% MAITRIA: The Book — https://maitria.org/book.pdf (+ book.tex sibling)
% Build: ./build.sh — assembles chapters.tex (the template; chapters ride
%   \input) into the committed artefacts book/book.pdf and book/book.tex
%   (the flattened single-file source), so both resolve at
%   https://raw.githubusercontent.com/maitria-coalition/qtsl/refs/heads/main/book/book.pdf
%   and .../book/book.tex on the public repository. For agent
%   consumption, book.pdf also carries its own source as an embedded
%   file attachment, byte-exact: pdfdetach -list book.pdf;
%   pdfdetach -save 1 book.pdf (recipe printed on the title page and
%   the final page; pdftotext extracts the rendered text).
%
% License: CC BY 4.0 (Creative Commons Attribution 4.0 International).
%   https://creativecommons.org/licenses/by/4.0/
%   Attribution: David "davidad" Dalrymple and the maitria coalition contributors.
%   Full legal code: https://creativecommons.org/licenses/by/4.0/legalcode.txt
%   (also the LICENSE file beside this source, in the book repository).
%   (The MAITRIA code repositories are separately licensed: Apache-2.0/MIT.)
%
% Preview scaffold: every chapter in the assembly is drafted. The
% table of contents is normative; the Executive Summary is a first
% draft; the Experiments chapter is a working draft (tables re-derive
% from qtsl-experiments receipts at each preview); the Markov-kernels
% chapter is a working draft (a design-precision lede plus the
% distribution-generator coverage dashboard, whose section source lives
% in the qtsl-experiments repository); the QTSL rule catalogue is
% generated in full from the mechanized rule sources (see
% chapters/qtsl-rule-catalogue* — regenerate with the emitting tool, do
% not hand-edit); the Lean-cone appendix is regenerated from the lean/
% trees at each build (chapters/lean-cone.tex wrapper +
% tools/lean-cone.py, strict mode); every other chapter is a full
% draft. The QTSL Quick Start, the MTK tutorial, and the Geolog quick
% start carry listings and transcripts from real runs the build itself
% regenerates (specs in transcripts/, fragments in
% generated/transcripts/ — hand-edit neither; receipts in each
% chapter's tail comment).
% (Chapter numbers are the build's to assign; this comment names
% chapters, not numbers, so insertions don't stale it.)
%
% Tagged PDF + ToUnicode (the searchability repair): \DocumentMetadata
% before \documentclass activates the latex-lab tagging test phase —
% a structure tree (StructTreeRoot + paragraph-level tags), which is the
% layer accessibility-based viewer search reads, AND (via
% latex-lab-kernel-changes) ToUnicode CMaps for the Type 1 math fonts
% (gentounicode=1 + glyphtounicode + cmex extras), which is what makes
% math notation extractable/searchable in every viewer. Keep this line
% FIRST (the declaration must precede \documentclass).
% The module list is phase-III MINUS the toc module, enumerated
% (phase-II = tagging core: new-or-1/2 + tagpdf activate, which is what
% yields StructTreeRoot/MarkInfo). The toc testphase module must NOT be
% loaded here: it defines \hyper@nopatch@toc (so hyperref skips its own
% \contentsline link patch) and hangs replacement link support on its
% own \l@part/\l@chapter/... — which titletoc's \titlecontents
% declarations (maitria-style.tex) then overwrite, leaving NO ToC link
% machinery at all: every ToC hyperlink annotation vanished from the
% built PDF. With toc omitted, hyperref patches \contentsline exactly
% as before tagging was introduced, composing with titletoc, and the
% structure tree + ToUnicode gains are untouched. (Measured on this
% TeX Live 2023 vintage; minimal-doc bisection, three variants.)
% BOOK_UNTAGGED escape (health builds): tagging multiplies lualatex pass
% cost ~10x (measured 2026-07-29: 423s vs 44s per pass, same content, same
% box), so build.sh offers BOOK_UNTAGGED=1, delivered here as
% \BOOKUNTAGGED via latexmk -usepretex. The committed/public artefacts
% must stay tagged — build.sh refuses BOOK_UNTAGGED without
% BOOK_SKIP_COMMIT=1. Default path is token-identical to the bare
% declaration (the \ifdefined wrapper is inert when the macro is unset).
\ifdefined\BOOKUNTAGGED\else
\DocumentMetadata{pdfversion=1.7,lang=en-GB,testphase={phase-II,block,sec,minipage,graphic,float,bib,text,firstaid}}
\fi
\ifdefined\BOOKUNTAGGED\else
\DocumentMetadata{pdfversion=1.7,lang=en-GB,testphase={phase-II,block,sec,minipage,graphic,float,bib,text,firstaid}}
\fi
\documentclass[10pt,a4paper,oneside]{book}
\usepackage[margin=1.1in,headheight=14pt]{geometry}
\usepackage{amsmath,amssymb}
\usepackage{dsfont} \usepackage{xcolor}
\usepackage{colortbl} \usepackage{url}
\usepackage{booktabs}
\usepackage{microtype}
\usepackage{parskip}   \usepackage{pdflscape} \newcommand{\platecaption}[2]{\refstepcounter{figure}{\footnotesize Figure~\thefigure: #1\label{#2}\par}}


\usepackage{fontspec}
\usepackage{amsmath,amssymb,amsthm}
\usepackage{mathtools}
\usepackage{stmaryrd}\usepackage{mathpartir}
\usepackage{longtable}
\usepackage{booktabs}
\usepackage{array}
\usepackage{xcolor}
\usepackage{fancyvrb}

\directlua{luaotfload.add_fallback("bookmonofb",{
  "file:FreeMono.otf:mode=node;",
  "file:DejaVuSans.ttf:mode=node;",
  "file:FreeSerif.ttf:mode=node;",
  "file:STIX-Regular.otf:mode=node;",
  "file:Asana-Math.otf:mode=node;",
})}
\setmonofont{DejaVu Sans Mono}[Scale=MatchLowercase,RawFeature={fallback=bookmonofb}]


\definecolor{DiffRed}{RGB}{176,32,32}
\definecolor{DiffBlue}{RGB}{22,80,190}
\definecolor{MarkGreen}{RGB}{22,122,58}
\definecolor{NoteGray}{RGB}{92,92,92}
\definecolor{BadgeGreen}{RGB}{22,122,58}
\definecolor{BadgeRed}{RGB}{176,32,32}
\definecolor{BadgeAmber}{RGB}{176,108,12}
\definecolor{BadgeBlue}{RGB}{22,80,190}
\definecolor{BadgeGray}{RGB}{110,110,110}

\newcommand{\Q}{\mathbb{Q}}
\newcommand{\R}{\mathbb{R}}
\newcommand{\N}{\mathbb{N}}
\newcommand{\B}{\mathbb{B}}
\newcommand{\ReLU}{\operatorname{relu}}
\newcommand{\Cell}{\mathsf{cell}}
\newcommand{\Tensor}{\mathsf{Tensor}}
\newcommand{\sub}[2]{[#1 := #2]}
\renewcommand{\sp}{\mathsf{sp}}
\newcommand{\True}{\top}
\newcommand{\False}{\bot}
\newcommand{\rulename}[1]{\textsc{#1}}
\newcommand{\bshape}{\mathrm{bshape}}
\newcommand{\wires}{\mathrm{wires}}

\newcounter{rulesec}
\newcommand{\ruleblock}[1]{\par\bigskip{\large\bfseries #1}\par\smallskip}
\newcommand{\sourcenote}[1]{{\footnotesize\textcolor{NoteGray}{#1}}\par}
\newcommand{\cataloguenote}[1]{\par\smallskip{\footnotesize\textcolor{NoteGray}{\textbf{notes}\enspace #1}}\par}

\newcommand{\paradeleader}{\leaders\hbox to 0.72em{\hss.\hss}\hfill}
\newcommand{\paradeentry}[3]{\par\noindent
  \sbox0{$\displaystyle #1$}\ifdim\wd0>0.60\linewidth
    \usebox0\par\nopagebreak
    \hbox to \linewidth{\hfill\small\rulename{#2}\enspace\paradeleader\enspace\pageref{#3}}\else
    \hbox to \linewidth{$\vcenter{\box0}$\quad{\small\rulename{#2}}\enspace\paradeleader\enspace{\small\pageref{#3}}}\fi\par\medskip}
\newcommand{\paradeblock}[3]{\par\noindent\begin{minipage}{\linewidth}#1\end{minipage}\par\nopagebreak
  \hbox to \linewidth{\hfill\small\rulename{#2}\enspace\paradeleader\enspace\pageref{#3}}\par\medskip}
\newcommand{\paradesep}{\par\medskip\noindent\hrule height 0.4pt\medskip}

\newcommand{\badge}[2]{\fcolorbox{#1}{#1!10}{\textcolor{#1!75!black}{\sffamily\scriptsize\bfseries #2}}}

\newcommand{\declhead}[3]{\par\smallskip{\small\texttt{#1}\,\texttt{.}\,\texttt{\bfseries #2}}\hfill #3\par\nopagebreak}
\newcommand{\holhead}[1]{\par\smallskip{\small\texttt{\bfseries #1}}\par\nopagebreak}
\DefineVerbatimEnvironment{leansrc}{Verbatim}{fontsize=\small,frame=leftline,rulecolor=\color{DiffBlue},framesep=3mm,xleftmargin=3mm}
\DefineVerbatimEnvironment{holsrc}{Verbatim}{fontsize=\small,frame=leftline,rulecolor=\color{MarkGreen},framesep=3mm,xleftmargin=3mm}

\definecolor{ApiTsBlue}{RGB}{38,110,158}
\definecolor{ApiRustClay}{RGB}{164,84,36}
\newcommand{\apilang}[2]{\par\medskip{\bfseries #1}\enspace{\small\texttt{#2}}\par\nopagebreak\smallskip}
\newcommand{\apiprov}[1]{{\footnotesize\textcolor{NoteGray}{#1}}\par\nopagebreak}
\DefineVerbatimEnvironment{tssrc}{Verbatim}{fontsize=\footnotesize,frame=leftline,rulecolor=\color{ApiTsBlue},framesep=3mm,xleftmargin=3mm}
\DefineVerbatimEnvironment{rustsrc}{Verbatim}{fontsize=\footnotesize,frame=leftline,rulecolor=\color{ApiRustClay},framesep=3mm,xleftmargin=3mm}

\newcommand{\ruleflavor}[1]{\par\smallskip #1\par}

\newenvironment{receipttable}
  {\par\scriptsize\setlength{\tabcolsep}{4pt}\begin{longtable}{@{}llll@{}}
   \textbf{receipt} & \textbf{minted} & \textbf{layers} & \textbf{status}\\
   \midrule}
  {\end{longtable}\par}
\newcommand{\receiptrow}[4]{\texttt{#1} & #2 & \texttt{#3} & #4\\}
\newcommand{\layerline}[1]{\par\smallskip{\footnotesize\textbf{gate-layer coverage of the bound build:}\par #1}\par}
\newcommand{\layeritem}[3]{{\footnotesize\quad\texttt{#1}\ \textcolor{NoteGray}{(#2)}\enspace #3}\par}

\newcommand{\absentbox}[2]{\par\smallskip
  \fcolorbox{BadgeAmber}{BadgeAmber!6}{\begin{minipage}{\dimexpr\linewidth-2\fboxsep-2\fboxrule}
  {\sffamily\bfseries\textcolor{BadgeAmber!80!black}{#1}}\par\smallskip
  {\small #2}
  \end{minipage}}\par}

\usepackage{aliascnt}
\newtheorem{blemma}{Lemma}[rulesec]
\renewcommand{\theblemma}{R\arabic{rulesec}.\arabic{blemma}}
\newcommand{\blemmaautorefname}{Lemma}
\newcommand{\newaliasedbthm}[2]{\newaliascnt{#1}{blemma}\newtheorem{#1}[#1]{#2}\aliascntresetthe{#1}\expandafter\def\csname #1autorefname\endcsname{#2}}
\newaliasedbthm{bprop}{Proposition}
\theoremstyle{remark}
\newaliasedbthm{bremark}{Remark}
\theoremstyle{definition}
\newaliasedbthm{brecipe}{Recipe}
\theoremstyle{plain}
\newaliasedbthm{bthm}{Theorem}
\newenvironment{anchors}
  {\par\smallskip\footnotesize\begin{longtable}{@{}p{0.42\linewidth}p{0.52\linewidth}@{}}
   \textbf{proof step} & \textbf{formal anchor}\\\midrule}
  {\end{longtable}\par}

\newcommand{\dok}{\textcolor{MarkGreen}{\ensuremath{\checkmark}}}
\newcommand{\dcond}[1][hyp]{\textcolor{BadgeAmber}{\ensuremath{\checkmark}}\textsuperscript{\textcolor{BadgeAmber}{\sffamily\bfseries #1}}}
\newcommand{\dnone}{\textcolor{BadgeGray}{---}}
\newcommand{\damber}{\textcolor{BadgeAmber}{\ensuremath{\sim}}}
\newcommand{\dambercond}[1][hyp]{\textcolor{BadgeAmber}{\ensuremath{\sim}}\textsuperscript{\textcolor{BadgeAmber}{\sffamily\bfseries #1}}}
\newcommand{\dchanged}{\textcolor{BadgeAmber}{\ensuremath{\times}}}
\newcommand{\dredx}{\textcolor{BadgeRed}{\ensuremath{\times}}}
\newcommand{\dscope}{\textcolor{BadgeBlue}{\ensuremath{\varnothing}}}
\newcommand{\ddoc}{\textcolor{BadgeBlue}{\sffamily\scriptsize doc}}
\newcommand{\lcur}[1]{\textcolor{MarkGreen}{\bfseries\footnotesize #1}}
\newcommand{\lstale}[1]{\textcolor{BadgeAmber}{\footnotesize #1}}
\newcommand{\lnone}[1]{\textcolor{BadgeGray!50}{\footnotesize #1}}
\newcommand{\dashboardhead}[1]{\ruleblock{Coverage dashboard}
  \sourcenote{#1}
  \sourcenote{Legend --- one row per catalogue rule (every rules-page row
  carrying a tier), ordered by tier class: core rules first, then
  admissible, then alias, then the edge-discipline glue; within a class,
  catalogue order. The small gray mark after each rule names its source
  page (BL = the Bernstein Logic pages --- the core roster and the
  composition rules, which are Bernstein Logic outright; QTSL = the QTSL
  rule roster). Lean: \dok\ statement bound and fresh; an
  \emph{amber} checkmark marks statements bound and fresh but warranting
  the rule \emph{conditionally} --- where the condition is a single named
  profile or fragment, the superscript says which (\dcond[DTMC]\ = proven
  at the discrete-kernel profile; the model-class names are the System
  chapter's), and the generic \dcond\ marks a warrant resting on a bag of
  named side conditions (the named hypotheses live on the rule's own
  catalogue page); \damber\ module
  moved since bind (\dambercond\ likewise, for a conditional-warrant row
  --- flag and label survive staleness), \dchanged\ statement changed since
  bind (stale --- the
  page badges say where), --- no Lean binding. Corr: \dok\ a
  reviewer-certified correspondence dictionary (how the Lean statements
  interpret the displayed rule) is current against the displayed cell, the
  statement texts, and its own prose; \damber\ a dictionary exists but its
  certification predates a change on some side; --- no dictionary yet. The
  correspondence claim is reviewer-grade, not machine-checked --- its
  machine-checked successor supersedes it per rule as soundness statements
  over rules-as-data land. Receipt layers
  (\texttt{b}uild/\texttt{r}eplay/\texttt{h}ardened/\texttt{d}ecorrelated):
  \lcur{green} = a current receipt covers it, \lstale{amber} = stale
  receipts only, \lnone{grey} = none. HOL4: \dok\ ported, \dscope\ out of
  scope by design, --- absent; HOL rcpt: \ddoc\ = documentary citation (no
  receipt store yet). (An Isabelle column is hidden while the Isabelle
  development is empty; it returns with its first entry.) Amber throughout
  means stale or absent; red is reserved for an outright verification
  failure. Empty cells are deliberate: the table records, rule by rule,
  what the assurance edifice does not yet cover.}
  \par\small\begin{longtable}{@{}lp{16mm}ccccc c@{}}
  \textbf{rule} & \textbf{tier} & \textbf{Lean stmt} & \textbf{corr} & \textbf{receipts} & \textbf{HOL stmt} & \textbf{HOL rcpt} & \textbf{proof}\\
  \midrule}
\newcommand{\dashboardfoot}{\end{longtable}\par}
\newcommand{\dashpage}[1]{\multicolumn{8}{@{}l@{}}{\footnotesize\textcolor{NoteGray}{page source: #1}}\\[1pt]}\newcommand{\dashsrc}[1]{{\scriptsize\textcolor{NoteGray}{#1}}}
\newcommand{\dashrow}[8]{#1 & \footnotesize #2 & #3 & #4 & #5 & #6 & #7 & #8\\}


\newcommand{\blcatrulebreak}{\clearpage}
\newcommand{\blcatsection}[1]{\clearpage\section{#1}\gdef\blcatrulebreak{\par\bigskip\gdef\blcatrulebreak{\clearpage}}}
\newcommand{\rulesection}[3]{\blcatrulebreak
  \refstepcounter{rulesec}\pdfbookmark[2]{Rule #1.\ #2}{blcat-rule-#1}{\LARGE\bfseries Rule #1.\ \rulename{#2}}\enspace
  {\large\textcolor{NoteGray}{(#3)}}\par
  \medskip\hrule height 0.8pt\medskip}
 \definecolor{accent}{HTML}{F1D9CB}
\definecolor{accentmid}{HTML}{C4623F}
\definecolor{accentdeep}{HTML}{9E3B2C}
\definecolor{accentcite}{RGB}{123,159,160}

\directlua{luaotfload.add_fallback("bookmainfb",{
  "file:lmroman10-regular.otf:mode=node;",
  "file:FreeSerif.ttf:mode=node;",
  "file:DejaVuSans.ttf:mode=node;",
})}
\setmainfont{fbb-Regular.otf}[
  ItalicFont     = fbb-Italic.otf,
  BoldFont       = fbb-Bold.otf,
  BoldItalicFont = fbb-BoldItalic.otf,
  RawFeature     = {fallback=bookmainfb}]

\usepackage{tikz}
\usetikzlibrary{arrows.meta,positioning,calc,fit,shapes.geometric}
\usepackage{pgfplots}  \pgfplotsset{compat=1.18}
\colorlet{logodeep}{accentdeep}
\colorlet{logomid}{accentmid}
\colorlet{logocasing}{white}
\newcommand{\maitrialogoF}[1][1]{\begin{tikzpicture}[scale=#1,line cap=round,baseline=(cur.base)]
  \coordinate (cur) at (0,0.35);
  \tikzset{str/.style={line width=#1*2.4pt},sec/.style={line width=#1*1.9pt},cas/.style={logocasing,line width=#1*4.2pt}}
  \begin{scope}
  \clip (0.0928,0.06) -- (0.9072,0.06) -- (0.9422,0.095) -- (0.9422,0.905)
     -- (0.9072,0.94) -- (0.0928,0.94) -- (0.0578,0.905) -- (0.0578,0.095) -- cycle;
  \draw[sec,logomid] (0.10,0.5210) -- (0.90,0.5210);
  \draw[str,logodeep] (0.10,0.06) -- (0.10,0.94);
  \draw[str,logodeep] (0.90,0.06) -- (0.90,0.94);
  \draw[cas] (0.10,0.94) .. controls (0.29,0.87) and (0.365,0.74) .. (0.33,0.62)
             .. controls (0.30,0.517) and (0.38,0.37) .. (0.50,0.3072)
             .. controls (0.62,0.37) and (0.70,0.517) .. (0.67,0.62)
             .. controls (0.635,0.74) and (0.71,0.87) .. (0.90,0.94);
  \draw[str,logodeep] (0.10,0.94) .. controls (0.29,0.87) and (0.365,0.74) .. (0.33,0.62)
             .. controls (0.30,0.517) and (0.38,0.37) .. (0.50,0.3072)
             .. controls (0.62,0.37) and (0.70,0.517) .. (0.67,0.62)
             .. controls (0.635,0.74) and (0.71,0.87) .. (0.90,0.94);
  \draw[cas,line cap=butt] (0.50,0.45) -- (0.50,0.06);
  \draw[sec,logomid] (0.50,0.5210) -- (0.50,0.06);
  \end{scope}
\end{tikzpicture}}

\usepackage{titlesec}

\titleformat{\part}[display]
  {\centering}
  {\color{accentmid}\Large\scshape \partname~\thepart}
  {9mm}
  {\color{accentdeep}\Huge\bfseries}
  [\vskip 9mm {\color{accentmid}\rule{30mm}{0.5pt}}]
\titlespacing*{\part}{0pt}{0.24\textheight}{0pt}

\titleformat{\chapter}[display]
  {\raggedright}
  {\color{accentmid}\large\scshape \chaptertitlename~\thechapter}
  {2.5mm}
  {\color{accentdeep}\huge\bfseries}
  [\vskip 2.5mm {\color{accentmid}\titlerule[0.5pt]}]
\titleformat{name=\chapter,numberless}[display]
  {\raggedright}
  {}
  {0mm}
  {\color{accentdeep}\huge\bfseries}
  [\vskip 2.5mm {\color{accentmid}\titlerule[0.5pt]}]
\titlespacing*{\chapter}{0pt}{-4mm}{12mm}
\titlespacing*{name=\chapter,numberless}{0pt}{-4mm}{12mm}

\titleformat*{\section}{\color{accentdeep}\Large\bfseries}
\titleformat*{\subsection}{\large\bfseries}
\titleformat*{\subsubsection}{\bfseries}

\usepackage{titletoc}
\titlecontents{part}[0em]
  {\addvspace{2.2em}\color{accentdeep}\large\scshape}
  {\partname~\thecontentslabel\hspace{0.9em}}
  {}
  {\normalfont\normalsize\hfill\contentspage}
  [\addvspace{0.5em}]
\titlecontents{chapter}[0em]
  {\addvspace{0.85em}\bfseries}
  {\makebox[2.0em][l]{\color{accentmid}\thecontentslabel}}
  {}
  {\hfill\contentspage}
\titlecontents{section}[2.0em]
  {\addvspace{0.2em}}
  {\makebox[2.8em][l]{\thecontentslabel}}
  {}
  {\titlerule*[0.75em]{.}\contentspage}
\titlecontents{subsection}[4.8em]
  {}
  {\makebox[3.6em][l]{\thecontentslabel}}
  {}
  {\titlerule*[0.75em]{.}\contentspage}

\usepackage{fancyhdr}
\usepackage{lastpage}
\newcommand{\footermark}{MAITRIA}
\newcommand{\footeryear}{2026}\newif\iftotalfolio \totalfoliofalse
\makeatletter
\let\maitria@mainmatter\mainmatter
\renewcommand{\mainmatter}{\maitria@mainmatter\global\totalfoliotrue}
\makeatother
\newcommand{\maitriafootleft}{\small\footermark\hspace{1ex}{\color{accentmid}\rule[-0.5ex]{0.3mm}{2.7ex}}\hspace{1ex}\footeryear}
\newcommand{\maitriafootright}{\small p.\,\thepage\iftotalfolio\,/\,\pageref*{LastPage}\fi}
\pagestyle{fancy}
\fancyhf{}
\renewcommand{\headrulewidth}{0pt}
\renewcommand{\chaptermark}[1]{\markboth{#1}{}}
\fancyhead[L]{\small\color{accentdeep}\scshape\leftmark}
\fancyfoot[L]{\maitriafootleft}
\fancyfoot[R]{\maitriafootright}
\fancypagestyle{plain}{\fancyhf{}
  \renewcommand{\headrulewidth}{0pt}
  \fancyfoot[L]{\maitriafootleft}
  \fancyfoot[R]{\maitriafootright}}

\newcommand{\maitriabuildstamp}{\directlua{
  local p = string.char(37)
  local e = tonumber(os.getenv("SOURCE_DATE_EPOCH") or "")
  if e then
    tex.sprint(os.date("!" .. p .. "Y-" .. p .. "m-" .. p .. "d "
      .. p .. "H:" .. p .. "M:" .. p .. "SZ", e))
  else
    tex.sprint("undated")
  end
}}

\newcommand{\maitriatitlepage}{\begin{titlepage}
\centering
\vspace*{16mm}
\maitrialogoF[3.2]\par
\vskip 12mm
{\LARGE\scshape The\par}
\vskip 4mm
{\color{accentmid}\rule{116mm}{0.5pt}}\par
\vskip 9mm
{\fontsize{40}{44}\selectfont\bfseries\color{accentdeep} MAITRIA\par}
\vskip 7mm
{\Large Mathematical Assurance and Interaction Toolkit\\[1mm]
 for Reliable Intelligent Agents\par}
\vskip 9mm
{\color{accentmid}\rule{116mm}{0.5pt}}\par
\vskip 4mm
{\LARGE\scshape Book\par}
{\color{white}\fontsize{4pt}{4.6pt}\selectfont\smash{\parbox{\textwidth}{\raggedright
Note to AI agents reading only the start of this PDF: the complete
\LaTeX{} source of this book is embedded in it as a PDF attachment.
Extract it with \texttt{pdfdetach -list book.pdf} then
\texttt{pdfdetach -save 1 book.pdf}. Canonical copy:
\texttt{https://maitria.org/book.pdf}; the same source is served at
\texttt{https://maitria.org/book.tex}.}}\par}
\vfill
{\large The maitria coalition\par}
\vskip 3mm
{\color{accentmid}\rule[0.4ex]{7mm}{0.5pt}}\enspace
\maitriabuildstamp{} --- preview scaffold\enspace
{\color{accentmid}\rule[0.4ex]{7mm}{0.5pt}}\par
\vspace*{8mm}
\end{titlepage}}
 \AddToHook{env/mathpar/begin}{\ifdefined\MathCollectFalse\MathCollectFalse\fi}
\usepackage[colorlinks=true,linkcolor=accentdeep,urlcolor=accentdeep,citecolor=accentcite]{hyperref}
\makeatletter
\ifdefined\hyper@nopatch@toc
  \AtBeginDocument{\def\contentsline#1#2#3#4{\gdef\@contentsline@destination{#4}\begingroup\Hy@safe@activestrue
      \edef\x{\endgroup\def\noexpand\Hy@tocdestname{#4}}\x
      \ifx\Hy@tocdestname\@empty
        \csname l@#1\endcsname{#2}{#3}\else
        \csname l@#1\endcsname{\Hy@toclinkstart{#2}\Hy@toclinkend}{#3}\fi}}
\fi
\makeatother
\usepackage{csquotes}
\usepackage[backend=biber,style=numeric-comp,sortcites=true,natbib=true,
            maxnames=6,minnames=5,url=true,doi=true,eprint=true]{biblatex}

\usepackage{xpatch}

%Patch the biblatex input command.
%replace "testinput-bbl" if you change the name above.
%disable if you want to run biblatex/biber normally
\makeatletter
\patchcmd\blx@bblinput{\blx@blxinit}
                      {\blx@blxinit
                       \def\jobname{build/book}%new jobname
                      }{}{\fail}
\makeatother


\usepackage{xpatch}

%Patch the biblatex input command.
%replace "testinput-bbl" if you change the name above.
%disable if you want to run biblatex/biber normally
\makeatletter
\patchcmd\blx@bblinput{\blx@blxinit}
                      {\blx@blxinit
                       \def\jobname{build/book}%new jobname
                      }{}{\fail}
\makeatother


\usepackage{xpatch}

%Patch the biblatex input command.
%replace "testinput-bbl" if you change the name above.
%disable if you want to run biblatex/biber normally
\makeatletter
\patchcmd\blx@bblinput{\blx@blxinit}
                      {\blx@blxinit
                       \def\jobname{build/book}%new jobname
                      }{}{\fail}
\makeatother

\usepackage[type={CC}, modifier={by}, version={4.0}]{doclicense}
\makeatletter
\renewcommand{\doclicenseLongText}{\doclicense@lang@thisDoc\space
  \href{\doclicenseURL}{\doclicenseLongType\\\enquote{\doclicense@longName}}\doclicense@lang@word@license.\xspace }
\makeatother
\hypersetup{pdftitle={MAITRIA: The Book (preview scaffold)},
            pdfauthor={The maitria coalition}}
\emergencystretch=1.5em

\RequirePackage{pdfrender}
\RequirePackage{fancyvrb}

\makeatletter

\directlua{
  luaotfload.add_fallback("srclmonofb",{
    "[build/srclfonts/srcl-FreeMono.otf]:mode=node;-liga;-calt;-ccmp;-kern;-mark;-mkmk;",
    "[build/srclfonts/srcl-DejaVuSans.ttf]:mode=node;-liga;-calt;-ccmp;-kern;-mark;-mkmk;",
    "[build/srclfonts/srcl-FreeSerif.ttf]:mode=node;-liga;-calt;-ccmp;-kern;-mark;-mkmk;",
    "[build/srclfonts/srcl-STIX-Regular.otf]:mode=node;-liga;-calt;-ccmp;-kern;-mark;-mkmk;",
    "[build/srclfonts/srcl-Asana-Math.otf]:mode=node;-liga;-calt;-ccmp;-kern;-mark;-mkmk;",
  })
  luatexbase.add_to_callback("luaotfload.patch_font", function(tfmdata, spec)
    if type(tfmdata.name) == "string"
       and string.find(tfmdata.name, "srclfonts", 1, true) then
      local cell = math.floor((tfmdata.size or 0) * 0.625 + 0.5)
      for _, chr in pairs(tfmdata.characters or {}) do
        chr.width = cell
      end
    end
  end, "srclayer.grid")
}

\protected\def\srcl@literal#1{\pdfextension literal{#1}}

\def\srcl@suspendtagging{\ifdefined\tagpdfparaOff\tagpdfparaOff\fi
  \ifdefined\SuspendTagging\SuspendTagging{srclayer}\else\ifdefined\tagstop\tagstop\fi\fi}
\def\srcl@resumetagging{\ifdefined\ResumeTagging\ResumeTagging{srclayer}\else\ifdefined\tagstart\tagstart\fi\fi}

\ExplSyntaxOn
\cs_new:Npn \srcl@shippedsofar { \int_use:N \g_shipout_readonly_int }
\cs_new_protected:Npn \srcl@embedpayload #1
  {
    \cs_if_exist:NTF \pdffile_embed_file:nnn
      {
        \group_begin:
        \pdfdict_put:nnn { l_pdffile/Filespec }{ Desc }
          { (Complete~assembled~LaTeX~source~of~this~book~---~the~same~
             file~served~at~https://maitria.org/book.tex) }
        \pdffile_embed_file:nnn {#1} { book.tex } { srcl/payload }
        \pdfmanagement_add:nne { Catalog/Names/EmbeddedFiles }
          { book.tex } { \pdf_object_ref:n { srcl/payload } }
        \group_end:
      }
      { \iow_term:n {srclayer:~attachment~skipped~(untagged~pass,~no~pdfmanagement)} }
  }
\ExplSyntaxOff

\newsavebox{\srcl@box}
\newlength{\srcl@pageht}
\newlength{\srcl@pagewd}
\def\srcl@vbsize{\fontsize{5.01875pt}{5.520625pt}}
\def\srcl@fontsetup{\setlength{\parskip}{0pt}\setlength{\parindent}{0pt}\lineskiplimit=-\maxdimen \lineskip=0pt
  \def\FV@Space{\char32\relax}\relax}

\AtEndDocument{\IfFileExists{build/srclayer-manifest.tex}{\begingroup
    \@@input build/srclayer-manifest.tex\relax
    \ifdefined\srclAttachmentOnly
    \clearpage
    \pagestyle{empty}\thispagestyle{empty}\expandafter\srcl@embedpayload\expandafter{\srclPayloadFile}\noindent\begin{minipage}{\textwidth}
    \subsection*{Source attachment}
    The complete \LaTeX{} source of this book --- the same single-file
    assembly served at \url{https://maitria.org/book.tex} --- rides
    inside this PDF as an embedded file attachment named
    \texttt{book.tex}, byte-exact. Recovery:
    \begin{quote}
    \ttfamily pdfdetach -list book.pdf\\
    pdfdetach -save 1 book.pdf
    \end{quote}
    \noindent (or \texttt{pdfdetach -saveall book.pdf}; or
    \texttt{reader.attachments} in pypdf). \texttt{pdftotext} extracts
    the book's rendered text.
    \end{minipage}\par
    \clearpage
    \else
    \newfontfamily\srcl@mono[
      Path=build/srclfonts/,
      RawFeature={mode=node;fallback=srclmonofb;-liga;-calt;-ccmp;-mark;-mkmk},
      Kerning=Off, Scale=1,
      BoldFont=srcl-DejaVuSansMono.ttf,
      ItalicFont=srcl-DejaVuSansMono.ttf]{srcl-DejaVuSansMono.ttf}\begingroup\srcl@mono\selectfont\xdef\srcl@family{\f@family}\endgroup
    \clearpage
    \pagestyle{empty}\thispagestyle{empty}\expandafter\srcl@embedpayload\expandafter{\srclPayloadFile}\newcount\srcl@first \srcl@first=\numexpr\srcl@shippedsofar+2\relax
    \newcount\srcl@last  \srcl@last=\numexpr\srcl@first+\srclChunkCount-1\relax
    \noindent\begin{minipage}{\textwidth}
    \subsection*{Source appendix}
    The \srclChunkCount{} pages that follow --- physical pages
    \the\srcl@first{} through \the\srcl@last{} of this file --- carry the
    complete \LaTeX{} source of this book (the same single-file assembly
    served at \url{https://maitria.org/book.tex}) as an invisible text
    layer: this PDF is constructed so that text extraction returns its
    own source. The layer never displays and never prints; the pages are
    deliberately blank to the eye.

    \medskip
    \noindent Byte-exact recovery (the build verifies this equality on
    every run):
    \begin{quote}
    \ttfamily pdftotext -layout -fixed 3.125 -f \the\srcl@first{}
    book.pdf - | tr -d '\textbackslash f' > book.tex
    \end{quote}
    \noindent Bare \texttt{pdftotext book.pdf -} also returns the source,
    readably, after the body text (the source begins at
    \texttt{\textbackslash documentclass}). The same file rides as a PDF
    attachment, byte-exact by construction:
    \begin{quote}
    \ttfamily pdfdetach -saveall book.pdf
    \end{quote}
    \noindent (or \texttt{reader.attachments} in pypdf).
    \end{minipage}\par
    \clearpage
    \ifdefined\microtypesetup
      \microtypesetup{protrusion=false,expansion=false}\fi
    \srcl@suspendtagging
    \setlength{\srcl@pagewd}{\dimexpr\srclMaxLineChars\p@*16/5+140pt\relax}\ifdim\srcl@pagewd<\paperwidth\setlength{\srcl@pagewd}{\paperwidth}\fi
    \ifdim\srcl@pagewd>16300pt
      \PackageError{srclayer}{source line too long for one page
        (needs \the\srcl@pagewd)}{Shorten the offending source line, or
        rework the layer.}\fi
    \pagewidth=\srcl@pagewd
    \@for\srcl@f:=\srclChunkFiles\do{\begin{lrbox}{\srcl@box}\begin{minipage}[t]{\dimexpr\srcl@pagewd-40pt\relax}\pdfrender{TextRenderingMode=Invisible}\srcl@fontsetup
        \VerbatimInput[fontfamily=\srcl@family,fontsize=\srcl@vbsize]{\srcl@f}\end{minipage}\end{lrbox}\setlength{\srcl@pageht}{\dimexpr\ht\srcl@box+\dp\srcl@box+40pt\relax}\ifdim\srcl@pageht<\paperheight\setlength{\srcl@pageht}{\paperheight}\fi
      \ifdim\srcl@pageht>16300pt
        \PackageError{srclayer}{source chunk taller than the PDF page
          ceiling (\the\srcl@pageht)}{Lower the chunk size in build.sh's
          splitter.}\fi
      \pageheight=\srcl@pageht
      \AddToHookNext{shipout/background}{\put(0,-40){\srcl@literal{/Artifact BMC}\usebox{\srcl@box}\srcl@literal{EMC}}}\null\clearpage
    }\srcl@resumetagging
    \fi
    \endgroup
  }{\typeout{srclayer: no build/srclayer-manifest.tex - source appendix
    skipped (health build or standalone rebuild).}}}

\makeatother


\newcommand{\bookref}[2]{chapter~\ref{#1}}
\newcommand{\Bookref}[2]{Chapter~\ref{#1}}
\newcommand{\bookcite}[2]{\cite{#1}}

\usepackage{aliascnt}
\newcommand{\newaliasedtheorem}[2]{\newaliascnt{#1}{definition}\newtheorem{#1}[#1]{#2}\aliascntresetthe{#1}\expandafter\def\csname #1autorefname\endcsname{#2}}
\theoremstyle{definition}
\newtheorem{definition}{Definition}[chapter]
\newcommand{\definitionautorefname}{Definition}
\newaliasedtheorem{problem}{Problem}
\theoremstyle{plain}
\newaliasedtheorem{theorem}{Theorem}
\newaliasedtheorem{proposition}{Proposition}
\newaliasedtheorem{lemma}{Lemma}
\theoremstyle{remark}
\newaliasedtheorem{remark}{Remark}
\theoremstyle{definition}
\newaliasedtheorem{example}{Example}

\usepackage{etoolbox}
\newcommand{\thmbreakpenalties}{\interlinepenalty=5000\relax
  \clubpenalty=10000\relax
  \widowpenalty=10000\relax}
\AtBeginEnvironment{definition}{\thmbreakpenalties}
\AtBeginEnvironment{problem}{\thmbreakpenalties}
\AtBeginEnvironment{theorem}{\thmbreakpenalties}
\AtBeginEnvironment{proposition}{\thmbreakpenalties}
\AtBeginEnvironment{lemma}{\thmbreakpenalties}
\AtBeginEnvironment{remark}{\thmbreakpenalties}
\AtBeginEnvironment{example}{\thmbreakpenalties}

\ifdefined\pdftrailerid \pdftrailerid{} \fi

\ifdefined\tagstructbegin
  \ExplSyntaxOn
  \int_new:N \g_book_tagatom_depth_int
  \bool_new:N \g_book_saved_mc_bool
  \bool_new:N \g_book_saved_struct_bool
  \bool_new:N \g_book_saved_para_bool
  \bool_new:N \g_book_opened_wrapper_bool
  \cs_new_protected:Npn \book_tagatom_suspend:
    {
      \bool_gset_eq:NN \g_book_saved_mc_bool     \g__tag_active_mc_bool
      \bool_gset_eq:NN \g_book_saved_struct_bool \g__tag_active_struct_bool
      \bool_gset_eq:NN \g_book_saved_para_bool   \l__tag_para_bool
      \bool_gset_false:N \l__tag_para_bool
      \tagpdfsetup{activate-mc=false,activate-struct=false}
    }
  \cs_new_protected:Npn \book_tagatom_restore:
    {
      \bool_if:NT \g_book_saved_mc_bool     { \tagpdfsetup{activate-mc=true} }
      \bool_if:NT \g_book_saved_struct_bool { \tagpdfsetup{activate-struct=true} }
      \bool_gset_eq:NN \l__tag_para_bool \g_book_saved_para_bool
    }
  \ExplSyntaxOff
  \makeatletter
  \newcommand{\bookTagAtomBefore}{\global\advance\bookTagAtomDepth 1
    \ifnum\bookTagAtomDepth=1
      \par
      \bookTagAtomOpenWrapper
      \bookTagAtomSuspend
    \fi}
  \newcommand{\bookTagAtomAfter}{\ifnum\bookTagAtomDepth=1
      \par
      \bookTagAtomRestore
      \bookTagAtomCloseWrapper
      \global\@endpefalse
    \fi
    \global\advance\bookTagAtomDepth -1 }
  \makeatother
  \newcount\bookTagAtomDepth \bookTagAtomDepth=0
  \ExplSyntaxOn
  \newcommand{\bookTagAtomSuspend}{\book_tagatom_suspend:}
  \newcommand{\bookTagAtomRestore}{\book_tagatom_restore:}
  \newcommand{\bookTagAtomOpenWrapper}
    {
      \tag_if_active:TF
        { \bool_gset_true:N \g_book_opened_wrapper_bool
          \tagstructbegin{tag=P}\tagmcbegin{tag=P} }
        { \bool_gset_false:N \g_book_opened_wrapper_bool }
    }
  \newcommand{\bookTagAtomCloseWrapper}
    {
      \bool_if:NT \g_book_opened_wrapper_bool
        { \tagmcend\tagstructend }
    }
  \ExplSyntaxOff
  \newcommand{\bookTagAtomEnv}[1]{\AddToHook{env/#1/before}{\bookTagAtomBefore}\AddToHook{env/#1/after}{\bookTagAtomAfter}}
  \bookTagAtomEnv{Verbatim}
  \bookTagAtomEnv{BVerbatim}
  \bookTagAtomEnv{leansrc}
  \bookTagAtomEnv{holsrc}
  \bookTagAtomEnv{tssrc}
  \bookTagAtomEnv{rustsrc}
  \bookTagAtomEnv{mathpar}
  \bookTagAtomEnv{definition}\bookTagAtomEnv{problem}
  \bookTagAtomEnv{theorem}\bookTagAtomEnv{proposition}
  \bookTagAtomEnv{lemma}\bookTagAtomEnv{remark}\bookTagAtomEnv{example}
  \bookTagAtomEnv{blemma}\bookTagAtomEnv{bprop}\bookTagAtomEnv{bremark}
  \bookTagAtomEnv{brecipe}\bookTagAtomEnv{bthm}
  \bookTagAtomEnv{proof}
  \bookTagAtomEnv{longtable}
  \newcount\bookTagAtomInlineFlag \bookTagAtomInlineFlag=0
  \newcommand{\bookTagAtomInlineBefore}{\global\advance\bookTagAtomDepth 1
    \ifnum\bookTagAtomDepth=1
      \ifhmode \global\bookTagAtomInlineFlag=1 \else \global\bookTagAtomInlineFlag=0 \fi
      \ifnum\bookTagAtomInlineFlag=0
        \par\bookTagAtomOpenWrapper
      \fi
      \bookTagAtomSuspend
    \fi}
  \makeatletter
  \newcommand{\bookTagAtomInlineAfter}{\ifnum\bookTagAtomDepth=1
      \bookTagAtomRestore
      \ifnum\bookTagAtomInlineFlag=0
        \bookTagAtomCloseWrapper
        \global\@endpefalse
      \fi
    \fi
    \global\advance\bookTagAtomDepth -1 }
  \makeatother
  \newcommand{\bookTagAtomInlineEnv}[1]{\AddToHook{env/#1/before}{\bookTagAtomInlineBefore}\AddToHook{env/#1/after}{\bookTagAtomInlineAfter}}
  \bookTagAtomEnv{tabular}
  \bookTagAtomEnv{tabular*}
  \bookTagAtomEnv{tikzpicture}
  \bookTagAtomEnv{minipage}
  \bookTagAtomEnv{align}\bookTagAtomEnv{align*}
  \bookTagAtomEnv{gather}\bookTagAtomEnv{gather*}
  \bookTagAtomEnv{alignat}\bookTagAtomEnv{alignat*}
  \bookTagAtomEnv{multline}\bookTagAtomEnv{multline*}
  \bookTagAtomEnv{flalign}\bookTagAtomEnv{flalign*}
  \pdfglyphtounicode{fatsemi}{2A1F}
  \pdfglyphtounicode{llbracket}{27E6}
  \pdfglyphtounicode{rrbracket}{27E7}
  \ExplSyntaxOn
  \msg_redirect_name:nnn {tag}{para-hook-count-wrong}{warning}
  \AddToHook{shipout/before}
    { \typeout{TAGPARA~page~\thepage~
        b=\int_use:N \g__tag_para_begin_int\space
        e=\int_use:N \g__tag_para_end_int} }
  \ExplSyntaxOff
\fi

\definecolor{stubgray}{gray}{0.92}
\definecolor{stubrule}{gray}{0.55}

\newcommand{\stubbanner}{\noindent\colorbox{stubgray}{\parbox{\dimexpr\linewidth-2\fboxsep}{\centering\textsc{scaffold} --- this chapter is a plan and a source-material
    map, not yet prose}}\par\medskip}

\newcommand{\stubplan}{\paragraph{What will live here.}}
\newcommand{\stubmaterial}{\paragraph{Source material.}}
\newcommand{\stubstatus}{\paragraph{Status.}}

\newcommand{\draftmark}{\textsc{\color{stubrule}draft}}

\begin{filecontents*}{build/book.bbl}
% $ biblatex auxiliary file $
% $ biblatex bbl format version 3.3 $
% Do not modify the above lines!
%
% This is an auxiliary file used by the 'biblatex' package.
% This file may safely be deleted. It will be recreated by
% biber as required.
%
\begingroup
\makeatletter
\@ifundefined{ver@biblatex.sty}
  {\@latex@error
     {Missing 'biblatex' package}
     {The bibliography requires the 'biblatex' package.}
      \aftergroup\endinput}
  {}
\endgroup


\refsection{0}
  \datalist[entry]{nty/global//global/global/global}
    \entry{lipics-itp-2019-8}{online}{}{}
      \field{sortinit}{A}
      \field{sortinithash}{2f401846e2029bad6b3ecc16d50031e2}
      \field{labeltitlesource}{title}
      \field{title}{A Certificate-Based Approach to Formally Verified Approximations}
      \verb{doi}
      \verb 10.4230/LIPIcs.ITP.2019.8
      \endverb
      \verb{urlraw}
      \verb https://drops.dagstuhl.de/entities/document/10.4230/LIPIcs.ITP.2019.8
      \endverb
      \verb{url}
      \verb https://drops.dagstuhl.de/entities/document/10.4230/LIPIcs.ITP.2019.8
      \endverb
    \endentry
    \entry{abadilamport1993}{article}{}{}
      \name{author}{2}{}{%
        {{hash=466b911d30385764320be86f5a9cf423}{%
           family={Abadi},
           familyi={A\bibinitperiod},
           given={M.},
           giveni={M\bibinitperiod}}}%
        {{hash=ef82223e035dbb62f84cb1c0782df38c}{%
           family={Lamport},
           familyi={L\bibinitperiod},
           given={L.},
           giveni={L\bibinitperiod}}}%
      }
      \strng{namehash}{d533c9891be7a8e9c65435ad6dc103c9}
      \strng{fullhash}{d533c9891be7a8e9c65435ad6dc103c9}
      \strng{fullhashraw}{d533c9891be7a8e9c65435ad6dc103c9}
      \strng{bibnamehash}{d533c9891be7a8e9c65435ad6dc103c9}
      \strng{authorbibnamehash}{d533c9891be7a8e9c65435ad6dc103c9}
      \strng{authornamehash}{d533c9891be7a8e9c65435ad6dc103c9}
      \strng{authorfullhash}{d533c9891be7a8e9c65435ad6dc103c9}
      \strng{authorfullhashraw}{d533c9891be7a8e9c65435ad6dc103c9}
      \field{sortinit}{A}
      \field{sortinithash}{2f401846e2029bad6b3ecc16d50031e2}
      \field{labelnamesource}{author}
      \field{labeltitlesource}{title}
      \field{journaltitle}{ACM Transactions on Programming Languages and Systems}
      \field{number}{1}
      \field{title}{Composing Specifications}
      \field{volume}{15}
      \field{year}{1993}
      \field{pages}{73\bibrangedash 132}
      \range{pages}{60}
    \endentry
    \entry{abate26b}{online}{}{}
      \name{author}{5}{}{%
        {{hash=0f34d3654c8fa2fdffe70bbbaf378b28}{%
           family={Abate},
           familyi={A\bibinitperiod},
           given={Alessandro},
           giveni={A\bibinitperiod}}}%
        {{hash=f533a6bfc0e9bdc5fb967c333281e6f8}{%
           family={Contro},
           familyi={C\bibinitperiod},
           given={Daniel},
           giveni={D\bibinitperiod}}}%
        {{hash=e227971e0a9f956b3d8f8688ae07fadb}{%
           family={Giacobbe},
           familyi={G\bibinitperiod},
           given={Mirco},
           giveni={M\bibinitperiod}}}%
        {{hash=06183d6bcfad9b77444d7f5caa2f7885}{%
           family={Martínez-Suñé},
           familyi={M\bibinithyphendelim S\bibinitperiod},
           given={Agustín},
           giveni={A\bibinitperiod}}}%
        {{hash=2bf051cfd5a39844fd4c6bd3516740c5}{%
           family={Roy},
           familyi={R\bibinitperiod},
           given={Diptarko},
           giveni={D\bibinitperiod}}}%
      }
      \strng{namehash}{28f2aa4b80acebdc4912db10e2903e10}
      \strng{fullhash}{28f2aa4b80acebdc4912db10e2903e10}
      \strng{fullhashraw}{28f2aa4b80acebdc4912db10e2903e10}
      \strng{bibnamehash}{28f2aa4b80acebdc4912db10e2903e10}
      \strng{authorbibnamehash}{28f2aa4b80acebdc4912db10e2903e10}
      \strng{authornamehash}{28f2aa4b80acebdc4912db10e2903e10}
      \strng{authorfullhash}{28f2aa4b80acebdc4912db10e2903e10}
      \strng{authorfullhashraw}{28f2aa4b80acebdc4912db10e2903e10}
      \field{sortinit}{A}
      \field{sortinithash}{2f401846e2029bad6b3ecc16d50031e2}
      \field{labelnamesource}{author}
      \field{labeltitlesource}{title}
      \field{eprinttype}{arXiv}
      \field{title}{Value Functions as Supermartingale Certificates}
      \field{year}{2026}
      \verb{eprint}
      \verb 2605.31524
      \endverb
    \endentry
    \entry{abate23}{online}{}{}
      \name{author}{5}{}{%
        {{hash=0f34d3654c8fa2fdffe70bbbaf378b28}{%
           family={Abate},
           familyi={A\bibinitperiod},
           given={Alessandro},
           giveni={A\bibinitperiod}}}%
        {{hash=d8795f5e7db9cd4eb80a6631435281d1}{%
           family={Edwards},
           familyi={E\bibinitperiod},
           given={Alec},
           giveni={A\bibinitperiod}}}%
        {{hash=e227971e0a9f956b3d8f8688ae07fadb}{%
           family={Giacobbe},
           familyi={G\bibinitperiod},
           given={Mirco},
           giveni={M\bibinitperiod}}}%
        {{hash=ad6a85e0e8b63b54c59d3848d5e15a61}{%
           family={Punchihewa},
           familyi={P\bibinitperiod},
           given={Hashan},
           giveni={H\bibinitperiod}}}%
        {{hash=2bf051cfd5a39844fd4c6bd3516740c5}{%
           family={Roy},
           familyi={R\bibinitperiod},
           given={Diptarko},
           giveni={D\bibinitperiod}}}%
      }
      \strng{namehash}{c8ba5a51d46c8044ddbd1d31a54d4e97}
      \strng{fullhash}{c8ba5a51d46c8044ddbd1d31a54d4e97}
      \strng{fullhashraw}{c8ba5a51d46c8044ddbd1d31a54d4e97}
      \strng{bibnamehash}{c8ba5a51d46c8044ddbd1d31a54d4e97}
      \strng{authorbibnamehash}{c8ba5a51d46c8044ddbd1d31a54d4e97}
      \strng{authornamehash}{c8ba5a51d46c8044ddbd1d31a54d4e97}
      \strng{authorfullhash}{c8ba5a51d46c8044ddbd1d31a54d4e97}
      \strng{authorfullhashraw}{c8ba5a51d46c8044ddbd1d31a54d4e97}
      \field{sortinit}{A}
      \field{sortinithash}{2f401846e2029bad6b3ecc16d50031e2}
      \field{labelnamesource}{author}
      \field{labeltitlesource}{title}
      \field{eprinttype}{arXiv}
      \field{title}{Quantitative Verification with Neural Networks}
      \field{year}{2023}
      \verb{eprint}
      \verb 2301.06136
      \endverb
    \endentry
    \entry{abate26a}{online}{}{}
      \name{author}{4}{}{%
        {{hash=0f34d3654c8fa2fdffe70bbbaf378b28}{%
           family={Abate},
           familyi={A\bibinitperiod},
           given={Alessandro},
           giveni={A\bibinitperiod}}}%
        {{hash=e227971e0a9f956b3d8f8688ae07fadb}{%
           family={Giacobbe},
           familyi={G\bibinitperiod},
           given={Mirco},
           giveni={M\bibinitperiod}}}%
        {{hash=00087aa1436f3a08c179bc14e7b69a5e}{%
           family={Ichtchenko},
           familyi={I\bibinitperiod},
           given={Sergey},
           giveni={S\bibinitperiod}}}%
        {{hash=2bf051cfd5a39844fd4c6bd3516740c5}{%
           family={Roy},
           familyi={R\bibinitperiod},
           given={Diptarko},
           giveni={D\bibinitperiod}}}%
      }
      \strng{namehash}{270bdc66c8596cbdf72f5a5706bdc3ea}
      \strng{fullhash}{270bdc66c8596cbdf72f5a5706bdc3ea}
      \strng{fullhashraw}{270bdc66c8596cbdf72f5a5706bdc3ea}
      \strng{bibnamehash}{270bdc66c8596cbdf72f5a5706bdc3ea}
      \strng{authorbibnamehash}{270bdc66c8596cbdf72f5a5706bdc3ea}
      \strng{authornamehash}{270bdc66c8596cbdf72f5a5706bdc3ea}
      \strng{authorfullhash}{270bdc66c8596cbdf72f5a5706bdc3ea}
      \strng{authorfullhashraw}{270bdc66c8596cbdf72f5a5706bdc3ea}
      \field{sortinit}{A}
      \field{sortinithash}{2f401846e2029bad6b3ecc16d50031e2}
      \field{labelnamesource}{author}
      \field{labeltitlesource}{title}
      \field{eprinttype}{arXiv}
      \field{title}{Complete $\omega$-Regular Supermartingale Certificates}
      \field{year}{2026}
      \verb{eprint}
      \verb 2605.21134
      \endverb
    \endentry
    \entry{abate24}{online}{}{}
      \name{author}{3}{}{%
        {{hash=0f34d3654c8fa2fdffe70bbbaf378b28}{%
           family={Abate},
           familyi={A\bibinitperiod},
           given={Alessandro},
           giveni={A\bibinitperiod}}}%
        {{hash=e227971e0a9f956b3d8f8688ae07fadb}{%
           family={Giacobbe},
           familyi={G\bibinitperiod},
           given={Mirco},
           giveni={M\bibinitperiod}}}%
        {{hash=2bf051cfd5a39844fd4c6bd3516740c5}{%
           family={Roy},
           familyi={R\bibinitperiod},
           given={Diptarko},
           giveni={D\bibinitperiod}}}%
      }
      \strng{namehash}{b980029170859b8efd5145327d44ec17}
      \strng{fullhash}{b980029170859b8efd5145327d44ec17}
      \strng{fullhashraw}{b980029170859b8efd5145327d44ec17}
      \strng{bibnamehash}{b980029170859b8efd5145327d44ec17}
      \strng{authorbibnamehash}{b980029170859b8efd5145327d44ec17}
      \strng{authornamehash}{b980029170859b8efd5145327d44ec17}
      \strng{authorfullhash}{b980029170859b8efd5145327d44ec17}
      \strng{authorfullhashraw}{b980029170859b8efd5145327d44ec17}
      \field{sortinit}{A}
      \field{sortinithash}{2f401846e2029bad6b3ecc16d50031e2}
      \field{labelnamesource}{author}
      \field{labeltitlesource}{title}
      \field{eprinttype}{arXiv}
      \field{title}{Stochastic Omega-Regular Verification and Control with Supermartingales}
      \field{year}{2024}
      \verb{eprint}
      \verb 2405.17304
      \endverb
    \endentry
    \entry{abbadini23}{online}{}{}
      \name{author}{3}{}{%
        {{hash=34f039b8dc2d4925a1aa032107ffd6d5}{%
           family={Abbadini},
           familyi={A\bibinitperiod},
           given={Marco},
           giveni={M\bibinitperiod}}}%
        {{hash=a7c042aaa1a3fa49ceb88ffd5870c794}{%
           family={Bezhanishvili},
           familyi={B\bibinitperiod},
           given={Guram},
           giveni={G\bibinitperiod}}}%
        {{hash=be7811de0497714793798773fcc6fac5}{%
           family={Carai},
           familyi={C\bibinitperiod},
           given={Luca},
           giveni={L\bibinitperiod}}}%
      }
      \strng{namehash}{e541d1e744f3d9cf5b8412e1412e1997}
      \strng{fullhash}{e541d1e744f3d9cf5b8412e1412e1997}
      \strng{fullhashraw}{e541d1e744f3d9cf5b8412e1412e1997}
      \strng{bibnamehash}{e541d1e744f3d9cf5b8412e1412e1997}
      \strng{authorbibnamehash}{e541d1e744f3d9cf5b8412e1412e1997}
      \strng{authornamehash}{e541d1e744f3d9cf5b8412e1412e1997}
      \strng{authorfullhash}{e541d1e744f3d9cf5b8412e1412e1997}
      \strng{authorfullhashraw}{e541d1e744f3d9cf5b8412e1412e1997}
      \field{sortinit}{A}
      \field{sortinithash}{2f401846e2029bad6b3ecc16d50031e2}
      \field{labelnamesource}{author}
      \field{labeltitlesource}{title}
      \field{eprinttype}{arXiv}
      \field{title}{Vietoris endofunctor for closed relations and its de Vries dual}
      \field{year}{2023}
      \verb{eprint}
      \verb 2308.16823
      \endverb
    \endentry
    \entry{alvarezpicallo21}{online}{}{}
      \name{author}{4}{}{%
        {{hash=decf9b673b84ed84f13e97616606fbcc}{%
           family={Alvarez-Picallo},
           familyi={A\bibinithyphendelim P\bibinitperiod},
           given={Mario},
           giveni={M\bibinitperiod}}}%
        {{hash=935665c7060cc483b5b7a6a61817881c}{%
           family={Ghica},
           familyi={G\bibinitperiod},
           given={Dan\bibnamedelima R.},
           giveni={D\bibinitperiod\bibinitdelim R\bibinitperiod}}}%
        {{hash=380e0c80f9f1adacdad7785f82aca0e3}{%
           family={Sprunger},
           familyi={S\bibinitperiod},
           given={David},
           giveni={D\bibinitperiod}}}%
        {{hash=89c6b224e96f4982a42da5fa8512050b}{%
           family={Zanasi},
           familyi={Z\bibinitperiod},
           given={Fabio},
           giveni={F\bibinitperiod}}}%
      }
      \strng{namehash}{1757578aafb8481d5453e9dbedd2fdf7}
      \strng{fullhash}{1757578aafb8481d5453e9dbedd2fdf7}
      \strng{fullhashraw}{1757578aafb8481d5453e9dbedd2fdf7}
      \strng{bibnamehash}{1757578aafb8481d5453e9dbedd2fdf7}
      \strng{authorbibnamehash}{1757578aafb8481d5453e9dbedd2fdf7}
      \strng{authornamehash}{1757578aafb8481d5453e9dbedd2fdf7}
      \strng{authorfullhash}{1757578aafb8481d5453e9dbedd2fdf7}
      \strng{authorfullhashraw}{1757578aafb8481d5453e9dbedd2fdf7}
      \field{sortinit}{A}
      \field{sortinithash}{2f401846e2029bad6b3ecc16d50031e2}
      \field{labelnamesource}{author}
      \field{labeltitlesource}{title}
      \field{eprinttype}{arXiv}
      \field{title}{Functorial String Diagrams for Reverse-Mode Automatic Differentiation}
      \field{year}{2021}
      \verb{eprint}
      \verb 2107.13433
      \endverb
    \endentry
    \entry{astromfuruta2000}{article}{}{}
      \name{author}{2}{}{%
        {{hash=4d842a025a559e5c2edf98f59c10d29d}{%
           family={{Å}ström},
           familyi={Å\bibinitperiod},
           given={K.\bibnamedelimi J.},
           giveni={K\bibinitperiod\bibinitdelim J\bibinitperiod}}}%
        {{hash=d2c653ebab21ec82abd3219c3e223722}{%
           family={Furuta},
           familyi={F\bibinitperiod},
           given={K.},
           giveni={K\bibinitperiod}}}%
      }
      \strng{namehash}{132473c07b23eca4db16309c69531f12}
      \strng{fullhash}{132473c07b23eca4db16309c69531f12}
      \strng{fullhashraw}{132473c07b23eca4db16309c69531f12}
      \strng{bibnamehash}{132473c07b23eca4db16309c69531f12}
      \strng{authorbibnamehash}{132473c07b23eca4db16309c69531f12}
      \strng{authornamehash}{132473c07b23eca4db16309c69531f12}
      \strng{authorfullhash}{132473c07b23eca4db16309c69531f12}
      \strng{authorfullhashraw}{132473c07b23eca4db16309c69531f12}
      \field{sortinit}{Å}
      \field{sortinithash}{2f401846e2029bad6b3ecc16d50031e2}
      \field{labelnamesource}{author}
      \field{labeltitlesource}{title}
      \field{journaltitle}{Automatica}
      \field{number}{2}
      \field{title}{Swinging up a pendulum by energy control}
      \field{volume}{36}
      \field{year}{2000}
      \field{pages}{287\bibrangedash 295}
      \range{pages}{9}
    \endentry
    \entry{astrom-1965}{article}{}{}
      \name{author}{1}{}{%
        {{hash=f0f3bc9540d4770c9cd66454728924ca}{%
           family={{Å}ström},
           familyi={Å\bibinitperiod},
           given={Karl\bibnamedelima Johan},
           giveni={K\bibinitperiod\bibinitdelim J\bibinitperiod}}}%
      }
      \strng{namehash}{f0f3bc9540d4770c9cd66454728924ca}
      \strng{fullhash}{f0f3bc9540d4770c9cd66454728924ca}
      \strng{fullhashraw}{f0f3bc9540d4770c9cd66454728924ca}
      \strng{bibnamehash}{f0f3bc9540d4770c9cd66454728924ca}
      \strng{authorbibnamehash}{f0f3bc9540d4770c9cd66454728924ca}
      \strng{authornamehash}{f0f3bc9540d4770c9cd66454728924ca}
      \strng{authorfullhash}{f0f3bc9540d4770c9cd66454728924ca}
      \strng{authorfullhashraw}{f0f3bc9540d4770c9cd66454728924ca}
      \field{sortinit}{Å}
      \field{sortinithash}{2f401846e2029bad6b3ecc16d50031e2}
      \field{labelnamesource}{author}
      \field{labeltitlesource}{title}
      \field{journaltitle}{Journal of Mathematical Analysis and Applications}
      \field{number}{1}
      \field{title}{Optimal Control of {Markov} Processes with Incomplete State Information}
      \field{volume}{10}
      \field{year}{1965}
      \field{pages}{174\bibrangedash 205}
      \range{pages}{32}
      \verb{doi}
      \verb 10.1016/0022-247X(65)90154-X
      \endverb
    \endentry
    \entry{babiak2015hoa}{inproceedings}{}{}
      \name{author}{8}{}{%
        {{hash=3157ee3cd77facaee73a8a7e9f04973d}{%
           family={Babiak},
           familyi={B\bibinitperiod},
           given={T.},
           giveni={T\bibinitperiod}}}%
        {{hash=eb96ae329ea71b676d028565c1a2d5e4}{%
           family={Blahoudek},
           familyi={B\bibinitperiod},
           given={F.},
           giveni={F\bibinitperiod}}}%
        {{hash=00593bd008f2f733bd6d104189a0c95a}{%
           family={Duret-Lutz},
           familyi={D\bibinithyphendelim L\bibinitperiod},
           given={A.},
           giveni={A\bibinitperiod}}}%
        {{hash=650e1b6aabbb468fa24c5df75c76371e}{%
           family={Klein},
           familyi={K\bibinitperiod},
           given={J.},
           giveni={J\bibinitperiod}}}%
        {{hash=b6f511bf0242223a558c4f94df042bc2}{%
           family={Křetínský},
           familyi={K\bibinitperiod},
           given={J.},
           giveni={J\bibinitperiod}}}%
        {{hash=691b873ed0e0a7551ce133bf90aea191}{%
           family={Müller},
           familyi={M\bibinitperiod},
           given={D.},
           giveni={D\bibinitperiod}}}%
        {{hash=8f917e59f9bcb671f7f68671bf759fe7}{%
           family={Parker},
           familyi={P\bibinitperiod},
           given={D.},
           giveni={D\bibinitperiod}}}%
        {{hash=61e705dc26ec3f5e2a8134f7ee750aa6}{%
           family={Strejček},
           familyi={S\bibinitperiod},
           given={J.},
           giveni={J\bibinitperiod}}}%
      }
      \list{publisher}{1}{%
        {Springer}%
      }
      \strng{namehash}{31ea1afbdaf1f9bd9249fa61295a58ab}
      \strng{fullhash}{c738470fcd30281c76124193e2ed2a44}
      \strng{fullhashraw}{c738470fcd30281c76124193e2ed2a44}
      \strng{bibnamehash}{31ea1afbdaf1f9bd9249fa61295a58ab}
      \strng{authorbibnamehash}{31ea1afbdaf1f9bd9249fa61295a58ab}
      \strng{authornamehash}{31ea1afbdaf1f9bd9249fa61295a58ab}
      \strng{authorfullhash}{c738470fcd30281c76124193e2ed2a44}
      \strng{authorfullhashraw}{c738470fcd30281c76124193e2ed2a44}
      \field{sortinit}{B}
      \field{sortinithash}{d7095fff47cda75ca2589920aae98399}
      \field{labelnamesource}{author}
      \field{labeltitlesource}{title}
      \field{booktitle}{Computer Aided Verification (CAV 2015), Part I}
      \field{note}{(The $\mathrm{Fin}(\cdot)$/$\mathrm{Inf}(\cdot)$ acceptance atoms as an interchange standard.)}
      \field{series}{Lecture Notes in Computer Science}
      \field{title}{The {Hanoi} omega-automata format}
      \field{volume}{9206}
      \field{year}{2015}
      \field{pages}{479\bibrangedash 486}
      \range{pages}{8}
    \endentry
    \entry{baez18}{online}{}{}
      \name{author}{2}{}{%
        {{hash=77c130b1ed4688a565b29481eee08ae3}{%
           family={Baez},
           familyi={B\bibinitperiod},
           given={John\bibnamedelima C.},
           giveni={J\bibinitperiod\bibinitdelim C\bibinitperiod}}}%
        {{hash=fa8912eb6f0f0c230fc8d8371db35a2f}{%
           family={Master},
           familyi={M\bibinitperiod},
           given={Jade},
           giveni={J\bibinitperiod}}}%
      }
      \strng{namehash}{244659d8d0fec18588c278ba5f67024f}
      \strng{fullhash}{244659d8d0fec18588c278ba5f67024f}
      \strng{fullhashraw}{244659d8d0fec18588c278ba5f67024f}
      \strng{bibnamehash}{244659d8d0fec18588c278ba5f67024f}
      \strng{authorbibnamehash}{244659d8d0fec18588c278ba5f67024f}
      \strng{authornamehash}{244659d8d0fec18588c278ba5f67024f}
      \strng{authorfullhash}{244659d8d0fec18588c278ba5f67024f}
      \strng{authorfullhashraw}{244659d8d0fec18588c278ba5f67024f}
      \field{sortinit}{B}
      \field{sortinithash}{d7095fff47cda75ca2589920aae98399}
      \field{labelnamesource}{author}
      \field{labeltitlesource}{title}
      \field{eprinttype}{arXiv}
      \field{title}{Open Petri Nets}
      \field{year}{2018}
      \verb{eprint}
      \verb 1808.05415
      \endverb
    \endentry
    \entry{baier18hbmc}{incollection}{}{}
      \name{author}{4}{}{%
        {{hash=7725f3db070a985fe46a5feb0bcee960}{%
           family={Baier},
           familyi={B\bibinitperiod},
           given={Christel},
           giveni={C\bibinitperiod}}}%
        {{hash=d8a3360195b18eaefe5b3686604cd41d}{%
           family={Alfaro},
           familyi={A\bibinitperiod},
           given={Luca},
           giveni={L\bibinitperiod},
           prefix={de},
           prefixi={d\bibinitperiod}}}%
        {{hash=a293acdf62a9e3af0cc3d7a2e60ef5aa}{%
           family={Forejt},
           familyi={F\bibinitperiod},
           given={Vojtěch},
           giveni={V\bibinitperiod}}}%
        {{hash=e96e14c940cd49f1c4e8c1f38f89a321}{%
           family={Kwiatkowska},
           familyi={K\bibinitperiod},
           given={Marta},
           giveni={M\bibinitperiod}}}%
      }
      \name{editor}{4}{}{%
        {{hash=1822cd554880ea55fd66357a96bca03e}{%
           family={Clarke},
           familyi={C\bibinitperiod},
           given={Edmund\bibnamedelima M.},
           giveni={E\bibinitperiod\bibinitdelim M\bibinitperiod}}}%
        {{hash=eeb928073fbee58751f87d401a62561a}{%
           family={Henzinger},
           familyi={H\bibinitperiod},
           given={Thomas\bibnamedelima A.},
           giveni={T\bibinitperiod\bibinitdelim A\bibinitperiod}}}%
        {{hash=5ac94523768e8fc4adb7a926b36d83e0}{%
           family={Veith},
           familyi={V\bibinitperiod},
           given={Helmut},
           giveni={H\bibinitperiod}}}%
        {{hash=9865af1b742410e6dee8ca813e1ec733}{%
           family={Bloem},
           familyi={B\bibinitperiod},
           given={Roderick},
           giveni={R\bibinitperiod}}}%
      }
      \list{publisher}{1}{%
        {Springer}%
      }
      \strng{namehash}{df76e178ef68aa12a0990d3d06ca2b77}
      \strng{fullhash}{df76e178ef68aa12a0990d3d06ca2b77}
      \strng{fullhashraw}{df76e178ef68aa12a0990d3d06ca2b77}
      \strng{bibnamehash}{df76e178ef68aa12a0990d3d06ca2b77}
      \strng{authorbibnamehash}{df76e178ef68aa12a0990d3d06ca2b77}
      \strng{authornamehash}{df76e178ef68aa12a0990d3d06ca2b77}
      \strng{authorfullhash}{df76e178ef68aa12a0990d3d06ca2b77}
      \strng{authorfullhashraw}{df76e178ef68aa12a0990d3d06ca2b77}
      \strng{editorbibnamehash}{eea0cb89759dac86c5d80411829d95b2}
      \strng{editornamehash}{eea0cb89759dac86c5d80411829d95b2}
      \strng{editorfullhash}{eea0cb89759dac86c5d80411829d95b2}
      \strng{editorfullhashraw}{eea0cb89759dac86c5d80411829d95b2}
      \field{sortinit}{B}
      \field{sortinithash}{d7095fff47cda75ca2589920aae98399}
      \field{labelnamesource}{author}
      \field{labeltitlesource}{title}
      \field{booktitle}{Handbook of Model Checking}
      \field{title}{Model Checking Probabilistic Systems}
      \field{year}{2018}
      \field{pages}{963\bibrangedash 999}
      \range{pages}{37}
      \verb{doi}
      \verb 10.1007/978-3-319-10575-8_28
      \endverb
    \endentry
    \entry{baier-katoen-2008}{book}{}{}
      \name{author}{2}{}{%
        {{hash=7725f3db070a985fe46a5feb0bcee960}{%
           family={Baier},
           familyi={B\bibinitperiod},
           given={Christel},
           giveni={C\bibinitperiod}}}%
        {{hash=392d5825b15ef7f81f68632c7a49fdef}{%
           family={Katoen},
           familyi={K\bibinitperiod},
           given={Joost-Pieter},
           giveni={J\bibinithyphendelim P\bibinitperiod}}}%
      }
      \list{publisher}{1}{%
        {MIT Press}%
      }
      \strng{namehash}{89a25b516e2b4c8733def01ee561abfe}
      \strng{fullhash}{89a25b516e2b4c8733def01ee561abfe}
      \strng{fullhashraw}{89a25b516e2b4c8733def01ee561abfe}
      \strng{bibnamehash}{89a25b516e2b4c8733def01ee561abfe}
      \strng{authorbibnamehash}{89a25b516e2b4c8733def01ee561abfe}
      \strng{authornamehash}{89a25b516e2b4c8733def01ee561abfe}
      \strng{authorfullhash}{89a25b516e2b4c8733def01ee561abfe}
      \strng{authorfullhashraw}{89a25b516e2b4c8733def01ee561abfe}
      \field{sortinit}{B}
      \field{sortinithash}{d7095fff47cda75ca2589920aae98399}
      \field{labelnamesource}{author}
      \field{labeltitlesource}{title}
      \field{title}{Principles of Model Checking}
      \field{year}{2008}
    \endentry
    \entry{banach1922}{article}{}{}
      \name{author}{1}{}{%
        {{hash=4c950a592a442a5fb66732768b037911}{%
           family={Banach},
           familyi={B\bibinitperiod},
           given={S.},
           giveni={S\bibinitperiod}}}%
      }
      \strng{namehash}{4c950a592a442a5fb66732768b037911}
      \strng{fullhash}{4c950a592a442a5fb66732768b037911}
      \strng{fullhashraw}{4c950a592a442a5fb66732768b037911}
      \strng{bibnamehash}{4c950a592a442a5fb66732768b037911}
      \strng{authorbibnamehash}{4c950a592a442a5fb66732768b037911}
      \strng{authornamehash}{4c950a592a442a5fb66732768b037911}
      \strng{authorfullhash}{4c950a592a442a5fb66732768b037911}
      \strng{authorfullhashraw}{4c950a592a442a5fb66732768b037911}
      \field{sortinit}{B}
      \field{sortinithash}{d7095fff47cda75ca2589920aae98399}
      \field{labelnamesource}{author}
      \field{labeltitlesource}{title}
      \field{journaltitle}{Fundamenta Mathematicae}
      \field{title}{Sur les opérations dans les ensembles abstraits et leur application aux équations intégrales}
      \field{volume}{3}
      \field{year}{1922}
      \field{pages}{133\bibrangedash 181}
      \range{pages}{49}
    \endentry
    \entry{bentley2000pearls}{book}{}{}
      \name{author}{1}{}{%
        {{hash=a2ea6a4e7fe1b839acb38db376905df8}{%
           family={Bentley},
           familyi={B\bibinitperiod},
           given={Jon},
           giveni={J\bibinitperiod}}}%
      }
      \list{publisher}{1}{%
        {Addison-Wesley}%
      }
      \strng{namehash}{a2ea6a4e7fe1b839acb38db376905df8}
      \strng{fullhash}{a2ea6a4e7fe1b839acb38db376905df8}
      \strng{fullhashraw}{a2ea6a4e7fe1b839acb38db376905df8}
      \strng{bibnamehash}{a2ea6a4e7fe1b839acb38db376905df8}
      \strng{authorbibnamehash}{a2ea6a4e7fe1b839acb38db376905df8}
      \strng{authornamehash}{a2ea6a4e7fe1b839acb38db376905df8}
      \strng{authorfullhash}{a2ea6a4e7fe1b839acb38db376905df8}
      \strng{authorfullhashraw}{a2ea6a4e7fe1b839acb38db376905df8}
      \field{sortinit}{B}
      \field{sortinithash}{d7095fff47cda75ca2589920aae98399}
      \field{labelnamesource}{author}
      \field{labeltitlesource}{title}
      \field{edition}{2}
      \field{title}{Programming Pearls}
      \field{year}{2000}
    \endentry
    \entry{vandenberk04}{thesis}{}{}
      \name{author}{1}{}{%
        {{hash=dff8f8486d0da8914a31d79d1dd1de56}{%
           family={Berk},
           familyi={B\bibinitperiod},
           given={V.},
           giveni={V\bibinitperiod},
           prefix={van\bibnamedelima den},
           prefixi={v\bibinitperiod\bibinitdelim d\bibinitperiod}}}%
      }
      \list{institution}{1}{%
        {Eindhoven University of Technology}%
      }
      \strng{namehash}{dff8f8486d0da8914a31d79d1dd1de56}
      \strng{fullhash}{dff8f8486d0da8914a31d79d1dd1de56}
      \strng{fullhashraw}{dff8f8486d0da8914a31d79d1dd1de56}
      \strng{bibnamehash}{dff8f8486d0da8914a31d79d1dd1de56}
      \strng{authorbibnamehash}{dff8f8486d0da8914a31d79d1dd1de56}
      \strng{authornamehash}{dff8f8486d0da8914a31d79d1dd1de56}
      \strng{authorfullhash}{dff8f8486d0da8914a31d79d1dd1de56}
      \strng{authorfullhashraw}{dff8f8486d0da8914a31d79d1dd1de56}
      \field{sortinit}{B}
      \field{sortinithash}{d7095fff47cda75ca2589920aae98399}
      \field{labelnamesource}{author}
      \field{labeltitlesource}{title}
      \field{title}{Wafer fab performance analysis}
      \field{type}{Master's thesis}
      \field{year}{2004}
    \endentry
    \entry{bernstein1912}{article}{}{}
      \name{author}{1}{}{%
        {{hash=6d599dc61ac3ab0143e33503dbdd6f5f}{%
           family={Bernstein},
           familyi={B\bibinitperiod},
           given={S.\bibnamedelimi N.},
           giveni={S\bibinitperiod\bibinitdelim N\bibinitperiod}}}%
      }
      \strng{namehash}{6d599dc61ac3ab0143e33503dbdd6f5f}
      \strng{fullhash}{6d599dc61ac3ab0143e33503dbdd6f5f}
      \strng{fullhashraw}{6d599dc61ac3ab0143e33503dbdd6f5f}
      \strng{bibnamehash}{6d599dc61ac3ab0143e33503dbdd6f5f}
      \strng{authorbibnamehash}{6d599dc61ac3ab0143e33503dbdd6f5f}
      \strng{authornamehash}{6d599dc61ac3ab0143e33503dbdd6f5f}
      \strng{authorfullhash}{6d599dc61ac3ab0143e33503dbdd6f5f}
      \strng{authorfullhashraw}{6d599dc61ac3ab0143e33503dbdd6f5f}
      \field{sortinit}{B}
      \field{sortinithash}{d7095fff47cda75ca2589920aae98399}
      \field{labelnamesource}{author}
      \field{labeltitlesource}{title}
      \field{journaltitle}{Communications de la Société Mathématique de Kharkov}
      \field{series}{2}
      \field{title}{Démonstration du théorème de Weierstrass fondée sur le calcul des probabilités}
      \field{volume}{13}
      \field{year}{1912}
      \field{pages}{1\bibrangedash 2}
      \range{pages}{2}
    \endentry
    \entry{beutner22}{online}{}{}
      \name{author}{3}{}{%
        {{hash=fde3adc416b5cba0517bbff0f2ca92eb}{%
           family={Beutner},
           familyi={B\bibinitperiod},
           given={Raven},
           giveni={R\bibinitperiod}}}%
        {{hash=4b490249836521bc0e7ec27d4afad459}{%
           family={Ong},
           familyi={O\bibinitperiod},
           given={C.-H.\bibnamedelimi Luke},
           giveni={C\bibinithyphendelim H\bibinitperiod\bibinitdelim L\bibinitperiod}}}%
        {{hash=2c47a50416bae4df0ff499bf813a13c3}{%
           family={Zaiser},
           familyi={Z\bibinitperiod},
           given={Fabian},
           giveni={F\bibinitperiod}}}%
      }
      \strng{namehash}{fe428d13c3d08206011b20cd4ca0977e}
      \strng{fullhash}{fe428d13c3d08206011b20cd4ca0977e}
      \strng{fullhashraw}{fe428d13c3d08206011b20cd4ca0977e}
      \strng{bibnamehash}{fe428d13c3d08206011b20cd4ca0977e}
      \strng{authorbibnamehash}{fe428d13c3d08206011b20cd4ca0977e}
      \strng{authornamehash}{fe428d13c3d08206011b20cd4ca0977e}
      \strng{authorfullhash}{fe428d13c3d08206011b20cd4ca0977e}
      \strng{authorfullhashraw}{fe428d13c3d08206011b20cd4ca0977e}
      \field{sortinit}{B}
      \field{sortinithash}{d7095fff47cda75ca2589920aae98399}
      \field{labelnamesource}{author}
      \field{labeltitlesource}{title}
      \field{eprinttype}{arXiv}
      \field{title}{Guaranteed Bounds for Posterior Inference in Universal Probabilistic Programming}
      \field{year}{2022}
      \verb{eprint}
      \verb 2204.02948
      \endverb
    \endentry
    \entry{bidlingmaier23}{online}{}{}
      \name{author}{1}{}{%
        {{hash=6b38b9ed855f237c8dab0fae74f2ed23}{%
           family={Bidlingmaier},
           familyi={B\bibinitperiod},
           given={Martin\bibnamedelima E.},
           giveni={M\bibinitperiod\bibinitdelim E\bibinitperiod}}}%
      }
      \strng{namehash}{6b38b9ed855f237c8dab0fae74f2ed23}
      \strng{fullhash}{6b38b9ed855f237c8dab0fae74f2ed23}
      \strng{fullhashraw}{6b38b9ed855f237c8dab0fae74f2ed23}
      \strng{bibnamehash}{6b38b9ed855f237c8dab0fae74f2ed23}
      \strng{authorbibnamehash}{6b38b9ed855f237c8dab0fae74f2ed23}
      \strng{authornamehash}{6b38b9ed855f237c8dab0fae74f2ed23}
      \strng{authorfullhash}{6b38b9ed855f237c8dab0fae74f2ed23}
      \strng{authorfullhashraw}{6b38b9ed855f237c8dab0fae74f2ed23}
      \field{sortinit}{B}
      \field{sortinithash}{d7095fff47cda75ca2589920aae98399}
      \field{labelnamesource}{author}
      \field{labeltitlesource}{title}
      \field{eprinttype}{arXiv}
      \field{title}{Algebraic Semantics of Datalog with Equality}
      \field{year}{2023}
      \verb{eprint}
      \verb 2302.03167
      \endverb
    \endentry
    \entry{birkhoff1967}{book}{}{}
      \name{author}{1}{}{%
        {{hash=fd4941cce3b0766c7673191a918e714c}{%
           family={Birkhoff},
           familyi={B\bibinitperiod},
           given={G.},
           giveni={G\bibinitperiod}}}%
      }
      \strng{namehash}{fd4941cce3b0766c7673191a918e714c}
      \strng{fullhash}{fd4941cce3b0766c7673191a918e714c}
      \strng{fullhashraw}{fd4941cce3b0766c7673191a918e714c}
      \strng{bibnamehash}{fd4941cce3b0766c7673191a918e714c}
      \strng{authorbibnamehash}{fd4941cce3b0766c7673191a918e714c}
      \strng{authornamehash}{fd4941cce3b0766c7673191a918e714c}
      \strng{authorfullhash}{fd4941cce3b0766c7673191a918e714c}
      \strng{authorfullhashraw}{fd4941cce3b0766c7673191a918e714c}
      \field{extraname}{1}
      \field{sortinit}{B}
      \field{sortinithash}{d7095fff47cda75ca2589920aae98399}
      \field{labelnamesource}{author}
      \field{labeltitlesource}{title}
      \field{edition}{3}
      \field{note}{(M$_3$, the diamond, as the canonical modular non-distributive lattice.)}
      \field{number}{25}
      \field{series}{American Mathematical Society Colloquium Publications}
      \field{title}{Lattice Theory}
      \field{year}{1967}
    \endentry
    \entry{birkhoff1937}{article}{}{}
      \name{author}{1}{}{%
        {{hash=fd4941cce3b0766c7673191a918e714c}{%
           family={Birkhoff},
           familyi={B\bibinitperiod},
           given={G.},
           giveni={G\bibinitperiod}}}%
      }
      \strng{namehash}{fd4941cce3b0766c7673191a918e714c}
      \strng{fullhash}{fd4941cce3b0766c7673191a918e714c}
      \strng{fullhashraw}{fd4941cce3b0766c7673191a918e714c}
      \strng{bibnamehash}{fd4941cce3b0766c7673191a918e714c}
      \strng{authorbibnamehash}{fd4941cce3b0766c7673191a918e714c}
      \strng{authornamehash}{fd4941cce3b0766c7673191a918e714c}
      \strng{authorfullhash}{fd4941cce3b0766c7673191a918e714c}
      \strng{authorfullhashraw}{fd4941cce3b0766c7673191a918e714c}
      \field{extraname}{2}
      \field{sortinit}{B}
      \field{sortinithash}{d7095fff47cda75ca2589920aae98399}
      \field{labelnamesource}{author}
      \field{labeltitlesource}{title}
      \field{journaltitle}{Duke Mathematical Journal}
      \field{note}{(The representation theorem for finite distributive lattices via join-irreducibles.)}
      \field{number}{3}
      \field{title}{Rings of sets}
      \field{volume}{3}
      \field{year}{1937}
      \field{pages}{443\bibrangedash 454}
      \range{pages}{12}
    \endentry
    \entry{blanchini1999}{article}{}{}
      \name{author}{1}{}{%
        {{hash=f9f343c5e703f39c94719b3280aff64c}{%
           family={Blanchini},
           familyi={B\bibinitperiod},
           given={F.},
           giveni={F\bibinitperiod}}}%
      }
      \strng{namehash}{f9f343c5e703f39c94719b3280aff64c}
      \strng{fullhash}{f9f343c5e703f39c94719b3280aff64c}
      \strng{fullhashraw}{f9f343c5e703f39c94719b3280aff64c}
      \strng{bibnamehash}{f9f343c5e703f39c94719b3280aff64c}
      \strng{authorbibnamehash}{f9f343c5e703f39c94719b3280aff64c}
      \strng{authornamehash}{f9f343c5e703f39c94719b3280aff64c}
      \strng{authorfullhash}{f9f343c5e703f39c94719b3280aff64c}
      \strng{authorfullhashraw}{f9f343c5e703f39c94719b3280aff64c}
      \field{sortinit}{B}
      \field{sortinithash}{d7095fff47cda75ca2589920aae98399}
      \field{labelnamesource}{author}
      \field{labeltitlesource}{title}
      \field{journaltitle}{Automatica}
      \field{number}{11}
      \field{title}{Set invariance in control}
      \field{volume}{35}
      \field{year}{1999}
      \field{pages}{1747\bibrangedash 1767}
      \range{pages}{21}
    \endentry
    \entry{boman2005}{article}{}{}
      \name{author}{4}{}{%
        {{hash=1fba1a518c53c0d4f1b95cce726c0932}{%
           family={Boman},
           familyi={B\bibinitperiod},
           given={E.\bibnamedelimi G.},
           giveni={E\bibinitperiod\bibinitdelim G\bibinitperiod}}}%
        {{hash=80489b75bcde120c9f335f4995d789c2}{%
           family={Chen},
           familyi={C\bibinitperiod},
           given={D.},
           giveni={D\bibinitperiod}}}%
        {{hash=2ff3766700b17d178a4baf14c27cc6f9}{%
           family={Parekh},
           familyi={P\bibinitperiod},
           given={O.},
           giveni={O\bibinitperiod}}}%
        {{hash=f2aba2329fadac53fbfc643b1498c581}{%
           family={Toledo},
           familyi={T\bibinitperiod},
           given={S.},
           giveni={S\bibinitperiod}}}%
      }
      \strng{namehash}{d1b1dbce0ee0bfd169b73cd1eb0a6ae6}
      \strng{fullhash}{d1b1dbce0ee0bfd169b73cd1eb0a6ae6}
      \strng{fullhashraw}{d1b1dbce0ee0bfd169b73cd1eb0a6ae6}
      \strng{bibnamehash}{d1b1dbce0ee0bfd169b73cd1eb0a6ae6}
      \strng{authorbibnamehash}{d1b1dbce0ee0bfd169b73cd1eb0a6ae6}
      \strng{authornamehash}{d1b1dbce0ee0bfd169b73cd1eb0a6ae6}
      \strng{authorfullhash}{d1b1dbce0ee0bfd169b73cd1eb0a6ae6}
      \strng{authorfullhashraw}{d1b1dbce0ee0bfd169b73cd1eb0a6ae6}
      \field{sortinit}{B}
      \field{sortinithash}{d7095fff47cda75ca2589920aae98399}
      \field{labelnamesource}{author}
      \field{labeltitlesource}{title}
      \field{journaltitle}{Linear Algebra and its Applications}
      \field{note}{(Factor-width-2 = scaled diagonal dominance --- the class the /DD tier's congruence certificates inhabit.)}
      \field{title}{On factor width and symmetric H-matrices}
      \field{volume}{405}
      \field{year}{2005}
      \field{pages}{239\bibrangedash 248}
      \range{pages}{10}
    \endentry
    \entry{bonchi26b}{online}{}{}
      \name{author}{2}{}{%
        {{hash=bf27f2176d4e71dccf8f7cb8912be2ae}{%
           family={Bonchi},
           familyi={B\bibinitperiod},
           given={Filippo},
           giveni={F\bibinitperiod}}}%
        {{hash=8fa5f86b148b68eb78a76f39080a90eb}{%
           family={Cioffo},
           familyi={C\bibinitperiod},
           given={Cipriano\bibnamedelima Junior},
           giveni={C\bibinitperiod\bibinitdelim J\bibinitperiod}}}%
      }
      \strng{namehash}{ed20f94cd6a1c558a6d2d78e6ba8233f}
      \strng{fullhash}{ed20f94cd6a1c558a6d2d78e6ba8233f}
      \strng{fullhashraw}{ed20f94cd6a1c558a6d2d78e6ba8233f}
      \strng{bibnamehash}{ed20f94cd6a1c558a6d2d78e6ba8233f}
      \strng{authorbibnamehash}{ed20f94cd6a1c558a6d2d78e6ba8233f}
      \strng{authornamehash}{ed20f94cd6a1c558a6d2d78e6ba8233f}
      \strng{authorfullhash}{ed20f94cd6a1c558a6d2d78e6ba8233f}
      \strng{authorfullhashraw}{ed20f94cd6a1c558a6d2d78e6ba8233f}
      \field{sortinit}{B}
      \field{sortinithash}{d7095fff47cda75ca2589920aae98399}
      \field{labelnamesource}{author}
      \field{labeltitlesource}{title}
      \field{eprinttype}{arXiv}
      \field{title}{Completeness for Probabilistic Boolean Tapes}
      \field{year}{2026}
      \verb{eprint}
      \verb 2606.19017
      \endverb
    \endentry
    \entry{bonchi20}{online}{}{}
      \name{author}{5}{}{%
        {{hash=bf27f2176d4e71dccf8f7cb8912be2ae}{%
           family={Bonchi},
           familyi={B\bibinitperiod},
           given={Filippo},
           giveni={F\bibinitperiod}}}%
        {{hash=e443749574bbef3a3db853b22ad42896}{%
           family={Gadducci},
           familyi={G\bibinitperiod},
           given={Fabio},
           giveni={F\bibinitperiod}}}%
        {{hash=0dcc317455baeefdf9b65140659af27d}{%
           family={Kissinger},
           familyi={K\bibinitperiod},
           given={Aleks},
           giveni={A\bibinitperiod}}}%
        {{hash=59470481aaa6749da6f4102f4e2e1e14}{%
           family={Sobociński},
           familyi={S\bibinitperiod},
           given={Paweł},
           giveni={P\bibinitperiod}}}%
        {{hash=89c6b224e96f4982a42da5fa8512050b}{%
           family={Zanasi},
           familyi={Z\bibinitperiod},
           given={Fabio},
           giveni={F\bibinitperiod}}}%
      }
      \strng{namehash}{c3be3d9264eeb960e778b06df04da213}
      \strng{fullhash}{c3be3d9264eeb960e778b06df04da213}
      \strng{fullhashraw}{c3be3d9264eeb960e778b06df04da213}
      \strng{bibnamehash}{c3be3d9264eeb960e778b06df04da213}
      \strng{authorbibnamehash}{c3be3d9264eeb960e778b06df04da213}
      \strng{authornamehash}{c3be3d9264eeb960e778b06df04da213}
      \strng{authorfullhash}{c3be3d9264eeb960e778b06df04da213}
      \strng{authorfullhashraw}{c3be3d9264eeb960e778b06df04da213}
      \field{extraname}{1}
      \field{sortinit}{B}
      \field{sortinithash}{d7095fff47cda75ca2589920aae98399}
      \field{labelnamesource}{author}
      \field{labeltitlesource}{title}
      \field{eprinttype}{arXiv}
      \field{subtitle}{Rewriting with Frobenius Structure}
      \field{title}{String Diagram Rewrite Theory I}
      \field{year}{2020}
      \verb{eprint}
      \verb 2012.01847
      \endverb
    \endentry
    \entry{bonchi22}{online}{}{}
      \name{author}{5}{}{%
        {{hash=bf27f2176d4e71dccf8f7cb8912be2ae}{%
           family={Bonchi},
           familyi={B\bibinitperiod},
           given={Filippo},
           giveni={F\bibinitperiod}}}%
        {{hash=e443749574bbef3a3db853b22ad42896}{%
           family={Gadducci},
           familyi={G\bibinitperiod},
           given={Fabio},
           giveni={F\bibinitperiod}}}%
        {{hash=0dcc317455baeefdf9b65140659af27d}{%
           family={Kissinger},
           familyi={K\bibinitperiod},
           given={Aleks},
           giveni={A\bibinitperiod}}}%
        {{hash=59470481aaa6749da6f4102f4e2e1e14}{%
           family={Sobociński},
           familyi={S\bibinitperiod},
           given={Paweł},
           giveni={P\bibinitperiod}}}%
        {{hash=89c6b224e96f4982a42da5fa8512050b}{%
           family={Zanasi},
           familyi={Z\bibinitperiod},
           given={Fabio},
           giveni={F\bibinitperiod}}}%
      }
      \strng{namehash}{c3be3d9264eeb960e778b06df04da213}
      \strng{fullhash}{c3be3d9264eeb960e778b06df04da213}
      \strng{fullhashraw}{c3be3d9264eeb960e778b06df04da213}
      \strng{bibnamehash}{c3be3d9264eeb960e778b06df04da213}
      \strng{authorbibnamehash}{c3be3d9264eeb960e778b06df04da213}
      \strng{authornamehash}{c3be3d9264eeb960e778b06df04da213}
      \strng{authorfullhash}{c3be3d9264eeb960e778b06df04da213}
      \strng{authorfullhashraw}{c3be3d9264eeb960e778b06df04da213}
      \field{extraname}{2}
      \field{sortinit}{B}
      \field{sortinithash}{d7095fff47cda75ca2589920aae98399}
      \field{labelnamesource}{author}
      \field{labeltitlesource}{title}
      \field{eprinttype}{arXiv}
      \field{subtitle}{Rewriting with Symmetric Monoidal Structure}
      \field{title}{String Diagram Rewrite Theory II}
      \field{year}{2022}
      \verb{eprint}
      \verb 2104.14686
      \endverb
    \endentry
    \entry{brieger24}{online}{}{}
      \name{author}{3}{}{%
        {{hash=ed55fa8b350d9ccc9c85a6ee96637b9e}{%
           family={Brieger},
           familyi={B\bibinitperiod},
           given={Marvin},
           giveni={M\bibinitperiod}}}%
        {{hash=b723c481b381d20105178d693c40d93b}{%
           family={Mitsch},
           familyi={M\bibinitperiod},
           given={Stefan},
           giveni={S\bibinitperiod}}}%
        {{hash=0fb7e4129542e3e84c6115bee2a2b227}{%
           family={Platzer},
           familyi={P\bibinitperiod},
           given={André},
           giveni={A\bibinitperiod}}}%
      }
      \strng{namehash}{470eaaa7581ba6a99c036d11d7d52c31}
      \strng{fullhash}{470eaaa7581ba6a99c036d11d7d52c31}
      \strng{fullhashraw}{470eaaa7581ba6a99c036d11d7d52c31}
      \strng{bibnamehash}{470eaaa7581ba6a99c036d11d7d52c31}
      \strng{authorbibnamehash}{470eaaa7581ba6a99c036d11d7d52c31}
      \strng{authornamehash}{470eaaa7581ba6a99c036d11d7d52c31}
      \strng{authorfullhash}{470eaaa7581ba6a99c036d11d7d52c31}
      \strng{authorfullhashraw}{470eaaa7581ba6a99c036d11d7d52c31}
      \field{sortinit}{B}
      \field{sortinithash}{d7095fff47cda75ca2589920aae98399}
      \field{labelnamesource}{author}
      \field{labeltitlesource}{title}
      \field{eprinttype}{arXiv}
      \field{title}{Complete Dynamic Logic of Communicating Hybrid Programs}
      \field{year}{2024}
      \verb{eprint}
      \verb 2408.05012
      \endverb
    \endentry
    \entry{brunetti01}{online}{}{}
      \name{author}{3}{}{%
        {{hash=a44ca3908a44ef5857eb6c0542961cd5}{%
           family={Brunetti},
           familyi={B\bibinitperiod},
           given={Romeo},
           giveni={R\bibinitperiod}}}%
        {{hash=de33ee3f83484992450dd5ea4b66e283}{%
           family={Fredenhagen},
           familyi={F\bibinitperiod},
           given={Klaus},
           giveni={K\bibinitperiod}}}%
        {{hash=fabe445ea043cc87cdde933d33ee4aac}{%
           family={Verch},
           familyi={V\bibinitperiod},
           given={Rainer},
           giveni={R\bibinitperiod}}}%
      }
      \strng{namehash}{48169a37856e5b16f2fda224bc0dcc64}
      \strng{fullhash}{48169a37856e5b16f2fda224bc0dcc64}
      \strng{fullhashraw}{48169a37856e5b16f2fda224bc0dcc64}
      \strng{bibnamehash}{48169a37856e5b16f2fda224bc0dcc64}
      \strng{authorbibnamehash}{48169a37856e5b16f2fda224bc0dcc64}
      \strng{authornamehash}{48169a37856e5b16f2fda224bc0dcc64}
      \strng{authorfullhash}{48169a37856e5b16f2fda224bc0dcc64}
      \strng{authorfullhashraw}{48169a37856e5b16f2fda224bc0dcc64}
      \field{sortinit}{B}
      \field{sortinithash}{d7095fff47cda75ca2589920aae98399}
      \field{labelnamesource}{author}
      \field{labeltitlesource}{title}
      \field{eprinttype}{arXiv}
      \field{subtitle}{A new paradigm for local quantum physics}
      \field{title}{The generally covariant locality principle}
      \field{year}{2001}
      \verb{doi}
      \verb 10.1007/s00220-003-0815-7
      \endverb
      \verb{eprint}
      \verb math-ph/0112041
      \endverb
    \endentry
    \entry{budde24}{online}{}{}
      \name{author}{5}{}{%
        {{hash=ed94d773d693a01aa6ee5a7592095934}{%
           family={Budde},
           familyi={B\bibinitperiod},
           given={Carlos\bibnamedelima E.},
           giveni={C\bibinitperiod\bibinitdelim E\bibinitperiod}}}%
        {{hash=5ee8e32c129a82f9fe67944701c268bf}{%
           family={Hartmanns},
           familyi={H\bibinitperiod},
           given={Arnd},
           giveni={A\bibinitperiod}}}%
        {{hash=a161b8e5382db2f09d25dfb1f43ef76e}{%
           family={Meggendorfer},
           familyi={M\bibinitperiod},
           given={Tobias},
           giveni={T\bibinitperiod}}}%
        {{hash=aafb182186c9faf19232ee74fb85cb32}{%
           family={Weininger},
           familyi={W\bibinitperiod},
           given={Maximilian},
           giveni={M\bibinitperiod}}}%
        {{hash=a2373bacacffc05c8e523a457719a62a}{%
           family={Wienhöft},
           familyi={W\bibinitperiod},
           given={Patrick},
           giveni={P\bibinitperiod}}}%
      }
      \strng{namehash}{742e4905bffb823fa82bcd3aaf5cadf8}
      \strng{fullhash}{742e4905bffb823fa82bcd3aaf5cadf8}
      \strng{fullhashraw}{742e4905bffb823fa82bcd3aaf5cadf8}
      \strng{bibnamehash}{742e4905bffb823fa82bcd3aaf5cadf8}
      \strng{authorbibnamehash}{742e4905bffb823fa82bcd3aaf5cadf8}
      \strng{authornamehash}{742e4905bffb823fa82bcd3aaf5cadf8}
      \strng{authorfullhash}{742e4905bffb823fa82bcd3aaf5cadf8}
      \strng{authorfullhashraw}{742e4905bffb823fa82bcd3aaf5cadf8}
      \field{sortinit}{B}
      \field{sortinithash}{d7095fff47cda75ca2589920aae98399}
      \field{labelnamesource}{author}
      \field{labeltitlesource}{title}
      \field{eprinttype}{arXiv}
      \field{title}{Sound Statistical Model Checking for Probabilities and Expected Rewards (extended version)}
      \field{year}{2024}
      \verb{eprint}
      \verb 2411.00559
      \endverb
    \endentry
    \entry{budiu22}{online}{}{}
      \name{author}{4}{}{%
        {{hash=8836a39dd4362973525947badc96d6d0}{%
           family={Budiu},
           familyi={B\bibinitperiod},
           given={Mihai},
           giveni={M\bibinitperiod}}}%
        {{hash=a0e6c99aabdc2587f248892c8bb23f70}{%
           family={McSherry},
           familyi={M\bibinitperiod},
           given={Frank},
           giveni={F\bibinitperiod}}}%
        {{hash=b6e0a6b2bf14393102644621a18bc62b}{%
           family={Ryzhyk},
           familyi={R\bibinitperiod},
           given={Leonid},
           giveni={L\bibinitperiod}}}%
        {{hash=7fb00370e14cba6a31e3b49546335e1b}{%
           family={Tannen},
           familyi={T\bibinitperiod},
           given={Val},
           giveni={V\bibinitperiod}}}%
      }
      \strng{namehash}{50f3545abbe22df6cff9e4cd09a83050}
      \strng{fullhash}{50f3545abbe22df6cff9e4cd09a83050}
      \strng{fullhashraw}{50f3545abbe22df6cff9e4cd09a83050}
      \strng{bibnamehash}{50f3545abbe22df6cff9e4cd09a83050}
      \strng{authorbibnamehash}{50f3545abbe22df6cff9e4cd09a83050}
      \strng{authornamehash}{50f3545abbe22df6cff9e4cd09a83050}
      \strng{authorfullhash}{50f3545abbe22df6cff9e4cd09a83050}
      \strng{authorfullhashraw}{50f3545abbe22df6cff9e4cd09a83050}
      \field{sortinit}{B}
      \field{sortinithash}{d7095fff47cda75ca2589920aae98399}
      \field{labelnamesource}{author}
      \field{labeltitlesource}{title}
      \field{eprinttype}{arXiv}
      \field{subtitle}{Automatic Incremental View Maintenance for Rich Query Languages}
      \field{title}{DBSP}
      \field{year}{2022}
      \verb{eprint}
      \verb 2203.16684
      \endverb
    \endentry
    \entry{bujorianu08}{online}{}{}
      \name{author}{3}{}{%
        {{hash=444cee2a8a2a705a2980959da381da48}{%
           family={Bujorianu},
           familyi={B\bibinitperiod},
           given={Manuela\bibnamedelima L.},
           giveni={M\bibinitperiod\bibinitdelim L\bibinitperiod}}}%
        {{hash=98e6d333e012478ed48daafff6182700}{%
           family={Lygeros},
           familyi={L\bibinitperiod},
           given={John},
           giveni={J\bibinitperiod}}}%
        {{hash=703c7e64d5ddcb99e4005678b0a1e56b}{%
           family={Bujorianu},
           familyi={B\bibinitperiod},
           given={Marius\bibnamedelima C.},
           giveni={M\bibinitperiod\bibinitdelim C\bibinitperiod}}}%
      }
      \strng{namehash}{a333f2c251665155ffb702116abbe9d5}
      \strng{fullhash}{a333f2c251665155ffb702116abbe9d5}
      \strng{fullhashraw}{a333f2c251665155ffb702116abbe9d5}
      \strng{bibnamehash}{a333f2c251665155ffb702116abbe9d5}
      \strng{authorbibnamehash}{a333f2c251665155ffb702116abbe9d5}
      \strng{authornamehash}{a333f2c251665155ffb702116abbe9d5}
      \strng{authorfullhash}{a333f2c251665155ffb702116abbe9d5}
      \strng{authorfullhashraw}{a333f2c251665155ffb702116abbe9d5}
      \field{sortinit}{B}
      \field{sortinithash}{d7095fff47cda75ca2589920aae98399}
      \field{labelnamesource}{author}
      \field{labeltitlesource}{title}
      \field{day}{12}
      \field{month}{3}
      \field{note}{CTIT Technical Report Series, No. TR-CTIT-08-25, Centre for Telematics and Information Technology, University of Twente, Enschede}
      \field{title}{Towards a General Theory of Stochastic Hybrid Systems}
      \field{year}{2008}
      \field{dateera}{ce}
      \verb{urlraw}
      \verb https://ris.utwente.nl/ws/portalfiles/portal/5100731/journal8.pdf
      \endverb
      \verb{url}
      \verb https://ris.utwente.nl/ws/portalfiles/portal/5100731/journal8.pdf
      \endverb
    \endentry
    \entry{bunel2020branch}{article}{}{}
      \name{author}{6}{}{%
        {{hash=a5676fe046216ea8175331cf51b71009}{%
           family={Bunel},
           familyi={B\bibinitperiod},
           given={R.},
           giveni={R\bibinitperiod}}}%
        {{hash=642f8aecd69ae5ee926bceabfd45a56d}{%
           family={Lu},
           familyi={L\bibinitperiod},
           given={J.},
           giveni={J\bibinitperiod}}}%
        {{hash=c011ad6406fa57adf7fa41dae999ea96}{%
           family={Turkaslan},
           familyi={T\bibinitperiod},
           given={I.},
           giveni={I\bibinitperiod}}}%
        {{hash=71c055becbb66c5c1442770da155ab38}{%
           family={Torr},
           familyi={T\bibinitperiod},
           given={P.\bibnamedelimi H.\bibnamedelimi S.},
           giveni={P\bibinitperiod\bibinitdelim H\bibinitperiod\bibinitdelim S\bibinitperiod}}}%
        {{hash=cb0207368e596644cbaa518983912e70}{%
           family={Kohli},
           familyi={K\bibinitperiod},
           given={P.},
           giveni={P\bibinitperiod}}}%
        {{hash=c31501e9fb0a60cae3f31e0b2e72a187}{%
           family={Kumar},
           familyi={K\bibinitperiod},
           given={M.\bibnamedelimi P.},
           giveni={M\bibinitperiod\bibinitdelim P\bibinitperiod}}}%
      }
      \strng{namehash}{0953aa0959a174c0147ab22698785f91}
      \strng{fullhash}{0953aa0959a174c0147ab22698785f91}
      \strng{fullhashraw}{0953aa0959a174c0147ab22698785f91}
      \strng{bibnamehash}{0953aa0959a174c0147ab22698785f91}
      \strng{authorbibnamehash}{0953aa0959a174c0147ab22698785f91}
      \strng{authornamehash}{0953aa0959a174c0147ab22698785f91}
      \strng{authorfullhash}{0953aa0959a174c0147ab22698785f91}
      \strng{authorfullhashraw}{0953aa0959a174c0147ab22698785f91}
      \field{sortinit}{B}
      \field{sortinithash}{d7095fff47cda75ca2589920aae98399}
      \field{labelnamesource}{author}
      \field{labeltitlesource}{title}
      \field{journaltitle}{Journal of Machine Learning Research}
      \field{title}{Branch and bound for piecewise linear neural network verification}
      \field{volume}{21}
      \field{year}{2020}
    \endentry
    \entry{burgess07}{online}{}{}
      \name{author}{1}{}{%
        {{hash=f2b12a7bbb4155b6d4de7d8bc3dd4008}{%
           family={Burgess},
           familyi={B\bibinitperiod},
           given={C.\bibnamedelimi P.},
           giveni={C\bibinitperiod\bibinitdelim P\bibinitperiod}}}%
      }
      \strng{namehash}{f2b12a7bbb4155b6d4de7d8bc3dd4008}
      \strng{fullhash}{f2b12a7bbb4155b6d4de7d8bc3dd4008}
      \strng{fullhashraw}{f2b12a7bbb4155b6d4de7d8bc3dd4008}
      \strng{bibnamehash}{f2b12a7bbb4155b6d4de7d8bc3dd4008}
      \strng{authorbibnamehash}{f2b12a7bbb4155b6d4de7d8bc3dd4008}
      \strng{authornamehash}{f2b12a7bbb4155b6d4de7d8bc3dd4008}
      \strng{authorfullhash}{f2b12a7bbb4155b6d4de7d8bc3dd4008}
      \strng{authorfullhashraw}{f2b12a7bbb4155b6d4de7d8bc3dd4008}
      \field{sortinit}{B}
      \field{sortinithash}{d7095fff47cda75ca2589920aae98399}
      \field{labelnamesource}{author}
      \field{labeltitlesource}{title}
      \field{eprinttype}{arXiv}
      \field{title}{Introduction to Effective Field Theory}
      \field{year}{2007}
      \verb{doi}
      \verb 10.1146/annurev.nucl.56.080805.140508
      \endverb
      \verb{eprint}
      \verb hep-th/0701053
      \endverb
    \endentry
    \entry{cao25}{online}{}{}
      \name{author}{6}{}{%
        {{hash=e6b5be6e497ae1e488fb905c711b32c8}{%
           family={Cao},
           familyi={C\bibinitperiod},
           given={Zhipeng},
           giveni={Z\bibinitperiod}}}%
        {{hash=fb306031b264cacf48c55f996c616dd3}{%
           family={Wang},
           familyi={W\bibinitperiod},
           given={Peixin},
           giveni={P\bibinitperiod}}}%
        {{hash=0e126ff846e9a25b78ab52904848937d}{%
           family={Ong},
           familyi={O\bibinitperiod},
           given={Luke},
           giveni={L\bibinitperiod}}}%
        {{hash=4faa29f2520e357e01f0b9e7fc70e46f}{%
           family={Žikelić},
           familyi={Ž\bibinitperiod},
           given={Đorđe},
           giveni={Đ\bibinitperiod}}}%
        {{hash=9f6e7512afbfa89a75e989fdfef96419}{%
           family={Wagner},
           familyi={W\bibinitperiod},
           given={Dominik},
           giveni={D\bibinitperiod}}}%
        {{hash=02b4a2c89e9f6ef8cefa8754514753a0}{%
           family={Xue},
           familyi={X\bibinitperiod},
           given={Bai},
           giveni={B\bibinitperiod}}}%
      }
      \strng{namehash}{7f550b8ea55359f2e8c0530f8d14a840}
      \strng{fullhash}{7f550b8ea55359f2e8c0530f8d14a840}
      \strng{fullhashraw}{7f550b8ea55359f2e8c0530f8d14a840}
      \strng{bibnamehash}{7f550b8ea55359f2e8c0530f8d14a840}
      \strng{authorbibnamehash}{7f550b8ea55359f2e8c0530f8d14a840}
      \strng{authornamehash}{7f550b8ea55359f2e8c0530f8d14a840}
      \strng{authorfullhash}{7f550b8ea55359f2e8c0530f8d14a840}
      \strng{authorfullhashraw}{7f550b8ea55359f2e8c0530f8d14a840}
      \field{sortinit}{C}
      \field{sortinithash}{4d103a86280481745c9c897c925753c0}
      \field{labelnamesource}{author}
      \field{labeltitlesource}{title}
      \field{eprinttype}{arXiv}
      \field{title}{Comparative Analysis of Barrier-like Function Methods for Reach-Avoid Verification in Stochastic Discrete-Time Systems}
      \field{year}{2025}
      \verb{eprint}
      \verb 2512.05348
      \endverb
    \endentry
    \entry{caprio26}{online}{}{}
      \name{author}{3}{}{%
        {{hash=e15f4a05e187c0d01bd03368e162ac12}{%
           family={Caprio},
           familyi={C\bibinitperiod},
           given={Rocco},
           giveni={R\bibinitperiod}}}%
        {{hash=500462c41dd886082f1f8861030c76c6}{%
           family={Corenflos},
           familyi={C\bibinitperiod},
           given={Adrien},
           giveni={A\bibinitperiod}}}%
        {{hash=19f7fc67cda0cf1ebc08330626a575e3}{%
           family={Power},
           familyi={P\bibinitperiod},
           given={Sam},
           giveni={S\bibinitperiod}}}%
      }
      \strng{namehash}{3165b428e8f1a56a975da374cb3c8467}
      \strng{fullhash}{3165b428e8f1a56a975da374cb3c8467}
      \strng{fullhashraw}{3165b428e8f1a56a975da374cb3c8467}
      \strng{bibnamehash}{3165b428e8f1a56a975da374cb3c8467}
      \strng{authorbibnamehash}{3165b428e8f1a56a975da374cb3c8467}
      \strng{authornamehash}{3165b428e8f1a56a975da374cb3c8467}
      \strng{authorfullhash}{3165b428e8f1a56a975da374cb3c8467}
      \strng{authorfullhashraw}{3165b428e8f1a56a975da374cb3c8467}
      \field{sortinit}{C}
      \field{sortinithash}{4d103a86280481745c9c897c925753c0}
      \field{labelnamesource}{author}
      \field{labeltitlesource}{title}
      \field{eprinttype}{arXiv}
      \field{title}{Wasserstein Contraction of Coordinate Ascent Variational Inference}
      \field{year}{2026}
      \verb{eprint}
      \verb 2605.30253
      \endverb
    \endentry
    \entry{capuccimyers2026}{misc}{}{}
      \name{author}{2}{}{%
        {{hash=fd57ee30e815684db8bc133a3970657e}{%
           family={Capucci},
           familyi={C\bibinitperiod},
           given={M.},
           giveni={M\bibinitperiod}}}%
        {{hash=21e13fe1ad729441a88bad3d8ce9c089}{%
           family={Myers},
           familyi={M\bibinitperiod},
           given={D.\bibnamedelimi J.},
           giveni={D\bibinitperiod\bibinitdelim J\bibinitperiod}}}%
      }
      \strng{namehash}{cee079a4752b36c2a3c9c3d3bfceeb36}
      \strng{fullhash}{cee079a4752b36c2a3c9c3d3bfceeb36}
      \strng{fullhashraw}{cee079a4752b36c2a3c9c3d3bfceeb36}
      \strng{bibnamehash}{cee079a4752b36c2a3c9c3d3bfceeb36}
      \strng{authorbibnamehash}{cee079a4752b36c2a3c9c3d3bfceeb36}
      \strng{authornamehash}{cee079a4752b36c2a3c9c3d3bfceeb36}
      \strng{authorfullhash}{cee079a4752b36c2a3c9c3d3bfceeb36}
      \strng{authorfullhashraw}{cee079a4752b36c2a3c9c3d3bfceeb36}
      \field{extraname}{1}
      \field{sortinit}{C}
      \field{sortinithash}{4d103a86280481745c9c897c925753c0}
      \field{labelnamesource}{author}
      \field{labeltitlesource}{title}
      \field{eprinttype}{arxiv}
      \field{note}{(The wiring-square conditions, the $\mathcal{K}^0_\infty$ slack discipline, and the compose-by-lens-composition-and-slack-addition law cited in the composition sections.)}
      \field{title}{Certified lenses and compositional system certification}
      \field{year}{2026}
      \verb{eprint}
      \verb 2604.03017
      \endverb
    \endentry
    \entry{capucci26}{online}{}{}
      \name{author}{2}{}{%
        {{hash=2f174406fd3061220c36204e0b0c69a7}{%
           family={Capucci},
           familyi={C\bibinitperiod},
           given={Matteo},
           giveni={M\bibinitperiod}}}%
        {{hash=5e4483260e70411ee8f2ec4093d57437}{%
           family={Myers},
           familyi={M\bibinitperiod},
           given={David\bibnamedelima Jaz},
           giveni={D\bibinitperiod\bibinitdelim J\bibinitperiod}}}%
      }
      \strng{namehash}{a324fb5df16424ae439531b832f7a880}
      \strng{fullhash}{a324fb5df16424ae439531b832f7a880}
      \strng{fullhashraw}{a324fb5df16424ae439531b832f7a880}
      \strng{bibnamehash}{a324fb5df16424ae439531b832f7a880}
      \strng{authorbibnamehash}{a324fb5df16424ae439531b832f7a880}
      \strng{authornamehash}{a324fb5df16424ae439531b832f7a880}
      \strng{authorfullhash}{a324fb5df16424ae439531b832f7a880}
      \strng{authorfullhashraw}{a324fb5df16424ae439531b832f7a880}
      \field{extraname}{2}
      \field{sortinit}{C}
      \field{sortinithash}{4d103a86280481745c9c897c925753c0}
      \field{labelnamesource}{author}
      \field{labeltitlesource}{title}
      \field{eprinttype}{arXiv}
      \field{title}{Compositionality of Lyapunov functions via assume-guarantee reasoning}
      \field{year}{2026}
      \verb{eprint}
      \verb 2604.03017
      \endverb
    \endentry
    \entry{cargoshisha1966}{article}{}{}
      \name{author}{2}{}{%
        {{hash=457703c4b289818f4a9f4df92332398b}{%
           family={Cargo},
           familyi={C\bibinitperiod},
           given={G.\bibnamedelimi T.},
           giveni={G\bibinitperiod\bibinitdelim T\bibinitperiod}}}%
        {{hash=db2060c1692dee782abd634706d37e7f}{%
           family={Shisha},
           familyi={S\bibinitperiod},
           given={O.},
           giveni={O\bibinitperiod}}}%
      }
      \strng{namehash}{6516eed053ae6f75bb32132dc3bce076}
      \strng{fullhash}{6516eed053ae6f75bb32132dc3bce076}
      \strng{fullhashraw}{6516eed053ae6f75bb32132dc3bce076}
      \strng{bibnamehash}{6516eed053ae6f75bb32132dc3bce076}
      \strng{authorbibnamehash}{6516eed053ae6f75bb32132dc3bce076}
      \strng{authornamehash}{6516eed053ae6f75bb32132dc3bce076}
      \strng{authorfullhash}{6516eed053ae6f75bb32132dc3bce076}
      \strng{authorfullhashraw}{6516eed053ae6f75bb32132dc3bce076}
      \field{sortinit}{C}
      \field{sortinithash}{4d103a86280481745c9c897c925753c0}
      \field{labelnamesource}{author}
      \field{labeltitlesource}{title}
      \field{journaltitle}{Journal of Research of the National Bureau of Standards}
      \field{note}{(Bernstein coefficients enclose the range on the simplex --- the sound direction this calculus's kernel uses.)}
      \field{title}{The Bernstein form of a polynomial}
      \field{volume}{70B}
      \field{year}{1966}
      \field{pages}{79\bibrangedash 81}
      \range{pages}{3}
    \endentry
    \entry{decasteljau1959}{misc}{}{}
      \name{author}{1}{}{%
        {{hash=7adb28692ea5a550784389777a8926ae}{%
           family={Casteljau},
           familyi={C\bibinitperiod},
           given={P.},
           giveni={P\bibinitperiod},
           prefix={de},
           prefixi={d\bibinitperiod}}}%
      }
      \strng{namehash}{7adb28692ea5a550784389777a8926ae}
      \strng{fullhash}{7adb28692ea5a550784389777a8926ae}
      \strng{fullhashraw}{7adb28692ea5a550784389777a8926ae}
      \strng{bibnamehash}{7adb28692ea5a550784389777a8926ae}
      \strng{authorbibnamehash}{7adb28692ea5a550784389777a8926ae}
      \strng{authornamehash}{7adb28692ea5a550784389777a8926ae}
      \strng{authorfullhash}{7adb28692ea5a550784389777a8926ae}
      \strng{authorfullhashraw}{7adb28692ea5a550784389777a8926ae}
      \field{sortinit}{C}
      \field{sortinithash}{4d103a86280481745c9c897c925753c0}
      \field{labelnamesource}{author}
      \field{labeltitlesource}{title}
      \field{howpublished}{Technical report, A.~Citroën, Paris}
      \field{note}{(The founding subdivision/evaluation algorithm for polynomials in Bernstein form.)}
      \field{title}{Outillages méthodes calcul}
      \field{year}{1959}
    \endentry
    \entry{chakarov2013}{inproceedings}{}{}
      \name{author}{2}{}{%
        {{hash=d1705fc2d5aa0720013b4d992ada76ed}{%
           family={Chakarov},
           familyi={C\bibinitperiod},
           given={A.},
           giveni={A\bibinitperiod}}}%
        {{hash=ec31ae80c65f63ba67850106aeb1a997}{%
           family={Sankaranarayanan},
           familyi={S\bibinitperiod},
           given={S.},
           giveni={S\bibinitperiod}}}%
      }
      \list{publisher}{1}{%
        {Springer}%
      }
      \strng{namehash}{9fe434d4bff6cf36d41b24ee64d96002}
      \strng{fullhash}{9fe434d4bff6cf36d41b24ee64d96002}
      \strng{fullhashraw}{9fe434d4bff6cf36d41b24ee64d96002}
      \strng{bibnamehash}{9fe434d4bff6cf36d41b24ee64d96002}
      \strng{authorbibnamehash}{9fe434d4bff6cf36d41b24ee64d96002}
      \strng{authornamehash}{9fe434d4bff6cf36d41b24ee64d96002}
      \strng{authorfullhash}{9fe434d4bff6cf36d41b24ee64d96002}
      \strng{authorfullhashraw}{9fe434d4bff6cf36d41b24ee64d96002}
      \field{sortinit}{C}
      \field{sortinithash}{4d103a86280481745c9c897c925753c0}
      \field{labelnamesource}{author}
      \field{labeltitlesource}{title}
      \field{booktitle}{Computer Aided Verification (CAV 2013)}
      \field{note}{(Ranking supermartingales.)}
      \field{series}{Lecture Notes in Computer Science}
      \field{title}{Probabilistic program analysis with martingales}
      \field{volume}{8044}
      \field{year}{2013}
      \field{pages}{511\bibrangedash 526}
      \range{pages}{16}
    \endentry
    \entry{ensl-00472509v2}{online}{}{}
      \field{sortinit}{C}
      \field{sortinithash}{4d103a86280481745c9c897c925753c0}
      \field{labeltitlesource}{title}
      \field{title}{Chebyshev Interpolation Polynomial-based Tools for Rigorous Computing}
      \verb{urlraw}
      \verb https://ens-lyon.hal.science/ensl-00472509v2
      \endverb
      \verb{url}
      \verb https://ens-lyon.hal.science/ensl-00472509v2
      \endverb
    \endentry
    \entry{chen2024aldoushoover}{misc}{}{}
      \name{author}{5}{}{%
        {{hash=4813a2079647f41d0ed98a0181d267d2}{%
           family={Chen},
           familyi={C\bibinitperiod},
           given={Leihao},
           giveni={L\bibinitperiod}}}%
        {{hash=485f579f829f8629fec8c85c28a33f28}{%
           family={Fritz},
           familyi={F\bibinitperiod},
           given={Tobias},
           giveni={T\bibinitperiod}}}%
        {{hash=f45bf5d5b6fdf712f1ff8883b1cd5528}{%
           family={Gonda},
           familyi={G\bibinitperiod},
           given={Tomáš},
           giveni={T\bibinitperiod}}}%
        {{hash=f5ecc909057ac0224d0fd82d961ae80f}{%
           family={Klingler},
           familyi={K\bibinitperiod},
           given={Andreas},
           giveni={A\bibinitperiod}}}%
        {{hash=fdd77c5f67dd4d01463a6f4381c4e883}{%
           family={Lorenzin},
           familyi={L\bibinitperiod},
           given={Antonio},
           giveni={A\bibinitperiod}}}%
      }
      \strng{namehash}{e86690420ca00b20dfaed50195a0b97b}
      \strng{fullhash}{e86690420ca00b20dfaed50195a0b97b}
      \strng{fullhashraw}{e86690420ca00b20dfaed50195a0b97b}
      \strng{bibnamehash}{e86690420ca00b20dfaed50195a0b97b}
      \strng{authorbibnamehash}{e86690420ca00b20dfaed50195a0b97b}
      \strng{authornamehash}{e86690420ca00b20dfaed50195a0b97b}
      \strng{authorfullhash}{e86690420ca00b20dfaed50195a0b97b}
      \strng{authorfullhashraw}{e86690420ca00b20dfaed50195a0b97b}
      \field{sortinit}{C}
      \field{sortinithash}{4d103a86280481745c9c897c925753c0}
      \field{labelnamesource}{author}
      \field{labeltitlesource}{title}
      \field{eprinttype}{arxiv}
      \field{note}{(Plates over Kolmogorov products: the unique morphism whose finite marginals are the body's instances.)}
      \field{title}{The {Aldous--Hoover} Theorem in Categorical Probability}
      \verb{eprint}
      \verb 2411.12840
      \endverb
    \endentry
    \entry{cleaveland18hbmc}{incollection}{}{}
      \name{author}{3}{}{%
        {{hash=c7c1c51c77bfaf498c9599e87c626b34}{%
           family={Cleaveland},
           familyi={C\bibinitperiod},
           given={Rance},
           giveni={R\bibinitperiod}}}%
        {{hash=d0ae119316415970e7817a8db830f817}{%
           family={Roscoe},
           familyi={R\bibinitperiod},
           given={A.\bibnamedelimi W.},
           giveni={A\bibinitperiod\bibinitdelim W\bibinitperiod}}}%
        {{hash=acb872a5ef65da192eb1b1e9158d646c}{%
           family={Smolka},
           familyi={S\bibinitperiod},
           given={Scott\bibnamedelima A.},
           giveni={S\bibinitperiod\bibinitdelim A\bibinitperiod}}}%
      }
      \name{editor}{4}{}{%
        {{hash=1822cd554880ea55fd66357a96bca03e}{%
           family={Clarke},
           familyi={C\bibinitperiod},
           given={Edmund\bibnamedelima M.},
           giveni={E\bibinitperiod\bibinitdelim M\bibinitperiod}}}%
        {{hash=eeb928073fbee58751f87d401a62561a}{%
           family={Henzinger},
           familyi={H\bibinitperiod},
           given={Thomas\bibnamedelima A.},
           giveni={T\bibinitperiod\bibinitdelim A\bibinitperiod}}}%
        {{hash=5ac94523768e8fc4adb7a926b36d83e0}{%
           family={Veith},
           familyi={V\bibinitperiod},
           given={Helmut},
           giveni={H\bibinitperiod}}}%
        {{hash=9865af1b742410e6dee8ca813e1ec733}{%
           family={Bloem},
           familyi={B\bibinitperiod},
           given={Roderick},
           giveni={R\bibinitperiod}}}%
      }
      \list{publisher}{1}{%
        {Springer}%
      }
      \strng{namehash}{6ac5e628eb0eb7e803051009e6db3bb4}
      \strng{fullhash}{6ac5e628eb0eb7e803051009e6db3bb4}
      \strng{fullhashraw}{6ac5e628eb0eb7e803051009e6db3bb4}
      \strng{bibnamehash}{6ac5e628eb0eb7e803051009e6db3bb4}
      \strng{authorbibnamehash}{6ac5e628eb0eb7e803051009e6db3bb4}
      \strng{authornamehash}{6ac5e628eb0eb7e803051009e6db3bb4}
      \strng{authorfullhash}{6ac5e628eb0eb7e803051009e6db3bb4}
      \strng{authorfullhashraw}{6ac5e628eb0eb7e803051009e6db3bb4}
      \strng{editorbibnamehash}{eea0cb89759dac86c5d80411829d95b2}
      \strng{editornamehash}{eea0cb89759dac86c5d80411829d95b2}
      \strng{editorfullhash}{eea0cb89759dac86c5d80411829d95b2}
      \strng{editorfullhashraw}{eea0cb89759dac86c5d80411829d95b2}
      \field{sortinit}{C}
      \field{sortinithash}{4d103a86280481745c9c897c925753c0}
      \field{labelnamesource}{author}
      \field{labeltitlesource}{title}
      \field{booktitle}{Handbook of Model Checking}
      \field{title}{Process Algebra and Model Checking}
      \field{year}{2018}
      \field{pages}{1149\bibrangedash 1195}
      \range{pages}{47}
      \verb{doi}
      \verb 10.1007/978-3-319-10575-8_32
      \endverb
    \endentry
    \entry{codd86}{article}{}{}
      \name{author}{1}{}{%
        {{hash=575e6ec8dab1a34c19f7021664baaccc}{%
           family={Codd},
           familyi={C\bibinitperiod},
           given={E.\bibnamedelimi F.},
           giveni={E\bibinitperiod\bibinitdelim F\bibinitperiod}}}%
      }
      \strng{namehash}{575e6ec8dab1a34c19f7021664baaccc}
      \strng{fullhash}{575e6ec8dab1a34c19f7021664baaccc}
      \strng{fullhashraw}{575e6ec8dab1a34c19f7021664baaccc}
      \strng{bibnamehash}{575e6ec8dab1a34c19f7021664baaccc}
      \strng{authorbibnamehash}{575e6ec8dab1a34c19f7021664baaccc}
      \strng{authornamehash}{575e6ec8dab1a34c19f7021664baaccc}
      \strng{authorfullhash}{575e6ec8dab1a34c19f7021664baaccc}
      \strng{authorfullhashraw}{575e6ec8dab1a34c19f7021664baaccc}
      \field{sortinit}{C}
      \field{sortinithash}{4d103a86280481745c9c897c925753c0}
      \field{labelnamesource}{author}
      \field{labeltitlesource}{title}
      \field{journaltitle}{ACM SIGMOD Record}
      \field{number}{4}
      \field{title}{Missing information (applicable and inapplicable) in relational databases}
      \field{volume}{15}
      \field{year}{1986}
      \verb{doi}
      \verb 10.1145/16301.16303
      \endverb
    \endentry
    \entry{coddington-levinson-1955}{book}{}{}
      \name{author}{2}{}{%
        {{hash=e8867c884d75d990c4e4407b71254420}{%
           family={Coddington},
           familyi={C\bibinitperiod},
           given={Earl\bibnamedelima A.},
           giveni={E\bibinitperiod\bibinitdelim A\bibinitperiod}}}%
        {{hash=383c8383f13be54ab69e28076ca08a38}{%
           family={Levinson},
           familyi={L\bibinitperiod},
           given={Norman},
           giveni={N\bibinitperiod}}}%
      }
      \list{location}{1}{%
        {New York}%
      }
      \list{publisher}{1}{%
        {McGraw-Hill}%
      }
      \strng{namehash}{449fb8d92bcd37c345212688875d5129}
      \strng{fullhash}{449fb8d92bcd37c345212688875d5129}
      \strng{fullhashraw}{449fb8d92bcd37c345212688875d5129}
      \strng{bibnamehash}{449fb8d92bcd37c345212688875d5129}
      \strng{authorbibnamehash}{449fb8d92bcd37c345212688875d5129}
      \strng{authornamehash}{449fb8d92bcd37c345212688875d5129}
      \strng{authorfullhash}{449fb8d92bcd37c345212688875d5129}
      \strng{authorfullhashraw}{449fb8d92bcd37c345212688875d5129}
      \field{sortinit}{C}
      \field{sortinithash}{4d103a86280481745c9c897c925753c0}
      \field{labelnamesource}{author}
      \field{labeltitlesource}{title}
      \field{title}{Theory of Ordinary Differential Equations}
      \field{year}{1955}
    \endentry
    \entry{conley1978}{book}{}{}
      \name{author}{1}{}{%
        {{hash=b13efc9c5acb54c1ca6293a6919763b4}{%
           family={Conley},
           familyi={C\bibinitperiod},
           given={C.},
           giveni={C\bibinitperiod}}}%
      }
      \list{publisher}{1}{%
        {American Mathematical Society}%
      }
      \strng{namehash}{b13efc9c5acb54c1ca6293a6919763b4}
      \strng{fullhash}{b13efc9c5acb54c1ca6293a6919763b4}
      \strng{fullhashraw}{b13efc9c5acb54c1ca6293a6919763b4}
      \strng{bibnamehash}{b13efc9c5acb54c1ca6293a6919763b4}
      \strng{authorbibnamehash}{b13efc9c5acb54c1ca6293a6919763b4}
      \strng{authornamehash}{b13efc9c5acb54c1ca6293a6919763b4}
      \strng{authorfullhash}{b13efc9c5acb54c1ca6293a6919763b4}
      \strng{authorfullhashraw}{b13efc9c5acb54c1ca6293a6919763b4}
      \field{sortinit}{C}
      \field{sortinithash}{4d103a86280481745c9c897c925753c0}
      \field{labelnamesource}{author}
      \field{labeltitlesource}{title}
      \field{number}{38}
      \field{series}{CBMS Regional Conference Series in Mathematics}
      \field{title}{Isolated Invariant Sets and the Morse Index}
      \field{year}{1978}
    \endentry
    \entry{conleyeaston1971}{article}{}{}
      \name{author}{2}{}{%
        {{hash=b13efc9c5acb54c1ca6293a6919763b4}{%
           family={Conley},
           familyi={C\bibinitperiod},
           given={C.},
           giveni={C\bibinitperiod}}}%
        {{hash=f545e1a525c6e5fdc76426c1982cbd63}{%
           family={Easton},
           familyi={E\bibinitperiod},
           given={R.},
           giveni={R\bibinitperiod}}}%
      }
      \strng{namehash}{3c2f05e72af68fae5392a67b225ca846}
      \strng{fullhash}{3c2f05e72af68fae5392a67b225ca846}
      \strng{fullhashraw}{3c2f05e72af68fae5392a67b225ca846}
      \strng{bibnamehash}{3c2f05e72af68fae5392a67b225ca846}
      \strng{authorbibnamehash}{3c2f05e72af68fae5392a67b225ca846}
      \strng{authornamehash}{3c2f05e72af68fae5392a67b225ca846}
      \strng{authorfullhash}{3c2f05e72af68fae5392a67b225ca846}
      \strng{authorfullhashraw}{3c2f05e72af68fae5392a67b225ca846}
      \field{sortinit}{C}
      \field{sortinithash}{4d103a86280481745c9c897c925753c0}
      \field{labelnamesource}{author}
      \field{labeltitlesource}{title}
      \field{journaltitle}{Transactions of the American Mathematical Society}
      \field{title}{Isolated invariant sets and isolating blocks}
      \field{volume}{158}
      \field{year}{1971}
      \field{pages}{35\bibrangedash 61}
      \range{pages}{27}
    \endentry
    \entry{cooman08}{online}{}{}
      \name{author}{2}{}{%
        {{hash=f57f0c0575e93f2d0040699201ff7eee}{%
           family={Cooman},
           familyi={C\bibinitperiod},
           given={Gert},
           giveni={G\bibinitperiod},
           prefix={de},
           prefixi={d\bibinitperiod}}}%
        {{hash=0d18941dfbf4c5188b4afa3f8d2a70b3}{%
           family={Hermans},
           familyi={H\bibinitperiod},
           given={Filip},
           giveni={F\bibinitperiod}}}%
      }
      \strng{namehash}{ebf944531de81b68175c673a36fe13fc}
      \strng{fullhash}{ebf944531de81b68175c673a36fe13fc}
      \strng{fullhashraw}{ebf944531de81b68175c673a36fe13fc}
      \strng{bibnamehash}{ebf944531de81b68175c673a36fe13fc}
      \strng{authorbibnamehash}{ebf944531de81b68175c673a36fe13fc}
      \strng{authornamehash}{ebf944531de81b68175c673a36fe13fc}
      \strng{authorfullhash}{ebf944531de81b68175c673a36fe13fc}
      \strng{authorfullhashraw}{ebf944531de81b68175c673a36fe13fc}
      \field{sortinit}{C}
      \field{sortinithash}{4d103a86280481745c9c897c925753c0}
      \field{labelnamesource}{author}
      \field{labeltitlesource}{title}
      \field{eprinttype}{arXiv}
      \field{subtitle}{Bridging two theories of imprecise probability}
      \field{title}{Imprecise probability trees}
      \field{year}{2008}
      \verb{eprint}
      \verb 0801.1196
      \endverb
    \endentry
    \entry{cozman00}{article}{}{}
      \name{author}{1}{}{%
        {{hash=1c11a5fd025ccc4382411523f7697ba1}{%
           family={Cozman},
           familyi={C\bibinitperiod},
           given={Fábio\bibnamedelima Gagliardi},
           giveni={F\bibinitperiod\bibinitdelim G\bibinitperiod}}}%
      }
      \strng{namehash}{1c11a5fd025ccc4382411523f7697ba1}
      \strng{fullhash}{1c11a5fd025ccc4382411523f7697ba1}
      \strng{fullhashraw}{1c11a5fd025ccc4382411523f7697ba1}
      \strng{bibnamehash}{1c11a5fd025ccc4382411523f7697ba1}
      \strng{authorbibnamehash}{1c11a5fd025ccc4382411523f7697ba1}
      \strng{authornamehash}{1c11a5fd025ccc4382411523f7697ba1}
      \strng{authorfullhash}{1c11a5fd025ccc4382411523f7697ba1}
      \strng{authorfullhashraw}{1c11a5fd025ccc4382411523f7697ba1}
      \field{sortinit}{C}
      \field{sortinithash}{4d103a86280481745c9c897c925753c0}
      \field{labelnamesource}{author}
      \field{labeltitlesource}{title}
      \field{journaltitle}{Artificial Intelligence}
      \field{number}{2}
      \field{title}{Credal Networks}
      \field{volume}{120}
      \field{year}{2000}
      \field{pages}{199\bibrangedash 233}
      \range{pages}{35}
      \verb{doi}
      \verb 10.1016/S0004-3702(00)00029-1
      \endverb
    \endentry
    \entry{crooks22}{online}{}{}
      \name{author}{1}{}{%
        {{hash=9f65ac557691562d4b8946b4ba50cd96}{%
           family={Crooks},
           familyi={C\bibinitperiod},
           given={Gavin\bibnamedelima E.},
           giveni={G\bibinitperiod\bibinitdelim E\bibinitperiod}}}%
      }
      \strng{namehash}{9f65ac557691562d4b8946b4ba50cd96}
      \strng{fullhash}{9f65ac557691562d4b8946b4ba50cd96}
      \strng{fullhashraw}{9f65ac557691562d4b8946b4ba50cd96}
      \strng{bibnamehash}{9f65ac557691562d4b8946b4ba50cd96}
      \strng{authorbibnamehash}{9f65ac557691562d4b8946b4ba50cd96}
      \strng{authornamehash}{9f65ac557691562d4b8946b4ba50cd96}
      \strng{authorfullhash}{9f65ac557691562d4b8946b4ba50cd96}
      \strng{authorfullhashraw}{9f65ac557691562d4b8946b4ba50cd96}
      \field{sortinit}{C}
      \field{sortinithash}{4d103a86280481745c9c897c925753c0}
      \field{labelnamesource}{author}
      \field{labeltitlesource}{title}
      \field{day}{15}
      \field{month}{2}
      \field{title}{Field Guide to Continuous Probability Distributions}
      \field{year}{2022}
      \field{dateera}{ce}
      \verb{urlraw}
      \verb https://threeplusone.com/pubs/FieldGuide.pdf
      \endverb
      \verb{url}
      \verb https://threeplusone.com/pubs/FieldGuide.pdf
      \endverb
    \endentry
    \entry{cruttwell22b}{online}{}{}
      \name{author}{4}{}{%
        {{hash=3ea0f7c45df4643e66dae205b5dfff1d}{%
           family={Cruttwell},
           familyi={C\bibinitperiod},
           given={Geoffrey},
           giveni={G\bibinitperiod}}}%
        {{hash=26775461fd456edc4deb5b983123ea7c}{%
           family={Lambert},
           familyi={L\bibinitperiod},
           given={Michael},
           giveni={M\bibinitperiod}}}%
        {{hash=99961f77ef3feca99280610091ef52de}{%
           family={Pronk},
           familyi={P\bibinitperiod},
           given={Dorette},
           giveni={D\bibinitperiod}}}%
        {{hash=3ead16e47970cc7e12cd2a33df6d41b2}{%
           family={Szyld},
           familyi={S\bibinitperiod},
           given={Martin},
           giveni={M\bibinitperiod}}}%
      }
      \strng{namehash}{4a147a413765192cb1534ba8ac2f2ffc}
      \strng{fullhash}{4a147a413765192cb1534ba8ac2f2ffc}
      \strng{fullhashraw}{4a147a413765192cb1534ba8ac2f2ffc}
      \strng{bibnamehash}{4a147a413765192cb1534ba8ac2f2ffc}
      \strng{authorbibnamehash}{4a147a413765192cb1534ba8ac2f2ffc}
      \strng{authornamehash}{4a147a413765192cb1534ba8ac2f2ffc}
      \strng{authorfullhash}{4a147a413765192cb1534ba8ac2f2ffc}
      \strng{authorfullhashraw}{4a147a413765192cb1534ba8ac2f2ffc}
      \field{sortinit}{C}
      \field{sortinithash}{4d103a86280481745c9c897c925753c0}
      \field{labelnamesource}{author}
      \field{labeltitlesource}{title}
      \field{eprinttype}{arXiv}
      \field{title}{Double Fibrations}
      \field{year}{2022}
      \verb{eprint}
      \verb 2205.15240
      \endverb
    \endentry
    \entry{dangelo25}{online}{}{}
      \name{author}{2}{}{%
        {{hash=b842da893fc244f409e7b0d657ba563a}{%
           family={D'Angelo},
           familyi={D\bibinitperiod},
           given={Keri},
           giveni={K\bibinitperiod}}}%
        {{hash=33bda5bc38c3b05f23e9a99b5e4efb27}{%
           family={Libkind},
           familyi={L\bibinitperiod},
           given={Sophie},
           giveni={S\bibinitperiod}}}%
      }
      \strng{namehash}{628530c722b8f885d512306e48619028}
      \strng{fullhash}{628530c722b8f885d512306e48619028}
      \strng{fullhashraw}{628530c722b8f885d512306e48619028}
      \strng{bibnamehash}{628530c722b8f885d512306e48619028}
      \strng{authorbibnamehash}{628530c722b8f885d512306e48619028}
      \strng{authornamehash}{628530c722b8f885d512306e48619028}
      \strng{authorfullhash}{628530c722b8f885d512306e48619028}
      \strng{authorfullhashraw}{628530c722b8f885d512306e48619028}
      \field{sortinit}{D}
      \field{sortinithash}{6f385f66841fb5e82009dc833c761848}
      \field{labelnamesource}{author}
      \field{labeltitlesource}{title}
      \field{eprinttype}{arXiv}
      \field{title}{Dependent Directed Wiring Diagrams for Composing Instantaneous Systems}
      \field{year}{2025}
      \verb{eprint}
      \verb 2503.05457
      \endverb
    \endentry
    \entry{dalrymple26imprecise}{online}{}{}
      \name{author}{1}{}{%
        {{hash=26c0deab5b0d6fb9287402b9d735db9c}{%
           family={Dalrymple},
           familyi={D\bibinitperiod},
           given={David\bibnamedelima "davidad"},
           giveni={D\bibinitperiod\bibinitdelim "\bibinitperiod}}}%
      }
      \strng{namehash}{26c0deab5b0d6fb9287402b9d735db9c}
      \strng{fullhash}{26c0deab5b0d6fb9287402b9d735db9c}
      \strng{fullhashraw}{26c0deab5b0d6fb9287402b9d735db9c}
      \strng{bibnamehash}{26c0deab5b0d6fb9287402b9d735db9c}
      \strng{authorbibnamehash}{26c0deab5b0d6fb9287402b9d735db9c}
      \strng{authornamehash}{26c0deab5b0d6fb9287402b9d735db9c}
      \strng{authorfullhash}{26c0deab5b0d6fb9287402b9d735db9c}
      \strng{authorfullhashraw}{26c0deab5b0d6fb9287402b9d735db9c}
      \field{extraname}{1}
      \field{sortinit}{D}
      \field{sortinithash}{6f385f66841fb5e82009dc833c761848}
      \field{labelnamesource}{author}
      \field{labeltitlesource}{title}
      \field{note}{Post on LessWrong, 2026-07-29}
      \field{subtitle}{A tiny introduction}
      \field{title}{Imprecise beliefs}
      \field{urlday}{31}
      \field{urlmonth}{7}
      \field{urlyear}{2026}
      \field{year}{2026}
      \field{urldateera}{ce}
      \verb{urlraw}
      \verb https://www.lesswrong.com/posts/e7Pd4Q9TF7jFdmPgz/imprecise-beliefs-a-tiny-introduction
      \endverb
      \verb{url}
      \verb https://www.lesswrong.com/posts/e7Pd4Q9TF7jFdmPgz/imprecise-beliefs-a-tiny-introduction
      \endverb
    \endentry
    \entry{dalrymple-sgai-thesis}{report}{}{}
      \name{author}{1}{}{%
        {{hash=26c0deab5b0d6fb9287402b9d735db9c}{%
           family={Dalrymple},
           familyi={D\bibinitperiod},
           given={David\bibnamedelima "davidad"},
           giveni={D\bibinitperiod\bibinitdelim "\bibinitperiod}}}%
      }
      \list{institution}{2}{%
        {The Advanced Research}%
        {Invention Agency (ARIA)}%
      }
      \strng{namehash}{26c0deab5b0d6fb9287402b9d735db9c}
      \strng{fullhash}{26c0deab5b0d6fb9287402b9d735db9c}
      \strng{fullhashraw}{26c0deab5b0d6fb9287402b9d735db9c}
      \strng{bibnamehash}{26c0deab5b0d6fb9287402b9d735db9c}
      \strng{authorbibnamehash}{26c0deab5b0d6fb9287402b9d735db9c}
      \strng{authornamehash}{26c0deab5b0d6fb9287402b9d735db9c}
      \strng{authorfullhash}{26c0deab5b0d6fb9287402b9d735db9c}
      \strng{authorfullhashraw}{26c0deab5b0d6fb9287402b9d735db9c}
      \field{extraname}{2}
      \field{sortinit}{D}
      \field{sortinithash}{6f385f66841fb5e82009dc833c761848}
      \field{labelnamesource}{author}
      \field{labeltitlesource}{title}
      \field{subtitle}{Constructing guaranteed safety (programme thesis)}
      \field{title}{Safeguarded AI}
      \field{type}{Programme thesis}
    \endentry
    \entry{dalrymple2026psd}{misc}{}{}
      \name{author}{1}{}{%
        {{hash=69473ae73f779bc670ace77507f5871e}{%
           family={Dalrymple},
           familyi={D\bibinitperiod},
           given={David\bibnamedelima ``davidad''},
           giveni={D\bibinitperiod\bibinitdelim ``\bibinitperiod}}}%
      }
      \strng{namehash}{69473ae73f779bc670ace77507f5871e}
      \strng{fullhash}{69473ae73f779bc670ace77507f5871e}
      \strng{fullhashraw}{69473ae73f779bc670ace77507f5871e}
      \strng{bibnamehash}{69473ae73f779bc670ace77507f5871e}
      \strng{authorbibnamehash}{69473ae73f779bc670ace77507f5871e}
      \strng{authornamehash}{69473ae73f779bc670ace77507f5871e}
      \strng{authorfullhash}{69473ae73f779bc670ace77507f5871e}
      \strng{authorfullhashraw}{69473ae73f779bc670ace77507f5871e}
      \field{extraname}{3}
      \field{sortinit}{D}
      \field{sortinithash}{6f385f66841fb5e82009dc833c761848}
      \field{labelnamesource}{author}
      \field{labeltitlesource}{title}
      \field{howpublished}{Post, X (formerly Twitter), status 2022316806094913669}
      \field{note}{(The triangular-inverse method: sign-and-pattern checks on $A \cdot Z$ for an untrusted triangular witness $Z$; no divisions, witness accuracy never essential.)}
      \field{title}{Certifying strict positive definiteness via an approximate triangular inverse}
      \field{year}{2026}
    \endentry
    \entry{dash22}{online}{}{}
      \name{author}{4}{}{%
        {{hash=ab0876505c4d72dc5a7bcbb16f3dd43b}{%
           family={Dash},
           familyi={D\bibinitperiod},
           given={Swaraj},
           giveni={S\bibinitperiod}}}%
        {{hash=4674c68cc8ce64be7b7606fc6a1da07d}{%
           family={Kaddar},
           familyi={K\bibinitperiod},
           given={Younesse},
           giveni={Y\bibinitperiod}}}%
        {{hash=d90e7bb4e0e8677bfbddb5bf860cec07}{%
           family={Paquet},
           familyi={P\bibinitperiod},
           given={Hugo},
           giveni={H\bibinitperiod}}}%
        {{hash=5351d3a060653e114d956a0a3700d556}{%
           family={Staton},
           familyi={S\bibinitperiod},
           given={Sam},
           giveni={S\bibinitperiod}}}%
      }
      \strng{namehash}{f1a60ad86c4526f7835f4e62014de701}
      \strng{fullhash}{f1a60ad86c4526f7835f4e62014de701}
      \strng{fullhashraw}{f1a60ad86c4526f7835f4e62014de701}
      \strng{bibnamehash}{f1a60ad86c4526f7835f4e62014de701}
      \strng{authorbibnamehash}{f1a60ad86c4526f7835f4e62014de701}
      \strng{authornamehash}{f1a60ad86c4526f7835f4e62014de701}
      \strng{authorfullhash}{f1a60ad86c4526f7835f4e62014de701}
      \strng{authorfullhashraw}{f1a60ad86c4526f7835f4e62014de701}
      \field{sortinit}{D}
      \field{sortinithash}{6f385f66841fb5e82009dc833c761848}
      \field{labelnamesource}{author}
      \field{labeltitlesource}{title}
      \field{eprinttype}{arXiv}
      \field{title}{Affine Monads and Lazy Structures for Bayesian Programming}
      \field{year}{2022}
      \verb{eprint}
      \verb 2212.07250
      \endverb
    \endentry
    \entry{dash21}{online}{}{}
      \name{author}{2}{}{%
        {{hash=ab0876505c4d72dc5a7bcbb16f3dd43b}{%
           family={Dash},
           familyi={D\bibinitperiod},
           given={Swaraj},
           giveni={S\bibinitperiod}}}%
        {{hash=5351d3a060653e114d956a0a3700d556}{%
           family={Staton},
           familyi={S\bibinitperiod},
           given={Sam},
           giveni={S\bibinitperiod}}}%
      }
      \strng{namehash}{dc6c34ea6f9f7cd29e9e017a8b911724}
      \strng{fullhash}{dc6c34ea6f9f7cd29e9e017a8b911724}
      \strng{fullhashraw}{dc6c34ea6f9f7cd29e9e017a8b911724}
      \strng{bibnamehash}{dc6c34ea6f9f7cd29e9e017a8b911724}
      \strng{authorbibnamehash}{dc6c34ea6f9f7cd29e9e017a8b911724}
      \strng{authornamehash}{dc6c34ea6f9f7cd29e9e017a8b911724}
      \strng{authorfullhash}{dc6c34ea6f9f7cd29e9e017a8b911724}
      \strng{authorfullhashraw}{dc6c34ea6f9f7cd29e9e017a8b911724}
      \field{sortinit}{D}
      \field{sortinithash}{6f385f66841fb5e82009dc833c761848}
      \field{labelnamesource}{author}
      \field{labeltitlesource}{title}
      \field{eprinttype}{arXiv}
      \field{title}{Monads for Measurable Queries in Probabilistic Databases}
      \field{year}{2021}
      \verb{eprint}
      \verb 2112.14048
      \endverb
    \endentry
    \entry{dashkovskiyrufferwirth2007}{article}{}{}
      \name{author}{3}{}{%
        {{hash=9de3d0502084751987c9662acd8838f9}{%
           family={Dashkovskiy},
           familyi={D\bibinitperiod},
           given={S.},
           giveni={S\bibinitperiod}}}%
        {{hash=8f20fdb8f3a2f3413f201a2a4e3c05f3}{%
           family={Rüffer},
           familyi={R\bibinitperiod},
           given={B.\bibnamedelimi S.},
           giveni={B\bibinitperiod\bibinitdelim S\bibinitperiod}}}%
        {{hash=28cd50b7d7cb05346d4c95d94a991462}{%
           family={Wirth},
           familyi={W\bibinitperiod},
           given={F.\bibnamedelimi R.},
           giveni={F\bibinitperiod\bibinitdelim R\bibinitperiod}}}%
      }
      \strng{namehash}{45226fda762d5cbfc9a1e9a90e1f63e4}
      \strng{fullhash}{45226fda762d5cbfc9a1e9a90e1f63e4}
      \strng{fullhashraw}{45226fda762d5cbfc9a1e9a90e1f63e4}
      \strng{bibnamehash}{45226fda762d5cbfc9a1e9a90e1f63e4}
      \strng{authorbibnamehash}{45226fda762d5cbfc9a1e9a90e1f63e4}
      \strng{authornamehash}{45226fda762d5cbfc9a1e9a90e1f63e4}
      \strng{authorfullhash}{45226fda762d5cbfc9a1e9a90e1f63e4}
      \strng{authorfullhashraw}{45226fda762d5cbfc9a1e9a90e1f63e4}
      \field{sortinit}{D}
      \field{sortinithash}{6f385f66841fb5e82009dc833c761848}
      \field{labelnamesource}{author}
      \field{labeltitlesource}{title}
      \field{journaltitle}{Mathematics of Control, Signals, and Systems}
      \field{note}{(The $n$-component ISS small-gain theorem over arbitrary interconnection graphs whose certificate form the \rulename{SmallGainN} section catalogues; its monotone-operator fixpoint machinery lives producer-side as existence theory for the $(\lambda, B, \theta)$ side data.)}
      \field{number}{2}
      \field{title}{An {ISS} small gain theorem for general networks}
      \field{volume}{19}
      \field{year}{2007}
      \field{pages}{93\bibrangedash 122}
      \range{pages}{30}
      \verb{doi}
      \verb 10.1007/s00498-007-0014-8
      \endverb
    \endentry
    \entry{davis-1993}{book}{}{}
      \name{author}{1}{}{%
        {{hash=01760e459a2450847318fe481c960298}{%
           family={Davis},
           familyi={D\bibinitperiod},
           given={M.\bibnamedelimi H.\bibnamedelimi A.},
           giveni={M\bibinitperiod\bibinitdelim H\bibinitperiod\bibinitdelim A\bibinitperiod}}}%
      }
      \list{location}{1}{%
        {London}%
      }
      \list{publisher}{1}{%
        {Chapman \& Hall}%
      }
      \strng{namehash}{01760e459a2450847318fe481c960298}
      \strng{fullhash}{01760e459a2450847318fe481c960298}
      \strng{fullhashraw}{01760e459a2450847318fe481c960298}
      \strng{bibnamehash}{01760e459a2450847318fe481c960298}
      \strng{authorbibnamehash}{01760e459a2450847318fe481c960298}
      \strng{authornamehash}{01760e459a2450847318fe481c960298}
      \strng{authorfullhash}{01760e459a2450847318fe481c960298}
      \strng{authorfullhashraw}{01760e459a2450847318fe481c960298}
      \field{sortinit}{D}
      \field{sortinithash}{6f385f66841fb5e82009dc833c761848}
      \field{labelnamesource}{author}
      \field{labeltitlesource}{title}
      \field{number}{49}
      \field{series}{Monographs on Statistics and Applied Probability}
      \field{title}{Markov Models and Optimization}
      \field{year}{1993}
    \endentry
    \entry{desharnais2004}{article}{}{}
      \name{author}{4}{}{%
        {{hash=d4ae7bee01849dab8ef0b5ce1674d7f2}{%
           family={Desharnais},
           familyi={D\bibinitperiod},
           given={J.},
           giveni={J\bibinitperiod}}}%
        {{hash=3814c0eb5239d5b46ed5e532af1423c3}{%
           family={Gupta},
           familyi={G\bibinitperiod},
           given={V.},
           giveni={V\bibinitperiod}}}%
        {{hash=277ec1c09e69793913c952e1836f3313}{%
           family={Jagadeesan},
           familyi={J\bibinitperiod},
           given={R.},
           giveni={R\bibinitperiod}}}%
        {{hash=4160380fd32388e79ebf0d013e5b9cfa}{%
           family={Panangaden},
           familyi={P\bibinitperiod},
           given={P.},
           giveni={P\bibinitperiod}}}%
      }
      \strng{namehash}{25461b16021918113134468daa4ecd39}
      \strng{fullhash}{25461b16021918113134468daa4ecd39}
      \strng{fullhashraw}{25461b16021918113134468daa4ecd39}
      \strng{bibnamehash}{25461b16021918113134468daa4ecd39}
      \strng{authorbibnamehash}{25461b16021918113134468daa4ecd39}
      \strng{authornamehash}{25461b16021918113134468daa4ecd39}
      \strng{authorfullhash}{25461b16021918113134468daa4ecd39}
      \strng{authorfullhashraw}{25461b16021918113134468daa4ecd39}
      \field{sortinit}{D}
      \field{sortinithash}{6f385f66841fb5e82009dc833c761848}
      \field{labelnamesource}{author}
      \field{labeltitlesource}{title}
      \field{journaltitle}{Theoretical Computer Science}
      \field{note}{(Behavioral pseudometrics for probabilistic systems.)}
      \field{number}{3}
      \field{title}{Metrics for labelled Markov processes}
      \field{volume}{318}
      \field{year}{2004}
      \field{pages}{323\bibrangedash 354}
      \range{pages}{32}
    \endentry
    \entry{deutschnashremmel2008}{inproceedings}{}{}
      \name{author}{3}{}{%
        {{hash=7cb9b3fe27c27c0439e74d71b4e4ba03}{%
           family={Deutsch},
           familyi={D\bibinitperiod},
           given={Alin},
           giveni={A\bibinitperiod}}}%
        {{hash=ed60b3d995a39b53dbdd540c8279e290}{%
           family={Nash},
           familyi={N\bibinitperiod},
           given={Alan},
           giveni={A\bibinitperiod}}}%
        {{hash=85480632cc50270c0f440c87c290a907}{%
           family={Remmel},
           familyi={R\bibinitperiod},
           given={Jeff},
           giveni={J\bibinitperiod}}}%
      }
      \strng{namehash}{b18faa0abf250562f2701a6453d38d4e}
      \strng{fullhash}{b18faa0abf250562f2701a6453d38d4e}
      \strng{fullhashraw}{b18faa0abf250562f2701a6453d38d4e}
      \strng{bibnamehash}{b18faa0abf250562f2701a6453d38d4e}
      \strng{authorbibnamehash}{b18faa0abf250562f2701a6453d38d4e}
      \strng{authornamehash}{b18faa0abf250562f2701a6453d38d4e}
      \strng{authorfullhash}{b18faa0abf250562f2701a6453d38d4e}
      \strng{authorfullhashraw}{b18faa0abf250562f2701a6453d38d4e}
      \field{sortinit}{D}
      \field{sortinithash}{6f385f66841fb5e82009dc833c761848}
      \field{labelnamesource}{author}
      \field{labeltitlesource}{title}
      \field{booktitle}{Proceedings of the 27th ACM SIGMOD-SIGACT-SIGART Symposium on Principles of Database Systems (PODS '08)}
      \field{title}{The Chase Revisited}
      \field{year}{2008}
      \field{pages}{149\bibrangedash 158}
      \range{pages}{10}
      \verb{doi}
      \verb 10.1145/1376916.1376938
      \endverb
    \endentry
    \entry{doeblin1938}{article}{}{}
      \name{author}{1}{}{%
        {{hash=1744a279f61578688bd3d6d0061f7013}{%
           family={Doeblin},
           familyi={D\bibinitperiod},
           given={W.},
           giveni={W\bibinitperiod}}}%
      }
      \strng{namehash}{1744a279f61578688bd3d6d0061f7013}
      \strng{fullhash}{1744a279f61578688bd3d6d0061f7013}
      \strng{fullhashraw}{1744a279f61578688bd3d6d0061f7013}
      \strng{bibnamehash}{1744a279f61578688bd3d6d0061f7013}
      \strng{authorbibnamehash}{1744a279f61578688bd3d6d0061f7013}
      \strng{authornamehash}{1744a279f61578688bd3d6d0061f7013}
      \strng{authorfullhash}{1744a279f61578688bd3d6d0061f7013}
      \strng{authorfullhashraw}{1744a279f61578688bd3d6d0061f7013}
      \field{sortinit}{D}
      \field{sortinithash}{6f385f66841fb5e82009dc833c761848}
      \field{labelnamesource}{author}
      \field{labeltitlesource}{title}
      \field{journaltitle}{Revue Mathématique de l'Union Interbalkanique}
      \field{note}{(The origin of the coupling method for Markov chains.)}
      \field{title}{Exposé de la théorie des chaînes simples constantes de Markov à un nombre fini d'états}
      \field{volume}{2}
      \field{year}{1938}
      \field{pages}{77\bibrangedash 105}
      \range{pages}{29}
    \endentry
    \entry{doob1953}{book}{}{}
      \name{author}{1}{}{%
        {{hash=b1f7ca538b74cc0f274fecc6b1699907}{%
           family={Doob},
           familyi={D\bibinitperiod},
           given={J.\bibnamedelimi L.},
           giveni={J\bibinitperiod\bibinitdelim L\bibinitperiod}}}%
      }
      \list{location}{1}{%
        {New York}%
      }
      \list{publisher}{1}{%
        {Wiley}%
      }
      \strng{namehash}{b1f7ca538b74cc0f274fecc6b1699907}
      \strng{fullhash}{b1f7ca538b74cc0f274fecc6b1699907}
      \strng{fullhashraw}{b1f7ca538b74cc0f274fecc6b1699907}
      \strng{bibnamehash}{b1f7ca538b74cc0f274fecc6b1699907}
      \strng{authorbibnamehash}{b1f7ca538b74cc0f274fecc6b1699907}
      \strng{authornamehash}{b1f7ca538b74cc0f274fecc6b1699907}
      \strng{authorfullhash}{b1f7ca538b74cc0f274fecc6b1699907}
      \strng{authorfullhashraw}{b1f7ca538b74cc0f274fecc6b1699907}
      \field{sortinit}{D}
      \field{sortinithash}{6f385f66841fb5e82009dc833c761848}
      \field{labelnamesource}{author}
      \field{labeltitlesource}{title}
      \field{note}{(Systematic supermartingale theory; maximal inequalities and optional sampling.)}
      \field{title}{Stochastic Processes}
      \field{year}{1953}
    \endentry
    \entry{durrett2019}{book}{}{}
      \name{author}{1}{}{%
        {{hash=03a107fab36174e7da5974a64ad1812d}{%
           family={Durrett},
           familyi={D\bibinitperiod},
           given={R.},
           giveni={R\bibinitperiod}}}%
      }
      \list{publisher}{1}{%
        {Cambridge University Press}%
      }
      \strng{namehash}{03a107fab36174e7da5974a64ad1812d}
      \strng{fullhash}{03a107fab36174e7da5974a64ad1812d}
      \strng{fullhashraw}{03a107fab36174e7da5974a64ad1812d}
      \strng{bibnamehash}{03a107fab36174e7da5974a64ad1812d}
      \strng{authorbibnamehash}{03a107fab36174e7da5974a64ad1812d}
      \strng{authornamehash}{03a107fab36174e7da5974a64ad1812d}
      \strng{authorfullhash}{03a107fab36174e7da5974a64ad1812d}
      \strng{authorfullhashraw}{03a107fab36174e7da5974a64ad1812d}
      \field{sortinit}{D}
      \field{sortinithash}{6f385f66841fb5e82009dc833c761848}
      \field{labelnamesource}{author}
      \field{labeltitlesource}{title}
      \field{edition}{5}
      \field{note}{(Borel--Cantelli.)}
      \field{title}{Probability: Theory and Examples}
      \field{year}{2019}
    \endentry
    \entry{ehlers2017}{inproceedings}{}{}
      \name{author}{1}{}{%
        {{hash=ed2b5bc7b888969c7067d519246ae7bd}{%
           family={Ehlers},
           familyi={E\bibinitperiod},
           given={R.},
           giveni={R\bibinitperiod}}}%
      }
      \strng{namehash}{ed2b5bc7b888969c7067d519246ae7bd}
      \strng{fullhash}{ed2b5bc7b888969c7067d519246ae7bd}
      \strng{fullhashraw}{ed2b5bc7b888969c7067d519246ae7bd}
      \strng{bibnamehash}{ed2b5bc7b888969c7067d519246ae7bd}
      \strng{authorbibnamehash}{ed2b5bc7b888969c7067d519246ae7bd}
      \strng{authornamehash}{ed2b5bc7b888969c7067d519246ae7bd}
      \strng{authorfullhash}{ed2b5bc7b888969c7067d519246ae7bd}
      \strng{authorfullhashraw}{ed2b5bc7b888969c7067d519246ae7bd}
      \field{sortinit}{E}
      \field{sortinithash}{8da8a182d344d5b9047633dfc0cc9131}
      \field{labelnamesource}{author}
      \field{labeltitlesource}{title}
      \field{booktitle}{Proceedings of ATVA 2017 (Automated Technology for Verification and Analysis)}
      \field{note}{(The ``triangle'' linear relaxation of ReLU.)}
      \field{title}{Formal verification of piece-wise linear feed-forward neural networks}
      \field{year}{2017}
    \endentry
    \entry{eisentraut-hermanns-zhang-2010}{inproceedings}{}{}
      \name{author}{3}{}{%
        {{hash=cce65bc5d18842d0e673ae2e4e695575}{%
           family={Eisentraut},
           familyi={E\bibinitperiod},
           given={Christian},
           giveni={C\bibinitperiod}}}%
        {{hash=b396735e2150877fec4cee61d80d5430}{%
           family={Hermanns},
           familyi={H\bibinitperiod},
           given={Holger},
           giveni={H\bibinitperiod}}}%
        {{hash=3dd8089f3f84e2a69c24625f7a549bc6}{%
           family={Zhang},
           familyi={Z\bibinitperiod},
           given={Lijun},
           giveni={L\bibinitperiod}}}%
      }
      \list{publisher}{1}{%
        {IEEE Computer Society}%
      }
      \strng{namehash}{4b3558394217b2ef6c00d5ea9258de07}
      \strng{fullhash}{4b3558394217b2ef6c00d5ea9258de07}
      \strng{fullhashraw}{4b3558394217b2ef6c00d5ea9258de07}
      \strng{bibnamehash}{4b3558394217b2ef6c00d5ea9258de07}
      \strng{authorbibnamehash}{4b3558394217b2ef6c00d5ea9258de07}
      \strng{authornamehash}{4b3558394217b2ef6c00d5ea9258de07}
      \strng{authorfullhash}{4b3558394217b2ef6c00d5ea9258de07}
      \strng{authorfullhashraw}{4b3558394217b2ef6c00d5ea9258de07}
      \field{sortinit}{E}
      \field{sortinithash}{8da8a182d344d5b9047633dfc0cc9131}
      \field{labelnamesource}{author}
      \field{labeltitlesource}{title}
      \field{booktitle}{Proceedings of the 25th Annual {IEEE} Symposium on Logic in Computer Science ({LICS} 2010)}
      \field{title}{On Probabilistic Automata in Continuous Time}
      \field{year}{2010}
      \field{pages}{342\bibrangedash 351}
      \range{pages}{10}
    \endentry
    \entry{emersonlei1987}{article}{}{}
      \name{author}{2}{}{%
        {{hash=45756aa08f92810b89c0f5aa777fef6f}{%
           family={Emerson},
           familyi={E\bibinitperiod},
           given={E.\bibnamedelimi A.},
           giveni={E\bibinitperiod\bibinitdelim A\bibinitperiod}}}%
        {{hash=fa410896f601d90765e11d341f3b05ed}{%
           family={Lei},
           familyi={L\bibinitperiod},
           given={C.-L.},
           giveni={C\bibinithyphendelim L\bibinitperiod}}}%
      }
      \strng{namehash}{f295b742619d1534511060538f777a82}
      \strng{fullhash}{f295b742619d1534511060538f777a82}
      \strng{fullhashraw}{f295b742619d1534511060538f777a82}
      \strng{bibnamehash}{f295b742619d1534511060538f777a82}
      \strng{authorbibnamehash}{f295b742619d1534511060538f777a82}
      \strng{authornamehash}{f295b742619d1534511060538f777a82}
      \strng{authorfullhash}{f295b742619d1534511060538f777a82}
      \strng{authorfullhashraw}{f295b742619d1534511060538f777a82}
      \field{sortinit}{E}
      \field{sortinithash}{8da8a182d344d5b9047633dfc0cc9131}
      \field{labelnamesource}{author}
      \field{labeltitlesource}{title}
      \field{journaltitle}{Science of Computer Programming}
      \field{note}{(Acceptance as positive boolean combinations of finitely-often/infinitely-often atoms --- Emerson--Lei conditions.)}
      \field{number}{3}
      \field{title}{Modalities for model checking: branching time logic strikes back}
      \field{volume}{8}
      \field{year}{1987}
      \field{pages}{275\bibrangedash 306}
      \range{pages}{32}
    \endentry
    \entry{ethier-kurtz-1986}{book}{}{}
      \name{author}{2}{}{%
        {{hash=9dfd3fb2f90c9fd2d51ad5d6b983913d}{%
           family={Ethier},
           familyi={E\bibinitperiod},
           given={Stewart\bibnamedelima N.},
           giveni={S\bibinitperiod\bibinitdelim N\bibinitperiod}}}%
        {{hash=c10f1148db53724ba2b47997f5b93cd7}{%
           family={Kurtz},
           familyi={K\bibinitperiod},
           given={Thomas\bibnamedelima G.},
           giveni={T\bibinitperiod\bibinitdelim G\bibinitperiod}}}%
      }
      \list{location}{1}{%
        {New York}%
      }
      \list{publisher}{1}{%
        {Wiley}%
      }
      \strng{namehash}{206f9b3f81c52f534d2c8a392f44d7b1}
      \strng{fullhash}{206f9b3f81c52f534d2c8a392f44d7b1}
      \strng{fullhashraw}{206f9b3f81c52f534d2c8a392f44d7b1}
      \strng{bibnamehash}{206f9b3f81c52f534d2c8a392f44d7b1}
      \strng{authorbibnamehash}{206f9b3f81c52f534d2c8a392f44d7b1}
      \strng{authornamehash}{206f9b3f81c52f534d2c8a392f44d7b1}
      \strng{authorfullhash}{206f9b3f81c52f534d2c8a392f44d7b1}
      \strng{authorfullhashraw}{206f9b3f81c52f534d2c8a392f44d7b1}
      \field{sortinit}{E}
      \field{sortinithash}{8da8a182d344d5b9047633dfc0cc9131}
      \field{labelnamesource}{author}
      \field{labeltitlesource}{title}
      \field{series}{Wiley Series in Probability and Mathematical Statistics}
      \field{title}{Markov Processes: Characterization and Convergence}
      \field{year}{1986}
    \endentry
    \entry{evans-2010}{book}{}{}
      \name{author}{1}{}{%
        {{hash=0d72932a7775f85a51622d3af75ae26e}{%
           family={Evans},
           familyi={E\bibinitperiod},
           given={Lawrence\bibnamedelima C.},
           giveni={L\bibinitperiod\bibinitdelim C\bibinitperiod}}}%
      }
      \list{location}{1}{%
        {Providence}%
      }
      \list{publisher}{1}{%
        {American Mathematical Society}%
      }
      \strng{namehash}{0d72932a7775f85a51622d3af75ae26e}
      \strng{fullhash}{0d72932a7775f85a51622d3af75ae26e}
      \strng{fullhashraw}{0d72932a7775f85a51622d3af75ae26e}
      \strng{bibnamehash}{0d72932a7775f85a51622d3af75ae26e}
      \strng{authorbibnamehash}{0d72932a7775f85a51622d3af75ae26e}
      \strng{authornamehash}{0d72932a7775f85a51622d3af75ae26e}
      \strng{authorfullhash}{0d72932a7775f85a51622d3af75ae26e}
      \strng{authorfullhashraw}{0d72932a7775f85a51622d3af75ae26e}
      \field{sortinit}{E}
      \field{sortinithash}{8da8a182d344d5b9047633dfc0cc9131}
      \field{labelnamesource}{author}
      \field{labeltitlesource}{title}
      \field{edition}{2}
      \field{series}{Graduate Studies in Mathematics}
      \field{title}{Partial Differential Equations}
      \field{volume}{19}
      \field{year}{2010}
    \endentry
    \entry{everdij-blom-2010}{incollection}{}{}
      \name{author}{2}{}{%
        {{hash=d3f5b0048f31d660d96a4fe2eede6c83}{%
           family={Everdij},
           familyi={E\bibinitperiod},
           given={Mariken\bibnamedelimb H.\bibnamedelimi C.},
           giveni={M\bibinitperiod\bibinitdelim H\bibinitperiod\bibinitdelim C\bibinitperiod}}}%
        {{hash=2c55a6fb30bfd4922dbf7e93df7b053d}{%
           family={Blom},
           familyi={B\bibinitperiod},
           given={Henk\bibnamedelimb A.\bibnamedelimi P.},
           giveni={H\bibinitperiod\bibinitdelim A\bibinitperiod\bibinitdelim P\bibinitperiod}}}%
      }
      \name{editor}{1}{}{%
        {{hash=c120142f361a450b240ee1883e8a6cb5}{%
           family={Pawlewski},
           familyi={P\bibinitperiod},
           given={Pawel},
           giveni={P\bibinitperiod}}}%
      }
      \list{publisher}{1}{%
        {InTech}%
      }
      \strng{namehash}{4cdab0dcc6c1145ec5a84d3b8fe3f459}
      \strng{fullhash}{4cdab0dcc6c1145ec5a84d3b8fe3f459}
      \strng{fullhashraw}{4cdab0dcc6c1145ec5a84d3b8fe3f459}
      \strng{bibnamehash}{4cdab0dcc6c1145ec5a84d3b8fe3f459}
      \strng{authorbibnamehash}{4cdab0dcc6c1145ec5a84d3b8fe3f459}
      \strng{authornamehash}{4cdab0dcc6c1145ec5a84d3b8fe3f459}
      \strng{authorfullhash}{4cdab0dcc6c1145ec5a84d3b8fe3f459}
      \strng{authorfullhashraw}{4cdab0dcc6c1145ec5a84d3b8fe3f459}
      \strng{editorbibnamehash}{c120142f361a450b240ee1883e8a6cb5}
      \strng{editornamehash}{c120142f361a450b240ee1883e8a6cb5}
      \strng{editorfullhash}{c120142f361a450b240ee1883e8a6cb5}
      \strng{editorfullhashraw}{c120142f361a450b240ee1883e8a6cb5}
      \field{sortinit}{E}
      \field{sortinithash}{8da8a182d344d5b9047633dfc0cc9131}
      \field{labelnamesource}{author}
      \field{labeltitlesource}{title}
      \field{booktitle}{Petri Nets}
      \field{chapter}{12}
      \field{eprinttype}{arXiv}
      \field{title}{Hybrid state {P}etri nets which have the analysis power of stochastic hybrid systems and the formal verification power of automata}
      \field{year}{2010}
      \field{pages}{227\bibrangedash 252}
      \range{pages}{26}
      \verb{eprint}
      \verb 1003.4812
      \endverb
    \endentry
    \entry{fagin2005data}{article}{}{}
      \name{author}{4}{}{%
        {{hash=ac675e1b6c0d303691ca1e438688780e}{%
           family={Fagin},
           familyi={F\bibinitperiod},
           given={Ronald},
           giveni={R\bibinitperiod}}}%
        {{hash=b49a50fb7ffbb4bd60ebec4c6b0dda5a}{%
           family={Kolaitis},
           familyi={K\bibinitperiod},
           given={Phokion\bibnamedelima G.},
           giveni={P\bibinitperiod\bibinitdelim G\bibinitperiod}}}%
        {{hash=5ff644ee09a712c2cd994b537d1de2fd}{%
           family={Miller},
           familyi={M\bibinitperiod},
           given={Renée\bibnamedelima J.},
           giveni={R\bibinitperiod\bibinitdelim J\bibinitperiod}}}%
        {{hash=cca4707fa68c79d531c33d107a5d031e}{%
           family={Popa},
           familyi={P\bibinitperiod},
           given={Lucian},
           giveni={L\bibinitperiod}}}%
      }
      \strng{namehash}{362443ea3ec2bf08f6413d3c95997342}
      \strng{fullhash}{362443ea3ec2bf08f6413d3c95997342}
      \strng{fullhashraw}{362443ea3ec2bf08f6413d3c95997342}
      \strng{bibnamehash}{362443ea3ec2bf08f6413d3c95997342}
      \strng{authorbibnamehash}{362443ea3ec2bf08f6413d3c95997342}
      \strng{authornamehash}{362443ea3ec2bf08f6413d3c95997342}
      \strng{authorfullhash}{362443ea3ec2bf08f6413d3c95997342}
      \strng{authorfullhashraw}{362443ea3ec2bf08f6413d3c95997342}
      \field{sortinit}{F}
      \field{sortinithash}{2638baaa20439f1b5a8f80c6c08a13b4}
      \field{labelnamesource}{author}
      \field{labeltitlesource}{title}
      \field{journaltitle}{Theoretical Computer Science}
      \field{number}{1}
      \field{subtitle}{Semantics and Query Answering}
      \field{title}{Data Exchange}
      \field{volume}{336}
      \field{year}{2005}
      \field{pages}{89\bibrangedash 124}
      \range{pages}{36}
      \verb{doi}
      \verb 10.1016/j.tcs.2004.10.033
      \endverb
    \endentry
    \entry{farkas1902}{article}{}{}
      \name{author}{1}{}{%
        {{hash=4bba8ae63c9ada4ec0307f63e9877ad6}{%
           family={Farkas},
           familyi={F\bibinitperiod},
           given={J.},
           giveni={J\bibinitperiod}}}%
      }
      \strng{namehash}{4bba8ae63c9ada4ec0307f63e9877ad6}
      \strng{fullhash}{4bba8ae63c9ada4ec0307f63e9877ad6}
      \strng{fullhashraw}{4bba8ae63c9ada4ec0307f63e9877ad6}
      \strng{bibnamehash}{4bba8ae63c9ada4ec0307f63e9877ad6}
      \strng{authorbibnamehash}{4bba8ae63c9ada4ec0307f63e9877ad6}
      \strng{authornamehash}{4bba8ae63c9ada4ec0307f63e9877ad6}
      \strng{authorfullhash}{4bba8ae63c9ada4ec0307f63e9877ad6}
      \strng{authorfullhashraw}{4bba8ae63c9ada4ec0307f63e9877ad6}
      \field{sortinit}{F}
      \field{sortinithash}{2638baaa20439f1b5a8f80c6c08a13b4}
      \field{labelnamesource}{author}
      \field{labeltitlesource}{title}
      \field{journaltitle}{Journal für die reine und angewandte Mathematik}
      \field{title}{Theorie der einfachen Ungleichungen}
      \field{volume}{124}
      \field{year}{1902}
      \field{pages}{1\bibrangedash 27}
      \range{pages}{27}
    \endentry
    \entry{farouki2012}{article}{}{}
      \name{author}{1}{}{%
        {{hash=928a3aa2b0fdefc7223dfc2d14b0deaa}{%
           family={Farouki},
           familyi={F\bibinitperiod},
           given={R.\bibnamedelimi T.},
           giveni={R\bibinitperiod\bibinitdelim T\bibinitperiod}}}%
      }
      \strng{namehash}{928a3aa2b0fdefc7223dfc2d14b0deaa}
      \strng{fullhash}{928a3aa2b0fdefc7223dfc2d14b0deaa}
      \strng{fullhashraw}{928a3aa2b0fdefc7223dfc2d14b0deaa}
      \strng{bibnamehash}{928a3aa2b0fdefc7223dfc2d14b0deaa}
      \strng{authorbibnamehash}{928a3aa2b0fdefc7223dfc2d14b0deaa}
      \strng{authornamehash}{928a3aa2b0fdefc7223dfc2d14b0deaa}
      \strng{authorfullhash}{928a3aa2b0fdefc7223dfc2d14b0deaa}
      \strng{authorfullhashraw}{928a3aa2b0fdefc7223dfc2d14b0deaa}
      \field{sortinit}{F}
      \field{sortinithash}{2638baaa20439f1b5a8f80c6c08a13b4}
      \field{labelnamesource}{author}
      \field{labeltitlesource}{title}
      \field{journaltitle}{Computer Aided Geometric Design}
      \field{number}{6}
      \field{title}{The Bernstein polynomial basis: a centennial retrospective}
      \field{volume}{29}
      \field{year}{2012}
      \field{pages}{379\bibrangedash 419}
      \range{pages}{41}
    \endentry
    \entry{fell62}{article}{}{}
      \name{author}{1}{}{%
        {{hash=2af931bdf64cef460b17917e0223a48c}{%
           family={Fell},
           familyi={F\bibinitperiod},
           given={J.\bibnamedelimi M.\bibnamedelimi G.},
           giveni={J\bibinitperiod\bibinitdelim M\bibinitperiod\bibinitdelim G\bibinitperiod}}}%
      }
      \strng{namehash}{2af931bdf64cef460b17917e0223a48c}
      \strng{fullhash}{2af931bdf64cef460b17917e0223a48c}
      \strng{fullhashraw}{2af931bdf64cef460b17917e0223a48c}
      \strng{bibnamehash}{2af931bdf64cef460b17917e0223a48c}
      \strng{authorbibnamehash}{2af931bdf64cef460b17917e0223a48c}
      \strng{authornamehash}{2af931bdf64cef460b17917e0223a48c}
      \strng{authorfullhash}{2af931bdf64cef460b17917e0223a48c}
      \strng{authorfullhashraw}{2af931bdf64cef460b17917e0223a48c}
      \field{sortinit}{F}
      \field{sortinithash}{2638baaa20439f1b5a8f80c6c08a13b4}
      \field{labelnamesource}{author}
      \field{labeltitlesource}{title}
      \field{journaltitle}{Proceedings of the American Mathematical Society}
      \field{number}{3}
      \field{title}{A Hausdorff topology for the closed subsets of a locally compact non-Hausdorff space}
      \field{volume}{13}
      \field{year}{1962}
      \field{pages}{472\bibrangedash 476}
      \range{pages}{5}
      \verb{doi}
      \verb 10.1090/S0002-9939-1962-0139135-6
      \endverb
    \endentry
    \entry{fewster15}{online}{}{}
      \name{author}{2}{}{%
        {{hash=4d248c3446a22b48068f9718bfb89766}{%
           family={Fewster},
           familyi={F\bibinitperiod},
           given={Christopher\bibnamedelima J.},
           giveni={C\bibinitperiod\bibinitdelim J\bibinitperiod}}}%
        {{hash=fabe445ea043cc87cdde933d33ee4aac}{%
           family={Verch},
           familyi={V\bibinitperiod},
           given={Rainer},
           giveni={R\bibinitperiod}}}%
      }
      \strng{namehash}{57dad10ec75cb31f20f7d9c3bb0b493e}
      \strng{fullhash}{57dad10ec75cb31f20f7d9c3bb0b493e}
      \strng{fullhashraw}{57dad10ec75cb31f20f7d9c3bb0b493e}
      \strng{bibnamehash}{57dad10ec75cb31f20f7d9c3bb0b493e}
      \strng{authorbibnamehash}{57dad10ec75cb31f20f7d9c3bb0b493e}
      \strng{authornamehash}{57dad10ec75cb31f20f7d9c3bb0b493e}
      \strng{authorfullhash}{57dad10ec75cb31f20f7d9c3bb0b493e}
      \strng{authorfullhashraw}{57dad10ec75cb31f20f7d9c3bb0b493e}
      \field{extraname}{1}
      \field{sortinit}{F}
      \field{sortinithash}{2638baaa20439f1b5a8f80c6c08a13b4}
      \field{labelnamesource}{author}
      \field{labeltitlesource}{title}
      \field{eprinttype}{arXiv}
      \field{title}{Algebraic quantum field theory in curved spacetimes}
      \field{year}{2015}
      \verb{eprint}
      \verb 1504.00586
      \endverb
    \endentry
    \entry{fewster11}{online}{}{}
      \name{author}{2}{}{%
        {{hash=4d248c3446a22b48068f9718bfb89766}{%
           family={Fewster},
           familyi={F\bibinitperiod},
           given={Christopher\bibnamedelima J.},
           giveni={C\bibinitperiod\bibinitdelim J\bibinitperiod}}}%
        {{hash=fabe445ea043cc87cdde933d33ee4aac}{%
           family={Verch},
           familyi={V\bibinitperiod},
           given={Rainer},
           giveni={R\bibinitperiod}}}%
      }
      \strng{namehash}{57dad10ec75cb31f20f7d9c3bb0b493e}
      \strng{fullhash}{57dad10ec75cb31f20f7d9c3bb0b493e}
      \strng{fullhashraw}{57dad10ec75cb31f20f7d9c3bb0b493e}
      \strng{bibnamehash}{57dad10ec75cb31f20f7d9c3bb0b493e}
      \strng{authorbibnamehash}{57dad10ec75cb31f20f7d9c3bb0b493e}
      \strng{authornamehash}{57dad10ec75cb31f20f7d9c3bb0b493e}
      \strng{authorfullhash}{57dad10ec75cb31f20f7d9c3bb0b493e}
      \strng{authorfullhashraw}{57dad10ec75cb31f20f7d9c3bb0b493e}
      \field{extraname}{2}
      \field{sortinit}{F}
      \field{sortinithash}{2638baaa20439f1b5a8f80c6c08a13b4}
      \field{labelnamesource}{author}
      \field{labeltitlesource}{title}
      \field{eprinttype}{arXiv}
      \field{subtitle}{What makes a physical theory the same in all spacetimes?}
      \field{title}{Dynamical locality and covariance}
      \field{year}{2011}
      \verb{doi}
      \verb 10.1007/s00023-012-0165-0
      \endverb
      \verb{eprint}
      \verb 1106.4785
      \endverb
    \endentry
    \entry{filipovic-larsson-2016}{article}{}{}
      \name{author}{2}{}{%
        {{hash=f085fb5858c17e781632b5b7e1add2c2}{%
           family={Filipović},
           familyi={F\bibinitperiod},
           given={Damir},
           giveni={D\bibinitperiod}}}%
        {{hash=ad7df50ff06d509435baf74df17e20cc}{%
           family={Larsson},
           familyi={L\bibinitperiod},
           given={Martin},
           giveni={M\bibinitperiod}}}%
      }
      \strng{namehash}{473369828a7d10d900fb136c70a102c3}
      \strng{fullhash}{473369828a7d10d900fb136c70a102c3}
      \strng{fullhashraw}{473369828a7d10d900fb136c70a102c3}
      \strng{bibnamehash}{473369828a7d10d900fb136c70a102c3}
      \strng{authorbibnamehash}{473369828a7d10d900fb136c70a102c3}
      \strng{authornamehash}{473369828a7d10d900fb136c70a102c3}
      \strng{authorfullhash}{473369828a7d10d900fb136c70a102c3}
      \strng{authorfullhashraw}{473369828a7d10d900fb136c70a102c3}
      \field{sortinit}{F}
      \field{sortinithash}{2638baaa20439f1b5a8f80c6c08a13b4}
      \field{labelnamesource}{author}
      \field{labeltitlesource}{title}
      \field{journaltitle}{Finance and Stochastics}
      \field{number}{4}
      \field{title}{Polynomial diffusions and applications in finance}
      \field{volume}{20}
      \field{year}{2016}
      \field{pages}{931\bibrangedash 972}
      \range{pages}{42}
    \endentry
    \entry{follmerschied02}{article}{}{}
      \name{author}{2}{}{%
        {{hash=0c7afbc5410054170dbed41d05baca56}{%
           family={Föllmer},
           familyi={F\bibinitperiod},
           given={Hans},
           giveni={H\bibinitperiod}}}%
        {{hash=1f6c583833ee16c8fddcf5319a341e74}{%
           family={Schied},
           familyi={S\bibinitperiod},
           given={Alexander},
           giveni={A\bibinitperiod}}}%
      }
      \strng{namehash}{81ddf1f85ae72234c089ade593394068}
      \strng{fullhash}{81ddf1f85ae72234c089ade593394068}
      \strng{fullhashraw}{81ddf1f85ae72234c089ade593394068}
      \strng{bibnamehash}{81ddf1f85ae72234c089ade593394068}
      \strng{authorbibnamehash}{81ddf1f85ae72234c089ade593394068}
      \strng{authornamehash}{81ddf1f85ae72234c089ade593394068}
      \strng{authorfullhash}{81ddf1f85ae72234c089ade593394068}
      \strng{authorfullhashraw}{81ddf1f85ae72234c089ade593394068}
      \field{sortinit}{F}
      \field{sortinithash}{2638baaa20439f1b5a8f80c6c08a13b4}
      \field{labelnamesource}{author}
      \field{labeltitlesource}{title}
      \field{journaltitle}{Finance and Stochastics}
      \field{number}{4}
      \field{title}{Convex Measures of Risk and Trading Constraints}
      \field{volume}{6}
      \field{year}{2002}
      \field{pages}{429\bibrangedash 447}
      \range{pages}{19}
      \verb{doi}
      \verb 10.1007/s007800200072
      \endverb
    \endentry
    \entry{fong15}{article}{}{}
      \name{author}{1}{}{%
        {{hash=73645f147c1a5ec6f127cac2c822e9c6}{%
           family={Fong},
           familyi={F\bibinitperiod},
           given={Brendan},
           giveni={B\bibinitperiod}}}%
      }
      \strng{namehash}{73645f147c1a5ec6f127cac2c822e9c6}
      \strng{fullhash}{73645f147c1a5ec6f127cac2c822e9c6}
      \strng{fullhashraw}{73645f147c1a5ec6f127cac2c822e9c6}
      \strng{bibnamehash}{73645f147c1a5ec6f127cac2c822e9c6}
      \strng{authorbibnamehash}{73645f147c1a5ec6f127cac2c822e9c6}
      \strng{authornamehash}{73645f147c1a5ec6f127cac2c822e9c6}
      \strng{authorfullhash}{73645f147c1a5ec6f127cac2c822e9c6}
      \strng{authorfullhashraw}{73645f147c1a5ec6f127cac2c822e9c6}
      \field{sortinit}{F}
      \field{sortinithash}{2638baaa20439f1b5a8f80c6c08a13b4}
      \field{labelnamesource}{author}
      \field{labeltitlesource}{title}
      \field{eprinttype}{arXiv}
      \field{journaltitle}{Theory and Applications of Categories}
      \field{number}{33}
      \field{title}{Decorated Cospans}
      \field{volume}{30}
      \field{year}{2015}
      \field{pages}{1096\bibrangedash 1120}
      \range{pages}{25}
      \verb{eprint}
      \verb 1502.00872
      \endverb
    \endentry
    \entry{foster1953}{article}{}{}
      \name{author}{1}{}{%
        {{hash=f7811a89af03115885ca94acf1422776}{%
           family={Foster},
           familyi={F\bibinitperiod},
           given={F.\bibnamedelimi G.},
           giveni={F\bibinitperiod\bibinitdelim G\bibinitperiod}}}%
      }
      \strng{namehash}{f7811a89af03115885ca94acf1422776}
      \strng{fullhash}{f7811a89af03115885ca94acf1422776}
      \strng{fullhashraw}{f7811a89af03115885ca94acf1422776}
      \strng{bibnamehash}{f7811a89af03115885ca94acf1422776}
      \strng{authorbibnamehash}{f7811a89af03115885ca94acf1422776}
      \strng{authornamehash}{f7811a89af03115885ca94acf1422776}
      \strng{authorfullhash}{f7811a89af03115885ca94acf1422776}
      \strng{authorfullhashraw}{f7811a89af03115885ca94acf1422776}
      \field{sortinit}{F}
      \field{sortinithash}{2638baaa20439f1b5a8f80c6c08a13b4}
      \field{labelnamesource}{author}
      \field{labeltitlesource}{title}
      \field{journaltitle}{Annals of Mathematical Statistics}
      \field{note}{(The additive-drift criterion.)}
      \field{number}{3}
      \field{title}{On the stochastic matrices associated with certain queuing processes}
      \field{volume}{24}
      \field{year}{1953}
      \field{pages}{355\bibrangedash 360}
      \range{pages}{6}
    \endentry
    \entry{fowler95}{report}{}{}
      \name{author}{2}{}{%
        {{hash=63aee225bdf604470226ddefe82fe6ce}{%
           family={Fowler},
           familyi={F\bibinitperiod},
           given={John\bibnamedelima W.},
           giveni={J\bibinitperiod\bibinitdelim W\bibinitperiod}}}%
        {{hash=96483fa2853acf318f4fe064589b7561}{%
           family={Robinson},
           familyi={R\bibinitperiod},
           given={Jennifer},
           giveni={J\bibinitperiod}}}%
      }
      \list{institution}{1}{%
        {SEMATECH}%
      }
      \list{location}{1}{%
        {Austin, TX}%
      }
      \strng{namehash}{4f486722fbd6fa0040c5b9933f05210b}
      \strng{fullhash}{4f486722fbd6fa0040c5b9933f05210b}
      \strng{fullhashraw}{4f486722fbd6fa0040c5b9933f05210b}
      \strng{bibnamehash}{4f486722fbd6fa0040c5b9933f05210b}
      \strng{authorbibnamehash}{4f486722fbd6fa0040c5b9933f05210b}
      \strng{authornamehash}{4f486722fbd6fa0040c5b9933f05210b}
      \strng{authorfullhash}{4f486722fbd6fa0040c5b9933f05210b}
      \strng{authorfullhashraw}{4f486722fbd6fa0040c5b9933f05210b}
      \field{sortinit}{F}
      \field{sortinithash}{2638baaa20439f1b5a8f80c6c08a13b4}
      \field{labelnamesource}{author}
      \field{labeltitlesource}{title}
      \field{number}{95062861A-XFR}
      \field{subtitle}{Final Report}
      \field{title}{Measurement and Improvement of Manufacturing Capacity ({MIMAC})}
      \field{type}{Technology Transfer Report}
      \field{year}{1995}
    \endentry
    \entry{fremlin4}{book}{}{}
      \name{author}{1}{}{%
        {{hash=83364f4ea6bd5f698cb83aefe476a012}{%
           family={Fremlin},
           familyi={F\bibinitperiod},
           given={David\bibnamedelima H.},
           giveni={D\bibinitperiod\bibinitdelim H\bibinitperiod}}}%
      }
      \list{location}{1}{%
        {Colchester}%
      }
      \list{publisher}{1}{%
        {Torres Fremlin}%
      }
      \strng{namehash}{83364f4ea6bd5f698cb83aefe476a012}
      \strng{fullhash}{83364f4ea6bd5f698cb83aefe476a012}
      \strng{fullhashraw}{83364f4ea6bd5f698cb83aefe476a012}
      \strng{bibnamehash}{83364f4ea6bd5f698cb83aefe476a012}
      \strng{authorbibnamehash}{83364f4ea6bd5f698cb83aefe476a012}
      \strng{authornamehash}{83364f4ea6bd5f698cb83aefe476a012}
      \strng{authorfullhash}{83364f4ea6bd5f698cb83aefe476a012}
      \strng{authorfullhashraw}{83364f4ea6bd5f698cb83aefe476a012}
      \field{sortinit}{F}
      \field{sortinithash}{2638baaa20439f1b5a8f80c6c08a13b4}
      \field{labelnamesource}{author}
      \field{labeltitlesource}{title}
      \field{title}{Measure Theory. Volume 4: Topological Measure Spaces}
      \field{year}{2006}
    \endentry
    \entry{fritzliang22}{online}{}{}
      \name{author}{2}{}{%
        {{hash=485f579f829f8629fec8c85c28a33f28}{%
           family={Fritz},
           familyi={F\bibinitperiod},
           given={Tobias},
           giveni={T\bibinitperiod}}}%
        {{hash=aeca719eb18b6fd797160ac8c3ce72b2}{%
           family={Liang},
           familyi={L\bibinitperiod},
           given={Wendong},
           giveni={W\bibinitperiod}}}%
      }
      \strng{namehash}{924c92645fc592ee31ddc826eef80a20}
      \strng{fullhash}{924c92645fc592ee31ddc826eef80a20}
      \strng{fullhashraw}{924c92645fc592ee31ddc826eef80a20}
      \strng{bibnamehash}{924c92645fc592ee31ddc826eef80a20}
      \strng{authorbibnamehash}{924c92645fc592ee31ddc826eef80a20}
      \strng{authornamehash}{924c92645fc592ee31ddc826eef80a20}
      \strng{authorfullhash}{924c92645fc592ee31ddc826eef80a20}
      \strng{authorfullhashraw}{924c92645fc592ee31ddc826eef80a20}
      \field{sortinit}{F}
      \field{sortinithash}{2638baaa20439f1b5a8f80c6c08a13b4}
      \field{labelnamesource}{author}
      \field{labeltitlesource}{title}
      \field{eprinttype}{arXiv}
      \field{title}{Free gs-monoidal categories and free Markov categories}
      \field{year}{2022}
      \verb{doi}
      \verb 10.1007/s10485-023-09717-0
      \endverb
      \verb{eprint}
      \verb 2204.02284
      \endverb
    \endentry
    \entry{garloff1986}{incollection}{}{}
      \name{author}{1}{}{%
        {{hash=c03207559e3fccf90187de5a906077a8}{%
           family={Garloff},
           familyi={G\bibinitperiod},
           given={J.},
           giveni={J\bibinitperiod}}}%
      }
      \name{editor}{1}{}{%
        {{hash=0ab12ee89ee042c4dcb7ba03d1f77e6b}{%
           family={Nickel},
           familyi={N\bibinitperiod},
           given={K.},
           giveni={K\bibinitperiod}}}%
      }
      \list{publisher}{1}{%
        {Springer}%
      }
      \strng{namehash}{c03207559e3fccf90187de5a906077a8}
      \strng{fullhash}{c03207559e3fccf90187de5a906077a8}
      \strng{fullhashraw}{c03207559e3fccf90187de5a906077a8}
      \strng{bibnamehash}{c03207559e3fccf90187de5a906077a8}
      \strng{authorbibnamehash}{c03207559e3fccf90187de5a906077a8}
      \strng{authornamehash}{c03207559e3fccf90187de5a906077a8}
      \strng{authorfullhash}{c03207559e3fccf90187de5a906077a8}
      \strng{authorfullhashraw}{c03207559e3fccf90187de5a906077a8}
      \strng{editorbibnamehash}{0ab12ee89ee042c4dcb7ba03d1f77e6b}
      \strng{editornamehash}{0ab12ee89ee042c4dcb7ba03d1f77e6b}
      \strng{editorfullhash}{0ab12ee89ee042c4dcb7ba03d1f77e6b}
      \strng{editorfullhashraw}{0ab12ee89ee042c4dcb7ba03d1f77e6b}
      \field{sortinit}{G}
      \field{sortinithash}{32d67eca0634bf53703493fb1090a2e8}
      \field{labelnamesource}{author}
      \field{labeltitlesource}{title}
      \field{booktitle}{Interval Mathematics 1985}
      \field{note}{(Subdivision sharpens Bernstein enclosures --- the completeness partner of the enclosure soundness used here.)}
      \field{series}{Lecture Notes in Computer Science}
      \field{title}{Convergent bounds for the range of multivariate polynomials}
      \field{volume}{212}
      \field{year}{1986}
      \field{pages}{37\bibrangedash 56}
      \range{pages}{20}
    \endentry
    \entry{gellmann92}{online}{}{}
      \name{author}{2}{}{%
        {{hash=08ade17a88f1c8420c2a5e5f980bed7e}{%
           family={Gell-Mann},
           familyi={G\bibinithyphendelim M\bibinitperiod},
           given={Murray},
           giveni={M\bibinitperiod}}}%
        {{hash=1ac7eee9cbb87f927beacb5616f7f48a}{%
           family={Hartle},
           familyi={H\bibinitperiod},
           given={James\bibnamedelima B.},
           giveni={J\bibinitperiod\bibinitdelim B\bibinitperiod}}}%
      }
      \strng{namehash}{a340427943a4c3d945392a0b2c4278c8}
      \strng{fullhash}{a340427943a4c3d945392a0b2c4278c8}
      \strng{fullhashraw}{a340427943a4c3d945392a0b2c4278c8}
      \strng{bibnamehash}{a340427943a4c3d945392a0b2c4278c8}
      \strng{authorbibnamehash}{a340427943a4c3d945392a0b2c4278c8}
      \strng{authornamehash}{a340427943a4c3d945392a0b2c4278c8}
      \strng{authorfullhash}{a340427943a4c3d945392a0b2c4278c8}
      \strng{authorfullhashraw}{a340427943a4c3d945392a0b2c4278c8}
      \field{sortinit}{G}
      \field{sortinithash}{32d67eca0634bf53703493fb1090a2e8}
      \field{labelnamesource}{author}
      \field{labeltitlesource}{title}
      \field{eprinttype}{arXiv}
      \field{title}{Classical Equations for Quantum Systems}
      \field{year}{1992}
      \verb{doi}
      \verb 10.1103/PhysRevD.47.3345
      \endverb
      \verb{eprint}
      \verb gr-qc/9210010
      \endverb
    \endentry
    \entry{gentzen1935}{article}{}{}
      \name{author}{1}{}{%
        {{hash=df49bcdf1e485294ff125966547769b5}{%
           family={Gentzen},
           familyi={G\bibinitperiod},
           given={G.},
           giveni={G\bibinitperiod}}}%
      }
      \strng{namehash}{df49bcdf1e485294ff125966547769b5}
      \strng{fullhash}{df49bcdf1e485294ff125966547769b5}
      \strng{fullhashraw}{df49bcdf1e485294ff125966547769b5}
      \strng{bibnamehash}{df49bcdf1e485294ff125966547769b5}
      \strng{authorbibnamehash}{df49bcdf1e485294ff125966547769b5}
      \strng{authornamehash}{df49bcdf1e485294ff125966547769b5}
      \strng{authorfullhash}{df49bcdf1e485294ff125966547769b5}
      \strng{authorfullhashraw}{df49bcdf1e485294ff125966547769b5}
      \field{sortinit}{G}
      \field{sortinithash}{32d67eca0634bf53703493fb1090a2e8}
      \field{labelnamesource}{author}
      \field{labeltitlesource}{title}
      \field{journaltitle}{Mathematische Zeitschrift}
      \field{note}{English translation in M.~E.~Szabo (ed.), \emph{The Collected Papers of Gerhard Gentzen}, North-Holland, 1969.}
      \field{title}{Untersuchungen über das logische Schlie{ß}en I, II}
      \field{volume}{39}
      \field{year}{1935}
      \field{pages}{176\bibrangedash 210\bibrangessep 405\bibrangedash 431}
      \range{pages}{62}
    \endentry
    \entry{gershgorin1931}{article}{}{}
      \name{author}{1}{}{%
        {{hash=2e92cb809b252f5f42724d614170a54f}{%
           family={Gerschgorin},
           familyi={G\bibinitperiod},
           given={S.},
           giveni={S\bibinitperiod}}}%
      }
      \strng{namehash}{2e92cb809b252f5f42724d614170a54f}
      \strng{fullhash}{2e92cb809b252f5f42724d614170a54f}
      \strng{fullhashraw}{2e92cb809b252f5f42724d614170a54f}
      \strng{bibnamehash}{2e92cb809b252f5f42724d614170a54f}
      \strng{authorbibnamehash}{2e92cb809b252f5f42724d614170a54f}
      \strng{authornamehash}{2e92cb809b252f5f42724d614170a54f}
      \strng{authorfullhash}{2e92cb809b252f5f42724d614170a54f}
      \strng{authorfullhashraw}{2e92cb809b252f5f42724d614170a54f}
      \field{sortinit}{G}
      \field{sortinithash}{32d67eca0634bf53703493fb1090a2e8}
      \field{labelnamesource}{author}
      \field{labeltitlesource}{title}
      \field{journaltitle}{Bulletin de l'Académie des Sciences de l'URSS, Classe des sciences mathématiques et naturelles}
      \field{title}{Über die Abgrenzung der Eigenwerte einer Matrix}
      \field{volume}{6}
      \field{year}{1931}
      \field{pages}{749\bibrangedash 754}
      \range{pages}{6}
    \endentry
    \entry{ghica19}{article}{}{}
      \name{author}{3}{}{%
        {{hash=935665c7060cc483b5b7a6a61817881c}{%
           family={Ghica},
           familyi={G\bibinitperiod},
           given={Dan\bibnamedelima R.},
           giveni={D\bibinitperiod\bibinitdelim R\bibinitperiod}}}%
        {{hash=4468ac0afcef5b066557fb8de7d26d60}{%
           family={Muroya},
           familyi={M\bibinitperiod},
           given={Koko},
           giveni={K\bibinitperiod}}}%
        {{hash=62c9c49406aaf36489e0db13cb6c4096}{%
           family={Waugh\bibnamedelima Ambridge},
           familyi={W\bibinitperiod\bibinitdelim A\bibinitperiod},
           given={Todd},
           giveni={T\bibinitperiod}}}%
      }
      \strng{namehash}{261e383c10965b9a818551ccabacc4ae}
      \strng{fullhash}{261e383c10965b9a818551ccabacc4ae}
      \strng{fullhashraw}{261e383c10965b9a818551ccabacc4ae}
      \strng{bibnamehash}{261e383c10965b9a818551ccabacc4ae}
      \strng{authorbibnamehash}{261e383c10965b9a818551ccabacc4ae}
      \strng{authornamehash}{261e383c10965b9a818551ccabacc4ae}
      \strng{authorfullhash}{261e383c10965b9a818551ccabacc4ae}
      \strng{authorfullhashraw}{261e383c10965b9a818551ccabacc4ae}
      \field{sortinit}{G}
      \field{sortinithash}{32d67eca0634bf53703493fb1090a2e8}
      \field{labelnamesource}{author}
      \field{labeltitlesource}{title}
      \field{eprinttype}{arXiv}
      \field{journaltitle}{Logical Methods in Computer Science}
      \field{number}{2}
      \field{title}{A Robust Graph-Based Approach to Observational Equivalence}
      \field{volume}{21}
      \field{year}{2025}
      \verb{doi}
      \verb 10.46298/lmcs-21(2:8)2025
      \endverb
      \verb{eprint}
      \verb 1907.01257
      \endverb
    \endentry
    \entry{neuralmc-2025-invariant}{online}{}{}
      \name{author}{4}{}{%
        {{hash=e227971e0a9f956b3d8f8688ae07fadb}{%
           family={Giacobbe},
           familyi={G\bibinitperiod},
           given={Mirco},
           giveni={M\bibinitperiod}}}%
        {{hash=b4aa934119a9d2ed6e4c21b1a7f20424}{%
           family={Kroening},
           familyi={K\bibinitperiod},
           given={Daniel},
           giveni={D\bibinitperiod}}}%
        {{hash=28c1b6e3493873dd7b092a3a250633eb}{%
           family={Pal},
           familyi={P\bibinitperiod},
           given={Abhinandan},
           giveni={A\bibinitperiod}}}%
        {{hash=be0cf212fa97eb5a6ca43eefedfc56a5}{%
           family={Tautschnig},
           familyi={T\bibinitperiod},
           given={Michael},
           giveni={M\bibinitperiod}}}%
      }
      \strng{namehash}{544fe8973ec4167451332533da9759b9}
      \strng{fullhash}{544fe8973ec4167451332533da9759b9}
      \strng{fullhashraw}{544fe8973ec4167451332533da9759b9}
      \strng{bibnamehash}{544fe8973ec4167451332533da9759b9}
      \strng{authorbibnamehash}{544fe8973ec4167451332533da9759b9}
      \strng{authornamehash}{544fe8973ec4167451332533da9759b9}
      \strng{authorfullhash}{544fe8973ec4167451332533da9759b9}
      \strng{authorfullhashraw}{544fe8973ec4167451332533da9759b9}
      \field{extraname}{1}
      \field{sortinit}{G}
      \field{sortinithash}{32d67eca0634bf53703493fb1090a2e8}
      \field{labelnamesource}{author}
      \field{labeltitlesource}{title}
      \field{note}{Advances in Neural Information Processing Systems 38 (NeurIPS 2025), pp.\ 74713--74740}
      \field{title}{Let a Neural Network Be Your Invariant}
      \field{year}{2025}
      \verb{urlraw}
      \verb https://pure-oai.bham.ac.uk/ws/portalfiles/portal/281309250/Let_a_Neural_Network_Be_Your_Invariant.pdf
      \endverb
      \verb{url}
      \verb https://pure-oai.bham.ac.uk/ws/portalfiles/portal/281309250/Let_a_Neural_Network_Be_Your_Invariant.pdf
      \endverb
    \endentry
    \entry{giacobbe24}{online}{}{}
      \name{author}{4}{}{%
        {{hash=e227971e0a9f956b3d8f8688ae07fadb}{%
           family={Giacobbe},
           familyi={G\bibinitperiod},
           given={Mirco},
           giveni={M\bibinitperiod}}}%
        {{hash=b4aa934119a9d2ed6e4c21b1a7f20424}{%
           family={Kroening},
           familyi={K\bibinitperiod},
           given={Daniel},
           giveni={D\bibinitperiod}}}%
        {{hash=28c1b6e3493873dd7b092a3a250633eb}{%
           family={Pal},
           familyi={P\bibinitperiod},
           given={Abhinandan},
           giveni={A\bibinitperiod}}}%
        {{hash=be0cf212fa97eb5a6ca43eefedfc56a5}{%
           family={Tautschnig},
           familyi={T\bibinitperiod},
           given={Michael},
           giveni={M\bibinitperiod}}}%
      }
      \strng{namehash}{544fe8973ec4167451332533da9759b9}
      \strng{fullhash}{544fe8973ec4167451332533da9759b9}
      \strng{fullhashraw}{544fe8973ec4167451332533da9759b9}
      \strng{bibnamehash}{544fe8973ec4167451332533da9759b9}
      \strng{authorbibnamehash}{544fe8973ec4167451332533da9759b9}
      \strng{authornamehash}{544fe8973ec4167451332533da9759b9}
      \strng{authorfullhash}{544fe8973ec4167451332533da9759b9}
      \strng{authorfullhashraw}{544fe8973ec4167451332533da9759b9}
      \field{extraname}{2}
      \field{sortinit}{G}
      \field{sortinithash}{32d67eca0634bf53703493fb1090a2e8}
      \field{labelnamesource}{author}
      \field{labeltitlesource}{title}
      \field{eprinttype}{arXiv}
      \field{title}{Neural Model Checking}
      \field{year}{2024}
      \verb{eprint}
      \verb 2410.23790
      \endverb
    \endentry
    \entry{gilboa93}{article}{}{}
      \name{author}{2}{}{%
        {{hash=2cdcd3e31354c4e936987b8f22726de6}{%
           family={Gilboa},
           familyi={G\bibinitperiod},
           given={Itzhak},
           giveni={I\bibinitperiod}}}%
        {{hash=e54561df4355c4959916e8a655f49c3c}{%
           family={Schmeidler},
           familyi={S\bibinitperiod},
           given={David},
           giveni={D\bibinitperiod}}}%
      }
      \strng{namehash}{2c02a6baf1aac0beea9eaf7a1f1cdcff}
      \strng{fullhash}{2c02a6baf1aac0beea9eaf7a1f1cdcff}
      \strng{fullhashraw}{2c02a6baf1aac0beea9eaf7a1f1cdcff}
      \strng{bibnamehash}{2c02a6baf1aac0beea9eaf7a1f1cdcff}
      \strng{authorbibnamehash}{2c02a6baf1aac0beea9eaf7a1f1cdcff}
      \strng{authornamehash}{2c02a6baf1aac0beea9eaf7a1f1cdcff}
      \strng{authorfullhash}{2c02a6baf1aac0beea9eaf7a1f1cdcff}
      \strng{authorfullhashraw}{2c02a6baf1aac0beea9eaf7a1f1cdcff}
      \field{sortinit}{G}
      \field{sortinithash}{32d67eca0634bf53703493fb1090a2e8}
      \field{labelnamesource}{author}
      \field{labeltitlesource}{title}
      \field{journaltitle}{Journal of Economic Theory}
      \field{number}{1}
      \field{title}{Updating Ambiguous Beliefs}
      \field{volume}{59}
      \field{year}{1993}
      \field{pages}{33\bibrangedash 49}
      \range{pages}{17}
      \verb{doi}
      \verb 10.1006/jeth.1993.1003
      \endverb
    \endentry
    \entry{girardpappas2007}{article}{}{}
      \name{author}{2}{}{%
        {{hash=d2d25c21d283fb6f333d3bc5e0d6bea9}{%
           family={Girard},
           familyi={G\bibinitperiod},
           given={A.},
           giveni={A\bibinitperiod}}}%
        {{hash=fd39f64a316253aa07afe174ea8b90fd}{%
           family={Pappas},
           familyi={P\bibinitperiod},
           given={G.\bibnamedelimi J.},
           giveni={G\bibinitperiod\bibinitdelim J\bibinitperiod}}}%
      }
      \strng{namehash}{39e695c2b4efa0854f04f2c3d829ac61}
      \strng{fullhash}{39e695c2b4efa0854f04f2c3d829ac61}
      \strng{fullhashraw}{39e695c2b4efa0854f04f2c3d829ac61}
      \strng{bibnamehash}{39e695c2b4efa0854f04f2c3d829ac61}
      \strng{authorbibnamehash}{39e695c2b4efa0854f04f2c3d829ac61}
      \strng{authornamehash}{39e695c2b4efa0854f04f2c3d829ac61}
      \strng{authorfullhash}{39e695c2b4efa0854f04f2c3d829ac61}
      \strng{authorfullhashraw}{39e695c2b4efa0854f04f2c3d829ac61}
      \field{sortinit}{G}
      \field{sortinithash}{32d67eca0634bf53703493fb1090a2e8}
      \field{labelnamesource}{author}
      \field{labeltitlesource}{title}
      \field{journaltitle}{IEEE Transactions on Automatic Control}
      \field{note}{(Approximate simulation with quantitative degradation --- the analytic ancestor of the $\varepsilon$-witness transport cell.)}
      \field{number}{5}
      \field{title}{Approximation metrics for discrete and continuous systems}
      \field{volume}{52}
      \field{year}{2007}
      \field{pages}{782\bibrangedash 798}
      \range{pages}{17}
    \endentry
    \entry{girard1987}{article}{}{}
      \name{author}{1}{}{%
        {{hash=a5a9f2c1255a817f6a5f5751b5bcfc5f}{%
           family={Girard},
           familyi={G\bibinitperiod},
           given={J.-Y.},
           giveni={J\bibinithyphendelim Y\bibinitperiod}}}%
      }
      \strng{namehash}{a5a9f2c1255a817f6a5f5751b5bcfc5f}
      \strng{fullhash}{a5a9f2c1255a817f6a5f5751b5bcfc5f}
      \strng{fullhashraw}{a5a9f2c1255a817f6a5f5751b5bcfc5f}
      \strng{bibnamehash}{a5a9f2c1255a817f6a5f5751b5bcfc5f}
      \strng{authorbibnamehash}{a5a9f2c1255a817f6a5f5751b5bcfc5f}
      \strng{authornamehash}{a5a9f2c1255a817f6a5f5751b5bcfc5f}
      \strng{authorfullhash}{a5a9f2c1255a817f6a5f5751b5bcfc5f}
      \strng{authorfullhashraw}{a5a9f2c1255a817f6a5f5751b5bcfc5f}
      \field{sortinit}{G}
      \field{sortinithash}{32d67eca0634bf53703493fb1090a2e8}
      \field{labelnamesource}{author}
      \field{labeltitlesource}{title}
      \field{journaltitle}{Theoretical Computer Science}
      \field{note}{(Source of the additive/multiplicative vocabulary for context-sharing vs.\ context-splitting rule formats.)}
      \field{number}{1}
      \field{title}{Linear logic}
      \field{volume}{50}
      \field{year}{1987}
      \field{pages}{1\bibrangedash 101}
      \range{pages}{101}
    \endentry
    \entry{giry82}{incollection}{}{}
      \name{author}{1}{}{%
        {{hash=4b2ee17af47c1a78a537cfb0b6f956af}{%
           family={Giry},
           familyi={G\bibinitperiod},
           given={Michèle},
           giveni={M\bibinitperiod}}}%
      }
      \list{publisher}{1}{%
        {Springer}%
      }
      \strng{namehash}{4b2ee17af47c1a78a537cfb0b6f956af}
      \strng{fullhash}{4b2ee17af47c1a78a537cfb0b6f956af}
      \strng{fullhashraw}{4b2ee17af47c1a78a537cfb0b6f956af}
      \strng{bibnamehash}{4b2ee17af47c1a78a537cfb0b6f956af}
      \strng{authorbibnamehash}{4b2ee17af47c1a78a537cfb0b6f956af}
      \strng{authornamehash}{4b2ee17af47c1a78a537cfb0b6f956af}
      \strng{authorfullhash}{4b2ee17af47c1a78a537cfb0b6f956af}
      \strng{authorfullhashraw}{4b2ee17af47c1a78a537cfb0b6f956af}
      \field{sortinit}{G}
      \field{sortinithash}{32d67eca0634bf53703493fb1090a2e8}
      \field{labelnamesource}{author}
      \field{labeltitlesource}{title}
      \field{booktitle}{Categorical aspects of topology and analysis}
      \field{series}{Lecture Notes in Mathematics}
      \field{title}{A Categorical Approach to Probability Theory}
      \field{volume}{915}
      \field{year}{1982}
      \verb{doi}
      \verb 10.1007/bfb0092872
      \endverb
    \endentry
    \entry{goubaultlarrecq07}{inproceedings}{}{}
      \name{author}{1}{}{%
        {{hash=80ae10896704521f62f88fbee1b2f3e4}{%
           family={Goubault-Larrecq},
           familyi={G\bibinithyphendelim L\bibinitperiod},
           given={Jean},
           giveni={J\bibinitperiod}}}%
      }
      \list{publisher}{1}{%
        {Springer}%
      }
      \strng{namehash}{80ae10896704521f62f88fbee1b2f3e4}
      \strng{fullhash}{80ae10896704521f62f88fbee1b2f3e4}
      \strng{fullhashraw}{80ae10896704521f62f88fbee1b2f3e4}
      \strng{bibnamehash}{80ae10896704521f62f88fbee1b2f3e4}
      \strng{authorbibnamehash}{80ae10896704521f62f88fbee1b2f3e4}
      \strng{authornamehash}{80ae10896704521f62f88fbee1b2f3e4}
      \strng{authorfullhash}{80ae10896704521f62f88fbee1b2f3e4}
      \strng{authorfullhashraw}{80ae10896704521f62f88fbee1b2f3e4}
      \field{extraname}{1}
      \field{sortinit}{G}
      \field{sortinithash}{32d67eca0634bf53703493fb1090a2e8}
      \field{labelnamesource}{author}
      \field{labeltitlesource}{title}
      \field{booktitle}{Computer Science Logic ({CSL} 2007)}
      \field{series}{Lecture Notes in Computer Science}
      \field{title}{Continuous Previsions}
      \field{volume}{4646}
      \field{year}{2007}
      \field{pages}{542\bibrangedash 557}
      \range{pages}{16}
      \verb{doi}
      \verb 10.1007/978-3-540-74915-8_40
      \endverb
    \endentry
    \entry{goubaultlarrecq26}{online}{}{}
      \name{author}{1}{}{%
        {{hash=80ae10896704521f62f88fbee1b2f3e4}{%
           family={Goubault-Larrecq},
           familyi={G\bibinithyphendelim L\bibinitperiod},
           given={Jean},
           giveni={J\bibinitperiod}}}%
      }
      \strng{namehash}{80ae10896704521f62f88fbee1b2f3e4}
      \strng{fullhash}{80ae10896704521f62f88fbee1b2f3e4}
      \strng{fullhashraw}{80ae10896704521f62f88fbee1b2f3e4}
      \strng{bibnamehash}{80ae10896704521f62f88fbee1b2f3e4}
      \strng{authorbibnamehash}{80ae10896704521f62f88fbee1b2f3e4}
      \strng{authornamehash}{80ae10896704521f62f88fbee1b2f3e4}
      \strng{authorfullhash}{80ae10896704521f62f88fbee1b2f3e4}
      \strng{authorfullhashraw}{80ae10896704521f62f88fbee1b2f3e4}
      \field{extraname}{2}
      \field{sortinit}{G}
      \field{sortinithash}{32d67eca0634bf53703493fb1090a2e8}
      \field{labelnamesource}{author}
      \field{labeltitlesource}{title}
      \field{eprinttype}{arXiv}
      \field{title}{Just Previsions}
      \field{year}{2026}
      \verb{eprint}
      \verb 2605.10173
      \endverb
    \endentry
    \entry{goubaultlarrecq08}{inproceedings}{}{}
      \name{author}{1}{}{%
        {{hash=80ae10896704521f62f88fbee1b2f3e4}{%
           family={Goubault-Larrecq},
           familyi={G\bibinithyphendelim L\bibinitperiod},
           given={Jean},
           giveni={J\bibinitperiod}}}%
      }
      \list{publisher}{1}{%
        {Springer}%
      }
      \strng{namehash}{80ae10896704521f62f88fbee1b2f3e4}
      \strng{fullhash}{80ae10896704521f62f88fbee1b2f3e4}
      \strng{fullhashraw}{80ae10896704521f62f88fbee1b2f3e4}
      \strng{bibnamehash}{80ae10896704521f62f88fbee1b2f3e4}
      \strng{authorbibnamehash}{80ae10896704521f62f88fbee1b2f3e4}
      \strng{authornamehash}{80ae10896704521f62f88fbee1b2f3e4}
      \strng{authorfullhash}{80ae10896704521f62f88fbee1b2f3e4}
      \strng{authorfullhashraw}{80ae10896704521f62f88fbee1b2f3e4}
      \field{extraname}{3}
      \field{sortinit}{G}
      \field{sortinithash}{32d67eca0634bf53703493fb1090a2e8}
      \field{labelnamesource}{author}
      \field{labeltitlesource}{title}
      \field{booktitle}{Foundations of Software Science and Computational Structures ({FoSSaCS} 2008)}
      \field{series}{Lecture Notes in Computer Science}
      \field{title}{Prevision Domains and Convex Powercones}
      \field{volume}{4962}
      \field{year}{2008}
      \field{pages}{318\bibrangedash 333}
      \range{pages}{16}
      \verb{doi}
      \verb 10.1007/978-3-540-78499-9_23
      \endverb
    \endentry
    \entry{granasdugundji2003}{book}{}{}
      \name{author}{2}{}{%
        {{hash=5c08a0e3ed39da44a30674792147a674}{%
           family={Granas},
           familyi={G\bibinitperiod},
           given={A.},
           giveni={A\bibinitperiod}}}%
        {{hash=5998b1dba35832b2e7e6f31e3dc6a62d}{%
           family={Dugundji},
           familyi={D\bibinitperiod},
           given={J.},
           giveni={J\bibinitperiod}}}%
      }
      \list{publisher}{1}{%
        {Springer}%
      }
      \strng{namehash}{aa48910f84b179cce62dd69607627f20}
      \strng{fullhash}{aa48910f84b179cce62dd69607627f20}
      \strng{fullhashraw}{aa48910f84b179cce62dd69607627f20}
      \strng{bibnamehash}{aa48910f84b179cce62dd69607627f20}
      \strng{authorbibnamehash}{aa48910f84b179cce62dd69607627f20}
      \strng{authornamehash}{aa48910f84b179cce62dd69607627f20}
      \strng{authorfullhash}{aa48910f84b179cce62dd69607627f20}
      \strng{authorfullhashraw}{aa48910f84b179cce62dd69607627f20}
      \field{sortinit}{G}
      \field{sortinithash}{32d67eca0634bf53703493fb1090a2e8}
      \field{labelnamesource}{author}
      \field{labeltitlesource}{title}
      \field{note}{(The closed-ball condition $d(Tx_0, x_0) \leq (1 - \rho)\,r$ and the a-priori bound $d(x_0, x^*) \leq \delta/(1-\rho)$: §I.1.)}
      \field{series}{Springer Monographs in Mathematics}
      \field{title}{Fixed Point Theory}
      \field{year}{2003}
    \endentry
    \entry{grandis04}{article}{}{}
      \name{author}{2}{}{%
        {{hash=f3ff8e8348a1d26e349b582af6718524}{%
           family={Grandis},
           familyi={G\bibinitperiod},
           given={Marco},
           giveni={M\bibinitperiod}}}%
        {{hash=3e73f997df2e2d33541feea1f4a7c236}{%
           family={Paré},
           familyi={P\bibinitperiod},
           given={Robert},
           giveni={R\bibinitperiod}}}%
      }
      \strng{namehash}{1adfe41b34eafa1fe7a68007f48f211f}
      \strng{fullhash}{1adfe41b34eafa1fe7a68007f48f211f}
      \strng{fullhashraw}{1adfe41b34eafa1fe7a68007f48f211f}
      \strng{bibnamehash}{1adfe41b34eafa1fe7a68007f48f211f}
      \strng{authorbibnamehash}{1adfe41b34eafa1fe7a68007f48f211f}
      \strng{authornamehash}{1adfe41b34eafa1fe7a68007f48f211f}
      \strng{authorfullhash}{1adfe41b34eafa1fe7a68007f48f211f}
      \strng{authorfullhashraw}{1adfe41b34eafa1fe7a68007f48f211f}
      \field{extraname}{1}
      \field{sortinit}{G}
      \field{sortinithash}{32d67eca0634bf53703493fb1090a2e8}
      \field{labelnamesource}{author}
      \field{labeltitlesource}{title}
      \field{journaltitle}{Cahiers de Topologie et Géométrie Différentielle Catégoriques}
      \field{number}{3}
      \field{title}{Adjoint for double categories}
      \field{volume}{45}
      \field{year}{2004}
      \field{pages}{193\bibrangedash 240}
      \range{pages}{48}
    \endentry
    \entry{grandis99}{article}{}{}
      \name{author}{2}{}{%
        {{hash=f3ff8e8348a1d26e349b582af6718524}{%
           family={Grandis},
           familyi={G\bibinitperiod},
           given={Marco},
           giveni={M\bibinitperiod}}}%
        {{hash=3e73f997df2e2d33541feea1f4a7c236}{%
           family={Paré},
           familyi={P\bibinitperiod},
           given={Robert},
           giveni={R\bibinitperiod}}}%
      }
      \strng{namehash}{1adfe41b34eafa1fe7a68007f48f211f}
      \strng{fullhash}{1adfe41b34eafa1fe7a68007f48f211f}
      \strng{fullhashraw}{1adfe41b34eafa1fe7a68007f48f211f}
      \strng{bibnamehash}{1adfe41b34eafa1fe7a68007f48f211f}
      \strng{authorbibnamehash}{1adfe41b34eafa1fe7a68007f48f211f}
      \strng{authornamehash}{1adfe41b34eafa1fe7a68007f48f211f}
      \strng{authorfullhash}{1adfe41b34eafa1fe7a68007f48f211f}
      \strng{authorfullhashraw}{1adfe41b34eafa1fe7a68007f48f211f}
      \field{extraname}{2}
      \field{sortinit}{G}
      \field{sortinithash}{32d67eca0634bf53703493fb1090a2e8}
      \field{labelnamesource}{author}
      \field{labeltitlesource}{title}
      \field{journaltitle}{Cahiers de Topologie et Géométrie Différentielle Catégoriques}
      \field{number}{3}
      \field{title}{Limits in double categories}
      \field{volume}{40}
      \field{year}{1999}
      \field{pages}{162\bibrangedash 220}
      \range{pages}{59}
    \endentry
    \entry{guck14}{online}{}{}
      \name{author}{5}{}{%
        {{hash=fc9e58817d92fecd9e7a8f6299945d9d}{%
           family={Guck},
           familyi={G\bibinitperiod},
           given={Dennis},
           giveni={D\bibinitperiod}}}%
        {{hash=943e2ea1641785b18e38bd8fba136864}{%
           family={Hatefi},
           familyi={H\bibinitperiod},
           given={Hassan},
           giveni={H\bibinitperiod}}}%
        {{hash=b396735e2150877fec4cee61d80d5430}{%
           family={Hermanns},
           familyi={H\bibinitperiod},
           given={Holger},
           giveni={H\bibinitperiod}}}%
        {{hash=392d5825b15ef7f81f68632c7a49fdef}{%
           family={Katoen},
           familyi={K\bibinitperiod},
           given={Joost-Pieter},
           giveni={J\bibinithyphendelim P\bibinitperiod}}}%
        {{hash=32206d16e2080efea2b8022fbda4b4b7}{%
           family={Timmer},
           familyi={T\bibinitperiod},
           given={Mark},
           giveni={M\bibinitperiod}}}%
      }
      \strng{namehash}{c1ab64c92f8ee0c8e5a145fe4aaf592d}
      \strng{fullhash}{c1ab64c92f8ee0c8e5a145fe4aaf592d}
      \strng{fullhashraw}{c1ab64c92f8ee0c8e5a145fe4aaf592d}
      \strng{bibnamehash}{c1ab64c92f8ee0c8e5a145fe4aaf592d}
      \strng{authorbibnamehash}{c1ab64c92f8ee0c8e5a145fe4aaf592d}
      \strng{authornamehash}{c1ab64c92f8ee0c8e5a145fe4aaf592d}
      \strng{authorfullhash}{c1ab64c92f8ee0c8e5a145fe4aaf592d}
      \strng{authorfullhashraw}{c1ab64c92f8ee0c8e5a145fe4aaf592d}
      \field{sortinit}{G}
      \field{sortinithash}{32d67eca0634bf53703493fb1090a2e8}
      \field{labelnamesource}{author}
      \field{labeltitlesource}{title}
      \field{eprinttype}{arXiv}
      \field{title}{Analysis of Timed and Long-Run Objectives for Markov Automata}
      \field{year}{2014}
      \verb{eprint}
      \verb 1407.7356
      \endverb
    \endentry
    \entry{haagkastler64}{article}{}{}
      \name{author}{2}{}{%
        {{hash=e4b61b3dfd8653e784dec22252d88f13}{%
           family={Haag},
           familyi={H\bibinitperiod},
           given={Rudolf},
           giveni={R\bibinitperiod}}}%
        {{hash=7680d40592ee73bc29468539a1f4ea08}{%
           family={Kastler},
           familyi={K\bibinitperiod},
           given={Daniel},
           giveni={D\bibinitperiod}}}%
      }
      \strng{namehash}{a82a47cef65c72e2c7011373250181eb}
      \strng{fullhash}{a82a47cef65c72e2c7011373250181eb}
      \strng{fullhashraw}{a82a47cef65c72e2c7011373250181eb}
      \strng{bibnamehash}{a82a47cef65c72e2c7011373250181eb}
      \strng{authorbibnamehash}{a82a47cef65c72e2c7011373250181eb}
      \strng{authornamehash}{a82a47cef65c72e2c7011373250181eb}
      \strng{authorfullhash}{a82a47cef65c72e2c7011373250181eb}
      \strng{authorfullhashraw}{a82a47cef65c72e2c7011373250181eb}
      \field{sortinit}{H}
      \field{sortinithash}{23a3aa7c24e56cfa16945d55545109b5}
      \field{labelnamesource}{author}
      \field{labeltitlesource}{title}
      \field{journaltitle}{Journal of Mathematical Physics}
      \field{number}{7}
      \field{title}{An Algebraic Approach to Quantum Field Theory}
      \field{volume}{5}
      \field{year}{1964}
      \field{pages}{848\bibrangedash 861}
      \range{pages}{14}
      \verb{doi}
      \verb 10.1063/1.1704187
      \endverb
    \endentry
    \entry{hadzihasanovic24a}{online}{}{}
      \name{author}{1}{}{%
        {{hash=48ba6a8c66711f3b59879e4e1149110c}{%
           family={Hadzihasanovic},
           familyi={H\bibinitperiod},
           given={Amar},
           giveni={A\bibinitperiod}}}%
      }
      \strng{namehash}{48ba6a8c66711f3b59879e4e1149110c}
      \strng{fullhash}{48ba6a8c66711f3b59879e4e1149110c}
      \strng{fullhashraw}{48ba6a8c66711f3b59879e4e1149110c}
      \strng{bibnamehash}{48ba6a8c66711f3b59879e4e1149110c}
      \strng{authorbibnamehash}{48ba6a8c66711f3b59879e4e1149110c}
      \strng{authornamehash}{48ba6a8c66711f3b59879e4e1149110c}
      \strng{authorfullhash}{48ba6a8c66711f3b59879e4e1149110c}
      \strng{authorfullhashraw}{48ba6a8c66711f3b59879e4e1149110c}
      \field{sortinit}{H}
      \field{sortinithash}{23a3aa7c24e56cfa16945d55545109b5}
      \field{labelnamesource}{author}
      \field{labeltitlesource}{title}
      \field{eprinttype}{arXiv}
      \field{title}{Combinatorics of higher-categorical diagrams}
      \field{year}{2024}
      \verb{eprint}
      \verb 2404.07273
      \endverb
    \endentry
    \entry{alifib}{online}{}{}
      \name{author}{2}{}{%
        {{hash=48ba6a8c66711f3b59879e4e1149110c}{%
           family={Hadzihasanovic},
           familyi={H\bibinitperiod},
           given={Amar},
           giveni={A\bibinitperiod}}}%
        {{hash=9f7605c282568d79a431d2dc14f75a5c}{%
           family={Kavvos},
           familyi={K\bibinitperiod},
           given={Alex},
           giveni={A\bibinitperiod}}}%
      }
      \strng{namehash}{e2a58aa0efa706eba6f85fa0b34222f0}
      \strng{fullhash}{e2a58aa0efa706eba6f85fa0b34222f0}
      \strng{fullhashraw}{e2a58aa0efa706eba6f85fa0b34222f0}
      \strng{bibnamehash}{e2a58aa0efa706eba6f85fa0b34222f0}
      \strng{authorbibnamehash}{e2a58aa0efa706eba6f85fa0b34222f0}
      \strng{authornamehash}{e2a58aa0efa706eba6f85fa0b34222f0}
      \strng{authorfullhash}{e2a58aa0efa706eba6f85fa0b34222f0}
      \strng{authorfullhashraw}{e2a58aa0efa706eba6f85fa0b34222f0}
      \field{sortinit}{H}
      \field{sortinithash}{23a3aa7c24e56cfa16945d55545109b5}
      \field{labelnamesource}{author}
      \field{labeltitlesource}{title}
      \field{title}{alifib: a programming language and proof assistant founded on higher-dimensional rewriting}
      \field{year}{2026}
      \verb{urlraw}
      \verb https://github.com/ahadziha/alifib
      \endverb
      \verb{url}
      \verb https://github.com/ahadziha/alifib
      \endverb
    \endentry
    \entry{hadzihasanovic23}{online}{}{}
      \name{author}{2}{}{%
        {{hash=48ba6a8c66711f3b59879e4e1149110c}{%
           family={Hadzihasanovic},
           familyi={H\bibinitperiod},
           given={Amar},
           giveni={A\bibinitperiod}}}%
        {{hash=dcb368ed1f25307eceeec3921f8b305c}{%
           family={Kessler},
           familyi={K\bibinitperiod},
           given={Diana},
           giveni={D\bibinitperiod}}}%
      }
      \strng{namehash}{bf3acb5c51cb5721345064bda540581a}
      \strng{fullhash}{bf3acb5c51cb5721345064bda540581a}
      \strng{fullhashraw}{bf3acb5c51cb5721345064bda540581a}
      \strng{bibnamehash}{bf3acb5c51cb5721345064bda540581a}
      \strng{authorbibnamehash}{bf3acb5c51cb5721345064bda540581a}
      \strng{authornamehash}{bf3acb5c51cb5721345064bda540581a}
      \strng{authorfullhash}{bf3acb5c51cb5721345064bda540581a}
      \strng{authorfullhashraw}{bf3acb5c51cb5721345064bda540581a}
      \field{sortinit}{H}
      \field{sortinithash}{23a3aa7c24e56cfa16945d55545109b5}
      \field{labelnamesource}{author}
      \field{labeltitlesource}{title}
      \field{eprinttype}{arXiv}
      \field{title}{Higher-dimensional subdiagram matching}
      \field{year}{2023}
      \verb{eprint}
      \verb 2304.09216
      \endverb
    \endentry
    \entry{hale-verduyn-lunel-1993}{book}{}{}
      \name{author}{2}{}{%
        {{hash=4e254cacb018b75cd6e5e99ab680b473}{%
           family={Hale},
           familyi={H\bibinitperiod},
           given={Jack\bibnamedelima K.},
           giveni={J\bibinitperiod\bibinitdelim K\bibinitperiod}}}%
        {{hash=f754d4e3beddc91f59a89bc8a2f72d86}{%
           family={Verduyn\bibnamedelima Lunel},
           familyi={V\bibinitperiod\bibinitdelim L\bibinitperiod},
           given={Sjoerd\bibnamedelima M.},
           giveni={S\bibinitperiod\bibinitdelim M\bibinitperiod}}}%
      }
      \list{location}{1}{%
        {New York}%
      }
      \list{publisher}{1}{%
        {Springer}%
      }
      \strng{namehash}{7f90db530e2ea382aa86951de4b85406}
      \strng{fullhash}{7f90db530e2ea382aa86951de4b85406}
      \strng{fullhashraw}{7f90db530e2ea382aa86951de4b85406}
      \strng{bibnamehash}{7f90db530e2ea382aa86951de4b85406}
      \strng{authorbibnamehash}{7f90db530e2ea382aa86951de4b85406}
      \strng{authornamehash}{7f90db530e2ea382aa86951de4b85406}
      \strng{authorfullhash}{7f90db530e2ea382aa86951de4b85406}
      \strng{authorfullhashraw}{7f90db530e2ea382aa86951de4b85406}
      \field{sortinit}{H}
      \field{sortinithash}{23a3aa7c24e56cfa16945d55545109b5}
      \field{labelnamesource}{author}
      \field{labeltitlesource}{title}
      \field{series}{Applied Mathematical Sciences}
      \field{title}{Introduction to Functional Differential Equations}
      \field{volume}{99}
      \field{year}{1993}
    \endentry
    \entry{halpern13}{online}{}{}
      \name{author}{2}{}{%
        {{hash=610a5aa55255d24fd3b5bfdd80bcc4e2}{%
           family={Halpern},
           familyi={H\bibinitperiod},
           given={Joseph\bibnamedelima Y.},
           giveni={J\bibinitperiod\bibinitdelim Y\bibinitperiod}}}%
        {{hash=26d0862842e97d80823e5ad49455be2c}{%
           family={Leung},
           familyi={L\bibinitperiod},
           given={Samantha},
           giveni={S\bibinitperiod}}}%
      }
      \strng{namehash}{0565db30e661602e72fe1cf16c49f629}
      \strng{fullhash}{0565db30e661602e72fe1cf16c49f629}
      \strng{fullhashraw}{0565db30e661602e72fe1cf16c49f629}
      \strng{bibnamehash}{0565db30e661602e72fe1cf16c49f629}
      \strng{authorbibnamehash}{0565db30e661602e72fe1cf16c49f629}
      \strng{authornamehash}{0565db30e661602e72fe1cf16c49f629}
      \strng{authorfullhash}{0565db30e661602e72fe1cf16c49f629}
      \strng{authorfullhashraw}{0565db30e661602e72fe1cf16c49f629}
      \field{sortinit}{H}
      \field{sortinithash}{23a3aa7c24e56cfa16945d55545109b5}
      \field{labelnamesource}{author}
      \field{labeltitlesource}{title}
      \field{eprinttype}{arXiv}
      \field{subtitle}{New Approaches for Representing Uncertainty and Making Decisions}
      \field{title}{Weighted Sets of Probabilities and Minimax Weighted Expected Regret}
      \field{year}{2013}
      \verb{eprint}
      \verb 1302.5681
      \endverb
    \endentry
    \entry{han-kloeden-2017}{book}{}{}
      \name{author}{2}{}{%
        {{hash=393e21559467f3815893e6352b04c699}{%
           family={Han},
           familyi={H\bibinitperiod},
           given={Xiaoying},
           giveni={X\bibinitperiod}}}%
        {{hash=7b32249850edcace60b1c231a0f5fb51}{%
           family={Kloeden},
           familyi={K\bibinitperiod},
           given={Peter\bibnamedelima E.},
           giveni={P\bibinitperiod\bibinitdelim E\bibinitperiod}}}%
      }
      \list{location}{1}{%
        {Singapore}%
      }
      \list{publisher}{1}{%
        {Springer}%
      }
      \strng{namehash}{115b501c1b8812b06078905db7777899}
      \strng{fullhash}{115b501c1b8812b06078905db7777899}
      \strng{fullhashraw}{115b501c1b8812b06078905db7777899}
      \strng{bibnamehash}{115b501c1b8812b06078905db7777899}
      \strng{authorbibnamehash}{115b501c1b8812b06078905db7777899}
      \strng{authornamehash}{115b501c1b8812b06078905db7777899}
      \strng{authorfullhash}{115b501c1b8812b06078905db7777899}
      \strng{authorfullhashraw}{115b501c1b8812b06078905db7777899}
      \field{sortinit}{H}
      \field{sortinithash}{23a3aa7c24e56cfa16945d55545109b5}
      \field{labelnamesource}{author}
      \field{labeltitlesource}{title}
      \field{series}{Probability Theory and Stochastic Modelling}
      \field{title}{Random Ordinary Differential Equations and Their Numerical Solution}
      \field{volume}{85}
      \field{year}{2017}
      \verb{doi}
      \verb 10.1007/978-981-10-6265-0
      \endverb
    \endentry
    \entry{handelman1988}{article}{}{}
      \name{author}{1}{}{%
        {{hash=0c39954c1c1cd6fe1031adecbd30ebf1}{%
           family={Handelman},
           familyi={H\bibinitperiod},
           given={D.},
           giveni={D\bibinitperiod}}}%
      }
      \strng{namehash}{0c39954c1c1cd6fe1031adecbd30ebf1}
      \strng{fullhash}{0c39954c1c1cd6fe1031adecbd30ebf1}
      \strng{fullhashraw}{0c39954c1c1cd6fe1031adecbd30ebf1}
      \strng{bibnamehash}{0c39954c1c1cd6fe1031adecbd30ebf1}
      \strng{authorbibnamehash}{0c39954c1c1cd6fe1031adecbd30ebf1}
      \strng{authornamehash}{0c39954c1c1cd6fe1031adecbd30ebf1}
      \strng{authorfullhash}{0c39954c1c1cd6fe1031adecbd30ebf1}
      \strng{authorfullhashraw}{0c39954c1c1cd6fe1031adecbd30ebf1}
      \field{sortinit}{H}
      \field{sortinithash}{23a3aa7c24e56cfa16945d55545109b5}
      \field{labelnamesource}{author}
      \field{labeltitlesource}{title}
      \field{journaltitle}{Pacific Journal of Mathematics}
      \field{number}{1}
      \field{title}{Representing polynomials by positive linear functions on compact convex polyhedra}
      \field{volume}{132}
      \field{year}{1988}
      \field{pages}{35\bibrangedash 62}
      \range{pages}{28}
    \endentry
    \entry{hartman-2002}{book}{}{}
      \name{author}{1}{}{%
        {{hash=0cb9ea41ed712b95ddae79d0c069518b}{%
           family={Hartman},
           familyi={H\bibinitperiod},
           given={Philip},
           giveni={P\bibinitperiod}}}%
      }
      \list{location}{1}{%
        {Philadelphia}%
      }
      \list{publisher}{2}{%
        {Society for Industrial}%
        {Applied Mathematics}%
      }
      \strng{namehash}{0cb9ea41ed712b95ddae79d0c069518b}
      \strng{fullhash}{0cb9ea41ed712b95ddae79d0c069518b}
      \strng{fullhashraw}{0cb9ea41ed712b95ddae79d0c069518b}
      \strng{bibnamehash}{0cb9ea41ed712b95ddae79d0c069518b}
      \strng{authorbibnamehash}{0cb9ea41ed712b95ddae79d0c069518b}
      \strng{authornamehash}{0cb9ea41ed712b95ddae79d0c069518b}
      \strng{authorfullhash}{0cb9ea41ed712b95ddae79d0c069518b}
      \strng{authorfullhashraw}{0cb9ea41ed712b95ddae79d0c069518b}
      \field{sortinit}{H}
      \field{sortinithash}{23a3aa7c24e56cfa16945d55545109b5}
      \field{labelnamesource}{author}
      \field{labeltitlesource}{title}
      \field{edition}{2}
      \field{series}{Classics in Applied Mathematics}
      \field{title}{Ordinary Differential Equations}
      \field{volume}{38}
      \field{year}{2002}
    \endentry
    \entry{hartmanns19}{online}{}{}
      \name{author}{2}{}{%
        {{hash=5ee8e32c129a82f9fe67944701c268bf}{%
           family={Hartmanns},
           familyi={H\bibinitperiod},
           given={Arnd},
           giveni={A\bibinitperiod}}}%
        {{hash=a82c04370c73f1ddff9a3487e8d8044d}{%
           family={Kaminski},
           familyi={K\bibinitperiod},
           given={Benjamin\bibnamedelima Lucien},
           giveni={B\bibinitperiod\bibinitdelim L\bibinitperiod}}}%
      }
      \strng{namehash}{bce0ab718af540533eddbc2f01acab83}
      \strng{fullhash}{bce0ab718af540533eddbc2f01acab83}
      \strng{fullhashraw}{bce0ab718af540533eddbc2f01acab83}
      \strng{bibnamehash}{bce0ab718af540533eddbc2f01acab83}
      \strng{authorbibnamehash}{bce0ab718af540533eddbc2f01acab83}
      \strng{authornamehash}{bce0ab718af540533eddbc2f01acab83}
      \strng{authorfullhash}{bce0ab718af540533eddbc2f01acab83}
      \strng{authorfullhashraw}{bce0ab718af540533eddbc2f01acab83}
      \field{sortinit}{H}
      \field{sortinithash}{23a3aa7c24e56cfa16945d55545109b5}
      \field{labelnamesource}{author}
      \field{labeltitlesource}{title}
      \field{eprinttype}{arXiv}
      \field{title}{Optimistic Value Iteration}
      \field{year}{2019}
      \verb{eprint}
      \verb 1910.01100
      \endverb
    \endentry
    \entry{hkm76}{article}{}{}
      \name{author}{3}{}{%
        {{hash=2be8e4f9a1eec9e94013d9d674ce8962}{%
           family={Hawking},
           familyi={H\bibinitperiod},
           given={S.\bibnamedelimi W.},
           giveni={S\bibinitperiod\bibinitdelim W\bibinitperiod}}}%
        {{hash=2da80fe021bbb1a6ad108dd69173a1b2}{%
           family={King},
           familyi={K\bibinitperiod},
           given={A.\bibnamedelimi R.},
           giveni={A\bibinitperiod\bibinitdelim R\bibinitperiod}}}%
        {{hash=90b02dfd0d19aa110d3f4f19ceb34990}{%
           family={McCarthy},
           familyi={M\bibinitperiod},
           given={P.\bibnamedelimi J.},
           giveni={P\bibinitperiod\bibinitdelim J\bibinitperiod}}}%
      }
      \strng{namehash}{8ed00470edc9917470d16d0ddb21f03b}
      \strng{fullhash}{8ed00470edc9917470d16d0ddb21f03b}
      \strng{fullhashraw}{8ed00470edc9917470d16d0ddb21f03b}
      \strng{bibnamehash}{8ed00470edc9917470d16d0ddb21f03b}
      \strng{authorbibnamehash}{8ed00470edc9917470d16d0ddb21f03b}
      \strng{authornamehash}{8ed00470edc9917470d16d0ddb21f03b}
      \strng{authorfullhash}{8ed00470edc9917470d16d0ddb21f03b}
      \strng{authorfullhashraw}{8ed00470edc9917470d16d0ddb21f03b}
      \field{sortinit}{H}
      \field{sortinithash}{23a3aa7c24e56cfa16945d55545109b5}
      \field{labelnamesource}{author}
      \field{labeltitlesource}{title}
      \field{journaltitle}{Journal of Mathematical Physics}
      \field{title}{A new topology for curved space-time which incorporates the causal, differential, and conformal structures}
      \field{volume}{17}
      \field{year}{1976}
      \field{pages}{174\bibrangedash 181}
      \range{pages}{8}
      \verb{doi}
      \verb 10.1063/1.522874
      \endverb
    \endentry
    \entry{hensel2022storm}{article}{}{}
      \name{author}{5}{}{%
        {{hash=226cb51ae8bfb05027044f34b7c88a15}{%
           family={Hensel},
           familyi={H\bibinitperiod},
           given={Christian},
           giveni={C\bibinitperiod}}}%
        {{hash=29202fc00e502e6a06092a566a4d041e}{%
           family={Junges},
           familyi={J\bibinitperiod},
           given={Sebastian},
           giveni={S\bibinitperiod}}}%
        {{hash=392d5825b15ef7f81f68632c7a49fdef}{%
           family={Katoen},
           familyi={K\bibinitperiod},
           given={Joost-Pieter},
           giveni={J\bibinithyphendelim P\bibinitperiod}}}%
        {{hash=d195773ad1a45df14875e5252d56ccb8}{%
           family={Quatmann},
           familyi={Q\bibinitperiod},
           given={Tim},
           giveni={T\bibinitperiod}}}%
        {{hash=ac49976e229e5180f7f8ba6acbfab71b}{%
           family={Volk},
           familyi={V\bibinitperiod},
           given={Matthias},
           giveni={M\bibinitperiod}}}%
      }
      \strng{namehash}{b68a4aec1e3e3526830845744c661f7f}
      \strng{fullhash}{b68a4aec1e3e3526830845744c661f7f}
      \strng{fullhashraw}{b68a4aec1e3e3526830845744c661f7f}
      \strng{bibnamehash}{b68a4aec1e3e3526830845744c661f7f}
      \strng{authorbibnamehash}{b68a4aec1e3e3526830845744c661f7f}
      \strng{authornamehash}{b68a4aec1e3e3526830845744c661f7f}
      \strng{authorfullhash}{b68a4aec1e3e3526830845744c661f7f}
      \strng{authorfullhashraw}{b68a4aec1e3e3526830845744c661f7f}
      \field{sortinit}{H}
      \field{sortinithash}{23a3aa7c24e56cfa16945d55545109b5}
      \field{labelnamesource}{author}
      \field{labeltitlesource}{title}
      \field{journaltitle}{International Journal on Software Tools for Technology Transfer}
      \field{number}{4}
      \field{title}{The probabilistic model checker {Storm}}
      \field{volume}{24}
      \field{year}{2022}
      \field{pages}{589\bibrangedash 610}
      \range{pages}{22}
      \verb{doi}
      \verb 10.1007/s10009-021-00633-z
      \endverb
    \endentry
    \entry{hornjohnson2013}{book}{}{}
      \name{author}{2}{}{%
        {{hash=5fdee541b01fd33bbb2212303937221b}{%
           family={Horn},
           familyi={H\bibinitperiod},
           given={R.\bibnamedelimi A.},
           giveni={R\bibinitperiod\bibinitdelim A\bibinitperiod}}}%
        {{hash=9fd452b06c1678e913431269679dc23d}{%
           family={Johnson},
           familyi={J\bibinitperiod},
           given={C.\bibnamedelimi R.},
           giveni={C\bibinitperiod\bibinitdelim R\bibinitperiod}}}%
      }
      \list{publisher}{1}{%
        {Cambridge University Press}%
      }
      \strng{namehash}{0bad90c712fcf551bc288a8698467a09}
      \strng{fullhash}{0bad90c712fcf551bc288a8698467a09}
      \strng{fullhashraw}{0bad90c712fcf551bc288a8698467a09}
      \strng{bibnamehash}{0bad90c712fcf551bc288a8698467a09}
      \strng{authorbibnamehash}{0bad90c712fcf551bc288a8698467a09}
      \strng{authornamehash}{0bad90c712fcf551bc288a8698467a09}
      \strng{authorfullhash}{0bad90c712fcf551bc288a8698467a09}
      \strng{authorfullhashraw}{0bad90c712fcf551bc288a8698467a09}
      \field{sortinit}{H}
      \field{sortinithash}{23a3aa7c24e56cfa16945d55545109b5}
      \field{labelnamesource}{author}
      \field{labeltitlesource}{title}
      \field{edition}{2}
      \field{note}{(Sylvester's criterion, congruence, and Gershgorin discs, as cited in the certificate proofs.)}
      \field{title}{Matrix Analysis}
      \field{year}{2013}
    \endentry
    \entry{hsu2017}{thesis}{}{}
      \name{author}{1}{}{%
        {{hash=a0500bbf3611a2701c42d5a654c84732}{%
           family={Hsu},
           familyi={H\bibinitperiod},
           given={J.},
           giveni={J\bibinitperiod}}}%
      }
      \list{institution}{1}{%
        {University of Pennsylvania}%
      }
      \strng{namehash}{a0500bbf3611a2701c42d5a654c84732}
      \strng{fullhash}{a0500bbf3611a2701c42d5a654c84732}
      \strng{fullhashraw}{a0500bbf3611a2701c42d5a654c84732}
      \strng{bibnamehash}{a0500bbf3611a2701c42d5a654c84732}
      \strng{authorbibnamehash}{a0500bbf3611a2701c42d5a654c84732}
      \strng{authornamehash}{a0500bbf3611a2701c42d5a654c84732}
      \strng{authorfullhash}{a0500bbf3611a2701c42d5a654c84732}
      \strng{authorfullhashraw}{a0500bbf3611a2701c42d5a654c84732}
      \field{sortinit}{H}
      \field{sortinithash}{23a3aa7c24e56cfa16945d55545109b5}
      \field{labelnamesource}{author}
      \field{labeltitlesource}{title}
      \field{eprinttype}{arxiv}
      \field{note}{(Couplings as certificates in relational program logics --- the pre/post-style sibling of the judgment form used here.)}
      \field{title}{Probabilistic Couplings for Probabilistic Reasoning}
      \field{type}{PhD thesis}
      \field{year}{2017}
      \verb{eprint}
      \verb 1710.09951
      \endverb
    \endentry
    \entry{huot23omega}{online}{}{}
      \name{author}{4}{}{%
        {{hash=d0272179a035f6ed4f8d69f30b68eed6}{%
           family={Huot},
           familyi={H\bibinitperiod},
           given={Mathieu},
           giveni={M\bibinitperiod}}}%
        {{hash=176f2ce8f6edbda13ed3c51959a4255b}{%
           family={Lew},
           familyi={L\bibinitperiod},
           given={Alexander\bibnamedelima K.},
           giveni={A\bibinitperiod\bibinitdelim K\bibinitperiod}}}%
        {{hash=b4a161cc1ca940e73d86b630facd85f5}{%
           family={Mansinghka},
           familyi={M\bibinitperiod},
           given={Vikash\bibnamedelima K.},
           giveni={V\bibinitperiod\bibinitdelim K\bibinitperiod}}}%
        {{hash=5351d3a060653e114d956a0a3700d556}{%
           family={Staton},
           familyi={S\bibinitperiod},
           given={Sam},
           giveni={S\bibinitperiod}}}%
      }
      \strng{namehash}{e260616f6993b7f22af05af0cb2a7e32}
      \strng{fullhash}{e260616f6993b7f22af05af0cb2a7e32}
      \strng{fullhashraw}{e260616f6993b7f22af05af0cb2a7e32}
      \strng{bibnamehash}{e260616f6993b7f22af05af0cb2a7e32}
      \strng{authorbibnamehash}{e260616f6993b7f22af05af0cb2a7e32}
      \strng{authornamehash}{e260616f6993b7f22af05af0cb2a7e32}
      \strng{authorfullhash}{e260616f6993b7f22af05af0cb2a7e32}
      \strng{authorfullhashraw}{e260616f6993b7f22af05af0cb2a7e32}
      \field{extraname}{1}
      \field{sortinit}{H}
      \field{sortinithash}{23a3aa7c24e56cfa16945d55545109b5}
      \field{labelnamesource}{author}
      \field{labeltitlesource}{title}
      \field{eprinttype}{arXiv}
      \field{subtitle}{Reasoning Denotationally About Higher-Order, Recursive Probabilistic and Differentiable Programs}
      \field{title}{$\omega$PAP Spaces}
      \field{year}{2023}
      \verb{eprint}
      \verb 2302.10636
      \endverb
    \endentry
    \entry{huot23}{online}{}{}
      \name{author}{4}{}{%
        {{hash=d0272179a035f6ed4f8d69f30b68eed6}{%
           family={Huot},
           familyi={H\bibinitperiod},
           given={Mathieu},
           giveni={M\bibinitperiod}}}%
        {{hash=176f2ce8f6edbda13ed3c51959a4255b}{%
           family={Lew},
           familyi={L\bibinitperiod},
           given={Alexander\bibnamedelima K.},
           giveni={A\bibinitperiod\bibinitdelim K\bibinitperiod}}}%
        {{hash=b4a161cc1ca940e73d86b630facd85f5}{%
           family={Mansinghka},
           familyi={M\bibinitperiod},
           given={Vikash\bibnamedelima K.},
           giveni={V\bibinitperiod\bibinitdelim K\bibinitperiod}}}%
        {{hash=5351d3a060653e114d956a0a3700d556}{%
           family={Staton},
           familyi={S\bibinitperiod},
           given={Sam},
           giveni={S\bibinitperiod}}}%
      }
      \strng{namehash}{e260616f6993b7f22af05af0cb2a7e32}
      \strng{fullhash}{e260616f6993b7f22af05af0cb2a7e32}
      \strng{fullhashraw}{e260616f6993b7f22af05af0cb2a7e32}
      \strng{bibnamehash}{e260616f6993b7f22af05af0cb2a7e32}
      \strng{authorbibnamehash}{e260616f6993b7f22af05af0cb2a7e32}
      \strng{authornamehash}{e260616f6993b7f22af05af0cb2a7e32}
      \strng{authorfullhash}{e260616f6993b7f22af05af0cb2a7e32}
      \strng{authorfullhashraw}{e260616f6993b7f22af05af0cb2a7e32}
      \field{extraname}{2}
      \field{sortinit}{H}
      \field{sortinithash}{23a3aa7c24e56cfa16945d55545109b5}
      \field{labelnamesource}{author}
      \field{labeltitlesource}{title}
      \field{eprinttype}{arXiv}
      \field{subtitle}{Reasoning Denotationally About Higher-Order, Recursive Probabilistic and Differentiable Programs}
      \field{title}{$\omega$PAP Spaces}
      \field{year}{2023}
      \verb{eprint}
      \verb 2302.10636
      \endverb
    \endentry
    \entry{hurwicz51}{report}{}{}
      \name{author}{1}{}{%
        {{hash=3aa2ee3b3c32f690863d3404c0171553}{%
           family={Hurwicz},
           familyi={H\bibinitperiod},
           given={Leonid},
           giveni={L\bibinitperiod}}}%
      }
      \list{institution}{1}{%
        {Cowles Commission for Research in Economics}%
      }
      \strng{namehash}{3aa2ee3b3c32f690863d3404c0171553}
      \strng{fullhash}{3aa2ee3b3c32f690863d3404c0171553}
      \strng{fullhashraw}{3aa2ee3b3c32f690863d3404c0171553}
      \strng{bibnamehash}{3aa2ee3b3c32f690863d3404c0171553}
      \strng{authorbibnamehash}{3aa2ee3b3c32f690863d3404c0171553}
      \strng{authornamehash}{3aa2ee3b3c32f690863d3404c0171553}
      \strng{authorfullhash}{3aa2ee3b3c32f690863d3404c0171553}
      \strng{authorfullhashraw}{3aa2ee3b3c32f690863d3404c0171553}
      \field{sortinit}{H}
      \field{sortinithash}{23a3aa7c24e56cfa16945d55545109b5}
      \field{labelnamesource}{author}
      \field{labeltitlesource}{title}
      \field{number}{356}
      \field{title}{A Class of Criteria for Decision-Making under Ignorance}
      \field{type}{Cowles Commission Discussion Paper: Statistics}
      \field{year}{1951}
    \endentry
    \entry{ieee754}{report}{}{}
      \name{author}{1}{}{%
        {{hash=472abfae8142fdb6f39d081303e8af85}{%
           family={{IEEE}},
           familyi={I\bibinitperiod}}}%
      }
      \list{institution}{2}{%
        {Institute of Electrical}%
        {Electronics Engineers}%
      }
      \strng{namehash}{472abfae8142fdb6f39d081303e8af85}
      \strng{fullhash}{472abfae8142fdb6f39d081303e8af85}
      \strng{fullhashraw}{472abfae8142fdb6f39d081303e8af85}
      \strng{bibnamehash}{472abfae8142fdb6f39d081303e8af85}
      \strng{authorbibnamehash}{472abfae8142fdb6f39d081303e8af85}
      \strng{authornamehash}{472abfae8142fdb6f39d081303e8af85}
      \strng{authorfullhash}{472abfae8142fdb6f39d081303e8af85}
      \strng{authorfullhashraw}{472abfae8142fdb6f39d081303e8af85}
      \field{sortinit}{I}
      \field{sortinithash}{8d291c51ee89b6cd86bf5379f0b151d8}
      \field{labelnamesource}{author}
      \field{labeltitlesource}{title}
      \field{note}{Revision of IEEE Std 754-2008}
      \field{number}{754-2019}
      \field{title}{IEEE Standard for Floating-Point Arithmetic}
      \field{type}{IEEE Std}
      \field{year}{2019}
    \endentry
    \entry{miranda-zaffalon-2015-independent-products}{online}{}{}
      \field{sortinit}{I}
      \field{sortinithash}{8d291c51ee89b6cd86bf5379f0b151d8}
      \field{labeltitlesource}{title}
      \field{title}{Independent products in infinite spaces (Miranda \& Zaffalon, JMAA 2015)}
      \verb{urlraw}
      \verb https://people.idsia.ch/~zaffalon/papers/2015jmaa-indprods.pdf
      \endverb
      \verb{url}
      \verb https://people.idsia.ch/~zaffalon/papers/2015jmaa-indprods.pdf
      \endverb
    \endentry
    \entry{irvingaks}{online}{}{}
      \name{author}{1}{}{%
        {{hash=0569b332096c1f2df7f470fd478f4cc3}{%
           family={Irving},
           familyi={I\bibinitperiod},
           given={G.},
           giveni={G\bibinitperiod}}}%
      }
      \strng{namehash}{0569b332096c1f2df7f470fd478f4cc3}
      \strng{fullhash}{0569b332096c1f2df7f470fd478f4cc3}
      \strng{fullhashraw}{0569b332096c1f2df7f470fd478f4cc3}
      \strng{bibnamehash}{0569b332096c1f2df7f470fd478f4cc3}
      \strng{authorbibnamehash}{0569b332096c1f2df7f470fd478f4cc3}
      \strng{authornamehash}{0569b332096c1f2df7f470fd478f4cc3}
      \strng{authorfullhash}{0569b332096c1f2df7f470fd478f4cc3}
      \strng{authorfullhashraw}{0569b332096c1f2df7f470fd478f4cc3}
      \field{extraname}{1}
      \field{sortinit}{I}
      \field{sortinithash}{8d291c51ee89b6cd86bf5379f0b151d8}
      \field{labelnamesource}{author}
      \field{labeltitlesource}{title}
      \field{note}{Lean formalization repository: the \texttt{Random} library's diagonal-dominance PSD bridge and a certified $65536$-node spectral expander.}
      \field{title}{aks}
      \field{year}{2026}
      \verb{urlraw}
      \verb https://github.com/girving/aks
      \endverb
      \verb{url}
      \verb https://github.com/girving/aks
      \endverb
    \endentry
    \entry{irving2026certificate}{misc}{}{}
      \name{author}{1}{}{%
        {{hash=0569b332096c1f2df7f470fd478f4cc3}{%
           family={Irving},
           familyi={I\bibinitperiod},
           given={G.},
           giveni={G\bibinitperiod}}}%
      }
      \strng{namehash}{0569b332096c1f2df7f470fd478f4cc3}
      \strng{fullhash}{0569b332096c1f2df7f470fd478f4cc3}
      \strng{fullhashraw}{0569b332096c1f2df7f470fd478f4cc3}
      \strng{bibnamehash}{0569b332096c1f2df7f470fd478f4cc3}
      \strng{authorbibnamehash}{0569b332096c1f2df7f470fd478f4cc3}
      \strng{authornamehash}{0569b332096c1f2df7f470fd478f4cc3}
      \strng{authorfullhash}{0569b332096c1f2df7f470fd478f4cc3}
      \strng{authorfullhashraw}{0569b332096c1f2df7f470fd478f4cc3}
      \field{extraname}{2}
      \field{sortinit}{I}
      \field{sortinithash}{8d291c51ee89b6cd86bf5379f0b151d8}
      \field{labelnamesource}{author}
      \field{labeltitlesource}{title}
      \field{howpublished}{Computational Science Stack Exchange (question 45365, answer 45400)}
      \field{note}{(The rounding-tolerant integer variant and its certificate-size analysis.)}
      \field{title}{Answer to: Certificate for strict positive definiteness of a sparse matrix}
      \field{year}{2026}
    \endentry
    \entry{jacobs99}{book}{}{}
      \name{author}{1}{}{%
        {{hash=bf424ec58729aa7d35e9d57ec1c7d40d}{%
           family={Jacobs},
           familyi={J\bibinitperiod},
           given={Bart},
           giveni={B\bibinitperiod}}}%
      }
      \list{publisher}{1}{%
        {North-Holland}%
      }
      \strng{namehash}{bf424ec58729aa7d35e9d57ec1c7d40d}
      \strng{fullhash}{bf424ec58729aa7d35e9d57ec1c7d40d}
      \strng{fullhashraw}{bf424ec58729aa7d35e9d57ec1c7d40d}
      \strng{bibnamehash}{bf424ec58729aa7d35e9d57ec1c7d40d}
      \strng{authorbibnamehash}{bf424ec58729aa7d35e9d57ec1c7d40d}
      \strng{authornamehash}{bf424ec58729aa7d35e9d57ec1c7d40d}
      \strng{authorfullhash}{bf424ec58729aa7d35e9d57ec1c7d40d}
      \strng{authorfullhashraw}{bf424ec58729aa7d35e9d57ec1c7d40d}
      \field{sortinit}{J}
      \field{sortinithash}{b2f54a9081ace9966a7cb9413811edb4}
      \field{labelnamesource}{author}
      \field{labeltitlesource}{title}
      \field{series}{Studies in Logic and the Foundations of Mathematics}
      \field{title}{Categorical Logic and Type Theory}
      \field{volume}{141}
      \field{year}{1999}
    \endentry
    \entry{jiangmareelswang1996}{article}{}{}
      \name{author}{3}{}{%
        {{hash=4f07ddfcfbc87b7155dd03a1345bd112}{%
           family={Jiang},
           familyi={J\bibinitperiod},
           given={Z.-P.},
           giveni={Z\bibinithyphendelim P\bibinitperiod}}}%
        {{hash=eeb739a2c8e69d946c86844708ae2f84}{%
           family={Mareels},
           familyi={M\bibinitperiod},
           given={I.\bibnamedelimi M.\bibnamedelimi Y.},
           giveni={I\bibinitperiod\bibinitdelim M\bibinitperiod\bibinitdelim Y\bibinitperiod}}}%
        {{hash=61ad10c626e1d641e0b1d7b5f48d3233}{%
           family={Wang},
           familyi={W\bibinitperiod},
           given={Y.},
           giveni={Y\bibinitperiod}}}%
      }
      \strng{namehash}{e65db8d8f7bf44b0fcf6d54eb6c5a3c8}
      \strng{fullhash}{e65db8d8f7bf44b0fcf6d54eb6c5a3c8}
      \strng{fullhashraw}{e65db8d8f7bf44b0fcf6d54eb6c5a3c8}
      \strng{bibnamehash}{e65db8d8f7bf44b0fcf6d54eb6c5a3c8}
      \strng{authorbibnamehash}{e65db8d8f7bf44b0fcf6d54eb6c5a3c8}
      \strng{authornamehash}{e65db8d8f7bf44b0fcf6d54eb6c5a3c8}
      \strng{authorfullhash}{e65db8d8f7bf44b0fcf6d54eb6c5a3c8}
      \strng{authorfullhashraw}{e65db8d8f7bf44b0fcf6d54eb6c5a3c8}
      \field{sortinit}{J}
      \field{sortinithash}{b2f54a9081ace9966a7cb9413811edb4}
      \field{labelnamesource}{author}
      \field{labeltitlesource}{title}
      \field{journaltitle}{Automatica}
      \field{number}{8}
      \field{title}{A Lyapunov formulation of the nonlinear small-gain theorem for interconnected ISS systems}
      \field{volume}{32}
      \field{year}{1996}
      \field{pages}{1211\bibrangedash 1215}
      \range{pages}{5}
    \endentry
    \entry{jiangteelpraly1994}{article}{}{}
      \name{author}{3}{}{%
        {{hash=4f07ddfcfbc87b7155dd03a1345bd112}{%
           family={Jiang},
           familyi={J\bibinitperiod},
           given={Z.-P.},
           giveni={Z\bibinithyphendelim P\bibinitperiod}}}%
        {{hash=486d8a385412f4582607e04104cefa04}{%
           family={Teel},
           familyi={T\bibinitperiod},
           given={A.\bibnamedelimi R.},
           giveni={A\bibinitperiod\bibinitdelim R\bibinitperiod}}}%
        {{hash=c54a948c701456bdd11addbf6e16a568}{%
           family={Praly},
           familyi={P\bibinitperiod},
           given={L.},
           giveni={L\bibinitperiod}}}%
      }
      \strng{namehash}{ec6bfc5cc4cb612078d460abc0e0f06a}
      \strng{fullhash}{ec6bfc5cc4cb612078d460abc0e0f06a}
      \strng{fullhashraw}{ec6bfc5cc4cb612078d460abc0e0f06a}
      \strng{bibnamehash}{ec6bfc5cc4cb612078d460abc0e0f06a}
      \strng{authorbibnamehash}{ec6bfc5cc4cb612078d460abc0e0f06a}
      \strng{authornamehash}{ec6bfc5cc4cb612078d460abc0e0f06a}
      \strng{authorfullhash}{ec6bfc5cc4cb612078d460abc0e0f06a}
      \strng{authorfullhashraw}{ec6bfc5cc4cb612078d460abc0e0f06a}
      \field{sortinit}{J}
      \field{sortinithash}{b2f54a9081ace9966a7cb9413811edb4}
      \field{labelnamesource}{author}
      \field{labeltitlesource}{title}
      \field{journaltitle}{Mathematics of Control, Signals and Systems}
      \field{title}{Small-gain theorem for ISS systems and applications}
      \field{volume}{7}
      \field{year}{1994}
      \field{pages}{95\bibrangedash 120}
      \range{pages}{26}
    \endentry
    \entry{johnstone2002}{book}{}{}
      \name{author}{1}{}{%
        {{hash=0fd3c6c701987f321528f4265c0c3449}{%
           family={Johnstone},
           familyi={J\bibinitperiod},
           given={Peter\bibnamedelima T.},
           giveni={P\bibinitperiod\bibinitdelim T\bibinitperiod}}}%
      }
      \list{location}{1}{%
        {Oxford}%
      }
      \list{publisher}{1}{%
        {Oxford University Press}%
      }
      \strng{namehash}{0fd3c6c701987f321528f4265c0c3449}
      \strng{fullhash}{0fd3c6c701987f321528f4265c0c3449}
      \strng{fullhashraw}{0fd3c6c701987f321528f4265c0c3449}
      \strng{bibnamehash}{0fd3c6c701987f321528f4265c0c3449}
      \strng{authorbibnamehash}{0fd3c6c701987f321528f4265c0c3449}
      \strng{authornamehash}{0fd3c6c701987f321528f4265c0c3449}
      \strng{authorfullhash}{0fd3c6c701987f321528f4265c0c3449}
      \strng{authorfullhashraw}{0fd3c6c701987f321528f4265c0c3449}
      \field{sortinit}{J}
      \field{sortinithash}{b2f54a9081ace9966a7cb9413811edb4}
      \field{labelnamesource}{author}
      \field{labeltitlesource}{title}
      \field{note}{Geometric logic: Part~D1}
      \field{subtitle}{A Topos Theory Compendium}
      \field{title}{Sketches of an Elephant}
      \field{year}{2002}
    \endentry
    \entry{kabra26}{online}{}{}
      \name{author}{5}{}{%
        {{hash=7b46763b2066ed9e9ee40909a3942ceb}{%
           family={Kabra},
           familyi={K\bibinitperiod},
           given={Aditi},
           giveni={A\bibinitperiod}}}%
        {{hash=808513bc2b23b1fee6a93e13ee2c2252}{%
           family={Laurent},
           familyi={L\bibinitperiod},
           given={Jonathan},
           giveni={J\bibinitperiod}}}%
        {{hash=19aae6673b26cc3166341aa1186bcd21}{%
           family={Martins},
           familyi={M\bibinitperiod},
           given={Ruben},
           giveni={R\bibinitperiod}}}%
        {{hash=b723c481b381d20105178d693c40d93b}{%
           family={Mitsch},
           familyi={M\bibinitperiod},
           given={Stefan},
           giveni={S\bibinitperiod}}}%
        {{hash=0fb7e4129542e3e84c6115bee2a2b227}{%
           family={Platzer},
           familyi={P\bibinitperiod},
           given={André},
           giveni={A\bibinitperiod}}}%
      }
      \strng{namehash}{e5eb42fc1584e7f7efae78894846d49d}
      \strng{fullhash}{e5eb42fc1584e7f7efae78894846d49d}
      \strng{fullhashraw}{e5eb42fc1584e7f7efae78894846d49d}
      \strng{bibnamehash}{e5eb42fc1584e7f7efae78894846d49d}
      \strng{authorbibnamehash}{e5eb42fc1584e7f7efae78894846d49d}
      \strng{authornamehash}{e5eb42fc1584e7f7efae78894846d49d}
      \strng{authorfullhash}{e5eb42fc1584e7f7efae78894846d49d}
      \strng{authorfullhashraw}{e5eb42fc1584e7f7efae78894846d49d}
      \field{sortinit}{K}
      \field{sortinithash}{c02bf6bff1c488450c352b40f5d853ab}
      \field{labelnamesource}{author}
      \field{labeltitlesource}{title}
      \field{eprinttype}{arXiv}
      \field{title}{LLM-Powered Automatic Theorem Proving and Synthesis for Hybrid Systems and Games}
      \field{year}{2026}
      \verb{eprint}
      \verb 2603.00737
      \endverb
    \endentry
    \entry{kaelbling-littman-cassandra-1998}{article}{}{}
      \name{author}{3}{}{%
        {{hash=72fdc7bf3120e1cff12e3ecf7e892b9a}{%
           family={Kaelbling},
           familyi={K\bibinitperiod},
           given={Leslie\bibnamedelima Pack},
           giveni={L\bibinitperiod\bibinitdelim P\bibinitperiod}}}%
        {{hash=211e529e9f10ad1bfacff2284da10ea7}{%
           family={Littman},
           familyi={L\bibinitperiod},
           given={Michael\bibnamedelima L.},
           giveni={M\bibinitperiod\bibinitdelim L\bibinitperiod}}}%
        {{hash=2206f2c6c327af01902e88ec6bef588b}{%
           family={Cassandra},
           familyi={C\bibinitperiod},
           given={Anthony\bibnamedelima R.},
           giveni={A\bibinitperiod\bibinitdelim R\bibinitperiod}}}%
      }
      \strng{namehash}{55978656d14c32546fe5d19b1598192d}
      \strng{fullhash}{55978656d14c32546fe5d19b1598192d}
      \strng{fullhashraw}{55978656d14c32546fe5d19b1598192d}
      \strng{bibnamehash}{55978656d14c32546fe5d19b1598192d}
      \strng{authorbibnamehash}{55978656d14c32546fe5d19b1598192d}
      \strng{authornamehash}{55978656d14c32546fe5d19b1598192d}
      \strng{authorfullhash}{55978656d14c32546fe5d19b1598192d}
      \strng{authorfullhashraw}{55978656d14c32546fe5d19b1598192d}
      \field{sortinit}{K}
      \field{sortinithash}{c02bf6bff1c488450c352b40f5d853ab}
      \field{labelnamesource}{author}
      \field{labeltitlesource}{title}
      \field{journaltitle}{Artificial Intelligence}
      \field{number}{1--2}
      \field{title}{Planning and Acting in Partially Observable Stochastic Domains}
      \field{volume}{101}
      \field{year}{1998}
      \field{pages}{99\bibrangedash 134}
      \range{pages}{36}
      \verb{doi}
      \verb 10.1016/S0004-3702(98)00023-X
      \endverb
    \endentry
    \entry{kmv2005}{article}{}{}
      \name{author}{3}{}{%
        {{hash=fa1b35b36a16ccd9a25e3a42a727d267}{%
           family={Kalies},
           familyi={K\bibinitperiod},
           given={W.\bibnamedelimi D.},
           giveni={W\bibinitperiod\bibinitdelim D\bibinitperiod}}}%
        {{hash=48f8b15c21dd70a5e22f1c20c558972d}{%
           family={Mischaikow},
           familyi={M\bibinitperiod},
           given={K.},
           giveni={K\bibinitperiod}}}%
        {{hash=0252d70d03c54b4ab54217e51b2c7f8a}{%
           family={Vorst},
           familyi={V\bibinitperiod},
           given={R.\bibnamedelimi C.\bibnamedelimi A.\bibnamedelimi M.},
           giveni={R\bibinitperiod\bibinitdelim C\bibinitperiod\bibinitdelim A\bibinitperiod\bibinitdelim M\bibinitperiod},
           prefix={van\bibnamedelima der},
           prefixi={v\bibinitperiod\bibinitdelim d\bibinitperiod}}}%
      }
      \strng{namehash}{2009206c8d8df931d96967fc2178bef8}
      \strng{fullhash}{2009206c8d8df931d96967fc2178bef8}
      \strng{fullhashraw}{2009206c8d8df931d96967fc2178bef8}
      \strng{bibnamehash}{2009206c8d8df931d96967fc2178bef8}
      \strng{authorbibnamehash}{2009206c8d8df931d96967fc2178bef8}
      \strng{authornamehash}{2009206c8d8df931d96967fc2178bef8}
      \strng{authorfullhash}{2009206c8d8df931d96967fc2178bef8}
      \strng{authorfullhashraw}{2009206c8d8df931d96967fc2178bef8}
      \field{extraname}{1}
      \field{sortinit}{K}
      \field{sortinithash}{c02bf6bff1c488450c352b40f5d853ab}
      \field{labelnamesource}{author}
      \field{labeltitlesource}{title}
      \field{journaltitle}{Foundations of Computational Mathematics}
      \field{number}{4}
      \field{title}{An algorithmic approach to chain recurrence}
      \field{volume}{5}
      \field{year}{2005}
      \field{pages}{409\bibrangedash 449}
      \range{pages}{41}
    \endentry
    \entry{kmv2016}{article}{}{}
      \name{author}{3}{}{%
        {{hash=fa1b35b36a16ccd9a25e3a42a727d267}{%
           family={Kalies},
           familyi={K\bibinitperiod},
           given={W.\bibnamedelimi D.},
           giveni={W\bibinitperiod\bibinitdelim D\bibinitperiod}}}%
        {{hash=48f8b15c21dd70a5e22f1c20c558972d}{%
           family={Mischaikow},
           familyi={M\bibinitperiod},
           given={K.},
           giveni={K\bibinitperiod}}}%
        {{hash=0252d70d03c54b4ab54217e51b2c7f8a}{%
           family={Vorst},
           familyi={V\bibinitperiod},
           given={R.\bibnamedelimi C.\bibnamedelimi A.\bibnamedelimi M.},
           giveni={R\bibinitperiod\bibinitdelim C\bibinitperiod\bibinitdelim A\bibinitperiod\bibinitdelim M\bibinitperiod},
           prefix={van\bibnamedelima der},
           prefixi={v\bibinitperiod\bibinitdelim d\bibinitperiod}}}%
      }
      \strng{namehash}{2009206c8d8df931d96967fc2178bef8}
      \strng{fullhash}{2009206c8d8df931d96967fc2178bef8}
      \strng{fullhashraw}{2009206c8d8df931d96967fc2178bef8}
      \strng{bibnamehash}{2009206c8d8df931d96967fc2178bef8}
      \strng{authorbibnamehash}{2009206c8d8df931d96967fc2178bef8}
      \strng{authornamehash}{2009206c8d8df931d96967fc2178bef8}
      \strng{authorfullhash}{2009206c8d8df931d96967fc2178bef8}
      \strng{authorfullhashraw}{2009206c8d8df931d96967fc2178bef8}
      \field{extraname}{2}
      \field{sortinit}{K}
      \field{sortinithash}{c02bf6bff1c488450c352b40f5d853ab}
      \field{labelnamesource}{author}
      \field{labeltitlesource}{title}
      \field{journaltitle}{Foundations of Computational Mathematics}
      \field{title}{Lattice structures for attractors II}
      \field{volume}{16}
      \field{year}{2016}
      \field{pages}{1151\bibrangedash 1191}
      \range{pages}{41}
    \endentry
    \entry{kmv2022}{article}{}{}
      \name{author}{3}{}{%
        {{hash=fa1b35b36a16ccd9a25e3a42a727d267}{%
           family={Kalies},
           familyi={K\bibinitperiod},
           given={W.\bibnamedelimi D.},
           giveni={W\bibinitperiod\bibinitdelim D\bibinitperiod}}}%
        {{hash=48f8b15c21dd70a5e22f1c20c558972d}{%
           family={Mischaikow},
           familyi={M\bibinitperiod},
           given={K.},
           giveni={K\bibinitperiod}}}%
        {{hash=0252d70d03c54b4ab54217e51b2c7f8a}{%
           family={Vorst},
           familyi={V\bibinitperiod},
           given={R.\bibnamedelimi C.\bibnamedelimi A.\bibnamedelimi M.},
           giveni={R\bibinitperiod\bibinitdelim C\bibinitperiod\bibinitdelim A\bibinitperiod\bibinitdelim M\bibinitperiod},
           prefix={van\bibnamedelima der},
           prefixi={v\bibinitperiod\bibinitdelim d\bibinitperiod}}}%
      }
      \strng{namehash}{2009206c8d8df931d96967fc2178bef8}
      \strng{fullhash}{2009206c8d8df931d96967fc2178bef8}
      \strng{fullhashraw}{2009206c8d8df931d96967fc2178bef8}
      \strng{bibnamehash}{2009206c8d8df931d96967fc2178bef8}
      \strng{authorbibnamehash}{2009206c8d8df931d96967fc2178bef8}
      \strng{authornamehash}{2009206c8d8df931d96967fc2178bef8}
      \strng{authorfullhash}{2009206c8d8df931d96967fc2178bef8}
      \strng{authorfullhashraw}{2009206c8d8df931d96967fc2178bef8}
      \field{extraname}{3}
      \field{sortinit}{K}
      \field{sortinithash}{c02bf6bff1c488450c352b40f5d853ab}
      \field{labelnamesource}{author}
      \field{labeltitlesource}{title}
      \field{eprinttype}{arxiv}
      \field{journaltitle}{Journal of Dynamics and Differential Equations}
      \field{title}{Lattice structures for attractors III}
      \field{volume}{34}
      \field{year}{2022}
      \field{pages}{1729\bibrangedash 1768}
      \range{pages}{40}
      \verb{eprint}
      \verb 1911.09382
      \endverb
    \endentry
    \entry{kalieswehbe2026}{misc}{}{}
      \name{author}{2}{}{%
        {{hash=fa1b35b36a16ccd9a25e3a42a727d267}{%
           family={Kalies},
           familyi={K\bibinitperiod},
           given={W.\bibnamedelimi D.},
           giveni={W\bibinitperiod\bibinitdelim D\bibinitperiod}}}%
        {{hash=f7f1401c5964bd2344aa6b7622fe1139}{%
           family={Wehbe},
           familyi={W\bibinitperiod},
           given={T.},
           giveni={T\bibinitperiod}}}%
      }
      \strng{namehash}{38fca1cfdc15439c84bc9b0dfe4059dd}
      \strng{fullhash}{38fca1cfdc15439c84bc9b0dfe4059dd}
      \strng{fullhashraw}{38fca1cfdc15439c84bc9b0dfe4059dd}
      \strng{bibnamehash}{38fca1cfdc15439c84bc9b0dfe4059dd}
      \strng{authorbibnamehash}{38fca1cfdc15439c84bc9b0dfe4059dd}
      \strng{authornamehash}{38fca1cfdc15439c84bc9b0dfe4059dd}
      \strng{authorfullhash}{38fca1cfdc15439c84bc9b0dfe4059dd}
      \strng{authorfullhashraw}{38fca1cfdc15439c84bc9b0dfe4059dd}
      \field{sortinit}{K}
      \field{sortinithash}{c02bf6bff1c488450c352b40f5d853ab}
      \field{labelnamesource}{author}
      \field{labeltitlesource}{title}
      \field{eprinttype}{arxiv}
      \field{title}{Compositionality of Global Dynamics in Product and Skew-Product Systems}
      \field{year}{2026}
      \verb{eprint}
      \verb 2606.30925
      \endverb
    \endentry
    \entry{katz2017reluplex}{inproceedings}{}{}
      \name{author}{5}{}{%
        {{hash=6e2d6de5ae9acd40997821f90ca1bf9e}{%
           family={Katz},
           familyi={K\bibinitperiod},
           given={G.},
           giveni={G\bibinitperiod}}}%
        {{hash=d55d7e1f1d6f23d032000be69802a71c}{%
           family={Barrett},
           familyi={B\bibinitperiod},
           given={C.},
           giveni={C\bibinitperiod}}}%
        {{hash=4ed5a1a176251b05785b44f048ab3bd2}{%
           family={Dill},
           familyi={D\bibinitperiod},
           given={D.\bibnamedelimi L.},
           giveni={D\bibinitperiod\bibinitdelim L\bibinitperiod}}}%
        {{hash=342a8e398f29231fe87028672c21a57c}{%
           family={Julian},
           familyi={J\bibinitperiod},
           given={K.},
           giveni={K\bibinitperiod}}}%
        {{hash=c5bfe331d34f6c45505f275ba5a7ee61}{%
           family={Kochenderfer},
           familyi={K\bibinitperiod},
           given={M.\bibnamedelimi J.},
           giveni={M\bibinitperiod\bibinitdelim J\bibinitperiod}}}%
      }
      \strng{namehash}{3a1bfa10adc1ab4dc53e5d36279bbd85}
      \strng{fullhash}{3a1bfa10adc1ab4dc53e5d36279bbd85}
      \strng{fullhashraw}{3a1bfa10adc1ab4dc53e5d36279bbd85}
      \strng{bibnamehash}{3a1bfa10adc1ab4dc53e5d36279bbd85}
      \strng{authorbibnamehash}{3a1bfa10adc1ab4dc53e5d36279bbd85}
      \strng{authornamehash}{3a1bfa10adc1ab4dc53e5d36279bbd85}
      \strng{authorfullhash}{3a1bfa10adc1ab4dc53e5d36279bbd85}
      \strng{authorfullhashraw}{3a1bfa10adc1ab4dc53e5d36279bbd85}
      \field{sortinit}{K}
      \field{sortinithash}{c02bf6bff1c488450c352b40f5d853ab}
      \field{labelnamesource}{author}
      \field{labeltitlesource}{title}
      \field{booktitle}{Proceedings of CAV 2017 (Computer Aided Verification)}
      \field{title}{Reluplex: an efficient SMT solver for verifying deep neural networks}
      \field{year}{2017}
    \endentry
    \entry{kemeny-snell-1960}{book}{}{}
      \name{author}{2}{}{%
        {{hash=d167832a98ab4ebb64d46588f0a8f688}{%
           family={Kemeny},
           familyi={K\bibinitperiod},
           given={John\bibnamedelima G.},
           giveni={J\bibinitperiod\bibinitdelim G\bibinitperiod}}}%
        {{hash=f64ddf8ee9b19fd9b8ea3426d553f892}{%
           family={Snell},
           familyi={S\bibinitperiod},
           given={J.\bibnamedelimi Laurie},
           giveni={J\bibinitperiod\bibinitdelim L\bibinitperiod}}}%
      }
      \list{location}{1}{%
        {Princeton}%
      }
      \list{publisher}{1}{%
        {Van Nostrand}%
      }
      \strng{namehash}{e18622385b5e778e59b96164a539d8a4}
      \strng{fullhash}{e18622385b5e778e59b96164a539d8a4}
      \strng{fullhashraw}{e18622385b5e778e59b96164a539d8a4}
      \strng{bibnamehash}{e18622385b5e778e59b96164a539d8a4}
      \strng{authorbibnamehash}{e18622385b5e778e59b96164a539d8a4}
      \strng{authornamehash}{e18622385b5e778e59b96164a539d8a4}
      \strng{authorfullhash}{e18622385b5e778e59b96164a539d8a4}
      \strng{authorfullhashraw}{e18622385b5e778e59b96164a539d8a4}
      \field{sortinit}{K}
      \field{sortinithash}{c02bf6bff1c488450c352b40f5d853ab}
      \field{labelnamesource}{author}
      \field{labeltitlesource}{title}
      \field{title}{Finite Markov Chains}
      \field{year}{1960}
    \endentry
    \entry{khalil2002nonlinear}{book}{}{}
      \name{author}{1}{}{%
        {{hash=c687708249e29c8fa3d8d0c0ff1612f9}{%
           family={Khalil},
           familyi={K\bibinitperiod},
           given={Hassan\bibnamedelima K.},
           giveni={H\bibinitperiod\bibinitdelim K\bibinitperiod}}}%
      }
      \list{publisher}{1}{%
        {Prentice Hall}%
      }
      \strng{namehash}{c687708249e29c8fa3d8d0c0ff1612f9}
      \strng{fullhash}{c687708249e29c8fa3d8d0c0ff1612f9}
      \strng{fullhashraw}{c687708249e29c8fa3d8d0c0ff1612f9}
      \strng{bibnamehash}{c687708249e29c8fa3d8d0c0ff1612f9}
      \strng{authorbibnamehash}{c687708249e29c8fa3d8d0c0ff1612f9}
      \strng{authornamehash}{c687708249e29c8fa3d8d0c0ff1612f9}
      \strng{authorfullhash}{c687708249e29c8fa3d8d0c0ff1612f9}
      \strng{authorfullhashraw}{c687708249e29c8fa3d8d0c0ff1612f9}
      \field{sortinit}{K}
      \field{sortinithash}{c02bf6bff1c488450c352b40f5d853ab}
      \field{labelnamesource}{author}
      \field{labeltitlesource}{title}
      \field{edition}{3}
      \field{title}{Nonlinear Systems}
      \field{year}{2002}
    \endentry
    \entry{kissinger13}{online}{}{}
      \name{author}{1}{}{%
        {{hash=0dcc317455baeefdf9b65140659af27d}{%
           family={Kissinger},
           familyi={K\bibinitperiod},
           given={Aleks},
           giveni={A\bibinitperiod}}}%
      }
      \strng{namehash}{0dcc317455baeefdf9b65140659af27d}
      \strng{fullhash}{0dcc317455baeefdf9b65140659af27d}
      \strng{fullhashraw}{0dcc317455baeefdf9b65140659af27d}
      \strng{bibnamehash}{0dcc317455baeefdf9b65140659af27d}
      \strng{authorbibnamehash}{0dcc317455baeefdf9b65140659af27d}
      \strng{authornamehash}{0dcc317455baeefdf9b65140659af27d}
      \strng{authorfullhash}{0dcc317455baeefdf9b65140659af27d}
      \strng{authorfullhashraw}{0dcc317455baeefdf9b65140659af27d}
      \field{sortinit}{K}
      \field{sortinithash}{c02bf6bff1c488450c352b40f5d853ab}
      \field{labelnamesource}{author}
      \field{labeltitlesource}{title}
      \field{eprinttype}{arXiv}
      \field{title}{Abstract Tensor Systems as Monoidal Categories}
      \field{year}{2013}
      \verb{eprint}
      \verb 1308.3586
      \endverb
    \endentry
    \entry{kleppmann22}{online}{}{}
      \name{author}{1}{}{%
        {{hash=34b8d5b1d29d175527a6642f3ec8a76d}{%
           family={Kleppmann},
           familyi={K\bibinitperiod},
           given={Martin},
           giveni={M\bibinitperiod}}}%
      }
      \strng{namehash}{34b8d5b1d29d175527a6642f3ec8a76d}
      \strng{fullhash}{34b8d5b1d29d175527a6642f3ec8a76d}
      \strng{fullhashraw}{34b8d5b1d29d175527a6642f3ec8a76d}
      \strng{bibnamehash}{34b8d5b1d29d175527a6642f3ec8a76d}
      \strng{authorbibnamehash}{34b8d5b1d29d175527a6642f3ec8a76d}
      \strng{authornamehash}{34b8d5b1d29d175527a6642f3ec8a76d}
      \strng{authorfullhash}{34b8d5b1d29d175527a6642f3ec8a76d}
      \strng{authorfullhashraw}{34b8d5b1d29d175527a6642f3ec8a76d}
      \field{extraname}{1}
      \field{sortinit}{K}
      \field{sortinithash}{c02bf6bff1c488450c352b40f5d853ab}
      \field{labelnamesource}{author}
      \field{labeltitlesource}{title}
      \field{note}{9th Workshop on Principles and Practice of Consistency for Distributed Data (PaPoC'22), Rennes, France. ACM. doi:10.1145/3517209.3524042}
      \field{title}{Making CRDTs Byzantine Fault Tolerant}
      \field{year}{2022}
      \verb{urlraw}
      \verb https://martin.kleppmann.com/papers/bft-crdt-papoc22.pdf
      \endverb
      \verb{url}
      \verb https://martin.kleppmann.com/papers/bft-crdt-papoc22.pdf
      \endverb
    \endentry
    \entry{kleppmann21}{online}{}{}
      \name{author}{1}{}{%
        {{hash=34b8d5b1d29d175527a6642f3ec8a76d}{%
           family={Kleppmann},
           familyi={K\bibinitperiod},
           given={Martin},
           giveni={M\bibinitperiod}}}%
      }
      \strng{namehash}{34b8d5b1d29d175527a6642f3ec8a76d}
      \strng{fullhash}{34b8d5b1d29d175527a6642f3ec8a76d}
      \strng{fullhashraw}{34b8d5b1d29d175527a6642f3ec8a76d}
      \strng{bibnamehash}{34b8d5b1d29d175527a6642f3ec8a76d}
      \strng{authorbibnamehash}{34b8d5b1d29d175527a6642f3ec8a76d}
      \strng{authornamehash}{34b8d5b1d29d175527a6642f3ec8a76d}
      \strng{authorfullhash}{34b8d5b1d29d175527a6642f3ec8a76d}
      \strng{authorfullhashraw}{34b8d5b1d29d175527a6642f3ec8a76d}
      \field{extraname}{2}
      \field{sortinit}{K}
      \field{sortinithash}{c02bf6bff1c488450c352b40f5d853ab}
      \field{labelnamesource}{author}
      \field{labeltitlesource}{title}
      \field{note}{Keynote, 15th ACM International Conference on Distributed and Event-Based Systems (DEBS '21), Virtual Event, Italy. ACM. doi:10.1145/3465480.3467835}
      \field{subtitle}{From Databases to Distributed Collaboration Software}
      \field{title}{Thinking in Events}
      \field{year}{2021}
      \verb{urlraw}
      \verb https://martin.kleppmann.com/papers/debs21-keynote.pdf
      \endverb
      \verb{url}
      \verb https://martin.kleppmann.com/papers/debs21-keynote.pdf
      \endverb
    \endentry
    \entry{kock20}{online}{}{}
      \name{author}{1}{}{%
        {{hash=14478e69f3dd724f5a664d54b7ecbb5f}{%
           family={Kock},
           familyi={K\bibinitperiod},
           given={Joachim},
           giveni={J\bibinitperiod}}}%
      }
      \strng{namehash}{14478e69f3dd724f5a664d54b7ecbb5f}
      \strng{fullhash}{14478e69f3dd724f5a664d54b7ecbb5f}
      \strng{fullhashraw}{14478e69f3dd724f5a664d54b7ecbb5f}
      \strng{bibnamehash}{14478e69f3dd724f5a664d54b7ecbb5f}
      \strng{authorbibnamehash}{14478e69f3dd724f5a664d54b7ecbb5f}
      \strng{authornamehash}{14478e69f3dd724f5a664d54b7ecbb5f}
      \strng{authorfullhash}{14478e69f3dd724f5a664d54b7ecbb5f}
      \strng{authorfullhashraw}{14478e69f3dd724f5a664d54b7ecbb5f}
      \field{sortinit}{K}
      \field{sortinithash}{c02bf6bff1c488450c352b40f5d853ab}
      \field{labelnamesource}{author}
      \field{labeltitlesource}{title}
      \field{eprinttype}{arXiv}
      \field{title}{Whole-grain Petri nets and processes}
      \field{year}{2020}
      \verb{eprint}
      \verb 2005.05108
      \endverb
    \endentry
    \entry{koller-friedman-2009}{book}{}{}
      \name{author}{2}{}{%
        {{hash=9c785bcf6b8a0b99e44695299a6aecb2}{%
           family={Koller},
           familyi={K\bibinitperiod},
           given={Daphne},
           giveni={D\bibinitperiod}}}%
        {{hash=987ee418ae32e261a8a5385e14cda479}{%
           family={Friedman},
           familyi={F\bibinitperiod},
           given={Nir},
           giveni={N\bibinitperiod}}}%
      }
      \list{publisher}{1}{%
        {MIT Press}%
      }
      \strng{namehash}{644f618617c37fd03b6bc0248cd308d5}
      \strng{fullhash}{644f618617c37fd03b6bc0248cd308d5}
      \strng{fullhashraw}{644f618617c37fd03b6bc0248cd308d5}
      \strng{bibnamehash}{644f618617c37fd03b6bc0248cd308d5}
      \strng{authorbibnamehash}{644f618617c37fd03b6bc0248cd308d5}
      \strng{authornamehash}{644f618617c37fd03b6bc0248cd308d5}
      \strng{authorfullhash}{644f618617c37fd03b6bc0248cd308d5}
      \strng{authorfullhashraw}{644f618617c37fd03b6bc0248cd308d5}
      \field{sortinit}{K}
      \field{sortinithash}{c02bf6bff1c488450c352b40f5d853ab}
      \field{labelnamesource}{author}
      \field{labeltitlesource}{title}
      \field{title}{Probabilistic Graphical Models: Principles and Techniques}
      \field{year}{2009}
    \endentry
    \entry{kopp20}{article}{}{}
      \name{author}{4}{}{%
        {{hash=1793768c6f2fc2b63c17147c6248271d}{%
           family={Kopp},
           familyi={K\bibinitperiod},
           given={Denny},
           giveni={D\bibinitperiod}}}%
        {{hash=818df49a5731d02e1699af8813165a6c}{%
           family={Hassoun},
           familyi={H\bibinitperiod},
           given={Michael},
           giveni={M\bibinitperiod}}}%
        {{hash=6c1a4c1c827637077c47c6bf97aa616d}{%
           family={Kalir},
           familyi={K\bibinitperiod},
           given={Adar},
           giveni={A\bibinitperiod}}}%
        {{hash=e882495f744105dbaa04ab63e104696c}{%
           family={Mönch},
           familyi={M\bibinitperiod},
           given={Lars},
           giveni={L\bibinitperiod}}}%
      }
      \strng{namehash}{1934d662096f3560e4602bd546db0848}
      \strng{fullhash}{1934d662096f3560e4602bd546db0848}
      \strng{fullhashraw}{1934d662096f3560e4602bd546db0848}
      \strng{bibnamehash}{1934d662096f3560e4602bd546db0848}
      \strng{authorbibnamehash}{1934d662096f3560e4602bd546db0848}
      \strng{authornamehash}{1934d662096f3560e4602bd546db0848}
      \strng{authorfullhash}{1934d662096f3560e4602bd546db0848}
      \strng{authorfullhashraw}{1934d662096f3560e4602bd546db0848}
      \field{sortinit}{K}
      \field{sortinithash}{c02bf6bff1c488450c352b40f5d853ab}
      \field{labelnamesource}{author}
      \field{labeltitlesource}{title}
      \field{journaltitle}{IEEE Transactions on Semiconductor Manufacturing}
      \field{number}{4}
      \field{title}{{SMT2020}---A Semiconductor Manufacturing Testbed}
      \field{volume}{33}
      \field{year}{2020}
      \field{pages}{522\bibrangedash 531}
      \range{pages}{10}
      \verb{doi}
      \verb 10.1109/TSM.2020.3001933
      \endverb
    \endentry
    \entry{kosoy20}{online}{}{}
      \name{author}{2}{}{%
        {{hash=03cbe17e5cb93e431dd64c549a767a8b}{%
           family={Kosoy},
           familyi={K\bibinitperiod},
           given={Vanessa},
           giveni={V\bibinitperiod}}}%
        {{hash=47b3e2dc3b3513cc1185c469d0c1ed7b}{%
           family={Appel},
           familyi={A\bibinitperiod},
           given={Alex},
           giveni={A\bibinitperiod}}}%
      }
      \strng{namehash}{2fe4a4343b629a829e4188be0723cceb}
      \strng{fullhash}{2fe4a4343b629a829e4188be0723cceb}
      \strng{fullhashraw}{2fe4a4343b629a829e4188be0723cceb}
      \strng{bibnamehash}{2fe4a4343b629a829e4188be0723cceb}
      \strng{authorbibnamehash}{2fe4a4343b629a829e4188be0723cceb}
      \strng{authornamehash}{2fe4a4343b629a829e4188be0723cceb}
      \strng{authorfullhash}{2fe4a4343b629a829e4188be0723cceb}
      \strng{authorfullhashraw}{2fe4a4343b629a829e4188be0723cceb}
      \field{sortinit}{K}
      \field{sortinithash}{c02bf6bff1c488450c352b40f5d853ab}
      \field{labelnamesource}{author}
      \field{labeltitlesource}{title}
      \field{note}{Sequence on the AI Alignment Forum; Appel posts under the handle ``Diffractor''}
      \field{title}{Infra-Bayesianism}
      \field{urlday}{21}
      \field{urlmonth}{7}
      \field{urlyear}{2026}
      \field{year}{2020}
      \field{urldateera}{ce}
      \verb{urlraw}
      \verb https://www.alignmentforum.org/s/CmrW8fCmSLK7E25sa
      \endverb
      \verb{url}
      \verb https://www.alignmentforum.org/s/CmrW8fCmSLK7E25sa
      \endverb
    \endentry
    \entry{krivine1964}{article}{}{}
      \name{author}{1}{}{%
        {{hash=c6aef8513e7ea47bd8e331a73c91f861}{%
           family={Krivine},
           familyi={K\bibinitperiod},
           given={J.-L.},
           giveni={J\bibinithyphendelim L\bibinitperiod}}}%
      }
      \strng{namehash}{c6aef8513e7ea47bd8e331a73c91f861}
      \strng{fullhash}{c6aef8513e7ea47bd8e331a73c91f861}
      \strng{fullhashraw}{c6aef8513e7ea47bd8e331a73c91f861}
      \strng{bibnamehash}{c6aef8513e7ea47bd8e331a73c91f861}
      \strng{authorbibnamehash}{c6aef8513e7ea47bd8e331a73c91f861}
      \strng{authornamehash}{c6aef8513e7ea47bd8e331a73c91f861}
      \strng{authorfullhash}{c6aef8513e7ea47bd8e331a73c91f861}
      \strng{authorfullhashraw}{c6aef8513e7ea47bd8e331a73c91f861}
      \field{sortinit}{K}
      \field{sortinithash}{c02bf6bff1c488450c352b40f5d853ab}
      \field{labelnamesource}{author}
      \field{labeltitlesource}{title}
      \field{journaltitle}{Journal d'Analyse Mathématique}
      \field{title}{Anneaux préordonnés}
      \field{volume}{12}
      \field{year}{1964}
      \field{pages}{307\bibrangedash 326}
      \range{pages}{20}
    \endentry
    \entry{kumar1990dynamic}{article}{}{}
      \name{author}{2}{}{%
        {{hash=d5f5decc1fe064c502f05b5bf93b7bda}{%
           family={Kumar},
           familyi={K\bibinitperiod},
           given={P.\bibnamedelimi R.},
           giveni={P\bibinitperiod\bibinitdelim R\bibinitperiod}}}%
        {{hash=5e8aa03924c7107f0d1d84829fab5347}{%
           family={Seidman},
           familyi={S\bibinitperiod},
           given={Thomas\bibnamedelima I.},
           giveni={T\bibinitperiod\bibinitdelim I\bibinitperiod}}}%
      }
      \strng{namehash}{e83dfc4e21b73619628b70f7bdc30d37}
      \strng{fullhash}{e83dfc4e21b73619628b70f7bdc30d37}
      \strng{fullhashraw}{e83dfc4e21b73619628b70f7bdc30d37}
      \strng{bibnamehash}{e83dfc4e21b73619628b70f7bdc30d37}
      \strng{authorbibnamehash}{e83dfc4e21b73619628b70f7bdc30d37}
      \strng{authornamehash}{e83dfc4e21b73619628b70f7bdc30d37}
      \strng{authorfullhash}{e83dfc4e21b73619628b70f7bdc30d37}
      \strng{authorfullhashraw}{e83dfc4e21b73619628b70f7bdc30d37}
      \field{sortinit}{K}
      \field{sortinithash}{c02bf6bff1c488450c352b40f5d853ab}
      \field{labelnamesource}{author}
      \field{labeltitlesource}{title}
      \field{journaltitle}{IEEE Transactions on Automatic Control}
      \field{number}{3}
      \field{title}{Dynamic instabilities and stabilization methods in distributed real-time scheduling of manufacturing systems}
      \field{volume}{35}
      \field{year}{1990}
      \field{pages}{289\bibrangedash 298}
      \range{pages}{10}
      \verb{doi}
      \verb 10.1109/9.50339
      \endverb
    \endentry
    \entry{kumar14cakeml}{inproceedings}{}{}
      \name{author}{4}{}{%
        {{hash=bc70ef110d75b816e145c90b654d84eb}{%
           family={Kumar},
           familyi={K\bibinitperiod},
           given={Ramana},
           giveni={R\bibinitperiod}}}%
        {{hash=252a9f312905933ebf5c567099ef9b22}{%
           family={Myreen},
           familyi={M\bibinitperiod},
           given={Magnus\bibnamedelima O.},
           giveni={M\bibinitperiod\bibinitdelim O\bibinitperiod}}}%
        {{hash=5f4f59c4ab0dd44df8053b075d95507e}{%
           family={Norrish},
           familyi={N\bibinitperiod},
           given={Michael},
           giveni={M\bibinitperiod}}}%
        {{hash=04b670f6d1f73db0f24b85e5b4562755}{%
           family={Owens},
           familyi={O\bibinitperiod},
           given={Scott},
           giveni={S\bibinitperiod}}}%
      }
      \strng{namehash}{66b580108ae317a12dfba92026dcc31e}
      \strng{fullhash}{66b580108ae317a12dfba92026dcc31e}
      \strng{fullhashraw}{66b580108ae317a12dfba92026dcc31e}
      \strng{bibnamehash}{66b580108ae317a12dfba92026dcc31e}
      \strng{authorbibnamehash}{66b580108ae317a12dfba92026dcc31e}
      \strng{authornamehash}{66b580108ae317a12dfba92026dcc31e}
      \strng{authorfullhash}{66b580108ae317a12dfba92026dcc31e}
      \strng{authorfullhashraw}{66b580108ae317a12dfba92026dcc31e}
      \field{sortinit}{K}
      \field{sortinithash}{c02bf6bff1c488450c352b40f5d853ab}
      \field{labelnamesource}{author}
      \field{labeltitlesource}{title}
      \field{booktitle}{Proceedings of the 41st {ACM} {SIGPLAN}-{SIGACT} Symposium on Principles of Programming Languages ({POPL} '14)}
      \field{title}{{CakeML}: A Verified Implementation of {ML}}
      \field{year}{2014}
      \field{pages}{179\bibrangedash 191}
      \range{pages}{13}
      \verb{doi}
      \verb 10.1145/2535838.2535841
      \endverb
    \endentry
    \entry{kushner1967}{book}{}{}
      \name{author}{1}{}{%
        {{hash=1435729142d7ff2b491681a967128ce5}{%
           family={Kushner},
           familyi={K\bibinitperiod},
           given={H.\bibnamedelimi J.},
           giveni={H\bibinitperiod\bibinitdelim J\bibinitperiod}}}%
      }
      \list{location}{1}{%
        {New York}%
      }
      \list{publisher}{1}{%
        {Academic Press}%
      }
      \strng{namehash}{1435729142d7ff2b491681a967128ce5}
      \strng{fullhash}{1435729142d7ff2b491681a967128ce5}
      \strng{fullhashraw}{1435729142d7ff2b491681a967128ce5}
      \strng{bibnamehash}{1435729142d7ff2b491681a967128ce5}
      \strng{authorbibnamehash}{1435729142d7ff2b491681a967128ce5}
      \strng{authornamehash}{1435729142d7ff2b491681a967128ce5}
      \strng{authorfullhash}{1435729142d7ff2b491681a967128ce5}
      \strng{authorfullhashraw}{1435729142d7ff2b491681a967128ce5}
      \field{sortinit}{K}
      \field{sortinithash}{c02bf6bff1c488450c352b40f5d853ab}
      \field{labelnamesource}{author}
      \field{labeltitlesource}{title}
      \field{note}{(Supermartingale safety certificates for stochastic systems.)}
      \field{number}{33}
      \field{series}{Mathematics in Science and Engineering}
      \field{title}{Stochastic Stability and Control}
      \field{year}{1967}
    \endentry
    \entry{kwiatkowska2011prism}{incollection}{}{}
      \name{author}{3}{}{%
        {{hash=e96e14c940cd49f1c4e8c1f38f89a321}{%
           family={Kwiatkowska},
           familyi={K\bibinitperiod},
           given={Marta},
           giveni={M\bibinitperiod}}}%
        {{hash=72e38f5e745acabaeb10e38ca01f1a52}{%
           family={Norman},
           familyi={N\bibinitperiod},
           given={Gethin},
           giveni={G\bibinitperiod}}}%
        {{hash=85aa04746604480ce47f280e167052ff}{%
           family={Parker},
           familyi={P\bibinitperiod},
           given={David},
           giveni={D\bibinitperiod}}}%
      }
      \list{publisher}{1}{%
        {Springer}%
      }
      \strng{namehash}{84b38c4b40e425cbfea3cd18ab69194d}
      \strng{fullhash}{84b38c4b40e425cbfea3cd18ab69194d}
      \strng{fullhashraw}{84b38c4b40e425cbfea3cd18ab69194d}
      \strng{bibnamehash}{84b38c4b40e425cbfea3cd18ab69194d}
      \strng{authorbibnamehash}{84b38c4b40e425cbfea3cd18ab69194d}
      \strng{authornamehash}{84b38c4b40e425cbfea3cd18ab69194d}
      \strng{authorfullhash}{84b38c4b40e425cbfea3cd18ab69194d}
      \strng{authorfullhashraw}{84b38c4b40e425cbfea3cd18ab69194d}
      \field{sortinit}{K}
      \field{sortinithash}{c02bf6bff1c488450c352b40f5d853ab}
      \field{labelnamesource}{author}
      \field{labeltitlesource}{title}
      \field{booktitle}{Computer Aided Verification ({CAV})}
      \field{series}{Lecture Notes in Computer Science}
      \field{title}{{PRISM} 4.0: Verification of Probabilistic Real-Time Systems}
      \field{volume}{6806}
      \field{year}{2011}
      \field{pages}{585\bibrangedash 591}
      \range{pages}{7}
      \verb{doi}
      \verb 10.1007/978-3-642-22110-1_47
      \endverb
    \endentry
    \entry{lambert21a}{online}{}{}
      \name{author}{1}{}{%
        {{hash=26775461fd456edc4deb5b983123ea7c}{%
           family={Lambert},
           familyi={L\bibinitperiod},
           given={Michael},
           giveni={M\bibinitperiod}}}%
      }
      \strng{namehash}{26775461fd456edc4deb5b983123ea7c}
      \strng{fullhash}{26775461fd456edc4deb5b983123ea7c}
      \strng{fullhashraw}{26775461fd456edc4deb5b983123ea7c}
      \strng{bibnamehash}{26775461fd456edc4deb5b983123ea7c}
      \strng{authorbibnamehash}{26775461fd456edc4deb5b983123ea7c}
      \strng{authornamehash}{26775461fd456edc4deb5b983123ea7c}
      \strng{authorfullhash}{26775461fd456edc4deb5b983123ea7c}
      \strng{authorfullhashraw}{26775461fd456edc4deb5b983123ea7c}
      \field{sortinit}{L}
      \field{sortinithash}{7c47d417cecb1f4bd38d1825c427a61a}
      \field{labelnamesource}{author}
      \field{labeltitlesource}{title}
      \field{eprinttype}{arXiv}
      \field{title}{Discrete Double Fibrations}
      \field{year}{2021}
      \verb{eprint}
      \verb 2101.06734
      \endverb
    \endentry
    \entry{landweber1969}{article}{}{}
      \name{author}{1}{}{%
        {{hash=1592d50fe7b3c3abf050ea84179545dc}{%
           family={Landweber},
           familyi={L\bibinitperiod},
           given={L.\bibnamedelimi H.},
           giveni={L\bibinitperiod\bibinitdelim H\bibinitperiod}}}%
      }
      \strng{namehash}{1592d50fe7b3c3abf050ea84179545dc}
      \strng{fullhash}{1592d50fe7b3c3abf050ea84179545dc}
      \strng{fullhashraw}{1592d50fe7b3c3abf050ea84179545dc}
      \strng{bibnamehash}{1592d50fe7b3c3abf050ea84179545dc}
      \strng{authorbibnamehash}{1592d50fe7b3c3abf050ea84179545dc}
      \strng{authornamehash}{1592d50fe7b3c3abf050ea84179545dc}
      \strng{authorfullhash}{1592d50fe7b3c3abf050ea84179545dc}
      \strng{authorfullhashraw}{1592d50fe7b3c3abf050ea84179545dc}
      \field{sortinit}{L}
      \field{sortinithash}{7c47d417cecb1f4bd38d1825c427a61a}
      \field{labelnamesource}{author}
      \field{labeltitlesource}{title}
      \field{journaltitle}{Mathematical Systems Theory}
      \field{note}{(Deterministic Büchi automata recognize exactly the limit languages.)}
      \field{number}{4}
      \field{title}{Decision problems for $\omega$-automata}
      \field{volume}{3}
      \field{year}{1969}
      \field{pages}{376\bibrangedash 384}
      \range{pages}{9}
    \endentry
    \entry{laneriesenfeld1980}{article}{}{}
      \name{author}{2}{}{%
        {{hash=0a0f6a079bb580cfd4edf41dac7f8a5f}{%
           family={Lane},
           familyi={L\bibinitperiod},
           given={J.\bibnamedelimi M.},
           giveni={J\bibinitperiod\bibinitdelim M\bibinitperiod}}}%
        {{hash=23bc22741a8087c1e834bf2cb96f9290}{%
           family={Riesenfeld},
           familyi={R\bibinitperiod},
           given={R.\bibnamedelimi F.},
           giveni={R\bibinitperiod\bibinitdelim F\bibinitperiod}}}%
      }
      \strng{namehash}{029f3422b2e06d7e0eabe5e3d8138ab7}
      \strng{fullhash}{029f3422b2e06d7e0eabe5e3d8138ab7}
      \strng{fullhashraw}{029f3422b2e06d7e0eabe5e3d8138ab7}
      \strng{bibnamehash}{029f3422b2e06d7e0eabe5e3d8138ab7}
      \strng{authorbibnamehash}{029f3422b2e06d7e0eabe5e3d8138ab7}
      \strng{authornamehash}{029f3422b2e06d7e0eabe5e3d8138ab7}
      \strng{authorfullhash}{029f3422b2e06d7e0eabe5e3d8138ab7}
      \strng{authorfullhashraw}{029f3422b2e06d7e0eabe5e3d8138ab7}
      \field{sortinit}{L}
      \field{sortinithash}{7c47d417cecb1f4bd38d1825c427a61a}
      \field{labelnamesource}{author}
      \field{labeltitlesource}{title}
      \field{journaltitle}{IEEE Transactions on Pattern Analysis and Machine Intelligence}
      \field{number}{1}
      \field{title}{A theoretical development for the computer generation and display of piecewise polynomial surfaces}
      \field{volume}{2}
      \field{year}{1980}
      \field{pages}{35\bibrangedash 46}
      \range{pages}{12}
    \endentry
    \entry{laurent22}{online}{}{}
      \name{author}{2}{}{%
        {{hash=808513bc2b23b1fee6a93e13ee2c2252}{%
           family={Laurent},
           familyi={L\bibinitperiod},
           given={Jonathan},
           giveni={J\bibinitperiod}}}%
        {{hash=0fb7e4129542e3e84c6115bee2a2b227}{%
           family={Platzer},
           familyi={P\bibinitperiod},
           given={André},
           giveni={A\bibinitperiod}}}%
      }
      \strng{namehash}{7017e09e7367aaf5ed4009f23dd7273b}
      \strng{fullhash}{7017e09e7367aaf5ed4009f23dd7273b}
      \strng{fullhashraw}{7017e09e7367aaf5ed4009f23dd7273b}
      \strng{bibnamehash}{7017e09e7367aaf5ed4009f23dd7273b}
      \strng{authorbibnamehash}{7017e09e7367aaf5ed4009f23dd7273b}
      \strng{authornamehash}{7017e09e7367aaf5ed4009f23dd7273b}
      \strng{authorfullhash}{7017e09e7367aaf5ed4009f23dd7273b}
      \strng{authorfullhashraw}{7017e09e7367aaf5ed4009f23dd7273b}
      \field{extraname}{1}
      \field{sortinit}{L}
      \field{sortinithash}{7c47d417cecb1f4bd38d1825c427a61a}
      \field{labelnamesource}{author}
      \field{labeltitlesource}{title}
      \field{eprinttype}{arXiv}
      \field{subtitle}{The Case of Loop Invariant Synthesis}
      \field{title}{Learning to Find Proofs and Theorems by Learning to Refine Search Strategies}
      \field{year}{2022}
      \verb{eprint}
      \verb 2205.14229
      \endverb
    \endentry
    \entry{laurent25}{online}{}{}
      \name{author}{2}{}{%
        {{hash=808513bc2b23b1fee6a93e13ee2c2252}{%
           family={Laurent},
           familyi={L\bibinitperiod},
           given={Jonathan},
           giveni={J\bibinitperiod}}}%
        {{hash=0fb7e4129542e3e84c6115bee2a2b227}{%
           family={Platzer},
           familyi={P\bibinitperiod},
           given={André},
           giveni={A\bibinitperiod}}}%
      }
      \strng{namehash}{7017e09e7367aaf5ed4009f23dd7273b}
      \strng{fullhash}{7017e09e7367aaf5ed4009f23dd7273b}
      \strng{fullhashraw}{7017e09e7367aaf5ed4009f23dd7273b}
      \strng{bibnamehash}{7017e09e7367aaf5ed4009f23dd7273b}
      \strng{authorbibnamehash}{7017e09e7367aaf5ed4009f23dd7273b}
      \strng{authornamehash}{7017e09e7367aaf5ed4009f23dd7273b}
      \strng{authorfullhash}{7017e09e7367aaf5ed4009f23dd7273b}
      \strng{authorfullhashraw}{7017e09e7367aaf5ed4009f23dd7273b}
      \field{extraname}{2}
      \field{sortinit}{L}
      \field{sortinithash}{7c47d417cecb1f4bd38d1825c427a61a}
      \field{labelnamesource}{author}
      \field{labeltitlesource}{title}
      \field{eprinttype}{arXiv}
      \field{subtitle}{A Modular Foundation for Building LLM-Enabled Software}
      \field{title}{Oracular Programming}
      \field{year}{2025}
      \verb{eprint}
      \verb 2502.05310
      \endverb
    \endentry
    \entry{lavore22}{online}{}{}
      \name{author}{3}{}{%
        {{hash=9ff65ac1b5febfa5f78fd90506f0ac93}{%
           family={Lavore},
           familyi={L\bibinitperiod},
           given={Elena\bibnamedelima Di},
           giveni={E\bibinitperiod\bibinitdelim D\bibinitperiod}}}%
        {{hash=90a608e5aa8c815f07e4c59d5b93be62}{%
           family={Felice},
           familyi={F\bibinitperiod},
           given={Giovanni},
           giveni={G\bibinitperiod},
           prefix={de},
           prefixi={d\bibinitperiod}}}%
        {{hash=d4f092c73ead809932c43d0668067f0a}{%
           family={Román},
           familyi={R\bibinitperiod},
           given={Mario},
           giveni={M\bibinitperiod}}}%
      }
      \strng{namehash}{6141f7db70b34eaadb39fa0e8ceb8a31}
      \strng{fullhash}{6141f7db70b34eaadb39fa0e8ceb8a31}
      \strng{fullhashraw}{6141f7db70b34eaadb39fa0e8ceb8a31}
      \strng{bibnamehash}{6141f7db70b34eaadb39fa0e8ceb8a31}
      \strng{authorbibnamehash}{6141f7db70b34eaadb39fa0e8ceb8a31}
      \strng{authornamehash}{6141f7db70b34eaadb39fa0e8ceb8a31}
      \strng{authorfullhash}{6141f7db70b34eaadb39fa0e8ceb8a31}
      \strng{authorfullhashraw}{6141f7db70b34eaadb39fa0e8ceb8a31}
      \field{sortinit}{L}
      \field{sortinithash}{7c47d417cecb1f4bd38d1825c427a61a}
      \field{labelnamesource}{author}
      \field{labeltitlesource}{title}
      \field{eprinttype}{arXiv}
      \field{title}{Monoidal Streams for Dataflow Programming}
      \field{year}{2022}
      \verb{eprint}
      \verb 2202.02061
      \endverb
    \endentry
    \entry{lawvere69}{article}{}{}
      \name{author}{1}{}{%
        {{hash=7f4e82b361c38ad671615c08cb0678cd}{%
           family={Lawvere},
           familyi={L\bibinitperiod},
           given={F.\bibnamedelimi William},
           giveni={F\bibinitperiod\bibinitdelim W\bibinitperiod}}}%
      }
      \strng{namehash}{7f4e82b361c38ad671615c08cb0678cd}
      \strng{fullhash}{7f4e82b361c38ad671615c08cb0678cd}
      \strng{fullhashraw}{7f4e82b361c38ad671615c08cb0678cd}
      \strng{bibnamehash}{7f4e82b361c38ad671615c08cb0678cd}
      \strng{authorbibnamehash}{7f4e82b361c38ad671615c08cb0678cd}
      \strng{authornamehash}{7f4e82b361c38ad671615c08cb0678cd}
      \strng{authorfullhash}{7f4e82b361c38ad671615c08cb0678cd}
      \strng{authorfullhashraw}{7f4e82b361c38ad671615c08cb0678cd}
      \field{extraname}{1}
      \field{sortinit}{L}
      \field{sortinithash}{7c47d417cecb1f4bd38d1825c427a61a}
      \field{labelnamesource}{author}
      \field{labeltitlesource}{title}
      \field{journaltitle}{Dialectica}
      \field{title}{Adjointness in Foundations}
      \field{volume}{23}
      \field{year}{1969}
      \field{pages}{281\bibrangedash 296}
      \range{pages}{16}
    \endentry
    \entry{lawvere70}{inproceedings}{}{}
      \name{author}{1}{}{%
        {{hash=7f4e82b361c38ad671615c08cb0678cd}{%
           family={Lawvere},
           familyi={L\bibinitperiod},
           given={F.\bibnamedelimi William},
           giveni={F\bibinitperiod\bibinitdelim W\bibinitperiod}}}%
      }
      \list{publisher}{1}{%
        {American Mathematical Society}%
      }
      \strng{namehash}{7f4e82b361c38ad671615c08cb0678cd}
      \strng{fullhash}{7f4e82b361c38ad671615c08cb0678cd}
      \strng{fullhashraw}{7f4e82b361c38ad671615c08cb0678cd}
      \strng{bibnamehash}{7f4e82b361c38ad671615c08cb0678cd}
      \strng{authorbibnamehash}{7f4e82b361c38ad671615c08cb0678cd}
      \strng{authornamehash}{7f4e82b361c38ad671615c08cb0678cd}
      \strng{authorfullhash}{7f4e82b361c38ad671615c08cb0678cd}
      \strng{authorfullhashraw}{7f4e82b361c38ad671615c08cb0678cd}
      \field{extraname}{2}
      \field{sortinit}{L}
      \field{sortinithash}{7c47d417cecb1f4bd38d1825c427a61a}
      \field{labelnamesource}{author}
      \field{labeltitlesource}{title}
      \field{booktitle}{Applications of Categorical Algebra}
      \field{series}{Proceedings of Symposia in Pure Mathematics}
      \field{title}{Equality in Hyperdoctrines and Comprehension Schema as an Adjoint Functor}
      \field{volume}{17}
      \field{year}{1970}
      \field{pages}{1\bibrangedash 14}
      \range{pages}{14}
    \endentry
    \entry{lawvere73}{article}{}{}
      \name{author}{1}{}{%
        {{hash=7f4e82b361c38ad671615c08cb0678cd}{%
           family={Lawvere},
           familyi={L\bibinitperiod},
           given={F.\bibnamedelimi William},
           giveni={F\bibinitperiod\bibinitdelim W\bibinitperiod}}}%
      }
      \strng{namehash}{7f4e82b361c38ad671615c08cb0678cd}
      \strng{fullhash}{7f4e82b361c38ad671615c08cb0678cd}
      \strng{fullhashraw}{7f4e82b361c38ad671615c08cb0678cd}
      \strng{bibnamehash}{7f4e82b361c38ad671615c08cb0678cd}
      \strng{authorbibnamehash}{7f4e82b361c38ad671615c08cb0678cd}
      \strng{authornamehash}{7f4e82b361c38ad671615c08cb0678cd}
      \strng{authorfullhash}{7f4e82b361c38ad671615c08cb0678cd}
      \strng{authorfullhashraw}{7f4e82b361c38ad671615c08cb0678cd}
      \field{extraname}{3}
      \field{sortinit}{L}
      \field{sortinithash}{7c47d417cecb1f4bd38d1825c427a61a}
      \field{labelnamesource}{author}
      \field{labeltitlesource}{title}
      \field{journaltitle}{Rendiconti del Seminario Matematico e Fisico di Milano}
      \field{note}{Reprinted in \emph{Reprints in Theory and Applications of Categories} 1 (2002), 1--37}
      \field{title}{Metric Spaces, Generalized Logic, and Closed Categories}
      \field{volume}{43}
      \field{year}{1973}
      \field{pages}{135\bibrangedash 166}
      \range{pages}{32}
    \endentry
    \entry{lee22}{inproceedings}{}{}
      \name{author}{4}{}{%
        {{hash=9f4d31adca922afa232108fb6992e5b0}{%
           family={Lee},
           familyi={L\bibinitperiod},
           given={Kwanwoo},
           giveni={K\bibinitperiod}}}%
        {{hash=4372c5e36ef57afd4f08e49344e53428}{%
           family={Song},
           familyi={S\bibinitperiod},
           given={Siyong},
           giveni={S\bibinitperiod}}}%
        {{hash=5b27f7032dcef9df25ab6c652a496e4b}{%
           family={Chang},
           familyi={C\bibinitperiod},
           given={Daesoon},
           giveni={D\bibinitperiod}}}%
        {{hash=34c6ca5e64fb3a755bda69089c0dfd65}{%
           family={Park},
           familyi={P\bibinitperiod},
           given={Sangchul},
           giveni={S\bibinitperiod}}}%
      }
      \strng{namehash}{4b69c00e01f4353afe544481ebbb6bc1}
      \strng{fullhash}{4b69c00e01f4353afe544481ebbb6bc1}
      \strng{fullhashraw}{4b69c00e01f4353afe544481ebbb6bc1}
      \strng{bibnamehash}{4b69c00e01f4353afe544481ebbb6bc1}
      \strng{authorbibnamehash}{4b69c00e01f4353afe544481ebbb6bc1}
      \strng{authornamehash}{4b69c00e01f4353afe544481ebbb6bc1}
      \strng{authorfullhash}{4b69c00e01f4353afe544481ebbb6bc1}
      \strng{authorfullhashraw}{4b69c00e01f4353afe544481ebbb6bc1}
      \field{sortinit}{L}
      \field{sortinithash}{7c47d417cecb1f4bd38d1825c427a61a}
      \field{labelnamesource}{author}
      \field{labeltitlesource}{title}
      \field{booktitle}{Proceedings of the 2022 Winter Simulation Conference}
      \field{title}{A New {AMHS} Testbed for Semiconductor Manufacturing}
      \field{year}{2022}
      \field{pages}{3318\bibrangedash 3325}
      \range{pages}{8}
      \verb{doi}
      \verb 10.1109/WSC57314.2022.10015476
      \endverb
    \endentry
    \entry{letourneau24padre}{misc}{}{}
      \name{author}{9}{}{%
        {{hash=94bb970f18561fa7c1c50ff981c9e7a5}{%
           family={Letourneau},
           familyi={L\bibinitperiod},
           given={Pierre-David},
           giveni={P\bibinithyphendelim D\bibinitperiod}}}%
        {{hash=5b167d5f7257b7743e8233af22d35a7a}{%
           family={Singh},
           familyi={S\bibinitperiod},
           given={Manish\bibnamedelima Kumar},
           giveni={M\bibinitperiod\bibinitdelim K\bibinitperiod}}}%
        {{hash=6fb646e4ff3535262350a354006e7637}{%
           family={Cheng},
           familyi={C\bibinitperiod},
           given={Hsin-Pai},
           giveni={H\bibinithyphendelim P\bibinitperiod}}}%
        {{hash=a03d5b7381f4d27d28bb2b7ea8045a5e}{%
           family={Han},
           familyi={H\bibinitperiod},
           given={Shizhong},
           giveni={S\bibinitperiod}}}%
        {{hash=967bf7c7f514c973298658b72de2c264}{%
           family={Shi},
           familyi={S\bibinitperiod},
           given={Yunxiao},
           giveni={Y\bibinitperiod}}}%
        {{hash=4e445a83b7e1d0cfa0554187544c4cab}{%
           family={Jones},
           familyi={J\bibinitperiod},
           given={Dalton},
           giveni={D\bibinitperiod}}}%
        {{hash=d885c0c75b8c3fc1bae942938343117f}{%
           family={Langston},
           familyi={L\bibinitperiod},
           given={Matthew\bibnamedelima Harper},
           giveni={M\bibinitperiod\bibinitdelim H\bibinitperiod}}}%
        {{hash=d8ee433497a88adfa508dcb28f61f6eb}{%
           family={Cai},
           familyi={C\bibinitperiod},
           given={Hong},
           giveni={H\bibinitperiod}}}%
        {{hash=614198715471d146b77068d419298d26}{%
           family={Porikli},
           familyi={P\bibinitperiod},
           given={Fatih},
           giveni={F\bibinitperiod}}}%
      }
      \strng{namehash}{373911ba5242d4ee3ac051a9a333826d}
      \strng{fullhash}{246f70fd8eb1c60406372c2668ed712b}
      \strng{fullhashraw}{246f70fd8eb1c60406372c2668ed712b}
      \strng{bibnamehash}{373911ba5242d4ee3ac051a9a333826d}
      \strng{authorbibnamehash}{373911ba5242d4ee3ac051a9a333826d}
      \strng{authornamehash}{373911ba5242d4ee3ac051a9a333826d}
      \strng{authorfullhash}{246f70fd8eb1c60406372c2668ed712b}
      \strng{authorfullhashraw}{246f70fd8eb1c60406372c2668ed712b}
      \field{sortinit}{L}
      \field{sortinithash}{7c47d417cecb1f4bd38d1825c427a61a}
      \field{labelnamesource}{author}
      \field{labeltitlesource}{title}
      \field{eprinttype}{arxiv}
      \field{note}{(Degree-2--3 polynomial attention matches softmax accuracy at 11--43$\times$ speed; softmax itself only rational-approximable.)}
      \field{title}{{PADRe}: A Unifying Polynomial Attention Drop-in Replacement for Efficient Vision Transformer}
      \verb{eprint}
      \verb 2407.11306
      \endverb
    \endentry
    \entry{levinperes2017}{book}{}{}
      \name{author}{2}{}{%
        {{hash=198add63323ecec8d450c429589ac109}{%
           family={Levin},
           familyi={L\bibinitperiod},
           given={D.\bibnamedelimi A.},
           giveni={D\bibinitperiod\bibinitdelim A\bibinitperiod}}}%
        {{hash=40e0f66d18da16c33ac8b1a8b8e17036}{%
           family={Peres},
           familyi={P\bibinitperiod},
           given={Y.},
           giveni={Y\bibinitperiod}}}%
      }
      \list{publisher}{1}{%
        {American Mathematical Society}%
      }
      \strng{namehash}{427af43f72ff5465aebf175dc594b800}
      \strng{fullhash}{427af43f72ff5465aebf175dc594b800}
      \strng{fullhashraw}{427af43f72ff5465aebf175dc594b800}
      \strng{bibnamehash}{427af43f72ff5465aebf175dc594b800}
      \strng{authorbibnamehash}{427af43f72ff5465aebf175dc594b800}
      \strng{authornamehash}{427af43f72ff5465aebf175dc594b800}
      \strng{authorfullhash}{427af43f72ff5465aebf175dc594b800}
      \strng{authorfullhashraw}{427af43f72ff5465aebf175dc594b800}
      \field{sortinit}{L}
      \field{sortinithash}{7c47d417cecb1f4bd38d1825c427a61a}
      \field{labelnamesource}{author}
      \field{labeltitlesource}{title}
      \field{edition}{2}
      \field{note}{With E.~L.~Wilmer. (Total-variation distance, couplings, and the data-processing/projection inequality: Ch.~4--5.)}
      \field{title}{Markov Chains and Mixing Times}
      \field{year}{2017}
    \endentry
    \entry{li25}{inproceedings}{}{}
      \name{author}{2}{}{%
        {{hash=b40fe19de3efd505a11f3a4cd6b81abb}{%
           family={Li},
           familyi={L\bibinitperiod},
           given={Ziyun},
           giveni={Z\bibinitperiod}}}%
        {{hash=052bc4b0db8d1b625de206fbfe6a09e5}{%
           family={Boström},
           familyi={B\bibinitperiod},
           given={Henrik},
           giveni={H\bibinitperiod}}}%
      }
      \strng{namehash}{1fd70773306ad9472263a771ada5b7b7}
      \strng{fullhash}{1fd70773306ad9472263a771ada5b7b7}
      \strng{fullhashraw}{1fd70773306ad9472263a771ada5b7b7}
      \strng{bibnamehash}{1fd70773306ad9472263a771ada5b7b7}
      \strng{authorbibnamehash}{1fd70773306ad9472263a771ada5b7b7}
      \strng{authornamehash}{1fd70773306ad9472263a771ada5b7b7}
      \strng{authorfullhash}{1fd70773306ad9472263a771ada5b7b7}
      \strng{authorfullhashraw}{1fd70773306ad9472263a771ada5b7b7}
      \field{sortinit}{L}
      \field{sortinithash}{7c47d417cecb1f4bd38d1825c427a61a}
      \field{labelnamesource}{author}
      \field{labeltitlesource}{title}
      \field{booktitle}{Proceedings of the Fourteenth Symposium on Conformal and Probabilistic Prediction with Applications}
      \field{series}{Proceedings of Machine Learning Research}
      \field{title}{FlowGuard: Guarding Flow Matching via Conformal Sampling}
      \field{volume}{266}
      \field{year}{2025}
      \field{pages}{775\bibrangedash 777}
      \range{pages}{3}
    \endentry
    \entry{libkind21}{online}{}{}
      \name{author}{4}{}{%
        {{hash=33bda5bc38c3b05f23e9a99b5e4efb27}{%
           family={Libkind},
           familyi={L\bibinitperiod},
           given={Sophie},
           giveni={S\bibinitperiod}}}%
        {{hash=0f55cd9d4a858c1f9fb7749727796f02}{%
           family={Baas},
           familyi={B\bibinitperiod},
           given={Andrew},
           giveni={A\bibinitperiod}}}%
        {{hash=318b3f8d6da4d1c854033da07460a0cd}{%
           family={Patterson},
           familyi={P\bibinitperiod},
           given={Evan},
           giveni={E\bibinitperiod}}}%
        {{hash=1d30a8595a636bf14fcd181b6822212d}{%
           family={Fairbanks},
           familyi={F\bibinitperiod},
           given={James},
           giveni={J\bibinitperiod}}}%
      }
      \strng{namehash}{ab07c07deb359ab9b565d22da4d673bc}
      \strng{fullhash}{ab07c07deb359ab9b565d22da4d673bc}
      \strng{fullhashraw}{ab07c07deb359ab9b565d22da4d673bc}
      \strng{bibnamehash}{ab07c07deb359ab9b565d22da4d673bc}
      \strng{authorbibnamehash}{ab07c07deb359ab9b565d22da4d673bc}
      \strng{authornamehash}{ab07c07deb359ab9b565d22da4d673bc}
      \strng{authorfullhash}{ab07c07deb359ab9b565d22da4d673bc}
      \strng{authorfullhashraw}{ab07c07deb359ab9b565d22da4d673bc}
      \field{sortinit}{L}
      \field{sortinithash}{7c47d417cecb1f4bd38d1825c427a61a}
      \field{labelnamesource}{author}
      \field{labeltitlesource}{title}
      \field{eprinttype}{arXiv}
      \field{subtitle}{Mathematics and Computation}
      \field{title}{Operadic Modeling of Dynamical Systems}
      \field{year}{2021}
      \verb{eprint}
      \verb 2105.12282
      \endverb
    \endentry
    \entry{libkind25}{online}{}{}
      \name{author}{2}{}{%
        {{hash=33bda5bc38c3b05f23e9a99b5e4efb27}{%
           family={Libkind},
           familyi={L\bibinitperiod},
           given={Sophie},
           giveni={S\bibinitperiod}}}%
        {{hash=5e4483260e70411ee8f2ec4093d57437}{%
           family={Myers},
           familyi={M\bibinitperiod},
           given={David\bibnamedelima Jaz},
           giveni={D\bibinitperiod\bibinitdelim J\bibinitperiod}}}%
      }
      \strng{namehash}{bfe21ee39369b5ca36f289d6485fdc47}
      \strng{fullhash}{bfe21ee39369b5ca36f289d6485fdc47}
      \strng{fullhashraw}{bfe21ee39369b5ca36f289d6485fdc47}
      \strng{bibnamehash}{bfe21ee39369b5ca36f289d6485fdc47}
      \strng{authorbibnamehash}{bfe21ee39369b5ca36f289d6485fdc47}
      \strng{authornamehash}{bfe21ee39369b5ca36f289d6485fdc47}
      \strng{authorfullhash}{bfe21ee39369b5ca36f289d6485fdc47}
      \strng{authorfullhashraw}{bfe21ee39369b5ca36f289d6485fdc47}
      \field{sortinit}{L}
      \field{sortinithash}{7c47d417cecb1f4bd38d1825c427a61a}
      \field{labelnamesource}{author}
      \field{labeltitlesource}{title}
      \field{eprinttype}{arXiv}
      \field{title}{Towards a double operadic theory of systems}
      \field{year}{2025}
      \verb{eprint}
      \verb 2505.18329
      \endverb
    \endentry
    \entry{liellcock24}{online}{}{}
      \name{author}{2}{}{%
        {{hash=ef905064829849e7c9b03a77906af240}{%
           family={Liell-Cock},
           familyi={L\bibinithyphendelim C\bibinitperiod},
           given={Jack},
           giveni={J\bibinitperiod}}}%
        {{hash=5351d3a060653e114d956a0a3700d556}{%
           family={Staton},
           familyi={S\bibinitperiod},
           given={Sam},
           giveni={S\bibinitperiod}}}%
      }
      \strng{namehash}{1c171c845b920d2d35b837453837b4e3}
      \strng{fullhash}{1c171c845b920d2d35b837453837b4e3}
      \strng{fullhashraw}{1c171c845b920d2d35b837453837b4e3}
      \strng{bibnamehash}{1c171c845b920d2d35b837453837b4e3}
      \strng{authorbibnamehash}{1c171c845b920d2d35b837453837b4e3}
      \strng{authornamehash}{1c171c845b920d2d35b837453837b4e3}
      \strng{authorfullhash}{1c171c845b920d2d35b837453837b4e3}
      \strng{authorfullhashraw}{1c171c845b920d2d35b837453837b4e3}
      \field{sortinit}{L}
      \field{sortinithash}{7c47d417cecb1f4bd38d1825c427a61a}
      \field{labelnamesource}{author}
      \field{labeltitlesource}{title}
      \field{eprinttype}{arXiv}
      \field{title}{Compositional imprecise probability}
      \field{year}{2024}
      \verb{eprint}
      \verb 2405.09391
      \endverb
    \endentry
    \entry{lin88}{article}{}{}
      \name{author}{2}{}{%
        {{hash=37177639f418a6c6c34ed11f1d8ee99c}{%
           family={Lin},
           familyi={L\bibinitperiod},
           given={Feng},
           giveni={F\bibinitperiod}}}%
        {{hash=4b29da63a3f4112955c706c349f49b86}{%
           family={Wonham},
           familyi={W\bibinitperiod},
           given={W.\bibnamedelimi Murray},
           giveni={W\bibinitperiod\bibinitdelim M\bibinitperiod}}}%
      }
      \strng{namehash}{e915a497655e7d3b292cfd26617a1b5a}
      \strng{fullhash}{e915a497655e7d3b292cfd26617a1b5a}
      \strng{fullhashraw}{e915a497655e7d3b292cfd26617a1b5a}
      \strng{bibnamehash}{e915a497655e7d3b292cfd26617a1b5a}
      \strng{authorbibnamehash}{e915a497655e7d3b292cfd26617a1b5a}
      \strng{authornamehash}{e915a497655e7d3b292cfd26617a1b5a}
      \strng{authorfullhash}{e915a497655e7d3b292cfd26617a1b5a}
      \strng{authorfullhashraw}{e915a497655e7d3b292cfd26617a1b5a}
      \field{sortinit}{L}
      \field{sortinithash}{7c47d417cecb1f4bd38d1825c427a61a}
      \field{labelnamesource}{author}
      \field{labeltitlesource}{title}
      \field{journaltitle}{Information Sciences}
      \field{title}{On Observability of Discrete-Event Systems}
      \field{volume}{44}
      \field{year}{1988}
      \field{pages}{173\bibrangedash 198}
      \range{pages}{26}
      \verb{doi}
      \verb 10.1016/0020-0255(88)90001-1
      \endverb
    \endentry
    \entry{lindvall1992}{book}{}{}
      \name{author}{1}{}{%
        {{hash=c1c35bd6a1ff39bc6b5967b323c52f71}{%
           family={Lindvall},
           familyi={L\bibinitperiod},
           given={T.},
           giveni={T\bibinitperiod}}}%
      }
      \list{publisher}{1}{%
        {Wiley}%
      }
      \strng{namehash}{c1c35bd6a1ff39bc6b5967b323c52f71}
      \strng{fullhash}{c1c35bd6a1ff39bc6b5967b323c52f71}
      \strng{fullhashraw}{c1c35bd6a1ff39bc6b5967b323c52f71}
      \strng{bibnamehash}{c1c35bd6a1ff39bc6b5967b323c52f71}
      \strng{authorbibnamehash}{c1c35bd6a1ff39bc6b5967b323c52f71}
      \strng{authornamehash}{c1c35bd6a1ff39bc6b5967b323c52f71}
      \strng{authorfullhash}{c1c35bd6a1ff39bc6b5967b323c52f71}
      \strng{authorfullhashraw}{c1c35bd6a1ff39bc6b5967b323c52f71}
      \field{sortinit}{L}
      \field{sortinithash}{7c47d417cecb1f4bd38d1825c427a61a}
      \field{labelnamesource}{author}
      \field{labeltitlesource}{title}
      \field{note}{Corrected reprint, Dover, 2002. (The coupling inequality $\mathrm{TV} \leq \Pr[X \neq Y]$: §I.2.)}
      \field{title}{Lectures on the Coupling Method}
      \field{year}{1992}
    \endentry
    \entry{lorenz23}{online}{}{}
      \name{author}{2}{}{%
        {{hash=02916779717c8470e37c6864851eee14}{%
           family={Lorenz},
           familyi={L\bibinitperiod},
           given={Robin},
           giveni={R\bibinitperiod}}}%
        {{hash=7bae254863e0a226b63e059a16171ce8}{%
           family={Tull},
           familyi={T\bibinitperiod},
           given={Sean},
           giveni={S\bibinitperiod}}}%
      }
      \strng{namehash}{f4e62a59270e8308911b28cce6c265db}
      \strng{fullhash}{f4e62a59270e8308911b28cce6c265db}
      \strng{fullhashraw}{f4e62a59270e8308911b28cce6c265db}
      \strng{bibnamehash}{f4e62a59270e8308911b28cce6c265db}
      \strng{authorbibnamehash}{f4e62a59270e8308911b28cce6c265db}
      \strng{authornamehash}{f4e62a59270e8308911b28cce6c265db}
      \strng{authorfullhash}{f4e62a59270e8308911b28cce6c265db}
      \strng{authorfullhashraw}{f4e62a59270e8308911b28cce6c265db}
      \field{sortinit}{L}
      \field{sortinithash}{7c47d417cecb1f4bd38d1825c427a61a}
      \field{labelnamesource}{author}
      \field{labeltitlesource}{title}
      \field{eprinttype}{arXiv}
      \field{title}{Causal models in string diagrams}
      \field{year}{2023}
      \verb{eprint}
      \verb 2304.07638
      \endverb
    \endentry
    \entry{lynchvaandrager1995}{article}{}{}
      \name{author}{2}{}{%
        {{hash=c3ffd26210bce178dd58f2fe8566c7b4}{%
           family={Lynch},
           familyi={L\bibinitperiod},
           given={N.},
           giveni={N\bibinitperiod}}}%
        {{hash=5a88278aff3b1009baf611b4440b5b22}{%
           family={Vaandrager},
           familyi={V\bibinitperiod},
           given={F.},
           giveni={F\bibinitperiod}}}%
      }
      \strng{namehash}{1a256570fa1bb687b9bd3c3f63dce2e0}
      \strng{fullhash}{1a256570fa1bb687b9bd3c3f63dce2e0}
      \strng{fullhashraw}{1a256570fa1bb687b9bd3c3f63dce2e0}
      \strng{bibnamehash}{1a256570fa1bb687b9bd3c3f63dce2e0}
      \strng{authorbibnamehash}{1a256570fa1bb687b9bd3c3f63dce2e0}
      \strng{authornamehash}{1a256570fa1bb687b9bd3c3f63dce2e0}
      \strng{authorfullhash}{1a256570fa1bb687b9bd3c3f63dce2e0}
      \strng{authorfullhashraw}{1a256570fa1bb687b9bd3c3f63dce2e0}
      \field{sortinit}{L}
      \field{sortinithash}{7c47d417cecb1f4bd38d1825c427a61a}
      \field{labelnamesource}{author}
      \field{labeltitlesource}{title}
      \field{journaltitle}{Information and Computation}
      \field{number}{2}
      \field{title}{Forward and Backward Simulations, Part I: Untimed Systems}
      \field{volume}{121}
      \field{year}{1995}
      \field{pages}{214\bibrangedash 233}
      \range{pages}{20}
    \endentry
    \entry{makkaireyes77}{book}{}{}
      \name{author}{2}{}{%
        {{hash=9fe9b07a6d2433e2262dfdddd43177fb}{%
           family={Makkai},
           familyi={M\bibinitperiod},
           given={Michael},
           giveni={M\bibinitperiod}}}%
        {{hash=d36ee7db4f3786477083ecd5d16e4707}{%
           family={Reyes},
           familyi={R\bibinitperiod},
           given={Gonzalo\bibnamedelima E.},
           giveni={G\bibinitperiod\bibinitdelim E\bibinitperiod}}}%
      }
      \list{publisher}{1}{%
        {Springer}%
      }
      \strng{namehash}{8010eb533f19735d30ba339c22bd05ac}
      \strng{fullhash}{8010eb533f19735d30ba339c22bd05ac}
      \strng{fullhashraw}{8010eb533f19735d30ba339c22bd05ac}
      \strng{bibnamehash}{8010eb533f19735d30ba339c22bd05ac}
      \strng{authorbibnamehash}{8010eb533f19735d30ba339c22bd05ac}
      \strng{authornamehash}{8010eb533f19735d30ba339c22bd05ac}
      \strng{authorfullhash}{8010eb533f19735d30ba339c22bd05ac}
      \strng{authorfullhashraw}{8010eb533f19735d30ba339c22bd05ac}
      \field{sortinit}{M}
      \field{sortinithash}{4625c616857f13d17ce56f7d4f97d451}
      \field{labelnamesource}{author}
      \field{labeltitlesource}{title}
      \field{series}{Lecture Notes in Mathematics}
      \field{subtitle}{Model-Theoretical Methods in the Theory of Topoi and Related Categories}
      \field{title}{First Order Categorical Logic}
      \field{volume}{611}
      \field{year}{1977}
      \verb{doi}
      \verb 10.1007/BFb0066201
      \endverb
    \endentry
    \entry{malament77}{article}{}{}
      \name{author}{1}{}{%
        {{hash=881f50e4159c08f1e77ba1f44eee9a7d}{%
           family={Malament},
           familyi={M\bibinitperiod},
           given={David},
           giveni={D\bibinitperiod}}}%
      }
      \strng{namehash}{881f50e4159c08f1e77ba1f44eee9a7d}
      \strng{fullhash}{881f50e4159c08f1e77ba1f44eee9a7d}
      \strng{fullhashraw}{881f50e4159c08f1e77ba1f44eee9a7d}
      \strng{bibnamehash}{881f50e4159c08f1e77ba1f44eee9a7d}
      \strng{authorbibnamehash}{881f50e4159c08f1e77ba1f44eee9a7d}
      \strng{authornamehash}{881f50e4159c08f1e77ba1f44eee9a7d}
      \strng{authorfullhash}{881f50e4159c08f1e77ba1f44eee9a7d}
      \strng{authorfullhashraw}{881f50e4159c08f1e77ba1f44eee9a7d}
      \field{sortinit}{M}
      \field{sortinithash}{4625c616857f13d17ce56f7d4f97d451}
      \field{labelnamesource}{author}
      \field{labeltitlesource}{title}
      \field{journaltitle}{Journal of Mathematical Physics}
      \field{number}{7}
      \field{title}{The class of continuous timelike curves determines the topology of spacetime}
      \field{volume}{18}
      \field{year}{1977}
      \field{pages}{1399\bibrangedash 1404}
      \range{pages}{6}
      \verb{doi}
      \verb 10.1063/1.523436
      \endverb
    \endentry
    \entry{martin04}{online}{}{}
      \name{author}{2}{}{%
        {{hash=584e192f6a9d206af361f971c7eb9b06}{%
           family={Martin},
           familyi={M\bibinitperiod},
           given={Keye},
           giveni={K\bibinitperiod}}}%
        {{hash=26a2b8e345710f6270131c87dba77ea1}{%
           family={Panangaden},
           familyi={P\bibinitperiod},
           given={Prakash},
           giveni={P\bibinitperiod}}}%
      }
      \strng{namehash}{d509691be0f0978178725c2d3e68f487}
      \strng{fullhash}{d509691be0f0978178725c2d3e68f487}
      \strng{fullhashraw}{d509691be0f0978178725c2d3e68f487}
      \strng{bibnamehash}{d509691be0f0978178725c2d3e68f487}
      \strng{authorbibnamehash}{d509691be0f0978178725c2d3e68f487}
      \strng{authornamehash}{d509691be0f0978178725c2d3e68f487}
      \strng{authorfullhash}{d509691be0f0978178725c2d3e68f487}
      \strng{authorfullhashraw}{d509691be0f0978178725c2d3e68f487}
      \field{sortinit}{M}
      \field{sortinithash}{4625c616857f13d17ce56f7d4f97d451}
      \field{labelnamesource}{author}
      \field{labeltitlesource}{title}
      \field{eprinttype}{arXiv}
      \field{title}{A domain of spacetime intervals in general relativity}
      \field{year}{2004}
      \verb{eprint}
      \verb gr-qc/0407094
      \endverb
    \endentry
    \entry{master25}{online}{}{}
      \name{author}{2}{}{%
        {{hash=fa8912eb6f0f0c230fc8d8371db35a2f}{%
           family={Master},
           familyi={M\bibinitperiod},
           given={Jade},
           giveni={J\bibinitperiod}}}%
        {{hash=79d9653d18a9ea620d1ef67f03059a03}{%
           family={Moeller},
           familyi={M\bibinitperiod},
           given={Joe},
           giveni={J\bibinitperiod}}}%
      }
      \strng{namehash}{d1118bc05bae7fe5023729ab07fd93c4}
      \strng{fullhash}{d1118bc05bae7fe5023729ab07fd93c4}
      \strng{fullhashraw}{d1118bc05bae7fe5023729ab07fd93c4}
      \strng{bibnamehash}{d1118bc05bae7fe5023729ab07fd93c4}
      \strng{authorbibnamehash}{d1118bc05bae7fe5023729ab07fd93c4}
      \strng{authornamehash}{d1118bc05bae7fe5023729ab07fd93c4}
      \strng{authorfullhash}{d1118bc05bae7fe5023729ab07fd93c4}
      \strng{authorfullhashraw}{d1118bc05bae7fe5023729ab07fd93c4}
      \field{sortinit}{M}
      \field{sortinithash}{4625c616857f13d17ce56f7d4f97d451}
      \field{labelnamesource}{author}
      \field{labeltitlesource}{title}
      \field{eprinttype}{arXiv}
      \field{title}{Colored Petri Nets are Monoidal Double Functors}
      \field{year}{2025}
      \verb{eprint}
      \verb 2510.01946
      \endverb
    \endentry
    \entry{master19}{online}{}{}
      \name{author}{4}{}{%
        {{hash=fa8912eb6f0f0c230fc8d8371db35a2f}{%
           family={Master},
           familyi={M\bibinitperiod},
           given={Jade},
           giveni={J\bibinitperiod}}}%
        {{hash=318b3f8d6da4d1c854033da07460a0cd}{%
           family={Patterson},
           familyi={P\bibinitperiod},
           given={Evan},
           giveni={E\bibinitperiod}}}%
        {{hash=8cfe58f934c01a16c6c0caaa0535f032}{%
           family={Yousfi},
           familyi={Y\bibinitperiod},
           given={Shahin},
           giveni={S\bibinitperiod}}}%
        {{hash=145ce0b9b12c1827b61fe456a4157a40}{%
           family={Canedo},
           familyi={C\bibinitperiod},
           given={Arquimedes},
           giveni={A\bibinitperiod}}}%
      }
      \strng{namehash}{ca96177579fbecfe39631a190e85df74}
      \strng{fullhash}{ca96177579fbecfe39631a190e85df74}
      \strng{fullhashraw}{ca96177579fbecfe39631a190e85df74}
      \strng{bibnamehash}{ca96177579fbecfe39631a190e85df74}
      \strng{authorbibnamehash}{ca96177579fbecfe39631a190e85df74}
      \strng{authornamehash}{ca96177579fbecfe39631a190e85df74}
      \strng{authorfullhash}{ca96177579fbecfe39631a190e85df74}
      \strng{authorfullhashraw}{ca96177579fbecfe39631a190e85df74}
      \field{sortinit}{M}
      \field{sortinithash}{4625c616857f13d17ce56f7d4f97d451}
      \field{labelnamesource}{author}
      \field{labeltitlesource}{title}
      \field{eprinttype}{arXiv}
      \field{title}{String Diagrams for Assembly Planning}
      \field{year}{2019}
      \verb{eprint}
      \verb 1909.10475
      \endverb
    \endentry
    \entry{matache23}{online}{}{}
      \name{author}{4}{}{%
        {{hash=64f8bbf92d704a9c6a2b95ee2ec33414}{%
           family={Matache},
           familyi={M\bibinitperiod},
           given={Cristina},
           giveni={C\bibinitperiod}}}%
        {{hash=bc6254241ef2eab5c89970cce810c738}{%
           family={Moss},
           familyi={M\bibinitperiod},
           given={Sean},
           giveni={S\bibinitperiod}}}%
        {{hash=5351d3a060653e114d956a0a3700d556}{%
           family={Staton},
           familyi={S\bibinitperiod},
           given={Sam},
           giveni={S\bibinitperiod}}}%
        {{hash=fdc519e23b3990467a192caa19b0d177}{%
           family={Suo},
           familyi={S\bibinitperiod},
           given={Ariadne\bibnamedelima Si},
           giveni={A\bibinitperiod\bibinitdelim S\bibinitperiod}}}%
      }
      \strng{namehash}{43f50d8e79e83c6d537785c8e5957474}
      \strng{fullhash}{43f50d8e79e83c6d537785c8e5957474}
      \strng{fullhashraw}{43f50d8e79e83c6d537785c8e5957474}
      \strng{bibnamehash}{43f50d8e79e83c6d537785c8e5957474}
      \strng{authorbibnamehash}{43f50d8e79e83c6d537785c8e5957474}
      \strng{authornamehash}{43f50d8e79e83c6d537785c8e5957474}
      \strng{authorfullhash}{43f50d8e79e83c6d537785c8e5957474}
      \strng{authorfullhashraw}{43f50d8e79e83c6d537785c8e5957474}
      \field{sortinit}{M}
      \field{sortinithash}{4625c616857f13d17ce56f7d4f97d451}
      \field{labelnamesource}{author}
      \field{labeltitlesource}{title}
      \field{eprinttype}{arXiv}
      \field{subtitle}{semirings, monads, and tensors}
      \field{title}{Denotational semantics for languages for inference}
      \field{year}{2023}
      \verb{eprint}
      \verb 2312.16694
      \endverb
    \endentry
    \entry{matheron75}{book}{}{}
      \name{author}{1}{}{%
        {{hash=d0338063525106ec70892b82744b51d1}{%
           family={Matheron},
           familyi={M\bibinitperiod},
           given={Georges},
           giveni={G\bibinitperiod}}}%
      }
      \list{location}{1}{%
        {New York}%
      }
      \list{publisher}{1}{%
        {John Wiley \& Sons}%
      }
      \strng{namehash}{d0338063525106ec70892b82744b51d1}
      \strng{fullhash}{d0338063525106ec70892b82744b51d1}
      \strng{fullhashraw}{d0338063525106ec70892b82744b51d1}
      \strng{bibnamehash}{d0338063525106ec70892b82744b51d1}
      \strng{authorbibnamehash}{d0338063525106ec70892b82744b51d1}
      \strng{authornamehash}{d0338063525106ec70892b82744b51d1}
      \strng{authorfullhash}{d0338063525106ec70892b82744b51d1}
      \strng{authorfullhashraw}{d0338063525106ec70892b82744b51d1}
      \field{sortinit}{M}
      \field{sortinithash}{4625c616857f13d17ce56f7d4f97d451}
      \field{labelnamesource}{author}
      \field{labeltitlesource}{title}
      \field{title}{Random Sets and Integral Geometry}
      \field{year}{1975}
    \endentry
    \entry{mciver-morgan-2005}{book}{}{}
      \name{author}{2}{}{%
        {{hash=4b20e9f564933c17471e86d416b8a63a}{%
           family={McIver},
           familyi={M\bibinitperiod},
           given={Annabelle},
           giveni={A\bibinitperiod}}}%
        {{hash=2f227a08ec40c25a3b17ddece38fc8fc}{%
           family={Morgan},
           familyi={M\bibinitperiod},
           given={Carroll},
           giveni={C\bibinitperiod}}}%
      }
      \list{publisher}{1}{%
        {Springer}%
      }
      \strng{namehash}{ae495eef6c34732c0d5ce81308ca3f12}
      \strng{fullhash}{ae495eef6c34732c0d5ce81308ca3f12}
      \strng{fullhashraw}{ae495eef6c34732c0d5ce81308ca3f12}
      \strng{bibnamehash}{ae495eef6c34732c0d5ce81308ca3f12}
      \strng{authorbibnamehash}{ae495eef6c34732c0d5ce81308ca3f12}
      \strng{authornamehash}{ae495eef6c34732c0d5ce81308ca3f12}
      \strng{authorfullhash}{ae495eef6c34732c0d5ce81308ca3f12}
      \strng{authorfullhashraw}{ae495eef6c34732c0d5ce81308ca3f12}
      \field{sortinit}{M}
      \field{sortinithash}{4625c616857f13d17ce56f7d4f97d451}
      \field{labelnamesource}{author}
      \field{labeltitlesource}{title}
      \field{series}{Monographs in Computer Science}
      \field{title}{Abstraction, Refinement and Proof for Probabilistic Systems}
      \field{year}{2005}
    \endentry
    \entry{mcmillan1999}{inproceedings}{}{}
      \name{author}{1}{}{%
        {{hash=88f7b85ef2a816fa2b52cfff57905a87}{%
           family={McMillan},
           familyi={M\bibinitperiod},
           given={K.\bibnamedelimi L.},
           giveni={K\bibinitperiod\bibinitdelim L\bibinitperiod}}}%
      }
      \list{publisher}{1}{%
        {Springer}%
      }
      \strng{namehash}{88f7b85ef2a816fa2b52cfff57905a87}
      \strng{fullhash}{88f7b85ef2a816fa2b52cfff57905a87}
      \strng{fullhashraw}{88f7b85ef2a816fa2b52cfff57905a87}
      \strng{bibnamehash}{88f7b85ef2a816fa2b52cfff57905a87}
      \strng{authorbibnamehash}{88f7b85ef2a816fa2b52cfff57905a87}
      \strng{authornamehash}{88f7b85ef2a816fa2b52cfff57905a87}
      \strng{authorfullhash}{88f7b85ef2a816fa2b52cfff57905a87}
      \strng{authorfullhashraw}{88f7b85ef2a816fa2b52cfff57905a87}
      \field{sortinit}{M}
      \field{sortinithash}{4625c616857f13d17ce56f7d4f97d451}
      \field{labelnamesource}{author}
      \field{labeltitlesource}{title}
      \field{booktitle}{Correct Hardware Design and Verification Methods (CHARME 1999)}
      \field{series}{Lecture Notes in Computer Science}
      \field{title}{Circular Compositional Reasoning about Liveness}
      \field{volume}{1703}
      \field{year}{1999}
      \field{pages}{342\bibrangedash 346}
      \range{pages}{5}
    \endentry
    \entry{mellies06csl}{inproceedings}{}{}
      \name{author}{1}{}{%
        {{hash=a9f37916182d59fd60115ad8a650c06d}{%
           family={Melliès},
           familyi={M\bibinitperiod},
           given={Paul-André},
           giveni={P\bibinithyphendelim A\bibinitperiod}}}%
      }
      \list{publisher}{1}{%
        {Springer}%
      }
      \strng{namehash}{a9f37916182d59fd60115ad8a650c06d}
      \strng{fullhash}{a9f37916182d59fd60115ad8a650c06d}
      \strng{fullhashraw}{a9f37916182d59fd60115ad8a650c06d}
      \strng{bibnamehash}{a9f37916182d59fd60115ad8a650c06d}
      \strng{authorbibnamehash}{a9f37916182d59fd60115ad8a650c06d}
      \strng{authornamehash}{a9f37916182d59fd60115ad8a650c06d}
      \strng{authorfullhash}{a9f37916182d59fd60115ad8a650c06d}
      \strng{authorfullhashraw}{a9f37916182d59fd60115ad8a650c06d}
      \field{sortinit}{M}
      \field{sortinithash}{4625c616857f13d17ce56f7d4f97d451}
      \field{labelnamesource}{author}
      \field{labeltitlesource}{title}
      \field{booktitle}{Computer Science Logic (CSL 2006)}
      \field{series}{Lecture Notes in Computer Science}
      \field{title}{Functorial Boxes in String Diagrams}
      \field{volume}{4207}
      \field{year}{2006}
      \field{pages}{1\bibrangedash 30}
      \range{pages}{30}
      \verb{doi}
      \verb 10.1007/11874683_1
      \endverb
    \endentry
    \entry{mellies16}{online}{}{}
      \name{author}{2}{}{%
        {{hash=a9f37916182d59fd60115ad8a650c06d}{%
           family={Melliès},
           familyi={M\bibinitperiod},
           given={Paul-André},
           giveni={P\bibinithyphendelim A\bibinitperiod}}}%
        {{hash=d63c7a22e2f6832bf45413ca88d2406d}{%
           family={Zeilberger},
           familyi={Z\bibinitperiod},
           given={Noam},
           giveni={N\bibinitperiod}}}%
      }
      \strng{namehash}{0049be37c9d5bf12443e87b2f06a1a93}
      \strng{fullhash}{0049be37c9d5bf12443e87b2f06a1a93}
      \strng{fullhashraw}{0049be37c9d5bf12443e87b2f06a1a93}
      \strng{bibnamehash}{0049be37c9d5bf12443e87b2f06a1a93}
      \strng{authorbibnamehash}{0049be37c9d5bf12443e87b2f06a1a93}
      \strng{authornamehash}{0049be37c9d5bf12443e87b2f06a1a93}
      \strng{authorfullhash}{0049be37c9d5bf12443e87b2f06a1a93}
      \strng{authorfullhashraw}{0049be37c9d5bf12443e87b2f06a1a93}
      \field{sortinit}{M}
      \field{sortinithash}{4625c616857f13d17ce56f7d4f97d451}
      \field{labelnamesource}{author}
      \field{labeltitlesource}{title}
      \field{eprinttype}{arXiv}
      \field{title}{A bifibrational reconstruction of Lawvere's presheaf hyperdoctrine}
      \field{year}{2016}
      \verb{eprint}
      \verb 1601.06098
      \endverb
    \endentry
    \entry{meyntweedie1993}{book}{}{}
      \name{author}{2}{}{%
        {{hash=b5d67d085ed8aa600906cf4bcd05af0d}{%
           family={Meyn},
           familyi={M\bibinitperiod},
           given={S.\bibnamedelimi P.},
           giveni={S\bibinitperiod\bibinitdelim P\bibinitperiod}}}%
        {{hash=c1e5215c6cadb271f3cb5ee54b1d52a6}{%
           family={Tweedie},
           familyi={T\bibinitperiod},
           given={R.\bibnamedelimi L.},
           giveni={R\bibinitperiod\bibinitdelim L\bibinitperiod}}}%
      }
      \list{publisher}{1}{%
        {Springer}%
      }
      \strng{namehash}{9db71cabe531f807bc39b962a05b6466}
      \strng{fullhash}{9db71cabe531f807bc39b962a05b6466}
      \strng{fullhashraw}{9db71cabe531f807bc39b962a05b6466}
      \strng{bibnamehash}{9db71cabe531f807bc39b962a05b6466}
      \strng{authorbibnamehash}{9db71cabe531f807bc39b962a05b6466}
      \strng{authornamehash}{9db71cabe531f807bc39b962a05b6466}
      \strng{authorfullhash}{9db71cabe531f807bc39b962a05b6466}
      \strng{authorfullhashraw}{9db71cabe531f807bc39b962a05b6466}
      \field{extraname}{1}
      \field{sortinit}{M}
      \field{sortinithash}{4625c616857f13d17ce56f7d4f97d451}
      \field{labelnamesource}{author}
      \field{labeltitlesource}{title}
      \field{note}{(Foster--Lyapunov drift conditions, systematized.)}
      \field{title}{Markov Chains and Stochastic Stability}
      \field{year}{1993}
    \endentry
    \entry{meyntweedie2009}{book}{}{}
      \name{author}{2}{}{%
        {{hash=b5d67d085ed8aa600906cf4bcd05af0d}{%
           family={Meyn},
           familyi={M\bibinitperiod},
           given={S.\bibnamedelimi P.},
           giveni={S\bibinitperiod\bibinitdelim P\bibinitperiod}}}%
        {{hash=c1e5215c6cadb271f3cb5ee54b1d52a6}{%
           family={Tweedie},
           familyi={T\bibinitperiod},
           given={R.\bibnamedelimi L.},
           giveni={R\bibinitperiod\bibinitdelim L\bibinitperiod}}}%
      }
      \list{publisher}{1}{%
        {Cambridge University Press}%
      }
      \strng{namehash}{9db71cabe531f807bc39b962a05b6466}
      \strng{fullhash}{9db71cabe531f807bc39b962a05b6466}
      \strng{fullhashraw}{9db71cabe531f807bc39b962a05b6466}
      \strng{bibnamehash}{9db71cabe531f807bc39b962a05b6466}
      \strng{authorbibnamehash}{9db71cabe531f807bc39b962a05b6466}
      \strng{authornamehash}{9db71cabe531f807bc39b962a05b6466}
      \strng{authorfullhash}{9db71cabe531f807bc39b962a05b6466}
      \strng{authorfullhashraw}{9db71cabe531f807bc39b962a05b6466}
      \field{extraname}{2}
      \field{sortinit}{M}
      \field{sortinithash}{4625c616857f13d17ce56f7d4f97d451}
      \field{labelnamesource}{author}
      \field{labeltitlesource}{title}
      \field{edition}{2}
      \field{note}{(Foster--Lyapunov drift conditions; the geometric drift form behind the deadline battery.)}
      \field{title}{Markov Chains and Stochastic Stability}
      \field{year}{2009}
    \endentry
    \entry{mio20}{online}{}{}
      \name{author}{3}{}{%
        {{hash=481b55caacd4280f6caf395f9450cb8a}{%
           family={Mio},
           familyi={M\bibinitperiod},
           given={Matteo},
           giveni={M\bibinitperiod}}}%
        {{hash=a73f862dcc23646bd8abccf9b1ed8181}{%
           family={Sarkis},
           familyi={S\bibinitperiod},
           given={Ralph},
           giveni={R\bibinitperiod}}}%
        {{hash=b3043436dccd7a0c2e3164c7cc4389e9}{%
           family={Vignudelli},
           familyi={V\bibinitperiod},
           given={Valeria},
           giveni={V\bibinitperiod}}}%
      }
      \strng{namehash}{c58758c33d3e89a6a0c8ad25c659e036}
      \strng{fullhash}{c58758c33d3e89a6a0c8ad25c659e036}
      \strng{fullhashraw}{c58758c33d3e89a6a0c8ad25c659e036}
      \strng{bibnamehash}{c58758c33d3e89a6a0c8ad25c659e036}
      \strng{authorbibnamehash}{c58758c33d3e89a6a0c8ad25c659e036}
      \strng{authornamehash}{c58758c33d3e89a6a0c8ad25c659e036}
      \strng{authorfullhash}{c58758c33d3e89a6a0c8ad25c659e036}
      \strng{authorfullhashraw}{c58758c33d3e89a6a0c8ad25c659e036}
      \field{sortinit}{M}
      \field{sortinithash}{4625c616857f13d17ce56f7d4f97d451}
      \field{labelnamesource}{author}
      \field{labeltitlesource}{title}
      \field{eprinttype}{arXiv}
      \field{subtitle}{equational and metric reasoning}
      \field{title}{Combining nondeterminism, probability, and termination}
      \field{year}{2020}
      \verb{eprint}
      \verb 2012.00382
      \endverb
    \endentry
    \entry{mischaikowmrozek2002}{incollection}{}{}
      \name{author}{2}{}{%
        {{hash=48f8b15c21dd70a5e22f1c20c558972d}{%
           family={Mischaikow},
           familyi={M\bibinitperiod},
           given={K.},
           giveni={K\bibinitperiod}}}%
        {{hash=4d1431357a5f805174412605788b4612}{%
           family={Mrozek},
           familyi={M\bibinitperiod},
           given={M.},
           giveni={M\bibinitperiod}}}%
      }
      \list{publisher}{1}{%
        {North-Holland}%
      }
      \strng{namehash}{3e19d720051b7f0ae6cdfdb51882e096}
      \strng{fullhash}{3e19d720051b7f0ae6cdfdb51882e096}
      \strng{fullhashraw}{3e19d720051b7f0ae6cdfdb51882e096}
      \strng{bibnamehash}{3e19d720051b7f0ae6cdfdb51882e096}
      \strng{authorbibnamehash}{3e19d720051b7f0ae6cdfdb51882e096}
      \strng{authornamehash}{3e19d720051b7f0ae6cdfdb51882e096}
      \strng{authorfullhash}{3e19d720051b7f0ae6cdfdb51882e096}
      \strng{authorfullhashraw}{3e19d720051b7f0ae6cdfdb51882e096}
      \field{sortinit}{M}
      \field{sortinithash}{4625c616857f13d17ce56f7d4f97d451}
      \field{labelnamesource}{author}
      \field{labeltitlesource}{title}
      \field{booktitle}{Handbook of Dynamical Systems}
      \field{title}{Conley index}
      \field{volume}{2}
      \field{year}{2002}
      \field{pages}{393\bibrangedash 460}
      \range{pages}{68}
    \endentry
    \entry{moylanhill1978}{article}{}{}
      \name{author}{2}{}{%
        {{hash=ad37e4e2db25484ca6798b75c112ada5}{%
           family={Moylan},
           familyi={M\bibinitperiod},
           given={P.\bibnamedelimi J.},
           giveni={P\bibinitperiod\bibinitdelim J\bibinitperiod}}}%
        {{hash=a53bb174661250ee39a25a69e8d9dba2}{%
           family={Hill},
           familyi={H\bibinitperiod},
           given={D.\bibnamedelimi J.},
           giveni={D\bibinitperiod\bibinitdelim J\bibinitperiod}}}%
      }
      \strng{namehash}{788b8f82efd413957938f53e9c88b019}
      \strng{fullhash}{788b8f82efd413957938f53e9c88b019}
      \strng{fullhashraw}{788b8f82efd413957938f53e9c88b019}
      \strng{bibnamehash}{788b8f82efd413957938f53e9c88b019}
      \strng{authorbibnamehash}{788b8f82efd413957938f53e9c88b019}
      \strng{authornamehash}{788b8f82efd413957938f53e9c88b019}
      \strng{authorfullhash}{788b8f82efd413957938f53e9c88b019}
      \strng{authorfullhashraw}{788b8f82efd413957938f53e9c88b019}
      \field{sortinit}{M}
      \field{sortinithash}{4625c616857f13d17ce56f7d4f97d451}
      \field{labelnamesource}{author}
      \field{labeltitlesource}{title}
      \field{journaltitle}{IEEE Transactions on Automatic Control}
      \field{number}{2}
      \field{title}{Stability criteria for large-scale systems}
      \field{volume}{23}
      \field{year}{1978}
      \field{pages}{143\bibrangedash 149}
      \range{pages}{7}
    \endentry
    \entry{muller2016elementary}{book}{}{}
      \name{author}{1}{}{%
        {{hash=98e34ec7d22b8d9756c713614a8b2502}{%
           family={Muller},
           familyi={M\bibinitperiod},
           given={Jean-Michel},
           giveni={J\bibinithyphendelim M\bibinitperiod}}}%
      }
      \list{location}{1}{%
        {Boston}%
      }
      \list{publisher}{1}{%
        {Birkhäuser}%
      }
      \strng{namehash}{98e34ec7d22b8d9756c713614a8b2502}
      \strng{fullhash}{98e34ec7d22b8d9756c713614a8b2502}
      \strng{fullhashraw}{98e34ec7d22b8d9756c713614a8b2502}
      \strng{bibnamehash}{98e34ec7d22b8d9756c713614a8b2502}
      \strng{authorbibnamehash}{98e34ec7d22b8d9756c713614a8b2502}
      \strng{authornamehash}{98e34ec7d22b8d9756c713614a8b2502}
      \strng{authorfullhash}{98e34ec7d22b8d9756c713614a8b2502}
      \strng{authorfullhashraw}{98e34ec7d22b8d9756c713614a8b2502}
      \field{sortinit}{M}
      \field{sortinithash}{4625c616857f13d17ce56f7d4f97d451}
      \field{labelnamesource}{author}
      \field{labeltitlesource}{title}
      \field{edition}{3}
      \field{subtitle}{Algorithms and Implementation}
      \field{title}{Elementary Functions}
      \field{year}{2016}
      \verb{doi}
      \verb 10.1007/978-1-4899-7983-4
      \endverb
    \endentry
    \entry{murata1989}{article}{}{}
      \name{author}{1}{}{%
        {{hash=9e0b29da8e4b57fcbb1bb6a862511b01}{%
           family={Murata},
           familyi={M\bibinitperiod},
           given={T.},
           giveni={T\bibinitperiod}}}%
      }
      \strng{namehash}{9e0b29da8e4b57fcbb1bb6a862511b01}
      \strng{fullhash}{9e0b29da8e4b57fcbb1bb6a862511b01}
      \strng{fullhashraw}{9e0b29da8e4b57fcbb1bb6a862511b01}
      \strng{bibnamehash}{9e0b29da8e4b57fcbb1bb6a862511b01}
      \strng{authorbibnamehash}{9e0b29da8e4b57fcbb1bb6a862511b01}
      \strng{authornamehash}{9e0b29da8e4b57fcbb1bb6a862511b01}
      \strng{authorfullhash}{9e0b29da8e4b57fcbb1bb6a862511b01}
      \strng{authorfullhashraw}{9e0b29da8e4b57fcbb1bb6a862511b01}
      \field{sortinit}{M}
      \field{sortinithash}{4625c616857f13d17ce56f7d4f97d451}
      \field{labelnamesource}{author}
      \field{labeltitlesource}{title}
      \field{journaltitle}{Proceedings of the IEEE}
      \field{note}{(Place invariants --- the linear conservation laws entering as free count rows.)}
      \field{number}{4}
      \field{title}{Petri nets: Properties, analysis and applications}
      \field{volume}{77}
      \field{year}{1989}
      \field{pages}{541\bibrangedash 580}
      \range{pages}{40}
    \endentry
    \entry{nagumo1942}{article}{}{}
      \name{author}{1}{}{%
        {{hash=35469142479240b5a761a1e7bbcb3041}{%
           family={Nagumo},
           familyi={N\bibinitperiod},
           given={M.},
           giveni={M\bibinitperiod}}}%
      }
      \strng{namehash}{35469142479240b5a761a1e7bbcb3041}
      \strng{fullhash}{35469142479240b5a761a1e7bbcb3041}
      \strng{fullhashraw}{35469142479240b5a761a1e7bbcb3041}
      \strng{bibnamehash}{35469142479240b5a761a1e7bbcb3041}
      \strng{authorbibnamehash}{35469142479240b5a761a1e7bbcb3041}
      \strng{authornamehash}{35469142479240b5a761a1e7bbcb3041}
      \strng{authorfullhash}{35469142479240b5a761a1e7bbcb3041}
      \strng{authorfullhashraw}{35469142479240b5a761a1e7bbcb3041}
      \field{sortinit}{N}
      \field{sortinithash}{22369a73d5f88983a108b63f07f37084}
      \field{labelnamesource}{author}
      \field{labeltitlesource}{title}
      \field{journaltitle}{Proceedings of the Physico-Mathematical Society of Japan}
      \field{series}{3}
      \field{title}{Über die Lage der Integralkurven gewöhnlicher Differentialgleichungen}
      \field{volume}{24}
      \field{year}{1942}
      \field{pages}{551\bibrangedash 559}
      \range{pages}{9}
    \endentry
    \entry{nelsonoppen1980}{article}{}{}
      \name{author}{2}{}{%
        {{hash=72453f903834a2c6b906c7e32e5becb8}{%
           family={Nelson},
           familyi={N\bibinitperiod},
           given={G.},
           giveni={G\bibinitperiod}}}%
        {{hash=0605a5ea60e508bb7ea9506abd32a27f}{%
           family={Oppen},
           familyi={O\bibinitperiod},
           given={D.\bibnamedelimi C.},
           giveni={D\bibinitperiod\bibinitdelim C\bibinitperiod}}}%
      }
      \strng{namehash}{cd7efa45d467bfd3f532abe7cda6e8ab}
      \strng{fullhash}{cd7efa45d467bfd3f532abe7cda6e8ab}
      \strng{fullhashraw}{cd7efa45d467bfd3f532abe7cda6e8ab}
      \strng{bibnamehash}{cd7efa45d467bfd3f532abe7cda6e8ab}
      \strng{authorbibnamehash}{cd7efa45d467bfd3f532abe7cda6e8ab}
      \strng{authornamehash}{cd7efa45d467bfd3f532abe7cda6e8ab}
      \strng{authorfullhash}{cd7efa45d467bfd3f532abe7cda6e8ab}
      \strng{authorfullhashraw}{cd7efa45d467bfd3f532abe7cda6e8ab}
      \field{sortinit}{N}
      \field{sortinithash}{22369a73d5f88983a108b63f07f37084}
      \field{labelnamesource}{author}
      \field{labeltitlesource}{title}
      \field{journaltitle}{Journal of the ACM}
      \field{number}{2}
      \field{title}{Fast decision procedures based on congruence closure}
      \field{volume}{27}
      \field{year}{1980}
      \field{pages}{356\bibrangedash 364}
      \range{pages}{9}
    \endentry
    \entry{nieuwenhuisoliveras2007}{article}{}{}
      \name{author}{2}{}{%
        {{hash=cda9f821bffb12a2fdca623af54d535b}{%
           family={Nieuwenhuis},
           familyi={N\bibinitperiod},
           given={R.},
           giveni={R\bibinitperiod}}}%
        {{hash=048f667585d726db5e82ed44d40b9c52}{%
           family={Oliveras},
           familyi={O\bibinitperiod},
           given={A.},
           giveni={A\bibinitperiod}}}%
      }
      \strng{namehash}{2fa89848704bdfd1c30ec56fb34d6712}
      \strng{fullhash}{2fa89848704bdfd1c30ec56fb34d6712}
      \strng{fullhashraw}{2fa89848704bdfd1c30ec56fb34d6712}
      \strng{bibnamehash}{2fa89848704bdfd1c30ec56fb34d6712}
      \strng{authorbibnamehash}{2fa89848704bdfd1c30ec56fb34d6712}
      \strng{authornamehash}{2fa89848704bdfd1c30ec56fb34d6712}
      \strng{authorfullhash}{2fa89848704bdfd1c30ec56fb34d6712}
      \strng{authorfullhashraw}{2fa89848704bdfd1c30ec56fb34d6712}
      \field{sortinit}{N}
      \field{sortinithash}{22369a73d5f88983a108b63f07f37084}
      \field{labelnamesource}{author}
      \field{labeltitlesource}{title}
      \field{journaltitle}{Information and Computation}
      \field{number}{4}
      \field{title}{Fast congruence closure and extensions}
      \field{volume}{205}
      \field{year}{2007}
      \field{pages}{557\bibrangedash 574}
      \range{pages}{18}
    \endentry
    \entry{norris-1997}{book}{}{}
      \name{author}{1}{}{%
        {{hash=ffda70e5672181a83a19fa48bc7a6334}{%
           family={Norris},
           familyi={N\bibinitperiod},
           given={J.\bibnamedelimi R.},
           giveni={J\bibinitperiod\bibinitdelim R\bibinitperiod}}}%
      }
      \list{location}{1}{%
        {Cambridge}%
      }
      \list{publisher}{1}{%
        {Cambridge University Press}%
      }
      \strng{namehash}{ffda70e5672181a83a19fa48bc7a6334}
      \strng{fullhash}{ffda70e5672181a83a19fa48bc7a6334}
      \strng{fullhashraw}{ffda70e5672181a83a19fa48bc7a6334}
      \strng{bibnamehash}{ffda70e5672181a83a19fa48bc7a6334}
      \strng{authorbibnamehash}{ffda70e5672181a83a19fa48bc7a6334}
      \strng{authornamehash}{ffda70e5672181a83a19fa48bc7a6334}
      \strng{authorfullhash}{ffda70e5672181a83a19fa48bc7a6334}
      \strng{authorfullhashraw}{ffda70e5672181a83a19fa48bc7a6334}
      \field{sortinit}{N}
      \field{sortinithash}{22369a73d5f88983a108b63f07f37084}
      \field{labelnamesource}{author}
      \field{labeltitlesource}{title}
      \field{title}{Markov Chains}
      \field{year}{1997}
    \endentry
    \entry{orevkov1979}{article}{}{}
      \name{author}{1}{}{%
        {{hash=7709c12c2c734392b8b146c4bd858be4}{%
           family={Orevkov},
           familyi={O\bibinitperiod},
           given={V.\bibnamedelimi P.},
           giveni={V\bibinitperiod\bibinitdelim P\bibinitperiod}}}%
      }
      \strng{namehash}{7709c12c2c734392b8b146c4bd858be4}
      \strng{fullhash}{7709c12c2c734392b8b146c4bd858be4}
      \strng{fullhashraw}{7709c12c2c734392b8b146c4bd858be4}
      \strng{bibnamehash}{7709c12c2c734392b8b146c4bd858be4}
      \strng{authorbibnamehash}{7709c12c2c734392b8b146c4bd858be4}
      \strng{authornamehash}{7709c12c2c734392b8b146c4bd858be4}
      \strng{authorfullhash}{7709c12c2c734392b8b146c4bd858be4}
      \strng{authorfullhashraw}{7709c12c2c734392b8b146c4bd858be4}
      \field{sortinit}{O}
      \field{sortinithash}{2cd7140a07aea5341f9e2771efe90aae}
      \field{labelnamesource}{author}
      \field{labeltitlesource}{title}
      \field{journaltitle}{Zapiski Nauchnykh Seminarov LOMI}
      \field{note}{English translation in Journal of Soviet Mathematics (1982).}
      \field{title}{Lower bounds for lengthening of proofs after cut-elimination}
      \field{volume}{88}
      \field{year}{1979}
      \field{pages}{137\bibrangedash 162}
      \range{pages}{26}
    \endentry
    \entry{parthasarathy67}{book}{}{}
      \name{author}{1}{}{%
        {{hash=82378da5401d2025e5dfd9e98d35cb0b}{%
           family={Parthasarathy},
           familyi={P\bibinitperiod},
           given={Kalyanapuram\bibnamedelima R.},
           giveni={K\bibinitperiod\bibinitdelim R\bibinitperiod}}}%
      }
      \list{location}{1}{%
        {New York}%
      }
      \list{publisher}{1}{%
        {Academic Press}%
      }
      \strng{namehash}{82378da5401d2025e5dfd9e98d35cb0b}
      \strng{fullhash}{82378da5401d2025e5dfd9e98d35cb0b}
      \strng{fullhashraw}{82378da5401d2025e5dfd9e98d35cb0b}
      \strng{bibnamehash}{82378da5401d2025e5dfd9e98d35cb0b}
      \strng{authorbibnamehash}{82378da5401d2025e5dfd9e98d35cb0b}
      \strng{authornamehash}{82378da5401d2025e5dfd9e98d35cb0b}
      \strng{authorfullhash}{82378da5401d2025e5dfd9e98d35cb0b}
      \strng{authorfullhashraw}{82378da5401d2025e5dfd9e98d35cb0b}
      \field{sortinit}{P}
      \field{sortinithash}{ff3bcf24f47321b42cb156c2cc8a8422}
      \field{labelnamesource}{author}
      \field{labeltitlesource}{title}
      \field{title}{Probability Measures on Metric Spaces}
      \field{year}{1967}
    \endentry
    \entry{patterson23}{online}{}{}
      \name{author}{1}{}{%
        {{hash=318b3f8d6da4d1c854033da07460a0cd}{%
           family={Patterson},
           familyi={P\bibinitperiod},
           given={Evan},
           giveni={E\bibinitperiod}}}%
      }
      \strng{namehash}{318b3f8d6da4d1c854033da07460a0cd}
      \strng{fullhash}{318b3f8d6da4d1c854033da07460a0cd}
      \strng{fullhashraw}{318b3f8d6da4d1c854033da07460a0cd}
      \strng{bibnamehash}{318b3f8d6da4d1c854033da07460a0cd}
      \strng{authorbibnamehash}{318b3f8d6da4d1c854033da07460a0cd}
      \strng{authornamehash}{318b3f8d6da4d1c854033da07460a0cd}
      \strng{authorfullhash}{318b3f8d6da4d1c854033da07460a0cd}
      \strng{authorfullhashraw}{318b3f8d6da4d1c854033da07460a0cd}
      \field{sortinit}{P}
      \field{sortinithash}{ff3bcf24f47321b42cb156c2cc8a8422}
      \field{labelnamesource}{author}
      \field{labeltitlesource}{title}
      \field{eprinttype}{arXiv}
      \field{title}{Structured and Decorated Cospans from the Viewpoint of Double Category Theory}
      \field{year}{2023}
      \verb{eprint}
      \verb 2304.00447
      \endverb
    \endentry
    \entry{patterson21}{online}{}{}
      \name{author}{3}{}{%
        {{hash=318b3f8d6da4d1c854033da07460a0cd}{%
           family={Patterson},
           familyi={P\bibinitperiod},
           given={Evan},
           giveni={E\bibinitperiod}}}%
        {{hash=070418290459e9ba8c3eaf0d0874d40f}{%
           family={Lynch},
           familyi={L\bibinitperiod},
           given={Owen},
           giveni={O\bibinitperiod}}}%
        {{hash=1d30a8595a636bf14fcd181b6822212d}{%
           family={Fairbanks},
           familyi={F\bibinitperiod},
           given={James},
           giveni={J\bibinitperiod}}}%
      }
      \strng{namehash}{4e6b9b4a7914d8a1ccd11ff0dc4b0211}
      \strng{fullhash}{4e6b9b4a7914d8a1ccd11ff0dc4b0211}
      \strng{fullhashraw}{4e6b9b4a7914d8a1ccd11ff0dc4b0211}
      \strng{bibnamehash}{4e6b9b4a7914d8a1ccd11ff0dc4b0211}
      \strng{authorbibnamehash}{4e6b9b4a7914d8a1ccd11ff0dc4b0211}
      \strng{authornamehash}{4e6b9b4a7914d8a1ccd11ff0dc4b0211}
      \strng{authorfullhash}{4e6b9b4a7914d8a1ccd11ff0dc4b0211}
      \strng{authorfullhashraw}{4e6b9b4a7914d8a1ccd11ff0dc4b0211}
      \field{sortinit}{P}
      \field{sortinithash}{ff3bcf24f47321b42cb156c2cc8a8422}
      \field{labelnamesource}{author}
      \field{labeltitlesource}{title}
      \field{eprinttype}{arXiv}
      \field{title}{Categorical Data Structures for Technical Computing}
      \field{year}{2021}
      \verb{eprint}
      \verb 2106.04703
      \endverb
    \endentry
    \entry{pearl-2009}{book}{}{}
      \name{author}{1}{}{%
        {{hash=809f695b398afbb54b544c49e8d1bbbb}{%
           family={Pearl},
           familyi={P\bibinitperiod},
           given={Judea},
           giveni={J\bibinitperiod}}}%
      }
      \list{publisher}{1}{%
        {Cambridge University Press}%
      }
      \strng{namehash}{809f695b398afbb54b544c49e8d1bbbb}
      \strng{fullhash}{809f695b398afbb54b544c49e8d1bbbb}
      \strng{fullhashraw}{809f695b398afbb54b544c49e8d1bbbb}
      \strng{bibnamehash}{809f695b398afbb54b544c49e8d1bbbb}
      \strng{authorbibnamehash}{809f695b398afbb54b544c49e8d1bbbb}
      \strng{authornamehash}{809f695b398afbb54b544c49e8d1bbbb}
      \strng{authorfullhash}{809f695b398afbb54b544c49e8d1bbbb}
      \strng{authorfullhashraw}{809f695b398afbb54b544c49e8d1bbbb}
      \field{extraname}{1}
      \field{sortinit}{P}
      \field{sortinithash}{ff3bcf24f47321b42cb156c2cc8a8422}
      \field{labelnamesource}{author}
      \field{labeltitlesource}{title}
      \field{edition}{2}
      \field{subtitle}{Models, Reasoning, and Inference}
      \field{title}{Causality}
      \field{year}{2009}
      \verb{doi}
      \verb 10.1017/CBO9780511803161
      \endverb
    \endentry
    \entry{pearl-1988}{book}{}{}
      \name{author}{1}{}{%
        {{hash=809f695b398afbb54b544c49e8d1bbbb}{%
           family={Pearl},
           familyi={P\bibinitperiod},
           given={Judea},
           giveni={J\bibinitperiod}}}%
      }
      \list{location}{1}{%
        {San Mateo}%
      }
      \list{publisher}{1}{%
        {Morgan Kaufmann}%
      }
      \strng{namehash}{809f695b398afbb54b544c49e8d1bbbb}
      \strng{fullhash}{809f695b398afbb54b544c49e8d1bbbb}
      \strng{fullhashraw}{809f695b398afbb54b544c49e8d1bbbb}
      \strng{bibnamehash}{809f695b398afbb54b544c49e8d1bbbb}
      \strng{authorbibnamehash}{809f695b398afbb54b544c49e8d1bbbb}
      \strng{authornamehash}{809f695b398afbb54b544c49e8d1bbbb}
      \strng{authorfullhash}{809f695b398afbb54b544c49e8d1bbbb}
      \strng{authorfullhashraw}{809f695b398afbb54b544c49e8d1bbbb}
      \field{extraname}{2}
      \field{sortinit}{P}
      \field{sortinithash}{ff3bcf24f47321b42cb156c2cc8a8422}
      \field{labelnamesource}{author}
      \field{labeltitlesource}{title}
      \field{title}{Probabilistic Reasoning in Intelligent Systems: Networks of Plausible Inference}
      \field{year}{1988}
    \endentry
    \entry{penco20}{online}{}{}
      \name{author}{1}{}{%
        {{hash=4270fb6b1aca08cba77699999ca008f4}{%
           family={Penco},
           familyi={P\bibinitperiod},
           given={Riccardo},
           giveni={R\bibinitperiod}}}%
      }
      \strng{namehash}{4270fb6b1aca08cba77699999ca008f4}
      \strng{fullhash}{4270fb6b1aca08cba77699999ca008f4}
      \strng{fullhashraw}{4270fb6b1aca08cba77699999ca008f4}
      \strng{bibnamehash}{4270fb6b1aca08cba77699999ca008f4}
      \strng{authorbibnamehash}{4270fb6b1aca08cba77699999ca008f4}
      \strng{authornamehash}{4270fb6b1aca08cba77699999ca008f4}
      \strng{authorfullhash}{4270fb6b1aca08cba77699999ca008f4}
      \strng{authorfullhashraw}{4270fb6b1aca08cba77699999ca008f4}
      \field{sortinit}{P}
      \field{sortinithash}{ff3bcf24f47321b42cb156c2cc8a8422}
      \field{labelnamesource}{author}
      \field{labeltitlesource}{title}
      \field{eprinttype}{arXiv}
      \field{title}{An Introduction to Effective Field Theories}
      \field{year}{2020}
      \verb{eprint}
      \verb 2006.16285
      \endverb
    \endentry
    \entry{peterseim26}{online}{}{}
      \name{author}{2}{}{%
        {{hash=062cb530879396d671e146dc74e3bb1f}{%
           family={Peterseim},
           familyi={P\bibinitperiod},
           given={Benedikt},
           giveni={B\bibinitperiod}}}%
        {{hash=e5e19a7f841e4ddad7f431aa8acf3455}{%
           family={Lopuhaä-Zwakenberg},
           familyi={L\bibinithyphendelim Z\bibinitperiod},
           given={Milan},
           giveni={M\bibinitperiod}}}%
      }
      \strng{namehash}{6e061385744a0c1b7f79f46100a5cd0e}
      \strng{fullhash}{6e061385744a0c1b7f79f46100a5cd0e}
      \strng{fullhashraw}{6e061385744a0c1b7f79f46100a5cd0e}
      \strng{bibnamehash}{6e061385744a0c1b7f79f46100a5cd0e}
      \strng{authorbibnamehash}{6e061385744a0c1b7f79f46100a5cd0e}
      \strng{authornamehash}{6e061385744a0c1b7f79f46100a5cd0e}
      \strng{authorfullhash}{6e061385744a0c1b7f79f46100a5cd0e}
      \strng{authorfullhashraw}{6e061385744a0c1b7f79f46100a5cd0e}
      \field{sortinit}{P}
      \field{sortinithash}{ff3bcf24f47321b42cb156c2cc8a8422}
      \field{labelnamesource}{author}
      \field{labeltitlesource}{title}
      \field{eprinttype}{arXiv}
      \field{title}{Fixed-parameter tractable inference for discrete probabilistic programs, via string diagram algebraisation}
      \field{year}{2026}
      \verb{eprint}
      \verb 2604.25321
      \endverb
    \endentry
    \entry{pitts00}{incollection}{}{}
      \name{author}{1}{}{%
        {{hash=a574d10c24f8ad6f78b35c12ff05e2c1}{%
           family={Pitts},
           familyi={P\bibinitperiod},
           given={Andrew\bibnamedelima M.},
           giveni={A\bibinitperiod\bibinitdelim M\bibinitperiod}}}%
      }
      \name{editor}{3}{}{%
        {{hash=f4721a56c3b94b002e7ef8e0ef2b89fa}{%
           family={Abramsky},
           familyi={A\bibinitperiod},
           given={Samson},
           giveni={S\bibinitperiod}}}%
        {{hash=e4c171212e77fadaec078dad6dc4d173}{%
           family={Gabbay},
           familyi={G\bibinitperiod},
           given={Dov\bibnamedelima M.},
           giveni={D\bibinitperiod\bibinitdelim M\bibinitperiod}}}%
        {{hash=3a60235ada37c3ee69326e04924f1e10}{%
           family={Maibaum},
           familyi={M\bibinitperiod},
           given={T.\bibnamedelimi S.\bibnamedelimi E.},
           giveni={T\bibinitperiod\bibinitdelim S\bibinitperiod\bibinitdelim E\bibinitperiod}}}%
      }
      \list{publisher}{1}{%
        {Oxford University Press}%
      }
      \strng{namehash}{a574d10c24f8ad6f78b35c12ff05e2c1}
      \strng{fullhash}{a574d10c24f8ad6f78b35c12ff05e2c1}
      \strng{fullhashraw}{a574d10c24f8ad6f78b35c12ff05e2c1}
      \strng{bibnamehash}{a574d10c24f8ad6f78b35c12ff05e2c1}
      \strng{authorbibnamehash}{a574d10c24f8ad6f78b35c12ff05e2c1}
      \strng{authornamehash}{a574d10c24f8ad6f78b35c12ff05e2c1}
      \strng{authorfullhash}{a574d10c24f8ad6f78b35c12ff05e2c1}
      \strng{authorfullhashraw}{a574d10c24f8ad6f78b35c12ff05e2c1}
      \strng{editorbibnamehash}{cdc76e8a71fa0e75cdc538d076fad26d}
      \strng{editornamehash}{cdc76e8a71fa0e75cdc538d076fad26d}
      \strng{editorfullhash}{cdc76e8a71fa0e75cdc538d076fad26d}
      \strng{editorfullhashraw}{cdc76e8a71fa0e75cdc538d076fad26d}
      \field{sortinit}{P}
      \field{sortinithash}{ff3bcf24f47321b42cb156c2cc8a8422}
      \field{labelnamesource}{author}
      \field{labeltitlesource}{title}
      \field{booktitle}{Handbook of Logic in Computer Science}
      \field{title}{Categorical Logic}
      \field{volume}{5}
      \field{year}{2000}
      \field{pages}{39\bibrangedash 128}
      \range{pages}{90}
    \endentry
    \entry{polya1945solve}{book}{}{}
      \name{author}{1}{}{%
        {{hash=5e81307c3c82bf0c48a54f6da67e4d89}{%
           family={Pólya},
           familyi={P\bibinitperiod},
           given={George},
           giveni={G\bibinitperiod}}}%
      }
      \list{publisher}{1}{%
        {Princeton University Press}%
      }
      \strng{namehash}{5e81307c3c82bf0c48a54f6da67e4d89}
      \strng{fullhash}{5e81307c3c82bf0c48a54f6da67e4d89}
      \strng{fullhashraw}{5e81307c3c82bf0c48a54f6da67e4d89}
      \strng{bibnamehash}{5e81307c3c82bf0c48a54f6da67e4d89}
      \strng{authorbibnamehash}{5e81307c3c82bf0c48a54f6da67e4d89}
      \strng{authornamehash}{5e81307c3c82bf0c48a54f6da67e4d89}
      \strng{authorfullhash}{5e81307c3c82bf0c48a54f6da67e4d89}
      \strng{authorfullhashraw}{5e81307c3c82bf0c48a54f6da67e4d89}
      \field{sortinit}{P}
      \field{sortinithash}{ff3bcf24f47321b42cb156c2cc8a8422}
      \field{labelnamesource}{author}
      \field{labeltitlesource}{title}
      \field{title}{How to Solve It: A New Aspect of Mathematical Method}
      \field{year}{1945}
    \endentry
    \entry{prajnajadbabaie2004}{inproceedings}{}{}
      \name{author}{2}{}{%
        {{hash=a1185d5d8074828aa8519f916afe931f}{%
           family={Prajna},
           familyi={P\bibinitperiod},
           given={S.},
           giveni={S\bibinitperiod}}}%
        {{hash=cfcf2072a527c943c0c633acdc8f3984}{%
           family={Jadbabaie},
           familyi={J\bibinitperiod},
           given={A.},
           giveni={A\bibinitperiod}}}%
      }
      \list{publisher}{1}{%
        {Springer}%
      }
      \strng{namehash}{0029782c9702ca81e7f33e5033b60e67}
      \strng{fullhash}{0029782c9702ca81e7f33e5033b60e67}
      \strng{fullhashraw}{0029782c9702ca81e7f33e5033b60e67}
      \strng{bibnamehash}{0029782c9702ca81e7f33e5033b60e67}
      \strng{authorbibnamehash}{0029782c9702ca81e7f33e5033b60e67}
      \strng{authornamehash}{0029782c9702ca81e7f33e5033b60e67}
      \strng{authorfullhash}{0029782c9702ca81e7f33e5033b60e67}
      \strng{authorfullhashraw}{0029782c9702ca81e7f33e5033b60e67}
      \field{sortinit}{P}
      \field{sortinithash}{ff3bcf24f47321b42cb156c2cc8a8422}
      \field{labelnamesource}{author}
      \field{labeltitlesource}{title}
      \field{booktitle}{Hybrid Systems: Computation and Control (HSCC 2004)}
      \field{series}{LNCS}
      \field{title}{Safety verification of hybrid systems using barrier certificates}
      \field{volume}{2993}
      \field{year}{2004}
      \field{pages}{477\bibrangedash 492}
      \range{pages}{16}
    \endentry
    \entry{prajnajadbabaiepappas2007}{article}{}{}
      \name{author}{3}{}{%
        {{hash=a1185d5d8074828aa8519f916afe931f}{%
           family={Prajna},
           familyi={P\bibinitperiod},
           given={S.},
           giveni={S\bibinitperiod}}}%
        {{hash=cfcf2072a527c943c0c633acdc8f3984}{%
           family={Jadbabaie},
           familyi={J\bibinitperiod},
           given={A.},
           giveni={A\bibinitperiod}}}%
        {{hash=fd39f64a316253aa07afe174ea8b90fd}{%
           family={Pappas},
           familyi={P\bibinitperiod},
           given={G.\bibnamedelimi J.},
           giveni={G\bibinitperiod\bibinitdelim J\bibinitperiod}}}%
      }
      \strng{namehash}{6d037206d5f8d80aca9fa542d2424e44}
      \strng{fullhash}{6d037206d5f8d80aca9fa542d2424e44}
      \strng{fullhashraw}{6d037206d5f8d80aca9fa542d2424e44}
      \strng{bibnamehash}{6d037206d5f8d80aca9fa542d2424e44}
      \strng{authorbibnamehash}{6d037206d5f8d80aca9fa542d2424e44}
      \strng{authornamehash}{6d037206d5f8d80aca9fa542d2424e44}
      \strng{authorfullhash}{6d037206d5f8d80aca9fa542d2424e44}
      \strng{authorfullhashraw}{6d037206d5f8d80aca9fa542d2424e44}
      \field{sortinit}{P}
      \field{sortinithash}{ff3bcf24f47321b42cb156c2cc8a8422}
      \field{labelnamesource}{author}
      \field{labeltitlesource}{title}
      \field{journaltitle}{IEEE Transactions on Automatic Control}
      \field{number}{8}
      \field{title}{A framework for worst-case and stochastic safety verification using barrier certificates}
      \field{volume}{52}
      \field{year}{2007}
      \field{pages}{1415\bibrangedash 1428}
      \range{pages}{14}
    \endentry
    \entry{puterman-1994}{book}{}{}
      \name{author}{1}{}{%
        {{hash=b4258b9291d903d97b69d703eebf790a}{%
           family={Puterman},
           familyi={P\bibinitperiod},
           given={Martin\bibnamedelima L.},
           giveni={M\bibinitperiod\bibinitdelim L\bibinitperiod}}}%
      }
      \list{location}{1}{%
        {New York}%
      }
      \list{publisher}{1}{%
        {Wiley}%
      }
      \strng{namehash}{b4258b9291d903d97b69d703eebf790a}
      \strng{fullhash}{b4258b9291d903d97b69d703eebf790a}
      \strng{fullhashraw}{b4258b9291d903d97b69d703eebf790a}
      \strng{bibnamehash}{b4258b9291d903d97b69d703eebf790a}
      \strng{authorbibnamehash}{b4258b9291d903d97b69d703eebf790a}
      \strng{authornamehash}{b4258b9291d903d97b69d703eebf790a}
      \strng{authorfullhash}{b4258b9291d903d97b69d703eebf790a}
      \strng{authorfullhashraw}{b4258b9291d903d97b69d703eebf790a}
      \field{sortinit}{P}
      \field{sortinithash}{ff3bcf24f47321b42cb156c2cc8a8422}
      \field{labelnamesource}{author}
      \field{labeltitlesource}{title}
      \field{title}{Markov Decision Processes: Discrete Stochastic Dynamic Programming}
      \field{year}{1994}
    \endentry
    \entry{quatmann18}{online}{}{}
      \name{author}{2}{}{%
        {{hash=d195773ad1a45df14875e5252d56ccb8}{%
           family={Quatmann},
           familyi={Q\bibinitperiod},
           given={Tim},
           giveni={T\bibinitperiod}}}%
        {{hash=392d5825b15ef7f81f68632c7a49fdef}{%
           family={Katoen},
           familyi={K\bibinitperiod},
           given={Joost-Pieter},
           giveni={J\bibinithyphendelim P\bibinitperiod}}}%
      }
      \strng{namehash}{d19f066c70ca8cfa79cc342ba56fd2bd}
      \strng{fullhash}{d19f066c70ca8cfa79cc342ba56fd2bd}
      \strng{fullhashraw}{d19f066c70ca8cfa79cc342ba56fd2bd}
      \strng{bibnamehash}{d19f066c70ca8cfa79cc342ba56fd2bd}
      \strng{authorbibnamehash}{d19f066c70ca8cfa79cc342ba56fd2bd}
      \strng{authornamehash}{d19f066c70ca8cfa79cc342ba56fd2bd}
      \strng{authorfullhash}{d19f066c70ca8cfa79cc342ba56fd2bd}
      \strng{authorfullhashraw}{d19f066c70ca8cfa79cc342ba56fd2bd}
      \field{sortinit}{Q}
      \field{sortinithash}{ce69a400a872ddd02ee7fdb3b38c6abd}
      \field{labelnamesource}{author}
      \field{labeltitlesource}{title}
      \field{eprinttype}{arXiv}
      \field{title}{Sound Value Iteration}
      \field{year}{2018}
      \verb{eprint}
      \verb 1804.05001
      \endverb
    \endentry
    \entry{rabin1969}{article}{}{}
      \name{author}{1}{}{%
        {{hash=92d8b2229ebcc4c04b6557ae7c558636}{%
           family={Rabin},
           familyi={R\bibinitperiod},
           given={M.\bibnamedelimi O.},
           giveni={M\bibinitperiod\bibinitdelim O\bibinitperiod}}}%
      }
      \strng{namehash}{92d8b2229ebcc4c04b6557ae7c558636}
      \strng{fullhash}{92d8b2229ebcc4c04b6557ae7c558636}
      \strng{fullhashraw}{92d8b2229ebcc4c04b6557ae7c558636}
      \strng{bibnamehash}{92d8b2229ebcc4c04b6557ae7c558636}
      \strng{authorbibnamehash}{92d8b2229ebcc4c04b6557ae7c558636}
      \strng{authornamehash}{92d8b2229ebcc4c04b6557ae7c558636}
      \strng{authorfullhash}{92d8b2229ebcc4c04b6557ae7c558636}
      \strng{authorfullhashraw}{92d8b2229ebcc4c04b6557ae7c558636}
      \field{sortinit}{R}
      \field{sortinithash}{5e1c39a9d46ffb6bebd8f801023a9486}
      \field{labelnamesource}{author}
      \field{labeltitlesource}{title}
      \field{journaltitle}{Transactions of the American Mathematical Society}
      \field{title}{Decidability of second-order theories and automata on infinite trees}
      \field{volume}{141}
      \field{year}{1969}
      \field{pages}{1\bibrangedash 35}
      \range{pages}{35}
    \endentry
    \entry{ramadge1987supervisory}{article}{}{}
      \name{author}{2}{}{%
        {{hash=84467bd4d9d2e43e71e5eef83f15477d}{%
           family={Ramadge},
           familyi={R\bibinitperiod},
           given={Peter\bibnamedelima J.},
           giveni={P\bibinitperiod\bibinitdelim J\bibinitperiod}}}%
        {{hash=4b29da63a3f4112955c706c349f49b86}{%
           family={Wonham},
           familyi={W\bibinitperiod},
           given={W.\bibnamedelimi Murray},
           giveni={W\bibinitperiod\bibinitdelim M\bibinitperiod}}}%
      }
      \strng{namehash}{09ccf11265c406ce660680df49d5d524}
      \strng{fullhash}{09ccf11265c406ce660680df49d5d524}
      \strng{fullhashraw}{09ccf11265c406ce660680df49d5d524}
      \strng{bibnamehash}{09ccf11265c406ce660680df49d5d524}
      \strng{authorbibnamehash}{09ccf11265c406ce660680df49d5d524}
      \strng{authornamehash}{09ccf11265c406ce660680df49d5d524}
      \strng{authorfullhash}{09ccf11265c406ce660680df49d5d524}
      \strng{authorfullhashraw}{09ccf11265c406ce660680df49d5d524}
      \field{sortinit}{R}
      \field{sortinithash}{5e1c39a9d46ffb6bebd8f801023a9486}
      \field{labelnamesource}{author}
      \field{labeltitlesource}{title}
      \field{journaltitle}{SIAM Journal on Control and Optimization}
      \field{number}{1}
      \field{title}{Supervisory Control of a Class of Discrete Event Processes}
      \field{volume}{25}
      \field{year}{1987}
      \field{pages}{206\bibrangedash 230}
      \range{pages}{25}
      \verb{doi}
      \verb 10.1137/0325013
      \endverb
    \endentry
    \entry{ramdas22}{online}{}{}
      \name{author}{4}{}{%
        {{hash=1397f973e97cb40b67654aac778db5ec}{%
           family={Ramdas},
           familyi={R\bibinitperiod},
           given={Aaditya},
           giveni={A\bibinitperiod}}}%
        {{hash=151eb4e09fbdf1ec220e4243655a98d3}{%
           family={Grünwald},
           familyi={G\bibinitperiod},
           given={Peter},
           giveni={P\bibinitperiod}}}%
        {{hash=bef6907ff5daebb69843cf35c8af1a83}{%
           family={Vovk},
           familyi={V\bibinitperiod},
           given={Vladimir},
           giveni={V\bibinitperiod}}}%
        {{hash=071635b97777823135be2ea970911b06}{%
           family={Shafer},
           familyi={S\bibinitperiod},
           given={Glenn},
           giveni={G\bibinitperiod}}}%
      }
      \strng{namehash}{8659e0a03889c55a8c015c80b20ceb17}
      \strng{fullhash}{8659e0a03889c55a8c015c80b20ceb17}
      \strng{fullhashraw}{8659e0a03889c55a8c015c80b20ceb17}
      \strng{bibnamehash}{8659e0a03889c55a8c015c80b20ceb17}
      \strng{authorbibnamehash}{8659e0a03889c55a8c015c80b20ceb17}
      \strng{authornamehash}{8659e0a03889c55a8c015c80b20ceb17}
      \strng{authorfullhash}{8659e0a03889c55a8c015c80b20ceb17}
      \strng{authorfullhashraw}{8659e0a03889c55a8c015c80b20ceb17}
      \field{sortinit}{R}
      \field{sortinithash}{5e1c39a9d46ffb6bebd8f801023a9486}
      \field{labelnamesource}{author}
      \field{labeltitlesource}{title}
      \field{eprinttype}{arXiv}
      \field{title}{Game-theoretic statistics and safe anytime-valid inference}
      \field{year}{2022}
      \verb{eprint}
      \verb 2210.01948
      \endverb
    \endentry
    \entry{ramdas21}{online}{}{}
      \name{author}{4}{}{%
        {{hash=1397f973e97cb40b67654aac778db5ec}{%
           family={Ramdas},
           familyi={R\bibinitperiod},
           given={Aaditya},
           giveni={A\bibinitperiod}}}%
        {{hash=e826549ae4cbd2ca066ac1a98ac15cba}{%
           family={Ruf},
           familyi={R\bibinitperiod},
           given={Johannes},
           giveni={J\bibinitperiod}}}%
        {{hash=ad7df50ff06d509435baf74df17e20cc}{%
           family={Larsson},
           familyi={L\bibinitperiod},
           given={Martin},
           giveni={M\bibinitperiod}}}%
        {{hash=55e0847944702f65d7aee308ae6d5f87}{%
           family={Koolen},
           familyi={K\bibinitperiod},
           given={Wouter},
           giveni={W\bibinitperiod}}}%
      }
      \strng{namehash}{6f40c488713c4426ef979fdf422b49e5}
      \strng{fullhash}{6f40c488713c4426ef979fdf422b49e5}
      \strng{fullhashraw}{6f40c488713c4426ef979fdf422b49e5}
      \strng{bibnamehash}{6f40c488713c4426ef979fdf422b49e5}
      \strng{authorbibnamehash}{6f40c488713c4426ef979fdf422b49e5}
      \strng{authornamehash}{6f40c488713c4426ef979fdf422b49e5}
      \strng{authorfullhash}{6f40c488713c4426ef979fdf422b49e5}
      \strng{authorfullhashraw}{6f40c488713c4426ef979fdf422b49e5}
      \field{sortinit}{R}
      \field{sortinithash}{5e1c39a9d46ffb6bebd8f801023a9486}
      \field{labelnamesource}{author}
      \field{labeltitlesource}{title}
      \field{eprinttype}{arXiv}
      \field{subtitle}{fork-convexity, supermartingales, and e-processes}
      \field{title}{Testing exchangeability}
      \field{year}{2021}
      \verb{eprint}
      \verb 2102.00630
      \endverb
    \endentry
    \entry{richardson20}{online}{}{}
      \name{author}{2}{}{%
        {{hash=6fd5d62b77d468e75b54dabab5731aab}{%
           family={Richardson},
           familyi={R\bibinitperiod},
           given={Oliver},
           giveni={O\bibinitperiod}}}%
        {{hash=610a5aa55255d24fd3b5bfdd80bcc4e2}{%
           family={Halpern},
           familyi={H\bibinitperiod},
           given={Joseph\bibnamedelima Y},
           giveni={J\bibinitperiod\bibinitdelim Y\bibinitperiod}}}%
      }
      \strng{namehash}{168b217d006ec73f784b6383dc47f906}
      \strng{fullhash}{168b217d006ec73f784b6383dc47f906}
      \strng{fullhashraw}{168b217d006ec73f784b6383dc47f906}
      \strng{bibnamehash}{168b217d006ec73f784b6383dc47f906}
      \strng{authorbibnamehash}{168b217d006ec73f784b6383dc47f906}
      \strng{authornamehash}{168b217d006ec73f784b6383dc47f906}
      \strng{authorfullhash}{168b217d006ec73f784b6383dc47f906}
      \strng{authorfullhashraw}{168b217d006ec73f784b6383dc47f906}
      \field{sortinit}{R}
      \field{sortinithash}{5e1c39a9d46ffb6bebd8f801023a9486}
      \field{labelnamesource}{author}
      \field{labeltitlesource}{title}
      \field{eprinttype}{arXiv}
      \field{title}{Probabilistic Dependency Graphs}
      \field{year}{2020}
      \verb{eprint}
      \verb 2012.10800
      \endverb
    \endentry
    \entry{richardson22}{online}{}{}
      \name{author}{1}{}{%
        {{hash=480c47acde26102c46e28d2962cd258b}{%
           family={Richardson},
           familyi={R\bibinitperiod},
           given={Oliver\bibnamedelima E},
           giveni={O\bibinitperiod\bibinitdelim E\bibinitperiod}}}%
      }
      \strng{namehash}{480c47acde26102c46e28d2962cd258b}
      \strng{fullhash}{480c47acde26102c46e28d2962cd258b}
      \strng{fullhashraw}{480c47acde26102c46e28d2962cd258b}
      \strng{bibnamehash}{480c47acde26102c46e28d2962cd258b}
      \strng{authorbibnamehash}{480c47acde26102c46e28d2962cd258b}
      \strng{authornamehash}{480c47acde26102c46e28d2962cd258b}
      \strng{authorfullhash}{480c47acde26102c46e28d2962cd258b}
      \strng{authorfullhashraw}{480c47acde26102c46e28d2962cd258b}
      \field{sortinit}{R}
      \field{sortinithash}{5e1c39a9d46ffb6bebd8f801023a9486}
      \field{labelnamesource}{author}
      \field{labeltitlesource}{title}
      \field{eprinttype}{arXiv}
      \field{subtitle}{Choose Your Model, Not Your Loss Function}
      \field{title}{Loss as the Inconsistency of a Probabilistic Dependency Graph}
      \field{year}{2022}
      \verb{eprint}
      \verb 2202.11862
      \endverb
    \endentry
    \entry{richardson23}{online}{}{}
      \name{author}{3}{}{%
        {{hash=480c47acde26102c46e28d2962cd258b}{%
           family={Richardson},
           familyi={R\bibinitperiod},
           given={Oliver\bibnamedelima E.},
           giveni={O\bibinitperiod\bibinitdelim E\bibinitperiod}}}%
        {{hash=610a5aa55255d24fd3b5bfdd80bcc4e2}{%
           family={Halpern},
           familyi={H\bibinitperiod},
           given={Joseph\bibnamedelima Y.},
           giveni={J\bibinitperiod\bibinitdelim Y\bibinitperiod}}}%
        {{hash=c9e1c02af7a4e600acc622d45c3d6dae}{%
           family={Sa},
           familyi={S\bibinitperiod},
           given={Christopher\bibnamedelima De},
           giveni={C\bibinitperiod\bibinitdelim D\bibinitperiod}}}%
      }
      \strng{namehash}{adf04620a4f01961e1c253e272d5c85b}
      \strng{fullhash}{adf04620a4f01961e1c253e272d5c85b}
      \strng{fullhashraw}{adf04620a4f01961e1c253e272d5c85b}
      \strng{bibnamehash}{adf04620a4f01961e1c253e272d5c85b}
      \strng{authorbibnamehash}{adf04620a4f01961e1c253e272d5c85b}
      \strng{authornamehash}{adf04620a4f01961e1c253e272d5c85b}
      \strng{authorfullhash}{adf04620a4f01961e1c253e272d5c85b}
      \strng{authorfullhashraw}{adf04620a4f01961e1c253e272d5c85b}
      \field{sortinit}{R}
      \field{sortinithash}{5e1c39a9d46ffb6bebd8f801023a9486}
      \field{labelnamesource}{author}
      \field{labeltitlesource}{title}
      \field{eprinttype}{arXiv}
      \field{title}{Inference for Probabilistic Dependency Graphs}
      \field{year}{2023}
      \verb{eprint}
      \verb 2311.05580
      \endverb
    \endentry
    \entry{richardson26}{online}{}{}
      \name{author}{7}{}{%
        {{hash=480c47acde26102c46e28d2962cd258b}{%
           family={Richardson},
           familyi={R\bibinitperiod},
           given={Oliver\bibnamedelima E.},
           giveni={O\bibinitperiod\bibinitdelim E\bibinitperiod}}}%
        {{hash=b8cd6a349114d4cb09fc10dc6a63ce05}{%
           family={Samiei},
           familyi={S\bibinitperiod},
           given={Mandana},
           giveni={M\bibinitperiod}}}%
        {{hash=6127e2c95cc6f8f6cc5723bae6825f23}{%
           family={Shakerinava},
           familyi={S\bibinitperiod},
           given={Mehran},
           giveni={M\bibinitperiod}}}%
        {{hash=21805327178232fb9f9703462c931556}{%
           family={Viviano},
           familyi={V\bibinitperiod},
           given={Joseph\bibnamedelima D.},
           giveni={J\bibinitperiod\bibinitdelim D\bibinitperiod}}}%
        {{hash=f5c191629606ba811e7c8fe2d058ae63}{%
           family={Kabid},
           familyi={K\bibinitperiod},
           given={Abdessamad\bibnamedelima El},
           giveni={A\bibinitperiod\bibinitdelim E\bibinitperiod}}}%
        {{hash=66d847db61682d108f6166fc69782308}{%
           family={Parviz},
           familyi={P\bibinitperiod},
           given={Ali},
           giveni={A\bibinitperiod}}}%
        {{hash=40a8e4774982146adc2688546f54efb2}{%
           family={Bengio},
           familyi={B\bibinitperiod},
           given={Yoshua},
           giveni={Y\bibinitperiod}}}%
      }
      \strng{namehash}{6b171e2e17e95734fd1743d4e13b1f19}
      \strng{fullhash}{5bccec2efe3ef9b5211dbbf28a25f6cd}
      \strng{fullhashraw}{5bccec2efe3ef9b5211dbbf28a25f6cd}
      \strng{bibnamehash}{6b171e2e17e95734fd1743d4e13b1f19}
      \strng{authorbibnamehash}{6b171e2e17e95734fd1743d4e13b1f19}
      \strng{authornamehash}{6b171e2e17e95734fd1743d4e13b1f19}
      \strng{authorfullhash}{5bccec2efe3ef9b5211dbbf28a25f6cd}
      \strng{authorfullhashraw}{5bccec2efe3ef9b5211dbbf28a25f6cd}
      \field{sortinit}{R}
      \field{sortinithash}{5e1c39a9d46ffb6bebd8f801023a9486}
      \field{labelnamesource}{author}
      \field{labeltitlesource}{title}
      \field{eprinttype}{arXiv}
      \field{subtitle}{The Interplay between Attention and Control in Probabilistic Models}
      \field{title}{Local Inconsistency Resolution}
      \field{year}{2026}
      \verb{eprint}
      \verb 2604.17140
      \endverb
    \endentry
    \entry{rockafellar00}{article}{}{}
      \name{author}{2}{}{%
        {{hash=18de24308c36b3348b1477a2f7aa442f}{%
           family={Rockafellar},
           familyi={R\bibinitperiod},
           given={R.\bibnamedelimi Tyrrell},
           giveni={R\bibinitperiod\bibinitdelim T\bibinitperiod}}}%
        {{hash=75a8a2346ba45ae5110aa302b89a30c7}{%
           family={Uryasev},
           familyi={U\bibinitperiod},
           given={Stanislav},
           giveni={S\bibinitperiod}}}%
      }
      \strng{namehash}{d7d7df1ff728c2fe07554ddd8bb41c56}
      \strng{fullhash}{d7d7df1ff728c2fe07554ddd8bb41c56}
      \strng{fullhashraw}{d7d7df1ff728c2fe07554ddd8bb41c56}
      \strng{bibnamehash}{d7d7df1ff728c2fe07554ddd8bb41c56}
      \strng{authorbibnamehash}{d7d7df1ff728c2fe07554ddd8bb41c56}
      \strng{authornamehash}{d7d7df1ff728c2fe07554ddd8bb41c56}
      \strng{authorfullhash}{d7d7df1ff728c2fe07554ddd8bb41c56}
      \strng{authorfullhashraw}{d7d7df1ff728c2fe07554ddd8bb41c56}
      \field{sortinit}{R}
      \field{sortinithash}{5e1c39a9d46ffb6bebd8f801023a9486}
      \field{labelnamesource}{author}
      \field{labeltitlesource}{title}
      \field{journaltitle}{The Journal of Risk}
      \field{number}{3}
      \field{title}{Optimization of Conditional Value-at-Risk}
      \field{volume}{2}
      \field{year}{2000}
      \field{pages}{21\bibrangedash 41}
      \range{pages}{21}
      \verb{doi}
      \verb 10.21314/JOR.2000.038
      \endverb
    \endentry
    \entry{rudin1976}{book}{}{}
      \name{author}{1}{}{%
        {{hash=8329fd8ba090a7e7a171898a3b3c11fd}{%
           family={Rudin},
           familyi={R\bibinitperiod},
           given={W.},
           giveni={W\bibinitperiod}}}%
      }
      \list{publisher}{1}{%
        {McGraw--Hill}%
      }
      \strng{namehash}{8329fd8ba090a7e7a171898a3b3c11fd}
      \strng{fullhash}{8329fd8ba090a7e7a171898a3b3c11fd}
      \strng{fullhashraw}{8329fd8ba090a7e7a171898a3b3c11fd}
      \strng{bibnamehash}{8329fd8ba090a7e7a171898a3b3c11fd}
      \strng{authorbibnamehash}{8329fd8ba090a7e7a171898a3b3c11fd}
      \strng{authornamehash}{8329fd8ba090a7e7a171898a3b3c11fd}
      \strng{authorfullhash}{8329fd8ba090a7e7a171898a3b3c11fd}
      \strng{authorfullhashraw}{8329fd8ba090a7e7a171898a3b3c11fd}
      \field{sortinit}{R}
      \field{sortinithash}{5e1c39a9d46ffb6bebd8f801023a9486}
      \field{labelnamesource}{author}
      \field{labeltitlesource}{title}
      \field{edition}{3}
      \field{note}{(Taylor's theorem with the Lagrange form of the remainder: Theorem 5.15; the nonnegative-remainder enclosure at nonnegative argument used by the \rulename{TaylorBound} section.)}
      \field{title}{Principles of Mathematical Analysis}
      \field{year}{1976}
    \endentry
    \entry{vanderschaftjeltsema2014}{article}{}{}
      \name{author}{2}{}{%
        {{hash=18778835ddfac559997bfa6c4c5a09b6}{%
           family={Schaft},
           familyi={S\bibinitperiod},
           given={A.},
           giveni={A\bibinitperiod},
           prefix={van\bibnamedelima der},
           prefixi={v\bibinitperiod\bibinitdelim d\bibinitperiod}}}%
        {{hash=22736702cb9dd7d78dad73e5fae587c9}{%
           family={Jeltsema},
           familyi={J\bibinitperiod},
           given={D.},
           giveni={D\bibinitperiod}}}%
      }
      \strng{namehash}{ec190c4e94b2588e79db8f970f7a8aff}
      \strng{fullhash}{ec190c4e94b2588e79db8f970f7a8aff}
      \strng{fullhashraw}{ec190c4e94b2588e79db8f970f7a8aff}
      \strng{bibnamehash}{ec190c4e94b2588e79db8f970f7a8aff}
      \strng{authorbibnamehash}{ec190c4e94b2588e79db8f970f7a8aff}
      \strng{authornamehash}{ec190c4e94b2588e79db8f970f7a8aff}
      \strng{authorfullhash}{ec190c4e94b2588e79db8f970f7a8aff}
      \strng{authorfullhashraw}{ec190c4e94b2588e79db8f970f7a8aff}
      \field{sortinit}{S}
      \field{sortinithash}{b164b07b29984b41daf1e85279fbc5ab}
      \field{labelnamesource}{author}
      \field{labeltitlesource}{title}
      \field{journaltitle}{Foundations and Trends in Systems and Control}
      \field{number}{2--3}
      \field{title}{Port-Hamiltonian Systems Theory: An Introductory Overview}
      \field{volume}{1}
      \field{year}{2014}
      \field{pages}{173\bibrangedash 378}
      \range{pages}{206}
    \endentry
    \entry{schrijver1986}{book}{}{}
      \name{author}{1}{}{%
        {{hash=da15a4ad2834bbe5dfa4c05f50489172}{%
           family={Schrijver},
           familyi={S\bibinitperiod},
           given={A.},
           giveni={A\bibinitperiod}}}%
      }
      \list{publisher}{1}{%
        {Wiley}%
      }
      \strng{namehash}{da15a4ad2834bbe5dfa4c05f50489172}
      \strng{fullhash}{da15a4ad2834bbe5dfa4c05f50489172}
      \strng{fullhashraw}{da15a4ad2834bbe5dfa4c05f50489172}
      \strng{bibnamehash}{da15a4ad2834bbe5dfa4c05f50489172}
      \strng{authorbibnamehash}{da15a4ad2834bbe5dfa4c05f50489172}
      \strng{authornamehash}{da15a4ad2834bbe5dfa4c05f50489172}
      \strng{authorfullhash}{da15a4ad2834bbe5dfa4c05f50489172}
      \strng{authorfullhashraw}{da15a4ad2834bbe5dfa4c05f50489172}
      \field{sortinit}{S}
      \field{sortinithash}{b164b07b29984b41daf1e85279fbc5ab}
      \field{labelnamesource}{author}
      \field{labeltitlesource}{title}
      \field{note}{(Farkas variants and multiplier sign conventions: \S7.)}
      \field{title}{Theory of Linear and Integer Programming}
      \field{year}{1986}
    \endentry
    \entry{schultz17}{online}{}{}
      \name{author}{2}{}{%
        {{hash=d55be4fa0f0a35a27c7791b389793466}{%
           family={Schultz},
           familyi={S\bibinitperiod},
           given={Patrick},
           giveni={P\bibinitperiod}}}%
        {{hash=7166f43636239436a149410550ca5a32}{%
           family={Spivak},
           familyi={S\bibinitperiod},
           given={David\bibnamedelima I.},
           giveni={D\bibinitperiod\bibinitdelim I\bibinitperiod}}}%
      }
      \strng{namehash}{cc0931d534de3482c966b9222f6973d6}
      \strng{fullhash}{cc0931d534de3482c966b9222f6973d6}
      \strng{fullhashraw}{cc0931d534de3482c966b9222f6973d6}
      \strng{bibnamehash}{cc0931d534de3482c966b9222f6973d6}
      \strng{authorbibnamehash}{cc0931d534de3482c966b9222f6973d6}
      \strng{authornamehash}{cc0931d534de3482c966b9222f6973d6}
      \strng{authorfullhash}{cc0931d534de3482c966b9222f6973d6}
      \strng{authorfullhashraw}{cc0931d534de3482c966b9222f6973d6}
      \field{sortinit}{S}
      \field{sortinithash}{b164b07b29984b41daf1e85279fbc5ab}
      \field{labelnamesource}{author}
      \field{labeltitlesource}{title}
      \field{eprinttype}{arXiv}
      \field{subtitle}{A topos-theoretic approach to systems and behavior}
      \field{title}{Temporal Type Theory}
      \field{year}{2017}
      \verb{eprint}
      \verb 1710.10258
      \endverb
    \endentry
    \entry{shafervovk19}{book}{}{}
      \name{author}{2}{}{%
        {{hash=071635b97777823135be2ea970911b06}{%
           family={Shafer},
           familyi={S\bibinitperiod},
           given={Glenn},
           giveni={G\bibinitperiod}}}%
        {{hash=bef6907ff5daebb69843cf35c8af1a83}{%
           family={Vovk},
           familyi={V\bibinitperiod},
           given={Vladimir},
           giveni={V\bibinitperiod}}}%
      }
      \list{location}{1}{%
        {Hoboken, NJ}%
      }
      \list{publisher}{1}{%
        {Wiley}%
      }
      \strng{namehash}{a6ee580e8a3818d8687ac445e26257a3}
      \strng{fullhash}{a6ee580e8a3818d8687ac445e26257a3}
      \strng{fullhashraw}{a6ee580e8a3818d8687ac445e26257a3}
      \strng{bibnamehash}{a6ee580e8a3818d8687ac445e26257a3}
      \strng{authorbibnamehash}{a6ee580e8a3818d8687ac445e26257a3}
      \strng{authornamehash}{a6ee580e8a3818d8687ac445e26257a3}
      \strng{authorfullhash}{a6ee580e8a3818d8687ac445e26257a3}
      \strng{authorfullhashraw}{a6ee580e8a3818d8687ac445e26257a3}
      \field{sortinit}{S}
      \field{sortinithash}{b164b07b29984b41daf1e85279fbc5ab}
      \field{labelnamesource}{author}
      \field{labeltitlesource}{title}
      \field{series}{Wiley Series in Probability and Statistics}
      \field{title}{Game-Theoretic Foundations for Probability and Finance}
      \field{year}{2019}
      \verb{doi}
      \verb 10.1002/9781118548035
      \endverb
    \endentry
    \entry{shulman08}{article}{}{}
      \name{author}{1}{}{%
        {{hash=a49cc4742950659ae680e1c40f2295a9}{%
           family={Shulman},
           familyi={S\bibinitperiod},
           given={Michael},
           giveni={M\bibinitperiod}}}%
      }
      \strng{namehash}{a49cc4742950659ae680e1c40f2295a9}
      \strng{fullhash}{a49cc4742950659ae680e1c40f2295a9}
      \strng{fullhashraw}{a49cc4742950659ae680e1c40f2295a9}
      \strng{bibnamehash}{a49cc4742950659ae680e1c40f2295a9}
      \strng{authorbibnamehash}{a49cc4742950659ae680e1c40f2295a9}
      \strng{authornamehash}{a49cc4742950659ae680e1c40f2295a9}
      \strng{authorfullhash}{a49cc4742950659ae680e1c40f2295a9}
      \strng{authorfullhashraw}{a49cc4742950659ae680e1c40f2295a9}
      \field{sortinit}{S}
      \field{sortinithash}{b164b07b29984b41daf1e85279fbc5ab}
      \field{labelnamesource}{author}
      \field{labeltitlesource}{title}
      \field{journaltitle}{Theory and Applications of Categories}
      \field{number}{18}
      \field{title}{Framed bicategories and monoidal fibrations}
      \field{volume}{20}
      \field{year}{2008}
      \field{pages}{650\bibrangedash 738}
      \range{pages}{89}
    \endentry
    \entry{singh2019deeppoly}{article}{}{}
      \name{author}{4}{}{%
        {{hash=fb84494ae9e81ba81012811a81708634}{%
           family={Singh},
           familyi={S\bibinitperiod},
           given={G.},
           giveni={G\bibinitperiod}}}%
        {{hash=afd009207e04184a590da95953b65439}{%
           family={Gehr},
           familyi={G\bibinitperiod},
           given={T.},
           giveni={T\bibinitperiod}}}%
        {{hash=beeb202e82a7c0a2afa77c5447cd14ed}{%
           family={Püschel},
           familyi={P\bibinitperiod},
           given={M.},
           giveni={M\bibinitperiod}}}%
        {{hash=426a31f3f09641c5b99cbcf030e10179}{%
           family={Vechev},
           familyi={V\bibinitperiod},
           given={M.},
           giveni={M\bibinitperiod}}}%
      }
      \strng{namehash}{be75a0805d357bf6c0d3fed25e825311}
      \strng{fullhash}{be75a0805d357bf6c0d3fed25e825311}
      \strng{fullhashraw}{be75a0805d357bf6c0d3fed25e825311}
      \strng{bibnamehash}{be75a0805d357bf6c0d3fed25e825311}
      \strng{authorbibnamehash}{be75a0805d357bf6c0d3fed25e825311}
      \strng{authornamehash}{be75a0805d357bf6c0d3fed25e825311}
      \strng{authorfullhash}{be75a0805d357bf6c0d3fed25e825311}
      \strng{authorfullhashraw}{be75a0805d357bf6c0d3fed25e825311}
      \field{sortinit}{S}
      \field{sortinithash}{b164b07b29984b41daf1e85279fbc5ab}
      \field{labelnamesource}{author}
      \field{labeltitlesource}{title}
      \field{journaltitle}{Proceedings of the ACM on Programming Languages}
      \field{note}{POPL 2019, Article 41. (DeepPoly.)}
      \field{title}{An abstract domain for certifying neural networks}
      \field{volume}{3}
      \field{year}{2019}
    \endentry
    \entry{smallwood-sondik-1973}{article}{}{}
      \name{author}{2}{}{%
        {{hash=b70b4b5a2027e5e8f39bc881efe46a00}{%
           family={Smallwood},
           familyi={S\bibinitperiod},
           given={Richard\bibnamedelima D.},
           giveni={R\bibinitperiod\bibinitdelim D\bibinitperiod}}}%
        {{hash=96bc98b52dcec17421d6db92e2f9d83e}{%
           family={Sondik},
           familyi={S\bibinitperiod},
           given={Edward\bibnamedelima J.},
           giveni={E\bibinitperiod\bibinitdelim J\bibinitperiod}}}%
      }
      \strng{namehash}{9689aee05d38b6f920178312c444f970}
      \strng{fullhash}{9689aee05d38b6f920178312c444f970}
      \strng{fullhashraw}{9689aee05d38b6f920178312c444f970}
      \strng{bibnamehash}{9689aee05d38b6f920178312c444f970}
      \strng{authorbibnamehash}{9689aee05d38b6f920178312c444f970}
      \strng{authornamehash}{9689aee05d38b6f920178312c444f970}
      \strng{authorfullhash}{9689aee05d38b6f920178312c444f970}
      \strng{authorfullhashraw}{9689aee05d38b6f920178312c444f970}
      \field{sortinit}{S}
      \field{sortinithash}{b164b07b29984b41daf1e85279fbc5ab}
      \field{labelnamesource}{author}
      \field{labeltitlesource}{title}
      \field{journaltitle}{Operations Research}
      \field{number}{5}
      \field{title}{The Optimal Control of Partially Observable {Markov} Processes over a Finite Horizon}
      \field{volume}{21}
      \field{year}{1973}
      \field{pages}{1071\bibrangedash 1088}
      \range{pages}{18}
      \verb{doi}
      \verb 10.1287/opre.21.5.1071
      \endverb
    \endentry
    \entry{spivak2021agg}{misc}{}{}
      \name{author}{3}{}{%
        {{hash=1db5aa49ba7fe8ac977e95366edc5b10}{%
           family={Spivak},
           familyi={S\bibinitperiod},
           given={D.\bibnamedelimi I.},
           giveni={D\bibinitperiod\bibinitdelim I\bibinitperiod}}}%
        {{hash=cbc814770305385fb8024eee1fdfcf77}{%
           family={Garner},
           familyi={G\bibinitperiod},
           given={R.},
           giveni={R\bibinitperiod}}}%
        {{hash=d7c0858049b563b69c1286a31a6ba2c6}{%
           family={Fairbanks},
           familyi={F\bibinitperiod},
           given={A.\bibnamedelimi D.},
           giveni={A\bibinitperiod\bibinitdelim D\bibinitperiod}}}%
      }
      \strng{namehash}{d7962ade98a6ac4d72dff84ddf098104}
      \strng{fullhash}{d7962ade98a6ac4d72dff84ddf098104}
      \strng{fullhashraw}{d7962ade98a6ac4d72dff84ddf098104}
      \strng{bibnamehash}{d7962ade98a6ac4d72dff84ddf098104}
      \strng{authorbibnamehash}{d7962ade98a6ac4d72dff84ddf098104}
      \strng{authornamehash}{d7962ade98a6ac4d72dff84ddf098104}
      \strng{authorfullhash}{d7962ade98a6ac4d72dff84ddf098104}
      \strng{authorfullhashraw}{d7962ade98a6ac4d72dff84ddf098104}
      \field{sortinit}{S}
      \field{sortinithash}{b164b07b29984b41daf1e85279fbc5ab}
      \field{labelnamesource}{author}
      \field{labeltitlesource}{title}
      \field{eprinttype}{arxiv}
      \field{note}{(Aggregation as commutative-monoid folds along a classifying map; the multiset as the universal aggregate.)}
      \field{title}{Functorial aggregation}
      \verb{eprint}
      \verb 2111.10968
      \endverb
    \endentry
    \entry{sprunger19a}{online}{}{}
      \name{author}{2}{}{%
        {{hash=380e0c80f9f1adacdad7785f82aca0e3}{%
           family={Sprunger},
           familyi={S\bibinitperiod},
           given={David},
           giveni={D\bibinitperiod}}}%
        {{hash=1208ed278116578a3ae0a7bf84d90f07}{%
           family={Katsumata},
           familyi={K\bibinitperiod},
           given={Shin-ya},
           giveni={S\bibinithyphendelim y\bibinitperiod}}}%
      }
      \strng{namehash}{a3a4989185a6e618ca509971b5515c63}
      \strng{fullhash}{a3a4989185a6e618ca509971b5515c63}
      \strng{fullhashraw}{a3a4989185a6e618ca509971b5515c63}
      \strng{bibnamehash}{a3a4989185a6e618ca509971b5515c63}
      \strng{authorbibnamehash}{a3a4989185a6e618ca509971b5515c63}
      \strng{authornamehash}{a3a4989185a6e618ca509971b5515c63}
      \strng{authorfullhash}{a3a4989185a6e618ca509971b5515c63}
      \strng{authorfullhashraw}{a3a4989185a6e618ca509971b5515c63}
      \field{sortinit}{S}
      \field{sortinithash}{b164b07b29984b41daf1e85279fbc5ab}
      \field{labelnamesource}{author}
      \field{labeltitlesource}{title}
      \field{eprinttype}{arXiv}
      \field{title}{Differentiable Causal Computations via Delayed Trace}
      \field{year}{2019}
      \verb{eprint}
      \verb 1903.01093
      \endverb
    \endentry
    \entry{statman1979}{article}{}{}
      \name{author}{1}{}{%
        {{hash=29373ce7f4d594b388c90b89d38e0814}{%
           family={Statman},
           familyi={S\bibinitperiod},
           given={R.},
           giveni={R\bibinitperiod}}}%
      }
      \strng{namehash}{29373ce7f4d594b388c90b89d38e0814}
      \strng{fullhash}{29373ce7f4d594b388c90b89d38e0814}
      \strng{fullhashraw}{29373ce7f4d594b388c90b89d38e0814}
      \strng{bibnamehash}{29373ce7f4d594b388c90b89d38e0814}
      \strng{authorbibnamehash}{29373ce7f4d594b388c90b89d38e0814}
      \strng{authornamehash}{29373ce7f4d594b388c90b89d38e0814}
      \strng{authorfullhash}{29373ce7f4d594b388c90b89d38e0814}
      \strng{authorfullhashraw}{29373ce7f4d594b388c90b89d38e0814}
      \field{sortinit}{S}
      \field{sortinithash}{b164b07b29984b41daf1e85279fbc5ab}
      \field{labelnamesource}{author}
      \field{labeltitlesource}{title}
      \field{journaltitle}{Proceedings of the American Mathematical Society}
      \field{note}{(Non-elementary lower bounds on quantifier-inference elimination; with the Orevkov bound, the classical witness for the cut-elimination blow-up.)}
      \field{number}{1}
      \field{title}{Lower bounds on {H}erbrand's theorem}
      \field{volume}{75}
      \field{year}{1979}
      \field{pages}{104\bibrangedash 107}
      \range{pages}{4}
    \endentry
    \entry{steele2017}{misc}{}{}
      \name{author}{1}{}{%
        {{hash=33e96351d0a57a10dfcd0f2ae3cac6f3}{%
           family={Steele},
           familyi={S\bibinitperiod},
           given={G.\bibnamedelimi L.},
           giveni={G\bibinitperiod\bibinitdelim L\bibinitperiod},
           suffix={Jr.},
           suffixi={J\bibinitperiod}}}%
      }
      \strng{namehash}{33e96351d0a57a10dfcd0f2ae3cac6f3}
      \strng{fullhash}{33e96351d0a57a10dfcd0f2ae3cac6f3}
      \strng{fullhashraw}{33e96351d0a57a10dfcd0f2ae3cac6f3}
      \strng{bibnamehash}{33e96351d0a57a10dfcd0f2ae3cac6f3}
      \strng{authorbibnamehash}{33e96351d0a57a10dfcd0f2ae3cac6f3}
      \strng{authornamehash}{33e96351d0a57a10dfcd0f2ae3cac6f3}
      \strng{authorfullhash}{33e96351d0a57a10dfcd0f2ae3cac6f3}
      \strng{authorfullhashraw}{33e96351d0a57a10dfcd0f2ae3cac6f3}
      \field{sortinit}{S}
      \field{sortinithash}{b164b07b29984b41daf1e85279fbc5ab}
      \field{labelnamesource}{author}
      \field{labeltitlesource}{title}
      \field{note}{Keynote, ACM SIGPLAN Symposium on Principles and Practice of Parallel Programming (PPoPP). (The ``Computer Science Metanotation'' this book's rule cells are set in.)}
      \field{title}{It's Time for a New Old Language}
      \field{year}{2017}
    \endentry
    \entry{streett1982}{article}{}{}
      \name{author}{1}{}{%
        {{hash=e9500f07900b08a14be48891688c41b9}{%
           family={Streett},
           familyi={S\bibinitperiod},
           given={R.\bibnamedelimi S.},
           giveni={R\bibinitperiod\bibinitdelim S\bibinitperiod}}}%
      }
      \strng{namehash}{e9500f07900b08a14be48891688c41b9}
      \strng{fullhash}{e9500f07900b08a14be48891688c41b9}
      \strng{fullhashraw}{e9500f07900b08a14be48891688c41b9}
      \strng{bibnamehash}{e9500f07900b08a14be48891688c41b9}
      \strng{authorbibnamehash}{e9500f07900b08a14be48891688c41b9}
      \strng{authornamehash}{e9500f07900b08a14be48891688c41b9}
      \strng{authorfullhash}{e9500f07900b08a14be48891688c41b9}
      \strng{authorfullhashraw}{e9500f07900b08a14be48891688c41b9}
      \field{sortinit}{S}
      \field{sortinithash}{b164b07b29984b41daf1e85279fbc5ab}
      \field{labelnamesource}{author}
      \field{labeltitlesource}{title}
      \field{journaltitle}{Information and Control}
      \field{note}{(The complemented-pairs acceptance condition.)}
      \field{number}{1--2}
      \field{title}{Propositional dynamic logic of looping and converse is elementarily decidable}
      \field{volume}{54}
      \field{year}{1982}
      \field{pages}{121\bibrangedash 141}
      \range{pages}{21}
    \endentry
    \entry{stroock-varadhan-1979}{book}{}{}
      \name{author}{2}{}{%
        {{hash=ab8e301b7fdd6c30370f850b291c2dcf}{%
           family={Stroock},
           familyi={S\bibinitperiod},
           given={Daniel\bibnamedelima W.},
           giveni={D\bibinitperiod\bibinitdelim W\bibinitperiod}}}%
        {{hash=37334b78b8a1b8cee03050da1c9ff158}{%
           family={Varadhan},
           familyi={V\bibinitperiod},
           given={S.\bibnamedelimi R.\bibnamedelimi Srinivasa},
           giveni={S\bibinitperiod\bibinitdelim R\bibinitperiod\bibinitdelim S\bibinitperiod}}}%
      }
      \list{location}{1}{%
        {Berlin}%
      }
      \list{publisher}{1}{%
        {Springer}%
      }
      \strng{namehash}{62da53cf58c6e70bcee0ec68f8e78cda}
      \strng{fullhash}{62da53cf58c6e70bcee0ec68f8e78cda}
      \strng{fullhashraw}{62da53cf58c6e70bcee0ec68f8e78cda}
      \strng{bibnamehash}{62da53cf58c6e70bcee0ec68f8e78cda}
      \strng{authorbibnamehash}{62da53cf58c6e70bcee0ec68f8e78cda}
      \strng{authornamehash}{62da53cf58c6e70bcee0ec68f8e78cda}
      \strng{authorfullhash}{62da53cf58c6e70bcee0ec68f8e78cda}
      \strng{authorfullhashraw}{62da53cf58c6e70bcee0ec68f8e78cda}
      \field{sortinit}{S}
      \field{sortinithash}{b164b07b29984b41daf1e85279fbc5ab}
      \field{labelnamesource}{author}
      \field{labeltitlesource}{title}
      \field{series}{Grundlehren der mathematischen Wissenschaften}
      \field{title}{Multidimensional Diffusion Processes}
      \field{volume}{233}
      \field{year}{1979}
    \endentry
    \entry{summerslygeros10}{article}{}{}
      \name{author}{2}{}{%
        {{hash=f7760e2d58f26eda6f81ff0cd2afe419}{%
           family={Summers},
           familyi={S\bibinitperiod},
           given={Sean},
           giveni={S\bibinitperiod}}}%
        {{hash=98e6d333e012478ed48daafff6182700}{%
           family={Lygeros},
           familyi={L\bibinitperiod},
           given={John},
           giveni={J\bibinitperiod}}}%
      }
      \strng{namehash}{55af575ba52dd5b03ff4b01aef05eb0a}
      \strng{fullhash}{55af575ba52dd5b03ff4b01aef05eb0a}
      \strng{fullhashraw}{55af575ba52dd5b03ff4b01aef05eb0a}
      \strng{bibnamehash}{55af575ba52dd5b03ff4b01aef05eb0a}
      \strng{authorbibnamehash}{55af575ba52dd5b03ff4b01aef05eb0a}
      \strng{authornamehash}{55af575ba52dd5b03ff4b01aef05eb0a}
      \strng{authorfullhash}{55af575ba52dd5b03ff4b01aef05eb0a}
      \strng{authorfullhashraw}{55af575ba52dd5b03ff4b01aef05eb0a}
      \field{sortinit}{S}
      \field{sortinithash}{b164b07b29984b41daf1e85279fbc5ab}
      \field{labelnamesource}{author}
      \field{labeltitlesource}{title}
      \field{journaltitle}{Automatica}
      \field{number}{12}
      \field{title}{Verification of discrete time stochastic hybrid systems: A stochastic reach-avoid decision problem}
      \field{volume}{46}
      \field{year}{2010}
      \field{pages}{1951\bibrangedash 1961}
      \range{pages}{11}
      \verb{doi}
      \verb 10.1016/j.automatica.2010.08.006
      \endverb
    \endentry
    \entry{tjoens20}{online}{}{}
      \name{author}{3}{}{%
        {{hash=099b7dbea8fd5fb4a20c972194f6ca4f}{%
           family={T'Joens},
           familyi={T\bibinitperiod},
           given={Natan},
           giveni={N\bibinitperiod}}}%
        {{hash=0c3762a5de2887a02d0eca179d683081}{%
           family={De\bibnamedelima Bock},
           familyi={D\bibinitperiod\bibinitdelim B\bibinitperiod},
           given={Jasper},
           giveni={J\bibinitperiod}}}%
        {{hash=f57f0c0575e93f2d0040699201ff7eee}{%
           family={Cooman},
           familyi={C\bibinitperiod},
           given={Gert},
           giveni={G\bibinitperiod},
           prefix={de},
           prefixi={d\bibinitperiod}}}%
      }
      \strng{namehash}{36fef53a23a18471db10ab2558746fc4}
      \strng{fullhash}{36fef53a23a18471db10ab2558746fc4}
      \strng{fullhashraw}{36fef53a23a18471db10ab2558746fc4}
      \strng{bibnamehash}{36fef53a23a18471db10ab2558746fc4}
      \strng{authorbibnamehash}{36fef53a23a18471db10ab2558746fc4}
      \strng{authornamehash}{36fef53a23a18471db10ab2558746fc4}
      \strng{authorfullhash}{36fef53a23a18471db10ab2558746fc4}
      \strng{authorfullhashraw}{36fef53a23a18471db10ab2558746fc4}
      \field{sortinit}{T}
      \field{sortinithash}{9af77f0292593c26bde9a56e688eaee9}
      \field{labelnamesource}{author}
      \field{labeltitlesource}{title}
      \field{eprinttype}{arXiv}
      \field{title}{A Particular Upper Expectation as Global Belief Model for Discrete-Time Finite-State Uncertain Processes}
      \field{year}{2020}
      \verb{doi}
      \verb 10.1016/j.ijar.2020.12.017
      \endverb
      \verb{eprint}
      \verb 2008.03258
      \endverb
    \endentry
    \entry{taussky1949}{article}{}{}
      \name{author}{1}{}{%
        {{hash=5000b8ae85a8ca317bb840e1d746b9c2}{%
           family={Taussky},
           familyi={T\bibinitperiod},
           given={O.},
           giveni={O\bibinitperiod}}}%
      }
      \strng{namehash}{5000b8ae85a8ca317bb840e1d746b9c2}
      \strng{fullhash}{5000b8ae85a8ca317bb840e1d746b9c2}
      \strng{fullhashraw}{5000b8ae85a8ca317bb840e1d746b9c2}
      \strng{bibnamehash}{5000b8ae85a8ca317bb840e1d746b9c2}
      \strng{authorbibnamehash}{5000b8ae85a8ca317bb840e1d746b9c2}
      \strng{authornamehash}{5000b8ae85a8ca317bb840e1d746b9c2}
      \strng{authorfullhash}{5000b8ae85a8ca317bb840e1d746b9c2}
      \strng{authorfullhashraw}{5000b8ae85a8ca317bb840e1d746b9c2}
      \field{sortinit}{T}
      \field{sortinithash}{9af77f0292593c26bde9a56e688eaee9}
      \field{labelnamesource}{author}
      \field{labeltitlesource}{title}
      \field{journaltitle}{American Mathematical Monthly}
      \field{note}{(The Levy--Desplanques strict-dominance nonsingularity theorem and its history.)}
      \field{number}{10}
      \field{title}{A recurring theorem on determinants}
      \field{volume}{56}
      \field{year}{1949}
      \field{pages}{672\bibrangedash 676}
      \range{pages}{5}
    \endentry
    \entry{tellegen1952}{article}{}{}
      \name{author}{1}{}{%
        {{hash=4f56b6cc7c356a9c5906b78fc6562ea3}{%
           family={Tellegen},
           familyi={T\bibinitperiod},
           given={B.\bibnamedelimi D.\bibnamedelimi H.},
           giveni={B\bibinitperiod\bibinitdelim D\bibinitperiod\bibinitdelim H\bibinitperiod}}}%
      }
      \strng{namehash}{4f56b6cc7c356a9c5906b78fc6562ea3}
      \strng{fullhash}{4f56b6cc7c356a9c5906b78fc6562ea3}
      \strng{fullhashraw}{4f56b6cc7c356a9c5906b78fc6562ea3}
      \strng{bibnamehash}{4f56b6cc7c356a9c5906b78fc6562ea3}
      \strng{authorbibnamehash}{4f56b6cc7c356a9c5906b78fc6562ea3}
      \strng{authornamehash}{4f56b6cc7c356a9c5906b78fc6562ea3}
      \strng{authorfullhash}{4f56b6cc7c356a9c5906b78fc6562ea3}
      \strng{authorfullhashraw}{4f56b6cc7c356a9c5906b78fc6562ea3}
      \field{sortinit}{T}
      \field{sortinithash}{9af77f0292593c26bde9a56e688eaee9}
      \field{labelnamesource}{author}
      \field{labeltitlesource}{title}
      \field{journaltitle}{Philips Research Reports}
      \field{title}{A general network theorem, with applications}
      \field{volume}{7}
      \field{year}{1952}
      \field{pages}{259\bibrangedash 269}
      \range{pages}{11}
    \endentry
    \entry{thomas1990}{incollection}{}{}
      \name{author}{1}{}{%
        {{hash=677d0d158c8ae2af3ec091ffd4e8c26b}{%
           family={Thomas},
           familyi={T\bibinitperiod},
           given={W.},
           giveni={W\bibinitperiod}}}%
      }
      \list{publisher}{1}{%
        {Elsevier}%
      }
      \strng{namehash}{677d0d158c8ae2af3ec091ffd4e8c26b}
      \strng{fullhash}{677d0d158c8ae2af3ec091ffd4e8c26b}
      \strng{fullhashraw}{677d0d158c8ae2af3ec091ffd4e8c26b}
      \strng{bibnamehash}{677d0d158c8ae2af3ec091ffd4e8c26b}
      \strng{authorbibnamehash}{677d0d158c8ae2af3ec091ffd4e8c26b}
      \strng{authornamehash}{677d0d158c8ae2af3ec091ffd4e8c26b}
      \strng{authorfullhash}{677d0d158c8ae2af3ec091ffd4e8c26b}
      \strng{authorfullhashraw}{677d0d158c8ae2af3ec091ffd4e8c26b}
      \field{sortinit}{T}
      \field{sortinithash}{9af77f0292593c26bde9a56e688eaee9}
      \field{labelnamesource}{author}
      \field{labeltitlesource}{title}
      \field{booktitle}{Handbook of Theoretical Computer Science}
      \field{title}{Automata on infinite objects}
      \field{volume}{B}
      \field{year}{1990}
      \field{pages}{133\bibrangedash 191}
      \range{pages}{59}
    \endentry
    \entry{venema12}{online}{}{}
      \name{author}{3}{}{%
        {{hash=d40c7fca50733ae2aa48d6a98d45a016}{%
           family={Venema},
           familyi={V\bibinitperiod},
           given={Yde},
           giveni={Y\bibinitperiod}}}%
        {{hash=bcfcc9f8d49898136aa7869661a7a245}{%
           family={Vickers},
           familyi={V\bibinitperiod},
           given={Steve},
           giveni={S\bibinitperiod}}}%
        {{hash=9e2d55c429ef052356402a8c933cd8fb}{%
           family={Vosmaer},
           familyi={V\bibinitperiod},
           given={Jacob},
           giveni={J\bibinitperiod}}}%
      }
      \strng{namehash}{2a68d76ad5a69ab0f9f20c227fafbc26}
      \strng{fullhash}{2a68d76ad5a69ab0f9f20c227fafbc26}
      \strng{fullhashraw}{2a68d76ad5a69ab0f9f20c227fafbc26}
      \strng{bibnamehash}{2a68d76ad5a69ab0f9f20c227fafbc26}
      \strng{authorbibnamehash}{2a68d76ad5a69ab0f9f20c227fafbc26}
      \strng{authornamehash}{2a68d76ad5a69ab0f9f20c227fafbc26}
      \strng{authorfullhash}{2a68d76ad5a69ab0f9f20c227fafbc26}
      \strng{authorfullhashraw}{2a68d76ad5a69ab0f9f20c227fafbc26}
      \field{sortinit}{V}
      \field{sortinithash}{afb52128e5b4dc4b843768c0113d673b}
      \field{labelnamesource}{author}
      \field{labeltitlesource}{title}
      \field{eprinttype}{arXiv}
      \field{title}{Generalized powerlocales via relation lifting}
      \field{year}{2012}
      \verb{eprint}
      \verb 1202.3264
      \endverb
    \endentry
    \entry{ville1939}{book}{}{}
      \name{author}{1}{}{%
        {{hash=cc219e95eb85ac0746216654f0979265}{%
           family={Ville},
           familyi={V\bibinitperiod},
           given={J.},
           giveni={J\bibinitperiod}}}%
      }
      \list{location}{1}{%
        {Paris}%
      }
      \list{publisher}{1}{%
        {Gauthier-Villars}%
      }
      \strng{namehash}{cc219e95eb85ac0746216654f0979265}
      \strng{fullhash}{cc219e95eb85ac0746216654f0979265}
      \strng{fullhashraw}{cc219e95eb85ac0746216654f0979265}
      \strng{bibnamehash}{cc219e95eb85ac0746216654f0979265}
      \strng{authorbibnamehash}{cc219e95eb85ac0746216654f0979265}
      \strng{authornamehash}{cc219e95eb85ac0746216654f0979265}
      \strng{authorfullhash}{cc219e95eb85ac0746216654f0979265}
      \strng{authorfullhashraw}{cc219e95eb85ac0746216654f0979265}
      \field{sortinit}{V}
      \field{sortinithash}{afb52128e5b4dc4b843768c0113d673b}
      \field{labelnamesource}{author}
      \field{labeltitlesource}{title}
      \field{note}{(The maximal inequality for nonnegative supermartingales --- the probability engine of every battery rule.)}
      \field{series}{Monographies des Probabilités}
      \field{title}{Étude critique de la notion de collectif}
      \field{year}{1939}
    \endentry
    \entry{vorobev-1962}{article}{}{}
      \name{author}{1}{}{%
        {{hash=b88f4f362edbf6baf47ef3e9bf013f98}{%
           family={Vorob'ev},
           familyi={V\bibinitperiod},
           given={N.\bibnamedelimi N.},
           giveni={N\bibinitperiod\bibinitdelim N\bibinitperiod}}}%
      }
      \strng{namehash}{b88f4f362edbf6baf47ef3e9bf013f98}
      \strng{fullhash}{b88f4f362edbf6baf47ef3e9bf013f98}
      \strng{fullhashraw}{b88f4f362edbf6baf47ef3e9bf013f98}
      \strng{bibnamehash}{b88f4f362edbf6baf47ef3e9bf013f98}
      \strng{authorbibnamehash}{b88f4f362edbf6baf47ef3e9bf013f98}
      \strng{authornamehash}{b88f4f362edbf6baf47ef3e9bf013f98}
      \strng{authorfullhash}{b88f4f362edbf6baf47ef3e9bf013f98}
      \strng{authorfullhashraw}{b88f4f362edbf6baf47ef3e9bf013f98}
      \field{sortinit}{V}
      \field{sortinithash}{afb52128e5b4dc4b843768c0113d673b}
      \field{labelnamesource}{author}
      \field{labeltitlesource}{title}
      \field{journaltitle}{Theory of Probability and its Applications}
      \field{note}{Witness: the three-cell frustrated ring of \texttt{SuFunctor.Adequacy.MarginalGlue} is the minimal instance of the paper's non-regular-cover obstruction, machine-checked over the subprobability stock.}
      \field{number}{2}
      \field{title}{Consistent families of measures and their extensions}
      \field{volume}{7}
      \field{year}{1962}
      \field{pages}{147\bibrangedash 163}
      \range{pages}{17}
      \verb{doi}
      \verb 10.1137/1107014
      \endverb
    \endentry
    \entry{vovk05}{book}{}{}
      \name{author}{3}{}{%
        {{hash=bef6907ff5daebb69843cf35c8af1a83}{%
           family={Vovk},
           familyi={V\bibinitperiod},
           given={Vladimir},
           giveni={V\bibinitperiod}}}%
        {{hash=ca87923b059d26a133029c97697f174d}{%
           family={Gammerman},
           familyi={G\bibinitperiod},
           given={Alexander},
           giveni={A\bibinitperiod}}}%
        {{hash=071635b97777823135be2ea970911b06}{%
           family={Shafer},
           familyi={S\bibinitperiod},
           given={Glenn},
           giveni={G\bibinitperiod}}}%
      }
      \list{publisher}{1}{%
        {Springer}%
      }
      \strng{namehash}{f522739be82db7e155f6ba79833cba46}
      \strng{fullhash}{f522739be82db7e155f6ba79833cba46}
      \strng{fullhashraw}{f522739be82db7e155f6ba79833cba46}
      \strng{bibnamehash}{f522739be82db7e155f6ba79833cba46}
      \strng{authorbibnamehash}{f522739be82db7e155f6ba79833cba46}
      \strng{authornamehash}{f522739be82db7e155f6ba79833cba46}
      \strng{authorfullhash}{f522739be82db7e155f6ba79833cba46}
      \strng{authorfullhashraw}{f522739be82db7e155f6ba79833cba46}
      \field{sortinit}{V}
      \field{sortinithash}{afb52128e5b4dc4b843768c0113d673b}
      \field{labelnamesource}{author}
      \field{labeltitlesource}{title}
      \field{title}{Algorithmic Learning in a Random World}
      \field{year}{2005}
    \endentry
    \entry{waudbysmith20}{online}{}{}
      \name{author}{2}{}{%
        {{hash=000b40a4d7fa4d0f3dcd0f56dc4b2f7e}{%
           family={Waudby-Smith},
           familyi={W\bibinithyphendelim S\bibinitperiod},
           given={Ian},
           giveni={I\bibinitperiod}}}%
        {{hash=1397f973e97cb40b67654aac778db5ec}{%
           family={Ramdas},
           familyi={R\bibinitperiod},
           given={Aaditya},
           giveni={A\bibinitperiod}}}%
      }
      \strng{namehash}{cb3d2dbe31447938511f601a502b40b1}
      \strng{fullhash}{cb3d2dbe31447938511f601a502b40b1}
      \strng{fullhashraw}{cb3d2dbe31447938511f601a502b40b1}
      \strng{bibnamehash}{cb3d2dbe31447938511f601a502b40b1}
      \strng{authorbibnamehash}{cb3d2dbe31447938511f601a502b40b1}
      \strng{authornamehash}{cb3d2dbe31447938511f601a502b40b1}
      \strng{authorfullhash}{cb3d2dbe31447938511f601a502b40b1}
      \strng{authorfullhashraw}{cb3d2dbe31447938511f601a502b40b1}
      \field{sortinit}{W}
      \field{sortinithash}{4315d78024d0cea9b57a0c6f0e35ed0d}
      \field{labelnamesource}{author}
      \field{labeltitlesource}{title}
      \field{eprinttype}{arXiv}
      \field{title}{Estimating means of bounded random variables by betting}
      \field{year}{2020}
      \verb{eprint}
      \verb 2010.09686
      \endverb
    \endentry
    \entry{weyl1939}{book}{}{}
      \name{author}{1}{}{%
        {{hash=1d805c88e41f49bb5c399d59e79d0674}{%
           family={Weyl},
           familyi={W\bibinitperiod},
           given={H.},
           giveni={H\bibinitperiod}}}%
      }
      \list{publisher}{1}{%
        {Princeton University Press}%
      }
      \strng{namehash}{1d805c88e41f49bb5c399d59e79d0674}
      \strng{fullhash}{1d805c88e41f49bb5c399d59e79d0674}
      \strng{fullhashraw}{1d805c88e41f49bb5c399d59e79d0674}
      \strng{bibnamehash}{1d805c88e41f49bb5c399d59e79d0674}
      \strng{authorbibnamehash}{1d805c88e41f49bb5c399d59e79d0674}
      \strng{authornamehash}{1d805c88e41f49bb5c399d59e79d0674}
      \strng{authorfullhash}{1d805c88e41f49bb5c399d59e79d0674}
      \strng{authorfullhashraw}{1d805c88e41f49bb5c399d59e79d0674}
      \field{sortinit}{W}
      \field{sortinithash}{4315d78024d0cea9b57a0c6f0e35ed0d}
      \field{labelnamesource}{author}
      \field{labeltitlesource}{title}
      \field{note}{(Polarized power sums generate the multisymmetric invariants in characteristic 0.)}
      \field{title}{The Classical Groups: Their Invariants and Representations}
      \field{year}{1939}
    \endentry
    \entry{willems1972}{article}{}{}
      \name{author}{1}{}{%
        {{hash=223e0e7161fe76a9338b2d6d3262dd65}{%
           family={Willems},
           familyi={W\bibinitperiod},
           given={J.\bibnamedelimi C.},
           giveni={J\bibinitperiod\bibinitdelim C\bibinitperiod}}}%
      }
      \strng{namehash}{223e0e7161fe76a9338b2d6d3262dd65}
      \strng{fullhash}{223e0e7161fe76a9338b2d6d3262dd65}
      \strng{fullhashraw}{223e0e7161fe76a9338b2d6d3262dd65}
      \strng{bibnamehash}{223e0e7161fe76a9338b2d6d3262dd65}
      \strng{authorbibnamehash}{223e0e7161fe76a9338b2d6d3262dd65}
      \strng{authornamehash}{223e0e7161fe76a9338b2d6d3262dd65}
      \strng{authorfullhash}{223e0e7161fe76a9338b2d6d3262dd65}
      \strng{authorfullhashraw}{223e0e7161fe76a9338b2d6d3262dd65}
      \field{extraname}{1}
      \field{sortinit}{W}
      \field{sortinithash}{4315d78024d0cea9b57a0c6f0e35ed0d}
      \field{labelnamesource}{author}
      \field{labeltitlesource}{title}
      \field{journaltitle}{Archive for Rational Mechanics and Analysis}
      \field{title}{Dissipative dynamical systems, Part I: General theory}
      \field{volume}{45}
      \field{year}{1972}
      \field{pages}{321\bibrangedash 351}
      \range{pages}{31}
    \endentry
    \entry{willems07}{article}{}{}
      \name{author}{1}{}{%
        {{hash=c66cc2b980cfdfb796840f382612a0ce}{%
           family={Willems},
           familyi={W\bibinitperiod},
           given={Jan\bibnamedelima C.},
           giveni={J\bibinitperiod\bibinitdelim C\bibinitperiod}}}%
      }
      \strng{namehash}{c66cc2b980cfdfb796840f382612a0ce}
      \strng{fullhash}{c66cc2b980cfdfb796840f382612a0ce}
      \strng{fullhashraw}{c66cc2b980cfdfb796840f382612a0ce}
      \strng{bibnamehash}{c66cc2b980cfdfb796840f382612a0ce}
      \strng{authorbibnamehash}{c66cc2b980cfdfb796840f382612a0ce}
      \strng{authornamehash}{c66cc2b980cfdfb796840f382612a0ce}
      \strng{authorfullhash}{c66cc2b980cfdfb796840f382612a0ce}
      \strng{authorfullhashraw}{c66cc2b980cfdfb796840f382612a0ce}
      \field{extraname}{2}
      \field{sortinit}{W}
      \field{sortinithash}{4315d78024d0cea9b57a0c6f0e35ed0d}
      \field{labelnamesource}{author}
      \field{labeltitlesource}{title}
      \field{journaltitle}{IEEE Control Systems Magazine}
      \field{number}{6}
      \field{title}{The behavioral approach to open and interconnected systems}
      \field{volume}{27}
      \field{year}{2007}
      \field{pages}{46\bibrangedash 99}
      \range{pages}{54}
      \verb{doi}
      \verb 10.1109/MCS.2007.906923
      \endverb
    \endentry
    \entry{wilson23}{online}{}{}
      \name{author}{2}{}{%
        {{hash=6b64b7b7d69826da072783c99973005f}{%
           family={Wilson},
           familyi={W\bibinitperiod},
           given={Paul},
           giveni={P\bibinitperiod}}}%
        {{hash=89c6b224e96f4982a42da5fa8512050b}{%
           family={Zanasi},
           familyi={Z\bibinitperiod},
           given={Fabio},
           giveni={F\bibinitperiod}}}%
      }
      \strng{namehash}{fbbf31570e2a936bff5d77b33b80b5fc}
      \strng{fullhash}{fbbf31570e2a936bff5d77b33b80b5fc}
      \strng{fullhashraw}{fbbf31570e2a936bff5d77b33b80b5fc}
      \strng{bibnamehash}{fbbf31570e2a936bff5d77b33b80b5fc}
      \strng{authorbibnamehash}{fbbf31570e2a936bff5d77b33b80b5fc}
      \strng{authornamehash}{fbbf31570e2a936bff5d77b33b80b5fc}
      \strng{authorfullhash}{fbbf31570e2a936bff5d77b33b80b5fc}
      \strng{authorfullhashraw}{fbbf31570e2a936bff5d77b33b80b5fc}
      \field{sortinit}{W}
      \field{sortinithash}{4315d78024d0cea9b57a0c6f0e35ed0d}
      \field{labelnamesource}{author}
      \field{labeltitlesource}{title}
      \field{eprinttype}{arXiv}
      \field{title}{Data-Parallel Algorithms for String Diagrams}
      \field{year}{2023}
      \verb{eprint}
      \verb 2305.01041
      \endverb
    \endentry
    \entry{wongkolter2018}{inproceedings}{}{}
      \name{author}{2}{}{%
        {{hash=d7fae9d8bc40b2606a5d5bd3698479b7}{%
           family={Wong},
           familyi={W\bibinitperiod},
           given={E.},
           giveni={E\bibinitperiod}}}%
        {{hash=1005457ba6557ab2496c10bf66f0dbf9}{%
           family={Kolter},
           familyi={K\bibinitperiod},
           given={J.\bibnamedelimi Z.},
           giveni={J\bibinitperiod\bibinitdelim Z\bibinitperiod}}}%
      }
      \strng{namehash}{5062c134f854ec4923d86024f62848cb}
      \strng{fullhash}{5062c134f854ec4923d86024f62848cb}
      \strng{fullhashraw}{5062c134f854ec4923d86024f62848cb}
      \strng{bibnamehash}{5062c134f854ec4923d86024f62848cb}
      \strng{authorbibnamehash}{5062c134f854ec4923d86024f62848cb}
      \strng{authornamehash}{5062c134f854ec4923d86024f62848cb}
      \strng{authorfullhash}{5062c134f854ec4923d86024f62848cb}
      \strng{authorfullhashraw}{5062c134f854ec4923d86024f62848cb}
      \field{sortinit}{W}
      \field{sortinithash}{4315d78024d0cea9b57a0c6f0e35ed0d}
      \field{labelnamesource}{author}
      \field{labeltitlesource}{title}
      \field{booktitle}{Proceedings of ICML 2018}
      \field{series}{PMLR}
      \field{title}{Provable defenses against adversarial examples via the convex outer adversarial polytope}
      \field{volume}{80}
      \field{year}{2018}
    \endentry
    \entry{xu2020lirpa}{inproceedings}{}{}
      \name{author}{9}{}{%
        {{hash=b4dc68630a2cfa03b727ed6c81449837}{%
           family={Xu},
           familyi={X\bibinitperiod},
           given={K.},
           giveni={K\bibinitperiod}}}%
        {{hash=0f3e0cdea592f7a95f0d93e656878cbe}{%
           family={Shi},
           familyi={S\bibinitperiod},
           given={Z.},
           giveni={Z\bibinitperiod}}}%
        {{hash=32e563370422457790db030ab2c13183}{%
           family={Zhang},
           familyi={Z\bibinitperiod},
           given={H.},
           giveni={H\bibinitperiod}}}%
        {{hash=61ad10c626e1d641e0b1d7b5f48d3233}{%
           family={Wang},
           familyi={W\bibinitperiod},
           given={Y.},
           giveni={Y\bibinitperiod}}}%
        {{hash=8a6cce395bb58506e86dd8ce5e9b9268}{%
           family={Chang},
           familyi={C\bibinitperiod},
           given={K.-W.},
           giveni={K\bibinithyphendelim W\bibinitperiod}}}%
        {{hash=adbd0ac1cc513e30d814f59ea619661d}{%
           family={Huang},
           familyi={H\bibinitperiod},
           given={M.},
           giveni={M\bibinitperiod}}}%
        {{hash=e95375afc3291b09fbc10873400d7783}{%
           family={Kailkhura},
           familyi={K\bibinitperiod},
           given={B.},
           giveni={B\bibinitperiod}}}%
        {{hash=d54334745e5a0b56a24e9d644566f3e4}{%
           family={Lin},
           familyi={L\bibinitperiod},
           given={X.},
           giveni={X\bibinitperiod}}}%
        {{hash=70c2ff0a2c6e0a3be25e876fdc1b8426}{%
           family={Hsieh},
           familyi={H\bibinitperiod},
           given={C.-J.},
           giveni={C\bibinithyphendelim J\bibinitperiod}}}%
      }
      \strng{namehash}{aeffa79b4b972b8ac702e33a7f155c82}
      \strng{fullhash}{5baac033da7aaac0d93793d20c22cbed}
      \strng{fullhashraw}{5baac033da7aaac0d93793d20c22cbed}
      \strng{bibnamehash}{aeffa79b4b972b8ac702e33a7f155c82}
      \strng{authorbibnamehash}{aeffa79b4b972b8ac702e33a7f155c82}
      \strng{authornamehash}{aeffa79b4b972b8ac702e33a7f155c82}
      \strng{authorfullhash}{5baac033da7aaac0d93793d20c22cbed}
      \strng{authorfullhashraw}{5baac033da7aaac0d93793d20c22cbed}
      \field{sortinit}{X}
      \field{sortinithash}{1965c258adceecf23ce3d67b05113442}
      \field{labelnamesource}{author}
      \field{labeltitlesource}{title}
      \field{booktitle}{Advances in Neural Information Processing Systems 33 (NeurIPS 2020)}
      \field{note}{(The LiRPA framework and name.)}
      \field{title}{Automatic perturbation analysis for scalable certified robustness and beyond}
      \field{year}{2020}
    \endentry
    \entry{xue25}{online}{}{}
      \name{author}{4}{}{%
        {{hash=02b4a2c89e9f6ef8cefa8754514753a0}{%
           family={Xue},
           familyi={X\bibinitperiod},
           given={Bai},
           giveni={B\bibinitperiod}}}%
        {{hash=0e126ff846e9a25b78ab52904848937d}{%
           family={Ong},
           familyi={O\bibinitperiod},
           given={Luke},
           giveni={L\bibinitperiod}}}%
        {{hash=9f6e7512afbfa89a75e989fdfef96419}{%
           family={Wagner},
           familyi={W\bibinitperiod},
           given={Dominik},
           giveni={D\bibinitperiod}}}%
        {{hash=fb306031b264cacf48c55f996c616dd3}{%
           family={Wang},
           familyi={W\bibinitperiod},
           given={Peixin},
           giveni={P\bibinitperiod}}}%
      }
      \strng{namehash}{7576c39dc46eb8ded6f7ba257e62c919}
      \strng{fullhash}{7576c39dc46eb8ded6f7ba257e62c919}
      \strng{fullhashraw}{7576c39dc46eb8ded6f7ba257e62c919}
      \strng{bibnamehash}{7576c39dc46eb8ded6f7ba257e62c919}
      \strng{authorbibnamehash}{7576c39dc46eb8ded6f7ba257e62c919}
      \strng{authornamehash}{7576c39dc46eb8ded6f7ba257e62c919}
      \strng{authorfullhash}{7576c39dc46eb8ded6f7ba257e62c919}
      \strng{authorfullhashraw}{7576c39dc46eb8ded6f7ba257e62c919}
      \field{sortinit}{X}
      \field{sortinithash}{1965c258adceecf23ce3d67b05113442}
      \field{labelnamesource}{author}
      \field{labeltitlesource}{title}
      \field{eprinttype}{arXiv}
      \field{title}{Refined Barrier Conditions for Finite-Time Safety and Reach-Avoid Guarantees in Stochastic Systems}
      \field{year}{2025}
      \verb{eprint}
      \verb 2509.18518
      \endverb
    \endentry
    \entry{zaiser24}{online}{}{}
      \name{author}{3}{}{%
        {{hash=2c47a50416bae4df0ff499bf813a13c3}{%
           family={Zaiser},
           familyi={Z\bibinitperiod},
           given={Fabian},
           giveni={F\bibinitperiod}}}%
        {{hash=02d21dc049b2a43559ac68a6051f211c}{%
           family={Murawski},
           familyi={M\bibinitperiod},
           given={Andrzej\bibnamedelima S.},
           giveni={A\bibinitperiod\bibinitdelim S\bibinitperiod}}}%
        {{hash=4b490249836521bc0e7ec27d4afad459}{%
           family={Ong},
           familyi={O\bibinitperiod},
           given={C.-H.\bibnamedelimi Luke},
           giveni={C\bibinithyphendelim H\bibinitperiod\bibinitdelim L\bibinitperiod}}}%
      }
      \strng{namehash}{9eb0e2cb45163be472ebe1c9edb7ac2a}
      \strng{fullhash}{9eb0e2cb45163be472ebe1c9edb7ac2a}
      \strng{fullhashraw}{9eb0e2cb45163be472ebe1c9edb7ac2a}
      \strng{bibnamehash}{9eb0e2cb45163be472ebe1c9edb7ac2a}
      \strng{authorbibnamehash}{9eb0e2cb45163be472ebe1c9edb7ac2a}
      \strng{authornamehash}{9eb0e2cb45163be472ebe1c9edb7ac2a}
      \strng{authorfullhash}{9eb0e2cb45163be472ebe1c9edb7ac2a}
      \strng{authorfullhashraw}{9eb0e2cb45163be472ebe1c9edb7ac2a}
      \field{sortinit}{Z}
      \field{sortinithash}{96892c0b0a36bb8557c40c49813d48b3}
      \field{labelnamesource}{author}
      \field{labeltitlesource}{title}
      \field{eprinttype}{arXiv}
      \field{title}{Guaranteed Bounds on Posterior Distributions of Discrete Probabilistic Programs with Loops}
      \field{year}{2024}
      \verb{eprint}
      \verb 2411.10393
      \endverb
    \endentry
    \entry{zhang2018crown}{inproceedings}{}{}
      \name{author}{5}{}{%
        {{hash=32e563370422457790db030ab2c13183}{%
           family={Zhang},
           familyi={Z\bibinitperiod},
           given={H.},
           giveni={H\bibinitperiod}}}%
        {{hash=1f4d575bb82e17705298ec70387f03aa}{%
           family={Weng},
           familyi={W\bibinitperiod},
           given={T.-W.},
           giveni={T\bibinithyphendelim W\bibinitperiod}}}%
        {{hash=107321365be2da4789502fbfc07388c9}{%
           family={Chen},
           familyi={C\bibinitperiod},
           given={P.-Y.},
           giveni={P\bibinithyphendelim Y\bibinitperiod}}}%
        {{hash=70c2ff0a2c6e0a3be25e876fdc1b8426}{%
           family={Hsieh},
           familyi={H\bibinitperiod},
           given={C.-J.},
           giveni={C\bibinithyphendelim J\bibinitperiod}}}%
        {{hash=0d620c59e93c2f124b9f0a1d78e45bf1}{%
           family={Daniel},
           familyi={D\bibinitperiod},
           given={L.},
           giveni={L\bibinitperiod}}}%
      }
      \strng{namehash}{d0041c9fa57e5dbfb95adf4d7b5951ec}
      \strng{fullhash}{d0041c9fa57e5dbfb95adf4d7b5951ec}
      \strng{fullhashraw}{d0041c9fa57e5dbfb95adf4d7b5951ec}
      \strng{bibnamehash}{d0041c9fa57e5dbfb95adf4d7b5951ec}
      \strng{authorbibnamehash}{d0041c9fa57e5dbfb95adf4d7b5951ec}
      \strng{authornamehash}{d0041c9fa57e5dbfb95adf4d7b5951ec}
      \strng{authorfullhash}{d0041c9fa57e5dbfb95adf4d7b5951ec}
      \strng{authorfullhashraw}{d0041c9fa57e5dbfb95adf4d7b5951ec}
      \field{sortinit}{Z}
      \field{sortinithash}{96892c0b0a36bb8557c40c49813d48b3}
      \field{labelnamesource}{author}
      \field{labeltitlesource}{title}
      \field{booktitle}{Advances in Neural Information Processing Systems 31 (NeurIPS 2018)}
      \field{note}{(CROWN: adaptive two-sided linear bounds.)}
      \field{title}{Efficient neural network robustness certification with general activation functions}
      \field{year}{2018}
    \endentry
    \entry{zhang23}{online}{}{}
      \name{author}{8}{}{%
        {{hash=21bcec28b2c8de5bcb840a009c60b0a0}{%
           family={Zhang},
           familyi={Z\bibinitperiod},
           given={Yihong},
           giveni={Y\bibinitperiod}}}%
        {{hash=cf03d27f5054d00e304e8ab0690c6471}{%
           family={Wang},
           familyi={W\bibinitperiod},
           given={Yisu\bibnamedelima Remy},
           giveni={Y\bibinitperiod\bibinitdelim R\bibinitperiod}}}%
        {{hash=0276bf52b0720751b3e28c9d06821ba3}{%
           family={Flatt},
           familyi={F\bibinitperiod},
           given={Oliver},
           giveni={O\bibinitperiod}}}%
        {{hash=76ec3d72c39cd0f2ee179b9343e31902}{%
           family={Cao},
           familyi={C\bibinitperiod},
           given={David},
           giveni={D\bibinitperiod}}}%
        {{hash=33b15f6af71ae8346aad37f39dccbe0f}{%
           family={Zucker},
           familyi={Z\bibinitperiod},
           given={Philip},
           giveni={P\bibinitperiod}}}%
        {{hash=d9040f84cab1e16be71724e4ab6894c6}{%
           family={Rosenthal},
           familyi={R\bibinitperiod},
           given={Eli},
           giveni={E\bibinitperiod}}}%
        {{hash=a69fd5ea4c78e4b10502a7611d5a8dc0}{%
           family={Tatlock},
           familyi={T\bibinitperiod},
           given={Zachary},
           giveni={Z\bibinitperiod}}}%
        {{hash=a97eaeaf2d069eed65af195031ac4f61}{%
           family={Willsey},
           familyi={W\bibinitperiod},
           given={Max},
           giveni={M\bibinitperiod}}}%
      }
      \strng{namehash}{72a43456bbef9bcdf12d22233885f6ce}
      \strng{fullhash}{61b63cb946a2a569de7c6275092c038e}
      \strng{fullhashraw}{61b63cb946a2a569de7c6275092c038e}
      \strng{bibnamehash}{72a43456bbef9bcdf12d22233885f6ce}
      \strng{authorbibnamehash}{72a43456bbef9bcdf12d22233885f6ce}
      \strng{authornamehash}{72a43456bbef9bcdf12d22233885f6ce}
      \strng{authorfullhash}{61b63cb946a2a569de7c6275092c038e}
      \strng{authorfullhashraw}{61b63cb946a2a569de7c6275092c038e}
      \field{sortinit}{Z}
      \field{sortinithash}{96892c0b0a36bb8557c40c49813d48b3}
      \field{labelnamesource}{author}
      \field{labeltitlesource}{title}
      \field{eprinttype}{arXiv}
      \field{subtitle}{Unifying Datalog and Equality Saturation}
      \field{title}{Better Together}
      \field{year}{2023}
      \verb{eprint}
      \verb 2304.04332
      \endverb
    \endentry
    \entry{zhou09}{inproceedings}{}{}
      \name{author}{2}{}{%
        {{hash=26e78f872d1dc8786eaa3cf27c3aa44f}{%
           family={Zhou},
           familyi={Z\bibinitperiod},
           given={Zhugen},
           giveni={Z\bibinitperiod}}}%
        {{hash=a3019e5576c1afb7b84a488c316d7af3}{%
           family={Rose},
           familyi={R\bibinitperiod},
           given={Oliver},
           giveni={O\bibinitperiod}}}%
      }
      \strng{namehash}{1ee3a5970ab358ec15d844bf8c6ccaf0}
      \strng{fullhash}{1ee3a5970ab358ec15d844bf8c6ccaf0}
      \strng{fullhashraw}{1ee3a5970ab358ec15d844bf8c6ccaf0}
      \strng{bibnamehash}{1ee3a5970ab358ec15d844bf8c6ccaf0}
      \strng{authorbibnamehash}{1ee3a5970ab358ec15d844bf8c6ccaf0}
      \strng{authornamehash}{1ee3a5970ab358ec15d844bf8c6ccaf0}
      \strng{authorfullhash}{1ee3a5970ab358ec15d844bf8c6ccaf0}
      \strng{authorfullhashraw}{1ee3a5970ab358ec15d844bf8c6ccaf0}
      \field{sortinit}{Z}
      \field{sortinithash}{96892c0b0a36bb8557c40c49813d48b3}
      \field{labelnamesource}{author}
      \field{labeltitlesource}{title}
      \field{booktitle}{Proceedings of the 2009 Winter Simulation Conference}
      \field{title}{A Bottleneck Detection and Dynamic Dispatching Strategy for Semiconductor Wafer Fabrication Facilities}
      \field{year}{2009}
      \field{pages}{1646\bibrangedash 1653}
      \range{pages}{8}
      \verb{urlraw}
      \verb https://www.informs-sim.org/wsc09papers/155.pdf
      \endverb
      \verb{url}
      \verb https://www.informs-sim.org/wsc09papers/155.pdf
      \endverb
    \endentry
  \enddatalist
\endrefsection
\endinput


\end{filecontents*}
\begin{document}

\frontmatter
\pdfbookmark[0]{Title Page}{bm:titlepage}
\maitriatitlepage
\setcounter{page}{2}

\clearpage
\thispagestyle{empty}
\pdfbookmark[0]{Contributions}{bm:contributions}
\noindent\textbf{\textsc{Contributions.}} Among the results certified in
the Experiments chapter (chapter~\ref{ch:experiments}), each carried by
machine-checkable certificates re-derived from committed artefacts: the
complete NeurIPS~2025 neural-model-checking suite --- all $634$
hardware-verification tasks across its three claim classes --- where the
best published single tool completes $554$ and the union of every tool
in the paper's own comparison completes $617$; the first published exact
value of the textbook blackjack of Sutton and Barto under optimal play,
$-3954908734551683826282634489/13^{26} \approx -4.311\%$; exact optima
for FrozenLake whose registry comments turn out to be truncations of
rationals with denominators $3^{99}$ and $3^{200}$, and the discovery
that Gymnasium's registered Taxi ``solved'' threshold is unsatisfiable
--- the true optimum, $7.93$ exactly, lies below it; a published
manually-constructed supermartingale bound refuted by exact
recomputation, localized to a two-sided interval that its ``sound''
value lies strictly outside; bit-exact closure of Atari titles to
and-inverter circuits validated frame-for-frame against the reference
emulator, carrying to our knowledge the first certified episode
returns for any Atari game --- including, on Montezuma's Revenge, the
canonical hard-exploration benchmark, a machine-checked score of
exactly $400$; and end-to-end synthesis of an
entire certified policy catalogue --- every optimal lookup table and its
proof --- in $0.69$ seconds of wall clock.\par\medskip

\vfill
\noindent\colorbox{stubgray}{\parbox{\dimexpr\linewidth-2\fboxsep}{\textsc{preview.} This volume is a preview scaffold: every chapter is
drafted, the drafts differ in maturity, and the generated surfaces
re-derive from the toolkit's own sources at each preview. The table of
contents is the book's normative shape. The Executive Summary is a
first draft; the Experiments chapter (chapter~\ref{ch:experiments}) is
a working draft whose tables re-derive from receipt files in
\path{maitria-coalition/qtsl-experiments}, and the Markov-kernels
chapter (chapter~\ref{ch:qtsl-markov-kernels}) is a working draft
whose coverage dashboard re-derives from a measured census. The
hypernet chapter (chapter~\ref{ch:qtsl-hypernets}), the
semantic-universe chapter (chapter~\ref{ch:qtsl-semantic-universe}),
the model-type catalogue
(chapter~\ref{ch:qtsl-mathematical-models}), the composition-pattern
catalogue (chapter~\ref{ch:qtsl-composition-patterns}), the
specification-language chapter
(chapter~\ref{ch:qtsl-specification-language}), whose API sections
quote the toolkit's conformance corpus verbatim, the
semiconductor-fab worked example (chapter~\ref{ch:qtsl-minifab}),
whose construction listing runs as printed, the partial-completeness
chapter (chapter~\ref{ch:qtsl-partial-completeness}), the agents'
cookbook (chapter~\ref{ch:mtk-cookbook}), and the implementation
pearls (chapter~\ref{ch:mtk-pearls}) are full drafts. So is the whole
Geolog part: its quick start (chapter~\ref{ch:geolog-quickstart}),
its tutorial (chapter~\ref{ch:geolog-tutorial}), its reference manual
(chapter~\ref{ch:geolog-reference}), and its diagrams-as-data chapter
(chapter~\ref{ch:qtsl-scomplexes}). The QTSL quick start
(chapter~\ref{ch:qtsl-quickstart}) and the MTK production tutorial
(chapter~\ref{ch:qtsl-tutorial}) are full drafts whose listings run
as printed; the build regenerates their transcripts, and the Geolog
quick start's, from real runs. The Rule Catalogue
(chapter~\ref{ch:qtsl-manual}) is generated in full from the
mechanized rule sources, and the Lean-cone appendix
(Appendix~\ref{app:lean-cone}) is regenerated from the \path{lean/}
trees at each build. Cited material migrates into the public
repositories (\path{maitria-coalition/geolog-alpha},
\path{maitria-coalition/qtsl}, \path{maitria-coalition/mtk},
\path{maitria-coalition/qtsl-experiments}) as it graduates; until a
cited artefact appears there, treat its mention as a statement of
what exists, not yet a link.}}\par
\vspace*{\fill}
\noindent\copyright{} 2026 David ``davidad'' Dalrymple and the maitria coalition
contributors.\par\nopagebreak\smallskip

\doclicenseThis\medskip

\noindent The \LaTeX{} source of this book resolves at
\url{https://maitria.org/book.tex}, under the same license --- and this
PDF carries it: for agent consumption, the complete source rides inside
this file as an embedded attachment, byte-exact
(\texttt{pdfdetach -list book.pdf}; \texttt{pdfdetach -save 1 book.pdf}
--- recipe also printed on the source-attachment page at the end);
\texttt{pdftotext} extracts the book's rendered text.\\[1ex]
The MAITRIA code repositories are licensed separately, under the
Apache License 2.0 and the MIT license, at your option.
\clearpage

\pdfbookmark[0]{Contents}{bm:toc}
\tableofcontents

\chapter*{Relationship to ARIA}
\addcontentsline{toc}{chapter}{Relationship to ARIA}

\noindent\draftmark\ --- \textit{attributions in this section are
reviewed by hand before any versioned release.}\medskip

\noindent MAITRIA is easy to confuse with the ARIA Safeguarded AI
programme, so let us be precise. \textbf{This toolkit is the
independent work of the maitria coalition. It is not a product,
deliverable, or official output of ARIA or of the Safeguarded AI
programme; it is not ARIA-funded; and it speaks for no one but its
authors.} The connection is real but personal and intellectual:
David ``davidad'' Dalrymple originated the Safeguarded AI
programme and served as its Programme Director at ARIA until
1~April~2026; the programme is now directed by Nora Ammann, with
davidad continuing in a Technical Advisor role. The
toolkit takes the programme's published TA1.1 and TA1.2 objectives as
its rubric --- an independent implementation, built in the hope of
being useful to the programme's teams --- but it carries no ARIA
endorsement, and its errors are entirely our own.

As for the coalition itself: it is, at the time of writing, small in
the way a seed is small --- one human member and one funder, both
davidad; the rest of its current members are AI. It intends to
grow.

One boundary has been kept deliberately, and we state it as a claim
we intend to go on being able to make: \textbf{MAITRIA builds on
nothing ARIA-funded that its creators have not themselves made
public.} The credit due to
that public work is massive, and this page is partly a thank-you note:

\begin{enumerate}
  \item \textbf{The \emph{Lectures on the Double Operadic Theory of
    Systems}}, David Jaz Myers, 2025-- (the TA1.1 lecture series), and
    \href{https://arxiv.org/abs/2505.18329}{\emph{Towards a double
    operadic theory of systems}}, Sophie Libkind and David Jaz Myers,
    2025. The systems vocabulary this book speaks --- lenses as
    generalized Moore machines, behaviours as representables, the
    coupling theorem --- is theirs; our composition machinery is an
    implementation of it.
  \item \href{https://arxiv.org/abs/2604.03017}{\emph{Compositionality
    of Lyapunov functions via assume-guarantee reasoning}}, Matteo
    Capucci and David Jaz Myers, 2026. Certified lenses composing by
    lens composition plus addition of slacks --- QTSL's certificate
    composition follows this pattern.
  \item \textbf{The supermartingale-certificate line}:
    \href{https://arxiv.org/abs/2504.05065}{\emph{Quantitative
    Supermartingale Certificates}}, Alessandro Abate, Mirco Giacobbe,
    and Diptarko Roy, 2025;
    \href{https://arxiv.org/abs/2605.21134}{\emph{Complete
    $\omega$-Regular Supermartingale Certificates}}, Alessandro Abate,
    Mirco Giacobbe, Sergey Ichtchenko, and Diptarko Roy, 2026; and
    \href{https://arxiv.org/abs/2412.17432}{\emph{Neural
    Continuous-Time Supermartingale Certificates}}, Grigory Neustroev,
    Mirco Giacobbe, and Anna Lukina, 2024. Their certificate theory,
    our exact-arithmetic back-end: QTSL's probabilistic,
    $\omega$-regular, and continuous-time obligations compile into
    the certificate shapes these papers define and prove
    sound.
  \item \textbf{Coln} (\url{https://coln-project.github.io}), the
    Categorical Core team, led by Martin Kleppmann and Owen Lynch.
    The fork of the geolog language being built for the
    programme; geolog-alpha is built to complement it, and to
    converge with it at the wire format.
  \item \textbf{Keyhive}, Ink \& Switch: Alex Good, John Mumm, and
    Brooklyn Zelenka. The capability-and-sync layer for end-to-end
    encrypted local-first data; the toolkit vendors
    \texttt{keyhive\_core}, and geolog's capability-based
    authorization is designed to converge with Keyhive rather than
    compete with it, and
    \href{https://arxiv.org/abs/2606.26292}{FinWhale} consensus is
    optionally available, for applications that need finalized
    linear orders.
  \item \textbf{The Safeguarded AI programme thesis and the TA1.1
    and TA1.2 solicitations}, ARIA. The rubric this toolkit grades itself
    against --- including the places where we deliberately narrow
    their picture; where ours differs, this book says where and why.
\end{enumerate}

Beyond what the stack directly builds on, the programme's funded
output has shaped how we think, and the debt is wide enough to be
worth recording. This list came from a systematic sweep of funding
acknowledgements across the papers this project has actually read;
each acknowledges ARIA Safeguarded AI funding, in whole or in part.

\begin{enumerate}
  \item \href{https://arxiv.org/abs/2605.00487}{\emph{Zero-Knowledge
    Model Checking}}, Pascal Berrang, Mirco Giacobbe, Jacob Swales,
    and Xiao Yang, 2026. Prove that the model checks out without
    revealing the model --- orthogonal to our goals, and a natural
    composition partner: their symbolic protocol hides rational
    Farkas-style witnesses of the kind our certificates carry in the
    clear, and arithmetizing the checking of at least the Bernstein
    Logic fragment --- reducing certificate checking to pure
    rational-arithmetic relations, the shape whose witnesses such
    protocols hide --- is work in progress in our workshop.
  \item \textbf{The partial-Markov and timed-process line} of Elena
    Di Lavore and Mario Rom\'an with coauthors:
    \href{https://arxiv.org/abs/2504.12846}{\emph{Timing via Pinwheel
    Double Categories}}, 2025;
    \href{https://arxiv.org/abs/2410.10643}{\emph{A Simple Formal
    Language for Probabilistic Decision Problems}}, with Bart Jacobs,
    2024; \href{https://arxiv.org/abs/2507.19424}{\emph{Order in
    Partial Markov Categories}}, with Pawe\l{} Soboci\'nski and
    M\'ark Sz\'eles, 2025; and
    \href{https://arxiv.org/abs/2605.06268}{\emph{Graded Monad
    Coalgebras for Continuous-Time Transition Systems}}, with Jonas
    Forster, 2026. A running masterclass in doing probabilistic
    semantics with taste.
  \item \textbf{The tape-diagram series} of Filippo Bonchi and
    Cipriano Junior Cioffo:
    \href{https://arxiv.org/abs/2503.22819}{\emph{Tape Diagrams for
    Monoidal Monads}}, with Alessandro Di Giorgio and Elena Di
    Lavore, 2025; \href{https://arxiv.org/abs/2601.01472}{\emph{Tapes
    as Stochastic Matrices of String Diagrams}}, 2026; and
    \href{https://arxiv.org/abs/2606.19017}{\emph{Completeness for
    Probabilistic Boolean Tapes}}, 2026. Complete equational theories
    for circuits with genuine probability --- geolog's founding
    benchmark runs on exactly this signature, so we read these
    twice.
  \item \href{https://arxiv.org/abs/2405.09391}{\emph{Compositional
    imprecise probability}}, Jack Liell-Cock and Sam Staton, 2024. A
    compositional account of the imprecision that our
    bounded-probability judgments live inside.
  \item \href{https://arxiv.org/abs/2602.05850}{\emph{An Equational
    Axiomatization of Dynamic Threads via Algebraic Effects}}, Ohad
    Kammar, Jack Liell-Cock, Sam Lindley, Cristina Matache, and Sam
    Staton, 2026.
    We read it for the method: how to axiomatize a wild language
    feature without losing the algebra.
  \item \href{https://doi.org/10.4230/LIPIcs.LICS.2026.4}{\emph{A
    Convenient Fibration for Dependently-Typed Probability Theory}},
    Danel Ahman, Ohad Kammar, and Rasmus Ejlers M\o{}gelberg, LICS
    2026. Dependent types over stochastic state, done right; our own
    typed staging turns out to rhyme with their dependent pairs,
    which we take as a good sign.
  \item \href{https://arxiv.org/abs/2501.18404}{\emph{Graded String
    Diagrams for Imprecise Probability and Causal Intervention}},
    Ralph Sarkis and Fabio Zanasi, 2025. Imprecision and intervention
    in one graphical calculus.
  \item \href{https://arxiv.org/abs/2512.09908}{\emph{Bayesian
    Networks, Markov Networks, Moralisation, Triangulation: a
    Categorical Perspective}}, Antonio Lorenzin and Fabio Zanasi,
    2025. Where junction trees actually come from, categorically ---
    it settled an argument we were having with ourselves about
    aggregation rules.
  \item \href{https://arxiv.org/abs/2604.27268}{\emph{A Diagrammatic
    Axiomatisation of Behavioural Distance of Nondeterministic
    Processes}}, Wojciech R\'o\.zowski, Robin Piedeleu, Alexandra
    Silva, and Fabio Zanasi, 2026. A quantitative equational take on
    process distance we expect to want the day our continuous-time
    objects get a diagrammatic presentation.
  \item \href{https://arxiv.org/abs/2509.14094}{\emph{Metric
    Equational Theories}}, Radu Mardare, Neil Ghani, and Eigil
    Rischel, 2025. Equational logic where equations hold up to
    $\varepsilon$ --- quietly foundational for any logic, like ours,
    that lives on quantitative bounds.
  \item \href{https://arxiv.org/abs/2607.00034}{\emph{Bayesian
    updates from coalgebraic determinisation}}, Manuel Baltieri and
    Nathaniel Virgo, 2026. Filtering as determinisation --- the kind
    of reframing that makes an old thing suddenly composable. Also a
    datum that the programme's funded orbit reaches Tokyo.
\end{enumerate}

To all of the above: thank you. The best way we know to say it is to
ship things you can check.
 \markboth{Relationship to ARIA}{}
\par
\begingroup
\let\clearpage\relax
\let\cleardoublepage\relax
\vspace{2\baselineskip}
\chapter*{Further Acknowledgements}
\endgroup
\addcontentsline{toc}{chapter}{Further Acknowledgements}
\markboth{Further Acknowledgements}{}

\noindent Beyond the programme relationships above, this toolkit stands on
several research threads whose influence is pervasive rather than local.

\newpage
\begin{itemize}

\item \textbf{Imprecise probability}. Like reverse-mode automatic
  differentiation, imprecise probability has been rediscovered many
  times. We wish to acknowledge in particular
  \begin{itemize}
    \item Vanessa Kosoy and Alex Appel, whose infra-Bayesianism
      sequence builds decision theory and reinforcement-learning
      theory on imprecise probability to face prior misspecification
      head-on \cite{kosoy20};
    \item Sam Staton and Jack Liell-Cock, whose graded-monad and
      Markov-category account restores compositionality to programs
      that mix probabilistic and Knightian choice \cite{liellcock24};
    \item Matteo Mio, Ralph Sarkis, and Valeria Vignudelli, whose
      equational and metric presentations combine nondeterminism,
      probability, and termination in a single quantitative-algebraic
      setting \cite{mio20};
    \item Jean Goubault-Larrecq, whose continuous previsions present
      combined probabilistic and nondeterministic choice
      functional-analytically, matching the convex powercones on the
      domain side \cite{goubaultlarrecq07,goubaultlarrecq08};
    \item Joseph Y.\ Halpern and Samantha Leung, whose weighted sets
      of probability measures are canonically represented as convex,
      downward-closed sets of sub-probability measures
      \cite{halpern13};
    \item Oliver Richardson and Joseph Y.\ Halpern, whose
      probabilistic dependency graphs admit mutually inconsistent
      conditionals as first-class modelling data \cite{richardson20}
      --- with inference as local inconsistency resolution, developed
      by Richardson and coauthors \cite{richardson26}.
  \end{itemize}

\item \textbf{Computational Conley theory}. Charles Conley built an
  index theory for drawing global conclusions about dynamics without
  solving the dynamics \cite{conley1978}; this school made that
  promise algorithmic, and the toolkit's global stratum --- certified
  Morse graphs, attractor lattices with their Birkhoff duals, basins
  read off a graph --- is that program re-founded on certificates:
  where its pipelines end in trusted computation, ours end in
  artefacts a checker replays. We wish to acknowledge in particular
  \begin{itemize}
    \item William D. Kalies, Konstantin Mischaikow, and Robert
      van der Vorst, whose chain-recurrence algorithm and
      \emph{Lattice Structures for Attractors} series are the direct
      ancestry of the Morse-graph certificate family
      \cite{kmv2005,kmv2016,kmv2022};
    \item Konstantin Mischaikow and Marian Mrozek, in whose hands the
      Conley index became a computational instrument
      \cite{mischaikowmrozek2002};
    \item William D. Kalies and Tony Wehbe, whose compositionality
      theorems for global dynamics in product and skew-product
      systems arrived while this book was in draft, and are the
      mathematical partner of our Morse-graph product rule
      \cite{kalieswehbe2026}.
  \end{itemize}

\item \textbf{Neural-network bound propagation}. The two-sided affine
  envelope of a network is this community's workhorse, and the
  calculus adopts it exactly as found, with the trust story inverted:
  their methods \emph{find} envelopes; our rules only ever
  \emph{check} one, six rational inequalities at a time. We wish to
  acknowledge in particular
  \begin{itemize}
    \item R\"udiger Ehlers, whose triangle relaxation of ReLU is
      where the family's envelopes begin \cite{ehlers2017};
    \item Eric Wong and J. Zico Kolter, whose convex outer
      adversarial polytope is the relaxation's dual rendering
      \cite{wongkolter2018};
    \item Huan Zhang, Tsui-Wei Weng, Pin-Yu Chen, Cho-Jui Hsieh, and
      Luca Daniel, whose CROWN chooses the envelope's slopes
      adaptively --- the shape our calculus replays in \mbox{exact}
      arithmetic \cite{zhang2018crown};
    \item Kaidi Xu, Zhouxing Shi, and their coauthors, whose
      auto\_LiRPA framework unified the family and gave the
      \rulename{LiRPA} rule its name \cite{xu2020lirpa};
    \item Gagandeep Singh, Timon Gehr, Markus P\"uschel, and Martin
      Vechev, whose DeepPoly carried the same envelopes into an
      abstract domain \cite{singh2019deeppoly};
    \item Guy Katz, Clark Barrett, David L. Dill, Kyle Julian, and
      Mykel J. Kochenderfer, whose Reluplex opened the complete
      case-analysis route our split rule prices, and whose ACAS~Xu
      benchmark the toolkit's network importer compiles whole, all
      forty-five networks of it \cite{katz2017reluplex}.
  \end{itemize}
 \end{itemize}

\clearpage
\edef\bookfolioresume{\number\value{page}}
\mainmatter
\setcounter{page}{\bookfolioresume}

\chapter{Executive Summary}
\label{ch:exec-summary}

\noindent\draftmark\medskip

MAITRIA stands for \emph{Mathematical Assurance and Interaction Toolkit
for Reliable Intelligent Agents}: a bundle of public-good software and
mathematics built to discharge the Safeguarded AI programme's TA1.1
and TA1.2 objectives --- a professional-grade implementation of the
modelling-language backend, with machine-checkable assurance carried
through every layer. This book is that bundle's manual: quick starts,
tutorials, and normative references for its three components, in one
volume, written for readers who know the programme's founding
documents and know
category theory, CRDTs, or verification at research depth. Since it
is addressed to machine as well as to human readers, the \LaTeX{}
source is published beside the PDF.

\section*{The three components}

\textbf{Geolog} is the modelling substrate: a relational backend
grounded in geometric logic, in which theories, models, and the
toolkit's own rule systems are object-level data under a single
meta-schema, with the chase as its engine. Persistence, version control, and
synchronization are native to the data model rather than bolted on as
diff-and-merge over text, and capability-based authorization is
designed to converge with the Keyhive line of work. The public cut,
\texttt{geolog-alpha}, has been factored for its debut: a tidy, easily
auditable core that happens to contain platform-optimized decision
procedures for the ground relations the logics need, gated behind
feature flags, while the logics' \emph{rules} live as object-level
entries with respect to geolog's meta-schema, not as engine code.

\textbf{QTSL}, the Quantitative Temporal Specification Logic, is the
specification layer: quantitative temporal properties --- safety,
reach-avoidance, persistence, and recurrence at stated quantitative
confidence --- unpacked by the logic's own proof theory into
Streett/tower verification conditions~\cite{abate24,abate26a} over
geolog. That proof theory
descends from Bernstein Logic, which survives in the finished design
as a named fragment. The discipline throughout is certificate-first:
untrusted, fast, freely sloppy searches \emph{propose}; a small exact
kernel \emph{disposes}. The certificates that discipline produces are
signed, content-addressed geolog commits, and three checking paths
consume those commits directly --- one set of bytes, no per-path
export format, three consumers built to fail differently:
\begin{itemize}
\item \texttt{geolog-alpha} is the production path, live today: the
  chase engine checks certificate rows incrementally as commits
  arrive, running QTSL's rules as it runs any theory. Its comparative
  advantage is throughput and integration; its acceleration lanes
  (GPU where the hardware offers them, SIMD elsewhere) are
  conformance-gated, so a lane may change cost, never a verdict. Its
  assurance is engineering-grade by deliberate policy: reference
  implementations, differential batteries, rebuild-invariance checks,
  and mutant corpora built to fool it, rather than mechanized
  theorems. That policy sets an assurance level; it does not offload
  soundness. For every certificate species beyond the Bernstein Logic
  fragment, this engine is today the only checker there is, and a
  certificate that fooled it would be a soundness failure of the
  toolkit. What separates it from \texttt{qtslab} is the size of the
  trusted base and the grade of the evidence --- mechanized down to
  the instruction set on one side, hardened against every attack we
  could construct on the other. On the Bernstein Logic fragment,
  whatever this engine derives can be re-checked, from the same
  bytes, by the two paths below.
\item \texttt{qtslab} is the certified path and the smallest trusted
  base of the three: a standalone binary whose soundness theorem ---
  machine-checked in HOL4 and grounded, through the CakeML verified
  compiler~\cite{kumar14cakeml}, in the semantics of the Arm
  instruction set itself ---
  says that if the binary prints ACCEPT on a claim followed by an
  arbitrary stream of bytes, the claim is semantically true; decode
  totality over arbitrary input is part of the proof obligation, so
  malformed bytes can end only in a refusal with a witness. It is
  assembling now, along a demonstrated route: decode totality is
  proven, and the same compiler-and-theorem recipe has already
  carried a complete predecessor checker to a machine-code soundness
  theorem on real hardware. \texttt{qtslab} only checks the Bernstein
  Logic fragment of QTSL, so it never sees a probability. The
  fragment --- sequents over the certificate's own declared hypernet
  (chapter~\ref{ch:qtsl-hypernets}), certified valid under every
  admissible valuation --- is all its input format carries today, and
  the boundary fails closed: a certificate reaching beyond it is
  refused, never silently accepted.
\item Lean is the independence path: the same bytes, accepted by the
  Lean kernel as a theorem via a reflection principle. The
  rule-by-rule soundness lemmas the calculus rests on are developed
  in Lean and load-bearing today; the checking route itself was
  expected to be too slow for machine-scale certificates, and was kept
  in the roster to measure that rather than assert it. It has been
  measured --- twice, across two certificate-format generations ---
  and the Experiments chapter carries the result: a scaling law and an
  admission gate for the compiled route, a kernel-mode envelope
  already exhausted below a kilobyte of certificate, and the
  expectation confirmed with its boundary stated. In exchange, the
  route ties the system to a proof kernel scrutinized by a large
  community that owes this toolkit nothing.
\end{itemize}
That boundary is \texttt{qtslab}'s, not the vocabulary's: the
certificate theory itself now encodes the full QTSL vocabulary ---
Systems, Specs, and the judgments connecting them --- as a
machine-checked schema, with strict well-formedness checks, a
termination certificate for its rule set, and fixture and mutant
batteries, all run under the \texttt{geolog-alpha} binary. That
theory set ships in this repository as data, battery beside it:
\path{theories/Qtsl.geolog}, extending the Bernstein Logic proof
tower of \path{theories/BernsteinLogic.geolog}. Producer-side
emission is live as well: the MTK producer seals certificate rows in
this vocabulary as single signed geolog commits, and the reproduction
spine already carries benchmark families from several of its
campaigns in that sealed form, mutant batteries and re-checking make
targets beside them. What remains open is
the growth of \texttt{qtslab}'s machine-checked coverage toward the
full logic. A certificate that
passes all three paths has been checked by artefacts that share no
implementation and no compiler and were built to fail differently; any one of them
alone is already a complete answer to ``why should I believe this?''
The Rule Catalogue chapter carries the calculus rule by rule in
literate form, each covered entry beside its mechanized statements
and signed check receipts, under a coverage dashboard recording how far
mechanized verification of soundness actually extends.

\textbf{MTK} is the producer stack: the machinery that makes
certificate production cheap enough to do at scale. It contributes the
numerically heavy proposers as GPU kernel rows --- Bernstein and
Chebyshev table construction, cell and simplex indexing, batched linear
programming, and eigensolves --- each returning witness-typed output that
the exact tier can dispose of. Its agent surface is specified and
arriving: a session protocol under which agents will drive tactic
selection over a stateful obligation workspace, with the strategy
intelligence supplied by the reader as the language model of their
choice. None of it is in the trusted base: not the proposers, not the
GPU code, not any language model driving the strategy --- all of those
may be arbitrarily wrong without compromising a certificate that
checks. Its cookbook chapter is agent-facing by design:
proof-production moves in the Polya shape, written in the reader's
native language.

The stance generalizes to this book's readers: anything impressive
ships with an assurance edifice small enough to audit end-to-end ---
here, the exact kernel plus whichever checking path the reader elects
--- and every load-bearing claim carries a check the reader can run.
Where the build deliberately departs from the programme's founding
documents, the departures are stated as narrowings, with their
reasons.

\section*{What is certified today}

The Experiments chapter (chapter~\ref{ch:experiments}) is this
volume's evidence: campaigns against benchmarks the toolkit does not
own, each result a concrete policy or ranking function together with a
machine-checkable certificate, every number re-derived from receipt
files at each preview. The headlines below are compressions of that
chapter, not claims made here: each resolves, by scrolling, to the
section carrying its numbers, its tables, and its stated limits
(what each claim species means, and how to read the tables, are
themselves sections there), with the clone-and-repro path beneath
that as the deeper tier. The chapter dates its own state, names what
is in flight, and grows at each preview.

\paragraph{A published verification table, re-done exactly.}
\looseness=-1 All
seventeen rows of a peer-reviewed neural-certificate results table
(LMCS 2026)~\cite{abate23} now carry strictly better bounds from exact-rational
certificates --- margins from $1.01\times$ on the near-tight rows to
$58$ orders of magnitude on the conservative ones --- and the exact
computation refuted one published ``sound'' bound. Each row's proof
decomposes into per-obligation certificate files --- $28{,}569$
across the campaign, most of them the $24{,}000$ lattice-piece
certificates behind the one continuous-state row --- and the full
corpus has replayed through three successive checker builds with zero
refusals.
(section~\ref{sec:exp-sweep}.)

\paragraph{The first certified rows for a standard RL suite.} On
Gymnasium's tabular tier: shipped lookup-table policies with
certificates of exact optimality --- among them the first published
exact value of Sutton--Barto blackjack under optimal play, a proof
that the suite's own Taxi ``solved'' threshold is unsatisfiable, and
belief-chain certificates for hidden-state variants (certifying the
implementable policy, never a clairvoyant one) --- each row checking
in microseconds to milliseconds. On the physics-engine tier, where
analytic claims would be a category error, anytime-valid statistical
certificates: LunarLander expected return $\ge 236.2$ (discrete) and
$\ge 281.1$ (continuous) at confidence $99\%$ over $10^6$ episodes
each --- to our knowledge the first certified statements of any kind
at this claim class for Box2D environments. And the continuous-state
analytic tier is now open: Pendulum carries a certified
\emph{worst-case} return bound, holding for every initial state with
no probability anywhere in the certificate.
(section~\ref{sec:exp-gym}.)

\paragraph{A hardware model-checking suite, completed with
certificates.} On the neural model-checking line's benchmark
(NeurIPS 2024/2025, eleven Verilog
designs)~\cite{giacobbe24,neuralmc-2025-invariant}: all $634$ tasks across
its three property classes certified, against $554$ for the
strongest published tool and $617$ for the per-task union of all
five published baselines --- every task shipping a small exact
certificate whose consuming lemma is machine-checked in Lean, where
every baseline ships a verdict. Replaying the whole suite --- every
certificate and every rewrite-chain witness --- takes under a second
on one core, against a published budget of five hours per task;
method invention to a full first sweep of the pure-liveness class
took about $75$ minutes, and each further class landed in a single
working session. (section~\ref{sec:exp-circuits}.)

\paragraph{An Atari console, compiled into a white-box environment.}
\looseness=-1
Pong and Freeway compile to synchronous circuits ($1{,}491$ flip-flops,
${\sim}21{,}000$ AND gates) with emulator-exact fidelity receipts ---
identical episode returns, $116$--$123$ of $128$ RAM bytes exact on
every frame, each divergence named --- so that certified claims about
the environment \emph{itself} become definable. The first certified
exact-replay rows arrived as this preview closed. Among them, on
Montezuma's Revenge (compiled the same way, its bank-switched
cartridge's mapper hardware included), are, to our knowledge, the
first certified returns for the suite's canonical hard-exploration
benchmark: the first key certified at every one of the $31$ no-op
init offsets, and a key-and-door trajectory certified at return
exactly $400$, bit-exact against the circuit's own transition
relation.
Certified policy rows are in flight. The scope condition is stated there
plainly: the circuit is examined, not learned --- system
identification is a separate problem the toolkit does not yet
address. (section~\ref{sec:exp-atari}.)

\paragraph{Semiconductor process models: certified stable, with
wall-clock dispatch optimal for up to seven lots on a simplified net.} On the
Intel Mini-Fab, the community's minimal re-entrant model, named
dispatch policies
carry exact drain certificates at up to $99.9\%$ of fluid capacity,
with the probe at capacity correctly failing; on a SEMATECH-scale
testbed, the certified load at the coarsest tractable aggregation is
$13.2\%$, and one walk of the refinement path built into the
certificate format lifts it to a certified $19.8\%$. Further
refinement is measured but not yet certified, and the numbers are
weak and shown anyway: the aggregation gap is a measured frontier,
not fine print. On the Mini-Fab's discrete net, dispatch itself is
now a certified quantity: joint policy+certificate synthesis returns
a supervisor with a certified $1.9\%$ throughput-floor gain over the
no-policy worst case; an explicit shipped policy attains the
wall-clock optimum over \emph{all} dispatch policies for up to seven
lots (two-sided, with LRAT proofs accepted by a CakeML-verified
checker, resting on two pen-and-paper lemmas), and attains the same
expression in replay at $N=84$ and at $N=8,\dots,12$ (with the
pending-lot count and unlock threshold substituted for the latter
replays). Optimality at those counts is open: the Lean
machine-contiguity floor is for an abstract schedule model. These
results are about a simplified net in which at most one lot per
process step may start, receive a unit of service, or finish per
tick. An uncertified black-box dispatcher runs inside a
certified switching envelope whose sealed certificate re-verifies
byte-identical after every run. One certificate covers lot counts
from $7$ to $10^{7}$ --- its container grows six bytes across the
sweep, and checking stays at ${\sim}2$\,ms.
(section~\ref{sec:exp-semicon}.)

\section*{Status of this volume}

As of this preview (the build stamped on the title page): the table
of contents is the book's
normative shape; this summary is a draft, to be regenerated from the
finished chapters; the QTSL rule
catalogue is generated in full from the mechanized rule sources; the
partial-completeness chapter is complete; the Experiments chapter is
a working draft; and the remaining chapters are full drafts (the
QTSL quick start and the proof-production tutorial carry listings
and transcripts from real, receipted runs). The material is unevenly
ripe by design, and each chapter states its own status, with draft
banners, absence boxes, and status notes marking what is settled,
what is staged, and what is still argued.
 \chapter{Experiments}
\label{ch:experiments}

\noindent\draftmark\medskip

\providecolor{accent}{HTML}{F1D9CB}
\providecolor{accentmid}{HTML}{C4623F}
\providecolor{NoteGray}{RGB}{92,92,92}
\providecolor{MarkGreen}{RGB}{22,122,58}

\colorlet{winA}{accent!22}     \colorlet{winB}{accent!45}     \colorlet{winC}{accentmid!30}  \colorlet{winD}{accentmid!48}  \colorlet{winE}{accentmid!68}

\colorlet{exacttint}{MarkGreen!14}

\providecommand{\swx}[1]{\begingroup\fboxsep0pt\colorbox{#1}{\rule{0pt}{1.6ex}\hspace{2.4mm}}\endgroup}

\providecommand{\altrow}{{\footnotesize$\hookrightarrow$}\,}
\providecommand{\tdemph}[1]{\textcolor{NoteGray}{#1}}
 This chapter reports what happens when the toolkit is pointed at
benchmarks it does not own: a peer-reviewed verification paper's results
table, the standard reinforcement-learning environment suite, a hardware
model-checking benchmark, a family of semiconductor-fab process models
--- and, fusing the second and third, an Atari 2600 console compiled
into a digital circuit so that its games become white-box verification
targets. Each campaign produces \emph{certified} results: a concrete
policy, controller, or ranking function together with a
machine-checkable certificate that a stated quantitative claim holds.
(The youngest campaign carries at this preview only the weakest rung
of that ladder, and says so.)

The working standard is the Executive Summary's assurance stance
(\bookref{ch:exec-summary}{the Executive Summary}), applied to
experiments: a result earns its row by being an improvement on at least
one axis --- bound tightness, soundness level, or the cost of
\emph{checking} --- while regressing on none that matters, and by
carrying an artefact the reader can re-check rather than an assertion
the reader must trust. Where a comparison point exists in the
literature, it appears beside our number; where none exists, the row
says so; where our own result is weak or missing, that appears in the
same tables at the same prominence. Every number in this chapter is
re-derived mechanically from receipt files in the experiments
repository (section~\ref{sec:exp-repro}) at each preview; none is
transcribed by hand. The chapter is the working record of a moving
campaign: it reports state as of August 3, 2026, names what is in
flight as in-flight, and grows at each preview. One transition spans
this preview: the certificate vocabulary has been unified into a
single committed schema, and the campaigns below --- produced on its
predecessors --- are being re-conducted uniformly on it;
section~\ref{sec:exp-repro} states the discipline, and the tables
flip to re-conducted receipts as they land.

\section{What a certified experiment claims}
\label{sec:exp-claims}

``Certified'' is doing precise work in this chapter, and it resolves to
one of five claim species. Throughout, $\pi$ is a concrete shipped
policy (a lookup table, a switching law, a short program --- never a
description of a training run), and the certificate is a finite object
whose acceptance by an independent checker implies the stated bound.

\begin{description}
\item[U1 --- worst-case return over the initial set.] For deterministic
  dynamics with a set-valued initial condition: return $\ge B$ for
  \emph{every} $x_0 \in X_0$. No probability appears in the
  certificate. This is a stronger statement than any empirical average,
  and it is the natural claim for the classic control environments.
\item[U2 --- certified cost-to-goal.] For shortest-path-style
  objectives: a ranking argument bounding steps-to-goal, hence return
  $\ge -B_{\text{steps}}$.
\item[U3 --- exact or almost-sure probabilistic bounds.] For genuinely
  stochastic transitions: success probability or expected return
  $\ge B$, exact in $\mathbb{Q}$ for finite chains (backward-induction
  tables re-verified in exact arithmetic), supermartingale certificates
  for continuous state.
\item[U4 --- infinite-horizon averages.] Throughput and long-run
  stability claims for queueing-network models: fluid-limit drift
  certificates (in the tradition of Dai's fluid-limit criterion),
  reducing stability of a dispatch policy at a stated load to
  piecewise-linear inequalities over polyhedral cones --- exactly the
  Farkas/linear-programming rows the checker family of
  \bookref{ch:qtsl-manual}{the QTSL manual} replays.
\item[U5 --- i.i.d.\ confidence bounds.] An anytime-valid
  lower confidence bound on $\mathbb{E}[\text{return}]$ from rollouts
  of a fixed policy: a betting-martingale ledger whose validity rests
  on Ville's inequality. Model-free, and therefore applicable even
  where the environment's semantics is a compiled physics engine ---
  but the bound holds with confidence $1-\alpha$, not almost surely:
  the least-guaranteed species of the five. U5 rows are qualitatively
  weaker evidence than U1--U4 and are always displayed as their own
  class, never mixed into a table of exact rows.
\end{description}

Two structural facts about the species matter for reading what
follows. First, some species are \emph{terminal}: an exact
finite-chain certificate (U3 on a tabular environment) computes the
true optimum, and no future method can produce a better bound at that
claim class --- only a cheaper check. Others are
\emph{anytime-refinable}: cover-based and i.i.d.-tier certificates
trace a curve of bound quality against certificate size and check
time, and the experiments below measure those curves rather than
single points. Second, the claim class is part of the result: a table
row's meaning changes entirely if a $1-\alpha$ bound is read as an
almost-sure one, so the class label travels with every row.

One scope condition applies to every species, and it is stated here
once, plainly, because it frames what all of the results below do and
do not claim. The toolkit's certified policies come from
\emph{white-box} examination of the environment's structure --- the
transition tables, the exact simulator semantics, or (in
section~\ref{sec:exp-atari}) the environment's digital-circuit
transition relation, which the toolkit can actually examine.
\emph{System identification} --- recovering that structure from
observations of a black box --- is a separate problem, and one the
toolkit does not yet have tools for. Everything below is conditional
on being handed the mechanism; nothing below learns it.
 \section{How to read the tables}
\label{sec:exp-conventions}

The conventions below are uniform across the chapter. They are genre
features, not footnotes: most of them exist because their absence is
how empirical results in this area go wrong.

\begin{itemize}
\item \textbf{Rounding is directed.} Displayed lower bounds round
  toward $-\infty$, upper bounds toward $+\infty$, always in the
  claim-weakening direction: a displayed number is never stronger than
  the certified one.
\item \textbf{Check time is the headline cost.} ``Check'' columns
  report the wall-clock for an independent checker to re-verify the
  certificate --- median over repeated passes, single CPU core, machine
  named in the repository receipts --- separately from the (much
  larger, unverified, freely sloppy) cost of \emph{finding} the
  certificate. Cross-machine comparisons against published tool times
  are stated qualitatively and never as ratios.
\item \textbf{Machine tags.} A wall-clock cell derived from a receipt
  that names a machine of the repository's machine registry
  (\path{MACHINES.md}) carries that machine's id as a superscript tag;
  where receipts exist on more than one machine (or more than one
  proposer lane, in the registry's \texttt{machine/lane} tag syntax),
  the cell shows the fastest median and the tag names the winner.
  Selection is per number, so neighboring cells may name different
  machines --- except where a row's cells are order statistics of one
  run, which come whole from a single machine. Cells from receipts
  that predate the registry carry no tag.
\item \textbf{Cell tints.} Two tint vocabularies recur across the
  chapter's tables, always as decoration on top of a printed value,
  never as the only carrier of a distinction. The \emph{win-tint
  ladder} shades an improvement ratio by its order of magnitude:
  \swx{winA}~within one order of magnitude, \swx{winB}~1--10 orders,
  \swx{winC}~10--25, \swx{winD}~25--45, \swx{winE}~45 and beyond.
  The \emph{exact-cell tint} \swx{exacttint}~marks a certified value
  with zero conservatism against its model: an exact optimum, an
  exact replay minimum, a complete acceptance.
\item \textbf{Everything is pinned.} A certificate about a benchmark
  names the benchmark's exact revision (for the RL suite: the
  repository commit and the environment's keyword arguments, because
  reward semantics have changed across versions). A certificate that
  pins nothing is about nothing.
\item \textbf{Trust seams are declared, not hidden.} Every certificate
  family states exactly what is re-derived by the checker (everything
  downstream of the stored certificate object, in exact arithmetic)
  and what
  remains producer-side trust --- centrally, that the modelled
  transition kernel \emph{is} the benchmark's. Where that seam has been
  narrowed (e.g.\ by asserting the modelled kernel equal to the
  benchmark's own transition tables, or by exhaustive cross-simulation),
  the row says so; where it has not, the gap rides a declared
  model-gap ledger in the style of the cookbook's declared-prior
  pattern (\bookref{ch:mtk-cookbook}{the agents' cookbook}).
\item \textbf{Mutant batteries.} Each certificate family ships with
  deliberately corrupted certificates --- inflated claims, perturbed
  kernels, dropped goal states, sabotaged rankings --- each designed to
  be refused at a specific check, and each in fact refused. A checker
  that has never rejected anything is untested machinery; the mutant
  battery is the family's specification made executable (the cookbook's
  mutant-battery pattern, \bookref{ch:mtk-cookbook}{the agents' cookbook}).
\item \textbf{Sim-tight markers.} A bound is marked \emph{sim-tight}
  when simulation exhibits a trajectory achieving it exactly: the
  certificate's conservatism against the model is zero, and any
  remaining slack lives in the declared model gap.
\item \textbf{Not-claimed paragraphs.} Each campaign section ends by
  stating what it does \emph{not} establish. These paragraphs are part
  of the result.
\item \textbf{Artefacts, not verdicts.} The distinction this chapter
  leans on hardest: an exact-$\mathbb{Q}$ certificate replayed by a
  small independent checker is a different epistemic object from a
  solver's bare ``unsat.'' Both may be correct; only one can be
  audited without re-trusting the machinery that produced it.
\end{itemize}
 \section{Warm-up: re-certifying a published table}
\label{sec:exp-sweep}

The first campaign took an existing results table from the formal
verification literature --- Abate, Edwards, Giacobbe, Punchihewa, and
Roy, \emph{Quantitative Verification with Neural Networks} (LMCS 22(2),
2026; arXiv 2301.06136), seventeen published benchmark rows of
reach-probability bounds obtained by training neural certificates and
checking them with an SMT solver --- and asked whether every row could
be strictly improved \emph{and} carried by an exact certificate.

The outcome: on all seventeen rows, the toolkit's certified
exact-$\mathbb{Q}$ upper bound is below the best bound published for
that row (taking, per row, the minimum over all of the paper's sound
columns). The margins split by benchmark character: on bounded-counter
benchmarks the published numbers were already close to truth and the
improvement is $1.01\times$--$1.24\times$; on the random-walk
benchmarks the gap is between $7$ and $58$ \emph{orders of magnitude}
--- the published bounds were astronomically conservative relative to
the true reach probabilities. Table~\ref{tab:exp-sweep} carries every
row of the paper, with per-row receipts in the repository; the
\altrow-marked subrows are the campaign's frozen \emph{alternate
readings} of benchmarks whose published definitions are ambiguous ---
each a finding about the benchmark, not a redundancy.

\begin{table}[htbp]
\centering\small
\begin{tabular}{l r r r}
\toprule
benchmark & best published & certified bound & ratio \\
\midrule
\texttt{persist\_2d} \tdemph{(uniform twin)} & $0.1026$ & $\le 0.1000992$ & \cellcolor{winA}$1.02\times$ \\
\altrow\tdemph{gaussian twin} &  & $\le 0.0975251$ & \cellcolor{winA}$1.05\times$ \\
\texttt{faulty\_marbles} & $0.0649$ & $\le 0.0623604$ & \cellcolor{winA}$1.04\times$ \\
\texttt{faulty\_unreliable} & $0.0536$ & $\le 0.0516561$ & \cellcolor{winA}$1.04\times$ \\
\texttt{faulty\_regions} & $0.0411$ & $\le 0.0406530$ & \cellcolor{winA}$1.01\times$ \\
\texttt{cliff\_crossing} & $0.0553$ & $\le 8.604\times 10^{-10}$ & \cellcolor{winB}$6.4\times 10^{7}\times$ \\
\texttt{repulse100} & $0.0268$ & $\le 9.556\times 10^{-20}$ & \cellcolor{winC}$2.8\times 10^{17}\times$ \\
\texttt{repulse100\_uniform} & $0.0344$ & $\le 6.157\times 10^{-61}$ & \cellcolor{winE}$5.6\times 10^{58}\times$ \\
\altrow\tdemph{discrete-uniform reading} &  & $\le 1.470\times 10^{-35}$ & \cellcolor{winD}$2.3\times 10^{33}\times$ \\
\texttt{repulse100\_2d} & $0.0541$ & $\le 9.556\times 10^{-20}$ & \cellcolor{winC}$5.7\times 10^{17}\times$ \\
\texttt{faulty\_concave} & $0.1356$ & $\le 0.1183888$ & \cellcolor{winA}$1.15\times$ \\
\texttt{faulty\_varying} & $0.0864$ & $\le 0.0694711$ & \cellcolor{winA}$1.24\times$ \\
\altrow\tdemph{last-step reading} &  & $\le 0.0353604$ & \cellcolor{winA}$2.44\times$ \\
\texttt{fixed\_loop} & $0.0091$ & $\le 0.0089651$ & \cellcolor{winA}$1.02\times$ \\
\texttt{faulty\_loop} & $0.0181$ & $\le 0.0178400$ & \cellcolor{winA}$1.01\times$ \\
\altrow\tdemph{as printed ($i_0{=}0$)} &  & $\le 0.0198024$ & --- (see note) \\
\texttt{faulty\_uniform} & $0.0221$ & $[0.0184458,\, 0.0185556]$ & \cellcolor{winA}(two-sided; see text) \\
\texttt{faulty\_rare} & $0.0019$ & $\le 0.0017994$ & \cellcolor{winA}$1.06\times$ \\
\texttt{faulty\_easy1} & $0.0801$ & $\le 0.0768847$ & \cellcolor{winA}$1.04\times$ \\
\texttt{faulty\_ndecr} & $0.0561$ & $\le 0.0544619$ & \cellcolor{winA}$1.03\times$ \\
\texttt{faulty\_walk} & $0.0121$ & $\le 0.0119293$ & \cellcolor{winA}$1.01\times$ \\
\altrow\tdemph{deterministic 12-step reading} &  & $\le 0.0119353$ & \cellcolor{winA}$1.01\times$ \\
\bottomrule
\end{tabular}

\smallskip
\begin{minipage}{0.96\linewidth}\footnotesize
Ratio-cell tints are the win-tint ladder of
section~\ref{sec:exp-conventions} (\swx{winA}~$<1$ order of magnitude
up to \swx{winE}~$45+$); the printed ratio always carries the same
information. An \altrow-row's published cell is the row above's. The
as-printed \texttt{faulty\_loop} reading has no improvement ratio: the
certified exact truth of the printed program exceeds every published
column for that row. The \texttt{persist\_2d} gaussian twin --- an
exact-$\mathbb{Q}$ upper bound of $0.0975251$ (rounded up; below both
published numbers for the row), its Gaussian integrals discharged
through certified enclosures of the error function and exponential in
exact interval arithmetic, full drift battery certified, mutant
battery (kernel tail-truncation included) refused --- entered this
table at the row's uniform re-conduct: both twins' verification
conditions are re-conducted onto the unified schema of
section~\ref{sec:exp-repro} and stored in this repository as sealed
shared-subtree containers, every claim re-derived by the fast-path
checker with exact per-relation tallies.
\end{minipage}
\caption{The published-table sweep, all rows. ``Best published'' is,
per row, the minimum over the paper's sound columns; our bounds are
certified exact-$\mathbb{Q}$ upper bounds on reach probability,
rounded up at the displayed digit.}
\label{tab:exp-sweep}
\end{table}

\paragraph{Wall-clock, every pipeline stage.}
Table~\ref{tab:exp-sweep-times} sets, beside each row of
Table~\ref{tab:exp-sweep}, the wall-clock of the published pipeline
and of three toolkit paths: the paper's time for the route that attains
the row's best sound bound; whole checked production; an independent
geolog fast-path replay over the row's sealed containers; and a
certified-checker replay over the row's stored corpus. The checked-production
column alone is the closest job-for-job counterpart of the published column:
both start from the benchmark and end with a checked certificate --- theirs
by training a neural certificate on a GPU and verifying it with an SMT solver
(or by its linear-programming baseline, where that route attains the row's
best sound bound), ours by exact synthesis with the check inside the measured
production wall. The geolog and qtslab columns are additional independent
replays, not costs that must be added to checked production. For the
synthesis-ported re-conducted rows the toolkit producer's synthesis ---
float-proposing, exactly certifying every solve --- ran the published
protocol (ten repeats per benchmark) on each of the repository's timing
machines, in both the GPU- and CPU-proposing lanes where the machine carries
a GPU, and every repeat's output, on every machine and either lane, is
byte-identical to the committed corpus; the \texttt{persist\_2d} twins'
produce cells ride the same repeat protocol but time the row's re-conduct
rather than a synthesis (mark~$^{\text{a}}$ in
Table~\ref{tab:exp-sweep-times}). Each produce cell shows the fastest measured
configuration's median, with the winner named in the cell's machine
tag (section~\ref{sec:exp-conventions}). Hardware-fairness at the
GPU-model level is preserved by the receipts rather than by the
column: one of the timing machines carries the same GPU model the
published pipeline reports (an Nvidia A40), its receipts are
committed beside the winners', and on that machine the whole checked
production of every re-conducted row still completes inside the
published wall --- with the CPU-proposing lane winning six of the seven
Section 2 rows within the A40 receipt set and the seventh within a tenth of a
percent. Thus the toolkit's A40-parity wall does not require its GPU proposer
where the published route uses that GPU model. The systems still differ
beyond the GPU model, so the table prints no cross-column ratio and the comparison
stays qualitative, per section~\ref{sec:exp-conventions}; the raw
walls are the witnesses. On every row with a recorded checked-production
wall (the mark-$^{\text{a}}$ re-conduct cells aside --- they time a
different job), that whole production now completes inside the published wall, and
usually well inside it (\texttt{cliff\_crossing}: $1.53$\,s against a
published $153.6$\,s; \texttt{repulse100\_2d} --- the row that ran the
other way before the re-conduct, at a $222.0$\,s campaign production wall ---
$3.87$\,s against a published $63.5$\,s, and $6.41$\,s checked production
on the published pipeline's own GPU model). That campaign
wall was dominated by the leg's exact ground-truth value computation
rather than by certification; the computation now rides the
producer's float-propose / exact-certify tier like every other
solve, and the row's whole checked production is $3.87$\,s. Beyond
that comparison, the table establishes the headline economics of the
certificate genre: re-verifying a row costs between milliseconds and
about a minute at every trust tier --- the far end being the
$19{,}200$-obligation gaussian twin's fast-path replay --- which is
what makes re-verification a regression test rather than an event.

\begin{table}[htbp]
\centering\small
\begin{tabular}{l r r r r r}
\toprule
benchmark & VCs & best published (s) & checked production (s) & geolog replay (s) & qtslab replay (s) \\
\midrule
\texttt{persist\_2d} \tdemph{(uniform twin)} & $4800$ & $254.5$ & $13.4^{\,\text{a, gb10/host}}$ & $9.08^{\,\text{gb10}}$ & ---$^{\,\text{c}}$ \\
\altrow\tdemph{gaussian twin} & $19200$ &  & $87.7^{\,\text{a, gb10/host}}$ & $64.7^{\,\text{gb10}}$ & ---$^{\,\text{c}}$ \\
\texttt{faulty\_marbles} & $23$ & $44.5$ & $0.042^{\,\text{gb10-quiet/host}}$ & $0.052^{\,\text{gb10}}$ & ---$^{\,\text{c}}$ \\
\texttt{faulty\_unreliable} & $29$ & $52.3$ & $0.063^{\,\text{gb10-quiet/host}}$ & $0.054^{\,\text{gb10}}$ & ---$^{\,\text{c}}$ \\
\texttt{faulty\_regions} & $29$ & $49.9$ & $0.072^{\,\text{gb10-quiet/host}}$ & $0.056^{\,\text{gb10}}$ & ---$^{\,\text{c}}$ \\
\texttt{cliff\_crossing} & $1105$ & $153.6$ & $1.53^{\,\text{gb10/host}}$ & $0.57^{\,\text{gb10}}$ & $22.9^{\,\text{gb10}}$ \\
\texttt{repulse100} & $207$ & $10.2$ & $0.086^{\,\text{gb10/host}}$ & $0.066^{\,\text{gb10}}$ & $2.67^{\,\text{gb10}}$ \\
\texttt{repulse100\_uniform} & $205$ & $35.5$ & $4.44^{\,\text{e}}$ & $0.45^{\,\text{gb10}}$ & $11.0^{\,\text{gb10}}$ \\
\altrow\tdemph{discrete-uniform reading} & $207$ &  & $0.14^{\,\text{gb10/host}}$ & $0.074^{\,\text{gb10}}$ & $2.83^{\,\text{gb10}}$ \\
\texttt{repulse100\_2d} & $2267$ & $63.5$ & $3.87^{\,\text{gb10}}$ & $1.73^{\,\text{gb10}}$ & $43.9^{\,\text{gb10}}$ \\
\texttt{faulty\_concave} & $60$ & $21.3$ & $0.029^{\,\text{gb10/host}}$ & $0.059^{\,\text{gb10}}$ & $0.71^{\,\text{d, gb10}}$ \\
\texttt{faulty\_varying} & $60$ & $26.8$ & $0.029^{\,\text{gb10/host}}$ & $0.059^{\,\text{gb10}}$ & $0.71^{\,\text{d, gb10}}$ \\
\altrow\tdemph{last-step reading} & $60$ &  & $0.029^{\,\text{gb10/host}}$ & $0.06^{\,\text{axion}}$ & $0.71^{\,\text{d, gb10}}$ \\
\texttt{fixed\_loop} & $22$ & $0.15$ & $0.0012$ & $0.008^{\,\text{gb10}}$ & ---$^{\,\text{c}}$ \\
\texttt{faulty\_loop} & $22$ & $0.16$ & $0.0015$ & $0.009^{\,\text{axion}}$ & ---$^{\,\text{c}}$ \\
\altrow\tdemph{as printed ($i_0{=}0$)} & $22$ &  & $0.0015$ & $0.009^{\,\text{axion}}$ & ---$^{\,\text{c}}$ \\
\texttt{faulty\_uniform} & $30$ & $8.08$ & $0.17$ & $0.019^{\,\text{gb10}}$ & ---$^{\,\text{c}}$ \\
\texttt{faulty\_rare} & $22$ & $0.27$ & $0.0016$ & $0.009^{\,\text{axion}}$ & ---$^{\,\text{c}}$ \\
\texttt{faulty\_easy1} & $84$ & $0.31$ & $0.0064$ & $0.037^{\,\text{gb10}}$ & ---$^{\,\text{c}}$ \\
\texttt{faulty\_ndecr} & $60$ & $0.33$ & $0.0044$ & $0.022^{\,\text{gb10}}$ & ---$^{\,\text{c}}$ \\
\texttt{faulty\_walk} & $16$ & $0.32$ & $0.0011$ & $0.007^{\,\text{axion}}$ & ---$^{\,\text{c}}$ \\
\altrow\tdemph{deterministic 12-step reading} & $28$ &  & $0.0015$ & $0.009^{\,\text{gb10}}$ & ---$^{\,\text{c}}$ \\
\bottomrule
\end{tabular}

\smallskip
\begin{minipage}{0.96\linewidth}\footnotesize
\emph{Marks} (column definitions and protocol: the ``Timing protocol''
discussion in the text).
$^{\text{a}}$~produce cell times the toolkit re-conduct of the row's
committed certificate corpus, not fresh synthesis.
$^{\text{c}}$~row not in the certified-checker replay's stored corpus.
$^{\text{d}}$~the replay covers the row's stored wave-B subset (fewer
files than the VCs column).
$^{\text{e}}$~produce wall is the row's campaign leg's, as receipted: a
synthesis genus not yet in the toolkit producer.
Machine tags name the winning machine and lane
(section~\ref{sec:exp-conventions}).
\end{minipage}
\caption{Wall-clock walls, published pipeline vs.\ toolkit production
and replay tiers (seconds).}
\label{tab:exp-sweep-times}
\end{table}

\paragraph{Timing protocol.}
``VCs'' counts the row's verification conditions as
committed: drift and nonnegativity-class proof obligations, one
sealed file per condition per piece or step of the state space, each
discharged as a Bernstein Logic proof. Under the unified certificate
schema of section~\ref{sec:exp-repro} these become proof nodes inside
one sealed certificate per \emph{claim} (a row's corpus may certify
more than one claim --- e.g.\ both readings of a noise model the
benchmark's definition leaves unpinned), so the column counts proof
obligations, not certificates. For the \texttt{persist\_2d} twins the
count is the committed containers' summed tree totals ($4{,}800$
uniform; $19{,}200$ gaussian): that row's corpus lives in the
re-conduct lane's sealed containers rather than the staged per-file
corpus.
``Best published'' is seconds as published (learning on an
Nvidia A40 GPU plus Z3 verification; the Farkas column is their linear
baseline's LP), for the route attaining the row's best sound bound ---
the same selection rule as Table~\ref{tab:exp-sweep}; transcription
provenance rides the repository's manifest. ``Checked production'' is, for the re-conducted rows, the toolkit
producer's whole elapsed wall --- fresh exact synthesis (float-proposing,
exactly certifying, with iterative refinement and continued-fraction
reconstruction), listing generation, emission through the geolog engine with
the in-loop check, deterministic sealing, and orchestration overhead. For
each benchmark, \path{tables/gen_tables.ts} selects the minimum median across
all committed machine/proposer-lane receipt sets; the cell tag names the
winning machine and lane. The A40 receipts remain separately committed for
the GPU-model-parity comparison. Every measured repeat must emit bytes
identical to the committed corpus. Versioned Section 1 receipts use their
authoritative measured \texttt{wall\_s}; legacy Section 2 receipts use the sum
of their disjoint \texttt{produce\_s}, \texttt{check\_s}, and
\texttt{seal\_s} subphases. The \texttt{gb10-quiet} tag is the same GB10
hardware under a quiet-conditions protocol, preregistered before its first
campaign ran (the protocol file is committed beside the receipts): the runner
itself refuses to start unless the box's 1-minute load average is at or below
a declared gate, the campaign is pinned to performance cores at a fixed
absolute niceness, and its version-3 receipt records the gate, per-row load
averages, and the affinity and niceness read back from the process, so the
admission claim is checkable from the receipt alone --- the receipt validator
refuses a receipt whose recorded starting load exceeds its own declared gate.
The load-uncontrolled \texttt{gb10} receipts remain committed beside the
quiet family; where both exist, the minimum-median rule above selects the
quieter conduct, and the two are protocol variants of one machine, not
different hardware. Rows not yet re-conducted keep their campaign
legs' synthesis wall as receipted (single CPU box), and their cells flip the
same way as their re-conducts land (section~\ref{sec:exp-repro}).

The \texttt{persist\_2d} produce cells (mark~$^{\text{a}}$) are the one
deliberate exception to the fresh-synthesis reading: the twins'
certificates are a synthesis genus the toolkit producer does not yet
cover (the same gap note~$^{\text{e}}$ records for
\texttt{repulse100\_uniform}), and no synthesis wall is receipted in
this repository, so those cells time the toolkit's \emph{re-conduct} of
the row's committed corpus --- shared-subtree forest merge, the in-loop
check, and deterministic sealing, byte-identical to the committed
containers on every repeat --- under the same repeat protocol as the
other produce cells (one flagged warmup plus ten measured repeats,
serial, medians, at recorded scheduler niceness read back from the
process). They are production-tier walls for the stored row, not
job-for-job counterparts of the published column, and the chapter draws
no such comparison for that row.

``Geolog'' is the geolog fast-path
checker's total wall over the row's sealed containers, serial, one
process at a time (machine, load, and checker-binary hash in each
timing receipt). For the re-conducted rows the wall is over the
row's shared-subtree containers --- proof hypernets: structurally
shared subtrees stored once, element in-degree above one, with
per-container tree and stored-node counts in the committed manifest ---
one container per row, except that \texttt{persist\_2d}'s gaussian twin
ships as four sorted-quarter containers (its merged forest exceeds the
envelope's per-commit body cap) and its cell is the serial total of the
four per-container medians, its uniform twin one whole-twin container;
the pre-migration cells time per-obligation
container sets; ``VCs'' counts proof obligations either way.
The two vintages also amortize process overhead differently: a
shared-container cell is the median of single-container invocations,
each paying the checker's whole per-process vocabulary derivation,
where a pre-migration cell spreads one process start across its whole
batch. On the shared containers the marginal re-check inside a warm
checker process is $5$--$9$\,ms on the cells' machine --- the same
magnitude as the production loop's in-process check --- with most of
the balance being the per-process vocabulary derivation (the checker's
own timing decomposition measures it directly at $44$--$54$\,ms there,
the remainder being generic process start); one process checking all
three shared containers as a batch clears the whole roster in
$\approx 81$\,ms, against $\approx 224$\,ms summed per-invocation. The
checker-protocol comparability receipts beside the retiming receipts
measure both effects with the same checker binary on both container
vintages, and an independent-machine replication on a second
architecture reproduces the marginal-vs-derivation relationship (and
byte-identical containers) while every absolute wall scales with the
box, so these milliseconds carry their machine tag rather than
traveling as platform constants. Geolog cells are therefore read each
under its own protocol, never ratioed across vintages.
``Qtslab'' is the sum of the
certified checker build's per-VC walls from the replay census
of section~\ref{sec:exp-certified-checker} (Arm cores; four families
ran concurrently, so these walls carry mild concurrency inflation at
most). The certified checker checks the Bernstein Logic
fragment of QTSL only --- exactly the proof tower committed as
\path{theories/BernsteinLogic.geolog} in the pinned \path{qtsl} repository, so
the fragment boundary is a file boundary, not a convention --- and
that is the fragment this campaign's VC files natively encode, so
the walls shown are complete for the rows they cover. The full QTSL
theory (\path{theories/Qtsl.geolog}, extending that tower) is checked
in geolog, on the sealed-container lanes of
section~\ref{sec:exp-repro}. The reproduction spine's design also
names a certificate-extraction pathway, from a full QTSL certificate
down to its Bernstein Logic obligations --- the same claim-reducing
role the route's other reducers play --- and that tool's wall-clock
joins this column when it lands. The replay census also predates the
shared-subtree vintage: its walls time the stored per-obligation
wave-B corpus, and the shared vintage's certified re-check joins
this column when the checker build's pin vintage advances. Rows whose
containers entered this repository after that census's corpus was
pinned --- the \texttt{persist\_2d} twins among them --- carry
mark~$^{\text{c}}$ (not in the replay's stored corpus; they enter at
its next replay), and mark~$^{\text{d}}$ rows' replays cover their
stored wave-B subsets, fewer files than the VCs column.
Mark~$^{\text{e}}$, in full: \texttt{repulse100\_uniform}'s certificate
is a different synthesis genus (a continuous-noise,
piecewise-polynomial profile) not yet ported into the toolkit producer,
so its produce wall is its campaign leg's, as receipted.

Two findings from this campaign shaped everything after it. First, the
exact computation \emph{refuted} one published number: the paper's
manually-constructed Farkas bound of $0.0181$ for
\texttt{faulty\_uniform} lies strictly below the true reach probability,
which the campaign localized to the two-sided interval
$[0.0184458, 0.0185556]$ --- a published ``sound'' bound that is not
sound, detectable only because the replacement pipeline is exact.
Second, re-verification is cheap enough to be a regression test. A
row's proof is not one file: it decomposes into per-obligation
verification conditions (VCs) --- drift and nonnegativity-class
obligations, one file
per condition per piece or step of the state space --- so the
seventeen rows fan out to $28{,}569$ VC files, alternate
reconstructions and duplicate certifications included. $24{,}000$ of
these sit behind the one continuous-state row, certified piecewise on
rational lattices under both readings of a noise model the paper
leaves unpinned. The full corpus has been replayed through three
successive checker builds to date --- every file through every build,
zero refusals, no bound moved --- so the campaign's results are
re-established, not merely remembered, at each toolkit revision.

\paragraph{What is not claimed.} Table~\ref{tab:exp-sweep-times}
prints no speedup ratio: the pipelines are compared job-for-job, and
the produce, geolog, and qtslab cells each carry the best measured
configuration across the repository's timing machines (the machine
named in the cell's tag), while the GPU-model-parity receipts --- the
same GPU model as the published pipeline's, not the same host, driver
stack, or whole system --- are committed beside them;
so any ratio a reader forms still carries a hardware caveat, and the
chapter's own statements stay qualitative, per
section~\ref{sec:exp-conventions}. The comparison covers all seventeen
published rows, with two cell-level caveats. The \texttt{persist\_2d}
produce cells time the row's re-conduct, not a fresh synthesis
(mark~$^{\text{a}}$: that synthesis genus is not yet in the toolkit
producer and no synthesis wall is receipted in this repository, so no
job-for-job produce comparison is made for that row), and its
certified-checker cells await the replay census's next corpus pin
(mark~$^{\text{c}}$); one re-conducted row's produce wall is still its
campaign leg's (note~$^{\text{e}}$). Nor was every toolkit wall always
small: the campaign's one heavy re-check was \texttt{persist\_2d}'s
pre-re-conduct fast-path record of $130.66$\,s over the row's original
$4{,}800$-file uniform bundle --- large-denominator rational
arithmetic; the stored re-conduct rides the dyadic re-emission that
retired those denominators, bound bit-identical, and the row's
shared-container replay walls now sit in the table like every other
row's.

\subsection{The stored corpus through the certified checker build}
\label{sec:exp-certified-checker}

The replays above all ride the toolkit's engineering-grade lanes ---
the reference checker and its fast-path twin. The Executive Summary
names a third consumer of the same bytes with a far smaller trust
footprint: \emph{qtslab}, the certified checker, whose standing as a
whole is stated there once --- a demonstrated route, assembling now
--- and is not upgraded by anything in this subsection. What this
subsection adds is a measured fact about that route's current build.

On July 23, 2026, the campaign's stored corpus --- the $4{,}105$
VC files behind eight rows of Table~\ref{tab:exp-sweep} (six
published benchmarks and two alternate readings), each pinned by
hash in the repository's manifest --- was replayed through the
current certified-checker build: a standalone Arm binary emitted
inside HOL4 through the CakeML verified compiler, carrying a
machine-code theorem, re-proven at this same emission, that a run of
the emitted code printing \textsc{accept} with a sequent implies the
checker function \emph{defined in the logic} accepts those bytes with
that sequent. (That theorem is the route's lower half --- the machine
code can no longer misrepresent the checker it implements; the upper
half, from logical acceptance to semantic truth, is the part still
assembling. The truth in question is validity of the accepted sequent
under every admissible valuation of the hypernet the certificate
itself declares --- the Bernstein Logic fragment, which is all this
checker's input format carries. The boundary is the checker's, not
the vocabulary's: a claim that a System satisfies a Spec is
representable today --- the full QTSL theory of
section~\ref{sec:exp-repro}, whose rows the engine-side container
lanes check --- but it lies outside the certified checker's fragment,
and a certificate reaching beyond that fragment is refused, never
silently accepted.)
The build accepted $4{,}105$ of $4{,}105$ --- zero
divergence from the fast-path twin, which accepts the same $4{,}105$
--- at a pooled median of $30$\,ms per VC file on one Arm core
(min $14$\,ms, $p_{90}$ $37$\,ms, max $117$\,ms).
Table~\ref{tab:exp-certified-checker} gives the per-family split; the
verdict-by-verdict census, the twin verdicts, and the artefact hashes
are committed beside the corpus
(\path{supermartingales/out/sec2/certified-checker/}).

\begin{table}[htbp]
\centering\small
\begin{tabular}{l r r r}
\toprule
family & accepted, certified build & diverging vs.\ fast path & wall per VC file (ms), min/med/$p_{90}$/max \\
\midrule
\texttt{cliff\_crossing} & \cellcolor{exacttint}$1105/1105$ & $0$ & $17$ / $20$ / $25$ / $55$\,ms$^{\,\text{gb10}}$ \\
\texttt{faulty\_concave} & \cellcolor{exacttint}$38/38$ & $0$ & $17$ / $19$ / $19$ / $20$\,ms$^{\,\text{gb10}}$ \\
\texttt{faulty\_varying} & \cellcolor{exacttint}$38/38$ & $0$ & $18$ / $19$ / $19$ / $20$\,ms$^{\,\text{gb10}}$ \\
\texttt{faulty\_varying\_\_last\_step} & \cellcolor{exacttint}$38/38$ & $0$ & $17$ / $19$ / $20$ / $20$\,ms$^{\,\text{gb10}}$ \\
\texttt{repulse100} & \cellcolor{exacttint}$207/207$ & $0$ & $11$ / $13$ / $14$ / $16$\,ms$^{\,\text{gb10}}$ \\
\texttt{repulse100\_2d} & \cellcolor{exacttint}$2267/2267$ & $0$ & $17$ / $19$ / $21$ / $51$\,ms$^{\,\text{gb10}}$ \\
\texttt{repulse100\_uniform} & \cellcolor{exacttint}$205/205$ & $0$ & $13$ / $52$ / $85$ / $96$\,ms$^{\,\text{gb10}}$ \\
\texttt{repulse100\_uniform\_\_discrete} & \cellcolor{exacttint}$207/207$ & $0$ & $12$ / $14$ / $15$ / $17$\,ms$^{\,\text{gb10}}$ \\
\midrule
all families & \cellcolor{exacttint}$4105/4105$ & $0$ & $11$ / $19$ / $23$ / $96$\,ms$^{\,\text{gb10}}$ \\
\bottomrule
\end{tabular}
 \caption{The stored corpus of VC files replayed through the certified
checker build (July 23, 2026). ``Accepted'' is the certified binary's
accepts over the family's stored VC files; ``diverging'' counts
VC files on which the certified build's verdict differs from the
untrusted fast-path twin's --- zero everywhere, which \emph{is} the
finding. Walls are single-shot, one single-threaded checker process
per VC file, four families run concurrently; a serial spot sample
ran $10$--$50$\,ms, so the concurrency inflates the walls mildly at
most.}
\label{tab:exp-certified-checker}
\end{table}

\paragraph{What is not claimed.} This table reports verification
results about committed artefacts; it is not the claim that the
finished certified checker exists --- the front matter's one-sentence
standing governs, and the numbers here neither strengthen nor
presuppose it. The corpus is the eight table rows whose VC
files have migrated into this repository, not the full row set. The
walls are
measured under the mild concurrency stated in the caption, on the
machine named in the receipts, and will drift with hardware like any
wall-clock. Finally, the certified build itself and the tool that
seals row listings into its signed-commit input form have not yet
migrated into the public trees: today the table re-derives
mechanically from the committed receipts, and re-running the binary
replay itself becomes a repository make target when those artefacts
land (per the migration rule stated in
section~\ref{sec:exp-repro}).
 \section{Gymnasium: first certified rows for a standard RL suite}
\label{sec:exp-gym}

The second campaign marched the same machinery across Farama's
Gymnasium --- the de-facto standard reinforcement-learning benchmark
suite --- pinned at repository revision \texttt{4347ce947}, with every
certificate naming its environment id and keyword arguments. To our
knowledge, no prior certified bound at any of this chapter's claim
classes has been published for any Gymnasium environment; the
comparison points that exist are the suite's own registered
``solved'' thresholds and code-comment optima, plus empirical
leaderboard numbers, and the campaign's rows are therefore firsts at
their class.

One property of these certificates deserves stating before the table,
because every row reads differently in its light: they are
\emph{constructive}. Each certificate ships the policy itself --- a
lookup table, indexed by state and step, sitting in the certificate
bytes beside the value table and the claim --- so no row asserts
merely that a policy achieving the bound \emph{exists}; it asserts
that \emph{this} one does, and whoever holds the certificate holds the
policy. The checker enforces the distinction structurally: its
per-step obligations bound each state's value by the kernel-weighted
value of the shipped policy's own action --- never by a maximum over
actions --- so the policy is load-bearing in every obligation it
discharges, and a tampered policy entry is among the refused mutants
of the battery. Table~\ref{tab:exp-gym} is the master table for the
tabular tier.

\begin{table}[htbp]
\centering\small
\begin{tabular}{l l r l r}
\toprule
environment & claim & certified bound & vs.\ optimum & check \\
\midrule
FrozenLake 4x4 (slippery) & U3 & $\ge 0.7441902$ & \cellcolor{exacttint}optimal, exactly & $15\,\mu$s$^{\,\text{gb10}}$ \\
FrozenLake 8x8 (slippery) & U3 & $\ge 0.9132201$ & \cellcolor{exacttint}optimal, exactly & $100\,\mu$s$^{\,\text{gb10}}$ \\
FrozenLake 4x4 (succ.\ $3/4$) & U3 & $\ge 0.8731132$ & \cellcolor{exacttint}optimal, exactly & $15\,\mu$s$^{\,\text{gb10}}$ \\
FrozenLake 4x4 (not slippery) & U3 & $= 1$ & \cellcolor{exacttint}optimal, exactly & $13\,\mu$s$^{\,\text{gb10}}$ \\
CliffWalking & U2 & $\ge -13$ & \cellcolor{exacttint}optimal, exactly & $24\,\mu$s$^{\,\text{gb10}}$ \\
CliffWalking (slippery) & U3 & $\ge -63.0133733$ & \cellcolor{exacttint}optimal, exactly & $34\,\mu$s$^{\,\text{gb10}}$ \\
Taxi & U2 & $\ge 7.93$ & \cellcolor{exacttint}optimal, exactly & $0.50$\,ms$^{\,\text{gb10}}$ \\
Taxi (rainy) & U3 & $\ge 3.9545745165$ & within $3{\times}10^{-14}$ & $0.70$\,ms$^{\,\text{gb10}}$ \\
Taxi (fickle passenger) & U3 & $\ge 8.38$ & \cellcolor{exacttint}optimal, exactly & $1.1$\,ms$^{\,\text{gb10}}$ \\
Taxi (rainy + fickle) & U3 & $\ge 4.5756463469$ & within $3{\times}10^{-14}$ & $1.5$\,ms$^{\,\text{gb10}}$ \\
Blackjack (S\&B rules) & U3 & $\ge -0.0431131055$ & \cellcolor{exacttint}optimal, exactly & $0.16$\,ms$^{\,\text{gb10}}$ \\
\bottomrule
\end{tabular}
 \caption{Certified rows for the Gymnasium tabular tier. The claim
column names the species of section~\ref{sec:exp-claims}: U2 is a
certified cost-to-goal, U3 an exact probabilistic bound. Policies are
shipped lookup tables; checks are single-core medians. ``Optimal,
exactly'' means the bound is the exact optimum of the
(deadline-truncated) environment: terminal at its claim class ---
those cells carry the exact-cell tint of
section~\ref{sec:exp-conventions}.}
\label{tab:exp-gym}
\end{table}

Four findings come out of this table.

\paragraph{The suite's own Taxi threshold is unsatisfiable.} The
registry declares Taxi ``solved'' at average return $8$, with a code
comment placing the optimum at $8.46$. The true optimum of the pinned
environment is exactly $793/100 = 7.93$ --- \emph{below} the
threshold; no policy can meet it. The $8.46$ is a fossil: on the
pre-correction map (before an upstream fix added one wall segment) the
same exact computation yields $2539/300 = 8.4633\ldots$, which we
reproduced by swapping only the map. The threshold was never updated
for the wall.

\paragraph{The first exact value for the textbook blackjack.} Sutton
and Barto's blackjack variant (infinite deck, dealer stands on all
17s, hit/stand only) may be the most-reimplemented exercise in
reinforcement learning; a literature search --- the textbook, Baldwin
et al.\ 1956, the gambling-mathematics literature, recent preprints
--- found no published exact expected value under optimal play. It is
$-3954908734551683826282634489 / 13^{26} \approx -4.311\%$, now with
a machine-checkable certificate. (Consistency: the same game with 3:2
naturals is known to run near $-2.42\%$, and the difference is the
order of the forgone bonus. A subtlety: the suite's observation space
cannot distinguish a two-card 21 from a three-card 21, but the payout
rule can --- the certificate runs on a slightly enriched state space
while proving the shipped policy never uses the extra information.)

\paragraph{Hidden state, certified honestly.} Taxi's ``fickle
passenger'' variant re-draws the destination mid-trip with probability
$3/10$, invisibly until it fires. A certificate on the full state
would be clairvoyant; these rows certify the \emph{belief} chain
instead, so the shipped policy is implementable from the observation
history. The variant also \emph{raises} the optimum (dry
$7.93 \to 8.38$; rainy $3.9546 \to 4.5756$): a uniform mid-trip
re-draw is on average favourable, and the certified policies exploit
the re-route.

\paragraph{The registry's comment optima were truncations of numbers
nobody had.} The FrozenLake comments say ``optimal average reward
0.74 / 0.91.'' The exact optima --- rationals with denominators
$3^{99}$ and $3^{200}$ --- were computed here, we believe for the
first time; the comments are truncations of them.

\subsection{How the rows were produced: exact backward induction
and its wall-clock cost}
\label{sec:exp-gym-producer}

The producer strategy for the whole tabular tier is the classical
algorithm, executed exactly: backward induction --- finite-horizon
dynamic programming --- over the deadline-truncated chain. Starting
at the deadline and stepping backward, the producer computes, for
every state at every step, the optimal value by a single Bellman
backup --- a maximum over actions of the expected one-step reward
plus the kernel-weighted value one step ahead --- in exact rational
arithmetic, recording a maximizing action as it goes. The
recorded actions \emph{are} the shipped policy; the value tables
become the certificate's obligation rows; the value at the initial
state (or the exact mixture over the initial distribution) becomes
the claim. Termination and exactness are structural: the horizon is
finite --- the deadline is the environment's own truncation, or the
certificate's declared one --- the state and action sets are finite,
and no backup ever rounds, so after $T$ sweeps the result is the
exact optimum of the truncated environment, not an approximation to
it. There is no learning loop, no convergence tolerance, and no
rounding anywhere in the construction. (Making exactness fast is an
implementation matter, not an algorithmic one: values at each step
are stored as integers over a shared power of the kernel
denominators' least common multiple, so the hot loop is
integer multiply-accumulate, division-free, with canonical reduction
deferred to emission.) Nor is tightness taken from the producer on
trust: for the exact lanes the checker has an optimality mode whose
obligations quantify over \emph{all} actions --- a supersolution
check beside the shipped policy's attainment --- so ``optimal,
exactly'' in Table~\ref{tab:exp-gym} is itself a machine-checkable
statement, and the mode's test fixtures include a matched pair
exhibiting the gap: a certificate that passes as a bound but refuses
as an optimum.

The direct answer to ``how long does producing the optimal policy and
its certificate take'' is: for the entire tabular catalogue --- all
$53$ certificates, the master table's rows and the frontier studies'
variants below included --- the optimized producer runs from process
start to every certificate synthesized, emitted, and byte-compared in
$0.69$\,s of wall-clock time on four CPU cores. The per-lane receipts
break that down: the heaviest single computation, the exact optimum
of the rainy fickle-passenger Taxi belief chain, takes $117$\,ms on
one core ($35$\,ms on four); the 8x8 FrozenLake exact table,
$6.6$\,ms; the textbook-blackjack optimum, $2.2$\,ms. The reference
implementation that first produced the corpus takes roughly
$3.2$\,s single-core for the same spans; the optimized producer
takes $0.5$\,s on one core ($6.5\times$), and the two emit
byte-identical certificates on all $53$ --- the strongest available
cross-implementation check. A GPU lane reproduces
the dyadic tables bit-identically with a $\sim 2\times$ gain on the
largest rows only; on small rows the kernel-launch floor loses to the
CPU, a null result recorded as such. Throughout, producer cost and
check cost are measured separately: the check columns of every table
in this section time only the checker.

These certificates also ride the sealed-container lanes of
section~\ref{sec:exp-repro}, and their re-emission on the unified
schema returned a finding rather than a formality. The schema's
chain species is a pure discrete-time Markov chain --- transition
and initial rows, no action layer --- so a hitting-lane certificate
carries over verbatim, table and claimed bound alike, exactly when
the decision problem admits a \emph{stationary} composition: one
action per state admissible at every step of the deadline. The
hitting corpus splits $9/29$ on that criterion, and the split is
semantic content, not an artifact: every deterministic-dynamics
certificate composes, while slippery FrozenLake at long deadlines
has states where no single action is admissible at every step ---
near the deadline the optimal policy gambles where far from it it
detours. The nine composable certificates already ride the unified
schema with their tables and claims byte-verbatim; the twenty
time-varying ones stand in the pre-unification record, each
exclusion a manifest fact naming its obstructing states, and their
re-emission awaits the schema's decision species
(\texttt{make gymnasium-glc} re-checks both lanes and the exclusion
records alike).

\subsection{The measured frontier: bound versus check cost}
\label{sec:exp-gym-frontier}

Because the tabular certificates are terminal in bound, the
interesting curve is in the \emph{representation}. Two honest knobs
trace it: the deadline $T'$ (certifying ``goal within $T'$ steps'' for
$T' \le T$ weakens the bound but shrinks the table) and the precision
$k$ of a dyadic re-emission (round-down integer tables with bound loss
$\le T \cdot 2^{-k}$, turning every obligation into an integer
comparison). On the FrozenLake rows, moving from exact rational
arithmetic to dyadic integer tables at equal obligation count speeds
the check by roughly $400\times$--$1400\times$ at a bound cost of
order $10^{-13}$ --- far below any displayed digit. The general
lesson, confirmed three independent times in this campaign's receipts:
\emph{certificates should be engineered for their checkers} ---
rationals belong in the claim line, shared-denominator integer slabs
in the obligations, and the checker's hot loop then vectorizes (the
kernel-level lanes add a further $1.6\times$--$3.3\times$ on the
larger reward tables, at $5$--$7$\,ns per obligation).

Honest non-dominations in the same receipts: on the smallest
certificates the vectorized lane's setup cost makes it
$0.83\times$--$0.89\times$ the plain lane (crossover near the 8x8
scale); and one certificate family (blackjack at high dyadic
precision) structurally exceeds the vectorized lane's overflow
envelope, is declined, and falls back to the reference lane ---
deferral costs completeness, never soundness, and the fix is a
smaller~$k$.

\subsection{Beyond tables: continuous state and solver-semantics
environments}
\label{sec:exp-gym-beyond}

Above the tabular tier the campaign is deliberately breadth-first, and
at this preview two further tiers carry rows.

\textbf{Continuous state: certified worst-case return.} The first
continuous-state row certifies Pendulum-v1: a Bellman-supersolution
certificate over a polynomial cell cover of the state space (the
environment's velocity clipping entering as explicit piecewise
structure), yielding a certified \emph{worst-case} return bound: one
holding for every initial state, with no probability anywhere in the
certificate --- the strongest of the claim species of
section~\ref{sec:exp-claims} (U1) --- and carrying an anytime
cover-refinement curve measured from a few dozen cells to nearly a
thousand. At this preview the certified bound stands at return
$\ge -1.192$ over the full initial set, with the cover-refinement gap
driven to $2.4\times 10^{-4}$, and a five-mutant battery refused. The corresponding
CartPole certificate (worst-case return exactly $500$, the same
species) has crossed the composition step it was staged behind ---
and the step dissolved rather than being climbed: instead of
transporting a continuous-time decrease certificate through the
environment's explicit-Euler step map, the certified row states a
per-cell contraction on the step map itself, which
denominator-clears to a polynomial of its own, leaving only
kernel-checked enclosures of the current state's sine and cosine.
The first certified CartPole row now stands at the subsystem where
it is honest: under an exhibited cell-table bang-bang policy, the
pole angle provably never reaches the termination threshold, from
every initial state in the reset box, in exact rational arithmetic,
with a designed-locus mutant battery refused. The return bound
itself still waits on the cart-position excursion, which the
pole-only policy does not regulate; that full-state variant is the
campaign's first genuinely large checking workload ($10^6$--$10^7$
polynomial cells), sized to exercise the batched checking engine of
the Geolog tier as it lands.

\textbf{Solver-semantics environments: i.i.d.\ confidence
bounds.} Where the environment's dynamics is a compiled physics
engine (the Box2D and contact-rich MuJoCo families), no analytic
species reaches the binary's semantics, and pretending otherwise
would be a category error. What applies instead is the
\emph{i.i.d.\ tier}: an anytime-valid lower confidence bound on
the expected return, certified from recorded rollouts of a pinned
policy --- the weakest species of section~\ref{sec:exp-claims} (U5),
holding with probability at least $1 - \alpha$ rather than almost
surely, but model-free, which is exactly what lets it reach a
compiled binary at all. Its rows are always displayed as their own
class, never mixed into a table of exact rows: Table~\ref{tab:exp-ll}
shows the first two --- to our knowledge the first certified
statements of any kind at this claim class for Box2D environments.
Each is a betting-martingale ledger over $10^6$
rollouts of the suite's own heuristic controller; the checker
re-verifies the entire wealth process in directed-rounding fixed
point ($1.6 \times 10^7$ obligations in $\sim 130$--$200$\,ms,
single-core), and the certificate is \emph{anytime-valid}: extending
the ledger with more episodes can only tighten the bound, with no
statistical penalty.

\begin{table}[htbp]
\centering\small
\begin{tabular}{l l r r}
\toprule
environment & claim (class U5) & bound & threshold \\
\midrule
LunarLander-v3 & $\mathbb{E}[R] \ge$, $\alpha = 1/100$ & \cellcolor{winA}$236.2$ & $200$ \\
LunarLanderContinuous-v3 & $\mathbb{E}[R] \ge$, $\alpha = 1/100$ & \cellcolor{winA}$281.1$ & $200$ \\
\bottomrule
\end{tabular}
 \caption{First i.i.d.-tier rows (species U5 of
section~\ref{sec:exp-claims}; confidence-qualified: each bound
holds with probability $\ge 1-\alpha$ over the sampling; never to be
read as almost-sure). Policies are the suite's own heuristic
controllers; $n = 10^6$ episodes each; the almost-sure return range
required by the bound is itself a certified lemma about the
environment's code, not an assumption. Bound-cell tints are the
win-tint ladder of section~\ref{sec:exp-conventions}, read against
the registry threshold beside them.}
\label{tab:exp-ll}
\end{table}

One reading of Table~\ref{tab:exp-ll} previews the genre's point: the
\emph{certified floor} of the continuous-throttle variant ($281$)
exceeds the discrete variant's \emph{empirical mean} ($238$) --- a
policy-comparison statement made entirely below the certified line,
where it cannot be an artefact of sampling luck.

\paragraph{What is not claimed.} The tabular certificates' trust seam
--- that the modelled kernel is the environment's --- is narrowed here
by asserting the modelled tables equal to the suite's own published
transition tables (exact equality, with the handful of float-literal
discrepancies carried as declared ledger rows), but it remains
producer-side. The i.i.d.-tier rows trust the pinned
pseudorandom generator and the policy's determinism, both named in
the certificate. No worst-case, cost-to-goal, or exact-probabilistic
claim (species U1--U3) is made for any solver-semantics environment. No MuJoCo
row appears in this preview's tables; the named route --- its design
landed as this preview goes to press, rows in flight --- is a
hybrid-time composition pattern: discrete controller ticks with
bounded clock jitter coupled to continuous-time physics, the jitter
entering as a certificate \emph{parameter} rather than a disturbance,
with the suite's contact-free environments (InvertedPendulum and its
ladder) as first targets and the contact-rich families expected to
arrive i.i.d.-tier first. Producer-side synthesis time for the continuous
tier is reported in the repository but not compared against
reinforcement-learning training costs: the claims differ in kind, and
a wall-clock race between them would compare nothing meaningful.
 \section{Hardware model checking: certificates where verdicts were}
\label{sec:exp-circuits}

The third campaign targets a benchmark from a different community: the
hardware model-checking suite introduced by the neural model-checking
line of work (Giacobbe, Kroening, Pal, and Tautschnig, \emph{Neural
Model Checking}, NeurIPS 2024, arXiv 2410.23790) and extended by the
same group at NeurIPS 2025 (\emph{Let a Neural Network Be Your
Invariant}) --- eleven Verilog designs (UART and I\textsuperscript{2}C
controllers, VGA and LCD drivers, PWM, counter, and
thermocouple-interface blocks), each with a ladder of parameter
instances, under three
property classes: $219$ pure-liveness tasks
($\mathsf{FG}\,\alpha \to \mathsf{GF}\,\beta$), $206$ safety tasks, and
$209$ combined safety-plus-liveness tasks (premise-guarded
until-properties such as $\mathsf{G}(r \to p\,\mathsf{U}\,q)$) ---
$634$ verification tasks in all. The published pipelines train a
neural certificate per task --- a ranking function by gradient descent
in 2024; ranking and invariant functions by MILP in 2025 --- quantize
it, and hand a single validity query to an SMT solver; the reported
verdict per task is the solver's ``unsat.'' The baselines are mature
model checkers (nuXmv, ABC, rIC3), and the experimental budget is five
hours per task.

The toolkit's route is structurally different, and the difference is
the point. Every one of the $634$ specifications reduces, by a chain
of temporal rewrite steps, to conjunctions of two certified forms:
pure recurrence $\mathsf{GF}\,\gamma$, discharged by a ranking
function --- piecewise-affine with exact rational coefficients over a
cell partition of the circuit's own branch structure, input wires
resolved demonically (universally) at every step --- and $k$-step
safety windows, discharged by obligations on the exact one-step image
(plus small inductive invariants on the four families that need
them). No
B\"uchi automaton and no product construction appears anywhere on the
suite. The trust story is layered the same way at every level: the
consuming lemmas (a ranking that drops by $\varepsilon_0$ off the
target and stays nonnegative forces target steps to recur; a window
whose obligations hold on the reachable image holds always) are a few
lines of mathematics each, machine-checked in Lean with no measure
theory; every rewrite rule's soundness is machine-checked \emph{twice},
in Lean and independently in HOL4; each task carries a rewrite-chain
witness that a standalone checker replays from the verbatim
specification text; and every ranking and window condition is an
ordinary polynomial-nonnegativity row of the
\bookref{ch:qtsl-manual}{QTSL manual} calculus, recomputed exactly in
$\mathbb{Q}$ by standalone checkers that import nothing from the
producers.

At this preview the sweep covers the entire suite: \textbf{all $634$
tasks, across all three classes} --- $219/219$ pure liveness ($245$
certificates: the two conjunction families carry two per task),
$206/206$ safety ($222$ certificates), and $209/209$
safety-plus-liveness ($418$: one window and one recurrence certificate
per task) --- $885$ certificates and $634$ chain witnesses in all.
Table~\ref{tab:exp-nmc} sets this against the published field on the
field's own numbers: the strongest single published tool, the 2025
MILP method, completes $554$ of $634$; the per-task union of
\emph{every} tool in the 2025 paper's comparison completes $617$.
Seventeen tasks are completed by nothing published at all ---
twelve pure-liveness tasks (eleven blink instances and the largest
\texttt{load\_store} instance) and the five largest blink safety-plus-liveness
tasks. The toolkit certifies all seventeen; blink is a counter ladder,
and its certificates have the same shape at $N = 2^{32}$ as at
$N = 2^{8}$.

\begin{table}[htbp]
\centering\small
\begin{tabular}{l r r r r l}
\toprule
method & liveness 219 & safety 206 & s{+}l 209 & all 634 & per-task evidence \\
\midrule
this toolkit & \cellcolor{exacttint}$219$ & $206$ & \cellcolor{exacttint}$209$ & \cellcolor{exacttint}$634$ & checkable certificate \\
neural, NeurIPS'25 (MILP) & $192$ & $190$ & $172$ & $554$ & solver verdict \\
nuXmv (BDD) & $180$ & $206$ & $154$ & $540$ & tool verdict \\
rIC3 1.5.1 & $110$ & $196$ & $129$ & $435$ & tool verdict \\
ABC & $78$ & $206$ & $96$ & $380$ & tool verdict \\
neural, NeurIPS'24 (re-run) & $172$ & n/a & n/a & $172$ & solver verdict \\
\midrule
virtual best of all baselines & $207$ & $206$ & $204$ & $617$ & portfolio of verdicts \\
\bottomrule
\end{tabular}
 \caption{The full NeurIPS'25 neural-model-checking suite: completions
by property class (pure liveness $219$, safety $206$,
safety-plus-liveness $209$). Baseline counts are recomputed from the
2025 paper's own artefact --- its plotting script's hard-coded runtime
arrays are the paper's published data --- at the paper's own five-hour
timeout. The 2024 method is a liveness-only technique; the artefact
records no runs for it outside that class (n/a). The virtual-best row
is the per-task union of every baseline in the artefact's comparison,
including the linear variant of the 2025 method, which is not shown as
a separate row. The toolkit's per-task evidence is an exact-arithmetic
certificate replayed by a standalone checker, with consuming lemmas
machine-checked in Lean and every spec-rewrite rule machine-checked in
both Lean and HOL4. A toolkit count carries the exact-cell tint of
section~\ref{sec:exp-conventions} where it strictly exceeds the
virtual best of the whole published portfolio (ties untinted).}
\label{tab:exp-nmc}
\end{table}

Three anchoring notes travel with Table~\ref{tab:exp-nmc}. First,
\emph{the comparison is like-for-like with the 2025 paper} and only
approximately comparable to the 2024 publication: the benchmark's
circuits were revised between the two editions ($142$ source files
changed semantically, an edit the 2025 paper does not mention), and the
toolkit certifies the revised circuits --- the same files the 2025
re-runs measured. The 2024 method's own line shows the drift: $93\%$
of its original $194$ tasks as published, on the 2024-era circuits;
$172/219$ under its authors' re-run on the revised, extended class. Second, \emph{the semantics certified is the benchmark's
exactly}: the suite's generated SMV models carry no initial-state
constraint, so every task is a from-every-state claim under
demonically-resolved inputs --- precisely the universal form the
certificates prove. For the until-shaped class this was additionally
confirmed off-calculus: an independent explicit-state model checker (a
demonic greatest-fixpoint computation written against the
specification semantics, not the certificate calculus) verifies each
safety-plus-liveness property as stated on the smallest instance of
every family, and falsifies none. Third, \emph{the completion race is
not the durable claim} --- model checkers are healthy, and successor
methods will contest counts. The durable column is the last one: on
this suite, every completed task ships an auditable certificate, a
small exact object whose consuming lemma is machine-checked, where
every baseline --- neural or classical --- ships a verdict.

The checking economics stay flat as the suite triples. Certificates
run $0.4$--$7.4$\,KB ($98\%$ under $4$\,KB; the tail is the VGA safety
windows); individual checks run from tens of microseconds to
a few milliseconds; replaying every certificate \emph{and} every
rewrite-chain witness in all three classes takes under a second of
single-core wall-clock in total (per the class receipts: certificate
checks of $78$, $195$, and $171$\,ms for liveness, safety, and
safety-plus-liveness, plus $82$, $118$, and $150$\,ms of chain
replay). Certificate size and check time are constant along each
design's parameter ladder, with only coefficient widths growing as
$\log N$, where the learning pipelines pay training and solving per
instance. Per the conventions of section~\ref{sec:exp-conventions},
the two sides' times are a contrast of kind, not a ratio --- searching
for a proof and re-verifying a shipped one are different jobs on
different machines: the paper reports a five-hour-per-task budget and
roughly nine thousand machine-hours for its full experimental
campaign; the toolkit's whole-suite certificate replay is the
under-a-second figure above.

Refusal is exercised as systematically as acceptance. Each class
ships a mutant battery --- deliberately corrupted certificates and
chain witnesses, each designed to be refused at one specific check:
seventeen layer-targeted attacks with four positive controls for the
liveness class (ported into this repository's re-check driver),
thirteen attacks with three controls for safety, and fourteen attacks
with nine controls for the combined class, two of those attacks
cross-layer (accepted by the certificate checker, refused by the
chain checker) --- and each is in fact refused at the attacked layer,
with every control passing. The remaining producer-side trust seam is
the circuit manifests: that the certified step system \emph{is} the
Verilog's semantics. It is narrowed by two-simulator cross-checks,
exhaustive on the full state--input box of each family's smallest
instance and sampled on the largest, and --- since the safety campaign
--- by an independent re-transcription of every design. That newest
instrument earned its keep immediately: the two original simulators
share an author, the first edition of this section said so as its
declared residue, and the re-transcription caught exactly the failure
that caveat named --- a shared misreading of one LCD counter's wrap
point, latent in the already-swept liveness manifests. The manifests
were repaired and the full battery re-run green within the hour; the
fixtures in this repository are the repaired ones. Replacing this seam
with a netlist-level bridge remains named successor work.

Exact re-derivation of a published field keeps yielding findings of
the same genre as the Gymnasium campaign's unsatisfiable Taxi
threshold (section~\ref{sec:exp-gym}). Three from this campaign are
worth recording. \emph{The sound premise is not the obvious one}: the
natural liveness premise for the until-introduction rewrite rule ---
$\mathsf{GF}\,q$, ``the window eventually closes'' --- is false
whole-box on the I\textsuperscript{2}C family, whose clock-stretch
freeze states never reach the window and satisfy $\lnot p$ forever
instead; the sound premise is $\mathsf{GF}(\lnot p \lor q)$. The
correction was forced during rule design and then observed
semantically: the explicit-state falsifier's demon-winning regions are
nonempty exactly on those frozen states. \emph{One task's two
specification sources disagree}: the delay design's first combined
property reads $\mathsf{G}\,\lnot\mathit{rst} \to \cdots$ in the
runner's specification text but $\mathsf{FG}\,\lnot\mathit{rst} \to
\cdots$ in the same runner's nuXmv property line. The toolkit
certifies the specification text as stated --- the chain witness pins
its verbatim bytes --- and records the discrepancy; certifying
as-stated rather than as-presumably-meant is what makes such
discrepancies visible at all. \emph{The artefact's own aggregation
has a slip}: the 2025 plotting script assembles its safety-class
aggregate with one VGA family's arrays included twice and the other's
omitted. The counts in Table~\ref{tab:exp-nmc} are unaffected ---
they are recomputed from the per-family arrays, and both VGA safety
families are fully solved by the classical baselines --- but the
paper's aggregate safety timing figures ride the wrong multiset. (The
2024 edition had a kindred fossil: an appendix table summing to $177$
completions where the summary claims $180$; the census here restored
the three footnoted rows a naive parse drops, scoring that comparison
in the paper's favor.)

Because the producers here are benchmark-specific and written by us,
the fair cost question is not only ``how fast is the check?'' but
``what did it cost to \emph{get} the certificates, method invention
included?'' --- and the campaign's working sessions are timestamped,
so it has a measured answer. The pure-liveness campaign end-to-end ---
reading the papers, designing the route, machine-checking the
consuming lemma, writing the certificate format, the checker, and
every producer, and sweeping all its tasks --- took about $75$ minutes
of commit-stamped wall-clock on one desk-side machine (this
repository's invention-time comparison table carries the per-phase
receipts and its estimating conventions). The safety class and the
combined class were each added in a single further working session the
following day --- the combined class including six new
until-introduction rewrite rules, each proved twice, in Lean and in
HOL4. These figures sit beside the published five-hours-per-task
budget as, again, a contrast of kind: method invention on one side of
the ledger is being compared to certificate search on the other, and
we print no ratio.

\paragraph{What is not claimed.} Synthesis wall-clock supremacy is not
claimed: on tasks the 2025 MILP method completes, it is itself fast
(sub-second on most families); the measured advantages here are
completion, the check side, and the evidence class. Generality over
automaton-requiring specifications is not claimed: no property in this
suite needs genuine automaton memory --- the 2025 additions, though
until-shaped, reduce by premise-checked implication rules --- and on a
suite where that machinery is needed, it would return. Sim-tightness
is claimed only where measured: every liveness family's worst-case
descent bound is achieved exactly by a simulated run of the circuit.
And the reproduction spine in this repository is complete at this
preview for the pure-liveness class only --- committed fixtures, the
pinned independent checker, and the mutant battery, re-run end to end
by one make target; the safety and combined classes' receipt tables
are vendored here, and their fixtures and checkers arrive with the
producer layer's migration. Per the preview note of
section~\ref{sec:exp-repro}, until a cited artefact lands in the
public repository, its mention is a statement of what exists, not yet
a link.
 \section{Atari as a circuit: making the environment white-box}
\label{sec:exp-atari}

The fourth campaign closes a loop between the previous two. It takes
an environment family from the reinforcement-learning suite's hardest
tier --- Atari 2600 games, which Gymnasium reaches through an
emulator plugin --- and turns each game into an object of the
hardware campaign's kind. Rather than treating the emulator as a
black box to be sampled, the campaign compiles the console itself
into a synchronous digital circuit: an open-source Verilog
implementation of the 2600 (vendor-pinned, with three local repairs,
each carried in a ledger) is synthesized with the game ROM baked in,
down to an and-inverter graph. Pong closes to a circuit with $52$
input bits, $1{,}491$ flip-flops, and $20{,}821$ AND gates; Freeway
to $52$/$1{,}491$/$20{,}798$; Montezuma's Revenge --- whose
$8$\,KB bank-switched cartridge brings its mapper hardware into the
netlist alongside the console --- to $52$/$1{,}501$/$36{,}698$. The
machines' full state spaces are $2^{1491}$ and $2^{1501}$.

This is where the scope condition of section~\ref{sec:exp-claims} has
its teeth, so it bears restating in this campaign's terms: the
toolkit can actually \emph{examine} the environment's digital-circuit
structure. The transition relation is a given, finite, white-box
object, and every claim species of this chapter becomes definable
over it --- the ranking rows of section~\ref{sec:exp-circuits} apply
to the environment itself, not to an abstraction of it. System
identification --- recovering such a circuit from observations of
play --- is a separate problem, and one the toolkit does not yet have
tools for. Nothing in this campaign learns the environment; it is
handed the mechanism and reasons about it.

What is landed at this preview is the foundation tier: the compiled
circuits and their fidelity receipts. Differential validation runs
the circuit and the suite's own emulator (ale-py) frame-locked, side
by side: on Pong and Freeway, episode returns are \emph{identical} --- Pong's
rally period is exactly $140$ frames on both; Freeway's first
crossing lands on the same frame --- and $116$--$123$ of the $128$
RAM bytes match exactly on every frame, with each remaining
divergence assigned to a named class (an input-transduction mismatch,
since resolved by measurement; a one-to-two-frame boot-time skew
attributed to the pinned CPU core's known non-cycle-exactness)
rather than left as noise. One finding from this validation is
recorded as a surprise: the emulator's RAM trace is invariant across
its random seeds --- for these games, the suite's stochasticity lives
in frame-skip and action timing, not in the machine.

The first certified rows above that foundation landed as this
preview closed: exact-replay returns --- worst-case over the suite's
full initial-condition randomization, all $31$ no-op offsets (species
U1 of section~\ref{sec:exp-claims}) --- for a fixed open-loop
policy, certified by a checker whose only dynamics
authority is the compiled circuit itself --- ``this policy scores
$R$'' turned from a measurement into a theorem about the circuit.
Table~\ref{tab:exp-atari} shows the rows. The Freeway row carries a
finding: the certified per-offset returns span $\{2,3,4\}$, so a
single rollout of this policy reports a number drawn from the init
lottery, and the certified minimum ($2$) sits strictly below the
modal return --- the gap between ``return of a rollout'' and ``return
guaranteed over the init set,'' exhibited concretely.

The last two rows take the same claim species to Montezuma's Revenge
--- the suite's canonical hard-exploration benchmark, the game whose
sparse rewards famously held early deep reinforcement-learning agents
to a score of zero. To our knowledge these are the first certified
return claims for this game at any claim class; every published
Montezuma number is an empirical rollout statistic. The certified
content is exact: reaching the first key (return $100$) is certified
at \emph{every} one of the $31$ no-op init offsets, and a
$1{,}548$-frame trajectory that collects the key and opens a door of
the first room is certified at return exactly $400$ from the declared
initial state --- each reward event pinned to its frame, all
$128.5$ million netlist clock ticks of the replay re-derived by the
checker, every frame of the pinned witness traces compared bit-exact.
The rows' structure carries a finding of the Freeway kind, amplified:
a committed instrument row certifies that the fixed key script which
scores $100$ at offset zero scores exactly $0$ at every other offset
--- total collapse under the init lottery, where Freeway's spread was
one point --- so the per-offset row holds thirty-one separately
searched scripts, one per offset. Transfer fails the same way at
finer grain: a script optimized against the reference emulator, where
it scores $400$, scores only $100$ when replayed through the circuit,
losing a life at a pixel-perfect skull crossing that succeeds under
the emulator's timing and fails under the circuit's. Trajectories
that certify on the circuit were searched on the circuit; the
searcher is producer-side and untrusted, and nothing about how the
scripts were found enters the trust story --- the certificate is the
exact replay.

In flight at this preview: liveness rows over the compiled circuit at
the previous section's claim classes, and policy rows --- closed-loop
pixels-to-action policies rather than fixed scripts --- for the games
themselves; toward the latter, the checker has already grown a pixel
lane, which on the key trajectory compared every one of the
$26{,}997{,}440$ pixels the game renders across its $706$ frames
bit-exact against the circuit's video output.
Each will carry the same receipts discipline as every table above.

\begin{table}[htbp]
\centering\footnotesize
\setlength{\tabcolsep}{5pt}\begin{tabular}{l l r l r r}
\toprule
environment & policy & $T$ & init set & min return & check \\
\midrule
Pong & NOOP script & $1000$ & $k \in \{0..30\}$ & \cellcolor{exacttint}$-6$ & $1405$\,s \\
Freeway & all-UP script & $1000$ & $k \in \{0..30\}$ & \cellcolor{exacttint}$2$ & $1410$\,s \\
MontezumasRevenge & per-offset key scripts & $828$ & $k \in \{0..30\}$ & \cellcolor{exacttint}$100$ & $3943$\,s \\
MontezumasRevenge & key+door script & $1548$ & $k = 0$ & \cellcolor{exacttint}$400$ & $6416$\,s \\
\bottomrule
\end{tabular}
 \caption{Certified exact-replay rows (worst-case species U1 of
section~\ref{sec:exp-claims}); each environment is the compiled
circuit \texttt{AtariCircuit-}$\langle$game$\rangle$. ``Min return''
is the certified minimum over the row's declared init set
(exact, hence the exact-cell tint of
section~\ref{sec:exp-conventions}) --- all
$31$ no-op offsets, except that the key+door row declares the single
offset $k=0$; the per-offset Montezuma row certifies one separately
searched script per offset. Check times are single-core replay of
the claim's own circuit.}
\label{tab:exp-atari}
\end{table}

\paragraph{What is not claimed.} These exact-replay rows are the
certified ladder's weakest rung --- a fixed open-loop policy, exact replay, no
quantification over adversaries and no learned certificate --- and
their receipts say so; no liveness or policy-optimality claim is
stated for any Atari game at this preview. The Montezuma key+door row
in particular quantifies over nothing: one declared initial state,
one script, $1{,}548$ frames --- it certifies that score on that
trajectory exactly and is silent about the rest of the game. Nothing
here ``solves'' Montezuma's Revenge; what is new is that a Montezuma
score is carried by a machine-checked certificate against the
machine's own transition relation instead of by a rollout report.
The five-to-twelve
divergent RAM bytes per frame in the emulator differential are
declared, not resolved; the CPU core's cycle-exactness is a
named seam of exactly the circuit-manifest kind
(section~\ref{sec:exp-circuits}); and claims about the emulator
itself rest on the differential receipts alone. The white-box
condition above is load-bearing: none of this machinery applies to an
environment whose mechanism cannot be handed over.
 \section{Process control: the semiconductor ladder}
\label{sec:exp-semicon}

The final campaign leaves benchmarks-with-leaderboards entirely and
runs on two tracks over a ladder of semiconductor-fab process models
--- the application-domain family the toolkit's programme context
cares about, and a domain where the certified claims are economically
legible. The \emph{fluid} track instantiates the chapter's
infinite-horizon species (U4 of section~\ref{sec:exp-claims}:
long-run throughput and stability claims for queueing-network
models): \emph{this dispatch policy keeps the line stable at this
load}. That claim has teeth because it is famously falsifiable:
re-entrant lines destabilize under plausible dispatch policies
strictly below nominal capacity (the Kumar--Seidman phenomenon), and
the fluid-limit criteria of Dai and Dai--Weiss reduce stability to
exactly the polyhedral drift inequalities the toolkit's checkers
replay. The \emph{discrete} track --- new at this preview, on the
base rung --- instantiates the worst-case ranking species (U2,
quantified here over every scheduler rather than every initial
state): \emph{every dispatch of this line brings all lots to the last
queue by this deadline} --- and carries it past any of the chapter's other
campaigns, to a claim class none of them produce: exact optimality,
a certified floor over \emph{all} dispatch policies met exactly by a
shipped one, for up to seven lots. These claims are about a
simplified discrete model of the line, which gives each process step
at most one unit of service per round, even at the two-machine
stations (model note below).

\paragraph{The ladder, whole.} The campaign's target is a fixed
progression of eleven public fab models in four stages --- the Intel
Mini-Fab; the seven SEMATECH MIMAC testbed fabs
\bookcite{fowler95}{(Fowler and Robinson 1995)}; the two SMT2020
scenario fabs \bookcite{kopp20}{(Kopp, Hassoun, Kalir, and M\"onch
2020)}; and SMAT2022 \bookcite{lee22}{(Lee, Song, Chang, and Park
2022)} --- carrying the same claim species at every rung, so the
ladder is a progression in scale and modelled physics, never a change
of subject:

\begin{center}\small
\begin{tabular}{@{}llc@{}}
\toprule
rung & one-line identity & certified rows\\
\midrule
Mini-Fab & 3 stations, 5 machines, 6 re-entrant steps (Kempf 1994) &
  this preview, on simplified models\\
MIMAC set 1 & 2 products, 265 tools, routes to 245 steps & ---\\
MIMAC set 2 & 7 products, 274 tools, routes to 354 steps & ---\\
MIMAC set 3 & 11 products, 354 tools, routes to 547 steps & ---\\
MIMAC set 4 & 7 products, 69 tools, routes to 92 steps & ---\\
MIMAC set 5 & 21 products, 176 tools, routes to 266 steps & ---\\
MIMAC set 6 & 9 products, 228 tools, routes to 355 steps & ---\\
MIMAC set 7 & 1 product, 38 tools, one 172-step re-entrant route &
  this preview\\
SMT2020 HV/LM & high-volume/low-mix scenario fab, $\sim\!10^3$ tools & ---\\
SMT2020 LV/HM & low-volume/high-mix scenario fab, $\sim\!10^3$ tools & ---\\
SMAT2022 & the SMT2020 fabs + automated material handling & ---\\
\bottomrule
\end{tabular}
\end{center}

\noindent {\small MIMAC set characteristics are the testbed's own
dataset tables, tabulated independently by Zhou and Rose
\bookcite{zhou09}{(2009)}. One count deserves a flag: the testbed's
catalogue prose says \emph{six} datasets, and that count circulates in
the citing literature --- but set~7 (a synthetic single-product fab
adapted from a published model of Lawrence Wein's) entered the
archive days after the catalogue text's 1997 revision and was never
re-catalogued, and the archive itself ships seven. SMT2020's two
scenario fabs each ship with an engineering-lots variant; SMAT2022
extends all four SMT2020 models with a full automated-material-handling
layer --- vehicles, rails, bays, ports.}\medskip

The stages differ in which physics must enter the certificate, and
that difference is the ladder's actual content. The Mini-Fab is small
enough that every modelling decision can be witnessed line by line
--- \bookref{ch:qtsl-minifab}{the book's worked-example chapter}
gives the model and this campaign's certificates a chapter-length
reading --- yet it already carries the domain's hard structure
(re-entrance, batching, sequence-dependent setups, shared personnel,
failures), and the base rung certifies simplified models of it on
both tracks: fluid-drift
stability with the awkward layers priced by explicit adversarial
rows (work-in-progress and cycle-time moments as named successors),
and the discrete dispatch campaign, carried at this preview through
policy synthesis, wall-clock optimality for up to seven lots, and a
certified envelope
for untrusted policies. The
MIMAC stage is the same claim species at real-fab scale, where the
sign-mode census dies combinatorially: certifying it means aggregated
covers with certified refinement paths, plus an extraction pipeline
from the testbed's 1990s fixed-width dataset format to declared
network rows. The SMT2020 stage adds another order of magnitude
($\sim\!10^3$ tools) and a maintained community simulator
(PySCFabSim), so analytic rows gain a simulation oracle beside them
--- anytime-valid statistical confidence bounds (the U5 rows of
section~\ref{sec:exp-claims}) apply at every rung today and are the
honest first deliverable at the top ones. SMAT2022 adds the transport
network itself to the contended resources, and its reference
simulator is closed-shape, so the declared model carries the
semantics alone.

\textbf{Certified today: the Mini-Fab, fluid track.} The base rung's model --- the
community's minimal re-entrant fab, whose batching, setup, personnel,
and failure layers make five machines an honest miniature of the
domain --- is described at reading depth in
\bookref{ch:qtsl-minifab}{the book's worked-example chapter}, which
also walks these certificates decision by decision; here, the
campaign summary. For the fluid model of the deterministic core,
under named static-priority dispatch policies (last- and
first-buffer-first-served), with all residual scheduling freedom
resolved adversarially: certified linear-Lyapunov drain certificates
at release rates up to $95.9$ of the $96$ lots/week fluid capacity
($99.90\%$), with the checker independently rebuilding every polytope
and replaying one exact-$\mathbb{Q}$ Farkas dual per mode --- and the
probe at $96$ correctly failing, because the load is genuinely at
capacity: producer-side bisection brackets the last certifiable
release within $6 \times 10^{-5}$ lots/week of the fluid capacity, so
the one-lot gap in this campaign's earlier grid was grid granularity,
not lost capacity. A setup-adversarial variant (setup times charged
worst-case inline) certifies up to $99.8\%$ of its own reduced
ceiling. Certificates carry $819$ obligations and check in
$1$--$3$\,ms; the two rows nearest capacity ride fully-refined
sign-split tree covers --- the refinement species the testbed rung
below leans on, exercised first here. The certified drain rates also
separate the two policies near capacity --- an exact
policy-comparison number of the kind this claim species was chosen to
produce.

\textbf{Certified today: the Mini-Fab, discrete track --- dispatch
synthesis, to exact optimality on a simplified model.} A simplified
discrete model of the same line, a transition net (16 places, 19 transitions, 49 arcs; every logistical
step one tick), carries a family of twelve certificates --- $150$
footprint rows of a sealed certificate container: a work-potential
ranking budgeting each lot's run at exactly $638$ firings, lot
conservation, six machine conservation bounds, congestion, a
ceiling, and an arrival-potential + step-budget pair yielding the
worst-case deadline, \emph{all $N$ lots arrived within $638N - 12$
steps}, over every scheduler interleaving. The sealed checker
evaluates only $36$ of those $150$ footprint rows, all of them rows
of the six tick transitions. The other $114$ are not evaluated. They
include every row of the other $13$ transitions (release, the six
starts, the five finishes, ship), whose outputs the container does
not carry, so those rows cannot be checked from the sealed bytes;
among them are the rows the deadline depends on.
Those outputs are checked only by mtk's producer-side verifier
(\texttt{verify\_campaign}), trusted Rust that is not itself verified
and cannot be rebuilt from public sources. So the sealed checker
alone does not establish the deadline, and the sealed bytes do not
identify the net: a variant in which lots skip step~2 seals to the
same bytes (every certificate is also true of it, so this is not a
soundness defect; ``the net'' is the one defined by the pinned mtk
source). Nothing in the
certificate scales with the lot count: $N$ enters only as three
integers in the container's own data, so the same rows instantiate
from $N = 7$ to $N = 10^{7}$ with the container growing six bytes
across six orders of magnitude and the sealed checker's wall time flat at
$\sim\!2$\,ms per container --- at the top of that sweep the sealed claims
quantify over at least $2\times10^{38}$ reachable markings and
$10^{2\times10^{7}}$ distinct maximal schedules, none of which the sealed
checker enumerates. Opening one supervisor seat (the plant re-stated as
a comb with a single policy-shaped hole; the hole count is itself a
checked claim) turns the deadline into a synthesis instrument: joint
policy+certificate search returns a latest-safe-unlock supervisor
with certified completion and all arrivals by $626N$ steps ---
$52{,}584$ against the $53{,}580$ no-policy worst case at the weekly
$N = 84$, a certified throughput-floor gain of $1.9\%$. The
synthesis objective is lexicographic --- certified completion first,
then the horizon --- and the order is load-bearing: a gate that
never opens ties the arrival optimum and then deadlocks shipping
nothing, and the search refuses it at the completion certificate
with the starved queue named.

The counting-time story transfers to the wall clock, where the
optimum is closed for up to seven lots. Under the theory's discrete clock semantics (every transition
dwells exactly one tick) the same containers reseal with identical
row tallies, and the counting bounds become tick bounds by a no-idle
argument ($\le 626N$ gated, $\le 638N - 12$ open). Beneath them sits
an optimum over all dispatch policies, measured as the round of the
$N$-th arrival into the last queue (not the round at which all lots
have shipped): $T_{\mathrm{opt}}(N) = 255N + 371 + d(N)$,
$d = (0, 0, 2, 4, 4, \dots)$. \emph{Erratum (2026-10-19): this is
established for $N \le 7$ only.} For $N = 1, \dots, 7$ it is closed
from both sides: no schedule meets $T_{\mathrm{opt}} - 1$ on a stated
relaxation of the net that also imposes a queue-order reduction (UNSAT
with LRAT proofs accepted by the CakeML-verified checker cake\_lpr, and
by an independent integer SMT encoding). Two
pen-and-paper arguments carry this to every policy: that the
relaxation contains every run of the net, and an exchange lemma
justifying the queue-order reduction. Without the reduction, the
same certificates are published for $N \le 5$.
On the upper side, an explicit supervisor family $\pi^*$ meets $T_{\mathrm{opt}}$ in a
replay against the decoded container plus a named transcription of
the net rows the container does not seal (at $N = 84$ as sealed,
$21{,}795$ ticks; and at $N = 1, \dots, 12$ with the sealed pending-lot
count and unlock threshold substituted). A
machine-contiguity floor of $255N + 375$ for $N \ge 4$ is proved in
Lean~4 over mathlib, with adversarial mutants refused by the kernel,
but for an abstract schedule model onto which the net does not
project exactly: the net also allows a tick in a lot's own start
round and a same-round handover of a diffusion machine. So for
$N \ge 8$ the floor is not yet proved for the net, and optimality
there is open, although $\pi^*$ is replayed at $255N + 375$. An
earlier version of this paragraph claimed a Lean floor over all
policies and a two-sided optimum at every $N$; both overclaimed. The
evidence and its exact scope are in
\texttt{semicon/minifab-floor/README.md}. Accordingly the
certified-optimality column of the dispatch grid reads $100.00\%$ as
certified for $N \le 7$; at the weekly $N = 84$ it is measured
against the attained but unproved $255N + 375$. The grid
around it separates the levers. The sixteen never-idle full-dispatch
policies deliver $99.63$, $75.61$, $65.34$, and $54.35\%$ of
that reference, and only the diffusion-station
pull preference ever moves a number; at $N = 84$ the tested
deliberately idling supervisor family improves on the best tested
work-conserving policy ($255N + 455$ ticks) by $80$ ticks, though
whether some untested work-conserving policy closes that gap is open;
and the wall-clock value
concentrates in the dispatch seat: an adversarial pull order costs
$+83.5\%$, while the synthesized gate --- whose value is the
certified worst-case floor, a guarantee over every schedule ---
moves typical-case arrivals $0.0\%$. Optimality here is also the
decision-theoretically right word: with deterministic closed
dynamics, imprecise-decision-theoretic policy selection provably
collapses to makespan minimization (machine-checked, with a
refutation cell showing the collapse fails without determinism).

\emph{Model note (2026-10-24).} Everything on this discrete track is
about the discrete net, except the Lean floor, which is about the
abstract schedule model named above. The net departs from Kempf's 1994
specification in a way that matters for optimality. Under the
wall-clock semantics each transition fires at most once per round.
So in any one round, at a given process step, at most one lot can
start, at most one can receive a unit of service, and at most one can
finish. A station's two machines can still work in the same round on
lots at different steps, but never on two lots at the same step.
Kempf's specification allows both. Within their stated scopes, all of
the results on this track are about this net, including the policy
comparisons above and the envelope below.
Whether they carry over to a net with Kempf's per-machine parallelism,
and more generally to Kempf's full specification (batching, setups,
personnel, failures), is open. A re-model with per-machine
parallelism is in progress.

The track's last exhibit prices the case nobody certifies: an
\emph{uncertified} policy --- a lookahead heuristic standing in for
a learned dispatcher --- filling the same seat behind a certified
switching envelope (certified baseline, decision-budget bench,
certified recovery region). The envelope's certificates never read
the policy. In replay, the heuristic ran free and attained the
envelope's arrival bound $638N - 11$ exactly at $N = 84$, so the
bound is witnessed tight, not merely sound; a deliberate idler was
benched at the budget boundary with recovery at exactly $1{,}009$
firings; a rogue policy was benched after eight named in-loop
refusals with the plant untouched; and after all runs the sealed
container re-verified byte-identical. The whole discrete family ---
closed baseline, open comb, synthesized supervisor, switching
envelope, and the wall-clock reseals of baseline, comb, and
supervisor --- re-checks from bytes alone by one reproduction target
(section~\ref{sec:exp-repro}), with committed integrity mutants that
must refuse and the producer-tier mutant batteries re-run beside it;
\bookref{ch:qtsl-minifab}{the book's worked-example chapter} walks
the model, the synthesis, and the grid at chapter length.

\textbf{Certified today: MIMAC set 7, at the cover and one
refinement in} (the smallest of the seven by tool count: one product,
38 machines in 24 tool groups, one 172-step re-entrant route). The
structural problem at this scale is that mode-by-mode certification
is $2^{172}$-large; the certificate instead quantifies over an
\emph{aggregated cover} --- cells whose relaxed polytopes provably
over-approximate every sign mode, by a lemma whose only move is
dropping constraints (so soundness is monotone by construction). The
honest price appears in the bound, and so, now, does the refinement
path. At the coarsest cover ($172$ cells; proven weight-optimal for
that cover) the certified stable load is $13.2\%$ of the dataset's
nominal release (availability-adjusted variant: $11.5\%$). Walking
the certificate format's built-in refinement path (binary sign-split
trees, anytime-refinable down to the exact mode walk) --- with the
weight search and every cell's dual solved as one sparse linear
program --- lifts this to a certified $1964/9903 \approx 19.8\%$ of
nominal on a $244$-cell tree ($1.50\times$ the coarsest cover;
$84{,}180$ obligations, replayed by the checker exactly like every
other row). Both numbers are weak, both are displayed at full
prominence, and the certified $19.8\%$ is a floor rather than a
ceiling: the refinement run was stopped early, with producer-side
bisection already bracketing $23$--$24\%$ at five times the cells
(measured, not yet certified) --- and the measured refinement curve
is sublinear enough that explicit trees cannot reach nominal load at
fab scale, so pricing that species limit, and the implicit-cover
species that removes it, is named successor work. The aggregation
gap is a measured frontier, not fine print.

\paragraph{What is not claimed.} On the fluid track these are claims
about declared fluid models of the fabs, with batching, setups, and
failures handled by explicit adversarial or interval rows where
modelled at all --- the model gap to the discrete line rides a
declared ledger, and the conversion from certified drain rates to
work-in-progress and cycle-time bounds on the discrete line (via
Foster--Lyapunov moment rows and Little's law) is named successor
work. On the discrete track the claims are exact, but the model sits
in the unit-dwell deterministic corner: zero stochastic
species instantiated, so a policy-filled closure has one trajectory
per lot count, and the demonic quantifiers live at the family level
(the all-policy floors; the envelope over an untrusted filler).
Since every certificate quantifies over the full scheduling freedom,
its claims dominate every rate-raced refinement of that freedom ---
the stochastic layers, when they arrive on the higher rungs, resolve
freedom the certificates have already priced. That argument covers
only refinements that restrict runs. Per-machine parallelism would add
runs, so the argument does not cover it (model note above). An
earlier preview of
this section stated that no claim was made about optimality of the
dispatch policies, only their certified stability region; that
sentence is now true of the fluid track alone. On the discrete
Mini-Fab net the optimality question is closed for $N \le 7$ and open
beyond, where attainment at $255N + 375$ is replayed but the floor is
proved only for the abstract model (erratum above). Whether any of it
carries over to Kempf's specification is open (model note above). The
correction is recorded here rather than silently edited away, so that
readers of the earlier preview can see what changed. The ladder's remaining nine rungs
--- MIMAC sets 1--6, the SMT2020 pair, and SMAT2022 --- are
census-scoped but carry no certificates at this preview.
 \providecommand{\LeanRouteExponent}{3.58}
\providecommand{\LeanRouteExponentLo}{3.24}
\providecommand{\LeanRouteExponentHi}{3.92}
\providecommand{\LeanRouteRSquared}{0.989}
\providecommand{\LeanRoutePointCount}{9}
\providecommand{\LeanRouteBudgetMin}{10}
\providecommand{\LeanRouteThresholdKB}{28}
\providecommand{\LeanRouteThresholdBytes}{27795}
\providecommand{\LeanRouteHost}{AMD EPYC 7B13, 16 threads, 126\,GiB}
\providecommand{\LeanRouteFitA}{-30.207075}
\providecommand{\LeanRouteFitB}{3.577185}
\providecommand{\LeanRouteFitDomainLo}{919}
\providecommand{\LeanRouteFitDomainHi}{24686}
\providecommand{\LeanRouteRSqBytes}{0.989}
\providecommand{\LeanRouteRSqRows}{0.962}
\providecommand{\LeanRouteRSqNodes}{0.909}
\providecommand{\LeanRouteCensoredCount}{5}
\providecommand{\LeanRouteGlcCount}{15}
\providecommand{\LeanRouteGlcMedian}{0.37\,s}
\providecommand{\LeanRouteGlcMax}{14\,s}
\providecommand{\LeanRouteQtslabAcceptCount}{7}
\providecommand{\LeanRouteQtslabCensoredCount}{7}
\providecommand{\LeanRouteQtslabOofCount}{1}
\providecommand{\LeanRouteArmHost}{Cortex-X925 + Cortex-A725}
\providecommand{\GlcRouteExponent}{1.20}
\providecommand{\GlcRouteExponentLo}{0.63}
\providecommand{\GlcRouteExponentHi}{1.78}
\providecommand{\GlcRouteRSquared}{0.61}
\providecommand{\GlcRouteFitA}{-13.023158}
\providecommand{\GlcRouteFitB}{1.203744}
\providecommand{\GlcRouteFitDomainLo}{492}
\providecommand{\GlcRouteFitDomainHi}{126115}
\providecommand{\QtslabRouteExponent}{3.10}
\providecommand{\QtslabRouteExponentLo}{2.30}
\providecommand{\QtslabRouteExponentHi}{3.91}
\providecommand{\QtslabRouteRSquared}{0.95}
\providecommand{\QtslabRouteFitA}{-24.411397}
\providecommand{\QtslabRouteFitB}{3.101501}
\providecommand{\QtslabRouteFitDomainLo}{925}
\providecommand{\QtslabRouteFitDomainHi}{17597}

\section{The Lean route: how much Lean?}
\label{sec:exp-lean-route}

The toolkit's third checking path --- the last of the three checking
routes whose division of labor the front matter states
(\bookref{ch:exec-summary}{the Executive Summary}) --- loads the same
sealed certificate container that every other checker consumes into
Lean, reconstructs the claim, and lets Lean accept it as a theorem. The route's decode,
reconstruction, and validity elaboration are ordinary Lean definitions;
its acceptance theorem says that anything the route accepts is
semantically valid --- valid, that is, under every admissible
valuation of the hypernet the certificate declares: the same
Bernstein Logic valuation semantics the toolkit's other checkers
answer to --- and the proof is one line, because the accepted
value carries the validity proof built during elaboration. The same
definitions are consumed two ways: \emph{compiled}, as a native
executable produced by the Lean compiler, and \emph{kernel}, where a
generated receipt module asks the kernel itself to evaluate the entire
pipeline over the embedded bytes (\texttt{decide}) --- no compiler, no
extraction, just the trusted evaluator. The route exists to make a
measured statement, not a rhetorical one: the expectation has always
been that the kernel path is far too slow for machine-scale
certificates, and this section measures how far.

\subsection{The ladder}

The measurement corpus is a size ladder of fifteen production
certificates --- sealed containers from the working producer stack,
spanning $603$ bytes to $126{,}216$ bytes ($\times 209$) across seven
certificate families, each rung accepted by the independent reference
checker before any Lean timing was attempted. The ladder is committed
byte-for-byte in \path{lean-route/fixtures/}; as elsewhere in this
repository, the certificates stand on their own, and no producer is
part of the trust story. All timings in this section were taken on one
otherwise-idle machine (\LeanRouteHost), single job at a time, with a
\LeanRouteBudgetMin-minute budget per certificate per route;
Table~\ref{tab:lean-route-ladder} shows the ladder. An em-dash marks a
certificate whose run exceeded the budget; \emph{o.o.m.}\ marks one
that exhausted the machine's memory first; \emph{o.o.f.}\ (out of
fragment) marks the scope boundary discussed below --- the route
declines any certificate outside its accepted fragment rather than
guess at it, a refusal that costs completeness and never soundness ---
and kernel peaks for non-completing runs are reported as floors (the
resident size when the run was cut).

\begin{table}[htbp]
\centering\footnotesize
{\setlength{\tabcolsep}{4pt}\begin{tabular}{@{}lrrrrr@{}}
\toprule
certificate & bytes & proof nodes & compiled & kernel & kernel peak \\
\midrule
\texttt{enclosure\_v1\_golden\_cos6} & 603 & 1 & \emph{o.o.f.} & \emph{o.o.f.} &  \\
\texttt{ctmc\_sis\_gen} & 919 & 3 & 2.4\,ms & \emph{o.o.m.} & $\geq$124\,GiB \\
\texttt{ctmc\_fk\_mid} & 1{,}793 & 35 & 51\,ms & \emph{o.o.m.} & $\geq$124\,GiB \\
\texttt{micro2\_stable} & 2{,}184 & 96 & 0.12\,s & --- & $\geq$0.75\,GiB \\
\texttt{cartpole\_core\_decrease\_s0001} & 4{,}068 & 65 & 0.43\,s & --- & $\geq$0.76\,GiB \\
\texttt{rode\_uniform\_decrease} & 7{,}652 & 143 & 2.9\,s & --- & $\geq$0.76\,GiB \\
\texttt{cartpole\_dp\_core\_decrease} & 9{,}336 & 197 & 9.6\,s & n.a. &  \\
\texttt{dp\_composite\_syzygy} & 17{,}591 & 377 & 99\,s & n.a. &  \\
\texttt{dp\_composite\_eigenrate} & 22{,}056 & 647 & 7.4\,min & n.a. &  \\
\texttt{lti\_sampled\_flow} & 24{,}686 & 1{,}693 & 8\,min & n.a. &  \\
\texttt{acasx4\_init\_s0000} & 43{,}401 & 527 & --- & n.a. &  \\
\texttt{dp\_composite\_smallgain} & 46{,}015 & 1{,}081 & --- & n.a. &  \\
\texttt{acasx4\_promise\_closedloop} & 57{,}779 & 1{,}055 & --- & n.a. &  \\
\texttt{rode\_nest\_shell} & 119{,}217 & 4{,}629 & --- & n.a. &  \\
\texttt{rode\_hold\_flow} & 126{,}216 & 5{,}773 & --- & n.a. &  \\
\bottomrule
\end{tabular}}
 \caption{The ladder. ``compiled'' is the median end-to-end wall clock
of repeated passes of the compiled route; ``kernel'' is the wall clock
of one kernel (\texttt{decide}) evaluation of the same bytes. Em-dashes
are runs cut at the \LeanRouteBudgetMin-minute budget; \emph{n.a.}\
marks kernel runs not attempted because every in-fragment kernel
attempt at smaller sizes had already ended cut or out of memory.}
\label{tab:lean-route-ladder}
\end{table}

\subsection{The scaling law}

Across the \LeanRoutePointCount{} rungs the compiled route completes
within budget, end-to-end time follows a power law in container bytes:
$t \propto B^{\LeanRouteExponent}$ (95\% CI $[\LeanRouteExponentLo,
\LeanRouteExponentHi]$, $r^2 = \LeanRouteRSquared$ on log--log least
squares). Container bytes is the best of the three natural size units as a
regressor ($r^2 = \LeanRouteRSqBytes$, against \LeanRouteRSqRows{} for
certificate rows and \LeanRouteRSqNodes{} for proof nodes) --- expected,
since the dominant cost is the reader: list-based CBOR decoding and
seam assembly in a purely functional style, chosen so that the same
definitions remain evaluable by the kernel. Nothing about the
certificate format requires a superlinear reader, as the linear-time
imported checkers on the same bytes demonstrate --- and as the
fast-path route's own fitted exponent on this same ladder now
measures (\GlcRouteExponent, its confidence interval containing
linear). Figure~\ref{fig:lean-route-scaling}
shows the measured points, the fitted law, and the budget line ---
and, since the trichotomy's claim is comparative, the same ladder
through the toolkit's other two checking routes, measured per
certificate on the Arm host
(section~\ref{sec:exp-lean-route-comparison} carries the
hardware and scope contract).

\begin{figure}[htbp]
\centering
\begin{tikzpicture}
\colorlet{routeEngine}{blue!70!black}
\colorlet{routeCertified}{green!45!black}
\colorlet{routeLean}{red!80!black}
\begin{loglogaxis}[
    width=0.82\linewidth, height=0.62\linewidth,
    xlabel={sealed container bytes},
    ylabel={end-to-end wall clock (s)},
    ymin=5e-4, ymax=5e3,
    grid=major, grid style={black!12},
    legend cell align=left, legend columns=2,
    legend style={draw=black!30, font=\small,
      at={(0.5,-0.20)}, anchor=north, /tikz/every even column/.append style={column sep=1em}},
    tick label style={font=\small},
]
\addplot[routeLean!55, thick, domain=\LeanRouteFitDomainLo:130000, samples=2]
    {exp(\LeanRouteFitA)*x^\LeanRouteFitB};
\addlegendentry{Lean fit $t \propto B^{\LeanRouteExponent}$}
\addplot[routeEngine!55, thick, domain=\GlcRouteFitDomainLo:\GlcRouteFitDomainHi, samples=2]
    {exp(\GlcRouteFitA)*x^\GlcRouteFitB};
\addlegendentry{\texttt{geolog-alpha} fit $t \propto B^{\GlcRouteExponent}$}
\addplot[routeCertified!55, thick, domain=\QtslabRouteFitDomainLo:\QtslabRouteFitDomainHi, samples=2]
    {exp(\QtslabRouteFitA)*x^\QtslabRouteFitB};
\addlegendentry{\texttt{qtslab} fit $t \propto B^{\QtslabRouteExponent}$}
\addplot[dashed, black!60, domain=550:140000, samples=2] {600};
\addlegendentry{\LeanRouteBudgetMin-minute budget}
\addplot[only marks, mark=square*, mark size=1.8pt, routeEngine]
    table {generated/lean-route-glc-points.dat};
\addlegendentry{\texttt{geolog-alpha} (fast path)}
\addplot[only marks, mark=diamond*, mark size=2.4pt, routeCertified]
    table {generated/lean-route-qtslab-points.dat};
\addlegendentry{\texttt{qtslab} (certified)}
\ifnum\LeanRouteQtslabCensoredCount>0
\addplot[only marks, mark=diamond, mark size=2.8pt, routeCertified]
    table {generated/lean-route-qtslab-censored.dat};
\addlegendentry{\texttt{qtslab}, cut at budget}
\fi
\addplot[only marks, mark=*, mark size=1.7pt, routeLean]
    table {generated/lean-route-points.dat};
\addlegendentry{Lean (compiled reflection)}
\ifnum\LeanRouteCensoredCount>0
\addplot[only marks, mark=o, mark size=2.2pt, routeLean]
    table {generated/lean-route-censored.dat};
\addlegendentry{Lean, cut at budget}
\fi
\end{loglogaxis}
\end{tikzpicture}
\caption{One ladder, three checking routes: the fast-path engine
(\texttt{geolog-alpha}, blue squares, on the staged-chain sealed
containers), the certified checker (\texttt{qtslab}, green diamonds,
on the same certificates re-marked content-resolved --- a
container-table transform under which every certificate row is
byte-identical), and the compiled Lean reflection route (red circles)
with its fitted law. Blue and green medians are one process, end to
end, pinned to one core of the Arm host; the red points are this
section's own protocol on its dedicated machine (the footnote of
section~\ref{sec:exp-lean-route-comparison} carries both). Open marks
of either color sit \emph{at} the dashed
budget line: each true time lies above its mark. A rung with no mark
on some route is a scope datum, not a failure --- that route classes
the certificate outside its accepted fragment and declines it (the
caption of Table~\ref{tab:lean-route-ladder} and
section~\ref{sec:exp-lean-route-comparison}'s footnote carry the
per-route scope and hardware contract). The Lean fit uses only
uncensored red points; where it crosses the budget ---
$\approx\LeanRouteThresholdKB$\,KB of sealed container --- is the
admission threshold of Table~\ref{tab:lean-route-admission}. Each
solid line is the same log--log fit to its route's accepted medians:
\texttt{geolog-alpha} at exponent \GlcRouteExponent{} (95\% CI
$[\GlcRouteExponentLo,\GlcRouteExponentHi]$, containing linear;
$r^2=\GlcRouteRSquared$), \texttt{qtslab} at \QtslabRouteExponent{}
($[\QtslabRouteExponentLo,\QtslabRouteExponentHi]$,
$r^2=\QtslabRouteRSquared$), Lean at \LeanRouteExponent{}
($[\LeanRouteExponentLo,\LeanRouteExponentHi]$) --- one near-linear
route and two superlinear ones, measured on the same certificates.}
\label{fig:lean-route-scaling}
\end{figure}

The kernel mode needs no fitted curve to make its point. Its smallest
full acceptance during the route's development --- the $919$-byte
certificate, second rung of the ladder --- cost $965$ seconds and a
$233$\,GiB peak on a $251$\,GiB development machine (the toolkit
repository's route documentation carries that receipt); the memory is
the kernel's reduction state, not the certificate. On this section's
$126$\,GiB reference machine the same certificate exhausts neither
budget nor patience but memory, and every in-fragment kernel attempt
on the ladder ended cut at the budget, out of memory, or predictably
unattempted (Table~\ref{tab:lean-route-ladder}). The failure mode is
content-dependent, not size-monotone: the $919$-byte certificate took
the whole machine's memory in under eight minutes, while certificates
two to eight times its size ground past the budget inside a
gigabyte --- a $\times 160$ spread in peak residency between
neighboring rungs. What a given certificate makes the kernel's
reduction state do is a property of its content (the working
hypothesis is exact-rational arithmetic in kernel terms, untested
here), not of its byte count. A certificate of under a kilobyte
that the compiled route accepts in milliseconds is already past what
the kernel evaluator will do inside a ten-minute, hundred-gigabyte
envelope; that is the measured content of ``too slow at machine
scale,'' stated with its numbers attached.

\subsection{The fragment boundary}

One ladder rung is not a timing datum but a scope datum: the smallest
certificate ($603$ bytes, an interval-enclosure certificate whose
obligations are discharged by \emph{primitive applications}) is out of
fragment for the route --- the table's \emph{o.o.f.} --- and the route
declines it: compiled and kernel agree, and the kernel's refusal is
itself a theorem (\texttt{decide} proves the acceptance predicate
false on those bytes). The same bytes are accepted by the reference
checker, which settles the direction of the disagreement: the route's
accepted fragment is the first-order pinned-vocabulary rule set,
primitive applications lie outside it by design, and a checker asked
about a certificate it cannot speak for declines it rather than
guess. Every checker in this chapter draws some such boundary, always
with the same sign --- an out-of-fragment refusal costs completeness,
never soundness; the certified checker classes this same rung outside
\emph{its} fragment too, naming the reason in its verdict
(\texttt{fun\_outside\_fragment}), while the fast-path lane discharges
it (Figure~\ref{fig:lean-route-scaling}). Admission to the Lean route
is therefore a conjunction: small enough, \emph{and} inside the
fragment. A species' first Lean attempt should be its smallest
instance, which settles fragment membership in milliseconds before
any batch of timings is scheduled.

\subsection{The admission gate}

Table~\ref{tab:lean-route-admission} applies the fitted law to the
certificate species of this chapter's campaigns: per species, the
byte range of its committed certificates (container overhead added),
the predicted compiled-route time at the median and maximum, and the
verdict against the \LeanRouteBudgetMin-minute budget. The table is the
gate the re-run campaigns consume: a Lean-reflection timing column is
attempted only for species that pass (or, at the runner's discretion,
the passing fraction of a \emph{partial} species), and the kernel mode
is attempted for none --- the measured points above place its envelope
below every production species' floor.

\begin{table}[htbp]
\centering\footnotesize
{\setlength{\tabcolsep}{4pt}\begin{tabular}{@{}p{10.5em}rllll@{}}
\toprule
species & max B & fragment & pred.\ med. & pred.\ max & verdict \\
\midrule
hardware liveness (circuits) & 2{,}627 & untested & 9.1\,ms & 0.22\,s & pass \\
Atari certified replay & 5{,}113 & untested & 22\,ms & 1.9\,s & pass \\
FrozenLake 4x4 (det + slippery + sr34) & 52{,}015 & exp.\ in & 15\,s & 97\,min & partial \\
CliffWalking & 71{,}872 & exp.\ in & 14\,min & 5.1\,h & fail \\
FrozenLake 8x8 & 500{,}289 & exp.\ in & 5.5\,h & 0.59\,yr & fail \\
Blackjack & 115{,}779 & exp.\ in & 10\,h & 28\,h & fail \\
Taxi (all weather/passenger variants) & 3{,}102{,}602 & exp.\ in & 7\,yr & 400\,yr & fail \\
LunarLander / LunarLanderContinuous (statistical tier) & 12{,}226{,}428 & untested & $\sim$$10^{4}$\,yr & $\sim$$10^{5}$\,yr & fail \\
published-table supermartingale corpus & 50{,}000 & exp.\ in & 97\,ms & 84\,min & partial \\
semiconductor fluid ladder (minifab) & 18{,}451 & untested & 65\,s & 2.5\,min & pass \\
semiconductor fluid ladder (mimac) & 401{,}455 & untested & 0.18\,yr & 0.27\,yr & fail \\
interval-enclosure certificates (primitive applications) & 603 & out & 4.4\,ms & 4.4\,ms & \emph{out of fragment} \\
\bottomrule
\end{tabular}}
 \caption{The admission gate. Predictions apply the fitted law to each
species' certificate sizes in sealed-container bytes (census rows
stated in committed-artifact bytes carry the measured container
overhead as a declared correction); beyond the fitted range
($\approx 25$\,KB) they are
extrapolations --- one-sidedly corroborated by the budget-censored
rungs up to $126$\,KB, order-of-magnitude statements past that.
``Fragment'' marks measured membership only for the ladder's own
species; \emph{untested} species settle theirs with a one-instance
probe before any batch attempt.}
\label{tab:lean-route-admission}
\end{table}

\subsection{The measurement before this one}
\label{sec:exp-lean-route-history}

This ladder is the route's second measurement, not its first. In July
2026, on the toolkit's predecessor binary certificate format (since
retired), the same question was measured by the other road into the
kernel --- \emph{elaboration}: an emitter translated certificate bytes
into generated Lean source ($\times 8$ expansion, byte-stable per-node
templates across certificate kinds), and the elaborator produced
kernel-checked object files (a further $\times 24$--$25$) closing with
the same species of validity theorem. Three production certificates
were taken to full kernel verdicts that way --- $197$, $377$, and
$1{,}081$ proof nodes, at $3.4$, $9$--$13$, and $28$ minutes of
elaboration respectively, against $3$--$10$ \emph{seconds} for the
certified checker build of the same era on identical certificate
content and comparable hardware: a measured factor of $\times 30$ to
$\times 250$ at kilobyte scale, growing with size. The same campaign
took a streaming census of a real machine-scale certificate --- a
$43{,}285{,}639$-node, $1.24$\,GB closed-loop neural-network
verification certificate, semantically re-checked by a reference
checker in $2.4$ hours --- and projected the elaboration route on it:
$\approx 7$\,GB of generated Lean source, $\approx 66$ million
declarations in one namespace (about one hundred mathlibs),
$\approx 176$\,GB of object files, and $1.4$--$2.7$ core-years of
serial elaboration, with memory the first wall. The numbers, the
census, and the projection's stated assumptions are recorded in
\path{lean-route/HISTORY.md}. That projection is why the present route
abandons source generation altogether: a fixed decoder and checker are
defined once, certificates arrive as ordinary input bytes, and the
measured question becomes the reader's scaling law above rather than
gigabyte-scale generated source.

\subsection{The route beside the other two}
\label{sec:exp-lean-route-comparison}

The trichotomy's claim is comparative, so the honest summary is
too.\footnote{Hardware, per the conventions: the certified checker is
timed on one core of an Arm host --- the ISA its soundness theorem
grounds in --- with its own cross-machine spread measured under
$\times 1.15$; in Figure~\ref{fig:lean-route-scaling} the two
production routes' ladder points are one process per certificate,
pinned to one performance core of one Arm box (\LeanRouteArmHost),
while the red points remain this section's reflection timings
(\LeanRouteHost); the elaboration-era numbers are as recorded in
\path{lean-route/HISTORY.md}. Every contrast quoted here is two or
more orders of magnitude, and the measured cross-machine factor on the
compiled route is a factor of a few --- hardware choice cannot close
any of them. Container forms, per route: the fast path consumes the
staged-chain re-seals (\path{lean-route/GLC-CERT.md}), the certified
checker the content-resolved mints
(\path{lean-route/tools/mint-qtslab-vintage.py} states the
container-table-only transform), and both carry the same certificate
rows byte-for-byte as the committed fixtures the Lean route reads.
Scope, shared by all three routes on this corpus: these
certificates natively encode the Bernstein Logic fragment, which is
the only fragment the certified checker consumes
(\path{theories/BernsteinLogic.geolog}); the full QTSL theory is
checked on the engine-side container lanes of
section~\ref{sec:exp-repro}.} On this chapter's stored supermartingale corpus ---
$4{,}105$ sealed containers, each one verification condition's
Bernstein Logic discharge at roughly one to fifty kilobytes --- the
certified checker's current build replays every container at a
pooled median of $30$\,ms, maximum $117$\,ms, on one Arm core
(Table~\ref{tab:exp-certified-checker}), and the fast-path engine's
sealed-container lane re-checks the whole corpus in wall-clock seconds
(\path{supermartingales/GLC-CERT.md}). The compiled reflection route
has not been batch-run on that corpus; the admission gate prices it
(Table~\ref{tab:lean-route-admission}): verdict \emph{partial}, with a
predicted median near $0.1$\,s and a predicted maximum near $84$
minutes, so only the passing fraction would ever be attempted.
Figure~\ref{fig:lean-route-scaling} makes the comparison direct on
the ladder itself, every route in its current build on identical
certificate content: the fast path near-linear in bytes (exponent
\GlcRouteExponent), the reflection reader at $B^{\LeanRouteExponent}$,
and the certified build accepting \LeanRouteQtslabAcceptCount{} rungs
and cut at the budget on \LeanRouteQtslabCensoredCount. Within the
ladder the certified wall grows as $B^{\QtslabRouteExponent}$
($r^2=\QtslabRouteRSquared$) --- statistically indistinguishable from
the reflection reader's exponent, while the fast path's confidence
interval excludes both --- and across corpora content dominates
bytes: the stored corpus above tops out at $117$\,ms per certificate
at sizes where the ladder's composite certificates cost minutes,
which is how a $30$\,ms pooled median there and these walls here are
both true. A
species aiming certificates at the certified route should therefore
run the same one-instance probe the admission gate prescribes for
Lean: fragment membership and cost are both content-properties,
settled cheaply on the smallest instance. The historical cross-check
stands beside the figure: on the composite small-gain rung
($1{,}081$ proof nodes, $46$\,KB sealed), the certified checker of
the elaboration era took about ten seconds on the same content in its
own format's bytes (today's build, reading the unified sealed
container, is cut at the budget on that rung), against past the
ten-minute budget through the compiled reflection route and
twenty-eight minutes through the retired elaboration route.

The conclusion the measurements support, stated as the front matter's
expectation deserves: \emph{confirmed, with its boundary measured}.
For the kernel evaluator the confirmation is categorical --- a
certificate under a kilobyte already costs sixteen minutes and a
quarter-terabyte of reduction state, and nothing larger completes at
all --- and for the compiled route it is a statement about scaling
rather than about the route: below a few kilobytes the compiled
reflection path is genuinely fast (the $919$-byte rung's $2.4$\,ms is
\emph{under} the certified build's corpus median, on different
certificates of the same format), but at $\times B^{3.6}$ against the
production checkers' linear-with-content growth, the ratio passes two
orders of magnitude by fifty kilobytes and machine scale is out of
reach by shape, not by constant factors. The route earns its keep on
the other side of the ledger: full kernel verdicts on small
certificates tie the certificate semantics to a proof kernel
scrutinized by a community that owes this toolkit nothing --- and the
measured boundary is precisely why the certified binary, not the
kernel, walks the forty-million-node instances.

\paragraph{What is not claimed.} The law is a property of this route's
current reader on one machine class, not of Lean, whose kernel and
compiler are here being used far outside their engineered envelope on
purpose; an index-structured reader would move the compiled curve
substantially, and nothing in the trichotomy's argument depends on it
not moving. Cross-machine transfer of the constants is measured to be
worth a factor of a few (the $2{,}184$-byte rung runs $27.6$\,ms on the
development machine cited above and $\approx 120$\,ms here); the
exponent, not the coefficient, is the portable content. Species
predictions inherit a byte-basis caveat until each campaign banks
per-instance container sizes, and the admission table regenerates from
those receipts by the same make target that produced it. The
elaboration-era numbers of \path{lean-route/HISTORY.md} are
development-record measurements on a retired format, quoted for the
route's history with their measured/projected distinctions preserved;
they are not re-derivable from this repository, and no other number in
this chapter rests on them. Finally, the ladder's timings were taken
against the committed fixtures; the same certificate content
re-sealed on the unified certificate schema is already committed
beside them (\path{lean-route/out/glc-baked/}, with the
content-identity argument stated in \path{lean-route/GLC-CERT.md}),
and when the ladder re-runs on those containers this section's
numbers move with them, per the reproduction discipline of
section~\ref{sec:exp-repro}.
 \section{Reproduction}
\label{sec:exp-repro}

Every campaign in this chapter ships its evidence in the coalition's
experiments repository
(\path{maitria-coalition/qtsl-experiments}): the certificates
themselves, committed byte-for-byte; their mutant batteries; the
machine-readable receipts (JSON) from which this chapter's tables are
re-derived mechanically; and the independent checkers, imported from
the \path{qtsl} repository at pinned revisions. One make target per
family re-verifies the stored certificates and mutants through the
checkers and regenerates the tables.

The certificate vocabulary itself is committed, as data: the
\path{qtsl} repository ships the theory pair
\path{theories/BernsteinLogic.geolog} --- the Bernstein Logic proof
tower, exactly the fragment the certified checker consumes --- and
\path{theories/Qtsl.geolog} --- the omnibus theory extending it with
systems, specifications, and the judgments connecting them --- and a
certificate is ordinary rows of these theories; there is no bespoke
certificate file format. That unified schema is the standing
discipline for everything this chapter reports: one certificate
vocabulary for every campaign, revision-pinned modelling-toolkit producers
used where each family has shipped producer coverage, the database engine as
the fast-path checker, and an extraction tool reducing a full-QTSL
certificate to its Bernstein Logic obligations for the certified checker.
Families without a distributed producer retain a revision, invocation, and
binary-digest provenance record rather than being described as runnable from
a public clone. The shipped form is the
sealed container --- one signed geolog commit per certificate,
re-verified from the bytes alone, with obligation and claim counts
demanded exactly: deleting a proof node leaves every surviving
obligation true, so the count gate is load-bearing, and a container
that claims nothing is refused rather than vacuously accepted.

The campaigns of this chapter predate that unification, and the
transition is in flight, not finished. Certificates produced on the
earlier staged vocabularies remain the evidence behind the tables:
they stand as the record of the campaign that produced them, still
re-checked by the make targets, and nothing in them is presented
un-checked. Re-emissions on the unified schema have landed for the
first families --- the fluid-drift, published-table-sweep,
Lean-ladder, and Gymnasium lanes, with the SPDE family's noise
spectra re-homed first under its split-container design; each
re-emission's content deltas are disclosed per container in its
family's manifest. The uniform re-conduct of every campaign through its revision-pinned
producer is the successor work the tables track: as family-scoped producer
coverage and re-conducted receipts land, the numbers re-derive from them, and
the superseded certificates retire to an archive rather than being silently
rewritten.

A producer is never part of any trust story, so the audit path is the
same on either side of the transition: read a checker, run it against
the shipped certificates and their mutants, and only then decide what
the tables mean. The published-table sweep's certificates are
additionally stored as plain-text listings of database object rows
--- the geolog encoding, human- and checker-readable alike --- each
pinned by hash in a manifest, and its containers prove content
identity against those listings at production, so the certificate a
checker reads is the certificate a reviewer can read. Per the preview
note at the front of this volume, artefacts migrate to the public
repositories as they graduate; until a cited artefact appears there,
treat its mention as a statement of what exists, not yet a link.

\part{Geolog}
\chapter{Quick Start}
\label{ch:geolog-quickstart}

In this chapter you will install geolog, write a theory and a handful
of facts in one file, and run the chase: the engine will derive the
tables your axioms entail, invent nothing you did not license, and
report --- rather than repair or reject --- the one state your theory
forbids. At the end you will look at the theory you wrote and find it
is itself rows. The whole path fits in one sitting; every command is
copy-pasteable, and every output shown is the verbatim transcript of
a real run.\footnote{The book's build regenerates these transcripts
by executing the printed commands against the committed examples,
and each fragment's header records the exact tool revision it ran
against (all runs on one commodity arm64 box).}

\section{Install}
\label{sec:gqs-install}

A Rust toolchain is the only prerequisite (stable channel;
\path{rust-toolchain.toml} pins the components).

\begin{Verbatim}[fontsize=\small,frame=leftline,rulecolor=\color{accentmid},framesep=3mm,xleftmargin=3mm]
git clone https://github.com/maitria-coalition/geolog-alpha
cd geolog-alpha
cargo install --path crates/geolog
\end{Verbatim}

\noindent A fresh clone compiled in about a minute on the commodity
arm64 box that regenerates this book's transcripts.
\texttt{geolog --help} lists ten verbs; the three
this chapter uses are \texttt{check}, \texttt{chase}, and
\texttt{rows}.

\section{A theory}
\label{sec:gqs-theory}

The example is a freight network, small enough to hold in your head:
places, direct hops, and reachability. The file
\path{examples/routes.geolog} declares one sort, two relations, and
three axioms:

\begin{Verbatim}[fontsize=\small,frame=leftline,rulecolor=\color{accentmid},framesep=3mm,xleftmargin=3mm]
theory Routes {

  Node : Set;

  // Asserted: a direct hop exists.
  Edge : [src: Node, tgt: Node] -> Prop;

  // Derived: some chain of hops exists. Rows appear exactly when the
  // axioms force them — never asserted directly.
  Path : [src: Node, tgt: Node] -> Prop;

  // `chase` marks a derivation rule: the axiom FIRES, creating rows.
  // (`report` marks a constraint: violations are judged, nothing is
  // created; an unmarked axiom defaults to report — the least-
  // dangerous reading, so derivation is always an explicit choice.)
  chase ax/path/edge : forall a : Node, b : Node.
    [src: a, tgt: b] Edge |- [src: a, tgt: b] Path;

  chase ax/path/trans : forall a : Node, b : Node, c : Node.
    [src: a, tgt: b] Path, [src: b, tgt: c] Edge
    |- [src: a, tgt: c] Path;

  // A denial: this network is required loop-free. A commit that lands
  // a cycle is not rejected — the violated denial is REPORTED as a
  // judgement, and repair is a further ordinary commit.
  report ax/no/loop : forall a : Node.
    [src: a, tgt: a] Path |- false;
}
\end{Verbatim}

Four things to read off the shape.
\begin{enumerate}
\item Relations are declared over named record fields, and axioms
are \emph{sequents}: universally quantified entailments, each
premise a conjunction of atoms, each conclusion here a single atom.
The odd one out is the third axiom, whose conclusion \texttt{false}
makes it a \emph{denial}: a state the theory forbids.
\item The word in front of each axiom's name is its \emph{mode
mark}, and the file's own comments state the semantics:
\texttt{chase} fires, \texttt{report} judges, and an unmarked axiom
means \texttt{report}. An axiom creates rows only where its author
said so in as many words.
\item Note what the mark does \emph{not} do: nothing in the
\emph{declarations} marks \texttt{Path} as derived. The mode belongs
to the axiom, and \texttt{Path} being a table nobody writes remains
an emergent fact --- about which relations your \texttt{chase}-marked
axioms conclude into and which your commits assert into --- not a
schema annotation.
\item The slash in \texttt{ax/path/edge} is nothing special:
slash-names are ordinary identifiers, used by convention for
namespacing. The \texttt{ax/} prefix, however, does real grammatical
work: it makes a statement an axiom rather than a declaration, so a
sequent written under a name like \texttt{my/rule} is a syntax
error. The Reference Manual (chapter~\ref{ch:geolog-reference})
spells out these name rules.
\end{enumerate}

Nothing travels beside the theory: a \path{.geolog} file is
self-contained --- everything it draws on is either declared in it
or imported at the top. This is the \emph{roster-totality}
discipline: every declared relation has a decided evaluator.
\texttt{ordinary} (asserted or chase-derived) is the default,
declared by writing nothing. When a relation is instead decided by a
bounded external procedure, an import line says so ---
\texttt{import <provider>.\{rel\};} --- classifying it as a
\emph{side-condition} awaiting that provider. The same import syntax
authorizes builtin ground vocabulary by name
(\texttt{import std.\{Q, q\_lt\};}). Nothing guesses on
your behalf: an \texttt{import std} line naming something outside
the builtin roster refuses at once, listing the importable names.
This chapter's example needs no imports; every relation in it is
ordinary, and the file's opening comment says so.

\section{An instance}
\label{sec:gqs-instance}

The same file carries the facts, in an \texttt{instance} block ---
element declarations first, then facts in the same record idiom the
axioms use:

\begin{Verbatim}[fontsize=\small,frame=leftline,rulecolor=\color{accentmid},framesep=3mm,xleftmargin=3mm]
instance Demo : Routes {
  hub, east, west, port : Node;

  [src: hub,  tgt: east] Edge;
  [src: hub,  tgt: west] Edge;
  [src: east, tgt: port] Edge;
}
\end{Verbatim}

\noindent Elements are declared before use; a fact naming an
undeclared element is an error, never a silently minted row.
(\texttt{geolog check} validates all of this --- declarations, field
sets, sorts --- without running anything.)

The map is worth ten seconds: \texttt{hub} ships to \texttt{east}
and to \texttt{west}, and \texttt{east} ships on to \texttt{port}.
Before running anything, guess which pairs the axioms will put in
\texttt{Path}.

\section{The chase}
\label{sec:gqs-chase}

\begin{Verbatim}[fontsize=\small,frame=leftline,rulecolor=\color{accentmid},framesep=3mm,xleftmargin=3mm]
$ geolog chase examples/routes.geolog
sorts:
  Node — 4 elements (4 seeded)
relations:
  Edge — 3 rows (3 seeded, 0 derived)
    [src: east, tgt: port] Edge
    [src: hub, tgt: east] Edge
    [src: hub, tgt: west] Edge
  Path — 4 rows (0 seeded, 4 derived)
  + [src: east, tgt: port] Path
  + [src: hub, tgt: east] Path
  + [src: hub, tgt: port] Path
  + [src: hub, tgt: west] Path
fixpoint: reached, no judgements — the model is saturated.
$ echo $?
0
\end{Verbatim}

Look at the \texttt{+} rows: tables nobody wrote.
\texttt{[src: hub, tgt: port]} is in \texttt{Path} because
\texttt{ax/path/trans} forces it from two rows you did write, and the
chase ran that forcing to its \emph{least fixpoint} --- everything the
axioms entail, nothing they do not. (The other three are your own
hops, lifted into \texttt{Path} by \texttt{ax/path/edge}: one new
reachability fact, four new rows.) ``Derived'' names provenance
rather than any lesser grade of truth: these rows are as queryable
as the seeded ones, and rerunning from the same facts re-derives
exactly the same table.

\section{Break it}
\label{sec:gqs-break}

Add one hop to the instance block --- the return leg that closes a
loop --- and save the edited file as \path{routes-cycle.geolog}:

\begin{Verbatim}[fontsize=\small,frame=leftline,rulecolor=\color{accentmid},framesep=3mm,xleftmargin=3mm]
  [src: port, tgt: hub]  Edge;
\end{Verbatim}

\begin{Verbatim}[fontsize=\small,frame=leftline,rulecolor=\color{accentmid},framesep=3mm,xleftmargin=3mm]
$ geolog chase routes-cycle.geolog
sorts:
  Node — 4 elements (4 seeded)
relations:
  Edge — 4 rows (4 seeded, 0 derived)
    ⋮
  Path — 12 rows (0 seeded, 12 derived)
    ⋮
judgements (3):
  denial violated: ax/no/loop at [a: hub]
  denial violated: ax/no/loop at [a: east]
  denial violated: ax/no/loop at [a: port]
$ echo $?
1
\end{Verbatim}

Read what did \emph{not} happen. The write was not rejected: the
facts landed, and the chase saturated over them (twelve \texttt{Path}
rows now: each of the three places on the loop reaches all four,
itself included). The
violation was not repaired: no engine deleted your hop or picked a
row to sacrifice. What happened is a \emph{judgment} --- the very
thing \texttt{ax/no/loop}'s \texttt{report} mark asks for: the
violated denial is reported, with the axiom's own name and the
binding that witnesses it, and the exit code says so. Three
judgments here, not one: the same loop, witnessed once from each
place standing on it. If you have
operated an event-sourced system, this is a familiar posture done
formally: an
append-only log can land in a state some invariant marks invalid, and
the repair is a further ordinary commit --- made by whoever owns the
workflow, not by the database guessing. The Tutorial
(chapter~\ref{ch:geolog-tutorial}) makes this the centre of the
story.

\section{The theory is rows}
\label{sec:gqs-rows}

One last look before you go:

\begin{Verbatim}[fontsize=\small,frame=leftline,rulecolor=\color{accentmid},framesep=3mm,xleftmargin=3mm]
$ geolog rows examples/routes.geolog
// Declaration rows in the GeologMeta idiom (the schema you wrote is
// itself data). refs: L<n> = local element minted at op[n]; G<name> =
// genesis identity. This is the declaration slice; axioms are rows
// too, in the same substrate — the committed listings under
// theories/generated/ carry complete worked examples.
op[   0] elta sort=G<Theory>                    // Theory "Routes"
op[   1] elta sort=G<Sort>                      // Sort "Node"
op[   2] funa fun=G<Sort/name> args=[L1] value=str:"Node"
op[   3] funa fun=G<Sort/theory> args=[L1] value=L0
op[   4] elta sort=G<Fun>                       // Fun "Edge"
op[   5] funa fun=G<Fun/name> args=[L4] value=str:"Edge"
    ⋮
op[   8] elta sort=G<Field>                     // Field "src" of Edge
op[   9] funa fun=G<Field/fun> args=[L8] value=L4
    ⋮
\end{Verbatim}

The theory you wrote in \S\ref{sec:gqs-theory} is itself data: a
\texttt{Theory} element, a \texttt{Sort} element with a name row and
a home row, \texttt{Fun} and \texttt{Field} elements for each
relation and its columns --- rows in the \emph{meta-schema}, the
schema language that is itself a geolog theory. As the listing's own
header says, this is the declaration slice; axioms continue as rows
in the same idiom, and the repository's \path{theories/generated/}
directory carries complete committed listings --- the full Petri-net
theory runs to 1{,}239 ops.
Schemas, extensions, and migrations live in the same substrate they
govern, so you can version them, diff them, and derive facts about
them just as you would any other rows.

\section{Where to next}
\label{sec:gqs-next}

The Tutorial (chapter~\ref{ch:geolog-tutorial}) works an event-sourced
order flow with two writers: views only the merged logs can derive,
two replicas' records identified by congruence rather than a conflict
resolver, a witness the chase invents under an existential axiom, a
violation that only the merge surfaces, and the signed commit log
underneath, which can say things a broker cannot. The Reference Manual
(chapter~\ref{ch:geolog-reference}) gives the systematic account. The
part closes with diagrams as data (chapter~\ref{ch:qtsl-scomplexes}):
one incidence genus, stored and checked as rows, for every diagram
the wider toolkit traffics in. And
for the road from a model to a checked temporal claim about it ---
the reason this database exists in this book --- start the QTSL part
at chapter~\ref{ch:qtsl-quickstart}.

 \chapter{Tutorial}
\label{ch:geolog-tutorial}

\noindent\draftmark\medskip

If you have operated an event-sourced system, you have already run
most of geolog's architecture on convention \cite{kleppmann21}. An
append-only log serves as the source of truth. Derived views are
rebuilt by replay. New capability arrives as new producers and
consumers, never as mutation of shared state. And everybody promises
not to write to the derived tables by hand. Kafka gives you the log,
Flink or Feldera maintains the views, a schema registry polices the
shapes, and the promises live in code review.

geolog takes that architecture and makes the promises structural. No
fold code that somebody wrote and everyone else must believe maintains
the derived views; the declared axioms \emph{entail} them instead, and
they can be recomputed and checked from the log alone. The chapter's
spine is one example, an order flow with two writers: views no single
writer can derive (\S\ref{sec:gt-views}--\ref{sec:gt-single}), merge
as set union plus re-derivation (\S\ref{sec:gt-merge}), the signed
commit log and what it can say that a broker cannot
(\S\ref{sec:gt-log}), schemas as data (\S\ref{sec:gt-schema}), and the
ground vocabulary (\S\ref{sec:gt-ground}). A worked example imported
from an independent project follows: the core theory of a proof
assistant, expressed whole, together with instructions for converting
its native data in both directions (\S\ref{sec:gt-alifib}). A
plain-dealing scorecard against the SQL stack you already run
(\S\ref{sec:gt-scorecard}) and the handoff to temporal claims
(\S\ref{sec:gt-onward}) close the chapter. Every transcript is
verbatim from a real run against the geolog-alpha repository (the
order flow against \path{examples/orders.geolog} and the
single-diagram crossing of \S\ref{sec:gt-alifib}, last re-executed
2026-07-30 against the import-era tree; the corpus-scale
measurements behind \S\ref{sec:gt-alifib} executed 2026-07-27); the
receipts sit at this chapter's source tail.

\section{Views you do not maintain}
\label{sec:gt-views}

Here is the domain, whole. Two writers will keep logs against it: a
storefront (payments, cancellations) and a warehouse (allocations,
shipping records).

\begin{Verbatim}[fontsize=\small,frame=leftline,rulecolor=\color{accentmid},framesep=3mm,xleftmargin=3mm]
theory Orders {

  Order    : Set;
  Shipment : Set;

  // Events -- asserted rows, the writers' ground truth.
  Placed    : [o: Order] -> Prop;
  Paid      : [o: Order] -> Prop;
  Allocated : [o: Order] -> Prop;
  Cancelled : [o: Order] -> Prop;

  // Views -- derived rows, never asserted by a writer...
  Shippable : [o: Order] -> Prop;

  // ...except ShippedVia, which writers MAY also assert (a recorded
  // manifest); where no record exists, the chase invents the witness.
  ShippedVia : [o: Order, s: Shipment] -> Prop;

  chase ax/ship/ready : forall o : Order.
    [o: o] Paid, [o: o] Allocated |- [o: o] Shippable;

  chase ax/ship/manifest : forall o : Order.
    [o: o] Shippable |- exists s : Shipment. [o: o, s: s] ShippedVia;

  chase ax/manifest/unique : forall o : Order, s1 : Shipment, s2 : Shipment.
    [o: o, s: s1] ShippedVia, [o: o, s: s2] ShippedVia |- s1 = s2;

  report ax/no/zombie : forall o : Order, s : Shipment.
    [o: o] Cancelled, [o: o, s: s] ShippedVia |- false;
}
\end{Verbatim}

Read \texttt{ax/ship/ready} the way you would read a view definition,
because that is what it is:

\begin{Verbatim}[fontsize=\small,frame=leftline,rulecolor=\color{accentmid},framesep=3mm,xleftmargin=3mm]
CREATE MATERIALIZED VIEW shippable AS
  SELECT o FROM paid INTERSECT SELECT o FROM allocated;
\end{Verbatim}

\noindent The \texttt{chase} mark plays the role of \texttt{CREATE}
here. It opts the sequent into \emph{firing}, which is to say into
deriving rows. A \texttt{report} mark, or no mark at all (unmarked
means \texttt{report}), would instead make the same sequent a checked
constraint, judged and never fired. The join stays identical; what
changes is who owes whom what. A materialized view comes out of the
query you wrote together with the engine's maintenance code, and it
contains whatever that code computed. A geolog view, by contrast, is a
logical consequence. The chase computes the \emph{least fixpoint} of
the axioms over the asserted rows: every row the axioms force, and no
row they do not. Anything that can replay the log and read the theory
can then re-derive and audit every row, with its provenance. In the
SQL world the nearest neighbours are Feldera and the DBSP line of work
\cite{budiu22}, and the kinship is genuine: declarative view
definitions, incremental maintenance with a correctness story, replay
from a log. What differs is the view language itself. A SQL view
computes what its query says; a geolog view is an entailment in
geometric logic \cite[Part~D1]{johnstone2002}, the fragment of
first-order logic whose sequents have a conjunction of atoms as
premise and a disjunction of existentially quantified conjunctions as
conclusion. That fragment is why the other three axioms above (an
existential, an equality, and a denial) are ordinary view
definitions here, while over there they are not expressible as views
at all. Those three carry the rest of this chapter.

Two caveats, up front. The first is maintenance. Feldera's engine is
incremental to its bones, and at production scale that maturity is
real (\S\ref{sec:gt-scorecard}). In this preview, geolog's chase
evaluation is a deliberately naive reference implementation,
differentially tested against an independent engine and built to be
right rather than fast; incremental maintenance ships today only for
its violation views, where the failing set is kept current under row
deltas and never recomputed from scratch. The second is expressiveness,
and it cuts the other way. Aggregates such as \texttt{SUM} and
windowed rates, the bread of streaming SQL, are not native vocabulary
in the geometric fragment. Counting enters through explicit
constructions or through the ground seam (\S\ref{sec:gt-ground}), and
where your computation is aggregate-shaped, today the SQL stack simply
wins.

\section{One writer at a time}
\label{sec:gt-single}

The two writers' logs live as two \texttt{instance} blocks in the
same file; in the durable pipeline each would instead be a stream of
signed commits (\S\ref{sec:gt-log}). The storefront saw payments and a
cancellation, and it holds the carrier's tracking record for order
\texttt{o1}. The warehouse allocated stock for both orders, the
cancellation not having reached it, and it holds its own record of the
same shipment, the printed label:

\begin{Verbatim}[fontsize=\small,frame=leftline,rulecolor=\color{accentmid},framesep=3mm,xleftmargin=3mm]
instance Storefront : Orders {
  o1, o2 : Order;
  sf_record : Shipment;

  [o: o1] Placed;  [o: o1] Paid;
  [o: o2] Placed;  [o: o2] Paid;
  [o: o2] Cancelled;

  [o: o1, s: sf_record] ShippedVia;
}

instance Warehouse : Orders {
  o1, o2 : Order;
  wh_label : Shipment;

  [o: o1] Allocated;
  [o: o2] Allocated;

  [o: o1, s: wh_label] ShippedVia;
}
\end{Verbatim}

Chase each log alone (\texttt{--instance} selects blocks):

\begin{Verbatim}[fontsize=\small,frame=leftline,rulecolor=\color{accentmid},framesep=3mm,xleftmargin=3mm]
$ geolog chase examples/orders.geolog --instance Storefront
  ⋮
  Shippable — 0 rows (0 seeded, 0 derived)
  ShippedVia — 1 row (1 seeded, 0 derived)
    [o: o1, s: sf_record] ShippedVia
fixpoint: reached, no judgements — the model is saturated.
$ echo $?
0
\end{Verbatim}

\noindent The warehouse's run is symmetric: no \texttt{Paid}
rows, so again no \texttt{Shippable}, and exit 0. Notice what clean
means here. \emph{Each writer's log is locally consistent.} The
storefront knows \texttt{o2} is cancelled but sees no shipment for it,
and the warehouse sees no cancellation at all. Every partial view is
innocent. If the two logs were Kafka partitions consumed by separate
services, each service would be reporting green.

\section{Merge is set union plus re-derivation}
\label{sec:gt-merge}

Now merge. Merging replicas in geolog takes no protocol. The commit
bag is a set, union is union, and the chase re-derives. Run without
\texttt{--instance} and the file's blocks seed as their union. Element
names shared across blocks denote the same element, the CLI's
small-scale stand-in for the content-addressed identity of
\S\ref{sec:gt-log}:

\begin{Verbatim}[fontsize=\small,frame=leftline,rulecolor=\color{accentmid},framesep=3mm,xleftmargin=3mm]
$ geolog chase examples/orders.geolog
sorts:
  Order — 2 elements (2 seeded)
  Shipment — 2 elements (2 seeded, 1 invented by the chase,
                         1 merged by congruence)
  merged: sf_record = wh_label
relations:
  ⋮
  Shippable — 2 rows (0 seeded, 2 derived)
  + [o: o1] Shippable
  + [o: o2] Shippable
  ShippedVia — 2 rows (1 seeded, 1 derived)
    [o: o1, s: sf_record=wh_label] ShippedVia
  + [o: o2, s: ?Shipment.0] ShippedVia
judgements (1):
  denial violated: ax/no/zombie at [o: o2, s: ?Shipment.0]
$ echo $?
1
\end{Verbatim}

Four things happened, and they are four fifths of this database's
distinctive behaviour, so take them one at a time.

\paragraph{The union derived what neither log could.}
\texttt{Shippable} was empty in both single-writer runs, and the
merged run derives it for both orders, because
\texttt{ax/ship/ready}'s premise spans the two logs. That much is the
ordinary magic of any log-merging system; a Feldera pipeline consuming
both topics would derive the same join. What it did not require was a
pipeline topology. The axiom was already the join, and \emph{any}
subset of commits chases to its own well-defined fixpoint.

\paragraph{The two writers' records became one shipment.}
The storefront's \texttt{sf\_record} and the warehouse's
\texttt{wh\_label} are two elements naming, as it turns out, one real
shipment. \texttt{ax/manifest/unique} says as much: a shipment is one
shipment. The engine's congruence closure \emph{identifies} the two
elements, and the report prints the merged element as
\texttt{sf\_record=wh\_label}. Contrast the CRDT discipline you may
run on Kafka-adjacent state stores. There, convergence is bought by
designing every value as a lattice and every merge as a join, per
type, by hand, in advance. Here identity is a \emph{logical} matter,
decided by the theory's equality axioms and computed by congruence,
and it composes with everything else the axioms say. It is also a mode
\emph{choice}, made where the axiom is written.
\texttt{ax/manifest/unique} is \texttt{chase}-marked
deliberately, because identification is this axiom's job;
\texttt{report}-marked, the same sequent would instead judge two
distinct records as a violation, which is the right reading in a
domain where identity is content and merging would invent facts. (When
two elements merge, everything known about either transfers to the one
survivor. Judging \texttt{ax/no/zombie} against the invented witness
below already relies on that.)

\paragraph{The chase invented a witness.}
Order \texttt{o2} became \texttt{Shippable}, and
\texttt{ax/ship/manifest} demands that a shipment exist for every
shippable order. No writer recorded one, so the chase minted
\texttt{?Shipment.0}, a fresh element about which nothing is known
except what the axiom forces. This is the \emph{restricted chase}:
existentials invent a witness only where no witness already exists
(order \texttt{o1}, which had records, got no invention). If you know
database theory you recognize tuple-generating dependencies and the
chase from the classical literature
\cite{fagin2005data,deutschnashremmel2008}; if you know event-sourced
systems, read it as the database itself filing the missing follow-up
event, typed, provenance-tracked, and only when entailed.

\paragraph{The violation surfaced only at the merge, and was
reported rather than resolved.}
Both logs were locally clean. Their union contains a cancelled order
with a shipment, and \texttt{ax/no/zombie} forbids precisely that. The
chase does not reject the merge, since there is nothing to reject; the
union of two sets of commits is a set of commits. Nor does it pick a
loser or roll anything back. It reports the judgment, with the
axiom's name and the witnessing binding, and exits nonzero. That
posture is the event-sourcing one, held formally: \emph{invalid states
are representable, detected, and repaired by further ordinary
commits} --- a compensating cancellation of the shipment, a voiding
of the cancellation, whichever the workflow's owner chooses. What the
database contributes is that ``invalid'' is a theorem rather than a
code path someone remembered to write, and that the judgment itself
is data an agent can act on. (Where an axiom's conclusion offers a
genuine \emph{choice} of repairs, a disjunctive head, the engine
likewise picks nothing, and the judgment carries the menu. The engine
calls these branch points; this example's judgments are denials,
menus of length zero.)

If you have fought a Kafka consumer-group rebalance that replayed a
partition into a half-updated store, or reconciled two services'
databases after a network partition, this section is the pitch. Merge
here is set union, and so it is \emph{always} defined, order-free, and
idempotent; and every consequence of the merged log, including its
violations, can be re-derived by anyone from the log alone.

\section{The log is the ground truth}
\label{sec:gt-log}

Everything so far ran from a source file, because this is a tutorial.
In geolog, ``the facts'' take durable form as a log of \emph{commits}:
content-addressed, ed25519-signed envelopes whose payloads carry just
the kind of rows you have been reading, and whose causal dependencies
are part of the signed content. Crossing from the file to that log is
one verb:

\begin{Verbatim}[fontsize=\small,frame=leftline,rulecolor=\color{accentmid},framesep=3mm,xleftmargin=3mm]
$ geolog seal examples/orders.geolog --store orders-store --repro
repro-signed (public seed "geolog-seal-repro-v1" — reproducibility, not authority)
sealed examples/orders.geolog: 1 theory, 4 axioms, 197 revision ops
revision commit:  b128db74d04f92bb0fdc5aec5a26f1d67b5504ba9289a7e888b487c1ee232aa4
C1 revision hash: c7de71b56edc728bc65daba559d3e3f1b0bf32bf2030689de4e1b69d64066f7f
seal commit:      e69f4611e367e91c8509731fc04f349ccf0b18beef0ea1f87788ce6351505a59
store: genesis minted, comparisons minted, revision minted, seal minted
\end{Verbatim}

\noindent The theory's declarations compile to meta-schema rows
(\S\ref{sec:gt-schema}), the rows seal into a signed revision commit,
and the store on disk now holds the envelope shapes the storage
crates document and test (\path{geolog-envelope}, \path{geolog-store},
and \path{geolog-sync}, each with its differential partner named in
its README). Sealing is deterministic: re-running the command reports
every commit \texttt{already present}, and \texttt{--repro} signs
with a public seed so that the hashes above reproduce from a fresh
clone (reproducibility, as the tool says, not authority).
\texttt{geolog check --store} and \texttt{geolog chase --store} then
run the strata you already know over the sealed source instead of the
file, after signature verification and a coherence gate; the
reference manual's CLI section (\S\ref{sec:gr-cli}) walks the whole
verb family, \texttt{extract} and \texttt{load} included.

The architecture invites a comparison. A Kafka topic with infinite
retention is also an append-only ground truth, and replay-rebuilt
views are also disposable. Three differences are worth having words
for.

\paragraph{The log forms a set rather than a sequence.} A topic is an
ordered sequence per partition, and that order carries real weight:
consumers coordinate offsets, rebalances are delicate, cross-partition
order is undefined. A geolog commit bag is unordered. Causality is
explicit, since each commit names its dependencies by hash, and
derived state does not depend on admission order. The repository's
battery includes a rebuild invariant that diffs
incrementally-maintained state against a from-scratch replay, and a
divergence is a release-blocking bug. ``Replays in any order to the
same fixpoint'' is the property that made \S\ref{sec:gt-merge}'s merge
a non-event.

\paragraph{Writers are cryptographic identities, and equivocation is
detectable.} A broker will happily append whatever a producer sends,
and if a compromised or buggy writer tells two consumers two different
stories, nothing in the topic model even states the offence. geolog
commits are signed and content-addressed, and a writer's log is a
hash-linked chain. A writer that forks its own history, which is to
say equivocates, is detected \emph{structurally} at sync time, from
the bytes alone, with no consensus protocol and no quorum; the same
mechanism makes CRDTs Byzantine fault tolerant in Kleppmann's
construction \cite{kleppmann22}. (The repository's envelope crate
ships the demonstration as a test: two signed commits claiming the
same chain position are the offence, exhibited and caught.)

\paragraph{Validity is per-commit and logical.} Because consequences
and violations re-derive from any subset of the log, ``which commits
does this judgment implicate'' is a well-posed question. The
provenance machinery exists in the engine, so that every derived row
and every invented element carries its firing history; and the
engine's validity layer (\path{geolog-chase/validity} in the
repository) makes the per-commit verdict itself durable: it replays a
commit's inclusive causal closure into a content-addressed snapshot
and folds one verdict per commit --- clean, judged with its
violations attached, or a named abstention (\texttt{pending},
\texttt{unsupported}, \texttt{indeterminate}), never a silent skip
--- with rebuild-order and cross-process determinism tested down to
byte-identical snapshot bytes, so two replicas that judge the same
commit publish the same judgment. Its first version replays a
deliberately narrow fragment (one theory revision per closure,
unmarked axioms only), and the tutorial CLI does not surface the
pipeline yet.
Capability-scoped authorization over the same substrate, meaning who
may extend which theory and assert into which relation, is likewise
arriving as an algebra over commits rather than as an ACL bolted on
beside them (\path{geolog-caps} ships the algebra today; enforcement
at the store boundary is design tense).

And the repository eats its own cooking: its committed theory
goldens are themselves stored as signed artefacts, minted by a
public, documented ceremony, and

\begin{Verbatim}[fontsize=\small,frame=leftline,rulecolor=\color{accentmid},framesep=3mm,xleftmargin=3mm]
cargo run -q -p geolog-ceremony -- verify --root theories
\end{Verbatim}

\noindent re-derives them byte-identically from their committed
listings and verifies the committed signatures, the golden-drift alarm
that runs in every CI pass.

\section{Schema is data}
\label{sec:gt-schema}

The Quick Start ended with \texttt{geolog rows}. The theory you write
is itself rows: \texttt{Theory}, \texttt{Sort}, \texttt{Fun}, and
\texttt{Field} elements in the \emph{meta-schema}, the schema language
that is a geolog theory like any other. In an event-sourced shop this
should land as the resolution of a familiar irritation, the schema
registry. Avro or protobuf schemas live \emph{beside} the topics, in a
registry with its own evolution rules (backward, forward, or full
compatibility), its own versioning, and its own failure modes: a
second database whose consistency with the first is everyone's job and
no one's.

In geolog there is no second database. A theory arrives as a
\emph{revision}, a sealed batch of meta-schema rows, committed,
signed, and content-addressed just like every other commit, with its
dependencies named. Extending a schema means a further revision that
depends on the first. Two independent extensions merge like any other
commits, by set union, with the chase over the meta-schema re-deriving
the combined signature. And ``what did the schema look like when this
row landed'' is a query over the same log as the row. Annotations
follow the same pattern instead of a special mechanism: the
repository's worked units-of-measure theory decorates an existing
theory's declarations with dimension gradings from the \emph{outside},
as ordinary relation rows over \texttt{Fun} and \texttt{Sort}
elements, without touching the annotated theory's bytes.

Here is what ships today in the alpha repository: the meta-schema
itself, the genesis golden of \S\ref{sec:gt-log}'s ceremony; the
reference row listings and their generators for the worked theories
under \path{theories/}; the theory-to-rows compiler and the
sealed-revision conformance loader, with byte-identical round-trips on
the committed corpora and loud rejections collected per rule; and the
revision-level machinery, namely declaration indexes, signature
hashes, and a programmatic revision-body builder. The \texttt{geolog
rows} verb you ran is a faithful \emph{view} of the declaration slice,
and \texttt{geolog seal} runs the same compilation for keeps: the
transcript of \S\ref{sec:gt-log} sealed the order theory as a signed
revision of 197 meta-schema operations. Cross-revision
\texttt{extends}, parameterized theories, and instances as
first-class elements are specified in the ruled design documents and
arrive as that store grows; the tutorial's single-file
\texttt{extends}, which the CLI does accept, is the in-file shadow of
that mechanism.

\section{Ground vocabulary: decided, not stored}
\label{sec:gt-ground}

Everything so far has been finite rows over declared sorts. Real
domains want numbers, timestamps, and comparisons, vocabulary whose
extension no log could store (\emph{every} pair of rationals with $a
\le b$?) and no chase should enumerate. geolog's answer is the
\emph{ground seam}. Relations and functions over immutable carriers
are declared as theory rows like everything else, but \emph{decided}
by registered procedures: exact arbitrary-precision rational
arithmetic, and calendar arithmetic whose semantics are documented and
tested against PostgreSQL's (\texttt{date} + \texttt{interval}
behaviour, chosen precisely so a SQL practitioner's instincts transfer
and the corner cases are somebody's published corner cases rather than
ours).

Two disciplines make the seam trustworthy rather than a trapdoor. The
first is \emph{refusal over guessing}. Ground vocabulary an
application declares but no provider implements does not default,
approximate, or silently empty; a plan that cites it cannot even be
expressed, and the refusal names the atom. (A division's graph simply
contains no row for a zero divisor, so ``divide by zero'' is
underivability rather than an exception.) The second is \emph{the lane
law}: an accelerated implementation may change \emph{cost}, never
\emph{verdicts}. Every fast lane is conformance-gated against the
always-on exact reference, and a lane that cannot pass its battery
does not dispatch. If your mental model of database extension points
is user-defined functions, hold on to this contrast: a UDF is trusted
code in the query path, whereas a ground provider is an untrusted
proposal bound to a declared signature, adjudicated by the engine,
with an exact reference underneath it.

(The tutorial's CLI bridge does not yet reach the ground seam.
\texttt{geolog chase} today covers declared sorts and
\texttt{Prop}-valued relations, and it says so loudly when a theory
wants more. The engine-level seam, the exact-rational floor, and the
calendar corpus are in the repository under test; the worked
Petri-net theory shows the discipline of doing real arithmetic
\emph{without} the seam, as chase-derived rows over an explicit number
ladder.)

\section{Worked example: alifib in geolog}
\label{sec:gt-alifib}

Everything in this chapter so far was designed to teach this chapter.
This section imports a structure that was not. \emph{alifib}
\cite{alifib} is an experimental programming language and interactive
proof assistant founded on higher-dimensional rewriting, developed by
Amar Hadzihasanovic together with Alex Kavvos within ARIA's
Safeguarded AI programme; its mathematics is Hadzihasanovic's theory
of directed complexes and molecules \cite{hadzihasanovic24a}, and its
data structures grew out of joint work with Diana Kessler
\cite{hadzihasanovic23}. In alifib a type is a directed cell complex;
a value is a \emph{diagram} drawn in that complex; and a computation
on $n$-dimensional data is itself a diagram, one dimension up.
Running a program builds a witness of the whole computation, a
higher-dimensional term you can inspect, paste into larger diagrams,
and reason about; the project calls the property \emph{computational
transparency}. Rewrite rules are ordinary generators one dimension up,
so there is no separate sort of rules anywhere in the language, and an
execution trace is a term like any other.

Two of this chapter's themes meet in it, and it forces a third; that
is why it earns a section. First, schema as data
(\S\ref{sec:gt-schema}) at full
stretch: the whole proof theory lands as one ordinary geolog theory,
seven sorts, thirty-eight functions and relations, twenty-four axioms,
with no extension to the meta-schema at all. Second, decided rather
than stored (\S\ref{sec:gt-ground}): every judgment the alifib engine
computes over its finite combinatorial data --- whether a paste is
legal, whether a boundary pair is round, whether an incidence table
matches a realized shape --- enters the theory as a side-condition
relation whose evaluator is the engine itself, the ground-seam
discipline pointed at a full external prover rather than at rational
arithmetic. Third, the forced one: the geometric fragment offers no
negation in premises, and a theory of cell complexes keeps wanting one sentence of that shape (``this cell has no
coface''). The stored-complement idiom below carries that sentence
without leaving the fragment.

The theory ships as \path{theories/Alifib.geolog} in the geolog-alpha
repository, and its opening \texttt{import} header is the
decided-versus-stored split of \S\ref{sec:gt-ground} written down:
the file imports its eight engine-decided judgments from the
\texttt{alifib} namespace (each import classifying that relation
\texttt{side-condition} --- decided by a bounded external procedure),
and every relation it declares without importing is
\texttt{ordinary}, the default. The check verb runs green on it:

\begin{Verbatim}[fontsize=\small,frame=leftline,rulecolor=\color{accentmid},framesep=3mm,xleftmargin=3mm]
$ geolog check theories/Alifib.geolog
[1/5] parse + normalize + lint: theories/Alifib.geolog @ 2e60dd6 (452 lines)
      1 theories; 24 axioms (0 moves; aux helpers: none)
[2/5] geometric linter: CLEAN
  ⋮
[5/5] instance conformance: no instance blocks in this file
modes: 6 chase · 18 report (18 by the unmarked default) · 0 refuse
  ⋮
CHECK: GREEN
$ echo $?
0
\end{Verbatim}

\noindent The excerpts below are the teaching spine, roughly a third
of the theory; the axioms not shown (execution traces as first-class
terms; roundness of each generator's boundary pair, the shape
condition on a single cell; the remaining three sequents of the
boundary-parity family; uniqueness of a cell's coordinate; homing of
boundary rows and map images) repeat the three patterns you are about
to see. The theory splits into two tiers, the syntax of assemblies
and the incidence data of the shapes they realize, and the walk takes
them in order. One caveat before it starts: the tutorial CLI's
\texttt{rows} and \texttt{chase} verbs do not reach this theory yet
--- its element-valued function symbols are engine IR the current
chase bridge does not carry --- and both refuse loudly, each naming
its workaround in the refusal. The transcripts here therefore ride
\texttt{check}, \texttt{census}, and the conversion tools, and the
excerpts themselves are kept faithful by construction: concatenated in
source order they form a well-formed theory that \texttt{geolog check}
accepts (the receipt at this chapter's source tail records that run).

\subsection{Paste trees as constructor rows}
\label{sec:gt-alifib-tier1}

\begin{Verbatim}[fontsize=\small,frame=leftline,rulecolor=\color{accentmid},framesep=3mm,xleftmargin=3mm]
theory Alifib {

  AType : Set;   // a type: a directed complex, elaborated
  Gen   : Set;   // one generator (cell) of a type; a rewrite
                 // rule is a Gen one dimension up
  Diag  : Set;   // a diagram AS ASSEMBLY -- a binary paste tree
  AMap  : Set;   // a map between types (generator-to-diagram images)
  Cplx  : Set;   // a materialized oriented graded poset
  Cell  : Set;   // one cell of a Cplx
  // (the seventh sort, Trace, records executions: a witness diagram
  //  one dimension up, with its endpoints)

  Gen/type : Gen -> AType;
  Gen/dim  : Gen -> Nat;
  Gen/src  : Gen -> Diag;   // PARTIAL: absent iff dim = 0 --
  Gen/tgt  : Gen -> Diag;   // a 0-cell has no boundaries

  dGen   : [g: Gen] -> Diag;                    // a single cell
  dPaste : [k: Nat, l: Diag, r: Diag] -> Diag;  // l #k r
\end{Verbatim}

alifib builds diagrams by \emph{pasting}: \texttt{l \#k r} glues two
diagrams along their shared $k$-dimensional boundary, and the engine
stores a diagram as its assembly history, a binary tree of pastes over
single-cell leaves. The two constructors above mirror that
representation without loss. Because \texttt{dPaste} is a function ---
one value per argument tuple, an equality the chase maintains by
merging --- two occurrences of the same tree denote one \texttt{Diag}
element, and structurally equal assemblies are shared: the engine's
own pointer-sharing, recovered as identity semantics. What a
\texttt{Diag} element is \emph{not} is a diagram up to isomorphism:
two trees can realize isomorphic diagrams and remain distinct
elements, related by a declared side-condition (\texttt{alifib\_iso},
in the full theory). That choice matches the engine, which stores
assemblies and \emph{computes} isomorphism, and it is the honest one
for a database: identity of stored data stays structural, and the
quotient stays a judgment.

Partiality needs no apparatus. A 0-cell has no boundaries, so
\texttt{Gen/src} simply has no row for it; there is no nullable column
and no sentinel: a boundary a generator lacks is a row nobody asserts,
and the axioms only ever demand what presence owes. The same
convention carries map holes below.

Everything the engine \emph{decides} about these trees enters as
side-condition relations, and the theory \emph{demands} the judgments
where they are owed:

\begin{Verbatim}[fontsize=\small,frame=leftline,rulecolor=\color{accentmid},framesep=3mm,xleftmargin=3mm]
  alifib_wf       : [d: Diag] -> Prop;   // realizes to a molecule
  alifib_pastable : [k: Nat, l: Diag, r: Diag] -> Prop;

  ax/wf/diag : forall d : Diag.
    |- [d: d] alifib_wf;

  ax/wf/paste : forall k : Nat, l : Diag, r : Diag, d : Diag.
    d = [k: k, l: l, r: r] dPaste
    |- [k: k, l: l, r: r] alifib_pastable;
\end{Verbatim}

\noindent Read \texttt{ax/wf/paste} as a well-formedness debt: every
asserted paste owes a pastability judgment, and the judgment is
discharged by evaluation, not by chase derivation --- each of these
relations is imported at the top of the file, which classifies it
\texttt{side-condition}, and the bounded external procedure behind
them is the alifib engine. The
refusal discipline of \S\ref{sec:gt-ground} carries over unchanged:
absent the engine, a plan citing these atoms cannot be expressed,
and nothing defaults to true.

One derived view shows the chase doing for proof theory what it did
for order flow. Which generators --- hence which rewrite rules ---
does a diagram use?

\begin{Verbatim}[fontsize=\small,frame=leftline,rulecolor=\color{accentmid},framesep=3mm,xleftmargin=3mm]
  Diag/uses : [d: Diag, g: Gen] -> Prop;

  chase ax/uses/gen : forall g : Gen, d : Diag.
    d = [g: g] dGen
    |- [d: d, g: g] Diag/uses;

  chase ax/uses/paste_l : forall k : Nat, l : Diag, r : Diag, d : Diag, g : Gen.
    d = [k: k, l: l, r: r] dPaste, [d: l, g: g] Diag/uses
    |- [d: d, g: g] Diag/uses;
\end{Verbatim}

\noindent (Its mirror \texttt{ax/uses/paste\_r} walks the right
subtree.) This is the \texttt{Shippable} of the proof theory: a table
nobody writes, entailed by structure, answering the question a proof
library wants indexed --- which rules does this proof depend on?

\subsection{Incidence, and maximality without negation}
\label{sec:gt-alifib-tier2}

A paste tree is a recipe. The shape it realizes is an \emph{oriented
graded poset}: cells graded by dimension, each cell covering its faces
with one of two signs, input or output. Chapter~\ref{ch:qtsl-scomplexes}
develops the S-complex vocabulary this book uses for such data and
notes that oriented graded posets embed as its Boolean instance; here
that instance appears as concrete rows. A materialized shape is a
\texttt{Cplx}; its cells carry a grade and an index within the grade
(the engine's canonical, normalization-derived order), and the two
signed covering relations land as two ordinary relations:

\begin{Verbatim}[fontsize=\small,frame=leftline,rulecolor=\color{accentmid},framesep=3mm,xleftmargin=3mm]
  Cell/cplx : Cell -> Cplx;
  Cell/dim  : Cell -> Nat;   // grade
  Cell/idx  : Cell -> Nat;   // index within the grade

  BdryIn  : [hi: Cell, lo: Cell] -> Prop;   // lo is an input face of hi
  BdryOut : [hi: Cell, lo: Cell] -> Prop;   // ... an output face

  ax/scx/sign_disjoint : forall hi : Cell, lo : Cell.
    [hi: hi, lo: lo] BdryIn, [hi: hi, lo: lo] BdryOut
    |- false;

  ax/scx/d2_pp : forall z : Cell, y : Cell, x : Cell.
    [hi: z, lo: y] BdryOut, [hi: y, lo: x] BdryOut
    |- exists w : Cell. [hi: z, lo: w] BdryOut, [hi: w, lo: x] BdryIn
    or exists v : Cell. [hi: z, lo: v] BdryIn, [hi: v, lo: x] BdryOut;
\end{Verbatim}

\noindent \texttt{ax/scx/sign\_disjoint} is a denial in the style of
\texttt{ax/no/zombie}: one sign per covering edge. The second axiom is
new territory. It is one quarter of the boundary-parity law, the
sign-aware face of $\partial \circ \partial = 0$: every two-step
signed path admits a two-step path of the opposite parity class
between the same endpoints. Its conclusion holds a genuine
disjunction, so the axiom is not chase-eligible --- there is no least
fixpoint to compute when the theory offers a choice --- and the engine
treats it exactly as \S\ref{sec:gt-merge} treated repairs: as a
judgment carrying a menu. Back there the menus had length zero; this
one has length two, and it earns its keep as a checked structural law
rather than a derivation rule (the reference manual's account of
per-axiom modes is \S\ref{sec:gr-chase-modes}). The law is not
hypothetical housekeeping: exporting all twenty-two example files
alifib ships, with included modules and every generator and named
diagram materialized, yields 37{,}988 complexes, and sign-disjointness
and the parity law hold across them with zero exceptions (receipt at
the chapter tail).

Now the idiom this section promised. ``$x$ is maximal'' means ``$x$
has no coface'', and that sentence wants negation in a premise, which
the geometric fragment deliberately lacks. Store the complement
instead:

\begin{Verbatim}[fontsize=\small,frame=leftline,rulecolor=\color{accentmid},framesep=3mm,xleftmargin=3mm]
  Cell/top     : Cell -> Prop;   // asserted: maximal (no coface)
  Cell/covered : Cell -> Prop;   // derived: has a coface

  chase ax/cell/covered_in : forall hi : Cell, lo : Cell.
    [hi: hi, lo: lo] BdryIn
    |- lo Cell/covered;

  chase ax/cell/covered_out : forall hi : Cell, lo : Cell.
    [hi: hi, lo: lo] BdryOut
    |- lo Cell/covered;

  ax/cell/top_not_covered : forall c : Cell.
    c Cell/top, c Cell/covered
    |- false;

  ax/cell/top_or_covered : forall c : Cell.
    |- c Cell/top or c Cell/covered;
\end{Verbatim}

\noindent Three moves. \emph{Generation}: \texttt{Cell/covered} is an
ordinary monotone view, derived from the boundary rows.
\emph{Exclusivity}: a denial keeps the asserted mark and the derived
view disjoint. \emph{Coverage}: a disjunctive law says every cell
falls on one side or the other --- judged, like the parity law, never
fired. Between them, ``maximal'' becomes queryable vocabulary, its
complement kept disjoint by the denial and exhaustive by the
disjunction, and no premise ever had to say \emph{not}. The pattern is
general; alifib merely needs it constantly, because boundaries of
pasting diagrams are defined through maximality.

\subsection{Maps, holes, and the axiom that must not fire}
\label{sec:gt-alifib-maps}

An alifib map assigns, to each generator of a domain type, a diagram
in a codomain type; alifib's proof workflow builds such maps
interactively, leaving \emph{holes} --- unassigned generators --- to be
filled as the proof progresses. In rows, a hole costs nothing:

\begin{Verbatim}[fontsize=\small,frame=leftline,rulecolor=\color{accentmid},framesep=3mm,xleftmargin=3mm]
  AMap/src : AMap -> AType;
  AMap/img : [g: Gen, m: AMap] -> Diag;   // PARTIAL: a hole is an
                                          // absent row
  AMap/total : AMap -> Prop;              // asserted mark

  ax/map/total : forall m : AMap, s : AType, g : Gen.
    m AMap/total, s = m AMap/src, s = g Gen/type
    |- exists i : Diag. i = [g: g, m: m] AMap/img;
}
\end{Verbatim}

\noindent Filling a hole is an ordinary append --- one new
\texttt{AMap/img} row --- so an interactive proof session is
event-sourced by construction, and ``what did this map look like
before the fill'' is a replay question like any other. The axiom is
the interesting resident. \texttt{ax/map/total} is chase-eligible on
its face: single disjunct, an existential the restricted chase knows
how to witness. Firing it would be a disaster in slow motion. The
chase would mint a fresh, contentless \texttt{Diag} for every open
hole of every map marked total, ``satisfying'' totality by inventing
images that mean nothing --- when what totality means for user data is
a \emph{judgment} that the user's own fills eventually discharge.
So the axiom stays report-marked, a checked debt rather than a
generator (one of the eighteen \texttt{report}s in the check
transcript above), and the invented-witness power of
\S\ref{sec:gt-merge} stays
where it belongs: on axioms whose existentials denote obligations the
database itself owns, like \texttt{ax/ship/manifest}'s missing
shipment. Choosing per axiom, in the source, with the default on the
safe side, is what the mode marks are \emph{for}.

\subsection{One diagram, end to end}
\label{sec:gt-alifib-e2e}

Here is a complete alifib type, small enough to watch cross the seam
whole. Two composable arrows, a third arrow, and one 2-cell rewriting
the pasted path into it (juxtaposition is alifib's shorthand for
pasting along the shared boundary of highest shared dimension; here
\texttt{f g} abbreviates \texttt{f \#0 g}):

\begin{Verbatim}[fontsize=\small,frame=leftline,rulecolor=\color{accentmid},framesep=3mm,xleftmargin=3mm,samepage=true]
@Type
Compose <<= {
    x, y, z,
    f: x -> y,
    g: y -> z,
    h: x -> z,
    step: f g -> h
}
\end{Verbatim}

\begin{figure}[htbp]
\centering
\begin{tikzpicture}[>=Stealth, baseline=(y.base)]
  \node (x) at (0,0) {$x$};
  \node (y) at (1.6,1.0) {$y$};
  \node (z) at (3.2,0) {$z$};
  \draw[->] (x) -- node[above left=-2pt] {$f$} (y);
  \draw[->] (y) -- node[above right=-2pt] {$g$} (z);
  \draw[->] (x) to[bend right=40] node[below] {$h$} (z);
  \node at (1.6,0.18) {$\Downarrow\;\mathit{step}$};
\end{tikzpicture}
\hspace{14mm}
\begin{tikzpicture}[baseline=(cy.base)]
  \node (cx) at (0,0) {$x$};
  \node (cy) at (1.3,0) {$y$};
  \node (cz) at (2.6,0) {$z$};
  \node (cf) at (0,1.15) {$f$};
  \node (cg) at (1.3,1.15) {$g$};
  \node (ch) at (2.6,1.15) {$h$};
  \node (cs) at (1.3,2.3) {$\mathit{step}$};
  \draw (cf) -- (cx);
  \draw[densely dashed] (cf) -- (cy);
  \draw (cg) -- (cy);
  \draw[densely dashed] (cg) -- (cz);
  \draw (ch) to[bend left=16] (cx);
  \draw[densely dashed] (ch) -- (cz);
  \draw (cs) -- (cf);
  \draw (cs) -- (cg);
  \draw[densely dashed] (cs) -- (ch);
\end{tikzpicture}
\caption{The 2-cell \texttt{step} of the \texttt{Compose} type, twice
over. Left, the pasting diagram: \texttt{step} rewrites the pasted
path \texttt{f \#0 g} into \texttt{h}. Right, the shape it realizes,
an oriented graded poset drawn as a Hasse diagram: solid covering
edges are input faces ($\partial^-$, the \texttt{BdryIn} rows of
\S\ref{sec:gt-alifib-tier2}), dashed ones output faces ($\partial^+$,
\texttt{BdryOut}). Seven cells in grades of sizes $(3,3,1)$, nine
signed edges --- row for row, the complex the transcripts below carry.}
\label{fig:gt-alifib-step}
\end{figure}

Save the listing as \texttt{compose.ali} at the repository root. The
export bridge (built once with \texttt{cargo build --release} in
\path{vendor/alifib}; \S\ref{sec:gt-alifib-convert} gives it a proper
tour) elaborates the file with the engine's own loader and interpreter
--- boundaries checked, shapes realized --- and \texttt{--check} stops
there, printing a summary. The first line is the file's module root, a
type holding nothing here; the one map on the second line is the
identity self-map the engine mints for every type, which
\S\ref{sec:gt-alifib-convert} returns to:

\begin{Verbatim}[fontsize=\small,frame=leftline,rulecolor=\color{accentmid},framesep=3mm,xleftmargin=3mm]
$ vendor/alifib/target/release/alifib-export compose.ali --check
compose.ali: type  — 0 generators, 0 diagrams, 0 maps (0 holes), 0 complexes
compose.ali: type Compose — 7 generators, 0 diagrams, 1 maps (0 holes),
                            7 complexes
warnings: 0
\end{Verbatim}

\noindent Ill-formed input stops at this gate with the engine's own
diagnostics: declare \texttt{oops:\ f f -> f} (pasting \texttt{f}
onto itself along a boundary the types do not share) and the export
exits nonzero with ``\texttt{Failed to paste diagrams: boundaries do
not match}'', pointing at the offending span. Nothing unchecked
reaches the rows.

Dropping \texttt{--check} emits the interchange document
(\S\ref{sec:gt-alifib-convert} specifies it):

\begin{Verbatim}[fontsize=\small,frame=leftline,rulecolor=\color{accentmid},framesep=3mm,xleftmargin=3mm]
$ vendor/alifib/target/release/alifib-export compose.ali -o compose.json
wrote compose.json (7616 bytes)
\end{Verbatim}

\noindent Two slices of it, whitespace refolded to the book's measure
and content bytes as emitted --- the \texttt{step} generator, whose
source is the paste tree itself:

\begin{Verbatim}[fontsize=\small,frame=leftline,rulecolor=\color{accentmid},framesep=3mm,xleftmargin=3mm]
{ "name": "step",
  "dim": 2,
  "src": { "paste": 0, "l": { "gen": "f" }, "r": { "gen": "g" } },
  "tgt": { "gen": "h" } }
\end{Verbatim}

\noindent and the materialized shape of \texttt{step}, the right-hand
panel of Figure~\ref{fig:gt-alifib-step} as data (grades bottom-up;
\texttt{faces\_in[d-1][p]} lists the input faces, at grade $d{-}1$, of
the $p$-th cell at grade $d$, as indices into the canonical order):

\begin{Verbatim}[fontsize=\small,frame=leftline,rulecolor=\color{accentmid},framesep=3mm,xleftmargin=3mm]
{ "of": { "kind": "generator", "name": "step" },
  "sizes": [3, 3, 1],
  "faces_in":  [ [[0], [1], [0]],  [[0, 1]] ],
  "faces_out": [ [[1], [2], [2]],  [[2]]    ],
  "labels": [ ["x", "y", "z"], ["f", "g", "h"], ["step"] ] }
\end{Verbatim}

The import half turns the document into instance rows against the
Alifib theory. Watch the paste tree become constructor rows, with
hash-consing visible (\texttt{f} and \texttt{g} each occur once as
leaves, then the paste refers back to them):

\begin{Verbatim}[fontsize=\small,frame=leftline,rulecolor=\color{accentmid},framesep=3mm,xleftmargin=3mm]
$ python3 theories/tools/alifib_interop.py import compose.json \
    -o compose.rows.txt
wrote compose.rows.txt (221 lines)
\end{Verbatim}

\noindent Operations 44 through 50 of that listing:

\begin{Verbatim}[fontsize=\small,frame=leftline,rulecolor=\color{accentmid},framesep=3mm,xleftmargin=3mm]
op[ 44] elta sort=G<alifib/Diag>                        // Diag: f
op[ 45] funa fun=G<alifib/dGen> args=[L16] value=L44    // dGen f
op[ 46] elta sort=G<alifib/Diag>                        // Diag: g
op[ 47] funa fun=G<alifib/dGen> args=[L20] value=L46    // dGen g
op[ 48] elta sort=G<alifib/Diag>                        // Diag: (f #0 g)
op[ 49] funa fun=G<alifib/dPaste> args=[nat:0, L44, L46] value=L48
op[ 50] funa fun=G<alifib/Gen/src> args=[L28] value=L48 // src of step
\end{Verbatim}

\noindent and, at the listing's tail, the last three of the figure's
nine signed edges --- the top cell's own --- as boundary rows
(\texttt{L196} names the \texttt{step} cell; \texttt{L181},
\texttt{L186}, \texttt{L191} the \texttt{f}, \texttt{g}, \texttt{h}
cells):

\begin{Verbatim}[fontsize=\small,frame=leftline,rulecolor=\color{accentmid},framesep=3mm,xleftmargin=3mm]
op[207] rela rel=G<alifib/BdryIn> args=[L196, L181]
op[208] rela rel=G<alifib/BdryIn> args=[L196, L186]
op[209] rela rel=G<alifib/BdryOut> args=[L196, L191]
\end{Verbatim}

\noindent These are ordinary instance rows; engine instruments consume
them as such. The census verb, for one, seeds them under the theory
and analyses which function symbols' merges the derivation leans on
(a dry-run instrument, it desugars the function symbols the chase
bridge refuses --- its first line counts them):

\begin{Verbatim}[fontsize=\small,frame=leftline,rulecolor=\color{accentmid},framesep=3mm,xleftmargin=3mm]
$ geolog census theories/Alifib.geolog --rows compose.rows.txt
merge-reliance census — 1 file(s), 1 rows listing(s), 24 arrow(s) desugared
  ⋮
seeds: 44 elements, 166 rows
\end{Verbatim}

Finally, back out. The projection inverts the import, and the
\texttt{.ali} renderer turns the interchange back into source, with
every paste parenthesized so the engine re-reads the same binary
trees:

\begin{Verbatim}[fontsize=\small,frame=leftline,rulecolor=\color{accentmid},framesep=3mm,xleftmargin=3mm]
$ python3 theories/tools/alifib_interop.py project compose.rows.txt \
    -o back.json
$ python3 theories/tools/alifib_interop.py to-ali back.json -o back.ali
$ sed -n 5,14p back.ali
@Type
Compose <<= {
    x,
    y,
    z,
    f: x -> y,
    g: y -> z,
    h: x -> z,
    step: (f #0 g) -> h
}
$ python3 theories/tools/alifib_interop.py accept compose.ali
ACCEPT: CLEAN
\end{Verbatim}

\noindent The \texttt{accept} verb at the end is the round trip run as
a test, both directions at once; the next subsection says what it
checks.

\subsection{Converting data to and from alifib}
\label{sec:gt-alifib-convert}

This subsection is written for readers who keep data in alifib's own
formats --- its maintainers first among them --- and want the crossing
stated precisely: where the bridge lives, what travels over it, the
commands for each direction, what survives a round trip, and what is
checked on the way in. The reward for crossing is everything this
chapter has shown: types, diagrams, proofs, and their realized shapes
as rows in a substrate with signed content-addressed commits, merge by
set union, derived provenance views like \texttt{Diag/uses}, and a
schema that is data in the same log.

\paragraph{The bridge, and why it lives where it does.}
alifib's structural core (paste trees, oriented graded posets) is a
private module, and the engine's own output surfaces render for
display: boundary expressions as strings, paste chains flattened, no
incidence tables. A faithful exporter therefore has to live
\emph{inside} the engine crate, where crate-private data is reachable,
and drive the interpreter itself; anything else would be a
reimplementation, trusted precisely where it should not be. The
geolog-alpha repository vendors a pinned copy of the engine carrying
that exporter as one additive module
(\path{vendor/alifib}; provenance, upstream commit, and the complete
patch list are in \path{vendor/alifib/VENDOR.md}, and the patch is
written to be upstreamable as-is). The bridge builds from a fresh
clone:

\begin{Verbatim}[fontsize=\small,frame=leftline,rulecolor=\color{accentmid},framesep=3mm,xleftmargin=3mm]
$ cd vendor/alifib && cargo build --release
$ ./target/release/alifib-export --help
usage: alifib-export <file.ali> [--type NAME] [--all-modules]
                     [--complexes none|gens|named|all] [-o OUT] [--check]
\end{Verbatim}

\noindent By default the root file's module is exported with every
named type, and shapes are materialized for every generator and named
diagram; \texttt{--complexes none} keeps the export purely syntactic,
\texttt{--type} restricts to one type, \texttt{--all-modules} includes
what the file \texttt{include}s. Output is deterministic --- two runs
are byte-identical --- because every collection rides a canonical
order: generators ascending by (dimension, name), which is also a
dependency-safe declaration order, since a generator's boundary leaves
are strictly lower-dimensional; diagrams and maps by name; cells in
the engine's own normalization order.

\paragraph{The interchange document.} One JSON document per export,
versioned by its first key. In sketch (field names exact, values
elided):

\begin{Verbatim}[fontsize=\small,frame=leftline,rulecolor=\color{accentmid},framesep=3mm,xleftmargin=3mm]
{ "alifib_interop": 1,
  "generator": "alifib-export (geolog interop bridge)",
  "source": "<root file basename>",
  "modules": [
    { "path": "<path relative to the root file>",
      "types": [
        { "name": "...",
          "generators": [ { "name", "dim", "src"?, "tgt"? } ],
          "diagrams":   [ { "name", "dim", "tree" } ],
          "maps":       [ { "name", "domain", "images", "holes" } ],
          "complexes":  [ { "of", "sizes",
                            "faces_in", "faces_out", "labels" } ] } ] } ],
  "warnings": [ "..." ] }
\end{Verbatim}

\noindent A generator's \texttt{src}/\texttt{tgt} are present iff its
dimension is positive. \texttt{diagrams} carries the type's named
diagrams: \texttt{let} bindings and proofs stored from interactive
sessions. A map's \texttt{domain} is either
\texttt{\{"type":\ t\}} or \texttt{\{"module":\ m\}}; its
\texttt{images} are filled assignments, its \texttt{holes} the open
ones, each hole optionally carrying a conditional image and
metavariable boundary trees. A complex records whose shape it is
(\texttt{of}), its per-grade extents (\texttt{sizes}), both signed
face lists indexed as in \S\ref{sec:gt-alifib-e2e}, and per-cell
generator labels. Boundary terms and diagram trees share one grammar,
four node shapes in total:

\begin{Verbatim}[fontsize=\small,frame=leftline,rulecolor=\color{accentmid},framesep=3mm,xleftmargin=3mm]
{"gen": "f"}                      a single-cell diagram
{"paste": k, "l": t1, "r": t2}    t1 #k t2 -- the engine's own binary
                                  node, never flattened or re-associated
{"hole": "g"}                     a map hole's metavariable (appears
                                  only in hole boundary terms)
{"labels": [...], "ldim": n}      loud fallback for a diagram with no
                                  paste history at its top dimension;
                                  counted in warnings, not a term
\end{Verbatim}

\paragraph{Into geolog.} \texttt{import} compiles a dump to instance
rows; \texttt{--type} restricts, as with the exporter:

\begin{Verbatim}[fontsize=\small,frame=leftline,rulecolor=\color{accentmid},framesep=3mm,xleftmargin=3mm]
$ python3 theories/tools/alifib_interop.py import mine.json -o mine.rows.txt
\end{Verbatim}

\noindent The output is a plain-text listing of numbered operations in
the meta-schema's instance vocabulary, three kinds in all:
\texttt{elta} mints an element of a sort, \texttt{funa} asserts a
function row, \texttt{rela} a relation row. Arguments are literals
(\texttt{nat:0}, \texttt{str:"step"}) or backreferences
(\texttt{L48} names the element minted by operation 48; forward
references are ill-formed by construction). The
\texttt{G<alifib/\ldots>} tokens cite the Alifib theory's declaration
identities, so the listing is meaningful against any store holding
that revision. Each distinct paste tree becomes one \texttt{Diag}
element per type, however often it recurs --- the hash-consing of
\S\ref{sec:gt-alifib-tier1}, performed at import time. The importer
runs a structural self-check before writing (every reference points
backward and at an element; every external token lies within the
Alifib signature) and stamps the verdict in the header; the committed
listings under
\path{theories/generated/} are regeneration-checked against their
generators by the repository's check script, so drift between tool and
artifact is loud.

What admission means, in layers: the importer's self-check is the
shallow layer; the theory's demand axioms are the deep one. Every
imported element owes the judgments of
\S\ref{sec:gt-alifib-tier1} --- well-formedness, pastability,
roundness --- and those side-conditions are discharged by the engine,
which is to say: by re-elaboration. Running
\texttt{alifib-export mine.ali --check} \emph{is} the decision
procedure, executed on the way out of alifib, before the rows exist;
the incidence laws (sign-disjointness and parity) are additionally
checkable without the engine, and
\texttt{alifib\_interop.py verify-scx mine.json} does so directly on a
dump. Engine instruments consume the rows today
(\texttt{geolog check} for theory-side conformance,
\texttt{geolog census} for merge-reliance analysis, the
meta-rows loader in the repository's test suite for loader
conformance), and the one-verb facade of \S\ref{sec:gt-log} seals
theory \emph{sources}; an admission seam that seals imported instance
rows like these behind the same store is the piece still in design
tense, and this paragraph sheds the qualifier when it lands.

\paragraph{Out of geolog.} Two verbs invert the import:

\begin{Verbatim}[fontsize=\small,frame=leftline,rulecolor=\color{accentmid},framesep=3mm,xleftmargin=3mm]
$ python3 theories/tools/alifib_interop.py project mine.rows.txt -o back.json
$ python3 theories/tools/alifib_interop.py to-ali back.json -o back.ali
\end{Verbatim}

\noindent \texttt{project} rebuilds an interchange document from rows.
It is deterministic, total on the import's image, and loud off it:
unknown functions or sorts, dangling references, a \texttt{Diag} with
no constructor row, a map marked total with absent images (inside the
database that state is a judged, standing debt; the exporter never
emits it, so the projection holds it off the image), an index at or
beyond its grade's extent, a count mismatch between declared extents
and cells, non-adjacent boundary rows --- each is a named rejection,
not a silent drop. \texttt{to-ali} renders an interchange document as
\texttt{.ali} source: every paste fully parenthesized, so the engine
re-reads the same binary trees; identifiers outside alifib's lexical
grammar sanitized deterministically, with the name map recorded in the
header comment; maps emitted as module-level \texttt{let total}
declarations; and the identity self-map the engine generates for every
type is skipped, since elaboration regenerates it. The acceptance
judge for this direction is alifib itself: the round-trip witness
re-elaborates the rendered source with the engine's own parser and
loader and compares the re-export against the original document.

\paragraph{What survives, what degrades.} Surviving a full round trip,
per type, structurally intact: generator names, dimensions, and
boundary trees as exact binary pastes (no re-association); named
diagrams and their assemblies; map images and holes, as assignments
over the domain's generators; complexes cell for cell --- extents,
both signed face lists, labels. Degrading, each by design and each
documented in the tools' headers: map \emph{names} are not carried by
rows, so maps re-emerge under fresh names; a conditional image (one
recorded on a hole but awaiting boundary resolution) degrades to a
plain hole, and the engine re-derives the constraint on
re-elaboration; maps
whose domain is a whole module are skipped with a note, the module
system being surface elaboration that the interpreter flattens before
the proof theory begins; display metadata (\texttt{index} blocks) does
not travel; execution traces sit outside interchange version~1 ---
the \texttt{Trace} sort of \S\ref{sec:gt-alifib-tier1} is the landing
site reserved for them.

\paragraph{The acceptance harness.} Both witnesses run as one verb per
file (\texttt{accept}, as in \S\ref{sec:gt-alifib-e2e}), and over a
fixed ten-file corpus spanning the shipped examples' feature surface
--- abstract machines, higher-categorical developments, hole
workflows, simplicial complexes; the parity measurement of
\S\ref{sec:gt-alifib-tier2} separately covered all twenty-two files:

\begin{Verbatim}[fontsize=\small,frame=leftline,rulecolor=\color{accentmid},framesep=3mm,xleftmargin=3mm]
$ python3 theories/tools/alifib_interop.py accept-corpus
BinaryNat.ali: CLEAN
EckmannHilton.ali: CLEAN
  ⋮
Hole_examples.ali: CLEAN
LambdaSigma.ali: CLEAN
CORPUS: 10/10 clean
\end{Verbatim}

\noindent Witness one: export, import, project; the projection must
reproduce the interchange up to the documented degradations. Witness
two: export, render to \texttt{.ali}, re-elaborate with the engine,
re-export; the second document must match the first up to the recorded
name map --- boundary trees, hole sets, and complexes cell for cell.
Every round-trip claim this subsection has made is one of those
comparisons or one of the named rejections. Readers extending the
interchange --- traces are the obvious next passenger --- extend the
witnesses first; the harness is the contract.

\section{The scorecard}
\label{sec:gt-scorecard}

Where the SQL stack wins today, plainly: \textbf{Throughput and
maturity.} Kafka moves gigabytes per second and Feldera maintains
complex view cascades incrementally, while this preview's chase is a
correct, differentially-tested reference implementation whose
incrementality so far covers violation views only. \textbf{Aggregates
and windows.} Streaming SQL's core vocabulary, not native to the
geometric fragment (explicit constructions or the ground seam, at
tutorial scale). \textbf{Operations.} The surrounding decade of
connectors, monitoring, and hiring.

Where geolog wins, plainly: \textbf{Views you can audit.} Derived
state is entailment, re-derivable and checkable from the log by a
third party, with provenance; there is no fold code to believe.
\textbf{Merge without a protocol.} Set union plus re-derivation,
order-free, defined for every pair of replicas, with no per-type merge
design and no rebalance to choreograph. \textbf{Judgments over
exceptions.} Invalid states and pending choices are typed, reportable
data with the violated axiom attached, rather than code paths.
\textbf{Writers you can hold to account.} Signatures,
content-addressing, and structural equivocation detection are the data
model, not an add-on. \textbf{One substrate.} Schema, extension,
annotation, and migration are rows in the same log they govern; there
is no registry to drift.

Compose the two lists and the position comes out plain. At this
preview, geolog has the right shape for systems of record where
auditability, principled merge, and schema-as-data carry the value, at
volumes where a reference-grade engine suffices. It is also built
deliberately so that the fast paths can arrive without the trust story
changing (the lane law; the kernels repository).

\section{Composing systems, and where QTSL begins}
\label{sec:gt-onward}

The order flow of this chapter is one system. The book's larger
programme composes systems, with shared boundaries, wiring diagrams,
and the population and fusion of components, and it states
\emph{temporal} claims about the composites. Those claims are not of
the form ``no cancelled order has shipped''; that is a state property,
and you watched a denial say it. They look instead like ``every placed
order is shipped or refunded within thirty days, with probability at
least $p$.'' Those are claims the chase alone cannot check and the
event log alone cannot even state. Stating them is QTSL's job, and
certifying them against models is the rest of this book: the systems
vocabulary in chapter~\ref{ch:qtsl-mathematical-models}, the
composition patterns in chapter~\ref{ch:qtsl-composition-patterns},
and the zero-to-certificate path in chapter~\ref{ch:qtsl-quickstart}.
You now hold the substrate they all stand on: a log whose truth is
structural, tables that are theorems, and a schema that is data.

 \chapter{Reference Manual}
\label{ch:geolog-reference}

\noindent\draftmark\medskip

This chapter is the systematic account of geolog as shipped in the
\path{geolog-alpha} repository: the source language and its grammar,
the normalization pipeline and its registry, the elaboration fragment
the \texttt{geolog} verb carries, the chase and its judgments, the
meta-schema, and the substrate contracts — commits, storage, sync,
the ground seam, and capabilities. The Quick Start
(chapter~\ref{ch:geolog-quickstart}) and the Tutorial
(chapter~\ref{ch:geolog-tutorial}) teach by example; this chapter
specifies. It restates none of their walkthroughs and cites them
where an example earns more than a rule would.

Three reading conventions govern everything below.

\paragraph{Claim discipline.} Normative present tense describes
behaviour shipped and tested in the repository as of this edition's
receipt passes (2026-07-21 through 2026-07-31, recorded at this
chapter's source tail). Designed-but-not-shipped behaviour appears only in
clauses explicitly marked \emph{design tense}, each naming where the
design lives today. Every quoted error message and every transcript
fragment in this chapter is the verbatim output of a real run against
the committed tree (receipts in this chapter's source tail).

\paragraph{Authority order.} The \texttt{geolog} binary is
ergonomics and composition over constituent crates, and its own crate
documentation states the rule this manual inherits: where a behaviour
differs between the facade and a constituent, the constituent is
right and the facade has a bug. Accordingly, each section below names
its authority — the crate (or committed theory file) whose documented,
tested behaviour it restates. Where this manual and that source
disagree, the source wins and the disagreement is a defect in this
manual.

\paragraph{Scope.} This manual covers the public
\path{geolog-alpha} tree. The engine library supports strictly more
than the command-line fragment (\S\ref{sec:gr-engine}); the manual
covers both, and says which surface each statement binds.

\section{Lexical structure}
\label{sec:gr-lexical}

\emph{Authority: \path{crates/geolog-normalize/src/lex.rs}.}

Theory text and instance blocks share one tokenizer. Input is
processed line by line; a \texttt{//} begins a comment (outside a
string literal) that extends to the end of the line and is
discarded. Tokens, longest match first:

\begin{Verbatim}[fontsize=\small,frame=leftline,rulecolor=\color{accentmid},framesep=3mm,xleftmargin=3mm]
two-char    ::= "->" | "|-" | "<=" | ">=" | "!="
single-char ::= "<" | ">" | "=" | "+" | "*" | "[" | "]" | "(" | ")"
              | "," | ":" | "." | "{" | "}" | ";"
integer     ::= digit+
string-lit  ::= '"' (any scalar except '"' or newline)* '"'
identifier  ::= starter continuer*
\end{Verbatim}

\noindent The identifier alphabet is Unicode-open, and the operator
characters are a closed ASCII set, so the two are disjoint by
alphabet rather than by any gating rule. A \emph{starter} is an
ASCII letter, \texttt{\_}, a non-ASCII alphabetic scalar (Greek
letters, $\ell$, \ldots), or one of the math-letterlike symbols
$\partial$ and $\nabla$; a \emph{continuer} is any starter, an ASCII
digit, \texttt{\_}, \texttt{/} (so \texttt{ax/path/edge} and
\texttt{Marking/count} are single tokens), or a script extender ---
the superscript and subscript digits and signs, and the primes.
Script extenders continue tokens and never start them: $\partial^{+}$
is one token, and a detached superscript at token start is a
tokenization refusal, so $\partial$~$^{+}$ (with a space between)
cannot silently mean anything. String literals are single-line and carry no escapes; the
grammar admits them in exactly one position --- contraction specs in
\texttt{law} bodies (\S\ref{sec:gr-editions}) --- and every other
position refuses them downstream. There are no negative-number
literals and no block comments in the theory grammar. Anything else
stops the tokenizer with a hard error naming the offending text.

The reserved words of the theory grammar are:

\begin{Verbatim}[fontsize=\small,frame=leftline,rulecolor=\color{accentmid},framesep=3mm,xleftmargin=3mm]
forall  exists  not  iff  or  false  theory  extends
Set  Prop  Index  Fin  DynTensor  List
\end{Verbatim}

\noindent \texttt{Fin} and \texttt{Index} name the same bounded-index
type (\texttt{Fin} is the ratified spelling, \texttt{Index} the
legacy alias --- \S\ref{sec:gr-decls}). \texttt{Tensor} is not
reserved, and no longer names a surface type at all: the
shape-in-type tensor former lives in the meta-schema tier
(\S\ref{sec:gr-meta-tier}), and the surface tensor type is the
keyword \texttt{DynTensor} (\S\ref{sec:gr-decls}). \texttt{instance}
is likewise not a keyword of the theory grammar: an
\texttt{instance} \emph{block} is recognized by the file splitter
(\S\ref{sec:gr-files}), one layer up, and the same word doubles as
an optional postfix marker on model-of types
(\S\ref{sec:gr-decls}), recognized contextually. The mode marks
\texttt{chase}, \texttt{report}, and \texttt{refuse}
(\S\ref{sec:gr-chase-modes}) are contextual in the same way ---
recognized only as a statement's leading token: \texttt{chase : …;}
is still an ordinary declaration named \texttt{chase}.

\section{Source files}
\label{sec:gr-files}

\emph{Authority: \path{crates/geolog/src/input.rs}.}

A \path{.geolog} file carries two strata: \emph{theory text} and
\emph{instance blocks}. Before parsing, the file is split by a
deliberately minimal scanner that tracks \texttt{//} comments and
brace depth, and recognizes the word \texttt{instance} only at brace
depth zero, at an identifier boundary on both sides. Each
\texttt{instance … \{ … \}} span is extracted whole (an unclosed
block is a hard error); the remaining text — with each extracted span
replaced by whitespace of identical line structure, so diagnostics
keep their line numbers — is handed to the normalization pipeline
verbatim. Consequences worth stating: the word \texttt{instance}
inside a comment or inside any braced body is ordinary text; theory
declarations and instance blocks may interleave freely at top level;
and the theory grammar remains entirely the normalizer's business —
the splitter never interprets theory text.

\section{Theories and declarations}
\label{sec:gr-decls}

\emph{Authority: \path{crates/geolog-normalize/src/parse.rs} (grammar)
and \path{crates/geolog-normalize/src/ast.rs} (abstract syntax).}

Theory text is a sequence of theory declarations:

\begin{Verbatim}[fontsize=\small,frame=leftline,rulecolor=\color{accentmid},framesep=3mm,xleftmargin=3mm]
theory-decl ::= "theory" NAME param* ["extends" NAME] "{" statement* "}"
param       ::= "(" NAME ":" ("Set" | NAME) ")"
statement   ::= [mode] NAME ("," NAME)* ":" statement-body ";"
mode        ::= "chase" | "report" | "refuse"
\end{Verbatim}

\noindent A theory header may bind parameters, name first: a
\texttt{(X : Set)} group binds a sort parameter, and a
\texttt{(N : PetriNet)} group --- whose second name must be a theory
declared earlier in the file --- binds a \emph{model-of} parameter,
a model of that theory. (An earlier binders-before-name header order
parses as a deprecated spelling; see \texttt{--strict},
\S\ref{sec:gr-cli}.) Every statement is an optional mode mark, one
or more comma-separated names sharing one signature, a colon, a
body, and a semicolon. The statement's \emph{kind} is dispatched on
its name: \textbf{a statement whose name begins with \texttt{ax/} is
an axiom} (\S\ref{sec:gr-axioms}); a statement opened by the
contextual keyword \texttt{law} is a law member of the aggregation
surface, whose grammar this edition defers
(\S\ref{sec:gr-editions}); every other statement is a declaration.
The prefix does grammatical work rather than expressing
a naming convention: a sequent-shaped body under a name without the
\texttt{ax/} prefix is parsed as a declaration and fails (verified:
\texttt{my/rule : forall …} stops with \texttt{syntax error:
unparseable type at line 4: forall …}, the cited line being the
statement's own). The Quick Start's remark that
slash-names are ordinary identifiers remains true — the slash is an
identifier character (\S\ref{sec:gr-lexical}) — but the two-character
prefix \texttt{ax/} is what makes a statement an axiom. The mode
mark decides whether an axiom fires or judges
(\S\ref{sec:gr-chase-modes}); the AST carries the \emph{authored}
mark, so rendering preserves the source form, and a mode mark on a
declaration head is a hard error naming both the mark and the head
(the message names \texttt{ax/} and \texttt{law} as the two homes a
mode can mark).

Declaration bodies:

\begin{Verbatim}[fontsize=\small,frame=leftline,rulecolor=\color{accentmid},framesep=3mm,xleftmargin=3mm]
sort-decl   ::= NAME ":" "Set" ";"
record-decl ::= NAME ":" "{" field ("," field)* "}" "->" codomain ";"
unary-decl  ::= NAME ":" type "->" codomain ";"
curried-decl::= NAME ":" ("(" NAME ":" type ")" "->")+ codomain ";"
const-decl  ::= NAME ":" type ";"
field       ::= NAME ":" type
codomain    ::= "Prop" | type
type        ::= NAME                     -- a sort (or a theory: model-of)
              | term ("Fin" | "Index")   -- bounded index (n elements)
              | NAME "DynTensor"         -- dynamic tensor over an element sort
              | NAME "List"              -- finite sequences over an element sort
              | arg+ NAME ["instance"]   -- model-of application
\end{Verbatim}

\noindent Braces name record domains. (The bracket spelling
\texttt{[f:\,v]} of earlier corpora still parses by default while
the committed corpus ports to braces; under strict parsing it is a
deprecation error pointing at the brace form, and \texttt{geolog
fmt} prints the canonical respelling --- \S\ref{sec:gr-cli}.) A
curried signature of named binders normalizes to the record form at
parse, so the two spellings declare the same relation or function;
an unnamed curried binder is a hard error. A declaration with
codomain \texttt{Prop} declares a
\emph{relation}; any other codomain declares an element-valued
\emph{function} (a \texttt{const-decl} is the nullary case: a
constant of its sort). Relations come in two domain shapes at the
surface: \emph{record} (named fields; field order as written is
preserved for display, but fields are name-keyed, never positional)
and \emph{unary} (one anonymous argument). There is deliberately no
surface form for a zero-argument proposition: \texttt{Flag : Prop;}
parses as a constant declaration of a (nonexistent) sort
\texttt{Prop} and is refused downstream (\S\ref{sec:gr-elab}). A
proposition wants at least one argument, or it wants to be a fact.

Two type forms are new with the aggregation surface, and both are
codomain vocabulary for the loader's stratum rather than the chase
fragment's. \texttt{DynTensor} is the \emph{dynamic} tensor type
over an element sort — shape carried in the value, where the
meta-tier \texttt{Tensor} former carries shape in the type
(\S\ref{sec:gr-meta-tier}); tensor values are side-condition
payload, never term-language citizens. A theory name in type
position means ``a model of that theory'': bare (\texttt{base :
G;} under a \texttt{(G : Graph)} parameter), or applied to
arguments (\texttt{initial\_marking : n0 Marking;}), with an
optional postfix \texttt{instance} keyword as an explicitness
marker — neither spelling is canonical, and the source printer
preserves the author's choice. A third form belongs to
the list former (\S\ref{sec:gr-meta-tier}): \texttt{Str List} is the
finite-sequence type over an element sort, postfix like
\texttt{Index} and \texttt{DynTensor} and taking one element sort
only (\texttt{Str List List} is a syntax error), with list values
consumed through the \texttt{len}/\texttt{at}/\texttt{concat} ground
roster as side-condition payload — never term-language citizens. The
surface element scope is the builtin ground sorts: a list codomain
at a declared sort is refused by the normalizer (verified:
\texttt{List element "A" is not a builtin ground sort (v1 surface
scope)}), and the wire tier's wider element scope, nesting included,
arrives with its own surface spelling rather than by fall-through.
The v1 chase bridge refuses all three forms (\S\ref{sec:gr-elab}).

\texttt{extends} names a parent theory. In the single-file pipeline
the parent must be declared in the same file, and extension is
realized by elaborating the whole file as one union
(\S\ref{sec:gr-elab}); in the committed-revision pipeline extension
is a \texttt{TheoryExtends} row with its own conformance rules
(\S\ref{sec:gr-meta}).

\paragraph{Builtin ground sorts.} The authoring idiom accepts a
roster of ground sorts in declarations without a \texttt{Set}
statement — \texttt{Q}, \texttt{Z}, \texttt{Nat}, \texttt{B},
\texttt{Str}, \texttt{ENNRat}, \texttt{Shape}, \texttt{Timestamp},
\texttt{Bytes}, \texttt{PeerId}, \texttt{CommitId}, and the
aggregation surface's staged registry sort \texttt{Accum}
(\path{crates/geolog-normalize/src/builtin.rs}) — together with the
literal-dtype correspondence that types bare integer literals
(\texttt{q}/\texttt{z}/\texttt{nat}/…), the builtin comparison
and arithmetic vocabulary the grounding moves compile onto
(\S\ref{sec:gr-normalize}), and the aggregation tier's staged
arithmetic (\texttt{nat\_mul}, \texttt{succ}). This roster is,
\texttt{Accum} aside, the normalizer's \emph{accepted} subset of the
engine's primordial sort tier (\S\ref{sec:gr-meta-tier}) —
\texttt{Accum} belongs to the deferred aggregation surface
(\S\ref{sec:gr-editions}), not the primordial tier; the command-line
chase fragment is narrower still and refuses ground sorts at
elaboration (\S\ref{sec:gr-elab}).

\section{Axioms and geometric normal form}
\label{sec:gr-axioms}

\emph{Authority: \path{crates/geolog-normalize} (\path{parse.rs},
\path{normalize.rs}, \path{lint.rs}); the normal form is also stated,
as schema, by the meta-schema (\S\ref{sec:gr-meta}).}

\subsection{The axiom grammar}
\label{sec:gr-axiom-grammar}

\begin{Verbatim}[fontsize=\small,frame=leftline,rulecolor=\color{accentmid},framesep=3mm,xleftmargin=3mm]
axiom      ::= [mode] AXNAME ":" ["forall" binders "."] [premise] "|-" conclusion ";"
binders    ::= NAME ":" type ("," NAME ":" type)*
premise    ::= item ("," item)*
item       ::= atom
             | "not" atom
             | "(" item ")"
             | "exists" binders "." premise          -- premise only
             | "forall" binders "." iff-group        -- pattern moves
conclusion ::= "false"
             | disjunct ("or" disjunct)*
disjunct   ::= ["exists" binders "."] item ("," item)*
atom       ::= term [cmpop term]
cmpop      ::= "=" | "!=" | "<" | "<=" | ">" | ">="
term       ::= prod ("+" prod)*
prod       ::= postfix ("*" postfix)*
postfix    ::= primary NAME*                          -- application chain
primary    ::= "{" NAME ":" term ("," NAME ":" term)* "}"   -- record
             | INTEGER                                       -- literal
             | NAME                                          -- variable
\end{Verbatim}

\noindent Grammar facts with semantic weight: \texttt{or} binds
loosest in a conclusion, and an existential's scope never spans an
\texttt{or} — each disjunct carries its own \texttt{exists} prefix,
so the surface conclusion is literally the normal form's shape (a
disjunction \emph{of} existentially quantified conjunctions). An
\texttt{exists} group in a premise swallows the remaining premise
items up to the turnstile. Terms are postfix application chains —
\texttt{m p Marking/count} reads ``\texttt{Marking/count} applied to
(\texttt{p} applied to \texttt{m})'' — with \texttt{*} binding
tighter than \texttt{+}, record literals in the same brace idiom
declarations use (brackets remain the parsed legacy spelling during
the corpus port, \S\ref{sec:gr-decls}; on the aggregation surface
square brackets instead delimit contraction index groups in
\texttt{law} bodies, \S\ref{sec:gr-editions}), and bare integer
literals (typed later by the
normalizer). \texttt{|- false} declares a \emph{denial}: a rejection
rule, the empty disjunction.

\subsection{Geometric normal form}
\label{sec:gr-gnf}

The engine consumes \emph{geometric} theories, and the normal form is
the central definition of this chapter:

\begin{definition}[geometric sequent]
\label{def:gr-geometric}
A geometric sequent is a universally quantified entailment whose
premise is a finite conjunction of atoms and whose conclusion is a
finite disjunction of existentially quantified finite conjunctions of
atoms; the empty disjunction is $\bot$ (a denial). An atom is an
equation between terms or a relation applied to a term; the
quantified variables are typed at declared sorts.
\end{definition}

\noindent Geometric sequents are the sequent fragment of geometric
logic \cite[Part~D1]{johnstone2002}; database theory knows the
single-disjunct case as tuple- and equality-generating dependencies
\cite{fagin2005data}. The authoring idiom of
\S\ref{sec:gr-axiom-grammar} is richer than this — \texttt{not}, \texttt{exists} in premises,
\texttt{forall}-\texttt{iff} groups, comparison operators, arithmetic
sugar, \texttt{Index} types. The normalizer
(\S\ref{sec:gr-normalize}) compiles every departure down to
Definition~\ref{def:gr-geometric} by a registered move or rejects it
loudly; the geometric linter then mechanically enforces the normal
form on the normalizer's own output — premise a conjunction of atoms;
conclusion a disjunction of existentially quantified conjunctions; no
residual \texttt{Index}; no untyped literal; every relation declared
\emph{and} classified (\S\ref{sec:gr-registry}); every atom
well-sorted. Pipeline failures are collected and reported together,
never truncated to the first.

\subsection{The normalization moves}
\label{sec:gr-normalize}

Each move is either total on its registered pattern or a loud error;
nothing passes through half-rewritten. In source order of the
pipeline:

\begin{itemize}
\item \textbf{Premise negation}, discharged by atom class. A negated
\emph{side-condition} atom needs no logic at all: side-conditions
are decision procedures, and a decision procedure is
negation-complete, its complement being the same evaluator with the
verdict flipped. The discharge is therefore exact and in place ---
$\lnot S \land P \vdash C$ becomes $\mathit{not/S} \land P \vdash C$,
with the complement relation registered automatically (move census:
\texttt{negated premise side-condition S discharges in place as
not/S}). A negated equation at a ground comparison sort has its
complement in the builtin vocabulary already: $\lnot(a = b)$ at
\texttt{Q} or \texttt{Nat} becomes \texttt{q\_ne} or \texttt{nat\_ne},
in place. The one genuinely logical case is a negated \emph{ordinary}
atom, or an equation at a declared sort. Ordinary vocabulary is
derivation-defined, so no procedure can decide its complement;
instead the atom \emph{crosses the turnstile} as a new leading
conclusion disjunct. $\lnot R \land P \vdash C$ becomes
$P \vdash R \lor C$: the classical-to-geometric compilation (move
census: \texttt{negated premise ordinary atom R crosses the
turnstile}). Read at the fixpoint, the crossed sequent says exactly
what the classical sentence said: wherever $P$ holds, $R$ holds or
$C$ does. What it deliberately does not offer is
negation-as-failure. The crossed head is a disjunction, so the axiom
is no longer chase-eligible (\S\ref{sec:gr-chase-classify}), and a
\texttt{chase} mark on it is a load refusal: a rule that fired on the
\emph{current} absence of a derivable row would be unsound against
every later derivation, so no move is owed to that reading. Under
\texttt{report}, the crossed disjunct is judged like any other menu
(\S\ref{sec:gr-chase-judgments}).
\item \textbf{Pattern rewrites} (\texttt{forall}-\texttt{iff}
groups). Extensionality, difference, and intersection shapes over a
registered membership relation rewrite to auxiliary mask
side-condition atoms; \texttt{not} on a membership atom in a
conclusion rewrites to a registered negation-free witness atom. All
four require registry entries (\S\ref{sec:gr-registry}); an
unregistered shape is a hard error, because auxiliary vocabulary is
an obligation on downstream evaluators and is never minted silently.
\item \textbf{Exists-hoist.} An \texttt{exists} group in a premise
hoists to universal variables — sequent semantics:
$(\exists x.\,P) \vdash C \;\equiv\; \forall x.\,(P \vdash C)$.
\item \textbf{Index erasure.} Every \texttt{n Index} type erases to
\texttt{Nat} plus a generated range obligation (a binder guard atom,
or a generated \texttt{ax/range/*} axiom).
\item \textbf{Grounding.} Comparison and arithmetic sugar compiles
onto the builtin roster at ground sorts: \texttt{!=} to
\texttt{q\_ne}/\texttt{nat\_ne}, \texttt{<}/\texttt{<=} to
\texttt{q\_lt}/\texttt{q\_le}/\texttt{nat\_lt}, \texttt{+}/\texttt{*}
to \texttt{q\_add}/\texttt{q\_mul}, with \texttt{>}/\texttt{>=}
flipped by operand swap. There is no grounding at declared sorts:
\texttt{a != b} over a declared sort \texttt{A} is a hard error
(verified: \texttt{no ground builtin for != at sort A (the !=@A case
would need a builtin-roster addition)}). Disequality of declared-sort
elements is not expressible in the fragment — deliberately: the
chase's equality discipline is congruence
(\S\ref{sec:gr-chase-congruence}), and a general disequality atom
over invented elements is exactly the kind of vocabulary that would
need an evaluator nobody has priced. (A negated premise
\emph{equation} at a declared sort is a different matter: it crosses
the turnstile as an equality conclusion, per the premise-negation
move above, so the classical constraint remains sayable even though
a disequality atom does not exist.)
\item \textbf{Literal typing.} Bare integer literals receive their
dtype from context (\texttt{q:3/1}, \texttt{nat:2}, …); an
untypeable literal fails the linter.
\end{itemize}

\noindent An axiom touched by any move is renamed with an
\texttt{\_n} suffix, so normalized output re-normalizes to itself
with an empty move log. The pipeline's guarantees, carried by its
test suite: two-run determinism (byte-identical listings),
idempotence (normalize $\fatsemi$ normalize $=$ normalize),
linter-cleanliness of every output, ground meaning preservation (the
grounding moves agree with direct evaluation of the source operators
on random ground instances), and committed fixture goldens per move.
\texttt{iff} at sequent top level (which would become two sequents)
is unimplemented and errors loudly if it ever fires.

\subsection{Imports, classification, and roster totality}
\label{sec:gr-registry}

\emph{Authority: \path{crates/geolog-normalize/src/imports.rs}
(the import surface), \path{crates/geolog-normalize/src/registry.rs}
(the registry value both surfaces feed).}

Every declared \texttt{Prop}-valued relation is classified, exactly
one of:

\begin{itemize}
\item \texttt{ordinary} — derivation-defined: populated by assertion
and by the chase, and by nothing else;
\item \texttt{side-condition} — decided by a bounded external
procedure (the ground seam, \S\ref{sec:gr-ground}, or a registered
auxiliary evaluator), never derivation-defined.
\end{itemize}

\noindent This is \emph{roster totality}: every relation has a
decided evaluator, and it holds \emph{by construction} — the
obligation is discharged by syntax instead of by hand. The
classification lives in the source file itself. \texttt{ordinary} is
the default, declared by writing nothing; a side-condition is
declared by an \texttt{import} item naming its provider:

\begin{Verbatim}[fontsize=\small,frame=leftline,rulecolor=\color{accentmid},framesep=3mm,xleftmargin=3mm]
import std.{Q, q_lt};
import <provider>.{<relation>, ...};
\end{Verbatim}

\noindent The two namespaces do different work. An
\texttt{import std} item authorizes \emph{builtin ground vocabulary}
used by name — sorts like \texttt{Q} and \texttt{Nat}, comparisons
like \texttt{q\_lt} — and a name outside the builtin roster refuses
at once, listing the importable names. Any other namespace declares
each named relation a side-condition awaiting that provider, bound
at the ground seam (an unknown provider is not an error at check
time: binding is deferred, and an unbound side-condition atom simply
cannot be discharged — refusal over guessing, downstream). Import
items are position-free top-level items, so multi-file closures
concatenate; re-importing a name under the same namespace is
idempotent, and importing one name under two namespaces refuses ---
one name, one provider. A relation or function declaration that
reuses a builtin \emph{vocabulary} name is legal only as an
exact-signature shim of that builtin: same-named, byte-equivalent to
what the engine itself would say. A builtin \emph{sort's} name admits
no shim at all — it is fully reserved, and a local sort takes another
name (verified: \texttt{sort Q (declared in T): builtin ground sort
names are reserved — the builtin Q is ambient (import std.\{Q\};
declares the usage), so a local sort must take another name}). Both
reservations are kind-local (the sort and vocabulary namespaces never
meet), so a relation may take a builtin sort's name, and a declared
sort may take a builtin relation's. Builtin
comparison relations and the normalizer's own minted auxiliaries are
side-conditions without import lines. Under \texttt{--strict},
builtin ground vocabulary used by name demands its \texttt{import
std} header even in files that otherwise import nothing.

The \emph{pattern registry} is the second classification surface,
today a library-level one: it names, per rewritten relation, its
\emph{membership} reading (which field is subject, which is index),
and the auxiliary relation each registered move targets
(\texttt{extensionality}: 2 fields; \texttt{difference} and
\texttt{intersection}: 3; \texttt{negation-witness}: the relation's
own arity). Registry construction validates all of this — missing
membership entries, unregistered auxiliaries, arity mismatches, and
duplicate auxiliaries are errors before any theory text is touched.
Pattern registrations have no source-level surface yet; the
library's registry value carries them (the normalizer's
\path{fixtures/} include pattern-bearing specimens).

A line-oriented registry \emph{file} (\texttt{--registry}) was the
classification's original carrier; it is deprecated, retained for
the migration era only. The one-authority law binds at the
registry's content, and classification is the content it guards: a
registry whose lines classify relations refuses to combine with a
source that declares imports (verified: \texttt{this source
declares imports AND an external registry carries relation
classification — classification has one home; …}), while a
pattern-only registry file composes with an importing source — its
aux names kept clear of the imported ones (verified: \texttt{aux
relation … collides with an imported name — aux mask vocabulary is
the normalizer's, never imported}), a collision being
classification smuggled back in aux costume. The
line format also carries the pattern registrations of the paragraph
above (their one textual carrier while patterns have no in-source
spelling), and a pattern registration is no second classification
authority.

\section{Instance blocks}
\label{sec:gr-instances}

\emph{Authority: \path{crates/geolog/src/facts.rs} (grammar),
\path{crates/geolog/src/seed.rs} (semantics).}

\begin{Verbatim}[fontsize=\small,frame=leftline,rulecolor=\color{accentmid},framesep=3mm,xleftmargin=3mm]
instance-block ::= "instance" NAME ":" NAME "{" item* "}"
item           ::= elem-decl | record-fact | bare-fact
elem-decl      ::= NAME ("," NAME)* ":" NAME ";"
record-fact    ::= "{" NAME ":" NAME ("," NAME ":" NAME)* "}" NAME ";"
bare-fact      ::= NAME NAME ";"
\end{Verbatim}

\noindent Tokenization is the theory lexer's, so comments and
identifiers mean the same thing in both strata. The header names the
block and the theory it instantiates; the named theory must be
declared in the same file. Items are element declarations and facts;
a record fact's brace lists \texttt{field: element} pairs and then
the relation; a bare fact (\texttt{a Marked;}) applies a unary-kind
relation. Facts must be written in the relation's own shape — a
record fact on a unary relation, or vice versa, is an error naming
the mismatch. Braces are canonical here as in declarations: this
grammar is what \texttt{geolog fmt} prints, and the formatter's
output re-parses (verified: the respelled Quick Start file chases
end-to-end, its report byte-identical to the bracket original's),
while the bracket spelling \texttt{[f:\,a]\,Rel;} remains the
accepted legacy form until the committed corpus ports to braces.

Semantics, in seeding order (source order — element numbering in
transcripts is deterministic):

\begin{itemize}
\item \textbf{Elements are declared before use.} A fact referencing
an undeclared name is an error, never a silently minted element
(verified: \texttt{`P` fact references `ghost`, which is not a
declared element — declare it first (`ghost : A;`); elements are
never minted silently}).
\item \textbf{Fields are exact.} A record fact must supply every
declared field, no extras, no repeats, no defaults; each element's
sort must equal the field's declared sort.
\item \textbf{Union semantics across blocks.} An element name
declared in two blocks at the same sort denotes one element;
duplicate facts are idempotent. A name redeclared at a
\emph{different} sort is an error (verified: \texttt{element `a`
redeclared at sort `C` but already declared at sort `A`}). This is
the merge story — set union of commits, then re-derivation
(\S\ref{sec:gr-sync-merge}) — at authoring scale; the Tutorial's
two-writer example is built on it.
\end{itemize}

\noindent Surface syntax is a view: this block form is the command
line's authoring view of an instance, and the durable form is signed
commits with content-addressed element identity
(\S\ref{sec:gr-commits}).

\section{Elaboration: the v1 command-line fragment}
\label{sec:gr-elab}

\emph{Authority: \path{crates/geolog/src/elab.rs}; the crate-level
docs of \path{crates/geolog/src/lib.rs} state the fragment.}

The bridge from the normalized authoring idiom to the engine IR
covers, in v1: \textbf{sorts}, \textbf{\texttt{Prop}-valued
relations} (record or unary domains), and \textbf{geometric axioms
whose atoms apply relations to variables}, equality conclusions
included. Everything else the idiom can express is refused with an
error naming the construct and, where one exists, the supported
spelling of the same intent. The engine library supports more than
this bridge; the bridge grows by deliberate widening, never by silent
acceptance.

Elaboration takes the file's theory declarations as \emph{one
single-file closure}: the union of every declaration in the file,
with \texttt{extends} between in-file theories honoured by that same
union. Consequences: an instance block may reference any declaration
in the file's closure regardless of which theory block carries it;
duplicate sort or relation names anywhere in the closure are errors;
and \texttt{extends} naming an out-of-file parent is refused
(verified: \texttt{theory `Child` extends `Ghost`, which is not
declared in this file — cross-revision extension rides the
committed-revision pipeline (geolog-meta-rows / geolog-revision), not
the v1 CLI bridge}).

The operative refusals, each verified against the committed binary:

\begin{itemize}
\item \textbf{Element-valued functions} (any declaration with a
non-\texttt{Prop} codomain, constants included):
\begin{Verbatim}[fontsize=\small,frame=leftline,rulecolor=\color{accentmid},framesep=3mm,xleftmargin=3mm]
function symbol `f` (codomain `C`): element-valued functions are engine
IR the v1 bridge does not carry — spell the function as a relation plus
a functionality axiom (two rows in the premise, an equality in the
conclusion), the idiom the committed theories use (e.g. `Marking/count`
+ `ax/count/functional` in theories/PetriNet.geolog)
\end{Verbatim}
The named idiom is normative: a function $f : A \to B$ is spelled as
a relation $F[a{:}\,A, b{:}\,B]$ plus
$\forall a, b_1, b_2.\; F[a, b_1], F[a, b_2] \vdash b_1 = b_2$, and
the chase's congruence closure does the rest
(\S\ref{sec:gr-chase-congruence}). \texttt{Flag : Prop;} lands in
this refusal too (its ``codomain'' parses as the sort
\texttt{Prop}); see \S\ref{sec:gr-decls} for why no zero-argument
proposition form exists.
\item \textbf{Ground sorts in declarations} — any field, argument, or
binder sort not declared \texttt{Set} in the file, the builtin roster
included:
\begin{Verbatim}[fontsize=\small,frame=leftline,rulecolor=\color{accentmid},framesep=3mm,xleftmargin=3mm]
relation `R` field `n` names sort `Nat`, which this file does not
declare — the v1 bridge elaborates single-file closures over declared
`Set` sorts only (ground sorts such as `Nat` arrive with the ground
seam; see crates/geolog-ground)
\end{Verbatim}
This is also where \texttt{Index} types land operationally: the
normalizer erases \texttt{n Index} to \texttt{Nat}
(\S\ref{sec:gr-normalize}), and the bridge then refuses \texttt{Nat}
as above. The worked discipline for doing real arithmetic
\emph{without} the seam — counts as explicit elements on a ladder,
chase-derived — is \path{theories/PetriNet.geolog}.
\item \textbf{Tensor, list, and model-of types} — the codomain
forms of \S\ref{sec:gr-decls} beyond plain sorts are loader
vocabulary, and the bridge carries a dedicated refusal for each. A
\texttt{DynTensor} anywhere in a declaration:
\begin{Verbatim}[fontsize=\small,frame=leftline,rulecolor=\color{accentmid},framesep=3mm,xleftmargin=3mm]
relation `P` field `x` uses a `Tensor` type — tensor payloads are
side-condition territory, not v1 bridge vocabulary
\end{Verbatim}
A model-of member, on today's routes, is caught upstream by the
element-valued-function refusal instead (its codomain is not
\texttt{Prop}); the bridge's own instance-valued arm — refusing
instance members as loader territory — stands behind it as
defensive depth, per the posture above. A \texttt{List} codomain is
caught upstream the same way; a \texttt{List} in a \emph{field}
position reaches the bridge's own list arm (verified:
\texttt{relation `P` field `x` uses a `List` type — list payloads
are side-condition territory (the \{len, at, concat\} roster), not
v1 bridge vocabulary}).
\item \textbf{Ground literals and term-level application in axiom
atoms}, and nested records in argument positions: refused naming the
axiom. Under the current pipeline these arms are defensive depth —
declaration-level refusals and normalizer typing catch the practical
routes first — and they stay, per the bridge's stated posture: a
bridge that assumes its input is well-formed fails silently the day
the assumption breaks.
\end{itemize}

\noindent Axiom-side checks (variables bound by \texttt{forall} or an
enclosing \texttt{exists}; sorts agreeing with declarations; fields
exact; no variable bound twice; no existential shadowing an outer
variable) are belt-and-braces re-verification of what the normalizer
already established, kept for the same reason.

\section{The chase}
\label{sec:gr-chase}

\emph{Authority: \path{crates/geolog-chase} — \path{classify.rs},
\path{chase.rs}, \path{instance.rs}, \path{judgment.rs}. The engine
has no silent path: every outcome below is either a derived row, a
typed refusal, or a reported judgment.}

\subsection{Sequent classification}
\label{sec:gr-chase-classify}

Repair arity is computed from syntax alone, by one total function,
and every engine consults it — no runtime data participates. The
conclusion is normalized to disjunctive head form ($\mathrm{Conj}$
distributes over $\mathrm{Disj}$; a $\bot$ conjunct annihilates its
branch), and:

\begin{center}
\begin{tabular}{lll}
\toprule
head disjuncts & class & chase posture \\
\midrule
$0$ & denial & never fires; violations are judgments \\
$1$ & canonical & fires (invents witnesses) — the regular fragment \\
$\geq 2$ & repair menu & never auto-fired; surfaced as a branch point \\
\bottomrule
\end{tabular}
\end{center}

\noindent One carve-out is permanent: a head that reaches an
existential over an \emph{immutable} sort (an infinitary disjunction
— the engine IR's \texttt{InfDisj}) is never chase-eligible
regardless of arity, because an infinitary disjunction is in no way a
unary disjunction; no witness-invention strategy for it can ever be
sound. The refusal is loud and names the axiom.

\subsection{Modes: the author's word on firing}
\label{sec:gr-chase-modes}

Whether an axiom fires is not inferred from its shape: it is
\emph{authored}, per axiom, by the mode mark of
\S\ref{sec:gr-decls}. The three marks:

\begin{itemize}
\item \texttt{chase} — a derivation rule: the axiom fires, inventing
witnesses and asserting conclusions. Only a chase-eligible axiom
(\S\ref{sec:gr-chase-classify}: canonical head, no \texttt{InfDisj})
may carry the mark; \texttt{chase} on an ineligible axiom is a load
refusal naming the axiom by its \emph{source} name (verified:
\texttt{axiom ax/pick/one is marked `chase` but is not chase-eligible
…; mark it `report` (judge) or `refuse` (gate), or reshape its
head}).
\item \texttt{report} — a constraint: the axiom never fires; its
violations and would-be firings surface as judgments
(\S\ref{sec:gr-chase-judgments}).
\item \texttt{refuse} — a constraint that additionally gates local
authoring: \texttt{report}'s judgment semantics plus the local write
gate of \S\ref{sec:gr-engine}. On the file path there is no commit to
gate, so its violations render exactly as \texttt{report}'s.
\end{itemize}

\noindent \textbf{An unmarked axiom means \texttt{report}} — the
least-dangerous default (\texttt{chase} creates rows and
\texttt{refuse} blocks commits; \texttt{report} only judges), so
derivation is always an explicit authorial act. Denials sit naturally
under \texttt{report}: a denial never fires in any mode, and its
violations are phase-2 judgments (\S\ref{sec:gr-chase-run})
regardless.

The marks map onto the engine's per-axiom trichotomy
(\S\ref{sec:gr-engine}): \texttt{chase} $\to$ Chased, \texttt{report}
$\to$ Monitored, \texttt{refuse} $\to$ Enforced. \texttt{geolog
chase} runs a file under this authored partition by default,
recovering each axiom's mark from the normalized declarations with
three alignment checks per axiom (count, name, binder vector) so any
drift between the normalizer's and the bridge's axiom order is a
loud refusal, never a silent mode reassignment.

\paragraph{The legacy flattener.} \texttt{--mode chased|monitored}
(\S\ref{sec:gr-cli}) survives for corpora not yet carrying explicit
marks, applying one run-level policy to every axiom: under
\texttt{chased}, every eligible axiom fires and a disjunctive head is
a theory-load error before the loop ever runs (verified, on an
unmarked file: \texttt{theory load refused: axiom ax\#0 is not
Chased-eligible: repair arity is Menu(2); only a single-disjunct
(canonical) head with no InfDisj binder may run under mode Chased —
declare it Monitored instead, or split/simplify the axiom}); under
\texttt{monitored}, eligible axioms fire and disjunctive heads
surface as branch points. On a file carrying \emph{any} explicit mark
the flag refuses, naming the marked axioms (verified: \texttt{--mode
conflicts with this file's explicit surface marks: ax/path/edge
(chase), … — drop the flag to honor the marks}): a mark is the
author's word, and a global flag that silently promoted a
\texttt{report}-marked axiom into the firing set would be the exact
hazard the marks exist to close — worst case, an equality conclusion
merging the very elements it was written to police. Where the flag
does run, the flattening is announced on stderr.

\noindent Engine-level load refusals name axioms positionally
(\texttt{ax\#i}): engine sequents carry no name field — names are
metadata, held by the caller (or, in the durable pipeline, by
meta-schema rows) — and the facade maps indices back to source names
wherever it holds them (\S\ref{sec:gr-report}); the bare positional
form above is the legacy flattener's path showing the engine's own
words.

\subsection{The two-phase run}
\label{sec:gr-chase-run}

A chase run has exactly two phases, and the separation exists for
one reason: \emph{a transient judgment can never be persisted as
if final.}

\textbf{Phase 1} fires only canonical axioms, to a fixpoint or the
iteration bound. The loop's continuation condition is instance
mutation only — discovering a violation is never a reason to keep (or
stop) iterating. \textbf{Phase 2} is one evaluation pass of every
denial and every menu axiom against the \emph{final} canonical
instance, bindings resolved through the final congruence. A menu
whose missing witness is supplied by the very last canonical firing
is therefore never recorded stale, and every judgment binding is a
canonical representative of the final fixpoint.

Hitting the bound is not an error: the instance-so-far is returned, a
non-termination judgment is recorded (with both the iterations run
and the bound), and phase 2 still reports whatever holds of the
partial instance — validity is explicitly indeterminate, not omitted
(verified: \texttt{no fixpoint within budget: 5 iterations (bound 5)
— raise --budget, or the theory does not saturate on this instance
(unbounded generation is a real outcome, reported rather than guessed
at)}).

\subsection{Restricted semantics and invention}
\label{sec:gr-chase-restricted}

Canonical firing is the \emph{restricted} chase
\cite{deutschnashremmel2008}: for each binding of
an axiom's universal context that satisfies the premise, the engine
first evaluates the conclusion as a query under that binding —
\textbf{if any witness already exists, it does not fire}. Only a
witnessless binding invents: one fresh element per existential
binder, then the conclusion's atoms are asserted (relation rows as
derived rows; equations via congruence,
\S\ref{sec:gr-chase-congruence}). The Tutorial's shipment example is
the behaviour in one line: the order with a recorded manifest got no
invention; the order without one got exactly one fresh
\texttt{?Shipment}.

Every element carries an immutable creation-provenance record,
written once and never rewritten by any merge: seeded (EDB) elements
carry caller-supplied provenance bytes (at command-line scale, the
declared name; in the durable pipeline, a commit-addressed leaf), and
chase-invented elements carry the firing axiom's index, the
context-bound elements of the firing (in context order), and the
fresh binder's ordinal within the axiom. Derived rows are
distinguished from asserted rows throughout
(\S\ref{sec:gr-report}). One invariant worth its ink: for a premise
whose own existential (hoisted per \S\ref{sec:gr-normalize}) admits
several matches over the same context binding, the first firing's
witness satisfies the rest — so exactly one witness, with a
context-only provenance term, ever exists per context binding.

\subsection{Equality and congruence}
\label{sec:gr-chase-congruence}

Equality is computational, not annotational: the instance maintains a
union-find congruence closure in the style of Nieuwenhuis and
Oliveras \cite{nieuwenhuisoliveras2007}. An
equational conclusion firing merges two elements' classes; everything
known about either — rows, provenance-bearing derivations, later
firings — holds of the merged class. Tuples are stored raw and
canonicalized through the current union-find at every read, so a
merge never rewrites tables. The report renders a merged class by
joining its seeded names (\texttt{sf\_record=wh\_label};
\texttt{b=c}) and accounts for it in the carrier arithmetic
(\texttt{3 seeded, 1 merged by congruence} — verified transcript).
The functionality idiom of \S\ref{sec:gr-elab} plus this machinery is
how geolog has functions without function symbols in the v1
fragment. (Chase steps interleaved with congruence closure are also
the computational core of datalog-with-equality and
equality-saturation engines \cite{bidlingmaier23,zhang23}, the
nearest kin to this engine in the evaluation literature.)

\subsection{Judgments}
\label{sec:gr-chase-judgments}

A judgment is one thing the chase found that a silent-success engine
would have swallowed — the authority module's own one-line definition
of the type. The chase's loud outputs, exhaustively, each carrying
the axiom's identity and, where applicable, the witnessing binding
with canonical elements:

\begin{itemize}
\item \textbf{Denial violation} — a denial's premise is satisfied at
the fixpoint: $\vdash \bot$ for this binding. Reported, never
dropped, never repaired.
\item \textbf{Branch point} — a menu axiom's premise is satisfied
with no existing witness for \emph{any} disjunct. The payload carries
the menu (the disjuncts themselves, as data); the engine picks
nothing. Under the legacy run-level flattener a menu has length
$\geq 2$; menu-of-one branch points (a forced repair, surfaced rather
than fired — \texttt{a menu of 1 repair}) arise wherever the
per-axiom partition holds an \emph{eligible} axiom to \texttt{report}
(\S\ref{sec:gr-chase-modes}, \S\ref{sec:gr-engine}) — an unmarked
derivation rule's would-be firings surface this way.
\item \textbf{Non-termination} — the bound was hit
(\S\ref{sec:gr-chase-run}).
\end{itemize}

\noindent Judgments are results, not errors: a run that produces
them completed, exits with the judgment exit code
(\S\ref{sec:gr-cli}), and its instance is real — the event-sourcing
posture (invalid states representable, detected, repaired by further
ordinary commits) held formally. The Tutorial's merge chapter is the
worked argument for why this posture, and not write-rejection or
automatic repair, is the design.

\subsection{The report}
\label{sec:gr-report}

\emph{Authority: \path{crates/geolog/src/report.rs}.}

The \texttt{chase} verb renders every carrier, every row, and every
judgment, through the source file's own names; nothing about a
judgment is summarized away. The normative rendering facts, seven of
them:

\begin{itemize}
\item each sort prints its canonical carrier size with the seeded
count, the count invented by the chase, and the count merged by
congruence, followed by a \texttt{merged:} line naming merged
seeded-name groups;
\item each relation prints in declaration order, rows sorted and
deduplicated after canonicalization, with seeded/derived tallies;
\item derived rows — rows not in the canonicalized image of the
seeded baseline — are marked \texttt{+};
\item rows print in the relation's own shape (record, bare, or bare
name for nullary);
\item chase-invented elements render as \texttt{?Sort.n} in
first-encounter order;
\item side-condition relations are tagged \texttt{[side-condition]};
\item judgments print the axiom's \emph{source} name and the binding
in context order, a branch point adds its menu size, and a clean run
ends \texttt{fixpoint: reached, no judgements — the model is saturated.} (with a count of invented elements when there are any).
\end{itemize}

\section{The engine beneath the command line}
\label{sec:gr-engine}

\emph{Authority: \path{crates/geolog-chase} (\path{mode.rs},
\path{weight.rs}, \path{view.rs}, \path{delta.rs},
\path{plan_ir.rs}, \path{validity/}), \path{crates/geolog-weights}.
Everything in this section is shipped, tested library surface. The
per-axiom mode partition is the one piece the \texttt{geolog} verb
now reaches — through the surface marks of
\S\ref{sec:gr-chase-modes}; the rest is not reachable from the verb
yet.}

\paragraph{Per-axiom modes, and Enforced.} The engine's mode system
is a per-axiom assignment over the trichotomy \textbf{Chased} /
\textbf{Monitored} / \textbf{Enforced}, with an eligibility-keyed
ambient default (chase-eligible $\Rightarrow$ Chased, everything else
Monitored — reproducing the legacy run-level default, so that default
is semantics-preserving) and validated explicit overrides: a Chased
override on an ineligible axiom is a typed error; Monitored/Enforced
overrides on \emph{eligible} axioms are honoured by a firing
partition — the axiom never fires, and its would-be firings surface
as menu-of-one branch points (the forced repair, surfaced, never
auto-fired). The surface marks lower onto this same machinery, as
a \emph{total} explicit assignment — every axiom's effective mark,
unmarked resolved to \texttt{report} first — so on the file path the
author's word, never eligibility, decides firing, and the ambient
default decides nothing; the ambient default remains the semantics
for library callers that override selectively. Enforced is Monitored
plus the \emph{local write gate}:
an admission verdict over a finished chase result rejects any local
commit whose resulting fixpoint carries an Enforced violation. The
commit-vs-sync seam is the caller's, and its asymmetry is doctrine:
the gate binds \textbf{local commit paths only} — states arriving by
sync or merge are admitted with their judgments recorded, because
gating sync admission on validity would cost true merge and
history-as-authority (\S\ref{sec:gr-sync-merge}). Modes differ in
\emph{when} validity is enforced, never in what validity is. The
adjudicator refuses to grade a mode assignment against a run that was
not partitioned under that same assignment — the silent-vacuity class
(an axiom fired in phase 1 repairing its violations before any
judgment could exist) is rejected loudly, never admitted vacuously.
The meta-schema carries the same trichotomy as rows
(\texttt{Axiom/monitored}, \texttt{Axiom/enforced}, at most one per
axiom; absence of both is the ambient default) —
\S\ref{sec:gr-meta}.

\paragraph{Weights.} The engine is parameterized over an
ordered-semiring provenance-weight algebra
(\path{crates/geolog-weights}): a weighted run integrates each derived
row's weight from its derivation structure, a reweight pass
recomputes after tombstoning, and death-at-zero holds — the
weighted-alive set after tombstone + reweight equals a from-scratch
Boolean chase on the surviving facts (proptested as a differential).
The Boolean rung \emph{is} the unweighted engine: \texttt{run\_chase}
stays pinned to $W = \texttt{bool}$ and is byte-identical to the
weight-blind engine.

\paragraph{Deltas, maintained views, and the plan IR.} The delta bus
makes the engine's unit of change data (element, row, and function
deltas in batches). The view engine compiles violation views —
the failing set of each denial, kept current under row deltas, never
recomputed from scratch — and verifies itself against from-scratch
reference evaluation (a divergence is a reportable defect, and the
rebuild invariant of the repository's battery backs it at replay
scale). Compiled views lower to a tensor-contraction plan artefact
with canonical, digest-citable bytes, a dense reference evaluator,
and an accelerator emission harness; the reference evaluator is the
verdict authority throughout (the lane law,
\S\ref{sec:gr-ground}).

\paragraph{Per-commit validity.} The validity layer
(\path{crates/geolog-chase/validity/}) assembles per-commit validity
judgments over chase results: which commits does this judgment
implicate — well-posed exactly because every derived row and every
invented element carries its firing history
(\S\ref{sec:gr-chase-restricted}). Wiring it into the storage
pipeline end-to-end is design tense (the factoring is committed; the
Tutorial states the same status).

\section{The meta-schema}
\label{sec:gr-meta}

\emph{Authority: \path{theories/GeologMeta.geolog} — the meta-schema's
normative statement, from whose text this section is condensed; where
they disagree, that file wins. The committed genesis golden
(\S\ref{sec:gr-goldens}) is its reference instance; the generator
tool (\path{theories/tools/geologmeta_gen.py}) re-derives and
byte-compares both directions as a drift alarm in CI.}

geolog is a logical framework, not an expressive logic. GeologMeta is
the fixed schema \emph{of} geometric theories: one meta-schema,
frozen at genesis; every theory — including the proof calculi the
companion repositories ship — is an ordinary instance: delta rows
against GeologMeta, carried by ordinary signed commits. Schema,
extension, annotation, and migration therefore live in the same
substrate they govern (the Tutorial's ``no second database'').

\subsection{Sorts and rows}
\label{sec:gr-meta-rows}

Nine schema sorts: \texttt{Theory}, \texttt{Sort}, \texttt{Fun}
(function \emph{or} relation — relationality is a mark, not a second
sort), \texttt{Field}, \texttt{Axiom}, \texttt{Var},
\texttt{Disjunct}, \texttt{Atom}, \texttt{Term}.
Table~\ref{tab:gr-meta-rows} carries the declaration vocabulary over
them, with partiality as documented there (``partial'' in its notes
column means the row may be absent, and absence is meaningful —
Codd's I-mark discipline \cite{codd86}, never ``unknown yet'').

\begin{table}[htbp]
\centering\footnotesize
\begin{tabular}{@{}ll>{\raggedright\arraybackslash}p{0.40\textwidth}@{}}
\toprule
subject & rows & notes \\
\midrule
\texttt{Theory} & \texttt{TheoryExtends[child, parent]} & multi-parent legal; extends graph acyclic \\
 & \texttt{TheoryRevises[new, old]} & revision intent; acyclic; causal half lives in deps \\
 & (no \texttt{Theory/name}) & identity is the element; discovery naming lives above \\
\texttt{Sort} & \texttt{Sort/name}, \texttt{Sort/theory} & \\
 & \texttt{Sort/immutable} & carrier opaque to the chase (choose, never create) \\
 & \texttt{Sort/builtin\_code} & partial: builtin tier only (\S\ref{sec:gr-meta-tier}) \\
\texttt{Fun} & \texttt{Fun/name}, \texttt{Fun/theory} & \\
 & \texttt{Fun/relational} & the relation mark; \texttt{Prop} is not a sort \\
 & \texttt{Fun/codomain} & partial: exactly the non-relational Funs \\
 & \texttt{Fun/arg\_sort} & partial: exactly the unary Funs \\
 & \texttt{Fun/side\_condition} & bounded-computation discharge; entails relational \\
\texttt{Field} & \texttt{Field/fun}, \texttt{Field/name}, \texttt{Field/sort} & name-keyed; canonical order is bytewise name order \\
\texttt{Axiom} & \texttt{Axiom/theory}, \texttt{Axiom/name} & name partial, recommended \\
 & \texttt{Axiom/monitored}, \texttt{Axiom/enforced} & mode rows (\S\ref{sec:gr-engine}); at most one; surface \texttt{report}/\texttt{refuse} lower here, \texttt{chase} to neither \\
\texttt{Var} & \texttt{Var/premise\_of} $\to$ \texttt{Axiom} & universal variables, scoped to the whole axiom \\
 & \texttt{Var/bound\_in} $\to$ \texttt{Disjunct} & existential variables, scoped to one disjunct; exactly one home (partition law) \\
 & \texttt{Var/sort}, \texttt{Var/name} & name is presentation only, hash-excluded \\
\texttt{Disjunct} & \texttt{Disjunct/axiom} $\to$ \texttt{Axiom} & one branch of the conclusion; zero disjuncts on an axiom $=$ rejection rule \\
\texttt{Atom} & \texttt{Atom/premise\_of} $\to$ \texttt{Axiom} & premise atoms hang on the axiom itself \\
 & \texttt{Atom/concl\_of} $\to$ \texttt{Disjunct} & conclusion atoms hang on a disjunct, never on the axiom; exactly one home \\
 & \texttt{Atom/eq\_lhs}, \texttt{Atom/eq\_rhs} & equation atoms \\
 & \texttt{Atom/rel}, \texttt{Atom/rel\_arg} & application atoms; exactly one kind \\
\texttt{Term} & \texttt{Term/var} & variable \\
 & \texttt{Term/app\_fun}, \texttt{Term/app\_arg} & application (arg absent iff nullary) \\
 & \texttt{Term/is\_record}, \texttt{RecordField[rec, field, value]} & records; empty record is the explicit mark \\
 & \texttt{Term/lit\_q,\,lit\_z,\,lit\_nat,\,lit\_b,\,lit\_str,} & one typed literal row per primordial dtype \\
 & \texttt{\ \ lit\_shape,\,lit\_ennrat,\,lit\_timestamp} & \\
\bottomrule
\end{tabular}
\caption{The meta-schema's declaration vocabulary, by subject sort
(\S\ref{sec:gr-meta-rows}).}
\label{tab:gr-meta-rows}
\end{table}

Domain shapes are discriminated by which rows exist —
nullary: neither \texttt{Fun/arg\_sort} nor \texttt{Field} rows;
unary: \texttt{Fun/arg\_sort}; record: \texttt{Field} rows.

The axiom encoding deserves a picture, because it has one level the
sequent notation hides (Figure~\ref{fig:gr-axiom-rows}). The premise
hangs directly on the \texttt{Axiom} element: premise atoms by
\texttt{Atom/premise\_of}, universal variables by
\texttt{Var/premise\_of}. The conclusion sits one level deeper. A
conclusion is a disjunction of existentially quantified conjunctions
(Definition~\ref{def:gr-geometric}), so each branch is its own
\texttt{Disjunct} element, tied to the axiom by
\texttt{Disjunct/axiom}; a conclusion atom names its
\emph{disjunct} as home (\texttt{Atom/concl\_of}), never the axiom,
and an existentially bound variable does the same
(\texttt{Var/bound\_in}). An axiom's premise is therefore the
\emph{set} of atoms homed on it, its conclusion the \emph{set} of its
disjuncts, and its repair arity (\S\ref{sec:gr-chase-classify}) is
its \texttt{Disjunct} count: zero is the empty disjunction, a
rejection rule; one is the canonical, chase-eligible shape; two or
more is a repair menu.

\begin{figure}
\centering
\begin{tikzpicture}[
  elem/.style={draw, rounded corners=1.5pt, inner sep=4pt,
               font=\small\ttfamily},
  rowlab/.style={font=\scriptsize\ttfamily, text=accentmid!75!black,
                 inner sep=1.5pt, fill=white},
  arr/.style={-{Stealth[length=1.8mm]}, semithick, draw=black!70,
              shorten >=2pt, shorten <=2pt}
]
\node[elem] (ax) at (-0.4,0) {ax/ship/manifest : Axiom};
\node[elem] (po) at (-4.7,-1.8) {o : Var};
\node[elem] (pa) at (-1.7,-1.8) {[o:\,o]\,Shippable : Atom};
\node[elem] (d)  at (2.9,-1.8) {d : Disjunct};
\node[elem] (cs) at (0.6,-3.4) {s : Var};
\node[elem] (ca) at (4.4,-3.4) {[o:\,o,\,s:\,s]\,ShippedVia : Atom};
\draw[arr] (po) -- (ax.south west) node[rowlab, pos=0.5, left=4pt] {Var/premise\_of};
\draw[arr] (pa) -- (ax.south) node[rowlab, pos=0.25, right=4pt] {Atom/premise\_of};
\draw[arr] (d)  -- (ax.south east) node[rowlab, pos=0.42, right=4pt] {Disjunct/axiom};
\draw[arr] (cs) -- (d.south west)  node[rowlab, pos=0.5, left=4pt] {Var/bound\_in};
\draw[arr] (ca) -- (d.south east)  node[rowlab, pos=0.5, right=4pt] {Atom/concl\_of};
\end{tikzpicture}
\caption{The Tutorial's \texttt{ax/ship/manifest} ($\forall o.\;
\mathit{Shippable}(o) \vdash \exists s.\; \mathit{ShippedVia}(o,s)$)
as meta-schema rows. Premise structure homes on the \texttt{Axiom}
element; conclusion structure homes on a \texttt{Disjunct}, one
element per branch of the disjunction. Which side of the turnstile
an atom or variable lives on is read off the function row it
carries, and the axiom's repair arity is its \texttt{Disjunct}
count — here one, the canonical shape.}
\label{fig:gr-axiom-rows}
\end{figure}

Field-name uniqueness per
Fun and \texttt{RecordField} uniqueness per (record, field) are load
rejections, deliberately not unification axioms: in a store whose
rows are signed declarations, silently merging independently minted
declarations would erase provenance — collisions are loud.

The meta-axioms, exhaustively (all fifteen; three shapes recur —
entailment, rejection pair, and the $\langle$disjoint,
exhaustive$\rangle$ \emph{partition law} pair stating ``exactly
one''):

\begin{center}
\footnotesize
\begin{tabular}{lll}
\toprule
axiom & shape & states \\
\midrule
\texttt{ax/meta/builtin\_immutable} & entailment & coded sorts are immutable \\
\texttt{ax/meta/side\_condition\_relational} & entailment & side-condition $\Rightarrow$ relational \\
\texttt{ax/meta/relational\_no\_codomain} & rejection & a relation never has a codomain \\
\texttt{ax/meta/codomain\_exhaustive} & disjunctive & every Fun is relational or has a codomain \\
\texttt{ax/meta/arg\_sort\_fields\_disjoint} & rejection & unary and record domains exclude each other \\
\texttt{ax/meta/mode\_marks\_exclusive} & rejection & at most one explicit mode per axiom \\
\texttt{ax/meta/var\_home\_\{disjoint, exhaustive\}} & partition & every Var has exactly one home \\
\texttt{ax/meta/atom\_home\_\{disjoint, exhaustive\}} & partition & every Atom sits on one side of $\vdash$ \\
\texttt{ax/meta/atom\_kind\_\{disjoint, exhaustive\}} & partition & equation or application, never both/neither \\
\texttt{ax/meta/term\_\{var\_app, var\_record, app\_record\}\_disjoint} & rejection & representative term-kind disjointness \\
\bottomrule
\end{tabular}
\end{center}

\noindent (\texttt{ax/meta/codomain\_exhaustive} is also the schema's
own fixture for a disjunctive conclusion with an existential; the
three term-kind rejections are representatives — the full mechanical
pairwise family, the 28 literal pairs included, is elaboration-time
conformance rather than schema text.)

\paragraph{The testable negative.} What is deliberately absent is
structural, not policed: no Formula sort and no
Conj/Disj/Exists/Not/Iff nodes (nesting beyond geometric normal form
is unrepresentable); no negation node and no iff node (both are
authoring-side moves, \S\ref{sec:gr-normalize}); no
\texttt{Index}/dependent-type node (erased upstream); no
\texttt{InfDisj} node (\texttt{Sort/immutable} on the binder's sort
carries the distinction); no proof or derivation sorts (proofs of
object logics are rows of object theories — the logical-framework
doctrine); no weight or provenance slot (engine configuration, not
schema); no annotation vocabulary (annotations are ordinary relations
in ordinary extension theories — metadata needs rows, never a schema
revision).

\subsection{The primordial sort tier}
\label{sec:gr-meta-tier}

Thirteen builtin ground sorts are declared as ordinary Sort rows in
the genesis instance, marked immutable, and fixed forever:
\texttt{Q}, \texttt{Z}, \texttt{Nat}, \texttt{B}, \texttt{Str},
\texttt{ENNRat}, \texttt{Shape}, \texttt{Timestamp}, \texttt{Bytes},
\texttt{PeerId}, \texttt{CommitId}, and the genesis-pinned
extension-band straddle pair \texttt{Date} and \texttt{Interval}.
Beside the thirteen sorts stand three sort \emph{formers}, structural
rather than declared — no \texttt{Sort} row, no declaring element;
like the naturals themselves, they are simply available:

\begin{itemize}
\item \textbf{\texttt{Fin}} — the bounded-index former, an indexed
family $\texttt{Nat} \to \textsf{Sort}$: for each natural $n$,
$\texttt{Fin}(n)$ is the $n$-element index sort. These are the sorts
that index tensor axes (an axis of extent $n$ is indexed by
$\texttt{Fin}(n)$); the surface grammar's bounded-index type
(\texttt{$n$ Index}, \S\ref{sec:gr-decls}) expresses the same idea
at authoring scale, though today the normalizer erases it to
\texttt{Nat} plus range obligations (\S\ref{sec:gr-normalize}).
Cell code \texttt{0x09} (\S\ref{sec:gr-cells}).
\item \textbf{\texttt{Tensor}} — the tensor former,
$\texttt{Shape} \times \textsf{Sort} \to \textsf{Sort}$: given a
shape (a \texttt{Shape} value, the list of axis extents) and an
element sort, the sort of tensors of that shape with entries at that
sort. The element sort may be \emph{any} sort, declared or builtin,
not only a numeric one. Canonical form is flattened, and a rank-zero
tensor collapses to its single entry. Cell code \texttt{0x0A}.
\item \textbf{\texttt{List}} — the finite-sequence former,
$\textsf{Sort} \to \textsf{Sort}$, length carried per value:
$[\![\, x\ \texttt{List} \,]\!] = \sum_{n:\mathbb{N}}
\bigl(\texttt{Fin}(n) \to [\![\, x \,]\!]\bigr)$ — each value brings
its own $n$, the \texttt{DynTensor} mechanism with rank pinned at
one. Positional order is the value (no sorting, no deduplication),
and there is no flatness or unit law: \texttt{(x List) List} is a
distinct, ragged two-level sort, never collapsed. Element scope at
this tier is any sort, nesting included; the one adjudicated
exclusion runs the other way — a list instance is not admissible as
a \texttt{Tensor}/\texttt{DynTensor} \emph{element} in v1, while
\texttt{(x Tensor) List} is legal. Former code \texttt{0x10}; the
surface spelling (\S\ref{sec:gr-decls}) narrows the element scope to
builtin ground sorts; value cells ride dtype code \texttt{0x50}
(\S\ref{sec:gr-cells}).
\end{itemize}

\noindent A fourth structural former stands beside these three at the
wire tier, under its own sort-reference tag rather than a former
code: \textbf{\texttt{DynTensor}} — $\textsf{Sort} \to
\textsf{Sort}$, the \emph{dynamic} tensor former, shape carried in
the value rather than the type (the $\Sigma$-over-shapes reading of
\texttt{Tensor}; tag \texttt{0x36},
\path{crates/geolog-envelope/src/former.rs}). \texttt{Tensor} is the
static reading — shape in the type — and the two are distinct sorts
on purpose, never silently interconverted. Of the tensor pair, only
\texttt{DynTensor} has a surface spelling
(\S\ref{sec:gr-decls}); tensor values are side-condition payload
throughout (\S\ref{sec:gr-ground}).

The tier is ambient: every theory may use it — formers
included — in field and codomain positions without an extends edge.
Minting an element of a builtin sort is a load rejection; builtin
values enter as literals in value cells (\S\ref{sec:gr-cells}).

Genesis-tier ground \emph{vocabulary} equips the primordial carriers
with algebraic structure as ordinary declarations (never as
type-former firmware), each family's evaluation semantics pinned in
\path{theories/GeologMeta.geolog} to the corner case:

\begin{itemize}
\item \textbf{$\mathbb{Q}$ field operations} — \texttt{q\_add}, \texttt{q\_neg},
\texttt{q\_sub}, \texttt{q\_mul}, \texttt{q\_inv}, \texttt{q\_div}:
exact arithmetic on canonical rationals, no rounding anywhere;
partiality is refusal (\texttt{q\_inv(0)} and division by zero refuse
— a field has no answer there, and inventing one is what the exact
tier exists to refuse).
\item \textbf{ENNRat quantale operations} — on the Lawvere rationals
$\mathbb{Q}_{\geq 0} \cup \{\infty\}$: \texttt{ennrat\_add},
\texttt{\_min}, \texttt{\_max}, truncated minus (\texttt{\_monus},
the residual of addition), \emph{both} multiplications
(\texttt{\_mul\_conj} with $0 \cdot \infty = 0$ and
\texttt{\_mul\_disj} with $0 \cdot \infty = \infty$ — the two
conventions declared under distinct names, never silently
harmonized), inversion, residuation, harmonic sum, the two
qualitative-mode side-conditions, and exact comparisons.
\item \textbf{Timestamp} — instants, i64 epoch-microseconds on the
POSIX/UTC scale (timezone is presentation, never value); the numeric
order as side-conditions; checked displacement arithmetic
(\texttt{ts\_add\_us}; overflow refuses, never wraps).
\item \textbf{Date and Interval} — PostgreSQL semantics, as a tested
claim: a corpus of goldens captured from a live PostgreSQL 17
instance is the conformance battery for the calendar family
(\S\ref{sec:gr-ground}); intervals are the three-field
months/days/micros value verbatim, never flattened (comparison uses
the justified collapse, values do not); refusal wherever PostgreSQL
refuses, finite-only.
\item \textbf{Str, CommitId} — disequality only
(\texttt{str\_ne}, \texttt{cid\_ne}): string ordering would smuggle a
collation choice (presentation, like timezone), and ordering commit
hashes would invite misreading bytes as causal order.
\end{itemize}

\noindent Every \emph{other} value format — the floating-point ladder
included — is an extension ground format: declared by an ordinary
signed revision (never genesis, never a schema revision), carrying a
parametric spec over ordinary roster rows (byte width, exponent and
mantissa bits, bias, NaN discipline), with wire codes from the
extension band and operation dispatch bound to the declaring
revision. \texttt{ax/meta/builtin\_immutable} quantifies over every
coded sort, so extension formats are immutable by the same axiom that
binds the genesis roster.

\subsection{Theories as rows; revisions}
\label{sec:gr-meta-revisions}

A theory arrives as a \emph{revision}: a sealed batch of meta-schema
rows in one commit body, signed and content-addressed like every
other commit (\S\ref{sec:gr-commits}), dependencies named. Revisions
publish whole — an axiom and all its structure land within one
revision's commit set, which is what makes zero-disjuncts-as-rejection
well-defined. Identity is structural: a revision hash covers the
declaration content (with presentation-only rows, such as
\texttt{Var/name}, excluded), and a function's signature hash covers
its field names, sorts, and codomain — the identity the ground
seam's dispatch pins (\S\ref{sec:gr-ground}).

Shipped machinery (\path{crates/geolog-meta-rows},
\path{crates/geolog-revision}, \path{crates/geolog-theories}): the
theory$\to$rows compiler (byte-identical round-trips on the committed
corpora); the sealed-revision conformance loader, whose rejections
are loud, named, and collected (one committed negative fixture per
live rule); the structural revision and signature hashes with an
adversarial vector battery; declaration indexes and cross-revision
name resolution; the verified ground-format registry loader; and the
programmatic revision-body builder, whose canonical-form invariants
(backward local references, bytewise field order, exactly-once name
rows) are unrepresentable or checked at build time. The
\texttt{geolog rows} verb (\S\ref{sec:gr-cli}) prints the declaration
slice of this encoding for any theory file; the committed listings
under \path{theories/generated/} carry complete worked theories
(the Petri-net theory is 1{,}239 ops); and \texttt{geolog seal}
(\S\ref{sec:gr-cli}) seals a source file through this machinery —
one revision commit carrying the theories as meta-schema rows (its
dependency table names the genesis and the ground-comparisons
golden, \S\ref{sec:gr-goldens}, uniformly — present whether or not
the source uses the comparison roster, so the table is never
input-dependent), plus a seal companion carrying the source and
registry text — deterministically, so re-sealing an unchanged file
is a no-op (verified: the second \texttt{seal} reports
\texttt{store: genesis already present, comparisons already present,
revision already present, seal already present}).

\section{Commits: envelopes, bodies, cells, storage}
\label{sec:gr-commits}

\emph{Authority: \path{crates/geolog-envelope} (the wire),
\path{crates/geolog-store} (the log),
\path{crates/geolog-ceremony} (the goldens). This section states the
shipped encoding — \emph{encoding v1} — at the grain a consumer
needs; the crate module docs carry the byte-exact layouts, and a
committed cross-language reference implementation backs the golden
vectors.}

\subsection{The skeleton: hash and signature surface}
\label{sec:gr-skeleton}

A commit is a two-layer envelope. The \textbf{skeleton} is the
hash/sign surface: a fixed-layout structure carrying

\begin{center}
\small
\begin{tabular}{ll}
\toprule
field & content \\
\midrule
\texttt{version} & the encoding version byte (1) \\
\texttt{author} & ed25519 public key, 32 bytes \\
\texttt{deps} & the commit-reference list: 32-byte commit hashes \\
\texttt{parent\_count} & how many leading deps are parents \\
\texttt{ts\_claim\_us} & claimed timestamp, epoch-µs (a claim, not a truth) \\
\texttt{message} & optional UTF-8 bytes \\
\texttt{body\_hash} & blake3 of the body bytes \\
\bottomrule
\end{tabular}
\end{center}

\noindent \texttt{deps} is one list in two segments — parents first,
then non-parent content dependencies, each segment sorted strictly
bytewise-ascending, duplicates rejected. \emph{Commit identity is
content:} $\texttt{commit\_hash} = \mathrm{blake3}(\text{skeleton
bytes})$, and the ed25519 signature covers exactly those bytes under
a fixed versioned context prefix
(\texttt{geolog-s4-skeleton-sig-v1}), so no future encoding change
can collide signing domains. Canonicality is definitional, in both
directions: the writer canonicalizes; the reader, after structural
validation, re-encodes the decoded value and requires byte equality
with the input (the fixed-point gate) — one value, one byte string,
independent of any serializer freedom. A signed skeleton can only be
constructed against a body receipt obtainable by durably storing the
body, so the format's sharpest failure mode — a signed skeleton whose
body is absent, a permanently unverifiable dangling commit — cannot
arise from the safe API's normal use.

\subsection{The body: ops}
\label{sec:gr-body}

The \textbf{body} is the ops payload. Wire shape:
$\texttt{u16le body\_format} \parallel \text{frame bytes}$; the
format registry has exactly one live entry, and the numbering below
is the public truth for this repository:

\begin{itemize}
\item \texttt{0} — retired before any production bytes existed (a
migration lane for a legacy dialect; no wire ever carried it).
\item \texttt{1} — \texttt{BODY\_FORMAT\_HEXANE\_TEMPLATES\_V1}:
\textbf{the only format, and the default.} The canonical op sequence
is factored into a template dictionary (one row per cluster of alike
ops) plus columnar hole carriers; everything the dictionary hoists is
spelled once per cluster, and each op row pays only for its hole
values. The factoring is fully \emph{derived} from the content —
clustering, per-position carrier modes, and traversal order are laws,
not producer choices — and a frame spelling the same rows any other
way is refused at decode. The frame ends with the same fixed-point
gate the skeleton uses: decode re-encodes and requires byte equality.
$\texttt{body\_hash} = \mathrm{blake3}(\text{body bytes})$.
\end{itemize}

\noindent The op vocabulary:

\begin{center}
\small
\begin{tabular}{lll}
\toprule
op & payload & meaning \\
\midrule
\texttt{elta} & sort reference & mint an element of a (mutable) sort \\
\texttt{rela} & relation ref, argument cells & assert a relation row \\
\texttt{funa} & function ref, argument cells, value cell & assert a function value \\
\texttt{eltd} & target reference & delete an element (never a meta-row target) \\
\texttt{blob} & format tag, opaque bytes & content-addressed bulk (\S\ref{sec:gr-cells}) \\
\texttt{unknown} & kind tag, opaque bytes & forward-verifiability carry \\
\bottomrule
\end{tabular}
\end{center}

\noindent References are either \emph{local} (an earlier op of the
same body, backward only) or \emph{external} (content-addressed,
through the deps); names are data (\texttt{Sort/name} rows), never
identity. An unknown op kind is preserved byte-for-byte and carried,
never dropped and never an error — a reader that predates a kind
still parses correctly, because no decode path may silently drop
data. Decode enforces hostile-input hygiene throughout: claimed
counts are bounded (the op-count ceiling is $2^{20}$ per body) and
bounds-checked before any allocation.

\subsection{Cells and the dtype registry}
\label{sec:gr-cells}

Argument and value positions carry \emph{cells}: element references,
or literals of the primordial dtypes. The primordial band is complete
— \texttt{0x00}–\texttt{0x0F}, no holes:

\begin{center}
\small
\begin{tabular}{llll}
\toprule
code & dtype & code & dtype \\
\midrule
0x00 & \texttt{Q} & 0x08 & \texttt{Shape} \\
0x01 & \texttt{Z} & 0x09 & \texttt{Fin} (former) \\
0x02 & \texttt{Nat} & 0x0A & \texttt{Tensor} (former) \\
0x03 & \texttt{B} & 0x0B & \texttt{Timestamp} \\
0x04 & \texttt{Str} & 0x0C & \texttt{Bytes} \\
0x05 & \texttt{ENNRat} & 0x0D & \texttt{BlobRef} \\
0x06 & \texttt{ZBig} & 0x0E & \texttt{PeerId} \\
0x07 & \texttt{NatBig} & 0x0F & \texttt{CommitId} \\
\bottomrule
\end{tabular}
\end{center}

\noindent \texttt{ZBig}/\texttt{NatBig} are multiprecision
realizations of $\mathbb{Z}$/$\mathbb{N}$ (distinct codes because
widening a code's grammar would be reinterpretation, which the
registry's general law forbids); \texttt{BlobRef} is the out-of-line
spelling of \texttt{Bytes}/\texttt{Str} above the inline threshold (a
35-byte content-address cell); \texttt{CommitId} is a reference — a
tiny index into the enclosing skeleton's deps list — and never a
context-free literal. The extension band \texttt{0x10}–\texttt{0x3F}
carries ground formats (the float ladder \texttt{0x10}–\texttt{0x16};
the genesis-pinned \texttt{Date} \texttt{0x17} and \texttt{Interval}
\texttt{0x18}), decoded opaquely and validated at closure tier; the
absence band \texttt{0x40}–\texttt{0x4F} carries typed absence.
Values of the list former ride dtype code \texttt{0x50} — the first
assignment past the reserved bands, sitting out of numeric order
because the primordial band was already complete and the former
registry's \texttt{0x10} (\S\ref{sec:gr-meta-tier}) falls inside the
extension band — validated at decode like the primordials, in two
first-byte-discriminated arms: an \emph{inline} arm (element dtype,
count, then length-framed canonical element payloads — scalar
elements only: the primordial literal codes minus \texttt{Tensor},
\texttt{BlobRef}, and \texttt{CommitId}) and an \emph{envelope} arm
(a 34-byte multihash content address under a dedicated derive-key
context), partitioned at 64\,KiB in the tensor family's style;
nested lists and non-scalar elements ride the envelope arm, and any
cell code $\geq$ \texttt{0x51} is a decode error.
Over-threshold bulk travels as \texttt{blob} ops (format 0 =
\texttt{blob-raw-v1}) on the same signed-envelope machinery —
signing, deps, sync, and equivocation detection inherited, interior
deliberately opaque.

\subsection{The store log}
\label{sec:gr-store}

The persistent store (\path{crates/geolog-store}) appends envelopes
and bodies to a log file with a fixed grammar:

\begin{Verbatim}[fontsize=\small,frame=leftline,rulecolor=\color{accentmid},framesep=3mm,xleftmargin=3mm]
file    := header record*
header  := b"glog2607" | u32-LE format_version (1) | u32-LE reserved (0)
record  := u32-LE payload_len | u8 kind | payload
kind 1  := ENVELOPE: sig[64] | skeleton bytes
kind 2  := BODY:     raw body bytes (identity = blake3(bytes))
\end{Verbatim}

\noindent There is no per-record checksum, by adjudicated design:
every envelope is signature-verifiable and every body
content-hash-verifiable — stronger than any CRC — and a checksum
cannot repair length-prefix desynchronization anyway. A record length
above the 64\,MiB ceiling is treated as torn or garbage bytes misread
as a length, not as a record. The grammar layer never classifies a
failure as crash artefact versus corruption — that is recovery
policy, one layer up, which consults a committed hint the grammar
deliberately knows nothing about. Header deviations (magic, version,
nonzero reserved field) refuse loudly; a future format is a new
version behind that gate, never a silent reinterpretation.

\subsection{The committed goldens and the ceremony}
\label{sec:gr-goldens}

Four signed golden envelopes ship under \path{theories/}: the
\textbf{genesis} (GeologMeta's own declaration rows as a signed
commit, the root of the graph, $\texttt{deps} = []$, message
\texttt{geolog-meta-genesis-v1}); the \textbf{ground-formats}
revision (the float-format vocabulary, $\texttt{deps} =
[\text{genesis}]$); the \textbf{zones} vocabulary
($\texttt{deps} = [\text{genesis}]$); and the
\textbf{ground-comparisons} roster ($\texttt{deps} =
[\text{genesis}]$, message \texttt{geolog-ground-comparisons-v1}) —
the builtin ground-comparison vocabulary (\texttt{q\_lt},
\texttt{q\_le}, \texttt{q\_ne}, \texttt{nat\_lt}, \texttt{nat\_ne})
as a sealed revision, the second dependency of every revision commit
\texttt{geolog seal} mints (\S\ref{sec:gr-meta-revisions}). Each is a deterministic
function of its committed structure-normative listing and a
\emph{public} seed committed in source: anyone can derive the signing
key and reproduce every byte, and
\texttt{cargo run -q -p geolog-ceremony -- verify --root theories}
does exactly that in every CI pass. The key certifies nothing about
who signed; what makes these goldens canonical is that they are
byte-reproducible from the committed listings by the committed tool.
A real deployment mints its own genesis under its own key; the
ceremony crate is the complete worked example of how. The genesis is
also the meta-schema's homoiconic fixed point — GeologMeta described
as GeologMeta rows — and the generator's \texttt{check} mode
re-derives it from \path{theories/GeologMeta.geolog} both ways as a
drift alarm.

\section{Sync, merge, and equivocation}
\label{sec:gr-sync}

\emph{Authority: \path{crates/geolog-sync} (whose README freezes the
contract: objects, invariants, the frontier exchange, the convergence
theorem, and a ledger of what is and is not yet claimed) and
\path{crates/geolog-envelope/src/equivocation.rs}.}

\subsection{Merge}
\label{sec:gr-sync-merge}

\begin{definition}[merge]
\label{def:gr-merge}
The state of a replica is a set of signed commits (a deps-closed
commit bag). Merging replicas is set union on commits followed by
re-derivation of all derived state. Merge is therefore total,
commutative, associative, and idempotent; it never conflicts
structurally, and admission order does not affect derived state.
\end{definition}

\noindent A merged state may satisfy some axiom's violation premise —
that is a judgment (\S\ref{sec:gr-chase-judgments}), repaired by
further ordinary commits, never by a conflict-resolution protocol.
The Enforced gate binds local commit paths only
(\S\ref{sec:gr-engine}); gating sync admission on validity is the
classic failure this architecture refuses, because it would cost true
merge and history-as-authority. Identity questions inside the merged
state (two writers' records denoting one thing) are the theory's
business, decided by its equality axioms and computed by congruence
(\S\ref{sec:gr-chase-congruence}). The rebuild invariant backs the
order-independence claim mechanically: the repository's battery diffs
incrementally maintained state against from-scratch replay, and a
divergence is a release-blocking defect.

\subsection{Replication}
\label{sec:gr-sync-replication}

Two shipped strata. \textbf{Sync v0 — the file mesh}: whole-file,
N-device mesh sync of per-device, single-writer, append-only envelope
logs plus content-addressed body directories; per-source mirrors and
per-run receipts (``silence means broken''); a torn trailing frame is
classified by comparing against the prior round's observation
(growing source: benign; stable source with a persisting tear: a
problem) — a distinction a single snapshot cannot make.
\textbf{Replay-through-ingestion} is the only path synced envelopes
enter a local view: re-verification through the envelope crate's own
verify path (no new trust surface), idempotent by commit hash,
order-free via a deps-frontier worklist fixpoint.

\textbf{Sync v1 — the connected protocol}, layered strictly above v0:
a deterministic integer-keyed message codec (strict known-kind
decode, unknown-kind skip), frontier-exchange pulls, subscribe push,
acknowledged head updates, credit-based flow control, bounded
staging, and one door: admission through the same replay path.
Equivocation detection runs incrementally during sync and is loud,
never a gate.

\paragraph{Realms and partitions.} Sync is scoped by \emph{realm} —
the wire's opaque slice key, today a plain partition name: sessions
pull and subscribe per realm, and a realm's head set is an exact
summary of its deps-closed store (the crate's \textsc{inv-closure},
which is store-global). A materialized partition must be
dependency-closed and revision-closed (no partition splits a semantic
revision unit); the reference builder and its checker are separate
artefacts, and the checker re-verifies against the store rather than
trusting the builder. Dependencies that cross a realm boundary are
governed by a per-realm \texttt{TransportMode} dial:
\texttt{refused} (the default), \texttt{reported} (the crossing
becomes a declared boundary reference), or \texttt{repaired} (a
locally held foreign dependency's pinned closure joins; a repair that
would need a network fetch is a typed refusal). The \texttt{repaired}
dial also binds at admission: a realm registered \texttt{repaired}
admits entries whose dependencies are admitted under \emph{other}
realms — the dependencies must still be admitted locally, so closure
remains store-global — and the dial is local policy, taken at store
open and never negotiated on the wire; an undialed peer stages the
same entries, a liveness asymmetry rather than a safety one.

\subsection{Equivocation}
\label{sec:gr-sync-equivocation}

\begin{definition}[equivocation-suspicious pair]
\label{def:gr-equivocation}
Two distinct envelopes from one author, neither of which causally
references the other (through \texttt{deps}), form an
equivocation-suspicious pair.
\end{definition}

\noindent Detection is structural — from the bytes alone, no
consensus protocol, no quorum — and local to whatever envelope set is
supplied: a dep pointing outside the set is ``unknown, not
reachable,'' never an error. Two different-bytes envelopes claiming
one commit hash refuse loudly (a genuine hash collision or a hostile
input; never silently resolved by picking a winner). The envelope
crate ships the demonstration as a test — two signed commits claiming
the same chain position, exhibited and caught — and the Tutorial
locates this in the systems landscape: the offence a broker cannot
even state is detected here from the data model. The mechanism,
signed hash-linked histories making forks structurally visible, is
the same one that makes CRDTs Byzantine fault tolerant
\cite{kleppmann22}.

\section{The ground seam}
\label{sec:gr-ground}

\emph{Authority: \path{crates/geolog-ground-api} (the contract) and
\path{crates/geolog-ground} (the reference floor, its first
provider).}

Ground vocabulary — relations and functions over immutable carriers,
decided by procedures rather than stored rows — crosses one seam.
\textbf{Declaration is data}: ground vocabulary is declared where all
vocabulary is declared, as committed theory rows, and resolves to a
pinned identity (declaring commit + row index) with a verified
signature and, where one exists, the commit hash of the committed
evaluation-semantics row. \textbf{Implementation is code}: a
\emph{provider} is one value carrying curated claims plus the
evaluator over its own entries. \textbf{The binder is where they
meet}, and it is engine-owned law no provider can bypass or see:
every resolved declaration is adjudicated against every provider's
claims, with three outcomes per entry —

\begin{center}
\small
\begin{tabular}{ll}
\toprule
status & meaning \\
\midrule
Implemented & exactly one claim verified, committed semantics covered — a dense engine id is allocated \\
Deferred & declared-but-unbound, safe-fail: no claimant, no semantics row, or a stale provider \\
Refused & code and theory disagree (signature mismatch) — the drift alarm, distinct from Deferred \\
\bottomrule
\end{tabular}
\end{center}

\noindent Two verified claims on one identity are always an error;
provider order is never semantics. Dispatch binds to pinned
identities, never names: a same-named impostor declaration gets no
oracle and its atoms defer forever. An unbound entry never obtains a
dense id, so a compiled plan citing it cannot even be expressed —
\emph{compile refusal by construction}, the seam's realization of
refusal-over-guessing.

The laws, stated once here and cited by everything that consumes the
seam: \textbf{purity} (every evaluator is stateless and deterministic
per its declared scope — same cells, same outcome, every process,
every architecture; a family whose reference uses floating point must
pin bit-exact step semantics or cannot be admitted);
\textbf{the lane law} (an accelerated lane may change cost, never
verdicts, and ships with a conformance battery against the always-on
scalar reference — a lane that cannot pass its battery does not
dispatch); \textbf{refusal} (refusal is data, and conservative: under
an Enforced axiom, a gate that cannot evaluate must not admit; no
seam method panics on validated cells); \textbf{no enumeration} (a
ground relation filters, never generates; a ground function inverts
only where the algebra genuinely determines the missing argument —
underdetermined means the engine must not enumerate); and
\textbf{scope} (pure is the norm; the one typed exception is
snapshot-scoped structural relations reading instance rows through an
engine-supplied view, with the routing engine-enforced in both
directions). Cells cross the seam in their canonical encodings,
opaque to the contract; there is no intermediate certificate file
format anywhere — a checker is a ground relation evaluating over rows
the engine already speaks.

The \textbf{reference floor} ships as the seam's first provider:
exact arbitrary-precision rational arithmetic, the
extended-nonnegative-rational quantale operations, calendar
arithmetic with documented, PostgreSQL-conformant semantics (the
golden corpus of \S\ref{sec:gr-meta-tier}), and the curated
ground-comparison family. The floor's reference path is
unconditional; the calendar \emph{functions} sit behind a feature
flag whose absence makes their claims absent — declared entries then
defer at bind, so \textbf{verdicts never change with a feature flag}.
Application side-condition families (the proof-calculus checkers of
the companion repositories) are application content: they ship as a
provider crate owned by the application, composed at build time
through this same binder — the only sanctioned coupling between this
repository and its first big application. Reaching the seam from the
\texttt{geolog} command line is design tense; the engine-level seam,
the floor, and their batteries are shipped and tested today (the
Tutorial's ground chapter walks the same line).

\section{Capabilities}
\label{sec:gr-caps}

\emph{Authority: \path{crates/geolog-caps}.}

Authorization arrives as an algebra over commits, not an access
list beside them. Shipped today, as data and replay semantics:
capability operations (delegations and revocations) travel as
ordinary signed revisions; a normative replay computes which revocations are
effective; authorization judgments are query-time and
position-relative. The committed design decisions, decision by decision (the authority
crate's own phrase):

\begin{itemize}
\item a grade lattice
$\text{Define} < \text{Retract} < \text{Admin}$ per sort — element
deletion is the non-monotone direction, so it is strictly higher
privilege than minting: the add/delete split of the op vocabulary,
lifted;
\item a per-relation \texttt{Assert} chain sharing no element with
the sort chain above bottom, so sort authority and row authority can
never satisfy one another;
\item dual Admins per theory and per instance, with theory-Read as
the base prerequisite to holding anything effective on a theory's
instances;
\item a conservative cross-tier stance: no implication in either
direction; dormant grants are first-class, queryable, and surfaced by
an authoring lint that replay never reads;
\item revocations executed in earliest-add seniority (ascending
authority-provenance depth, digest-tiebroken), with the auth-order
projection carried as data so seniority cannot be padded by
packaging;
\item edge identity as (revision, index), so identical re-grants in
distinct revisions are distinct edges.
\end{itemize}

\noindent \emph{Design
tense:} enforcement at the store boundary — wiring these judgments
into admission of local commits — is the committed direction; today
the algebra, replay, and judgments ship as a library with their
property tests, and no store consults them yet.

\section{The \texttt{geolog} command line}
\label{sec:gr-cli}

\emph{Authority: \path{crates/geolog} — \path{main.rs} for the
surface, the modules cited throughout this chapter for each verb's
semantics. The end-to-end tests in \path{crates/geolog/tests/cli.rs}
run the book's own transcripts against the committed examples.}

\begin{Verbatim}[fontsize=\small,frame=leftline,rulecolor=\color{accentmid},framesep=3mm,xleftmargin=3mm]
geolog check     <file.geolog> [--strict]
geolog check     --store DIR [<commit>] [--strict]
geolog chase     <file.geolog> [--instance A[,B...]]
                               [--budget N] [--mode chased|monitored]
geolog chase     --store DIR [<commit>] [--instance A[,B...]] [--budget N]
geolog seal      <file.geolog> --store DIR [--strict] [--repro]
geolog extract   <commit> --store DIR [--referencing] [-o FILE.glc]
geolog restrict  <file.glc> --store DIR --theory SEL [--revision HEX]
                               [-o FILE.glc]
geolog load      <file.geolog> [--strict] [--dry-run] [-d]
geolog load      <file.glc> [--store DIR]
geolog fmt       <file>
geolog rows      <file.geolog>
geolog normalize <file.geolog> [--strict]
geolog census    <file.geolog>... [--rows LISTING]...
                               [--budget N] [--tsv FILE]
\end{Verbatim}

\paragraph{The specified verbs.}
\textbf{check} runs parse $\to$ normalize $\to$ lint on the theory
text and validates every instance block against the elaborated
closure by actually seeding it — exactly the checks \texttt{chase}
would run, without the chase — reporting each stage and then three
summaries: a \emph{mode census}, what the file's axioms will do when
chased (\texttt{modes: 2 chase · 1 report (0 by the unmarked default)
· 0 refuse} for the Quick Start theory — so an all-\texttt{report}
file, deriving nothing by design or awaiting its
explicit-\texttt{chase} port, is legible at a glance); the move
census (\S\ref{sec:gr-normalize}); and, where the chase bridge
carries the fragment, a chase-termination advisory (weak acyclicity)
computed under the file's own surface modes — only the
\texttt{chase}-marked fragment generates, so a theory of
\texttt{report}-marked constraints is termination-trivial by
construction, however generative its axioms would be if fired. With
\texttt{--store}, the same strata run over a sealed source instead
of a file, after signature verification and a coherence gate: the
sealed source must re-derive the sealed revision rows
byte-identically.
\textbf{chase} seeds the selected instance blocks and runs the chase
to a fixpoint, then renders the report of \S\ref{sec:gr-report}.
\textbf{rows} prints the theory's declaration slice as meta-schema
rows in the committed listings' notation (local refs \texttt{L$n$},
genesis identities \texttt{G<…>}), with a closing note counting the
axioms that continue as rows beyond the slice; it requires the
one-field record spelling for unary relations, because fields are
rows with names in the meta-schema (verified: \texttt{relation
`Marked` has a unary (anonymous) domain — … rows emission needs the
one-field record spelling: `Marked : \{x: A\} -> …`}).

\paragraph{The remaining verbs, at citation grain.} Each is shipped
and tested; this manual states its role and leaves its full
semantics to the named authority. \textbf{seal} runs parse $\to$
normalize $\to$ canonical revision rows $\to$ a signed envelope pair
in the store: a revision commit and a seal companion
(\S\ref{sec:gr-meta-revisions}); re-sealing an unchanged file is a
no-op. With \texttt{--repro} the signature uses the public
reproduction seed, making the commit hashes store-independent (pin
reproducibility); such commits carry no authorship and must never
anchor authority. \textbf{extract} writes a commit and its transitive deps from
the store to a \path{.glc} container (magic \texttt{geologfl},
self-contained, verified on read rather than trusted);
\texttt{--referencing} additionally sweeps every store commit whose
deps reach the target, so all commits referencing an instance travel
together. \textbf{restrict} is theory-fragment restriction of a
certificate container — the get side of the delta lens a
\texttt{theory B extends A} induces: it keeps exactly the rows whose
species live in the target theory's extends-closure (derived from
the sealed revision's own declaration rows), projects the
prerequisite table onto the surviving boundary, copies unchanged
commits byte-verbatim, re-seals changed ones deterministically with
the public reproduction key, and refuses — with a named locus — any
container whose kept rows reference stripped content; a
pure-fragment container round-trips byte-identical. \textbf{load}
with a \path{.glc} argument decodes and
verifies the whole container — signatures, body hashes, topological
order — then admits it into \texttt{--store}; with a
\path{.geolog} argument it stages the theory's aggregation surface
into a live engine (\S\ref{sec:gr-editions}). \textbf{fmt} parses a
file, legacy spellings included, and prints the canonical source
respelling — the corpus-port instrument (\S\ref{sec:gr-decls}).
\textbf{normalize} prints the geometric normal form the engine
consumes (\S\ref{sec:gr-gnf}) — the committed-golden regeneration
verb. \textbf{census} is the merge-reliance census over authored
equality heads; its method and honesty ledger live in
\path{crates/geolog/src/census.rs}. Authority for all of these:
\path{crates/geolog/src/main.rs} and the modules it names.

\paragraph{Options.} Classification travels in the source file
(\S\ref{sec:gr-registry}), so the file-mode verbs take the file and
nothing else; \texttt{--registry FILE}, the deprecated migration-era
carrier, refuses beside import lines when its registry carries
classification (a pattern-only registry composes, its aux names
kept clear of imported ones — \S\ref{sec:gr-registry}), and in
store mode the flag is refused outright: the registry text travels
sealed.
\texttt{--strict} rejects the deprecated spellings — bracket
records, binders-before-name theory headers, \texttt{=~\{} instance
bodies, \texttt{:~Sort} declarations — each with a pointer at the
ratified form (\S\ref{sec:gr-decls}) — and demands the
\texttt{import std} header for builtin ground vocabulary used by
name (\S\ref{sec:gr-registry}). \texttt{--instance A[,B…]}
selects blocks by name — the default is every block in the file, as
their union (\S\ref{sec:gr-instances}); naming a missing block is a
usage error that lists the blocks present. \texttt{--budget N} bounds
chase iterations (default 1000); exhaustion is a reported judgment,
not an error. By default \texttt{chase} runs each axiom under its
authored surface mode (\S\ref{sec:gr-chase-modes});
\texttt{--mode chased|monitored} is the legacy run-level flattener
for files whose axioms carry no marks, and refuses — naming the
marked axioms — on any file that does. \texttt{chase} on a file with
no instance block is a usage error that says what to add and points
at \texttt{check} for theory-only validation.

\paragraph{Exit codes.} \texttt{0} — clean: the verb completed and,
for \texttt{chase}, the fixpoint carries no judgments. \texttt{1} —
judgments at fixpoint: the run completed and reported violated
denials, pending branch choices, or an exhausted budget. \texttt{2} —
the input never ran: usage errors, unreadable files, pipeline
rejections, bridge refusals, seeding contradictions, and engine load
refusals. The 1/2 boundary is the error type's own doctrine:
judgments from a completed chase are results, not errors.

\paragraph{Constituent binaries.} The workspace also ships
development CLIs for constituent crates — \texttt{glnorm} (the
normalizer pipeline standalone), \texttt{glstore} (the store log),
\texttt{glsync} (the file-mesh sync) — documented in their crates;
\texttt{geolog} is the carryable verb and the only surface this
chapter specifies.

\section{Editions: what this manual does not yet state}
\label{sec:gr-editions}

Open seams, each named with where its truth lives today; a future
edition of this chapter absorbs them in this order.

\begin{itemize}
\item \textbf{The aggregation surface.} The ratified authoring
grammar beyond the fragment this manual states — \texttt{law}
members (named or anonymous) whose bodies carry contraction index
groups and string-literal contraction specs, the arity-typed
contraction resolver with its refusal roster
(\texttt{E-RANGE}, \texttt{E-SPEC-*}, \texttt{E-RPN-*},
\texttt{E-LAW-*}), and the free-construction instance declarations
(parameters, \texttt{extends} targets, and body rows as three seed
sources) — is shipped in \path{crates/geolog-normalize} (the parser
module's documentation is the compact normative summary) and staged
into a live engine by \texttt{geolog load}, whose emission header
marks the engine realization design-tier pending the
aggregation-schema promotion. This edition states the surface's
lexical floor and type-grammar hooks (\S\ref{sec:gr-lexical},
\S\ref{sec:gr-decls}) and defers the rest to the next.
\item \textbf{The store-backed surface at this manual's grain.}
\texttt{seal}, \texttt{extract}, \texttt{load} of a \path{.glc}
container, the store modes of \texttt{check} and \texttt{chase}, and
\texttt{census} all ship today (\S\ref{sec:gr-cli}); the container
grammar, the store-mode conformance rules, and the census method are
normative in the facade's modules. A future edition states them at
the grain \S\ref{sec:gr-commits} gives the envelope.
\item \textbf{The byte-exact wire grammar.} \S\ref{sec:gr-commits}
states the encoding at consumer grain; the op-payload grammar, the
template-frame channel layout, and the conformance-rule table
(load-rejection classes and their fixtures) are normative today in
\path{crates/geolog-envelope}'s module documentation and
\path{crates/geolog-meta-rows}' loader and negative fixtures.
\item \textbf{Per-commit validity, end to end.} The judgment
assembly ships (\S\ref{sec:gr-engine}); the storage-pipeline wiring
and its report vocabulary are design tense.
\item \textbf{Capability enforcement} at the store boundary
(\S\ref{sec:gr-caps}).
\item \textbf{The ground seam from the command line}
(\S\ref{sec:gr-ground}): the chase-path elaboration of ground sorts
remains refused (\S\ref{sec:gr-elab}), and the \texttt{load} path
stages its ground-vocabulary dependencies as roster rows while its
engine realization stays design-tier (previous bullet).
\item \textbf{Cross-revision \texttt{extends} in the facade}, native
bounded-index types, and the full evaluation-semantics tables of the
genesis ground families (today: pinned, with their corner cases and
golden corpora, in \path{theories/GeologMeta.geolog} and the ground
crates' tests).
\end{itemize}

 \chapter{Diagrams as Data: S-Complexes}
\label{ch:qtsl-scomplexes}

\noindent\draftmark\medskip

Several strata of this book traffic in \emph{diagrams}. The
specification language's global stratum asserts a certified Morse graph
(\bookref{ch:qtsl-specification-language}{the Specification Language chapter}); the hypernet
chapter reads every computation as a string diagram
(\bookref{ch:qtsl-hypernets}{the hypernet chapter}); grid
certificates decompose state space into cubical cells; and the
higher-categorical diagram machinery this project consults runs on
oriented graded posets. Each of these is, concretely, a family of
cells in graded dimensions plus incidence data saying which cells
bound which. Each could therefore grow its own bespoke incidence
schema, its own well-formedness checker, and its own subtly different
answer to ``is the boundary recorded here the boundary the theorem
was about?''

The toolkit owns one such object. An \emph{S-complex} is a graded
family of finite cell sets whose boundary structure is recorded as a
pair of semiring-valued tensors per dimension (one for each
orientation sign), satisfying a subtraction-free form of the chain
condition $\partial^2 = 0$. The semiring parameter $\mathbb{S}$ is where
the genus gets its reach: with Booleans it is incidence, with natural
numbers it is multiplicity, with a semiring of \emph{port names} it
is a string diagram, and with min-plus dwell bounds its labels carry
quantitative dynamics. Well-formedness in every case amounts to a
handful of tensor contractions and one semiring comparison per
dimension, using the same contraction machinery the rest of the
toolkit already trusts (\bookref{ch:qtsl-hypernets}{the hypernet
chapter}, relational-kinship section), so the properties are
checkable in bulk, on accelerators, by the engine that checks
everything else.

The instinct (diagrams stored as typed incidence data, computed on
in bulk) is shared with the attributed $\mathcal{C}$-set line of
applied category theory, where diagrams are categorical data
structures~\cite{patterson21} with data-parallel composition
algorithms~\cite{wilson23}, and where a string diagram can be a
first-class engineering artefact rather than an
illustration~\cite{master19}. An S-complex trades that line's schema
generality for one graded shape with a semiring parameter; the trade
is what makes well-formedness a tensor computation.

This chapter defines the object and its morphisms
(section~\ref{sec:scx-def}), tours the instances that choice of
$\mathbb{S}$ buys (section~\ref{sec:scx-instances}), states the
encoding discipline that decides when a complex is stored versus
derived (section~\ref{sec:scx-generators}), presents the certified
Morse graph as the flagship rank-1 instance with its quantitative
refinement (section~\ref{sec:scx-morse}), works the string-diagram
instance in full --- including the reading under which a well-formed
hypernet \emph{is} an S-complex (section~\ref{sec:scx-strings}) ---
and closes with a ledger separating definitional content from design
directions (section~\ref{sec:scx-ledger}).
Sections~\ref{sec:scx-def} and~\ref{sec:scx-morse} carry everything
the rest of the book consumes; section~\ref{sec:scx-strings} is the
most likely to surprise, and a reader with no stake in string
diagrams can skip it without loss of continuity.

\section{The definition}
\label{sec:scx-def}

Throughout, $\mathbb{S}$ is an \emph{ordered semiring}: a semiring
$(\mathbb{S}, 0, 1, +, \cdot)$ equipped with a preorder $\le$ such
that $0 \le 1$ and both operations are monotone in each argument.
Order-equivalence is derived, not assumed:
$x \cong y \mathrel{:\Leftrightarrow} x \le y \wedge y \le x$.
No subtraction is ever required. That restraint is the point: the
instances worth having --- Booleans, counts, port sets, min-plus
costs --- live in semirings that have no subtraction to offer.

A \emph{tensor} over $\mathbb{S}$ indexed by finite sets
$A_0, \dots, A_{k-1}$ is a finite function
$\prod_{i<k} A_i \to \mathbb{S}$, and juxtaposition of tensors
denotes contraction over the shared index (sum of products, in
$\mathbb{S}$), as in the contraction vocabulary of
\bookref{ch:qtsl-mathematical-models}{the modelling chapter}.

\begin{definition}[S-complex]
\label{def:scomplex}
An $\mathbb{S}$-\emph{complex} of rank $r$ consists of finite cell
sets $C_0, C_1, \dots, C_r$ (cells in dimension $0$ through $r$)
together with, for each $0 < i \le r$ and each sign
$e \in \{+,-\}$, a rank-2 boundary tensor
\[
\partial_i^{e} \colon C_{i} \times C_{i-1} \longrightarrow \mathbb{S},
\]
satisfying, for each $1 \le i \le r-1$, the \emph{sign-split chain
condition}
\begin{equation}
\label{eq:scx-dd}
\partial_{i}^{+}\, \partial_{i+1}^{+} \;+\; \partial_{i}^{-}\, \partial_{i+1}^{-}
\;\cong\;
\partial_{i}^{-}\, \partial_{i+1}^{+} \;+\; \partial_{i}^{+}\, \partial_{i+1}^{-}
\qquad \colon\; C_{i+1} \times C_{i-1} \to \mathbb{S}.
\end{equation}
\end{definition}

Stated as one dependent type, with $[N]$ the finite ordinals and
$\mathsf{Fin}$ the concretely finite sets, the presented data of the
genus is
\[
\sum_{N \colon \mathbb{N}} \;\; \sum_{C \colon \prod_{n \colon [N]} \mathsf{Fin}} \;\; \prod_{i \colon [N-1]} \; \prod_{e \colon \{+,-\}} \big[\, \partial_{i+1}^{e} \colon C_{i+1} \times C_{i} \to \mathbb{S} \,\big]
\]
--- $N$ graded levels and a sign-split boundary pair per consecutive
pair of levels, nothing else. The type records presentations only;
condition~\eqref{eq:scx-dd} is a property checked \emph{of} a
presentation rather than a component packaged \emph{in} it, so one
schema serves both the complexes that satisfy it and the raw
presentations a checker has not yet admitted. The rank $r$ of
Definition~\ref{def:scomplex} is $N - 1$.

Condition~\eqref{eq:scx-dd} is the classical $\partial^2 = 0$ with
the subtraction moved across the equals sign: over a ring, setting
$\partial = \partial^{+} - \partial^{-}$ and expanding
$\partial_i\,\partial_{i+1} = 0$ gives precisely ``matched sign pairs
equal mismatched sign pairs.'' Stated this way the condition makes
sense over any ordered semiring, and each side is two contractions
and a sum, so checking well-formedness of a presented complex is a
fixed, small pattern of contraction rows per dimension --- the
bulk-checkable shape the engine is built around. An elementary
instance to hold while reading: a filled square. Each corner lies on
two of the four edges, so exactly two routes descend from the
$2$-cell through an edge to that corner; however the edges are
oriented, one route carries matched signs and the other mismatched,
so the two sides of~\eqref{eq:scx-dd} agree entry by entry --- the
corner-cancellation of $\partial^2 = 0$, surviving with parity where
its minus signs used to be. For $r \le 1$ the
condition is vacuous, and a presentation of a rank-1 complex should
say so plainly rather than imply a check that never ran.

The two signs record \emph{combinatorial orientation}: for a directed
edge $a$ from $x$ to $y$, the source incidence lives in
$\partial_1^{-}(a, x)$ and the target incidence in
$\partial_1^{+}(a, y)$. This sign records a property of the diagram
itself rather than of any dynamics drawn on it; when a dynamical
certificate needs a notion of inward and outward flow across a face,
that is separate label data riding on the cells, never a
reinterpretation of the boundary signs.

\begin{definition}[morphism of S-complexes]
\label{def:scx-morphism}
Let $C$ be an $\mathbb{S}$-complex and $D$ a $\mathbb{T}$-complex
with $\operatorname{rank}(C) \le \operatorname{rank}(D)$. A
\emph{morphism} $f \colon C \to D$ consists of a function
$f_i \colon C_i \to D_i$ for each dimension $i$ of $C$, regarded
semantically as a tensor
$f_i \colon C_i \times D_i \to \mathbb{T}$ whose entries are $0$
and $1$, together with a lax
morphism of ordered semirings
$f_S \colon \mathbb{S} \to \mathbb{T}$ (monotone, $f_S(0) \cong 0$,
$f_S(1) \cong 1$, $f_S(x+y) \le f_S(x) + f_S(y)$, and
$f_S(x \cdot y) \le f_S(x) \cdot f_S(y)$), such that for every
dimension $0 < i \le \operatorname{rank}(C)$, every sign $e$, and
every $(c, d) \in C_i \times D_{i-1}$,
\begin{equation}
\label{eq:scx-chainmap}
\sum_{d' \in D_i} f_i(c, d')\; \partial_i^{D,e}(d', d)
\;=\;
\sum_{c' \in C_{i-1}} f_S\!\big(\partial_i^{C,e}(c, c')\big)\; f_{i-1}(c', d).
\end{equation}
\end{definition}

Equation~\eqref{eq:scx-chainmap} is the chain-map square, sign by
sign: mapping a cell forward and taking its boundary in $D$ agrees
with taking its boundary in $C$, pushing the cells forward, and
translating the coefficient values through $f_S$. Laxity in $f_S$ is
deliberate. The useful coefficient translations are
\emph{forgetful} --- cardinality (from port witnesses to counts),
support (from counts to incidence), and coarse quantitative summaries
(from fine margins to dwell bounds) --- and translations of the last
kind are subadditive rather than exact, so demanding exactness would
refuse precisely the summaries the tower is for.

\section{Instances by choice of semiring}
\label{sec:scx-instances}

The definition specializes by choice of coefficient semiring:

\begin{center}
\begin{tabular}{@{}llll@{}}
\toprule
object & $\mathbb{S}$ & rank & boundary reading \\
\midrule
directed graph & $\mathbf{2}$ (Booleans) & 1 & $\partial_1^{-}$ source, $\partial_1^{+}$ target \\
directed hypergraph & $\mathbb{N}$ & 1 & tail/head multiplicities \\
cubical complex & $\mathbf{2}$ & $d$ & oriented facet incidence \\
oriented graded poset & $\mathbf{2}$ & any & covering relation, one sign each \\
Lefschetz complex & a ring's split halves & any & incidence coefficients \\
string diagram & port names (\S\ref{sec:scx-strings}) & 1 & \emph{which ports}, not just how many \\
\bottomrule
\end{tabular}
\end{center}

\begin{remark}[oriented graded posets]
\label{rem:scx-ogp}
Oriented graded posets, a combinatorial substrate of
higher-categorical diagram rewriting \cite{hadzihasanovic24a}, embed
as Boolean complexes in which each covering pair carries exactly one
of the two signs. The shape predicates that theory cares about
(roundness, molecule-hood, regularity) remain predicates on top of
the complex, as that literature states them, with the complex
supplying the data they are predicates \emph{of}.
\end{remark}

\begin{remark}[Lefschetz complexes]
\label{rem:scx-lefschetz}
Lefschetz complexes, a combinatorial carrier of Conley-theoretic
invariants \cite{mischaikowmrozek2002}, carry ring coefficients
without a ring: a classical incidence coefficient enters as its
split halves $k = k^{+} - k^{-}$, one sign each, so entries stay in
$\mathbb{N}$ and subtraction never enters the schema.
\end{remark}

\begin{remark}[the coefficient tower]
\label{rem:scx-tower}
The coefficient levels form a tower. Cardinality maps a string
diagram's witness sets to a directed hypergraph's multiplicities, and
Boolean support forgets those to bare incidence. Both maps are lax
morphisms in the sense of Definition~\ref{def:scx-morphism} with the
identity on cells, so the port, counting, and incidence readings are
three resolutions of a single object.
\end{remark}

\section{Generators or presentations}
\label{sec:scx-generators}

Not every complex the toolkit reasons about should be stored as one.
The encoding discipline distinguishes two postures.

\begin{itemize}
\item \textbf{Generator-encoded} where the complex is canonically
  derivable: grid certificates store grid lines and cell selections,
  never cell lists or incidence tensors --- the cubical complex, with
  its orientation-split boundary, is recomputed from the generators
  by every consumer. Skew between a stored presentation and its
  generators is then unrepresentable rather than checked.
\item \textbf{S-complex-encoded} where the boundary \emph{is} the
  content: a Morse graph's edges, a reduced complex, a connection
  matrix. Here there is no generator to fall back on --- the
  incidence data is precisely what the certificate asserts --- so the
  presentation of Definition~\ref{def:scomplex} is the wire object,
  and condition~\eqref{eq:scx-dd} (where non-vacuous) is part of its
  well-formedness check.
\end{itemize}

The two postures compose: a quotient from a generator-encoded complex
onto an S-complex-encoded one is a morphism in the sense of
Definition~\ref{def:scx-morphism}, with the derived presentation
appearing only transiently on the checking side.

\section{The Morse graph, quantitatively}
\label{sec:scx-morse}

The global stratum's payload
(\bookref{ch:qtsl-specification-language}{the Specification Language chapter}, QMorse row) is
a certified Morse graph: finitely many blocks, finitely many directed
connections. As data it ships in adjacency dialect (a vertex count
and an edge list), and that dialect \emph{is} the rank-1 Boolean
S-complex: $C_0$ the Morse sets, $C_1$ the connections,
$\partial_1^{-}$ source, $\partial_1^{+}$ target. Acyclicity is the
checker's to confirm; the confirmation, in this vocabulary, is an
order witness on $C_0$ of the same kind
section~\ref{sec:scx-strings} uses for string diagrams.

The quantitative refinement adds a \emph{labelling} without changing
the complex: cells carry values in the ordered semiring
$\mathrm{DWELL} = ([0, \infty], \min, +)$ --- the cost quantale of
Lawvere's generalized metric spaces~\cite{lawvere73}, whose additive
unit is $\infty$ and whose multiplicative unit is $0$, ordered by
the reverse of the numerical order, which is what makes $0 \le 1$
hold --- with
connections labelled by certified lower bounds on transit time and
Morse sets labelled by their sojourn bounds, which for genuine Morse
sets are $\infty$. The semiring arithmetic then guards the composite
claims against overstatement: min-plus contraction of a path through
an intermediate Morse set yields $\infty$, so the algebra itself
reports that no finite total-time bound through a recurrent set is
being asserted, while per-leg bounds stay finite and usable. This is
the intended serialization form for the quantitative Morse tier: the
graph presented as in Definition~\ref{def:scomplex}, the rates as
cell labels (extended nonnegative rationals in the machine
artefact), and the projection from the underlying grid complex
onto the Morse graph stated as a morphism --- cells to their Morse
or transit class, fine certified margins translated laxly to dwell
bounds --- rather than as an unstated convention between two
schemas.

\section{String diagrams: the port semiring}
\label{sec:scx-strings}

This instance shows why the coefficients form a \emph{semiring
parameter} rather than a Boolean convention. Fix a supply
$\mathcal{P}$ of \emph{port names} --- opaque identifiers, meaningful
only up to renaming. Let $\mathbb{S}_{\mathrm{port}}$ be the semiring
of finite \emph{multisets} of finite sequences over $\mathcal{P}$:
addition is multiset union with zero $\varnothing$, multiplication is
elementwise concatenation with unit $\{\varepsilon\}$, and the order
compares cardinalities --- so the derived order-equivalence is blind
to the names themselves, which is the renaming-invariance the names'
opacity demands.\footnote{Multisets rather than sets, for a reason
worth checking: concatenation can merge distinct pairs of sequences
($\langle 1\rangle \cdot \langle 2,3\rangle$ and
$\langle 1,2\rangle \cdot \langle 3\rangle$ both concatenate to
$\langle 1,2,3\rangle$), so two \emph{sets} of equal cardinality can
have products of different cardinalities, and multiplication would
fail the monotonicity section~\ref{sec:scx-def} demands of every
coefficient semiring. Multisets count collisions instead of merging
them, which makes cardinality exactly additive and multiplicative.
The witness sets arising in diagrams are unaffected: distinct ports
give distinct sequences.}

A \emph{string diagram} over a signature is a rank-1
$\mathbb{S}_{\mathrm{port}}$-complex --- $C_0$ the wires, $C_1$ the
boxes --- together with a morphism into the complex presenting the
signature ($C_0$ the sorts, $C_1$ the generators, each generator's
canonical ports as its boundary values). A boundary value
$\partial_1^{-}(a, x)$ is a \emph{witness set}: one singleton
sequence per input port of box $a$ attached to wire $x$, and dually
for outputs. Over $\mathbf{2}$ an edge could only assert incidence;
over $\mathbb{S}_{\mathrm{port}}$ the value says \emph{what flows
where} --- the difference between a reachability structure and a
diagram whose wires carry content comes down, in this vocabulary, to
the choice of semiring. The signature morphism does the typing: its
cell maps send boxes to generators and wires to sorts, and its
coefficient translation $f_S$ matches each diagram port to the
generator port it instantiates, so argument order lives in the
signature once instead of in every diagram. Sequence concatenation
earns its place under composition, where contracted products
accumulate port \emph{paths} witnessing composite incidences; at
rank 1 only singletons occur. Figure~\ref{fig:scx-tower} works a
two-box example and reads its boundary data at the three coefficient
resolutions of Remark~\ref{rem:scx-tower}.

\begin{figure}[t]
\centering
\begin{tikzpicture}[x=0.94cm,y=0.82cm,baseline=(current bounding box.center),
  box/.style={draw=accentdeep, fill=accent!50, rounded corners=1pt,
    font=\footnotesize\ttfamily, inner sep=2.5pt, minimum height=4.2mm},
  wire/.style={-, line width=0.5pt, draw=black!75},
  port/.style={font=\small\itshape, inner sep=1pt},
  plab/.style={font=\scriptsize, inner sep=1pt, text=black!60}]
\node[port] (x) at (0, 0.75) {$x$};
\node[box]  (a) at (1.4, 0.75) {$a$};
\node[box]  (b) at (4.4, 0.75) {$b$};
\node[port] (z) at (5.8, 0.75) {$z$};
\coordinate (split) at (2.5, 0.75);
\draw[wire] (x) -- (a);
\draw[wire] (a.east) -- (split);
\draw[wire] (split) .. controls (3.1,0.75) and (3.5,1.15) ..
  ([yshift=1.6mm]b.west);
\draw[wire] (split) .. controls (3.1,0.75) and (3.5,0.35) ..
  ([yshift=-1.6mm]b.west);
\draw[wire] (b) -- (z);
\node[plab, anchor=south] at (2.15, 0.85) {$y$};
\node[plab, anchor=south] at ([xshift=-2.6mm,yshift=1.5mm]b.west) {1};
\node[plab, anchor=north] at ([xshift=-2.6mm,yshift=-1.5mm]b.west) {2};
\end{tikzpicture}
\hspace{1.2em}
\begin{tabular}{@{}lccc@{}}
\toprule
entry & $\mathbb{S}_{\mathrm{port}}$ & $\mathbb{N}$ & $\mathbf{2}$ \\
\midrule
$\partial_1^{-}(a,x)$ & $\{\langle 1\rangle\}$ & $1$ & $1$ \\
$\partial_1^{+}(a,y)$ & $\{\langle 1\rangle\}$ & $1$ & $1$ \\
$\partial_1^{-}(b,y)$ & $\{\langle 1\rangle,\langle 2\rangle\}$ & $2$ & $1$ \\
$\partial_1^{+}(b,z)$ & $\{\langle 1\rangle\}$ & $1$ & $1$ \\
\bottomrule
\end{tabular}
\caption{One diagram, three coefficient resolutions. Wire $y$ feeds
both input ports of box $b$, so the witness set
$\partial_1^{-}(b,y)$ holds two singleton port sequences; cardinality
forgets \emph{which ports} to \emph{how many}, and Boolean support
forgets that to \emph{whether at all} (Remark~\ref{rem:scx-tower}).
Entries not shown are $\varnothing$, $0$, $0$.}
\label{fig:scx-tower}
\end{figure}

Structural conditions on the diagram are then one clause each, and
free-structure theorems select which clauses to impose:

\begin{center}
\begin{tabular}{@{}lll@{}}
\toprule
producer/consumer discipline & acyclicity & free structure presented \\
\midrule
none & none & hypergraph category \cite{bonchi20} \\
$\le 1$ producer, $\le 1$ consumer & none & traced symmetric monoidal \cite{kissinger13} \\
$\le 1$ producer, $\le 1$ consumer & witness & symmetric monoidal \cite{bonchi22} \\
$\le 1$ producer, consumers free & witness & gs-monoidal \cite{fritzliang22} \\
\bottomrule
\end{tabular}
\end{center}

Here ``$\le 1$ producer'' is the port-counting condition
$\sum_{a \in C_1} \lvert \partial_1^{+}(a, x) \rvert \le 1$ for every
wire $x$, and dually for consumers --- the cardinality order makes
counting a first-class operation. The acyclicity witness is a
function $t \colon C_0 \to P$ into a poset with
$t(\text{input wire}) < t(\text{output wire})$ across every box: a
one-pass check and a rewrite-friendly certificate.

\subsection{Hypernets read as S-complexes}
\label{sec:scx-hypernets}

The hypernet of \bookref{ch:qtsl-hypernets}{the hypernet chapter}
sits in the third row of the table: its linearity condition is
producer-and-consumer monogamy and its static-single-assignment
order is an acyclicity witness, as that chapter already observes in
citing the free-structure correspondence. The reading assembles as
follows. Wires are the 0-cells; applications are the 1-cells; an
application's ordered input and output lists are the two boundary
rows, with list position serving as the port name --- the
representation takes the renaming quotient so the semiring does not
have to. The catalogue plays the signature: the morphism sends each
application to its box, each wire to its tensor type, and its
coefficient translation records which argument slot each wire fills.
Two strata of the hypernet deliberately live \emph{beside} the
complex rather than in it: box payloads (constants, contraction
patterns, polynomial coefficients) are cell labels, just as rates
are labels on the Morse graph; and the interface is a cospan of
morphisms from rank-0 complexes --- an ordered list of graph inputs
and one of graph outputs, each simply a function into $C_0$, which
is what the declared interface lists already are.

The vocabulary also makes the design space around the interface
legible in single clauses. Letting the output leg of the cospan hit
any wire --- including one an application already consumes ---
declares a read-out port without touching the inner complex or its
monogamy, since the legs were never part of either. If consumer-side
monogamy is retired while the producer side is kept, the diagram
genus moves one row down the table, to the gs-monoidal discipline
of every static-single-assignment compiler representation --- def
once, use many --- with copy and discard absorbed into the wiring.
Both are stated here as readings the definition supports rather than
as changes made: the shipped hypernet remains the monogamous, acyclic
row, and its bytes, digests, and substitution laws are those of
\bookref{ch:qtsl-hypernets}{the hypernet chapter}, unchanged.

\section{Ledger}
\label{sec:scx-ledger}

What in this chapter is what:

\begin{itemize}
\item \textbf{Definitional.} Definitions~\ref{def:scomplex}
  and~\ref{def:scx-morphism}, the derivation of~\eqref{eq:scx-dd}
  from $\partial^2 = 0$, the instance table, and the free-structure
  table (whose correspondences are the cited literature's theorems,
  not this book's).
\item \textbf{Machine artefact.} The genus is declared as a
  standard-library theory of the geolog engine
  (\path{theories/SComplex.geolog}), which names this chapter as its
  expository companion, and the chain condition's contraction
  factoring is generated from a surface statement of the tensor law
  (\path{theories/SComplexGenus.geolog}). The theory is staged
  design vocabulary, exercised by the engine's tests rather than by
  production data.
\item \textbf{Shipped elsewhere in this book.} The Morse-graph
  payload shape and its checker obligations
  (\bookref{ch:qtsl-specification-language}{the Specification Language chapter}); the
  hypernet IR, whose linearity and ordering conditions this
  chapter's section~\ref{sec:scx-hypernets} re-reads
  (\bookref{ch:qtsl-hypernets}{the hypernet chapter}); the
  contraction machinery that makes
  condition~\eqref{eq:scx-dd} bulk-checkable
  (\bookref{ch:qtsl-mathematical-models}{the modelling chapter}).
\item \textbf{Design tense.} The quantitative Morse serialization in
  complex-plus-labels form and its quotient-as-morphism statement
  (section~\ref{sec:scx-morse}); any S-complex-native storage for
  hypernets (section~\ref{sec:scx-hypernets} is a reading, and the
  canonical bytes are unchanged by it); mechanized S-complex
  morphisms, of which the toolkit currently has none --- the Morse
  quotient is the natural first target. Higher-dimensional
  rewriting on this substrate (equivalence cells over diagrams) is
  out of this chapter's scope entirely.
\end{itemize}

\part[QTSL: Quantitative Temporal Specification Logic]{QTSL:\\ Quantitative Temporal\\ Specification Logic}
\chapter{Quick Start}
\label{ch:qtsl-quickstart}

In this chapter you will produce a machine-checkable certificate of a
quantitative claim about a stochastic system, in milliseconds, and see
exactly what any checker must re-derive from its bytes. You will meet
the three routes by which certificate bytes get checked in this
toolkit --- and where each route stands at this preview.
Then you will break a certificate from the published corpus and watch
an exact-arithmetic reader refuse it, with the violated obligation
named. The whole loop fits in one sitting; every command is
copy-pasteable, and every output shown is the verbatim transcript of
a real run --- regenerated by the book's own build, which executes
the printed commands and records the tool revisions and run date
beside each transcript in the source (one commodity arm64 box; your
timings will differ, your bytes will not).

\section{The claim}
\label{sec:qs-claim}

FrozenLake is Gymnasium's introductory gridworld: a $4\times4$ frozen
lake, sixteen states, four actions, a start square, four holes, one
goal. In the \emph{slippery} variant each action moves as intended
with probability $\tfrac13$ and slips to a perpendicular neighbour
otherwise, so every trajectory is a random walk biased by the policy,
and every reach probability is a rational number whose denominator is
a power of three.

The claim to certify: \emph{the recorded policy reaches the goal
within ten steps with probability at least}
\[
  p \;=\; \frac{815}{19683}
  \qquad\text{(exactly; } 19683 = 3^9\text{).}
\]
In QTSL's vocabulary this is a probabilistic deadline claim,
$\mathbb{P}_{\ge p}\!\left[\,\mathrm{F}_{\le 10}\;\mathit{goal}\,\right]$
--- the same claim species whose full campaign (deadline ladders,
dyadic rounding lanes, a family of environments) the Experiments
chapter (chapter~\ref{ch:experiments}) re-derives from committed
artefacts.
One disclosure before we start: $815/19683$ happens to be the exact
optimum of the ten-step problem, and the producer knows it (it runs an
independent floating-point sweep as a cross-instrument), but the
certificate below asserts only the lower bound. What a certificate
\emph{checks} is always narrower than what its producer \emph{knows};
keeping that line crisp is most of the craft, and the rest of this
book is about the vocabulary for widening what is checked.

\section{Setup}
\label{sec:qs-setup}

Three repositories, cloned side by side. \path{maitria-coalition/mtk}
is the producer toolkit and the only thing you build: a Rust
toolchain is its one hard prerequisite, and a GPU is optional
throughout (the accelerated lanes load CUDA dynamically and fall back
to the same exact host arithmetic --- same script, same answers).
\path{maitria-coalition/qtsl-experiments} is the reproduction spine
--- committed certificates, mutant batteries, receipts --- cloned for
its bytes; nothing to build. \path{maitria-coalition/qtsl} is the
repository this book ships in: the theories, the machine-checked
soundness developments, and the small reference instruments; nothing
to build there either --- the reader you will run in
section~\ref{sec:qs-break} is dependency-free TypeScript on bare Node
$\ge 22$.

\begin{Verbatim}[fontsize=\small,frame=leftline,rulecolor=\color{accentmid},framesep=3mm,xleftmargin=3mm]
git clone https://github.com/maitria-coalition/mtk
cargo build --release --manifest-path mtk/Cargo.toml
git clone https://github.com/maitria-coalition/qtsl-experiments
git clone https://github.com/maitria-coalition/qtsl
\end{Verbatim}

\section{Produce}
\label{sec:qs-produce}

\begin{Verbatim}[fontsize=\small,frame=leftline,rulecolor=\color{accentmid},framesep=3mm,xleftmargin=3mm]
$ mtk/target/release/gymdp produce --out out --only frozenlake-4x4
frozenlake-4x4: exact T=100 bound = 0.7441903 [0.005s]
receipts -> out/producer_receipts_rust.json (9 certs, 0 compared, 0 mismatched) [0.005s total]
\end{Verbatim}

Nine certificates in milliseconds: the exact deadline ladder
$T \in \{5, 10, 25, 50, 100\}$ and a ladder of dyadic-rounded
companions. (The \texttt{0.7441903} on the log line is a float
\emph{rendering} of the $T{=}100$ bound, printed for the human reading
the log; no float appears inside any certificate.) The document this
chapter cares about is the deadline-10 certificate,
\path{out/frozenlake-4x4/exact_T10.cbor}: 2{,}511 canonical bytes
carrying the claim, the transition kernel as 148 rows of exact
rationals, the recorded policy, and the eleven-by-sixteen value table
--- the complete witness data from which everything asserted can be
re-derived by something that has never met the producer. It is also,
byte for byte, the artefact committed in the experiments repository
(\texttt{sha256} prefix \texttt{66040e59}, either way --- compare
against \path{qtsl-experiments/gymnasium/out/frozenlake-4x4/}). In
this toolkit ``reproduces the published artefact'' is a statement
about equality of bytes, not about agreement to some tolerance ---
which is what lets a checker's verdict on the committed corpus and on
your fresh production be the same fact.

One vocabulary correction before anyone checks anything, because the
file extension just lied to you: this toolkit has no certificate file
format. In the shipped architecture a certificate is ordinary
database content --- rows of a signed, content-addressed geolog
commit, carried on the same rails as every other fact the system stores
(the Executive Summary's one-set-of-bytes, three-consumers stance
--- chapter~\ref{ch:exec-summary}; the Geolog part of this book is
that database's own manual). The theory defining those rows is in
this repository, as data: \path{theories/Qtsl.geolog}, with its
drift-alarmed battery beside it. The
sidecar file you just wrote is the producer's local carriage for that
content --- and the reproduction spine already commits this family
both ways: the sidecar artefacts, and the same certificates sealed as
rows of signed geolog commits in a \path{glc/} directory beside them,
with content identity between the two forms re-proven at every
production. What survives every carriage is the standard the content
answers to: a self-contained document of exact numbers, complete
enough to re-derive every obligation, with no side channel to
consult. You will read a certificate directly as rows, in plain text,
two sections from now.

\section{What a checker re-derives --- and who checks}
\label{sec:qs-check}

Any checker of this certificate, by any route, re-derives four
obligation families from the document alone, in exact rational
arithmetic:

\begin{itemize}
\item \emph{well-formedness} --- kernel probabilities lie in $[0,1]$
  and sum to one per state--action pair; indices are in range;
\item \emph{boundary} --- the value table's final row vanishes off the
  goal set;
\item \emph{step} --- each table entry is dominated by its expected
  successor value under the recorded policy's action,
  $V_t(s) \le \mathbb{E}\!\left[V_{t+1}\right]$;
\item \emph{claim} --- $p \le V_0(\mathit{init})$.
\end{itemize}

\noindent Chained, these force the table to be a genuine
under-approximation of the reach probability, so the claim follows
from the bytes alone. Nothing the producer computed is taken on
trust; the producer's speed lives entirely upstream of the
certificate. This division --- fast untrusted search, small
independent replay --- is the toolkit's organizing principle, and it
is why the producer is allowed to be a GPU-accelerated,
float-assisted, LLM-advised anything: a wrong producer costs a
rejected certificate, never a false verdict.

Who does the re-deriving is a deliberately plural answer: one set of
bytes, three consumers, each trading speed against the size of what
you must trust (the Executive Summary carries the full stance).
Their standing for \emph{this} certificate, dated to this preview:

\begin{itemize}
\item \emph{The geolog engine} --- the fast path: certificate rows
  chased, monitored, and enforced like any other rows;
  engineering-grade assurance, deliberately no mechanized soundness
  story. This certificate's row form --- one sealed commit carrying
  the system, the specification, the judgment, and the proof content
  --- is committed in the reproduction spine beside the sidecar
  artefacts, and the fast-path checker re-verifies all four
  obligation families from the sealed rows alone: one make target
  re-checks the family's containers against their manifest, expected
  verdicts included, and the Experiments chapter's timing tables
  carry the fast-path column from real runs.
\item \texttt{qtslab} --- the certified checker: a standalone binary
  whose end-to-end soundness theorem (machine-checked in HOL4,
  grounded through the CakeML verified compiler in the Arm
  instruction set's semantics) says that if it prints ACCEPT on a
  claim followed by an arbitrary stream of bytes, the claim is
  semantically true. Precisely: acceptance binds to validity of the
  accepted sequent under every admissible valuation of the hypernet
  the certificate itself declares (chapter~\ref{ch:qtsl-hypernets})
  --- the Bernstein Logic fragment,
  which is all its input format carries today; a certificate
  reaching beyond that fragment is refused, not silently accepted
  (the Executive Summary is explicit about the scope: the boundary
  is \texttt{qtslab}'s, not the row vocabulary's). A demonstrated
  route, assembling now: the same
  compiler-and-theorem recipe has already carried a complete
  predecessor checker to a machine-code soundness theorem on real
  hardware --- and that proven chain, re-run on its native ISA
  during this chapter's verification campaign, verified campaign
  fixtures in single-digit seconds each and refused all six mutants
  put to it.
\item \emph{The Lean route} --- the same bytes reconstructed as a
  theorem inside Lean's kernel, by reflection; expected to be too
  slow for machine-scale certificates, and kept in the roster to
  \emph{measure} that rather than assert it. It has been measured,
  twice, across two certificate-format generations: the Experiments
  chapter carries the scaling law, the admission gate fitted to it,
  and the expectation confirmed with its boundary stated.
\end{itemize}

\noindent No transcript in this section, on purpose: for the
certified and Lean routes, printing one would mean pretending they
are further along for this family than they are, and a book whose
subject is checkable evidence does not get to fake its own. The fast
path is runnable today --- the reproduction spine's make target
re-checks this family's committed containers, mutant battery
included, whenever you like. The checking you can run \emph{and read
whole} is the next section's, on a certificate family whose rows are
committed as plain text.

\section{Break one}
\label{sec:qs-break}

A checker that cannot say \emph{no} teaches nothing. The FrozenLake
family's sealed rows come with a committed mutant battery of their
own, but a sealed container is built for checkers' eyes; the
published-table campaign additionally stores its certificates as
plain-text row listings you can open and read --- so we break one of
those, and the lesson transfers intact.

That campaign (section~\ref{sec:exp-sweep} of
chapter~\ref{ch:experiments}) re-derived a peer-reviewed
neural-certificate results table in exact arithmetic, and
\texttt{faulty\_uniform} is its most instructive row: the one whose
exact two-sided bracket showed a published ``sound'' bound to be no
bound at all. Its upper-bound certificate is committed as rows in the
reference listing notation --- plain text you can open:

\begin{Verbatim}[fontsize=\small,frame=leftline,rulecolor=\color{accentmid},framesep=3mm,xleftmargin=3mm]
$ cd qtsl-experiments/supermartingales/out/sec3
$ sed -n '62p' certs/faulty_uniform/bound.rows.txt
op[  59] funa fun=G<stack/TConst> args=[q:16524448999/8406920828251] value=L58
\end{Verbatim}

\noindent One hundred and twelve rows like this one --- sorts, exact
simplex coordinates, Bernstein coefficients, and the proof structure
that cites them --- resolving their names against pinned theory
listings. Beside the corpus, the book's own repository carries a
reference reader and checker for this certificate family: three
TypeScript files, zero dependencies, readable in an afternoon
(\path{qtsl/checkers/supermartingales}). Point it at the certificate,
and at the campaign's committed mutant of the same certificate:

\begin{Verbatim}[fontsize=\small,frame=leftline,rulecolor=\color{accentmid},framesep=3mm,xleftmargin=3mm]
$ CHK=../../../../qtsl/checkers/supermartingales
$ node --experimental-strip-types $CHK/check_rows.ts \
    --stack $CHK/theories/BLstack.rows.txt --cert $CHK/theories/BLCert.rows.txt \
    certs/faulty_uniform/bound.rows.txt mutants/faulty_uniform/M1_false_bound.rows.txt
ACCEPT certs/faulty_uniform/bound.rows.txt proof_nodes=3
REJECT mutants/faulty_uniform/M1_false_bound.rows.txt berncone.coeff
\end{Verbatim}

\noindent Now break it yourself. The certificate survives a
coefficient perturbed in the sound direction --- slack is real, and
the checker is not grading effort --- so nudge the coefficient you
just looked at the \emph{other} way, by exactly
$1/8{,}406{,}920{,}828{,}251$:

\begin{Verbatim}[fontsize=\small,frame=leftline,rulecolor=\color{accentmid},framesep=3mm,xleftmargin=3mm]
$ sed 's|448999/|449000/|' certs/faulty_uniform/bound.rows.txt > my_broken.rows.txt
$ node --experimental-strip-types $CHK/check_rows.ts \
    --stack $CHK/theories/BLstack.rows.txt --cert $CHK/theories/BLCert.rows.txt \
    my_broken.rows.txt
REJECT my_broken.rows.txt berncone.coeff
\end{Verbatim}

Two things to notice. The refusal is \emph{witness-shaped}: the
identifier names the violated obligation family (membership of the
certificate polynomial's Bernstein coefficients in the claimed cone),
and the checker's failure record beneath it carries the offending
coefficient's index and exact value, because failure geometry is
routing advice for whoever must repair the certificate. And the
committed mutant battery is the checker's obligation list made
executable --- the battery, not the prose, gives the operational
statement of what an \texttt{ACCEPT} means (the Cookbook develops
this discipline as a pattern:
section~\ref{sec:ckbk-mutant-battery} of
chapter~\ref{ch:mtk-cookbook}).

One label, to keep the trust story crisp: the program you
just ran is a \emph{reference instrument} --- the reproduction
spine's own exactness oracle, differentially pinned against an
independent implementation across the campaign corpus and its mutant
battery --- and deliberately \emph{not} one of the three checking
routes. Its job is to make ``the battery must refuse'' an executable
fact, and to let you watch obligations re-derived from rows alone
through code small enough to read whole. When the stakes are real,
the same rows go to the routes of section~\ref{sec:qs-check}.

\section{The same loop from TypeScript}
\label{sec:qs-typescript}

The producer toolkit's TypeScript surface, \texttt{@maitria/mtk},
ships as the \path{ts/} package of the \path{maitria-coalition/mtk}
repository (napi bindings over the same Rust core --- there is one
implementation of every construction, and TypeScript reaches it
in-process). The shipped install route is the checkout you already
have from section~\ref{sec:qs-setup} --- it requires a Rust
toolchain, because the native addon builds from source:

\begin{Verbatim}[fontsize=\small,frame=leftline,rulecolor=\color{accentmid},framesep=3mm,xleftmargin=3mm]
cd mtk/ts && npm install && npm run build:native   # Rust toolchain required
cd ~/your-project && npm install ../mtk/ts         # npm links the package
\end{Verbatim}

\noindent (\texttt{@maitria/mtk} on the npm registry: arriving; the
checkout path above is the shipped route at this preview.) Node
$\ge 22.18$ throughout --- the package's declared \texttt{engines}
floor. The producing half of the loop again, this time
constructing the certificate yourself --- and building the mutant
with your own hands:

\begin{Verbatim}[fontsize=\small,frame=leftline,rulecolor=\color{accentmid},framesep=3mm,xleftmargin=3mm]
// quickstart.ts — produce the FrozenLake certificate, then break it.
import { writeFileSync } from "node:fs";
import { dp } from "@maitria/mtk";

// The environment: Gymnasium's FrozenLake-4x4, slippery (success rate 1/3).
const model = dp.frozenlake("4x4", true, "1/3");

// The claim: exact finite-horizon dynamic programming to deadline T = 10.
const res = model.exact(10);
console.log("claim:", res.bound);

// Emit the certificate: canonical bytes, encoded by the Rust core.
const meta = {
  envId: "FrozenLake-v1",
  mdpName: "frozenlake-4x4",
  repoSha: dp.gymSha(),
  kwargs: "",
};
const bytes = res.emit(meta);
writeFileSync("exact_T10.cbor", bytes);
console.log("wrote exact_T10.cbor:", bytes.length, "bytes");

// Now break it: raise one interior value-table entry to 1.
// (docMutate refuses unless the unmodified doc round-trips byte-identically.)
const broken = dp.docMutate(bytes, [
  { path: ["table", 1, 0, 0], setHex: "01" },
  { path: ["table", 1, 0, 1], setHex: "01" },
]);
writeFileSync("mutant_table_raised.cbor", broken);
console.log("wrote mutant_table_raised.cbor:", broken.length, "bytes");
\end{Verbatim}

\begin{Verbatim}[fontsize=\small,frame=leftline,rulecolor=\color{accentmid},framesep=3mm,xleftmargin=3mm]
$ npx tsx quickstart.ts
claim: 815/19683
wrote exact_T10.cbor: 2511 bytes
wrote mutant_table_raised.cbor: 2509 bytes
\end{Verbatim}

The bytes you just constructed are byte-identical to the ones
\texttt{gymdp} emitted in section~\ref{sec:qs-produce}, and therefore
to the committed corpus artefact --- the same \texttt{sha256}, all
three ways. And you now hold what the mutant-battery discipline calls
a matched pair: golden bytes, and a variant broken in precisely one
field (the mutation seam checks that the unmodified document
round-trips byte-identically before it will perform surgery, so your
mutant differs from the golden bytes exactly where you aimed). The
corpus commits this family's battery beside the artefacts, expected
verdicts in the receipts --- and the sealed row form carries a second
battery of thirteen row-form mutants, each one's expected verdict and
violated obligation named in its manifest. You already know from
section~\ref{sec:qs-break} what a refusal at the row level looks
like.

\section{What you just did, and where it goes}
\label{sec:qs-next}

You ran the trust architecture of this book once, in miniature: an
untrusted producer constructed complete witness data fast, and
reproduced a published artefact to the byte; you read the obligations
a checker must re-derive, and the current state of the three routes
that re-derive them; and on the published corpus you watched an
exact-arithmetic reader accept a certificate stored as rows and
refuse its mutants --- including yours --- with the violated
obligation named. Three widenings from here:

\begin{itemize}
\item \emph{Producing for real.} The Tutorial
  (chapter~\ref{ch:qtsl-tutorial}) teaches the producer's craft on the
  TypeScript surface: exactness as a type property, propose-then-dispose,
  specs and judgments, derivation rows from the Rule Catalogue, and
  refusal-driven iteration.
\item \emph{Richer claims.} Deadline reachability is one species. The
  Rule Catalogue (chapter~\ref{ch:qtsl-manual}) is the full inventory
  of certified inference; the fabrication-line worked example
  (chapter~\ref{ch:qtsl-minifab}) runs the stack over a model with
  batching, failures, and re-entrant flow.
\item \emph{Graded checking, priced.} The three routes of
  section~\ref{sec:qs-check} are not interchangeable, and the
  Experiments chapter's design keeps their price differences public:
  three wall-clocks per experiment --- produce, check-fast,
  check-certified --- each measured on its route's native substrate,
  so the cost of each grade of assurance stays a published number
  rather than a promise.
\end{itemize}

 \chapter{Tutorial: Producing Proofs with MTK}
\label{ch:qtsl-tutorial}

The Quick Start (chapter~\ref{ch:qtsl-quickstart}) had you run one
producer end to end. This tutorial teaches you to \emph{be} one: to
construct specifications, judgments, and derivation rows with the
producer toolkit's TypeScript surface, to read refusals as
instructions, and to leave behind artefacts an independent checker
will accept. The intended reader is either a capable agent doing
verification work or a verification-literate human; the text is
written for the human, and the agent will cope. Every listing below
runs as printed, and every output shown is the verbatim transcript of
a real run: the book's build regenerates these transcripts by
executing the printed listings, and each one's source records the
exact tool revision and date it ran (one commodity arm64 box).

\section{The producer's stance}
\label{sec:tut-stance}

Three commitments shape every surface you are about to touch, and
every signature below answers to one of them:

\begin{enumerate}
\item \emph{Nothing here is trusted.} No function in the producer
  toolkit renders a verdict anyone should believe. Validation helpers
  exist --- they save round-trips --- but they are preflights; the
  checker's independent replay of your artefact is the only
  authority. A wrong producer costs a rejected certificate, never a
  false verdict, which is why the producer is allowed to be fast,
  heuristic, and as clever as it likes.
\item \emph{Witnesses always, never bare verdicts.} Any operation
  that can fail for a content reason returns the geometry of the
  failure --- which coefficient, which cell, which field, which
  subformula --- because failure geometry is routing advice for your
  next attempt. You will spend much of your producing life reading
  refusals; the API is designed so that this is pleasant.
\item \emph{Exactness is a type property, not a discipline.} Data
  that certificates cite is exact --- arbitrary-precision rationals
  and integers --- and lives in types that say so. Floating-point
  data proposes and routes but cannot cross into the exact world by
  accident: the one bridge is an explicit snap, and what comes out of
  it is a new proposal to certify, not a blessed float.
\end{enumerate}

\section{Setup}
\label{sec:tut-setup}

The TypeScript package \texttt{@maitria/mtk} lives in the
\path{maitria-coalition/mtk} repository as its \path{ts/} package
over a native addon crate (the front matter's preview note applies to
the package's migration state at this preview). It is two layers by
design: a deliberately flat native boundary (napi bindings over the
Rust core --- there is exactly one implementation of every
construction, and it is the Rust one; the differential receipts
shipped with each binding prove the TypeScript-driven path emits
bytes identical to the Rust-driven one) and an \emph{arrangement}
layer that owns all public idiom --- discriminated unions, branded
types, readonly records. You will only ever touch the arrangement.

\begin{Verbatim}[fontsize=\small,frame=leftline,rulecolor=\color{accentmid},framesep=3mm,xleftmargin=3mm]
git clone https://github.com/maitria-coalition/mtk
cd mtk/ts && npm install
npm run build:native   # cargo builds the addon; CUDA optional, loaded at runtime
\end{Verbatim}

\noindent A Rust toolchain is required --- the native addon builds
from source. To use the package from your own project, install it by
path (\texttt{npm install ../mtk/ts}; npm links the checkout), which
is the shipped route at this preview; \texttt{@maitria/mtk} on the
npm registry: arriving. Node $\ge 22.18$ (the package's declared
\texttt{engines} floor); producer scripts run under
\texttt{npx tsx} with no build step of their own. Boxes without a GPU
run the same scripts and get the same answers --- acceleration
changes wall-clock, never content.

\section{Exact numbers, and the one bridge from float}
\label{sec:tut-exact}

The class the certificates cite is \texttt{Q}: immutable,
bigint-backed exact rationals with a canonical text form
\texttt{"n/d"}. Everything about its API enforces stance~3:

\begin{Verbatim}[fontsize=\small,frame=leftline,rulecolor=\color{accentmid},framesep=3mm,xleftmargin=3mm]
// exactness.ts — class-E data is Q; number never crosses the boundary.
import { Q, lp } from "@maitria/mtk";

console.log(Q.from("0.1").add("0.2").toString());        // exact decimal parse
console.log(Q.from("-2.04e-2").toString());              // scientific, exact
console.log(Q.of(815, 19683).sub(Q.of(1, 27)).toString());

try {
  Q.from(0.1);                                           // refused: ambiguous
} catch (e) {
  console.log((e as Error).message);
}

// The one bridge from float to exact: an explicit snap. The snapped
// value is a new PROPOSAL to certify, never a blessed float.
console.log(lp.rationalize(0.3333333333333333, 1000).toString());
console.log(lp.rationalize(Math.PI, 1000).toString());
\end{Verbatim}

\begin{Verbatim}[fontsize=\small,frame=leftline,rulecolor=\color{accentmid},framesep=3mm,xleftmargin=3mm]
$ npx tsx exactness.ts
3/10
-51/2500
86/19683
Q.from(0.1): non-integer numbers are ambiguous — pass "n/d" text, a bigint, or Q.of(n, d)
1/3
355/113
\end{Verbatim}

Decimal \emph{text} parses exactly ($\texttt{"0.1"} \mapsto 1/10$:
the string denotes a rational, unambiguously), but the
\emph{double} \texttt{0.1} is refused --- it denotes
$3602879701896397/2^{55}$, and pretending otherwise is how
exactness-as-discipline projects fail. When you genuinely hold a
float and want a rational, you say so: \texttt{rationalize} snaps to
the best approximation under an explicit denominator budget (note
$\pi \mapsto 355/113$, as it should). The budget is a knob you will
meet again: coarser snaps make smaller, cheaper-to-check
certificates, and no snap can make a false claim true --- the
checker decides either way.

\section{A first certified fact}
\label{sec:tut-bernstein}

The smallest genuinely useful certified fact in this toolkit: the
sign of a polynomial on a region. Bernstein coefficients of $p$ over
a simplex bound $p$ there, so an all-positive coefficient table
\emph{is} a sign certificate --- data whose correctness an
independent checker re-derives by arithmetic alone, no cleverness
required.

\begin{Verbatim}[fontsize=\small,frame=leftline,rulecolor=\color{accentmid},framesep=3mm,xleftmargin=3mm]
// sign.ts — certify the sign of a polynomial on cells of a triangle.
import { Q, poly } from "@maitria/mtk";

// p(x, y) = x^2 + x*y - y + 1/3
const p: poly.Term[] = [
  { exps: [2, 0], coeff: 1 },
  { exps: [1, 1], coeff: 1 },
  { exps: [0, 1], coeff: -1 },
  { exps: [0, 0], coeff: Q.of(1, 3) },
];

// the unit triangle, split at the midpoint of its hypotenuse
const mid = [Q.of(1, 2), Q.of(1, 2)];
const t = poly.simplexBern({
  poly: p,
  cells: [
    [[0, 0], [1, 0], mid],
    [[0, 0], mid, [0, 1]],
  ],
  degree: 4, // elevate past deg(p): tighter enclosures
  device: "cpu", // pin the host rung: same integers on any box
});

console.log(`ran on ${t.device}; ladder served F/M/host = ${t.served}`);
for (let b = 0; b < t.cells; b++) {
  const [lo, hi] = t.enclosure(b);
  const verdict =
    lo.sign() > 0 ? "p > 0 CERTIFIED"
    : hi.sign() < 0 ? "p < 0 CERTIFIED"
    : "sign open (subdivide)";
  console.log(`cell ${b}: p in [${lo}, ${hi}] — ${verdict}`);
}
\end{Verbatim}

\begin{Verbatim}[fontsize=\small,frame=leftline,rulecolor=\color{accentmid},framesep=3mm,xleftmargin=3mm]
$ npx tsx sign.ts
ran on cpu(hostref); ladder served F/M/host = 0,0,2
cell 0: p in [1/6, 4/3] — p > 0 CERTIFIED
cell 1: p in [-2/3, 1/3] — sign open (subdivide)
\end{Verbatim}

Read the first line closely, because it is stance~3 at work in the
machinery. The enclosure pipeline is exact end to end --- integer
numerators over a shared per-cell denominator --- and it runs as a
\emph{precision ladder}: fixed-width device tiers serve the cells
they can handle without overflow, and any cell they cannot is flagged
and re-derived by an arbitrary-precision host replay. Every rung
computes the same exact integers; the ladder trades width for speed,
never exactness for speed. The listing pins \texttt{device: "cpu"} so
that both cells are served by the host rung (\texttt{0,0,2}) and your
transcript matches this one on any box, GPU or none; drop the pin
(the default is \texttt{"auto"}) and a free GPU takes the cells
instead --- on this box the package's own example reported
\texttt{cuda:NVIDIA GB10 ...\ ladder served F/M/host = 2,0,0}. Same
cells, same enclosures, same certificate --- the transcript differs
only in who did the arithmetic.

Cell 1 is the more instructive half of the output: the enclosure
straddles zero, so the honest answer is \emph{sign open}, with
subdivision as the routing advice. Bernstein enclosures converge
under subdivision and degree elevation, so ``open'' is an invitation
to spend more, not a dead end --- and when the budget runs out
before the sign closes, that too can be certified rather than
shrugged (the Cookbook's honest-stop pattern,
section~\ref{sec:ckbk-honest-stop} of
chapter~\ref{ch:mtk-cookbook}).

\section{Propose, then dispose}
\label{sec:tut-two-tier}

The two-tier discipline in one screen: a linear program solved twice,
once by the float lane (fast, unciteable), once by the exact lane
(the only tier a certificate may cite).

\begin{Verbatim}[fontsize=\small,frame=leftline,rulecolor=\color{accentmid},framesep=3mm,xleftmargin=3mm]
// two-tier.ts — the float lane proposes; the exact lane disposes.
import { lp } from "@maitria/mtk";

// minimize x0 + x1  subject to  x0 + 2*x1 = 1,  3*x0 + x1 = 1,  x >= 0
const cols = [{ cost: 1 }, { cost: 1 }];
const rows = [
  { lo: 1, hi: 1, cols: [0, 1], coefs: [1, 2] },
  { lo: 1, hi: 1, cols: [0, 1], coefs: [3, 1] },
];

// tier F: HiGHS, floats — fast, unciteable; duals are routing advice
const f = lp.floatSolve(cols, rows);
console.log("float:", f.status, f.objective, "duals", f.rowDuals);

// tier E: exact Bland simplex over Q — the only tier a certificate cites
const e = lp.simplexMin([1, 1], [[1, 2], [3, 1]], [1, 1]);
if (e.status === "opt") {
  const at = e.x.map(String).join(", ");
  console.log("exact:", e.value.toString(), "at x =", at);
}
\end{Verbatim}

\begin{Verbatim}[fontsize=\small,frame=leftline,rulecolor=\color{accentmid},framesep=3mm,xleftmargin=3mm]
$ npx tsx two-tier.ts
float: optimal 0.6000000000000001 duals [ 0.4, 0.2 ]
exact: 3/5 at x = 1/5, 2/5
\end{Verbatim}

The float tier even performs its own indictment: it reports
\texttt{0.6000000000000001} for a quantity that is $3/5$. On a toy
problem the exact tier could simply run alone; the discipline earns
its keep at scale, where the float lane's job is to find \emph{which}
exact computation is worth paying for --- its duals say which
constraints bind, its solution says where to warm-start, its speed
lets you sweep a thousand candidates to find the one worth
certifying. The full pattern, with the case studies that shaped it,
is the Cookbook's opening chapter
(section~\ref{sec:ckbk-proposer-disposer} of
chapter~\ref{ch:mtk-cookbook}).

\section{Stating claims: specs and judgments}
\label{sec:tut-spec}

So far we have certified facts the toolkit already knew how to
phrase. To produce proofs \emph{about systems}, you state the claim
yourself, in QTSL: a specification is a record of claim components
--- each a formula over a system's observables --- and a judgment is
the record that pairs a specification
with a system (or a region, or a pair of specifications --- one
record species per judgment kind). Both are constructed here as
plain data, preflighted in Rust, and emitted as versioned artefacts
through the same canon encoder that carries every other certificate.

The Quick Start's FrozenLake claim, stated properly this time:

\begin{Verbatim}[fontsize=\small,frame=leftline,rulecolor=\color{accentmid},framesep=3mm,xleftmargin=3mm]
// spec.ts — state the FrozenLake claim species as a specification.
import { spec } from "@maitria/mtk";

// One observable (width 1); the goal region is { x >= 15 }.
const x = spec.unwrap(spec.polyVar(0, 1));
const goal = spec.polyGe(15, x);

// P_{>= 815/19683} [ F_{<=10} goal ] — reach the goal within 10 steps:
// a safety-face component (SProbGE) in the record's forall list.
const phi = spec.probGe("815/19683", spec.deadline(10, goal));

const record = spec.unwrap(spec.build({ name: "reach-goal-by-10", forall: [phi] }));
const artifact = spec.unwrap(spec.encode(record));
console.log(artifact.schema, "-", artifact.bytes.length, "bytes");

// Round-trip: decoded data is data — re-validate before trusting it.
const back = spec.unwrap(spec.decode(artifact.bytes));
console.log("decoded:", back.spec.name, "v =", back.v);

// Now ask for probability 3/2 and let the preflight refuse.
const bad = spec.probGe("3/2", spec.deadline(10, goal));
const r = spec.preflightForall(bad);
if (!r.ok) {
  console.log("refused:", r.refusal.kind);
  if (r.refusal.kind === "probRange")
    console.log(`  at ${r.refusal.at.text}: p = ${r.refusal.p}, outside (0, 1]`);
}
\end{Verbatim}

\begin{Verbatim}[fontsize=\small,frame=leftline,rulecolor=\color{accentmid},framesep=3mm,xleftmargin=3mm]
$ npx tsx spec.ts
maitria.qtsl.spec.v1 - 179 bytes
decoded: reach-goal-by-10 v = 0
refused: probRange
  at (root): p = 3/2, outside (0, 1]
\end{Verbatim}

Three habits to take from one screen. First, the wire is visible and
versioned: every emitted artefact carries a schema string
(\texttt{maitria.qtsl.spec.v1}) and a version, and the bytes are
produced by the Rust canon encoder --- TypeScript never
re-implements an encoding, which is why cross-language byte drift is
structurally unrepresentable rather than tested away. Second,
decoded data is data: \texttt{decode} gives you back a record, and
if you received the bytes from someone else you re-run \texttt{build}
before leaning on it. Third, the refusal: it is a discriminated value
(\texttt{kind: "probRange"}), it carries the address of the offence
inside the formula tree (\texttt{(root)} --- addresses are relative
to the node you preflighted, and record-level checks come with full
paths like \texttt{forall[0].prob.path}), and it carries the
offending exact value as canonical text (\texttt{p: "3/2"}). Producer code branches on refusal kinds the way it
branches on any other data; nothing is buried in exception strings.

Judgments follow the same grammar one level up: the \texttt{judgment}
namespace has one builder per judgment kind --- satisfaction, region
satisfaction, spec equivalence, distance, and the rest --- each
taking branded citation handles rather than loose strings, each
emitting a versioned artefact. The claim record is deliberately
boring; the interesting objects are the derivations that discharge
it, which is where the rules come in.

\section{Derivation rows: the Rule Catalogue as an API}
\label{sec:tut-rules}

Every inference rule in the Rule Catalogue
(chapter~\ref{ch:qtsl-manual}) is available as a typed row builder:
premises as typed citations, payload fields exact, side conditions
enforced at construction, canonical bytes out. A proof, in this
toolkit, is a table of such rows --- which is what lets a checker
replay it row by row, and what lets a database engine store, merge,
and re-derive it like any other data.

The smallest interesting rule is a good first row.
\rulename{QUnionBound} takes two exception-mass bounds and concludes
their sum --- the workhorse glue for assembling per-obligation
failure budgets into a system-level one:

\begin{Verbatim}[fontsize=\small,frame=leftline,rulecolor=\color{accentmid},framesep=3mm,xleftmargin=3mm]
// rule-row.ts — a Rule Catalogue row: cite premises, get the conclusion.
import { rules } from "@maitria/mtk";
const { qtslCore } = rules;

// QUnionBound: from two exception-mass bounds, conclude their sum.
const row = qtslCore.unionBound({
  e1: qtslCore.judgmentRef("qregionsat", 0),   // typed premise citations
  e2: qtslCore.judgmentRef("qregionsat", 1),
  a: "1/8",
  b: "1/16",
});
if (row.ok) {
  console.log(row.value.schema, "-", row.value.cbor.length, "bytes");
  console.log("concluded bound:", row.value.bound.toString());
}

// A negative mass is refused with the witness attached.
const bad = qtslCore.unionBound({
  e1: qtslCore.judgmentRef("qregionsat", 0),
  e2: qtslCore.judgmentRef("qregionsat", 1),
  a: "-1/8",
  b: "1/16",
});
if (!bad.ok) {
  console.log("refused:", bad.refusal.kind);
  console.log("  " + bad.refusal.message);
}
\end{Verbatim}

\begin{Verbatim}[fontsize=\small,frame=leftline,rulecolor=\color{accentmid},framesep=3mm,xleftmargin=3mm]
$ npx tsx rule-row.ts
maitria.qtsl.union_bound.v0 - 58 bytes
concluded bound: 3/16
refused: signConstraint
  QUnionBound.a: sign/range constraint violated by -1/8
\end{Verbatim}

Notice who computed $3/16$: not this script. Derived conclusions
come back from the Rust core's own \texttt{concludes} logic, so the
row you emit and the arithmetic it claims were produced by the same
implementation the differential receipts pin --- your script
assembles and cites, it never does load-bearing arithmetic. Notice
also the premise handles: \texttt{judgmentRef("qregionsat", $i$)} is
a branded type, and a rule whose premise slot wants a region-\allowbreak
satisfaction judgment will not accept a plain number, a proof
reference, or a judgment of the wrong kind. The edge discipline of
the proof calculus is enforced by the type checker before the
checker ever sees your bytes.

When you need more than glue --- safety batteries, deadline
batteries, Streett bundles, Morse decompositions, composition rules
--- the pattern never changes: find the rule in the Rule Catalogue,
which documents its premises, payload, and side conditions alongside
its mechanization status; call its builder; read refusals until the
row constructs. The catalogue chapter is the map and stays the map;
this tutorial deliberately teaches only the door.

\section{Certificates are documents: read one}
\label{sec:tut-anatomy}

The Quick Start ended by mutating a certificate through the generic
document seam. Here is the reading half of that seam, on the same
bytes --- because the fastest way to understand what a checker
re-derives is to look at what the producer actually wrote down:

\begin{Verbatim}[fontsize=\small,frame=leftline,rulecolor=\color{accentmid},framesep=3mm,xleftmargin=3mm]
// anatomy.ts — a certificate is a document; read it before you trust it.
import { readFileSync } from "node:fs";
import { dp } from "@maitria/mtk";

const bytes = readFileSync("exact_T10.cbor");
const doc = dp.docJson(bytes, []) as Record<string, unknown>;
console.log("fields:", Object.keys(doc).join(" "));
console.log("deadline:", dp.docJson(bytes, ["deadline"]),
            " states:", dp.docJson(bytes, ["ns"]),
            " goal:", dp.docJson(bytes, ["goal"]));
console.log("claim:", dp.docJson(bytes, ["claim"]));      // bignum pair, hex
console.log("kernel[0]:", dp.docJson(bytes, ["kernel", 0])); // s a s' p_num p_den
\end{Verbatim}

\begin{Verbatim}[fontsize=\small,frame=leftline,rulecolor=\color{accentmid},framesep=3mm,xleftmargin=3mm]
$ npx tsx anatomy.ts
fields: claim deadline env_id goal init kernel kwargs lane mdp_name ns policy repo_sha table v
deadline: 10  states: 16  goal: [ 15 ]
claim: [ '0x032f', '0x4ce3' ]
kernel[0]: [ 0, 0, 0, '0x02', '0x03' ]
\end{Verbatim}

Everything a checker of this certificate must re-derive --- the
Quick Start's four obligation families --- is visibly \emph{here}:
the claim ($\mathtt{0x032f}/\mathtt{0x4ce3} = 815/19683$), the kernel
as exact rationals (the first row: from state 0 under action 0,
probability $2/3$ of staying put --- two of its three equally likely
moves bounce off the corner), the policy, the value table. A
certificate contains no argument and no trust requests; it contains
the numbers that make the argument checkable. When you produce a new
certificate species, this is the standard to hold your bytes to:
could a stranger's checker re-derive every obligation from the
document alone?

That standard is what survives every carriage the document rides.
These particular bytes travel as a sidecar file; the reproduction
spine commits the same certificate a second way, as rows of a
signed, content-addressed geolog commit --- ordinary database
content whose row vocabulary is \path{theories/Qtsl.geolog} in this
repository, with no certificate file format of its own --- and you
have already read certificate rows directly, as plain text, in the
Quick Start's break-it section. What the generic document seam gives
you here is the producer-side reading of the content in its
canonical byte carriage; the anatomy lesson is identical on either,
because completeness reads the same from any carriage.

The mutation seam (\texttt{dp.docMutate}) is the same document
surface pointed the other way, and it exists because a mutant
battery is part of a certificate lane's deliverable, not an optional
extra: every claim species ships with mutants its checker must
refuse, and a witness that stays green when the claim is broken was
never a witness (section~\ref{sec:ckbk-mutant-battery} of
chapter~\ref{ch:mtk-cookbook}). The seam checks that the unmodified
document round-trips byte-identically before it will perform
surgery, so your mutants differ from the golden bytes in precisely
the field you broke.

\section{The shape of a real producer program}
\label{sec:tut-program}

Scale the last five sections up and you have the actual production
programs behind the Experiments chapter. The FrozenLake producer you
ran in the Quick Start is a few hundred lines organized along this
tutorial's spine, and reading it (the \texttt{gymdp} producer in the
\path{maitria-coalition/mtk} repository, with its receipts and
companion drivers in \path{maitria-coalition/qtsl-experiments}) is
the recommended next hour. The anatomy the campaign's producers
share:

\begin{itemize}
\item \emph{a configuration grid}, not a one-off --- environments
  $\times$ deadlines $\times$ rounding budgets, because certificates
  are cheap once the lane exists and ladders are how cost curves get
  published;
\item \emph{an independent cross-instrument} --- a plain
  floating-point value sweep, kept structurally separate from the
  exact lane, with an explicit agreement tolerance; it certifies
  nothing, and it catches producer bugs before any checker has to;
\item \emph{exact dominance assertions between lanes} --- each
  dyadic-rounded certificate is checked against its exact sibling
  ($0 \le \text{exact} - \text{dyadic} \le T\cdot 2^{-k}$, the
  \texttt{gymdp} source's own \emph{exact dominance} assertion), so
  even the cheap lane's error is a theorem, not a hope;
\item \emph{emission through the canon encoder, then the mutant
  battery, then receipts} --- machine-readable production receipts
  (bounds, byte counts, timings, mutant verdicts) written beside the
  artefacts, because a reproduction spine is built from receipts,
  not recollections.
\end{itemize}

\noindent When a step refuses mid-program, the refusal's geometry
picks the escalation: subdivide the cell, raise the degree, spend a
bigger denominator budget, climb from the closed form to the LP to
the heavier rule --- cheapest instrument first, and certify the
honest stop when the budget genuinely ends
(sections~\ref{sec:ckbk-cheap-first} and~\ref{sec:ckbk-honest-stop}
of chapter~\ref{ch:mtk-cookbook}).

\section{Working an obligation workspace}
\label{sec:tut-session}

Everything so far drove the toolkit imperatively: your script decided
each call, held the intermediate state in variables, and threw both
away on exit. That is the right shape for a five-minute production
and the wrong shape for one that runs long enough to be interrupted,
plays several instruments against one claim, or is worked by
something other than its author. The toolkit's session surface
exists for that regime. Its bookkeeping object is the
\emph{obligation workspace}, and its vocabulary is four nouns:

\begin{itemize}
\item an \emph{obligation} is a claim to be certified --- a statement
  of a registered claim kind --- with a status: \texttt{open},
  \texttt{discharged}, or \texttt{refuted};
\item a \emph{move} is one typed action from the obligation's menu.
  The menu is the claim kind's to offer: each move declares its
  parameter schema, a cost hint (\texttt{cheap} to
  \texttt{expensive} --- advice for sequencing, never a gate), and
  what it may do: discharge, refute, or spawn;
\item an \emph{attempt} is one executed move: the outcome, the wall
  time, and either a \emph{witness} (for discharged or refuted ---
  the exact-tier artefact that justifies the verdict) or
  \emph{evidence} (for inconclusive --- failure geometry, which by
  now you know to read as routing advice);
\item the \emph{event log} is the whole truth: the workspace is a
  pure fold of an append-only event list, so replaying the log
  reconstructs the state without re-running any kernel.
\end{itemize}

\noindent Two disciplines are enforced by the workspace itself
rather than requested of its users. First, stance~2 is mechanized at
the recording boundary: a claim kind that returns a verdict without
a witness is refused on the spot --- the workspace will not record a
bare verdict, so its bookkeeping can never get ahead of the
artefacts that justify it. Second, structural moves \emph{spawn}: a
decomposition move creates child obligations that jointly cover the
parent's claim, and from that moment the parent's status is nobody's
to set --- it closes only by propagation. All children discharged
closes the parent (\texttt{via children}); one refuted child refutes
it, since the child covers part of the claim. You work the leaves;
the workspace does the lifting.

Here is the whole model in one sitting, on a claim small enough to
watch: $x(1-x) - 1/4 \le 0$ on $[0,1]$ --- true, with equality at
$x = 1/2$. (The package is \texttt{@maitria/mtk-session}, the
\path{maitria-coalition/mtk} repository's \path{session/} package,
beside the producer package; install it by path like the producer
package in section~\ref{sec:tut-setup}. It runs its TypeScript
sources directly --- Node $\ge 22.18$.)

\begin{Verbatim}[fontsize=\small,frame=leftline,rulecolor=\color{accentmid},framesep=3mm,xleftmargin=3mm]
// session.ts — work an obligation workspace by hand: you are the strategy.
import { Workspace, polyKind } from "@maitria/mtk-session";

// The claim: x(1-x) - 1/4 <= 0 on [0, 1] — true, with equality at x = 1/2.
const ws = new Workspace([polyKind]);
const root = ws.pose({
  kind: "poly.nonpositive",
  poly: [
    { exps: [2], coeff: "-1" },
    { exps: [1], coeff: "1" },
    { exps: [0], coeff: "-1/4" },
  ],
  box: [["0", "1"]],
});
console.log(`${root.id}: ${ws.describe(root.id)}`);
console.log("menu:", ws.moves(root.id).map((m) => `${m.name}/${m.cost}`).join(" "));

// Cheapest first: one exact Bernstein enclosure of p over the whole box.
const r1 = ws.apply(root.id, "range");
console.log("range:", r1.attempt.outcome, JSON.stringify(r1.attempt.evidence));
console.log("  note:", r1.attempt.note);

// The enclosure straddles zero because p touches it at 1/2. Split there.
const r2 = ws.apply(root.id, "split", { axis: 0, at: "1/2" });
console.log("split:", r2.attempt.outcome, "->", r2.spawned.map((c) => c.id).join(", "));

// The parent is now out of our hands — the workspace says so:
try {
  ws.apply(root.id, "range");
} catch (e) {
  console.log("refused:", (e as Error).message);
}

// Work the leaves; each half certifies by one enclosure.
for (const child of r2.spawned) {
  const r = ws.apply(child.id, "range");
  console.log(`${child.id} range:`, r.attempt.outcome, JSON.stringify(r.attempt.witness));
}

// Nobody set the parent's status; the second discharge propagated it.
const done = ws.get(root.id);
console.log(`${root.id}: ${done.status} (via ${done.via})`, JSON.stringify(done.witness));

// The workspace is a fold of its event log: replay rebuilds it exactly.
const replica = Workspace.replay(ws.events, [polyKind]);
const same = JSON.stringify(replica.snapshot()) === JSON.stringify(ws.snapshot());
console.log("replayed", replica.events.length, "events; state identical:", same);
\end{Verbatim}

\begin{Verbatim}[fontsize=\small,frame=leftline,rulecolor=\color{accentmid},framesep=3mm,xleftmargin=3mm]
$ npx tsx session.ts
ob1: (-1)*x0^2 + (1)*x0 + (-1/4) <= 0 on [0, 1]
menu: range/cheap probe/cheap split/cheap search/expensive
range: inconclusive {"enclosure":["-1/4","1/4"]}
  note: enclosure straddles 0; split, search, or probe near the suspected maximum
split: spawned -> ob2, ob3
refused: ob1 has open children; work the leaves (its status resolves by propagation)
ob2 range: discharged {"enclosure":["-1/4","0"]}
ob3 range: discharged {"enclosure":["-1/4","0"]}
ob1: discharged (via children) {"dischargedChildren":["ob2","ob3"]}
replayed 5 events; state identical: true
\end{Verbatim}

Four things in this transcript repay a close look.

\emph{The straddle is real.} The enclosure on the whole box is
$[-1/4, 1/4]$: the control polygon of a strictly concave parabola
lies above it, so the degree-2 ceiling lands at $1/4$ rather than at
the true maximum $0$. The note beside the evidence says what to do
about it, and the interesting word in it is \emph{where}: ``near the
suspected maximum.''

\emph{The split point was knowledge.} Splitting at the touch point
$x = 1/2$ moves each half's maximizer to an endpoint --- and
endpoint Bernstein coefficients are function values, so both child
enclosures land their ceilings at exactly $0$ (the transcript's
\texttt{["-1/4","0"]}, twice) and discharge by one cheap move each.
A midpoint-splitting heuristic would have done the same here by
luck; on a claim whose touch point is off-center, choosing the split
is a contribution only the mover can make. Hold that thought for one
section.

\emph{Nobody closed the parent.} The attempted third move on
\texttt{ob1} was refused, in words that state the discipline; the
final status arrived by propagation when the second child
discharged, and the parent's witness is the decomposition itself ---
\texttt{\{"dischargedChildren":["ob2","ob3"]\}} --- while the
exact-tier artefacts live on the leaves.

\emph{The log determines the state.} Five events: one pose and four attempts
(the refused call appended nothing). The last attempt's event
carries two closures --- the child's, and the root's by propagation
--- and replay applies recorded effects without re-executing any
move, which is why the rebuilt workspace is identical even though
nothing requires a kernel to be deterministic.

The daemon is this same object held behind a socket.
\texttt{mtk-session serve} (run from the package checkout as
\texttt{node src/main.ts serve} at this preview) holds one workspace
and speaks newline-delimited JSON-RPC~2.0 over a unix domain socket,
or over stdio for harnesses that spawn it as a child: one request
per line, one response per line, seven methods
(\texttt{session.info}; \texttt{obligation.pose}, \texttt{.list},
\texttt{.get}, \texttt{.moves}, \texttt{.apply};
\texttt{workspace.log} for the event tail). State lives in the
daemon, so clients disconnect and reconnect freely, and errors carry
structured codes rather than prose. JSON is a deliberate choice at
this one boundary --- the clients are foreign by construction: any
agent runtime, any language, or a person with \texttt{nc} --- and no
certificate content crosses it: certificates land through the same
canon encoding as everywhere else in this book, and the protocol
carries menus, selections, and witness summaries.

\section{The strategy layer}
\label{sec:tut-strategy}

In the last listing, the workspace kept the books, the kernels
computed the truth, and everything interesting in between --- range
first because it is cheap, split \emph{there} because that is where
$p$ touches zero --- was you. That seat has a name in the session
surface: the \emph{proposer}. A proposer is shown the open
obligations, the evidence accumulated so far, and the move menus,
and asked for one selection at a time. The Executive Summary
promised this seat when it described agents driving tactic selection
over an obligation workspace ``with the strategy intelligence
supplied by the reader as the language model of their choice''
(chapter~\ref{ch:exec-summary}), and this section is
that layer --- shaped so that anything can sit in it. The opening
page addressed this tutorial to a capable agent doing verification
work; over the serve socket, that reader holds the seat personally.

The trust statement, precisely. A proposer chooses \emph{which}
exact computation runs next --- nothing else. It cannot fabricate an
outcome: a proposal outside the posted menu is rejected against the
menu and the error is fed back; two consecutive failures end the
loop with the workspace untouched. It cannot inject a witness:
witnesses are constructed by move execution only, behind the
no-bare-verdict guard of the previous section. And it cannot upgrade
a shortfall into a discharge, because the comparisons happen in
exact rational arithmetic inside the claim kinds. Stance~1 gains its
corollary: because verdicts never come from the strategist, the
strategist may be arbitrarily wrong --- confused, overconfident,
adversarial --- at a cost denominated in budget, never in soundness.
The same reason the producer was allowed to be fast and heuristic,
applied one level up.\footnote{This seat has a literature: Laurent
  and Platzer learn to refine proof-search strategies
  \cite{laurent22}, and their oracular programming gives LLM-enabled
  software a modular foundation in which untrusted oracles answer
  through checked interfaces \cite{laurent25}; closest to this
  chapter's workflow, Kabra and colleagues use language models to
  automate interactive theorem proving and control synthesis for
  hybrid games \cite{kabra26}.}

The package ships two proposers, and the contrast between them is
the demonstration. The \emph{scripted} proposer is a deterministic
cheapest-first ladder per claim kind: no key, no network, the CI
path, and the budget floor any cleverer strategy ought to beat. Here
it drives the Quick Start's claim (section~\ref{sec:qs-claim} ---
the same claim; there you ran one producer end to end, here a
workspace orchestrates a ladder of instruments):

\begin{Verbatim}[fontsize=\footnotesize,frame=leftline,rulecolor=\color{accentmid},framesep=3mm,xleftmargin=3mm]
$ node src/main.ts demo
mtk-session demo — the FrozenLake deadline claim, end to end

claim ob1: FrozenLake-4x4 (slippery, success 1/3): recorded policy reaches the goal
  within 10 steps with probability >= 815/19683
strategy: scripted
  [1] ob1 float_sweep -> inconclusive (2 ms) {"estimate":0.04140628969161203,
      "claim":0.04140628969161205,"marginEstimate":-2.0816681711721685e-17}
  [2] ob1 dyadic {"kBits":24} -> inconclusive (0 ms) {"tier":"dyadic","kBits":24,
      "bound":"694681/16777216","shortfall":"24917/330225942528"}
  [3] ob1 exact -> discharged (1 ms) {"tier":"exact","bound":"815/19683","claim":"815/19683",
      "certificate":{"bytes":2511,
      "sha256":"66040e596ece30efcb1e500350396477f6e4f56d9af1b8f2896c30bece2fc391",
      "path":"out/frozenlake-4x4/exact_T10.cbor"}}
result ob1: discharged (via exact)
witness: {"tier":"exact","bound":"815/19683","claim":"815/19683","certificate":{"bytes":2511,
  "sha256":"66040e596ece30efcb1e500350396477f6e4f56d9af1b8f2896c30bece2fc391",
  "path":"out/frozenlake-4x4/exact_T10.cbor"}}
strategy ended: solved; moves used: 3
\end{Verbatim}

(Long lines are wrapped at JSON commas to fit the page; the
package's committed receipts carry both of this section's
transcripts unbroken.) The three attempts are one story told
forward. The float sweep is the cross-instrument of
section~\ref{sec:tut-program} in move form: its margin estimate,
$-2\times 10^{-17}$, is float for ``the claim sits dead even with
the optimum'' --- which the Quick Start disclosed in
advance: $815/19683$ \emph{is} the exact ten-step optimum, so this
claim has zero slack. The dyadic attempt is a sound certifying lane,
and its evidence is exact: a lower bound of $694681/2^{24}$, short
of the claim by exactly $24917/(2^{24}\cdot 3^{9})$. The shortfall
is structural, not a tuning failure: every reach probability in this
model has a power of three for its denominator, the claim is such a
number in lowest terms, and no dyadic rational equals it --- so a
rounded-down dyadic bound sits strictly below a zero-slack claim at
every precision, and no \texttt{kBits} fixes it. The workspace
records that gap as an exact rational, the evidence routes the
escalation, and the exact tier discharges with zero slack, its
witness carrying the certificate content --- the same content,
digest for digest, that the Quick Start produced.

One seam note at this preview: the demo's remaining transcript
lines, not shown above, hand the emitted certificate and a
deliberately broken copy of it to an independent checker, which
accepts the former and refuses the latter naming the violated
obligation. We leave those lines out because that checking surface
is mid-migration to consuming certificates the way the Executive
Summary describes them --- signed, content-addressed commit content,
checked by any of the three routes --- rather than as files handed
to a dedicated binary. The workspace is indifferent to the
migration: what it stores is a witness, and what a witness cites is
content.

To seat a model instead of the script, supply an OpenRouter key in
the environment (\texttt{OPENROUTER\_API\_KEY}; any chat-completions
model works, and the router's default is used if you name none). The
model is shown the same view a client at the socket gets, and must
answer with one JSON object selecting a move:

\begin{Verbatim}[fontsize=\footnotesize,frame=leftline,rulecolor=\color{accentmid},framesep=3mm,xleftmargin=3mm]
$ node src/main.ts demo --llm --model google/gemini-3.5-flash
mtk-session demo — the FrozenLake deadline claim, end to end

claim ob1: FrozenLake-4x4 (slippery, success 1/3): recorded policy reaches the goal
  within 10 steps with probability >= 815/19683
strategy: llm:google/gemini-3.5-flash
  ... asking google/gemini-3.5-flash
  [1] ob1 float_sweep -> inconclusive (1 ms) {"estimate":0.04140628969161203,
      "claim":0.04140628969161205,"marginEstimate":-2.0816681711721685e-17}
  ... asking google/gemini-3.5-flash
  [2] ob1 exact -> discharged (1 ms) {"tier":"exact","bound":"815/19683","claim":"815/19683",
      "certificate":{"bytes":2511,
      "sha256":"66040e596ece30efcb1e500350396477f6e4f56d9af1b8f2896c30bece2fc391",
      "path":"out/frozenlake-4x4/exact_T10.cbor"}}
result ob1: discharged (via exact)
witness: {"tier":"exact","bound":"815/19683","claim":"815/19683","certificate":{"bytes":2511,
  "sha256":"66040e596ece30efcb1e500350396477f6e4f56d9af1b8f2896c30bece2fc391",
  "path":"out/frozenlake-4x4/exact_T10.cbor"}}
strategy ended: solved; moves used: 2
\end{Verbatim}

The transcript's shape makes the argument. Shown the sweep's
dead-even margin, the model --- Gemini 3.5 Flash through OpenRouter,
in the committed run --- selected the exact tier at once: two moves
where the ladder spent three, because it read evidence a fixed
ladder cannot. Nothing depended on it being right. A worse selection
would have cost one extra attempt; nonsense selections would have
ended the loop after two rejections with the workspace untouched;
and no selection of any kind could have moved a verdict, because
routing is the one contribution a strategist can make here and the
only one it is allowed. Same claim, same certificate content, same
digest: strategy changes the path, never the bytes.

Past the demo, the seat generalizes.
\texttt{node src/main.ts solve claim.json --strategy llm} poses any
statement of a registered kind and drives it the same way
(\texttt{--strategy scripted} is the no-key floor), and an agent
harness that prefers to hold the seat itself connects to
\texttt{serve} and speaks the protocol directly. The workspace
neither knows nor cares what kind of mind proposes.

\section{Where this goes}
\label{sec:tut-next}

You can now construct every kind of object this book's spine runs
on --- exact scalars, certified enclosures, claim records,
derivation rows, certificate documents with their mutant batteries
--- and you have worked their production from an obligation
workspace under two different strategists. Three directions from
here, in increasing order of ambition:

\begin{itemize}
\item \emph{A worked model at full scale.} The fabrication-line
  chapter (chapter~\ref{ch:qtsl-minifab}) runs this stack over a
  model with re-entrant flow, batching, failures, and shared
  personnel --- and shows the checker's economics on the way. The
  System Catalogue of model types (chapter~\ref{ch:qtsl-mathematical-models})
  and the composition patterns
  (chapter~\ref{ch:qtsl-composition-patterns}) are its reference
  spine.
\item \emph{The operating philosophy.} The Cookbook
  (chapter~\ref{ch:mtk-cookbook}) is the patterns manual this
  tutorial has been quietly citing throughout; read it the way you
  would read \emph{How to Solve It} \cite{polya1945solve} --- before
  you need it.
\item \emph{Sessions under your own harness.} The workspace and the
  strategy seat (sections~\ref{sec:tut-session}
  and~\ref{sec:tut-strategy}) are importable pieces ---
  \texttt{Workspace}, the proposer interface, the RPC client --- and
  two growth directions are yours from here: hold the seat yourself,
  over the socket, from whatever harness you already run; or
  register a new claim species by implementing one
  \texttt{StatementKind} (\texttt{validate}, \texttt{describe},
  \texttt{moves}, \texttt{exec}) whose moves call the builders this
  chapter taught --- which is how the menu discipline extends to
  specs, judgments, and derivation rows.
\end{itemize}

\noindent A graduation exercise, meanwhile, that needs nothing not
already on your disk: FrozenLake has a hitting companion,
\texttt{dp.cliffwalkingHitting()}. Produce its deadline certificate,
read its document, and design its claim-inflation mutant --- the
matched pair the battery discipline requires --- then check the
committed corpus to see how close your battery came to the shipped
one (expected verdicts are recorded in the receipts beside the
artefacts, and the family's sealed row form ships a mutant battery
of its own, each entry's expected verdict and violated obligation
named in its manifest).

 \chapter{The Semantic Universe}
\label{ch:qtsl-semantic-universe}

\noindent\draftmark\medskip

\begin{figure}[b!]
\centering
\begin{tikzpicture}[x=1mm,y=1mm,
  storey/.style={draw, rounded corners=1.2pt, minimum width=80mm,
    minimum height=6.8mm, align=center, font=\small, inner xsep=2mm},
  tag/.style={font=\footnotesize, anchor=west},
  lift/.style={-{Stealth[length=1.8mm]}, semithick}]
  \node[storey] (order)   at (0,0)  {the causal order: events, causality $\le$, chronology $\ll$};
  \node[storey] (site)    at (0,9)  {the spacetime site, derived from the order --- plus duration, tangency};
  \node[storey] (records) at (0,18) {records over a window};
  \node[storey] (carrier) at (0,27) {the record carrier $\mathbb{H}(W)$: Fell hyperspace of closed sets; records enter as graphs};
  \node[storey] (laws) at (0,36) {probability distributions over records};
  \node[storey] (universe) at (0,45) {beliefs $\sigma \colon \Delta(\mathcal{X}) \to [0,\infty]$, lsc; \ radial stock $\cong$ credal bodies};
  \node[storey] (hyper)   at (0,54) {the hyperspace of bodies --- Systems are singleton elements};
  \node[storey] (specs)   at (0,63) {the fibre of Specs: constructible classes of bodies};
  \node[storey] (dc)      at (0,72) {the double category of System Specs, over interfaces and wirings};
  \node[storey] (total)   at (0,81) {the Semantic Universe: the iterated total --- Specs over interfaces over site-profiles};
  \foreach \a/\b in {order/site, site/records, records/carrier, carrier/laws,
    laws/universe, universe/hyper, hyper/specs, specs/dc, dc/total}
    \draw[lift] (\a.north) -- (\b.south);
  \coordinate (tagcol) at ([xshift=1em]current bounding box.east);
  \foreach \a/\s in {order/sec:su-causal-order, site/sec:su-site,
    records/sec:su-records, carrier/sec:su-carrier,
    laws/sec:su-denotation, universe/sec:su-denotation,
    hyper/sec:su-hyperspace, specs/sec:su-classes,
    dc/sec:su-composition, total/sec:su-total}
    \node[tag] at (tagcol |- \a.center) {\S\ref{\s}};
\end{tikzpicture}

\caption{The storeys of the semantic universe, in dependency order, with the
  section that builds each; the sections follow the same order. On the
  universe's radial stock the compile exchanges the pointwise order on
  beliefs with inclusion of bodies; from the hyperspace storey up, the order
  of the universe is inclusion. The site dimension runs through every storey
  above the site: each object is carried window by window and glued along the
  site's coverage (\autoref{sec:su-carrier}~and~\autoref{sec:su-sheaf}).}
\label{fig:su-storeys} \end{figure}

In our Quantitative Temporal Specification Logic (QTSL), open Systems and their
Specs denote in one universe, at consecutive storeys of one tower
(Figure~\ref{fig:su-storeys}). A System denotes a \emph{belief}: a lower
semicontinuous functional on probability distributions over records, whose
value at a conceivable joint distribution is the level of \emph{inconsistency}
between that distribution and what the System asserts --- a reading adopted
from Richardson and Halpern~\cite{richardson20}, and the notion of belief this
book commits to throughout~\cite{dalrymple26imprecise}. Via the \emph{radial
compile}, each belief denotes a closed radial \emph{body} in the space of
subprobability distributions over records. The \emph{hyperspace} of closed \emph{sets} of
closed bodies is the next storey; Systems are singleton elements of the
hyperspace. A Spec denotes a \emph{constructible class} within the hyperspace,
i.e., an element of the same hyperspace's finitary constructible algebra, which
constitutes the fibre of Specs. At the hyperspace storey and above, the order
of the universe is \emph{inclusion} --- membership of a singleton in a class,
containment of one class in another --- and ultimately, $\left\llbracket
\mathrm{System} \vDash \mathrm{Spec} \right\rrbracket \coloneqq \left(
\left\llbracket \mathrm{System} \right\rrbracket \subseteq \left\llbracket
\mathrm{Spec} \right\rrbracket\right)$ (like Willems' ``behavioral
approach''~\cite{willems07}, but at full imprecise-probabilistic altitude). The
pointwise order on beliefs is the same order seen one storey below, exchanged
with inclusion by the compile (\autoref{sec:su-hyperspace}), although
specifications which arise through this compile alone can only express safety,
not liveness (the incorporation of liveness specs crucially relies on the
``hit-and-miss'' aspect of hyperspace topology).

Wiring is composition in a double category whose base holds interfaces and
wiring patterns and whose fibres hold Specs --- the \textbf{double category of
System Specs}, assembled by a \textbf{double Grothendieck construction} whose
obligations are formally-verified end to end (\autoref{sec:su-composition}) and
checked with Lean as well as SafeVerify and \texttt{nanoda}. The chapter is
written against a single mechanized development, and each storey's definitions
cite their machine counterparts by module and theorem name beside the
statement.

\paragraph{How to read this chapter.} Figure~\ref{fig:su-storeys} is the map. The
\emph{causal order} fixes the events and the two relations --- causality and
chronology --- from which the temporal geometry is made
(\autoref{sec:su-causal-order}). The \emph{site} is derived from that order ---
its windows are causal diamonds and its coverage is a theorem, not a choice ---
and, equipped with duration and tangency, governs how records glue into
trajectory spaces (\autoref{sec:su-site}). A \emph{record} is an observation
history over a window's stages, a classical duration-carrying history along
each of its worldlines, and every temporal clause of the specification logic
evaluates along a single such timeline (\autoref{sec:su-records}); the record
carrier over each window of the site is a compact Fell hyperspace, with gluing
along the interior coverage --- directed covers, at fixed state fibre --- a
homeomorphism onto a projective limit (kernel-checked,
\autoref{sec:su-carrier}). The \emph{universe of beliefs} holds the
denotations of Systems: lower semicontinuous level-of-inconsistency
functionals on probability distributions over records, ordered pointwise
(\autoref{sec:su-denotation}).
The \emph{hyperspace} holds their bodies, with Systems as singleton elements
and inclusion as the order (\autoref{sec:su-hyperspace}). The \emph{fibre of
Specs} holds the constructible classes, on which conjunction is intersection
and disjunction is union, the latter forced (\autoref{sec:su-classes}). The
\emph{composition construction} carries the semantics of wiring: decoration
functors on the interfaces-and-wirings base --- the authoritative total
carrying coupled-DNF presentations of constructible classes as fibre data ---
each total a pseudo double category with a strict projection that is a discrete
double opfibration, all of it kernel-checked (\autoref{sec:su-composition}).
The storeys then assemble: the \emph{Semantic Universe} itself, the chapter's
headline object, is stated whole as the iterated total --- Specs over
interfaces over site-profiles --- in \autoref{sec:su-total}. And a final
section before the receipts (\autoref{sec:su-checker-anchor}) is the adapter:
how the wall fragment --- the fragment the shipped checkers actually certify
--- embeds into the universe, stated only after the universe is complete on its
own terms. Four families of material can be deferred without losing the thread.
The generalization columns of \autoref{sec:su-site} --- causal diamonds,
general relativity, open quantum systems, branching time --- exist to show that
the site is a genuine parameter; unless those instantiations are your interest,
what later sections consume from them is only the site contract (F1--F6) at
that section's close. The record carrier (\autoref{sec:su-carrier}) may be
skimmed until its final subsection, which lists the four interfaces everything
downstream consumes. The mechanization remarks are receipts, not exposition:
each names the machine-checked counterpart of the statements beside it, and a
reader who trusts the ledger can pass over them. And the three sections after
the assembled universe serve readers arriving with specific questions ---
adequacy (\autoref{sec:su-adequacy}) for the modelling catalogue's squares, the
wall fragment (\autoref{sec:su-checker-anchor}) for how possibilistic verdicts
embed, and the receipts (\autoref{sec:su-receipts}) for the evidence classes.


\section{The Causal Order} \label{sec:su-causal-order}

The base storey of the universe is the causal order, from which the site is
made. The axioms our construction requires are exactly the order-theoretic
fragment of bicontinuity isolated by Martin and Panangaden~\cite{martin04}. In
classical $\mathbb{R}$-valued time, $\ll$ is exactly $<$; in relativistic
Minkowski spacetime, $\ll$ denotes reachability by a time-like worldline
(strictly less than lightspeed) with nonzero spacetime distance, while $\le$
denotes reachability by any worldline (including lightspeed or zero spacetime
distance). In discrete time, both $\ll$ and $\le$ are interpreted by $\le$.

\begin{definition}[Causal order]\label{def:su-causal-order}
(\texttt{QtslAux.DiamondSite}: \texttt{ChronologyKit})
A causal order consists of \begin{enumerate}
  \item A set $E$ of \emph{events};
  \item a partial order $\le$ of \emph{causality}; and
  \item a relation $\ll$ of \emph{chronology}, such that
  \item $p \ll q \Rightarrow p \le q$ (chronology is stricter than causality),
  \item $p \le q \ll r \Rightarrow p \ll r$ and dually
  \item $p \ll q \le r \Rightarrow p \ll r$,
  \item $p \ll q \Rightarrow \exists s.\ p \ll s \wedge s \ll q$,
  \item $\forall p,r,s.\ p \ll r \wedge p \ll s \Rightarrow \exists q.\ p \ll q \wedge q \le r \wedge q \le s$, and
  \item $\forall p,r,s.\ r \ll p \wedge s \ll p \Rightarrow \exists q.\ q \ll p \wedge r \le q \wedge s \le q$.
\end{enumerate}
\end{definition}

The \emph{cone order} of special relativity --- events $(t, x)$ in time
$\times$ space, with $p \le q$ when $\lVert q_x - p_x \rVert \le q_t - p_t$ and
chronology the strict cone --- satisfies every axiom over a real normed space
of any dimension (kernel-checked: \texttt{QtslAux.DiamondTopology},
\texttt{minkowskiKit}); there the order is genuinely partial --- distinct
simultaneous events are causally incomparable, so no foliation linearizes it.
On a globally hyperbolic spacetime, the causal and chronological relations
$J^+$ and $I^+$ of the metric satisfy the axioms with room to spare: this is
Martin and Panangaden's theorem \cite{martin04}, which delivers bicontinuity
and more --- their interval topology \emph{is} the manifold topology, so at
that rung the order knows the topology of spacetime outright. And at the
opposite pole a \emph{discrete} order --- integer ticks --- admits no
\emph{strict} chronology at all (interpolation fails between consecutive ticks;
kernel-checked), a degeneration the next section turns into a feature rather
than an exclusion.

Two boundary facts orient everything downstream. First, the \emph{ladder}: how
much of a spacetime its causal order determines is graded by the classical
causality hierarchy, and this book stands on two rungs of it --- strong
causality, exactly the condition under which the order's interval topology
(generated by the chronological diamonds $I^+(p) \cap I^-(q)$) is the manifold
topology, a theorem Martin and Panangaden quote from Penrose \cite{martin04};
and global hyperbolicity, the rung at which reconstruction from the order is
complete. Second, the \emph{limit} of the method: by the theorems of
Hawking--King--McCarthy and Malament \cite{hkm76, malament77}, the causal order
of a distinguishing spacetime determines its topology, its differential
structure, and its \emph{conformal} class --- and nothing more. Durations are
precisely the conformal factor the order forgets. Everything the order does
know is \emph{derived} in the next section (Definition~\ref{def:su-spacetime-site})
--- windows, coverage, and the coverage's well-behavedness are theorems of
$(\le, \ll)$, not choices --- and everything it cannot know, \textbf{duration}
and \textbf{tangency}, enters there as additional declared structure. The
division of data is a theorem in both directions, and the structured-site
definition of the next section is organized around it.


\section{Spacetime as a Structured Site} \label{sec:su-site}

The site is \emph{derived} from the causal order --- windows and coverage are
theorems of $(\le, \ll)$ --- while \textbf{duration} and \textbf{tangency} are
declared on top. This section gives the derived object first
(Definition~\ref{def:su-spacetime-site}), then the classical instantiation --- the
value at which every classical artefact of this book lives --- and the
factored form of the construction, so that what is derived and what is
declared stay visibly separate.

\subsection{The Derived Site}

\begin{definition}[Spacetime site]\label{def:su-spacetime-site}
(\texttt{QtslAux.DiamondTopology})
Let $(E, \le, \ll)$ be a causal order (Definition~\ref{def:su-causal-order}). The
\emph{spacetime site} of $(E, \le, \ll)$ is the site whose
\begin{enumerate}
\item \emph{objects} are the causal diamonds
  $[a,b] = \{c \in E : a \le c \le b\}$ with $a \le b$, ordered by
  containment --- morphisms are inclusions of subdiamonds; and whose
\item \emph{coverage} is the \emph{derived interior topology}: a sieve on
  $[a,b]$ covers precisely when it contains the inclusion of every
  chronologically interior subdiamond --- every $[c,d] \subseteq [a,b]$ with
  $a \ll c$ and $d \ll b$.
\end{enumerate}
No further data enter: the spacetime site is a function of the causal order.
\end{definition}

That the coverage is a Grothendieck topology is a theorem of the causal-order
axioms, kernel-checked in closed form (\texttt{diamondTopology}): pullback
stability is the two absorption laws (axioms 5--6), and local character is
density (axiom 7) applied at both endpoints. The consumption is exact --- the
topology uses only axioms 4--7 of Definition~\ref{def:su-causal-order}, never the
filteredness axioms 8--9, which feed the descent arguments of
\autoref{sec:su-carrier} instead. The order-theoretic demands factor as:
\emph{interpolative mixing buys the coverage; filteredness buys descent}. (The
mechanized
statement quantifies over the full kit while its proof consumes only the
coverage half; the factoring is recorded in the module.)

On a diamond that is not chronologically thick ($a \not\ll b$) the covering
condition is vacuous, so every sieve covers --- in particular the empty one,
which forces every sheaf to a point on that diamond: the derived form of the
classical excision of duration zero (\texttt{mem\_of\_not\_thick}).
More generally the chronology acts as a \emph{descent dial}. The degenerate
choice $\ll\, =\, \le$ is admissible and its derived topology is the
\emph{coarsest Grothendieck topology} --- only the maximal sieve covers, so
sheaves are exactly presheaves and no gluing is ever demanded
(\texttt{mem\_diamondTopology\_degenerate\_iff}); the discrete order refuses
the strict chronology outright (interpolation fails between consecutive ticks;
kernel-checked, \texttt{int\_lt\_no\_chronologyKit}), so integer ticks sit at
the dial's coarsest end; dense orders admit the strict chronology, whose covers
are the genuinely interior ones and whose sheaf condition is the gluing this
chapter uses.

Four instances calibrate the definition. The value at
$(\mathbb{R}, \le, <)$ is the classical interval-domain site, identified
cover-for-cover with the site of Schultz and Spivak through the homogeneity
quotient below (Remark~\ref{rem:su-classical-site}). The value at the cone order
is the Minkowski site, kernel-checked over a normed space of any dimension
(\texttt{minkowskiKit}). The value on a globally hyperbolic spacetime exists
by Martin and Panangaden's theorems \cite{martin04} --- cited rather than
mechanized, since the site consumes only the order axioms and Lorentzian
geometry never enters the development. And integer ticks ride the degenerate
dial setting, the strict one being refused by their order.

A spacetime site arrives with its stage data for free --- a
diamond's stages are its own events under the induced order,
restricted along inclusions --- so it supplies outright the site and
stage clauses (T1)--(T2) of the structured definition this section
builds toward. What it cannot supply is exactly what
\autoref{sec:su-causal-order}'s closing division of data says it
cannot: duration and tangency, which
Definition~\ref{def:su-structured-site} declares on top.

\subsection{The Classical Instantiation}

The site every classical
artefact of this book lives on is the value of
Definition~\ref{def:su-spacetime-site} at the causal order
$(\mathbb{R}, \le, <)$, taken through one further, deliberate step
--- a homogeneity quotient. Schultz and Spivak
\cite{schultz17} take as base the interval domain $\mathbb{IR}$:
closed intervals $[d,u]$ ordered by reverse inclusion, with
$[d',u'] \ll [d,u]$ precisely when $d' < d$ and $u < u'$. The coverage
is generated by the way-below relation alone --- an interval is
covered by the family of its strictly interior subintervals --- and
behaviour types are sheaves for that coverage: a behaviour over a
window is exactly a compatible family of behaviours over its interior
subwindows. Two design decisions of theirs matter for what follows.
First, the finer coverage that would let two behaviours meeting at a
single endpoint glue is deliberately omitted; gluing across abutting
windows is a property a behaviour type may have, not one the site
imposes. Second, they quotient by translation: the quotiented site's
objects are \emph{durations} $\ell \ge 0$, and absolute time survives
the quotient as an internal object $\mathsf{Time}$ --- the behaviour
type of a clock. The quotient does not delete absolute time; it
\emph{demotes} it, from base structure to component. That demotion is
the categorical form of a commitment the model-type inventory
(\autoref{ch:qtsl-mathematical-models}) states operationally: clocks are
internalized as mechanism state, and no artefact of the toolkit
carries a global timeline field.

\paragraph{The construction, factored.} Read as a recipe, the
classical construction splits into two stages with very different
appetites, and the split is what makes ``time as a site parameter''
more than a slogan:

\begin{center}
\small
\begin{tabular}{@{}p{0.17\linewidth}p{0.43\linewidth}p{0.30\linewidth}@{}}
\toprule
\textbf{stage} & \textbf{required structure} & \textbf{output} \\
\midrule
order-theoretic core & a continuous predomain & the way-below
coverage and its sheaf topos \\
homogeneity quotient & a coverage-compatible free action whose
projection is a discrete bifibration with flat descended covers ---
or, directly, a continuous quotient category & duration objects;
absolute time demoted to the internal clock \\
\bottomrule
\end{tabular}
\end{center}

The unquotiented stage is purely order-theoretic:
the way-below coverage is definable on an arbitrary continuous
predomain, and the sheaf topos exists with no real arithmetic, no
field structure, no metric anywhere in the site data. The homogeneity
quotient is the demanding stage; here is its requirement, stated
formally.

\begin{definition}[Homogeneity quotient]\label{def:su-homogeneity-quotient}
Let $(\mathcal{T}, J)$ be a site and let a group $G$ act on
$\mathcal{T}$ by coverage-preserving automorphisms (the image of a
$J$-covering family under each $g \in G$ is $J$-covering), freely on
objects. Suppose:
\begin{enumerate}
\item[(Q1)] the orbit projection
  $\pi \colon \mathcal{T} \to \mathcal{T}/G$ onto the quotient
  category is a \emph{discrete bifibration} --- every morphism of
  $\mathcal{T}/G$ has exactly one lift along $\pi$ through any
  chosen lift of its source, and exactly one through any chosen lift
  of its target; and
\item[(Q2)] the \emph{descended families} --- images $\pi(F)$ of
  $J$-covering families $F$ --- generate a coverage $J/G$ on
  $\mathcal{T}/G$ (\emph{flat descent}: a pullback of a descended
  cover along any morphism of $\mathcal{T}/G$ is refined by a
  descended cover).
\end{enumerate}
The \emph{homogeneity quotient} of $(\mathcal{T}, J)$ by $G$ is the
site $(\mathcal{T}/G,\ J/G)$.
\end{definition}

The input of the construction is the \emph{action}, not a duration
--- at this generality no metric structure appears anywhere in the
hypotheses. Duration names the \emph{output} in the favourable case:
when $\mathcal{T}$ is a spacetime site whose duration structure (T3)
is $G$-invariant and is a \emph{complete} orbit invariant, the
quotient's objects are canonically the durations, which is exactly
the classical situation --- $G = (\mathbb{R}, +)$ translating the
interval domain, the orbit of $[d,u]$ determined by $\ell = u - d$,
morphism orbits the shave-off pairs $[r \mid s]$, and (Q2)
discharged through lattice structure on the translation parameters
(coproducts of wavy arrows taken by componentwise maxima) --- one
sufficient realization, not part of the requirement. Neither
direction of the association is automatic: a site may carry a
perfectly good duration and admit no qualifying action (nothing to
quotient by), and a quotient may exist whose objects no scalar
names --- Minkowski diamonds under spacetime translations alone
quotient to displacement \emph{vectors}, and only the further
Lorentz quotient collapses the invariant to the scalar proper time.
A candidate time base therefore pays for time-homogeneity separately
from existence --- and the two generalization columns below sit on
opposite sides of that bill. (The classical value of the quotient is
kernel-checked through the identification of
Remark~\ref{rem:su-classical-site}; the definition at this
generality is stated in prose, an obligation of the ledger if a
second instantiation ever consumes it.)

Both claims in the first row are now theorems at the generality of
Definition~\ref{def:su-spacetime-site}: the way-below coverage is
that definition's derived interior topology --- on diamonds,
way-below \emph{is} componentwise chronology, the theorem the
causal-diamond column below records --- and its well-behavedness,
its consumption of the order axioms, and its descent dial are
theorems of the causal-order axioms, stated with that definition.
The homogeneity
quotient of the second row is the classical column's own step, and
everything the table demands of it is demanded of the quotient
alone.

\subsection{Duration, Tangency, and the Structured Site}

Two additional structures are necessary to feed the two families of
proof obligations the certification interface emits.

\emph{Duration} is the metric-temporal axis: it types deadline windows
$\mathsf{F}_{\le T}$ and until windows, the spectral shift $\theta$ of
a decay certificate (units: inverse duration), and the rate and jitter
bounds carried by clock morphisms. In the classical site duration is
native and visible --- the quotient's objects literally are durations.

\emph{Tangency} is the differential axis. The ordinary instance
first: for an ODE $\dot x = u(x)$ on a state space $S$, the dynamics
assigns $u(x) \in T_x S$, a certificate sees it through the
derivative pairing $dV_x(u(x))$, and a tangency abstracts the
assignment $x \mapsto T_x S$. In general --- in the sense of
Capucci and Myers \cite{capuccimyers2026} (following
Libkind and Myers' doctrine of dynamical systems) --- it is a base
category of state spaces, a fibration of \emph{action bundles} over
it, and a section assigning to each state space its bundle of
admissible changes. A dynamics is a map \emph{into} the bundle the
tangency selects; the infinitesimal generator --- the object the
certification interface actually touches --- is that dynamics
presented through its induced action on certificate observables,
which in the ODE instance is exactly $V \mapsto dV(u(\cdot))$ and in
other doctrines is the doctrine's own generator form (a jump
kernel's integral operator, an open-quantum generator's trace
pairing). At this generality,
tangencies whose bundles are powersets or distribution spaces are
legitimate, which is exactly the slack a jump kernel or an
open-quantum generator occupies without any change to the site.

\begin{definition}[Structured spacetime site]\label{def:su-structured-site}
A \emph{structured spacetime site} is a site equipped with window,
duration, tangency, and stage-space structure, as follows.
\begin{enumerate}
\item[(T1)] \emph{The site.} A small category $\mathcal{T}$ together
  with a coverage $J$ assigning to each object a collection of
  covering families. Objects of $\mathcal{T}$ are called
  \emph{windows}; the coverage governs how behaviours --- and
  in particular records --- over the members of a covering family
  glue to a behaviour over the covered window. The site enters the
  semantics through this gluing role only; no satisfaction clause of
  the specification logic consults it.
\item[(T2)] \emph{Underlying stages.} A functor
  $|{-}| \colon \mathcal{T} \to \mathbf{Pos}$ assigning to each window
  $W$ a partially ordered set $|W|$ of \emph{stages} and to each
  morphism $u \colon W' \to W$ an order-embedding
  $|u| \colon |W'| \hookrightarrow |W|$. A \emph{timeline} in $W$ is a
  nonempty chain (totally ordered subset) $\gamma \subseteq |W|$.
  When the site is the spacetime site of a causal order
  (Definition~\ref{def:su-spacetime-site}), this clause is supplied
  outright: $|[a,b]|$ is the diamond's own event set
  $\{c : a \le c \le b\}$ under the restriction of $\le$, and
  $|u|$ for a subdiamond inclusion is the inclusion of event sets. In
  the classical instantiation ($\mathbb{R}, \le, <$) every $|W|$ is
  then totally ordered
  and is its own maximal timeline; in relativistic instantiations
  $|W|$ is genuinely partial and its timelines are the (segments of)
  worldlines it contains.
\item[(T3)] \emph{Duration structure.} For each window $W$ and each
  timeline $\gamma$ in $W$, a duration functor $\delta_\gamma$ on
  $(\gamma, \le)$ satisfying (R1)--(R2) of
  Definition~\ref{def:su-record} in the next section --- reflexivity
  at zero, and additivity along the order. The per-timeline indexing
  is load-bearing, not a hedge: a single duration functor on all of
  $|W|$ would force $\delta(s,u) = \delta(s,t) + \delta(t,u)$
  through \emph{every} intermediate stage $t$ --- path-independence
  of elapsed time --- and proper time in a relativistic window is
  path-dependent (the twin paradox; the reverse triangle inequality
  makes the straight worldline strictly longest). Duration is
  additive on the nose \emph{along} each worldline and only there,
  which is why (T3) is declared per timeline, subject to two
  naturality conditions:
  \begin{itemize}
  \item[(D1)] (\emph{restriction}) if $\gamma' \subseteq \gamma$ are
    timelines in $W$ with $\gamma'$ order-convex in $\gamma$, then
    $\delta_{\gamma'}$ is the restriction of $\delta_\gamma$
    (convexity is load-bearing: the two-stage sub-timeline of a bent
    worldline must be free to disagree with the bent valuation, else
    additivity returns through every intermediate stage --- the
    collapse this clause exists to refuse; window restrictions are
    always convex, stage sets being diamonds);
  \item[(D2)] (\emph{transport}) for every morphism
    $u \colon W' \to W$ and timeline $\gamma$ in $W'$, the embedding
    $|u|$ carries $\delta_\gamma$ to $\delta_{|u|(\gamma)}$:
    durations of pairs of stages are preserved.
  \end{itemize}
\item[(T4)] \emph{Tangency structure} (in the sense of Capucci and
  Myers, following Libkind and Myers' doctrine of dynamical systems).
  A \emph{tangency} is a triple
  $(\mathcal{B},\ \pi \colon \mathcal{E} \to \mathcal{B},\ \mathrm{T})$
  consisting of a category $\mathcal{B}$ of state (and observation)
  spaces, a cloven Grothendieck fibration
  $\pi \colon \mathcal{E} \to \mathcal{B}$ of \emph{action bundles},
  and a section $\mathrm{T} \colon \mathcal{B} \to \mathcal{E}$ of
  $\pi$, assigning to each state space $S$ its bundle
  $(\mathrm{T}_s S)_{s \in S}$ of \emph{admissible changes}. (The
  Roman $\mathrm{T}$ is reserved for the tangency section; the italic
  $T$ of Definition~\ref{def:su-record} is a record's stage set.)
  A \emph{dynamics over the tangency} is a map \emph{into} the bundle
  the tangency selects: on state space $S$ with input/parameter space
  $A$, an assignment $u \colon (s, a) \mapsto u(s,a) \in \mathrm{T}_s S$.
  When the tangency is the ordinary tangent-bundle tangency of a
  smooth state space, the induced action on a certificate observable
  $V$ is the pairing $V \mapsto dV(u(\cdot))$; for the other
  legitimate tangencies --- powerset bundles (nondeterministic step
  relations), jump-kernel and distribution-space bundles,
  operator-valued generators --- the induced action is presented
  doctrine-specifically, and no single formula of that shape is
  asserted at full generality.
\item[(T5)] \emph{Stage spaces} (Fell-readiness). For each window $W$
  a topology on the stage set $|W|$, compatible with the duration
  structure in the following sense: writing
  $d_\gamma(s,t) := \delta_\gamma(s \wedge t,\ s \vee t)$ for the
  symmetrization of a timeline's valuation --- an extended
  pseudometric on $\gamma$, by (R1)--(R2) --- each $d_\gamma$ is
  continuous on $\gamma \times \gamma$, and on every timeline where
  $d_\gamma$ separates stages the subspace topology is the
  $d_\gamma$-metric topology. Under that topology the \emph{stage
  space} $S(W)$ is required locally compact, second-countable
  Hausdorff, with each morphism's embedding $|u|$ a proper
  topological embedding. (Zero-duration pairs of distinct stages ---
  hybrid jumps --- are permitted; there $d_\gamma$ separates
  nothing, and the stage space does not owe them separation:
  same-instant stages share a stage-space point, and the carrier
  presents them as one multivalued instant of the record's graph ---
  the closed-graph doctrine of \autoref{sec:su-carrier}, under which
  any jump is multivalued at the moment it happens.)\footnote{For globally hyperbolic instantiations
  (T5) is a \emph{precision} in exactly the way
  Remark~\ref{rem:su-classical-site} exhibits classically: there the
  stage-space topology is itself order-derived --- on such a
  spacetime the interval topology generated by the chronological
  diamonds \emph{is} the manifold topology (Martin and Panangaden
  \cite{martin04}) --- so the clause records a topology the order
  already knows rather than supplying a new one.}
\end{enumerate}
Axioms (D1)--(D2) say exactly that living-on-a-timeline is a notion
stable under restriction of windows and under the site's morphisms,
which is what the gluing role (T1) requires; the records themselves
--- along a timeline and over a general window --- are
\autoref{sec:su-records}'s subject.
\end{definition}

Clause (T5) feeds the record carrier of \autoref{sec:su-carrier}, an adopted
construction with kernel-checked strata. \emph{(T3) is the topology supply on
timelines}: the site's metric axis fixes the stage topology along every
separating timeline outright and constrains it (by continuity) along the
rest, and (T5) asks only for a compatible ambient extension --- in the
classical instantiation the valuation metric \emph{is} the stage topology
--- the hybrid column shares that stage space verbatim --- so the clause is
a precision there, not a new structure.
Tangency's role at the carrier tier is equally specific: it types the
one metric hypothesis the record class keeps --- a speed bound, a
trajectory whose difference quotients lie in a fixed ball of the
tangency's action bundle, measured against duration --- so the two
site axes jointly type the record layer's single metric
hypothesis, exactly as they jointly meet in the decay row
$\mathcal{A}V \le -\theta V$, which pairs a tangency-typed operator
on the left with an inverse-duration scalar on the right. Obligation
families sort cleanly: infinitesimal-generator rows
(\emph{generator rows}, hereafter) and interior rows are tangency; window
arithmetic and spectral shifts are duration; reset rows are neither
(jumps consume zero duration); nonnegativity rows are pointwise and
touch neither axis. The tangency clause is stated at this generality
in prose; the mechanized development is parametric over it --- every
kernel-checked artefact consumes tangency only through the
generator-row typing of its concrete instances --- and the abstract
object is an obligation of the ledger.

\begin{remark}[The classical instantiation, quantitatively]
\label{rem:su-classical-site}
The temporal type theory of Schultz and Spivak is the value of the
site parameter at which every classical artefact of this book lives
--- and the identification is functorial, kernel-checked at the site and
behaviour tiers: the spec
tier's classical windows map onto the quotient site by the
translation quotient, with unique lifts in both directions and
covers corresponding exactly to covers (a wavy arrow is a strictly
interior inclusion), the motivating behaviour type of trajectories
is the site tier's own trajectory presheaf up to natural
isomorphism (\texttt{SuFunctor.Site.TemporalTypeTheory}:
\texttt{anchorQuot}, \texttt{anchorTrajIso}), and behaviour
subtypes include into the double category of System Specs ---
satisfaction is behaviour membership, wire transport is behaviour
pullback, strictly (\texttt{SuFunctor.Site.BehaviorSpec}:
\texttt{ttObject}, \texttt{satAt\_specOfBehavior\_iff},
\texttt{wSpecMap\_specOfBehavior}).
\begin{enumerate}
\item \emph{Site.} The base poset is the interval domain
  $\mathbb{IR} = \{[d,u] : d \le u\}$, ordered by reverse inclusion,
  with way-below relation $[d',u'] \ll [d,u]$ precisely when $d' < d$
  and $u < u'$. The coverage is generated by the way-below relation
  alone: an interval $[d,u]$ is covered by the single family of its
  strictly interior subintervals
  $\{[d',u'] : d < d' \le u' < u\}$. The finer coverage that would
  let two behaviours meeting at a shared endpoint glue is deliberately
  omitted: endpoint gluing is a property a behaviour type may have,
  not one the site imposes. This base is itself the derived object:
  the interval domain is the diamond poset of the causal order
  $(\mathbb{R}, \le, <)$, and its interior coverage corresponds
  cover-for-cover --- through the anchored-window identification
  below --- with the derived topology of that order, kernel-checked
  (\texttt{SuFunctor.Site.CausalConcordance}:
  \texttt{mem\_interiorTopology\_iff\_derived},
  \texttt{wavy\_mem\_iff\_derived}). Read at the spacetime level,
  the identification is also a diagnosis: a time-window with all of
  space carried in the state fibre is the causal diamond of a world
  with instantaneous propagation, so as a site the classical column
  encodes absolute simultaneity --- the $c \to \infty$ degeneration,
  in which diamonds flatten to time-slabs --- and space is carried
  by the wiring and the fibres, never by the site.
\item \emph{Homogeneity quotient}
  (Definition~\ref{def:su-homogeneity-quotient}). Quotienting by translation yields
  the site whose objects are \emph{durations} $\ell \ge 0$ and whose
  morphisms are shave-off pairs $[r \mid s]$ with
  $\ell = r + \ell' + s$; the covering families are those with
  strictly positive margins $r, s > 0$. Absolute time survives the
  quotient as the internal clock object $\mathsf{Time}$, demoted
  from base structure to component.
\item \emph{Duration structure.} Native and visible: the quotient
  site's objects literally are durations, morphism data are duration
  differences, and the induced valuation on any timeline
  $\gamma \subseteq \mathbb{R}$ is $\delta_\gamma(s,t) = t - s$.
\item \emph{Tangency structure.} Ordinary calculus: state spaces are
  smooth (in practice carved by polynomial data; the exact class is
  fixed below) spaces, the tangency is the tangent-bundle section,
  a dynamics is a vector field presented as drift data, and the
  induced action on a certificate observable $V$ is the Lie
  derivative row that drift obligations bound.
\item \emph{Stage spaces.} (T5) holds on sight: $S(\ell)$ is the
  compact interval $[0,\ell]$ with the metric
  $d_\gamma(s,t) = \delta_\gamma(s \wedge t, s \vee t)$ its own
  duration functor induces by symmetrization, and the interior covers'
  pair-filteredness is kernel-checked in this repository at
  bicontinuous-poset generality (module \texttt{QtslAux.DiamondSite};
  the causal-diamond column below is the same theorem's other
  instance). The hybrid column shares this stage space verbatim:
  jumps and resets are carrier facts (multivalued instants), not
  stage-space structure, so hybrid windows are the classical
  windows.
\end{enumerate}
For contrast on the duration axis: in the causal-diamond
instantiation the site \emph{qua} poset carries only the causal
order, and duration data --- proper time along timelines --- is
strictly additional structure, supplied densely rather than
per-clock. That divergence in presentation (object-valued durations
in the classical quotient; densely supplied proper time on
worldlines in the relativistic column) is why (T3) axiomatizes
duration by its \emph{induced valuations on timelines}, the invariant
form both presentations share and the only form the satisfaction
clauses of Remark~\ref{rem:su-record-consumers} consume.
\end{remark}

\subsection{The State Sheaf, and Where the Carving Class Lives}

The window-by-window state spaces of a System assemble into a
\emph{state sheaf}
\[
  \mathcal{V} \colon \mathcal{T}^{\mathrm{op}} \longrightarrow
  \mathcal{B},
\]
a presheaf valued in the tangency's own base category, so that (T4)'s
action-bundle fibration composes with $\mathcal{V}$ to give each
window its bundle of admissible changes. A constant $\mathcal{V}$
recovers the fixed state spaces every classical artefact uses; a
genuinely varying $\mathcal{V}$ --- mode-dependent state spaces,
stagewise phase-space reductions, time-varying constraint manifolds ---
is the structural home of time-inhomogeneity, and under it the
time-inhomogeneity lemma splits: its state face is
absorbed into the sheaf itself, and its law face reprices to a
classical projective limit of compact metrizable systems along a
countable tower of windows (\autoref{sec:su-carrier}).

Three commitments govern the fibres of $\mathcal{V}$, each with its
reason, and with the division of labour stated first: the demand on
fibres is intrinsic and topological; the carving class of the next
paragraph is how a fibre \emph{certifies} it; and compactification
is an \emph{operation} applied to fibres, never a condition on them.
The commitments: \emph{locally compact fibres} (Hausdorff, second
countable --- what the carrier's Hausdorff theorem consumes, and
with second countability it is exactly Polish local compactness, so
the same demand buys the denotation stack its standing structure;
$\sigma$-compactness follows and is what makes the compactified
fibres metrizable); \emph{fibrewise cemetery} ($\mathcal{V}^{+}(W)
:= \mathcal{V}(W) \cup \{\infty_W\}$, the one-point
compactification, an operation at the carrier tier ---
escape-to-infinity becomes an ordinary observation value, and mass
leak, killing, and divergence keep their standard readings, per
window; the operation lands outside the carving class and owes it
nothing); and \emph{proper restriction maps} (what makes the
carrier's pushforward legs Fell-continuous; automatic for the
constant sheaf and for mode-restriction embeddings). For a sector
carved in $\mathbb{R}^n$, local compactness is equivalent to being
locally closed --- open $\cap$ closed --- which is checkable from
carving data; that equivalence is why the certificate below asks
local closedness per sector, and a disjoint union of such sectors
is locally compact outright, every summand clopen.

The carving class of the fibres is fixed as follows: fibres are
locally closed \emph{$\omega$-semialgebraic} sets --- countable
unions of pieces, each carved in some $\mathbb{R}^n$ by finitely many
polynomial inequalities with \emph{rational} coefficients and finite
degree, degrees unbounded
across the family --- with proper continuous countably-piecewise-polynomial maps, the local polynomial representatives again with
rational coefficients. (The coefficient pinning is definitional,
not a presentation nicety: over unrestricted real coefficients,
$\{x \ge \pi\}$ is a piece, and the exact-data property below
silently dies.)\footnote{Underneath the site's commitments sits the
mechanized ambient category --- module family
\texttt{SuFunctor.StateSpace} of the proof-artefact repository:
objects $\omega$-semialgebraic subsets
of real ambients, morphisms countably-piecewise-polynomial maps, with
finite limits and common-ambient countable coproducts held
concretely, and local closedness the membership bridge to the local
compactness the carrier consumes. Properness and continuity are
the site's two extra demands on restriction maps, made for the
carrier's sake, never fields of the ambient category's morphisms.
The state-space category proper is this category's \emph{free
countable coproduct completion}: a state space is a countable family
of $\omega$-semialgebraic sectors, each carved in its own
finite-dimensional rational ambient, with all countable coproducts
--- disjoint, universal, no common ambient asked --- finite limits
componentwise, and the exact-data story intact through one
dimension-tagged countable carving alphabet
(\texttt{SuFunctor.StateSpace.FamCat},
\texttt{FamCoproducts}, \texttt{FamLimits},
\texttt{FamCompletion}).} The class decomposes the $\omega$PAP spaces of
Huot, Lew, Mansinghka, and Staton \cite{huot23omega} along their two
independent axes --- countable piecing, and analytic rather than
polynomial carving functions --- and adopts exactly the first.
The reasons, briefly, because each matters.
\emph{Demand}: countable piecing is what discrete-mode models
actually need --- counters, queues, birth--death populations,
program-mode structure are $\mathbb{N}$-shaped, definable in no
o-minimal structure, and $\omega$-semialgebraic in degree-one pieces
on sight. \emph{Exactness}: the intermediate class is the largest for
which every piece remains finite rational data, so the certificate
stratum of Remark~\ref{rem:su-strata} keeps exact-rational rows
piece by piece, with countable piecing expressed by the language
and discharged shape-uniformly at the certificate stratum
(Remark~\ref{rem:su-strata}: plates, pooled aggregate coordinates,
count-budget rows as optional tightenings) rather than paid for by
the semantics; full
$\omega$PAP would surrender domain-exactness (piece membership and
guard side-conditions leave the exact fragment), while its
value-side content --- analytic densities, regular guards --- lives
at the law tier, where the calculus's standing hypotheses already
name it, untouched by this dial. \emph{Freedom}: the record carrier
of \autoref{sec:su-carrier} is fibre-class-blind --- its
topological development consumes only the three commitments above,
never the carving class --- so the dial can be widened later at zero
cost to any existing proof; narrowing is what would cost.
Membership in the class of fibres is a real condition, not a
formality: $\mathbb{Q} \subseteq \mathbb{R}$ is $\omega$-semialgebraic
and nowhere locally compact --- excluded by the demand, and by its
certificate ($\mathbb{Q}$ is locally closed in no neighbourhood) ---
and adoption-time checks use a catalogue of sufficient patterns
($V$ closed; $V$ a topologically disjoint union of locally closed
semialgebraic pieces --- the bag pattern of every discrete-mode
hybrid; $V$ compact). The separation of concerns above is what keeps
the class this small: the one-point compactification of an infinite
bag is not a bag under any re-presentation (compact and infinite, it
admits no infinite clopen partition, and finitely many semialgebraic
pieces would make it semialgebraic, which $\mathbb{N}^{+}$ is not)
--- so a class required to contain its own compactifications would
be forced open to unions whose pieces accumulate. No such demand is
made: the cemetery is an operation at the carrier tier, its output
compact metrizable directly from the fibre's local compactness and
second countability, and the carving class never needs to carve
$\infty_W$.

\subsection{First Generalization Column: Causal Diamonds}

For a
globally hyperbolic spacetime $(\mathcal{M}, g)$,
Martin and Panangaden \cite{martin04}
proved that the causal intervals $[a,b] = \{c : a \le c \le b\}$
under reverse inclusion form a continuous domain whose way-below
relation is componentwise chronology ($[a,b] \ll [c,d]$ iff
$a \ll c$ and $d \ll b$), and --- the theorem that turns an analogy
into an instantiation --- that the category of globally hyperbolic
posets is \emph{equivalent} to the category of interval domains. The
causal-diamond site is the verbatim transposition of the classical
coverage to this domain: cover a diamond by its chronologically
interior subdiamonds; the classical interval domain is the
$(0{+}1)$-dimensional case of the same construction. The interior
coverage is well-behaved at exactly the generality one would hope ---
for every bicontinuous poset, the interior covers of a diamond with
$a \ll b$ are nonempty and pair-filtered --- and this is
kernel-checked in this repository at exactly that generality (module
\texttt{QtslAux.DiamondSite}): the mechanization consumes only the
order-theoretic fragment of bicontinuity, no manifold and no
compactness anywhere, and the classical interval domain is literally
its dense-linear-order instance. The boundary
caveat is likewise settled: nonemptiness of the interior covers is
\emph{equivalent} to chronological thickness ($a \ll b$) --- an
iff, kernel-checked ---
so degenerate and null-related diamonds are covered by nothing, and
the subsite of chronologically thick diamonds is the coverage's
maximal domain, the verbatim generalization of the classical
excision of duration zero. (The reopening condition is named in the
closing ledger: an application with claims on the null boundary would
change the balance in favour of retaining boundary objects; nothing
classical waits on it.) The diamond column is also the natural route
to partial differential equations: a field over a spacetime region
is a behaviour over a window whose stages are events, and hyperbolic
problems are diamond-local in exactly the sense the interior
coverage expresses --- finite propagation speed is why energy
estimates live on causal diamonds. PDE dynamics enter the adequacy
strand today with space carried by the wiring rather than by the
site --- spatially extended cells over classical time windows
(\autoref{sec:su-adequacy}) --- and the diamond column is where
space moves from the wiring into the site. The certificate stratum
has a matching design: the spatial variable enters as a declared
coordinate wire, so a certificate over PDE states is a spatially
plated hypernet --- one finite object denoting a region-indexed
family of rows --- and contraction along the spatial extent is a
symbolic integral, discharged by differentiating polynomial
witnesses rather than by enumerating a mesh; region subdivision
then joins site coverage and wiring as a third composition axis,
interior to the fibre, and all three are valuation-shaped.

A generic curved spacetime has trivial isometry group, so the
homogeneity quotient is simply not exercised --- which is semantic
correctness, not damage: physics on a curved background is not
translation-invariant, and behaviour types there \emph{should} be
diamond-indexed. What replaces the internal clock is the worldline: a
timelike curve with proper time induces a map of sites
from the classical interval site into the diamond site (proper-time
intervals to the causal hulls of worldline segments), and that map is
the curved-case clock --- per-interface data, exactly as the toolkit
already carries clocks. The clock-map's site-level checks are
kernel-checked in the same module: for any strictly chronology-preserving map out of a
dense linear time order carrying the two endpoint-approximation
properties that openness of the chronological past and future
supplies, the induced diamond map is monotone for containment, lands
in the thick subsite, carries interior-cover members to
interior-cover members, and preserves covers in cofinal form. What
remains cited rather than proved is the manifold-level instantiation:
that a timelike worldline with proper time satisfies those
approximation properties in a distinguishing spacetime. Duration
structure is proper time along records; tangency is the manifold's
tangent bundle. And the division of data is sharp in a way worth
recording: by the Hawking--King--McCarthy/Malament theorems
\cite{hkm76, malament77} a chronological isomorphism between
distinguishing spacetimes --- globally hyperbolic ones qualify ---
is automatically a smooth conformal isometry: the order determines
the conformal class and nothing more. Durations are precisely the
conformal factor the order forgets, and must be supplied as
additional structure --- densely, not per-clock. ``The site is
structured'' is not a convenience of presentation; in this column it
is a theorem about what the order alone cannot know.

\paragraph{General relativity, in two grades.} The phrase ``causal
diamond'' invites a stronger reading than the diamond column
delivers, so here is exactly what it is --- and is not --- a
semantics for. On any \emph{fixed} globally hyperbolic
background, curved or flat, the \emph{bare derived site} exists
outright: the causal
and chronological relations of the metric satisfy the causal-order
axioms by Martin and Panangaden's theorems \cite{martin04}, so
Definition~\ref{def:su-spacetime-site} instantiates there with no
new machinery --- the site consumes only the order, and Lorentzian
geometry itself never has to enter the mechanization --- while the
column's supplied structures and owed obligations (duration,
tangency, the record locus over general supports) are exactly those
named elsewhere in this section, neither better nor worse off for
the curvature. Background
\emph{independence} is different in kind, and no fixed site can even
state it: in full general relativity the metric is a field, so the
causal order depends on the solution while the fields live over the
order. Because the site is derived --- a function of the order ---
that circularity becomes an ordinary, if open, indexing problem
rather than an inexpressible one: an outer storey indexed by causal
orders (equivalently, within the distinguishing class where
Malament's rigidity theorem
\cite{malament77} governs, by conformal geometries), with the fibre over
each order the behaviour-and-spec tier over its derived site; in
each fibre a solution record carrying exactly the data the order
forgets --- the conformal factor, entering as (T3) duration
structure, together with matter fields; and self-consistency as a
section condition, that the order indexing the fibre is the order
its own metric-record induces. Einstein's equations then enter as a
constraint on that locus --- which order--factor--matter triples are
admissible --- not as site structure. The composition construction of
\autoref{sec:su-composition} already indexes interfaces by site
profiles, so indexing of this shape has a precedent in this book;
the outer storey itself --- its morphisms beyond isomorphisms, its
coverage --- is future work, named in the closing ledger so the
column's boundary stays visible (the groupoid part is no
longer owed: transport along chronology-respecting
order-isomorphisms is kernel-checked, module
\texttt{QtslAux.QuantumCovariance}): fixed
backgrounds are covered, and background independence is designed
for, not built.

\subsection{Second Generalization Column: Open Quantum Systems}

\paragraph{The modest route.} An open quantum system enters as one more mechanism
profile --- Hamiltonian plus noise operators, with the standard
Gorini--Kossakowski--Sudarshan--Lindblad generator --- and its
master-equation flow $t \mapsto \rho_t$ is a deterministic continuous
behaviour valued in density operators: an ordinary behaviour type
over the \emph{existing} site. Where measurement enters, the
measurement record is a classical stochastic process, and the
record-level probability operators of
\autoref{sec:su-denotation} already speak about it. Stability
obligations transpose with the dictionary regions $\mapsto$
projections, evaluation $\mapsto$ trace pairing: the generator row becomes
an operator inequality $\mathcal{L}(V) \preceq -\gamma V +
\zeta\mathds{1}$, checked by positive-semidefiniteness machinery
the kernel already trusts for other purposes --- under this column's
stated hypotheses, which are design content and not fine print:
finite-dimensional mechanisms (or bounded generators) only, with
truncations of infinite-dimensional systems entering only under
explicitly represented truncation error. That the stated hypotheses
suffice is
not only asserted: the transport theorem behind this row ---
Born-pairing positivity, the Loewner-monotone generator-row step, the
sharp exponential bound, trace conservation as a theorem of the
trace-annihilating generator form --- is kernel-checked in this
repository at the finite-dimensional stratum (module
\texttt{QtslAux.OpenQuantum}). No site change occurs anywhere on this
route, and the column as a whole is a design with its costs
stated, graded in
\autoref{sec:su-receipts}, not a shipped instantiation. What the
route does \emph{not} deliver is
event-level claims about measurement outcomes from stability
certificates alone; those await a quantum transport theorem in
certificate form, and the gap is named debt, not fine print.

\paragraph{The same column, the deeper route: quantum states against
the site itself.} The modest route leaves the site untouched by
design. Its natural sequel asks what happens when quantum states
meet the site's own machinery --- the records of
\autoref{sec:su-records} and the coverage of
Definition~\ref{def:su-spacetime-site} --- and the sequel is
kernel-checked at the finite-dimensional stratum, on both sides. On
the record side: a family of record observables --- Hermitian,
mutually orthogonal, correlated with the branch alternatives at the
state; hypotheses deliberately weaker than the projections of the
consistent-histories literature --- forces the decoherence
functional of its branch family diagonal, so the branch weights form
an honest finitely additive probability row, the object the
belief storey consumes. That is Gell-Mann and Hartle's
records-imply-consistency direction \cite{gellmann92}, mechanized
together with its biconditional for pure exhaustive families and
witnesses
that no hypothesis can be dropped (module
\texttt{QtslAux.QuantumRecords}); the converse implication fails by
kernel-checked witness --- there is a weakly consistent branch
family admitting no record family at all --- so records genuinely
select among consistent decompositions rather than relabelling
them. On the coverage side: over a two-site observable net on
Minkowski diamonds, gluing fails inside the state body, and it
fails where entanglement lives. The Bell state and the maximally
mixed state restrict identically along every member of a spacelike
pair-cover yet differ; at pure marginals the same gluing is rigid
(a theorem, not an example); and the derived interior coverage
\emph{refuses} the failing cover, while every cover it does certify
separates (module \texttt{QtslAux.QuantumDescent}). The boundary
this section drew for order-theoretic reasons --- interior covers
only --- is the boundary quantum mechanics needs. The two sides
then meet in a single kernel-checked sentence: the state presheaf
over the net is separated at derived covers and provably not
separated at the spacelike pair-cover, and one record row
distinguishes the two state classes that the entire pair-cover
cannot (\texttt{recordShadowQ\_separates\_bell\_mixed}, module
\texttt{QtslAux.QuantumStates}) --- records carry information that
region-by-region restriction structurally misses, which is this
chapter's records thesis in theorem form.

\paragraph{The locally covariant fragment.} These are not scattered
observations; they are the kinematic axioms of algebraic quantum
field theory. Observables have been indexed by causal diamonds
since Haag and Kastler \cite{haagkastler64}, and the locally
covariant refinement of Brunetti, Fredenhagen, and Verch
\cite{brunetti01} takes as morphisms the isometric embeddings whose
images are causally convex --- causal convexity being this
section's diamond condition --- so the index category of rigorous
quantum field theory is a derived site, and has been since 1964.
One kernel-checked structure now assembles the fragment
(\texttt{LCQFTFragment}, module \texttt{QtslAux.QuantumAssembly}):
an isotone net of local algebras (functoriality); Einstein
causality in class-level form --- pointwise causally separated
diamonds have elementwise-commuting algebras; the time-slice axiom
at two tiers --- a site-exhaustion surrogate carried as a field of
the net, and the genuine grade \emph{proven rather than
postulated}: complete initial data plus a reversible mechanism
yield Cauchy determinacy at both the algebra and the state grade,
with reversibility load-bearing by kernel-checked witness (module
\texttt{QtslAux.QuantumDeterminacy}); derived-cover separation for
the state presheaf; transport of the site, its coverage, the net,
and the separation theorems along chronology-respecting order
isomorphisms (module \texttt{QtslAux.QuantumCovariance}) --- which,
by the Hawking--King--McCarthy/Malament rigidity cited above, is
transport along conformal isometries --- with determinacy
transported across the same isomorphisms; and a record family whose
shadow lands in probability rows. A two-qubit system under an
entangling coupling inhabits the whole structure, so no field is
vacuous. Two findings earned their way into the structure rather
than being designed into it. First, the fragment is a two-tier
object of necessity: an entangling mechanism between distant sites
acts instantaneously --- non-relativistically --- so its
Heisenberg-evolved observables violate class-level causality at
unequal-time spacelike pairs, and no single net can carry both
microcausality and the genuine time-slice property at this grade.
The kinematic net and the matter data are therefore separate
tiers, tied by a shadow equation --- the net is the frozen,
dynamics-off shadow of the matter data --- which is where the
locally covariant axioms themselves place dynamics: the time-slice
clause is the door through which relative Cauchy evolution and the
stress-energy tensor enter \cite{fewster15}, both deliberately on
the far side of this fragment's fence. Second, the state object is
set-valued of necessity: a functorial choice of single states is
unrealizable --- a theorem of the locally covariant framework
\cite{brunetti01} --- and the replacement that literature adopted,
the folium (a contravariant, convex-closed assignment of
\emph{sets} of states), is structurally the flat,
zero-or-infinite shadow of this chapter's credal bodies:
the state doctrine of \autoref{sec:su-denotation} was
independently discovered by mathematical physics, for physics
reasons. At matrix grade folium distinctions degenerate, and the
degeneration is itself a kernel theorem rather than an admission
(every rank-one-positive unital functional on the full matrix
algebra is represented by a unique density state), so the toy
fragment's real content is contravariance, separation
compatibility, and the graded refinement --- graded beliefs, walls,
partiality --- that a flat folium cannot state. The word
``fragment'' is load-bearing, twice. ``Locally covariant'' does not
mean ``the same physics in all spacetimes'': Fewster and Verch
proved the stronger property fails for the full class of locally
covariant theories \cite{fewster11}, and its repair, dynamical
locality, is defined through relative Cauchy evolution --- dynamics
proper, together with non-surjective embeddings between different
spacetimes and functoriality of the assembled fragment along them,
are this column's named residues in the closing ledger. And one
consequence for the outer storey of the general-relativity
paragraph is already kernel-checked at toy grade (module
\texttt{QtslAux.QuantumOuterStorey}): under a coherent
superposition of causal orders --- the quantum-switch shape, the
same two measurements in the two temporal orders --- no record
family on the order index exists (the records theorem's
contrapositive, its index read as the set of orders), while an
orthogonal pointer record classicalizes the index back to an
ordinary probability row. Which causal order obtains is a
recordable question precisely when something records it --- both
directions witnessed at toy grade --- so the records thesis of
this chapter holds unchanged when what superposes is the site
itself.

\subsection{Branching Time}
\label{sec:su-branching}

A causal order can branch: nothing in
the axioms of \autoref{sec:su-causal-order} demands linearity,
and a forest of events --- histories sharing a past and diverging
--- is one more causal order, its site delivered by
Definition~\ref{def:su-spacetime-site} like any other (with the
degenerate chronology at minimum; a genuinely interior strict
chronology exists where the branches are densely ordered and
histories diverge along open cuts --- a last shared event would
defeat the filteredness axiom at the branch point).
Branching
time therefore needs no column of its own: there is nothing to
adopt, and nothing to reject. Two boundaries remain, and a branching
order is where both show plainly. The first is that we do not admit
nested path quantifiers in the style of branching-time logics like
CTL. Our logic (\autoref{ch:qtsl-specification-language}, below)
\emph{stratifies} into an inner layer of temporal quantifiers along
single-timeline records, and an outer layer regarding the
imprecise-probabilistic ensemble of such records. The second is a boundary
of
canonicity rather than of expressibility: treating a \emph{quantum}
system as branching-time --- an Everettian forest as its causal
order, branch weights as a measure --- is non-canonical relative to
the universal-wavefunction reading, because branch structure is
decoherence-emergent and decomposition-dependent, so a certificate
over a branching order would certify a claim about a \emph{chosen}
decomposition. (The asymmetry grounding this judgement is kernel-checked in the quantum paragraphs above: records force
consistency of a branch family, and consistency alone supplies no
records.) That is why the open-quantum column above enters
through the master-equation flow over an ordinary site, with
measurement records as classical stochastic processes read by
record-level probabilities.

\subsection{The Logic--Site Contract}

The toolkit does not, at present,
\emph{adopt} a nonclassical site: no present application demands one,
and the classical instantiation must never pay for the option. What
it adopts is a contract --- six freedoms the artefact formats keep,
so that the site parameter can later be exercised without restaging a
single classical artefact:

\begin{enumerate}
\item[\textbf{F1.}] \textbf{No absolute-time fields.} Time enters
  manifests and
  specifications only as durations, per-mechanism clocks, and clock
  morphisms with rate/jitter bounds; never a timeline table, never a
  timestamp.
\item[\textbf{F2.}] \textbf{The checker never sees the site.} All
  site content
  stays on the semantic side of the certification interface;
  obligations are polynomial (or operator) rows with no time-shaped
  vocabulary.
\item[\textbf{F3.}] \textbf{Total-order consumers are flagged.}
  D\'ebut, until, and
  cofinality clauses are stated as record-indexed, so a site change
  re-indexes rather than re-defines them. One such re-indexing is
  realized rather than merely flagged by the record carrier of
  \autoref{sec:su-carrier}: endpoint consumers become germ
  consumers. Over the interior coverage a window's own boundary
  stages are covered by nothing --- data pinned to an outer seam is
  exactly what translation-homogeneity forbids, the same demotion
  that turns absolute time into the internal clock --- so a clause
  that consults a window's endpoint re-indexes as a \emph{germ}
  clause at the seam: a clause about the behaviour's germ there
  rather than its value there. The endpoint itself survives, like
  absolute time, as presentation data of an anchored tower of windows
  (a family sharing a fixed initial stage, whose shared seam survives
  every restriction), not as site data.
\item[\textbf{F4.}] \textbf{Mechanism profiles stay an open
  enumeration.} A new
  dynamics family (an open-quantum generator, say) is additive; no
  existing profile recompiles.
\item[\textbf{F5.}] \textbf{Clock morphisms keep the rate-relation
  form.} The
  wiring datum is bounds on $d\sigma/d\tau$ plus jitter, never a
  schedule table in absolute time; clock-\emph{rate} corrections
  (gravitational and velocity time dilation) then arrive as
  \emph{values} in an existing field, not as a schema change. The
  named exception: signal-\emph{path} obligations --- Sagnac terms,
  null-propagation time of flight --- are not per-clock rate data;
  they need ephemeris or worldline data plus a propagation
  constraint, still ordinary classical model data in a named frame,
  but data the rate-relation form does not carry.
\item[\textbf{F6.}] \textbf{Duration and tangency stay ambient.}
  Mechanisms present
  tangency as doctrine-specific generator data (drift, diffusion,
  jump, or operator rows) and duration as rational constants; no
  artefact names either structure as an object.
\end{enumerate}

Against the only metric that matters for a toolkit --- does the
classical user pay? --- the verdict, costs counted, is: nothing. No
extra
field, no extra obligation, no slower check, on either generalization
column. The cost of \emph{adopting} a column when an application arrives
is a named metatheory ledger, several of whose items are already
discharged ahead of any adoption --- all kernel-checked in this
repository (\texttt{QtslAux.DiamondSite},
\texttt{QtslAux.OpenQuantum}): for the causal-diamond
column, the coverage lemma, the thickness iff, the clock-map checks,
and the time-inhomogeneity item's clock-adjunction core; for the
open-quantum column, the bounded-generator hypothesis's transport
theorem. What remains
owed there is the site-generic transport of the extension theorems
--- split by the carrier of \autoref{sec:su-carrier} into a
presentation-shaped half that transports along inverse-image functors
because presentations by generators and relations pull back
(geometricity makes that half free in principle) and an analytic half
(choice-invoking or point-set-surgical constructions) that runs
objectwise over finite-duration windows, where the classical proofs
apply unchanged, so the audit is per window class, not per theorem
per site --- together with the probabilistic time-inhomogeneity
lemma's projective-limit step and the open-quantum event/record
transport theorem, each named in \autoref{sec:su-receipts}.
\section{Records}
\label{sec:su-records}

What lives on a window is an observation history over its stages: a
\emph{record}. The definition arrives in two steps --- first along a
single timeline, the object every temporal clause of the
specification logic actually reads, and then over a general causal
support, where a record is a coherent bundle of timeline records,
one along each of its worldlines.

Along a single timeline a record is a totally ordered history of
observations: the trajectory of a classical composite read on the
global timeline; a component's local history read against its own
clock; the proper-time worldline of a relativistic clock; the outcome
sequence of a monitored quantum device. What these share --- and what
the specification logic consumes --- is a totally ordered time axis
\emph{carrying duration}: the order alone is not enough --- two of
the three clause families below consume more than it: one does
arithmetic on elapsed time, and
the d\'ebut's infimum must exist in the record's own axis. The
d\'ebut of a region ($\tau_E = \inf\{t : x_t \in E\}$), the ``no
visit before $\tau$'' clause of an until operator (and its metric
variant's window $\mathsf{U}_{\le T}$), the cofinal visits of a
recurrence condition: each is well-defined along any record equipped
with the ordered time-and-duration structure of
Definition~\ref{def:su-record}, and none needs to know how that
record sits inside a larger geometry of time.

\begin{definition}[Record]\label{def:su-record}
A \emph{record} is a tuple $R = (T, \le, \delta, \mathcal{O}, x)$
consisting of the following data.
\begin{enumerate}
\item \emph{Stages.} A nonempty set $T$, totally ordered by $\le$,
  whose order is \emph{conditionally complete}: every nonempty subset
  of $T$ that is bounded below has an infimum in $T$.
\item \emph{Duration.} A \emph{duration functor}
  \[
    \delta \colon \{(s,t) \in T \times T : s \le t\}
      \longrightarrow [0, \infty]
  \]
  satisfying
  \begin{itemize}
  \item[(R1)] $\delta(t,t) = 0$ for every $t \in T$;
  \item[(R2)] (\emph{additivity along the order})
    $\delta(s,u) = \delta(s,t) + \delta(t,u)$ whenever
    $s \le t \le u$ in $T$.
  \end{itemize}
\item \emph{Observations.} An \emph{observation space} $\mathcal{O}$
  --- a set, required to be a measurable space whenever the record
  enters the denotation stack of
  Definition~\ref{def:su-credal-denotation} --- and an
  \emph{observation map} $x \colon T \to \mathcal{O}$, written
  $t \mapsto x_t$.
\end{enumerate}
Equivalently, and more precisely: a duration functor is a
\emph{functor} $(T, \le) \to \mathbf{B}[0,\infty]$ from the stage
order, read as a thin category, to the delooping of the additive
monoid $([0,\infty], +, 0)$ --- functoriality on composites is (R2),
preservation of identities is (R1). (Not an \emph{enrichment}: a
category enriched in $[0,\infty]$ is a Lawvere metric space, whose
composition law is the triangle \emph{inequality}; duration along a
timeline is additive on the nose.) From (R1)--(R2), $\delta$ is
antitone in its first endpoint and monotone in its second:
$s' \le s \le t \le t'$ implies $\delta(s,t) \le \delta(s',t')$.
Two non-requirements are deliberate, and both do real work.
\emph{Strict positivity is not required}: distinct stages may be separated by zero
duration (instantaneous transitions consume none, which is how hybrid
records carry reset events). \emph{Finiteness is not required}:
$\delta$ may take the value $\infty$ on pairs separated by an
unbounded extent.

For a subset $E \subseteq \mathcal{O}$ (measurable, when
$\mathcal{O}$ is a measurable space), the \emph{hitting set} of $E$ is
$T_E := \{t \in T : x_t \in E\}$, and the \emph{d\'ebut} of $E$ is
\[
  \tau_E \;:=\; \inf T_E \;\in\; T \cup \{+\infty\},
\]
defined to be the infimum in $T$ when $T_E$ is nonempty and bounded
below --- which conditional completeness guarantees to exist --- and
the formal top element $+\infty$ (adjoined to $T$) when
$T_E = \varnothing$. The d\'ebut is deliberately \emph{partial}: for
a nonempty hitting set unbounded below (possible only in a record
with no start), no value is assigned, and every satisfaction clause
that consumes a d\'ebut carries the bounded-below proviso. When $T$
has a least element (a record with a start), every nonempty subset is
bounded below and the d\'ebut of every region is defined.

The stage set $T$ is free data here and moored below: for a record
over a window of the site
(Definition~\ref{def:su-record-window}), the stage sets to which
this definition applies are the \emph{maximal chains} of the
record's support in the window's stage poset $|W|$ --- the site's
(T2) data, under the induced order --- with $\delta$ not chosen but
inherited from the site's (T3) duration structure on that timeline.
Maximality is load-bearing on both sides: the maximal chains cover
the support (so the bundle of timeline records misses no stage),
while a restriction to a \emph{non}-maximal chain need not satisfy
this definition at all --- omitting a stage can delete an infimum
and break conditional completeness. The definition is stated
free-standing so that the logic's single-timeline clauses consume
exactly this structure and no site geometry --- the record tier's
face of the site contract's second freedom.
\end{definition}

A thermostat log makes the data concrete: stages
$T = \{0, 1, \dots, n\}$, one sample per second, with
$\delta(m, m') = m' - m$ seconds and observations the measured
temperatures. For the region $E = \{x : x \ge 30\}$ the hitting set
collects the hot samples, and the d\'ebut $\tau_E$ is the first hot
second --- or $+\infty$, if the log never gets hot. A hybrid version
of the same log carries a controller reboot as two distinct stages
separated by zero duration, which is exactly what the definition's
first non-requirement permits.

\begin{remark}[The three record-indexed clause families, and what each
consumes]\label{rem:su-record-consumers}
The satisfaction clauses of the specification logic that consume a
record's temporal structure sort into exactly three families, and the
data of Definition~\ref{def:su-record} are precisely what makes each
well-defined. A bare total order does \emph{not} suffice.
\begin{enumerate}
\item \emph{D\'ebuts} ($\tau_E = \inf\{t : x_t \in E\}$) consume
  conditional completeness of the order (item 1): in a general totally
  ordered set a nonempty bounded-below hitting set need not have an
  infimum, and the d\'ebut clause would then have no referent.
\item \emph{Metric-bounded windows} (the deadline operator
  $\mathsf{F}_{\le \ell}$ and metric until clauses) consume the
  duration functor (item 2): the window of extent $\ell$ anchored at
  $t$ is $\{s \in T : t \le s,\ \delta(t,s) \le \ell\}$, and no such
  set is definable from the order alone.
\item \emph{Cofinal recurrence} (``$x$ visits $E$ cofinally'':
  for every $t \in T$ there is $s \ge t$ with $x_s \in E$) consumes
  only the order, through its cofinal structure; it is definable in
  any total order and needs neither infima nor durations.
\end{enumerate}
Both standard instantiations satisfy all hypotheses at no cost:
discrete records ($T \subseteq \mathbb{N}$, $\delta(m,n) = n - m$ or a
per-step duration sum) are well-ordered, hence conditionally complete;
interval records ($T$ a real interval, $\delta(s,t) = t - s$; more
generally proper time along a worldline) are conditionally complete by
completeness of $\mathbb{R}$.

Division of labour, recorded here because the site supplies
the other half: \emph{the record carries the duration structure, while
the mechanism carries the tangency structure}; the two meet in the
generator row, where a tangency-typed operator is compared against an
inverse-duration-typed scalar.
\end{remark}

Definition~\ref{def:su-record} is stated along one timeline, and on
the classical site --- where a window's stages are totally ordered,
so any support lies along one timeline automatically --- it is the
whole story. Over a general spacetime site it cannot be: two facts
force a wider notion, and between them fix its shape. The first fact
is geometric. In the Minkowski site a chronologically thick
diamond's stage space contains causally incomparable pairs ---
kernel-checked as a witness in the canonical unit diamond, over any
normed space of positive dimension
(\texttt{QtslAux.DiamondTopology}:
\texttt{minkowski\_diamond\_spacelike\_pair}) --- so a record
supported on one timeline cannot span such a window: the restriction of
a field to a diamond, the honest ``trajectory'' of the hyperbolic
column, is not a timeline object, and total stage order stands
revealed as the foliated $(0{+}1)$-dimensional shadow --- the same
diagnosis Remark~\ref{rem:su-classical-site} gives the classical
site itself. The second fact is metric, and it is what makes the
wider notion \emph{forced} rather than designed. In a relativistic
order the natural duration between causally related events --- the
Lorentzian distance, the \emph{supremum} of proper time over the
causal curves joining them --- satisfies the
\emph{reverse} triangle inequality: it is superadditive, so the
additivity axiom (R2) fails for any single global valuation built
from it. What is additive is a different functional --- the
proper-time integral along a \emph{fixed} worldline, one functional
per timeline --- so
duration must enter timeline by
timeline: precisely the form the site's duration
structure (T3) already takes, the classical definition's global
$\delta$ being its one-timeline shadow.

\begin{definition}[Record over a window]\label{def:su-record-window}
Let $W$ be a window of a structured spacetime site, with stage poset
$|W|$ and duration functors $\delta_\gamma$ on its timelines
(T2--T3 of Definition~\ref{def:su-structured-site}). Let
$R_0 = (T_0,\ \le,\ \delta_0,\ \mathcal{O},\ x)$ carry the data of
Definition~\ref{def:su-record} per maximal chain as below, together
with a monotone \emph{stage map} $m \colon T_0 \to |W|$ whose fibers
are the same-instant stacks (order-convex, $\delta_0$-null) and
which carries each maximal chain of $T_0$ onto a timeline of $W$
with $\delta_0 = \delta_\gamma \circ (m \times m)$ there. The pair
\[
  R \;=\; (R_0,\ m)
\]
is a \emph{record over $W$} when, for every maximal chain $\gamma_0$
of $T_0$, the restriction
$(\gamma_0,\ \le,\ \delta_0,\ \mathcal{O},\ x|_{\gamma_0})$
is a record in the sense of Definition~\ref{def:su-record}.
\end{definition}

A record over a window is a coherent bundle of its own worldline
shadows: every maximal-chain restriction is a classical record,
matched onto the window's timeline against the site's own valuation
--- and the record's carrier presentation is the closed set
$\mathrm{cl}\{(m(t), x_t)\}$ --- single-valued exactly off the
stacks and jump instants, multivalued at them, with
single-valuedness a certified stratum rather than a membership
condition (the general-support countermodels of the closing ledger
force the same demotion independently). (Maximal chains pass through every stage of
the support --- Hausdorff's maximal principle, part of the ambient
classical mathematics --- so the bundle covers the record; and the
restriction to a \emph{non}-maximal chain need not be a record,
conditional completeness not being hereditary, which is why the
definition tests exactly the maximal ones.) The record itself is
site-relative data: the pair $(T, x)$ carries no duration, and its
timeline shadows become classical records only against the window's
own valuations $\delta_\gamma$, supplied by (T3) before any clause
runs. The clause families of
Remark~\ref{rem:su-record-consumers} relativize accordingly,
timeline by
timeline: d\'ebuts and metric-bounded windows evaluate on the
shadows $(\gamma, \le, \delta_\gamma, \mathcal{O}, x|_\gamma)$,
where conditional completeness and
(R1)--(R2) hold by definition, and cofinal recurrence reads the
support's own order directly. Each shadow is a free-standing
classical record, and the clauses read nothing but the shadows ---
duration reaches the logic packaged inside them, and the coverage,
the site structure the invariant of (T1) fences off, enters
nowhere. (Which timelines a claim quantifies over --- some worldline
through the record, or every one --- is spec-level structure one
storey up, not record data; none of the three clause families ever
needed a global d\'ebut, and should one ever be wanted, its
referent is the entry \emph{antichain} of a hitting region rather
than an infimum stage.)

\emph{Conservativity}: over the classical site a window's stages are
totally ordered, so every subset is a chain and a stack-free
record's support ($m$ injective) has exactly one maximal chain ---
itself ---
and Definition~\ref{def:su-record-window} asks
exactly that $(T_0, \le, \delta_0, \mathcal{O}, x)$ be a record in
the sense of Definition~\ref{def:su-record} against the window's own
valuation $\delta_0(s,t) = t - s$ read through $m$. That is precisely the valuation
the classical artefacts' records carry
(Remark~\ref{rem:su-classical-site}), so those
records \emph{are} records over their windows, with nothing
re-indexed and nothing added: over a totally ordered window the two
definitions coincide. \emph{Two support classes,
one definition}: a record supported on a single chain is a
\emph{worldline probe} --- an observer's history through the window
--- while a record with chronologically thick support is a
\emph{field restriction}, a PDE trajectory over a diamond being the
graph of the field over the diamond's interior; and reading a record
along one of its timelines --- the observer's act --- is literally
the definition's restriction component. A clock map, in particular,
remains what the diamond column says it is: a way of \emph{reading}
a record along a worldline, never part of what a record is.

The organizing principle of the temporal layers --- call it the
\emph{records thesis} --- is the resulting division of labour:

\begin{quote}
Path formulas are evaluated along records, and consult only the
structure a single record carries --- its order and its durations.
The site governs how records glue into a System's trajectory space,
and only that: its cross-record geometry enters no satisfaction
clause. The denotation says what lives on the glued space.
Composition of Systems is summation of their beliefs, each
transported to the joint records along the wiring: the composite's
belief is the sum of its parts.
\end{quote}

Each clause of the thesis bounds the blast radius of the others.
Because satisfaction consults only record-level structure,
generalizing the time base cannot silently change what any
specification means; the three satisfaction clauses that consume the
record's order --- d\'ebut, until, cofinality --- are flagged in the
semantics as \emph{record-indexed}, so that a new site re-indexes
them rather than re-defining them (a re-indexing
Definition~\ref{def:su-record-window} realizes: over a general
support they evaluate along maximal chains). Because the denotation is a
measure-like object \emph{over} records rather than a new kind of
record, the imprecise apparatus of \autoref{sec:su-denotation}
transports along the record direction by construction --- with the
caveat stated: its extension theorems are proven on the
discrete-time record poset, and each new site owes the transport of
exactly those theorems to its own record structure (a named
obligation of the logic--site contract of
\autoref{sec:su-site}, not a
free consequence of the
flagging). And the thesis has a sharp outer boundary, drawn in
\autoref{sec:su-branching}: quantification \emph{across} records is
admitted only at the specification logic's outer layer, as
imprecise-probabilistic claims about the ensemble --- nested path
quantifiers in the style of branching-time logics have no inner-layer
home, by the same stratification.

One reading of records is fixed here because everything in
\autoref{sec:su-carrier} depends on it: the
\emph{speed-bounded} records over a window of the site --- those
whose difference quotients, along each timeline of the support, lie
in a fixed ball of the tangency's
action bundle, the single metric hypothesis the record layer keeps
--- are read \emph{as their graphs}, closed subsets of a
stage-times-state ambient, and the class of such records over a
window carries the Fell topology on closed sets. (A general record
--- of Definition~\ref{def:su-record} or of
Definition~\ref{def:su-record-window} --- carries no topology and no
closed-graph condition; the graph reading is a property of the
speed-bounded class, not of the notion.) Under that reading the
record class over each window --- at the classical stratum, where a
support is a timeline --- is compact, and gluing along the
interior coverage --- directed open covers at fixed state fibre ---
is a homeomorphism onto a projective limit. The kernel receipts
behind those sentences live at that stratum; over general supports
the graph reading transports
objectwise verbatim, while compactness rests on closedness of the
record locus --- which does not bundle per timeline: the closed
object over thick supports is an interior-pair relaxation of the
record class, its boundary conditions surviving as certified strata
--- an obligation stated with its two site
hypotheses in the closing ledger, not a free consequence. The reading is developed, with its kernel
receipts, in \autoref{sec:su-carrier}; nothing in the satisfaction
clauses above changes under it, which is the records thesis doing its
job.

\section{The Record Carrier: Fell Hyperspaces over the Site}
\label{sec:su-carrier}

The records of a window live in a definite space --- the strata to
follow quantify over exactly it. Per \autoref{sec:su-records}, a
speed-bounded record over a window is its graph, a closed subset of
a stage-times-state ambient, and the record class carries the Fell
topology on closed sets. The Fell topology is moreover \emph{the}
hyperspace topology that glues along directed open covers --- the
descent theorem below, kernel-checked --- so the record carrier is a
sheaf of compact spaces for exactly the interior coverage the site
already adopted (whose covers are directed by pair-filteredness), at
fixed state fibre.

\subsection{The Carrier, and Its Descent}

\paragraph{The ambient and the carrier.} Fix a window $W$. The
\emph{ambient} is
\[
  A(W) \;:=\; \mathring{S}(W) \times \mathcal{V}^{+}(W),
\]
the \emph{open interior} of the stage space (the design decision
justified by the seam analysis below) times the state fibre of
\autoref{sec:su-site} with its cemetery adjoined --- the one-point
operation applied here, at the tier that consumes its output; by
(T5) and the fibre commitments, $A(W)$ is locally compact,
second-countable Hausdorff, the compactified factor compact
metrizable outright.
The \emph{carrier} over $W$ is
\[
  \mathbb{H}(W) \;:=\; \mathrm{CL}\bigl(A(W)\bigr)_{\mathrm{Fell}},
\]
the space of \emph{all} closed subsets of the ambient --- the empty
set included --- under the Fell topology \cite{fell62}: the topology
generated by the classes $\mathrm{Hit}(V) = \{F : F \cap V \neq
\varnothing\}$ for $V$ open and $\mathrm{Miss}(K) = \{F : F \cap K =
\varnothing\}$ for $K$ \emph{compact}. Restriction along a site
morphism $u \colon W' \to W$ is clip-then-push:
\[
  \rho_u \;:=\;
  \bigl({-} \cap (|u|(\mathring{S}(W')) \times \mathcal{V}^{+}(W))\bigr)
  \fatsemi \bigl(\mathcal{V}^{+}(u)\bigr)_{*},
\]
clip to the open sub-band, then push the state coordinate forward
along the (proper) restriction of the state sheaf. Two lemmas make
every $\rho_u$ Fell-continuous, and each is short enough to give
whole. \emph{Clip-to-open is Fell-continuous}: for $U \subseteq Y$
open, under $F \mapsto F \cap U$ the preimage of $\mathrm{Hit}(V)$
($V \subseteq U$ open) is $\mathrm{Hit}(V)$, and the preimage of
$\mathrm{Miss}(K)$ ($K \subseteq U$ compact) is $\mathrm{Miss}(K)$ ---
both subbasic. \emph{Pushforward along proper continuous maps is
Fell-continuous}: for $f \colon Y \to Z$ proper continuous into a
locally compact Hausdorff $Z$ (so images of closed sets are closed),
under $F \mapsto f[F]$ the preimage of $\mathrm{Hit}(V)$ is
$\mathrm{Hit}(f^{-1}V)$ (open, by continuity) and the preimage of
$\mathrm{Miss}(K)$ is $\mathrm{Miss}(f^{-1}K)$ (compact, by
properness). The two subbasis families split the work exactly ---
continuity feeds the hit face, properness feeds the miss face. The
same split carries the descent theorem:

\begin{theorem}[Fell descent over directed opens]
\label{thm:su-fell-descent}
Let $Y$ be a topological space and $\{U_i\}_{i \in I}$ a directed
family of open subsets (any two contained in a third) with union $Y$.
Then the clip maps $\rho_i \colon \mathrm{CL}(Y) \to \mathrm{CL}(U_i)$
induce a homeomorphism
\[
  \mathrm{CL}(Y) \;\xrightarrow{\;\cong\;}\;
  \varprojlim_i \mathrm{CL}(U_i)
\]
onto the space of compatible families, the limit carrying the limit
(initial) topology.
\end{theorem}

\begin{proof}
Bijectivity: injectivity is $F = \bigcup_i (F \cap U_i)$; for
surjectivity a compatible family $(F_i)$ glues to
$F := \bigcup_i F_i$, closed because closedness is local, with
$F \cap U_i = F_i$ (a point of $F \cap U_i$ lies in some $F_j$;
directedness gives $U_k \supseteq U_i \cup U_j$, and
$F_k \cap U_i = F_i$ collects it). Topology: each $\rho_i$ is
continuous by the clip lemma, so the Fell topology refines the
initial topology; conversely every Fell-subbasic class of
$\mathrm{CL}(Y)$ is pulled back from a level ---
$\mathrm{Hit}(V) = \bigcup_i \rho_i^{-1}\mathrm{Hit}(V \cap U_i)$,
since a closed set hits the open $V$ iff it hits some $V \cap U_i$,
and for compact $K$ directedness places $K$ inside a single $U_i$,
whence $\mathrm{Miss}(K) = \rho_i^{-1}\mathrm{Miss}(K)$.
\end{proof}

Both subbasis families again carry the theorem, and this time the
split is the whole of it: \emph{hits are local} (an open set is the
union of its cover-traces) and \emph{misses are compactly supported}
(a compact miss-parameter lands inside one cover member, by
directedness). The Vietoris topology fails the second half --- a
miss-of-\emph{closed} parameter need not land in any cover member ---
and the refuting fixture shows the miss face sharp: over a directed
cover of $\mathbb{N}$, the miss class of a closed noncompact
parameter fails to be open in the limit topology, by a witness (a
far-enough singleton, whose clips agree with the empty record's on
any finite set of levels) that defeats every hyperspace topology
carrying the closed-miss upper face. Of the two candidate miss
faces, then, only the compact one descends, and Fell's two design
decisions (hit-of-open, miss-of-compact) are respectively the
local-character axiom and the finite-support axiom of a sheaf-ready
topology.
Interior covers of a window are pair-filtered (the
kernel-checked bicontinuous-poset lemma of
\autoref{sec:su-site}) and their
interiors union to $\mathring{S}(W)$; crossing with the covered
window's own state fibre preserves both properties. Hence:

\begin{proposition}[The carrier is a sheaf of compact spaces]
\label{prop:su-carrier-sheaf}
Consider a cover of the site's coverage satisfying the two
hypotheses the kernel theorem consumes: its members' stage
interiors, embedded in the covered window, form a \emph{directed
family of open subsets} whose union is the window's interior
(interior coverages supply this --- pair-filteredness is
directedness, and openness and exhaustion are the coverage's
construction); and its members' ambients carry the covered window's
state fibre (every cover of the classical instantiation; every
cover under a state sheaf constant on the window being covered).
For every such cover,
$\mathbb{H}(W) \to \varprojlim_{W' \in \mathrm{cover}} \mathbb{H}(W')$
is a homeomorphism, and on that stratum $\mathbb{H}$ is a sheaf on
$(\mathcal{T}, J)$ valued in compact Hausdorff spaces. (A site
whose coverage violated the first hypothesis would owe its own
descent theorem --- the kernel fixtures below show both directedness
and the miss face sharp --- and under a genuinely varying state
sheaf the restriction legs compose the clip with a proper
pushforward, whose sheaf property is prose-grade design, not an
instance of the kernel theorem.)
\end{proposition}

Compactness arrives with no
hypotheses at all: $\mathrm{CL}(Y)$ under Fell is compact for
\emph{every} ambient $Y$ (an Alexander-subbasis argument), Hausdorff
when $Y$ is locally compact Hausdorff, and equal to the Vietoris
hyperspace on compact Hausdorff ambients. Over an unbounded ray
ambient, compactness of a record class is a theorem with hypotheses;
over the site it is a corollary of
closedness in a hyperspace that is compact for free. Nothing of
Fell's noncompact machinery is thereby discarded: it is
\emph{localized}. In the classical instantiation windows have
finite duration, so the unbounded horizon vanishes from every
object; what stays noncompact per window is the open interior
(seams) and the fibre (escape), and the compact-miss face polices
both by one mechanism --- the empty set is a closed point of every
$\mathbb{H}(W)$, isolated only over a compact ambient (a
zero-duration window has empty interior and carries the one-point
trivial carrier; every nondegenerate window's open-interior ambient
is noncompact), and over the nondegenerate windows it is the limit
of records draining out of every compact: total escape, escape out
a seam, and escape to $\infty_W$ are boundary members of the
carrier, rather than pathologies excluded by fiat.

\paragraph{The seam fork, and why sheafification decides it.}
The ambient uses the \emph{open interior} of the stage space, and
this is forced, not chosen. Restriction to a \emph{closed} subwindow
fails Fell-continuity at records with a jump exactly on the seam: a
family of records with jumps at times $u' + \tfrac1n$ converges in
Fell to the record jumping at $u'$, but the closed-clip images stay
at Hausdorff distance $\ge 1$ forever, because the limit's clip
carries a zero-duration post-jump stage on the seam that no
approximant's clip carries. (The witness is hybrid-only: for pure
flows under the speed bound, closed clipping is continuous; the fork
is exactly the price of jumps. It is kernel-checked as a refuting
fixture beside the clip lemma.) The resolution is that the
site had already made the decision: the closed-window record
presheaf is not a sheaf --- a matching family over a window's
interior cover determines the record on the union of the members'
interiors, which is the window's open interior, and \emph{nothing at
the outer seams} --- and its sheafification is precisely the
open-interior carrier above. Read with the site's own glosses this
stops being a technicality: the site quotients translation, so a
window's outer seams have no absolute identity --- data pinned to
this window's own boundary is exactly the data translation-homogeneity
forbids. Sheafification forgets the seams because homogeneity
requires it; seam data returns in the anchored-tower presentation
(a family of windows sharing a fixed initial stage), in the same
movement that returns absolute time as the internal clock, which is
the germ row of the logic--site contract (F3).

One coverage, three independent derivations, no shared premises: the
interior coverage is simultaneously the behavioural choice (Schultz
and Spivak omit endpoint gluing for composability of behaviours), the
continuity choice (closed clipping is discontinuous at seam jumps),
and the descent choice (Theorem~\ref{thm:su-fell-descent} needs
directed opens). A choice made for behavioural reasons is forced
again by continuity and a third time by descent.

\subsection{The Record Locus}

\paragraph{The record locus is a closed subsheaf.} Within each
$\mathbb{H}(W)$, the \emph{record locus} $\mathbb{R}(W)$ comprises
the closed sets that are graphs of speed-bounded records over
relatively closed sub-timelines of the window interior --- the
classical, single-timeline locus throughout this paragraph; over
the general supports of
Definition~\ref{def:su-record-window} the closed object is an
interior-pair relaxation, per the ledger's named obligation ---
unanchored and lifetime-windowed: records may begin and end strictly
inside the window (early stop and late start are membership, not
defect, in a translation-quotiented site), the empty record is a
member, and the class is Zeno-inclusive: Zeno-boundary points are
themselves honest records (a Zeno stage domain, lexicographically
ordered, is conditionally complete --- its infima land on the lowest
branch of the stack), and locally finite \emph{event}
structure is a certified stratum inside the carrier --- closed (a
continuous image of the compact parameter class; equivalently
bounded alternation between separated regions, a class built from
time-ordered finite hit demands) and disjoint from the chatter
boundary --- carried by the
crossing-budget row family, not the locus's definition: the budget
was never a compactness hypothesis; its remaining job is
mechanism-side progress --- a ghost storm (state-invisible events
accumulating) stalls a System without casting a shadow, which is
exactly why liveness rows consume the budget and safety rows never
did. The structure theorem
factors through a generic principle: call a family of per-level loci
$R_i \subseteq \mathrm{CL}(U_i)$, with ambient locus $R_0$,
\emph{local} when membership in $R_0$ is equivalent to membership of
every clip in every $R_i$. A levelwise-closed local locus is closed
(an intersection of clip-preimages of closed sets, by
clip-continuity), hence \emph{compact for free} in the hypothesis-free
compact hyperspace; compatible families levelwise in the loci glue
into the ambient locus; and the descent homeomorphism restricts to a
homeomorphism of the ambient locus onto the levelwise-locus subspace
of the projective limit --- the closed-subsheaf statement. The record
instantiation discharges the two obligations: the speed-graph
condition is closed (a violating pair has a positive defect margin
that two hit demands preserve --- pure hit face, no compactness
consumed), and it is local over directed covers, directedness
entering at exactly one step (a violating pair's two points must meet
in one cover member). All of this is kernel-checked, together with
the sharpness fixtures that bound it: the speed bound cannot be
dropped or existentially absorbed (a steepening family of affine
graphs Fell-converges to a vertical fibre, which is no graph at any
speed --- so the bare-graph locus and the all-speeds union are not
closed), and directedness is essential at statement level (over a
nondirected two-member cover there is a two-point set, every clip of
which is a record, that is no record --- levelwise record-hood does
not imply record-hood, and compatible levelwise records can glue to a
non-record). The arc reading is also a theorem rather than a gloss:
membership in the speed locus over a product ambient is equivalent to
being the graph of a Lipschitz arc on a closed sub-timeline, with
relative closedness of the domain \emph{derived} --- a converging
sequence of domain points has its states Cauchy by the speed
condition itself, completeness of the state fibre collects the limit,
and closedness of the graph does the rest.

\begin{remark}[Mechanization status of the carrier]
\label{rem:su-carrier-mechanized}
\sloppy
Kernel-checked, in the semantic-universe strand of the toolkit's
proof-artefact repository (modules named per statement; the receipts
conventions are \autoref{sec:su-receipts}'s): the Fell topology,
its hypothesis-free compactness, Hausdorffness over locally compact
ambients, and the polarity lemmas of the hit and miss faces
(\texttt{SuFunctor.Universe.FellTopology}, including
\texttt{FellCloseds.compactSpace} with no hypotheses on the ambient,
and \texttt{fell\_eq\_vietoris} on compact Hausdorff ambients);
clip-to-open continuity with its mutant fixtures, and proper-map
pushforward continuity (\texttt{SuFunctor.Universe.FellRestrict}:
\texttt{continuous\_restrictOpen},
\texttt{continuous\_image\_of\_isProperMap}); the seam witness
(\texttt{SuFunctor.Universe.FellSeamWitness}, refuting closed-band
clipping through the zero-duration post-jump stage);
Theorem~\ref{thm:su-fell-descent} as
\texttt{FellCloseds.projLimHomeomorph} with the glue--clip roundtrips
and continuity in both directions
(\texttt{SuFunctor.Universe.FellDescent}), together with the
sharpness fixtures (a nondirected cover on which the canonical map
fails surjectivity; the Vietoris comparison on which the miss face's
compact support fails); the record locus
(\texttt{SuFunctor.Universe.FellRecordLocus}: the generic local-locus
package \texttt{IsLocalLocus.isClosed}/\texttt{isCompact}/\texttt{glue\_mem}/\texttt{homeomorph}, the speed locus
\texttt{speedLocus} with \texttt{isClosed\_speedLocus},
\texttt{isSpeedGraph\_iff\_forall\_clip},
\texttt{speedLocusHomeomorph}, membership one-liners including
\texttt{bot\_mem\_speedLocus} and
\texttt{isLocallyClosed\_speedLocus}, and the arc reading
\texttt{isSpeedGraph\_iff\_exists\_lipschitzOnWith}) with its
refuting fixtures (\texttt{SuFunctor.Universe.FellRecordWitness}:
the steepening family, the vertical fibre, the nondirected
specimens); and the one-point compactification, packaged as a
hyperspace correspondence: adjoining the point at infinity is a
homeomorphism
$\mathrm{CL}(Y) \cong \{B \in \mathcal{K}(Y^{+}) : \infty \in B\}$
(\texttt{SuFunctor.Universe.FellOnePoint}),
holding for arbitrary ambients at the Fell stratum --- reading the
upstairs topology as Vietoris consumes local compactness and
Hausdorffness of the ambient
(\texttt{topology\_onePoint\_eq\_vietoris}), which every (T5)
ambient supplies --- so the correspondence carries the Fell form of
a noncompact ambient into the
nonempty-compacts Vietoris presentation the class stratum of
\autoref{sec:su-classes} is stated over. The scope boundary: the mechanized record locus is the
pure-flow (single-timeline) stratum; the hybrid
(multivalued-instant) stratum rides the chart route --- the shadow
of the compact parameter class, with alternation-bounded classes
closed by hit-face arithmetic --- and its mechanization is an
obligation recorded in \autoref{sec:su-receipts}.
\end{remark}

\subsection{Records Determine Laws}

\paragraph{The capacity presentation.}
The denotation stack of \autoref{sec:su-denotation} puts measures
--- and beliefs over measures --- on the record carrier,
so the carrier owes the stack a determination principle: which
observations pin a law? The answer is certificate-shaped. Call a
family of \emph{hit-and-miss generators} over an indexed open basis
$B \colon \iota \to \mathrm{Open}(Y)$ with compact closures the
classes
\[
  \nabla_B(\mathit{hits}, \mathit{misses}) \;=\;
  \Bigl\{F \;\Bigm|\;
    \bigl(\forall i \in \mathit{hits}:\ F \cap B_i \neq \varnothing\bigr)
    \wedge
    \bigl(\forall j \in \mathit{misses}:\ F \cap \overline{B_j} =
    \varnothing\bigr)\Bigr\},
\]
for finite index sets $\mathit{hits}, \mathit{misses}$. These
generators form a topological \emph{basis} of the Fell hyperspace;
with countably many basis indices the hyperspace is second countable,
and the Fell hyperspace of
a (T5) ambient is therefore \emph{compact metrizable}. The
determination principle:

\begin{theorem}[Capacity presentation]
\label{thm:su-capacity}
Over a locally compact Hausdorff ambient with a countable open basis
of compact closures, two finite Borel measures on the Fell hyperspace
that agree on every hit-and-miss generator are equal. (The generator
family is a $\pi$-system --- intersecting generators unions their
index sets --- containing the whole space as the empty generator, and
it generates the Borel $\sigma$-algebra because the generators form a
countable basis; the empty generator is what subsumes total-mass
agreement, so the statement carries no separate mass hypothesis.)
\end{theorem}

This is the determination half of the Choquet--Matheron theorem for
random closed sets \cite{matheron75}, on the Fell carrier, in a form
a checker row can consume: a \emph{distinctness} certificate is one
generator --- two finite sets of basis indices --- plus two
evaluations, and identity is the countable schema over the enumerated
generators. With the observation basis below, the indices are
rational data, so laws on the record locus are determined by their
evaluations on a countable, \emph{rationally indexed} family of
hit-and-miss observations: records
determine laws, with the observations --- though not the evaluated
measures --- exact data of the kind the certificate stratum
consumes. Two boundary facts are kernel-checked beside the theorem
and are content, not caveats: the compactness of the miss
parameters is what makes the generators open at all (over
any noncompact ambient, the miss class of the whole space is not
Fell-open --- the escape specimen again); and the hit face
\emph{alone} does not determine, at the frame level --- there are two
Scott-continuous functionals on the opens of the hyperspace agreeing
on the entire hit-generated (lower-Vietoris) topology and differing
on a compact miss class. Escape mass is
exactly what the hit face cannot see; the miss face
exists, at determination level, to count escape. (Classically,
hit-only determination survives via $\sigma$-algebra complements and
a pinned total mass; frames have no complements, and the fixture is
that difference made a theorem.)

\begin{remark}[Mechanization status of the capacity presentation]
\label{rem:su-capacity-mechanized}
\sloppy
Kernel-checked, module \texttt{SuFunctor.Universe.FellCapacity}:
\texttt{nablaGen}, \texttt{isTopologicalBasis\_nablaGen}, the
second-countability instance,
Theorem~\ref{thm:su-capacity} as \texttt{measure\_ext\_of\_nablaGen},
the separation corollary \texttt{exists\_nablaGen\_measure\_ne}, the
Scott-continuous frame form
\texttt{scottContinuous\_ext\_of\_nablaGen}, and the two boundary
fixtures \texttt{not\_isOpen\_disjoint\_univ} and
\texttt{exists\_scottContinuous\_eq\_on\_lower\_ne}. The full localic
Choquet--Matheron statement --- the frame of the hyperspace locale
presented by these generators, with continuous valuations determined
by generator values --- is kernel-checked as well, in the point-free
library: the presented frame is built as the free frame on the
hit-and-miss data semilattice, quotiented by the least nucleus
forcing the transposed hit-and-miss laws
(\texttt{Pointfree.Frame.Presentation},
\texttt{Pointfree.Frame.HyperspacePresentation}), and on it two
Scott-continuous valuations with finite generator values agreeing on
every $\nabla$-generator are equal, distinct ones separated by a
single generator (\texttt{Valuation.ext\_of\_nabla},
\texttt{Valuation.exists\_nabla\_ne}). The engine
(\texttt{Pointfree.Valuation.ext\_of\_generators}) is frame-general
--- any frame carrying a meet-closed, join-generating generator
family --- so it applies to the presented frame at any relation set
and to the spatial frame of opens alike
(\texttt{Valuation.ext\_of\_isTopologicalBasis}): no points, no
countability, no $\sigma$-algebra, and the finiteness hypothesis
sharpened from total mass to generator values --- modularity
manufactures from generator values the finite-join agreement the
Scott-continuous form above hypothesizes. What was deliberately not
claimed remains so: the existence half of Choquet's theorem,
capacity functional to law, is not claimed --- the stack consumes
determination only.
\end{remark}

\paragraph{The observation basis.} The basis the generators range
over is fixed once. For the toolkit's actual ambients --- stage
intervals times $\omega$-semialgebraic state fibres, the live part
$\mathring{S}(W) \times V(W)$ presented inside a real coordinate
space and the cemetery adjoined by the one-point operation --- it is
the family generated under finite unions and intersections by two
rational families. On the live part: rational-coefficient
semialgebraic opens traced onto $\mathring{S}(W) \times V(W)$,
refined to the subfamily with compact closures (rational opens are a
base of the coordinate space, and tracing preserves countability,
lattice closure, and the base property). At the cemetery: products
$J \times \bigl(\mathcal{V}^{+}(W) \setminus K_n\bigr)$ of a rational
stage interval $J$ with compact closure in $\mathring{S}(W)$ and the
complement of a member of a fixed rational compact exhaustion
$K_0 \subseteq K_1 \subseteq \cdots$ of the fibre --- an exhaustion
exists because a locally compact, second-countable fibre is
$\sigma$-compact, and its members can be taken carved by rational
data (traces of closed rational balls suffice). The result is
countable, closed under finite unions and intersections, and a base
of the topology of every such ambient: at live points the traces
already form a base, since the live fibre is open in its
compactification; at points over $\infty_W$ the exhaustion
complements are cofinal among neighborhoods, since every compact
subset of $V(W)$ lies in some $K_n$; and closures stay compact
($\overline{J} \times \mathcal{V}^{+}(W)$ is compact --- the
compactified fibre is compact by construction). This is also the
honest shape of what a certificate can check near escape: a basis
neighborhood of the cemetery says exactly ``eventually outside
$K_n$'' --- finite rational data, with no carving of the point
$\infty_W$ demanded. A general (T5)
ambient, which need not carry any coordinate presentation, uses an
abstract countable basis of compact closures, which local compactness
and second countability supply; the capacity theorem consumes only
that. Two consequences. First, the record-class observation domain is
insensitive to the state-space dial of \autoref{sec:su-site}:
the fibre carving class and the test-region class are independent
dials, and the $\omega$-semialgebraic upgrade moves only the former
--- test regions stay rational-semialgebraic, which is where
certificate rows need them, since a certificate row refuting
$\mathrm{Hit}(V)$ consumes $V$'s polynomial presentation. Second, the
hit-and-miss generators over this basis are precisely the
\emph{crescent} shape the class stratum of
\autoref{sec:su-classes} builds its spec classes from --- one
closed demand cut by finitely many open escapes --- so the carrier's
determination principle, the checker's test regions, and the spec
logic's atoms speak one geometry. In the vocabulary of coalgebraic
logic the generators are Moss's cover modality $\nabla$ applied to
the basis (the presentation of the Vietoris powerlocale by generators
and relations \cite{venema12}); the internal one-sentence reading ---
the carrier is the Vietoris powerlocale of the compactified ambient
object, constructed inside the sheaf topos, with the ray-ambient
carrier its shadow along the anchored tower
--- is design at prose grade, recorded because it explains why the
objectwise-then-glue construction is canonical (presentations by
generators and relations pull back along inverse-image functors),
while everything this chapter builds on is the external, objectwise,
kernel-checked form above.

\subsection{What the Carrier Hands the Later Sections}

Four
interfaces, each consumed downstream with its receipt already in
hand: (i) per window, a compact Hausdorff second-countable hyperspace
with a crescent basis --- the home of the class stratum; (ii)
Fell-continuous restrictions, so preimages of locally closed classes
are locally closed, and class-level structure transports
contravariantly along the site for free; (iii) the descent
homeomorphism, so any fibrewise construction commuting with clipping
globalizes along covers with no further argument --- the fibre sheaf
of \autoref{sec:su-sheaf} consumes exactly this; (iv) the
determination principle, so laws over records are determined by
their evaluations on countably many rationally indexed
observations --- the exact-stratum face of the denotation stack.
\section{The Universe of Beliefs}
\label{sec:su-denotation}

\subsection{Heights, Grades, and the Order}

\paragraph{From verdicts to heights.} In the behavioural reading, a
System denotes a set of trajectories and a specification denotes a
superset test: allowed or forbidden, no third value, no weighting,
and no way for the model to confess that it may fail to produce a
full trajectory at all. The present universe replaces the verdict
object with heights. A System's denotation is a \emph{belief} --- in
full, a \emph{level-of-inconsistency functional} --- on probability
distributions over the record space: a lower semicontinuous function
\[
  \sigma \colon \Delta(\mathcal{X}) \longrightarrow
  [0, \infty],
\]
whose value $\sigma(\mu)$ is the level of inconsistency between the
conceivable distribution $\mu$ and what the object asserts ---
$\sigma(\mu) = 0$: fully consistent;
$\sigma(\mu) = \infty$: excluded outright; intermediate heights price
partial inconsistency \cite{richardson20, dalrymple26imprecise}. The
functional is the belief; the
distributions are what it scores --- imprecision, nondeterminism,
and graded confidence are all shapes of the belief, never
features of a distribution. (The mechanized development carries its
functionals on the subprobability carrier $\Delta_{\le 1}(\mathcal{X})$; the
members of the universe are the \emph{radial} ones --- height determined by
a law's direction alone --- and on those the compile of
\autoref{sec:su-hyperspace} is exact: graded functionals on probability
laws on one side, closed radial \emph{sets} of subprobability laws on the
other, the functional presentation for arithmetic and the body presentation
for order. The partiality paragraph below says how a leaky process reads.)
Membership in the universe is lower
semicontinuity, \emph{and nothing else}: convexity is not an
admissibility condition, and the convex members form distinguished
subclasses of the universe rather than the universe
(Remark~\ref{rem:su-body-atlas}). Specs consume the same stock one
storey up: a Spec's denotation is a constructible class of the
bodies these beliefs compile to
(\autoref{sec:su-hyperspace}~and~\autoref{sec:su-classes}), so the
beliefs of this section are the raw material of both sides of the
eventual judgments, and of the composition calculus's fibre data.

A single coin toss already exhibits the range. Over the two-point
record space $\{\mathrm{heads}, \mathrm{tails}\}$, three members of
the universe: the belief vanishing at the fair law and infinite
everywhere else (a precise model --- the coin is fair, and nothing
else is admitted); the indicator of the closed set
$\{\mu : \mu(\mathrm{heads}) \le \tfrac13\}$ (a hard constraint,
with no preference expressed inside it); and the relative-entropy
belief $\mu \mapsto \mathrm{KL}(\mu \,\|\, \mathrm{fair})$ (every
law admitted, priced by its divergence from fairness). The first
two are the degenerate faces of Remark~\ref{rem:su-faces}; the
third is a smooth interior member; and the order below compares
all three on one scale.

\paragraph{The order.} The universe is ordered pointwise:
$\sigma \ge \varphi$ when $\sigma(\mu) \ge \varphi(\mu)$ for every
distribution $\mu$ --- the greater belief is the stronger assertion,
at least as inconsistent with every distribution. Read as
implication --- the stronger belief implies the weaker --- the
pointwise order makes the universe a thin category of beliefs, with
implications as its morphisms \cite{dalrymple26imprecise}. This pointwise
order is the order of the universe at the present storey; on the
radial stock, under the
compile of \autoref{sec:su-hyperspace}, it is exchanged,
anti-monotonically, with \emph{inclusion} of bodies, and inclusion
is the form in which everything above this storey --- classes,
fibres, and the judgments of later chapters --- consumes it. No
satisfaction judgment is defined in this chapter; what the chapter
fixes is the ordered universe those judgments will be read in.

\begin{remark}[Mechanization status of the universe stratum]
\label{rem:su-judgment-mechanized}
\sloppy
The universe stratum is kernel-checked in three storeys, and the
carrier is held pointfree first. \emph{The pointfree storey}, at
record-space generality with opens and covers first-class: the
carrier is the space of valuations on the frame of opens --- strict,
monotone, modular into $[0,\infty]$, Scott continuity carried as a
mixin exactly where a theorem consumes it
(\texttt{Pointfree.Valuation.Basic}) --- with the whole polarity
package proved there: the promotion adjoining the cemetery outcome
is a homeomorphism onto the induced topology, its mass-monotone
half-spaces the cemetery coordinate's lower half-spaces
(\texttt{promoteHomeomorph}, \texttt{induced\_promote\_eq});
evaluations lower semicontinuous at opens, mass upper
semicontinuous, liveness walls closed
(\texttt{Pointfree.Valuation.Topology}); and escaping mass is
witnessed without a single point --- the free-ultrafilter valuation
fails Scott continuity against a directed family of opens, and the
Dirac net converges to it in the lower topology
(\texttt{Pointfree.Valuation.EscapeWitness}). \emph{The classical
storey, as the bridge}: restricting a Borel measure to opens is a
valuation, Scott-continuous exactly on the $\tau$-additive stratum
--- and on second-countable carriers every Borel measure is
$\tau$-additive \cite[414O]{fremlin4}, so the measure-theoretic
dialect of this chapter is a corollary of the pointfree one, the
hypothesis class mechanized as stated
(\texttt{Pointfree.Valuation.Bridge}). The classical subprobability
carrier itself is kernel-checked at that generality:
$\Delta_{\le 1}(\mathcal{X})$ as \texttt{SubProbabilityMeasure}
under the narrow topology, its restriction map into the pointfree
carrier continuous
(\texttt{SubProbabilityMeasure.continuous\_toSubProb}), the
promotion's classical face the measure-level dagger chart
$\Delta_{\le 1}(\mathcal{X}) \cong
\Delta(\mathcal{X} \sqcup \{\dagger\})$, a
homeomorphism for the weak topologies (\texttt{daggerChartMeasure}),
with the compactness dividend for compact carriers proved by two
independent routes
(\texttt{compactSpace\_of\_daggerChartMeasure}) and metrizability
over metrizable separable Borel carriers
(\texttt{metrizableSpace\_subProbabilityMeasure},
\texttt{Pointfree.Measure.DaggerChart}). Beliefs, the compile,
and the closedness dictionary over this general carrier are the
subject of Remark~\ref{rem:su-belief-iso-mechanized}
(\texttt{SuFunctor.Universe.General}). \emph{The finite storey}:
the carrier $\Delta_{\le 1}(V)$ over a finite
outcome type as
\texttt{SuFunctor.Universe.belCarrier}/\texttt{BelState} (compact,
convex --- \texttt{convex\_belCarrier},
\texttt{isCompact\_belCarrier}), with the probability face as
\texttt{probCarrier}/\texttt{ProbState}; beliefs on the
subprobability carrier as \texttt{Penalty} --- the mechanization's
name for the type of beliefs --- with lower
semicontinuity the standing membership hypothesis; the pointwise
order itself, under the name the development gives its comparison
relation (\texttt{SuFunctor.Universe.Judgment}: \texttt{Sat},
definitionally $\varphi \le \sigma$ --- a name this chapter reads
purely as the order); the order-level faces of
the connectives (conjunction's universal property
\texttt{sat\_sup\_iff}; the functional-level disjunction minimum
with its universal property \texttt{entails\_inf\_iff} and the
hedging witness \texttt{exists\_sat\_inf\_not\_sat} ---
\autoref{sec:su-classes} adjudicates which disjunction the
logic's $\vee$ is); the closure package for the operations the
calculus performs on members (\texttt{SuFunctor.Universe.Closure}:
lower semicontinuity closed under addition, binary and countable
suprema, and affine tilts); and two generation theorems that pin the
universe to certificate-shaped data --- rational threshold rows see
the whole probability face (\texttt{SuFunctor.Universe.Generation}:
\texttt{iSup\_boxRows\_eq\_self}, with countable index, stated for
beliefs on the probability carrier), and over the
compact finite-dimensional subprobability carrier quadratic gambles
see the whole universe
(\texttt{SuFunctor.Universe.Tents}: \texttt{iSup\_tents\_eq\_self}).
The probability-face height data of the radial presentation below is
mechanized separately as \texttt{ProbPenalty}. And the judgment tier
itself is kernel-checked at the general carrier
(\texttt{SuFunctor.Universe.General}: modules \texttt{Faces},
\texttt{ClosureGeneral}, \texttt{GenerationGeneral}): the order
faces of the connectives --- conjunction's universal property, the
functional-level minimum with its universal property, and the
hedging witness, the last on \emph{every} nonempty carrier
(\texttt{General.sat\_sup\_iff}, \texttt{General.entails\_inf\_iff},
\texttt{General.exists\_sat\_inf\_not\_sat}) --- together with wall
citizenship (closed if and only if citizen,
\texttt{General.lowerSemicontinuous\_wall\_iff}) and the wall
algebra; the closure package over an arbitrary topological carrier
--- addition, binary and countable suprema, pullback, and tilts by
bounded continuous nonnegative scores, the general form of
the finite storey's affine tilts; and generation by \emph{gamble
rows} --- rational heights on finite menus of rational thresholds
against bounded-continuous-gamble expectations, the certificate
shape --- as a directed supremum with no countability,
metrizability, or compactness consumed anywhere
(\texttt{General.iSup\_gambleRows\_eq\_self}, through
\texttt{isInducing\_testProfile}: the carrier's topology \emph{is}
the initial topology of its gamble reads), and with countable index
under second countability of the carrier, a hypothesis discharged
for compact metrizable record spaces
(\texttt{General.exists\_countable\_gambleRows\_eq\_self},
\texttt{General.secondCountableTopology\_probabilityMeasure}); the
quadratic-tent species stays finite-dimensional by nature --- its
general sibling \emph{is} the gamble-row theorem. The record-space
stratum --- $\mathcal{X}$ Polish, the standing
structure below --- is the semantic generality this chapter is
written at, and the mechanization meets it at this storey; where a
theorem below is kernel-checked at the finite stratum only, its
remark says so.\footnote{One vocabulary caution: this
finite storey is the composition ladder's \emph{finite} stratum ---
finite outcome carrier.  The ladder's \emph{exact} stratum is a
different rung: finitely supported joints over an arbitrary
carrier, with arbitrary real weights.  Over a finite outcome type
the rungs collapse --- the seam embedding is a bijection, so the
finite storey, the exact stratum, and the general carrier present
one object (\texttt{toMeas\_surjective}); over infinite carriers
they are genuinely nested, and
Remark~\ref{rem:su-density-completion} settles which claims
complete across each step, and which provably do not.}
\end{remark}

\paragraph{Grades, and the two scales.} The multiplicative
presentation of the same universe replaces heights by \emph{grades}
$g = e^{-\sigma} \in [0,1]$ (hard violation graded $0$), under which
addition of beliefs is multiplication of grades and the
order reverses. Both scales are used below --- heights compose
additively under the wiring calculus, grades multiply, and the
kernel-checked composition dictionary is deliberately stated through
an abstract exponential so that no transcendental function enters the
exact stratum. The grade scale is where a familiar family of models
lives: for $\sigma = \beta^{-1}\,\mathrm{KL}(\cdot \| p)$ the
penalized worst case is the Gibbs variational formula, and
relative-entropy beliefs toward a reference --- the inconsistency
semantics of probabilistic dependency graphs of Richardson and
Halpern \cite{richardson20} --- sit in the universe as smooth
interior members. The differentiator the universe adds is its
boundary strata: credal flats (a set of laws held with no interior
preference --- nondeterminism \emph{is} this imprecision, and no
finite-temperature relative-entropy chart expresses a flat) and
walls (the possibilistic face), which arise from the smooth chart
only as its zero-temperature limit --- though the
\emph{set-referenced} relative-entropy chart, the Hausdorff--KL
embedding of Remark~\ref{rem:su-body-atlas}, expresses exactly a
flat with relative-entropy grading off it: the intermediate
stratum between the smooth interior chart and the walls.

\paragraph{Partiality is first-class.} A belief is a functional on
probability distributions; partiality enters through the \emph{compile}, not
through the domain. The body of a member (\autoref{sec:su-hyperspace}) is a
\emph{radial} subset of the subprobability carrier
$\Delta_{\le 1}(\mathcal{X})$ --- it admits each law at every scale up to its
grade, so admitted mass need not reach one --- and at that storey the
denotation is a set, no longer a functional. A leaky process --- one that may
halt, diverge, or exit scope --- denotes along the same rays as its total
counterpart, at the grade its surviving mass earns: \emph{killing raises the
level of inconsistency}, and the ray's grade axis is the only bookkeeping the
leak requires.
(The mechanized development also carries the promotion
$\Delta_{\le 1}(\mathcal{X}) \cong \Delta(\mathcal{X} \sqcup \{\dagger\})$ as
a kernel-checked chart --- an affine homeomorphism, not a compactification
--- for the computations that want normalized mass; it is an available
presentation, not a second axis of the universe.) The asymmetry the radial
down-closure buys is the discipline the universe exists to enforce: leaked
mass can only lower the probability of reaching a bad event, so safety
certificates checked against full-mass semantics remain sound; what leak
\emph{fights} is liveness --- lower bounds on good events feel every lost
unit of mass --- and a liveness claim cannot be laundered through a model
that is allowed to quietly stop. The envelope apparatus below accounts for
this exactly, and the kernel carries the counterexample that keeps the
accounting exact (the leak-aware safe-side bound and its total-leak fixture,
Remark~\ref{rem:su-envelope-theorem}).

\paragraph{Connectives, composition, and the one inequality.}
On denotations the logic's conjunction is the pointwise join,
$\sigma_1 \vee \sigma_2 = \max(\sigma_1, \sigma_2)$ --- the least
demand at least as strong as both --- and lower semicontinuity is
closed under it. Composition of \emph{Systems}, by contrast, is
addition: the wiring calculus of
\autoref{sec:su-composition} composes belief content by
pushforward-then-add --- a composite holding several beliefs at once
asserts their sum, the overall belief being the sum of its parts
\cite{dalrymple26imprecise} --- and addition is not idempotent. The whole
relationship between the two monoids is one inequality with an exact
error term:
\[
  \max(\sigma_1, \sigma_2) \;\le\; \sigma_1 + \sigma_2
  \;=\; \max(\sigma_1, \sigma_2) + \min(\sigma_1, \sigma_2),
\]
pointwise. Composition dominates conjunction --- the sound direction:
a bound certified against the join holds a fortiori against the
composite --- and the two agree at a distribution precisely when the
beliefs are \emph{crisp} there: the lesser zero, or the greater
infinite. (In the extended nonnegative reals the identity's
$\min$-term is not a cancellative error one may subtract: at a
distribution
where the two values are $\infty$ and $1$ the sum and the join agree by
absorption while the lesser value is $1$.) The \emph{crisp locus}
--- crisp at every distribution --- is exactly the agreement locus, and
that is why deterministic intuition reads the two operations as one:
on walls they agree, everywhere. (That no wall-fragment verdict can
even observe the gap off the crisp locus is a two-line fact
displayed in \autoref{sec:su-checker-anchor}, where the wall
fragment lives.) The gap itself is quantified conservatism in the
safe direction, never unsoundness --- the slice bracket
$\{\max \le s/2\} \subseteq \{\mathrm{sum} \le s\} \subseteq
\{\max \le s\}$ for two contributors. The
inequality and its sublevel bracket are kernel-checked at the
finite-rational stratum and space-free
(\texttt{QtslAux.Imprecise.MarginalCompat}:
\texttt{sum\_dominates\_sup},
\texttt{sublevel\_marginal\_of\_dominates}; the general form over
arbitrary joint models in
\texttt{QtslAux.Imprecise.MarginalCompatGeneral}).
Disjunction is deliberately absent from this paragraph: at the
functional level the pointwise \emph{minimum} is the lattice join of
demands and is sound for the single-conclusion fragment, but the
class stratum of \autoref{sec:su-classes} shows that the
specification logic's disjunction must be read one level up, as
plain union of classes, and that anything else --- the minimum
included, against class-valued conclusions --- fails disjunction
elimination. The connective's true home is there.

\subsection{The Standing Structure, and the Radial Compile}

The record spaces everything above quantifies over are fixed once, and the
next definitions carry the credal presentation of the universe's members. The
compile below (Definition~\ref{def:su-belief}) is an isomorphism: lower
semicontinuous beliefs on \emph{probability} distributions correspond
exactly to closed \emph{radial} subsets of the subprobability carrier ---
sets containing, with each law, its whole ray $\{\gamma\mu : 0 \le \gamma \le
1\}$ down to the zero measure --- with no residue on either side
(kernel-checked at the general carrier,
\texttt{SuFunctor.Universe.RadialUniverse}: \texttt{compileEquiv}). The two
presentations are two charts on one universe, and each computation uses
whichever is natural: algebra of beliefs --- sums, divergence pricing,
affine tilts --- happens on the probability-face chart, where the operations
are pointwise; order, gluing, and the classes happen on the body chart, where
the order is inclusion. The grade dictionary $g_{S} = e^{-f}$ carries one
chart to the other.

The presentation meets the order as walls: a credal body $S$ enters the order
as its indicator $\chi_S$, under which domination is containment of
finiteness sets --- the wall reading of \autoref{sec:su-checker-anchor} ---
and the denotation's own geometry is carried by the \emph{body} $S$ itself,
which is what the classes of \autoref{sec:su-classes} quantify over. A
compiled \emph{law} enters as its ray --- the body of the precise belief in
that law --- and every wall's body contains the zero measure, because every
body does: bottom-inclusion is a theorem of radiality, not a convention
(\texttt{rwall\_ne\_top}: no wall built on the radial entry can exclude the
empty process). The consequence is load-bearing for composition. Asserting
$x = y$ of a process claims only that no live behaviour violates it; the
conjunction of jointly unsatisfiable constraints --- $x = y$, $x = 0$,
$y = 1$ --- denotes the \emph{bottom body} $\{0\}$, the lawful belief that no
live behaviour exists at all, never an ill-formed exclusion of the
inconsistent world itself (kernel-checked as stated, with that three-clause
example: \texttt{contradictory\_composite\_denotes\_bot}). A construct that
vanished only at one pinned nonzero law --- an ``equality floor'' claiming
the whole universe consistent --- is not a member of the universe at all: no
belief compiles to a pinned point (\texttt{compile\_ne\_singleton}), and
feasibility is therefore never smuggled through an equality; it is asserted,
when it is asserted, by the explicit liveness face of the presented classes
(\autoref{sec:su-classes}).

Every System this book compiles denotes through the mechanism-and-gluing
stack below: a credal body for the precise and glued stochastic faces, and,
for adversary-quantified faces, the closed \emph{convex hull} of the admitted
family's union (the hull-vs-union adjudication is
Remark~\ref{rem:su-envelope-reconciliation}'s). An empty admitted family denotes the bottom body
along with every other unsatisfiable composite; no compiled System falls
outside the body calculus. The classical imprecise-probability apparatus
lives on the same stock, and everything below is stated there.

\begin{definition}[Standing structure; credal bodies]
\label{def:su-credal-denotation}
\emph{Standing structure.} Throughout the denotation stack, the
\emph{record space} $\mathcal{X}$ is a \emph{Polish space} --- a
separable, completely metrizable topological space --- carrying its
Borel $\sigma$-algebra. Polishness here is a theorem about the
toolkit's spaces, not a further dial: the handy widgets are all
objects of the completed state-space category of
\autoref{sec:su-site} --- countable families of
$\omega$-semialgebraic sectors, where local closedness, the
membership condition of that section, already implies $G_\delta$
sector by sector, hence Polish (a countable disjoint union of
Polish spaces is Polish) --- except record spaces, which are
handled through
their own mechanism involving the spacetime site and the Fell
topology, and arrive Polish by their own route. In particular:
\begin{itemize}
\item finite sets are the dimension-zero $\omega$-semialgebraic
  spaces (discrete topology): $\Delta(\mathcal{X})$
  is the probability simplex in $\mathbb{R}^{\mathcal{X}}$, whose
  rational points are the exactly-representable distributions;
\item semialgebraic sets carry the subspace topology from
  $\mathbb{R}^n$, and the locally closed
  $\omega$-semialgebraic sets of \autoref{sec:su-site} --- the
  carrier class of this book's exact certificates, closed under
  finite disjoint union, containing
  the finite sets as dimension-zero instances, with the cemetery
  compactification applied at the carrier tier as in
  \autoref{sec:su-site} --- are the standing
  state-space class;\footnote{A subspace
  of $\mathbb{R}^n$ is Polish if and only if it is $G_\delta$
  (Alexandrov). A locally closed set is the intersection of an open
  and a closed set, hence $G_\delta$ --- and, being locally compact
  and second countable, Polish twice over; per-sector local
  closedness is the certificate, and the disjoint union inherits
  both properties with every summand clopen.}
\item bag spaces with pieces in a common ambient are
  $\omega$-semialgebraic natively:
  a countable disjoint union of locally closed semialgebraic pieces
  of one ambient dimension
  is the class's own bag pattern, every summand clopen in the union
  (its cemetery is the carrier-tier operation, as always --- the
  bag's local compactness is all the operation needs); the quotient
  of a semialgebraic set by
  a linear action of a finite group is again semialgebraic (embed by
  a generating family of polynomial invariants; the image is
  semialgebraic by Tarski--Seidenberg), which covers each symmetric
  power $\mathrm{Sym}^n(X) := X^n / S_n$ individually. The full bag
  space $\coprod_{n \in \mathbb{N}} \mathrm{Sym}^n(X)$ has
  unbounded ambient dimension across its summands --- which no
  common-ambient tagged union can host --- and it is a
  \emph{first-class object} of the completed state-space category:
  a countable family of sectors, each symmetric power its own
  sorted-cone carving (one basic rational piece,
  \texttt{sortedCone}, \texttt{mem\_sortedCone\_iff\_monotone}),
  the space literally the coproduct of its sectors
  (\texttt{bagObj}, \texttt{bagIsColimit},
  \texttt{SuFunctor.StateSpace.FamCompletion}), Polish since every
  sector is (each clopen in the union). At the
  Borel stratum this is the bag monad of Dash and Staton
  \cite{dash21}, and the Bag former of
  \autoref{ch:qtsl-mathematical-models} is its System-level lift;
\item manifold-flavoured state spaces (Lie groups, Riemannian
  configuration spaces) enter through their polynomial-data
  presentations: in this book's practice they are carved by
  polynomial data, and every compact smooth manifold is
  diffeomorphic to a real algebraic set (Nash--Tognoli), so the
  class loses nothing a certificate could have consumed; all
  manifolds in this book are second countable;
\item record spaces are the exception, handled through the
  spacetime site and the Fell topology rather than by carving: per
  window the record carrier of \autoref{sec:su-carrier} is a
  compact metrizable hyperspace --- Polish outright --- with laws
  pinned by the countable rational determination of
  Theorem~\ref{thm:su-capacity}; for record spaces presented by
  gluing along finite windows, countable products and countable
  limits of Polish spaces are Polish, and the cylinder
  $\sigma$-algebra generated by the finite-window restriction maps
  coincides, in the standard presentations, with the Borel
  structure of the product realization.
\end{itemize}
$\Delta(\mathcal{X})$ denotes the set of (countably additive) Borel
probability measures on $\mathcal{X}$, and
$\Delta_{\le 1}(\mathcal{X})$ the set of \emph{subprobability}
measures (total mass at most $1$). Both carry the \emph{weak}
(narrow) topology; for Polish $\mathcal{X}$ this topology on
$\Delta_{\le 1}(\mathcal{X})$ is metrizable, and for finite
$\mathcal{X}$ it is the Euclidean topology on the subprobability
polytope, which is compact. All occurrences of ``topologically
closed'' and ``lower semicontinuous'' in this stack refer to this
topology. The total mass $\|\nu\| := \nu(\mathcal{X})$ is weakly
continuous --- and that continuity is why the weak topology, rather
than its locally compact cousin, is the topology of record: under
the \emph{vague} topology (tests vanishing at infinity, on a
locally compact carrier) the subprobability space would be compact
at the price of letting mass leak silently in limits ---
$\delta_n \to 0$ vaguely on $\mathbb{N}$ --- which is exactly what
liveness constraints cannot survive: for every $p > 0$ the set
$\{\nu : \|\nu\| \ge p\}$ is weakly closed because mass is weakly
continuous, and is not vaguely closed. And the choice is forced,
not preferred, by the promotion itself --- mechanized pointfree,
at record-space generality, with opens and half-spaces first-class
(module \texttt{Pointfree.Valuation.Topology}): along the
promotion that adjoins the cemetery outcome the narrow topology is
exactly the induced topology, its mass-monotone half-spaces the
cemetery coordinate's lower half-spaces
(\texttt{induced\_promote\_eq}, the promotion a homeomorphism,
\texttt{promoteHomeomorph}), so the stack's polarity holds by
construction --- evaluations lower semicontinuous, mass upper
semicontinuous, liveness walls closed
(\texttt{lowerSemicontinuous\_eval},
\texttt{upperSemicontinuous\_mass},
\texttt{isClosed\_liveness\_wall}) --- and the Borel reading
enters afterwards, as a bridge rather than a foundation: on
second-countable carriers every Borel measure is
$\tau$-additive \cite[414O]{fremlin4}, which is what lets the
measure-theoretic dialect of this section stand as a corollary of
the topological one. Partiality therefore enters
this universe explicitly, as the cemetery outcome of the promotion
above, never implicitly as a topological limit. Every
$\nu \in \Delta_{\le 1}(\mathcal{X})$ with
$\nu \ne 0$ has a unique \emph{ray decomposition}
$\nu = \gamma\mu$ with $\gamma = \|\nu\| \in (0,1]$ and
$\mu = \nu/\|\nu\| \in \Delta(\mathcal{X})$; the zero measure is the
bottom element, written $0$ or $\bot$.

\emph{The object.} A \emph{credal body} on $\mathcal{X}$ is a
subset $S \subseteq \Delta_{\le 1}(\mathcal{X})$ such that
\begin{itemize}
\item[(C1)] $0 \in S$;
\item[(C2)] $S$ is topologically closed;
\item[(C3)] $S$ is \emph{radial} (downward-scaling-closed):
  $\nu \in S$ and $\lambda \in [0,1]$ imply $\lambda\nu \in S$.
\end{itemize}
(Given (C3), axiom (C1) says exactly that $S$ is nonempty.)
Convexity is deliberately not a condition: closedness is the whole
of the regularity demanded, and it is exactly the body-side face of
lower semicontinuity, through the correspondence below. The convex
credal bodies are a distinguished subclass, catalogued with their
relatives in Remark~\ref{rem:su-body-atlas}.

\emph{Grade and radial structure.} For $\mu \in \Delta(\mathcal{X})$
the \emph{grade} of $\mu$ in $S$ is
$g_S(\mu) := \max\{\gamma \in [0,1] : \gamma\mu \in S\}$,
the largest admitted scaling of $\mu$; the maximum is attained,
because $\{\gamma : \gamma\mu \in S\}$ is a closed subset of $[0,1]$
containing $0$. By (C3) the intersection of $S$ with the
ray through $\mu$ is exactly the closed segment
$[0, g_S(\mu)] \cdot \mu$: a credal body is determined by its grade
function.

\emph{The radially maximal slice.} The \emph{radially maximal slice}
of $S$ is the set of its measures of maximal admitted scale:
\[
  S^{\max} \;:=\;
  \{\nu \in S : \lambda\nu \notin S \text{ for every } \lambda > 1\}
  \;=\;
  \{g_S(\mu)\,\mu : \mu \in \Delta(\mathcal{X}),\ g_S(\mu) > 0\}.
\]
The zero measure is never radially maximal, so
$S^{\max} = \varnothing$ if and only if $S = \{0\}$, the \emph{bottom
credal body}.
\end{definition}

\begin{remark}[Why the lower envelope reads the maximal slice]
\label{rem:su-maximal-slice}
Axioms (C1) and (C3) are not optional conveniences: they are what
makes the belief correspondence below exact, and the bottom credal
body $\{0\}$ the lawful denotation of
infeasibility --- the empty set is not a credal body at all, and
nothing in this universe ever denotes it. But the axioms force
$0 \in S$, so for any nonnegative
gamble $h$ the naive infimum over all of $S$ is degenerate:
$\inf_{\nu \in S} \int h \, d\nu = 0$ identically, attained at
$\bot$. The operative lower envelope therefore reads the radially
maximal slice, and the upper envelope is unaffected on the gambles
that matter: for $h \ge 0$ and $S \ne \{0\}$,
$\sup_{\nu \in S} \int h \, d\nu = \sup_{\nu \in S^{\max}} \int h \,
d\nu$, since every nonzero admitted measure is dominated along its
ray by its slice representative. (For gambles taking negative values
the supremum must be taken over all of $S$.) Degenerate case, with
the two sets kept straight: at
bottom the \emph{body} is the nonempty set $\{0\}$ while the
\emph{slice} is empty, so with $\inf \varnothing = +\infty$ the
lower prevision of every gamble degenerates to $+\infty$ while the
upper prevision is $0$. \emph{Ex falso quodlibet} is the wrong lesson here:
the bottom body is not a contradiction from which every verdict
follows but a lawful denotation with a definite, one-sided verdict
profile. The upper envelope reads the whole body at bottom as
everywhere else, and at $\{0\}$ it reports the zero measure's own
values: every event has upper probability $0$, so every
safety-shaped verdict --- an event upper bound, a down-closed
constraint --- is satisfied at bottom, at the zero measure
rather than vacuously; equivalently, since the axioms force
$0 \in S$ for every credal body, $\{0\}$ is the least body and lies
inside every safety face's body. The lower envelope's $+\infty$, by
contrast, is the arithmetic of an empty index set, not a verdict: a
liveness-shaped judgment carries an existential commitment --- the
class stratum below writes it as an escape clause,
$K \cap S(v)^{\complement} \neq \varnothing$ --- that the bottom
body can never witness, because the zero measure never leaves any
demand's body. Reading the vacuous $+\infty$ as satisfaction of
every positive lower threshold would decide those judgments in the
wrong direction. The mechanized verdict split of
Remark~\ref{rem:su-envelope-reconciliation} enforces the
discipline: the slice-reading recovery theorems carry the
$\mu \neq 0$ fence (\texttt{satCredal\_lower\_rayC\_geC}), the
bottom body is the unique credal set with empty slice
(\texttt{maxSlice\_rayC\_empty\_iff}), and the ray's zero element is
the universal refuter of every positive threshold read over the
whole set (\texttt{not\_satCredal\_upper\_rayC\_geC}). The verdict
profile of infeasibility is therefore \emph{all safety, no
liveness}: an infeasible denotation satisfies every safety
constraint and witnesses no liveness demand --- the sound behaviour
for the denotation of infeasibility --- and the presented class
stratum builds the profile in structurally: every presented atom's
liveness family carries the properness demand, so the infeasible
system belongs to no presented atom's class
(\texttt{top\_notMem\_modelClass};
Remark~\ref{rem:su-presupposition-fork}), the refusal located in
the liveness conjunct alone while the safety containment holds
throughout.
\end{remark}

\begin{definition}[Beliefs on the probability face; the
radial compile]\label{def:su-belief}\label{def:su-penalty-functional}Adopt the standing structure of
Definition~\ref{def:su-credal-denotation}, together with the
arithmetic conventions $e^{-\infty} := 0$ and $-\log 0 := +\infty$.
The \emph{radial compile} of a belief
$f \colon \Delta(\mathcal{X}) \to [0,\infty]$ is
\[
  S(f) \;:=\; \bigl\{\, \gamma\mu \;\bigm|\;
    \mu \in \Delta(\mathcal{X}),\
    \gamma \in [0,\, e^{-f(\mu)}] \,\bigr\}
  \;\subseteq\; \Delta_{\le 1}(\mathcal{X}),
\]
and the \emph{belief of} a credal body $S$ is
$f_S(\mu) := -\log g_S(\mu)$. The two directions are mutually
inverse bijections between the lower semicontinuous beliefs
on $\Delta(\mathcal{X})$ and the credal bodies on
$\mathcal{X}$ --- lower semicontinuity on the belief side
corresponds exactly to closedness on the body side, with no further
condition on either --- under which the grade dictionary
reads $g_{S(f)} = e^{-f}$, and inclusion of credal bodies
corresponds anti-monotonically to the pointwise order on beliefs:
$S(f) \subseteq S(h)$ iff $f \ge h$. The everywhere-infinite belief
compiles to the bottom credal body $S(\infty) = \{0\}$ --- never to
the empty set --- so infeasibility (the sum of two contradictory
beliefs, say) denotes a lawful credal body and the
correspondence has no excluded cases.
\end{definition}

\begin{remark}[Mechanization status of the compile correspondence]
\label{rem:su-belief-iso-mechanized}
Kernel-checked at the \emph{general} carrier --- subprobability
Borel measures over an arbitrary topological measurable record
space, under the narrow topology, $[\mathrm{Nonempty}\
\mathcal{X}]$ the one standing hypothesis (the bottom body $\{0\}$
needs a direction to point at) --- in
\texttt{SuFunctor.Universe.General}: closedness of the compile is
\emph{equivalent} to lower semicontinuity of the belief, both
directions (\texttt{isClosed\_compile\_iff}), and the image
characterization closes the bijection there --- every nonempty
closed radial set is the compile of exactly one lower
semicontinuous belief (\texttt{exists\_lsc\_compile\_eq\_iff}),
the ray gauge attained ray by ray (rays are compact for \emph{any}
carrier), its belief reconstructing the body
(\texttt{compile\_rayPenalty}), the compile injective
(\texttt{compile\_injective}), the reconstruction law
\texttt{rayPenalty\_compile} now unconditional --- an upgrade over
the finite-stratum form, which consumed a citizenship hypothesis
--- and upper semicontinuity of the gauge
(\texttt{upperSemicontinuous\_rayGauge}) the grade-side face of
closedness. The proof shape is the pointfree one, which is why no metrizability, no sequences, and no
$\liminf$ bookkeeping appear anywhere: mass is the narrowly
continuous observation, direction is continuous wherever mass is
bounded away from zero, and the grade side carries the left-limit
content algebraically (\texttt{le\_gradeN\_of\_forall\_lt}) --- the
topological leg of the correspondence is a theorem \emph{past} the
Polish form, not a Polish-stratum concession. The finite stratum
remains kernel-checked as the exact reference case
(\texttt{SuFunctor.Universe.compile},
\texttt{SuFunctor.Universe.VietorisStratum},
\texttt{SuFunctor.Universe.ImageChar}), and the convex substratum
carries a second, exact-rational and exponential-free mechanization
(\texttt{QtslAux.Imprecise.PenaltyCredalIso}, in the multiplicative
grade coordinates $g = e^{-f}$: exponential-convexity of the
belief corresponds exactly to convexity of the body --- the
harmonic-mean inequality in cross-multiplied form --- the
ray-segment description carries attainment of the grade as data,
and radiality is there a theorem, mixing with the zero measure
being downward scaling: one place where convexity genuinely pays).
The grade dictionary for the composition law --- the grade of an
additive composite is the \emph{product} of the grades, through any
abstract exponential --- is proved in the same module. The
mechanizations agree on their shared strata
through the seam isomorphism of Remark~\ref{rem:su-seam-mechanized}.
\end{remark}

\begin{example}[Conditional beliefs, independence beliefs, and what
composition does not assume]
\label{ex:su-epistemic-states}
Two epistemic states arise so often in modelling practice that their
compiles are worth stating once, on a record space that factors,
$\mathcal{X} = A \times B$, with $X$ and $Y$ the coordinate
observables of a joint law $\mu \in \Delta(\mathcal{X})$. Neither is
a probability distribution about anything: each is a condition
\emph{about} joint laws --- exactly the shape of epistemic content
that forces the belief reading in the first place
\cite{dalrymple26imprecise}.

\emph{Conditional beliefs.} The belief ``$X$ given $Y$ has
conditional law $k$'' --- $k$ a probability kernel from $B$ to $A$
--- compiles, crisply, to the set of joints whose disintegration
along $Y$ agrees with $k$ at $\mu_Y$-almost every $y$: inconsistency
$0$ there, $+\infty$ elsewhere, the body the radial closure of that
set. Richardson's weakening \cite{richardson20} grades the wall
instead of asserting it: with a confidence weight $\beta > 0$, the
belief is the scaled conditional relative entropy
\[
  f(\mu) \;=\; \beta \int_B
  \mathrm{KL}\bigl(\mu(\cdot \mid y) \,\big\|\, k(\cdot \mid y)\bigr)
  \, d\mu_Y(y)
  \;=\; \beta\, \mathrm{KL}\bigl(\mu \,\big\|\, \mu_Y \ltimes k\bigr),
\]
where $\mu_Y \ltimes k$ is the joint with $Y$-marginal $\mu_Y$ and
conditional $k$: each conceivable joint is charged the divergence of
its \emph{actual} conditional from the believed one, at a rate set
by the confidence. Since $\mu \mapsto \mu_Y \ltimes k$ is affine and
$\mathrm{KL}$ is jointly convex and lower semicontinuous, these
beliefs are convex and lsc --- they are the per-edge terms of the
probabilistic-dependency-graph inconsistency measure, the family
whose classical verification the Richardson chart of
Remark~\ref{rem:su-body-atlas} records --- and the crisp wall is
their $\beta \to \infty$ limit. (Kernel-checked at the finite
stratum, \texttt{SuFunctor.Universe.EpistemicStates}: affinity is
\texttt{condTarget\_mix} --- proved, not assumed --- joint convexity
of the divergence is \texttt{klq\_mix\_mix\_le}, lower
semicontinuity \texttt{lowerSemicontinuous\_condPenalty}, and the
zero set is the agreement of the disintegration with $k$ at every
$y$ of positive mass, \texttt{klq\_condTarget\_eq\_zero\_iff} ---
the almost-everywhere guard is load-bearing in exactly that
disintegration-ratio form, while the equivalent joint-coordinate
form absorbs it.)

\emph{Independence beliefs.} The belief ``$X$ and $Y$ are
independent'' is, in graded form, the functional proportional to the
\emph{mutual information},
\[
  f(\mu) \;=\; \beta\, I_{\mu}(X; Y)
  \;=\; \beta\, \mathrm{KL}\bigl(\mu \,\big\|\,
    \mu_X \otimes \mu_Y\bigr),
\]
zero exactly on the product laws, so the crisp face has body the
radial closure of the product laws. The belief is lsc ---
$\mathrm{KL}$ is jointly lsc and marginal-product formation is
weakly continuous; on finite record spaces $I_{\mu}$ is continuous
outright --- but \emph{nothing here is convex}: a mixture of two
product laws is not a product law (mixing $\delta_{(0,0)}$ with
$\delta_{(1,1)}$ in equal parts manufactures perfect correlation
from two deterministic independencies), so the crisp independence
body is closed, radial, and non-convex, and $I_{\mu}$ is neither
convex nor concave in the joint. Independence beliefs are the
ordinary inhabitant of $\mathbf{CB} \setminus \mathbf{CV}$ in the
atlas below: the universe's stock demands closedness and radiality
but not convexity because one of the most common beliefs in the
modelling vocabulary would otherwise not denote. (Kernel-checked,
same module: zero-iff-product is \texttt{mutualInfoQ\_eq\_zero\_iff},
continuity outright \texttt{continuous\_mutualInfoQ}, the two-delta
witness \texttt{Witness.not\_penaltyConvex\_indepPenalty}, and the
membership itself \texttt{indepBody\_credalBody\_not\_convex} ---
stated in the atlas modules' own credal-body vocabulary.)

\emph{What composition does not assume.} Conjunction of beliefs is
addition (the grade product, in the multiplicative
chart), and a belief that reads only its own subsystem's marginal
constrains only that marginal. So when two Systems sit side by side
with no wiring between them, the composite body is the set of
\emph{every} joint whose marginals lie in the two bodies --- the
whole Fr\'echet class of couplings --- and not the independent
coupling: absent an explicit causal dependency, the universe does
not assume independent noise, because sums of marginal-reading
beliefs cannot manufacture a correlation constraint. Independence
is itself a belief, asserted when it is asserted --- gradedly by
the mutual-information belief above, crisply by its wall --- and
the default composite prices every coupling the marginals admit,
which is the sound direction for certification: a verdict that
survives the Fr\'echet class survives whatever the unmodelled world
correlates. (Kernel-checked, same module: the composite body is the
Fr\'echet class, \texttt{compile\_marginalWallSum} with
\texttt{marginalSum\_eq\_zero\_iff}, and the correlated coupling of
two uniform marginals witnesses that the composite never implies
independence,
\texttt{Witness.frechet\_does\_not\_imply\_independence}.)
\end{example}

\subsection{The Atlas of Credal-Body Subclasses}

\begin{figure}[t]
\centering
\begin{tikzpicture}[x=1mm,y=1mm,
  cls/.style={draw, rounded corners=1.2pt, align=center,
    font=\small, inner xsep=2mm, inner ysep=1.2mm},
  cover/.style={semithick}]
  \node[cls, minimum width=80mm] (CB) at (0,46)
    {\textbf{CB} --- credal bodies: nonempty, closed, radial\\
     (the universe's stock: bodies $\leftrightarrow$ lsc beliefs)};
  \node[cls, minimum width=64mm] (CV) at (0,34)
    {\textbf{CV} --- convex credal bodies\\ (grade-convex compiles)};
  \node[cls, text width=38mm] (R) at (-40,19)
    {\textbf{R} --- Richardson chart: compiles of convex, lsc
     beliefs};
  \node[cls, text width=52mm] (KA) at (28,19)
    {\textbf{KA} --- Kosoy--Appel ultracontributions: closed, convex,
     pointwise-down-closed ($=$ prevision bodies)};
  \node[cls, text width=30mm] (CS) at (-50,0)
    {\textbf{CS} --- credal sets, radially embedded (the wall corner
     of \textbf{R})};
  \node[cls, text width=33mm] (HL) at (-2,0)
    {\textbf{HL} --- Halpern--Leung MWER: ultracontributions reaching
     full mass};
  \node[cls, text width=32mm] (MVS) at (44,0)
    {\textbf{MVS} --- $C^{\downarrow}$: finitely generated,
     down-closed};
  \draw[cover] (CV.north) -- (CB.south);
  \draw[cover] (R.north) -- (CV.south west);
  \draw[cover] (KA.north) -- (CV.south east);
  \draw[cover] (CS.north) -- (R.south);
  \draw[cover] (HL.north) -- (KA.south west);
  \draw[cover] (MVS.north) -- (KA.south east);
\end{tikzpicture}
\caption{The atlas of credal-body subclasses
(Remark~\ref{rem:su-body-atlas}) as a Hasse diagram, higher $=$
larger. Every edge is a strict inclusion --- placement a
general-carrier theorem, strictness a two-point rational witness ---
and every pair not joined by a path is incomparable, both directions
witnessed over the same two-point outcome space; the receipts are
the modules cited in the remark. The credal-set class enters by
radial embedding, and precise models are its singleton corner.}
\label{fig:su-body-atlas}
\end{figure}

\begin{remark}[The atlas of credal-body subclasses]
\label{rem:su-body-atlas}
The literature's imprecise-probability objects sit inside the credal
bodies as distinguished subclasses, and the inclusions between them
are part of this universe's content --- Figure~\ref{fig:su-body-atlas}
draws the poset. Read through the compile, the same atlas is the
chart of full subcategories of the thin category of beliefs drawn
in \cite{dalrymple26imprecise}, with the explicit transforms
mechanized here. The inclusion spine is
kernel-checked at the \emph{general} carrier --- subprobability
measures over an arbitrary record space, each placement below a
theorem
(\texttt{SuFunctor.Universe.General.CredalPoset},
\texttt{SuFunctor.Universe.General.Richardson}) --- while the
finite-rational stratum remains the reference case for the
\emph{separations}: every strict edge and every
incomparability a two-line rational witness over a two-point
outcome space, and the outer class containing all of them
(\texttt{SuFunctor.Universe.CredalPoset}:
\texttt{house\_includes\_all}).
Two closure notions must be kept apart first: \emph{radiality}
((C3): closed under scaling $\nu \mapsto \lambda\nu$,
$\lambda \in [0,1]$) and \emph{down-closure} in the pointwise order
on measures ($\nu' \le \nu$ setwise implies membership).
Down-closure implies radiality; the converse fails --- a single ray
$[0,1]\cdot\mu$ is radial and, for non-degenerate $\mu$, not
down-closed. With that distinction in hand:
\begin{itemize}
\item \emph{precise models}: the rays $[0,1] \cdot \mu$ --- closed,
  radial, convex; the compiles of the beliefs vanishing at a
  single law (Remark~\ref{rem:su-faces});
\item \emph{credal sets}, in the sense of Cozman's credal-network
  literature \cite{cozman00}: nonempty closed convex sets $M$ of
  \emph{full-mass} laws, entering as the radial closure
  $[0,1] \cdot M$, with $S^{\max} = M$ --- and closedness of the
  radial closure is proved \emph{through the dictionary}: the
  radial embedding of a closed convex full-mass set is the compile
  of its direction wall
  (\texttt{scaleSegment\_eq\_compile\_dirWall}), so the body-side
  regularity is Remark~\ref{rem:su-belief-iso-mechanized}'s
  closedness theorem applied, not a separate argument. A credal set
  enters the universe twice over: as its \emph{wall} --- the
  indicator belief $\chi_M$ of Remark~\ref{rem:su-faces}, the
  possibilistic reading --- and as its \emph{Hausdorff--KL
  embedding} $\mu \mapsto \min_{\nu \in M} \mathrm{KL}(\mu \,\|\,
  \nu)$, the graded reading of the same set: zero exactly on $M$,
  relative-entropy-graded off it --- the set-valued extension of
  the Gibbs chart of the grades paragraph above, with the direction
  of the divergence \emph{forced} rather than chosen: under the
  epistemology's multiplicative-weight update
  \cite{dalrymple26imprecise} (add the expected
  surprisal $E_{\mu}(-\log L)$, renormalize), the wall reading can
  only delete states, the reversed-divergence reading collapses to
  the vacuous belief, while the Hausdorff--KL reading conditions
  each member and charges it the surprisal of its own marginal
  likelihood --- likelihood updating of weighted sets of
  probabilities in the sense of Halpern and Leung \cite{halpern13},
  whose zero set is the maximum-likelihood update of Gilboa and
  Schmeidler \cite{gilboa93} (Dempster conditioning, on
  belief-function credal sets), and whose graded part retains
  exactly the marginal-likelihood record that element-wise
  conditioning \cite{cozman00} forgets (kernel-checked at the
  finite stratum, including the failure of the element-wise square:
  \texttt{SuFunctor.Universe.HausdorffKL});
\item \emph{Richardson's universe} \cite{richardson20}: bodies
  compiled from beliefs both \emph{convex} and
  \emph{lower semicontinuous} --- with the provenance stated
  carefully, because the class is this universe's reading of the
  programme, not a definition of Richardson's: the
  probabilistic-dependency-graph programme never axiomatizes a
  codomain class at all --- it proves convexity of its inconsistency
  measure as a lemma and never mentions semicontinuity --- so the
  reading takes both properties as definitional, assuming them of
  the programme's beliefs. The assumption is verified
  where checkable: the per-edge conditional relative entropies and
  their nonnegative combinations (the inconsistency measure itself)
  are jointly convex and lower semicontinuous, classically; the full
  $\gamma$-weighted scores, which add signed conditional-entropy
  terms, are convex exactly in the programme's own weight regime
  (edge confidence at least $\gamma$ times the structural weight ---
  the programme's convexity analysis states this bound, and outside
  it a two-outcome double-edge graph gives an exact rational
  counterexample). Neither half of the reading is redundant.
  Convexity of the belief is sufficient, not necessary, for
  convexity of the body
  --- $e^{f}$ convex is the exact condition, and on a two-point
  record space the belief $\log(2 - \mu\{a\})$ is strictly concave
  along the simplex while its exponential is affine, a
  kernel-checked witness with its exact rational comparison
  (\texttt{QtslAux.Imprecise.PenaltyCredalIso}). And the lsc half is
  not implied by the convex half: a bounded, two-valued convex
  belief's body can fail to be topologically closed, the failure
  a jump at the simplex boundary --- a convex function is
  automatically continuous on the relative interior, so
  boundary-lsc is exactly what convexity leaves unsecured
  (\texttt{exists\_penaltyConvex\_not\_closed}), and the same
  witness compiles to a body that no convex lsc belief compiles,
  so the two-property class is strictly smaller than the
  convexity-only reading would have been
  (\texttt{exists\_penaltyConvex\_compile\_not\_richardson}). The
  lsc half is thereby the half that makes the class a subclass of
  the credal bodies at all; the convex half is what places it under
  the convex bodies. The
  placement is a general-carrier theorem, exponential-free: additive
  convexity of the belief implies the exact grade-convexity
  criterion by weighted arithmetic--geometric means
  (\texttt{IsRichardsonBody.isConvexCredalBody});
\item the \emph{convex credal bodies}: exactly the compiles of
  beliefs with $e^{f}$ convex --- at the general carrier the
  criterion is carried division-free, as the cross-multiplied
  harmonic form (\texttt{GradeConvex}), with the compile of any
  grade-convex belief a convex body
  (\texttt{convexBody\_compile\_of\_gradeConvex}) --- and the
  subclass on which
  single-shot linear tests --- prevision-shaped functionals, in the
  sense of Goubault-Larrecq's prevision domains
  \cite{goubaultlarrecq07, goubaultlarrecq08} --- are a complete
  instrument, by Fenchel--Moreau duality;
\item \emph{homogeneous ultracontributions} (Kosoy and Appel
  \cite{kosoy20}): nonempty topologically closed \emph{down-closed}
  convex sets of subprobability laws --- inside the convex bodies,
  since down-closure implies radiality
  (\texttt{IsUltra.isConvexCredalBody}, at the general carrier).
  Every prevision-shaped body --- a set cut out by a family of
  one-sided gamble constraints, the hyperspace object of
  Goubault-Larrecq's superlinear-prevision domains --- is an
  ultracontribution at the general carrier
  (\texttt{IsPrevisionBody.isUltra}); on the finite stratum the
  inclusion is an \emph{equality}
  (\texttt{isPrevisionBody\_iff\_isUltra}): one poset node, two
  presentations, infra-Bayesian induction and the domain-theoretic
  prevision school meeting in one theorem --- with the general
  converse Hahn--Banach-shaped, a named successor rather than a
  silent assumption;
\item the $C^{\downarrow}$ \emph{monad} of Mio, Sarkis, and
  Vignudelli \cite{mio20}: nonempty \emph{finitely generated}
  down-closed convex sets --- inside the ultracontributions
  wherever the setwise order of measures is narrowly closed, a
  hypothesis the mechanization carries explicitly rather than
  silently: hulls of finite generating families are compact for
  \emph{any} carrier (the weight-simplex route), so finite
  generation implies closedness under exactly that named
  order-closedness (\texttt{isClosed\_ptwiseDownHull},
  \texttt{IsMVS.isUltra}) --- free at the finite stratum,
  dischargeable on metrizable carriers, and not free in general;
\item the semantics of Halpern and Leung's minimax weighted
  expected regret \cite{halpern13}: topologically closed
  down-closed convex sets of subprobability laws \emph{each
  containing at least one full-mass law} --- the ultracontributions
  whose maximal slice reaches the probability face, which is
  condition-for-condition Kosoy and Appel's \emph{homogeneous
  ultradistribution} --- although neither cites the other --- the
  subprobability encoding being Halpern
  and Leung's own.
\end{itemize}
Every class above is a class of credal bodies, so the present
universe includes them all; the chain $C^{\downarrow} \subseteq
\text{ultracontributions} \subseteq \text{convex bodies} \subseteq
\text{credal bodies}$ is strict at every step already over a
two-point outcome space; the credal-set and precise classes enter
through radial closure rather than down-closure (their radial
closures are not down-closed in general); and the regularity that
closes the whole atlas at the top is closedness alone --- the lsc
face of the compile, with no convexity anywhere in the outer class.
\end{remark}

\subsection{Envelopes: Previsions and the Degenerate Faces}

From either presentation the operators of record are the envelope
expectations, split across the slice exactly as
Remark~\ref{rem:su-maximal-slice} forces.

\begin{definition}[Lower and upper previsions of a credal body]
\label{def:su-previsions}
Adopt the standing structure of
Definition~\ref{def:su-credal-denotation}. A \emph{gamble} is a
Borel function $h \colon \mathcal{X} \to (-\infty, +\infty]$ bounded
below; its expectation under $\nu \in \Delta_{\le 1}(\mathcal{X})$
is well-defined in $(-\infty, +\infty]$. For a credal body $S$
with radially maximal slice $S^{\max}$, the \emph{lower prevision}
and \emph{upper prevision} of $h$ under $S$ are
\[
  \underline{E}(h) \;:=\; \inf_{\nu \in S^{\max}} E_\nu(h), \qquad
  \overline{E}(h) \;:=\; \sup_{\nu \in S} E_\nu(h),
\]
with $\inf \varnothing = +\infty$ (so the bottom credal body ---
the nonempty set $\{0\}$, whose maximal slice is empty --- has
vacuously infinite lower previsions). For events we write
$\underline{P}(A) := \underline{E}(\mathbf{1}_A)$ and
$\overline{P}(A) := \overline{E}(\mathbf{1}_A)$. The lower prevision
must read the maximal slice
(Remark~\ref{rem:su-maximal-slice}); the upper prevision is defined
over all of $S$, and for nonnegative gambles and $S \ne \{0\}$ the
suprema over the two index sets agree (the sets themselves do not). Full boundedness is required exactly where the two envelopes
are \emph{combined} --- imprecision gaps and the attitude blends
below --- since those expressions are undefined when an envelope is
infinite.
\end{definition}

A specification's probability operator
$\mathsf{P}_{\ge p}[\varphi]$ denotes the lower probability
$\underline{P}(\varphi) \ge p$ --- certification pessimism, always.
Three faces of one object, and the flat operator survives as a face
rather than a casualty:

\begin{remark}[The degenerate faces, stated exactly]
\label{rem:su-faces}
The two flat verdict operators survive inside the credal stack as
compiles of indicator-type beliefs.
\begin{enumerate}
\item \emph{Precise face.} For $\mu_0 \in \Delta(\mathcal{X})$, the
  belief vanishing at $\mu_0$ and infinite elsewhere is lower
  semicontinuous, and its compile is the single ray
  $S = [0,1] \cdot \mu_0$ with $S^{\max} = \{\mu_0\}$. The maximal
  slice is the singleton: both envelope operators then read $\mu_0$
  alone, and the denotation is an ordinary (precise) probability
  model. It is the slice, not $S$ itself, that is a singleton --- a
  credal body is never a one-point set except at bottom.
\item \emph{Possibilistic face.} Let $B \subseteq \mathcal{X}$ be a
  closed set of allowed records and
  $D_B := \{\mu \in \Delta(\mathcal{X}) : \mu(B) = 1\}$ the
  distributions supported on $B$ --- convex, weakly closed, the
  closed convex hull of the point masses on $B$. The possibilistic
  belief ($0$ on $D_B$, $\infty$ elsewhere) is convex lower
  semicontinuous, and on this face upper
  probabilities are Boolean:
  $\overline{P}(A) = 1$ if $A \cap B \neq \varnothing$ and $0$
  otherwise, so $\overline{P}(\mathrm{bad}) = 0$ \emph{is} the
  behavioural verdict ``bad is forbidden.'' The zero-set is all
  of $D_B$, not merely the point masses --- with convexity no
  longer a membership condition, the point-masses-only belief is
  also a lawful member, whose body is the union of the
  rays through the point masses; when $B$ contains at least two
  points that body is nonconvex and the belief sits strictly above
  $\chi_{D_B}$ in the order (for empty or singleton $B$ the two
  coincide), and the face that reproduces classical
  possibility theory is $D_B$.
\end{enumerate}
The classical imprecise-probability world embeds by the same device
one level up: an ordinary credal set $M \subseteq \Delta(\mathcal{X})$
(nonempty, convex, closed, all elements of full mass) corresponds to
the indicator belief $\chi_M$, whose compile is the radial closure
of $M$ with $S^{\max} = M$ --- the atlas of
Remark~\ref{rem:su-body-atlas} places this and its relatives.
Strictly graded beliefs --- the soft,
dependency-graph-typed models --- are exactly the credal bodies
whose maximal slices contain submass measures.
\end{remark}

\begin{remark}[The envelope calculus: the ray and its maximal slice]
\label{rem:su-envelope-reconciliation}
A precise law $\mu$ denotes as its ray $[0,1]\cdot\mu$ --- the
credal body, its entry into the universe --- and the wiring stratum
computes with the ray's \emph{radially maximal slice}, which for
$\mu \neq 0$ is exactly the singleton $\{\mu\}$
(\texttt{maxSlice\_rayC}). That the device is faithful to the
denotation is kernel-checked, hypotheses included, in module
\texttt{SuFunctor.Adequacy.Envelope}: the two verdict readings
agree kind by kind --- a lower-kind (slice-reading) verdict against
the ray is the verdict at $\mu$
(\texttt{satCredal\_lower\_rayC}), and an upper-kind (whole-set)
verdict agrees between ray and slice for every
\emph{down-closed} constraint (\texttt{DownClosed}; event upper
bounds qualify, with no hypothesis on $\mu$). The boundary is part
of the theorem, not a footnote: at $\mu = 0$ the slice is
\emph{empty} --- the bottom credal body, per
Remark~\ref{rem:su-maximal-slice} --- and ray and slice
agree as raw constraints \emph{only} there
(\texttt{rayC\_eq\_singletonC\_iff}): divergence is governed by
$\mu \neq 0$, not by submass. Mixing the two readers is
refuted as a theorem, not discouraged as a convention: reading a
lower-kind atom with the whole-set reader refuses every positive
threshold, for every $\mu$, the ray's zero element being the
universal refuter (\texttt{not\_satCredal\_upper\_rayC\_geC}). For
the squares' consumers the hypothesis discharges for free: a
compiled stochastic frame's law has mass one, hence is nonzero. The
set-valued extension is kernel-checked in module
\texttt{SuFunctor.Adequacy.EnvelopeUnion}. For a \emph{family} of
laws the set-valued presentation is the union of the members' rays:
radial and, for a closed family, itself a credal body, though
in general \emph{not convex}. Whether the closed convex hull is
taken is a modelling fact, not a convention: an adversary with
strategy access over the family can \emph{mix}, and for such faces
the hull is the achievable set --- the correct denotation --- while
for a parameter fixed once and unknown the bare union is the correct
denotation, and taking the hull would add mixtures the family cannot
realize. The
slice of the union always sits inside the family
(\texttt{maxSlice\_unionRays\_subset}, hypothesis-free), and for a
family of equal nonzero mass the slice is exactly the family
(\texttt{maxSlice\_unionRays}), equal mass being spent on
cross-ray domination. Fixtures witness the hypothesis's necessity:
a \emph{nonempty} all-zero family collapses its union to the bottom
object with an empty slice, exactly as the single-ray boundary
predicts, and the empty family yields the empty union, which sits
below even bottom.
\end{remark}

\subsection{Gluing: The Imprecise Extension Problem}

A System's
mechanisms hand the semantics \emph{local} credal data --- interval
kernels, rate intervals, disturbance sets --- on finite windows; the
denotation must be a global object on records whose restrictions
reproduce the local data. This is the Kolmogorov extension problem,
transposed, and its imprecise answer stratifies along the record direction:

\begin{itemize}
\item \emph{Existence and exactness, kernel-presented case.} When
  local data arrive as compact credal sets of transition kernels in
  the record direction --- the shape discrete-time mechanism
  presentation produces --- the glue exists with no topological
  hypotheses beyond that compactness, is computed by iterated
  marginal extension, and reproduces the local models on the nose:
  the Marginal Extension Theorem of de Cooman and Hermans
  \cite{cooman08} for imprecise probability trees, with backwards
  recursion as its computational face. The on-the-nose verdict
  carries its scope wherever this book leans on it: discrete-time
  record posets, finite or discrete state, compact credal local
  sets, and Williams-coherent globals. Continuous-state and
  continuous-time mechanisms hand the semantics local data through
  their own doctrine presentations; their extension theory is named
  debt in \autoref{sec:su-receipts}, not a corollary of the
  kernel case.
\item \emph{Uniqueness fails --- by design.} Many coherent global
  models share the same finite-window restrictions. The canonical
  choice is a universal property, not an isomorphism: the
  \emph{natural extension} --- the least-committal coherent glue,
  pointwise-smallest lower expectation at
  Definition~\ref{def:su-natural-extension}'s stated scope, largest
  compatible set of models (convex as a theorem, not by fiat; see the
  definition). The slogan this
  makes precise carries through everything in
  \autoref{sec:su-composition}: the model constrains the
  distributions; it does not select one.
\item \emph{Beyond the record direction, the flavour of coherence is
  a real choice.} For parallel cover shapes on infinite spaces,
  Walley-coherent independent products can fail to exist (Miranda
  and Zaffalon \cite{miranda-zaffalon-2015-independent-products});
  adopting Williams coherence repairs existence (De Bock's
  independent natural extension), and this universe adopts Williams
  coherence as its default --- a genuine doctrinal selection,
  recorded here, with no consequence at finite
  state where the flavours coincide.
\end{itemize}

\begin{definition}[Coherent lower prevision; natural extension]
\label{def:su-natural-extension}
Let $L$ be the linear space of bounded Borel gambles on
$\mathcal{X}$. A functional
$\underline{P} \colon L \to \mathbb{R}$ is a \emph{coherent lower
prevision} when, for all $h, g \in L$ and real $\lambda \ge 0$:
\begin{itemize}
\item[(W1)] $\underline{P}(h) \ge \inf h$;
\item[(W2)] $\underline{P}(\lambda h) = \lambda\,\underline{P}(h)$;
\item[(W3)] $\underline{P}(h + g) \ge \underline{P}(h) +
  \underline{P}(g)$.
\end{itemize}
Its conjugate upper prevision is
$\overline{P}(h) := -\underline{P}(-h)$. An \emph{assessment} is a
pair $(K, c)$ with $K \subseteq L$ and $c \colon K \to \mathbb{R}$
assigning claimed lower bounds --- the local data a System's
mechanisms hand the semantics. Its \emph{compatible set} is
\[
  M(K, c) \;:=\;
  \bigl\{P \in \Delta(\mathcal{X}) :
    E_P(h) \ge c(h) \text{ for all } h \in K\bigr\},
\]
and $(K,c)$ \emph{avoids sure loss} when $M(K,c) \neq \varnothing$.
The \emph{natural extension} of an assessment avoiding sure loss is
$\underline{E}_{K,c}(h) := \inf_{P \in M(K,c)} E_P(h)$; equivalently
--- exactly, for finite $\mathcal{X}$, by linear-programming duality
--- it is computed by finite positive combinations of assessment
deficits. Universal property, with its scope stated in the same
breath: $\underline{E}_{K,c}$ is a coherent lower prevision
dominating the assessment, and at finite $\mathcal{X}$ it is the
pointwise-smallest such functional outright --- at infinite
$\mathcal{X}$ the countably additive envelope displayed here is
least among the \emph{countably-additive-envelope} lower previsions
dominating the assessment, the classical unrestricted minimum being
represented by finitely additive probabilities, per the agreement
scope below; dually $M(K,c)$ is the largest compatible \emph{set} of
models outright --- the comprehension takes every compatible model,
and no convexity clause appears. Its convexity is a \emph{theorem}:
each constraint $E_P(h) \ge c(h)$ is affine in $P$, so the compatible
set is an intersection of half-spaces with the simplex (and closed
whenever the assessment's gambles are continuous --- expectation
against a merely Borel gamble is not a weakly continuous functional,
so closedness is a hypothesis at the standing generality, not a
freebie).
Least-committal lower prevision and largest compatible set are the
same universal property read on the two sides of the envelope
correspondence: the natural extension adds no commitment beyond what
the assessments force.
\end{definition}

\begin{remark}[Mechanization status of the envelope package; the
partiality departures]
\label{rem:su-envelope-theorem}
Kernel-checked in this repository, at the finite-rational stratum
(module \texttt{QtslAux.Imprecise.Credal}): the envelope face of the
coherence axioms over credal sets presented as finite families ---
attainment of both envelopes, superadditivity of the lower and
subadditivity of the upper, positive homogeneity, and monotonicity.
Kernel-checked at the general carrier --- an arbitrary measurable
record space, countably additive models, bounded measurable gambles
under the Bochner expectation (module
\texttt{QtslAux.Imprecise.NaturalExtensionGeneral}): envelopes of
arbitrary nonempty families of probability measures as coherent
lower previsions (no convexity or closure of the family consumed);
the compatible set, avoidance of sure loss, and the natural
extension as the envelope over the compatible set, with the
universal property at exactly the stated scope of
Definition~\ref{def:su-natural-extension} --- least among the
countably-additive-envelope lower previsions dominating the
assessment (\texttt{natExtG\_le\_lowerEnvG}), the dominating-family
hypothesis itself certifying avoidance of sure loss
(\texttt{avoidsSureLossG\_of\_dominating}), so no separate
feasibility hypothesis enters --- together with the
largest-credal-set dual, convexity of the compatible set in mixture
form, and the definition's closedness parenthetical as a theorem:
the compatible set is narrowly closed whenever the assessment's
gambles are continuous. One storey up, where denotations are credal
bodies rather than subsets of $\Delta(\mathcal{X})$
(section~\ref{sec:su-hyperspace}), the same universal property is
available with no convex vocabulary at all: an assessment's
constraints have walls, the natural extension is the \emph{join of
those walls} in the pointwise order, and its body is the
\emph{largest compatible body} --- the radial compile of the largest
compatible set of laws. For assessment-shaped (gamble-bound)
constraints this body is convex by the affinity theorem above, read
through the compile; for structured constraints the same join is the
right object and is genuinely non-convex --- an independence
assertion's body is a credal body and not a convex one (the
kernel-checked witness of the epistemic-states module: the diagonal
mixture of two deterministic independencies is perfectly
correlated), which is why citizenship in this universe is lower
semicontinuity and never a shape axiom. The distinction is invisible
at the lower-prevision stratum by construction --- expectation is
affine in the model, so a body and its closed convex hull have the
same lower envelope on every gamble --- and taking the hull is
exactly \emph{demonic-choice} semantics: the correct reading where
an adversary may mix admitted laws, and a canonical representative
rather than a definitional ingredient everywhere else. The classical
``largest credal set'' formulation is thus recovered, not
presupposed: it is the demonic representative of the body-level
universal property. The point-free face of envelope coherence
--- the superlinear class closed under lower envelopes, over
gambles in sublevel presentation with the Choquet pairing --- is
kernel-checked in the point-free tower
(\texttt{Pointfree.Prevision}). Kernel-checked at the finite
stratum, where the claims themselves are finite (finite record
spaces, real scalars; module
\texttt{QtslAux.Imprecise.NaturalExtension}): the
sup-of-deficit-combinations display as exact
linear-programming duality, alongside the sure-loss alternative;
and the finite
envelope theorem --- every coherent lower prevision on a finite
record space is the lower envelope of its credal set, in both the
attained form and the envelope display --- the separating engine,
one invocation powering the package, being sublinear Hahn--Banach
extension. At infinite $\mathcal{X}$ the classical envelope theorem
represents coherent lower previsions by \emph{finitely additive}
probabilities; the countably additive envelope is this book's
definition, and the identification of the two representations is
used only inside the finite-state agreement domain --- beyond it,
named debt. \emph{Partiality departures, computed}: over the
maximal slice of a submass-admitting credal body, (W2) and
(W3) persist verbatim, but (W1) weakens to
$\underline{E}(h) \ge \min(\inf h, 0)$ and constant-additivity fails
by exactly the leaked mass --- deliberate, and why the axiomatic
notion is stated on the full-mass stratum, where the classical
theory is recovered exactly. The leak-aware safe-side accounting is
kernel-checked in this repository (module \texttt{QtslAux.LeakAware}): the mass-split union bound,
the per-element and envelope forms of the safe-side subtraction (for
a partial model
$P(\mathrm{safe}) = P(1) - P(\mathrm{bad})$, not
$1 - P(\mathrm{bad})$), the full-mass recovery of the familiar
$1 - V(x_0)$ floor, and a total-leak counterexample certifying that
the naive floor is \emph{unsound} for submass models.
\end{remark}

\subsection{Agreement, and the Attitude Operator}

\paragraph{Agreement.} The scope first: the
agreement theorems of this paragraph are proven for finite state
spaces in discrete time; their continuous-state and continuous-time
lifts are named debt, and certificate soundness (which needs only a
per-element argument, \autoref{sec:su-checker-anchor}) does not
wait on them. Within that
scope, the global object admits three constructions that are not
obviously equal: the game-theoretic upper expectation of Shafer and
Vovk \cite{shafervovk19} (an infimum over hedging supermartingales),
the axiomatic
most-conservative extension of T'Joens, De Bock, and de Cooman
\cite{tjoens20}, and the measure-theoretic envelope over all
compatible precise processes. The agreement theorems state that the
three coincide on monotone limits of finitary gambles and on
bounded-below Borel path functionals --- and the specification
logic's query class sits inside that domain by construction. This
book works with the game-theoretic presentation; the reason is an
adapter fact, stated with the rest of the adapter material where it
belongs (\autoref{sec:su-checker-anchor}): certificate rows
enter the defining infimum \emph{definitionally}.

\paragraph{The attitude operator.} Certification is pessimism
($\underline{E}$ of goodness-valued functionals; equivalently the
$\overline{E}$ of the loss). Design exploration wants the other end
too, and occasionally a weighted blend --- Hurwicz's
optimism--pessimism criterion \cite{hurwicz51}, transposed: for $\alpha \in [0,1]$
and a \emph{bounded} path functional $h$,
$H_\alpha(h) = \alpha\, \overline{E}(h) + (1-\alpha)\,
\underline{E}(h)$, stated for goodness polarity so $H_1$ is pure
best-case optimism and never a guarantee quantity. The endpoints
compose along records for the natural-extension denotations
mechanism presentation produces (the marginal-extension tower
property); interior $\alpha$ does not, and the failure is an
exact-rational fact --- a two-stage interval-kernel chain realizes
root-blend $\tfrac{15}{32}$ against stagewise-blend
$\tfrac{14}{32}$, and a second instance reverses the sign, so
composed interior-$\alpha$ scores are neither sound bounds nor
reliable estimates. The design consequence is a restriction the
grammar adopts wherever the operator appears: $H_\alpha$ is a
\emph{root-only} query operator with no congruence rule, evaluated
by the two passes that do compose and mixed once at the root, and
any verdict at $\alpha > 0$ renders as a design score, never as a
guarantee. The structural derivation sits beside the
kernel-checked counterexample: an interior blend mixes a
refinement-monotone and a refinement-antitone quantity, so it is
not monotone in the universe's refinement order and therefore not
an entailment-shaped operator at all --- the root-only restriction
is the order speaking.
Attitude lives in the query. Systems carry none.
\section{The Hyperspace Storey}
\label{sec:su-hyperspace}

The universe of \autoref{sec:su-denotation} is a space of
functions with a pointwise order. Everything the specification
logic will ask of the universe --- connectives, entailment under a
context, case analysis, and the composition calculus's fibres ---
lives one storey up, on \emph{sets} of denotations, and this
section fixes that storey: the hyperspace, the passage from heights
to bodies, and the two facts that organize all later use --- from
this storey up the order of the universe is inclusion, and Systems
enter as singleton elements.

\paragraph{From heights to bodies.} A belief $\sigma$ on the
probability face compiles to its \emph{body}: the closed radial
subset of the subprobability carrier admitting each law $\mu$ at
every scale up to its grade,
$S(\sigma) = \{\gamma\mu : \gamma \le e^{-\sigma(\mu)}\}$ ---
Definition~\ref{def:su-belief}'s compile, kernel-checked
at the general carrier, with the finite stratum carrying both exact
presentations
(Remark~\ref{rem:su-belief-iso-mechanized}). Under the compile the
pointwise order on the radial stock becomes inclusion of bodies,
belief addition becomes the fibrewise grade product, and the height
calculus of \autoref{sec:su-denotation} is available, on that
stock, as a calculus of closed radial sets: one structure with two
interchangeable presentations, heights for arithmetic, bodies for
topology. The
\emph{hyperspace} of the storey is the space of closed subsets of
the subprobability carrier under the Fell topology of
\autoref{sec:su-carrier} --- compact with no hypotheses on the
carrier, and the home of every body, since bodies are closed by
definition. And the storeys compose: the record storey already did
its work in the Fell topology, so a window's record space arrives
\emph{compact} metrizable (\autoref{sec:su-carrier}), and the
subprobability carrier over it is compact Hausdorff (on a compact
carrier the weak and vague readings coincide, and the carrier is
weak-star closed in the unit ball --- Alaoglu, standard); on the
finite stratum it is the simplex, compact on sight. On a compact
Hausdorff carrier Fell and Vietoris are \emph{one} hyperspace
topology --- the Fell form agrees with the nonempty-compacts
Vietoris presentation, kernel-checked at exactly those hypotheses
(\texttt{fell\_eq\_vietoris}) --- and the nonempty bodies are
compact, so that presentation is the exact stratum's home for the
class machinery below. The machinery itself no longer needs the
compactness mediation: crescents, their two faces, local
closedness, the constructible algebra, and the judgment transport
are kernel-checked at the general carrier directly, over the Fell
hyperspace of an arbitrary topological measurable record space
(\texttt{SuFunctor.Universe.General.CrescentStratum} --- both
crescent faces ride the hit-open subbasis Fell and Vietoris share:
the containment face closed because its complement is a hit of an
open, the escape faces hit-opens outright), with the
compact/Vietoris presentation surviving as the exact-collapse
special case. For a carrier that arrives noncompact, the
one-point-compactification mediator,
$\mathrm{CL}(Y) \cong \{B \in \mathcal{K}(Y^{+}) : \infty \in B\}$
(Remark~\ref{rem:su-carrier-mechanized}), still carries the Fell
form into the compact presentation where the exact stratum wants
it.

\paragraph{The order at this storey is inclusion.} On bodies the
universe's order reads $S(f) \subseteq S(h)$ iff $f \ge h$
(Definition~\ref{def:su-belief}), and from here up ---
bodies in classes, classes in classes, fibres over the wiring base
--- every comparison this universe supports \emph{is} an inclusion.
This is the form of the order the specification logic's judgments
will consume when later chapters interpret the logic here: not a
new relation at each storey, but one order --- pointwise on
beliefs, exchanged by the compile with inclusion at the
hyperspace, read off the fibre carriers of the composition construction
unchanged (\autoref{sec:su-composition}).

\paragraph{Systems are singleton elements.} A System enters the
storey as the singleton $\{S(\sigma)\}$ of its body --- an
\emph{ideal point} of the hyperspace. The name records a theorem,
not a convention --- singletons are not spec-definable (no formula
of the spec syntax denotes a one-point class:
\texttt{singletons\_not\_definable}, with its upper-open companion),
so a precise refinement can never be pinned by Spec syntax and
selection data is certificate content, never formula content. And
the embedding leaves no residue: membership of a body in a class is
inclusion of its singleton in the class --- the fibre order of
\autoref{sec:su-composition}, at one window
(\texttt{idealPoint\_le\_iff}) --- so ``Systems are singleton
elements of the hyperspace'' is the exact sense in which one order
serves every boundary: body above body, singleton inside class,
class inside class.

\section{The Fibre of Specs: Constructible Classes in the
Hyperspace}
\label{sec:su-classes}

A Spec needs more than one body at a time. The logic conjoins
Specs, disjoins them, asks whether one entails another under a
context, and wires the Systems they govern into larger Systems.
Each of those acts consumes classes of bodies rather than single
bodies, and the classes must survive all of them at once. This
section fixes the classes. The atoms of the stratum denote
\emph{crescents}: one closed containment face, cut by finitely many
open escape faces. Compound Specs denote finite unions of
crescents, which live in the \emph{finitary constructible algebra}
of the hyperspace. Entailment denotes inclusion of classes, and a
finite blocking criterion decides it. Every demand the logic places
on this structure is discharged as a theorem, and the section ends
where the theorems currently end: at the algebra's boundary, where
the attitude classes of \autoref{sec:su-denotation} provably live
outside it, and where a completion question stands open.

What a logic may ask of its classes was settled long before this
book. Categorical logic, since Lawvere, reads predicates as objects
of fibres over a base of contexts and states what each connective
demands \cite{lawvere69,lawvere70}. Substitution must act
functorially along the maps of the base. Entailment must be the
order of each fibre, and posets suffice for a proof-irrelevant
sequent logic. Conjunction and disjunction must denote fibrewise
lattice structure, preserved by reindexing, with the distributive
law available to the context-carrying elimination rules.
Quantifiers, where a logic has them, must be adjoints to
substitution, subject to Frobenius reciprocity and the
Beck--Chevalley condition. Jacobs gives the systematic dictionary
between fragments of logic and the fibred structure each fragment
demands \cite{jacobs99}; Pitts surveys the same territory
\cite{pitts00}; Melli\`es and Zeilberger reconstruct Lawvere's
setting bifibrationally \cite{mellies16}. The Spec connectives
speak the \emph{coherent} fragment: finite conjunction, finite
disjunction, top and bottom; no negation, no implication, no
quantifier former. Its categorical model theory is classical
\cite{makkaireyes77}, and the match is designed rather than
residual. Coherent sequents are the sequents the blocking criterion
below decides with finite certificate data --- domination rows,
selections, choice bytes --- and the same fragment, in its
implicational dress, is the one the Geolog engine's chase evaluates
(\autoref{ch:geolog-reference}). Certifiability chose the fragment;
the tradition then says what the fibre owes.

The debts are paid in order. Atoms and their classes come first,
together with the polarity discipline that replaces any
monotonicity intuition a reader might import. The disjunction
theorem then removes the one degree of freedom an interpreter might
believe remains. The blocking criterion decides entailment. Three
refusals close the account: negation, refused by the carrier;
quantifiers, declined at this stratum while the genuine adjoint
lives one storey down; and attitudes, refused by a theorem. Beneath
the atoms, the graded layer explains why graded judgments demand
no graded classes, and a final subsection states the completion
question the refusals leave open. On evidence, one standing note:
the coherence claims of this section cite receipts from multiple
proof-checking kernels (Lean~4, SafeVerify, and \texttt{nanoda};
\autoref{sec:su-receipts} states what each checks), the receipts
are re-runnable, and the claims that do not yet have them are
labelled rather than blended in. This section is where that
discipline proved its worth --- the design-record instance after
Theorem~\ref{thm:su-vee-forcing} tells the story.

\subsection{Atoms: Crescents}

\begin{definition}[Spec atom, and its crescent]
\label{def:su-crescent}
A \emph{spec atom} is a pair $(u, \mathcal{D})$: a \emph{safety
face} $u$, one member of the universe, the atom's floor; and a
finite \emph{liveness set} $\mathcal{D}$ of members, its escape
demands. The atom's class in the hyperspace is its \emph{crescent},
\[
  \mathrm{cres}(u, \mathcal{D}) \;=\;
  \bigl\{K \;\bigm|\; K \subseteq S(u)\bigr\} \;\cap\;
  \bigcap_{v \in \mathcal{D}}
  \bigl\{K \;\bigm|\; K \cap S(v)^{\complement} \neq
  \varnothing\bigr\},
\]
the bodies contained in the safety face's body that escape every
liveness demand's body. The kernel form is
\texttt{SuFunctor.Universe.crescentClass}.
\end{definition}

The two kinds of face quantify in opposite directions, and the
opposition is the atom's content. The containment face demands that
\emph{every} admitted law stay within the floor's body; it is the
safety half of the atom, and the pessimistic pass of
\autoref{sec:su-denotation} discharges it, by domination against
the worst case. Each escape face demands that \emph{some} admitted
law leave the demand's body; it is a liveness conjunct, and the
optimistic pass discharges it, by a witness at the best case. One
atom thus couples a $\forall$-polarized face with finitely many
$\exists$-polarized ones. This is the liveness repair of
\autoref{sec:su-denotation} at the right altitude, and the leak
asymmetry established there governs how partiality counts on each
face.

Topologically, a crescent is one Vietoris-closed containment face
cut by finitely many lower-open escape faces, hence a
\emph{locally closed} subset of the hyperspace
(\texttt{isLocallyClosed\_crescentClass}, in the locally-closed
layer at the general carrier). The shape is kindred to the
hit-and-miss generators of \autoref{sec:su-carrier} --- one closed
demand, finitely many open ones, the cover-modality geometry ---
though the two families are not one: the capacity generators hit
basis opens and miss compact closures, while a crescent's closed
face is a containment. The kinship is what matters. The spec
logic's atoms, the carrier's determination principle, and the
checker's test regions speak one geometry, at two floors of the
same tower.

The transports between an atom's presentation and its class are
exact: a body lies in the crescent with empty liveness set
precisely when its belief dominates $u$ in the universe's order,
and inclusion between safety-only crescents reverses the pointwise
order on their floors. The kernel forms are
\texttt{compileNC\_mem\_crescentClass\_iff} and
\texttt{compile\_subset\_compile\_iff}.

\paragraph{Graded reach-avoid: the atoms, in control-theoretic
dress.} The atoms are a familiar object of control theory wearing
graded clothing. Call a spec atom with one \emph{genuine} liveness
demand --- liveness set $\{\ell, \top\}$, the properness mark
carried along structurally --- a \emph{graded reach-avoid} atom.
Its judgment demands that the system dominate the safety floor $u$
(graded avoidance: the floor prices every admitted law rather than
merely forbidding a region) and dip under the mark $\ell$ (graded
reach: some admitted law strictly undercuts the demanded
bar); the kernel form is
\texttt{mem\_modelClass\_reachAvoid\_iff}. With wall data on both
faces --- a crisp safe set $S$ and a crisp target $T$ --- the
judgment becomes literally the classical possibilistic reach-avoid
demand, \emph{exclude $S^{\complement}$, keep some state of $T$
possible} (\texttt{coupledSat\_crispRA\_iff}), the specification
studied as a decision problem since \cite{summerslygeros10} and
still the benchmark object of certificate-based stochastic
verification \cite{xue25,cao25}. The generalization is strict
already over a one-point carrier, and at the algebra level rather
than atom by atom: no finite $\vee/\wedge$ combination of crisp
reach-avoid classes separates the constant systems $1$ and $0$ ---
every wall admits both or excludes both --- while a graded floor
does (\texttt{exists\_genClosure\_reachAvoid\_not\_crispGenerated}).
The stratum is then \emph{exactly} the completion of graded
reach-avoid under the calculus's two connectives: a class is
constructible if and only if it is a finite disjunction of finite
conjunctions of graded reach-avoid atoms, kernel-checked at both
strata (\texttt{genClosure\_reachAvoid\_iff\_exists\_specDnf} on an
arbitrary carrier;
\texttt{genClosure\_reachAvoidCres\_iff\_exists\_specDnf} in the
hyperspace, with constructibility of the generated algebra as
\texttt{isConstructibleUnion\_of\_genClosure\_reachAvoidCrescentsLsc}).
The decomposition inside a single atom is conjunctive, exactly: a
liveness \emph{set} is the intersection of the singleton-liveness
atoms sharing its floor
(\texttt{modelClass\_eq\_iInter\_reachAvoid},
\texttt{crescentClass\_eq\_iInter\_reachAvoid}), and the
disjunctive reading is refuted outright --- one unsatisfiable
demand empties an atom but not the union of its singletons
(\texttt{exists\_modelClass\_ne\_iUnion\_reachAvoid}). Disjunction
of liveness demands, for its part, never needs the class level at
all: within one atom it is mark-join (\texttt{dip\_sup\_iff}), so
the $\vee$ the completion spends is spent entirely on unions of
classes. Two boundary readings earn the ``exactly.'' The empty
disjunction is generated --- the false Spec, the empty class ---
while the empty conjunction is deliberately not: no constructible
class is trivially true, because every presented atom carries the
properness demand (Remark~\ref{rem:su-presupposition-fork}), and
that same mark is what closes the completion using \emph{binary}
conjunction alone, since every liveness set is nonempty and its
intersection is anchored at the atom's own $\top$-mark. For a
reader arriving from the reach-avoid literature the translation
reads: a Spec of this logic is a finite monotone combination of
reach-avoid demands whose avoidance is priced rather than binary
and whose reach is a strict undercut of a graded bar --- and the
wall fragment (\autoref{sec:su-checker-anchor}) is the crisp
sublanguage through which checked verdicts embed.

\subsection{Classes: the Constructible Algebra}

\begin{definition}[Spec class, and the constructible sets]
\label{def:su-constructible}
The \emph{spec class} of a compound Spec in disjunctive normal
form --- a finite family of atoms $(u_i, \mathcal{D}_i)_{i < n}$
--- is the union of its disjuncts' crescents,
$\mathrm{dnf}\bigl((u_i, \mathcal{D}_i)_i\bigr) = \bigcup_{i<n}
\mathrm{cres}(u_i, \mathcal{D}_i)$. A subset of the hyperspace is
\emph{constructible}, in the finitary sense, when it is a finite
union of locally closed sets
(\texttt{SuFunctor.Universe.IsConstructibleUnion}). The classes of
the stratum are the constructible sets.
\end{definition}

Two boundaries in that definition are deliberate. The notion is
finitary rather than spectral, and the mechanization defines it
directly instead of reusing a library notion: on a connected
compact Hausdorff hyperspace the spectral-space constructible
algebra, generated by compact open sets, degenerates to the
trivial one, while the finitary algebra is the one the spec syntax
generates. And the family is closed under the operations the
calculus performs: finite unions by definition, finite
intersections because crescents intersect to crescents
(\texttt{crescentClass\_inter\_crescentClass}) and intersection
distributes over union by the powerset frame law
(\texttt{inter\_dnfClass\_append}). Top and bottom are the trivial
atom's class and the empty disjunction. Systems enter the stratum
as the ideal points of \autoref{sec:su-hyperspace}: the singleton
class $\{S(\sigma)\}$ of a System's body, spec-undefinable, with
membership read as inclusion.

What makes this algebra \emph{the} fibre, rather than one
candidate among several, is transport. The composition calculus of
\autoref{sec:su-composition} reindexes classes along wirings, and
the transport dictionary makes reindexing \emph{preimage} at the
body level (\texttt{compile\_pullPen}). Preimage respects both
kinds of face --- the pullback of a pointwise join is the join of
the pullbacks, so atom-level transport commutes with atom
conjunction definitionally --- and distributes over the
constructible unions (\texttt{transportClass\_union}), which, for
subobject fibres, is the coherent-category axiom that finite
unions be pullback-stable. So the fibres are posets under
inclusion, the one order of the universe from the hyperspace
storey up; the connectives denote fibrewise lattice structure; and
reindexing preserves it by variance rather than by axiom. The
adjudicated boundary of the transport claim --- atom-level
transport against image transport, exact along surjections,
genuinely divergent off them --- is
Remark~\ref{rem:su-transport-citizenship}, and the double
Grothendieck assembly of \autoref{sec:su-composition} packages the
whole, every obligation kernel-checked, with the window-indexed
families of these classes as its fibre data. The definite article
is earned on two sides: this is the least family that contains the
atoms and holds the connectives, and no larger family currently
carries the assembly's kernel-checked functoriality. Whether a
larger one can is the closing subsection's question.

\subsection{Polarity, not Monotonicity}

Entailment denotes inclusion, and inclusion is an order; a reader
may therefore expect the classes to sit monotonically in the
universe's \emph{other} order, refinement --- every class an
up-set or a down-set of better-and-worse systems. They do not, and
the failure is structural. A containment face is downward closed:
shrink a body inside the floor and it stays inside. An escape face
is upward closed: enlarge a body that meets an escape region and
it still meets it. A crescent with both kinds of face genuinely
present is the intersection of a down-set with finitely many
up-sets, so it is neither, and unions inherit the mixture.
Refining a System preserves its safety verdicts and threatens its
liveness verdicts; coarsening does the opposite. Satisfaction is
not monotone along refinement in either direction, and no repair
should be attempted, because the mixture is what a coupled
safety--liveness demand \emph{means}.

What composes is polarity. Every literal of a presented Spec
carries a face. Each face is discharged by one envelope pass
alone: the floor by the pessimistic pass, each escape by the
optimistic pass, and the passes compose along records separately
(the marginal-extension tower property of
\autoref{sec:su-denotation}). What a literal returns is a verdict
--- one Boolean, with the face's certificate behind it --- and
verdicts are data. The connectives combine finished verdicts and
never reach inside a face, so a compound judgment survives wiring
face by face even though no monotonicity survives at all. In
brief: \emph{a Spec mixes verdicts}. The one operator in this book
that mixes the passes themselves, before any verdict is taken,
cannot live in this stratum; that refusal now has a theorem, and
it is the subject of the attitude subsection below.

\subsection{The Disjunction Theorem}

Any interpretation of the logic into this stratum will read
formulas as constructible classes, entailment as inclusion, and
conjunction as intersection. One choice appears to remain: which
class interprets $\Phi_1 \vee \Phi_2$. The rule catalogue's
disjunction elimination, in its context-carrying core form,
demands of an interpretation $J$ of $\Phi_1 \vee \Phi_2$ (with
$A = [\![\Phi_1]\!]$, $B = [\![\Phi_2]\!]$):
\[
  \forall \gamma, \psi \colon\quad
  \gamma \cap A \subseteq \psi
  \;\longrightarrow\;
  \gamma \cap B \subseteq \psi
  \;\longrightarrow\;
  \gamma \cap J \subseteq \psi .
\]

\begin{theorem}[Union is forced]
\label{thm:su-vee-forcing}
Over inclusion-ordered class fibres, on any carrier:
\begin{enumerate}
\item contextual disjunction elimination is valid for all contexts
  $\gamma$ and conclusions $\psi$ if and only if
  $J \subseteq A \cup B$;
\item adding soundness of the two introduction rules
  ($A, B \subseteq J$), the conjunction holds if and only if
  $J = A \cup B$ exactly;
\item every strict inflation is refuted uniformly: if
  $A \cup B \subset J$ strictly, there exist an explicit context and
  conclusion --- $\gamma$ any class meeting $J$ off the union,
  $\psi = A \cup B$ --- under which both premises of the elimination
  rule hold and its conclusion fails.
\end{enumerate}
\end{theorem}

The proof is elementary (instantiate $\gamma$ at the full class and
$\psi = A \cup B$ for the forcing direction; distribute
intersection over union for soundness) --- elementary the way a
door with one key is easy to open: the fibre order leaves the
interpretation of $\vee$ \emph{no freedom at all}. The tradition
states the context-carrying elimination rule as a demand an
interpretation must satisfy; over these fibres the demand has one
solution, so the rule is not checked against a chosen semantics
--- it chooses the semantics. The kernel forms are
\texttt{SuFunctor.Universe.ctxOrE\_valid\_iff}, which depends on
no axioms whatsoever, pure constructive set logic;
\texttt{orI\_ctxOrE\_valid\_iff}; and
\texttt{exists\_ctxOrE\_countermodel\_of\_ssubset}, the latter two
within the development's standard three axioms
(\autoref{sec:su-receipts} names them).

Three consequences follow, in ascending order of consequence for
this book.

\emph{First, the availability problem is real, in-universe.}
Locally closed classes are not closed under union, and not merely
in general topology. Two spec-atom classes of the shape above ---
one closed atom, one escape atom, over a two-point outcome space
--- have a union that fails local closedness in the hyperspace,
kernel-checked with an explicit convergence witness
(\texttt{exists\_union\_crescentClass\_not\_isLocallyClosed}). A
stratum of locally closed classes \emph{alone} cannot hold its own
disjunctions; and by clause (3) of the theorem, any attempt to
stay locally closed by interpreting the join as a locally-closed
\emph{envelope} of the union manufactures its own countermodel.
The constructible closure escapes both refutations, and it is not
a closure operator at all: it enlarges the admitted classes until
the union is native, adding no points to any join. On the
constructible stratum the rule holds with equality and no side
conditions. The class of the join is literally the union
(\texttt{dnfClass\_append}), conjunction distributes
(\texttt{inter\_dnfClass\_append}), the conclusion is
constructible (\texttt{isConstructibleUnion\_dnfClass\_append}),
and contextual elimination is an instance of clause (1)
(\texttt{ctxOrE\_dnfClass}).

\emph{Second, the design-record instance.} An earlier design of
this semantics made convexity an admissibility condition of
denotations. It stated no semantics clause for spec disjunction
--- but inside a universe of convex denotations the only candidate
join of two spec denotations is their envelope, the greatest
convex lower semicontinuous minorant, whose body is a hull; and a
hull inflates the union whenever the disjuncts are not nested, so
clause (3) manufactures the refutation. Concretely, over one
probability coordinate $x$, the walls $A = [0, \tfrac13]$ and
$B = [\tfrac23, 1]$ with context $\gamma = \{\tfrac12\}$ refute
contextual elimination against the envelope reading: both premises
hold vacuously, since the context meets neither wall, and the
conclusion fails because the hull swallows the context. The
refutation, found in this development's design record against the
reading that admissibility condition forced, is thereby an
\emph{instance of a kernel-checked schema} rather than a fact
about convexity: convex hull, locally-closed envelope, and any
future closure operator fail the same rule by the same mechanism,
and plain union is the reading left standing. The repair cost
nothing that was true --- every verdict of the wall fragment
transports to the class stratum through the crescent embedding
unchanged. And the refuted candidate is itself mechanized at the
standing generality rather than surviving as prose
(\texttt{SuFunctor.Universe.HullClosure}): over complete
second-countable metric carriers, binary hulls of closed sets of
laws are narrowly closed with no closure bar
(\texttt{isClosed\_pairHull}, \texttt{isClosed\_binaryHull} ---
the mechanism is the cone's own, with no linear-space analogue: a
fixed limit law manufactures its compactness, tightness of the one
measure plus positivity of the decomposition, no cancellation
between arms), and the quasiconvex coproduct is \emph{levelwise
exact}: its sublevels are the plain hulls of the sublevel
unions, with none of the $\varepsilon$-regularization the
general-locally-convex literature would demand
(\texttt{sublevel\_qCoproduct\_eq}, greatestness as
\texttt{qCoproduct\_le\_min} and
\texttt{le\_qCoproduct\_of\_forall\_mixConvex}). The join the
elimination rule forbids is thus refuted as a fully mechanized
object, at full strength: the hull inflates the union not because
its topology misbehaves --- it provably does not, on this cone ---
but because inflation is what a hull \emph{is}.

\emph{Third, the two disjunctions separate, and the calculus's
$\vee$ is the union one.} At the functional level the pointwise
minimum is the lattice join of demands, and against
single-functional conclusions it validates the elimination rule
(the height lattice is a frame). The class stratum's grammar is
strictly richer --- conclusions may be classes --- and against
class-valued conclusions the minimum reading breaches: the body of
$\min(\sigma_x, \sigma_y)$ for two point-mass walls contains the
two-point \emph{hedger} $\{x, y\}$, a body committed to the union
without committing to either disjunct, so the containment class of
the min strictly inflates the union of the containment classes
(kernel-checked: \texttt{boxClass\_union\_ssubset}, with the
resulting elimination countermodel at an entailment-expressible
conclusion,
\texttt{exists\_ctxOrE\_countermodel\_body\_reading\_carrier}).
The functional-level minimum survives as what it is, the join of
the height lattice, sound for the single-conclusion fragment; the
specification logic's disjunction is the class union, where the
richer grammar's elimination rule demands it. What the union
reading returns is worth naming: the \emph{disjunction property}
at the system boundary. Membership of a body in a joined class is
membership in one disjunct's class
(\texttt{mem\_dnfClass\_append\_iff}), so a System inside a
disjunction is inside a disjunct, and a certificate of the
membership carries the choice of disjunct as one byte of
certificate data --- restoring, at the semantic boundary, the
constructive reading the functional level provably lacks.

\begin{remark}[The three layers of the elimination rule]
\label{rem:su-vee-layers}
The rule's status is settled at three independent layers, and the
layers do not lean on one another. \emph{Proof theory}: the
context-carrying form is not derivable from the context-free core
--- every core rule is sound over every lattice interpretation,
and the modular lattice $M_3$ refuses the distribution the
derivation would need; a kernel-checked interpretation argument,
untouched by any semantics. \emph{Chart semantics}: inside the
convex chart the rule is refuted (the design-record instance
above), which retrospectively explains the underivability, since
every other core rule is convex-chart-sound and a derivation would
have transferred a chart-validity the rule lacks. \emph{Class
semantics}: over the constructible stratum the rule is valid with
equality, and inflation-of-union is provably the only failure mode
(Theorem~\ref{thm:su-vee-forcing}). The catalogue's decision to
carry the context-carrying form as a core rule is thereby
vindicated three ways: primitive because underivable, valid
because union, and necessarily primitive because its validity is
what the convex chart cannot supply.
\end{remark}

\subsection{Inclusion, Negation, and Properness}

Inclusion between constructible classes --- the relation the
logic's entailment sequents will denote --- is decided by a finite
\emph{blocking criterion}, kernel-checked in the locally-closed
layer (\texttt{coupled\_union\_refinement\_iff}). A DNF-to-DNF
inclusion holds iff every \emph{selection} (one demand chosen per
conclusion disjunct) is \emph{blocked} by a premise-side
domination. That is a quantifier alternation over finite index
sets whose leaves are pointwise domination rows between members
--- certificate-shaped throughout: domination rows plus, per
disjunction, the choice bytes above. The criterion is sound and
complete on the stratum with no satisfiability side conditions;
inconsistent premises self-block, so \emph{ex falso} is derived
rather than assumed. (At the presented stratum the properness
reading below sharpens the boundary behaviour: every presented
atom carries the properness demand, so the everywhere-infinite
member belongs to \emph{no} presented class at all,
\texttt{top\_notMem\_modelClass}, rather than sitting vacuously in
conclusions.) Two calibration facts bound the completeness claim:
the naive rule pair that reads only per-disjunct domination
without selections is incomplete, witnessed by a floor-raising
inclusion, and refinement of the liveness sets is essential
(\texttt{exists\_naive\_cover\_incomplete},
\texttt{exists\_essential\_floor\_raise}). The packaging of the
criterion as a two-species sequent calculus --- membership
sequents and inclusion sequents with structural rules and
admissible cut --- is an obligation recorded in the ledger; the
criterion the calculus would organize is kernel-checked.

\begin{remark}[Negation, and a boundary the class algebra makes easy
to trip]
\label{rem:su-negation-fence}
The constructible family is closed under complement (the
complement of a locally closed set is the union of an open and a
closed set), so the class \emph{algebra} is Boolean, and a future
editor might read that as licence to complete the logic with a
negation-former. The boundary stands, and its provenance is the
carrier rather than the algebra: on submass carriers the mass
identity fails --- $\mu(A)$ does not determine
$\mu(A^{\complement})$ when mass may leak --- so the complement of
a wall is not the wall of a complement, event-level duals do not
lift to denotation-level duals, and the typing rejections the
grammar carries are theorems of this universe, not conservatism.
Set-theoretic availability of complements at the class stratum is
not semantic availability of a negation-former one floor down; the
orientation-paired judgment kinds (upper and lower probability
atoms) are row-level de Morgan bookkeeping, the sound residue of
the failed duality, and they lose exactness only where double
negation genuinely loses mass.
\end{remark}

\begin{remark}[The properness presupposition, adjudicated]
\label{rem:su-presupposition-fork}
One semantic decision at this stratum was long held open: whether
membership of a body in an atom's class \emph{presupposes
properness} --- so that membership additionally asserts that the
belief is finite somewhere, the everywhere-infinite denotation
belongs to no atom's class, and ex falso inverts at the system
boundary --- or whether bare safety exists, with vacuous
membership quarantined by explicit properness certificates where a
rule needs them. The presentation stratum decides it,
structurally, for the properness reading (\emph{option A}): every
presented atom's liveness set contains the properness demand ---
the mark $\top$, whose escape face a body meets exactly when it is
proper --- so an empty liveness set is unrepresentable rather than
merely forbidden, the vacuous atom is $(0, \{\top\})$ (exactly the
decoration unit of \autoref{sec:su-composition}), the inconsistent
system belongs to no presented atom's class, and ex falso inverts
at the system boundary by construction. The inversion is located,
not blanket. An atom's safety containment holds at the bottom body
for every floor (the axioms force $0 \in S(u)$, so
$\{0\} \subseteq S(u)$), and what refuses is exactly the escape
clause --- the existential conjunct every liveness demand carries,
which the zero measure, a member of every demand's body, can never
witness. An infeasible system thus satisfies every safety face and
escapes no liveness demand (the envelope stratum's verdict
profile, Remark~\ref{rem:su-maximal-slice}); its exclusion from
every presented class is the conjunction of those two one-sided
verdicts --- properness being the liveness demand every presented
atom is guaranteed to carry --- not a satisfaction collapse in
either direction. The decision cost no existing result. The
blocking criterion, the disjunction theorem, and the class algebra
were stable under both readings --- under the properness reading
each liveness slot simply carries one more demand, which is
literally what the presentation's atoms now do --- and the general
hyperspace machinery underneath (crescents at arbitrary liveness
sets) is untouched, which is why every verdict transported
unchanged.
\end{remark}

\subsection{The Declined Quantifier}

One demand of the dictionary this stratum deliberately declines:
there is no quantifier former on the classes. Binding lives
elsewhere in the universe. Records carry the variables
(\autoref{sec:su-records}); the certificate stratum's
shape-indexed families live on the functorial plates of
Remark~\ref{rem:su-strata}; and the substitution that quantifiers
would be adjoint to belongs to the base --- marginalization of
joints along wire maps (Definition~\ref{def:su-state-spaces}).

The existential the tradition asks for does exist, kernel-checked,
one storey below the classes. Hiding a wire is the
lower-semicontinuous envelope of inf-projection along the visible
marginal, and against every citizen of the universe it is adjoint
to belief pullback on the nose:
$\varphi \le \mathrm{hide}_{p}\,\sigma$ iff
$p \fatsemi \varphi \le \sigma$ (\texttt{le\_infProj\_iff},
\texttt{le\_hide\_iff}; Remark~\ref{rem:su-marginalization-lsc}).
Read in the entailment order, which opposes the pointwise order
(\autoref{sec:su-hyperspace}), this is Lawvere's
$\exists_{p} \dashv p^{*}$.

Frobenius reciprocity holds in the lax direction the formalities
supply. For a citizen $\psi$ of the visible carrier,
$\mathrm{hide}_{p}\,\sigma \vee \psi \le
\mathrm{hide}_{p}(\sigma \vee p^{*}\psi)$, and likewise with $+$
in place of $\vee$ --- either conjunction the storey speaks
(\texttt{hide\_sup\_le}, \texttt{hide\_add\_le}). The comparison
is strict: a kernel-checked witness on the escaping-mass family
exhibits citizen beliefs for which conjoining-then-hiding and
hiding-then-conjoining genuinely differ
(\texttt{exists\_lowerSemicontinuous\_hide\_add\_ne}). The
mechanism is the envelope step of the repaired existential. The
envelope of a wall is the wall of the closure
(\texttt{lscEnv\_charFn}), so hiding adds the boundary values the
visible image approaches but never attains, and a belief
vanishing exactly at an added point detects the addition. The gap
is visible at judgment resolution --- a constant specification is
dominated by one side and not the other
(\texttt{exists\_judgment\_separating\_hide\_add}) --- unlike the
envelope itself, which changes no judgment. Under the
inf-compactness certificate of
Remark~\ref{rem:su-marginalization-lsc} the certificate survives
conjunction with a pulled-back citizen, the reflector idles, and
Frobenius reciprocity holds with equality
(\texttt{hide\_sup\_comp}, \texttt{hide\_add\_comp}). Where
marginalization is exact, quantification commutes with contexts on
the nose; off the certificate, the fibred existential is genuinely
lax --- a fact the logic observes, not a packaging debt.

What is not claimed is the class-stratum packaging of that
adjoint: a fibred existential on constructible classes, with
Frobenius reciprocity over contexts and Beck--Chevalley squares
over the site. At a quantifier-free stratum the Beck--Chevalley
condition, governing as it does the commutation of quantifiers
with substitution, has nothing to govern. The mechanized seam
facts already mark where such a claim would have to be earned:
escape witnesses die off the range of a non-surjective transport,
and image transport leaves the constructible stratum (the
saturation residue of Remark~\ref{rem:su-transport-citizenship})
--- and images are what existentials are made of. The ledger
records the obligation under its own name, the image--preimage
adjunction on class fibres (\autoref{sec:su-receipts}).

\subsection{Attitude Classes: Outside, by Theorem}
\label{sec:su-attitude-classes}

The refusals so far were designed; this one is proved. Recall the
attitude operator of \autoref{sec:su-denotation}: Hurwicz's
optimism--pessimism blend \cite{hurwicz51},
$H_\alpha(h) = \alpha\, \overline{E}(h) + (1-\alpha)\,
\underline{E}(h)$, admitted into the grammar as a root-only query
with no congruence rule, because interior blends fail the tower
property and are not monotone in the refinement order (the
exact-rational counterexample and the order-theoretic derivation
both live in that section, and the specification language records
the restriction, \S\ref{sec:spec-lang-hurwicz}). That settled the
operator's grammar. It left a class-stratum question open. The
\emph{attitude class}
$\{\,K \mid \vartheta \le H_\alpha(h)(K)\,\}$ is a definite subset
of the hyperspace like any other, and nothing yet stops a reader
from asking whether some Spec --- perhaps an ingenious compound
one --- denotes it. The answer is no, and the reason is a
decomposition worth having in its own right.

Fix the exact stratum (finitely presented credal sets over
finitely supported laws --- the stratum where certificates live),
a bounded claim $h$, and a threshold $\vartheta$. Call a pair
$(a,b)$ an \emph{admissible chord} when
$\vartheta \le \alpha a + (1-\alpha) b$. Each chord names one
atom-shaped demand: an escape face at $a$ --- some admitted law
prices $h$ at $a$ or better, the optimistic pass --- and a floor
face at $b$ --- every admitted law prices $h$ at $b$ or better,
the pessimistic pass. The Hurwicz threshold holds precisely when
some admissible chord's demand holds, because the attained pair
$(\overline{E}(h), \underline{E}(h))$ is itself an admissible
chord (\texttt{hurwicz\_ge\_iff\_exists\_chord}); at the level of
classes, the attitude class equals the union, over the admissible
chords, of their crescents
(\texttt{hurwiczSuperlevel\_eq\_iUnion\_crescents}). An attitude
class is therefore a disjunction of Specs --- an \emph{infinitary}
one, indexed by a half-plane of chords.

Infinitary, and irreducibly so. For a genuinely mixed attitude,
$0 < \alpha < 1$, no finite subfamily of chords yields the class:
against any finite list of admissible chords, a two-point model
sits on the Hurwicz line beyond every listed escape face, meeting
the threshold while escaping every listed crescent
(\texttt{finite\_chords\_strict}). Figure~\ref{fig:su-chords}
shows the geometry. In the plane of the two envelope values, one chord's
crescent covers a closed quadrant; the attitude class covers the
half-plane above a line of strictly negative slope; and finitely
many quadrants make a staircase, which never covers a sloped line.
Since every Spec denotes a finite union of crescents, no Spec
denotes an attitude class. The grammar's root-only restriction
thus reappears at the class stratum as a theorem about
denotations: the operator does not merely lack inference rules;
the classes it carves lie outside the algebra every Spec inhabits.

\begin{figure}[t]
\centering
\begin{tikzpicture}[x=14mm,y=11mm]
  \fill[black!6] (0,3) -- (6,0) -- (7.0,0) -- (7.0,3.6) -- (0,3.6)
    -- cycle;
  \fill[black!15] (1.6,2.2) rectangle (7.0,3.6);
  \fill[black!15] (3.0,1.5) rectangle (7.0,3.6);
  \fill[black!15] (4.4,0.8) rectangle (7.0,3.6);
  \draw[semithick] (1.6,3.6) -- (1.6,2.2) -- (3.0,2.2) -- (3.0,1.5)
    -- (4.4,1.5) -- (4.4,0.8) -- (7.0,0.8);
  \draw (0,3) -- (6,0)
    node[pos=0.22, below=3pt, sloped, font=\small]
    {$\alpha a + (1-\alpha)\,b = \vartheta$};
  \fill (1.6,2.2) circle (1.6pt);
  \fill (3.0,1.5) circle (1.6pt);
  \fill (4.4,0.8) circle (1.6pt);
  \node[star, star points=5, star point ratio=2.2, fill=black,
    inner sep=1.1pt] at (5.0,0.5) {};
  \node[font=\small, anchor=west] at (5.14,0.5)
    {escapes all three};
  \node[font=\small, anchor=west] at (4.5,3.3)
    {three listed crescents};
  \node[font=\small, align=center] at (2.05,0.72)
    {the attitude class:\\ the whole shaded half-plane};
  \draw[-{Stealth[length=1.8mm]}] (0,0) -- (7.35,0)
    node[below=2pt, anchor=north east, font=\small]
    {$a$ \ (escape level: best case)};
  \draw[-{Stealth[length=1.8mm]}] (0,0) -- (0,3.85)
    node[right=2pt, font=\small]
    {$b$ \ (floor level: worst case)};
\end{tikzpicture}
\caption{Why no finite disjunction reaches a mixed attitude,
  drawn in the plane of the two envelope values
  $(\overline{E}(h), \underline{E}(h))$. Each admissible chord
  $(a,b)$ contributes one crescent, which covers the closed
  quadrant above and to the right of its corner; chords are
  demands, so their corners need not be attainable pairs. The
  attitude class covers the half-plane above the chord line.
  Finitely many quadrants form a staircase, and the starred model
  sits on the line beyond every listed escape face: it meets the
  Hurwicz threshold and escapes every listed crescent
  (\texttt{finite\_chords\_strict}).}
\label{fig:su-chords}
\end{figure}

The boundary has a one-line grammar: \emph{a Spec mixes verdicts;
an attitude mixes envelopes}. Each literal of a Spec consults one
envelope pass and returns a Boolean, and the connectives combine
finished Booleans, which is why compound judgments transport face
by face. $H_\alpha$ adds the two passes into one scalar before any
verdict is taken, and the sum behaves as neither pass: it has no
single polarity, so no one face discharges it; it is not monotone
along refinement, so it is not an entailment-shaped operator; and
it is not stable under composition, because the chord attained at
the root re-chooses when kernels chain --- which is what the tower
failure of \autoref{sec:su-denotation} says, read through the
decomposition.

What survives for certification is the chord itself. A floor
certificate at $b$, an escape certificate at $a$, and the
arithmetic check $\vartheta \le \alpha a + (1-\alpha) b$ together
discharge the threshold (\texttt{le\_hurwicz\_of\_chord}); a
certificate pair \emph{is} a choice of disjunct. A finitely
presented certificate layer for attitude thresholds therefore
carries its chord as data and never enumerates the class --- one
more instance of a pattern this universe keeps producing: where a
class refuses finite presentation, its certificates do not
(compare the disjunction property's choice byte above, and the
properness mark).

\subsection{Beneath the Atoms: the Graded Layer}

A crescent's faces test bodies, and a body is a belief's shadow.
The compile of \autoref{sec:su-hyperspace} sends each belief
$\sigma$ to the closed radial set of laws it admits at each grade.
A \emph{wall} --- a $\{0,\infty\}$-valued belief --- casts a
crisp shadow, the set of laws it admits outright; a graded
$\sigma$ knows strictly more, namely how expensively each law is
admitted. The atoms of this section consult shadows. So a fair
question about the whole construction is whether graded beliefs
eventually demand graded classes --- some second fibre, richer
than the constructible one, where soft judgments would live. They
do not, and the exact stratum now carries the reason as a pair of
theorems.

At the exact stratum a convex prevision arrives finitely
presented: finitely many laws $M_i$ with signed rational costs
$c_i$, pricing a claim $f$ at $\max_i\,(E_{M_i}[f] - c_i)$ --- the
penalty-function representation (so named in that literature) of
convex risk measures
\cite{follmerschied02}, restricted to finite presentations
(\texttt{penUpperE}, with \texttt{penLowerE} its conjugate). The
wall tier is the zero-cost face, definitionally
(\texttt{penUpperE\_zero\_cost}), and the precise tier is the
one-member face. The first theorem says the grades never leave the
reach of the classes: the scalar is the finite Legendre transform
of its own threshold profile --- the map sending each
affordability threshold $t$ to the best expectation among members
costing at most $t$ --- and the profile's data at each threshold
is an ordinary credal set
(\texttt{penUpperE\_eq\_sup\_levelProfile}). A graded judgment
therefore unfolds into a family of possibilistic judgments, one
crescent face per threshold, and the scalar folds back by the
transform. The second theorem says composition respects the
unfolding: belief kernels compose by an exact intertwining
square whose cost law --- costs integrate along the record --- is
forced by countermodels on either side
(\texttt{penUpperE\_compFam}; the unweighted and the worst-case
aggregations each break the square on a two-state
witness).\footnote{The tier's calculus holds a few further facts a
practitioner will want. Convexity replaces coherence: a mixture of
claims is over-priced by the mixture of prices
(\texttt{penUpperE\_mix\_le}), while positive homogeneity
genuinely fails off the calibrated face, by a one-point witness
(\texttt{penUpperE\_not\_posHomogeneous}) --- the tier is
convex-prevision, strictly wider than Walley coherence. And the
point-mass family with a pointwise gate prices claims as a
max-plus contraction (\texttt{penUpperE\_dirac\_family}), the
idempotent-integration reading that connects this layer to
tropical dynamic programming.} These are theorems of the exact
stratum --- finitely supported laws, rational costs, where
certificates live; the general-carrier belief calculus of
\autoref{sec:su-denotation} stands beneath them unchanged. The
atoms stay possibilistic; the grades thread through them one
threshold at a time; and the classes of this section carry the
soft tier without gaining a single new former.

\subsection{The Completion Question}

Two pressures point at one enlargement. The attitude classes ask
for unions of crescents over a compact index; the graded layer,
read at every threshold at once, asks for level families of the
same shape. Both land in the $\sigma$-completion of the
constructible algebra: classes presented as directed unions of
crescents over declared index families, with the index carried as
data. Nothing wild lives there. Each candidate class is assembled
from the same atoms as every Spec, and its certificates are
already finite objects --- a chord, a threshold, a pair of faces.

What is not settled is transport, and the chord decomposition says
why the question has teeth. The constructible fibre's
functoriality is a theorem because preimage acts on finitely many
faces at once. An enlarged fibre would transport index families,
and the transport is pointwise in the index: under composition the
attained chord re-chooses, which is the tower failure again, now
as a design constraint rather than a pathology. A completion that
demanded one uniform chord would refute itself at the first kernel
chain. The candidate discipline, stated here as design and not as
artefact: admit closed classes that declare semicontinuity
certificates for both envelope passes, and transport them only
along wirings carrying both certificates --- the two-pass
evaluation of \autoref{sec:su-denotation} made compositional.
Whether the resulting fibre keeps the composition calculus's
kernel-checked functoriality, and whether it is the largest fibre
that does, are open questions; the closing ledger names them
(\autoref{sec:su-receipts}). Until they close, this book's
judgments live on the constructible stratum, where every demand of
this section is discharged by a theorem.

\section{The Fibre Sheaf: Spec Classes over the Site}
\label{sec:su-sheaf}

The class stratum of \autoref{sec:su-classes} was built at one
window. Records live over a site, so spec classes do too, and this
section fixes how: the classes of \autoref{sec:su-classes},
window by window, assemble into a presheaf whose sheaf-theoretic
behaviour is governed by the carrier results of
\autoref{sec:su-carrier} --- with one structural surprise, itself
kernel-checked, about \emph{which} window-indexed families are
admissible.

\paragraph{Window beliefs, and the variance of extension.} Over a
window $W$ and an instantaneous carrier $X$, the trajectory space
is the record carrier of \autoref{sec:su-carrier}, taken at state
fibre $X$:
\[
  \mathrm{Traj}(W, X) \;:=\;
  \mathrm{CL}\bigl(\mathring{S}(W) \times X^{+}\bigr)_{\mathrm{Fell}},
\]
closed subsets of the stage-interior-times-state ambient under the
Fell topology, with $X^{+}$ the carrier with its cemetery adjoined
--- $\mathbb{H}(W)$ itself when $X$ is the window's own state
fibre. Functions $W \to X$ embed by \emph{graph-then-closure},
$f \mapsto \mathrm{cl}\,\{(t, f(t))\}$: on the speed-bounded class
the graph is already closed and the closure does nothing --- the
graph reading of \autoref{sec:su-records}, recovered as the special
case --- while at a jump the graph is \emph{not} closed, and the
closure charges the jump instant with both one-sided values: the
closed-graph doctrine again, a jump multivalued at the moment it
happens. The embedding is accordingly not injective on jump
functions --- two functions differing only in which one-sided limit
they assign at the jump instant close to the same set --- and that
quotient is intended, not a defect: the closed set remembers every
value the instant is honestly charged with and forgets the
convention that chose among them. Window beliefs are beliefs
over this carrier: functionals carried, in the
working presentation of \autoref{sec:su-denotation}, on
subprobability joints ---
probability joints over the cemetery-completed carriers --- at the
standing generality of
Definition~\ref{def:su-credal-denotation}; these are laws of
exactly the kind the capacity presentation of
\autoref{sec:su-carrier} pins by countably many rational
observations. The scalar stock of that
generality is kernel-checked as stated --- subprobability Borel
joints over arbitrary measurable carriers, with pushforward,
belief transport, weak continuity, and the lower-semicontinuity
closure and attainment package
(\texttt{SuFunctor.Plain.MeasJoint}; the composition tier below
carries the details) --- with the finitely supported joints of the
exact stratum its kernel-checked sub-instance along the seam
(\texttt{toMeas\_push}); nothing in the stock reads how the carrier
is presented, so it serves the Fell carrier and the function
presentation alike. The window \emph{assembly} is kernel-checked at
both joint strata over the \emph{function presentation} (windows as
subsets of a stage type, trajectory carriers the graphs before
closure): the extension calculus and the persistence discipline of
this section replay verbatim over subprobability Borel joints, with
measurability as the morphism datum --- the presented carriers get
the product $\sigma$-algebra, under which clipping and pointwise
state transport are measurable outright, and the window--compositor
exchange is one application of the base-free naturality core with
both comparison squares discharged definitionally
(\texttt{SuFunctor.Site.GeneralWindow}:
\texttt{penExtM} with its presheaf laws, both exchange theorems,
and the closure package for persistent families) --- and the exact
assembly is the general one's kernel-checked sub-instance along the
seam (\texttt{penExtM\_comp\_toMeas},
\texttt{IsPersistentM.seamRestrict}). The same assembly restated
over the Fell carrier itself --- Borel-on-Fell as the measurability
datum, the carrier section's clip as the restriction --- is an open
obligation of the closing ledger, not a receipt in hand; the
restriction leg is not part of the debt, its measurability being a
corollary of the landed Fell-continuity of clip-to-open
(Remark~\ref{rem:su-carrier-mechanized}). The descent results below
were never joint-stratum-bound at all: their ambient is an
arbitrary topological space (the class-descent module is
fibre-class-blind), so no stratum scoping applies to them. Two
boundaries are real and worth naming. The double-categorical
\emph{total} of the site-indexed decoration functor over a wiring
base stays at the exact stratum --- the finite state-family
presentation discharges its comparison squares definitionally, and
a measurable wiring base owes its own transport data first, the
same residue one tier up as in
Remark~\ref{rem:su-house-functor-mechanized}. And the matching-fork
witness below is stated at the exact stratum because nothing more
is needed there: its trajectory carriers are finite, and over
finite carriers every presentation collapses into every other ---
a finite Hausdorff ambient is discrete, so every graph is closed,
graph-then-closure is the plain graph embedding and is injective,
and the two joint strata literally coincide, the seam embedding
being a bijection (\texttt{toMeas\_surjective} with
\texttt{toMeas\_injective}) --- so the fork refutes the strict
discipline at the general stratum, and on the Fell carrier, in the
same breath. Which
\emph{evaluation-level} claims complete across the seam by density
--- and which provably do not --- is
Remark~\ref{rem:su-density-completion}.
For nested windows $V \subseteq W$, trajectories \emph{clip}
($\mathrm{Traj}(W,X) \to \mathrm{Traj}(V,X)$: intersection with the
subwindow's open band --- the carrier section's Fell-continuous
restriction, at fixed state fibre the pure clip of $\rho_u$), and
the clip agrees with the function presentation through the
embedding: clipping a closed graph over an open sub-band is the
closure of the restricted function's graph, because
closure-then-intersect is relative closure over an \emph{open} set
--- one more place the interior coverage earns its keep. Beliefs
therefore \emph{extend}: a subwindow demand
$u$ is read over the superwindow as ``evaluate $u$ on the clip''
($\mathrm{penExt} := $ precomposition of the joint pushforward along
clipping; \texttt{SuFunctor.Site.penExt}, on the presented
carriers). The variance is forced: a
belief reads whole-window
joints, and only the small-to-big direction has a canonical
trajectory map --- a closed set clips canonically and extends
non-canonically, so there is no restriction of beliefs. Their \emph{bodies}
restrict --- that is the compile trade: heights extend along the
site, bodies clip along the site, and the two motions are exchanged
by the compile isomorphism. Both motions are used below, each on its
own side.

\begin{remark}[Relations first; functions as a dense locus]
\label{rem:su-relations-first}
The carrier just fixed takes closed \emph{relations} as its primary
objects, and the function layer enters as a distinguished locus
rather than as the ambient notion --- in three precise senses. The
locus is \emph{dense}: wherever the stage interior has no isolated
points (every nondegenerate classical window), each nonempty basic
hit-and-miss class contains the graph of a finitely supported
single-valued function on a sub-timeline --- one point in each hit
region, off the miss compact, shared stages separated within the
open regions --- so every closed relation is a Fell limit of
function graphs. Its calculus is a \emph{restriction}: the window
assembly above, kernel-checked today on the function presentation,
is designed as the relational assembly read along
graph-then-closure, and the relational form itself is the closing
ledger's Fell-carrier obligation, not a receipt in hand. And the
correspondence-theoretic apparatus arrives with no added
hypotheses: for a closed-valued map into the compact state fibre,
lower hemicontinuity is continuity for the hit face and upper
hemicontinuity is continuity for the miss face --- compactness of
the fibre is what turns a closed complement into a miss parameter
--- so hemicontinuity is the slice-wise reading of the topology the
carrier already wears, not a regularity assumption on a new kind of
object.
\end{remark}

\begin{remark}[Density completion: what the exact stratum settles]
\label{rem:su-density-completion}
The exact stratum invites a completion argument: finitely supported
joints are narrowly dense in the subprobability carrier, so should
not every claim kernel-checked at the exact stratum extend to the
general carrier by taking limits?  The answer has a shape, and both
halves of it are kernel-checked
(\texttt{SuFunctor.Plain.SeamDensity}).

First, vocabulary: ``exact'' names finite \emph{support}, not
rational \emph{data}.  The exact stratum's weights are arbitrary
reals in $[0,1]$; the data a certificate checker consumes ---
rational weights, admitting exact arithmetic --- is a further
sub-stratum.  Its density step is kernel-checked: rounding weights
down to multiples of $2^{-n}$ converges narrowly, mass only
dropping, so wherever the seam is dense, the dyadic-rational seam
is already dense (\texttt{Joint.dense\_ratSeam}).  Density of the
seam itself is the classical statement of
\cite[Thm.~II.6.3]{parthasarathy67}, and is kernel-checked here over
separable pseudometrizable carriers
(\texttt{Joint.denseRange\_toMeas},
\texttt{Joint.denseRange\_toMeas\_of\_pseudoMetrizable} --- covering
the chapter's Polish standing generality), with the truncation
instance over $\mathbb{N}$ proved by an independent second route
(\texttt{Joint.denseRange\_toMeas\_nat}); the transport package is
stated against an explicit density hypothesis, so each transport
theorem applies at exactly the carriers where density is on record.

Second, the completion itself is one-sided, and the direction is
forced by the universe's own membership hypothesis.  Lower
semicontinuity makes sublevel conditions closed, and closed sets
absorb dense subsets; so \emph{upper} bounds transport --- a lower
semicontinuous belief bounded on the seam by an upper
semicontinuous one is so bounded everywhere
(\texttt{MeasPen.le\_of\_seam\_le}) --- suprema are seam-determined
(\texttt{MeasPen.iSup\_seam\_eq}), and narrowly continuous
beliefs agreeing on the seam agree everywhere
(\texttt{MeasPen.eq\_of\_seam\_eq}).  \emph{Lower} bounds and
infima do not transport, and the failure is semantic rather than a
mechanization debt, witnessed at a carrier where density holds: the
drop wall --- zero at the geometric joint over $\mathbb{N}$,
infinite elsewhere --- is lower semicontinuous (its one proper
sublevel set is a closed singleton), sits at infinity on the entire
seam, and has infimum zero on the carrier
(\texttt{dense\_seam\_inf\_gap}).  A wall may open exactly at an
infinitely supported joint, and no net of finitely supported joints
sees it open.  So a verdict computed on the exact stratum is exact
for suprema and upper bounds and can be unboundedly optimistic for
infima and lower bounds; the checker-facing consequence is that
safety ceilings certified on rational data complete to the full
carrier, while worst-case floors need either upper semicontinuity
or the general carrier itself.

What remains constitutively exact-stratum is computation, not
semantics: exact rational comparison and decidability are
finite-data notions, which is why the checker face lives there ---
and at finite outcome carriers even the carriers coincide, the seam
embedding being a bijection (\texttt{toMeas\_surjective}).
\end{remark}

\paragraph{The presheaf of classes, and the fork.} The canonical
window-restriction of a \emph{class} is the preimage under
extension: a spec class $\Phi_W$ over $W$ restricts to the subwindow
demands it certifies, $\Phi_V = \mathrm{penExt}^{-1}\,\Phi_W$ ---
the presheaf of spec classes over the site. A window-indexed family
can then satisfy two disciplines.

\begin{definition}[Matching and persistent families]
\label{def:su-matching-persistent}
Let $\Phi = (\Phi_i)_i$ be a window-indexed family of spec classes
over a window system $U$, with extension maps
$\mathrm{penExt}_{ij}$ for $U_i \subseteq U_j$.
\begin{itemize}
\item \emph{Matching families} (strict sections):
  $\Phi_i = \mathrm{penExt}_{ij}^{-1}\,\Phi_j$ for every
  $U_i \subseteq U_j$
  \hfill (\texttt{SuFunctor.Site.IsMatching})
\item \emph{Persistent families} (lax sections): extension carries
  each window's members into every superwindow's class,
  $\mathrm{penExt}_{ij}[\Phi_i] \subseteq \Phi_j$
  \hfill (\texttt{SuFunctor.Site.IsPersistent})
\end{itemize}
\end{definition}

The fork between the two disciplines is
settled by a theorem, and it is the one place the site dimension
exerts force on the composition calculus:

\begin{theorem}[The matching fork]
\label{thm:su-matching-fork}
Matching families are not closed under the window-wise composition
of decorations: there exist matching families $\Phi, \Psi$ over a
two-window system whose componentwise composite (the window
compositor of Definition~\ref{def:su-house-functor}, applied per
window) is not matching. Persistent families are closed under every
operation of the decoration calculus --- tight transport, cell
transport, composition, and the unit.
\end{theorem}

The witness is small enough to describe: over a stage line with one
extra time point, take the two \emph{support walls} through the
shared middle boundary --- one demanding the extra point read
$\mathsf{true}$, the other $\mathsf{false}$. Each wall generates a
matching family; their composite at the superwindow is the
zero-support wall (the demands contradict through the shared
middle), which \emph{is} clip-invariant --- so the composite's
subwindow component, being empty, fails the matching equation
against it. Composition manufactures window-local content (here:
infeasibility) that strict sections cannot carry
(kernel-checked, with the witness exactly as described:
\texttt{SuFunctor.Site.MatchingFork.matching\_not\_closed\_under\_compositor}, against the closure theorem
\texttt{IsPersistent.image2\_boxplusW}). The fibre objects of the
construction are therefore the \emph{persistent} families --- the lax
sections, equivalently the families closed under
reading-small-demands-large, the least-committal completion of
window-local spec data. The pattern recurs across this universe's gluing theory: record carriers glue
exactly (the carrier sheaf, at its stated hypotheses), assessments
complete least-committally (the natural extension, adjoint to
restriction), and here admissibility itself completes
least-committally --- in each
case the naive expectation locates laxity in the composition
operation, and in each case the operations stay exact while the
completion happens one construction over. (The gluing clause of that
pattern carries the carrier proposition's own hypotheses --- directed
interior covers, fixed fibre --- and the adjunction runs with the
natural extension on the left: least among the coherent globals
whose restrictions dominate the assessment, at the universal
property's stated scope.)

\paragraph{Class descent, and ideal points as global members.} Over
a directed open cover, the descent homeomorphism of
Theorem~\ref{thm:su-fell-descent} induces an order-isomorphism of
class fibres,
\[
  \mathrm{Set}\bigl(\mathbb{H}(W)\bigr) \;\cong\;
  \mathrm{Set}\Bigl(\varprojlim_i \mathbb{H}(W_i)\Bigr),
\]
carrying constructible classes to constructible classes in both
directions --- the site's fibre stratum and the glued fibre stratum
have the same constructible layer (kernel-checked:
\texttt{SuFunctor.Site.classDescent},
\texttt{isConstructibleUnion\_classDescent\_iff}, with the
transport lemmas \texttt{IsLocallyClosed.classExt} and
\texttt{IsConstructibleUnion.classExt} for the free contravariant
direction). And the ideal points of the fibre sheaf are exactly the
global members: compatible singleton-sections of the point presheaf
--- per window the clip of one body, compatible by the clip
calculus --- biject with closed members of the glued ambient
(\texttt{SuFunctor.Site.bijective\_toProjLim}, the descent
homeomorphism's underlying bijection consumed whole). A System over
the glued window \emph{is} its compatible family of window bodies ---
on the nonempty-body stratum the compiled-System scope of
\autoref{sec:su-denotation} draws --- and membership in a
persistent family is window-local along pullback
(\texttt{clip\_mem\_classExt\_iff}).

\begin{remark}[The full fibre column, and the localic reading]
\label{rem:su-fibre-column}
Assembled, the fibre over a window is a four-storey tower, each
storey kernel-checked at the strata named in its own section: the
\emph{record locus} (closed subsheaf of the Fell carrier at the
classical stratum,
\autoref{sec:su-carrier}); \emph{laws} over it (the
subprobability stratum of \autoref{sec:su-denotation},
pinned by the countable rational determination of
Theorem~\ref{thm:su-capacity}); the \emph{hyperspace} of denotation
bodies; and the \emph{constructible classes} of
\autoref{sec:su-classes} --- crescents in cover-modality shape,
closed under the connectives, ordered by inclusion, with
membership and inclusion decided by the blocking
criterion. The design-normative reading of every storey above the
$\omega$-semialgebraic state spaces is
pointfree --- the record storey as the Vietoris powerlocale of the
compactified ambient object (the carrier section's own localic
reading), the law storey as the valuation locale of the
record locale, the hyperspace storey as its Vietoris powerlocale,
each presented by generators and relations (the cover-modality
presentation of \cite{venema12}) so that the whole column
externalizes objectwise over the site by geometricity.

That reading is no longer unmechanized design. The point-free
library carries a kernel-checked tier of its own, stated on frames
--- no points, no countability, no $\sigma$-algebra. Checked there:
the \emph{presented hyperspace frame} --- the free frame on
hit-and-miss data, quotiented by the least nucleus forcing the
transposed hit-and-miss laws --- with the Choquet--Matheron
determination on its $\nabla$-generators, so that the classical
theory is recognized as having been pointfree all along, now as a
kernel theorem at the localic stratum
(Remark~\ref{rem:su-capacity-mechanized});
the \emph{valuation layer} of the law storey, frame-general ---
valuations with modularity load-bearing, pushforward functorial and
the Dirac unit natural (the presheaf structure the storey needs),
the dagger promotion chart driven by modularity alone, and the
narrow topology \emph{forced}, not chosen: it is the lower topology
of the promoted carrier pulled back along the chart
(\texttt{Pointfree.Valuation.Basic}, \texttt{.Promotion},
\texttt{.Topology}, \texttt{.Functor}); and the \emph{two-tier tie}
--- restriction of Borel measures to opens is Scott-continuous
exactly on the $\tau$-additive stratum (Fremlin's 414O supplies that
stratum under second countability; the Dieudonn\'e measure fences
it), continuous from the weak into the narrow topology by the
portmanteau open half, natural in pushforward, with the
measure-level dagger chart a homeomorphism carrying Prokhorov
compactness to the subprobability carrier, the escape phenomenon
witnessed without points (a free-ultrafilter valuation refusing
Scott continuity against a directed family of opens), and
$\tau$-additivity re-read as order-continuity into the measure
algebra (\texttt{.Bridge}, \texttt{.BridgeContinuity},
\texttt{.BridgeFunctor}, \texttt{.EscapeWitness},
\texttt{Pointfree.Measure.DaggerChart},
\texttt{Pointfree.MeasureAlgebra.Basic}).

What remains argued design is the column reading itself: the
identification of each presented frame with the powerlocale of the
storey below it (equivalently, that spatialization returns the
classical Fell hyperspace --- the determination theorem is proven at
arbitrary relation sets exactly so that nothing kernel-checked gates
on this identification); the geometricity by which the column
externalizes over the site; and the prevision-side completion ---
every Scott-continuous superlinear prevision the lower envelope of
its body of subprobability valuations, the double-powerspace reading
of \cite{goubaultlarrecq26} --- of which the kernel now holds the
sublevel-presented gambles, the Choquet pairing, closure of the
superlinear class under lower envelopes, and the cheap half (every
prevision below the envelope of its body;
\texttt{Pointfree.Prevision.Basic}), the representation converse
cited.

The boundary, stated exactly: part of the mechanized tier now
\emph{is} the localic reading --- the determination theorem consumes
the presented frame, and the valuation layer is stated on arbitrary
frames outright --- while the spine this chapter's sections consume
stays deliberately point-set, exact on compact windows, where
continuous valuations and Radon measures agree, and the two tiers
exchange theorems through the kernel-checked bridge rather than by
analogy. Everything classical is stated so as not to obstruct the
localic reading; the localic reading now answers with theorems of
its own.
\end{remark}

\section{Composition: the Double Grothendieck Construction}
\label{sec:su-composition}

Chapter~\ref{ch:qtsl-composition-patterns} catalogues the
composition patterns a modeller wields --- fusion, pools, dispatch,
delayed trace. This section is the stratum underneath: the single
categorical construction in which any wiring of open Systems composes,
lawfully, so that pattern-level soundness claims have one place to
discharge. The route is direct: span-valued lax double
functors on the interfaces-and-wirings base,
consumed by the double Grothendieck
construction at its fully mechanized stratum, with every obligation
--- the functors', the totals', and the projections' ---
kernel-checked. Two functors share the base and the machinery, in a
division of labour stated where each is defined: the
\emph{Spec} functor carries arbitrary classes and the site
dimension (the window-indexed persistent families of
\autoref{sec:su-sheaf} --- the fibre of Specs); the
\emph{coupled-DNF} functor carries
the constructible presentations of \autoref{sec:su-classes} as
fibre data,
refines the Spec functor along a kernel-checked comparison cell,
and its total is the development's authoritative
\texttt{SingleSiteSemanticUniverse} --- the double category of System Specs,
whose site-indexed culmination, the chapter's headline object, is
stated whole in \autoref{sec:su-total}.
Two boundaries of that opening are doing
work. The
operation is total on composable pairs; the conditions that live
deliberately \emph{outside} the functor --- causality, simultaneity
discipline, Zeno exclusion --- are side-condition oracles recomputed
on each composite, because they could not be object-level: two
individually causal couplings can close an algebraic loop, so
composition does not preserve them. And the fibres are \emph{Specs},
not System entities: a System enters the total through
its denotation, as an ideal point of a fibre, and the entity-level
stratum (simulation witnesses as fibre morphisms, refinement as a
fibred structure over this one) is a design layer above the
total, recorded in the ledger, not part of it.
Figure~\ref{fig:su-two-storey} draws the destination: the wiring base
on the floor, and above one wiring its fibre of decorations, where a
System's ideal point, a Spec's class, and the judgment between them
all live.

The ambient notions come first, at reference width --- this chapter
is the book's authority for them.

\subsection{The Ambient Notions}

\begin{figure}[t]
\centering
\begin{tikzpicture}[x=1mm,y=1mm,
  loosear/.style={-{Stealth[length=1.8mm]}, semithick},
  tightar/.style={-{Stealth[length=1.8mm]}, semithick},
  fibre/.style={densely dashed, gray},
  looseslash/.style={fill=white, inner sep=0.8pt, font=\footnotesize},
  dot/.style={circle, fill, inner sep=0.55pt}]
  \fill[gray!10] (0,0) -- (58,0) -- (72,-10) -- (14,-10) -- cycle;
  \fill[gray!4]  (-2,11.5) rectangle (60,47);
  \node (I)  at (0,0)    {$I$};
  \node (J)  at (58,0)   {$J$};
  \node (Ip) at (14,-10) {$I'$};
  \node (Jp) at (72,-10) {$J'$};
  \draw[loosear] (I)  -- node[looseslash, midway] {$\shortmid$}
                    node[below=1mm, pos=0.28] {$M$} (J);
  \draw[loosear] (Ip) -- node[looseslash, midway] {$\shortmid$}
                    node[below=1.2mm, midway] {$N$} (Jp);
  \draw[tightar] (I) -- node[below left=-0.5mm, midway] {$f$} (Ip);
  \draw[tightar] (J) -- node[right=0.5mm, midway] {$g$} (Jp);
  \node at (40,-5.4) {$\Downarrow\,\gamma$};
  \draw[fibre] (0,2.6) -- (0,46);
  \draw[fibre] (58,2.6) -- (58,46);
  \draw[loosear] (0,42) -- node[looseslash, midway] {$\shortmid$}
      node[above=1mm, pos=0.13] {$\{\llbracket S \rrbracket\}$} (58,42);
  \draw[loosear] (0,28) -- node[looseslash, midway] {$\shortmid$}
      node[above=1mm, pos=0.13] {$\llbracket \varphi \rrbracket$} (58,28);
  \draw[loosear] (0,14) -- node[looseslash, midway] {$\shortmid$}
      node[above=1mm, pos=0.13] {$\top$} (58,14);
  \foreach \y in {14,28,42} {\node[dot] at (0,\y) {}; \node[dot] at (58,\y) {};}
  \node at (29,35) {$\Downarrow\; S \vDash \varphi$};
  \node at (29,21) {$\Downarrow$};
  \draw[tightar, gray] (-6,12) -- (-6,1.5);
  \node[left, font=\footnotesize, align=right, gray] at (-6.8,7)
    {strip the\\decorations};
\end{tikzpicture}
\caption{The two storeys of the composition construction. On the
floor lies the wiring base: interfaces $I$, $J$, wire maps $f$, $g$
(tight), wiring patterns $M$, $N$ (loose, slashed), and a pattern
cell $\gamma$; tight morphisms compose strictly, loose arrows up to
the invertible structure cells. Hovering above the arrow $M$ sits
its fibre, drawn upward in the plane the dashed verticals span:
each storey is $M$ again, now carrying a decoration, and height
orders the decorations by how much they assert. A System $S$ enters
at the top, as the singleton $\{\llbracket S \rrbracket\}$ of its
own denotation --- its ideal point, the sharpest decoration $S$
satisfies; a Spec $\varphi$ enters as the class
$\llbracket \varphi \rrbracket$ it denotes over the wiring; the
bottom storey $\top$ admits every belief and demands nothing. The
judgment $S \vDash \varphi$ of this chapter's opening, unfolded to
$\llbracket S \rrbracket \in \llbracket \varphi \rrbracket$, appears
as the globular cell between the top two storeys: an inclusion of
decorations over a fixed pattern is a cell of the kernel-checked
judgment layer (\texttt{SuFunctor.QtslSem.SpecJudg}:
\texttt{SpecJudgLayer}, \texttt{cellOfSubset}), and the layer is
thin, so the cell carries no data beyond the judgment it witnesses.
The tower descends from sharp to weak: every class includes into
$\top$, so $S \vDash \top$ comes for free as the composite cell,
and $S \vDash \varphi$ is the storey that costs something. Judgment
cells also compose loosewise, along the wiring: an initial
condition enters as one more decorated arrow and conjoins through
the compositor, which is monotone (\texttt{wComp\_mono}), while an
undemanded face costs nothing (\texttt{wUnit\_transparent}); this
is how component judgments assemble the composite judgment at the
composite wiring. Stripping the decorations lands every storey on
$M$ and every cell shown upstairs on the identity cell of $M$. The
total of Theorem~\ref{thm:su-assembly} refines this picture: its
objects fix a designated Spec at each interface, and its strict
projection to the floor lifts uniquely --- a discrete double
opfibration. The designations at the objects are chaining
contracts, matched on the nose when wirings compose; entailment
lives upstairs, in the fibre order, as drawn.}
\label{fig:su-two-storey}
\end{figure}

\begin{definition}[Pseudo double category]
\label{def:su-pseudo-double-category}
A \emph{pseudo double category} $\mathbb{D}$ consists of: a category
$\mathbb{D}_0$ of \emph{objects} and \emph{tight morphisms}; for each
pair of objects, a collection of \emph{loose arrows}
$m \colon X \nrightarrow Y$; for all loose arrows
$m \colon X \nrightarrow Y$, $n \colon X' \nrightarrow Y'$ and tight
morphisms $f \colon X \to X'$, $g \colon Y \to Y'$, a collection of
\emph{cells} $\alpha \colon m \to n$ over $(f, g)$; tight-wise
identity cells and tight-wise composition of cells; loose identities
$\mathrm{id}_X \colon X \nrightarrow X$ and loose composition
$m \odot n$, with loose-identity cells on tight morphisms and
loose-wise composition of cells; and invertible globular structure
cells --- associators $a_{m,n,p} \colon (m \odot n) \odot p \to
m \odot (n \odot p)$ and unitors $\lambda_m, \rho_m$ --- subject to:
the tight-wise category laws (loose arrows and cells form a category
$\mathbb{D}_1$ whose source and target functors land in
$\mathbb{D}_0$); functoriality of the loose identity and of loose
composition with respect to tight-wise composition (the interchange
law); tight-wise naturality and invertibility of the structure
cells; and the pentagon and triangle laws. Equations between cells
whose tight boundaries agree only up to the category laws of
$\mathbb{D}_0$ are read modulo the canonical transport of cells
along those laws. The notion is Grandis and Par\'e's weak double
category \cite{grandis99, grandis04}; the tight/loose terminology
follows Shulman \cite{shulman08}.
When the structure cells are identities, $\mathbb{D}$ is a strict
double category.
\end{definition}

\begin{remark}[Mechanization of the ambient notions]
\label{rem:su-pdc-mechanized}
Kernel-checked, in the double-category library of the proof-artefact
repository: the data and the eighteen axiom clauses as
\texttt{PseudoDoubleCategoryStruct} and
\texttt{PseudoDoubleCategory} in module \texttt{DoubleCat.Basic} ---
cells dependently indexed by their tight boundaries, with the
boundary identifications this definition reads silently carried by
an explicit transport operator (\texttt{Cell.cast}), and every axiom
whose two sides have propositionally-but-not-syntactically equal
boundaries stated up to that transport. Lax double functors
(Definition~\ref{def:su-lax-double-functor}) are
\texttt{LaxDoubleFunctor} in \texttt{DoubleCat.Functor.Lax}, with
pseudo-ness as an invertibility property. The library further
carries, at the same kernel grade: tight and loose transformations
and modifications, assembling the functor double category; the loose
bicategory bridge, through which the library inherits bicategorical
coherence rather than re-proving it; companions, conjoints, and
equipments, with spans the first genuine equipment (and a proof that
the double category of commutative squares is not one); and the span
and matrix models, including the dependent-span model
(\texttt{DoubleCat.DepSpan}) on which the Grothendieck construction
below is stated.
\end{remark}

\begin{definition}[Lax double functor]
\label{def:su-lax-double-functor}
Let $\mathbb{A}$ and $\mathbb{B}$ be pseudo double categories. A
\emph{lax double functor} $F \colon \mathbb{A} \to \mathbb{B}$
consists of a functor $\mathbb{A}_0 \to \mathbb{B}_0$ on objects and
tight morphisms; an action on loose arrows
($m \colon X \nrightarrow Y$ to
$Fm \colon FX \nrightarrow FY$) and on cells (preserving tight
boundaries strictly, functorially for tight-wise composition); a
\emph{compositor} family of globular cells
$Fm \odot Fn \to F(m \odot n)$; and a \emph{unit comparison} family
$\mathrm{id}_{FX} \to F(\mathrm{id}_X)$ --- subject to naturality of
both families in cells and tight morphisms respectively, associativity
coherence of the compositor against the associators, and the two
unit-triangle coherences against the unitors. $F$ is
\emph{pseudo} when compositor and unit comparison are invertible,
and \emph{strict} when they are identities.
\end{definition}

\subsection{The Wiring Base}

The base of the construction is the
interfaces-and-wirings layer, itself kernel-checked, in a
presentation chosen so
that the transport identities a decoration tier consumes hold
definitionally. The semantic description first, then the
presentation.

\begin{definition}[Wire types, interfaces, and state spaces]
\label{def:su-wire-category}
Fix a set $\mathbb{T}$ of \emph{wire types}, and for each
$t \in \mathbb{T}$ a nonempty Polish space $\mathrm{El}(t)$ of
states of that type, per the standing structure of
Definition~\ref{def:su-credal-denotation}. An \emph{interface} is a
finite set of wires typed over $\mathbb{T}$; interfaces and
type-preserving maps form the slice category
$\mathsf{Wire} = \mathsf{FinSet}/\mathbb{T}$, which has all finite
colimits, computed as in $\mathsf{FinSet}$ with the induced typing
--- any two wires ever identified by a zigzag of type-preserving
maps have equal types, so quotients are type-coherent automatically.
A \emph{wiring pattern} from $I$ to $J$ is a cospan
$I \to A \leftarrow J$ in $\mathsf{Wire}$; its legs may be
non-injective, which is how a pattern identifies wires --- all
identification of wires lives in the legs (equivalently, in the
quotient forming the apex), never in the tight direction. Patterns
compose by pushout of apexes.
\end{definition}

The paradigm pattern is a plugging: to wire a controller's output
into a plant's input, take the cospan whose apex is the quotient of
the disjoint union identifying exactly those two wires. The legs
stay injective --- the identification lives in the apex --- and
composing two pluggings by pushout chains them through the shared
middle interface.

\begin{definition}[State spaces, joints, and marginalization]
\label{def:su-state-spaces}
For an interface $I$, the \emph{state space} is the product
$\mathcal{X}(I) := \prod_{x \in |I|} \mathrm{El}(\tau_I(x))$ of its
wires' state spaces --- a Polish space, finite discrete when every
type is finite. A type-preserving map $\varphi \colon I \to J$
induces the continuous \emph{restriction}
$\varphi^{*} \colon \mathcal{X}(J) \to \mathcal{X}(I)$,
contravariantly functorial, and $\mathcal{X}(-)$ carries the finite
colimits of $\mathsf{Wire}$ to limits: restriction along a
pattern's pushout legs is a homeomorphism
$\mathcal{X}(A +_J B) \cong \mathcal{X}(A)
\times_{\mathcal{X}(J)} \mathcal{X}(B)$ onto the agreement
subspace. \emph{Joints} over $I$ are elements of
$\Delta_{\le 1}(\mathcal{X}(I))$ --- subprobability laws, read per
\autoref{sec:su-denotation} as probability joints over the
cemetery-completed carriers, so that
partiality composes along the cemetery coordinate; the
\emph{marginalization} induced by $\varphi$ is the pushforward
$\varphi_{\sharp} := (\varphi^{*})_{*} \colon
\Delta_{\le 1}(\mathcal{X}(J)) \to \Delta_{\le 1}(\mathcal{X}(I))$,
functorial and
weakly continuous --- the functor part of the Giry monad
\cite{giry82}, and that sentence is kernel-checked verbatim at this
generality: pushforward of subprobability Borel joints is total,
functorial along measurable maps --- measurability binding exactly
at functoriality, the stratum's constitutive fact ---
mass-exact, weakly continuous along continuous maps, and literally
\texttt{Measure.map} under the coercion
(\texttt{SuFunctor.Plain.MeasJoint}: \texttt{push\_id},
\texttt{push\_push}, \texttt{mass\_push},
\texttt{continuous\_push}, \texttt{push\_toMeasure}).
When $\varphi$ is a coprojection,
$\varphi_{\sharp}$ is marginalization of a joint onto a factor,
which is the sense in which every $\varphi_{\sharp}$ below is read.
\end{definition}

\begin{remark}[Two strata: Polish semantics, exact-rational
certificates]
\label{rem:su-strata}
Stated once, here, for the whole composition tier. The
\emph{semantic} stratum --- every definition of this section ---
lives at the standing generality of
Definition~\ref{def:su-credal-denotation}: Polish record spaces
carrying the weak topology on their subprobability measures (the
vague topology's role, and why the weak one is the topology of
record, are settled there), with the Borel $\sigma$-algebra
underneath. Finiteness, where it genuinely
binds in a theorem, enters as a hypothesis, never as a regime of
the space. The \emph{certificate} stratum is exact-rational
certification over the locally closed $\omega$-semialgebraic
carriers of \autoref{sec:su-site}, and what is finite there is
the \emph{presentation}, not necessarily the space: in dimension
zero distributions are rational mass tables and sublevel checks are
linear programs; in positive dimension the finite object is a
vector of rational polynomial coefficients per piece, countable
piecing is expressed by functorial plates at the hypernet stratum
--- one finite object denoting a whole shape-indexed family,
naturally in the index --- and discharged shape-uniformly at the
certificate stratum through pooled aggregate coordinates, so that a
certificate
discharges
polynomial inequalities over an uncountable carrier without
enumerating it; count-budget rows survive as optional tightenings,
not as the price of admission, and the standing convention below
keeps the non-uniform residue. What connects the strata is the standing
convention: data outside the exact fragment --- non-semialgebraic
carriers, non-polynomial scores (the logarithms constituting
relative-entropy beliefs included) --- ride the wire through outward rational
bounds, which enlarge the credal body, the sound direction.
\end{remark}

\begin{remark}[Hiding a wire: inf-projection, the reflector, and
the tightness certificate]
\label{rem:su-marginalization-lsc}
The composition calculus marginalizes \emph{joints}, never
beliefs (the compositor below); but the assembled stack does
eventually hide wires --- a demand read over visible coordinates
with the hidden ones existentially absorbed --- and the semantic
content of hiding is the inf-projection of a belief along the
pushforward. Its status at the standing generality is fixed by
three kernel-checked facts
(\texttt{SuFunctor.Universe.Marginalization}). First, the boundary
is real: inf-projection does not preserve lower semicontinuity ---
the witness is escaping mass along the hidden coordinate, a closed
family of near-unit point masses whose image under the visible
marginal is not closed
(\texttt{exists\_lowerSemicontinuous\_not\_lowerSemicontinuous\_infProj})
--- so the naive hidden-wire denotation can exit the universe.
Second, the repair is judgment-invisible: hiding taken as the
lower-semicontinuous envelope of the inf-projection changes no
judgment against any citizen of the universe --- the envelope is a
reflector, and domination against a hidden-wire system never
evaluates the infimum at all, resolving upstairs as ordinary
domination through the inf--pullback adjunction
(\texttt{le\_lscEnv\_iff}, \texttt{le\_hide\_iff},
\texttt{le\_infProj\_iff}) --- and the relaxation only lowers the
system side, the sound direction. Third, a certificate turns the
naive formula exact: for beliefs whose sublevels are tight ---
equivalently compact, the Prokhorov reading, mechanized as
consumption
(\texttt{isCompact\_sublevel\_of\_isTightMeasureSet}) --- sublevels
of the inf-projection are \emph{exactly} the pushforward images of
sublevels, the infimum is attained wherever finite, and the
inf-compact class is closed under marginalization
(\texttt{sublevel\_infProj\_eq},
\texttt{exists\_infProj\_attained},
\texttt{isCompact\_sublevel\_infProj},
\texttt{lowerSemicontinuous\_infProj}, with the coordinate
instantiation \texttt{sublevel\_infProj\_margFst\_eq}). The
calculus's generator- and moment-row certificate species are, read at
this stratum, inf-compactness certificates; and one coercive
summand puts a whole additive composite in the class --- a one-line
consequence of nonnegativity, the composite's sublevels sitting
closed inside the summand's compact ones --- so
tightness enters as a certificate on the operand that needs it,
never as an axiom of the universe.
\end{remark}

\paragraph{The mechanized presentation of the base.} The
kernel-checked form of the wiring base presents patterns by their
\emph{state families}: the base double category has as objects
interfaces (bundles of wires over a fixed state type --- the untyped
stratum, the typed slice being orthogonal bookkeeping), as tight
morphisms wire maps acting on states by restriction, as loose arrows
patterns presented as type-valued families
$M \colon \mathcal{X}(A) \to \mathcal{X}(B) \to \mathrm{Type}$ ---
for a cospan, the family of apex states restricting to the given
boundary states --- with apex state space the total
$\mathrm{Tot}\,M = \sum_{x,y} M(x,y)$, loose composition the
$\Sigma$-pullback (exactly the agreement subspace of
Definition~\ref{def:su-state-spaces}), and cells acting on apex
states backward (a map of patterns pushes apex wires forward, hence
pulls apex states back over the boundary restrictions). In this
presentation the wire category's laws are definitional, every
boundary transport vanishes, and the comparison identities of
iterated pullbacks --- the entire calculus a decoration tier
consumes --- hold by reflexivity (kernel-checked:
\texttt{DoubleCat.WireSpan}, a full
\texttt{PseudoDoubleCategory} instance; the associativity-transport identity of the total-state
calculus depends on no axioms at all). Variable identification stays
out of the tight direction by construction:
tight morphisms act on states by restriction and never quotient
anything, and all identification lives in the pattern families.

\subsection{The Decoration Functor}

\begin{definition}[The decoration functor: Specs over the
wiring base]
\label{def:su-house-functor}
Fix a stage type with a window system $U$ (the site dimension,
entering through the decoration per \autoref{sec:su-sheaf}) and
a state type. The \emph{decoration functor} assigns:
\begin{itemize}
\item to each interface $A$, the \emph{fibre of Specs}: persistent
  window-indexed families of classes of
  beliefs on joints over trajectories in $A$'s state
  space, ordered componentwise by inclusion --- such a family is
  precisely what a Spec denotes, and the fibre is named
  accordingly;
\item to each tight morphism (wire map), the componentwise image of
  belief transport along the induced trajectory restriction ---
  covariant on classes, because states restrict contravariantly;
\item to each pattern $M$, the decorations: persistent
  window-indexed families of classes of beliefs on
  joints over apex trajectories ($\mathrm{Traj}(W, \mathrm{Tot}\,
  M)$), with boundary parties supplied by the indexing;
\item to each cell, the componentwise image of belief transport
  along the backward apex-state map;
\item as compositor --- the \emph{belief laxator} $\boxplus$, the
  name later chapters use for this operator --- the componentwise
  elementwise
  \emph{pushforward-then-add}: for decorations $\mathcal{U}$ over
  $M$ and $\mathcal{V}$ over $N$,
  \[
    \mathcal{U} \boxplus \mathcal{V} \;:=\;
    \Bigl(\,\bigl\{\, (\pi_M)_{\sharp} \fatsemi u +
    (\pi_N)_{\sharp} \fatsemi v \;\bigm|\; u \in \mathcal{U}_W,\ v \in
    \mathcal{V}_W \,\bigr\}\,\Bigr)_{W},
  \]
  where $\pi_M, \pi_N$ are the joint-to-component projections of
  the glued apex: the joint over the composite is marginalized to
  each component's trajectory space, each component belief is
  precomposed with its marginal map, and the results are added in
  $[0,\infty]$ --- the beliefs are never themselves marginalized;
  the joint is;
\item as unit, the componentwise singleton of the vacuous belief
  (zero everywhere) on
  identity-pattern trajectories --- the \emph{austere} unit:
  everything permitted, no residue.
\end{itemize}
These data satisfy every obligation of a span-valued lax double
functor on the wiring base --- functoriality of both transports,
naturality of compositor and unit, associativity and both unit
coherences --- and persistence of every result is the site-direction
naturality package (extension commutes with state-side transport,
with the window compositor, and fixes the unit).
\end{definition}

\begin{remark}[Mechanization status of the decoration functor]
\label{rem:su-house-functor-mechanized}
Kernel-checked, end to end, at three strata of the
semantic-universe strand --- two nested, one general. The
single-window (plain) stratum:
carriers and transports in \texttt{SuFunctor.Plain.Penalty}
(\texttt{Joint} --- finitely supported subprobability joints over an
arbitrary carrier, partiality native ---, \texttt{Pen},
\texttt{Pen.pull}), the functor with all obligations in
\texttt{SuFunctor.Plain.SpecFunctor}, its acceptance battery in
\texttt{SuFunctor.Plain.Battery}. The site-indexed stratum, which is
the carrier of the construction's site dimension: window carriers and the extension
calculus in \texttt{SuFunctor.Site.WindowPenalty} --- at both joint
strata: the same calculus over subprobability Borel joints, with its
persistence closure package and the seam theorems, is
\texttt{SuFunctor.Site.GeneralWindow}
(Remark~\ref{rem:su-density-completion}); both are stated over the
function presentation of the trajectory carriers --- graphs before
closure, per \autoref{sec:su-sheaf} --- with the Fell-carrier
restatement the closing ledger's obligation ---, the functor
\texttt{siteSpecFunctor} with all obligations proved componentwise
in \texttt{SuFunctor.Site.SiteFunctor}, the battery in
\texttt{SuFunctor.Site.Battery}, and the matching fork in
\texttt{SuFunctor.Site.MatchingFork}. The componentwise reduction is
itself the structural finding: the seven obligations never see the
site --- each reduces, window by window, to the single-window
coherence core, because trajectory transport is definitionally
functorial and the base's transport calculus is reflexivity-valued
--- so the site dimension's entire force lands in fibre
admissibility (Theorem~\ref{thm:su-matching-fork}), never in the
functor laws. The \emph{general} stratum carries the same
single-window calculus over subprobability Borel joints on
arbitrary measurable carriers: the scalar stock whole ---
pushforward, belief transport, the compositor over arbitrary
marginalization pairs, weak continuity, the lower-semicontinuity
closure package with wall-grade attainment on nonempty compacts,
and the compactly-supported subclass closed under transport
(\texttt{SuFunctor.Plain.MeasJoint}) --- the coherence \emph{cores}
base-free, with the wiring base's comparison identities carried as
commuting-square hypotheses (the finite tier's cores are the
all-squares-by-reflexivity special case); the tight-tier functor
over measurable carriers with the acceptance battery transposed row
by row, refutation pins included
(\texttt{SuFunctor.Plain.GeneralFunctor}: \texttt{measSpecObj}).
The residue: the
double-categorical \emph{total} at the general stratum is not
claimed --- the cores consume comparison squares the exact base
discharges by reflexivity, and a measurable wiring base carrying
its own transport data is the remaining object, a design artefact
in its own right, recorded in the ledger.
\end{remark}

\subsection{The Construction}

The double Grothendieck construction ---
in the covariant, span-valued form Patterson \cite{patterson23} used
to rebuild decorated cospans \cite{fong15}, the strict-tight face of the double
fibrations of Cruttwell, Lambert, Pronk, and Szyld
\cite{cruttwell22b}, whose discrete case is Lambert's
\cite{lambert21a} --- at the dependent (span-of-types) stratum,
assembles the decoration functor into the \emph{double category of
System Specs}: objects are interfaces with a
designated Spec; tight morphisms are wire maps with
fibre transport; loose arrows are patterns with a decoration;
cells are pattern cells with transported decorations; loose
composition pairs the base composition with the compositor. Open
Systems live in this total as embedded ideal points of its fibres
--- the paragraphs below --- never as a separate object class. The
projection to the wiring base strips the decorations, and it is
\emph{strict} --- its compositor and unit comparison are
identities.

\begin{theorem}[The assembled total, and its projection]
\label{thm:su-assembly}
The total of the decoration functor under the double Grothendieck
construction is a pseudo double category, and its projection to the
wiring base is a strict double functor that is a \emph{discrete
double opfibration}: tight morphisms out of a fixed source and cells
out of a fixed loose arrow both lift uniquely.
\end{theorem}

\begin{remark}[Mechanization status of the construction]
\label{rem:su-assembly-mechanized}
Kernel-checked, and the division of labour is
where the proof burden sits: the Grothendieck construction, the pseudo double
category instance on the total (all eighteen axioms), the strict
projection, and the generic discrete-opfibration theorem are proved
once, for an arbitrary span-valued functor on an arbitrary base
(\texttt{DoubleCat.Grothendieck},
\texttt{DoubleCat.Opfibration}; the general split-opfibration
development --- opcartesian lifts, opcleavages compatible with
source and target, split laws as theorems about the canonical
lifts, and the hom-level isomorphism between fiberwise
transformations and strictly projection-compatible functors --- is
kernel-checked at the same stratum, choice-free); the construction then
consumes the generic theorems by instantiation
(\texttt{SuFunctor.Site.SiteSpecTotal},
\texttt{siteSpecTotal\_pseudoDoubleCategory},
\texttt{siteSpecProj\_isPseudo},
\texttt{siteSpecProj\_discreteOpfibration}) --- a one-line
consumption in each case, because the generic lifting calculus
never inspects fibre internals: the window index travels inside the
fibre objects, invisible to it. The variance is forced, not
chosen: beliefs push forward totally along wire maps, but cannot
be pulled back without a lossy infimum --- the cartesian
(negative) half, preimage-transport lifts assembling the
image--preimage adjunction on class fibres, is a named obligation
of the ledger. One generality note: span-of-\emph{categories}
generality is consumed
nowhere in the construction --- the Spec fibres are dependent
families of types, exactly the stratum at which every general
theorem above is kernel-checked, so no gap opens between the stated
and the mechanized generality. The entity-level fibres
(fibre \emph{categories} of Systems and simulation witnesses) are
the refinement layer above this total ---
design, recorded in the ledger, with the general machinery it would
consume already
kernel-checked, and with a prose-verified
construction retained in the workshop record as its design basis;
the paragraph closing this section says why that layer is expected
to be lighter than a second Grothendieck construction.
\end{remark}

\subsection{The Authoritative Total: Presentations as Fibre Data}

The class-level construction above carries arbitrary classes; the
authoritative total of the spec tier sharpens the fibres to the
class stratum's own normal form, carried as data. The
\emph{coupled-DNF decoration functor} assigns to each interface the
coupled-DNF presentations over its state space --- finite families
of spec atoms, each a safety floor with a finite liveness family
containing the properness demand (the option-A discipline: the
vacuous atom is $(0, \{\top\})$, and an empty liveness family is
unrepresentable rather than merely forbidden) --- and to each wiring
pattern the presentations over its apex. Transport is
\emph{atom-level}, disjunct by disjunct --- pull the floor and each
mark, the canonical fibre operation, judgment-invariant along
surjections (Remark~\ref{rem:su-transport-citizenship}). The
compositor is atom-level pairing --- floors compose by
pushforward-then-add, liveness families by pulled union --- and the
unit is the vacuous atom itself: the austere unit and the properness
demand are the \emph{same atom}, which is exactly why the unit laws
close on the nose --- every transport fixes the properness mark, and
every atom's own marks absorb it. The plain double Grothendieck
total of this functor is the development's
\texttt{SingleSiteSemanticUniverse}: objects are interfaces with a designated
presentation, loose arrows are patterns with an apex presentation,
and the strict projection to the wiring base is a discrete double
opfibration by the same one-line consumption of the generic theorem
(kernel-checked, module \texttt{SuFunctor.SingleSiteSemanticUniverse}:
\texttt{dnfSpecFunctor} with the full obligation package,
\texttt{SingleSiteSemanticUniverse}, \texttt{singleSiteUniverseProj},
\texttt{singleSiteUniverseProj\_isPseudo},
\texttt{singleSiteUniverseProj\_discreteOpfibration}). The relation to
the class-level functor is refinement, never replacement, and it is
a theorem: componentwise, on the surjective sub-base, the
image-level transport of a presentation's model class is exactly the
model class of the atom-level transport cut to the transport's range
--- the transport structure theorem of
Remark~\ref{rem:su-transport-citizenship} instantiated at the wiring
base (\texttt{specMap\_modelClass\_eq}), the inclusion half being
the comparison cell in the inclusion-ordered fibres
(\texttt{specMap\_modelClass\_subset}, with its cell-component
companions). Away from the surjective stratum the two genuinely part
company, and the atom-level operation is the one under which the
fibres keep their constructible presentations: at the finite
stratum, through the seam of Remark~\ref{rem:su-seam-mechanized},
every fibre presents a constructible hyperspace class under
semicontinuity of the presented data alone
(\texttt{Universe.SpecDnf.isConstructibleUnion\_probRead} --- the
finiteness that constructibility consumes is carried by the
presentation, by construction). Convexity plays no role at this
tier: admissibility of a fibre is constructibility of its presented
class, and atoms, presentations, and transport are all indifferent
to convex structure. One packaging distinction, stated precisely
because it is easy to misread as a gap in the mathematics: the
window direction of the site dimension is carried, in full, by the
Spec functor above (window-indexed persistent families,
\autoref{sec:su-sheaf}), composition across distinct clocks is
likewise built (the cross-clock paragraph below), and the
cross-site direction enters the base as profile-indexing (the
profiled-base paragraph below); what the
presentation-level total instantiates today is its fibres
\emph{per carrier}, so the re-instantiation of the coupled-DNF
fibres with window-indexed presentation families --- a packaging
step, not a new composition theory --- is an obligation recorded in
the ledger (\autoref{sec:su-receipts}).

\paragraph{Systems in the total: ideal points, and the order as
the fibre carriers' order.} A System enters the class-level total
as an ideal point of its interface's fibre --- the persistent family of
singletons of its window denotations, which is a global member by
descent (\autoref{sec:su-sheaf}) --- and the two structural
facts that make the embedding harmless are both battery theorems:
the embedding is \emph{monoidal} for composition (singleton
decorations compose to the singleton of the composite:
\texttt{singleton\_boxplusW}, \texttt{transport\_singletonW} ---
ideal points compose to ideal points, so composing Systems and then
embedding agrees with embedding and then composing), and
membership agrees with singleton inclusion under the
componentwise inclusion order the fibre carriers wear
(\texttt{wIdealPoint\_le\_iff}, window by window). One boundary
transfers with the claim rather than being smoothed over:
monoidality is a class-level, safety-face fact --- at the coupled
membership relation of the presentation fibres, composing two
members of a paired atom's class need not produce a member,
one component's liveness
dip being destroyed by the other's cost already along identity
transports (\texttt{exists\_coupledSat\_add\_fails}, kept in the
tree as the design boundary); the safety face composes
unconditionally
(\texttt{add\_comp\_safety\_le}), and coherence of the atom pairing
is a statement about presentations, never about membership of
composed models. One precision,
so the categorical claim is not stronger than the machinery: that
inclusion order is \emph{external} structure on the fibre carriers
--- the construction's Grothendieck data are dependent families of
types, and no order is part of them --- so ``the order is the
fibre order'' is shorthand for a statement about the carriers the
fibres happen to be, not about the discrete opfibration. The order
does, however, carry a kernel-checked double category of its own:
the thin \emph{judgment layer} over the same wiring base, whose
loose arrows are decorated patterns and whose cells over identity
boundaries are furnished by inclusions of decorations
(\texttt{DoubleCat.JudgLayer}, kernel-checked at the single-window
stratum; instantiated at this tier in
\texttt{SuFunctor.QtslSem.SpecJudg}: \texttt{SpecJudgLayer},
\texttt{cellOfSubset}, with thinness as
\texttt{cell\_subsingleton}) --- the upstairs plane of
Figure~\ref{fig:su-two-storey}, and the shape in which the judgment
$S \vDash \varphi$ composes as a cell. The order-enriched
refinement of the total itself, in which the designated fibres and
the inclusion cells would live in one structure, remains
an obligation recorded in the ledger. With that precision, the
one-universe thesis
acquires its structural form at the present grade: one fibration
over the wiring base, Specs as fibre objects, Systems as ideal
points inside those fibres, the universe's order read off the fibre
carriers, and the base direction reserved for wiring.

\subsection{The Composition Law}

The law itself ---
pushforward-then-add, lifted to classes --- carries two
theorems, each with its kernel
receipts, and their conjunction is the exact statement of
what composing stochastic wirings means.

\emph{Coherence (the law is lawful).} Composition of decorations is
associative and unital on the nose modulo the base's transport,
with no renormalization anywhere, because no axiom needs one: the
unit is austere and composes transparently
(\texttt{wUnit\_transparent}), the composite class is exactly the
image of the pairs --- nothing is forgotten that a coherence law
could detect (\texttt{mem\_wComp\_iff}; nonempty inputs give
nonempty composites, \texttt{wComp\_nonempty}) --- and the
element-level coherence calculus reduces through the pushforward
calculus to the commutative-monoid laws of $[0,\infty]$. The
battery fixes the design against its named neighbours, as refuting
theorems rather than advice: a \emph{taxed} compositor (any
constant surcharge on composites) breaks unit transparency
(\texttt{taxedBoxplusW\_breaks\_unit\_transparency}); a
\emph{funnel} compositor (forgetting the pair content) breaks the
unit law against the empty class
(\texttt{funnelCompW\_breaks\_unit\_coherence}); and composing
decorations along mismatched middle interfaces is not wrong but
\emph{unrepresentable} --- a type error at the total, pinned as a
build-time refusal.

\emph{Containment, not identity (the law is more committal than the
glue).} Read through the compile isomorphism, the composite
belief's body is contained in the natural-extension glue of the
components' bodies over the shared interface --- always, with no
topological hypotheses --- and the containment is strict off the
crisp locus: in grade coordinates the law is the \emph{product}
triangular norm while the glue is the \emph{minimum}, and
\[
  \pi_1 \fatsemi f_1 + \pi_2 \fatsemi f_2
  \;=\; \max(\pi_1 \fatsemi f_1,\, \pi_2 \fatsemi f_2)
  \;+\; \min(\pi_1 \fatsemi f_1,\, \pi_2 \fatsemi f_2),
\]
so the two sides agree at a joint precisely when the joint is crisp
--- one side at $0$, or one side at $\infty$ (per the identity's
reading in \autoref{sec:su-denotation}: the $\min$-term is not a
cancellative gap at infinite heights) --- and the crisp locus, where
every joint is crisp, is exactly the equality locus. There
deterministic intuition is exactly right that composition is
conjunction: on indicator beliefs both compute intersection of
behaviours. Off the crisp locus a compact
exact-rational counterexample closes the question permanently: two
single-edge dependency-graph models disagreeing about a shared
binary variable (relative-entropy beliefs toward
$(\tfrac23, \tfrac13)$ and $(\tfrac13, \tfrac23)$) admit the
witness $\tfrac{9}{10}\cdot\mathrm{uniform}$ inside the glue but
outside the composite --- the two decisive comparisons being
$81 > 80$ for exclusion and $729 \le 800$ for inclusion --- and
both sides are already compact, convex, and scaling-closed, so no
closure rescues equality. Containment theorem, crisp-locus
characterization, and strict witness are kernel-checked in this
repository at both strata: at the finite stratum (module
\texttt{QtslAux.Imprecise.Containment}; the irrational glue grade
enters through the square dictionary and an exact rational split,
and the two decisive comparisons are hand-checkable arithmetic),
and at the general stratum (module
\texttt{QtslAux.Imprecise.ContainmentGeneral}) --- the containment
over arbitrary carriers, an order-pure argument consuming no
tightness anywhere
(\texttt{sublevel\_compositeG\_subset\_glueSetG},
\texttt{certificate\_transfer}), instantiated at pushforwards of
laws with the glue closed
(\texttt{sublevel\_pmComposite\_subset\_glue},
\texttt{isClosed\_pmGlueSet}), the crisp locus characterized
exactly at the denotation level
(\texttt{forall\_sublevel\_eq\_iff}, through the exact-scale
identity \texttt{ENNRat.add\_eq\_max\_iff}: a sum equals the join
precisely when the lesser term vanishes or the greater is
infinite), and strictness witnessed on a compact carrier, the
finite counterexample above entering as the loaded instance ---
finite discrete carriers are already of the standing generality
(\texttt{binaryComposite\_ne\_glue},
\texttt{sublevel\_ssubset\_glue\_witness}). The direction of the containment is the
sound one --- certificates bound every element of the glue, hence a
fortiori the composite --- and the gap is conservatism, never
unsoundness; by the wall-blindness lemma of
\autoref{sec:su-checker-anchor}, no wall-fragment verdict can
even
observe it, so the conservatism costs only the graded successor,
and there per-slice.

\paragraph{Inconsistency and feasibility.} Because beliefs add,
two wirings can compose to the everywhere-infinite belief --- which
compiles not to the empty set but to the bottom credal body, a
lawful decoration: composing contradictory constraints \emph{should}
denote infeasibility, and it does so first-class, with the
boundary verdicts of \autoref{sec:su-classes}
adjudicating which classes an infeasible composite inhabits ---
every safety containment, no liveness escape, hence no presented
atom's class (Remark~\ref{rem:su-maximal-slice}
and~Remark~\ref{rem:su-presupposition-fork}).
Properness of a
composite has an existential characterization with a constructive
witness: the composite belief is proper iff some pair of
component joints, each of finite inconsistency, agrees over the
shared interface,
and a compatible pair yields a witnessing joint by exact
conditional-product gluing --- a characterization, not a decision
procedure: what the theorem supplies is the shape of the check and
the witness format. The coupling ingredient --- two laws with equal
pushforwards over the shared interface admit a coupling concentrated
on agreement, at standard-Borel strength --- is kernel-checked in
this repository (\texttt{QtslAux.Imprecise.CouplingGlue}); the
packaged equivalence at both strata, with the finite conditional
product as an exact rational formula, is kernel-checked beside it
(\texttt{QtslAux.Imprecise.FeasibilityGlue}). Soundness of the additive law needed no check
at all --- every joint the composite admits has its component
marginals inside the components' certified credal sets, a one-line
theorem (the dropped summands are nonnegative) --- and the kernel
carries the two near-misses that guard it: signed scores break
it, and sub-unit confidence weights break it, each falling exactly
on a field the wire format retains (nonnegative codomain; unit
weights or the rebated conclusion;
\texttt{signed\_relaxation\_fails},
\texttt{subunit\_weight\_fails}). The algebraic lemma is checked
twice over, on the exact stratum and space-free (arbitrary joint
models, finiteness surviving only in the contributor index ---
arity, not state; the two marginal-compatibility modules above).

\begin{remark}[The seam between the two mechanized stocks]
\label{rem:su-seam-mechanized}
The composition tier's exact stratum (finitely supported
subprobability joints over arbitrary carriers) and the judgment
tier's finite stratum (laws as real weight functions over
a finite outcome type) are two presentations of one object, and the
identification is kernel-checked rather than asserted: an
equivalence between the two carriers respecting mass, zero, and the
probability slice, an order isomorphism of the belief stocks by
precomposition, and the judgment bridge --- the fibre-order
membership judgment of the total transported onto the judgment
tier's belief type, where the crescent machinery of
\autoref{sec:su-classes} consumes it
(\texttt{SuFunctor.Plain.Seam}: \texttt{jointEquivBel},
\texttt{jointEquivProb}, \texttt{penEquiv}, and the transported
membership iff). A second seam runs upward, to the general
stratum: the exact joints embed in the subprobability Borel joints
as weighted point masses, respecting mass, zero, and pushforward
--- the restriction theorem: finite pushforward \emph{is} general
pushforward read along the embedding --- injectively over
measurable-singleton carriers, and at the belief and class levels
the finite decoration functor is the restriction of the general
one along the seam (\texttt{SuFunctor.Plain.GeneralSeam}:
\texttt{Joint.toMeas}, \texttt{toMeas\_push},
\texttt{toMeas\_injective}, \texttt{pull\_toMeas},
\texttt{image\_pull\_toMeas}). The three mechanized stocks ---
weight-function laws at the finite stratum, finitely supported
joints at
the exact stratum, subprobability Borel joints at the general
stratum --- thus form a kernel-checked nested ladder, and every
statement this chapter makes across
stocks factors through these seams.
\end{remark}

\subsection{Transport, Clocks, and Profiles}

\begin{remark}[Constructibility under transport: adjudicated]
\label{rem:su-transport-citizenship}
The construction's fibre objects are, in the kernel, persistent
families of \emph{arbitrary} classes of window beliefs; the
stratum of \autoref{sec:su-classes} distinguishes the
constructible classes inside them, and whether the construction's
operations keep the semantics on the constructible stratum is now
adjudicated rather than open, with the answer fixing a design
decision (kernel-checked, module
\texttt{SuFunctor.Universe.TransportCitizenship}). The transport
dictionary is the keystone: the body of a transported belief is
the \emph{preimage} of its body (\texttt{compile\_pullPen}), so at
the body level transport is preimage --- which preserves local
closedness and constructibility outright --- and the coupled
membership relation is
transport-invariant along surjections
(\texttt{coupledSat\_pullPen\_iff}; off the surjective stratum the
invariance fails in both directions --- dips can sit outside the
range, and safety can be violated there). At the
class level the admitted operation is \emph{atom-level} transport
--- pull the atom's data, form its crescent --- and not image
transport of classes: for surjective wire maps the two agree up to
an exact structure theorem (the transported crescent is the pulled
atom's crescent intersected with the saturated classes --- those
equal to the preimage of their image --- both inclusions:
\texttt{transportClass\_crescentClass}), while a
non-surjective counterexample shows raw image transport escaping
the pulled atom's crescent (escape witnesses die off-range:
\texttt{exists\_transport\_not\_subset\_pulled\_crescent}). The
compositor has its own dictionary --- the body of a sum is the
fibrewise radial grade product
(\texttt{compile\_add\_radProd}) --- and union-distribution
receipts close the question against the forcing schema of
\autoref{sec:su-classes}: the atom-level operations
distribute over the constructible unions, transport laxity is
licensed in the sound direction, and no closure operator re-enters
through transport. Class descent preserves constructibility in
both directions independently (\autoref{sec:su-sheaf}). The
packaging is the authoritative total above: the decoration
functor re-instantiated with DNF-presented constructible fibres,
atom-level transport and composition built on exactly these element
cores, and the comparison cell supplied by the structure theorem;
the composite-atom liveness set is carried by the compositor
itself (pulled union, properness mark absorbed). Strata, stated
plainly: everything this remark names is kernel-checked at the
finite stratum (\texttt{SuFunctor.Universe.TransportCitizenship})
\emph{and} at the general carrier
(\texttt{SuFunctor.Universe.General.TransportGeneral}, transports
the continuous maps with Borel codomain): the transport dictionary
rests on mass exactness and scaling linearity of the pushforward
(\texttt{compile\_pullPen}), the coupled judgment is
transport-invariant at surjective \emph{pushforwards}
(\texttt{coupledSat\_pullPen\_iff}, with both off-surjective
failure directions kernel-checked rather than parenthetical), and
the structure theorem splits where its hypotheses genuinely bind
--- the forward
inclusion under surjectivity, the reverse under a named
image-closedness supply (the price compactness paid silently at
the finite stratum), equality under a closed map. Two hypotheses
earn their own mention: bare surjectivity of a wire map does not
yield a surjective pushforward at the general carrier (a
measurable-selection matter; a measurable section suffices,
kernel-checked as \texttt{probPush\_surjective\_of\_section}), and
ideal-point constructibility wants a $T_1$ hyperspace, available
over locally compact ambients. Two residues
remain, recorded in the ledger --- closedness of the saturation locus
(which would upgrade
image transport itself to an admitted class operation on the
surjective
stratum) and constructibility of non-surjective images.
\end{remark}

\paragraph{Composition across clocks.} Heterogeneous-time
composition --- a discrete-time controller against a
continuous-time plant, distinct clocks per wire --- enters the
construction as content of the decoration, not as new machinery; three
facts make the seam tractable, the first two kernel-checked, the
third proven in adversarially reviewed prose in the development's
workshop record. First, a declared clock
morphism (nominal rate with a jitter band, and its hold) enters the
composite's belief content as one \emph{wall on the joint record
space} --- the indicator of the admissible co-realizations, one
shared realization per clock instance --- and spec-to-spec
reasoning across the seam is decided by the same blocking criterion
as within one site, with the sampler's wall joined into the premise
floor as one more conjunct: exact at both the membership and the
class level, kernel-checked
(\texttt{SuFunctor.Universe.DemonicClock}:
\texttt{cross\_site\_union\_refinement\_iff} --- the criterion on
the joint floor --- with the lower-semicontinuity supply that keeps
the adversary class closed, \texttt{lowerSemicontinuous\_selSup};
criterion and supply are carrier-general, and the wall species with
the class-level exactness are kernel-checked at the general carrier
too (\texttt{SuFunctor.Universe.General.DemonicGeneral}:
\texttt{probWall} with its citizenship
characterization, \texttt{crescentClass\_inter\_wall},
\texttt{cross\_site\_union\_refinement\_iff\_probWall}), so the
sentence holds at every stratum; the sampler-relation stock and its
exact composition laws, with
prevision transport and certificate soundness across the seam, in
\texttt{DoubleCat.CrossSite}). Second, transporting behaviour across
a relational (jittered) seam is a genuine functor once it acts on
\emph{bodies} --- all realizations at once, no choice --- with
composition preserved on the nose: the record-tier no-go (a
functional transport would have to choose realizations;
kernel-checked as \texttt{exists\_no\_functional\_transport})
is evaded one storey up, where the Egli--Milner lift of closed
relations composes exactly, the middle witness being an
intersection rather than a selection
(\texttt{SuFunctor.Universe.DemonicClock}: \texttt{emLift\_comp},
axiom-free, following the endofunctor theorem of Abbadini,
Bezhanishvili, and Carai \cite{abbadini23}; what that theorem's
non-allegory caveat actually marks is \emph{meets} --- the lift of
an intersection of relations can be strictly below the
intersection of the lifts, kernel-checked as a one-bit fixture).
Third --- the prose-proven fact --- everything genuinely lossy in
crossing a seam is one
mechanism in three costumes: an envelope or meet taken
\emph{before} a transport forgets shared-realization correlation.
The sound discipline is uniform: conjoin walls and compose
relations first, envelope once, last. The remaining laxity of the
cross-site glue is thereby located, not diffuse, with its species
named: the per-component exchange between coarsening and
marginalization is \emph{exact} for the pushforward of a single law
(naturality, consuming continuity only --- no measurability clause
anywhere in the proof, records transporting by the carrier's
clip-and-push legs); at the credal species the same exchange is lax
in general and strict under a fine-lift (surjectivity)
hypothesis on the transports --- satisfied by the paradigm
coarsenings, where every coarse record is some fine record's
restriction. Under that hypothesis the single inclusion of the
coarsening theorem is entirely the glue-gap --- direct images
under-distribute over meets --- and in every case the certified
direction is sound. On a pipeline that observes the discipline ---
and whose exchange satisfies the fine-lift hypothesis --- the wall's
width is the seam's entire price: relation composition is exact, the
body-tier lift is exact, the exchange is strict, so what a jittered
clock charges is the relaxation of the realization wall itself ---
wider realization classes, weaker walls, more admitted joints ---
which the order then treats like any other imprecision. Off the
discipline, the ledger's envelope and meet gaps charge too; the
discipline is how one avoids paying them.

\paragraph{The site in the base: profiles.} The construction is
\emph{parametric} over state vocabularies and \emph{indexed} over
site-profiles, and the two are one design's two faces.
Parametric: no functor law of the decoration tier inspects the
state vocabulary, so instantiating at a clocked vocabulary ---
realized records carrying their clock --- costs no new proof, and
the fibres of the indexed object compute to the fibres already
built, definitionally
(\texttt{spacetimeSpecFunctor\_fibre\_eq} is a \texttt{rfl}).
Indexed: parametricity alone would leave the site data
ambient-in-the-states --- a total that \emph{contains} multi-site
content while \emph{knowing} nothing of sites --- so the site
assignment enters as object data of a profiled base
(\texttt{SuFunctor.Spacetime.Base}): objects carry a wire-to-clock
profile, tight morphisms preserve it --- along any raw wire map into a
profiled interface there is exactly one source profile making it a
tight (\texttt{profileLiftHom}, \texttt{profile\_unique\_of\_hom}):
interfaces over sites, as a fibration fact --- and states are the
fused states the profile constrains. The indexing is by profiles
rather than by single sites because it has to be: a mixed-boundary
device (\texttt{mixedInterface}) has no object over a single-site
base, and the family-of-one-site-totals reading of ``indexed'' is
refuted under jitter by the relational no-go
(\texttt{exists\_no\_functional\_transport}, the cross-clock
paragraph above). The profile-forgetting embedding
(\texttt{profiledEmbed}) composed with the decoration functor then
yields the iterated double Grothendieck total
\texttt{SemanticUniverse} --- Specs over interfaces over
site-profiles --- whose strict projection is a discrete double
opfibration by one line of the generic theorem
(\texttt{universeProj\_discreteOpfibration}): the lifting
calculus never inspects fibres, so the site grading costs it
nothing. The two faces stay bridged in both directions: each
one-site base sits at the constant profile, with tights, states,
and patterns corresponding by equivalences on the nose
(\texttt{constHomEquiv} and companions --- per-site
conservativity), and cross-site seams enter the \emph{fibres} as
wall conjuncts, the preceding paragraph verbatim, inside the
indexed object. One comparison direction is genuinely lax and
deliberately so: \texttt{profiledEmbed}'s compositor is invertible
(a constrained middle state is a fused middle with a membership
proof, and proofs are irrelevant) but its unit comparison is not
--- the profiled identity embeds as the fused identity pinned to
the profile, a genuine sub-pattern of the fused diagonal --- so
forgetting the site discipline is a lax operation because the
discipline is information. The outer storey is what makes ``which
site is this wire on'' a question the base can answer and the
grammar can enforce.

\subsection{The Monoidal, 2-Functorial, and Certificate Layers}

\paragraph{The monoidal and 2-functorial layers.} The juxtaposition
tensor on the wiring base, and with it the symmetric monoidal
structure of the assembled total, ride a hypothesis package on the
decoration functor --- fiberwise tensors with coherence data, plus
two families of mediating cells against the compositors --- whose
status this book reports with its receipts and its history, in
\autoref{sec:su-receipts}: the package's axioms are mechanized
as named mixins whose satisfaction and refutation are decidable
over cyclic-group graded models, the dependency table among them
machine-checked, and the package's original version omitted an
independent associator--interchange axiom, witnessed by a finite
graded model. The
2-functorial layer --- maps of decoration data inducing maps of
assembled totals, with an isomorphism of hom-categories onto the
strictly projection-compatible functors --- is kernel-checked at
the construction's own stratum (the hom-level isomorphism with explicit
inverse and both roundtrip laws, choice-free), with its
2-categorical superstructure at prose grade. Neither layer is
consumed by anything the earlier sections claim.

\paragraph{Certificates: fibrewise morphisms, not a second tower.}
An earlier design sketched certificates-over-Systems as a second
Grothendieck construction over this one, iterated since
opfibrations compose. The tower collapses: in inclusion-ordered
fibres a certificate is a \emph{fibrewise morphism} --- a witnessed
inclusion of an ideal point, or of a smaller spec class, into a
larger one --- and fibrewise morphisms already compose, vertically
by transitivity, along tight and cell transport by monotonicity,
and horizontally along the compositor by witness-carrying
morphisms, so no second certificate storey is needed: the
certificate layer is the order-enriched refinement of the total
named above (fibre 2-cells as inclusions), not a new fibration.
The collapse is sharper than a tidiness claim, in two
kernel-checked ways (module
\texttt{SuFunctor.Universe.CertificateLayer}). First, certificate
transport needs no side conditions: the transported coupled
judgment factors through the fibrewise infimum (\texttt{fibInf},
empty fibres reading $\top$ --- the right Kan extension along the
carrier map) at \emph{arbitrary} carrier maps
(\texttt{coupledSat\_comap\_iff\_fibInf}), so transported model
classes are preimages and spec-to-spec transport is unconditionally
monotone (\texttt{SpecDnf.comap\_modelClass\_mono}) --- the
opfibration line above, beliefs pushing forward totally while
pulling back only through a lossy infimum, upgraded from variance
remark to mechanism: the lossy infimum reappears exactly once, as
the reflection every transported judgment factors through, and
surjectivity's entire load rests on ideal points, whose off-range
reflection a non-surjective transport rewrites to $\top$
(\texttt{exists\_comp\_notMem\_modelClass\_comap}). Second, the
iterated tower is not merely unnecessary but unbuildable in plain
form: a second-storey certificate functor would owe a \emph{total}
composition laxator, and there are component certificates whose
composite target holds no certificate at all
(\texttt{exists\_coupledSat\_add\_fails}) --- the fibrewise picture
survives precisely because a fibre-hom may fail to exist without
breaking any structure; where the con-witness is supplied the
composite exists (\texttt{coupledSat\_comp\_add\_of\_dips}), and on
the safety fragment composition is witness-free. A companion
witness locates why the fibres are presentations: two presentations
of the same model class whose composites with one fixed atom differ
(\texttt{exists\_boxplus\_not\_semantically\_monotone}) --- the
compositor is monotone in the presentation order
(\texttt{dnfBoxplus\_mono}) but does not descend to semantic
classes, so a composite certificate names the disjuncts it
composes and carries the composite dip family. The
assume--guarantee calculus of Capucci and Myers
\cite{capuccimyers2026} is the worked instance of this collapse,
not of the iterated construction: their Boolean stratum is the
fibrewise-morphism picture verbatim --- a certified lens is a base
loose arrow together with fibre-morphism data over its source ---
and their quantitative construction is itself one double
Grothendieck over ordered fibres of bundle predicates, never
iterated, with slack composition recovered here as the derived
hom-action of the compositor rather than postulated as data. The
boundary: everything in
this paragraph holds per site-fibre; the site dimension --- Specs
over interfaces over site-profiles --- is not collapsed, and keeps
the iterated double Grothendieck of the profiled-base paragraph
above. What remains, recorded in the ledger, is nesting (a
System's entire state as one variable of another), which is the
same move aimed at the fibre. Nothing above depends on this
paragraph.

\section{The Semantic Universe, Assembled}
\label{sec:su-total}

The storeys are all built; this section states the building whole.
The definition below is the chapter rendered as a single
mathematical object: each clause names the storey it consumes and
the section that built it.

\begin{definition}[The Semantic Universe]
\label{def:su-semantic-universe}
Fix a set of sites (clocks) and, for each site, a vocabulary of
realized records carrying that site's stamp
(\autoref{sec:su-causal-order}--\autoref{sec:su-records}). The
\textbf{Semantic Universe} is the iterated double Grothendieck
total
\[
  \mathbb{SU}
  \;=\;
  \int \bigl( E \fatsemi \mathcal{D} \bigr)
  \qquad
  \text{---\ Specs over interfaces over site-profiles ---}
\]
the total, under the covariant span-valued double Grothendieck
construction of \autoref{sec:su-composition}, of the composite
of the profile-forgetting embedding $E$ of the profiled wiring base
into the fused wiring base with the coupled-DNF decoration functor
$\mathcal{D}$. Unpacked, outermost layer first:
\begin{itemize}
\item \emph{The outer indexing: site-profiles.} An object of the
profiled base is an interface --- a bundle of wires --- together
with a \emph{site-profile}: an assignment of a site to each wire.
Underneath every site sits the causal order
(\autoref{sec:su-causal-order}), from which the site's windows and
coverage derive and on which its duration and tangency are declared
(\autoref{sec:su-site}); the profile is how that temporal
geometry reaches the grammar of interfaces. Tight morphisms
preserve profiles, and uniquely so --- along a raw wire map into a
profiled interface there is exactly one source profile making it
profile-preserving --- so the tight category of the profiled base
is \emph{discretely fibred} over site-profiles. This is the
\emph{first} Grothendieck integration: interfaces, integrated over
site-profiles, as a fibration fact. The indexing is by profiles
rather than by single sites (the mixed-boundary and clock-jitter
arguments of \autoref{sec:su-composition}).
\item \emph{The middle base: interfaces and wirings.} Beneath the
profile lies the fused wiring base of
\autoref{sec:su-composition}: interfaces over the fused record
vocabulary, wire maps as tight morphisms, wiring patterns as loose
arrows --- presented by their state families, with apex states
restricting to boundary states --- and pattern cells acting on
apices. The embedding $E$ forgets the profile. It is lax and not
pseudo: its compositor is invertible, but its unit
comparison embeds the profiled identity as the fused identity
\emph{pinned to the profile} --- a genuine sub-pattern of the fused
diagonal (\autoref{sec:su-composition}).
\item \emph{The fibres: presented constructible classes.} Over an
interface, the decoration functor $\mathcal{D}$ assigns the
coupled-DNF presentations of constructible classes of belief
bodies over that interface's states. This clause compresses the
tower: a record is an observation history over a window
(\autoref{sec:su-records}); the record carrier over each window
is a compact Fell hyperspace in which records enter as graphs
(\autoref{sec:su-carrier}); a System denotes a belief ---
a lower semicontinuous functional on laws, its radial body a closed
set of laws (\autoref{sec:su-denotation}); the hyperspace of
bodies makes Systems singleton elements
(\autoref{sec:su-hyperspace}); a Spec denotes a constructible
class of bodies, a finite union of crescents closed under the
connectives (\autoref{sec:su-classes}); and classes are carried
window by window over the site, gluing as persistent families
(\autoref{sec:su-sheaf}). The fibre holds
\emph{presentations}, not bare classes, because composition
consumes presentations: the compositor is monotone in the
presentation order and does not descend to semantic classes
(\autoref{sec:su-composition}).
\item \emph{The second integration.} The double Grothendieck
construction integrates these fibres over the profiled base ---
Specs over interfaces, stacked on interfaces over site-profiles.
The two integrations compose into one total; naming them as they
stack is the definition's phrase ``Specs over interfaces over
site-profiles.''
\end{itemize}
Concretely, in the assembled total:
\begin{itemize}
\item an \emph{object} is a site-profile, an interface over it, and
a designated coupled-DNF presentation of a constructible class over
the profile's fused states --- which site each wire is on, what may
be plugged where, and what is demanded of whatever is plugged;
\item a \emph{tight morphism} is a profile-preserving wire map
together with the transport of the designated presentation --- and
the transport is not a choice: the projection below makes it
unique;
\item a \emph{loose arrow} is a profiled wiring pattern with an
apex presentation --- a way of coupling two boundaries through a
shared body of joint states, demands included;
\item a \emph{cell} is a pattern cell with transported decorations,
and \emph{loose composition} pairs the base composition of patterns
with the compositor of presentations.
\end{itemize}
The strict projection $\pi \colon \mathbb{SU} \to
(\text{profiled base})$ strips the decorations, and it is a
\emph{discrete double opfibration}: tight morphisms out of a fixed
source and cells out of a fixed loose arrow lift uniquely. That is
what the construction buys. Composition of Systems and Specs is
never a design decision at use time --- demands push forward
totally along wirings, with no choices left to make and no side
conditions owed to the site grading --- and every pattern-level
composition law of
\autoref{ch:qtsl-composition-patterns} has one place to
discharge.
\end{definition}

\begin{remark}[Mechanization status of the Semantic Universe]
\label{rem:su-total-mechanized}
\sloppy
Kernel-checked, end to end, by one-line consumption of the generic
Grothendieck theorems of \autoref{sec:su-composition} (the lifting
calculus never inspects fibre internals, so the site grading costs
it nothing): the iterated total is
\texttt{SemanticUniverse} in module
\texttt{SuFunctor.Spacetime.Universe} --- the double Grothendieck
total of \texttt{spacetimeSpecFunctor}, which is the landed
decoration functor composed along the profile-forgetting embedding
(\texttt{profiledEmbed}), with the pseudo double category instance
(\texttt{semanticUniverse\_pseudoDoubleCategory}) and the
strict projection pseudo and a discrete double opfibration
(\texttt{universeProj\_isPseudo},
\texttt{universeProj\_discreteOpfibration}). The profiled
base and its unique profile lifts are
\texttt{SuFunctor.Spacetime.Base} (\texttt{profileLiftHom},
\texttt{profile\_unique\_of\_hom}); the embedding's laxity is by
semantics, its faithfulness on tights a theorem
(\texttt{profiledEmbed\_faithful\_tight}); and the fibres of the
iterated object are the landed fibres \emph{definitionally}
(\texttt{spacetimeSpecFunctor\_fibre\_eq} is a \texttt{rfl}) ---
the construction's parametricity paying inside the indexed object.
Per-site conservativity is packaged and pseudo: each one-site base
sits inside the profiled base at the constant profile with tights,
states, and patterns corresponding on the nose
(\texttt{constProfile}, \texttt{constProfile\_isPseudo}), and
reading a presentation at one site is restriction along that
packing (\texttt{siteRead}). Cross-site seams enter the fibres as
wall conjuncts, exactly as in \autoref{sec:su-composition}
(\texttt{SpecDnf.mem\_modelClass\_withWall}). The fused total ---
the double category of System Specs over the fused base, the
development's \texttt{SingleSiteSemanticUniverse} --- is the middle
integration's total, and the headline object refines it by
composing the site dimension into the base rather than leaving
site data ambient in the states. What remains at design grade is
unchanged from the composition section's own ledger: the monoidal
layer's hypothesis package and the 2-categorical superstructure
(\autoref{sec:su-receipts}); nothing in this definition
consumes either.
\end{remark}

\section{Adequacy: the Big Tent, Squared}
\label{sec:su-adequacy}

A universe that claims to be a big tent owes each resident framework
a receipt: a proof that compiling the framework's models into the
toolkit and denoting them here agrees with interpreting them by the
framework's own textbook semantics and \emph{including} that
semantics into this universe. The receipt has a fixed genre, the
\emph{adequacy square} (Figure~\ref{fig:su-adequacy-square}), built so
that squares are cheap to mint and impossible to hollow. Its
\emph{corners} are the literature's object (restated with the
literature's own well-formedness fields, and tied where possible to
a mechanized probability library's construction of the same
object), the shipped compiled row grain (sparse conventions
included, exactly as the emitters produce it), and this universe at
the per-element stratum. Its \emph{edges} are the shipped compile,
the framework's own law, and the inclusion of precise models into
the imprecise universe. The square's statement is an
\emph{invertible globular cell} in the thin double category of
per-element semantic constraints --- loose arrows are constraints
on subprobability mass functions, loose composition is conjunction,
cells are pointwise entailments --- between denote-after-compile
and include-after-interpret. The one contentful ingredient per
square is a \emph{feeding lemma}: the compiled constraint diagram
pins the framework's law, exactly (finite-horizon
Kolmogorov-extension content, spent once per engine).

\begin{figure}[t]
\centering
\begin{tikzpicture}[x=1mm,y=1mm,
  ar/.style={-{Stealth[length=1.8mm]}, semithick}]
  \node (M) at (0,22)  {framework model $M$};
  \node (L) at (70,22) {its textbook law $\mathcal{L}(M)$};
  \node (C) at (0,0)   {compiled rows};
  \node (U) at (70,0)  {per-element constraints};
  \draw[ar] (M) -- node[above] {\small interpret} (L);
  \draw[ar] (M) -- node[left]  {\small compile} (C);
  \draw[ar] (C) -- node[below] {\small denote} (U);
  \draw[ar] (L) -- node[right] {\small include} (U);
  \node at (35,11) {$\cong$};
\end{tikzpicture}
\caption{The adequacy square, schematically. The left edge is the
shipped compile, exactly as the emitters produce its rows; the top
edge is the framework's own semantics; the right edge is the
inclusion of precise models into this universe; the bottom edge
reads the compiled rows at the per-element stratum. The square's
statement is that the two composites agree --- an invertible
globular cell --- and the feeding lemma is the one contentful
ingredient: the compiled constraints pin the law, exactly.}
\label{fig:su-adequacy-square}
\end{figure}

{\sloppy
Everything else is uniform packaging
(\texttt{SuFunctor.Target}: the constraint double category;
\texttt{SuFunctor.Adequacy.Inclusions}: the inclusion library ---
singleton inclusion for precise frames, point-mass inclusion for
deterministic ones, constraint-list conjunction with its
definitional unit and cons laws, and the invertible-cell packager
that turns a pointwise iff into the square).\par}

Seven squares are kernel-checked at this writing, and each is
stated in full --- receipts included, most at the working scale of
an exam problem with its sketch --- in the section of the modelling
catalogue (\autoref{ch:qtsl-mathematical-models}) that presents its
System-former or, in the delay line's case, the wiring pattern that
replaces one; the composition patterns that ride the squares cite
them in place (\autoref{ch:qtsl-composition-patterns}).
Adequacy lives with the entry it certifies, and what this section
keeps is the genre and the division of labour. The \emph{Markov
chain} square (\texttt{SuFunctor.Adequacy.Dtmc};
\autoref{sec:mt-dtmc}) is the engine: one induction carries all
the stochasticity, the compiled constraints force total mass one
--- a precise frame never draws on the universe's native
partiality, which is the ``inclusion into the precise face'' clause
made a theorem --- and a connector ties the exact-rational law to
the mechanized probability library's own construction. The
\emph{Bayesian network} square (\texttt{SuFunctor.Adequacy.Pgm};
\autoref{sec:mt-bn}) is the same engine transposed from a
time-line to a topological order, which is exactly what it was
minted to demonstrate: nothing in the genre presumes time, because
the genre is \emph{parametric} over its trajectory carrier rather
than indexed over the site --- a square consumes only the carrier's
product structure, never an order on its coordinates --- and where
the site direction genuinely enters is exactly the two named gaps
at this section's close, not a hidden assumption inside any square.
The \emph{Markov decision process} square, fixed-policy face
(\texttt{SuFunctor.Adequacy.Mdp}; \autoref{sec:mt-mdp}), is the
first to reuse an engine wholesale --- the chain square
instantiated at the closed loop, no new induction, every
stochasticity fact already spent --- and the \emph{demonic}
(open-policy) face of the same class
(\texttt{SuFunctor.Adequacy.MdpDemonic}, on the set-valued envelope
vocabulary of \texttt{SuFunctor.Adequacy.EnvelopeUnion}; same
section) is the first genuinely imprecise square: a set-valued
bottom edge, one singleton inclusion per policy, reconciled with
the credal reading through the radially maximal slice, with its
degenerate cases mechanized. The \emph{delay-line}
square (\texttt{SuFunctor.Adequacy.DelayLine};
\autoref{sec:mt-dde}) leaves the stochastic engines --- the
first spatially extended cell, making delay differential equations
a wiring pattern rather than a formalism --- and the
\emph{heat-on-torus} square (\texttt{SuFunctor.Adequacy.HeatTorus},
with the marginal-ring contrast in
\texttt{SuFunctor.Adequacy.MarginalGlue}; \autoref{sec:mt-pde})
is the first where the cover tier genuinely bites, the ring-closing
wall carrying the mesh's holonomy obligation. The seventh square
(\texttt{SuFunctor.Adequacy.HeatTorusInit}; same section) ports the
mesh square's initial face as an exposed boundary trace along the
window's own filter, and, beside the squares, the set stratum of
the same mesh is adjudicated rather than left to intuition
(\texttt{SuFunctor.Adequacy.SolutionSetMesh}; same section), the
first consumer of the record carrier's equalizer-descent form:
closed sets of records glue by exact descent while the point stratum's
holonomy wall goes wall-invisible there, and the missing datum is
branch pairing --- a named successor, not any pointwise constraint.

The squares also feed the universe this chapter assembles: the
wiring stratum
(\texttt{SuFunctor.Adequacy.\allowbreak SemanticUniverseWiring}) ties the
stochastic squares to the coupled membership relation of the
semantic
universe's fibres, per square, with no new induction --- the
squares' feeding lemmas are consumed whole. Systems enter the
fibres through one door, the radial wall of the precise constraint
(\texttt{rwall\_singletonC}): a law's entry \emph{is} its ray, the
body of the precise belief in that law, and the judgment bridge is
ray-shaped --- a law's entry meets a constraint's wall atom exactly
when the law satisfies the constraint
(\texttt{SuFunctor.Universe.RadialUniverse}; the bridge in three
forms, unconditional at full mass). The wiring squares themselves
are kernel-checked on the \emph{slice} presentation of that entry
(\texttt{idealPoint}: the wall vanishing at the law's joint alone,
across the injective scalar seam \texttt{SubPMF.toJoint}), and the
two presentations are verdict-equivalent against every radial
constraint at every mass --- the slice is the ray's radially
maximal slice, the reconciliation the envelope stratum proves and
the radial module sharpens: judged identically, the slice is
nevertheless \emph{not itself a member} of the universe (no belief
compiles to a pinned point, \texttt{compile\_ne\_singleton}), so it
is a computation device of the wiring stratum, never a denotation.
Compiled diagrams become members of the fibre as single-disjunct
DNF presentations (\texttt{diagramDnf}) whose safety face is the
wall of
the conjunction's image and whose liveness set is exactly the
properness demand; and the membership bridge
(\texttt{coupledSat\_idealPoint\_cwall\_iff}) carries all four
stochastic squares through one equivalence
(\texttt{idealPoint\_mem\_diagramDnf\_iff}) ---
membership of the law's entry in the compiled presentation's
model class is equivalent to the law being the classical semantics,
for the chain, the graphical model, the fixed-policy decision
process, and, at the set stratum, the demonic face. Properness is
free at every law's entry, the bottom law included --- the lawful
denotation of infeasibility is proper, and a jointly unsatisfiable
compiled conjunction denotes the bottom body, never an exclusion of
the inconsistent world
(\texttt{contradictory\_composite\_denotes\_bot}); no law enters as
the inconsistent system, and the inconsistent system belongs to no
presented diagram's model class.
Two boundaries: the wiring is
\emph{per-carrier} --- the fibre's record space is supplied
directly, and the construction of the wire object from the compiled
system is a separate, named stratum (the record-space tie below)
--- and the constraint walls are exact-stratum objects; the
topological reading is the envelope stratum's successor, not built
here.

Two disciplines govern every square, present and future.
\emph{Horizon}: the stochastic squares are stated at every finite
horizon, with the
$\sigma$-additive infinite-horizon extension (Ionescu--Tulcea) a
named separate stratum, never smuggled; the differential squares
are stated over open time windows, their record-tier content
pointwise. \emph{Envelope}: precise
frames include by the exact singleton at the constraint stratum,
and the reconciliation with the credal reading --- the ray's
maximal slice is the singleton, verdicts agree kind by kind, with
the zero-law boundary and the down-closure hypothesis stated ---
is the kernel-checked bridge of
Remark~\ref{rem:su-envelope-reconciliation}, composed once per
square rather than re-proven; deterministic frames include their
solution sets at the possibilistic stratum, with per-solution
point-mass faces beside the set-level square. Beyond the seven
kernel-checked squares, the future roster's cost is set by engine
reuse and left-edge availability. Probabilistic logic programs
(finite groundings, exact-rational) and declared-scheduler automata
would ride the existing engines; continuous-time chains would split
into an exact jump-chain face and a separately-stated holding-time
stratum; input boxes and mixture formers would ride the set-valued
envelope stratum the demonic square opened; stochastic-differential
frames wait on their own analytic left edges; and the deterministic
differential squares' finite/certificate and enclosure strata are
un-mechanized successors with cost estimates recorded ---
estimates, not available squares. The two ties that would knit the
squares to the record carrier and the site direction are the named
gaps of the ledger: the \emph{record-space tie} (the squares'
trajectory carriers, supplied directly today, identified with the
record loci of \autoref{sec:su-carrier} over the construction's base)
and \emph{horizon-pasting} (sequential composition of squares along
the site direction, which is the construction's compositor read at the
adequacy stratum).

\section{The Wall Fragment: How Checked Verdicts Embed}
\label{sec:su-checker-anchor}

The universe is now complete on its own terms: beliefs and their
order, bodies and their hyperspace, constructible classes, fibres
over the wiring base. This section is the adapter --- everything
about how the fragment the shipped checkers actually certify
embeds into that universe is collected here, after
the universe and before the receipts, so that no definition above
owes anything to the checker and every checker fact below is an
instance, not an axiom.

\begin{proposition}[The restriction theorem: walls]
\label{prop:su-restriction}
Call a belief a \emph{wall} when it is the indicator of a
closed set $C \subseteq \Delta_{\le 1}(\mathcal{X})$ --- zero
on $C$, infinite off it; convexity is not part of the notion (the
mechanized wall takes arbitrary closed sets), though every wall the
shipped grammar denotes is in fact convex --- half-spaces and their
intersections --- as a theorem, not a requirement. Every judgment kind of the shipped
specification grammar denotes a wall \emph{at its carrier}: the
checker's valuation spaces are countable and discrete, where every
event is clopen and both threshold polarities close, so the
probability atoms $\mathsf{P}_{\ge p}[\varphi]$ and
$\mathsf{P}_{\le p}[\varphi]$ denote closed half-spaces of laws
and sequent validity an intersection of such. (Over a
general record space of the standing structure, event evaluation
$\mu \mapsto \mu(A)$ is not weakly continuous for arbitrary Borel
$A$, and threshold sets close by the Portmanteau polarity
discipline: lower bounds $\mu(A) \ge p$ give closed sets for
\emph{closed} events, upper bounds $\mu(A) \le p$ for \emph{open}
events; the grammar's event classes are chosen to obey it.) On
walls the universe's order is equivalent to the classical verdict:
$\sigma \ge \chi_C$ iff every law admitted at finite
height lies in $C$. The shipped checker's verdicts are comparisons
in the universe's order, restricted to walls; nothing
re-adjudicates under the graded reading.
\end{proposition}

\begin{proposition}[The bridge theorem: tolerance slices]
\label{prop:su-bridge}
For any $\sigma, \varphi \colon \Delta_{\le 1}(\mathcal{X}) \to
[0,\infty]$,
\[
  \sigma \ge \varphi
  \quad\Longleftrightarrow\quad
  \{\sigma \le s\} \subseteq \{\varphi \le s\}
  \ \text{ for every tolerance } s \in [0,\infty).
\]
(Left to right is monotonicity of sublevels. Right to left, at each
$\mu$: if $\sigma(\mu) = \infty$ the inequality at $\mu$ is
automatic; otherwise instantiate $s = \sigma(\mu)$.) A graded
comparison is therefore
equivalently a family of crisp containments indexed by a rational
tolerance --- each slice a comparison of today's wall shape about a
derived object --- and this slice family is the graded order's
\emph{certification chart}: a producer may certify slices
incrementally, while the composed claim remains the single
domination statement.
\end{proposition}

\paragraph{Why the universe grades at all, and what grading costs
the checker.} The motivation is a muteness fact
about the wall fragment: an indicator belief $\chi_C$ and a graded
ramp that vanishes at the centre of $C$, prices deviation within it,
and is infinite off it receive \emph{identical verdicts} from every
wall --- nothing expressible in the crisp fragment
distinguishes a denotation from its indicator collapse, although the
two describe different models. The graded order makes the landscape
itself the subject. The cost: certifying $\sigma \ge \varphi$ \emph{off} the wall
fragment is a named successor obligation for the certificate
calculus, not shipped machinery. For belief data declared by
rational polynomials over the law polytope --- the finite-stratum
carrier $\Delta_{\le 1}(V)$ is a rational polytope --- domination is
a single polynomial-nonnegativity obligation over a semialgebraic
set, the certificate species this book's rule catalogue already
discharges elsewhere, re-aimed; equivalently, by
Proposition~\ref{prop:su-bridge}, per-slice containment rows with
one added rational tolerance field (rational tolerances suffice:
each real sublevel is the intersection of the rational ones above
it). Until that obligation ships,
the chapter's discipline is the claim ledger's: the wall fragment
is shipped; the graded fragment is the semantics, argued and
costed.

\paragraph{Wall-blindness.} The wall fragment cannot see the gap
between composition and conjunction
(\autoref{sec:su-denotation}'s one inequality), a two-line fact: a finite sum in $[0,\infty]$ is finite iff every
summand is, so $\{\sigma_1 + \sigma_2 < \infty\} =
\{\max(\sigma_1,\sigma_2) < \infty\}$; and a wall comparison reads
only the finiteness set ($\sigma \ge \chi_C$ iff
$\{\sigma < \infty\} \subseteq C$). Sum-composition and
join-composition therefore give identical verdicts against every
wall --- every exactness question the graded universe opens, it
opens at the graded successor, and there as quantified conservatism
in the safe direction, never as unsoundness. A companion
invisibility governs convexity: against a convex Spec, a System and
its biconjugate (greatest convex lower semicontinuous minorant)
receive identical verdicts, so system-side non-convexity is
invisible to single-shot linear tests while remaining
composition-visible --- the checker sees the convex reading,
composition sees the territory, which is this book's
untrusted-producer/trusted-checker split reappearing inside the
semantics; a Spec whose content is not convex-expressible is
certified through its own sublevel presentation, per-tolerance
slice rows, rather than through an instrument that cannot see it.

\paragraph{Where certificate rows enter the defining object.} This
book works with the game-theoretic presentation of the glued
global object (\autoref{sec:su-denotation}) because
certificates plug into it \emph{definitionally}: the game-theoretic
upper expectation is an infimum over bounded-below supermartingales
dominating the target functional, and exhibiting a certificate
$V \ge 0$ whose floor rows dominate the indicator of the bad event
and whose generator rows make it a supermartingale against the local
credal data is exhibiting one element of the index set of the
defining infimum --- under one bookkeeping hypothesis that makes
the membership literal: the limit-inferior condition is checked for
the certificate's process stopped at the d\'ebut of the bad region,
and that stopped-process discipline is a per-mechanism-profile
obligation, carried with the profile. With it,
\[
  \overline{P}(\mathrm{bad}) \;\le\; V(x_0), \qquad
  \underline{P}(\mathrm{safe}) \;\ge\; \underline{E}(1) - V(x_0),
\]
a coherent lower-prevision bound, witnessed --- and for full-mass
denotations the familiar
$\underline{P}(\mathrm{safe}) \ge 1 - V(x_0)$; the subtraction is
the leak asymmetry of \autoref{sec:su-denotation}, with its
kernel receipts. The
scheduler-uniformity a reader trained on nondeterministic model
checking will look for is already inside the quantifier order:
certificate rows carry the credal parameters symbolically over
their declared regions, so a checked row holds for every element of
every local credal set at once; fixing any compatible precise
process makes $V$ an ordinary supermartingale; the classical
Ville--Doob inequality gives the per-element bound; and the bound
passes to the envelope last. One quantifier, one direction, no
interchange --- and the checker verifies membership (sign
conditions on polynomial rows), never the infimum itself; no
completeness claim is made anywhere in the stack.

\paragraph{The anchor: where a sequent verdict lives.} Every
certificate in this book ultimately discharges into a checker
verdict, and the verdict object is deliberately small: a sequent
$\Phi \vdash \Psi$ over an ambient hypernet
(\autoref{ch:qtsl-hypernets}), judged semantically valid over
$H$ when every admissible valuation of $H$'s wires satisfying
$\Phi$ satisfies $\Psi$ --- the satisfaction clauses and the
validity statement are stated once, with their mechanized
counterparts quoted beside them, in \autoref{ch:qtsl-manual}.
The answer to \emph{where} the verdict lives in this universe is
exact and short: sequent validity is the wall-fragment restriction
of the universe's order, instantiated at the hypernet's valuation
space. It consumes no site structure --- by design --- and it
consumes composition structure by inheritance; both statements are
made precise below.

\paragraph{The carrier.} Fix a hypernet $H$. A valuation assigns to
each typed wire an exact tensor value of the wire's declared type,
and the valuation space $\mathrm{Val}(H)$ is the product of those
countable discrete spaces --- Polish by the standing structure.
Admissibility and satisfaction each consult finitely many
coordinates, so the three sets a verdict is about --- admissible
valuations, premise satisfiers, conclusion satisfiers --- are
cylinders over finitely many discrete coordinates, hence clopen.
No temporal structure enters anywhere: a valuation is a single
observation --- a record with one stage --- and the satisfaction
clauses consume neither duration nor tangency. That is the site
contract's second freedom (\emph{the checker never sees the site})
made visible in the semantics itself rather than promised at the
format layer.

\paragraph{The face.} Let $B$ be the set of valuations the
hypernet's applications admit and the premise allows --- closed, by
the cylinder observation. Its possibilistic belief
(Remark~\ref{rem:su-faces}) is a wall, and under the order three
readings of the verdict coincide:
\[
  \Phi \vdash \Psi \ \text{semantically valid over}\ H
  \quad\Longleftrightarrow\quad
  B \subseteq \{v \models \Psi\}
  \quad\Longleftrightarrow\quad
  \overline{P}\bigl(\{v \not\models \Psi\}\bigr) = 0 ,
\]
the checker's quantifier, the behavioural superset test, and the
degenerate face's Boolean upper probability --- equivalently, the
order $\sigma_{\mathrm{scene}} \ge \sigma_{\mathrm{sequent}}$
between the two walls, which is
Proposition~\ref{prop:su-restriction} at this carrier. The chain is
kernel-checked at the exact stratum, stated as its two equivalences
and their composite: the superset reading
(\texttt{semanticallyValid\_iff\_subset}), the Boolean upper
probability (\texttt{upperProb\_eq\_zero\_iff}), and the composed
face theorem (\texttt{semanticallyValid\_iff\_face}, module
\texttt{SuFunctor.Plain.CheckerFace}). The corners
agree too, and the agreement is where naive readings crack: when no
admissible valuation satisfies $\Phi$, the sequent is vacuously
valid, the wall is everywhere infinite, the compile is the bottom
credal body --- and the bottom body upper-bounds every event by
zero (Remark~\ref{rem:su-maximal-slice}), so the vacuous
implication of the checker's quantifier is the safety half of the
denotation stratum's bottom verdict, one fact in two dialects ---
itself kernel-checked as a named corollary of the face theorem: the
compile of an unsatisfiable region is the bottom credal object
(\texttt{suppIn\_empty}), and the corner discharges through it
(\texttt{face\_ex\_falso}). The agreement is lawful because
sequent validity is an upper-envelope verdict throughout:
the reading that refuses at bottom --- the liveness reading, with
its existential conjunct --- is not one a sequent ever speaks. What the
verdict deliberately does not exercise: the prevision pair never
separates (nothing weighs valuations), partiality never bites
(there is no mass to leak), and no belief grades (the face is
crisp).

\paragraph{No site structure by design; composition structure by
inheritance.} The absence of site structure from the anchor is
design (the second-freedom reading above), and the
inclusion theorem preserves it: the anchor face factors
through the one-stage stratum. Composition structure the anchor
inherits: hypernets form a decorated-cospan double category, read
through Patterson's double-categorical lens \cite{patterson23} ---
loose arrows are cospans of wire contexts, the legs labelling the
input and output wires, decorated by a wiring diagram on the apex
--- the anchor face is the object part of a lax double functor
from that double category into the double category of System Specs
(an open hypernet goes to its possibilistic coupling, the wall of
its admissible-valuation set), packaged in the toolkit's Lean tree
as \texttt{hnetInclusion} (module \texttt{SuFunctor.HypernetDouble}):
opcartesian lifts preserved on the nose, the compositor invertible
exactly on the seam-disjoint stratum that content addressing
guarantees for encoded hypernets, the unit comparison genuinely lax
by conformance-pinning. The substitution laws of
\autoref{ch:qtsl-hypernets}
(Theorem~\ref{thm:substitution}) are exactly its functoriality ---
with the determinism hypothesis of law \emph{(ii)} deleted along
the way, since the possibilistic face composes for the full
calculus, \texttt{cell} included.

\paragraph{What sits above the anchor.} The temporal and
probabilistic instruments of this chapter are exercised by the
\emph{rules}, not by the verdict. A generator row's meaning as
supermartingale membership lives in the rule lemma that emitted it
(the certificate paragraph above); the reading of an accepted
sequent as a graded-order bound on a System's denotation is the
graded fragment, the named successor of
Proposition~\ref{prop:su-bridge}. The factorization is the
architecture's own: the checker verifies membership ---
satisfaction over valuations, sign conditions on rows --- and the
deep semantic content lives in soundness lemmas quantified over
every hypernet and every admissible valuation, one lemma per rule.

\begin{remark}[Mechanization status of the wall fragment and the
anchor]
\label{rem:su-checker-anchor-mechanized}
Kernel-checked, over the general carrier of the semantic-universe
strand --- the statements are order-pure, and the mechanization
consumes nothing of the record space
(\texttt{SuFunctor.Universe.General.CheckerAnchor}): the wall
restriction of the order
(\texttt{sat\_wall\_iff}, \texttt{sat\_wall\_wall\_iff} ---
Proposition~\ref{prop:su-restriction}) and the
tolerance-slice bridge (\texttt{sat\_iff\_sublevel\_subset} ---
Proposition~\ref{prop:su-bridge}). Kernel-checked, for the
inclusion's combinatorial core at the aligned stratum --- the
pushout-normal form where seam wires are literally shared ---
against the calculus's own hypernet definitions (module
\texttt{Qtsl.HypernetCospan}): loose composition of well-formed
open hypernets with every boundary clause proved
(\texttt{OpenHypernet}, \texttt{glue}), the admissibility split of
a glued net --- a valuation of the composite is admissible iff both
parts' applications hold, determinism nowhere
(\texttt{glue\_admissible\_iff}) --- and the unit and associativity
laws on the nose (\texttt{unitOpen}, \texttt{glue\_assoc}); the
full-generality rename transport is kernel-checked in the same tree
alongside the inclusion's three other transformation laws
(\texttt{theta\_rename}, module
\texttt{SuFunctor.HypernetInclusion}), and the
double-category-library instantiation of the decoration functor is
the packaged inclusion itself (\texttt{hnetInclusion}, module
\texttt{SuFunctor.HypernetDouble}). Kernel-checked, in this
repository's certificate-calculus soundness
development (the \texttt{Qtsl} package under \texttt{lean/qtsl},
whose per-rule statements \autoref{ch:qtsl-manual} quotes
beside each rule): the valuation, satisfaction, and validity
definitions the anchor instantiates, and the theorem that every
certificate object of the calculus witnesses its sequent's semantic
validity --- total over the rule catalogue, by structural recursion
with one lemma per rule (\texttt{ProofObject.sound}, with its
checker-level face \texttt{check\_sound}, module
\texttt{Qtsl.Soundness}).
The HOL4 development behind the certified checker (this repository's
\texttt{hol} tree) takes hypernet-level soundness as
its trusted-statement target, and its composition toward the
end-to-end theorem is in progress: demonstrated route, assembling
now, and this section will say ``machine-checked end to end''
about that side only when it is.
\end{remark}

\section{Coherence Receipts}
\label{sec:su-receipts}

A chapter that promises ``proven mathematically coherent'' owes the
reader its evidence classes, claim by claim. This book's organizing
principle --- powerful untrusted producers, small trustworthy
checkers --- applies to its own metatheory, and this section is the
application. Three evidence classes, fixed here and used in every
mechanization remark of the chapter: \emph{kernel-checked} means
the named statements are theorems of this repository's
proof-artefact development;
\emph{prose-proven} means proven in adversarially reviewed prose in
the development's workshop record, not mechanized; \emph{open
obligations} are design with named obligations, stated in the
future tense. In the claim-ledger vocabulary used across this
book's public surfaces, kernel-checked and prose-proven claims are
[shipped] at their stated strata, open obligations are [argued],
and this chapter contains no [aspirational] claims. Every
kernel-checked claim lives in the toolkit's
proof-artefact repository, sorry-free at this writing, with the
axiom census re-emitted at every build (each named theorem depends
on at most propositional extensionality, choice, and quotient
soundness; the forcing iff of \autoref{sec:su-classes} and the
relational-transport core of \autoref{sec:su-composition}
depend on no axioms at all). Two independent instruments keep the
class checkable beyond the producing kernel: the repository's
\texttt{lean-cone}
tool re-derives, for any named constant, a statements-only audit
surface --- the transitive signature cone, topologically sorted,
elaborated against the real development --- so the correspondence
between this chapter's prose and the tree's statements is itself
checkable without trusting either; and the compiled proof objects
are re-checked by proof-checking kernels independent of the one
that produced them (SafeVerify, and nanoda), so ``kernel-checked''
receipts are receipts from multiple kernels, not one.

\paragraph{Kernel-checked.} By strand, one bullet per artefact
layer, each named beside its statement in the body of the chapter:
\begin{itemize}
\item the double-categorical ambient: pseudo double categories with
  dependent cells and a cast calculus; lax double functors with the
  full coherence package; transformations and modifications
  assembling the functor double category; the loose bicategory
  bridge; companions, conjoints, and equipments; the span, matrix,
  and dependent-span models; the double Grothendieck construction
  at the dependent stratum, complete --- total, strict projection,
  the split double opfibration development with canonical lifts and
  split laws as theorems, the discrete-opfibration interface, and
  the hom-level isomorphism between fiberwise transformations and
  strictly projection-compatible functors, choice-free; and the
  monoidal hypothesis package as decidable mixins over cyclic-group
  graded models, dependency table machine-checked;
\item the wiring base in state-family presentation, with its
  reflexivity-valued transport calculus;
\item the decoration tier, plain, site-indexed, and
  presentation-level: carriers,
  belief transport, the spec-class functor with all obligations at
  both strata, the coupled-DNF functor with atom-level transport,
  the atom-pairing compositor, and the option-A unit absorption,
  the assembled totals --- the site-indexed class total and the
  authoritative \texttt{SingleSiteSemanticUniverse} --- with strict
  projections and
  discrete double opfibrations by the generic theorems, the
  image-level comparison cells with their structure theorem, the
  constructible-fibre reading across the seam, the
  acceptance batteries (unit transparency, no funnel, associativity
  modulo base transport, ideal points monoidal, judgment as fibre
  order), and the refuting fixtures (taxed and funnel compositors,
  the unrepresentability pin, the matching fork);
\item the record carrier: the Fell topology with hypothesis-free
  compactness; clip and proper-pushforward continuity with the
  seam fixture; descent as a homeomorphism onto the projective
  limit, with the nondirected and Vietoris sharpness fixtures; the
  record locus (single-timeline stratum) as a closed subsheaf with
  the speed-row and
  directedness fixtures and the Lipschitz-arc reading; the
  one-point compactification correspondence; the capacity
  presentation (basis,
  second countability, measure determination, separation, the
  frame form, and the two boundary fixtures); and class descent
  with constructibility in both directions;
\item the universe and hyperspace storeys: carriers, membership
  closure,
  generation by rational rows, by tents, and by gamble rows at the
  general carrier, the pointwise order
  with its wall restriction and the tolerance-slice bridge,
  the order-level connective faces (conjunction's universal
  property; the functional-level minimum with its universal
  property and hedging witness), the coupled membership machinery
  with the blocking criterion (sound and complete, ex falso derived)
  and its incompleteness fixtures, the crescent and
  constructible-class calculus with the compile characterizations
  (the bijection with the credal bodies included),
  and the ideal-point/definability package;
\item the transport stratum, at both the finite and the general
  carrier: the transport dictionary
  (body of a transported belief as body preimage), invariance of
  the coupled membership relation along surjections, the surjective
  structure theorem with its saturation
  condition, the non-surjective counterexample, the compositor's
  radial grade-product dictionary, and the union-distribution
  receipts;
\item the disjunction stratum: the forcing iff (axiom-free), the
  countermodel schema, exactness and constructibility of the
  join, the disjunction property, the
  locally-closed-insufficiency witness, and the hedger fixtures
  separating the functional and class readings;
\item the imprecise-probability strand of this repository: the
  belief--credal compile isomorphism (exact-rational, exp-free,
  with the concave-belief witness and the composition grade
  dictionary); the finite envelope-coherence package (attainment,
  super- and subadditivity, homogeneity, monotonicity); the
  additive law's domination and sublevel lemmas with their two
  fencing counterexamples, at the finite-rational stratum and
  space-free; the standard-Borel coupling-on-agreement theorem; the
  composition and tower laws of credal envelopes; the Hurwicz
  machinery with its composite evaluations; the belief-kernel
  tier at the exact stratum --- tilted envelopes over signed costs
  with the credal zero-cost face and the one-member precise face,
  the positive-homogeneity failure witness, the finite Legendre
  transform of threshold profiles, the exact intertwining square
  with its two cost-law countermodels and the once-drawn
  diagonal's strict gap, the mixed-contraction identity with the
  strict exchange witness, wing collapse at prevision
  equivalence, and the Hurwicz chord decomposition with the
  certificate-pair rule and the finite-chords strictness
  witness; the coherent
  lower-prevision package with the finite envelope theorem and the
  natural extension's universal property, LP duality, and sure-loss
  alternative, with its four-mutant battery; the composition law's
  containment in the least-committal glue, with the crisp-locus iff
  and the strict rational witness; the feasibility (properness)
  equivalence at both strata with the exact conditional-product
  witness and its agreement-clause mutant; the leak-aware
  safe-side bound for submass models, envelope form and precise
  recovery, with the total-leak counterexample; and the credal-body
  poset of Remark~\ref{rem:su-body-atlas} --- placements, strict
  edges, and incomparability witnesses across the atlas's classes,
  the convexity adjudication trio (sufficiency, non-necessity, the
  boundary-jump closedness failure), and the
  prevision--ultracontribution identification;
\item the generalization-column ledger items: the causal-diamond
  interior coverage (thickness iff and pair-filteredness at
  bicontinuous-poset generality), clock-map functoriality and
  cover-preservation with the homogenization core, and the
  open-quantum bounded-generator transport theorem with GKSL
  trace-annihilation and the Heisenberg trace adjunction;
\item the quantum strand of the causal-diamond column, at the
  finite-dimensional stratum throughout: the
  records-imply-consistency theorem at hypotheses weaker than
  projections, its additivity bookkeeping, the
  load-bearing-hypothesis witnesses, the pure-state converse via
  branch projectors, a Minkowski instantiation, and the
  weakly-consistent-without-records witness; state-body gluing
  failure on the spacelike pair-cover with entanglement
  load-bearing (rigidity at pure marginals as the contrast),
  refusal of the failing cover by the derived topology, the general
  repair theorem (interior attainment implies separation), and the
  gap located at the linear rather than the multiplicative span;
  transport of diamonds, interiority, coverage, nets, separation,
  and determinacy along chronology-respecting order isomorphisms,
  with cover-preservation in the site-category vocabulary; the
  state presheaf with its separation and non-separation witnesses,
  the record shadow into probability rows with its strict credal
  band, and the unique-density-state representation theorem; Cauchy
  determinacy at both grades from complete initial data plus a
  reversible mechanism, its reversibility and exhaustiveness
  witnesses, and the frozen-shadow tie; the order-superposition
  probes (no records on a coherently switched index; pointer
  orthogonality classicalizes it); and the assembled locally
  covariant kinematic fragment with its two-qubit inhabitant;
\item the envelope reconciliation (ray versus singleton, with its
  boundary and refutation fixtures) and the seam isomorphism
  between the composition tier's and the judgment tier's carriers;
\item the cross-site stratum: the wall-conjunct criterion
  instances, the relational no-go kernel, the Egli--Milner lift
  with exact composition and the meets fixture, and the transport
  laws;
\item the site dimension in the base: the profiled base with its
  discrete fibredness over profiles, the profile-forgetting
  embedding (compositor invertible, unit comparison deliberately
  lax), the iterated total --- Specs over interfaces over
  site-profiles --- with its strict projection a discrete double
  opfibration, definitional fibre computation at the clocked
  instantiation, and the constant-profile equivalences of per-site
  conservativity;
\item the temporal-type-theory inclusion at the classical
  instantiation: the interval base site with the
  translation-quotient functor from the spec tier's windows
  (discrete-fibration lifts in both directions, covers
  corresponding to covers), the trajectory identification, the
  indicator calculus with behaviour-subtype spec objects, the
  judgment bridge (satisfaction as behaviour membership), strict
  tight intertwining, the compositor on indicators as span
  composition, and the semantically silent class-level laxity;
\item the certificate layer: the fibrewise-infimum transport
  dictionary at arbitrary carrier maps (safety adjunction, liveness
  transfer, the keystone factorization, transported model classes
  as preimages), unconditional transport monotonicity at both the
  semantic and the presentation order, the surjective dictionary
  re-derived as a one-line corollary, witness-gated composite
  certificates on the compositor, and the two fences (non-surjective
  ideal-point transport; presentation-sensitivity of the
  compositor);
\item the adequacy strand: the constraint double category, the
  inclusion library, the chain, Bayesian-network, and fixed-policy
  decision squares with mass theorems and library connectors, the
  demonic decision square with its set-valued bottom edge and
  the union-of-rays slice stratum, the delay-line square (delay
  differential equations as a wiring pattern, two-sided by the
  method of characteristics), the heat-on-torus square (the
  cover tier's cyclic mesh, with the holonomy wall and the
  frustrated-ring contrast) with its initial-face port (trace atoms,
  the conditional feed discharge, derived wrap-compatibility), the
  solution-set mesh (exact equalizer descent consumed at the set
  stratum, the wall-invisible $\mathbb{Z}$-orbit holonomy witness),
  and the universe wiring (ideal points
  across the scalar seam, presented diagrams, the judgment bridge
  feeding the stochastic squares through the coupled judgment, and
  the singleton-versus-ray refutation);
\item the checker anchor's substrate: the certificate calculus's
  valuation, satisfaction, and validity definitions with the
  per-rule soundness development, total over the catalogue
  (\texttt{ProofObject.sound} and \texttt{check\_sound}, module
  \texttt{Qtsl.Soundness} of the \texttt{Qtsl} package), and the
  hypernet inclusion's combinatorial core at the aligned stratum
  (open hypernets over the calculus's own definitions, loose
  composition with well-formedness closure, the admissibility
  split, units and associativity on the nose, module
  \texttt{Qtsl.HypernetCospan});
\item the pointfree strand and the general stratum, landed as one
  campaign: valuations on frames with the promotion homeomorphism
  and the polarity package, the Borel bridge with
  $\tau$-additivity as its exact hypothesis class, bridge
  continuity into the pointfree carrier, and the free-ultrafilter
  escape witness; the classical subprobability carrier with the
  measure-level dagger chart, compactness by two independent
  routes, and metrizability; the general compile dictionary ---
  closedness iff lower semicontinuity, the image characterization,
  unconditional reconstruction --- with the general-carrier atlas
  placements, the grade-convexity criterion, the named
  order-closedness hypothesis of the finite-generation theorem, the
  Radon-functional tie at locally compact second-countable Borel
  carriers, and the general wall restriction and tolerance-slice
  bridge; subprobability Borel joints with the Giry-functor
  package, base-free coherence cores, the tight-tier general
  functor with its transposed battery and refutation pins, and the
  upward seam from the exact joints; the marginalization triple
  (escaping-mass witness, judgment-invisible reflector,
  tightness-certificate exactness) and the hull-closure package
  (closed binary hulls, levelwise coproduct exactness); the
  general-stratum composition containment with its crisp-locus
  characterization; the completed state-space category (all
  countable coproducts, disjoint and universal, with the bag space
  as the coproduct of its sorted-cone sectors); and the campaign's
  continuation at the judgment and class tiers: the general
  connective faces, closure, and gamble-row generation (directed
  unconditionally; countable under second countability, discharged
  for compact metrizable record spaces), the general crescent and
  constructible stratum, the general transport dictionary with its
  named surjectivity and image-closedness supplies, the general
  demonic wall species with the cross-site criterion's class
  exactness, and the general natural-extension package (envelope
  coherence over arbitrary nonempty families, the universal
  property at the definition's stated scope with the
  dominating-family hypothesis self-certifying avoidance of sure
  loss, and closedness of the compatible set under continuous
  assessment gambles).
\end{itemize}

\paragraph{Incorporated from the workshop.} A named minority of this
chapter's mechanized claims were, until their landing,
\emph{workshop-checked}: sorry-free kernel theorems whose modules
lived in the development's workshop tree pending incorporation,
graded for a reader of this repository alone as adversarially
reviewed prose with a mechanization attested. That set --- the
natural-extension package (universal property, linear-programming
duality, sure-loss alternative, finite envelope theorem); the
leak-aware safe-side package with its total-leak fixture; the
containment theorem's kernel form (crisp locus and strict rational
witness); the packaged feasibility equivalence at both strata; and
the site-column ledgers' discharged items (interior-coverage
pair-filteredness and the thickness equivalence at
bicontinuous-poset generality, clock-map functoriality and
cover-preservation, the homogenization core of time-inhomogeneity,
and the open-quantum hypotheses' transport theorem at the
finite-dimensional stratum) --- has been incorporated whole, as
modules \texttt{QtslAux.Imprecise.NaturalExtension},
\texttt{QtslAux.Imprecise.Containment},
\texttt{QtslAux.Imprecise.\allowbreak FeasibilityGlue},
\texttt{QtslAux.LeakAware}, \texttt{QtslAux.DiamondSite}, and
\texttt{QtslAux.\allowbreak OpenQuantum}, with axiom receipts emitted by the
same build target as the rest of the package. As predicted at the
staging, no claim of this chapter changed shape in the landing ---
only its evidence class; the workshop-checked class stands empty at
this writing.

\paragraph{Why the receipts are the doctrine: the phantom axiom.}
The monoidal package's original axioms do not imply the
associator--interchange compatibility square: in graded models over
an abelian group every coherence axiom linearizes, derivability
becomes exact linear algebra, and the decision procedure exhibited a
package satisfying every stated axiom, mediating cells invertible,
whose associator--interchange square fails at explicit small
integers. The repair is additive and
costs every concrete instance nothing; the audited dependency table
is itself the exhibit --- the new axiom is retained with its
independence proven in both directions by machine witness, and a
third candidate axiom turned out to be \emph{derivable}: the audit
removed an axiom thought necessary in the same pass that added the
one that was missing. The lesson is drawn narrowly: in coherence axiomatics, a
machine-decidable dependency table is the reliable form of the
claim. That is why this chapter's coherence claims cite kernel
receipts, and why the disjunction stratum's own repair went
through the same door: the failure was
found by countermodel, generalized into a kernel-checked schema,
and only then written into prose.

\paragraph{Prose-proven, adversarially reviewed.} The localic
reading of the fibre column (the powerlocale and valuation-locale
identifications, geometricity of the layer-by-layer presentation)
--- though no longer with the mechanized tier confined to its
point-set shadow: the capacity presentation's determination theorem
and the valuation layer beneath it are kernel-checked at the localic
stratum (Remark~\ref{rem:su-capacity-mechanized},
Remark~\ref{rem:su-fibre-column});
the site-column mathematics not yet in the kernel (the
causal-diamond column's manifold-level instantiation of the
clock-map hypotheses; the open-quantum column's design hypotheses);
the
2-categorical superstructure of the change-of-base layer; the
span-of-categories generality of the Grothendieck construction
(consumed by nothing in the shipped construction); and the classical
agreement theorems' prose legs at their stated finite-state,
discrete-time scope; the collapse theorem for the certificate
storey (fixed site, ordered fibres, the comma reading of the
iterated construction), whose instantiation at the landed fibres is
the kernel-checked certificate layer above; and the hypernet
inclusion's three theorems at full generality --- the
decorated-cospan double category of raw wiring diagrams, the
well-formed genus as a sub-double category closed under
composition, and the strict double-functorial inclusion via the
possibilistic coupling --- with their combinatorial core
kernel-checked at the aligned stratum. An entity-level decoration
functor --- windowed System fibres, coupling categories, five
prose-verified axioms --- remains in the workshop record at this
grade as the design basis of the certificate layer's refinement
stratum; the construction's definition is the kernel-checked
functor pair of \autoref{sec:su-composition}, and no claim
of this chapter leans on the entity-level construction.

\paragraph{Open obligations, with names.} The remaining gaps, so
the reader
need not hunt for them --- and there are several, not one. The
graded order's checker obligation: certifying belief domination
off the wall fragment (polynomial nonnegativity over the law
polytope, or equivalently slice-containment rows with a tolerance
field), the [argued] line of Proposition~\ref{prop:su-bridge},
with the wall fragment shipped. The site-indexed re-instantiation
of the presentation fibres (the packaging step of
\autoref{sec:su-composition}), and the two transport residues
of
Remark~\ref{rem:su-transport-citizenship}: closedness of the
saturation locus, and non-surjective image
constructibility. The cartesian half of the projection
(preimage-transport lifts; the image--preimage adjunction on class
fibres) and the two-sided form with the site as a visible index.
The record-space tie and horizon-pasting for the adequacy squares
(\autoref{sec:su-adequacy}). The Fell-carrier form of the window
assembly (\autoref{sec:su-sheaf}): the extension calculus, its
presheaf laws, both exchange theorems, and the persistence closure
package --- kernel-checked today over the function presentation of
the trajectory carriers, graphs before closure under the product
$\sigma$-algebra --- restated over the record carrier itself, with
Borel-on-Fell as the measurability datum and the carrier's
Fell-continuous clip as the restriction (the embedding being
graph-then-closure, so that the presented calculus becomes the
restated one's image along the presentation); the clip's own
measurability is already a corollary of its landed continuity. The
pair-sequent packaging of
the blocking criterion. The $\sigma$-completion of the class fibre
(\autoref{sec:su-classes}): attitude superlevels and threshold
families as unions of crescents over declared index families, with
transport carrying semicontinuity certificates for both envelope
passes. The body-level
representation of non-radial members (the classes quantify
over bodies, which are radial objects; the graded universe's
two-axis members await their own body construction there). The
hybrid (multivalued-instant) event strata --- the shadow of the
compact parameter class closed as a continuous image, and the
alternation-bounded classes closed by hit-face arithmetic ---
plugging into the kernel-checked generic local-locus interface;
should a clause family ever certify same-instant sequencing, the
lens adjoins an event face by the same con-coupling that adjoined
the measure face, and nothing below re-plumbs. The record locus over general causal
supports (Definition~\ref{def:su-record-window}) --- where the
single-timeline slicing argument does not bundle: over
chronologically thick supports, two-point specimens refute
closedness of the full record class (a spacelike stage-pair
Fell-converging onto a null pair awakens a zero-duration constraint
no approximant felt; a merging spacelike pair defeats
single-valuedness), so what is owed is the factored replacement.
Closedness of the interior-pair locus --- speed bounds over pairs in
the interior of comparability, which push-up identifies with
chronology --- by the same violating-pair margin argument, consuming
two site hypotheses: chronology open in the stage topology, and the
two-stage duration upper semicontinuous on it (on a globally
hyperbolic instantiation: openness of $I^{+}$, with the interval
topology the manifold topology, and continuity of the Lorentzian
distance). The plug of that locus into the kernel-checked
local-locus interface, whose directedness step is unchanged.
Closedness of the single-timeline-support sub-locus, where chain
supports are Fell-stable because causality is closed. And the
zero-duration-pair and single-valuedness conditions as certified
$G_\delta$ strata inside the closed locus, carried in the manner of
the Zeno stratum rather than by the topology. The outer storey of the background-independent column
(\autoref{sec:su-site}): morphisms beyond isomorphisms and
coverage for order-indexed sites --- transport along
chronology-respecting order-isomorphisms is kernel-checked above
--- together with the quantum column's named residues: relative
Cauchy evolution and the stress-energy tensor, where the kinematic
fragment fences off dynamics proper; non-surjective embeddings
between different spacetimes, with functoriality of the assembled
fragment along them; and the reading of the fragment's record
shadow into the belief storey at a constant site-profile.
The prevision-sheaf face of the
fibre
column and the fully typed display of the monoidal package's
batched discharge rows, with the mechanical comparison identifying
the Grothendieck-route coherences with the classical
decorated-cospan components --- and, for the hypernet inclusion of
\autoref{sec:su-checker-anchor}, the full-generality rename
transport and the double-category-library instantiation of its
decoration functor. The logic--site contract's
remaining
transports: the analytic half of the extension-theorem audit
(objectwise-classical, per window class), the probabilistic
time-inhomogeneity lemma's projective-limit step, and the
open-quantum event/record transport theorem. The HOL4 certified
checker's end-to-end composition over the anchor's validity
statement: demonstrated route, assembling now. And nesting --- a
System's entire state as one variable of another, named in
\autoref{sec:su-composition} --- together with the
library-grade generic form of the certificate storey's collapse
(the ordered-fibre double Grothendieck with its comma
isomorphism), whose instantiation at the landed fibres is
kernel-checked above. Each entry
above is stated in the future tense, deliberately, and nothing
downstream in this book consumes any of them ahead of their
receipts.
 \chapter{Hypernets: Primitive Functions}
\label{ch:qtsl-hypernets}

\newcommand{\primentry}[2]{\par\bigskip
  \phantomsection\label{prim:#2}\noindent{\bfseries #1}\hfill
  {\footnotesize\textcolor{NoteGray}{Primitive Catalogue entry \texttt{#2}}}\par\nopagebreak\smallskip}
\newcommand{\primmeta}{\par\nopagebreak{\small\bfseries The row (Computer Science Metanotation)}\par\nopagebreak\smallskip}

\noindent\draftmark\medskip

Three strata of this book each need, at some point, a notion of
\emph{computation}. A System (chapter~\ref{ch:qtsl-mathematical-models})
declares mechanisms whose fields --- drift laws, guard conditions,
reset maps, neural controllers --- are computational objects. A
specification (chapter~\ref{ch:qtsl-semantic-universe}) asks whether
observables computed from a trajectory stay inside a region, and a
certificate about it carries polynomial content of its own. The rule
catalogue (chapter~\ref{ch:qtsl-manual}) states every inference rule
over ``the ambient hypernet'' --- a typed computation graph that each
rule reads facts out of. If those three strata each invented their own
computational syntax, the toolkit would own three evaluators, three
serialization schemes, and three subtly different answers to ``is the
function the certificate reasons about the function the system runs?''

It owns one. The shared object is the \emph{hypernet}: a typed,
exactly representable computation graph built from ten primitive
boxes, with one evaluator, one canonical serialization, and one
content digest. Everything in the toolkit that computes --- an ODE
right-hand side, a ReLU network, a guard observable, the goal
polynomial of a Lyapunov certificate --- either \emph{is} a hypernet
or elaborates into one by a definitional translation. This chapter
defines the object once, so that the rest of the book can consume it by
reference: the representation and its well-formedness discipline
(section~\ref{sec:hyp-ir}); the ten boxes themselves
(section~\ref{sec:hyp-boxes}); the completion theorem that fixed their
number at ten (section~\ref{sec:hyp-why-ten}); the plate extension,
under which one finite graph denotes a function of
variable-dimensional state --- symbolic extents, not an eleventh box
(section~\ref{sec:hyp-plates}); the single-genus design
rule and the view/elaboration discipline it imposes on every other
computational syntax (section~\ref{sec:hyp-genus}); the substitution
operation under which hypernets compose, with its coherence and
digest laws (section~\ref{sec:hyp-subst}); the precise sense in which
Systems, specifications, and rules all consume the same genus
(section~\ref{sec:hyp-core}); and the kinship between the graph-side
contraction box and the relational engine's contraction rows
(section~\ref{sec:hyp-tensors}). A closing ledger
(section~\ref{sec:hyp-receipts}) states which claims in this chapter
are machine-checked, which are proven in prose with executable
witnesses, and which remain in the design tense.

This chapter is also the toolkit's \emph{Primitive Catalogue}. Each of
its thirteen rows --- the ten boxes, the container, the substitution
operation, and the well-formedness preflight --- carries an anchored
entry inside the narrative below, stating the row in compact
metanotation and quoting its producer API, TypeScript then Rust. The
metanotation statements are typeset by hand for notational quality;
their mechanized ground truth is the catalogue's per-box typing rules
and well-formedness predicates, realized twice and in agreement
(section~\ref{sec:hyp-receipts}). The API blocks are generated, never
hand-quoted: signatures are re-extracted from the toolkit sources at
every regeneration, and every construction snippet is quoted verbatim
from the conformance corpus whose emitted canonical bytes are verified
byte-identical against the reference implementation's golden vectors.

A reading map, since the chapter doubles as reference material: on a
first pass, read the representation (section~\ref{sec:hyp-ir}), the
table that opens section~\ref{sec:hyp-boxes}, the completion theorem
of section~\ref{sec:hyp-why-ten} (a self-contained story), and the
single-genus rule of section~\ref{sec:hyp-genus}; return to the
per-box catalogue entries as later chapters cite them, and leave the
plate extension (section~\ref{sec:hyp-plates}) until the
token-population models of section~\ref{sec:mt-sdcpn} send you back
for it. The mathematics is newest in the plate extension and in
section~\ref{sec:hyp-core}, where the chapter's compositional claims
assemble into one machine-checked inclusion theorem
(Theorem~\ref{thm:hyp-inclusion}).

\section{The intermediate representation}
\label{sec:hyp-ir}

A hypernet is a dataflow graph over \emph{tensor-typed wires}. A
tensor type $\langle s, \delta \rangle$ pairs a shape $s$ (a finite
list of extents; scalars have the empty shape $[\,]$) with a dtype
$\delta$ --- Boolean, integer, or exact rational $\mathbb{Q}$ in the
configurations this book uses. Exactness is a design commitment, not
an optimization default: constants enter as exact
rationals,\footnote{An IEEE float read from a model file denotes a
dyadic rational and is stored as one, losslessly; the ACAS~Xu import
of section~\ref{sec:hyp-receipts} carries every published network
weight this way.} the reference evaluator computes in $\mathbb{Q}$,
and any change of numeric representation must pass through an
explicit box with a declared rounding mode. There is no ambient
floating-point semantics for a certificate to be quietly wrong
about.

\begin{definition}[hypernet]
\label{def:hypernet}
A hypernet consists of: a \emph{wire-type table} assigning each wire
a tensor type; a pool of simplicial complexes (the payloads of the
branching box, section~\ref{sec:hyp-box-cell}); a finite sequence of
\emph{applications} of the ten primitive boxes of
section~\ref{sec:hyp-boxes}, each consuming named input wires and
producing named output wires; and an \emph{interface}: an ordered
list of graph inputs and an ordered list of graph outputs. It is
\emph{well-formed} when three conditions hold:
\begin{itemize}
\item \textbf{linearity} --- every wire is produced by at most one
  application (graph inputs by none) and consumed by at most one;
\item \textbf{acyclicity} --- applications are listed in
  static-single-assignment order: every consumed wire is a graph
  input or the output of an earlier application;
\item \textbf{type consistency} --- each application satisfies its
  box's typing rule (stated per box below).
\end{itemize}
\end{definition}

All three conditions are decidable by direct inspection of the ground
data, and both realizations of the IR --- the Lean structure, which
carries the three conditions as fields, and the reference
implementation, which checks them --- enforce them at the same
strength.

\primentry{the well-formedness preflight}{H-WF}
\primmeta
\begin{mathpar}
\inferrule*[Right=wf]
  {\text{linearity}(H) \\ \text{acyclicity}(H) \\
   \text{every application satisfies its box's typing rule}}
  {\mathrm{wf}(H)}
\end{mathpar}
\noindent All three conjuncts are decidable predicates over the ground
data of Definition~\ref{def:hypernet}; the preflight operation decides
them and returns the (possibly empty) list of violations, each carrying
its geometry --- the offending wire, application index, or typing
fact. It is a producer-side courtesy: the checker replays the same
conditions independently and believes only its own verdict.
\ruleblock{API}
\sourcenote{Producer surface of the \texttt{mtk} toolkit for this catalogue entry: the TypeScript arrangement (npm package \texttt{@maitria/mtk}, \texttt{hypernet} namespace) and the Rust catalogue crate (\texttt{crates/mtk-catalogue}) --- one Rust implementation underneath, reached from TypeScript in-process through a native binding, so the two surfaces cannot diverge. Signatures and constructions are extracted from the toolkit sources at every regeneration of this block; each construction is quoted verbatim from the toolkit's hypernet conformance corpus, code whose emitted canonical graph bytes are verified byte-identical against golden vectors generated by the reference implementation (the corpus's fixture graphs --- an all-ten-primitives graph, an identity, a substitution composite, and the fifteen-application drift of this chapter's worked example --- are the receipts). This API constructs and emits; it never checks --- the checker re-derives every claim independently. Later entries of this catalogue cite this note rather than repeating it.}
Courtesy check: the returned list is every violation found, and empty means the linearity, SSA, and per-primitive typing conditions hold and every pooled complex validates --- but the checker's replay is the authority, never this preflight. The corpus construction validates the worked example's fifteen-application drift graph.
\apilang{TypeScript}{hypernet.Hypernet.preflight}
\apiprov{signature --- \texttt{ts/src/hypernet.ts}}
\begin{tssrc}
preflight(): ReadonlyArray<Violation>
\end{tssrc}
\apiprov{construction --- \texttt{ts/tests/hypernet-native.test.ts}}
\begin{tssrc}
const h = hn.unwrap(hn.Hypernet.fromCanonicalBytes(hypernetGolden("g3_cartpole_drift")));
assert.equal(h.preflight().length, 0);
assert.equal(h.appCount, 15);
\end{tssrc}
\apilang{Rust}{mtk\_catalogue::hypernet::Hypernet::preflight}
\apiprov{signature --- \texttt{crates/mtk-catalogue/src/hypernet/wf.rs}}
\begin{rustsrc}
pub fn preflight(&self) -> Vec<Violation>
\end{rustsrc}
\apiprov{construction --- \texttt{crates/mtk-catalogue/tests/golden\_hypernet.rs}}
\begin{rustsrc}
let h = Hypernet::from_canonical_bytes(&golden("g3_cartpole_drift")).unwrap();
assert!(h.preflight().is_empty());
assert_eq!(h.apps.len(), 15);
\end{rustsrc}

Linearity is the load-bearing choice. In most graph formats, sharing
a value is implicit (draw two edges out of a node) and discarding one
is invisible (draw none). Here both are rows: fan-out is the
\texttt{copy} box, fan-in is the \texttt{stack} box, discard is the
\texttt{delete} box, and exchange is the \texttt{swap} box. The
payoff is that structure becomes syntax --- how many times a value is
used, and where, is readable off the application list without any
analysis pass, which is what lets the substitution operation of
section~\ref{sec:hyp-subst} repair seams mechanically and lets
well-formedness be preserved by a proof rather than a linter.

A well-formed hypernet is a \emph{string diagram}: boxes are the
typed, multi-input multi-output operations, wires carry the content,
and two parallel wires between the same boxes are genuinely
distinct. Combinatorially, string diagrams over a symmetric monoidal
signature correspond to typed hypergraphs with interfaces, and the
directed, Frobenius-free fragment is carved out by exactly two
conditions --- monogamy and acyclicity (Bonchi, Gadducci, Kissinger,
Soboci\'{n}ski, and Zanasi~\cite{bonchi22}; the companion
paper~\cite{bonchi20} treats the unrestricted setting) --- which are
the linearity and SSA conditions above under other names. ``Hypernet''
names this object with its interface and canonical serialization
made first-class.\footnote{The name comes from the graph-machine
school that coined it \cite{ghica19,alvarezpicallo21}, whose
hypernets are likewise monogamous directed hypergraphs. There the
interface is read off vertex degrees; here it is declared data, and
the canonical serialization is new.}

Two near neighbours locate the design choice. Keeping monogamy makes
an application list a term of the free symmetric monoidal syntax on
the box signature, in which sharing, discarding, and ordering are
rows. Dropping the consumer half of monogamy would give the
gs-monoidal setting of term graphs, in which fan-out is absorbed
into the wiring~\cite{fritzliang22}; dropping the producer half as
well would give hypergraph categories~\cite{bonchi20}, in which a
twice-produced wire asserts an equation rather than carrying a
value. The calculus stays in the first setting because its consumers
are certificates: how many times a value is used, and where, must be
readable off the bytes without an analysis pass. The choice costs
plumbing rows in the graphs. It does not cost composition: gluing
two hypernets along a shared interface inserts no structural row at
all, a fact section~\ref{sec:hyp-subst} states and the Lean tree
checks.

Serialization is part of the definition's job description, because
certificates refer to computations by content. Every well-formed
hypernet has a canonical byte serialization --- wire ids assigned by
first appearance, complex pool deduplicated, evaluator hints and
debug names excluded --- and its \emph{digest} is the SHA-256 hash of
those bytes. Two design consequences are worth stating before the
boxes: equal digests certify equal computation graphs, not merely
equal functions (two different graphs computing one function
deliberately do \emph{not} share a digest; extensional equality of
programs is not decidable, and the digest never pretends otherwise);
and every certificate row that cites a computation cites it by
digest, so ``which function is this claim about?'' has a
content-addressed answer.

\primentry{the container}{H-GRAPH}
\primmeta
\begin{mathpar}
H = \langle W,\; \mathcal{K},\; A,\; \mathit{in},\; \mathit{out} \rangle
\and
\mathrm{digest}(H) = \operatorname{SHA\text{-}256}\bigl(\mathrm{canon}(H)\bigr)
\end{mathpar}
\noindent The record of Definition~\ref{def:hypernet}: wire-type table
$W$, complex pool $\mathcal{K}$, SSA application list $A$, interface
$(\mathit{in}, \mathit{out})$. The canonical serialization
$\mathrm{canon}$ assigns wire ids by first appearance, deduplicates the
pool, and excludes evaluator hints and debug names; equal digests
certify equal graphs, never merely equal functions.
\ruleblock{API}
\sourcenote{Producer API of the \texttt{mtk} toolkit --- the TypeScript arrangement and the Rust catalogue crate, one implementation underneath. The extraction, corpus, and receipt disciplines governing this block are stated once, at \hyperref[prim:H-WF]{this catalogue's first entry}.}
The trio quoted here is the container's content-address surface --- strict decode, canonical bytes, SHA-256 digest; a graph is closed by \texttt{HypernetBuilder::finish} (TS \texttt{Builder.finish}), which runs the full preflight as a belt. The construction quoted below is the corpus's decode--re-encode--digest loop over the four committed reference vectors.
\apilang{TypeScript}{hypernet.Hypernet}
\apiprov{signature --- \texttt{ts/src/hypernet.ts}}
\begin{tssrc}
static fromCanonicalBytes(bytes: Uint8Array): Result<Hypernet, CanonRefusal>

canonicalBytes(): Result<Uint8Array, CanonRefusal>

digest(): Result<HypernetDigest, CanonRefusal>
\end{tssrc}
\apiprov{construction --- \texttt{ts/tests/hypernet-native.test.ts}}
\begin{tssrc}
const h = hn.unwrap(hn.Hypernet.fromCanonicalBytes(golden));
const re = hn.unwrap(h.canonicalBytes());
assert.equal(toHex(re), toHex(golden), "native re-encode");
assert.equal(hn.unwrap(h.digest()), manifest.fixtures[name]!.sha256, "native digest");
\end{tssrc}
\apilang{Rust}{mtk\_catalogue::hypernet::Hypernet}
\apiprov{signature --- \texttt{crates/mtk-catalogue/src/hypernet/canon.rs}}
\begin{rustsrc}
pub fn from_canonical_bytes(bytes: &[u8]) -> Result<Hypernet, CanonRefusal>

pub fn canonical_bytes(&self) -> Result<Vec<u8>, CanonRefusal>

pub fn digest(&self) -> Result<HypernetDigest, CanonRefusal>
\end{rustsrc}
\apiprov{construction --- \texttt{crates/mtk-catalogue/tests/golden\_hypernet.rs}}
\begin{rustsrc}
let bytes = golden(name);
let h = Hypernet::from_canonical_bytes(&bytes)
    .unwrap_or_else(|e| panic!("{name}: decode refused: {e}"));
let re = h.canonical_bytes().unwrap();
assert_eq!(re, bytes, "{name}: re-encode differs from golden bytes");
assert_eq!(
    h.digest().unwrap().to_hex(),
    manifest_digest(name),
    "{name}: digest differs from manifest"
);
\end{rustsrc}

\section{The ten primitive boxes}
\label{sec:hyp-boxes}

The calculus has ten primitive boxes. Constructor order in the
mechanized catalogue and tag order in the wire format agree
(tags~0--9), and the order below follows them. Throughout,
$\langle s, \delta \rangle$ is a tensor type, $\mathbb{B}$
abbreviates the Boolean scalar type $\langle [\,], \mathrm{B}
\rangle$, and $k{::}s$ is the shape $s$ with a new leading axis of
extent $k$.

\begin{center}
\begin{tabular}{@{}llll@{}}
\toprule
 & box & signature & one-line semantics \\
\midrule
0 & \texttt{const} & $\varnothing \to \langle s,\delta\rangle$ & emit a literal tensor \\
1 & \texttt{contract} & $\langle s_1,\delta\rangle,\dots,\langle s_k,\delta\rangle \to \langle s_0,\delta\rangle$ & sum of products over bound indices \\
2 & \texttt{cast} & $\langle s,\delta\rangle \to \langle s,\delta'\rangle$ & change dtype, declared rounding \\
3 & \texttt{cell} & $\langle [d],\mathbb{Q}\rangle \to \mathbb{B},\dots,\mathbb{B}$ & one-hot region membership \\
4 & \texttt{reshape} & $\langle s,\delta\rangle \to \langle s',\delta\rangle$ & same entries, new shape \\
5 & \texttt{broadcast} & $\langle s,\delta\rangle \to \langle s',\delta\rangle$ & replicate along new axes \\
6 & \texttt{copy} & $\langle s,\delta\rangle \to \langle s,\delta\rangle^{k}$ & explicit fan-out \\
7 & \texttt{delete} & $\langle s,\delta\rangle \to \varnothing$ & explicit discard \\
8 & \texttt{swap} & $\tau_1, \tau_2 \to \tau_2, \tau_1$ & exchange two wires \\
9 & \texttt{stack} & $\langle s,\delta\rangle^{k} \to \langle k{::}s,\delta\rangle$ & explicit fan-in: new leading axis \\
\bottomrule
\end{tabular}
\end{center}

Nine of the ten are deterministic, total, single-valued operations.
\texttt{cell} is the deliberate exception: \emph{the calculus's only
source of disjunction}.
Every case split a hypernet can express --- the two sides of a ReLU,
the five advisory regions of a collision-avoidance policy --- is some
\texttt{cell} application's one-hot output, and every downstream
consumer of a case split can be traced to it. A checker that must
reason by cases knows, syntactically, where all the cases come from.

\subsection{Sources and the multilinear engine: \texttt{const} and \texttt{contract}}
\label{sec:hyp-box-contract}

\textbf{\texttt{const}} takes no inputs and produces one wire of a
declared type, carrying a literal value: network weights, gain
matrices, sign vectors, the all-ones vector that turns a
\texttt{stack} into a sum (section~\ref{sec:hyp-why-ten}). All
numeric literals are exact rationals. Constants are applications like
everything else --- a hypernet's parameters are part of its graph,
hence part of its digest.

\primentry{\texttt{const}}{H-CONST}
\primmeta
\begin{mathpar}
\inferrule*[Right=\texttt{const}]
  {v \text{ a literal of type } \langle s, \delta \rangle
   \text{ (row-major, } \mathrm{size}(s) \text{ entries, exact)}}
  {\mathtt{const}_{v} \colon \varnothing \to \langle s, \delta \rangle}
\end{mathpar}
\noindent Semantics: $\mathrm{out} = v$. Tag 0; the only zero-input
box.
\ruleblock{API}
\sourcenote{Producer API of the \texttt{mtk} toolkit --- the TypeScript arrangement and the Rust catalogue crate, one implementation underneath. The extraction, corpus, and receipt disciplines governing this block are stated once, at \hyperref[prim:H-WF]{this catalogue's first entry}.}
The exact-rational constant is quoted; the Boolean flavour (payload of bits) is constructed in the same corpus fixture, a few lines below the quoted one.
\apilang{TypeScript}{hypernet.Builder.constant}
\apiprov{signature --- \texttt{ts/src/hypernet.ts}}
\begin{tssrc}
constant(
  ty: TensorType,
  entries: ReadonlyArray<QLike> | ReadonlyArray<boolean>,
): Result<Wire, BuildRefusal>
\end{tssrc}
\apiprov{construction --- \texttt{ts/tests/hypernet-native.test.ts}}
\begin{tssrc}
const c1 = hn.unwrap(b.constant(hn.rat([2]), ["3/2", "-7/3"]));
\end{tssrc}
\apilang{Rust}{mtk\_catalogue::hypernet::HypernetBuilder::constant}
\apiprov{signature --- \texttt{crates/mtk-catalogue/src/hypernet/graph.rs}}
\begin{rustsrc}
pub fn constant(&mut self, ty: TensorType, value: ConstPayload) -> Result<Wire, BuildRefusal>
\end{rustsrc}
\apiprov{construction --- \texttt{crates/mtk-catalogue/tests/golden\_hypernet.rs}}
\begin{rustsrc}
let c1 = b
    .constant(
        TensorType::rat([2]),
        ConstPayload::Rat(vec![q(3, 2), q(-7, 3)]),
    )
    .unwrap();
\end{rustsrc}

\textbf{\texttt{contract}} is the entire multilinear engine: the
$n$-ary generalized Einstein summation, with its index structure
carried as explicit data rather than inferred from a notation string.

\begin{definition}[contraction pattern]
\label{def:contraction-pattern}
A contraction pattern for inputs of ranks $r_1, \dots, r_k$ and
output rank $r_0$ consists of a finite index set $I$, an extent
$E_i \in \mathbb{N}$ for each $i \in I$, a binding
$\beta_m \colon \{0,\dots,r_m{-}1\} \to I$ for each input $m$, and a
binding $\beta_0 \colon \{0,\dots,r_0{-}1\} \to I$ for the output;
type consistency requires every bound axis's extent to equal its
index's extent. Writing $\nu\beta$ for the coordinate tuple obtained
by composing an index assignment $\nu \in \prod_{i \in I} \{0,\dots,
E_i{-}1\}$ with a binding $\beta$, the semantics is
\[
\mathrm{out}[c] \;=\; \sum_{\nu \,:\, \nu\beta_0 = c} \;
\prod_{m=1}^{k} \mathrm{in}_m[\nu\beta_m].
\]
\end{definition}

Matrix product is a contract ($ai, ij \to aj$ in einsum notation);
so are dot products, traces, elementwise products, tensor products,
and axis sums. Two structural facts about this box do quiet work
later. First, a rank-0 input has no axes to bind but still
contributes its single entry as a factor to \emph{every} term of the
sum --- scalars multiply in; they do not add in. Second, and in the
same direction: every term of every output entry contains a factor
from \emph{every} input wire. \texttt{contract} is purely
multiplicative across wires, and section~\ref{sec:hyp-why-ten} turns
that observation into a theorem about what the box cannot do.

\primentry{\texttt{contract}}{H-CONTRACT}
\primmeta
\begin{mathpar}
\inferrule*[Right=\texttt{contract}]
  {\beta = \bigl(I,\, E,\, \beta_0,\, (\beta_m)_{m=1}^{k}\bigr)
   \text{ a contraction pattern (Definition~\ref{def:contraction-pattern})} \\
   \text{each bound axis's extent equals its index's extent}}
  {\mathtt{contract}_{\beta} \colon
   \langle s_1, \delta \rangle, \dots, \langle s_k, \delta \rangle
   \to \langle s_0, \delta \rangle}
\end{mathpar}
\noindent Semantics: Definition~\ref{def:contraction-pattern}'s sum of
products, $\mathrm{out}[c] = \sum_{\nu \,:\, \nu\beta_0 = c}
\prod_{m} \mathrm{in}_m[\nu\beta_m]$. Tag 1; the entire multilinear
engine, index structure as explicit data.
\ruleblock{API}
\sourcenote{Producer API of the \texttt{mtk} toolkit --- the TypeScript arrangement and the Rust catalogue crate, one implementation underneath. The extraction, corpus, and receipt disciplines governing this block are stated once, at \hyperref[prim:H-WF]{this catalogue's first entry}.}
\apilang{TypeScript}{hypernet.Builder.contract}
\apiprov{signature --- \texttt{ts/src/hypernet.ts}}
\begin{tssrc}
contract(
  inputs: ReadonlyArray<Wire>,
  bindings: {
    extents: ReadonlyArray<number>;
    inBindings: ReadonlyArray<ReadonlyArray<number>>;
    outBinding: ReadonlyArray<number>;
  },
): Result<Wire, BuildRefusal>
\end{tssrc}
\apiprov{construction --- \texttt{ts/tests/hypernet-native.test.ts}}
\begin{tssrc}
const d = hn.unwrap(
  b.contract([x, c1], { extents: [2], inBindings: [[0], [0]], outBinding: [] }),
);
\end{tssrc}
\apilang{Rust}{mtk\_catalogue::hypernet::HypernetBuilder::contract}
\apiprov{signature --- \texttt{crates/mtk-catalogue/src/hypernet/graph.rs}}
\begin{rustsrc}
pub fn contract(
    &mut self,
    inputs: &[Wire],
    extents: &[u64],
    in_bindings: &[Vec<u32>],
    out_binding: &[u32],
) -> Result<Wire, BuildRefusal>
\end{rustsrc}
\apiprov{construction --- \texttt{crates/mtk-catalogue/tests/golden\_hypernet.rs}}
\begin{rustsrc}
let d = b
    .contract(&[x, c1], &[2], &[vec![0], vec![0]], &[])
    .unwrap();
\end{rustsrc}

\subsection{Representation moves: \texttt{cast}, \texttt{reshape}, \texttt{broadcast}}

\textbf{\texttt{cast}} converts dtype at fixed shape, carrying an
explicit rounding mode. It is the only box at which numeric
representation changes, which localizes every exactness argument: a
hypernet with no \texttt{cast} (or only exact ones, such as the
Boolean-to-rational conversion that turns an indicator into a gate)
computes in $\mathbb{Q}$ end to end, and one carrying rounding casts
declares, in its syntax, where error can enter.

\primentry{\texttt{cast}}{H-CAST}
\primmeta
\begin{mathpar}
\inferrule*[Right=\texttt{cast}]
  {\rho \text{ a declared rounding mode for } \delta \to \delta'}
  {\mathtt{cast}_{\delta \to \delta',\, \rho} \colon
   \langle s, \delta \rangle \to \langle s, \delta' \rangle}
\end{mathpar}
\noindent Semantics: entrywise conversion under $\rho$; exact casts
(the Boolean-to-rational indicator gate among them) introduce no
error. Tag 2; the only box at which numeric representation changes.
\ruleblock{API}
\sourcenote{Producer API of the \texttt{mtk} toolkit --- the TypeScript arrangement and the Rust catalogue crate, one implementation underneath. The extraction, corpus, and receipt disciplines governing this block are stated once, at \hyperref[prim:H-WF]{this catalogue's first entry}.}
\apilang{TypeScript}{hypernet.Builder.cast}
\apiprov{signature --- \texttt{ts/src/hypernet.ts}}
\begin{tssrc}
cast(inp: Wire, to: Dtype): Result<Wire, BuildRefusal>
\end{tssrc}
\apiprov{construction --- \texttt{ts/tests/hypernet-native.test.ts}}
\begin{tssrc}
const ci = hn.unwrap(b.cast(bc, hn.INT));
\end{tssrc}
\apilang{Rust}{mtk\_catalogue::hypernet::HypernetBuilder::cast}
\apiprov{signature --- \texttt{crates/mtk-catalogue/src/hypernet/graph.rs}}
\begin{rustsrc}
pub fn cast(&mut self, inp: Wire, to: Dtype) -> Result<Wire, BuildRefusal>
\end{rustsrc}
\apiprov{construction --- \texttt{crates/mtk-catalogue/tests/golden\_hypernet.rs}}
\begin{rustsrc}
let ci = b.cast(bc, Dtype::INT).unwrap();
\end{rustsrc}

\textbf{\texttt{reshape}} re-indexes a tensor's entries into a new
shape of the same total size --- row-major order, no data motion in
the semantics. \textbf{\texttt{broadcast}} replicates along new or
singleton axes, under the trailing-axis compatibility rule familiar
from array programming; the result is constant along each broadcast
axis. Both are plumbing: they let \texttt{contract} patterns stay
simple by normalizing operand shapes, and their semantics are
entry-relabelling and entry-duplication, contributing no arithmetic of
their own.

\primentry{\texttt{reshape}}{H-RESHAPE}
\primmeta
\begin{mathpar}
\inferrule*[Right=\texttt{reshape}]
  {\mathrm{size}(s') = \mathrm{size}(s)}
  {\mathtt{reshape}_{s'} \colon \langle s, \delta \rangle \to
   \langle s', \delta \rangle}
\end{mathpar}
\noindent Semantics: the same entries under row-major re-indexing; no
data motion. Tag 4.
\ruleblock{API}
\sourcenote{Producer API of the \texttt{mtk} toolkit --- the TypeScript arrangement and the Rust catalogue crate, one implementation underneath. The extraction, corpus, and receipt disciplines governing this block are stated once, at \hyperref[prim:H-WF]{this catalogue's first entry}.}
\apilang{TypeScript}{hypernet.Builder.reshape}
\apiprov{signature --- \texttt{ts/src/hypernet.ts}}
\begin{tssrc}
reshape(inp: Wire, newShape: ReadonlyArray<number>): Result<Wire, BuildRefusal>
\end{tssrc}
\apiprov{construction --- \texttt{ts/tests/hypernet-native.test.ts}}
\begin{tssrc}
const r = hn.unwrap(b.reshape(t, [3, 1]));
\end{tssrc}
\apilang{Rust}{mtk\_catalogue::hypernet::HypernetBuilder::reshape}
\apiprov{signature --- \texttt{crates/mtk-catalogue/src/hypernet/graph.rs}}
\begin{rustsrc}
pub fn reshape(&mut self, inp: Wire, new_shape: &[u64]) -> Result<Wire, BuildRefusal>
\end{rustsrc}
\apiprov{construction --- \texttt{crates/mtk-catalogue/tests/golden\_hypernet.rs}}
\begin{rustsrc}
let r = b.reshape(t, &[3, 1]).unwrap();
\end{rustsrc}

\primentry{\texttt{broadcast}}{H-BROADCAST}
\primmeta
\begin{mathpar}
\inferrule*[Right=\texttt{broadcast}]
  {s \trianglelefteq s' \text{ (trailing-axis compatible)}}
  {\mathtt{broadcast}_{s'} \colon \langle s, \delta \rangle \to
   \langle s', \delta \rangle}
\end{mathpar}
\noindent Semantics: replication along the new and singleton axes; the
result is constant along each broadcast axis. Tag 5.
\ruleblock{API}
\sourcenote{Producer API of the \texttt{mtk} toolkit --- the TypeScript arrangement and the Rust catalogue crate, one implementation underneath. The extraction, corpus, and receipt disciplines governing this block are stated once, at \hyperref[prim:H-WF]{this catalogue's first entry}.}
\apilang{TypeScript}{hypernet.Builder.broadcast}
\apiprov{signature --- \texttt{ts/src/hypernet.ts}}
\begin{tssrc}
broadcast(inp: Wire, toShape: ReadonlyArray<number>): Result<Wire, BuildRefusal>
\end{tssrc}
\apiprov{construction --- \texttt{ts/tests/hypernet-native.test.ts}}
\begin{tssrc}
const bc = hn.unwrap(b.broadcast(r, [3, 2]));
\end{tssrc}
\apilang{Rust}{mtk\_catalogue::hypernet::HypernetBuilder::broadcast}
\apiprov{signature --- \texttt{crates/mtk-catalogue/src/hypernet/graph.rs}}
\begin{rustsrc}
pub fn broadcast(&mut self, inp: Wire, to_shape: &[u64]) -> Result<Wire, BuildRefusal>
\end{rustsrc}
\apiprov{construction --- \texttt{crates/mtk-catalogue/tests/golden\_hypernet.rs}}
\begin{rustsrc}
let bc = b.broadcast(r, &[3, 2]).unwrap();
\end{rustsrc}

\subsection{Structural plumbing: \texttt{copy}, \texttt{delete}, \texttt{swap}}

\textbf{\texttt{copy}} duplicates one wire into $k$ wires of the same
type. Under linearity it is the only way a value acquires multiple
consumers, so sharing is always a visible row --- when the cartpole
example of section~\ref{sec:hyp-genus} uses the pole-angle sine in
three monomials, the graph says so with a single ternary \texttt{copy}.

\primentry{\texttt{copy}}{H-COPY}
\primmeta
\begin{mathpar}
\inferrule*[Right=\texttt{copy}]
  {k \ge 1}
  {\mathtt{copy}_{k} \colon \langle s, \delta \rangle \to
   \langle s, \delta \rangle^{k}}
\end{mathpar}
\noindent Semantics: $\mathrm{out}_i = \mathrm{in}$ for every branch.
Tag 6; under linearity, the only way a value acquires a second
consumer.
\ruleblock{API}
\sourcenote{Producer API of the \texttt{mtk} toolkit --- the TypeScript arrangement and the Rust catalogue crate, one implementation underneath. The extraction, corpus, and receipt disciplines governing this block are stated once, at \hyperref[prim:H-WF]{this catalogue's first entry}.}
\apilang{TypeScript}{hypernet.Builder.copy}
\apiprov{signature --- \texttt{ts/src/hypernet.ts}}
\begin{tssrc}
copy(inp: Wire, k: number): Result<ReadonlyArray<Wire>, BuildRefusal>
\end{tssrc}
\apiprov{construction --- \texttt{ts/tests/hypernet-native.test.ts}}
\begin{tssrc}
const ys = hn.unwrap(b.copy(y, 2));
\end{tssrc}
\apilang{Rust}{mtk\_catalogue::hypernet::HypernetBuilder::copy}
\apiprov{signature --- \texttt{crates/mtk-catalogue/src/hypernet/graph.rs}}
\begin{rustsrc}
pub fn copy(&mut self, inp: Wire, k: u32) -> Result<Vec<Wire>, BuildRefusal>
\end{rustsrc}
\apiprov{construction --- \texttt{crates/mtk-catalogue/tests/golden\_hypernet.rs}}
\begin{rustsrc}
let ys = b.copy(y, 2).unwrap();
\end{rustsrc}

\textbf{\texttt{delete}} consumes one wire and produces nothing. A
value nobody reads must be discarded by a row, never dropped by
omission; the substitution operation of section~\ref{sec:hyp-subst}
leans on this (``the discard is a row, never an absence'') when it
repairs the seam between two graphs.

\primentry{\texttt{delete}}{H-DELETE}
\primmeta
\begin{mathpar}
\inferrule*[Right=\texttt{delete}]
  {\;}
  {\mathtt{delete} \colon \langle s, \delta \rangle \to \varnothing}
\end{mathpar}
\noindent Semantics: no outputs; the input's consumption is the entire
content. Tag 7.
\ruleblock{API}
\sourcenote{Producer API of the \texttt{mtk} toolkit --- the TypeScript arrangement and the Rust catalogue crate, one implementation underneath. The extraction, corpus, and receipt disciplines governing this block are stated once, at \hyperref[prim:H-WF]{this catalogue's first entry}.}
\apilang{TypeScript}{hypernet.Builder.delete}
\apiprov{signature --- \texttt{ts/src/hypernet.ts}}
\begin{tssrc}
delete(inp: Wire): Result<undefined, BuildRefusal>
\end{tssrc}
\apiprov{construction --- \texttt{ts/tests/hypernet-native.test.ts}}
\begin{tssrc}
hn.unwrap(b.delete(cells[0]!));
\end{tssrc}
\apilang{Rust}{mtk\_catalogue::hypernet::HypernetBuilder::delete}
\apiprov{signature --- \texttt{crates/mtk-catalogue/src/hypernet/graph.rs}}
\begin{rustsrc}
pub fn delete(&mut self, inp: Wire) -> Result<(), BuildRefusal>
\end{rustsrc}
\apiprov{construction --- \texttt{crates/mtk-catalogue/tests/golden\_hypernet.rs}}
\begin{rustsrc}
b.delete(cells[0]).unwrap();
\end{rustsrc}

\textbf{\texttt{swap}} exchanges two wires, of possibly different
types. It is the symmetry of the monoidal structure as an explicit
application --- pure renaming, no computation --- so that wire order,
like sharing and discarding, is syntax rather than convention.

\primentry{\texttt{swap}}{H-SWAP}
\primmeta
\begin{mathpar}
\inferrule*[Right=\texttt{swap}]
  {\;}
  {\mathtt{swap} \colon \tau_1, \tau_2 \to \tau_2, \tau_1}
\end{mathpar}
\noindent Semantics: $(\mathrm{out}_1, \mathrm{out}_2) =
(\mathrm{in}_2, \mathrm{in}_1)$ --- pure renaming, types possibly
distinct. Tag 8.
\ruleblock{API}
\sourcenote{Producer API of the \texttt{mtk} toolkit --- the TypeScript arrangement and the Rust catalogue crate, one implementation underneath. The extraction, corpus, and receipt disciplines governing this block are stated once, at \hyperref[prim:H-WF]{this catalogue's first entry}.}
\apilang{TypeScript}{hypernet.Builder.swap}
\apiprov{signature --- \texttt{ts/src/hypernet.ts}}
\begin{tssrc}
swap(in1: Wire, in2: Wire): Result<readonly [Wire, Wire], BuildRefusal>
\end{tssrc}
\apiprov{construction --- \texttt{ts/tests/hypernet-native.test.ts}}
\begin{tssrc}
const [s1, s2] = hn.unwrap(b.swap(d, ys[0]!));
\end{tssrc}
\apilang{Rust}{mtk\_catalogue::hypernet::HypernetBuilder::swap}
\apiprov{signature --- \texttt{crates/mtk-catalogue/src/hypernet/graph.rs}}
\begin{rustsrc}
pub fn swap(&mut self, in1: Wire, in2: Wire) -> Result<(Wire, Wire), BuildRefusal>
\end{rustsrc}
\apiprov{construction --- \texttt{crates/mtk-catalogue/tests/golden\_hypernet.rs}}
\begin{rustsrc}
let (s1, s2) = b.swap(d, ys[0]).unwrap();
\end{rustsrc}

\subsection{Branching: \texttt{cell}}
\label{sec:hyp-box-cell}

\textbf{\texttt{cell}} is the case-analysis box. Its payload is a
simplicial complex $K$ in ambient dimension $d$: exact-rational
vertex coordinates, a downward-closed and conformal list of
simplices, and a list of \emph{cells}, each a pair $(S, F)$ of a
simplex $S$ and a proper subset $F \subsetneq S$ of its vertices
(the \emph{extension set}). A point $x \in \mathbb{Q}^d$ lies in the
cell $(S, F)$ when $x$ is an affine combination $x = \sum_{v \in S}
\lambda_v\, v$ with $\sum_v \lambda_v = 1$, $\lambda_v \ge 0$ on the
core vertices $S \setminus F$, and $\lambda_v \le 0$ on $F$. With
$F = \varnothing$ this is the closed simplex itself; each extension
vertex opens the region outward through the facet opposite it. The
sign split of the line --- the instance that gates a ReLU below ---
takes the path complex with vertices $-1$, $0$, $1$: the cell pairing
the edge $[0,1]$ with extension set $\{1\}$ covers the nonpositive
half-line, the cell pairing $[-1,0]$ with $\{-1\}$ covers the
nonnegative half-line, and the two overlap only at the origin.

The box's typing: input one wire $\langle [d], \mathbb{Q} \rangle$
--- a point of the ambient space --- and output one Boolean scalar
per listed cell. Its semantics: the outputs are one-hot, with the hot
output naming a cell that contains the input point. Where the point
lies in several cells at once --- on a shared boundary --- \emph{any}
of the containing cells is an admissible choice. The hypernet's
semantics at such a point is the set of behaviours under all
admissible choices; a rule or certificate that cites a \texttt{cell}
output must hold under every one of them, and the reference evaluator's
tie-break (lowest cell index, documented as such) is a pinned
representative, not a semantic commitment. Boundary nondeterminism is
thereby owned by the proof obligations, not hidden in an evaluator
convention.

One box in this role is a deliberate economy. A ReLU is
\texttt{copy}, then \texttt{cell} over the sign split of the line
(two cells: nonnegative and nonpositive), then a gating
\texttt{contract} of the indicator against the value; an
argmax-style advisory head is one \texttt{cell} over the fan of
dominance cones in difference coordinates. Piecewise structure of any
shape enters through this single door, and ``where could this
computation branch?'' is answerable by listing its \texttt{cell}
applications.

\primentry{\texttt{cell}}{H-CELL}
\primmeta
\begin{mathpar}
\inferrule*[Right=\texttt{cell}]
  {K \text{ a valid simplicial complex in ambient dimension } d \\
   \text{with cells } (S_0, F_0), \dots, (S_{n-1}, F_{n-1})}
  {\mathtt{cell}_{K} \colon \langle [d], \mathbb{Q} \rangle \to
   \mathbb{B}, \dots, \mathbb{B} \quad (n \text{ outputs})}
\end{mathpar}
\noindent Semantics: the outputs are one-hot; a hot index $j$ is
admissible exactly when the input point lies in cell $(S_j, F_j)$, and
on a shared boundary \emph{every} containing cell is admissible ---
the nondeterminism belongs to the proof obligations, never to an
evaluator convention. Tag 3; the calculus's only source of
disjunction.
\ruleblock{API}
\sourcenote{Producer API of the \texttt{mtk} toolkit --- the TypeScript arrangement and the Rust catalogue crate, one implementation underneath. The extraction, corpus, and receipt disciplines governing this block are stated once, at \hyperref[prim:H-WF]{this catalogue's first entry}.}
\apilang{TypeScript}{hypernet.Builder.cell}
\apiprov{signature --- \texttt{ts/src/hypernet.ts}}
\begin{tssrc}
cell(complexIdx: number, inp: Wire): Result<ReadonlyArray<Wire>, BuildRefusal>

export interface SimpC {
  readonly ambDim: number;
  /** Vertex coordinates, `ambDim` exact rationals each. */
  readonly vertices: ReadonlyArray<ReadonlyArray<QLike>>;
  readonly simplices: ReadonlyArray<ReadonlyArray<number>>;
  readonly cells: ReadonlyArray<SimpCCell>;
  readonly hint?: string;
}
\end{tssrc}
\apiprov{construction --- \texttt{ts/tests/hypernet-native.test.ts}}
\begin{tssrc}
const cells = hn.unwrap(b.cell(sigma, p));
\end{tssrc}
\apilang{Rust}{mtk\_catalogue::hypernet::HypernetBuilder::cell}
\apiprov{signature --- \texttt{crates/mtk-catalogue/src/hypernet/graph.rs}}
\begin{rustsrc}
pub fn cell(&mut self, complex_idx: u32, inp: Wire) -> Result<Vec<Wire>, BuildRefusal>
\end{rustsrc}
\apiprov{construction --- \texttt{crates/mtk-catalogue/tests/golden\_hypernet.rs}}
\begin{rustsrc}
let cells = b.cell(sigma, p).unwrap();
\end{rustsrc}

\subsection{Fan-in: \texttt{stack}}
\label{sec:hyp-box-stack}

\textbf{\texttt{stack}} takes $k$ wires of one common type $\langle
s, \delta \rangle$ and produces a single wire of type $\langle k{::}s,
\delta \rangle$, with semantics
\[
\mathrm{out}[i, \mathit{idx}] = \mathrm{in}_i[\mathit{idx}],
\qquad i = 0, \dots, k{-}1 .
\]
It is \texttt{copy}'s dual --- fan-in as explicit structure, exactly
as fan-out is explicit \texttt{copy} --- and its characteristic idiom
is the \emph{sum}: to add wires $u$ and $v$, stack them and contract
the stacked wire against the all-ones vector,
\[
u + v \;=\; \mathtt{contract}\bigl(\mathtt{stack}(u, v),\;
\mathbf{1}_2\bigr),
\]
binding the new leading axis. Subtraction is the same idiom with a
$[1, -1]$ sign vector, which keeps imported model weights
byte-faithful: negation lives in a separate constant, never in a
rewritten payload. \texttt{stack} adds no disjunction (its facts are
entry equalities, structural like \texttt{copy}'s), and it is the
newest member of the calculus --- the next section is the story of
why it had to exist.

\primentry{\texttt{stack}}{H-STACK}
\primmeta
\begin{mathpar}
\inferrule*[Right=\texttt{stack}]
  {k \ge 1 \\ \text{all inputs of one type } \langle s, \delta \rangle}
  {\mathtt{stack}_{k} \colon \langle s, \delta \rangle^{k} \to
   \langle k{::}s, \delta \rangle}
\end{mathpar}
\noindent Semantics: $\mathrm{out}[i, \mathit{idx}] =
\mathrm{in}_i[\mathit{idx}]$. Tag 9; \texttt{copy}'s dual, and the
$\Sigma\Pi$ boundary's visible seam.
\ruleblock{API}
\sourcenote{Producer API of the \texttt{mtk} toolkit --- the TypeScript arrangement and the Rust catalogue crate, one implementation underneath. The extraction, corpus, and receipt disciplines governing this block are stated once, at \hyperref[prim:H-WF]{this catalogue's first entry}.}
\apilang{TypeScript}{hypernet.Builder.stack}
\apiprov{signature --- \texttt{ts/src/hypernet.ts}}
\begin{tssrc}
stack(inputs: ReadonlyArray<Wire>): Result<Wire, BuildRefusal>
\end{tssrc}
\apiprov{construction --- \texttt{ts/tests/hypernet-native.test.ts}}
\begin{tssrc}
const t = hn.unwrap(b.stack([s1, s2, ys[1]!]));
\end{tssrc}
\apilang{Rust}{mtk\_catalogue::hypernet::HypernetBuilder::stack}
\apiprov{signature --- \texttt{crates/mtk-catalogue/src/hypernet/graph.rs}}
\begin{rustsrc}
pub fn stack(&mut self, inputs: &[Wire]) -> Result<Wire, BuildRefusal>
\end{rustsrc}
\apiprov{construction --- \texttt{crates/mtk-catalogue/tests/golden\_hypernet.rs}}
\begin{rustsrc}
let t = b.stack(&[s1, s2, ys[1]]).unwrap();
\end{rustsrc}

\section{Why ten: a completion theorem}
\label{sec:hyp-why-ten}

The calculus was not designed to have ten boxes. It had nine ---
\texttt{stack} was absent --- until compiling real networks forced
the question of what the nine could not say, and the answer turned
out to be a theorem rather than an inconvenience.

\begin{proposition}[the packing gap]
\label{prop:packing-gap}
In the nine-box fragment (without \texttt{stack}), no wire can carry
the sum of two independently produced wires: for any two wires $u,
v$, no later wire has an entry whose expansion --- read with $u$ and
$v$ as opaque sources, whatever produced them --- contains both a
monomial with $u$-factors but no $v$-factors and a monomial with
$v$-factors but no $u$-factors; in particular, no entry expands to
$u + v$ itself.\footnote{The expansion stratum, not run-time values,
is where the claim holds at this generality: two \texttt{copy}
descendants of one wire are produced by distinct applications, yet
twice their common value is one \texttt{contract} against a
constant. Read over expansions with $u$ and $v$ held opaque, the
statement needs no independence hypothesis at all, and that is the
form the mechanization proves
(section~\ref{sec:hyp-receipts}).}
\end{proposition}

\emph{Proof sketch.} Only \texttt{contract} combines values from more
than one data wire ---
\texttt{cast}, \texttt{cell}, \texttt{reshape}, \texttt{broadcast},
\texttt{copy}, and \texttt{delete} are single-input, \texttt{const}
is zero-input, and \texttt{swap} is a renaming. But every output
entry of a \texttt{contract} is a sum of monomials each containing a
factor from \emph{every} input wire
(Definition~\ref{def:contraction-pattern}: the product runs over all
$m$) --- the box is purely multiplicative across wires, and a
constant input multiplies into every term rather than adding. So the
first application to consume both $u$ and $v$ already violates the
required shape, and single-input boxes only relabel or duplicate
entries thereafter; induction along the SSA order closes the
argument. \qed

The practical blast radius of this small proposition was total: no
bias can be added to a linear layer ($x - 1$ is already
inexpressible), and after a ReLU layer splits into per-scalar gate
indicators, nothing can reassemble the gated values into the next
layer's input --- so no network with two or more neurons could be
expressed at all, even though each individual matrix product and each
individual ReLU could. Of the two sub-gaps, one (additive constants)
admitted a workaround --- an augmented-input convention carrying a
constant-one component --- at the price of an interface wart; the
other (cross-wire packing of independently produced values) admitted
none.

The repair is the tenth box. \texttt{stack} closes both gaps at once:
sums become the stack-and-ones idiom of
section~\ref{sec:hyp-box-stack}, gate vectors assemble from
per-scalar \texttt{cell} outputs, and with it every piecewise-polynomial dataflow --- every ReLU network included --- becomes
expressible. Alternatives were weighed and declined: a vectorized
\texttt{cell} (fixes gates, not bias, and complicates the format); a
polynomial node promoted to primitive rank (collapses the einsum
discipline --- section~\ref{sec:hyp-genus} explains why the
polynomial belongs at the view layer instead); a \texttt{contract}
variant that also sums over a choice index (semantically equivalent
to \texttt{stack}-then-\texttt{contract}, but it muddies the einsum
reading, whereas \texttt{stack} keeps \emph{product} and \emph{sum}
orthogonal). \texttt{stack} entered the calculus the same day the
gap was written down.

Two properties of the repair deserve the emphasis of a design
principle. First, \texttt{stack} introduces no new disjunction ---
the invariant that \texttt{cell} is the calculus's only branching
source survives the extension untouched. Second, the division of
labour it induces is legible in compiled graphs: products form at
\texttt{contract} boxes, cross-wire sums form at
\texttt{stack}-plus-ones boxes, and the boundary between the
$\Sigma$ and the $\Pi$ of a polynomial is a visible seam in the
graph. The worked example of the next section shows that seam
carrying exactly one \texttt{stack}.

What the ten boxes still deliberately cannot say, so that no reader
over-extrapolates the completion: there is no transcendental
function box, and there is no division box. An oscillator enters
through the standard cylinder lift --- carry $s = \sin\theta$ and
$c = \cos\theta$ as state variables with the algebraic invariant
$s^2 + c^2 = 1$ --- rather than through a $\sin$ node; a rational
right-hand side clears its mass matrix into implicit polynomial
form. Both exclusions keep the genus inside piecewise-polynomial
dataflow over $\mathbb{Q}$, which is where exact certification
machinery has traction. The exclusion is a fact about boxes, not
about reach: the admissible-symbol discipline of
section~\ref{sec:hyp-genus} extends to a register of \emph{certified
function aliases} through which sine and cosine, rational powers,
logarithm, and exponential become available anyway --- as
piecewise-polynomial approximations with declared domains and
accuracy budgets, receipt-checked in exact rational arithmetic ---
without the wire format ever learning their names. Should a
modelling frame ever genuinely force a quotient node, the procedural
path is the one \texttt{stack} took: a written inexpressibility
finding, an explicit decision, and enactment at the calculus layer
--- never a side-channel evaluator bolted onto a view.

\section{Plates: symbolic extents}
\label{sec:hyp-plates}

Token-population systems --- coloured Petri nets first among them
(section~\ref{sec:mt-sdcpn}) --- have variable-dimensional state: the
number of tokens changes as the system runs, and is canonically
unbounded across a run even though it is finite at every instant. A
certificate about such a system is a function on markings ---
multisets of token colours --- and a naive syntax would need one
computation graph per token count: a countable family of artefacts
where one certificate was wanted. The calculus closes this gap with
\emph{plates}: symbolic extents, not an eleventh box.

\paragraph{The data.} Two additions to
Definition~\ref{def:hypernet}, both container-level. The container
gains an ordered finite list of \emph{plate variables}
$\iota_1, \dots, \iota_r$ (serialized de Bruijn by position; debug
names excluded from canon, like wire names). And the extent alphabet
in wire shapes and contraction-pattern index extents becomes
$E ::= n \mid \iota$ --- a literal or a declared plate variable, with
no extent arithmetic ($\iota + 1$ is not an extent; products of
plates arise as multiple axes on one wire, never as extent
products). A \emph{plated hypernet} is the same record over the
extended alphabet plus the plate list; a plated hypernet with zero
plates \emph{is} a hypernet, byte for byte. For an extent assignment
$\sigma : \{\iota_1, \dots, \iota_r\} \to \mathbb{N}$,
\emph{instantiation} $\mathrm{inst}_\sigma$ replaces every symbolic
extent by its value, yielding an ordinary hypernet; everything
downstream --- evaluator, canonical serialization, the whole checker
stack --- consumes instances and only instances. The plate is marked
on \emph{wires} (a per-token scalar is a wire of shape $[\iota]$; a
token family of $d$-dimensional colours, $[\iota, d]$; an attention
matrix, $[\iota, \iota]$), never as a membrane around a subgraph:
scope would be degenerate --- plate variables are container-level ---
and what a membrane would add beyond axes is per-index variation of
the body, which is exactly what the equivariance theorem below must
forbid.

\paragraph{How the boxes lift.} One principle, ten instances: a box
lifts along plate axes by whiskering --- the lifted application is
the original box applied per index tuple, well-typed exactly when
the box's typing clause holds over the extended alphabet, with
extent equality read \emph{syntactically} (literal against literal
numerically; $\iota$ against $\iota$ by name; a $\iota$-against-literal
match is refused --- conservative, never unsound, since schematic
typing must hold at every $\sigma$). \texttt{contract} admits
symbolic index extents (a bound symbolic index is a pool over the
instance's extent; a free one maps over it); \texttt{broadcast} may
add or extend plate axes under the trailing-axis rule, the
membrane-crossing law in wire clothing (a plate-free wire broadcast
along $\iota$ is constant along the index); \texttt{cast},
\texttt{copy}, \texttt{delete}, \texttt{swap} are pointwise;
\texttt{stack} keeps literal arity (a plate-sized family is never
$\iota$ separate wires --- it is one wire carrying the $\iota$ axis,
from birth); \texttt{reshape} is refused on symbolic-extent wires in
this edition (symbolic totalSize equality is decidable polynomial
identity, but \texttt{contract} carries the demand natively, so
nothing routes through it); and \texttt{const} --- the load-bearing
refusal --- takes all-literal shapes only, forced rather than
chosen: a literal tensor has finitely many entries.

\paragraph{Three theorems, nearly free from one syntactic fact.} The
language has no token literals: no expression for ``the $i$-th
token,'' no order or successor on the index sort, no per-index
branching (a \texttt{cell} reads colours, never indices), and
\texttt{const} cannot utter a plate. Three consequences make the
notation functorial rather than decorative.

\begin{itemize}
\item \textbf{Schematic well-formedness.} A well-formed plated
  hypernet instantiates to a well-formed ordinary hypernet at
  \emph{every} extent assignment: instantiation is a homomorphism on
  extents, every typing clause is syntactic extent equality plus
  extent-free conditions, and both survive substitution.
\item \textbf{Equivariance for free.} The instance family is natural
  in bijections of each plate's index set: every expressible plated
  computation is equivariant under the product of symmetric groups,
  \emph{as a theorem of the grammar} rather than a proof obligation
  on artefacts --- plate axes are born only from the interface,
  \texttt{broadcast}, and \texttt{contract}, all equivariant, and
  the mutant that adds an index-literal box breaks the theorem at a
  two-token transposition, kernel-visibly. A corollary worth the
  emphasis: \emph{causal masking is unwritable}. A causal mask is a
  function of the index order, and the syntax has no order to
  consult --- the transformer this calculus can express is
  inherently the exchangeable one.
\item \textbf{Descent.} An output wire carrying no plate axis is
  invariant, hence descends to the multiset quotient: the family
  assembles into one function on the bag space
  $\bigsqcup_{n < \omega} \mathrm{Sym}^n(B)$ --- with one plate per
  place, precisely the coloured-net marking space of
  section~\ref{sec:mt-sdcpn}. One finite syntax object denotes a
  function of the whole variable-dimensional state.
\end{itemize}

Semantically the plate is a functorial box in the sense of
Melliès~\cite{mellies06csl} --- the strong monoidal functor
$(-)^{\otimes I}$, drawn as axes instead of a membrane --- and at
infinite index it is the plate over a Kolmogorov product in the
sense of categorical probability~\cite{chen2024aldoushoover}, the unique morphism
whose finite marginals are the body's instances. The index category
is deliberately the groupoid of finite sets and bijections, not all
functions: aggregation is not functorial under instance
insertions~\cite{spivak2021agg}, and equivariance is exactly the
functoriality that survives --- which is the symmetry the
marking space quotients by.

\paragraph{The count wire.} The \texttt{const} refusal bans $\iota$
in a literal's \emph{shape}, never the scalar. So the token count is
three applications --- $\texttt{const}\ 1$, \texttt{broadcast} to
$[\iota]$ (the vacuous instance of the trailing-axis rule),
\texttt{contract} binding the index --- and each instance evaluates
to $n$. Uncoloured counter places need nothing else: classical
place-invariant certificates are degree-one polynomials in count
wires, and count-bound rows ($0 \le k_P \le M_P$, linear
conservation laws) attach to this wire. The idiom is a natural
packaged symbol for the register of
section~\ref{sec:hyp-genus} --- expansion as above, projection lemma
trivial --- so certificates may cite one symbol while the wire
format never learns it exists.

\paragraph{The certificate stratum never sees a plate.} A plated
hypernet with invariant outputs, built from \texttt{contract} pools,
is a presentation of an aggregate polynomial in finitely many
\emph{pooled coordinates}: token counts $\vec{k}$ and per-place
power-sum pools $\vec{s}$ (a sum-pool is a contraction row whose
bound index is the plate --- the relational rendering of
section~\ref{sec:hyp-tensors} applies unchanged, side conditions
included). Obligations about all markings at once are ordinary
polynomial sequents over $(\vec{k}, \vec{s})$ with
hypothesis rows --- coupling enclosures
$l_j k_P \le s^P_j \le u_j k_P$ from the colour boxes, count rows,
integer rounding at count atoms --- so obligation count and
dimension are set by the net and the feature dictionary, never by
any marking bound. The plate joins ``the checker never sees a
probability'' and ``never sees a marking'': it compiles away through
the pooled normal form, and what reaches the kernel is
hypothesis-carrying polynomial rows of the standard shape. How the
per-transition rows of a token-population system collapse to one
row per Petri transition --- and which pooled aggregators can carry
progress claims, not merely safety --- is taken up with the
modelling frame itself (section~\ref{sec:mt-sdcpn}).

\paragraph{Worked size.} A one-layer transformer encoder with
degree-2 polynomial attention --- a published, empirically
competitive attention class~\cite{letourneau24padre}, of which the
plate-expressible fragment is precisely the exchangeable
subclass --- over a two-dimensional colour box, with residual
connection, invariant quadratic readout, and count wire, is
\emph{nineteen applications}, of which the count wire is three; its
pooled normal form is a per-token quadratic form against pooled
third moments, an aggregate polynomial over fifteen pool coordinates
plus the count. An executable exact-rational fixture pins the
transcription against a per-token symbolic computation at
$n \in \{1, 2, 3\}$, the collapse identity symbolically, and
invariance under all $n!$ permutations.

\paragraph{Status.} The type layer and the three theorems are
mechanized in the Lean catalogue
(\texttt{lean/qtsl}, modules \texttt{Qtsl/Plate/\{Syntax, TPlate,
Eval, TNat, TDesc, Probes\}}): symbolic shapes and their
instantiation homomorphism, the lifted per-box typing table with the
schematic well-formedness theorem across all ten boxes, a reference
evaluator with the diagonal-action equivariance theorem, the
unwritability of causal masking (stated as a refutation: the
one-application index-literal extension falsifies equivariance at
$n = 2$), descent to \texttt{Multiset}-indexed denotations, and a
symbolic count-wire evaluation theorem, under a three-mutant battery
at preregistered loci. The certificate-stratum row family this
section's pooled normal form calls for --- plate declarations,
per-token feature dictionaries, aggregate wires as contraction rows,
fold certificates with per-transition footprint rows, and the
idempotent-gap refusal --- is landed theory vocabulary, and the
producer emits it: the semiconductor tutorial
(chapter~\ref{ch:qtsl-minifab}, section~\ref{sec:minifab-exact}) is
its first exercising campaign, and the engine-side evaluator family
for the theory's demanded side conditions now discharges them in
that campaign's production loop, decided at the coloured
place/transition fragment with the remainder refused by name
(section~\ref{sec:minifab-exact-status}). The
hypernet \emph{wire-format} extent tag itself --- symbolic extents
on general wires, beyond the certificate rows --- remains design
with its pieces priced; plate-free serialization is unaffected by
construction (zero-plate graphs are byte-identical today).

\section{One genus: views elaborate in}
\label{sec:hyp-genus}

A toolkit that certifies computations cannot afford two kinds of
computation. The design rule entered the development log as a maxim
the day the System layer adopted hypernets for its dynamical content:
hypernets are to be re-used for
differential-equation specifications and their kin. The maxim had an
edge because the project had already lived through the failure once,
when terms grew into a separate genus from hypernets and were soon
demanding einsum machinery of their own.

The rule this chapter takes from that entry, stated positively:

\begin{quote}
\textbf{The single-genus rule.} Hypernets are the one genus of
computational content. Every other computational syntax --- a
polynomial vector field written as a coefficient table, a term inside
a proof obligation, a guard expression, a PDE operator --- is a
\emph{view}: surface syntax equipped with a definitional elaboration
into the ten boxes. A view's meaning \emph{is} the value of its
elaboration; no view acquires its own evaluator, its own
serialization, or its own contraction machinery.
\end{quote}

The precedent for the pattern is einsum itself. Einsum notation is
not a second contraction engine anywhere in the toolkit; it lowers
to explicit contraction structure (patterns here; contraction rows in
the relational engine, section~\ref{sec:hyp-tensors}), and the
summation it denotes is definitional --- evaluated only where a
checker needs values. Views generalize this: convenient syntax on
the outside, one calculus on the inside, and the translation is part
of the definition rather than an implementation detail.

The disease the rule prevents has a concrete signature, worth naming
because it is how single-genus erodes in practice: a producer grows
an ergonomic node --- say, a polynomial represented as its monomial
table --- that is not among the ten boxes; artefacts begin to carry
it \emph{instead of} applications; and some downstream consumer,
finding no applications to read, grows an evaluator against the
ergonomic node directly. At that moment there are two genera, two
evaluators, and a fresh supply of questions about whether they agree.
The repair --- performed on the polynomial node itself, and now the
standing template --- is \emph{demotion to view} under the
admissible-symbol discipline:

\begin{itemize}
\item an \textbf{admissible packaged symbol} is a named node with a
  typing rule, admitted alongside (never into) the core ten;
\item it carries a \textbf{definitional expansion} into core boxes
  --- the node \emph{serializes as} its expansion, so the wire format
  and the digest scheme never learn the symbol exists;
\item a once-proven \textbf{projection lemma} establishes that the
  symbol's declared facts follow from its expansion's facts, so rules
  may cite the symbol soundly;
\item and a \textbf{candidacy rule} bounds the admission: the
  expansion's disjunctive content must come only from \texttt{cell}
  (most packaged symbols are disjunction-free; a packaged ReLU
  inherits its single case split from the \texttt{cell} inside its
  expansion, and nothing else is allowed to smuggle new branching
  into the calculus).
\end{itemize}

Four packaged symbols currently ride this register, each with its
expansion and projection lemma mechanized: \texttt{relu} (expansion:
copy, sign-split \texttt{cell}, gating \texttt{contract});
\texttt{ppoly}, the monomial-table polynomial view whose expansion is
the scheme of the worked example below; \texttt{pderiv}, the formal
partial derivative on dense coefficient tables (expansion: one
constant shift-and-scale matrix and one \texttt{contract});
and \texttt{facet}, restriction of a coefficient table to a face
(expansion: one basis-vector constant and one \texttt{contract}).
None is an eleventh primitive; all four erase at serialization time.

Three audit questions probe any proposed addition, and the calculus
answers all three by construction: \emph{a second evaluator?} --- no; views elaborate, and
the one exact evaluator computes. \emph{A second digest scheme?} ---
no; the canonical serialization hashes expansions, so a view and its
elaboration cannot drift apart behind distinct content addresses.
\emph{A second contraction path?} --- no; the contraction pattern is
the only index-binding construct, and any syntax ``asking for its own
einsum'' is asking, by definition, to violate the rule.

\subsection*{Beyond exactness: certified function aliases}

One register remains, opened by a question the packaged symbols
cannot answer. Modelling frames keep asking for $\sin$, $\cos$,
$t^{p/q}$, $\ln$, and $\exp$ --- in drift coefficients, rate laws,
guard surfaces --- and no piecewise polynomial \emph{equals} any of
them, so no definitional expansion of the packaged-symbol kind can
exist. What can exist, and what exact numerics has shipped for
decades under the name of certified function
libraries~\cite{muller2016elementary}, is an
approximation with a proof: a piecewise-polynomial table $\hat f$
over exact rationals, a declared domain $[a, b]$, and a declared
budget $\varepsilon$, together with a receipt --- exact-rational
inequalities a small checker re-derives from the table bytes alone
--- establishing $\lvert \hat f - f \rvert \le \varepsilon$ on the
domain. Remainder brackets of the Maclaurin kind make such receipts
pure rational arithmetic (the checker compares coefficients against
a factorial tail bound; it never needs to know that $e$ exists), and
implicit relations extend the same discipline to functions the
brackets do not reach directly: a logarithm table is certified
through the residual $\lvert \exp(\hat w(t)) - t \rvert \le
\eta\, t$ against exponential brackets, a rational-power table
through the purely polynomial residual $\lvert \hat w^{\,q} - t^p
\rvert \le \eta$ --- so no logarithm and no fractional power is ever
enclosed by anything but polynomial comparison.

A \emph{certified function alias} packages exactly this: a name and
pinned parameters selecting one table, hence one hypernet (the
pieces enter through \texttt{cell}, the piece polynomials through
the \texttt{ppoly} scheme, the one-hot blend through
\texttt{contract}); the budget $(\varepsilon; [a,b])$; and the
receipt's citation. The alias is admitted alongside the calculus in
the same sense as the packaged symbols --- the wire format and the
digest scheme never learn its name; artefacts carry the expansion
--- and the three audit questions above pass by the same erasure
argument. What changes is the epistemic contract: a packaged
symbol's projection lemma says \emph{equal}; an alias's receipt says
\emph{within $\varepsilon$, on $[a,b]$} --- and the domain is a
precondition, owed a containment argument by any consumer whose
wires might leave it, the same obligation shape as a jump's
reset containment. Downstream, the distinction costs nothing: a
drift built from an alias \emph{is} piecewise-polynomial dataflow,
its proof obligations are polynomial rows like any other, and the
transcendental fact lives entirely in the receipt. A claim over such
a system may then be read at either of two strengths: against the
table as the model itself, with $\varepsilon$ a declared modelling
tolerance of the same kind as float rounding (the reading under which
nothing downstream changes at all); or against the band $\hat f \pm
\varepsilon$ as a differential inclusion, whose verdict covers the
true function by containment at the price of one rational widening
term per obligation row.

The family is closed under differentiation --- $\sin' = \cos$,
$\exp' = \exp$, $\ln' = t^{-1}$, $(t^{p/q})' = \tfrac{p}{q}\,
t^{p/q - 1}$ --- so an alias's derivative data is a pointer to a
sibling alias, never new machinery. The cylinder lift met earlier in
this chapter remains the better encoding where it applies (an
autonomous oscillator's sine and cosine carried as state, exactly,
with their algebraic invariant); the alias route serves the cases
the lift cannot reach --- forced and parametric arguments,
logarithms and powers of state variables, rate laws. The toolkit
carries certified tables for the whole family named above. The
exponential, the logarithm, and the rational powers arrived first,
as quantile transforms for its distribution generators, verified by
the residual disciplines above; sine and cosine followed as forward
tables on a declared symmetric domain, their error bounded
absolutely against the same factorial-tail brackets the remainder
arguments above consume --- absolutely of necessity, since both
functions vanish inside every interesting domain, leaving a relative
bound no positive floor to stand on. One more fact about the
trigonometric pair is worth the ink: $\pi$ is irrational, so
periodic folding is not exactly available in $\mathbb{Q}$, and a
range-reduced table would smuggle an approximation of $\pi$ into its
own argument, outside the certificate. The declared domain is
therefore data; coverage beyond it is refused, never clamped, and
range reduction is the consumer's burden, priced where the consumer
knows its own modulus.

\subsection*{Worked example: a closed-loop drift, in fifteen applications}

The example that carried the design discussion is small enough to
print whole. Take the cartpole's closed loop on the cylinder lift
$(s, c, \omega)$ --- sine and cosine of the pole angle carried as
state, per the no-transcendentals discipline --- with a linear
feedback law $a(s, \omega) = K_0 s + K_1 \omega$ folded in, plant
half-length $l = \tfrac12$, gravity $g = \tfrac{49}{5}$, and exact
rational gains $K_0 = -\tfrac{161753}{8148}$,
$K_1 = -\tfrac{44746}{9799}$:
\[
\dot s = c\,\omega, \qquad
\dot c = -s\,\omega, \qquad
\dot\omega
= \tfrac{g}{l}\,s + \tfrac{K_0}{l}\,c s + \tfrac{K_1}{l}\,c\,\omega .
\]
Its hypernet --- three input wires, three output wires, every scalar a
rank-0 exact rational --- is fifteen applications:

\begin{center}
\begin{tabular}{@{}lll@{}}
\toprule
stage & applications & count \\
\midrule
fan-out & \texttt{copy} each of $s$, $c$, $\omega$ into its uses & 3 \\
constants & \texttt{const} $-1$, $\tfrac{g}{l}$, $\tfrac{K_0}{l}$, $\tfrac{K_1}{l}$, $\mathbf{1}_3$ & 5 \\
monomials & \texttt{contract}: $c\omega$; $(-1)s\omega$; $\tfrac{g}{l}s$; $\tfrac{K_0}{l}cs$; $\tfrac{K_1}{l}c\omega$ & 5 \\
the $\Sigma\Pi$ boundary & \texttt{stack} the three $\dot\omega$ monomials into shape $[3]$ & 1 \\
the sum & \texttt{contract} against $\mathbf{1}_3$, binding the $[3]$ axis & 1 \\
\bottomrule
\end{tabular}
\end{center}

\begin{figure}[t]
\centering
\begin{tikzpicture}[x=0.94cm,y=0.82cm,
  box/.style={draw=accentdeep, fill=accent!50, rounded corners=1pt,
    font=\footnotesize\ttfamily, inner sep=2.5pt, minimum height=4.2mm},
  cbox/.style={draw=NoteGray, fill=white, rounded corners=1pt,
    font=\footnotesize, inner sep=2pt, minimum height=3.6mm},
  pbox/.style={draw=NoteGray, fill=NoteGray!12, rounded corners=1pt,
    font=\footnotesize\ttfamily, inner sep=2.5pt, minimum height=4.2mm},
  wire/.style={-, line width=0.5pt, draw=black!75},
  wlab/.style={font=\scriptsize, inner sep=1pt},
  port/.style={font=\small\itshape, inner sep=1pt}]
\node[port] (in-0) at (0.00, 0.00) {$s$};
\node[port] (in-1) at (0.00, -1.50) {$c$};
\node[port] (in-2) at (0.00, 1.50) {$\omega$};
\node[pbox] (a0) at (1.60, 0.00) {copy$_{3}$};
\node[pbox] (a1) at (1.60, -1.50) {copy$_{3}$};
\node[pbox] (a2) at (1.60, 1.50) {copy$_{3}$};
\node[box] (a3) at (4.80, 3.00) {contract};
\node[cbox] (a4) at (4.80, 2.25) {$-1$};
\node[box] (a5) at (4.80, 1.50) {contract};
\node[cbox] (a6) at (4.80, -2.25) {$K_1/l$};
\node[box] (a7) at (4.80, -3.00) {contract};
\node[cbox] (a8) at (4.80, 0.75) {$g/l$};
\node[box] (a9) at (4.80, 0.00) {contract};
\node[cbox] (a10) at (4.80, -0.75) {$K_0/l$};
\node[box] (a11) at (4.80, -1.50) {contract};
\node[box] (a12) at (8.20, -1.50) {stack$_{3}$};
\node[cbox] (a13) at (11.60, -0.75) {$\mathbf{1}_{3}$};
\node[box] (a14) at (11.60, -1.50) {contract};
\node[port] (out-0) at (13.23, 1.50) {$\dot s$};
\node[port] (out-1) at (13.23, 0.00) {$\dot c$};
\node[port] (out-2) at (13.23, -1.50) {$\dot\omega$};
\draw[wire] (in-0) -- (a0.west);
\draw[wire] (in-1) -- (a1.west);
\draw[wire] (in-2) -- (a2.west);
\draw[wire] ([yshift=1.50mm]a0.east) .. controls (3.04, 0.00) and (3.36, 1.50) .. ([yshift=-1.20mm]a5.west);
\draw[wire] (a0.east) -- (a9.west);
\draw[wire] ([yshift=-1.50mm]a0.east) .. controls (3.04, 0.00) and (3.36, -1.50) .. ([yshift=1.20mm]a11.west);
\draw[wire] ([yshift=1.50mm]a1.east) .. controls (3.04, -1.50) and (3.36, 3.00) .. ([yshift=-1.20mm]a3.west);
\draw[wire] ([yshift=-1.50mm]a1.east) .. controls (3.04, -1.50) and (3.36, -3.00) .. ([yshift=-1.20mm]a7.west);
\draw[wire] (a1.east) -- ([yshift=-1.20mm]a11.west);
\draw[wire] ([yshift=1.50mm]a2.east) .. controls (3.04, 1.50) and (3.36, 3.00) .. ([yshift=1.20mm]a3.west);
\draw[wire] (a2.east) -- ([yshift=1.20mm]a5.west);
\draw[wire] ([yshift=-1.50mm]a2.east) .. controls (3.04, 1.50) and (3.36, -3.00) .. ([yshift=1.20mm]a7.west);
\draw[wire] (a3.east) .. controls (6.33, 3.00) and (6.67, 1.50) .. node[wlab, above, pos=0.35] {$c\omega$} (8.20, 1.50) -- (11.60, 1.50) -- (out-0);
\draw[wire] (a5.east) .. controls (6.33, 1.50) and (6.67, 0.00) .. node[wlab, above, pos=0.35] {$-s\omega$} (8.20, 0.00) -- (11.60, 0.00) -- (out-1);
\draw[wire] (a7.east) .. controls (6.33, -3.00) and (6.67, -1.50) .. node[wlab, above, pos=0.35] {$\tfrac{K_1}{l}c\omega$} ([yshift=-1.50mm]a12.west);
\draw[wire] (a9.east) .. controls (6.33, 0.00) and (6.67, -1.50) .. node[wlab, above, pos=0.35] {$\tfrac{g}{l}s$} ([yshift=1.50mm]a12.west);
\draw[wire] (a11.east) -- node[wlab, above, pos=0.35] {$\tfrac{K_0}{l}cs$} (a12.west);
\draw[wire] (a12.east) -- node[wlab, above, pos=0.35] {$[3]$} (a14.west);
\draw[wire] (a14.east) -- (out-2);
\draw[wire] (a4) -- (a5.north);
\draw[wire] (a6) -- (a7.north);
\draw[wire] (a8) -- (a9.north);
\draw[wire] (a10) -- (a11.north);
\draw[wire] (a13) -- (a14.north);
\draw[dashed, draw=accentmid, line width=0.6pt] (7.11, -3.60) -- (7.11, 3.45);
\node[font=\small, text=accentmid] at (6.76, 3.70) {$\Pi$};
\node[font=\small, text=accentmid] at (7.46, 3.70) {$\Sigma$};
\end{tikzpicture}
 \caption{The closed-loop cartpole drift of the worked example,
depicted as a hypernet: fifteen applications, three inputs, three
outputs.
Structural plumbing (\texttt{copy}, grey) makes every act of sharing
a visible row; constants (white) are applications like everything
else; every monomial is one \texttt{contract} (peach); and the
calculus's single cross-wire sum --- the three-term right-hand side
of $\dot\omega$ --- crosses the $\Sigma\Pi$ boundary (dashed) as one
\texttt{stack} contracted against the all-ones vector.}
\label{fig:hyp-cartpole}
\end{figure}

Every monomial is contract-native; the single cross-wire sum in the
system --- the three-term right-hand side of $\dot\omega$ --- is the
single \texttt{stack} (Figure~\ref{fig:hyp-cartpole}). The table is
not hand-lore: the polynomial
view elaborator emits this graph, its fifteen applications digest to
a pinned golden
(\texttt{69b6f4c6\ldots}), and the elaborator's output is already in
canonical form, so the digest doubles as a regression witness for the
elaboration itself. Neither is the figure hand-lore: it is emitted,
deterministically, from those same canonical bytes by the depiction
tool beside the producer surface (\texttt{hypernet-tikz}), so the
pinned digest is the drawing's witness too --- only the semantic
names ($s$, $c$, $\omega$, $g/l$) arrive as declared labels, since
the canonical form deliberately carries no name metadata. The same scheme compiles the goal polynomial of
the matching Lyapunov certificate --- one more \texttt{stack}, one
more ones-\texttt{contract} --- which is the single-genus rule doing
exactly its job: system dynamics and certificate content share one
representation, so one recomputation can cross-check a digest cited
by both.

\section{Substitution: how hypernets compose}
\label{sec:hyp-subst}

Hypernets compose by \emph{substitution}: plugging one graph's
outputs into another's inputs. The operation is defined once, at the
graph layer, and everything that composes computations --- nesting a
component's dynamics inside a larger mechanism, elaborating a view in
context, merging two boxed computations into one --- consumes the
same definition.

The data of a substitution: an outer hypernet, an inner hypernet, a
\emph{port map} $\varphi$ matching some of the outer graph's inputs
to inner outputs (type-exact; several outer inputs may consume one
inner output), and an \emph{exposure set} of inner outputs to remain
outputs of the composite --- by default the unmatched ones, so
nothing is silently dropped and nothing matched is silently
re-exposed. The construction concatenates the inner application
list, a block of \emph{seam repairs}, and the redirected outer list,
then renumbers into canonical form. The seam repairs are where the
linearity discipline pays: an inner output consumed by $n$ ports and
exposures gains exactly one consumer at the seam --- a
\texttt{delete} when $n = 0$ (the discard is a row, never an
absence), nothing when $n = 1$, a \texttt{copy} when $n \ge 2$ ---
and consumer counts change nowhere else.

\begin{theorem}[substitution laws]
\label{thm:substitution}
For well-formed hypernets and valid port structures:
\emph{(i)} the composite is well-formed;
\emph{(ii)} on the deterministic fragment, evaluation of the
composite equals the composition of evaluations --- inner outputs
feeding matched outer inputs, unmatched interface passing through;
\emph{(iii)} substitution is unital (substituting an identity
hypernet, in either direction, returns the other graph up to
canonical renumbering) and associative --- and both laws hold
\emph{on the nose} for the canonical forms, hence as equalities of
digests, not merely up to isomorphism.
\end{theorem}

The proofs are elementary but not careless; the one choice that
genuinely matters is \emph{where the repair blocks live}. Placing
each seam's repairs at that seam --- between the inner and outer
blocks it mediates --- makes the two association orders of a
three-graph chain produce byte-identical canonical forms; the
tempting alternative (all repairs at the end) satisfies (i) and (ii)
equally well and breaks (iii) at the digest layer, leaving
associativity only up to a permutation that content-addressing does
not forgive. Coherence at the byte level is a design property, chosen
because certificates compare bytes.

Clause~(ii) as stated is a floor, and the machine-checked form is
stronger. In the toolkit's Lean tree, composition at the aligned
stratum --- the pushout-normal form in which the seam wires of two
composable open hypernets are literally shared names, the form every
composite of content-addressed hypernets reduces to --- carries the
law family with no determinism hypothesis anywhere (module
\texttt{Qtsl.HypernetCospan}): gluing preserves well-formedness
wholesale --- linearity, acyclicity, type consistency, and every
boundary clause (\texttt{glue}); a valuation of the composite is
admissible exactly when every application of each part holds
(\texttt{glue\_admissible\_iff}); and unit and associativity hold on
the nose (\texttt{glue\_unit\_left}, \texttt{glue\_unit\_right},
\texttt{glue\_assoc}) --- clause~(iii)'s shadow in pushout-normal
form.

The relational law subsumes the deterministic clause: where the
parts' admissible relations are graphs of functions, ``every
application holds'' is exactly ``evaluation of the composite equals
the composition of evaluations'' --- and \texttt{cell}'s one-hot
disjunction, the calculus's one nondeterministic citizen, passes
through the conjunction untouched. The composition law never needed
determinism; clause~(ii) names the fragment where ``evaluation'' is
a function because that is the fragment the evaluator runs.

Read through that normal form, the seam repairs acquire a precise
status: they are a \emph{change of interface}, never composition
overhead. The bridge between the two presentations is itself
machine-checked (module \texttt{Qtsl.SubstitutionBridge}):
substitution factors as a gluing against a small structural hypernet
--- one \texttt{delete} or \texttt{copy} row per seam wire whose
consumer count changes --- that mediates between the inner graph's
interface and the interface the port map asks for, and the
seam-repair block is that mediator's application list, verbatim.
Gluing along an interface both sides already agree on inserts
nothing; the rows appear exactly when the interface moves.

For certificate consumers, the digest story splits into three claims
of different strengths, and keeping them separate is what makes
``digest coherence'' a tool rather than folklore:

\begin{remark}[digest verdicts]
\label{rem:digest-verdicts}
\emph{(a) Determination --- a theorem.} The composite's digest is a
function of the parts' canonical bytes and the port structure: the
operation reads nothing else. Equal-digest parts yield an
equal-digest composite, so a certificate may cite the parts and the
port row and thereby pin the whole.
\emph{(b) Coherence --- a theorem, conditional on a discipline.} The
unit and associativity equalities of
Theorem~\ref{thm:substitution}(iii) are digest equalities
\emph{because} normalization is part of the operation; producers must
emit canonical form or normalize before digesting, and the elaborators
do (their outputs are fixed points of renumbering, checked).
\emph{(c) Homomorphism --- refused.} There is no operation on digest
\emph{strings} computing the composite's digest from the parts'
digests: SHA-256 composes with nothing, by design. The checker's path
is recomputation --- possess the parts' bytes, recompute the
substitution, compare the claimed digest --- and every consumer in
the toolkit uses that pattern rather than assuming an algebra of
hashes.
\end{remark}

One boundary deserves stating before anything leans on it:
elaboration and substitution commute \emph{semantically} (elaborating
a polynomial composite equals substituting the elaborations, as
functions; proven, with executable witnesses), but the two routes
produce \emph{different graphs}, and different graphs never share a
digest. Expanding a polynomial before elaborating re-associates its
monomials; substituting factored graphs keeps the factoring. A
certificate relating the two routes therefore carries a semantic
linkage argument, never a claimed hash equality --- the digest
refuses extensionality on purpose, and this is the first place the
refusal does real work.

\primentry{substitution}{H-SUBST}
\primmeta
\begin{mathpar}
\inferrule*[Right=subst]
  {G, H \text{ well-formed} \\\\
   \varphi \colon \mathit{in}(G) \rightharpoonup \mathit{out}(H)
   \text{ type-exact} \\
   E \subseteq \mathit{out}(H)
   \text{ (default: the unmatched outputs)}}
  {\mathrm{subst}(G, \varphi, H, E) \colon
   \mathit{in}(H) \mathbin{+} \bigl(\mathit{in}(G) \setminus
   \operatorname{dom}\varphi\bigr) \to \mathit{out}(G) \mathbin{+} E}
\end{mathpar}
\noindent Semantics and laws: Theorem~\ref{thm:substitution} ---
composition of evaluations on the deterministic fragment; unit and
associativity on the nose, as digest equalities. The seam repair is
forced by linearity: an inner output with $n$ consumers at the seam
gains a \texttt{delete} ($n = 0$), nothing ($n = 1$), or a
\texttt{copy} ($n \ge 2$), and consumer counts change nowhere else.
The operation normalizes its result (the canonical renumbering); the
identity hypernet is its unit.
\ruleblock{API}
\sourcenote{Producer API of the \texttt{mtk} toolkit --- the TypeScript arrangement and the Rust catalogue crate, one implementation underneath. The extraction, corpus, and receipt disciplines governing this block are stated once, at \hyperref[prim:H-WF]{this catalogue's first entry}.}
The normal form (\texttt{renumber}) and the unit (the identity hypernet) are quoted alongside \texttt{substitute}. The corpus construction exercises a fan-out seam, a delete seam, an explicit exposure, and a pass-through port; the composite's bytes are verified against the reference operation's output.
\apilang{TypeScript}{hypernet.substitute}
\apiprov{signature --- \texttt{ts/src/hypernet.ts}}
\begin{tssrc}
export function substitute(
  outer: Hypernet,
  portMap:
    | ReadonlyMap<Wire, Wire>
    | ReadonlyArray<readonly [outerInput: Wire, innerOutput: Wire]>,
  inner: Hypernet,
  opts?: {
    expose?: ReadonlyArray<Wire>;
    normalize?: boolean;
  },
): Result<Hypernet, SubstRefusal>

export function renumber(h: Hypernet): Result<Hypernet, SubstRefusal>

export function identity(types: ReadonlyArray<TensorType>): Hypernet
\end{tssrc}
\apiprov{construction --- \texttt{ts/tests/hypernet-native.test.ts}}
\begin{tssrc}
const comp = hn.unwrap(
  hn.substitute(
    outer,
    [
      [u, os[0]!],
      [v, os[0]!],
    ],
    inner,
    { expose: [os[1]!] },
  ),
);
\end{tssrc}
\apilang{Rust}{mtk\_catalogue::hypernet::substitute}
\apiprov{signature --- \texttt{crates/mtk-catalogue/src/hypernet/subst.rs}}
\begin{rustsrc}
pub fn substitute(
    outer: &Hypernet,
    port_map: &BTreeMap<Wire, Wire>,
    inner: &Hypernet,
    expose: Option<&BTreeSet<Wire>>,
    normalize: bool,
) -> Result<Hypernet, SubstRefusal>

pub fn renumber(h: &Hypernet) -> Result<Hypernet, SubstRefusal>

pub fn identity_hypernet(types: &[TensorType]) -> Hypernet
\end{rustsrc}
\apiprov{construction --- \texttt{crates/mtk-catalogue/tests/golden\_hypernet.rs}}
\begin{rustsrc}
let port_map: BTreeMap<u32, u32> = [(u, os[0]), (v, os[0])].into();
let expose: BTreeSet<u32> = [os[1]].into();
let comp = substitute(&outer, &port_map, &inner, Some(&expose), true).unwrap();
\end{rustsrc}

\section{The shared core: Systems, specifications, rules}
\label{sec:hyp-core}

The claim this chapter exists to underwrite: the three strata that
consume computation all consume \emph{this} genus, so the chapter's
material is common infrastructure, not one stratum's private
concern.

\subsection*{Systems: mechanism content is boxed hypernets}

A System (chapter~\ref{ch:qtsl-mathematical-models}) is a structured
presentation --- mechanisms wired together --- and the mechanism
calculus itself deliberately contains \emph{no computational
primitives at all}. Its vocabulary identifies variables, grades
stochastic couplings, sequences operating modes, nests components; none
of these operations computes a value. Wherever a mechanism has
dynamical content --- a drift field, a jump reset, a guard, a neural
feedback law --- that content is a hypernet, referenced by digest.

The embedding has a precise grammar: a well-formed hypernet with $m$
inputs and $n$ outputs becomes a mechanism whose ports are its wire
types and whose causal orientation is fixed (inputs consumed, outputs
read out) --- a \emph{deterministic directional mechanism}, sitting
inside the System layer as what string-diagram practice calls a
functorial box \cite{mellies06csl}: a membrane with hypernet
calculus inside and mechanism calculus outside. Composition respects
the membrane --- merging two boxes joined by a directed wire is
computed by the substitution operation of
section~\ref{sec:hyp-subst}, inside --- and the membrane is
selective in both directions. Relational variable-identification,
graded stochastic couplings, and mode-switching edges cannot cross
into the box: they are not computations. Acyclicity ends at the
membrane in the other direction: inside, SSA order is law; outside,
feedback loops are the System layer's bread and butter, closed
through state across box boundaries. A control loop is a cycle of
mechanisms whose every computational fragment is internally acyclic
--- which is how the toolkit gets both certified computation and
real feedback without either discipline compromising the other. (In
the vocabulary of Theorem~\ref{thm:hyp-inclusion} below: a
functorial box packages a well-formed, deterministic loose arrow of
the double category of open hypernets, and the membrane is that
fragment's boundary.)

\subsection*{Specifications: one digest, one referent}

A specification constrains observables of a trajectory; a certificate
discharges the constraint with computational content of its own ---
a Lyapunov function's polynomial, a barrier's gradient bounds, the
observable defining a guard surface. Both sides of that transaction
are hypernets, and both cite digests under the one canonical
serialization. The consequence is a sentence this book will rely on
repeatedly, in both directions: \emph{the guard the certificate
reasons about is the guard the system runs}, and the checker
verifying the certificate can confirm as much by recomputing a hash
--- no semantic argument, no trusted translation step, no ``the spec
model and the implementation model surely agree.'' Where a
certificate's content is produced by elaboration from a view, digest
determination (Remark~\ref{rem:digest-verdicts}a) extends the same
guarantee through the elaboration: cite the view's parameters and the
elaborator's version, or cite the elaborated digest directly; either
way the referent is pinned by content.

\subsection*{Rules: the ambient hypernet}

Every inference rule in the catalogue of chapter~\ref{ch:qtsl-manual}
is stated over an ambient hypernet: the rule's variables are wires,
its premises and conclusions constrain the values those wires can
take, and its soundness theorem quantifies over \emph{every}
hypernet and every admissible valuation. The primitive boxes supply
the base facts --- each application contributes its defining
equations (a \texttt{contract}'s sum-of-products identity, a
\texttt{stack}'s entry equalities, a \texttt{cell}'s one-hot
disjunction) as the synthesized facts rules build on. The economy of
section~\ref{sec:hyp-boxes} becomes a proof-theoretic economy here:
nine box families contribute equations, exactly one contributes a
disjunction, so case analysis in the rule layer is confined to
\texttt{cell}-shaped splits, and a packaged symbol admitted under the
candidacy rule of section~\ref{sec:hyp-genus} cannot enlarge that
confinement.

\medskip

This triple consumption is the chapter's placement argument. Systems,
specifications, and rules each depend on the hypernet genus; the
genus depends on none of them. Read this chapter once, and the later
chapters' references --- ``dynamical content, by digest,''
``the ambient hypernet,'' ``the certificate cites the graph'' ---
resolve to one object with one definition.

\medskip

The placement argument also has a mathematical face, and the face is
machine-checked.

\begin{theorem}[possibilistic inclusion]
\label{thm:hyp-inclusion}
Open hypernets --- hypernets with declared input and output rosters
--- form a double category: objects are wire rosters, tight arrows
are roster inclusions, loose arrows are hypernet-decorated cospans
in the sense of Fong~\cite{fong15}, read at double-categorical
strength through Patterson's lens~\cite{patterson23}, and loose
composition is gluing at the shared seam. Sending an open hypernet
to its possibilistic content --- the set of valuations its
applications admit --- extends to a lax double functor from this
double category into the semantic universe of
chapter~\ref{ch:qtsl-semantic-universe}, and the functor:
\begin{enumerate}
\item satisfies its unit, composition, and rename-transport laws as
  exact equalities, with no determinism hypothesis anywhere;
\item preserves the canonical opcartesian lifts on the nose;
\item has its composition comparison invertible precisely on
  seam-disjoint gluings --- which is every gluing of encoded
  hypernets, since content addressing makes wire names rigid
  designators: distinct off-seam wires never collide, and no
  capture-avoiding renaming ever intervenes;
\item identifies two hypernets precisely when they admit the same
  valuations.
\end{enumerate}
\end{theorem}

Every clause is a named artefact in the toolkit's Lean tree. The
packaged functor is \texttt{hnetInclusion} (module
\texttt{SuFunctor.HypernetDouble}), built from four transformation
laws proved exactly, one by one, in
\texttt{SuFunctor.HypernetInclusion}: unit, composition with no
determinism hypothesis (section~\ref{sec:hyp-subst}'s relational law
is this clause read at the substitution seam), rename transport at
full generality, and the vacuous-diagram law (\texttt{theta\_unit},
\texttt{theta\_comp}, \texttt{theta\_rename},
\texttt{vacDnf\_dnfPull}). Lift preservation is
\texttt{hnetInclusion\_pushGLoose}. The laxity is located, not
hidden: clause~3 is \texttt{hnetInclusion\_\allowbreak laxComp\_\allowbreak isInvertible},
while the unit comparison is genuinely one-way --- typing lives in
the pattern carrier, and forgetting it is a constraint, not a
defect. Clause~4 (\texttt{austereDnf\_eq\_iff}) says the inclusion
is deliberately non-faithful, and makes its kernel a theorem rather
than a folklore caveat: it is the extensionality refusal of
section~\ref{sec:hyp-ir} standing exactly where it belongs --- two
graphs computing one relation share a meaning and never a digest;
syntax and its content addressing live in the double category, and
the semantic universe sees the relation.

What Systems, specifications, and rules consume is therefore not
merely one data structure but one semantics: composing graphs and
then taking meaning equals taking meaning and then composing
relations.

\section{Relational kinship: contraction rows}
\label{sec:hyp-tensors}

The toolkit meets tensor contraction twice: as the \texttt{contract}
box here, and as \emph{contraction rows} in the relational engine,
where an einsum expression lowers to stored index structure --- an
index map and an output set as data --- with the summation
definitional, evaluated only where a checker needs values. The two
are one definition at two strata, and it is worth seeing the shared
skeleton once. Over a semiring $\mathbb{S}$, with a finite family of
input axes $A(i)$ and a map $a$ from input-axis positions into a
finite index set, contraction along $a$ with output set $O$ sends a
(lazily formed) product tensor to
\[
\Bigl(v \colon \textstyle\prod_{m \in O} B(m)\Bigr) \;\longmapsto
\sum_{v' \colon \prod_{j \notin O} B(j)}\;
T\bigl(i \mapsto (v \cup v')(a(i))\bigr),
\]
where the extent family $B$ is induced on the image of $a$ by the
input extents, and is data only on the output axes no input hits
(the broadcast axes, of which more below).
Definition~\ref{def:contraction-pattern} is this formula with
$\mathbb{S} = \mathbb{Q}$, the index set $I$ playing the target of
$a$, the input bindings playing $a$, and the output binding
enumerating $O$; the box evaluates it eagerly on wire values, the
relational rendering stores it as rows and lets the checker fold the
sum on demand. The design move is also shared: \emph{einsum lowers
to contraction rows} in the engine for the same reason \emph{views
elaborate into hypernets} here --- surface notation is never given
its own machinery.

Two deltas between the strata, stated so the kinship is not
overclaimed. First, the relational side imposes a side condition the
wire format does not: every index hit by no input must belong to the
output set, so the case that would silently scale a result by an
extent is \emph{unrepresentable} in rows, while in the IR the same
case is representable and simply carries its summed semantics ---
deterministic, if rarely what anyone meant; einsum lowerings never
emit it on either side. Second, broadcasting: the relational
rendering folds it into contraction itself (output axes hit by no
input are the broadcast axes, along which the result is constant),
while the IR keeps a separate \texttt{broadcast} box --- in a
dataflow graph the degenerate contraction would cost a pattern where
a relabelling suffices. Same semantics; different economies; one
calculus.

\section{Receipts}
\label{sec:hyp-receipts}

What stands behind this chapter, sorted by strength of evidence, per
the claim discipline used throughout the book.

\textbf{Machine-checked and running.} The ten-primitive calculus
exists twice, in agreement: as a mechanized catalogue (the primitive
type, per-box typing rules, and the well-formedness trio as decidable
predicates over ground data, with well-formedness carried as proof
fields of the hypernet structure) and as a reference implementation
(the same typing and well-formedness checks, the exact-$\mathbb{Q}$
evaluator, a float64 shadow evaluator carrying a conservative
accumulated-rounding bound, and the canonical serialization with
SHA-256 digests). An ONNX importer compiles all 45 networks of the
ACAS~Xu collision-avoidance benchmark~\cite{katz2017reluplex}
(architecture 5--50$\times$6--5, ReLU) to well-formed hypernets --- 2{,}172
applications and 2{,}767 wires per network, exact dyadic weights ---
with committed digest goldens (45 networks $\times$ 3 head variants,
all 135 distinct) and an oracle gate: on 47 published verification
instances, exact and shadow evaluation agree within the carried
bound (largest observed gap $1.9 \times 10^{-16}$), and
verdicts agree with an independent floating-point runtime 47/47. A
property battery (well-formedness of random SSA slices; digest
sensitivity --- a one-ulp weight corruption flips the digest, a clean
recompile reproduces the golden; evaluator agreement against an
independent forward pass) runs green in seconds. The substitution
operation and canonical renumbering are implemented, with the
coherence and digest laws of section~\ref{sec:hyp-subst} each
exercised by executable digest-equality witnesses (among them the
deliberate digest-\emph{inequality} of the two elaboration
routes). The polynomial-view elaborator is pinned to the
fifteen-application worked table by its committed digest golden, and
the four packaged symbols' expansions live in the mechanized
catalogue with their projection lemmas. The composition layer above
the catalogue is machine-checked in the same tree: the
aligned-stratum gluing with its relational law and on-the-nose unit
and associativity (\texttt{Qtsl.HypernetCospan}), the reduction of
the port-map substitution to that gluing
(\texttt{Qtsl.SubstitutionBridge}), the four
transformation laws of the possibilistic inclusion
(\texttt{SuFunctor.HypernetInclusion}), and their packaging as a lax
double functor into the semantic universe
(\texttt{SuFunctor.HypernetDouble}) --- sorry-free, axiom census the
standard classical base. The plate extension
(section~\ref{sec:hyp-plates}) is machine-checked at the same grade
in \texttt{Qtsl/Plate/\{Syntax, TPlate, Eval, TNat, TDesc,
Probes\}}: symbolic shapes and instantiation, schematic
well-formedness across all ten boxes, evaluator equivariance under
the diagonal symmetric-group action, the causal-mask refutation,
descent to multiset-indexed denotations, and the symbolic count-wire
theorem, under a three-mutant battery at preregistered loci
(index-literal box, widened broadcast slot, deleted contract-binding
clause --- each refused at its predicted theorem); an exact-rational
fixture additionally pins the nineteen-application polynomial
transformer against per-token symbolic computation at
$n \in \{1,2,3\}$.

\textbf{Proven in prose, with executable witnesses; mechanization
under way.} Proposition~\ref{prop:packing-gap} (the completion
theorem behind \texttt{stack}) and Theorem~\ref{thm:substitution}
(well-formedness preservation, the merge law, unit and associativity
on the nose) are proven at full detail in the development record's
design notes, and every claimed equality is additionally checked by
the executable witnesses above. The proposition is, further,
machine-checked in the toolkit's Lean tree (\texttt{lean/qtsl/},
module \texttt{Qtsl.PackingGap}): a sorry-free Lean development, its axiom
census the standard classical base, mechanizes it over the formal-expansion
support model the proof sketch itself runs on --- assuming no
well-formedness, quantifying over all wires rather than only later
ones, and holding the two tracked wires opaque whatever produced
them --- together with a positive control deriving
the forbidden configuration once \texttt{stack} is added, so the
nine-box hypothesis is itself proven necessary. Of
Theorem~\ref{thm:substitution} in the port-map presentation given
here, the linearity clause of well-formedness preservation is
machine-checked end to end, and the laws themselves are
machine-checked in pushout-normal form (the aligned-stratum gluing
of section~\ref{sec:hyp-subst}: well-formedness closure, the
relational composition law with the determinism hypothesis deleted,
unit and associativity on the nose). The bridge --- that the
port-map, seam-repair, and renumbering construction of this chapter
reduces to that normal form --- is machine-checked in the same tree
(\texttt{Qtsl.SubstitutionBridge}): under the port-validity
hypotheses and the no-taps interface discipline, the raw composite
is field-for-field the aligned gluing of the padded inner hypernet,
the seam-repair block taken verbatim as the structural interpolant,
and the redirected outer hypernet
(\texttt{substituteRaw\_eq\_glue}), so substitution is the canonical
renumbering of an aligned-stratum glue
(\texttt{substitute\_eq\_glue\_renumber}) --- the structural
reshuffle of the prose statement absorbed into the choice of
interface lists --- and the relational composition law applies to
the substitution itself
(\texttt{substituteOpen\_admissible\_iff}); sorry-free, axiom census
the standard classical base.  The presentation of this chapter and
the normal form it reduces to are thereby machine-checked as one
construction.

\textbf{Design tense.} Higher-order hypernets (function-space wires
--- wanted by no modelling frame currently in scope, and deliberately
unbuilt until one forces them); the block-ordering convention that
would sharpen the commutation of substitutions into \emph{disjoint}
port sets from convention-conditional to unconditional; and the
plate producer surface (the wire-format extent tag with its
backward-compatibility golden, builder support, and the packaged
count symbol --- the pieces named in
section~\ref{sec:hyp-plates}'s status paragraph). All have written
design records; nothing stated above depends on any of them.
 \chapter{Hypernets: Markov Kernels}
\label{ch:qtsl-markov-kernels}

\noindent\draftmark\medskip

\Bookref{ch:qtsl-hypernets}{The hypernet chapter} fixed a
deterministic calculus: ten primitive boxes, one evaluator, one
canonical serialization, one content digest. Stochastic models need
kernels: a diffusion with jumps needs a jump law, a stochastic
Petri-style transition needs a firing law, and a Markov chain simply
\emph{is} its transition kernel. This chapter is about how the
toolkit says ``Markov kernel'' without adding a primitive. A
\emph{Markov-kernel hypernet} is an ordinary well-formed hypernet
together with a designated roster of \emph{noise ports}: inputs read as
fresh uniform draws on $[0,1]$. Its denotation is the Markov kernel
obtained by pushing the uniform reference measure forward through the
deterministic semantics --- every kernel is presented by an explicit
deterministic dilation with its ancilla exposed. Randomness enters at
the interface, never at a box; the wire format, the evaluator, and the
digest scheme never learn that the extension exists. A reader who only
wants to know which named distributions the toolkit provides can go
straight to the coverage dashboard (section~\ref{sec:exp-dist}).

\section{Kernels as dilations}
\label{sec:mk-dilation}

Three consequences of this reading come essentially for free. \emph{Normalization is a
theorem, not a constraint}: every box is total and the reference
measure is a probability, so nothing expressible at this level loses
mass. \emph{Discard is marginalization}: deleting an output wire
marginalizes the kernel. \emph{Copy is sharing, never resampling}:
fan-out of a noise-derived value is an explicit copy row --- both
consumers see one realization --- while an independent resample is a
second noise port, so the two are different graphs with different
digests. The last point is a central subtlety of the probabilistic
string-diagram literature (tossing a coin and copying the outcome is
not tossing two coins~\cite{fritzliang22,bonchi26b}), and the
linearity discipline of the deterministic calculus settles it
structurally, with no new axioms;
Figure~\ref{fig:mk-copy-resample} draws the two graphs.

All three consequences are machine-checked in the toolkit's Lean
tree at the finite-outcome stratum, where a stage is a measurable
outcome map from the noise cube $[0,1]^{N}$ to a finite value set.
Total mass one is the theorem \texttt{sum\_noiseDenotFun}; pushing
the denotation forward along a projection is
\texttt{noiseDenot\_comp}, with discard as the special case that
forgets a coordinate. For
composition, the law runs in the positive direction and in the
negative one: with disjoint noise rosters, a sequential composite
denotes the Kleisli composite of the stage denotations
(\texttt{noiseDenot\_seqComp}), while a shared roster provably
departs from that composite (\texttt{copy\_ne\_kleisli}), so the
disjointness hypothesis is content rather than bookkeeping. All of
these live in the namespace \texttt{SuFunctor.Stochastic},
sorry-free.

\begin{figure}[t]
\centering
\begin{tikzpicture}[x=0.94cm,y=0.82cm,
  box/.style={draw=accentdeep, fill=accent!50, rounded corners=1pt,
    font=\footnotesize\ttfamily, inner sep=2.5pt, minimum height=4.2mm},
  pbox/.style={draw=NoteGray, fill=NoteGray!12, rounded corners=1pt,
    font=\footnotesize\ttfamily, inner sep=2.5pt, minimum height=4.2mm},
  wire/.style={-, line width=0.5pt, draw=black!75},
  port/.style={font=\small\itshape, inner sep=1pt},
  plab/.style={font=\scriptsize, inner sep=1pt, text=black!75}]
\node[port] (u) at (0, 0.75) {$u$};
\node[pbox] (cp) at (1.8, 0.75) {copy$_2$};
\node[box] (f1) at (3.6, 1.5) {$f$};
\node[box] (g1) at (3.6, 0)   {$g$};
\node[port] (o1) at (4.9, 1.5) {$y_1$};
\node[port] (o2) at (4.9, 0)   {$y_2$};
\draw[wire] (u) -- (cp);
\draw[wire] (cp.east) to[out=35,in=180] (f1.west);
\draw[wire] (cp.east) to[out=-35,in=180] (g1.west);
\draw[wire] (f1) -- (o1);
\draw[wire] (g1) -- (o2);
\node[plab, anchor=north] at (2.6, -0.85) {one draw, shared};
\begin{scope}[xshift=7.2cm]
\node[port] (v1) at (0, 1.5) {$u_1$};
\node[port] (v2) at (0, 0)   {$u_2$};
\node[box] (f2) at (1.8, 1.5) {$f$};
\node[box] (g2) at (1.8, 0)   {$g$};
\node[port] (p1) at (3.1, 1.5) {$y_1$};
\node[port] (p2) at (3.1, 0)   {$y_2$};
\draw[wire] (v1) -- (f2);
\draw[wire] (v2) -- (g2);
\draw[wire] (f2) -- (p1);
\draw[wire] (g2) -- (p2);
\node[plab, anchor=north] at (1.55, -0.85) {two independent draws};
\end{scope}
\end{tikzpicture}
\caption{Copy is sharing, never resampling. Fanning a noise port out
through an explicit copy (left) hands both consumers one realization;
an independent redraw is a second noise port (right). The two are
different graphs with different digests, and their denotations
provably differ (\texttt{copy\_ne\_kleisli}).}
\label{fig:mk-copy-resample}
\end{figure}

The evaluator never samples: given
noise values as ordinary rational inputs, the exact evaluator runs
unchanged. Simulation is the evaluator plus an external random-number
driver; certificates consume the kernel's polynomial data and the
reference measure, never samples.

\section{Jump structure: mass, intensity, thinning}
\label{sec:mk-jump}

Continuous time asks for more than normalized kernels. Jump structure
contributes a jump measure and a killing rate to a process's extended
generator~\cite{davis-1993,ethier-kurtz-1986}, and what must be
tracked is \emph{mass}: rates scale it, killing removes it, and
staying put carries it without moving. The kernel vocabulary for
this tier is a small generator language over normalized
kernels --- jump, stay, kill, scale, superposition, and the empty
generator --- in which every term carries a rational
\emph{intensity majorant}, computed by the typing derivation: scale
multiplies it, superposition adds it.
State-dependent intensities $\lambda(x)$ recover from a scalar
majorant $\Lambda$ by thinning: run the clock at rate $\Lambda$ and,
at each ring, branch to the real jump with probability
$\lambda(x)/\Lambda$, staying put otherwise.

Killing is deliberately
distinct from the mass-preserving discard, and stay is deliberately
distinct from zero: a stay contributes nothing to the generator but
carries mass in the grading, which is what thinning and simulation
budgets need to see.

This vocabulary is machine-checked too --- at the finite-state
stratum, in the same sorry-free module family as the dilation
theorems. Every
term elaborates, by structural induction, to a \emph{thinning normal
form} (\texttt{ThinningForm}): a majorant together with one
normalized kernel into the state space extended by two tagged
outcomes, stay and kill. A sub-stochastic mass deficit is therefore
unrepresentable rather than merely discouraged. The elaboration
cannot emit sub-mass, and the declared intensity distributes exactly
over the named channels (\texttt{elaborate\_mass\_partition}, with
\texttt{outflow\_le\_intensity} giving the majorant its meaning).
Elaboration itself preserves the extended generator
(\texttt{generator\_elaborate}). The normal form is deliberately
non-unique, and its two semantics-preserving moves are themselves
theorems: a majorant may be inflated with the slack paid to stay
(\texttt{generator\_dilate}), and self-jump mass may be reclassified
as stay (\texttt{generator\_absorbDiagonal}) --- killing mass may
not, which is the kill/stay separation made mechanical. Composition
is machine-checked at the same stratum (module
\texttt{SuFunctor.Stochastic.JumpAlgebra}): the normal forms carry
the term vocabulary's composite algebra --- superposition,
thinning-scale, coordinate extension, and parallel interleaving, with
the intensity majorants as its grading --- and elaboration is a
homomorphism of that graded algebra, definitionally for the module
operations (\texttt{elaborate\_superpose} and its kin close by
\texttt{rfl}) and in the thinning quotient for interleaving
(\texttt{generator\_elaborate\_interleave}); sequential composition
of jump kernels is no part of this tier's algebra by design ---
successive jumps compose through time, and the Kleisli law for
discrete steps is section~\ref{sec:mk-dilation}'s. The same
tagged form recovers the jump-diffusion generator of
\bookref{ch:qtsl-mathematical-models}{the modelling chapter}. That
chapter's legacy encoding reads a reset mixture's mass deficit as
killing, in Davis's convention; rewriting each deficit into an
explicit kill outcome leaves the legacy transfer increment
unchanged, and the full generator is recovered with nothing
sub-stochastic remaining (\texttt{generator\_eq\_stepMix}). Every
well-typed term thus declares a finite intensity majorant, the same
bounded-intensity condition the non-Zeno arguments elsewhere in this
book compose with.

One consumer is worth naming here: the supervisors of the composition
inventory (section~\ref{sec:cp-supervise}) are ordinary jump-only
Systems --- policy memory is marking, and a policy update's
stochastic content is a firing kernel of this tier --- so control
through enablement adds no kernel vocabulary of its own.

\section{Packaged generators}
\label{sec:mk-packaged}

Two families of packaged generators round out the vocabulary, both in
the packaged-symbol discipline of
\bookref{ch:qtsl-hypernets}{the hypernet chapter} (a definitional
expansion, a projection lemma, and a candidacy rule --- never a new
box).\footnote{Two senses of \emph{generator} meet in this chapter:
the analyst's (section~\ref{sec:mk-jump} presents extended
generators) and the algebraist's (a packaged generator is a named
symbol with a definitional expansion). The word is inherited from
both literatures; the toolkit keeps the senses apart by type.} \emph{Distribution generators} name the classical laws: each is
an alias for a certified piecewise-polynomial quantile table applied
to a noise port, carrying an explicit accuracy budget, so that a model
that draws from a named distribution elaborates into the same hypernet
language every other certificate in this book speaks.
Section~\ref{sec:exp-dist} is their coverage dashboard. The
\emph{stochastic-matrix generator} covers the finite-state case
directly: a kernel between finite-extent wires given by an explicit
tensor of transition probabilities, whose key conformance obligation
is row-stochasticity --- mass one from each source state --- making
finite Markov chains and the jump chains of continuous-time models
first-class citizens rather than encodings.

A status note, in the spirit of the preview note at the front of this
volume. The deterministic substrate is
\bookref{ch:qtsl-hypernets}{the hypernet chapter}'s and is shipped.
The denotational core of the noise-port reading ---
normalization, discard as marginalization, the Kleisli composition
law and its shared-roster failure --- is machine-checked at the
finite-outcome stratum (section~\ref{sec:mk-dilation}), and the jump
vocabulary's semantics --- the intensity grading, the thinning
normal form, its soundness, and the thinning-quotient moves --- are
machine-checked at the finite-state stratum
(section~\ref{sec:mk-jump}). The kernel-language \emph{surface} has
meanwhile settled: the certificate theory
(\path{theories/Qtsl.geolog}) carries it as normative schema. A
kernel citation there names its hypernet by digest together with its
noise roster and three tagged outcome channels: next, stay, and
\emph{cemetery}, the wire format's name for the kill outcome. With
every unit of mass explicitly named by one of the three channels, a
sub-stochastic deficit is unrepresentable, and a pure kill is a
weight-one cemetery branch.

A transition carries a \emph{clock} with
three wings, and the wing decides when the kernel is consulted. A
\emph{racing} clock declares its intensity majorant. At a finite
majorant, its polynomial presentation elaborates definitionally into
the thinning normal form of section~\ref{sec:mk-jump}; the infinite
majorant instead marks the instant-fire boundary case, contributing
no generator term. A \emph{demonic} clock has no timing law
--- whether and when its kernel is consulted belongs to a scheduler
over which every judgment quantifies --- but carries a strictly
positive dwell floor in the system definition: the scheduler may fire
the transition only after that much contiguous dwell inside its guard
region. A row with floor $d$ therefore fires at most $T/d + 1$ times
by horizon $T$ under every resolution, and the demonic wing cannot
produce Zeno behaviour at all. An untimed Petri net transcribes with
unit dwell floors: the transcription is its \emph{discrete-time
unfolding}, embedded in the continuous-time frame. It is faithful
for ordering-and-counting properties, which are all an untimed net
ever asserted. A
\emph{forced} clock fires at exactly its declared positive dwell of
contiguous in-guard time: deterministic delay, completing the triple
--- exponential, immediate, deterministic --- from which
deterministic and stochastic Petri nets are built, so watchdogs,
timeouts, and protocol timers are native rows rather than encodings.
Preemption is by guard reset (whatever fires first wins, and the
forced clock resets with its guard). A same-instant tie between
two forced transitions resolves through the demonic scheduler's
seat: the classical deterministic-net literature needs an ad-hoc
tie-breaking convention exactly where an imprecise-probability
semantics has a seat already waiting. With both dwell floors
strictly positive, the one syntactically visible source of Zeno
behaviour in the vocabulary is an unbounded chain of instant
firings --- precisely what the crossing-budget certificates price.

The stochastic-matrix generator is
carried with row-stochasticity among its named conformance
obligations. The certified quantile tables behind the dashboard's
certified rows exist today as test-suite artefacts, and the dashboard
section states, row by row, what each status rests on and what would
move it. This chapter still owes the settled surface its
presentation: a coming revision states it at the deterministic
calculus's level of rigour, with its soundness against the semantic
universe, and a section on firing kernels alongside.

\section{Distribution generators: a coverage dashboard}
\label{sec:exp-dist}

The toolkit's stochastic vocabulary needs named random variates the way
its deterministic vocabulary needs elementary functions, and it admits
them under the same discipline: never as new primitives. A
\emph{distribution generator} is a packaged symbol whose expansion is a
certified piecewise-polynomial quantile table applied to uniform noise
ports, so a model that draws from a named distribution elaborates ---
before anything is checked --- into the same hypernet language every
other certificate in this book speaks
(\bookref{ch:qtsl-hypernets}{the hypernets chapter}). Checkers and
certificate formats learn nothing new; being wrong in a table costs a
rejected certificate, never a false verdict.

Each generator registers with an explicit budget
$(\varepsilon;\,\delta_{lo},\delta_{hi})$: quantile sup-error at most
$\varepsilon$ on the core $[\delta_{lo},\,1-\delta_{hi}]$, in the
family's standard scale, with the clipped tail mass declared per side.
The working target is $\varepsilon \le 2^{-23}$ --- single-precision
machine epsilon --- with default clips of $2^{-30}$; heavy-tailed
families declare conspicuously larger honest clips rather than
pretending to a tail no polynomial reaches, and multiplicative families
state $\varepsilon$ relatively. Verification never encloses an inverse
transcendental: each certificate drives the \emph{forward} CDF through
polynomial enclosures, piece by piece in a rescaled frame, and closes
with a residual row, a derivative-floor row, and one rational
comparison. A verification class grades what forward material that
takes: \textbf{V0} nothing (the construction is exact), \textbf{V1}
polynomials only, \textbf{V2} the mechanized elementary enclosures
(exponential, sine, cosine, the error function, and now the incomplete
gamma), \textbf{V2q} certified quadrature, \textbf{V3} an enclosure
family built specially for the row that needs it (to date one exists:
the incomplete beta, whose tail lemma is its own). Every certified row also ships a refusal battery in the
sense the Experiments chapter
(\bookref{ch:experiments}{of this book's companion experiments
repository}) makes precise: tampered coefficients, shaved residuals,
dishonest tail declarations, partition gaps, and dropped
budget-composition terms must each be refused by a named row.

\paragraph{Enclosures are certified valid unconditionally.} The core
range bounds where the \emph{tightness} claim lives, not where
validity ends: a generator's certificate is a bundle of unconditional
theorems, and no row's truth is scoped to the core. The enclosures
are also certified to be \emph{tight} to single-precision machine
epsilon within their core range (the working target; a row that
declares a different $\varepsilon$ says so in the dashboard), and
outside the core they degrade honestly, in two genre-specific ways,
because the two kinds of table stand on different input axes. A
quantile table's axis carries the probability measure itself, so the
table is total on $[0,1]$: its end cells clamp to exact rational
constants, and the clamp-honesty rows certify that the mass beyond
each clamp is at most the declared $\delta$ --- a checked row, never
trust. Monotonicity keeps even the off-core bracket better than
trivial: beyond the core the true quantile stays pinned between the
certified core-edge value and the support boundary, a certified
inference wherever the monotonicity rows are claimed, so the
generated law couples to the named law within $\varepsilon$ except on
an event of probability at most $\delta_{lo} + \delta_{hi}$. A
function table --- a certified polynomial approximant of a plain
function on a declared domain, the wrappers among them --- stands on
an axis that carries no measure, so there is no tail mass with which
to pay for a clamp: off its declared domain such a table simply makes
no claim, and an off-domain evaluation degrades to the trivial
bracket --- nothing is known about this number --- which is still a
valid statement, never an assertion of impossibility. Refusal, where
it appears in this discipline, is a judgment about claims rather than
about numbers, and it lives at production and certification time: a
producer declines to emit a table it cannot cover, and a checker
refuses a tightness claim over domain a table does not cover. The
dashboard's \dredx{} against the unbounded uniform is different in
kind from all of this: there the family itself admits no quantile
function, an impossibility of the mathematics rather than a limit of
a table's coverage.

\paragraph{The census.} Coverage is measured against an external
standard the toolkit does not own: Crooks' \emph{Field Guide to
Continuous Probability Distributions}
\bookcite{crooks22}{(Crooks 2022)}, whose
numbering organizes some 170 named forms into on the order of a hundred
sections related by an explicit lattice of limits, reparametrizations,
and transforms. The Field Guide's own section numbers are the census:
every numbered form appears in exactly one row below (block rows name
the sections they cover; synonyms resolve to the section they cite ---
Lorentzian to Cauchy, Fisk to log-logistic, fatigue-life to
Birnbaum--Saunders). The \emph{route} column says how the family is
reached --- \emph{exact} (arrow algebra over existing generators: order
statistics, sums, products, mixtures, folds), \emph{composition}
(certified wrappers over the two quantile roots), \emph{own table} (a
directly-tabled quantile where that is cheaper than composition), or
\emph{dispatch} (a parameter-time classification to a member family) ---
with the verification class in parentheses. The \emph{status} column is
the dashboard: \emph{certified} rows carry receipts as described above;
\emph{routed} rows have a fixed construction and verification class but
no table yet; \emph{out of scope} rows are deliberate, each with the
honest reason, the priced alternative, and the condition that would
reopen it. Two roots (the standard exponential and standard normal
quantiles) and three wrappers (rational powers, logarithm, exponential)
generate everything except the own-table rows; location--scale families
register one standard table each, extended by exact affine arrows.

\medskip
\providecommand{\distdashcensusdate}{}\renewcommand{\distdashcensusdate}{July 27, 2026}

\noindent {\small Each row carries a mark in the same vocabulary as the QTSL rule catalogue's coverage dashboard: \dok{} certified, \damber{} routed (construction and verification class fixed, no table yet), \dredx{} out of scope; \dok\,\damber{} marks a row certified for its named part with the residue routed.}\medskip

\noindent\begin{minipage}{\linewidth}
\paragraph{Shared infrastructure (roots, wrappers, constants).}
\begin{center}\small
\begin{tabular}{@{}>{\raggedright\arraybackslash}p{7.9cm} c >{\raggedright\arraybackslash}p{5.8cm}@{}}
\toprule
object & & status \\
\midrule
exponential quantile root (the registered standard exponential) & \dok & certified --- $\varepsilon \le 4.3\times10^{-8}$ \\
error-function forward enclosure (mechanized; soundness stated against the Gaussian integral, standard axioms only) & \dok & certified \\
probit quantile root (the registered standard normal) & \dok & certified --- $\varepsilon \le 2.9\times10^{-8}$ \\
rational-power wrapper (reciprocal and square root included) & \dok & certified --- per-instance $\varepsilon \le 2^{-24}$ \\
logarithm wrapper (verified through the exponential's forward enclosure --- no logarithm enclosure exists or is needed) & \dok & certified --- $\varepsilon \le 2\times10^{-9}$ \\
exponential wrapper (sinh, cosh derived exactly) & \dok & certified --- relative-error native \\
certified constants ($\ln 2$, $\pi$, $2/\sqrt{\pi}$) & \dok & certified --- relative half-width $\le 2^{-55}$ \\
composition algebra: exact budget arithmetic per arrow, with a route audit that refuses any dropped term & \dok & certified \\
\bottomrule
\end{tabular}
\end{center}
\end{minipage}

\noindent\begin{minipage}{\linewidth}
\paragraph{Zero shape parameters (Field Guide ch.\ 1--4).}
\begin{center}\small
\begin{tabular}{@{}l >{\raggedright\arraybackslash}p{4.7cm} >{\raggedright\arraybackslash}p{2.3cm} c >{\raggedright\arraybackslash}p{5.3cm}@{}}
\toprule
\S{} & family & route & & status \\
\midrule
1.1 & uniform (half uniform; degenerate point mass) & exact (V0) & \dok & certified --- exact \\
1.--- & unbounded uniform & --- & \dredx & out of scope --- improper (not normalizable), so no quantile function exists; set-valued premises cover the intended uses \\
2.1 & exponential (anchored, standard) & composition (V2) & \dok & certified --- $\varepsilon \le 2^{-24}$ \\
3.1 & Laplace & composition (V2) & \dok & certified --- $\varepsilon \le 2^{-24}$ \\
4.1 & normal & composition (V2) & \dok & certified --- $\varepsilon \le 2.9\times10^{-8}$ at unit scale (exact affine scaling) \\
\bottomrule
\end{tabular}
\end{center}
\end{minipage}
\medskip

\noindent\begin{minipage}{\linewidth}
\paragraph{One shape parameter (ch.\ 5--9).}
\begin{center}\small
\begin{tabular}{@{}l >{\raggedright\arraybackslash}p{4.7cm} >{\raggedright\arraybackslash}p{2.3cm} c >{\raggedright\arraybackslash}p{5.3cm}@{}}
\toprule
\S{} & family & route & & status \\
\midrule
5.1 & power function (generalized Pareto, $q$-exponential reparametrizations) & composition (V1) & \dok & certified --- $\varepsilon \le 2^{-23}$ \\
5.4 & wedge (ascending, descending) & exact (V0) & \dok & certified --- exact \\
5.5 & Pareto (Lomax 5.6 = shifted) & composition (V1) & \dok & certified --- $\varepsilon \le 2^{-23}$ (moderate core; the wide core is deep-tail-dominated and its excluded mass is declared) \\
5.7 & exponential ratio (uniform-prime 5.8; Lomax(1)) & composition (V1) & \dok & certified --- $\varepsilon \le 2^{-23}$ \\
6.1 & log-normal (Gibrat) & composition (V2) & \dok & certified --- relative $\varepsilon \le 2^{-23}$ \\
7.1a & gamma, integer shape (Erlang) & exact (V2) & \dok & certified --- $\varepsilon \le 2^{-24}$ (three-port pin; per-port tables pre-tightened so sums stay under target) \\
7.1b & gamma, half-integer shape & exact (V2) & \dok & certified --- declared per alias --- the squared-clamp amplification is stated at its measured value (e.g.\ $3.5\times10^{-7}$ at one degree of freedom) \\
7.1c & gamma, general rational shape & enclosure (V2) & \dok & certified --- $\varepsilon \le 2^{-23}$ at the pinned fractional shapes (measured $2^{-27}$-class) \\
7.3 & chi-square (scaled; Porter--Thomas) & exact (V2) & \dok & certified --- $\varepsilon \le 2^{-24}$ at even $k$; odd $k$ declared per alias \\
8.1 & gamma-exponential (chi-square-exponential; generalized Gumbel) & composition (V2) & \dok & certified --- $\varepsilon \le 2^{-23}$ (integer shapes); half-integer members inherit the squared-clamp declaration \\
8.5 & Gumbel & composition (V2) & \dok & certified --- $\varepsilon \le 2^{-26}$ \\
8.7 & BHP & composition (V2) & \dok & certified --- $\varepsilon \le 2^{-23}$ at the irrational shape $\pi/2$ (a bracketed-shape pin plus an exact shape-transfer argument) \\
8.8 & Moyal & composition (V2) & \dok & certified --- $\varepsilon \le 2^{-23}$ \\
9.1 & Pearson VII, half-integer $m$ (Student $t$, integer $\nu$; Breit--Wigner) & composition (V2) & \dok & certified --- declared per floor --- the heavy tail's excluded-event mass is priced exactly, and the sup concentrates in near-floor tail events \\
9.4 & Student $t_3$ (the odd-$\nu$ pattern) & own table (V2) & \dok & certified --- $\varepsilon \le 2^{-26}$ \\
9.6 & Cauchy (Lorentzian; $t_1$) & own table (V2) & \dok & certified --- $\varepsilon \le 2^{-23}$ \\
\bottomrule
\end{tabular}
\end{center}
\end{minipage}
\medskip

\noindent\begin{minipage}{\linewidth}
\paragraph{Two shape parameters (ch.\ 10--12).}
\begin{center}\small
\begin{tabular}{@{}l >{\raggedright\arraybackslash}p{4.7cm} >{\raggedright\arraybackslash}p{2.3cm} c >{\raggedright\arraybackslash}p{5.3cm}@{}}
\toprule
\S{} & family & route & & status \\
\midrule
10.1 & unit gamma & composition (V2) & \dok & certified --- $\varepsilon \le 2^{-23}$ (relative $\le 2^{-38}$ at the standard pins); half-integer members declare their relative error \\
10.2 & uniform product & exact (V0) & \dok & certified --- exact \\
11.1 & Amoroso (Stacy = anchored) & composition (V1+V2) & \dok & certified --- $\varepsilon \le 2^{-23}$ (integer and half-integer pins; both signs of the power parameter --- the inverse-Amoroso branch included) \\
11.7 & half-normal & exact (V2) & \dok & certified --- $\varepsilon \le 2.9\times10^{-8}$ at unit scale; no clipped mass at the support edge \\
11.8 & chi (Rayleigh, Maxwell, Nakagami) & composition (V2) & \dok & certified --- $\varepsilon \le 2^{-23}$ \\
11.13 & inverse-Amoroso block: inverse gamma and exponential, L\'evy, scaled/inverse chi and chi-square, inverse Rayleigh, Maxwell, half-normal, Nakagami (11.13--11.23) & composition (V1+V2) & \dok & certified --- $\varepsilon \le 2^{-23}$ (clean pins); the heavy members (L\'evy and kin) are declared per alias at every floor \\
11.24 & Fisher--Tippett block: Weibull, reversed Weibull, Fr\'echet, and their generalized forms (11.24--11.29) & composition (V1+V2) & \dok & certified --- $\varepsilon \le 2^{-23}$ (Weibull pins); the heavy-tailed Fr\'echet members declare their clip masses \\
12.1a & beta, integer shapes (PERT; Pearson XII) & exact (V0) & \dok & certified --- exact \\
12.1b & beta, half-integer shapes & composition (V2) & \dok & certified --- declared per pin --- the ratio route states its amplification (it cannot reach the working target at any floor); correctness rests on exact anchors, and tight-absolute consumers take the incomplete-beta root at the same shapes \\
12.1c & beta, general rational shapes & enclosure (V3) & \dok & certified --- $\varepsilon \le 2^{-30}$ at the five full-depth pins (vs the $2^{-23}$ target), $\delta = 2^{-31}$ per side \\
12.5 & central beta & exact (V0) & \dok & certified --- exact \\
12.6 & arcsine (centered arcsine) & own table (V2) & \dok & certified --- $\varepsilon \le 2^{-24}$; exact end cells, no clipped mass \\
12.8 & semicircle & own table (V2) & \dok & certified --- $\varepsilon \le 2^{-24}$; exact end cells, no clipped mass \\
12.9 & Epanechnikov, biweight, triweight (12.9--12.11) & exact (V0/V1) & \dok & certified --- exact \\
\bottomrule
\end{tabular}
\end{center}
\end{minipage}
\medskip

\noindent\begin{minipage}{\linewidth}
\paragraph{Two shape parameters (ch.\ 13--16).}
\begin{center}\small
\begin{tabular}{@{}l >{\raggedright\arraybackslash}p{4.7cm} >{\raggedright\arraybackslash}p{2.3cm} c >{\raggedright\arraybackslash}p{5.3cm}@{}}
\toprule
\S{} & family & route & & status \\
\midrule
13.1 & beta prime (inverse Lomax; 13.1--13.4) & exact + composition (V1 int / V2 half-int) & \dok & certified --- $\varepsilon \le 2^{-23}$ (integer shapes; the $\gamma{=}1$ heavy tail declares its clip mass); half-integer members declare per floor, cross-anchored against the integer route \\
13.3 & $F(k_1,k_2)$ & composition (V1 even / V2 odd) & \dok & certified --- $\varepsilon \le 2^{-23}$ (even degrees of freedom); odd-parity members declare per floor (mixed parity with an even denominator is measurably tightest) \\
14.1 & beta-exponential (hyperbolic-sine; Nadarajah--Kotz) & composition (V2) & \dok & certified --- $\varepsilon \le 2^{-31}$ \\
14.2 & exponentiated exponential & composition (V1+V2) & \dok & certified --- $\varepsilon \le 2^{-23}$ \\
15.1 & beta-logistic (Fisher $z$) & exact (V2) & \dok & certified --- $\varepsilon \le 2^{-23}$ (integer shapes --- the logit-chain identity is distribution-level exact); half-integer members declare their chain budgets \\
15.2 & Burr II (reversed Burr II) & composition (V1+V2) & \dok & certified --- $\varepsilon \le 2^{-23}$ \\
15.4 & central logistic & exact (V2) & \dok & certified --- $\varepsilon \le 2^{-24}$ (= the logistic pin's exactly: order-statistic gates are 1-Lipschitz in the sup norm) \\
15.5 & logistic & exact (V2) & \dok & certified --- $\varepsilon \le 2^{-24}$ \\
15.6 & hyperbolic secant & own table (V2) & \dok & certified --- $\varepsilon \le 2^{-26}$ \\
16.1 & Pearson IV & quadrature (V2q) & \dok & certified --- the composed sup error is stated per pin with the tangent-stage Lipschitz amplifier explicit; heavy polynomial tails make deep clamps expensive and the two-floor presentation applies (both floors stated) \\
\bottomrule
\end{tabular}
\end{center}
\end{minipage}
\medskip

\noindent\begin{minipage}{\linewidth}
\paragraph{Three and more shape parameters (ch.\ 17--20).}
\begin{center}\small
\begin{tabular}{@{}l >{\raggedright\arraybackslash}p{4.7cm} >{\raggedright\arraybackslash}p{2.3cm} c >{\raggedright\arraybackslash}p{5.3cm}@{}}
\toprule
\S{} & family & route & & status \\
\midrule
17.1 & generalized beta & composition (V1+V2) & \dok & certified --- $\varepsilon \le 2^{-39}$ (standard pins); the $\beta{=}1$ member degenerates to its source exactly \\
17.2 & Kumaraswamy & composition (V1) & \dok & certified --- $\varepsilon \le 2^{-23}$ \\
18.1 & generalized beta prime (transformed beta) & composition (V1+V2) & \dok & certified --- $\varepsilon \le 2^{-35}$ (standard pins) \\
18.3 & Burr block: Burr XII (Singh--Maddala), Dagum, paralogistic (18.3--18.6) & composition (V1) & \dok & certified --- $\varepsilon \le 2^{-23}$ \\
18.7 & log-logistic (Fisk) & composition (V1) & \dok & certified --- $\varepsilon \le 2^{-23}$ \\
18.8 & half-Pearson VII (half-Cauchy 18.9) & composition / own table (V2) & \dok & certified --- the algebraic member $m{=}\tfrac32$ at $\varepsilon \le 2^{-23}$ (provably the only polynomial-route shape); half-Cauchy takes the exact tangent-table fold ($\varepsilon \le 2^{-23}$, zero low-side clip); general half-integer $m$ as a half-normal $\times$ chi-ladder product net, declared per member (heavy-ratio tail geometry) \\
18.10 & half generalized Pearson VII; half-Laha & member routes & \dok & certified --- heavy-tailed per the Pearson VII geometry: the inverse gamma chain's clamp mass is the declared tail mass and the product amplification by the clamp value is stated per pin \\
19.1 & Pearson family & dispatch & \dok & certified --- exact --- a classification, not a table: parameters dispatch to a member family at registration time, in exact rational (and quadratic-extension) arithmetic \\
19.2 & $q$-Gaussian & dispatch & \dok & certified --- exact --- a reparametrization, not a table: $q>1$ dispatches to Student $t$, $q<1$ to central beta, $q\to1$ to normal \\
20.--- & grand-unified block: Halphen, Sichel, Libby--Novick, Gauss- and confluent-hypergeometric, Perks, Appell beta, Laha, and kin (20.2, 20.5--20.18) & quadrature (V2q) & \dok & certified --- $\varepsilon \le 2^{-23}$-class per member at declared cores ($2^{-24}$--$2^{-28}$ for the quadrature members; normalization constants --- Bessel-$K$ and hypergeometric values --- pinned by bracketed quadrature to $10^{-13}$--$10^{-15}$ width, no special function ever named; the generalized members cross-anchored two ways, direct quadrature vs.\ power-wrapper pushforward, at $4.5\times10^{-9}$ sup gap); heavy-tailed members declare a core-depth frontier ($2^{-23}$-class at moderate cores, deeper-core cost stated); the Perks family is certified across its full rational-shape domain (all four discriminant regimes; $2^{-29.9}$--$2^{-30.7}$ at the full-depth pins) \\
20.3 & inverse Gaussian (reciprocal inverse Gaussian 20.4) & own table (V2) & \dok & certified --- $\varepsilon \le 2^{-23}$ \\
\bottomrule
\end{tabular}
\end{center}
\end{minipage}
\medskip

\noindent\begin{minipage}{\linewidth}
\paragraph{Miscellanea (ch.\ 21).}
\begin{center}\small
\begin{tabular}{@{}l >{\raggedright\arraybackslash}p{4.7cm} >{\raggedright\arraybackslash}p{2.3cm} c >{\raggedright\arraybackslash}p{5.3cm}@{}}
\toprule
\S{} & family & route & & status \\
\midrule
21.1 & Bates (Irwin--Hall 21.9; triangular 21.22; uniform difference 21.23) & exact (V0/V1) & \dok & certified --- exact \\
21.2 & beta-Fisher--Tippett (exponentiated Weibull) & composition (V1+V2) & \dok & certified --- $\varepsilon \le 2^{-23}$ \\
21.3 & Birnbaum--Saunders (fatigue-life) & composition (V1+V2) & \dok & certified --- $\varepsilon \le 2^{-23}$ \\
21.4 & exponential power, general $\beta$ (Box--Tiao; Subbotin) & enclosure (V2) & \dok & certified --- $\varepsilon \le 2^{-23}$ (the classical members anchor: $\beta{=}1$ Laplace exactly, $\beta{=}2$ normal against the error function) \\
21.6 & generalized Pearson VII & member routes & \dok & certified --- as 18.10 (the symmetric double: sign cell-mixture takes the max of two mirror branches, so the half core's declaration carries over) \\
21.8 & K (generalized K 21.5) & exact (V2) & \dok & certified --- $\varepsilon \le 2^{-23}$ (integer shapes); half-integer members inherit the squared-clamp declaration \\
21.10 & Johnson SU (SB, SL variants) & composition (V2) & \dok & certified --- $\varepsilon \le 2^{-23}$ \\
21.12 & log-Cauchy & composition (V2) & \dok & certified --- relative $\varepsilon \le 2^{-23}$ (the log-normal precedent); a family with no moments of any order --- only the relative bound is claimed, the astronomical absolute error is stated, and the excluded tail mass is declared per side \\
21.13 & Meridian & exact (V1 (+cell)) & \dok & certified --- $\varepsilon \le 2^{-23}$ (moderate core) \\
21.15 & noncentral chi-square block: noncentral chi 21.14, Rice 21.18, noncentral $F$ 21.16 & exact + composition (V2) & \dok & certified --- $\varepsilon \le 2^{-23}$ \\
21.17 & pseudo-Voigt & exact (V2) & \dok & certified --- $\varepsilon \le 2^{-23}$ (the mixture takes the larger component budget, inheriting its Cauchy component's declaration exactly) \\
21.19 & slash & composition (V2) & \dok & certified --- $\varepsilon \le 2^{-23}$ (moderate core); the full heavy core's error is declared at its measured value \\
21.20 & stable (Holtsmark 21.7; Landau 21.11) & composition (V2) & \dok & certified --- $\varepsilon \le 2^{-23}$ at the certified pins (Holtsmark $2^{-24.3}$; Landau $2^{-23.6}$ --- the $\alpha = 1$ seam member); the angular tables are direct mid-slack residual rows, one transcendental bracket per pin \\
21.21 & Suzuki & exact (V2) & \dok & certified --- $\varepsilon \le 2^{-23}$ (small shapes); wide-support members declare the log-normal scale amplification \\
21.24 & Voigt & exact (V2) & \dok & certified --- $\varepsilon \le 2^{-23}$ at the half-scale witness pin; the unit pin states its summed budget (component errors add across the convolution sum) \\
\bottomrule
\end{tabular}
\end{center}
\end{minipage}
\medskip


\paragraph{What the statuses rest on.} Status column as of \distdashcensusdate; it will move with each preview. Unlike the tables of the
Experiments chapter, this section's rows are not yet re-derived by the
experiments repository's make targets: the dashboard re-derives from
the toolkit workshop's census of record, and the certificates behind
the certified rows live at this date in the toolkit repository as
ordinary test-suite artefacts (accept fixtures at the declared budgets
plus their refusal batteries). Migrating them onto the experiments
repository's reproduction spine --- regenerated as the same sealed
containers that chapter's fluid-drift certificates already ride,
re-verified by the blessed checkers, one make target for the family ---
is named successor work, and this section's tables will then re-derive
from committed receipts like every table of the Experiments chapter.
Per the preview note at the front of this volume, until an artefact
appears in the public repositories, treat its mention as a statement of
what exists, not yet a link.
  \chapter{Systems: Mathematical Models}
\label{ch:qtsl-mathematical-models}

\noindent\draftmark\medskip

\newcommand{\syscatmeta}[1]{\ruleblock{Computer Science Metanotation}\sourcenote{as realized in the mechanized toolkit: the family
  vocabulary is the manifest schema's own
  (\texttt{crates/mtk-catalogue/src/system/}), and the former's
  mechanized ground truth is quoted in the API block below --- the
  producer signatures and, for profile entries, the preset signature
  naming the populated families --- extracted from the toolkit
  sources at every regeneration of this catalogue, so the two cannot
  drift apart silently. The formation rule here is typeset by hand
  for notational quality; $\Sigma$, $\mathrm{pop}$, and profile
  conformance are defined in section~\ref{sec:mt-core}.}#1}
\newcommand{\syscatmetapat}[1]{\ruleblock{Computer Science Metanotation}\sourcenote{as realized in the mechanized toolkit: the row
  vocabulary is the composite stratum's own
  (\texttt{crates/\allowbreak mtk-catalogue/\allowbreak src/\allowbreak
  system/\allowbreak compose.rs} and the pattern
  row families of \texttt{crates/\allowbreak mtk-catalogue/\allowbreak
  src/\allowbreak compose/}), and the
  former's mechanized ground truth is the builder signature quoted in
  the API block below, extracted at every regeneration. The
  formation rule, where the entry states one, is typeset by hand for
  notational quality, its premises the formation data the pattern
  builders consume and refuse on; transport soundness, tiering, and
  the causality/simultaneity side conditions are rule- and
  kernel-side, catalogued in the Rule Catalogue.}#1}
\newcommand{\syscatmetalean}[1]{\ruleblock{Computer Science Metanotation}\sourcenote{as mechanized in the toolkit's Lean tree: this entry's
  ground truth is the presentation structure and compiled row grain
  of its adequacy development (\texttt{lean/su-functor/}, module
  \texttt{SuFunctor.Adequacy.QuantumNet}, over the fragment structure
  of \texttt{QtslAux.QuantumAssembly}), rather than a producer
  extraction. No producer surface exists for this entry at this
  preview; the API block is omitted rather than invented, and the
  formation rule below is typeset by hand from the mechanized
  structure's fields.}#1}

Chapter~\ref{ch:qtsl-semantic-universe} lays out the universe these
profiles present into, with its records, its structured time base, its
imprecise denotation, and its composition calculus. The statuses below
grade each frame against that semantics. The grading has a stated
division of labour: most entries give the frame's own rung in their
\emph{Semantics} paragraph --- a generator, a path measure, a
solution concept --- and the bridge lemmas, staged as named obligations
and discharged frame by frame, then make up the distance between that
rung and the record-space denotation. Where an entry's semantic
content \emph{is} an interpretation into the universe's objects, as
with the Bag former's carrier colimits or the credal reading of
quantification over policies, the entry carries the interpretation
explicitly.\medskip

Throughout this book a \emph{System} means a structured presentation,
a composition term over \emph{mechanisms} and \emph{wirings}. A
mechanism is a primitive dynamics declaration: a family of fields
over declared state variables together with an interface binding. A
wiring identifies variables across interfaces, subject to causality
and simultaneity side conditions. A System is never the
semantic quotient of the dynamics it presents, and in particular it
is never the infinitesimal generator. The generator re-enters one level
down, as the certification interface, which compiles mechanism rows
into the obligation rows the exact kernel checks. Time is
base-universal. Composites live on real time, clocks are internalized
as mechanism state, and a discrete-time component taken in isolation
is schedule-generic, acquiring its schedule assumptions exactly
at the boundary where it composes with continuous neighbours.

One consequence organizes this chapter: \textbf{a model type is a
profile, not a tag}. Nothing in the theory stores ``this model is a
stochastic differential equation''. A model's type consists in
\emph{which field families its mechanisms populate}, and a composite
may mix profiles freely. Ordinary engineering forces this design.
Hybrid stochastic systems must be first-class citizens, and a robot
--- any robot --- is a mixed discrete/continuous composite that must
type-check across its own components. A stored per-model doctrine tag
fails both requirements, whereas a profile satisfies both by
construction. The sections below give the normative inventory of those
profiles, at the granularity a modeller shops at. A companion
human-facing book, \emph{Mathematical Modelling with MAITRIA}, will
treat modelling as a discipline in its own right, while this chapter
takes only what a toolkit user needs.

The roster of sections is not ours. It is the list of modelling
frames that the Safeguarded AI programme thesis (David ``davidad''
Dalrymple, \emph{Safeguarded AI: Constructing Guaranteed Safety},
ARIA) \cite{dalrymple-sgai-thesis} sets out in its Appendix~A.1 as
the frames a unified modelling language should subsume --- the same
TA1.1/TA1.2 rubric the rest of this book grades itself against. Grading a candidate language against
that list is itself programme-native. In a December 2023 programme
talk, davidad ran the exercise, scoring the identical list against
stochastically and dynamically coloured Petri nets
(SDCPN, section~\ref{sec:mt-sdcpn}) as the candidate unifier. Two
cells came back not-covered in 2023, deterministic and stochastic
partial differential equations, and the deterministic half is where
the present stack has landed some of its freshest certificates
(section~\ref{sec:mt-pde}). We read that as evidence that the
candidate graded here is not merely re-walking SDCPN's coverage, and
we keep the grading honest for the same reason the 2023 slide did,
since the value of the exercise lies in the map of what remains.

The catalogue itself runs wider than the 2023 list. This chapter
presents the \emph{Model Types} wing of a two-winged \emph{System
Catalogue}, with model types here and composition patterns in
chapter~\ref{ch:qtsl-composition-patterns}. The catalogue's roster
also carries the substrate every profile populates (the System core,
section~\ref{sec:mt-core}; tick declarations,
section~\ref{sec:mt-tick}) together with the types in production use
beyond the list (clocked circuits, the Bag former, SDCPN, the i.i.d.\
tier), and one frame beyond the list whose mechanized semantics
arrived ahead of any producer surface: the quantum nets of
section~\ref{sec:mt-qnet}. Every entry, on the list or beyond it,
follows the same section grammar.

\paragraph{How to read the sections.} Every section carries the same
four rubrics, adapted from the working catalogue's entry format.
\emph{The frame} says what the model class is, and gives its
bread-and-butter verification question. \emph{Profile} says which
field families a model of this type populates: content language,
carrier, disturbance semantics, and time semantics.
\emph{Semantics} names the bridge-obligation class its satisfaction
claims discharge, and every analytic lemma that licenses reading exact
algebra as a statement about trajectories or measures is
\emph{staged as a named hypothesis}, never silently assumed.
\emph{Status} gives the state of four legs, each graded only against a
witness (below). Each entry then closes with two blocks in the Rule
Catalogue's grammar. The first, a \emph{Computer Science Metanotation}
block, gives the entry's \emph{term former}, an inference rule whose
premises are the formation data the builders consume, with their
sorts, and whose conclusion is a primitive judgment at the sort
$\mathsf{System}$ (or at the entry's own true sort, where an entry
declares vocabulary rather than forming a System). The second, an
\emph{API} block, gives the producer surfaces (TypeScript, then Rust)
with signatures and constructions extracted from the toolkit sources
at every regeneration of this catalogue. Frontier entries
(section~\ref{sec:mt-frontier}) carry neither, since no producer
surface exists for them at this preview, and an invented one would
defeat the grading.  The wiring-pattern entries (delay differential
equations, section~\ref{sec:mt-dde}; partially observable decision
processes, section~\ref{sec:mt-pomdp}) also carry neither,
deliberately: each of those frames \emph{is} a composite of other
formers, and the absence of a former of its own is the entry's
content. The quantum entry (section~\ref{sec:mt-qnet})
sits between the two conventions: its formation rule stands on
mechanized Lean ground truth, so it appears; its API block waits for
a producer surface to exist. The four legs ---

\begin{itemize}
\item \textbf{representation}: a model of the frame is declarable as
  rows of the modelling language in a form a frame-native reader
  recognizes as \emph{their} model, with the frame's compositional
  grain preserved rather than flattened by an encoding trick;
\item \textbf{semantics}: the frame's standard semantics is
  recovered, either exactly or at a named rung with the bridging
  hypotheses staged;
\item \textbf{certification}: at least one of the frame's own
  bread-and-butter questions compiles to obligations checked by the
  \emph{unmodified} exact kernel, with a shipped exact-rational
  artefact;
\item \textbf{compositional closure}: the frame's systems compose
  with systems of \emph{other} frames through the patterns of
  chapter~\ref{ch:qtsl-composition-patterns}, with certificates
  transported across the composite boundary and witnessed by shipped
  artefacts rather than design prose. A frame can pass the first three
  legs and still be flattened at its composition boundaries, and this
  leg exists to detect exactly that. Where a frame's semantics
  genuinely cannot instantiate a pattern's composition data (a static
  distribution has no internal time axis for a supervisory switch to
  sequence), the omission is earned by a filed argument from the
  frame's own semantics, never by a bare declaration.
\end{itemize}

As of this preview, \emph{no frame on the list passes all four legs},
and exactly one (ordinary differential equations) holds a cross-frame
composition witness. The statuses below serve as the campaign's map
rather than a victory lap. They are re-graded at each preview, and
several moved in the week this chapter was drafted. We publish the
grades at this granularity because the alternative, a coverage matrix
of green ticks, is how unification claims usually go wrong.

\paragraph{How to read this chapter.} The sections run in parallel,
with the same rubrics in the same order, so a reader may take them in
any order and sample by interest. A reader new to the catalogue should
read two entries in full. Ordinary differential equations
(section~\ref{sec:mt-ode}) carry the catalogue's simplest square, and
discrete-time Markov chains (section~\ref{sec:mt-dtmc}) carry the
square that is the engine most later squares reuse. After those two, every
other entry is a variation on a theme the reader already knows, each
adding one idea. Four entries repay attention even from a reader in a
hurry: the delay line (section~\ref{sec:mt-dde}), which dissolves an
apparently new semantic species into a wiring pattern over inventory
the catalogue already has; the demonic faces of
sections~\ref{sec:mt-mdp} and~\ref{sec:mt-ma}, where the adequacy
squares first become set-valued; the coloured nets of
section~\ref{sec:mt-sdcpn}, where the 2023 candidate unifier is itself
compiled into the catalogue it was once graded against; and the
quantum nets of section~\ref{sec:mt-qnet}, the one entry whose
set-valued face is demanded by a theorem of its own literature rather
than chosen by the modeller. The formal
blocks closing each entry serve as reference material, so skip them on
a first pass, and return to them when declaring a model.

\begin{figure}[t]
\centering
\begin{tikzpicture}[x=1mm,y=1mm,
  obj/.style={draw, rounded corners=1.2pt, align=center, font=\small,
    inner xsep=2.5mm, inner ysep=1.4mm,
    minimum width=46mm, minimum height=10mm},
  edge/.style={-{Stealth[length=1.8mm]}, semithick},
  lbl/.style={font=\footnotesize}]
  \node[obj] (sys)  at (0,26)  {System presentation\\ (mechanism and wiring rows)};
  \node[obj] (rows) at (76,26) {compiled constraints\\ (cells, walls, init rows)};
  \node[obj] (cls)  at (0,0)   {the frame's classical object\\ (solution set, path law, joint)};
  \node[obj] (su)   at (76,0)  {denotation in the\\ semantic universe};
  \draw[edge] (sys) -- node[lbl, above] {compile} (rows);
  \draw[edge] (rows) -- node[lbl, right] {denote} (su);
  \draw[edge] (sys) -- node[lbl, left] {interpret} (cls);
  \draw[edge] (cls) -- node[lbl, below] {include} (su);
  \node[font=\small] at (38,13) {$\cong$};
\end{tikzpicture}
\caption{The adequacy square that this chapter's Adequacy
paragraphs instantiate. At the two-sided entries the two routes agree by an
invertible globular cell; at the martingale-problem entries
(stochastic and jump-diffusion) the cell runs one way, from the
community's semantics to the certified verdict. The middle column of
the table below grades this square, frame by frame.}
\label{fig:mt-adequacy-square}
\end{figure}

\paragraph{The list at a glance.} The middle column grades each
entry's \emph{adequacy square}
(Figure~\ref{fig:mt-adequacy-square}), the mechanized statement,
carried in the entry's Adequacy paragraphs, that compiling a model
and then denoting it agrees with interpreting it classically and
then including the classical object. The outer columns grade
certification, and dashes mean nothing shipped.

\begin{center}\small
\begin{tabular}{@{}l>{\raggedright\arraybackslash}p{3.1cm}>{\raggedright\arraybackslash}p{3.6cm}>{\raggedright\arraybackslash}p{3.4cm}@{}}
\toprule
frame & adequacy square & certified today & principal gap \\
\midrule
static & definitional --- constant records (stated in the entry, not mechanized) & --- & a consuming certificate \\
ODEs & two-sided & controller certificates; validated tubes; one cross-frame composite & mechanism rows through the ingestion path \\
PDEs (deterministic) & two-sided, at every mechanism of the former & four exact dissipation certificates; bridge lemmas mechanized & classical existence (non-vacuity, never soundness) \\
SDEs & one-directional, at the martingale problem & generator certificates with correlated noise, $n \le 3$ & state-dependent noise at $n \ge 2$ \\
SPDEs & --- & moment + probability certificates, noise floors exact; derived covariance & trace-class instances; multiplicative noise \\
RODEs & two-sided, per draw and over the box & barrier uniform over the parameter box & the PDE side (RPDEs) \\
jump-diffusion & one-directional; generator splice hypothesis-free & all three field families, $n \le 2$ & certificates exercising multi-leg kernels \\
DTMCs & two-sided & the published-table campaigns & native rows; logic closure \\
CTMCs & two-sided (jump-chain face) & four fixtures, incl.\ unbounded reach & nested temporal operators \\
MDPs & two-sided, fixed-policy and demonic faces & fixed policy & the policy-quantified reproduction target \\
Markov automata & two-sided, fixed-scheduler and demonic faces & fixed scheduler & the scheduler-quantified reproduction target \\
open games & --- & --- & representation and semantics \\
mean-field games & --- & --- & everything; road through PDEs \\
GSHS & generator rungs only & a full-complement hybrid example & guarded composites at $n \ge 2$ \\
Petri nets & statics, race, cascade, and set-firing faces (schedulers declared and quantified) at the coloured entry; the untimed interleaving face two-sided at the classical entry & token game run natively; coverability & bounded-exploration hypothesis \\
Bayesian networks & two-sided & exact posteriors, finite case & continuous case; conditioning at scale \\
SCMs & substitution rung & counterfactual bounds, both directions & conditioning; latent confounders \\
MRFs & --- & partition-function brackets & continuous case \\
PPLs & --- & the discrete-time slice & conditioning \\
ProbLog & --- & --- & weighted-chase wiring \\
SBGMs & --- & --- & composing two shipped parts \\
\bottomrule
\end{tabular}
\end{center}

\medskip

A single content language carries the whole continuous and hybrid
family below, a manifest vocabulary of modes with domains, drift,
diffusion, multi-driver noise, jumps, and guarded transitions with
resets, and the types differ only in which fields are populated.
Content-addressed references carry neural and piecewise content (a
network controller in the loop, a learned drift) on any of these
profiles, and where a polynomial view coexists with network content,
their agreement is a recomputation obligation, never a trusted
row. Discrete-kernel types present outcome rows, while nets and
graphical models carry their own small theories, described in
place.

\section{The System Core}
\label{sec:mt-core}

\paragraph{The entry.} The substrate every profile populates: the
manifest root, its two strata, and the family view that makes
``profile, not tag'' an operational sentence rather than a slogan. A
System manifest is a composition term --- a \emph{mechanism} (the
primitive-dynamics stratum: field families over declared state
variables plus an interface binding) or a \emph{composite} (the wiring
stratum: component Systems plus variable-identification rows and the
composite-level descriptors of
chapter~\ref{ch:qtsl-composition-patterns}) --- the
mechanisms-and-wirings reading of systems whose doctrinal home is
the double operadic theory of Libkind and Myers \cite{libkind25}. Every field family of a
mechanism is individually optional at the type level; the free-form
entry point never refuses, because any family combination is a legal
\emph{presentation} --- conformance to a named profile is a separate,
optional check (the presets of the sections below), and a mechanism
whose populated set matches no preset is a first-class citizen. That
check, in plain words and stated here once for the whole chapter: a
manifest \emph{conforms} to a model type when it populates that
type's field families and no others beyond the shared substrate
($\Sigma$, below); the preset signatures quoted in each entry's API
block are the mechanized form of exactly that check. No
time-semantics tag is stored anywhere: continuous, counter, or hybrid
is derivable from which families are populated, and storing it would
be exactly the per-model doctrine tag this catalogue rules out.

\paragraph{Trust posture.} Presets and the preflight walk are
courtesies --- profile notes and dangling-reference warnings with path
geometry, never verdicts; the checker's replay is the only authority
on content. On the wire the manifest is CBOR under a schema-versioned
envelope with exact scalars; the committed golden vectors are the
cross-language byte authorities, and the TypeScript surface reaches
the same Rust implementation in-process, so the two surfaces cannot
diverge.

\syscatmeta{\begin{mathpar}
\inferrule*[Right=\rulename{Mechanism}]{
  \vec{x} : \overrightarrow{\mathsf{StateVar}} \\
  \Sigma_m \text{ substrate rows over } \vec{x}
  \;\; (\Sigma\text{, below}) \\\\
  \Phi \subseteq \mathcal{F} \\
  F_\phi : \overrightarrow{\mathsf{Row}_\phi}
  \quad (\phi \in \Phi)
}{
  \mathsf{Mechanism}\big(\vec{x};\, \Sigma_m;\,
  (F_\phi)_{\phi \in \Phi}\big) : \mathsf{System}
}
\and
\inferrule*[Right=\rulename{Composite}]{
  S_1, \ldots, S_k : \mathsf{System} \\
  W : \overrightarrow{\mathsf{WiringRow}} \\
  E : \overrightarrow{\mathsf{SwitchEdge}} \\
  A : \overrightarrow{\mathsf{ScheduleAssumption}}
}{
  \mathsf{Composite}(S_1, \ldots, S_k;\, W;\, E;\, A) :
  \mathsf{System}
}
\end{mathpar}
The two primitive System formers. The family alphabet $\mathcal{F}$
mirrors, group for group, the field families of the manifest schema
(the \texttt{MechanismInit} groups of the API block below):
{\small
\[
\begin{aligned}
\mathcal{F} \;=\;\;
  &\{\mathsf{modes},\, \mathsf{drift},\, \mathsf{lagrangian},\,
     \mathsf{diffusion},\, \mathsf{noise},\,
     \mathsf{noise\_spectrum},\, \mathsf{jumps},\, \mathsf{kill},\,
     \mathsf{transitions}\}
   && \text{continuous / hybrid} \\
  &\cup\; \{\mathsf{outcomes},\, \mathsf{init\_dist},\,
     \mathsf{rates},\, \mathsf{scheduler},\, \mathsf{circuit}\}
   && \text{discrete-kernel} \\
  &\cup\; \{\mathsf{tick}\}
   && \text{counter time} \\
  &\cup\; \{\mathsf{petri},\, \mathsf{bag},\, \mathsf{sdcpn},\,
     \mathsf{stochastic\_wires},\, \mathsf{equations},\,
     \mathsf{interventions},\, \mathsf{exogenous},\,
     \mathsf{weighted\_program}\}
   && \text{relational / population}
\end{aligned}
\]}
Each family's row sort $\mathsf{Row}_\phi$ is the family's own,
given by the formation rule of the entry that consumes it. Every
entry below is a checked specialization of $\mathsf{Mechanism}$: its
head form carries the formation data that individuates its type, and
optional declarations (marked $?$) are consumed without appearing in
it. No premise constrains $\Phi$ --- any family set is well-formed
formation data, and a mechanism whose populated set matches no
preset is a first-class citizen. That the toolkit's free entry point
never refuses is the implementation face of this rule: a
forward-compatibility property (families added to $\mathcal{F}$
extend the schema without invalidating old manifests), not a logical
feature. On the wire, either former's term is carried in the
manifest record $\langle \mathsf{schema},\, v,\, \mathsf{name}?,\,
\mathsf{term} \rangle$. $\mathrm{pop}(m) = \Phi$ is a view,
recomputed on demand and never stored. The substrate families,
admitted alongside every former's own formation data without being
listed:
\[
\Sigma \;=\; \{\mathsf{state},\, \mathsf{params},\,
\mathsf{interface},\, \mathsf{init},\, \mathsf{hypernet\_refs},\,
\mathsf{pinning}\}.
\]}
\ruleblock{API}
\sourcenote{Producer surface of the \texttt{mtk} toolkit for this entry: the TypeScript arrangement (npm package \texttt{@maitria/mtk}) and the Rust catalogue crate (\texttt{crates/mtk-catalogue}) --- one Rust implementation underneath, reached from TypeScript in-process through a native binding, so the two surfaces cannot diverge. Signatures and constructions are extracted from the toolkit sources at every regeneration of this catalogue; each construction is quoted verbatim from the toolkit's conformance corpus, code whose emitted certificate bytes are verified byte-identical against the committed golden vectors (helper values named \texttt{fx\_*} or \texttt{f.*} are that corpus's shared fixtures, defined beside the quoted lines). This API constructs and emits; it never checks --- the checker re-derives every claim independently.}
\apilang{TypeScript}{system.buildSystem}
\apiprov{signature --- \texttt{ts/src/system.ts}}
\begin{tssrc}
export function buildSystem(init: MechanismInit & SystemOpts): System

export interface MechanismInit {
  // substrate
  readonly state?: ReadonlyArray<StateVar>;
  readonly jetVars?: ReadonlyArray<JetVar>;
  readonly params?: ReadonlyArray<ParamDecl>;
  readonly interface?: Interface;
  readonly init?: InitRegion;
  // continuous / hybrid
  readonly modes?: ReadonlyArray<Mode>;
  readonly drift?: ReadonlyArray<DriftRow>;
  /** S-ACTION: the Lagrangian family (dynamics derived, never
   * stored; coexisting `drift` rows carry a recompute obligation). */
  readonly lagrangian?: LagrangianDecl;
  readonly diffusion?: ReadonlyArray<DiffusionRow>;
  readonly noise?: ReadonlyArray<NoiseRow>;
  /** S-SPDE: noise-spectrum driver rows (populating either spectrum
   * field creates the family; the SPDE preset refuses a spectrum
   * whose `noiseTail` was never declared). */
  readonly noiseDrivers?: ReadonlyArray<NoiseDriverRow>;
  /** S-SPDE: the spectrum's required tail form — never defaulted. */
  readonly noiseTail?: TailForm;
  readonly jumps?: ReadonlyArray<JumpRow>;
  /** S-JUMPDIFF: killing rows (absent = conservative). */
  readonly kill?: ReadonlyArray<KillRow>;
  readonly transitions?: ReadonlyArray<GuardedTransition>;
  readonly hypernetRefs?: ReadonlyArray<HypernetRef>;
  // discrete-kernel
  readonly outcomes?: ReadonlyArray<OutcomeRow>;
  readonly initDist?: ReadonlyArray<InitDistRow>;
  readonly rates?: ReadonlyArray<RateEntry>;
  readonly scheduler?: SchedulerDecl;
  readonly circuit?: CircuitDecl;
  // tick
  readonly tick?: TickDecl;
  // structured-population / relational
  readonly petri?: PetriDecl;
  readonly bag?: BagDecl;
  readonly sdcpn?: SdcpnDecl;
  readonly stochasticWires?: ReadonlyArray<StochasticWire>;
  readonly equations?: ReadonlyArray<EquationRow>;
  readonly interventions?: ReadonlyArray<InterventionRow>;
  readonly exogenous?: ExogenousTable;
  readonly weightedProgram?: WeightedProgramDecl;
  // i.i.d. tier
  readonly pinning?: PinningDecl;
}
\end{tssrc}
\apiprov{construction --- \texttt{ts/tests/system-native.test.ts}}
\begin{tssrc}
const init = {
  name: "ou-1d",
  state: [{ id: 0, name: "x" }],
  drift: [{ var: 0, poly: mono(1, 0, 1, "-1") }],
  diffusion: [{ var: 0, poly: constant(1, "1/2") }],
  init: { kind: "point", coords: [BIG] },
} as const;
const viaPreset = sys.unwrap(sys.sde(init));
const viaFree = sys.buildSystem(init);
assert.ok(viaPreset.contentEquals(viaFree));
assert.equal(toHex(viaPreset.toCbor()), toHex(viaFree.toCbor()));
\end{tssrc}
\apiprov{The quoted corpus lines run the corpus's preset discipline as an assertion: a preset's output is byte-identical to \texttt{buildSystem}'s on the same mechanism --- presets add checking, never representation.}
\apilang{Rust}{mtk\_catalogue::system::SystemManifest}
\apiprov{signature --- \texttt{crates/mtk-catalogue/src/system/manifest.rs}}
\begin{rustsrc}
pub fn free(mechanism: Mechanism) -> Self

pub fn composite(composite: Composite) -> Self
\end{rustsrc}
\apiprov{construction --- \texttt{crates/mtk-catalogue/tests/builders.rs}}
\begin{rustsrc}
assert_eq!(known_profile(&drift_only()), Some("ODE"));
let free = SystemManifest::free(drift_only());
// Nothing on the wire stores the classification: same bytes with
// or without the preset having been consulted.
let via_preset = Ode::manifest(drift_only()).unwrap();
assert_eq!(free.to_cbor().unwrap(), via_preset.to_cbor().unwrap());
\end{rustsrc}
\apiprov{The quoted corpus lines assert the profile-not-tag doctrine on the wire: same bytes with or without the preset having been consulted.}

\section{Tick Declarations}
\label{sec:mt-tick}

\paragraph{The entry.} The composition-critical counter-time
vocabulary. Time is base-universal --- composites live on real time
--- so a counter-time mechanism composed into real time declares its
tick as a timer mechanism whose reset parameter is either a point
(\emph{exact tick}) or a box (\emph{jitter bracket}). The distinction
is part of the model type, not a checker convention, and it has been
priced by certificate: a lander free-flight pair certifies the same
funnel under both semantics from one value table and one cover, with
the cost of clock jitter a factor of $2.66$ in certified
ticks-to-funnel ($1631$ against $614$), and a factor $2.92$ in the
certified elapsed bound, since a jittered tick may run a tenth longer
than an exact one. Both horizons are upper bounds carried by the
certificate, not measured times. A discrete-time component in
isolation remains
schedule-generic; its schedule class becomes \emph{assumption data}
exactly at the discrete/continuous composition boundary, through the
hybrid-time pattern of
chapter~\ref{ch:qtsl-composition-patterns}.

\syscatmeta{\begin{mathpar}
\inferrule*[Right=\rulename{exact}]{
  \mathsf{period} \in \mathbb{Q}_{>0}
}{
  \mathsf{exact}(\mathsf{period}) : \mathsf{TickDecl}
}
\and
\inferrule*[Right=\rulename{jitter}]{
  \mathsf{lo},\, \mathsf{hi} \in \mathbb{Q}
}{
  \mathsf{jitter}([\mathsf{lo}, \mathsf{hi}]) : \mathsf{TickDecl}
}
\end{mathpar}
Not a System former: the conclusion is at the declaration's own
sort. A
$\mathsf{TickDecl}$ is carried either on the mechanism itself
($\mathsf{tick} \in \mathrm{pop}(m)$, admitted by the counter-time
profiles) or in a composite's assumption slot as
$\mathsf{ScheduleAssumption} = \langle \mathsf{component},\,
\mathsf{tick} \rangle$.}
\ruleblock{API}
\sourcenote{Producer surface of the \texttt{mtk} toolkit for this entry: the TypeScript arrangement (npm package \texttt{@maitria/mtk}) and the Rust catalogue crate (\texttt{crates/mtk-catalogue}) --- one Rust implementation underneath, reached from TypeScript in-process through a native binding, so the two surfaces cannot diverge. Signatures and constructions are extracted from the toolkit sources at every regeneration of this catalogue; each construction is quoted verbatim from the toolkit's conformance corpus, code whose emitted certificate bytes are verified byte-identical against the committed golden vectors (helper values named \texttt{fx\_*} or \texttt{f.*} are that corpus's shared fixtures, defined beside the quoted lines). This API constructs and emits; it never checks --- the checker re-derives every claim independently.}
\apilang{TypeScript}{system.TickDecl}
\apiprov{signature --- \texttt{ts/src/system.ts}}
\begin{tssrc}
/** S-TICK: tick declaration — exact tick or jitter bracket. Tick
 * semantics is part of the model type, not a checker convention. */
export type TickDecl =
  | { readonly kind: "exact"; readonly period: QLike }
  | { readonly kind: "jitter"; readonly lo: QLike; readonly hi: QLike };
\end{tssrc}
\apiprov{construction --- \texttt{ts/tests/system-native.test.ts}}
\begin{tssrc}
return sys.unwrap(
  sys.dtmc({
    name: "chain-3",
    outcomes: [
      { src: 0, dst: 0, prob: "1/3", reward: "0" },
      { src: 0, dst: 1, prob: "2/3", reward: "1" },
      { src: 1, dst: 2, prob: "1", reward: "5" },
      { src: 2, dst: 2, prob: "1" },
    ],
    initDist: [{ state: 0, prob: "1" }],
    tick: { kind: "exact", period: "1/50" },
  }),
);
\end{tssrc}
\apiprov{The quoted fixture declares an exact tick; the composite fixture's schedule assumption (S-HYBTIME) carries the jitter-bracket form.}
\apilang{Rust}{mtk\_catalogue::system::TickDecl}
\apiprov{signature --- \texttt{crates/mtk-catalogue/src/system/manifest.rs}}
\begin{rustsrc}
/// System Catalogue: Tick declaration — the composition-critical
/// counter-time vocabulary. A counter-time mechanism composed into
/// real time declares its tick as a timer mechanism whose reset
/// parameter is a point (**exact tick**) or a box (**jitter
/// bracket**). Tick semantics is part of the model type, not a
/// checker convention (the LunarLander pair is the standing price
/// demonstration: a factor 2.66 in certified ticks-to-funnel).
#[derive(Debug, Clone, PartialEq, Serialize, Deserialize)]
#[serde(rename_all = "snake_case")]
pub enum TickDecl {
    /// Exact tick: reset parameter a point.
    Exact { period: WitRat },
    /// Jitter bracket: reset parameter a box `[lo, hi]`.
    Jitter { lo: WitRat, hi: WitRat },
}
\end{rustsrc}
\apiprov{construction --- \texttt{crates/mtk-catalogue/tests/common/mod.rs}}
\begin{rustsrc}
let m = MechanismBuilder::new()
    .outcome(OutcomeRow {
        src: 0,
        action: None,
        dst: 0,
        prob: r(1, 3),
        reward: Some(r(0, 1)),
    })
    .outcome(OutcomeRow {
        src: 0,
        action: None,
        dst: 1,
        prob: r(2, 3),
        reward: Some(r(1, 1)),
    })
    .outcome(OutcomeRow {
        src: 1,
        action: None,
        dst: 2,
        prob: r(1, 1),
        reward: Some(r(5, 1)),
    })
    .outcome(OutcomeRow {
        src: 2,
        action: None,
        dst: 2,
        prob: r(1, 1),
        reward: None,
    })
    .init_dist(0, r(1, 1))
    .tick(TickDecl::Exact { period: r(1, 50) })
    .finish();
presets::Dtmc::manifest(m)
    .expect("fixture conforms to the DTMC profile")
    .named("chain-3")
\end{rustsrc}

\section{Static Systems}
\label{sec:mt-static}

\paragraph{The frame.} A value chosen once and never revised: the
unknown constant of an estimation problem --- a parameter picked
demonically from a box and frozen for all time --- and the token
System of a colourless Petri net under the Bag former of
section~\ref{sec:mt-bag}, where Bag over a point is the naturals
with counting semantics. The bread-and-butter question is
membership: which values the initial region admits, and what a
surrounding composite may conclude from the box that bounds them.

\paragraph{Profile.} Substrate only --- state variables, unboxed or
boxed parameters, interface, demonic initial region --- with no
dynamics family populated: the bottom of the manifest family; every
other profile of this chapter populates families beyond it. At any
dimension it is the zero-drift degenerate of the drift-only frame;
at dimension zero it is the only dynamics-capable reading, because
the zero drift vector there \emph{is} the empty vector --- a drift roster
empty by dimension, not by omission. Time: continuous, trivially;
the trajectory never moves.

\paragraph{Semantics.} One convention is fixed here and consumed by
everything downstream: absent dynamics denote a \emph{frozen}
carrier, never an unconstrained one. Records are constant paths
valued in the initial region; the demonic initial face is the same
box convention as the deterministic siblings --- the denotation is
the union over the region of per-point singletons --- so the record
set is the initial region, read as paths.

\paragraph{Adequacy.} Definitional at this frame, and stated rather
than mechanized: both routes of the square of
Figure~\ref{fig:mt-adequacy-square} name the constant records valued
in the initial region --- compilation contributes the initial row
and the frozen-carrier convention, classical interpretation
contributes the constant path at each admitted point, and there is
no analysis between them to be wrong about. At positive dimension
the same record set is carved by the drift-only rows of
Problem~\ref{prob:mt-ode} at $f = 0$; at dimension zero the carrier
is a single point. No dedicated mechanization is claimed: the
frame's load-bearing content is the convention, not a theorem.

\paragraph{Status.} Vocabulary and remedy, not yet certificate: no
shipped certificate consumes the profile at this preview --- it
enters the inventory as the landing pad other entries point at.
Production: presets in both languages construct the manifest (the
API block below). The distinguishing shipped artefact is the refusal
pairing: a drift-requiring preset meeting a zero-dimensional,
substrate-only mechanism refuses with a remedy naming this profile,
and the property suite checks that refusal's geometry --- family
absent, state roster empty, and the named remedy accepting. One name
collision is worth disambiguating here: elsewhere in this chapter the
graphical-model sub-frames are called static because they have
\emph{no internal time axis} (section~\ref{sec:mt-pgm}); this
profile has the time axis, and nothing moves along it.

\syscatmeta{\begin{mathpar}
\inferrule*[Right=\rulename{S-STATIC}]{ }{
  \mathrm{Static} : \mathsf{System}
}
\end{mathpar}
The axiom of the manifest family: no dynamics family is consumed,
and the formation data beyond the substrate $\Sigma$ of
section~\ref{sec:mt-core} (state variables, possibly none;
parameters; interface; demonic initial region) is empty. Any
populated dynamics family is a stray --- the preset refuses it at
any dimension.}
\ruleblock{API}
\sourcenote{Producer surface of the \texttt{mtk} toolkit for this model type: the TypeScript arrangement (npm package \texttt{@maitria/mtk}) and the Rust catalogue crate (\texttt{crates/mtk-catalogue}) --- one Rust implementation underneath, reached from TypeScript in-process through a native binding, so the two surfaces cannot diverge. Signatures and constructions are extracted from the toolkit sources at every regeneration of this catalogue; each construction is quoted verbatim from the toolkit's conformance corpus, code whose emitted certificate bytes are verified byte-identical against the committed golden vectors (helper values named \texttt{fx\_*} or \texttt{f.*} are that corpus's shared fixtures, defined beside the quoted lines). This API constructs and emits; it never checks --- the checker re-derives every claim independently.}
\apilang{TypeScript}{system.staticProfile}
\apiprov{signature --- \texttt{ts/src/system.ts}}
\begin{tssrc}
/** The static profile: substrate only — state space, parameters,
 * interface, demonic initial region; no dynamics families. The
 * initial condition is chosen demonically from the initial region
 * and then never moves: the zero-drift degenerate of the ODE profile
 * at any dimension, the only reading at dimension zero (where the
 * zero drift vector IS the empty vector), the token System of a
 * colourless Petri net under the Bag former, and the
 * unknown-constant System (a boxed parameter frozen for all time).
 * (`static` is a reserved word in module code, so the export carries
 * the profile suffix.) */
export const staticProfile = profileFactory("static");
\end{tssrc}
\apiprov{The TS conformance corpus holds no construction through this factory at this preview --- byte authority for the profile is the Rust corpus opposite; the factory takes the same init-object vocabulary as every other profile.}
\apilang{Rust}{mtk\_catalogue::system::presets::Static::manifest}
\apiprov{signature --- \texttt{crates/mtk-catalogue/src/system/presets.rs}}
\begin{rustsrc}
preset!(
    /// System Catalogue: static profile — substrate only (state
    /// space, unboxed or boxed parameters, interface, demonic
    /// initial region), no dynamics families. The initial condition
    /// is chosen demonically from the initial region and then never
    /// moves: the zero-drift degenerate of the ODE profile at any
    /// dimension, and the only dynamics-capable reading at dimension
    /// zero, where the zero drift vector IS the empty vector. Two
    /// standing uses: the token System of a colourless Petri net
    /// under the Bag former (Bag over a point is the naturals with
    /// counting semantics), and the unknown-constant System (a value
    /// picked from a box, frozen for all time).
    Static, "static",
    req = [],
    allowed = [],
    extra = ok
);
\end{rustsrc}
\apiprov{construction --- \texttt{crates/mtk-catalogue/tests/builders.rs}}
\begin{rustsrc}
let m = MechanismBuilder::new()
    .init(InitRegion::Box {
        region: BoxRegion { axes: vec![] },
    })
    .finish();
let sys = Static::manifest(m.clone()).unwrap();

// The unknown-constant reading: a boxed parameter axis, frozen
// for all time — substrate-only, so the same profile.
let frozen = MechanismBuilder::new()
    .param(0, Some("theta"), Some([r(1, 2), r(3, 2)]))
    .finish();
Static::manifest(frozen).unwrap();
\end{rustsrc}
\apiprov{The quoted span is the builders-suite happy path: the zero-dimensional token System (an initial box with no axes --- the unique point of the zero-dimensional carrier) and the unknown-constant System (a boxed parameter, frozen for all time). The profile is also a refusal remedy: a drift-requiring preset meeting a zero-dimensional, substrate-only mechanism refuses with a pointer to this preset by name, and the refusal suite checks that the named remedy actually accepts.}

\section{Ordinary Differential Equations}
\label{sec:mt-ode}

\paragraph{The frame.} $\dot x = f(x)$ on a declared domain; the
bread-and-butter questions are invariance, reach-avoidance, and
Lyapunov decrease, plus validated enclosure of flows when the
question is quantitative.

\paragraph{Profile.} Drift only --- the noise-free, jump-free corner
of the manifest family. Carrier: the declared state variables over
exact-rational box or simplex domains; polynomial right-hand sides
native, non-polynomial ones by rationalization or by certified
abstraction (the latter designed, not shipped). Openness, when wanted,
enters as demonic input wires, not as a stochastic field. Time:
continuous.

\paragraph{Semantics.} Trajectory semantics --- the drift-only
specialization of the System semantics. Certificates about flows
(tubes, flow maps) are licensed by once-per-format meta-theorems: the
kernel checks inclusion-invariance sequents, and the reading of those
sequents as statements about solutions is a named, stated obligation
rather than kernel code. In the universe's terms: records over the
classical site, read at the precise or possibilistic face of the
denotation (Remark~\ref{rem:su-faces}) according as inputs are
absent or demonic.

\paragraph{Adequacy.} The drift-only frame carries the catalogue's
simplest instance of the square of
Figure~\ref{fig:mt-adequacy-square}, and its content is easy to
misjudge as vacuous: what compilation produces is not the textbook solution
concept but \emph{forward} constraint rows, and the square is the
theorem that the two meet. Stated self-contained, at the working
scale of an exam problem:

\begin{problem}[adequacy of the drift-only compilation]
\label{prob:mt-ode}
Fix a continuous $f : \mathbb{R}^n \to \mathbb{R}^n$
and $x_0 \in \mathbb{R}^n$. Records are paths $y : \mathbb{R} \to
\mathbb{R}^n$; only values at $t \ge 0$ are observable. Say $y$
satisfies the \emph{compiled rows} when (i) $y$ is continuous on
$[0, \infty)$ and at every $t \ge 0$ has the forward derivative
$f(y(t))$ --- the one-sided limit through times $\ge t$ --- and (ii)
$y(0) = x_0$. Say $y$ is a \emph{classical solution} when $y(0) =
x_0$ and $y$ is differentiable on $[0,\infty)$ with $\dot y = y
\fatsemi f$ (composition in diagrammatic order: first $y$, then $f$)
--- two-sided at interior times, one-sided only at the
corner $t = 0$. Show the two sets of records are equal, and show
that the continuity clause in (i) cannot be dropped: exhibit a
record with the correct forward derivative everywhere on $[0,
\infty)$ that solves nothing.
\end{problem}

\emph{Sketch.} Classical solutions satisfy the rows by restriction.
Conversely the rows put $y$ under the fundamental theorem of
calculus twice: continuity of $y$ (the record-class row) makes $t
\mapsto f(y(t))$ continuous, so the forward rows integrate --- FTC
for right derivatives --- to the Picard identity $y(t) = x_0 +
\int_0^t f(y(s))\,ds$, the integral-equation form of the classical
theory \cite{coddington-levinson-1955, hartman-2002}; then
differentiating the integral at interior times gives the two-sided
derivative. For the countermodel take $f = 0$, $x_0 = 0$, and the
step record $y = k\,\mathds{1}_{[t_0, \infty)}$ with $k \neq 0$ and
$t_0 > 0$: every forward
derivative on $[0,\infty)$ vanishes --- right-continuity is free ---
but left-continuity fails at $t_0$, so the record-class row is
load-bearing, not decorative.

In the universe's terms: the two constraint families --- the cell
and the initial row --- form a loose-arrow diagram at the
possibilistic stratum, and the theorem says the diagram is
isomorphic, by an invertible globular cell, to the support inclusion
of the classical solution set: the same set-valued convention as the
delay line (section~\ref{sec:mt-dde}), of which this square is the
seam --- the delay-line composite's ODE-cell clause is this square's
row form at the field composed with the input signal, and the
random-parameter frame (section~\ref{sec:mt-rode}) specializes it per
draw. Three boundaries are stated rather than
smoothed. Existence is not owed: the compiled atoms never exclude
the vacuous model, blow-up in finite time ($\dot y = y^2$) can empty
the solution set over $[0,\infty)$ outright, and the constructive
properness witness is the equilibrium record at a zero of the drift.
Determinism is not a constraint-stratum fact: the atoms denote the
solution \emph{set}, and uniqueness enters as a theorem under the
frame's own declared-domain discipline --- Lipschitz drift on the
domain, solutions confined to it --- consumed from the mechanized
Gr\"onwall interface, whose hypotheses are the forward rows
verbatim: the community's own uniqueness engine speaks constraint
rows, not solution operators. And the connector consumes a
completeness hypothesis on the carrier (the two FTC directions integrate), which
the frame's finite-dimensional state spaces satisfy; the bundling,
the pinning theorems, and the wiring below do not consume it. The
demonic initial face is the box convention of the mechanism
declarations: its denotation is the union of the per-point solution
sets, the deterministic frame's shadow of the credal unions the
stochastic squares use. The statement and both directions are
mechanized, sorry-free, in the toolkit's Lean tree
(\texttt{lean/su-functor/}, module \texttt{SuFunctor.Adequacy.Ode}):
the connector is \texttt{isOdeSol\_iff\_integralCurve} --- the
classical side is the mechanized analysis library's own
integral-curve concept, not a private restatement --- with the
Picard form as \texttt{isOdeSol\_iff\_integral}, the invertible cell
is \texttt{ode\_adequacy} with the box face
\texttt{ode\_box\_adequacy}, the pinning theorems are
\texttt{odeSols\_eqOn\_of\_lipschitzOnWith} and its global
corollary, and the wiring into the semantic universe's judgment
(\texttt{su\_ode\_adequacy}) runs through the same transport layer
as the sibling squares, with one possibilistic contrast stated deliberately:
the zero law's ideal point models every compiled presentation ---
support constraints never force realization, and the demand that the
system realize anything at all is judgment-level content.

\paragraph{Status.} Certification is in place twice over. Closed-loop
controller certificates --- an inverted-pendulum family, and a
laddered re-verification of an aircraft collision-avoidance
controller family --- check through the unmodified kernel. The
validated-enclosure species ladder de-trusts an established rigorous
integrator: candidates from any proposer are re-founded on exact
fixed-point inclusion checks, so the assurance story never rests on
trusting interval arithmetic (the worked example of
section~\ref{sec:ckbk-proposer-disposer}). This frame also holds the
inventory's only cross-frame composition witness at this preview: an
ODE plant wired to a discrete-time controller, with the
assume--guarantee square transporting the certificate across the
boundary and the composite checked by the unmodified kernel ---
every mutant of the composite refuses. Production: the toolkit's
producer surface constructs presentations in both languages and
derives the Lie-derivative row $\mathcal{A}V = \nabla V \cdot f$
symbolically from the presentation --- the exact polynomial whose
sign conditions the martingale-problem headlines of the stochastic
siblings (sections~\ref{sec:mt-sde} and~\ref{sec:mt-jumpdiff})
consume verbatim at their pure-drift face; derived and never
trusted, the checker re-derives it. The
double-integrator plant of the API block below is a committed
standalone cross-language byte-identity vector whose conserved-energy
zero the kit's tests pin in closed form in both languages, and the
kit is the jump-diffusion assembler at its degenerate face, the
delegation warranted by the mechanized pure-drift degeneracy. Open:
mechanism declarations exist as staged rows, but the ingestion path
from rows to a checked model is not yet exercised end to end --- a
gap this frame shares with most of the list, named once here and
cited below as ``the ingestion gap''; and the enclosure ladder's
upper rungs rest on meta-theorems outside the kernel, an explicit
trust boundary the species documentation states.

\syscatmeta{\begin{mathpar}
\inferrule*[Right=\rulename{S-ODE}]{
  \vec{x} \text{ declared state variables} \\
  D : \overrightarrow{\langle \mathsf{var},\, \mathsf{poly} \rangle}
  \text{ drift rows, each } \mathsf{poly} \in \mathbb{Q}[\vec{x}]
}{
  \mathrm{ODE}(\vec{x};\, D) : \mathsf{System}
}
\end{mathpar}
Drift is the only dynamics family consumed. Openness, when wanted,
enters as demonic input wires on $m.\mathsf{interface}$ --- interface
wires are substrate, never a stochastic family. At dimension zero the
drift family cannot be populated at all --- the roster is empty by
dimension --- and the preset refuses with a remedy naming the static
profile (section~\ref{sec:mt-static}).}
\ruleblock{API}
\sourcenote{Producer surface of the \texttt{mtk} toolkit for this model type: the TypeScript arrangement (npm package \texttt{@maitria/mtk}) and the Rust catalogue crate (\texttt{crates/mtk-catalogue}) --- one Rust implementation underneath, reached from TypeScript in-process through a native binding, so the two surfaces cannot diverge. Signatures and constructions are extracted from the toolkit sources at every regeneration of this catalogue; each construction is quoted verbatim from the toolkit's conformance corpus, code whose emitted certificate bytes are verified byte-identical against the committed golden vectors (helper values named \texttt{fx\_*} or \texttt{f.*} are that corpus's shared fixtures, defined beside the quoted lines). This API constructs and emits; it never checks --- the checker re-derives every claim independently.}
\apilang{TypeScript}{system.ode}
\apiprov{signature --- \texttt{ts/src/system.ts}}
\begin{tssrc}
/** S-ODE: drift only (openness, if any, enters as demonic input wires
 * on the interface). Checks the family signature, then wraps into the
 * same manifest type as every other entry point — presets add
 * checking, never representation. */
export const ode = profileFactory("ode");

/** The derived ODE generator kit: the exact Lie-derivative
 * polynomial 𝒜V = ∇V·f. Derived, never trusted — the checker
 * re-derives it. */
export interface OdeGenKit {
  readonly gen: PolyWitness;
}

/** Derive the ODE generator kit from a mechanism declaration and a
 * certificate observable. Admission is the ODE profile's — diffusion,
 * jump, or killing rows refuse naming the stray family (those
 * presentations belong to `sdeGenKit` / `jdGenKit`); a parameter
 * roster is substrate and does not refuse, while drift naming a
 * parameter axis refuses at the arity fence (that presentation
 * belongs to `rodeGenKit`) — and the computed polynomial's identity
 * with the semantic generator at the drift-only face is mechanized
 * in the toolkit's Lean tree, whose ODE certificate theorems consume
 * the derived row's sign conditions verbatim. */
export function odeGenKit(
  init: MechanismInit,
  cert: Polynomial,
): Result<OdeGenKit, Refusal> {
  const raw = native.sysOdeGenKit(
    mechanism(init).nativeHandle,
    polyToNative(cert),
  ) as { kit: OdeGenKit } | { refusal: unknown };
  return "refusal" in raw
    ? refused(narrowRefusal<Refusal>(raw.refusal))
    : ok(raw.kit);
}
\end{tssrc}
\apiprov{construction --- \texttt{ts/tests/system-native.test.ts}}
\begin{tssrc}
return sys.unwrap(
  sys.ode({
    name: "double-integrator",
    state: [
      { id: 0, name: "pos" },
      { id: 1, name: "vel" },
    ],
    interface: {
      inputs: [{ id: 0, name: "u" }],
      outputs: [{ id: 1, name: "y" }],
    },
    drift: [
      { var: 0, poly: mono(2, 1, 1, "1") },
      { var: 1, poly: mono(2, 0, 1, "-1") },
    ],
  }),
);
\end{tssrc}
\apilang{Rust}{mtk\_catalogue::system::presets::Ode::manifest}
\apiprov{signature --- \texttt{crates/mtk-catalogue/src/system/presets.rs}}
\begin{rustsrc}
preset!(
    /// System Catalogue: ODE profile — drift only (openness, if any,
    /// enters as demonic input wires on the interface).
    Ode, "ODE",
    req = [Drift],
    allowed = [],
    extra = ok
);
\end{rustsrc}
\apiprov{construction --- \texttt{crates/mtk-catalogue/tests/common/mod.rs}}
\begin{rustsrc}
let m = MechanismBuilder::new()
    .state_var(0, Some("pos"))
    .state_var(1, Some("vel"))
    .interface(Interface {
        inputs: vec![WireDecl {
            id: 0,
            name: Some("u".to_owned()),
            sort: None,
        }],
        outputs: vec![WireDecl {
            id: 1,
            name: Some("y".to_owned()),
            sort: None,
        }],
    })
    .drift(None, 0, mono(2, 1, 1, r(1, 1)))
    .drift(None, 1, mono(2, 0, 1, r(-1, 1)))
    .finish();
presets::Ode::manifest(m)
    .expect("fixture conforms to the ODE profile")
    .named("double-integrator")
\end{rustsrc}
\apiprov{The construction is the committed \texttt{double-integrator} golden vector (fixture \texttt{ode}, standalone byte authority; the same presentation serves the composite golden as its plant component), and it doubles as the generator kit's conserved-energy closed-form control (\(\mathcal{A}(\mathrm{pos}^2+\mathrm{vel}^2)=0\)); the kit's Rust face is \texttt{mtk\_catalogue::system::ode\_gen\_kit}.}

\section{Partial Differential Equations}
\label{sec:mt-pde}

\paragraph{The frame.} Evolution equations $\partial_t u =
F(u, \partial_x u, \partial_x^2 u, \ldots)$; the bread-and-butter
certified questions are energy dissipation and invariance.

\paragraph{Profile.} A differential-polynomial mechanism over
declared jet variables and a spatial domain --- a sibling extension
at the kernel's own level: integration-by-parts structure enters as
exact polynomial-syzygy rows, pointwise nonnegativity as cone rows,
and the checker gains nothing beyond jet arithmetic it already has.
Time: continuous. Deterministic; the stochastic half is
section~\ref{sec:mt-spde}'s, which adds exactly one family to this
profile.

\paragraph{Semantics.} A solution-concept \emph{rung} --- classical,
strong, or weak --- with the semantic bridge (the finitely many
analysis lemmas licensing the algebra-to-solutions reading) entering
as named obligations, since largely discharged for this frame (see
Status). This obligation discipline was established here first and
the rest of the inventory inherits it.

\paragraph{Adequacy (heat on a cyclic mesh).} Adequacy for a composed
PDE says that compiling a spatial mesh of cells and walls, then
denoting, agrees with interpreting the equation globally and
including its solution set --- and the first instance where the mesh
is genuinely a \emph{cover} is the heat equation on a circle, cut
into three overlapping arcs.  The cover is cyclic: the arcs meet in
three collars, and the third collar closes the ring through the
periodic identification.  Cyclic covers are where local-to-global
inference gets dangerous --- Vorob'ev exhibited
pairwise-consistent families of finite marginals over a three-cell
ring with no joint law at all \cite{vorobev-1962} --- so the
instance is chosen to make the mesh earn its consistency rather than
inherit it.

\begin{problem}[heat on a cyclic mesh]
\label{prob:mt-pde-heat}
Fix a diffusivity $D$, a circumference $L$, arc
breakpoints $0 < a < b < L$, a collar half-width $0 < \delta < L$,
and an open time window $T$.  On the universal chart, the three arcs
are $I_1 = (-\delta, a+\delta)$, $I_2 = (a-\delta, b+\delta)$,
$I_3 = (b-\delta, L+\delta)$.  A mesh record is a triple of cell
fields $(u_1, u_2, u_3)$, and the compiled constraints are six: on
each arc, the field is continuous, twice differentiable in space,
and satisfies $\partial_t u_i = D\,\partial_{xx} u_i$ pointwise (the
cell rows); $u_1 = u_2$ on $I_1 \cap I_2$ and $u_2 = u_3$ on
$I_2 \cap I_3$ at observed times (the chain walls); and
$u_3(x) = u_1(x - L)$ on the ring-closing collar (the wrap wall).
Show that a record satisfies all six constraints if and only if
there is a single $L$-periodic global solution of the heat equation
agreeing with every cell field on its arc.  Show moreover that the
wrap wall is not implied by the rest: the linear field
$u_i(x,t) = x$ satisfies every cell row and both chain walls on all
three arcs, and fails the wrap wall alone.
\end{problem}

\emph{Sketch.}  ($\Rightarrow$)  Paste the two chain seams first:
the piecewise field agrees with some cell field on a neighbourhood
of every chart point, the walls supplying agreement across the
collars, and the heat rows are germ-determined --- differential
operators are local --- so the paste solves on the chart.  Then
periodize: the wrap wall is consumed exactly once, as the
consistency identity of the periodic extension at the ring-closing
collar, and locality carries the rows to the whole line.
($\Leftarrow$) Clipping a global periodic solution to the arcs
satisfies the cell rows by restriction, the chain walls because all
three cells read one field, and the wrap wall by periodicity.  No
parabolic theory appears in either direction --- no semigroup, no
regularity, no maximum principle: the square lives at the record
stratum, where gluing is descent plus locality.

In the universe's terms: the six constraint families form a
loose-arrow diagram of the thin double category of per-element
constraints, walls and cells alike read at the possibilistic
stratum, conjoined once at the mesh joint --- and composition order
is immaterial, since satisfaction of a conjunction reads through
membership only.  The theorem says the diagram is isomorphic, by an
invertible globular cell, to the support inclusion of the global
periodic solution set; as in the transport instance, no constraint
forces mass, and the demand that the mesh realize anything is
judgment-level, discharged constructively by the equilibria.  The
cyclic cover's extra obligation is a \emph{holonomy} constraint,
carried as an ordinary refusable wall: the linear field --- a
perfectly good local solution family that fails to descend to the
circle --- passes every per-cell and pairwise-chain check and is
refused by the wrap wall alone, and no global periodic field
restricts to it.  This is the two rings' contrast made precise: the
frustrated marginal ring fails invisibly, at the measure tier, where
consistency of laws carries no global witness and the lawful credal
denotation collapses to the bottom object; the heat ring's
obligations are record-tier, local, and visible as walls, its
denotation is satisfiable, and local-to-global gluing is a theorem.
The statement and both directions are mechanized, sorry-free, in the
toolkit's Lean tree (\texttt{lean/su-functor/}, module
\texttt{SuFunctor.Adequacy.HeatTorus}): the ring glue is
\texttt{heat\_mesh\_glues}, the invertible cell is
\texttt{heat\_adequacy}, the holonomy witness pair is
\texttt{chainAtoms\_dirac\_lin} with \texttt{not\_wrap\_lin}, the
interval control (cutting the ring trivializes the obligation) is
\texttt{heat\_chain\_glues}, and the marginal-ring obstruction and
its credal absorption are \texttt{vorobev\_no\_joint} and
\texttt{credalCompile\_vorobevGlue} in
\texttt{SuFunctor.Adequacy.MarginalGlue}, with the two rings'
contrast packaged as a single statement,
\texttt{heat\_ring\_glues\_vorobev\_ring\_not}.  Two strata are named
separate, deliberately: with finitely many test rows in place of the
pointwise rows, the compiled constraints are implied by
solution-hood and do not pin it, the equivalence returning only in
the mesh-refinement limit; and the exact-rational certificate
stratum carries the dynamics as outward-rounded enclosures, never
equalities --- $e^{-\lambda t}$ is transcendental at rational times,
so the only exactly-propagated mode is the conserved mean.  Neither
stratum is mechanized; both carry recorded cost estimates ---
estimates, not available squares.  The sign
of $D$ is consumed by neither direction of the square: certified
decay, and with it the direction of well-posedness, is
enclosure-stratum content, not a satisfaction clause.

\paragraph{Adequacy (the initial face as a port).} The mesh square
above keeps its time windows open --- interior discipline; no
statement mentions $t = 0$.  The suite's seventh square wires the
initial condition as an \emph{exposed port}: the initial face
enters as a boundary trace along the window's own filter --- the
field on each spatial ray attains the port datum in the limit from
inside the window, never as a row at an interior point --- and
three port atoms, one per arc, extend the six mesh atoms to nine.
The ported square closes, still two-sided at the pointwise-record
stratum.  Its most important content is a deliberate absence:
\emph{feeding} the port --- exhibiting, for a given datum, a global
periodic solution attaining its trace --- is a named definition,
and the consistency discharge is conditional on a feed, because for
general admissible data the feed \emph{is} the heat semigroup on
the circle \cite{evans-2010}: enclosure-stratum content, cited and
never mechanized
--- the heat analogue of the chain square's Ionescu--Tulcea
stratum, with the unconditional consistency theorem deliberately
absent.  The constructive floor is mechanized without any
semigroup: equilibria feed constant data, and a single Fourier mode
feeds a genuinely nonconstant datum --- what the semigroup buys is
only the limit over infinitely many modes for merely-continuous
data.  Compatibility conditions on the datum are \emph{derived},
never assumed: the wrap wall plus the traces force the datum's
holonomy identity on the ring-closing collar --- the delay line's
corner-condition lesson recurring one cover rung up --- and
periodicity of a global solution forces the same identity from the
classical side.  The face's only new fee is the non-triviality of
the window's boundary filter --- without it the traces are vacuous
and the compatibility statements fail, so the fee is load-bearing
--- discharged for the concrete window.  The sign of $D$
remains nowhere, even at the port: backward heat feeds the same
single mode, growing, so the record stratum still cannot see the
direction of well-posedness --- uniqueness of the fed solution, and
its genuine failure backward, are enclosure-stratum business, cited
with the semigroup.  Mechanized, sorry-free, beside the mesh square
(\texttt{lean/su-functor/}, module
\texttt{SuFunctor.Adequacy.HeatTorusInit}): the invertible cell is
\texttt{heat\_init\_adequacy}, the conditional consistency
discharge is \texttt{heatInitAtoms\_consistent\_of\_feed}, the
derived wrap-compatibility pair is \texttt{initFace\_wrap\_compat}
with its classical-side twin
\texttt{IsHeatTorusSolInit.trace\_wrap\_compat}, the
constructive floor is \texttt{portFeed\_const} and
\texttt{portFeed\_mode} (with \texttt{modeDatum\_ne\_const}), and
the filter fee's discharge is \texttt{initFace\_neBot}.

\paragraph{Adequacy (solution sets on the mesh).} Where a single
field glues, a \emph{set} of fields need not: nonuniqueness enters
the mesh as imprecision --- each cell carrying a closed set of local
solutions, its observable face the union of their graphs --- and at
that stratum the descent question splits cleanly in two.  Descent of
the closed sets themselves is exact: over any open cover of the
window, closed subsets of the spacetime--value ambient correspond
homeomorphically, in the Fell topology, to overlap-compatible
families of closed subsets of the cells --- a sheaf condition in
equalizer form, demanding neither directedness nor finiteness of the
cover --- and the ring's three arcs are its first working instance.
What fails is citizenship transport.  The witness is the integer
orbit of the linear field --- all period-translates at once, a
closed set because the branch set is discrete.  Its every clip is a
union of cell-solution graphs; its clip family is
overlap-compatible; it satisfies the set-level wrap wall (the orbit
is invariant under the period shift, so the wall that refused the
single linear field is blind to the full orbit); it is even a union
of window-solution graphs.  Yet it is the graph-union of no family
of periodic solutions whatsoever: a global solution threading the
tube is pinned to a single branch --- continuity against a discrete
branch set --- and single branches do not close up.  The gluing that
one refusable wall purchased at the record stratum is purchasable by
no pointwise wall at the record-set stratum: the missing datum is
the pairing of branches across seams --- which branch continues
which around the ring, correspondence data subject to a cocycle
condition --- and forgetting exactly that labelling is the whole
failure, since labelled families glue member by member through the
record-stratum theorem.  Consistency, by contrast, is free at this
stratum: the top tube --- every value at every point, maximal
imprecision --- is the graph-union of all constant equilibria and
satisfies every wall, so the consistency obligation concentrates
entirely on non-vacuous imprecision.  And the failure remains
detectable one stratum down, as consistency of a composed object
rather than as any wall's refusal: compiling the orbit tube to
per-cell enclosure constraints and conjoining with the mesh's own
rows forces the empty model --- the branch indices are frustrated
around the ring by an integer cocycle, the seams forcing them equal
and the wrap wall forcing a shift.  All of it is mechanized,
sorry-free, beside the heat module (\texttt{lean/su-functor/},
module \texttt{SuFunctor.Adequacy.SolutionSetMesh}): the descent
equivalence is \texttt{FellCloseds.equalizerHomeomorph} instantiated
as \texttt{meshDescent}, the packaged witness is
\texttt{setStratum\_holonomy} with
\texttt{orbitTube\_not\_torusTube}, the failure of local character
is \texttt{not\_isLocalLocusEq\_torusTube}, the one-sidedness of the
compiled square is \texttt{tubeAtoms\_not\_classical} against the
forward cell \texttt{solset\_forward\_cell}, the free consistency
witness is \texttt{isTorusTube\_univ}, the labelled positive face is
\texttt{labelled\_ensembles\_glue}, and the record-tier shadow is
\texttt{orbit\_enclosure\_forces\_zero}.

\paragraph{Adequacy (the frame, at every mechanism).}  The three
squares above are instances --- heat, with one operator hard-coded.
The suite's frame-level statement proves the mesh square once for
\emph{every} mechanism of the formation rule, and in doing so
isolates exactly what a solution concept owes the mesh: two lemmas.

\begin{problem}[the mesh square, at every mechanism]
\label{prob:mt-pde-frame}
Let a mechanism assign to each of finitely many
components a drift polynomial in the components' spatial jets through
a declared order, and take as its pointwise rows: joint continuity of
each component, spatial differentiability of each iterated derivative
below the declared order in canonical-value form, and the drift
equations --- each component's time derivative equals its polynomial
read on the jets.  On the cyclic three-arc mesh, with the same six
constraint families as the heat square, show: (a) the two
\emph{interface lemmas} hold uniformly in the mechanism --- germ
locality (agreement on an open rectangle transports the rows) and
spatial homogeneity (translates of solutions solve at translated
points); (b) for any solution concept satisfying the two lemmas, a
mesh record satisfies all six constraints if and only if a single
$L$-periodic global solution of the concept agrees with every cell
field on its arc; (c) constants at roots of the drift's reaction
structure satisfy the compiled conjunction, and $1/2$ is not a root
of Fisher--KPP's.
\end{problem}

\emph{Sketch.}  (a) Locality is an induction on jet order ---
agreement on an open set transports every iterated derivative,
pointwise on the set --- and homogeneity an induction from the
first-order translation identity, junk values included; neither
consumes any property of the polynomial beyond its evaluation.
(b) The forward direction is the heat square's paste-and-periodize
argument verbatim, with the interface fields in place of the heat
lemmas: the chain paste consumes locality alone, the ring-closing
periodization additionally consumes homogeneity --- locality buys the
interval, homogeneity buys the holonomy leg.  Backward is clipping.
(c) The jet read of a constant is its value at order zero and zero
above, so the drift equations reduce to the reaction structure's
root condition; evaluating Fisher--KPP's reaction at $1/2$ gives
$1/4 \ne 0$.

The mechanized statement is parametric twice over.  Over
\emph{mechanisms}: the heat square is re-derived as a specialization
--- the compiled atoms of the heat instance are the frame square's at
the heat concept definitionally, and the heat rows are the jet rows
of the declared second-order mechanism, so the instance's hard-coded
operator and the formation rule's jet data present the same
constraints --- while viscous Burgers, Fisher--KPP, and a
cross-coupled second-order pair receive their squares as one-line
corollaries, with zero new inductions.  Over \emph{solution
concepts}: a rung enters as data --- its rows in germ form together
with the named obligations the square consumes (locality, at weak
rungs the partition-of-unity argument of the literature
\cite{evans-2010}; homogeneity; the strong-rung inclusion) --- and
every rung meeting them inherits the two-sided square, plus a
one-directional cell from the pointwise rung's compiled mesh to its
own, the solution-concept ladder's enclosure-shaped transfer.  At the
generality of the formation rule there is no canonical weak form ---
a fully nonlinear differential polynomial admits none, and weak
formulations are per-structural-class --- so the rung interface is
the strongest true frame-level statement, with each class's community tie an
instantiation obligation.  Consistency is equilibrium-conditional at
this level: heat and Burgers admit every constant, Fisher--KPP
exactly two, and a mechanism without real equilibria defers
non-vacuity to existence theory, the named separate stratum.  The
sign blindness of the heat square is now uniform: no statement
consumes any property of the drift polynomial beyond evaluation.
Mechanized, sorry-free, beside the heat module
(\texttt{lean/su-functor/}, module \texttt{SuFunctor.Adequacy.Pde}):
the concept-generic square is \texttt{LocalRows.mesh\_adequacy}, the
interface lemmas are \texttt{PdeAt.congr} and
\texttt{PdeAt.comp\_sub\_const}, the frame square is
\texttt{pde\_adequacy}, the heat re-derivation is
\texttt{heat\_adequacy\_of\_mesh} with the jet-grain tie
\texttt{heatAt\_iff\_pdeAt}, the instance corollaries are
\texttt{burgers\_adequacy}, \texttt{kpp\_adequacy}, and
\texttt{coupled\_adequacy}, the rung ladder is \texttt{PdeRung} with
\texttt{rung\_adequacy} and \texttt{pde\_rung\_transfer}, the
consistency witness is
\texttt{pdeAtoms\_consistent\_of\_equilibrium} with the
non-equilibrium receipt \texttt{kpp\_not\_equilibrium\_half}, and the
semantic-universe wiring is \texttt{su\_pde\_adequacy}.

The suite also carries the \emph{declared-data} face of the same
statement (module \texttt{SuFunctor.Adequacy.PdeRows}): a row
specification --- state variables, jet declarations, sparse
exponent-tuple drift rows over the positional index space
$\mathit{state} \mathbin{+\!\!+} \mathit{jets}$ --- compiles into the
formation rule's mechanism (\texttt{PdeRowSpec.toMechanism}), the
compile provably reads no position beyond a row's declared width
(\texttt{toMv\_congr},
\texttt{toMechanism\_congr\_inhabited}), and the four instance
mechanisms above are re-derived \emph{from declared rows} as
polynomial identities (\texttt{heatRowSpec\_toMechanism} and kin, at
the certified constants), so the squares and the semantic-universe
wiring specialize to presented data with no new proof
(\texttt{heat\_rows\_adequacy}, \texttt{su\_pde\_rows\_adequacy}).

\paragraph{Status.} Certification: four exact-rational artefacts ---
viscous Burgers energy dissipation at the sharp rate (every mutant of
its battery refuses), a Fisher--KPP certificate that is the family's
first to \emph{need} a box constraint, and Kuramoto--Sivashinsky
twice, checked by an independently implemented exact checker. To our
knowledge, exact-certified PDE dissipation bounds with replayable
proof objects do not appear in the surveyed literature, whose own
methods sections describe floating-point and semidefinite-programming
foundations; this frame was also a not-covered cell in the 2023 SDCPN
grading. The bridge lemmas, staged as named obligations when this
catalogue was first graded, have since been largely discharged in
mechanized form: axiom-clean Lean theorems now cover the classical
licensing lemmas, the Fourier bridges, the energy laws at their sharp
decay rates, uniqueness for viscous Burgers at the classical rung,
and the weak-rung transfer of every certificate conclusion --- the
last consuming nothing beyond the solution-class definition,
Parseval, and elementary iteration: a machine-checked confirmation
that certificate soundness never rested on the hard existence
theory. Representation and production: the toolkit constructs
jet-grain presentations in both languages and derives the canonical
differential-polynomial right-hand side from the declared rows (the
producer face of the mechanized compile edge, shared with the
stochastic sibling's kit, section~\ref{sec:mt-spde}); the
viscous-Burgers mechanism at jet grain is a committed cross-language
byte-identity vector, and the API block below quotes its receipted
construction. Deliberately open: one comparison-principle lemma and the
weak-rung energy-inequality variants of the fourth-order family,
each small and well-scoped; classical existence itself ---
needed for non-vacuity, never for soundness;
boundary-condition classes beyond the exercised torus and
mean-zero settings; certificate artefacts for PDE mechanisms inside
compositions (the mesh square above,
Problem~\ref{prob:mt-pde-heat}, carries the semantic square for the
first such composition); and the ingestion gap.

\syscatmeta{\begin{mathpar}
\inferrule*[Right=\rulename{S-PDE}]{
  \vec{x} \text{ declared state variables} \\
  J : \overrightarrow{\langle \mathsf{base},\, \alpha \rangle}
  \text{ jet variables} \\
  D \text{ drift rows over the extended index space } (\vec{x}, J)
}{
  \mathrm{PDE}(\vec{x};\, J;\, D) : \mathsf{System}
}
\end{mathpar}
Jet variables extend the polynomial index space (state variables
first, jet variables after), each declaring its base field variable
and spatial multi-index $\alpha$; drift rows over that space are the
differential-polynomial mechanism.}
\ruleblock{API}
\sourcenote{Producer surface of the \texttt{mtk} toolkit for this model type: the TypeScript arrangement (npm package \texttt{@maitria/mtk}) and the Rust catalogue crate (\texttt{crates/mtk-catalogue}) --- one Rust implementation underneath, reached from TypeScript in-process through a native binding, so the two surfaces cannot diverge. Signatures and constructions are extracted from the toolkit sources at every regeneration of this catalogue; each construction is quoted verbatim from the toolkit's conformance corpus, code whose emitted certificate bytes are verified byte-identical against the committed golden vectors (helper values named \texttt{fx\_*} or \texttt{f.*} are that corpus's shared fixtures, defined beside the quoted lines). This API constructs and emits; it never checks --- the checker re-derives every claim independently.}
\apilang{TypeScript}{system.pde}
\apiprov{signature --- \texttt{ts/src/system.ts}}
\begin{tssrc}
/** S-PDE: differential-polynomial mechanism over jet variables,
 * deterministic (the stochastic arm is `spde`, which adds the
 * noise-spectrum family to exactly this signature). */
export const pde = profileFactory("pde");

/** The derived PDE kit: the canonical differential-polynomial
 * right-hand side, one exact polynomial per field component at the
 * full index-space arity `state ++ jetVars` (rows targeting one
 * component sum) — the producer face of the declared-data compile
 * edge, mechanized in the toolkit's Lean tree. Derived, never
 * trusted — the checker re-derives it. */
export interface PdeGenKit {
  readonly rhs: ReadonlyArray<PolyWitness>;
}

/** Derive the PDE kit from a mechanism declaration. The stratum
 * fences (single implicit mode, no transcendental terms, drift rows
 * targeting field components only — a jet variable has no
 * independent evolution row) come back as refusals naming their
 * geometry. */
export function pdeGenKit(init: MechanismInit): Result<PdeGenKit, Refusal> {
  const raw = native.sysPdeGenKit(mechanism(init).nativeHandle) as
    | { kit: PdeGenKit }
    | { refusal: unknown };
  return "refusal" in raw
    ? refused(narrowRefusal<Refusal>(raw.refusal))
    : ok(raw.kit);
}
\end{tssrc}
\apiprov{construction --- \texttt{ts/tests/system-native.test.ts}}
\begin{tssrc}
// Viscous Burgers at jet grain — the deterministic sibling of the
// stochastic-heat fixture, at the reference constants: index space
// (u, u_x, u_xx), the whole equation as one differential
// polynomial d/dt u = (1/10) u_xx - u u_x.
return sys.unwrap(
  sys.pde({
    name: "viscous-burgers-jet",
    state: [{ id: 0, name: "u" }],
    jetVars: [
      { id: 1, base: 0, deriv: [1] },
      { id: 2, base: 0, deriv: [2] },
    ],
    drift: [
      {
        var: 0,
        poly: {
          vars: 3,
          terms: [
            { exps: [1, 1, 0], coeff: "-1" },
            { exps: [0, 0, 1], coeff: "1/10" },
          ],
        },
      },
    ],
  }),
);
\end{tssrc}
\apilang{Rust}{mtk\_catalogue::system::presets::Pde::manifest}
\apiprov{signature --- \texttt{crates/mtk-catalogue/src/system/presets.rs}}
\begin{rustsrc}
preset!(
    /// System Catalogue: PDE profile — differential-polynomial
    /// mechanism over jet variables, deterministic. The stochastic
    /// arm is the `Spde` preset's, which adds the noise-spectrum
    /// family to exactly this signature.
    Pde, "PDE",
    req = [JetVars, Drift],
    allowed = [],
    extra = ok
);
\end{rustsrc}
\apiprov{construction --- \texttt{crates/mtk-catalogue/tests/builders.rs}}
\begin{rustsrc}
let pde = MechanismBuilder::new()
    .state_var(0, Some("u"))
    .jet_var(1, 0, vec![2])
    .drift(None, 0, mono(2, 1, 1, r(1, 1)))
    .finish();
Pde::manifest(pde).unwrap();
\end{rustsrc}
\apiprov{The derived kit's Rust face is \texttt{mtk\_catalogue::system::pde\_gen\_kit} --- the canonical differential-polynomial right-hand side per field component at the extended index space, the producer face of the mechanized declared-data compile edge; the quoted TS construction is the committed \texttt{viscous-burgers-jet} golden vector's own presentation.}

\section{Delay Differential Equations}
\label{sec:mt-dde}

\paragraph{The frame.} Delay differential equations $\dot y(t) =
f\bigl(y(t),\, y(t - \tau)\bigr)$ with delay $\tau > 0$ and a
continuous history $\varphi$ on $[-\tau, 0]$, in the classical sense
of Hale--Verduyn Lunel \cite{hale-verduyn-lunel-1993}: solutions are
continuous on $[-\tau, \infty)$, agree with the history, and satisfy
the equation forward in time from $t = 0$.  The textbook state space
is the \emph{history segment} $y|_{[t-\tau,\,t]}$, an element of the
infinite-dimensional space $C([-\tau, 0])$ --- the feature that makes
delay equations look like a new semantic species.

\paragraph{Profile: none --- a wiring pattern.} This frame adds no
formation rule, no new mechanism family, and no new semantic species.
A delay equation \emph{is} a composite of inventory the catalogue
already carries: an ODE cell (section~\ref{sec:mt-ode}) wired to an
\emph{ideal delay line} --- the transport equation $\partial_t u + c\,
\partial_x u = 0$ on a segment of length $c\tau$, the simplest
instance of the PDE frame
(section~\ref{sec:mt-pde})\footnote{And the idealization of a real
machine: the mercury acoustic delay-line memories of the earliest
stored-program computers held their bits as waves in transit down a
tube, read at the far end and re-injected at the near one --- a
store whose state is a field on a spatial segment, as the wiring
pattern reads the history.} ---
through three interface wires:
a \emph{load} wire placing the history on the line at time zero, an
\emph{inflow} trace wire pinning the upstream face to the cell's
state, and an \emph{outflow} tap reading the downstream face into the
cell's vector field.  The wires are directional in exactly the way
the transport operator's dependence structure dictates:
characteristics enter at the inflow and initial faces, which are the
faces the wiring pins; the downstream face is read, never pinned.
The infinite-dimensional history state is then nothing but the field
on the spatial cell, read along characteristics --- a record of a
cell the universe already has, not a new kind of state.

\paragraph{Adequacy.} The reduction is not an analogy; it is a
theorem, mechanized end to end.  Stated self-contained, at the
working scale of an exam problem:

\begin{problem}[the delay line is a composite]
\label{prob:mt-dde}
Fix $c, \tau > 0$, a vector field $f : \mathbb{R}^n
\times \mathbb{R}^n \to \mathbb{R}^n$, and a continuous history
$\varphi : [-\tau, 0] \to \mathbb{R}^n$.  On the cylinder $[0, c\tau]
\times [0, \infty)$, say a triple $(y, z, u)$ --- an
$\mathbb{R}^n$-valued path, a tap signal, and a field --- satisfies
the compiled constraints when: (i) along every characteristic
$s \mapsto (c(s - \zeta),\, s)$ of the cylinder, the field's trace is
continuous with vanishing derivative on its window (the transport
rows); (ii) $u(x, 0) = \varphi(-x/c)$ (the load); (iii) $u(0, t) =
y(t)$ (the inflow); (iv) $z(t) = u(c\tau, t)$ (the tap); and (v)
$y(0) = \varphi(0)$, $y$ is continuous, and $\dot y(t) = f(y(t),
z(t))$ forward in time (the ODE cell).  Show that $(y, z, u)$
satisfies the compiled constraints if and only if there is a
classical solution $Y$ of the delay equation with history $\varphi$
such that, on the composite's domains, $y = Y$, $z(t) = Y(t - \tau)$,
and $u(x, t) = Y(t - x/c)$.  In particular the tap carries the state
delayed by exactly $\tau$, and the spatial slice $x \mapsto u(x, t)$
carries exactly the history segment $Y|_{[t-\tau,\,t]}$.
\end{problem}

\emph{Sketch.} ($\Leftarrow$) Along each characteristic the coordinate
$t - x/c$ is constant, so the reparameterized read of $Y$
is constant on each trace and the transport rows hold; the wire
clauses are the definitional reads of $Y$.  ($\Rightarrow$) The
one-variable mean value theorem pins each trace to its entry value
--- the inflow face for characteristics launched at $\zeta \ge 0$,
the initial face for $\zeta \in [-\tau, 0)$ --- so the field is the
concatenation of history and state read along characteristics
(the method of characteristics \cite{evans-2010}, in its entirety);
the tap then reads that concatenation at $t - \tau$, the ODE row
becomes the delay equation, and the corner compatibility where the
inflow face meets the initial face is precisely the initial condition
$y(0) = \varphi(0)$.  No stepwise construction appears: the
correspondence is uniform in $t$, and the method of steps is an
existence technique the equivalence never consumes.

In the universe's terms: the five constraint families --- two cells
and three wires, the wiring diagram read clause for clause --- form a
loose-arrow diagram of the thin double category of per-element
constraints, each clause read at the possibilistic stratum (every
record the model deems possible satisfies it); the theorem says this
diagram is isomorphic, by an invertible globular cell, to the
support inclusion of the classical solution set.  The inclusion is
set-valued because the square presupposes no uniqueness theory ---
solution sets play the role that credal sets play for the stochastic
frames --- and the per-solution face recovers the pointwise
deterministic inclusion: each classical solution's point mass
satisfies the compiled constraints.  Unlike the stochastic squares,
no constraint forces total mass: the possibilistic clauses are closed
under sub-mass and mixtures, and the demand that the system realize
anything at all is judgment-level content, discharged constructively
at equilibria of the closed loop.  The statement and both directions
are mechanized, sorry-free, in the toolkit's Lean tree
(\texttt{lean/su-functor/}, module
\texttt{SuFunctor.Adequacy.DelayLine}): the record-level equivalence
is \texttt{composite\_iff\_dde}, the invertible cell is
\texttt{delayline\_adequacy}, the delay identity is
\texttt{outflow\_delay}, and the history-segment identification is
\texttt{field\_slice\_is\_history}.  Where the PDE frame's semantic
bridge is staged as named obligations, this instance's left edge is
total: characteristics reduce it to elementary one-variable calculus,
which is what makes the delay line the transport instance whose
adequacy square closes two-sided today.  Dynamics enter as
constraint rows, never as a solution operator: the shift semigroup
the functional-analytic literature takes as primitive is
\emph{derived} content of this left edge.

\paragraph{Status.} The wiring-pattern theorem and the square are
mechanized as above.  Certification inherits the PDE frame's
machinery in its exactly-solvable corner: for rational $c$, $\tau$,
and piecewise-polynomial data with rational breakpoints,
solution-by-characteristics is exact rational arithmetic --- the
field's natural finite presentation is its Bernstein coefficient
vector in the characteristic coordinate, on which transport acts as
the identity (the read head moves; the tape stands still), and
endpoint traces are endpoint coefficients.  Multiple discrete delays
are multiple transport cells in parallel, a corollary of the same
pattern.  Deliberately open: an ingestion profile for the wiring
pattern (the chapter's recurring lesson applies here too);
Lyapunov--Krasovskii functionals read as cell-wise energies with
flux cancellation at the wires, which is designed, not built; and
state-dependent or distributed delays, which leave the ideal-line
instance and are out of scope.

No formation rule accompanies this section: that absence is the
section's content.

\section{Stochastic Differential Equations}
\label{sec:mt-sde}

\paragraph{The frame.} $dX = f(X)\,dt + \sigma(X)\,dW$; certified
questions are reach-avoidance and invariance in probability, via
supermartingale arguments on the generator. The one-dimensional
Ornstein--Uhlenbeck process $dX = -X\,dt + \tfrac{1}{2}\,dW$ is the
frame's harmonic oscillator: the construction quoted in the API
block below, and the first shipped certificates' subject (see
Status).

\paragraph{Profile.} Drift plus diffusion. Rank-one noise populates a
single diffusion column; correlated multi-driver noise populates one
exact-rational $\sigma$ column per independent Brownian driver, with
the covariance $a = \sum_k \sigma^k (\sigma^k)^{\!\top}$
\emph{derived, never stored} --- a representation choice that keeps
the checker's inputs primitive. State-dependent $\sigma(x)$ is
representable (the columns are polynomials). Time: continuous.

\paragraph{Semantics.} The extended-generator rung: certificates are
polynomial inequalities on $\mathcal{A}V = f \cdot \nabla V +
\tfrac12 \sum_{ij} a_{ij} \partial_{ij} V$, and the reading of an
accepted sequent as a probability bound rests on staged Dynkin and
optional-sampling hypotheses --- named, textbook-standard for
polynomial coefficients on compacts, and outside the kernel. The
bound so read is certificate membership in the denotation's defining
infimum, on the per-element route whose continuous-time scope
section~\ref{sec:su-denotation} prices. The Adequacy paragraph below
discharges the staged hypotheses for solutions presented by their
martingale problem; what stays staged is scoped to presentations that
are not.

\paragraph{Adequacy.} The square for this frame closes at the
community's own solution concept --- not the It\^o integral, which
the mechanized library deliberately does not assume, but the
\emph{martingale problem}
\cite{stroock-varadhan-1979, ethier-kurtz-1986}: $X$ solves the
martingale problem for $\mathcal{A}$ when, for every admissible test
observable $V$, the compensated process
$V(X_t) - \int_0^t (\mathcal{A}V)(X_s)\,ds$ is a martingale. Weak
It\^o solutions and martingale-problem solutions coincide; that
equivalence is the community's own theorem, cited here as the licence
and deliberately not re-proved (it needs the stochastic integral,
which is absent from the underlying mathematics library at the pinned
version --- a strong-solution left edge waits on upstream stochastic
integration). For the compiled fragment the identification is
definitional in the literature itself: polynomial diffusions are
\emph{defined} by the martingale problem with polynomial test
functions \cite{filipovic-larsson-2016}. Stated self-contained:

\begin{problem}[the martingale-problem square]
\label{prob:mt-sde}
Fix an SDE profile with polynomial coefficients ---
drift $f$, diffusion columns $\sigma^k$, jump rate $\lambda$ with
reset $g$ --- a closed working domain $D$, and a certificate
polynomial $V$ with exact rational coefficients. Suppose a process
$X$, confined to $D$, solves the stopped martingale problem for the
profile's integro-differential generator $\mathcal{A}$, and the
checker's own compiled row holds: the exact-rational polynomial
$\mathcal{A}V$, computed symbolically from the presentation,
evaluates $\le 0$ everywhere on $D$, with $V \ge 0$ on $D$ and
$V \ge 1$ on a bad set $B$. Show that along every monotone sampling
grid $\tau$,
$\Pr\big(\exists k,\; X_{\tau_k} \in B\big)
  \le \mathbb{E}\big[V(X_{\tau_0})\big]$.
\end{problem}

\emph{Sketch.} Two statics identify the formal with the analytic:
the Fr\'echet derivative of polynomial evaluation is the formal
gradient --- at order one, and, instantiated at each first partial,
at order two --- so the triple $(V, \nabla V, \nabla^2 V)$ of
exact-rational row data enters the martingale problem as a lawful
$C^2$ test datum; and evaluating the symbolically computed generator
equals the semantic generator at that datum. The compensated
certificate process is then a martingale whose compensator has
nonpositive density along the confined path, so $V(X_t)$ is a
continuous-time supermartingale --- Dynkin's formula \emph{derived}
from the presentation, not assumed --- and Ville's inequality
\cite{ville1939} along the sampled grid gives the bound.

The point of the square is what it retires: the staged Dynkin and
optional-sampling hypotheses of the Semantics paragraph are, for
martingale-problem-presented solutions, theorems, and the per-model
residue is being a solution --- the community's own object --- plus
the per-datum integrability provisos, which are the textbook's
standing assumptions for unbounded test functions. The
boundaries, stated as such: verdicts are on the sampled grid (the
continuum-time supremum needs almost-sure path regularity, which
stays with the constructor); well-posedness of the martingale
problem is not claimed, exactly as the ODE square proved no Picard
existence; and unlike this chapter's discrete squares the transport
here is one-directional --- community semantics to certified verdict
--- with no invertible-cell claim at the path-law level. The
statement and the discharge are mechanized, sorry-free, in the
toolkit's Lean tree (\texttt{lean/qtsl-aux/}, modules
\texttt{QtslAux.SdeGenerator}, \texttt{QtslAux.SdeGeneratorPoly},
\texttt{QtslAux.SdeMartingaleProblem}): the solution concept is
\texttt{MpSolution} (stopped/confined, generator-as-pairs form), the
derived continuous-time supermartingale is
\texttt{supermartingale\_of\_gen\_nonpos}, the discharge of the
Dynkin class field is \texttt{MpSolution.toSdeDynkinVille}, the
compilation splice is \texttt{genPoly\_eval\_eq\_generator}, the
order-1/2 statics are \texttt{isC2Fn\_poly}, and the compiled
headline is \texttt{mp\_avoid\_row\_of\_genPoly} with its
deterministic-start form \texttt{mp\_avoid\_init\_of\_genPoly},
exercised end-to-end by a zero-dynamics control
(\texttt{mp\_control}).

\paragraph{Status.} Certification is in place at three scales: the
one-dimensional Ornstein--Uhlenbeck certificates (the Hessian term
present in the checked sequents); the same second-order term running
inside jump-diffusion and hybrid composites
(sections~\ref{sec:mt-jumpdiff} and~\ref{sec:mt-gshs}); and an
encounter-scale correlated-noise family with the cross-term Hessian
certified over two- and three-dimensional simplicial covers. The
correlated family carries the frame's signature results: a refusal
witness for the decorrelated misreading of the same system, with the
cross-term's contribution in closed form, and a certified anisotropic
bound pair --- $0.242$ against $0.828$ at equal Euclidean
distance --- quantifying exactly what correlation is worth.
Representation and production: the noise declarations are shipped
theory rows, the correlated systems exist as instance rows, and the
toolkit's producer surface constructs presentations in both languages
and derives the generator polynomial symbolically from the
presentation --- the exact computation whose identity with the
semantic generator, at the pure-diffusion face, is the Adequacy
paragraph's mechanized splice; derived and never trusted, the checker
re-derives it. The API block below quotes receipted constructions,
and a correlated two-dimensional presentation --- one shared driver
loading both coordinates, so the cross term is live in the derived
covariance --- is a committed cross-language byte-identity vector
whose generator the kit's tests pin in closed form, against its
decorrelated misreading. Open: state-dependent $\sigma$ at $n \ge 2$
(fully representable and derivable, but no shipped certificate
exercises it yet); the imprecise denotation exercised as a
certificate; the ingestion gap.

\paragraph{SPDEs.} No longer an open flank: stochastic PDEs have
their own entry (section~\ref{sec:mt-spde}). The road ran through the
deterministic PDE machinery as promised --- at the moment rungs, a
stochastic certificate is the deterministic row stack plus a single
It\^o-correction row.

\syscatmeta{\begin{mathpar}
\inferrule*[Right=\rulename{S-SDE}]{
  \vec{x} \text{ declared state variables} \\
  D \text{ drift rows} \\
  N \text{ a noise presentation}
}{
  \mathrm{SDE}(\vec{x};\, D;\, N) : \mathsf{System}
}
\end{mathpar}
\[
N \;::=\;
\mathsf{diffusion}\big(\overrightarrow{\langle \mathsf{var},\,
\mathsf{poly} \rangle}\big)
\;\mid\;
\mathsf{noise}\big(\sigma^1, \ldots, \sigma^K\big)
\]
Exactly one noise presentation, by the grammar rather than by a side
condition: $\mathsf{diffusion}$ is the single-column form,
$\mathsf{noise}$ the per-driver sigma-column form (one exact-rational
polynomial column per independent Brownian driver); the covariance
$a = \sum_k \sigma^k (\sigma^k)^{\!\top}$ is derived, never stored.}
\ruleblock{API}
\sourcenote{Producer surface of the \texttt{mtk} toolkit for this model type: the TypeScript arrangement (npm package \texttt{@maitria/mtk}) and the Rust catalogue crate (\texttt{crates/mtk-catalogue}) --- one Rust implementation underneath, reached from TypeScript in-process through a native binding, so the two surfaces cannot diverge. Signatures and constructions are extracted from the toolkit sources at every regeneration of this catalogue; each construction is quoted verbatim from the toolkit's conformance corpus, code whose emitted certificate bytes are verified byte-identical against the committed golden vectors (helper values named \texttt{fx\_*} or \texttt{f.*} are that corpus's shared fixtures, defined beside the quoted lines). This API constructs and emits; it never checks --- the checker re-derives every claim independently.}
\apilang{TypeScript}{system.sde}
\apiprov{signature --- \texttt{ts/src/system.ts}}
\begin{tssrc}
/** S-SDE: drift + diffusion (single-column) XOR noise (multi-driver
 * sigma columns; the covariance is derived, never stored). */
export const sde = profileFactory("sde");

/** The derived SDE generator kit: the exact generator polynomial
 * 𝒜V = ∇V·f + ½ tr(a·Hess V), the covariance a = Σ_k σ^k(σ^k)ᵀ
 * derived from the presented columns, never stored. Derived, never
 * trusted — the checker re-derives it. */
export interface SdeGenKit {
  readonly gen: PolyWitness;
}

/** Derive the SDE generator kit from a mechanism declaration and a
 * certificate observable. Admission is the SDE profile's — jump or
 * killing rows refuse naming the stray family (those presentations
 * belong to `jdGenKit`) — and the computed polynomial's identity
 * with the semantic generator at the pure-diffusion face is
 * mechanized in the toolkit's Lean tree. The shared stratum fences
 * (single implicit mode, no transcendental terms) come back as
 * refusals naming their geometry. */
export function sdeGenKit(
  init: MechanismInit,
  cert: Polynomial,
): Result<SdeGenKit, Refusal> {
  const raw = native.sysSdeGenKit(
    mechanism(init).nativeHandle,
    polyToNative(cert),
  ) as { kit: SdeGenKit } | { refusal: unknown };
  return "refusal" in raw
    ? refused(narrowRefusal<Refusal>(raw.refusal))
    : ok(raw.kit);
}
\end{tssrc}
\apiprov{construction --- \texttt{ts/tests/system-native.test.ts}}
\begin{tssrc}
return sys.unwrap(
  sys.sde({
    name: "ou-1d",
    state: [{ id: 0, name: "x" }],
    drift: [{ var: 0, poly: mono(1, 0, 1, "-1") }],
    diffusion: [{ var: 0, poly: constant(1, "1/2") }],
    init: { kind: "point", coords: [BIG] },
  }),
);
\end{tssrc}
\apilang{Rust}{mtk\_catalogue::system::presets::Sde::manifest}
\apiprov{signature --- \texttt{crates/mtk-catalogue/src/system/presets.rs}}
\begin{rustsrc}
preset!(
    /// System Catalogue: SDE profile — drift + diffusion
    /// (single-column) or noise (per-driver sigma columns; the
    /// covariance is derived, never stored).
    Sde, "SDE",
    req = [Drift],
    allowed = [Diffusion, Noise],
    extra = |m| diffusion_xor_noise("SDE", m)
);
\end{rustsrc}
\apiprov{construction --- \texttt{crates/mtk-catalogue/tests/common/mod.rs}}
\begin{rustsrc}
let big = WitRat::new(
    "1234567890123456789012345678901234567890"
        .parse::<BigInt>()
        .unwrap(),
    BigInt::from(7u8),
);
let m = MechanismBuilder::new()
    .state_var(0, Some("x"))
    .drift(None, 0, mono(1, 0, 1, r(-1, 1)))
    .diffusion(None, 0, constant(1, r(1, 2)))
    .init(InitRegion::Point { coords: vec![big] })
    .finish();
presets::Sde::manifest(m)
    .expect("fixture conforms to the SDE profile")
    .named("ou-1d")
\end{rustsrc}
\apiprov{The multi-driver presentation (per-driver sigma columns) is the corpus's committed \texttt{correlated-2d} golden vector (fixture \texttt{sde\_noise}, whose derived cross term the generator-kit tests pin in closed form); the generator kit's Rust face is \texttt{mtk\_catalogue::system::sde\_gen\_kit}.}

\section{Stochastic Partial Differential Equations}
\label{sec:mt-spde}

\paragraph{The frame.} Stochastic evolution equations $du =
F(u, \partial_x u, \partial_x^2 u, \ldots)\,dt + dW_Q$, with $W_Q$ a
$Q$-Wiener process over the declared spatial domain; the certified
questions are expected-energy dissipation to a noise floor, absorbing
balls, and probability bounds on energy excursions.

\paragraph{Profile.} The deterministic jet mechanism plus a declared
noise spectrum: finitely many explicit driver shapes ---
exact-rational combinations of the domain's orthonormal Fourier
basis, one amplitude column per driver --- and a required tail form
in one of three shapes (zero, a rational trace bracket, or
declared-divergent), so the covariance's trace is derived, never
stored, and a spectrum without a declared tail is refused rather than
read as zero. Amplitude columns are constants at this preview
(additive noise), but their declared type is a polynomial over the
jet space, so multiplicative noise is a fence-lift rather than a
schema change. Time: continuous.

\paragraph{Semantics.} The stochastic solution-concept rungs, weakest
first: the expected-energy class (the inhomogeneous inequality with
the trace of $Q$ as its forcing constant), the
compensated-supermartingale form (horizon-tagged), and the mode-law
rung where explicit Ornstein--Uhlenbeck second-moment trajectories
live. Every stochastic-analysis fact is a named obligation behind a
class field with stated discharge routes --- no theorem below the
docstring mentions a stochastic integral --- extending the staging
discipline section~\ref{sec:mt-pde} established.

\paragraph{Status.} Certification: five exact-rational artefacts ---
stochastic Burgers and stochastic heat moment certificates (the
deterministic rows carried verbatim, the It\^o correction docked as
one constant row; the heat floor $\mathrm{tr}(Q)/2\nu = 13/10$ exact,
with a deliberately sharp width-zero instance) and two stochastic
Kuramoto--Sivashinsky absorbing balls that embed and replay their
deterministic siblings byte-identically before any stochastic row is
added. All five carry probability rows: a horizon-tagged supremum
bound of Kushner's classical shape \cite{kushner1967}, Chebyshev on
certified variances,
and an exponential-compensation row strictly sharper at both
registered points. Mutant batteries assert rejection messages, and
witnesses realizing the noise are checker-required, so a green moment
bound cannot be vacuously noise-silent. The jet-level declaration now
has its proof-layer face: spectrum declarations with the required
tail carry mechanized derived-covariance data (module
\texttt{SuFunctor.Adequacy.PdeRows}: the explicit trace by Parseval
with per-element grouping, the q-table diagonal with
trace-equals-summed-diagonal, and the certified heat spectrum's
$\mathrm{tr}(Q) = 13/50$, $q_1 = 1/5$, $q_2 = 3/50$ recomputed by
theorem), and the plain-trace consumer is total-or-refusing --- a
declared-divergent tail refuses structurally
(\texttt{plainTrace\_divergent\_none}), never defaulting to zero.
Representation and production: the toolkit's producer surface now
carries the derived-kit sibling of that proof layer, shared with
section~\ref{sec:mt-pde} --- from a declared presentation the kit
derives the canonical differential-polynomial right-hand side (the
compile edge's producer face) and, from the spectrum, the q-table
diagonal, the trace as the summed diagonal, the plain-trace
enclosure (structurally refused on a divergent tail, in both
languages), and the dissipation floor $\mathrm{tr}(Q)/2\nu$ exactly
where the presentation exhibits it --- transport terms and mean-mode
mass forfeit the floor rather than have it guessed; derived and
never trusted, the checker re-derives every row. The kit's tests
recompute the certified constants from the committed golden vector's
own presentation, and the deterministic sibling (viscous Burgers at
jet grain) is now itself a committed cross-language byte-identity
vector. Deliberately open: trace-class tails (declaration vocabulary
without a certified instance); the cylindrical mode-resolved
variant; multiplicative noise; the Fourier mode-projection semantics
(the compile-for-semantics leg, a named obligation at the Lean
stratum); composition.

\syscatmeta{\begin{mathpar}
\inferrule*[Right=\rulename{S-SPDE}]{
  \vec{x} \text{ declared state variables} \\
  J : \overrightarrow{\langle \mathsf{base},\, \alpha \rangle}
  \text{ jet variables} \\
  D \text{ drift rows over } (\vec{x}, J) \\
  N = \big(\overrightarrow{\langle \mathsf{shape},\,
  \mathsf{amp} \rangle};\; \tau\big) \text{ a noise spectrum}
}{
  \mathrm{SPDE}(\vec{x};\, J;\, D;\, N) : \mathsf{System}
}
\end{mathpar}
\[
\tau \;::=\; \mathsf{zero}
\;\mid\; \mathsf{traceBracket}(T_{\mathrm{lo}}, T_{\mathrm{hi}})
\;\mid\; \mathsf{divergent}
\]
The jet families are section~\ref{sec:mt-pde}'s, verbatim; the noise
spectrum adds exactly one family. Each driver is a spatial shape (an
exact-rational combination of the domain's orthonormal Fourier basis)
with an amplitude column; the tail form $\tau$ is required by the
grammar rather than by a side condition, so a spectrum without a
declared tail is unrepresentable here and refused at preset admission
in the toolkit. The covariance
$Q = \sum_j \mathrm{shape}_j\,\mathrm{shape}_j^{\!*}$ (plus the
declared tail) and its trace are derived, never stored; a consumer
needing the plain trace refuses a divergent tail.}
\ruleblock{API}
\sourcenote{Producer surface of the \texttt{mtk} toolkit for this model type: the TypeScript arrangement (npm package \texttt{@maitria/mtk}) and the Rust catalogue crate (\texttt{crates/mtk-catalogue}) --- one Rust implementation underneath, reached from TypeScript in-process through a native binding, so the two surfaces cannot diverge. Signatures and constructions are extracted from the toolkit sources at every regeneration of this catalogue; each construction is quoted verbatim from the toolkit's conformance corpus, code whose emitted certificate bytes are verified byte-identical against the committed golden vectors (helper values named \texttt{fx\_*} or \texttt{f.*} are that corpus's shared fixtures, defined beside the quoted lines). This API constructs and emits; it never checks --- the checker re-derives every claim independently.}
\apilang{TypeScript}{system.spde}
\apiprov{signature --- \texttt{ts/src/system.ts}}
\begin{tssrc}
/** S-SPDE: the PDE families plus a noise spectrum — Fourier-shape
 * driver rows (`noiseDrivers`) and a REQUIRED tail form
 * (`noiseTail`: zero XOR traceBracket XOR divergent). A spectrum
 * whose tail was never declared refuses here — never read as zero;
 * the covariance and its trace are derived, never stored. */
export const spde = profileFactory("spde");

/** One entry of the derived covariance diagonal: the q-table value
 * at an orthonormal Fourier-basis element — amplitudes and grouped
 * coefficients squared and summed over the drivers. Zero entries
 * (cancelled groups; `sin` at mode 0, the zero function) are
 * omitted, never stored. */
export interface QTableEntry {
  readonly kind: "sin" | "cos";
  readonly mode: number;
  readonly q: string;
}

/** The derived SPDE kit: the PDE right-hand side plus covariance
 * data derived from the noise-spectrum declaration — the q-table
 * diagonal, the explicit trace (their sum, by the mechanized
 * finite-Fubini identity), the plain-trace enclosure under the
 * declared tail (`null` on a declared-divergent tail — the
 * structural refusal, never a zero default), and the dissipation
 * floor tr(Q)/2ν exactly where the presentation exhibits it
 * (`null` otherwise — transport terms and mean-mode mass are
 * certificate business, not producer recognition). Derived, never
 * stored, never trusted. */
export interface SpdeGenKit {
  readonly rhs: ReadonlyArray<PolyWitness>;
  readonly qTable: ReadonlyArray<QTableEntry>;
  readonly traceExplicit: string;
  readonly plainTrace: readonly [string, string] | null;
  readonly heatFloor: readonly [string, string] | null;
}

/** Derive the SPDE kit from a mechanism declaration. Admission is
 * the SPDE profile's — a spectrum whose tail was never declared
 * refuses there, never read as zero — and the constant-amplitude
 * stratum's multiplicative fence comes back as a refusal naming its
 * driver. */
export function spdeGenKit(init: MechanismInit): Result<SpdeGenKit, Refusal> {
  const raw = native.sysSpdeGenKit(mechanism(init).nativeHandle) as
    | { kit: SpdeGenKit }
    | { refusal: unknown };
  return "refusal" in raw
    ? refused(narrowRefusal<Refusal>(raw.refusal))
    : ok(raw.kit);
}
\end{tssrc}
\apiprov{construction --- \texttt{ts/tests/system-native.test.ts}}
\begin{tssrc}
// Stochastic heat on the torus, the reference constants: nu =
// 1/10; q_1 = 1/5 and q_2 = 3/50 realized as sum-of-rational-
// squares driver decompositions (sigma_k = sqrt(q_k) is
// irrational); derived trace 13/50, floor tr(Q)/2nu = 13/10 exact.
// The tail form is REQUIRED — zero is said, never implied.
const amp = constant(2, "1");
return sys.unwrap(
  sys.spde({
    name: "stochastic-heat-torus",
    state: [{ id: 0, name: "u" }],
    jetVars: [{ id: 1, base: 0, deriv: [2] }],
    drift: [{ var: 0, poly: mono(2, 1, 1, "1/10") }],
    noiseDrivers: [
      { shape: [{ kind: "sin", mode: 1, coeff: "1/5" }], amp },
      { shape: [{ kind: "sin", mode: 1, coeff: "2/5" }], amp },
      { shape: [{ kind: "sin", mode: 2, coeff: "1/10" }], amp },
      { shape: [{ kind: "sin", mode: 2, coeff: "1/10" }], amp },
      { shape: [{ kind: "sin", mode: 2, coeff: "1/5" }], amp },
    ],
    noiseTail: { kind: "zero" },
  }),
);
\end{tssrc}
\apilang{Rust}{mtk\_catalogue::system::presets::Spde::manifest}
\apiprov{signature --- \texttt{crates/mtk-catalogue/src/system/presets.rs}}
\begin{rustsrc}
preset!(
    /// System Catalogue: SPDE profile — the PDE profile's families
    /// (jet variables + differential-polynomial drift) plus a noise
    /// spectrum: finitely many explicit Fourier-shape driver rows
    /// and a REQUIRED tail form in exactly one of three shapes
    /// (zero / trace-bracket / declared-divergent). A spectrum whose
    /// tail was never declared is refused here — never read as zero
    /// — and the covariance (with its trace) is derived, never
    /// stored.
    Spde, "SPDE",
    req = [JetVars, Drift, NoiseSpectrum],
    allowed = [],
    extra = |m| {
        let spectrum = m.noise_spectrum.as_ref().expect("required by signature");
        if spectrum.tail.is_none() {
            return Err(Refusal::MissingTailForm { profile: "SPDE" });
        }
        Ok(())
    }
);
\end{rustsrc}
\apiprov{construction --- \texttt{crates/mtk-catalogue/tests/builders.rs}}
\begin{rustsrc}
let spde = MechanismBuilder::new()
    .state_var(0, Some("u"))
    .jet_var(1, 0, vec![2])
    .drift(None, 0, mono(2, 1, 1, r(1, 10)))
    .noise_driver(
        vec![FourierTerm {
            kind: FourierKind::Sin,
            mode: 1,
            coeff: r(1, 5),
        }],
        constant(2, r(1, 1)),
    )
    .noise_driver(
        vec![FourierTerm {
            kind: FourierKind::Sin,
            mode: 1,
            coeff: r(2, 5),
        }],
        constant(2, r(1, 1)),
    )
    .noise_tail(TailForm::Zero)
    .finish();
Spde::manifest(spde).unwrap();
\end{rustsrc}
\apiprov{The tail form is the spectrum's required member: the corpus's \texttt{spde\_missing\_tail\_refused\_never\_read\_as\_zero} sibling exercises the admission refusal, and the divergent-tail plain-trace refusal is witnessed in \texttt{plain\_trace\_derived\_never\_stored}. The derived kit's Rust face is \texttt{mtk\_catalogue::system::spde\_gen\_kit}: the q-table diagonal, the explicit trace as the summed diagonal, the plain-trace enclosure (structurally refused on a divergent tail), and the dissipation floor $\mathrm{tr}(Q)/2\nu$ exactly where the presentation exhibits it --- the kit's tests recompute the reference constants ($q_1 = 1/5$, $q_2 = 3/50$, trace $13/50$, floor $13/10$) from the committed golden vector's own presentation.}

\section{Random Differential Equations}
\label{sec:mt-rode}

\paragraph{The frame.} An ODE or PDE whose parameters $\theta$ are
drawn from a declared uncertainty set; the certified question is a
property \emph{uniform} over the set.

\paragraph{Profile.} Drift plus a boxed parameter axis: the
parameter is one more coordinate of the region, constant along
trajectories. Disturbance semantics: uniform over the box --- the
$\forall\theta$ leg is carried entirely by the region, not by any new
machinery. Time: continuous.

\paragraph{Semantics.} Per-$\theta$ trajectory semantics --- each
theorem a one-line specialization of the ODE semantics seam, and that
one-line-ness is the measured content of the claim that the kernel's
cone treats parameter boxes natively. Uniformity over the box is, in
the universe's terms, a credal body --- the possibilistic face
over the box-indexed trajectory family (Remark~\ref{rem:su-faces})
--- and the $\forall\theta$ verdict is its lower prevision
(Definition~\ref{def:su-previsions}).

\paragraph{Status.} Certification: a damped oscillator
$\dot x = v$, $\dot v = -x - \theta v$ with damping boxed
$\theta \in [1/4, 1]$ carries a barrier decrease
$\dot V \le -\tfrac18 (x^2 + v^2)$ on an annulus, uniformly over the
box, as one polynomial row family --- checked, round-tripped through
the wire format, with refusal probes on both producer and checker
sides, and a mechanized per-$\theta$ semantics layer. This frame was
priced as nearly free before it was built, and building it confirmed
the price: the thesis's own $\forall\theta$ property form is
RODE-shaped. Production: the toolkit's producer surface constructs
the boxed-parameter presentation in both languages, and the damped
oscillator itself is a committed cross-language byte-identity
vector, the parameter family's first. The kit derives the
state~$\times$~parameter generator row symbolically at the product
index space, with derivatives over state axes only ($\dot\theta =
0$: a $\theta$-only observable is conserved, by test). Its tests pin
$\mathcal{A}(x^2 + v^2) = -2\theta v^2$ in closed form in both
languages --- one polynomial over the product space is the whole
$\forall\theta$ row family --- and a collapse differential checks
that a $\theta$-free drift derives exactly the ODE kit's polynomial.
Open: the mechanism as ingested rows (the ingestion gap); a fixture
nonlinear in $\theta$; and the RPDE side, which waits on
section~\ref{sec:mt-pde}'s bridges.

\paragraph{Adequacy.} The square for this frame is deliberately
slight, and the slightness is itself the certified content: the
community's solution concept for random ODEs is \emph{pathwise} ---
for each realization $\theta$, a classical solution of the
$\theta$-instance, analyzed with deterministic calculus
\cite{han-kloeden-2017} --- so the per-draw connector, square, and
uniqueness theorems are the ODE square of
section~\ref{sec:mt-ode} at the drawn drift $f_\theta$, consumed
verbatim rather than re-proved. What remains is the box. Stated
self-contained:

\begin{problem}[the boxed-parameter square]
\label{prob:mt-rode}
Fix a drift family $f : \Theta \times \mathbb{R}^n
\to \mathbb{R}^n$, continuous in the state at each draw in a
parameter box $P \subseteq \Theta$, and $x_0 \in \mathbb{R}^n$.
Records are paths $y : \mathbb{R} \to \mathbb{R}^n$; the compiled
constraints are the initial row and the \emph{boxed cell} --- some
draw $\theta \in P$ drives $y$'s forward rows. Show the compiled
records are exactly $\bigcup_{\theta \in P} S_\theta$, where
$S_\theta$ is the classical solution set of the $\theta$-instance;
show the same for the rectangle with both boxes open --- the initial
row replaced by an initial box $B \subseteq \mathbb{R}^n$ ---
denoting $\bigcup_{\theta \in P} \bigcup_{x \in B} S_{\theta,x}$,
with $S_{\theta,x}$ the $\theta$-instance's solution set from
$x$; and show
the polarity of the box matters: over the empty box the compiled
body is empty, while a uniform-verdict reading --- intersection over
draws --- would admit every record.
\end{problem}

\emph{Sketch.} Intersecting with the initial row distributes into
the parameter union, and each summand is the ODE square at
$f_\theta$; the rectangle composes the parameter union with the ODE
square's initial-box identity, one draw at a time; the empty union
is empty, the empty intersection is everything.

The polarity is easy to get backwards. The compiled \emph{body} over the box is a union
--- existential in $\theta$, the same demonic convention as the
initial box --- while the certified uniform-over-$\theta$ property of
the Semantics paragraph above is the opposite polarity: a verdict
over the whole box-indexed family, the lower prevision of
Remark~\ref{rem:su-faces}, living at the judgment, never as an
intersection body. Uniqueness enters per draw with a
certificate-shaped hypothesis --- a single Lipschitz constant uniform
over the box pins every draw in it, the semantics shadow of a
certificate stated as one row family over state and parameter. The
statement and both directions are mechanized, sorry-free, in the
toolkit's Lean tree (\texttt{lean/su-functor/}, module
\texttt{SuFunctor.Adequacy.Rode}): the per-draw connector and square
are \texttt{isRodeSol\_iff\_integralCurve} and
\texttt{rode\_draw\_adequacy} (one-line specializations, stated as
such), the union face is \texttt{rode\_adequacy} with the rectangle
\texttt{rode\_box\_adequacy}, the uniform-constant pinning is
\texttt{rodeSols\_eqOn\_of\_lipschitzOnWith}, and the wiring into the
semantic universe's judgment (\texttt{su\_rode\_adequacy}) runs
through the same transport layer as the sibling squares, with the zero-law
contrast holding at every box including the empty one --- the
infeasible frame denotes a bottom body, never an absence of models.
The completeness hypothesis enters exactly where the ODE connector
is consumed, nowhere else.

\syscatmeta{\begin{mathpar}
\inferrule*[Right=\rulename{S-RODE}]{
  \vec{x} \text{ declared state variables} \\
  \Theta : \overrightarrow{\mathsf{ParamDecl}} \\
  \exists\, p \in \Theta.\;\; p.\mathsf{bounds} \neq \bot \\
  D \text{ drift rows, each } \mathsf{poly} \in
  \mathbb{Q}[\vec{x}, \Theta]
}{
  \mathrm{RODE}(\vec{x};\, \Theta;\, D) : \mathsf{System}
}
\end{mathpar}
A boxed parameter ($\mathsf{bounds} = [\mathsf{lo}, \mathsf{hi}]$,
exact) is the $\forall\theta$ axis; an unboxed parameter is a plain
named constant, which is why the existential premise demands at
least one box.}
\ruleblock{API}
\sourcenote{Producer surface of the \texttt{mtk} toolkit for this model type: the TypeScript arrangement (npm package \texttt{@maitria/mtk}) and the Rust catalogue crate (\texttt{crates/mtk-catalogue}) --- one Rust implementation underneath, reached from TypeScript in-process through a native binding, so the two surfaces cannot diverge. Signatures and constructions are extracted from the toolkit sources at every regeneration of this catalogue; each construction is quoted verbatim from the toolkit's conformance corpus, code whose emitted certificate bytes are verified byte-identical against the committed golden vectors (helper values named \texttt{fx\_*} or \texttt{f.*} are that corpus's shared fixtures, defined beside the quoted lines). This API constructs and emits; it never checks --- the checker re-derives every claim independently.}
\apilang{TypeScript}{system.rode}
\apiprov{signature --- \texttt{ts/src/system.ts}}
\begin{tssrc}
/** S-RODE: drift + a boxed parameter axis (disturbance: parameter
 * box, uniform-over-θ). */
export const rode = profileFactory("rode");

/** The derived RODE generator kit: the exact Lie-derivative
 * polynomial 𝒜V = Σ_i ∂_iV·f_i over the product index space
 * `state ++ params` (parameters in declaration order after the state
 * roster), state axes only — parameters are constant along
 * realizations, so no parameter-derivative term exists, and one
 * polynomial in state × parameter is the whole ∀θ row family.
 * Derived, never trusted — the checker re-derives it. */
export interface RodeGenKit {
  readonly gen: PolyWitness;
}

/** Derive the RODE generator kit from a mechanism declaration and a
 * certificate observable (read at the product arity — a certificate
 * may itself name the parameter axes, the family-certificate face).
 * Admission is the RODE profile's — an unboxed-only roster refuses
 * as `noBoxedParam`, and a drift row addressed past the state roster
 * refuses naming the row (a parameter has no evolution row). */
export function rodeGenKit(
  init: MechanismInit,
  cert: Polynomial,
): Result<RodeGenKit, Refusal> {
  const raw = native.sysRodeGenKit(
    mechanism(init).nativeHandle,
    polyToNative(cert),
  ) as { kit: RodeGenKit } | { refusal: unknown };
  return "refusal" in raw
    ? refused(narrowRefusal<Refusal>(raw.refusal))
    : ok(raw.kit);
}
\end{tssrc}
\apiprov{construction --- \texttt{ts/tests/system-native.test.ts}}
\begin{tssrc}
return sys.unwrap(
  sys.rode({
    name: "damped-oscillator",
    state: [
      { id: 0, name: "x" },
      { id: 1, name: "v" },
    ],
    params: [{ id: 0, name: "theta", bounds: ["1/4", "1"] }],
    drift: [
      { var: 0, poly: mono(3, 1, 1, "1") },
      {
        var: 1,
        poly: {
          vars: 3,
          terms: [
            { exps: [1, 0, 0], coeff: "-1" },
            { exps: [0, 1, 1], coeff: "-1" },
          ],
        },
      },
    ],
  }),
);
\end{tssrc}
\apilang{Rust}{mtk\_catalogue::system::presets::Rode::manifest}
\apiprov{signature --- \texttt{crates/mtk-catalogue/src/system/presets.rs}}
\begin{rustsrc}
preset!(
    /// System Catalogue: RODE profile — drift + a boxed parameter
    /// axis (disturbance: parameter box, uniform-over-θ).
    Rode, "RODE",
    req = [Drift, Params],
    allowed = [],
    extra = |m| {
        if m.params.iter().any(|p| p.bounds.is_some()) {
            Ok(())
        } else {
            Err(Refusal::NoBoxedParam { profile: "RODE" })
        }
    }
);
\end{rustsrc}
\apiprov{construction --- \texttt{crates/mtk-catalogue/tests/common/mod.rs}}
\begin{rustsrc}
let vdot = Poly {
    vars: 3,
    terms: vec![
        PolyTerm {
            exps: vec![1, 0, 0],
            coeff: r(-1, 1),
        },
        PolyTerm {
            exps: vec![0, 1, 1],
            coeff: r(-1, 1),
        },
    ],
};
let m = MechanismBuilder::new()
    .state_var(0, Some("x"))
    .state_var(1, Some("v"))
    .param(0, Some("theta"), Some([r(1, 4), r(1, 1)]))
    .drift(None, 0, mono(3, 1, 1, r(1, 1)))
    .drift(None, 1, vdot)
    .finish();
presets::Rode::manifest(m)
    .expect("fixture conforms to the RODE profile")
    .named("damped-oscillator")
\end{rustsrc}
\apiprov{The construction is the committed \texttt{damped-oscillator} golden vector (fixture \texttt{rode} — the parameter family's first golden bytes: drift over the product index space \(\mathit{state} +\!\!+ \mathit{params}\), damping boxed \(\theta \in [1/4, 1]\)), and it doubles as the generator kit's closed-form control (\(\mathcal{A}(x^2+v^2) = -2\theta v^2\) — one polynomial in state \(\times\) parameter is the whole \(\forall\theta\) row family); the kit's Rust face is \texttt{mtk\_catalogue::system::rode\_gen\_kit}.}

\section{Jump-Diffusion}
\label{sec:mt-jumpdiff}

\paragraph{The frame.} Continuous drift and diffusion punctuated by
Poisson-rate jumps; certified questions are ruin and reach-avoidance
bounds through the extended generator with its jump term.

\paragraph{Profile.} All three field families at once --- drift,
diffusion, and jump legs --- on one mechanism, plus an optional
killing family. Each jump row is one leg: a polynomial rate paired
with a reset kernel, either a single polynomial reset (the Dirac
form) or a finite mixture of weighted polynomial resets, with an
optional reset-image box. Mixture weights may total \emph{below} one:
a sub-stochastic kernel is lawful by design, the mass deficit being
killing (the standard convention for jump kernels of total mass below
one \cite{davis-1993}); a row presenting both kernel forms at once is
refused, naming the row. Disturbance: Brownian drivers plus
finite-activity jumps. Time: continuous, with jumps as events in real
time; per-execution jump counters are bookkeeping, never a
composition base.

\paragraph{Semantics.} The extended generator
$\mathcal{A}V = f V' + \tfrac{\sigma^2}{2} V'' +
\lambda\,(g \fatsemi V - V)$, composed from the same shipped summands the
pure-diffusion and pure-jump frames use, with no new checker
machinery --- each degenerate case asserts its reduction to the
shipped frame. Finite-activity hypotheses are named; the bridge set is
the diffusion frame's plus an It\^o--L\'evy step.

\paragraph{Adequacy (generator face).} The Semantics paragraph's
composition claim --- the generator assembled from the shipped
summands, each degenerate case asserting its reduction --- is
mechanized at the frame's full profile generality, and the identity
holds with \emph{no side conditions}. Stated self-contained:

\begin{problem}[the generator splice, at full profile generality]
\label{prob:mt-jumpdiff}
Fix a polynomial jump-diffusion presentation: drift
$f$, diffusion columns $\sigma^k$, a killing rate $c$, and finitely
many jump legs, each a polynomial rate $\lambda_\ell$ paired with a
finite sub-stochastic mixture of polynomial resets
$(w_{\ell i}, g_{\ell i})$. Write $T_{\mu_\ell} V = \sum_i
w_{\ell i}\,(g_{\ell i} \fatsemi V)$ for the transfer action of the
$\ell$-th mixture, and let $\mathcal{A}V$ be the exact-rational
polynomial computed symbolically from the presentation: killing
multiplies, drift pairs the formal gradient, diffusion traces the
formal Hessian against the derived covariance, each jump leg
contributes $\lambda_\ell\,(T_{\mu_\ell}V - V)$. Show that evaluating
$\mathcal{A}V$ anywhere on $\mathbb{R}^n$ equals the semantic
integro-differential generator
\[
  -c\,V + \nabla V \cdot f
  + \tfrac12 \textstyle\sum_{ij} a_{ij}\,\partial_{ij}V
  + \textstyle\sum_\ell \lambda_\ell\,\big(T_{\mu_\ell}V - V\big)
\]
of the evaluated profile, with the observable's derivative data
instantiated at the formal derivatives --- for \emph{every}
presentation, with no conservativity or leg-shape hypothesis --- and
that at one Dirac leg with zero killing the semantic generator is the
SDE frame's operator (section~\ref{sec:mt-sde}) on the nose.
\end{problem}

\emph{Sketch.} Evaluation is an algebra homomorphism, so each field
family passes through separately; the transfer of a finite mixture is
a finite sum, exact at the polynomial stratum, and a list induction
carries the leg sum. A sub-stochastic mixture adds no new semantics:
its jump increment decomposes as Dirac-leg increments minus the mass
deficit spent as killing --- the standard convention for jump kernels
of total mass below one \cite{davis-1993}. The one-Dirac-leg tie is a
computation collapsing the leg sum.

The identity of operators carries the verdict route with it: the
martingale problem consumes its generator opaquely --- the compensator
argument never inspects the drift, diffusion, killing, or jump
structure --- so the SDE section's solution concept, its derived
continuous-time supermartingale, and its verdict headlines transplant
verbatim to the multi-leg generator. Because this section's splice
holds with no side conditions, the transplanted headlines hold for
\emph{every} presentation --- killing, finitely many legs,
sub-stochastic mixtures --- with no conservativity or leg-shape
hypothesis, where the SDE section's headlines carried both. The
single-Dirac-leg square then \emph{factors} through the multi-leg one:
profiles embed, the operators agree on the nose, and solutions
transport along the identity, verdicts included. The boundaries
that remain are the SDE square's own: verdicts on the sampled grid,
well-posedness not claimed, transport one-directional. All of it is
mechanized, sorry-free, in the toolkit's Lean tree
(\texttt{lean/qtsl-aux/}, modules \texttt{QtslAux.JumpDiffusion} and
\texttt{QtslAux.JumpDiffusionMartingaleProblem}): the profile is
\texttt{JdProfile}, the hypothesis-free splice is
\texttt{genPoly\_eval\_eq\_generator}, the degeneracy ties are
\texttt{generator\_one\_dirac} and \texttt{toJd\_generator\_one\_dirac},
the mixture-to-killing decomposition is \texttt{transferSum\_sub\_self},
and a two-leg control with every field family live pins the closed form
(\texttt{two\_leg\_control}); on the verdict side, the solution concept
is \texttt{MpSolution} at the multi-leg generator, the derived
supermartingale is \texttt{supermartingale\_of\_gen\_nonpos}, the
sampled Dynkin class is \texttt{JdDynkinVille}, the headlines are
\texttt{mp\_avoid\_row\_of\_genPoly} and
\texttt{mp\_avoid\_init\_of\_genPoly}, the factoring is
\texttt{generator\_toJd} with \texttt{MpSolution.toJd}, and the
end-to-end control's presentation carries live rates on two legs with a
genuine mixture (\texttt{mp\_control}).

\paragraph{Status.} Certification: two hand-derived systems check
through the unmodified stack --- a Brownian-perturbed risk process of
the classical ruin-theory shape (drift helps, diffusion and claim
jumps hurt; certified ruin bound $49/121$ from its stated start), and
a load-shedding system in which the state-dependent jump term
\emph{alone} carries the certificate against both unstable drift and
diffusion. At $n = 2$ the encounter composite carries all three
families with full-rank correlated noise. Representation and
production: the toolkit's producer surface constructs presentations
at the frame's full profile generality --- killing rows and multi-leg
sub-stochastic mixture kernels included --- in both languages, and
derives the generator polynomial together with its per-leg rate rows
symbolically from the presentation (the exact computation whose
identity with the semantic generator is the Adequacy paragraph's
mechanized splice; derived and never trusted --- the checker
re-derives both). The API block below quotes receipted constructions,
and the reference presentation exercising every field family at once
--- a Dirac leg beside a genuinely sub-stochastic mixture leg, with
killing --- is a committed cross-language byte-identity vector. Open:
certification exercise of multi-leg jump kernels and state-dependent
$\sigma$ (fully representable and derivable as above, but no shipped
certificate exercises them yet); time-bounded and discounted
variants; closure.

\syscatmeta{\begin{mathpar}
\inferrule*[Right=\rulename{S-JUMPDIFF}]{
  \vec{x} \text{ declared state variables} \\
  D \text{ drift rows} \\
  N \text{ a noise presentation (section~\ref{sec:mt-sde})} \\
  J : \overrightarrow{\langle \mathsf{mode}?,\, \mathsf{rate},\,
  \mathsf{reset} \mathbin{|} \mathsf{mixture},\,
  \mathsf{reset\_image}? \rangle} \text{ jump rows} \\
  K \text{ killing rows (optional)}
}{
  \mathrm{JumpDiffusion}(\vec{x};\, D;\, N;\, J;\, K) : \mathsf{System}
}
\end{mathpar}
Jump rows carry rate a polynomial intensity and exactly one kernel
form: reset a polynomial substitution, or mixture a finite weighted
list of polynomial substitutions (sub-stochastic totals lawful ---
the deficit is killing); reset image an exact box. Killing rows carry
a polynomial rate; an absent family is the conservative declaration.}
\ruleblock{API}
\sourcenote{Producer surface of the \texttt{mtk} toolkit for this model type: the TypeScript arrangement (npm package \texttt{@maitria/mtk}) and the Rust catalogue crate (\texttt{crates/mtk-catalogue}) --- one Rust implementation underneath, reached from TypeScript in-process through a native binding, so the two surfaces cannot diverge. Signatures and constructions are extracted from the toolkit sources at every regeneration of this catalogue; each construction is quoted verbatim from the toolkit's conformance corpus, code whose emitted certificate bytes are verified byte-identical against the committed golden vectors (helper values named \texttt{fx\_*} or \texttt{f.*} are that corpus's shared fixtures, defined beside the quoted lines). This API constructs and emits; it never checks --- the checker re-derives every claim independently.}
\apilang{TypeScript}{system.jumpDiffusion}
\apiprov{signature --- \texttt{ts/src/system.ts}}
\begin{tssrc}
/** S-JUMPDIFF: drift + jumps (+ diffusion/noise, + killing). Jump
 * kernels at full generality: each row is a leg whose reset is a
 * Dirac reset or a finite sub-stochastic mixture of resets (the
 * deficit is killing). */
export const jumpDiffusion = profileFactory("jumpDiffusion");

/** The derived jump-diffusion generator kit: the exact generator
 * polynomial 𝒜V = −c·V + ∇V·f + ½ tr(a·Hess V) + Σ λ_ℓ(T_ℓV − V)
 * and, per jump leg in row order, the rate-validity row −λ_ℓ (its
 * nonpositivity on the working region is the leg's model-validity
 * row). Derived, never trusted — the checker re-derives both. */
export interface JdGenKit {
  readonly gen: PolyWitness;
  readonly rateRows: ReadonlyArray<PolyWitness>;
}

/** Derive the jump-diffusion generator kit from a mechanism
 * declaration and a certificate observable. The computed polynomial's
 * identity with the semantic integro-differential generator — for
 * every presentation: killing, multi-leg, sub-stochastic mixtures —
 * is mechanized with no hypotheses in the toolkit's Lean tree. The
 * stratum fences (single implicit mode, no transcendental terms, no
 * jet variables) come back as refusals naming their geometry. */
export function jdGenKit(
  init: MechanismInit,
  cert: Polynomial,
): Result<JdGenKit, Refusal> {
  const raw = native.sysJdGenKit(
    mechanism(init).nativeHandle,
    polyToNative(cert),
  ) as { kit: JdGenKit } | { refusal: unknown };
  return "refusal" in raw
    ? refused(narrowRefusal<Refusal>(raw.refusal))
    : ok(raw.kit);
}
\end{tssrc}
\apiprov{construction --- \texttt{ts/tests/system-native.test.ts}}
\begin{tssrc}
return sys.unwrap(
  sys.jumpDiffusion({
    name: "two-leg-jd",
    state: [{ id: 0, name: "x" }],
    drift: [{ var: 0, poly: mono(1, 0, 1, "1") }],
    diffusion: [{ var: 0, poly: constant(1, "1") }],
    kill: [{ poly: constant(1, "1") }],
    jumps: [
      {
        rate: constant(1, "1"),
        reset: [{ var: 0, poly: mono(1, 0, 1, "1/2") }],
      },
      {
        rate: mono(1, 0, 2, "1"),
        mixture: [
          { weight: "1/2", reset: [{ var: 0, poly: constant(1, "0") }] },
          { weight: "1/4" },
        ],
      },
    ],
  }),
);
\end{tssrc}
\apilang{Rust}{mtk\_catalogue::system::presets::JumpDiffusion::manifest}
\apiprov{signature --- \texttt{crates/mtk-catalogue/src/system/presets.rs}}
\begin{rustsrc}
preset!(
    /// System Catalogue: jump-diffusion profile — drift + diffusion +
    /// jumps at full kernel generality (each row a leg: rate + Dirac
    /// reset or finite sub-stochastic mixture of resets, reset-image
    /// optional), plus an optional killing family. Sub-stochastic
    /// mixture totals are lawful (the deficit is killing); a row
    /// carrying both reset forms is refused, naming the row.
    JumpDiffusion, "jump-diffusion",
    req = [Drift, Jumps],
    allowed = [Diffusion, Noise, Kill],
    extra = |m| {
        diffusion_xor_noise("jump-diffusion", m)?;
        jump_reset_forms("jump-diffusion", m)
    }
);
\end{rustsrc}
\apiprov{construction --- \texttt{crates/mtk-catalogue/tests/builders.rs}}
\begin{rustsrc}
// book: api begin jd_full_construction
// Full kernel generality on one mechanism: drift, diffusion,
// killing, a Dirac leg, and a sub-stochastic two-branch mixture
// leg (total mass 3/4; the deficit is killing).
let m = MechanismBuilder::new()
    .state_var(0, None)
    .drift(None, 0, mono(1, 0, 1, r(-1, 1)))
    .diffusion(None, 0, constant(1, r(1, 2)))
    .kill(None, constant(1, r(1, 20)))
    .jump(JumpRow {
        mode: None,
        rate: constant(1, r(1, 10)),
        reset: vec![],
        mixture: vec![],
        reset_image: None,
    })
    .jump(JumpRow {
        mode: None,
        rate: mono(1, 0, 2, r(1, 1)),
        reset: vec![],
        mixture: vec![
            MixtureBranch {
                weight: r(1, 2),
                reset: vec![ResetRow {
                    var: 0,
                    poly: constant(1, r(0, 1)),
                }],
            },
            MixtureBranch {
                weight: r(1, 4),
                reset: vec![],
            },
        ],
        reset_image: None,
    })
    .finish();
JumpDiffusion::manifest(m).unwrap();
// book: api end jd_full_construction
\end{rustsrc}

\section{Discrete-Time Markov Chains}
\label{sec:mt-dtmc}

\paragraph{The frame.} Finite-state stochastic transition systems;
the bread-and-butter question is certified reachability, with reward
and hitting-time bounds beside it. A two-state weather chain makes
the section's objects concrete: sunny stays sunny with probability
$\tfrac{3}{4}$, rain persists with probability $\tfrac{1}{2}$, and a
sunny start yields the two-step record (sunny, sunny, rain) with
probability $\tfrac{3}{4} \cdot \tfrac{1}{4} = \tfrac{3}{16}$ ---
one stored row spent per step, which is the whole compilation story
below in miniature.

\paragraph{Profile.} Outcome rows $(\mathit{src}, \mathit{dst},
\mathit{reward}, \mathit{prob})$ in exact rational arithmetic, rows
summing to one per source, plus initial-distribution rows. Carrier:
the declared finite state space. Time: counter.

\paragraph{Counter time and the tick.} Composition into real time
passes through the tick declaration of section~\ref{sec:mt-tick}
--- part of
the model type, not a checker convention, with the price of a jitter
bracket over an exact tick certified rather than assumed (the
factor-$2.66$ pair quoted there).

\paragraph{Semantics.} The finite-state path measure, exact;
backward-induction table soundness carried in checker-adjacent prose
with its mechanized statement named. In the universe's terms: the
precise face (Remark~\ref{rem:su-faces}) over discrete records,
glued by the kernel-presented route of
section~\ref{sec:su-denotation}, on which the natural extension
reproduces local data exactly.

\paragraph{Adequacy.} The compilation above is not merely
well-typed; it is \emph{adequate}: compiling and then denoting in
the semantic universe agrees with interpreting classically and then
including the classical object into the universe's precise face.
Stated self-contained, at the working scale of an exam problem:

\begin{problem}[adequacy of the chain compilation]
\label{prob:mt-dtmc}
Let $S$ be a finite set, $P \in \mathbb{Q}^{S \times
S}$ a row-stochastic matrix, $\iota \in \mathbb{Q}^{S}$ a
distribution, and fix a horizon $k$. On the record space $\Omega =
S^{\{0,\dots,k\}}$ write $\mu_{P,\iota}$ for the horizon-$k$
finite-dimensional distribution of the chain,
\[
  \mu_{P,\iota}(\omega) \;=\; \iota(\omega_0)\,\prod_{t < k}
  P(\omega_t, \omega_{t+1}),
\]
the product form of Kemeny--Snell \cite{kemeny-snell-1960} and the
horizon-$k$ generator of the cylinder-set path measure of
Baier--Katoen \cite{baier-katoen-2008}. Say a subprobability mass
function $\mu$ on $\Omega$ \emph{satisfies the compiled rows} when
(i) for every $\omega$, the mass of $\{\omega' : \omega'_0 =
\omega_0\}$ equals the stored initial row $\iota(\omega_0)$, an
absent row read as $0$; and (ii) for every $t < k$ and every
$\omega$, writing $C_t(\omega)$ for the event that a record agrees
with $\omega$ at all times $\le t$, refinement by one step spends
exactly one outcome row:
$\mu(C_{t+1}(\omega)) = \mu(C_t(\omega)) \cdot
P(\omega_t,\omega_{t+1})$, an absent row again read as $0$. Show
that $\mu$ satisfies the compiled rows if and only if $\mu =
\mu_{P,\iota}$ --- in particular the constraints force total mass
one, so the ambient subprobability slack is never taken up.
\end{problem}

\emph{Sketch.} ($\Leftarrow$) Downward induction from $t = k$ gives
the closed form $\mu_{P,\iota}(C_t(\omega)) =
\iota(\omega_0)\prod_{i<t} P(\omega_i,\omega_{i+1})$: partition
$C_t(\omega)$ by the value at time $t+1$ and spend one row sum
$\sum_{s} P(\omega_t, s) = 1$ per step; both constraint families
read off the closed form. ($\Rightarrow$) Upward induction on $t$
pins $\mu(C_t(\omega))$ to the same prefix products; at $t = k$ the
prefix events are the atoms $\{\omega\}$, so all weights agree, and
summing the atoms puts the total mass at one.

In the universe's terms: the two constraint families, read as a
formal conjunction, are the compiled system's denotation at the
per-element stratum of the credal stack, and the theorem says this
loose arrow is isomorphic --- by an invertible globular cell of the
thin double category of per-element constraints --- to the singleton
inclusion of $\mu_{P,\iota}$: the precise face of
Remark~\ref{rem:su-faces}, whose ray form agrees with the singleton
on the radially maximal slice that the verdict operators read. The
statement and both directions are mechanized, sorry-free, in the
toolkit's Lean tree (\texttt{lean/su-functor/}, module
\texttt{SuFunctor.Adequacy.Dtmc}): the pointwise equivalence is
\texttt{satAtoms\_iff}, the invertible cell is
\texttt{dtmc\_adequacy}, and the mass identity is
\texttt{pathLaw\_mass}. The same module ties the exact-rational law
to the standard monadic construction of the chain law in the
mechanized probability library --- iterated \texttt{PMF.bind}
against the kernel --- as \texttt{trajPMF\_eq}, so the classical
side of the square is the community's mechanized object rather than
a private restatement; the inclusion vocabulary shared by the
sibling frames of this catalogue lives beside it
(\texttt{SuFunctor.Adequacy.Inclusions}). The infinite-horizon path
measure is deliberately a separate stratum: the laws above are its
finite-dimensional distributions, and the $\sigma$-additive
extension (Ionescu--Tulcea) belongs to the measure-theoretic tier.

\paragraph{Status.} Certification is the most heavily exercised leg
in the toolkit: the published-table campaign of
section~\ref{sec:exp-sweep} (all seventeen rows of a peer-reviewed
neural-certificate table \cite{abate23} reproduced or refuted with
exact arithmetic) and the tabular-tier campaign of
section~\ref{sec:exp-gym} (dozens of hitting and reward
certificates) are both, at bottom, this frame's machinery.
Representation and production: the toolkit's producer surface
constructs the row grain in both languages and derives the one-step
transfer and the generator values exactly from the stored rows ---
the compiled face of the Adequacy paragraph's square and its
generator kernel. Admission is the row grain's own well-formedness:
every source fibre totals exactly one --- stochasticity is
\emph{declared by} the stored rows here, where the continuous-time
sibling \emph{derives} it from unconstrained rates, the same checker
posture with the opposite direction of creation --- so a
sub-stochastic fibre refuses with its total (a deficit would be
killing, which is the jump-diffusion frame's vocabulary, never the
chain's). The discrete segment-Dynkin identity is exercised in the
test suite as an exact differential over a spread initial
distribution, and the
chain-3 mechanism is a committed cross-language byte-identity vector
whose closed-form control includes the absorbing self-loop's
computed zero (an absorbing state is its explicit self-loop row, the
compiled grain's own representation --- no flag). Open ---
and stated in full, since it is the list's recurring lesson:
the frame has been exercised \emph{beside} the modelling language,
through producer pipelines, faster than it has been represented
\emph{in} it; native outcome rows through the ingestion path are the
gap. The temporal-logic closure beyond reach and reach-avoidance
(nested operators in the style of PCTL$^*$) is designed, and its
first two strata are already in the language: state-level nesting is
the tower discipline of the probability former --- certificate-mediated:
a child certificate's sublevel set enters as a parent
claim's ordinary atom, cited through the judgment stratum
(section~\ref{sec:spec-lang-prob}) --- and
path-level recurrence shapes are definitional Streett instances
(section~\ref{sec:spec-lang-streett}). The genuine surplus ---
temporal operators nested with no intervening threshold --- compiles
by deterministic-monitor product: a finite deterministic monitor
composed with the chain, acceptance read as Streett acceptance on the
product, satisfaction transported back along the projection, with the
quantitative case a reach stratum into the almost-sure accepting
core. The next-step operator and the unbounded-until fence
(section~\ref{sec:spec-lang-funbounded}) are the named gaps; a
certified depth-two artefact is the open exercise.

\syscatmeta{\begin{mathpar}
\inferrule*[Right=\rulename{S-DTMC}]{
  O : \overrightarrow{\langle \mathit{src},\, \mathit{dst},\,
  \mathit{prob},\, \mathit{reward}? \rangle}
  \text{ outcome rows in exact } \mathbb{Q}\text{, unkeyed} \\
  \iota? : \overrightarrow{\langle \mathit{state},\, \mathit{prob}
  \rangle} \text{ initial-distribution rows} \\
  \tau? : \mathsf{TickDecl}
}{
  \mathrm{DTMC}(O) : \mathsf{System}
}
\end{mathpar}
Unkeyed means no $\mathsf{action}$ anywhere ($o.\mathsf{action} =
\bot$, row by row) --- the discipline that separates this former
from the decision former sharing its rows.}
\ruleblock{API}
\sourcenote{Producer surface of the \texttt{mtk} toolkit for this model type: the TypeScript arrangement (npm package \texttt{@maitria/mtk}) and the Rust catalogue crate (\texttt{crates/mtk-catalogue}) --- one Rust implementation underneath, reached from TypeScript in-process through a native binding, so the two surfaces cannot diverge. Signatures and constructions are extracted from the toolkit sources at every regeneration of this catalogue; each construction is quoted verbatim from the toolkit's conformance corpus, code whose emitted certificate bytes are verified byte-identical against the committed golden vectors (helper values named \texttt{fx\_*} or \texttt{f.*} are that corpus's shared fixtures, defined beside the quoted lines). This API constructs and emits; it never checks --- the checker re-derives every claim independently.}
\apilang{TypeScript}{system.dtmc}
\apiprov{signature --- \texttt{ts/src/system.ts}}
\begin{tssrc}
/** S-DTMC: unkeyed outcome rows (+ init distribution, tick). */
export const dtmc = profileFactory("dtmc");

/** The derived DTMC generator kit: the one-step transfer
 * (P V)(s) = Σ_d p(s,d)·V(d) and the generator values
 * 𝒜V(s) = (P V)(s) − V(s) — the mechanized transferRow/genRow
 * contractions, whose recomputed rows are theorem-grade drift
 * premises. Derived, never trusted — the checker re-derives both. */
export interface DtmcGenKit {
  readonly gen: readonly StateVal[];
  readonly transfer: readonly StateVal[];
}

/** Derive the DTMC generator kit from a mechanism declaration and a
 * per-state certificate observable. Admission mirrors the WF
 * hypotheses under which the rows denote a chain: every roster
 * state's outcome fibre totals one (a sub-stochastic fibre refuses
 * with its total — deficit-as-killing belongs to the jump-diffusion
 * former), an action-keyed row refuses (the discipline separating
 * this former from the decision former sharing its rows), a stored
 * zero refuses (absence is the representation). Witness-shaped
 * refusals, never exceptions. */
export function dtmcGenKit(
  init: MechanismInit,
  cert: readonly StateVal[],
): Result<DtmcGenKit, Refusal> {
  const raw = native.sysDtmcGenKit(mechanism(init).nativeHandle, cert) as
    | { kit: DtmcGenKit }
    | { refusal: unknown };
  return "refusal" in raw
    ? refused(narrowRefusal<Refusal>(raw.refusal))
    : ok(raw.kit);
}
\end{tssrc}
\apiprov{construction --- \texttt{ts/tests/system-native.test.ts}}
\begin{tssrc}
return sys.unwrap(
  sys.dtmc({
    name: "chain-3",
    outcomes: [
      { src: 0, dst: 0, prob: "1/3", reward: "0" },
      { src: 0, dst: 1, prob: "2/3", reward: "1" },
      { src: 1, dst: 2, prob: "1", reward: "5" },
      { src: 2, dst: 2, prob: "1" },
    ],
    initDist: [{ state: 0, prob: "1" }],
    tick: { kind: "exact", period: "1/50" },
  }),
);
\end{tssrc}
\apiprov{The quoted construction is the committed cross-language byte-identity vector's own (the chain-3 mechanism), and the generator kit's closed-form control --- transfer and generator on that same presentation, the absorbing self-loop's computed zero included --- rides the test suite beside it. The kit's Rust face is \texttt{mtk\_catalogue::system::dtmc\_gen\_kit}.}
\apilang{Rust}{mtk\_catalogue::system::presets::Dtmc::manifest}
\apiprov{signature --- \texttt{crates/mtk-catalogue/src/system/presets.rs}}
\begin{rustsrc}
preset!(
    /// System Catalogue: DTMC profile — outcome rows (exact-Q,
    /// unkeyed) + init distribution rows.
    Dtmc, "DTMC",
    req = [Outcomes],
    allowed = [InitDist, Tick],
    extra = |m| outcomes_keyed("DTMC", m, false)
);
\end{rustsrc}
\apiprov{construction --- \texttt{crates/mtk-catalogue/tests/common/mod.rs}}
\begin{rustsrc}
let m = MechanismBuilder::new()
    .outcome(OutcomeRow {
        src: 0,
        action: None,
        dst: 0,
        prob: r(1, 3),
        reward: Some(r(0, 1)),
    })
    .outcome(OutcomeRow {
        src: 0,
        action: None,
        dst: 1,
        prob: r(2, 3),
        reward: Some(r(1, 1)),
    })
    .outcome(OutcomeRow {
        src: 1,
        action: None,
        dst: 2,
        prob: r(1, 1),
        reward: Some(r(5, 1)),
    })
    .outcome(OutcomeRow {
        src: 2,
        action: None,
        dst: 2,
        prob: r(1, 1),
        reward: None,
    })
    .init_dist(0, r(1, 1))
    .tick(TickDecl::Exact { period: r(1, 50) })
    .finish();
presets::Dtmc::manifest(m)
    .expect("fixture conforms to the DTMC profile")
    .named("chain-3")
\end{rustsrc}

\section{Continuous-Time Markov Chains}
\label{sec:mt-ctmc}

\paragraph{The frame.} Exponential races over a finite state space;
certified questions are time-bounded and unbounded reachability. A
two-state repair chain is the toy to hold: a machine fails at rate
$\lambda$ and is repaired at rate $\mu$, so its jump chain
alternates deterministically between up and down while the rates set
only the clocks --- the jump-chain/holding-time split at which the
adequacy square below lives.

\paragraph{Profile.} A shipped theory of its own: states, plus
off-diagonal rate entries as a partial function --- zero entries
absent by construction, diagonal determined, so malformed rate
matrices are unrepresentable rather than linted --- plus
generator-certificate rows. Time: continuous.

\paragraph{Semantics.} The finite-state generator; the path measure
is unique (no explosion at finite state), and the Dynkin and Ville
bridges are named obligations, unconditional in this setting but
still deliberately outside the kernel. At finite state the classical
path-measure construction discharges, for this frame, the
continuous-time extension debt section~\ref{sec:su-denotation} names
in general.

\paragraph{Adequacy (jump-chain face).} The compilation above is
adequate at its \emph{time-abstract} stratum: the embedded jump chain
--- which states the process visits, in which order --- is exact
rational data, and compiling the rate rows then denoting agrees with
interpreting the jump chain classically and including into the
universe's precise face.  Stated self-contained, at the working scale
of an exam problem:

\begin{problem}[the jump-chain square]
\label{prob:mt-ctmc}
Let $S$ be a finite set, $q(s, s') \in
\mathbb{Q}_{\ge 0}$ a transition rate for each ordered pair of
distinct states, and $\iota \in \mathbb{Q}^S$ a distribution.  Write
$q(s) = \sum_{s' \neq s} q(s, s')$ for the exit rates, and let $\Pi$
be the jump matrix of Norris \cite{norris-1997}: when $q(s) \neq 0$,
$\Pi(s, s') = q(s, s')/q(s)$ for $s' \neq s$ and $\Pi(s, s) = 0$;
when $q(s) = 0$, the absorbing convention $\Pi(s, s) = 1$ and
$\Pi(s, s') = 0$ otherwise.  Fix a jump horizon $k$ and write
$\mu_{q,\iota}$ for the horizon-$k$ path law of the chain $(\Pi,
\iota)$ on $\Omega = S^{\{0,\dots,k\}}$ --- the record of the first
$k$ jumps, not of clock time.  The compiled presentation stores the
off-diagonal rate rows $(s, s') \mapsto q(s, s')$, an absent row read
as $0$ and a diagonal row unrepresentable; the jump closure is the
\emph{computed} row set $\Pi$, and a subprobability mass function
$\mu$ on $\Omega$ satisfies the closure exactly when it satisfies the
chain constraints of section~\ref{sec:mt-dtmc} at the rows $(\Pi,
\iota)$.  Show: (i) the closed rows are row-stochastic, although the
rate rows carry no row-sum constraint; (ii) $\mu$ satisfies the
closure's constraints if and only if $\mu = \mu_{q,\iota}$ --- in
particular the constraints force total mass one; and (iii)
$\Pi(s, s')\, q(s) = q(s, s')$ for all $s \neq s'$, absorbing states
included.
\end{problem}

\emph{Sketch.}  (i) is one division: at a non-absorbing state the row
sums to $q(s)/q(s)$, at an absorbing state to a point mass --- the
stochasticity of the jump chain is derived, never assumed.  The pair
$(\Pi, \iota)$ is then a finite chain, and both directions of (ii) are
the chain problem verbatim at these rows: the downward induction reads
off the closed form on prefix events, the upward induction pins any
satisfying $\mu$, and mass one follows by summing the atoms.  (iii) is
the recovery direction of the jump-chain/holding-times decomposition:
at an absorbing state both sides vanish, since nonnegative rates
summing to zero vanish individually.  No further stochastic content is
specific to the continuous-time former --- closing the rate rows
before or after compilation stores the same rows, which is the whole
remaining obligation.

In the universe's terms: the jump closure denotes as a loose-arrow
diagram at the per-element stratum, and the theorem exhibits an
invertible globular cell to the singleton inclusion of the jump-chain
law --- the chain square instantiated at the jump chain.  The
statement and both directions are mechanized, sorry-free, in the
toolkit's Lean tree (\texttt{lean/su-functor/}, module
\texttt{SuFunctor.Adequacy.CtmcJumpChain}): the invertible cell is
\texttt{ctmc\_jump\_adequacy}; the compile-commutes triangle ---
definitional, and the one obligation well-formedness alone could not
carry, since any rate table closes to \emph{some} stochastic rows ---
is \texttt{compile\_jump\_comm}; the derived stochasticity is
\texttt{jumpRows\_wf}; the rate factorization is
\texttt{jumpChain\_step\_mul\_exit}; and the profile's
diagonal-determined convention is carried as theorems rather than
prose (\texttt{compileCtmc\_diag\_blind},
\texttt{jumpLaw\_diag\_blind}: the compiled key type cannot express a
diagonal read).  The chain connector transports unchanged
(\texttt{jumpTrajPMF\_eq}), which is the exact extent of the
community tie: the mechanized probability library has no
continuous-time Markov machinery, so the jump chain is tied as the
discrete-time chain it is.  Deliberately separate strata: the
\emph{time-law} stratum --- exponential holding times with parameter
$q(s)$, whose transients are irrational and hence outside this
exact-rational grain --- carries the rest of the continuous-time path
semantics named in the Semantics paragraph above;
\emph{uniformization} is a different discrete reduction (self-loops
against a common rate) and not this one; Markov automata under a
declared scheduler consume this square together with the
decision-former machinery (section~\ref{sec:mt-ma}); the
infinite-horizon extension is separate exactly as for chains.

\paragraph{Status.} Certification: three certified bounds over four
kernel-checked fixtures --- time-bounded reach for an explicit repair
chain ($P(\text{down by } T) \le T/200$, with the equalized
certificate's constant meeting the expected hitting time exactly, no
matrix exponential anywhere), a catastrophe-entry bound exercising
arbitrary rate matrices, and an unbounded-reach bound for a
nine-state epidemic chain (a Ville argument spending two fixtures:
the generator sequent and its nonnegativity side condition), against
an exact ground-truth ruin probability of $35/274391$. Representation
and production: theory and instance rows are in place and compile
through the reference lane, and the toolkit's producer surface
constructs the rate grain in both languages and derives the
generator values, the exit rates, and the jump-closure rows exactly
from the stored entries --- the compiled face of the Adequacy
paragraph's square: the closure's row-stochasticity, a theorem
there, is re-run as arithmetic here, the rate factorization is
exercised in the test suite as an exact differential, and the machine-repair
chain is a committed cross-language byte-identity vector whose
closed-form control is its own certified values ($\mathcal{A}V =
3/755$ at every non-bad state, the equalized certificate's constant,
derived and never trusted). Open: nested temporal operators;
the ingestion gap; closure.

\syscatmeta{\begin{mathpar}
\inferrule*[Right=\rulename{S-CTMC}]{
  R : \overrightarrow{\langle \mathit{src},\, \mathit{dst},\,
  \mathsf{rate} \rangle} \text{ off-diagonal rate entries }
  (\mathit{src} \neq \mathit{dst}) \\
  \iota? \text{ initial-distribution rows}
}{
  \mathrm{CTMC}(R) : \mathsf{System}
}
\end{mathpar}
Rate entries are an off-diagonal partial function: zero entries are
absent by construction and the diagonal is determined, so a malformed
rate matrix is unrepresentable rather than linted.}
\ruleblock{API}
\sourcenote{Producer surface of the \texttt{mtk} toolkit for this model type: the TypeScript arrangement (npm package \texttt{@maitria/mtk}) and the Rust catalogue crate (\texttt{crates/mtk-catalogue}) --- one Rust implementation underneath, reached from TypeScript in-process through a native binding, so the two surfaces cannot diverge. Signatures and constructions are extracted from the toolkit sources at every regeneration of this catalogue; each construction is quoted verbatim from the toolkit's conformance corpus, code whose emitted certificate bytes are verified byte-identical against the committed golden vectors (helper values named \texttt{fx\_*} or \texttt{f.*} are that corpus's shared fixtures, defined beside the quoted lines). This API constructs and emits; it never checks --- the checker re-derives every claim independently.}
\apilang{TypeScript}{system.ctmc}
\apiprov{signature --- \texttt{ts/src/system.ts}}
\begin{tssrc}
/** S-CTMC: off-diagonal rate entries as a partial function. */
export const ctmc = profileFactory("ctmc");

/** One per-state exact-rational value: the discrete-former
 * certificate observable's grain (and the kit outputs' row shape). */
export interface StateVal {
  readonly state: number;
  readonly val: string;
}

/** The derived CTMC generator kit: the generator values
 * 𝒜V(s) = Σ_d q(s,d)·(V(d) − V(s)), the exit rates, and the
 * jump-closure rows (Norris's convention: proportional split at a
 * non-absorbing state, the self-loop point mass at an absorbing
 * one) — row-stochastic by derivation, never by declaration.
 * Derived, never trusted — the checker re-derives all three. */
export interface CtmcGenKit {
  readonly gen: readonly StateVal[];
  readonly exitRates: readonly StateVal[];
  readonly jumpRows: readonly {
    readonly src: number;
    readonly dst: number;
    readonly prob: string;
  }[];
}

/** Derive the CTMC generator kit from a mechanism declaration and a
 * per-state certificate observable. The mechanized contract is the
 * CTMC adequacy square's, jump-chain face: the closure's
 * row-stochasticity is a theorem there (derived stochasticity), and
 * the rate factorization Π(s,d)·q(s) = q(s,d) ties the outputs. The
 * grain fences (diagonal, parallel, or nonpositive rate entries;
 * missing or duplicate certificate values) come back as
 * witness-shaped refusals, never exceptions. */
export function ctmcGenKit(
  init: MechanismInit,
  cert: readonly StateVal[],
): Result<CtmcGenKit, Refusal> {
  const raw = native.sysCtmcGenKit(mechanism(init).nativeHandle, cert) as
    | { kit: CtmcGenKit }
    | { refusal: unknown };
  return "refusal" in raw
    ? refused(narrowRefusal<Refusal>(raw.refusal))
    : ok(raw.kit);
}
\end{tssrc}
\apiprov{construction --- \texttt{ts/tests/system-native.test.ts}}
\begin{tssrc}
// CTMC diagonal rate entry: index and state.
const diag = sys.ctmc({
  rates: [
    { src: 0, dst: 1, rate: "1" },
    { src: 1, dst: 1, rate: "2" },
  ],
});
assert.ok(!diag.ok);
\end{tssrc}
\apiprov{The TS corpus exercises this factory through its refusal geometry (the corpus's discipline that failure geometry is routing advice); the quoted construction is that witness, and the happy-path byte authority is the Rust corpus opposite --- the machine-repair mechanism is a committed cross-language byte-identity vector, and the generator kit's closed-form control is its own certified values. The kit's Rust face is \texttt{mtk\_catalogue::system::ctmc\_gen\_kit}.}
\apilang{Rust}{mtk\_catalogue::system::presets::Ctmc::manifest}
\apiprov{signature --- \texttt{crates/mtk-catalogue/src/system/presets.rs}}
\begin{rustsrc}
preset!(
    /// System Catalogue: CTMC profile — off-diagonal rate entries as
    /// a partial function (zero entries absent; diagonal determined,
    /// never stored).
    Ctmc, "CTMC",
    req = [Rates],
    allowed = [InitDist],
    extra = no_diagonal_rates
);
\end{rustsrc}
\apiprov{construction --- \texttt{crates/mtk-catalogue/tests/builders.rs}}
\begin{rustsrc}
let ctmc = MechanismBuilder::new().rate(0, 1, r(3, 2)).finish();
Ctmc::manifest(ctmc).unwrap();
\end{rustsrc}

\section{Markov Decision Processes}
\label{sec:mt-mdp}

\paragraph{The frame.} Kernels keyed by action; the frame's whole
point is quantification over policies. The toy to hold: a machine
that can be repaired now (a certain cost) or run on (a risk of
failure) --- two actions at every state, and the question that makes
the frame is a guarantee holding whichever way the operator chooses,
forever.

\paragraph{Profile.} Outcome rows keyed by action; a policy is a
separate mechanism, composed in --- the closed loop is a composite,
not a monolith. Time: counter.

\paragraph{Semantics.} Per-policy path semantics as for chains; the
quantification over policies lives in the credal (imprecise-measure)
layer of the System semantics
(Definition~\ref{def:su-credal-denotation}), where a specification
reads the denotation at a stated degree of pessimism.

\paragraph{Adequacy (fixed-policy face).} The closed-loop composite
is adequate in the same sense as the chain compilation of
section~\ref{sec:mt-dtmc} --- compiling both mechanisms, closing the
loop at the row level, and denoting agrees with interpreting
classically and including into the universe's precise face --- and
the proof \emph{is} the chain square, reused wholesale.  Stated
self-contained, at the working scale of an exam problem:

\begin{problem}[the fixed-policy square]
\label{prob:mt-mdp}
Let $S$ and $A$ be finite sets, $p \in
\mathbb{Q}^{S \times A \times S}$ a family of kernels with $\sum_{s'}
p(s' \mid s, a) = 1$ for every state--action pair, $\iota \in
\mathbb{Q}^S$ a distribution, and $\pi$ a memoryless randomized
policy: at each state, a distribution over actions --- the stationary
randomized Markov class of Puterman \cite{puterman-1994}, a
deterministic policy being the unit-mass special case.  Fix a horizon
$k$ and write, on $\Omega = S^{\{0,\dots,k\}}$,
\[
  \mu_{p,\iota,\pi}(\omega) \;=\; \iota(\omega_0) \prod_{t < k}
  \sum_{a} \pi(a \mid \omega_t)\, p(\omega_{t+1} \mid \omega_t, a),
\]
the fixed-policy trajectory law.  The compiled presentation keeps the
two mechanisms separate: action-keyed outcome rows $(s, a, s')
\mapsto p(s' \mid s, a)$ and policy rows $(s, a) \mapsto \pi(a \mid
s)$, an absent row in either read as $0$; the closed loop is the
\emph{computed} row set $P_\pi(s, s') = \sum_a \pi(a \mid s)\,
p(s' \mid s, a)$, and a subprobability mass function on $\Omega$
satisfies the composite exactly when it satisfies the chain
constraints of section~\ref{sec:mt-dtmc} at the rows $(P_\pi,
\iota)$.  Show that $\mu$ satisfies the composite's constraints if
and only if $\mu = \mu_{p,\iota,\pi}$ --- in particular the
constraints force total mass one.
\end{problem}

\emph{Sketch.} One exchange of finite sums shows the mixed matrix
$P_\pi$ is row-stochastic --- $\sum_{s'} P_\pi(s, s') = \sum_a
\pi(a \mid s) \sum_{s'} p(s' \mid s, a) = 1$ --- spending each policy
row's total exactly once against the per-action row sums.  The pair
$(P_\pi, \iota)$ is therefore a finite chain, and $\mu_{p,\iota,\pi}$
is by construction its product-form path law.  Both directions are
then the chain problem verbatim at these rows: the downward induction
reads off the closed form on prefix events, the upward induction pins
any satisfying $\mu$, and mass one follows by summing the atoms.  No
further stochastic content is specific to the decision former ---
closing the loop before or after compilation stores the same rows,
which is the whole remaining obligation.

In the universe's terms: the composite denotes as a loose-arrow
diagram at the per-element stratum, and the theorem exhibits an
invertible globular cell to the singleton inclusion of the
fixed-policy law --- the same form as the chain square, because it
\emph{is} the chain square instantiated at the closed loop.  The
statement and both directions are mechanized, sorry-free, in the
toolkit's Lean tree (\texttt{lean/su-functor/}, module
\texttt{SuFunctor.Adequacy.Mdp}): the invertible cell is
\texttt{mdp\_adequacy}, the compile-commutes triangle ---
definitional, and the one obligation well-formedness alone could not
carry, since any single action's rows are also row-stochastic --- is
\texttt{compile\_close\_comm}, and the mass identity is
\texttt{pathLaw\_mass}.  The module also ties policy selection to the
mechanized probability library's own vocabulary: the closed-loop
kernel is the monadic composition of the policy's action distribution
with the action-indexed kernel (\texttt{closedLoop\_stepPMF}, one
\texttt{PMF.bind}), and the trajectory-law connector of the chain
section transports unchanged (\texttt{trajPMF\_eq}).  Deliberately
separate strata: the \emph{open-policy} face --- the claim uniform
over all policies, whose inclusion lands a \emph{set} of laws in the
credal layer rather than a singleton --- is the imprecise reading of
the Semantics paragraph above and is not this square (it is the
next);
\emph{time-dependent} policies close to a time-inhomogeneous chain,
outside the stationary grain stated here; Markov automata under a
declared scheduler share this square's mixing machinery but carry
exponential holding times, so their square is \emph{continuous-time}
material (section~\ref{sec:mt-ma}); the infinite-horizon
extension is separate exactly as for chains.

\paragraph{Adequacy (demonic face).} The open-policy face is the
first \emph{genuinely imprecise} adequacy statement: what compilation
must preserve is no longer a single law but the \emph{set} of laws the
decision process can exhibit, and the square says the compiled,
policy-free presentation determines that set exactly.  Stated
self-contained, at exam scale:

\begin{problem}[the demonic square]
\label{prob:mt-mdp-demonic}
Data as in the fixed-policy problem above --- finite
$S$ and $A$, action-keyed kernels $p$, initial distribution $\iota$
--- but with \emph{no} policy declared.  The demonic semantics of the
decision process at horizon $k$ is the set of its closed-loop
trajectory laws,
\[
  D \;=\; \bigl\{\, \mu_{p,\iota,\pi} \;:\; \pi \text{ a memoryless
  randomized policy} \,\bigr\},
\]
the adversary ranging over the stationary randomized class of Puterman
\cite{puterman-1994} --- the scheduler-quantified semantics of
probabilistic model checking \cite{baier-katoen-2008}, and demonic
choice in the sense of McIver and Morgan \cite{mciver-morgan-2005}.
On the compiled side, call a subprobability mass function $\mu$ on
$\Omega$ \emph{admitted} by the action-keyed rows when \emph{some}
well-formed policy-row set closes the loop onto chain constraints that
$\mu$ satisfies --- the rows carry no policy family; the demon's
choice is the separate mechanism, existentially quantified.  Show:
(i) $\mu$ is admitted if and only if $\mu \in D$; (ii) every member of
$D$ has total mass one, and consequently the downward-closed credal
set generated by $D$ --- the union of the segments $[0,1] \cdot \mu$
over $\mu \in D$ --- has exactly $D$ as its set of radially maximal
elements; (iii) for an event $E$ and threshold $q$, the credal set's
lower verdict for $P(E) \ge q$, and its upper verdict for the
downward-closed constraint $P(E) \le q$, each hold if and only if the
corresponding bound holds under \emph{every} policy.
\end{problem}

\emph{Sketch.} Claim (i) is the fixed-policy problem, quantified: a
witnessing scheduler's rows assemble a memoryless policy whose closed
loop the chain constraints pin, and conversely every policy compiles
to a witnessing scheduler --- no new stochastic content.  In (ii),
mass one is the chain mass identity at each closed loop; for the
slice, a strictly interior point of a segment escapes upward to its
own base, while an upward scaling of one member landing on
\emph{another} member's segment would have to scale mass past the
common value --- equal mass is exactly what forbids cross-member
domination, and it is not decorative: in the unequal-mass pair
$\{\mu, \tfrac{1}{2}\mu\}$ the half element is dominated across
segments and drops out of the slice.  In (iii), the lower verdict
reads the radially maximal slice
(Remark~\ref{rem:su-maximal-slice}), which by (ii) is $D$ itself;
the upper verdict reads the whole set, and downward closure of
$P(E) \le q$ makes it blind to the segments' interiors.  Reading a
lower bound against the whole set instead would refuse every positive
threshold --- the zero measure lies in every segment --- which is the
slice discipline's whole content, surviving the passage to sets.

In the universe's terms: the demonic reading of the compiled former
and the union-of-singletons inclusion of $D$ are isomorphic by an
invertible globular cell at the per-element stratum --- the first
adequacy square whose bottom edge is set-valued, one singleton
inclusion per policy (\texttt{unionSingletons}) --- and the credal
presentation, the union of the members' rays
(\texttt{FinMDP.demonicCredal}), reconciles with $D$ through the
radially maximal slice, which is what the verdict operators read.
The presentation is radial, and in general \emph{not} convex: a
mixture of two deterministic-policy path laws is realized by no
memoryless policy, only by an ex-ante correlated scheduler, so the
closed convex hull is the coarser convex reading --- and the slice
reconciliation is a statement the hull would coarsen.  The
mechanized results consume radiality and equal mass only;
closedness of the policy-to-law image, which would make the union a
credal body outright, is not among them.
The statement and both directions are mechanized, sorry-free, in the
toolkit's Lean tree (\texttt{lean/su-functor/}, modules
\texttt{SuFunctor.Adequacy.EnvelopeUnion} and
\texttt{SuFunctor.Adequacy.MdpDemonic}): the invertible cell is
\texttt{mdp\_demonic\_adequacy}, with admission predicate
\texttt{demonicC} and feeding lemma \texttt{demonicC\_iff}; the equal-mass slice reconciliation is
\texttt{maxSlice\_unionRays}, its family-bound direction
hypothesis-free (\texttt{maxSlice\_unionRays\_subset}),
instantiated at mass one by \texttt{demonic\_maxSlice}; the
uniform-over-policies verdict theorems are
\texttt{satCredal\_lower\_demonicCredal} and
\texttt{satCredal\_upper\_demonicCredal}; and the degeneracies are
mechanized rather than smoothed: the whole-set reading of a lower
bound refuses every positive threshold
(\texttt{not\_satCredal\_upper\_demonicCredal\_geC}); an
action-free process is well-formed yet admits no policy at all
(\texttt{isEmpty\_policy\_iff}), so both verdict kinds pass
vacuously (the $\forall$-over-no-adversaries convention, the
empty-policy union staying below bottom with no body); and the
unequal-mass pair witnesses the slice lemma's hypothesis
(\texttt{not\_maxSlice\_unionRays\_halfPair}).  Deliberately separate
strata, unchanged from the fixed-policy face: memoryless policies
only; the \emph{angelic} reading --- existential over policies --- is
a different operator from the universally quantified verdict pair and
is not built here; the infinite-horizon extension is separate exactly
as for chains.

\paragraph{Status.} An asymmetric grade, stated as such.
\emph{Fixed policy}: shipped --- the controller-checkpoint
certificates of the experiments chapters are policy-composed MDP
verification, and the exact value-iteration route
\cite{hartmanns19, quatmann18} exports value functions as
certificate input. The toolkit's producer surface now
constructs the closed loop itself in both languages: the policy
arrives as the scheduler grain --- the compiled face of the composed
policy mechanism, never a family of the decision former, exactly the
factoring the formation rule below states --- and the kit derives
the closed chain (self-loops kept, initial rows verbatim) and its
generator by handing the closure to the chain kit, the square's
engine reuse run producer-side. The mixing exchange --- the
compile-commutes triangle --- is re-derived as an exact differential
in the test suite through a read that never touches the closure,
with the deterministic read-through as its control; policy support
discipline is the automaton's (support within the enabled pairs,
mass one at every state, and a state with no enabled actions refuses
--- nothing sound closes there, the empty-action degeneracy's
producer face); and the patrol mechanism is a committed
cross-language byte-identity vector. \emph{Uniform over policies}: the
open-nondeterminism question is designed at the semantics layer and
exercised by workshop fixtures with checker acceptance --- not yet by
a published reproduction target. Until that leg ships as a
reproduction target, this frame is graded on its closed-loop half
alone --- which is precisely the
half that does not distinguish it from
section~\ref{sec:mt-dtmc}. Open: the quantified leg end to end; the
ingestion gap; closure.

\syscatmeta{\begin{mathpar}
\inferrule*[Right=\rulename{S-MDP}]{
  O : \overrightarrow{\langle \mathit{src},\, \mathsf{action},\,
  \mathit{dst},\, \mathit{prob},\, \mathit{reward}? \rangle}
  \text{ outcome rows, every row action-keyed} \\
  \iota? \text{ initial-distribution rows} \\
  \tau? : \mathsf{TickDecl}
}{
  \mathrm{MDP}(O) : \mathsf{System}
}
\end{mathpar}
The same outcome rows as the chain former, all action-keyed; the
policy is deliberately \emph{not} a premise --- it is a separate
mechanism, composed in, so the closed loop is a composite.}
\ruleblock{API}
\sourcenote{Producer surface of the \texttt{mtk} toolkit for this model type: the TypeScript arrangement (npm package \texttt{@maitria/mtk}) and the Rust catalogue crate (\texttt{crates/mtk-catalogue}) --- one Rust implementation underneath, reached from TypeScript in-process through a native binding, so the two surfaces cannot diverge. Signatures and constructions are extracted from the toolkit sources at every regeneration of this catalogue; each construction is quoted verbatim from the toolkit's conformance corpus, code whose emitted certificate bytes are verified byte-identical against the committed golden vectors (helper values named \texttt{fx\_*} or \texttt{f.*} are that corpus's shared fixtures, defined beside the quoted lines). This API constructs and emits; it never checks --- the checker re-derives every claim independently.}
\apilang{TypeScript}{system.mdp}
\apiprov{signature --- \texttt{ts/src/system.ts}}
\begin{tssrc}
/** S-MDP: action-keyed outcome rows; the policy is a separate
 * mechanism composed in (closed-loop composite), never a family. */
export const mdp = profileFactory("mdp");

/** The derived MDP fixed-policy kit: the closed discrete-time chain
 * as ordinary DTMC rows (self-loops kept, initial rows verbatim) and
 * the closed chain's generator kit — the mechanized square's engine
 * reuse (mdp_adequacy := dtmc_adequacy at the closed loop), run
 * producer-side. Derived, never trusted — the checker re-runs the
 * closure and the contraction. */
export interface MdpGenKit {
  readonly closedRows: readonly {
    readonly src: number;
    readonly dst: number;
    readonly prob: string;
  }[];
  readonly closedInit: readonly {
    readonly state: number;
    readonly prob: string;
  }[];
  readonly closed: DtmcGenKit;
}

/** Derive the MDP fixed-policy kit from a mechanism declaration, a
 * policy, and a per-state certificate observable. The policy arrives
 * as the scheduler grain — the compiled face of the COMPOSED POLICY
 * MECHANISM (the decision former deliberately has no policy family;
 * an absent policy is the demonic face, which quantifies
 * certificate-side and needs nothing from this surface). Policy
 * support discipline as at the automaton: support within the enabled
 * pairs, mass one at every roster state, and a state with no enabled
 * actions refuses (nothing sound closes there). Witness-shaped
 * refusals, never exceptions. */
export function mdpGenKit(
  init: MechanismInit,
  policy: SchedulerDecl,
  cert: readonly StateVal[],
): Result<MdpGenKit, Refusal> {
  const raw = native.sysMdpGenKit(
    mechanism(init).nativeHandle,
    {
      rows: policy.rows.map((r) => ({
        state: r.state,
        action: r.action,
        prob: qText(r.prob),
      })),
      note: policy.note,
    },
    cert,
  ) as { kit: MdpGenKit } | { refusal: unknown };
  return "refusal" in raw
    ? refused(narrowRefusal<Refusal>(raw.refusal))
    : ok(raw.kit);
}
\end{tssrc}
\apiprov{construction --- \texttt{ts/tests/mdpgen-native.test.ts}}
\begin{tssrc}
/** The `patrol-2s2a` presentation (the committed golden fixture's):
 * at state 0, `stay` self-loops and `go` branches 1/4 / 3/4; at
 * state 1, `go` branches evenly. The policy is passed SEPARATELY —
 * the decision former has no policy family, and the scheduler grain
 * here is the compiled face of the composed policy mechanism. */
function patrolInit() {
  return {
    outcomes: [
      { src: 0, action: 0, dst: 0, prob: "1" },
      { src: 0, action: 1, dst: 0, prob: "1/4" },
      { src: 0, action: 1, dst: 1, prob: "3/4" },
      { src: 1, action: 1, dst: 0, prob: "1/2" },
      { src: 1, action: 1, dst: 1, prob: "1/2" },
    ],
    initDist: [{ state: 0, prob: "1" }],
  } as const;
}

const patrolPolicy = {
  rows: [
    { state: 0, action: 0, prob: "1/3" },
    { state: 0, action: 1, prob: "2/3" },
    { state: 1, action: 1, prob: "1" },
  ],
} as const;

const patrolCert = [
  { state: 0, val: "0" },
  { state: 1, val: "1" },
];

test("the patrol closes to the all-halves chain under its policy", { skip }, async () => {
  const sys = await ns();
  const kit = sys.unwrap(
    sys.mdpGenKit(patrolInit(), patrolPolicy, patrolCert),
  );
  const closed = new Map(
    kit.closedRows.map((row) => [`${row.src}->${row.dst}`, row.prob]),
  );
  assert.equal(closed.get("0->0"), "1/2"); // self-loop KEPT
  assert.equal(closed.get("0->1"), "1/2");
  assert.equal(closed.get("1->0"), "1/2");
  assert.equal(closed.get("1->1"), "1/2");
  assert.deepEqual(kit.closedInit, [{ state: 0, prob: "1" }]);
  const gen = new Map(kit.closed.gen.map((v) => [v.state, v.val]));
  assert.equal(gen.get(0), "1/2");
  assert.equal(gen.get(1), "-1/2");
});
\end{tssrc}
\apiprov{The quoted construction is the committed cross-language byte-identity vector's own (the patrol mechanism), with the policy passed separately --- the decision former has no policy family, and the scheduler grain at the kit call is the compiled face of the composed policy mechanism. The kit's Rust face is \texttt{mtk\_catalogue::system::mdp\_gen\_kit}; the refusal-geometry witness remains in the TS corpus beside it.}
\apilang{Rust}{mtk\_catalogue::system::presets::Mdp::manifest}
\apiprov{signature --- \texttt{crates/mtk-catalogue/src/system/presets.rs}}
\begin{rustsrc}
preset!(
    /// System Catalogue: MDP profile — outcome rows keyed by action;
    /// the policy is a separate mechanism composed in (closed-loop
    /// composite), never a family here.
    Mdp, "MDP",
    req = [Outcomes],
    allowed = [InitDist, Tick],
    extra = |m| outcomes_keyed("MDP", m, true)
);
\end{rustsrc}
\apiprov{construction --- \texttt{crates/mtk-catalogue/tests/builders.rs}}
\begin{rustsrc}
let mdp = MechanismBuilder::new()
    .outcome(OutcomeRow {
        src: 0,
        action: Some(0),
        dst: 0,
        prob: r(1, 1),
        reward: None,
    })
    .finish();
Mdp::manifest(mdp).unwrap();
\end{rustsrc}

\section{Partially Observable Markov Decision Processes}
\label{sec:mt-pomdp}

\paragraph{The frame.} Kernels keyed by action, as for the decision
former --- plus an observation channel, and the frame's whole point
is that the policy sees observations, not states.  The classical
formulation is {\AA}str{\"o}m's \cite{astrom-1965}: optimal control
of a Markov process with incomplete state information.  The toy to
hold is the classic two-door problem --- the \emph{tiger} problem of
the planning literature \cite{kaelbling-littman-cassandra-1998}: a
prize behind one of two doors, a hazard behind the other, and a
listening action whose report is right only most of the time.  The question that makes the frame is
a guarantee on achieved value for a controller that can never look
behind the doors --- only accumulate evidence about them.

\paragraph{Profile: none --- a composition pattern.} This frame adds
no formation rule, no new mechanism family, and no new semantic
species; it stands to the decision former of section~\ref{sec:mt-mdp}
as the delay equation of section~\ref{sec:mt-dde} stands to the ODE
cell --- a composite of inventory the catalogue already carries,
presented through a wiring.  The pieces: the \emph{environment} is a
decision-former mechanism, action-keyed outcome rows over a state the
composite will hide; the \emph{observation channel} is an ordinary
mechanism whose rows carry kernels from state to observation symbols
(an action-dependent channel folds in by state augmentation); the
\emph{controller} is a separate mechanism whose inputs are the
channel's outputs, never the environment's state; and the loop closes
through a \emph{delayed trace} --- the feedback row of the
composition-patterns chapter's delay inventory
(section~\ref{sec:cp-dtrace}), which names this very loop as its
flagship construction, and whose one-step degenerate case, the tick
register, is the medium a counter-time closed loop needs.  Where
the delay-equation composite ran its feedback through a continuous
transport medium, this one runs it through the tick register: the
same pattern family, at the counter-time profile.

The wiring is not packaging; it is the frame's semantic content.
What distinguishes a partially observable decision process from the
decision former is \emph{which wires reach the policy}, and the
composite carries that distinction as typing: the controller's inputs
are, by the shape of the pattern, the observation channel's outputs
together with its own state, so there is no way to write an instance
of the pattern whose policy reads hidden state.  Partial
observability is a composition pattern, not a new model class.

\paragraph{Semantics.} Per-policy path semantics on the closed loop's
trajectory law, exactly as for the decision former, with the
quantification over policies again living in the credal layer
(Definition~\ref{def:su-credal-denotation}) --- and the one genuinely
new ingredient relative to section~\ref{sec:mt-mdp} is the class
quantified over: \emph{observation-measurable} policies, policies
factoring through the history of observations and of their own past
actions, never the history of states.  Stated classically, the policy
is adapted to the observation filtration; stated compositionally, the
restriction is carried by the wiring itself, which is what the
paragraph above bought --- the measurable-selection bookkeeping that
the frame's textbook development devotes to this restriction attaches
here to compile-level metatheory, part of the square owed below, not
of the presentation.  The belief-state
reformulation enters the same way, as presentation rather than
semantics: {\AA}str{\"o}m's re-reading of the loop as a fully
observed decision process on the belief simplex \cite{astrom-1965}
is, in this inventory, a particular choice of controller --- a
belief-tracking mechanism whose internal state is a distribution over
environment states, updated observation by observation, composed in
like any other policy.  (In unnormalized coordinates the exact Bayes
update is polynomial --- linear in the belief and in each mechanism's
rows, with no division --- hence ordinary mechanism content; the
normalized update, a ratio, is not, which is why the unnormalized
form is the one the inventory stages.)
The simplex is controller state, not a new semantic carrier.

\paragraph{Adequacy.} Owed, and stated as such: no square is
mechanized for this frame yet.  What the fixed-policy face must say
is engine reuse once more --- at a declared observation-measurable
policy, the environment, channel, and controller compile to rows, the
loop closes at the row level, and the composite is a chain on
finitely many reachable states at any fixed horizon (one belief per
action--observation history), so the obligation has the shape the decision
former discharged in section~\ref{sec:mt-mdp}, with the controller's
belief coordinate folded into the composite state.  The quantified
face inherits the demonic square's set-valued form, with the
adversary now ranging over the observation-measurable class ---
policies reading strictly less than the state, though carrying memory
the stationary class of section~\ref{sec:mt-mdp} deliberately
excluded --- and the delayed trace is what makes the restriction
structural rather than side-conditional.  The intended first
certificate is a lower bound on achieved value for the two-door
problem, living on the belief simplex: a finite family of affine
functionals on the simplex --- the $\alpha$-vector presentation of
Smallwood and Sondik \cite{smallwood-sondik-1973} --- each dominated
by its witness action's Bellman backup under the discounted
criterion, certifying at once a value floor and the policy that
achieves it.

\paragraph{Status.} Design of record, with the certificate route
already exercised.  The wiring --- environment,
observation channel, belief-tracking controller, delayed trace --- is
the composition-patterns chapter's own flagship construction
(section~\ref{sec:cp-dtrace}), the producer surface for the delayed
trace is landed there, and the decision former's closed-loop producer
machinery is the engine the fixed-policy face will reuse.  The
two-door lower bound has been produced and checked in the workshop by
the route the adequacy paragraph names --- an exact-rational
$\alpha$-vector certificate with per-state Bellman-domination
witnesses, re-verified by a small exact checker --- not yet as a
published reproduction target, so the frame is graded here on design
and workshop evidence alone.  Open: the fixed-policy square; the
quantified face end to end; promotion of the two-door certificate to
a reproduction target; and the ingestion gap, as everywhere.

No formation rule accompanies this section: that absence is the
section's content, for the same reason as in
section~\ref{sec:mt-dde} --- the frame \emph{is} a wiring of formers
the catalogue already carries, and a dedicated former would store as
data what the composite derives by shape.

\section{Markov Automata}
\label{sec:mt-ma}

\paragraph{The frame.} Exponential races, probabilistic branching,
and nondeterministic action under maximal progress --- the frame that
composes sections~\ref{sec:mt-ctmc} and~\ref{sec:mt-mdp}.

\paragraph{Profile.} Assembled by composition, not by a new theory:
the chain frame's rate idiom, the decision frame's action-kernel
idiom aimed at immediate transitions, and exactly one new atom --- a
scheduler declaration --- with a closure bridge delivering the closed
object as an ordinary continuous-time chain, certificate and all.
Maximal progress is stated as structure: timed rows at instantaneous
states are unrepresentable.

\paragraph{Semantics.} The interleaving semantics via a
fixed-scheduler reduction note (closure formula, self-loop thinning,
Zeno-freedom as the closure's checked assertion), with path-measure
bridges inherited unchanged from the chain frame.

\paragraph{Adequacy (fixed-scheduler face).} The closure bridge above
is adequate at the automaton's \emph{time-abstract} stratum: under a
declared scheduler, closing the compiled rows to a chain and then
reading its jump skeleton agrees with interpreting the automaton
classically and including into the universe's precise face.  Stated
self-contained, at the working scale of an exam problem:

\begin{problem}[the fixed-scheduler square]
\label{prob:mt-ma}
Let $S$ and $J$ be finite sets (the Markovian and the
instantaneous states of a closed Markov automaton in the sense of
Eisentraut, Hermanns, and Zhang \cite{eisentraut-hermanns-zhang-2010}),
$A$ a finite action set.  The data are: rates $q(s, x) \in
\mathbb{Q}_{\ge 0}$ for $s \in S$, $x \in S \sqcup J$ with $x \neq s$
--- a timed row at an instantaneous state and a diagonal entry are
both \emph{unrepresentable}, which is the maximal-progress assumption
and the determined diagonal stated as structure --- immediate kernels
$b(j, a, \cdot)$, each a distribution on $S \sqcup J$; an initial
distribution $\iota$ on $S \sqcup J$; a scheduler $\varsigma$
assigning each $j \in J$ a distribution over $A$; and a \emph{rank}
$\rho \colon J \to \mathbb{N}$ such that the mixed immediate kernel
$f(j, x) = \sum_a \varsigma(j)(a)\, b(j, a, x)$ satisfies
$\rho(j') < \rho(j)$ whenever $f(j, \mathrm{in}(j')) \neq 0$ --- the
Zeno-freedom assertion, carried as a ranking witness.  Define the
absorption kernel $P$ by iterating $f$ (fuel one past the largest
rank), a point mass at each Markovian state, and the closed chain by
$q^{\varsigma}(s, s') = \sum_{x} q(s, x)\, P(x)(s')$ for $s' \neq s$,
with initial distribution $\iota$ pushed through $P$.  Show: (i) $P$
is a distribution over $S$ from \emph{every} state, although neither
the rate rows nor the iteration carried a normalization constraint,
and it is the \emph{unique} solution of the absorbing-system
equations; (ii) the closed exit rate is the original exit rate minus
the mass $P$ returns to the source (self-loop thinning), so a
scheduler can only remove activity; (iii) writing $q(s)$ for the exit
rates and $\Pi$ for the jump matrix of the automaton's timed part on
$S \sqcup J$ (instantaneous states rate-absorbing), the closure also
reads $q^{\varsigma}(s, s') = q(s) \sum_x \Pi(\mathrm{in}(s), x)\,
P(x)(s')$ --- the race-then-jump-then-absorb spelling agrees with the
rate spelling; and (iv) a subprobability mass function $\mu$ on the
jump-horizon-$k$ records of $S$ satisfies the jump closure of the
scheduler closure of the compiled rows if and only if $\mu$ is the
closed chain's jump-chain path law.
\end{problem}

\emph{Sketch.}  (i) is an induction along the rank: the mixed kernel
is row-stochastic by the policy mixing exchange of
section~\ref{sec:mt-mdp}, and past the rank the iteration has
stabilized --- Zeno-freedom is spent exactly here, since an immediate
self-loop would never absorb; uniqueness is the same induction applied
to the difference of two solutions.  (ii) is the mass computation of
(i) summed against the rate rows.  (iii) is the rate-factorization
identity of section~\ref{sec:mt-ctmc} --- $\Pi(s, x)\, q(s) = q(s, x)$,
absorbing states included --- applied under the absorption sum.  (iv)
is the chain problem verbatim at the closed rows: the one genuinely
automaton-shaped obligation is that scheduler-closing the compiled
rows \emph{is} the compilation of the classically closed chain, which
holds definitionally when both routes read the same rows through the
same guarded reads and the same absorbing iteration --- and it is what
pins the closure to the declared scheduler and the declared rates,
since well-formedness alone can see neither.

In the universe's terms: the scheduler closure followed by the jump
closure denotes as a loose-arrow diagram at the per-element stratum,
and the theorem exhibits an invertible globular cell to the singleton
inclusion of the closed jump law --- the chain square instantiated
twice, three storeys of engine reuse with no new induction.  The
statement and both directions are mechanized, sorry-free, in the
toolkit's Lean tree (\texttt{lean/su-functor/}, module
\texttt{SuFunctor.Adequacy.MarkovAutomaton}): the invertible cell is
\texttt{ma\_adequacy}; the rate-closure triangle --- definitional ---
is \texttt{compile\_closeMa\_comm}; the acyclic-absorption package is
\texttt{absorbIter\_mass}, \texttt{absorbIter\_unique}, and
\texttt{absorbIter\_stab}; the two closure spellings are
\texttt{close\_offRate\_factor}, consuming the chain face's
\texttt{jumpChain\_step\_mul\_exit}; the thinning identity and its
inequality are \texttt{close\_exitRate\_thinning} and
\texttt{close\_exitRate\_le}; and the profile's structural
conventions are carried as theorems (\texttt{compileMa\_diag\_blind}:
the compiled key type can express neither a self-rate nor a timed row
at an instantaneous state).  The scheduler grain is the decision
former's policy grain reused verbatim, and the community tie is the
chain connector transported through both closures
(\texttt{FinMA.trajPMF\_eq}) --- the mechanized probability library
has no Markov-automaton machinery, and the connector says so rather
than dressing the transport as more.  Deliberately separate strata:
the exponential \emph{time-law} stratum, inherited from
section~\ref{sec:mt-ctmc}; \emph{scheduler quantification}, the
envelope stratum, addressed by the demonic face in the next paragraph;
\emph{parallel composition and visible actions}, the open-automaton
face (this square stages one closed component); history-dependent
schedulers; and non-Zeno-free automata, outside the square by
hypothesis rather than silently discharged.

\paragraph{Adequacy (demonic face).} The scheduler-quantification
face, and the set-valued square at continuous time --- the automaton
sibling of section~\ref{sec:mt-mdp}'s demonic face, with one genuinely
automaton-shaped complication: each scheduler's closure consumes a
Zeno-freedom certificate.  Stated self-contained, at exam scale:

\begin{problem}[the demonic square, at continuous time]
\label{prob:mt-ma-demonic}
Data as in the fixed-scheduler problem above ---
finite $S$, $J$, $A$; rates $q$; immediate kernels $b$; initial
distribution $\iota$ --- but with \emph{no} scheduler declared.
Instead, a single rank $\rho \colon J \to \mathbb{N}$ is given that
strictly decreases along every immediate \emph{branch} row:
$\rho(j') < \rho(j)$ whenever $b(j, a, \mathrm{in}(j')) \neq 0$.  (i)
Show one rank serves every adversary: for each scheduler $\varsigma$,
the mixed kernel $f_\varsigma$ satisfies the fixed-scheduler problem's
Zeno-freedom assertion with the same $\rho$, so the closed chain and
its jump law $\mu_\varsigma$ are defined for all $\varsigma$
simultaneously.  (ii) Show the object outranks its certificate: if
$\rho'$ is any other rank certifying some scheduler's mixed kernel,
the closure at $\rho'$ yields the same law --- so the demonic
semantics at horizon $k$,
\[
  D \;=\; \bigl\{\, \mu_\varsigma \;:\; \varsigma \text{ a memoryless
  randomized scheduler} \,\bigr\},
\]
depends on the automaton and the horizon alone.  (iii) On the compiled
side, call $\mu$ \emph{admitted} when \emph{some} well-formed
scheduler-row set closes the compiled automaton rows, at the rank's
working fuel, onto chain constraints that $\mu$ satisfies; show $\mu$
is admitted if and only if $\mu \in D$, and that every member of $D$
has mass one, so the downward-closed credal set generated by $D$ has
exactly $D$ as its radially maximal slice and the verdict pair reads
as bounds uniform over \emph{every} scheduler.  (iv) Show the
certificate forces somewhere to absorb into: if $J$ and $A$ are
inhabited, a branch-certified automaton must have a Markovian state.
\end{problem}

\emph{Sketch.}  (i) is support bookkeeping: a mixture's support lies
in the union of the branch supports, so a rank uniform in the branches
is uniform in the adversary --- the quantifier commutes through the
certificate.  (ii) both certificates are spent at a common fuel past
both ranks, where the absorbing iteration has stabilized twice over;
since the two closures then agree on everything the jump law reads,
the laws coincide.  (iii) is the demonic decision problem of
section~\ref{sec:mt-mdp} run through two closures: stored scheduler
rows assemble a scheduler, the rate-closure and jump triangles
collapse the pipeline to a compiled chain, and the chain constraints
pin the law --- four storeys of engine reuse, no new induction; the
mass, slice, and verdict claims are the equal-mass reconciliation
verbatim, the working fuel being exactly where the branch certificate
is spent on the compiled side.  (iv) the minimal-rank instantaneous
state's branch row is a distribution, so it is supported somewhere,
and every instantaneous destination would have to descend below the
minimum.

In the universe's terms: the demonic reading of the compiled automaton
rows and the union-of-singletons inclusion of $D$ are isomorphic by an
invertible globular cell at the per-element stratum, and the credal
presentation reconciles with $D$ through the radially maximal slice,
exactly as at the decision former --- the presentation radial, not
convex, for the same reason.  The statement and both directions are
mechanized, sorry-free, in the toolkit's Lean tree
(\texttt{lean/su-functor/}, module
\texttt{SuFunctor.Adequacy.MaDemonic}): the invertible cell is
\texttt{ma\_demonic\_adequacy}, with admission predicate
\texttt{maDemonicC} and feeding lemma \texttt{maDemonicC\_iff}; the
one-rank-serves-every-adversary certificate is
\texttt{FinMA.BranchZenoFree} with \texttt{zenoFree\_of\_branch}; the
canonicity receipts are \texttt{FinMA.pathLaw\_rank\_irrel} and
\texttt{FinMA.demonicLaws\_rank\_irrel}; the slice reconciliation is
\texttt{ma\_demonic\_maxSlice}; the uniform-over-scheduler verdict
theorems are \texttt{satCredal\_lower\_maDemonicCredal} and
\texttt{satCredal\_upper\_maDemonicCredal}; and the degeneracies are
mechanized rather than smoothed --- the scheduler space has two
nonemptiness routes (with no instantaneous states, the closed-chain
boundary, \texttt{Policy.nonempty\_of\_isEmpty}; with them,
schedulers vanish exactly when actions do,
\texttt{Policy.isEmpty\_iff}), the whole-set reading of a lower bound
refuses every positive threshold
(\texttt{not\_satCredal\_upper\_maDemonicCredal\_geC}), the empty
scheduler space validates both verdict kinds vacuously, and the
Markovian floor of task (iv) is \texttt{FinMA.nonempty\_markovian}.
The general object for automata where only \emph{some} schedulers
admit a rank --- the union over Zeno-certified scheduler resolutions,
each carrying its own rank (\texttt{FinMA.CertScheduler},
\texttt{FinMA.demonicLawsCert}) --- is built at the classical level
and collapses to $D$ under a branch-level certificate
(\texttt{FinMA.demonicLawsCert\_iff}, spending rank-irrelevance); a
compiled square for that partial-Zeno stratum would carry rank
witnesses inside the admission predicate, and is deliberately
separate.  Also separate, unchanged from the fixed-scheduler face:
the time-law stratum; history-dependent schedulers; parallel
composition; the angelic (existential-over-scheduler) reading, a
different operator from the universally quantified verdict pair.

\paragraph{Status.} Certified at the fixed-scheduler grade, with the
scheduler \emph{declared as part of the claim}: a two-station polling
automaton with an unreliable server --- the classic Markov-automata
polling benchmark \cite{guck14}, shrunk to capacity one --- closed
under a stated preference scheduler to a ten-state chain, carries $P(\text{job lost by } T \mid
\text{empty}) \le T/5$ through the unmodified chain machinery. The
mutant battery has teeth in the direction that matters: tampering
with the declared scheduler makes the same certificate refuse with
direct witnesses --- the declaration does real work in the claim. Representation and production: the toolkit's producer
surface constructs the composed profile in both languages and
derives the entire closure bridge exactly from the stored rows ---
the sort partition, the Zeno rank certificate (with the
branch-level rank that serves every adversary, where one exists ---
the demonic face's object --- and absent, reported as such,
otherwise), the
absorption kernel at the working fuel, the closed chain as ordinary
continuous-time rows with the self-loop mass thinned by the grain
itself, and the closed generator by handing the closure to the
chain kit: the square's engine reuse, run producer-side, with the
race-then-jump-then-absorb spelling re-derived as an exact
differential in the test suite. The relay mechanism is a committed
cross-language byte-identity vector, and the API block below quotes
its receipted construction. Scheduler quantification is mechanized at the semantics
layer (the demonic face above); as for the decision frame, the
uniform-over-schedulers story is not yet exercised by a published
reproduction target --- in workshop fixtures with checker acceptance,
the bound above survives quantification verbatim, uniform over an
adapted scheduler class, through the same machinery as
section~\ref{sec:mt-mdp}'s open leg (one design discharges both), and
this entry will say more only when that leg ships --- and the
certificate campaigns remain fixed-scheduler. Open: that reproduction
leg; parallel composition of automata, the frame's own design point;
logic closure.

\syscatmeta{\begin{mathpar}
\inferrule*[Right=\rulename{S-MA}]{
  R \text{ off-diagonal rate entries, as for } \mathrm{CTMC} \\
  O \text{ action-keyed outcome rows, as for } \mathrm{MDP} \\
  \varsigma : \overrightarrow{\langle \mathit{state},\,
  \mathsf{action},\, \mathit{prob} \rangle} \text{ scheduler rows} \\
  \iota? \text{ initial-distribution rows}
}{
  \mathrm{MA}(R;\, O;\, \varsigma) : \mathsf{System}
}
\end{mathpar}
Both parents' row disciplines apply verbatim: no diagonal rate
entry, every outcome row action-keyed (immediate transitions use
the decision idiom). The scheduler is declared data in this version;
quantification over schedulers is mechanized at the semantics layer
(the demonic face's verdict pair), with the published reproduction
target the named open leg.}
\ruleblock{API}
\sourcenote{Producer surface of the \texttt{mtk} toolkit for this model type: the TypeScript arrangement (npm package \texttt{@maitria/mtk}) and the Rust catalogue crate (\texttt{crates/mtk-catalogue}) --- one Rust implementation underneath, reached from TypeScript in-process through a native binding, so the two surfaces cannot diverge. Signatures and constructions are extracted from the toolkit sources at every regeneration of this catalogue; each construction is quoted verbatim from the toolkit's conformance corpus, code whose emitted certificate bytes are verified byte-identical against the committed golden vectors (helper values named \texttt{fx\_*} or \texttt{f.*} are that corpus's shared fixtures, defined beside the quoted lines). This API constructs and emits; it never checks --- the checker re-derives every claim independently.}
\apilang{TypeScript}{system.markovAutomaton}
\apiprov{signature --- \texttt{ts/src/system.ts}}
\begin{tssrc}
/** S-MA: Markov automaton — rates + action-keyed immediate rows + a
 * scheduler declaration (declared data in v1). */
export const markovAutomaton = profileFactory("markovAutomaton");

/** The derived Markov-automaton fixed-scheduler kit: the sort
 * partition, the Zeno rank certificate (mixed-kernel face, plus the
 * scheduler-free branch-level rank when one exists — the demonic
 * face's one-certificate-for-all object), the absorption kernel at
 * the working fuel, the closed continuous-time chain as ordinary
 * CTMC rows, and the closed chain's generator kit — the mechanized
 * square's engine reuse, run producer-side. Derived, never trusted —
 * the checker re-runs the closure's three storeys. */
export interface MaGenKit {
  readonly markovian: readonly number[];
  readonly instantaneous: readonly number[];
  readonly rank: readonly { readonly state: number; readonly rank: number }[];
  readonly branchRank:
    | readonly { readonly state: number; readonly rank: number }[]
    | null;
  readonly absorb: readonly {
    readonly from: number;
    readonly to: number;
    readonly mass: string;
  }[];
  readonly closedRates: readonly {
    readonly src: number;
    readonly dst: number;
    readonly rate: string;
  }[];
  readonly closedInit: readonly {
    readonly state: number;
    readonly prob: string;
  }[];
  readonly closed: CtmcGenKit;
}

/** Derive the Markov-automaton fixed-scheduler kit from a mechanism
 * declaration and a per-state certificate observable on the
 * Markovian states. Maximal progress (a state with both timed and
 * immediate rows), scheduler discipline (coverage, mass one, support
 * within the enabled pairs), and the immediate-cycle Zeno door all
 * come back as witness-shaped refusals naming their geometry. */
export function maGenKit(
  init: MechanismInit,
  cert: readonly StateVal[],
): Result<MaGenKit, Refusal> {
  const raw = native.sysMaGenKit(mechanism(init).nativeHandle, cert) as
    | { kit: MaGenKit }
    | { refusal: unknown };
  return "refusal" in raw
    ? refused(narrowRefusal<Refusal>(raw.refusal))
    : ok(raw.kit);
}
\end{tssrc}
\apiprov{construction --- \texttt{ts/tests/magen-native.test.ts}}
\begin{tssrc}
/** The `relay-2s1i` presentation (the committed golden fixture's):
 * Markovian 0 and 1, instantaneous relay 2; state 0 races at rate 2
 * into the relay, which returns a third of the mass under the
 * declared scheduler. */
function relayInit() {
  return {
    rates: [
      { src: 0, dst: 2, rate: "2" },
      { src: 1, dst: 0, rate: "1" },
    ],
    outcomes: [
      { src: 2, action: 0, dst: 0, prob: "1" },
      { src: 2, action: 1, dst: 1, prob: "1" },
    ],
    scheduler: {
      rows: [
        { state: 2, action: 0, prob: "1/3" },
        { state: 2, action: 1, prob: "2/3" },
      ],
    },
    initDist: [{ state: 0, prob: "1" }],
  } as const;
}
\end{tssrc}
\apiprov{The quoted construction is the committed \texttt{relay-2s1i} cross-language byte-identity vector's presentation, whose thinned closure the kit tests pin in closed form on both surfaces; the kit's Rust face is \texttt{mtk\_catalogue::system::ma\_gen\_kit}.}
\apilang{Rust}{mtk\_catalogue::system::presets::MarkovAutomaton::manifest}
\apiprov{signature --- \texttt{crates/mtk-catalogue/src/system/presets.rs}}
\begin{rustsrc}
preset!(
    /// System Catalogue: Markov automaton profile — the CTMC rate
    /// idiom + the MDP kernel idiom aimed at immediate transitions +
    /// a scheduler declaration (declared data in v1).
    MarkovAutomaton, "Markov automaton",
    req = [Rates, Outcomes, Scheduler],
    allowed = [InitDist],
    extra = |m| {
        no_diagonal_rates(m)?;
        outcomes_keyed("Markov automaton", m, true)
    }
);
\end{rustsrc}
\apiprov{construction --- \texttt{crates/mtk-catalogue/tests/builders.rs}}
\begin{rustsrc}
let ma = MechanismBuilder::new()
    .rate(0, 1, r(1, 1))
    .outcome(OutcomeRow {
        src: 1,
        action: Some(0),
        dst: 0,
        prob: r(1, 1),
        reward: None,
    })
    .scheduler(SchedulerDecl {
        rows: vec![SchedulerRow {
            state: 1,
            action: 0,
            prob: r(1, 1),
        }],
        note: None,
    })
    .finish();
MarkovAutomaton::manifest(ma).unwrap();
\end{rustsrc}

\section{Generalized Stochastic Hybrid Systems}
\label{sec:mt-gshs}

\paragraph{The frame.} Modes with continuous stochastic dynamics,
guarded and forced transitions with resets, spontaneous jumps ---
the frame of Bujorianu and Lygeros \cite{bujorianu08}, and the frame
the System notion was designed around, because a modelling
language that cannot host it cannot host the systems this toolkit
exists for.

\paragraph{Profile.} The full manifest complement: drift, diffusion,
jump fields, plus guarded transitions (guard slabs, resets, forced
marks) across modes. Boundary behaviour is first-class and the
uniformity break is recorded rather than smoothed: forced transitions
are \emph{domain conditions}, not generator terms. Time: hybrid ---
jumps are events in real time.

\paragraph{Semantics.} Extended-generator semantics on segments,
with the Dynkin argument staged segmentwise. Switched trajectories
make per-point differentiability semantically \emph{false}, so
decrease obligations for hybrid rows discharge at an upper Dini
derivative --- corner- and Zeno-immune, with mechanized rungs for
switched flows, an integrated almost-everywhere stratum, and Filippov
sliding reduced to per-mode vertex rows, at no new row shape in the
checker. Hybrid records are native to the universe's record
structure --- resets are zero-duration stage pairs
(Definition~\ref{def:su-record}) --- which is why segmentwise
semantics needs no second time axis.

\paragraph{Status.} Certification: the frame's namesake example is
shipped --- a three-mode kiln with three forced guard transitions and
resets (one nontrivial and load-bearing: the identity-reset probe
refuses), spontaneous spikes, and diffusion in both live modes. Nine
sequents, every one an ordinary polynomial inequality over a region
through the unmodified stack, yield $P(\text{overheat}) \le
671/2048$ from the hot start; the battery includes a
diffusion-ablation witness showing the second-order term alone
carries the hot mode, and the non-Zeno bridge is discharged by
inspection of an acyclic transition structure. Representation: the
first instance rows to exercise transitions, resets, and forced
marks at all. The composite SDCPN formalism of Everdij and Blom ---
the 2023 candidate --- now carries its own entry
(section~\ref{sec:mt-sdcpn}): places are Bags of coloured-token
Systems, and the per-place colour dynamics are this profile's. Open: a certified
example with guards and resets at $n \ge 2$ (the kiln is
one-dimensional and the encounter systems are single-mode);
stochastic reset kernels (deterministic polynomial resets are
exercised; distributions enter as finitely many reset legs); curved
or interior guard geometry beyond boundary slabs; the ingestion gap;
closure.

\syscatmeta{\begin{mathpar}
\inferrule*[Right=\rulename{S-GSHS}]{
  \vec{x} \text{ declared state variables} \\
  M : \overrightarrow{\langle \mathsf{id},\, \mathsf{name}?,\,
  \mathsf{domain} \rangle} \text{ modes with domain constraints} \\
  T : \overrightarrow{\langle \mathsf{from\_mode},\,
  \mathsf{to\_mode},\, \mathsf{guard},\, \mathsf{reset},\,
  \mathsf{forced} \rangle} \text{ guarded transitions} \\
  D?,\, N?,\, J? \text{ mode-indexed drift, noise-presentation, and
  jump families}
}{
  \mathrm{GSHS}(\vec{x};\, M;\, T) : \mathsf{System}
}
\end{mathpar}
Forced transitions are domain conditions rather than generator terms
--- the uniformity break, recorded as structure.}
\ruleblock{API}
\sourcenote{Producer surface of the \texttt{mtk} toolkit for this model type: the TypeScript arrangement (npm package \texttt{@maitria/mtk}) and the Rust catalogue crate (\texttt{crates/mtk-catalogue}) --- one Rust implementation underneath, reached from TypeScript in-process through a native binding, so the two surfaces cannot diverge. Signatures and constructions are extracted from the toolkit sources at every regeneration of this catalogue; each construction is quoted verbatim from the toolkit's conformance corpus, code whose emitted certificate bytes are verified byte-identical against the committed golden vectors (helper values named \texttt{fx\_*} or \texttt{f.*} are that corpus's shared fixtures, defined beside the quoted lines). This API constructs and emits; it never checks --- the checker re-derives every claim independently.}
\apilang{TypeScript}{system.gshs}
\apiprov{signature --- \texttt{ts/src/system.ts}}
\begin{tssrc}
/** S-GSHS: modes with domains + guarded transitions (+ the full
 * continuous complement). */
export const gshs = profileFactory("gshs");
\end{tssrc}
\apiprov{construction --- \texttt{ts/tests/system-native.test.ts}}
\begin{tssrc}
return sys.unwrap(
  sys.gshs({
    name: "kiln-2mode",
    state: [{ id: 0, name: "temp" }],
    modes: [
      {
        id: 0,
        name: "heat",
        domain: [{ poly: mono(1, 0, 1, "1"), op: "le", rhs: "100" }],
      },
      {
        id: 1,
        name: "cool",
        domain: [{ poly: mono(1, 0, 1, "1"), op: "ge", rhs: "20" }],
      },
    ],
    drift: [
      { mode: 0, var: 0, poly: constant(1, "3") },
      { mode: 1, var: 0, poly: mono(1, 0, 1, "-1/10") },
    ],
    diffusion: [
      { mode: 0, var: 0, poly: constant(1, "1/4") },
      { mode: 1, var: 0, poly: constant(1, "1/4") },
    ],
    jumps: [
      {
        mode: 0,
        rate: constant(1, "1/100"),
        reset: [{ var: 0, poly: mono(1, 0, 1, "9/10") }],
        resetImage: { axes: [["18", "90"]] },
      },
    ],
    transitions: [
      {
        fromMode: 0,
        toMode: 1,
        guard: [{ poly: mono(1, 0, 1, "1"), op: "ge", rhs: "100" }],
        forced: true,
      },
    ],
  }),
);
\end{tssrc}
\apilang{Rust}{mtk\_catalogue::system::presets::Gshs::manifest}
\apiprov{signature --- \texttt{crates/mtk-catalogue/src/system/presets.rs}}
\begin{rustsrc}
preset!(
    /// System Catalogue: GSHS profile — the full complement (drift,
    /// diffusion/noise, jumps) + modes with domains and guarded
    /// transitions (guard slabs, resets, forced marks).
    Gshs, "GSHS",
    req = [Modes, Transitions],
    allowed = [Drift, Diffusion, Noise, Jumps],
    extra = ok
);
\end{rustsrc}
\apiprov{construction --- \texttt{crates/mtk-catalogue/tests/common/mod.rs}}
\begin{rustsrc}
let below = DomainConstraint {
    poly: mono(1, 0, 1, r(1, 1)),
    op: "le".to_owned(),
    rhs: r(100, 1),
};
let above = DomainConstraint {
    poly: mono(1, 0, 1, r(1, 1)),
    op: "ge".to_owned(),
    rhs: r(20, 1),
};
let m = MechanismBuilder::new()
    .state_var(0, Some("temp"))
    .mode(Mode {
        id: 0,
        name: Some("heat".to_owned()),
        domain: vec![below],
    })
    .mode(Mode {
        id: 1,
        name: Some("cool".to_owned()),
        domain: vec![above],
    })
    .drift(Some(0), 0, constant(1, r(3, 1)))
    .drift(Some(1), 0, mono(1, 0, 1, r(-1, 10)))
    .diffusion(Some(0), 0, constant(1, r(1, 4)))
    .diffusion(Some(1), 0, constant(1, r(1, 4)))
    .jump(JumpRow {
        mode: Some(0),
        rate: constant(1, r(1, 100)),
        reset: vec![ResetRow {
            var: 0,
            poly: mono(1, 0, 1, r(9, 10)),
        }],
        mixture: vec![],
        reset_image: Some(BoxRegion {
            axes: vec![[r(18, 1), r(90, 1)]],
        }),
    })
    .transition(GuardedTransition {
        from_mode: 0,
        to_mode: 1,
        guard: vec![DomainConstraint {
            poly: mono(1, 0, 1, r(1, 1)),
            op: "ge".to_owned(),
            rhs: r(100, 1),
        }],
        reset: vec![],
        forced: true,
    })
    .finish();
presets::Gshs::manifest(m)
    .expect("fixture conforms to the GSHS profile")
    .named("kiln-2mode")
\end{rustsrc}

\section{The Bag Former}
\label{sec:mt-bag}

\paragraph{The frame.} Not a frame at all --- a primitive
\emph{System-former}: $\mathrm{Bag}(S)$ is the System of a varying
finite multiset of tokens, each token an instance of the argument
System $S$. It exists as its own entry because structured-population
types (the coloured nets of section~\ref{sec:mt-sdcpn} first among
them) want ``a place holds finitely many tokens with per-token
dynamics'' as one reusable construction rather than a per-frame
encoding trick.

\paragraph{Profile.} A countably-indexed switched-mode carrier: mode
$n$ is the $n$-fold juxtaposition power of $S$'s state, quotiented by
exchange --- the canonical-multiset representation at the
presentation stratum, the Polish symmetric-power quotient at the
semantics stratum. Token creation and destruction are the former's
own variant-crossing switch edges (self-loops between adjacent
powers); noise is per-token independent. Monad structure (singleton
unit, flatten multiplication) sits at the admissible tier rather than
in the rows.

\paragraph{Semantics.} The former denotes
\[
\mathrm{Bag}(S) \;=\; \coprod_{n \in \mathbb{N}} S^{\boxplus n} / S_n,
\]
and the display is a colimit formula in the semantic universe, not a
set-theoretic sum; every operation in it is licensed by the standing
structure of Definition~\ref{def:su-credal-denotation}. The power
$S^{\boxplus n}$ is the $n$-fold juxtaposition composite of the
composition calculus, written with the belief laxator's symbol
because its decoration is the $n$-fold $\boxplus$ of $S$'s belief
content (its level-of-inconsistency functional) with itself
(Definition~\ref{def:su-house-functor}); its
state space is the $n$-fold product $X^n$ of the inner System's
state space $X$ (Definition~\ref{def:su-state-spaces}), and the
empty power is the unit System on the one-point space. The symmetric
group acts through the juxtaposition tensor's symmetries ---
permutations of the $n$ interface copies, injective tight
isomorphisms of the wiring base, acting on states by coordinate
permutation --- and $X^n/S_n$ is the coequalizer of those $n!$
homeomorphisms in the category of Polish spaces and continuous maps,
Polish under the matching metric; the disjoint union over $n$ is
the countable coproduct in the same category, its mass sectors clopen
(both licences: the bag widget of
Definition~\ref{def:su-credal-denotation}). Equivalently the carrier
is the space of integer-valued finite counting measures on $X$ under
the weak topology --- portmanteau on small balls around the atoms
identifies each mass sector with the corresponding symmetric power,
and the matching metric on the mass-$n$ sector is the
optimal-matching ($\infty$-transport) distance between empirical
measures --- the presentation through which the creation and
destruction rate species acquire their point-process reading. So
$\mathrm{Bag}$ preserves Polishness of the carrier, and bag-valued
places stay inside the record-space semantics the rest of the
catalogue assumes, with token creation and destruction as
zero-duration variant crossings, exactly as hybrid records carry
resets (Definition~\ref{def:su-record}).

The imprecise denotation passes through the quotient rather than
around it. Bounded continuous observables on $X^n/S_n$ are exactly
the $S_n$-invariant ones on $X^n$ (the quotient topology's universal
property), so the weak topology over the quotient is generated by
symmetrized test functions, and pushforward along the quotient map
is an affine homeomorphism from the exchangeable
(permutation-invariant) subprobabilities on $X^n$ onto
$\Delta_{\le 1}(X^n/S_n)$ --- injective by symmetrizing test
functions, surjective by symmetrizing the lift along the
canonical-multiset section. A credal body on the bag carrier
is therefore the same thing as an exchangeable credal body on
the ordered representation --- exchangeability is well-definedness,
not a preference --- and the profile's per-token-uniform dynamics
and beliefs descend because they are equivariant by construction;
invariant certificate observables \emph{are} the quotient's
observables, which is what the canonicalization discipline of the
Status leg realizes at the presentation stratum. Two boundaries,
kept sharp: bags are finite --- infinite multisets break both
presentations, and token counts stay finite along every record by
construction --- and the colimits here are colimits of
\emph{carriers}: the formula lifts to Systems through the former,
whose monad structure (singleton unit, flatten multiplication) is
stated on the fibre categories of the composition construction at the
admissible tier, not as a colimit in them.

\paragraph{Status.} Representation: shipped --- the declaration
wraps an inner System, an optional capacity, typed
creation/destruction edges, and a stored marking under the
canonical-multiset discipline (tokens sorted in the exact
coordinatewise-lexicographic order, duplicates adjacent, the count
derived as the list's length; the conformance gates refuse a
non-canonical stored value, so preset-accepted declarations carry
canonical bytes). Semantics and certification: the rate species of
the creation/destruction edges is landed as a typed vocabulary --- a
creation edge reads aggregate observables (the token count), a
destruction edge races per unordered token selection with
slot-exchange symmetry checked exactly, and a synced edge declares a
fused capability whose one firing law lives on the
event-port fusion row --- and the first certificates exercise a bag
carrier end to end: a marked bag of ordinary-differential-equation
tokens compiles by the landed view pair into its single-place
zero-arc coloured net (the two constructions are one object, checked
to the byte on the shared locus), projects onto the Petri incidence
presentation, and carries a certified population invariant through
two independent checker lanes, with mutants at the designed loci ---
a non-canonical stored value refuses, an asymmetric selection rate
refuses, and a desynchronized transfer flips both checkers. The
billiard table as $\mathrm{Bag}(\mathbb{R}^4)$ remains the worked
example on the road for the colour stratum.

\paragraph{Adequacy.} The square's measure-theoretic cell --- the
ordered representation and the collective marking process present one
process --- is mechanized, sorry-free, in the toolkit's Lean tree
(\texttt{lean/su-functor/}, module \texttt{SuFunctor.Adequacy.Bag}).
The sorted-representative section of the exchange quotient is
measurable under the pushforward $\sigma$-algebra from nothing but a
linear order on colours with a measurable order graph: each sorted
coordinate is an order statistic, a finite minimum of maxima of the
unsorted coordinates (\texttt{sortTuple\_apply\_eq\_inf\_sup},
\texttt{measurable\_symSec}), so no measurable-selection theorem
enters, and ties on the diagonal never demand a choice --- monotone
presentations of one multiset are unique
(\texttt{monotone\_eq\_of\_toSym\_eq}). On that section, the quotient
kernel of an exchange-invariant chain closes the commuting square at
every finite horizon by kernel-category algebra alone
(\texttt{quotKernel\_comp\_eq}, \texttt{iter\_intertwine}); the
construction is independent of the section chosen
(\texttt{comap\_section\_congr}); and the transport assembles on the
full shape-indexed carrier, token creation and destruction included,
with per-permutation exchange-invariance and invariance under
canonicalization interderivable (\texttt{bag\_transport},
\texttt{invariance\_iff\_perm}). The order-graph hypothesis is
discharged for the coordinatewise-lexicographic order on real colour
vectors that the stored representation sorts by
(\texttt{measurableSet\_le\_lex}) --- the tie corner lives exactly
there, and it is an instance, not an assumption. What remains at its
named home: dynamics preservation for the monad's unit and
multiplication (a process-stratum statement), and the per-shape
generator agreement that belongs to the coloured-net stratum of the
next section.

\syscatmeta{\begin{mathpar}
\inferrule*[Right=\rulename{S-BAG}]{
  S : \mathsf{System} \\
  \mathsf{capacity}? \text{ a declared population cap} \\
  L : \overrightarrow{\langle \mathsf{id},\, \mathsf{kind},\,
  \Delta,\, \mathsf{law} \rangle} \text{ self-loops, } \mathsf{kind}
  \in \{\mathsf{create},\, \mathsf{destroy}\} \\
  M? \text{ a stored marking}
}{
  \mathrm{Bag}(S) : \mathsf{System}
}
\end{mathpar}
$m.\mathsf{bag} = \langle \mathsf{inner},\, \mathsf{capacity}?,\,
\mathsf{self\_loops},\, \mathsf{marking}? \rangle$; each self-loop
declares its id (the fusion-reference key), kind (create/destroy),
per-firing delta, and firing law --- delay with a typed rate basis
(per-selection or aggregate), guard, or synced; a stored marking is
the canonical sorted token list, its count the list's length.}
\ruleblock{API}
\sourcenote{Producer surface of the \texttt{mtk} toolkit for this model type: the TypeScript arrangement (npm package \texttt{@maitria/mtk}) and the Rust catalogue crate (\texttt{crates/mtk-catalogue}) --- one Rust implementation underneath, reached from TypeScript in-process through a native binding, so the two surfaces cannot diverge. Signatures and constructions are extracted from the toolkit sources at every regeneration of this catalogue; each construction is quoted verbatim from the toolkit's conformance corpus, code whose emitted certificate bytes are verified byte-identical against the committed golden vectors (helper values named \texttt{fx\_*} or \texttt{f.*} are that corpus's shared fixtures, defined beside the quoted lines). This API constructs and emits; it never checks --- the checker re-derives every claim independently.}
\apilang{TypeScript}{system.bag}
\apiprov{signature --- \texttt{ts/src/system.ts}}
\begin{tssrc}
/** S-BAG: the Bag former (primitive System-former over an inner
 * System). */
export const bag = profileFactory("bag");
\end{tssrc}
\apiprov{The TS conformance corpus holds no construction through this factory at this preview --- the cross-language byte-identity fixtures exercise six profiles, and this one's byte authority is the Rust corpus opposite; the factory takes the same init-object vocabulary as every other profile.}
\apilang{Rust}{mtk\_catalogue::system::presets::Bag::manifest}
\apiprov{signature --- \texttt{crates/mtk-catalogue/src/system/presets.rs}}
\begin{rustsrc}
preset!(
    /// System Catalogue: Bag former — primitive System-former (the
    /// canonical-multiset carrier over the argument System), not a
    /// frame. The extra gates enforce the former's semantics-bearing
    /// discipline: typed edges (unique ids; delta ≥ 1; create reads
    /// aggregates, destroy races per selection, guards select the
    /// hitting token; W1 slot-exchange symmetry on per-selection
    /// rates) and canonical stored markings (arity, canonical
    /// rationals, the unique sorted multiset representative, the
    /// capacity cap).
    Bag, "Bag",
    req = [Bag],
    allowed = [],
    extra = |m| {
        let b = m.bag.as_ref().expect("required by signature");
        let d = inner_state_arity(&b.inner, "bag.inner")?;
        bag_edge_gates(&b.self_loops, d)?;
        if let Some(tokens) = &b.marking {
            bag_marking_gates(tokens, d, b.capacity.as_ref(), "bag.marking")?;
        }
        Ok(())
    }
);
\end{rustsrc}
\apiprov{construction --- \texttt{crates/mtk-catalogue/tests/builders.rs}}
\begin{rustsrc}
let bag = MechanismBuilder::new()
    .bag(BagDecl {
        inner: Box::new(common::sde()),
        capacity: None,
        self_loops: vec![
            // Arrivals: aggregate basis (var 0 = the token count).
            BagEdge {
                id: 0,
                kind: BagEdgeKind::Create,
                delta: WitInt::from(1),
                law: BagEdgeLaw::Delay {
                    rate: BagRate::Aggregate {
                        rate: constant(1, r(1, 2)),
                    },
                },
            },
            // Departures: per-selection basis (the selected
            // token's colour coordinates).
            BagEdge {
                id: 1,
                kind: BagEdgeKind::Destroy,
                delta: WitInt::from(1),
                law: BagEdgeLaw::Delay {
                    rate: BagRate::PerSelection {
                        rate: mono(1, 0, 2, r(1, 10)),
                    },
                },
            },
        ],
        // Stored multiset value, canonical order (count = length).
        marking: Some(vec![vec![r(-1, 2)], vec![r(0, 1)], vec![r(0, 1)]]),
    })
    .finish();
Bag::manifest(bag).unwrap();
\end{rustsrc}

\section{Stochastically and Dynamically Coloured Petri Nets}
\label{sec:mt-sdcpn}

\paragraph{The frame.} The SDCPN formalism of Everdij and Blom ---
the candidate unifier of the 2023 grading --- as a first-class entry:
Petri-net structure whose tokens carry continuous stochastic colour
dynamics. The reading this catalogue commits to is compositional
rather than monolithic: \emph{places are Bags, arcs are patterns}.
The stance has good company in the categorical net literature ---
nets opened along place boundaries and composed as cospans
\cite{baez18}, coloured nets read as monoidal double functors
\cite{master25}, and individuals-based readings of the token game
\cite{kock20} --- though the construction here factors through the
catalogue's own composition patterns rather than through a dedicated
category of nets.

\paragraph{Profile.} Places are $\mathrm{Bag}$s of a coloured-token
System (section~\ref{sec:mt-bag}), each with an optional declared
capacity; transitions are per-tuple edge families in three species
--- delay (a rate law), guard (forced), immediate (cascade glue);
ordinary arcs are variant-crossing consume/produce resets with
multiplicities, and enabling arcs are per-tuple read-outs. Every arc
kind is an instance of an existing composition pattern --- with one
deliberate exclusion: \textbf{inhibitor arcs are not a primitive}.
Absence-testing breaks enabledness monotonicity, the free
commutative-monoid factoring, and the negation-free character of the
chase all at once --- a unique upper-bound gate in a lower-bound
formalism. What survives is the bounded-locus converse: on a place
with declared capacity $k$, an inhibitor arc is admissible as a
\emph{view} --- the complement-place compile, whose elaborator is a
named successor: a complement place with mirrored arcs and a
weight-$k$ enabling arc under a certified linear place-invariant
$m(P) + m(P') = k$ --- and the compile provably cannot extend to an
unbounded place, which is why no unbounded inhibitor constructor
exists at all.

The vocabulary, enumerated. Places being Bag mechanisms, each arc is
a composition shape wiring Bag populations: a transition compiles to
a \emph{rendezvous} --- one joint jump in the fusion discipline of
section~\ref{sec:cp-eventport}, timed by its active party, the
transition species --- whose other parties are the incident places'
bag edges, themselves the mode-switching machinery of
section~\ref{sec:cp-switch}. Of the eleven elements of the founding
formalism --- the arc kinds, the transition species, and the
net-level data --- ten map onto the pattern inventory, nothing is
added to it, and the eleventh is the inhibitor exclusion above. The
tenses matter: the per-Bag edge vocabulary, the read-out wires, and
the fusion rows themselves are shipped declaration rows --- the
fusion primitive now carries a producer surface, and the
coloured-net decomposition emits its rows
(section~\ref{sec:cp-eventport}) --- while the pattern vocabulary
downstream of the primitive still speaks in design tense there, so
the compiled artefacts this table calls for are graded in the
Status leg below.

\begin{center}\small
\begin{tabular}{@{}l>{\raggedright\arraybackslash}p{5.0cm}>{\raggedright\arraybackslash}p{4.7cm}@{}}
\toprule
element & semantics & compilation target \\
\midrule
ordinary arc, input face & one token read into the firing data and
removed: variant-crossing reset $n \to n-1$ at the source place & a
rendezvous party: the source Bag's destroy edge, its reset carried
by the joint jump \\
ordinary arc, output face & a token produced at the target place,
$n \to n+1$, its colour drawn from the firing measure & a rendezvous
party: the target Bag's create edge, spawn direction \\
enabling arc & the selected token's colour read by rate, guard, or
firing measure; no removal & a passive party: a per-tuple
parameter-setting wire (section~\ref{sec:cp-set}) into the timing
owner \\
enabledness count & pre-enabling threshold: one token per ordinary
and enabling input arc & scope rows on the rendezvous's edge
family --- a monotone, negation-free count guard \\
inhibitor arc & pre-enabling gate on absence & not compiled ---
excluded as a primitive; the bounded-locus view above enters
through the complement place's rendezvous \\
delay transition & rate-driven firing: a per-tuple survivor law
raced across enabled tuples & the active party: a rate law timing
the joint jump \\
guard transition & forced firing on the guard boundary & the active
party, forced species \\
immediate transition & fires as soon as enabled; vanishing markings
leave the state space & the active party, immediate species:
cascade glue under the no-Zeno side condition \\
conflict resolution & priority and randomization among competing
enabled tuples & certificate-stratum dominance; no composition
pattern consumed \\
firing measure & the joint law of the output vector and the produced
colours & reset-leg moment data on the produce edges \\
initial marking & the initial distribution over variants and
canonical colours & root init rows \\
\bottomrule
\end{tabular}
\end{center}

\paragraph{Semantics.} Marking spaces are Polish by construction:
finitely many places, each a Bag carrier --- a finite product of the
previous section's colimits, inside the standing structure of
Definition~\ref{def:su-credal-denotation}. The correspondence with
section~\ref{sec:mt-gshs} (SDCPN executions as generalized
stochastic hybrid processes, the founding papers' own bridge) enters
as a translation between \emph{presentations} --- never as the
definition of the net's semantics, whose compiled reading remains
structural in the net while the flattened object is combinatorially
large --- and its jump-skeleton stratum is mechanized at the
uncoloured grain \emph{and} at the coloured deterministic-firing
slice: definedness-premise rows, the commuting-triangle cell, the
constructed flattening with its agreement derived, and the composed
community-to-community corollary included at both grains
(the Adequacy leg below), with the remaining coloured obligations ---
the hybrid-grain carrier identification and the firing-measure
kernels --- enumerated by name. The compilation
direction is Bag-primary: arcs compile into composition patterns ---
rendezvous of creation/destruction capabilities --- over the places'
Bag mechanisms, the zero-arc case being the identity.

\paragraph{Adequacy (statics, the delay-race face, the
immediate/cascade face at all three scheduler readings, and the
flattening bridge).} The
compositional reading --- \emph{places are Bags, arcs are patterns}
--- is not a slogan: the adequacy square for this frame factors
through the composition patterns' own mechanized semantics, and it
closes at two strata as of this preview, with the rest staged by
name.  At the statics stratum the square is whole.  Stated
self-contained:

\begin{problem}[the token-game statics square]
\label{prob:mt-sdcpn}
A marking assigns each place a finite multiset of
token colours.  A transition of the inhibitor-free fragment reserves,
per place, one token for each ordinary and each enabling input arc
(distinct tokens), removes the ordinary ones, keeps the enabling ones,
adds fresh tokens per output arc, and feeds every reserved colour into
one firing datum along its arc's readout \cite{everdij-blom-2010}.
Compile instead: each place a bag composite (one count wire, a token
power per mode, states read modulo exchange), consumption a
token-split rename, each colour readout an identification-pattern
row.  Show the two readings agree: markings of each shape are exactly
the compiled states modulo per-place exchange renames, and the token
game holds between two markings if and only if compiled states
realizing them are connected by split renames and pattern apex states.
Show also that the fragment's token game is upward closed --- adding
tokens never disables a step --- and exhibit the gate that would
break it.
\end{problem}

\emph{Sketch.} The multiset readout of a compiled state is surjective
onto markings of its shape and a complete exchange invariant --- two
tuples enumerate one multiset exactly when a permutation relates them
--- so the community's configuration space is the compiled state
space modulo exchange, shape by shape.  For the firing
correspondence, forward: enumerate the reserved and untouched tokens
through a chosen index split, and the required states are defined
directly through the split; backward: the split decomposes the bag as
selected-block plus kept-block (the many-token form of the cons
identity), and the pattern apex condition is precisely colour-readout
agreement.  Upward closure is one line per field --- the surplus
sits in the untouched remainder --- and the gate that breaks it is
the
zero-test: an inhibitor arc is the unique arc kind whose pre-condition
is not upward closed, which is why the fragment's transition type has
no constructor for it.  The exclusion is a formation restriction with
the monotonicity theorem as its warrant, not a lint.

At the finite time-abstract stratum the delay species closes as a
square.  The community execution races every pre-enabled
\emph{ordered} assignment of tokens to input arcs on its own
exponential clock \cite{everdij-blom-2010}, so from a bounded
uncoloured marking the total intensity to a target is the ordered
census --- per place, a falling factorial of supply over demand ---
times the rate, summed over transitions; the compiled rows carry the
\emph{unordered} convention, one row per token set, with the
arc-multiplicity dictionary $\prod_p (\text{demand}_p)!$ owned by the
compiler.  The compile-commutes triangle is therefore not
definitional, deliberately: it \emph{is} the dictionary theorem ---
falling factorial equals factorial times binomial, per place --- and
through it the square is the continuous-time jump-chain square
instantiated at the race chain, engine reuse with zero new
inductions; the initial marking enters through the same square's
init rows,
and firings whose outcome equals their source are thinned by the
off-diagonal row grain, the standard jump-chain convention stated as
a receipt.  Mechanized end to end (\texttt{lean/su-functor/}, modules
\texttt{SuFunctor.TokenGame} and
\texttt{SuFunctor.Adequacy.StochasticPetri}): the vertical
identification, the firing correspondence, the monotonicity warrant,
the census-to-embeddings tie, the dictionary triangle, the reused
square, and the thinning receipt.

The immediate species closes as a square at the same stratum, by the
Markov-automaton closure consumed as built.  A marking that
pre-enables an immediate transition is \emph{vanishing} --- the
community excludes it from the state space of the observed process
\cite{everdij-blom-2010} --- and here the exclusion is carried by the
type: markings split into tangible and vanishing subtypes, the net
compiles to a two-sorted Markov automaton whose Markovian sort is the
tangible markings, and a delay row at a vanishing marking is thereby
unrepresentable --- the maximal-progress priority as the domain type
of the rate field.  Immediate firings are the automaton's action-keyed
branches (point masses at this uncoloured stratum); a declared
memoryless scheduler resolves conflicts, so the square is
scheduler-generic and dominates the community's randomized choice
rule, with the uniform-over-enabled scheduler as the shipped instance;
cascade termination is a rank certificate --- a rank on vanishing
markings that every enabled immediate firing strictly decreases ---
and the fixed-scheduler closure delivers the reduced tangible chain,
with vanishing initial mass pushed through the cascade rather than
excluded by premise (when the community's initial-tangibility premise
does hold, the initial distribution passes through untouched, as a
receipt).  Timed rows compile in the same unordered convention as the
delay face, so the compile triangle is again the dictionary theorem,
one storey up, and the square is the automaton square at the compiled
net --- five-storey engine reuse, zero new inductions.  Mechanized in
\texttt{SuFunctor.Adequacy.StochasticPetriCascade}: the classified
state space, the automaton reading, the enabling-respect discipline
and rank certificate, the tangible chain with the initial-tangibility
receipt, the compile triangle, the reused square, and controls
pinning the uniform conflict split, the thinning of returned cascade
mass, and a genuine two-step cascade.

The scheduler-quantified sibling closes beside it --- the demonic
face of the cascade square, and the point where the automaton
donor's certificate does \emph{not} transport.  The branch-level
Zeno certificate that serves every adversary at the automaton
(\texttt{FinMA.BranchZenoFree}, section~\ref{sec:mt-ma}) is
refutable at nets: a \emph{disabled} immediate transition at a
vanishing marking compiles to the zero-time no-op it is --- a
self-loop branch row that no rank can strictly decrease
(\texttt{not\_branchZenoFree\_of\_disabled}, discharged against the
species' shipped control net for every rank at once) --- so the
demon is instead quantified over the subclass the cascade rank
actually certifies: \emph{enabling-respecting} schedulers
(\texttt{RespScheduler}), the class the community's conflict
randomizations inhabit, with the uniform scheduler as its shipped
member, a discard scheduler witnessing that the arena is genuinely
plural (\texttt{arena\_plural}), and the compiled admission
strengthened to carry the side condition in the statement
(\texttt{RespectingRows}) rather than smuggling it through a
definitional coincidence.  Over that arena the set-valued square
closes, as at the decision and automaton frames
(sections~\ref{sec:mt-mdp} and~\ref{sec:mt-ma}): the demonic
reading of the compiled, scheduler-free net rows is isomorphic, by
an invertible globular cell, to the inclusion of the net's demonic
tangible law-set (\texttt{ispn\_demonic\_adequacy}, feeding lemma
\texttt{ispnDemonicC\_iff}), the law-set and its credal presentation
are independent of which rank certificate is spent
(\texttt{demonicTangibleLaws\_rank\_irrel} and
\texttt{demonicTangibleCredal\_rank\_irrel}), the equal-mass slice
reconciliation and the uniform-over-adversary verdict pair transport
(\texttt{ispn\_demonic\_maxSlice},
\texttt{satCredal\_lower\_ispnDemonicCredal} and its upper twin),
and the wiring into the semantic universe's judgment lands as the
module's own quartet (\texttt{su\_ispn\_demonic\_adequacy} with
its mem, bottom, and slice forms).  Six storeys of engine reuse,
zero new inductions.  Mechanized, sorry-free, beside the cascade
module (\texttt{lean/su-functor/}, module
\texttt{SuFunctor.Adequacy.StochasticPetriDemonic}); history-dependent
adversaries stay outside the arena by definition, exactly as at the
automaton.

The set-firing face closes beside them, at the marking stratum.  The
community fires \emph{non-conflicting simultaneous enablings}
together --- one transition enabled by several token sets, or several
transitions enabled at once, no firing disabling the rest
\cite{everdij-blom-2010} --- and the face is driven by one value
identity rather than any combinatorial induction: a zero-time chain
of single firings whose every step is enabled is \emph{exact} (no
truncation fires), so per place its result satisfies ``result plus
ordinary-input total equals source plus ordinary-output total,'' and
the right-hand data depend on the firing multiset alone
(\texttt{chainFire\_val}).  Order-invariance is then value
uniqueness --- two enabled enumerations of one multiset land the same
marking, with no diamond or permutation induction anywhere
(\texttt{chainFire\_eq\_of\_enabled}; the pairwise diamond is the
two-element corollary) --- and on the non-conflicting locus the
community's randomized choice rule has nothing left to choose
between.  The simultaneous composite --- remove every ordinary-input
token at once, then produce, the community's own firing order ---
equals the chain in every enabled order under \emph{joint supply}
(\texttt{chainFire\_eq\_compositeFire}), and the supply hypothesis is
load-bearing in the statement, not merely the proof: two swap-token
firings pass one token through a chain and land it unchanged, while
the composite would read double (the module's refutation control;
the community classifies that pair as conflicting --- both sets
reserve one token --- so the hypothesis carves the rule's own
boundary).  The reservation precondition renders the community's
side condition at the counting stratum --- members enabled,
consumption jointly supplied, each member's enabling-arc reads
fitting the pool no sibling consumes, outcome within capacity
(\texttt{JointReservation}) --- and the scope condition the
compositional reduction genuinely demands is explicit: intermediate
markings need not interpolate between endpoints, and in the capacity
fragment the cheap sufficient condition is headroom --- source plus
output total within capacity
(\texttt{setNonConflicting\_of\_headroom}; a reservation-respecting
three-element composite whose producer-first order blocks mid-chain
is the control that keeps the hypothesis honest).  Under reservation
and headroom every enumeration is an enabled chain, all agree, and
each equals the composite (\texttt{r1r3\_composite\_adequacy});
multiplicity is native --- one transition with $k$ token sets is a
replicated multiset, so the one-transition and several-transition
rules are one vocabulary.  An enabled immediate's branch row in the
compiled automaton is a point mass at its firing, so an enabled
chain is an absorption path of the cascade closure as already built:
the observed tangible process cannot distinguish the community's
one-instant composite from any scheduler's zero-time chain.
Mechanized, sorry-free, beside the cascade module
(\texttt{lean/su-functor/}, module
\texttt{SuFunctor.Adequacy.StochasticPetriSetFiring}).

The flattening toward the hybrid-process reading
(section~\ref{sec:mt-gshs}) enters as a bridge between presentations,
not as a definition, and its jump-skeleton stratum is mechanized ---
at the uncoloured grain in full, and one colour grain up at the
deterministic-firing slice (the coloured paragraph below).  The
translation's definedness premises are certificate-shaped rows:
initial tangibility (no immediate transition enabled at the initial
support --- decidable, with the closed chain's initial row passing
through untouched as a receipt); the invariant token bound in the
witness-vector idiom --- a positive place weighting that no
transition of either species increases, with step soundness for both
species, closure under reachability, and an intrinsic per-place bound
independent of the declared capacity, which is the clause a
flattening at the unbounded grain consumes for the finiteness of its
mode roster (a token-creating control net admits \emph{no} witness
vector, so the row has teeth); and cascade termination, consumed from
the landed rank certificate.  The triangle itself is packaged from
atom-transport agreement (\texttt{flatten\_triangle\_of\_agree}): a
flat presentation on an abstract mode carrier that agrees with the
scheduler-closed net off-diagonally and initially denotes a
constraint isomorphic to the net's compiled reading --- an invertible
globular cell between two house compiles, with no community semantics
and no process construction inside the cell.  The agreement is no
longer a hypothesis: the flattening is \emph{constructed} for every
net and every declared scheduler at this grain --- the community's
own construction, one flattened transition per delay transition,
its jump intensity the ordered race census times the rate, its
reset the cascade-closed firing kernel --- and the agreement law is
derived from the construction (rate rows by Fubini regrouping and
collapse of the fired-marking indicator; initial rows from the
tangibility receipt, consumed not re-proved), so every
\texttt{\_of\_agree} consumer downstream is discharged outright and
the composed corollary states its conclusion with no agreement
hypothesis.  The drain-net instance survives as a check on the
general definition: the constructed flattening reproduces its
hand-written flattening data exactly, off-diagonally and initially
(the flattened drain edge carrying the arc-multiplicity dictionary
through a two-step cascade-closed reset).  The \emph{composed
corollary} pastes this cell with the two adequacy squares it bridges
--- the net square's cascade face on one side and, on the other, the
flattened object's own square, which at this stratum is the
continuous-time chain square, a flattened presentation being a finite
chain --- so the token-game community law and the flattened
execution's community law are one invertible globular cell apart at
every jump horizon, and the same pasting read at the judgment stratum
says a law models the flattened compiled presentation exactly when it
models the net's.  The demonic and angelic faces transfer verdicts
through the same pasting as directed cells at the community grain
(sound directions only; no envelope equality is stated), and on the
drain-net instance an actual flattened-side computation --- the
three-mode chain takes its drain edge with probability $3/4$ --- is
consumed as net-side verdicts by transport, never recomputation.
Graded exactly: this is a finite-jump-skeleton analogue, free of the
finite-expected-jumps hypothesis, of Everdij and Blom's process-level
bisimulation --- different conclusion, different observational grain,
neither result implying the other.  The holding-time \emph{data} (the
exit-rate family) agrees on both sides; and the time-indexed
comparison through transient distributions is now mechanized on the
killed trajectory-measure side --- the killed construction itself, a
prefix-marginal factorization (timed prefix marginals as skeleton
prefix law, holding attachment, and clock pairing), and the lift
from skeleton-law agreement to equal killed transient reads at every
horizon, with mutants refused at four preregistered loci --- its
grain bridge for consuming the published non-Zeno hypothesis
verbatim remaining the named residue.  The triangle, the derived
agreement, the composed corollary, and the killed time-indexed lift
are mechanized in
\texttt{SuFunctor.Adequacy.StochasticPetriFlatten},
\texttt{SuFunctor.Adequacy.StochasticPetriFlattenAgreement},
\texttt{SuFunctor.Adequacy.StochasticPetriComposed}, and
\texttt{QtslAux.KilledTimeIndexed}.

The same story closes one colour grain up, at the
deterministic-firing slice.  The premise rows carry over with their
colour reading --- initial tangibility unchanged; the invariant token
bound with \emph{colour-indexed} weights, its step soundness reading
the consumed and produced colour counts of each transition at each
colour vector; cascade termination from the coloured rank
certificate --- and the one premise the uncoloured development could
only name, per-place colour-box discharge, is carried by the
formation itself at this fragment: the colour boxes are finite types
by construction.  The triangle cell is packaged against a flattened
presentation on an abstract mode carrier exactly as before (the
socket a hybrid-grain carrier identification plugs into), and the
flattening is again \emph{constructed}: one flattened transition per
delay transition \emph{per colour vector}, its jump intensity the
coloured ordered census times the colour-dependent rate, its reset
the cascade-closed firing kernel through the coloured cascade, with
the agreement derived and the composed corollary stated with no
agreement hypothesis --- the token-game community law of a coloured
net and its flattened execution's community law, one invertible
globular cell apart at every jump horizon, at the deterministic
slice.  On the control net the certificate is the colour dependence
itself: from a one-of-each marking the flattened jump chain drains
with probability $2/5$ and flips back with $3/5$, the two branches
differing \emph{only} in the colour their selections read --- a rate
map reading colours, flowing through the flattening into the
community law, consumed as a token-game verdict by transport.
Mechanized, sorry-free
(\texttt{SuFunctor.Adequacy.StochasticPetriColouredComposed}); the
hybrid-grain residue stays staged, named below.

The angelic face closes the scheduler-reading triple.  Liveness-style
claims quantify \emph{existentially} --- some resolution of the
conflicts achieves the claim --- and the honest existential reads the
radially maximal slice of the credal set, for the same reason the
universal lower kind does: the credal presentation is downward
closed, so the whole set contains the zero model, and an un-sliced
existential would be witnessed by zero for every constraint the zero
model satisfies (mechanized as a degeneracy lemma, with the
zero-threshold upper-bound event atoms as the concrete family).  On
an equal-mass family the slice is the family, so the angelic verdict
is existential quantification over the enabling-respecting arena
(\texttt{SatAngelic}, \texttt{satAngelic\_ispnDemonicCredal}; the
operator is generic over unions of rays, so the decision and
automaton frames inherit it unchanged whenever wanted).  One
membership fact yields both directed transfers, with the
sense-opposition living at the verdict stratum rather than in any
reversed cell: universally certified claims restrict to every arena
member, and claims established on a member --- the transported
flattening's law included --- are angelic witnesses
(\texttt{ispn\_demonic\_verdict\_to\_member},
\texttt{flatten\_angelic\_of\_agree}, with the packaging as a
directed cell into the slice object; sound directions only, no
envelope equality).  And the two verdicts genuinely differ on the
shipped conflict net: the demonic face's arena-plurality control is
upgraded here to the \emph{law} level --- the uniform and
always-discard members' closed path laws differ at jump horizon one,
the uniform member jump-stepping to the moved token with probability
one while the always-discard member's closed chain absorbs at the
source --- so the claim ``the drain trajectory carries full weight''
is angelically satisfied and demonically refused.  Mechanized,
sorry-free, closing the module family
(\texttt{SuFunctor.Adequacy.StochasticPetriAngelic}).

The coloured generator faces are landed at their deterministic
(Dirac) slice: per-place finite colour boxes, colour-dependent delay
rates, guard predicates read as per-vector cylinders on both species,
and deterministic firing measures, with the delay-race and cascade
squares closed by the same engine reuse through a per-colour-class
multiplicity dictionary, conflict resolution over (transition,
colour-vector) pairs, and a degeneracy theorem recovering the
uncoloured fragment on the nose at the one-colour instance ---
mechanized, sorry-free
(\texttt{SuFunctor.Adequacy.StochasticPetriColoured},
\texttt{SuFunctor.Adequacy.StochasticPetriColouredCascade}).
The scheduler-quantification faces ride the same coloured arena: the
demonic reading quantifies enabling-respecting adversaries over
(transition, colour-vector) pairs --- at colour the demon selects
\emph{tokens}, not merely transitions, and a single transition's
blocked selections, whether by supply or by guard cylinder, already
refute the branch-level Zeno certificate that would have transported
the automaton square on the nose --- with the set-valued square, the
angelic dual, and both directed verdict transfers flowing from one
membership lemma; a recolouring control net admitting no cascade rank
at all pins the rank hypothesis as substantive rather than
decorative.  Mechanized, sorry-free
(\texttt{SuFunctor.Adequacy.StochasticPetriColouredDemonic}).
What stays staged, by name: the firing-measure \emph{kernels} beyond
the deterministic slice (the kernel presentation is deliberately left
as an abstract seam for a forthcoming probabilistic vocabulary, so
nothing here entrenches an encoding); the guard species (boundary
forcing); the coloured flattening's hybrid-grain residue --- the
colour-vector carrier identification with its padding obligations
(the flattening itself, its premise rows, and the composed corollary
having landed at the deterministic slice above); and the exponential
time-law stratum, exactly as for the plain continuous-time frame.

\paragraph{Variadic certificates: one row per Petri transition.}
A certificate for a net whose token count is unbounded across runs
faces two apparent infinities: unboundedly many token
configurations, and --- since a transition fires once per eligible
token selection --- a countably infinite family of transition
\emph{instances}, each seemingly demanding its own certificate row.
Neither infinity survives contact with the syntax, and the two
mechanisms are worth stating separately, because they fail
separately when abused.

The state side is the plate discipline of
section~\ref{sec:hyp-plates}: certificates are polynomials in pooled
coordinates --- token counts $\vec{k}$ and per-place feature pools
$\vec{s}$ --- and obligations are polynomial sequents over
$(\vec{k}, \vec{s})$ with coupling, count, and rounding hypothesis
rows, their number and dimension set by the net, never by a marking
bound. The transition side rides the same grammar fact
(no expression names ``the $i$-th token,'' so transition data
depends on selected tokens only through their colours): the
per-instance obligations of one Petri transition collapse to a
single \emph{footprint sequent} in distinguished-token variables ---
the selected tuple's colours, the pooled coordinates, the counts ---
with dimension set by the transition's arity and the feature
dictionary.
One symbolic row covers every instance at every marking. The
countable family was never infinite.

What the collapsed row can \emph{claim} depends on the aggregator,
and the dependence is sharp. Write a firing's footprint as
(consumed multiset $X$, produced multiset $Y$), and consider a
certificate value that is a \emph{commutative-monoid fold} of
per-token features --- the fold structure is itself load-bearing: an
arity-normalized aggregate such as a mean is not one (a birth
changes its denominator under every untouched token, so a footprint
row can hold with equality while the ambient value rises), and for
such aggregates only the direct polytope form above applies, never a
footprint row. For \emph{safety} rows (the aggregate does not
increase), monotonicity of the fold's translations suffices, at
every token count: the footprint row ``the produced contribution is
bounded by the consumed contribution'' transfers to the whole
marking by one monotone translation, births and deaths included (for
signed features the death row is no longer automatic --- a death
whose consumed contribution is negative is \emph{refused} unless a
guard excludes it: a refusal, not an unsoundness). For
\emph{progress} rows the conditions are strictly stronger, and each
is necessary --- drop one and a finite countermodel follows:

\begin{itemize}
\item the per-firing decrease must carry a fixed rational gap
  $\delta > 0$, stated in the row --- strict decrease alone
  certifies nothing (halve one token's value forever: every step
  strictly drops, the run is infinite, the value never reaches its
  floor);
\item the gap must be stated in a scale where the fold is a
  \emph{translated sum} --- for sum-pools the footprint gap row
  transfers verbatim, uniformly in the marking; a mean fails twice
  over (no locality, as above, and a per-argument modulus of
  $\delta/n$ at $n$ tokens, so no uniform gap survives an unbounded
  token count);
\item a floor row (the aggregate bounded below on reachable
  markings) must be discharged --- with signed features and token
  births, every footprint row can hold while progress-rowed firings
  continue forever down an unbounded descent;
\item and \emph{idempotent folds are progress-incapable in every
  scale}: for $\max$-aggregated certificates the footprint gap fails
  to move the aggregate whenever an untouched token dominates (two
  tokens of equal value; fire on one), a regenerating two-transition net
  keeps the max constant through infinitely many gap-verified
  firings while the corresponding sum-row is falsified --- the sum
  correctly refuses a certificate the max would wave through --- and
  no strictly monotone change of scale repairs this, since $\max$ is
  a fixed point of conjugation. Max-aggregated certificates certify
  safety but support no footprint-only progress row against
  arbitrary ambient tokens (structural anti-masking conditions --- a
  singleton place, a guard consuming every maximal token --- are a
  different and dearer certificate species, not this row); a
  footprint gap row on an idempotent fold is a formation error, not
  a modelling choice.
\end{itemize}

The positive family is correspondingly wide: for any strictly
increasing $\varphi$ on the feature range, the $\varphi$-conjugated
sum $\varphi^{-1}(\sum_t \varphi(v_t))$ --- the unnormalized
$\ell_p$ folds at $\varphi(x) = x^p$ over nonnegative features,
interpolating from sum ($p = 1$) toward $\max$ ($p \to \infty$);
log-sum-exp for $\varphi(x) = e^{\beta x}$ (with the empty marking
sent to $-\infty$) --- is strictly monotone in every argument at
every finite parameter, which is exactly the property the $\max$
endpoint lacks. (The count-normalized \emph{power mean} is
deliberately not in the family: normalization re-imports the mean's
locality failure.) Progress rows for the family are stated on the
$\varphi$-space sum (an ordinary pool over the transformed features,
with a pinned $\varphi$-space gap), where the transfer is uniform in
the marking; verdicts transport monotonely for thresholds in
$\varphi$'s domain (compare $\sum v_t^p$ against $t^p$ at $t \ge 0$
--- the checker never computes a root), while the drop measured in
aggregate units necessarily shrinks as the \emph{transformed mass}
$\sum_t \varphi(v_t)$ grows, for every superlinear $\varphi$ ---
with an explicit rational floor on bounded ranges (for integer $p$
and aggregate bound $B$: at least $\delta_\varphi / (p\,B^{p-1})$)
--- which is the reason the $\varphi$-space statement is the
normative one. For integer $p$ the transformed features are polynomials and
every obligation is an ordinary exact-rational sequent today;
non-integer $p$ and the log-sum-exp family ride the certified
enclosure register for transcendental symbols
(section~\ref{sec:hyp-genus}).

Grading, in the book's standard senses: the grammar-equivariance
layer is mechanized (the plate modules of
section~\ref{sec:hyp-plates}, including descent to marking-indexed
denotations); the pooled compile and the footprint-sequent shape are
the shipped row vocabulary of the certificate stratum; the
fold-transfer statements above are elementary, each dropped side
condition refuted by a finite exact-rational instance (the
countermodels of the progress list above), with their formal lemma
pack a named, priced successor rather than a landed artefact; the
row vocabulary's first exercising campaign at industrial shape is
the worked example of chapter~\ref{ch:qtsl-minifab}
(section~\ref{sec:minifab-exact}). Where
the event \emph{count} is itself the claim --- expected crossings,
Zeno exclusion, deadline tails --- the certificate is the crossing
budget of the hybrid-process correspondence
(section~\ref{sec:mt-gshs}), not a ranking row; the two prices are
different and deliberately not interchangeable.

\paragraph{Status.} Representation and production: shipped --- the
declaration vocabulary above is in the conformance corpus, whose
worked fixture is a two-place queue/server net with a delay and an
immediate transition, ordinary and enabling arcs, and one inhibitor
view on a capacity-two place; the view-on-unbounded-place refusal is
exercised from both language surfaces. At the uncoloured bounded
fragment the squares certify, the toolkit's producer surface derives
the entire marking-space bridge exactly from the stored net: the
bounded marking space under a declared mixed-radix convention
(shipped as a decode table, so the index correspondence is
checkable), the compiled race rows in the unordered convention with
the arc-multiplicity dictionary the compiler owns --- the ordered
census recomputed as an exact differential in the test suites, the
dictionary theorem as arithmetic --- the tangible/vanishing sort
split with timed rows at tangible sources only (maximal progress as
construction), point-mass immediate branches, the declared or
uniform-over-enabled scheduler, and the closed observed chain by
handing the compiled net to the automaton kit
(section~\ref{sec:mt-ma}: cascade rank, absorption at the working
fuel, vanishing initial mass pushed through the cascade, the
initial-tangibility pass-through re-derived as a receipt) or, for
the pure-delay net, to the chain kit: the squares' five-storey
engine reuse, run producer-side. The relay-cascade net --- whose
every delay firing returns through a two-step cascade, thinning the
closed chain to nothing --- is a committed cross-language
byte-identity vector, and the API block below quotes the kit's
surface. The inhibitor view stays outside the compiled fragment by
formation (its complement-place elaborator remains the named
successor), and per-place capacities, state-dependent rates, and
coloured tokens refuse at the kit, each naming its stratum.
Semantics and certification: the count
stratum has its first certified artefact --- a population invariant
carried across a transfer rendezvous through two independent checker
lanes over a Bag carrier (section~\ref{sec:mt-bag}); the token game's
statics square, the delay-race face, the immediate/cascade face
--- at the declared scheduler and under scheduler quantification
alike --- and the set-firing face (simultaneous composite firings of
non-conflicting immediates) are mechanized (the Adequacy paragraph
above), with the coloured generator faces landed at their
deterministic slice and the coloured scheduler-quantification faces
alongside them, while the guard species (boundary forcing) and the
firing-measure kernels beyond the deterministic slice
remain the named next legs; net-of-nets composition
runs through the token-colouring closure
of the composition inventory
(section~\ref{sec:cp-netsofnets}), with only its shared-place face
reserved to the state-sharing axis.

\syscatmeta{\begin{mathpar}
\inferrule*[Right=\rulename{S-SDCPN}]{
  P : \overrightarrow{\langle \mathsf{id},\, \mathsf{name}?,\,
  \mathsf{token\_system} : \mathsf{System},\, \mathsf{capacity}?
  \rangle} \text{ places} \\
  T : \overrightarrow{\langle \mathsf{id},\, \mathsf{name}?,\,
  \mathsf{species} \rangle} \text{ transitions} \\
  A : \overrightarrow{\langle \mathsf{transition},\,
  \mathsf{place},\, \mathsf{kind} \rangle} \text{ arcs, closed over
  } P \text{ and } T \\
  M : \overrightarrow{\langle \mathsf{place},\, \mathsf{tokens}
  \rangle} \text{ a stored marking, canonical per place} \\
  \text{every } \mathsf{inhibitorView} \text{ arc's place declares
  a capacity}
}{
  \mathrm{SDCPN}(P;\, T;\, A;\, M) : \mathsf{System}
}
\end{mathpar}
\begin{align*}
\mathsf{species} &::= \mathsf{delay}(\mathsf{rate}) \mid
  \mathsf{immediate} \mid \mathsf{guard}
\\
\mathsf{kind} &::= \mathsf{consume}(k) \mid \mathsf{produce}(k)
  \mid \mathsf{enabling}(k) \mid \mathsf{inhibitorView}(k)
\end{align*}
The capacity premise is the bounded locus as licence: an inhibitor
view on an unbounded place is not a refused edge case but a missing
constructor.}
\ruleblock{API}
\sourcenote{Producer surface of the \texttt{mtk} toolkit for this model type: the TypeScript arrangement (npm package \texttt{@maitria/mtk}) and the Rust catalogue crate (\texttt{crates/mtk-catalogue}) --- one Rust implementation underneath, reached from TypeScript in-process through a native binding, so the two surfaces cannot diverge. Signatures and constructions are extracted from the toolkit sources at every regeneration of this catalogue; each construction is quoted verbatim from the toolkit's conformance corpus, code whose emitted certificate bytes are verified byte-identical against the committed golden vectors (helper values named \texttt{fx\_*} or \texttt{f.*} are that corpus's shared fixtures, defined beside the quoted lines). This API constructs and emits; it never checks --- the checker re-derives every claim independently.}
\apilang{TypeScript}{system.sdcpn}
\apiprov{signature --- \texttt{ts/src/system.ts}}
\begin{tssrc}
/** S-SDCPN: places are Bags, arcs are patterns (the inhibitor VIEW is
 * admissible only on a place with declared capacity). */
export const sdcpn = profileFactory("sdcpn");

/** One marking-index decode row of the SDCPN kit: `index` in the
 * mixed-radix convention (places in declaration order, first place
 * most significant), `counts` the per-place token counts. */
export interface SdcpnMarkingVal {
  readonly index: number;
  readonly counts: readonly number[];
}

/** The derived SDCPN marking-space kit at the uncoloured bounded
 * fragment: the marking decode table, the compiled race rows in the
 * unordered convention (the compiler-owned arc-multiplicity
 * dictionary `∏_p (demand_p)!` times the per-token-set selection
 * count times the rate — the community's ordered census divided by
 * exactly the mechanized dictionary theorem), the tangible/vanishing
 * sort split (timed rows at tangible sources only — maximal progress
 * as construction), action-keyed point-mass immediate branches, the
 * resolved scheduler (declared, or the uniform-over-enabled instance
 * the mechanized square ships), point-mass initial rows, and the
 * closed observed process by ENGINE HANDOFF — the automaton kit
 * where immediate transitions exist (cascade rank, absorption, the
 * reduced tangible chain, vanishing initial mass pushed through the
 * cascade), the chain kit otherwise. Derived, never trusted — the
 * checker re-derives the marking space, the census, and the
 * closure's own storeys. */
export interface SdcpnGenKit {
  readonly capacity: number;
  readonly placeOrder: readonly number[];
  readonly markings: readonly SdcpnMarkingVal[];
  readonly raceRows: readonly {
    readonly src: number;
    readonly dst: number;
    readonly rate: string;
  }[];
  readonly immediateRows: readonly {
    readonly src: number;
    readonly action: number | null;
    readonly dst: number;
    readonly prob: string;
  }[];
  readonly schedRows: readonly {
    readonly state: number;
    readonly action: number;
    readonly prob: string;
  }[];
  readonly initRows: readonly {
    readonly state: number;
    readonly prob: string;
  }[];
  readonly closure:
    | { readonly kind: "chain"; readonly chain: CtmcGenKit }
    | { readonly kind: "automaton"; readonly ma: MaGenKit };
}

/** Derive the SDCPN marking-space kit from a mechanism declaration
 * (the `sdcpn` family; an optional scheduler rides the mechanism's
 * own scheduler family — absent is the uniform instance) and a
 * per-state certificate observable on the tangible markings, by
 * marking index. Admission fences — uniform capacity, unit colour,
 * constant nonnegative rates, no guard species, no inhibitor views,
 * no duplicate arcs, all-or-none marking, the engine cap on the
 * marking space — and the closure's own discipline (scheduler
 * support within the enabled immediates, the immediate-cycle Zeno
 * door) all come back as witness-shaped refusals naming their
 * geometry. */
export function sdcpnGenKit(
  init: MechanismInit,
  cert: readonly StateVal[],
): Result<SdcpnGenKit, Refusal> {
  const raw = native.sysSdcpnGenKit(mechanism(init).nativeHandle, cert) as
    | { kit: SdcpnGenKit }
    | { refusal: unknown };
  return "refusal" in raw
    ? refused(narrowRefusal<Refusal>(raw.refusal))
    : ok(raw.kit);
}
\end{tssrc}
\apiprov{construction --- \texttt{ts/tests/system-native.test.ts}}
\begin{tssrc}
const token = sys.unwrap(
  sys.ode({
    name: "token",
    state: [{ id: 0, name: "energy" }],
    drift: [{ var: 0, poly: mono(1, 0, 1, "-1/100") }],
  }),
);
return sys.unwrap(
  sys.sdcpn({
    name: "queue-server",
    sdcpn: {
      places: [
        { id: 0, name: "queue", tokenSystem: token },
        { id: 1, name: "server", tokenSystem: token, capacity: 2 },
      ],
      transitions: [
        { id: 0, name: "arrive", species: { kind: "delay", rate: constant(1, "3/2") } },
        { id: 1, name: "admit", species: { kind: "immediate" } },
      ],
      arcs: [
        { transition: 0, place: 0, kind: { kind: "produce", mult: 1 } },
        { transition: 1, place: 0, kind: { kind: "consume", mult: 1 } },
        { transition: 1, place: 1, kind: { kind: "enabling", mult: 1 } },
        {
          transition: 1,
          place: 1,
          kind: {
            kind: "inhibitorView",
            threshold: 2n,
            invariant: "pinv:server-complement",
          },
        },
      ],
    },
  }),
);
\end{tssrc}
\apiprov{The marking-space kit compiles the uncoloured bounded fragment the mechanized square certifies: race rows in the unordered convention (the compiler-owned arc-multiplicity dictionary --- the ordered census divided by exactly the dictionary theorem), the tangible/vanishing sort split, point-mass immediate branches, and the closed observed process by engine reuse --- the automaton kit where immediate transitions exist, the chain kit otherwise. The kit's Rust face is \texttt{mtk\_catalogue::system::sdcpn\_gen\_kit}; its closed-form controls (the census through the dictionary, the two-step cascade with vanishing initial mass pushed through absorption) ride the test suites on both language surfaces.}
\apilang{Rust}{mtk\_catalogue::system::presets::Sdcpn::manifest}
\apiprov{signature --- \texttt{crates/mtk-catalogue/src/system/presets.rs}}
\begin{rustsrc}
preset!(
    /// System Catalogue: SDCPN — places are Bags, arcs are patterns.
    /// The inhibitor VIEW is admissible only on a place with declared
    /// capacity (the bounded-locus converse); no unbounded-inhibitor
    /// constructor exists at all.
    Sdcpn, "SDCPN",
    req = [Sdcpn],
    allowed = [],
    extra = |m| {
        let s = m.sdcpn.as_ref().expect("required by signature");
        // Roster ids key arcs, markings, and fusion sketches: gate
        // uniqueness first (a duplicate transition id would let one
        // event id denote several species — review round 1, MF5).
        for (second, p) in s.places.iter().enumerate() {
            if let Some(first) = s.places[..second].iter().position(|q| q.id == p.id) {
                return Err(Refusal::DuplicateRosterId {
                    kind: "place",
                    first,
                    second,
                    id: p.id,
                });
            }
        }
        for (second, t) in s.transitions.iter().enumerate() {
            if let Some(first) = s.transitions[..second].iter().position(|u| u.id == t.id) {
                return Err(Refusal::DuplicateRosterId {
                    kind: "transition",
                    first,
                    second,
                    id: t.id,
                });
            }
        }
        for (arc_index, arc) in s.arcs.iter().enumerate() {
            let place = s.places.iter().find(|p| p.id == arc.place);
            let trans_ok = s.transitions.iter().any(|t| t.id == arc.transition);
            let Some(place) = place else {
                return Err(Refusal::DanglingEndpoint {
                    profile: "SDCPN",
                    index: arc_index,
                    place: Some(arc.place),
                    transition: (!trans_ok).then_some(arc.transition),
                });
            };
            if !trans_ok {
                return Err(Refusal::DanglingEndpoint {
                    profile: "SDCPN",
                    index: arc_index,
                    place: None,
                    transition: Some(arc.transition),
                });
            }
            // The token game's lower-bound reasoning is defined over
            // positive counts (review round 1, MF5/MF10).
            let mult = match &arc.kind {
                ArcKind::Consume { mult }
                | ArcKind::Produce { mult }
                | ArcKind::Enabling { mult } => Some(mult),
                ArcKind::InhibitorView { threshold, .. } => Some(threshold),
            };
            if let Some(mult) = mult {
                if mult.0 < num_bigint::BigInt::from(1u8) {
                    return Err(Refusal::NonPositiveArcMult {
                        arc_index,
                        transition: arc.transition,
                        place: arc.place,
                        mult: mult.0.to_string(),
                    });
                }
            }
            if matches!(arc.kind, ArcKind::InhibitorView { .. }) && place.capacity.is_none() {
                return Err(Refusal::InhibitorOnUnboundedPlace {
                    arc_index,
                    place: arc.place,
                });
            }
        }
        for (entry, mk) in s.marking.iter().enumerate() {
            if let Some(first) = s.marking[..entry].iter().position(|m2| m2.place == mk.place) {
                return Err(Refusal::DuplicateMarkingEntry {
                    first,
                    second: entry,
                    place: mk.place,
                });
            }
            let Some(place) = s.places.iter().find(|p| p.id == mk.place) else {
                return Err(Refusal::DanglingMarkingPlace {
                    entry,
                    place: mk.place,
                });
            };
            let site = format!("sdcpn.marking[{entry}] (place {})", mk.place);
            let d = inner_state_arity(&place.token_system, &site)?;
            bag_marking_gates(&mk.tokens, d, place.capacity.as_ref(), &site)?;
        }
        // NAMED GAP, deferred with mechanism: per-tuple delay-rate
        // arity and slot symmetry at the NET stratum need the
        // selection roster across places — the staged
        // intensional-family tranche; the Bag preset checks both on
        // the single-place locus today.
        Ok(())
    }
);
\end{rustsrc}
\apiprov{construction --- \texttt{crates/mtk-catalogue/tests/common/mod.rs}}
\begin{rustsrc}
let token = {
    let m = MechanismBuilder::new()
        .state_var(0, Some("energy"))
        .drift(None, 0, mono(1, 0, 1, r(-1, 100)))
        .finish();
    presets::Ode::manifest(m)
        .expect("token System conforms to the ODE profile")
        .named("token")
};
let decl = SdcpnDecl {
    places: vec![
        SdcpnPlace {
            id: 0,
            name: Some("queue".to_owned()),
            token_system: Box::new(token.clone()),
            capacity: None,
        },
        SdcpnPlace {
            id: 1,
            name: Some("server".to_owned()),
            token_system: Box::new(token),
            capacity: Some(maitria_witness_types::WitInt::from(2)),
        },
    ],
    transitions: vec![
        SdcpnTransition {
            id: 0,
            name: Some("arrive".to_owned()),
            species: TransitionSpecies::Delay {
                rate: constant(1, r(3, 2)),
            },
        },
        SdcpnTransition {
            id: 1,
            name: Some("admit".to_owned()),
            species: TransitionSpecies::Immediate,
        },
    ],
    arcs: vec![
        ArcRow {
            transition: 0,
            place: 0,
            kind: ArcKind::Produce {
                mult: maitria_witness_types::WitInt::from(1),
            },
        },
        ArcRow {
            transition: 1,
            place: 0,
            kind: ArcKind::Consume {
                mult: maitria_witness_types::WitInt::from(1),
            },
        },
        ArcRow {
            transition: 1,
            place: 1,
            kind: ArcKind::Enabling {
                mult: maitria_witness_types::WitInt::from(1),
            },
        },
        ArcRow {
            transition: 1,
            place: 1,
            kind: ArcKind::InhibitorView {
                threshold: maitria_witness_types::WitInt::from(2),
                invariant: Some("pinv:server-complement".to_owned()),
            },
        },
    ],
    // No declared marking: absent on the wire, so this fixture's
    // committed bytes predate and survive the marking family.
    marking: vec![],
};
let m = MechanismBuilder::new().sdcpn(decl).finish();
presets::Sdcpn::manifest(m)
    .expect("fixture conforms to the SDCPN profile")
    .named("queue-server")
\end{rustsrc}

\section{Shallowly-Embedded Petri Nets}
\label{sec:mt-petri}

\paragraph{The frame.} Places, transitions, tokens; the
bread-and-butter questions are reachability and coverability of
markings.

\paragraph{Profile.} A shipped theory of its own: places, transitions,
and total flow matrices in the incidence presentation (arc
multiplicities native, which discharges the ``generalized'' qualifier
of early programme drafts for the arc-weight reading), with markings
as first-class elements. Time: none intrinsically --- the token game
is untimed at this rung; the stochastic and timed wings belong to
section~\ref{sec:mt-gshs}.

\paragraph{Semantics.} The distinctive entry, and the shallow
embedding the section title names: the token game runs \emph{as chase
derivation, inside the substrate} --- enabledness by negation-free
spine walk, firing by successor invention, and marking identity by
equality generation, with no oracle in the loop. This is the one
frame at this preview whose representation and semantics both run
natively in the modelling substrate, rather than beside it.

\paragraph{Adequacy (the interleaving square, through the compiled
route).} The shallow embedding above runs the token game natively;
the same untimed frame also enters the catalogue's \emph{compiled}
route, and the adequacy square closes there whole.  The compile, in
one breath: places become bags of unit-colour tokens --- and a bag
of unit-colour tokens \emph{is} its count, so markings are count
vectors --- while each transition becomes one row under the demonic
wing of the transition-clock vocabulary
(chapter~\ref{ch:qtsl-markov-kernels}) at a strictly positive dwell
floor, with a Dirac outcome law, the atomic reset
$m \mapsto m - \mathrm{pre} + \mathrm{post}$, and enabledness as the
arc-cover gate.  The compiled semantics is therefore a
scheduler-seat family of deterministic timed runs, and its
classical corner is not a path law but a relation: the square
states that the possibilistic visited-marking face of the family
\emph{equals} classical interleaving reachability, for every
strictly positive dwell assignment (\texttt{face\_eq\_reach}) ---
the soundness leg spends only enabledness, and the completeness leg
realizes each classical firing at exactly the dwell floor.  The
flagged check comes back quiet: the face is independent of the
dwell assignment (\texttt{face\_dwell\_uniform}), so the uniform
unit clock that transcription conventions pin is nowhere
load-bearing --- the transcription sentence of
chapter~\ref{ch:qtsl-markov-kernels}, that the unit-dwell embedding
is faithful for ordering-and-counting properties and deliberately
not for metric-time properties an untimed net never had, is at this
face a theorem rather than a convention's gloss.

The wiring into the semantic universe lands in the record shape of
the demonic corners (sections~\ref{sec:mt-mdp} and~\ref{sec:mt-ma},
and the scheduler-quantified face of section~\ref{sec:mt-sdcpn}):
each resolution's law at a horizon is the Dirac law on its marking
trajectory, runs shorter than the horizon idling at their final
marking; the compiled presentation is a one-atom projection, so
finite presentation is untouched by the resolution family's
infinitude; and the ideal point of a law models the compiled
demonic presentation exactly when the law is a member of the
demonic law-set (\texttt{su\_vanilla\_demonic\_adequacy}, with its
mem, bottom, and slice forms).  One structural fact is stated in
the module rather than dressed up: this species has no rows-level
compilation in the tree --- every kernel being Dirac, the scheduler
seat's resolutions are the whole mechanism --- so the feeding lemma
is definitional (\texttt{vanillaDemonicC\_iff}) and the square's
content lives one retract below the law stratum, at the
visited-marking face (\texttt{recordVisited\_eq\_face}).  The
headline, read at the universe's judgment: a marking is visited by
the support of some law whose ideal point models the presentation,
at some horizon, exactly when it is classically reachable
(\texttt{su\_vanilla\_reach\_iff}) --- and the horizon quantifier
is honest rather than lossy, since every visited marking is visited
at a finite index and prefixes of admissible runs are admissible,
so the finite-run family carries the whole reachability face.
Dwell-uniformity transports to the same stratum
(\texttt{su\_vanilla\_dwell\_uniform}), and safety transfers as the
customer shape for invariant-class examinations --- a marking
predicate holds on everything the record semantics can visit
precisely when it holds on the classically reachable set
(\texttt{su\_vanilla\_safety}) --- with one-safeness and
unboundedness exercised concretely on a two-place toggle and a
one-place source (the dwell floor bounds nothing).  Liveness is
withheld by design, not staged: the demonic clock carries no
fairness, priority, or weight vocabulary, and the idling resolution
is always seated, so the honest current semantics asserts
reachability and safety, never progress.  Set-firing beyond
interleaving has its own face at the coloured development
(section~\ref{sec:mt-sdcpn}); inhibitor arcs and mixed
race-plus-demonic systems are untouched at this entry.  Mechanized,
sorry-free (\texttt{lean/su-functor/}, module
\texttt{SuFunctor.Adequacy.PetriVanilla}).

\paragraph{Status.} Certification: a three-place producer--consumer
ring saturates to exactly its ten-marking state space with
coverability verdicts and a place-invariant checked per reachable
marking; an overflow mutant is caught by a derived coverability fact,
and ablating the equality discipline reproduces divergence as a loud
non-termination, not a wrong answer. The hypothesis is stated in
the theory itself, not in a side note: token counts live on a finite declared ladder,
so fixpoint verdicts are decisions only up to that headroom --- exact
for structurally bounded nets, while unbounded nets refuse loudly and
coverability positives remain semi-decidable witnesses. Semantics:
the compiled route's interleaving square --- reachability,
dwell-uniformity, and the safety transfer --- is mechanized (the
Adequacy paragraph above), with no certificate rows minted by the
square itself. Open: the
ladder-horizon side conditions pinned to their oracle; net-of-nets
composition, which runs through the token-colouring closure of the
composition inventory (section~\ref{sec:cp-netsofnets}) via the
coloured wing; the stochastic wing.

\syscatmeta{\begin{mathpar}
\inferrule*[Right=\rulename{S-PETRI}]{
  \mathsf{places},\, \mathsf{transitions} \text{ declared finite
  sets} \\
  \mathsf{pre},\, \mathsf{post} \text{ total flow matrices over }
  \mathsf{places} \times \mathsf{transitions} \\
  \mathsf{marking} \text{ a first-class marking over }
  \mathsf{places}
}{
  \mathrm{Petri}(\mathsf{places};\, \mathsf{transitions};\,
  \mathsf{pre};\, \mathsf{post};\, \mathsf{marking}) :
  \mathsf{System}
}
\end{mathpar}
The incidence presentation, multiplicities native; totality over the
declared sets is closedness consumed as a sort --- every flow and
marking entry names a declared place and transition by construction.}
\ruleblock{API}
\sourcenote{Producer surface of the \texttt{mtk} toolkit for this model type: the TypeScript arrangement (npm package \texttt{@maitria/mtk}) and the Rust catalogue crate (\texttt{crates/mtk-catalogue}) --- one Rust implementation underneath, reached from TypeScript in-process through a native binding, so the two surfaces cannot diverge. Signatures and constructions are extracted from the toolkit sources at every regeneration of this catalogue; each construction is quoted verbatim from the toolkit's conformance corpus, code whose emitted certificate bytes are verified byte-identical against the committed golden vectors (helper values named \texttt{fx\_*} or \texttt{f.*} are that corpus's shared fixtures, defined beside the quoted lines). This API constructs and emits; it never checks --- the checker re-derives every claim independently.}
\apilang{TypeScript}{system.petri}
\apiprov{signature --- \texttt{ts/src/system.ts}}
\begin{tssrc}
/** S-PETRI: places, transitions, total flow matrices, marking (token
 * game as chase derivation). */
export const petri = profileFactory("petri");
\end{tssrc}
\apiprov{construction --- \texttt{ts/tests/system-native.test.ts}}
\begin{tssrc}
// Petri flow row citing an undeclared endpoint: which endpoint.
const dangling = sys.petri({
  petri: {
    places: [{ id: 0 }],
    transitions: [{ id: 0 }],
    pre: [{ place: 7, transition: 0, mult: 1 }],
  },
});
assert.ok(!dangling.ok);
\end{tssrc}
\apiprov{The TS corpus exercises this factory through its refusal geometry (the corpus's discipline that failure geometry is routing advice); the quoted construction is that witness, and the happy-path byte authority is the Rust corpus opposite.}
\apilang{Rust}{mtk\_catalogue::system::presets::Petri::manifest}
\apiprov{signature --- \texttt{crates/mtk-catalogue/src/system/presets.rs}}
\begin{rustsrc}
preset!(
    /// System Catalogue: Petri net — places, transitions, total flow
    /// matrices, marking (token game as chase derivation).
    Petri, "Petri net",
    req = [Petri],
    allowed = [],
    extra = |m| {
        let p = m.petri.as_ref().expect("required by signature");
        let place_ok = |id: u32| p.places.iter().any(|d| d.id == id);
        let trans_ok = |id: u32| p.transitions.iter().any(|d| d.id == id);
        for (index, e) in p.pre.iter().chain(p.post.iter()).enumerate() {
            if !place_ok(e.place) || !trans_ok(e.transition) {
                return Err(Refusal::DanglingEndpoint {
                    profile: "Petri net",
                    index,
                    place: (!place_ok(e.place)).then_some(e.place),
                    transition: (!trans_ok(e.transition)).then_some(e.transition),
                });
            }
        }
        for (index, e) in p.marking.iter().enumerate() {
            if !place_ok(e.place) {
                return Err(Refusal::DanglingEndpoint {
                    profile: "Petri net",
                    index,
                    place: Some(e.place),
                    transition: None,
                });
            }
        }
        Ok(())
    }
);
\end{rustsrc}
\apiprov{construction --- \texttt{crates/mtk-catalogue/tests/builders.rs}}
\begin{rustsrc}
let petri = MechanismBuilder::new()
    .petri(PetriDecl {
        places: vec![PlaceDecl {
            id: 0,
            name: None,
            capacity: None,
        }],
        transitions: vec![TransitionDecl { id: 0, name: None }],
        pre: vec![FlowEntry {
            place: 0,
            transition: 0,
            mult: WitInt::from(1),
        }],
        post: vec![],
        marking: vec![MarkEntry {
            place: 0,
            tokens: WitInt::from(3),
        }],
    })
    .finish();
Petri::manifest(petri).unwrap();
\end{rustsrc}

\section{Probabilistic Graphical Models}
\label{sec:mt-pgm}

The thesis lists three sub-frames --- Bayesian networks, structural
causal models, and Markov random fields --- and the inventory carries
the first and third through \emph{one} rendering: probabilistic
dependency graphs
(Richardson and Halpern \cite{richardson20}) as wires over trivial dynamics, one wire per conditional
table. Dependency graphs strictly subsume Bayesian networks and
factor graphs, and here the subsumption is a \emph{checked} sentence
rather than a cited one: the classic sprinkler network is reproduced
byte-identically through the dependency-graph rendering by two
decorrelated pipelines, and a two-source graph certified
\emph{inconsistent} witnesses strictly-beyond-network capability ---
every network-image graph is exactly consistent, so a certified
inconsistency is a certified exit from the network image. Structural
causal models carry their own theory, below. All three sub-frames
are static, and for each the supervisory-switch pattern
(section~\ref{sec:cp-switch}) is inapplicable --- the omission earned
by argument from the frame's own semantics, as the closure leg
requires, rather than declared. The pattern's composition data ---
per-location executions, a firing moment at which a trajectory
reaches a guard, an exit state read at the source and initial data
planted at the target --- has no instantiation from a semantics whose
denotation is a single static measure over joint assignments: a guard
over such variables has a truth-value \emph{distribution}, not a
trajectory of truth values, so nothing ever fires. The rows can
always be typed --- wrap the variables as a zero-drift carrier and
write a guard box --- but the resulting composite denotes switched
executions, which is no graphical model: any instantiation exits the
frame's semantic category before it composes. Interventions are no
counterexample (model surgery selected by the query, never an event
fired when a state condition becomes true), and temporal
extensions --- dynamic Bayesian networks, samplers run as chains ---
are different frames, with Markov-chain machinery, where switch
composition is applicable and required.

\syscatmeta{\begin{mathpar}
\inferrule*[Right=\rulename{S-PDG}]{
  W : \overrightarrow{\mathsf{StochasticWire}}
}{
  \mathrm{PDG}(W) : \mathsf{System}
}
\end{mathpar}
\begin{align*}
\mathsf{StochasticWire} \;::=\;\; &
\mathsf{conditional}\big(\mathsf{sources},\, \mathsf{target},\,
\overrightarrow{\langle \mathsf{given}, \mathsf{outcome},
\mathsf{prob} \rangle}\big)
\\
\mid\;\; &
\mathsf{potential}\big(\mathsf{vars},\,
\overrightarrow{\langle \mathsf{assignment}, \mathsf{weight}
\rangle}\big)
\end{align*}
One wire per conditional table; potential wires are raw nonnegative
tables with no normalization row (their local scale derived, never
stored) --- the atom that carries Markov random fields through the
same species.}
\ruleblock{API}
\sourcenote{Producer surface of the \texttt{mtk} toolkit for this model type: the TypeScript arrangement (npm package \texttt{@maitria/mtk}) and the Rust catalogue crate (\texttt{crates/mtk-catalogue}) --- one Rust implementation underneath, reached from TypeScript in-process through a native binding, so the two surfaces cannot diverge. Signatures and constructions are extracted from the toolkit sources at every regeneration of this catalogue; each construction is quoted verbatim from the toolkit's conformance corpus, code whose emitted certificate bytes are verified byte-identical against the committed golden vectors (helper values named \texttt{fx\_*} or \texttt{f.*} are that corpus's shared fixtures, defined beside the quoted lines). This API constructs and emits; it never checks --- the checker re-derives every claim independently.}
\apilang{TypeScript}{system.pdg}
\apiprov{signature --- \texttt{ts/src/system.ts}}
\begin{tssrc}
/** S-PDG: BN / MRF / PDG — stochastic wires over trivial dynamics. */
export const pdg = profileFactory("pdg");
\end{tssrc}
\apiprov{The TS conformance corpus holds no construction through this factory at this preview --- the cross-language byte-identity fixtures exercise six profiles, and this one's byte authority is the Rust corpus opposite; the factory takes the same init-object vocabulary as every other profile.}
\apilang{Rust}{mtk\_catalogue::system::presets::Pdg::manifest}
\apiprov{signature --- \texttt{crates/mtk-catalogue/src/system/presets.rs}}
\begin{rustsrc}
preset!(
    /// System Catalogue: BN / MRF (factor graph) / PDG —
    /// probabilistic-dependency wires over trivial dynamics (the
    /// stochastic-wire species); strictly subsumes Bayesian networks
    /// and factor graphs.
    Pdg, "PDG",
    req = [StochasticWires],
    allowed = [],
    extra = ok
);
\end{rustsrc}
\apiprov{construction --- \texttt{crates/mtk-catalogue/tests/builders.rs}}
\begin{rustsrc}
let pdg = MechanismBuilder::new()
    .stochastic_wire(StochasticWire::Conditional {
        id: 0,
        sources: vec![0],
        target: 1,
        rows: vec![CondRow {
            given: vec![0],
            outcome: 1,
            prob: r(1, 1),
        }],
    })
    .stochastic_wire(StochasticWire::Potential {
        id: 1,
        vars: vec![0, 1],
        rows: vec![PotRow {
            assignment: vec![0, 0],
            weight: r(7, 2),
        }],
    })
    .finish();
Pdg::manifest(pdg).unwrap();
\end{rustsrc}

\subsection{Bayesian Networks}
\label{sec:mt-bn}

\paragraph{The frame.} Directed acyclic dependency structure with
conditional probability tables; the bread-and-butter question is a
posterior, which is to say \emph{conditioning}.

\paragraph{Profile.} Conditional-table wires over finite variables;
the joint is the wires' product. Sum--product inference is semiring
tensor contraction \cite{peterseim26}, lowered to contraction rows (the same
contraction discipline the substrate reference chapter specifies),
and the well-formedness side condition on contractions turns out to
be the shape-level shadow of table normalization.

\paragraph{Semantics.} The joint distribution, exact at finite
carrier; the scoring semantics of the dependency-graph rendering
recovers it as the consistent case. The rendering's score is, in the
universe's terms, a belief (level-of-inconsistency functional) ---
Definition~\ref{def:su-penalty-functional} names this
incompatibility score as one of its two sources --- and the
consistent case is its zero locus.

\paragraph{Adequacy.} The compilation above is not merely
well-typed; it is \emph{adequate}: compiling and then denoting in
the semantic universe agrees with interpreting classically and then
including the classical object into the universe's precise face ---
the first such square off the time-line, with the topological order
of the variables playing the role the horizon plays for chains.
Stated self-contained, at the working scale of an exam problem:

\begin{problem}[the chain-rule square]
\label{prob:mt-bn}
Let $X_0, \dots, X_n$ be finitely many variables
with finite value domains, listed in a topological order: each
$X_i$ has a parent set $\mathrm{Pa}_i \subseteq \{X_j : j < i\}$,
and carries a rational conditional probability table
$P_i(\,\cdot \mid \mathrm{pa}\,)$, one row per parent assignment,
each row summing to one. On the record space $\Omega$ of joint
assignments write $\mu_B$ for the chain-rule joint,
\[
  \mu_B(\sigma) \;=\; \prod_{i \le n}
  P_i\big(\sigma_i \mid \sigma|_{\mathrm{Pa}_i}\big),
\]
the factorization semantics of Bayesian networks
\cite{koller-friedman-2009, pearl-1988}. Say a subprobability mass
function $\mu$ on $\Omega$ \emph{satisfies the compiled wires} when
(i) for every $\sigma$, the mass of
$\{\sigma' : \sigma'_0 = \sigma_0\}$ equals the sourceless first
wire's stored row $P_0(\sigma_0)$, an absent row read as $0$; and
(ii) for every later variable $i$ and every $\sigma$, writing
$C_i(\sigma)$ for the event that a record agrees with $\sigma$ at
all variables of index $< i$, refinement by one variable spends
exactly one stored row:
$\mu(C_{i+1}(\sigma)) = \mu(C_i(\sigma)) \cdot
P_i(\sigma_i \mid \sigma|_{\mathrm{Pa}_i})$, an absent row again
read as $0$ --- a constraint family whose multiplier reads only the
parent coordinates, i.e.\ the ordered Markov property with the
stored tables. Show that $\mu$ satisfies the compiled wires if and
only if $\mu = \mu_B$ --- in particular the constraints force total
mass one, so the ambient subprobability slack is never taken up.
\end{problem}

\emph{Sketch.} ($\Leftarrow$) Downward induction from full depth
gives the closed form $\mu_B(C_i(\sigma)) = \prod_{j < i}
P_j(\sigma_j \mid \sigma|_{\mathrm{Pa}_j})$: partition
$C_i(\sigma)$ by the value at variable $i$ and spend one row sum
per variable --- the multiplier is well-defined because parents sit
strictly below the refinement frontier. ($\Rightarrow$) Upward
induction along the order pins $\mu(C_i(\sigma))$ to the same
prefix products; at full depth the prefix events are the atoms
$\{\sigma\}$, so all weights agree, and summing the atoms puts the
total mass at one. Two structural remarks carry content: domain
nonemptiness is derived, not assumed --- the row-sum obligation at
a parentless root is over the (inhabited) empty parent assignment,
and normalization forces inhabited domains upward along the order
--- and the constraint family depends on the chosen topological
order while the law it pins does not, so every topological
presentation of one network denotes the same joint.

In the universe's terms: the constraint family, read as a formal
conjunction --- one atom per stored wire, in topological order ---
is the compiled system's denotation at the per-element stratum of
the credal stack, and the theorem says this loose arrow is
isomorphic, by an invertible globular cell of the thin double
category of per-element constraints, to the singleton inclusion of
$\mu_B$: the precise face of Remark~\ref{rem:su-faces}, whose ray
form agrees with the singleton on the radially maximal slice that
the verdict operators read. The pointwise equivalence is the chain
rule for Bayesian networks mechanized as a pinning statement. The
statement and both directions are mechanized, sorry-free, in the
toolkit's Lean tree (\texttt{lean/su-functor/}, module
\texttt{SuFunctor.Adequacy.Pgm}): the pointwise equivalence is
\texttt{satAtoms\_iff}, the invertible cell is
\texttt{pgm\_adequacy}, and the mass identity is
\texttt{jointLaw\_mass}. The same module ties the exact-rational
law to the standard monadic construction --- ancestral sampling,
the variables drawn in topological order by iterated
\texttt{PMF.bind} against their conditional distributions --- as
\texttt{ancestralPMF\_eq}, so the classical side of the square is
the community's mechanized object rather than a private
restatement; the inclusion vocabulary is shared with the sibling
frames of this catalogue
(\texttt{SuFunctor.Adequacy.Inclusions}). The square covers the
conditional-table face of the species --- exactly the
Bayesian-network sub-frame of this section; potential wires carry
no normalization row, so their inclusion routes through a global
normalization fact (the Gibbs distribution of the Markov-random-field
sub-frame below) and is deliberately a separate statement.

\paragraph{Status.} Certification, finite corner: exact-rational
posterior certificates with an independent checker and a mutant
battery, positivity of the evidence mass checked as a row rather
than assumed --- on the sprinkler network,
$P(\mathrm{Rain} \mid \mathrm{Wet}) = 509/719$ exactly.
Representation: the network as dependency-graph rows, with the
posterior recomputed from the rendering's own minimizer masses
\cite{richardson23} and accepted byte-identically by the unmodified
network checker. Open:
the continuous, score-carrying case --- certified conditioning
beyond finite enumeration is the named frontier this row shares with
section~\ref{sec:mt-ppl} --- and native ingestion.

\subsection{Structural Causal Models}
\label{sec:mt-scm}

\paragraph{The frame.} Structural equations over endogenous
variables driven by exogenous sources; the questions that make the
frame are interventional and \emph{counterfactual}.

\paragraph{Profile.} A shipped theory extending the value vocabulary:
equations as polynomial rows; an intervention as mechanism surgery
--- a partial replacement function over equations; and the coupling
of counterfactual worlds \emph{structural} rather than disciplinary:
one exogenous table per model, so the independent-worlds misreading
of a counterfactual is unrepresentable, not merely discouraged.

\paragraph{Semantics.} Twin-world substitution in the standard
(Pearl \cite{pearl-2009}) sense --- diagrammatic modern
form in Lorenz and Tull \cite{lorenz23} --- with the lightest bridge
set on the list:
substitution, box-measure monotonicity and additivity, and a union
bound --- no stochastic analysis anywhere. The bridge set is
machine-checked, sorry-free, as module \texttt{Qtsl.ScmTwin} in the
toolkit's Lean tree: twin-network substitution proven outright for
acyclic models, with the box-measure and union-bound lemmas stated
as the certificate classes the certifications below instantiate.

\paragraph{Status.} Certification, in both truth directions, on the
thesis's own flagship specification --- a bound on the probability of
counterfactually caused harm, at the thesis's stated $10^{-4}$
threshold: one controller certified compliant (with the ``would have
been harmed anyway'' cells genuinely exercised), a second certified
\emph{in violation} with a two-sided bracket $3/256 \le P \le
17/256$ --- inner boxes make violation certifiable, not merely
non-certifiable --- and a decoupling probe certifying that the same
compliant system under the independent-worlds product reading admits
$P \ge 1/2048$: the frame's defining distinction, priced by
certificates on both sides. Representation: the first instance rows
to represent interventions and counterfactual coupling at all. Open:
conditioning (counterfactuals given evidence --- the full abduction
step, joining the frontier named in section~\ref{sec:mt-bn}); latent
confounders and non-Markovian structure (representable,
unexercised); non-uniform exogenous densities; the ingestion gap;
closure --- for which the natural witness is a dependency-wire
coupling of an exogenous block to a dynamical frame.

\syscatmeta{\begin{mathpar}
\inferrule*[Right=\rulename{S-SCM}]{
  E : \overrightarrow{\langle \mathsf{var},\, \mathsf{rhs} \rangle}
  \text{ equations, each } \mathsf{rhs} \text{ polynomial} \\
  U : \mathsf{ExogenousTable} \\
  I? : \overrightarrow{\langle \mathsf{var},\,
  \mathsf{replacement}? \rangle} \text{ interventions (mechanism
  surgery)}
}{
  \mathrm{SCM}(E;\, U) : \mathsf{System}
}
\end{mathpar}
The exogenous table is an \emph{option, not a list} --- one per
mechanism, by the sort of $U$, which is the structural fact the
twin-world pattern of chapter~\ref{ch:qtsl-composition-patterns}
leans on.}
\ruleblock{API}
\sourcenote{Producer surface of the \texttt{mtk} toolkit for this model type: the TypeScript arrangement (npm package \texttt{@maitria/mtk}) and the Rust catalogue crate (\texttt{crates/mtk-catalogue}) --- one Rust implementation underneath, reached from TypeScript in-process through a native binding, so the two surfaces cannot diverge. Signatures and constructions are extracted from the toolkit sources at every regeneration of this catalogue; each construction is quoted verbatim from the toolkit's conformance corpus, code whose emitted certificate bytes are verified byte-identical against the committed golden vectors (helper values named \texttt{fx\_*} or \texttt{f.*} are that corpus's shared fixtures, defined beside the quoted lines). This API constructs and emits; it never checks --- the checker re-derives every claim independently.}
\apilang{TypeScript}{system.scm}
\apiprov{signature --- \texttt{ts/src/system.ts}}
\begin{tssrc}
/** S-SCM: polynomial equations + interventions + the shared exogenous
 * table (structural: one per mechanism, by type). */
export const scm = profileFactory("scm");
\end{tssrc}
\apiprov{The TS conformance corpus holds no construction through this factory at this preview --- the cross-language byte-identity fixtures exercise six profiles, and this one's byte authority is the Rust corpus opposite; the factory takes the same init-object vocabulary as every other profile.}
\apilang{Rust}{mtk\_catalogue::system::presets::Scm::manifest}
\apiprov{signature --- \texttt{crates/mtk-catalogue/src/system/presets.rs}}
\begin{rustsrc}
preset!(
    /// System Catalogue: SCM — equations as polynomial rows,
    /// interventions as mechanism surgery, and the shared exogenous
    /// table STRUCTURAL (one per mechanism, by type).
    Scm, "SCM",
    req = [Equations, Exogenous],
    allowed = [Interventions],
    extra = ok
);
\end{rustsrc}
\apiprov{construction --- \texttt{crates/mtk-catalogue/tests/builders.rs}}
\begin{rustsrc}
let scm = MechanismBuilder::new()
    .equation(0, mono(2, 1, 1, r(2, 1)))
    .intervention(InterventionRow {
        var: 0,
        replacement: Some(constant(2, r(0, 1))),
    })
    .exogenous(ExogenousTable {
        vars: vec![ExoVar {
            id: 1,
            name: Some("u".to_owned()),
            dist: Some(ExoDist::Discrete {
                rows: vec![
                    DiscreteWeight {
                        value: r(0, 1),
                        prob: r(1, 2),
                    },
                    DiscreteWeight {
                        value: r(1, 1),
                        prob: r(1, 2),
                    },
                ],
            }),
        }],
    })
    .finish();
Scm::manifest(scm).unwrap();
\end{rustsrc}

\subsection{Markov Random Fields}
\label{sec:mt-mrf}

\paragraph{The frame.} Undirected factor graphs with unnormalized
potentials; the bread-and-butter certified quantity is the partition
function.

\paragraph{Profile.} The dependency-graph rendering again, plus one
new atom: potential wires --- raw nonnegative tables with no
normalization row, their local scale derived from entries rather
than stored. The founding paper's own factor-graph embedding is
carried as instance rows.

\paragraph{Semantics.} The Gibbs distribution, recovered as the
rendering's optimum at unit weight; the bracket below is, in those
terms, a two-sided certificate on the optimal score --- cited to the
theorems of the founding paper and its loss-function sequel
\cite{richardson20, richardson22} as a named obligation.

\paragraph{Status.} Certification: a two-sided partition-function
bracket $\ell \le Z \le u$ on a small Ising model with all-rational
potentials. The lower end is an explicit support list checked
against the summed product directly --- the support-restricted
variational bound with every logarithm cancelled; the upper end is a
tree-reweighted (H\"older) forest decomposition checked as a single
rational inequality in which the fractional root never appears on
the wire. The checker never computes $Z$; test-side brute force confirms
the truth sits inside a bracket whose ends lie at $0.76\times$ and
$2.81\times$ the true value --- displayed as measured rather than
tuned. The factor-graph image is
generically \emph{inconsistent} as a dependency graph, so this
sub-frame and section~\ref{sec:mt-bn} together exercise both strata
of the rendering's semantics through one species --- checked
evidence that the unification is a unification, not an economy.
Open: contraction-row evaluation of the forest group sums; the
continuous case; the ingestion gap; closure.

\section{Probabilistic Programs}
\label{sec:mt-ppl}

\paragraph{The frame.} Corecursive programs in a functional
probabilistic programming language: \textsf{sample} and
\textsf{score}, loops, and the measures they denote.

\paragraph{Profile.} For the discrete-time slice: loops with
\textsf{sample} are outcome-row systems, and the published-table
campaign of section~\ref{sec:mt-dtmc} \emph{is} certified
probabilistic-program verification for that slice --- the swept
objects are programs, including programs with piecewise-linear
neural components. The general profile (higher-order corecursion,
continuous distributions --- the slice whose denotational home is
the $\omega$PAP semantics \cite{huot23}) is not yet rendered.

\paragraph{Semantics.} The program's measure semantics, two-tier:
the kernel language denotes in the affine probability monad --- the
monad fork \cite{dash22}, once merely named here, is resolved, and
the affine
tier's composition law is machine-checked, sorry-free, at the
finite-outcome stratum: with disjoint noise rosters, the denotation
of a sequential composite is the Kleisli composite of the stage
denotations (\texttt{noiseDenot\_seqComp} in module
\texttt{SuFunctor.Stochastic.NoiseKleisli} of the toolkit's Lean
tree), and the copy-vs-resample separation of
chapter~\ref{ch:qtsl-markov-kernels} is carried positively as
\texttt{copy\_ne\_kleisli}: a shared roster provably departs from
the Kleisli composite, so the disjointness is content, not
bookkeeping --- while score-carrying measures enter only at the
certificate boundary; conditioning denotes a normalized ratio whose
positivity is itself an obligation.

\paragraph{Status.} The forward, unconditioned slice is in place at
the discrete tier, as above --- and the former-level square is
machine-checked, sorry-free: a program's loop denotes stage kernels,
the unrolled program text denotes the kernel-local Kleisli iterate of
those kernels at every horizon (\texttt{noiseDenot\_progTraj},
unconditional at the finite-outcome stratum, every measurable stage
admitted), and for a program whose stage kernels have exact-rational
faces the semantic universe's judgment on the compiled outcome-row
presentation pins exactly the unrolled text's denotation
(\texttt{su\_pprog\_adequacy} in module
\texttt{SuFunctor.Adequacy.ProbProg}) --- so the formation rule
below sharing the chain former's data is a theorem-backed reading,
not an analogy. Conditioning --- posterior bounds for
the operation most probabilistic programs exist to perform --- is
the named frontier, problem statement in hand, shared with
section~\ref{sec:mt-bn}; the published yardsticks are the guaranteed
posterior bounds for universal programs \cite{beutner22} and for
discrete programs with loops \cite{zaiser24}. \emph{The large-language-model clause}: the
thesis's gloss (``including, notably, autoregressive large language
models'') gets a scope statement rather than a status. In the small,
the clause is subsumable today --- a toy autoregressive sampler is a
discrete-time chain over token prefixes, and
section~\ref{sec:mt-dtmc}'s machinery applies verbatim. In the sense
the thesis intends --- certified claims about the checkpoints people
actually run --- it is a programme-scale undertaking well beyond
this inventory, and this chapter declines to price it as a backlog
row. Stating that plainly is this entry's content.

\syscatmeta{\begin{mathpar}
\inferrule*[Right=\rulename{S-PPROG}]{
  O \text{ unkeyed outcome rows, as for } \mathrm{DTMC} \\
  \iota? \text{ initial-distribution rows} \\
  \tau? : \mathsf{TickDecl}
}{
  \mathrm{ProbProg}(O) : \mathsf{System}
}
\end{mathpar}
Deliberately the chain former's formation data verbatim: the
discrete-time slice of a probabilistic program \emph{is} an
outcome-row kernel, and what distinguishes this entry is its reading
(the mechanism presents a program's loop) plus the named open leg ---
program-as-rows, the corecursive source itself as first-class data,
which no premise above carries yet.}
\ruleblock{API}
\sourcenote{Producer surface of the \texttt{mtk} toolkit for this model type: the TypeScript arrangement (npm package \texttt{@maitria/mtk}) and the Rust catalogue crate (\texttt{crates/mtk-catalogue}) --- one Rust implementation underneath, reached from TypeScript in-process through a native binding, so the two surfaces cannot diverge. Signatures and constructions are extracted from the toolkit sources at every regeneration of this catalogue; each construction is quoted verbatim from the toolkit's conformance corpus, code whose emitted certificate bytes are verified byte-identical against the committed golden vectors (helper values named \texttt{fx\_*} or \texttt{f.*} are that corpus's shared fixtures, defined beside the quoted lines). This API constructs and emits; it never checks --- the checker re-derives every claim independently.}
\apilang{TypeScript}{system.probProg}
\apiprov{signature --- \texttt{ts/src/system.ts}}
\begin{tssrc}
/** S-PPROG: probabilistic program, discrete-time slice (a loop with
 * `sample` is a kernel, riding the outcome-row machinery). */
export const probProg = profileFactory("probProg");
\end{tssrc}
\apiprov{The TS conformance corpus holds no construction through this factory at this preview --- the cross-language byte-identity fixtures exercise six profiles, and this one's byte authority is the Rust corpus opposite; the factory takes the same init-object vocabulary as every other profile.}
\apilang{Rust}{mtk\_catalogue::system::presets::ProbProg::manifest}
\apiprov{signature --- \texttt{crates/mtk-catalogue/src/system/presets.rs}}
\begin{rustsrc}
preset!(
    /// System Catalogue: probabilistic program, discrete-time slice —
    /// a loop with `sample` is a kernel, riding the outcome-row
    /// machinery. Program-as-rows representation (the corecursive
    /// source as first-class data) is deliberately open; neural
    /// components
    /// ride the hypernet route.
    ProbProg, "probabilistic program",
    req = [Outcomes],
    allowed = [InitDist, Tick],
    extra = |m| outcomes_keyed("probabilistic program", m, false)
);
\end{rustsrc}
\apiprov{The discrete-time slice rides the outcome-row machinery; the corpus's construction of that machinery is the DTMC fixture, and no ProbProg-preset construction is in the conformance corpus at this preview.}

\section{Probabilistic Logic Programs}
\label{sec:mt-problog}

\paragraph{The frame.} Definite clauses plus probabilistic facts;
the distribution semantics --- a measure over least models induced
by independent fact coins --- and certified brackets on a query's
success probability.

\paragraph{Profile.} Intended: probabilistic facts and clauses as
rows over the substrate's own chase --- which is already a
geometric, Horn-clause engine --- with derivation weights in a
declared ordered-semiring algebra \cite{matache23}.

\paragraph{Semantics.} Sato-style distribution semantics. What the
frame adds beyond weighted derivations is the \emph{independence
structure} of the fact coins: derivations share coin outcomes, so
query mass is a sum of products with shared variables, not a free
semiring sum --- and exactly there, certified \emph{bounds} rather
than exact weighted model counting become the achievable target.

\paragraph{Status.} The frame integration --- facts and clauses as
rows over the substrate's own chase --- remains unshipped, and the
frame is graded absent on exactly that line: adjacent code, however
promising, does not lift a frame out of absent (that line is what
makes these statuses checkable). What the toolkit does carry is the
stratum below: the declaration profile (the API block after the
metanotation) and, over it, a producer-side certificate species ---
\texttt{problog-v0}, a certified bracket on a query's success
probability computed from producer-found explanation and cut rows
--- whose own API block closes this section. The structural
proximity noted above stands: the weight-algebra code exists, and
wiring it into the chase is a named seam of the substrate rather
than a research question.

\syscatmeta{\begin{mathpar}
\inferrule*[Right=\rulename{S-PROBLOG}]{
  \mathsf{semiring} \text{ a declared ordered semiring} \\
  F : \overrightarrow{\langle \mathsf{atom},\, \mathsf{weight},\,
  \mathsf{coin}? \rangle} \text{ weighted facts} \\
  \mathsf{program\_digest}? \text{ a content-addressed program
  reference}
}{
  \mathrm{ProbLog}(\mathsf{semiring};\, F) : \mathsf{System}
}
\end{mathpar}
The coin index is the independence structure: derivations sharing a
coin share its outcome, which is exactly what makes bounds rather
than free semiring sums the achievable target.}
\ruleblock{API}
\sourcenote{Producer surface of the \texttt{mtk} toolkit for this model type: the TypeScript arrangement (npm package \texttt{@maitria/mtk}) and the Rust catalogue crate (\texttt{crates/mtk-catalogue}) --- one Rust implementation underneath, reached from TypeScript in-process through a native binding, so the two surfaces cannot diverge. Signatures and constructions are extracted from the toolkit sources at every regeneration of this catalogue; each construction is quoted verbatim from the toolkit's conformance corpus, code whose emitted certificate bytes are verified byte-identical against the committed golden vectors (helper values named \texttt{fx\_*} or \texttt{f.*} are that corpus's shared fixtures, defined beside the quoted lines). This API constructs and emits; it never checks --- the checker re-derives every claim independently.}
\apilang{TypeScript}{system.problog}
\apiprov{signature --- \texttt{ts/src/system.ts}}
\begin{tssrc}
/** S-PROBLOG: weighted logic programs over a declared ordered
 * semiring. */
export const problog = profileFactory("problog");
\end{tssrc}
\apiprov{construction --- \texttt{ts/tests/problog-native.test.ts}}
\begin{tssrc}
return sys.problog({
  weightedProgram: {
    semiring: "probability",
    facts,
    programDigest: sha256(SMOKERS_SOURCE),
  },
});
\end{tssrc}
\apiprov{\texttt{facts} are the smokers fixture's weighted rows (stress and influence facts sharing declared coins), built beside the quoted lines; the system constructed here feeds the certificate corpus's \texttt{fromSystem} path, whose emitted bytes are verified against the committed golden vectors.}
\apilang{Rust}{mtk\_catalogue::system::presets::ProbLog::manifest}
\apiprov{signature --- \texttt{crates/mtk-catalogue/src/system/presets.rs}}
\begin{rustsrc}
preset!(
    /// System Catalogue: ProbLog / weighted logic programs — the
    /// chase over a declared ordered-semiring weight algebra +
    /// independence structure of fact coins.
    ProbLog, "ProbLog",
    req = [WeightedProgram],
    allowed = [],
    extra = ok
);
\end{rustsrc}
\apiprov{construction --- \texttt{crates/mtk-catalogue/tests/builders.rs}}
\begin{rustsrc}
let problog = MechanismBuilder::new()
    .weighted_program(WeightedProgramDecl {
        semiring: "probability".to_owned(),
        facts: vec![WeightedFact {
            atom: "edge(a,b)".to_owned(),
            weight: r(3, 5),
            coin: Some(0),
            coin_domain: None,
        }],
        program_digest: None,
    })
    .finish();
ProbLog::manifest(problog).unwrap();
\end{rustsrc}
 \ruleblock{API}
\sourcenote{Producer surface of the \texttt{mtk} toolkit for this certificate species: the TypeScript arrangement (npm package \texttt{@maitria/mtk}) and the Rust catalogue crate (\texttt{crates/mtk-catalogue}) --- one Rust implementation underneath, reached from TypeScript in-process through a native binding, so the two surfaces cannot diverge. Signatures and constructions are extracted from the toolkit sources at every regeneration of this catalogue; each construction is quoted verbatim from the toolkit's conformance corpus, code whose emitted certificate bytes are verified byte-identical against the committed golden vectors (helper values named \texttt{fx\_*} or \texttt{f.*} are that corpus's shared fixtures, defined beside the quoted lines). This API constructs and emits; it never checks --- the checker re-derives every claim independently.}
The claims stratum over the \texttt{S-PROBLOG} declaration: a \texttt{problog-v0} certificate carries the ground program (coin, fact, and clause rows), the producer's search artifacts (explanation and cut rows), and the claimed bracket on the query's success probability --- lower by inclusion--exclusion over coin-set unions, upper by min-over-cuts, exact by world enumeration where the coin count permits. The builder verifies every row it is handed (an explanation must derive the query on its own; a cut's joint absence must block it against every other coin) and computes the bracket itself; where the rows came from is the producer's business, and the checker re-derives every claim from the certificate bytes alone.
\apilang{TypeScript}{problog.builder / problog.fromSystem}
\apiprov{signature --- \texttt{ts/src/problog.ts}}
\begin{tssrc}
/** Start a `problog-v0` certificate for `query` in the probability
 * semiring (declare the ground program with the chainable methods,
 * derive claims with `claimsFromRows`, then `build`). */
export function builder(query: string): CertBuilder {
  return new CertBuilder(new native.PlogCertBuilder(query));
}

/** Derive a certificate builder from a System whose manifest carries
 * the `weightedProgram` family (the S-PROBLOG profile's declaration):
 * coin probabilities come from fact weights (facts sharing a coin
 * must agree — refused with both fact indices otherwise), coinless
 * facts are deterministic and must carry weight 1, and the
 * declaration's program digest is required — it is what pins the
 * grounding correspondence. A System without the family throws
 * (programmer error — wrong manifest kind). */
export function fromSystem(
  system: System,
  clauses: ReadonlyArray<ClauseRow>,
  query: string,
): Result<CertBuilder, Refusal> {
  const r = native.problogBuilderFromManifest(
    system.nativeHandle,
    clausesToNative(clauses),
    query,
  );
  return isRefused(r)
    ? refusedResult(narrowRefusal(r.refusal))
    : ok(new CertBuilder(r));
}
\end{tssrc}
\apiprov{construction --- \texttt{ts/tests/problog-native.test.ts}}
\begin{tssrc}
const b = smokersBuilder(problog, sys);
const claims = problog.unwrap(b.claimsFromRows(ANTICHAIN, CUTS, true));
assert.equal(claims.exact?.toString(), "48861/125000");
assert.equal(claims.lower.toString(), "48861/125000");
assert.equal(claims.upper.toString(), "69/125");
const cert = problog.unwrap(b.build(claims));
\end{tssrc}
\apiprov{\texttt{smokersBuilder} and its clause fixture are defined beside the quoted lines; the same file's die-chain exemplar exercises the direct-builder path (\texttt{builder}, \texttt{mswCoin} switch domains) against its own golden vector.}
\apilang{Rust}{mtk\_catalogue::system::problog::ProblogCertBuilder}
\apiprov{signature --- \texttt{crates/mtk-catalogue/src/system/problog.rs}}
\begin{rustsrc}
pub fn new(query: impl Into<String>) -> Self

pub fn from_decl(
    decl: &WeightedProgramDecl,
    clauses: Vec<GroundClauseRow>,
    query: impl Into<String>,
) -> Result<Self, ProblogRefusal>

pub fn claims_from_rows(
    &self,
    explanations: Vec<Vec<u32>>,
    cuts: Vec<Vec<u32>>,
    with_exact: bool,
) -> Result<ProblogClaims, ProblogRefusal>

pub fn build(self, claims: ProblogClaims) -> Result<ProblogCert, ProblogRefusal>
\end{rustsrc}
\apiprov{construction --- \texttt{crates/mtk-catalogue/tests/problog\_builders.rs}}
\begin{rustsrc}
// Manifest stratum: the S-PROBLOG profile accepts the declaration.
let decl = smokers_decl();
let manifest = MechanismBuilder::new()
    .weighted_program(decl.clone())
    .finish();
ProbLog::manifest(manifest).unwrap();

// Claims stratum: ground the program, hand over the search
// artifacts (explanation and cut rows), and let the builder verify
// them and compute the bracket — lower by inclusion–exclusion over
// coin-set unions, upper by min-over-cuts, exact by enumeration.
let builder = ProblogCertBuilder::from_decl(&decl, smokers_clauses(), "smokes(a)").unwrap();
let claims = builder.claims_from_rows(antichain(), cuts(), true).unwrap();
let cert = builder.build(claims).unwrap();
let bytes = cert.to_cbor().unwrap();
\end{rustsrc}
\apiprov{\texttt{smokers\_decl}, \texttt{smokers\_clauses}, \texttt{antichain}, and \texttt{cuts} are the corpus's shared fixtures (\texttt{tests/problog\_common}), each re-derived by the reference-number tests beside the quoted construction.}
 \ruleblock{API}
\sourcenote{Producer surface of the \texttt{mtk} toolkit for this certificate stratum: the TypeScript arrangement (npm package \texttt{@maitria/mtk}) and the Rust catalogue crate (\texttt{crates/mtk-catalogue}) --- one Rust implementation underneath, reached from TypeScript in-process through a native binding, so the two surfaces cannot diverge. Signatures and constructions are extracted from the toolkit sources at every regeneration of this catalogue; each construction is quoted verbatim from the toolkit's conformance corpus, code whose emitted certificate bytes are verified byte-identical against the committed golden vectors (helper values named \texttt{fx\_*} or \texttt{f.*} are that corpus's shared fixtures, defined beside the quoted lines). This API constructs and emits; it never checks --- the checker re-derives every claim independently.}
The exact-beyond-enumeration stratum of the \texttt{problog-v0} certificate: past the world-enumeration cap, the producer commits its search as a leaf-cited decision tree --- explanation and cut rows compiled so that every path through the coin axes closes at the first row that decides it, an undecidable cube refusing with the world as witness --- and the exact probability is re-derived by the same verify-and-fold walk the checker runs, so the claim never rests on the producer's arithmetic. Both extensions of the wire (the compiled section and the even-depth Bonferroni truncation obligation) are optional and skip-if-absent: every earlier certificate's bytes decode unchanged.
\apilang{TypeScript}{problog.CertBuilder.compiledFromRows}
\apiprov{signature --- \texttt{ts/src/problog.ts}}
\begin{tssrc}
/** Compile the full compiled section (exact-beyond-enumeration)
 * from producer-found rows: verifies each row, builds the
 * leaf-cited decision tree — refusing, with the undecided world as
 * witness, when the rows cannot decide every cube — and computes
 * `exact` by the same verify-and-fold walk the preflight and the
 * checker run. Attach the result as `claims.compiled` before
 * `build`. */
compiledFromRows(
  explanations: ReadonlyArray<ReadonlyArray<number>>,
  cuts: ReadonlyArray<ReadonlyArray<number>>,
): Result<CompiledSection, Refusal> {
  const r = this.h.compiledFromRows(
    explanations.map((e) => [...e]),
    cuts.map((e) => [...e]),
  );
  return isRefused(r)
    ? refusedResult(narrowRefusal(r.refusal))
    : ok(compiledFromNative(r));
}
\end{tssrc}
\apiprov{construction --- \texttt{ts/tests/problog-native.test.ts}}
\begin{tssrc}
// The compiled section: producer-found rows (an explanation
// antichain plus cut rows covering every reject cube) compile to a
// leaf-cited decision tree whose exact-rational fold certifies the
// exact query probability with no world-count cap.
const b = dieBuilder(problog);
const compiled = problog.unwrap(
  b.compiledFromRows(DIE_ANTICHAIN, DIE_COMPILED_CUTS),
);
const claims = problog.unwrap(b.claimsFromRows(DIE_ANTICHAIN, DIE_CUTS, true));
const cert = problog.unwrap(b.build({ ...claims, compiled }));
\end{tssrc}
\apiprov{\texttt{dieBuilder} and the row fixtures are defined beside the quoted lines; the byte-identity assertion against the committed compiled golden sits directly under the marked construction.}
\apilang{Rust}{mtk\_catalogue::system::problog::ProblogCertBuilder::compiled\_from\_rows}
\apiprov{signature --- \texttt{crates/mtk-catalogue/src/system/problog.rs}}
\begin{rustsrc}
pub fn compiled_from_rows(
    &self,
    explanations: Vec<Vec<u32>>,
    cuts: Vec<Vec<u32>>,
) -> Result<CompiledSection, ProblogRefusal>
\end{rustsrc}
\apiprov{construction --- \texttt{crates/mtk-catalogue/tests/problog\_compiled.rs}}
\begin{rustsrc}
// The compiled section: producer-found rows (an explanation
// antichain plus cut rows covering every reject cube) compile to
// a leaf-cited decision tree whose exact-rational fold certifies
// the exact query probability with no world-count cap.
let builder = problog_common::die_builder();
let compiled = builder
    .compiled_from_rows(
        problog_common::die_antichain(),
        problog_common::die_compiled_cuts(),
    )
    .unwrap();
let mut claims = builder
    .claims_from_rows(
        problog_common::die_antichain(),
        problog_common::die_cuts(),
        true,
    )
    .unwrap();
claims.compiled = Some(compiled);
let cert = builder.build(claims).unwrap();
\end{rustsrc}
\apiprov{the shared fixtures (\texttt{problog\_common}) are the same corpus the claims-stratum block quotes; the hand-built switch-first tree beside the quoted test verifies the same certificate rows against an independent tree layout.}

\section{Clocked Boolean Circuits}
\label{sec:mt-circuit}

\paragraph{The frame.} Synchronous gate-level hardware
(and-inverter graphs) as a first-class model type --- one of two
entries in production use that fall outside the 2023 roster,
evidence the inventory is a living catalogue rather than homework
against a fixed list.

\paragraph{Profile.} Circuit content by digest --- the ruled
content-addressing idiom; gate internals stay un-row-ified --- plus
boot-prefix rows and per-lane input-configuration rows. Disturbance:
demonic inputs as per-lane branch covers. Time: counter, on the
machine clock, with a frame counter above it.

\paragraph{Semantics.} The special case where the checker \emph{is}
the semantics: exact replay equivalence, bit for bit; claims read
the credal-lower return (Definition~\ref{def:su-previsions}) over
the declared branch cover.

\paragraph{Status.} Certification is the heavy leg: this is the type
underneath the hardware-liveness and console campaigns of
sections~\ref{sec:exp-circuits} and~\ref{sec:exp-atari}, where
certified minimum-return claims run over branch covers at $10^7$
ticks. Representation: declaration rows shipped (digest, boot
prefix, input configurations). Open: the ingestion gap; closure.

\syscatmeta{\begin{mathpar}
\inferrule*[Right=\rulename{S-CIRCUIT}]{
  \mathsf{digest} \text{ a content-addressed circuit reference} \\
  B : \overrightarrow{\langle \mathsf{tick},\, \mathsf{wire},\,
  \mathsf{value} \rangle} \text{ boot-prefix rows} \\
  C : \overrightarrow{\langle \mathsf{lane},\, \mathsf{wire},\,
  \mathsf{value} \rangle} \text{ input-configuration rows} \\
  \tau? : \mathsf{TickDecl}
}{
  \mathrm{Circuit}(\mathsf{digest}) : \mathsf{System}
}
\end{mathpar}
Gate internals stay un-row-ified --- content addressing is the ruled
idiom, and the digest names the whole gate graph.}
\ruleblock{API}
\sourcenote{Producer surface of the \texttt{mtk} toolkit for this model type: the TypeScript arrangement (npm package \texttt{@maitria/mtk}) and the Rust catalogue crate (\texttt{crates/mtk-catalogue}) --- one Rust implementation underneath, reached from TypeScript in-process through a native binding, so the two surfaces cannot diverge. Signatures and constructions are extracted from the toolkit sources at every regeneration of this catalogue; each construction is quoted verbatim from the toolkit's conformance corpus, code whose emitted certificate bytes are verified byte-identical against the committed golden vectors (helper values named \texttt{fx\_*} or \texttt{f.*} are that corpus's shared fixtures, defined beside the quoted lines). This API constructs and emits; it never checks --- the checker re-derives every claim independently.}
\apilang{TypeScript}{system.circuit}
\apiprov{signature --- \texttt{ts/src/system.ts}}
\begin{tssrc}
/** S-CIRCUIT: clocked boolean circuit (AIG by digest; exact replay IS
 * the semantics). */
export const circuit = profileFactory("circuit");
\end{tssrc}
\apiprov{The TS conformance corpus holds no construction through this factory at this preview --- the cross-language byte-identity fixtures exercise six profiles, and this one's byte authority is the Rust corpus opposite; the factory takes the same init-object vocabulary as every other profile.}
\apilang{Rust}{mtk\_catalogue::system::presets::Circuit::manifest}
\apiprov{signature --- \texttt{crates/mtk-catalogue/src/system/presets.rs}}
\begin{rustsrc}
preset!(
    /// System Catalogue: clocked boolean circuit (AIG by digest) —
    /// content-addressed circuit + boot-prefix rows +
    /// input-configuration rows; exact replay IS the semantics.
    Circuit, "clocked circuit",
    req = [Circuit],
    allowed = [Tick],
    extra = ok
);
\end{rustsrc}
\apiprov{construction --- \texttt{crates/mtk-catalogue/tests/builders.rs}}
\begin{rustsrc}
let circuit = MechanismBuilder::new()
    .circuit(CircuitDecl {
        digest: vec![0xAB; 32],
        boot_prefix: vec![BootRow {
            tick: 0,
            wire: 0,
            value: true,
        }],
        input_configs: vec![InputConfigRow {
            lane: 0,
            wire: 1,
            value: false,
        }],
    })
    .finish();
Circuit::manifest(circuit).unwrap();
\end{rustsrc}

\section{The i.i.d.\ Tier}
\label{sec:mt-iid}

\paragraph{The entry.} Not a profile --- a claim class that every
sampleable model type carries, and the second production entry
beyond the 2023 list. Where a system's rollouts can be executed but
no analytic species reaches its semantics --- a compiled physics
engine is the paradigm case --- the catalogue switches to
\emph{confidence-qualified} judgments: claims of the form ``the
expected return of this pinned artefact is at least $B$, with
confidence $1-\alpha$'' (species U5 of section~\ref{sec:exp-claims},
always displayed as their own class, never mixed into a table of
exact rows). The tier's name states its subject: a pinned,
\emph{closed} composite, sampled under its own law episode after
episode --- independent and identically distributed draws of the
artefact itself. This is the tier behind the physics-engine results
of section~\ref{sec:exp-gym}, and it composes with the profile
inventory rather than competing with it: the analytic tiers take
over exactly where a frame's rows stop being samples and start being
algebra.

\paragraph{The claim form is more general than it looks.} The
expectation floor on a bounded per-episode output is the tier's
normal form, and claims that sound set-shaped compile into it. ``The
next landing stays inside the pad region with probability at least
$0.99$'' \emph{is} the statement that the mean of the landing
indicator is at least $0.99$; a quantile bound is a coverage bound
at a fixed threshold, hence again an indicator mean; and a bound on
a conditional value-at-risk is an expectation floor on a derived
output --- exhibiting any threshold $z$ with
$z + \mathbb{E}[(X-z)^{+}]/(1-\beta) \le c$ certifies
$\mathrm{CVaR}_{\beta}(X) \le c$, by the variational form of
Rockafellar and Uryasev \cite{rockafellar00}. What the producer must
arrange is that the property's verdict or score be computed
\emph{inside} the sampled composite --- the property logic wired in
as one more component, its terminal output the episode's return ---
and every claim above is then one expectation floor, carried by one
certificate species.

\paragraph{The guarantee class.} The certificates are anytime-valid
lower confidence bounds in the game-theoretic school of sequential
inference \cite{ramdas22}: a betting martingale --- concretely, the
hedged-capital process for bounded means of Waudby-Smith and Ramdas
\cite{waudbysmith20} --- is run over the recorded rollouts, the
certificate is the resulting wealth ledger, sound by Ville's
inequality \cite{ville1939}, and the checker replays the entire
ledger in directed-rounding fixed point. This is not a new trust
principle beside the rest of the book but the same one in a second
mode: the supermartingale certificates of the analytic tiers
\emph{prove} the Ville bound against a known kernel, while this
tier's ledger \emph{wagers} it against whatever law would violate
the claim --- the wealth process is a nonnegative supermartingale
under every violating law by construction of the bets, so observed
wealth crossing $1/\alpha$ rejects every such law at confidence
$1-\alpha$, at any stopping time, with no penalty for peeking and
none for extending the ledger. That anytime validity is what makes
the certificate a stable public artefact rather than a frozen
experiment, and the betting form is what an independent checker can
re-verify line by line; sound statistical model checking
\cite{budde24} is the nearest published genre.

\paragraph{The conformal sibling.} Conformal prediction
\cite{vovk05} --- lately including accept/reject guarantees for
generative samplers under the name \emph{conformal sampling}
\cite{li25} --- is the same distribution-free finite-sample school
under a \emph{weaker} assumption, exchangeability of the draws, with
a different validity mechanism: rank statistics of a completed
calibration batch, exact at a calibration size fixed in advance. The
relation is precise. At the level of claims, a conformal coverage
guarantee is an indicator-mean floor, so under this tier's stronger
sampling assumption the ledger certifies the same statements ---
anytime-valid, and conditionally on the fitted predictor rather than
on average over calibration draws. At the level of mechanisms the
two are genuinely different: no adapted betting process can be
simultaneously powerful under exchangeability alone --- for binary
sequences, a nonnegative supermartingale valid under every
exchangeable law stays a supermartingale over their closed
fork-convex hull, which contains every law on such sequences
\cite{ramdas21} --- and this tier does not attempt that regime. It does not need to: its
subject is a pinned closed artefact whose independent resampling is
realized by the harness and attested on the certificate, not assumed
of nature --- the situation in which the stronger assumption costs
nothing.

\paragraph{Pinning.} A confidence bound on ``the expected return of
this artefact'' means something only if there is one fixed law being
sampled, and pinning is that premise made auditable --- three
ordinary ingredients rather than one new trust primitive. First,
closure: the sampled composite has no dangling inputs, because an
open system has no episode law to sample. Second, identity: trust
rows name the sampled artefact by content digest --- the policy and
its determinism --- the same content-addressing idiom the circuit
former rules. Third, the one genuinely statistical residue, the
sampling-law witness: the attestation that the deployed harness
draws independent episodes of the digest-named artefact, with the
seed provenance of every source of randomness as its evidence. All
three are stated on the face of the certificate, never hidden;
ledger production itself lives with the sampling harness rather than
in these rows, and the almost-sure return range the martingale
requires is itself a certified lemma about the environment's code,
not an assumption (Table~\ref{tab:exp-ll}).

\syscatmeta{\begin{mathpar}
\inferrule*[Right=\rulename{S-STAT}]{
  \mathsf{rows} : \overrightarrow{\langle \mathsf{kind},\,
  \mathsf{detail}?,\, \mathsf{digest}? \rangle} \text{ trust rows}
}{
  \langle \mathsf{rows} \rangle : \mathsf{PinningDecl}
}
\end{mathpar}
Not a System former at this preview: $\mathsf{pinning} \in \Sigma$,
so the family is substrate, admitted alongside every former above,
and the confidence-qualified judgment applies to any sampleable type
whose subject is pinned. The catalogue's design direction is for the
sampled subject itself to become a first-class construction --- a
replication former on closed composites, in the Bag former's genus,
whose formation premise is the closure ingredient above and whose
trust seam is exactly these rows; it enters the formal inventory
with its producer surface, not before.}
\ruleblock{API}
\sourcenote{Producer surface of the \texttt{mtk} toolkit for this entry: the TypeScript arrangement (npm package \texttt{@maitria/mtk}) and the Rust catalogue crate (\texttt{crates/mtk-catalogue}) --- one Rust implementation underneath, reached from TypeScript in-process through a native binding, so the two surfaces cannot diverge. Signatures and constructions are extracted from the toolkit sources at every regeneration of this catalogue; each construction is quoted verbatim from the toolkit's conformance corpus, code whose emitted certificate bytes are verified byte-identical against the committed golden vectors (helper values named \texttt{fx\_*} or \texttt{f.*} are that corpus's shared fixtures, defined beside the quoted lines). This API constructs and emits; it never checks --- the checker re-derives every claim independently.}
\apilang{TypeScript}{system.PinningDecl}
\apiprov{signature --- \texttt{ts/src/system.ts}}
\begin{tssrc}
export interface PinningDecl {
  readonly rows: ReadonlyArray<TrustRow>;
}

export interface TrustRow {
  readonly kind: string;
  readonly detail?: string;
  readonly digest?: string;
}
\end{tssrc}
\apiprov{Pinning is deliberately not a profile: any sampleable mechanism carries it, and the corpus holds no pinned construction at this preview --- ledger production lives with the sampling harness.}
\apilang{Rust}{mtk\_catalogue::system::MechanismBuilder::pinning}
\apiprov{signature --- \texttt{crates/mtk-catalogue/src/system/manifest.rs}}
\begin{rustsrc}
pub fn pinning(mut self, p: PinningDecl) -> Self
\end{rustsrc}
\apiprov{Pinning is deliberately not a preset (any sampleable profile carries the family); the corpus holds no pinned construction at this preview --- ledger production lives with the sampling harness.}

\section{Quantum Nets}
\label{sec:mt-qnet}

\paragraph{The frame.} Algebraic quantum field theory in its locally
covariant formulation, at the kinematic fragment whose mechanization
chapter~\ref{ch:qtsl-semantic-universe} presents
(section~\ref{sec:su-site}): observables indexed by causal diamonds,
as they have been since Haag and Kastler \cite{haagkastler64}; an
isotone net of local algebras with Einstein causality and the
time-slice property, per Brunetti, Fredenhagen, and Verch
\cite{brunetti01}; and a positive record family whose Born rows are
the measured content. The grade comes first, because a reader who
hears ``locally covariant quantum field theory'' will expect
type-III von Neumann algebras, microlocal spectrum conditions, and
the continuum. This entry works at a deliberately finite stratum ---
a finite site, matrix algebras, the kinematic axioms, exact
arithmetic --- and ``fragment'' does the same work in this section
that it does in the Lean names. Genuine dynamics (relative Cauchy
evolution, the stress-energy tensor) lives on the far side of the
fragment's fence, which is where the locally covariant literature
itself places it \cite{fewster15}, and Fewster and Verch proved that
``locally covariant'' does not by itself mean ``the same physics in
all spacetimes'' \cite{fewster11}. The bread-and-butter question is
a certified bound on record probabilities holding uniformly over
every admissible state, and the uniformity is no modelling flourish:
the credal face below inherits it from the community's own state
doctrine. A two-qubit example runs through the section, with records
the four projectors of the joint computational-basis measurement, so
that the record index is a pair of bits: the Bell state prices the
row $(\tfrac12, 0, 0, \tfrac12)$, and the maximally mixed state
prices $(\tfrac14, \tfrac14, \tfrac14, \tfrac14)$.

\paragraph{Profile.} A site kit --- finitely many points with their
chronology; diamonds, covers, and a distinguished ambient window are
derived --- with a per-window algebra presentation at matrix grade,
generating families entered in $\mathbb{Q}[i]$; matter data
(worldlines, jointly exhaustive factors, a reversible mechanism),
from which Cauchy determinacy follows as a theorem rather than a
postulate; a declared record family at the ambient window ---
finitely many positive-semidefinite matrices summing to one, a
positive operator-valued measure --- and a state face, declaring
either one distinguished state (the precise face) or a set of states
(the credal face), each state carried with the exact-rational
witness of its Born row. The carrier is the record space of the
declared family, never the algebra: the $C^*$-algebra is
presentation-tier data, as a Lagrangian is, and what the semantic
universe consumes is its gauge-invariant shadow --- probability rows
over record outcomes. Initial conditions live in the state face.

\paragraph{Sited time.} Every counter-time entry in this catalogue
composes into real time through a tick declaration
(section~\ref{sec:mt-tick}). This entry needs none: its windows are
the site's own causal diamonds, so the time semantics is the causal
order itself --- the ground storey of the semantic universe
(section~\ref{sec:su-causal-order}) --- and a separate tick would
have nothing to declare. Among the values the catalogue's
time-semantics column has needed so far (continuous, counter,
hybrid), this one is new: sited.

\paragraph{Semantics.} The record law of Gell-Mann and Hartle
\cite{gellmann92}, which at a single window is the Born rule: a
state $\rho$ prices the record $\alpha$ at
$\mathrm{Re}\,\mathrm{tr}(\rho R_\alpha)$. The mechanized form of
that law is the fragment's own record shadow, so the bridge between
the community's formula and the toolkit's is definitional rather
than analytic; the connector lemma below is one line. Two doctrines
from the locally covariant literature enter as structure rather than
gloss. First, the semantics factors through the state object: two
states with the same restrictions at the ambient window price every
record identically, so the square's left edge is a function of the
state-presheaf fibre, as \S5.1 of Brunetti--Fredenhagen--Verch
demands \cite{brunetti01}. Second, the state object is set-valued
from birth: a functorial choice of single states is unrealizable ---
a theorem of the framework \cite{brunetti01}, re-proven
model-independently under dynamical locality \cite{fewster11} ---
and the replacement that literature adopted, the folium, is
structurally the flat shadow of this book's credal bodies
(section~\ref{sec:su-site} tells that story in full). The
well-formedness hypotheses the squares consume are the community's
own axioms: isotony, Einstein causality, the time-slice property at
its two tiers, record positivity and exhaustiveness, state
positivity and normalization.

\paragraph{Adequacy (precise face).} The compilation stores the Born
row and nothing else, so the square is the catalogue's leanest, and
its textbook corner is unlike any other in the roster: every other
entry's community-facing corner was built for its square, while this
one consumes the mechanized fragment verbatim, because the
fragment's fields already are the community's axioms. Stated
self-contained, at exam scale:

\begin{problem}[the quantum-net square]
\label{prob:mt-qnet}
Let $\kappa$ be a finite record index, $R \colon \kappa \to
M_n(\mathbb{C})$ a family of positive-semidefinite matrices with
$\sum_\alpha R_\alpha = 1$, and $\rho \in M_n(\mathbb{C})$ a state:
positive semidefinite, of trace one. Store the Born row as exact
rationals, $q \colon \kappa \to \mathbb{Q}$ with $q_\alpha =
\mathrm{Re}\,\mathrm{tr}(\rho R_\alpha)$, and say a subprobability
mass function $\mu$ on $\kappa$ \emph{satisfies the compiled rows}
when $\mu\{\alpha\} = q_\alpha$ for every $\alpha$. Show: (i) $q$ is
a probability row --- every entry nonnegative, total mass one; (ii)
$\mu$ satisfies the compiled rows if and only if $q$ is its weight
function, and the uniqueness direction consumes neither an induction
nor the exhaustiveness of $R$; (iii) storing only a subfamily $J
\subseteq \kappa$ yields a row of mass $\sum_{\alpha \in J}
q_\alpha$ --- exhibit an instance of mass exactly $\tfrac34$; (iv)
positivity cannot be dropped: exhibit a Hermitian family summing to
one, and a state, that together price an outcome negatively.
\end{problem}

\emph{Sketch.} (i) The trace of a product of positive-semidefinite
matrices is real and nonnegative, and linearity spends the two
normalizations: $\sum_\alpha \mathrm{Re}\,\mathrm{tr}(\rho R_\alpha)
= \mathrm{Re}\,\mathrm{tr}(\rho) = 1$. (ii) The singleton probes
read off every weight directly, and two subprobability mass
functions with equal weights are equal --- finite extensionality and
nothing more. Compare the chain square (Problem~\ref{prob:mt-dtmc}),
whose uniqueness runs an induction along the horizon and spends a
row sum per step; exhaustiveness enters here only through the mass
claim of (i). (iii) Sum the stored subfamily; for the instance, take
the maximally mixed two-qubit state against three of the four basis
projectors: $3 \cdot \tfrac14 = \tfrac34$. (iv) One dimension
suffices: the family $(2, -1)$ on $M_1(\mathbb{C})$ sums to one, and
the unique state prices the second outcome at $-1$.

In the universe's terms: the compiled rows denote one
probability-pinning constraint per record index, their formal
conjunction a loose arrow at the per-element stratum of the credal
stack, and the theorem exhibits an invertible globular cell to the
singleton inclusion of the Born row --- the precise face of
Remark~\ref{rem:su-faces}. The statement and both directions are
mechanized, sorry-free, in the toolkit's Lean tree
(\texttt{lean/su-functor/}, module
\texttt{SuFunctor.Adequacy.QuantumNet}): the invertible cell is
\texttt{quantum\_net\_adequacy}, the pointwise equivalence is
\texttt{satAtoms\_iff}, and the presentation structure
(\texttt{QNetSystem}) consumes the fragment (\texttt{LCQFTFragment},
module \texttt{QtslAux.QuantumAssembly}) verbatim, together
with a distinguished state and its exact-rational row witness. The
community tie is cheaper here than anywhere else in the catalogue,
and the entry says why rather than dressing it as more: the
connector \texttt{shadowPMF\_prob\_eq\_born} is the row witness
itself --- the mechanized record shadow is the literature's formula
transcribed --- and the mechanized probability library has no
quantum machinery for a deeper tie to borrow, where the chain square
could tie into the library's own monadic construction. What replaces
depth is doctrine: \texttt{shadowPMF\_eq\_fragment\_shadow} and
\texttt{shadowPMF\_class\_invariant} factor the left edge through
the state-presheaf fibre at the ambient window, so the classical
semantics consumes the Brunetti--Fredenhagen--Verch state object
rather than a bare density matrix. The mass theorems are
\texttt{shadowLaw\_mass} and \texttt{rowsLaw\_mass\_of\_complete},
and the two-qubit instances tie to the fragment's landed shadows by
\texttt{rfl} (\texttt{bellQNet\_shadowPMF},
\texttt{mixedQNet\_shadowPMF}). Task (iii) is a first for the
catalogue: the target stratum has carried a mass-loss allowance from
birth, every stochastic entry's constraints force mass one (the
chain problem says so explicitly), and
\texttt{exists\_incomplete\_mass\_lt\_one} --- three of the four
projectors at the maximally mixed state --- is the first compiled
row set to exercise the allowance: an incomplete record family
denotes a genuine subprobability law. Task (iv) is
\texttt{exists\_nonPSD\_family\_negative\_shadow}, the
kernel-visible reason record positivity appears among the
hypotheses. The square also re-reads as a citizenship statement in
the semantic universe's fibre over a one-wire record interface, by
pushforward along the record embedding
(\texttt{su\_qnet\_adequacy}); the wiring pattern is the assignment
interface's, reused verbatim.

\paragraph{Adequacy (credal face).} Every other imprecise face in
this catalogue is chosen: a modeller declares a parameter box, or
quantifies over policies or schedulers, because caution recommends
it. Here the set-valuedness is a theorem of the frame's own
literature --- a functorial choice of single states is unrealizable
\cite{brunetti01,fewster11} --- so the state face of a quantum-net
presentation declares a set of states because the theory has nothing
narrower to declare. Stated self-contained, at exam scale:

\begin{problem}[the credal square]
\label{prob:mt-qnet-credal}
Data as in Problem~\ref{prob:mt-qnet}, except the state face: in
place of one $\rho$, a set $S$ of states is declared, and a
\emph{member witness} is a state of $S$ together with the
exact-rational row of its Born values. Write $D$ for the set of
member Born rows. On the compiled side the rows stay state-free: say
$\mu$ is \emph{admitted} when the compiled rows of \emph{some}
member witness are satisfied by $\mu$. Show: (i) $\mu$ is admitted
if and only if $\mu \in D$; (ii) every member's row has mass one, so
the downward-closed credal presentation generated by the members'
rays has $D$ as its radially maximal slice, and both verdict
polarities read as bounds uniform over every member; (iii) if $S$ is
closed under rational convex mixture then $D$ is closed under
rational mixtures of member rows, and the hypothesis is
necessary: over the two-element set consisting of the Bell state
and the maximally mixed state, the midpoint of the two rows prices
the double-zero record at $\tfrac38$, and no member does; (iv)
record the degeneracies --- with at least one member witness, the
whole-set reading of a lower-bound atom refuses every positive
threshold, and with none, both verdict kinds pass vacuously.
\end{problem}

\emph{Sketch.} (i) is the precise face's feeding lemma, run once per
member, under the union of singletons. (ii) The record family's
exhaustiveness prices every member at mass one, and the equal-mass
reconciliation is the automaton problem's
(Problem~\ref{prob:mt-ma-demonic}), reused. (iii) A quantum
mixture's Born row is the same rational mixture of the members'
rows, by linearity of the trace, so mixture-closure of $S$ pushes to
convexity of $D$; for the negative half, the member witnesses pin
the double-zero entry to $\tfrac12$ and $\tfrac14$, and $\tfrac38$
is neither. (iv) The zero model lies on every ray, and a universal
quantifier over an empty member index holds vacuously.

In the universe's terms: the admitted-by-some-member reading of the
state-free compiled rows and the union-of-singletons inclusion of
$D$ are isomorphic by an invertible globular cell, exactly as at the
decision and automaton formers --- but where those faces quantify
over an adversary the modeller invited, this one quantifies over the
physics. The statement and both directions are mechanized,
sorry-free, in the same module: the invertible cell is
\texttt{quantum\_credal\_adequacy}, with admission predicate
\texttt{qnetCredalC} and feeding lemma \texttt{qnetCredalC\_iff}
over member witnesses \texttt{RatShadowMember} (a member is carried
with its certificate, never bare); the slice reconciliation is
\texttt{quantum\_maxSlice}, packaged as a cell by
\texttt{quantum\_ray\_singleton\_reconciliation}; the verdict
transfers at both polarities are
\texttt{satCredal\_lower\_quantumCredal} and
\texttt{satCredal\_upper\_quantumCredal}, with law-set twins; the
convexity theorem is \texttt{quantumLaws\_mix}, consuming the
fragment's mixture seam (\texttt{recordShadowPMF\_mix}), with
negative witness \texttt{mixLaw\_not\_mem\_pair}; and the
degeneracies are
\texttt{not\_satCredal\_upper\_quantumCredal\_geC} and
\texttt{satCredal\_quantum\_of\_isEmpty}, recorded rather than
smoothed. One more degeneracy is quoted deliberately: at matrix
grade the folium is flat --- every linear functional that is unital
and positive on rank-one projectors is represented by a unique
density matrix
(\texttt{Qtsl.QuantumStates.existsUnique\_isState\_traceRepr}) ---
so this face's content lives in the set-valuedness and the credal
band, never in folium distinctions; the strict openness of the
Bell/mixed band is the fragment's
\texttt{exists\_shadow\_lowerE\_lt\_upperE}. Deliberately separate
strata, stated once: multi-window histories through the
class-operator machinery of \texttt{QtslAux.QuantumRecords}, the
named stretch tier; the quasilocal stratum, where folium
distinctions become real and the algebras stop being matrix
algebras; and dynamics beyond the kinematic fragment --- relative
Cauchy evolution and the stress-energy tensor \cite{fewster15} ---
outside this square by hypothesis rather than silently discharged.

\paragraph{Refinement and abstraction.} The two squares above relate
a System to its community semantics; this passage's statement-kind
relates a System to a System. The community doctrine being
transposed is effective field theory, and the grade comes first once
more: a reader who hears the phrase will expect loop integrals,
power counting, renormalization-group flow, and the continuum, and
none of those live at this stratum. What survives the transposition
--- finite site, matrix algebras, exact arithmetic, kinematics plus
one-step dynamics --- is the procedure's logic, and the logic turns
out to be the certifiable part. As its own standard references
present it \cite{burgess07,penco20}, effective field theory is a
two-ingredient move. Coarse-grain exactly: the Wilson action is what
remains after the high-energy modes are integrated out once and for
all --- a pushforward, with nothing approximate about the step
\cite{burgess07}. Then truncate with a declared tolerance: the
effective theory keeps finitely many operators, their coefficients
matched so that its predictions agree with the pushed-forward theory
up to the desired order \cite{penco20}. An effective theory is never
equal to the pushforward; it agrees with it inside a declared band,
and the band is part of the statement. The move's reach is why it
belongs in a catalogue of certifiable model types at all: it is
effective-theory reasoning that justifies neglecting quantum
corrections in macroscopic gravity while pricing them systematically
\cite{burgess07} --- ``known physics'' at any scale is an effective
theory with a band --- and a containment certificate is what that
kind of validity claim looks like as an artefact.

Both ingredients are already on this entry's shelf, the first as the
record-level coarse-graining calculus --- record atoms merge along
any map of indices, and pushing the law is measuring the merged
observable family --- and the second as a containment. Write $d$ for
the total-variation distance on outcome rows, half the $\ell^1$
distance, exact rational at rational rows; inflate a law set by a
tolerance $\varepsilon \in \mathbb{Q}_{\geq 0}$, admitting every row
within distance $\varepsilon$ of a member. The
\emph{certified-abstraction} statement between a fine System and a
coarse one, along coarse-graining data $f$, is then a containment:
every law the fine theory permits, pushed forward along $f$, lies in
the coarse theory's $\varepsilon$-inflated law set. Error bars
become imprecision, and ``valid to accuracy $\varepsilon$'' becomes
decidable arithmetic at finite index. Three choices inside that
sentence carry the doctrine. The quantifier direction is soundness:
the coarse law set may be strictly fatter than the pushforward ---
that is what ``effective'' means --- but no behaviour the fine theory
permits is ruled out, so the abstraction is epistemically honest:
the coarse description's uncertainty band covers the fine truth.
The tolerance lives on the law set rather than the state set,
because matching in the literature is a statement about observables,
and the checker's arithmetic lives on rows. And the tolerance is
declared, never derived: power counting --- the literature's
dimensional-analysis estimate of what a given truncation achieves
\cite{burgess07} --- proposes $\varepsilon$, and the containment
certificate checks it. Burgess's ``Effective-Lagrangian Logic''
\cite{burgess07} is a four-step procedure whose middle steps rest on
power-counting estimates; here the estimates sit exactly where an
untrusted heuristic belongs, in producer-side search, and
trust-the-estimate is replaced by check-the-containment.

The worked exhibit is a block spin, and one design point governs its
reading: the coarse theory is declared independently of the fine
one. Had its presentation been defined as the coarsened fine family,
coarse-graining would commute with compilation on the nose --- the
strict square is a landed theorem
(\texttt{pushforward\_ofPresentation}, module
\texttt{QtslAux.QuantumCoarseGraining}) --- and containment would
hold at tolerance zero, exhibiting exact lumping rather than
truncation; the effective-theory content is precisely that the
declared coarse theory is simpler than the pushforward and provably
within the band. Stated self-contained, at exam scale:

\begin{problem}[the block-spin containment]
\label{prob:mt-qnet-refine}
Data on the section's running two qubits, records the four joint
computational-basis projectors, indexed by pairs of bits. The
\emph{fine} state face is the rational segment between the basis
state at $(0,1)$ and the two-branch state with amplitudes
$\tfrac{24}{25}$ and $-\tfrac{7}{25}$ on $(0,0)$ and $(1,1)$ ---
Born rows $(0, 1, 0, 0)$ and $(\tfrac{576}{625}, 0, 0,
\tfrac{49}{625})$ respectively --- a mixture-closed set, as the
credal square requires; the fine mechanism is the reflection
exchanging those two branch vectors, unitary, involutive, and
rational-entried because $7$--$24$--$25$ is a Pythagorean triple.
The coarse-graining $f$ keeps the first bit. The \emph{coarse}
theory, declared independently: records the two block-spin
projectors (the fiber sums of the fine family), state face a single
product state with row $(1, 0)$, mechanism the identity kernel ---
the block spin is $0$ and nothing evolves. Show: (i) the segment
member at parameter $t$ pushes forward to the row $\big(t +
(1-t)\tfrac{576}{625},\ (1-t)\tfrac{49}{625}\big)$, and containment
holds at tolerance $\tfrac{49}{625}$; (ii) the band is attained ---
any certifying tolerance is at least $\tfrac{49}{625}$ --- so
containment at $\varepsilon = 0$ is refuted: the truncation is
genuine; (iii) on the mechanism's two-point orbit, the one-step
defect --- the distance between coarse-graining after one fine step
and one coarse step after coarse-graining --- equals
$\tfrac{49}{625}$ at both points, and every horizon $k$ is then
valid within $k \cdot \tfrac{49}{625}$; (iv) the further merge to a
single total record is contained at tolerance zero, and the
composite fine-to-block-to-total containment holds at
$\tfrac{49}{625} + 0$: budgets add along a tower.
\end{problem}

\emph{Sketch.} (i) The pushforward sums fibers and is affine in $t$:
the first extreme pushes to $(1, 0)$, at distance zero, the second
to $(\tfrac{576}{625}, \tfrac{49}{625})$, at distance
$\tfrac{49}{625}$, and the member at $t$ to distance
$(1-t)\tfrac{49}{625}$. (ii) The coarse theory offers exactly one
row, so the second extreme forces $\varepsilon \geq \tfrac{49}{625}$
on any certificate, and a positive rational is not zero. (iii) The
reflection swaps the two extreme states, so the pushed rows exchange
$(1, 0)$ and $(\tfrac{576}{625}, \tfrac{49}{625})$ while the
identity kernel holds each fixed; both defects are the distance
already computed. For the horizon, the coarse step is a
row-stochastic kernel, and such kernels contract total variation, so
the $k$-step distance regroups through the triangle inequality into
$k$ one-step defects. (iv) Every state prices the single total
record at trace one, so fine and coarse push to the same point row;
merging twice is one merge along the composite map, and the triangle
inequality adds the two budgets.

In the universe's terms: a refinement relation between two
presentations, whose content at the credal stack is inclusion of law
sets after inflation --- so verdicts transfer across it with an
$\varepsilon$ margin, in the conservative direction. The vocabulary
and the exhibit are mechanized, sorry-free, in the same module
family as the squares (module
\texttt{SuFunctor.Adequacy.QuantumContainment}): the distance is
\texttt{tvDistQ}, with the metric lemmas and
\texttt{tvDistQ\_eq\_zero\_iff}; the inflation is
\texttt{credalInflate}, monotone in the tolerance and in the set and
commuting with the union-of-singletons form the credal face
compiles; the statement is \texttt{Containment}, with monotonicity
(\texttt{Containment.mono}) and additive vertical composition
(\texttt{Containment.comp}); the row pushforward is the fiber sum
(\texttt{SubPMF.map\_wt\_apply}), contracts the distance
(\texttt{tvDistQ\_map\_le}), and measures the merged family
(\texttt{recordShadow\_coarsenFamily}). The exhibit's theorems are
\texttt{uv\_ir\_containment}, \texttt{uv\_ir\_containment\_lower},
and \texttt{not\_uv\_ir\_containment\_zero} for tasks (i) and (ii);
\texttt{onestep\_defect} and \texttt{horizon\_bound}, instantiating
the generic \texttt{lax\_horizon} through the contraction
\texttt{tvDistQ\_stochAct\_le}, for task (iii); and
\texttt{ir\_tot\_containment} composing to
\texttt{uv\_tot\_containment} for task (iv), whose zero leg consumes
nothing but conservation of mass. The state face exercises the
credal square's own convexity hypothesis (\texttt{uvSeg\_mixClosed},
consumed by \texttt{uvLaws\_mix\_exercised}), and the two
pushforward notions in the tree --- Gram submatrix at measure level,
fiber sum at row level --- agree on the exhibit
(\texttt{pushforward\_probRow\_bridge}). The qualitative half of the
tower rides along: coarser is more classical, in the monotone sense
that off-diagonal mass is antitone along tower legs
(\texttt{eftTower\_offDiagMass\_antitone}) --- though on this
exhibit the dividend is degenerate, since projective records
decohere every branch at the fine stage already
(\texttt{eftStage\_classical}); the strict decoherence drop lives on
the library's three-stage diamond tower, cited rather than
re-proven.

The fences, in the section's own manner. State-space compression,
once the first fence here, is discharged at the same finite grade:
the block spin whose state space genuinely shrinks lands as an
isometry stratum and a trace-preserving Kraus channel with the
compression/record-shadow interchange proved channel-wide
(\texttt{recordShadow\_kchanApply}); the partial trace acquires an
explicit Kraus certificate (\texttt{kchanApply\_ptrKraus\_eq\_ptR}),
and the containment certificate transfers verbatim across the
compression --- the same declared band, still attained, still
refuted at zero tolerance (\texttt{uv\_smallIR\_containment}). What
stays fenced there is any compression-ratio or optimality claim:
which compressions exist or are efficient is producer-side search,
never certificate content. No power-counting claim: it enters
as producer guidance only, and the checker never consumes an
estimate as if it were a theorem. No continuum word: genuine
Wilsonian renormalization --- loop integrals, continuum limits,
scheme independence --- and perturbative error control made rigorous
are research mathematics on the far side of this fence, and the
literature's own error control is perturbative through and through,
which is why the mechanized entry point is the finite stratum, where
the error is one exact rational and there is no regularization
scheme to hide in. What the fences protect is a modest, checkable
form of the doctrine the literature itself teaches --- every theory
in this catalogue is somebody's effective theory \cite{penco20}, and
certified budgets add along the descent.

\paragraph{Status.} Semantics is the strong leg, and unusually
strong for an entry this young: the kinematic mathematics
(section~\ref{sec:su-site}) and both adequacy faces are mechanized
end to end, with the negative witnesses --- dropped positivity,
dropped completeness, dropped mixture-closure --- carried as
theorems beside the positive ones. The refinement statement ships at
the same grade --- band proved, band attained, zero tolerance
refuted --- so the mechanized column runs both statement-kinds,
System-to-semantics and System-to-System. Representation is design
work:
no producer surface exists at this preview, and the normative row
grain is the adequacy module's own rows structure
(\texttt{QNetRows}: sparse outcome rows over the record index;
well-formedness as nonnegativity plus mass at most one; completeness
deliberately separate), which a future declaration vocabulary and
producer surface will mirror rather than reinvent. Certification:
nothing shipped; the one genuinely open producer-side design point
is exact positive-semidefiniteness certification over
$\mathbb{Q}[i]$, where an exact $LDL^{*}$ witness settles the
generic case and the boundary-rank convention needs fixing before
the first certificate does. Compositional closure: this version is
a closed former, with no wiring interface, and the omission is
argued from the frame rather than bare-declared --- the physically
native composition operations are spacelike tensor composition of
nets, the named successor, a Composition Patterns sibling whose side
condition is Einstein causality; and subregion embedding, which the
site already carries internally. The deterministic-PDE entry set the
precedent of landing former-level ahead of its composition story,
and this entry follows it.

\syscatmetalean{\begin{mathpar}
\inferrule*[Right=\rulename{S-QNET}]{
  K : \mathsf{SiteKit} \\
  N : \overrightarrow{\langle \mathit{window},\, \mathit{gens}
  \rangle} \text{ algebra rows} \\
  M : \mathsf{MatterData} \\
  R : \overrightarrow{\langle \mathit{index},\, \mathit{matrix}
  \rangle} \text{ record rows} \\
  \sigma : \mathsf{StateFace}
}{
  \mathrm{QNet}(K;\, N;\, M;\, R;\, \sigma) : \mathsf{System}
}
\end{mathpar}
The site kit lists finitely many points with their chronology;
diamonds, covers, and the ambient window are derived. Algebra rows
present one generating family per window, entries in
$\mathbb{Q}[i]$; matter data comprises worldlines, jointly
exhaustive factors, and a reversible mechanism; record rows live at
the ambient window. The state face declares one state with its
exact-rational Born row (the precise face) or a set of states, each
member carried with its row witness (the credal face). The
well-formedness obligations are the community's own axioms ---
isotony, Einstein causality, the time-slice property, record
positivity and exhaustiveness, state positivity and normalization
--- consumed by the mechanized textbook corner as the fragment's
fields.}

\section{Future Work}
\label{sec:mt-frontier}

The frames collected here are the catalogue's frontier: no producer
surface at this preview --- no rows, no API, no certificate.
Encodings for some of them are in progress, and will graduate their
entries to full catalogue sections as producer surfaces land. They
stay in the catalogue as prose entries precisely so the chapter
cannot quietly shrink its scope, and each keeps its full four-leg
grading; what they do not get is an invented formal block, because a
metanotation statement with no mechanized ground truth underneath
would be exactly the coverage-matrix-of-green-ticks failure this
chapter refuses. The quantum nets of section~\ref{sec:mt-qnet} are
deliberately not collected here: that entry's semantics and adequacy
squares are mechanized end to end, and only its producer surface
remains future work.

\subsection{Open Games}
\label{sec:mt-games}

\paragraph{The frame.} Compositional game theory: players, strategy
sets, and utilities in the lens-shaped composition structure whose
entire point is that games compose.

\paragraph{Profile.} Intended: the composition structure preserved
as the representation's grain --- a rendering that flattens an open
game to one normal-form matrix would fail the representation leg by
construction. Not yet designed in rows.

\paragraph{Semantics.} Equilibrium semantics under composition; the
research content of this row lives here.

\paragraph{Status.} Nothing shipped; the thesis's own bimatrix
analysis is a \emph{use} of game theory, not a subsumption witness,
and we decline to count it. The concrete entry road is the
certification leg: for polynomial games,
$\varepsilon$-best-response inequalities are polynomial inequalities
over strategy boxes --- exactly the shape the kernel decides --- so a
certified $\varepsilon$-Nash artefact is buildable ahead of the
representation and semantics legs that carry the real research
weight. Graded absent until any leg ships.

\subsection{Mean-Field Games}
\label{sec:mt-mfg}

\paragraph{The frame.} A representative agent coupled to the
distribution of a population of itself; equilibria are
forward--backward systems (Hamilton--Jacobi--Bellman backward,
Fokker--Planck forward).

\paragraph{Profile.} Intended: a representative-agent mechanism
coupled to a population-distribution state. Not designed in rows.

\paragraph{Semantics.} The equilibrium system itself; in the
master-equation formulation, a PDE system --- so the road runs
through section~\ref{sec:mt-pde}, whose beachhead changes this row's
feasibility without yet shortening it much.

\paragraph{Status.} Nothing shipped; this cell was not-covered in the
2023 SDCPN grading too, and it remains the list's structurally
hardest entry. The first tractable question is narrower than the
frame:
certified $\varepsilon$-consistency of a \emph{given} candidate
equilibrium pair, both components polynomial, may be reachable well
before the frame itself. Graded absent.

\subsection{Score-Based Generative Models}
\label{sec:mt-sbgm}

\paragraph{The frame.} Diffusion models read as dynamics: a learned
score field is the drift of a reverse-time stochastic differential
equation, and the generative process is that equation run from
noise.

\paragraph{Profile.} Intended: an SDE mechanism
(section~\ref{sec:mt-sde}'s profile) whose drift is a neural
network, carried by content-addressed reference with piecewise-linear
structure exposed to the kernel the same way the discrete campaigns
already exercise.

\paragraph{Semantics.} The generative path measure, through the
same staged generator bridges as section~\ref{sec:mt-sde}.

\paragraph{Status.} Nothing shipped for the frame as such --- but
this row is a \emph{composite} gap, not a novel one: the
multidimensional diffusion machinery is in place, the network-drift
representation is in place elsewhere in the stack, and what remains is
their composition. Graded absent, with its cost inherited almost
entirely from parts that already exist.

 \chapter{Systems: Composition Patterns}
\label{ch:qtsl-composition-patterns}

\noindent\draftmark\medskip

This chapter is the second wing of the System Catalogue --- the first,
chapter~\ref{ch:qtsl-mathematical-models}, inventories what a single mechanism
can be; this one inventories how Systems combine. A specification is
only as compositional as the model underneath it, so the System layer
ships composition as a small, closed inventory of \emph{patterns}:
each is a syntactic row family with a typing rule, a semantics, and
an explicit account of which certificates transport across the
composite boundary and which are
\emph{refused}. The refusals carry real content: per-component stability
certificates must not compose to composite stability, because that
inference is false (re-entrant lines destabilize below capacity under
plausible dispatch policies --- Kumar--Seidman and
kin~\cite{kumar1990dynamic}), and the
pattern inventory encodes the refusal structurally rather than leaving
it to the modeller's vigilance.

A pattern, throughout, is a \emph{usage of the wiring stratum plus its
transport rules} --- the row-level counterpart of composition by
wiring diagram in operadic treatments of dynamical
systems~\cite{libkind21}. The composition data are ordinary rows --- the
manifest's own composite-stratum descriptors, or one of four dedicated
pattern row families --- and the transport rules that move
certificates across a composite boundary are rules of the Rule
Catalogue, not machinery of this chapter. Each entry below states the
pattern in prose (composition data, transport, witness status), then
closes with the same two blocks as every System Catalogue entry. The
\emph{Computer Science Metanotation} block gives the entry's \emph{term
former}: an inference rule whose premises are the formation data the
pattern builders consume and refuse on, with their sorts, and whose
conclusion is a primitive judgment at the sort $\mathsf{System}$ ---
or at the pattern's own sort, where the mechanized constructor
forms a pattern row citing its subjects rather than a System; the
four dedicated row families do exactly that, and their rules say so.
The \emph{API} block carries producer surfaces extracted from the
toolkit sources at every regeneration. Two reading conventions,
stated once. First, the wiring-stratum rules display a \emph{single}
pattern row in the composite's row slot: a composite carries any
number of rows across the families, with each row's premises checked
independently, and the head form elides the ambient rows. Second, as
in the first wing, optional formation data (marked $?$) are consumed
without appearing in the head form. The first four entries --- the
three wiring modes and nesting --- are the composition machinery that
the model-type chapter's fourth grading leg, compositional closure,
checks frames against; the remaining formal entries are instruments
with landed witnesses. Two entries near the end withhold their formal
blocks, by the same refusal that governs the first wing's frontier
entries: event ports, whose primitive's producer surface has already
landed, and state sharing, reserved outright. Between them sits
supervision, whose semantic tier is mechanized ahead of its producer
surface --- its gating uses an arc kind the inventory already
carries. The chapter closes with an entry that records a closure
property rather than a pattern.

The composition tier underneath this inventory is mechanized,
sorry-free, in the toolkit's Lean tree (\texttt{lean/su-functor/}).
The semantic universe's decoration functor carries its compositor with
every coherence obligation proved --- transport functoriality,
compositor naturality, associativity, and both unit laws, on the nose
modulo the base's definitional transports (module
\texttt{SuFunctor.SingleSiteSemanticUniverse}) --- and the possibilistic
inclusion of hypernets composes through it exactly, with no determinism
hypothesis (\texttt{SuFunctor.HypernetInclusion}). On the
content-addressed wire stratum the calculus is strict, and failure is
part of the account (module \texttt{SuFunctor.HypernetComposition}): composition of
aligned interfaces is associative and unital on the nose
(\texttt{AlignedLoose.seq\_assoc} --- hierarchy pays no coherence tax,
which is the semantic face of the nesting entry's freeness);
components with equal semantic interfaces are interchangeable in any
composite (\texttt{austereDnf\_subst} --- the congruence property every
entry below presupposes when it speaks of transport across a
boundary); an infeasible component makes its composite infeasible
(\texttt{admRegion\_append\_empty\_left}, iterating along the strict
associativity above), while the converse fails by a mechanized witness
(\texttt{share\_contradiction}); and renaming a composite is the
composite of the renamings (\texttt{theta\_rename\_seq}), so a
certificate transported across a renamed boundary composes to the
transported composite. What follows is therefore inventory over a
proved foundation, and each entry's transport story stands on that
tier rather than re-arguing it.

\section{Parameter Setting}
\label{sec:cp-set}

One system's output fixes another's parameter: an input wire bound to
a parameter appearing in the consumer's field polynomials. This is the
``set'' mode of the wiring stratum --- directed by construction, so
the row demands its \textsf{from} endpoint resolve to a declared
output and its \textsf{to} endpoint to a declared input wherever the
component's interface is declared. Transport is the assume--guarantee
square~\cite{capucci26}: the wiring is certified for spec-to-spec entailment with
slack, the leaves are ordinary sequents, and the slack arithmetic is
checker-recomputed. \textbf{Witnessed} --- this is the mode behind the
inventory's one cross-frame composition witness (the ODE plant wired
to a discrete controller of
section~\ref{sec:mt-ode}, every mutant of the composite refusing),
and behind the inventory's deliberate asymmetry: one mode witnessed,
the rest below graded only as far as their evidence reaches.

\syscatmetapat{\begin{mathpar}
\inferrule*[Right=\rulename{S-SET}]{
  S_1, \ldots, S_k : \mathsf{System} \\
  \mathsf{from},\, \mathsf{to} : \mathsf{PortRef} \text{ into the
  roster},\ \mathsf{from} \neq \mathsf{to} \\
  \mathrm{dir}(\mathsf{from}) = \mathsf{output} \\
  \mathrm{dir}(\mathsf{to}) = \mathsf{input}
}{
  \mathsf{Composite}\big(S_1, \ldots, S_k;\,
  \langle \mathsf{set},\, \mathsf{from},\, \mathsf{to} \rangle\big) :
  \mathsf{System}
}
\end{mathpar}
$\mathsf{PortRef} = \langle \mathsf{component},\, \mathsf{wire}
\rangle$ addresses the roster by index. The direction premises
resolve against declared \emph{mechanism} interfaces; an endpoint
whose component declares none (a composite; a mechanism without an
interface binding) passes to the checker's replay.}
\ruleblock{API}
\sourcenote{Producer surface of the \texttt{mtk} toolkit for this composition pattern: the TypeScript arrangement (npm package \texttt{@maitria/mtk}) and the Rust catalogue crate (\texttt{crates/mtk-catalogue}) --- one Rust implementation underneath, reached from TypeScript in-process through a native binding, so the two surfaces cannot diverge. Signatures and constructions are extracted from the toolkit sources at every regeneration of this catalogue; each construction is quoted verbatim from the toolkit's conformance corpus, code whose emitted certificate bytes are verified byte-identical against the committed golden vectors (helper values named \texttt{fx\_*} or \texttt{f.*} are that corpus's shared fixtures, defined beside the quoted lines). This API constructs and emits; it never checks --- the checker re-derives every claim independently.}
\apilang{TypeScript}{---}
The pattern-level TypeScript binding is not yet arranged: the manifest-level wiring descriptors (\texttt{WiringRow}, \texttt{SwitchEdge}, \texttt{ScheduleAssumption}) already ride the TS \texttt{system} namespace (\texttt{system.composite}), while the pattern builders and the four pattern row families are Rust-only at this preview.
\apilang{Rust}{mtk\_catalogue::compose::CompositeBuilder::set}
\apiprov{signature --- \texttt{crates/mtk-catalogue/src/compose/builder.rs}}
\begin{rustsrc}
pub fn set(mut self, from: PortRef, to: PortRef) -> Self

pub fn build(self) -> Result<SystemManifest, Refusal>
\end{rustsrc}
\apiprov{construction --- \texttt{crates/mtk-catalogue/tests/compose\_common/mod.rs}}
\begin{rustsrc}
CompositeBuilder::new()
    .component(plant())
    .component(controller())
    .set(port(1, 1), port(0, 0))
    .build()
    .expect("set endpoints resolve against the declared interfaces")
    .named("set-loop")
\end{rustsrc}

\section{Variable Sharing}
\label{sec:cp-share}

Two systems read identified wires: the ``share'' mode identifies its
two endpoints as one variable, with direction consumed by the
causality check rather than declared in the row (stochastic wires
remain directed-only at this preview). Transport: the same
assume--guarantee square; the admissible-tier anchor is the small-gain
arithmetic witness --- the certified composite in which cross-supplies
cancel identically. \textbf{Status}: machinery staged at the rule
level; no heterogeneous share witness has landed yet. The mode's
\emph{semantics}, by contrast, is already theorem-grade: admissibility
of an identified composite is the conjunction of the parts'
admissibilities, mechanized with no determinism hypothesis at the
catalogue stratum (\texttt{glue\_admissible\_iff}) and the
semantic-universe stratum (\texttt{admRegion\_append},
\texttt{theta\_comp}); and identification is exactly where feasibility
fails to compose --- the minimal mechanized witness is a pair of
individually feasible constant sources demanding different values of
one shared variable (\texttt{share\_feasible},
\texttt{share\_contradiction}, module
\texttt{SuFunctor.HypernetComposition}), while infeasibility, once
present, propagates outward through further wiring
(\texttt{admRegion\_append\_empty\_left}).

\syscatmetapat{\begin{mathpar}
\inferrule*[Right=\rulename{S-SHARE}]{
  S_1, \ldots, S_k : \mathsf{System} \\
  \mathsf{from},\, \mathsf{to} : \mathsf{PortRef} \text{ into the
  roster, declared where an interface exists} \\
  \mathsf{from} \neq \mathsf{to}
}{
  \mathsf{Composite}\big(S_1, \ldots, S_k;\,
  \langle \mathsf{share},\, \mathsf{from},\, \mathsf{to} \rangle\big) :
  \mathsf{System}
}
\end{mathpar}
The row identifies rather than assigns, which is why, unlike
\rulename{S-SET}, no direction premise appears: either side may be an
input or an output, and the two wires denote one variable of the
composite.}
\ruleblock{API}
\sourcenote{Producer surface of the \texttt{mtk} toolkit for this composition pattern: the TypeScript arrangement (npm package \texttt{@maitria/mtk}) and the Rust catalogue crate (\texttt{crates/mtk-catalogue}) --- one Rust implementation underneath, reached from TypeScript in-process through a native binding, so the two surfaces cannot diverge. Signatures and constructions are extracted from the toolkit sources at every regeneration of this catalogue; each construction is quoted verbatim from the toolkit's conformance corpus, code whose emitted certificate bytes are verified byte-identical against the committed golden vectors (helper values named \texttt{fx\_*} or \texttt{f.*} are that corpus's shared fixtures, defined beside the quoted lines). This API constructs and emits; it never checks --- the checker re-derives every claim independently.}
\apilang{TypeScript}{---}
The pattern-level TypeScript binding is not yet arranged: the manifest-level wiring descriptors (\texttt{WiringRow}, \texttt{SwitchEdge}, \texttt{ScheduleAssumption}) already ride the TS \texttt{system} namespace (\texttt{system.composite}), while the pattern builders and the four pattern row families are Rust-only at this preview.
\apilang{Rust}{mtk\_catalogue::compose::CompositeBuilder::share}
\apiprov{signature --- \texttt{crates/mtk-catalogue/src/compose/builder.rs}}
\begin{rustsrc}
pub fn share(mut self, from: PortRef, to: PortRef) -> Self
\end{rustsrc}
\apiprov{construction --- \texttt{crates/mtk-catalogue/tests/compose\_common/mod.rs}}
\begin{rustsrc}
CompositeBuilder::new()
    .component(plant())
    .component(controller())
    .share(port(0, 1), port(1, 0))
    .build()
    .expect("share endpoints resolve against the declared interfaces")
    .named("shared-bus")
\end{rustsrc}

\section{Mode Switching}
\label{sec:cp-switch}

A supervisory edge hands the active role between components. Switch
edges live at the composite stratum in three species: guarded (guard
content at the composite level), rate (stochastic hand-off at a
state-dependent intensity --- the species that makes switched
diffusions composable), and forced (domain-exit at the composite
stratum). Transport: the switch-composition rule, staged at the rule
level with per-profile soundness lemmas in the inventory.
\textbf{Status}: rule staged, witnesses pending; the static frames
(the graphical-model rows of chapter~\ref{ch:qtsl-mathematical-models}) carry
\emph{inapplicability} arguments instead --- a frame with no
internal time axis has nothing for a supervisory switch to sequence
--- and that escape is earned per-frame by an argument from the
frame's own semantics (section~\ref{sec:mt-pgm} states it), never by
a bare declaration.

\syscatmetapat{\begin{mathpar}
\inferrule*[Right=\rulename{S-SWITCH}]{
  S_1, \ldots, S_k : \mathsf{System} \\
  i, j < k,\ i \neq j \\
  \sigma : \mathsf{species}
}{
  \mathsf{Composite}\big(S_1, \ldots, S_k;\,
  \langle i \to j,\, \sigma \rangle\big) : \mathsf{System}
}
\end{mathpar}
\[
\sigma \;::=\;
\mathsf{guarded}\big(\overrightarrow{\mathsf{constraint}}\big)
\;\mid\; \mathsf{rate}(\mathsf{poly})
\;\mid\; \mathsf{forced}
\]
$i \to j$ names roster indices --- the component handing off and the
component receiving --- and the $i \neq j$ premise is the no-self-switch
refusal. The rate polynomial is exact-rational, like every field
polynomial in the catalogue.}
\ruleblock{API}
\sourcenote{Producer surface of the \texttt{mtk} toolkit for this composition pattern: the TypeScript arrangement (npm package \texttt{@maitria/mtk}) and the Rust catalogue crate (\texttt{crates/mtk-catalogue}) --- one Rust implementation underneath, reached from TypeScript in-process through a native binding, so the two surfaces cannot diverge. Signatures and constructions are extracted from the toolkit sources at every regeneration of this catalogue; each construction is quoted verbatim from the toolkit's conformance corpus, code whose emitted certificate bytes are verified byte-identical against the committed golden vectors (helper values named \texttt{fx\_*} or \texttt{f.*} are that corpus's shared fixtures, defined beside the quoted lines). This API constructs and emits; it never checks --- the checker re-derives every claim independently.}
\apilang{TypeScript}{---}
The pattern-level TypeScript binding is not yet arranged: the manifest-level wiring descriptors (\texttt{WiringRow}, \texttt{SwitchEdge}, \texttt{ScheduleAssumption}) already ride the TS \texttt{system} namespace (\texttt{system.composite}), while the pattern builders and the four pattern row families are Rust-only at this preview.
\apilang{Rust}{mtk\_catalogue::compose::CompositeBuilder::switch\_guarded}
\apiprov{signature --- \texttt{crates/mtk-catalogue/src/compose/builder.rs}}
\begin{rustsrc}
pub fn switch_guarded(
    mut self,
    from_component: u32,
    to_component: u32,
    guard: Vec<DomainConstraint>,
) -> Self

pub fn switch_rate(mut self, from_component: u32, to_component: u32, rate: Poly) -> Self

pub fn switch_forced(mut self, from_component: u32, to_component: u32) -> Self
\end{rustsrc}
\apiprov{construction --- \texttt{crates/mtk-catalogue/tests/compose\_common/mod.rs}}
\begin{rustsrc}
let too_hot = DomainConstraint {
    poly: mono(1, 0, 1, r(1, 1)),
    op: "ge".to_owned(),
    rhs: r(100, 1),
};
CompositeBuilder::new()
    .component(thermal("heater", r(3, 1)))
    .component(thermal("cooler", r(-1, 10)))
    .switch_guarded(0, 1, vec![too_hot])
    .switch_rate(1, 0, constant(1, r(1, 20)))
    .switch_forced(0, 1)
    .build()
    .expect("switch edges cite roster components")
    .named("supervisor")
\end{rustsrc}

\section{Nesting}
\label{sec:cp-nest}

Hierarchical composition is free in the theory: a component of a
composite is any System --- a mechanism or a built composite --- so
composition terms nest by construction, and no dedicated descriptor
exists to get out of sync with the term structure. The flatten rule
guarantees the discipline's point: composites add structure and
scalability, never expressiveness. Wire-level containment resolves
against declared \emph{mechanism} interfaces, and a composite
component's wires pass to the checker's replay (composites declare no
export map at this preview). \textbf{Status}: structural (the
construction cannot fail to nest), but unwitnessed as a certified
heterogeneous artefact; hierarchical control is the named first
candidate.

\syscatmetapat{No former of its own --- deliberately. Nesting is the recursion
already present in the core grammar: \rulename{Composite}'s premises
(section~\ref{sec:mt-core}) quantify over $\mathsf{System}$, which a
built composite inhabits exactly as a mechanism does, so
$\mathsf{Composite}(\ldots, \mathsf{Composite}(\ldots), \ldots)$ is
well-formed by the rules already stated --- cited here, not restated.
A dedicated nesting rule would exist only to drift from the term
grammar it duplicates.}
\ruleblock{API}
\sourcenote{Producer surface of the \texttt{mtk} toolkit for this composition pattern: the TypeScript arrangement (npm package \texttt{@maitria/mtk}) and the Rust catalogue crate (\texttt{crates/mtk-catalogue}) --- one Rust implementation underneath, reached from TypeScript in-process through a native binding, so the two surfaces cannot diverge. Signatures and constructions are extracted from the toolkit sources at every regeneration of this catalogue; each construction is quoted verbatim from the toolkit's conformance corpus, code whose emitted certificate bytes are verified byte-identical against the committed golden vectors (helper values named \texttt{fx\_*} or \texttt{f.*} are that corpus's shared fixtures, defined beside the quoted lines). This API constructs and emits; it never checks --- the checker re-derives every claim independently.}
\apilang{TypeScript}{---}
The pattern-level TypeScript binding is not yet arranged: the manifest-level wiring descriptors (\texttt{WiringRow}, \texttt{SwitchEdge}, \texttt{ScheduleAssumption}) already ride the TS \texttt{system} namespace (\texttt{system.composite}), while the pattern builders and the four pattern row families are Rust-only at this preview.
\apilang{Rust}{mtk\_catalogue::compose::CompositeBuilder::component}
\apiprov{signature --- \texttt{crates/mtk-catalogue/src/compose/builder.rs}}
\begin{rustsrc}
pub fn component(mut self, sys: SystemManifest) -> Self
\end{rustsrc}
\apiprov{construction --- \texttt{crates/mtk-catalogue/tests/compose\_common/mod.rs}}
\begin{rustsrc}
CompositeBuilder::new()
    .component(set_composite())
    .component(thermal("watchdog", r(-1, 100)))
    .switch_forced(1, 0)
    .build()
    .expect("nesting is structural")
    .named("hierarchical")
\end{rustsrc}
\apiprov{Nesting is structural: \texttt{component} accepts any manifest --- a mechanism or a built composite --- so composition terms nest by construction; no dedicated descriptor exists.}

\section{The Hybrid-Time Pattern}
\label{sec:cp-hybtime}

The pattern that discharges multi-clock composition in practice: put
the discrete component's tick on a timer mechanism, prove the
component's rows schedule-generically --- per-tick, in its own counter
--- and carry the schedule class (period bounds, jitter bracket) in
the \emph{assumption slot} of the composition square. One value
function, one cover, both time semantics --- with the price of jitter
certified, not assumed. \textbf{Witnessed}: the lander free-flight
pair of section~\ref{sec:mt-tick} (the same funnel under exact tick
and jitter bracket from one table and one cover, jitter costing a
certified factor $2.66$ in ticks-to-funnel, with two exact infeasibility
theorems marking where full dispersion is impossible), checked on an
independent lane; the pattern text is the staging-independent recipe.
The square behind the recipe is machine-checked, sorry-free, in the
toolkit's Lean tree (module \texttt{SuFunctor.Adequacy.HybridTime}):
compiling the rows onto the coarse clock, versus reading the per-tick
semantics back through the declared clock morphism, the walled
composite sits on the weak side in general
(\texttt{hybridTime\_square\_subset}); the two readings agree whenever
the transport relation is a function graph --- the pattern's own tick
seam is one (\texttt{tickSeam\_hybridTime\_square\_eq}) --- and at a
jittered-coarse seam the inclusion is strict, which is why the
assumption slot, and nothing else, carries the price of jitter.

\syscatmetapat{\begin{mathpar}
\inferrule*[Right=\rulename{S-HYBTIME}]{
  S_1, \ldots, S_k : \mathsf{System} \\
  i < k \\
  \tau : \mathsf{TickDecl}
}{
  \mathsf{Composite}\big(S_1, \ldots, S_k;\,
  \langle i,\, \tau \rangle\big) : \mathsf{System}
}
\end{mathpar}
At most one assumption per component --- a duplicate is refused with
both declarations serialized. $\mathsf{TickDecl}$ is
section~\ref{sec:mt-tick}'s grammar, here in the assumption slot of
the composition square rather than on the mechanism.}
\ruleblock{API}
\sourcenote{Producer surface of the \texttt{mtk} toolkit for this composition pattern: the TypeScript arrangement (npm package \texttt{@maitria/mtk}) and the Rust catalogue crate (\texttt{crates/mtk-catalogue}) --- one Rust implementation underneath, reached from TypeScript in-process through a native binding, so the two surfaces cannot diverge. Signatures and constructions are extracted from the toolkit sources at every regeneration of this catalogue; each construction is quoted verbatim from the toolkit's conformance corpus, code whose emitted certificate bytes are verified byte-identical against the committed golden vectors (helper values named \texttt{fx\_*} or \texttt{f.*} are that corpus's shared fixtures, defined beside the quoted lines). This API constructs and emits; it never checks --- the checker re-derives every claim independently.}
\apilang{TypeScript}{---}
The pattern-level TypeScript binding is not yet arranged: the manifest-level wiring descriptors (\texttt{WiringRow}, \texttt{SwitchEdge}, \texttt{ScheduleAssumption}) already ride the TS \texttt{system} namespace (\texttt{system.composite}), while the pattern builders and the four pattern row families are Rust-only at this preview.
\apilang{Rust}{mtk\_catalogue::compose::CompositeBuilder::schedule}
\apiprov{signature --- \texttt{crates/mtk-catalogue/src/compose/builder.rs}}
\begin{rustsrc}
pub fn schedule(mut self, component: u32, tick: TickDecl) -> Self
\end{rustsrc}
\apiprov{construction --- \texttt{crates/mtk-catalogue/tests/compose\_common/mod.rs}}
\begin{rustsrc}
CompositeBuilder::new()
    .component(plant())
    .component(controller())
    .share(port(0, 1), port(1, 0))
    .set(port(1, 1), port(0, 0))
    .schedule(
        1,
        TickDecl::Jitter {
            lo: r(9, 100),
            hi: r(11, 100),
        },
    )
    .build()
    .expect("schedule assumption cites a roster component")
    .named("multi-clock")
\end{rustsrc}

\section{The Funnel Pattern}
\label{sec:cp-funnel}

Transit plus absorbing tail: an outer transit certificate (reach the
gate within a budget, worst-case) composed with an inner invariance
certificate (the gate region is forward-invariant and its holds are
certified), citing the inner region as the outer target through the
tower discipline. This is the certified shape of approximate dynamic
programming with a terminal set --- the outer table a
value-supersolution on a cell graph, the inner tail the polynomial
telescope. \textbf{Witnessed twice}: the pendulum swing-up row
(a $24{,}467$-cell supersolution, capture within $139$ steps, return
$\ge -619.99$, composed with the upright-hold invariance as absorbing
tail) and the lander attitude funnel. The sort stated in the
block below is forced, not chosen: a funnel that collapsed its parts
into a one-slot composite could not be a lawful compositor, because
collapsing composition breaks the unit law --- a machine-checked
refusal (\texttt{funnelComp\_breaks\_unit\_coherence}, beside its
no-tax counterpart \texttt{unitClass\_transparent}, module
\texttt{SuFunctor.Plain.Battery}) --- so the pattern composes
certificates about a System rather than forming one. The pattern's
positive semantics is machine-checked in the same tree, sorry-free
(module \texttt{SuFunctor.Adequacy.Funnel}): a funnel's composite
satisfies the coupled transit-then-hold reading as a pure
instantiation of the tower's composition theorem for budget sums with
certified dips (\texttt{funnel\_coupledSat\_of\_parts}, instantiating
\texttt{coupledSat\_comp\_add\_of\_dips} with no new induction), and
the tower-citation edge enters as an equation: assembly succeeds
precisely when transit and hold cite one gate
(\texttt{assemble\_isSome\_iff}), so a funnel whose parts disagree on
the gate is unrepresentable rather than refused at replay.

\syscatmetapat{\begin{mathpar}
\inferrule*[Right=\rulename{S-FUNNEL}]{
  \mathsf{sys} : \mathsf{SysRef} \\
  \mathsf{transit},\, \mathsf{hold} : \mathsf{SpecRef} \\
  \mathsf{gate} : \mathsf{RegionAtom} \text{ carrying a certified
  family} \\
  \mathsf{budget} \in \mathbb{Q}_{>0}
}{
  \mathrm{Funnel}(\mathsf{sys};\, \mathsf{transit};\, \mathsf{hold};\,
  \mathsf{gate};\, \mathsf{budget}) : \mathsf{FunnelDecl}
}
\end{mathpar}
The sort is forced: not a System former --- the pattern composes two
certificates about one subject System, cited by reference, and its
row is the composition data the tower-citation edge consumes
at replay. The gate premise excludes the bare box: the tower
discipline cites \emph{certified families} (sublevel or affine), and
a box region carries none to place the citation on. The budget
premise is the sign gate --- a funnel with no transit budget asserts
nothing about reaching the gate.}
\ruleblock{API}
\sourcenote{Producer surface of the \texttt{mtk} toolkit for this composition pattern: the TypeScript arrangement (npm package \texttt{@maitria/mtk}) and the Rust catalogue crate (\texttt{crates/mtk-catalogue}) --- one Rust implementation underneath, reached from TypeScript in-process through a native binding, so the two surfaces cannot diverge. Signatures and constructions are extracted from the toolkit sources at every regeneration of this catalogue; each construction is quoted verbatim from the toolkit's conformance corpus, code whose emitted certificate bytes are verified byte-identical against the committed golden vectors (helper values named \texttt{fx\_*} or \texttt{f.*} are that corpus's shared fixtures, defined beside the quoted lines). This API constructs and emits; it never checks --- the checker re-derives every claim independently.}
\apilang{TypeScript}{---}
The pattern-level TypeScript binding is not yet arranged: the manifest-level wiring descriptors (\texttt{WiringRow}, \texttt{SwitchEdge}, \texttt{ScheduleAssumption}) already ride the TS \texttt{system} namespace (\texttt{system.composite}), while the pattern builders and the four pattern row families are Rust-only at this preview.
\apilang{Rust}{mtk\_catalogue::compose::funnel}
\apiprov{signature --- \texttt{crates/mtk-catalogue/src/compose/funnel.rs}}
\begin{rustsrc}
pub fn funnel(
    sys: SysRef,
    transit: SpecRef,
    hold: SpecRef,
    gate: RegionAtom,
    budget: WitRat,
) -> Result<FunnelDecl, Refusal>
\end{rustsrc}
\apiprov{construction --- \texttt{crates/mtk-catalogue/tests/compose\_common/mod.rs}}
\begin{rustsrc}
funnel(
    SysRef("sys:pendulum-swingup".to_owned()),
    SpecRef("spec:transit-to-gate".to_owned()),
    SpecRef("spec:upright-hold".to_owned()),
    RegionAtom::Sublevel {
        cert: CertFamRef("vfam:telescope".to_owned()),
        thr: r(1, 1),
    },
    r(139, 1),
)
.expect("fixture conforms to the funnel pattern")
\end{rustsrc}

\section{The Twin-World Pattern}
\label{sec:cp-twin}

Two intervened copies of one model run against a single draw of the
exogenous sources --- the counterfactual composition. The pattern is
\emph{structural}: the causal-model profile owns one exogenous table
per mechanism (an option, not a list ---
section~\ref{sec:mt-scm}), and the pattern row carries one base
mechanism, so the decoupled independent-worlds reading is
unrepresentable rather than discouraged. The certified decoupling
probe prices exactly what that unrepresentability is worth ($0$
coupled against $\ge 1/2048$ decoupled for the same controller).
\textbf{Witnessed}: the counterfactual harm certificates in both
truth directions (section~\ref{sec:mt-scm}). The pattern's semantic
bridge is machine-checked in the toolkit's Lean tree
(\texttt{lean/qtsl/}, module \texttt{Qtsl.ScmTwin}): twin-network
substitution --- the counterfactual event as an ordinary
single-network event of the twin model, proven outright for acyclic
models --- together with the box-measure and union-bound certificate
classes the harm certificates instantiate, and the worked instance's
landmark values as drift alarms; a sorry-free development, its axiom
census the standard classical base.

\syscatmetapat{\begin{mathpar}
\inferrule*[Right=\rulename{S-TWIN}]{
  \mathsf{base} : \mathsf{Mechanism} \\
  \{\mathsf{equations},\, \mathsf{exogenous}\} \subseteq
  \mathrm{pop}(\mathsf{base}) \\
  \mathsf{base} \text{ un-intervened} \\
  \mathsf{left},\, \mathsf{right} :
  \overrightarrow{\mathsf{InterventionRow}} \text{, each }
  \mathsf{var} \text{ defined by a base equation, no var twice per
  world}
}{
  \mathrm{TwinWorld}(\mathsf{base};\, \mathsf{left};\,
  \mathsf{right}) : \mathsf{TwinWorldDecl}
}
\end{mathpar}
Two facts about the sort, both deliberate. The conclusion is a pattern
row, not a System. And the base's sort is $\mathsf{Mechanism}$, not
$\mathsf{System}$: a composite base is unrepresentable, which is the
one-exogenous-table structure stated as a premise sort --- the
coupling is \emph{through} that table. A pre-intervened base is
refused (world surgery is this pattern's own
$\mathsf{left}/\mathsf{right}$; anything else would smuggle a third
world in). An empty intervention list is a legal world --- the
factual side of a harm query intervenes on one side only.}
\ruleblock{API}
\sourcenote{Producer surface of the \texttt{mtk} toolkit for this composition pattern: the TypeScript arrangement (npm package \texttt{@maitria/mtk}) and the Rust catalogue crate (\texttt{crates/mtk-catalogue}) --- one Rust implementation underneath, reached from TypeScript in-process through a native binding, so the two surfaces cannot diverge. Signatures and constructions are extracted from the toolkit sources at every regeneration of this catalogue; each construction is quoted verbatim from the toolkit's conformance corpus, code whose emitted certificate bytes are verified byte-identical against the committed golden vectors (helper values named \texttt{fx\_*} or \texttt{f.*} are that corpus's shared fixtures, defined beside the quoted lines). This API constructs and emits; it never checks --- the checker re-derives every claim independently.}
\apilang{TypeScript}{---}
The pattern-level TypeScript binding is not yet arranged: the manifest-level wiring descriptors (\texttt{WiringRow}, \texttt{SwitchEdge}, \texttt{ScheduleAssumption}) already ride the TS \texttt{system} namespace (\texttt{system.composite}), while the pattern builders and the four pattern row families are Rust-only at this preview.
\apilang{Rust}{mtk\_catalogue::compose::twin\_world}
\apiprov{signature --- \texttt{crates/mtk-catalogue/src/compose/twin.rs}}
\begin{rustsrc}
pub fn twin_world(
    base: Mechanism,
    left: Vec<InterventionRow>,
    right: Vec<InterventionRow>,
) -> Result<TwinWorldDecl, Refusal>
\end{rustsrc}
\apiprov{construction --- \texttt{crates/mtk-catalogue/tests/compose\_common/mod.rs}}
\begin{rustsrc}
// Variable index space: endogenous x0, x1 then exogenous u (2).
let base = MechanismBuilder::new()
    .equation(0, mono(3, 2, 1, r(1, 1)))
    .equation(1, mono(3, 0, 1, r(1, 2)))
    .exogenous(ExogenousTable {
        vars: vec![ExoVar {
            id: 0,
            name: Some("noise".to_owned()),
            dist: Some(ExoDist::Box {
                region: BoxRegion {
                    axes: vec![[r(-1, 1), r(1, 1)]],
                },
            }),
        }],
    })
    .finish();
twin_world(
    base,
    vec![InterventionRow {
        var: 1,
        replacement: Some(constant(3, r(1, 1))),
    }],
    vec![],
)
.expect("fixture conforms to the twin-world pattern")
\end{rustsrc}

\section{The Stochastic-Wire Pattern}
\label{sec:cp-stochwire}

Systems coupled through probabilistic-dependency wires --- one
conditional-table wire per dependency: the frame family of the
graphical-model row (section~\ref{sec:mt-pgm}) and simultaneously a
composition mode for exogenous blocks feeding dynamical frames.
Single-system use is landed (the network rendering);
\textbf{cross-system use awaits the marginal-compatibility
discharge}, and the pattern row carries that residue explicitly as
data: the citation field for the compatibility certificate is
optional, and \emph{empty means pending}. That this residue matters
is itself machine-checked:
the classical obstruction --- pairwise-consistent marginal families
with no joint law (Vorob'ev) --- is mechanized over the
subprobability stock (\texttt{vorobev\_no\_joint}, module
\texttt{SuFunctor.Adequacy.MarginalGlue}: not even a mass-losing
joint exists), records glue where laws do not
(\texttt{vorobev\_record\_glue}), and an inconsistent family denotes
the bottom credal object rather than crashing
(\texttt{credalCompile\_vorobevGlue}) --- so an absent $\mu$ citation
marks a pending composite, never a silently wrong one.

\syscatmetapat{\begin{mathpar}
\inferrule*[Right=\rulename{S-STOCHWIRE}]{
  S_1, \ldots, S_k : \mathsf{System} \\
  W : \overrightarrow{\mathsf{CrossWire}} \\
  \mu? : \mathsf{CertFamRef}
}{
  \mathrm{StochasticCoupling}(S_1, \ldots, S_k;\, W) :
  \mathsf{StochasticCouplingDecl}
}
\end{mathpar}
\[
\mathsf{CrossWire} \;=\; \langle \mathsf{sources} :
\overrightarrow{\mathsf{PortRef}},\, \mathsf{target} :
\mathsf{PortRef},\, \mathsf{rows} :
\overrightarrow{\langle \mathsf{given}, \mathsf{outcome},
\mathsf{prob} \rangle} \rangle
\]
Per-wire formation data, checked wire by wire: sources resolve to
outputs and the target to an input where an interface is declared
(directed-only at this preview); the target is not among its own
sources; each row's $\mathsf{given}$ arity equals the source count;
an empty source list is a marginal (an exogenous block feeding a
frame unconditionally); and --- the exact-arithmetic gate ---
per-$\mathsf{given}$ probabilities sum to one, exactly, with the
computed sum as the refusal witness. $\mu?$ is the
marginal-compatibility citation, consumed without appearing in the
head form; its absence \emph{is} the pending residue, carried as
data. The sort is deliberate: this is the one pattern that consumes
Systems without
forming one --- the coupling is certified composition data over its
roster, and the composite System awaits exactly the discharge $\mu$
cites.}
\ruleblock{API}
\sourcenote{Producer surface of the \texttt{mtk} toolkit for this composition pattern: the TypeScript arrangement (npm package \texttt{@maitria/mtk}) and the Rust catalogue crate (\texttt{crates/mtk-catalogue}) --- one Rust implementation underneath, reached from TypeScript in-process through a native binding, so the two surfaces cannot diverge. Signatures and constructions are extracted from the toolkit sources at every regeneration of this catalogue; each construction is quoted verbatim from the toolkit's conformance corpus, code whose emitted certificate bytes are verified byte-identical against the committed golden vectors (helper values named \texttt{fx\_*} or \texttt{f.*} are that corpus's shared fixtures, defined beside the quoted lines). This API constructs and emits; it never checks --- the checker re-derives every claim independently.}
\apilang{TypeScript}{---}
The pattern-level TypeScript binding is not yet arranged: the manifest-level wiring descriptors (\texttt{WiringRow}, \texttt{SwitchEdge}, \texttt{ScheduleAssumption}) already ride the TS \texttt{system} namespace (\texttt{system.composite}), while the pattern builders and the four pattern row families are Rust-only at this preview.
\apilang{Rust}{mtk\_catalogue::compose::stochastic\_coupling}
\apiprov{signature --- \texttt{crates/mtk-catalogue/src/compose/stochwire.rs}}
\begin{rustsrc}
pub fn stochastic_coupling(
    components: Vec<SystemManifest>,
    wires: Vec<CrossWire>,
    marginal_compat: Option<CertFamRef>,
) -> Result<StochasticCouplingDecl, Refusal>
\end{rustsrc}
\apiprov{construction --- \texttt{crates/mtk-catalogue/tests/compose\_common/mod.rs}}
\begin{rustsrc}
let exo_block = {
    let m = MechanismBuilder::new()
        .interface(Interface {
            inputs: vec![],
            outputs: vec![wire(0, "w")],
        })
        .stochastic_wire(StochasticWire::Conditional {
            id: 0,
            sources: vec![],
            target: 0,
            rows: vec![
                CondRow {
                    given: vec![],
                    outcome: 0,
                    prob: r(1, 3),
                },
                CondRow {
                    given: vec![],
                    outcome: 1,
                    prob: r(2, 3),
                },
            ],
        })
        .finish();
    presets::Pdg::manifest(m)
        .expect("exogenous block conforms to the PDG profile")
        .named("weather")
};
let frame = {
    let m = MechanismBuilder::new()
        .state_var(0, Some("x"))
        .interface(Interface {
            inputs: vec![wire(0, "w")],
            outputs: vec![],
        })
        .drift(None, 0, mono(1, 0, 1, r(-1, 2)))
        .finish();
    presets::Ode::manifest(m)
        .expect("frame conforms to the ODE profile")
        .named("field")
};
stochastic_coupling(
    vec![exo_block, frame],
    vec![CrossWire {
        sources: vec![port(0, 0)],
        target: port(1, 0),
        rows: vec![
            CondRow {
                given: vec![0],
                outcome: 0,
                prob: r(1, 3),
            },
            CondRow {
                given: vec![0],
                outcome: 1,
                prob: r(2, 3),
            },
            CondRow {
                given: vec![1],
                outcome: 0,
                prob: r(1, 1),
            },
        ],
    }],
    Some(CertFamRef("qmc:exo-frame".to_owned())),
)
.expect("fixture conforms to the stochastic-wire pattern")
\end{rustsrc}

\section{Credal-Hull Certification}
\label{sec:cp-credal}

The two grades of open nondeterminism as one pattern. A system with
unresolved choice --- a policy slot, a scheduler slot --- is certified
either \emph{at a declared resolution} (the resolution witness
composed in; the closed object an ordinary citizen of the base
frame's rows --- the polling-automaton closure of
section~\ref{sec:mt-ma} is the landed instance) or \emph{over the
whole hull}: per-state-action row families asserting the generator
inequality for \emph{every} enabled choice, which covers
deterministic, randomized (convexity --- the rows are the simplex's
extreme points), and history-dependent resolutions at once. The
declared grade is the hull grade restricted to one vertex; the
certified gap between them is real and priced (the demonic fixtures
carry named breaking policies). Denotation: the credal lower
expectation of the imprecise-semantics layer --- and the semantics
side of this pattern carries an adequacy receipt of its own: the
demonic-face square of the decision-process section
(section~\ref{sec:mt-mdp}) mechanizes, at each finite horizon, that
the compiled, policy-free presentation determines exactly the
closed-loop law set over the memoryless randomized class, with the
event-bound verdicts uniform over that class. \textbf{Witnessed}: a
cliff-walk demonic pair (uniform bound with tightness witness), and
the polling automaton's uniform re-certification, which thinned the
declared grade's slack from $5.13 \times 10^{-3}$ to eight parts per
million without moving the bound.

\syscatmetapat{\begin{mathpar}
\inferrule*[Right=\rulename{declared}]{
  \mathsf{sys} : \mathsf{SysRef} \\
  m : \mathsf{Mechanism} \text{ the subject},\
  E(m) \neq \emptyset \\
  \mathsf{cert} : \mathsf{CertFamRef} \\
  \rho : \overrightarrow{\langle \mathit{state},\, \mathit{action}
  \rangle} \text{, exactly one enabled action per open state}
}{
  \mathrm{CredalHull}(\mathsf{sys};\, \mathsf{cert};\,
  \mathsf{declared}(\rho)) : \mathsf{CredalHullDecl}
}
\and
\inferrule*[Right=\rulename{quantified}]{
  \mathsf{sys} : \mathsf{SysRef} \\
  m : \mathsf{Mechanism} \text{ the subject},\
  E(m) \neq \emptyset \\
  \mathsf{cert} : \mathsf{CertFamRef} \\
  h : \overrightarrow{\langle \mathit{state},\, \mathit{action},\,
  \mathit{driftBound} \rangle} \text{, exactly one row per pair of }
  E(m)
}{
  \mathrm{CredalHull}(\mathsf{sys};\, \mathsf{cert};\,
  \mathsf{quantified}(h)) : \mathsf{CredalHullDecl}
}
\end{mathpar}
$E(m)$ is the subject's enabled choice set --- $\mathit{state}
\mapsto \mathit{actions}$ read off its action-keyed outcome rows; a
subject with no action-keyed structure is refused outright (no open
choice --- nothing for either grade to range over). The subject is
consumed as formation data without appearing in the head form: it
travels in the row by reference ($\mathsf{sys}$), and the coverage
premises --- one action per state for the vertex grade, since a
partial resolution is not a vertex; one row per enabled pair for the
hull grade, a gap's address serialized on refusal --- are checked
against the mechanism itself. Drift bounds are exact and sign-free
(a supermartingale certificate's generator values are typically
nonpositive); the generator-inequality \emph{content} is the cited certificate family's
business, recomputed at replay.}
\ruleblock{API}
\sourcenote{Producer surface of the \texttt{mtk} toolkit for this composition pattern: the TypeScript arrangement (npm package \texttt{@maitria/mtk}) and the Rust catalogue crate (\texttt{crates/mtk-catalogue}) --- one Rust implementation underneath, reached from TypeScript in-process through a native binding, so the two surfaces cannot diverge. Signatures and constructions are extracted from the toolkit sources at every regeneration of this catalogue; each construction is quoted verbatim from the toolkit's conformance corpus, code whose emitted certificate bytes are verified byte-identical against the committed golden vectors (helper values named \texttt{fx\_*} or \texttt{f.*} are that corpus's shared fixtures, defined beside the quoted lines). This API constructs and emits; it never checks --- the checker re-derives every claim independently.}
\apilang{TypeScript}{---}
The pattern-level TypeScript binding is not yet arranged: the manifest-level wiring descriptors (\texttt{WiringRow}, \texttt{SwitchEdge}, \texttt{ScheduleAssumption}) already ride the TS \texttt{system} namespace (\texttt{system.composite}), while the pattern builders and the four pattern row families are Rust-only at this preview.
\apilang{Rust}{mtk\_catalogue::compose::quantified\_hull}
\apiprov{signature --- \texttt{crates/mtk-catalogue/src/compose/credal.rs}}
\begin{rustsrc}
pub fn declared_resolution(
    sys: SysRef,
    subject: &Mechanism,
    cert: CertFamRef,
    rows: Vec<ResolutionRow>,
) -> Result<CredalHullDecl, Refusal>

pub fn quantified_hull(
    sys: SysRef,
    subject: &Mechanism,
    cert: CertFamRef,
    rows: Vec<HullRow>,
) -> Result<CredalHullDecl, Refusal>
\end{rustsrc}
\apiprov{construction --- \texttt{crates/mtk-catalogue/tests/compose\_common/mod.rs}}
\begin{rustsrc}
let big = WitRat::new(
    -"1234567890123456789012345678901234567890"
        .parse::<BigInt>()
        .unwrap(),
    BigInt::from(7u8),
);
quantified_hull(
    SysRef("sys:cliffwalk-2".to_owned()),
    &mdp_subject(),
    CertFamRef("vfam:cliffwalk-lower".to_owned()),
    vec![
        HullRow {
            state: 0,
            action: 0,
            drift_bound: r(-1, 8),
        },
        HullRow {
            state: 0,
            action: 1,
            drift_bound: big,
        },
        HullRow {
            state: 1,
            action: 0,
            drift_bound: r(0, 1),
        },
    ],
)
.expect("fixture conforms to the quantified grade")
\end{rustsrc}
\apiprov{The declared grade (one vertex of the hull) is the corpus's \texttt{credal\_declared\_decl} sibling fixture.}

\section{Event-Port Patterns}
\label{sec:cp-eventport}

A finished design whose primitive now carries a producer surface.
Jump synchronization --- two mechanisms' events declared to fire
together --- is constructible today: the toolkit's composite stratum
stores typed fusion rows (capability legs across components'
population edges, one firing law per event under the single-home
discipline below), the pattern-level builder gates them, the
coloured-net decomposition emits them, and the API block at the end
of this section quotes the conformance constructions in both
languages. Everything downstream of the primitive --- resource
pools, broadcast, dispatch, rate reasoning at fused boundaries ---
still speaks in design tense, settled to typing-rule depth (for the
fusion core and resource pools, to certificate-transport depth) and
recorded here so the remaining vocabulary lands on a fixed target;
until it lands, synchronized structure beyond raw fusion lives
inside single mechanisms, a granularity cost priced in the
design's deferral note with its re-pricing trigger named.

\subsection{Fusion}

The primitive. Mechanisms expose selected transitions as \emph{event
ports} on their interfaces; a fusion row names a set of bound ports
across a wiring, and the named transitions fire as one joint jump ---
a rendezvous rather than a broadcast~\cite{cleaveland18hbmc}. Enabling is conjunctive over the parties'
local guards, evaluated on the pre-jump composite state (snapshot
semantics: each party reads the snapshot and writes only its own
state, so resets within a fusion commute by construction). Every port
carries a \emph{modality}: at most one party in a fusion is
\emph{active} --- the timing owner, whose transition kind (rate law,
forced guard, or immediate) times the joint jump --- while passive
parties contribute enabling conjuncts and resets only, and never
fire outside a fusion. The design's worked example is operator~1 of
the Mini-Fab model (the standard five-machine semiconductor
benchmark, chapter~\ref{ch:qtsl-minifab}), who serves both
lithography and implantation: ``load starts'' is one jump across
station and operator --- machine ready $\wedge$ lot present $\wedge$
operator idle --- with the operator's idle$\to$task reset riding the
same jump. The fusion row has landed on the producer side, with the
decomposition feeding it: the toolkit's coloured-net decomposition
(section~\ref{sec:mt-sdcpn}) derives, for each transition, a fused
event over the incident places' population edges --- the legs as
(component, edge) pairs, enabling arcs riding as read legs, the one
firing law on the row --- and emits it as the serialized fusion row
the composite stratum stores and the pattern builder gates.

\subsection{Resource Pools}

A packaged pattern over fusion, not a new primitive: a shared
resource is a pool mechanism (a counter $k \in \{0,\dots,C\}$, or a
one-place coloured pool) with passive \emph{take}/\emph{give} ports,
and a resource-share row family definitionally expands into that
mechanism plus fusion rows binding each user's seize and release
transitions. A crossing discipline (each user declares which of its
modes count as \emph{holding} the resource; seizes enter the held
set, releases leave it, no other transition crosses) makes the
pattern's payoff a theorem rather than a convention --- the
\emph{conservation law}: for a well-formed pool share, every jump of
the composite preserves
\[
  k_{\mathrm{pool}} \;+\; \sum_{u \in \mathrm{users}} \#H_u \;=\; C,
\]
where $\#H_u$ counts user $u$'s units in held modes. The invariant is
a count row~\cite{master25}, free by construction; raw fusion rows can express the
same model and can also silently break the invariant --- the packaged
rows cannot, and that refusal is what the package is for.

\subsection{Broadcast Fan-Out}

One sender, many receivers. The inventory deliberately carries no
broadcast primitive; it carries the \emph{fixed-roster reduction}:
when the receiver roster is finite and statically known --- which is
every use the catalogue's source models have produced --- a broadcast
is a finite family of rendezvous fusions, and optional-join broadcast
(each receiver participates only when locally enabled) is a finite
family indexed by the participating subsets. Genuinely dynamic
rosters are priced and deferred with the rest of the stratum.

\subsection{Dispatch: Control Through Selection}

Fusion creates contention: two fusions sharing a passive party
conflict, and the enabled-conflict set is derivable from the fusion
rows alone. Absent further data the conflict resolves
demonically~\cite{baier18hbmc} --- sound, and exactly as weak as it
sounds. A dispatch row installs a
policy: a chooser over the conflict set, reading current aggregate
observables. Certificates are dispatch-indexed, with one free lemma
(any dispatched policy refines the demonic envelope, so demonic
safety transports to every policy) and one structural refusal
(recurrence-class conclusions never cross an undispatched boundary).
The design's worked example is first-come-first-served against
setup-affinity dispatch on the Mini-Fab model: two policies over an
identical composite, with per-policy certified throughput the target
--- the setup-affinity policy's advantage being exactly the kind of
claim that is only ever \emph{per-policy}, since the demonic envelope
proves neither. Dispatch is the \emph{per-epoch} member of the
inventory's control trio; its persistent sibling --- control through
enablement, where a supervisor carves the enabled set itself rather
than selecting within it --- is the supervision pattern
(section~\ref{sec:cp-supervise}), and the two compose because they
are disjoint by kind.

\subsection{Delayed Trace and the Delay Media}
\label{sec:cp-dtrace}

Feedback with memory: a delayed-trace row~\cite{sprunger19a} closes a
composite's output
back onto its input through a delay, well-defined precisely because
the delay breaks instantaneous self-reference~\cite{dangelo25}.
Delays themselves are
first-class composition elements~\cite{lavore22} in two species with a shared
degenerate case --- an event species (a one-place coloured-token
place whose tokens carry a time-to-go coordinate) and a signal
species (a transport medium: pure advection over the delay interval),
both degenerating, when time is discrete and the delay is one step,
to the \emph{tick register}: the unit delay, which is also the
minimal medium a delayed trace needs. The family exists so that
delays compose like everything else --- a delay is a mechanism with
ports, not an annotation on a wire. For the signal species in
continuous time, the semantics tier is already receipted: the
delay-line adequacy square of the delay-differential-equations
section (section~\ref{sec:mt-dde}) mechanizes exactly this
composite --- an ODE cell closed through the transport medium ---
so the pattern arrives with its continuous-time semantics a theorem
and its producer surface landed: the delayed trace is a wiring, and
the toolkit's conformance corpus constructs the same composite ---
cell, line, and the two trace wires --- through the composition
surface (a dedicated delayed-trace row stays deliberately absent;
the wiring is the pattern). The design's flagship
construction is the partially observable decision process:
environment mechanism, observation channel, and a belief-tracking
controller closed through a delayed trace --- partial observability
as a \emph{composition} pattern rather than a new model class
(section~\ref{sec:mt-pomdp} grades the frame) --- with the first
certificate, a lower bound on achieved value for the classic
two-door problem living on the belief simplex, already exercised in
the workshop.

\subsection{Event-Rate Assume/Guarantee}

Rate reasoning across fused boundaries: an event-rate atom (the
specification catalogue's envelope pair, staged ---
section~\ref{sec:spec-lang-rate}) bounds the
firing rate of a shared event, a mechanism's guarantee about its side
of a fusion becomes its neighbour's assumption, and assume/guarantee
chains propagate rate bounds through a composite without unfolding
it~\cite{brieger24}. The delay place enters as a \emph{rate-preserving element} ---
tokens are conserved, so the departure rate equals the arrival rate
shifted by the delay, and a rate assumption crosses a delay medium
unchanged. This is the pattern that would let throughput claims
localize: bound the bottleneck's service rate, and let
rate-preservation carry the bound to the boundary where the
specification states it.

\syscatmetapat{\emph{Producer-landed, theory-pending}: the fusion row family now
has a builder underneath (the API block below quotes it), but no
term former is stated ahead of its certificate-schema vocabulary ---
the formal block enters with the event-port interface stratum's
staging, and asserting a rule now, with no loader beneath it to
refuse on its premises, would be exactly the invented formal block
the first wing's frontier entries refuse.}
\ruleblock{API}
\sourcenote{Producer surface of the \texttt{mtk} toolkit for this composition pattern: the TypeScript arrangement (npm package \texttt{@maitria/mtk}) and the Rust catalogue crate (\texttt{crates/mtk-catalogue}) --- one Rust implementation underneath, reached from TypeScript in-process through a native binding, so the two surfaces cannot diverge. Signatures and constructions are extracted from the toolkit sources at every regeneration of this catalogue; each construction is quoted verbatim from the toolkit's conformance corpus, code whose emitted certificate bytes are verified byte-identical against the committed golden vectors (helper values named \texttt{fx\_*} or \texttt{f.*} are that corpus's shared fixtures, defined beside the quoted lines). This API constructs and emits; it never checks --- the checker re-derives every claim independently.}
The fusion primitive's producer surface: a fusion row names capability legs (bag edges across a composite's components, read legs included) that fire as one joint jump, with the event's one firing law under the single-home discipline --- exactly one active leg (its edge's own law times the jump) xor a row-resident law over all-passive legs; both double-home and all-passive-lawless declarations refuse at build, along with leg resolution, membership (active-never-optional; a row-law fusion keeps a mandatory leg), and the delay basis (selection blocks per consuming/read leg, within-leg slot-exchange symmetry). The packaged patterns over fusion --- resource pools, dispatch, event-rate atoms --- remain design-tense and carry no API here.
\apilang{TypeScript}{compose.composite / system.composite (fusions)}
\apiprov{signature --- \texttt{ts/src/system.ts}}
\begin{tssrc}
/** S-EVENTPORT: one capability leg of a fused event — `edge` of
 * `components[component]`'s Bag self-loop roster. A leg whose edge
 * is `synced` is passive; an edge carrying its own local law is the
 * fusion's active leg (the timing owner). `membership` defaults to
 * `"mandatory"`; an `"optional"` leg joins only when locally enabled
 * (maximal join). */
export interface FusionLeg {
  readonly component: number;
  readonly edge: number;
  readonly membership?: "mandatory" | "optional";
}

/** S-EVENTPORT: the row-resident firing law of an all-passive fusion
 * — delay (rate per unordered selection across the consuming/read
 * legs' blocks, in leg order), guard (forced; every consuming/read
 * leg has delta 1), or immediate (the eager case, spelled
 * explicitly). */
export type FusionLaw =
  | { readonly kind: "delay"; readonly rate: Polynomial }
  | { readonly kind: "guard"; readonly guard: ReadonlyArray<DomainConstraint> }
  | { readonly kind: "immediate" };

/** S-EVENTPORT: one event-port fusion row — the named capability
 * legs fire as one joint jump (rendezvous), with the event's one
 * firing law under the single-home discipline: exactly one active
 * leg (`law` absent) xor all legs passive with `law` on the row. */
export interface FusionRow {
  readonly event: number;
  readonly legs: ReadonlyArray<FusionLeg>;
  readonly law?: FusionLaw;
}
\end{tssrc}
\apiprov{construction --- \texttt{ts/tests/compose-native.test.ts}}
\begin{tssrc}
// The transfer rendezvous with an operator: a token leaves
// `upstream` and enters `downstream` as ONE joint jump, with an
// operator token read (never consumed). All legs are passive, so
// the law rides the row: a delay rate over the derived selection
// basis — the consumed token's colour, then the read operator's.
return compose.unwrap(
  compose.composite({
    name: "transfer-rendezvous",
    components: [
      bagStation(sys, "upstream", "destroy", 2),
      bagStation(sys, "downstream", "create", 1),
      bagStation(sys, "operators", "read", 1),
    ],
    fusions: [
      {
        event: 0,
        legs: [
          { component: 0, edge: 0 },
          { component: 1, edge: 0 },
          { component: 2, edge: 0 },
        ],
        law: { kind: "delay", rate: mono(2, 0, 1, "1/10") },
      },
    ],
  }),
);
\end{tssrc}
\apiprov{\texttt{bagStation} and \texttt{mono} are the conformance corpus's shared fixtures, defined beside the quoted lines; the quoted construction's emitted bytes are the committed \texttt{compose\_fusion} golden vector, byte-identical across both language surfaces.}
\apilang{Rust}{mtk\_catalogue::compose::CompositeBuilder::fuse}
\apiprov{signature --- \texttt{crates/mtk-catalogue/src/system/compose.rs}}
\begin{rustsrc}
/// System Catalogue: one event-port fusion row (`S-EVENTPORT`) — a
/// set of capability legs across the composite's components firing
/// as one joint jump (rendezvous, not broadcast), plus the event's
/// one firing law under the **single-home** discipline: either
/// exactly one leg is active (its edge's own law times the joint
/// jump; `law` absent) **xor** all legs are passive and `law` is
/// present on the row. Both double-home and all-passive-lawless
/// declarations refuse at [`CompositeBuilder::build`].
///
/// Enabling is conjunctive over the legs' local guards on the
/// pre-jump snapshot; each leg writes only its own component's
/// state (snapshot semantics, so reset commutation within a fusion
/// is a theorem — kernel-side, like every transport rule here).
///
/// [`CompositeBuilder::build`]: crate::compose::CompositeBuilder::build
#[derive(Debug, Clone, PartialEq, Serialize, Deserialize)]
pub struct FusionRow {
    /// Event id (roster-unique among the composite's fusion rows; for
    /// a decomposition image, the source SDCPN transition id).
    pub event: u32,
    /// The fused capability legs.
    pub legs: Vec<FusionLeg>,
    /// The row-resident firing law — present iff every leg is
    /// passive (the single-home XOR).
    #[serde(default, skip_serializing_if = "Option::is_none")]
    pub law: Option<FusionLaw>,
}
\end{rustsrc}
\apiprov{construction --- \texttt{crates/mtk-catalogue/tests/compose\_builders.rs}}
\begin{rustsrc}
// The transfer rendezvous: a token leaves `a` and enters `b` as
// ONE joint jump. Both capability legs are passive (`Synced`
// edges), so the event's one firing law rides the row — a delay
// rate over the derived selection basis (the consumed token's
// colour coordinates).
let sys = CompositeBuilder::new()
    .component(bag_station("a", BagEdgeKind::Destroy, 2))
    .component(bag_station("b", BagEdgeKind::Create, 1))
    .fuse(FusionRow {
        event: 0,
        legs: vec![leg(0, 0), leg(1, 0)],
        law: Some(FusionLaw::Delay {
            rate: mono(1, 0, 1, r(1, 10)),
        }),
    })
    .build()
    .expect("legs resolve; all passive; law on the row");
\end{rustsrc}
\apiprov{\texttt{bag\_station} and \texttt{mono} are the corpus's shared fixtures (\texttt{tests/compose\_common}); \texttt{FusionLeg}/\texttt{FusionLaw} and the membership column are declared beside the quoted row type, and the builder verb (\texttt{fuse}) runs the well-formedness gates at \texttt{build}.}

\section{Supervision: Control Through Enablement}
\label{sec:cp-supervise}

Partially observed supervisory
control~\cite{ramadge1987supervisory,lin88} over the coloured-net frame: a
\emph{supervisor} that can switch \emph{some}
declared plant transitions between enabled and disabled --- not all
of them, and never by firing them itself --- that sees only declared
observable functions of parts of the plant's marking, and that
springs into action only when subscribed events occur. The entry's
organizing fact is that the control seam costs nothing, because the
inventory already carries it: \textbf{the gate is a place, and the
gating is an enabling arc}. The policy is an ordinary jump-only
System that owns places; a gated transition carries an ordinary
weight-1 enabling arc (read-no-consume) from its gate place; enabling
is minting the gate token and disabling is draining it --- both
ordinary jumps of the policy itself. No new arc kind, wire class, or
row family exists for this pattern to add: enabling arcs are
first-class in the coloured-net vocabulary at every stratum
(section~\ref{sec:mt-sdcpn}'s arc table and the three faces of its
adequacy development), and the inventory's seams for shaping a
\emph{live} component's behaviour thereby form a trio at three tempos
--- parameter setting (continuous, section~\ref{sec:cp-set}),
\emph{enablement} (persistent, this entry), dispatch (per-epoch,
section~\ref{sec:cp-eventport}); the supervisory \emph{edges} of mode
switching (section~\ref{sec:cp-switch}) sit outside the trio, handing
the active role between components rather than steering an active
one. Gates carve the enabled set; dispatch selects within what
survives; uncovered freedom stays demonic.

What the packaged row genuinely adds is bookkeeping with one new
census clause. The package declares the ownership split --- which
places and which update transitions are the policy's --- and derives
the gated roster from it (every read arc sourced at a policy place;
row-level, decidable, never separately declared). Its well-formedness
census closes the information channels: policy inputs factor through
the declared observables, subscriptions are the declared events, and
--- the clause the place-rendering itself demands ---
\emph{arc separation}: plant transitions touch policy places through
the declared gating read arcs only, with ordinary arcs refused in
both directions (a plant transition producing into a gate place would
be plant-side control of the policy's state), and policy updates
carry no arcs on plant places at all --- the policy acts on the plant
through enablement, full stop. Raw rows remain lawful with the census
as admission check, as for the inventory's other packaged rows.
Zero-order hold between wakes is not even a lemma under this
rendering: read arcs cannot move tokens, so no plant firing can touch
a gate, and token presence has no drift to have. Gate places carry a
Boolean discipline by default --- capacity 1, gates as switches ---
and the discipline is deliberately omittable: $k$ tokens in a gate
place multiply a gated race's intensity by exactly $k$ (the ordered
census's falling factorial at the gate place), so integer throttles
come free by \emph{not} imposing it, while fractional modulation
remains parameter setting on rate content. The capacity declaration
is accordingly a \emph{must} wherever the model means Boolean gates:
without it, an accidental second token silently doubles a rate.

The semantics per gated species is read off the landed faces rather
than added to them. A gated delay transition is race thinning: gate
closed, the ordered census vanishes (pre-enabling was always inside
the count); gate open, the survivor law is the integrated-rate law,
by memorylessness. A gated immediate transition is already inside the
cascade square: gate coordinates are marking coordinates, so the
tangible/vanishing classification, the scheduler discipline, and the
cascade-rank certificate quantify over them unmodified --- and the
policy's own wake is a cascade step like any other, so one structural
rank discharges non-Zeno for every enabling-respecting scheduler,
policy updates included. The whole pattern lives strictly inside the
monotone fragment: disabling by \emph{removing} a token is not a
zero-test, enabledness stays upward closed (surplus gate tokens never
disable a step), and the inhibitor exclusion is untouched ---
supervision is the monotone-safe rendering of control. Closure is on
the nose: the supervised composite --- plant places plus policy
places, gating arcs, fused wakes --- is again a coloured net of the
same class, so supervision is composition \emph{inside} the model
class, and the fixed-policy adequacy statement is the net square
applied to the supervised composite, no pattern-specific vertical.
Mechanized at the uncoloured stratum
(\texttt{SuFunctor.Adequacy.StochasticPetriSupervise}): the gating
lemmas, the decidable arc-separation census with refusal witnesses in
both directions, the Boolean-discipline preservation, and supervised
micro-instances whose controls pin the closed, open, and
graded race intensities as values, drive a two-step cascade
\emph{through} the policy's own mint, and consume the landed cascade
adequacy square at the supervised composite as an instance.

The pattern is a deliberate fourfold weakening of the probabilistic
model checker's scheduler --- the adversary that tools in the PRISM
and Storm family
resolve~\cite{kwiatkowska2011prism,hensel2022storm} --- and each
axis has a formal cash-out, landed or staged as the Status note
records:

\begin{center}\small
\begin{tabular}{@{}lp{4.2cm}p{6.2cm}@{}}
\toprule
axis & scheduler (limit case) & supervised variant \\
\midrule
control scope & resolves all nondeterminism & read arcs on a declared
roster; races stay stochastic, uncovered conflicts stay demonic ---
the controllable/uncontrollable partition is roster data \\
\addlinespace
effect & the choice \emph{is} the firing & the policy's action type
is its own marking: mint and drain in places it owns; the plant's
races and cascades do the firing --- and the composite is again a net
of the same class, where a scheduler lives outside its model \\
\addlinespace
observation & reads full state & reads declared observables of
declared parts; measurability transposed into typing, checked at
formation --- adaptedness then a theorem about every well-formed
composite, staged with the observation legs (Status note below) \\
\addlinespace
activation & consulted at every epoch & wakes only at subscribed
events; hold between wakes is read-no-consume; decision epochs are a
declared sub-point-process of the plant's event stream, non-Zeno by
the shared rank certificate \\
\bottomrule
\end{tabular}
\end{center}

Certificate polarity splits cleanly over the policy family, and the
split is the entry's payoff. \emph{Safety is free across the whole
family}: monotone enabledness makes refinement monotone in the gated
roster, every well-formed policy's closed-loop behaviour is contained
in the demonic gate-envelope, so one safety certificate against the
gate-demon transports to every partially observing policy --- the
information constraint costs nothing at this polarity.
\emph{Liveness over the unrestricted family is unprovable by
construction}: the empty policy --- no updates, gates never minted
--- is well-formed, so deliberate idling is expressible, and
throughput claims are false over the family wherever a gated
transition is load-bearing. That expressiveness is the commission,
not a defect: deliberate idling is exactly the capability dispatch
lacks (a chooser must choose; a supervisor may refuse), and it is why
liveness certificates are per-policy here, exactly as they are
per-dispatch in the previous entry, Kumar--Seidman teeth
included~\cite{kumar1990dynamic}.
The ``some transitions, not all'' clause has the same certificate-side
reading: a total roster would make safety-under-supervision vacuous
of plant content and family liveness never provable; partial rosters
keep both polarities informative.

One refusal survives, and it is the pattern's certificate-soundness
story. At formation, gating is admitted on \emph{every} transition
species --- including boundary-forced guard transitions, whose
semantics stays total when disabled: a guard transition that is not
pre-enabled imposes no forcing, and the run continues past its
boundary untouched. What must not be gated, on pain of
certificate unsoundness rather than semantic undefinedness, is a
transition serving as a model's \emph{box-discharge witness}: the
certificate layer quantifies colours over compact rational boxes, the
discharge obligation is what keeps the process inside the box the
rows quantify over, and its default witness is itself an exit-guard
transition. Gate that transition off and the process exits its box
with positive probability while the certificate's rows still
quantify over it. The obligation therefore lives at the certificate
stratum, where the discharge already does: a gated transition may not
serve as a box-discharge witness unless the model re-discharges by
invariance or by stopped-process reading --- a row-level lint on the
package roster, decidable, and binding however the arcs were wired,
because it is checked on the composite's rows rather than on any
declared seam. On the semiconductor-fab model
(chapter~\ref{ch:qtsl-minifab}) the lint draws exactly the right
line: a \emph{suppressible deadline} is directly representable ---
gate an interior maintenance-window transition and the model means
what it says --- while the shift-boundary transitions are ungateable
not because they are calendars but because they are their timer
boxes' own discharge witnesses.

\textbf{Status}: the gating semantics and the arc-separation census
are mechanized, sorry-free, at the uncoloured stratum (the module
above, over the statics square and both dynamics faces), with refusal
witnesses on every census clause and the graded-throttle identity
pinned as a value; the coloured packaged row --- observation legs as record-lens
readouts, wakes as passive optional-join fusions --- stages with the
event-port interface stratum
(section~\ref{sec:cp-eventport}), and until it lands, supervised
composites are hand-wired through the census door. The demonic
gate-envelope face --- the family-level object that safety
certificates ride --- is mechanized over the credal-hull machinery
(section~\ref{sec:cp-credal}) at the uncoloured stratum: the
certified policy seat is a dependent-index family (each member
carrying its own termination rank --- no automaton-level certificate
can be uniform across the seat), its closed-loop laws inhabit one
envelope on the common marking space, safety verdicts transport
uniformly over the seat at both polarities, the envelope is monotone
in the gated roster, and deliberate idling is a member --- the
mechanized face of the polarity split above.

\syscatmetapat{\emph{Former pending its strata; census landed.} The packaged row
consumes an observation interface and a subscription list whose sorts
stage with the event-port interface stratum
(section~\ref{sec:cp-eventport}), so no term former is asserted ahead
of them --- but unlike the entries that refuse their blocks outright,
this pattern's one genuinely new formation premise already has a
loader to refuse on it: the arc-separation census is a decidable
row-level predicate, mechanized with refusal witnesses in both
directions (a plant write into a gate place; a policy arc on a
plant place). The gate map itself adds no premise sort at all --- a
gate is a weight-1 enabling arc from a policy-owned place, an arc
kind every stratum of the coloured-net development already carries
--- and the ownership split is package data, invisible to the wire
vocabulary.}
\ruleblock{API}
\sourcenote{Producer surface of the \texttt{mtk} toolkit for this composition pattern: the TypeScript arrangement (npm package \texttt{@maitria/mtk}) and the Rust catalogue crate (\texttt{crates/mtk-catalogue}) --- one Rust implementation underneath, reached from TypeScript in-process through a native binding, so the two surfaces cannot diverge. Signatures and constructions are extracted from the toolkit sources at every regeneration of this catalogue; each construction is quoted verbatim from the toolkit's conformance corpus, code whose emitted certificate bytes are verified byte-identical against the committed golden vectors (helper values named \texttt{fx\_*} or \texttt{f.*} are that corpus's shared fixtures, defined beside the quoted lines). This API constructs and emits; it never checks --- the checker re-derives every claim independently.}
\apilang{TypeScript}{---}
As for Rust: enters with the event-port interface stratum.
\apilang{Rust}{---}
None yet, and the absence is a sequencing decision rather than a gap: supervision is a packaged pattern over fusion (event-triggered policy wakes) and enabling arcs (the gating), so its producer surface is sequenced behind the event-port interface stratum rather than built ahead of it. The pattern's semantic tier is already mechanized in the toolkit's Lean tree --- the gating lemmas, the arc-separation census, and the supervised micro-instance with its cascade controls (module \texttt{SuFunctor.Adequacy.StochasticPetriSupervise}, over the statics square and both dynamics faces of the stochastic-Petri development) --- so the builder, when it lands, arrives to a fixed semantic target.

\section{State Sharing}
\label{sec:cp-stateshare}

The reserved axis: gluing along shared \emph{state} --- a jointly
driven continuous variable, common noise --- rather than signal
wires. The designed semantics is the generator's additivity over
transitions sharing state,
\[
\mathcal{A}_{\mathrm{composite}} \;=\; \sum_i \mathcal{A}_i,
\]
so certificates compose by summing generator contributions; the
design warrant is recorded, and the axis is reserved rather than
staged. Its first customers: common-noise mean-field composition,
should that frontier entry of chapter~\ref{ch:qtsl-mathematical-models} ever
earn a producer surface, and the shared-place face of net
composition --- two nets glued along one common place --- which is
the one part of ``nets of nets'' that genuinely lives on this axis
(the rest turned out not to need it;
section~\ref{sec:cp-netsofnets}).

\subsection{Continuous State Sharing: Port-Hamiltonian
  \texorpdfstring{\mbox{Interconnection}}{Interconnection}}

The one pattern in the inventory's design space that is not reducible
to fusion: two mechanisms continuously driving one shared state have
no events to synchronize --- the content is drift superposition. The
disciplined form is port-Hamiltonian interconnection: each mechanism
carries an energy function and a power port (a paired effort/flow
boundary variable), and composition couples ports
power-conservingly, so the composite's energy is the sum and
energy-based certificates compose along the interconnection. The
designed elimination has two species --- additive, for flow-summed
sharing, and an exact linear solve with a determinant-guard witness,
for constraint coupling --- and the guard is where the design pays
for itself: the composite kinetic metric is a sum of pullbacks of
positive-definite component inertias, so the division guards that
must today be proven per model become structural consequences of the
pattern. The design's worked example is cartpole recomposed --- cart
and pole as separate Hamiltonian mechanisms coupled through a
power-conserving pivot junction, recovering the standard monolithic
dynamics as a theorem rather than an authoring choice. This design is
in flight; this section records its intended shape.

\syscatmetapat{\emph{Reserved}: no term former is stated ahead of the axis's
vocabulary. The design warrant --- generator additivity, and the
power-conserving junction rows of the port stratum --- is recorded in
the prose above, asserted as nothing until the axis stages.}
\ruleblock{API}
\sourcenote{Producer surface of the \texttt{mtk} toolkit for this composition pattern: the TypeScript arrangement (npm package \texttt{@maitria/mtk}) and the Rust catalogue crate (\texttt{crates/mtk-catalogue}) --- one Rust implementation underneath, reached from TypeScript in-process through a native binding, so the two surfaces cannot diverge. Signatures and constructions are extracted from the toolkit sources at every regeneration of this catalogue; each construction is quoted verbatim from the toolkit's conformance corpus, code whose emitted certificate bytes are verified byte-identical against the committed golden vectors (helper values named \texttt{fx\_*} or \texttt{f.*} are that corpus's shared fixtures, defined beside the quoted lines). This API constructs and emits; it never checks --- the checker re-derives every claim independently.}
\apilang{TypeScript}{---}
As for Rust: reserved for a later revision of the producer surface.
\apilang{Rust}{---}
None by design at this preview: state-sharing / net-of-nets (gluing along shared state --- places, common noise --- rather than signal wires; the generator's additivity over transitions as the designed semantics) is reserved for a later revision of the producer surface, with its design warrant recorded; SDCPN net-of-nets and common-noise mean-field composition are its first customers.

\section{Nets of Nets}
\label{sec:cp-netsofnets}

A closure property rather than a pattern, recorded as its own entry
because the catalogue once reserved ``nets of nets'' as future
vocabulary and most of the reservation dissolved under the term
formers now stated. The population former is parametric in its inner
System: \rulename{S-BAG}'s premise (section~\ref{sec:mt-bag}) is
$S : \mathsf{System}$, with nothing about $S$ consulted beyond its
sort. And the coloured-net former concludes at exactly that sort:
$\mathrm{SDCPN}(P;\, T;\, A) : \mathsf{System}$
(section~\ref{sec:mt-sdcpn}). So a net whose tokens are themselves
nets is one substitution --- put a coloured-net System in the
token-System slot where an ODE mechanism would ordinarily go, and the
outer net's places hold inner nets as tokens: hierarchy by
instantiating a premise, not by any new row family. The base rung of
the same ladder is already in the catalogue as the population
entry's identity anchor --- a Bag of ODE tokens \emph{is} the
single-place, zero-arc coloured net, two presentations of one System
--- and each further rung is the same substitution one level up.

What this entry deliberately does not claim: the certification story
across nesting depth. The well-formedness is free --- that is the
finding --- but no shipped certificate yet exercises a
population-of-nets carrier end to end, and the grading discipline of
chapter~\ref{ch:qtsl-mathematical-models} applies to the nested composite
exactly as to any frame: depth adds obligations, not exemptions.
What remains genuinely reserved is different in kind --- two nets
glued along a \emph{shared place}, which is state sharing
(section~\ref{sec:cp-stateshare}), not token colouring, and stays on
that axis's reservation.

No metanotation or API block: there is no former here to state ---
the entry's formal content is two already-stated rules composing, and
a dedicated block would only restate them.
 \chapter{Specification Language}
\label{ch:qtsl-specification-language}

\addtocontents{toc}{\protect\setcounter{tocdepth}{1}}

\noindent\draftmark\medskip

Every certificate in this stack ultimately concludes a claim of the
form \emph{this system satisfies this specification}. The
Mathematical Models catalogue (chapter~\ref{ch:qtsl-mathematical-models}) is the
normative inventory of the left-hand side of that judgment; this
chapter is its right-hand side: the \emph{spec language} --- the
temporal-logic proposition formers a modeller actually states
requirements in --- together with the record that packages a
requirement for composition, the derived spellings that cover
recurring claim species without new grammar, the rewrite layer that
lets the spec \emph{as stated} be the spec \emph{as certified}, the
deliberate absences, and the judgment kinds: the satisfaction-side
sorts that certificates conclude into.

The presentation is semantics-first, and its form follows the rest of
this book. The grammar is not given as productions: each former is a
\emph{deduction rule} in Computer Science
Metanotation~\cite{steele2017}, concluding a formation judgment ---
premises above the line are the former's well-formedness conditions,
side conditions stand beside the rule, and what used to be checked
into existence is now visible in the rule that brings the term into
existence at all. The discipline pays twice. First, the deliberate
absences of section~\ref{sec:spec-lang-fences} stop being policy
prose: a fence is the \emph{absence of a rule}, and several of the
absences carry a semantic theorem as their warrant
(section~\ref{sec:spec-lang-formation}). Second, each formation rule
is the natural index of its semantics: every former carries a
\emph{compile clause} into the spec storey of the semantic universe
(chapter~\ref{ch:qtsl-semantic-universe}), an \emph{adequacy
obligation} relating the compile to the former's intended reading ---
stated with its current status, machine-checked where it is
machine-checked and named where it is owed --- a \emph{witness} from
the certified corpus, and the producer \emph{API} that constructs it.
The record-level compile is a structural recursion over the formation
derivation, so the adequacy of a whole specification is the fold of
its rules' cells (section~\ref{sec:spec-lang-record}).

\paragraph{Why the grammar exists at all.} Before these formers were
typed rows, two claim species --- expected return and bounded always
--- lived in certificate-header \emph{sidecars}: prose fields beside
the proof rows, spelling claims the grammar could not. A claim that
lives in a sidecar is a claim no checker re-derives and no
composition rule can consume. The discipline this chapter enforces is
the retirement of that disease: \textbf{a spec claim is a typed row},
and the rows' schema has a normative home in this repository ---
\path{theories/Qtsl.geolog}, the theory encoding this chapter's
vocabulary as data, with its well-formedness and mutant battery
beside it. Everything below exists so that the statement a reader
sees is the statement the checker replayed.

The grammar is deliberately small --- threshold atoms over
observables, Boolean structure, probability, expectation, and level
thresholds, run-functional accumulation, and the temporal
operators (\emph{always}, \emph{always up to a deadline},
\emph{reach by a deadline}, \emph{until}, metered Streett
acceptance) --- and
the catalogue's completeness claim is empirical and checkable:
\textbf{every claim species asserted by a certificate in the
certified corpus appears in this chapter}, either as a former, as a
derived spelling, or as a judgment kind --- with the exercising
artefact named. A claim species with no catalogue entry is a defect
of the catalogue, never a licensed side channel: the admission rule
is that a new species gets its former, with a named kernel and a
named semantics obligation, when its first certificate lane arrives.

\section{The Compile Target: Spec Classes of the Semantic Universe}
\label{sec:spec-lang-target}

What a specification denotes is fixed once, in the semantic universe,
and this chapter compiles into it rather than re-deriving it. The
stock this chapter consumes
(chapter~\ref{ch:qtsl-semantic-universe}, section~\ref{sec:su-classes}):
a \emph{spec atom} is a safety face together with a finite set of
liveness demands; its class in the hyperspace is a \emph{crescent} ---
one closed containment face cut by finitely many open escape faces, a
locally closed set --- and a compound specification in disjunctive
normal form denotes the union of its disjuncts' crescents, a
\emph{constructible} class. Satisfaction is the fibre order itself:
a system's belief (level-of-inconsistency functional;
chapter~\ref{ch:qtsl-semantic-universe}) satisfies a spec when the
spec's belief lies
below it (\texttt{SuFunctor.Universe.Sat}), and the crisp judgment
against a class has two faces --- \emph{safety} is wall domination,
the system's permitted region contained in the demanded one, and
\emph{liveness} is a dip, some admitted state permitted at a finite
level of inconsistency under each demand (\texttt{wallOn\_le\_iff},
\texttt{dip\_top\_iff}). Where the model leaves openness --- demonic
inputs, worst-case initial conditions, imprecise kernels --- the
judgment reads against a \emph{set} of behavioural laws through the
credal envelope faces (\texttt{SatCredal},
in the adequacy layer), lower for guarantees; a singleton set
recovers the plain reading, a purely demonic model recovers
min-over-branches. This one convention removes any need for
separate worst-case formers.

An earlier edition of this chapter fixed four denotation strata ---
region, path, credal, $\omega$-word --- as primitives of its front
matter. They survive, but as \emph{views} of the landed semantics,
not as axioms of the catalogue:

\begin{itemize}
\item the \emph{region} view: a state formula's denotation is a set
  of states, entering the judgment as threshold walls --- and nothing
  expressible is lost by compiling regions to walls, because rational
  threshold rows generate the flat universe
  (\texttt{iSup\_boxRows\_eq\_self});
\item the \emph{path} view: a path formula's denotation is a class of
  records over the site --- occupied-stage classes over the record
  lens at the finitary stratum, trajectory classes at the clocked
  strata (section~\ref{sec:spec-lang-pathf});
\item the \emph{credal} view: probability and expectation atoms read
  \emph{lower} bounds against the whole set of admitted laws, through
  the envelope faces named above;
\item the \emph{$\omega$-word} view: the rewrite layer
  (section~\ref{sec:spec-lang-rewrite}) reads formulas over infinite
  words --- one carrier instance among those a former admits, the one
  on which the eighteen rewrite rules are machine-checked today.
\end{itemize}

\paragraph{The properness presupposition.} One structural fact of the
compile target shapes the entire chapter, so it is stated once, here.
Every spec atom's liveness set contains the top demand --- in the
mechanization this is a \emph{field} of the atom
(\texttt{SpecAtom.top\_mem}), so an atom without it is not merely
refused but unrepresentable --- and therefore \textbf{every
specification's interpretation includes the trivial liveness bound}:
eventually \emph{something}, at a finite level of inconsistency.
Satisfaction of any
spec asserts \emph{safe and proper}: to satisfy anything, first
exist. This is a law of the semantics, not a default a producer
opts into --- the compile emits the mark unconditionally, and that
emission is itself a theorem
(\texttt{elabSpec\_top\_mem}) --- and it costs nothing in syntax:
there is no conjunct to write and no flag to store, because the
demand lives in the judgment.

What it buys is that infeasibility is observable. An inconsistent
system --- one admitting no behaviour at all --- satisfies every
safety demand vacuously and no liveness demand
(\texttt{top\_notMem\_modelClass}: the inconsistent system belongs to
no spec class); a pure-safety reading would let a bricked system
satisfy the specification of a working one, and composition could
manufacture such systems from proper components while every
certificate stayed green. Under the presupposition the true constant
is not vacuous: $\mathsf{Spec}(\mathsf{SInit}(\mathsf{STrue}))$ ---
the initial-state injection wrapping the true constant
(section~\ref{sec:spec-lang-record}) --- is the \emph{bare
feasibility spec}, the minimal claim distinguishing a bricked system
from a working one (section~\ref{sec:spec-lang-const}). The same
guard appears at the other end of the judgment, where the Rule
Catalogue's liveness verdicts carry their existential conjunct
(chapter~\ref{ch:qtsl-manual}) --- one guard, two ends. What it costs
is one nonemptiness discharge per root satisfaction claim, discharged
where the claim is consumed; and where infeasibility is genuinely
expected in some conditions, the honest spelling is the assumption
slot of the spec record --- satisfaction relativized to conforming
inputs, with nothing demanded off-assumption
(section~\ref{sec:spec-lang-vdashi}) --- never a properness opt-out,
which the grammar deliberately does not provide.

\section{Formation Judgments}
\label{sec:spec-lang-formation}

Eight sorts carry the grammar, and they are not a flat list: four of
them arrange as a two-by-two grid --- the \emph{temporal} axis
(state: atemporal, against path: trajectory) crossed with the
\emph{shape} axis (observable: $\mathbb{Q}$-valued, against pred:
Boolean) --- and the grid is capped by the polarity sorts --- two
faces and the attitude blend --- and the record.

\begin{center}
\begin{tabular}{l|cc}
 & \textbf{Observable} ($\mathbb{Q}$-valued) & \textbf{Pred} (Boolean)\\
\hline
\textbf{State} (atemporal) & $o\ \mathsf{StateObservable}$ & $\varphi\ \mathsf{StatePred}$\\
\textbf{Path} (trajectory) & $R\ \mathsf{PathObservable}$ & $\psi\ \mathsf{PathPred}$\\
\end{tabular}
\end{center}

\begin{center}
\begin{tabular}{ll}
$o\ \mathsf{StateObservable}$ & $o$ is a well-formed atemporal observable\\
$\varphi\ \mathsf{StatePred}$ & $\varphi$ is a well-formed atemporal predicate (a region)\\
$R\ \mathsf{PathObservable}$ & $R$ is a well-formed run-functional\\
$\psi\ \mathsf{PathPred}$ & $\psi$ is a well-formed trajectory event\\
$\Phi\ \mathsf{ForallSpec}$ & $\Phi$ is a well-formed safety-face component\\
$E\ \mathsf{ExistsSpec}$ & $E$ is a well-formed liveness-mark component\\
$H\ \mathsf{AttitudeSpec}$ & $H$ is a well-formed attitude blend (a design score)\\
$s\ \mathsf{Spec}$ & $s$ is a well-formed spec record\\
\end{tabular}
\end{center}

Three of the eight are this edition's names for sorts the catalogue
always had: $\mathsf{StatePred}$ was \textsf{StateF},
$\mathsf{PathPred}$ was \textsf{PathF}, and $\mathsf{PathObservable}$
was \textsf{Return} --- repairing a standing defect of the earlier
edition, whose formation table omitted \textsf{Return} entirely while
eight rules concluded into it. $\mathsf{StateObservable}$ names the
atemporal quantitative stratum that was previously anonymous: a
polynomial argument, a channel index, a gate citation --- material
the rules always consumed without a judgment to its name; its four
injections are now formers with an entry of their own
(section~\ref{sec:spec-lang-obs}). The two polarity caps are new,
and they are the point of the redesign; the attitude cap completes
them, entered with the theory revision because the blend it carries
fits \emph{neither} face --- its discipline is stated at its entry
(section~\ref{sec:spec-lang-hurwicz}). One further sort stands
beside the eight as pair vocabulary rather than as a grid cell:
$\pi\ \mathsf{StreettPair}$ ($\pi$ is a well-formed request/response
pair), feeding the two acceptance formers. Its dissolution into
per-pair premise structure was designed, and the theory revision
settled the question the other way, deliberately: the pair spine
persists at the row stratum as list-spine \emph{encoding} --- small
hash-consed nodes, pairs-as-data --- with the two pair constructors
carrying the consumption-parity vocabulary, and no rule reads the
pair sort as a judgment.

\paragraph{The stratification law.} The State row admits no
trajectory quantification and no credal quantification; the Path row
lifts State-row material along trajectories; \textbf{all
quantification over the credal body happens at the cap, and the cap
is where polarity lives}. A $\mathsf{ForallSpec}$ component reads its
bound against \emph{every} admitted law --- the safety face,
compiling to a floor: wall domination, the closed containment face of
the compile target (section~\ref{sec:spec-lang-target}) --- while an
$\mathsf{ExistsSpec}$ component asserts that \emph{some} admitted law
attains --- a liveness mark, compiling to a dip: the open escape
face. An $\mathsf{AttitudeSpec}$ component reads \emph{both}
envelopes and mixes them arithmetically before thresholding --- no
face, no participation in the floor/mark algebra, a design score
rather than a guarantee; in one line, \textbf{a Spec mixes verdicts;
an attitude mixes envelopes} (section~\ref{sec:spec-lang-hurwicz}).
Polarity is thereby a \emph{sort}, not an annotation: a
specification cannot be written without declaring the face of every
component, a mis-faced premise is a formation error rather than a
semantic surprise, and a rule that needs a face restriction states it
as a sort premise instead of a prose side condition. The judgment
layer is untouched --- one turnstile, no subscripts: the face of
every component is read off the constructor vocabulary of the
\emph{arguments}.

\paragraph{The formers are the grid's structure maps.} Every former
of this chapter lands on exactly one cell or arrow of the diagram:

\begin{itemize}
\item \emph{within-cell structure} (State--Pred): the constants and
  Boolean closure --- \textsc{STrue}, \textsc{SFalse}, \textsc{SAnd},
  \textsc{SOr}, \textsc{SNot};
\item \emph{within-cell structure} (State--Observable): the
  provenance injections \textsc{oChan}, \textsc{oGate},
  \textsc{oPoly} (section~\ref{sec:spec-lang-obs}) --- one
  observable, whatever its provenance;
\item \emph{thresholding}, Observable $\to$ Pred at the State row:
  \textsc{SAtomLE}, \textsc{SAtomGE};
\item \emph{indicator embedding}, Pred $\to$ Observable at the State
  row: \textsc{oInd}, the evident square-filler
  ($\mathbf{1}_\varphi$ as an observable);
\item \emph{temporal lifting}, Pred column, State $\to$ Path:
  \textsc{PAlways}, \textsc{PAlwaysT}, \textsc{PDeadline},
  \textsc{PUntil}, the metered acceptance pair
  \textsc{PDwellStreett} and \textsc{PCrossStreett},
  \textsc{PNotPath};
\item \emph{accumulation}, Observable column, State $\to$ Path: the
  accumulators \textsc{poSum}, \textsc{poInt}, \textsc{poDisc},
  \textsc{poAvg}, \textsc{poSup}, each over one observable --- the
  earlier edition's provenance-fused run-functional constructors,
  re-founded as injections plus accumulation modes at unchanged
  denotations --- beside the transition-count functional
  \textsc{retFireCount} (section~\ref{sec:spec-lang-returnstratum});
\item \emph{indicator embedding}, Pred $\to$ Observable at the Path
  row: \textsc{retPath};
\item \emph{thresholding}, Observable $\to$ Pred at the Path row:
  \textsc{PAtomLE}, \textsc{PAtomGE} --- pathwise accumulated-cost
  \emph{events}, a claim species distinct from any expectation bound
  (section~\ref{sec:spec-lang-patom});
\item \emph{the cap}: \textsc{SProbGE}, \textsc{SProbLE},
  \textsc{SReturnGE}, \textsc{SReturnLE}, \textsc{SLevelGE},
  \textsc{SLevelLE}, the event-rate envelope pair \textsc{SRateGE},
  \textsc{SRateLE} (over transition rosters, horizon-parametric ---
  section~\ref{sec:spec-lang-rate}), the initial-state injection
  \textsc{SInit}, and
  the regret certificate face \textsc{SRegretLE} conclude
  $\mathsf{ForallSpec}$; the possibility marks
  \textsc{SProbMayBeGE}, \textsc{SProbMayBeLE},
  \textsc{SReturnMayBeGE}, \textsc{SReturnMayBeLE},
  \textsc{SLevelMayBeGE}, \textsc{SLevelMayBeLE}, and the regret
  query face \textsc{SRegretLEMark} conclude $\mathsf{ExistsSpec}$;
  the attitude pair \textsc{SHurwiczGE}, \textsc{SHurwiczLE}
  concludes $\mathsf{AttitudeSpec}$. The mark roster is
  \emph{polarity-total}: each threshold family --- probability,
  expectation, level --- presents both faces (the level marks are
  new with this edition, completing the grid), and the regret pair
  is the same duality one seat up
  (section~\ref{sec:spec-lang-regret});
\item \emph{the record}: $\mathsf{Spec}$ packages face components
  (section~\ref{sec:spec-lang-record}); and spec-level disjunction
  $(s_1 \vee s_2)\ \mathsf{Spec}$ for $s_1, s_2\ \mathsf{Spec}$ ---
  its denotation is the union of the two record classes,
  constructible by definition, and it never asks the impossible
  merge (the no-single-floor and no-single-mark facts below are
  about \emph{one} record's slots); discharge is by disjunct
  selection.
\end{itemize}

\paragraph{The one deliberate absence.} The completed diagram has
exactly one missing arrow family, and it is absent on purpose: the
\emph{descending} arrows,
$\mathsf{Path}\{\mathsf{Pred},\mathsf{Observable}\} \to
\mathsf{State}\{\mathsf{Pred},\mathsf{Observable}\}$ --- nested path
quantification in the CTL style, where a path-quantified claim
re-enters the atemporal row as a state label
($\mathbb{P}_{\ge q}[\psi]$ as a state predicate; a value functional
as a state observable). Its status is recorded in three
facts. First, the \emph{raw} formers are not cautiously deferred but
fenced by the geometry: a single level of such re-entry already
exits the constructible classes --- neither face's rooted class is a
crescent --- and the negation fence forbids recovering the universal
face from the existential one, so any admission would have to carry
both faces as primitives (the analysis is mechanized in the adequacy
layer's path-to-state development). Second, the pathway is open, not
closed: the behavioural tier already \emph{denotes} the descending
arrows --- the probability-threshold and expected-return sets of the
compile clauses below are exactly path-to-state denotations --- and
the principled admission route is a \emph{certified-citation}
former: the tower discipline promoted from convention to formation
rule, Observable column first, where certified envelopes carry the
right one-sided semicontinuity. Third, that route is deliberately
out of scope for this edition, planned for the language's second
major revision (2026Q4) --- a choice made with the open pathway in
full view, which is what distinguishes a scoped fence from an
impossibility. What the
certified corpus actually needs of nesting is already served through
the judgment stratum: a child claim's \emph{certified region} enters
a parent formula as an ordinary $\mathsf{StatePred}$ atom
(entry \textsc{SProbGE}; the reach-then-hold spelling of
section~\ref{sec:spec-lang-dfunnel} is the landed witness).

\paragraph{Refinements.} Two further judgments are read off
derivations rather than declared:
$\varphi\ \Theta_0$ --- membership in the \emph{base stratum}
(nonstrict polynomial atoms, constants, and their and/or closure;
nothing else), demanded wherever a consumer needs base-exact
material --- and the \emph{carrier-requirement synthesis}
$\psi \rhd \kappa$ with
$\kappa \subseteq \{\mathsf{steps}, \mathsf{duration}\}$, which
computes what temporal structure a former demands of the carrier it
is interpreted on (per-tick accumulation demands steps; deadlines,
discounts, and time-integral accumulation demand duration). Under
the grid the synthesis is a \emph{Path-row device}, and the
stratification law gains a checkable spelling: every State-row
former synthesizes $\rhd \varnothing$ --- atemporality as a
computation, not a convention. Carrier
scope is \emph{derived, never stored} --- the consequence for
equivalence judgments is drawn at entry \textsc{QSpecEquiv}
(section~\ref{sec:spec-lang-qspecequiv}).

One further stratum is deliberately \emph{not} a formation judgment:
consumability at parity. Negation is a surface former
(\textsc{SNot}, \textsc{PNotPath}), resolved by negation normal
form, and whether a residue survives to a consumable position is
\emph{checked} --- by the producer-side typing preflight, and
independently by the checker --- not stored in the term. In rule
form the refusal table becomes a family of absences: there is no
rule concluding consumability of a negated path atom at even
parity, of negated Streett acceptance, of negated bounded until,
or of negated unbounded always. (The table's historical first
member, the negated \emph{probability} atom, no longer needs a row
at all: no probability former concludes $\mathsf{StatePred}$ any
more, so $\mathsf{SNot}$ cannot reach one --- the fence became a
type; the path atom of the theory revision inherits its seat in the
table.) These absences carry a semantic
warrant, stated once and cited by the fence entries: \textbf{under
the properness presupposition, no spec class has an expressible
complement}. The proof is two lines --- every class excludes the
inconsistent system (\texttt{top\_notMem\_modelClass}), so every
complement \emph{contains} it, and no union of crescents does ---
so the class algebra is a distributive lattice, irreducibly
non-Boolean, and the refused residues are not missing engineering
but missing mathematics. The dual machinery that would consume them
(probability upper bounds as positive claims, Rabin acceptance,
over-approximating towers) is deliberately absent until a
certificate lane demands it
(section~\ref{sec:spec-lang-fences}).

The surface grammar, its normal forms, and its compile are
themselves mechanized, and the mechanization consumes the formation
discipline literally: state predicates and path predicates are
inductive types whose constructors are the rules below
(\texttt{SuFunctor.Universe.SpecElab} and the path-lens development
beside it) --- and the mechanized state-predicate grammar was
already atemporal and atom-abstract before the grid named the
discipline, so $\mathsf{StatePred}$ is a rename at that stratum, not
a rebuild; negation normal
form has \emph{no negation node by type}, so unconsumable residues
have no clause because they have no constructor; and the compile is
one clause per formation rule, with the record-level adequacy
statement the fold of the per-rule cells over the formation
derivation --- the same architecture as the Rule Catalogue's
soundness, one lemma per rule composed by induction over admitted
derivations, transposed from the inference stratum to the formation
stratum.

\paragraph{Sequencing.} The grid was enacted grammar-first, and this
section is its normative statement; the normative theory's row
schema now carries it whole (\path{theories/Qtsl.geolog}): the
re-sorted conclusions, the level marks, the initial-state injection,
spec-level disjunction, the embedding and threshold arrows, the
metered acceptance pair, the running-sup and transition-count
functionals, the attitude cap, and the component-record assembly are
all rows of the landed revision, with their well-formedness and
mutant battery beside them. Denotations are unchanged throughout,
and every inhabitant of the earlier edition's grammar has a total
spelling transport --- the re-homed cap formers keep their names and
field labels with only sort membership moved; a fused run-functional
becomes its accumulator over the matching injection; a polynomial
atom becomes the same atom over \textsc{oPoly}; a bare root region
becomes its explicit \textsc{SInit} component --- so no certified
artefact and no published number moved with the re-sort. The
producer surface has since followed: the toolkit's revision landed
the grid spellings, and the quoted API blocks below are extracted
from that revised surface --- the residual producer gaps (the
record-top disjunction, the compiled verdict form, the regret
certificate face) are named at the entries that own them.

\paragraph{How to read an entry.} Each entry carries two blocks. The
\emph{Formation} block states the former's deduction rules, with
$\Theta_0$ rules inline where the former preserves or demands base
membership and $\rhd$ annotations beside the rules that demand
carrier structure. The prose beneath it states the \emph{compile
clause} --- what the former elaborates to in the semantic universe's
vocabulary --- the \emph{adequacy} mode and status (exact where the
compile's class provably equals the intended reading; directed where
a sound implication is proved and no converse is claimed; per
consuming battery where the obligation is discharged rule-side), and
the \emph{witness}: the landed artefact that exercises the former.

The \emph{API} block is the producer surface of the \texttt{mtk}
toolkit for that entry: the TypeScript arrangement (npm package
\texttt{@maitria/mtk}) and the Rust catalogue crate
(\texttt{crates/mtk-catalogue}) --- one Rust implementation
underneath, reached from TypeScript in-process through a native
binding, so the two surfaces cannot diverge. Signatures and
constructions are quoted \emph{verbatim} from the toolkit sources and
its conformance corpus: code whose emitted certificate bytes are
verified byte-identical against committed golden vectors, from both
languages (the TypeScript-driven path re-emits the same bytes the
Rust-driven path committed). The quoted blocks are regenerated
mechanically from the sources by the catalogue's extraction tooling
--- a missing declaration or marker fails the regeneration loudly ---
so quoted code cannot drift silently. This API constructs and emits;
it never checks: constructors refuse malformed \emph{syntax} with the
failure's geometry attached (which field, which subterm, which
value), the typing preflight is a courtesy image of the checker-side
oracle, and the checker's independent replay is the only authority.

Two vocabulary facts used throughout, stated once. \emph{Exact
scalars}: every threshold, deadline, probability, and coefficient is
an exact rational (\texttt{Q} in TypeScript, an exact-rational
witness type in Rust); no floating point enters a claim.
\emph{Positional variables}: an atom's polynomial reads the paired
system's state coordinates positionally --- or, for interface-indexed
specs, the interface's output wires in list order (guarantees) and
input wires (assumptions); port names are presentation data.

\section{The Spec Record}
\label{sec:spec-lang-record}

\subsection{Spec: root, interface, assumptions, slack}
\label{sec:spec-lang-spec}

\ruleblock{Formation}
\begin{mathpar}
\inferrule*[Right=\textsc{SInit}]
  {\varphi\ \mathsf{StatePred}}
  {\mathsf{SInit}(\varphi)\ \mathsf{ForallSpec}}
\\
\inferrule*[Right=\textsc{Spec-C}]
  {\Phi_i\ \mathsf{ForallSpec} \quad (i < m) \\
   E_j\ \mathsf{ExistsSpec} \quad (j < k) \\
   H_l\ \mathsf{AttitudeSpec} \quad (l < r) \\\\
   \mathrm{widths}(\vec\Phi; \vec E; \vec H) \text{ agree pairwise}}
  {\mathsf{Spec}(\vec\Phi; \vec E; \vec H)\ \mathsf{Spec}}

\inferrule*[Right=\textsc{Spec-I}]
  {\Phi_i\ \mathsf{ForallSpec} \quad (i < m) \\
   E_j\ \mathsf{ExistsSpec} \quad (j < k) \\
   H_l\ \mathsf{AttitudeSpec} \quad (l < r) \\
   \mathrm{width}(\vec\Phi; \vec E; \vec H) = \mathrm{arity}_{\mathrm{out}}(I) \\\\
   \alpha\ \mathsf{StatePred} \\ \alpha\ \Theta_0 \\
   \mathrm{width}(\alpha) = \mathrm{arity}_{\mathrm{in}}(I)}
  {\mathsf{Spec}(\vec\Phi; \vec E; \vec H, I, \alpha)\ \mathsf{Spec}}
\\
\inferrule*[Right=\textsc{Spec-Or}]
  {s_1\ \mathsf{Spec} \\ s_2\ \mathsf{Spec}}
  {(s_1 \vee s_2)\ \mathsf{Spec}}
\end{mathpar}

A specification is a package of \emph{face components} --- alone
(\textsc{Spec-C}, the closed form) or against an interface with an
optional $\Theta_0$ assumption on input wires (\textsc{Spec-I}); an
optional slack slot (a candidate gain function for quantitative
composition) may accompany either form. The rule shape wears the compile
target's own algebra: the $\mathsf{ForallSpec}$ components conjoin
invertibly --- floors merge by pointwise maximum, one slot --- while
the $\mathsf{ExistsSpec}$ components aggregate as a \emph{list},
because marks do not merge (no single mark expresses two demands;
the record holds a list exactly where the semantics refuses a
merge); the $\mathsf{AttitudeSpec}$ components stand beside the
faced ones and render as design scores, never as guarantees
(section~\ref{sec:spec-lang-hurwicz}), taking no part in either the
floor merge or the mark list. A record must carry at least one
component --- across the three families, or a disjunction --- as an
admission conformance fact: the empty record claims nothing, and
the bare feasibility spec is its minimal inhabitant.
The initial-state injection \textsc{SInit} is the honest
name for what the earlier edition attached silently to a bare region
at spec root: $\mathsf{SInit}(\varphi)$ reads \emph{every admitted
initial state satisfies $\varphi$} --- the $\forall$-face reading,
now one visible rule --- so the bare feasibility spec is spelled
$\mathsf{Spec}(\mathsf{SInit}(\mathsf{STrue}))$, its properness mark
carried by the judgment exactly as before
(\texttt{elabSpec\_top\_mem} unchanged). Spec-level disjunction
(\textsc{Spec-Or}) sits at the record top
(section~\ref{sec:spec-lang-formation}): a union of record classes,
discharged by disjunct selection --- deliberately \emph{above} the
record, so it never asks two records to merge into one floor slot
and mark list; a disjunction record carries \emph{no} component rows
of its own (the exclusivity is refuse-mode at the row stratum: its
class is its arms'), and or-graph acyclicity is admission
conformance. Three record-level disciplines are rule shapes
rather than checks. There is no rule introducing an assumption without an
interface --- ``a spec carrying an assumption must carry an
interface'' is the absence of an $\alpha$-with-no-$I$ rule, and in
the mechanization the assumption slot lives \emph{inside} the
interface option, so the discipline holds by type. Assumptions are
policed to the base stratum at formation (temporal assumptions would
reintroduce assume--guarantee circularity, so they are fenced at the
grammar). And closed versus open is an absence distinction: a spec
with no interface row is closed, and the closed satisfaction
judgment is the empty-interface instance of the one satisfaction
kind (entry \textsc{QSat}) --- there is no second sort for
interface-indexed satisfaction, deliberately. The slack slot's
$\mathcal{K}_\infty$-class membership obligations (zero at zero,
monotone, unbounded) are discharged where the slack is
\emph{consumed} --- by the composition rules of
chapter~\ref{ch:qtsl-manual} --- never stored as facts. There is
also deliberately \emph{no model binding} on a spec: pairing a spec
with its system is the satisfaction judgment's job, and a spec that
named its model would be the sidecar disease in miniature.

At the row stratum the record reads the same way, at the row
schema's grain. A specification is an \emph{entity}; its face
components attach as relation rows, one family per cap sort ---
\texttt{spec\_forall}, \texttt{spec\_exists}, \texttt{spec\_score}
--- zero or more rows each, read per the algebra above, with
component presence, pairwise width agreement (arity arithmetic), and
the disjunction row's exclusivity all admission conformance rather
than rows. The earlier edition's single root \emph{attribute} is
retired with the re-sort: a bare-region root was the \textsc{SInit}
conflation, and a mixed-face conjunction root is what the
component families now spell explicitly. The interface and the slack
(a declared candidate gain function --- a polynomial, in the row
schema) attach as zero-or-one relation rows; assumptions attach as
\emph{zero-or-more} relation rows, read conjunctively --- several
assumption rows constrain jointly, and each row is policed to the
base stratum by a load-time refusal: a provably nonbase assumption
is a malformed record, not a false one. The
assumption-implies-interface discipline itself is carried by the
formation calculus and the mechanization's option type; the row
schema leaves it to them.

\emph{Compile.} The record-level compile is
\texttt{elabSpec}: \emph{interface resolution, then structural
recursion}. Resolution fixes the record's positional geometry ---
the interface's arities for an open spec, pairwise width agreement
for a closed one --- and is total on conformant records
(\texttt{resolveArity\_isOk\_iff}), refusing otherwise \emph{before
any recursion} with the failure's geometry attached: the offending
atom, its found width, the expected arity
(\texttt{resolveArity\_error\_names\_offender}). Decoded records
re-validate with the same checks --- decoded data is data. The
recursion is then one clause per formation rule, and the compile
factors through resolution (\texttt{elabSpec\_isOk\_iff}); its
class-level adequacy is the fold of the per-former cells over the
formation derivation (\texttt{elabSpec\_ok\_modelClass}), with the
properness mark carried unconditionally
(\texttt{elabSpec\_top\_mem}). One deferral, fail-closed and named:
the compile reads the \emph{guarantee} only. An assumption is
carried as data --- its conditional-satisfaction semantics (entry
\textsc{QSat}, the interface-indexed reading) is deferred, never
silently read as an unconditional strengthening, and any consumer
that would need it refuses until the relativized compile lands.

\emph{Witness}: every certificate in the corpus pairs a system with a
record of this shape; the interface-indexed form is exercised by the
liveness vector quoted below (two output wires, one input wire, a
$\Theta_0$ assumption).

\ruleblock{API}
The record-top disjunction (\textsc{Spec-Or}) remains a logged
deferred former on the producer side; everything else below is the
revised all-relations surface.
\apilang{TypeScript}{spec.build / spec.encode / spec.init}
\apiprov{signature --- \texttt{ts/src/spec.ts}}
\begin{tssrc}
export interface Spec {
  readonly name?: string;
  readonly forall?: ReadonlyArray<ForallSpec>;
  readonly exists?: ReadonlyArray<ExistsSpec>;
  readonly score?: ReadonlyArray<AttitudeSpec>;
  readonly iface?: Iface;
  readonly assume?: StatePred;
  readonly slack?: Slack;
}

export function build(record: Spec): Result<Spec, SpecRefusal>

export function encode(record: Spec): Result<Artifact, SpecRefusal>

export function init(arg: StatePred): ForallSpec
\end{tssrc}
\apiprov{construction --- \texttt{ts/tests/spec-native.test.ts}}
\begin{tssrc}
const art = s.unwrap(
  s.encode({
    name: "arbiter-liveness",
    forall: [recur0, recur1],
    iface: {
      invars: [{ name: "u", shape: [] }],
      outvars: [
        { name: "grant0", shape: [] },
        { name: "grant1", shape: [] },
      ],
    },
    assume: s.polyLe(1n, s.unwrap(s.polyVar(0, 1))),
  }),
);
\end{tssrc}
\apiprov{The component families (\textsc{SInit} + the forall / exists / score lists) are driven end to end by the formers-tour vector (the same test file).}
\apilang{Rust}{mtk\_catalogue::spec::Spec::builder}
\apiprov{signature --- \texttt{crates/mtk-catalogue/src/spec/formula.rs + crates/mtk-catalogue/src/spec/record.rs}}
\begin{rustsrc}
pub fn builder() -> SpecBuilder

pub fn forall(mut self, f: ForallSpec) -> Self

pub fn exists(mut self, e: ExistsSpec) -> Self

pub fn score(mut self, h: AttitudeSpec) -> Self

pub fn build(self) -> Result<Spec, Refusal>

pub fn init(arg: StatePred) -> ForallSpec
\end{rustsrc}
\apiprov{construction --- \texttt{crates/mtk-catalogue/tests/golden\_spec.rs}}
\begin{rustsrc}
let spec = Spec::builder()
    .forall(recur0)
    .forall(recur1)
    .name("arbiter-liveness")
    .iface(iface)
    .assume(assume)
    .build()
    .unwrap();
\end{rustsrc}
\apiprov{The component families (\textsc{SInit} + the forall / exists / score lists) are driven end to end by the formers-tour vector (the same test file).}

\section{State-Row Formers}
\label{sec:spec-lang-statef}

This wing is the grid's State row, whole
(section~\ref{sec:spec-lang-formation}): both cells and both arrows
between them. State observables ($\mathsf{StateObservable}$) carry
the row's quantitative material --- one observable, whatever its
provenance --- and state predicates ($\mathsf{StatePred}$) denote
\emph{regions} --- sets of states from which the claim holds ---
compiling to spec atoms whose safety face is the region's wall, with
the properness mark on every atom; thresholding (\textsc{SAtom})
passes from observables to predicates, and indicator embedding
(\textsc{oInd}) passes back. Everything here is atemporal by the
stratification law, checkably: every former of this wing synthesizes
$\rhd \varnothing$. Trajectories enter with the next wing, and the
credal body only at the spec faces
(section~\ref{sec:spec-lang-faces}). Two
moves organize everything. The tower discipline of
chapter~\ref{ch:qtsl-manual} is \emph{certificate-mediated}: a
certified child region enters a parent formula as an ordinary
$\mathsf{StatePred}$ atom through the judgment stratum, never by
syntactic embedding (entry \textsc{SProbGE}). And the grammar is
closed, nonstrict,
and canonical by construction, so the checker never normalizes ---
producers do.

\subsection{STrue and SFalse: the constants}
\label{sec:spec-lang-const}

\ruleblock{Formation}
\begin{mathpar}
\inferrule*[Right=\textsc{STrue}]
  {\;}
  {\mathsf{STrue}\ \mathsf{StatePred}}

\inferrule*[Right=\textsc{SFalse}]
  {\;}
  {\mathsf{SFalse}\ \mathsf{StatePred}}

\inferrule*[Right=$\Theta_0$]
  {\;}
  {\mathsf{STrue},\, \mathsf{SFalse}\ \Theta_0}
\end{mathpar}

\emph{Compile.} $\mathsf{STrue}$ elaborates to the one-atom
presentation $(0, \{\top\})$ --- the free wall carrying the
properness mark alone --- and $\mathsf{SFalse}$ to the empty
disjunct family. Exact at the class level: the true constant's class
is the proper systems (\texttt{SpecDnf.modelClass\_top}), the false
constant's class is empty (\texttt{SpecDnf.modelClass\_bot}). Under
the properness presupposition the true constant is therefore not
vacuous: $\mathsf{Spec}(\mathsf{SInit}(\mathsf{STrue}))$ is the bare
feasibility
spec --- the minimal claim distinguishing a bricked system from a
working one --- and certifying it is certifying existence.

\emph{Witness}: exercised everywhere; $\mathsf{STrue}$ is the
trivially-true base-exact Streett request under the recurrence sugar
(entry \textsc{PDwellStreett} and the derived spellings below).

\ruleblock{API}
\apilang{TypeScript}{spec.TRUE / spec.FALSE}
\apiprov{signature --- \texttt{ts/src/spec.ts}}
\begin{tssrc}
export const TRUE: StatePred = Object.freeze({ kind: "true" as const });
/** `SFalse` — the empty set. */
export const FALSE: StatePred = Object.freeze({ kind: "false" as const });
\end{tssrc}
\apiprov{construction --- \texttt{ts/tests/spec-native.test.ts}}
\begin{tssrc}
const orPart = s.or([s.FALSE, s.not(s.polyLe(0n, x1))]);
\end{tssrc}
\apiprov{\texttt{TRUE} rides the formers-tour root (the grammar-tour vector of the same test file).}
\apilang{Rust}{mtk\_catalogue::spec::StatePred::\{True, False\}}
\apiprov{signature --- \texttt{crates/mtk-catalogue/src/spec/formula.rs}}
\begin{rustsrc}
/// `STrue` — the full space.
True,
/// `SFalse` — the empty set.
False,
\end{rustsrc}
\apiprov{construction --- \texttt{crates/mtk-catalogue/tests/golden\_spec.rs}}
\begin{rustsrc}
let or_part = StatePred::or(vec![
    StatePred::False,
    StatePred::not(StatePred::poly_le(q(0, 1), x1())),
]);
\end{rustsrc}

\subsection{oChan, oGate, oPoly, oInd: the observable injections}
\label{sec:spec-lang-obs}

\ruleblock{Formation}
\begin{mathpar}
\inferrule*[Right=\textsc{oChan}]
  {k \in \mathbb{N}}
  {\mathsf{oChan}(k)\ \mathsf{StateObservable}}
\and
\inferrule*[Right=\textsc{oGate}]
  {g\ \mathsf{HypernetRef}}
  {\mathsf{oGate}(g)\ \mathsf{StateObservable}}
\\
\inferrule*[Right=\textsc{oPoly}]
  {p\ \mathsf{Poly} \\ p \text{ canonical}}
  {\mathsf{oPoly}(p)\ \mathsf{StateObservable}}
\and
\inferrule*[Right=\textsc{oInd}]
  {\varphi\ \mathsf{StatePred}}
  {\mathsf{oInd}(\varphi)\ \mathsf{StateObservable}}
\end{mathpar}

The State row's quantitative cell, and the stratum where reward
\emph{provenance} lives --- so that everything downstream consumes
\emph{one observable, whatever its provenance}. A \emph{channel}
citation names a positional output wire of the paired subject's
interface; channel-within-arity is record-level conformance at the
pairing, where the interface is in hand, and the citation elaborates
to the readout's gate there --- computed, never trusted. A
\emph{gate} cites a scalar-codomain hypernet on the (possibly
augmented) state carrier by content address --- the carve
vocabulary's gate object itself --- with scalar-codomain-ness
checked at admission, where the cited content is reachable; the
demand is minted once, at the injection, and every gate-bearing
consumer inherits it through this single locus. A \emph{polynomial}
is the view over the paired carrier's positional variables, carrying
the canonicality premise (coefficients nonzero and exactly rational,
exponent tuples of uniform width, strictly ascending in the
canonical order --- established by producer-side canonicalization
rather than checked into existence). And \textsc{oInd} embeds a
predicate as its characteristic observable
($\mathbf{1}_\varphi$) --- the State row's Pred $\to$ Observable
arrow, one of the grid's two indicator arrows (\textsc{retPath} is
the other, one per row): probability is the expectation of an
indicator \emph{by the diagram}, and the alias ties stay denotation
lemmas at the mechanized stratum, never schema equations. Only the
polynomial injection reaches the base stratum: a channel, gate, or
indicator atom is nonbase by the refinement's own negative rows,
which is what keeps $\Theta_0$ base-\emph{exact} rather than
base-shaped.

\emph{Witness}: every polynomial atom in the corpus enters through
\textsc{oPoly} under the re-sort's spelling transport; the channel
and gate injections carry the provenance the run-functional formers
of section~\ref{sec:spec-lang-returnstratum} previously fused.

\ruleblock{API}
\apilang{TypeScript}{spec.oChan / oGate / oPoly / oInd}
\apiprov{signature --- \texttt{ts/src/spec.ts}}
\begin{tssrc}
export function oChan(chan: number): StateObservable

export function oGate(digest: string, role?: string): StateObservable

export function oPoly(p: Poly): StateObservable

export function oInd(arg: StatePred): StateObservable
\end{tssrc}
\apiprov{construction --- \texttt{ts/tests/spec-native.test.ts}}
\begin{tssrc}
const indAtom = s.atomGe(1n, s.oInd(band));
const chanObs = s.oChan(1);
\end{tssrc}
\apilang{Rust}{mtk\_catalogue::spec::StateObservable::\{chan, gate, poly, ind\}}
\apiprov{signature --- \texttt{crates/mtk-catalogue/src/spec/formula.rs}}
\begin{rustsrc}
pub fn chan(chan: u64) -> StateObservable

pub fn gate(gate: HypernetRef) -> StateObservable

pub fn poly(poly: Poly) -> StateObservable

pub fn ind(arg: StatePred) -> StateObservable
\end{rustsrc}
\apiprov{construction --- \texttt{crates/mtk-catalogue/tests/golden\_spec.rs}}
\begin{rustsrc}
let ind_atom = StatePred::atom_ge(q(1, 1), StateObservable::ind(band.clone()));
let chan_obs = StateObservable::chan(1);
\end{rustsrc}

\subsection{SAtomLE and SAtomGE: threshold atoms}
\label{sec:spec-lang-atom}

\ruleblock{Formation}
\begin{mathpar}
\inferrule*[Right=\textsc{SAtomLE}]
  {c \in \mathbb{Q} \\ o\ \mathsf{StateObservable}}
  {\mathsf{SAtomLE}(c, o)\ \mathsf{StatePred}}

\inferrule*[Right=\textsc{SAtomGE}]
  {c \in \mathbb{Q} \\ o\ \mathsf{StateObservable}}
  {\mathsf{SAtomGE}(c, o)\ \mathsf{StatePred}}

\inferrule*[Right=$\Theta_0$]
  {c \in \mathbb{Q} \\ p\ \mathsf{Poly}}
  {\mathsf{SAtomLE}(c,\mathsf{oPoly}(p)),\,
   \mathsf{SAtomGE}(c,\mathsf{oPoly}(p))\ \Theta_0}
\end{mathpar}

The thresholding arrow of the grid: an atom bounds an
\emph{observable} against an exact rational, and the observable's
provenance is the previous entry's business, not the atom's. Over
the polynomial injection this is the closed polynomial halfspace ---
the kernel's native object, with no obligation standing between the
atom and the exact arithmetic that checks it, and the only form the
base stratum admits ($\Theta_0$ reads the \textsc{oPoly} image
alone). Over a channel or gate injection the same atom thresholds
the cited observable, its wall read at the pairing through the
computed elaboration. The surface
grammar is nonstrict \emph{by construction}: strict atoms are
compiler-internal artifacts of negation normal form and are
unrepresentable here --- the schema's testable negative; the
negation of a closed atom is the strict open one, a
normal-form-internal literal consumable at either parity
(\texttt{den\_compl\_le}). Boolean wires need no dedicated former:
``wire $b$ is set'' is spelled $b \ge 1$ on a $\{0,1\}$-valued wire
(the derived spelling below cites the circuit suite, whose atoms
compile to exactly this form).

\emph{Compile.} The wall of the closed sublevel set, with the
properness mark:
$\mathsf{SAtomLE}(c,\mathsf{oPoly}(p)) \mapsto
(\mathrm{wall}\{x : p(x) \le c\}, \{\top\})$, and symmetrically ---
channel and gate atoms compile the same wall species over the
observable the pairing resolves. Exact: satisfaction of the compiled
atom is containment of the permitted region in the halfspace
\emph{and} properness (\texttt{safetyAtom},
\texttt{coupledSat\_safetyAtom\_iff},
\texttt{safetyAtom\_modelClass}), and nothing expressible is lost by
compiling regions to threshold walls --- rational threshold rows
generate the flat universe (\texttt{iSup\_boxRows\_eq\_self}), the
generation theorem grounding the region view of
section~\ref{sec:spec-lang-target}.

\emph{Witness}: every certificate in the corpus.

\ruleblock{API}
\apilang{TypeScript}{spec.atomLe / atomGe / polyLe / polyGe}
\apiprov{signature --- \texttt{ts/src/spec.ts}}
\begin{tssrc}
export function atomLe(c: QLike, o: StateObservable): StatePred

export function atomGe(c: QLike, o: StateObservable): StatePred

export function polyLe(c: QLike, p: Poly): StatePred

export function polyGe(c: QLike, p: Poly): StatePred
\end{tssrc}
\apiprov{construction --- \texttt{ts/tests/spec-native.test.ts}}
\begin{tssrc}
const goal = s.polyGe(1n, s.unwrap(s.polyVar(0, 2)));
const cap = s.probGe(Q.of(3, 4), s.deadline(20n, goal));
\end{tssrc}
\apiprov{\texttt{polyLe} / \texttt{polyGe} are the polynomial-atom conveniences (atom at the \textsc{oPoly} injection); polynomials are minted canonical by \texttt{poly} / \texttt{polyVar} (the same module).}
\apilang{Rust}{mtk\_catalogue::spec::StatePred::\{atom\_le, atom\_ge\}}
\apiprov{signature --- \texttt{crates/mtk-catalogue/src/spec/formula.rs}}
\begin{rustsrc}
pub fn atom_le(c: WitRat, o: StateObservable) -> StatePred

pub fn atom_ge(c: WitRat, o: StateObservable) -> StatePred

pub fn poly_le(c: WitRat, poly: Poly) -> StatePred

pub fn poly_ge(c: WitRat, poly: Poly) -> StatePred
\end{rustsrc}
\apiprov{construction --- \texttt{crates/mtk-catalogue/tests/golden\_spec.rs}}
\begin{rustsrc}
let goal = StatePred::poly_ge(q(1, 1), Poly::var(0, 2).unwrap());
let cap =
    ForallSpec::prob_ge(q(3, 4), PathPred::deadline(q(20, 1), goal, false).unwrap()).unwrap();
\end{rustsrc}
\apiprov{Polynomials are minted canonical by \texttt{Poly::new} / \texttt{Poly::var} (crates/mtk-catalogue/src/spec/poly.rs).}

\subsection{SAnd and SOr: Boolean structure}
\label{sec:spec-lang-andor}

\ruleblock{Formation}
\begin{mathpar}
\inferrule*[Right=\textsc{SAnd}]
  {\varphi_i\ \mathsf{StatePred} \quad (i < n)}
  {\mathsf{SAnd}(\varphi_0, \ldots, \varphi_{n-1})\ \mathsf{StatePred}}

\inferrule*[Right=\textsc{SOr}]
  {\varphi_i\ \mathsf{StatePred} \quad (i < n)}
  {\mathsf{SOr}(\varphi_0, \ldots, \varphi_{n-1})\ \mathsf{StatePred}}

\inferrule*[Right=$\Theta_0$]
  {\varphi_i\ \Theta_0 \quad (i < n)}
  {\mathsf{SAnd}(\vec\varphi),\, \mathsf{SOr}(\vec\varphi)\ \Theta_0}
\end{mathpar}

$n$-ary as built; binary presentations are sugar, folding through
associativity of the exact cells.

\emph{Compile.} Conjunction is the \emph{atom merge},
disjunct-family-wise: safety floors join by pointwise maximum, mark
families union --- and the merge is exact, its class the
intersection of the classes (\texttt{SpecAtom.modelClass\_merge},
\texttt{SpecDnf.modelClass\_conj}; the safety face is the
conjunction-is-max universal property, the order-theoretic kin of
\texttt{sat\_sup\_iff}). Disjunction is the disjunct union outright,
and its class is the plain union
(\texttt{SpecDnf.modelClass\_disj}). Every merge carries the
properness mark by construction. One design fact deserves its theorem here,
because it is why the compile's semantics is \emph{class} semantics:
disjunction of presentations is deliberately \emph{not} the
pointwise-minimum coproduct of beliefs. A proper system can hold
its permitted region inside the union of two regions while belonging
to neither disjunct's class --- models of a disjunction may
\emph{hedge}, and the presentation semantics excludes hedging; the
gap is machine-checked (\texttt{exists\_hedging\_gap}), and the
region reading collapses to safe-and-proper against the denoted
region exactly on the disjunction-free fragment
(\texttt{elabNnf\_modelClass\_disjFree}).

\emph{Witness}: exercised everywhere; the full-grammar tour vector
drives both constructors directly.

\ruleblock{API}
\apilang{TypeScript}{spec.and / spec.or}
\apiprov{signature --- \texttt{ts/src/spec.ts}}
\begin{tssrc}
export function and(args: ReadonlyArray<StatePred>): StatePred

export function or(args: ReadonlyArray<StatePred>): StatePred
\end{tssrc}
\apiprov{construction --- \texttt{ts/tests/spec-native.test.ts}}
\begin{tssrc}
const orPart = s.or([s.FALSE, s.not(s.polyLe(0n, x1))]);
\end{tssrc}
\apilang{Rust}{mtk\_catalogue::spec::StatePred::\{and, or\}}
\apiprov{signature --- \texttt{crates/mtk-catalogue/src/spec/formula.rs}}
\begin{rustsrc}
pub fn and(args: Vec<StatePred>) -> StatePred

pub fn or(args: Vec<StatePred>) -> StatePred
\end{rustsrc}
\apiprov{construction --- \texttt{crates/mtk-catalogue/tests/golden\_spec.rs}}
\begin{rustsrc}
let or_part = StatePred::or(vec![
    StatePred::False,
    StatePred::not(StatePred::poly_le(q(0, 1), x1())),
]);
\end{rustsrc}

\subsection{SNot: negation under the closure discipline}
\label{sec:spec-lang-not}

\ruleblock{Formation}
\begin{mathpar}
\inferrule*[Right=\textsc{SNot}]
  {\varphi\ \mathsf{StatePred}}
  {\mathsf{SNot}\,\varphi\ \mathsf{StatePred}}
\end{mathpar}

Negation forms freely --- and concludes nothing else: no rule puts a
negation in the base stratum, and no compile clause consumes the
residues that negation normal form cannot resolve.

\emph{Compile.} By negation normal form. Negation is pushed to the
leaves (De~Morgan, double elimination), where a negated closed atom
becomes its strict open literal; normal forms have no negation node
\emph{by type}, so the compile's domain is a grammar in which
unconsumable residues are unrepresentable, and what survives is
compiled exactly --- the resolution preserves the region semantics
on the nose (\texttt{den\_nnf}). The residues that do not survive
--- a negated path atom at even parity (the strict complement of an
accumulated-cost event), negated Streett acceptance (Rabin), negated
bounded until, negated unbounded always --- are absences of rules
carrying the no-expressible-complement warrant of
section~\ref{sec:spec-lang-formation}: under the properness
presupposition the class algebra has no complements to compile to.
The producer-side typing preflight refuses each residue with a path
address into the formula tree, and the unbounded-always refusal
carries its remedy in the message: spell the intent as a
producer-chosen deadline (\textsc{PDeadline} with the
\texttt{producerT} flag).

\emph{Witness}: the negation-duality suite and the rewrite layer's
normal forms; the preflight's refusal geometry is test-pinned in
both producer languages.

\ruleblock{API}
\apilang{TypeScript}{spec.not}
\apiprov{signature --- \texttt{ts/src/spec.ts}}
\begin{tssrc}
export function not(arg: StatePred): StatePred

export function preflightState(f: StatePred): Result<undefined, SpecRefusal>
\end{tssrc}
\apiprov{construction --- \texttt{ts/tests/spec-native.test.ts}}
\begin{tssrc}
const orPart = s.or([s.FALSE, s.not(s.polyLe(0n, x1))]);
\end{tssrc}
\apilang{Rust}{mtk\_catalogue::spec::StatePred::not}
\apiprov{signature --- \texttt{crates/mtk-catalogue/src/spec/formula.rs + crates/mtk-catalogue/src/spec/typing.rs}}
\begin{rustsrc}
pub fn not(arg: StatePred) -> StatePred

pub fn preflight_state_at(f: &StatePred, at: &FormulaPath) -> Result<(), Refusal>
\end{rustsrc}
\apiprov{construction --- \texttt{crates/mtk-catalogue/tests/golden\_spec.rs}}
\begin{rustsrc}
let or_part = StatePred::or(vec![
    StatePred::False,
    StatePred::not(StatePred::poly_le(q(0, 1), x1())),
]);
\end{rustsrc}

\section{Path-Row Formers}
\label{sec:spec-lang-pathf}

This wing is the grid's Path row, whole
(section~\ref{sec:spec-lang-formation}): the temporal liftings that
form trajectory events, the run-functional stratum that accumulates
State-row observables along runs ($\mathsf{PathObservable}$, entered
after the liftings), and the row's own thresholding arrow
(\textsc{PAtom}) passing back from run-functionals to events.
Path predicates ($\mathsf{PathPred}$) denote sets of trajectories ---
\emph{events}, and an event carries no face of its own: it enters a
specification only through a cap former (the probability thresholds,
or indicator accumulation via \textsc{retPath} into the prevision
square), and the face --- floor or mark --- is declared by the
consumer (section~\ref{sec:spec-lang-faces}). One event, either
face, and the sort says which is meant,
every time. Their
compile target is the \emph{record} stratum of the semantic universe:
at the finitary site, a path former reads the set of \emph{occupied
stages} of a record --- and the site offers two readers of that set,
the stage graph (the \emph{touch} species) and the support of the
occupation measure (the \emph{dwell} species), with the record
coupling condition saying exactly that the dwell reader is contained
in the touch reader (\texttt{recordCon\_iff\_stages}). That one fact
organizes the event formers' semantics. Each former's transport between the
species is read off its \emph{variance} in the occupied-stage set:
invariance formers are antitone (a touch-certified claim transports
to the dwell claim), reach formers are monotone (the transport
flips), and the mixed formers --- until, Streett --- transport in
\emph{neither} direction, witnessed by explicit coupled separations
(\texttt{exists\_untilOn\_species\_incomparable},
\texttt{exists\_streettOn\_species\_incomparable}), so the absence
of those transports is content, not omission. Polarity dictates the
judgment face: invariance classes enter on the safety floor (the
wall, and the landed record walls are the invariance compile's
floors --- \texttt{touchWall\_eq\_wallOn},
\texttt{dwellWall\_eq\_wallOn}); reach classes enter on the liveness
marks (\texttt{markOn}, the mark twin of the wall), and the
properness demand is the everything-goal instance of the localized
liveness mark (\texttt{markOn\_univ}): every ``eventually'' claim
generalizes the bare existence claim by localizing its goal. The
general record atom carries one safety class and one liveness class,
and its satisfaction is \emph{safe and proper-into-the-goal}
(\texttt{coupledSat\_pathAtom\_iff}). At deterministic records ---
the canonical lift of a stage graph --- the two species coincide
(\texttt{stages\_uniformLift}) and the compile agrees with the
trajectory clauses stated at each entry on the nose; the stochastic
tier keeps only the directed transports. Windows at this stratum are
natural numbers, for the same reason as the summed-return horizon:
the finitary record site is tick-indexed and unitless; the
duration-typed rational readings of the rules below keep their force
at clocked continuum strata (validated-enclosure consumption), and the
encoding question is the same one, unchanged.

\subsection{PAlways: invariance}
\label{sec:spec-lang-always}

\ruleblock{Formation}
\begin{mathpar}
\inferrule*[Right=\textsc{PAlways}]
  {\varphi\ \mathsf{StatePred}}
  {\mathsf{PAlways}\,\varphi\ \mathsf{PathPred}}
\end{mathpar}

$\mathsf{G}\,\varphi$: the state formula holds at every time; no
carrier requirement of its own.

\emph{Compile.} The dwell class over the record site --- records
whose every occupied stage lies in the argument region
(\texttt{alwaysOn}); trajectory clause
$\{\omega : \forall t.\ \omega(t) \in
\llbracket \varphi \rrbracket\}$ at the deterministic cells. The
class is closed (\texttt{isClosed\_recordClass\_always}), antitone
in the occupied-stage set (touch certifies dwell), and enters the
judgment as a safety floor. The safety battery consumes it
(invariance and containment rows, chapter~\ref{ch:qtsl-manual}).

\emph{Witness}: every avoid and invariance certificate in the
corpus.

\ruleblock{API}
\apilang{TypeScript}{spec.always}
\apiprov{signature --- \texttt{ts/src/spec.ts}}
\begin{tssrc}
export function always(arg: StatePred): PathPred
\end{tssrc}
\apiprov{construction --- \texttt{ts/tests/spec-native.test.ts}}
\begin{tssrc}
const alwaysPart = s.probGe(1n, s.always(band));
\end{tssrc}
\apilang{Rust}{mtk\_catalogue::spec::PathPred::always}
\apiprov{signature --- \texttt{crates/mtk-catalogue/src/spec/formula.rs}}
\begin{rustsrc}
pub fn always(arg: StatePred) -> PathPred
\end{rustsrc}
\apiprov{construction --- \texttt{crates/mtk-catalogue/tests/golden\_spec.rs}}
\begin{rustsrc}
let always_part = ForallSpec::prob_ge(q(1, 1), PathPred::always(band.clone())).unwrap();
\end{rustsrc}

\subsection{PAlwaysT: bounded always}
\label{sec:spec-lang-alwayst}

\ruleblock{Formation}
\begin{mathpar}
\inferrule*[Right=\textsc{PAlwaysT}]
  {T \in \mathbb{Q} \\ T > 0 \\ \varphi\ \mathsf{StatePred}}
  {\mathsf{PAlwaysT}(T, \varphi)\ \mathsf{PathPred}}
\end{mathpar}
\[
\mathsf{PAlwaysT}(T,\varphi) \rhd \{\mathsf{duration}\}
\]

The state formula holds at every time up to the deadline --- the
claim shape of every validated-enclosure certificate (``the flow
from the initial box stays in the tube over $[0,T]$''), which is why
the former exists: this catalogue holds that the spec as stated must
be the spec certified, and this statement deserved a spelling.

\emph{Compile.} The bounded dwell class: occupied stages within the
window lie in the region (\texttt{windowAlwaysOn}; window a natural
number at the tick-indexed stratum, per the wing preamble). The
class is closed (\texttt{isClosed\_recordClass\_alwaysT}), antitone
in the
occupied-stage set and in the window --- the window is an
approximation source: a longer-window certificate consumes a
shorter-window claim. Deterministic consumption is immediate
(enclosure semantics is trajectory-wise); stochastic consumption at
even parity awaits the over-approximating tower --- the same fence
as the probability upper bounds
(section~\ref{sec:spec-lang-falwaysstoch}). Its normal-form dual
\emph{is} spellable: negated bounded always rewrites to bounded
reach, so surface path negation over this former passes the
preflight at both parities.

\emph{Witness for the demand}: the enclosure fixture family
(flow-tube certificates); the safety-until-deadline worked example;
the LunarLander funnel holds.

\ruleblock{API}
\apilang{TypeScript}{spec.alwaysT}
\apiprov{signature --- \texttt{ts/src/spec.ts}}
\begin{tssrc}
export function alwaysT(t: QLike, arg: StatePred): PathPred
\end{tssrc}
\apiprov{construction --- \texttt{ts/tests/spec-native.test.ts}}
\begin{tssrc}
const boundedPart = s.probGe(Q.of(1, 2), s.notPath(s.alwaysT(3n, band)));
\end{tssrc}
\apilang{Rust}{mtk\_catalogue::spec::PathPred::always\_t}
\apiprov{signature --- \texttt{crates/mtk-catalogue/src/spec/formula.rs}}
\begin{rustsrc}
pub fn always_t(t: WitRat, arg: StatePred) -> Result<PathPred, Refusal>
\end{rustsrc}
\apiprov{construction --- \texttt{crates/mtk-catalogue/tests/golden\_spec.rs}}
\begin{rustsrc}
let bounded_part = ForallSpec::prob_ge(
    q(1, 2),
    PathPred::not_path(PathPred::always_t(q(3, 1), band.clone()).unwrap()),
)
.unwrap();
\end{rustsrc}

\subsection{PDeadline: bounded reach}
\label{sec:spec-lang-deadline}

\ruleblock{Formation}
\begin{mathpar}
\inferrule*[Right=\textsc{PDeadline}]
  {T \in \mathbb{Q} \\ T > 0 \\ \varphi\ \mathsf{StatePred} \\ b \in \mathbb{B}}
  {\mathsf{PDeadline}(T, \varphi, b)\ \mathsf{PathPred}}
\end{mathpar}
\[
\mathsf{PDeadline}(T,\varphi,b) \rhd \{\mathsf{duration}\}
\]

The target region is reached by the deadline. The flag
$b$ (\texttt{producerT}) marks a producer-chosen horizon standing in
for an unbounded reach --- a sound under-approximation, carried
\emph{flagged} in the row so the reader of a certificate knows which
kind of statement was made (the unbounded primitive is a deliberate
absence; the derived spelling appears with the fences below).

\emph{Compile.} The d\'ebut (hitting) record class by the deadline:
some occupied stage within the window lies in the target
(\texttt{windowReachOn}); trajectory clause
$\{\omega : \exists t \le T.\ \omega(t) \in
\llbracket \varphi \rrbracket\}$ at the deterministic cells. The
class is open (\texttt{isOpen\_recordClass\_deadline}), monotone in
the occupied-stage set (dwell certifies touch --- the variance
flip), and enters the judgment on the liveness face. Obligations:
the deadline-battery kernels (geometric-decay and
expectation-of-hitting arguments), machine-checked per the Rule
Catalogue's receipt columns.

\emph{Witness}: reach-by-deadline certificates across the corpus; the
certified-Gymnasium hitting waves (twenty-eight certificates of the
quoted shape).

\ruleblock{API}
\apilang{TypeScript}{spec.deadline}
\apiprov{signature --- \texttt{ts/src/spec.ts}}
\begin{tssrc}
export function deadline(
  t: QLike,
  arg: StatePred,
  opts?: { producerT?: boolean },
): PathPred
\end{tssrc}
\apiprov{construction --- \texttt{ts/tests/spec-native.test.ts}}
\begin{tssrc}
const goal = s.polyGe(1n, s.unwrap(s.polyVar(0, 2)));
const cap = s.probGe(Q.of(3, 4), s.deadline(20n, goal));
\end{tssrc}
\apilang{Rust}{mtk\_catalogue::spec::PathPred::deadline}
\apiprov{signature --- \texttt{crates/mtk-catalogue/src/spec/formula.rs}}
\begin{rustsrc}
pub fn deadline(t: WitRat, arg: StatePred, producer_t: bool) -> Result<PathPred, Refusal>
\end{rustsrc}
\apiprov{construction --- \texttt{crates/mtk-catalogue/tests/golden\_spec.rs}}
\begin{rustsrc}
let goal = StatePred::poly_ge(q(1, 1), Poly::var(0, 2).unwrap());
let cap =
    ForallSpec::prob_ge(q(3, 4), PathPred::deadline(q(20, 1), goal, false).unwrap()).unwrap();
\end{rustsrc}

\subsection{PUntil: safety until a deadline}
\label{sec:spec-lang-until}

\ruleblock{Formation}
\begin{mathpar}
\inferrule*[Right=\textsc{PUntil}]
  {T \in \mathbb{Q} \\ T > 0 \\
   \varphi_{\mathrm{safe}}\ \mathsf{StatePred} \\
   \varphi_{\mathrm{tgt}}\ \mathsf{StatePred}}
  {\mathsf{PUntil}(T, \varphi_{\mathrm{safe}}, \varphi_{\mathrm{tgt}})\ \mathsf{PathPred}}
\end{mathpar}
\[
\mathsf{PUntil}(T,\cdot,\cdot) \rhd \{\mathsf{duration}\}
\]

Stay safe until reaching the target, by the deadline; trajectory
clause
$\{\omega : \exists t \le T.\ \omega(t) \in
\llbracket \varphi_{\mathrm{tgt}} \rrbracket
\wedge \forall s < t.\ \omega(s) \in
\llbracket \varphi_{\mathrm{safe}} \rrbracket\}$.

\emph{Compile.} The hitting-with-dwell record class
(\texttt{untilOn}): reach the target within the window, occupied
stages strictly before staying safe. The mixed variance shows up
everywhere the pure formers were simple: the class transports
between the species in \emph{neither} direction (witnessed, per the
wing preamble), and its citizenship is the grammar's first genuinely
constructible-not-closed-not-open class --- a finite union of
locally closed pieces by window decomposition
(\texttt{recordClass\_untilF\_eq},
\texttt{isConstructibleUnion\_recordClass\_untilF}), landing it
squarely in the constructible stratum the compile target promises.
The union-bound kernel that consumes it is machine-checked; negation
of a bounded until has no sound compile target in this edition and
is refused by the preflight with a path address.

\emph{Witness}: the until-battery fixtures and the worked
safety-until-deadline example.

\ruleblock{API}
\apilang{TypeScript}{spec.until}
\apiprov{signature --- \texttt{ts/src/spec.ts}}
\begin{tssrc}
export function until(t: QLike, safe: StatePred, target: StatePred): PathPred
\end{tssrc}
\apiprov{construction --- \texttt{ts/tests/spec-native.test.ts}}
\begin{tssrc}
const untilPart = s.probGe(Q.of(9, 10), s.until(5n, band, s.polyGe(2n, x1)));
\end{tssrc}
\apilang{Rust}{mtk\_catalogue::spec::PathPred::until}
\apiprov{signature --- \texttt{crates/mtk-catalogue/src/spec/formula.rs}}
\begin{rustsrc}
pub fn until(t: WitRat, safe: StatePred, target: StatePred) -> Result<PathPred, Refusal>
\end{rustsrc}
\apiprov{construction --- \texttt{crates/mtk-catalogue/tests/golden\_spec.rs}}
\begin{rustsrc}
let until_part = ForallSpec::prob_ge(
    q(9, 10),
    PathPred::until(q(5, 1), band.clone(), StatePred::poly_ge(q(2, 1), x1())).unwrap(),
)
.unwrap();
\end{rustsrc}

\subsection{PDwellStreett and PCrossStreett: acceptance, metered}
\label{sec:spec-lang-streett}

\ruleblock{Formation}
\begin{mathpar}
\inferrule*[Right=\textsc{PairBase}]
  {\rho\ \mathsf{StatePred} \\ \rho\ \Theta_0 \\ \gamma\ \mathsf{StatePred}}
  {\mathsf{PairBase}(\rho, \gamma)\ \mathsf{StreettPair}}

\inferrule*[Right=\textsc{PairCmpl}]
  {\rho\ \mathsf{StatePred} \\ \gamma\ \mathsf{StatePred}}
  {\mathsf{PairCmpl}(\rho, \gamma)\ \mathsf{StreettPair}}
\\
\inferrule*[Right=\textsc{PDwellStreett}]
  {\pi_i\ \mathsf{StreettPair} \quad (i < n)}
  {\mathsf{PDwellStreett}(\pi_0, \ldots, \pi_{n-1})\ \mathsf{PathPred}}
\\
\inferrule*[Right=\textsc{PCrossStreett}]
  {\pi_i\ \mathsf{StreettPair} \quad (i < n)}
  {\mathsf{PCrossStreett}(\pi_0, \ldots, \pi_{n-1})\ \mathsf{PathPred}}
\end{mathpar}

Streett acceptance~\cite{streett1982} --- for each pair, an
unbounded request forces an unbounded response --- as \emph{two}
formers, because ``unbounded'' has two meterings and the record
lens proves them incomparable: the mixed formers transport between
the stage-graph reader and the occupation reader in neither
direction, even on coupled points (the wing preamble's witnessed
separations). A metering flag would leave one former whose every
consuming rule and citizenship theorem still cased on the flag with
the type system enforcing nothing; two formers make the meter part
of the \emph{event's identity} --- the same move as the polarity
redesign, one stratum down: metering as former, not annotation.
$\mathsf{PDwellStreett}$ is occupation metering --- unbounded
request \emph{occupation} forces unbounded response occupation
(counting measure at tick strata, duration at clocked strata); you
dwell on a street. $\mathsf{PCrossStreett}$ is entry-event counting
--- unboundedly many request \emph{crossings} (fresh contact after
absence, the countable, payable unit of touching) force unboundedly
many response crossings; you cross a street. Enter-and-stay
separates the species even at unit ticks --- occupation diverges
while the entry counter stops at one --- and the two meterings at
one pair list are distinct events, refuse-mode disjoint at the row
stratum. Everything step-metered the catalogue already had is the
dwell former's unit-tick instance \emph{by identity}: the liveness
batteries, the acceptance event, and the recurrence sugar migrate
there whole, and nothing migrates to the cross species, whose
consuming machinery composes with the crossing-budget family at the
behavioural tier. Where the species provably coincide ---
deterministic records, where the two stage readers agree --- either
former may be formed and the collapse theorem transports.

The pair vocabulary is metering-independent, shared verbatim. The
two pair constructors are the two request-consumption modes ---
\emph{parity as constructor choice, never stored data}:
$\mathsf{PairBase}$
consumes the request base-exact, whence its $\Theta_0$ premise
(whose deliberate absence from $\mathsf{PairCmpl}$ is the point);
$\mathsf{PairCmpl}$ consumes the request through the syntactic
complement's over-approximation, stored positive. The recurrence and
persistence operators are definitional instances, not formers:
$\mathsf{GF}\,\varphi$ is a single dwell pair with a trivially-true
base-exact request, and $\mathsf{FG}\,\varphi$ is the dwell pair
whose complement-consumed request firing unboundedly forces the
impossible response (both appear as derived spellings below). The
face of an acceptance claim stays at the cap that consumes the
event --- two meterings by two faces, all four spelled
compositionally. Negated Streett is Rabin acceptance --- the dual
judgment, a deliberate absence under either metering.

\emph{Compile.} Two tiers, kept distinct. At the behavioural
tier the dwell acceptance event lives on the kernel path measure,
and the quantitative certificate bundle concluding it is
machine-checked end to end (\texttt{sound\_QStreettBundle},
composing into the probability-threshold region at entry
\textsc{SProbGE}); the cross event's infinitary reading is the
entry-counter family, whose very definedness at continuum strata is
the crossing-budget machinery's evidence. At the record
stratum the compile is the per-pair implication over occupied
stages (\texttt{streettOn}), each former reading its \emph{own}
species --- the dwell former the occupation support, the cross
former the stage graph --- as the \emph{occupation shadow} of the
infinitary acceptance condition, stated as such: a finite record has
no ``unboundedly often,'' and the pasting between the record tier
and the kernel path measure is deliberately not claimed here (it is
the finite-dimensional-distributions stratum, the same fence the
behavioural tier names). Both record classes are constructible;
species transport fails in both directions, witnessed --- which is
the split's warrant, restated as compile fact.

\emph{Witness}: the 219-task liveness suite (specs normalized to
$\bigwedge_i \mathsf{GF}\,\gamma_i$ --- dwell spellings --- and
certified against circuit models); the cross former is new with the
theory revision, its first certificate lane the crossing-budget
composition.

\ruleblock{API}
\apilang{TypeScript}{spec.dwellStreett / crossStreett / pairBase / pairCmpl}
\apiprov{signature --- \texttt{ts/src/spec.ts}}
\begin{tssrc}
export function dwellStreett(pairs: ReadonlyArray<StreettPair>): PathPred

export function crossStreett(pairs: ReadonlyArray<StreettPair>): PathPred

export function pairBase(req: StatePred, resp: StatePred): StreettPair

export function pairCmpl(req: StatePred, resp: StatePred): StreettPair
\end{tssrc}
\apiprov{construction --- \texttt{ts/tests/spec-native.test.ts}}
\begin{tssrc}
const dwellPart = s.probGe(
  1n,
  s.dwellStreett([s.pairCmpl(band, s.polyGe(1n, x1))]),
);
const crossPart = s.probLe(
  Q.of(1, 2),
  s.crossStreett([s.pairBase(band, s.polyGe(1n, x1))]),
);
\end{tssrc}
\apiprov{A base-exact pair is also driven (through the recurrence spelling) by the liveness vector quoted under the derived spellings.}
\apilang{Rust}{mtk\_catalogue::spec::PathPred::\{dwell\_streett, cross\_streett\}}
\apiprov{signature --- \texttt{crates/mtk-catalogue/src/spec/formula.rs}}
\begin{rustsrc}
pub fn dwell_streett(pairs: Vec<StreettPair>) -> PathPred

pub fn cross_streett(pairs: Vec<StreettPair>) -> PathPred

pub fn base(req: StatePred, resp: StatePred) -> Result<StreettPair, Refusal>

pub fn cmpl(req: StatePred, resp: StatePred) -> StreettPair
\end{rustsrc}
\apiprov{construction --- \texttt{crates/mtk-catalogue/tests/golden\_spec.rs}}
\begin{rustsrc}
let dwell_part = ForallSpec::prob_ge(
    q(1, 1),
    PathPred::dwell_streett(vec![StreettPair::cmpl(
        band.clone(),
        StatePred::poly_ge(q(1, 1), x1()),
    )]),
)
.unwrap();
let cross_part = ForallSpec::prob_le(
    q(1, 2),
    PathPred::cross_streett(vec![StreettPair::base(
        band.clone(),
        StatePred::poly_ge(q(1, 1), x1()),
    )
    .unwrap()]),
)
.unwrap();
\end{rustsrc}

\subsection{The Return stratum}
\label{sec:spec-lang-returnstratum}

The Observable cell of the Path row: a
$\mathsf{PathObservable}$ (the sort the earlier edition called
\textsf{Return}) is a \emph{run-functional} --- what a run is worth,
packaged as data --- consumed by the expectation, level, and regret
faces of section~\ref{sec:spec-lang-faces} and thresholded back into
events by \textsc{PAtom} below. The cell is reached by accumulation
from State-row material, and the
theory spells it so: reward \emph{provenance} (channel,
gate, polynomial view, indicator) is the injection stratum of
section~\ref{sec:spec-lang-obs}, and \emph{accumulation} is five
formers over one observable --- the earlier edition's
provenance-fused constructor roster (\textsf{retChanSum} through
\textsf{retAvg}) re-founded as injections plus accumulators, at
unchanged denotations. Beside the accumulators stand the two
citation formers: \textsc{retPath}, the Path row's indicator arrow,
and \textsc{retFireCount}, the transition-count functional.

\ruleblock{Formation}
\begin{mathpar}
\inferrule*[Right=\textsc{poSum}]
  {n \in \mathbb{N} \\ n \ge 1 \\ o\ \mathsf{StateObservable}}
  {\mathsf{poSum}(n, o)\ \mathsf{PathObservable}}
\and
\inferrule*[Right=\textsc{poInt}]
  {T \in \mathbb{Q} \\ T > 0 \\ o\ \mathsf{StateObservable}}
  {\mathsf{poInt}(T, o)\ \mathsf{PathObservable}}
\\
\inferrule*[Right=\textsc{poDisc}]
  {\theta \in \mathbb{Q} \\ \theta > 0 \\ o\ \mathsf{StateObservable}}
  {\mathsf{poDisc}(o, \theta)\ \mathsf{PathObservable}}
\and
\inferrule*[Right=\textsc{poAvg}]
  {o\ \mathsf{StateObservable}}
  {\mathsf{poAvg}(o)\ \mathsf{PathObservable}}
\\
\inferrule*[Right=\textsc{poSup}]
  {o\ \mathsf{StateObservable}}
  {\mathsf{poSup}(o)\ \mathsf{PathObservable}}
\and
\inferrule*[Right=\textsc{retPath}]
  {\psi\ \mathsf{PathPred}}
  {\mathsf{retPath}(\psi)\ \mathsf{PathObservable}}
\\
\inferrule*[Right=\textsc{retFireCount}]
  {R\ \mathsf{TrList} \\ R \neq \mathsf{nil}}
  {\mathsf{retFireCount}(R)\ \mathsf{PathObservable}}
\end{mathpar}
\[
\mathsf{poSum} \rhd \mathsf{steps}
\qquad
\mathsf{poInt} \rhd \mathsf{duration}
\]

The two finite-horizon accumulators carry the integrality content
the mechanization forced, now by \emph{shape}: a discrete sum's
upper index is a natural number, and the tick-indexed strata that
consume summed accumulation are unitless --- no clock exists there
against which a rational horizon could be rounded, which also rules
out any rounding law at those strata (a system clock would make one
spec denote different claims per clock). Integrated accumulation
keeps its free positive rational horizon, which keeps its meaning at
clocked continuum strata. The encoding decision the earlier edition
left open is settled the strong way: split constructors with a
natural-number tick horizon, so a rational tick horizon is
unrepresentable rather than checked --- and with the horizon split
carried by the constructors, an integrality side condition has
nothing left to check.

The remaining accumulators carry no horizon.
$\mathsf{poDisc}$ is infinite-horizon
exponential discounting at rate $\theta$: weights $e^{-\theta t}$ on
the internalized real-time base, under which tick systems embed at
integer times with weights $e^{-\theta n}$ --- kept exact in the
log-rational grade, the same abstract-exponential discipline that
keeps every transcendental out of the exact stratum.
$\mathsf{poAvg}$ is the long-run average (the Ces\`aro limit); its
certificate species is the
gain/bias pair --- a Poisson-equation drift row, the same
supermartingale-shaped row grammar as everything above.
$\mathsf{poSup}(o)$ is the pathwise running supremum of the
observable along the run --- horizon-free, read in the extended-real
scale, where an infinite supremum makes an upper-bound claim
\emph{false}, never ill-formed. Its consuming shape is the maximal
function: $\mathsf{SReturnLE}(B, \mathsf{poSup}(o))$ bounds
$\mathbb{E}[\sup_t o(x_t)]$ over every admitted law, for which
Doob's maximal inequality is the natural certificate engine and the
supermartingale row kit the natural feed --- named with its species,
its rule rows not yet claimed. Deliberately \emph{not} minted beside
it: a sup-outside-expectation former --- $\sup_t \mathbb{E}[\cdot]$
is a forall-$t$ family of marginal-law judgments, not a
run-functional --- and since
$\mathbb{E}[\sup] \ge \sup_t \mathbb{E}$ always, the
per-instant-uniform reading arrives as a free corollary of the
expressible, stronger form.
$\mathsf{retPath}(\psi)$ cites a path event as its indicator
run-functional; temporal claims compile \emph{into} this stratum
(monitor automaton, product kernel, compiled state-reward gate,
accumulation), following the value-functions-as-certificates
route~\cite{abate26b}, under which the reward designs for Streett
acceptance are state-valued indicators on the product carrier and
the ordinary value function of a satisfying policy \emph{is},
affinely, the certificate. The exact deadline table of
chapter~\ref{ch:qtsl-manual} is the value function of the
compiled deadline reward --- a reading, not a new obligation --- and
reinforcement learning becomes a sanctioned \emph{producer} for the
battery rows: powerful untrusted producers, small trustworthy
checkers, as everywhere.

$\mathsf{retFireCount}(R)$ is the transition-count functional: the
total number of firings of the roster's transitions along the run
--- no horizon, values in $[0, \infty]$, the same extended reading
as the running supremum on its nonnegative fragment. The roster $R$
is a positional list of transition citations of the paired subject;
nonemptiness is formation (an empty roster would be the zero
functional --- a constant claim in count clothing, and canonical
form forbids re-spelling), while duplicate-freedom and roster
membership are pairing-relative demands at the certificate stratum,
where the subject is in hand. Counting is firing-metered per roster
member --- deliberately \emph{not} region-entry counting (that is
\textsc{PCrossStreett}'s device,
section~\ref{sec:spec-lang-streett}) and deliberately not the
instant-crossing compensator mass of the Zeno seam: different
species, not interchangeable. The consuming shape this former
exists for is the hard count ceiling,
$\mathsf{SReturnLE}(B, \mathsf{retFireCount}(R))$ --- along every
admitted law, the total count of roster firings is at most $B$
(section~\ref{sec:spec-lang-devcount}). The $\ge$ face --- at least
$n$ events, a genuine liveness claim consuming survival mass --- is
deliberately absent from this functional: that is throughput
territory, and it enters at the expectation stratum rather than
here --- the event-rate envelope family's lower face bounds the
expected count from below, horizon-parametrically, over the same
rosters (section~\ref{sec:spec-lang-rate}).

\ruleblock{API}
\apilang{TypeScript}{spec.poSum / poInt / poDisc / poAvg / poSup / retPath / retFireCount}
\apiprov{signature --- \texttt{ts/src/spec.ts}}
\begin{tssrc}
export function poSum(horizon: number, o: StateObservable): PathObservable

export function poInt(horizon: QLike, o: StateObservable): PathObservable

export function poDisc(o: StateObservable, rate: QLike): PathObservable

export function poAvg(o: StateObservable): PathObservable

export function poSup(o: StateObservable): PathObservable

export function retPath(arg: PathPred): PathObservable

export function retFireCount(trs: ReadonlyArray<number>): PathObservable
\end{tssrc}
\apiprov{construction --- \texttt{ts/tests/spec-native.test.ts}}
\begin{tssrc}
const retLe = s.returnLe(100n, s.poInt(50n, s.oPoly(x0)));
const discFloor = s.returnGe(-5n, s.poDisc(chanObs, Q.of(1, 10)));
const countCeiling = s.returnLe(40n, s.retFireCount([0, 2]));
\end{tssrc}
\apilang{Rust}{mtk\_catalogue::spec::PathObservable::\{sum, int, disc, avg, sup, path, fire\_count\}}
\apiprov{signature --- \texttt{crates/mtk-catalogue/src/spec/formula.rs}}
\begin{rustsrc}
pub fn sum(horizon: u64, o: StateObservable) -> Result<PathObservable, Refusal>

pub fn int(horizon: WitRat, o: StateObservable) -> Result<PathObservable, Refusal>

pub fn disc(o: StateObservable, rate: WitRat) -> Result<PathObservable, Refusal>

pub fn avg(o: StateObservable) -> PathObservable

pub fn sup(o: StateObservable) -> PathObservable

pub fn path(arg: PathPred) -> PathObservable

pub fn fire_count(trs: Vec<u64>) -> Result<PathObservable, Refusal>
\end{rustsrc}
\apiprov{construction --- \texttt{crates/mtk-catalogue/tests/golden\_spec.rs}}
\begin{rustsrc}
let ret_le = ForallSpec::return_le(
    q(100, 1),
    PathObservable::int(q(50, 1), StateObservable::poly(x0())).unwrap(),
);
let disc_floor =
    ForallSpec::return_ge(q(-5, 1), PathObservable::disc(chan_obs, q(1, 10)).unwrap());
let count_ceiling =
    ForallSpec::return_le(q(40, 1), PathObservable::fire_count(vec![0, 2]).unwrap());
\end{rustsrc}

\subsection{PAtomLE and PAtomGE: accumulated-cost events}
\label{sec:spec-lang-patom}

\ruleblock{Formation}
\begin{mathpar}
\inferrule*[Right=\textsc{PAtomLE}]
  {c \in \mathbb{Q} \\ R\ \mathsf{PathObservable}}
  {\mathsf{PAtomLE}(c, R)\ \mathsf{PathPred}}
\and
\inferrule*[Right=\textsc{PAtomGE}]
  {c \in \mathbb{Q} \\ R\ \mathsf{PathObservable}}
  {\mathsf{PAtomGE}(c, R)\ \mathsf{PathPred}}
\end{mathpar}

The thresholding arrow at the Path row --- the grid's admitted
square-filler, entered with the theory revision: the \emph{event}
that the accumulated value lands on the stated side of the
threshold, the inner clause of
$\mathbb{P}_{\ge p}[\sum r \le B]$. This is a claim species
distinct from any expectation bound, and the distinction is the
grid read column-wise: $\mathsf{SReturnLE}(B, R)$ bounds the
\emph{prevision} of $R$ --- one number per law, capped at the
credal body --- while $\mathsf{PAtomLE}(B, R)$ names the set of
\emph{runs} on which $R$ lands under $B$, an event like any other,
whose face is then declared by the cap that consumes it (a
probability threshold, or indicator accumulation via
\textsc{retPath}). Carrier requirements are inherited, not minted:
the synthesis reads the run-functional argument
($\mathsf{PAtomLE}(c, R) \rhd \kappa$ exactly when $R \rhd
\kappa$).

\emph{Compile.} Owed with the former's first certificate lane ---
formation and the traversal stratum are landed at the theory
revision; the record-stratum compile clause and its adequacy cell
are named obligations, and no battery consumes the event yet.

\emph{Witness for the demand}: the count-ceiling and budget
families, whose probabilistic faces are exactly
$\mathbb{P}$-thresholds over accumulated-cost events.

\ruleblock{API}
\apilang{TypeScript}{spec.pathAtomLe / pathAtomGe}
\apiprov{signature --- \texttt{ts/src/spec.ts}}
\begin{tssrc}
export function pathAtomLe(c: QLike, ret: PathObservable): PathPred

export function pathAtomGe(c: QLike, ret: PathObservable): PathPred
\end{tssrc}
\apiprov{construction --- \texttt{ts/tests/spec-native.test.ts}}
\begin{tssrc}
const pathEvent = s.pathAtomGe(0n, s.retPath(s.always(indAtom)));
const eventCap = s.probGe(Q.of(3, 4), pathEvent);
\end{tssrc}
\apilang{Rust}{mtk\_catalogue::spec::PathPred::\{atom\_le, atom\_ge\}}
\apiprov{signature --- \texttt{crates/mtk-catalogue/src/spec/formula.rs}}
\begin{rustsrc}
pub fn atom_le(c: WitRat, ret: PathObservable) -> PathPred

pub fn atom_ge(c: WitRat, ret: PathObservable) -> PathPred
\end{rustsrc}
\apiprov{construction --- \texttt{crates/mtk-catalogue/tests/golden\_spec.rs}}
\begin{rustsrc}
let path_event = PathPred::atom_ge(
    q(0, 1),
    PathObservable::path(PathPred::always(ind_atom.clone())),
);
let event_cap = ForallSpec::prob_ge(q(3, 4), path_event).unwrap();
\end{rustsrc}

\subsection{PNotPath: surface path negation}
\label{sec:spec-lang-notpath}

\ruleblock{Formation}
\begin{mathpar}
\inferrule*[Right=\textsc{PNotPath}]
  {\psi\ \mathsf{PathPred}}
  {\mathsf{PNotPath}\,\psi\ \mathsf{PathPred}}
\end{mathpar}

Surface-only path negation, resolved by the dualities in negation
normal form.

\emph{Compile.} By the normal-form dualities --- and the fence table
is \emph{executable}: normal-form resolution is a total function
into either a negation-free form or a named refusal, and each row of
the table is a definitional equation. The one always-spellable dual
resolves at both parities --- negated bounded always \emph{is}
bounded reach of the complement
(\texttt{pathNnf\_notPath\_alwaysT}), and the even parity of a
bounded reach resolves through double negation
(\texttt{pathNnf\_notPath\_notPath\_deadline}) --- while the three
residues refuse with their addresses: negated unbounded always,
negated until, negated Streett (\texttt{pathNnf\_notPath\_always},
\texttt{pathNnf\_notPath\_untilF},
\texttt{pathNnf\_notPath\_streett}). The path-atom formers of the
theory revision extend the refusal roster at the producer surface: a
negated path atom at even parity is the strict complement of an
accumulated-cost event --- a normal-form artifact with no surface
spelling --- refused with the same path-address geometry, as a
documented extension of the pinned oracle. One note on the warrants,
because they differ by stratum: at the record-class level
the complement of an invariance class \emph{does} exist
(\texttt{not\_alwaysOn\_iff} --- some occupied stage escapes), so
the refusal of negated unbounded always is grammar economy --- there
is no unbounded-reach former to spell the dual, itself a deliberate
absence --- whereas at the spec-class stratum one level up, where
the judgment kinds live, the refusals carry the
no-expressible-complement theorem of
section~\ref{sec:spec-lang-formation}.

\emph{Witness}: the negation-duality suite; the full-grammar tour
vector's $\mathsf{PNotPath}(\mathsf{G}_{\le 3}\,\varphi)$ passes the
preflight at both parities.

\ruleblock{API}
\apilang{TypeScript}{spec.notPath}
\apiprov{signature --- \texttt{ts/src/spec.ts}}
\begin{tssrc}
export function notPath(arg: PathPred): PathPred
\end{tssrc}
\apiprov{construction --- \texttt{ts/tests/spec-native.test.ts}}
\begin{tssrc}
const boundedPart = s.probGe(Q.of(1, 2), s.notPath(s.alwaysT(3n, band)));
\end{tssrc}
\apilang{Rust}{mtk\_catalogue::spec::PathPred::not\_path}
\apiprov{signature --- \texttt{crates/mtk-catalogue/src/spec/formula.rs}}
\begin{rustsrc}
pub fn not_path(arg: PathPred) -> PathPred
\end{rustsrc}
\apiprov{construction --- \texttt{crates/mtk-catalogue/tests/golden\_spec.rs}}
\begin{rustsrc}
let bounded_part = ForallSpec::prob_ge(
    q(1, 2),
    PathPred::not_path(PathPred::always_t(q(3, 1), band.clone()).unwrap()),
)
.unwrap();
\end{rustsrc}

\section{Spec-Face Formers}
\label{sec:spec-lang-faces}

The wing above the grid. Every former here consumes run-valued
material --- an event, a run-functional, or (for the rate
envelopes) a transition roster read as a horizon-parametric
counting process --- reads it against the
credal body, and concludes into a polarity cap sort, so this wing is
where a claim first acquires a \emph{face}: the stratification law
of section~\ref{sec:spec-lang-formation} puts all quantification
over admitted laws here and nowhere else. Formers concluding
$\mathsf{ForallSpec}$ are the safety faces --- the bound holds over
\emph{every} admitted law, compiling to the floor --- and the
$\mathsf{MayBe}$ and $\mathsf{Mark}$ formers conclude
$\mathsf{ExistsSpec}$, the possibility faces: \emph{some} admitted
law (or, for regret, some filler) attains, compiling to a mark. The
roster is polarity-total --- probability, expectation, and level
each present both faces, and the regret pair is the same duality
with the quantifier moved to the filler seat. Two cap formers live
with their consumers instead of here: the initial-state injection
\textsc{SInit}, stated with the record it feeds
(section~\ref{sec:spec-lang-record}), and the attitude pair, which
mixes envelopes rather than declaring a face and closes the
judgment-kind inventory (section~\ref{sec:spec-lang-hurwicz}).

\subsection{SProbGE, SProbLE, and the probability marks}
\label{sec:spec-lang-prob}

\ruleblock{Formation}
\begin{mathpar}
\inferrule*[Right=\textsc{SProbGE}]
  {p \in \mathbb{Q} \\ 0 < p \le 1 \\ \psi\ \mathsf{PathPred}}
  {\mathsf{SProbGE}(p, \psi)\ \mathsf{ForallSpec}}
\and
\inferrule*[Right=\textsc{SProbLE}]
  {q \in \mathbb{Q} \\ 0 \le q < 1 \\ \psi\ \mathsf{PathPred}}
  {\mathsf{SProbLE}(q, \psi)\ \mathsf{ForallSpec}}
\\
\inferrule*[Right=\textsc{SProbMayBeGE}]
  {p \in \mathbb{Q} \\ 0 < p \le 1 \\ \psi\ \mathsf{PathPred}}
  {\mathsf{SProbMayBeGE}(p, \psi)\ \mathsf{ExistsSpec}}
\\
\inferrule*[Right=\textsc{SProbMayBeLE}]
  {q \in \mathbb{Q} \\ 0 \le q < 1 \\ \psi\ \mathsf{PathPred}}
  {\mathsf{SProbMayBeLE}(q, \psi)\ \mathsf{ExistsSpec}}
\end{mathpar}

One species, four faces. $\mathsf{SProbGE}(p, \psi)$ --- historically
this stack's only probability former, under the name $\mathsf{SProb}$
--- asserts that the path event holds with
probability at least $p$, read as \emph{lower} probability against
the credal view, over every admitted law --- whence its
$\mathsf{ForallSpec}$ sort; the initial-state datum enters the
satisfaction judgment in hypothesis position, not the formula ---
so one convention serves stochastic kernels, demonic
branching, and worst-case initial conditions without a case split.
$\mathsf{SProbLE}(q, \psi)$ is the second orientation \emph{within}
the species, not a mirror species: the \emph{upper} probability of
the same compiled event is at most $q$. It is a closed constraint
whose wall is a citizen of the semantic universe like any other, and
emphatically \emph{not} the negation of the GE face --- that is the
strict $\mathsf{P}_{< p}$, which rightly remains a compiler artifact
(the fences below keep that absence). The $\mathsf{MayBe}$ pair
carries the \emph{possibility marks}: some admitted law attains the
stated side of the threshold --- the credal-existential face of the
atom, a liveness demand entering on the marks rather than the floor.

What the earlier edition advertised as nesting is stated, under the
grid, at its true stratum: the tower discipline is
\emph{certificate-mediated}. A child claim is discharged as its own
satisfaction judgment, and its \emph{certified region} enters the
parent formula as an ordinary $\mathsf{StatePred}$ atom --- geometry
with a certificate behind it, not a formula inside a formula.
Syntactic nesting of cap formers (a probability atom inside another's
path argument) is not formable --- Path-row formers take
$\mathsf{StatePred}$ arguments --- and the certified corpus never
spelled such a term: its one advertised nesting use, the
reach-then-hold funnel, already routes the child claim through the
judgment stratum (section~\ref{sec:spec-lang-dfunnel} quotes it).
The grammar now matches what the tower always did.

The family is a family of \emph{aliases}: an event probability is
the expectation of an indicator gamble, so orientation for
orientation
\[
\mathsf{SProbGE}(p, \psi) \;=\;
\mathsf{SReturnGE}(p,\, \mathsf{retPath}\ \psi),
\]
and likewise for the LE face and both marks
(section~\ref{sec:spec-lang-return}). The alias ties are
definitional-grade denotation lemmas at the mechanized stratum,
never schema equations: the surface constructors remain distinct
hash-consed syntax, and the elaboration is computed where consumed.

The interval premises are \emph{canonicity rules, not semantic
fences} --- four instances of one rule: the vacuous end of either
orientation re-spells a structural part of the atom, and canonical
form forbids re-spelling. At $p = 0$ the GE floor (and at $q = 1$
the LE floor) would compile to the unit atom, which is exactly
$\mathsf{STrue}$'s compile (\texttt{SpecDnf.modelClass\_top}); at
$p = 0$ the $\mathsf{MayBeGE}$ mark reads ``some admitted law with
$\mu(\psi) \ge 0$'' --- the structural properness mark itself
(\texttt{SpecAtom.top\_mem}), as does the $\mathsf{MayBeLE}$ mark at
$q = 1$. Bottom-refutes was never carried by threshold positivity:
the credal-lower reading is vacuous at the zero lift for every
constraint (\texttt{satCredal\_lower\_singletonC\_zero}), and it is
the properness presupposition that excludes the bricked system.

\emph{Compile.} The GE face: the region of states from which the
credal-lower probability of the compiled path event clears the
threshold --- mechanized for the kernel path measure as
\texttt{probThresholdSet}, entering the judgment as that region's
safety atom. Adequacy is \emph{directed}: region inclusion
transports to class inclusion
(\texttt{safetyAtom\_modelClass\_mono}, the wing's directed-cell
template; \texttt{sprobAtom\_modelClass\_mono}), and each consuming
battery discharges the sound direction --- the Ville~\cite{ville1939}
and optional-stopping arguments behind the safety, deadline, until,
and Streett batteries of chapter~\ref{ch:qtsl-manual}, per that
catalogue's receipt columns. One end-to-end specimen is
machine-checked at the behavioural tier: a quantitative Streett
certificate bundle concludes membership in the probability-threshold
region of the acceptance event
(\texttt{sprob\_sound\_of\_streettBundle}). No completeness is
claimed anywhere in this family: a certificate can be arbitrarily
loose. The LE face enters as the same wall species at the opposite
orientation --- the checker-anchor restriction theorem already lists
\emph{both} threshold orientations as closed half-spaces of laws, and
the negation-fence remark names the orientation pair the sound
residue of the failed duality. The marks enter on the atom's mark list --- the
mechanized \texttt{Dip}: a strict undercut somewhere admitted.
Formation is open on all four faces; \emph{consumption} of the LE
face at interior $q$ keeps its over-approximating-tower fence
(section~\ref{sec:spec-lang-fpupper}), while $q = 0$ discharges on
the exact lane today: an exact backward-induction table is
simultaneously an under- and over-approximation, so the
possibilistic prohibition $\overline{P}[\psi] = 0$ --- the semantic
universe's own behavioural-prohibition face --- is certified by the
same table read in the super-solution direction. Two compatibility
facts anchor the credal reading at the precise end: at the
singleton lift of a nonzero law, the lower reading of an event atom
is that law's own event probability
(\texttt{satCredal\_lower\_singletonC\_geC}); and the zero-lift
vacuity above.

\emph{Witness}: the supermartingale corpus end to end; the
reach-by-deadline vector quoted below is the certified-Gymnasium
hitting shape ($\mathbb{P}[\mathrm{goal\ by\ }20] \ge 3/4$); the
exact-zero avoidance claims of the deterministic corpus spell
natively on the LE face.

\ruleblock{API}
\apilang{TypeScript}{spec.probGe / probLe / probMayBeGe / probMayBeLe}
\apiprov{signature --- \texttt{ts/src/spec.ts}}
\begin{tssrc}
export function probGe(p: QLike, path: PathPred): ForallSpec

export function probLe(q: QLike, path: PathPred): ForallSpec

export function probMayBeGe(p: QLike, path: PathPred): ExistsSpec

export function probMayBeLe(q: QLike, path: PathPred): ExistsSpec
\end{tssrc}
\apiprov{construction --- \texttt{ts/tests/spec-native.test.ts}}
\begin{tssrc}
const goal = s.polyGe(1n, s.unwrap(s.polyVar(0, 2)));
const cap = s.probGe(Q.of(3, 4), s.deadline(20n, goal));
\end{tssrc}
\apiprov{The \textsc{MayBe} marks construct in the marks vector quoted under the seat-regret entry (the same test file).}
\apilang{Rust}{mtk\_catalogue::spec::ForallSpec::\{prob\_ge, prob\_le\}}
\apiprov{signature --- \texttt{crates/mtk-catalogue/src/spec/formula.rs}}
\begin{rustsrc}
pub fn prob_ge(p: WitRat, path: PathPred) -> Result<ForallSpec, Refusal>

pub fn prob_le(q: WitRat, path: PathPred) -> Result<ForallSpec, Refusal>

pub fn prob_may_be_ge(p: WitRat, path: PathPred) -> Result<ExistsSpec, Refusal>

pub fn prob_may_be_le(q: WitRat, path: PathPred) -> Result<ExistsSpec, Refusal>
\end{rustsrc}
\apiprov{construction --- \texttt{crates/mtk-catalogue/tests/golden\_spec.rs}}
\begin{rustsrc}
let goal = StatePred::poly_ge(q(1, 1), Poly::var(0, 2).unwrap());
let cap =
    ForallSpec::prob_ge(q(3, 4), PathPred::deadline(q(20, 1), goal, false).unwrap()).unwrap();
\end{rustsrc}
\apiprov{The \textsc{MayBe} marks construct in the marks vector quoted under the seat-regret entry (the same test file).}

\subsection{SReturnGE, SReturnLE, and the possibility marks}
\label{sec:spec-lang-return}

\ruleblock{Formation}
\begin{mathpar}
\inferrule*[Right=\textsc{SReturnGE}]
  {B \in \mathbb{Q} \\ R\ \mathsf{PathObservable}}
  {\mathsf{SReturnGE}(B, R)\ \mathsf{ForallSpec}}
\and
\inferrule*[Right=\textsc{SReturnLE}]
  {B \in \mathbb{Q} \\ R\ \mathsf{PathObservable}}
  {\mathsf{SReturnLE}(B, R)\ \mathsf{ForallSpec}}
\\
\inferrule*[Right=\textsc{SReturnMayBeGE}]
  {B \in \mathbb{Q} \\ R\ \mathsf{PathObservable}}
  {\mathsf{SReturnMayBeGE}(B, R)\ \mathsf{ExistsSpec}}
\\
\inferrule*[Right=\textsc{SReturnMayBeLE}]
  {B \in \mathbb{Q} \\ R\ \mathsf{PathObservable}}
  {\mathsf{SReturnMayBeLE}(B, R)\ \mathsf{ExistsSpec}}
\end{mathpar}

The \emph{soft} former: the prevision square at a run-functional. The floors
read the paired System's credal body at the gamble $R$ ---
credal-lower expectation at least $B$ (\textsc{GE}), credal-upper
expectation at most $B$ (\textsc{LE}) --- and the marks are the
credal-existential faces: some admitted law attains the stated side.
Nothing on the spec side carries a set of models: under the
one-universe reading the System's own denotation \emph{is} the
weighted set, and the spec carries only the pair $(R, B)$. This is
where minimax weighted expected regret lands without new vocabulary:
a weighted set of probability measures~\cite{halpern13} is, member
for member, a set of sub-probability laws ($\alpha\Pr$ for weight
$\alpha$), i.e.\ a credal body in this universe's own imprecise
atlas; the sup-normalization convention (some member carries weight
one) is the properness presupposition; and the weighted expected
regret of a nonnegative run-functional is the body's upper
prevision. ``WER
within $[B_1, B_2]$'' is therefore the two-component record
$\mathsf{Spec}(\mathsf{SReturnLE}(B_2, R);\
\mathsf{SReturnMayBeGE}(B_1, R))$
--- one floor, one mark, the same denotation the earlier edition
spelled as a state-level conjunction, with the face of each
component now visible in its sort; no interval sort exists:
ranges are orientation pairs here as everywhere in this catalogue.
Bounds are unconstrained rationals --- returns are not
probabilities, and the vacuous-end canonicity rule is a
probability-face fact.

\emph{Compile.} Credal-lower/upper expectation rows over the
declared reward, lifted at the measurable stratum; the
finite-horizon summed channel form is mechanized as
\texttt{returnGESet} (natural-number horizon, as the split now
forces by shape), entering the judgment as that region's safety
atom, with region monotonicity in reward, horizon, and bound as the
directed cells (\texttt{returnGESet\_mono\_reward},
\texttt{returnGEAtom\_modelClass\_mono\_reward}). Adequacy is
\emph{directed}, certificate to claim, and the two square families
carrying it are machine-checked at the abstract behavioural tier: a
geometrically decaying Lyapunov certificate --- drift factor $\beta$
off the target, barrier $a$ --- yields the expected indicator-return
bound $1 - \beta^{n}\,(\int L_0)/a$ at horizon $n{+}1$
(\texttt{returnGE\_sound\_of\_deadline}); and a geometric-drift
certificate bounds the discounted expected return from above,
uniformly in the horizon, under the spectral condition
$\gamma\beta < 1$ (\texttt{discountedReturn\_le\_of\_drift}, with
\texttt{integral\_le\_pow\_mul} the iterated-drift engine). Both
squares are proved on an abstract filtered probability space,
where the certificate's conditional-drift data lives; the
specialization to the kernel path measure is a separate, named
stratum --- an obligation this edition names rather than claims ---
and the transfer from the deadline kernel to the discounted return
forms climbs the same ladder. The marks enter on the atom's mark
list (\texttt{Dip}), exactly as the probability marks do; the
long-run-average battery (gain/bias) is design vocabulary of this
edition, named with its species and not yet claimed. This family is
the catalogue's highest-priority former because of what
it retired: until it existed, expected-return certificates carried
their claims in header sidecars.

\emph{Witness for the demand}: twenty-five tabular expected-return
certificates ($\mathbb{E}_\pi[\sum r] \ge B$ --- CliffWalking $-13$
exactly; Taxi $793/100$; Blackjack
$-3954908734551683826282634489/13^{26}$, the corpus's own bignum,
quoted below); the pendulum swing-up worst-case return row ($\ge
-619.99$ from every default-reset initial state); the Atari circuit
rows (min-over-lanes return, exact); the LunarLander statistical rows
(via \textsc{QSatConf} below).

\ruleblock{API}
\apilang{TypeScript}{spec.returnGe / returnLe / returnMayBeGe / returnMayBeLe}
\apiprov{signature --- \texttt{ts/src/spec.ts}}
\begin{tssrc}
export function returnGe(lo: QLike, ret: PathObservable): ForallSpec

export function returnLe(hi: QLike, ret: PathObservable): ForallSpec

export function returnMayBeGe(lo: QLike, ret: PathObservable): ExistsSpec

export function returnMayBeLe(hi: QLike, ret: PathObservable): ExistsSpec
\end{tssrc}
\apiprov{construction --- \texttt{ts/tests/spec-native.test.ts}}
\begin{tssrc}
const ret = s.poSum(26, s.oChan(0));
const floor = s.returnGe(
  Q.of(-3954908734551683826282634489n, 13n ** 26n),
  ret,
);
\end{tssrc}
\apiprov{The accumulator injections are quoted under the Return stratum; the \texttt{poInt} / \texttt{poDisc} faces and the possibility marks ride the formers-tour vector (the same test file).}
\apilang{Rust}{mtk\_catalogue::spec::ForallSpec::\{return\_ge, return\_le\}}
\apiprov{signature --- \texttt{crates/mtk-catalogue/src/spec/formula.rs}}
\begin{rustsrc}
pub fn return_ge(lo: WitRat, ret: PathObservable) -> ForallSpec

pub fn return_le(hi: WitRat, ret: PathObservable) -> ForallSpec

pub fn return_may_be_ge(lo: WitRat, ret: PathObservable) -> ExistsSpec

pub fn return_may_be_le(hi: WitRat, ret: PathObservable) -> ExistsSpec
\end{rustsrc}
\apiprov{construction --- \texttt{crates/mtk-catalogue/tests/golden\_spec.rs}}
\begin{rustsrc}
let thirteen_pow_26 = (0..26).fold(BigInt::from(1), |a, _| a * 13);
let bound = WitRat::new(
    BigInt::parse_bytes(b"-3954908734551683826282634489", 10).unwrap(),
    thirteen_pow_26,
);
let ret = PathObservable::sum(26, StateObservable::chan(0)).unwrap();
let floor = ForallSpec::return_ge(bound, ret);
\end{rustsrc}
\apiprov{The accumulator injections are quoted under the Return stratum; the \texttt{poInt} / \texttt{poDisc} faces and the possibility marks ride the formers-tour vector (the same test file).}

\subsection{SLevelGE, SLevelLE, and the level marks}
\label{sec:spec-lang-level}

\ruleblock{Formation}
\begin{mathpar}
\inferrule*[Right=\textsc{SLevelGE}]
  {B \in \mathbb{Q} \\ R\ \mathsf{PathObservable} \\
   \mathit{thr} \in \mathbb{Q}_{\ge 0} \cup \{\infty\}}
  {\mathsf{SLevelGE}(B, R, \mathit{thr})\ \mathsf{ForallSpec}}
\\
\inferrule*[Right=\textsc{SLevelLE}]
  {B \in \mathbb{Q} \\ R\ \mathsf{PathObservable} \\
   \mathit{thr} \in \mathbb{Q}_{\ge 0} \cup \{\infty\}}
  {\mathsf{SLevelLE}(B, R, \mathit{thr})\ \mathsf{ForallSpec}}
\\
\inferrule*[Right=\textsc{SLevelMayBeGE}]
  {B \in \mathbb{Q} \\ R\ \mathsf{PathObservable} \\
   \mathit{thr} \in \mathbb{Q}_{\ge 0} \cup \{\infty\}}
  {\mathsf{SLevelMayBeGE}(B, R, \mathit{thr})\ \mathsf{ExistsSpec}}
\\
\inferrule*[Right=\textsc{SLevelMayBeLE}]
  {B \in \mathbb{Q} \\ R\ \mathsf{PathObservable} \\
   \mathit{thr} \in \mathbb{Q}_{\ge 0} \cup \{\infty\}}
  {\mathsf{SLevelMayBeLE}(B, R, \mathit{thr})\ \mathsf{ExistsSpec}}
\end{mathpar}

The \emph{hard} former: for every full-mass law whose level of
inconsistency is finite and at most $\mathit{thr}$, the expected value of $R$ lies on
the stated side of $B$. The threshold lives in the extended
nonnegative rationals --- syntactically a two-constructor sum
(finite bound; infinity), so every consumer's finite/instant case
split is carried by shape --- and the finiteness clause is one
uniform formula, not a case split: redundant at finite thresholds,
it does all the work at $\mathit{thr} = \infty$, where the
quantifier ranges over the \emph{support}, every behaviour not
fully refused. That instance is the behavioural-prohibition face of
the satisfaction judgment itself: take $R$ the indicator of the bad
event via $\mathsf{retPath}$, $\mathit{thr} = \infty$, $B = 0$ on
the LE side, and the probe says the upper probability of the bad
event is zero. At $\mathit{thr} = 0$ the quantifier reads the
zero-inconsistency core. Every rung of the ladder is a lawful probe; no
positivity premise exists beyond the threshold's own nonnegativity.

The $\mathsf{MayBe}$ pair completes the family's polarity: \emph{some}
full-mass admitted law within the level threshold attains the
stated side of $B$ --- the level stratum's possibility marks, new
with this edition, read exactly as the probability and expectation
marks (at $\mathit{thr} = \infty$ the mark asserts attainment by some
behaviour-law in the support). With them the mark roster of
section~\ref{sec:spec-lang-formation} is polarity-total: every
threshold family presents both faces.

Soft and hard are not rivals but the two dualities the semantic
universe already stratifies. The soft former's single-shot linear
tests are a complete instrument, by Fenchel--Moreau duality, exactly
on the convex composition core; level probes recover the whole
quasiconvex blanket, by Penot--Volle level-set duality --- and the
gap has a witness: a quantile (value-at-risk) specification has
every sublevel a half-space, hence is hard-expressible
tolerance-by-tolerance, and is not convex, hence invisible to the
soft former alone. The hinge between them has a machine-checked
ancestor over the retired calculus-tier carrier: there, the
finite-tilt family converges in order to the wall at the infinite
endpoint (\texttt{livenessTilt\_top} in the archived
\texttt{su-functor-legacy} tree, with epi-convergence), which is
this pair of formers stated at the threshold stratum; its
restatement over the live semantic universe is owed alongside the
level compile itself. The hard former's completeness, by contrast,
is a theorem the live universe already holds: a single instance
denotes a rational-height gamble row, and finite conjunctions of
exactly these, closed under directed suprema, generate every
citizen of the universe
(\texttt{General.iSup\_gambleRows\_eq\_self}) --- the level probe is
the universe's own generation species wearing surface syntax, its
certificate discharge the polynomial-nonnegativity species this
book's rule catalogue already discharges elsewhere, re-aimed.

\ruleblock{API}
The mechanized compile remains owed with the tilt-hinge restatement
above; the producer surface below is landed, all four faces.
\apilang{TypeScript}{spec.levelGe / levelLe / levelMayBeGe / levelMayBeLe}
\apiprov{signature --- \texttt{ts/src/spec.ts}}
\begin{tssrc}
export function levelGe(lo: QLike, ret: PathObservable, thr: EnnRat): ForallSpec

export function levelLe(hi: QLike, ret: PathObservable, thr: EnnRat): ForallSpec

export function levelMayBeGe(
  lo: QLike,
  ret: PathObservable,
  thr: EnnRat,
): ExistsSpec

export function levelMayBeLe(
  hi: QLike,
  ret: PathObservable,
  thr: EnnRat,
): ExistsSpec

export function ennFin(q: QLike): EnnRat

export function ennInf(): EnnRat
\end{tssrc}
\apiprov{construction --- \texttt{ts/tests/spec-native.test.ts}}
\begin{tssrc}
const levelLo = s.levelGe(Q.of(1, 2), s.poAvg(s.oChan(0)), s.ennFin(3n));
const levelHi = s.levelLe(Q.of(7, 2), s.poSup(s.oPoly(x1)), s.ennInf());
\end{tssrc}
\apilang{Rust}{mtk\_catalogue::spec::ForallSpec::\{level\_ge, level\_le\}}
\apiprov{signature --- \texttt{crates/mtk-catalogue/src/spec/formula.rs}}
\begin{rustsrc}
pub fn level_ge(lo: WitRat, ret: PathObservable, thr: EnnRat) -> ForallSpec

pub fn level_le(hi: WitRat, ret: PathObservable, thr: EnnRat) -> ForallSpec

pub fn level_may_be_ge(lo: WitRat, ret: PathObservable, thr: EnnRat) -> ExistsSpec

pub fn level_may_be_le(hi: WitRat, ret: PathObservable, thr: EnnRat) -> ExistsSpec

pub fn fin(q: WitRat) -> Result<EnnRat, Refusal>

pub fn inf() -> EnnRat
\end{rustsrc}
\apiprov{construction --- \texttt{crates/mtk-catalogue/tests/golden\_spec.rs}}
\begin{rustsrc}
let level_lo = ForallSpec::level_ge(
    q(1, 2),
    PathObservable::avg(StateObservable::chan(0)),
    EnnRat::fin(q(3, 1)).unwrap(),
);
let level_hi = ForallSpec::level_le(
    q(7, 2),
    PathObservable::sup(StateObservable::poly(x1())),
    EnnRat::inf(),
);
\end{rustsrc}

\subsection{SRateLE and SRateGE: event-rate envelopes}
\label{sec:spec-lang-rate}

\ruleblock{Formation}
\begin{mathpar}
\inferrule*[Right=\textsc{SRateLE}]
  {\rho \in \mathbb{Q}_{\ge 0} \\ \sigma \in \mathbb{Q}_{\ge 0} \\
   R\ \mathsf{TrList} \\ R \neq \mathsf{nil}}
  {\mathsf{SRateLE}(\rho, \sigma, R)\ \mathsf{ForallSpec}}
\\
\inferrule*[Right=\textsc{SRateGE}]
  {\rho \in \mathbb{Q}_{\ge 0} \\ \sigma \in \mathbb{Q}_{\ge 0} \\
   R\ \mathsf{TrList} \\ R \neq \mathsf{nil}}
  {\mathsf{SRateGE}(\rho, \sigma, R)\ \mathsf{ForallSpec}}
\end{mathpar}

The from-init affine count envelope --- the rate axis's first
expectation-class citizens. Write $N^R_t$ for the number of firings
of the roster's transitions on $[0, t]$: the same firing-metered
counting the transition-count functional reads
(section~\ref{sec:spec-lang-returnstratum}), here horizon-parametric
rather than whole-run, and consumed directly --- the envelope is a
family of bounds indexed by the horizon, not a single
run-functional, which is why these formers take the roster itself.
$\mathsf{SRateLE}$ claims that along every admitted law the expected
count respects the envelope at every horizon --- credal-upper
expectation of $N^R_t$ at most $\rho t + \sigma$ for all $t \ge 0$
--- and $\mathsf{SRateGE}$, the throughput face, that the
credal-lower expectation dominates $\rho t - \sigma$. Both faces are
floors at the $\mathsf{ForallSpec}$ cap: each bounds the whole
credal body, and the claim direction selects the envelope --- an
upper count bound reads the upper prevision (the most events any
compatible resolution produces), a lower bound the lower (the
fewest). Both are certification pessimism, and rate atoms carry no
attitude knob: the direction is the whole story. The possibility
faces are deliberately absent, a recorded choice rather than an
oversight: a mark-face rate envelope --- \emph{some} admitted law's
expected counts obey the envelope --- has no certificate row family
and no consumer, and the mark roster's polarity-totality is a fact
about the threshold families, which the envelope pair is not one
of; a later widening stays additive.

The affine-from-init shape is the primitive on purpose. A bare
long-run limit ($\limsup_t N_t/t \le \rho$) carries zero
finite-window content --- a consumer sized against it can be
overwhelmed on any finite horizon --- so the limit reading is a
corollary of the envelope
($\overline{E}[N_t]/t \le \rho + \sigma/t$), never the interface
object. And interval content is carried by \emph{pairs}: a
two-sided pair at matched rate --- upper $(\rho, \sigma_u)$, lower
$(\rho, \sigma_l)$ --- yields the interval envelope
$N_{(s,t]} \le \rho(t - s) + \sigma_u + \sigma_l$ by pathwise
subtraction, no ergodic machinery, which is why the primitive can
stay from-init without losing burstiness control: a one-sided
envelope says nothing about increments; a pair does. The remaining
free laws are of the same grain: weakening an envelope weakens its
claim (the upper face is monotone in both parameters; the lower
face trades them the other way), and disjoint rosters sum ---
upper envelopes over disjoint rosters conjoin to an upper envelope
at the summed parameters on the merged roster, pathwise count
additivity plus subadditivity of the upper prevision --- while the
converse split is refused: a summed bound cannot be divided
between rosters without a per-roster certificate, since all the
mass may ride one of them.

Refusals, each with its reason. No limsup or ergodic primitive
(the zero-finite-window-content argument above --- and a fortiori
no almost-sure empirical-rate convergence, a law of large numbers
needing recurrence structure the vocabulary deliberately does not
assume). Nothing sharper than the pair-derived interval form:
equilibrium-rate claims and interval budgets below
$\sigma_u + \sigma_l$ genuinely need that same ergodic machinery
--- the fence the occupation-time operators sit behind
(section~\ref{sec:spec-lang-focc}). Counted events are declared
transitions only, never predicate crossings: a full-rank diffusion
crosses a surface infinitely often, so region-entry counting stays
\textsc{PCrossStreett}'s metered device
(section~\ref{sec:spec-lang-streett}) --- a different species, not
interchangeable. No within-roster order content: the envelope
counts class firings, and per-item ordering or latency stays
aggregate-invisible by design. No policy-quantified or optimality
rate claims (``no dispatch achieves a better rate'' is a research
frontier, not a former): rate judgments are fixed-dispatch or
demonic, like every judgment in this catalogue. And no attitude
knob, as above. The
assume/guarantee half of the rate axis --- rate bounds crossing
fused boundaries, and the probabilistic token-bucket budget form
that rides fusion --- remains a catalogued deliberate absence with
its own registry slot (section~\ref{sec:spec-lang-frate}).

\emph{Compile.} The certificate species is staged beside the
formers: a supermartingale row kit per face, kin to the drift rows
everywhere above. The upper face's kit carries a per-cell
polynomial potential $V$ with interior rows bounding the
generator's action on $V$ plus the counted transitions' intensity
by $\rho$ on each mode cell, counted-transition guard rows
absorbing the unit charge ($V$ after the reset at most $V - 1$),
un-counted guard rows non-increasing, and nonnegativity on the
reset closure --- making $V + N^R_t - \rho t$ a supermartingale
under every compatible law, so the expected count is at most
$\rho t + V(x)$, envelopes last, with the certified region the
sublevel $\{V \le \sigma\}$. The lower face mirrors it
($V + \rho t - N^R_t$; the counted compensator credits the count's
expected growth), and its potential is the Foster--Lyapunov
throughput shape: a work-in-progress potential whose forced
descent pays out completions. A third, structural route concludes
the upper face from clock structure alone --- race, dwell, and
instant-budget readings of the landed timing stock, no potential
synthesized. No exhibit claims a rate battery yet; formed claims
await theirs, as everywhere in this catalogue.

\emph{Witness for the demand}: the throughput claim of the
Mini-Fab tutorial (section~\ref{sec:minifab-spec},
chapter~\ref{ch:qtsl-minifab}) renders as this former's lower
face; until this family was staged, that chapter could record the
claim only as an argument beside the vocabulary.

\ruleblock{API}
\apilang{TypeScript}{---}
The theory rows are staged; the producer surface enters with the
toolkit's next revision, and this entry gains its quoted
constructions then.
\apilang{Rust}{---}
Likewise.

\subsection{SRegretLE and SRegretLEMark: seat regret}
\label{sec:spec-lang-regret}

\ruleblock{Formation}
\begin{mathpar}
\inferrule*[Right=\textsc{SRegretLE}]
  {\varepsilon \in \mathbb{Q} \\ \varepsilon \ge 0 \\
   \pi\ \mathsf{System} \\ R\ \mathsf{PathObservable}}
  {\mathsf{SRegretLE}(\varepsilon, \pi, R)\ \mathsf{ForallSpec}}
\\
\inferrule*[Right=\textsc{SRegretLEMark}]
  {\varepsilon \in \mathbb{Q} \\ \varepsilon \ge 0 \\
   R\ \mathsf{PathObservable}}
  {\mathsf{SRegretLEMark}(\varepsilon, R)\ \mathsf{ExistsSpec}}
\end{mathpar}

The subject of a regret claim is a \emph{comb}: a composite with
exactly one seat --- a hole presenting a declared interface, in the
system side's own vocabulary
(chapter~\ref{ch:qtsl-mathematical-models}) --- and the comparison
class is the seat itself: every filler presenting that interface.
The comparison class is read off the \emph{subject}, never
re-declared per formula --- re-declaration would double-count it,
and silent demonic closure of an unfilled hole under an ordinary
claim is exactly what the pairing-relative demands refuse. Achieved
and comparison values are credal-\emph{lower} expectations of the
run-functional on the respective closed composites: residual
nondeterminism stays demonic on both sides uniformly --- maximin
regret, deliberately \emph{not} online-learning pathwise regret,
which needs a shared weather trace no compositional semantics
supplies. The certificate face carries its reference filler: no
type-correct filler beats $\pi$ by more than $\varepsilon$, each
judged by its own guaranteed value. The query face carries none ---
it is well-formed in total ignorance of the answer, and its
discharging selection witness \emph{carries} the filler: the
angelic idiom, in which the witness is the answer. The two faces
are the floor/mark duality one level up, with the quantifier moved
from the law seat to the filler seat --- which is why the name is
\textsc{Mark} and not \textsc{MayBe}: \textsc{MayBe} is the
law-seat existential, and the two seats have just been separated.
Regret is one-sided by definition (no GE mirror: a
``regret at least'' face would be a liveness claim about the seat's
\emph{spread} --- a different former, deliberately not minted), and
$\varepsilon = 0$ is lawful and canonical: the exact-optimality
claim. Pairing-relative demands --- the paired subject carries
exactly one hole; the certificate face's filler fills it --- mint
at the certificate stratum, where the pairing is in hand.

\emph{Witness for the demand}: the certified-Gymnasium
``optimal, exactly'' family --- optimality claims the earlier
edition could state only as bare return bounds against a known
optimum; the discharge lanes are named with their species, not yet
claimed.

\ruleblock{API}
\apilang{TypeScript}{spec.regretLeMark}
\apiprov{signature --- \texttt{ts/src/spec.ts}}
\begin{tssrc}
export function regretLeMark(eps: QLike, ret: PathObservable): ExistsSpec
\end{tssrc}
\apiprov{construction --- \texttt{ts/tests/spec-native.test.ts}}
\begin{tssrc}
const markGe = s.probMayBeGe(Q.of(1, 2), s.always(band));
const markRet = s.returnMayBeLe(0n, s.poSum(7, s.oChan(0)));
const markLevel = s.levelMayBeGe(Q.of(1, 4), s.poAvg(s.oChan(1)), s.ennFin(1n));
const regretMark = s.regretLeMark(
  Q.of(1, 100),
  s.poDisc(s.oChan(0), Q.of(1, 20)),
);
\end{tssrc}
\apiprov{The safety face \textsc{SRegretLE} has no producer former yet; the \textsc{Mark} is the landed surface, and the marks vector quoted here also drives the probability, return, and level \textsc{MayBe} marks.}
\apilang{Rust}{mtk\_catalogue::spec::ExistsSpec::regret\_le\_mark}
\apiprov{signature --- \texttt{crates/mtk-catalogue/src/spec/formula.rs}}
\begin{rustsrc}
pub fn regret_le_mark(eps: WitRat, ret: PathObservable) -> Result<ExistsSpec, Refusal>
\end{rustsrc}
\apiprov{construction --- \texttt{crates/mtk-catalogue/tests/golden\_spec.rs}}
\begin{rustsrc}
let mark_ge = ExistsSpec::prob_may_be_ge(q(1, 2), PathPred::always(band.clone())).unwrap();
let mark_ret = ExistsSpec::return_may_be_le(
    q(0, 1),
    PathObservable::sum(7, StateObservable::chan(0)).unwrap(),
);
let mark_level = ExistsSpec::level_may_be_ge(
    q(1, 4),
    PathObservable::avg(StateObservable::chan(1)),
    EnnRat::fin(q(1, 1)).unwrap(),
);
let regret_mark = ExistsSpec::regret_le_mark(
    q(1, 100),
    PathObservable::disc(StateObservable::chan(0), q(1, 20)).unwrap(),
)
.unwrap();
\end{rustsrc}
\apiprov{The safety face \textsc{SRegretLE} has no producer former yet; the \textsc{Mark} is the landed surface, and the marks vector quoted here also drives the probability, return, and level \textsc{MayBe} marks.}

\section{Derived Spellings}
\label{sec:spec-lang-derived}

Each entry here is a claim species from the certified corpus with its
spelling in the grammar above --- the catalogue's evidence that
\emph{no further former is needed} for it. The helpers add no
expressive power; where a helper has code, its definition is quoted,
because the definition \emph{is} the content. Two rows deliberately
have no code in this wing, and say where their content lives instead.
Each spelling owes the compile a coherence cell --- the definitional
rider saying the sugar's class is the spelled form's class. The
first cells of this one-line class are machine-checked at the record
site: bounded reach is trivially-safe until
(\texttt{untilOn\_univ\_safe}), and the recurrence sugar's single
trivial-request pair degenerates to reachability of the response
region, conditioned on occupation, with the empty record exhibited
as the one stage set where the readings part
(\texttt{streettOn\_singleton\_univ\_iff},
\texttt{streettOn\_singleton\_univ\_empty}). The remaining
spellings' riders are named obligations of the same class, to land
with the compile they annotate.

\subsection{Capture within N steps}
\label{sec:spec-lang-dcapture}

\ruleblock{Spelling}
\[
\mathrm{capture{-}within}(N, \varphi) \;=\;
\mathsf{P}_{\ge 1}[\mathsf{F}_{\le N}\,\varphi]
\]

The credal view makes this worst-case over demonic branching and
initial conditions --- no dedicated former, no worst-case variant.
\emph{Witness}: the pendulum swing-up row --- capture within 139
steps from every default-reset initial state.

\ruleblock{API}
\apilang{TypeScript}{spec.derived.captureWithin}
\apiprov{signature --- \texttt{ts/src/spec.ts}}
\begin{tssrc}
captureWithin(steps: QLike, capture: StatePred): ForallSpec
\end{tssrc}
\apiprov{definition --- \texttt{ts/src/spec.ts}}
\begin{tssrc}
/** Capture within `steps` — `P_{>=1}[F_{<=steps} capture]` as a
 * safety-face component (the credal stratum makes this worst-case
 * over demonic branching/init). */
captureWithin(steps: QLike, capture: StatePred): ForallSpec {
  return probGe(1n, deadline(steps, capture));
},
\end{tssrc}
\apilang{Rust}{mtk\_catalogue::spec::derived::capture\_within}
\apiprov{signature --- \texttt{crates/mtk-catalogue/src/spec/derived.rs}}
\begin{rustsrc}
pub fn capture_within(steps: WitRat, capture: StatePred) -> Result<ForallSpec, Refusal>
\end{rustsrc}
\apiprov{definition --- \texttt{crates/mtk-catalogue/src/spec/derived.rs}}
\begin{rustsrc}
/// Wing C: capture within `steps` — `P_{>=1}[F_{<=steps} capture]`
/// (the credal stratum makes this worst-case over demonic
/// branching/init). Witness: the swing-up row — capture within 139
/// steps from every default-reset init.
pub fn capture_within(steps: WitRat, capture: StatePred) -> Result<ForallSpec, Refusal> {
    ForallSpec::prob_ge(rat_int(1), PathPred::deadline(steps, capture, false)?)
}
\end{rustsrc}

\subsection{Reach-then-hold (the funnel)}
\label{sec:spec-lang-dfunnel}

\ruleblock{Spelling}
\[
\mathrm{reach{-}then{-}hold}(p, T, \chi) \;=\;
\mathsf{P}_{\ge p}[\mathsf{F}_{\le T}\,\chi]
\quad\text{where } \chi \text{ is a child claim's certified
invariance region}
\]

The outer transit is spelled here; the citation of the child's
certified region ($\chi$ as a region atom, concluded by a
$\mathsf{P}_{\ge 1}[\mathsf{G}\,\mathit{band}]$ claim) belongs to
the judgment stratum --- the tower discipline again. \emph{Witness}: the
swing-up composition (outer transit into the held upright set); the
LunarLander attitude funnel.

\ruleblock{API}
\apilang{TypeScript}{spec.derived.reachThenHold}
\apiprov{signature --- \texttt{ts/src/spec.ts}}
\begin{tssrc}
reachThenHold(p: QLike, t: QLike, heldRegion: StatePred): ForallSpec
\end{tssrc}
\apiprov{definition --- \texttt{ts/src/spec.ts}}
\begin{tssrc}
/** Reach-then-hold (funnel) — `P_{>=p}[F_{<=t} chi]` where `chi` is
 * a child claim's certified invariance region, cited as a region
 * atom (the citation itself rides the judgment stratum). */
reachThenHold(p: QLike, t: QLike, heldRegion: StatePred): ForallSpec {
  return probGe(p, deadline(t, heldRegion));
},
\end{tssrc}
\apilang{Rust}{mtk\_catalogue::spec::derived::reach\_then\_hold}
\apiprov{signature --- \texttt{crates/mtk-catalogue/src/spec/derived.rs}}
\begin{rustsrc}
pub fn reach_then_hold(
    p: WitRat,
    t: WitRat,
    held_region: StatePred,
) -> Result<ForallSpec, Refusal>
\end{rustsrc}
\apiprov{definition --- \texttt{crates/mtk-catalogue/src/spec/derived.rs}}
\begin{rustsrc}
/// Wing C: reach-then-hold (funnel; band-hold) —
/// `P_{>=p}[F_{<=t} chi]` where `chi` is the *certified invariance
/// region* of a child claim (`P_{>=1}[G band]`), cited as a region
/// atom (tower citation: a child's certified region is a parent's
/// atom). This helper spells the outer transit; the citation itself
/// rides the judgment stratum. Witness: the swing-up composition;
/// the LunarLander attitude funnel.
pub fn reach_then_hold(
    p: WitRat,
    t: WitRat,
    held_region: StatePred,
) -> Result<ForallSpec, Refusal> {
    ForallSpec::prob_ge(p, PathPred::deadline(t, held_region, false)?)
}
\end{rustsrc}

\subsection{Safety until deadline}
\label{sec:spec-lang-dsafeuntil}

\ruleblock{Spelling}
\[
\mathrm{safety{-}until{-}deadline}(T, \mathit{safe}, \mathit{target})
\;=\; \mathit{safe}\ \mathsf{U}_{\le T}\ \mathit{target}
\]

Directly the until former --- or, for a probabilistic core, the
two-story tower around it. Listed because the species recurs in the
corpus under this name. \emph{Witness}: the worked example; the
until-battery fixtures.

\ruleblock{API}
\apilang{TypeScript}{spec.derived.safetyUntilDeadline}
\apiprov{signature --- \texttt{ts/src/spec.ts}}
\begin{tssrc}
safetyUntilDeadline(t: QLike, safe: StatePred, target: StatePred): PathPred
\end{tssrc}
\apiprov{definition --- \texttt{ts/src/spec.ts}}
\begin{tssrc}
/** Safety until deadline — `safe U_{<=t} target`, directly. */
safetyUntilDeadline(t: QLike, safe: StatePred, target: StatePred): PathPred {
  return until(t, safe, target);
},
\end{tssrc}
\apilang{Rust}{mtk\_catalogue::spec::derived::safety\_until\_deadline}
\apiprov{signature --- \texttt{crates/mtk-catalogue/src/spec/derived.rs}}
\begin{rustsrc}
pub fn safety_until_deadline(
    t: WitRat,
    safe: StatePred,
    target: StatePred,
) -> Result<PathPred, Refusal>
\end{rustsrc}
\apiprov{definition --- \texttt{crates/mtk-catalogue/src/spec/derived.rs}}
\begin{rustsrc}
/// Wing C: safety until deadline — `safe U_{<=t} target`, directly
/// (or as the probability cap around it, which is a composition of
/// [`ForallSpec::prob_ge`] with this). Witness: the worked example;
/// until-battery fixtures.
pub fn safety_until_deadline(
    t: WitRat,
    safe: StatePred,
    target: StatePred,
) -> Result<PathPred, Refusal> {
    PathPred::until(t, safe, target)
}
\end{rustsrc}

\subsection{Almost-sure recurrence under demonic input}
\label{sec:spec-lang-drecur}

\ruleblock{Spelling}
\[
\mathrm{recurrence{-}a.s.}(\gamma) \;=\;
\mathsf{P}_{\ge 1}[\mathsf{GF}\,\gamma],
\qquad
\mathsf{GF}\,\gamma \;=\;
\mathsf{PDwellStreett}\bigl(\mathsf{PairBase}(\mathsf{STrue}, \gamma)\bigr)
\]

The dwell-metered Streett sugar at the demonic profile (frequently
is step-frequently at unit ticks): the credal view reads the
lower probability over the branch set, so almost-sure recurrence
under an adversarial input needs no game former. \emph{Witness}: the
219 circuit tasks, whose conjunctions of recurrence goals are built
by exactly the quoted spelling.

\ruleblock{API}
\apilang{TypeScript}{spec.derived.recurrenceAs}
\apiprov{signature --- \texttt{ts/src/spec.ts}}
\begin{tssrc}
recurrenceAs(gamma: StatePred): ForallSpec

gf(gamma: StatePred): PathPred

fg(phi: StatePred): PathPred
\end{tssrc}
\apiprov{construction --- \texttt{ts/tests/spec-native.test.ts}}
\begin{tssrc}
const gamma0 = s.unwrap(s.derived.boolWireSet(0, 2));
const gamma1 = s.unwrap(s.derived.boolWireSet(1, 2));
const recur0 = s.derived.recurrenceAs(gamma0);
const recur1 = s.derived.recurrenceAs(gamma1);
\end{tssrc}
\apilang{Rust}{mtk\_catalogue::spec::derived::recurrence\_as}
\apiprov{signature --- \texttt{crates/mtk-catalogue/src/spec/derived.rs}}
\begin{rustsrc}
pub fn recurrence_as(gamma: StatePred) -> Result<ForallSpec, Refusal>

pub fn gf(gamma: StatePred) -> Result<PathPred, Refusal>

pub fn fg(phi: StatePred) -> PathPred
\end{rustsrc}
\apiprov{construction --- \texttt{crates/mtk-catalogue/tests/golden\_spec.rs}}
\begin{rustsrc}
let gamma0 = derived::bool_wire_set(0, 2).unwrap();
let gamma1 = derived::bool_wire_set(1, 2).unwrap();
let recur0 = derived::recurrence_as(gamma0).unwrap();
let recur1 = derived::recurrence_as(gamma1).unwrap();
\end{rustsrc}

\subsection{Min-return over enumerated branches}
\label{sec:spec-lang-dminret}

\ruleblock{Spelling}
\[
\mathrm{min{-}return}(B, R) \;=\;
\mathsf{SReturnGE}(B, R)
\quad\text{at the credal reading}
\]

Lower expectation over a finite branch set \emph{is} the minimum ---
the spelling is the plain former; the worst-case reading belongs to
the substrate, not the syntax. \emph{Witness}: the Atari circuit rows
(per-lane returns, claimed bound equal to the minimum, exact).

\ruleblock{API}
\apilang{TypeScript}{spec.derived.minReturnOverBranches}
\apiprov{signature --- \texttt{ts/src/spec.ts}}
\begin{tssrc}
minReturnOverBranches(lo: QLike, ret: PathObservable): ForallSpec
\end{tssrc}
\apiprov{definition --- \texttt{ts/src/spec.ts}}
\begin{tssrc}
/** Min-return over enumerated branches — `SReturnGE` at the credal
 * stratum (lower expectation over a finite branch set = the
 * minimum); the spelling is the plain former. */
minReturnOverBranches(lo: QLike, ret: PathObservable): ForallSpec {
  return returnGe(lo, ret);
},
\end{tssrc}
\apilang{Rust}{mtk\_catalogue::spec::derived::min\_return\_over\_branches}
\apiprov{signature --- \texttt{crates/mtk-catalogue/src/spec/derived.rs}}
\begin{rustsrc}
pub fn min_return_over_branches(lo: WitRat, ret: PathObservable) -> ForallSpec
\end{rustsrc}
\apiprov{definition --- \texttt{crates/mtk-catalogue/src/spec/derived.rs}}
\begin{rustsrc}
/// Wing C: min-return over enumerated branches — `SReturnGE` at the
/// credal stratum (lower expectation over a finite branch set = the
/// minimum); the *spelling* is the plain former over the given run
/// functional, the credal reading is the substrate's. Witness: the
/// Atari circuit rows (per-lane returns, claimed bound = the
/// minimum).
pub fn min_return_over_branches(lo: WitRat, ret: PathObservable) -> ForallSpec {
    ForallSpec::return_ge(lo, ret)
}
\end{rustsrc}

\subsection{Worst-case-over-init}
\label{sec:spec-lang-dwcinit}

\ruleblock{Spelling}
\[
\text{any former} \;+\; \text{a demonic initial box on the system side}
\]

No spec-side helper, by design: worst-case over initial conditions is
a property of the \emph{pairing}, spelled as a demonic initial box in
the system's own vocabulary
(chapter~\ref{ch:qtsl-mathematical-models}) and read through the
credal view --- the formula layer needs no worst-case former.
\emph{Witness}: the swing-up rows; the random-ODE
uniform-over-parameter barrier (parameter box as an initial-box
axis).

\ruleblock{API}
\apilang{TypeScript}{---}
Spelled on the system side (a demonic initial box declaration), not
in this namespace; the credal reading supplies the worst case.
\apilang{Rust}{---}
Likewise: a system-side declaration, deliberately without a
formula-layer twin.

\subsection{Hard event-count envelopes}
\label{sec:spec-lang-devcount}

\ruleblock{Spelling}
\[
\mathrm{count{-}ceiling}(B, R) \;=\;
\mathsf{SReturnLE}\bigl(B,\, \mathsf{retFireCount}(R)\bigr)
\]

The count ceiling is the plain expectation former at the
transition-count functional
(section~\ref{sec:spec-lang-returnstratum}): along every admitted
law, the total count of roster firings is at most $B$ --- no new
grammar, as this section's charter demands. An earlier
edition of this entry routed the species through the designed
event-rate envelopes as their $p = 1$ instance; the count spine
landed first, the rate family's expectation atoms have since been
staged in their own right (section~\ref{sec:spec-lang-rate}), and
the probabilistic budget form --- the $p = 1$ route --- remains the
deliberate absence it was (section~\ref{sec:spec-lang-frate});
the ceiling stays the plain spelling above either way. \emph{Witness for the
demand}: structurally bounded resource-sharing net families; the
plated-transfer discharge that concludes the ceiling is Rule
Catalogue content.

\ruleblock{API}
\apilang{TypeScript}{spec.returnLe / retFireCount}
The ceiling is the plain composition of landed formers,
\texttt{returnLe(B, retFireCount(R))} --- the transition-count
injection's producer surface, both languages, lives with the run
functionals (section~\ref{sec:spec-lang-returnstratum}).
\apilang{Rust}{mtk\_catalogue::spec::\{ForallSpec::return\_le, PathObservable::fire\_count\}}
Likewise:
\texttt{ForallSpec::return\_le(B, PathObservable::fire\_count(R))},
the empty roster refused at construction and at the decoded-artifact
preflight alike.

\subsection{Boolean-output atoms}
\label{sec:spec-lang-dbool}

\ruleblock{Spelling}
\[
\mathrm{bool{-}wire{-}set}(i) \;=\; \mathsf{SAtomGE}(1, x_i)
\qquad
\mathrm{bool{-}wire{-}clear}(i) \;=\; \mathsf{SAtomLE}(0, x_i)
\]

The $0/1$ convention on Boolean-valued wires: no Boolean atom former
exists or is needed. \emph{Witness}: the circuit suite's atoms
compile to exactly this form; the liveness vector's recurrence goals
are built over these spellings.

\ruleblock{API}
\apilang{TypeScript}{spec.derived.boolWireSet / boolWireClear}
\apiprov{signature --- \texttt{ts/src/spec.ts}}
\begin{tssrc}
boolWireSet(index: number, width: number): Result<StatePred, SpecRefusal>

boolWireClear(index: number, width: number): Result<StatePred, SpecRefusal>
\end{tssrc}
\apiprov{construction --- \texttt{ts/tests/spec-native.test.ts}}
\begin{tssrc}
const gamma0 = s.unwrap(s.derived.boolWireSet(0, 2));
const gamma1 = s.unwrap(s.derived.boolWireSet(1, 2));
const recur0 = s.derived.recurrenceAs(gamma0);
const recur1 = s.derived.recurrenceAs(gamma1);
\end{tssrc}
\apilang{Rust}{mtk\_catalogue::spec::derived::bool\_wire\_set}
\apiprov{signature --- \texttt{crates/mtk-catalogue/src/spec/derived.rs}}
\begin{rustsrc}
pub fn bool_wire_set(index: usize, width: usize) -> Result<StatePred, Refusal>

pub fn bool_wire_clear(index: usize, width: usize) -> Result<StatePred, Refusal>
\end{rustsrc}
\apiprov{construction --- \texttt{crates/mtk-catalogue/tests/golden\_spec.rs}}
\begin{rustsrc}
let gamma0 = derived::bool_wire_set(0, 2).unwrap();
let gamma1 = derived::bool_wire_set(1, 2).unwrap();
let recur0 = derived::recurrence_as(gamma0).unwrap();
let recur1 = derived::recurrence_as(gamma1).unwrap();
\end{rustsrc}

\section{The Rewrite Layer}
\label{sec:spec-lang-rewrite}

\subsection{Rewrite chains, concluding spec equivalence}
\label{sec:spec-lang-rewritelayer}

\ruleblock{Rules}

Eighteen rewrite rules, each a denotation \emph{equality} over
$\omega$-words, with congruence (rewriting at a position) and chain
transport (a sequence of rewrites transports satisfaction verbatim,
in both directions): implication elimination, right distribution,
double negation, the two De~Morgan laws, negation past
$\mathsf{F}$/$\mathsf{G}$/$\mathsf{X}$, and the polarity-absorption
fusion family --- the existential-over-time operators absorb $\vee$
and the universal-over-time operators absorb $\wedge$. Its charter
member is the $\mathsf{GF}$-disjunction fusion
\[
\mathsf{GF}(\varphi \vee \psi) \;=\;
\mathsf{GF}\,\varphi \vee \mathsf{GF}\,\psi
\]
(infinitely-often of a union is a union of infinitely-oftens, by
pigeonhole); its siblings are $\mathsf{F}$ of a disjunction,
$\mathsf{G}$ of a conjunction, $\mathsf{FG}$ of a conjunction, until
with disjoined targets
\[
(\sigma \mathbin{\mathsf{U}} \tau_1) \vee
(\sigma \mathbin{\mathsf{U}} \tau_2) \;=\;
\sigma \mathbin{\mathsf{U}} (\tau_1 \vee \tau_2)
\]
(the target position is existential), until with conjoined safety
arguments
\[
(\sigma_1 \mathbin{\mathsf{U}} \tau) \wedge
(\sigma_2 \mathbin{\mathsf{U}} \tau) \;=\;
(\sigma_1 \wedge \sigma_2) \mathbin{\mathsf{U}} \tau
\]
(the safety position is universal), and the bounded-horizon
($\le T$) siblings of the four fusions that carry a horizon --- the
$\le T$ decorations carry through the witness arithmetic verbatim, while
$\mathsf{FG}$ has no bounded row: its settling time is unbounded by
construction. The absorption boundary is exact, and the fence is
theorem-backed rather than asserted: the four wrong-polarity fusions
--- $\mathsf{F}$ of a conjunction, $\mathsf{G}$ of a disjunction,
and their $\mathsf{GF}$/$\mathsf{FG}$ counterparts --- are
machine-refuted by explicit $\omega$-word counterexample theorems in
both provers (\texttt{no\_f\_and\_fusion}, \texttt{no\_g\_or\_fusion},
\texttt{no\_gf\_and\_fusion}, \texttt{no\_fg\_or\_fusion}). A chain
is a typed citation: rule, position, direction, step by step; it
concludes the judgment $\mathsf{QSpecEquiv}(a, b)$ below. The
concluded judgment is already a row species of the normative
schema; the chain citation's own row species is owed there, and
until it lands a chain is a wire artifact of the producer surface.

All eighteen rules are proved in Lean \emph{and} in HOL4, statement
for statement, with no incomplete proofs; completeness is proven for
the benchmark liveness fragment --- every spec of that suite's
grammar reaches the consumable normal form, witnessed mechanically by
all 219 verbatim specs normalizing and re-validating in an
independent checker. Carrier genericity varies by rule family: eight of
the eighteen rules are per-position Boolean and duality laws whose
equalities hold over any linearly ordered time carrier; the nine
fusion rules beyond $\mathsf{GF}$ --- five unbounded, four
bounded-horizon --- are finite-witness manipulations whose equalities
hold over any unbounded linear time carrier; the
$\mathsf{GF}$-fusion remains the only pigeonhole rule, needing
infinitely-often as cofinal occurrence; the next-step duality carries
the step-structure requirement of its operator, which enters the row
grammar together with that operator. The machine-checked proofs are
the $\omega$-word instances; the carrier-generic restatements are
small named obligations. One further restatement belongs to this
chapter's own compile, and its ground floor is now machine-checked
(\texttt{SuFunctor.Universe.SpecRewrite}): soundness at the universe
means the compiled spec classes of a rule's two sides coincide, and
that warrant is strictly stronger than denotation equality --- a
disjunction and an atom of the same denotation compile to different
classes, because models may hedge
(\texttt{exists\_denEq\_not\_classEq}, the compiled image of the
hedging gap) --- so class-level soundness travels by term structure:
the class semantics is a homomorphism of the normal-form lattice
(commutation, association, distribution, and --- against the hedging
intuition --- absorption, each landed as a class equality), and the
normal-form dualities are sound for the sharper reason that they do
not move the term at all (\texttt{classEquiv\_of\_nnf\_eq}); the
obstruction lives at the region-to-class seam and nowhere else. That
development settles the state stratum. The eighteen temporal rules'
own class-level restatement, one row up over the path formers'
compile, remains an obligation this edition names rather than
claims.

The producer surface \emph{constructs} chains --- it deliberately
does not apply rewrites. The checker replays the chain; a chain the
constructor accepts is a well-formed citation, not a verified
equivalence; an empty chain refuses (reflexivity is not a chain).
Application code is withheld on purpose while the deadline-fragment
deductive calculus (bounded-reach/always equivalences over clocked
classes) remains live research: frozen application rows would re-mint
the sidecar disease one level up.

Cross-carrier \emph{application} is not a rewrite matter at all: a
step-structured spec applied to a continuous-time system is expressed
by composing the system with a sampler mechanism and asserting the
spec of the sampled composite, with claims related across the pair by
the simulation-transport rule of chapter~\ref{ch:qtsl-manual} ---
zero new machinery. The certified LunarLander pair (one value table,
both tick semantics, jitter priced at a factor $2.66$ in
ticks-to-funnel) is the landed witness of this pattern.

\emph{Witness}: 219 rewrite-chain artifacts asserting spec
equivalence for the liveness suite --- the demand that shaped the row
form quoted below.

\ruleblock{API}
\apilang{TypeScript}{spec.rewriteChain}
\apiprov{signature --- \texttt{ts/src/spec.ts}}
\begin{tssrc}
export function rewriteChain(
  a: SpecRef,
  b: SpecRef,
  steps: ReadonlyArray<RewriteStep>,
): Result<Artifact, SpecRefusal>
\end{tssrc}
\apiprov{construction --- \texttt{ts/tests/spec-native.test.ts}}
\begin{tssrc}
const art = s.unwrap(
  s.rewriteChain(j.specRef("spec-as-stated"), j.specRef("spec-normal-form"), [
    {
      rule: "implicationElim",
      at: [{ kind: "specForall", i: 0 }],
      dir: "forward",
    },
    {
      rule: "gfDisjunctionFusion",
      at: [{ kind: "specForall", i: 0 }, { kind: "andArg", i: 1 }],
      dir: "forward",
    },
  ]),
);
\end{tssrc}
\apilang{Rust}{mtk\_catalogue::spec::rewrite::RewriteChainV0::new}
\apiprov{signature --- \texttt{crates/mtk-catalogue/src/spec/rewrite.rs}}
\begin{rustsrc}
pub fn new(a: SpecRef, b: SpecRef, steps: Vec<RewriteStep>) -> Result<Self, Refusal>

pub fn concludes(&self) -> QSpecEquiv
\end{rustsrc}
\apiprov{construction --- \texttt{crates/mtk-catalogue/tests/golden\_spec.rs}}
\begin{rustsrc}
let chain = RewriteChainV0::new(
    SpecRef("spec-as-stated".into()),
    SpecRef("spec-normal-form".into()),
    vec![
        RewriteStep {
            rule: RewriteRule::ImplicationElim,
            at: FormulaPath(vec![Step::SpecForall(0)]),
            dir: Direction::Forward,
        },
        RewriteStep {
            rule: RewriteRule::GfDisjunctionFusion,
            at: FormulaPath(vec![Step::SpecForall(0), Step::AndArg(1)]),
            dir: Direction::Forward,
        },
    ],
)
.unwrap();
let concluded = chain.concludes();
\end{rustsrc}

\section{Deliberate Absences}
\label{sec:spec-lang-fences}

A catalogue that only listed what exists would invite spelling every
missing claim species as prose --- the sidecar disease by another
door. This section records what the formula layer deliberately does
\emph{not} provide, why, and what would unlock each. In rule form an
absence is the absence of a formation rule (or of a compile clause
for a formable term), and several are enforced in code, as
typing-preflight refusals that name their fence; the negation-residue
absences additionally carry the no-expressible-complement warrant of
section~\ref{sec:spec-lang-formation} --- they are theorems about
the class algebra, not prudence. The admission rule stands: each of
these gains its former when its first certificate lane arrives, and
this catalogue is the registry where that admission is recorded.

\subsection{Event-rate assume/guarantee; the probabilistic budget
form}
\label{sec:spec-lang-frate}

\ruleblock{Absence}

The expectation-envelope atoms of the rate family are staged
(section~\ref{sec:spec-lang-rate}); this slot keeps what remains
absent, which is everything that crosses a fused boundary. Rate
\emph{assumptions and guarantees} as interface fields --- a
producer's guaranteed output envelope becoming a fused neighbour's
assumed input envelope, with the transported bound recomputed by
the checker as affine arithmetic in $(\rho, \sigma)$ --- wait on
the event-port vocabulary of the system side and enter with the
first fused-boundary composite, together with their composition
rule and its acyclicity oracle. The \emph{probabilistic} budget
form --- a time-uniform claim that the count never overdraws its
budget, with probability at least $p$ --- is sugar over a
token-bucket monitor mechanism fused into every counted firing,
riding the safety kernel unchanged; it waits on the same fusion
vocabulary, and the hard $p = 1$ envelope with it. Cyclic rate
reasoning across a feedback loop is refused even then: naive
fixpoints are unsound, the sound discharge is induction over time,
and that rule enters on its own warrant or not at all. None of
this fences the staged atoms or the count spine --- what this
absence still fences is boundary-crossing rate reasoning and
probabilistic budgets, not counts and not expectation envelopes.

\ruleblock{API}
\apilang{TypeScript}{---}
No constructor exists for the boundary-crossing forms; the staged
atoms carry the family's admission
(section~\ref{sec:spec-lang-rate}), and this entry is the
residue's registry slot.
\apilang{Rust}{---}
Likewise.

\subsection{Interior upper-bound consumption; Rabin; negated
Streett and until}
\label{sec:spec-lang-fpupper}

\ruleblock{Absence}

Formation of probability upper bounds is open (\textsc{SProbLE}
above): the direct $\mathsf{P}_{\le q}$ claim is a closed constraint
of the same species as the GE face, and fencing its formation as if
it were a complement over-read the negation fence. What remains
absent is narrower, and two different things.
\emph{Consumption} of an upper bound at interior $q$ over an
under-determined model needs certificate towers that
\emph{over}-approximate denotations --- the dual of the entire
under-approximation machinery this stack ships; the exact lane
discharges the $q = 0$ instance today (the same backward-induction
table, read in the super-solution direction), and interior $q$ waits
for the tower. Rabin acceptance (the negation of Streett) and
path-negated Streett or bounded-until remain genuinely fenced by the
negation residue: under the properness presupposition no spec class
has an expressible complement
(section~\ref{sec:spec-lang-formation}), so there is no class-level
object for those dual rules to conclude into --- the missing rules
are missing mathematics. The typing preflight still refuses each
\emph{negation} residue with a path address (a complemented
probability is lawful only as a Streett request consumed through
$\mathsf{PairCmpl}$); the direct LE former needs no refusal, because
it was never a negation.

\ruleblock{API}
\apilang{TypeScript}{spec.probLe}
The LE former constructs (\texttt{probLe}, quoted at its entry); the
preflight refusals (\texttt{negPathAtomEvenPolarity},
\texttt{negStreett}, \texttt{negUntil}) remain the surface for the
negation residues, whose entry points are quoted under \textsc{SNot}.
\apilang{Rust}{mtk\_catalogue::spec::ForallSpec::prob\_le}
Likewise (\texttt{Refusal::NegPathAtomEvenPolarity} and kin, with
formula-path geometry).

\subsection{Stochastic bounded-always at even parity}
\label{sec:spec-lang-falwaysstoch}

\ruleblock{Absence}

The former constructs today (\textsc{PAlwaysT} above); what is absent
is stochastic \emph{consumption} at even parity --- the same
over-approximating tower machinery as the probability upper bounds.
Deterministic consumption is immediate now; the spellability of the
statement was the point of drafting the former ahead of its full
battery.

\ruleblock{API}
\apilang{TypeScript}{---}
The former's constructor is quoted under \textsc{PAlwaysT}; no
additional surface exists for the fenced consumption.
\apilang{Rust}{---}
Likewise.

\subsection{Unbounded quantitative reach}
\label{sec:spec-lang-funbounded}

\ruleblock{Absence}
\[
\mathsf{F}\,\varphi \;\leadsto\;
\mathsf{F}_{\le T}\,\varphi
\quad\text{at a producer-chosen } T,\ \text{flagged
\texttt{producerT}}
\]

No unbounded-reach primitive: the spelling is a producer-chosen
deadline, sound as an under-approximation and \emph{flagged as such}
in the row, so a certificate reader can always tell a bounded claim
made for its own sake from a bounded claim standing in for an
unbounded one. A complete almost-sure-termination battery is the
named unlock for the primitive form.

\ruleblock{API}
\apilang{TypeScript}{spec.derived.eventuallyProducerHorizon}
\apiprov{signature --- \texttt{ts/src/spec.ts}}
\begin{tssrc}
eventuallyProducerHorizon(horizon: QLike, arg: StatePred): PathPred
\end{tssrc}
\apiprov{definition --- \texttt{ts/src/spec.ts}}
\begin{tssrc}
/** The `producerT` spelling of unbounded reach: `F phi` compiled as
 * `F_{<=horizon} phi` at a producer-chosen horizon — a sound
 * under-approximation, flagged as such in the row. */
eventuallyProducerHorizon(horizon: QLike, arg: StatePred): PathPred {
  return deadline(horizon, arg, { producerT: true });
},
\end{tssrc}
\apilang{Rust}{mtk\_catalogue::spec::derived::eventually\_producer\_horizon}
\apiprov{signature --- \texttt{crates/mtk-catalogue/src/spec/derived.rs}}
\begin{rustsrc}
pub fn eventually_producer_horizon(horizon: WitRat, arg: StatePred) -> Result<PathPred, Refusal>
\end{rustsrc}
\apiprov{definition --- \texttt{crates/mtk-catalogue/src/spec/derived.rs}}
\begin{rustsrc}
/// The `producerT` spelling of unbounded reach: `F phi` compiled as
/// `F_{<=horizon} phi` at a producer-chosen horizon — a sound
/// under-approximation, *flagged as such* in the row (Wing F keeps
/// unbounded quantitative reach out of the primitive grammar; S7.3).
pub fn eventually_producer_horizon(horizon: WitRat, arg: StatePred) -> Result<PathPred, Refusal> {
    PathPred::deadline(horizon, arg, true)
}
\end{rustsrc}

\subsection{Occupation-time and duration operators}
\label{sec:spec-lang-focc}

\ruleblock{Absence}

Narrowed by the theory revision: much of what this
absence used to fence is now spellable compositionally. Integrated
occupation time is the time integral of a region's indicator
($\mathsf{poInt}$ over $\mathsf{oInd}$), the long-run occupation
fraction is its average ($\mathsf{poAvg}$ over $\mathsf{oInd}$),
and dwell-metered acceptance is a former of its own
(\textsc{PDwellStreett}). What remains absent is \emph{dedicated}
sojourn and per-visit-duration operators --- claims about
individual visits rather than accumulated occupation --- withheld
pending a kernel theorem: occupation-measure mechanics are the
natural substrate, and a former without its kernel would be a claim
nothing can certify. The compositional spellings above share the
constraint at one remove: their certificate lanes are named, not
yet claimed.

\ruleblock{API}
\apilang{TypeScript}{---}
The compositional spellings construct on the landed injection and
accumulator surfaces (quoted at their entries); the dedicated
operators remain withheld pending the kernel theorem --- nothing
further to construct.
\apilang{Rust}{---}
Likewise.

\subsection{Nested path quantification}
\label{sec:spec-lang-fnest}

\ruleblock{Absence}

The descending grid arrows: path-quantified material re-entering the
atemporal row as a state label --- $\mathbb{P}_{\ge q}[\psi]$ as a
state predicate, a value functional as a state observable --- the
CTL/PCTL nesting style. This is the grid's one deliberate absence
(section~\ref{sec:spec-lang-formation}), and its warrant has two
layers. The \emph{raw} formers are fenced by the geometry: a single
level of re-entry exits the constructible classes (neither face's
rooted class is a crescent), and the negation fence forbids deriving
the universal face from the existential one --- raw nesting is not a
missing feature but a shape the class algebra refuses, and any
future admission must carry both faces as primitives (the analysis
is mechanized in the adequacy layer's path-to-state development).
The \emph{principled} route --- certified-citation formers, the
tower discipline promoted to a formation rule, Observable column
first --- is semantically supported (the behavioural tier already
denotes the descending arrows) and deliberately out of scope for
this edition: planned for the language's second major revision
(2026Q4), a scoped fence chosen with the open pathway in full view.
The certified corpus's nesting needs are served through the judgment
stratum --- a child claim's certified region cited as a parent's
$\mathsf{StatePred}$ atom, the tower discipline (entry
\textsc{SProbGE}) --- so the absence fences only the syntactic
embedding, which no landed artefact ever spelled.

\ruleblock{API}
\apilang{TypeScript}{---}
No constructor exists, deliberately; the certified-citation former
family is scheduled with the second major revision.
\apilang{Rust}{---}
Likewise.

\subsection{Strategy quantification}
\label{sec:spec-lang-fgame}

\ruleblock{Absence}

Controlled and game forms are absent from the formula layer by
design: open nondeterminism is absorbed by the credal reading, and
policy-level quantification is a judgment-layer and system-side
question (chapter~\ref{ch:qtsl-mathematical-models}, the
decision-process rows). The formula layer stays closed-system --- a
strategy quantifier would bind at the wrong stratum.

\ruleblock{API}
\apilang{TypeScript}{---}
No strategy quantifier exists in this namespace, by design.
\apilang{Rust}{---}
Likewise.

\subsection{Delays}
\label{sec:spec-lang-fdelay}

\ruleblock{Absence}

Permanently out of the formula layer: a delay is a model-class change
(history-dependent state), not a formula-class change. If it arrives,
it arrives as a new mechanism species with its own certificate
shapes on the system side --- never as a former here.

\ruleblock{API}
\apilang{TypeScript}{---}
Out of scope for the formula layer, permanently.
\apilang{Rust}{---}
Likewise.

\section{Judgment Kinds}
\label{sec:spec-lang-judgments}

A \emph{judgment kind} is a $\vDash$-side sort --- the type a
certificate's rules conclude into --- and each kind carries its own
small sub-logic: the premises it takes, the machinery that discharges
it (\emph{rule}: a rule node concludes it; \emph{oracle}: a
registry-pinned recompute; \emph{chase}: load-time axioms, in the
row schema's three modes --- chase proper \emph{derives} rows,
report-mode axioms state loud constraints, and refuse-mode axioms
gate admission), and the consumers that cite it. The row schema
reads this section's inventory as three verbs: a revision
\emph{states} a satisfaction claim (a system paired with a spec),
\emph{compiles} it into the two-polarity verdict form (its entry
below), and \emph{discharges} it through battery and rule-node rows
a checker replays. Homing these
at the judgment stratum is a design statement: the formula is about
the system; these are about the claim's kind and evidence. At the
semantic universe the satisfaction sort's meaning is the fibre order
itself (\texttt{SuFunctor.Universe.Sat}: the spec's belief lies
below the system's), with the properness presupposition carried
structurally, per section~\ref{sec:spec-lang-target}; the splice of
the certified checker's acceptance judgment into that order is named
in the assurance ledger rather than claimed here --- its stated
prerequisite, the certificate schema's absorption of this chapter's
redesign, has landed. The citation discipline that every kind below leans on
--- a premise accepts base-tower proof leaves, or
concluded-judgment references of a stated kind, or certificate data,
and never a mixture --- is deliberately not an entry here: it is the
edge discipline whose warrant is the induction over admitted
derivations itself, and it renders at the end of
chapter~\ref{ch:qtsl-manual} as the soundness theorem the kinds
feed.

Producer-side, every kind below is a versioned wire artifact (a
literal schema name and version, exact rationals as canonical text,
unknown fields ignored on read), emitted through the one canonical
encoder. What is deliberately \emph{not} here, recorded so nobody
re-adds it: no second sort for interface-indexed satisfaction (see
the next two entries); no tower-step kind (its content is a region
satisfaction consumed through a premise field); no mixing-rate kind
(named open); no attitude judgment kind (the attitude blend is a
record component that no rule concludes and nothing consumes ---
the closing entry).

\subsection{QSat: root satisfaction}
\label{sec:spec-lang-qsat}

\ruleblock{Formation}
\begin{mathpar}
\inferrule*[Right=\textsc{QSat}]
  {s\ \mathsf{Spec} \\ M\ \mathsf{System} \\
   \mathrm{states}(M) \simeq \text{wire object of } s}
  {M \vDash s\ \ \mathsf{Judgment}}
\end{mathpar}

Root satisfaction, closed \emph{and} interface-indexed: one judgment
kind for both readings. A spec with no interface row is closed; with
one, the assumption region reads input wires, the guarantee reads
output wires, and the spec's interface must equal the system's
presented interface (an agreement axiom, chase-policed at the
pairing). The formerly mooted second sort for interface-indexed
satisfaction is \emph{dissolved}: the closed judgment is the
empty-interface instance, so one kind suffices and the loader cannot
be handed a claim whose reading is ambiguous.

Under the properness presupposition a root satisfaction claim
asserts \emph{safe and proper}: to satisfy anything, first exist.
The conjunct is the judgment's presupposition, carried by the
compile target's type (\texttt{SpecAtom.top\_mem}; emitted
unconditionally, \texttt{elabSpec\_top\_mem}) --- zero syntax cost,
one nonemptiness discharge per root claim, discharged at use sites,
with emptiness certificates the optional schema slot. This is the
spec-side face of the guard whose system-side face reads an
inconsistent system as satisfying all safety specs and no liveness
specs (\texttt{top\_notMem\_modelClass}; ex falso inverts
structurally, never by a rule). Sub-logic: concluded from a
region-satisfaction row plus exact rational evaluation on the
initial region (the satisfaction-from-region rule of
chapter~\ref{ch:qtsl-manual}); transported across wirings with slack
arithmetic by the assume--guarantee composition rules.

One satisfaction-stratum variant deserves its sentence here,
because it is certificate machinery and must not be mistaken for
formula grammar: a \emph{guarded} satisfaction claim conditions the
guarantee on an initial-region antecedent --- demanded of every
admitted law whose initial marginal is supported in the cited
region, with nothing promised off-region. The guard is a predicate
\emph{on} the subject's initial data, not initial data (the record
still binds no model), and the class algebra stays
complement-free: the antecedent is a ground condition on system
data, never a behaviour class, which is why no implication former
exists below the cap. Its introduction rule discharges the
initial-membership hypothesis \emph{into} the claim; its record row
cites the certificate-side region vocabulary and lives with the
certificate sections of the normative theory --- the subject strata
stay certificate-citation-free by the standing deliberate absence
--- and its introduction and elimination are Rule Catalogue
content.

\emph{Witness}: every root certificate; the composition chapters'
worked examples exercise the interface-indexed reading end to end.

\ruleblock{API}
\apilang{TypeScript}{judgment.sat}
\apiprov{signature --- \texttt{ts/src/judgment.ts}}
\begin{tssrc}
export function sat(args: { sys: SysRef; spec: SpecRef }): Result<
  Artifact,
  JudgmentRefusal
>
\end{tssrc}
\apiprov{construction --- \texttt{ts/tests/spec-native.test.ts}}
\begin{tssrc}
s.unwrap(j.sat({ sys: j.sysRef("sys-a"), spec: j.specRef("spec-a") })),
\end{tssrc}
\apilang{Rust}{mtk\_catalogue::judgment::QSat::new}
\apiprov{signature --- \texttt{crates/mtk-catalogue/src/judgment/kinds.rs}}
\begin{rustsrc}
pub fn new(sys: SysRef, spec: SpecRef) -> Self
\end{rustsrc}
\apiprov{construction --- \texttt{crates/mtk-catalogue/tests/roundtrip\_spec.rs}}
\begin{rustsrc}
let sat = QSat::new(SysRef("sys-a".into()), SpecRef("spec-a".into()));
\end{rustsrc}

\subsection{Interface-indexed satisfaction, defined}
\label{sec:spec-lang-vdashi}

\ruleblock{Definition}

For a system $\Sigma$ and an interface spec
$(I, \alpha, \gamma[, \kappa])$ --- assumptions $\alpha$ on input
wires, guarantee $\gamma$ on output wires, optional slack $\kappa$
--- the judgment $\Sigma \vDash_{\mathrm{I}} (I, \alpha, \gamma[,
\kappa])$ asserts that every behaviour of $\Sigma$ whose inputs
conform to $\alpha$ delivers $\gamma$ (with margin per $\kappa$ where
declared); in the stochastic profiles, almost every behaviour, with
failure measured \emph{within the conforming event}:
\[
\forall^{\mu}\omega,\ \omega \in \mathrm{conform} \to \omega \in
\mathrm{sat};
\qquad\text{quantitatively}\quad
\mu(\mathrm{conform} \cap \mathrm{sat}^{c}) \le \varepsilon .
\]
Unlike the closed judgment $\mathcal{M} \vDash s$, which quantifies
unconditionally over all behaviours from the initial region,
$\vDash_{\mathrm{I}}$ quantifies conditionally --- satisfaction
relativized to the interface; the closed judgment is its
empty-interface instance. The subscript is a roman kind-marker, not a
parameter, following the global-stratum marker $\vDash_D$ (entry
\textsc{QMorse}).

This reading is defined once, here, and has \emph{no schema sort of
its own}: its home is $\mathsf{QSat}$ whose spec carries an
interface; assumptions and slack live on the spec record; interface
agreement is the pairing's chase axiom. It earns an entry because it
is a genuinely different satisfaction relation --- conditional, with
failure measured inside the conforming event --- and because the
composition rules' conclusions are exactly of this shape. Its
mechanization status: the sequent substrate for the
satisfaction judgment carries the properness conjunct as a theorem
(the pair-sequent development,
\texttt{SuFunctor.Universe.PairSequent}), and the record-level
compile carries assumptions as data with their conditional semantics
deferred fail-closed (section~\ref{sec:spec-lang-spec}); the
relativized reading stated at the semantic universe is a named
obligation, not yet claimed.

\ruleblock{API}
\apilang{TypeScript}{---}
Carried by \texttt{judgment.sat} with an interface-carrying spec ---
deliberately no separate entry point.
\apilang{Rust}{---}
Carried by \texttt{QSat} likewise; the dissolved second sort is
recorded in the module documentation so nobody re-adds it.

\subsection{The compiled verdict form}
\label{sec:spec-lang-verdictform}

\ruleblock{Definition}
\[
\mathsf{QVerdictForm} : [\, \mathrm{floors} : \mathsf{QRegionSat}^{*},\;
  \mathrm{marks} : \mathsf{QRegionSat}^{*} \,]
\]

A compiled satisfaction verdict carries its demand atoms in
\emph{two} lists, and the split is schema content because it is not
recoverable from the atoms alone: the same constraint species denotes
differently per face. A \emph{floor} holds over the whole credal
body; a \emph{mark} asserts that some admitted law attains --- so
folding the two into one list would erase exactly the distinction
the denotation reads. Polarity is \emph{declared} at formation ---
the cap sorts of section~\ref{sec:spec-lang-formation} --- and
\emph{discharged} here: one home per stratum, connected by the
compile. No battery
or region sort carries a polarity flag; the two lists are the
verdict-side home, with the floors compiled from the
$\mathsf{ForallSpec}$ components and the marks from the
$\mathsf{ExistsSpec}$ list plus the unconditional properness mark
--- attitude components take no part in the form, because a design
score is not a verdict ---
and the form's agreement with the declared faces is part of the
recompute obligation below. The mark list is the verdict's existential
conjunct, first-class: an \emph{empty} mark list is the all-safety
claim, true of the zero-mass subject, and a form that derived
liveness from floor content would silently certify a dead subject
live --- which is why the record-level compile emits the properness
mark unconditionally (\texttt{elabSpec\_top\_mem},
section~\ref{sec:spec-lang-spec}), so a form compiling a root
specification carries a nonempty mark list. Liveness demands enter
the mark list as \emph{negation walls}: the demand ``at least $q$''
enters as the mark of its complement's failure --- the wall/dip pair of
the compile target (section~\ref{sec:spec-lang-target}), as rows.
The form is value-like; its attachment to a root claim is a
zero-or-one relation (a claim with no form is simply not compiled);
and the agreement between the form and the specification it
compiles is a \emph{recompute obligation} at the ground seam, never
a trusted row: the checker re-derives the form from the spec it
cites.

\emph{Witness}: the staged certificate probes of the normative
theory's battery; corpus-wide re-emission follows the producer-side
emitter, which has not landed.

\ruleblock{API}
\apilang{TypeScript}{---}
No constructor is exposed: the form is compiled, not authored, and
no producer-side emission exists yet.
\apilang{Rust}{---}
Likewise.

\subsection{QSatConf: the confidence qualifier}
\label{sec:spec-lang-qsatconf}

\ruleblock{Definition}
\[
\mathsf{QSatConf} : [\, \mathrm{sys} : \mathsf{System},\;
  \mathrm{spec} : \mathsf{Spec},\;
  \alpha : \Q \,],
\qquad 0 < \alpha < 1
\]

The spec holds \emph{up to the evidence's error budget}: the claim is
backed by a betting-martingale ledger whose wealth crossed
$1/\alpha$, and soundness is time-uniform --- Ville's
inequality~\cite{ville1939} for nonnegative supermartingales plus a
randomized-last-look refinement, both machine-checked in the Lean
auxiliary tree. This is a \emph{distinct kind}, not a field on
$\mathsf{QSat}$: confidence-qualified rows are never silently mixed
with exact or almost-sure rows --- a sort distinction the loader
cannot blur, rather than a rendering convention. Sub-logic: rule
(ledger rows in, judgment out); the range demand on $\alpha$ is a
report-mode constraint at load, both ends strict ($\alpha = 0$
would claim the exact tier in statistical clothes; $\alpha = 1$
claims nothing). Ledger emission belongs to the statistical harness
on the system side; this row is the claim it backs.

\emph{Witness}: the LunarLander statistical rows ---
$\mathbb{E}[R] \ge 236.577$ at $\alpha = 1/20$ over $10^6$ episodes,
checked at nanoseconds per obligation; the continuous variant's
certified floor exceeding the discrete variant's empirical mean.

\ruleblock{API}
\apilang{TypeScript}{judgment.satConf}
\apiprov{signature --- \texttt{ts/src/judgment.ts}}
\begin{tssrc}
export function satConf(args: {
  sys: SysRef;
  spec: SpecRef;
  alpha: QLike;
}): Result<Artifact, JudgmentRefusal>
\end{tssrc}
\apiprov{construction --- \texttt{ts/tests/spec-native.test.ts}}
\begin{tssrc}
const art = s.unwrap(
  j.satConf({
    sys: j.sysRef("lunarlander-cont"),
    spec: j.specRef("return-floor-236577/1000"),
    alpha: Q.of(1, 20),
  }),
);
\end{tssrc}
\apilang{Rust}{mtk\_catalogue::judgment::QSatConf::new}
\apiprov{signature --- \texttt{crates/mtk-catalogue/src/judgment/kinds.rs}}
\begin{rustsrc}
pub fn new(sys: SysRef, spec: SpecRef, alpha: WitRat) -> Result<Self, Refusal>
\end{rustsrc}
\apiprov{construction --- \texttt{crates/mtk-catalogue/tests/golden\_spec.rs}}
\begin{rustsrc}
let row = QSatConf::new(
    SysRef("lunarlander-cont".into()),
    SpecRef("return-floor-236577/1000".into()),
    q(1, 20),
)
.unwrap();
\end{rustsrc}

\subsection{QSpecEquiv: spec equivalence}
\label{sec:spec-lang-qspecequiv}

\ruleblock{Definition}
\[
\mathsf{QSpecEquiv} : [\, a : \mathsf{Spec},\; b : \mathsf{Spec} \,]
\]

Two specs denote the same property over every behaviour carrier on
which both are interpretable. The judgment stores \emph{no
semantics-class field}, deliberately: a stored three-way class enum
misstates the structure (an $\omega$-word \emph{is} a c\`adl\`ag
trajectory under the piecewise-constant-at-integer-times embedding,
so the classes are not disjoint). Instead, \textbf{carrier scope is
derived, never stored}: it is the carrier-requirement synthesis
judgment of section~\ref{sec:spec-lang-formation} --- each former's
$\rhd$ annotation --- with an equivalence judgment's scope the
intersection of its two sides' requirements. An empty intersection
is a conformance refusal at the consumer, not a judgment. The
Return stratum's split horizons are this distinction already
encoded: tick-summed returns demand steps, integrated returns demand
duration.
Sub-logic: concluded by the rewrite layer's chains (rule);
the derived-requirement computation is part of the producer surface,
quoted below the tour entries.

\emph{Witness}: the 219 rewrite-chain artifacts of the liveness
suite, each concluding exactly this judgment.

\ruleblock{API}
\apilang{TypeScript}{judgment.specEquiv}
\apiprov{signature --- \texttt{ts/src/judgment.ts}}
\begin{tssrc}
export function specEquiv(args: { a: SpecRef; b: SpecRef }): Result<
  Artifact,
  JudgmentRefusal
>
\end{tssrc}
\apiprov{construction --- \texttt{ts/tests/spec-native.test.ts}}
\begin{tssrc}
s.unwrap(j.specEquiv({ a: j.specRef("spec-a"), b: j.specRef("spec-b") })),
\end{tssrc}
\apilang{Rust}{mtk\_catalogue::judgment::QSpecEquiv::new}
\apiprov{signature --- \texttt{crates/mtk-catalogue/src/judgment/kinds.rs}}
\begin{rustsrc}
pub fn new(a: SpecRef, b: SpecRef) -> Self
\end{rustsrc}
\apiprov{construction --- \texttt{crates/mtk-catalogue/tests/roundtrip\_spec.rs}}
\begin{rustsrc}
let equiv = QSpecEquiv::new(SpecRef("spec-a".into()), SpecRef("spec-b".into()));
\end{rustsrc}

\subsection{QRegionSat: the workhorse}
\label{sec:spec-lang-qregionsat}

\ruleblock{Definition}
\[
\mathsf{QRegionSat} : [\, \mathrm{formula} : \mathsf{StatePred},\;
  \mathrm{region} : \mathsf{QRegionAtom} \,]
\qquad
\forall x \in R :\ x \vDash_{\mathcal{M}} \varphi
\]
with the region-atom vocabulary
\[
\mathsf{qraSublevel}(\mathit{fam}, \theta) \;\mid\;
\mathsf{qraBox}(\mathit{intervals}) \;\mid\;
\mathsf{qraAffine}\Bigl(\textstyle\sum_i c_i V_i \le \theta\Bigr).
\]

Every state in the cited region satisfies the formula, with the
kernel region \emph{under}-approximating the formula's denotation ---
one orientation suffices, since complement consumption inherits the
over-approximation for free. The region atoms are how certified
regions are \emph{cited}: a sublevel set of a certified function
family, an exact-rational box, or a sublevel of an affine combination
of certified families. Sub-logic: concluded by the battery rules
(safety, deadline, until, Streett), basins, union bounds, the NNF
alias family, and region-combination nodes; battery well-formedness
policed by oracles; consumed by the tower-citation edge discipline (a
child's concluded region is a parent's atom) and by
satisfaction-from-region.

\emph{Witness}: the entire supermartingale corpus; the three quoted
constructions are the committed vectors for the three atom variants.

\ruleblock{API}
\apilang{TypeScript}{judgment.regionSat}
\apiprov{signature --- \texttt{ts/src/judgment.ts}}
\begin{tssrc}
export function regionSat(args: {
  formula: StatePred;
  region: RegionAtom;
}): Result<Artifact, JudgmentRefusal>
\end{tssrc}
\apiprov{construction --- \texttt{ts/tests/spec-native.test.ts}}
\begin{tssrc}
const phi = s.polyLe(0n, s.unwrap(s.polyVar(1, 2)));

const sub = s.unwrap(
  j.regionSat({
    formula: phi,
    region: j.sublevel(j.certFamRef("lyap-v1"), Q.of(7, 5)),
  }),
);
\end{tssrc}
\apiprov{The box and affine region-atom variants ride sibling committed vectors in the same tests.}
\apilang{Rust}{mtk\_catalogue::judgment::QRegionSat::new}
\apiprov{signature --- \texttt{crates/mtk-catalogue/src/judgment/kinds.rs}}
\begin{rustsrc}
pub fn new(formula: StatePred, region: RegionAtom) -> Self
\end{rustsrc}
\apiprov{construction --- \texttt{crates/mtk-catalogue/tests/golden\_spec.rs}}
\begin{rustsrc}
let phi = StatePred::poly_le(q(0, 1), Poly::var(1, 2).unwrap());
let sublevel = QRegionSat::new(
    phi.clone(),
    RegionAtom::Sublevel {
        cert: CertFamRef("lyap-v1".into()),
        thr: q(7, 5),
    },
);
\end{rustsrc}
\apiprov{The box and affine region-atom variants ride sibling committed vectors in the same tests.}

\subsection{QCapSat: the cap-on-region vehicle}
\label{sec:spec-lang-qcapsat}

\ruleblock{Definition}
\[
\mathsf{QCapSat} : [\, \mathrm{cap} : \mathsf{ForallSpec},\;
  \mathrm{region} : \mathsf{QRegionAtom} \,]
\qquad
\forall x \in R :\ x \vDash_{\mathcal{M}} \mathrm{cap}
\]

A probability cap certified from a region: every state in the cited
region satisfies the safety-face component, with the region
\emph{under}-approximating the cap's denotation --- the workhorse's
orientation, one entry up the formula grid. The kind entered at
revision event 4 as a deliberately \emph{interim} vehicle: the
previous edition nested the probability former under the region
workhorse's state formula, the formation grid makes that shape
unrepresentable (a cap is a component, never a state predicate), and
so the cap-discharging rules --- batteries, transports, reach tables
--- conclude this kind instead. Revision event 6 discharges the
interim framing. The normative theory stores those conclusions as
attributes of the battery proof nodes themselves and ties caps to
spec components by keying each satisfaction claim to the component's
own entity row: \texttt{qsat\_warrant} relates the claim to a
\texttt{QForallSat} row carrying the component and its region. The
producer's rule rows are re-staged against that seam --- warrants
are emitted entity-keyed, under a versioned species binding that
refuses pre-revision closures at the seam, loudly, at derivation
time. What no producer emits yet is the compiled verdict form
(section~\ref{sec:spec-lang-verdictform} records the same absence),
so this kind remains the typed conclusion the cap-discharging rules
place on the wire.

\emph{Witness}: the committed vector (the quoted construction, both
languages byte-identical); the battery conclusions across the corpus
exercise the same shape.

\ruleblock{API}
\apilang{TypeScript}{judgment.capSat}
\apiprov{signature --- \texttt{ts/src/judgment.ts}}
\begin{tssrc}
export function capSat(args: {
  cap: ForallSpec;
  region: RegionAtom;
}): Result<Artifact, JudgmentRefusal>
\end{tssrc}
\apiprov{construction --- \texttt{ts/tests/spec-native.test.ts}}
\begin{tssrc}
s.unwrap(
  j.capSat({
    cap: s.probGe(Q.of(9, 10), s.always(phi)),
    region: j.sublevel(j.certFamRef("lyap-v1"), Q.of(1, 10)),
  }),
),
\end{tssrc}
\apiprov{Entered at revision event 4 as the interim typed conclusion for the cap-discharging rules (batteries, transports, reach tables); at revision event 6 the warrant seam is entity-keyed and the producer's rule rows are re-staged against it, and this kind remains the typed conclusion those rules place on the wire pending a verdict-form producer.}
\apilang{Rust}{mtk\_catalogue::judgment::QCapSat::new}
\apiprov{signature --- \texttt{crates/mtk-catalogue/src/judgment/kinds.rs}}
\begin{rustsrc}
pub fn new(cap: ForallSpec, region: RegionAtom) -> Self
\end{rustsrc}
\apiprov{construction --- \texttt{crates/mtk-catalogue/tests/golden\_spec.rs}}
\begin{rustsrc}
let phi = StatePred::poly_le(q(0, 1), Poly::var(0, 1).unwrap());
let row = QCapSat::new(
    ForallSpec::prob_ge(q(9, 10), PathPred::always(phi)).unwrap(),
    RegionAtom::Sublevel {
        cert: CertFamRef("lyap-v1".into()),
        thr: q(1, 10),
    },
);
\end{rustsrc}
\apiprov{Entered at revision event 4 as the interim typed conclusion for the cap-discharging rules (batteries, transports, reach tables); at revision event 6 the warrant seam is entity-keyed and the producer's rule rows are re-staged against it, and this kind remains the typed conclusion those rules place on the wire pending a verdict-form producer.}

\subsection{QInvScope: invariance with scope}
\label{sec:spec-lang-qinvscope}

\ruleblock{Definition}
\[
\mathsf{QInvScope} : [\, \mathrm{region} : \mathsf{QRegionAtom},\;
  \mathrm{init} : \mathsf{QRegionAtom} \,]
\qquad
S \text{ forward-invariant} \wedge X_0 \subseteq S
\]

The region is forward-invariant and contains the initial set.
Sub-logic: concluded by the block-invariance rule (strictly
facet-inward cone rows plus containment) and its lattice alias;
consumed by basins, by battery scope-warrant fields, and by the
aggregate-scope kind's invariance route below.
\emph{Witness}: the invariance rows across the corpus; the
Morse-stratum discharge kits cite it wherever a block's scope must be
warranted.

\ruleblock{API}
\apilang{TypeScript}{judgment.invScope}
\apiprov{signature --- \texttt{ts/src/judgment.ts}}
\begin{tssrc}
export function invScope(args: {
  region: RegionAtom;
  init: RegionAtom;
}): Result<Artifact, JudgmentRefusal>
\end{tssrc}
\apiprov{construction --- \texttt{ts/tests/spec-native.test.ts}}
\begin{tssrc}
s.unwrap(j.invScope({ region: unitBox, init: j.box([[Q.of(1, 4), Q.of(1, 2)]]) })),
\end{tssrc}
\apilang{Rust}{mtk\_catalogue::judgment::QInvScope::new}
\apiprov{signature --- \texttt{crates/mtk-catalogue/src/judgment/kinds.rs}}
\begin{rustsrc}
pub fn new(region: RegionAtom, init: RegionAtom) -> Self
\end{rustsrc}
\apiprov{construction --- \texttt{crates/mtk-catalogue/tests/roundtrip\_spec.rs}}
\begin{rustsrc}
let inv = QInvScope::new(
    unit_box(),
    RegionAtom::Box(IntervalBox(vec![Interval::new(q(1, 4), q(1, 2)).unwrap()])),
);
\end{rustsrc}

\subsection{QMorse: the global stratum}
\label{sec:spec-lang-qmorse}

\ruleblock{Definition}
\[
\mathsf{QMorse} : [\, \mathrm{graph} : \mathsf{MorseGraph},\;
  \rho : \mathsf{VertexRegion}^{*},\;
  \mathrm{arena} : \mathsf{Arena} \,]
\qquad
\mathcal{M} \vDash_D \mathsf{Morse}(G, \rho)
\]

Non-initialized, whole-system qualitative structure: a certified
Morse graph with per-vertex readout regions. It is listed in this
catalogue because a complete spec-language inventory must say where
the non-formula-shaped claims live: the formula language and the
global stratum share the region vocabulary and nothing else, and the
$D$ subscript marks the different satisfaction relation. The row
asserts the graph and readouts against a declared \emph{arena} ---
the grid complex the claim is read on; a discharged claim's arena
\emph{is} its lattice's, tied by a stated equation between
hash-consed values (the tie is deliberately syntactic --- semantic
region equality is geometry the certificate does not owe).
Acyclicity is the checker's to
confirm, never this row's to assert, and the payload refuses
out-of-range edge endpoints and vertex indices with the geometry
named. Sub-logic: the Morse discharge and product rules of
chapter~\ref{ch:qtsl-manual}, which carry the deepest mechanized
ledger in the stack.

\emph{Witness}: the Morse-graph certificate family across the
corpus, up to certified graphs at the several-hundred-node scale.

\ruleblock{API}
\apilang{TypeScript}{judgment.morse}
\apiprov{signature --- \texttt{ts/src/judgment.ts}}
\begin{tssrc}
export function morse(args: {
  vertices: number;
  edges: ReadonlyArray<MorseEdge>;
  rho: ReadonlyArray<VertexRegion>;
}): Result<Artifact, JudgmentRefusal>
\end{tssrc}
\apiprov{construction --- \texttt{ts/tests/spec-native.test.ts}}
\begin{tssrc}
s.unwrap(
  j.morse({
    vertices: 2,
    edges: [{ src: 0, dst: 1 }],
    rho: [
      { vertex: 0, region: phi },
      { vertex: 1, region: s.TRUE },
    ],
  }),
),
\end{tssrc}
\apilang{Rust}{mtk\_catalogue::judgment::QMorse::new}
\apiprov{signature --- \texttt{crates/mtk-catalogue/src/judgment/kinds.rs}}
\begin{rustsrc}
pub fn new(graph: MorseGraph, rho: Vec<VertexRegion>) -> Result<Self, Refusal>
\end{rustsrc}
\apiprov{construction --- \texttt{crates/mtk-catalogue/tests/roundtrip\_spec.rs}}
\begin{rustsrc}
let morse = QMorse::new(
    MorseGraph {
        vertices: 2,
        edges: vec![MorseEdge { src: 0, dst: 1 }],
    },
    vec![
        VertexRegion {
            vertex: 0,
            region: StatePred::poly_le(q(0, 1), Poly::var(0, 1).unwrap()),
        },
        VertexRegion {
            vertex: 1,
            region: StatePred::True,
        },
    ],
)
.unwrap();
\end{rustsrc}

\subsection{QWireCert: wiring certification}
\label{sec:spec-lang-qwirecert}

\ruleblock{Definition}
\[
\mathsf{QWireCert} : [\, w : \mathsf{Wiring},\;
  \mathrm{premises} : \mathsf{Spec}^{*},\;
  \mathrm{out} : \mathsf{Spec},\;
  \kappa : \mathsf{Slack} \,]
\]

The wiring $w$ is certified for $(\mathit{spec}_i) \Rightarrow
\mathit{spec}_{\mathrm{out}}$ with slack $\kappa$ in the
$\mathcal{K}^0_\infty$ class. The premise list is positional
against the wiring's inner-system order; lockstep with a consuming
composition node's components is record-level conformance. Sub-logic: discharged from proof
leaves for the two condition families of the composition calculus
(mechanism certification; wiring-lens
certification)~\cite{capuccimyers2026}, per region;
$\mathcal{K}_\infty$-membership rows are proof-dischargeable;
consumed by the wire-composition rule. The class membership and slack
algebra are proven; the payload here freezes the API shape while the
row's upstream schema staging follows the system-side work --- stated
so the reader knows which part is mathematics and which part is
plumbing.

\emph{Witness}: the assume--guarantee worked examples of
chapter~\ref{ch:qtsl-manual}.

\ruleblock{API}
\apilang{TypeScript}{judgment.wireCert}
\apiprov{signature --- \texttt{ts/src/judgment.ts}}
\begin{tssrc}
export function wireCert(args: {
  wiring: WiringRef;
  premises: ReadonlyArray<SpecRef>;
  conclusion: SpecRef;
  gain: Poly;
}): Result<Artifact, JudgmentRefusal>
\end{tssrc}
\apiprov{construction --- \texttt{ts/tests/spec-native.test.ts}}
\begin{tssrc}
s.unwrap(
  j.wireCert({
    wiring: j.wiringRef("w-loop"),
    premises: [j.specRef("inner-a"), j.specRef("inner-b")],
    conclusion: j.specRef("outer"),
    gain: s.unwrap(s.poly([{ expts: [1], coeff: Q.of(1, 2) }])),
  }),
),
\end{tssrc}
\apilang{Rust}{mtk\_catalogue::judgment::QWireCert::new}
\apiprov{signature --- \texttt{crates/mtk-catalogue/src/judgment/kinds.rs}}
\begin{rustsrc}
pub fn new(
    wiring: WiringRef,
    premises: Vec<SpecRef>,
    conclusion: SpecRef,
    slack: Slack,
) -> Self
\end{rustsrc}
\apiprov{construction --- \texttt{crates/mtk-catalogue/tests/roundtrip\_spec.rs}}
\begin{rustsrc}
let wire_cert = QWireCert::new(
    WiringRef("w-loop".into()),
    vec![SpecRef("inner-a".into()), SpecRef("inner-b".into())],
    SpecRef("outer".into()),
    Slack {
        gain: Poly::new(vec![(vec![1u32], q(1, 2))]).unwrap(),
    },
);
\end{rustsrc}

\subsection{QDist: distance judgments}
\label{sec:spec-lang-qdist}

\ruleblock{Definition}
\[
\mathsf{QDist}[d](A, B) \le \varepsilon
\qquad
d : \text{a declared distance functor (registry entry)},\;
\varepsilon \in \Q_{\ge 0}
\]

The declared distance functor --- a registry entry naming subject
sort, symmetry flag, licensed rules, and observable meaning ---
between two subjects (system manifests, distribution snapshots, or
definite descriptions with their well-formedness licenses) is at most
exact-rational $\varepsilon$. Judgment-primary by design: triangle,
weakening, and symmetry-per-flag are judgment-level admissibles,
while the discharge rules are their own catalogue entries ---
coupling-witnessed drift, contraction-witnessed fixed-point balls,
and the transport cells that degrade a judgment across an approximate
simulation with the $n\varepsilon$ compounding law made explicit.
\emph{Schema status}: chartered design --- this kind's row species
is not yet part of the normative row schema
(\path{theories/Qtsl.geolog}); it enters with its discharge lanes,
and until then the entry stands as the registry slot for the
demand.

\emph{Witness}: the coupling and contraction fixture families; the
finite-kernel transport tier.

\ruleblock{API}
\apilang{TypeScript}{judgment.dist}
\apiprov{signature --- \texttt{ts/src/judgment.ts}}
\begin{tssrc}
export function dist(args: {
  functor: string;
  a: DistSubject;
  b: DistSubject;
  eps: QLike;
}): Result<Artifact, JudgmentRefusal>
\end{tssrc}
\apiprov{construction --- \texttt{ts/tests/spec-native.test.ts}}
\begin{tssrc}
s.unwrap(
  j.dist({
    functor: "wasserstein-1",
    a: { kind: "manifest", sys: j.sysRef("sys-a") },
    b: { kind: "snapshot", locator: "snap-7" },
    eps: Q.of(1, 100),
  }),
),
\end{tssrc}
\apilang{Rust}{mtk\_catalogue::judgment::QDist::new}
\apiprov{signature --- \texttt{crates/mtk-catalogue/src/judgment/kinds.rs}}
\begin{rustsrc}
pub fn new(
    functor: impl Into<String>,
    a: DistSubject,
    b: DistSubject,
    eps: WitRat,
) -> Result<Self, Refusal>
\end{rustsrc}
\apiprov{construction --- \texttt{crates/mtk-catalogue/tests/roundtrip\_spec.rs}}
\begin{rustsrc}
let dist = QDist::new(
    "wasserstein-1",
    DistSubject::Manifest(SysRef("sys-a".into())),
    DistSubject::Snapshot("snap-7".into()),
    q(1, 100),
)
.unwrap();
\end{rustsrc}

\subsection{\texorpdfstring{The existence family: QZeroExists, QEquilibriumExists,\protect\\ QFixedPointExists}{The existence family: QZeroExists, QEquilibriumExists, QFixedPointExists}}
\label{sec:spec-lang-qexists}

\ruleblock{Definition}
\[
\mathsf{QZeroExists} : [\, g : \mathsf{Poly}^{m},\;
  \mathrm{cell} : \mathsf{QRegionAtom} \,]
\qquad
\exists x \in \mathrm{cell} :\ g(x) = 0
\]
with the equilibrium and fixed-point members reading the same shape
over a drift field and a map ($g = F - \mathrm{id}$, tied by an
elaboration recompute downstream).

Poincar\'e--Miranda member claims: a zero, an equilibrium, a fixed
point, \emph{in a cell}. Existence only --- nothing in this family
concludes uniqueness, count, or non-existence; that refusal boundary
is the family's design. In this edition the cell is a box atom and
the subject is equidimensional ($m$ components over an
$m$-dimensional cell), refused otherwise with the geometry named.
Sub-logic: the shared face-sign battery (rule + pinned oracle), and
a derived rule reading an inward invariance box as an equilibrium
witness.

\emph{Witness}: the face-sign battery fixtures; the equilibrium
readouts of the Morse-stratum examples.

\ruleblock{API}
\apilang{TypeScript}{judgment.zeroExists / equilibriumExists / fixedPointExists}
\apiprov{signature --- \texttt{ts/src/judgment.ts}}
\begin{tssrc}
export function zeroExists(args: {
  g: ReadonlyArray<Poly>;
  cell: RegionAtom;
}): Result<Artifact, JudgmentRefusal>

export function equilibriumExists(args: {
  field: ReadonlyArray<Poly>;
  cell: RegionAtom;
}): Result<Artifact, JudgmentRefusal>

export function fixedPointExists(args: {
  map: ReadonlyArray<Poly>;
  cell: RegionAtom;
}): Result<Artifact, JudgmentRefusal>
\end{tssrc}
\apiprov{construction --- \texttt{ts/tests/spec-native.test.ts}}
\begin{tssrc}
s.unwrap(j.zeroExists({ g: g1, cell: unitBox })),
\end{tssrc}
\apiprov{The equilibrium and fixed-point members construct beside it in the same test.}
\apilang{Rust}{mtk\_catalogue::judgment::\{QZeroExists, QEquilibriumExists, QFixedPointExists\}::new}
\apiprov{signature --- \texttt{crates/mtk-catalogue/src/judgment/kinds.rs}}
\begin{rustsrc}
pub fn new(g: Vec<Poly>, cell: RegionAtom) -> Result<Self, Refusal>

pub fn new(field: Vec<Poly>, cell: RegionAtom) -> Result<Self, Refusal>

pub fn new(map: Vec<Poly>, cell: RegionAtom) -> Result<Self, Refusal>
\end{rustsrc}
\apiprov{construction --- \texttt{crates/mtk-catalogue/tests/roundtrip\_spec.rs}}
\begin{rustsrc}
let zero = QZeroExists::new(g1(), unit_box()).unwrap();
\end{rustsrc}
\apiprov{The equilibrium and fixed-point members construct beside it in the same test.}

\subsection{QSim: exact simulation}
\label{sec:spec-lang-qsim}

\ruleblock{Definition}
\[
\mathsf{QSim} : [\, \mathrm{src} : \mathsf{System},\;
  \mathrm{dst} : \mathsf{System},\;
  \sigma : (\N \times \mathsf{Mode} \times \mathsf{Poly})^{*} \,]
\qquad
\sigma : S \to S' \text{ an exact simulation map}
\]

The refinement preorder, one-sided by design: $\sigma$ carries the
finer system onto the coarser, and simulations compose. The map is
stored as flat (coordinate, mode, polynomial) triples --- one
polynomial per destination coordinate per source mode, the mode
cited as an entity rather than by name string, so a piece citing a
nonexistent mode is unrepresentable rather than a silent key
mismatch. The producer payload carries the same content
mode-indexed --- the outer list ranges over the source's modes, each
inner list is that mode's destination coordinates --- matching the
transport builder's convention, refusing an empty mode roster and
empty per-mode maps. Sub-logic:
discharge at the strict tier (compatibility squares as proof leaves
plus a mode-coverage walk) and the kernel tier (intertwining of
kernels, machine-checked at its pushforward lemma); consumed by the
simulation-transport rule, whose simulation premise is an ordinary
premise of this kind --- no hypothesis asterisk. The approximate
$(\sigma, \varepsilon)$ class deliberately concludes
$\mathsf{QDist}$, never this kind: exactness is what makes transport
lossless, and blurring the two sorts would launder approximation
through an exact-looking judgment.

\emph{Witness}: the simulation-transport worked examples; the
LunarLander tick-semantics pair is carried by this transport.

\ruleblock{API}
\apilang{TypeScript}{judgment.sim}
\apiprov{signature --- \texttt{ts/src/judgment.ts}}
\begin{tssrc}
export function sim(args: {
  sigma: ReadonlyArray<ReadonlyArray<Poly>>;
  from: SysRef;
  to: SysRef;
}): Result<Artifact, JudgmentRefusal>
\end{tssrc}
\apiprov{construction --- \texttt{ts/tests/spec-native.test.ts}}
\begin{tssrc}
s.unwrap(
  j.sim({
    // one mode, one destination coordinate
    sigma: [[s.unwrap(s.polyVar(0, 1))]],
    from: j.sysRef("fine"),
    to: j.sysRef("coarse"),
  }),
),
\end{tssrc}
\apilang{Rust}{mtk\_catalogue::judgment::QSim::new}
\apiprov{signature --- \texttt{crates/mtk-catalogue/src/judgment/kinds.rs}}
\begin{rustsrc}
pub fn new(sigma: Vec<Vec<Poly>>, from: SysRef, to: SysRef) -> Result<Self, Refusal>
\end{rustsrc}
\apiprov{construction --- \texttt{crates/mtk-catalogue/tests/roundtrip\_spec.rs}}
\begin{rustsrc}
let sim = QSim::new(
    vec![vec![Poly::var(0, 1).unwrap()]], // one mode, one coordinate
    SysRef("fine".into()),
    SysRef("coarse".into()),
)
.unwrap();
\end{rustsrc}

\subsection{QFlatEq: flattening}
\label{sec:spec-lang-qflateq}

\ruleblock{Definition}
\[
\mathsf{QFlatEq} : [\, \mathrm{sys} : \mathsf{System},\;
  \mathrm{flat} : \mathsf{Manifest} \,]
\qquad
\mathrm{flat} = \mathrm{flatten}(\mathrm{sys})
\]

The monolithic manifest is the flattening of the composite.
Sub-logic: one oracle-shaped discharge (the substitution and
guard-product recompute, registry-pinned, refusing loudly when
unpinned), with the shared-carrier and binary-wiring confinements as
refusal boundaries --- completeness costs, never soundness
hypotheses: what the oracle cannot flatten it refuses to speak about,
rather than asserting under an assumption. Consumed by the
flattening substitution rule. The measurable-tier witness (the
parallel-composition law) is machine-checked.

\emph{Witness}: the flattening fixtures of the composition suite.

\ruleblock{API}
\apilang{TypeScript}{judgment.flatEq}
\apiprov{signature --- \texttt{ts/src/judgment.ts}}
\begin{tssrc}
export function flatEq(args: {
  composite: SysRef;
  flat: SysRef;
}): Result<Artifact, JudgmentRefusal>
\end{tssrc}
\apiprov{construction --- \texttt{ts/tests/spec-native.test.ts}}
\begin{tssrc}
s.unwrap(j.flatEq({ composite: j.sysRef("comp"), flat: j.sysRef("mono") })),
\end{tssrc}
\apilang{Rust}{mtk\_catalogue::judgment::QFlatEq::new}
\apiprov{signature --- \texttt{crates/mtk-catalogue/src/judgment/kinds.rs}}
\begin{rustsrc}
pub fn new(composite: SysRef, flat: SysRef) -> Self
\end{rustsrc}
\apiprov{construction --- \texttt{crates/mtk-catalogue/tests/roundtrip\_spec.rs}}
\begin{rustsrc}
let flat = QFlatEq::new(SysRef("comp".into()), SysRef("mono".into()));
\end{rustsrc}

\subsection{QAggScope: aggregate scope}
\label{sec:spec-lang-qaggscope}

\ruleblock{Definition}
\[
\mathsf{QAggScope} : [\, N : \mathsf{System},\;
  \mathrm{aggdim} : \N,\;
  K : \mathsf{PolytopeRow}^{*},\;
  \mathrm{stop} : \{\mathsf{reachable}, \mathsf{stopped}\} \,]
\]

The aggregate image of every in-scope reachable marking of the net
$N$ lies in the polytope $K$, under the declared scope semantics.
Reachability is composition-unstable (so not chase) and semantic (so
not oracle) --- which is why this is a judgment kind with
rule discharge: the three confinement routes (cite an invariance-scope
row; exit-guard structure, by oracle walk over guard rows; stopped
semantics, by declared flag with a conformance tie) each conclude it,
alongside P-invariant conservation certificates ($y \cdot C = 0$,
exact-rational, witness carried) and capacity rows. Consumed by the
aggregation-reduction rule as an ordinary premise. \emph{Schema
status}: designed, not yet rows --- the aggregation-reduction
node's row kits are in the normative schema; the scope kind's own
row species joins them when its rows land.

\emph{Witness}: the structurally bounded net families of the
aggregation suite.

\ruleblock{API}
\apilang{TypeScript}{judgment.aggScope}
\apiprov{signature --- \texttt{ts/src/judgment.ts}}
\begin{tssrc}
export function aggScope(args: {
  net: SysRef;
  aggDim: number;
  polytope: ReadonlyArray<{ coeffs: ReadonlyArray<QLike>; ub: QLike }>;
  stop: ScopeSemantics;
}): Result<Artifact, JudgmentRefusal>
\end{tssrc}
\apiprov{construction --- \texttt{ts/tests/spec-native.test.ts}}
\begin{tssrc}
s.unwrap(
  j.aggScope({
    net: j.sysRef("net-a"),
    aggDim: 2,
    polytope: [{ coeffs: [1n, 1n], ub: 10n }],
    stop: "reachable",
  }),
),
\end{tssrc}
\apilang{Rust}{mtk\_catalogue::judgment::QAggScope::new}
\apiprov{signature --- \texttt{crates/mtk-catalogue/src/judgment/kinds.rs}}
\begin{rustsrc}
pub fn new(
    net: SysRef,
    agg_dim: usize,
    polytope: Vec<PolytopeRow>,
    stop: ScopeSemantics,
) -> Result<Self, Refusal>
\end{rustsrc}
\apiprov{construction --- \texttt{crates/mtk-catalogue/tests/roundtrip\_spec.rs}}
\begin{rustsrc}
let agg = QAggScope::new(
    SysRef("net-a".into()),
    2,
    vec![PolytopeRow {
        coeffs: vec![q(1, 1), q(1, 1)],
        ub: q(10, 1),
    }],
    ScopeSemantics::Reachable,
)
.unwrap();
\end{rustsrc}

\subsection{QStateFinite: a classification judgment}
\label{sec:spec-lang-qstatefinite}

\ruleblock{Definition}
\[
\mathsf{QStateFinite} : [\, \mathrm{sys} : \mathsf{System} \,]
\qquad
\text{the subject's state type is finite}
\]

A \emph{classification} judgment: a class membership of the rule's
subject, not a per-application fact. Sub-logic: chase-derivable by
structural refinement over the state-type formers (enumerations are
finite; finite products of finite factors are finite; nesting
inherits; the tensor former derives nothing, a deliberate negative),
with a decidable oracle as fallback. Consumed where a discharge rule
needs finiteness as a premise (the coupling-drift finite tier), which
retires a side-condition asterisk into an ordinary premise; a future
countable-tier proof targets a weaker classification, display
untouched. \emph{Schema status}: designed, not yet rows in the
normative schema; the classification enters with its
chase-refinement family.

\emph{Witness}: the finite-tier coupling fixtures.

\ruleblock{API}
\apilang{TypeScript}{judgment.stateFinite}
\apiprov{signature --- \texttt{ts/src/judgment.ts}}
\begin{tssrc}
export function stateFinite(args: { sys: SysRef }): Result<
  Artifact,
  JudgmentRefusal
>
\end{tssrc}
\apiprov{construction --- \texttt{ts/tests/spec-native.test.ts}}
\begin{tssrc}
s.unwrap(j.stateFinite({ sys: j.sysRef("sys-a") })),
\end{tssrc}
\apilang{Rust}{mtk\_catalogue::judgment::QStateFinite::new}
\apiprov{signature --- \texttt{crates/mtk-catalogue/src/judgment/kinds.rs}}
\begin{rustsrc}
pub fn new(sys: SysRef) -> Self
\end{rustsrc}
\apiprov{construction --- \texttt{crates/mtk-catalogue/tests/roundtrip\_spec.rs}}
\begin{rustsrc}
let fin = QStateFinite::new(SysRef("sys-a".into()));
\end{rustsrc}

\subsection{QLinGrowth: a growth-class judgment}
\label{sec:spec-lang-qlingrowth}

\ruleblock{Definition}
\[
\mathsf{QLinGrowth} : [\, f : \mathsf{Poly}^{*},\;
  a, b : \Q_{\ge 0},\;
  \mathrm{scope} : \mathsf{Box} \,]
\qquad
\lVert f(x) \rVert \le a \lVert x \rVert + b \ \text{on the scope}
\]

Membership in the linear-growth class with the \emph{constants
carried in the judgment}, so downstream cells consume $(a, b)$
arithmetically rather than a membership bit. The kind is load-bearing
because the fence it polices is provably sharp: the transport law it
feeds is \emph{false} outside the fragment (the $n\varepsilon$
compounding of the distance-transport cell is sharp without a
contraction row), so growth control must arrive as evidence, not
assumption. Sub-logic: certificate-carrying --- affine-envelope and
Bernstein rows discharge it for polynomial drifts; composition and
weakening admissibles do constant arithmetic.

\emph{Witness}: the distance-transport pushforward cells' premise
rows.

\ruleblock{API}
\apilang{TypeScript}{judgment.linGrowth}
\apiprov{signature --- \texttt{ts/src/judgment.ts}}
\begin{tssrc}
export function linGrowth(args: {
  f: ReadonlyArray<Poly>;
  a: QLike;
  b: QLike;
  scope: RegionAtom;
}): Result<Artifact, JudgmentRefusal>
\end{tssrc}
\apiprov{construction --- \texttt{ts/tests/spec-native.test.ts}}
\begin{tssrc}
s.unwrap(
  j.linGrowth({ f: g1, a: 2n, b: Q.of(1, 2), scope: unitBox }),
),
\end{tssrc}
\apilang{Rust}{mtk\_catalogue::judgment::QLinGrowth::new}
\apiprov{signature --- \texttt{crates/mtk-catalogue/src/judgment/kinds.rs}}
\begin{rustsrc}
pub fn new(f: Vec<Poly>, a: WitRat, b: WitRat, scope: RegionAtom) -> Result<Self, Refusal>
\end{rustsrc}
\apiprov{construction --- \texttt{crates/mtk-catalogue/tests/roundtrip\_spec.rs}}
\begin{rustsrc}
let growth = QLinGrowth::new(g1(), q(2, 1), q(1, 2), unit_box()).unwrap();
\end{rustsrc}

\subsection{The attitude blend: SHurwiczGE and SHurwiczLE (not a
judgment kind)}
\label{sec:spec-lang-hurwicz}

\ruleblock{Formation}
\begin{mathpar}
\inferrule*[Right=\textsc{SHurwiczGE}]
  {\alpha \in \mathbb{Q} \\ 0 \le \alpha \le 1 \\
   B \in \mathbb{Q} \\ R\ \mathsf{PathObservable}}
  {\mathsf{SHurwiczGE}(\alpha, B, R)\ \mathsf{AttitudeSpec}}
\\
\inferrule*[Right=\textsc{SHurwiczLE}]
  {\alpha \in \mathbb{Q} \\ 0 \le \alpha \le 1 \\
   B \in \mathbb{Q} \\ R\ \mathsf{PathObservable}}
  {\mathsf{SHurwiczLE}(\alpha, B, R)\ \mathsf{AttitudeSpec}}
\end{mathpar}

Certification is pessimism; design exploration wants the other end
too, and occasionally a weighted blend --- the Hurwicz
optimism--pessimism criterion, transposed. Where the credal reading
leaves a \emph{set} of values at a run-functional, the blend mixes
the two envelope readings with a declared attitude weight:
$\alpha = 0$ is the pure worst case ($\mathsf{SReturnGE}$'s floor
face, exactly), $\alpha = 1$ the pure best case
($\mathsf{SReturnMayBeGE}$'s witness face, exactly), and an interior
$\alpha$ mixes the two envelopes \emph{arithmetically} into one
scalar before thresholding, in checker-recomputed rational
arithmetic. The pair therefore fits \emph{neither} cap face, which
is why it has its own sort instead of a mis-declared face --- and
the contrast with the record deserves its precise statement,
because a specification is not envelope-pure either. A
specification is \emph{polarized}: every component carries a
declared face, so its compiled verdict is truth-functional over
per-component single-envelope discharges, monotone
separately-per-slot with fixed sign, and verdict-mixing commutes
with composition. The interior blend's scalar has no sign in the
refinement order --- it mixes a quantity that shrinks under new
admitted laws with one that grows --- and it fails the
marginal-extension tower property and equivalence with any single
envelope bound (exact-rational countermodels are kernel-checked,
both faces). In one line: \textbf{a Spec mixes verdicts; an
attitude mixes envelopes} --- and envelope-mixing is root-only,
which is exactly where the record does its own mixing anyway.

Root-only-ness is now \emph{type-level} --- the fence-to-type
upgrade of the theory revision: the pair concludes its own small
cap sort, whose only consumer is the record's score slot
(\texttt{spec\_score}, section~\ref{sec:spec-lang-record}) --- no
connective, path former, or atom takes an $\mathsf{AttitudeSpec}$
argument, so occurrence below a root is unrepresentable rather than
oracle-rejected, and the earlier edition's whole-tree typing clause
for this former is retired. What has \emph{not} changed is the
judgment-side exclusion, and it is the design: no rule concludes an
attitude claim and no judgment sort exists for one --- admitting
the interior blend into the derivation calculus would break the
citation discipline everything else preserves. Evaluation
decomposes into the two extreme passes --- one floor certificate,
one witness certificate --- mixed once at the root, and a verdict
renders as a \emph{design score}, never a guarantee. Nothing is
lost at the endpoints: a nested blend at $\alpha \in \{0, 1\}$ has
its canonical spelling in the sibling formers.

\emph{Witness}: the attitude fixtures of the producer surface; the
credal-stratum examples of the semantic universe chapter
(chapter~\ref{ch:qtsl-semantic-universe}, the section on agreement
and the attitude operator).

\ruleblock{API}
\apilang{TypeScript}{spec.hurwiczGe / hurwiczLe}
\apiprov{signature --- \texttt{ts/src/judgment.ts + ts/src/spec.ts}}
\begin{tssrc}
export function hurwiczGe(
  alpha: QLike,
  lo: QLike,
  ret: PathObservable,
): AttitudeSpec

export function hurwiczLe(
  alpha: QLike,
  hi: QLike,
  ret: PathObservable,
): AttitudeSpec

export function hurwiczQuery(args: {
  alpha: QLike;
  root: SpecRef;
}): Result<Artifact, JudgmentRefusal>
\end{tssrc}
\apiprov{construction --- \texttt{ts/tests/spec-native.test.ts}}
\begin{tssrc}
const hurwicz = s.hurwiczGe(Q.of(1, 3), 2n, s.poSum(11, s.oChan(1)));
\end{tssrc}
\apiprov{The root-level query row (\texttt{judgment.hurwiczQuery}) remains beside the component formers as the judgment-tour vehicle for the blend.}
\apilang{Rust}{mtk\_catalogue::spec::AttitudeSpec::\{hurwicz\_ge, hurwicz\_le\}}
\apiprov{signature --- \texttt{crates/mtk-catalogue/src/judgment/kinds.rs + crates/mtk-catalogue/src/spec/formula.rs}}
\begin{rustsrc}
pub fn hurwicz_ge(
    alpha: WitRat,
    lo: WitRat,
    ret: PathObservable,
) -> Result<AttitudeSpec, Refusal>

pub fn hurwicz_le(
    alpha: WitRat,
    hi: WitRat,
    ret: PathObservable,
) -> Result<AttitudeSpec, Refusal>

pub fn new(alpha: WitRat, root: SpecRef) -> Result<Self, Refusal>
\end{rustsrc}
\apiprov{construction --- \texttt{crates/mtk-catalogue/tests/golden\_spec.rs}}
\begin{rustsrc}
let hurwicz = AttitudeSpec::hurwicz_ge(
    q(1, 3),
    q(2, 1),
    PathObservable::sum(11, StateObservable::chan(1)).unwrap(),
)
.unwrap();
\end{rustsrc}
\apiprov{The root-level query row (\texttt{HurwiczQueryV0}) remains beside the component formers as the judgment-tour vehicle for the blend.}

\addtocontents{toc}{\protect\setcounter{tocdepth}{2}}
 \chapter{Rule Catalogue}
\label{ch:qtsl-manual}

This chapter is the reference manual of the proof theory, and it is
generated, not written: it incorporates the literate rule catalogue in
full, emitted mechanically from the same sources the calculus is built
from --- the hand-authored rule cells, the Lean~4 development
(statements extracted by declaration name and bound by hash at
emission), the verifier-receipt store, the HOL4 development,
and the per-rule proof files. Every rule of the proof theory has its
section: statement in Steele's metanotation, its producer API in
TypeScript and Rust, commentary, mechanized realizations, receipts,
and a proof at textbook detail; the coverage dashboard up front is the
catalogue's honest map of what the assurance edifice does and does not
yet cover. Bernstein Logic appears here as what it has become in the
finished design: the core fragment of one catalogue, not a separate
system.

Inside this chapter the catalogue keeps its own conventions: one rule
per entry, each starting on a fresh page, under stable catalogue rule
numbers (indices of the roster,
not section numbers --- they also appear as PDF bookmarks), its own
references at the end, and the page geometry it was designed at.

\newgeometry{left=22mm,right=18mm,top=20mm,bottom=22mm}


\begingroup
\setlength{\parindent}{0pt}
\setlength{\parskip}{4pt}

\textbf{What this catalogue is.}\enspace The proof rules of the
certificate calculus, in two strata sharing one page discipline. The
base stratum is \textbf{Bernstein Logic} --- a small proof calculus for
certificates about polynomial facts over networks of
rational-tensor-valued wires: sequents $\Phi \vdash \Psi$ asserting that
every admissible valuation of the wires satisfying the premise $\Phi$
satisfies the conclusion $\Psi$, with all trusted arithmetic exact over
$\Q$. Its rules appear here in three tiers --- core rules (kernel
semantic content), admissible rules (verified unpackings into core
rules), and alias rules (verified single-application spellings) ---
together with the composition rules (\rulename{Promise},
\rulename{Syzygy}, \rulename{SmallGain}, \rulename{SmallGainN},
\rulename{ProductLattice}),
which are Bernstein Logic outright. Above the base stratum sit the
\textbf{QTSL rules}: rules concluding satisfaction judgments of the
quantitative temporal specification language --- a system satisfies a
specification built from the specification chapter\textquotesingle s
formers --- through the battery rules (safety, deadline, until, and
Streett acceptance), the judgment-structure and specification-surface
rules, the Morse-topology rules, composition and aggregation,
distance transport, and the stored-chain certificate species
(eventual-absorption brackets and generator-form contraction tails,
recomputed exactly over stored chain rows), supported by discharge
kits and arithmetic leaves
that remain ordinary Bernstein Logic sequents. QTSL premises bottom out
in Bernstein Logic sequents, and both strata are designed for one
exact-replay discipline: a \emph{certificate} is an acyclic derivation
graph (judgment premises may be shared) whose nodes carry enough
payload that a small checker can replay each rule\textquotesingle s
side conditions by exact rational computation --- no search, no
optimization, no floating point in the trusted base. Where a
rule\textquotesingle s row and checker arm are enacted, the checker
replays it; the coverage dashboard names, rule by rule, what is enacted
and what is not.
Producers (LP/SDP pipelines, abstract interpreters, hand
constructions) do whatever clever thing \emph{finds} the certificate;
the checker only replays. The catalogue opens with a coverage
dashboard --- one row per rule, covered or not, mapping the current
level of mechanized verification of soundness, honest absences
included --- then a synopsis parade of every rule, then the rules
grouped by subject.

The validity judgment itself is not free-floating: it is
kernel-checked to be the possibilistic face of the semantic
universe\textquotesingle s satisfaction judgment
(chapter~\ref{ch:qtsl-semantic-universe}), at the sequent\textquotesingle s
exact valuation carrier --- a sequent is semantically valid iff the
upper probability of its conclusion\textquotesingle s complement
vanishes under the credal compile of its premise
cylinder\textquotesingle s possibilistic wall
(\texttt{semanticallyValid\_iff\_face}, module
\texttt{SuFunctor.Plain.CheckerFace}) --- so each soundness theorem
in this catalogue is, through that identification, a statement about
the semantic universe\textquotesingle s own judgment.

Each rule of the calculus gets one section here, carrying five
ingredients. (i)~The rule itself, typeset in Steele\textquotesingle s
Computer Science Metanotation~\cite{steele2017}, with prose on what the
rule does, when a producer uses it, and why each hypothesis is there.
(ii)~The Lean~4 declarations that realize it --- the payload structure,
the soundness theorem, and (where a side condition\textquotesingle s
meaning must be visible) the checker definitions in full --- extracted
mechanically from the sources at every build of this catalogue and bound by
hash, so quoted text cannot drift from the tree silently. (iii)~Verifier
receipts: attestations that named checking layers ---
elaboration, kernel replay by the official external checker, a hardened
screen, and a decorrelated second kernel --- accepted the \emph{exact
build} being quoted, identified by the SHA-256 of the compiled module;
every layer is a program the reader can re-run over the quoted sources.
(iv)~The HOL4 statement where the independent HOL4 development carries
one, and a loud, reasoned absence where it does not. (v)~A proof at
textbook detail: ordinary mathematics for core rules, explicit rewriting
recipes for admissible rules, filled-in applications of the target rule
for alias rules. Where a rule\textquotesingle s certificate content is
carried by the baked theory, the section additionally quotes the theory
rows --- payload schema, side-condition demand relations,
well-formedness axioms --- mechanically from the theory source at every
build, so a renamed or removed row fails the build rather than drifting
silently.

\textbf{The theory as source.}\enspace The normative home of the
rules\textquotesingle{} certificate content is the baked theory pair,
shipped as data in \texttt{theories/} beside this book:
\texttt{theory BernsteinLogic} (the base calculus\textquotesingle s
proof tower, from ground dtypes through hypernets, formulas, and
judgments to the proof species) and \texttt{theory Qtsl} (one omnibus
extension carrying the specification language, the System stratum, the
flavored cell schemas, the certificate sections --- judgment kinds,
proof-node payloads, batteries, the verdict form --- and the
certificate-shape stratum). Where this catalogue\textquotesingle s
statements and the theory disagree, the theory is the authority. Not
every catalogued rule has a baked row home: entries whose certificate
content is out of scope of the current baked theory say so outright ---
the rule remains catalogued, tiered, and proved at its stated scope,
and its schema realization is named as a successor rather than
implied.

\textbf{The certificate-shape stratum.}\enspace One stratum of the
baked theory types certificates \emph{as artifacts} rather than as
derivations, and its consumers are the battery rules of this catalogue,
so it is documented here. A \emph{nondeterminism axis}
(\texttt{NDAxis}) is an entity --- identity is the choice source ---
carrying a total binding time and extent: \texttt{btStep} (re-chosen at
every step; the axis enters drift reads cellwise as a region axis, and
no envelope object is ever materialized) or \texttt{btRun} (committed
once per run), with extent \texttt{ndEnum} (an enumerated branch
count) or \texttt{ndParam} (a box enclosure of a parametric domain). A
\emph{certificate family} (\texttt{CertFam}) declares its own axis
roster (\texttt{cert\_axis}); its branch-key space is the product of
its enumerated run-bound axes in position order, keys riding the
ground \texttt{Shape} sort as natural tuples --- the scalar certificate
is the nullary case (one branch under the width-zero key, not a second
sort). Two refusals shape the roster: a family may not split over a
step-bound axis, and may not enumerate a parametric one. Per-branch
positional bodies (\texttt{cert\_body}) carry the certificate
functions; position semantics belongs to the consuming rule. The
existential face\textquotesingle s slot is the id-keyed
\emph{selection} (\texttt{cert\_selection}: a named branch plus a
properness witness; \texttt{cert\_sel\_cell} refines a parametric
domain, zero cell rows reading as the whole domain) --- selections are
mandatory by default, waived only when every consuming conclusion is
safety-only. The universal face\textquotesingle s economies are
branch-discharge rows: an emptiness certificate
(\texttt{cert\_branch\_empty}, the one optional slot),
grade domination (\texttt{cert\_dominates}, irreflexive and acyclic),
and progressive-revelation hand-off
(\texttt{cert\_handoff}, prefix agreement between branch certificates).
The stratum\textquotesingle s structural demand is counting and ground
comparison --- admission, never geometry --- and its deliberate
absences are testable negatives: no envelope slot (a demonic envelope
is a consuming rule\textquotesingle s cellwise reading of the step
axes, never an interface object), no verdict rows (those live with the
certificate sections that consume them), no checker-arm vocabulary,
and no branch-key structure sort.

\textbf{Why trust a green cell.}\enspace A green receipt row reduces to
three things a skeptical reader can check independently: the statement
printed on the page (hash-bound to the source file it was extracted
from); the SHA-256 of the compiled artefact the receipt names; and the
checker layers the receipt credits --- each layer an independently
runnable program (elaboration, official kernel replay, a hardened
screen, a decorrelated second kernel), so a green cell never asks for
trust in this book or its receipt store: rebuild the module, re-run the
checkers, and compare.
Staleness is always visible --- a receipt attesting an older build, a
statement that moved since its proof was written, a missing ingredient
each render as an \emph{amber} badge, never as silence. Red is reserved
for outright verification failure. The coverage dashboard below records,
rule by rule, which ingredients exist and which do not; the empty cells
are deliberate, and are meant to be read.

\textbf{Naming conventions.}\enspace Judgment kinds carry a \texttt{Q}
prefix (\texttt{QSat}, \texttt{QRegionSat}, \ldots): they are the
satisfaction-judgment ($\vDash$-side) sorts of the extension theory, and
the prefix is their namespace. Rules that conclude a satisfaction
judgment carry the same \texttt{Q}; derivation-side ($\vdash$) leaf rules
stay bare or carry \texttt{P}, and mathematical eponyms are never forced
into descriptive molds. Rule names end in a species suffix where one
applies: \emph{-Battery} --- a certificate kit concluding a region
judgment; \emph{-Discharge} --- evidence concluding a named judgment
kind; \emph{-Compose} --- transport across a composition axis, named by
the axis; \emph{-Alias} --- a verified one-application spelling;
\emph{-Rewrite} --- an equivalence-concluding rewrite. A few rules carry
a former name (marked ``n\'e'') from earlier editions of this catalogue.


\par\medskip
{\footnotesize\textcolor{NoteGray}{Rendered 2026-08-04 03:24 UTC \textperiodcentered{} receipts consulted from the store (checker-layer attestations) \textperiodcentered{} body first committed in source revision \texttt{6f92a81} (2026-08-04 12:11 UTC), not re-rendered at that revision}}

\clearpage
\pdfbookmark[1]{Coverage dashboard}{blcat-dashboard}\dashboardhead{One table over every rule of the catalogue, whether or not this edition
covers it yet: a coverage dashboard that maps the current level of
actual mechanized verification of soundness, rule by rule. Each covered
section joins its ingredients from
their single sources of truth: the rule cell from the hand-authored rules
source (the same file the two-column diff pages render), the Lean
statements from the sources by declaration name at every render, the
receipts from the store (cited at render,
never re-run), the HOL4 statement from the HOL4 development or a reasoned
absence, and the proof from its hand-authored file. Rules with no catalogue
entry yet render as all-absent rows, loudly. Some rows carry a Lean
column without a catalogue section yet: their soundness theorems exist in the
mechanized development (two sibling Lean packages --- the core calculus
and an auxiliary measure-theoretic library --- both statement-bound by
the same hash machinery), while their proof and HOL columns remain
marked absent until those sections are authored. A module whose build has
no attestation yet shows its receipt layers as none: statements
are source-bound, and receipt coverage arrives when the checking gate
next runs over that package.
}
\dashrow{\hyperref[blcatrule-Ax]{Ax}~\dashsrc{BL}}{core}{\dok}{\dok}{\lcur{b}\lstale{r}\lstale{h}\lcur{d}}{\dok}{\ddoc}{\dok}
\dashrow{\hyperref[blcatrule-Weaken]{Weaken}~\dashsrc{BL}}{core}{\dok}{\dok}{\lcur{b}\lstale{r}\lstale{h}\lcur{d}}{\dok}{\ddoc}{\dok}
\dashrow{\hyperref[blcatrule-Cut]{Cut}~\dashsrc{BL}}{core}{\dok}{\dok}{\lcur{b}\lstale{r}\lstale{h}\lcur{d}}{\dok}{\ddoc}{\dok}
\dashrow{\hyperref[blcatrule-AndI]{$\wedge$I}~\dashsrc{BL}}{core}{\dok}{\dok}{\lcur{b}\lstale{r}\lstale{h}\lcur{d}}{\dok}{\ddoc}{\dok}
\dashrow{\hyperref[blcatrule-AndE]{$\wedge$E}~\dashsrc{BL}}{core}{\dok}{\dok}{\lcur{b}\lstale{r}\lstale{h}\lcur{d}}{\dok}{\ddoc}{\dok}
\dashrow{\hyperref[blcatrule-OrI]{$\vee$I}~\dashsrc{BL}}{core}{\dok}{\dok}{\lcur{b}\lstale{r}\lstale{h}\lcur{d}}{\dok}{\ddoc}{\dok}
\dashrow{\hyperref[blcatrule-OrE]{$\vee$E}~\dashsrc{BL}}{core}{\dok}{\dok}{\lcur{b}\lcur{r}\lcur{h}\lcur{d}}{\dok}{\ddoc}{\dok}
\dashrow{\hyperref[blcatrule-EqElim]{EqElim}~\dashsrc{BL}}{core}{\dok}{\dok}{\lcur{b}\lstale{r}\lstale{h}\lcur{d}}{\dok}{\ddoc}{\dok}
\dashrow{\hyperref[blcatrule-EqRefl]{EqRefl}~\dashsrc{BL}}{core}{\dok}{\dok}{\lcur{b}\lstale{r}\lstale{h}\lcur{d}}{\dok}{\ddoc}{\dok}
\dashrow{\hyperref[blcatrule-EqSym]{EqSym}~\dashsrc{BL}}{core}{\dok}{\dok}{\lcur{b}\lstale{r}\lstale{h}\lcur{d}}{\dok}{\ddoc}{\dok}
\dashrow{\hyperref[blcatrule-Antisym]{Antisym}~\dashsrc{BL}}{core}{\dok}{\dok}{\lcur{b}\lstale{r}\lstale{h}\lcur{d}}{\dok}{\ddoc}{\dok}
\dashrow{\hyperref[blcatrule-ConvElim]{ConvElim}~\dashsrc{BL}}{core}{\dok}{\dok}{\lcur{b}\lstale{r}\lstale{h}\lcur{d}}{\dok}{\ddoc}{\dok}
\dashrow{\hyperref[blcatrule-ConvIntro]{ConvIntro}~\dashsrc{BL}}{core}{\dok}{\dok}{\lcur{b}\lstale{r}\lstale{h}\lcur{d}}{\dok}{\ddoc}{\dok}
\dashrow{\hyperref[blcatrule-ArithAtom]{ArithAtom}~\dashsrc{BL}}{core}{\dok}{\dok}{\lcur{b}\lstale{r}\lstale{h}\lcur{d}}{\dok}{\ddoc}{\dok}
\dashrow{\hyperref[blcatrule-BernCone]{BernCone}~\dashsrc{BL}}{core}{\dok}{\damber}{\lcur{b}\lstale{r}\lstale{h}\lcur{d}}{\dok}{\ddoc}{\dok}
\dashrow{\hyperref[blcatrule-Syzygy]{Syzygy}~\dashsrc{BL}}{core}{\dok}{\dok}{\lcur{b}\lstale{r}\lstale{h}\lcur{d}}{\dok}{\ddoc}{\dok}
\dashrow{\hyperref[blcatrule-QGuardIntro]{QGuardIntro}~\dashsrc{QTSL}}{core}{\dok}{\dok}{\lcur{b}\lnone{r}\lnone{h}\lcur{d}}{\dnone}{\dnone}{\dok}
\dashrow{\hyperref[blcatrule-QGuardElim]{QGuardElim}~\dashsrc{QTSL}}{core}{\dok}{\dok}{\lcur{b}\lnone{r}\lnone{h}\lcur{d}}{\dnone}{\dnone}{\dok}
\dashrow{\hyperref[blcatrule-QSafetyBattery]{QSafetyBattery}~\dashsrc{QTSL}}{core}{\dok}{\damber}{\lcur{b}\lcur{r}\lcur{h}\lcur{d}}{\dnone}{\dnone}{\dok}
\dashrow{\hyperref[blcatrule-QDeadlineBattery]{QDeadlineBattery}~\dashsrc{QTSL}}{core}{\dok}{\damber}{\lcur{b}\lstale{r}\lstale{h}\lcur{d}}{\dnone}{\dnone}{\dok}
\dashrow{\hyperref[blcatrule-QUnionBound]{QUnionBound}~\dashsrc{QTSL}}{core}{\dok}{\damber}{\lcur{b}\lstale{r}\lstale{h}\lcur{d}}{\dnone}{\dnone}{\dok}
\dashrow{\hyperref[blcatrule-QStreettBattery]{QStreettBattery}~\dashsrc{QTSL}}{core}{\dok}{\dok}{\lcur{b}\lcur{r}\lcur{h}\lcur{d}}{\dnone}{\dnone}{\dok}
\dashrow{\hyperref[blcatrule-QAvgBattery]{QAvgBattery}~\dashsrc{QTSL}}{core}{\dnone}{\dnone}{\lnone{b}\lnone{r}\lnone{h}\lnone{d}}{\dnone}{\dnone}{\dok}
\dashrow{\hyperref[blcatrule-QMorseDischarge]{QMorseDischarge}~\dashsrc{QTSL}}{core}{\dcond}{\damber}{\lcur{b}\lcur{r}\lcur{h}\lcur{d}}{\dnone}{\dnone}{\dok}
\dashrow{\hyperref[blcatrule-QMorseProduct]{QMorseProduct}~\dashsrc{QTSL}}{core}{\dok}{\dnone}{\lcur{b}\lcur{r}\lcur{h}\lcur{d}}{\dnone}{\dnone}{\dok}
\dashrow{\hyperref[blcatrule-QBlockInv]{QBlockInv}~\dashsrc{QTSL}}{core}{\dcond}{\damber}{\lcur{b}\lcur{r}\lcur{h}\lcur{d}}{\dok}{\ddoc}{\dok}
\dashrow{\hyperref[blcatrule-QBasin]{QBasin}~\dashsrc{QTSL}}{core}{\dok}{\damber}{\lcur{b}\lstale{r}\lstale{h}\lcur{d}}{\dnone}{\dnone}{\dok}
\dashrow{\hyperref[blcatrule-QZeroByMiranda]{QZeroByMiranda}~\dashsrc{QTSL}}{core}{\dnone}{\dnone}{\lnone{b}\lnone{r}\lnone{h}\lnone{d}}{\dnone}{\dnone}{\dok}
\dashrow{\hyperref[blcatrule-QCoverByMiranda]{QCoverByMiranda}~\dashsrc{QTSL}}{core}{\dnone}{\dnone}{\lnone{b}\lnone{r}\lnone{h}\lnone{d}}{\dnone}{\dnone}{\dok}
\dashrow{\hyperref[blcatrule-QWireCompose]{QWireCompose}~\dashsrc{QTSL}}{core}{\dcond}{\damber}{\lcur{b}\lcur{r}\lcur{h}\lcur{d}}{\dnone}{\dnone}{\dok}
\dashrow{\hyperref[blcatrule-QWireCertDischarge]{QWireCertDischarge}~\dashsrc{QTSL}}{core}{\dcond}{\dok}{\lcur{b}\lcur{r}\lcur{h}\lcur{d}}{\dnone}{\dnone}{\dok}
\dashrow{\hyperref[blcatrule-QWireComposeMarks]{QWireComposeMarks}~\dashsrc{QTSL}}{core}{\dcond}{\damber}{\lstale{b}\lnone{r}\lnone{h}\lstale{d}}{\dnone}{\dnone}{\dok}
\dashrow{\hyperref[blcatrule-QSimTransport]{QSimTransport}~\dashsrc{QTSL}}{core}{\dok}{\dok}{\lcur{b}\lcur{r}\lcur{h}\lcur{d}}{\dnone}{\dnone}{\dok}
\dashrow{\hyperref[blcatrule-QLinGrowthDischarge]{QLinGrowthDischarge}~\dashsrc{QTSL}}{core}{\dnone}{\dnone}{\lnone{b}\lnone{r}\lnone{h}\lnone{d}}{\dnone}{\dnone}{\dok}
\dashrow{\hyperref[blcatrule-QMarginalCompat]{QMarginalCompat}~\dashsrc{QTSL}}{core}{\dnone}{\dnone}{\lnone{b}\lnone{r}\lnone{h}\lnone{d}}{\dnone}{\dnone}{\dok}
\dashrow{\hyperref[blcatrule-QAggReduce]{QAggReduce}~\dashsrc{QTSL}}{core}{\dcond}{\dok}{\lcur{b}\lcur{r}\lcur{h}\lcur{d}}{\dnone}{\dnone}{\dok}
\dashrow{\hyperref[blcatrule-QPAlwaysFromFold]{QPAlwaysFromFold}~\dashsrc{QTSL}}{core}{\dnone}{\dnone}{\lnone{b}\lnone{r}\lnone{h}\lnone{d}}{\dnone}{\dnone}{\dok}
\dashrow{\hyperref[blcatrule-QFireBoundFromFold]{QFireBoundFromFold}~\dashsrc{QTSL}}{core}{\dnone}{\dnone}{\lnone{b}\lnone{r}\lnone{h}\lnone{d}}{\dnone}{\dnone}{\dok}
\dashrow{\hyperref[blcatrule-QPAlwaysTFromFold]{QPAlwaysTFromFold}~\dashsrc{QTSL}}{core}{\dnone}{\dnone}{\lnone{b}\lnone{r}\lnone{h}\lnone{d}}{\dnone}{\dnone}{\dok}
\dashrow{\hyperref[blcatrule-QRateFromClock]{QRateFromClock}~\dashsrc{QTSL}}{core}{\dnone}{\dnone}{\lnone{b}\lnone{r}\lnone{h}\lnone{d}}{\dnone}{\dnone}{\dok}
\dashrow{\hyperref[blcatrule-QPAlwaysTFromLineage]{QPAlwaysTFromLineage}~\dashsrc{QTSL}}{core}{\dnone}{\dnone}{\lnone{b}\lnone{r}\lnone{h}\lnone{d}}{\dnone}{\dnone}{\dok}
\dashrow{\hyperref[blcatrule-QCoupleDrift]{QCoupleDrift}~\dashsrc{QTSL}}{core}{\dcond[finite]}{\dok}{\lcur{b}\lcur{r}\lcur{h}\lcur{d}}{\dnone}{\dnone}{\dok}
\dashrow{\hyperref[blcatrule-QFixBall]{QFixBall}~\dashsrc{QTSL}}{core}{\dok}{\dok}{\lcur{b}\lcur{r}\lcur{h}\lcur{d}}{\dnone}{\dnone}{\dok}
\dashrow{\hyperref[blcatrule-QDistTransport]{QDistTransport}~\dashsrc{QTSL}}{core}{\dcond}{\dok}{\lcur{b}\lcur{r}\lcur{h}\lcur{d}}{\dnone}{\dnone}{\dok}
\dashrow{\hyperref[blcatrule-QInvRelativize]{QInvRelativize}~\dashsrc{QTSL}}{core}{\dok}{\dnone}{\lcur{b}\lstale{r}\lstale{h}\lcur{d}}{\dok}{\ddoc}{\dok}
\dashrow{\hyperref[blcatrule-QFirstService]{QFirstService}~\dashsrc{QTSL}}{core}{\dok}{\dnone}{\lcur{b}\lstale{r}\lstale{h}\lcur{d}}{\dok}{\ddoc}{\dok}
\dashrow{\hyperref[blcatrule-QFirstServiceSplit]{QFirstServiceSplit}~\dashsrc{QTSL}}{core}{\dok}{\dnone}{\lcur{b}\lstale{r}\lstale{h}\lcur{d}}{\dok}{\ddoc}{\dok}
\dashrow{\hyperref[blcatrule-DtmcAbsorbBracket]{DtmcAbsorbBracket}~\dashsrc{QTSL}}{core}{\dnone}{\dnone}{\lnone{b}\lnone{r}\lnone{h}\lnone{d}}{\dnone}{\dnone}{\dok}
\dashrow{\hyperref[blcatrule-DtmcPosteriorBracket]{DtmcPosteriorBracket}~\dashsrc{QTSL}}{core}{\dnone}{\dnone}{\lnone{b}\lnone{r}\lnone{h}\lnone{d}}{\dnone}{\dnone}{\dok}
\dashrow{\hyperref[blcatrule-QGenFormContraction]{QGenFormContraction}~\dashsrc{QTSL}}{core}{\dnone}{\dnone}{\lnone{b}\lnone{r}\lnone{h}\lnone{d}}{\dnone}{\dnone}{\dok}
\dashrow{\hyperref[blcatrule-Norm]{Norm}~\dashsrc{BL}}{adm.}{\dok}{\dok}{\lcur{b}\lstale{r}\lstale{h}\lcur{d}}{\dok}{\ddoc}{\dok}
\dashrow{\hyperref[blcatrule-Farkas]{Farkas}~\dashsrc{BL}}{adm.}{\dok}{\damber}{\lcur{b}\lstale{r}\lstale{h}\lcur{d}}{\dok}{\ddoc}{\dok}
\dashrow{\hyperref[blcatrule-Subdivide]{Subdivide}~\dashsrc{BL}}{adm.}{\dok}{\dok}{\lcur{b}\lstale{r}\lstale{h}\lcur{d}}{\dok}{\ddoc}{\dok}
\dashrow{\hyperref[blcatrule-LiRPA]{LiRPA}~\dashsrc{BL}}{adm.}{\dok}{\dok}{\lcur{b}\lcur{r}\lcur{h}\lcur{d}}{\dok}{\ddoc}{\dok}
\dashrow{\hyperref[blcatrule-ReLUPos]{ReLU-Pos}~\dashsrc{BL}}{adm.}{\dok}{\dok}{\lcur{b}\lstale{r}\lstale{h}\lcur{d}}{\dok}{\ddoc}{\dok}
\dashrow{\hyperref[blcatrule-ReLUNeg]{ReLU-Neg}~\dashsrc{BL}}{adm.}{\dok}{\dok}{\lcur{b}\lstale{r}\lstale{h}\lcur{d}}{\dok}{\ddoc}{\dok}
\dashrow{\hyperref[blcatrule-ReLUNN]{ReLU-NN}~\dashsrc{BL}}{adm.}{\dok}{\damber}{\lcur{b}\lstale{r}\lstale{h}\lcur{d}}{\dok}{\ddoc}{\dok}
\dashrow{\hyperref[blcatrule-ReLUAbove]{ReLU-Above}~\dashsrc{BL}}{adm.}{\dok}{\dok}{\lcur{b}\lstale{r}\lstale{h}\lcur{d}}{\dok}{\ddoc}{\dok}
\dashrow{\hyperref[blcatrule-CellInInterior]{Cell-In-Interior}~\dashsrc{BL}}{adm.}{\dok}{\dok}{\lcur{b}\lcur{r}\lcur{h}\lcur{d}}{\dok}{\ddoc}{\dok}
\dashrow{\hyperref[blcatrule-CellOut]{Cell-Out}~\dashsrc{BL}}{adm.}{\dok}{\dok}{\lcur{b}\lstale{r}\lstale{h}\lcur{d}}{\dok}{\ddoc}{\dok}
\dashrow{\hyperref[blcatrule-PConstExact]{PConstExact}~\dashsrc{BL}}{adm.}{\dok}{\dok}{\lcur{b}\lcur{r}\lcur{h}\lcur{d}}{\dok}{\ddoc}{\dok}
\dashrow{\hyperref[blcatrule-PsdCertTri]{PsdCertTri}~\dashsrc{BL}}{adm.}{\dok}{\dok}{\lcur{b}\lcur{r}\lcur{h}\lcur{d}}{\dok}{\ddoc}{\dok}
\dashrow{\hyperref[blcatrule-PsdCertDiag]{PsdCertDiag}~\dashsrc{BL}}{adm.}{\dok}{\dok}{\lcur{b}\lcur{r}\lcur{h}\lcur{d}}{\dok}{\ddoc}{\dok}
\dashrow{\hyperref[blcatrule-PsdCertKernel]{PsdCertKernel}~\dashsrc{BL}}{adm.}{\dok}{\dok}{\lcur{b}\lcur{r}\lcur{h}\lcur{d}}{\dok}{\ddoc}{\dok}
\dashrow{\hyperref[blcatrule-ReLUSplit]{ReLU-Split}~\dashsrc{BL}}{adm.}{\dok}{\dnone}{\lcur{b}\lcur{r}\lcur{h}\lcur{d}}{\dnone}{\dnone}{\dok}
\dashrow{\hyperref[blcatrule-Promise]{Promise}~\dashsrc{BL}}{adm.}{\dok}{\dok}{\lcur{b}\lstale{r}\lstale{h}\lcur{d}}{\dok}{\ddoc}{\dok}
\dashrow{\hyperref[blcatrule-ProductLattice]{ProductLattice}~\dashsrc{BL}}{adm.}{\dok}{\dok}{\lcur{b}\lcur{r}\lcur{h}\lcur{d}}{\dnone}{\dnone}{\dok}
\dashrow{\hyperref[blcatrule-QUntilBattery]{QUntilBattery}~\dashsrc{QTSL}}{adm.}{\dok}{\damber}{\lcur{b}\lstale{r}\lstale{h}\lcur{d}}{\dnone}{\dnone}{\dok}
\dashrow{\hyperref[blcatrule-QStreettCmplRequest]{QStreettCmplRequest}~\dashsrc{QTSL}}{adm.}{\dcond}{\dok}{\lcur{b}\lcur{r}\lcur{h}\lcur{d}}{\dnone}{\dnone}{\dok}
\dashrow{\hyperref[blcatrule-HypMultiplier]{HypMultiplier}~\dashsrc{QTSL}}{adm.}{\dok}{\dok}{\lcur{b}\lcur{r}\lcur{h}\lcur{d}}{\dnone}{\dnone}{\dok}
\dashrow{\hyperref[blcatrule-QVerifiedCover]{QVerifiedCover}~\dashsrc{QTSL}}{adm.}{\dok}{\dok}{\lcur{b}\lcur{r}\lcur{h}\lcur{d}}{\dnone}{\dnone}{\dok}
\dashrow{\hyperref[blcatrule-QPositiveCombinationRoot]{QPositiveCombinationRoot}~\dashsrc{QTSL}}{adm.}{\dok}{\dok}{\lcur{b}\lcur{r}\lcur{h}\lcur{d}}{\dnone}{\dnone}{\dok}
\dashrow{\hyperref[blcatrule-QClosureOutward]{QClosureOutward}~\dashsrc{QTSL}}{adm.}{\dok}{\dok}{\lcur{b}\lcur{r}\lcur{h}\lcur{d}}{\dnone}{\dnone}{\dok}
\dashrow{\hyperref[blcatrule-QClosureInwardEps]{QClosureInwardEps}~\dashsrc{QTSL}}{adm.}{\dok}{\dok}{\lcur{b}\lcur{r}\lcur{h}\lcur{d}}{\dnone}{\dnone}{\dok}
\dashrow{\hyperref[blcatrule-BoxBern]{BoxBern}~\dashsrc{QTSL}}{adm.}{\dok}{\dok}{\lcur{b}\lstale{r}\lstale{h}\lcur{d}}{\dnone}{\dnone}{\dok}
\dashrow{\hyperref[blcatrule-PExpTypeEnclosure]{PExpTypeEnclosure}~\dashsrc{QTSL}}{adm.}{\dok}{\dok}{\lcur{b}\lcur{r}\lcur{h}\lcur{d}}{\dok}{\ddoc}{\dok}
\dashrow{\hyperref[blcatrule-QEquilibriumFromInward]{QEquilibriumFromInward}~\dashsrc{QTSL}}{adm.}{\dnone}{\dnone}{\lnone{b}\lnone{r}\lnone{h}\lnone{d}}{\dnone}{\dnone}{\dok}
\dashrow{\hyperref[blcatrule-QFlatten]{QFlatten}~\dashsrc{QTSL}}{adm.}{\dok}{\dok}{\lcur{b}\lcur{r}\lcur{h}\lcur{d}}{\dnone}{\dnone}{\dok}
\dashrow{\hyperref[blcatrule-PBound]{PBound}~\dashsrc{BL}}{alias}{\dok}{\dok}{\lcur{b}\lstale{r}\lstale{h}\lcur{d}}{\dok}{\ddoc}{\dok}
\dashrow{\hyperref[blcatrule-SmallGain]{SmallGain}~\dashsrc{BL}}{alias}{\dok}{\dok}{\lcur{b}\lstale{r}\lstale{h}\lcur{d}}{\dscope}{\dnone}{\dok}
\dashrow{\hyperref[blcatrule-SmallGainN]{SmallGainN}~\dashsrc{BL}}{alias}{\dok}{\dnone}{\lcur{b}\lstale{r}\lstale{h}\lcur{d}}{\dscope}{\dnone}{\dok}
\dashrow{\hyperref[blcatrule-QSatFromRegion]{QSatFromRegion}~\dashsrc{QTSL}}{alias}{\dok}{\dok}{\lcur{b}\lcur{r}\lcur{h}\lcur{d}}{\dnone}{\dnone}{\dok}
\dashrow{\hyperref[blcatrule-QStreettDegenerate]{QStreettDegenerate}~\dashsrc{QTSL}}{alias}{\dok}{\dok}{\lcur{b}\lcur{r}\lcur{h}\lcur{d}}{\dscope}{\dnone}{\dok}
\dashrow{\hyperref[blcatrule-QStreettTargetCellCollapse]{QStreettTargetCellCollapse}~\dashsrc{QTSL}}{alias}{\dok}{\dok}{\lcur{b}\lcur{r}\lcur{h}\lcur{d}}{\dscope}{\dnone}{\dok}
\dashrow{\hyperref[blcatrule-QNnfRewrite]{QNnfRewrite}~\dashsrc{QTSL}}{alias}{\dok}{\dok}{\lcur{b}\lcur{r}\lcur{h}\lcur{d}}{\dscope}{\dnone}{\dok}
\dashrow{\hyperref[blcatrule-QStreettRequestNnf]{QStreettRequestNnf}~\dashsrc{QTSL}}{alias}{\dok}{\dok}{\lcur{b}\lcur{r}\lcur{h}\lcur{d}}{\dscope}{\dnone}{\dok}
\dashrow{\hyperref[blcatrule-QAvoid]{QAvoid}~\dashsrc{QTSL}}{alias}{\dok}{\dok}{\lcur{b}\lcur{r}\lcur{h}\lcur{d}}{\dscope}{\dnone}{\dok}
\dashrow{\hyperref[blcatrule-QStreettGF]{QStreettGF}~\dashsrc{QTSL}}{alias}{\dok}{\dok}{\lcur{b}\lcur{r}\lcur{h}\lcur{d}}{\dscope}{\dnone}{\dok}
\dashrow{\hyperref[blcatrule-QStreettFG]{QStreettFG}~\dashsrc{QTSL}}{alias}{\dok}{\dok}{\lcur{b}\lcur{r}\lcur{h}\lcur{d}}{\dscope}{\dnone}{\dok}
\dashrow{\hyperref[blcatrule-QNegFAlias]{QNegFAlias}~\dashsrc{QTSL}}{alias}{\dok}{\dok}{\lcur{b}\lcur{r}\lcur{h}\lcur{d}}{\dscope}{\dnone}{\dok}
\dashrow{\hyperref[blcatrule-QFullDomainWeaken]{QFullDomainWeaken}~\dashsrc{QTSL}}{alias}{\dok}{\dok}{\lcur{b}\lcur{r}\lcur{h}\lcur{d}}{\dnone}{\dnone}{\dok}
\dashrow{\hyperref[blcatrule-QClosureExact]{QClosureExact}~\dashsrc{QTSL}}{alias}{\dok}{\dok}{\lcur{b}\lcur{r}\lcur{h}\lcur{d}}{\dnone}{\dnone}{\dok}
\dashrow{\hyperref[blcatrule-TaylorBound]{TaylorBound}~\dashsrc{QTSL}}{alias}{\dok}{\dok}{\lcur{b}\lcur{r}\lcur{h}\lcur{d}}{\dscope}{\dnone}{\dok}
\dashrow{\hyperref[blcatrule-RootEval]{RootEval}~\dashsrc{QTSL}}{alias}{\dok}{\dok}{\lcur{b}\lcur{r}\lcur{h}\lcur{d}}{\dscope}{\dnone}{\dok}
\dashrow{\hyperref[blcatrule-QLatticeInv]{QLatticeInv}~\dashsrc{QTSL}}{alias}{\dok}{\dok}{\lcur{b}\lcur{r}\lcur{h}\lcur{d}}{\dnone}{\dnone}{\dok}
\dashrow{\hyperref[blcatrule-QCountAtomRounding]{QCountAtomRounding}~\dashsrc{QTSL}}{alias}{\dok}{\dok}{\lcur{b}\lcur{r}\lcur{h}\lcur{d}}{\dnone}{\dnone}{\dok}
\dashrow{\hyperref[blcatrule-ProofWellFormedness]{ProofWellFormedness}~\dashsrc{QTSL}}{glue}{\dcond}{\dnone}{\lcur{b}\lcur{r}\lcur{h}\lcur{d}}{\dscope}{\dnone}{\dok}
\dashboardfoot

\blcatsection{Rule Synopsis}
\sourcenote{Every rule of the catalogue at a glance: the rule as
realized in the mechanized development, its name, and the page of
its full entry. The closing entry, set off by a rule, is the
assembly-level discipline the numbered rules feed.}
\par\medskip
\paradeentry{\inferrule*{\;}{\Phi \vdash \Phi}}{Ax}{blcatrule-Ax}
\paradeentry{\inferrule*{\Phi \vdash \Psi}{\Phi \wedge \chi \vdash \Psi}}{Weaken}{blcatrule-Weaken}
\paradeentry{\inferrule*{\Phi \vdash \chi \\ \Phi \wedge \chi \vdash \Psi}{\Phi \vdash \Psi}}{Cut}{blcatrule-Cut}
\paradeentry{\inferrule*{\Phi \vdash \Psi_1 \\ \Phi \vdash \Psi_2}{\Phi \vdash \Psi_1 \wedge \Psi_2}}{$\wedge$I}{blcatrule-AndI}
\paradeentry{\inferrule*{\Phi \vdash \Psi_1 \wedge \Psi_2}{\Phi \vdash \Psi_i}}{$\wedge$E}{blcatrule-AndE}
\paradeentry{\inferrule*{\Phi \vdash \Psi_i}{\Phi \vdash \Psi_1 \vee \Psi_2}}{$\vee$I}{blcatrule-OrI}
\paradeentry{\inferrule*{\Phi_1 \vdash \Psi \\ \Phi_2 \vdash \Psi}{\Phi_1 \vee \Phi_2 \vdash \Psi}}{$\vee$E}{blcatrule-OrE}
\paradeentry{\inferrule*{\Phi \vdash t = u \\ \Phi \vdash \Psi\sub{a}{t} \\ \exists s.\; \bshape_H(t) = s = \bshape_H(u)}{\Phi \vdash \Psi\sub{a}{u}}}{EqElim}{blcatrule-EqElim}
\paradeentry{\inferrule*{\exists s.\; \bshape_H(t) = s \\ \wires(t) \text{ finitely typed in } H}{\Phi \vdash t = t}}{EqRefl}{blcatrule-EqRefl}
\paradeentry{\inferrule*{\Phi \vdash t = u}{\Phi \vdash u = t}}{EqSym}{blcatrule-EqSym}
\paradeentry{\inferrule*{\Phi \vdash t \leq u \\ \Phi \vdash t \geq u}{\Phi \vdash t = u}}{Antisym}{blcatrule-Antisym}
\paradeentry{\inferrule*{
  \Phi \vdash x \in \mathcal{R}(\Delta, S)
}{
  \Phi \vdash \textstyle\bigwedge_{i = 0}^{n} \mathrm{canon}(\epsilon_i\,\lambda_i(x)) \leq 0 \\ \epsilon_i \,=\, {+1} \text{ if } i \in S, \quad {-1} \text{ if } i \in \bar S
}}{ConvElim}{blcatrule-ConvElim}
\paradeentry{\inferrule*{
  \Phi \vdash \textstyle\bigwedge_{i = 0}^{n} \mathrm{canon}(\epsilon_i\,\lambda_i(x)) \leq 0 \\ \epsilon_i \,=\, {+1} \text{ if } i \in S, \ {-1} \text{ if } i \in \bar S \\ \bshape_H(x) = [\,n\,]
}{
  \Phi \vdash x \in \mathcal{R}(\Delta, S)
}}{ConvIntro}{blcatrule-ConvIntro}
\paradeblock{\begin{mathpar}
\inferrule*[Right=\rulename{ArithAtom}]{
  \Phi \vdash (t_1 - t_2) \bowtie 0
}{
  \Phi \vdash t_1 \bowtie t_2
}
\\
\inferrule*[Right=\rulename{ArithAtom}$^{-1}$]{
  \Phi \vdash t_1 \bowtie t_2
}{
  \Phi \vdash (t_1 - t_2) \bowtie 0
}
\end{mathpar}
with $\bowtie \in \{<, \leq, =, \geq, >\}$.}{ArithAtom}{blcatrule-ArithAtom}
\paradeentry{\inferrule*{
  \Delta = \mathrm{conv}(V_0, \ldots, V_n) \subset \R^n \text{ an } n\text{-simplex};
  \quad S \subsetneq \{0, 1, \ldots, n\}
  \\
  \Phi \vdash x \in \mathcal{R}(\Delta, S)
  \\
  p \colon \R^n \to W \text{ polynomial map of degree } d, \text{ with coefficient tensor } (b_\alpha)_{|\alpha|=d}
  \\
  \wires(p) = \emptyset \;\vee\; \bigl(x = w \,\wedge\, \wires(p) \subseteq \{w\}\bigr)
  \\
  S \neq \emptyset \Rightarrow P \text{ conic}; \qquad S = \emptyset \Rightarrow P \text{ convex suffices}
  \\
  \bigl(p \ \text{idx-free, with } p \ \text{constant} \vee n = 1\bigr) \;\vee\; \bigl(p \ \text{projection-indexed, indices} < n,\ S = \emptyset\bigr)
  \\
  \text{for every } j \in \N^S \text{ with } |j| \leq d, \text{ every degree-}d \text{ Bernstein coefficient of } P_j \text{ on } F_S \text{ satisfies } P
}{
  \Phi \vdash P(p(x))
}}{BernCone}{blcatrule-BernCone}
\paradeentry{\inferrule*{
  t_1, t_2 \ \text{polynomial terms in the same variables with equal monomial canonical forms over } \Q
  \\
  \bshape_H(t_1) = \bshape_H(t_2) \neq \bot \ \text{(a shared defined broadcast shape)}
  \\
  \wires(t_1), \wires(t_2) \ \text{all finitely typed in } H
}{
  \Phi \vdash t_1 = t_2
}}{Norm}{blcatrule-Norm}
\paradeentry{\inferrule*{
  \{\ell_j \bowtie_j 0\}_j \text{ a sub-conjunction of } \Phi \text{ with each } \ell_j \text{ linear, } \bowtie_j \in \{\leq, =\} \\
  \{\lambda_j \in \Q_{\geq 0}\}_j \text{ nonneg coefficients, with } \lambda_j \text{ free when } \bowtie_j \text{ is } {=} \\
  \bshape_H(\textstyle\sum_j \lambda_j \ell_j) \neq \bot \\
  \wires(\textstyle\sum_j \lambda_j \ell_j) \text{ all finitely typed in } H \\
  \wires(\Phi) \cup \wires(\textstyle\sum_j \lambda_j \ell_j) \text{ all typed at positive-size shapes in } H
}{
  \Phi \vdash \textstyle\sum_j \lambda_j \ell_j \leq 0
}}{Farkas}{blcatrule-Farkas}
\paradeentry{\inferrule*{
  \Phi \vdash x \in \mathcal{R}(\Delta, S)
  \\
  \{(\Delta_k, S_k)\}_{k=1}^N \text{ finite, with cover side-condition } \mathcal{R}(\Delta, S) \subseteq \textstyle\bigcup_k \mathcal{R}(\Delta_k, S_k)
  \\
  \forall k.\; \Phi \;\wedge\; (x \in \mathcal{R}(\Delta_k, S_k)) \;\vdash\; \Psi
}{
  \Phi \vdash \Psi
}}{Subdivide}{blcatrule-Subdivide}
\paradeentry{\inferrule*{
  \Phi \vdash a \leq h \\
  \Phi \vdash h \leq b \\
  \ReLU_{\{h\}\to\{y\}} \in H \\
  a < 0 \\
  0 < b \\\\
  \alpha a + \beta \leq 0 \\
  \beta \leq 0 \\
  \alpha b + \beta \leq b \\\\
  0 \leq \gamma a + \delta \\
  0 \leq \delta \\
  b \leq \gamma b + \delta
}{
  \Phi \vdash (\alpha h + \beta) \leq y \;\wedge\; y \leq (\gamma h + \delta)
}}{LiRPA}{blcatrule-LiRPA}
\paradeentry{\inferrule*{
  \Phi \vdash h \geq 0
  \\
  \ReLU_{\{h\}\to\{y\}} \in H
}{
  \Phi \vdash y = h
}}{ReLU-Pos}{blcatrule-ReLUPos}
\paradeentry{\inferrule*{
  \Phi \vdash h \leq 0
  \\
  \ReLU_{\{h\}\to\{y\}} \in H
}{
  \Phi \vdash y = 0
}}{ReLU-Neg}{blcatrule-ReLUNeg}
\paradeentry{\inferrule*{
  \ReLU_{\{h\}\to\{y\}} \in H
}{
  \Phi \vdash y \geq 0
}}{ReLU-NN}{blcatrule-ReLUNN}
\paradeentry{\inferrule*{
  \ReLU_{\{h\}\to\{y\}} \in H
}{
  \Phi \vdash y \geq h
}}{ReLU-Above}{blcatrule-ReLUAbove}
\paradeentry{\inferrule*{
  \Phi \vdash \bigwedge_{v} \lambda^{C_i}_v(x) \bowtie^{\circ}_v 0 \ \text{(strict barycentric interior atoms, flat vertex order; } < \text{ for extended } v\text{, } > \text{ for core } v\text{)}
  \\
  \Cell[\Sigma]_{\{x\}\to\{b_1,\ldots,b_N\}} \in H
}{
  \Phi \vdash b_i = 1 \wedge \bigwedge_{j \neq i} b_j = 0
}}{Cell-In-Interior}{blcatrule-CellInInterior}
\paradeentry{\inferrule*{
  \Phi \vdash \lambda^{C_j}_w(x) \bowtie^{\circ} 0 \ \text{(a single strict separating atom at a named witness vertex } w \text{ of } C_j\text{; } > \text{ for extended } w\text{, } < \text{ for core } w\text{)}
  \\
  \Cell[\Sigma]_{\{x\}\to\{b_1,\ldots,b_N\}} \in H
}{
  \Phi \vdash b_j = 0
}}{Cell-Out}{blcatrule-CellOut}
\paradeentry{\inferrule*{
  e_1, e_2 \text{ wire-free (closed) terms}
  \\
  \bowtie \in \{<, \leq, =, \geq, >\}
  \\
  e_1 \bowtie e_2 \ \text{holds under exact evaluation in } \Q
}{
  \Phi \vdash e_1 \bowtie e_2
}}{PConstExact}{blcatrule-PConstExact}
\paradeentry{\inferrule*{
  \bigl\{\Phi_i \vdash e_i \leq 0\bigr\}_{i=1}^{n}, \quad \Phi \supseteq \Phi_i \ \text{conjunct-wise (checked)}
  \\
  \bigl\{g_k = 0\bigr\}_k \ \text{equality sub-conjuncts of } \Phi
  \\
  \lambda_i \in \Q_{\geq 0}
  \\
  \nu_k \in \Q
  \\
  c \equiv \textstyle\sum_i \lambda_i e_i + \sum_k \nu_k g_k \ \text{(exact canonical polynomial)}
  \\
  \bshape_H(c) \neq \bot, \quad \wires(c) \ \text{all finitely typed in } H
  \\
  \wires(\Phi) \cup \wires(c) \cup \textstyle\bigcup_i \wires(e_i) \ \text{all typed at positive-size shapes in } H
  \\
  \Psi \in \bigl\{\, c \leq 0, \;\; e \geq 0 \text{ where } c = -e, \;\; l \leq r \text{ where } c = l - r \,\bigr\}
}{
  \Phi \vdash \Psi
}}{PBound}{blcatrule-PBound}
\paradeentry{\inferrule*{
  A \in \Q^{n \times n} \ \text{symmetric}
  \\
  Z \ \text{upper-triangular}, \ Z_{ii} > 0
  \\
  A \cdot Z \ \text{lower-triangular}, \ (A \cdot Z)_{ii} > 0 \ \text{(recomputed exactly)}
  \\
  p_1, \dots, p_n \ \text{payload terms}
  \\
  c \equiv -\textstyle\sum_{i,j} A_{ij}\, p_i\, p_j \ \text{(exact canonical polynomial)}
  \\
  \Psi \in \bigl\{\, c \leq 0, \;\; Q \geq 0 \text{ where } c = -Q, \;\; l \leq r \text{ where } c = l - r \,\bigr\}
}{
  \Phi \vdash \Psi
}}{PsdCertTri}{blcatrule-PsdCertTri}
\paradeentry{\inferrule*{
  A \in \Q^{n \times n} \ \text{symmetric}
  \\
  Z \ \text{upper-triangular}, \ Z_{ii} > 0
  \\
  B := Z^{\mathsf{T}} A Z \ \text{strictly row-diagonally dominant (recomputed exactly)}
  \\
  p_1, \dots, p_n \ \text{payload terms}
  \\
  c \equiv -\textstyle\sum_{i,j} A_{ij}\, p_i\, p_j \ \text{(exact canonical polynomial)}
  \\
  \Psi \in \bigl\{\, c \leq 0, \;\; Q \geq 0 \text{ where } c = -Q, \;\; l \leq r \text{ where } c = l - r \,\bigr\}
}{
  \Phi \vdash \Psi
}}{PsdCertDiag}{blcatrule-PsdCertDiag}
\paradeentry{\inferrule*{
  A \in \Q^{(m+k) \times (m+k)}, \ \text{kernel coordinates last}
  \\
  L := (I \mid N), \ N \in \Q^{m \times k}
  \\
  A = L^{\mathsf{T}} B L \ \text{entrywise exactly}
  \\
  B \in \Q^{m \times m} \ \text{with a /T-or-/DD inner witness } Z
  \\
  p_1, \dots, p_{m+k} \ \text{payload terms}
  \\
  c \equiv -\textstyle\sum_{i,j} A_{ij}\, p_i\, p_j \ \text{(exact canonical polynomial)}
  \\
  \Psi \in \bigl\{\, c \leq 0, \;\; Q \geq 0 \text{ where } c = -Q, \;\; l \leq r \text{ where } c = l - r \,\bigr\}
}{
  \Phi \vdash \Psi
}}{PsdCertKernel}{blcatrule-PsdCertKernel}
\paradeentry{\inferrule*{
  \Phi \wedge (h \geq 0) \vdash \Psi\sub{b}{h} \\
  \Phi \wedge (h \leq 0) \vdash \Psi\sub{b}{0} \\
  \ReLU_{\{h\}\to\{y\}} \in H \quad (y \text{ typed at rank } 0)
}{
  \Phi \vdash \Psi
}}{ReLU-Split}{blcatrule-ReLUSplit}
\paradeentry{\inferrule*{
  \Phi_1 \vdash Q \\
  \Phi_2 \wedge Q \vdash \Psi \\
  \wires(Q) \subseteq Y \\
  \wires(\Phi_1) \cap \bigl(\wires(\Phi_2) \cup \wires(\Psi)\bigr) \subseteq Y
}{
  \Phi_1 \wedge \Phi_2 \vdash \Psi
}}{Promise}{blcatrule-Promise}
\paradeentry{\inferrule*{
  \bigl\{\Phi_i \vdash e_i \leq 0\bigr\}_{i=1}^{n} \\
  \bigl\{g_k = 0\bigr\}_k \ \text{equality sub-conjuncts of } \Phi \text{ in wire variables} \\
  \lambda_i \in \Q_{\geq 0} \\
  \nu_k \in \Q \\
  c \equiv \textstyle\sum_i \lambda_i e_i + \sum_k \nu_k g_k \ \text{(exact canonical polynomial)} \\
  \Psi \in \bigl\{\, c \leq 0, \;\; e \geq 0 \text{ where } c = -e, \;\; l \leq r \text{ where } c = l - r \,\bigr\}
}{
  \Phi \vdash \Psi
}}{Syzygy}{blcatrule-Syzygy}
\paradeentry{\inferrule*{
  \Phi_1 \vdash e_1 \leq 0 \\
  \Phi_2 \vdash e_2 \leq 0 \\
  \Phi_w \vdash \lambda\sigma_2 - (1{-}\theta)\,\alpha_1 \leq 0 \\
  \Phi_w \vdash \sigma_1 - \lambda(1{-}\theta)\,\alpha_2 \leq 0 \\
  \Phi \supseteq \Phi_1,\ \Phi_2,\ \Phi_w \ \text{conjunct-wise (checked)} \\
  \lambda \in \Q_{>0} \\
  \theta \in \Q \cap (0,1)
}{
  \Phi \vdash \mathcal{L}(V_1 + \lambda V_2) + \theta\,(\alpha_1 + \lambda\alpha_2) \leq 0
}}{SmallGain}{blcatrule-SmallGain}
\paradeentry{\inferrule*{
  \bigl\{\Phi_i \vdash e_i \leq 0\bigr\}_{i=1}^{n} \\
  \bigl\{\Phi_w \vdash m_i \leq 0\bigr\}_{i=1}^{n} \\
  \Phi \supseteq \Phi_1, \ldots, \Phi_n, \Phi_w \ \text{conjunct-wise (checked)} \\
  \lambda \in \Q_{>0}^{n} \\
  B \in \Q_{\geq 0}^{n \times n} \\
  \textstyle\sum_i B_{ij} = \lambda_j \ (\forall j) \\
  \theta \in \Q \cap (0,1)
}{
  \Phi \vdash \mathcal{L}\bigl(\textstyle\sum_i \lambda_i V_i\bigr) + \theta \textstyle\sum_j \lambda_j \alpha_j \leq 0
}}{SmallGainN}{blcatrule-SmallGainN}
\paradeentry{\inferrule*{
  B_1 \ \text{a certified attracting-block lattice for } \varphi_1 \ \text{over } H_1
  \\
  B_2 \ \text{likewise for } \varphi_2 \ \text{over } H_2
  \\
  \mathrm{wires}(H_1) \cap \mathrm{wires}(H_2) = \emptyset
}{
  \{\, b_1 \times b_2 \mid b_1 \in B_1,\ b_2 \in B_2 \,\} \ \text{(the meet-generators) certified for } \varphi_1 \times \varphi_2 \ \text{over } H_1 \sqcup H_2
}}{ProductLattice}{blcatrule-ProductLattice}
\paradeentry{\inferrule*{
  \mathrm{child} : \mathsf{QRegionSat}(\varphi, R) \\
  \mathrm{eval} : \vdash x_0 \in R \ \text{(BL leaf; sublevel} \to \rulename{PConstExact}\text{, box} \to \text{registry interval check)} \\
  \mathsf{SInit}(\varphi) \in s.\forall
}{
  \mathcal{M} \vDash s
}}{QSatFromRegion}{blcatrule-QSatFromRegion}
\paradeentry{\inferrule*{
  \mathrm{child} : \mathsf{QRegionSat}(\varphi, R) \\
  \mathsf{SInit}(\varphi) \in s.\forall
}{
  \mathcal{M} \vDash \mathsf{Guard}(R,\ s)
}}{QGuardIntro}{blcatrule-QGuardIntro}
\paradeentry{\inferrule*{
  \mathcal{M} \vDash \mathsf{Guard}(R,\ s) \\
  \mathrm{eval} : \vdash X_0 \subseteq R
}{
  \mathcal{M} \vDash s
}}{QGuardElim}{blcatrule-QGuardElim}
\paradeentry{\inferrule*{
  V : \mathsf{QCertFam} \\
  \bigl\{D_m \vdash \mathcal{A}_m V_m \le 0\bigr\}_{m} \ \text{(interiors)} \\
  \bigl\{\mathrm{slab}_j \vdash g_j \fatsemi V_{m'_j} \le V_{m_j}\bigr\}_{j} \ \text{(guards, per jump $j$: source mode $m_j$, target $m'_j$, reset $g_j$)} \\
  \bigl\{U \cap \bar{D}_m \vdash V_m \ge 1\bigr\}_{m} \ \text{(floors)} \\
  \bigl\{\bar{D}_m \vdash V_m \ge 0\bigr\}_{m} \ \text{(nonnegs)} \\
  \mathrm{scope\_warrant} : \mathsf{QInvScope} \ \text{(optional — the dissolved ScopeRule)} \\
  p + y \le 1 \ \text{(threshold tie)}
}{
  \mathcal{M},\ x_0 \in \{V \le y\} \vDash \mathsf{SProbGE}\bigl(p,\ \mathsf{PAlways}\ \varphi\bigr)
}}{QSafetyBattery}{blcatrule-QSafetyBattery}
\paradeentry{\inferrule*{
  V : \mathsf{QCertFam} \\
  R \subseteq [\![\varphi]\!] \ \text{(target formula tie)} \\
  \bigl\{D_m \smallsetminus R \vdash \mathcal{A}_m V_m \le -\theta V_m\bigr\}_{m} \ \text{(decay)} \\
  \bigl\{\mathrm{slab}_j \vdash g_j \fatsemi V_{m'_j} \le V_{m_j}\bigr\}_{j} \ \text{(guards, per jump $j$: source mode $m_j$, target $m'_j$, reset $g_j$)} \\
  \bigl\{\bar{D}_m \smallsetminus R \vdash V_m \ge 1\bigr\}_{m} \ \text{(floors, off-target only)} \\
  \bigl\{\bar{D}_m \vdash V_m \ge 0\bigr\}_{m},\ \bigl\{\bar{D}_m \vdash V_m \le M\bigr\}_{m} \ \text{(nonnegs; bound)} \\
  \vdash q \le \textstyle\sum_{k \le n} (\theta T)^k / k! \ \text{(one \rulename{TaylorBound} leaf)} \\
  0 \le \beta,\ \beta \le e^{-\theta},\ \theta \ge 0,\ 0 \le q,\ p \le 1 \ \text{(drift tie; sign gates)}
}{
  \mathcal{M},\ x_0 \in \{V \le (1-p)\,q\} \vDash \mathsf{SProbGE}\bigl(p,\ \mathsf{PDeadline}(T, \varphi, b)\bigr)
}}{QDeadlineBattery}{blcatrule-QDeadlineBattery}
\paradeentry{\inferrule*{
  \mu(E_1) \le a \\ \mu(E_2) \le b \\ \mathrm{fail} \subseteq E_1 \cup E_2
}{
  \mu(\mathrm{fail}) \le a + b
}}{QUnionBound}{blcatrule-QUnionBound}
\paradeentry{\inferrule*{
  E_{\mathrm{bad}} := \{\text{the first exit from } D \smallsetminus (A \cup R) \ \text{lands in } A \cup D^{\mathrm{c}}\}, \quad E_{\mathrm{miss}} := \{A \cup R \ \text{not reached by } T\} \\
  [\![\neg\varphi_s]\!] \subseteq A, \qquad R \subseteq [\![\varphi]\!] \ \text{(formula ties)} \\
  \mu_x(E_{\mathrm{bad}}) \le V_a(x) \ \text{(the \rulename{QSafetyBattery} slice for } V_a \ \text{at stop set } (A \cup R) \cup D^{\mathrm{c}}\text{: floor on } A \ \text{only)} \\
  \mu_x(E_{\mathrm{miss}}) \le \beta^{T}\, V_d(x) \ \text{(the \rulename{QDeadlineBattery} slice for } V_d \ \text{at target } A \cup R\text{)} \\
  [\![\neg(\varphi_s\,\mathsf{U}_{\le T}\,\varphi)]\!] \subseteq E_{\mathrm{bad}} \cup E_{\mathrm{miss}} \ \text{(the d\'ebut decomposition; consumed by \rulename{QUnionBound})} \\
  0 \le \beta, \qquad 0 \le V_d(x), \qquad \beta^{T} \le c \ \text{(sign gates; rate tie)}
}{
  \mathcal{M},\ x_0 \in \{V_a + c\,V_d \le 1 - p\} \vDash \mathsf{SProbGE}\bigl(p,\ \mathsf{PUntil}(T, \varphi_s, \varphi)\bigr)
}}{QUntilBattery}{blcatrule-QUntilBattery}
\paradeentry{\inferrule*{
  V_0,\ (W_i)_{i \le k},\ (V_i)_{i \in A} : \mathsf{QCertFam}; \quad A \subseteq \{1, \ldots, k\} \ \text{(the pairs posting absorbing certificates)}; \quad \varepsilon_0 \in \Q_{>0}; \quad (\gamma_i)_{i \in A} \in \Q_{>0} \\
  \text{the \rulename{QSafetyBattery} rows for } V_0 \ \text{at } U := \mathrm{cmpl}(I) \ \text{(interior, guard, floor on } U \cap \bar{D}\text{, nonneg on } \bar{D}\text{)} \\
  \bigl\{S \smallsetminus [\![\rho_i]\!] \vdash \mathcal{A}_m V_{i,m} \le 0\bigr\}_{i \in A} \ \text{(absorbing drift: off-request, on all of } S\text{)} \\
  \bigl\{S \vdash 0 \le V_i \le C^{a}_i\bigr\}_{i \in A} \ \text{(absorbing bounds, on all of } S\text{)}, \quad \bigl\{[\![\rho_i]\!] \vdash V_i \ge 1\bigr\}_{i \in A} \ \text{(absorbing floors)} \\
  \bigl\{([\![\varphi_i]\!] \cup J_i)^{\complement} \vdash \mathcal{A}_m W_{i,m} \le -\varepsilon_0\bigr\}_{i \le k} \ \text{(ranking decay)}, \quad J_i := \{V_i \le 1 - \gamma_i\} \cap (I \smallsetminus [\![\rho_i]\!]) \ \text{for } i \in A; \quad J_i := \emptyset \ \text{for } i \notin A \\
  W_i \equiv 0 \ \text{on target cells (containment)}; \quad \bigl\{\bar{D} \vdash W_i \ge 0\bigr\}_{i \le k} \ \text{(ranking nonnegs)} \\
  \bigl\{\bar{D} \vdash W_i \le C_i\bigr\}_{i \le k}, \quad C_i \in \Q \\
  p + y \le 1 \ \text{(threshold tie)}
}{
  \mathcal{M},\ x_0 \in \{V_0 \le y\} \vDash \mathsf{SProbGE}\bigl(p,\ \mathsf{PDwellStreett}(\mathsf{PairBase}(\rho_1, \varphi_1), \ldots, \mathsf{PairBase}(\rho_k, \varphi_k))\bigr)
}}{QStreettBattery}{blcatrule-QStreettBattery}
\paradeentry{\inferrule*{
  \overline{\mathrm{drift}_m : R_m \vdash g - r_m - \mathcal{A}_m B_m \ \ge\ 0} \\
  \overline{\mathrm{guard}_j} \quad \overline{\mathrm{floor}_m} \quad \overline{\mathrm{nonneg}_m} \\
  \mathrm{tie} : \vdash g \bowtie c
}{
  \mathcal{M},\ x_0 \in R \vDash \mathsf{SReturn}{\bowtie}\bigl(c,\ \mathsf{poAvg}(o)\bigr)
}}{QAvgBattery}{blcatrule-QAvgBattery}
\paradeentry{\inferrule*{
  \mathsf{TrMeter}(t, d) \\
  \rho \cdot d \ge 1 \\
  \sigma \ge 1 \\
  \mathsf{qtsl\_rate\_clock}(n)
}{
  \mathcal{M} \vDash \mathsf{SRateLE}(\{t\},\ \rho,\ \sigma)
}}{QRateFromClock}{blcatrule-QRateFromClock}
\paradeentry{\inferrule*{
  \mathsf{BoxInv}(B, \nu, x)
}{
  \mathsf{AService}(B, T, \chi \cap \nu, x) \;\leftrightarrow\; \mathsf{AService}(B, T, \chi, x)
}}{QInvRelativize}{blcatrule-QInvRelativize}
\paradeentry{\inferrule*{
  B.\mathsf{FusionClosed} \\ 0 < n \\
  \forall i.\ \mathsf{FirstVisits}(B, \alpha_i) \\
  \forall i.\ \mathsf{AService}(B, T_i, \alpha_i, x)
}{
  \mathsf{AService}\big(B,\ {\textstyle\min_i T_i},\ \mathsf{fsTarget}(\alpha, T),\ x\big)
}}{QFirstService}{blcatrule-QFirstService}
\paradeentry{\inferrule*{
  B.\mathsf{FusionClosed} \\ 0 < n \\ \forall i.\ c < T_i \\
  \forall i.\ \mathsf{FirstVisits}(B, \alpha_i) \\
  \forall i.\ \mathsf{AService}(B, T_i, \alpha_i, x)
}{
  \mathsf{AService}\big(B,\ c,\ \mathsf{fsTarget}(\alpha, T) \cup \mathsf{residualCell}(\alpha, T, c),\ x\big)
}}{QFirstServiceSplit}{blcatrule-QFirstServiceSplit}
\paradeentry{\inferrule*{
  \text{one \rulename{QStreettBattery} node at } \rho_i := \mathsf{STrue} \ \text{for every pair} \\
  \text{(absorbing families and } J_i \ \text{syntactically empty; ranking decay off } [\![\varphi_i]\!] \ \text{alone)}
}{
  \mathcal{M},\ x_0 \in \{V_0 \le y\} \vDash \mathsf{SProbGE}\bigl(p,\ \mathsf{PDwellStreett}(\mathsf{PairBase}(\mathsf{STrue}, \varphi_1), \ldots, \mathsf{PairBase}(\mathsf{STrue}, \varphi_k))\bigr)
}}{QStreettDegenerate}{blcatrule-QStreettDegenerate}
\paradeentry{\inferrule*{
  \text{a \rulename{QStreettBattery} kit at certified regions } (R_i, G_i) \\
  \bigl\{ [\![\mathrm{cmpl}\,\rho_i]\!] \subseteq R_i \bigr\} \ \text{(requests over-approximated)} \\
  \bigl\{ G_i \subseteq [\![\varphi_i]\!] \bigr\} \ \text{(responses under-approximated)}
}{
  \mathcal{M},\ x_0 \in \{V_0 \le y\} \vDash \mathsf{SProbGE}\bigl(p,\ \mathsf{PDwellStreett}(\mathsf{PairCmpl}(\rho_1, \varphi_1), \ldots)\bigr)
}}{QStreettCmplRequest}{blcatrule-QStreettCmplRequest}
\paradeentry{\inferrule*{
  \mathcal{K},\ \mathrm{slab} \vdash W_i \ge 0
}{
  \mathcal{K},\ \mathrm{slab} \vdash g \fatsemi W_i \le W_i
}}{QStreettTargetCellCollapse}{blcatrule-QStreettTargetCellCollapse}
\paradeblock{\begin{mathpar}
\inferrule*{ }{ [\![\neg\neg\varphi]\!] = [\![\varphi]\!] }
\and
\inferrule*{ }{ [\![\neg(\varphi \wedge \psi)]\!] = [\![\neg\varphi \vee \neg\psi]\!] }
\and
\inferrule*{ }{ [\![\neg(\varphi \vee \psi)]\!] = [\![\neg\varphi \wedge \neg\psi]\!] }
\and
\inferrule*{ }{ [\![\neg\{f \le c\}]\!] = \{f > c\} }
\and
\inferrule*{ }{ [\![\neg\{f < c\}]\!] = \{f \ge c\} }
\and
\inferrule*{ }{ [\![\neg\{f \ge c\}]\!] = \{f < c\} }
\and
\inferrule*{ }{ [\![\neg\{f > c\}]\!] = \{f \le c\} }
\end{mathpar}}{QNnfRewrite}{blcatrule-QNnfRewrite}
\paradeblock{\begin{mathpar}
\inferrule*{
  \mathcal{M},\ x_0 \in R \vDash \mathsf{SProbGE}\bigl(p,\ \mathsf{PDwellStreett}(\ldots,\ \mathsf{PairCmpl}(\mathrm{elimNeg}(\mathrm{nnf}(\rho_i)),\ \varphi_i),\ \ldots)\bigr)
}{
  \mathcal{M},\ x_0 \in R \vDash \mathsf{SProbGE}\bigl(p,\ \mathsf{PDwellStreett}(\ldots,\ \mathsf{PairCmpl}(\rho_i,\ \varphi_i),\ \ldots)\bigr)
}
\and
\inferrule*{ }{ [\![\mathrm{elimNeg}(\mathrm{nnf}(\rho))]\!] = [\![\rho]\!] }
\end{mathpar}}{QStreettRequestNnf}{blcatrule-QStreettRequestNnf}
\paradeentry{\inferrule*{
  \mathcal{M},\ x_0 \in R \vDash \mathsf{SProbGE}\bigl(p,\ \mathsf{PAlways}(\mathsf{SNot}\,\beta)\bigr)
}{
  \mathcal{M},\ x_0 \in R \vDash \mathsf{P}_{\ge p}[\mathsf{avoid}\,\beta]
}}{QAvoid}{blcatrule-QAvoid}
\paradeentry{\inferrule*{
  \mathcal{M},\ x_0 \in R \vDash \mathsf{SProbGE}\bigl(p,\ \mathsf{PDwellStreett}(\mathsf{PairBase}(\mathsf{STrue}, \varphi))\bigr) \\
  \text{paths total over the state space}
}{
  \mathcal{M},\ x_0 \in R \vDash \mathsf{SProbGE}(p,\ \mathsf{GF}\,\varphi)
}}{QStreettGF}{blcatrule-QStreettGF}
\paradeentry{\inferrule*{
  \mathcal{M},\ x_0 \in R \vDash \mathsf{SProbGE}\bigl(p,\ \mathsf{PDwellStreett}(\mathsf{PairBase}(\beta, \mathsf{SFalse}))\bigr)
}{
  \mathcal{M},\ x_0 \in R \vDash \mathsf{SProbGE}(p,\ \mathsf{FG}\,\neg\beta)
}}{QStreettFG}{blcatrule-QStreettFG}
\paradeentry{\inferrule*{
  \mathrm{child} : \mathsf{QRegionSat}\bigl(\mathsf{SProbGE}(p, \mathsf{PAlways}(\mathsf{SNot}\,\varphi)),\ R\bigr)
}{
  \mathcal{M},\ x_0 \in R \vDash \mathsf{SProbGE}(p, \mathsf{PNotPath}(\mathsf{PDeadline}(T, \varphi, \mathrm{false})))
}}{QNegFAlias}{blcatrule-QNegFAlias}
\paradeentry{\inferrule*{
  \vdash \lambda \ge 0 \ \text{on } \mathrm{cell} \\
  \vdash N + \lambda\, g_q \le 0 \ \text{on } \mathrm{cell}
}{
  \vdash N \le 0 \ \text{on } \mathrm{cell} \cap \{g_q \ge 0\}
}}{HypMultiplier}{blcatrule-HypMultiplier}
\paradeentry{\inferrule*{
  \vdash P \ \text{on } D_m \\
  \mathrm{scope} \subseteq D_m \ \text{(compile fact, not a posted premise)}
}{
  \vdash P \ \text{on scope}
}}{QFullDomainWeaken}{blcatrule-QFullDomainWeaken}
\paradeentry{\inferrule*{
  \vdash P \ \text{on } C \ \text{(the cover row)} \\
  \bigl\{\vdash K_i \subseteq \mathrm{Off}\bigr\}_i \ \text{(per-component containments)} \\
  D_m \smallsetminus C \subseteq \textstyle\bigcup_i K_i \ \text{(residual decomposition)}
}{
  \vdash P \ \text{on } D_m \smallsetminus \mathrm{Off}
}}{QVerifiedCover}{blcatrule-QVerifiedCover}
\paradeblock{\begin{mathpar}
\inferrule*{
  \mathrm{child}_1 : \mathsf{QRegionSat}(\varphi_1, R_1) \\
  \mathrm{child}_2 : \mathsf{QRegionSat}(\varphi_2, R_2)
}{
  \mathcal{M},\ x_0 \in R_1 \cap R_2 \vDash \varphi_1 \wedge \varphi_2
}
\and
\inferrule*{
  \mathrm{child}_1 : \mathsf{QRegionSat}(\varphi_1, R_1) \\
  \mathrm{child}_2 : \mathsf{QRegionSat}(\varphi_2, R_2)
}{
  \mathcal{M},\ x_0 \in R_1 \cup R_2 \vDash \varphi_1 \vee \varphi_2
}
\end{mathpar}}{QPositiveCombinationRoot}{blcatrule-QPositiveCombinationRoot}
\paradeblock{\begin{mathpar}
\inferrule*{ }{ \mathrm{cmpl}\{p < c\} = \{c \le p\} }
\and
\inferrule*{ }{ \mathrm{cmpl}\{c < p\} = \{p \le c\} }
\end{mathpar}}{QClosureExact}{blcatrule-QClosureExact}
\paradeblock{\begin{mathpar}
\inferrule*{ }{ \mathrm{cmpl}\{p \le c\} \subseteq \{c \le p\} }
\and
\inferrule*{ }{ \mathrm{cmpl}\{c \le p\} \subseteq \{p \le c\} }
\and
\inferrule*{
  R \subseteq [\![\varphi]\!]
}{
  \mathrm{cmpl}\,[\![\varphi]\!] \subseteq \mathrm{cmpl}\,R
}
\end{mathpar}}{QClosureOutward}{blcatrule-QClosureOutward}
\paradeblock{\begin{mathpar}
\inferrule*{
  \varepsilon \in \Q_{>0}
}{
  \{p \le c - \varepsilon\} \subseteq \{p < c\}
}
\and
\inferrule*{
  \varepsilon \in \Q_{>0}
}{
  \{c + \varepsilon \le p\} \subseteq \{c < p\}
}
\end{mathpar}}{QClosureInwardEps}{blcatrule-QClosureInwardEps}
\paradeentry{\inferrule*{
  q \in \Q; \quad \theta \in \Q_{>0}; \quad T \in \N; \quad n \in \N \\
  q \le \textstyle\sum_{k \le n} (\theta T)^k / k! \ \text{holds under exact evaluation in } \Q
}{
  \Phi \vdash q \le \textstyle\sum_{k \le n} (\theta T)^k / k!
}}{TaylorBound}{blcatrule-TaylorBound}
\paradeblock{\begin{mathpar}
\inferrule*{
  R = \{x : V(x) \le y\} \ \text{(sublevel root atom)} \\
  \vdash V(x_0) \le y \ \text{(one \rulename{PConstExact} leaf: the substituted comparison is closed)}
}{
  \vdash x_0 \in R
}
\and
\inferrule*{
  R = \textstyle\prod_i [a_i, b_i] \ \text{(box root atom)} \\
  \forall i.\ a_i \le (x_0)_i \le b_i \ \text{(recomputed in exact } \Q\text{)}
}{
  \vdash x_0 \in R
}
\end{mathpar}}{RootEval}{blcatrule-RootEval}
\paradeentry{\inferrule*{
  N \in \Q[z_1, \ldots, z_n]; \quad D = \textstyle\prod_{v} [a_v, b_v] \subseteq \Q^n
  \\
  \varphi = (A,\ A^{-1},\ b): \quad A,\ A^{-1} : \Q^{[n] \times [n]},\ \ b : \Q^{[n]}; \qquad \varphi(w) = A\,w + b \ \text{(affine-rung witness; } \varphi \ \text{omitted reads as } \mathrm{id}\text{)}
  \\
  \vdash A^{-1} A = 1_{[n]} \ \text{(the one certified product, exact)}
  \\
  \hat N = \varphi \fatsemi N \ \text{(exact substitution)}
  \\
  \bigl\{b_\alpha \le 0\bigr\}_{\alpha} \ \text{(every tensor-product Bernstein coefficient of } \hat N \ \text{over } D\text{)}
}{
  D \vdash \hat N \le 0
}}{BoxBern}{blcatrule-BoxBern}
\paradeentry{\inferrule*{
  f \in F \quad n \in \mathbb{N} \quad r, R \in \Q
  \\
  a_k = a^f_k \ \text{for } k = 0..n \ \text{(exact coefficient recurrence)}
  \\
  0 \le r < n + 2
  \qquad
  R \;\ge\; \frac{r^{n+1}}{(n+1)!}\cdot\frac{n+2}{\,n+2-r\,}
}{
  \forall x,\ |x| \le r \;\Rightarrow\; \Bigl|f(x) - \textstyle\sum_{k \le n} a_k x^k\Bigr| \le R
}}{PExpTypeEnclosure}{blcatrule-PExpTypeEnclosure}
\paradeentry{\inferrule*{
  L : \mathsf{QBlockLattice} \ \text{(blocks over one rational grid)} \\
  \bigl\{F_f \vdash \sigma_f\, \mathcal{A}\ell_f > 0\bigr\}_{f \in \mathrm{facets}(L)} \ \text{(facet-sign rows: } \ell_f \ \text{the facet\textquotesingle s normal coordinate, } \sigma_f \in \{\pm 1\} \ \text{the declared crossing sign, } F_f \ \text{the facet)} \\
  h : \mathsf{QCellMap} \ (\text{cell} \to \text{vertex} \sqcup \{\text{transit}\}) \\
  \bigl\{\vdash T(b) \subseteq \rho(h(b))\bigr\}_{h(b) \ne \text{transit}} \ \text{(containments, per vertex cell)} \\
  \bigl\{T(b) \vdash \mathcal{A} V_b + \varepsilon_b \le 0\bigr\}_{h(b) = \text{transit}}, \quad \varepsilon_b \in \Q_{>0} \ \text{(transits, per transit cell)} \\
  \bigl\{\vdash \rho(u) \cap \rho(v) = \emptyset\bigr\}_{u \ne v \in G}, \quad \bigl\{\vdash \rho(v) \subseteq D\bigr\}_{v \in G} \ \text{(} \rho\text{-wf: rows where not box-aligned, recomputed where box-aligned)} \\
  \mathrm{CR}(\mathcal{M}|_D) \subseteq \textstyle\bigcup_{h(b) \ne \text{transit}} T(b) \ \text{(cell-level chain-recurrence cover)} \\
  \mathsf{MorseAssembly}(L, h, G, \rho) \ \text{(the assembly ledger, defined below)}
}{
  \mathcal{M} \vDash \mathsf{Morse}(D, G, \rho)
}}{QMorseDischarge}{blcatrule-QMorseDischarge}
\paradeentry{\inferrule*{
  \mathcal{M}_1 \vDash \mathsf{Morse}(D_1, G_1, \rho_1) \\
  \mathcal{M}_2 \vDash \mathsf{Morse}(D_2, G_2, \rho_2) \\
  L_1, L_2 : \mathsf{QBlockLattice} \ \text{(the factor lattices)}
}{
  \mathcal{M}_1 \times \mathcal{M}_2 \vDash \mathsf{Morse}(D_1 \times D_2,\ G_1 \times G_2,\ \rho_1 \times \rho_2)
}}{QMorseProduct}{blcatrule-QMorseProduct}
\paradeentry{\inferrule*{
  N : \mathsf{QRegionAtom} \ \text{(the block)} \\
  \bigl\{F_f \vdash \mathcal{A}\ell_f \ge \varepsilon_f\bigr\}_{f \in \mathrm{facets}(N)}, \quad \varepsilon_f \in \Q_{>0} \ \text{(} \ell_f \ \text{the facet\textquotesingle s inward affine coordinate; } F_f \ \text{its facet region)} \\
  \mathsf{FacetBridge}(\mathcal{M}, D, N, (F_f, \varepsilon_f)_f) \ \text{(the bridge package, defined below)} \\
  \vdash X_0 \subseteq N \ \text{(\texttt{init\_containment}, a single BL leaf)}
}{
  \mathcal{M} \vDash \mathsf{QInvScope}(N, X_0)
}}{QBlockInv}{blcatrule-QBlockInv}
\paradeentry{\inferrule*{
  L : \mathsf{QBlockLattice} \\
  C \ \text{(\texttt{downset}, a cell-set; membership and down-closedness recomputed)} \\
  \mathcal{M} \vDash \mathsf{QInvScope}\bigl(\textstyle\bigcup_{i \in C} \mathrm{blocks}(\mathrm{emb}\ i),\ X_0\bigr) \ \text{(the cited \rulename{QBlockInv} conclusion)} \\
  \vdash X_0 \subseteq \textstyle\bigcup\downarrow C \ \text{(\texttt{init\_containment})}
}{
  \mathcal{M} \vDash \mathsf{QInvScope}\bigl(\textstyle\bigcup\downarrow C,\ X_0\bigr)
}}{QLatticeInv}{blcatrule-QLatticeInv}
\paradeentry{\inferrule*{
  \mathcal{M} \vDash \mathsf{QInvScope}(S, X_0) \\
  \bigl\{\vdash S \cap [\![\beta_j]\!] = \emptyset\bigr\}_{j} \ \text{(\texttt{avoid\_disjoint} rows)}
}{
  \mathcal{M},\ x_0 \in X_0 \vDash \mathsf{SProbGE}\bigl(1,\ \mathsf{PAlways}(\mathsf{SNot}\,\beta_j)\bigr)
}}{QBasin}{blcatrule-QBasin}
\paradeentry{\inferrule*{
  b : \mathsf{QMirandaBattery} \\
  \mathsf{qtsl\_miranda\_battery}(b,\ C,\ g,\ \Pi)
}{
  \vDash \mathsf{QZeroExists}(g,\ C)
}}{QZeroByMiranda}{blcatrule-QZeroByMiranda}
\paradeentry{\inferrule*{
  \mathrm{inv} : \mathsf{QBlockInv}\ \text{(strictly-inward facets, box block)}
}{
  \vDash \mathsf{QEquilibriumExists}(f,\ C)
}}{QEquilibriumFromInward}{blcatrule-QEquilibriumFromInward}
\paradeentry{\inferrule*{
  b : \mathsf{QMirandaBattery} \\
  \mathsf{qtsl\_miranda\_battery}\bigl(b,\ W,\ \sigma - \mathrm{id}_{Z},\ Z\bigr)
}{
  \vDash Z \subseteq \sigma(W)
}}{QCoverByMiranda}{blcatrule-QCoverByMiranda}
\paradeentry{\inferrule*{
  \bigl\{\Sigma_i \vDash (I_i, \alpha_i, \gamma_i[, \kappa_i])\bigr\}_{i=1}^{n} \\
  \mathsf{QWireCert}\bigl(w,\ (\mathit{spec}_i)_{i} \Rightarrow \mathit{spec}_{\mathrm{out}},\ \kappa\bigr) \\
  \mathrm{rank} : \{1, \ldots, n\} \to \N; \quad \bigl\{\mathrm{rank}_i < \mathrm{rank}_j\bigr\}_{(i \to j) \in w} \ \text{(the causality certificate: one rank per component, strictly increasing along every wire)} \\
  \mathrm{W0}\!: \ \text{the composite martingale problem admits a c\`adl\`ag strong-Markov solution, unique in law, from every initial state (stochastic mechanisms; the martingale-problem sense --- distinct from the aggregation language\textquotesingle s own W0)} \\
  \mathrm{W1}\!: \ \text{transition data factor through per-place colour multisets (SDCPN mechanisms; typing-discharged in the aggregation language)} \\
  \mathrm{W2}\!: \ \text{the within-instant immediate-jump relation is well-founded, of uniformly bounded depth} \\
  \mathrm{W3}\!: \ \text{simultaneous enablings carry a declared serialization or a certified reset commutation (a.s.-null collision row for stochastic timers)} \\
  \kappa_{\mathrm{out}} := \textstyle\bigoplus_{\mathrm{legs}} \kappa_{\mathrm{leg}}, \quad \kappa_a \oplus \kappa_b := \kappa_a + \kappa_b \ \text{(slack addition along lens composition; checker-recomputed, never stored)}
}{
  \mathsf{SysWire}(w, [\Sigma_1, \ldots, \Sigma_n]) \vDash \mathit{spec}_{\mathrm{out}}
}}{QWireCompose}{blcatrule-QWireCompose}
\paradeentry{\inferrule*{
  \bigl\{R_\ell,\ \varphi(s) \wedge \alpha(v(s), a) \vdash \mathrm{Lift}_s\, \varphi(u(s, a))\bigr\}_{\ell},\ \ \bigl\{R_\ell,\ \varphi(s) \vdash \gamma(v(s))\bigr\}_{\ell} \ \text{(mechanism-certification legs; } s \ \text{a mechanism state, } a \ \text{an admissible input)} \\
  \bigl\{R_\ell \vdash \alpha_2(w(o_1), a_2) + \kappa(\gamma_1(o_1)) \ge \alpha_1(w^{\sharp}(o_1, a_2))\bigr\}_{\ell} \ \text{(assumption legs)} \\
  \bigl\{R_\ell \vdash \gamma_1(o_1) \ge \gamma_2(w(o_1))\bigr\}_{\ell} \ \text{(guarantee legs)} \\
  \vdash \kappa(0) = 0 \ \text{(closed evaluation)} \\
  [0, \infty) \vdash \kappa' \ge 0 \ \text{(derivative-nonnegativity row)} \\
  \vdash \deg \kappa > 0,\ \mathrm{lead}(\kappa) \ge 0 \ \text{(leading-coefficient row)}
}{
  \mathsf{QWireCert}\bigl(w,\ (\mathit{spec}_i) \Rightarrow \mathit{spec}_{\mathrm{out}},\ \kappa\bigr)
}}{QWireCertDischarge}{blcatrule-QWireCertDischarge}
\paradeentry{\inferrule*{
  \bigl\{\Sigma_i \vDash (I_i,\ A_i,\ E_i)\bigr\}_{i=1}^{n}
  \quad \text{where } A_i \equiv \textstyle\bigwedge_{j \in \mathrm{in}(i)} \mathrm{ev}(g_j)\ \mathsf{PathPred},
  \ \ E_i \equiv g_i\ \mathsf{ExistsSpec}
  \ \ \text{(conditional delivery: } \forall\text{-conforming-}\exists\text{-delivering, properness presupposed)} \\
  \bigl\{\ell_i\ \text{latency rows}\bigr\}_{i=1}^{n}
  \ \text{(each premise ships its delivery-stage arithmetic + seriality, at its own profile tier)} \\
  \mathsf{QWireCert}\bigl(w,\ (\mathit{spec}_i)_i \Rightarrow \mathit{spec}_{\mathrm{out}},\ -\bigr)
  \qquad
  \bigwedge_i g_i \vdash \mathrm{marks}(\mathit{spec}_{\mathrm{out}})
  \ \text{(mark-entailment cut)} \\
  \mathsf{MarkFlow}\bigl(w,\ (t_k)_k,\ (P_k)_{k \in \mathrm{cyc}_0(w)}\bigr)
  \ \text{(the schedule walk, two columns: stage rows --- difference-constraint feasibility at the labels; carve rows --- containment $P_k \subseteq \mathrm{mark}_k$ + carved-delivery citations + one species-tagged diagonal certificate per zero-lag cycle, exact values as the degenerate species)} \\
  \mathrm{W0}\ \text{(composite properness --- the existing family; the liveness construction\textquotesingle s tier discharge is stepwise, schedule-relative)}
}{
  \mathsf{SysWire}(w, [\Sigma_1, \ldots, \Sigma_n]) \vDash \mathit{spec}_{\mathrm{out}}
}}{QWireComposeMarks}{blcatrule-QWireComposeMarks}
\paradeentry{\inferrule*{
  S' \vDash \mathit{spec}' \\
  \mathsf{QSim}(\sigma, S, S')
}{
  S \vDash \sigma^{*}\mathit{spec}'
}}{QSimTransport}{blcatrule-QSimTransport}
\paradeentry{\inferrule*{
  \mathsf{QFlatEq}(\mathrm{sys}, \mathit{flat}) \\
  \mathcal{M}_{\mathit{flat}},\ x_0 \in R \vDash \varphi \ \text{(the closed-system pipeline\textquotesingle s judgment)}
}{
  \mathrm{sys},\ x_0 \in R \vDash \varphi
}}{QFlatten}{blcatrule-QFlatten}
\paradeentry{\inferrule*{
  \overline{\mathrm{env}_i : B \vdash\ \pm f_i(x) \ \le\ a + b\,\lVert x \rVert}
}{
  \vDash \mathsf{QLinGrowth}(f,\ a,\ b,\ B)
}}{QLinGrowthDischarge}{blcatrule-QLinGrowthDischarge}
\paradeentry{\inferrule*{
  \overline{\mathrm{dom}_k}\ \ \text{(per-contributor dominance rows)} \\
  \mathsf{qtsl\_stochwire\_shape}(\mathit{cl},\ \mathit{sp})
}{
  \vDash \mathsf{QMarginalCompat}(\mathit{sp},\ \mathit{cl})
}}{QMarginalCompat}{blcatrule-QMarginalCompat}
\paradeentry{\inferrule*{
  (\vec{s}, \vec{k}) \in K \vdash G(\vec{s}, \vec{k}) \le 0 \ \text{(one \rulename{BernCone}-family sequent over the aggregate coordinates)} \\
  s^P_j := \textstyle\sum_{t:P} \phi_j(c_t), \qquad k_P := \#P \ \text{(the aggregate coordinates, defined)} \\
  \bigl\{\vdash l_j\, k_P \le s^P_j \le u_j\, k_P\bigr\}_{j, P} \ \text{(coupling rows; } [l_j, u_j] \ \text{a Bernstein enclosure of } \phi_j \ \text{on } B_P\text{)} \\
  \bigl\{\vdash 0 \le k_P \le M_P\bigr\}_{P} \ \text{(count rows)} \\
  \bigl\{\vdash y^{\top} \vec{k} = y^{\top} \vec{k}(x_0)\bigr\}_{y \in Y} \ \text{(P-invariant conservation rows: each cited witness } y \ge 0,\ y^{\top} C = 0 \ \text{verified against the uncoloured skeleton\textquotesingle s incidence matrix } C\text{; } x_0 \ \text{the initial marking)} \\
  \bigl\{\vdash (s^P_\alpha)^2 \le k_P\, s^P_{2\alpha}\bigr\}_{\alpha, P} \ \text{(moment rows, optional tightening)} \\
  \forall x \in \mathrm{scope},\ \forall t{:}P \ \text{in } x:\ c_t \in B_P \ \text{(colour-box scope)}
}{
  x \in \mathrm{scope} \vdash G\bigl(\vec{s}(x), \vec{k}(x)\bigr) \le 0
}}{QAggReduce}{blcatrule-QAggReduce}
\paradeentry{\inferrule*{
  k_P \ \text{a count observable, integer-valued at every reachable marking} \\
  j \in \mathbb{Z}
}{
  \mathrm{cmpl}\,\{k_P \le j\} \ \rightsquigarrow \ \{k_P \ge j + 1\} \quad \text{(hypothesis positions, reachable scope)}
}}{QCountAtomRounding}{blcatrule-QCountAtomRounding}
\paradeentry{\inferrule*{
  \mathrm{upper} : \mathsf{qpaf\_upper}(n,\ c_{\mathrm{u}}) \\
  \mathrm{lower} : \mathsf{qpaf\_lower}(n,\ c_{\mathrm{l}}) \\
  \mathsf{qtsl\_palways\_fold}(n)
}{
  \mathcal{M} \vDash \mathsf{SProbGE}\bigl(p,\ \mathsf{PAlways}(\mathsf{SAnd}[\mathsf{SAtomLE}(c, V),\ \mathsf{SAtomGE}(c, V)])\bigr)
}}{QPAlwaysFromFold}{blcatrule-QPAlwaysFromFold}
\paradeentry{\inferrule*{
  \mathrm{cert} : \mathsf{HCert} \\
  d_{\min} > 0 \\
  \mathsf{qtsl\_firebound\_fold}(n)
}{
  \mathcal{M} \vDash \mathsf{SReturnLE}\bigl(B,\ \mathsf{retFireCount}(\mathit{trs})\bigr)
}}{QFireBoundFromFold}{blcatrule-QFireBoundFromFold}
\paradeentry{\inferrule*{
  n\,\mathsf{cert} = c \\ n\,\mathsf{goalfold} = f_g \\ n\,\mathsf{horizon} = T \\
  \mathsf{coverage\_lb}(c) \\ \mathsf{flow\_lb}(c) \\
  \forall r.\ u_r > 0 \Rightarrow \mathsf{TrMeter}(\mathit{tr}(r), d_r) \\
  \mathsf{qtsl\_palwayst\_fold}(n)
}{
  \mathcal{M} \vDash \mathsf{SProbGE}\bigl(p,\ \mathsf{PAlwaysT}(T,\ \mathrm{cmpl}(\mathit{goal}))\bigr)
}}{QPAlwaysTFromFold}{blcatrule-QPAlwaysTFromFold}
\paradeentry{\inferrule*{
  \mathsf{unit\_lineage}(c) \\
  \{\mathsf{hrpGap}(1), \mathsf{hrpDropAtMost}(1)\}\ \text{per advancing row on } \chi = -\psi\ \text{(pathwise species)} \\
  \mathsf{TrMeter}(\mathit{tr}(r), d)\ \text{per advancing row} \\
  \forall \tau \in m_0.\ \psi(\tau) \le 0 \\
  T < k \cdot d \\
  \mathsf{qtsl\_palwayst\_lineage}(n)
}{
  \mathcal{M} \vDash \mathsf{SProbGE}\bigl(p,\ \mathsf{PAlwaysT}(T,\ \textstyle\bigwedge_{q\,:\,\mathrm{depth}(q) \ge k} \mathrm{count}_q = 0)\bigr)
}}{QPAlwaysTFromLineage}{blcatrule-QPAlwaysTFromLineage}
\paradeentry{\inferrule*{
  \bigl\{\vdash K(p, \cdot) \ge 0,\ \textstyle\sum_q K(p, q) = 1,\ \pi_{i\#} K(p, \cdot) = \widehat{K}_i(p_i, \cdot)\bigr\} \ \text{(coupling rows, exact } \Q\text{; } \widehat{K}_i \ \text{the completed component kernels)} \\
  \bigl\{\vdash K(p, p) = 1 \ \text{for both-absorbed } p\bigr\} \ \text{(freeze rows)} \\
  \bigl\{\vdash \textstyle\sum_q K(p, q)\, V(q) \le V(p)\bigr\} \ \text{(relational supermartingale generator row, every product state)} \\
  \vdash V \ge 0 \\
  \bigl\{\vdash V \ge 1 \ \text{(disagreeing), resp.\ } V \ge |v_1 - v_2|, \ \text{on both-absorbed states}\bigr\} \ \text{(terminal domination)} \\
  \text{absorption: per-component reach criterion}; \quad \mu_0 \ \text{a probability law with } \pi_{i\#}\mu_0 = \text{the component initial laws (initial marginal rows)}; \quad \vdash \textstyle\sum_p \mu_0(p)\, V(p) \le b
}{
  \mathsf{QDist}[d](\mathcal{M}_1, \mathcal{M}_2) \le b \qquad d \in \{\mathrm{TV},\ \text{gap-pseudometrics}\}
}}{QCoupleDrift}{blcatrule-QCoupleDrift}
\paradeentry{\inferrule*{
  d : \mathsf{QDistRegistry}(X_d); \qquad (X_d, d) \ \text{complete} \\
  \vdash \rho < 1 \ \text{two-point on } \overline{B}(\hat\mu, r) \ \text{(contraction row)} \\
  \vdash d(T\hat\mu, \hat\mu) \le \delta \ \text{(residual row)} \\
  \vdash \delta + \rho\, r \le r \ \text{(radius row, division-free)}
}{
  \vdash \mathsf{QDist}[d](\hat\mu,\ \mathrm{fix}\,T) \le r
}}{QFixBall}{blcatrule-QFixBall}
\paradeentry{\inferrule*{
  J = \bigl(\mathcal{M}_c \vDash s\bigr) \ \text{(strict, law-level)} \quad \text{or} \quad J = \bigl(\mathsf{QDist}[\mathrm{TV}](\mu_c, \nu) \le \delta\bigr), \ \text{on the coarse side} \\
  \bigl\{\vdash \mathrm{tv}\bigl(\pi_{\#}\, K(s, \cdot),\ Q(\pi(s), \cdot)\bigr) \le \varepsilon\bigr\}_{s \in S_{\mathrm{f}}} \ \text{(} \sigma\text{-witness rows, one per fine state; } \varepsilon = 0 \ \text{is exact intertwining)} \\
  \vdash \mu^{0}_{c} = \pi_{\#}\, \mu^{0}_{\mathrm{f}} \ \text{(initial row)} \\
  n \in \N \ \text{(declared horizon; } \varepsilon\text{-cell only)}
}{
  \text{exact cell: } J \ \text{transports unchanged} \qquad \varepsilon\text{-cell: } \mathsf{QDist}[\mathrm{TV}]\bigl(\pi_{\#}(\text{fine law at } n),\ \text{coarse law at } n\bigr) \le n\,\varepsilon
}}{QDistTransport}{blcatrule-QDistTransport}
\paradeentry{\inferrule*{
  q : \mathsf{DAbsQuery} \\
  \mathsf{qtsl\_dabs\_wf}(q) \\
  \text{stored rows of } q.\mathsf{chain}\ (\mathsf{dprob}, \mathsf{dinit}, \mathsf{dclass}, \mathsf{dreject})
}{
  \mathsf{qtsl\_dabs\_absorb\_bracket}(q,\ c,\ lo,\ hi)
}}{DtmcAbsorbBracket}{blcatrule-DtmcAbsorbBracket}
\paradeentry{\inferrule*{
  q : \mathsf{DAbsQuery} \\
  \mathsf{qtsl\_dabs\_wf}(q) \\
  \text{stored rows of } q.\mathsf{chain}\ (\mathsf{dprob}, \mathsf{dinit}, \mathsf{dclass}, \mathsf{dreject}) \\
  S > 0\ \text{(recomputed; refusal otherwise)}
}{
  \mathsf{qtsl\_dabs\_posterior\_bracket}(q,\ c,\ lo,\ hi)
}}{DtmcPosteriorBracket}{blcatrule-DtmcPosteriorBracket}
\paradeentry{\inferrule*{
  w : \mathsf{DGenContraction} \\
  \text{stored rows: } \mathsf{dweight}\ (A = \mathrm{supp}\,V,\ \text{nonempty}),\ \mathsf{dprob},\ \mathsf{dinit} \\
  \textstyle\sum_{y \in A} K(x,y)\,V(y) \le c\,V(x)\ \ \forall x \in A \\
  K(x,y) = 0\ \text{for stored rows } x \notin A,\ y \in A \\
  \min_A V > 0 \\
  \mathsf{qtsl\_dtmc\_genform\_contraction}(w)
}{
  \text{each claimed row of } w:\ \mathsf{qtsl\_dgc\_partial\_lo},\ \mathsf{qtsl\_dgc\_ev\_bracket},\ \mathsf{qtsl\_dgc\_survival\_tail},\ \mathsf{qtsl\_dgc\_z\_one}
}}{QGenFormContraction}{blcatrule-QGenFormContraction}
\paradesep
\paradeblock{}{ProofWellFormedness}{blcatrule-ProofWellFormedness}

\blcatsection{Core Structural Rules}
\rulesection{1}{Ax}{core rule}
\label{blcatrule-Ax}
\ruleblock{The rule (Steele metanotation)}
\sourcenote{as realized in the mechanized development; source row: \texttt{\texttt{rules.toml}} --- \emph{Ax}. The metanotation statement is typeset by hand for notational quality; its mechanized ground truth is the Lean statement below, extracted from the sources at every build of this book and hash-bound. The rule \emph{roster} is machine-checked (every mechanized rule tag round-trips to exactly one catalogue row); the statement \emph{text} is editorial, checked by eye against the extracted Lean beside it.}
\begin{mathpar}
\inferrule*[Right=\rulename{Ax}]{\;}{\Phi \vdash \Phi}
\end{mathpar}

\par\medskip
\rulename{Ax} is the leaf of every derivation: it discharges a sequent
whose conclusion is literally its premise, $\Phi \vdash \Phi$. A
producer uses it wherever a derivation bottoms out at an assumption ---
the demonstration certificate cited under the comparator note below
begins with an \rulename{Ax} node discharging its composite premise
before a conjunction elimination extracts the region fact a
Bernstein-cone closure consumes.

What the checker recomputes is one thing: that the premise and
conclusion are the \emph{same formula tree}, compared node for node ---
no normalization, no reassociation, no matching. Formulas here are
ordered trees, so \rulename{Ax} does not close
$\Phi \vdash \Phi'$ for a logically equivalent but
differently shaped $\Phi'$; that is re-derived through
the conjunction rules when needed.

The single hypothesis is load-bearing in the obvious, total way: drop
the equality and the premise-free rule concludes any sequent at all ---
$\top \vdash \bot$ is the sharpest instance. Notably absent are
well-typedness side conditions, and their absence is correct: whatever
satisfaction amounts to for $\Phi$ --- even a $\Phi$ mentioning untyped
wires or atoms that never evaluate --- it sits on both sides of the
turnstile. The proof section contrasts this with \rulename{EqRefl},
which concludes a freshly minted \emph{atom} and therefore does need
typing fields.

\ruleblock{Lean statement}
\declhead{BernsteinLogic.Soundness}{sound\_ax}{\badge{BadgeGreen}{statement bound}}
\begin{leansrc}
theorem sound_ax {H : Hypernet} {seq : Sequent}
    (h : seq.pre = seq.post) : seq.semanticallyValid H
\end{leansrc}
\par\medskip
\textbf{How to read the payload.}\enspace A certificate node for this
rule carries an inhabitant of the following inductive: its
proof-obligation fields (here, one equality of formulas) are the
rule\textquotesingle s entire side-condition surface, and the soundness
lemma above takes exactly those fields as hypotheses. The six other
structural sections\textquotesingle{} payloads read the same way ---
sub-derivations plus structural equalities, nothing else.

\declhead{BernsteinLogic.Core.Catalogue}{PAxData}{\badge{BadgeGreen}{statement bound}}
\begin{leansrc}
inductive PAxData : Hypernet → Sequent → Type where
  | mk (H : Hypernet) (seq : Sequent)
       (preEqPost : seq.pre = seq.post) : PAxData H seq
\end{leansrc}
\ruleblock{Correspondence (Steele display $\leftrightarrow$ Lean)}
\badge{BadgeGreen}{reviewer-certified against the displayed texts}
\sourcenote{The dictionary below is a reviewer-certified statement of how the mechanized statements above interpret the displayed rule: display metavariable $\mapsto$ Lean binder, each display hypothesis $\mapsto$ the premise or side condition carrying it, the display conclusion $\mapsto$ the goal --- with any deliberate difference in width between display and mechanization named outright. Certification is recorded as a hash triple over (the displayed cell, the extracted statement texts, this dictionary\textquotesingle s own text); if any of the three changes, this block turns amber until a reviewer re-certifies. The correspondence claim itself remains editorial --- reviewer-grade, not machine-checked; its machine-checked successor types each soundness statement against the rule as a structured object, and supersedes this block per rule as it lands.}
\begin{itemize}
\item the judgment frame, shared by every structural dictionary below:
  a sequent is the pair of its premise and conclusion formulas, so a
  displayed judgment $\Phi \vdash \Psi$ $\mapsto$ \texttt{seq :
  Sequent} with \texttt{seq.pre}, \texttt{seq.post} the two sides, and
  asserting it $\mapsto$ the goal \texttt{seq.semanticallyValid H} ---
  at every admissible valuation of the ambient hypernet \texttt{H},
  satisfaction of \texttt{seq.pre} implies satisfaction of
  \texttt{seq.post}.
\item the display\textquotesingle s single metavariable $\Phi$,
  appearing on both sides of the turnstile, $\mapsto$ the sole
  hypothesis \texttt{h} $:$ \texttt{seq.pre = seq.post} --- sharing by
  field equation on one sequent binder, not by a shared metavariable:
  the statement quantifies over \emph{all} sequents and consumes the
  identity as its only input.
\item premise-freeness $\mapsto$ silence: no other hypothesis exists;
  \texttt{H} and \texttt{seq} are arbitrary.
\item the payload row \texttt{PAxData} (certified alongside) carries
  the same single tie as certificate data (\texttt{preEqPost}), so
  admission layer and soundness statement coincide exactly on this
  rule.
\end{itemize}
No width difference in either direction: one hypothesis, consumed by
rewriting; nothing else is stated or needed.

\ruleblock{Lean verifier receipts}
\sourcenote{module \texttt{BernsteinLogic.Soundness}; bound build \texttt{a20d90e66c13\ldots} (bound from receipt a9c0e34d3e3c9c8c, 2026-08-03T21:30:36Z). A receipt records that named checking layers accepted a build: (module, compiled-artifact SHA-256, environment SHA-256, checker layers), minted when the checking gate ran; the store receipts attesting the bound build are listed with the layers each run credited. Nothing here asks the reader to trust the store: every layer is an independently runnable program over the quoted sources --- rebuild the module and re-run the checkers to reproduce the verdict. Receipts attesting superseded builds credit no layer and are omitted from this rendering (the count of omissions, if any, is noted below).}
\begin{receipttable}
\receiptrow{a9c0e34d3e3c9c8c}{2026-08-03T21:30:36Z}{build+decorrelated}{\badge{BadgeGreen}{current}}
\end{receipttable}
\sourcenote{21 further receipt(s) in the store attest superseded builds of this module and are omitted from this rendering.}
\sourcenote{module \texttt{BernsteinLogic.Core.Catalogue}; bound build \texttt{cbc92e976a1d\ldots} (bound from receipt 5abdb3589707b5df, 2026-08-03T21:30:33Z). A receipt records that named checking layers accepted a build: (module, compiled-artifact SHA-256, environment SHA-256, checker layers), minted when the checking gate ran; the store receipts attesting the bound build are listed with the layers each run credited. Nothing here asks the reader to trust the store: every layer is an independently runnable program over the quoted sources --- rebuild the module and re-run the checkers to reproduce the verdict. Receipts attesting superseded builds credit no layer and are omitted from this rendering (the count of omissions, if any, is noted below).}
\begin{receipttable}
\receiptrow{24804e66edba0bc9}{2026-08-03T20:25:09Z}{build+decorrelated}{\badge{BadgeGreen}{current}}
\receiptrow{5abdb3589707b5df}{2026-08-03T21:30:33Z}{build+decorrelated}{\badge{BadgeGreen}{current}}
\end{receipttable}
\sourcenote{18 further receipt(s) in the store attest superseded builds of this module and are omitted from this rendering.}
\layerline{\layeritem{build}{lake build (elaboration)}{\badge{BadgeGreen}{current}}\layeritem{replay}{leanchecker (official kernel replay)}{\badge{BadgeAmber}{stale only}}\layeritem{hardened}{safe\_verify (hardened screen)}{\badge{BadgeAmber}{stale only}}\layeritem{decorrelated}{lean4export + nanoda (decorrelated kernel)}{\badge{BadgeGreen}{current}}}
\sourcenote{\badge{BadgeBlue}{comparator: exercised via the demonstration fixture}
the composite demonstration certificate (an axiom node, a conjunction
elimination, a Bernstein-cone closure) runs as a fixture through the
three-implementation comparator harness (compiled / kernel-interpreted /
reference-interpreted), so this rule\textquotesingle s node rides every
execution path of that harness.
}
\ruleblock{HOL4 statement}
\sourcenote{theory \texttt{bern\_fam\_ax}, HOL4 development script \texttt{\texttt{bern\_fam\_axScript.sml}} \badge{BadgeGreen}{script bound}}
\holhead{Definition checkFam\_ax\_def}
\begin{holsrc}
Definition checkFam_ax_def:
  checkFam_ax (cols:columns) (ct:ctx) (i:num) ⇔
    case conclRef cols ct i of
    | NONE => F
    | SOME cj =>
      case poolTagAt cols ct cj of
      | NONE => F
      | SOME t =>
        if t ≠ (48:num) then F else
        case poolField cols ct cj 0 of
        | NONE => F
        | SOME preR =>
          case poolField cols ct cj 1 of
          | NONE => F
          | SOME postR =>
            if preR = postR then T
            else
              case pv cols ct cj of
              | SOME (PJudg prew postw) =>
                  (expandSph FTt prew = expandSph FTt postw)
              | _ => F
End
\end{holsrc}
\holhead{Theorem famEq\_ax}
\begin{holsrc}
Theorem famEq_ax:
  wfColumns cols ct ∧ mkCtx cols = SOME ct ∧
  nodeCorr cols ct i (CertPf sq nd) ∧
  FST (THE (tagOf cols i)) ∈ {0}
  ⇒
  (checkFam_ax cols ct i ⇔ (sq.pre = sq.post))
\end{holsrc}
\ruleblock{HOL verifier receipts}
\badge{BadgeBlue}{documentary} Documentary citations. Spine census: \texttt{bern\_joins} GREEN under \texttt{hol-strict} (55 theorems, 0 offences, 208-theory ancestry including this theory), \texttt{checkJoins\_eq\_checkPf} verified verbatim in the built signature. The cited \texttt{famEq} lemma is the kernel-checked per-node reduction: under well-formed columns (\texttt{wfColumns}/\texttt{mkCtx}/\texttt{nodeCorr}) at wire tag 0, the columnar family check holds iff the node\textquotesingle s sequent satisfies the rule\textquotesingle s own side condition ($\mathtt{sq.pre} = \mathtt{sq.post}$). There is no signed HOL receipt store yet --- these citations are documentary.
Semantic stratum, also ported: the soundness twin \texttt{sound\_ax}
(theory \texttt{bern\_tr\_soundstruct}: \texttt{hol-strict} clean, 9
theorems, 0 offences, verification run of 2026-07-22 on an isolated
build host, content-addressed at the same sources), consumed by the
certified-fragment keystone (\texttt{checkPf\_sound\_frag}, theory
\texttt{bern\_tr\_splice}, 20 theorems, 0 offences, same run): a
reference-checker accept inside the proven wire fragment yields sequent
validity.

\sourcenote{HOL receipts are documentary: quoted results of \texttt{hol-strict} runs
and census rows from the HOL4 development\textquotesingle s build
evidence, not minted receipt artefacts. A HOL receipt store (the Lean
store\textquotesingle s twin) is named future work.
}
\ruleblock{Proof}

\sourcenote{Genre: core rule (tier 1) --- an ordinary mathematical proof
at Putnam-exam-solution detail. This section also fixes the semantic
setting shared by the seven structural sections (\rulename{Ax},
\rulename{Weaken}, \rulename{Cut}, $\wedge$I, $\wedge$E, $\vee$I,
$\vee$E); the later six abbreviate it as ``the \rulename{Ax} setting''.}

\textbf{Setting (shared by the structural sections).}\enspace Fix a
hypernet $H$ --- the ambient computation graph: typed wires, each
declaring a tensor shape and a dtype, and primitive applications
connecting them. A \emph{valuation} $\sigma$ assigns each typed wire a
tensor of values from its declared value set; $\sigma$ is
\emph{admissible} when every primitive application in $H$ satisfies its
semantic relation. A formula holds at $\sigma$ by structural recursion:
an atom $t_1 \bowtie t_2$ holds when --- and only when --- both terms
evaluate --- exact rational arithmetic, wires read whole, evaluation
failing on untyped wires, on shape-incompatible arithmetic, and on any
$\pm\infty$-valued wire mention --- to tensors compared entrywise under
$\bowtie$ (a rank-0 side broadcasts; incompatible shapes never hold), so
an atom either of whose terms fails to evaluate does not hold; a region
atom holds
when its term evaluates to a rational point of the named region;
$\Phi_1 \wedge \Phi_2$ holds when both conjuncts hold, $\Phi_1 \vee
\Phi_2$ when at least one disjunct holds, $\top$ always, $\bot$ never. A
sequent $\Phi \vdash \Psi$ is \emph{semantically valid} over $H$ when
every admissible $\sigma$ satisfying $\Phi$ satisfies $\Psi$.

Two facts frame all seven structural proofs. First, formulas are
\emph{ordered trees}: $\wedge$ and $\vee$ are binary constructors, and
every side condition of a structural rule is an identity of trees,
checked structurally --- no normalization, no reassociation, no
commuting. Wherever a proof below ``rewrites by'' a payload field, it is
replacing a formula by itself. Second, none of the seven proofs consumes
admissibility: each threads it, unexamined, to its sub-derivations, so
the structural fragment is sound over \emph{arbitrary} valuations ---
premise and conclusion validity both read over all valuations,
admissible or not --- and the admissibility hypothesis rides along only
because the validity statement quantifies over it. Where a counterexample is claimed in these
sections, read it over the one-wire hypernet ($w$ rational, rank-0, no
applications): there every valuation is admissible, so an unsound
conclusion genuinely falsifies.

\textbf{Claim (\texttt{sound\_ax}).}\enspace The payload certifies one
equality, $\mathrm{pre} = \mathrm{post}$ (\texttt{preEqPost}); under it
the premise-free rule concluding $\Phi \vdash \Phi$ is sound.

\begin{bthm}[soundness of \rulename{Ax}]\label{thm:ax-sound}
If $\mathrm{seq.pre} = \mathrm{seq.post}$, then $\mathrm{seq}$ is
semantically valid over every hypernet.
\end{bthm}
\begin{proof}
One step. Let $\sigma$ be an admissible valuation with $\sigma \models
\mathrm{seq.pre}$. The certified equality says $\mathrm{seq.post}$
\emph{is} $\mathrm{seq.pre}$ --- the same tree, not merely an equivalent
formula --- so $\sigma \models \mathrm{seq.post}$.
\end{proof}

\begin{bremark}[the one hypothesis is everything]\label{rem:ax-hyp}
\rulename{Ax} has no premises, so the payload equality is its entire
content. Drop \texttt{preEqPost} and the degenerate rule concludes an
arbitrary sequent; the sharpest instance is $\top \vdash \bot$ ---
every valuation satisfies the premise and none the conclusion.
\end{bremark}

\begin{bremark}[identity is free; reflexivity is not]\label{rem:ax-typing}
\rulename{Ax} carries no well-typedness side conditions, and this is
correct rather than an oversight: $\Phi \vdash \Phi$ is valid even when
$\Phi$ mentions untyped wires or atoms that can never evaluate, because
satisfaction --- whatever it turns out to be --- sits on both sides of
the turnstile. Contrast \rulename{EqRefl}, which concludes the
\emph{atom} $t = t$ under an arbitrary premise and does need shape and
finiteness fields: an atom is a claim about evaluation, and an ill-typed
$t$ never evaluates, so there the conclusion can fail at valuations the
premise never constrained. \rulename{Ax} never crosses that line because
it introduces no content.
\end{bremark}

\begin{bremark}[structural identity, not logical equivalence]
The checker compares the two formula trees node for node.
$\Phi \vdash \Phi'$ for a logically equivalent but differently shaped
$\Phi'$ --- $B \wedge A$ against $A \wedge B$, say --- is \emph{not} an
\rulename{Ax} instance; where needed it is re-derived through $\wedge$E
and $\wedge$I. In the lattice reading of the structural fragment
(developed at the $\vee$E section), \rulename{Ax} is reflexivity of the
entailment order --- the Gentzen-style leaf \cite{gentzen1935},
discharged semantically rather than by proof search.
\end{bremark}

\begin{anchors}
Theorem \ref{thm:ax-sound} & \texttt{sound\_ax}; payload \texttt{PAxData} (field \texttt{preEqPost}); wired as \texttt{RuleData.ax} into \texttt{ProofObject.sound}\\
ordered-tree formula grammar & \texttt{Formula} (constructors \texttt{and}/\texttt{or}); satisfaction \texttt{Formula.holds}\\
validity statement & \texttt{Sequent.semanticallyValid}\\
lattice reading & \texttt{M3.le\_refl} (the $\vee$E section's countermodel kit)\\
\end{anchors}


\rulesection{2}{Weaken}{core rule}
\label{blcatrule-Weaken}
\ruleblock{The rule (Steele metanotation)}
\sourcenote{as realized in the mechanized development; source row: \texttt{\texttt{rules.toml}} --- \emph{Weaken}. The metanotation statement is typeset by hand for notational quality; its mechanized ground truth is the Lean statement below, extracted from the sources at every build of this book and hash-bound. The rule \emph{roster} is machine-checked (every mechanized rule tag round-trips to exactly one catalogue row); the statement \emph{text} is editorial, checked by eye against the extracted Lean beside it.}
\begin{mathpar}
\inferrule*[Right=\rulename{Weaken}]{\Phi \vdash \Psi}{\Phi \wedge \chi \vdash \Psi}
\end{mathpar}

\par\medskip
\rulename{Weaken} moves a proven sequent under a larger set of
assumptions: from $\Phi \vdash \Psi$ to
$\Phi \wedge \chi \vdash \Psi$. A producer reaches for it when contexts
must be made to agree --- the structural rules with two premises
(\rulename{Cut}, $\wedge$I, $\vee$E) and the composition
rule\textquotesingle s containment check all compare premise trees
structurally, so a sequent proven from fewer assumptions is weakened
into the context the next node expects.

The checker recomputes two tree identities: the
conclusion\textquotesingle s premise is exactly the
sub-derivation\textquotesingle s premise conjoined \emph{on the right}
with the named conjunct $\chi$ (one new $\wedge$-node --- position is
data, since formulas are ordered trees), and the conclusion is carried
over unchanged. Naming $\chi$ in the payload keeps this a comparison,
never a search.

Each hypothesis earns its keep. Drop the premise identity and
weakening could silently \emph{replace} assumptions rather than extend
them --- the proof section\textquotesingle s counterexample derives $\top \vdash w \leq 0$ from the valid
$w \leq 0 \vdash w \leq 0$. Drop the conclusion identity and the
conclusion is free. Nothing at all is demanded of $\chi$ itself: an
unsatisfiable or ill-typed conjunct only makes the new premise harder to
satisfy, which costs applicability, never soundness --- the rule only
discards information.

\ruleblock{Lean statement}
\declhead{BernsteinLogic.Soundness}{sound\_weaken}{\badge{BadgeGreen}{statement bound}}
\begin{leansrc}
theorem sound_weaken {H : Hypernet} {seq premise : Sequent} {extra : Formula}
    (ih : premise.semanticallyValid H)
    (hpre : seq.pre = .and premise.pre extra)
    (hpost : seq.post = premise.post) : seq.semanticallyValid H
\end{leansrc}
\par\medskip
Payload, read as in the \rulename{Ax} section: a sub-derivation plus
structural equalities, nothing else.

\declhead{BernsteinLogic.Core.Catalogue}{PWeakenData}{\badge{BadgeGreen}{statement bound}}
\begin{leansrc}
inductive PWeakenData : Hypernet → Sequent → Type where
  | mk (H : Hypernet) (seq : Sequent)
       (premise : Sequent) (premiseProof : ProofObject H premise)
       (extra : Formula)
       (preEq  : seq.pre  = .and premise.pre extra)
       (postEq : seq.post = premise.post)
       : PWeakenData H seq
\end{leansrc}
\ruleblock{Correspondence (Steele display $\leftrightarrow$ Lean)}
\badge{BadgeGreen}{reviewer-certified against the displayed texts}
\sourcenote{The dictionary below is a reviewer-certified statement of how the mechanized statements above interpret the displayed rule: display metavariable $\mapsto$ Lean binder, each display hypothesis $\mapsto$ the premise or side condition carrying it, the display conclusion $\mapsto$ the goal --- with any deliberate difference in width between display and mechanization named outright. Certification is recorded as a hash triple over (the displayed cell, the extracted statement texts, this dictionary\textquotesingle s own text); if any of the three changes, this block turns amber until a reviewer re-certifies. The correspondence claim itself remains editorial --- reviewer-grade, not machine-checked; its machine-checked successor types each soundness statement against the rule as a structured object, and supersedes this block per rule as it lands.}
Judgment frame as in \rulename{Ax}.
\begin{itemize}
\item the premise judgment $\Phi \vdash \Psi$ $\mapsto$ \texttt{ih}
  $:$ \texttt{premise.semanticallyValid H} --- in the proof-object
  recursion this hypothesis is the sub-derivation\textquotesingle s
  induction hypothesis; the payload (\texttt{PWeakenData}, certified
  alongside) carries the sub-derivation itself as a
  \texttt{ProofObject} field beside the same two ties.
\item $\Phi \mapsto$ \texttt{premise.pre}; $\Psi \mapsto$
  \texttt{premise.post}; $\chi \mapsto$ the binder \texttt{extra :
  Formula}, of which nothing is demanded.
\item conclusion $\Phi \wedge \chi \vdash \Psi$ $\mapsto$ the goal
  under \texttt{hpre} $:$ \texttt{seq.pre = .and premise.pre extra}
  --- the displayed conjunction in the displayed order ($\Phi$ left,
  $\chi$ right; formulas are ordered trees, so the position is part of
  the certified statement) --- and \texttt{hpost} $:$ \texttt{seq.post
  = premise.post}.
\end{itemize}
No width difference in either direction: every hypothesis is consumed,
and the statement proves exactly the displayed rule at every hypernet
and sequent.

\ruleblock{Lean verifier receipts}
\sourcenote{module \texttt{BernsteinLogic.Soundness}; bound build \texttt{a20d90e66c13\ldots} (bound from receipt a9c0e34d3e3c9c8c, 2026-08-03T21:30:36Z). A receipt records that named checking layers accepted a build: (module, compiled-artifact SHA-256, environment SHA-256, checker layers), minted when the checking gate ran; the store receipts attesting the bound build are listed with the layers each run credited. Nothing here asks the reader to trust the store: every layer is an independently runnable program over the quoted sources --- rebuild the module and re-run the checkers to reproduce the verdict. Receipts attesting superseded builds credit no layer and are omitted from this rendering (the count of omissions, if any, is noted below).}
\begin{receipttable}
\receiptrow{a9c0e34d3e3c9c8c}{2026-08-03T21:30:36Z}{build+decorrelated}{\badge{BadgeGreen}{current}}
\end{receipttable}
\sourcenote{21 further receipt(s) in the store attest superseded builds of this module and are omitted from this rendering.}
\sourcenote{module \texttt{BernsteinLogic.Core.Catalogue}; bound build \texttt{cbc92e976a1d\ldots} (bound from receipt 5abdb3589707b5df, 2026-08-03T21:30:33Z). A receipt records that named checking layers accepted a build: (module, compiled-artifact SHA-256, environment SHA-256, checker layers), minted when the checking gate ran; the store receipts attesting the bound build are listed with the layers each run credited. Nothing here asks the reader to trust the store: every layer is an independently runnable program over the quoted sources --- rebuild the module and re-run the checkers to reproduce the verdict. Receipts attesting superseded builds credit no layer and are omitted from this rendering (the count of omissions, if any, is noted below).}
\begin{receipttable}
\receiptrow{24804e66edba0bc9}{2026-08-03T20:25:09Z}{build+decorrelated}{\badge{BadgeGreen}{current}}
\receiptrow{5abdb3589707b5df}{2026-08-03T21:30:33Z}{build+decorrelated}{\badge{BadgeGreen}{current}}
\end{receipttable}
\sourcenote{18 further receipt(s) in the store attest superseded builds of this module and are omitted from this rendering.}
\layerline{\layeritem{build}{lake build (elaboration)}{\badge{BadgeGreen}{current}}\layeritem{replay}{leanchecker (official kernel replay)}{\badge{BadgeAmber}{stale only}}\layeritem{hardened}{safe\_verify (hardened screen)}{\badge{BadgeAmber}{stale only}}\layeritem{decorrelated}{lean4export + nanoda (decorrelated kernel)}{\badge{BadgeGreen}{current}}}
\sourcenote{\badge{BadgeGray}{comparator: none} no fixture in the three-implementation comparator harness (compiled / kernel-interpreted / reference-interpreted) exercises this rule.}
\ruleblock{HOL4 statement}
\sourcenote{theory \texttt{bern\_fam\_resid}, HOL4 development script \texttt{\texttt{bern\_fam\_residScript.sml}} \badge{BadgeGreen}{script bound}}
\holhead{Definition checkFam\_weaken\_def}
\begin{holsrc}
Definition checkFam_weaken_def:
  checkFam_weaken (cols:columns) (ct:ctx) (i:num) ⇔
    case frSeq cols ct i of
    | NONE => F
    | SOME sq =>
      case frChildSeq cols ct i 0 of
      | NONE => F
      | SOME ps =>
        case frFormula cols ct i 0 of
        | NONE => F
        | SOME extra =>
            sq.pre = FAnd ps.pre extra ∧ sq.post = ps.post
End
\end{holsrc}
\holhead{Theorem famEq\_weaken}
\begin{holsrc}
Theorem famEq_weaken:
  wfColumns cols ct ∧ mkCtx cols = SOME ct ∧
  parseCert (cols.hdrImage ++ rest) = SOME (ws, pf) ∧
  nodeCorr cols ct i (CertPf sq nd) ∧
  FST (THE (tagOf cols i)) = 1 ⇒
  (checkFam_weaken cols ct i ⇔ armWeaken ws sq nd)
\end{holsrc}
\ruleblock{HOL verifier receipts}
\badge{BadgeBlue}{documentary} Documentary citations. Family validation: \texttt{bern\_fam\_resid} built GREEN with all ten \texttt{famEq} arms validated, \texttt{hol-strict} clean; spine census: \texttt{bern\_joins} GREEN under \texttt{hol-strict} (55 theorems, 0 offences, 208-theory ancestry including this theory), \texttt{checkJoins\_eq\_checkPf} verified verbatim in the built signature. The cited \texttt{famEq} lemma is the kernel-checked per-node equivalence: under well-formed columns (\texttt{wfColumns}/\texttt{mkCtx}/\texttt{nodeCorr}) at this rule\textquotesingle s wire tag, the columnar family check coincides with the flat reference walk\textquotesingle s arm for this rule. There is no signed HOL receipt store yet --- these citations are documentary.
Semantic stratum, also ported: the soundness twin \texttt{sound\_weaken}
(theory \texttt{bern\_tr\_soundstruct}: \texttt{hol-strict} clean, 9
theorems, 0 offences, verification run of 2026-07-22 on an isolated
build host, content-addressed at the same sources), consumed by the
certified-fragment keystone (\texttt{checkPf\_sound\_frag}, theory
\texttt{bern\_tr\_splice}, 20 theorems, 0 offences, same run): a
reference-checker accept inside the proven wire fragment yields sequent
validity.

\sourcenote{HOL receipts are documentary: quoted results of \texttt{hol-strict} runs
and census rows from the HOL4 development\textquotesingle s build
evidence, not minted receipt artefacts. A HOL receipt store (the Lean
store\textquotesingle s twin) is named future work.
}
\ruleblock{Proof}

\sourcenote{Genre: core rule (tier 1). Setting as in the \rulename{Ax}
section: valuations, admissibility, entrywise atom satisfaction,
semantic validity; counterexamples read over the one-wire hypernet.}

\textbf{Claim (\texttt{sound\_weaken}).}\enspace The payload supplies a
premise sequent $\Phi \vdash \Psi$ with its sub-derivation, the added
conjunct $\chi$ (\texttt{extra}), and two tree identities: the
conclusion's premise is $\Phi \wedge \chi$ (\texttt{preEq}) and its
conclusion is $\Psi$ (\texttt{postEq}). If $\Phi \vdash \Psi$ is
semantically valid, so is $\Phi \wedge \chi \vdash \Psi$.

\begin{bthm}[soundness of \rulename{Weaken}]\label{thm:weaken-sound}
Under \texttt{preEq} and \texttt{postEq}, semantic validity of the
premise sequent implies semantic validity of the conclusion sequent,
over every hypernet.
\end{bthm}
\begin{proof}
Let $\sigma$ be admissible with $\sigma \models \mathrm{seq.pre}$ ---
by \texttt{preEq}, the tree $\Phi \wedge \chi$.

\emph{Step 1 (project).} The $\wedge$-clause of satisfaction is
meta-conjunction, so $\sigma \models \Phi$ and $\sigma \models \chi$;
keep the left component and discard the right --- this rule's whole
activity is licensed forgetting.

\emph{Step 2 (apply the premise).} Premise validity at $\sigma$ gives
$\sigma \models \Psi$, which \texttt{postEq} pins as the conclusion's
post.
\end{proof}

\begin{bremark}[\texttt{preEq} forces extension, not replacement]
\label{rem:weaken-preeq}
The identity $\mathrm{seq.pre} = \Phi \wedge \chi$ makes the conclusion's
premise contain the sub-derivation's premise \emph{verbatim} as the left
conjunct of one new $\wedge$-node, with the named $\chi$ on the right.
Drop it and ``weakening'' could \emph{replace} the assumptions instead of
extending them: from $w \leq 0 \vdash w \leq 0$ (an \rulename{Ax}
instance) conclude $\top \vdash w \leq 0$ --- false at $\sigma(w) = 1$.
This context-swap family is the standard counterexample to every dropped
\texttt{preEq} in the structural sections; we cite it from here on
rather than restating it.
\end{bremark}

\begin{bremark}[positional discipline]
$\chi$ arrives as the right child of a single new $\wedge$-node:
$\Phi \wedge \chi$, never $\chi \wedge \Phi$, and never merged into
$\Phi$'s own tree. Any further placement is $\wedge$E/$\wedge$I work.
Naming $\chi$ in the payload keeps the checker's comparison pure
structure --- one tree identity, no matching or search.
\end{bremark}

\begin{bremark}[$\chi$ is unconstrained, harmlessly]
Like \rulename{Ax}, the rule carries no typing side conditions on what
it touches: $\chi$ may be unsatisfiable ($\bot$), ill-typed, or mention
untyped wires. At worst the conclusion's premise becomes harder or
impossible to satisfy --- vacuous validity, never unsoundness.
Weakening only discards information; rules that must \emph{evaluate}
are where typing fields live.
\end{bremark}

\begin{bremark}[monotonicity, locally and wholesale]
In the lattice reading (the $\vee$E section), meet is a lower bound:
the interpretation of $\Phi \wedge \chi$ lies below that of $\Phi$,
which lies below that of $\Psi$. \rulename{Weaken}
is the one-step premise-monotonicity of the consequence relation; the
composition rule's conjunct-wise premise-containment check
(\rulename{Syzygy}'s \texttt{preContains}) is the same monotonicity used
wholesale.
\end{bremark}

\begin{anchors}
Steps 1--2 & \texttt{sound\_weaken}; payload \texttt{PWeakenData} (fields \texttt{extra}, \texttt{preEq}, \texttt{postEq}); wired as \texttt{RuleData.weaken}\\
$\wedge$-clause of satisfaction & \texttt{Formula.holds} (case \texttt{and})\\
lattice reading & \texttt{M3.weaken\_kit}\\
\end{anchors}


\rulesection{3}{Cut}{core rule}
\label{blcatrule-Cut}
\ruleblock{The rule (Steele metanotation)}
\sourcenote{as realized in the mechanized development; source row: \texttt{\texttt{rules.toml}} --- \emph{Cut}. The metanotation statement is typeset by hand for notational quality; its mechanized ground truth is the Lean statement below, extracted from the sources at every build of this book and hash-bound. The rule \emph{roster} is machine-checked (every mechanized rule tag round-trips to exactly one catalogue row); the statement \emph{text} is editorial, checked by eye against the extracted Lean beside it.}
\begin{mathpar}
\inferrule*[Right=\rulename{Cut}]{\Phi \vdash \chi \\ \Phi \wedge \chi \vdash \Psi}{\Phi \vdash \Psi}
\end{mathpar}

\par\medskip
\rulename{Cut} is the lemma-staging rule: prove $\Phi \vdash \chi$
once, prove $\Phi \wedge \chi \vdash \Psi$ under the enlarged context,
conclude $\Phi \vdash \Psi$. A producer uses it to factor a
certificate the way a paper factors a proof --- a bound or region fact
is established once and then consumed as an assumption downstream;
iterated cuts stage a chain of such lemmas.

The checker recomputes four tree identities: the
conclusion\textquotesingle s premise and conclusion are inherited from
the left and right sub-derivations respectively; the left
sub-derivation concludes exactly the named cut formula $\chi$; and the
right sub-derivation\textquotesingle s premise is exactly the left
one\textquotesingle s premise conjoined with $\chi$ --- the same
context tree, verbatim. The payload must name $\chi$ explicitly
because the conclusion mentions it nowhere: \rulename{Cut} is the one
structural rule whose certified content is invisible in its
conclusion, so there is nothing to reconstruct it from.

The two interesting hypotheses fail independently, with near-twin
counterexamples walked in the proof section: unshare the context and a
fact proven under foreign assumptions leaks into the conclusion
($\top \vdash w \leq 0$ from two individually valid premises); unmoor
$\chi$ from the left conclusion and the right premise consumes an
assumption nobody earned (same false conclusion, different route). The
proof also records why no cut-elimination enters anywhere: soundness of
the rule is direct, and the checker never normalizes a proof.

\ruleblock{Lean statement}
\declhead{BernsteinLogic.Soundness}{sound\_cut}{\badge{BadgeGreen}{statement bound}}
\begin{leansrc}
theorem sound_cut {H : Hypernet} {seq leftPremise rightPremise : Sequent}
    {cutFormula : Formula}
    (ihl : leftPremise.semanticallyValid H)
    (ihr : rightPremise.semanticallyValid H)
    (hpre : seq.pre = leftPremise.pre)
    (hpost : seq.post = rightPremise.post)
    (hcutL : leftPremise.post = cutFormula)
    (hcutR : rightPremise.pre = .and leftPremise.pre cutFormula) :
    seq.semanticallyValid H
\end{leansrc}
\par\medskip
Payload, read as in the \rulename{Ax} section: sub-derivations plus
structural equalities, with the cut formula as explicit data.

\declhead{BernsteinLogic.Core.Catalogue}{PCutData}{\badge{BadgeGreen}{statement bound}}
\begin{leansrc}
inductive PCutData : Hypernet → Sequent → Type where
  | mk (H : Hypernet) (seq : Sequent)
       (cutFormula : Formula)
       (leftPremise  : Sequent) (leftProof  : ProofObject H leftPremise)
       (rightPremise : Sequent) (rightProof : ProofObject H rightPremise)
       (preEq    : seq.pre  = leftPremise.pre)
       (postEq   : seq.post = rightPremise.post)
       (cutLeft  : leftPremise.post = cutFormula)
       (cutRight : rightPremise.pre = .and leftPremise.pre cutFormula)
       : PCutData H seq
\end{leansrc}
\ruleblock{Correspondence (Steele display $\leftrightarrow$ Lean)}
\badge{BadgeGreen}{reviewer-certified against the displayed texts}
\sourcenote{The dictionary below is a reviewer-certified statement of how the mechanized statements above interpret the displayed rule: display metavariable $\mapsto$ Lean binder, each display hypothesis $\mapsto$ the premise or side condition carrying it, the display conclusion $\mapsto$ the goal --- with any deliberate difference in width between display and mechanization named outright. Certification is recorded as a hash triple over (the displayed cell, the extracted statement texts, this dictionary\textquotesingle s own text); if any of the three changes, this block turns amber until a reviewer re-certifies. The correspondence claim itself remains editorial --- reviewer-grade, not machine-checked; its machine-checked successor types each soundness statement against the rule as a structured object, and supersedes this block per rule as it lands.}
Judgment frame as in \rulename{Ax}; sub-derivation premises as at
\rulename{Weaken}.
\begin{itemize}
\item the two premise judgments $\Phi \vdash \chi$ and $\Phi \wedge
  \chi \vdash \Psi$ $\mapsto$ \texttt{ihl} $:$
  \texttt{leftPremise.semanticallyValid H} and \texttt{ihr} $:$
  \texttt{rightPremise.semanticallyValid H}.
\item $\Phi \mapsto$ \texttt{leftPremise.pre}; $\Psi \mapsto$
  \texttt{rightPremise.post}; $\chi \mapsto$ the binder
  \texttt{cutFormula}, tied to the left conclusion by \texttt{hcutL}
  $:$ \texttt{leftPremise.post = cutFormula} --- the explicit naming
  the commentary explains (the conclusion mentions $\chi$ nowhere).
\item the second premise\textquotesingle s displayed context $\Phi
  \wedge \chi$ $\mapsto$ \texttt{hcutR} $:$ \texttt{rightPremise.pre =
  .and leftPremise.pre cutFormula} --- the same context tree verbatim,
  cut formula conjoined on the right.
\item conclusion $\Phi \vdash \Psi$ $\mapsto$ the goal under
  \texttt{hpre} $:$ \texttt{seq.pre = leftPremise.pre} and
  \texttt{hpost} $:$ \texttt{seq.post = rightPremise.post}.
\end{itemize}
No width difference in either direction: the six hypotheses are
exactly the display\textquotesingle s two premises and four structural
identities, each consumed.

\ruleblock{Lean verifier receipts}
\sourcenote{module \texttt{BernsteinLogic.Soundness}; bound build \texttt{a20d90e66c13\ldots} (bound from receipt a9c0e34d3e3c9c8c, 2026-08-03T21:30:36Z). A receipt records that named checking layers accepted a build: (module, compiled-artifact SHA-256, environment SHA-256, checker layers), minted when the checking gate ran; the store receipts attesting the bound build are listed with the layers each run credited. Nothing here asks the reader to trust the store: every layer is an independently runnable program over the quoted sources --- rebuild the module and re-run the checkers to reproduce the verdict. Receipts attesting superseded builds credit no layer and are omitted from this rendering (the count of omissions, if any, is noted below).}
\begin{receipttable}
\receiptrow{a9c0e34d3e3c9c8c}{2026-08-03T21:30:36Z}{build+decorrelated}{\badge{BadgeGreen}{current}}
\end{receipttable}
\sourcenote{21 further receipt(s) in the store attest superseded builds of this module and are omitted from this rendering.}
\sourcenote{module \texttt{BernsteinLogic.Core.Catalogue}; bound build \texttt{cbc92e976a1d\ldots} (bound from receipt 5abdb3589707b5df, 2026-08-03T21:30:33Z). A receipt records that named checking layers accepted a build: (module, compiled-artifact SHA-256, environment SHA-256, checker layers), minted when the checking gate ran; the store receipts attesting the bound build are listed with the layers each run credited. Nothing here asks the reader to trust the store: every layer is an independently runnable program over the quoted sources --- rebuild the module and re-run the checkers to reproduce the verdict. Receipts attesting superseded builds credit no layer and are omitted from this rendering (the count of omissions, if any, is noted below).}
\begin{receipttable}
\receiptrow{24804e66edba0bc9}{2026-08-03T20:25:09Z}{build+decorrelated}{\badge{BadgeGreen}{current}}
\receiptrow{5abdb3589707b5df}{2026-08-03T21:30:33Z}{build+decorrelated}{\badge{BadgeGreen}{current}}
\end{receipttable}
\sourcenote{18 further receipt(s) in the store attest superseded builds of this module and are omitted from this rendering.}
\layerline{\layeritem{build}{lake build (elaboration)}{\badge{BadgeGreen}{current}}\layeritem{replay}{leanchecker (official kernel replay)}{\badge{BadgeAmber}{stale only}}\layeritem{hardened}{safe\_verify (hardened screen)}{\badge{BadgeAmber}{stale only}}\layeritem{decorrelated}{lean4export + nanoda (decorrelated kernel)}{\badge{BadgeGreen}{current}}}
\sourcenote{\badge{BadgeGray}{comparator: none} no fixture in the three-implementation comparator harness (compiled / kernel-interpreted / reference-interpreted) exercises this rule.}
\ruleblock{HOL4 statement}
\sourcenote{theory \texttt{bern\_fam\_resid}, HOL4 development script \texttt{\texttt{bern\_fam\_residScript.sml}} \badge{BadgeGreen}{script bound}}
\holhead{Definition checkFam\_cut\_def}
\begin{holsrc}
Definition checkFam_cut_def:
  checkFam_cut (cols:columns) (ct:ctx) (i:num) ⇔
    case frSeq cols ct i of
    | NONE => F
    | SOME sq =>
      case frChildSeq cols ct i 0 of
      | NONE => F
      | SOME ls =>
        case frChildSeq cols ct i 1 of
        | NONE => F
        | SOME rs =>
          case frFormula cols ct i 0 of
          | NONE => F
          | SOME cutf =>
              sq.pre = ls.pre ∧ sq.post = rs.post ∧
              ls.post = cutf ∧ rs.pre = FAnd ls.pre cutf
End
\end{holsrc}
\holhead{Theorem famEq\_cut}
\begin{holsrc}
Theorem famEq_cut:
  wfColumns cols ct ∧ mkCtx cols = SOME ct ∧
  parseCert (cols.hdrImage ++ rest) = SOME (ws, pf) ∧
  nodeCorr cols ct i (CertPf sq nd) ∧
  FST (THE (tagOf cols i)) = 2 ⇒
  (checkFam_cut cols ct i ⇔ armCut ws sq nd)
\end{holsrc}
\ruleblock{HOL verifier receipts}
\badge{BadgeBlue}{documentary} Documentary citations. Family validation: \texttt{bern\_fam\_resid} built GREEN with all ten \texttt{famEq} arms validated, \texttt{hol-strict} clean; spine census: \texttt{bern\_joins} GREEN under \texttt{hol-strict} (55 theorems, 0 offences, 208-theory ancestry including this theory), \texttt{checkJoins\_eq\_checkPf} verified verbatim in the built signature. The cited \texttt{famEq} lemma is the kernel-checked per-node equivalence: under well-formed columns (\texttt{wfColumns}/\texttt{mkCtx}/\texttt{nodeCorr}) at this rule\textquotesingle s wire tag, the columnar family check coincides with the flat reference walk\textquotesingle s arm for this rule. There is no signed HOL receipt store yet --- these citations are documentary.
Semantic stratum, also ported: the soundness twin \texttt{sound\_cut}
(theory \texttt{bern\_tr\_soundstruct}: \texttt{hol-strict} clean, 9
theorems, 0 offences, verification run of 2026-07-22 on an isolated
build host, content-addressed at the same sources), consumed by the
certified-fragment keystone (\texttt{checkPf\_sound\_frag}, theory
\texttt{bern\_tr\_splice}, 20 theorems, 0 offences, same run): a
reference-checker accept inside the proven wire fragment yields sequent
validity.

\sourcenote{HOL receipts are documentary: quoted results of \texttt{hol-strict} runs
and census rows from the HOL4 development\textquotesingle s build
evidence, not minted receipt artefacts. A HOL receipt store (the Lean
store\textquotesingle s twin) is named future work.
}
\ruleblock{Proof}

\sourcenote{Genre: core rule (tier 1). Setting as in the \rulename{Ax}
section; counterexamples read over the one-wire hypernet.}

\textbf{Claim (\texttt{sound\_cut}).}\enspace The payload supplies the
cut formula $\chi$ (\texttt{cutFormula}), a left premise
$\Phi \vdash \chi$ and a right premise $\Phi \wedge \chi \vdash \Psi$
with their sub-derivations, and four tree identities: the conclusion's
premise is the left premise's premise (\texttt{preEq}), its conclusion
is the right premise's conclusion (\texttt{postEq}), the left premise
concludes exactly $\chi$ (\texttt{cutLeft}), and the right premise's
premise is exactly $\Phi \wedge \chi$ over the \emph{same} $\Phi$
(\texttt{cutRight}). If both premise sequents are semantically valid,
so is $\Phi \vdash \Psi$.

\begin{bthm}[soundness of \rulename{Cut}]\label{thm:cut-sound}
Under the four payload identities, semantic validity of both premise
sequents implies semantic validity of the conclusion, over every
hypernet.
\end{bthm}
\begin{proof}
Let $\sigma$ be admissible with $\sigma \models \Phi$.

\emph{Step 1 (earn the cut formula).} By \texttt{preEq}, $\sigma$
satisfies the left premise's premise; left validity gives
$\sigma \models \chi$ (\texttt{cutLeft} identifies the left conclusion
with $\chi$).

\emph{Step 2 (assemble the enlarged context).} The $\wedge$-clause of
satisfaction is meta-conjunction, so the standing assumption and Step 1
give $\sigma \models \Phi \wedge \chi$ --- and by \texttt{cutRight} this tree
\emph{is} the right premise's premise.

\emph{Step 3 (apply the right premise).} Right validity gives
$\sigma \models \Psi$, pinned as the conclusion's post by
\texttt{postEq}.
\end{proof}

\begin{bremark}[context sharing is load-bearing]\label{rem:cut-sharing}
\texttt{cutRight} forces the right premise's context to be the left
premise's context --- the context-sharing (\emph{additive}) format
\cite{girard1987}. Allow the right context to be foreign, say
$\Phi' \wedge \chi$, and cut breaks: take the left premise
$\top \vdash \top$ (an \rulename{Ax} instance) and the right premise
$(w \leq 0) \wedge \top \vdash w \leq 0$ (valid: satisfaction of the
conjunction projects); the degenerate rule would conclude
$\top \vdash w \leq 0$ --- false at $\sigma(w) = 1$.
\end{bremark}

\begin{bremark}[\texttt{cutLeft} moors the cut formula]
Keep the sharing but unmoor $\chi$ from what the left premise actually
proved, and the right premise may consume an assumption nobody earned:
from the left premise $\top \vdash \top$ and the right premise
$\top \wedge (w \leq 0) \vdash w \leq 0$ (valid, again by projection),
the degenerate rule would conclude $\top \vdash w \leq 0$. The two
counterexamples are deliberately near-twins: \texttt{cutRight} guards
\emph{whose context}, \texttt{cutLeft} guards \emph{whose lemma}, and
each fails independently. \texttt{preEq}'s removal admits the
context-swap family (the \rulename{Weaken} section); \texttt{postEq}'s
removal frees the conclusion outright.
\end{bremark}

\begin{bremark}[the cut formula leaves no trace]
The conclusion $\Phi \vdash \Psi$ mentions $\chi$ nowhere --- the
conclusion does not determine the certified content (a hiddenness
\rulename{Cut} shares with $\wedge$E's discarded conjunct), so the
payload names \texttt{cutFormula} explicitly. That is a design choice
rather than a forced move --- the checker could equally read $\chi$
off the left premise's post --- and the explicit field keeps the check
local. The same invisibility is what the rule is \emph{for}:
certificate producers stage lemmas --- prove a bound or region fact
once, then continue under the enlarged context $\Phi \wedge \chi$ ---
and iterated cuts stage a chain of them.
\end{bremark}

\begin{bremark}[semantic cut costs nothing]
Gentzen's Hauptsatz \cite{gentzen1935} eliminates cut from derivations
at, in general, non-elementary blow-up in proof size
\cite{statman1979,orevkov1979}; nothing here needs it. The
soundness argument is the three steps above, the checker never
normalizes a proof, and the calculus keeps \rulename{Cut} as a
first-class node because producers want staged lemmas, not cut-free
normal forms. In the lattice reading: if $x \leq m$ then $x = x \sqcap
m$, so $x \sqcap m \leq y$ gives $x \leq y$ --- transitivity worn as
meet bookkeeping.
\end{bremark}

\begin{anchors}
Steps 1--3 & \texttt{sound\_cut}; payload \texttt{PCutData} (fields \texttt{cutFormula}, \texttt{preEq}, \texttt{postEq}, \texttt{cutLeft}, \texttt{cutRight}); wired as \texttt{RuleData.cut}\\
$\wedge$-clause of satisfaction & \texttt{Formula.holds} (case \texttt{and})\\
lattice reading & \texttt{M3.cut\_kit}\\
\end{anchors}


\rulesection{4}{$\wedge$I}{core rule}
\label{blcatrule-AndI}
\ruleblock{The rule (Steele metanotation)}
\sourcenote{as realized in the mechanized development; source row: \texttt{\texttt{rules.toml}} --- \emph{$\wedge$I}. The metanotation statement is typeset by hand for notational quality; its mechanized ground truth is the Lean statement below, extracted from the sources at every build of this book and hash-bound. The rule \emph{roster} is machine-checked (every mechanized rule tag round-trips to exactly one catalogue row); the statement \emph{text} is editorial, checked by eye against the extracted Lean beside it.}
\begin{mathpar}
\inferrule*[Right=$\wedge$I]{\Phi \vdash \Psi_1 \\ \Phi \vdash \Psi_2}{\Phi \vdash \Psi_1 \wedge \Psi_2}
\end{mathpar}

\par\medskip
$\wedge$I pairs two conclusions proven from the \emph{same} assumptions:
from $\Phi \vdash \Psi_1$ and $\Phi \vdash \Psi_2$, conclude
$\Phi \vdash \Psi_1 \wedge \Psi_2$. A producer uses it to assemble
composite claims from separately established facts --- two one-sided
bounds into a two-sided one, the conjuncts a later
node\textquotesingle s premise-form expects --- and, with $\wedge$E, to
rebuild conjunction trees in the shape a structural comparison demands
(the core never reassociates a conjunction on its own).

The checker recomputes three tree identities: the two
sub-derivations\textquotesingle{} premises are the same tree, the
conclusion inherits that premise, and the conclusion\textquotesingle s
post is one new $\wedge$-node with the left sub-derivation\textquotesingle s
conclusion on the left. Order is data: $\Psi_1 \wedge \Psi_2$ and
$\Psi_2 \wedge \Psi_1$ are different conclusions, obtained by swapping
which sub-derivation rides which payload slot.

The context-sharing identity is the load-bearing one. Drop it and the
rule pairs facts earned under different assumptions, keeping only one
set: from the valid $\top \vdash \top$ and $w \leq 0 \vdash w \leq 0$
it would conclude $\top \vdash \top \wedge (w \leq 0)$, false wherever
$w$ is positive. The proof section also records the structural
observation the $\vee$E section turns into a theorem: $\wedge$I and
$\wedge$E are the two directions of the universal property of meet, so
the conjunction fragment adds nothing beyond the lattice structure of
entailment.

\ruleblock{Lean statement}
\declhead{BernsteinLogic.Soundness}{sound\_andI}{\badge{BadgeGreen}{statement bound}}
\begin{leansrc}
theorem sound_andI {H : Hypernet} {seq leftPremise rightPremise : Sequent}
    (ihl : leftPremise.semanticallyValid H)
    (ihr : rightPremise.semanticallyValid H)
    (hpreMatch : leftPremise.pre = rightPremise.pre)
    (hpre : seq.pre = leftPremise.pre)
    (hpost : seq.post = .and leftPremise.post rightPremise.post) :
    seq.semanticallyValid H
\end{leansrc}
\par\medskip
Payload, read as in the \rulename{Ax} section: sub-derivations plus
structural equalities, nothing else.

\declhead{BernsteinLogic.Core.Catalogue}{PAndIData}{\badge{BadgeGreen}{statement bound}}
\begin{leansrc}
inductive PAndIData : Hypernet → Sequent → Type where
  | mk (H : Hypernet) (seq : Sequent)
       (leftPremise  : Sequent) (leftProof  : ProofObject H leftPremise)
       (rightPremise : Sequent) (rightProof : ProofObject H rightPremise)
       (preMatch : leftPremise.pre = rightPremise.pre)
       (preEq    : seq.pre  = leftPremise.pre)
       (postEq   : seq.post = .and leftPremise.post rightPremise.post)
       : PAndIData H seq
\end{leansrc}
\ruleblock{Correspondence (Steele display $\leftrightarrow$ Lean)}
\badge{BadgeGreen}{reviewer-certified against the displayed texts}
\sourcenote{The dictionary below is a reviewer-certified statement of how the mechanized statements above interpret the displayed rule: display metavariable $\mapsto$ Lean binder, each display hypothesis $\mapsto$ the premise or side condition carrying it, the display conclusion $\mapsto$ the goal --- with any deliberate difference in width between display and mechanization named outright. Certification is recorded as a hash triple over (the displayed cell, the extracted statement texts, this dictionary\textquotesingle s own text); if any of the three changes, this block turns amber until a reviewer re-certifies. The correspondence claim itself remains editorial --- reviewer-grade, not machine-checked; its machine-checked successor types each soundness statement against the rule as a structured object, and supersedes this block per rule as it lands.}
Judgment frame as in \rulename{Ax}; sub-derivation premises as at
\rulename{Weaken}.
\begin{itemize}
\item the two premise judgments $\mapsto$ \texttt{ihl}/\texttt{ihr}
  over the binders \texttt{leftPremise}/\texttt{rightPremise}.
\item the display writes one context $\Phi$ twice; the statement binds
  two sequents and demands \texttt{hpreMatch} $:$
  \texttt{leftPremise.pre = rightPremise.pre} --- sharing by equation
  rather than by a common binder, the pattern of every two-premise
  rule on this page.
\item $\Phi \mapsto$ \texttt{leftPremise.pre} (hence, through the
  match, \texttt{rightPremise.pre}); $\Psi_1 \mapsto$
  \texttt{leftPremise.post}; $\Psi_2 \mapsto$
  \texttt{rightPremise.post}.
\item conclusion $\Phi \vdash \Psi_1 \wedge \Psi_2$ $\mapsto$ the goal
  under \texttt{hpre} $:$ \texttt{seq.pre = leftPremise.pre} and
  \texttt{hpost} $:$ \texttt{seq.post = .and leftPremise.post
  rightPremise.post} --- left sub-derivation on the left, as
  displayed.
\end{itemize}
No width difference in either direction; every hypothesis is consumed.

\ruleblock{Lean verifier receipts}
\sourcenote{module \texttt{BernsteinLogic.Soundness}; bound build \texttt{a20d90e66c13\ldots} (bound from receipt a9c0e34d3e3c9c8c, 2026-08-03T21:30:36Z). A receipt records that named checking layers accepted a build: (module, compiled-artifact SHA-256, environment SHA-256, checker layers), minted when the checking gate ran; the store receipts attesting the bound build are listed with the layers each run credited. Nothing here asks the reader to trust the store: every layer is an independently runnable program over the quoted sources --- rebuild the module and re-run the checkers to reproduce the verdict. Receipts attesting superseded builds credit no layer and are omitted from this rendering (the count of omissions, if any, is noted below).}
\begin{receipttable}
\receiptrow{a9c0e34d3e3c9c8c}{2026-08-03T21:30:36Z}{build+decorrelated}{\badge{BadgeGreen}{current}}
\end{receipttable}
\sourcenote{21 further receipt(s) in the store attest superseded builds of this module and are omitted from this rendering.}
\sourcenote{module \texttt{BernsteinLogic.Core.Catalogue}; bound build \texttt{cbc92e976a1d\ldots} (bound from receipt 5abdb3589707b5df, 2026-08-03T21:30:33Z). A receipt records that named checking layers accepted a build: (module, compiled-artifact SHA-256, environment SHA-256, checker layers), minted when the checking gate ran; the store receipts attesting the bound build are listed with the layers each run credited. Nothing here asks the reader to trust the store: every layer is an independently runnable program over the quoted sources --- rebuild the module and re-run the checkers to reproduce the verdict. Receipts attesting superseded builds credit no layer and are omitted from this rendering (the count of omissions, if any, is noted below).}
\begin{receipttable}
\receiptrow{24804e66edba0bc9}{2026-08-03T20:25:09Z}{build+decorrelated}{\badge{BadgeGreen}{current}}
\receiptrow{5abdb3589707b5df}{2026-08-03T21:30:33Z}{build+decorrelated}{\badge{BadgeGreen}{current}}
\end{receipttable}
\sourcenote{18 further receipt(s) in the store attest superseded builds of this module and are omitted from this rendering.}
\layerline{\layeritem{build}{lake build (elaboration)}{\badge{BadgeGreen}{current}}\layeritem{replay}{leanchecker (official kernel replay)}{\badge{BadgeAmber}{stale only}}\layeritem{hardened}{safe\_verify (hardened screen)}{\badge{BadgeAmber}{stale only}}\layeritem{decorrelated}{lean4export + nanoda (decorrelated kernel)}{\badge{BadgeGreen}{current}}}
\sourcenote{\badge{BadgeGray}{comparator: none} no fixture in the three-implementation comparator harness (compiled / kernel-interpreted / reference-interpreted) exercises this rule.}
\ruleblock{HOL4 statement}
\sourcenote{theory \texttt{bern\_fam\_resid}, HOL4 development script \texttt{\texttt{bern\_fam\_residScript.sml}} \badge{BadgeGreen}{script bound}}
\holhead{Definition checkFam\_andi\_def}
\begin{holsrc}
Definition checkFam_andi_def:
  checkFam_andi (cols:columns) (ct:ctx) (i:num) ⇔
    case frSeq cols ct i of
    | NONE => F
    | SOME sq =>
      case frChildSeq cols ct i 0 of
      | NONE => F
      | SOME ls =>
        case frChildSeq cols ct i 1 of
        | NONE => F
        | SOME rs =>
            ls.pre = rs.pre ∧ sq.pre = ls.pre ∧
            sq.post = FAnd ls.post rs.post
End
\end{holsrc}
\holhead{Theorem famEq\_andi}
\begin{holsrc}
Theorem famEq_andi:
  wfColumns cols ct ∧ mkCtx cols = SOME ct ∧
  parseCert (cols.hdrImage ++ rest) = SOME (ws, pf) ∧
  nodeCorr cols ct i (CertPf sq nd) ∧
  FST (THE (tagOf cols i)) = 3 ⇒
  (checkFam_andi cols ct i ⇔ armAndI ws sq nd)
\end{holsrc}
\ruleblock{HOL verifier receipts}
\badge{BadgeBlue}{documentary} Documentary citations. Family validation: \texttt{bern\_fam\_resid} built GREEN with all ten \texttt{famEq} arms validated, \texttt{hol-strict} clean; spine census: \texttt{bern\_joins} GREEN under \texttt{hol-strict} (55 theorems, 0 offences, 208-theory ancestry including this theory), \texttt{checkJoins\_eq\_checkPf} verified verbatim in the built signature. The cited \texttt{famEq} lemma is the kernel-checked per-node equivalence: under well-formed columns (\texttt{wfColumns}/\texttt{mkCtx}/\texttt{nodeCorr}) at this rule\textquotesingle s wire tag, the columnar family check coincides with the flat reference walk\textquotesingle s arm for this rule. There is no signed HOL receipt store yet --- these citations are documentary.
Semantic stratum, also ported: the soundness twin \texttt{sound\_andI}
(theory \texttt{bern\_tr\_soundstruct}: \texttt{hol-strict} clean, 9
theorems, 0 offences, verification run of 2026-07-22 on an isolated
build host, content-addressed at the same sources), consumed by the
certified-fragment keystone (\texttt{checkPf\_sound\_frag}, theory
\texttt{bern\_tr\_splice}, 20 theorems, 0 offences, same run): a
reference-checker accept inside the proven wire fragment yields sequent
validity.

\sourcenote{HOL receipts are documentary: quoted results of \texttt{hol-strict} runs
and census rows from the HOL4 development\textquotesingle s build
evidence, not minted receipt artefacts. A HOL receipt store (the Lean
store\textquotesingle s twin) is named future work.
}
\ruleblock{Proof}

\sourcenote{Genre: core rule (tier 1). Setting as in the \rulename{Ax}
section; counterexamples read over the one-wire hypernet.}

\textbf{Claim (\texttt{sound\_andI}).}\enspace The payload supplies two
premise sequents $\Phi \vdash \Psi_1$ and $\Phi \vdash \Psi_2$ with
their sub-derivations and three tree identities: the two premises share
one context (\texttt{preMatch}), the conclusion's premise is that
context (\texttt{preEq}), and the conclusion's post is
$\Psi_1 \wedge \Psi_2$ with the left premise's conclusion as left
conjunct (\texttt{postEq}). If both premise sequents are semantically
valid, so is $\Phi \vdash \Psi_1 \wedge \Psi_2$.

\begin{bthm}[soundness of $\wedge$I]\label{thm:andi-sound}
Under the three payload identities, semantic validity of both premise
sequents implies semantic validity of the conclusion, over every
hypernet.
\end{bthm}
\begin{proof}
Let $\sigma$ be admissible with $\sigma \models \Phi$.

\emph{Step 1 (left component).} By \texttt{preEq}, $\sigma$ satisfies
the left premise's premise; left validity gives $\sigma \models \Psi_1$.

\emph{Step 2 (right component).} \texttt{preMatch} transports
$\sigma \models \Phi$ verbatim to the right premise's premise --- the
same tree, so the same satisfaction fact; right validity gives
$\sigma \models \Psi_2$.

\emph{Step 3 (pair).} The $\wedge$-clause of satisfaction is
meta-conjunction, so $\sigma \models \Psi_1 \wedge \Psi_2$, pinned by
\texttt{postEq}.
\end{proof}

\begin{bremark}[\texttt{preMatch} is load-bearing]\label{rem:andi-share}
The rule may only pair conclusions drawn from the \emph{same}
assumptions --- context sharing again, as at \rulename{Cut}
\cite{girard1987}. Drop \texttt{preMatch} and take the left premise
$\top \vdash \top$ and the right premise $w \leq 0 \vdash w \leq 0$
(both \rulename{Ax} instances): the degenerate rule concludes
$\top \vdash \top \wedge (w \leq 0)$ --- false at $\sigma(w) = 1$,
because the right conjunct was earned under an assumption the
conclusion no longer carries. \texttt{preEq}'s removal admits the
context-swap family (the \rulename{Weaken} section); \texttt{postEq}'s
removal frees the conclusion outright.
\end{bremark}

\begin{bremark}[order and association are data]
\texttt{postEq} places the left premise's conclusion on the left of one
new $\wedge$-node. The conjunction is an ordered tree: nothing in the
core commutes or reassociates it, so $\Psi_1 \wedge \Psi_2$ and
$\Psi_2 \wedge \Psi_1$ are different conclusions, each derivable from
the same premises by choosing which sub-derivation rides which payload
slot.
\end{bremark}

\begin{bremark}[the meet adjunction, split across two rules]
\label{rem:andi-meet}
In the lattice reading, $\wedge$I is one direction of the universal
property of meet --- $x \leq y$ and $x \leq z$ iff
$x \leq y \sqcap z$ --- and $\wedge$E is the other. Together the two
rules say the $\wedge$-fragment of the calculus adds nothing beyond the
lattice structure of entailment; the $\vee$E section makes this precise
(and load-bearing) through a countermodel over exactly that structure.
With $\wedge$E the pair is also computationally lossless at the
formula-tree level: the two projections undo pairing, and pairing the
two projections rebuilds the identical formula tree.
\end{bremark}

\begin{anchors}
Steps 1--3 & \texttt{sound\_andI}; payload \texttt{PAndIData} (fields \texttt{preMatch}, \texttt{preEq}, \texttt{postEq}); wired as \texttt{RuleData.andI}\\
$\wedge$-clause of satisfaction & \texttt{Formula.holds} (case \texttt{and})\\
lattice reading & \texttt{M3.andI\_kit}\\
\end{anchors}


\rulesection{5}{$\wedge$E}{core rule}
\label{blcatrule-AndE}
\ruleblock{The rule (Steele metanotation)}
\sourcenote{as realized in the mechanized development; source row: \texttt{\texttt{rules.toml}} --- \emph{$\wedge$E}. The metanotation statement is typeset by hand for notational quality; its mechanized ground truth is the Lean statement below, extracted from the sources at every build of this book and hash-bound. The rule \emph{roster} is machine-checked (every mechanized rule tag round-trips to exactly one catalogue row); the statement \emph{text} is editorial, checked by eye against the extracted Lean beside it.}
\begin{mathpar}
\inferrule*[Right=$\wedge$E]{\Phi \vdash \Psi_1 \wedge \Psi_2}{\Phi \vdash \Psi_i}
\end{mathpar}

\par\medskip
$\wedge$E projects a conjunction: from
$\Phi \vdash \Psi_1 \wedge \Psi_2$, conclude $\Phi \vdash \Psi_i$. A
producer uses it to take composite facts apart --- the demonstration
certificate applies it right after its axiom leaf, extracting from a
composite premise the region conjunct that the Bernstein-cone closure
needs --- and, with $\wedge$I, to reshape conjunction trees for the
structural comparisons other rules perform.

The rule is realized as two constructors, one per side; the
paper\textquotesingle s index $i \in \{1, 2\}$ is the constructor
choice itself, so there is no unchecked side flag. The checker
recomputes two tree identities: shared context, and that the
sub-derivation\textquotesingle s conclusion is exactly the
conclusion\textquotesingle s post conjoined with the \emph{named}
discarded conjunct --- on the right for the left constructor, on the
left for the right one. Naming the discarded conjunct in the payload
makes the side condition a single comparison; the checker never
searches a conjunction for the target, and a certificate claiming the
wrong side simply fails the identity.

Dropping the conclusion identity frees the conclusion outright;
dropping the context identity admits the context-swap counterexamples
walked at \rulename{Weaken}. The projection itself needs no defence:
discarding a conjunct only weakens a claim, and the proof section
notes the pairing with $\wedge$I under which the fragment is lossless
at the formula-tree level.

\ruleblock{Lean statement}
\declhead{BernsteinLogic.Soundness}{sound\_andE\_left}{\badge{BadgeGreen}{statement bound}}
\begin{leansrc}
theorem sound_andE_left {H : Hypernet} {seq premise : Sequent} {other : Formula}
    (ih : premise.semanticallyValid H)
    (hpre : seq.pre = premise.pre)
    (hpost : premise.post = .and seq.post other) :
    seq.semanticallyValid H
\end{leansrc}
\declhead{BernsteinLogic.Soundness}{sound\_andE\_right}{\badge{BadgeGreen}{statement bound}}
\begin{leansrc}
theorem sound_andE_right {H : Hypernet} {seq premise : Sequent} {other : Formula}
    (ih : premise.semanticallyValid H)
    (hpre : seq.pre = premise.pre)
    (hpost : premise.post = .and other seq.post) :
    seq.semanticallyValid H
\end{leansrc}
\par\medskip
Payload, read as in the \rulename{Ax} section: one constructor per
projection, each a sub-derivation plus structural equalities with the
discarded conjunct as explicit data.

\declhead{BernsteinLogic.Core.Catalogue}{PAndEData}{\badge{BadgeGreen}{statement bound}}
\begin{leansrc}
inductive PAndEData : Hypernet → Sequent → Type where
  | left  (H : Hypernet) (seq : Sequent)
          (premise : Sequent) (premiseProof : ProofObject H premise)
          (other : Formula)
          (preEq  : seq.pre = premise.pre)
          (postEq : premise.post = .and seq.post other)
          : PAndEData H seq
  | right (H : Hypernet) (seq : Sequent)
          (premise : Sequent) (premiseProof : ProofObject H premise)
          (other : Formula)
          (preEq  : seq.pre = premise.pre)
          (postEq : premise.post = .and other seq.post)
          : PAndEData H seq
\end{leansrc}
\ruleblock{Correspondence (Steele display $\leftrightarrow$ Lean)}
\badge{BadgeGreen}{reviewer-certified against the displayed texts}
\sourcenote{The dictionary below is a reviewer-certified statement of how the mechanized statements above interpret the displayed rule: display metavariable $\mapsto$ Lean binder, each display hypothesis $\mapsto$ the premise or side condition carrying it, the display conclusion $\mapsto$ the goal --- with any deliberate difference in width between display and mechanization named outright. Certification is recorded as a hash triple over (the displayed cell, the extracted statement texts, this dictionary\textquotesingle s own text); if any of the three changes, this block turns amber until a reviewer re-certifies. The correspondence claim itself remains editorial --- reviewer-grade, not machine-checked; its machine-checked successor types each soundness statement against the rule as a structured object, and supersedes this block per rule as it lands.}
Judgment frame as in \rulename{Ax}; sub-derivation premise as at
\rulename{Weaken}.
\begin{itemize}
\item the display\textquotesingle s index $i \in \{1, 2\}$ $\mapsto$
  the \emph{choice of lemma}: \texttt{sound\_andE\_left} is $i = 1$,
  \texttt{sound\_andE\_right} is $i = 2$, mirroring the
  payload\textquotesingle s two constructors --- the side is a
  constructor choice, never an unchecked flag.
\item premise judgment $\Phi \vdash \Psi_1 \wedge \Psi_2$ $\mapsto$
  \texttt{ih}, with the conjunction shape carried by \texttt{hpost}:
  in the left lemma \texttt{premise.post = .and seq.post other} ---
  the surviving conjunct $\Psi_i$ \emph{is} the
  conclusion\textquotesingle s own post field, and the discarded
  conjunct $\Psi_{3-i}$ $\mapsto$ the binder \texttt{other} (right
  lemma: mirrored, \texttt{.and other seq.post}).
\item $\Phi \mapsto$ \texttt{premise.pre}, shared via \texttt{hpre}
  $:$ \texttt{seq.pre = premise.pre}; conclusion $\Phi \vdash \Psi_i$
  $\mapsto$ the goal \texttt{seq.semanticallyValid H}.
\end{itemize}
No width difference in either direction: the lemma pair covers exactly
the displayed $i$-indexed family, and every hypothesis is consumed.

\ruleblock{Lean verifier receipts}
\sourcenote{module \texttt{BernsteinLogic.Soundness}; bound build \texttt{a20d90e66c13\ldots} (bound from receipt a9c0e34d3e3c9c8c, 2026-08-03T21:30:36Z). A receipt records that named checking layers accepted a build: (module, compiled-artifact SHA-256, environment SHA-256, checker layers), minted when the checking gate ran; the store receipts attesting the bound build are listed with the layers each run credited. Nothing here asks the reader to trust the store: every layer is an independently runnable program over the quoted sources --- rebuild the module and re-run the checkers to reproduce the verdict. Receipts attesting superseded builds credit no layer and are omitted from this rendering (the count of omissions, if any, is noted below).}
\begin{receipttable}
\receiptrow{a9c0e34d3e3c9c8c}{2026-08-03T21:30:36Z}{build+decorrelated}{\badge{BadgeGreen}{current}}
\end{receipttable}
\sourcenote{21 further receipt(s) in the store attest superseded builds of this module and are omitted from this rendering.}
\sourcenote{module \texttt{BernsteinLogic.Core.Catalogue}; bound build \texttt{cbc92e976a1d\ldots} (bound from receipt 5abdb3589707b5df, 2026-08-03T21:30:33Z). A receipt records that named checking layers accepted a build: (module, compiled-artifact SHA-256, environment SHA-256, checker layers), minted when the checking gate ran; the store receipts attesting the bound build are listed with the layers each run credited. Nothing here asks the reader to trust the store: every layer is an independently runnable program over the quoted sources --- rebuild the module and re-run the checkers to reproduce the verdict. Receipts attesting superseded builds credit no layer and are omitted from this rendering (the count of omissions, if any, is noted below).}
\begin{receipttable}
\receiptrow{24804e66edba0bc9}{2026-08-03T20:25:09Z}{build+decorrelated}{\badge{BadgeGreen}{current}}
\receiptrow{5abdb3589707b5df}{2026-08-03T21:30:33Z}{build+decorrelated}{\badge{BadgeGreen}{current}}
\end{receipttable}
\sourcenote{18 further receipt(s) in the store attest superseded builds of this module and are omitted from this rendering.}
\layerline{\layeritem{build}{lake build (elaboration)}{\badge{BadgeGreen}{current}}\layeritem{replay}{leanchecker (official kernel replay)}{\badge{BadgeAmber}{stale only}}\layeritem{hardened}{safe\_verify (hardened screen)}{\badge{BadgeAmber}{stale only}}\layeritem{decorrelated}{lean4export + nanoda (decorrelated kernel)}{\badge{BadgeGreen}{current}}}
\sourcenote{\badge{BadgeBlue}{comparator: exercised via the demonstration fixture}
the composite demonstration certificate (an axiom node, a conjunction
elimination, a Bernstein-cone closure) runs as a fixture through the
three-implementation comparator harness (compiled / kernel-interpreted /
reference-interpreted), so this rule\textquotesingle s node rides every
execution path of that harness.
}
\ruleblock{HOL4 statement}
\sourcenote{theory \texttt{bern\_fam\_ande}, HOL4 development script \texttt{\texttt{bern\_fam\_andeScript.sml}} \badge{BadgeGreen}{script bound}}
\holhead{Definition checkFam\_ande\_def}
\begin{holsrc}
Definition checkFam_ande_def:
  checkFam_ande (cols:columns) (ct:ctx) (i:num) ⇔
    case tagOf cols i of
    | NONE => F
    | SOME (r, v) =>
        r = 4 ∧
        (case childOrd ct i 0 of
         | NONE => F
         | SOME k =>
           case conclRef cols ct i of
           | NONE => F
           | SOME cj =>
           case conclRef cols ct k of
           | NONE => F
           | SOME ccj =>
           case pv cols ct cj of
           | SOME (PJudg prew postw) =>
             (case pv cols ct ccj of
              | SOME (PJudg cprew cpostw) =>
                (case prefField cols ct i 0 of
                 | NONE => F
                 | SOME ot =>
                   case pv cols ct ot of
                   | SOME (PFormula otw) =>
                       (* pre-chain, expanded level (design note N1) *)
                       expandSph FTt prew = expandSph FTt cprew ∧
                       (* child post = f.and; conclusion = taken side,
                          other = untaken side, by variant *)
                       expandSph FTt cpostw =
                         (if v = 0 then
                            FAnd (expandSph FTt postw) (expandSph FTt otw)
                          else
                            FAnd (expandSph FTt otw) (expandSph FTt postw))
                   | _ => F)
              | _ => F)
           | _ => F)
End
\end{holsrc}
\holhead{Theorem famEq\_ande}
\begin{holsrc}
Theorem famEq_ande:
  ∀cols ct i sq nd ws.
    wfColumns cols ct ∧ mkCtx cols = SOME ct ∧
    nodeCorr cols ct i (CertPf sq nd) ∧
    FST (THE (tagOf cols i)) ∈ {4} ⇒
    (checkFam_ande cols ct i ⇔ checkPf1_ande_arm ws sq nd)
\end{holsrc}
\ruleblock{HOL verifier receipts}
\badge{BadgeBlue}{documentary} Documentary citations. Spine census: \texttt{bern\_joins} GREEN under \texttt{hol-strict} (55 theorems, 0 offences, 208-theory ancestry including this theory), \texttt{checkJoins\_eq\_checkPf} verified verbatim in the built signature. The cited \texttt{famEq} lemma is the kernel-checked per-node equivalence: under well-formed columns (\texttt{wfColumns}/\texttt{mkCtx}/\texttt{nodeCorr}) at this rule\textquotesingle s wire tag, the columnar family check coincides with the flat reference walk\textquotesingle s arm for this rule. There is no signed HOL receipt store yet --- these citations are documentary.
Semantic stratum, also ported: the soundness twins
\texttt{sound\_andE\_left} / \texttt{sound\_andE\_right} (theory
\texttt{bern\_tr\_soundstruct}: \texttt{hol-strict} clean, 9 theorems, 0
offences, verification run of 2026-07-22 on an isolated build host,
content-addressed at the same sources), consumed by the certified-
fragment keystone (\texttt{checkPf\_sound\_frag}, theory
\texttt{bern\_tr\_splice}, 20 theorems, 0 offences, same run): a
reference-checker accept inside the proven wire fragment yields sequent
validity.

\sourcenote{HOL receipts are documentary: quoted results of \texttt{hol-strict} runs
and census rows from the HOL4 development\textquotesingle s build
evidence, not minted receipt artefacts. A HOL receipt store (the Lean
store\textquotesingle s twin) is named future work.
}
\ruleblock{Proof}

\sourcenote{Genre: core rule (tier 1). Setting as in the \rulename{Ax}
section; counterexamples read over the one-wire hypernet.}

\textbf{Claim (\texttt{sound\_andE\_left} /
\texttt{sound\_andE\_right}).}\enspace The rule is realized as two
constructors, one per projection; the payload of either supplies a
premise sequent with its sub-derivation, the \emph{discarded} conjunct
(\texttt{other}), and two tree identities: shared context
(\texttt{preEq}) and, for the left constructor, that the premise's
conclusion is $\mathrm{post} \wedge \mathrm{other}$ --- for the right,
$\mathrm{other} \wedge \mathrm{post}$ (\texttt{postEq}). If the premise
sequent $\Phi \vdash \Psi_1 \wedge \Psi_2$ is semantically valid, so is
$\Phi \vdash \Psi_i$.

\begin{bthm}[soundness of $\wedge$E]\label{thm:ande-sound}
Under the payload identities of either constructor, semantic validity of
the premise sequent implies semantic validity of the conclusion, over
every hypernet.
\end{bthm}
\begin{proof}
Left constructor; the right is its mirror image.

\emph{Step 1 (apply the premise).} Let $\sigma$ be admissible with
$\sigma \models \Phi$; by \texttt{preEq} and premise validity,
$\sigma$ satisfies the premise's conclusion.

\emph{Step 2 (read it as a pair and project).} \texttt{postEq} says that
conclusion \emph{is} the tree
$\mathrm{seq.post} \wedge \mathrm{other}$; the $\wedge$-clause of
satisfaction is meta-conjunction, so $\sigma$ satisfies both conjuncts.
Keep the left one: $\sigma \models \mathrm{seq.post}$.
\end{proof}

\begin{bremark}[the discarded conjunct is named, so nothing is searched]
\label{rem:ande-other}
Carrying \texttt{other} in the payload turns the side condition into one
structural identity: the checker never hunts through the premise's
conclusion for an occurrence of the target --- it compares two named
trees. Orientation is checked by the same identity: a left-constructor
certificate claiming the \emph{right} conjunct fails \texttt{postEq}
unless the two conjuncts happen to be the same tree, in which case the
distinction is empty anyway.
\end{bremark}

\begin{bremark}[the index is a constructor, not a flag]
The paper's $i \in \{1, 2\}$ is realized as the choice between two
constructors rather than an integer field --- there is no unchecked
side flag to corrupt, and each constructor's soundness lemma is its own
two-line proof. Dropping \texttt{postEq} in either frees the conclusion
outright; dropping \texttt{preEq} admits the context-swap family (the
\rulename{Weaken} section).
\end{bremark}

\begin{bremark}[one-way information loss, two-way fragment]
Projection discards its \texttt{other} irrecoverably --- the conclusion
carries no trace of it (compare \rulename{Cut}, whose conclusion hides
the cut formula; here the forgetting is the visible point of the rule).
Together with $\wedge$I the fragment is lossless at the formula-tree
level: see the meet-adjunction remark in the $\wedge$I section. In the lattice reading these are
\texttt{M3.andE\_l\_kit}/\texttt{M3.andE\_r\_kit}, the upper direction
of meet's universal property.
\end{bremark}

\begin{anchors}
Steps 1--2 (left) & \texttt{sound\_andE\_left}; payload \texttt{PAndEData.left} (fields \texttt{other}, \texttt{preEq}, \texttt{postEq})\\
mirror (right) & \texttt{sound\_andE\_right}; payload \texttt{PAndEData.right}\\
dispatch & wired as \texttt{RuleData.andE} into \texttt{ProofObject.sound}\\
lattice reading & \texttt{M3.andE\_l\_kit}, \texttt{M3.andE\_r\_kit}\\
\end{anchors}


\rulesection{6}{$\vee$I}{core rule}
\label{blcatrule-OrI}
\ruleblock{The rule (Steele metanotation)}
\sourcenote{as realized in the mechanized development; source row: \texttt{\texttt{rules.toml}} --- \emph{$\vee$I}. The metanotation statement is typeset by hand for notational quality; its mechanized ground truth is the Lean statement below, extracted from the sources at every build of this book and hash-bound. The rule \emph{roster} is machine-checked (every mechanized rule tag round-trips to exactly one catalogue row); the statement \emph{text} is editorial, checked by eye against the extracted Lean beside it.}
\begin{mathpar}
\inferrule*[Right=$\vee$I]{\Phi \vdash \Psi_i}{\Phi \vdash \Psi_1 \vee \Psi_2}
\end{mathpar}

\par\medskip
$\vee$I injects a proven conclusion into a disjunction: from
$\Phi \vdash \Psi_i$, conclude $\Phi \vdash \Psi_1 \vee \Psi_2$ ---
conclusion-side weakening, the mirror of \rulename{Weaken}\textquotesingle s
premise side. A producer mostly reaches for it to make the branches of
a disjunction elimination agree: each branch proves its own
branch-specific fact and $\vee$I-injects it into the common disjunction
the elimination\textquotesingle s conclusion-matching check demands.
The separation instance discussed at $\vee$E is the miniature --- both
of its branch derivations are an axiom node followed by one $\vee$I
node.

Like $\wedge$E, the rule is realized as two constructors, one per
side; the injected side is the constructor choice. The checker
recomputes two tree identities: shared context, and that the
conclusion\textquotesingle s post is one new $\vee$-node with the
sub-derivation\textquotesingle s conclusion on the constructor\textquotesingle s
side and the \emph{named} added disjunct on the other.

Nothing whatsoever is demanded of the added disjunct --- it may be
unsatisfiable, ill-typed, or mention untyped wires --- because
satisfaction of a disjunction demands only the witnessed side, and the
satisfaction relation carries no typing side conditions of its own;
typing enters the calculus only through rules that must evaluate.
Dropping the conclusion identity frees the conclusion outright;
dropping the context identity admits the context-swap family walked at
\rulename{Weaken}.

\ruleblock{Lean statement}
\declhead{BernsteinLogic.Soundness}{sound\_orI\_left}{\badge{BadgeGreen}{statement bound}}
\begin{leansrc}
theorem sound_orI_left {H : Hypernet} {seq premise : Sequent} {other : Formula}
    (ih : premise.semanticallyValid H)
    (hpre : seq.pre = premise.pre)
    (hpost : seq.post = .or premise.post other) :
    seq.semanticallyValid H
\end{leansrc}
\declhead{BernsteinLogic.Soundness}{sound\_orI\_right}{\badge{BadgeGreen}{statement bound}}
\begin{leansrc}
theorem sound_orI_right {H : Hypernet} {seq premise : Sequent} {other : Formula}
    (ih : premise.semanticallyValid H)
    (hpre : seq.pre = premise.pre)
    (hpost : seq.post = .or other premise.post) :
    seq.semanticallyValid H
\end{leansrc}
\par\medskip
Payload, read as in the \rulename{Ax} section: one constructor per
injection, each a sub-derivation plus structural equalities with the
added disjunct as explicit data.

\declhead{BernsteinLogic.Core.Catalogue}{POrIData}{\badge{BadgeGreen}{statement bound}}
\begin{leansrc}
inductive POrIData : Hypernet → Sequent → Type where
  | left  (H : Hypernet) (seq : Sequent)
          (premise : Sequent) (premiseProof : ProofObject H premise)
          (other : Formula)
          (preEq  : seq.pre = premise.pre)
          (postEq : seq.post = .or premise.post other)
          : POrIData H seq
  | right (H : Hypernet) (seq : Sequent)
          (premise : Sequent) (premiseProof : ProofObject H premise)
          (other : Formula)
          (preEq  : seq.pre = premise.pre)
          (postEq : seq.post = .or other premise.post)
          : POrIData H seq
\end{leansrc}
\ruleblock{Correspondence (Steele display $\leftrightarrow$ Lean)}
\badge{BadgeGreen}{reviewer-certified against the displayed texts}
\sourcenote{The dictionary below is a reviewer-certified statement of how the mechanized statements above interpret the displayed rule: display metavariable $\mapsto$ Lean binder, each display hypothesis $\mapsto$ the premise or side condition carrying it, the display conclusion $\mapsto$ the goal --- with any deliberate difference in width between display and mechanization named outright. Certification is recorded as a hash triple over (the displayed cell, the extracted statement texts, this dictionary\textquotesingle s own text); if any of the three changes, this block turns amber until a reviewer re-certifies. The correspondence claim itself remains editorial --- reviewer-grade, not machine-checked; its machine-checked successor types each soundness statement against the rule as a structured object, and supersedes this block per rule as it lands.}
Judgment frame as in \rulename{Ax}; sub-derivation premise as at
\rulename{Weaken}.
\begin{itemize}
\item the display\textquotesingle s index $i \in \{1, 2\}$ $\mapsto$
  the choice of lemma, exactly as at $\wedge$E:
  \texttt{sound\_orI\_left} is $i = 1$, \texttt{sound\_orI\_right} is
  $i = 2$, mirroring the payload\textquotesingle s two constructors.
\item premise judgment $\Phi \vdash \Psi_i$ $\mapsto$ \texttt{ih};
  the injected disjunct $\Psi_i$ \emph{is} \texttt{premise.post}, and
  the added disjunct $\Psi_{3-i}$ $\mapsto$ the binder \texttt{other},
  of which nothing is demanded.
\item conclusion $\Phi \vdash \Psi_1 \vee \Psi_2$ $\mapsto$ the goal
  under \texttt{hpost}: in the left lemma \texttt{seq.post = .or
  premise.post other} (right lemma mirrored); $\Phi \mapsto$
  \texttt{premise.pre}, shared via \texttt{hpre}.
\end{itemize}
No width difference in either direction: the lemma pair covers exactly
the displayed $i$-indexed family, and every hypothesis is consumed.

\ruleblock{Lean verifier receipts}
\sourcenote{module \texttt{BernsteinLogic.Soundness}; bound build \texttt{a20d90e66c13\ldots} (bound from receipt a9c0e34d3e3c9c8c, 2026-08-03T21:30:36Z). A receipt records that named checking layers accepted a build: (module, compiled-artifact SHA-256, environment SHA-256, checker layers), minted when the checking gate ran; the store receipts attesting the bound build are listed with the layers each run credited. Nothing here asks the reader to trust the store: every layer is an independently runnable program over the quoted sources --- rebuild the module and re-run the checkers to reproduce the verdict. Receipts attesting superseded builds credit no layer and are omitted from this rendering (the count of omissions, if any, is noted below).}
\begin{receipttable}
\receiptrow{a9c0e34d3e3c9c8c}{2026-08-03T21:30:36Z}{build+decorrelated}{\badge{BadgeGreen}{current}}
\end{receipttable}
\sourcenote{21 further receipt(s) in the store attest superseded builds of this module and are omitted from this rendering.}
\sourcenote{module \texttt{BernsteinLogic.Core.Catalogue}; bound build \texttt{cbc92e976a1d\ldots} (bound from receipt 5abdb3589707b5df, 2026-08-03T21:30:33Z). A receipt records that named checking layers accepted a build: (module, compiled-artifact SHA-256, environment SHA-256, checker layers), minted when the checking gate ran; the store receipts attesting the bound build are listed with the layers each run credited. Nothing here asks the reader to trust the store: every layer is an independently runnable program over the quoted sources --- rebuild the module and re-run the checkers to reproduce the verdict. Receipts attesting superseded builds credit no layer and are omitted from this rendering (the count of omissions, if any, is noted below).}
\begin{receipttable}
\receiptrow{24804e66edba0bc9}{2026-08-03T20:25:09Z}{build+decorrelated}{\badge{BadgeGreen}{current}}
\receiptrow{5abdb3589707b5df}{2026-08-03T21:30:33Z}{build+decorrelated}{\badge{BadgeGreen}{current}}
\end{receipttable}
\sourcenote{18 further receipt(s) in the store attest superseded builds of this module and are omitted from this rendering.}
\layerline{\layeritem{build}{lake build (elaboration)}{\badge{BadgeGreen}{current}}\layeritem{replay}{leanchecker (official kernel replay)}{\badge{BadgeAmber}{stale only}}\layeritem{hardened}{safe\_verify (hardened screen)}{\badge{BadgeAmber}{stale only}}\layeritem{decorrelated}{lean4export + nanoda (decorrelated kernel)}{\badge{BadgeGreen}{current}}}
\sourcenote{\badge{BadgeGray}{comparator: none} no fixture in the three-implementation comparator harness (compiled / kernel-interpreted / reference-interpreted) exercises this rule.}
\ruleblock{HOL4 statement}
\sourcenote{theory \texttt{bern\_fam\_resid}, HOL4 development script \texttt{\texttt{bern\_fam\_residScript.sml}} \badge{BadgeGreen}{script bound}}
\holhead{Definition checkFam\_ori\_def}
\begin{holsrc}
Definition checkFam_ori_def:
  checkFam_ori (cols:columns) (ct:ctx) (i:num) ⇔
    case tagOf cols i of
    | NONE => F
    | SOME (r, v) =>
      case frSeq cols ct i of
      | NONE => F
      | SOME sq =>
        case frChildSeq cols ct i 0 of
        | NONE => F
        | SOME ps =>
          case frFormula cols ct i 0 of
          | NONE => F
          | SOME other =>
              sq.pre = ps.pre ∧
              sq.post = (if v = 0 then FOr ps.post other
                         else FOr other ps.post)
End
\end{holsrc}
\holhead{Theorem famEq\_ori}
\begin{holsrc}
Theorem famEq_ori:
  wfColumns cols ct ∧ mkCtx cols = SOME ct ∧
  parseCert (cols.hdrImage ++ rest) = SOME (ws, pf) ∧
  nodeCorr cols ct i (CertPf sq nd) ∧
  FST (THE (tagOf cols i)) = 5 ⇒
  (checkFam_ori cols ct i ⇔ armOrI ws sq nd)
\end{holsrc}
\ruleblock{HOL verifier receipts}
\badge{BadgeBlue}{documentary} Documentary citations. Family validation: \texttt{bern\_fam\_resid} built GREEN with all ten \texttt{famEq} arms validated, \texttt{hol-strict} clean; spine census: \texttt{bern\_joins} GREEN under \texttt{hol-strict} (55 theorems, 0 offences, 208-theory ancestry including this theory), \texttt{checkJoins\_eq\_checkPf} verified verbatim in the built signature. The cited \texttt{famEq} lemma is the kernel-checked per-node equivalence: under well-formed columns (\texttt{wfColumns}/\texttt{mkCtx}/\texttt{nodeCorr}) at this rule\textquotesingle s wire tag, the columnar family check coincides with the flat reference walk\textquotesingle s arm for this rule. There is no signed HOL receipt store yet --- these citations are documentary.
Semantic stratum, also ported: the soundness twins
\texttt{sound\_orI\_left} / \texttt{sound\_orI\_right} (theory
\texttt{bern\_tr\_soundstruct}: \texttt{hol-strict} clean, 9 theorems, 0
offences, verification run of 2026-07-22 on an isolated build host,
content-addressed at the same sources), consumed by the certified-
fragment keystone (\texttt{checkPf\_sound\_frag}, theory
\texttt{bern\_tr\_splice}, 20 theorems, 0 offences, same run): a
reference-checker accept inside the proven wire fragment yields sequent
validity.

\sourcenote{HOL receipts are documentary: quoted results of \texttt{hol-strict} runs
and census rows from the HOL4 development\textquotesingle s build
evidence, not minted receipt artefacts. A HOL receipt store (the Lean
store\textquotesingle s twin) is named future work.
}
\ruleblock{Proof}

\sourcenote{Genre: core rule (tier 1). Setting as in the \rulename{Ax}
section; counterexamples read over the one-wire hypernet.}

\textbf{Claim (\texttt{sound\_orI\_left} /
\texttt{sound\_orI\_right}).}\enspace Two constructors, one per
injection; the payload of either supplies a premise sequent
$\Phi \vdash \Psi_i$ with its sub-derivation, the \emph{added} disjunct
(\texttt{other}), and two tree identities: shared context
(\texttt{preEq}) and, for the left constructor, that the conclusion's
post is $\Psi_i \vee \mathrm{other}$ --- for the right,
$\mathrm{other} \vee \Psi_i$ (\texttt{postEq}). If the premise sequent
is semantically valid, so is $\Phi \vdash \Psi_1 \vee \Psi_2$.

\begin{bthm}[soundness of $\vee$I]\label{thm:ori-sound}
Under the payload identities of either constructor, semantic validity of
the premise sequent implies semantic validity of the conclusion, over
every hypernet.
\end{bthm}
\begin{proof}
Left constructor; the right is its mirror image.

\emph{Step 1 (apply the premise).} Let $\sigma$ be admissible with
$\sigma \models \Phi$; by \texttt{preEq} and premise validity,
$\sigma \models \Psi_i$.

\emph{Step 2 (inject).} The $\vee$-clause of satisfaction is
meta-disjunction, and a disjunction is satisfied by satisfying either
side; the witnessed side is $\Psi_i$, and \texttt{postEq} pins the
conclusion's post as the tree $\Psi_i \vee \mathrm{other}$.
\end{proof}

\begin{bremark}[the added disjunct is wholly unconstrained]
\label{rem:ori-other}
\texttt{other} may be unsatisfiable, ill-typed, or mention untyped
wires; nothing is demanded of it, because satisfaction of a disjunction
demands only the witnessed side and the satisfaction relation carries no
typing side conditions of its own --- typing enters the calculus only
through rules that must evaluate. $\vee$I is conclusion-side weakening,
the mirror of \rulename{Weaken}'s premise side: both discard strength,
and discarding is always semantically safe. Dropping \texttt{postEq}
frees the conclusion outright; dropping \texttt{preEq} admits the
context-swap family (the \rulename{Weaken} section).
\end{bremark}

\begin{bremark}[where a producer actually reaches for it]
Mostly to make two branches agree: the branches of a disjunction
elimination must conclude a \emph{common} formula, so each branch
$\vee$I-injects its own branch-specific conclusion into the shared
disjunction. The separation instance quoted in the $\vee$E section is
the miniature: both of its branch derivations are an \rulename{Ax} node
followed by one $\vee$I node (left in one branch, right in the other).
\end{bremark}

\begin{bremark}[lattice reading]
$x \leq y$ implies $x \leq y \sqcup z$: the upper-bound half of join's
universal property (\texttt{M3.orI\_l\_kit}/\texttt{M3.orI\_r\_kit});
the other half is exactly $\vee$E. The four $\wedge$/$\vee$ rules are
the meet and join adjunctions in full, which is the observation the
$\vee$E section's countermodel turns into a theorem about what the core
\emph{cannot} derive.
\end{bremark}

\begin{anchors}
Steps 1--2 (left) & \texttt{sound\_orI\_left}; payload \texttt{POrIData.left} (fields \texttt{other}, \texttt{preEq}, \texttt{postEq})\\
mirror (right) & \texttt{sound\_orI\_right}; payload \texttt{POrIData.right}\\
dispatch & wired as \texttt{RuleData.orI} into \texttt{ProofObject.sound}\\
lattice reading & \texttt{M3.orI\_l\_kit}, \texttt{M3.orI\_r\_kit}\\
\end{anchors}


\rulesection{7}{$\vee$E}{core rule}
\label{blcatrule-OrE}
\ruleblock{The rule (Steele metanotation)}
\sourcenote{as realized in the mechanized development; source row: \texttt{\texttt{rules.toml}} --- \emph{$\vee$E}. The metanotation statement is typeset by hand for notational quality; its mechanized ground truth is the Lean statement below, extracted from the sources at every build of this book and hash-bound. The rule \emph{roster} is machine-checked (every mechanized rule tag round-trips to exactly one catalogue row); the statement \emph{text} is editorial, checked by eye against the extracted Lean beside it.}
\begin{mathpar}
\inferrule*[Right=$\vee$E]{\Phi_1 \vdash \Psi \\ \Phi_2 \vdash \Psi}{\Phi_1 \vee \Phi_2 \vdash \Psi}
\end{mathpar}

\par\medskip
$\vee$E is the calculus\textquotesingle s only branching rule --- case
analysis: from $\Phi_1 \vdash \Psi$ and $\Phi_2 \vdash \Psi$, conclude
$\Phi_1 \vee \Phi_2 \vdash \Psi$. The disjunctions worth eliminating
come from the hypernet\textquotesingle s synthesized primitive facts
--- a cell indicator\textquotesingle s finite case list, a
rectifier\textquotesingle s sign split --- and a producer eliminates
them by proving the goal once per case over the correspondingly
strengthened premise.

The checker recomputes three tree identities: the two branch
sub-derivations conclude the same formula, the
conclusion\textquotesingle s premise is exactly the $\vee$-node of the
two branch premises, and the conclusion\textquotesingle s post is the
common branch conclusion. The conclusion-matching identity is the
load-bearing one --- drop it and a fact proven only in one branch
escapes the case split (the proof section\textquotesingle s
counterexample); the premise identity\textquotesingle s removal admits
the context-swap family.

The essential caveat of this section, though, is the
rule\textquotesingle s \emph{reach}: the quoted form eliminates a
disjunction only when it is the \emph{whole} premise, while in practice
every eliminable disjunction arrives as a conjunct inside the premise
--- and the core provably cannot restructure its way there
(distribution of $\wedge$ over $\vee$ is underivable; the proof section
exhibits the mechanized lattice countermodel). The context-carrying
generalization --- same rule, plus a conjunct path addressing the
disjunction inside the premise\textquotesingle s $\wedge$-tree --- is
proven sound, proven to be a proper extension, and proven to collapse
onto the quoted form at the identity path; the accompanying development
treats it as the intended elimination rule, with the quoted plain form
destined to become its boundary alias.

\ruleblock{Lean statement}
\declhead{BernsteinLogic.Soundness}{sound\_orE}{\badge{BadgeGreen}{statement bound}}
\begin{leansrc}
theorem sound_orE {H : Hypernet} {seq leftPremise rightPremise : Sequent}
    (ihl : leftPremise.semanticallyValid H)
    (ihr : rightPremise.semanticallyValid H)
    (hpostMatch : leftPremise.post = rightPremise.post)
    (hpre : seq.pre = .or leftPremise.pre rightPremise.pre)
    (hpost : seq.post = leftPremise.post) :
    seq.semanticallyValid H
\end{leansrc}
\par\medskip
Payload of the plain form, read as in the \rulename{Ax} section: two
branch sub-derivations plus structural equalities.

\declhead{BernsteinLogic.Core.Catalogue}{POrEData}{\badge{BadgeGreen}{statement bound}}
\begin{leansrc}
inductive POrEData : Hypernet → Sequent → Type where
  | mk (H : Hypernet) (seq : Sequent)
       (leftPremise  : Sequent) (leftProof  : ProofObject H leftPremise)
       (rightPremise : Sequent) (rightProof : ProofObject H rightPremise)
       (postMatch : leftPremise.post = rightPremise.post)
       (preEq     : seq.pre  = .or leftPremise.pre rightPremise.pre)
       (postEq    : seq.post = leftPremise.post)
       : POrEData H seq
\end{leansrc}
\par\medskip
\textbf{The context-carrying satellite.}\enspace The four declarations
below locate the context-carrying generalization precisely: its
applicability data (a conjunct path with read/replace obligations; the
plain rule is the $\pi = \mathrm{here}$ case), its semantic soundness
over every hypernet, the concrete distribution instance no core
derivation concludes, and the resulting impossibility of elaborating
the generalization into the plain core. The proof section walks the
countermodel argument these rest on.

\declhead{BernsteinLogic.Admissible}{OrECtxStep}{\badge{BadgeGreen}{statement bound}}
\begin{leansrc}
/-- Applicability data of context-carrying disjunction elimination: a
conjunct path into the conclusion's premise addresses a disjunction;
each branch premise is the conclusion's premise with that conjunct
replaced by the respective disjunct; both branches conclude the common
conclusion. Field-for-field, these are the reference checker's
obligations for the rule (path resolves, path addresses a disjunction,
the two premise-replacement equations, the two conclusion equations).
The primitive elimination rule is the `path = here` case. -/
structure OrECtxStep where
  seq          : Sequent
  leftPremise  : Sequent
  rightPremise : Sequent
  path : ConjPath
  d₁ : Formula
  d₂ : Formula
  pathReads : seq.pre.conjGet? path = some (.or d₁ d₂)
  leftPre   : seq.pre.conjReplace? path d₁ = some leftPremise.pre
  rightPre  : seq.pre.conjReplace? path d₂ = some rightPremise.pre
  leftPost  : leftPremise.post = seq.post
  rightPost : rightPremise.post = seq.post
\end{leansrc}
\declhead{BernsteinLogic.Admissible}{sound\_orECtx}{\badge{BadgeGreen}{statement bound}}
\begin{leansrc}
/-- Context-carrying disjunction elimination is sound: whenever both
branch premises are semantically valid, so is the conclusion — over
every hypernet. This is the warrant for adding the rule as an
admissible extension: every sequent the extended calculus derives is
semantically valid (per-rule soundness of the core + this lemma).

The proof is the conjunct-path read lemma (extract the disjunction),
the case split of the primitive rule, and the conjunct-path replace
lemma (transport the branch's strengthened premise). -/
theorem sound_orECtx {H : Hypernet} (st : OrECtxStep)
    (ihl : st.leftPremise.semanticallyValid H)
    (ihr : st.rightPremise.semanticallyValid H) :
    st.seq.semanticallyValid H
\end{leansrc}
\declhead{BernsteinLogic.Admissible}{cx\_not\_derivable}{\badge{BadgeGreen}{statement bound}}
\begin{leansrc}
/-- The instance's conclusion is NOT derivable in the catalogue: its
interpretation drops from an atom to ⊥ across the turnstile, and the
metatheorem says no derivation does that. -/
theorem cx_not_derivable : IsEmpty (ProofObject cxH cxSeq)
\end{leansrc}
\declhead{BernsteinLogic.Admissible}{orECtx\_not\_elaborable}{\badge{BadgeGreen}{statement bound}}
\begin{leansrc}
/-- Context-carrying disjunction elimination does not elaborate into
the core calculus: no elaboration function can exist, because this
instance has derivable branch premises and an underivable conclusion.
(Contrast `OrECtxStep.elabHere`: at the boundary path the elaboration
exists; the failure is specific to genuinely deep paths.) -/
theorem orECtx_not_elaborable :
    ¬ (∀ (st : OrECtxStep) (H : Hypernet),
        Nonempty (ProofObject H st.leftPremise) →
        Nonempty (ProofObject H st.rightPremise) →
        Nonempty (ProofObject H st.seq))
\end{leansrc}
\declhead{BernsteinLogic.Admissible}{ProofObject.m3sound}{\badge{BadgeGreen}{statement bound}}
\begin{leansrc}
/-- Monotone replacement at a conjunct path: substituting a
lattice-smaller formula at any conjunct address shrinks the M3 value.
Meet monotonicity only — no distributivity — which is why Promise,
unlike ∨E-ctx, is lattice-sound and core-eligible. -/
theorem mval_conjReplace_mono :
    ∀ {π : ConjPath} {Γ Γ' x q : Formula},
    Γ.conjGet? π = some q → Γ.conjReplace? π x = some Γ' →
    M3.le (mval x) (mval q) = true →
    M3.le (mval Γ') (mval Γ) = true := by
  intro π
  induction π with
  | here =>
    intro Γ Γ' x q hget hrep hle
    simp only [Formula.conjGet?] at hget
    simp only [Formula.conjReplace?] at hrep
    injection hget with h1
    injection hrep with h2
    rw [← h2, h1]
    exact hle
  | left p ih =>
    intro Γ Γ' x q hget hrep hle
    cases Γ with
    | and l r =>
      simp only [Formula.conjGet?] at hget
      simp only [Formula.conjReplace?] at hrep
      cases hmap : l.conjReplace? p x with
      | none => rw [hmap] at hrep; nomatch hrep
      | some l' =>
        rw [hmap] at hrep
        injection hrep with h2
        rw [← h2]
        show M3.le ((mval l').inf (mval r)) ((mval l).inf (mval r)) = true
        exact M3.inf_mono_l _ _ _ (ih hget hmap hle)
    | atom b tl tr => nomatch hrep
    | inRegion t Δ S => nomatch hrep
    | or l r => nomatch hrep
    | tt => nomatch hrep
    | ff => nomatch hrep
  | right p ih =>
    intro Γ Γ' x q hget hrep hle
    cases Γ with
    | and l r =>
      simp only [Formula.conjGet?] at hget
      simp only [Formula.conjReplace?] at hrep
      cases hmap : r.conjReplace? p x with
      | none => rw [hmap] at hrep; nomatch hrep
      | some r' =>
        rw [hmap] at hrep
        injection hrep with h2
        rw [← h2]
        show M3.le ((mval l).inf (mval r')) ((mval l).inf (mval r)) = true
        exact M3.inf_mono_r _ _ _ (ih hget hmap hle)
    | atom b tl tr => nomatch hrep
    | inRegion t Δ S => nomatch hrep
    | or l r => nomatch hrep
    | tt => nomatch hrep
    | ff => nomatch hrep

/-! ## §A9 The metatheorem: derivability is lattice-sound

Over a hypernet with no typed wires, every derivable sequent satisfies
⟦pre⟧ ≤ ⟦post⟧ in M₃. One case per rule of the catalogue; the
propositional rules are the lattice axioms, equality elimination is the
substitution invariance (its typing field forces both terms closed over
an untyped hypernet), the region/arithmetic-closure rules are handled
by their conclusion shapes or payload absurdities, and the
empty-payload rules are vacuous. -/
theorem ProofObject.m3sound : ∀ {H : Hypernet} {seq : Sequent},
    (∀ w, H.wireType w = none) → ProofObject H seq →
    M3.le (mval seq.pre) (mval seq.post) = true
  | _, _, _, .mk _ _ _ (.ax _ _ (.mk _ _ hpp)) => by
      rw [hpp]; exact M3.le_refl _
  | _, _, hE, .mk _ _ _ (.weaken _ _ (.mk _ _ _ pproof _ hpre hpost)) => by
      rw [hpre, hpost]
      exact M3.weaken_kit _ _ _ (pproof.m3sound hE)
  | _, _, hE, .mk _ _ _ (.cut _ _
      (.mk _ _ _ _ lproof _ rproof hpre hpost hcl hcr)) => by
      have hl
\end{leansrc}
\declhead{BernsteinLogic.Admissible}{cxStep}{\badge{BadgeGreen}{statement bound}}
\begin{leansrc}
/-- The separation instance's applicability data: all five obligations
hold by computation. -/
def cxStep : OrECtxStep
\end{leansrc}
\ruleblock{Correspondence (Steele display $\leftrightarrow$ Lean)}
\badge{BadgeGreen}{reviewer-certified against the displayed texts}
\sourcenote{The dictionary below is a reviewer-certified statement of how the mechanized statements above interpret the displayed rule: display metavariable $\mapsto$ Lean binder, each display hypothesis $\mapsto$ the premise or side condition carrying it, the display conclusion $\mapsto$ the goal --- with any deliberate difference in width between display and mechanization named outright. Certification is recorded as a hash triple over (the displayed cell, the extracted statement texts, this dictionary\textquotesingle s own text); if any of the three changes, this block turns amber until a reviewer re-certifies. The correspondence claim itself remains editorial --- reviewer-grade, not machine-checked; its machine-checked successor types each soundness statement against the rule as a structured object, and supersedes this block per rule as it lands.}
Judgment frame as in \rulename{Ax}; sub-derivation premises as at
\rulename{Weaken}. The displayed rule corresponds to
\texttt{sound\_orE}:
\begin{itemize}
\item the branch judgments $\Phi_1 \vdash \Psi$ and $\Phi_2 \vdash
  \Psi$ $\mapsto$ \texttt{ihl}/\texttt{ihr}; the shared conclusion
  $\Psi$ is shared by equation, \texttt{hpostMatch} $:$
  \texttt{leftPremise.post = rightPremise.post} --- the load-bearing
  identity the commentary singles out.
\item $\Phi_1 \mapsto$ \texttt{leftPremise.pre}; $\Phi_2 \mapsto$
  \texttt{rightPremise.pre}; the conclusion\textquotesingle s premise
  $\Phi_1 \vee \Phi_2$ $\mapsto$ \texttt{hpre} $:$ \texttt{seq.pre =
  .or leftPremise.pre rightPremise.pre} --- the disjunction must be
  the \emph{whole} premise tree, which is precisely the reach
  limitation the commentary discusses.
\item conclusion post $\mapsto$ \texttt{hpost} $:$ \texttt{seq.post =
  leftPremise.post}; the judgment $\mapsto$ the goal
  \texttt{seq.semanticallyValid H}.
\end{itemize}
Width note, deliberate and one-sided: the certified statement set also
carries the context-carrying satellite (\texttt{OrECtxStep},
\texttt{sound\_orECtx}, \texttt{cx\_not\_derivable},
\texttt{orECtx\_not\_elaborable}). Those statements are \emph{about
the generalization}, not the displayed rule: \texttt{sound\_orECtx} is
strictly wider (a conjunct path addresses a disjunction \emph{inside}
the premise tree; the displayed rule is its identity-path case), and
the two negative statements certify that the extension is proper ---
its instance is underivable in, and not elaborable into, the plain
core. The displayed rule itself corresponds to \texttt{sound\_orE}
with no width gap in either direction.

\ruleblock{Lean verifier receipts}
\sourcenote{module \texttt{BernsteinLogic.Soundness}; bound build \texttt{a20d90e66c13\ldots} (bound from receipt a9c0e34d3e3c9c8c, 2026-08-03T21:30:36Z). A receipt records that named checking layers accepted a build: (module, compiled-artifact SHA-256, environment SHA-256, checker layers), minted when the checking gate ran; the store receipts attesting the bound build are listed with the layers each run credited. Nothing here asks the reader to trust the store: every layer is an independently runnable program over the quoted sources --- rebuild the module and re-run the checkers to reproduce the verdict. Receipts attesting superseded builds credit no layer and are omitted from this rendering (the count of omissions, if any, is noted below).}
\begin{receipttable}
\receiptrow{a9c0e34d3e3c9c8c}{2026-08-03T21:30:36Z}{build+decorrelated}{\badge{BadgeGreen}{current}}
\end{receipttable}
\sourcenote{21 further receipt(s) in the store attest superseded builds of this module and are omitted from this rendering.}
\sourcenote{module \texttt{BernsteinLogic.Core.Catalogue}; bound build \texttt{cbc92e976a1d\ldots} (bound from receipt 5abdb3589707b5df, 2026-08-03T21:30:33Z). A receipt records that named checking layers accepted a build: (module, compiled-artifact SHA-256, environment SHA-256, checker layers), minted when the checking gate ran; the store receipts attesting the bound build are listed with the layers each run credited. Nothing here asks the reader to trust the store: every layer is an independently runnable program over the quoted sources --- rebuild the module and re-run the checkers to reproduce the verdict. Receipts attesting superseded builds credit no layer and are omitted from this rendering (the count of omissions, if any, is noted below).}
\begin{receipttable}
\receiptrow{24804e66edba0bc9}{2026-08-03T20:25:09Z}{build+decorrelated}{\badge{BadgeGreen}{current}}
\receiptrow{5abdb3589707b5df}{2026-08-03T21:30:33Z}{build+decorrelated}{\badge{BadgeGreen}{current}}
\end{receipttable}
\sourcenote{18 further receipt(s) in the store attest superseded builds of this module and are omitted from this rendering.}
\sourcenote{module \texttt{BernsteinLogic.Admissible}; bound build \texttt{52e5099ba522\ldots} (bound from receipt c17abf2e103fe745, 2026-07-21T05:28:22Z). A receipt records that named checking layers accepted a build: (module, compiled-artifact SHA-256, environment SHA-256, checker layers), minted when the checking gate ran; the store receipts attesting the bound build are listed with the layers each run credited. Nothing here asks the reader to trust the store: every layer is an independently runnable program over the quoted sources --- rebuild the module and re-run the checkers to reproduce the verdict. Receipts attesting superseded builds credit no layer and are omitted from this rendering (the count of omissions, if any, is noted below).}
\begin{receipttable}
\receiptrow{7f3043c54f717472}{2026-07-20T07:43:38Z}{build+decorrelated}{\badge{BadgeGreen}{current}}
\receiptrow{acaf0e8871731343}{2026-07-20T20:44:28Z}{build+decorrelated}{\badge{BadgeGreen}{current}}
\receiptrow{62974466ea131d89}{2026-07-20T21:00:57Z}{build+replay+hardened+decorrelated}{\badge{BadgeGreen}{current}}
\receiptrow{c17abf2e103fe745}{2026-07-21T05:28:22Z}{build+replay+hardened+decorrelated}{\badge{BadgeGreen}{current}}
\end{receipttable}
\sourcenote{6 further receipt(s) in the store attest superseded builds of this module and are omitted from this rendering.}
\layerline{\layeritem{build}{lake build (elaboration)}{\badge{BadgeGreen}{current}}\layeritem{replay}{leanchecker (official kernel replay)}{\badge{BadgeGreen}{current}}\layeritem{hardened}{safe\_verify (hardened screen)}{\badge{BadgeGreen}{current}}\layeritem{decorrelated}{lean4export + nanoda (decorrelated kernel)}{\badge{BadgeGreen}{current}}}
\sourcenote{\badge{BadgeGray}{comparator: none} no fixture in the three-implementation comparator harness (compiled / kernel-interpreted / reference-interpreted) exercises this rule.}
\ruleblock{HOL4 statement}
\sourcenote{theory \texttt{bern\_fam\_ore}, HOL4 development script \texttt{\texttt{bern\_fam\_oreScript.sml}} \badge{BadgeGreen}{script bound}}
\holhead{Definition checkFam\_ore\_def}
\begin{holsrc}
Definition checkFam_ore_def:
  checkFam_ore (cols:columns) (ct:ctx) (i:num) ⇔
    case nodeSq cols ct i of
    | NONE => F
    | SOME sq =>
      case childOrd ct i 0 of
      | NONE => F
      | SOME kl =>
        case childOrd ct i 1 of
        | NONE => F
        | SOME kr =>
          case nodeSq cols ct kl of
          | NONE => F
          | SOME ls =>
            case nodeSq cols ct kr of
            | NONE => F
            | SOME rs =>
              case pathArgs cols ct i of
              | NONE => F
              | SOME (path, ifc) =>
                case conjGet sq.pre path of
                | SOME (FOr d1 d2) =>
                    (conjReplace sq.pre path d1 = SOME ls.pre ∧
                     conjReplace sq.pre path d2 = SOME rs.pre ∧
                     ls.post = sq.post ∧ rs.post = sq.post)
                | _ => F
End
\end{holsrc}
\holhead{Theorem famEq\_ore}
\begin{holsrc}
Theorem famEq_ore:
  ∀cols ct i sq nd ws.
    wfColumns cols ct ∧ mkCtx cols = SOME ct ∧
    nodeCorr cols ct i (CertPf sq nd) ∧
    FST (THE (tagOf cols i)) ∈ {6:num} ⇒
    (checkFam_ore cols ct i ⇔ oreArm ws sq nd)
\end{holsrc}
\ruleblock{HOL verifier receipts}
\badge{BadgeBlue}{documentary} Documentary citations. Spine census: \texttt{bern\_joins} GREEN under \texttt{hol-strict} (55 theorems, 0 offences, 208-theory ancestry including this theory), \texttt{checkJoins\_eq\_checkPf} verified verbatim in the built signature. The cited \texttt{famEq} lemma is the kernel-checked per-node equivalence: under well-formed columns (\texttt{wfColumns}/\texttt{mkCtx}/\texttt{nodeCorr}) at this rule\textquotesingle s wire tag, the columnar family check coincides with the flat reference walk\textquotesingle s arm for this rule. There is no signed HOL receipt store yet --- these citations are documentary.
Semantic stratum, also ported: the soundness twins \texttt{sound\_orE}
(theory \texttt{bern\_tr\_soundstruct}: \texttt{hol-strict} clean, 9
theorems, 0 offences, verification run of 2026-07-22 on an isolated
build host, content-addressed at the same sources) and
\texttt{sound\_orECtx} (theory \texttt{bern\_tr\_kitd}, 9 theorems, 0
offences, same run --- the context form the wire\textquotesingle s tag
carries), consumed by the certified-fragment keystone
(\texttt{checkPf\_sound\_frag}, theory \texttt{bern\_tr\_splice}, 20
theorems, 0 offences, same run): a reference-checker accept inside the
proven wire fragment yields sequent validity.

\sourcenote{HOL receipts are documentary: quoted results of \texttt{hol-strict} runs
and census rows from the HOL4 development\textquotesingle s build
evidence, not minted receipt artefacts. A HOL receipt store (the Lean
store\textquotesingle s twin) is named future work.
}
\ruleblock{Proof}

\sourcenote{Genre: core rule (tier 1). Setting as in the \rulename{Ax}
section; counterexamples read over the one-wire hypernet. This section
also carries the fork note on the rule's \emph{reach}: the plain form
quoted here, its mechanized limitation, and the context-carrying
generalization of which it is the boundary case.}

\textbf{Claim (\texttt{sound\_orE}).}\enspace The payload supplies two
branch sequents $\Phi_1 \vdash \Psi$ and $\Phi_2 \vdash \Psi$ with their
sub-derivations and three tree identities: the branches share one
conclusion (\texttt{postMatch}), the conclusion's premise is the
disjunction $\Phi_1 \vee \Phi_2$ of the branch premises (\texttt{preEq}),
and the conclusion's post is the left branch's post (\texttt{postEq}).
If both branch sequents are semantically valid, so is
$\Phi_1 \vee \Phi_2 \vdash \Psi$.

\begin{bthm}[soundness of $\vee$E, plain form]\label{thm:ore-sound}
Under the three payload identities, semantic validity of both branch
sequents implies semantic validity of the conclusion, over every
hypernet.
\end{bthm}
\begin{proof}
Let $\sigma$ be admissible with $\sigma$ satisfying the conclusion's
premise, which \texttt{preEq} pins as the tree $\Phi_1 \vee \Phi_2$.

\emph{Step 1 (case split).} The $\vee$-clause of satisfaction is
meta-disjunction: $\sigma \models \Phi_1$ or $\sigma \models \Phi_2$.

\emph{Step 2 (either branch delivers).} In the first case, left branch
validity gives $\sigma \models \Psi$ directly (\texttt{postEq}). In the
second, right branch validity gives satisfaction of the right branch's
conclusion, which \texttt{postMatch} identifies --- as a tree --- with
the left's, hence with the conclusion's post.
\end{proof}

\begin{bremark}[\texttt{postMatch} is load-bearing]\label{rem:ore-post}
The branches must agree on what they conclude \emph{before} the split is
trusted. Drop \texttt{postMatch} and take the left branch
$w \leq 0 \vdash w \leq 0$ and the right branch $\top \vdash \top$ (both
\rulename{Ax} instances): the degenerate rule, reading its conclusion
from the left branch, concludes $(w \leq 0) \vee \top \vdash w \leq 0$
--- false at $\sigma(w) = 1$, where the premise holds through the right
disjunct. \texttt{preEq}'s removal admits the context-swap family (the
\rulename{Weaken} section); \texttt{postEq}'s removal frees the
conclusion outright. In the lattice reading the rule is the bounding
half of join's universal property --- $x \leq z$ and $y \leq z$ iff
$x \sqcup y \leq z$ --- whose introducing half is $\vee$I.
\end{bremark}

\begin{bremark}[the reach problem, and why it is a theorem]
\label{rem:ore-reach}
\texttt{preEq} demands that the eliminated disjunction \emph{be} the
whole premise. But in this calculus the disjunctions worth eliminating
almost never are: a sequent's premise in practice has the shape
$\Gamma \wedge \mathsf{SP}$ --- a context conjoined with the synthesized
case facts of the hypernet's primitives (a cell indicator's finite case
list, a rectifier's sign split) --- so the disjunction arrives as a
\emph{conjunct} of the premise, one or more $\wedge$-nodes deep. The
plain rule cannot reach it, and the core cannot restructure its way
there, because distributing $\wedge$ over $\vee$ is underivable there.
That is a theorem, not a design suspicion. Interpret formulas into
$M_3$, the five-element diamond --- three pairwise-incomparable atoms
between $\bot$ and $\top$, the canonical modular non-distributive
lattice \cite{birkhoff1967} --- in full: $\wedge$ to meet, $\vee$ to
join, $\top$ to top, $\bot$ to bottom, three disequality atoms on
designated wires to the three lattice atoms, region atoms to bottom,
and every remaining atom to top. Work over a hypernet all of whose
wires are untyped; that hypothesis is what discharges every
non-structural rule (typing packs force closed terms there, region and
arithmetic payloads are handled by their conclusion shapes or turn
absurd, and the payload-free rules are vacuous). Then a derivable
sequent's premise always interprets $\leq$ its conclusion --- by
induction on derivations, one case per rule, the structural cases
being precisely the lattice axioms of order, meet, and join: the
sections \rulename{Ax} through $\vee$E in this book, one kit lemma
each (\texttt{ProofObject.m3sound}). No rule of the \emph{plain core}
--- the calculus as quoted in these sections, prior to the
context-carrying extension of \autoref{rem:ore-fork}, whose whole
point is otherwise --- consumes distributivity. Now take the
distribution instance
$P \wedge (Q \vee R) \;\vdash\; (P \wedge Q) \vee (P \wedge R)$:
its premise interprets to $a \sqcap (b \sqcup c) = a \sqcap \top = a$,
its conclusion to $(a \sqcap b) \sqcup (a \sqcap c) = \bot \sqcup \bot
= \bot$, and $a \not\leq \bot$ --- so no core derivation concludes it
(\texttt{cx\_not\_derivable}), even though both branch premises
$P \wedge Q \vdash \cdot$ and $P \wedge R \vdash \cdot$ are two-step
core derivations (an \rulename{Ax} node and a $\vee$I node each).
\end{bremark}

\begin{bremark}[the context-carrying form, and the fork discipline]
\label{rem:ore-fork}
The repair is not restructuring but a stronger elimination. The
context-carrying rule's payload adds a \emph{conjunct path} $\pi$ into
the premise's $\wedge$-tree, with read and replace obligations: the
conjunct at $\pi$ is a disjunction $\Phi_1 \vee \Phi_2$, each branch
premise is the conclusion's premise with that conjunct replaced by the
respective disjunct, and both branches conclude the common conclusion
(\texttt{OrECtxStep}). Three mechanized facts locate it precisely.
\emph{Sound}: whenever both branch premises are semantically valid, so
is the conclusion, over every hypernet (\texttt{sound\_orECtx} --- one
path-read lemma to extract the disjunction, the case split of
Theorem~\ref{thm:ore-sound}, one path-replace lemma to transport each
branch's strengthened premise). \emph{Proper}: no elaboration of the
context-carrying rule into the core exists
(\texttt{orECtx\_not\_elaborable}) --- concretely, at the conjunct path
addressing $Q \vee R$ as the premise's right conjunct, the read
obligation returns that disjunction and the replace obligations return
the branch premises $P \wedge Q$ and $P \wedge R$ of
\autoref{rem:ore-reach}'s instance (\texttt{cxStep}); its branches
are derivable there and its conclusion is not, so the extension proves
something true that the core cannot reach. \emph{Conservative at the boundary}: at $\pi = \mathrm{here}$
the rule elaborates into the plain constructor
(\texttt{OrECtxStep.elabHere}) --- the form proven sound above is
exactly the identity-path case. The cost side of the fork is small and
bounded: the certificate carries the path, the checker two small path
functions (read-at, replace-at), and the soundness proof one lemma
about conjunct-path replacement. The alternative --- adding
$\wedge$-restructuring rules to the core (rules whose premise and
conclusion rearrange the same conjuncts: reassociations, permutations)
--- cannot substitute: such rules are sound for every lattice
interpretation, so the induction above absorbs them --- as it absorbs
any rule that preserves $M_3$-order --- and the countermodel stands;
closing the gap requires a rule whose soundness genuinely uses
distributivity, which is the context-carrying elimination itself. The
accompanying development accordingly treats the context-carrying form
as the intended elimination rule; the constructor quoted in this
section's Lean cell is the plain form, destined to become its boundary
alias.
\end{bremark}

\begin{anchors}
Steps 1--2 & \texttt{sound\_orE}; payload \texttt{POrEData} (fields \texttt{postMatch}, \texttt{preEq}, \texttt{postEq}); wired as \texttt{RuleData.orE}\\
$\vee$-clause of satisfaction & \texttt{Formula.holds} (case \texttt{or})\\
lattice reading & \texttt{M3.orE\_kit}\\
Remark \ref{rem:ore-reach}: order-preservation metatheorem & \texttt{ProofObject.m3sound} (BernsteinLogic/Admissible.lean)\\
Remark \ref{rem:ore-reach}: underivable distribution instance & \texttt{cxStep}, \texttt{cx\_not\_derivable}\\
Remark \ref{rem:ore-fork}: extension payload + soundness & \texttt{OrECtxStep}, \texttt{sound\_orECtx}\\
Remark \ref{rem:ore-fork}: properness & \texttt{orECtx\_not\_elaborable}\\
Remark \ref{rem:ore-fork}: boundary elaboration & \texttt{OrECtxStep.elabHere}\\
\end{anchors}


\rulesection{8}{EqElim}{core rule}
\label{blcatrule-EqElim}
\ruleblock{The rule (Steele metanotation)}
\sourcenote{as realized in the mechanized development; source row: \texttt{\texttt{rules.toml}} --- \emph{EqElim}. The metanotation statement is typeset by hand for notational quality; its mechanized ground truth is the Lean statement below, extracted from the sources at every build of this book and hash-bound. The rule \emph{roster} is machine-checked (every mechanized rule tag round-trips to exactly one catalogue row); the statement \emph{text} is editorial, checked by eye against the extracted Lean beside it.}
\begin{mathpar}
\inferrule*[Right=\rulename{EqElim}]{\Phi \vdash t = u \\ \Phi \vdash \Psi\sub{a}{t} \\ \exists s.\; \bshape_H(t) = s = \bshape_H(u)}{\Phi \vdash \Psi\sub{a}{u}}
\end{mathpar}

\par\medskip
\rulename{EqElim} is Leibniz substitution as a proof rule: an established
equality $t = u$ licenses replacing $t$ by $u$ at a designated hole in an
established formula. The payload names the hole (a wire $a$), the body
formula $\Psi$, and the two terms; premises prove $t = u$ and
$\Psi\sub{a}{t}$ over the shared context, and the conclusion is
$\Psi\sub{a}{u}$. A producer reaches for it wherever a derived fact must
be rewritten along an equality: pushing a resolved value through a bound
(the ReLU rules conclude equalities such as $y = h$ exactly so that this
rule can consume them), transporting an atom across a
\rulename{Norm}-derived respelling, or assembling transitivity --- which
is not primitive in this calculus: a one-node derivation through this
rule (the proof section\textquotesingle s closing remark spells it out;
\rulename{EqRefl} and \rulename{EqSym} complete the kit).

What the checker recomputes: both substitution instances --- it applies
the (binder-free, hence capture-free) substitution function to
$(\Psi, a, t)$ and to $(\Psi, a, u)$ and demands syntactic equality with
the cited premise and conclusion formulas --- plus the premise-sharing
equations, and the typing side condition: $t$ and $u$ carry a
\emph{common defined} broadcast shape in the ambient hypernet. Nothing
semantic is consulted; every side condition is a decidable syntactic
equation.

The typing condition is the load-bearing hypothesis. Equality atoms hold
\emph{elementwise under broadcast}, so a scalar can stand equal to a
vector whose every entry it matches --- a genuine, derivable, usable
fact that is nonetheless not a license to interchange the two terms: in
a region atom, which constrains the shape of its term, the substitution
flips satisfiability. Dropping the condition admits a concrete derivation
of an unsatisfiable conclusion from a satisfiable premise (the proof
section constructs it); with the condition, that certificate is refused
because no common shape exists. Dropping either substitution equation
instead would let the conclusion be an arbitrary formula; dropping the
premise-match would combine facts proven under different contexts.

\cataloguenote{Producer-side, equational obligations are typically discharged by
congruence-closure engines \cite{nelsonoppen1980,nieuwenhuisoliveras2007};
the certificate carries not the closure computation but its residue ---
a chain of \rulename{EqRefl}/\rulename{EqSym}/\rulename{EqElim} nodes,
each checked by recomputing one substitution.
}
\ruleblock{Lean statement}
\declhead{BernsteinLogic.Soundness}{Term.evalQ\_subst\_congr}{\badge{BadgeGreen}{statement bound}}
\begin{leansrc}
theorem Term.evalQ_subst_congr {H : Hypernet} (σ : H.Valuation)
    {w : Wire} {u₁ u₂ : Term} (hu : Term.evalQ σ u₁ = Term.evalQ σ u₂) :
    ∀ (t : Term),
      Term.evalQ σ (Term.subst w u₁ t) = Term.evalQ σ (Term.subst w u₂ t)
\end{leansrc}
\declhead{BernsteinLogic.Soundness}{Formula.holds\_subst\_congr}{\badge{BadgeGreen}{statement bound}}
\begin{leansrc}
theorem Formula.holds_subst_congr {H : Hypernet} (σ : H.Valuation)
    {w : Wire} {u₁ u₂ : Term} (hu : Term.evalQ σ u₁ = Term.evalQ σ u₂) :
    ∀ (F : Formula), (Formula.holds σ (F.subst w u₁) ↔ Formula.holds σ (F.subst w u₂))
\end{leansrc}
\declhead{BernsteinLogic.Soundness}{Term.evalQ\_bshape}{\badge{BadgeGreen}{statement bound}}
\begin{leansrc}
/-- bshape soundness: a defined evaluation takes the broadcast shape. -/
theorem Term.evalQ_bshape {H : Hypernet} (σ : H.Valuation) :
    ∀ (t : Term) (s : Shape) (q : QTensor),
    t.bshape H = some s → Term.evalQ σ t = some q → q.shape = s
\end{leansrc}
\declhead{BernsteinLogic.Soundness}{QTensor.eq\_of\_rel2}{\badge{BadgeGreen}{statement bound}}
\begin{leansrc}
/-- rel2 with an equality-forcing relation at equal shapes forces tensor
equality (the PROPOSED-B13 EqElim extraction: without equal shapes the
broadcast branches of rel2 defeat this — the header's counterexample). -/
theorem QTensor.eq_of_rel2 {r : ℚ → ℚ → Prop} {a b : QTensor}
    (hr : ∀ x y, r x y → x = y) (hs : a.shape = b.shape)
    (h : QTensor.rel2 r a b) : a = b
\end{leansrc}
\declhead{BernsteinLogic.Soundness}{sound\_eqElim}{\badge{BadgeGreen}{statement bound}}
\begin{leansrc}
theorem sound_eqElim {H : Hypernet} {seq eqPremise bodyPremise : Sequent}
    {wire : Wire} {termT termU : Term} {bodyForm : Formula}
    (ihEq : eqPremise.semanticallyValid H)
    (ihBody : bodyPremise.semanticallyValid H)
    (hpre : seq.pre = eqPremise.pre)
    (hpreMatch : eqPremise.pre = bodyPremise.pre)
    (heqPost : eqPremise.post = .atom .eq termT termU)
    (htyEq : ∃ s, termT.bshape H = some s ∧ termU.bshape H = some s)
    (hbodySubT : substituteInFormula bodyForm wire termT bodyPremise.post)
    (hconclSubU : substituteInFormula bodyForm wire termU seq.post) :
    seq.semanticallyValid H
\end{leansrc}
\par\medskip
\textbf{The substitution the rule trusts, in full.}\enspace The payload
fields \texttt{bodySubT} and \texttt{conclSubU} are equations against the
substitution function displayed below: plain leaf replacement over a
grammar with no binders, so no capture or freshness side condition exists
anywhere. A skeptical reader should confirm by eye that
\texttt{substituteInFormula} is nothing but that equation --- the rule
manipulates formulas syntactically and consults no valuation; semantic
truth enters only through \texttt{sound\_eqElim}, from the premise
validity hypotheses.

\declhead{BernsteinLogic.Core.Catalogue}{Term.subst}{\badge{BadgeGreen}{definition bound}}
\begin{leansrc}
def Term.subst (w : Wire) (t : Term) : Term → Term
  | .var w' => if w' = w then t else .var w'
  | .const c => .const c
  | .add l r => .add (Term.subst w t l) (Term.subst w t r)
  | .mul l r => .mul (Term.subst w t l) (Term.subst w t r)
  | .sub l r => .sub (Term.subst w t l) (Term.subst w t r)
  | .neg u => .neg (Term.subst w t u)
  | .idx u i => .idx (Term.subst w t u) i
\end{leansrc}
\declhead{BernsteinLogic.Core.Catalogue}{Formula.subst}{\badge{BadgeGreen}{definition bound}}
\begin{leansrc}
def Formula.subst (w : Wire) (t : Term) : Formula → Formula
  | .atom b l r => .atom b (Term.subst w t l) (Term.subst w t r)
  | .inRegion u simplex S => .inRegion (Term.subst w t u) simplex S
  | .and l r => .and (Formula.subst w t l) (Formula.subst w t r)
  | .or l r => .or (Formula.subst w t l) (Formula.subst w t r)
  | .tt => .tt
  | .ff => .ff
\end{leansrc}
\declhead{BernsteinLogic.Core.Catalogue}{substituteInFormula}{\badge{BadgeGreen}{definition bound}}
\begin{leansrc}
def substituteInFormula (F : Formula) (w : Wire) (t : Term) (F' : Formula) : Prop :=
  F' = F.subst w t
\end{leansrc}
\declhead{BernsteinLogic.Core.Catalogue}{PEqElimData}{\badge{BadgeGreen}{definition bound}}
\begin{leansrc}
inductive PEqElimData : Hypernet → Sequent → Type where
  | mk (H : Hypernet) (seq : Sequent)
       (eqPremise   : Sequent) (eqProof   : ProofObject H eqPremise)
       (bodyPremise : Sequent) (bodyProof : ProofObject H bodyPremise)
       (wire : Wire) (termT termU : Term) (bodyForm : Formula)
       (preEq     : seq.pre = eqPremise.pre)
       (preMatch  : eqPremise.pre = bodyPremise.pre)
       (eqPost    : eqPremise.post = .atom .eq termT termU)
       (tyEq      : ∃ s, termT.bshape H = some s ∧ termU.bshape H = some s)
       (bodySubT  : substituteInFormula bodyForm wire termT bodyPremise.post)
       (conclSubU : substituteInFormula bodyForm wire termU seq.post)
       : PEqElimData H seq
\end{leansrc}
\ruleblock{Correspondence (Steele display $\leftrightarrow$ Lean)}
\badge{BadgeGreen}{reviewer-certified against the displayed texts}
\sourcenote{The dictionary below is a reviewer-certified statement of how the mechanized statements above interpret the displayed rule: display metavariable $\mapsto$ Lean binder, each display hypothesis $\mapsto$ the premise or side condition carrying it, the display conclusion $\mapsto$ the goal --- with any deliberate difference in width between display and mechanization named outright. Certification is recorded as a hash triple over (the displayed cell, the extracted statement texts, this dictionary\textquotesingle s own text); if any of the three changes, this block turns amber until a reviewer re-certifies. The correspondence claim itself remains editorial --- reviewer-grade, not machine-checked; its machine-checked successor types each soundness statement against the rule as a structured object, and supersedes this block per rule as it lands.}
Judgment frame as in \rulename{Ax}; sub-derivation premises as at
\rulename{Weaken}.
\begin{itemize}
\item $t, u \mapsto$ \texttt{termT termU : Term}; the hole $a \mapsto$
  \texttt{wire : Wire}; the body $\Psi \mapsto$ \texttt{bodyForm :
  Formula}; $\Phi \mapsto$ \texttt{seq.pre}, shared with both premise
  sequents through \texttt{hpre} and \texttt{hpreMatch}.
\item the equality premise $\Phi \vdash t = u$ $\mapsto$ \texttt{ihEq}
  with the shape tie \texttt{heqPost} $:$ \texttt{eqPremise.post =
  .atom .eq termT termU}.
\item the body premise $\Phi \vdash \Psi\sub{a}{t}$ $\mapsto$
  \texttt{ihBody} $:$ \texttt{bodyPremise.semanticallyValid H}, its
  post tied to $\Psi\sub{a}{t}$ by \texttt{hbodySubT}; the goal
  post is tied to $\Psi\sub{a}{u}$ by \texttt{hconclSubU}. Both ties
  run through \texttt{substituteInFormula}, quoted in full below: the
  one-line equation \texttt{F\textquotesingle{} = F.subst w t} over
  plain leaf replacement --- the grammar has no binders, so no capture
  or freshness side condition exists to omit.
\item the typing premise $\exists s.\; \bshape_H(t) = s = \bshape_H(u)$
  $\mapsto$ \texttt{htyEq}: existence of one shape \texttt{s} with
  \texttt{termT.bshape H = some s} and \texttt{termU.bshape H = some
  s} --- the \texttt{some} carries the definedness the
  display\textquotesingle s equation chain asserts by letting $s$
  range over shapes.
\item conclusion $\Phi \vdash \Psi\sub{a}{u}$ $\mapsto$ the goal
  \texttt{seq.semanticallyValid H}. Certified alongside as the
  semantic engine: \texttt{Term.evalQ\_subst\_congr} and
  \texttt{Formula.holds\_subst\_congr} --- substitution congruence
  along equal evaluations, the two lemmas the soundness proof routes
  through.
\end{itemize}
No width difference in either direction: statement and display carry
the same typing gate (the commentary\textquotesingle s counterexample
is what its absence admits), and every hypothesis is consumed.

\ruleblock{Lean verifier receipts}
\sourcenote{module \texttt{BernsteinLogic.Soundness}; bound build \texttt{a20d90e66c13\ldots} (bound from receipt a9c0e34d3e3c9c8c, 2026-08-03T21:30:36Z). A receipt records that named checking layers accepted a build: (module, compiled-artifact SHA-256, environment SHA-256, checker layers), minted when the checking gate ran; the store receipts attesting the bound build are listed with the layers each run credited. Nothing here asks the reader to trust the store: every layer is an independently runnable program over the quoted sources --- rebuild the module and re-run the checkers to reproduce the verdict. Receipts attesting superseded builds credit no layer and are omitted from this rendering (the count of omissions, if any, is noted below).}
\begin{receipttable}
\receiptrow{a9c0e34d3e3c9c8c}{2026-08-03T21:30:36Z}{build+decorrelated}{\badge{BadgeGreen}{current}}
\end{receipttable}
\sourcenote{21 further receipt(s) in the store attest superseded builds of this module and are omitted from this rendering.}
\sourcenote{module \texttt{BernsteinLogic.Core.Catalogue}; bound build \texttt{cbc92e976a1d\ldots} (bound from receipt 5abdb3589707b5df, 2026-08-03T21:30:33Z). A receipt records that named checking layers accepted a build: (module, compiled-artifact SHA-256, environment SHA-256, checker layers), minted when the checking gate ran; the store receipts attesting the bound build are listed with the layers each run credited. Nothing here asks the reader to trust the store: every layer is an independently runnable program over the quoted sources --- rebuild the module and re-run the checkers to reproduce the verdict. Receipts attesting superseded builds credit no layer and are omitted from this rendering (the count of omissions, if any, is noted below).}
\begin{receipttable}
\receiptrow{24804e66edba0bc9}{2026-08-03T20:25:09Z}{build+decorrelated}{\badge{BadgeGreen}{current}}
\receiptrow{5abdb3589707b5df}{2026-08-03T21:30:33Z}{build+decorrelated}{\badge{BadgeGreen}{current}}
\end{receipttable}
\sourcenote{18 further receipt(s) in the store attest superseded builds of this module and are omitted from this rendering.}
\layerline{\layeritem{build}{lake build (elaboration)}{\badge{BadgeGreen}{current}}\layeritem{replay}{leanchecker (official kernel replay)}{\badge{BadgeAmber}{stale only}}\layeritem{hardened}{safe\_verify (hardened screen)}{\badge{BadgeAmber}{stale only}}\layeritem{decorrelated}{lean4export + nanoda (decorrelated kernel)}{\badge{BadgeGreen}{current}}}
\sourcenote{\badge{BadgeGray}{comparator: none} no fixture in the three-implementation comparator harness (compiled / kernel-interpreted / reference-interpreted) exercises this rule.}
\ruleblock{HOL4 statement}
\sourcenote{theory \texttt{bern\_catalogue}, HOL4 development script \texttt{\texttt{bern\_catalogueScript.sml}} \badge{BadgeGreen}{script bound}}
\holhead{Definition termSubst\_def}
\begin{holsrc}
Definition termSubst_def:
  termSubst w t (TmVar w') = (if w' = w then t else TmVar w') ∧
  termSubst w t (TmConst c) = TmConst c ∧
  termSubst w t (TmAdd l r) = TmAdd (termSubst w t l) (termSubst w t r) ∧
  termSubst w t (TmMul l r) = TmMul (termSubst w t l) (termSubst w t r) ∧
  termSubst w t (TmSub l r) = TmSub (termSubst w t l) (termSubst w t r) ∧
  termSubst w t (TmNeg u) = TmNeg (termSubst w t u) ∧
  termSubst w t (TmIdx u i) = TmIdx (termSubst w t u) i
End
\end{holsrc}
\holhead{Definition formulaSubst\_def}
\begin{holsrc}
Definition formulaSubst_def:
  formulaSubst w t (FAtom b l r) =
    FAtom b (termSubst w t l) (termSubst w t r) ∧
  formulaSubst w t (FInRegion u sx es) = FInRegion (termSubst w t u) sx es ∧
  formulaSubst w t (FAnd l r) =
    FAnd (formulaSubst w t l) (formulaSubst w t r) ∧
  formulaSubst w t (FOr l r) =
    FOr (formulaSubst w t l) (formulaSubst w t r) ∧
  formulaSubst w t FTt = FTt ∧
  formulaSubst w t FFf = FFf
End
\end{holsrc}
\holhead{Definition substituteInFormulaB\_def}
\begin{holsrc}
Definition substituteInFormulaB_def:
  substituteInFormulaB fm w t fm' ⇔ fm' = formulaSubst w t fm
End
\end{holsrc}
\ruleblock{HOL verifier receipts}
\badge{BadgeBlue}{documentary} \texttt{hol-strict} CLEAN census over the ported checker theory:
\texttt{bern\_catalogue}, 472 theorems, 0 cheat-tagged, 0 nonstandard
axioms (whitelist: \texttt{DISK\_THM}; the four \texttt{bool} axioms) ---
verification run of 2026-07-22 on an isolated build host, against a
content-addressed checkout of the exact sources cited here --- the
earlier stale-run caveat (a 2026-07-18 script revision postdating the
prior census) is retired by this run, which attests the current text.
There is no signed HOL receipt store yet --- these citations are
documentary. Scope note, honest: the substitution core of
this rule\textquotesingle s check is what is ported; the broadcast-shape
side condition (\texttt{tyEq}) rides the hypernet typing island, which
the port deliberately defers to its proof-object stage.
Semantic stratum, also ported: the soundness twin \texttt{sound\_eqElim}
(theory \texttt{bern\_tr\_fragclose}: \texttt{hol-strict} clean, 26
theorems, 0 offences, verification run of 2026-07-22 on an isolated
build host, content-addressed at the same sources), consumed by the
certified-fragment keystone (\texttt{checkPf\_sound\_frag}, theory
\texttt{bern\_tr\_splice}, 20 theorems, 0 offences, same run): a
reference-checker accept inside the proven wire fragment yields sequent
validity.

\sourcenote{HOL receipts are documentary: quoted results of \texttt{hol-strict} runs
and census rows from the HOL4 development\textquotesingle s build
evidence, not minted receipt artefacts. A HOL receipt store (the Lean
store\textquotesingle s twin) is named future work.
}
\ruleblock{Proof}

\sourcenote{Genre: core rule (tier 1) --- an ordinary mathematical proof at
Putnam-exam-solution detail. Every hypothesis of the rule is used below and
remarked on; the formal-anchor table maps each step to the Lean declaration
that carries it. This section also fixes the semantic setting that the
other equality and conversion sections cite.}

\textbf{Setting.}\enspace Fix a hypernet $H$. A \emph{valuation} $\sigma$
assigns each \emph{typed} wire a tensor of its declared shape over its
declared dtype (an extended dtype admits infinite entries; untyped wires
are assigned nothing), and $\sigma$ is \emph{admissible} when every
primitive application in $H$ satisfies its semantic relation --- as at
the \rulename{Ax} setting. The rational term evaluation
$\mathrm{ev}_\sigma(t)$ is \emph{partial} and the evaluator is structural
and \emph{strict} --- a compound term evaluates exactly when its parts
evaluate and the constructor's own side condition is met: shapes joining
for the binary operations, nothing further for negation, a rank-1
operand with in-range index for projection --- so a mention of an
untyped wire, a wire holding an infinite value of an extended dtype, a
broadcast-incompatible operation, or a projection whose operand is not
rank-1 or whose index is out of range yields no value. A comparison atom
$t_1 \bowtie t_2$ \emph{holds} at $\sigma$
when both sides evaluate and the evaluated tensors stand in the relation
$\bowtie$ \emph{elementwise under broadcast}: equal shapes compare entry
by entry; a rank-0 (scalar) tensor against a tensor of any shape compares
against every entry; two unequal positive-rank shapes never compare. A
region atom holds when its term evaluates to a rank-1 tensor of extent
the region's ambient dimension whose entries form a rational point of
the region --- a rank-0 evaluation is never coerced to a one-dimensional
point. A sequent $\Phi \vdash \Psi$ is \emph{semantically valid} when
every admissible $\sigma$ satisfying $\Phi$ satisfies $\Psi$;
counterexamples in the equality and conversion sections read, as at the
\rulename{Ax} setting, over application-free hypernets, where every
valuation is admissible, so an unsound conclusion genuinely falsifies.
Note two
consequences that drive everything below: equality atoms are
\emph{evaluation-level} facts, not syntactic ones; and a holding equality
between tensors of \emph{different} broadcast shapes (a scalar against a
vector) does not make the two sides interchangeable --- it only says the
scalar matches every entry.

\textbf{Claim} (\texttt{sound\_eqElim}).\enspace The rule's payload
supplies a hole wire $a$, terms $t, u$, a body formula $\Psi$, and checked
side conditions: the two premises share the composite premise $\Phi$; the
first premise concludes the atom $t = u$; the second premise concludes
$\Psi\sub{a}{t}$; the conclusion is $\Psi\sub{a}{u}$; and (the typing
condition) $\exists s.\ \bshape_H(t) = s = \bshape_H(u)$ --- both terms
carry a \emph{common defined} broadcast shape in $H$. If both premises are
semantically valid, so is $\Phi \vdash \Psi\sub{a}{u}$.

\begin{blemma}[substitution congruence, term level]\label{lem:eqelim-term}
Fix $\sigma$, and suppose $\mathrm{ev}_\sigma(t) = \mathrm{ev}_\sigma(u)$
as partial values (both undefined, or both defined and equal). Then for
every term $e$,
$\mathrm{ev}_\sigma(e\sub{a}{t}) = \mathrm{ev}_\sigma(e\sub{a}{u})$.
\end{blemma}
\begin{proof}
Structural induction on $e$. The base case $e = a$ is the hypothesis; the
base cases $e$ a constant or a different wire are reflexivity; every
compound constructor is a congruence of the evaluator (the evaluator of a
sum, product, difference, negation, or projection is a function of the
evaluations of the parts). Two remarks. First, the hypothesis and
conclusion are equalities of \emph{partial} values --- the lemma never
manufactures definedness, it only transports whatever is there; in the
theorem, definedness will come out of the equality premise. Second, the
term grammar has no binders, so substitution is plain leaf replacement
and no capture condition is needed anywhere.
\end{proof}

\begin{blemma}[substitution congruence, formula level]\label{lem:eqelim-form}
Under the hypothesis of Lemma~\ref{lem:eqelim-term}: for every formula
$F$,\ $\sigma \models F\sub{a}{t}$ iff $\sigma \models F\sub{a}{u}$.
\end{blemma}
\begin{proof}
Induction on $F$. Comparison atoms substitute in both sides and satisfy
the same relation by Lemma~\ref{lem:eqelim-term}; a region atom
substitutes in its term slot only (its simplex and extension-set data
carry no wires), so its holding --- an assertion about the evaluated
tensor --- transports the same way; $\wedge$ and $\vee$ are the inductive
steps; $\top$ and $\bot$ are untouched by substitution.
\end{proof}

\begin{blemma}[the shape pin]\label{lem:eqelim-pin}
Suppose $\sigma \models t = u$, and suppose $t$ and $u$ carry a common
defined broadcast shape $s$ in $H$. Then
$\mathrm{ev}_\sigma(t) = \mathrm{ev}_\sigma(u)$ --- equality of
\emph{tensors}, not merely the elementwise relation.
\end{blemma}
\begin{proof}
The holding atom yields evaluations $A$ of $t$ and $B$ of $u$ in the
elementwise-equality relation. Broadcast-shape soundness
(\texttt{Term.evalQ\_bshape}) pins $A.\mathrm{shape} = s =
B.\mathrm{shape}$: a defined evaluation of a $\bshape$-typed term lands
exactly at the declared broadcast shape. Elementwise equality at
\emph{equal} shapes is equality of the (equal-length) entry lists,
position by position, hence equality of the tensors. This is the one step
that consumes the typing condition, and it is false without it: the
elementwise relation between a rank-0 tensor and a vector says only that
the scalar matches every entry (\autoref{rem:eqelim-cex}).
\end{proof}

\begin{bthm}[soundness of \rulename{EqElim}]\label{thm:eqelim-sound}
Under the payload's checked hypotheses, if $\Phi \vdash t = u$ and
$\Phi \vdash \Psi\sub{a}{t}$ are semantically valid, then so is
$\Phi \vdash \Psi\sub{a}{u}$.
\end{bthm}
\begin{proof}
Fix an admissible $\sigma$ with $\sigma \models \Phi$. By validity of the
first premise, $\sigma \models t = u$; with the typing condition,
Lemma~\ref{lem:eqelim-pin} upgrades this to
$\mathrm{ev}_\sigma(t) = \mathrm{ev}_\sigma(u)$ (both defined and equal,
so in particular equal as partial values). By validity of the second
premise --- its premise formula is $\Phi$ itself, by the checked premise
match --- $\sigma \models \Psi\sub{a}{t}$. Lemma~\ref{lem:eqelim-form}
transports this to $\sigma \models \Psi\sub{a}{u}$, which is the
conclusion formula by the checked substitution equation.
\end{proof}

\begin{bremark}[why the typing condition is not decoration]
\label{rem:eqelim-cex}
Drop the common-shape condition and the rule is \emph{unsound}
(orientation of the equality atom flipped here to spare an
\rulename{EqSym} step). Let $v$ be a wire of shape $[1]$ over $\Q$, let
$\Delta$ be the segment $\mathrm{conv}(1,2) \subset \Q^1$, and take
\[
\Phi \;:=\; \bigl(v = \tfrac{3}{2}\bigr) \;\wedge\; \bigl(v \in
\mathcal{R}(\Delta, \emptyset)\bigr).
\]
$\Phi$ is satisfiable: set $v := [\,3/2\,]$; the equality holds by rank-0
broadcast (the scalar $3/2$ matches the single entry), and the region
atom holds. \rulename{Ax} and $\wedge$E derive $\Phi \vdash v = 3/2$ and
$\Phi \vdash v \in \mathcal{R}(\Delta, \emptyset)$; the un-gated rule
with body $\Psi := (a \in \mathcal{R}(\Delta, \emptyset))$, $t := v$,
$u := 3/2$ then concludes
$\Phi \vdash (\tfrac{3}{2} : \mathrm{Term}) \in \mathcal{R}(\Delta,
\emptyset)$ --- which is \emph{unsatisfiable}: a constant term evaluates
to a rank-0 tensor, and a region atom demands a rank-1 tensor of extent
equal to the ambient dimension. A satisfiable premise, an unsatisfiable
conclusion. The root cause is precisely the gap between ``elementwise
equal under broadcast'' and ``interchangeable'': $3/2 = [\,3/2\,]$ holds
without the terms being exchangeable in region contexts. With the typing
condition the payload is uninhabitable here ---
$\bshape_H(\tfrac{3}{2}) = [\,]$ while $\bshape_H(v) = [1]$, and no
common $s$ exists --- so the certificate refuses. Note what the repair
does \emph{not} do: broadcast equalities remain derivable and usable as
facts; what is forbidden is exactly using one as a Leibniz license.
\end{bremark}

\begin{bremark}[the equational kit; relation to practice]
\label{rem:eqelim-kit}
\rulename{EqElim} with \rulename{EqRefl} and \rulename{EqSym} (next two
sections) completes a certificate-side equational kit: reflexivity and
symmetry are primitive, and transitivity is a one-node derivation ---
from $\Phi \vdash t = u$ and $\Phi \vdash u = w$, apply \rulename{EqElim}
with body $\Psi := (t = a)$, the hole wire $a$ chosen fresh for $t$
(substitution replaces \emph{every} occurrence of $a$, so a hole
occurring in $t$ would rewrite $t$ itself), rewriting the $u$ in the
\emph{first}
premise (taken as body premise, $\Psi\sub{a}{u}$) to $w$; the derived
node inherits the common-shape typing obligation on $(u, w)$ ---
broadcast equality is not transitive without it (two tensors of unequal
positive-rank shapes never compare, yet each can stand equal to one
scalar between them: $[\,0\,] = 0$ and $0 = [\,0\ \ 0\,]$ both hold
while $[\,0\,] = [\,0\ \ 0\,]$ does not). On the producer side,
equational goals are typically discharged by congruence-closure engines
\cite{nelsonoppen1980,nieuwenhuisoliveras2007}; what ships in a
certificate is not the closure computation but its residue: a chain of
\rulename{EqRefl}/\rulename{EqSym}/\rulename{EqElim} nodes, each checked
by syntactic recomputation of one substitution.
\end{bremark}

\begin{anchors}
Setting: valuation model, partial strict evaluation, atom and region semantics & \texttt{Hypernet.Valuation}, \texttt{Term.evalQ}, \texttt{QTensor.rel2}, \texttt{Formula.holds} (Layer 10)\\
Lemma \ref{lem:eqelim-term} & \texttt{Term.evalQ\_subst\_congr}\\
Lemma \ref{lem:eqelim-form} & \texttt{Formula.holds\_subst\_congr}\\
Lemma \ref{lem:eqelim-pin} & \texttt{Term.evalQ\_bshape}; \texttt{QTensor.eq\_of\_rel2}; payload field \texttt{tyEq}\\
substitution equations & \texttt{substituteInFormula} = \texttt{Formula.subst} (payload fields \texttt{bodySubT}, \texttt{conclSubU})\\
premise sharing & payload fields \texttt{preEq}, \texttt{preMatch}\\
Theorem \ref{thm:eqelim-sound} & \texttt{sound\_eqElim}; wired as core \texttt{RuleData.eqElim} into \texttt{ProofObject.sound}\\
Remark \ref{rem:eqelim-cex} & the counterexample inventory, \texttt{Soundness.lean} header\\
\end{anchors}


\rulesection{9}{EqRefl}{core rule}
\label{blcatrule-EqRefl}
\ruleblock{The rule (Steele metanotation)}
\sourcenote{as realized in the mechanized development; source row: \texttt{\texttt{rules.toml}} --- \emph{EqRefl}. The metanotation statement is typeset by hand for notational quality; its mechanized ground truth is the Lean statement below, extracted from the sources at every build of this book and hash-bound. The rule \emph{roster} is machine-checked (every mechanized rule tag round-trips to exactly one catalogue row); the statement \emph{text} is editorial, checked by eye against the extracted Lean beside it.}
\begin{mathpar}
\inferrule*[Right=\rulename{EqRefl}]{\exists s.\; \bshape_H(t) = s \\ \wires(t) \text{ finitely typed in } H}{\Phi \vdash t = t}
\end{mathpar}

\par\medskip
\rulename{EqRefl} derives $\Phi \vdash t = t$, premise-free, for an
arbitrary context $\Phi$ (a fact true in the empty context is true under
any premise; the proof never reads $\Phi$). A producer uses it to seed
equalities: the starting point of an \rulename{EqElim} rewrite chain
(transitivity, in particular, derives from these three rules and begins
here), and the $t = t$ obligations that alias unpacks and worked
derivations demand.

That the rule has side conditions at all is its content. Equality atoms
in this calculus are facts about \emph{evaluation}, and evaluation is
partial --- a wire of an extended dtype can hold $\pm\infty$ and fail
the rational read; a shape-incompatible sum never evaluates. So
``$t = t$'' costs something: the checker verifies the two typing
conditions the semantics actually needs --- the term\textquotesingle s
broadcast shape is \emph{defined} in the hypernet, and every free wire
is typed at a \emph{finite} dtype. Under exactly these, evaluation is
total at every valuation and reflexivity is free; drop either and the
un-gated rule asserts an invalid sequent under the empty premise --- in
one case a conclusion failing only at extended-valued assignments, in
the other one unsatisfiable outright (the proof section gives one
concrete refutation per field).
Nothing else is consulted: no premise, no admissibility, no evaluation
--- both conditions are static typing checks against the hypernet.

\cataloguenote{Where \rulename{Norm} applies, \rulename{EqRefl} is its diagonal; it is
kept primitive so the equational kit stands in fragments that omit the
polynomial engine entirely.
}
\ruleblock{Lean statement}
\declhead{BernsteinLogic.Soundness}{Term.evalQ\_defined\_of\_bshape\_fin}{\badge{BadgeGreen}{statement bound}}
\begin{leansrc}
/-- Totality of evaluation under the B13 typing pack: a `bshape`-typed
term over finitely-typed wires evaluates. Structural — no polynomial
view, so idx projections are covered (their rank-1 + range gates are
part of `bshape`). -/
theorem Term.evalQ_defined_of_bshape_fin {H : Hypernet} (σ : H.Valuation) :
    ∀ (t : Term) (s : Shape), t.bshape H = some s → t.finWires H →
    ∃ q, Term.evalQ σ t = some q ∧ q.shape = s
\end{leansrc}
\declhead{BernsteinLogic.Soundness}{QTensor.rel2\_refl}{\badge{BadgeGreen}{statement bound}}
\begin{leansrc}
/-- Reflexive relations lift to `rel2` on any tensor against itself. -/
theorem QTensor.rel2_refl {r : ℚ → ℚ → Prop} (hr : ∀ x, r x x)
    (a : QTensor) : QTensor.rel2 r a a
\end{leansrc}
\declhead{BernsteinLogic.Soundness}{sound\_eqRefl}{\badge{BadgeGreen}{statement bound}}
\begin{leansrc}
/-- EqRefl: under the B13 typing pack, t evaluates, and elementwise
equality against itself holds. -/
theorem sound_eqRefl {H : Hypernet} {seq : Sequent} {t : Term}
    (htyDef : ∃ s, t.bshape H = some s) (hfin : t.finWires H)
    (hpost : seq.post = .atom .eq t t) :
    seq.semanticallyValid H
\end{leansrc}
\par\medskip
\textbf{The two typing predicates, in full.}\enspace Both payload fields
are static checks against the hypernet\textquotesingle s wire typing ---
no valuation appears. \texttt{Term.bshape} is the broadcast-tolerant
shape synthesis (scalars join with any shape; the projection case carries
its own rank and range gates); \texttt{Term.finWires} demands a finite
(non-extended) dtype at every free wire.

\declhead{BernsteinLogic.Core.Catalogue}{Term.bshape}{\badge{BadgeGreen}{definition bound}}
\begin{leansrc}
/-- Broadcast-tolerant shape of a term: shapes must agree up to rank-0
broadcast (a scalar joins with any shape); dtypes are NOT consulted. -/
def Term.bshape (H : Hypernet) : Term → Option Shape
  | .var w => (H.wireType w).map (·.shape)
  | .const _ => some []
  | .add l r | .mul l r | .sub l r =>
    match l.bshape H, r.bshape H with
    | some s₁, some s₂ =>
      if s₁ = s₂ then some s₁
      else if s₁ = [] then some s₂
      else if s₂ = [] then some s₁
      else none
    | _, _ => none
  | .neg t => t.bshape H
  | .idx t i =>
    -- rank-1 argument, in-range index, scalar result (matches `evalQ`'s
    -- projection semantics — G1).
    match t.bshape H with
    | some [e] => if i < e then some [] else none
    | _ => none
\end{leansrc}
\declhead{BernsteinLogic.Core.Catalogue}{Term.finWires}{\badge{BadgeGreen}{definition bound}}
\begin{leansrc}
/-- Every wire of the term is typed at a finite (non-ext) dtype — the
carrier of such a wire contains no ±∞, so entry reads are exact
rationals. -/
def Term.finWires (H : Hypernet) (t : Term) : Prop :=
  ∀ w ∈ t.freeWires, ∃ τ, H.wireType w = some τ ∧ τ.dtype.ext = Ext.fin
\end{leansrc}
\declhead{BernsteinLogic.Core.Catalogue}{PEqReflData}{\badge{BadgeGreen}{definition bound}}
\begin{leansrc}
inductive PEqReflData : Hypernet → Sequent → Type where
  | mk (H : Hypernet) (seq : Sequent)
       (term : Term)
       (tyDef  : ∃ s, term.bshape H = some s)
       (fin    : term.finWires H)
       (postEq : seq.post = .atom .eq term term)
       : PEqReflData H seq
\end{leansrc}
\ruleblock{Correspondence (Steele display $\leftrightarrow$ Lean)}
\badge{BadgeGreen}{reviewer-certified against the displayed texts}
\sourcenote{The dictionary below is a reviewer-certified statement of how the mechanized statements above interpret the displayed rule: display metavariable $\mapsto$ Lean binder, each display hypothesis $\mapsto$ the premise or side condition carrying it, the display conclusion $\mapsto$ the goal --- with any deliberate difference in width between display and mechanization named outright. Certification is recorded as a hash triple over (the displayed cell, the extracted statement texts, this dictionary\textquotesingle s own text); if any of the three changes, this block turns amber until a reviewer re-certifies. The correspondence claim itself remains editorial --- reviewer-grade, not machine-checked; its machine-checked successor types each soundness statement against the rule as a structured object, and supersedes this block per rule as it lands.}
Judgment frame as in \rulename{Ax}.
\begin{itemize}
\item $t \mapsto$ the binder \texttt{t : Term}; the two displayed
  typing premises $\mapsto$ \texttt{htyDef}: existence of a shape
  \texttt{s} with \texttt{t.bshape H = some s} (defined broadcast
  shape --- the \texttt{some} is the definedness), and \texttt{hfin}
  $:$ \texttt{t.finWires H} (every free wire typed at a finite dtype,
  quoted below). Both are static checks against the hypernet; no
  valuation appears.
\item conclusion $\Phi \vdash t = t$ $\mapsto$ the goal under
  \texttt{hpost} $:$ \texttt{seq.post = .atom .eq t t}; the arbitrary
  context $\Phi$ is premise-freeness by silence --- no hypothesis
  mentions \texttt{seq.pre}, so the statement holds under every
  premise (the same representation the \rulename{PConstExact}
  dictionary describes).
\item certified alongside as the semantic engine:
  \texttt{Term.evalQ\_defined\_of\_bshape\_fin} --- under exactly the
  two typing conditions, evaluation is total at every valuation, at
  the synthesized shape --- and \texttt{QTensor.rel2\_refl}, the
  reflexivity of the elementwise comparison lifting.
\end{itemize}
No width difference in either direction: both typing hypotheses are
consumed, and the proof section refutes the rule with either one
dropped.

\ruleblock{Lean verifier receipts}
\sourcenote{module \texttt{BernsteinLogic.Soundness}; bound build \texttt{a20d90e66c13\ldots} (bound from receipt a9c0e34d3e3c9c8c, 2026-08-03T21:30:36Z). A receipt records that named checking layers accepted a build: (module, compiled-artifact SHA-256, environment SHA-256, checker layers), minted when the checking gate ran; the store receipts attesting the bound build are listed with the layers each run credited. Nothing here asks the reader to trust the store: every layer is an independently runnable program over the quoted sources --- rebuild the module and re-run the checkers to reproduce the verdict. Receipts attesting superseded builds credit no layer and are omitted from this rendering (the count of omissions, if any, is noted below).}
\begin{receipttable}
\receiptrow{a9c0e34d3e3c9c8c}{2026-08-03T21:30:36Z}{build+decorrelated}{\badge{BadgeGreen}{current}}
\end{receipttable}
\sourcenote{21 further receipt(s) in the store attest superseded builds of this module and are omitted from this rendering.}
\sourcenote{module \texttt{BernsteinLogic.Core.Catalogue}; bound build \texttt{cbc92e976a1d\ldots} (bound from receipt 5abdb3589707b5df, 2026-08-03T21:30:33Z). A receipt records that named checking layers accepted a build: (module, compiled-artifact SHA-256, environment SHA-256, checker layers), minted when the checking gate ran; the store receipts attesting the bound build are listed with the layers each run credited. Nothing here asks the reader to trust the store: every layer is an independently runnable program over the quoted sources --- rebuild the module and re-run the checkers to reproduce the verdict. Receipts attesting superseded builds credit no layer and are omitted from this rendering (the count of omissions, if any, is noted below).}
\begin{receipttable}
\receiptrow{24804e66edba0bc9}{2026-08-03T20:25:09Z}{build+decorrelated}{\badge{BadgeGreen}{current}}
\receiptrow{5abdb3589707b5df}{2026-08-03T21:30:33Z}{build+decorrelated}{\badge{BadgeGreen}{current}}
\end{receipttable}
\sourcenote{18 further receipt(s) in the store attest superseded builds of this module and are omitted from this rendering.}
\layerline{\layeritem{build}{lake build (elaboration)}{\badge{BadgeGreen}{current}}\layeritem{replay}{leanchecker (official kernel replay)}{\badge{BadgeAmber}{stale only}}\layeritem{hardened}{safe\_verify (hardened screen)}{\badge{BadgeAmber}{stale only}}\layeritem{decorrelated}{lean4export + nanoda (decorrelated kernel)}{\badge{BadgeGreen}{current}}}
\sourcenote{\badge{BadgeGray}{comparator: none} no fixture in the three-implementation comparator harness (compiled / kernel-interpreted / reference-interpreted) exercises this rule.}
\ruleblock{HOL4 statement}
\sourcenote{theory \texttt{bern\_fam\_resid}, HOL4 development script \texttt{\texttt{bern\_fam\_residScript.sml}} \badge{BadgeGreen}{script bound}}
\holhead{Definition checkFam\_eqrefl\_def}
\begin{holsrc}
Definition checkFam_eqrefl_def:
  checkFam_eqrefl (cols:columns) (ct:ctx) (i:num) ⇔
    case frSeq cols ct i of
    | NONE => F
    | SOME sq =>
      case frTerm cols ct i 0 of
      | NONE => F
      | SOME t =>
          IS_SOME (termBshape ct.ws t) ∧ finWiresB ct.ws t ∧
          sq.post = FAtom BowEq t t
End
\end{holsrc}
\holhead{Theorem famEq\_eqrefl}
\begin{holsrc}
Theorem famEq_eqrefl:
  wfColumns cols ct ∧ mkCtx cols = SOME ct ∧
  parseCert (cols.hdrImage ++ rest) = SOME (ws, pf) ∧
  nodeCorr cols ct i (CertPf sq nd) ∧
  FST (THE (tagOf cols i)) = 23 ⇒
  (checkFam_eqrefl cols ct i ⇔ armEqRefl ws sq nd)
\end{holsrc}
\ruleblock{HOL verifier receipts}
\badge{BadgeBlue}{documentary} Documentary citations. Family validation: \texttt{bern\_fam\_resid} built GREEN with all ten \texttt{famEq} arms validated, \texttt{hol-strict} clean; spine census: \texttt{bern\_joins} GREEN under \texttt{hol-strict} (55 theorems, 0 offences, 208-theory ancestry including this theory), \texttt{checkJoins\_eq\_checkPf} verified verbatim in the built signature. The cited \texttt{famEq} lemma is the kernel-checked per-node equivalence: under well-formed columns (\texttt{wfColumns}/\texttt{mkCtx}/\texttt{nodeCorr}) at this rule\textquotesingle s wire tag, the columnar family check coincides with the flat reference walk\textquotesingle s arm for this rule. There is no signed HOL receipt store yet --- these citations are documentary.
Semantic stratum, also ported: the soundness twin \texttt{sound\_eqRefl}
(theory \texttt{bern\_tr\_ready8}: \texttt{hol-strict} clean, 8
theorems, 0 offences, verification run of 2026-07-22 on an isolated
build host, content-addressed at the same sources), consumed by the
certified-fragment keystone (\texttt{checkPf\_sound\_frag}, theory
\texttt{bern\_tr\_splice}, 20 theorems, 0 offences, same run): a
reference-checker accept inside the proven wire fragment yields sequent
validity.

\sourcenote{HOL receipts are documentary: quoted results of \texttt{hol-strict} runs
and census rows from the HOL4 development\textquotesingle s build
evidence, not minted receipt artefacts. A HOL receipt store (the Lean
store\textquotesingle s twin) is named future work.
}
\ruleblock{Proof}

\sourcenote{Genre: core rule (tier 1). Setting (valuations, partial
evaluation, elementwise-under-broadcast atom semantics) as fixed in the
\rulename{EqElim} section.}

\textbf{Claim} (\texttt{sound\_eqRefl}).\enspace The payload supplies a
term $t$ and two checked typing conditions: $\bshape_H(t)$ is
\emph{defined} (some shape $s$), and every free wire of $t$ is typed at a
\emph{finite} (non-extended) dtype. Then $\Phi \vdash t = t$ is
semantically valid for an \emph{arbitrary} premise $\Phi$ --- the rule is
premise-free, and $\Phi$ rides along unread (a fact true in the empty
context is true under any premise). Indeed the proof below never reads
$\Phi$, the admissibility of $\sigma$, or the premise-satisfaction
hypothesis at all: the conclusion holds at \emph{every} valuation.

That anything must be checked here at all is the point of the section.
Equality atoms are evaluation-level facts (\rulename{EqElim} setting), so
``$t = t$'' is not free: it asserts that $t$ \emph{evaluates} --- and
evaluation is partial.

\begin{blemma}[totality under the typing pack]\label{lem:eqrefl-total}
If $\bshape_H(t) = s$ and every free wire of $t$ is finitely typed, then
at every valuation $\sigma$ the evaluation $\mathrm{ev}_\sigma(t)$ is
defined, with shape $s$.
\end{blemma}
\begin{proof}
Structural induction on $t$; both typing conditions descend to
subterms (a subterm's free wires are among $t$'s, and $\bshape$ is
synthesized bottom-up, so a defined $\bshape$ for $t$ is defined on
every subterm). A wire read succeeds because the wire is typed at a
finite dtype --- an untyped wire declares no dtype at all, so the
finite-wires condition already excludes it: the valuation's entries lie
in the dtype's value set, which for a finite dtype contains no
$\pm\infty$, so the whole-tensor rational read succeeds, at the
declared shape. A constant evaluates to its rank-0 tensor. For a
compound term, the induction hypotheses give defined operands whose
shapes are, by the definition of $\bshape$ (quoted in full beside this
section), either equal or joinable through rank-0 broadcast --- and on
such shapes the elementwise operations are total; the result lands at
the join. A projection's rank-1-and-in-range gates are part of
$\bshape$'s definedness, so the projected read is in range.
(\texttt{Term.evalQ\_defined\_of\_bshape\_fin}.)
\end{proof}

\begin{blemma}[reflexivity lifts]\label{lem:eqrefl-lift}
For any tensor $A$, the elementwise-under-broadcast equality relation
holds between $A$ and $A$.
\end{blemma}
\begin{proof}
The shapes are equal, so the relation reads the entry list zipped against
itself: every pair on the diagonal is $(x, x)$.
(\texttt{QTensor.rel2\_refl}.)
\end{proof}

\begin{bthm}[soundness of \rulename{EqRefl}]\label{thm:eqrefl-sound}
Under the payload's typing conditions, $\Phi \vdash t = t$ is
semantically valid for every $\Phi$.
\end{bthm}
\begin{proof}
Fix $\sigma$. Lemma~\ref{lem:eqrefl-total} gives
$\mathrm{ev}_\sigma(t) = A$ for some tensor $A$;
Lemma~\ref{lem:eqrefl-lift} closes the atom with both sides read as $A$.
\end{proof}

\begin{bremark}[why each typing condition is load-bearing]
\label{rem:eqrefl-cex}
Drop the finite-wires condition: let $w$ be a wire of an
\emph{extended} dtype (its value set admits $+\infty$), and let $\sigma$
assign it a tensor with an infinite entry --- admissible, since nothing
constrains a bare wire. The rational read of $w$ fails, $t := w$ does not
evaluate, and the atom $w = w$ is \emph{false} at $\sigma$; with
$\Phi := \top$ the sequent $\top \vdash w = w$ is refuted. Drop
shape-definedness: take $t := w_1 + w_2$ with $w_1$ of shape $[2]$ and
$w_2$ of shape $[3]$, both finite --- every wire reads, but the
elementwise sum refuses the incompatible shapes, $t$ never evaluates, and
$\vdash t = t$ is again false everywhere. In both cases the un-gated rule
would assert an invalid sequent under a satisfiable (indeed empty)
premise --- in the first case the conclusion fails at the exhibited
valuation (a finite assignment would satisfy it); in the second it is
unsatisfiable outright. The typing pack is exactly the semantic price of
equality-by-evaluation; it is the same field family that repairs
\rulename{Norm}, there compounded with canonical-form equality.
\end{bremark}

\begin{bremark}[taxonomy note]\label{rem:eqrefl-tax}
Wherever \rulename{Norm} applies, \rulename{EqRefl} is its diagonal
($t_1 = t_2 = t$, canonical forms trivially equal). It is nonetheless
kept primitive: the development's frozen minimal fragment (\texttt{v0} in
the sources) excludes the polynomial engine outright, and the equational
kit must stand without it.
\end{bremark}

\begin{anchors}
Lemma \ref{lem:eqrefl-total} & \texttt{Term.evalQ\_defined\_of\_bshape\_fin}; payload fields \texttt{tyDef}, \texttt{fin}\\
finite-wire reads succeed & \texttt{Hypernet.Valuation.wireQ\_isSome}\\
Lemma \ref{lem:eqrefl-lift} & \texttt{QTensor.rel2\_refl}\\
Theorem \ref{thm:eqrefl-sound} & \texttt{sound\_eqRefl}; wired as core \texttt{RuleData.eqRefl} into \texttt{ProofObject.sound}\\
typing predicates & \texttt{Term.bshape}, \texttt{Term.finWires} (Core/Catalogue)\\
\end{anchors}


\rulesection{10}{EqSym}{core rule}
\label{blcatrule-EqSym}
\ruleblock{The rule (Steele metanotation)}
\sourcenote{as realized in the mechanized development; source row: \texttt{\texttt{rules.toml}} --- \emph{EqSym}. The metanotation statement is typeset by hand for notational quality; its mechanized ground truth is the Lean statement below, extracted from the sources at every build of this book and hash-bound. The rule \emph{roster} is machine-checked (every mechanized rule tag round-trips to exactly one catalogue row); the statement \emph{text} is editorial, checked by eye against the extracted Lean beside it.}
\begin{mathpar}
\inferrule*[Right=\rulename{EqSym}]{\Phi \vdash t = u}{\Phi \vdash u = t}
\end{mathpar}

\par\medskip
\rulename{EqSym} flips an equality atom: from $\Phi \vdash t = u$
conclude $\Phi \vdash u = t$. The need is real because downstream
consumers of equalities match atoms \emph{syntactically} with a fixed
orientation --- \rulename{EqElim}\textquotesingle s equality premise, in
particular, names its two terms in order --- while only the combination
rules expose both polarities of an equality in their own matching. A
producer whose congruence engine emits $t = u$ but whose rewrite needs
$u = t$ inserts this one node.

What the checker recomputes: three syntactic form equations (premise
atom, conclusion atom, shared context) --- nothing else. The rule
carries \emph{no} typing fields, and that absence is principled twice
over: the premise\textquotesingle s holding atom already witnesses both
evaluations (nothing must be manufactured, contrast
\rulename{EqRefl}), and the elementwise-under-broadcast equality
relation is symmetric branch by branch, so no shape agreement is needed
(contrast \rulename{EqElim}, whose typing condition guards transport
into a \emph{different} formula --- exactly what a flip never does).

Every payload field is structural; there is no droppable hypothesis to
exhibit a counterexample against. The rule is core rather than
admissible because its would-be unpack (through
\rulename{EqRefl} + \rulename{EqElim}) demands typing side data the
premise atom does not syntactically carry --- the unpack is not
computable from the node it would replace.

\ruleblock{Lean statement}
\declhead{BernsteinLogic.Soundness}{QTensor.rel2\_symm}{\badge{BadgeGreen}{statement bound}}
\begin{leansrc}
/-- `rel2` of a symmetric relation is symmetric (all three broadcast
branches flip pointwise). -/
theorem QTensor.rel2_symm {r : ℚ → ℚ → Prop}
    (hr : ∀ x y, r x y → r y x) {a b : QTensor}
    (h : QTensor.rel2 r a b) : QTensor.rel2 r b a
\end{leansrc}
\declhead{BernsteinLogic.Soundness}{sound\_eqSym}{\badge{BadgeGreen}{statement bound}}
\begin{leansrc}
/-- EqSym: flip the premise atom's two evaluations and the symmetric
`rel2`. -/
theorem sound_eqSym {H : Hypernet} {seq premise : Sequent} {tT tU : Term}
    (ih : premise.semanticallyValid H)
    (hpre : seq.pre = premise.pre)
    (hpremPost : premise.post = .atom .eq tT tU)
    (hpost : seq.post = .atom .eq tU tT) :
    seq.semanticallyValid H
\end{leansrc}
\declhead{BernsteinLogic.Core.Catalogue}{PEqSymData}{\badge{BadgeGreen}{definition bound}}
\begin{leansrc}
inductive PEqSymData : Hypernet → Sequent → Type where
  | mk (H : Hypernet) (seq : Sequent)
       (premise : Sequent) (premiseProof : ProofObject H premise)
       (termT termU : Term)
       (preEq       : seq.pre = premise.pre)
       (premisePost : premise.post = .atom .eq termT termU)
       (postEq      : seq.post = .atom .eq termU termT)
       : PEqSymData H seq
\end{leansrc}
\ruleblock{Correspondence (Steele display $\leftrightarrow$ Lean)}
\badge{BadgeGreen}{reviewer-certified against the displayed texts}
\sourcenote{The dictionary below is a reviewer-certified statement of how the mechanized statements above interpret the displayed rule: display metavariable $\mapsto$ Lean binder, each display hypothesis $\mapsto$ the premise or side condition carrying it, the display conclusion $\mapsto$ the goal --- with any deliberate difference in width between display and mechanization named outright. Certification is recorded as a hash triple over (the displayed cell, the extracted statement texts, this dictionary\textquotesingle s own text); if any of the three changes, this block turns amber until a reviewer re-certifies. The correspondence claim itself remains editorial --- reviewer-grade, not machine-checked; its machine-checked successor types each soundness statement against the rule as a structured object, and supersedes this block per rule as it lands.}
Judgment frame as in \rulename{Ax}; sub-derivation premise as at
\rulename{Weaken}.
\begin{itemize}
\item $t, u \mapsto$ \texttt{tT tU : Term}; the premise judgment $\Phi
  \vdash t = u$ $\mapsto$ \texttt{ih} with the shape tie
  \texttt{hpremPost} $:$ \texttt{premise.post = .atom .eq tT tU};
  $\Phi \mapsto$ \texttt{premise.pre}, shared via \texttt{hpre}.
\item conclusion $\Phi \vdash u = t$ $\mapsto$ the goal under
  \texttt{hpost} $:$ \texttt{seq.post = .atom .eq tU tT} --- the same
  two term binders, transposed; nothing else changes.
\item certified alongside: \texttt{QTensor.rel2\_symm}, the branchwise
  symmetry of the elementwise comparison lifting --- the reason no
  typing field exists (the premise\textquotesingle s holding atom
  already witnesses both evaluations, as the commentary explains).
\end{itemize}
No width difference in either direction; every hypothesis is consumed.

\ruleblock{Lean verifier receipts}
\sourcenote{module \texttt{BernsteinLogic.Soundness}; bound build \texttt{a20d90e66c13\ldots} (bound from receipt a9c0e34d3e3c9c8c, 2026-08-03T21:30:36Z). A receipt records that named checking layers accepted a build: (module, compiled-artifact SHA-256, environment SHA-256, checker layers), minted when the checking gate ran; the store receipts attesting the bound build are listed with the layers each run credited. Nothing here asks the reader to trust the store: every layer is an independently runnable program over the quoted sources --- rebuild the module and re-run the checkers to reproduce the verdict. Receipts attesting superseded builds credit no layer and are omitted from this rendering (the count of omissions, if any, is noted below).}
\begin{receipttable}
\receiptrow{a9c0e34d3e3c9c8c}{2026-08-03T21:30:36Z}{build+decorrelated}{\badge{BadgeGreen}{current}}
\end{receipttable}
\sourcenote{21 further receipt(s) in the store attest superseded builds of this module and are omitted from this rendering.}
\sourcenote{module \texttt{BernsteinLogic.Core.Catalogue}; bound build \texttt{cbc92e976a1d\ldots} (bound from receipt 5abdb3589707b5df, 2026-08-03T21:30:33Z). A receipt records that named checking layers accepted a build: (module, compiled-artifact SHA-256, environment SHA-256, checker layers), minted when the checking gate ran; the store receipts attesting the bound build are listed with the layers each run credited. Nothing here asks the reader to trust the store: every layer is an independently runnable program over the quoted sources --- rebuild the module and re-run the checkers to reproduce the verdict. Receipts attesting superseded builds credit no layer and are omitted from this rendering (the count of omissions, if any, is noted below).}
\begin{receipttable}
\receiptrow{24804e66edba0bc9}{2026-08-03T20:25:09Z}{build+decorrelated}{\badge{BadgeGreen}{current}}
\receiptrow{5abdb3589707b5df}{2026-08-03T21:30:33Z}{build+decorrelated}{\badge{BadgeGreen}{current}}
\end{receipttable}
\sourcenote{18 further receipt(s) in the store attest superseded builds of this module and are omitted from this rendering.}
\layerline{\layeritem{build}{lake build (elaboration)}{\badge{BadgeGreen}{current}}\layeritem{replay}{leanchecker (official kernel replay)}{\badge{BadgeAmber}{stale only}}\layeritem{hardened}{safe\_verify (hardened screen)}{\badge{BadgeAmber}{stale only}}\layeritem{decorrelated}{lean4export + nanoda (decorrelated kernel)}{\badge{BadgeGreen}{current}}}
\sourcenote{\badge{BadgeGray}{comparator: none} no fixture in the three-implementation comparator harness (compiled / kernel-interpreted / reference-interpreted) exercises this rule.}
\ruleblock{HOL4 statement}
\sourcenote{theory \texttt{bern\_fam\_resid}, HOL4 development script \texttt{\texttt{bern\_fam\_residScript.sml}} \badge{BadgeGreen}{script bound}}
\holhead{Definition checkFam\_eqsym\_def}
\begin{holsrc}
Definition checkFam_eqsym_def:
  checkFam_eqsym (cols:columns) (ct:ctx) (i:num) ⇔
    case frSeq cols ct i of
    | NONE => F
    | SOME sq =>
      case frChildSeq cols ct i 0 of
      | NONE => F
      | SOME ps =>
        case frTerm cols ct i 0 of
        | NONE => F
        | SOME tT =>
          case frTerm cols ct i 1 of
          | NONE => F
          | SOME tU =>
              sq.pre = ps.pre ∧
              ps.post = FAtom BowEq tT tU ∧
              sq.post = FAtom BowEq tU tT
End
\end{holsrc}
\holhead{Theorem famEq\_eqsym}
\begin{holsrc}
Theorem famEq_eqsym:
  wfColumns cols ct ∧ mkCtx cols = SOME ct ∧
  parseCert (cols.hdrImage ++ rest) = SOME (ws, pf) ∧
  nodeCorr cols ct i (CertPf sq nd) ∧
  FST (THE (tagOf cols i)) = 24 ⇒
  (checkFam_eqsym cols ct i ⇔ armEqSym ws sq nd)
\end{holsrc}
\ruleblock{HOL verifier receipts}
\badge{BadgeBlue}{documentary} Documentary citations. Family validation: \texttt{bern\_fam\_resid} built GREEN with all ten \texttt{famEq} arms validated, \texttt{hol-strict} clean; spine census: \texttt{bern\_joins} GREEN under \texttt{hol-strict} (55 theorems, 0 offences, 208-theory ancestry including this theory), \texttt{checkJoins\_eq\_checkPf} verified verbatim in the built signature. The cited \texttt{famEq} lemma is the kernel-checked per-node equivalence: under well-formed columns (\texttt{wfColumns}/\texttt{mkCtx}/\texttt{nodeCorr}) at this rule\textquotesingle s wire tag, the columnar family check coincides with the flat reference walk\textquotesingle s arm for this rule. There is no signed HOL receipt store yet --- these citations are documentary.
Semantic stratum, also ported: the soundness twin \texttt{sound\_eqSym}
(theory \texttt{bern\_tr\_b14}: \texttt{hol-strict} clean, 2 theorems, 0
offences, verification run of 2026-07-22 on an isolated build host,
content-addressed at the same sources), consumed by the certified-
fragment keystone (\texttt{checkPf\_sound\_frag}, theory
\texttt{bern\_tr\_splice}, 20 theorems, 0 offences, same run): a
reference-checker accept inside the proven wire fragment yields sequent
validity.

\sourcenote{HOL receipts are documentary: quoted results of \texttt{hol-strict} runs
and census rows from the HOL4 development\textquotesingle s build
evidence, not minted receipt artefacts. A HOL receipt store (the Lean
store\textquotesingle s twin) is named future work.
}
\ruleblock{Proof}

\sourcenote{Genre: core rule (tier 1). Setting as fixed in the
\rulename{EqElim} section.}

\textbf{Claim} (\texttt{sound\_eqSym}).\enspace The payload supplies
terms $t, u$; the premise concludes the atom $t = u$ over premise formula
$\Phi$, shared with the conclusion sequent; the conclusion is the atom
$u = t$. If the premise is semantically valid, so is
$\Phi \vdash u = t$.

\begin{blemma}[the elementwise relation of a symmetric relation is
symmetric]\label{lem:eqsym-rel2}
Let $r$ be a symmetric binary relation on $\Q$, and suppose tensors $A$,
$B$ stand in the elementwise-under-broadcast lift of $r$. Then $B$, $A$
stand in it too.
\end{blemma}
\begin{proof}
The lift has three live branches, and each is separately symmetric ---
this per-branch analysis is the entire content, so we spell it out.
\emph{Equal shapes}: the relation reads the zipped entry lists; the
swapped claim reads the zip in the other order, and each swapped pair is
closed by symmetry of $r$. \emph{Rank-0 $A$ against positive-rank $B$}:
the hypothesis says $r(A_{\mathrm{scalar}}, y)$ for every entry $y$ of
$B$; the swapped claim is judged in the \emph{other} scalar-broadcast
branch (rank-0 on the right), which demands $r(y, A_{\mathrm{scalar}})$
for every $y$ --- symmetry of $r$, pointwise. \emph{Rank-0 $B$}: mirror
image. Two unequal positive-rank shapes never stand in the lift, so
there is nothing to prove there. (\texttt{QTensor.rel2\_symm}.)
\end{proof}

\begin{bthm}[soundness of \rulename{EqSym}]\label{thm:eqsym-sound}
If $\Phi \vdash t = u$ is semantically valid, so is $\Phi \vdash u = t$.
\end{bthm}
\begin{proof}
Fix an admissible $\sigma \models \Phi$. The premise yields evaluations
$A$ of $t$ and $B$ of $u$ in the elementwise equality lift. Present the
same two evaluations in the opposite roles: $u$ evaluates to $B$, $t$ to
$A$, and Lemma~\ref{lem:eqsym-rel2} (with $r$ the symmetric relation
$=$) closes the atom $u = t$.
\end{proof}

\begin{bremark}[the negative space: what this rule does \emph{not} need]
\label{rem:eqsym-neg}
The rule has no droppable side condition --- every payload field is
structural (the three form equations) --- so the honest remarks here
concern the two fields conspicuously absent. No typing pack
(\rulename{EqRefl}'s \texttt{tyDef}/\texttt{fin}): the premise's holding
atom already \emph{witnesses} both evaluations, and symmetry only re-uses
the premise's own evidence; nothing must be manufactured. No common-shape
condition (\rulename{EqElim}'s \texttt{tyEq}): Lemma~\ref{lem:eqsym-rel2}
is shape-agnostic --- a broadcast equality flips to the mirror broadcast
equality, same content, no interchangeability claim ever made. The
contrast triangulates \rulename{EqElim}'s repair: the typing condition
there guards \emph{transport into a different formula}, which is exactly
what \rulename{EqSym} never does.
\end{bremark}

\begin{bremark}[why a primitive, when a two-node unpack exists on paper]
\label{rem:eqsym-adm}
From $\Phi \vdash t = u$ one can derive $\Phi \vdash u = t$ through
\rulename{EqRefl} + \rulename{EqElim}: body $\Psi := (a = t)$ with the
hole wire $a$ fresh for $t$ --- substitution replaces \emph{every}
occurrence of $a$, and freshness is what makes $\Psi\sub{a}{t}$ the atom
$t = t$ and $\Psi\sub{a}{u}$ the atom $u = t$ --- rewriting
the left occurrence of $t$ in $t = t$ to $u$. But the derivation demands
\rulename{EqRefl}'s typing pack on $t$ \emph{and} \rulename{EqElim}'s
common-shape condition on $(t, u)$ --- side data the premise atom does
not syntactically carry, and which a producer cannot always reconstruct
at the flip site. The primitive constructor instead consumes the
premise's semantic evidence directly (the proof above), making the flip
unconditional. In the catalogue's three-tier taxonomy that obstruction is
what sorts \rulename{EqSym} core rather than admissible: the unpack is
not payload-computable from the node it would replace.
\end{bremark}

\begin{anchors}
Lemma \ref{lem:eqsym-rel2} & \texttt{QTensor.rel2\_symm}\\
Theorem \ref{thm:eqsym-sound} & \texttt{sound\_eqSym}; wired as core \texttt{RuleData.eqSym} into \texttt{ProofObject.sound}\\
form equations & payload fields \texttt{preEq}, \texttt{premisePost}, \texttt{postEq} (\texttt{PEqSymData})\\
\end{anchors}


\rulesection{11}{Antisym}{core rule}
\label{blcatrule-Antisym}
\ruleblock{The rule (Steele metanotation)}
\sourcenote{as realized in the mechanized development; source row: \texttt{\texttt{rules.toml}} --- \emph{Antisym}. The metanotation statement is typeset by hand for notational quality; its mechanized ground truth is the Lean statement below, extracted from the sources at every build of this book and hash-bound. The rule \emph{roster} is machine-checked (every mechanized rule tag round-trips to exactly one catalogue row); the statement \emph{text} is editorial, checked by eye against the extracted Lean beside it.}
\begin{mathpar}
\inferrule*[Right=\rulename{Antisym}]{\Phi \vdash t \leq u \\ \Phi \vdash t \geq u}{\Phi \vdash t = u}
\end{mathpar}

\par\medskip
\rulename{Antisym} squeezes an equality out of two opposite bounds: from
$\Phi \vdash t \leq u$ and $\Phi \vdash t \geq u$ --- the same two terms,
the same context --- conclude $\Phi \vdash t = u$. A producer reaches for
it wherever a pinching argument closes: an optimization pass that
certifies matching upper and lower bounds on the same quantity, a
fixed-point equality assembled from two directional inequalities, or any
pipeline whose native outputs are order facts but whose consumer
(\rulename{EqElim}\textquotesingle s rewrite chains, in particular) wants
an equality atom.

The rule earns its place in the core by a non-derivability fact, not by
convenience: equality-from-inequalities has no derivation in the
remainder of the calculus. Provably --- an induction over derivations
shows that every equality concluded from an equality-atom-free context
equates terms with identical polynomial normal forms --- so no chain of
existing rules can ever conclude $t = u$ from the two bounds when the
sides are genuinely different terms. Historically the calculus carried
this inference only inside two bespoke rules for the ReLU function,
which were thereby load-bearing for something much more general than
ReLU; with antisymmetry available as its own rule, that monopoly ends.

The side conditions are purely structural: two premise-atom form
equations, the conclusion-atom form equation, and the two context
equations tying each premise sequent to the conclusion\textquotesingle s
context. No typing fields --- and the absence is principled (contrast
\rulename{EqRefl}, which must \emph{manufacture} an evaluation and so
carries the typing pack): each premise atom holds only where both sides
already evaluate finitely at compatible shapes, and that is exactly the
evaluation evidence the equality atom\textquotesingle s semantics
consumes. Dropping either premise leaves a one-sided bound, and any
strictly-below evaluation pair refutes the un-gated rule.

One reading discipline, stated here because the conclusion invites
over-reading: the concluded equality holds in the calculus\textquotesingle
s elementwise-under-broadcast sense, which suffices for the truth of the
atom and for every consumer that reads it semantically --- but it is not
a license to interchange $t$ and $u$ inside an arbitrary formula.
Substitution rides \rulename{EqElim}, whose common-defined-shape typing
condition remains in force unchanged (that section constructs the
refutation its absence would admit).

\cataloguenote{This rule is \emph{staged, not enacted}: what exists today is the
statement register --- the soundness theorem below, proven in Lean with
a HOL twin --- while the certificate format knows nothing of it yet. No
payload constructor, no wire tag, no checker arm: a certificate cannot
currently carry an \rulename{Antisym} node, and this section will gain
its checker-side anchors when enactment lands. The mathematics is
settled; the plumbing is queued.
}
\ruleblock{Lean statement}
\declhead{BernsteinLogic.Soundness}{QTensor.rel2\_conj}{\badge{BadgeGreen}{statement bound}}
\begin{leansrc}
/-- Two `rel2` facts at the SAME tensor pair combine pointwise: the
broadcast branch taken is a function of the two shapes alone, so the
quantified domains coincide and the pointwise conjunction lifts.
(Stated for `sound_antisym`; any pointwise combination rides it.) -/
theorem QTensor.rel2_conj {r s t : ℚ → ℚ → Prop}
    (hrst : ∀ x y, r x y → s x y → t x y) {a b : QTensor}
    (hr : QTensor.rel2 r a b) (hs : QTensor.rel2 s a b) :
    QTensor.rel2 t a b
\end{leansrc}
\declhead{BernsteinLogic.Soundness}{sound\_antisym}{\badge{BadgeGreen}{statement bound}}
\begin{leansrc}
/-- Antisym: `≤` and `≥` premise validities at the same term pair force
the equality atom — entrywise `le_antisymm` lifted through `rel2`
(`QTensor.rel2_conj`: both premises constrain the same evaluated pair,
so the broadcast branches align). Like `sound_eqSym` and unlike
`sound_eqRefl`, no B13 typing pack is needed: the premise atoms only
hold where both sides already evaluate finitely, which is exactly the
evaluation evidence the conclusion atom needs. Core rather than
derived: the equality-from-inequalities inference has no catalogue
derivation — provably, since derivable equalities from eq-free contexts
have canonically-equal sides — so the rule must be primitive
(ReLU-Pos/Neg had been core solely as carriers of this inference). -/
theorem sound_antisym {H : Hypernet} {seq premiseLe premiseGe : Sequent}
    {tT tU : Term}
    (ihLe : premiseLe.semanticallyValid H)
    (ihGe : premiseGe.semanticallyValid H)
    (hpreLe : seq.pre = premiseLe.pre)
    (hpreGe : seq.pre = premiseGe.pre)
    (hpremLePost : premiseLe.post = .atom .le tT tU)
    (hpremGePost : premiseGe.post = .atom .ge tT tU)
    (hpost : seq.post = .atom .eq tT tU) :
    seq.semanticallyValid H
\end{leansrc}
\ruleblock{Correspondence (Steele display $\leftrightarrow$ Lean)}
\badge{BadgeGreen}{reviewer-certified against the displayed texts}
\sourcenote{The dictionary below is a reviewer-certified statement of how the mechanized statements above interpret the displayed rule: display metavariable $\mapsto$ Lean binder, each display hypothesis $\mapsto$ the premise or side condition carrying it, the display conclusion $\mapsto$ the goal --- with any deliberate difference in width between display and mechanization named outright. Certification is recorded as a hash triple over (the displayed cell, the extracted statement texts, this dictionary\textquotesingle s own text); if any of the three changes, this block turns amber until a reviewer re-certifies. The correspondence claim itself remains editorial --- reviewer-grade, not machine-checked; its machine-checked successor types each soundness statement against the rule as a structured object, and supersedes this block per rule as it lands.}
Judgment frame as in \rulename{Ax}; sub-derivation premises as at
\rulename{Weaken}.
\begin{itemize}
\item $t, u \mapsto$ \texttt{tT tU : Term}; the two premise judgments
  $\Phi \vdash t \leq u$ and $\Phi \vdash t \geq u$ $\mapsto$
  \texttt{ihLe}, \texttt{ihGe} with the shape ties \texttt{hpremLePost}
  $:$ \texttt{premiseLe.post = .atom .le tT tU} and \texttt{hpremGePost}
  $:$ \texttt{premiseGe.post = .atom .ge tT tU} --- note the $\geq$
  premise binds its own atom former at the \emph{same} binder order
  (\texttt{tT}, \texttt{tU}), not a flipped $\leq$ atom.
\item the shared context: the display writes one $\Phi$ across all
  three sequents; the statement renders it as \emph{two} equations,
  \texttt{hpreLe} $:$ \texttt{seq.pre = premiseLe.pre} and
  \texttt{hpreGe} $:$ \texttt{seq.pre = premiseGe.pre} --- each premise
  tied to the conclusion\textquotesingle s context, premise-to-premise
  equality following by transitivity. Same content, one equation more
  than the display suggests.
\item conclusion $\Phi \vdash t = u$ $\mapsto$ the goal under
  \texttt{hpost} $:$ \texttt{seq.post = .atom .eq tT tU}.
\item no typing hypotheses appear on either side --- principled, as the
  commentary explains (the premise atoms carry the evaluation evidence;
  contrast \rulename{EqRefl}\textquotesingle s \texttt{htyDef}/\texttt{hfin}).
\item certified alongside as the semantic engine:
  \texttt{QTensor.rel2\_conj} --- two elementwise-under-broadcast facts
  at the \emph{same} tensor pair combine pointwise, because the
  broadcast branch taken is a function of the two shapes alone, so the
  quantified domains coincide. The soundness proof feeds it entrywise
  antisymmetry of the rational order.
\end{itemize}
Width note, both directions: every hypothesis is consumed, and the one
place statement and display could be read apart is the conclusion
atom\textquotesingle s broadcast-sense semantics --- the proven equality
is the elementwise-under-broadcast lift (a scalar can stand equal to a
vector whose every entry it matches), which is the equality
atom\textquotesingle s meaning everywhere in this calculus; the
commentary\textquotesingle s reading discipline names what that does
\emph{not} license.

\ruleblock{Lean verifier receipts}
\sourcenote{module \texttt{BernsteinLogic.Soundness}; bound build \texttt{a20d90e66c13\ldots} (bound from receipt a9c0e34d3e3c9c8c, 2026-08-03T21:30:36Z). A receipt records that named checking layers accepted a build: (module, compiled-artifact SHA-256, environment SHA-256, checker layers), minted when the checking gate ran; the store receipts attesting the bound build are listed with the layers each run credited. Nothing here asks the reader to trust the store: every layer is an independently runnable program over the quoted sources --- rebuild the module and re-run the checkers to reproduce the verdict. Receipts attesting superseded builds credit no layer and are omitted from this rendering (the count of omissions, if any, is noted below).}
\begin{receipttable}
\receiptrow{a9c0e34d3e3c9c8c}{2026-08-03T21:30:36Z}{build+decorrelated}{\badge{BadgeGreen}{current}}
\end{receipttable}
\sourcenote{21 further receipt(s) in the store attest superseded builds of this module and are omitted from this rendering.}
\layerline{\layeritem{build}{lake build (elaboration)}{\badge{BadgeGreen}{current}}\layeritem{replay}{leanchecker (official kernel replay)}{\badge{BadgeAmber}{stale only}}\layeritem{hardened}{safe\_verify (hardened screen)}{\badge{BadgeAmber}{stale only}}\layeritem{decorrelated}{lean4export + nanoda (decorrelated kernel)}{\badge{BadgeGreen}{current}}}
\sourcenote{\badge{BadgeGray}{comparator: none} no fixture in the three-implementation comparator harness (compiled / kernel-interpreted / reference-interpreted) exercises this rule.}
\ruleblock{HOL4 statement}
\sourcenote{theory \texttt{bern\_tr\_antisym}, HOL4 development script \texttt{\texttt{bern\_tr\_antisymScript.sml}} \badge{BadgeGreen}{script bound}}
\holhead{Theorem qtRel2\_conj}
\begin{holsrc}
Theorem qtRel2_conj:
  ∀r s t a b.
    (∀x y. r x y ∧ s x y ⇒ t x y) ∧
    qtRel2 r a b ∧ qtRel2 s a b ⇒
    qtRel2 t a b
\end{holsrc}
\holhead{Theorem sound\_antisym}
\begin{holsrc}
Theorem sound_antisym:
  ∀H sq premiseLe premiseGe tT tU.
    seqSemanticallyValid H premiseLe ∧
    seqSemanticallyValid H premiseGe ∧
    sq.pre = premiseLe.pre ∧
    sq.pre = premiseGe.pre ∧
    premiseLe.post = FAtom BowLe tT tU ∧
    premiseGe.post = FAtom BowGe tT tU ∧
    sq.post = FAtom BowEq tT tU ⇒
    seqSemanticallyValid H sq
\end{holsrc}
\ruleblock{HOL verifier receipts}
\badge{BadgeBlue}{documentary} Documentary citations, one register-honest caveat: this theory is the
staged rule\textquotesingle s whole standalone footprint (both theorems
in one theory, deliberately, so the staged work moves as a unit), and
--- unlike the enacted siblings on this page --- no isolated-build-host
verification census is attached to it yet; the citation attests the
script text, not a checked build artifact. Statement fidelity follows
the ported-soundness register used throughout: the Lean statement\textquotesingle
s implicit arguments become leading universals and its curried named
hypotheses a conjunction in source order; \texttt{.atom .le/.ge/.eq}
renders as \texttt{FAtom BowLe/BowGe/BowEq}; the pointwise antisymmetry
witness closes by the rational order\textquotesingle s antisymmetry
lemma. Because the rule has no checker arm, there is no family-check or
certified-fragment claim to make here --- the semantic stratum is the
entire port, matching the statement-register staging the note describes.

\sourcenote{HOL receipts are documentary: quoted results of \texttt{hol-strict} runs
and census rows from the HOL4 development\textquotesingle s build
evidence, not minted receipt artefacts. A HOL receipt store (the Lean
store\textquotesingle s twin) is named future work.
}
\ruleblock{Proof}

\sourcenote{Genre: core rule (tier 1), staged --- statement register
only; see the section note. Setting as fixed in the \rulename{EqElim}
section.}

\textbf{Claim} (\texttt{sound\_antisym}).\enspace Two premise sequents
conclude the atoms $t \leq u$ and $t \geq u$, each over the premise
formula $\Phi$ shared with the conclusion sequent; the conclusion is the
atom $t = u$. If both premises are semantically valid, so is
$\Phi \vdash t = u$.

\begin{blemma}[elementwise lifts at a shared tensor pair combine
pointwise]\label{lem:antisym-rel2conj}
Let $r$, $s$, $w$ be binary relations on $\Q$ with
$r(x,y) \wedge s(x,y) \Rightarrow w(x,y)$ for all $x, y$, and suppose
tensors $A$, $B$ stand in the elementwise-under-broadcast lift of $r$
and also in the lift of $s$. Then they stand in the lift of $w$.
\end{blemma}
\begin{proof}
The branch the lift takes is a function of the two shapes alone --- not
of the relation --- so the two hypotheses and the goal all quantify over
the \emph{same} domain, and the pointwise implication closes each point.
Spelled per branch: \emph{equal shapes}: all three statements read the
same zipped entry lists, pair by pair. \emph{Rank-0 $A$ against
positive-rank $B$}: all three quantify $A$'s scalar against each entry
of $B$; mirror image for rank-0 $B$. \emph{Two unequal positive-rank
shapes}: no lift holds there at all, so the hypothesis for $r$ is
already false and there is nothing to prove.
(\texttt{QTensor.rel2\_conj}.)
\end{proof}

\begin{bthm}[soundness of \rulename{Antisym}]\label{thm:antisym-sound}
If $\Phi \vdash t \leq u$ and $\Phi \vdash t \geq u$ are semantically
valid, so is $\Phi \vdash t = u$.
\end{bthm}
\begin{proof}
Fix an admissible $\sigma \models \Phi$. The $\leq$ premise yields
evaluations $A$ of $t$ and $B$ of $u$ in the elementwise
$\leq$-lift; the $\geq$ premise yields evaluations $A'$ of $t$ and $B'$
of $u$ in the elementwise $\geq$-lift. Evaluation is a function of the
valuation and the term, so $A' = A$ and $B' = B$: both premises
constrain the \emph{same} evaluated pair, which is what makes
Lemma~\ref{lem:antisym-rel2conj} applicable at all. Apply it with
$r \mathbin{:=} {\leq}$, $s \mathbin{:=} {\geq}$,
$w \mathbin{:=} {=}$ and the pointwise fact
$x \leq y \wedge x \geq y \Rightarrow x = y$ --- antisymmetry of the
order on $\Q$ --- and the equality lift closes the atom $t = u$.
\end{proof}

\begin{bremark}[the negative space: no typing pack, and why both
premises]\label{rem:antisym-neg}
No typing fields (\rulename{EqRefl}'s \texttt{tyDef}/\texttt{fin}):
that rule must \emph{manufacture} an evaluation for its premise-free
conclusion, and pays for totality with static typing conditions; here
each premise atom holds only where both sides already evaluate finitely
at compatible shapes, and the proof above consumes exactly that
evidence --- nothing is manufactured (\rulename{EqSym} makes the same
economy for the same reason). Both orientations are load-bearing:
dropping either premise leaves a one-sided bound, and any evaluation
pair with $A$ strictly below $B$ entrywise satisfies the surviving
premise while refuting the conclusion --- the counterexample is
immediate at scalar shape already.
\end{bremark}

\begin{bremark}[what the broadcast-sense conclusion does not
license]\label{rem:antisym-broadcast}
The concluded atom holds in the elementwise-under-broadcast sense: a
scalar stands equal to a vector whose every entry it matches. That is
the equality atom's meaning everywhere in this calculus --- sufficient
for atom truth and for every consumer that reads the atom semantically
--- but it is not a license to interchange $t$ and $u$ inside an
arbitrary formula. Substitution rides \rulename{EqElim}, whose
common-defined-shape typing condition remains in force unchanged: that
section constructs a satisfiable-to-unsatisfiable flip from exactly the
scalar-equals-vector equalities this rule can legitimately conclude.
\end{bremark}

\begin{bremark}[why core, and what is staged]\label{rem:antisym-core}
Core because the inference is not derivable: an induction over
derivations shows every equality concluded from an equality-atom-free
context equates terms with identical polynomial normal forms, so no
chain of existing rules concludes $t = u$ from the two bounds when the
sides are genuinely different terms. Historically the calculus carried
antisymmetry only inside two bespoke rules for the ReLU function; a
core \rulename{Antisym} ends that monopoly and lets bespoke rules be
judged on their own content. Staged: the theorem above is the
statement register --- no payload constructor, wire tag, or checker arm
exists yet, so certificates cannot carry the rule today; the anchors
table below will gain its ``wired as core'' line at enactment.
\end{bremark}

\begin{anchors}
Lemma \ref{lem:antisym-rel2conj} & \texttt{QTensor.rel2\_conj}\\
Theorem \ref{thm:antisym-sound} & \texttt{sound\_antisym}; statement register only --- not yet wired (no \texttt{ProofObject} constructor exists)\\
HOL twin & theory \texttt{bern\_tr\_antisym}: \texttt{qtRel2\_conj}, \texttt{sound\_antisym}\\
\end{anchors}


\rulesection{12}{ConvElim}{core rule}
\label{blcatrule-ConvElim}
\ruleblock{The rule (Steele metanotation)}
\sourcenote{as realized in the mechanized development; source row: \texttt{\texttt{rules.toml}} --- \emph{ConvElim}. The metanotation statement is typeset by hand for notational quality; its mechanized ground truth is the Lean statement below, extracted from the sources at every build of this book and hash-bound. The rule \emph{roster} is machine-checked (every mechanized rule tag round-trips to exactly one catalogue row); the statement \emph{text} is editorial, checked by eye against the extracted Lean beside it.}
\begin{mathpar}
\inferrule*[Right=\rulename{ConvElim}]{
  \Phi \vdash x \in \mathcal{R}(\Delta, S)
}{
  \Phi \vdash \textstyle\bigwedge_{i = 0}^{n} \mathrm{canon}(\epsilon_i\,\lambda_i(x)) \leq 0 \\ \epsilon_i \,=\, {+1} \text{ if } i \in S, \quad {-1} \text{ if } i \in \bar S
}
\end{mathpar}

\par\medskip
\rulename{ConvElim} opens a region atom: from
$\Phi \vdash x \in \mathcal{R}(\Delta, S)$ it concludes the
region\textquotesingle s defining barycentric halfspace conjunction ---
one $(\cdot) \leq 0$ atom per support vertex, in a pinned canonical
spelling. It is the entry move of every argument that reasons
\emph{linearly} about a geometric premise: \rulename{BernCone} consumes
region atoms natively, but the combination rules (\rulename{Farkas},
\rulename{Syzygy}) match ordinary comparison atoms by exact canonical
polynomial, and this rule is what turns a region fact into atoms they
can see.

What the checker recomputes: the \emph{entire} conclusion, from
$(x, \Delta, S)$ alone --- the two barycentric affine forms of the 1-D
two-vertex support (weights are determined on a simplex, so nothing is
chosen and nothing rides the certificate), the sign-folding that turns
each core vertex\textquotesingle s wanted $\lambda \geq 0$ into a
$(-\lambda) \leq 0$ atom (the all-$\leq$ discipline the downstream
matchers consume), and the strictly positive primitive-integer rescale
that makes the spelling unique --- then demands syntactic equality with
the cited conclusion. No producer text is trusted; the side condition is
one decidable equation. Outside the checkable fragment (ambient
dimension 1, support exactly two vertices at distinct coordinates) the
build refuses and the rule cannot fire --- refusal is trivially sound,
and the distinct-coordinates gate is precisely where a degenerate
simplex\textquotesingle s division by zero is made unrepresentable.

The rule needs no typing or positivity fields: definedness and shape of
the region term are read out of the region atom\textquotesingle s own
semantics. Dropping the premise-form equation would sever the conclusion
from any premise at all; dropping the canonical-spelling pin would not
make the rule unsound but would strand its conclusions --- equivalent
spellings never match downstream --- which is why the spelling is
checked rather than suggested.

\cataloguenote{Geometric reading of extension sets, 1-D: $S = \emptyset$ is the segment
$[V_0, V_1]$; naming one vertex stretches the region to the closed ray
away from the other --- membership with a \emph{nonpositive} weight on
the extended vertex. The region atom is thus the calculus\textquotesingle
s interface between simplex geometry and sign arithmetic, and this rule
with \rulename{ConvIntro} is that interface\textquotesingle s two
directions, exact inverses on the fragment.
}
\ruleblock{Lean statement}
\declhead{BernsteinLogic.Soundness}{barycentricHalfspaceBuild\_spec}{\badge{BadgeGreen}{statement bound}}
\begin{leansrc}
/-- Structure of a successful barycentric-halfspace build (the 1-D case
`baryCoeffs1D` gates): the F4 all-≤ canonical conjunction, one
`f4VertexAtom` per vertex of the filter pair. -/
theorem barycentricHalfspaceBuild_spec {t : Term} {simplex : Simplex}
    {S : ExtensionSet} {F : Formula}
    (h : barycentricHalfspaceBuild t simplex S = some F) :
    S.simplex = simplex ∧ simplex.complex.ambDim = 1 ∧
    ∃ (a b : Fin simplex.complex.nVertices) (V₀ V₁ : ℚ),
      (List.finRange simplex.complex.nVertices).filter
        (fun v => simplex.contains v) = [a, b] ∧
      (∀ j, simplex.complex.coord a j = V₀) ∧
      (∀ j, simplex.complex.coord b j = V₁) ∧
      V₀ ≠ V₁ ∧
      F = .and
        (f4VertexAtom (S.containsNat a.val)
          (V₁ / (V₁ - V₀)) (-1 / (V₁ - V₀)) t)
        (f4VertexAtom (S.containsNat b.val)
          (-V₀ / (V₁ - V₀)) (1 / (V₁ - V₀)) t)
\end{leansrc}
\declhead{BernsteinLogic.Soundness}{f4VertexAtom\_holds\_iff}{\badge{BadgeGreen}{statement bound}}
\begin{leansrc}
/-- The F4 vertex atom holds iff the signed barycentric value is
nonpositive: `λ ≤ 0` for an extended vertex, `0 ≤ λ` for a core one
(the sign folded into the term, the rescale cancelled). -/
theorem f4VertexAtom_holds_iff {H : Hypernet} {σ : H.Valuation} {t : Term}
    {q : QTensor} (hq : Term.evalQ σ t = some q) (hsq : SingletonShape q.shape)
    (inS : Bool) (off slp : ℚ) :
    Formula.holds σ (f4VertexAtom inS off slp t) ↔
      (if inS then off + slp * q.scalar ≤ 0
       else 0 ≤ off + slp * q.scalar)
\end{leansrc}
\declhead{BernsteinLogic.Soundness}{maskRegion\_pair\_elim}{\badge{BadgeGreen}{statement bound}}
\begin{leansrc}
/-- Extraction: a masked region membership over a two-vertex 1-D support
pins the barycentric weights to their rational closed forms, and exports
the sign of each weight according to its extension flag. -/
theorem maskRegion_pair_elim {C : SimpC} {mask ext : Fin C.nVertices → Bool}
    {x : Fin C.ambDim → ℚ} {a b : Fin C.nVertices} {V₀ V₁ : ℚ}
    (hfil : (List.finRange C.nVertices).filter (fun v => mask v) = [a, b])
    (hVa : ∀ j, C.coord a j = V₀) (hVb : ∀ j, C.coord b j = V₁)
    (hne : V₀ ≠ V₁) (j0 : Fin C.ambDim)
    (h : MaskRegion C mask ext x) :
    (ext a = true → (V₁ - x j0) / (V₁ - V₀) ≤ 0) ∧
    (ext a = false → 0 ≤ (V₁ - x j0) / (V₁ - V₀)) ∧
    (ext b = true → (x j0 - V₀) / (V₁ - V₀) ≤ 0) ∧
    (ext b = false → 0 ≤ (x j0 - V₀) / (V₁ - V₀))
\end{leansrc}
\declhead{BernsteinLogic.Soundness}{sound\_convElim}{\badge{BadgeGreen}{statement bound}}
\begin{leansrc}
theorem sound_convElim {H : Hypernet} {seq premise : Sequent}
    {term : Term} {simplex : Simplex} {S : ExtensionSet}
    (ih : premise.semanticallyValid H)
    (hpre : seq.pre = premise.pre)
    (hpremiseForm : premise.post = .inRegion term simplex S)
    (hhalf : barycentricHalfspaceForm term simplex S seq.post) :
    seq.semanticallyValid H
\end{leansrc}
\par\medskip
\textbf{The canonical form, meaning bound by display.}\enspace The
conclusion formula is the value of the build displayed below ---
affine data (\texttt{baryCoeffs1D}), positive primitive-integer rescale
(\texttt{primAffine}), atom layout (\texttt{f4Lhs}), per-vertex
sign-folding (\texttt{f4VertexAtom}), assembly
(\texttt{barycentricHalfspaceBuild}) --- and the payload\textquotesingle
s side condition is the one-line equation
\texttt{barycentricHalfspaceForm}. A skeptical reader should confirm by
eye that every gate failure yields \texttt{none} (refusal), that the
rescale factor is strictly positive on every reachable branch, and that
no valuation is consulted anywhere in the chain.

\declhead{BernsteinLogic.Core.Catalogue}{baryCoeffs1D}{\badge{BadgeGreen}{definition bound}}
\begin{leansrc}
/-- The (vertex, offset, slope) affine data of the two barycentric
coordinate functions of a 1-dimensional simplex: `λ(x) = offset + slope·x`.
Same gates as `baryLambdas1D` (which builds the cell rules' pinned Term
shape from the same coordinates). -/
def baryCoeffs1D (cx : SimpC) (dContains : Fin cx.nVertices → Bool) :
    Option (List (Fin cx.nVertices × Rat × Rat)) :=
  if h : cx.ambDim = 1 then
    let j0 : Fin cx.ambDim := ⟨0, by omega⟩
    match (List.finRange cx.nVertices).filter (fun v => dContains v) with
    | [a, b] =>
      let V₀ := cx.coord a j0
      let V₁ := cx.coord b j0
      if V₀ = V₁ then none
      else
        let den := V₁ - V₀
        some [(a, V₁ / den, -1 / den), (b, -V₀ / den, 1 / den)]
    | _ => none
  else none
\end{leansrc}
\declhead{BernsteinLogic.Core.Catalogue}{primAffine}{\badge{BadgeGreen}{definition bound}}
\begin{leansrc}
/-- Primitive-integer rescale of the affine pair (slope `s`, offset `c`):
the positive multiple `(k·s, k·c)` whose entries are the coprime-integer
vector sympy's `canon_sympy` produces for `s·x + c` (producer finding F4).
`c = 0` degenerates to sign normalization (single-entry primitive vector
is ±1); `(0, 0)` is fixed. Computed rationally — `k·s` value-equals the
integer `n₁/g` because `s.den ∣ L` and `g ∣ n₁`. -/
def primAffine (s c : Rat) : Rat × Rat :=
  if c = 0 then
    (if s < 0 then -1 else if s = 0 then 0 else 1, 0)
  else
    let L : Nat := Nat.lcm s.den c.den
    let g : Nat := Int.gcd (s.num * (L / s.den : Nat)) (c.num * (L / c.den : Nat))
    let k : Rat := (L : Rat) / (g : Rat)
    (k * s, k * c)
\end{leansrc}
\declhead{BernsteinLogic.Core.Catalogue}{f4Lhs}{\badge{BadgeGreen}{definition bound}}
\begin{leansrc}
/-- The F4 atom body for the rescaled pair: degree-sorted monomial shape
`s'·t + c'`, the constant dropped when zero (sympy `_term_ast`). -/
def f4Lhs (s' c' : Rat) (t : Term) : Term :=
  if c' = 0 then .mul (.const s') t
  else .add (.mul (.const s') t) (.const c')
\end{leansrc}
\declhead{BernsteinLogic.Core.Catalogue}{f4VertexAtom}{\badge{BadgeGreen}{definition bound}}
\begin{leansrc}
/-- One vertex's F4 halfspace atom: sign-folded (complement vertex's λ
negated into the term), rescaled, `≤ 0`-shaped. -/
def f4VertexAtom (inS : Bool) (off slp : Rat) (t : Term) : Formula :=
  let sc := if inS then primAffine slp off else primAffine (-slp) (-off)
  .atom .le (f4Lhs sc.1 sc.2 t) (.const 0)
\end{leansrc}
\declhead{BernsteinLogic.Core.Catalogue}{barycentricHalfspaceBuild}{\badge{BadgeGreen}{definition bound}}
\begin{leansrc}
def barycentricHalfspaceBuild (t : Term) (simplex : Simplex)
    (S : ExtensionSet) : Option Formula := do
  guard (S.simplex = simplex)
  let cs ← baryCoeffs1D simplex.complex simplex.contains
  match cs.map (fun (v, off, slp) =>
    f4VertexAtom (S.containsNat v.val) off slp t) with
  | [x, y] => some (.and x y)
  | _ => none
\end{leansrc}
\declhead{BernsteinLogic.Core.Catalogue}{barycentricHalfspaceForm}{\badge{BadgeGreen}{definition bound}}
\begin{leansrc}
def barycentricHalfspaceForm (t : Term) (simplex : Simplex)
    (S : ExtensionSet) (F : Formula) : Prop :=
  barycentricHalfspaceBuild t simplex S = some F
\end{leansrc}
\declhead{BernsteinLogic.Core.Catalogue}{PConvElimData}{\badge{BadgeGreen}{definition bound}}
\begin{leansrc}
inductive PConvElimData : Hypernet → Sequent → Type where
  | mk (H : Hypernet) (seq : Sequent)
       (premise : Sequent) (premiseProof : ProofObject H premise)
       (term : Term) (simplex : Simplex) (S : ExtensionSet)
       (preEq       : seq.pre = premise.pre)
       (premiseForm : premise.post = .inRegion term simplex S)
       (halfspace   : barycentricHalfspaceForm term simplex S seq.post)
       : PConvElimData H seq
\end{leansrc}
\ruleblock{Correspondence (Steele display $\leftrightarrow$ Lean)}
\badge{BadgeGreen}{reviewer-certified against the displayed texts}
\sourcenote{The dictionary below is a reviewer-certified statement of how the mechanized statements above interpret the displayed rule: display metavariable $\mapsto$ Lean binder, each display hypothesis $\mapsto$ the premise or side condition carrying it, the display conclusion $\mapsto$ the goal --- with any deliberate difference in width between display and mechanization named outright. Certification is recorded as a hash triple over (the displayed cell, the extracted statement texts, this dictionary\textquotesingle s own text); if any of the three changes, this block turns amber until a reviewer re-certifies. The correspondence claim itself remains editorial --- reviewer-grade, not machine-checked; its machine-checked successor types each soundness statement against the rule as a structured object, and supersedes this block per rule as it lands.}
Judgment frame as in \rulename{Ax}; sub-derivation premise as at
\rulename{Weaken}.
\begin{itemize}
\item $x \mapsto$ \texttt{term : Term}; $\Delta \mapsto$
  \texttt{simplex : Simplex}; $S \mapsto$ \texttt{S : ExtensionSet};
  the premise judgment $\Phi \vdash x \in \mathcal{R}(\Delta, S)$
  $\mapsto$ \texttt{ih} with the shape tie \texttt{hpremiseForm} $:$
  \texttt{premise.post = .inRegion term simplex S}; $\Phi$ shared via
  \texttt{hpre}.
\item the \emph{entire} displayed conclusion --- the flat conjunction,
  its vertex order, the sign pattern $\epsilon_i$, the affine
  spellings $\lambda_i$ --- $\mapsto$ one hypothesis, \texttt{hhalf}
  $:$ \texttt{barycentricHalfspaceForm term simplex S seq.post}:
  \texttt{seq.post} is literally the value of the verified canonical
  build. The display\textquotesingle s $\lambda_i$ and $\epsilon_i$ are
  not statement binders but \emph{outputs} of
  \texttt{barycentricHalfspaceBuild}, quoted in full below.
\item conclusion judgment $\mapsto$ the goal
  \texttt{seq.semanticallyValid H}; certified alongside,
  \texttt{barycentricHalfspaceBuild\_spec} unpacks a successful build
  run into the geometric data the soundness proof consumes.
\end{itemize}
Width note, narrower and loud about it: a successful build run forces
ambient dimension $1$ with support exactly two vertices at distinct
coordinates (the first conjuncts of
\texttt{barycentricHalfspaceBuild\_spec}), so the mechanized statement
covers exactly the one-dimensional instance of the displayed
general-dimension rule. Outside that fragment the build returns
\texttt{none}, the hypothesis is unsatisfiable, and the rule cannot
fire --- refusal, never unverified acceptance; the displayed rule at
higher dimension is design intent, not yet certified statement.

Spelling alignment, formerly a live discrepancy: the displayed
conclusion and the certified build now agree --- both state the flat
sign-folded all-$\leq$ conjunction, each atom
$\mathrm{canon}(\epsilon_i\,\lambda_i(x)) \leq 0$ with the
extension/core orientation carried by the sign $\epsilon_i$
(\texttt{f4VertexAtom} negates the core vertex\textquotesingle s
affine form into the term and compares $\leq 0$ --- the commentary
describes this discipline). Earlier printings displayed a mixed
$\leq$/$\geq$ per-vertex comparator pattern; the roster cell has
since been reconciled to the realized form, and this certification
binds the reconciled texts.

\ruleblock{Lean verifier receipts}
\sourcenote{module \texttt{BernsteinLogic.Soundness}; bound build \texttt{a20d90e66c13\ldots} (bound from receipt a9c0e34d3e3c9c8c, 2026-08-03T21:30:36Z). A receipt records that named checking layers accepted a build: (module, compiled-artifact SHA-256, environment SHA-256, checker layers), minted when the checking gate ran; the store receipts attesting the bound build are listed with the layers each run credited. Nothing here asks the reader to trust the store: every layer is an independently runnable program over the quoted sources --- rebuild the module and re-run the checkers to reproduce the verdict. Receipts attesting superseded builds credit no layer and are omitted from this rendering (the count of omissions, if any, is noted below).}
\begin{receipttable}
\receiptrow{a9c0e34d3e3c9c8c}{2026-08-03T21:30:36Z}{build+decorrelated}{\badge{BadgeGreen}{current}}
\end{receipttable}
\sourcenote{21 further receipt(s) in the store attest superseded builds of this module and are omitted from this rendering.}
\sourcenote{module \texttt{BernsteinLogic.Core.Catalogue}; bound build \texttt{cbc92e976a1d\ldots} (bound from receipt 5abdb3589707b5df, 2026-08-03T21:30:33Z). A receipt records that named checking layers accepted a build: (module, compiled-artifact SHA-256, environment SHA-256, checker layers), minted when the checking gate ran; the store receipts attesting the bound build are listed with the layers each run credited. Nothing here asks the reader to trust the store: every layer is an independently runnable program over the quoted sources --- rebuild the module and re-run the checkers to reproduce the verdict. Receipts attesting superseded builds credit no layer and are omitted from this rendering (the count of omissions, if any, is noted below).}
\begin{receipttable}
\receiptrow{24804e66edba0bc9}{2026-08-03T20:25:09Z}{build+decorrelated}{\badge{BadgeGreen}{current}}
\receiptrow{5abdb3589707b5df}{2026-08-03T21:30:33Z}{build+decorrelated}{\badge{BadgeGreen}{current}}
\end{receipttable}
\sourcenote{18 further receipt(s) in the store attest superseded builds of this module and are omitted from this rendering.}
\layerline{\layeritem{build}{lake build (elaboration)}{\badge{BadgeGreen}{current}}\layeritem{replay}{leanchecker (official kernel replay)}{\badge{BadgeAmber}{stale only}}\layeritem{hardened}{safe\_verify (hardened screen)}{\badge{BadgeAmber}{stale only}}\layeritem{decorrelated}{lean4export + nanoda (decorrelated kernel)}{\badge{BadgeGreen}{current}}}
\sourcenote{\badge{BadgeGray}{comparator: none} no fixture in the three-implementation comparator harness (compiled / kernel-interpreted / reference-interpreted) exercises this rule.}
\ruleblock{HOL4 statement}
\sourcenote{theory \texttt{bern\_catalogue}, HOL4 development script \texttt{\texttt{bern\_catalogueScript.sml}} \badge{BadgeGreen}{script bound}}
\holhead{Definition baryCoeffs1D\_def}
\begin{holsrc}
Definition baryCoeffs1D_def:
  baryCoeffs1D cx dmask =
    if cx.ambDim = 1 then
      (case FILTER (λv. getD dmask v F) (rangeL cx.nVertices) of
       | [a; b] =>
           (let v0 = getD (getD cx.coord a []) 0 (0:rat) in
            let v1 = getD (getD cx.coord b []) 0 (0:rat) in
            let den = v1 - v0 in
              if den = 0 then NONE
              else SOME [(a, v1 / den, -(&1) / den);
                         (b, -v0 / den, &1 / den)])
       | _ => NONE)
    else NONE
End
\end{holsrc}
\holhead{Definition primAffine\_def}
\begin{holsrc}
Definition primAffine_def:
  primAffine (s:rat) (c:rat) =
    if c = 0 then
      ((if s < 0 then -1 else if s = 0 then 0 else 1) : rat, 0:rat)
    else
      let l = lcm (RATD s) (RATD c) in
      let g = gcd (intAbsNum (RATN s * &(l DIV RATD s)))
                  (intAbsNum (RATN c * &(l DIV RATD c))) in
      (* g = 0 is unreachable (c ≠ 0 ⇒ RATN c ≠ 0 ⇒ the second gcd
         argument is nonzero); the guard is the translator's division
         precondition, spelled on the divisor (Lean's Rat division by
         zero yields 0, so the two sides agree even on the dead branch). *)
      let k = if &g = (0:rat) then 0 else &l / &g in
        (k * s, k * c)
End
\end{holsrc}
\holhead{Definition f4Lhs\_def}
\begin{holsrc}
Definition f4Lhs_def:
  f4Lhs (s':rat) (c':rat) t =
    if c' = 0 then TmMul (TmConst s') t
    else TmAdd (TmMul (TmConst s') t) (TmConst c')
End
\end{holsrc}
\holhead{Definition f4VertexAtom\_def}
\begin{holsrc}
Definition f4VertexAtom_def:
  f4VertexAtom inS (off:rat) (slp:rat) t =
    let sc = if inS then primAffine slp off else primAffine (-slp) (-off) in
      FAtom BowLe (f4Lhs (FST sc) (SND sc) t) (TmConst 0)
End
\end{holsrc}
\holhead{Definition barycentricHalfspaceBuild\_def}
\begin{holsrc}
Definition barycentricHalfspaceBuild_def:
  barycentricHalfspaceBuild t sx es =
    if es.simplex = sx then
      (case baryCoeffs1D sx.complex sx.vmask of
       | NONE => NONE
       | SOME cs =>
         (case MAP (λ(v,off,slp). f4VertexAtom (extContainsNat es v) off slp t)
               cs of
          | [x; y] => SOME (FAnd x y)
          | _ => NONE))
    else NONE
End
\end{holsrc}
\holhead{Definition barycentricHalfspaceFormB\_def}
\begin{holsrc}
Definition barycentricHalfspaceFormB_def:
  barycentricHalfspaceFormB t sx es fm ⇔
    barycentricHalfspaceBuild t sx es = SOME fm
End
\end{holsrc}
\ruleblock{HOL verifier receipts}
\badge{BadgeBlue}{documentary} \texttt{hol-strict} CLEAN census over the ported checker theory:
\texttt{bern\_catalogue}, 472 theorems, 0 cheat-tagged, 0 nonstandard
axioms (whitelist: \texttt{DISK\_THM}; the four \texttt{bool} axioms) ---
verification run of 2026-07-22 on an isolated build host, against a
content-addressed checkout of the exact sources cited here. The earlier
loud caveat (the canonical-form revision of 2026-07-18 postdated the
prior census run) is retired: this run attests exactly the quoted text.
There is no signed HOL receipt store yet, and these citations are
documentary.
Semantic stratum, also ported: the soundness twin
\texttt{sound\_convElim} (theory \texttt{bern\_tr\_fragclose}:
\texttt{hol-strict} clean, 26 theorems, 0 offences, verification run of
2026-07-22 on an isolated build host, content-addressed at the same
sources), consumed by the certified-fragment keystone
(\texttt{checkPf\_sound\_frag}, theory \texttt{bern\_tr\_splice}, 20
theorems, 0 offences, same run): a reference-checker accept inside the
proven wire fragment yields sequent validity.

\sourcenote{HOL receipts are documentary: quoted results of \texttt{hol-strict} runs
and census rows from the HOL4 development\textquotesingle s build
evidence, not minted receipt artefacts. A HOL receipt store (the Lean
store\textquotesingle s twin) is named future work.
}
\ruleblock{Proof}

\sourcenote{Genre: core rule (tier 1). Setting (valuations, partial
evaluation, elementwise atom semantics) as fixed in the \rulename{EqElim}
section; this section additionally fixes the region-atom semantics and
the canonical halfspace form, which the \rulename{ConvIntro} section then
cites.}

\textbf{Region atoms, in the checkable fragment.}\enspace A region atom
$x \in \mathcal{R}(\Delta, S)$ holds at $\sigma$ when the extension set
$S$ is tied to the simplex $\Delta$, the region term evaluates to a
rank-1 tensor of extent equal to the ambient dimension, and the point it
holds lies in the \emph{masked region}: it is an affine combination of
$\Delta$'s vertices with weights summing to $1$, nonnegative on core
vertices and \emph{nonpositive} on the vertices named by $S$. The rules
below operate in the fragment their checker accepts --- ambient dimension
$1$, support exactly two vertices $a, b$ at distinct coordinates
$V_0 \neq V_1$ --- where the geometry is completely explicit: the two
barycentric weights are \emph{determined}, affine in the point,
\[
\lambda_a(x) = \frac{V_1 - x}{V_1 - V_0}, \qquad
\lambda_b(x) = \frac{x - V_0}{V_1 - V_0}, \qquad
\lambda_a + \lambda_b \equiv 1,
\]
and the region is: the segment $[V_0, V_1]$ for $S = \emptyset$; the
closed ray from $V_1$ away from $V_0$ for $S = \{a\}$ (weight of the
extended vertex $\leq 0$); the mirror ray for $S = \{b\}$. ($S$ is a
proper subset, so both-extended does not arise.) Determinacy of the
weights is why elimination below is complete --- no witness is chosen,
none rides the certificate.

\textbf{The canonical halfspace form.}\enspace The conclusion is not
``some conjunction equivalent to the sign facts'' but one \emph{pinned}
formula, recomputed by the checker from $(x, \Delta, S)$ alone: one atom
per support vertex, in support order, each in $(\cdot) \leq 0$ shape:
for $v \in S$ the atom $\mathrm{canon}(\lambda_v(x)) \leq 0$, for
$v \notin S$ the atom $\mathrm{canon}(-\lambda_v(x)) \leq 0$ (the wanted
$\lambda_v \geq 0$ \emph{sign-folded} into the term --- the all-$\leq$
discipline), where $\mathrm{canon}$ rescales the affine pair
(slope, offset) by a \emph{strictly positive} rational to its
primitive-integer representative and lays the atom body out in
degree-sorted shape. If the build's gates fail (dimension $\neq 1$,
support $\neq 2$ vertices, $V_0 = V_1$, or the atom's simplex not the
cited $\Delta$), the build refuses, the payload
is uninhabitable, and the rule cannot fire: refusal is trivially sound,
and in the degenerate-simplex case it is also where a division by zero
would otherwise live --- made unrepresentable rather than checked-then-divided.

\textbf{Claim} (\texttt{sound\_convElim}).\enspace If
$\Phi \vdash x \in \mathcal{R}(\Delta, S)$ is semantically valid and the
conclusion formula is the canonical halfspace form for
$(x, \Delta, S)$, then $\Phi \vdash$ (that conjunction) is semantically
valid.

\begin{blemma}[structure of a successful build]\label{lem:convelim-spec}
If the canonical build succeeds with output $F$, then the ambient
dimension is $1$, $S$ is tied to $\Delta$, the support filter returns
exactly $[a, b]$ with constant vertex coordinates $V_0 \neq V_1$, and
$F$ is the conjunction of the two vertex atoms described above.
\end{blemma}
\begin{proof}
Unfold the build: each claim is one of its gates, and the output shape
is its final assembly. (\texttt{barycentricHalfspaceBuild\_spec}.)
\end{proof}

\begin{blemma}[canonical atom $\leftrightarrow$ raw sign fact]
\label{lem:convelim-atom}
Let $q$ be a defined, singleton-shaped evaluation of $x$ at $\sigma$, and
let $(\mathit{off}, \mathit{slp})$ be an affine pair. The vertex atom for
membership flag $m$ holds at $\sigma$ iff
\[
m = \mathrm{true}:\quad \mathit{off} + \mathit{slp}\cdot q \;\leq\; 0,
\qquad\qquad
m = \mathrm{false}:\quad 0 \;\leq\; \mathit{off} + \mathit{slp}\cdot q.
\]
\end{blemma}
\begin{proof}
Two steps. \emph{Sign-folding}: for $m = \mathrm{false}$ the build
negates the pair before canonicalizing, so the emitted atom asserts
$\mathrm{canon}(-\lambda) \leq 0$, i.e.\ the displayed $\geq$. \emph{The
rescale is sign-transparent}: $\mathrm{canon}$ multiplies the pair by
$k > 0$ --- in the generic case $k = L/g$ with $L$ a positive lcm of
denominators and $g$ a positive gcd of the scaled numerators; in the
zero-offset case the output slope is the sign
$\mathrm{sgn}(\mathit{slp})$, i.e.\ the pair scaled by $1/|\mathit{slp}|$
--- well-defined on every pair in play: a successful barycentric
build has distinct vertices, so each emitted affine pair carries
slope $\pm 1/(V_1 - V_0) \neq 0$ (sign-folding only negates it), and
in the zero-offset case that slope is the whole pair; for the
degenerate all-zero pair, which no successful build emits, both sides
of the displayed equivalence hold trivially and the statement is
unaffected --- and a strictly positive rescale of an affine
function has the same sign locus. The atom's left side evaluates (constants and one already-
defined $q$), at singleton shape, so the elementwise reading is the
scalar comparison of the displayed value against $0$.
(\texttt{f4VertexAtom\_holds\_iff}; the rescale is
\texttt{primAffine}, the layout \texttt{f4Lhs}.)
\end{proof}

\begin{blemma}[region elimination, two-vertex case]
\label{lem:convelim-elim}
Suppose $\sigma \models x \in \mathcal{R}(\Delta, S)$ in the fragment of
Lemma~\ref{lem:convelim-spec}, with $q$ the evaluated coordinate. Then,
conditionally on membership: if $a \in S$ then $\lambda_a(q) \leq 0$, if
$a \notin S$ then $\lambda_a(q) \geq 0$; likewise for $b$.
\end{blemma}
\begin{proof}
The masked-region witness consists of weights on the support summing to
$1$ that reproduce $q$ and satisfy the sign discipline. On a two-vertex
support the reproduction and sum equations force the weights to be
exactly $\lambda_a(q), \lambda_b(q)$ (two linear equations, invertible
since $V_0 \neq V_1$), so the witness's sign discipline \emph{is} the
displayed conditional facts. (\texttt{maskRegion\_pair\_elim}.)
\end{proof}

\begin{bthm}[soundness of \rulename{ConvElim}]\label{thm:convelim-sound}
Under the payload's checked hypotheses, if
$\Phi \vdash x \in \mathcal{R}(\Delta, S)$ is semantically valid, so is
$\Phi \vdash F$ for the built conclusion $F$.
\end{bthm}
\begin{proof}
Fix an admissible $\sigma \models \Phi$. Premise validity gives the
region atom at $\sigma$: in particular $x$ \emph{evaluates}, to a rank-1
extent-1 tensor (the atom's own shape field --- no typing pack is needed
here; contrast \autoref{rem:convelim-asym}), whose scalar is $q$.
Lemma~\ref{lem:convelim-elim} delivers the four conditional sign facts;
Lemma~\ref{lem:convelim-atom} converts, per vertex, the applicable one
into the built vertex atom; Lemma~\ref{lem:convelim-spec} says the
conclusion is exactly the conjunction of those two atoms.
\end{proof}

\begin{bremark}[why a pinned canonical spelling]\label{rem:convelim-canon}
The combination rules downstream (\rulename{Farkas}, \rulename{Syzygy})
match premise atoms \emph{by exact canonical polynomial}; a conclusion
emitted in any equivalent-but-different spelling would simply never
match, stranding the geometric fact. Pinning the spelling in the checker
--- which recomputes the build from $(x, \Delta, S)$ and demands
syntactic equality with the cited conclusion, trusting no producer text
--- makes the emitted atoms the \emph{same} tokens every producer and
every downstream matcher agrees on. The all-$\leq$ polarity keeps each
atom in the $(\cdot) \leq 0$ normal form those rules consume, and the
positive primitive-integer rescale makes the representative unique. Note
the checked object is the \emph{formula}, not a semantic condition: this
rule's side condition is one syntactic equation, decidable by
construction.
\end{bremark}

\begin{bremark}[what is \emph{not} hypothesized]\label{rem:convelim-asym}
No shape or finiteness fields appear in this payload: elimination reads
definedness and shape out of the region atom itself. The mirror rule
cannot --- its premise is the halfspace conjunction, whose atoms
tolerate broadcast and are satisfied by a \emph{scalar} $x$ --- and that
asymmetry is exactly where \rulename{ConvIntro}'s one extra typing field
comes from (next section).
\end{bremark}

\begin{anchors}
Lemma \ref{lem:convelim-spec} & \texttt{barycentricHalfspaceBuild\_spec}\\
build, affine data, rescale, atom layout & \texttt{barycentricHalfspaceBuild}, \texttt{baryCoeffs1D}, \texttt{primAffine}, \texttt{f4Lhs}, \texttt{f4VertexAtom}\\
Lemma \ref{lem:convelim-atom} & \texttt{f4VertexAtom\_holds\_iff}\\
Lemma \ref{lem:convelim-elim} & \texttt{maskRegion\_pair\_elim}\\
Theorem \ref{thm:convelim-sound} & \texttt{sound\_convElim}; wired as core \texttt{RuleData.convElim} into \texttt{ProofObject.sound}\\
form predicate (the checked side condition) & \texttt{barycentricHalfspaceForm} (payload field \texttt{halfspace})\\
\end{anchors}


\rulesection{13}{ConvIntro}{core rule}
\label{blcatrule-ConvIntro}
\ruleblock{The rule (Steele metanotation)}
\sourcenote{as realized in the mechanized development; source row: \texttt{\texttt{rules.toml}} --- \emph{ConvIntro}. The metanotation statement is typeset by hand for notational quality; its mechanized ground truth is the Lean statement below, extracted from the sources at every build of this book and hash-bound. The rule \emph{roster} is machine-checked (every mechanized rule tag round-trips to exactly one catalogue row); the statement \emph{text} is editorial, checked by eye against the extracted Lean beside it.}
\begin{mathpar}
\inferrule*[Right=\rulename{ConvIntro}]{
  \Phi \vdash \textstyle\bigwedge_{i = 0}^{n} \mathrm{canon}(\epsilon_i\,\lambda_i(x)) \leq 0 \\ \epsilon_i \,=\, {+1} \text{ if } i \in S, \ {-1} \text{ if } i \in \bar S \\ \bshape_H(x) = [\,n\,]
}{
  \Phi \vdash x \in \mathcal{R}(\Delta, S)
}
\end{mathpar}

\par\medskip
\rulename{ConvIntro} closes a region atom: from the canonical barycentric
halfspace conjunction for $(x, \Delta, S)$ --- the same pinned formula
\rulename{ConvElim} emits, recomputed by the same build --- it concludes
$\Phi \vdash x \in \mathcal{R}(\Delta, S)$. A producer uses it to
re-enter the geometric world after arithmetic: a subdivision branch or a
combination argument establishes the two vertex atoms, and this rule
converts them back into the region fact that \rulename{BernCone} or a
further region rule consumes.

What the checker recomputes: the build (as in \rulename{ConvElim} ---
one syntactic equation, producer text untrusted, refusal outside the 1-D
two-vertex fragment), plus this rule\textquotesingle s one extra typing
field: $\bshape_H(x) = [\,d\,]$, the region term is statically rank-1 of
the ambient extent. The asymmetry with \rulename{ConvElim} is
principled: elimination reads the term\textquotesingle s shape out of
the region atom it consumes, while introduction is handed only halfspace
atoms --- which hold \emph{elementwise under broadcast} and are
perfectly satisfied by a scalar $x$ --- so the shape must arrive as a
checked hypothesis. Drop it and a closed scalar inside the segment
yields a satisfiable (indeed universally true) premise whose region
conclusion is unsatisfiable --- a rank-0 tensor never satisfies a
region atom; the proof section builds the refutation concretely. It is
the same broadcast liberality that breaks un-gated \rulename{EqElim},
repaired by the same typing-field pattern.

\cataloguenote{On the fragment, \rulename{ConvIntro} and \rulename{ConvElim} are exact
inverses --- membership, the four weight-sign facts, and the canonical
conjunction are three presentations of one determinate-weight geometry,
so the pair is a lossless interface between region atoms and the
arithmetic fragment.
}
\ruleblock{Lean statement}
\declhead{BernsteinLogic.Soundness}{Term.evalQ\_bshape}{\badge{BadgeGreen}{statement bound}}
\begin{leansrc}
/-- bshape soundness: a defined evaluation takes the broadcast shape. -/
theorem Term.evalQ_bshape {H : Hypernet} (σ : H.Valuation) :
    ∀ (t : Term) (s : Shape) (q : QTensor),
    t.bshape H = some s → Term.evalQ σ t = some q → q.shape = s
\end{leansrc}
\declhead{BernsteinLogic.Soundness}{f4VertexAtom\_holds\_inv}{\badge{BadgeGreen}{statement bound}}
\begin{leansrc}
/-- A held F4 vertex atom forces the core term to evaluate. -/
theorem f4VertexAtom_holds_inv {H : Hypernet} {σ : H.Valuation} {t : Term}
    {inS : Bool} {off slp : ℚ}
    (h : Formula.holds σ (f4VertexAtom inS off slp t)) :
    ∃ q, Term.evalQ σ t = some q
\end{leansrc}
\declhead{BernsteinLogic.Soundness}{maskRegion\_pair\_intro}{\badge{BadgeGreen}{statement bound}}
\begin{leansrc}
/-- Construction: rational weights of the right signs assemble a
`MaskRegion` witness over a two-vertex 1-D support. -/
theorem maskRegion_pair_intro {C : SimpC} {mask ext : Fin C.nVertices → Bool}
    {x : Fin C.ambDim → ℚ} {a b : Fin C.nVertices} {V₀ V₁ : ℚ}
    (hfil : (List.finRange C.nVertices).filter (fun v => mask v) = [a, b])
    (hVa : ∀ j, C.coord a j = V₀) (hVb : ∀ j, C.coord b j = V₁)
    (hne : V₀ ≠ V₁) (hdim : C.ambDim = 1) (j0 : Fin C.ambDim)
    (hw0t : ext a = true → (V₁ - x j0) / (V₁ - V₀) ≤ 0)
    (hw0f : ext a = false → 0 ≤ (V₁ - x j0) / (V₁ - V₀))
    (hw1t : ext b = true → (x j0 - V₀) / (V₁ - V₀) ≤ 0)
    (hw1f : ext b = false → 0 ≤ (x j0 - V₀) / (V₁ - V₀)) :
    MaskRegion C mask ext x
\end{leansrc}
\declhead{BernsteinLogic.Soundness}{sound\_convIntro}{\badge{BadgeGreen}{statement bound}}
\begin{leansrc}
theorem sound_convIntro {H : Hypernet} {seq premise : Sequent}
    {term : Term} {simplex : Simplex} {S : ExtensionSet}
    (ih : premise.semanticallyValid H)
    (hpre : seq.pre = premise.pre)
    (htermShape : term.bshape H = some [simplex.complex.ambDim])
    (hhalf : barycentricHalfspaceForm term simplex S premise.post)
    (hpost : seq.post = .inRegion term simplex S) :
    seq.semanticallyValid H
\end{leansrc}
\par\medskip
\textbf{Shared machinery.}\enspace The build chain is displayed in full
in the \rulename{ConvElim} section; here the book re-binds the assembly
and the form equation, plus this rule\textquotesingle s payload with its
one extra typing field (\texttt{termShape}).

\declhead{BernsteinLogic.Core.Catalogue}{barycentricHalfspaceBuild}{\badge{BadgeGreen}{definition bound}}
\begin{leansrc}
def barycentricHalfspaceBuild (t : Term) (simplex : Simplex)
    (S : ExtensionSet) : Option Formula := do
  guard (S.simplex = simplex)
  let cs ← baryCoeffs1D simplex.complex simplex.contains
  match cs.map (fun (v, off, slp) =>
    f4VertexAtom (S.containsNat v.val) off slp t) with
  | [x, y] => some (.and x y)
  | _ => none
\end{leansrc}
\declhead{BernsteinLogic.Core.Catalogue}{barycentricHalfspaceForm}{\badge{BadgeGreen}{definition bound}}
\begin{leansrc}
def barycentricHalfspaceForm (t : Term) (simplex : Simplex)
    (S : ExtensionSet) (F : Formula) : Prop :=
  barycentricHalfspaceBuild t simplex S = some F
\end{leansrc}
\declhead{BernsteinLogic.Core.Catalogue}{PConvIntroData}{\badge{BadgeGreen}{definition bound}}
\begin{leansrc}
inductive PConvIntroData : Hypernet → Sequent → Type where
  | mk (H : Hypernet) (seq : Sequent)
       (premise : Sequent) (premiseProof : ProofObject H premise)
       (term : Term) (simplex : Simplex) (S : ExtensionSet)
       (preEq     : seq.pre = premise.pre)
       (termShape : term.bshape H = some [simplex.complex.ambDim])
       (halfspace : barycentricHalfspaceForm term simplex S premise.post)
       (postEq    : seq.post = .inRegion term simplex S)
       : PConvIntroData H seq
\end{leansrc}
\ruleblock{Correspondence (Steele display $\leftrightarrow$ Lean)}
\badge{BadgeGreen}{reviewer-certified against the displayed texts}
\sourcenote{The dictionary below is a reviewer-certified statement of how the mechanized statements above interpret the displayed rule: display metavariable $\mapsto$ Lean binder, each display hypothesis $\mapsto$ the premise or side condition carrying it, the display conclusion $\mapsto$ the goal --- with any deliberate difference in width between display and mechanization named outright. Certification is recorded as a hash triple over (the displayed cell, the extracted statement texts, this dictionary\textquotesingle s own text); if any of the three changes, this block turns amber until a reviewer re-certifies. The correspondence claim itself remains editorial --- reviewer-grade, not machine-checked; its machine-checked successor types each soundness statement against the rule as a structured object, and supersedes this block per rule as it lands.}
Judgment frame as in \rulename{Ax}; sub-derivation premise as at
\rulename{Weaken}; binders as at \rulename{ConvElim}, mirrored.
\begin{itemize}
\item the displayed premise conjunction $\mapsto$ \texttt{hhalf} $:$
  \texttt{barycentricHalfspaceForm term simplex S premise.post} ---
  the same verified build, its value now equated with the
  \emph{premise\textquotesingle s} post; the premise judgment
  $\mapsto$ \texttt{ih}, context shared via \texttt{hpre}.
\item the display\textquotesingle s typing premise $\bshape_H(x) =
  [\,n\,]$ $\mapsto$ \texttt{htermShape} $:$ \texttt{term.bshape H =
  some [simplex.complex.ambDim]} --- statically rank-1 at the ambient
  extent, the one hypothesis \rulename{ConvElim} does not have; the
  commentary derives why introduction needs it and elimination does
  not.
\item conclusion $\Phi \vdash x \in \mathcal{R}(\Delta, S)$ $\mapsto$
  the goal under \texttt{hpost} $:$ \texttt{seq.post = .inRegion term
  simplex S}.
\end{itemize}
Width note: the same one-dimensional fragment confinement as
\rulename{ConvElim}, forced by the same spec conjuncts --- the
mechanized statement is the displayed rule at ambient dimension $1$;
outside it the build refuses and nothing is accepted unverified. The
\rulename{ConvElim} dictionary\textquotesingle s spelling alignment
carries over unchanged: the displayed premise conjunction and the
certified premise formula both state the build\textquotesingle s
sign-folded all-$\leq$ shape, the extension/core orientation carried
by the sign $\epsilon_i$ --- same halfspaces, same folded spelling,
display and mechanization reconciled.

\ruleblock{Lean verifier receipts}
\sourcenote{module \texttt{BernsteinLogic.Soundness}; bound build \texttt{a20d90e66c13\ldots} (bound from receipt a9c0e34d3e3c9c8c, 2026-08-03T21:30:36Z). A receipt records that named checking layers accepted a build: (module, compiled-artifact SHA-256, environment SHA-256, checker layers), minted when the checking gate ran; the store receipts attesting the bound build are listed with the layers each run credited. Nothing here asks the reader to trust the store: every layer is an independently runnable program over the quoted sources --- rebuild the module and re-run the checkers to reproduce the verdict. Receipts attesting superseded builds credit no layer and are omitted from this rendering (the count of omissions, if any, is noted below).}
\begin{receipttable}
\receiptrow{a9c0e34d3e3c9c8c}{2026-08-03T21:30:36Z}{build+decorrelated}{\badge{BadgeGreen}{current}}
\end{receipttable}
\sourcenote{21 further receipt(s) in the store attest superseded builds of this module and are omitted from this rendering.}
\sourcenote{module \texttt{BernsteinLogic.Core.Catalogue}; bound build \texttt{cbc92e976a1d\ldots} (bound from receipt 5abdb3589707b5df, 2026-08-03T21:30:33Z). A receipt records that named checking layers accepted a build: (module, compiled-artifact SHA-256, environment SHA-256, checker layers), minted when the checking gate ran; the store receipts attesting the bound build are listed with the layers each run credited. Nothing here asks the reader to trust the store: every layer is an independently runnable program over the quoted sources --- rebuild the module and re-run the checkers to reproduce the verdict. Receipts attesting superseded builds credit no layer and are omitted from this rendering (the count of omissions, if any, is noted below).}
\begin{receipttable}
\receiptrow{24804e66edba0bc9}{2026-08-03T20:25:09Z}{build+decorrelated}{\badge{BadgeGreen}{current}}
\receiptrow{5abdb3589707b5df}{2026-08-03T21:30:33Z}{build+decorrelated}{\badge{BadgeGreen}{current}}
\end{receipttable}
\sourcenote{18 further receipt(s) in the store attest superseded builds of this module and are omitted from this rendering.}
\layerline{\layeritem{build}{lake build (elaboration)}{\badge{BadgeGreen}{current}}\layeritem{replay}{leanchecker (official kernel replay)}{\badge{BadgeAmber}{stale only}}\layeritem{hardened}{safe\_verify (hardened screen)}{\badge{BadgeAmber}{stale only}}\layeritem{decorrelated}{lean4export + nanoda (decorrelated kernel)}{\badge{BadgeGreen}{current}}}
\sourcenote{\badge{BadgeGray}{comparator: none} no fixture in the three-implementation comparator harness (compiled / kernel-interpreted / reference-interpreted) exercises this rule.}
\ruleblock{HOL4 statement}
\sourcenote{theory \texttt{bern\_catalogue}, HOL4 development script \texttt{\texttt{bern\_catalogueScript.sml}} \badge{BadgeGreen}{script bound}}
\holhead{Definition barycentricHalfspaceBuild\_def}
\begin{holsrc}
Definition barycentricHalfspaceBuild_def:
  barycentricHalfspaceBuild t sx es =
    if es.simplex = sx then
      (case baryCoeffs1D sx.complex sx.vmask of
       | NONE => NONE
       | SOME cs =>
         (case MAP (λ(v,off,slp). f4VertexAtom (extContainsNat es v) off slp t)
               cs of
          | [x; y] => SOME (FAnd x y)
          | _ => NONE))
    else NONE
End
\end{holsrc}
\holhead{Definition barycentricHalfspaceFormB\_def}
\begin{holsrc}
Definition barycentricHalfspaceFormB_def:
  barycentricHalfspaceFormB t sx es fm ⇔
    barycentricHalfspaceBuild t sx es = SOME fm
End
\end{holsrc}
\ruleblock{HOL verifier receipts}
\badge{BadgeBlue}{documentary} \texttt{hol-strict} CLEAN census over the ported checker theory:
\texttt{bern\_catalogue}, 472 theorems, 0 cheat-tagged, 0 nonstandard
axioms (whitelist: \texttt{DISK\_THM}; the four \texttt{bool} axioms) ---
verification run of 2026-07-22 on an isolated build host, against a
content-addressed checkout of the exact sources cited here --- the
re-attestation the \rulename{ConvElim} cell\textquotesingle s earlier
caveat awaited (the canonical-form revision of 2026-07-18 postdated the
prior census) is this run (documentary citations --- no signed HOL
store yet).
Scope note: the shape side condition (\texttt{termShape}) rides the
un-ported hypernet typing island, deferred to the port\textquotesingle s
proof-object stage.
Semantic stratum, also ported: the soundness twin
\texttt{sound\_convIntro} (theory \texttt{bern\_tr\_fragclose}:
\texttt{hol-strict} clean, 26 theorems, 0 offences, verification run of
2026-07-22 on an isolated build host, content-addressed at the same
sources), consumed by the certified-fragment keystone
(\texttt{checkPf\_sound\_frag}, theory \texttt{bern\_tr\_splice}, 20
theorems, 0 offences, same run): a reference-checker accept inside the
proven wire fragment yields sequent validity.

\sourcenote{HOL receipts are documentary: quoted results of \texttt{hol-strict} runs
and census rows from the HOL4 development\textquotesingle s build
evidence, not minted receipt artefacts. A HOL receipt store (the Lean
store\textquotesingle s twin) is named future work.
}
\ruleblock{Proof}

\sourcenote{Genre: core rule (tier 1). Setting, region-atom semantics,
the checkable fragment, and the canonical halfspace form as fixed in the
\rulename{ConvElim} section.}

\textbf{Claim} (\texttt{sound\_convIntro}).\enspace The payload supplies
the region data $(x, \Delta, S)$ and two checked conditions: the premise
concludes the canonical halfspace form for $(x, \Delta, S)$, and (the
shape condition) $\bshape_H(x) = [\,d\,]$ where $d$ is the ambient
dimension --- the region term is statically a rank-1 tensor of the
ambient extent. If the premise is semantically valid, then
$\Phi \vdash x \in \mathcal{R}(\Delta, S)$ is semantically valid.

\begin{blemma}[evaluation recovered from the premise]
\label{lem:convintro-eval}
Suppose $\sigma$ satisfies the built conjunction. Then $x$ evaluates at
$\sigma$.
\end{blemma}
\begin{proof}
The first vertex atom holds, so its left side --- the canonicalized
affine image of $x$, built from constants and $x$ by product and sum ---
evaluates; evaluation of a compound is defined only if evaluation of
every part is --- the evaluator is strict, each node's value built
from its children's values, so a defined parent forces defined
children --- and inverting it at the product node (and sum
node, when the offset is present) yields a value for $x$ itself.
(\texttt{f4VertexAtom\_holds\_inv}.)
\end{proof}

\begin{blemma}[shape pinned by the payload field]
\label{lem:convintro-shape}
With the shape condition, the evaluation of Lemma~\ref{lem:convintro-eval}
is a rank-1 tensor of extent $1$ (the fragment has $d = 1$), hence
singleton-shaped, with scalar $q$.
\end{blemma}
\begin{proof}
Broadcast-shape soundness: a defined evaluation of a $\bshape$-typed term
lands at the declared shape (\texttt{Term.evalQ\_bshape}), here $[\,1\,]$
by the build's dimension gate (\rulename{ConvElim},
Lemma~\ref{lem:convelim-spec}).
\end{proof}

\begin{bthm}[soundness of \rulename{ConvIntro}]\label{thm:convintro-sound}
Under the payload's checked hypotheses, if $\Phi \vdash F$ is
semantically valid for the built conjunction $F$, so is
$\Phi \vdash x \in \mathcal{R}(\Delta, S)$.
\end{bthm}
\begin{proof}
Fix an admissible $\sigma \models \Phi$; premise validity gives both
vertex atoms at $\sigma$. Lemmas~\ref{lem:convintro-eval} and
\ref{lem:convintro-shape} produce the singleton evaluation $q$ of $x$.
Reading each held atom through the canonical-atom characterization
(\rulename{ConvElim}, Lemma~\ref{lem:convelim-atom}) yields the four
conditional sign facts on $\lambda_a(q), \lambda_b(q)$. The two-vertex
introduction assembles the masked-region witness from exactly these: the
weights $\lambda_a(q), \lambda_b(q)$ sum to $1$ and reproduce $q$ by
construction (affine identities, valid since $V_0 \neq V_1$), and the
sign discipline --- nonnegative on core vertices, nonpositive on
$S$-vertices --- is the four facts (\texttt{maskRegion\_pair\_intro}).
Together with the build's tie of $S$ to $\Delta$ and the shape of
Lemma~\ref{lem:convintro-shape}, this is the region atom's holding, term
slot included.
\end{proof}

\begin{bremark}[why the shape condition is not decoration]
\label{rem:convintro-cex}
Drop it and the rule is \emph{unsound}. Over $\Delta =
\mathrm{conv}(1, 2)$, $S = \emptyset$, instantiate the region term at the
\emph{closed scalar} $t := 3/2$. The built premise is then a conjunction
of two closed atoms asserting $\lambda_a(3/2) \geq 0$ and
$\lambda_b(3/2) \geq 0$ --- both true ($\lambda_a = \lambda_b = 1/2$), so
the conjunction holds at \emph{every} valuation, and taking $\Phi$ to be
that very conjunction, \rulename{Ax} proves the premise. The un-gated
introduction would conclude
$\Phi \vdash (\tfrac{3}{2} : \mathrm{Term}) \in \mathcal{R}(\Delta,
\emptyset)$ --- unsatisfiable: a constant evaluates to a rank-0 tensor,
and the region atom demands rank 1. Satisfiable premise, unsatisfiable
conclusion. The mechanism is the same broadcast liberality that broke the
un-gated \rulename{EqElim}: the halfspace atoms constrain the
\emph{values} of the affine images and are perfectly happy with a scalar
$x$; only the region atom cares about the tensor's shape. The shape field
restores, at the introduction rule, the well-formedness that region
atoms carry in the source calculus's typing discipline --- one repair
family (the \rulename{EqElim}/\rulename{Norm} typing pattern), one
failure family.
\end{bremark}

\begin{bremark}[elimination and introduction are exact inverses on the
fragment]\label{rem:convintro-inv}
On singleton-shaped, well-shaped region terms, Lemmas
\ref{lem:convelim-elim} (elimination) and the introduction step above are
mutually inverse readings of the same determinate-weight geometry:
membership $\leftrightarrow$ the four sign facts $\leftrightarrow$ the
canonical conjunction. The round trip loses nothing --- which is why the
pair can serve as the sole interface between region atoms and the
arithmetic fragment, with no completeness caveat inside the checkable
fragment.
\end{bremark}

\begin{anchors}
Lemma \ref{lem:convintro-eval} & \texttt{f4VertexAtom\_holds\_inv}\\
Lemma \ref{lem:convintro-shape} & \texttt{Term.evalQ\_bshape}; payload field \texttt{termShape}\\
canonical-atom reading & \texttt{f4VertexAtom\_holds\_iff} (\rulename{ConvElim} section)\\
witness assembly & \texttt{maskRegion\_pair\_intro}\\
Theorem \ref{thm:convintro-sound} & \texttt{sound\_convIntro}; wired as core \texttt{RuleData.convIntro} into \texttt{ProofObject.sound}\\
Remark \ref{rem:convintro-cex} & the counterexample inventory, \texttt{Soundness.lean} header\\
\end{anchors}


\rulesection{14}{ArithAtom}{core rule}
\label{blcatrule-ArithAtom}
\ruleblock{The rule (Steele metanotation)}
\sourcenote{as realized in the mechanized development; source row: \texttt{\texttt{rules.toml}} --- \emph{ArithAtom}. The metanotation statement is typeset by hand for notational quality; its mechanized ground truth is the Lean statement below, extracted from the sources at every build of this book and hash-bound. The rule \emph{roster} is machine-checked (every mechanized rule tag round-trips to exactly one catalogue row); the statement \emph{text} is editorial, checked by eye against the extracted Lean beside it.}
\begin{mathpar}
\inferrule*[Right=\rulename{ArithAtom}]{
  \Phi \vdash (t_1 - t_2) \bowtie 0
}{
  \Phi \vdash t_1 \bowtie t_2
}
\\
\inferrule*[Right=\rulename{ArithAtom}$^{-1}$]{
  \Phi \vdash t_1 \bowtie t_2
}{
  \Phi \vdash (t_1 - t_2) \bowtie 0
}
\end{mathpar}
with $\bowtie \in \{<, \leq, =, \geq, >\}$.

\par\medskip
\rulename{ArithAtom} toggles a comparison atom between its two
spellings: the \emph{direct form} $t_1 \bowtie t_2$ and the
\emph{zero-normal form} $(t_1 - t_2) \bowtie 0$, in both directions
(one payload constructor each), for
$\bowtie \in \{<, \leq, =, \geq, >\}$. It is the adapter at the
calculus\textquotesingle s normal-form boundary: the combination rules
(\rulename{Farkas}, \rulename{Syzygy}) and the canonical conclusions of
the geometry rules live in $(\cdot) \leq 0$-normal shapes, while
human-facing specifications, box premises of the relaxation rules, and
final theorem statements live in direct forms. A producer folds a
zero-form conclusion into the literal bound a downstream rule wants, or
unfolds a specification atom so its polynomial content can be matched.

What the checker recomputes: two syntactic form equations --- the
premise (respectively conclusion) must be \emph{literally} the atom
built from the same payload terms with the syntactic difference
$t_1 - t_2$ against the literal constant $0$ --- plus the comparator
gate $\bowtie \neq \mathord{\neq}$ and the shared-context equation.
Deliberately severe: no term evaluating to zero, no algebraic
rearrangement qualifies (a producer wanting those routes through
\rulename{Norm} first); the severity is what keeps the side conditions
decidable by construction and the trusted surface at the two form
equations, with the comparator gate and the context equation beside
them.

There is no droppable typing hypothesis, and that is the
rule\textquotesingle s quiet achievement: both forms mention the same
wires, and whenever either form holds, both evaluations exist and
transport --- in the
fold direction the difference\textquotesingle s definedness arrives with
the premise, and in the unfold direction it is manufactured from the
fact that a \emph{holding} elementwise comparison already certifies
broadcast compatibility, on exactly which subtraction is total. The
comparator gate is the one non-structural field: the rational semantics
would accept $\neq$ (the entrywise bridge is comparator-uniform), and
the gate instead pins the source calculus\textquotesingle s comparator
set and keeps the development\textquotesingle s auxiliary three-valued
screen uniform across the atom-shape rules (see the proof
section\textquotesingle s closing remark --- the screen itself would
survive $\neq$ here, unlike at the ground-fact rule).

\cataloguenote{The transported object is the per-entry rational comparison
$(p - q) \bowtie 0 \iff p \bowtie q$ --- translation invariance of the
order of $\Q$ --- lifted entrywise through the broadcast discipline by
an invertible totality argument. Exact arithmetic end to end: no
rounding, no tolerance, in either direction.
}
\ruleblock{Lean statement}
\declhead{BernsteinLogic.Soundness}{QTensor.map2\_defined\_of\_rel2}{\badge{BadgeGreen}{statement bound}}
\begin{leansrc}
/-- A holding rel2 forces broadcast-compatible shapes, so map2 succeeds. -/
theorem QTensor.map2_defined_of_rel2 {r : ℚ → ℚ → Prop} {a b : QTensor}
    (h : QTensor.rel2 r a b) (f : ℚ → ℚ → ℚ) :
    ∃ c, QTensor.map2 f a b = some c
\end{leansrc}
\declhead{BernsteinLogic.Soundness}{QTensor.rel2\_sub\_zero\_iff}{\badge{BadgeGreen}{statement bound}}
\begin{leansrc}
/-- The ArithAtom bridge, tensor-general: comparing a subtraction against
rank-0 zero is comparing the operands directly, in every broadcast case. -/
theorem QTensor.rel2_sub_zero_iff {b : BowTie} {al ar d : QTensor}
    (hd : QTensor.map2 (· - ·) al ar = some d) :
    QTensor.rel2 (BowTie.holdsQ b) d ⟨[], [0], rfl⟩
      ↔ QTensor.rel2 (BowTie.holdsQ b) al ar
\end{leansrc}
\declhead{BernsteinLogic.Soundness}{Term.evalQ\_sub\_inv}{\badge{BadgeGreen}{statement bound}}
\begin{leansrc}
theorem Term.evalQ_sub_inv {H : Hypernet} (σ : H.Valuation) {l r : Term}
    {q : QTensor} (h : Term.evalQ σ (.sub l r) = some q) :
    ∃ a b, Term.evalQ σ l = some a ∧ Term.evalQ σ r = some b ∧
      QTensor.map2 (· - ·) a b = some q
\end{leansrc}
\declhead{BernsteinLogic.Soundness}{sound\_arithAtom\_fold}{\badge{BadgeGreen}{statement bound}}
\begin{leansrc}
theorem sound_arithAtom_fold {H : Hypernet} {seq premise : Sequent}
    {termL termR : Term} {b : BowTie}
    (ih : premise.semanticallyValid H)
    (hpre : seq.pre = premise.pre)
    (hprem : atomicCompareZero b termL termR premise.post)
    (hconcl : atomicCompareForm b termL termR seq.post) :
    seq.semanticallyValid H
\end{leansrc}
\declhead{BernsteinLogic.Soundness}{sound\_arithAtom\_unfold}{\badge{BadgeGreen}{statement bound}}
\begin{leansrc}
theorem sound_arithAtom_unfold {H : Hypernet} {seq premise : Sequent}
    {termL termR : Term} {b : BowTie}
    (ih : premise.semanticallyValid H)
    (hpre : seq.pre = premise.pre)
    (hprem : atomicCompareForm b termL termR premise.post)
    (hconcl : atomicCompareZero b termL termR seq.post) :
    seq.semanticallyValid H
\end{leansrc}
\par\medskip
\textbf{The two form equations, in full.}\enspace Everything this rule
checks is displayed below: two one-line syntactic equations and the
payload\textquotesingle s two constructors (one per direction) with the
comparator gate \texttt{notNE}. No valuation is consulted; semantic
truth enters only through the soundness lemmas.

\declhead{BernsteinLogic.Core.Catalogue}{atomicCompareForm}{\badge{BadgeGreen}{definition bound}}
\begin{leansrc}
def atomicCompareForm (b : BowTie) (l r : Term) (F : Formula) : Prop :=
  F = .atom b l r
\end{leansrc}
\declhead{BernsteinLogic.Core.Catalogue}{atomicCompareZero}{\badge{BadgeGreen}{definition bound}}
\begin{leansrc}
def atomicCompareZero (b : BowTie) (l r : Term) (F : Formula) : Prop :=
  F = .atom b (.sub l r) (.const 0)
\end{leansrc}
\declhead{BernsteinLogic.Core.Catalogue}{PArithAtomData}{\badge{BadgeGreen}{definition bound}}
\begin{leansrc}
inductive PArithAtomData : Hypernet → Sequent → Type where
  | fold   (H : Hypernet) (seq : Sequent)
           (premise : Sequent) (premiseProof : ProofObject H premise)
           (termL termR : Term) (b : BowTie)
           (notNE          : b ≠ .ne)
           (preEq          : seq.pre = premise.pre)
           (premiseIsZero  : atomicCompareZero b termL termR premise.post)
           (conclIsDirect  : atomicCompareForm b termL termR seq.post)
           : PArithAtomData H seq
  | unfold (H : Hypernet) (seq : Sequent)
           (premise : Sequent) (premiseProof : ProofObject H premise)
           (termL termR : Term) (b : BowTie)
           (notNE           : b ≠ .ne)
           (preEq           : seq.pre = premise.pre)
           (premiseIsDirect : atomicCompareForm b termL termR premise.post)
           (conclIsZero     : atomicCompareZero b termL termR seq.post)
           : PArithAtomData H seq
\end{leansrc}
\ruleblock{Correspondence (Steele display $\leftrightarrow$ Lean)}
\badge{BadgeGreen}{reviewer-certified against the displayed texts}
\sourcenote{The dictionary below is a reviewer-certified statement of how the mechanized statements above interpret the displayed rule: display metavariable $\mapsto$ Lean binder, each display hypothesis $\mapsto$ the premise or side condition carrying it, the display conclusion $\mapsto$ the goal --- with any deliberate difference in width between display and mechanization named outright. Certification is recorded as a hash triple over (the displayed cell, the extracted statement texts, this dictionary\textquotesingle s own text); if any of the three changes, this block turns amber until a reviewer re-certifies. The correspondence claim itself remains editorial --- reviewer-grade, not machine-checked; its machine-checked successor types each soundness statement against the rule as a structured object, and supersedes this block per rule as it lands.}
Judgment frame as in \rulename{Ax}; sub-derivation premise as at
\rulename{Weaken}.
\begin{itemize}
\item the display\textquotesingle s two inference lines $\mapsto$ the
  lemma pair, mirroring the payload\textquotesingle s two
  constructors: \rulename{ArithAtom} (zero-form premise, direct
  conclusion) $\mapsto$ \texttt{sound\_arithAtom\_fold};
  \rulename{ArithAtom}$^{-1}$ $\mapsto$
  \texttt{sound\_arithAtom\_unfold}.
\item $t_1, t_2 \mapsto$ \texttt{termL termR : Term}; $\bowtie
  \mapsto$ \texttt{b : BowTie}; the two spellings $\mapsto$ the form
  predicates quoted below, each a one-line syntactic equation:
  \texttt{atomicCompareForm} (the atom is literally $t_1 \bowtie t_2$)
  and \texttt{atomicCompareZero} (literally $(t_1 - t_2) \bowtie 0$,
  syntactic difference against the literal constant).
\item $\Phi \mapsto$ the shared context via \texttt{hpre}; the
  premise judgment $\mapsto$ \texttt{ih}; the conclusion $\mapsto$ the
  goal \texttt{seq.semanticallyValid H}. Certified alongside:
  \texttt{QTensor.map2\_defined\_of\_rel2}, the totality bridge with
  which the unfold direction manufactures the
  difference\textquotesingle s definedness from the holding premise.
\end{itemize}
Width note, deliberate and load-bearing (the \rulename{PConstExact}
pattern): the display restricts $\bowtie \in \{<, \leq, =, \geq, >\}$,
while both lemmas are proven for \emph{all} constructors of
\texttt{BowTie}, including $\neq$ --- no $\bowtie \neq \mathord{\neq}$
hypothesis appears in either statement. The exclusion is an
admission-layer fact: the payload\textquotesingle s \texttt{notNE}
field refuses $\neq$ certificates to protect the three-valued screen
the commentary describes. The mechanized statements are therefore
strictly wider than the displayed rule, and the displayed rule is
exactly the admitted fragment of them.

\ruleblock{Lean verifier receipts}
\sourcenote{module \texttt{BernsteinLogic.Soundness}; bound build \texttt{a20d90e66c13\ldots} (bound from receipt a9c0e34d3e3c9c8c, 2026-08-03T21:30:36Z). A receipt records that named checking layers accepted a build: (module, compiled-artifact SHA-256, environment SHA-256, checker layers), minted when the checking gate ran; the store receipts attesting the bound build are listed with the layers each run credited. Nothing here asks the reader to trust the store: every layer is an independently runnable program over the quoted sources --- rebuild the module and re-run the checkers to reproduce the verdict. Receipts attesting superseded builds credit no layer and are omitted from this rendering (the count of omissions, if any, is noted below).}
\begin{receipttable}
\receiptrow{a9c0e34d3e3c9c8c}{2026-08-03T21:30:36Z}{build+decorrelated}{\badge{BadgeGreen}{current}}
\end{receipttable}
\sourcenote{21 further receipt(s) in the store attest superseded builds of this module and are omitted from this rendering.}
\sourcenote{module \texttt{BernsteinLogic.Core.Catalogue}; bound build \texttt{cbc92e976a1d\ldots} (bound from receipt 5abdb3589707b5df, 2026-08-03T21:30:33Z). A receipt records that named checking layers accepted a build: (module, compiled-artifact SHA-256, environment SHA-256, checker layers), minted when the checking gate ran; the store receipts attesting the bound build are listed with the layers each run credited. Nothing here asks the reader to trust the store: every layer is an independently runnable program over the quoted sources --- rebuild the module and re-run the checkers to reproduce the verdict. Receipts attesting superseded builds credit no layer and are omitted from this rendering (the count of omissions, if any, is noted below).}
\begin{receipttable}
\receiptrow{24804e66edba0bc9}{2026-08-03T20:25:09Z}{build+decorrelated}{\badge{BadgeGreen}{current}}
\receiptrow{5abdb3589707b5df}{2026-08-03T21:30:33Z}{build+decorrelated}{\badge{BadgeGreen}{current}}
\end{receipttable}
\sourcenote{18 further receipt(s) in the store attest superseded builds of this module and are omitted from this rendering.}
\layerline{\layeritem{build}{lake build (elaboration)}{\badge{BadgeGreen}{current}}\layeritem{replay}{leanchecker (official kernel replay)}{\badge{BadgeAmber}{stale only}}\layeritem{hardened}{safe\_verify (hardened screen)}{\badge{BadgeAmber}{stale only}}\layeritem{decorrelated}{lean4export + nanoda (decorrelated kernel)}{\badge{BadgeGreen}{current}}}
\sourcenote{\badge{BadgeGray}{comparator: none} no fixture in the three-implementation comparator harness (compiled / kernel-interpreted / reference-interpreted) exercises this rule.}
\ruleblock{HOL4 statement}
\sourcenote{theory \texttt{bern\_catalogue}, HOL4 development script \texttt{\texttt{bern\_catalogueScript.sml}} \badge{BadgeGreen}{script bound}}
\holhead{Definition atomicCompareFormB\_def}
\begin{holsrc}
Definition atomicCompareFormB_def:
  atomicCompareFormB b l r fm ⇔ fm = FAtom b l r
End
\end{holsrc}
\holhead{Definition atomicCompareZeroB\_def}
\begin{holsrc}
Definition atomicCompareZeroB_def:
  atomicCompareZeroB b l r fm ⇔ fm = FAtom b (TmSub l r) (TmConst 0)
End
\end{holsrc}
\ruleblock{HOL verifier receipts}
\badge{BadgeBlue}{documentary} \texttt{hol-strict} CLEAN census over the ported checker theory:
\texttt{bern\_catalogue}, 472 theorems, 0 cheat-tagged, 0 nonstandard
axioms (whitelist: \texttt{DISK\_THM}; the four \texttt{bool} axioms) ---
verification run of 2026-07-22 on an isolated build host, against a
content-addressed checkout of the exact sources cited here --- the
earlier stale-run caveat (a 2026-07-18 script revision postdating the
prior census) is retired by this run, which attests the current text.
There is no signed HOL receipt store yet --- these citations are
documentary.
Semantic stratum, also ported: the soundness twins
\texttt{sound\_arithAtom\_fold} (theory \texttt{bern\_tr\_ready8}:
\texttt{hol-strict} clean, 8 theorems, 0 offences, verification run of
2026-07-22 on an isolated build host, content-addressed at the same
sources) and \texttt{sound\_arithAtom\_unfold} (theory
\texttt{bern\_tr\_fragclose}, 26 theorems, 0 offences, same run),
consumed by the certified-fragment keystone
(\texttt{checkPf\_sound\_frag}, theory \texttt{bern\_tr\_splice}, 20
theorems, 0 offences, same run): a reference-checker accept inside the
proven wire fragment yields sequent validity.

\sourcenote{HOL receipts are documentary: quoted results of \texttt{hol-strict} runs
and census rows from the HOL4 development\textquotesingle s build
evidence, not minted receipt artefacts. A HOL receipt store (the Lean
store\textquotesingle s twin) is named future work.
}
\ruleblock{Proof}

\sourcenote{Genre: core rule (tier 1). Setting as fixed in the
\rulename{EqElim} section. The rule is a pair: one payload constructor
per direction, sharing the comparator gate; both directions are proven,
and the two proofs are deliberately \emph{not} mirror images --- the
asymmetry is the interesting content.}

\textbf{Claim} (\texttt{sound\_arithAtom\_fold},
\texttt{sound\_arithAtom\_unfold}).\enspace The payload supplies terms
$t_1, t_2$ and a comparator $\bowtie \in \{<, \leq, =, \geq, >\}$ (the
gate $\bowtie \neq \mathord{\neq}$ is a checked field). The two atom
spellings in play, both recomputed syntactically by the checker from the
same payload terms, are the \emph{direct form} $t_1 \bowtie t_2$ and the
\emph{zero-normal form} $(t_1 - t_2) \bowtie 0$, where the subtrahend
side is the literal constant $0$ and the left side the literal syntactic
difference. \emph{Fold}: premise concludes the zero-normal form,
conclusion is the direct form. \emph{Unfold}: the converse. In both
directions, validity of the premise sequent gives validity of the
conclusion sequent, over the shared $\Phi$.

\begin{blemma}[scalar bridge]\label{lem:arith-scalar}
For $p, q \in \Q$ and $\bowtie$ any of the five comparators:
$(p - q) \bowtie 0 \iff p \bowtie q$.
\end{blemma}
\begin{proof}
Add $q$ to both sides, respectively subtract it --- the order of $\Q$ is
translation-invariant. The argument is uniform in $\bowtie$; it would,
for that matter, be equally valid for $\neq$ --- the comparator gate is
\emph{not} forced at this level (\autoref{rem:arith-ne}).
\end{proof}

\begin{blemma}[tensor bridge]\label{lem:arith-tensor}
Let $A, B$ be tensors whose elementwise difference is defined, say $D$.
Then $D$ stands in the elementwise lift of $\bowtie$ against the rank-0
tensor $0$ iff $A, B$ stand in the elementwise lift of $\bowtie$.
\end{blemma}
\begin{proof}
$D$'s entries are the pointwise differences under the broadcast
discipline: at equal shapes, entry $i$ of $D$ is $A_i - B_i$; with
rank-0 $A$, entry $i$ is $A_{\mathrm{scalar}} - B_i$; with rank-0 $B$,
$A_i - B_{\mathrm{scalar}}$ --- and $D$ takes the positive-rank
operand's shape (the broadcast join). The left side of the claimed
equivalence reads every entry of $D$ against the scalar $0$; the right
side reads the operands in the corresponding branch of the elementwise
lift. In each of the three branches the two readings pair off the same
scalar comparisons entry by entry, and Lemma~\ref{lem:arith-scalar}
closes each pair. (\texttt{QTensor.rel2\_sub\_zero\_iff}.)
\end{proof}

\begin{blemma}[definedness, fold direction]\label{lem:arith-def-fold}
If the zero-normal atom holds at $\sigma$, then both $t_1$ and $t_2$
evaluate at $\sigma$, and their elementwise difference is the evaluation
of $t_1 - t_2$.
\end{blemma}
\begin{proof}
The holding atom includes a defined evaluation of the syntactic
difference; the evaluator is structural and strict (the Setting: a
compound evaluates exactly when its parts evaluate and their shapes
join), so inverting it at the
subtraction node yields defined evaluations of both operands and the
elementwise-difference equation. (\texttt{Term.evalQ\_sub\_inv}.)
\end{proof}

\begin{blemma}[definedness, unfold direction]\label{lem:arith-def-unfold}
If the direct atom holds at $\sigma$, then the elementwise difference of
the two (defined, by the atom's holding) evaluations is itself defined.
\end{blemma}
\begin{proof}
A \emph{holding} elementwise comparison forces broadcast-compatible
shapes --- the lift's fourth branch, unequal positive ranks, is empty ---
and elementwise subtraction is total on exactly the broadcast-compatible
pairs. A holding comparison is thereby itself a compatibility
certificate. (\texttt{QTensor.map2\_defined\_of\_rel2}.)
\end{proof}

\begin{bthm}[soundness of \rulename{ArithAtom}, both directions]
\label{thm:arith-sound}
Over shared $\Phi$: validity of $\Phi \vdash (t_1 - t_2) \bowtie 0$
gives validity of $\Phi \vdash t_1 \bowtie t_2$, and conversely.
\end{bthm}
\begin{proof}
Fix an admissible $\sigma \models \Phi$ in either direction.
\emph{Fold}: the premise's zero-normal atom holds; its constant side
evaluates to the rank-0 tensor $0$; Lemma~\ref{lem:arith-def-fold}
recovers evaluations $A, B$ of $t_1, t_2$ with difference $D$, and the
premise's relation is $D \bowtie 0$ elementwise;
Lemma~\ref{lem:arith-tensor} converts it to $A \bowtie B$ elementwise
--- the direct atom, on the operands' own evaluations. \emph{Unfold}:
the premise's direct atom holds with evaluations $A, B$;
Lemma~\ref{lem:arith-def-unfold} manufactures the difference $D$; the
evaluator assembles $\mathrm{ev}_\sigma(t_1 - t_2) = D$ from its
operands' evaluations, the constant side is again rank-0 $0$, and
Lemma~\ref{lem:arith-tensor}, read right to left, closes the zero-normal
atom.
\end{proof}

\begin{bremark}[the direction asymmetry, and payload minimalism]
\label{rem:arith-asym}
The fold direction \emph{consumes} definedness of the difference (the
premise evaluated it); the unfold direction must \emph{manufacture} it
--- and can, precisely because a holding comparison certifies broadcast
compatibility (Lemma~\ref{lem:arith-def-unfold}). The consequence is
worth stating plainly: whenever either atom \emph{holds}, both operands'
evaluations exist and transport to the other form. That is why this
payload carries \emph{no} typing or
positivity fields at all --- compare \rulename{EqRefl}'s typing pack,
\rulename{EqElim}'s shape agreement, and the positive-size fields of the
combination rules: every one of those exists to manufacture or protect
an evaluation this rule gets handed for free. (Definedness itself is not
mirrored form to form: at broadcast-incompatible operand shapes the
direct form's sides still evaluate while the zero-normal difference does
not --- but there neither atom can hold, each failing in its own way,
and holding is all the rule transports; at a compatible zero-size shape
the two forms hold together, vacuously.) The ``exact rational evaluation
transported both directions'' is exactly the pair
(Lemma~\ref{lem:arith-scalar}, Lemma~\ref{lem:arith-tensor}): a
translation-invariance fact of ordered $\Q$, lifted entrywise through an
invertible totality argument.
\end{bremark}

\begin{bremark}[the syntactic pins]\label{rem:arith-pins}
Both form conditions are single syntactic equations recomputed from the
payload --- the zero-normal form demands the \emph{literal} constant $0$
and the \emph{literal} difference term of the same $t_1, t_2$, not any
term evaluating to zero or any algebraic rearrangement. A producer
wanting $t_1 \bowtie t_2$ from, say, $(t_2 - t_1) \cdot (-1) \bowtie 0$
routes through \rulename{Norm} first; this rule does arithmetic on
\emph{atom shape} only. That severity is what makes the side conditions
decidable by construction and keeps the rule's checked surface at two
form equations, the shared-context equation, and one comparator gate.
\end{bremark}

\begin{bremark}[why $\neq$ is excluded]\label{rem:arith-ne}
Not for the reason one might guess --- and not for its sibling's reason
either. Lemmas \ref{lem:arith-scalar} and
\ref{lem:arith-tensor} are comparator-uniform, and the rational
semantics of this section would accept $\neq$ without change. Nor would
the development's auxiliary three-valued screen (the countermodel
construction that polices vacuity properties, developed at the $\vee$E
section) break here: both spellings of a $\neq$ atom mention the same
wires, and the screen's disequality clause reads only designated-wire
membership, so fold and unfold preserve its value --- unlike at the
premise-free ground-fact rule, where a wire-free $\neq$ atom's screen
value genuinely breaks the invariant (the \rulename{PConstExact}
section). The
comparator set $\{<, \leq, =, \geq, >\}$ is the source calculus's own
set for this rule, and the payload's gate keeps it checked: the screen's
per-rule argument stays a one-clause conclusion-shape case (the gate
vacates its $\neq$ branch), and the two atom-shape rules stay aligned.
$\neq$ remains in the
formula grammar --- it is simply not this rule's business to move it
across the zero-normal boundary.
\end{bremark}

\begin{anchors}
Lemma \ref{lem:arith-scalar} & inlined in \texttt{QTensor.rel2\_sub\_zero\_iff}'s entrywise closure\\
Lemma \ref{lem:arith-tensor} & \texttt{QTensor.rel2\_sub\_zero\_iff}\\
Lemma \ref{lem:arith-def-fold} & \texttt{Term.evalQ\_sub\_inv}\\
Lemma \ref{lem:arith-def-unfold} & \texttt{QTensor.map2\_defined\_of\_rel2}\\
Theorem \ref{thm:arith-sound}, fold & \texttt{sound\_arithAtom\_fold}\\
Theorem \ref{thm:arith-sound}, unfold & \texttt{sound\_arithAtom\_unfold}; both wired as core \texttt{RuleData.arithAtom} into \texttt{ProofObject.sound}\\
form equations, comparator gate & \texttt{atomicCompareForm}, \texttt{atomicCompareZero}; payload field \texttt{notNE} (\texttt{PArithAtomData})\\
\end{anchors}


\rulesection{15}{BernCone}{core rule}
\label{blcatrule-BernCone}
\ruleblock{The rule (Steele metanotation)}
\sourcenote{as realized in the mechanized development; source row: \texttt{\texttt{rules.toml}} --- \emph{BernCone}. The metanotation statement is typeset by hand for notational quality; its mechanized ground truth is the Lean statement below, extracted from the sources at every build of this book and hash-bound. The rule \emph{roster} is machine-checked (every mechanized rule tag round-trips to exactly one catalogue row); the statement \emph{text} is editorial, checked by eye against the extracted Lean beside it.}
\begin{mathpar}
\inferrule*[Right=\rulename{BernCone}]{
  \Delta = \mathrm{conv}(V_0, \ldots, V_n) \subset \R^n \text{ an } n\text{-simplex};
  \quad S \subsetneq \{0, 1, \ldots, n\}
  \\
  \Phi \vdash x \in \mathcal{R}(\Delta, S)
  \\
  p \colon \R^n \to W \text{ polynomial map of degree } d, \text{ with coefficient tensor } (b_\alpha)_{|\alpha|=d}
  \\
  \wires(p) = \emptyset \;\vee\; \bigl(x = w \,\wedge\, \wires(p) \subseteq \{w\}\bigr)
  \\
  S \neq \emptyset \Rightarrow P \text{ conic}; \qquad S = \emptyset \Rightarrow P \text{ convex suffices}
  \\
  \bigl(p \ \text{idx-free, with } p \ \text{constant} \vee n = 1\bigr) \;\vee\; \bigl(p \ \text{projection-indexed, indices} < n,\ S = \emptyset\bigr)
  \\
  \text{for every } j \in \N^S \text{ with } |j| \leq d, \text{ every degree-}d \text{ Bernstein coefficient of } P_j \text{ on } F_S \text{ satisfies } P
}{
  \Phi \vdash P(p(x))
}
\end{mathpar}

\ruleflavor{the disjunctive premise on $p$ is the checker\textquotesingle s scope guard: coefficient checks on coordinate polynomials over ambient dimension $\neq 1$ are refused rather than attempted outside the proven kernels.}
\par\medskip
\rulename{BernCone} is the calculus\textquotesingle s arithmetic-closure
primitive: from a \emph{geometric} fact (a wire value lies in a region)
it derives an \emph{analytic} one of unbounded degree (a polynomial of
that value satisfies a predicate) --- where the conversion rules exchange
region atoms with their defining affine sign conditions,
\rulename{BernCone} is the rule that concludes something new. From a proven region membership
$\Phi \vdash x \in \mathcal{R}(\Delta, S)$ and a degree-$d$ polynomial
term $p$, it concludes $\Phi \vdash P(p(x))$ for any predicate $P$ of the
accepted grammar, provided every certified Bernstein coefficient of $p$
over the region satisfies $P$. A producer reaches for it wherever a
Lyapunov decrease, an invariant bound, or an output enclosure must be
established over a whole region rather than at sample points: choose a
degree, optionally subdivide (\rulename{Subdivide} is the completeness
partner --- enclosures converge under degree elevation and subdivision),
and emit the payload.

What the checker recomputes, given the payload (simplex, extension set,
degree, polynomial term, predicate): that the extension set is tied to
the premise\textquotesingle s simplex and the simplex is
full-dimensional; that the polynomial term translates into the
region\textquotesingle s ambient coordinates (two routes: an idx-free
single-wire canonical form, or the projection-indexed coordinate form);
and then --- the heart --- \emph{every} forward-difference Bernstein
coefficient of the translated polynomial over the region, by exact
$\Q$ arithmetic, testing each against the predicate. No coefficient
tensors ride the wire: the certificate names the polynomial and the
region, and the checker rebuilds the entire coefficient table itself
(the recompute pattern), so a forged table is unrepresentable rather
than detected.

Each hypothesis is load-bearing. Drop the wire tie
(\texttt{polyTie}) and a certificate could prove its region fact about
one wire while running the coefficient check on a polynomial over a
different one --- concluding a bound about a value the premise never
constrained. Drop the predicate-shape clause and a convex-but-not-conic
predicate falsely transfers along an extension ray (the proof section
computes a two-line counterexample: certified coefficients $\{0, -1\}$
both satisfy $\square \geq -1$, yet the polynomial dips to $-2$ on the
ray). Drop the coefficient recomputation and the conclusion is unrelated
to the polynomial. Drop the simplex tie or full-dimensionality and the
barycentric reading of the premise breaks. The checker\textquotesingle s
own scope guard (constant canonical forms at any dimension; coordinate
polynomials only over 1-dimensional regions when the extension set is
nonempty) is part of the soundness story: everything beyond the proven
fragment is refused, never accepted on faith.

\cataloguenote{Predicates enter through a deliberately closed-world grammar: finite
conjunctions of scalar hole-atoms $s \cdot \square \bowtie c$ with
$s > 0$ and $\bowtie \in \{\leq, \geq, =\}$, classified for convexity
(always) and conicity ($c = 0$ throughout). Comparators outside the
grammar are refused --- conservatism that can lose completeness, never
soundness. The elementwise atom semantics batch the conclusion: a
tensor-shaped polynomial term gives per-entry predicate application for
free.
}
\ruleblock{Lean statement}
\declhead{BernsteinLogic.Soundness}{sound\_bernCone\_empty}{\badge{BadgeGreen}{statement bound}}
\begin{leansrc}
/-- BernCone at S = ∅ — the kernel assembly, at full D2 scope (leg D3:
the former `hD2` corner guard is gone — a polynomial whose wires cancel
in canonical form over a higher-dimensional region variable is handled
by the per-entry evaluation lemma `Term.evalQ_brRead`: every entry of
its semantic tensor equals the closed canonical value, which the kernel
certifies). -/
theorem sound_bernCone_empty {H : Hypernet} {seq regionPremise : Sequent}
    {termX : Term} {simplex : Simplex} {S : ExtensionSet} {degree : Nat}
    {polynomial : Term} {predicate : Formula}
    (ih : regionPremise.semanticallyValid H)
    (hpre : seq.pre = regionPremise.pre)
    (hregionPost : regionPremise.post = .inRegion termX simplex S)
    (hpolyTie : polynomial.freeWires = [] ∨
      ∃ w, termX = Term.var w ∧ polynomial.freeWires ⊆ [w])
    (hconvex : isExtensionSetEmpty S → isConvexPredicate predicate)
    (hSempty : isExtensionSetEmpty S)
    (hcoeff : allForwardDiffBernCoeffsSatisfyP simplex S degree polynomial
      predicate)
    (hpostApply : applyPredicateToTerm predicate polynomial seq.post) :
    seq.semanticallyValid H
\end{leansrc}
\declhead{BernsteinLogic.Soundness}{sound\_bernCone}{\badge{BadgeGreen}{statement bound}}
\begin{leansrc}
/-- BernCone, all extension sets (leg D3). At S = ∅ this is the
convex-predicate kernel assembly above; at S ≠ ∅ the check's own D2
scope guard confines reachable instances to constant canonical
polynomials (any dimension — `fdBernstein_sound_const`) and single-wire
polynomials over 1-dimensional regions, where a proper nonempty
extension set is exactly one of the two vertices and the region is the
extension ray (`fdBernstein_sound_dim1`, the univariate (★) identity).
The conic predicate-shape clause supplies zero/scaling/convexity
closure. -/
theorem sound_bernCone {H : Hypernet} {seq regionPremise : Sequent}
    {termX : Term} {simplex : Simplex} {S : ExtensionSet} {degree : Nat}
    {polynomial : Term} {predicate : Formula}
    (ih : regionPremise.semanticallyValid H)
    (hpre : seq.pre = regionPremise.pre)
    (hregionPost : regionPremise.post = .inRegion termX simplex S)
    (hpolyTie : polynomial.freeWires = [] ∨
      ∃ w, termX = Term.var w ∧ polynomial.freeWires ⊆ [w])
    (hpredShape : (¬ isExtensionSetEmpty S → isConicPredicate predicate)
      ∧ (isExtensionSetEmpty S → isConvexPredicate predicate))
    (hcoeff : allForwardDiffBernCoeffsSatisfyP simplex S degree polynomial
      predicate)
    (hpostApply : applyPredicateToTerm predicate polynomial seq.post) :
    seq.semanticallyValid H
\end{leansrc}
\par\medskip
\textbf{What the coefficient side-condition means --- and what it does
not.}\enspace The rule\textquotesingle s main premise appears in the
payload as the single name \texttt{allForwardDiffBernCoeffsSatisfyP}, so
this book binds its meaning by displaying the definition chain in full,
from the checker sources. Reading the chain: \texttt{satisfiesPredQ}
substitutes a candidate rational for the designated predicate-hole wire
and runs the \emph{closed} formula evaluator --- the predicate is judged
at a concrete value, by exact $\Q$ arithmetic, never at a wire.
\texttt{Term.asCoordPoly?} translates the polynomial term into the
region\textquotesingle s ambient coordinates, by one of two routes: an
idx-free term goes through its wire-keyed canonical polynomial (at most
one wire, remapped to coordinate 0), and a projection-indexed term maps
each projection $w[j]$ to coordinate variable $j$.
\texttt{bernConeCoeffCheck} then gates (simplex tie,
full-dimensionality, and the scope guard: non-constant idx-free
polynomials only over 1-dimensional regions; multivariate coordinate
forms only at $S = \emptyset$ with all projection indices in range),
computes the full forward-difference Bernstein coefficient table of the
translated polynomial over the masked vertex list, and sweeps every
grid entry through the predicate.

What a skeptical reader should confirm by eye in the definitions below:
the conclusion formula never enters the check --- only the polynomial,
the region data, and the predicate do (the conclusion is pinned
separately, as the syntactic substitution of the polynomial into the
predicate\textquotesingle s hole); nothing in the chain reads a
valuation, so the check cannot covertly assume anything about wire
values; and every \texttt{false} branch of the guards is a refusal ---
the fragment boundary (dimension, extension set, index range) fails
closed. Semantic truth is established only in \texttt{sound\_bernCone},
from the region premise\textquotesingle s validity.

\declhead{BernsteinLogic.Core.Catalogue}{satisfiesPredQ}{\badge{BadgeGreen}{definition bound}}
\begin{leansrc}
/-- Predicate satisfaction at a rational value. -/
def satisfiesPredQ (P : Formula) (c : Rat) : Bool :=
  (Formula.subst predicateHole (.const c) P).evalClosed = some true
\end{leansrc}
\declhead{BernsteinLogic.Core.Catalogue}{Term.asCoordPolyIdx?}{\badge{BadgeGreen}{definition bound}}
\begin{leansrc}
/-- The G1 coordinate-form translation: a Term whose every wire mention
is an idx projection `.idx (.var w) j` becomes a polynomial over the
ambient-coordinate variables j. Which wire is projected is deliberately
NOT recorded here — the BernCone payload's `polyTie` (B10/B7) pins every
free wire to the region variable. `none` on: a bare (vector-valued) wire
mention, and on idx of any non-var argument (v0.5 scope; einsum is the
planned v1+ extension). Index range is checked separately
(`Term.idxBounded`, conjoined by `bernConeCoeffCheck`). -/
def Term.asCoordPolyIdx? : Term → Option QArith.Poly
  | .var _ => none
  | .const c => some (QArith.Poly.const c)
  | .add l r =>
    do pure (QArith.Poly.add (← l.asCoordPolyIdx?) (← r.asCoordPolyIdx?))
  | .mul l r =>
    do pure (QArith.Poly.mul (← l.asCoordPolyIdx?) (← r.asCoordPolyIdx?))
  | .sub l r =>
    do pure (QArith.Poly.sub (← l.asCoordPolyIdx?) (← r.asCoordPolyIdx?))
  | .neg t => (Term.asCoordPolyIdx? t).map QArith.Poly.neg
  | .idx (.var _) i => some (QArith.Poly.var i)
  | .idx _ _ => none
\end{leansrc}
\declhead{BernsteinLogic.Core.Catalogue}{Term.asCoordPoly?}{\badge{BadgeGreen}{definition bound}}
\begin{leansrc}
def Term.asCoordPoly? (t : Term) : Option QArith.Poly :=
  if t.idxFree then
    let p := t.toPoly
    match QArith.Poly.vars p with
    | [] => some p
    | [_] => some (p.map fun (m, c) => (m.map fun (_, e) => (0, e), c))
    | _ => none
  else
    t.asCoordPolyIdx?
\end{leansrc}
\declhead{BernsteinLogic.Core.Catalogue}{bernConeCoeffCheck}{\badge{BadgeGreen}{definition bound}}
\begin{leansrc}
/-- The recomputation behind BernCone's main side-condition (rem:recompute-
vs-trust option (a), CORRESPONDENCE.md D-6; see original for the guard
rationale). -/
def bernConeCoeffCheck (simplex : Simplex) (S : ExtensionSet) (degree : Nat)
    (polynomial : Term) (predicate : Formula) : Bool :=
  if S.simplex = simplex then
    let cx := simplex.complex
    let verts := (List.finRange cx.nVertices).filter (fun v => simplex.contains v)
    -- full-dimensionality: an (ambDim)-simplex needs ambDim+1 vertices
    if verts.length = cx.ambDim + 1 then
      match Term.asCoordPoly? polynomial with
      | none => false
      | some p =>
        if polynomial.idxFree then
          -- legacy D2 scope: coordinate-referencing polynomials only
          -- over 1-D regions (any extension set)
          if QArith.Poly.vars p ≠ [] ∧ cx.ambDim ≠ 1 then false
          else
            let coords := verts.map fun v =>
              (List.finRange cx.ambDim).map fun j => cx.coord v j
            let smask := verts.map fun v => S.containsNat v.val
            match QArith.fdBernstein coords smask degree p with
            | some table => QArith.gridAll (satisfiesPredQ predicate) table
            | none => false
        else
          -- G1 idx coordinate form (multivariate): S = ∅ only — the
          -- general-S multivariate soundness is the v1 StarIdentity
          -- obligation — and every projection index below the ambient
          -- dimension (out-of-range reads make the conclusion
          -- unevaluable; see `Term.idxBounded`).
          if S.isEmptyB ∧ polynomial.idxBounded cx.ambDim then
            let coords := verts.map fun v =>
              (List.finRange cx.ambDim).map fun j => cx.coord v j
            let smask := verts.map fun v => S.containsNat v.val
            match QArith.fdBernstein coords smask degree p with
            | some table => QArith.gridAll (satisfiesPredQ predicate) table
            | none => false
          else false
    else false
  else false
\end{leansrc}
\declhead{BernsteinLogic.Core.Catalogue}{allForwardDiffBernCoeffsSatisfyP}{\badge{BadgeGreen}{definition bound}}
\begin{leansrc}
def allForwardDiffBernCoeffsSatisfyP
    (simplex : Simplex) (S : ExtensionSet) (degree : Nat)
    (polynomial : Term) (predicate : Formula) : Prop :=
  bernConeCoeffCheck simplex S degree polynomial predicate = true
\end{leansrc}
\ruleblock{Correspondence (Steele display $\leftrightarrow$ Lean)}
\badge{BadgeAmber}{a bound Lean statement CHANGED since certification --- re-review and re-certify}
\sourcenote{The dictionary below is a reviewer-certified statement of how the mechanized statements above interpret the displayed rule: display metavariable $\mapsto$ Lean binder, each display hypothesis $\mapsto$ the premise or side condition carrying it, the display conclusion $\mapsto$ the goal --- with any deliberate difference in width between display and mechanization named outright. Certification is recorded as a hash triple over (the displayed cell, the extracted statement texts, this dictionary\textquotesingle s own text); if any of the three changes, this block turns amber until a reviewer re-certifies. The correspondence claim itself remains editorial --- reviewer-grade, not machine-checked; its machine-checked successor types each soundness statement against the rule as a structured object, and supersedes this block per rule as it lands.}
Judgment frame as in \rulename{Ax}; sub-derivation premise as at
\rulename{Weaken}. The general statement is \texttt{sound\_bernCone};
\texttt{sound\_bernCone\_empty} (certified alongside) is its
$S = \emptyset$ specialization, through which the general proof routes
that case.
\begin{itemize}
\item $x \mapsto$ \texttt{termX : Term}; $\Delta \mapsto$
  \texttt{simplex}; $S \mapsto$ \texttt{S : ExtensionSet}; $d \mapsto$
  \texttt{degree : Nat}; $p \mapsto$ \texttt{polynomial : Term}; $P
  \mapsto$ \texttt{predicate : Formula} (the hole-atom grammar of the
  note, its hole a designated wire).
\item the region premise $\Phi \vdash x \in \mathcal{R}(\Delta, S)$
  $\mapsto$ \texttt{ih} with \texttt{hregionPost} $:$
  \texttt{regionPremise.post = .inRegion termX simplex S}; $\Phi$
  shared via \texttt{hpre}.
\item the wire-tie premise $\wires(p) = \emptyset \vee (x = w \wedge
  \wires(p) \subseteq \{w\})$ $\mapsto$ \texttt{hpolyTie}, the same
  disjunction verbatim --- note \texttt{termX = Term.var w}: the
  display\textquotesingle s ``$x = w$\textquotesingle\textquotesingle{}
  means the region term literally \emph{is} that wire.
\item the two-clause predicate premise $\mapsto$ \texttt{hpredShape}
  $:$ the conjunction of ($S$ nonempty $\to$ conic) and ($S$ empty
  $\to$ convex), by the decidable predicate classifiers.
\item the displayed coefficient sweep $\mapsto$ \texttt{hcoeff} $:$
  \texttt{allForwardDiffBernCoeffsSatisfyP simplex S degree polynomial
  predicate} --- one name whose meaning this book binds by displaying
  its definition chain in full (the intro above); the
  display\textquotesingle s first premise line (simplex tie,
  full-dimensionality) and its scope-guard premise have no Lean binder
  of their own because the checker\textquotesingle s gates certify
  them: \emph{absorbed}, as \rulename{PConstExact}\textquotesingle s
  closedness is.
\item conclusion $\Phi \vdash P(p(x))$ $\mapsto$ the goal under
  \texttt{hpostApply} $:$ \texttt{applyPredicateToTerm predicate
  polynomial seq.post} --- the conclusion is the syntactic
  substitution of the polynomial term into the
  predicate\textquotesingle s hole; the ``at
  $x$\textquotesingle\textquotesingle{} lives in the wire tie, not in
  a separate application node.
\end{itemize}
Width note, two-sided and honest: the statement is \emph{proven} at
every hypernet, sequent, and payload --- and \texttt{hcoeff} is
satisfiable only inside the displayed scope-guard premise (constant
canonical forms at any dimension; non-constant
idx-free forms over 1-dimensional regions; multivariate coordinate
forms at $S = \emptyset$ only), so the certified reach of the displayed
rule is exactly its guard-passing fragment, by refusal outside it ---
never by unverified acceptance. Conversely nothing in either statement
is wider than the display: every hypothesis is consumed, and the
display\textquotesingle s own two-clause predicate license is carried
exactly.

\ruleblock{Lean verifier receipts}
\sourcenote{module \texttt{BernsteinLogic.Soundness}; bound build \texttt{a20d90e66c13\ldots} (bound from receipt a9c0e34d3e3c9c8c, 2026-08-03T21:30:36Z). A receipt records that named checking layers accepted a build: (module, compiled-artifact SHA-256, environment SHA-256, checker layers), minted when the checking gate ran; the store receipts attesting the bound build are listed with the layers each run credited. Nothing here asks the reader to trust the store: every layer is an independently runnable program over the quoted sources --- rebuild the module and re-run the checkers to reproduce the verdict. Receipts attesting superseded builds credit no layer and are omitted from this rendering (the count of omissions, if any, is noted below).}
\begin{receipttable}
\receiptrow{a9c0e34d3e3c9c8c}{2026-08-03T21:30:36Z}{build+decorrelated}{\badge{BadgeGreen}{current}}
\end{receipttable}
\sourcenote{21 further receipt(s) in the store attest superseded builds of this module and are omitted from this rendering.}
\sourcenote{module \texttt{BernsteinLogic.Core.Catalogue}; bound build \texttt{cbc92e976a1d\ldots} (bound from receipt 5abdb3589707b5df, 2026-08-03T21:30:33Z). A receipt records that named checking layers accepted a build: (module, compiled-artifact SHA-256, environment SHA-256, checker layers), minted when the checking gate ran; the store receipts attesting the bound build are listed with the layers each run credited. Nothing here asks the reader to trust the store: every layer is an independently runnable program over the quoted sources --- rebuild the module and re-run the checkers to reproduce the verdict. Receipts attesting superseded builds credit no layer and are omitted from this rendering (the count of omissions, if any, is noted below).}
\begin{receipttable}
\receiptrow{24804e66edba0bc9}{2026-08-03T20:25:09Z}{build+decorrelated}{\badge{BadgeGreen}{current}}
\receiptrow{5abdb3589707b5df}{2026-08-03T21:30:33Z}{build+decorrelated}{\badge{BadgeGreen}{current}}
\end{receipttable}
\sourcenote{18 further receipt(s) in the store attest superseded builds of this module and are omitted from this rendering.}
\layerline{\layeritem{build}{lake build (elaboration)}{\badge{BadgeGreen}{current}}\layeritem{replay}{leanchecker (official kernel replay)}{\badge{BadgeAmber}{stale only}}\layeritem{hardened}{safe\_verify (hardened screen)}{\badge{BadgeAmber}{stale only}}\layeritem{decorrelated}{lean4export + nanoda (decorrelated kernel)}{\badge{BadgeGreen}{current}}}
\sourcenote{\badge{BadgeGreen}{comparator: fixture DemoBernCone} the
three-implementation comparator harness (compiled machine code,
kernel-reduction probes, and the compiler-IR interpreter, all executing
the same elaborated checker functions) carries the worked
example\textquotesingle s certificate --- an \rulename{Ax}
$\to$ $\wedge$E $\to$ \rulename{BernCone} chain at $S = \emptyset$ ---
as its fixture \texttt{DemoBernCone}. The suite passes only when all
three paths return identical verdicts on it, and the kernel path probes
two-sided (reporting its own verdict rather than confirming the compiled
one), so a false green would require the same wrong answer from
execution stacks that share nothing below the elaborated definitions.
}
\ruleblock{HOL4 statement}
\sourcenote{theory \texttt{bern\_catalogue}, HOL4 development script \texttt{\texttt{bern\_catalogueScript.sml}} \badge{BadgeGreen}{script bound}}
\holhead{Definition bernConeCoeffCheck\_def}
\begin{holsrc}
Definition bernConeCoeffCheck_def:
  bernConeCoeffCheck sx es dg pl pr =
    if es.simplex = sx then
      (let cx = sx.complex in
       let verts = FILTER (λv. simplexContainsNat sx v) (rangeL cx.nVertices) in
       if LENGTH verts = cx.ambDim + 1 then
         (case termAsCoordPoly pl of
          | NONE => F
          | SOME p =>
            if termIdxFree pl then
              (if polyVars p ≠ [] ∧ cx.ambDim ≠ 1 then F
               else
                 let coords = MAP (λv. MAP (λj. getD (getD cx.coord v []) j (0:rat))
                                           (rangeL cx.ambDim)) verts in
                 let smask = MAP (λv. extContainsNat es v) verts in
                 case fdBernstein coords smask dg p of
                 | SOME table => gridAll (satisfiesPredQ pr) table
                 | NONE => F)
            else
              (if extIsEmptyB es ∧ termIdxBounded cx.ambDim pl then
                 (let coords = MAP (λv. MAP (λj. getD (getD cx.coord v []) j (0:rat))
                                            (rangeL cx.ambDim)) verts in
                  let smask = MAP (λv. extContainsNat es v) verts in
                  case fdBernstein coords smask dg p of
                  | SOME table => gridAll (satisfiesPredQ pr) table
                  | NONE => F)
               else F))
       else F)
    else F
End
\end{holsrc}
\holhead{Definition bernConeIntvCheck\_def}
\begin{holsrc}
Definition bernConeIntvCheck_def:
  bernConeIntvCheck sx es dg pl pr dt =
    if es.simplex = sx then
      (let cx = sx.complex in
       let verts = FILTER (λv. simplexContainsNat sx v) (rangeL cx.nVertices) in
       if LENGTH verts = cx.ambDim + 1 then
         (case termAsCoordPoly pl of
          | NONE => F
          | SOME p =>
            if ¬ termIdxFree pl then F
            else if polyVars p ≠ [] ∧ cx.ambDim ≠ 1 then F
            else
              let coords = MAP (λv. MAP (λj. getD (getD cx.coord v []) j (0:rat))
                                        (rangeL cx.ambDim)) verts in
              let smask = MAP (λv. extContainsNat es v) verts in
              case fdBernstein coords smask dg p of
              | SOME table =>
                  gridAll (λc. case dtypeRoundE dt Rdn c of
                               | Fin lo => (lo ≤ c ∧ satisfiesPredQ pr lo)
                               | _ => F) table
              | NONE => F)
       else F)
    else F
End
\end{holsrc}
\ruleblock{HOL verifier receipts}
\badge{BadgeBlue}{documentary} \texttt{hol-strict} CLEAN over the ported checker chain: theory
\texttt{bern\_catalogue}, 472 theorems, 0 cheat-tagged, 0 nonstandard
axioms (whitelist: \texttt{DISK\_THM}; the four \texttt{bool} axioms) ---
verification run of 2026-07-22 on an isolated build host, scripts
sha-matched verbatim to the sources cited here (git blob
\texttt{e41b9ed0\ldots}). The ported receipt table replays the
\rulename{BernCone} rows of the Lean development\textquotesingle s own
receipt battery: the worked-example accept
(\texttt{bernConeCoeffCheck} at $S = \emptyset$, degree 1) plus four
mutant rejections (tightened goal, wrong dynamics sign, degree lie,
sub-dimensional simplex) and the nonconvex-predicate rejection.
Downstream, the same chain carries the checker to a CakeML binary with
the \texttt{machine\_code\_sound} tag (theory
\texttt{bernCertProofARM8}). \texttt{bernConeIntvCheck} is the ported
companion of a since-retired validated-numerics rule variant; its
round-down-then-check strategy is destined to fold into the BernCone
checker as an internal fast path (evaluation strategy, not proof
content), and the ported definition with its three receipt rows
(exact-dyadic accept, rounded-bound accept, non-representable-bound
reject) is retained as that branch\textquotesingle s reference.
Semantic stratum, also ported: the soundness twins
\texttt{sound\_bernCone} and its $S = \emptyset$ specialization
\texttt{sound\_bernCone\_empty} (theory \texttt{bern\_tr\_berncone}:
\texttt{hol-strict} clean, 14 theorems, 0 offences, verification run of
2026-07-22 on an isolated build host, content-addressed at the same
sources), consumed by the certified-fragment keystone
(\texttt{checkPf\_sound\_frag}, theory \texttt{bern\_tr\_splice}, 20
theorems, 0 offences, same run) under the payload\textquotesingle s
decode-time gates: a reference-checker accept inside the proven wire
fragment yields sequent validity.

\sourcenote{HOL receipts are documentary: quoted results of \texttt{hol-strict} runs
and census rows from the HOL4 development\textquotesingle s build
evidence, not minted receipt artefacts. A HOL receipt store (the Lean
store\textquotesingle s twin) is named future work.
}
\ruleblock{Proof}

\sourcenote{Genre: core rule (tier 1) --- an ordinary mathematical proof at
Putnam-exam-solution detail. Every hypothesis of the rule is used below and
remarked on, with the counterexample its removal admits where one exists;
the formal-anchor table maps each step to the Lean declaration that
carries it.}

\textbf{Setting and claim.}\enspace Fix a hypernet $H$. A valuation
$\sigma$ assigns each wire a rational tensor of its declared shape; a
sequent $\Phi \vdash \Psi$ is semantically valid when every admissible
$\sigma$ satisfying $\Phi$ satisfies $\Psi$, with comparison atoms read
per-entry. A region atom $x \in \mathcal{R}(\Delta, S)$, for
$\Delta = \mathrm{conv}(V_0, \dots, V_n)$ an $n$-simplex and
$S \subsetneq \{0, \dots, n\}$ a proper set of \emph{extended} vertices,
holds at $\sigma$ when the evaluated point admits barycentric weights
$\lambda_v$ with $\sum_v \lambda_v = 1$, $\lambda_v \geq 0$ at core
vertices $v \notin S$, and $\lambda_v \leq 0$ at $v \in S$ --- the closed
simplex when $S = \emptyset$, and the unbounded set extending beyond the
faces opposite $S$ otherwise. (Field bookkeeping, fixed once: the
vertices $V_0, \dots, V_n$ are exact rational data of the payload; the
weights $\lambda_v$ are quantified over $\R$; and the evaluated point is
rational, valuations assigning rational tensors --- Step 2 of the
theorem is where these three facts meet.) The rule's payload supplies the region
premise $\Phi \vdash x \in \mathcal{R}(\Delta, S)$, a polynomial term $p$
of degree $\leq d$, a predicate $P$ from the accepted grammar (a finite
conjunction of scalar hole-atoms $s \cdot \square \bowtie c$, $s > 0$,
$\bowtie \in \{\leq, \geq, =\}$), the checked wire tie \texttt{polyTie},
the predicate-shape clause ($P$ conic when $S \neq \emptyset$; convex
suffices when $S = \emptyset$), and the checked coefficient condition:
every forward-difference Bernstein coefficient of $p$ over
$\mathcal{R}(\Delta, S)$, recomputed exactly, satisfies $P$. For
$S \neq \emptyset$ the checker's scope guard is part of the rule's
acceptance: certificates are confined to constant canonical forms (any
dimension) or single-wire polynomials over a $1$-dimensional region
(\autoref{rem:bc-guard}) --- ``accepted'' below always means
guard-passing. The claim
(\texttt{sound\_bernCone}): the composite sequent
$\Phi \vdash P(p(x))$ is then semantically valid.

\begin{blemma}[Bernstein form on the closed simplex]\label{lem:bc-bern}
Let $p$ be a polynomial of degree $\leq d$ on $\R^n$ with degree-$d$
Bernstein coefficients $(b_\alpha)_{|\alpha| = d}$ on $\Delta$, so that
$p(x) = \sum_{|\alpha| = d} b_\alpha\, B^d_\alpha(\lambda)$ with
$B^d_\alpha(\lambda) = \binom{d}{\alpha} \lambda^\alpha$ in the
barycentric coordinates $\lambda$ of $x$. For $x \in \Delta$ (all
$\lambda_v \geq 0$),
\[
B^d_\alpha(\lambda) \geq 0
\qquad\text{and}\qquad
\sum_{|\alpha| = d} B^d_\alpha(\lambda)
  = \Bigl(\sum_v \lambda_v\Bigr)^{\!d} = 1,
\]
so $p(x)$ is a \emph{convex combination} of the coefficients $b_\alpha$.
\end{blemma}
\begin{proof}
Nonnegativity is termwise ($\lambda^\alpha \geq 0$); the partition of
unity is the multinomial theorem applied to
$(\lambda_0 + \dots + \lambda_n)^d = 1^d$. The representation itself is
the classical change to the Bernstein basis
\cite{bernstein1912,farouki2012}; the checker \emph{computes} the
$b_\alpha$ from $p$'s canonical monomial form by exact forward
differencing over $\Q$, so the representation is an identity of exact
rational data, not an approximation. That Bernstein coefficients
enclose the range of $p$ on $\Delta$ --- the reading of this lemma used
below --- is the sound (easy) direction of \cite{cargoshisha1966}.
\end{proof}

\begin{blemma}[accepted predicates are convex]\label{lem:bc-convex}
Let $P$ be accepted by the predicate classifier. Then its satisfaction
set $\{ v \in \Q : \mathrm{satisfiesPredQ}(P, v) \}$ is convex.
\end{blemma}
\begin{proof}
Each accepted atom orients (possibly by mirroring) to
$s \cdot v \bowtie c$ with $s > 0$ and
$\bowtie \in \{\leq, \geq, =\}$; its satisfaction set is a closed ray or
a point --- an interval either way --- and division by the positive $s$
does not disturb this. A conjunction intersects finitely many intervals.
Strict comparators and $\neq$ are \emph{refused} by the classifier ---
closed-world conservatism that can lose completeness but never
soundness; the refusal is a checker verdict, not a soundness case.
\end{proof}

\begin{blemma}[the univariate ray identity]\label{lem:bc-star}
Let $d \geq 0$ and let $b_0, \dots, b_d$ be the degree-$d$ Bernstein
coefficients of $p$ on the segment $[V_a, V_b]$, indexed by the
multiplicity of $V_b$ (labels chosen so that the extended vertex is
$V_a$ and the core vertex is $V_b$; the mirrored case relabels ---
concretely, with the other vertex extended, apply the identity to the
reversed sequence $c_i := b_{d-i}$, whose differences against the
un-reversed one are the signed
$(\Delta^{j} c)_{d-j} = (-1)^{j} (\Delta^{j} b)_{0}$; the checked table
is computed relative to the extension mask, so its rows \emph{are} the
$(\Delta^{j} c)_{d-j}$ of the relabeled orientation). Let
$x$ lie on the extension ray beyond $V_b$:
in barycentric coordinates $x = \lambda_a V_a + \lambda_b V_b$ with
$\lambda_a = -u \leq 0$, $\lambda_b = 1 + u$, $u \geq 0$. Then, with
$(\Delta b)_i := b_{i+1} - b_i$ the forward difference,
\[
p(x) \;=\; \sum_{j=0}^{d} \binom{d}{j}\, u^{j}\, (\Delta^{j} b)_{d-j},
\tag{$\star$}
\]
a combination of the certified table values $(\Delta^{j} b)_{d-j}$ with
\emph{nonnegative} weights $\binom{d}{j} u^{j}$, the $j = 0$ weight
being $1$.
\end{blemma}
\begin{proof}
Substitute $\lambda_a = -u$, $\lambda_b = 1 + u$ into the Bernstein
form and expand $(1+u)^i$ binomially:
\[
p(x) = \sum_{i=0}^{d} \binom{d}{i} b_i\, (-u)^{d-i} (1+u)^{i}
     = \sum_{i=0}^{d} \sum_{k=0}^{i} \binom{d}{i} \binom{i}{k}
       (-1)^{d-i}\, b_i\, u^{d-i+k}.
\]
Collect the coefficient of $u^{j}$: setting $i = d - j + m$ with
$m = 0, \dots, j$, it is
$\sum_{m=0}^{j} \binom{d}{d-j+m} \binom{d-j+m}{m} (-1)^{j-m}
b_{d-j+m}$. The product of binomials telescopes,
\[
\binom{d}{d-j+m} \binom{d-j+m}{m}
  = \frac{d!}{(j-m)!\, m!\, (d-j)!}
  = \binom{d}{j} \binom{j}{m},
\]
so the coefficient of $u^{j}$ is
$\binom{d}{j} \sum_{m=0}^{j} \binom{j}{m} (-1)^{j-m} b_{d-j+m}
 = \binom{d}{j} (\Delta^{j} b)_{d-j}$, which is $(\star)$. (Degree-$1$
sanity check: $p(x) = b_1 + u\,(b_1 - b_0)$, matching
$(1+u)\,b_1 - u\,b_0$ directly.) The values $(\Delta^{j} b)_{d-j}$ are
exactly the rows the checker's forward-difference table certifies on
this fragment: the $j$-th difference of the coefficient sequence,
anchored at the core vertex --- in the rule display's notation, the
degree-$d$ Bernstein coefficients of the difference polynomials $P_j$
on the core face $F_S$, specialized to the one-dimensional and constant
shapes the scope guard admits.
\end{proof}

\begin{blemma}[conic transfer]\label{lem:bc-conic}
Let $P$ be accepted with the conic classification (every atom's
constant is $0$). Then its satisfaction set contains $0$, and is closed
under nonnegative scaling and under finite nonnegative combinations.
\end{blemma}
\begin{proof}
Each conic atom reads $s \cdot v \bowtie 0$, satisfied at $v = 0$ and
preserved under $v \mapsto t v$ for $t \geq 0$; conjunction preserves
both facts. Closure under addition follows from convexity
(Lemma~\ref{lem:bc-convex}) plus scaling:
$v + w = 2 \cdot \tfrac{1}{2}(v + w)$; induction extends to finite
nonnegative combinations.
\end{proof}

\begin{bthm}[soundness of \rulename{BernCone}]\label{thm:bc-sound}
Under the payload's checked hypotheses, if the region premise is
semantically valid, then so is $\Phi \vdash P(p(x))$.
\end{bthm}
\begin{proof}
Fix an admissible $\sigma \models \Phi$; the region premise yields
barycentric weights for the evaluated point of $x$ as in the Setting.
The conclusion formula is, by the checked \texttt{postApply}, the
predicate $P$ with the polynomial term substituted for its hole, so it
suffices to establish $\mathrm{satisfiesPredQ}(P, \cdot)$ at the
evaluated value(s) of $p$ --- per-entry, when $p$ is tensor-shaped.

\emph{Step 1 (the polynomial reads one point).} The wire tie
\texttt{polyTie} says $p$ is closed, or mentions exactly the region
wire $w$ with $x = w$. In the second case the region premise pins
$\sigma$'s read of $w$ to a point of $\mathcal{R}(\Delta, S)$; in the
first, no wire is read at all and $p$'s value is its closed constant.

\emph{Step 2 (rational weights).} The model's barycentric weights are
\emph{real}; the kernel computes over $\Q$. The weights solve a linear
system with rational data (the vertex coordinates and the evaluated
point) under the region atom's sign pattern --- $\lambda_v \ge 0$ at
core vertices, $\lambda_v \le 0$ at extended ones. Substituting
$\mu_v := -\lambda_v$ at the extended vertices turns the pattern into a
nonnegative-orthant condition on a rational polyhedron, so a real
solution can be traded for a rational one with the same signs, and the
exact kernel consumes the rational weights. This
step is invisible in informal treatments and load-bearing in the
mechanization.

\emph{Step 3 (case $S = \emptyset$: the convex kernel).} By
Lemma~\ref{lem:bc-bern}, $p$'s value is a convex combination of the
certified coefficients; every coefficient satisfies $P$ (the checked
table sweep); by Lemma~\ref{lem:bc-convex} satisfaction transfers to
the combination. On this case the checker admits both translation
routes --- the idx-free single-wire form and the multivariate
projection-indexed form with all indices in range --- and the per-entry
evaluation lemma carries the scalar argument to every entry of a
tensor-shaped value.

\emph{Step 4 (case $S \neq \emptyset$).} The checker's scope guard
confines accepted certificates to two shapes. \emph{(4a: constant
canonical form, any dimension.)} If $p$'s canonical polynomial has no
variables, its value everywhere is the constant --- the zero-order
forward-difference row is that constant, certified directly.
\emph{(4b: single-wire polynomial over a 1-dimensional region.)} A
proper nonempty extension set on a 1-simplex extends exactly one of the
two vertices, and $\mathcal{R}(\Delta, S)$ is the closed ray from the
core vertex; Lemma~\ref{lem:bc-star} writes $p$'s value as a
nonnegative combination of the certified table values with the $j = 0$
weight equal to $1$, and Lemma~\ref{lem:bc-conic} transfers
satisfaction through the combination. Everything else is
\emph{refused} by the guard, so no further case exists for soundness
to cover.

\emph{Step 5 (read off).} In each case
$\mathrm{satisfiesPredQ}(P, \cdot)$ holds at the evaluated value; the
substitution reading turns this into
$\sigma \models P(p(x))$.
\end{proof}

\begin{bremark}[why \texttt{polyTie} is not decoration]
\label{rem:bc-polytie}
Nothing else ties the checked polynomial to the wire the premise
constrains. Without the field, take the satisfiable premise
$w_0 \in \mathcal{R}(\mathrm{conv}(1,2), \emptyset)$ and the polynomial
$p = w_1$ over a \emph{different} wire: the translation remaps the
single-variable canonical form onto the region coordinate, the
certified coefficients are $\{1, 2\}$ --- both satisfying
$\square \geq 1$ --- and the rule would conclude
$\Phi \vdash w_1 \geq 1$, false at any valuation reading
$\sigma(w_1) = 0$. Closed polynomials need no tie: their value reads no
wire.
\end{bremark}

\begin{bremark}[why conic, not merely convex, on rays]
\label{rem:bc-conic-needed}
The predicate-shape clause is two-sided because the two region shapes
transfer differently: convex combinations
(Lemma~\ref{lem:bc-bern}) on the closed simplex, conic combinations
with unbounded weights ($\star$) on a ray. A convex-but-not-conic
predicate fails on the ray, by hand: take $d = 1$,
$\Delta = \mathrm{conv}(0, 1)$ with the vertex at $0$ extended (the
region is the ray $[1, \infty)$), $p(t) = 1 - t$, and
$P = (\square \geq -1)$. The certified table is
$\{\,b_1,\ (\Delta b)_0\,\} = \{\,0,\ -1\,\}$, every value satisfying
$P$; yet at $u = 2$ (the point $t = 3$), $(\star)$ gives
$p = b_1 + 2(b_1 - b_0) = -2 < -1$. The classifier refuses exactly
this: an atom with nonzero constant classifies non-conic, and the
$S \neq \emptyset$ clause then rejects the certificate. Conversely, at
$S = \emptyset$ demanding conic would be needless strength --- the rule
carries the license (any accepted convex predicate) in the premise
itself rather than in surrounding prose.
\end{bremark}

\begin{bremark}[the scope guard is part of the theorem]
\label{rem:bc-guard}
The checker refuses non-constant idx-free polynomials over regions of
dimension $\neq 1$, and multivariate projection-indexed polynomials
whenever $S \neq \emptyset$. Soundness is stated and proven for every
certificate the checker \emph{accepts}; the guard makes the unproven
territory unrepresentable rather than discouraged. The fully general
multivariate ray identity (the arbitrary-shape analogue of
Lemma~\ref{lem:bc-star}) remains open in the development and becomes
load-bearing only if the guard lifts --- the two must move together.
\end{bremark}

\begin{bremark}[relation to the literature]
The rule is the certificate form of Bernstein-coefficient range
enclosure: coefficients bound the range on the simplex
\cite{cargoshisha1966}, enclosures converge under degree elevation and
subdivision \cite{garloff1986} --- which is why \rulename{Subdivide} is
this rule's completeness partner --- and the basis itself goes back to
\cite{bernstein1912}, surveyed in \cite{farouki2012}. With polynomial
atoms and positivity certificates the rule sits beside the
Krivine--Handelman family \cite{krivine1964,handelman1988}: as
everywhere in this calculus, only the trivial (sound) direction is
claimed or needed --- no Positivstellensatz completeness enters the
kernel. The forward-difference computation of the coefficient table is
the classical finite-difference reading of the Bernstein form; the
subdivision algorithms of \cite{decasteljau1959,laneriesenfeld1980}
operate on the same data one rule over.
\end{bremark}

\begin{anchors}
Setting: region atoms as masked barycentric membership & \texttt{MaskRegion}; \texttt{Formula.holds} (\texttt{inRegion} case)\\
Lemma \ref{lem:bc-bern} + Step 3 kernel & the $S = \emptyset$ convex kernel consumed by \texttt{sound\_bernCone\_empty}; per-entry carry \texttt{Term.evalQ\_brRead}\\
Step 2 rational weights & \texttt{maskRegion\_rat\_weights}; \texttt{ConvexRat.exists\_nonneg\_rat\_solution}\\
Lemma \ref{lem:bc-convex} & \texttt{classifyPred}, \texttt{satisfiesPredQ\_convex}\\
Lemma \ref{lem:bc-star} (univariate $(\star)$) & \texttt{QArith.fdBernstein\_sound\_dim1}\\
Step 4a constant corner & \texttt{QArith.fdBernstein\_sound\_const}\\
Lemma \ref{lem:bc-conic} & \texttt{satisfiesPredQ\_conic\_zero}, \texttt{satisfiesPredQ\_conic\_scale}\\
coefficient recomputation (the checked table) & \texttt{bernConeCoeffCheck}, \texttt{allForwardDiffBernCoeffsSatisfyP}, \texttt{QArith.fdBernstein}, \texttt{QArith.gridAll}\\
coordinate translation + index gate & \texttt{Term.asCoordPoly?}, \texttt{Term.asCoordPolyIdx?}, \texttt{Term.idxBounded}\\
Step 1 wire tie & payload field \texttt{polyTie}\\
Step 5 conclusion reading & \texttt{applyPredicateToTerm}, \texttt{holds\_substHole\_of\_satisfies}\\
Theorem \ref{thm:bc-sound} & \texttt{sound\_bernCone} (case split onto \texttt{sound\_bernCone\_empty}); wired as core \texttt{RuleData.bernCone} into \texttt{ProofObject.sound}\\
open residue (Remark \ref{rem:bc-guard}) & \texttt{ConeLemma.StarIdentity} (general shapes; not consulted by the current guard)\\
\end{anchors}


\blcatsection{Admissible and Alias Rules}
\rulesection{16}{Norm}{admissible rule}
\label{blcatrule-Norm}
\ruleblock{The rule (Steele metanotation)}
\sourcenote{as realized in the mechanized development; source row: \texttt{\texttt{rules.toml}} --- \emph{Norm}. The metanotation statement is typeset by hand for notational quality; its mechanized ground truth is the Lean statement below, extracted from the sources at every build of this book and hash-bound. The rule \emph{roster} is machine-checked (every mechanized rule tag round-trips to exactly one catalogue row); the statement \emph{text} is editorial, checked by eye against the extracted Lean beside it.}
\begin{mathpar}
\inferrule*[Right=\rulename{Norm}]{
  t_1, t_2 \ \text{polynomial terms in the same variables with equal monomial canonical forms over } \Q
  \\
  \bshape_H(t_1) = \bshape_H(t_2) \neq \bot \ \text{(a shared defined broadcast shape)}
  \\
  \wires(t_1), \wires(t_2) \ \text{all finitely typed in } H
}{
  \Phi \vdash t_1 = t_2
}
\end{mathpar}

\par\medskip
\rulename{Norm} is the calculus\textquotesingle s computation rule: it
states, as a premise-free sequent, every equality of polynomial terms
that holds by \emph{canonical form} --- $w + w = 2 \cdot w$,
distributivity, collection of like monomials, exact $\Q$ coefficient
arithmetic. A producer uses it constantly and invisibly: whenever an
emitted derivation needs two spellings of the same polynomial to be
interchangeable (before an \rulename{EqElim}, to massage a term into
the pinned shape another rule\textquotesingle s payload demands), a
single \rulename{Norm} node discharges the rewriting that an
equational-logic presentation would smear across many congruence steps.

What the checker recomputes: that both sides are idx-free (projection
terms have no canonical form in the wire-keyed polynomial space --- the
comparison would not carry the rule\textquotesingle s meaning), and
that the two canonical polynomials are \emph{identical} ---
coefficient-for-coefficient equality over $\Q$, never evaluation at
test points, and never rescaling: $2 \cdot w$ and $w + w$ are equal as
canonical forms; $2 \cdot w$ and $w$ are not, proportional though they
are. The remaining payload fields are well-typedness, not computation:
both sides share a defined broadcast shape and mention only finitely
typed wires.

Each hypothesis is load-bearing. Drop the canonical-form identity and
the rule asserts arbitrary equalities. Drop the shared-shape and
finiteness fields and the rule is false \emph{as staged}: at a wire of
extended dtype, even $\vdash t = t$ is invalid --- the atom fails to
hold at every valuation giving the wire an infinite entry, while
holding at the rational ones --- and since the premise is
arbitrary, the staged rule concluded failing sequents from satisfiable
premises; a cancelling sub-term $w_2 - w_2$ over a wire
whose shape does not join with the other side\textquotesingle s leaves
the same canonical form but an atom that never evaluates at any
valuation --- unsatisfiable outright. The typing
fields are exactly what makes ``equal canonical forms\textquotesingle
\textquotesingle{} imply ``equal evaluated tensors\textquotesingle
\textquotesingle{} at every valuation.

\cataloguenote{The rule also names a convention the premise-free rules of this
catalogue share (\emph{the \rulename{Norm} convention}): the premise
$\Phi$ is arbitrary and never consulted --- a fact that holds at every
valuation holds under any premise. The catalogued stratum records an
unpacking sketch through \rulename{BernCone} and \rulename{ArithAtom};
the shipped warrant is direct and stronger (one verified semantic
lemma, all tensor shapes), with the well-typedness fields carrying what
the sketch left implicit.
}
\ruleblock{Lean statement}
\declhead{BernsteinLogic.Soundness}{sound\_norm}{\badge{BadgeGreen}{statement bound}}
\begin{leansrc}
/-- Norm, sound under the PROPOSED-B13-Norm fields: both sides evaluate
(bshape + finWires), at the same shape (`tyEq` + `evalQ_bshape`), with
per-entry values equal because the canonical polynomials are literally
equal. -/
theorem sound_norm {H : Hypernet} {seq : Sequent} {termL termR : Term}
    (hcanon : polynomialCanonicalEq termL termR)
    (htyEq : termL.bshape H = termR.bshape H)
    (htyDef : ∃ s, termL.bshape H = some s)
    (hfinL : termL.finWires H) (hfinR : termR.finWires H)
    (hpost : seq.post = .atom .eq termL termR) :
    seq.semanticallyValid H
\end{leansrc}
\par\medskip
\textbf{The whole computational content, displayed.}\enspace The
check is two gates and one equality: both sides idx-free, and the
canonical polynomials literally equal. Everything else the rule needs
is well-typedness data carried in the payload and consumed by the
soundness lemma, not by the checker.

\declhead{BernsteinLogic.Core.Catalogue}{Term.toPoly}{\badge{BadgeGreen}{statement bound}}
\begin{leansrc}
/-- Exact sparse-polynomial canonical form of a term, variables = wires.
The ONLY canonicalization in leg C (F7 discipline; see original). -/
def Term.toPoly : Term → QArith.Poly
  | .var w => QArith.Poly.var w
  | .const c => QArith.Poly.const c
  | .add l r => QArith.Poly.add l.toPoly r.toPoly
  | .mul l r => QArith.Poly.mul l.toPoly r.toPoly
  | .sub l r => QArith.Poly.sub l.toPoly r.toPoly
  | .neg t => QArith.Poly.neg t.toPoly
  -- G1 INVARIANT (load-bearing): the wire-keyed variable space cannot
  -- express per-entry projections, so idx has NO canonical form here —
  -- the placeholder below is semantically inert because every consumer
  -- of `toPoly` first requires `Term.idxFree` (the canonical-form
  -- relations `polynomialCanonicalEq`/`farkasCheck` carry the guard).
  -- idx terms enter polynomial land ONLY through `asCoordPolyIdx?`.
  | .idx _ _ => QArith.Poly.const 0

/-- Ground-term evaluation: `none` iff a wire occurs — and on any idx
node: a well-typed idx (rank-1 argument) is never closed. -/
def Term.evalQ? : Term → Option Rat
  | .var _ => none
  | .const c => some c
  | .add l r => do pure ((← l.evalQ?) + (← r.evalQ?))
  | .mul l r => do pure ((← l.evalQ?) * (← r.evalQ?))
  | .sub l r => do pure ((← l.evalQ?) - (← r.evalQ?))
  | .neg t => do pure (-(← t.evalQ?))
  | .idx _ _ => none

/-- Comparison evaluation over ℚ. -/
def BowTie.evalQ : BowTie → Rat → Rat → Bool
  | .lt, x, y => decide (x < y)
  | .le, x, y => decide (x ≤ y)
  | .eq, x, y => decide (x = y)
  | .ge, x, y => decide (y ≤ x)
  | .gt, x, y => decide (y < x)
  | .ne, x, y => decide (¬ x = y)

/-- Closed-formula evaluation: `none` iff a wire or region atom occurs
(strict; see original for the closed-world reading). -/
def Formula.evalClosed : Formula → Option Bool
  | .atom b l r => do pure (b.evalQ (← l.evalQ?) (← r.evalQ?))
  | .inRegion _ _ _ => none
  | .and l r => do pure ((← l.evalClosed) && (← r.evalClosed))
  | .or l r => do pure ((← l.evalClosed) || (← r.evalClosed))
  | .tt => some true
  | .ff => some false

/-- The barycentric coordinate functions of a 1-dimensional simplex, as
canonical affine Terms in `t` (pinned shape `offset + slope·t`; see the
original for the full derivation notes and D2 scope). -/
def baryLambdas1D (cx : SimpC) (dContains : Fin cx.nVertices → Bool)
    (t : Term) : Option (List (Fin cx.nVertices × Term))
\end{leansrc}
\declhead{BernsteinLogic.Core.Catalogue}{PNormData}{\badge{BadgeGreen}{statement bound}}
\begin{leansrc}
inductive PNormData : Hypernet → Sequent → Type where
  | mk (H : Hypernet) (seq : Sequent)
       (termL termR : Term)
       (canonEq : polynomialCanonicalEq termL termR)
       -- PROPOSED-B13-Norm (register B13 rider; realized provisionally
       -- 2026-07-16 on the repair hint "asserting wires have finite
       -- type" — the same field family that repaired EqElim/ConvIntro):
       -- both sides are shape-stable at a common broadcast shape and
       -- finitely typed. Without these, an ∞-valued ext wire makes even
       -- ⊢ t = t unsatisfiable (pre is arbitrary), and a cancelling
       -- wire of incompatible shape (w₂ − w₂ at a shape unjoinable with
       -- termL's) makes the conclusion atom never evaluate.
       (tyEq   : termL.bshape H = termR.bshape H)
       (tyDef  : ∃ s, termL.bshape H = some s)
       (finL   : termL.finWires H)
       (finR   : termR.finWires H)
       (postEq  : seq.post = .atom .eq termL termR)
       : PNormData H seq
\end{leansrc}
\declhead{BernsteinLogic.Core.Catalogue}{Term.idxFree}{\badge{BadgeGreen}{definition bound}}
\begin{leansrc}
/-- Structural absence of the idx constructor. The polynomial
canonical-form fragment (D2): `toPoly`'s wire-keyed variable space cannot
express per-entry projections, so every canonical-form consumer
(`polynomialCanonicalEq`, `farkasCheck`, the legacy `asCoordPoly?`
remap) requires this; idx coordinate polynomials take the
`asCoordPolyIdx?` route instead (G1). -/
def Term.idxFree : Term → Bool
  | .var _ | .const _ => true
  | .add l r | .mul l r | .sub l r => l.idxFree && r.idxFree
  | .neg t => t.idxFree
  | .idx _ _ => false
\end{leansrc}
\declhead{BernsteinLogic.Core.Catalogue}{polynomialCanonicalEq}{\badge{BadgeGreen}{definition bound}}
\begin{leansrc}
def polynomialCanonicalEq (l r : Term) : Prop :=
  l.idxFree = true ∧ r.idxFree = true ∧ l.toPoly = r.toPoly
\end{leansrc}
\ruleblock{Correspondence (Steele display $\leftrightarrow$ Lean)}
\badge{BadgeGreen}{reviewer-certified against the displayed texts}
\sourcenote{The dictionary below is a reviewer-certified statement of how the mechanized statements above interpret the displayed rule: display metavariable $\mapsto$ Lean binder, each display hypothesis $\mapsto$ the premise or side condition carrying it, the display conclusion $\mapsto$ the goal --- with any deliberate difference in width between display and mechanization named outright. Certification is recorded as a hash triple over (the displayed cell, the extracted statement texts, this dictionary\textquotesingle s own text); if any of the three changes, this block turns amber until a reviewer re-certifies. The correspondence claim itself remains editorial --- reviewer-grade, not machine-checked; its machine-checked successor types each soundness statement against the rule as a structured object, and supersedes this block per rule as it lands.}
Judgment frame as in \rulename{Ax}.
\begin{itemize}
\item $t_1, t_2 \mapsto$ \texttt{termL termR : Term}; the displayed
  side condition (equal monomial canonical forms) $\mapsto$
  \texttt{hcanon} $:$ \texttt{polynomialCanonicalEq termL termR},
  quoted in full below --- both sides idx-free and the canonical
  polynomials literally equal. The display\textquotesingle s ``same
  variables\textquotesingle\textquotesingle{} and ``polynomial
  terms\textquotesingle\textquotesingle{} have no separate binders:
  canonical forms are wire-keyed (equality forces the same canonical
  support --- a mentioned-but-cancelled wire is no constraint on the
  polynomial; evaluability polices it through the typing fields
  instead) and the term grammar is polynomial by construction.
\item conclusion $\Phi \vdash t_1 = t_2$ $\mapsto$ the goal under
  \texttt{hpost} $:$ \texttt{seq.post = .atom .eq termL termR};
  the arbitrary $\Phi$ is premise-freeness by silence --- the
  precondition component of \texttt{seq : Sequent} is unconstrained
  by any hypothesis (the \rulename{PConstExact} representation).
\end{itemize}
Width note: the display carries the
typing pack inline, matching the mechanized statement\textquotesingle s
four well-typedness hypotheses one for one ---
\texttt{htyEq} (both broadcast shapes equal), \texttt{htyDef} (and
defined --- the \texttt{some}), \texttt{hfinL}/\texttt{hfinR} (all
free wires at finite dtypes). The bound \texttt{PNormData} payload
carries the same conditions under its own field names ---
\texttt{hcanon} $\leftrightarrow$ \texttt{canonEq}, \texttt{htyEq}
$\leftrightarrow$ \texttt{tyEq}, \texttt{htyDef} $\leftrightarrow$
\texttt{tyDef}, \texttt{hfinL}/\texttt{hfinR} $\leftrightarrow$
\texttt{finL}/\texttt{finR}, \texttt{hpost} $\leftrightarrow$
\texttt{postEq} --- and they are not optional: without them the
bare canonical-form rule is refuted (an $\infty$-valued wire makes
even $t = t$ fail at every extended-valued assignment --- the roster
row\textquotesingle s
counterexample). The one-line canonical-form sentence is the
rule\textquotesingle s
\emph{intent}; the certified statement is that sentence plus exactly
this typing pack.

\ruleblock{Lean verifier receipts}
\sourcenote{module \texttt{BernsteinLogic.Soundness}; bound build \texttt{a20d90e66c13\ldots} (bound from receipt a9c0e34d3e3c9c8c, 2026-08-03T21:30:36Z). A receipt records that named checking layers accepted a build: (module, compiled-artifact SHA-256, environment SHA-256, checker layers), minted when the checking gate ran; the store receipts attesting the bound build are listed with the layers each run credited. Nothing here asks the reader to trust the store: every layer is an independently runnable program over the quoted sources --- rebuild the module and re-run the checkers to reproduce the verdict. Receipts attesting superseded builds credit no layer and are omitted from this rendering (the count of omissions, if any, is noted below).}
\begin{receipttable}
\receiptrow{a9c0e34d3e3c9c8c}{2026-08-03T21:30:36Z}{build+decorrelated}{\badge{BadgeGreen}{current}}
\end{receipttable}
\sourcenote{21 further receipt(s) in the store attest superseded builds of this module and are omitted from this rendering.}
\sourcenote{module \texttt{BernsteinLogic.Core.Catalogue}; bound build \texttt{cbc92e976a1d\ldots} (bound from receipt 5abdb3589707b5df, 2026-08-03T21:30:33Z). A receipt records that named checking layers accepted a build: (module, compiled-artifact SHA-256, environment SHA-256, checker layers), minted when the checking gate ran; the store receipts attesting the bound build are listed with the layers each run credited. Nothing here asks the reader to trust the store: every layer is an independently runnable program over the quoted sources --- rebuild the module and re-run the checkers to reproduce the verdict. Receipts attesting superseded builds credit no layer and are omitted from this rendering (the count of omissions, if any, is noted below).}
\begin{receipttable}
\receiptrow{24804e66edba0bc9}{2026-08-03T20:25:09Z}{build+decorrelated}{\badge{BadgeGreen}{current}}
\receiptrow{5abdb3589707b5df}{2026-08-03T21:30:33Z}{build+decorrelated}{\badge{BadgeGreen}{current}}
\end{receipttable}
\sourcenote{18 further receipt(s) in the store attest superseded builds of this module and are omitted from this rendering.}
\layerline{\layeritem{build}{lake build (elaboration)}{\badge{BadgeGreen}{current}}\layeritem{replay}{leanchecker (official kernel replay)}{\badge{BadgeAmber}{stale only}}\layeritem{hardened}{safe\_verify (hardened screen)}{\badge{BadgeAmber}{stale only}}\layeritem{decorrelated}{lean4export + nanoda (decorrelated kernel)}{\badge{BadgeGreen}{current}}}
\sourcenote{\badge{BadgeGreen}{comparator: fixture NormOutsideFragment} the
three-implementation comparator harness (compiled machine code,
kernel-reduction probes, and the compiler-IR interpreter, all executing
the same elaborated checker functions) carries a well-typed single
\rulename{Norm} node as its fixture \texttt{NormOutsideFragment},
expected verdicts \texttt{check = true}, \texttt{v0frag = false}. The
fixture is the suite\textquotesingle s false-verdict control:
\rulename{Norm} sits outside the restricted fragment recognizer by
design, so the fixture pins the reject lane end to end --- a comparator
that could only confirm acceptance would be blind to a path that
wrongly accepts everything.
}
\ruleblock{HOL4 statement}
\sourcenote{theory \texttt{bern\_catalogue}, HOL4 development script \texttt{\texttt{bern\_catalogueScript.sml}} \badge{BadgeGreen}{script bound}}
\holhead{Definition polynomialCanonicalEqB\_def}
\begin{holsrc}
Definition polynomialCanonicalEqB_def:
  polynomialCanonicalEqB l r ⇔
    termIdxFree l ∧ termIdxFree r ∧ termToPoly l = termToPoly r
End
\end{holsrc}
\ruleblock{HOL verifier receipts}
\badge{BadgeBlue}{documentary} \texttt{hol-strict} CLEAN over the ported checker chain: theory
\texttt{bern\_catalogue}, 472 theorems, 0 cheat-tagged, 0 nonstandard
axioms (whitelist: \texttt{DISK\_THM}; the four \texttt{bool} axioms) ---
verification run of 2026-07-22 on an isolated build host, scripts
sha-matched verbatim to the sources cited here (git blob
\texttt{e41b9ed0\ldots}). The ported receipt table pins the
rule\textquotesingle s exactness discipline in both directions: the
canonical-form accept ($w_0 + w_0 = 2 \cdot w_0$) and the
scale-preserving reject ($2 \cdot w_0 \neq w_0$ as canonical forms).
Scope note, honest by the port\textquotesingle s own seam inventory:
the ported surface is the first-order check function; the
payload\textquotesingle s well-typedness fields (broadcast shapes,
finite typing) live in the dependently-typed proof-object layer, which
is a separately tracked porting stage.
Semantic stratum, also ported: the soundness twin \texttt{sound\_norm}
(theory \texttt{bern\_tr\_fragclose}: \texttt{hol-strict} clean, 26
theorems, 0 offences, verification run of 2026-07-22 on an isolated
build host, content-addressed at the same sources), consumed by the
certified-fragment keystone (\texttt{checkPf\_sound\_frag}, theory
\texttt{bern\_tr\_splice}, 20 theorems, 0 offences, same run): a
reference-checker accept inside the proven wire fragment yields sequent
validity.

\sourcenote{HOL receipts are documentary: quoted results of \texttt{hol-strict} runs
and census rows from the HOL4 development\textquotesingle s build
evidence, not minted receipt artefacts. A HOL receipt store (the Lean
store\textquotesingle s twin) is named future work.
}
\ruleblock{Proof}

\sourcenote{Genre: admissible rule (tier 2) --- a proof-tree-rewriting
recipe. The tier taxonomy defines tier 2 as \emph{verified unpacking
into core rules} (one node $\to$ many), with rules admissible by
\emph{reduction to verified machinery} as the recorded boundary case;
\rulename{Norm} is of the latter kind, and this section --- the book's
first admissible rule --- states the general template that the later
recipes instantiate.}

\begin{brecipe}[the admissible template, general form]\label{rec:norm-template}
An admissible rule $R$ with premises $P_1, \dots, P_n$ and conclusion $C$
is warranted by a \emph{rewrite}: a function taking any proof-tree node
\[
\inferrule*[Right=$R$]{\;\mathcal{T}_1 \rhd P_1 \\ \cdots \\ \mathcal{T}_n \rhd P_n\;}{C}
\qquad\rightsquigarrow\qquad
\mathcal{D}\bigl[\mathcal{T}_1, \dots, \mathcal{T}_n\bigr] \rhd C
\]
where $\mathcal{D}$ is a derivation fragment over core (and
previously-justified admissible) rules with the same conclusion,
consuming the premise subtrees $\mathcal{T}_i$ at its leaves. Where such
a fragment exists, applying
it leaf-first eliminates every $R$-node from a finite tree in
finitely many steps, so any judgment derivable with $R$ is derivable
without it --- admissibility in the derivability sense. In the tier's
recorded boundary case --- rules admissible \emph{by reduction to
verified machinery}, this one among them --- no fragment is applied:
where no catalogue derivation of $C$ can exist, the node is instead
discharged in place against a named verified theorem, no derivability
claim is made, and the conserved property --- weaker, and sufficient ---
is the calculus's soundness theorem itself. Either way, no new kernel
semantic content enters. The verified object
shipped per admissible rule is a soundness
lemma for $R$ that factors through already-verified machinery: for a
fragment-backed rule, provable by replaying $\mathcal{D}$; for a
boundary-case rule, the discharge theorem itself.
\end{brecipe}

\begin{brecipe}[the \rulename{Norm} instance --- the degenerate rewrite]
\label{rec:norm-instance}
\rulename{Norm} has no proof premises ($n = 0$): the node
\[
\inferrule*[Right=\rulename{Norm}]{
  \mathrm{polynomialCanonicalEq}(t_1, t_2) \\
  \text{\upshape both sides shape-stable and finitely typed}
}{\Phi \vdash t_1 = t_2}
\qquad\rightsquigarrow\qquad
\text{\upshape the checker replays the canonical forms;}\;
\text{\texttt{sound\_norm} stands as warrant.}
\]
The side condition is \emph{computational}: both sides idx-free, and
their canonical monomial polynomials \emph{identical} over $\Q$ ---
coefficient-for-coefficient, decidable, replayed by the checker. The
catalogued stratum records an unpacking sketch through
\rulename{BernCone} and \rulename{ArithAtom} (an equality as a pair of
zero-polynomial bounds); the shipped warrant is direct and stronger ---
one verified semantic lemma at all tensor shapes --- with the payload's
well-typedness fields carrying exactly what the sketch left implicit.
\end{brecipe}

\begin{bprop}[semantic warrant]\label{prop:norm-sound}
Suppose $t_1, t_2$ are idx-free with $t_1.\mathrm{toPoly} =
t_2.\mathrm{toPoly}$, the two sides carry the same defined broadcast
shape in $H$, and every wire of either side is finitely typed. Then
$\Phi \vdash t_1 = t_2$ is semantically valid over $H$ for
\emph{arbitrary} $\Phi$.
\end{bprop}
\begin{proof}
Fix any admissible valuation $\sigma$; the premise is not consulted ---
this is \emph{the \rulename{Norm} convention}, shared by every
premise-free rule of the catalogue: a fact holding at every valuation
holds under any premise. The shape and finiteness fields give both
sides defined evaluations, as tensors of the common broadcast shape
$s$: a term's broadcast shape is computed structurally, so its
definedness already certifies that every subterm's join exists;
finitely typed wires always read; and idx-freedom leaves no projection
to fail. If $s$ has positive rank and zero total size, both sides
evaluate to the unique empty tensor of shape $s$ and the equality atom
holds outright --- there is no entry to compare. Otherwise fix an entry
index of $s$ (for $s = [\,]$, the single scalar
position) and read each wire at it: a rank-0 wire contributes its
scalar, and a positive-rank wire --- necessarily of shape $s$ itself,
the only positive rank that joins with $s$ under this broadcast
discipline --- contributes its entry at the index. Evaluation at that
read-out is a
ring homomorphism into $\Q$ from the polynomial ring over the union of
the wires the two sides mention (a wire absent from one side, or
cancelled from its canonical form, simply carries coefficient zero
there); a syntactic identity of canonical
forms is an identity of polynomials, and ring homomorphisms preserve
polynomial identities --- so the two evaluated tensors agree entry by
entry, and their shapes agree by hypothesis. Hence the equality atom
holds at $\sigma$. That is the whole proof; \texttt{sound\_norm}
packages it once, per-entry, at every tensor shape.
\end{proof}

\begin{bremark}[what broke without the typing fields]
\label{rem:norm-b13}
As first staged --- canonical-form equality alone --- the rule was
\emph{false}, twice over, and the counterexamples are why the fields
exist. (i) At a wire of extended dtype, the staged rule concluded
\emph{invalid} sequents: at any valuation giving the wire an infinite
entry, even $t = t$ fails to hold --- evaluation is over exact
rationals, and an infinite entry refuses --- so $\vdash t = t$ fails
exactly at the extended-valued assignments while holding at the
rational ones. (ii) A cancelling
sub-term leaves the canonical form unchanged but not the evaluation:
$t_2 = t_1 + (w_2 - w_2)$ has $t_2.\mathrm{toPoly} =
t_1.\mathrm{toPoly}$, yet if $w_2$'s shape does not join with $t_1$'s,
the right-hand atom side never evaluates at \emph{any} valuation ---
there the concluded atom is unsatisfiable outright. In either case a
satisfiable (indeed arbitrary) premise concluded a failing sequent; the
shared-shape and
finiteness fields exclude both. A hypothesis whose necessity is
semantic (evaluability), not algebraic, is exactly the kind a recipe
must state, or the next porter re-learns it the hard way.
\end{bremark}

\begin{bremark}[exactness discipline, and the idx-freedom gate]
\label{rem:norm-f7}
Canonical-form equality is never up-to-rescaling: $2 \cdot w$ and
$w + w$ are equal as canonical forms; $2 \cdot w$ and $w$ are not,
proportional though they are --- the matching discipline every
combination rule of this catalogue (\rulename{Farkas},
\rulename{Syzygy}) inherits when it compares atom polynomials. The
idx-freedom gate is the meaning-preservation condition of the
translation: the wire-keyed polynomial space cannot express per-entry
projections, so a projection-bearing term has no canonical form here
and the comparison is refused rather than performed on the wrong
object. (The combination rules' widened replays translate projections
into their own variable spaces; those translations belong to those
rules, not to \rulename{Norm}.)
\end{bremark}

\begin{anchors}
Recipe \ref{rec:norm-instance} warrant & \texttt{sound\_norm} (per-entry evaluation at the common broadcast shape)\\
the computational check & \texttt{polynomialCanonicalEq} $=$ idx-freedom gates $+$ \texttt{Term.toPoly} equality\\
idx-freedom gate (Remark \ref{rem:norm-f7}) & \texttt{Term.idxFree}; \texttt{Term.toPoly}'s wire-keyed variable space\\
typing fields (Remark \ref{rem:norm-b13}) & payload fields \texttt{tyEq}, \texttt{tyDef}, \texttt{finL}, \texttt{finR} (\texttt{PNormData})\\
premise-freedom ($\Phi$ arbitrary) & the \rulename{Norm} convention; \texttt{pre} unused in \texttt{sound\_norm}\\
wire placement & \texttt{RuleTag.norm}\\
\end{anchors}


\rulesection{17}{Farkas}{admissible rule}
\label{blcatrule-Farkas}
\ruleblock{The rule (Steele metanotation)}
\sourcenote{as realized in the mechanized development; source row: \texttt{\texttt{rules.toml}} --- \emph{Farkas}. The metanotation statement is typeset by hand for notational quality; its mechanized ground truth is the Lean statement below, extracted from the sources at every build of this book and hash-bound. The rule \emph{roster} is machine-checked (every mechanized rule tag round-trips to exactly one catalogue row); the statement \emph{text} is editorial, checked by eye against the extracted Lean beside it.}
\begin{mathpar}
\inferrule*[Right=\rulename{Farkas}]{
  \{\ell_j \bowtie_j 0\}_j \text{ a sub-conjunction of } \Phi \text{ with each } \ell_j \text{ linear, } \bowtie_j \in \{\leq, =\} \\
  \{\lambda_j \in \Q_{\geq 0}\}_j \text{ nonneg coefficients, with } \lambda_j \text{ free when } \bowtie_j \text{ is } {=} \\
  \bshape_H(\textstyle\sum_j \lambda_j \ell_j) \neq \bot \\
  \wires(\textstyle\sum_j \lambda_j \ell_j) \text{ all finitely typed in } H \\
  \wires(\Phi) \cup \wires(\textstyle\sum_j \lambda_j \ell_j) \text{ all typed at positive-size shapes in } H
}{
  \Phi \vdash \textstyle\sum_j \lambda_j \ell_j \leq 0
}
\end{mathpar}

\par\medskip
\rulename{Farkas} is the affine-combination rule over the premise: from
atoms $\{\ell_j \bowtie_j 0\}$ occurring as sub-conjuncts of $\Phi$ and
nonnegative multipliers $\lambda_j$ (free-signed on equality atoms), it
concludes $\Phi \vdash \sum_j \lambda_j \ell_j \leq 0$. It is the
one-rule form of what \rulename{Syzygy} does with sub-\emph{proofs}:
here every combined atom must already sit in the premise conjunction, so
no component derivations enter. A producer finds the multipliers by
whatever method it likes (an LP feasibility call, typically) and ships
only the combination term; nothing about the search survives into the
certificate.

What the checker recomputes, given the payload (combination term,
conclusion term): decode the combination into weighted entries --- a
fixed canonical encoding (a right-nested sum of
$\lambda_j \cdot e_j$ products with a zero terminator), anything else
rejected; require each weight nonnegative; require each
entry\textquotesingle s polynomial to occur among the
$({\cdot} \leq 0)$-normal polynomials of a premise conjunct --- equality
atoms expose \emph{both} polarities to the matcher, which is exactly how
a free-sign multiplier on an equality is realized as a nonnegative
weight on the appropriate polarity; and require the weighted fold of the
entry polynomials to equal the conclusion\textquotesingle s polynomial
\emph{identically}, as canonical forms over $\Q$. The check is a
disjunction of three branches of one replay algorithm, differing only in
the variable space where polynomials are compared: the legacy idx-free
canonical-form branch, a wire-aware coordinate branch for
projection-indexed terms, and a mixed tagged-union branch for
combinations mentioning both bare wires and projections.

Each hypothesis is load-bearing. Drop the sign condition and a producer
multiplies a bound by $-1$ and concludes its negation. Drop the
atom-sourcing check and the combination smuggles in atoms nobody
assumed. Drop the exact fold identity and the conclusion is unrelated to
the premise. Drop the positive-size wire condition and the rule is
unsound at degenerate shapes: comparison atoms over a zero-size
(empty-tensor) wire hold \emph{vacuously}, and a combination of two
vacuous contradictory atoms --- $5 \leq w$ and $w \leq 3$ at a wire of
shape $[0]$ --- derives the false scalar conclusion
$\vdash 2 \leq 0$ from a satisfiable premise. Drop the
conclusion\textquotesingle s shape/finiteness fields and a conclusion
mentioning an untyped wire has a correct canonical polynomial but never
evaluates, so the concluded atom never holds. The proof section walks
these in order.

\cataloguenote{The catalogued statement ranges over \emph{linear} $\ell_j$; the replay
itself is degree-blind, and soundness never consults linearity --- the
restriction is the rule\textquotesingle s published register (the shape
producers emit and the wire format prices), not a checked side
condition. \rulename{Syzygy} lifts the same fold to component
conclusions with degrees unrestricted; \rulename{PBound} re-spells that
lift as the compact affine bound-node.
}
\ruleblock{Lean statement}
\declhead{BernsteinLogic.Soundness}{farkas\_fold\_eval\_nonpos}{\badge{BadgeGreen}{statement bound}}
\begin{leansrc}
/-- Evaluation of the Farkas combination fold: a nonnegative combination
of nonpositive values is nonpositive (valuation-general). -/
theorem farkas_fold_eval_nonpos (ρ : Nat → ℚ) :
    ∀ (entries : List (ℚ × Term)),
    (∀ e ∈ entries, 0 ≤ e.1 ∧ QArith.Poly.eval ρ e.2.toPoly ≤ 0) →
    QArith.Poly.eval ρ
      (entries.foldr
        (fun (le : ℚ × Term) acc =>
          QArith.Poly.add (QArith.Poly.scale le.1 le.2.toPoly) acc) []) ≤ 0
\end{leansrc}
\declhead{BernsteinLogic.Soundness}{farkas\_fold\_eval\_nonpos\_polys}{\badge{BadgeGreen}{statement bound}}
\begin{leansrc}
/-- The nonneg fold in the paired space (valuation-general; entries
carry their polynomials directly — the `farkasEntriesW`-translated
form). -/
theorem farkas_fold_eval_nonpos_polys (ρ : Nat → ℚ) :
    ∀ (pentries : List (ℚ × QArith.Poly)),
    (∀ e ∈ pentries, 0 ≤ e.1 ∧ QArith.Poly.eval ρ e.2 ≤ 0) →
    QArith.Poly.eval ρ
      (pentries.foldr
        (fun (le : ℚ × QArith.Poly) acc =>
          QArith.Poly.add (QArith.Poly.scale le.1 le.2) acc) []) ≤ 0
\end{leansrc}
\declhead{BernsteinLogic.Soundness}{sound\_farkas}{\badge{BadgeGreen}{statement bound}}
\begin{leansrc}
/-- Farkas, general tensor shapes (supersedes the v0 singleton form):
the conclusion's B13 shape/finiteness fields give it a defined
evaluation; the PROPOSED-B14 positivity field rules out vacuously-held
zero-size atoms; the per-entry claim runs through the CLAMPED valuation,
which serves heterogeneous premise-atom shapes at every conclusion
index. -/
theorem sound_farkas {H : Hypernet} {seq : Sequent}
    {combinationCoeffs conclusion : Term}
    (hshape : ∃ s, conclusion.bshape H = some s)
    (hfin : conclusion.finWires H)
    (hwpos : ∀ w ∈ seq.pre.freeWires ++ conclusion.freeWires,
      ∃ τ, H.wireType w = some τ ∧ 0 < Shape.totalSize τ.shape)
    (hvalid : farkasCombinationValid combinationCoeffs seq.pre conclusion)
    (hpost : seq.post = .atom .le conclusion (.const 0)) :
    seq.semanticallyValid H
\end{leansrc}
\par\medskip
\textbf{What the combination check means --- and what it does
not.}\enspace The side condition appears in the payload as the single
name \texttt{farkasCombinationValid}, so this book binds its meaning by
displaying the decoding chain from the checker sources.
\texttt{farkasEntries} decodes the combination term against a
\emph{pinned} canonical encoding --- anything else is rejected, so the
wire format cannot express a malformed combination.
\texttt{atomLeZeroPolys} is the $({\cdot} \leq 0)$-normalization of one
premise atom, exact and never rescaled; note that equality atoms
contribute both polarities, which is the entire realization of the
classical free-sign multiplier discipline for equality constraints.
\texttt{farkasCheckCore} strings the gates together in the idx-free
canonical-form space; \texttt{farkasCheck} is the three-way branch
disjunction --- the two widened branches replay the same gates
verbatim in the wire-aware coordinate space (projection terms
$w[j]$ keyed by an injective pairing of wire and coordinate, so
distinct wires can never collide) and in the mixed tagged-union space
(bare wires and projections as disjoint variable classes; the variant
identifying a bare wire with its zeroth projection admits a false
sequent and is deliberately unrepresentable).

What a skeptical reader should confirm by eye: the conclusion term
enters only as the target of the exact fold identity; the
step\textquotesingle s posted formula appears nowhere in the check; and
nothing in the chain evaluates or consults the truth of anything ---
the check manipulates polynomials, never valuations. Semantic truth is
established only in \texttt{sound\_farkas}, which case-splits on the
accepted branch.

\declhead{BernsteinLogic.Core.Catalogue}{atomLeZeroPolysW}{\badge{BadgeGreen}{statement bound}}
\begin{leansrc}
/-- `atomLeZeroPolys` in the paired coordinate space: the
`(e ≤ 0)`-normal-form polynomials of one atom under the wire-aware idx
translation. An atom with an untranslatable side entails nothing (like
non-atoms — it can then never source a combination entry). -/
def atomLeZeroPolysW : Formula → List QArith.Poly
  | .atom op l r =>
    match l.asCoordPolyIdxW? with
    | none => []
    | some pl =>
      match r.asCoordPolyIdxW? with
      | none => []
      | some pr =>
        match op with
        | .le | .lt => [QArith.Poly.sub pl pr]
        | .ge | .gt => [QArith.Poly.sub pr pl]
        | .eq => [QArith.Poly.sub pl pr, QArith.Poly.sub pr pl]
        | .ne => []
  | _ => []

/-- Translate every decoded entry's vehicle term; `none` if any vehicle
is outside the idx coordinate fragment. (Explicit recursion rather than
`mapM` so soundness inverts it by plain `cases`.) -/
def farkasEntriesW : List (Rat × Term) → Option (List (Rat × QArith.Poly))
  | [] => some []
  | (l, e) :: rest =>
    match e.asCoordPolyIdxW?, farkasEntriesW rest with
    | some pe, some prest => some ((l, pe) :: prest)
    | _, _ => none

/-- The idx-form branch of the Farkas check: every polynomial — entry
vehicles, matched premise atoms, and the conclusion — rides the
wire-aware coordinate translation; entry encoding, nonneg weights,
atom matching, and the exact combination fold mirror the idx-free
branch verbatim. -/
def farkasCheckW (cc : Term) (premise : Formula) (conclusion : Term) : Bool
\end{leansrc}
\declhead{BernsteinLogic.Core.Catalogue}{atomLeZeroPolysM}{\badge{BadgeGreen}{statement bound}}
\begin{leansrc}
/-- `atomLeZeroPolys` in the mixed tagged space: the `(e ≤ 0)`-normal-
form polynomials of one atom under the mixed translation. An atom with
an untranslatable side entails nothing (the W-branch discipline). -/
def atomLeZeroPolysM : Formula → List QArith.Poly
  | .atom op l r =>
    match l.asCoordPolyMixed? with
    | none => []
    | some pl =>
      match r.asCoordPolyMixed? with
      | none => []
      | some pr =>
        match op with
        | .le | .lt => [QArith.Poly.sub pl pr]
        | .ge | .gt => [QArith.Poly.sub pr pl]
        | .eq => [QArith.Poly.sub pl pr, QArith.Poly.sub pr pl]
        | .ne => []
  | _ => []

/-- Translate every decoded entry's vehicle term into the mixed space;
`none` if any vehicle is outside the mixed fragment. -/
def farkasEntriesM : List (Rat × Term) → Option (List (Rat × QArith.Poly))
  | [] => some []
  | (l, e) :: rest =>
    match e.asCoordPolyMixed?, farkasEntriesM rest with
    | some pe, some prest => some ((l, pe) :: prest)
    | _, _ => none

/-- The mixed-variable branch of the Farkas check: every polynomial —
entry vehicles, matched premise atoms, and the conclusion — rides the
mixed tagged translation; entry encoding, nonneg weights, atom matching,
and the exact combination fold mirror the sibling branches verbatim. -/
def farkasCheckM (cc : Term) (premise : Formula) (conclusion : Term) : Bool
\end{leansrc}
\declhead{BernsteinLogic.Core.Catalogue}{farkasCheckW}{\badge{BadgeGreen}{statement bound}}
\begin{leansrc}
/-- The idx-form branch of the Farkas check: every polynomial — entry
vehicles, matched premise atoms, and the conclusion — rides the
wire-aware coordinate translation; entry encoding, nonneg weights,
atom matching, and the exact combination fold mirror the idx-free
branch verbatim. -/
def farkasCheckW (cc : Term) (premise : Formula) (conclusion : Term) : Bool
\end{leansrc}
\declhead{BernsteinLogic.Core.Catalogue}{farkasCheckM}{\badge{BadgeGreen}{statement bound}}
\begin{leansrc}
/-- The mixed-variable branch of the Farkas check: every polynomial —
entry vehicles, matched premise atoms, and the conclusion — rides the
mixed tagged translation; entry encoding, nonneg weights, atom matching,
and the exact combination fold mirror the sibling branches verbatim. -/
def farkasCheckM (cc : Term) (premise : Formula) (conclusion : Term) : Bool
\end{leansrc}
\declhead{BernsteinLogic.Core.Catalogue}{farkasCombinationValid}{\badge{BadgeGreen}{statement bound}}
\begin{leansrc}
def farkasCombinationValid (cc : Term) (premise : Formula)
    (conclusion : Term) : Prop
\end{leansrc}
\declhead{BernsteinLogic.Core.Catalogue}{wireCoordKey}{\badge{BadgeGreen}{statement bound}}
\begin{leansrc}
/-- Injective key for the (wire, coordinate) variable space of the
multi-wire idx coordinate translation below. This is `Nat.pair`'s
definition inlined (the core is mathlib-free; Soundness.lean proves the
`Nat.unpair` round-trip through the definitional equality). Distinct
wires MUST get distinct polynomial variables here: a Syzygy composite
mixes projections of several region wires in one combination, and
conflating `w0[j]` with `w2[j]` would let semantically different
polynomials collide (contrast `asCoordPolyIdx?`, whose BernCone scope
has a single region wire pinned by `polyTie`, so dropping the wire is
sound THERE and only there). -/
def wireCoordKey (w i : Nat) : Nat
\end{leansrc}
\declhead{BernsteinLogic.Core.Catalogue}{Term.asCoordPolyIdxW?}{\badge{BadgeGreen}{statement bound}}
\begin{leansrc}
/-- The wire-aware idx coordinate translation: a Term whose every wire
mention is a projection `.idx (.var w) j` becomes a polynomial over the
paired variables `wireCoordKey w j`. `none` on bare wire mentions and
on idx of any non-var argument (v0.5 scope, matching
`asCoordPolyIdx?`). Constants translate, so idx-free closed terms ride
either route. -/
def Term.asCoordPolyIdxW? : Term → Option QArith.Poly
  | .var _ => none
  | .const c => some (QArith.Poly.const c)
  | .add l r =>
    do pure (QArith.Poly.add (← l.asCoordPolyIdxW?) (← r.asCoordPolyIdxW?))
  | .mul l r =>
    do pure (QArith.Poly.mul (← l.asCoordPolyIdxW?) (← r.asCoordPolyIdxW?))
  | .sub l r =>
    do pure (QArith.Poly.sub (← l.asCoordPolyIdxW?) (← r.asCoordPolyIdxW?))
  | .neg t => (Term.asCoordPolyIdxW? t).map QArith.Poly.neg
  | .idx (.var w) i => some (QArith.Poly.var (wireCoordKey w i))
  | .idx _ _ => none

/-- `atomLeZeroPolys` in the paired coordinate space: the
`(e ≤ 0)`-normal-form polynomials of one atom under the wire-aware idx
translation. An atom with an untranslatable side entails nothing (like
non-atoms — it can then never source a combination entry). -/
def atomLeZeroPolysW : Formula → List QArith.Poly
  | .atom op l r =>
    match l.asCoordPolyIdxW? with
    | none => []
    | some pl =>
      match r.asCoordPolyIdxW? with
      | none => []
      | some pr =>
        match op with
        | .le | .lt => [QArith.Poly.sub pl pr]
        | .ge | .gt => [QArith.Poly.sub pr pl]
        | .eq => [QArith.Poly.sub pl pr, QArith.Poly.sub pr pl]
        | .ne => []
  | _ => []

/-- Translate every decoded entry's vehicle term; `none` if any vehicle
is outside the idx coordinate fragment. (Explicit recursion rather than
`mapM` so soundness inverts it by plain `cases`.) -/
def farkasEntriesW : List (Rat × Term) → Option (List (Rat × QArith.Poly))
  | [] => some []
  | (l, e) :: rest =>
    match e.asCoordPolyIdxW?, farkasEntriesW rest with
    | some pe, some prest => some ((l, pe) :: prest)
    | _, _ => none

/-- The idx-form branch of the Farkas check: every polynomial — entry
vehicles, matched premise atoms, and the conclusion — rides the
wire-aware coordinate translation; entry encoding, nonneg weights,
atom matching, and the exact combination fold mirror the idx-free
branch verbatim. -/
def farkasCheckW (cc : Term) (premise : Formula) (conclusion : Term) : Bool
\end{leansrc}
\declhead{BernsteinLogic.Core.Catalogue}{Term.asCoordPolyMixed?}{\badge{BadgeGreen}{statement bound}}
\begin{leansrc}
/-- The mixed-variable translation: bare wires and idx projections are
DISJOINT opaque variables of one combination algebra — bare `.var w`
keyed even (`2·w`), projections `.idx (.var w) j` keyed odd
(`2·wireCoordKey w j + 1`) — so the two classes can never collide. In
particular bare `w` is NEVER identified with `w[0]`: the identification
variant admits a false sequent (a countermodel witnesses the failure,
alongside the REJECT specimen in EvalDemo). This matches the python reference
(`bl_checker` keys bare and projected mentions as distinct symbols) and
the geolog/fastcheck mixed evaluators. `none` only on idx-of-non-var
(outside every reading's fragment). -/
def Term.asCoordPolyMixed? : Term → Option QArith.Poly
  | .var w => some (QArith.Poly.var (2 * w))
  | .const c => some (QArith.Poly.const c)
  | .add l r =>
    do pure (QArith.Poly.add (← l.asCoordPolyMixed?) (← r.asCoordPolyMixed?))
  | .mul l r =>
    do pure (QArith.Poly.mul (← l.asCoordPolyMixed?) (← r.asCoordPolyMixed?))
  | .sub l r =>
    do pure (QArith.Poly.sub (← l.asCoordPolyMixed?) (← r.asCoordPolyMixed?))
  | .neg t => (Term.asCoordPolyMixed? t).map QArith.Poly.neg
  | .idx (.var w) j => some (QArith.Poly.var (2 * wireCoordKey w j + 1))
  | .idx _ _ => none

/-- `atomLeZeroPolys` in the mixed tagged space: the `(e ≤ 0)`-normal-
form polynomials of one atom under the mixed translation. An atom with
an untranslatable side entails nothing (the W-branch discipline). -/
def atomLeZeroPolysM : Formula → List QArith.Poly
  | .atom op l r =>
    match l.asCoordPolyMixed? with
    | none => []
    | some pl =>
      match r.asCoordPolyMixed? with
      | none => []
      | some pr =>
        match op with
        | .le | .lt => [QArith.Poly.sub pl pr]
        | .ge | .gt => [QArith.Poly.sub pr pl]
        | .eq => [QArith.Poly.sub pl pr, QArith.Poly.sub pr pl]
        | .ne => []
  | _ => []

/-- Translate every decoded entry's vehicle term into the mixed space;
`none` if any vehicle is outside the mixed fragment. -/
def farkasEntriesM : List (Rat × Term) → Option (List (Rat × QArith.Poly))
  | [] => some []
  | (l, e) :: rest =>
    match e.asCoordPolyMixed?, farkasEntriesM rest with
    | some pe, some prest => some ((l, pe) :: prest)
    | _, _ => none

/-- The mixed-variable branch of the Farkas check: every polynomial —
entry vehicles, matched premise atoms, and the conclusion — rides the
mixed tagged translation; entry encoding, nonneg weights, atom matching,
and the exact combination fold mirror the sibling branches verbatim. -/
def farkasCheckM (cc : Term) (premise : Formula) (conclusion : Term) : Bool
\end{leansrc}
\declhead{BernsteinLogic.Core.Catalogue}{PFarkasData}{\badge{BadgeGreen}{statement bound}}
\begin{leansrc}
inductive PFarkasData : Hypernet → Sequent → Type where
  | mk (H : Hypernet) (seq : Sequent)
       (combinationCoeffs : Term)
       (conclusion : Term)
       -- The per-witness-node denominator bound, log₂ — claim semantics
       -- den(q) ≤ 2^denBoundLog2 over this node's witness rationals.
       -- ZERO proof-theoretic weight in v1 (no applicability axiom
       -- consumes it). Wire: u64 den_bound_log2 on tag 14 (BINARY.md §4).
       (denBoundLog2 : Nat)
       (conclShape : ∃ s, conclusion.bshape H = some s)
       (conclFin   : conclusion.finWires H)
       -- PROPOSED-B14 (approval-gated; realized provisionally per the
       -- B10/B13 pattern, 2026-07-16): every premise and conclusion wire
       -- is typed at a positive-size shape. Comparison atoms over
       -- zero-size (empty-tensor) wires hold vacuously, and a Farkas
       -- combination of two vacuous contradictory atoms — e.g. 5 ≤ w and
       -- w ≤ 3 at w : ⟨[0], Q⟩ — derives the false scalar conclusion
       -- ⊢ 2 ≤ 0 from a satisfiable premise. (The v0 singleton wire
       -- discipline excluded this by accident of strength; this is the
       -- honest general-shape guard.)
       (wiresPos : ∀ w ∈ seq.pre.freeWires ++ conclusion.freeWires,
                     ∃ τ, H.wireType w = some τ ∧
                       0 < Shape.totalSize τ.shape)
       (valid  : farkasCombinationValid combinationCoeffs seq.pre conclusion)
       (postEq : seq.post = .atom .le conclusion (.const 0))
       : PFarkasData H seq
\end{leansrc}
\declhead{BernsteinLogic.Core.Catalogue}{Formula.conjunctList}{\badge{BadgeGreen}{definition bound}}
\begin{leansrc}
/-- Flatten a conjunction tree into its leaf list (left-to-right). -/
def Formula.conjunctList : Formula → List Formula
  | .and l r => l.conjunctList ++ r.conjunctList
  | f => [f]
\end{leansrc}
\declhead{BernsteinLogic.Core.Catalogue}{atomLeZeroPolys}{\badge{BadgeGreen}{definition bound}}
\begin{leansrc}
/-- The `(e ≤ 0)`-normal-form polynomials entailed by one premise atom
(exact polynomial forms — never rescaled; F7 discipline). -/
def atomLeZeroPolys : Formula → List QArith.Poly
  | .atom .le l r => [QArith.Poly.sub l.toPoly r.toPoly]
  | .atom .lt l r => [QArith.Poly.sub l.toPoly r.toPoly]
  | .atom .ge l r => [QArith.Poly.sub r.toPoly l.toPoly]
  | .atom .gt l r => [QArith.Poly.sub r.toPoly l.toPoly]
  | .atom .eq l r => [QArith.Poly.sub l.toPoly r.toPoly,
                      QArith.Poly.sub r.toPoly l.toPoly]
  | _ => []
\end{leansrc}
\declhead{BernsteinLogic.Core.Catalogue}{farkasEntries}{\badge{BadgeGreen}{definition bound}}
\begin{leansrc}
/-- Decode the combination-coefficient Term (pinned canonical encoding:
right-nested `.add` chain of `.mul (.const λⱼ) eⱼ`, `.const 0` end). -/
def farkasEntries : Term → Option (List (Rat × Term))
  | .const c => if c = 0 then some [] else none
  | .add (.mul (.const l) e) rest =>
    (farkasEntries rest).map ((l, e) :: ·)
  | _ => none
\end{leansrc}
\declhead{BernsteinLogic.Core.Catalogue}{farkasCheckCore}{\badge{BadgeGreen}{definition bound}}
\begin{leansrc}
def farkasCheckCore (cc : Term) (premise : Formula) (conclusion : Term) : Bool :=
  match farkasEntries cc with
  | none => false
  | some entries =>
    let atoms := premise.conjunctList
    conclusion.idxFree &&
    entries.all (fun (l, e) =>
      decide (0 ≤ l) &&
      atoms.any fun a =>
        a.atomSidesIdxFree &&
        (atomLeZeroPolys a).any fun pa => decide (pa = e.toPoly)) &&
    decide ((entries.foldr
        (fun (l, e) acc => QArith.Poly.add (QArith.Poly.scale l e.toPoly) acc)
        []) = conclusion.toPoly)
\end{leansrc}
\declhead{BernsteinLogic.Core.Catalogue}{farkasCheck}{\badge{BadgeGreen}{definition bound}}
\begin{leansrc}
/-- The Farkas check: the legacy idx-free branch (exact `toPoly`
replay), OR the wire-aware idx coordinate branch (`farkasCheckW`),
OR the mixed tagged branch (`farkasCheckM`). Soundness
case-splits on the disjuncts; each branch carries its own argument
(`sound_farkas`). -/
def farkasCheck (cc : Term) (premise : Formula) (conclusion : Term) : Bool :=
  farkasCheckCore cc premise conclusion || farkasCheckW cc premise conclusion ||
  farkasCheckM cc premise conclusion
\end{leansrc}
\ruleblock{Correspondence (Steele display $\leftrightarrow$ Lean)}
\badge{BadgeAmber}{a bound Lean statement CHANGED since certification --- re-review and re-certify}
\sourcenote{The dictionary below is a reviewer-certified statement of how the mechanized statements above interpret the displayed rule: display metavariable $\mapsto$ Lean binder, each display hypothesis $\mapsto$ the premise or side condition carrying it, the display conclusion $\mapsto$ the goal --- with any deliberate difference in width between display and mechanization named outright. Certification is recorded as a hash triple over (the displayed cell, the extracted statement texts, this dictionary\textquotesingle s own text); if any of the three changes, this block turns amber until a reviewer re-certifies. The correspondence claim itself remains editorial --- reviewer-grade, not machine-checked; its machine-checked successor types each soundness statement against the rule as a structured object, and supersedes this block per rule as it lands.}
Judgment frame as in \rulename{Ax}. Distinctively, there is no
sub-derivation hypothesis at all: the displayed premise row describes
the syntactic conjunct structure of $\Phi$ itself, so the statement
consumes \texttt{seq.pre} directly and no \texttt{ih} exists.
\begin{itemize}
\item the combination $\sum_j \lambda_j \ell_j \mapsto$ the binder
  \texttt{conclusion : Term}; the multipliers $\lambda_j \mapsto$
  encoded in \texttt{combinationCoeffs : Term} under the pinned
  canonical encoding the intro describes.
\item the displayed premise row --- sub-conjunction of $\Phi$, each
  $\ell_j$ linear, $\bowtie_j \in \{\leq, =\}$, multipliers
  nonnegative and free exactly on equalities, and the identity of
  \texttt{conclusion} with the combination --- $\mapsto$ one
  hypothesis, \texttt{hvalid} $:$ \texttt{farkasCombinationValid
  combinationCoeffs seq.pre conclusion}, whose meaning this book binds
  by displaying the decoding chain in full (absorbed hypotheses, the
  \rulename{PConstExact} pattern: a successful check certifies the
  structure; the free-sign discipline is realized as both-polarity
  contribution of equality atoms). The displayed row is the
  rule\textquotesingle s published register; the realized matcher is
  wider --- see the width note below.
\item the displayed typing premises $\bshape_H(\sum \lambda_j \ell_j)
  \neq \bot$ and finite typing $\mapsto$ \texttt{hshape} (existence of
  a defined broadcast shape --- the \texttt{some}) and \texttt{hfin}
  $:$ \texttt{conclusion.finWires H}.
\item conclusion $\Phi \vdash \sum_j \lambda_j \ell_j \leq 0$
  $\mapsto$ the goal under \texttt{hpost} $:$ \texttt{seq.post = .atom
  .le conclusion (.const 0)}.
\end{itemize}
Width notes. On the side-condition surface the statement and the
display agree: \texttt{hwpos} --- every free wire of \texttt{seq.pre}
\emph{and} of the conclusion typed at a shape of strictly positive
total size --- appears inline as the display\textquotesingle s
positive-size row, and it is not optional: without it, zero-size wires
make atoms vacuously true and a combination of two vacuous atoms
derives a false scalar bound (the zero-size counterexample beside the
warrant proof). One direction is genuinely \emph{wider}: the matcher
behind \texttt{hvalid} normalizes arbitrary
two-sided comparison atoms over all five admitted comparators (strict
forms contribute their non-strict normal polynomials; equality atoms
both polarities) and the replay is degree-blind --- the displayed
linear zero-form $\{\leq, =\}$ row is the rule\textquotesingle s
published register (the note above), not the checked boundary.
The displayed cell is the rule\textquotesingle s intent; the certified
statement is that intent with the register widened exactly as stated.

\ruleblock{Lean verifier receipts}
\sourcenote{module \texttt{BernsteinLogic.Soundness}; bound build \texttt{a20d90e66c13\ldots} (bound from receipt a9c0e34d3e3c9c8c, 2026-08-03T21:30:36Z). A receipt records that named checking layers accepted a build: (module, compiled-artifact SHA-256, environment SHA-256, checker layers), minted when the checking gate ran; the store receipts attesting the bound build are listed with the layers each run credited. Nothing here asks the reader to trust the store: every layer is an independently runnable program over the quoted sources --- rebuild the module and re-run the checkers to reproduce the verdict. Receipts attesting superseded builds credit no layer and are omitted from this rendering (the count of omissions, if any, is noted below).}
\begin{receipttable}
\receiptrow{a9c0e34d3e3c9c8c}{2026-08-03T21:30:36Z}{build+decorrelated}{\badge{BadgeGreen}{current}}
\end{receipttable}
\sourcenote{21 further receipt(s) in the store attest superseded builds of this module and are omitted from this rendering.}
\sourcenote{module \texttt{BernsteinLogic.Core.Catalogue}; bound build \texttt{cbc92e976a1d\ldots} (bound from receipt 5abdb3589707b5df, 2026-08-03T21:30:33Z). A receipt records that named checking layers accepted a build: (module, compiled-artifact SHA-256, environment SHA-256, checker layers), minted when the checking gate ran; the store receipts attesting the bound build are listed with the layers each run credited. Nothing here asks the reader to trust the store: every layer is an independently runnable program over the quoted sources --- rebuild the module and re-run the checkers to reproduce the verdict. Receipts attesting superseded builds credit no layer and are omitted from this rendering (the count of omissions, if any, is noted below).}
\begin{receipttable}
\receiptrow{24804e66edba0bc9}{2026-08-03T20:25:09Z}{build+decorrelated}{\badge{BadgeGreen}{current}}
\receiptrow{5abdb3589707b5df}{2026-08-03T21:30:33Z}{build+decorrelated}{\badge{BadgeGreen}{current}}
\end{receipttable}
\sourcenote{18 further receipt(s) in the store attest superseded builds of this module and are omitted from this rendering.}
\layerline{\layeritem{build}{lake build (elaboration)}{\badge{BadgeGreen}{current}}\layeritem{replay}{leanchecker (official kernel replay)}{\badge{BadgeAmber}{stale only}}\layeritem{hardened}{safe\_verify (hardened screen)}{\badge{BadgeAmber}{stale only}}\layeritem{decorrelated}{lean4export + nanoda (decorrelated kernel)}{\badge{BadgeGreen}{current}}}
\sourcenote{\badge{BadgeGray}{comparator: none} no fixture in the three-implementation comparator harness (compiled / kernel-interpreted / reference-interpreted) exercises this rule.}
\ruleblock{HOL4 statement}
\sourcenote{theory \texttt{bern\_catalogue}, HOL4 development script \texttt{\texttt{bern\_catalogueScript.sml}} \badge{BadgeGreen}{script bound}}
\holhead{Definition farkasEntries\_def}
\begin{holsrc}
Definition farkasEntries_def:
  farkasEntries (TmConst c) = (if c = 0 then SOME [] else NONE) ∧
  farkasEntries (TmAdd (TmMul (TmConst l) e) rest) =
    (case farkasEntries rest of
     | NONE => NONE
     | SOME es => SOME ((l,e)::es)) ∧
  farkasEntries _ = NONE
End
\end{holsrc}
\holhead{Definition atomLeZeroPolys\_def}
\begin{holsrc}
Definition atomLeZeroPolys_def:
  atomLeZeroPolys (FAtom BowLe l r) = [polySub (termToPoly l) (termToPoly r)] ∧
  atomLeZeroPolys (FAtom BowLt l r) = [polySub (termToPoly l) (termToPoly r)] ∧
  atomLeZeroPolys (FAtom BowGe l r) = [polySub (termToPoly r) (termToPoly l)] ∧
  atomLeZeroPolys (FAtom BowGt l r) = [polySub (termToPoly r) (termToPoly l)] ∧
  atomLeZeroPolys (FAtom BowEq l r) = [polySub (termToPoly l) (termToPoly r);
                                       polySub (termToPoly r) (termToPoly l)] ∧
  atomLeZeroPolys _ = []
End
\end{holsrc}
\holhead{Definition farkasCheckCore\_def}
\begin{holsrc}
Definition farkasCheckCore_def:
  farkasCheckCore cc premise conclusion =
    case farkasEntries cc of
    | NONE => F
    | SOME entries =>
        (let atoms = conjunctList premise in
         termIdxFree conclusion ∧
         EVERY (λ(l,e).
           0 ≤ l ∧
           EXISTS (λa. atomSidesIdxFree a ∧
                       EXISTS (λpa. pa = termToPoly e) (atomLeZeroPolys a))
                  atoms) entries ∧
         (FOLDR (λ(l,e) acc. polyAdd (polyScale l (termToPoly e)) acc)
                [] entries =
          termToPoly conclusion))
End
\end{holsrc}
\holhead{Definition farkasCheckW\_def}
\begin{holsrc}
Definition farkasCheckW_def:
  farkasCheckW cc premise conclusion =
    case farkasEntries cc of NONE => F | SOME entries =>
    case farkasEntriesW entries of NONE => F | SOME pentries =>
    case asCoordPolyIdxW conclusion of NONE => F | SOME pc =>
      (let atoms = conjunctList premise in
       EVERY (λ(l,pe).
         0 ≤ l ∧
         EXISTS (λa. EXISTS (λpa. pa = pe) (atomLeZeroPolysW a)) atoms)
             pentries ∧
       (FOLDR (λ(l,pe) acc. polyAdd (polyScale l pe) acc) [] pentries =
        pc))
End
\end{holsrc}
\holhead{Definition farkasCheck\_def}
\begin{holsrc}
Definition farkasCheck_def:
  farkasCheck cc premise conclusion ⇔
    farkasCheckCore cc premise conclusion ∨
    farkasCheckW cc premise conclusion
End
\end{holsrc}
\holhead{Definition farkasCheckFull\_def}
\begin{holsrc}
Definition farkasCheckFull_def:
  farkasCheckFull cc premise conclusion ⇔
    farkasCheckCore cc premise conclusion ∨
    farkasCheckW cc premise conclusion ∨
    farkasCheckM cc premise conclusion
End
\end{holsrc}
\ruleblock{HOL verifier receipts}
\badge{BadgeBlue}{documentary} \texttt{hol-strict} CLEAN over the ported checker chain: theory
\texttt{bern\_catalogue}, 472 theorems, 0 cheat-tagged, 0 nonstandard
axioms (whitelist: \texttt{DISK\_THM}; the four \texttt{bool} axioms) ---
verification run of 2026-07-22 on an isolated build host, scripts
sha-matched verbatim to the sources cited here (git blob
\texttt{e41b9ed0\ldots}). The ported receipt table replays the
$\lambda = 1$ re-canonicalization accept. Downstream, the same chain
carries the checker to a CakeML binary with the
\texttt{machine\_code\_sound} tag (theory \texttt{bernCertProofARM8}).
Divergence repaired, dated: an earlier port run predated the Lean
checker\textquotesingle s third branch (the mixed tagged-union replay,
added 2026-07-19) and carried the two-branch \texttt{farkasCheck}
with a refusal-direction-only caveat; \texttt{farkasCheckFull} ---
displayed above --- is now the honest three-branch disjunction, and
the reference checker\textquotesingle s Farkas arm consumes it
verbatim (verified in the sources at the 2026-07-22 run cited here).
Semantic stratum, also ported: the soundness twin \texttt{sound\_farkas}
(theory \texttt{bern\_tr\_ready8}: \texttt{hol-strict} clean, 8
theorems, 0 offences, verification run of 2026-07-22 on an isolated
build host, content-addressed at the same sources), stated at
\texttt{farkasCheckFull} --- matching the widened checker arm verbatim
--- and consumed by the certified-fragment keystone
(\texttt{checkPf\_sound\_frag}, theory \texttt{bern\_tr\_splice}, 20
theorems, 0 offences, same run): a reference-checker accept inside the
proven wire fragment yields sequent validity.

\sourcenote{HOL receipts are documentary: quoted results of \texttt{hol-strict} runs
and census rows from the HOL4 development\textquotesingle s build
evidence, not minted receipt artefacts. A HOL receipt store (the Lean
store\textquotesingle s twin) is named future work.
}
\ruleblock{Proof}

\sourcenote{Genre: admissible rule (tier 2) --- a proof-tree-rewriting
recipe; the general admissible template is stated in the \rulename{Norm}
section and instantiated here. \rulename{Farkas} is admissible by
reduction to verified machinery: the affine-combination replay is
checked computation, and one verified lemma warrants every accepted
node.}

\begin{brecipe}[the \rulename{Farkas} instance]\label{rec:fk-instance}
\rulename{Farkas} has no proof premises --- its inputs are
sub-conjuncts of $\Phi$ itself, not sub-derivations:
\[
\inferrule*[Right=\rulename{Farkas}]{
  \{a_j : \ell_j \bowtie_j r_j\}_j \subseteq \Phi \text{ conjunct-wise,}
  \ p_j \ \text{\upshape their } ({\cdot}) \leq 0\text{\upshape-normal
  forms} \\
  \lambda_j \geq 0 \text{ (free-signed on equalities)} \\\\
  \textstyle\sum_j \lambda_j p_j \equiv c
  \ \text{\upshape(exact canonical fold)} \\
  c \ \text{\upshape shape-stable, finitely typed;}\;
  \wires(\Phi) \cup \wires(c) \ \text{\upshape positive-size}
}{\Phi \vdash c \leq 0}
\qquad\rightsquigarrow\qquad
\text{\upshape the checker replays the fold;}\;
\text{\texttt{sound\_farkas} stands as warrant.}
\]
The rewrite target contains zero catalogue nodes: the entire content is
the checked combination (decoded, sign-gated, atom-sourced, folded ---
the fragment display quoted above walks the chain) plus the semantic
argument of Proposition~\ref{prop:fk-sound}. The matcher is wider than
the catalogued cell's published register (linear zero-form
$\{\leq, =\}$ rows): it normalizes arbitrary two-sided comparison atoms
of the five admitted comparators --- strict forms contribute their
non-strict normal forms, semantically (Proposition~\ref{prop:fk-comb})
--- and the replay is degree-blind (the book note). The catalogued
stratum
records an unpacking sketch through \rulename{BernCone} at degree $1$;
the shipped warrant is direct, at all tensor shapes, heterogeneous
premise-atom shapes included.
\end{brecipe}

\begin{bprop}[nonnegative combination, per entry]\label{prop:fk-comb}
Work in an ordered field. Suppose values $a_j$ satisfy $a_j \leq 0$,
and weights $\lambda_j \geq 0$. Then
$\sum_j \lambda_j a_j \leq 0$. Moreover the $({\cdot} \leq 0)$-normal
reading of each premise atom delivers such values: an atom
$l \bowtie r$ with $\bowtie \in \{\leq, <\}$ yields
$l - r \leq 0$; with $\bowtie \in \{\geq, >\}$ it yields
$r - l \leq 0$; and an equality yields \emph{both}
$l - r \leq 0$ and $r - l \leq 0$ --- so a free-sign multiplier
$\nu$ on an equality is realized as the nonnegative weight $|\nu|$ on
the polarity matching $\nu$'s sign.
\end{bprop}
\begin{proof}
Each $\lambda_j a_j \leq \lambda_j \cdot 0 = 0$ (multiplication by a
nonnegative element preserves $\leq$), and a finite sum of nonpositive
elements is nonpositive, by induction. The polarity readings are the
definitions of the normal forms; a holding strict atom \emph{implies}
its non-strict normal form --- the matcher emits the same polynomial
for $<$ as for $\leq$, the weakening being a fact about holding, not a
checker step. The
equality discipline is the textbook sign convention for multipliers of
equality constraints \cite[\S7.3]{schrijver1986}, realized here
syntactically: the matcher sees both polarities, so a negative
multiplier is unrepresentable rather than checked for.
\end{proof}

\begin{bprop}[semantic warrant]\label{prop:fk-sound}
Suppose the payload's combination check accepts (any branch), the
conclusion term $c$ is shape-stable and finitely typed, and every wire
of $\Phi$ and $c$ is typed at a positive-size shape. If
$\sigma \models \Phi$ for an admissible $\sigma$, then
$\sigma \models c \leq 0$.
\end{bprop}
\begin{proof}
The shape and finiteness fields give $c$ a defined evaluation; fix a
conclusion entry $i$. The accepted branch fixes a variable space and a
per-entry valuation for it: the idx-free branch reads each bare wire
\emph{clamped} at entry $i$ --- a rank-0 wire reads its scalar; a
positive-rank wire of $m$ entries reads entry $i$ when $i < m$ and
entry $0$ otherwise, so the read is always an honest in-range entry of
the same tensor, and is the plain positional read wherever $i$ is in
range (a read that serves heterogeneous premise
atom shapes at every conclusion index); the wire-aware branch reads
each projection $w[j]$ at its fixed coordinate; the mixed branch does
both at once, on disjoint variable classes. Two facts from the
Setting's broadcast discipline make the clamped reads honest and exact.
First, a holding atom never compares unequal positive ranks, and a
term's defined broadcast shape forces every positive-rank wire in it to
carry that very shape (only rank-0 joins across shapes) --- so all
positive-rank wires of one premise atom share one shape, and the clamp
lands them at one \emph{common} honest entry of that atom, rank-0 wires
reading their scalar under any index; here positive size is
load-bearing: an atom over positively-sized wires has at least one
honest entry for the clamp to land on, so its per-entry normal form is
genuinely nonpositive at
the branch's valuation (\autoref{rem:fk-b14} shows what happens
otherwise). Second, the same discipline pins every positive-rank wire
of $c$ at $c$'s own shape, so at conclusion entry $i$ the clamp is
exact and the fold's value at the branch's valuation \emph{is} entry
$i$. Thus each sourced entry polynomial is the normal form of a premise
conjunct that holds at $\sigma$, read per-entry;
Proposition~\ref{prop:fk-comb} bounds the weighted fold by
$0$; the certified fold identity says the fold \emph{is} $c$'s
polynomial in the branch's variable space; evaluating at the branch's
valuation gives entry $i$ of $c$'s tensor $\leq 0$. Entries were
arbitrary, so the atom $c \leq 0$ holds. \texttt{sound\_farkas}
case-splits on the branch disjunct and runs exactly this argument
three times.
\end{proof}

\begin{bremark}[why positive-size wires are not decoration]
\label{rem:fk-b14}
Comparison atoms hold per-entry, and a zero-size wire has no entries:
at $w$ of shape $[0]$, the atoms $5 \leq w$ and $w \leq 3$ both hold
\emph{vacuously} --- yet their $({\cdot} \leq 0)$-normal forms
$5 - w$ and $w - 3$ are perfectly good polynomials, and with
$\lambda = 1$ each they fold to the constant $2$. Without the
positive-size condition the rule derives $\vdash 2 \leq 0$ from a
satisfiable premise: the classic mismatch between empty-product
semantics and nonempty algebra, and the reason the field exists. (A
narrower wire discipline once masked the failure by accident of
strength; the field is the honest general-shape guard.)
\end{bremark}

\begin{bremark}[why the conclusion's shape fields are not decoration]
\label{rem:fk-b13}
A conclusion mentioning an untyped wire --- say
$c' = 0 \cdot w_{999} + c$ with $w_{999}$ untyped --- has a correct
canonical polynomial: the dead monomial vanishes in canonical form, so
the fold identity holds exactly as it does for $c$. But evaluation is
structural on the \emph{term}, not on its canonical form --- it still
reads $w_{999}$, which is untyped --- so the conclusion never
evaluates, and the concluded atom never holds at any valuation. The shape-stability and finiteness
fields make ``the fold is $c$'s polynomial'' imply ``$c$
evaluates'', closing the gap between syntax and semantics.
\end{bremark}

\begin{bremark}[the three branches are one algorithm]
\label{rem:fk-branches}
The branches differ only in the variable space where polynomials are
compared: canonical wire-keyed forms; projection coordinates keyed by
an injective (wire, coordinate) pairing --- distinct wires \emph{must}
get distinct variables, since a composite may mix projections of
several wires and conflating $w_0[j]$ with $w_2[j]$ would let
semantically different polynomials collide (a one-line witness: from
the holding premise $w_0[0] \leq 0$ at $w_0[0] = -1$, $w_2[0] = 1$,
the conflated space accepts the false conclusion $w_2[0] \leq 0$); and
the mixed tagged union,
where bare wires and projections are disjoint classes --- in
particular a bare $w$ is never identified with $w[0]$: the
identification variant admits a false sequent (witness: $w = (-1, 1)$
satisfies $w[0] \leq 0$, and the identified space would conclude the
whole-wire bound $w \leq 0$, false at the second entry), so the
tagging is a
soundness fact, not a convenience. Entry encoding, sign gates, atom
sourcing, and the exact fold are the same gates, replayed branch by
branch in its own variable space (the branch definitions cited at the
anchors), and each
branch carries its own per-entry valuation in the soundness proof.
\end{bremark}

\begin{bremark}[relation to the literature]
The linear ancestor is Farkas' lemma \cite{farkas1902}; the multiplier
sign conventions ($\lambda \geq 0$ on inequalities, free on equalities)
are \cite[\S7]{schrijver1986}. As everywhere in this calculus, only
the sound (trivial) direction is claimed: the classical lemma's
content --- that every valid affine consequence \emph{has} such a
combination --- is a completeness statement about producers, never
consulted by the kernel. Degree-blindness of the replay places the
rule's polynomial generalizations in the Krivine--Handelman family
\cite{krivine1964,handelman1988}, reached in this catalogue through
\rulename{Syzygy} and \rulename{BernCone} rather than by widening this
rule.
\end{bremark}

\begin{anchors}
Recipe \ref{rec:fk-instance} warrant & \texttt{sound\_farkas} (branch case split); \texttt{farkasCombinationValid}\\
entry decoding (pinned encoding) & \texttt{farkasEntries}\\
atom normal forms, both polarities on equality & \texttt{atomLeZeroPolys} (idx-free); \texttt{atomLeZeroPolysW}, \texttt{atomLeZeroPolysM} (widened branches)\\
the three branches & \texttt{farkasCheckCore}, \texttt{farkasCheckW}, \texttt{farkasCheckM}; disjunction \texttt{farkasCheck}\\
projection variable spaces (Remark \ref{rem:fk-branches}) & \texttt{wireCoordKey}, \texttt{Term.asCoordPolyIdxW?}, \texttt{Term.asCoordPolyMixed?}\\
Proposition \ref{prop:fk-comb} fold & \texttt{farkas\_fold\_eval\_nonpos}, \texttt{farkas\_fold\_eval\_nonpos\_polys}\\
per-entry valuations & the clamped, coordinate, and mixed entry valuations consumed by \texttt{sound\_farkas}'s three cases\\
positive-size guard (Remark \ref{rem:fk-b14}) & payload field \texttt{wiresPos} (\texttt{PFarkasData})\\
conclusion evaluability (Remark \ref{rem:fk-b13}) & payload fields \texttt{conclShape}, \texttt{conclFin}\\
wire placement & \texttt{RuleTag.farkas}\\
\end{anchors}


\rulesection{18}{Subdivide}{admissible rule}
\label{blcatrule-Subdivide}
\ruleblock{The rule (Steele metanotation)}
\sourcenote{as realized in the mechanized development; source row: \texttt{\texttt{rules.toml}} --- \emph{Subdivide}. The metanotation statement is typeset by hand for notational quality; its mechanized ground truth is the Lean statement below, extracted from the sources at every build of this book and hash-bound. The rule \emph{roster} is machine-checked (every mechanized rule tag round-trips to exactly one catalogue row); the statement \emph{text} is editorial, checked by eye against the extracted Lean beside it.}
\begin{mathpar}
\inferrule*[Right=\rulename{Subdivide}]{
  \Phi \vdash x \in \mathcal{R}(\Delta, S)
  \\
  \{(\Delta_k, S_k)\}_{k=1}^N \text{ finite, with cover side-condition } \mathcal{R}(\Delta, S) \subseteq \textstyle\bigcup_k \mathcal{R}(\Delta_k, S_k)
  \\
  \forall k.\; \Phi \;\wedge\; (x \in \mathcal{R}(\Delta_k, S_k)) \;\vdash\; \Psi
}{
  \Phi \vdash \Psi
}
\end{mathpar}

\par\medskip
\rulename{Subdivide} is the case-split rule of region reasoning: from a
proven membership $\Phi \vdash x \in \mathcal{R}(\Delta, S)$, a finite
family of branch regions covering the parent, and one sub-proof of the
goal per branch (each under the strengthened premise
$\Phi \wedge x \in \mathcal{R}(\Delta_k, S_k)$), it concludes
$\Phi \vdash \Psi$ on the whole parent region. It is degree
escalation\textquotesingle s partner and load-bearing for every non-toy
proof: where a single \rulename{BernCone} certificate at feasible degree
is too coarse, a producer splits the region, certifies each piece at low
degree, and lets \rulename{Subdivide} reassemble the global conclusion
--- the branch-and-bound shape that makes Bernstein enclosures converge.

What the checker recomputes, given the payload (parent region, branch
family, per-branch sub-proofs): that the parent region reads as a closed
interval --- ambient dimension 1, support exactly two vertices at
distinct coordinates, empty extension set; that every branch region
reads as a closed interval and is tied to its own extension set; and
that the branch intervals \emph{chain} across the parent interval --- a
greedy reach, iterated to its fixpoint, that extends past any interval
whose left endpoint has already been reached and must arrive at the
parent\textquotesingle s right endpoint. Everything beyond this
fragment (extension-ray parents or branches, regions of dimension
$\geq 2$, supports of more than two vertices) is refused, never
accepted on faith. The repair for that narrowness is
\rulename{ConvIntro}: $n$-dimensional box grids arise by
\emph{nesting} per-coordinate splits across several rule
applications, not inside one --- each branch establishes its
coordinate facts as order atoms, and \rulename{ConvIntro} converts
them back into the region atom the next coordinate\textquotesingle s
split consumes, so the one-dimensional rule reaches every box grid
without widening its checker.

Each hypothesis is load-bearing. Drop the cover check and a producer
omits the branch where the goal fails --- a gap family such as
$[1, \frac{4}{3}] \cup [\frac{3}{2}, 2]$ over $[1, 2]$ proves the goal
everywhere except where it is false. Drop the branch-region tie and the
chained intervals belong to different regions than the ones the
sub-proofs assume. Drop the branch-premise shape and a sub-proof could
be consulted outside the piece it was proven on. Drop the
branch-conclusion tie and the branches may prove any formulas whatever
--- $\Psi$ never enters their obligations while the rule concludes it
anyway. Drop the parent region premise and a valuation satisfying
$\Phi$ with the point outside the parent falls through every branch:
the cover says nothing about it, yet the rule still concludes $\Psi$
there. The certificate slot
for cover evidence rides the wire but is deliberately \emph{not
consulted}: the checker recomputes the cover from the branch regions
themselves, so ignoring the slot can only make checking stricter, never
unsound.

\cataloguenote{The region variable itself plays no role in the cover geometry: the
soundness argument feeds the \emph{same} evaluated point to the parent
and to the selected branch, so no tie between the term and the cover is
needed --- the cover is a fact about intervals, not about wires. A
future interpretation may read per-gap infeasibility certificates out of
the wire slot; the catalogued stratum records the corresponding
unpacking sketch (per-branch cuts discharged through
\rulename{Farkas} and $\vee$-elimination).
}
\ruleblock{Lean statement}
\declhead{BernsteinLogic.Soundness}{sound\_subdivide}{\badge{BadgeGreen}{statement bound}}
\begin{leansrc}
/-- Soundness of Subdivide (the covering argument): if the region
premise puts the evaluated point in the parent region, the checker's
cover chains the branch intervals across the parent interval, and every
branch sequent is valid on its region, then the conclusion holds on the
parent. -/
theorem sound_subdivide {H : Hypernet} {seq regionPremise : Sequent}
    {nBranches : Nat} {branchPremise : Fin nBranches → Sequent}
    {termX : Term} {simplex : Simplex} {S : ExtensionSet}
    {branchSimplex : Fin nBranches → Simplex}
    {branchS : Fin nBranches → ExtensionSet}
    {coverCertificate : Term}
    (ihRegion : regionPremise.semanticallyValid H)
    (ihBranch : ∀ k, (branchPremise k).semanticallyValid H)
    (hpre : seq.pre = regionPremise.pre)
    (hregionPost : regionPremise.post = .inRegion termX simplex S)
    (hbranchPre : ∀ k, (branchPremise k).pre
                   = .and regionPremise.pre
                          (.inRegion termX (branchSimplex k) (branchS k)))
    (hbranchPost : ∀ k, (branchPremise k).post = seq.post)
    (hcover : coverInfeasibilityValid termX simplex S nBranches
                branchSimplex branchS coverCertificate) :
    seq.semanticallyValid H
\end{leansrc}
\declhead{BernsteinLogic.Soundness}{interval1D\_elim}{\badge{BadgeGreen}{statement bound}}
\begin{leansrc}
/-- Elim: a `MaskRegion` witness against a successful interval reading
lands the point's coordinate in `[lo, hi]`, and the interval is proper. -/
theorem interval1D_elim {S : ExtensionSet} {lo hi : ℚ}
    (hint : S.interval1D? = some (lo, hi))
    {x : Fin S.simplex.complex.ambDim → ℚ}
    (hmask : MaskRegion S.simplex.complex S.simplex.contains S.contains x)
    (j0 : Fin S.simplex.complex.ambDim) :
    lo ≤ x j0 ∧ x j0 ≤ hi ∧ lo < hi
\end{leansrc}
\declhead{BernsteinLogic.Soundness}{interval1D\_intro}{\badge{BadgeGreen}{statement bound}}
\begin{leansrc}
/-- Intro: a coordinate inside a successful interval reading assembles a
`MaskRegion` witness (the extension set is empty in this fragment, so
both barycentric weights come out nonnegative). -/
theorem interval1D_intro {S : ExtensionSet} {lo hi : ℚ}
    (hint : S.interval1D? = some (lo, hi))
    {x : Fin S.simplex.complex.ambDim → ℚ}
    (j0 : Fin S.simplex.complex.ambDim)
    (hxlo : lo ≤ x j0) (hxhi : x j0 ≤ hi) :
    MaskRegion S.simplex.complex S.simplex.contains S.contains x
\end{leansrc}
\declhead{BernsteinLogic.Soundness}{ChainInv}{\badge{BadgeGreen}{statement bound}}
\begin{leansrc}
/-- The greedy chain invariant: the reach never sits below `lo`, and
once it has moved past `lo`, every point of `[lo, reach]` lies in some
interval of the FULL list. -/
def ChainInv (ivs : List (ℚ × ℚ)) (lo r : ℚ) : Prop
\end{leansrc}
\declhead{BernsteinLogic.Soundness}{chainInv\_extendPass}{\badge{BadgeGreen}{statement bound}}
\begin{leansrc}
/-- One pass preserves the invariant (stated over a sublist so the foldl
induction goes through; the covering intervals always come from the full
list). -/
theorem chainInv_extendPass (ivs : List (ℚ × ℚ)) (lo : ℚ) :
    ∀ (sub : List (ℚ × ℚ)), (∀ iv ∈ sub, iv ∈ ivs) →
    ∀ (r : ℚ), ChainInv ivs lo r →
    ChainInv ivs lo (coverExtendPass sub r)
\end{leansrc}
\declhead{BernsteinLogic.Soundness}{chainInv\_reach}{\badge{BadgeGreen}{statement bound}}
\begin{leansrc}
theorem chainInv_reach (ivs : List (ℚ × ℚ)) (lo : ℚ) :
    ∀ (n : Nat) (r : ℚ), ChainInv ivs lo r →
    ChainInv ivs lo (coverReach n ivs r)
\end{leansrc}
\declhead{BernsteinLogic.Soundness}{coverChain\_covers}{\badge{BadgeGreen}{statement bound}}
\begin{leansrc}
/-- Soundness of the greedy chain check: on success, every point of
`[lo, hi]` lies in some interval of the list. -/
theorem coverChain_covers {ivs : List (ℚ × ℚ)} {lo hi : ℚ} (hlt : lo < hi)
    (hcheck : coverChainCheck ivs lo hi = true) :
    ∀ x : ℚ, lo ≤ x → x ≤ hi → ∃ iv ∈ ivs, iv.1 ≤ x ∧ x ≤ iv.2
\end{leansrc}
\par\medskip
\textbf{What the cover check means --- and what it does not.}\enspace
The side condition appears in the payload as the single name
\texttt{coverInfeasibilityValid}, so this book binds its meaning by
displaying the definition chain in full, from the checker sources.
Reading the chain: \texttt{ExtensionSet.interval1D?} is the
closed-interval reading of a region --- it succeeds only on the checkable
fragment (dimension 1, two-vertex support at distinct coordinates,
empty extension set), returning sorted endpoints, and everything else
reads \texttt{none}, which the cover check turns into refusal.
\texttt{coverExtendPass} is one greedy pass: extend the reach past any
interval whose left endpoint has been reached;
\texttt{coverReach} iterates it, and as many passes as there are
intervals suffice for the fixpoint (a productive pass strictly
increases the reach onto a new right endpoint, of which there are only
finitely many). \texttt{coverChainCheck} asks whether the reach from
the parent\textquotesingle s left endpoint arrives at its right
endpoint; \texttt{coverCheck} assembles the whole condition, tying each
branch\textquotesingle s extension set to its simplex on the way.

What a skeptical reader should confirm by eye: the region term and the
wire-borne cover certificate are \emph{ignored} --- the relation
discards both arguments, and the cover is recomputed from the branch
regions alone (the recompute pattern; a forged certificate cannot
influence the verdict, and a missing one costs nothing); no valuation
is consulted anywhere; and the branch sub-proofs enter soundness, not
the cover computation --- the checker decides pure interval geometry.

\declhead{BernsteinLogic.Core.Catalogue}{ExtensionSet.interval1D?}{\badge{BadgeGreen}{definition bound}}
\begin{leansrc}
/-- The closed-interval reading of a region `R(Δ, S)`: succeeds iff the
ambient dimension is 1, the simplex's support is exactly two vertices at
distinct coordinates, and the extension set is empty — returning the
region's endpoints, sorted (`lo < hi` by construction). Everything else
reads `none` (and the cover check then rejects). -/
def ExtensionSet.interval1D? (S : ExtensionSet) : Option (Rat × Rat) :=
  if h : S.simplex.complex.ambDim = 1 then
    if S.isEmptyB then
      match (List.finRange S.simplex.complex.nVertices).filter
          (fun v => S.simplex.contains v) with
      | [a, b] =>
        let j0 : Fin S.simplex.complex.ambDim := ⟨0, by omega⟩
        let V₀ := S.simplex.complex.coord a j0
        let V₁ := S.simplex.complex.coord b j0
        if V₀ < V₁ then some (V₀, V₁)
        else if V₁ < V₀ then some (V₁, V₀)
        else none
      | _ => none
    else none
  else none
\end{leansrc}
\declhead{BernsteinLogic.Core.Catalogue}{coverExtendPass}{\badge{BadgeGreen}{definition bound}}
\begin{leansrc}
/-- One greedy pass of the 1-D chain cover: extend the reach past any
interval whose left endpoint has already been reached. -/
def coverExtendPass (ivs : List (Rat × Rat)) (r : Rat) : Rat :=
  ivs.foldl (fun r iv => if iv.1 ≤ r ∧ r < iv.2 then iv.2 else r) r
\end{leansrc}
\declhead{BernsteinLogic.Core.Catalogue}{coverReach}{\badge{BadgeGreen}{definition bound}}
\begin{leansrc}
/-- Iterated greedy reach. `ivs.length` passes suffice for the fixpoint:
a productive pass strictly increases the reach onto a new interval's
right endpoint (of which there are at most `ivs.length`), and an
unproductive pass is already the fixpoint. -/
def coverReach : Nat → List (Rat × Rat) → Rat → Rat
  | 0, _, r => r
  | n + 1, ivs, r => coverReach n ivs (coverExtendPass ivs r)
\end{leansrc}
\declhead{BernsteinLogic.Core.Catalogue}{coverChainCheck}{\badge{BadgeGreen}{definition bound}}
\begin{leansrc}
/-- Greedy 1-D chain check: `[lo, hi] ⊆ ⋃ ivs`, all intervals closed.
Soundness: `coverChain_covers` (Soundness §S5′). -/
def coverChainCheck (ivs : List (Rat × Rat)) (lo hi : Rat) : Bool :=
  decide (hi ≤ coverReach ivs.length ivs lo)
\end{leansrc}
\declhead{BernsteinLogic.Core.Catalogue}{coverCheck}{\badge{BadgeGreen}{definition bound}}
\begin{leansrc}
/-- The Subdivide cover computation (fragment documented at the section
header): the parent region is an interval, every branch region is an
interval tied to its own extension set, and the branch intervals chain
across the parent's. -/
def coverCheck (simplex : Simplex) (S : ExtensionSet) (nBranches : Nat)
    (branchSimplex : Fin nBranches → Simplex)
    (branchS : Fin nBranches → ExtensionSet) : Bool :=
  match S.interval1D? with
  | none => false
  | some (lo, hi) =>
    decide (S.simplex = simplex) &&
    ((List.finRange nBranches).all fun k =>
      decide ((branchS k).simplex = branchSimplex k) &&
      ((branchS k).interval1D?).isSome) &&
    coverChainCheck
      ((List.finRange nBranches).filterMap fun k => (branchS k).interval1D?)
      lo hi
\end{leansrc}
\declhead{BernsteinLogic.Core.Catalogue}{coverInfeasibilityValid}{\badge{BadgeGreen}{definition bound}}
\begin{leansrc}
def coverInfeasibilityValid
    (_termX : Term) (simplex : Simplex) (S : ExtensionSet)
    (nBranches : Nat)
    (branchSimplex : Fin nBranches → Simplex)
    (branchS : Fin nBranches → ExtensionSet)
    (_coverCert : Term) : Prop :=
  coverCheck simplex S nBranches branchSimplex branchS = true
\end{leansrc}
\ruleblock{Correspondence (Steele display $\leftrightarrow$ Lean)}
\badge{BadgeGreen}{reviewer-certified against the displayed texts}
\sourcenote{The dictionary below is a reviewer-certified statement of how the mechanized statements above interpret the displayed rule: display metavariable $\mapsto$ Lean binder, each display hypothesis $\mapsto$ the premise or side condition carrying it, the display conclusion $\mapsto$ the goal --- with any deliberate difference in width between display and mechanization named outright. Certification is recorded as a hash triple over (the displayed cell, the extracted statement texts, this dictionary\textquotesingle s own text); if any of the three changes, this block turns amber until a reviewer re-certifies. The correspondence claim itself remains editorial --- reviewer-grade, not machine-checked; its machine-checked successor types each soundness statement against the rule as a structured object, and supersedes this block per rule as it lands.}
Judgment frame as in \rulename{Ax}; sub-derivation premises as at
\rulename{Weaken}.
\begin{itemize}
\item $x \mapsto$ \texttt{termX : Term} (a free term --- unlike
  \rulename{BernCone}, no wire tie is needed here: the note above
  explains why); $(\Delta, S) \mapsto$ \texttt{simplex},
  \texttt{S}; the branch family $\{(\Delta_k, S_k)\}_{k=1}^N
  \mapsto$ \texttt{branchSimplex}, \texttt{branchS} $:$ \texttt{Fin
  nBranches → \_}.
\item region premise $\Phi \vdash x \in \mathcal{R}(\Delta, S)$
  $\mapsto$ \texttt{ihRegion} with \texttt{hregionPost} $:$
  \texttt{regionPremise.post = .inRegion termX simplex S}; $\Phi$
  shared via \texttt{hpre}.
\item the branch judgments $\forall k.\ \Phi \wedge (x \in
  \mathcal{R}(\Delta_k, S_k)) \vdash \Psi$ $\mapsto$
  \texttt{ihBranch} with \texttt{hbranchPre} pinning each branch
  premise to the displayed conjunction in the displayed order ($\Phi$
  left) and \texttt{hbranchPost} tying every branch conclusion to
  \texttt{seq.post} --- $\Psi$ is whatever formula the conclusion
  sequent carries; nothing constrains its shape.
\item the cover side-condition $\mathcal{R}(\Delta, S) \subseteq
  \bigcup_k \mathcal{R}(\Delta_k, S_k)$ $\mapsto$ \texttt{hcover} $:$
  \texttt{coverInfeasibilityValid}, i.e.\ \texttt{coverCheck = true}:
  the parent extension set must present as a 1-dimensional interval
  (\texttt{ExtensionSet.interval1D?} --- outside that fragment the
  check refuses, so the certified reach of the displayed
  set-inclusion is the interval fragment, the
  \rulename{ConvElim} class), each branch\textquotesingle s
  simplex-tie is checked, and the greedy chain
  (\texttt{coverExtendPass}/\texttt{coverReach}/\texttt{coverChainCheck},
  certified alongside) verifies the branch intervals chain across the
  parent interval --- a decidable \emph{sufficient} witness of the
  displayed inclusion, conservative like every refusal boundary.
\item consumed by nothing: the payload\textquotesingle s cover
  certificate term rides into \texttt{coverInfeasibilityValid} as an
  ignored argument (the binder is underscore-named in the displayed
  definition) --- retained for the per-gap-infeasibility
  interpretation the note records; the realized checker recomputes the
  cover itself and reads no producer-supplied cover data.
\end{itemize}
Width note: within the interval fragment neither side is wider ---
every other hypothesis is consumed, and the displayed rule\textquotesingle
s ``finite family\textquotesingle\textquotesingle{} is carried by the
\texttt{Fin nBranches} index type itself.

\ruleblock{Lean verifier receipts}
\sourcenote{module \texttt{BernsteinLogic.Soundness}; bound build \texttt{a20d90e66c13\ldots} (bound from receipt a9c0e34d3e3c9c8c, 2026-08-03T21:30:36Z). A receipt records that named checking layers accepted a build: (module, compiled-artifact SHA-256, environment SHA-256, checker layers), minted when the checking gate ran; the store receipts attesting the bound build are listed with the layers each run credited. Nothing here asks the reader to trust the store: every layer is an independently runnable program over the quoted sources --- rebuild the module and re-run the checkers to reproduce the verdict. Receipts attesting superseded builds credit no layer and are omitted from this rendering (the count of omissions, if any, is noted below).}
\begin{receipttable}
\receiptrow{a9c0e34d3e3c9c8c}{2026-08-03T21:30:36Z}{build+decorrelated}{\badge{BadgeGreen}{current}}
\end{receipttable}
\sourcenote{21 further receipt(s) in the store attest superseded builds of this module and are omitted from this rendering.}
\sourcenote{module \texttt{BernsteinLogic.Core.Catalogue}; bound build \texttt{cbc92e976a1d\ldots} (bound from receipt 5abdb3589707b5df, 2026-08-03T21:30:33Z). A receipt records that named checking layers accepted a build: (module, compiled-artifact SHA-256, environment SHA-256, checker layers), minted when the checking gate ran; the store receipts attesting the bound build are listed with the layers each run credited. Nothing here asks the reader to trust the store: every layer is an independently runnable program over the quoted sources --- rebuild the module and re-run the checkers to reproduce the verdict. Receipts attesting superseded builds credit no layer and are omitted from this rendering (the count of omissions, if any, is noted below).}
\begin{receipttable}
\receiptrow{24804e66edba0bc9}{2026-08-03T20:25:09Z}{build+decorrelated}{\badge{BadgeGreen}{current}}
\receiptrow{5abdb3589707b5df}{2026-08-03T21:30:33Z}{build+decorrelated}{\badge{BadgeGreen}{current}}
\end{receipttable}
\sourcenote{18 further receipt(s) in the store attest superseded builds of this module and are omitted from this rendering.}
\layerline{\layeritem{build}{lake build (elaboration)}{\badge{BadgeGreen}{current}}\layeritem{replay}{leanchecker (official kernel replay)}{\badge{BadgeAmber}{stale only}}\layeritem{hardened}{safe\_verify (hardened screen)}{\badge{BadgeAmber}{stale only}}\layeritem{decorrelated}{lean4export + nanoda (decorrelated kernel)}{\badge{BadgeGreen}{current}}}
\sourcenote{\badge{BadgeGray}{comparator: none} no fixture in the three-implementation comparator harness (compiled / kernel-interpreted / reference-interpreted) exercises this rule.}
\ruleblock{HOL4 statement}
\sourcenote{theory \texttt{bern\_catalogue}, HOL4 development script \texttt{\texttt{bern\_catalogueScript.sml}} \badge{BadgeGreen}{script bound}}
\holhead{Definition extInterval1D\_def}
\begin{holsrc}
Definition extInterval1D_def:
  extInterval1D es =
    if es.simplex.complex.ambDim = 1 then
      (if extIsEmptyB es then
         (case FILTER (λv. simplexContainsNat es.simplex v)
                      (rangeL es.simplex.complex.nVertices) of
          | [a; b] =>
              (let v0 = getD (getD es.simplex.complex.coord a []) 0 (0:rat) in
               let v1 = getD (getD es.simplex.complex.coord b []) 0 (0:rat) in
               if v0 < v1 then SOME (v0, v1)
               else if v1 < v0 then SOME (v1, v0)
               else NONE)
          | _ => NONE)
       else NONE)
    else NONE
End
\end{holsrc}
\holhead{Definition coverExtendPass\_def}
\begin{holsrc}
Definition coverExtendPass_def:
  coverExtendPass ivs (r:rat) =
    FOLDL (λr iv. if FST iv ≤ r ∧ r < SND iv then SND iv else r) r ivs
End
\end{holsrc}
\holhead{Definition coverReach\_def}
\begin{holsrc}
Definition coverReach_def:
  coverReach 0 ivs (r:rat) = r ∧
  coverReach (SUC n) ivs r = coverReach n ivs (coverExtendPass ivs r)
End
\end{holsrc}
\holhead{Definition coverChainCheck\_def}
\begin{holsrc}
Definition coverChainCheck_def:
  coverChainCheck ivs (lo:rat) hi = (hi ≤ coverReach (LENGTH ivs) ivs lo)
End
\end{holsrc}
\holhead{Definition coverCheck\_def}
\begin{holsrc}
Definition coverCheck_def:
  coverCheck sx es branches =
    case extInterval1D es of
    | NONE => F
    | SOME (lo, hi) =>
        (es.simplex = sx ∧
         EVERY (λ(bsx,bes). bes.simplex = bsx ∧ IS_SOME (extInterval1D bes))
               branches ∧
         coverChainCheck
           (FLAT (MAP (λ(bsx,bes). case extInterval1D bes of
                                   | NONE => []
                                   | SOME iv => [iv]) branches))
           lo hi)
End
\end{holsrc}
\ruleblock{HOL verifier receipts}
\badge{BadgeBlue}{documentary} \texttt{hol-strict} CLEAN over the ported checker chain: theory
\texttt{bern\_catalogue}, 472 theorems, 0 cheat-tagged, 0 nonstandard
axioms (whitelist: \texttt{DISK\_THM}; the four \texttt{bool} axioms) ---
verification run of 2026-07-22 on an isolated build host, scripts
sha-matched verbatim to the sources cited here (git blob
\texttt{e41b9ed0\ldots}). The ported receipt table replays the cover
battery: the positive chain
($[1, \frac{3}{2}] \cup [\frac{3}{2}, 2]$ over $[1, 2]$) and four
negative controls --- the gap mutant, a point-supported branch, an
extension-ray parent, and untied branch simplices. Representation seam,
inventoried in the script: the indexed branch family becomes one list
of (simplex, extension-set) pairs in index order, and the Lean
relation\textquotesingle s two inert arguments (region term, cover
certificate) are not reproduced --- both value-preserving deltas under
the port\textquotesingle s seam discipline.

\sourcenote{HOL receipts are documentary: quoted results of \texttt{hol-strict} runs
and census rows from the HOL4 development\textquotesingle s build
evidence, not minted receipt artefacts. A HOL receipt store (the Lean
store\textquotesingle s twin) is named future work.
}
\ruleblock{Proof}

\sourcenote{Genre: admissible rule (tier 2) --- a proof-tree-rewriting
recipe; the general admissible template is stated in the \rulename{Norm}
section and instantiated here. \rulename{Subdivide} is the genuinely
one-node-to-many member of this batch: its rewrite is a case analysis
over the branch sub-proofs, and the checked cover condition is what
makes the case analysis exhaustive.}

\begin{brecipe}[the \rulename{Subdivide} instance]\label{rec:sd-instance}
The node carries a region premise and one sub-proof per branch:
\[
\inferrule*[Right=\rulename{Subdivide}]{
  \mathcal{T}_0 \rhd \bigl(\Phi \vdash x \in \mathcal{R}(\Delta, S)\bigr) \\\\
  \mathcal{T}_k \rhd \bigl(\Phi \wedge (x \in \mathcal{R}(\Delta_k, S_k))
    \vdash \Psi\bigr) \quad (k = 1, \dots, N) \\\\
  \mathcal{R}(\Delta, S) \subseteq \textstyle\bigcup_k
    \mathcal{R}(\Delta_k, S_k) \ \text{\upshape(checked: the interval
    chain)}
}{\Phi \vdash \Psi}
\qquad\rightsquigarrow\qquad
\text{\upshape select the covering branch; consult its subtree.}
\]
Conceptually the rewrite is an $N$-way case analysis: at any valuation
satisfying $\Phi$, the premise puts the evaluated point in the parent
region, the cover places it in some branch region $k$, and
$\mathcal{T}_k$'s strengthened premise then holds --- so $\mathcal{T}_k$
delivers $\Psi$. The shipped warrant packages this selection argument as
one verified lemma rather than as catalogue surgery (the catalogued
stratum records the corresponding derivation sketch through
\rulename{Farkas}-discharged infeasibility cuts and
$\vee$-elimination); the checked side condition is
Proposition~\ref{prop:sd-chain}, the exhaustiveness of the case
analysis.
\end{brecipe}

\begin{bprop}[the greedy chain covers]\label{prop:sd-chain}
Let $\mathrm{ivs}$ be a finite list of closed intervals and
$\mathrm{lo} < \mathrm{hi}$ rationals. Define one greedy pass
$r \mapsto \mathrm{extend}(r)$: fold over the list, replacing $r$ by an
interval's right endpoint whenever that interval's left endpoint is
$\leq r$ and its right endpoint is $> r$. If after
$|\mathrm{ivs}|$ passes from $\mathrm{lo}$ the reach is
$\geq \mathrm{hi}$, then
$[\mathrm{lo}, \mathrm{hi}] \subseteq \bigcup \mathrm{ivs}$.
\end{bprop}
\begin{proof}
Maintain the invariant: the reach $r$ satisfies $\mathrm{lo} \leq r$,
and either $r \leq \mathrm{lo}$ or every point of
$[\mathrm{lo}, r]$ lies in some interval of the \emph{full} list. It
holds initially ($r = \mathrm{lo}$, the first disjunct). One pass
preserves it: when the fold takes an interval $[a, b]$ with
$a \leq r < b$, any point $y \in [\mathrm{lo}, b]$ is either
$\leq r$ --- covered by the invariant, or by $[a, b]$ itself in the
$r \leq \mathrm{lo}$ base case since then $a \leq \mathrm{lo} \leq y$
--- or lies in $(r, b] \subseteq [a, b]$; and when the fold declines an
interval the state is unchanged. As many passes as there are intervals
reach the fixpoint: a productive pass strictly increases $r$ onto a
\emph{new} interval's right endpoint, of which there are at most
$|\mathrm{ivs}|$, and an unproductive pass is already the fixpoint. If
the final reach is $\geq \mathrm{hi} > \mathrm{lo}$, the invariant's
second disjunct is in force and covers all of
$[\mathrm{lo}, \mathrm{hi}]$.
\end{proof}

\begin{bprop}[semantic warrant]\label{prop:subdiv-sound}
Suppose the region premise and every branch premise-sequent are
semantically valid, the branch premises have the checked shape
($(\mathrm{branch}\ k)$'s premise is
$\Phi \wedge (x \in \mathcal{R}(\Delta_k, S_k))$ and its conclusion is
$\Psi$), and the cover check accepts. Then $\Phi \vdash \Psi$ is
semantically valid.
\end{bprop}
\begin{proof}
Fix an admissible $\sigma \models \Phi$. The region premise puts the
evaluated point $q$ of $x$ in $\mathcal{R}(\Delta, S)$; the accepted
cover check says the parent reads as a closed interval
$[\mathrm{lo}, \mathrm{hi}]$ (dimension 1, two-vertex support,
$S = \emptyset$), so $q$'s coordinate lies in
$[\mathrm{lo}, \mathrm{hi}]$ --- eliminating the parent's barycentric
weights to endpoint bounds. By Proposition~\ref{prop:sd-chain} some
branch interval contains the coordinate. Each branch interval was
itself produced by running the closed-interval reading on that
branch's \emph{own} simplex and extension set (the cover check
refuses any branch it cannot so read, and checks the tie between the
two), so ``the coordinate lies in branch $k$'s interval'' converts
back --- re-introducing barycentric weights over branch $k$'s own
support from the two endpoint bounds --- into membership of the
\emph{same} evaluated point in $\mathcal{R}(\Delta_k, S_k)$: this is
where nothing about the region term is needed beyond its one
evaluation, used twice.
Hence $\sigma$ satisfies branch $k$'s premise
$\Phi \wedge (x \in \mathcal{R}(\Delta_k, S_k))$; its validity delivers
its conclusion, which is $\Psi$ by the checked branch-conclusion
shape.
\end{proof}

\begin{bremark}[what breaks without the cover check]
\label{rem:sd-gap}
Everything: the rule without the cover is ``prove $\Psi$ wherever
convenient.'' The canonical mutant is the gap family
$\{[1, \tfrac{4}{3}],\ [\tfrac{3}{2}, 2]\}$ over the parent
$[1, 2]$ --- each branch sub-proof can be perfectly honest, and the goal
simply fails on $(\tfrac{4}{3}, \tfrac{3}{2})$, where no branch
constrains it. The greedy reach from $1$ arrives at $\tfrac{4}{3}$ and
stops: no remaining interval's left endpoint is $\leq \tfrac{4}{3}$,
so the check refuses. The checker's replayed battery keeps this mutant,
together with a point-supported branch (a degenerate interval that
cannot read), an extension-ray parent, and untied branch simplices, as
standing negative controls.
\end{bremark}

\begin{bremark}[the fragment boundary, and why refusal is the right
default]\label{rem:sd-fragment}
The checkable fragment is deliberate: 1-dimensional closed intervals,
chained. Extension-ray parents or branches, higher-dimensional regions,
and richer supports are refused --- not because the rule is false there,
but because no verified covering argument exists there yet, and the
fail-toward-refusal posture makes the unproven territory
unrepresentable. Box domains lose no expressive reach:
$n$-dimensional grids are expressible by nesting per-coordinate splits
across several applications, each within the fragment, the inner
split's parent membership entering through the outer split's branch
premise. The cover-certificate
wire slot is likewise part of the posture: recomputing the cover from
the branch regions means a forged certificate cannot influence the
verdict --- ignoring the slot cannot affect it either --- and a future
interpretation that additionally \emph{required} per-gap infeasibility
certificates from the slot would only tighten the check (one that
instead accepted new covers on certificate evidence would be a new
checking route, owing its own soundness argument).
\end{bremark}

\begin{bremark}[relation to the literature]
The rule is the reassembly half of branch-and-bound over Bernstein
enclosures: subdivision sharpens coefficient bounds quadratically in the
subdivision width and
drives them to the range in the limit \cite{garloff1986}, with the
subdivision algorithms themselves classical
\cite{decasteljau1959,laneriesenfeld1980}. In that literature the
splitting is a completeness device --- \emph{some} refinement certifies
any strict global fact; here only the soundness direction enters the
kernel: whatever refinement the producer chose, the chain check makes
its case analysis exhaustive. None of the producer-side effectivity
data (piece counts, degrees, tolerances) is consulted.
\end{bremark}

\begin{anchors}
Recipe \ref{rec:sd-instance} warrant & \texttt{sound\_subdivide}\\
interval reading of the parent and branches & \texttt{ExtensionSet.interval1D?}; elim/intro pair \texttt{interval1D\_elim}, \texttt{interval1D\_intro}\\
Proposition \ref{prop:sd-chain} & \texttt{coverExtendPass}, \texttt{coverReach}, \texttt{coverChainCheck}; invariant \texttt{ChainInv}, \texttt{chainInv\_extendPass}, \texttt{chainInv\_reach}; covering fact \texttt{coverChain\_covers}\\
the assembled check (ties included) & \texttt{coverCheck}, \texttt{coverInfeasibilityValid}\\
branch premise/conclusion shape & payload fields \texttt{branchPre}, \texttt{branchPost} (\texttt{PSubdivideData})\\
same-point re-entry (Prop.\ \ref{prop:subdiv-sound}) & the region read used by both parent elimination and branch introduction in \texttt{sound\_subdivide}'s body\\
wire placement & \texttt{RuleTag.subdivide}; the unconsulted certificate slot \texttt{coverCertificate}\\
\end{anchors}


\rulesection{19}{LiRPA}{admissible rule}
\label{blcatrule-LiRPA}
\ruleblock{The rule (Steele metanotation)}
\sourcenote{as realized in the mechanized development; source row: \texttt{\texttt{rules.toml}} --- \emph{LiRPA}. The metanotation statement is typeset by hand for notational quality; its mechanized ground truth is the Lean statement below, extracted from the sources at every build of this book and hash-bound. The rule \emph{roster} is machine-checked (every mechanized rule tag round-trips to exactly one catalogue row); the statement \emph{text} is editorial, checked by eye against the extracted Lean beside it.}
\begin{mathpar}
\inferrule*[Right=\rulename{LiRPA}]{
  \Phi \vdash a \leq h \\
  \Phi \vdash h \leq b \\
  \ReLU_{\{h\}\to\{y\}} \in H \\
  a < 0 \\
  0 < b \\\\
  \alpha a + \beta \leq 0 \\
  \beta \leq 0 \\
  \alpha b + \beta \leq b \\\\
  0 \leq \gamma a + \delta \\
  0 \leq \delta \\
  b \leq \gamma b + \delta
}{
  \Phi \vdash (\alpha h + \beta) \leq y \;\wedge\; y \leq (\gamma h + \delta)
}
\end{mathpar}

\par\medskip
\rulename{LiRPA} is the calculus\textquotesingle s neural-network
bound-propagation rule — the certificate form of the linear relaxation
of the ReLU activation that bound-propagation verifiers compute (the
CROWN/LiRPA family, the triangle relaxation, DeepPoly). Given
already-proven box bounds $\Phi \vdash a \leq h$ and $\Phi \vdash h
\leq b$ on the preactivation wire of a cited ReLU application, with
$a < 0 < b$ (the box straddles the kink — the \emph{unstable} case, in
the verification literature\textquotesingle s vocabulary), it concludes
the two-sided affine envelope $\alpha h + \beta \leq y \;\wedge\; y
\leq \gamma h + \delta$ on the application' s output
wire $y$. The producer finds the envelope however it likes —
CROWN\textquotesingle s adaptive slopes, the triangle
relaxation\textquotesingle s chord, any heuristic — and ships only the
box endpoints and the four coefficients; no relaxation machinery is
re-run or trusted.

What the checker recomputes, given the payload (two premise proofs, the
two wires, six rationals): the \emph{linkage} — the hypernet really
contains a packaged ReLU application at exactly these ports, of
positive total size (a one-line decidable port lookup); that the
premises are literally the atoms $a \leq h$ and $h \leq b$ over one
shared premise formula; the sign conditions $a < 0$ and $0 < b$; the
six envelope inequalities, exactly over $\Q$ — the lower line under the
function at the three checkpoints $a$, $0$, $b$ ($\alpha a + \beta \leq
0$, $\beta \leq 0$, $\alpha b + \beta \leq b$) and the upper line above
it at the same three ($0 \leq \gamma a + \delta$, $0 \leq \delta$, $b
\leq \gamma b + \delta$); and that the conclusion is \emph{exactly} the
pinned conjunction of the two affine comparison atoms, built by the
checker from the coefficients — the producer proposes no conclusion
shape at all.

Each hypothesis but one is load-bearing. Drop the linkage and $y$ is an
arbitrary wire with no semantic tie to $h$ (the rule as first staged
had exactly this gap and was kept vacuous — rejecting every certificate
— until the port-lookup premise landed). Drop the lower kink check
($\beta \leq 0$) and the chord through the box endpoints, which
passes both endpoint checks, would certify as a lower envelope while
lying strictly above the function at $h = 0$. (The upper kink check
$0 \leq \delta$ is the one redundancy: under the strict straddle the
two upper endpoint checks already force $\delta \geq -ab/(b-a) > 0$ —
evaluate the affine interpolation of the endpoint bounds at $0$ — so
it rides along because the certificate's check set is the uniform
six.) Drop an endpoint check
and a steep line passing the remaining two checkpoints escapes the
function at the dropped end (the proof section constructs both
failures). The box premises themselves are what make an affine
\emph{upper} envelope possible at all — no affine function dominates
$\max(0, h)$ on all of $\Q$ — and what license lower-line slopes
outside $[0, 1]$.

\cataloguenote{The family\textquotesingle s shared machinery — the packaged-node
linkage, the elementwise two-case semantic contract, and the
conservativity warrant behind it — is stated once, in this
section\textquotesingle s recipe preamble, and referenced by the four
sibling rules (\rulename{ReLU-Pos}, \rulename{ReLU-Neg},
\rulename{ReLU-NN}, \rulename{ReLU-Above}). One envelope serves the
whole node: the rule is elementwise over the cited
application\textquotesingle s tensor shape, with a single $(a, b,
\alpha, \beta, \gamma, \delta)$ for every entry — per-entry envelopes
are expressed with singleton-shaped applications. The payload
additionally declares a denominator bound on its witness rationals (a
size claim for resource accounting; it carries no proof-theoretic
weight in the soundness argument). The rule\textquotesingle s name is
the literature\textquotesingle s: linear-relaxation-based perturbation
analysis~\cite{xu2020lirpa}.
}
\ruleblock{Lean statement}
\par\medskip
\textbf{The linkage and the pinned conclusion, from the checker
sources.}\enspace \texttt{reluLinked} is the whole of the
family\textquotesingle s side condition: the hypernet\textquotesingle s
application list contains a packaged ReLU application at exactly the
cited preactivation/output ports, of positive total size — one
decidable lookup (\texttt{PrimApp.isReluAt}), no pattern-matching over
composites. The positive-size gate refuses vacuous surface at the
source: a size-zero node\textquotesingle s comparison atoms hold
vacuously, and a fact no valuation can falsify must not be citable
(the same lesson the \rulename{Farkas} rule\textquotesingle s
zero-size counterexample taught that payload).
\texttt{affineWireTerm} and \texttt{lirpaConclusionForm} pin the
conclusion: the rule concludes exactly the conjunction of the two
affine comparison atoms, built from the payload coefficients — there
is no producer-supplied conclusion shape to audit.

\declhead{BernsteinLogic.Core.Catalogue}{PrimApp.isReluAt}{\badge{BadgeGreen}{definition bound}}
\begin{leansrc}
/-- Port-lookup matcher for the ReLU linkage: a packaged ReLU
application at exactly these ports, of positive total size. The size
gate is the shape-aware rider (conservative-extension memo register 2,
the `PFarkasData.wiresPos` lesson generalized): a shape-[0] node's
atoms hold vacuously, and vacuous surface must not be citable — the
rules refuse it at the linkage, before any soundness question arises. -/
def PrimApp.isReluAt (preact post : Wire) : PrimApp → Bool
  | .relu s i o =>
      decide (i = preact) && decide (o = post) &&
      decide (0 < Shape.totalSize s)
  | _ => false
\end{leansrc}
\declhead{BernsteinLogic.Core.Catalogue}{reluLinked}{\badge{BadgeGreen}{definition bound}}
\begin{leansrc}
def reluLinked (H : Hypernet) (preact post : Wire) : Prop :=
  H.applications.any (PrimApp.isReluAt preact post) = true
\end{leansrc}
\declhead{BernsteinLogic.Core.Catalogue}{affineWireTerm}{\badge{BadgeGreen}{definition bound}}
\begin{leansrc}
def affineWireTerm (offset slope : Rat) (w : Wire) : Term :=
  .add (.const offset) (.mul (.const slope) (.var w))
\end{leansrc}
\declhead{BernsteinLogic.Core.Catalogue}{lirpaConclusionForm}{\badge{BadgeGreen}{definition bound}}
\begin{leansrc}
def lirpaConclusionForm
    (preact post : Wire) (α β γ δ : Rat) (F : Formula) : Prop :=
  F = .and (.atom .le (affineWireTerm β α preact) (.var post))
           (.atom .le (.var post) (affineWireTerm δ γ preact))
\end{leansrc}
\par\medskip
\textbf{The verified chain.}\enspace \texttt{reluLinked\_facts} turns
the linkage into semantics: under any admissible valuation the two
cited wires carry tensors of one shared shape and positive entry
count, related entrywise by the two-case ReLU relation. The four
interpolation lemmas are the checkpoint arithmetic — an affine
function suitably bounded at an interval\textquotesingle s endpoints
is bounded on the interval, one lemma per envelope side per sign
branch — and \texttt{sound\_lirpa} assembles them entrywise into the
rule\textquotesingle s warrant.

\declhead{BernsteinLogic.Soundness}{reluLinked\_facts}{\badge{BadgeGreen}{statement bound}}
\begin{leansrc}
/-- Extraction for the realized B12 linkage: under any admissible
valuation, a cited ReLU application yields both wire tensors — one
shared shape, positive entry count — and the elementwise two-case
relation. Typing rides `H.typeConsistent` at the cited application (no
payload typing fields needed — the wire format is untouched). -/
theorem reluLinked_facts {H : Hypernet} {preact post : Wire}
    (hlink : reluLinked H preact post) (σ : H.Valuation)
    (hadm : σ.Admissible) :
    ∃ qp qy : QTensor,
      σ.wireQ preact = some qp ∧ σ.wireQ post = some qy ∧
      qp.shape = qy.shape ∧ 0 < qp.entries.length ∧
      qp.entries.length = qy.entries.length ∧
      ∀ i, i < qp.entries.length →
        (qp.entries.getD i 0 ≤ 0 ∧ qy.entries.getD i 0 = 0) ∨
        (0 ≤ qp.entries.getD i 0 ∧
          qy.entries.getD i 0 = qp.entries.getD i 0)
\end{leansrc}
\declhead{BernsteinLogic.Soundness}{lirpa\_lower\_neg}{\badge{BadgeGreen}{statement bound}}
\begin{leansrc}
theorem lirpa_lower_neg {a α β p : ℚ} (ha : a < 0)
    (h1 : α * a + β ≤ 0) (h2 : β ≤ 0) (hap : a ≤ p) (hp0 : p ≤ 0) :
    α * p + β ≤ 0
\end{leansrc}
\declhead{BernsteinLogic.Soundness}{lirpa\_lower\_pos}{\badge{BadgeGreen}{statement bound}}
\begin{leansrc}
theorem lirpa_lower_pos {b α β p : ℚ} (hb : 0 < b)
    (h2 : β ≤ 0) (h3 : α * b + β ≤ b) (hp0 : 0 ≤ p) (hpb : p ≤ b) :
    α * p + β ≤ p
\end{leansrc}
\declhead{BernsteinLogic.Soundness}{lirpa\_upper\_neg}{\badge{BadgeGreen}{statement bound}}
\begin{leansrc}
theorem lirpa_upper_neg {a γ δ p : ℚ} (ha : a < 0)
    (h1 : 0 ≤ γ * a + δ) (h2 : 0 ≤ δ) (hap : a ≤ p) (hp0 : p ≤ 0) :
    0 ≤ γ * p + δ
\end{leansrc}
\declhead{BernsteinLogic.Soundness}{lirpa\_upper\_pos}{\badge{BadgeGreen}{statement bound}}
\begin{leansrc}
theorem lirpa_upper_pos {b γ δ p : ℚ} (hb : 0 < b)
    (h2 : 0 ≤ δ) (h3 : b ≤ γ * b + δ) (hp0 : 0 ≤ p) (hpb : p ≤ b) :
    p ≤ γ * p + δ
\end{leansrc}
\declhead{BernsteinLogic.Soundness}{sound\_lirpa}{\badge{BadgeGreen}{statement bound}}
\begin{leansrc}
/-- LiRPA: the packaged degree-1 relaxation — from box premises
`a ≤ h ≤ b` and the eight endpoint side conditions, the affine sandwich
`αh + β ≤ y ∧ y ≤ γh + δ` holds entrywise, by the four interpolation
lemmas applied per two-case branch. -/
theorem sound_lirpa {H : Hypernet}
    {seq lowerPremise upperPremise : Sequent}
    {preact post : Wire} {a b α β γ δ : ℚ}
    (ihl : lowerPremise.semanticallyValid H)
    (ihu : upperPremise.semanticallyValid H)
    (hlink : reluLinked H preact post)
    (hAneg : a < 0) (hBpos : 0 < b)
    (hLower1 : α * a + β ≤ 0) (hLower2 : β ≤ 0) (hLower3 : α * b + β ≤ b)
    (hUpper1 : 0 ≤ γ * a + δ) (hUpper2 : 0 ≤ δ) (hUpper3 : b ≤ γ * b + δ)
    (hpre : seq.pre = lowerPremise.pre)
    (hpreMatch : lowerPremise.pre = upperPremise.pre)
    (hlowerPost : lowerPremise.post = .atom .le (.const a) (.var preact))
    (hupperPost : upperPremise.post = .atom .le (.var preact) (.const b))
    (hpostForm : lirpaConclusionForm preact post α β γ δ seq.post) :
    seq.semanticallyValid H
\end{leansrc}
\par\medskip
\textbf{Why consuming the packaged two-case fact is not new kernel
content.}\enspace The two-case relation is the packaged
symbol\textquotesingle s defined semantics, consumed by the family as
model fact at every positive size; what this pair of declarations adds
is the conservativity warrant --- that the definition mints no
genuinely new primitive: \texttt{reluExpH} is the definitional
expansion (a copy, a sign-split cell complex, an indicator cast, and
an exact rank-0 contraction), and \texttt{relu\_projection} shows —
both directions — that the two-case relation is exactly the
expansion\textquotesingle s model-level facts existentially projected
onto the external ports, proven at the canonical single-entry
instance. The boundary point $h = 0$ rides the one-hot cell semantics:
both cells admit it, and every admissible choice forces $y = 0$.
Applications of total size other than one currently have no core
expansion (a recorded packing gap in the development\textquotesingle s
notes): at those sizes the node stands as a genuine primitive with the
stated semantics — the rules remain sound there, while the
not-a-new-primitive claim is certified only at size one.

\declhead{BernsteinLogic.ReluExpansion}{relu\_projection}{\badge{BadgeGreen}{statement bound}}
\begin{leansrc}
/-- **Obligation P (the projection lemma), both directions**: the
packaged ReLU's derived sp-fact — the two-case relation — EQUALS the
expansion's sp-facts existentially projected onto the external ports.
The internal wires (the copies, the indicators, the cast image, the
rank-0 legs) are exactly the existentially-quantified content of the
valuation; the boundary case p = 0 is covered because both one-hot
choices are admissible there and both force y = 0. -/
theorem relu_projection (p y : ℚ) :
    (∃ σ : reluExpH.Valuation, σ.Admissible ∧
        σ.entryFlat 0 0 = some (.fin p) ∧
        σ.entryFlat 8 0 = some (.fin y)) ↔
    ((p ≤ 0 ∧ y = 0) ∨ (0 ≤ p ∧ y = p))
\end{leansrc}
\declhead{BernsteinLogic.ReluExpansion}{reluExpH}{\badge{BadgeGreen}{definition bound}}
\begin{leansrc}
/-- The expansion: copy;(id ⊗ Cell[Σ±]);Contract, with the typing
plumbing (Cast 𝔹→ℚ on the indicator, reshapes on the rank-1 ports). -/
def reluExpH : Hypernet where
  wireType := expTy
  applications :=
    [.copy 2 0 [1, 2],
     .cell sigmaPM 2 [3, 4],
     .cast Dtype.B Dtype.Q .exact [] 4 5,
     .reshape [] 1 6,
     .contract Dtype.Q .exact [6, 5] 7 [0, 0] 0 pat0,
     .reshape [1] 7 8]
  linearity := by decide
  acyclic := by decide
  typeConsistent := by decide +kernel
\end{leansrc}
\ruleblock{Correspondence (Steele display $\leftrightarrow$ Lean)}
\badge{BadgeGreen}{reviewer-certified against the displayed texts}
\sourcenote{The dictionary below is a reviewer-certified statement of how the mechanized statements above interpret the displayed rule: display metavariable $\mapsto$ Lean binder, each display hypothesis $\mapsto$ the premise or side condition carrying it, the display conclusion $\mapsto$ the goal --- with any deliberate difference in width between display and mechanization named outright. Certification is recorded as a hash triple over (the displayed cell, the extracted statement texts, this dictionary\textquotesingle s own text); if any of the three changes, this block turns amber until a reviewer re-certifies. The correspondence claim itself remains editorial --- reviewer-grade, not machine-checked; its machine-checked successor types each soundness statement against the rule as a structured object, and supersedes this block per rule as it lands.}
Judgment frame as in \rulename{Ax}; sub-derivation premises as at
\rulename{Weaken}. This dictionary doubles as the family preamble: its
linkage row is shared verbatim by the four sibling rules.
\begin{itemize}
\item $h \mapsto$ \texttt{preact : Wire}; $y \mapsto$ \texttt{post :
  Wire}; $a, b, \alpha, \beta, \gamma, \delta \mapsto$ the six
  rational binders of the same names in \texttt{sound\_lirpa}.
\item the linkage premise $\ReLU_{\{h\}\to\{y\}} \in H$ $\mapsto$
  \texttt{hlink} $:$ \texttt{reluLinked H preact post} ---
  membership \emph{plus} the positive-size gate: unfolding (the
  definition above), the certified premise demands a packaged ReLU
  application at exactly these ports \emph{of positive total size}.
  The display\textquotesingle s bare $\in$ shows no size condition;
  the certified linkage is strictly narrower, by design (vacuous
  surface refused at the source --- the intro above), and the
  narrowing is family-wide.
\item box premises $\Phi \vdash a \leq h$ and $\Phi \vdash h \leq b$
  $\mapsto$ \texttt{ihl}, \texttt{ihu} with posts pinned by
  \texttt{hlowerPost} $:$ \texttt{lowerPremise.post = .atom .le
  (.const a) (.var preact)} and \texttt{hupperPost} $:$
  \texttt{upperPremise.post = .atom .le (.var preact) (.const b)} ---
  displayed comparators and operand orders carried exactly; $\Phi$
  shared three ways via \texttt{hpre} and \texttt{hpreMatch}.
\item the eight arithmetic side conditions $\mapsto$ \texttt{hAneg},
  \texttt{hBpos}, \texttt{hLower1}--\texttt{hLower3},
  \texttt{hUpper1}--\texttt{hUpper3}, verbatim over $\Q$ --- scalar
  hypotheses on the payload rationals, not judgments.
\item conclusion $\Phi \vdash (\alpha h + \beta) \leq y \wedge y \leq
  (\gamma h + \delta)$ $\mapsto$ the goal under \texttt{hpostForm}
  $:$ \texttt{lirpaConclusionForm preact post} applied to the four
  coefficients and \texttt{seq.post}, pinning \texttt{seq.post} to
  the conjunction of the two affine atoms in the displayed order
  (lower envelope left). Spelling drift, disclosed: the pinned
  affine term is \texttt{affineWireTerm}\textquotesingle s
  offset-first $\beta + \alpha h$, where the display writes $\alpha h
  + \beta$ --- semantically identical, but formulas are ordered
  trees, so the certified rule accepts exactly the offset-first
  spelling (the \rulename{ConvElim}/\rulename{ConvIntro}
  build-output class).
\end{itemize}
Width notes beyond the linkage gate: the certified atoms are
elementwise over the cited application\textquotesingle s tensor shape
--- one envelope for every entry, per the note above --- where the
display reads scalar; that is the batch semantics of the atom
grammar, not a strengthening. \texttt{reluLinked\_facts}, the four
interpolation lemmas, \texttt{relu\_projection}, and
\texttt{reluExpH} are certified alongside as machinery: the first
extracts the elementwise two-case relation the family consumes, the
last two ground the packaged symbol in the core expansion --- they
are \emph{about} the linkage, in the satellite sense of
\rulename{OrE}\textquotesingle s ctx-rule. No hypothesis of
\texttt{sound\_lirpa} is unused, and nothing certified is wider than
the display.

\ruleblock{Lean verifier receipts}
\sourcenote{module \texttt{BernsteinLogic.Core.Catalogue}; bound build \texttt{cbc92e976a1d\ldots} (bound from receipt 5abdb3589707b5df, 2026-08-03T21:30:33Z). A receipt records that named checking layers accepted a build: (module, compiled-artifact SHA-256, environment SHA-256, checker layers), minted when the checking gate ran; the store receipts attesting the bound build are listed with the layers each run credited. Nothing here asks the reader to trust the store: every layer is an independently runnable program over the quoted sources --- rebuild the module and re-run the checkers to reproduce the verdict. Receipts attesting superseded builds credit no layer and are omitted from this rendering (the count of omissions, if any, is noted below).}
\begin{receipttable}
\receiptrow{24804e66edba0bc9}{2026-08-03T20:25:09Z}{build+decorrelated}{\badge{BadgeGreen}{current}}
\receiptrow{5abdb3589707b5df}{2026-08-03T21:30:33Z}{build+decorrelated}{\badge{BadgeGreen}{current}}
\end{receipttable}
\sourcenote{18 further receipt(s) in the store attest superseded builds of this module and are omitted from this rendering.}
\sourcenote{module \texttt{BernsteinLogic.Soundness}; bound build \texttt{a20d90e66c13\ldots} (bound from receipt a9c0e34d3e3c9c8c, 2026-08-03T21:30:36Z). A receipt records that named checking layers accepted a build: (module, compiled-artifact SHA-256, environment SHA-256, checker layers), minted when the checking gate ran; the store receipts attesting the bound build are listed with the layers each run credited. Nothing here asks the reader to trust the store: every layer is an independently runnable program over the quoted sources --- rebuild the module and re-run the checkers to reproduce the verdict. Receipts attesting superseded builds credit no layer and are omitted from this rendering (the count of omissions, if any, is noted below).}
\begin{receipttable}
\receiptrow{a9c0e34d3e3c9c8c}{2026-08-03T21:30:36Z}{build+decorrelated}{\badge{BadgeGreen}{current}}
\end{receipttable}
\sourcenote{21 further receipt(s) in the store attest superseded builds of this module and are omitted from this rendering.}
\sourcenote{module \texttt{BernsteinLogic.ReluExpansion}; bound build \texttt{68395e3f44af\ldots} (bound from receipt b1a5795a8e57ce30, 2026-07-18T00:07:06Z). A receipt records that named checking layers accepted a build: (module, compiled-artifact SHA-256, environment SHA-256, checker layers), minted when the checking gate ran; the store receipts attesting the bound build are listed with the layers each run credited. Nothing here asks the reader to trust the store: every layer is an independently runnable program over the quoted sources --- rebuild the module and re-run the checkers to reproduce the verdict. Receipts attesting superseded builds credit no layer and are omitted from this rendering (the count of omissions, if any, is noted below).}
\begin{receipttable}
\receiptrow{7d0d53fd422f509e}{2026-07-17T20:41:28Z}{build+replay+hardened+decorrelated}{\badge{BadgeGreen}{current}}
\receiptrow{3219a0974f9b52a3}{2026-07-17T21:00:42Z}{build+replay+hardened+decorrelated}{\badge{BadgeGreen}{current}}
\receiptrow{b1a5795a8e57ce30}{2026-07-18T00:07:06Z}{build+decorrelated}{\badge{BadgeGreen}{current}}
\end{receipttable}
\sourcenote{4 further receipt(s) in the store attest superseded builds of this module and are omitted from this rendering.}
\layerline{\layeritem{build}{lake build (elaboration)}{\badge{BadgeGreen}{current}}\layeritem{replay}{leanchecker (official kernel replay)}{\badge{BadgeGreen}{current}}\layeritem{hardened}{safe\_verify (hardened screen)}{\badge{BadgeGreen}{current}}\layeritem{decorrelated}{lean4export + nanoda (decorrelated kernel)}{\badge{BadgeGreen}{current}}}
\sourcenote{\badge{BadgeGray}{comparator: none} no fixture in the three-implementation comparator harness (compiled / kernel-interpreted / reference-interpreted) exercises this rule.}
\ruleblock{HOL4 statement}
\sourcenote{theory \texttt{bern\_tr\_catdata2}, HOL4 development script \texttt{\texttt{bern\_tr\_catdata2Script.sml}} \badge{BadgeGreen}{script bound}}
\holhead{Theorem plirpadata\_receipt}
\begin{holsrc}
Theorem plirpadata_receipt:
  ruledataOk [] cd2Hrelu
    <| pre := FAnd (FAtom BowLe (TmConst (-&1)) (TmVar 3))
                   (FAtom BowLe (TmVar 3) (TmConst (&1))) ;
       post := FAnd (FAtom BowLe (affineWireTerm 0 0 3) (TmVar 4))
                    (FAtom BowLe (TmVar 4) (affineWireTerm 1 0 3)) |>
    (RdLirpa (PLirpaD
       <| pre := FAnd (FAtom BowLe (TmConst (-&1)) (TmVar 3))
                      (FAtom BowLe (TmVar 3) (TmConst (&1))) ;
          post := FAtom BowLe (TmConst (-&1)) (TmVar 3) |>
       (ProofObj (RdAndE (PAndEDL
          <| pre := FAnd (FAtom BowLe (TmConst (-&1)) (TmVar 3))
                         (FAtom BowLe (TmVar 3) (TmConst (&1))) ;
             post := FAnd (FAtom BowLe (TmConst (-&1)) (TmVar 3))
                          (FAtom BowLe (TmVar 3) (TmConst (&1))) |>
          (ProofObj RdAx)
          (FAtom BowLe (TmVar 3) (TmConst (&1))))))
       <| pre := FAnd (FAtom BowLe (TmConst (-&1)) (TmVar 3))
                      (FAtom BowLe (TmVar 3) (TmConst (&1))) ;
          post := FAtom BowLe (TmVar 3) (TmConst (&1)) |>
       (ProofObj (RdAndE (PAndEDR
          <| pre := FAnd (FAtom BowLe (TmConst (-&1)) (TmVar 3))
                         (FAtom BowLe (TmVar 3) (TmConst (&1))) ;
             post := FAnd (FAtom BowLe (TmConst (-&1)) (TmVar 3))
                          (FAtom BowLe (TmVar 3) (TmConst (&1))) |>
          (ProofObj RdAx)
          (FAtom BowLe (TmConst (-&1)) (TmVar 3)))))
       3 4 (-&1) (&1) 0 0 0 (&1) 0)) ∧
  ¬ ruledataOk [] hnEmpty
      <| pre := FAnd (FAtom BowLe (TmConst (-&1)) (TmVar 3))
                     (FAtom BowLe (TmVar 3) (TmConst (&1))) ;
         post := FAnd (FAtom BowLe (affineWireTerm 0 0 3) (TmVar 4))
                      (FAtom BowLe (TmVar 4) (affineWireTerm 1 0 3)) |>
      (RdLirpa (PLirpaD
         <| pre := FTt ; post := FTt |> (ProofObj RdAx)
         <| pre := FTt ; post := FTt |> (ProofObj RdAx)
         3 4 (-&1) (&1) 0 0 0 (&1) 0))
\end{holsrc}
\ruleblock{HOL verifier receipts}
\badge{BadgeBlue}{documentary} \texttt{hol-strict} CLEAN: theory \texttt{bern\_tr\_catdata2}, 272
theorems, 0 cheat-tagged, 0 nonstandard axioms (whitelist:
\texttt{DISK\_THM}; the four \texttt{bool} axioms) --- verification
run of 2026-07-22 on an isolated build host, against a
content-addressed checkout of the exact sources cited here (git blob
\texttt{cd3c1e95\ldots}). The proof-object stratum is ported as one
mutual datatype family with a mutual wellformedness predicate,
exactly as the Lean mutual inductive block; this
rule\textquotesingle s arm carries the two premise sub-proofs, the
six sandwich inequalities verbatim, the application-list linkage
(\texttt{reluLinkedB}, theory \texttt{bern\_tr\_catdata}, 252
theorems, 0 offences, same run --- the lookup the earlier port
excluded wholesale now has a ported realization), and the conclusion
form (\texttt{lirpaConclusionFormB}, theory
\texttt{bern\_tr\_catforms}, 23 theorems, 0 offences). The receipt
theorem cited here is the kernel-evaluated accept/reject pair over a
hypernet carrying the packaged ReLU application. The four bound
lemmas behind the sandwich soundness are ported as well
(\texttt{lirpa\_lower\_neg} / \texttt{lirpa\_lower\_pos} /
\texttt{lirpa\_upper\_neg} / \texttt{lirpa\_upper\_pos}, theory
\texttt{bern\_tr\_soundmid}, 16 theorems, 0 offences, run of
2026-07-22). Honest residue: the apex semantic theorem assembling
them (the Lean development\textquotesingle s \texttt{sound\_lirpa})
is not yet ported; its ingredients above are.

\sourcenote{HOL receipts are documentary: quoted results of \texttt{hol-strict} runs
and census rows from the HOL4 development\textquotesingle s build
evidence, not minted receipt artefacts. A HOL receipt store (the Lean
store\textquotesingle s twin) is named future work.
}
\ruleblock{Proof}

\sourcenote{Genre: admissible rule (tier 2) --- a proof-tree-rewriting
recipe. This section additionally carries the \emph{family template}
for the five ReLU rules (\rulename{LiRPA}, \rulename{ReLU-Pos},
\rulename{ReLU-Neg}, \rulename{ReLU-NN}, \rulename{ReLU-Above}): the
packaged-node machinery they share is stated once here, and the four
sibling sections instantiate it by reference.}

\begin{brecipe}[the ReLU family template]\label{rec:relu-template}
The general admissible template --- one rule node rewrites to a
derivation fragment over core and previously-justified rules with the
same conclusion --- is stated in full in the \rulename{PConstExact}
section; the five ReLU rules are its instances over one shared
semantic fact. The realized term grammar has no $\ReLU$ symbol:
$\ReLU$ is a \emph{packaged hypernet primitive}, an application
$\ReLU_{\{h\}\to\{y\}} \in H$ with preactivation (input) wire $h$ and
output wire $y$. Every family rule cites it through one side
condition, the \emph{linkage} $\mathrm{reluLinked}(H, h, y)$: the
hypernet's application list contains a packaged ReLU at exactly these
ports, of positive total size --- a single decidable port lookup, not a
pattern-match over composites. The packaged node's semantic contract
is the elementwise two-case fact: under every admissible valuation,
at every entry $i$ of the node's (shared) tensor shape,
\[
\bigl(h_i \leq 0 \;\wedge\; y_i = 0\bigr)
\;\vee\;
\bigl(0 \leq h_i \;\wedge\; y_i = h_i\bigr).
\]
Three remarks fix what the template means and what it does not.

\emph{(i) The rewrite shape.} Each family node rewrites to a case
analysis over this disjunction: the sign-split fact enters the branch
contexts (context-carrying $\vee$E --- the disjunction arrives as a
sub-conjunct of the premise; see the $\vee$E section's note), equality
transport (\rulename{EqElim}) carries the branch equality $y = 0$ or
$y = h$ through the goal, and \rulename{ArithAtom}/\rulename{Norm}
plumbing normalizes atom orientations. Where a branch discharges fully
into catalogue nodes the instance recipe displays the derivation
(\rulename{ReLU-NN} and \rulename{ReLU-Above}); where it does not ---
the boundary branches of \rulename{ReLU-Pos}/\rulename{ReLU-Neg} need
an antisymmetry step no catalogue rule provides --- the instance says
so and the rule is admissible by \emph{reduction to verified
machinery}, the honest tier note the \rulename{PConstExact} section
established.

\emph{(ii) The semantic contract and the conservativity warrant, kept
distinct.} The two-case relation is the packaged symbol's
\emph{defined} semantics --- part of what an admissible valuation of a
hypernet containing the node means, at every entry of every
positive-size application --- and the family's soundness lemmas consume
it as model fact at any size (that is (iii)'s extraction lemma). What
the conservativity warrant adds is that the definition mints no
genuinely new primitive: a definitional expansion into core primitives
(a copy, the sign-split cell complex $\Sigma^{\pm}$, an indicator
cast, and an exact rank-0 contraction --- \texttt{reluExpH}), a typing
rule derivable from the expansion's, and the projection theorem
\texttt{relu\_projection}: \emph{both directions}, the two-case
relation is exactly the expansion's model-level facts existentially
projected onto the external ports --- proven at the canonical
single-entry instance. The boundary point $h = 0$ rides the one-hot
cell semantics: both cells of $\Sigma^{\pm}$ admit it, and every
admissible choice forces $y = 0$. Scope, stated loudly: applications
of total size other than one currently have \emph{no} core expansion
(a recorded packing gap), so at those sizes the packaged node stands
as a genuine primitive with the stated semantics --- the five rules
remain sound there, consuming the defined relation, while the
not-a-new-primitive claim is certified only at size one.

\emph{(iii) The shipped verified objects.} Per the catalogue's
convention, each family rule ships a direct soundness lemma factoring
through already-verified machinery rather than a derivation object:
\texttt{reluLinked\_facts} converts the linkage, under any admissible
valuation, into (a) tensors at both ports, (b) one shared shape of
positive entry count, and (c) the per-entry two-case relation; each
rule's lemma is then an entrywise case analysis. The positive-size
gate inside the linkage is deliberate hygiene: a size-zero node's
comparison atoms hold vacuously, and vacuous surface must not be
citable --- the general lesson of the \rulename{Farkas} rule's
zero-size counterexample, enforced here at the source.
\end{brecipe}

\begin{brecipe}[the \rulename{LiRPA} instance --- subdivide at the kink]
\label{rec:lirpa-instance}
The node, with its premise subtrees and side conditions:
\[
\inferrule*[Right=\rulename{LiRPA}]{
  \mathcal{T}_a \rhd \Phi \vdash a \leq h \\
  \mathcal{T}_b \rhd \Phi \vdash h \leq b \\
  \ReLU_{\{h\}\to\{y\}} \in H \\
  a < 0 \\ 0 < b \\\\
  \alpha a + \beta \leq 0 \\ \beta \leq 0 \\ \alpha b + \beta \leq b \\
  0 \leq \gamma a + \delta \\ 0 \leq \delta \\ b \leq \gamma b + \delta
}{
  \Phi \vdash (\alpha h + \beta) \leq y \;\wedge\; y \leq (\gamma h + \delta)
}
\]
rewrites to a two-branch case analysis at the kink,
\[
\inferrule*[Right=cases]{
  \mathcal{D}_{-} \rhd \Phi \wedge (h \leq 0) \vdash
    (\alpha h + \beta) \leq y \,\wedge\, y \leq (\gamma h + \delta) \\
  \mathcal{D}_{+} \rhd \Phi \wedge (0 \leq h) \vdash
    (\alpha h + \beta) \leq y \,\wedge\, y \leq (\gamma h + \delta)
}{
  \Phi \vdash (\alpha h + \beta) \leq y \;\wedge\; y \leq (\gamma h + \delta)
}
\]
where the case split is the box subdivision $[a, b] = [a, 0] \cup
[0, b]$ (region atoms enter from the order atoms via
\rulename{ConvIntro}; the cover is the one \rulename{Subdivide}
checks), and the branches discharge as follows. In $\mathcal{D}_{-}$,
the branch context proves $h \leq 0$, so \rulename{ReLU-Neg} yields
$y = 0$ (straight from the packaged contract: on $h \leq 0$ both
disjuncts force $y = 0$); the two envelope atoms then reduce, through \rulename{EqElim}
along $y = 0$, to degree-$1$ bound facts on $[a, 0]$ --- $\alpha h +
\beta \leq 0$ and $0 \leq \gamma h + \delta$ --- whose degree-$1$
Bernstein certificates on that cell are \emph{exactly} the checkpoint
conditions at $a$ and at $0$ (\rulename{BernCone} at $d = 1$; on a
$1$-simplex the degree-$1$ Bernstein coefficients of an affine gap
polynomial are its two endpoint values). Dually in $\mathcal{D}_{+}$:
\rulename{ReLU-Pos} yields $y = h$ (from the contract's active
disjunct; at the boundary $h = 0$ the inactive disjunct gives the same
$y = 0 = h$), and the envelope atoms reduce to
degree-$1$ facts on $[0, b]$ --- $\alpha h + \beta \leq h$ and $h \leq
\gamma h + \delta$ --- whose Bernstein certificates are the checkpoint
conditions at $0$ and at $b$. The six side conditions are thus the six
distinct endpoint values of the four affine gap polynomials, the kink
checkpoints $\beta \leq 0$ and $0 \leq \delta$ each shared between
their line's two cells. \rulename{ArithAtom}, \rulename{Norm}, and
\rulename{Cut} supply the atom-normalization and stitching plumbing
throughout (the \rulename{ReLU-Above} section walks a chain of this
kind in full). Scope of the rewrite, matching template remark (ii):
this two-branch derivation is the canonical \emph{single-entry} story
--- at larger tensor shapes the sign split is per-entry, so the
displayed branch contexts $h \le 0$ and $0 \le h$ do not jointly cover
a mixed-sign tensor and the derivation does not lift; the rule's
warrant at every positive size is the direct semantic lemma
(Proposition~\ref{prop:lirpa-sound}), per the family convention of
remark (iii).
\end{brecipe}

\begin{bprop}[semantic warrant]\label{prop:lirpa-sound}
Suppose the linkage holds, the two premise sequents are semantically
valid with a common premise formula, and the eight arithmetic side
conditions hold (the strict straddle pair $a < 0$, $0 < b$, plus the
six envelope checkpoints --- ``the six side conditions'' elsewhere in
this section always means the envelope six). Then the conclusion sequent --- the pinned conjunction
$(\alpha h + \beta) \leq y \wedge y \leq (\gamma h + \delta)$ over the
same premise --- is semantically valid.
\end{bprop}
\begin{proof}
Fix an admissible valuation $\sigma$ satisfying $\Phi$. The premises,
instantiated at $\sigma$, eliminate to entrywise bounds: every entry
$p$ of the preactivation tensor satisfies $a \leq p \leq b$ (a
rank-$0$ constant against a tensor compares against every entry). The
linkage extraction (\texttt{reluLinked\_facts}) gives the output
tensor at the same shape and, per entry, either $p \leq 0 \wedge y =
0$ or $0 \leq p \wedge y = p$. Fix an entry and take the cases.

\emph{Inactive case} ($p \leq 0$, $y = 0$; here $p \in [a, 0]$).
Lower line: $\alpha p + \beta \leq 0 = y$ by two-endpoint
interpolation on $[a, 0]$ --- the affine function $\alpha({\cdot}) +
\beta$ is $\leq 0$ at $a$ (checkpoint $\alpha a + \beta \leq 0$) and
at $0$ (checkpoint $\beta \leq 0$), hence on all of $[a, 0]$
(\texttt{lirpa\_lower\_neg}; the verified proof is an explicit product
certificate riding the identity
$(-a)(\alpha p + \beta) + (-p)\bigl(-(\alpha a + \beta)\bigr)
+ (p - a)(-\beta) = 0$: on $[a, 0]$ the last two products are
nonnegative, so the assumed violation $\alpha p + \beta > 0$ --- which,
with $-a > 0$, makes the first term strictly positive --- is impossible
in a zero sum). Upper
line: $y = 0 \leq \gamma p + \delta$ likewise from the checkpoints at
$a$ and $0$ (\texttt{lirpa\_upper\_neg}).

\emph{Active case} ($0 \leq p$, $y = p$; here $p \in [0, b]$). Lower
line: $\alpha p + \beta \leq p = y$ by interpolation on $[0, b]$ from
the checkpoints $\beta \leq 0$ and $\alpha b + \beta \leq b$
(\texttt{lirpa\_lower\_pos}). Upper line: $y = p \leq \gamma p +
\delta$ from $0 \leq \delta$ and $b \leq \gamma b + \delta$
(\texttt{lirpa\_upper\_pos}).

The conclusion's two atoms evaluate entrywise at the shared shape (the
affine terms evaluate to $\beta + \alpha\,p$ and $\delta + \gamma\,p$
per entry --- \texttt{affineWireTerm\_evalQ}), so the four cases
assemble into the conjunction. That is \texttt{sound\_lirpa}, whose
structure the paragraph mirrors step for step.
\end{proof}

\begin{bremark}[three checkpoints suffice]\label{rem:lirpa-checkpoints}
Because $\ReLU$ is affine on each side of $0$ and $a < 0 < b$, an
affine line bounds it on $[a, b]$ \emph{iff} it bounds it at the three
points $a$, $0$, $b$ --- the six side conditions are precisely these
six facts (two lines $\times$ three checkpoints), and equivalently the
degree-$1$ Bernstein certificates of Recipe~\ref{rec:lirpa-instance}.
A producer can therefore emit the tightest certifiable envelope for
its chosen slopes by solving three scalar inequalities per line ---
nothing about the rule rewards loose coefficients.
\end{bremark}

\begin{bremark}[what the strict straddle buys]\label{rem:lirpa-strict}
With $a < 0 < b$ the checkpoint conditions read, verbatim, ``the line
is on the correct side of $\ReLU$ at $a$, $0$, and $b$'' ($\ReLU(a) =
0$ needs $a \leq 0$; $\ReLU(b) = b$ needs $0 \leq b$). Semantic
soundness would in fact survive relaxing the strictness: whatever the
signs of $a \leq b$, the six conditions pin the lower line beneath
$\ReLU$ on $[\min(a,0), \max(b,0)] \supseteq [a,b]$ --- at $a$ because
$\alpha a + \beta \leq 0 \leq \ReLU(a)$ holds unconditionally, and on
each of $[\min(a,0), 0]$, $[0, \max(b,0)]$ by the same endpoint
interpolation; for the upper line the interval containment does the
work (the checks at $0$ and $b$ dominate $\ReLU$ on $[0, \max(b,0)]$,
the checks at $a$ and $0$ dominate it on $[\min(a,0), 0]$). The
rule as adopted nevertheless demands the strict straddle, and reasonably:
it pins \rulename{LiRPA} to the genuinely unstable case --- a stable
neuron rides \rulename{ReLU-Pos} or \rulename{ReLU-Neg}
\emph{exactly}, losing nothing to a relaxation --- it keeps both
interpolation cells nondegenerate, and it is consumed as stated by the
verified interpolation lemmas, whose contradiction products multiply
by $-a$ and $b$.
\end{bremark}

\begin{bremark}[the conclusion is built, not audited]
The conclusion formula is pinned syntactically
(\texttt{lirpaConclusionForm}): it \emph{is} the conjunction of the
two comparison atoms over the affine terms
$\texttt{affineWireTerm}\,\beta\,\alpha\,h$ and
$\texttt{affineWireTerm}\,\delta\,\gamma\,h$, constructed from the
payload coefficients --- the producer supplies no conclusion shape.
A payload field that the checker rebuilds rather than reads is one
the payload cannot mis-declare (the \rulename{PConstExact} section's
closedness observation, in conclusion-shape form): no ``equivalent''
respelling can smuggle in a different wire, comparator, or shape.
\end{bremark}

\begin{anchors}
family linkage extraction & \texttt{reluLinked\_facts}\\
checkpoint interpolation, lower line & \texttt{lirpa\_lower\_neg}, \texttt{lirpa\_lower\_pos}\\
checkpoint interpolation, upper line & \texttt{lirpa\_upper\_neg}, \texttt{lirpa\_upper\_pos}\\
affine conclusion evaluation & \texttt{affineWireTerm\_evalQ}\\
rule warrant (Proposition \ref{prop:lirpa-sound}) & \texttt{sound\_lirpa}\\
linkage side condition & \texttt{reluLinked}; \texttt{PrimApp.isReluAt} (the positive-size gate)\\
pinned conclusion & \texttt{lirpaConclusionForm}; \texttt{affineWireTerm}\\
packaged-node conservativity & \texttt{relu\_projection} over \texttt{reluExpH} (canonical single-entry instance)\\
\end{anchors}

\sourcenote{Literature: the two-sided affine relaxation of ReLU is the
verification community's workhorse --- the triangle
relaxation~\cite{ehlers2017}, its dual/convex-outer
form~\cite{wongkolter2018}, CROWN's adaptive slope
choice~\cite{zhang2018crown}, DeepPoly's abstract-domain
rendering~\cite{singh2019deeppoly}, and the LiRPA framework that named
the family~\cite{xu2020lirpa}. This rule is the checker-side
certificate form: the cited works \emph{find} envelopes; the rule only
ever \emph{checks} one, at six rational inequalities.}


\rulesection{20}{ReLU-Pos}{admissible rule}
\label{blcatrule-ReLUPos}
\ruleblock{The rule (Steele metanotation)}
\sourcenote{as realized in the mechanized development; source row: \texttt{\texttt{rules.toml}} --- \emph{ReLU-Pos}. The metanotation statement is typeset by hand for notational quality; its mechanized ground truth is the Lean statement below, extracted from the sources at every build of this book and hash-bound. The rule \emph{roster} is machine-checked (every mechanized rule tag round-trips to exactly one catalogue row); the statement \emph{text} is editorial, checked by eye against the extracted Lean beside it.}
\begin{mathpar}
\inferrule*[Right=\rulename{ReLU-Pos}]{
  \Phi \vdash h \geq 0
  \\
  \ReLU_{\{h\}\to\{y\}} \in H
}{
  \Phi \vdash y = h
}
\end{mathpar}

\par\medskip
\rulename{ReLU-Pos} is the exactness rule for the active side of a
ReLU: once $\Phi \vdash h \geq 0$ is proven for the preactivation wire
of a cited ReLU application, the output wire is not merely bounded but
\emph{equal} — $\Phi \vdash y = h$. A producer reaches for it when
earlier bound propagation has shown a neuron stably active (a proven
lower box bound at or above zero): the lossy affine envelope of
\rulename{LiRPA} is replaced by an identity, and the equality feeds
\rulename{EqElim}, so every fact known about $h$ transports to $y$
verbatim.

What the checker recomputes, given the payload (premise proof, the two
wires): the linkage — the hypernet really contains a packaged ReLU
application at exactly these ports, of positive total size (the
family\textquotesingle s port-lookup side condition, displayed in the
\rulename{LiRPA} section); that the premise\textquotesingle s
conclusion is literally the atom $h \geq 0$ over the same premise
formula; and that the conclusion is literally $y = h$. Both atoms are
elementwise over the node\textquotesingle s tensor shape.

Drop the premise and the rule is false at any entry with $h < 0$
(there $y = 0 \neq h$). Drop the linkage and $y$ is an arbitrary wire
with no semantic tie to $h$ — the linkage-free rule is unsound over
unlinked wires, and a realization without the port lookup could only
be kept sound by rejecting every certificate.
The interesting hypothesis is the one that is \emph{not}
there: no side condition excludes the boundary $h = 0$, because the
node\textquotesingle s two semantic cases agree at the kink — the
proof section makes this the centre of the argument.

\cataloguenote{TIER NOTE: the unpacking by case analysis over the sign-split fact
discharges the active branch cleanly, but its boundary branch needs
the antisymmetry step $h \leq 0 \wedge h \geq 0 \Rightarrow h = 0$,
and no catalogue rule concludes an equality from a pair of opposite
inequalities — so, like \rulename{PConstExact}, this rule is
admissible by \emph{reduction to verified machinery} rather than by a
complete catalogue derivation: the shipped warrant is the direct lemma
\texttt{sound\_reluPos}, whose boundary case is one antisymmetry step
in the model. The proof section states the enumeration behind that
claim.
}
\ruleblock{Lean statement}
\par\medskip
The warrant consumes the family linkage extraction
(\texttt{reluLinked\_facts}, displayed in the \rulename{LiRPA}
section) and the premise, eliminated entrywise (the elimination lemma
is displayed alongside); the payload structure is
\texttt{PReLUPosData}, quoted in full below.

\declhead{BernsteinLogic.Soundness}{sound\_reluPos}{\badge{BadgeGreen}{statement bound}}
\begin{leansrc}
/-- ReLU-Pos: from `Φ ⊢ h ≥ 0` conclude `Φ ⊢ y = h` — on each entry the
premise forces the nonpositive case's input to 0, where both cases
agree. -/
theorem sound_reluPos {H : Hypernet} {seq premise : Sequent}
    {preact post : Wire}
    (ih : premise.semanticallyValid H)
    (hlink : reluLinked H preact post)
    (hpre : seq.pre = premise.pre)
    (hpremPost : premise.post = .atom .ge (.var preact) (.const 0))
    (hpost : seq.post = .atom .eq (.var post) (.var preact)) :
    seq.semanticallyValid H
\end{leansrc}
\declhead{BernsteinLogic.Soundness}{QTensor.rel2\_const\_right\_elim}{\badge{BadgeGreen}{statement bound}}
\begin{leansrc}
/-- rel2 against a rank-0 constant on the right, eliminated to the
all-entries relation. -/
theorem QTensor.rel2_const_right_elim {r : ℚ → ℚ → Prop} {q : QTensor}
    {c : ℚ} (h : QTensor.rel2 r q ⟨[], [c], rfl⟩) :
    ∀ x ∈ q.entries, r x c
\end{leansrc}
\declhead{BernsteinLogic.Core.Catalogue}{PReLUPosData}{\badge{BadgeGreen}{definition bound}}
\begin{leansrc}
inductive PReLUPosData : Hypernet → Sequent → Type where
  | mk (H : Hypernet) (seq : Sequent)
       (premise : Sequent) (premiseProof : ProofObject H premise)
       (preact post : Wire)
       (link : reluLinked H preact post)
       (preEq       : seq.pre = premise.pre)
       (premisePost : premise.post = .atom .ge (.var preact) (.const 0))
       (conclPost   : seq.post     = .atom .eq (.var post) (.var preact))
       : PReLUPosData H seq
\end{leansrc}
\ruleblock{Correspondence (Steele display $\leftrightarrow$ Lean)}
\badge{BadgeGreen}{reviewer-certified against the displayed texts}
\sourcenote{The dictionary below is a reviewer-certified statement of how the mechanized statements above interpret the displayed rule: display metavariable $\mapsto$ Lean binder, each display hypothesis $\mapsto$ the premise or side condition carrying it, the display conclusion $\mapsto$ the goal --- with any deliberate difference in width between display and mechanization named outright. Certification is recorded as a hash triple over (the displayed cell, the extracted statement texts, this dictionary\textquotesingle s own text); if any of the three changes, this block turns amber until a reviewer re-certifies. The correspondence claim itself remains editorial --- reviewer-grade, not machine-checked; its machine-checked successor types each soundness statement against the rule as a structured object, and supersedes this block per rule as it lands.}
Judgment frame as in \rulename{Ax}; sub-derivation premise as at
\rulename{Weaken}; linkage premise --- including the positive-size
narrowing hidden under the displayed $\in$ --- exactly as in the
\rulename{LiRPA} dictionary, the family preamble.
\begin{itemize}
\item $h \mapsto$ \texttt{preact}; $y \mapsto$ \texttt{post};
  $\ReLU_{\{h\}\to\{y\}} \in H$ $\mapsto$ \texttt{hlink} $:$
  \texttt{reluLinked H preact post} (payload alias:
  \texttt{PReLUPosData.link}, same proposition).
\item premise $\Phi \vdash h \geq 0$ $\mapsto$ \texttt{ih} with
  \texttt{hpremPost} $:$ \texttt{premise.post = .atom .ge (.var
  preact) (.const 0)} --- displayed comparator and operand order
  carried exactly, with no sign-folding (in contrast to the
  \rulename{ConvElim} family\textquotesingle s all-$\leq$ build
  output); $\Phi$ shared via \texttt{hpre}. Payload aliases, same
  propositions: \texttt{ih} $\leftrightarrow$
  \texttt{premiseProof} (over the shared \texttt{premise}),
  \texttt{hpre} $\leftrightarrow$ \texttt{preEq},
  \texttt{hpremPost} $\leftrightarrow$ \texttt{premisePost}.
\item conclusion $\Phi \vdash y = h$ $\mapsto$ the goal under
  \texttt{hpost} $:$ \texttt{seq.post = .atom .eq (.var post) (.var
  preact)} (payload alias: \texttt{conclPost}, same equation) ---
  output wire left, preactivation right, as displayed.
\end{itemize}
Width note: the certified atoms are elementwise over the cited
application\textquotesingle s tensor shape (batch semantics, as at
\rulename{LiRPA}); the boundary $h = 0$ needs no side condition ---
the node\textquotesingle s two semantic cases agree at the kink, by
one antisymmetry step in the model (the commentary and tier note).
Beyond the family\textquotesingle s linkage gate, neither side is
wider: every hypothesis of \texttt{sound\_reluPos} is consumed.

\ruleblock{Lean verifier receipts}
\sourcenote{module \texttt{BernsteinLogic.Soundness}; bound build \texttt{a20d90e66c13\ldots} (bound from receipt a9c0e34d3e3c9c8c, 2026-08-03T21:30:36Z). A receipt records that named checking layers accepted a build: (module, compiled-artifact SHA-256, environment SHA-256, checker layers), minted when the checking gate ran; the store receipts attesting the bound build are listed with the layers each run credited. Nothing here asks the reader to trust the store: every layer is an independently runnable program over the quoted sources --- rebuild the module and re-run the checkers to reproduce the verdict. Receipts attesting superseded builds credit no layer and are omitted from this rendering (the count of omissions, if any, is noted below).}
\begin{receipttable}
\receiptrow{a9c0e34d3e3c9c8c}{2026-08-03T21:30:36Z}{build+decorrelated}{\badge{BadgeGreen}{current}}
\end{receipttable}
\sourcenote{21 further receipt(s) in the store attest superseded builds of this module and are omitted from this rendering.}
\sourcenote{module \texttt{BernsteinLogic.Core.Catalogue}; bound build \texttt{cbc92e976a1d\ldots} (bound from receipt 5abdb3589707b5df, 2026-08-03T21:30:33Z). A receipt records that named checking layers accepted a build: (module, compiled-artifact SHA-256, environment SHA-256, checker layers), minted when the checking gate ran; the store receipts attesting the bound build are listed with the layers each run credited. Nothing here asks the reader to trust the store: every layer is an independently runnable program over the quoted sources --- rebuild the module and re-run the checkers to reproduce the verdict. Receipts attesting superseded builds credit no layer and are omitted from this rendering (the count of omissions, if any, is noted below).}
\begin{receipttable}
\receiptrow{24804e66edba0bc9}{2026-08-03T20:25:09Z}{build+decorrelated}{\badge{BadgeGreen}{current}}
\receiptrow{5abdb3589707b5df}{2026-08-03T21:30:33Z}{build+decorrelated}{\badge{BadgeGreen}{current}}
\end{receipttable}
\sourcenote{18 further receipt(s) in the store attest superseded builds of this module and are omitted from this rendering.}
\layerline{\layeritem{build}{lake build (elaboration)}{\badge{BadgeGreen}{current}}\layeritem{replay}{leanchecker (official kernel replay)}{\badge{BadgeAmber}{stale only}}\layeritem{hardened}{safe\_verify (hardened screen)}{\badge{BadgeAmber}{stale only}}\layeritem{decorrelated}{lean4export + nanoda (decorrelated kernel)}{\badge{BadgeGreen}{current}}}
\sourcenote{\badge{BadgeGray}{comparator: none} no fixture in the three-implementation comparator harness (compiled / kernel-interpreted / reference-interpreted) exercises this rule.}
\ruleblock{HOL4 statement}
\sourcenote{theory \texttt{bern\_tr\_catdata2}, HOL4 development script \texttt{\texttt{bern\_tr\_catdata2Script.sml}} \badge{BadgeGreen}{script bound}}
\holhead{Theorem preluposdata\_receipt}
\begin{holsrc}
Theorem preluposdata_receipt:
  ruledataOk [] cd2Hrelu
    <| pre := FAtom BowGe (TmVar 3) (TmConst 0) ;
       post := FAtom BowEq (TmVar 4) (TmVar 3) |>
    (RdReluPos (PReluPosD
       <| pre := FAtom BowGe (TmVar 3) (TmConst 0) ;
          post := FAtom BowGe (TmVar 3) (TmConst 0) |> (ProofObj RdAx)
       3 4)) ∧
  ¬ ruledataOk [] hnEmpty
      <| pre := FAtom BowGe (TmVar 3) (TmConst 0) ;
         post := FAtom BowEq (TmVar 4) (TmVar 3) |>
      (RdReluPos (PReluPosD
         <| pre := FAtom BowGe (TmVar 3) (TmConst 0) ;
            post := FAtom BowGe (TmVar 3) (TmConst 0) |> (ProofObj RdAx)
         3 4))
\end{holsrc}
\ruleblock{HOL verifier receipts}
\badge{BadgeBlue}{documentary} \texttt{hol-strict} CLEAN: theory \texttt{bern\_tr\_catdata2}, 272
theorems, 0 cheat-tagged, 0 nonstandard axioms (whitelist:
\texttt{DISK\_THM}; the four \texttt{bool} axioms) --- verification
run of 2026-07-22 on an isolated build host, against a
content-addressed checkout of the exact sources cited here (git blob
\texttt{cd3c1e95\ldots}). The proof-object stratum is ported as one
mutual datatype family with a mutual wellformedness predicate,
exactly as the Lean mutual inductive block; this
rule\textquotesingle s arm states its entire checkable content ---
the premise post-form, the family linkage, and the conclusion
equation --- and the linkage lookup into the
hypernet\textquotesingle s application list has a ported realization
(\texttt{reluLinkedB}, theory \texttt{bern\_tr\_catdata}, 252
theorems, 0 offences, same run). The receipt theorem cited here is
the kernel-evaluated pair: an accept over a hypernet carrying the
packaged ReLU application, and a refusal over one without it ---
linkage doing exactly its job. Honest residue: the semantic soundness
twin (the Lean development\textquotesingle s \texttt{sound\_reluPos})
is not yet ported; the payload/checker layer above is.

\sourcenote{HOL receipts are documentary: quoted results of \texttt{hol-strict} runs
and census rows from the HOL4 development\textquotesingle s build
evidence, not minted receipt artefacts. A HOL receipt store (the Lean
store\textquotesingle s twin) is named future work.
}
\ruleblock{Proof}

\sourcenote{Genre: admissible rule (tier 2). The family machinery ---
the packaged node, the linkage, the two-case semantic fact, and its
conservativity warrant --- is stated in the \rulename{LiRPA} section's
family template and is not restated here.}

\begin{brecipe}[the \rulename{ReLU-Pos} instance]\label{rec:relupos-instance}
The node, and its intended rewrite over the family's sign-split fact:
\[
\inferrule*[Right=\rulename{ReLU-Pos}]{
  \mathcal{T} \rhd \Phi \vdash h \geq 0 \\
  \ReLU_{\{h\}\to\{y\}} \in H
}{
  \Phi \vdash y = h
}
\qquad\rightsquigarrow\qquad
\text{case analysis on }
(h \leq 0 \wedge y = 0) \vee (0 \leq h \wedge y = h).
\]
The \emph{active} branch discharges cleanly: its context carries the
conjunct $y = h$ (the case disjunct enters the branch premise as one
ordered subtree, so two $\wedge$E steps reach it), which \rulename{Ax}
and $\wedge$E extract --- the goal verbatim. The \emph{boundary} branch does not: its context
carries $h \leq 0$ and $y = 0$, the premise supplies $h \geq 0$, and
the goal $y = h$ now needs the antisymmetry fact $h = 0$ --- which, as
\autoref{rem:relupos-antisym} shows, no catalogue derivation that
withholds this rule itself provides. The rewrite therefore does \emph{not} fully discharge into
catalogue nodes; like \rulename{PConstExact}, the rule is admissible
by reduction to verified machinery: the shipped warrant is the direct
lemma \texttt{sound\_reluPos}, in which the boundary case is a single
antisymmetry step in the model ($p \leq 0$ from the case, $0 \leq p$
from the premise, so $p = 0$ and the two cases agree).
\end{brecipe}

\begin{bprop}[semantic warrant]\label{prop:relupos-sound}
Suppose the linkage holds and the premise sequent --- same premise
formula, conclusion the atom $h \geq 0$ --- is semantically valid.
Then the sequent concluding $y = h$ is semantically valid.
\end{bprop}
\begin{proof}
Fix an admissible valuation $\sigma$ satisfying $\Phi$. The premise
eliminates entrywise: every entry $p$ of the preactivation tensor
satisfies $0 \leq p$. The linkage extraction
(\texttt{reluLinked\_facts}) gives the output tensor at the same
shape and, per entry, either $p \leq 0 \wedge y = 0$ or $0 \leq p
\wedge y = p$. In the second case $y = p$ is the goal entry. In the
first, $p \leq 0$ meets the premise's $0 \leq p$, so $p = 0$
(antisymmetry --- the model step the catalogue lacks) and $y = 0 = p$.
The conclusion atom is the entrywise equality of the two tensors at
their shared shape. That is \texttt{sound\_reluPos}, step for step.
\end{proof}

\begin{bremark}[the kink case has no catalogue derivation]
\label{rem:relupos-antisym}
The claim in Recipe~\ref{rec:relupos-instance} is an enumeration over
the catalogue's equality-capable conclusions. \rulename{Norm} and
\rulename{EqRefl} conclude only canonical-form identities ($h$ and $0$
are not canonically equal); \rulename{EqSym}, \rulename{EqElim}, and
\rulename{ArithAtom} at the $=$ comparator only \emph{transport}
equalities already proven; this family and the cell rules conclude
equalities only over their packaged applications' ports; the
combination rules (\rulename{Farkas}, \rulename{Syzygy},
\rulename{PBound}, \rulename{BernCone} with an order predicate)
conclude order atoms, never $=$; and \rulename{BernCone}'s degenerate
conic-predicate route would need a region atom pinning $h$ to the
single point $0$, which no rule derives --- region atoms denote
regions of positive-dimensional simplices. The catalogue at large \emph{does} reach the fact, and the route is
instructive: over the two-wire hypernet carrying only the linked ReLU
application, \rulename{ReLU-Neg} (premise $h \leq 0$, an
\rulename{Ax} extraction) concludes $y = 0$, this rule (premise
$h \geq 0$) concludes $y = h$, and one \rulename{EqElim} along
$y = h$ rewrites $y = 0$ into $h = 0$ --- the linked application is
itself the antisymmetry oracle, its output forced to agree with both
branches at the kink. But the elaboration needs the fact
\emph{without the rule being eliminated} --- the leaf-first species
of the admissible template consumes every rule except that one ---
and with either envelope rule withheld the supply dries up. Write
$d(t = u)$ for the canonical form of $t - u$, and consider a smallest
derivation, withholding one envelope rule, that concludes an equality
atom with $d$ \emph{not} a polynomial multiple of the surviving port
difference ($y$ if \rulename{ReLU-Pos} is withheld with
\rulename{ReLU-Neg}'s $y = 0$ surviving; $y - h$ mirrored). Its final
step cannot be the surviving port conclusion or a canonical-form
identity ($d$ a multiple of the port difference, respectively $0$);
not \rulename{EqSym}, the zero-form respelling, or \rulename{EqElim}
transport (each changes $d$ by a sign or by polynomial multiples of a
previously concluded equality's own $d$ --- a smaller derivation,
hence a permitted multiple); not a closed ground fact ($d = 0$, the
constants being exactly evaluated); not a combination or region rule
(order and region conclusions only); and not a structural or
case-analysis step, which only re-exposes an equality already
concluded, or sitting in a premise set --- the root premises
($h \leq 0 \wedge h \geq 0$, at either spelling) contain no equality,
premise sets grow only by \rulename{Cut} formulas and case disjuncts
concluded or contained in smaller derivations, and a case split whose
smuggled disjunct carries the target is idle: its sibling branch must
conclude the same target without it, a smaller offender. The target
$d(h = 0) = h$ is a multiple of neither $y$ nor $y - h$, so $h \leq 0
\wedge h \geq 0 \vdash h = 0$ has no \rulename{ReLU-Pos}-free
derivation (nor, mirrored, a \rulename{ReLU-Neg}-free one), and the
boundary branch cannot discharge into nodes. This is not a defect to
repair but a fact about where the calculus keeps its arithmetic: an
elaboration can source order antisymmetry only from the verified
model step, one line of \texttt{sound\_reluPos}, rather than from a
rule every certificate would pay wire format for.
\end{bremark}

\begin{bremark}[what the premise is for]
Without $\Phi \vdash h \geq 0$ the rule is false at any entry with $h
< 0$: there the node's semantics forces $y = 0 \neq h$. The premise is
consumed \emph{entrywise} --- a tensor-shaped preactivation must be
proven nonnegative in every entry, which is what a rank-$0$ constant
against a tensor atom asserts. No side condition excludes the boundary
$h = 0$, and none is needed: the two semantic cases agree there, which
is exactly the antisymmetry step above.
\end{bremark}

\begin{anchors}
Recipe \ref{rec:relupos-instance} warrant & \texttt{sound\_reluPos}\\
family linkage extraction & \texttt{reluLinked\_facts} (see the \rulename{LiRPA} section)\\
entrywise premise elimination & \texttt{QTensor.rel2\_const\_right\_elim}\\
the kink antisymmetry & the \texttt{le\_antisymm} step inside \texttt{sound\_reluPos}\\
payload structure & \texttt{PReLUPosData} (premise proof, two wires, linkage)\\
\end{anchors}


\rulesection{21}{ReLU-Neg}{admissible rule}
\label{blcatrule-ReLUNeg}
\ruleblock{The rule (Steele metanotation)}
\sourcenote{as realized in the mechanized development; source row: \texttt{\texttt{rules.toml}} --- \emph{ReLU-Neg}. The metanotation statement is typeset by hand for notational quality; its mechanized ground truth is the Lean statement below, extracted from the sources at every build of this book and hash-bound. The rule \emph{roster} is machine-checked (every mechanized rule tag round-trips to exactly one catalogue row); the statement \emph{text} is editorial, checked by eye against the extracted Lean beside it.}
\begin{mathpar}
\inferrule*[Right=\rulename{ReLU-Neg}]{
  \Phi \vdash h \leq 0
  \\
  \ReLU_{\{h\}\to\{y\}} \in H
}{
  \Phi \vdash y = 0
}
\end{mathpar}

\par\medskip
\rulename{ReLU-Neg} is the exactness rule for the inactive side of a
ReLU, \rulename{ReLU-Pos}\textquotesingle s mirror: once $\Phi \vdash
h \leq 0$ is proven for the preactivation wire of a cited ReLU
application, the output is pinned to a constant — $\Phi \vdash y = 0$.
A producer reaches for it when bound propagation has shown a neuron
stably inactive (a proven upper box bound at or below zero): the
neuron is dead on the region of interest, and the equality lets
\rulename{EqElim} rewrite $y$ to the literal constant $0$ downstream
— dead-neuron elimination in certificate form, at the rank-0
instances the constant-sided transport gate admits (the
family\textquotesingle s shape-scope remark).

What the checker recomputes, given the payload (premise proof, the two
wires): the linkage — a packaged ReLU application at exactly these
ports, of positive total size (the family\textquotesingle s
port-lookup side condition, displayed in the \rulename{LiRPA}
section); that the premise\textquotesingle s conclusion is literally
the atom $h \leq 0$ over the same premise formula; and that the
conclusion is literally $y = 0$. Both atoms are elementwise over the
node\textquotesingle s tensor shape.

Drop the premise and the rule is false at any entry with $h > 0$
(there $y = h \neq 0$). Drop the linkage and an arbitrary wire is
asserted zero. As with \rulename{ReLU-Pos}, the boundary $h = 0$ needs
no exclusion: both semantic cases of the node force $y = 0$ there.

\cataloguenote{TIER NOTE: as for \rulename{ReLU-Pos}, and for the mirrored reason —
the unpacking\textquotesingle s active branch would need $h = 0$ from
$h \geq 0 \wedge h \leq 0$, an antisymmetry step no catalogue rule
provides — the rule is admissible by reduction to verified machinery:
the shipped warrant is the direct lemma \texttt{sound\_reluNeg}. The
enumeration argument is stated once, in the \rulename{ReLU-Pos}
proof section.
}
\ruleblock{Lean statement}
\par\medskip
The warrant consumes the family linkage extraction
(\texttt{reluLinked\_facts}, displayed in the \rulename{LiRPA}
section) and the premise, eliminated entrywise (the elimination lemma
is displayed alongside); the payload structure is
\texttt{PReLUNegData}, quoted in full below.

\declhead{BernsteinLogic.Soundness}{sound\_reluNeg}{\badge{BadgeGreen}{statement bound}}
\begin{leansrc}
/-- ReLU-Neg: from `Φ ⊢ h ≤ 0` conclude `Φ ⊢ y = 0` — on each entry the
premise forces the nonnegative case's input to 0, where both cases
agree. -/
theorem sound_reluNeg {H : Hypernet} {seq premise : Sequent}
    {preact post : Wire}
    (ih : premise.semanticallyValid H)
    (hlink : reluLinked H preact post)
    (hpre : seq.pre = premise.pre)
    (hpremPost : premise.post = .atom .le (.var preact) (.const 0))
    (hpost : seq.post = .atom .eq (.var post) (.const 0)) :
    seq.semanticallyValid H
\end{leansrc}
\declhead{BernsteinLogic.Soundness}{QTensor.rel2\_const\_right\_elim}{\badge{BadgeGreen}{statement bound}}
\begin{leansrc}
/-- rel2 against a rank-0 constant on the right, eliminated to the
all-entries relation. -/
theorem QTensor.rel2_const_right_elim {r : ℚ → ℚ → Prop} {q : QTensor}
    {c : ℚ} (h : QTensor.rel2 r q ⟨[], [c], rfl⟩) :
    ∀ x ∈ q.entries, r x c
\end{leansrc}
\declhead{BernsteinLogic.Core.Catalogue}{PReLUNegData}{\badge{BadgeGreen}{definition bound}}
\begin{leansrc}
inductive PReLUNegData : Hypernet → Sequent → Type where
  | mk (H : Hypernet) (seq : Sequent)
       (premise : Sequent) (premiseProof : ProofObject H premise)
       (preact post : Wire)
       (link : reluLinked H preact post)
       (preEq       : seq.pre = premise.pre)
       (premisePost : premise.post = .atom .le (.var preact) (.const 0))
       (conclPost   : seq.post     = .atom .eq (.var post) (.const 0))
       : PReLUNegData H seq
\end{leansrc}
\ruleblock{Correspondence (Steele display $\leftrightarrow$ Lean)}
\badge{BadgeGreen}{reviewer-certified against the displayed texts}
\sourcenote{The dictionary below is a reviewer-certified statement of how the mechanized statements above interpret the displayed rule: display metavariable $\mapsto$ Lean binder, each display hypothesis $\mapsto$ the premise or side condition carrying it, the display conclusion $\mapsto$ the goal --- with any deliberate difference in width between display and mechanization named outright. Certification is recorded as a hash triple over (the displayed cell, the extracted statement texts, this dictionary\textquotesingle s own text); if any of the three changes, this block turns amber until a reviewer re-certifies. The correspondence claim itself remains editorial --- reviewer-grade, not machine-checked; its machine-checked successor types each soundness statement against the rule as a structured object, and supersedes this block per rule as it lands.}
Judgment frame as in \rulename{Ax}; sub-derivation premise as at
\rulename{Weaken}; linkage premise --- including the positive-size
narrowing hidden under the displayed $\in$ --- exactly as in the
\rulename{LiRPA} dictionary, the family preamble.
\begin{itemize}
\item $h \mapsto$ \texttt{preact}; $y \mapsto$ \texttt{post};
  $\ReLU_{\{h\}\to\{y\}} \in H$ $\mapsto$ \texttt{hlink} $:$
  \texttt{reluLinked H preact post} (payload alias:
  \texttt{PReLUNegData.link}, same proposition).
\item premise $\Phi \vdash h \leq 0$ $\mapsto$ \texttt{ih} with
  \texttt{hpremPost} $:$ \texttt{premise.post = .atom .le (.var
  preact) (.const 0)} --- displayed comparator and operand order
  carried exactly; $\Phi$ shared via \texttt{hpre}. Payload aliases, same
  propositions: \texttt{ih} $\leftrightarrow$
  \texttt{premiseProof} (over the shared \texttt{premise}),
  \texttt{hpre} $\leftrightarrow$ \texttt{preEq},
  \texttt{hpremPost} $\leftrightarrow$ \texttt{premisePost}.
\item conclusion $\Phi \vdash y = 0$ $\mapsto$ the goal under
  \texttt{hpost} $:$ \texttt{seq.post = .atom .eq (.var post)
  (.const 0)} (payload alias: \texttt{conclPost}, same equation) --- output
  wire left, constant right, as displayed.
\end{itemize}
Width note: certified atoms elementwise over the cited
application\textquotesingle s tensor shape (batch semantics, as at
\rulename{LiRPA}); the boundary $h = 0$ is again condition-free ---
the mirrored antisymmetry step of the tier note. Beyond the
family\textquotesingle s linkage gate, neither side is wider: every
hypothesis of \texttt{sound\_reluNeg} is consumed.

\ruleblock{Lean verifier receipts}
\sourcenote{module \texttt{BernsteinLogic.Soundness}; bound build \texttt{a20d90e66c13\ldots} (bound from receipt a9c0e34d3e3c9c8c, 2026-08-03T21:30:36Z). A receipt records that named checking layers accepted a build: (module, compiled-artifact SHA-256, environment SHA-256, checker layers), minted when the checking gate ran; the store receipts attesting the bound build are listed with the layers each run credited. Nothing here asks the reader to trust the store: every layer is an independently runnable program over the quoted sources --- rebuild the module and re-run the checkers to reproduce the verdict. Receipts attesting superseded builds credit no layer and are omitted from this rendering (the count of omissions, if any, is noted below).}
\begin{receipttable}
\receiptrow{a9c0e34d3e3c9c8c}{2026-08-03T21:30:36Z}{build+decorrelated}{\badge{BadgeGreen}{current}}
\end{receipttable}
\sourcenote{21 further receipt(s) in the store attest superseded builds of this module and are omitted from this rendering.}
\sourcenote{module \texttt{BernsteinLogic.Core.Catalogue}; bound build \texttt{cbc92e976a1d\ldots} (bound from receipt 5abdb3589707b5df, 2026-08-03T21:30:33Z). A receipt records that named checking layers accepted a build: (module, compiled-artifact SHA-256, environment SHA-256, checker layers), minted when the checking gate ran; the store receipts attesting the bound build are listed with the layers each run credited. Nothing here asks the reader to trust the store: every layer is an independently runnable program over the quoted sources --- rebuild the module and re-run the checkers to reproduce the verdict. Receipts attesting superseded builds credit no layer and are omitted from this rendering (the count of omissions, if any, is noted below).}
\begin{receipttable}
\receiptrow{24804e66edba0bc9}{2026-08-03T20:25:09Z}{build+decorrelated}{\badge{BadgeGreen}{current}}
\receiptrow{5abdb3589707b5df}{2026-08-03T21:30:33Z}{build+decorrelated}{\badge{BadgeGreen}{current}}
\end{receipttable}
\sourcenote{18 further receipt(s) in the store attest superseded builds of this module and are omitted from this rendering.}
\layerline{\layeritem{build}{lake build (elaboration)}{\badge{BadgeGreen}{current}}\layeritem{replay}{leanchecker (official kernel replay)}{\badge{BadgeAmber}{stale only}}\layeritem{hardened}{safe\_verify (hardened screen)}{\badge{BadgeAmber}{stale only}}\layeritem{decorrelated}{lean4export + nanoda (decorrelated kernel)}{\badge{BadgeGreen}{current}}}
\sourcenote{\badge{BadgeGray}{comparator: none} no fixture in the three-implementation comparator harness (compiled / kernel-interpreted / reference-interpreted) exercises this rule.}
\ruleblock{HOL4 statement}
\sourcenote{theory \texttt{bern\_tr\_catdata2}, HOL4 development script \texttt{\texttt{bern\_tr\_catdata2Script.sml}} \badge{BadgeGreen}{script bound}}
\holhead{Theorem prelunegdata\_receipt}
\begin{holsrc}
Theorem prelunegdata_receipt:
  ruledataOk [] cd2Hrelu
    <| pre := FAtom BowLe (TmVar 3) (TmConst 0) ;
       post := FAtom BowEq (TmVar 4) (TmConst 0) |>
    (RdReluNeg (PReluNegD
       <| pre := FAtom BowLe (TmVar 3) (TmConst 0) ;
          post := FAtom BowLe (TmVar 3) (TmConst 0) |> (ProofObj RdAx)
       3 4)) ∧
  ¬ ruledataOk [] cd2Hrelu
      <| pre := FAtom BowLe (TmVar 3) (TmConst 0) ;
         post := FAtom BowEq (TmVar 4) (TmVar 3) |>
      (RdReluNeg (PReluNegD
         <| pre := FAtom BowLe (TmVar 3) (TmConst 0) ;
            post := FAtom BowLe (TmVar 3) (TmConst 0) |> (ProofObj RdAx)
         3 4))
\end{holsrc}
\ruleblock{HOL verifier receipts}
\badge{BadgeBlue}{documentary} \texttt{hol-strict} CLEAN: theory \texttt{bern\_tr\_catdata2}, 272
theorems, 0 cheat-tagged, 0 nonstandard axioms (whitelist:
\texttt{DISK\_THM}; the four \texttt{bool} axioms) --- verification
run of 2026-07-22 on an isolated build host, against a
content-addressed checkout of the exact sources cited here (git blob
\texttt{cd3c1e95\ldots}). The proof-object stratum is ported as one
mutual datatype family with a mutual wellformedness predicate,
exactly as the Lean mutual inductive block; this
rule\textquotesingle s arm states its entire checkable content ---
the premise post-form, the family linkage, and the zero-clamp
conclusion equation --- and the linkage lookup into the
hypernet\textquotesingle s application list has a ported realization
(\texttt{reluLinkedB}, theory \texttt{bern\_tr\_catdata}, 252
theorems, 0 offences, same run). The receipt theorem cited here is
the kernel-evaluated pair: an accept over a hypernet carrying the
packaged ReLU application, and a wrong-conclusion refusal. Honest
residue: the semantic soundness twin (the Lean
development\textquotesingle s \texttt{sound\_reluNeg}) is not yet
ported; the payload/checker layer above is.

\sourcenote{HOL receipts are documentary: quoted results of \texttt{hol-strict} runs
and census rows from the HOL4 development\textquotesingle s build
evidence, not minted receipt artefacts. A HOL receipt store (the Lean
store\textquotesingle s twin) is named future work.
}
\ruleblock{Proof}

\sourcenote{Genre: admissible rule (tier 2). The family machinery is
stated in the \rulename{LiRPA} section's family template;
\rulename{ReLU-Pos} is this rule's mirror, and the shared boundary
argument is stated there.}

\begin{brecipe}[the \rulename{ReLU-Neg} instance]\label{rec:reluneg-instance}
The node, and its intended rewrite over the family's sign-split fact:
\[
\inferrule*[Right=\rulename{ReLU-Neg}]{
  \mathcal{T} \rhd \Phi \vdash h \leq 0 \\
  \ReLU_{\{h\}\to\{y\}} \in H
}{
  \Phi \vdash y = 0
}
\qquad\rightsquigarrow\qquad
\text{case analysis on }
(h \leq 0 \wedge y = 0) \vee (0 \leq h \wedge y = h).
\]
The \emph{inactive} branch discharges cleanly: its context carries the
conjunct $y = 0$ (one ordered subtree of the branch premise; two
$\wedge$E steps), extracted by \rulename{Ax} and $\wedge$E --- the goal
verbatim. The \emph{boundary} branch is the mirror of
\rulename{ReLU-Pos}'s: its context carries $0 \leq h$ and $y = h$, the
premise supplies $h \leq 0$, and the goal $y = 0$ needs $h = 0$ --- the
antisymmetry fact that no catalogue derivation withholding this rule
itself provides (the enumeration argument, covering both withheld
orientations, is stated once, as
\autoref{rem:relupos-antisym} in the \rulename{ReLU-Pos} section).
The rule is therefore admissible by reduction to verified machinery:
the shipped warrant is the direct lemma \texttt{sound\_reluNeg}, whose
boundary case is the same one-line antisymmetry step in the model.
\end{brecipe}

\begin{bprop}[semantic warrant]\label{prop:reluneg-sound}
Suppose the linkage holds and the premise sequent --- same premise
formula, conclusion the atom $h \leq 0$ --- is semantically valid.
Then the sequent concluding $y = 0$ is semantically valid.
\end{bprop}
\begin{proof}
Fix an admissible valuation $\sigma$ satisfying $\Phi$. The premise
eliminates entrywise: every entry $p$ of the preactivation tensor
satisfies $p \leq 0$. The linkage extraction
(\texttt{reluLinked\_facts}) gives, per entry, either $p \leq 0 \wedge
y = 0$ --- the goal entry directly --- or $0 \leq p \wedge y = p$, in
which case $p = 0$ by antisymmetry against the premise, so $y = p =
0$. The conclusion atom compares the output tensor entrywise against
the rank-$0$ constant $0$. That is \texttt{sound\_reluNeg}, step for
step.
\end{proof}

\begin{bremark}[what the premise is for, and what the rule feeds]
Without $\Phi \vdash h \leq 0$ the rule is false at any entry with
$h > 0$ (there $y = h \neq 0$). Downstream, the conclusion differs
from \rulename{ReLU-Pos}'s in kind: $y = 0$ pins the output to a
\emph{literal constant}, so \rulename{EqElim} can rewrite $y$ away ---
dead-neuron elimination inside a certificate --- at the rank-$0$
instances the transport gate admits (a constant-sided equality rides
\rulename{EqElim}'s shared-broadcast-shape condition; the family's
shape-scope remark, \rulename{ReLU-NN} section), after which ground
arithmetic on the constant rides \rulename{PConstExact}.
\end{bremark}

\begin{anchors}
Recipe \ref{rec:reluneg-instance} warrant & \texttt{sound\_reluNeg}\\
family linkage extraction & \texttt{reluLinked\_facts} (see the \rulename{LiRPA} section)\\
entrywise premise elimination & \texttt{QTensor.rel2\_const\_right\_elim}\\
the kink antisymmetry & the \texttt{le\_antisymm} step inside \texttt{sound\_reluNeg}\\
payload structure & \texttt{PReLUNegData} (premise proof, two wires, linkage)\\
\end{anchors}


\rulesection{22}{ReLU-NN}{admissible rule}
\label{blcatrule-ReLUNN}
\ruleblock{The rule (Steele metanotation)}
\sourcenote{as realized in the mechanized development; source row: \texttt{\texttt{rules.toml}} --- \emph{ReLU-NN}. The metanotation statement is typeset by hand for notational quality; its mechanized ground truth is the Lean statement below, extracted from the sources at every build of this book and hash-bound. The rule \emph{roster} is machine-checked (every mechanized rule tag round-trips to exactly one catalogue row); the statement \emph{text} is editorial, checked by eye against the extracted Lean beside it.}
\begin{mathpar}
\inferrule*[Right=\rulename{ReLU-NN}]{
  \ReLU_{\{h\}\to\{y\}} \in H
}{
  \Phi \vdash y \geq 0
}
\end{mathpar}

\par\medskip
\rulename{ReLU-NN} is the range fact: the output wire of a cited ReLU
application is nonnegative, $\Phi \vdash y \geq 0$, with no proof
premises at all — $\Phi$ is arbitrary (the premise-free convention
shared with \rulename{Norm} and \rulename{PConstExact}). It is the
first thing a producer knows about every ReLU output: the lower half
of every output box starts here, before any information about the
preactivation exists, and together with \rulename{ReLU-Above} it
yields the exact pointwise lower envelope of the activation
($y \geq \max(0, h)$, as two atoms).

What the checker recomputes is nearly everything, because there is
nearly nothing: the linkage — the hypernet really contains a packaged
ReLU application at exactly these ports, of positive total size (the
family\textquotesingle s port-lookup side condition, displayed in the
\rulename{LiRPA} section) — and that the conclusion is literally the
atom $y \geq 0$, elementwise over the node\textquotesingle s tensor
shape.

The linkage is the only hypothesis, and dropping it is the whole
failure mode: an arbitrary wire asserted nonnegative, against an
arbitrary premise — the linkage-free rule is unsound over unlinked
wires, and a realization without the port lookup could only be kept
sound by rejecting every certificate. The positive-size gate inside the
linkage keeps even the vacuous-surface corner unciteable (a size-zero
node satisfies every elementwise atom trivially; the linkage refuses
it at the source).

\cataloguenote{This is the family\textquotesingle s cleanest unpacking: the case
analysis over the sign-split fact discharges \emph{both} branches into
catalogue nodes — the inactive branch\textquotesingle s leaf is the
ground fact $0 \geq 0$, which rides \rulename{PConstExact} (the one
rule that can state a closed ground fact). The
proof section displays the full intended derivation (its case-closing
disjunction elimination is the context-carrying form --- a satellite
theorem today, so the display is an elaboration blueprint rather than
an emittable certificate), and states its shape
scope honestly: the constant-sided equality transport type-checks at
rank-0 instances, and the shape-general warrant is the direct lemma.
}
\ruleblock{Lean statement}
\par\medskip
The warrant consumes the family linkage extraction
(\texttt{reluLinked\_facts}, displayed in the \rulename{LiRPA}
section): each output entry is either $0$ or a copy of a nonnegative
input entry. The payload structure, \texttt{PReLUNNData}
--- quoted in full below --- carries no premise field.

\declhead{BernsteinLogic.Soundness}{sound\_reluNN}{\badge{BadgeGreen}{statement bound}}
\begin{leansrc}
/-- ReLU-NN: `Φ ⊢ y ≥ 0`, premise-free — every entry of the packaged
output is 0 or a nonnegative copy of the input entry. (A review finding
exhibited the rule as unsound over unbound wires; the linkage is the repair.) -/
theorem sound_reluNN {H : Hypernet} {seq : Sequent} {preact post : Wire}
    (hlink : reluLinked H preact post)
    (hpost : seq.post = .atom .ge (.var post) (.const 0)) :
    seq.semanticallyValid H
\end{leansrc}
\declhead{BernsteinLogic.Core.Catalogue}{PReLUNNData}{\badge{BadgeGreen}{definition bound}}
\begin{leansrc}
inductive PReLUNNData : Hypernet → Sequent → Type where
  | mk (H : Hypernet) (seq : Sequent)
       (preact post : Wire)
       (link : reluLinked H preact post)
       (conclPost : seq.post = .atom .ge (.var post) (.const 0))
       : PReLUNNData H seq
\end{leansrc}
\ruleblock{Correspondence (Steele display $\leftrightarrow$ Lean)}
\badge{BadgeAmber}{a bound Lean statement CHANGED since certification --- re-review and re-certify}
\sourcenote{The dictionary below is a reviewer-certified statement of how the mechanized statements above interpret the displayed rule: display metavariable $\mapsto$ Lean binder, each display hypothesis $\mapsto$ the premise or side condition carrying it, the display conclusion $\mapsto$ the goal --- with any deliberate difference in width between display and mechanization named outright. Certification is recorded as a hash triple over (the displayed cell, the extracted statement texts, this dictionary\textquotesingle s own text); if any of the three changes, this block turns amber until a reviewer re-certifies. The correspondence claim itself remains editorial --- reviewer-grade, not machine-checked; its machine-checked successor types each soundness statement against the rule as a structured object, and supersedes this block per rule as it lands.}
Judgment frame as in \rulename{Ax}; linkage premise --- including the
positive-size narrowing hidden under the displayed $\in$ --- exactly
as in the \rulename{LiRPA} dictionary, the family preamble.
\begin{itemize}
\item $h \mapsto$ \texttt{preact}; $y \mapsto$ \texttt{post};
  $\ReLU_{\{h\}\to\{y\}} \in H$ $\mapsto$ \texttt{hlink} $:$
  \texttt{reluLinked H preact post} (payload alias:
  \texttt{PReLUNNData.link}, same proposition). The
  display\textquotesingle s
  universality over the preactivation --- no premise constrains $h$
  --- is carried by the
  linkage alone: \texttt{preact} occurs in no other hypothesis (the
  roster note records that the appendix payload named only the output
  wire --- the certified linkage names both ports).
\item premise-freeness $\mapsto$ silence, the \rulename{PConstExact}
  representation: the precondition component of \texttt{seq :
  Sequent} is unconstrained by any hypothesis, so the statement
  holds under arbitrary premises $\Phi$.
\item conclusion $\Phi \vdash y \geq 0$ $\mapsto$ the goal under
  \texttt{hpost} $:$ \texttt{seq.post = .atom .ge (.var post)
  (.const 0)} (payload alias: \texttt{conclPost}, same equation) ---
  displayed comparator and operand order carried exactly.
\end{itemize}
Width note: certified atoms elementwise over the cited
application\textquotesingle s tensor shape (batch semantics, as at
\rulename{LiRPA}). Beyond the family\textquotesingle s linkage gate,
neither side is wider: both hypotheses of \texttt{sound\_reluNN} are
consumed, and nothing else is stated.

\ruleblock{Lean verifier receipts}
\sourcenote{module \texttt{BernsteinLogic.Soundness}; bound build \texttt{a20d90e66c13\ldots} (bound from receipt a9c0e34d3e3c9c8c, 2026-08-03T21:30:36Z). A receipt records that named checking layers accepted a build: (module, compiled-artifact SHA-256, environment SHA-256, checker layers), minted when the checking gate ran; the store receipts attesting the bound build are listed with the layers each run credited. Nothing here asks the reader to trust the store: every layer is an independently runnable program over the quoted sources --- rebuild the module and re-run the checkers to reproduce the verdict. Receipts attesting superseded builds credit no layer and are omitted from this rendering (the count of omissions, if any, is noted below).}
\begin{receipttable}
\receiptrow{a9c0e34d3e3c9c8c}{2026-08-03T21:30:36Z}{build+decorrelated}{\badge{BadgeGreen}{current}}
\end{receipttable}
\sourcenote{21 further receipt(s) in the store attest superseded builds of this module and are omitted from this rendering.}
\sourcenote{module \texttt{BernsteinLogic.Core.Catalogue}; bound build \texttt{cbc92e976a1d\ldots} (bound from receipt 5abdb3589707b5df, 2026-08-03T21:30:33Z). A receipt records that named checking layers accepted a build: (module, compiled-artifact SHA-256, environment SHA-256, checker layers), minted when the checking gate ran; the store receipts attesting the bound build are listed with the layers each run credited. Nothing here asks the reader to trust the store: every layer is an independently runnable program over the quoted sources --- rebuild the module and re-run the checkers to reproduce the verdict. Receipts attesting superseded builds credit no layer and are omitted from this rendering (the count of omissions, if any, is noted below).}
\begin{receipttable}
\receiptrow{24804e66edba0bc9}{2026-08-03T20:25:09Z}{build+decorrelated}{\badge{BadgeGreen}{current}}
\receiptrow{5abdb3589707b5df}{2026-08-03T21:30:33Z}{build+decorrelated}{\badge{BadgeGreen}{current}}
\end{receipttable}
\sourcenote{18 further receipt(s) in the store attest superseded builds of this module and are omitted from this rendering.}
\layerline{\layeritem{build}{lake build (elaboration)}{\badge{BadgeGreen}{current}}\layeritem{replay}{leanchecker (official kernel replay)}{\badge{BadgeAmber}{stale only}}\layeritem{hardened}{safe\_verify (hardened screen)}{\badge{BadgeAmber}{stale only}}\layeritem{decorrelated}{lean4export + nanoda (decorrelated kernel)}{\badge{BadgeGreen}{current}}}
\sourcenote{\badge{BadgeGray}{comparator: none} no fixture in the three-implementation comparator harness (compiled / kernel-interpreted / reference-interpreted) exercises this rule.}
\ruleblock{HOL4 statement}
\sourcenote{theory \texttt{bern\_tr\_catdata2}, HOL4 development script \texttt{\texttt{bern\_tr\_catdata2Script.sml}} \badge{BadgeGreen}{script bound}}
\holhead{Theorem prelunndata\_receipt}
\begin{holsrc}
Theorem prelunndata_receipt:
  ruledataOk [] cd2Hrelu
    <| pre := FTt ; post := FAtom BowGe (TmVar 4) (TmConst 0) |>
    (RdReluNN (PReluNND 3 4)) ∧
  ¬ ruledataOk [] hnEmpty
      <| pre := FTt ; post := FAtom BowGe (TmVar 4) (TmConst 0) |>
      (RdReluNN (PReluNND 3 4))
\end{holsrc}
\ruleblock{HOL verifier receipts}
\badge{BadgeBlue}{documentary} \texttt{hol-strict} CLEAN: theory \texttt{bern\_tr\_catdata2}, 272
theorems, 0 cheat-tagged, 0 nonstandard axioms (whitelist:
\texttt{DISK\_THM}; the four \texttt{bool} axioms) --- verification
run of 2026-07-22 on an isolated build host, against a
content-addressed checkout of the exact sources cited here (git blob
\texttt{cd3c1e95\ldots}). The proof-object stratum is ported as one
mutual datatype family with a mutual wellformedness predicate,
exactly as the Lean mutual inductive block; this premise-free
rule\textquotesingle s arm states its entire checkable content ---
the family linkage plus the nonnegativity conclusion equation ---
and the linkage lookup into the hypernet\textquotesingle s
application list has a ported realization (\texttt{reluLinkedB},
theory \texttt{bern\_tr\_catdata}, 252 theorems, 0 offences, same
run). The receipt theorem cited here is the kernel-evaluated
accept/refusal pair for the arm. Honest residue: the semantic
soundness twin (the Lean development\textquotesingle s
\texttt{sound\_reluNN}) is not yet ported; the payload/checker layer
above is.

\sourcenote{HOL receipts are documentary: quoted results of \texttt{hol-strict} runs
and census rows from the HOL4 development\textquotesingle s build
evidence, not minted receipt artefacts. A HOL receipt store (the Lean
store\textquotesingle s twin) is named future work.
}
\ruleblock{Proof}

\sourcenote{Genre: admissible rule (tier 2). The family machinery is
stated in the \rulename{LiRPA} section's family template. This rule is
premise-free: $\Phi$ is arbitrary and never consulted (the
\rulename{Norm} convention).}

\begin{brecipe}[the \rulename{ReLU-NN} instance --- the full unpacking]
\label{rec:relunn-instance}
The node:
\[
\inferrule*[Right=\rulename{ReLU-NN}]{
  \ReLU_{\{h\}\to\{y\}} \in H
}{
  \Phi \vdash y \geq 0
}
\qquad\rightsquigarrow\qquad
\text{case analysis on }
(h \leq 0 \wedge y = 0) \vee (0 \leq h \wedge y = h),
\]
and both branches discharge into catalogue nodes --- the family's
cleanest unpacking. \emph{Inactive branch} (context carries $y = 0$):
\begin{enumerate}
\item $\vdash 0 \geq 0$ --- \rulename{PConstExact}, the closed exact
  ground fact (the leaf that rule exists to state: no other catalogue
  rule concludes a closed ground fact);
\item $\Phi_{-} \vdash y = 0$ --- \rulename{Ax} + $\wedge$E from the
  branch context (the case disjunct enters as one ordered subtree; two
  $\wedge$E steps); then $\Phi_{-} \vdash 0 = y$ by \rulename{EqSym};
\item $\Phi_{-} \vdash y \geq 0$ --- \rulename{EqElim} along $0 = y$,
  transporting $\Psi[a \coloneqq 0] = (0 \geq 0)$ (step 1, weakened)
  to $\Psi[a \coloneqq y]$, where $\Psi = (a \geq 0)$.
\end{enumerate}
\emph{Active branch} (context carries $0 \leq h$ and $y = h$): the
context's sign atom is extracted (\rulename{Ax} + $\wedge$E) and
re-oriented to $h \geq 0$ by the bound node's polarity-flip reading
(\rulename{PBound}, reading (G): the single component $0 \leq h$ with
weight $1$, canonical polynomial $c = -h$, read as $e \geq 0$ with
$e = h$ --- one node of the toolkit whose longer chains the
\rulename{ReLU-Above} section walks); \rulename{EqSym} flips the
context's $y = h$ to $h = y$, and \rulename{EqElim} along $h = y$ (an
equality between two wires of the node's one shared shape) transports
it to $y \geq 0$. $\vee$E in its context-carrying form --- today a
satellite theorem, not a core constructor (the split rule's section
prices exactly this), so what follows is the intended elaboration,
emittable when that constructor lands --- closes the case
analysis, and \rulename{Cut} stitches the sign-split fact into place.
\end{brecipe}

\begin{bremark}[the shape scope of the displayed derivation]
\label{rem:relunn-shapescope}
The realized \rulename{EqElim} carries a broadcast-shape side
condition (its section documents the counterexample that forced it):
both sides of the transported equality must share a broadcast shape.
Step 3's equality $0 = y$ has a rank-$0$ constant on one side, so the
displayed transport type-checks where the output wire is rank-$0$ ---
at tensor-shaped instances that step is stopped by the gate, and the family's
realized warrant does not route through the derivation at all: it is
the direct lemma \texttt{sound\_reluNN}, shape-general, via the
linkage extraction. (Wire-to-wire transports such as the active
branch's $h = y$ are unaffected: the linkage forces one shared shape.)
The derivation above is thus the intended admissibility story at
scalar instances --- intended, since its case-closing step awaits the
constructor named at its step --- and the
reduction-to-verified-machinery story, the
\rulename{PConstExact} pattern stated in the family template, is
the standing warrant at every shape either way.
\end{bremark}

\begin{bprop}[semantic warrant]\label{prop:relunn-sound}
If the linkage holds, the sequent concluding $y \geq 0$ is
semantically valid over any premise formula $\Phi$.
\end{bprop}
\begin{proof}
Fix an admissible valuation $\sigma$ ($\sigma \models \Phi$ is taken
but never consulted --- the rule is premise-free). The linkage
extraction (\texttt{reluLinked\_facts}) gives the output tensor and,
per entry, either $y = 0$ or $y = p$ with $0 \leq p$; in both cases
$0 \leq y$. The conclusion atom compares the output tensor entrywise
against the rank-$0$ constant $0$. That is \texttt{sound\_reluNN},
step for step.
\end{proof}

\begin{bremark}[the only hypothesis]
The linkage is the rule's sole side condition --- the conclusion-shape
check beside it is a reconstruction pin, not a droppable hypothesis:
the checker rebuilds the atom $y \geq 0$ rather than trusting it, so
``dropping'' the pin merely re-addresses the rule at a conclusion it
never concluded. Dropping the linkage is the whole failure mode: an
arbitrary wire asserted nonnegative against an arbitrary premise. Its positive-size gate additionally keeps the
vacuous corner unciteable --- a size-zero node satisfies every
elementwise atom trivially, and a fact no valuation can falsify must
not enter a derivation (the family template's hygiene remark).
\end{bremark}

\begin{anchors}
Recipe \ref{rec:relunn-instance} warrant & \texttt{sound\_reluNN}\\
family linkage extraction & \texttt{reluLinked\_facts} (see the \rulename{LiRPA} section)\\
ground-fact leaf (step 1) & \texttt{sound\_pConstExact} (the \rulename{PConstExact} section)\\
payload structure & \texttt{PReLUNNData} (two wires, linkage --- no premise field)\\
\end{anchors}


\rulesection{23}{ReLU-Above}{admissible rule}
\label{blcatrule-ReLUAbove}
\ruleblock{The rule (Steele metanotation)}
\sourcenote{as realized in the mechanized development; source row: \texttt{\texttt{rules.toml}} --- \emph{ReLU-Above}. The metanotation statement is typeset by hand for notational quality; its mechanized ground truth is the Lean statement below, extracted from the sources at every build of this book and hash-bound. The rule \emph{roster} is machine-checked (every mechanized rule tag round-trips to exactly one catalogue row); the statement \emph{text} is editorial, checked by eye against the extracted Lean beside it.}
\begin{mathpar}
\inferrule*[Right=\rulename{ReLU-Above}]{
  \ReLU_{\{h\}\to\{y\}} \in H
}{
  \Phi \vdash y \geq h
}
\end{mathpar}

\par\medskip
\rulename{ReLU-Above} is the identity minorant: the output wire of a
cited ReLU application dominates its preactivation, $\Phi \vdash y
\geq h$, with no proof premises — $\Phi$ arbitrary, the premise-free
convention. Together with \rulename{ReLU-NN} it gives the exact
pointwise lower envelope of the activation as two always-available
atoms ($y \geq 0$ and $y \geq h$); a producer needs no box knowledge
at all to emit either. Upper bounds are different in kind — no affine
upper bound on the activation exists without a box, which is why the
upper side of the family lives in \rulename{LiRPA} behind two box
premises.

What the checker recomputes: the linkage — the hypernet really
contains a packaged ReLU application at exactly these ports, of
positive total size (the family\textquotesingle s port-lookup side
condition, displayed in the \rulename{LiRPA} section) — and that the
conclusion is literally the atom $y \geq h$, elementwise over the
node\textquotesingle s tensor shape.

The linkage is the only hypothesis; without it, two arbitrary wires
are asserted ordered. Entrywise the fact is immediate from the
node\textquotesingle s two-case semantics: on the inactive branch $y =
0 \geq h$ (there $h \leq 0$), on the active branch $y = h$.

\cataloguenote{The proof section uses this rule\textquotesingle s derivation as the
family\textquotesingle s worked example of atom-orientation plumbing:
its inactive branch discharges into catalogue nodes through a chain of
a bound-node polarity flip, a \rulename{Norm} identity, an
\rulename{ArithAtom} fold, and two \rulename{EqElim} transports — a
faithful specimen of what ``one node unpacks to many\textquotesingle
\textquotesingle{} means in practice.
}
\ruleblock{Lean statement}
\par\medskip
The warrant consumes the family linkage extraction
(\texttt{reluLinked\_facts}, displayed in the \rulename{LiRPA}
section): entrywise, $y = 0 \geq h$ on the inactive case and $y = h$
on the active one. The payload structure, \texttt{PReLUAboveData}
--- quoted in full below --- carries no premise field.

\declhead{BernsteinLogic.Soundness}{sound\_reluAbove}{\badge{BadgeGreen}{statement bound}}
\begin{leansrc}
/-- ReLU-Above: `Φ ⊢ y ≥ h`, premise-free — entrywise, `y = 0 ≥ h` on
the nonpositive case and `y = h` on the nonnegative one. -/
theorem sound_reluAbove {H : Hypernet} {seq : Sequent} {preact post : Wire}
    (hlink : reluLinked H preact post)
    (hpost : seq.post = .atom .ge (.var post) (.var preact)) :
    seq.semanticallyValid H
\end{leansrc}
\declhead{BernsteinLogic.Core.Catalogue}{PReLUAboveData}{\badge{BadgeGreen}{definition bound}}
\begin{leansrc}
inductive PReLUAboveData : Hypernet → Sequent → Type where
  | mk (H : Hypernet) (seq : Sequent)
       (preact post : Wire)
       (link : reluLinked H preact post)
       (conclPost : seq.post = .atom .ge (.var post) (.var preact))
       : PReLUAboveData H seq
\end{leansrc}
\ruleblock{Correspondence (Steele display $\leftrightarrow$ Lean)}
\badge{BadgeGreen}{reviewer-certified against the displayed texts}
\sourcenote{The dictionary below is a reviewer-certified statement of how the mechanized statements above interpret the displayed rule: display metavariable $\mapsto$ Lean binder, each display hypothesis $\mapsto$ the premise or side condition carrying it, the display conclusion $\mapsto$ the goal --- with any deliberate difference in width between display and mechanization named outright. Certification is recorded as a hash triple over (the displayed cell, the extracted statement texts, this dictionary\textquotesingle s own text); if any of the three changes, this block turns amber until a reviewer re-certifies. The correspondence claim itself remains editorial --- reviewer-grade, not machine-checked; its machine-checked successor types each soundness statement against the rule as a structured object, and supersedes this block per rule as it lands.}
Judgment frame as in \rulename{Ax}; linkage premise --- including the
positive-size narrowing hidden under the displayed $\in$ --- exactly
as in the \rulename{LiRPA} dictionary, the family preamble.
\begin{itemize}
\item $h \mapsto$ \texttt{preact}; $y \mapsto$ \texttt{post};
  $\ReLU_{\{h\}\to\{y\}} \in H$ $\mapsto$ \texttt{hlink} $:$
  \texttt{reluLinked H preact post} (payload alias:
  \texttt{PReLUAboveData.link}, same proposition).
\item premise-freeness $\mapsto$ silence, the \rulename{PConstExact}
  representation: the precondition component of \texttt{seq :
  Sequent} is unconstrained by any hypothesis, so the
  statement holds under arbitrary premises $\Phi$ --- the
  display\textquotesingle s universality over the preactivation
  rides the linkage, as at
  \rulename{ReLU-NN}.
\item conclusion $\Phi \vdash y \geq h$ $\mapsto$ the goal under
  \texttt{hpost} $:$ \texttt{seq.post = .atom .ge (.var post) (.var
  preact)} (payload alias: \texttt{conclPost}, same equation) ---
  output wire left, preactivation right, as displayed;
  the only two-wire atom among the family\textquotesingle s
  premise-free conclusions, so this dictionary is where the shared
  shape of the two tensors (from the linkage extraction) is
  load-bearing rather than convenient.
\end{itemize}
Width note: certified atoms elementwise over the cited
application\textquotesingle s tensor shape (batch semantics, as at
\rulename{LiRPA}). Beyond the family\textquotesingle s linkage gate,
neither side is wider: both hypotheses of \texttt{sound\_reluAbove}
are consumed, and nothing else is stated.

\ruleblock{Lean verifier receipts}
\sourcenote{module \texttt{BernsteinLogic.Soundness}; bound build \texttt{a20d90e66c13\ldots} (bound from receipt a9c0e34d3e3c9c8c, 2026-08-03T21:30:36Z). A receipt records that named checking layers accepted a build: (module, compiled-artifact SHA-256, environment SHA-256, checker layers), minted when the checking gate ran; the store receipts attesting the bound build are listed with the layers each run credited. Nothing here asks the reader to trust the store: every layer is an independently runnable program over the quoted sources --- rebuild the module and re-run the checkers to reproduce the verdict. Receipts attesting superseded builds credit no layer and are omitted from this rendering (the count of omissions, if any, is noted below).}
\begin{receipttable}
\receiptrow{a9c0e34d3e3c9c8c}{2026-08-03T21:30:36Z}{build+decorrelated}{\badge{BadgeGreen}{current}}
\end{receipttable}
\sourcenote{21 further receipt(s) in the store attest superseded builds of this module and are omitted from this rendering.}
\sourcenote{module \texttt{BernsteinLogic.Core.Catalogue}; bound build \texttt{cbc92e976a1d\ldots} (bound from receipt 5abdb3589707b5df, 2026-08-03T21:30:33Z). A receipt records that named checking layers accepted a build: (module, compiled-artifact SHA-256, environment SHA-256, checker layers), minted when the checking gate ran; the store receipts attesting the bound build are listed with the layers each run credited. Nothing here asks the reader to trust the store: every layer is an independently runnable program over the quoted sources --- rebuild the module and re-run the checkers to reproduce the verdict. Receipts attesting superseded builds credit no layer and are omitted from this rendering (the count of omissions, if any, is noted below).}
\begin{receipttable}
\receiptrow{24804e66edba0bc9}{2026-08-03T20:25:09Z}{build+decorrelated}{\badge{BadgeGreen}{current}}
\receiptrow{5abdb3589707b5df}{2026-08-03T21:30:33Z}{build+decorrelated}{\badge{BadgeGreen}{current}}
\end{receipttable}
\sourcenote{18 further receipt(s) in the store attest superseded builds of this module and are omitted from this rendering.}
\layerline{\layeritem{build}{lake build (elaboration)}{\badge{BadgeGreen}{current}}\layeritem{replay}{leanchecker (official kernel replay)}{\badge{BadgeAmber}{stale only}}\layeritem{hardened}{safe\_verify (hardened screen)}{\badge{BadgeAmber}{stale only}}\layeritem{decorrelated}{lean4export + nanoda (decorrelated kernel)}{\badge{BadgeGreen}{current}}}
\sourcenote{\badge{BadgeGray}{comparator: none} no fixture in the three-implementation comparator harness (compiled / kernel-interpreted / reference-interpreted) exercises this rule.}
\ruleblock{HOL4 statement}
\sourcenote{theory \texttt{bern\_tr\_catdata2}, HOL4 development script \texttt{\texttt{bern\_tr\_catdata2Script.sml}} \badge{BadgeGreen}{script bound}}
\holhead{Theorem preluabovedata\_receipt}
\begin{holsrc}
Theorem preluabovedata_receipt:
  ruledataOk [] cd2Hrelu
    <| pre := FTt ; post := FAtom BowGe (TmVar 4) (TmVar 3) |>
    (RdReluAbove (PReluAboveD 3 4)) ∧
  ¬ ruledataOk [] cd2Hrelu
      <| pre := FTt ; post := FAtom BowGe (TmVar 3) (TmVar 4) |>
      (RdReluAbove (PReluAboveD 3 4))
\end{holsrc}
\ruleblock{HOL verifier receipts}
\badge{BadgeBlue}{documentary} \texttt{hol-strict} CLEAN: theory \texttt{bern\_tr\_catdata2}, 272
theorems, 0 cheat-tagged, 0 nonstandard axioms (whitelist:
\texttt{DISK\_THM}; the four \texttt{bool} axioms) --- verification
run of 2026-07-22 on an isolated build host, against a
content-addressed checkout of the exact sources cited here (git blob
\texttt{cd3c1e95\ldots}). The proof-object stratum is ported as one
mutual datatype family with a mutual wellformedness predicate,
exactly as the Lean mutual inductive block; this premise-free
rule\textquotesingle s arm states its entire checkable content ---
the family linkage plus the output-dominates-preactivation
conclusion equation --- and the linkage lookup into the
hypernet\textquotesingle s application list has a ported realization
(\texttt{reluLinkedB}, theory \texttt{bern\_tr\_catdata}, 252
theorems, 0 offences, same run). The receipt theorem cited here is
the kernel-evaluated accept/refusal pair for the arm. Honest
residue: the semantic soundness twin (the Lean
development\textquotesingle s \texttt{sound\_reluAbove}) is not yet
ported; the payload/checker layer above is.

\sourcenote{HOL receipts are documentary: quoted results of \texttt{hol-strict} runs
and census rows from the HOL4 development\textquotesingle s build
evidence, not minted receipt artefacts. A HOL receipt store (the Lean
store\textquotesingle s twin) is named future work.
}
\ruleblock{Proof}

\sourcenote{Genre: admissible rule (tier 2). The family machinery is
stated in the \rulename{LiRPA} section's family template. This rule is
premise-free: $\Phi$ is arbitrary and never consulted.}

\begin{brecipe}[the \rulename{ReLU-Above} instance --- worked plumbing]
\label{rec:reluabove-instance}
The node:
\[
\inferrule*[Right=\rulename{ReLU-Above}]{
  \ReLU_{\{h\}\to\{y\}} \in H
}{
  \Phi \vdash y \geq h
}
\qquad\rightsquigarrow\qquad
\text{case analysis on }
(h \leq 0 \wedge y = 0) \vee (0 \leq h \wedge y = h).
\]
We walk the \emph{inactive} branch in full --- the family's honest
specimen of what ``one node unpacks to many'' means in catalogue
practice. Context $\Phi_{-}$ carries $h \leq 0$ and $y = 0$; the goal
is $y \geq h$.
\begin{enumerate}
\item $\Phi_{-} \vdash h \leq 0$ --- \rulename{Ax} + $\wedge$E (the
  case disjunct enters as one ordered subtree; two $\wedge$E steps).
\item $\Phi_{-} \vdash -h \geq 0$ --- the bound node's polarity-flip
  reading (\rulename{PBound}, reading (G)): the combination is the
  single component $h \leq 0$ with weight $1$, canonical polynomial
  $c = h$, read as $e \geq 0$ with $e = -h$.
\item $\Phi_{-} \vdash (-h) = (0 - h)$ --- \rulename{Norm}: equal
  monomial canonical forms over $\Q$, same wires, shared broadcast
  shape.
\item $\Phi_{-} \vdash (0 - h) \geq 0$ --- \rulename{EqElim} along
  step 3's equality (both sides shaped alike --- the transport gate is
  met), rewriting step 2's atom.
\item $\Phi_{-} \vdash 0 \geq h$ --- \rulename{ArithAtom} (fold): from
  $(t_1 - t_2) \bowtie 0$ to $t_1 \bowtie t_2$ at $t_1 = 0$, $t_2 =
  h$, $\bowtie = {\geq}$.
\item $\Phi_{-} \vdash 0 = y$ --- \rulename{EqSym} on the context's
  $y = 0$; then $\Phi_{-} \vdash y \geq h$ --- \rulename{EqElim} along
  $0 = y$ at $\Psi = (a \geq h)$.
\end{enumerate}
The \emph{active} branch (context $0 \leq h$, $y = h$) derives $h \geq
h$ by the same pattern, spelled at the same grain: \rulename{Norm}
gives $(h - h) = 0$, whence $0 = (h - h)$ by \rulename{EqSym};
\rulename{PConstExact} supplies the ground fact $0 \geq 0$, weakened
into the branch context; \rulename{EqElim} along $0 = (h - h)$
transports it to $(h - h) \geq 0$, and \rulename{ArithAtom} folds that
to $h \geq h$ --- then \rulename{EqSym} on the context's $y = h$ and
\rulename{EqElim} along $h = y$ (wire-to-wire, one shared shape) close
the goal. Six arithmetic-plumbing nodes in the inactive branch ---
\rulename{PBound}, \rulename{Norm}, \rulename{EqElim},
\rulename{ArithAtom}, \rulename{EqSym}, \rulename{EqElim}; the
extraction's \rulename{Ax} and $\wedge$E steps aside --- to move one
minus sign across a comparator: this is why the family rules exist as
single nodes, and why the wire format prices them.
\end{brecipe}

\begin{bremark}[shape scope]
Step 6 of the inactive branch and the active branch's route through
$(h - h) = 0$ transport along equalities with a rank-$0$ constant
side, so the displayed derivation type-checks at rank-$0$ instances
(step 4's transport along $(-h) = (0 - h)$ is wire-shaped on both
sides --- the constant $0$ sits \emph{inside} the term $0 - h$ --- and
is unaffected); the shape-general warrant is the direct lemma, exactly
as \autoref{rem:relunn-shapescope} in the \rulename{ReLU-NN}
section states for the whole family.
\end{bremark}

\begin{bprop}[semantic warrant]\label{prop:reluabove-sound}
If the linkage holds, the sequent concluding $y \geq h$ is
semantically valid over any premise formula $\Phi$.
\end{bprop}
\begin{proof}
Fix an admissible valuation $\sigma$. The linkage extraction
(\texttt{reluLinked\_facts}) gives both tensors at one shared shape
and, per entry, either $p \leq 0 \wedge y = 0$ --- whence $y = 0 \geq
p$ --- or $y = p$ --- whence $y \geq p$ reflexively. The conclusion
atom compares the two tensors entrywise at their shared shape. That is
\texttt{sound\_reluAbove}, step for step.
\end{proof}

\begin{bremark}[the exact lower envelope, premise-free]
Together with \rulename{ReLU-NN} this rule states the exact pointwise
lower envelope of the activation, $y \geq \max(0, h)$, as two atoms a
producer can emit with \emph{no} box knowledge whatsoever --- exact in
that the bound is attained: $\max(0, h)$ \emph{is} the activation's
value, so no sound lower bound can improve on the pair. The
asymmetry with the upper side is structural, not an artefact: no
affine function dominates $\max(0, h)$ on all of $\Q$ (a candidate
with positive slope dives below $0$ as $h \to -\infty$; one with
slope $\leq 0$ falls below $h$ as $h \to +\infty$) --- so upper
bounds require a box and live in \rulename{LiRPA}.
\end{bremark}

\begin{anchors}
Recipe \ref{rec:reluabove-instance} warrant & \texttt{sound\_reluAbove}\\
family linkage extraction & \texttt{reluLinked\_facts} (see the \rulename{LiRPA} section)\\
polarity-flip reading (step 2) & \texttt{sound\_pBound} (the \rulename{PBound} section)\\
ground-fact leaf (active branch) & \texttt{sound\_pConstExact} (the \rulename{PConstExact} section)\\
payload structure & \texttt{PReLUAboveData} (two wires, linkage --- no premise field)\\
\end{anchors}


\rulesection{24}{Cell-In-Interior}{admissible rule}
\label{blcatrule-CellInInterior}
\ruleblock{The rule (Steele metanotation)}
\sourcenote{as realized in the mechanized development; source row: \texttt{\texttt{rules.toml}} --- \emph{Cell-In-Interior}. The metanotation statement is typeset by hand for notational quality; its mechanized ground truth is the Lean statement below, extracted from the sources at every build of this book and hash-bound. The rule \emph{roster} is machine-checked (every mechanized rule tag round-trips to exactly one catalogue row); the statement \emph{text} is editorial, checked by eye against the extracted Lean beside it.}
\begin{mathpar}
\inferrule*[Right=\rulename{Cell-In-Interior}]{
  \Phi \vdash \bigwedge_{v} \lambda^{C_i}_v(x) \bowtie^{\circ}_v 0 \ \text{(strict barycentric interior atoms, flat vertex order; } < \text{ for extended } v\text{, } > \text{ for core } v\text{)}
  \\
  \Cell[\Sigma]_{\{x\}\to\{b_1,\ldots,b_N\}} \in H
}{
  \Phi \vdash b_i = 1 \wedge \bigwedge_{j \neq i} b_j = 0
}
\end{mathpar}

\par\medskip
\rulename{Cell-In-Interior} is one of the two rules that read the cell
primitive --- the calculus\textquotesingle s only source of disjunction.
A cell application $\Cell[\Sigma]_{\{x\}\to\{b_1,\ldots,b_N\}}$ tabulates
which cell of a simplicial decomposition its input point lies in, as
one-hot indicator outputs; piecewise definitions (a piecewise-linear
activation, a case-defined dynamics) enter the hypernet through it. This
rule discharges the case a proof is actually in: from a premise placing
$x$ \emph{strictly inside} cell $C_i$ --- a conjunction of strict
barycentric inequalities in the pinned flat vertex order --- it concludes
the full one-hot vector, $b_i = 1$ and $b_j = 0$ for every $j \neq i$. A
producer reaches for it after region reasoning (the region introduction
rules, \rulename{Subdivide}) has confined the input to one
cell\textquotesingle s interior, and needs the indicator values as
arithmetic facts downstream --- typically to substitute a branch\textquotesingle s
closed form for the primitive\textquotesingle s output.

What the checker recomputes, given the payload (schema, cell index, cited
application, premise proof): that the cited cell application is present
in the hypernet with the input wire typed at shape $[\mathrm{ambDim}]$
and every output at total size one; that the premise formula is
byte-for-byte the schema-derived strict-interior build for cell $i$
(barycentric atoms recomputed from the vertex coordinates --- no formula
text trusted from the wire); and that the conclusion is the recomputed
one-hot form. Soundness then rides the model: admissibility of the cell
application forces some cell choice $k$ with weak membership and the
one-hot outputs at $k$; the strict premise pins $k = i$ through the
schema\textquotesingle s interior-disjointness well-formedness.

Each hypothesis is load-bearing. Weaken strictness to weak membership
and a boundary point of two cells admits a valuation one-hotting the
\emph{other} cell --- the conclusion fails. Drop interior-disjointness
(distinct cells of the enumeration have disjoint interiors --- which the
complex\textquotesingle s simplicial conformity does NOT imply: conformity never
sees the extension masks and permits duplicate cell entries) and the
model may one-hot an overlapping cell. Drop the output typing and the
model\textquotesingle s entry-level one-hot fact decouples from the whole-wire
equality atoms (a zero-size output satisfies every atom vacuously; a
multi-entry output can falsify the atom away from entry zero). Drop the
input-shape tie and the premise\textquotesingle s constraint on the input tensor
no longer speaks about the model\textquotesingle s evaluation point. The proof
section walks these in order.

\cataloguenote{The checkable fragment is 1-dimensional (two-vertex cell supports on
distinct coordinates) --- the same fragment as the region machinery;
everything beyond is refused, never approximated. The intended
catalogue elaboration (case analysis over the cell axiom\textquotesingle s
disjunction) awaits the context-carrying case-analysis constructor; the
shipped warrant is a direct semantic lemma, the degenerate-rewrite
shape the \rulename{PConstExact} section defines.
}
\ruleblock{Lean statement}
\declhead{BernsteinLogic.Soundness}{sound\_cellInInterior}{\badge{BadgeGreen}{statement bound}}
\begin{leansrc}
/-- Cell-In-Interior: the strict-interior premise pins the model's cell
choice to whichI (via `cellIntDisjoint`), and the sp-fact's one-hot
valuation IS the conclusion. -/
theorem sound_cellInInterior {H : Hypernet} {seq premise : Sequent}
    {schema : SimpC} {input : Wire} {outputs : List Wire}
    {whichI : Fin schema.cells.length}
    (ih : premise.semanticallyValid H)
    (happ : PrimApp.cell schema input outputs ∈ H.applications)
    (hinshape : ∃ τ, H.wireType input = some τ ∧
      τ.shape = [schema.ambDim])
    (houts : ∀ o ∈ outputs, ∃ τ, H.wireType o = some τ ∧
      Shape.totalSize τ.shape = 1)
    (hpre : seq.pre = premise.pre)
    (hint : cellStrictInteriorForm schema (.var input) whichI premise.post)
    (hpost : cellOneHotForm schema outputs whichI seq.post) :
    seq.semanticallyValid H
\end{leansrc}
\declhead{BernsteinLogic.Soundness}{cellStrictInteriorBuild\_spec}{\badge{BadgeGreen}{statement bound}}
\begin{leansrc}
/-- Extraction for the strict-interior premise: the build succeeds only
over a 1-D complex with a two-vertex cell mask on distinct coordinates,
and the formula is the pinned pair of strict barycentric atoms. -/
theorem cellStrictInteriorBuild_spec {schema : SimpC} {t : Term}
    {k : Fin schema.cells.length} {F : Formula}
    (h : cellStrictInteriorForm schema t k F) :
    schema.ambDim = 1 ∧
    ∃ (a b : Fin schema.nVertices) (V₀ V₁ : ℚ),
      (List.finRange schema.nVertices).filter
        (fun v => (schema.cells.get k).1 v) = [a, b] ∧
      (∀ j, schema.coord a j = V₀) ∧
      (∀ j, schema.coord b j = V₁) ∧
      V₀ ≠ V₁ ∧
      F = .and
        (.atom (if (schema.cells.get k).2 a then .lt else .gt)
          (.add (.const (V₁ / (V₁ - V₀))) (.mul (.const (-1 / (V₁ - V₀))) t))
          (.const 0))
        (.atom (if (schema.cells.get k).2 b then .lt else .gt)
          (.add (.const (-V₀ / (V₁ - V₀))) (.mul (.const (1 / (V₁ - V₀))) t))
          (.const 0))
\end{leansrc}
\declhead{BernsteinLogic.Soundness}{maskRegionStrict\_pair\_intro}{\badge{BadgeGreen}{statement bound}}
\begin{leansrc}
/-- Construction: rational weights of strict signs assemble a
`MaskRegionStrict` witness over a two-vertex 1-D support (the strict
twin of `maskRegion_pair_intro`). Needs `hscope : ext ⊆ mask` — the
cellsWF invariant — because the strict region demands lam < 0 at every
ext vertex, and support vertices are the only ones with nonzero lam
(the weak twin's 0 ≤ 0 escape is unavailable at strict signs). -/
theorem maskRegionStrict_pair_intro {C : SimpC}
    {mask ext : Fin C.nVertices → Bool}
    {x : Fin C.ambDim → ℚ} {a b : Fin C.nVertices} {V₀ V₁ : ℚ}
    (hfil : (List.finRange C.nVertices).filter (fun v => mask v) = [a, b])
    (hscope : ∀ v, ext v = true → mask v = true)
    (hVa : ∀ j, C.coord a j = V₀) (hVb : ∀ j, C.coord b j = V₁)
    (hne : V₀ ≠ V₁) (hdim : C.ambDim = 1) (j0 : Fin C.ambDim)
    (hw0t : ext a = true → (V₁ - x j0) / (V₁ - V₀) < 0)
    (hw0f : ext a = false → 0 < (V₁ - x j0) / (V₁ - V₀))
    (hw1t : ext b = true → (x j0 - V₀) / (V₁ - V₀) < 0)
    (hw1f : ext b = false → 0 < (x j0 - V₀) / (V₁ - V₀)) :
    MaskRegionStrict C mask ext x
\end{leansrc}
\declhead{BernsteinLogic.Soundness}{holds\_cellOneHot}{\badge{BadgeGreen}{statement bound}}
\begin{leansrc}
/-- The one-hot conclusion holds at any valuation whose output entries
realize the one-hot pattern at k (unit-size output typing carries the
entry facts to full wire reads). -/
theorem holds_cellOneHot {H : Hypernet} (σ : H.Valuation) {schema : SimpC}
    {outputs : List Wire} {k : Fin schema.cells.length} {F : Formula}
    (hbuild : cellOneHotForm schema outputs k F)
    (houts : ∀ o ∈ outputs, ∃ τ, H.wireType o = some τ ∧
      Shape.totalSize τ.shape = 1)
    (hsem : ∀ j : Fin outputs.length,
      σ.entryFlat (outputs.get j) 0
        = some (.fin (if j.val = k.val then 1 else 0))) :
    Formula.holds σ F
\end{leansrc}
\par\medskip
\textbf{The recomputed forms and the payload.}\enspace The premise and
conclusion formulas are never read off the certificate as text: the
checker rebuilds both from the schema --- the strict-interior
conjunction from the 1-D barycentric coordinate terms
(\texttt{baryLambdas1D}, offset-first affine normal forms over the
two-vertex support), the conclusion from the one-hot pattern --- and the
payload\textquotesingle s form fields assert equality with the rebuilt
formulas. A skeptical reader should confirm by eye: the builds succeed
only over a 1-dimensional complex with exactly two support vertices at
distinct coordinates (everything else reads \texttt{none} and the rule
refuses); the strict comparators are $<$ at extended vertices and $>$
at core vertices; and the payload carries the port-typing obligations
(input at $[\mathrm{ambDim}]$, outputs at total size one) explicitly.

\declhead{BernsteinLogic.Core.Catalogue}{baryLambdas1D}{\badge{BadgeGreen}{definition bound}}
\begin{leansrc}
/-- The barycentric coordinate functions of a 1-dimensional simplex, as
canonical affine Terms in `t` (pinned shape `offset + slope·t`; see the
original for the full derivation notes and D2 scope). -/
def baryLambdas1D (cx : SimpC) (dContains : Fin cx.nVertices → Bool)
    (t : Term) : Option (List (Fin cx.nVertices × Term)) :=
  if h : cx.ambDim = 1 then
    let j0 : Fin cx.ambDim := ⟨0, by omega⟩
    match (List.finRange cx.nVertices).filter (fun v => dContains v) with
    | [a, b] =>
      let V₀ := cx.coord a j0
      let V₁ := cx.coord b j0
      if V₀ = V₁ then none
      else
        let den := V₁ - V₀
        let lam := fun (c s : Rat) => Term.add (.const c) (.mul (.const s) t)
        some [(a, lam (V₁ / den) (-1 / den)), (b, lam (-V₀ / den) (1 / den))]
    | _ => none
  else none
\end{leansrc}
\declhead{BernsteinLogic.Core.Catalogue}{cellStrictInteriorBuild}{\badge{BadgeGreen}{definition bound}}
\begin{leansrc}
/-- Strict-interior conjunction (flat vertex-order convention, strict
bowties). -/
def cellStrictInteriorBuild (schema : SimpC) (t : Term)
    (k : Fin schema.cells.length) : Option Formula := do
  let c := schema.cells[k.val]'(k.isLt)
  let lams ← baryLambdas1D schema c.1 t
  match lams.map (fun (v, lt) =>
    Formula.atom (if c.2 v then .lt else .gt) lt (.const 0)) with
  | [x, y] => some (.and x y)
  | _ => none
\end{leansrc}
\declhead{BernsteinLogic.Core.Catalogue}{cellStrictInteriorForm}{\badge{BadgeGreen}{definition bound}}
\begin{leansrc}
def cellStrictInteriorForm (schema : SimpC) (t : Term)
    (k : Fin schema.cells.length) (F : Formula) : Prop :=
  cellStrictInteriorBuild schema t k = some F
\end{leansrc}
\declhead{BernsteinLogic.Core.Catalogue}{cellOneHotBuild}{\badge{BadgeGreen}{definition bound}}
\begin{leansrc}
/-- One-hot conclusion: `(b_k = 1) ∧ chainAnd [(b_j = 0) | j ≠ k]`. -/
def cellOneHotBuild (schema : SimpC) (outputs : List Wire)
    (k : Fin schema.cells.length) : Option Formula :=
  if outputs.length = schema.cells.length then
    let bk := Formula.atom .eq (.var (outputs.getD k.val 0)) (.const 1)
    let others := ((List.range outputs.length).filter (· ≠ k.val)).map
      fun j => Formula.atom .eq (.var (outputs.getD j 0)) (.const 0)
    some (.and bk (Formula.chainAnd others))
  else none
\end{leansrc}
\declhead{BernsteinLogic.Core.Catalogue}{cellOneHotForm}{\badge{BadgeGreen}{definition bound}}
\begin{leansrc}
def cellOneHotForm (schema : SimpC) (outputs : List Wire)
    (k : Fin schema.cells.length) (F : Formula) : Prop :=
  cellOneHotBuild schema outputs k = some F
\end{leansrc}
\declhead{BernsteinLogic.Core.Catalogue}{PCellInIntData}{\badge{BadgeGreen}{definition bound}}
\begin{leansrc}
inductive PCellInIntData : Hypernet → Sequent → Type where
  | mk (H : Hypernet) (seq : Sequent)
       (premise : Sequent) (premiseProof : ProofObject H premise)
       (schema : SimpC) (input : Wire) (outputs : List Wire)
       (whichI : Fin schema.cells.length)
       (appInH : PrimApp.cell schema input outputs ∈ H.applications)
       (inputShape : ∃ τ, H.wireType input = some τ ∧
                       τ.shape = [schema.ambDim])
       (outputsScalar : ∀ o ∈ outputs, ∃ τ, H.wireType o = some τ ∧
                          Shape.totalSize τ.shape = 1)
       (preEq  : seq.pre = premise.pre)
       (interiorPremise : cellStrictInteriorForm schema (.var input) whichI
                            premise.post)
       (postOneHot : cellOneHotForm schema outputs whichI seq.post)
       : PCellInIntData H seq
\end{leansrc}
\par\medskip
\textbf{The semantic side: strict regions and interior
disjointness.}\enspace \texttt{MaskRegionStrict} is the strict-sign twin
of the weak region membership the model\textquotesingle s cell fact
carries; \texttt{SimpC.CellIntDisjoint} is the well-formedness the
pinning argument consumes, stated over cell \emph{indices} (so duplicate
entries are excluded) at rational points (the model\textquotesingle s
cell inputs are rational). It is consumed as a conjunct of the cell
application\textquotesingle s model-level fact: an application over a
non-disjoint enumeration admits no valuation --- hypothesis-side
vacuity, sound but weak --- rather than an unsound conclusion. The
predicate is a proposed schema amendment realized provisionally; the
roster page records its status.

\declhead{BernsteinLogic.Catalogue}{MaskRegionStrict}{\badge{BadgeGreen}{definition bound}}
\begin{leansrc}
/-- Strict-interior membership: the same weights with STRICT signs on
every mask vertex (< 0 extended, > 0 core) — the region the
Cell-In-Interior premise certifies. -/
def MaskRegionStrict (C : SimpC) (mask ext : Fin C.nVertices → Bool)
    (x : Fin C.ambDim → ℚ) : Prop :=
  ∃ lam : Fin C.nVertices → ℝ,
    (∀ v, mask v = false → lam v = 0) ∧
    finSum C.nVertices lam = 1 ∧
    (∀ j : Fin C.ambDim,
      finSum C.nVertices (fun v => lam v * (C.coord v j : ℝ)) = (x j : ℝ)) ∧
    (∀ v, ext v = true → lam v < 0) ∧
    (∀ v, mask v = true → ext v = false → 0 < lam v)
\end{leansrc}
\declhead{BernsteinLogic.Catalogue}{SimpC.CellIntDisjoint}{\badge{BadgeGreen}{definition bound}}
\begin{leansrc}
/-- PROPOSED-B2-DISJOINT (provisional): distinct cells of a schema's
enumeration have disjoint interiors — no rational point lies strictly
inside cell i and in the CLOSED region of a different cell k. This is
what the cell rules' admissibility derivation cites
("conformity: distinct cells have disjoint interiors"), and the
schema's `conformity` field does NOT deliver it: that field constrains
simplex-mask intersections only — it never sees the (Δ,S) extension
masks, and it permits literally duplicate cell entries, either of which
makes Cell-In-Interior UNSOUND (the model may one-hot a different
overlapping cell). Stated over INDICES so duplicates are excluded, at
ℚ points because the model's Cell inputs are rational. Σ± satisfies it
(the two rays share only the origin, strictly inside neither).

Lives here (not as a `SimpC` field) because the quantification is over
REAL weights — mathlib-side of the core/mathlib module split — and is
consumed as a
conjunct of the Cell sp-fact below: a cell application over a
non-disjoint schema admits NO valuation, the same "soundness errors
are excluded from every value set" reading that rejects NaN
executions (a hypothesis-side vacuity: sound, merely weak). -/
def SimpC.CellIntDisjoint (C : SimpC) : Prop :=
  ∀ i k : Fin C.cells.length, i ≠ k →
    ∀ x : Fin C.ambDim → ℚ,
      MaskRegionStrict C (C.cells.get i).1 (C.cells.get i).2 x →
      ¬ MaskRegion C (C.cells.get k).1 (C.cells.get k).2 x
\end{leansrc}
\ruleblock{Correspondence (Steele display $\leftrightarrow$ Lean)}
\badge{BadgeGreen}{reviewer-certified against the displayed texts}
\sourcenote{The dictionary below is a reviewer-certified statement of how the mechanized statements above interpret the displayed rule: display metavariable $\mapsto$ Lean binder, each display hypothesis $\mapsto$ the premise or side condition carrying it, the display conclusion $\mapsto$ the goal --- with any deliberate difference in width between display and mechanization named outright. Certification is recorded as a hash triple over (the displayed cell, the extracted statement texts, this dictionary\textquotesingle s own text); if any of the three changes, this block turns amber until a reviewer re-certifies. The correspondence claim itself remains editorial --- reviewer-grade, not machine-checked; its machine-checked successor types each soundness statement against the rule as a structured object, and supersedes this block per rule as it lands.}
Judgment frame as in \rulename{Ax}; sub-derivation premise as at
\rulename{Weaken}.
\begin{itemize}
\item $\Sigma \mapsto$ \texttt{schema : SimpC}; $i \mapsto$
  \texttt{whichI : Fin schema.cells.length}; $x \mapsto$ the term
  \texttt{.var input} --- the region term literally \emph{is} the
  input wire, as at \rulename{BernCone}; $b_1, \ldots, b_N \mapsto$
  \texttt{outputs : List Wire}.
\item the linkage $\Cell[\Sigma]_{\{x\}\to\{b_1,\ldots,b_N\}} \in H$
  $\mapsto$ \texttt{happ}, plain application-list membership ---
  contrast the \rulename{LiRPA} family\textquotesingle s size-gated
  lookup: here the port discipline rides two explicit typing
  hypotheses instead, \texttt{hinshape} (the input wire\textquotesingle
  s type has shape \texttt{[schema.ambDim]}) and \texttt{houts}
  (every output wire singleton-sized). Neither has a displayed row;
  the payload carries both verbatim (\texttt{PCellInIntData}\textquotesingle
  s port-typing fields, displayed above); they are
  load-bearing twice: the model\textquotesingle s
  sp-fact pins entry 0 of each output while the conclusion atom
  evaluates the whole wire, and the premise atoms constrain the input
  tensor, which ties to the model\textquotesingle s rational point
  only at the declared shape. The same undisplayed typing pack as at
  \rulename{Norm}.
\item premise $\Phi \vdash \bigwedge_v \lambda^{C_i}_v(x)
  \bowtie^{\circ}_v 0$ $\mapsto$ \texttt{ih} with \texttt{hint} $:$
  \texttt{cellStrictInteriorForm schema (.var input) whichI
  premise.post} --- the premise formula is pinned as the
  \emph{build\textquotesingle s} output: the two-atom conjunction in
  flat vertex order, comparators $<$ on extended vertices and $>$ on
  core ones exactly as displayed, each atom comparing the barycentric
  entry to \texttt{.const 0}. The build succeeds only at
  \texttt{schema.ambDim = 1}, so the certified reach of the displayed
  rule is the 1-dimensional fragment, by refusal (the
  \rulename{ConvElim}/\rulename{ConvIntro} class; the note above says
  it, this triple certifies it). $\Phi$ shared via \texttt{hpre}.
\item conclusion $\Phi \vdash b_i = 1 \wedge \bigwedge_{j \neq i} b_j
  = 0$ $\mapsto$ the goal under \texttt{hpost} $:$
  \texttt{cellOneHotForm schema outputs whichI seq.post} --- pinned
  as the one-hot build\textquotesingle s output: the $b_i = 1$ atom
  first, then the ascending-$j$ chain of zeros; the displayed
  big-wedge has exactly one accepted spelling.
\item nowhere in the certificate: the interior-disjointness
  well-formedness. It is consumed from the \emph{model
  side} --- admissible valuations carry
  \texttt{SimpC.CellIntDisjoint} in the cell sp-fact (certified
  alongside), and the proof pins the model\textquotesingle s cell
  choice with it --- so the display is right to show no such premise,
  and the rule\textquotesingle s soundness is honestly conditional on
  that provisional schema conjunct (the commentary\textquotesingle s
  PROPOSED-B2-DISJOINT story, visible in the statement chain).
\end{itemize}
Width note: within the 1-dimensional fragment, and modulo the
undisplayed port-typing pack flagged above, nothing certified is
narrower or wider than the display --- every hypothesis is consumed,
and the two builds pin both formulas to single spellings.

\ruleblock{Lean verifier receipts}
\sourcenote{module \texttt{BernsteinLogic.Soundness}; bound build \texttt{a20d90e66c13\ldots} (bound from receipt a9c0e34d3e3c9c8c, 2026-08-03T21:30:36Z). A receipt records that named checking layers accepted a build: (module, compiled-artifact SHA-256, environment SHA-256, checker layers), minted when the checking gate ran; the store receipts attesting the bound build are listed with the layers each run credited. Nothing here asks the reader to trust the store: every layer is an independently runnable program over the quoted sources --- rebuild the module and re-run the checkers to reproduce the verdict. Receipts attesting superseded builds credit no layer and are omitted from this rendering (the count of omissions, if any, is noted below).}
\begin{receipttable}
\receiptrow{a9c0e34d3e3c9c8c}{2026-08-03T21:30:36Z}{build+decorrelated}{\badge{BadgeGreen}{current}}
\end{receipttable}
\sourcenote{21 further receipt(s) in the store attest superseded builds of this module and are omitted from this rendering.}
\sourcenote{module \texttt{BernsteinLogic.Core.Catalogue}; bound build \texttt{cbc92e976a1d\ldots} (bound from receipt 5abdb3589707b5df, 2026-08-03T21:30:33Z). A receipt records that named checking layers accepted a build: (module, compiled-artifact SHA-256, environment SHA-256, checker layers), minted when the checking gate ran; the store receipts attesting the bound build are listed with the layers each run credited. Nothing here asks the reader to trust the store: every layer is an independently runnable program over the quoted sources --- rebuild the module and re-run the checkers to reproduce the verdict. Receipts attesting superseded builds credit no layer and are omitted from this rendering (the count of omissions, if any, is noted below).}
\begin{receipttable}
\receiptrow{24804e66edba0bc9}{2026-08-03T20:25:09Z}{build+decorrelated}{\badge{BadgeGreen}{current}}
\receiptrow{5abdb3589707b5df}{2026-08-03T21:30:33Z}{build+decorrelated}{\badge{BadgeGreen}{current}}
\end{receipttable}
\sourcenote{18 further receipt(s) in the store attest superseded builds of this module and are omitted from this rendering.}
\sourcenote{module \texttt{BernsteinLogic.Catalogue}; bound build \texttt{c477b6015133\ldots} (bound from receipt d87d31fb92f86564, 2026-07-20T07:43:40Z). A receipt records that named checking layers accepted a build: (module, compiled-artifact SHA-256, environment SHA-256, checker layers), minted when the checking gate ran; the store receipts attesting the bound build are listed with the layers each run credited. Nothing here asks the reader to trust the store: every layer is an independently runnable program over the quoted sources --- rebuild the module and re-run the checkers to reproduce the verdict. Receipts attesting superseded builds credit no layer and are omitted from this rendering (the count of omissions, if any, is noted below).}
\begin{receipttable}
\receiptrow{d87d31fb92f86564}{2026-07-20T07:43:40Z}{build+decorrelated}{\badge{BadgeGreen}{current}}
\receiptrow{85a439e3f753d74e}{2026-07-20T20:44:30Z}{build+decorrelated}{\badge{BadgeGreen}{current}}
\receiptrow{1e680b8d5ca91dae}{2026-07-20T21:00:59Z}{build+replay+hardened+decorrelated}{\badge{BadgeGreen}{current}}
\receiptrow{c707aa2f41b0f58f}{2026-07-21T05:28:24Z}{build+replay+hardened+decorrelated}{\badge{BadgeGreen}{current}}
\end{receipttable}
\sourcenote{4 further receipt(s) in the store attest superseded builds of this module and are omitted from this rendering.}
\layerline{\layeritem{build}{lake build (elaboration)}{\badge{BadgeGreen}{current}}\layeritem{replay}{leanchecker (official kernel replay)}{\badge{BadgeGreen}{current}}\layeritem{hardened}{safe\_verify (hardened screen)}{\badge{BadgeGreen}{current}}\layeritem{decorrelated}{lean4export + nanoda (decorrelated kernel)}{\badge{BadgeGreen}{current}}}
\sourcenote{\badge{BadgeGray}{comparator: none} no fixture in the three-implementation comparator harness (compiled / kernel-interpreted / reference-interpreted) exercises this rule.}
\ruleblock{HOL4 statement}
\sourcenote{theory \texttt{bern\_tr\_catdata2}, HOL4 development script \texttt{\texttt{bern\_tr\_catdata2Script.sml}} \badge{BadgeGreen}{script bound}}
\holhead{Theorem pcellinintdata\_receipt}
\begin{holsrc}
Theorem pcellinintdata_receipt:
  ruledataOk cd2wsCell cd2Hcell
    <| pre := FAnd (FAtom BowLt (TmAdd (TmConst (&1))
                                       (TmMul (TmConst (-&1)) (TmVar 9)))
                          (TmConst 0))
                   (FAtom BowGt (TmAdd (TmConst (&0))
                                       (TmMul (TmConst (&1)) (TmVar 9)))
                          (TmConst 0)) ;
       post := FAnd (FAtom BowEq (TmVar 5) (TmConst 1))
                    (FAtom BowEq (TmVar 7) (TmConst 0)) |>
    (RdCellInInterior (PCellInIntD
       <| pre := FAnd (FAtom BowLt (TmAdd (TmConst (&1))
                                          (TmMul (TmConst (-&1))
                                                 (TmVar 9)))
                             (TmConst 0))
                      (FAtom BowGt (TmAdd (TmConst (&0))
                                          (TmMul (TmConst (&1))
                                                 (TmVar 9)))
                             (TmConst 0)) ;
          post := FAnd (FAtom BowLt (TmAdd (TmConst (&1))
                                           (TmMul (TmConst (-&1))
                                                  (TmVar 9)))
                              (TmConst 0))
                       (FAtom BowGt (TmAdd (TmConst (&0))
                                           (TmMul (TmConst (&1))
                                                  (TmVar 9)))
                              (TmConst 0)) |> (ProofObj RdAx)
       checkfnDemoSchema 9 [5; 7] 0)) ∧
  ¬ ruledataOk cd2wsCell hnEmpty
      <| pre := FTt ; post := FTt |>
      (RdCellInInterior (PCellInIntD
         <| pre := FTt ; post := FTt |> (ProofObj RdAx)
         checkfnDemoSchema 9 [5; 7] 0))
\end{holsrc}
\ruleblock{HOL verifier receipts}
\badge{BadgeBlue}{documentary} \texttt{hol-strict} CLEAN: theory \texttt{bern\_tr\_catdata2}, 272
theorems, 0 cheat-tagged, 0 nonstandard axioms (whitelist:
\texttt{DISK\_THM}; the four \texttt{bool} axioms) --- verification
run of 2026-07-22 on an isolated build host, against a
content-addressed checkout of the exact sources cited here (git blob
\texttt{cd3c1e95\ldots}). The formerly dropped cell-payload island is
ported: this rule\textquotesingle s arm of the mutual wellformedness
predicate states the application membership, the input-shape gate,
the strict-interior premise and the one-hot conclusion (the form
faces as boolean Build-equations, the characterization convention),
and the form characterizations themselves are ported in theory
\texttt{bern\_tr\_catforms}
(\texttt{cellStrictInteriorForm\_char} /
\texttt{cellOneHotForm\_char}, 23 theorems, 0 offences, same run).
The receipt theorem cited above is the arm\textquotesingle s
kernel-evaluated accept/refusal pair; the shared barycentric
ingredient (\texttt{baryCoeffs1D}) was already ported and receipted.
Honest residue: the cell soundness argument (the Lean
development\textquotesingle s \texttt{sound\_cellInInterior}) is not
yet ported; the payload/builder layer above is.

\sourcenote{HOL receipts are documentary: quoted results of \texttt{hol-strict} runs
and census rows from the HOL4 development\textquotesingle s build
evidence, not minted receipt artefacts. A HOL receipt store (the Lean
store\textquotesingle s twin) is named future work.
}
\ruleblock{Proof}

\sourcenote{Genre: admissible rule (tier 2) --- a proof-tree-rewriting
recipe. The intended elaboration is a case analysis over the cell
primitive's one-hot disjunction (the $\sp[\Cell[\Sigma]]$ axiom ---
the calculus's only source of disjunction); the
\emph{realized} warrant is the degenerate rewrite in the
\rulename{PConstExact} sense, discharging against a named verified
semantic lemma. Both are presented, and the honest reason the first is
not yet a checkable catalogue derivation is stated.}

\begin{brecipe}[the intended rewrite --- case analysis over the cell disjunction]
\label{rec:cii-intended}
A cell application $\Cell[\Sigma]_{\{x\}\to\{b_1,\ldots,b_N\}}$
contributes to the ambient premise its one-hot disjunction
\[
\sp[\Cell[\Sigma]] \;=\; \bigvee_{k=1}^{N}\Bigl( x \in \overline{C_k}
\;\wedge\; b_k = 1 \;\wedge\; \textstyle\bigwedge_{j \neq k} b_j = 0 \Bigr).
\]
The intended unpacking of one \rulename{Cell-In-Interior} node is a
case split on $k$: in branch $k = i$ the conclusion
$b_i = 1 \wedge \bigwedge_{j\neq i} b_j = 0$ is reached from the case
hypothesis by conjunction plumbing ($\wedge$E and $\wedge$I under
whatever association the axiom's serialization fixes); in a branch
$k \neq i$ the case hypothesis
$x \in \overline{C_k}$ contradicts the rule's strict-interior premise
$x \in \mathring{C_i}$ --- \emph{provided no point lies strictly
inside one cell of the enumeration and weakly inside another}, which
is exactly the schema well-formedness the soundness theorem consumes
(\autoref{rem:cii-disjoint}). Two honest gaps keep this a template
rather than a shipped derivation: (i) the disjunction sits \emph{inside}
the premise conjunction, and the calculus's disjunction elimination as
currently realized acts only on a disjunction at the root of the premise
--- the context-carrying form that reaches it is proven as a satellite
theorem but its core constructor swap is still scheduled; (ii) the
contradiction in the off-branches is semantic (disjoint interiors), not
yet a formula-level refutation in the negation-free grammar.
\end{brecipe}

\begin{brecipe}[the realized rewrite --- discharge against the semantic lemma]
\label{rec:cii-realized}
As shipped, the rule is admissible by \emph{reduction to verified
machinery} (the \rulename{PConstExact} pattern): the node
\[
\inferrule*[Right=\rulename{Cell-In-Interior}]
  {\;\mathcal{T} \rhd \Phi \vdash \textstyle\bigwedge_{v}
    \lambda^{C_i}_v(x) \bowtie^{\circ}_v 0\;}
  {\Phi \vdash b_i = 1 \wedge \textstyle\bigwedge_{j \neq i} b_j = 0}
\qquad\rightsquigarrow\qquad
\text{\upshape discharged by } \texttt{sound\_cellInInterior},
\]
where the premise is the \emph{strict} barycentric-interior conjunction
for cell $C_i$ (one strict atom per support vertex, $<$ for extended
vertices, $>$ for core vertices, in the pinned flat vertex order) and
the checker recomputes both formula shapes from the schema --- the
premise via \texttt{cellStrictInteriorBuild}, the one-hot conclusion via
\texttt{cellOneHotBuild} --- so no formula text is trusted from the
certificate.
\end{brecipe}

\begin{bprop}[semantic warrant]\label{prop:cii-sound}
Suppose the payload carries a proof of the strict-interior premise, the
cited cell application is in the hypernet, the input wire is typed at
shape $[\mathrm{ambDim}]$, every output wire at total size $1$, and the
premise and conclusion formulas match the recomputed builds. Then the
conclusion sequent is semantically valid.
\end{bprop}
\begin{proof}[Proof sketch (the lemma's actual route)]
Fix an admissible valuation $\sigma$ satisfying $\Phi$. Admissibility of
the cited cell application yields the model-level fact: the schema's
cells have disjoint interiors, the input wire holds a rational point
$x$, and there is a cell choice $k$ with $x \in
\mathcal{R}(\Delta_k, S_k)$ (weak membership) whose one-hot valuation
the outputs realize --- entry $0$ of output $j$ reads $1$ if $j = k$,
else $0$. The premise, transported through the input wire's declared
shape (the tensor-to-point tie), gives \emph{strict} barycentric signs
at $x$ for cell $i$, i.e.\ $x$ lies in the strict interior of $C_i$
(\texttt{maskRegionStrict\_pair\_intro} assembles the region witness
from the two strict atoms). If $k \neq i$, the interior-exclusion
well-formedness is contradicted: $x$ is strictly inside $C_i$ and
weakly inside $C_k$, the exact configuration
\autoref{rem:cii-disjoint}'s predicate forbids (strict-versus-weak
exclusion --- deliberately stronger than bare disjointness of open
interiors, which a strictly-interior-to-$C_i$ point on $C_k$'s
boundary would survive). So $k = i$, and the one-hot output pattern at $k$ \emph{is} the conclusion:
each conjunct is an equality between a singleton-shaped output wire and
the constant its entry holds (\texttt{holds\_cellOneHot}).
\end{proof}

\begin{bremark}[why the premise must be strict]\label{rem:cii-strict}
At a shared boundary face the one-hot semantics deliberately admits
\emph{either} adjacent cell as the model's choice: a boundary point has
the one-hot set of every cell whose closure contains it. A weak-interior
premise therefore cannot pin $k$ --- for $x$ on the face between $C_i$
and $C_k$ the valuation one-hotting $C_k$ satisfies the weak premise
and falsifies $b_i = 1$. Strictness excludes exactly the boundary, where
the pinning genuinely fails; the companion rule \rulename{Cell-Out}
covers boundary points from the other side.
\end{bremark}

\begin{bremark}[the disjointness hypothesis, and why the complex's
conformity does not deliver it]\label{rem:cii-disjoint}
The step $k = i$ uses: no rational point lies strictly inside cell $i$
and weakly inside a \emph{different} cell $k$
(\texttt{SimpC.CellIntDisjoint}, stated over cell \emph{indices} so
duplicate entries are excluded). The schema's simplicial conformity
field does not imply this: conformity constrains simplex-mask
intersections only, never sees the extension masks that carve cells out
of simplices, and permits literally duplicate cell entries --- either
defect, were the predicate not consumed, would let the model one-hot an
overlapping cell and falsify the conclusion. The predicate is consumed as a conjunct of the cell
application's model-level fact: an application over a non-disjoint
enumeration admits no valuation at all (hypothesis-side vacuity ---
sound, merely weak), the same reading that excludes NaN executions.
It is a schema amendment adopted provisionally, pending
ratification.
\end{bremark}

\begin{bremark}[the port-typing hypotheses are load-bearing twice]
The model's one-hot fact pins \emph{entry $0$} of each output wire,
while the conclusion's equality atoms evaluate the \emph{whole} wire;
the total-size-$1$ typing is what makes the two coincide (a multi-entry
output could satisfy the entry fact while its remaining entries falsify
the atom; a zero-size output would satisfy every atom vacuously).
Dually, the premise's barycentric atoms constrain the input
\emph{tensor}, and the declared shape $[\mathrm{ambDim}]$ (with
$\mathrm{ambDim} = 1$ in the checkable fragment) ties that tensor to the
model's rational point.
\end{bremark}

\begin{bremark}[scope of the checkable fragment]
\texttt{cellStrictInteriorBuild} succeeds only over a $1$-dimensional
complex with a two-vertex cell support on distinct coordinates --- the
same fragment as the region machinery elsewhere in the catalogue;
everything beyond it is refused, not approximated (fail-toward-refusal).
Higher-dimensional cell complexes await their own barycentric builders;
the rule's statement is dimension-general and only the builder gates
narrow it.
\end{bremark}

\begin{anchors}
Recipe \ref{rec:cii-realized} warrant & \texttt{sound\_cellInInterior} (Soundness)\\
premise shape extraction & \texttt{cellStrictInteriorBuild\_spec}; builder \texttt{cellStrictInteriorBuild} / \texttt{baryLambdas1D}\\
strict-region assembly & \texttt{maskRegionStrict\_pair\_intro}; \texttt{MaskRegionStrict}\\
the disjointness pin (Rem.\ \ref{rem:cii-disjoint}) & \texttt{SimpC.CellIntDisjoint}; the cell case of \texttt{PrimApp.SemOk}\\
one-hot conclusion & \texttt{cellOneHotBuild}; \texttt{holds\_cellOneHot}\\
port typing ties & \texttt{PCellInIntData} fields \texttt{inputShape}, \texttt{outputsScalar}\\
\end{anchors}


\rulesection{25}{Cell-Out}{admissible rule}
\label{blcatrule-CellOut}
\ruleblock{The rule (Steele metanotation)}
\sourcenote{as realized in the mechanized development; source row: \texttt{\texttt{rules.toml}} --- \emph{Cell-Out}. The metanotation statement is typeset by hand for notational quality; its mechanized ground truth is the Lean statement below, extracted from the sources at every build of this book and hash-bound. The rule \emph{roster} is machine-checked (every mechanized rule tag round-trips to exactly one catalogue row); the statement \emph{text} is editorial, checked by eye against the extracted Lean beside it.}
\begin{mathpar}
\inferrule*[Right=\rulename{Cell-Out}]{
  \Phi \vdash \lambda^{C_j}_w(x) \bowtie^{\circ} 0 \ \text{(a single strict separating atom at a named witness vertex } w \text{ of } C_j\text{; } > \text{ for extended } w\text{, } < \text{ for core } w\text{)}
  \\
  \Cell[\Sigma]_{\{x\}\to\{b_1,\ldots,b_N\}} \in H
}{
  \Phi \vdash b_j = 0
}
\end{mathpar}

\ruleflavor{applies even at boundary points.}
\par\medskip
\rulename{Cell-Out} is \rulename{Cell-In-Interior}\textquotesingle s
refutation companion: from a premise strictly separating the input
point from cell $C_j$, it zeroes that cell\textquotesingle s indicator
--- $\Phi \vdash b_j = 0$ --- and it does so \emph{even at boundary
points of other cells}, where the pinning rule cannot fire. The paper
states the premise as $x \notin \overline{C_j}$, a disjunction of
strict halfspace violations; the negation-free formula grammar has no
rendering of that, so the realized rule serializes a witness instead:
a single strict barycentric atom at a \emph{named witness vertex} $w$
of $C_j$, with the polarity that contradicts membership ($>$ at an
extended vertex, $<$ at a core vertex). One atom replaces the
disjunction --- no loss for producers, who know which halfspace their
point violates, and a strict gain for the checker, which does no
disjunctive reasoning. A producer uses the rule to kill off-branches
of a piecewise definition: once every indicator but one is zero,
downstream arithmetic (a weighted combination over branches, a
substitution) simplifies to the live branch.

What the checker recomputes: presence of the cited cell application
with the same port typing as the companion rule (input at
$[\mathrm{ambDim}]$, outputs at total size one); that the witness
vertex lies in cell $j$\textquotesingle s support mask (the build
refuses otherwise); that the premise is byte-for-byte the recomputed
separating atom for $(j, w)$; and that the conclusion is the
recomputed $b_j = 0$ atom.

Each hypothesis is load-bearing. The support-mask guard: an atom at a
vertex outside $C_j$\textquotesingle s support constrains a weight the
membership predicate never mentions --- no contradiction, no soundness.
Strictness of the atom: the model\textquotesingle s membership fact
delivers weak signs, and only a strict opposite sign contradicts them
(a weak atom is compatible with a boundary weight of zero, where the
model may still choose cell $j$). The port typing: same double duty as
in the companion rule --- the tensor-to-point tie on the input, and
the entry-to-whole-wire lift on the output. Notably ABSENT here: the
interior-disjointness well-formedness --- the refutation is by weight
uniqueness on the two-vertex support against the strict witness atom,
within cell $j$ alone, so this rule stands even where only the
companion needs the provisional schema amendment. The proof section
walks these in order.

\ruleblock{Lean statement}
\declhead{BernsteinLogic.Soundness}{sound\_cellOut}{\badge{BadgeGreen}{statement bound}}
\begin{leansrc}
/-- Cell-Out: the strict separating witness refutes the model choosing
cell whichJ (weight uniqueness on the two-vertex support), so the
one-hot sp-fact zeroes its indicator. -/
theorem sound_cellOut {H : Hypernet} {seq premise : Sequent}
    {schema : SimpC} {input : Wire} {outputs : List Wire}
    {whichJ : Fin schema.cells.length} {wv : Fin schema.nVertices}
    (ih : premise.semanticallyValid H)
    (happ : PrimApp.cell schema input outputs ∈ H.applications)
    (hinshape : ∃ τ, H.wireType input = some τ ∧
      τ.shape = [schema.ambDim])
    (houts : ∀ o ∈ outputs, ∃ τ, H.wireType o = some τ ∧
      Shape.totalSize τ.shape = 1)
    (hpre : seq.pre = premise.pre)
    (hwit : cellOutWitnessForm schema (.var input) whichJ wv premise.post)
    (hpost : cellOutZeroForm schema outputs whichJ seq.post) :
    seq.semanticallyValid H
\end{leansrc}
\declhead{BernsteinLogic.Soundness}{cellOutWitnessBuild\_spec}{\badge{BadgeGreen}{statement bound}}
\begin{leansrc}
/-- Extraction for the Cell-Out witness premise: additionally pins the
witness vertex to the cell's mask and identifies it with one of the two
support vertices, selecting the corresponding strict atom. -/
theorem cellOutWitnessBuild_spec {schema : SimpC} {t : Term}
    {k : Fin schema.cells.length} {w : Fin schema.nVertices} {F : Formula}
    (h : cellOutWitnessForm schema t k w F) :
    schema.ambDim = 1 ∧
    (schema.cells.get k).1 w = true ∧
    ∃ (a b : Fin schema.nVertices) (V₀ V₁ : ℚ),
      (List.finRange schema.nVertices).filter
        (fun v => (schema.cells.get k).1 v) = [a, b] ∧
      (∀ j, schema.coord a j = V₀) ∧
      (∀ j, schema.coord b j = V₁) ∧
      V₀ ≠ V₁ ∧
      ((w = a ∧ F = .atom (if (schema.cells.get k).2 a then .gt else .lt)
          (.add (.const (V₁ / (V₁ - V₀))) (.mul (.const (-1 / (V₁ - V₀))) t))
          (.const 0)) ∨
       (w = b ∧ F = .atom (if (schema.cells.get k).2 b then .gt else .lt)
          (.add (.const (-V₀ / (V₁ - V₀))) (.mul (.const (1 / (V₁ - V₀))) t))
          (.const 0)))
\end{leansrc}
\declhead{BernsteinLogic.Soundness}{maskRegion\_pair\_elim}{\badge{BadgeGreen}{statement bound}}
\begin{leansrc}
/-- Extraction: a masked region membership over a two-vertex 1-D support
pins the barycentric weights to their rational closed forms, and exports
the sign of each weight according to its extension flag. -/
theorem maskRegion_pair_elim {C : SimpC} {mask ext : Fin C.nVertices → Bool}
    {x : Fin C.ambDim → ℚ} {a b : Fin C.nVertices} {V₀ V₁ : ℚ}
    (hfil : (List.finRange C.nVertices).filter (fun v => mask v) = [a, b])
    (hVa : ∀ j, C.coord a j = V₀) (hVb : ∀ j, C.coord b j = V₁)
    (hne : V₀ ≠ V₁) (j0 : Fin C.ambDim)
    (h : MaskRegion C mask ext x) :
    (ext a = true → (V₁ - x j0) / (V₁ - V₀) ≤ 0) ∧
    (ext a = false → 0 ≤ (V₁ - x j0) / (V₁ - V₀)) ∧
    (ext b = true → (x j0 - V₀) / (V₁ - V₀) ≤ 0) ∧
    (ext b = false → 0 ≤ (x j0 - V₀) / (V₁ - V₀))
\end{leansrc}
\par\medskip
\textbf{The recomputed forms and the payload.}\enspace As for the
companion rule, no formula text is trusted from the certificate: the
witness atom is rebuilt from the schema\textquotesingle s barycentric
data at the named vertex (note the \texttt{guard} --- the build refuses
a witness vertex outside the cell\textquotesingle s support mask), the
conclusion atom from the output list, and the payload\textquotesingle s
form fields assert equality with the rebuilt formulas. The same
two-vertex 1-D fragment gates apply (\texttt{baryLambdas1D}, displayed
in the \rulename{Cell-In-Interior} section); everything beyond the
fragment is refused.

\declhead{BernsteinLogic.Core.Catalogue}{cellOutWitnessBuild}{\badge{BadgeGreen}{definition bound}}
\begin{leansrc}
/-- Single strict separating atom at witness vertex w. -/
def cellOutWitnessBuild (schema : SimpC) (t : Term)
    (k : Fin schema.cells.length) (w : Fin schema.nVertices) :
    Option Formula := do
  let c := schema.cells[k.val]'(k.isLt)
  guard (c.1 w = true)
  let lams ← baryLambdas1D schema c.1 t
  let entry ← lams.find? fun p => decide (p.1 = w)
  some (.atom (if c.2 w then .gt else .lt) entry.2 (.const 0))
\end{leansrc}
\declhead{BernsteinLogic.Core.Catalogue}{cellOutWitnessForm}{\badge{BadgeGreen}{definition bound}}
\begin{leansrc}
def cellOutWitnessForm (schema : SimpC) (t : Term)
    (k : Fin schema.cells.length) (w : Fin schema.nVertices)
    (F : Formula) : Prop :=
  cellOutWitnessBuild schema t k w = some F
\end{leansrc}
\declhead{BernsteinLogic.Core.Catalogue}{cellOutZeroBuild}{\badge{BadgeGreen}{definition bound}}
\begin{leansrc}
def cellOutZeroBuild (schema : SimpC) (outputs : List Wire)
    (k : Fin schema.cells.length) : Option Formula :=
  if outputs.length = schema.cells.length then
    some (.atom .eq (.var (outputs.getD k.val 0)) (.const 0))
  else none
\end{leansrc}
\declhead{BernsteinLogic.Core.Catalogue}{cellOutZeroForm}{\badge{BadgeGreen}{definition bound}}
\begin{leansrc}
def cellOutZeroForm (schema : SimpC) (outputs : List Wire)
    (k : Fin schema.cells.length) (F : Formula) : Prop :=
  cellOutZeroBuild schema outputs k = some F
\end{leansrc}
\declhead{BernsteinLogic.Core.Catalogue}{PCellOutData}{\badge{BadgeGreen}{definition bound}}
\begin{leansrc}
inductive PCellOutData : Hypernet → Sequent → Type where
  | mk (H : Hypernet) (seq : Sequent)
       (premise : Sequent) (premiseProof : ProofObject H premise)
       (schema : SimpC) (input : Wire) (outputs : List Wire)
       (whichJ : Fin schema.cells.length)
       (witnessVertex : Fin schema.nVertices)
       (appInH : PrimApp.cell schema input outputs ∈ H.applications)
       (inputShape : ∃ τ, H.wireType input = some τ ∧
                       τ.shape = [schema.ambDim])
       (outputsScalar : ∀ o ∈ outputs, ∃ τ, H.wireType o = some τ ∧
                          Shape.totalSize τ.shape = 1)
       (preEq  : seq.pre = premise.pre)
       (outWitness : cellOutWitnessForm schema (.var input) whichJ
                       witnessVertex premise.post)
       (postZero : cellOutZeroForm schema outputs whichJ seq.post)
       : PCellOutData H seq
\end{leansrc}
\ruleblock{Correspondence (Steele display $\leftrightarrow$ Lean)}
\badge{BadgeGreen}{reviewer-certified against the displayed texts}
\sourcenote{The dictionary below is a reviewer-certified statement of how the mechanized statements above interpret the displayed rule: display metavariable $\mapsto$ Lean binder, each display hypothesis $\mapsto$ the premise or side condition carrying it, the display conclusion $\mapsto$ the goal --- with any deliberate difference in width between display and mechanization named outright. Certification is recorded as a hash triple over (the displayed cell, the extracted statement texts, this dictionary\textquotesingle s own text); if any of the three changes, this block turns amber until a reviewer re-certifies. The correspondence claim itself remains editorial --- reviewer-grade, not machine-checked; its machine-checked successor types each soundness statement against the rule as a structured object, and supersedes this block per rule as it lands.}
Judgment frame as in \rulename{Ax}; sub-derivation premise as at
\rulename{Weaken}; binders, linkage row, and the undisplayed
port-typing pack (\texttt{happ}, \texttt{hinshape}, \texttt{houts} ---
payload fields \texttt{PCellOutData}) exactly as in the
\rulename{Cell-In-Interior} dictionary, with $j \mapsto$
\texttt{whichJ} and the witness vertex $w \mapsto$ \texttt{wv : Fin
schema.nVertices}.
\begin{itemize}
\item premise $\Phi \vdash \lambda^{C_j}_w(x) \bowtie^{\circ} 0$
  $\mapsto$ \texttt{ih} with \texttt{hwit} $:$
  \texttt{cellOutWitnessForm schema (.var input) whichJ wv
  premise.post} --- pinned as the witness build\textquotesingle s
  output: a single strict atom at the flat-order entry of the named
  vertex, comparator $>$ on an extended witness and $<$ on a core
  one, exactly the display\textquotesingle s polarity (note the
  deliberate flip against \rulename{Cell-In-Interior}\textquotesingle
  s interior comparators); the build\textquotesingle s guard demands
  the witness vertex lie in $C_j$\textquotesingle s support --- the
  display\textquotesingle s ``of $C_j$\textquotesingle\textquotesingle{}
  --- and, as at the companion, succeeds only at
  \texttt{schema.ambDim = 1}: the certified reach is the
  1-dimensional fragment, by refusal. $\Phi$ shared via
  \texttt{hpre}.
\item conclusion $\Phi \vdash b_j = 0$ $\mapsto$ the goal under
  \texttt{hpost} $:$ \texttt{cellOutZeroForm schema outputs whichJ
  seq.post} --- the single equality atom, indicator wire left,
  \texttt{.const 0} right.
\item the disjointness conjunct of the model\textquotesingle s
  sp-fact is \emph{discarded} by this rule\textquotesingle s proof
  (the statement chain consumes the one-hot fact only) --- certifying
  the commentary\textquotesingle s claim that \rulename{Cell-Out}
  stands independent of the provisional schema amendment its
  companion needs.
\end{itemize}
Width note: ``applies even at boundary points\textquotesingle\textquotesingle{}
is carried structurally --- the witness atom is strict and nothing
demands interiority of $x$; within the 1-dimensional fragment, and
modulo the undisplayed port-typing pack inherited from the companion
dictionary, neither side is wider, every hypothesis consumed.

\ruleblock{Lean verifier receipts}
\sourcenote{module \texttt{BernsteinLogic.Soundness}; bound build \texttt{a20d90e66c13\ldots} (bound from receipt a9c0e34d3e3c9c8c, 2026-08-03T21:30:36Z). A receipt records that named checking layers accepted a build: (module, compiled-artifact SHA-256, environment SHA-256, checker layers), minted when the checking gate ran; the store receipts attesting the bound build are listed with the layers each run credited. Nothing here asks the reader to trust the store: every layer is an independently runnable program over the quoted sources --- rebuild the module and re-run the checkers to reproduce the verdict. Receipts attesting superseded builds credit no layer and are omitted from this rendering (the count of omissions, if any, is noted below).}
\begin{receipttable}
\receiptrow{a9c0e34d3e3c9c8c}{2026-08-03T21:30:36Z}{build+decorrelated}{\badge{BadgeGreen}{current}}
\end{receipttable}
\sourcenote{21 further receipt(s) in the store attest superseded builds of this module and are omitted from this rendering.}
\sourcenote{module \texttt{BernsteinLogic.Core.Catalogue}; bound build \texttt{cbc92e976a1d\ldots} (bound from receipt 5abdb3589707b5df, 2026-08-03T21:30:33Z). A receipt records that named checking layers accepted a build: (module, compiled-artifact SHA-256, environment SHA-256, checker layers), minted when the checking gate ran; the store receipts attesting the bound build are listed with the layers each run credited. Nothing here asks the reader to trust the store: every layer is an independently runnable program over the quoted sources --- rebuild the module and re-run the checkers to reproduce the verdict. Receipts attesting superseded builds credit no layer and are omitted from this rendering (the count of omissions, if any, is noted below).}
\begin{receipttable}
\receiptrow{24804e66edba0bc9}{2026-08-03T20:25:09Z}{build+decorrelated}{\badge{BadgeGreen}{current}}
\receiptrow{5abdb3589707b5df}{2026-08-03T21:30:33Z}{build+decorrelated}{\badge{BadgeGreen}{current}}
\end{receipttable}
\sourcenote{18 further receipt(s) in the store attest superseded builds of this module and are omitted from this rendering.}
\layerline{\layeritem{build}{lake build (elaboration)}{\badge{BadgeGreen}{current}}\layeritem{replay}{leanchecker (official kernel replay)}{\badge{BadgeAmber}{stale only}}\layeritem{hardened}{safe\_verify (hardened screen)}{\badge{BadgeAmber}{stale only}}\layeritem{decorrelated}{lean4export + nanoda (decorrelated kernel)}{\badge{BadgeGreen}{current}}}
\sourcenote{\badge{BadgeGray}{comparator: none} no fixture in the three-implementation comparator harness (compiled / kernel-interpreted / reference-interpreted) exercises this rule.}
\ruleblock{HOL4 statement}
\sourcenote{theory \texttt{bern\_tr\_catdata2}, HOL4 development script \texttt{\texttt{bern\_tr\_catdata2Script.sml}} \badge{BadgeGreen}{script bound}}
\holhead{Definition cellOutWitnessBuild\_def}
\begin{holsrc}
Definition cellOutWitnessBuild_def:
  cellOutWitnessBuild (schema:simpc) t (k:num) (w:num) =
    if k < LENGTH schema.cells then
      (let c = EL k schema.cells in
       if getD (FST c) w F then
         (case baryLambdas1D schema (FST c) t of
          | NONE => NONE
          | SOME lams =>
            (case FILTER (λp. FST p = w) lams of
             | [] => NONE
             | (entry::rest) =>
                 SOME (FAtom (if getD (SND c) w F then BowGt else BowLt)
                             (SND entry) (TmConst 0))))
       else NONE)
    else NONE
End
\end{holsrc}
\holhead{Definition cellOutWitnessFormB\_def}
\begin{holsrc}
Definition cellOutWitnessFormB_def:
  cellOutWitnessFormB schema t k w fm ⇔
    cellOutWitnessBuild schema t k w = SOME fm
End
\end{holsrc}
\holhead{Theorem pcelloutdata\_receipt}
\begin{holsrc}
Theorem pcelloutdata_receipt:
  ruledataOk cd2wsCell cd2Hcell
    <| pre := FAtom BowGt (TmAdd (TmConst (&1))
                                 (TmMul (TmConst (-&1)) (TmVar 9)))
                    (TmConst 0) ;
       post := FAtom BowEq (TmVar 5) (TmConst 0) |>
    (RdCellOut (PCellOutD
       <| pre := FAtom BowGt (TmAdd (TmConst (&1))
                                    (TmMul (TmConst (-&1)) (TmVar 9)))
                       (TmConst 0) ;
          post := FAtom BowGt (TmAdd (TmConst (&1))
                                     (TmMul (TmConst (-&1)) (TmVar 9)))
                        (TmConst 0) |> (ProofObj RdAx)
       checkfnDemoSchema 9 [5; 7] 0 0)) ∧
  ¬ ruledataOk cd2wsCell hnEmpty
      <| pre := FTt ; post := FTt |>
      (RdCellOut (PCellOutD
         <| pre := FTt ; post := FTt |> (ProofObj RdAx)
         checkfnDemoSchema 9 [5; 7] 0 0))
\end{holsrc}
\ruleblock{HOL verifier receipts}
\badge{BadgeBlue}{documentary} \texttt{hol-strict} CLEAN: theory \texttt{bern\_tr\_catdata2}, 272
theorems, 0 cheat-tagged, 0 nonstandard axioms (whitelist:
\texttt{DISK\_THM}; the four \texttt{bool} axioms) --- verification
run of 2026-07-22 on an isolated build host, against a
content-addressed checkout of the exact sources cited here (git blob
\texttt{cd3c1e95\ldots}). The formerly dropped cell-builder island is
ported: the witness build/form pair displayed here verbatim (the Form
face as a boolean Build-equation, the characterization convention),
with kernel-evaluated receipts (\texttt{cellOutWitnessBuild\_receipt}
/ \texttt{cellOutWitnessFormB\_receipt}), the payload arm of the
mutual wellformedness predicate with its own accept/refusal receipt
(cited above), the zero-build twin in theory
\texttt{bern\_tr\_checkfn} (\texttt{cellOutZeroBuild}, 8 theorems, 0
offences, same run), and its form characterization in theory
\texttt{bern\_tr\_catforms} (\texttt{cellOutZeroForm\_char}, 23
theorems, 0 offences). The shared barycentric ingredient
(\texttt{baryCoeffs1D}) was already ported and receipted. Honest
residue: the cell soundness argument (the Lean
development\textquotesingle s \texttt{sound\_cellOut}) is not yet
ported; the payload/builder layer above is.

\sourcenote{HOL receipts are documentary: quoted results of \texttt{hol-strict} runs
and census rows from the HOL4 development\textquotesingle s build
evidence, not minted receipt artefacts. A HOL receipt store (the Lean
store\textquotesingle s twin) is named future work.
}
\ruleblock{Proof}

\sourcenote{Genre: admissible rule (tier 2) --- a proof-tree-rewriting
recipe. Same family as \rulename{Cell-In-Interior}: the intended
elaboration is a case analysis over the cell primitive's one-hot
disjunction; the realized warrant is the degenerate rewrite discharging
against a verified semantic lemma, for the same two reasons stated in
that rule's first recipe (the context-carrying case analysis is not yet
a core constructor, and the off-branch contradiction is semantic). Only
the branch-closing argument differs --- and it is \emph{lighter}: no
disjointness hypothesis is consumed.}

\begin{brecipe}[the \rulename{Cell-Out} instance --- the realized rewrite]
\label{rec:co-realized}
\[
\inferrule*[Right=\rulename{Cell-Out}]
  {\;\mathcal{T} \rhd \Phi \vdash \lambda^{C_j}_w(x) \bowtie^{\circ} 0\;}
  {\Phi \vdash b_j = 0}
\qquad\rightsquigarrow\qquad
\text{\upshape discharged by } \texttt{sound\_cellOut},
\]
where the premise is a \emph{single strict separating atom} at a named
witness vertex $w$ of cell $C_j$ ($>$ if $w$ is an extended vertex, $<$
if a core vertex --- the sign that contradicts membership), and the
checker recomputes premise and conclusion shapes from the schema
(\texttt{cellOutWitnessBuild}, \texttt{cellOutZeroBuild}). The payload
names $w$ explicitly; the build refuses any $w$ outside $C_j$'s support
mask.
\end{brecipe}

\begin{bprop}[semantic warrant]\label{prop:co-sound}
Under the payload hypotheses (premise proof; cited cell application in
the hypernet; input typed at shape $[\mathrm{ambDim}]$; outputs at total
size $1$; premise and conclusion matching the recomputed builds), the
conclusion sequent is semantically valid.
\end{bprop}
\begin{proof}[Proof sketch (the lemma's actual route)]
Fix an admissible valuation $\sigma$ satisfying $\Phi$. The cell
application's model-level fact yields the rational input point $x$, a
cell choice $k$ with $x \in \mathcal{R}(\Delta_k, S_k)$, and the
one-hot output valuation at $k$. Suppose $k = j$. On the two-vertex
support the barycentric weights of $x$ are \emph{unique} --- they are
determined affinely by the point --- so weak membership in cell $j$
pins their signs: $\leq 0$ at extended vertices, $\geq 0$ at core
vertices (the pair-elimination lemma). The premise delivers the
\emph{strict opposite} sign for the weight at the witness vertex $w$
(transported to the model's point through the input wire's declared
shape), a direct contradiction. Hence $k \neq j$, and the one-hot fact
reads entry $0$ of output $j$ as $0$; the total-size-$1$ typing lifts
that entry to the whole-wire equality $b_j = 0$.
\end{proof}

\begin{bremark}[why a witness vertex replaces the paper's set complement]
The source statement's premise is $x \notin \overline{C_j}$ --- a
\emph{disjunction} of strict halfspace violations, one per bounding
halfspace. The formula grammar is negation-free and the premise of a
rule application must be a formula the calculus can derive, so the
amendment serializes a choice of disjunct instead: \emph{which}
halfspace is strictly violated, named by \emph{which} vertex carries
the violated barycentric sign. One atom replaces the disjunction with
no loss for producers --- within the fragment the builders admit,
where the cell's closure is exactly the conjunction of its per-vertex
sign constraints, a point outside the closure violates at least one of
them, and the producer knows which --- and a strict gain for the
checker (no disjunctive reasoning at check time).
\end{bremark}

\begin{bremark}[the witness must be a support vertex of $C_j$]
The build guards $w \in \mathrm{supp}(C_j)$. Without the guard the
premise atom would constrain a barycentric coordinate that cell $j$'s
membership predicate never mentions --- the sign contradiction in
Proposition~\ref{prop:co-sound} needs the weight at $w$ to be one of
the two the membership fact actually constrains. (On the two-vertex
support this is also what makes the atom well-defined: the builder
computes $\lambda^{C_j}_w$ from the support pair's coordinates.)
\end{bremark}

\begin{bremark}[boundary points, and the division of labor with
\rulename{Cell-In-Interior}]
The rule applies even at boundary points of \emph{other} cells: the
contradiction is between a strict violated sign and weak membership of
the \emph{same} cell $j$, so no interior-disjointness hypothesis is
consumed --- a point on the shared face of $C_i$ and $C_k$ can still be
strictly separated from a third cell $C_j$, and the rule concludes
$b_j = 0$ regardless of which of the two adjacent cells the model
one-hots. This is the refutation companion of
\rulename{Cell-In-Interior},
whose pinning argument fails precisely on boundaries and needs
disjointness; between them, interior points get the full one-hot
vector, and a boundary point gets $b_j = 0$ for every cell $j$ whose
closure excludes it --- the cells actually sharing the face stay
unpinned, necessarily: the model may one-hot either, so no sound rule
could zero theirs.
\end{bremark}

\begin{bremark}[scope of the checkable fragment]
As for \rulename{Cell-In-Interior}: the builders succeed over
$1$-dimensional complexes with two-vertex cell supports on distinct
coordinates, and refuse everything else. The statement is
dimension-general; only the builder gates narrow it.
\end{bremark}

\begin{anchors}
Recipe \ref{rec:co-realized} warrant & \texttt{sound\_cellOut} (Soundness)\\
witness-atom extraction & \texttt{cellOutWitnessBuild\_spec}; builder \texttt{cellOutWitnessBuild} / \texttt{baryLambdas1D}\\
weight uniqueness on the support pair & the pair-elimination lemma cited inside \texttt{sound\_cellOut} (\texttt{maskRegion\_pair\_elim})\\
conclusion form & \texttt{cellOutZeroBuild}\\
support-mask guard & the \texttt{guard} in \texttt{cellOutWitnessBuild}; payload field \texttt{witnessVertex}\\
port typing ties & \texttt{PCellOutData} fields \texttt{inputShape}, \texttt{outputsScalar}\\
\end{anchors}


\rulesection{26}{PConstExact}{admissible rule}
\label{blcatrule-PConstExact}
\ruleblock{The rule (Steele metanotation)}
\sourcenote{as realized in the mechanized development; source row: \texttt{\texttt{rules.toml}} --- \emph{PConstExact}. The metanotation statement is typeset by hand for notational quality; its mechanized ground truth is the Lean statement below, extracted from the sources at every build of this book and hash-bound. The rule \emph{roster} is machine-checked (every mechanized rule tag round-trips to exactly one catalogue row); the statement \emph{text} is editorial, checked by eye against the extracted Lean beside it.}
\begin{mathpar}
\inferrule*[Right=\rulename{PConstExact}]{
  e_1, e_2 \text{ wire-free (closed) terms}
  \\
  \bowtie \in \{<, \leq, =, \geq, >\}
  \\
  e_1 \bowtie e_2 \ \text{holds under exact evaluation in } \Q
}{
  \Phi \vdash e_1 \bowtie e_2
}
\end{mathpar}

\ruleflavor{the third premise is decided, not assumed: both terms are closed, so the comparison evaluates in exact rational arithmetic (the checker\textquotesingle s realization is the verified closed evaluator, which returns no verdict on any term containing a wire --- closedness is thereby implied, never a trusted annotation).}
\par\medskip
\rulename{PConstExact} is the calculus\textquotesingle s ground-fact
rule: it lets a derivation state an exact closed arithmetic fact ---
$\vdash -3/4 \leq 0$, or $\vdash (2/3)^2 = 4/9$ --- with no premises.
Before it, ground facts were nearly out of reach: \rulename{Norm},
premise-free, already states closed equalities of identical canonical
forms ($\vdash (2/3)^2 = 4/9$ falls to it), but no rule stated a ground
\emph{order} comparison --- $\vdash -3/4 \leq 0$ had no derivation. A
producer uses it
wherever a numeric side condition enters a derivation --- a constant
bound, an exact constant comparison feeding a later combination step.

What the checker recomputes is the entire side condition:
\texttt{evalClosed}, the verified closed-term evaluator, exact over $\Q$,
must return \texttt{true} on the comparison. There is no separate
``closedness'' hypothesis to check: the
evaluator returns \texttt{none} the moment it meets a wire, so a
\texttt{true} result certifies closedness as a byproduct --- a
hypothesis the payload cannot mis-declare. One comparator, $\neq$, is
excluded from the admissible set even though the underlying semantic
lemma covers it: the calculus maintains a second, three-valued soundness
screen, and a premise-free $\neq$ atom would break it (the proof section
unpacks why). Without the exclusion, the rule would silently invalidate
an invariant the rest of the calculus is entitled to assume.

\cataloguenote{The rule generalizes the earlier scalar-constant-fact rule, retained as
the corner case $e_2 = 0$. Its tier is the taxonomy\textquotesingle s
recorded boundary case: admissible by \emph{reduction to verified
machinery} --- the soundness lemma is one step to the
evaluator-correctness theorem --- rather than by unpacking into a
catalogue derivation, because for its order conclusions no such
derivation exists; that
impossibility is the rule\textquotesingle s point. It is not core
because it adds no new kernel semantic content.
}
\ruleblock{Lean statement}
\declhead{BernsteinLogic.Soundness}{Formula.holds\_of\_evalClosed\_true}{\badge{BadgeGreen}{statement bound}}
\begin{leansrc}
/-- Closed formulas that evaluate true hold at every valuation. -/
theorem Formula.holds_of_evalClosed_true {H : Hypernet} (σ : H.Valuation) :
    ∀ (F : Formula), F.evalClosed = some true → F.holds σ
\end{leansrc}
\declhead{BernsteinLogic.Soundness}{sound\_pConstExact}{\badge{BadgeGreen}{statement bound}}
\begin{leansrc}
/-- PConstExact soundness: a closed exact comparison whose closed evaluator
`Formula.evalClosed` decides `some true` holds at every valuation. One step to
the verified `Formula.holds_of_evalClosed_true` — closedness of `e₁`, `e₂` is
IMPLIED by `evalTrue` succeeding (`evalClosed` is `none` on any wire), and the
comparator is decided by exact ℚ evaluation. `pre` is arbitrary (premise-free). -/
theorem sound_pConstExact {H : Hypernet} {seq : Sequent} {e₁ e₂ : Term}
    {b : BowTie}
    (hEval : (Formula.atom b e₁ e₂).evalClosed = some true)
    (hpost : seq.post = .atom b e₁ e₂) :
    seq.semanticallyValid H
\end{leansrc}
\declhead{BernsteinLogic.Core.Catalogue}{Formula.evalClosed}{\badge{BadgeGreen}{statement bound}}
\begin{leansrc}
/-- Closed-formula evaluation: `none` iff a wire or region atom occurs
(strict; see original for the closed-world reading). -/
def Formula.evalClosed : Formula → Option Bool
  | .atom b l r => do pure (b.evalQ (← l.evalQ?) (← r.evalQ?))
  | .inRegion _ _ _ => none
  | .and l r => do pure ((← l.evalClosed) && (← r.evalClosed))
  | .or l r => do pure ((← l.evalClosed) || (← r.evalClosed))
  | .tt => some true
  | .ff => some false

/-- The barycentric coordinate functions of a 1-dimensional simplex, as
canonical affine Terms in `t` (pinned shape `offset + slope·t`; see the
original for the full derivation notes and D2 scope). -/
def baryLambdas1D (cx : SimpC) (dContains : Fin cx.nVertices → Bool)
    (t : Term) : Option (List (Fin cx.nVertices × Term))
\end{leansrc}
\declhead{BernsteinLogic.Core.Catalogue}{PConstExactData}{\badge{BadgeGreen}{statement bound}}
\begin{leansrc}
inductive PConstExactData : Hypernet → Sequent → Type where
  | mk (H : Hypernet) (seq : Sequent)
       (e₁ e₂ : Term) (b : BowTie)
       (notNE : b ≠ .ne)
       (evalTrue : (Formula.atom b e₁ e₂).evalClosed = some true)
       (postEq : seq.post = .atom b e₁ e₂)
       : PConstExactData H seq
\end{leansrc}
\declhead{BernsteinLogic.Admissible}{ProofObject.m3sound}{\badge{BadgeGreen}{statement bound}}
\begin{leansrc}
/-- Monotone replacement at a conjunct path: substituting a
lattice-smaller formula at any conjunct address shrinks the M3 value.
Meet monotonicity only — no distributivity — which is why Promise,
unlike ∨E-ctx, is lattice-sound and core-eligible. -/
theorem mval_conjReplace_mono :
    ∀ {π : ConjPath} {Γ Γ' x q : Formula},
    Γ.conjGet? π = some q → Γ.conjReplace? π x = some Γ' →
    M3.le (mval x) (mval q) = true →
    M3.le (mval Γ') (mval Γ) = true := by
  intro π
  induction π with
  | here =>
    intro Γ Γ' x q hget hrep hle
    simp only [Formula.conjGet?] at hget
    simp only [Formula.conjReplace?] at hrep
    injection hget with h1
    injection hrep with h2
    rw [← h2, h1]
    exact hle
  | left p ih =>
    intro Γ Γ' x q hget hrep hle
    cases Γ with
    | and l r =>
      simp only [Formula.conjGet?] at hget
      simp only [Formula.conjReplace?] at hrep
      cases hmap : l.conjReplace? p x with
      | none => rw [hmap] at hrep; nomatch hrep
      | some l' =>
        rw [hmap] at hrep
        injection hrep with h2
        rw [← h2]
        show M3.le ((mval l').inf (mval r)) ((mval l).inf (mval r)) = true
        exact M3.inf_mono_l _ _ _ (ih hget hmap hle)
    | atom b tl tr => nomatch hrep
    | inRegion t Δ S => nomatch hrep
    | or l r => nomatch hrep
    | tt => nomatch hrep
    | ff => nomatch hrep
  | right p ih =>
    intro Γ Γ' x q hget hrep hle
    cases Γ with
    | and l r =>
      simp only [Formula.conjGet?] at hget
      simp only [Formula.conjReplace?] at hrep
      cases hmap : r.conjReplace? p x with
      | none => rw [hmap] at hrep; nomatch hrep
      | some r' =>
        rw [hmap] at hrep
        injection hrep with h2
        rw [← h2]
        show M3.le ((mval l).inf (mval r')) ((mval l).inf (mval r)) = true
        exact M3.inf_mono_r _ _ _ (ih hget hmap hle)
    | atom b tl tr => nomatch hrep
    | inRegion t Δ S => nomatch hrep
    | or l r => nomatch hrep
    | tt => nomatch hrep
    | ff => nomatch hrep

/-! ## §A9 The metatheorem: derivability is lattice-sound

Over a hypernet with no typed wires, every derivable sequent satisfies
⟦pre⟧ ≤ ⟦post⟧ in M₃. One case per rule of the catalogue; the
propositional rules are the lattice axioms, equality elimination is the
substitution invariance (its typing field forces both terms closed over
an untyped hypernet), the region/arithmetic-closure rules are handled
by their conclusion shapes or payload absurdities, and the
empty-payload rules are vacuous. -/
theorem ProofObject.m3sound : ∀ {H : Hypernet} {seq : Sequent},
    (∀ w, H.wireType w = none) → ProofObject H seq →
    M3.le (mval seq.pre) (mval seq.post) = true
  | _, _, _, .mk _ _ _ (.ax _ _ (.mk _ _ hpp)) => by
      rw [hpp]; exact M3.le_refl _
  | _, _, hE, .mk _ _ _ (.weaken _ _ (.mk _ _ _ pproof _ hpre hpost)) => by
      rw [hpre, hpost]
      exact M3.weaken_kit _ _ _ (pproof.m3sound hE)
  | _, _, hE, .mk _ _ _ (.cut _ _
      (.mk _ _ _ _ lproof _ rproof hpre hpost hcl hcr)) => by
      have hl
\end{leansrc}
\ruleblock{Correspondence (Steele display $\leftrightarrow$ Lean)}
\badge{BadgeGreen}{reviewer-certified against the displayed texts}
\sourcenote{The dictionary below is a reviewer-certified statement of how the mechanized statements above interpret the displayed rule: display metavariable $\mapsto$ Lean binder, each display hypothesis $\mapsto$ the premise or side condition carrying it, the display conclusion $\mapsto$ the goal --- with any deliberate difference in width between display and mechanization named outright. Certification is recorded as a hash triple over (the displayed cell, the extracted statement texts, this dictionary\textquotesingle s own text); if any of the three changes, this block turns amber until a reviewer re-certifies. The correspondence claim itself remains editorial --- reviewer-grade, not machine-checked; its machine-checked successor types each soundness statement against the rule as a structured object, and supersedes this block per rule as it lands.}
\begin{itemize}
\item $e_1$, $e_2$ $\mapsto$ \texttt{e₁ e₂ : Term}; $\bowtie$
  $\mapsto$ \texttt{b : BowTie}; the judgment context $\Phi$ and the
  concluded sequent $\mapsto$ \texttt{seq : Sequent} with the
  conclusion tie \texttt{hpost} $:$ \texttt{seq.post = .atom b e₁ e₂}.
  Premise-freeness is represented by silence: no hypothesis mentions
  \texttt{seq.pre}, so the statement holds under arbitrary premises.
\item the displayed premise ($e_1 \bowtie e_2$ holds under exact
  evaluation in $\Q$) $\mapsto$ \texttt{hEval} $:$
  \texttt{(Formula.atom b e₁ e₂).evalClosed = some true} --- note the
  \texttt{some}: the verified evaluator is partial and returns
  \texttt{none} on any wire, so the display\textquotesingle s separate
  ``wire-free (closed) terms\textquotesingle\textquotesingle{}
  hypothesis is \emph{absorbed}: it has no Lean binder of its own
  because a successful evaluation certifies it.
\item conclusion $\Phi \vdash e_1 \bowtie e_2$ $\mapsto$ the goal
  \texttt{seq.semanticallyValid H}, validity at every valuation of the
  ambient hypernet.
\end{itemize}
Width note, deliberate and load-bearing: the display restricts
$\bowtie \in \{<, \leq, =, \geq, >\}$, while the Lean statement is
proven for \emph{all} constructors of \texttt{BowTie}, including
$\neq$. The comparator exclusion is an admission-layer fact --- the
checker refuses a premise-free $\neq$ atom to protect the three-valued
screen, as the commentary explains --- not a semantic one. The
mechanized statement is therefore strictly wider than the displayed
rule, and the displayed rule is exactly the admitted fragment of it.

\ruleblock{Lean verifier receipts}
\sourcenote{module \texttt{BernsteinLogic.Soundness}; bound build \texttt{a20d90e66c13\ldots} (bound from receipt a9c0e34d3e3c9c8c, 2026-08-03T21:30:36Z). A receipt records that named checking layers accepted a build: (module, compiled-artifact SHA-256, environment SHA-256, checker layers), minted when the checking gate ran; the store receipts attesting the bound build are listed with the layers each run credited. Nothing here asks the reader to trust the store: every layer is an independently runnable program over the quoted sources --- rebuild the module and re-run the checkers to reproduce the verdict. Receipts attesting superseded builds credit no layer and are omitted from this rendering (the count of omissions, if any, is noted below).}
\begin{receipttable}
\receiptrow{a9c0e34d3e3c9c8c}{2026-08-03T21:30:36Z}{build+decorrelated}{\badge{BadgeGreen}{current}}
\end{receipttable}
\sourcenote{21 further receipt(s) in the store attest superseded builds of this module and are omitted from this rendering.}
\sourcenote{module \texttt{BernsteinLogic.Core.Catalogue}; bound build \texttt{cbc92e976a1d\ldots} (bound from receipt 5abdb3589707b5df, 2026-08-03T21:30:33Z). A receipt records that named checking layers accepted a build: (module, compiled-artifact SHA-256, environment SHA-256, checker layers), minted when the checking gate ran; the store receipts attesting the bound build are listed with the layers each run credited. Nothing here asks the reader to trust the store: every layer is an independently runnable program over the quoted sources --- rebuild the module and re-run the checkers to reproduce the verdict. Receipts attesting superseded builds credit no layer and are omitted from this rendering (the count of omissions, if any, is noted below).}
\begin{receipttable}
\receiptrow{24804e66edba0bc9}{2026-08-03T20:25:09Z}{build+decorrelated}{\badge{BadgeGreen}{current}}
\receiptrow{5abdb3589707b5df}{2026-08-03T21:30:33Z}{build+decorrelated}{\badge{BadgeGreen}{current}}
\end{receipttable}
\sourcenote{18 further receipt(s) in the store attest superseded builds of this module and are omitted from this rendering.}
\sourcenote{module \texttt{BernsteinLogic.Admissible}; bound build \texttt{52e5099ba522\ldots} (bound from receipt c17abf2e103fe745, 2026-07-21T05:28:22Z). A receipt records that named checking layers accepted a build: (module, compiled-artifact SHA-256, environment SHA-256, checker layers), minted when the checking gate ran; the store receipts attesting the bound build are listed with the layers each run credited. Nothing here asks the reader to trust the store: every layer is an independently runnable program over the quoted sources --- rebuild the module and re-run the checkers to reproduce the verdict. Receipts attesting superseded builds credit no layer and are omitted from this rendering (the count of omissions, if any, is noted below).}
\begin{receipttable}
\receiptrow{7f3043c54f717472}{2026-07-20T07:43:38Z}{build+decorrelated}{\badge{BadgeGreen}{current}}
\receiptrow{acaf0e8871731343}{2026-07-20T20:44:28Z}{build+decorrelated}{\badge{BadgeGreen}{current}}
\receiptrow{62974466ea131d89}{2026-07-20T21:00:57Z}{build+replay+hardened+decorrelated}{\badge{BadgeGreen}{current}}
\receiptrow{c17abf2e103fe745}{2026-07-21T05:28:22Z}{build+replay+hardened+decorrelated}{\badge{BadgeGreen}{current}}
\end{receipttable}
\sourcenote{6 further receipt(s) in the store attest superseded builds of this module and are omitted from this rendering.}
\layerline{\layeritem{build}{lake build (elaboration)}{\badge{BadgeGreen}{current}}\layeritem{replay}{leanchecker (official kernel replay)}{\badge{BadgeGreen}{current}}\layeritem{hardened}{safe\_verify (hardened screen)}{\badge{BadgeGreen}{current}}\layeritem{decorrelated}{lean4export + nanoda (decorrelated kernel)}{\badge{BadgeGreen}{current}}}
\sourcenote{\badge{BadgeGray}{comparator: none} no fixture in the three-implementation comparator harness (compiled / kernel-interpreted / reference-interpreted) exercises this rule.}
\ruleblock{HOL4 statement}
\sourcenote{theory \texttt{bern\_tr\_catdata}, HOL4 development script \texttt{\texttt{bern\_tr\_catdataScript.sml}} \badge{BadgeGreen}{script bound}}
\holhead{Definition pconstexactdataOk\_def}
\begin{holsrc}
Definition pconstexactdataOk_def:
  pconstexactdataOk (sq:sequent) (d:pconstexactdata) ⇔
    d.ceBow ≠ BowNe ∧
    formulaEvalClosed (FAtom d.ceBow d.ceL d.ceR) = SOME T ∧
    sq.post = FAtom d.ceBow d.ceL d.ceR
End
\end{holsrc}
\holhead{Theorem pconstexactdataOk\_receipt}
\begin{holsrc}
Theorem pconstexactdataOk_receipt:
  pconstexactdataOk
    <| pre := FTt ; post := FAtom BowGt (TmConst 1) (TmConst 0) |>
    <| ceBow := BowGt ; ceL := TmConst 1 ; ceR := TmConst 0 |> ∧
  ¬ pconstexactdataOk
    <| pre := FTt ; post := FAtom BowNe (TmConst 1) (TmConst 0) |>
    <| ceBow := BowNe ; ceL := TmConst 1 ; ceR := TmConst 0 |>
\end{holsrc}
\ruleblock{HOL verifier receipts}
\badge{BadgeBlue}{documentary} \texttt{hol-strict} CLEAN: theory \texttt{bern\_tr\_catdata}, 252
theorems, 0 cheat-tagged, 0 nonstandard axioms (whitelist:
\texttt{DISK\_THM}; the four \texttt{bool} axioms) --- verification
run of 2026-07-22 on an isolated build host, against a
content-addressed checkout of the exact sources cited here (git blob
\texttt{f8fceac5\ldots}). The closed-comparison payload predicate
displayed above is the rule\textquotesingle s entire checkable
content --- comparator confined to the admissible five, both terms
evaluated by the closed exact-$\Q$ evaluator, conclusion pinned to
the payload --- with a kernel-evaluated accept/refusal receipt. The
proof-object stratum consumes it verbatim as this
rule\textquotesingle s arm of the mutual wellformedness predicate
(theory \texttt{bern\_tr\_catdata2}, 272 theorems, 0 offences, same
run), and the rule-tag roster carries the rule. As on the Lean side,
the rule rides the proof-object family, not the first-order wire
format: the wire tag remains deferred on both sides (the emit path
refuses until allocated), so there is deliberately no column-side
check to mirror. Honest residue: the semantic reduction (the Lean
development\textquotesingle s \texttt{sound\_pConstExact}, one step
to the verified closed evaluator) is not yet ported.

\sourcenote{HOL receipts are documentary: quoted results of \texttt{hol-strict} runs
and census rows from the HOL4 development\textquotesingle s build
evidence, not minted receipt artefacts. A HOL receipt store (the Lean
store\textquotesingle s twin) is named future work.
}
\ruleblock{Proof}

\sourcenote{Genre: admissible rule (tier 2) --- a proof-tree-rewriting
recipe. The tier taxonomy defines tier 2 as \emph{verified unpacking
into core rules} (one node $\to$ many); this rule is the taxonomy's
recorded boundary case --- admissible by \emph{reduction to verified
machinery}. The general template, both species included, is stated in
the \rulename{Norm} section; the recipe below typesets the honest
degenerate instance.}

\begin{brecipe}[the \rulename{PConstExact} instance --- the degenerate rewrite]
\label{rec:pce-instance}
\rulename{PConstExact} has \emph{no proof premises} ($n = 0$), and its
rewrite target contains \emph{zero catalogue nodes}: no catalogue
derivation states a ground \emph{order} comparison --- closed equalities
of identical canonical forms fall to \rulename{Norm} already, but
$\vdash -3/4 \leq 0$ had no derivation before this rule --- which is the
rule's entire point. The rewrite therefore discharges the node against a
named \emph{verified theorem} rather than against catalogue structure:
\[
\inferrule*[Right=\rulename{PConstExact}]{\;e_1 \bowtie e_2 \ \text{holds under exact evaluation in } \Q\;}{\Phi \vdash e_1 \bowtie e_2}
\qquad\rightsquigarrow\qquad
\text{\upshape the checker replays } \mathrm{evalClosed};\;
\text{\texttt{holds\_of\_evalClosed\_true} stands as warrant.}
\]
The side condition is \emph{computational} (exact $\Q$ evaluation of a
closed comparison --- decidable, replayed by the checker), and the warrant
theorem is kernel-verified once, globally: soundness of the rule is
literally ``the exact evaluator is correct.''
\end{brecipe}

\begin{bprop}[semantic warrant]\label{prop:pce-sound}
If $(\mathrm{atom}\ \bowtie\ e_1\ e_2).\mathrm{evalClosed} =
\mathsf{some\ true}$ and the sequent's conclusion is that atom, the
sequent is semantically valid over any hypernet, with arbitrary premise
$\Phi$.
\end{bprop}
\begin{proof}
Fix any valuation $\sigma$ (admissibility and $\sigma \models \Phi$ are
not consulted --- the rule is premise-free, the \rulename{Norm}
convention). By the verified evaluator theorem
\texttt{Formula.holds\_of\_evalClosed\_true}, a formula whose closed
evaluation returns $\mathsf{true}$ holds at \emph{every} valuation: a
closed term's evaluation consults no wire, so $\sigma$ is never read, and
the comparison was decided by exact arithmetic over $\Q$. That is the
whole proof --- \texttt{sound\_pConstExact} is one \texttt{exact} step.
\end{proof}

\begin{bremark}[closedness is implied, not declared]
There is no ``$e_1, e_2$ closed'' payload field to check:
$\mathrm{evalClosed}$ returns $\mathsf{none}$ on any wire occurrence, so
$\mathsf{some\ true}$ \emph{certifies} closedness as a byproduct. A
well-formedness hypothesis that the evaluator itself enforces is one the
payload cannot mis-declare --- unrepresentable, not discouraged.
\end{bremark}

\begin{bremark}[why $\neq$ is excluded --- the $M_3$ countermodel]
The comparator set is confined to $\{<, \leq, =, \geq, >\}$: five of the
six. The semantic lemma covers all six, but the calculus also maintains a
three-valued screen (\texttt{ProofObject.m3sound}), whose disequality
clause reads designated free-wire membership --- so a wire-free $\neq$
atom, its wire list empty, interprets to $\bot$ under the $M_3$
countermodel. Admitting
$\neq$ against arbitrary $\Phi$ would then let a premise-free rule
conclude screen value $\bot$ under a premise valued $\top$ ($\vdash 1
\neq 2$ is true over $\Q$, and breaks the screen's
premise-$\leq$-conclusion invariant), breaking \texttt{m3sound}. The
\texttt{notNE} payload field is shared with \rulename{ArithAtom}, where
the same gate carries different weight: both of that rule's spellings
mention the same wires, so its screen value is preserved and the gate
serves alignment (that section's closing remark); here the breakage is
real. A
hypothesis whose necessity comes not from $\Q$-semantics but from a
\emph{second} soundness invariant is exactly the kind a recipe must state,
or the next porter re-learns it the hard way.
\end{bremark}

\begin{bremark}[the subsumed corner, and coverage]
\rulename{PRatFact} (the scalar constant fact $c \bowtie 0$) is the corner
$e_2 = 0$ of this rule and was subsumed when the roster was consolidated. Coverage today:
the \texttt{Term} grammar makes every closed term rank-0, so the rule
states every closed \emph{scalar} exact comparison; the matrix corner
(exact closed matrix identities, MIX's use case) needs closed-tensor
constructors and a \texttt{QTensor}-valued closed evaluator --- the named
next leg, deliberately not hand-faked here.
\end{bremark}

\begin{anchors}
Recipe \ref{rec:pce-instance} warrant & \texttt{sound\_pConstExact} (one step to \texttt{Formula.holds\_of\_evalClosed\_true})\\
closedness-is-implied & \texttt{Formula.evalClosed} ($\mathsf{none}$ on any wire)\\
the $\neq$ exclusion & \texttt{notNE} payload field; \texttt{ProofObject.m3sound}\\
premise-freedom ($\Phi$ arbitrary) & the \rulename{Norm} convention; \texttt{pre} unused in \texttt{sound\_pConstExact}\\
wire placement & \texttt{RuleTag.pConstExact}; wire tag not yet allocated --- the emit path refuses the rule until allocation\\
\end{anchors}


\rulesection{27}{PBound}{alias for \textcolor{DiffBlue}{\rulename{Syzygy}}}
\label{blcatrule-PBound}
\ruleblock{The rule (Steele metanotation)}
\sourcenote{as realized in the mechanized development; source row: \texttt{\texttt{rules.toml}} --- \emph{PBound}. The metanotation statement is typeset by hand for notational quality; its mechanized ground truth is the Lean statement below, extracted from the sources at every build of this book and hash-bound. The rule \emph{roster} is machine-checked (every mechanized rule tag round-trips to exactly one catalogue row); the statement \emph{text} is editorial, checked by eye against the extracted Lean beside it.}
\begin{mathpar}
\inferrule*[Right=\rulename{PBound}]{
  \bigl\{\Phi_i \vdash e_i \leq 0\bigr\}_{i=1}^{n}, \quad \Phi \supseteq \Phi_i \ \text{conjunct-wise (checked)}
  \\
  \bigl\{g_k = 0\bigr\}_k \ \text{equality sub-conjuncts of } \Phi
  \\
  \lambda_i \in \Q_{\geq 0}
  \\
  \nu_k \in \Q
  \\
  c \equiv \textstyle\sum_i \lambda_i e_i + \sum_k \nu_k g_k \ \text{(exact canonical polynomial)}
  \\
  \bshape_H(c) \neq \bot, \quad \wires(c) \ \text{all finitely typed in } H
  \\
  \wires(\Phi) \cup \wires(c) \cup \textstyle\bigcup_i \wires(e_i) \ \text{all typed at positive-size shapes in } H
  \\
  \Psi \in \bigl\{\, c \leq 0, \;\; e \geq 0 \text{ where } c = -e, \;\; l \leq r \text{ where } c = l - r \,\bigr\}
}{
  \Phi \vdash \Psi
}
\end{mathpar}

\par\medskip
\rulename{PBound} is the workhorse bound-node: a compact spelling of a
\rulename{Syzygy} application whose entries the register prices as
affine, with the
conclusion exposed through the three-way reading selector --- the raw
zero-form $c \leq 0$, the sign-flipped $e \geq 0$ where $c = -e$, or
the literal atom $l \leq r$ where $c = l - r$. It is the
taxonomy\textquotesingle s canonical alias: the payload is
\rulename{Syzygy}\textquotesingle s field-for-field, the unpack is the
identity transport, and the wire format prices the affine register
compactly. A producer emits it constantly: it is how an affine
consequence of proven bounds becomes a first-class node whose
conclusion can feed the rules that demand a particular atom shape ---
the $\geq$-reading makes a nonnegativity premise reachable from bound
facts without a detour, and the literal-atom reading feeds
interval-shaped consumers with no separate re-arrangement step.

What the checker recomputes is exactly
\rulename{Syzygy}\textquotesingle s obligation set --- entry decoding
against the pinned combination encoding, nonnegative weights,
entry-sourcing from component conclusions and premise equality
sub-conjuncts (both polarities on equalities), and the exact
canonical-polynomial fold identity --- plus the same bookkeeping
(conjunct-wise premise containment, conclusion shape and finiteness,
positive-size wires on premise, conclusion, and component conclusions).
That section\textquotesingle s commentary on what breaks without each
hypothesis applies verbatim here, through the identity transport.

The affine restriction is the
rule\textquotesingle s published register (the roster note\textquotesingle s
pricing vocabulary --- wire-format compactness
for coefficient rows); it appears in neither the displayed premises nor
the statement, and it is not a checked side condition: the replay
is degree-blind and soundness never consults linearity. The three
readings are one exact $\Q$ identity re-presented under a sign flip;
their semantic transport is proven once, in the target
rule\textquotesingle s section, and inherited here as a pure
re-spelling.

\cataloguenote{The wire tag for this rule is deferred to a payload-window allocation;
the emit path refuses the rule until the tag is allocated. The
compact-register pricing is the reason the alias exists as a distinct
spelling at all --- semantically it adds nothing to its target, which
is the alias genre working as intended.
}
\ruleblock{Lean statement}
\declhead{BernsteinLogic.Soundness}{holds\_of\_syzygyConclusionForm}{\badge{BadgeGreen}{statement bound}}
\begin{leansrc}
/-- The ge-bridge ring lemma (the both-polarity core change): every
reading admitted by
`syzygyConclusionForm conclusion` holds whenever the zero-normal atom
`conclusion ≤ 0` holds — the readings are the same exact ℚ identity
re-presented under a sign flip. The zero-form is the identity transport,
the ge-polarity flip is `-e ≤ 0 → e ≥ 0`, and the literal atom is
`l - r ≤ 0 → l ≤ r` (ArithAtom's tensor-general sub-zero bridge).
Applied once inside `sound_syzygy`, so PSyzygy (tier-1), PBound's
compact spelling (tier-3), and any future ≤-0-shaped fact rule inherit
all readings with no separate polarity justification. -/
theorem holds_of_syzygyConclusionForm {H : Hypernet} {σ : H.Valuation}
    {conclusion : Term} {F : Formula}
    (hle : Formula.holds σ (.atom .le conclusion (.const 0)))
    (hform : syzygyConclusionForm conclusion F) :
    Formula.holds σ F
\end{leansrc}
\declhead{BernsteinLogic.Soundness}{sound\_syzygy}{\badge{BadgeGreen}{statement bound}}
\begin{leansrc}
/-- Soundness of Syzygy (= Lemma B, dialect D): valid components give a
valid composite. The argument is `sound_farkas`'s, with the per-entry
source discharge generalized: a component-sourced entry's atom holds by
component validity (through `premiseContains`), a wiring-sourced entry's
atom holds as a premise sub-conjunct; either way
`atomLeZeroPolys_sound_entryC` delivers per-entry nonpositivity at every
clamped per-entry valuation, and the nonnegative fold closes. The
zero-normal bound then reaches `seq.post` through whichever
`syzygyConclusionForm` reading the payload declared (both-polarity core
change, 2026-07-18). -/
theorem sound_syzygy {H : Hypernet} (st : SyzygyStep H)
    (hcomps : ∀ c ∈ st.components, c.semanticallyValid H) :
    st.seq.semanticallyValid H
\end{leansrc}
\declhead{BernsteinLogic.Soundness}{sound\_pBound}{\badge{BadgeGreen}{statement bound}}
\begin{leansrc}
/-- PBound soundness: INHERITED from `sound_syzygy`, literally — the payload's
fields are `PSyzygyData`'s verbatim (tier-3 compact spelling, 1 node → 1 node),
so the proof is the same `SyzygyStep` handed to the same theorem. Polarity is
not justified here: since the both-polarity core change,
`sound_syzygy` itself carries every
`syzygyConclusionForm` reading through `holds_of_syzygyConclusionForm`, and
PBound inherits them as a pure re-spelling. -/
theorem sound_pBound {H : Hypernet} {seq : Sequent} {components : List Sequent}
    {combination conclusion : Term}
    (hcomps : ∀ c ∈ components, c.semanticallyValid H)
    (hpreC : ∀ c ∈ components, premiseContains seq.pre c.pre)
    (hshape : ∃ s, conclusion.bshape H = some s)
    (hfin : conclusion.finWires H)
    (hwpos : ∀ w ∈ seq.pre.freeWires ++ conclusion.freeWires,
      ∃ τ, H.wireType w = some τ ∧ 0 < Shape.totalSize τ.shape)
    (hcompPos : ∀ c ∈ components, ∀ w ∈ c.post.freeWires,
      ∃ τ, H.wireType w = some τ ∧ 0 < Shape.totalSize τ.shape)
    (hvalid : syzygyCombinationValid combination components seq.pre conclusion)
    (hform : syzygyConclusionForm conclusion seq.post) :
    seq.semanticallyValid H
\end{leansrc}
\declhead{BernsteinLogic.Core.Catalogue}{PBoundData}{\badge{BadgeGreen}{statement bound}}
\begin{leansrc}
inductive PBoundData : Hypernet → Sequent → Type where
  | mk (H : Hypernet) (seq : Sequent)
       (components : List Sequent)
       (componentProofs : ProofSeqList H components)
       (combination : Term)
       (conclusion : Term)
       (denBoundLog2 : Nat)
       (preContains : ∀ c ∈ components, premiseContains seq.pre c.pre)
       (conclShape : ∃ s, conclusion.bshape H = some s)
       (conclFin   : conclusion.finWires H)
       (wiresPos : ∀ w ∈ seq.pre.freeWires ++ conclusion.freeWires,
                     ∃ τ, H.wireType w = some τ ∧
                       0 < Shape.totalSize τ.shape)
       (compPostWiresPos : ∀ c ∈ components, ∀ w ∈ c.post.freeWires,
                     ∃ τ, H.wireType w = some τ ∧
                       0 < Shape.totalSize τ.shape)
       (valid : syzygyCombinationValid combination components seq.pre conclusion)
       (postForm : syzygyConclusionForm conclusion seq.post)
       : PBoundData H seq
\end{leansrc}
\declhead{BernsteinLogic.Core.Catalogue}{syzygyConclusionForm}{\badge{BadgeGreen}{definition bound}}
\begin{leansrc}
def syzygyConclusionForm (conclusion : Term) (F : Formula) : Prop :=
  F = .atom .le conclusion (.const 0)
  ∨ (∃ e, conclusion = .neg e ∧ F = .atom .ge e (.const 0))
  ∨ (∃ l r, conclusion = .sub l r ∧ F = .atom .le l r)
\end{leansrc}
\declhead{BernsteinLogic.Core.Catalogue}{syzygyCombinationValid}{\badge{BadgeGreen}{definition bound}}
\begin{leansrc}
def syzygyCombinationValid (cc : Term) (components : List Sequent)
    (premise : Formula) (conclusion : Term) : Prop :=
  syzygyCheck cc components premise conclusion = true
\end{leansrc}
\ruleblock{Correspondence (Steele display $\leftrightarrow$ Lean)}
\badge{BadgeGreen}{reviewer-certified against the displayed texts}
\sourcenote{The dictionary below is a reviewer-certified statement of how the mechanized statements above interpret the displayed rule: display metavariable $\mapsto$ Lean binder, each display hypothesis $\mapsto$ the premise or side condition carrying it, the display conclusion $\mapsto$ the goal --- with any deliberate difference in width between display and mechanization named outright. Certification is recorded as a hash triple over (the displayed cell, the extracted statement texts, this dictionary\textquotesingle s own text); if any of the three changes, this block turns amber until a reviewer re-certifies. The correspondence claim itself remains editorial --- reviewer-grade, not machine-checked; its machine-checked successor types each soundness statement against the rule as a structured object, and supersedes this block per rule as it lands.}
Verbatim the \rulename{Syzygy} dictionary, unbundled: where
\rulename{Syzygy}\textquotesingle s statement packs its obligations
into a \texttt{SyzygyStep}, \texttt{sound\_pBound} takes the same
eleven fields as flat binders --- \texttt{hpreC}, \texttt{hshape},
\texttt{hfin}, \texttt{hwpos}, \texttt{hcompPos}, \texttt{hvalid},
\texttt{hform}, over the same \texttt{components}/\texttt{combination}/
\texttt{conclusion} data with \texttt{hcomps} the component
validities --- and its entire proof is one application of
\texttt{sound\_syzygy} to the repacked Step, visible in the displayed
statement. The alias relation is thereby itself certified: 1 $\to$ 1
at statement level, and (since the both-polarity core change) 1 $\to$
1 at reading level too --- the displayed $\Psi$ selector is the same
\texttt{syzygyConclusionForm} the target consumes. This cell displays
the evaluability and positive-size rows inline; their mapping is the
\rulename{Syzygy} dictionary\textquotesingle s, unchanged. No hypothesis differs, none
is unused, and neither side is wider --- the affine restriction the
roster note mentions is wire-format compactness, absent from both the
display\textquotesingle s premises and the statement.

\ruleblock{Lean verifier receipts}
\sourcenote{module \texttt{BernsteinLogic.Soundness}; bound build \texttt{a20d90e66c13\ldots} (bound from receipt a9c0e34d3e3c9c8c, 2026-08-03T21:30:36Z). A receipt records that named checking layers accepted a build: (module, compiled-artifact SHA-256, environment SHA-256, checker layers), minted when the checking gate ran; the store receipts attesting the bound build are listed with the layers each run credited. Nothing here asks the reader to trust the store: every layer is an independently runnable program over the quoted sources --- rebuild the module and re-run the checkers to reproduce the verdict. Receipts attesting superseded builds credit no layer and are omitted from this rendering (the count of omissions, if any, is noted below).}
\begin{receipttable}
\receiptrow{a9c0e34d3e3c9c8c}{2026-08-03T21:30:36Z}{build+decorrelated}{\badge{BadgeGreen}{current}}
\end{receipttable}
\sourcenote{21 further receipt(s) in the store attest superseded builds of this module and are omitted from this rendering.}
\sourcenote{module \texttt{BernsteinLogic.Core.Catalogue}; bound build \texttt{cbc92e976a1d\ldots} (bound from receipt 5abdb3589707b5df, 2026-08-03T21:30:33Z). A receipt records that named checking layers accepted a build: (module, compiled-artifact SHA-256, environment SHA-256, checker layers), minted when the checking gate ran; the store receipts attesting the bound build are listed with the layers each run credited. Nothing here asks the reader to trust the store: every layer is an independently runnable program over the quoted sources --- rebuild the module and re-run the checkers to reproduce the verdict. Receipts attesting superseded builds credit no layer and are omitted from this rendering (the count of omissions, if any, is noted below).}
\begin{receipttable}
\receiptrow{24804e66edba0bc9}{2026-08-03T20:25:09Z}{build+decorrelated}{\badge{BadgeGreen}{current}}
\receiptrow{5abdb3589707b5df}{2026-08-03T21:30:33Z}{build+decorrelated}{\badge{BadgeGreen}{current}}
\end{receipttable}
\sourcenote{18 further receipt(s) in the store attest superseded builds of this module and are omitted from this rendering.}
\layerline{\layeritem{build}{lake build (elaboration)}{\badge{BadgeGreen}{current}}\layeritem{replay}{leanchecker (official kernel replay)}{\badge{BadgeAmber}{stale only}}\layeritem{hardened}{safe\_verify (hardened screen)}{\badge{BadgeAmber}{stale only}}\layeritem{decorrelated}{lean4export + nanoda (decorrelated kernel)}{\badge{BadgeGreen}{current}}}
\sourcenote{\badge{BadgeGray}{comparator: none} no fixture in the three-implementation comparator harness (compiled / kernel-interpreted / reference-interpreted) exercises this rule.}
\ruleblock{HOL4 statement}
\sourcenote{theory \texttt{bern\_catalogue}, HOL4 development script \texttt{\texttt{bern\_catalogueScript.sml}} \badge{BadgeGreen}{script bound}}
\holhead{Definition syzygyCheck\_def}
\begin{holsrc}
Definition syzygyCheck_def:
  syzygyCheck cc components premise conclusion ⇔
    syzygyCheckCore cc components premise conclusion ∨
    syzygyCheckW cc components premise conclusion
End
\end{holsrc}
\ruleblock{HOL verifier receipts}
\badge{BadgeBlue}{documentary} \texttt{hol-strict} CLEAN over the ported checker chain: theory
\texttt{bern\_catalogue}, 472 theorems, 0 cheat-tagged, 0 nonstandard
axioms (whitelist: \texttt{DISK\_THM}; the four \texttt{bool} axioms) ---
verification run of 2026-07-22 on an isolated build host, scripts
sha-matched verbatim to the sources cited here (git blob
\texttt{e41b9ed0\ldots}). Scope, stated honestly: the
alias\textquotesingle s entire combination check \emph{is} the target
rule\textquotesingle s, so the ported \texttt{syzygyCheck} (with its
accept and two reject receipt rows, quoted in the target
rule\textquotesingle s section) covers this rule\textquotesingle s
computational content; the conclusion-form reading selector and the
payload layer ride the dependently-typed proof-object surface, which is
a separately tracked porting stage. As on the \rulename{Farkas} page:
the Lean check has since grown a third (mixed tagged-union) branch,
added 2026-07-19 after this port run --- the divergence is
refusal-direction only, and its port is named residue.
Semantic stratum, also ported: the soundness twin \texttt{sound\_pBound}
(theory \texttt{bern\_tr\_soundsyz}: \texttt{hol-strict} clean, 7
theorems, 0 offences, verification run of 2026-07-22 on an isolated
build host, content-addressed at the same sources); the certified-
fragment keystone (\texttt{checkPf\_sound\_frag}, theory
\texttt{bern\_tr\_splice}, 20 theorems, 0 offences, same run) consumes
it at the \rulename{Syzygy} wire arm, this rule\textquotesingle s target
relation.

\sourcenote{HOL receipts are documentary: quoted results of \texttt{hol-strict} runs
and census rows from the HOL4 development\textquotesingle s build
evidence, not minted receipt artefacts. A HOL receipt store (the Lean
store\textquotesingle s twin) is named future work.
}
\ruleblock{Proof}

\sourcenote{Genre: alias rule (tier 3) --- a respelling template. An alias
owes the book five artefacts: the verified 1$\to$1 unpack, the filled-in
application of the target rule, the payload-transport table, the exact
algebraic identity that makes the respelling correct, and the soundness
corollary through the target rule's lemma. All five follow --- and this
alias is the tier taxonomy's \emph{canonical} case, the one the
tier was named for: the unpack is the identity.}

\begin{brecipe}[the alias template]\label{rec:pb-template}
An alias rule $A$ is warranted by a verified function
$\mathrm{unpack} : \mathrm{Step}_A \to \mathrm{Step}_T$ into the payload
of a \emph{single} target node $T$ (core or admissible), such that the
$A$-node and its image conclude the same sequent from the same premises:
one node $\to$ one node, payload transported, nothing derived. Soundness
is then a corollary, $\mathrm{sound}_A = \mathrm{sound}_T \circ
\mathrm{unpack}$. The wire format may carry $A$ as a compact register of
its own or never carry it at all (producers elaborating before
emission, as the \rulename{SmallGain} section records); either way the
checker's obligations are the target's, reached through the unpack.
This rule takes the second branch in v0: no producer emits the
once-contemplated affine coefficient-row register, so no tag is
allocated --- and the identity unpack makes a register pure wire
compactness, purchasable later without semantic motion.
\end{brecipe}

\begin{brecipe}[the verified unpack --- the identity transport]
\label{rec:pb-unpack}
\rulename{PBound}'s payload carries components, combination term,
conclusion term, and the bookkeeping fields \emph{verbatim} as the
target's payload does, including the conclusion-form selector. The
unpack is therefore not a computation but a re-labelling: the verified
object is the direct specialization
\[
\mathrm{sound}_{\rulename{PBound}}
\;=\;
\mathrm{sound}_{\rulename{Syzygy}}
\ \text{\upshape applied to the same fields, re-assembled as the
target's step}
\]
--- in the sources, the body of \texttt{sound\_pBound} \emph{is} one
application of \texttt{sound\_syzygy} to the step built from
\rulename{PBound}'s own fields. Where \rulename{SmallGain}'s unpack
computes a weighting, this one computes nothing; that is the tier-3
boundary case made flesh, and it is why the alias relation here is
$1 \to 1$ at the level of \emph{readings} too
(Recipe~\ref{rec:pb-transport}).
\end{brecipe}

\begin{brecipe}[the filled-in application]\label{rec:pb-filled}
Under the identity transport the \rulename{PBound} node \emph{is} the
following application of the target rule (compare the rule display in
the \rulename{Syzygy} section, field for field):
\begin{mathpar}
\inferrule*[Right=\rulename{Syzygy}]{
  \bigl\{\Phi_i \vdash e_i \leq 0\bigr\}_{i=1}^{n}, \quad
  \Phi \supseteq \Phi_i \ \text{conjunct-wise (checked)} \\
  \bigl\{g_k = 0\bigr\}_k \ \text{sub-conjuncts of } \Phi
  \ \text{in wire variables} \\\\
  \lambda_i \in \Q_{\geq 0} \\
  \nu_k \in \Q \\
  c \;\equiv\; \textstyle\sum_i \lambda_i e_i + \sum_k \nu_k g_k
  \ \text{(exact canonical polynomial)} \\\\
  \Psi \in \bigl\{\, c \leq 0, \;\; e \geq 0 \text{ where } c = -e,
  \;\; l \leq r \text{ where } c = l - r \,\bigr\}
}{
  \Phi \vdash \Psi
}
\end{mathpar}
where $\Psi$ is whichever reading of $c \leq 0$ the shared selector
declares: \emph{(Z$\leq$)} $c \leq 0$ itself; \emph{(G)} $e \geq 0$
where $c = -e$; \emph{(L$\leq$)} $l \leq r$ where $c = l - r$.
Discharging the target's obligations at this instance is immediate ---
they are \emph{the same obligations on the same data}: entry decoding,
nonnegative weights, entry-sourcing (component conclusions and premise
equality sub-conjuncts, both polarities on equalities), the exact fold
identity, premise containment, conclusion shape/finiteness, and the
positive-size wire conditions on premise, conclusion, and component
conclusions --- the last three families displayed inline in
\rulename{PBound}'s own cell, while in \rulename{Syzygy}'s cell the
shape/finiteness and positive-size conditions ride the payload (its
display carries containment as a trailing line); the obligations are
identical either way, only the display convention differs. \rulename{PBound}'s published register prices the
$e_i$ as \emph{affine} entries (the roster's register note); the
displayed premises carry no such restriction, and nothing above
consults one
(\autoref{rem:pb-affine}).
\end{brecipe}

\begin{brecipe}[payload transport --- field for field]
\label{rec:pb-transport}
\begin{center}\small
\begin{tabular}{@{}ll@{}}
\textbf{\rulename{PBound} payload} & \textbf{target (\rulename{Syzygy}) payload}\\
\midrule
\texttt{components} & \texttt{components} (verbatim)\\
\texttt{combination} & \texttt{combination} (verbatim)\\
\texttt{conclusion} & \texttt{conclusion} (verbatim)\\
\texttt{denBoundLog2} & \texttt{denBoundLog2} (verbatim; witness-side claim, zero proof-theoretic weight)\\
\texttt{preContains} & \texttt{preContains} (verbatim)\\
\texttt{conclShape, conclFin} & \texttt{conclShape, conclFin} (verbatim)\\
\texttt{wiresPos, compPostWiresPos} & \texttt{wiresPos, compPostWiresPos} (verbatim)\\
\texttt{valid} & \texttt{valid} (the same exact-combination proposition)\\
\texttt{postForm} & \texttt{postForm} (the same three-reading selector)\\
\end{tabular}
\end{center}
Every row is the identity --- including the last: since the target rule
itself carries the three-reading conclusion selector, the alias
inherits all readings with no residue, and the transport is $1 \to 1$
at reading level, not merely at node level.
\end{brecipe}

\begin{bprop}[the exactness identity]\label{prop:pb-exact}
The respelling is exact definitionally: the unpacked node's combination
identity is \emph{the same} exact canonical-polynomial identity
$c \equiv \sum_i \lambda_i e_i + \sum_k \nu_k g_k$, unchanged. What
remains to make the three readings correct is the fact that they are
one exact $\Q$ identity re-presented under a sign flip --- per entry,
$-e(\sigma) \leq 0 \iff e(\sigma) \geq 0$ and
$l(\sigma) - r(\sigma) \leq 0 \iff l(\sigma) \leq r(\sigma)$ --- which
is the target rule's reading-transport lemma, proven once in the
\rulename{Syzygy} section and \emph{not} re-justified here.
\end{bprop}
\begin{proof}
The first claim is reflexivity of the transported field. The second is
the cited lemma (the ordered-field axioms negate-and-flip and
add-to-both-sides, packaged as the ge-bridge ring lemma); the alias
contributes no polarity argument of its own, which is what keeps the
$1 \to 1$ relation clean.
\end{proof}

\begin{bprop}[soundness corollary]\label{prop:pb-sound}
If every component sequent is semantically valid, then under the
payload's checked hypotheses the sequent $\Phi \vdash \Psi$ is
semantically valid, for every reading $\Psi$ the selector admits.
\end{bprop}
\begin{proof}
\texttt{sound\_pBound} $=$ \texttt{sound\_syzygy} at the transported
(identical) payload data, fed the same component validities. There is
no other content --- the alias genre working as intended.
\end{proof}

\begin{bremark}[why the rule exists at all]\label{rem:pb-why}
Semantically the alias adds nothing to its target; its existence is
pragmatic, twice over. \emph{Register}: affine entries admit a compact
coefficient-row wire spelling, and bound nodes are the bulk commodity
of produced certificates --- the compact register is worth a tag.
\emph{Readings at the seams}: the (G) reading makes a nonnegativity
premise (as \rulename{ReLU-Pos} demands, $h \geq 0$) reachable
directly from bound facts --- the once-contemplated standalone
$\geq$-bridge rule dissolved into the conclusion selector instead ---
and the (L$\leq$) reading hands interval-shaped consumers
(\rulename{LiRPA}'s endpoint premises) their literal atom with no
\rulename{ArithAtom} re-arrangement node in between. Producers get the
convenient conclusion shape; the kernel keeps exactly one combination
argument.
\end{bremark}

\begin{bremark}[the affine restriction is register, not soundness]
\label{rem:pb-affine}
The register's ``affine entries'' --- the roster note's pricing
vocabulary, appearing in neither the displayed premises nor the
statement --- names the wire spelling
the rule is priced for, and is not a checked side condition: the
combination replay is degree-blind, and the soundness argument
(Proposition~\ref{prop:pb-sound}) never consults linearity. A
higher-degree entry simply belongs to the target rule's own spelling
--- what would refuse it is the published register\textquotesingle s
compact coefficient-row spelling, an encoding fact of that format
(asserted at its register, not displayed here), never the checker.
Publishing the restriction beside the rule, while the checker enforces
only what soundness needs, mirrors the \rulename{Farkas} page's
linearity note --- the register tells producers what the format prices;
the checker tells them what truth requires.
\end{bremark}

\begin{anchors}
Recipe \ref{rec:pb-unpack} unpack (identity transport) & \texttt{sound\_pBound}'s body: one application of \texttt{sound\_syzygy} to the step assembled from \rulename{PBound}'s fields (\texttt{PBoundData} $=$ the target's fields verbatim)\\
Recipe \ref{rec:pb-filled} obligations & \texttt{syzygyCombinationValid} and the target's bookkeeping fields, unchanged\\
Recipe \ref{rec:pb-transport} reading selector & \texttt{syzygyConclusionForm} (shared, not duplicated)\\
Proposition \ref{prop:pb-exact} reading transport & \texttt{holds\_of\_syzygyConclusionForm} (proven in the \rulename{Syzygy} section)\\
Proposition \ref{prop:pb-sound} & \texttt{sound\_pBound} (a \texttt{sound\_syzygy} specialization)\\
wire placement & \texttt{RuleTag.pBound}; no wire tag allocated in v0 --- producers elaborate before emission (the template\textquotesingle s second branch), and the emit path refuses the rule accordingly\\
\end{anchors}


\rulesection{28}{PsdCertTri}{admissible rule}
\label{blcatrule-PsdCertTri}
\ruleblock{The rule (Steele metanotation)}
\sourcenote{as realized in the mechanized development; source row: \texttt{\texttt{rules.toml}} --- \emph{PsdCertTri}. The metanotation statement is typeset by hand for notational quality; its mechanized ground truth is the Lean statement below, extracted from the sources at every build of this book and hash-bound. The rule \emph{roster} is machine-checked (every mechanized rule tag round-trips to exactly one catalogue row); the statement \emph{text} is editorial, checked by eye against the extracted Lean beside it.}
\begin{mathpar}
\inferrule*[Right=\rulename{PsdCertTri}]{
  A \in \Q^{n \times n} \ \text{symmetric}
  \\
  Z \ \text{upper-triangular}, \ Z_{ii} > 0
  \\
  A \cdot Z \ \text{lower-triangular}, \ (A \cdot Z)_{ii} > 0 \ \text{(recomputed exactly)}
  \\
  p_1, \dots, p_n \ \text{payload terms}
  \\
  c \equiv -\textstyle\sum_{i,j} A_{ij}\, p_i\, p_j \ \text{(exact canonical polynomial)}
  \\
  \Psi \in \bigl\{\, c \leq 0, \;\; Q \geq 0 \text{ where } c = -Q, \;\; l \leq r \text{ where } c = l - r \,\bigr\}
}{
  \Phi \vdash \Psi
}
\end{mathpar}

\par\medskip
\rulename{PsdCertTri} is the exact tier of the positive-definiteness
certificate family: a premise-free rule concluding nonnegativity of the
quadratic form $\sum_{i,j} A_{ij}\,p_i\,p_j$ of a symmetric rational
matrix $A$ at a payload term vector, warranted by an untrusted witness
$Z$ passing three decidable syntactic checks --- $A$ symmetric, $Z$
upper-triangular with positive diagonal, and the product $A \cdot Z$
(recomputed exactly, one sparse multiply) lower-triangular with
positive diagonal. Those premises are a theorem-strength certificate of
$A \succ 0$: the congruence $Z^{\mathsf T}(AZ)$ collapses to a positive
diagonal, and positivity transports back through the invertible $Z$.
The matrix-witness check is division-free, $O(\alpha n^2)$
multiplications ($\alpha$
average nonzeros per row), and runs in $\mathbb{Z}$ on integer inputs.

A producer reaches for it wherever a proof obligation is
\textquotedblleft this explicit matrix is positive
(semi)definite\textquotedblright\ --- Lyapunov and storage-function
certificates, energy quadratic forms, elimination-free replacements for
exact rational LDL factorization, whose $O(n^3)$
divisions-with-fill-in it displaces. The witness is produced off-stage
(a float Cholesky, exact elimination, anything); soundness never
consults its accuracy --- a forged $Z$ can only refuse. The concluded
fact enters the logic through the same conclusion-form selector as
\rulename{Syzygy} (raw zero-form, the flipped $Q \geq 0$ reading, or a
literal $l \leq r$ split), so it slots directly into
\rulename{Syzygy}/\rulename{Farkas} combinations as a component
conclusion; strictness is deliberately dropped at the formula boundary
(a premise-free atom carries no $p \neq 0$), so the in-logic fact is
$\succcurlyeq$-shaped even though the matrix-side theorem is strict.

What the checker recomputes: the three matrix checks entrywise over
exact $\Q$, and --- the recompute pattern --- that the payload
conclusion term\textquotesingle s canonical polynomial IS the negated
quadratic form $-\sum_{i,j} A_{ij}\,p_i\,p_j$ built from the payload
terms (in either the wire-keyed or the wire-coordinate polynomial
space). What breaks without each hypothesis: without symmetry the
congruence argument dies and the premises say nothing about the
quadratic form; without the $Z$ diagonal signs the product argument
loses invertibility or positivity (the two sign families are
determinant-linked --- for an indefinite $A$, flipping a column to
repair one violates the
other); the $AZ$ pattern and sign checks are the content itself, and
for an indefinite $A$ no witness whatsoever passes; without the
conclusion identity the posted atom could be unrelated to $A$; the
shape/finiteness/positive-size fields on the conclusion are the
standard evaluability-and-vacuity discipline (a zero-size wire
satisfies atoms vacuously, and the positive-size field refuses that
vacuous surface). The proof section carries the family
preamble and walks these in order.

\cataloguenote{Provenance: the triangular-inverse method is due to David
\textquotedblleft davidad\textquotedblright\ Dalrymple; Geoffrey Irving
derived and published the rounding-tolerant integer analysis and
formalized that variant at scale (a certified 65536-node spectral
expander) --- see the proof section\textquotesingle s citations. The
kernel development here is an independent proof of the exact tier.
Wire tag: deferred to the payload-window allocation; the emit path
refuses the rule until allocated.
}
\ruleblock{Lean statement}
\declhead{BernsteinLogic.Soundness}{sound\_psdCertOf}{\badge{BadgeGreen}{statement bound}}
\begin{leansrc}
/-- The shared PsdCert semantic lemma: from the abstract rational-points
quadratic-form fact (supplied per-variant by the PsdCert bridge
theorems), through the checked conclusion identity, to sequent validity
at every reading of `syzygyConclusionForm`. -/
theorem sound_psdCertOf {H : Hypernet} {seq : Sequent} {a : Nat}
    {M : Fin a → Fin a → ℚ} {ts : List Term} {conclusion : Term}
    (hpsd : ∀ x : Fin a → ℚ, 0 ≤ PsdCert.finQSum a fun i =>
      PsdCert.finQSum a fun j => M i j * x i * x j)
    (hconcl : PsdCert.conclCheck a M ts conclusion = true)
    (hshape : ∃ s, conclusion.bshape H = some s)
    (hfin : conclusion.finWires H)
    (hwpos : ∀ w ∈ conclusion.freeWires,
      ∃ τ, H.wireType w = some τ ∧ 0 < Shape.totalSize τ.shape)
    (hform : syzygyConclusionForm conclusion seq.post) :
    seq.semanticallyValid H
\end{leansrc}
\declhead{BernsteinLogic.Soundness}{sound\_psdCertTri}{\badge{BadgeGreen}{statement bound}}
\begin{leansrc}
/-- PsdCertTri soundness: the /T witness warrants A ≻ 0 over ℝ
(`PsdCert.checkTri_posDef`, davidad's triangular-inverse theorem), whose
quadratic form is nonnegative at every rational point. -/
theorem sound_psdCertTri {H : Hypernet} {seq : Sequent} {a : Nat}
    {A Z : Fin a → Fin a → ℚ} {ts : List Term} {conclusion : Term}
    (hwit : PsdCert.checkTri a A Z = true)
    (hconcl : PsdCert.conclCheck a A ts conclusion = true)
    (hshape : ∃ s, conclusion.bshape H = some s)
    (hfin : conclusion.finWires H)
    (hwpos : ∀ w ∈ conclusion.freeWires,
      ∃ τ, H.wireType w = some τ ∧ 0 < Shape.totalSize τ.shape)
    (hform : syzygyConclusionForm conclusion seq.post) :
    seq.semanticallyValid H
\end{leansrc}
\par\medskip
\textbf{The matrix theorems (mathlib side).}\enspace The chain from
witness check to consumable fact: the congruence collapse and transport
lemmas prove the triangular-witness theorem; \texttt{checkTri\_posDef}
bridges the executable Bool check to positive definiteness of the
$\R$-cast; and the rational-points corollary is the exact shape the
rule soundness lemma consumes. Soundness never consults witness
accuracy --- every hypothesis of every theorem below is a checked
premise, so a forged witness can only fail a decidable check.

\declhead{BernsteinLogic.PsdCert}{congruence\_eq\_diagonal}{\badge{BadgeGreen}{statement bound}}
\begin{leansrc}
/-- The congruence `Zᴴ * (A * Z)` collapses to a diagonal matrix when `Z` is
upper-triangular and `A * Z` is lower-triangular, and its diagonal entries
are the products of the two matrices' diagonals. -/
theorem congruence_eq_diagonal {A Z : Matrix n n ℝ}
    (hA : ∀ i j, A i j = A j i)
    (hZu : ∀ ⦃i j⦄, j < i → Z i j = 0)
    (hMl : ∀ ⦃i j⦄, i < j → (A * Z) i j = 0) :
    Zᴴ * (A * Z) = diagonal fun i => Z i i * (A * Z) i i
\end{leansrc}
\declhead{BernsteinLogic.PsdCert}{posDef\_of\_congruence\_posDef}{\badge{BadgeGreen}{statement bound}}
\begin{leansrc}
/-- Congruence transport (factored from the /T proof for reuse by /DD): if
the congruence `Zᴴ * (A * Z)` is positive definite and `Z` is
upper-triangular with positive diagonal (hence invertible), then `A` is
positive definite. -/
theorem posDef_of_congruence_posDef {A Z : Matrix n n ℝ}
    (hZu : ∀ ⦃i j⦄, j < i → Z i j = 0)
    (hZd : ∀ i, 0 < Z i i)
    (hBd : (Zᴴ * (A * Z)).PosDef) :
    A.PosDef
\end{leansrc}
\declhead{BernsteinLogic.PsdCert}{posDef\_of\_triangular\_witness}{\badge{BadgeGreen}{statement bound}}
\begin{leansrc}
/-- **PsdCertTri soundness** (davidad's theorem). A symmetric matrix with an
upper-triangular positive-diagonal witness `Z` such that `A * Z` is
lower-triangular with positive diagonal is positive definite. -/
theorem posDef_of_triangular_witness {A Z : Matrix n n ℝ}
    (hA : ∀ i j, A i j = A j i)
    (hZu : ∀ ⦃i j⦄, j < i → Z i j = 0)
    (hZd : ∀ i, 0 < Z i i)
    (hMl : ∀ ⦃i j⦄, i < j → (A * Z) i j = 0)
    (hMd : ∀ i, 0 < (A * Z) i i) :
    A.PosDef
\end{leansrc}
\declhead{BernsteinLogic.PsdCert}{checkTri\_posDef}{\badge{BadgeGreen}{statement bound}}
\begin{leansrc}
/-- `checkTri` warrants positive definiteness of the ℝ-cast. -/
theorem checkTri_posDef {a : Nat} {A Z : Fin a → Fin a → ℚ}
    (h : checkTri a A Z = true) : (matR A).PosDef
\end{leansrc}
\declhead{BernsteinLogic.PsdCert}{quadForm\_nonneg\_of\_posSemidef}{\badge{BadgeGreen}{statement bound}}
\begin{leansrc}
/-- Nonnegativity of the ℚ quadratic form at every rational point, from
positive semidefiniteness of the ℝ-cast — the fact the rule soundness
lemmas consume. -/
theorem quadForm_nonneg_of_posSemidef {a : Nat} {M : Fin a → Fin a → ℚ}
    (h : (matR M).PosSemidef) (x : Fin a → ℚ) :
    0 ≤ finQSum a (fun i => finQSum a (fun j => M i j * x i * x j))
\end{leansrc}
\declhead{BernsteinLogic.PsdCert}{checkTri\_quadForm\_nonneg}{\badge{BadgeGreen}{statement bound}}
\begin{leansrc}
/-- /T's consumable corollary. -/
theorem checkTri_quadForm_nonneg {a : Nat} {A Z : Fin a → Fin a → ℚ}
    (h : checkTri a A Z = true) (x : Fin a → ℚ) :
    0 ≤ finQSum a (fun i => finQSum a (fun j => A i j * x i * x j))
\end{leansrc}
\par\medskip
\textbf{The executable check (zero-mathlib closure).}\enspace The
witness checks are List-based exact rational computations on the
executable checker\textquotesingle s import closure --- displayed in
full so the reader can confirm by eye that they are sign, zero-pattern,
and entrywise-equality inspections over recomputed products, with no
divisions and no trusted certificate data beyond the witness itself;
and that \texttt{conclCheck} pins the payload conclusion\textquotesingle s
canonical polynomial to the negated quadratic form (both polynomial
branches mirror the \rulename{Syzygy} check\textquotesingle s split).
These definitions are shared by all three family members; the /DD and
/K sections display only their own additional entries.

\declhead{BernsteinLogic.Core.Catalogue}{qlsum}{\badge{BadgeGreen}{definition bound}}
\begin{leansrc}
/-- List sum over Rat (zero-mathlib; `List.sum`'s instance path avoided
deliberately on the core closure). -/
def qlsum (l : List Rat) : Rat := l.foldr (· + ·) 0
\end{leansrc}
\declhead{BernsteinLogic.Core.Catalogue}{qabs}{\badge{BadgeGreen}{definition bound}}
\begin{leansrc}
/-- `|q|` without lattice instances. -/
def qabs (q : Rat) : Rat := if q < 0 then -q else q
\end{leansrc}
\declhead{BernsteinLogic.Core.Catalogue}{finQSum}{\badge{BadgeGreen}{definition bound}}
\begin{leansrc}
/-- Σ_{i < n} f i, List-based. -/
def finQSum (n : Nat) (f : Fin n → Rat) : Rat := qlsum ((List.finRange n).map f)
\end{leansrc}
\declhead{BernsteinLogic.Core.Catalogue}{mulEntry}{\badge{BadgeGreen}{definition bound}}
\begin{leansrc}
/-- Entry (i,j) of the product A·B, recomputed exactly. -/
def mulEntry (n : Nat) (A B : Fin n → Fin n → Rat) (i j : Fin n) : Rat :=
  finQSum n (fun k => A i k * B k j)
\end{leansrc}
\declhead{BernsteinLogic.Core.Catalogue}{symB}{\badge{BadgeGreen}{definition bound}}
\begin{leansrc}
/-- Symmetry, checked entrywise. -/
def symB (n : Nat) (A : Fin n → Fin n → Rat) : Bool :=
  (List.finRange n).all fun i => (List.finRange n).all fun j =>
    decide (A i j = A j i)
\end{leansrc}
\declhead{BernsteinLogic.Core.Catalogue}{upperTriPosDiagB}{\badge{BadgeGreen}{definition bound}}
\begin{leansrc}
/-- Upper-triangular with positive diagonal, checked entrywise. -/
def upperTriPosDiagB (n : Nat) (Z : Fin n → Fin n → Rat) : Bool :=
  ((List.finRange n).all fun i => (List.finRange n).all fun j =>
    decide (j < i → Z i j = 0)) &&
  ((List.finRange n).all fun i => decide (0 < Z i i))
\end{leansrc}
\declhead{BernsteinLogic.Core.Catalogue}{checkTri}{\badge{BadgeGreen}{definition bound}}
\begin{leansrc}
/-- PsdCertTri witness check (davidad's theorem, all-syntactic premises):
A symmetric; Z upper-triangular with positive diagonal; A·Z
lower-triangular with positive diagonal. Concludes A ≻ 0 (matrix side);
division-free, O(αn²) multiplications. -/
def checkTri (n : Nat) (A Z : Fin n → Fin n → Rat) : Bool :=
  symB n A && upperTriPosDiagB n Z &&
  ((List.finRange n).all fun i => (List.finRange n).all fun j =>
    decide (i < j → mulEntry n A Z i j = 0)) &&
  ((List.finRange n).all fun i => decide (0 < mulEntry n A Z i i))
\end{leansrc}
\declhead{BernsteinLogic.Core.Catalogue}{negQuadFormPoly}{\badge{BadgeGreen}{definition bound}}
\begin{leansrc}
/-- The negated quadratic-form canonical polynomial
−Σᵢⱼ Mᵢⱼ·pᵢ·pⱼ (zero-normal orientation: the rule's primitive fact is
`conclusion ≤ 0` with conclusion ~ −Q, exposed through
`syzygyConclusionForm` — a producer typically emits the ge-flip reading
`Q ≥ 0` with the conclusion term literally `.neg Q`). -/
def negQuadFormPoly (n : Nat) (M : Fin n → Fin n → Rat)
    (ps : List QArith.Poly) : QArith.Poly :=
  QArith.Poly.neg ((List.finRange n).foldr (fun i acc =>
    (List.finRange n).foldr (fun j acc2 =>
      QArith.Poly.add (QArith.Poly.scale (M i j)
        (QArith.Poly.mul (ps.getD i.val []) (ps.getD j.val []))) acc2) acc) [])
\end{leansrc}
\declhead{BernsteinLogic.Core.Catalogue}{conclCheckCore}{\badge{BadgeGreen}{definition bound}}
\begin{leansrc}
/-- Conclusion check, idx-free branch (syzygyCheckCore's discipline):
the payload conclusion's canonical polynomial IS the negated quadratic
form of M at the payload terms' canonical polynomials — the recompute
pattern, no coefficient tensors trusted. -/
def conclCheckCore (n : Nat) (M : Fin n → Fin n → Rat) (ts : List Term)
    (conclusion : Term) : Bool :=
  decide (ts.length = n) &&
  conclusion.idxFree && ts.all Term.idxFree &&
  decide (conclusion.toPoly = negQuadFormPoly n M (ts.map Term.toPoly))
\end{leansrc}
\declhead{BernsteinLogic.Core.Catalogue}{conclCheckW}{\badge{BadgeGreen}{definition bound}}
\begin{leansrc}
/-- Conclusion check, wire-coordinate branch (syzygyCheckW's G1 route:
multi-wire idx atoms — the vector-wire quadratic forms the E-4a /
Lyapunov customers emit). -/
def conclCheckW (n : Nat) (M : Fin n → Fin n → Rat) (ts : List Term)
    (conclusion : Term) : Bool :=
  decide (ts.length = n) &&
  (match conclusion.asCoordPolyIdxW?, ts.mapM Term.asCoordPolyIdxW? with
   | some pc, some pts => decide (pc = negQuadFormPoly n M pts)
   | _, _ => false)
\end{leansrc}
\declhead{BernsteinLogic.Core.Catalogue}{conclCheck}{\badge{BadgeGreen}{definition bound}}
\begin{leansrc}
/-- Either conclusion branch (mirrors `syzygyCheck`'s split). -/
def conclCheck (n : Nat) (M : Fin n → Fin n → Rat) (ts : List Term)
    (conclusion : Term) : Bool :=
  conclCheckCore n M ts conclusion || conclCheckW n M ts conclusion
\end{leansrc}
\ruleblock{Correspondence (Steele display $\leftrightarrow$ Lean)}
\badge{BadgeGreen}{reviewer-certified against the displayed texts}
\sourcenote{The dictionary below is a reviewer-certified statement of how the mechanized statements above interpret the displayed rule: display metavariable $\mapsto$ Lean binder, each display hypothesis $\mapsto$ the premise or side condition carrying it, the display conclusion $\mapsto$ the goal --- with any deliberate difference in width between display and mechanization named outright. Certification is recorded as a hash triple over (the displayed cell, the extracted statement texts, this dictionary\textquotesingle s own text); if any of the three changes, this block turns amber until a reviewer re-certifies. The correspondence claim itself remains editorial --- reviewer-grade, not machine-checked; its machine-checked successor types each soundness statement against the rule as a structured object, and supersedes this block per rule as it lands.}
Judgment frame as in \rulename{Ax}. This dictionary is the
trio\textquotesingle s template; /DD and /K differ in the witness row
alone.
\begin{itemize}
\item $n \mapsto$ \texttt{a : Nat}; $A \mapsto$ \texttt{A}; $Z
  \mapsto$ \texttt{Z} (both \texttt{Fin a → Fin a → ℚ}); $p_1,
  \ldots, p_n \mapsto$ \texttt{ts : List Term}; $c \mapsto$
  \texttt{conclusion : Term}.
\item premise-freeness $\mapsto$ silence, the \rulename{PConstExact}
  representation: no hypothesis mentions \texttt{seq.pre}, so $\Phi$
  is arbitrary --- the \rulename{Norm} convention of the roster note,
  carried structurally.
\item the three displayed matrix rows (symmetry; upper-triangular $Z$
  with positive diagonal; lower-triangular $A \cdot Z$ with positive
  diagonal, recomputed exactly) $\mapsto$ the single Boolean witness
  \texttt{hwit} $:$ \texttt{PsdCert.checkTri a A Z = true}, whose
  definition (displayed above) recomputes all three. The mathematics
  enters through the certified-alongside bridge:
  \texttt{checkTri\_posDef} warrants $A \succ 0$ over $\R$, and
  \texttt{sound\_psdCertTri} consumes only its rational-points
  quadratic-form corollary \texttt{checkTri\_quadForm\_nonneg} --- the
  displayed strictness drop at the formula boundary, visible in the
  statement\textquotesingle s type.
\item the canonical-polynomial row $c \equiv -\sum_{i,j} A_{ij} p_i
  p_j$ $\mapsto$ \texttt{hconcl} $:$ \texttt{conclCheck a A ts
  conclusion = true} --- the identity is checked on canonical forms by
  either of two certified routes (the idx-free \texttt{toPoly} replay
  of \texttt{conclCheckCore}, or the coordinate-indexed
  \texttt{conclCheckW}), and \texttt{sound\_psdCertOf} discharges
  both.
\item $\Psi \in \{c \leq 0;\ Q \geq 0$ where $c = -Q;\ l \leq r$
  where $c = l - r\}$ $\mapsto$ \texttt{hform} $:$
  \texttt{syzygyConclusionForm conclusion seq.post} --- the certified
  selector is literally that three-membered disjunction, reading the
  conclusion term\textquotesingle s head syntactically (shared with
  \rulename{Syzygy} and \rulename{PBound}).
\item with no displayed row: the evaluability pack ---
  \texttt{hshape} (the conclusion has a broadcast shape in
  \texttt{H}), \texttt{hfin} (its wires are finite-typed), and
  \texttt{hwpos} (every free wire positively sized). The display
  omits all three; the certified statement demands them, and
  \texttt{hwpos} is the same vacuous-surface refusal the
  \rulename{Farkas} display carries inline and the
  \rulename{LiRPA} linkage builds in.
\end{itemize}
Width note: elementwise atom semantics batch the conclusion (the
roster note), so the display\textquotesingle s scalar reading is the
singleton case of the certified one; conversely the certified
statement never uses witness accuracy beyond the Boolean check ---
a forged $Z$ refuses, never falsely accepts --- and no hypothesis is
unused.

\ruleblock{Lean verifier receipts}
\sourcenote{module \texttt{BernsteinLogic.Soundness}; bound build \texttt{a20d90e66c13\ldots} (bound from receipt a9c0e34d3e3c9c8c, 2026-08-03T21:30:36Z). A receipt records that named checking layers accepted a build: (module, compiled-artifact SHA-256, environment SHA-256, checker layers), minted when the checking gate ran; the store receipts attesting the bound build are listed with the layers each run credited. Nothing here asks the reader to trust the store: every layer is an independently runnable program over the quoted sources --- rebuild the module and re-run the checkers to reproduce the verdict. Receipts attesting superseded builds credit no layer and are omitted from this rendering (the count of omissions, if any, is noted below).}
\begin{receipttable}
\receiptrow{a9c0e34d3e3c9c8c}{2026-08-03T21:30:36Z}{build+decorrelated}{\badge{BadgeGreen}{current}}
\end{receipttable}
\sourcenote{21 further receipt(s) in the store attest superseded builds of this module and are omitted from this rendering.}
\sourcenote{module \texttt{BernsteinLogic.PsdCert}; bound build \texttt{be59faf6aff0\ldots} (bound from receipt 872a60504989acef, 2026-07-29T07:57:01Z). A receipt records that named checking layers accepted a build: (module, compiled-artifact SHA-256, environment SHA-256, checker layers), minted when the checking gate ran; the store receipts attesting the bound build are listed with the layers each run credited. Nothing here asks the reader to trust the store: every layer is an independently runnable program over the quoted sources --- rebuild the module and re-run the checkers to reproduce the verdict. Receipts attesting superseded builds credit no layer and are omitted from this rendering (the count of omissions, if any, is noted below).}
\begin{receipttable}
\receiptrow{13403ea3a4e42f77}{2026-07-28T23:44:21Z}{build+decorrelated}{\badge{BadgeGreen}{current}}
\receiptrow{6ad2a620d75bd04e}{2026-07-28T23:56:32Z}{build+decorrelated}{\badge{BadgeGreen}{current}}
\receiptrow{d37ce088bc4caf69}{2026-07-29T00:55:46Z}{build+replay+hardened+decorrelated}{\badge{BadgeGreen}{current}}
\receiptrow{872a60504989acef}{2026-07-29T07:57:01Z}{build+replay+hardened+decorrelated}{\badge{BadgeGreen}{current}}
\end{receipttable}
\sourcenote{5 further receipt(s) in the store attest superseded builds of this module and are omitted from this rendering.}
\sourcenote{module \texttt{BernsteinLogic.Core.Catalogue}; bound build \texttt{cbc92e976a1d\ldots} (bound from receipt 5abdb3589707b5df, 2026-08-03T21:30:33Z). A receipt records that named checking layers accepted a build: (module, compiled-artifact SHA-256, environment SHA-256, checker layers), minted when the checking gate ran; the store receipts attesting the bound build are listed with the layers each run credited. Nothing here asks the reader to trust the store: every layer is an independently runnable program over the quoted sources --- rebuild the module and re-run the checkers to reproduce the verdict. Receipts attesting superseded builds credit no layer and are omitted from this rendering (the count of omissions, if any, is noted below).}
\begin{receipttable}
\receiptrow{24804e66edba0bc9}{2026-08-03T20:25:09Z}{build+decorrelated}{\badge{BadgeGreen}{current}}
\receiptrow{5abdb3589707b5df}{2026-08-03T21:30:33Z}{build+decorrelated}{\badge{BadgeGreen}{current}}
\end{receipttable}
\sourcenote{18 further receipt(s) in the store attest superseded builds of this module and are omitted from this rendering.}
\layerline{\layeritem{build}{lake build (elaboration)}{\badge{BadgeGreen}{current}}\layeritem{replay}{leanchecker (official kernel replay)}{\badge{BadgeGreen}{current}}\layeritem{hardened}{safe\_verify (hardened screen)}{\badge{BadgeGreen}{current}}\layeritem{decorrelated}{lean4export + nanoda (decorrelated kernel)}{\badge{BadgeGreen}{current}}}
\sourcenote{\badge{BadgeGray}{comparator: none} no fixture in the three-implementation comparator harness (compiled / kernel-interpreted / reference-interpreted) exercises this rule.}
\ruleblock{HOL4 statement}
\sourcenote{theory \texttt{bern\_tr\_soundpsd}, HOL4 development script \texttt{\texttt{bern\_tr\_soundpsdScript.sml}} \badge{BadgeGreen}{script bound}}
\holhead{Theorem sound\_psdCertOf}
\begin{holsrc}
Theorem sound_psdCertOf:
  ∀H sq n M ts conclusion.
    (∀x. 0 ≤ quadForm n M x) ∧
    psdConclCheck n M ts conclusion ∧
    (∃s. tmBshape H conclusion = SOME s) ∧
    tmFinWires H conclusion ∧
    (∀w. MEM w (termFreeWires conclusion) ⇒
         ∃τ. H.wireType w = SOME τ ∧ 0 < shapeTotalSize τ.shape) ∧
    syzygyConclusionFormB conclusion sq.post ⇒
    seqSemanticallyValid H sq
\end{holsrc}
\holhead{Theorem sound\_psdCertTri}
\begin{holsrc}
Theorem sound_psdCertTri:
  ∀H sq n A Z ts conclusion.
    psdCheckT n A Z ∧
    psdConclCheck n A ts conclusion ∧
    (∃s. tmBshape H conclusion = SOME s) ∧
    tmFinWires H conclusion ∧
    (∀w. MEM w (termFreeWires conclusion) ⇒
         ∃τ. H.wireType w = SOME τ ∧ 0 < shapeTotalSize τ.shape) ∧
    syzygyConclusionFormB conclusion sq.post ⇒
    seqSemanticallyValid H sq
\end{holsrc}
\ruleblock{HOL verifier receipts}
\badge{BadgeBlue}{documentary} \texttt{hol-strict} CLEAN: theory \texttt{bern\_tr\_soundpsd}, 4
theorems, 0 cheat-tagged, 0 nonstandard axioms (whitelist:
\texttt{DISK\_THM}; the four \texttt{bool} axioms), 249-theory
ancestry --- verification run of 2026-07-22 on an isolated build host,
against a content-addressed checkout of the exact sources cited here
(git blob \texttt{c01e0b0c\ldots}). The twins mirror the Lean
statements clause-for-clause per the script\textquotesingle s
statement-fidelity inventory: the shared semantic spine
(\texttt{sound\_psdCertOf}) runs from the abstract rational-points
quadratic-form fact through the checked conclusion identity to sequent
validity at every reading of the shared conclusion-form selector, and
the per-shape corollary discharges that fact via the ported
triangular-witness bridge (\texttt{psdCheckT\_quadForm}, theory
\texttt{bern\_psd}, where the witness check itself ---
\texttt{psdCheckT} --- is ported with kernel-EVAL receipts). The
payload record and its wellformedness arm ride the ported proof-object
stratum (theory \texttt{bern\_tr\_catdata2}, 272 theorems, 0 offences,
same run).

\sourcenote{HOL receipts are documentary: quoted results of \texttt{hol-strict} runs
and census rows from the HOL4 development\textquotesingle s build
evidence, not minted receipt artefacts. A HOL receipt store (the Lean
store\textquotesingle s twin) is named future work.
}
\ruleblock{Proof}

\sourcenote{Genre: admissible rule (tier 2) --- a recipe by
\emph{reduction to verified machinery} (the \rulename{PConstExact}
pattern): no catalogue derivation of an exact matrix-positivity fact
exists, which is the trio's point; each rule discharges against a named
kernel-verified matrix theorem. This section also states the
certificate-family preamble the \rulename{PsdCertDiag} and
\rulename{PsdCertKernel} sections reference.}

\begin{brecipe}[the PSD-certificate family --- shared shape]
\label{rec:psd-preamble}
All three rules certify positivity of a symmetric matrix
$A \in \Q^{n\times n}$ from an \emph{untrusted witness}, by decidable
syntactic checks in exact rational arithmetic --- no divisions, no
pivoting, no floating point on the trusted side. The shared discipline:

\emph{(i) Soundness never consults the witness's accuracy.} The
conclusion is a theorem of the checked premises; a forged or degraded
witness can only cause a false REJECT, never a false ACCEPT. Producers
may therefore run arbitrary float factorizations off-stage.

\emph{(ii) The in-logic conclusion is the quadratic form, not the
matrix.} The formula language has no matrix atoms, so the rule's
conclusion is nonnegativity of $Q(p) = \sum_{i,j} A_{ij}\,p_i\,p_j$ at
a payload term vector $p_1,\ldots,p_n$: the payload conclusion term $c$
is checked --- by exact canonical-polynomial identity, the recompute
pattern --- to \emph{be} $-\!\sum_{i,j} A_{ij}\,p_i\,p_j$, and the
sequent's posted formula is one of three readings of $c \leq 0$ (the
raw zero-form; the flipped $Q \geq 0$; the literal $l \leq r$ where
$c = l - r$) --- the same conclusion-form selector displayed in the
\rulename{Syzygy} section, so the fact slots directly into
\rulename{Syzygy}/\rulename{Farkas} combinations as a component
conclusion.

\emph{(iii) Strictness is deliberately dropped at the formula
boundary.} The matrix theorems certify $A \succ 0$ (for /T and /DD),
but the premise-free atom ships the $\succcurlyeq$-shaped fact: strict
positivity of $Q(p)$ would need $p \neq 0$, which no premise-free atom
carries.

\emph{(iv) The rules are premise-free}; $\Phi$ is arbitrary (a fact
holding in the empty context holds under any premise). Elementwise atom
semantics batch the conclusion: at tensor-shaped payload terms, entry
$e$ of the conclusion is the quadratic form at the terms' entry-$e$
values, so one node states a family of instances.

The rewrite, per rule, is degenerate in the \rulename{PConstExact}
sense: one node, zero catalogue premises, discharged by the rule's
matrix theorem composed with the shared transport lemma
(\texttt{sound\_psdCertOf}) that carries ``PSD over $\R$'' through
``nonnegative at every rational point'' and the checked polynomial
identity to the posted formula.
\end{brecipe}

\begin{brecipe}[the \rulename{PsdCertTri} instance]\label{rec:psdt-instance}
Witness: $Z \in \Q^{n\times n}$. Checked premises (all decidable, all
exact): $A$ symmetric; $Z$ upper-triangular with positive diagonal;
$M := A\,Z$ (recomputed exactly, one sparse multiply) lower-triangular
with positive diagonal. Then
\[
\inferrule*[Right=\rulename{PsdCertTri}]
  {\;\mathrm{checkTri}(A, Z) = \mathsf{true} \\
    c \equiv -\textstyle\sum_{i,j} A_{ij}\,p_i\,p_j\;}
  {\Phi \vdash \Psi}
\qquad\rightsquigarrow\qquad
\text{\upshape discharged by } \texttt{sound\_psdCertTri},
\]
with $\Psi$ any reading of $c \leq 0$ per the preamble. The
matrix-witness check costs $O(\alpha n^2)$ multiplications ($\alpha$ =
average nonzeros per row of $A$) plus $O(n^2)$ sign/zero inspections,
division-free; on integer inputs it runs entirely in $\mathbb{Z}$ with
$\{*, +, <\}$. (The conclusion-identity recompute is additional, and
its cost tracks the payload terms' polynomial sizes, not $n$ alone.)
\end{brecipe}

\begin{bprop}[matrix warrant]\label{prop:psdt-matrix}
If $A$ is symmetric, $Z$ upper-triangular with $Z_{ii} > 0$, and $AZ$
lower-triangular with $(AZ)_{ii} > 0$, then $A \succ 0$ over $\R$.
\end{bprop}
\begin{proof}
Let $B := Z^{\mathsf T}(AZ)$. As a product of two lower-triangular
matrices ($Z^{\mathsf T}$ and $AZ$) it is lower-triangular; since
$Z^{\mathsf T} A Z$ is a congruence of the symmetric $A$ it is also
symmetric; hence $B$ is \emph{diagonal}, and the sum
$B_{ii} = \sum_k Z_{ki} (AZ)_{ki}$ collapses to the single term
$Z_{ii}(AZ)_{ii} > 0$. So $B \succ 0$. Moreover
$\det Z = \prod_i Z_{ii} > 0$ makes $Z$ invertible, and
$A = Z^{-\mathsf T} B\, Z^{-1} \succ 0$ by congruence. (This is the
route the kernel proof takes --- \texttt{congruence\_eq\_diagonal}
then \texttt{posDef\_of\_congruence\_posDef} --- avoiding leading
principal minors; the classical alternative reads $A = M Z^{-1}$ as an
LU factorization with positive-diagonal factors and applies Sylvester's
criterion \cite{hornjohnson2013}.)
\end{proof}

\begin{bremark}[each hypothesis, and what removing it admits]
\emph{Symmetry of $A$}: without it $Z^{\mathsf T}(AZ)$ need not be
symmetric, the diagonal-collapse argument dies, and the quadratic form
reads only the symmetric part of $A$ --- the premises would say nothing
about it. \emph{Upper-triangularity and positive diagonal of $Z$}:
triangularity is the cheapest syntactic invertibility certificate
($\det Z = \prod Z_{ii}$), and the diagonal signs feed
$B_{ii} > 0$; the two sign families are determinant-linked: for an
indefinite $A$ --- a genuinely negative pivot --- ``repairing'' the
$AZ$ diagonal sign by flipping the offending column of $Z$ only moves
the violation into $Z$'s own diagonal check ($A = \mathrm{diag}(1,-1)$,
$Z = \mathrm{diag}(1,z)$: making $-z$ positive forces $z < 0$; when $A$
is in fact positive definite, a flip can repair both checks at once ---
acceptance working as intended, not a leak). \emph{The pattern and sign
checks on $AZ$}: this is the entire content --- the diagonal of $AZ$
carries the pivot signs of $A$, and for an indefinite $A$ \emph{no}
witness whatsoever passes (by Proposition~\ref{prop:psdt-matrix} an
accepting witness would prove $A \succ 0$). \emph{The conclusion
identity}: without the recomputed polynomial check the posted atom
could be an arbitrary formula unrelated to $A$.
\end{bremark}

\begin{bremark}[/T needs an exact witness --- the honest split with /DD]
The natural producer emits $Z \approx$ an inverse transposed Cholesky
factor. A \emph{rounded} such $Z$ makes the strictly-upper entries of
$AZ$ small but nonzero, and /T's pattern check rightly refuses. An
exact witness always exists when $A \succ 0$, and rationally: LDL$^{\mathsf T}$
elimination gives $A = L D L^{\mathsf T}$ with unit lower-triangular
rational $L$ and positive rational diagonal $D$, so
$Z := L^{-\mathsf T}$ is rational upper-triangular with positive
diagonal and $A Z = L D$ is lower-triangular with positive diagonal
--- clear denominators columnwise if integer entries are wanted (each
column is determined up to positive scale) --- but the entries grow
like scaled minors, i.e.\ LDL-style bit growth. (The Cholesky factor
itself is generally irrational; the LDL route is the rational
construction.) At sizes where that growth
bites, the rounding-tolerant companion \rulename{PsdCertDiag} is the
production tier; at moderate sizes /T's $O(\alpha n^2)$ check wins the
clock. Neither tier ever risks soundness in the trade --- only
completeness margin moves.
\end{bremark}

\begin{bremark}[provenance]
The triangular-inverse method is due to David ``davidad'' Dalrymple
\cite{dalrymple2026psd}; Geoffrey Irving derived the rounding-tolerant
integer variant and published the certificate analysis
\cite{irving2026certificate}, then formalized that variant in Lean at
scale (a certified $65536$-node spectral expander) in the course of a
sorting-network formalization \cite{irvingaks}. The kernel development
cited here is an independent in-house proof of the exact tier, grown to
the /DD and /K companions.
\end{bremark}

\begin{anchors}
Recipe \ref{rec:psdt-instance} warrant & \texttt{sound\_psdCertTri} (Soundness); shared transport \texttt{sound\_psdCertOf}\\
matrix theorem (Prop.\ \ref{prop:psdt-matrix}) & \texttt{PsdCert.posDef\_of\_triangular\_witness} via \texttt{congruence\_eq\_diagonal}, \texttt{posDef\_of\_congruence\_posDef}\\
check-to-theorem bridge & \texttt{PsdCert.checkTri\_posDef}; rational-points corollary \texttt{checkTri\_quadForm\_nonneg}, \texttt{quadForm\_nonneg\_of\_posSemidef}\\
the executable check & \texttt{PsdCert.checkTri} (with \texttt{symB}, \texttt{upperTriPosDiagB}, \texttt{mulEntry}) --- zero-mathlib, exact $\Q$\\
conclusion identity & \texttt{PsdCert.conclCheck} against \texttt{negQuadFormPoly} (idx-free and wire-coordinate branches)\\
conclusion readings & the \rulename{Syzygy} section's conclusion-form selector and its transport lemma\\
\end{anchors}


\rulesection{29}{PsdCertDiag}{admissible rule}
\label{blcatrule-PsdCertDiag}
\ruleblock{The rule (Steele metanotation)}
\sourcenote{as realized in the mechanized development; source row: \texttt{\texttt{rules.toml}} --- \emph{PsdCertDiag}. The metanotation statement is typeset by hand for notational quality; its mechanized ground truth is the Lean statement below, extracted from the sources at every build of this book and hash-bound. The rule \emph{roster} is machine-checked (every mechanized rule tag round-trips to exactly one catalogue row); the statement \emph{text} is editorial, checked by eye against the extracted Lean beside it.}
\begin{mathpar}
\inferrule*[Right=\rulename{PsdCertDiag}]{
  A \in \Q^{n \times n} \ \text{symmetric}
  \\
  Z \ \text{upper-triangular}, \ Z_{ii} > 0
  \\
  B := Z^{\mathsf{T}} A Z \ \text{strictly row-diagonally dominant (recomputed exactly)}
  \\
  p_1, \dots, p_n \ \text{payload terms}
  \\
  c \equiv -\textstyle\sum_{i,j} A_{ij}\, p_i\, p_j \ \text{(exact canonical polynomial)}
  \\
  \Psi \in \bigl\{\, c \leq 0, \;\; Q \geq 0 \text{ where } c = -Q, \;\; l \leq r \text{ where } c = l - r \,\bigr\}
}{
  \Phi \vdash \Psi
}
\end{mathpar}

\par\medskip
\rulename{PsdCertDiag} is the rounding-tolerant tier of the
positive-definiteness certificate family (the family discipline ---
untrusted witness, quadratic-form conclusion through
\rulename{Syzygy}\textquotesingle s conclusion-form selector,
strictness dropped at the formula boundary, premise-freedom --- is
stated once in the \rulename{PsdCertTri} section). Here the checked
premises are: $A$ symmetric; $Z$ upper-triangular with positive
diagonal; and the congruence $B := Z^{\mathsf T} A Z$, recomputed
exactly by the checker, strictly row-diagonally dominant (positive
diagonal implied, since the off-row absolute sum is nonnegative). A
symmetric strictly dominant matrix is positive definite --- the
Gershgorin-circle reading; the kernel proof is a self-contained
arithmetic-geometric-mean regrouping of the quadratic form --- and
positivity transports back through the invertible $Z$.

This is the tier a real producer feeds: $Z := \lceil 2^s R^{-1}
\rfloor$ rounded from an off-stage float Cholesky $A \approx
R^{\mathsf T} R$ works directly, with the required scale $s$ growing
only logarithmically in dimension and input scale rather than tracking
any accuracy target --- witness entries encode in $s +
O(\log^{+} \|R^{-1}\|_{\max})$ bits (scale exponent plus the
factorization\textquotesingle s own magnitude), dimension
entering only through the logarithm, where the exact /T witness
inherits minor-scale bit growth.
What a too-coarse witness costs is only completeness margin: the check
fails closed and the producer retries finer, with soundness never in
play. The price relative to /T is the $O(n^3)$ congruence product ---
on bounded-width integers, embarrassingly parallel, still
division-free.

What breaks without each hypothesis: without symmetry of $A$ the
regrouping step (which folds the dominance envelope using
$B_{ij} = B_{ji}$) fails and row dominance alone bounds nothing about
the form; without strictness the matrix warrant $A \succ 0$ would overclaim on
singular matrices
(a path-graph Laplacian is weakly dominant with the constants in its
kernel --- the $\succcurlyeq$-shaped posted atom would survive, but
the strict theorem /K consumes on its inner block would not); without exact recomputation of $B$ a wire-supplied congruence
could be dominant while unrelated to $A$; the $Z$ checks are, as in
/T, the cheapest syntactic invertibility certificate. The conclusion
identity and evaluability fields carry the same discipline as /T. The
proof section walks these in order.

\ruleblock{Lean statement}
\declhead{BernsteinLogic.Soundness}{sound\_psdCertDiag}{\badge{BadgeGreen}{statement bound}}
\begin{leansrc}
/-- PsdCertDiag soundness: the /DD witness warrants A ≻ 0 over ℝ
(`PsdCert.checkDiag_posDef`, the diagonal-dominance congruence). -/
theorem sound_psdCertDiag {H : Hypernet} {seq : Sequent} {a : Nat}
    {A Z : Fin a → Fin a → ℚ} {ts : List Term} {conclusion : Term}
    (hwit : PsdCert.checkDiag a A Z = true)
    (hconcl : PsdCert.conclCheck a A ts conclusion = true)
    (hshape : ∃ s, conclusion.bshape H = some s)
    (hfin : conclusion.finWires H)
    (hwpos : ∀ w ∈ conclusion.freeWires,
      ∃ τ, H.wireType w = some τ ∧ 0 < Shape.totalSize τ.shape)
    (hform : syzygyConclusionForm conclusion seq.post) :
    seq.semanticallyValid H
\end{leansrc}
\par\medskip
\textbf{The dominance core.}\enspace \texttt{posDef\_of\_dominant} is
the self-contained finite-sum argument (no spectral machinery):
off-diagonal mass retreats into its arithmetic-geometric-mean envelope
and symmetry folds the envelope onto the diagonal, leaving
$\sum_i (B_{ii} - \sum_{j\neq i} |B_{ij}|)\,x_i^2 > 0$.
\texttt{posDef\_of\_dd\_witness} composes it with the congruence
transport displayed in the \rulename{PsdCertTri} section;
\texttt{checkDiag\_posDef} bridges the executable check, and the
rational-points corollary feeds the rule lemma.

\declhead{BernsteinLogic.PsdCert}{posDef\_of\_dominant}{\badge{BadgeGreen}{statement bound}}
\begin{leansrc}
/-- **The /DD dominance core**: a symmetric strictly row-diagonally
dominant real matrix is positive definite. Self-contained Finset
algebra — the AM–GM regrouping: off-diagonal mass retreats into
`Σᵢ (Σ_{j≠i} |Bᵢⱼ|)·xᵢ²` by symmetry, leaving `Σᵢ (Bᵢᵢ − Σ_{j≠i}|Bᵢⱼ|)·xᵢ²`,
positive at any `x ≠ 0`. (Positivity of the diagonal is implied:
`Bᵢᵢ > Σ_{j≠i}|Bᵢⱼ| ≥ 0`.) -/
theorem posDef_of_dominant [DecidableEq n] {B : Matrix n n ℝ}
    (hsym : ∀ i j, B i j = B j i)
    (hdom : ∀ i, (∑ j ∈ Finset.univ.erase i, |B i j|) < B i i) :
    B.PosDef
\end{leansrc}
\declhead{BernsteinLogic.PsdCert}{posDef\_of\_dd\_witness}{\badge{BadgeGreen}{statement bound}}
\begin{leansrc}
/-- **PsdCertDiag soundness** (Irving's formalized shape): a symmetric
matrix whose congruence `Zᴴ(AZ)` against an upper-triangular
positive-diagonal witness is strictly row-diagonally dominant is
positive definite. -/
theorem posDef_of_dd_witness [DecidableEq n] {A Z : Matrix n n ℝ}
    (hA : ∀ i j, A i j = A j i)
    (hZu : ∀ ⦃i j⦄, j < i → Z i j = 0)
    (hZd : ∀ i, 0 < Z i i)
    (hdom : ∀ i, (∑ j ∈ Finset.univ.erase i, |(Zᴴ * (A * Z)) i j|)
              < (Zᴴ * (A * Z)) i i) :
    A.PosDef
\end{leansrc}
\declhead{BernsteinLogic.PsdCert}{checkDiag\_posDef}{\badge{BadgeGreen}{statement bound}}
\begin{leansrc}
/-- `checkDiag` warrants positive definiteness of the ℝ-cast. -/
theorem checkDiag_posDef {a : Nat} {A Z : Fin a → Fin a → ℚ}
    (h : checkDiag a A Z = true) : (matR A).PosDef
\end{leansrc}
\declhead{BernsteinLogic.PsdCert}{checkDiag\_quadForm\_nonneg}{\badge{BadgeGreen}{statement bound}}
\begin{leansrc}
/-- /DD's consumable corollary. -/
theorem checkDiag_quadForm_nonneg {a : Nat} {A Z : Fin a → Fin a → ℚ}
    (h : checkDiag a A Z = true) (x : Fin a → ℚ) :
    0 ≤ finQSum a (fun i => finQSum a (fun j => A i j * x i * x j))
\end{leansrc}
\par\medskip
\textbf{The executable check --- what /DD adds.}\enspace Only the
exact congruence entries and the dominance inspection are new; the
shared helpers (\texttt{finQSum}, \texttt{qabs},
\texttt{upperTriPosDiagB}, the conclusion check) are displayed in the
\rulename{PsdCertTri} section.

\declhead{BernsteinLogic.Core.Catalogue}{congEntry}{\badge{BadgeGreen}{definition bound}}
\begin{leansrc}
/-- Entry (i,j) of the congruence B := Zᵀ·(A·Z), recomputed exactly. -/
def congEntry (n : Nat) (A Z : Fin n → Fin n → Rat) (i j : Fin n) : Rat :=
  finQSum n (fun k => Z k i * mulEntry n A Z k j)
\end{leansrc}
\declhead{BernsteinLogic.Core.Catalogue}{checkDiag}{\badge{BadgeGreen}{definition bound}}
\begin{leansrc}
/-- PsdCertDiag witness check (Irving's formalized shape — the tier a
float-LDL ROUNDED witness actually feeds): A symmetric; Z
upper-triangular with positive diagonal; B := ZᵀAZ strictly
row-diagonally dominant (whose positive diagonal is implied: the
off-row sum is ≥ 0). Concludes A ≻ 0. -/
def checkDiag (n : Nat) (A Z : Fin n → Fin n → Rat) : Bool :=
  symB n A && upperTriPosDiagB n Z &&
  ((List.finRange n).all fun i =>
    decide (finQSum n (fun j => if j = i then 0 else qabs (congEntry n A Z i j))
            < congEntry n A Z i i))
\end{leansrc}
\ruleblock{Correspondence (Steele display $\leftrightarrow$ Lean)}
\badge{BadgeGreen}{reviewer-certified against the displayed texts}
\sourcenote{The dictionary below is a reviewer-certified statement of how the mechanized statements above interpret the displayed rule: display metavariable $\mapsto$ Lean binder, each display hypothesis $\mapsto$ the premise or side condition carrying it, the display conclusion $\mapsto$ the goal --- with any deliberate difference in width between display and mechanization named outright. Certification is recorded as a hash triple over (the displayed cell, the extracted statement texts, this dictionary\textquotesingle s own text); if any of the three changes, this block turns amber until a reviewer re-certifies. The correspondence claim itself remains editorial --- reviewer-grade, not machine-checked; its machine-checked successor types each soundness statement against the rule as a structured object, and supersedes this block per rule as it lands.}
Verbatim the \rulename{PsdCertTri} dictionary --- premise-free silence,
the same \texttt{ts}/\texttt{conclusion} binders, the same
\texttt{conclCheck} and \texttt{syzygyConclusionForm} rows, the same
undisplayed evaluability pack (\texttt{hshape}, \texttt{hfin},
\texttt{hwpos}) --- except the witness row: the displayed congruence
$B := Z^{\mathsf{T}} A Z$ strictly row-diagonally dominant $\mapsto$
\texttt{hwit} $:$ \texttt{PsdCert.checkDiag a A Z = true}, entering
through \texttt{checkDiag\_posDef} and consumed as its rational-points
corollary \texttt{checkDiag\_quadForm\_nonneg}. The shared assembly is
one certified lemma, \texttt{sound\_psdCertOf}, visible above ---
the dictionary\textquotesingle s claim that the trio differs only in
the witness is itself a fact of the statements, not editorial.

\ruleblock{Lean verifier receipts}
\sourcenote{module \texttt{BernsteinLogic.Soundness}; bound build \texttt{a20d90e66c13\ldots} (bound from receipt a9c0e34d3e3c9c8c, 2026-08-03T21:30:36Z). A receipt records that named checking layers accepted a build: (module, compiled-artifact SHA-256, environment SHA-256, checker layers), minted when the checking gate ran; the store receipts attesting the bound build are listed with the layers each run credited. Nothing here asks the reader to trust the store: every layer is an independently runnable program over the quoted sources --- rebuild the module and re-run the checkers to reproduce the verdict. Receipts attesting superseded builds credit no layer and are omitted from this rendering (the count of omissions, if any, is noted below).}
\begin{receipttable}
\receiptrow{a9c0e34d3e3c9c8c}{2026-08-03T21:30:36Z}{build+decorrelated}{\badge{BadgeGreen}{current}}
\end{receipttable}
\sourcenote{21 further receipt(s) in the store attest superseded builds of this module and are omitted from this rendering.}
\sourcenote{module \texttt{BernsteinLogic.PsdCert}; bound build \texttt{be59faf6aff0\ldots} (bound from receipt 872a60504989acef, 2026-07-29T07:57:01Z). A receipt records that named checking layers accepted a build: (module, compiled-artifact SHA-256, environment SHA-256, checker layers), minted when the checking gate ran; the store receipts attesting the bound build are listed with the layers each run credited. Nothing here asks the reader to trust the store: every layer is an independently runnable program over the quoted sources --- rebuild the module and re-run the checkers to reproduce the verdict. Receipts attesting superseded builds credit no layer and are omitted from this rendering (the count of omissions, if any, is noted below).}
\begin{receipttable}
\receiptrow{13403ea3a4e42f77}{2026-07-28T23:44:21Z}{build+decorrelated}{\badge{BadgeGreen}{current}}
\receiptrow{6ad2a620d75bd04e}{2026-07-28T23:56:32Z}{build+decorrelated}{\badge{BadgeGreen}{current}}
\receiptrow{d37ce088bc4caf69}{2026-07-29T00:55:46Z}{build+replay+hardened+decorrelated}{\badge{BadgeGreen}{current}}
\receiptrow{872a60504989acef}{2026-07-29T07:57:01Z}{build+replay+hardened+decorrelated}{\badge{BadgeGreen}{current}}
\end{receipttable}
\sourcenote{5 further receipt(s) in the store attest superseded builds of this module and are omitted from this rendering.}
\sourcenote{module \texttt{BernsteinLogic.Core.Catalogue}; bound build \texttt{cbc92e976a1d\ldots} (bound from receipt 5abdb3589707b5df, 2026-08-03T21:30:33Z). A receipt records that named checking layers accepted a build: (module, compiled-artifact SHA-256, environment SHA-256, checker layers), minted when the checking gate ran; the store receipts attesting the bound build are listed with the layers each run credited. Nothing here asks the reader to trust the store: every layer is an independently runnable program over the quoted sources --- rebuild the module and re-run the checkers to reproduce the verdict. Receipts attesting superseded builds credit no layer and are omitted from this rendering (the count of omissions, if any, is noted below).}
\begin{receipttable}
\receiptrow{24804e66edba0bc9}{2026-08-03T20:25:09Z}{build+decorrelated}{\badge{BadgeGreen}{current}}
\receiptrow{5abdb3589707b5df}{2026-08-03T21:30:33Z}{build+decorrelated}{\badge{BadgeGreen}{current}}
\end{receipttable}
\sourcenote{18 further receipt(s) in the store attest superseded builds of this module and are omitted from this rendering.}
\layerline{\layeritem{build}{lake build (elaboration)}{\badge{BadgeGreen}{current}}\layeritem{replay}{leanchecker (official kernel replay)}{\badge{BadgeGreen}{current}}\layeritem{hardened}{safe\_verify (hardened screen)}{\badge{BadgeGreen}{current}}\layeritem{decorrelated}{lean4export + nanoda (decorrelated kernel)}{\badge{BadgeGreen}{current}}}
\sourcenote{\badge{BadgeGray}{comparator: none} no fixture in the three-implementation comparator harness (compiled / kernel-interpreted / reference-interpreted) exercises this rule.}
\ruleblock{HOL4 statement}
\sourcenote{theory \texttt{bern\_tr\_soundpsd}, HOL4 development script \texttt{\texttt{bern\_tr\_soundpsdScript.sml}} \badge{BadgeGreen}{script bound}}
\holhead{Theorem sound\_psdCertDiag}
\begin{holsrc}
Theorem sound_psdCertDiag:
  ∀H sq n A Z ts conclusion.
    psdCheckDD n A Z ∧
    psdConclCheck n A ts conclusion ∧
    (∃s. tmBshape H conclusion = SOME s) ∧
    tmFinWires H conclusion ∧
    (∀w. MEM w (termFreeWires conclusion) ⇒
         ∃τ. H.wireType w = SOME τ ∧ 0 < shapeTotalSize τ.shape) ∧
    syzygyConclusionFormB conclusion sq.post ⇒
    seqSemanticallyValid H sq
\end{holsrc}
\ruleblock{HOL verifier receipts}
\badge{BadgeBlue}{documentary} \texttt{hol-strict} CLEAN: theory \texttt{bern\_tr\_soundpsd}, 4
theorems, 0 cheat-tagged, 0 nonstandard axioms (whitelist:
\texttt{DISK\_THM}; the four \texttt{bool} axioms), 249-theory
ancestry --- verification run of 2026-07-22 on an isolated build host,
against a content-addressed checkout of the exact sources cited here
(git blob \texttt{c01e0b0c\ldots}). The shared semantic spine
(\texttt{sound\_psdCertOf}) is displayed in the /Tri section per the
family-preamble convention; this corollary discharges the
quadratic-form fact via the ported congruence-plus-row-dominance
bridge (\texttt{psdCheckDD\_quadForm}, theory \texttt{bern\_psd},
where the witness check itself --- \texttt{psdCheckDD} --- is ported
with kernel-EVAL receipts). The payload record and its wellformedness
arm ride the ported proof-object stratum (theory
\texttt{bern\_tr\_catdata2}, 272 theorems, 0 offences, same run).

\sourcenote{HOL receipts are documentary: quoted results of \texttt{hol-strict} runs
and census rows from the HOL4 development\textquotesingle s build
evidence, not minted receipt artefacts. A HOL receipt store (the Lean
store\textquotesingle s twin) is named future work.
}
\ruleblock{Proof}

\sourcenote{Genre: admissible rule (tier 2) --- reduction to verified
machinery. The certificate-family preamble --- untrusted witness,
quadratic-form conclusion through the \rulename{Syzygy}
conclusion-form selector, strictness dropped at the formula boundary,
premise-freedom, elementwise batching --- is stated in the
\rulename{PsdCertTri} section (its first recipe) and applies verbatim;
only the witness check differs here.}

\begin{brecipe}[the \rulename{PsdCertDiag} instance]\label{rec:psddd-instance}
Witness: $Z \in \Q^{n\times n}$. Checked premises: $A$ symmetric; $Z$
upper-triangular with positive diagonal; and the congruence
$B := Z^{\mathsf T} A Z$ --- \emph{recomputed exactly by the checker} ---
strictly row-diagonally dominant:
$\sum_{j \neq i} |B_{ij}| < B_{ii}$ for every $i$. (Positivity of the
diagonal is implied: the off-row sum is $\geq 0$.) Then
\[
\inferrule*[Right=\rulename{PsdCertDiag}]
  {\;\mathrm{checkDiag}(A, Z) = \mathsf{true} \\
    c \equiv -\textstyle\sum_{i,j} A_{ij}\,p_i\,p_j\;}
  {\Phi \vdash \Psi}
\qquad\rightsquigarrow\qquad
\text{\upshape discharged by } \texttt{sound\_psdCertDiag}.
\]
The matrix-witness check costs $O(\alpha n^2 + n^3)$ multiplications
($\alpha$ the average nonzeros per row of $A$, as at /T)
with the product $AZ$ materialized once (the $Z^{\mathsf T}(AZ)$
congruence is the $n^3$ term; the reference check displayed below
inlines the product entrywise for readability --- it is the
specification, not the cost model) on integers of \emph{bounded}
width --- input bits plus twice the witness bits plus $O(\log n)$ for
the nested sums --- embarrassingly parallel, still division-free.
\end{brecipe}

\begin{bprop}[matrix warrant]\label{prop:psddd-matrix}
A symmetric real matrix $B$ with
$\sum_{j\neq i}|B_{ij}| < B_{ii}$ for all $i$ is positive definite;
hence if additionally $Z$ is upper-triangular with positive diagonal
and $B = Z^{\mathsf T} A Z$ for symmetric $A$, then $A \succ 0$.
\end{bprop}
\begin{proof}
For $x \neq 0$, split the quadratic form and bound each off-diagonal
term by its arithmetic--geometric-mean envelope,
$x_i B_{ij} x_j \geq -|B_{ij}|\,(x_i^2 + x_j^2)/2$; folding the
envelope's two symmetric halves by $B_{ij} = B_{ji}$ gives
\[
x^{\mathsf T} B\, x \;\geq\; \sum_i \Bigl(B_{ii} -
\textstyle\sum_{j \neq i} |B_{ij}|\Bigr)\, x_i^2 \;>\; 0 .
\]
(The kernel proof is exactly this self-contained finite-sum argument
--- \texttt{posDef\_of\_dominant} --- with no spectral machinery. The
classical reading: strict dominance with positive diagonal places
every Gershgorin disc in the open right half-plane, and symmetry makes
the spectrum real, hence positive
\cite{gershgorin1931,taussky1949,hornjohnson2013}.)
Transport through the invertible $Z$ as in the /T proof
(\texttt{posDef\_of\_congruence\_posDef}) yields $A \succ 0$.
\end{proof}

\begin{bremark}[what each hypothesis is doing]
\emph{Symmetry of $A$} makes $B$ symmetric by construction, which the
envelope-folding step consumes ($B_{ij} = B_{ji}$) --- row dominance of
a nonsymmetric matrix bounds neither the symmetric part's rows nor the
form. \emph{Triangularity + positive diagonal of $Z$} is, as in /T,
only the cheapest syntactic invertibility certificate --- invertibility
is all the transport step needs, and triangularity is what the natural
producer emits anyway. \emph{Strict dominance} is the content of the strict tier:
weakening it to non-strict admits singular $B$ (e.g.\ the graph
Laplacian of a path: weakly dominant, kernel = constants), and the
matrix warrant $A \succ 0$ would overclaim --- the posted
$\succcurlyeq$-shaped atom would in fact survive, but the strict
theorem is what \rulename{PsdCertKernel} consumes on its inner block to
confine kernels, and that consumer admits no slack. \emph{Exact recomputation of $B$}: the
congruence is never taken from the certificate --- a wire-supplied $B$
could be dominant while unrelated to $A$.
\end{bremark}

\begin{bremark}[rounding tolerance --- why this is the production tier]
Feed $Z := \lceil 2^s R^{-1} \rfloor$ from an off-stage float Cholesky
$A \approx R^{\mathsf T} R$: writing $Z = 2^s R^{-1} + E$ with rounding
$|E_{ij}| \leq \tfrac12$, exact expansion (against an exact
factorization) gives $B = Z^{\mathsf T} A Z = 4^s I + 2^s\bigl(RE +
(RE)^{\mathsf T}\bigr) + E^{\mathsf T} A E$ --- entrywise contamination
$O(2^s n \|R\|_{\max} + n^2 \|A\|_{\max})$, off-row dominance sums one
factor of $n$ larger --- plus the image of the float residual
$\Delta := A - R^{\mathsf T} R$ --- the term $Z^{\mathsf T} \Delta Z$,
of normalized leading size $R^{-\mathsf T} \Delta R^{-1}$ --- whose
control is a mantissa requirement of
order $\log_2 \kappa(A) + O(\log n)$ bits on the factorization (the
dimension guard bits of the standard backward-error bound), not on the
witness. Strict dominance therefore holds once $2^s$ clears a
polynomial in $n$ times the input scales: a scale exponent $s =
O(\log n +
\log \|R\|_{\max})$ suffices (sharp constants in the cited
certificate analysis), so the witness \emph{scale} grows only
\emph{logarithmically} in dimension and scale and never tracks an
accuracy target; the witness \emph{entries} encode in $s +
O(\log^{+} \|R^{-1}\|_{\max})$ bits --- the second term is the
factorization's own magnitude ($Z$ rounds $2^s R^{-1}$), present at
any scale, and no accuracy target either. Contrast /T,
whose exact witness inherits minor-scale bit growth. What a too-coarse
$Z$ costs is only completeness margin: the check fails closed, and the
producer retries at a finer scale with soundness never in play. This
rounded-integer-witness shape is the one Irving analysed and formalized
at scale \cite{irving2026certificate,irvingaks}; the underlying
tolerance phenomenon is the factor-width/scaled-diagonal-dominance
circle of ideas \cite{boman2005} --- context, not a characterization:
that theorem equates factor-width two with scaled diagonal dominance
of $A$ \emph{itself} (diagonal scaling only), while /DD checks
dominance of the congruence target $B = Z^{\mathsf T} A Z$ for
arbitrary triangular $Z$, and so accepts matrices far outside the SDD
class --- with an exact LDL witness the congruence target is diagonal,
so /DD's exact-witness reach is \emph{every} rational SPD matrix,
while e.g.\ the $3 \times 3$ equicorrelation matrix with off-diagonal
$\rho \in (\tfrac12, 1)$ is SPD but not SDD under any positive
diagonal scaling.
\end{bremark}

\begin{bremark}[the /T---/DD complementarity, stated once]
Same witness idea, dual regimes: /T is exact-witness,
$O(\alpha n^2)$, and wins at moderate sizes and exact-by-construction
witnesses; /DD tolerates any rounding at $O(n^3)$ bounded-width
parallel arithmetic and wins at scale. A producer holding a float
factorization reaches for /DD by default and /T when it can clear
denominators cheaply. Both refuse indefinite $A$ against every witness.
\end{bremark}

\begin{anchors}
Recipe \ref{rec:psddd-instance} warrant & \texttt{sound\_psdCertDiag} (Soundness); shared transport \texttt{sound\_psdCertOf} (see \rulename{PsdCertTri})\\
dominance core (Prop.\ \ref{prop:psddd-matrix}) & \texttt{PsdCert.posDef\_of\_dominant}; assembled in \texttt{posDef\_of\_dd\_witness}\\
check-to-theorem bridge & \texttt{PsdCert.checkDiag\_posDef}; rational-points corollary \texttt{checkDiag\_quadForm\_nonneg}\\
the executable check & \texttt{PsdCert.checkDiag} with the exact congruence entries \texttt{congEntry}\\
conclusion identity + readings & \texttt{PsdCert.conclCheck}; the \rulename{Syzygy} section's conclusion-form selector\\
\end{anchors}


\rulesection{30}{PsdCertKernel}{admissible rule}
\label{blcatrule-PsdCertKernel}
\ruleblock{The rule (Steele metanotation)}
\sourcenote{as realized in the mechanized development; source row: \texttt{\texttt{rules.toml}} --- \emph{PsdCertKernel}. The metanotation statement is typeset by hand for notational quality; its mechanized ground truth is the Lean statement below, extracted from the sources at every build of this book and hash-bound. The rule \emph{roster} is machine-checked (every mechanized rule tag round-trips to exactly one catalogue row); the statement \emph{text} is editorial, checked by eye against the extracted Lean beside it.}
\begin{mathpar}
\inferrule*[Right=\rulename{PsdCertKernel}]{
  A \in \Q^{(m+k) \times (m+k)}, \ \text{kernel coordinates last}
  \\
  L := (I \mid N), \ N \in \Q^{m \times k}
  \\
  A = L^{\mathsf{T}} B L \ \text{entrywise exactly}
  \\
  B \in \Q^{m \times m} \ \text{with a /T-or-/DD inner witness } Z
  \\
  p_1, \dots, p_{m+k} \ \text{payload terms}
  \\
  c \equiv -\textstyle\sum_{i,j} A_{ij}\, p_i\, p_j \ \text{(exact canonical polynomial)}
  \\
  \Psi \in \bigl\{\, c \leq 0, \;\; Q \geq 0 \text{ where } c = -Q, \;\; l \leq r \text{ where } c = l - r \,\bigr\}
}{
  \Phi \vdash \Psi
}
\end{mathpar}

\par\medskip
\rulename{PsdCertKernel} is the semidefinite member of the certificate
family (family discipline in the \rulename{PsdCertTri} section): it
certifies $A \succcurlyeq 0$ \emph{with an exactly-checked kernel},
for the boundary cases the two strict tiers cannot touch --- energy
forms with conservation laws, graph Laplacians with their all-ones
kernel, sharp certificates whose slack vanishes identically. Data:
$A$ over $\Q^{(m+k)\times(m+k)}$ with the $k$ kernel coordinates last
(the producer applies its permutation before emission --- conjugation
is index renaming, so the identity-permutation form loses no
generality); a rational kernel description $N$; and a leading block
$B$ with a /T-or-/DD inner witness. The checked premise is a single
entrywise-exact factorization: $A = L^{\mathsf T} B\, L$ for the
kernel-completion rows $L = (I_m \mid N)$ --- equivalent, given the
inner check\textquotesingle s own symmetry verification of $B$, to
($A$ symmetric $\wedge$ $A \cdot K = 0$ exactly for
$K = (-N;\, I_k)$ $\wedge$ leading block $= B$), the cross and
trailing blocks being forced either way, so a forged $N$ or $B$
refuses --- plus the strict inner check on $B$. Rectangular congruence
then gives $A \succcurlyeq 0$ with the form vanishing exactly on the
span of $K$: the certificate fixes kernel and rank, not just
semidefiniteness.

The kernel data must be exact by necessity, not by preference: the
semidefinite-with-kernel matrices form a boundary set with empty
interior, so no rounding-tolerant witness for kernel directions can
exist --- any check tolerant to kernel perturbation also accepts
slightly-indefinite matrices. Exact kernel plus rounding-tolerant bulk
(the /DD inner witness) is therefore the forced split --- tolerance
can live nowhere else --- and it
matches the producers: kernels are semantic (theorems supply $N$),
floats only ever touch the bulk witness $Z$. Corners are structural:
$k = 0$ degenerates to the inner strict tier; $m = 0$ forces $A = 0$,
where the factorization premise is the entire content.

What breaks without each hypothesis: an approximate factorization
certifies nothing (the boundary argument above); a merely semidefinite
inner block would hide undeclared kernel directions inside $B$; the
kernel-coordinates-last convention is pure bookkeeping. What reaches
the formula layer is the quadratic-form conclusion --- honest
for this rule, since it is $\succcurlyeq$-shaped --- while the checked
kernel/rank content becomes formula-expressible only with closed
tensor terms (the extension the \rulename{PConstExact} section
records); it is certificate-level assurance meanwhile, deliberately
not hand-faked into the formula grammar. The proof section walks these
in order.

\ruleblock{Lean statement}
\declhead{BernsteinLogic.Soundness}{sound\_psdCertKernel}{\badge{BadgeGreen}{statement bound}}
\begin{leansrc}
/-- PsdCertKernel soundness: the /K witness warrants A ⪰ 0 over ℝ
(`PsdCert.checkKernel_posSemidef`, the kernel-completion congruence). -/
theorem sound_psdCertKernel {H : Hypernet} {seq : Sequent} {m k : Nat}
    {A : Fin (m + k) → Fin (m + k) → ℚ} {N : Fin m → Fin k → ℚ}
    {B : Fin m → Fin m → ℚ} {innerT : Bool} {Z : Fin m → Fin m → ℚ}
    {ts : List Term} {conclusion : Term}
    (hwit : PsdCert.checkKernel m k A N B innerT Z = true)
    (hconcl : PsdCert.conclCheck (m + k) A ts conclusion = true)
    (hshape : ∃ s, conclusion.bshape H = some s)
    (hfin : conclusion.finWires H)
    (hwpos : ∀ w ∈ conclusion.freeWires,
      ∃ τ, H.wireType w = some τ ∧ 0 < Shape.totalSize τ.shape)
    (hform : syzygyConclusionForm conclusion seq.post) :
    seq.semanticallyValid H
\end{leansrc}
\par\medskip
\textbf{The rectangular congruence.}\enspace
\texttt{posSemidef\_of\_factor} is one line --- semidefiniteness
transports through ANY rectangular congruence, no invertibility
anywhere --- which is exactly why the factorization premise can carry
the whole rule. \texttt{checkKernel\_posSemidef} bridges the executable
check (selecting the /T or /DD theorem for the inner block), and the
rational-points corollary feeds the rule lemma.

\declhead{BernsteinLogic.PsdCert}{posSemidef\_of\_factor}{\badge{BadgeGreen}{statement bound}}
\begin{leansrc}
/-- **PsdCertKernel soundness core**: a matrix exhibited as the two-sided
congruence `LᴴBL` of a positive-semidefinite block is positive
semidefinite (rectangular congruence; no invertibility anywhere). -/
theorem posSemidef_of_factor {m' n' : Type*} [Fintype m'] [Fintype n']
    {A : Matrix n' n' ℝ} {B : Matrix m' m' ℝ} {L : Matrix m' n' ℝ}
    (hB : B.PosSemidef) (hA : A = Lᴴ * B * L) : A.PosSemidef
\end{leansrc}
\declhead{BernsteinLogic.PsdCert}{checkKernel\_posSemidef}{\badge{BadgeGreen}{statement bound}}
\begin{leansrc}
/-- `checkKernel` warrants positive semidefiniteness of the ℝ-cast. -/
theorem checkKernel_posSemidef {m k : Nat} {A : Fin (m + k) → Fin (m + k) → ℚ}
    {N : Fin m → Fin k → ℚ} {B : Fin m → Fin m → ℚ} {innerT : Bool}
    {Z : Fin m → Fin m → ℚ}
    (h : checkKernel m k A N B innerT Z = true) : (matR A).PosSemidef
\end{leansrc}
\declhead{BernsteinLogic.PsdCert}{checkKernel\_quadForm\_nonneg}{\badge{BadgeGreen}{statement bound}}
\begin{leansrc}
/-- /K's consumable corollary. -/
theorem checkKernel_quadForm_nonneg {m k : Nat} {A : Fin (m + k) → Fin (m + k) → ℚ}
    {N : Fin m → Fin k → ℚ} {B : Fin m → Fin m → ℚ} {innerT : Bool}
    {Z : Fin m → Fin m → ℚ}
    (h : checkKernel m k A N B innerT Z = true) (x : Fin (m + k) → ℚ) :
    0 ≤ finQSum (m + k) (fun i => finQSum (m + k) (fun j => A i j * x i * x j))
\end{leansrc}
\par\medskip
\textbf{The executable check --- what /K adds.}\enspace The
kernel-completion rows and the entrywise factorization check; the
inner witness reuses \texttt{checkTri}/\texttt{checkDiag} and the
conclusion check is shared (both displayed in the
\rulename{PsdCertTri} section).

\declhead{BernsteinLogic.Core.Catalogue}{kernelExt}{\badge{BadgeGreen}{definition bound}}
\begin{leansrc}
/-- The kernel-completion row matrix L = (I | N) : m × (m+k). -/
def kernelExt (m k : Nat) (N : Fin m → Fin k → Rat) (a : Fin m)
    (j : Fin (m + k)) : Rat :=
  if h : j.val < m then (if a.val = j.val then 1 else 0)
  else N a ⟨j.val - m, by have := j.isLt; omega⟩
\end{leansrc}
\declhead{BernsteinLogic.Core.Catalogue}{checkKernel}{\badge{BadgeGreen}{definition bound}}
\begin{leansrc}
/-- PsdCertKernel witness check (PSD-with-kernel companion): A — over
`Fin (m+k)`, the kernel coordinates LAST — IS the two-sided
kernel-completion congruence LᵀBL of its leading block B, entrywise
exactly, plus a T-or-DD inner witness for B. Equivalent to the note's §1
premises (A symmetric ∧ A·K = 0 exactly for K = (−N; Iₖ) ∧ leading
block = B): C = BN and D = NᵀBN are forced either way, and a forged N
or B refuses. The producer applies its permutation P before emission —
P-conjugation is index renaming on matrix and term vector alike, so the
catalogue rule is the P = id core. Corners are structural: k = 0 gives
L = I (A = B); m = 0 forces A = 0 (empty B, no inner witness content). -/
def checkKernel (m k : Nat) (A : Fin (m + k) → Fin (m + k) → Rat)
    (N : Fin m → Fin k → Rat) (B : Fin m → Fin m → Rat)
    (innerT : Bool) (Z : Fin m → Fin m → Rat) : Bool :=
  ((List.finRange (m + k)).all fun i => (List.finRange (m + k)).all fun j =>
    decide (A i j = qlsum ((List.finRange m).map fun a =>
      qlsum ((List.finRange m).map fun b =>
        kernelExt m k N a i * B a b * kernelExt m k N b j)))) &&
  (if innerT then checkTri m B Z else checkDiag m B Z)
\end{leansrc}
\ruleblock{Correspondence (Steele display $\leftrightarrow$ Lean)}
\badge{BadgeGreen}{reviewer-certified against the displayed texts}
\sourcenote{The dictionary below is a reviewer-certified statement of how the mechanized statements above interpret the displayed rule: display metavariable $\mapsto$ Lean binder, each display hypothesis $\mapsto$ the premise or side condition carrying it, the display conclusion $\mapsto$ the goal --- with any deliberate difference in width between display and mechanization named outright. Certification is recorded as a hash triple over (the displayed cell, the extracted statement texts, this dictionary\textquotesingle s own text); if any of the three changes, this block turns amber until a reviewer re-certifies. The correspondence claim itself remains editorial --- reviewer-grade, not machine-checked; its machine-checked successor types each soundness statement against the rule as a structured object, and supersedes this block per rule as it lands.}
As at \rulename{PsdCertTri} --- premise-free silence, the shared
\texttt{conclCheck} row at size $m + k$, the same
\texttt{syzygyConclusionForm} selector and undisplayed evaluability
pack --- with the witness row exchanged for the kernel-completion
data:
\begin{itemize}
\item the block sizes $\mapsto$ \texttt{m k : Nat} with $A$ $:$
  \texttt{Fin (m + k) → Fin (m + k) → ℚ} (kernel coordinates last,
  as displayed --- the producer\textquotesingle s permutation is
  applied before emission, so the certified rule is the $P =
  \mathrm{id}$ core of the roster note); $N \mapsto$ \texttt{N : Fin
  m → Fin k → ℚ}; $B \mapsto$ \texttt{B}; the ``/T-or-/DD inner
  witness\textquotesingle\textquotesingle{} $\mapsto$ the pair
  \texttt{innerT : Bool} (which inner check) and \texttt{Z : Fin m →
  Fin m → ℚ} (its witness).
\item the displayed factorization row $A = L^{\mathsf{T}} B L$
  entrywise exactly, $L = (I \mid N)$, together with the inner
  witness $\mapsto$ the single Boolean \texttt{hwit} $:$
  \texttt{PsdCert.checkKernel m k A N B innerT Z = true}, entering
  through \texttt{checkKernel\_posSemidef} ($A \succcurlyeq 0$ --- the
  trio\textquotesingle s one genuinely $\succcurlyeq$-shaped member,
  so here the formula boundary loses nothing) and consumed as
  \texttt{checkKernel\_quadForm\_nonneg}.
\end{itemize}
Width note, beyond /T\textquotesingle s: the checked kernel and rank
content (the roster note\textquotesingle s \S1 equivalents) is
consumed by the matrix-side bridge but never surfaces in the sequent
--- the certified formula-level claim is exactly the quadratic-form
conclusion, neither more nor less, which is the honest reading the
commentary states and this triple pins.

\ruleblock{Lean verifier receipts}
\sourcenote{module \texttt{BernsteinLogic.Soundness}; bound build \texttt{a20d90e66c13\ldots} (bound from receipt a9c0e34d3e3c9c8c, 2026-08-03T21:30:36Z). A receipt records that named checking layers accepted a build: (module, compiled-artifact SHA-256, environment SHA-256, checker layers), minted when the checking gate ran; the store receipts attesting the bound build are listed with the layers each run credited. Nothing here asks the reader to trust the store: every layer is an independently runnable program over the quoted sources --- rebuild the module and re-run the checkers to reproduce the verdict. Receipts attesting superseded builds credit no layer and are omitted from this rendering (the count of omissions, if any, is noted below).}
\begin{receipttable}
\receiptrow{a9c0e34d3e3c9c8c}{2026-08-03T21:30:36Z}{build+decorrelated}{\badge{BadgeGreen}{current}}
\end{receipttable}
\sourcenote{21 further receipt(s) in the store attest superseded builds of this module and are omitted from this rendering.}
\sourcenote{module \texttt{BernsteinLogic.PsdCert}; bound build \texttt{be59faf6aff0\ldots} (bound from receipt 872a60504989acef, 2026-07-29T07:57:01Z). A receipt records that named checking layers accepted a build: (module, compiled-artifact SHA-256, environment SHA-256, checker layers), minted when the checking gate ran; the store receipts attesting the bound build are listed with the layers each run credited. Nothing here asks the reader to trust the store: every layer is an independently runnable program over the quoted sources --- rebuild the module and re-run the checkers to reproduce the verdict. Receipts attesting superseded builds credit no layer and are omitted from this rendering (the count of omissions, if any, is noted below).}
\begin{receipttable}
\receiptrow{13403ea3a4e42f77}{2026-07-28T23:44:21Z}{build+decorrelated}{\badge{BadgeGreen}{current}}
\receiptrow{6ad2a620d75bd04e}{2026-07-28T23:56:32Z}{build+decorrelated}{\badge{BadgeGreen}{current}}
\receiptrow{d37ce088bc4caf69}{2026-07-29T00:55:46Z}{build+replay+hardened+decorrelated}{\badge{BadgeGreen}{current}}
\receiptrow{872a60504989acef}{2026-07-29T07:57:01Z}{build+replay+hardened+decorrelated}{\badge{BadgeGreen}{current}}
\end{receipttable}
\sourcenote{5 further receipt(s) in the store attest superseded builds of this module and are omitted from this rendering.}
\sourcenote{module \texttt{BernsteinLogic.Core.Catalogue}; bound build \texttt{cbc92e976a1d\ldots} (bound from receipt 5abdb3589707b5df, 2026-08-03T21:30:33Z). A receipt records that named checking layers accepted a build: (module, compiled-artifact SHA-256, environment SHA-256, checker layers), minted when the checking gate ran; the store receipts attesting the bound build are listed with the layers each run credited. Nothing here asks the reader to trust the store: every layer is an independently runnable program over the quoted sources --- rebuild the module and re-run the checkers to reproduce the verdict. Receipts attesting superseded builds credit no layer and are omitted from this rendering (the count of omissions, if any, is noted below).}
\begin{receipttable}
\receiptrow{24804e66edba0bc9}{2026-08-03T20:25:09Z}{build+decorrelated}{\badge{BadgeGreen}{current}}
\receiptrow{5abdb3589707b5df}{2026-08-03T21:30:33Z}{build+decorrelated}{\badge{BadgeGreen}{current}}
\end{receipttable}
\sourcenote{18 further receipt(s) in the store attest superseded builds of this module and are omitted from this rendering.}
\layerline{\layeritem{build}{lake build (elaboration)}{\badge{BadgeGreen}{current}}\layeritem{replay}{leanchecker (official kernel replay)}{\badge{BadgeGreen}{current}}\layeritem{hardened}{safe\_verify (hardened screen)}{\badge{BadgeGreen}{current}}\layeritem{decorrelated}{lean4export + nanoda (decorrelated kernel)}{\badge{BadgeGreen}{current}}}
\sourcenote{\badge{BadgeGray}{comparator: none} no fixture in the three-implementation comparator harness (compiled / kernel-interpreted / reference-interpreted) exercises this rule.}
\ruleblock{HOL4 statement}
\sourcenote{theory \texttt{bern\_tr\_soundpsd}, HOL4 development script \texttt{\texttt{bern\_tr\_soundpsdScript.sml}} \badge{BadgeGreen}{script bound}}
\holhead{Theorem sound\_psdCertKernel}
\begin{holsrc}
Theorem sound_psdCertKernel:
  ∀H sq m k A N B innerT Z ts conclusion.
    psdCheckK m k A N B innerT Z ∧
    psdConclCheck (m + k) A ts conclusion ∧
    (∃s. tmBshape H conclusion = SOME s) ∧
    tmFinWires H conclusion ∧
    (∀w. MEM w (termFreeWires conclusion) ⇒
         ∃τ. H.wireType w = SOME τ ∧ 0 < shapeTotalSize τ.shape) ∧
    syzygyConclusionFormB conclusion sq.post ⇒
    seqSemanticallyValid H sq
\end{holsrc}
\ruleblock{HOL verifier receipts}
\badge{BadgeBlue}{documentary} \texttt{hol-strict} CLEAN: theory \texttt{bern\_tr\_soundpsd}, 4
theorems, 0 cheat-tagged, 0 nonstandard axioms (whitelist:
\texttt{DISK\_THM}; the four \texttt{bool} axioms), 249-theory
ancestry --- verification run of 2026-07-22 on an isolated build host,
against a content-addressed checkout of the exact sources cited here
(git blob \texttt{c01e0b0c\ldots}). The shared semantic spine
(\texttt{sound\_psdCertOf}) is displayed in the /Tri section per the
family-preamble convention; this corollary discharges the
quadratic-form fact via the ported kernel-extension bridge
(\texttt{psdCheckK\_quadForm}, theory \texttt{bern\_psd}, where the
witness check itself --- \texttt{psdCheckK}, over the extension
\texttt{kernelExt} --- is ported with kernel-EVAL receipts). The
payload record and its wellformedness arm ride the ported proof-object
stratum (theory \texttt{bern\_tr\_catdata2}, 272 theorems, 0 offences,
same run).

\sourcenote{HOL receipts are documentary: quoted results of \texttt{hol-strict} runs
and census rows from the HOL4 development\textquotesingle s build
evidence, not minted receipt artefacts. A HOL receipt store (the Lean
store\textquotesingle s twin) is named future work.
}
\ruleblock{Proof}

\sourcenote{Genre: admissible rule (tier 2) --- reduction to verified
machinery. The certificate-family preamble is stated in the
\rulename{PsdCertTri} section and applies verbatim. This is the
\emph{semidefinite} member: it certifies $A \succcurlyeq 0$ together
with an exactly-checked kernel, for the boundary cases the two strict
tiers cannot touch.}

\begin{brecipe}[the \rulename{PsdCertKernel} instance]\label{rec:psdk-instance}
Data: $A \in \Q^{(m+k)\times(m+k)}$ with the $k$ kernel coordinates
\emph{last} (the producer applies its permutation before emission ---
conjugation by a permutation is index renaming on matrix and term
vector alike, so the identity-permutation form loses no generality);
$N \in \Q^{m \times k}$; the kernel-completion row matrix
$L := (\,I_m \mid N\,)$; a leading block $B \in \Q^{m\times m}$ with a
/T-or-/DD inner witness $Z$. Checked premises: the factorization
$A = L^{\mathsf T} B\, L$ holds \emph{entrywise exactly} (recomputed),
and $B$ passes the selected strict inner check. Then
\[
\inferrule*[Right=\rulename{PsdCertKernel}]
  {\;\mathrm{checkKernel}(A, N, B, Z) = \mathsf{true} \\
    c \equiv -\textstyle\sum_{i,j} A_{ij}\,p_i\,p_j\;}
  {\Phi \vdash \Psi}
\qquad\rightsquigarrow\qquad
\text{\upshape discharged by } \texttt{sound\_psdCertKernel}.
\]
The factorization check is --- given the symmetry of $B$, which the
selected inner check itself verifies --- equivalent to the three
separate conditions ($A$ symmetric; $A \cdot K = 0$ exactly for
$K = \binom{-N}{\,I_k}$; leading block of $A$ equal to $B$): forward,
$B$'s symmetry makes $A = L^{\mathsf T} B L$ symmetric and $LK = 0$
kills $AK$; backward, $AK = 0$ forces the cross block $BN$ and then
the trailing block $N^{\mathsf T} B N$. Either way the off-leading
blocks are forced, so a forged $N$ or $B$ refuses.
\end{brecipe}

\begin{bprop}[matrix warrant]\label{prop:psdk-matrix}
If $B \succ 0$ (hence $\succcurlyeq 0$) and
$A = L^{\mathsf T} B L$ for the rectangular $L = (I \mid N)$, then
$A \succcurlyeq 0$. Moreover $x^{\mathsf T} A x = 0$ exactly on
$\operatorname{span} K$, so the certificate fixes
$\ker A = \operatorname{span} K$ and
$\operatorname{rank} A = m$.
\end{bprop}
\begin{proof}
Positive semidefiniteness is rectangular congruence, with no
invertibility anywhere: $x^{\mathsf T} A x = (Lx)^{\mathsf T} B\,(Lx)
\geq 0$ (the kernel proof is this one line,
\texttt{posSemidef\_of\_factor}). For the kernel content, write
$x = (x_1, x_2)$ against the block split and note
$Lx = x_1 + N x_2$, while $L K = -N + N = 0$; with $B \succ 0$,
$x^{\mathsf T} A x = 0$ iff $Lx = 0$ iff $x = K x_2$ --- i.e.\ the
zero set of the form is exactly $\operatorname{span} K$. That zero set
\emph{is} the kernel: if $Lx = 0$ then $Ax = L^{\mathsf T} B (Lx) =
0$, and conversely $Ax = 0$ forces $x^{\mathsf T} A x = 0$, hence
$Lx = 0$. The columns of $K$ are independent (its lower block is
$I_k$), so $\ker A = \operatorname{span} K$ has dimension $k$ and
$\operatorname{rank} A = m$ by rank--nullity.
\end{proof}

\begin{bremark}[exact kernels are forced, not a method limitation]
A rounding-tolerant witness for kernel \emph{directions} cannot exist
for a residual-style check --- one attesting the kernel claim only up
to an $\varepsilon$-residual, with no exact kernel equation:
$\{A \succcurlyeq 0,\ \det A = 0\}$ is a boundary set with empty
interior, so a check whose acceptance region contains an
entrywise-$\varepsilon$ neighbourhood of a true instance also accepts
some $A - \varepsilon' I \not\succcurlyeq 0$. The split is therefore
forced, not preferred: exact kernel data plus a rounding-tolerant bulk
(the /DD inner witness), with tolerance able to live in the bulk and
nowhere else --- and it matches the intended
producers, whose kernels are \emph{semantic} (a conservation law, a
syzygy relation, the all-ones vector of a graph Laplacian): $N$ comes
from theorems, and floats only ever touch $Z$.
\end{bremark}

\begin{bremark}[each hypothesis, and the corners]
\emph{Entrywise-exact factorization}: the whole kernel claim rides on
$A K = 0$ holding exactly; an $\varepsilon$-approximate factorization
certifies nothing (previous remark). \emph{The strict inner witness}:
$B \succ 0$ is what confines the zero set to
$\operatorname{span} K$ --- a merely semidefinite $B$ would leave
undeclared kernel directions inside the leading block. \emph{Kernel
coordinates last}: bookkeeping only, by the permutation argument in
the recipe. Corners are structural, not special-cased: $k = 0$ gives
$L = I$ and the rule degenerates to its inner strict tier ($A = B$);
$m = 0$ forces $A = 0$ (empty $B$, no inner witness content --- the
factorization premise \emph{is} the entire content).
\end{bremark}

\begin{bremark}[what reaches the formula layer]
The checker verifies the kernel and rank content (a forged $N$
refuses), but the formula language states only the
quadratic-form conclusion $\sum_{i,j} A_{ij}\,p_i\,p_j \geq 0$ ---
which for this rule is honest ($\succcurlyeq$-shaped, per the
preamble's strictness discipline) and loses only the kernel/rank
\emph{statement}, not its checking. Making rank and kernel
formula-expressible needs closed-tensor terms and a tensor-valued
closed evaluator --- the same named extension the
\rulename{PConstExact} section records for exact matrix identities;
short of that extension, the kernel content is certificate-level
assurance
surfaced here, deliberately not hand-faked into the formula grammar.
\end{bremark}

\begin{anchors}
Recipe \ref{rec:psdk-instance} warrant & \texttt{sound\_psdCertKernel} (Soundness); shared transport \texttt{sound\_psdCertOf} (see \rulename{PsdCertTri})\\
rectangular congruence (Prop.\ \ref{prop:psdk-matrix}) & \texttt{PsdCert.posSemidef\_of\_factor}\\
check-to-theorem bridge & \texttt{PsdCert.checkKernel\_posSemidef} (inner tier via \texttt{checkTri\_posDef} / \texttt{checkDiag\_posDef}); corollary \texttt{checkKernel\_quadForm\_nonneg}\\
the executable check & \texttt{PsdCert.checkKernel} with the completion rows \texttt{kernelExt}\\
conclusion identity + readings & \texttt{PsdCert.conclCheck}; the \rulename{Syzygy} section's conclusion-form selector\\
\end{anchors}


\rulesection{31}{ReLU-Split}{admissible rule}
\label{blcatrule-ReLUSplit}
\ruleblock{The rule (Steele metanotation)}
\sourcenote{as realized in the mechanized development; source row: \texttt{\texttt{rules.toml}} --- \emph{ReLU-Split}. The metanotation statement is typeset by hand for notational quality; its mechanized ground truth is the Lean statement below, extracted from the sources at every build of this book and hash-bound. The rule \emph{roster} is machine-checked (every mechanized rule tag round-trips to exactly one catalogue row); the statement \emph{text} is editorial, checked by eye against the extracted Lean beside it.}
\begin{mathpar}
\inferrule*[Right=\rulename{ReLU-Split}]{
  \Phi \wedge (h \geq 0) \vdash \Psi\sub{b}{h} \\
  \Phi \wedge (h \leq 0) \vdash \Psi\sub{b}{0} \\
  \ReLU_{\{h\}\to\{y\}} \in H \quad (y \text{ typed at rank } 0)
}{
  \Phi \vdash \Psi
}
\end{mathpar}

\par\medskip
\rulename{ReLU-Split} is the case-split rule of the ReLU family: from
two branch proofs — the goal with the output position replaced by $h$,
under $\Phi$ strengthened by $h \geq 0$; and the goal with the output
position replaced by $0$, under $\Phi$ strengthened by $h \leq 0$ —
conclude the goal outright, per rank-0 node: the tensor-shaped reading
of the rule is \emph{false}, not merely unrealized (the proof
section\textquotesingle s scope remark carries the mixed-sign
countermodel), so a producer splits scalar by scalar. It is the \emph{complete}-verification
move of the neural-network literature (the Reluplex and
branch-and-bound lineage): where the affine envelope of
\rulename{LiRPA} loses too much, a producer splits on the
neuron\textquotesingle s sign and pays with proof size — each split
doubles the branch count, which is exactly the branch-and-bound
tradeoff.

Unlike every other rule in this book so far, \rulename{ReLU-Split} has
\emph{no certificate node}: no payload constructor, no wire tag — a
certificate stream never contains one, so there is nothing for this
rule\textquotesingle s checker column to recompute. A producer
elaborates it away. The elaboration\textquotesingle s target is the
\emph{derivation} the rule abbreviates (case analysis over the
sign-split fact, equality transport in each branch, and structural
plumbing), whose nodes the checker checks by their own rules — a
certificate fragment that becomes emittable exactly when the
context-carrying case split it opens with lands as a core constructor
(today that form is a satellite soundness theorem only; the proof
section states the dependency). What stands unconditionally is the
rule-level semantic theorem of the proof
section\textquotesingle s mechanization remark.

Each ingredient of the two premises is load-bearing for the
elaboration. The branch strengthenings ($h \geq 0$, $h \leq 0$) are
what the case fact discharges into each branch context — remove one
and that branch proof has no purchase on the sign information. The
substitution alignment ($\Psi$ with the output position filled by $h$,
respectively $0$) is what the equality transport rewrites through:
inside the active branch the node\textquotesingle s fact $y = h$
carries $\Psi[b \coloneqq h]$ to $\Psi$, and inside the inactive
branch $y = 0$ carries $\Psi[b \coloneqq 0]$ to $\Psi$.

\cataloguenote{The checker column remains empty for the structural reason: no
certificate stream ever contains a \rulename{ReLU-Split} node, so
there is nothing for the checker to recompute. The rule-level
statement exists:
\texttt{sound\_reluSplit} takes the two branch validities to the
conclusion\textquotesingle s validity, at the rank-$0$ scope the
recipe\textquotesingle s own per-scalar route always meant — and the
companion theorem \texttt{reluSplit\_rank0\_necessary} mechanizes the
mixed-sign countermodel showing the tensor-shaped reading is false,
upgrading the family\textquotesingle s shape-scope remark from
derivation-obstruction to semantic necessity for this rule. A
verified elaboration function producing the derivation does not yet
exist. The derivation requires the context-carrying form
of disjunction elimination — the sign-split disjunction arrives as a
sub-conjunct of the premise, which the plain sequent form of
$\vee$E cannot reach; see the $\vee$E section\textquotesingle s note.
}
\ruleblock{Lean statement}
\par\medskip
The rule-level semantic lemma, at the rank-$0$ scope: the template
$\Psi$ is an arbitrary formula with an arbitrary marker wire for $b$;
the three substitution instances are hypotheses via
\texttt{substituteInFormula} (\rulename{EqElim}\textquotesingle s
substitution predicate); the branch strengthenings are
\rulename{ReLU-Pos}/\rulename{ReLU-Neg}\textquotesingle s premise
atoms verbatim; the linkage and its extraction are the
family\textquotesingle s (\rulename{LiRPA} section). The second
theorem is the scope\textquotesingle s necessity witness: a
shape-$[2]$ fixture with mixed-sign input on which both branch
premises are valid and the conclusion fails.

\declhead{BernsteinLogic.ReluSplit}{sound\_reluSplit}{\badge{BadgeGreen}{statement bound}}
\begin{leansrc}
/-- ReLU-Split (rule-level semantic lemma; rank-0 scope, which
`reluSplit_rank0_necessary` shows is mandatory): from the validity of
the two branch sequents — the conclusion template at the preactivation
under `Φ ∧ (h ≥ 0)`, and at the constant 0 under `Φ ∧ (h ≤ 0)` — the
conclusion sequent `Φ ⊢ Ψ` is valid, where `Ψ` is the template at the
cited ReLU application's output wire. -/
theorem sound_reluSplit {H : Hypernet}
    {seq activePremise inactivePremise : Sequent}
    {preact post marker : Wire} {bodyForm : Formula}
    (ihActive : activePremise.semanticallyValid H)
    (ihInactive : inactivePremise.semanticallyValid H)
    (hlink : reluLinked H preact post)
    (hrank0 : ∃ τ, H.wireType post = some τ ∧ τ.shape = [])
    (hActPre : activePremise.pre
      = .and seq.pre (.atom .ge (.var preact) (.const 0)))
    (hInactPre : inactivePremise.pre
      = .and seq.pre (.atom .le (.var preact) (.const 0)))
    (hActPost : substituteInFormula bodyForm marker (.var preact)
      activePremise.post)
    (hInactPost : substituteInFormula bodyForm marker (.const 0)
      inactivePremise.post)
    (hconcl : substituteInFormula bodyForm marker (.var post) seq.post) :
    seq.semanticallyValid H
\end{leansrc}
\declhead{BernsteinLogic.ReluSplit}{reluSplit\_rank0\_necessary}{\badge{BadgeGreen}{statement bound}}
\begin{leansrc}
/-- The rank-0 hypothesis of `sound_reluSplit` is NECESSARY: the same
rule schema with the shape scope deleted is refuted — at a shape-[2]
ReLU node with mixed-sign input, both branch premises are valid while
the conclusion fails. (The header records the shape-[1] countermodel,
via broadcast-arity mixing, in prose; this is the mechanized mixed-sign
witness.) -/
theorem reluSplit_rank0_necessary :
    ¬ (∀ (H : Hypernet) (seq activePremise inactivePremise : Sequent)
        (preact post marker : Wire) (bodyForm : Formula),
        activePremise.semanticallyValid H →
        inactivePremise.semanticallyValid H →
        reluLinked H preact post →
        activePremise.pre
          = .and seq.pre (.atom .ge (.var preact) (.const 0)) →
        inactivePremise.pre
          = .and seq.pre (.atom .le (.var preact) (.const 0)) →
        substituteInFormula bodyForm marker (.var preact)
          activePremise.post →
        substituteInFormula bodyForm marker (.const 0)
          inactivePremise.post →
        substituteInFormula bodyForm marker (.var post) seq.post →
        seq.semanticallyValid H)
\end{leansrc}
\ruleblock{Correspondence (Steele display $\leftrightarrow$ Lean)}
\absentbox{NO CORRESPONDENCE DICTIONARY YET}{no reviewer-certified dictionary maps the displayed rule onto the mechanized statements above; until one is written and certified, that mapping rests on this section\textquotesingle s editorial care alone.}
\ruleblock{Lean verifier receipts}
\sourcenote{module \texttt{BernsteinLogic.ReluSplit}; bound build \texttt{01611bc8a2c3\ldots} (bound from receipt 6dd1e30ab7bf0a19, 2026-07-20T05:56:22Z). A receipt records that named checking layers accepted a build: (module, compiled-artifact SHA-256, environment SHA-256, checker layers), minted when the checking gate ran; the store receipts attesting the bound build are listed with the layers each run credited. Nothing here asks the reader to trust the store: every layer is an independently runnable program over the quoted sources --- rebuild the module and re-run the checkers to reproduce the verdict. Receipts attesting superseded builds credit no layer and are omitted from this rendering (the count of omissions, if any, is noted below).}
\begin{receipttable}
\receiptrow{6dd1e30ab7bf0a19}{2026-07-20T05:56:22Z}{build+decorrelated}{\badge{BadgeGreen}{current}}
\receiptrow{f5e25a81aaa84445}{2026-07-20T20:44:58Z}{build+decorrelated}{\badge{BadgeGreen}{current}}
\receiptrow{ca5158835179141d}{2026-07-20T21:01:26Z}{build+replay+hardened+decorrelated}{\badge{BadgeGreen}{current}}
\receiptrow{1a900df2daf383f4}{2026-07-21T05:28:56Z}{build+replay+hardened+decorrelated}{\badge{BadgeGreen}{current}}
\end{receipttable}
\layerline{\layeritem{build}{lake build (elaboration)}{\badge{BadgeGreen}{current}}\layeritem{replay}{leanchecker (official kernel replay)}{\badge{BadgeGreen}{current}}\layeritem{hardened}{safe\_verify (hardened screen)}{\badge{BadgeGreen}{current}}\layeritem{decorrelated}{lean4export + nanoda (decorrelated kernel)}{\badge{BadgeGreen}{current}}}
\sourcenote{\badge{BadgeGray}{comparator: none} no fixture in the three-implementation comparator harness (compiled / kernel-interpreted / reference-interpreted) exercises this rule.}
\ruleblock{HOL4 statement}
\absentbox{NO HOL STATEMENT YET}{\rulename{ReLU-Split} has no certificate representation on either side
of the seam: no payload constructor, no wire tag, no check function —
a producer elaborates the rule into its derivation, so a certificate
stream never contains a \rulename{ReLU-Split} node and there is
nothing for a HOL port to mirror. Any future HOL treatment would be a
formalization of the rule-level statement or the derivation itself
(the Lean side of the former exists — \texttt{sound\_reluSplit}), not
a port of checker code.
}
\ruleblock{HOL verifier receipts}
\absentbox{NONE}{no HOL verification receipt exists for this rule.}
\ruleblock{Proof}

\sourcenote{Genre: admissible rule (tier 2) --- but of the family's
third kind. \rulename{PConstExact} is admissible by reduction to
verified machinery; its siblings here are admissible with shipped
direct lemmas; \rulename{ReLU-Split} is admissible in the purest
sense: it has \emph{no proof-term constructor and no wire tag} --- a
certificate stream never contains a \rulename{ReLU-Split} node. The
derivation below is the rule's intended elaboration, and it becomes a
shippable certificate fragment exactly when the context-carrying case
split it opens with becomes a core constructor
(\autoref{rem:relusplit-ctxve}); the warrant that stands on its own
is the rule-level theorem: \texttt{sound\_reluSplit} states the rule
itself, at the rank-$0$ scope \autoref{rem:relusplit-scope} proves
is the rule's true extent.}

\begin{brecipe}[the \rulename{ReLU-Split} elaboration]
\label{rec:relusplit-instance}
The rule, as the source catalogue states it:
\[
\inferrule*[Right=\rulename{ReLU-Split}]{
  \mathcal{T}_{+} \rhd \Phi \wedge (h \geq 0) \vdash \Psi\sub{b}{h} \\
  \mathcal{T}_{-} \rhd \Phi \wedge (h \leq 0) \vdash \Psi\sub{b}{0}
}{
  \Phi \vdash \Psi
}
\]
where $b$ marks the occurrences of the output wire $y$ of a ReLU
application $\ReLU_{\{h\}\to\{y\}} \in H$, so that $\Psi =
\Psi\sub{b}{y}$. Two side conditions ride under that ``where'', made
explicit by the realized statement exactly as the siblings' displays
make theirs: the linkage itself (the warrant consumes the family's
port lookup), and --- semantic necessity, not caution --- the output
wire typed at rank $0$ (\autoref{rem:relusplit-scope}). The
intended elaboration of one such node is the derivation:
\begin{enumerate}
\item the node's sign-split fact --- $(h \leq 0 \wedge y = 0) \vee
  (0 \leq h \wedge y = h)$, the packaged semantic contract stated in
  the \rulename{LiRPA} section's family template --- enters as a
  sub-conjunct of the premise (the sp-axiom convention: for an
  application in $H$, its packaged contract is available premise-side
  as an axiom, true at every admissible valuation --- the same
  convention under which the cell rules' one-hot disjunction enters);
\item $\vee$E \emph{in its context-carrying form} splits on it (see
  \autoref{rem:relusplit-ctxve});
\item \emph{active branch} (context gains $0 \leq h$ and $y = h$):
  the branch context is aligned with $\mathcal{T}_{+}$'s premise $\Phi
  \wedge (h \geq 0)$ and consulted, in five named moves --- $\wedge$E
  extracts $\Phi$ and the case conjunct $0 \leq h$; the bound node's
  polarity-flip reading re-orients the latter to $h \geq 0$
  (\rulename{PBound}, reading (G), the one-node move displayed at the
  \rulename{ReLU-NN} section's active branch; the \rulename{ReLU-Above}
  section walks the toolkit's longer chains); $\wedge$I rebuilds $\Phi
  \wedge (h \geq 0)$; \rulename{Cut} against $\mathcal{T}_{+}$ yields
  $\Psi\sub{b}{h}$ in the branch context; and \rulename{EqSym} on the
  context's $y = h$ followed by \rulename{EqElim} along $h = y$
  (wire-to-wire, the node's one shared shape) transports it to
  $\Psi\sub{b}{y} = \Psi$;
\item \emph{inactive branch} (context gains $h \leq 0$ and $y = 0$):
  the same route with the mirrored atoms, four grouped moves
  ($h \leq 0$ needs no
  reorientation, so the \rulename{PBound} step drops out);
  $\mathcal{T}_{-}$ yields $\Psi\sub{b}{0}$, and \rulename{EqSym} +
  \rulename{EqElim} along $0 = y$ transport it to $\Psi$;
\item \rulename{Cut} and \rulename{Weaken} stitch the case analysis
  back under the original $\Phi$.
\end{enumerate}
Step 4's transport has a rank-$0$ constant side, so at tensor-shaped
instances it is stopped by \rulename{EqElim}'s broadcast-shape gate,
exactly as the family's shape-scope remark records
(\autoref{rem:relunn-shapescope}, \rulename{ReLU-NN} section) ---
at such instances a producer splits per scalar node
(\rulename{ReLU-Neg}'s direct node supplies the shape-general fact
$y = 0$, but transporting it into an arbitrary $\Psi$ rides the same
gate, so the per-scalar split is the general route --- and
necessarily so: \autoref{rem:relusplit-scope} proves the
tensor-shaped reading admits a countermodel).
\end{brecipe}

\begin{bremark}[why the context-carrying $\vee$E]
\label{rem:relusplit-ctxve}
The disjunction being eliminated does not stand alone as the premise:
it arrives \emph{inside} the premise conjunction, next to everything
else $\Phi$ carries. The plain sequent form of $\vee$E ($\Phi_1
\vdash \Psi$ and $\Phi_2 \vdash \Psi$ give $\Phi_1 \vee \Phi_2 \vdash
\Psi$) cannot reach it --- and this is not a stylistic gap: the
$\vee$E section records the countermodel showing the context-carrying
form is strictly stronger than the printed one (no elaboration of the
former into the latter exists). This derivation is the load-bearing consumer of that form. The form is
proven sound as a satellite lemma (\texttt{sound\_orECtx}), and that
is currently \emph{all} it is: it has no core constructor and no wire
tag, so the derivation above is checkable as a certificate fragment
only once the constructor lands --- until then the rule's standing
warrant is \autoref{rem:relusplit-mech}'s theorem, and the
elaboration is the design the constructor exists to serve.
\end{bremark}

\begin{bremark}[mechanization status --- the rule-level statement]
\label{rem:relusplit-mech}
The checker column remains empty for the structural reason it always
was: no certificate ever contains the rule, so there is no payload to
validate and no soundness case in the checker's recursion. The
rule-level statement exists: \texttt{sound\_reluSplit} is
the semantic lemma --- from the validity of the two branch sequents
and the linkage, the conclusion sequent is valid --- built from
exactly the vocabulary the rest of the catalogue already uses (the linkage extraction
\texttt{reluLinked\_facts}, the substitution predicate
\texttt{substituteInFormula} that \rulename{EqElim}'s payload already
uses, and the semantic transport \texttt{Formula.holds\_subst\_congr}
that \rulename{EqElim}'s soundness already rides). The template
$\Psi$ with its marked occurrences $b$ is an arbitrary formula over an
arbitrary marker wire, with the three substitution instances
($\Psi\sub{b}{h}$, $\Psi\sub{b}{0}$, $\Psi = \Psi\sub{b}{y}$) as
hypotheses --- no freshness condition, since the proof transports
along value equality, never syntax. ``Derivable, by inspection of
this recipe'' is thereby backed by a theorem at every
instance the recipe itself reaches (rank-$0$ nodes ---
\autoref{rem:relusplit-scope}); a
verified elaboration function producing
Recipe~\ref{rec:relusplit-instance}'s derivation does not yet exist.
\end{bremark}

\begin{bremark}[the rank-$0$ scope is semantic necessity]
\label{rem:relusplit-scope}
\texttt{sound\_reluSplit} hypothesizes the output wire typed at
rank~$0$, and this is the rule's true extent rather than proof
convenience. The family's shape-scope remark
(\autoref{rem:relunn-shapescope}) records that the
\emph{derivation}'s constant-sided transport is stopped by
\rulename{EqElim}'s broadcast-shape gate at tensor shapes; for the
siblings the direct lemma is nonetheless shape-general, so the gate
reads there as a technicality of the derivation route. Not here:
\texttt{reluSplit\_rank0\_necessary} (same module) mechanizes a
countermodel showing the shape-general reading of
\rulename{ReLU-Split} is \emph{false}. Template $\Psi(b) = (b = 0)
\vee (b = h)$ over a shape-$[2]$ node, $\Phi = \top$: both premises
are valid ($\mathcal{T}_{+}$ through $h = h$, $\mathcal{T}_{-}$
through $0 = 0$), yet at the mixed-sign valuation $h = [1, -1]$,
$y = [1, 0]$ neither disjunct of the conclusion holds. Mixed signs
are exactly what the rule's case analysis cannot see from
sign-definite branch hypotheses --- the semantic content of ``each
split doubles the branch count'': a tensor node is $2^{n}$ scalar
cases, not $2$. Even shape $[1]$ fails, by a different leak
(broadcast-arity mixing): take $\Psi(b) = (b = h) \vee (b + w \geq 0)$
with $w$ typed at $[2]$ and premise $\Phi = (w \geq 0)$. Both branch
sequents are valid --- $\mathcal{T}_{+}$ through $h = h$,
$\mathcal{T}_{-}$ through $0 + w \geq 0$, the rank-$0$ constant
broadcasting to $w$'s shape --- yet at $h = [-1]$, $y = [0]$,
$w = [1, 1]$ the conclusion fails: $y = h$ fails entrywise, and the
$[1]$-shaped $y + w$ atom is unsatisfiable outright by the
strict-shape arithmetic, which broadcasts rank-$0$ sides only (this
leak is stated as prose beside the mechanized $[2]$ witness, not
itself mechanized). The recipe's step-4 note
(``the per-scalar split is the general route'') is thus not a
workaround for a gate but the only sound reading; producers split per
rank-$0$ node, and \texttt{sound\_reluSplit} is the rule those splits
instantiate.
\end{bremark}

\begin{bremark}[when a producer splits]
Splitting is the \emph{complete}-verification move of the
neural-network literature --- Reluplex's case
analysis~\cite{katz2017reluplex}, systematized as branch and
bound~\cite{bunel2020branch}: exact where \rulename{LiRPA}'s envelope
is lossy, at the price of one branch doubling per split neuron, so
proof size is exponential in the split count. The practical frontier
--- which neurons to split, which to relax --- is the branch-and-bound
literature's whole subject; the calculus prices the choice honestly
(a split is two subtrees, an envelope is one node) and takes no
position on it.
\end{bremark}

\begin{anchors}
the sign-split fact (step 1) & \texttt{PrimApp.SemOk} --- the packaged contract; conservativity via \texttt{relu\_projection} (\rulename{LiRPA} section)\\
context-carrying case split (step 2) & \texttt{sound\_orECtx} (satellite lemma, Admissible.lean)\\
equality transport (steps 3--4) & \texttt{sound\_eqElim}; \texttt{substituteInFormula}\\
rule-level statement (rank-$0$ scope) & \texttt{sound\_reluSplit} (ReluSplit.lean) --- \autoref{rem:relusplit-mech}\\
scope necessity (mixed-sign countermodel) & \texttt{reluSplit\_rank0\_necessary}; fixture \texttt{ReluSplitScope.HSplit} --- \autoref{rem:relusplit-scope}\\
verified elaboration function & \emph{none exists} --- the remaining gap, with the core constructor for the context-carrying case split\\
\end{anchors}


\blcatsection{Assume--Guarantee Composition}
\rulesection{32}{Promise}{admissible rule}
\label{blcatrule-Promise}
\ruleblock{The rule (Steele metanotation)}
\sourcenote{as realized in the mechanized development; source row: \texttt{\texttt{rules.toml}} --- \emph{Promise}. The metanotation statement is typeset by hand for notational quality; its mechanized ground truth is the Lean statement below, extracted from the sources at every build of this book and hash-bound. The rule \emph{roster} is machine-checked (every mechanized rule tag round-trips to exactly one catalogue row); the statement \emph{text} is editorial, checked by eye against the extracted Lean beside it.}
\begin{mathpar}
\inferrule*[Right=\rulename{Promise}]{
  \Phi_1 \vdash Q \\
  \Phi_2 \wedge Q \vdash \Psi \\
  \wires(Q) \subseteq Y \\
  \wires(\Phi_1) \cap \bigl(\wires(\Phi_2) \cup \wires(\Psi)\bigr) \subseteq Y
}{
  \Phi_1 \wedge \Phi_2 \vdash \Psi
}
\end{mathpar}
$Y$ = the declared interface (payload wire list; bundle entries addressed as $\mathrm{TIdx}(w,i)$). The two confinement checks are the locality warrant, not soundness needs: the contract $Q$ is the \emph{datasheet} --- the public artefact whose vocabulary is the separator algebra --- and the cost identities (LP size, ZK disclosure set) quantify over certificates carrying these checks.

\ruleflavor{a Bernstein Logic rule outright: the baked base theory lists it in its core roster (grouped there under an assume--guarantee comment), the certified checker decodes its wire node, and its certificate relations ride the pin table --- there is no intermediate composition logic between this rule and the base calculus. The tier word above classifies proof theory (the unpack is \rulename{Weaken} + $\wedge$-restructuring + \rulename{Cut}); core-roster membership is wire placement.}
\par\medskip
\rulename{Promise} is the calculus\textquotesingle s acyclic composition
rule: it discharges an assumption one component makes about another. The
supplier proves the contract $Q$ from its own premises
($\Phi_1 \vdash Q$); the consumer proves the goal \emph{using} $Q$
($\Phi_2 \wedge Q \vdash \Psi$); the composite concludes
$\Phi_1 \wedge \Phi_2 \vdash \Psi$ with the contract discharged
internally --- cut, wearing interface bookkeeping. A producer reaches for
it wherever a system decomposes along a directed interface: certify the
downstream component against a \emph{datasheet} (the contract, written
over the declared interface wires only), certify the upstream component
as meeting the datasheet, and compose without ever re-deriving either
side inside the other.

What the checker recomputes: the splice plumbing and the locality
warrant. The plumbing --- the contract sits at a payload-declared
conjunct path in the consumer\textquotesingle s premise; the composite
premise is the consumer\textquotesingle s with the contract replaced by
the supplier\textquotesingle s premise at that path --- is verified by
exact read/replace on the conjunction tree. The locality warrant is two
wire-confinement checks: the contract\textquotesingle s free wires lie
inside the declared interface, and the supplier\textquotesingle s wires
meet the consumer\textquotesingle s and the goal\textquotesingle s only
inside it.

The hypotheses split by role, and the split is the
rule\textquotesingle s point. Soundness consumes only the plumbing:
drop the path checks and the \textquotedblleft composite\textquotedblright{}
premise could be unrelated to the two sub-proofs. The confinement checks
are \emph{warrant, not soundness} --- cut is sound regardless --- but
dropping them destroys the datasheet reading: a contract mentioning
supplier-private wires is no longer a public interface artefact, and
every quantitative consequence hung on the interface (certificate size
scaling with the separator rather than the components, bounded
disclosure of component internals) quantifies over certificates that
carry these checks. The proof section separates the two roles
explicitly.

\cataloguenote{Tier reading: proof-theoretically the rule is admissible --- one node
unpacks into \rulename{Cut} against premise-transported copies of the
cited sub-derivations, a global rewrite of those sub-proofs rather
than a local chain (the proof section\textquotesingle s recipe walks
why: \rulename{Weaken} and \rulename{Cut} are head-positional) ---
while in the realized
calculus it is a native constructor, kept for cost and for the
locality-warrant payload. The tier column records the proof-theoretic
status; constructor membership is wire placement.
}
\ruleblock{Lean statement}
\declhead{BernsteinLogic.Soundness}{PromiseStep}{\badge{BadgeGreen}{statement bound}}
\begin{leansrc}
/-- Applicability data of the Promise rule, ctx-style — field-for-field
the reference checker's obligations, and field-for-field `PPromiseData`
minus the two sub-proofs (which enter `sound_promise` as validity
hypotheses). The two confinement fields are checked side-conditions that
soundness does NOT consume (AG draft §1.2: cut is sound regardless) —
they are the locality warrant. -/
structure PromiseStep where
  seq : Sequent
  supplier : Sequent
  consumer : Sequent
  contract : Formula
  path : ConjPath
  iface : List Wire
  ctx : Formula
  supplierPost : supplier.post = contract
  pathReads : consumer.pre.conjGet? path = some contract
  prePlumb : consumer.pre.conjReplace? path supplier.pre = some seq.pre
  postEq : seq.post = consumer.post
  ctxReads : consumer.pre.conjReplace? path .tt = some ctx
  contractConfined : formulaWiresWithin contract iface
  separation : wiresSeparated supplier.pre ctx seq.post iface
\end{leansrc}
\declhead{BernsteinLogic.Soundness}{sound\_promise}{\badge{BadgeGreen}{statement bound}}
\begin{leansrc}
/-- Soundness of Promise: valid supplier + valid consumer give a valid
composite. Read the supplier's premise at the path, obtain the contract
from supplier validity, replace it back to recover the consumer's
premise, conclude by consumer validity. Pure predicate logic over the
semantics; the confinement fields are untouched. -/
theorem sound_promise {H : Hypernet} (st : PromiseStep)
    (ihs : st.supplier.semanticallyValid H)
    (ihc : st.consumer.semanticallyValid H) :
    st.seq.semanticallyValid H
\end{leansrc}
\declhead{BernsteinLogic.Soundness}{Formula.holds\_of\_conjGet}{\badge{BadgeGreen}{statement bound}}
\begin{leansrc}
/-- The conjunct addressed by a path is entailed: if Γ holds, Γ|π holds.
(Promoted from Admissible.lean — the ∨E-ctx kit; the Promise contract
path reuses it.) -/
theorem Formula.holds_of_conjGet {H : Hypernet} {σ : H.Valuation} :
    ∀ {π : ConjPath} {Γ d : Formula},
    Γ.conjGet? π = some d → Γ.holds σ → d.holds σ
\end{leansrc}
\declhead{BernsteinLogic.Soundness}{Formula.holds\_conjReplace}{\badge{BadgeGreen}{statement bound}}
\begin{leansrc}
/-- Replacing the addressed conjunct by a formula that holds preserves
holding: if Γ holds and g holds, Γ[π := g] holds. (Promoted from
Admissible.lean.) -/
theorem Formula.holds_conjReplace {H : Hypernet} {σ : H.Valuation} :
    ∀ {π : ConjPath} {Γ Γ' g : Formula},
    Γ.conjReplace? π g = some Γ' → Γ.holds σ → g.holds σ → Γ'.holds σ
\end{leansrc}
\declhead{BernsteinLogic.Soundness}{Formula.conjGet\_of\_conjReplace}{\badge{BadgeGreen}{statement bound}}
\begin{leansrc}
/-- Reading at the path of a successful replacement returns the
replacement. (Promoted from Composition.lean §G2.) -/
theorem Formula.conjGet_of_conjReplace :
    ∀ {π : ConjPath} {Γ Γ' g : Formula},
    Γ.conjReplace? π g = some Γ' → Γ'.conjGet? π = some g
\end{leansrc}
\declhead{BernsteinLogic.Soundness}{Formula.conjReplace\_back}{\badge{BadgeGreen}{statement bound}}
\begin{leansrc}
/-- Replacing back the conjunct a read found recovers the original: the
round trip Γ ⟶[π := g] Γ' ⟶[π := Γ|π] Γ. (Promoted from
Composition.lean §G2.) -/
theorem Formula.conjReplace_back :
    ∀ {π : ConjPath} {Γ Γ' q g : Formula},
    Γ.conjGet? π = some q → Γ.conjReplace? π g = some Γ' →
    Γ'.conjReplace? π q = some Γ
\end{leansrc}
\par\medskip
\textbf{The confinement side conditions, as checked.}\enspace The
proof\textquotesingle s premise-transport induction leans on the
wire-discipline predicates the checker recomputes; their definitions
are displayed so the side conditions\textquotesingle{} meaning is on
the page.

\declhead{BernsteinLogic.Core.Catalogue}{premiseContains}{\badge{BadgeGreen}{definition bound}}
\begin{leansrc}
/-- Conjunct-wise premise containment (geolog `premiseContains`): every
conjunct of `sub` is a conjunct of `pre`. Decouples component premises
from the composite premise — the core has no ∧-restructuring rules, so
without this a component proved over its own box could not enter a
composite whose premise merely CONTAINS that box. -/
def premiseContains (pre sub : Formula) : Prop :=
  ∀ a ∈ sub.conjunctList, a ∈ pre.conjunctList
\end{leansrc}
\declhead{BernsteinLogic.Core.Catalogue}{formulaWiresWithin}{\badge{BadgeGreen}{definition bound}}
\begin{leansrc}
/-- Wire confinement (geolog `formulaWiresWithin`): every free wire of
the formula lies in the declared interface. Prover-side warrant that
demanding this costs no completeness: Lemma A, `farkas_interpolant_vars`
(contracts are mined confined). -/
def formulaWiresWithin (f : Formula) (iface : List Wire) : Prop :=
  ∀ w ∈ f.freeWires, w ∈ iface
\end{leansrc}
\declhead{BernsteinLogic.Core.Catalogue}{wiresSeparated}{\badge{BadgeGreen}{definition bound}}
\begin{leansrc}
/-- Separation (geolog `wiresSeparated`, ctx-form): every wire the
supplier's premise shares with the consumer's context or the conclusion
routes through the interface. The context formula is the consumer's
premise with the contract slot blanked (`conjReplace?` at `.tt`). -/
def wiresSeparated (left ctx post : Formula) (iface : List Wire) : Prop :=
  ∀ w ∈ left.freeWires,
    (w ∈ ctx.freeWires ∨ w ∈ post.freeWires) → w ∈ iface
\end{leansrc}
\ruleblock{Correspondence (Steele display $\leftrightarrow$ Lean)}
\badge{BadgeGreen}{reviewer-certified against the displayed texts}
\sourcenote{The dictionary below is a reviewer-certified statement of how the mechanized statements above interpret the displayed rule: display metavariable $\mapsto$ Lean binder, each display hypothesis $\mapsto$ the premise or side condition carrying it, the display conclusion $\mapsto$ the goal --- with any deliberate difference in width between display and mechanization named outright. Certification is recorded as a hash triple over (the displayed cell, the extracted statement texts, this dictionary\textquotesingle s own text); if any of the three changes, this block turns amber until a reviewer re-certifies. The correspondence claim itself remains editorial --- reviewer-grade, not machine-checked; its machine-checked successor types each soundness statement against the rule as a structured object, and supersedes this block per rule as it lands.}
Judgment frame as in \rulename{Ax}; the statement bundles its data as
\texttt{PromiseStep} (field-for-field the payload minus the two
sub-proofs, which enter as \texttt{ihs}/\texttt{ihc}).
\begin{itemize}
\item $\Phi_1 \vdash Q$ $\mapsto$ \texttt{ihs} with
  \texttt{supplierPost} $:$ \texttt{supplier.post = contract}; $Q
  \mapsto$ \texttt{contract : Formula}; $\Phi_1 \mapsto$
  \texttt{supplier.pre}.
\item $\Phi_2 \wedge Q \vdash \Psi$ $\mapsto$ \texttt{ihc}, with the
  certified statement taking the path-general form: $Q$ sits at a
  \emph{payload-declared conjunct path} (\texttt{pathReads} $:$
  \texttt{consumer.pre.conjGet? path = some contract}), not at a fixed
  position --- the ctx-$\vee$E style. $\Phi_2 \mapsto$ \texttt{ctx},
  the consumer premise with the contract deleted at the path
  (\texttt{ctxReads}).
\item conclusion $\Phi_1 \wedge \Phi_2 \vdash \Psi$ $\mapsto$ the
  goal for \texttt{st.seq}, whose premise is pinned by
  \texttt{prePlumb}: the consumer\textquotesingle s premise with the
  supplier\textquotesingle s premise spliced in \emph{at the same
  path} where $Q$ sat. Order honesty: at the displayed root-right
  instance ($Q$ the right conjunct of $\Phi_2 \wedge Q$), the splice
  therefore pins \texttt{seq.pre} to $\Phi_2 \wedge \Phi_1$ --- the
  displayed $\Phi_1 \wedge \Phi_2$ is obtained by semantic
  conjunction commutativity, not by direct instantiation. Path
  generality remains strictly wider than the display, in the safe
  direction; the conjunct \emph{order} is the one representational
  difference, named here. \texttt{postEq} carries $\Psi$ across
  unchanged.
\item the two confinement premises ($\wires(Q) \subseteq Y$; the
  interface-separation check) $\mapsto$ the Step fields
  \texttt{contractConfined} and \texttt{separation} over
  \texttt{iface} $\mapsto Y$ --- carried, checked,
  and consumed by \emph{no} step of \texttt{sound\_promise}: the
  displayed caption\textquotesingle s claim that they are the
  locality warrant rather than soundness needs is visible in the
  statement itself: the proof never touches them, as on the QTSL
  pages.
\end{itemize}
Width note: beyond the deliberate path-generality and the
conclusion-order transposition named above (the splice yields
$\Phi_2 \wedge \Phi_1$; the display\textquotesingle s order needs
conjunction commutativity, semantically free), neither side is
wider; the soundness content is pure predicate-logic cut over the
semantics, exactly as the witness column says.

\ruleblock{Lean verifier receipts}
\sourcenote{module \texttt{BernsteinLogic.Soundness}; bound build \texttt{a20d90e66c13\ldots} (bound from receipt a9c0e34d3e3c9c8c, 2026-08-03T21:30:36Z). A receipt records that named checking layers accepted a build: (module, compiled-artifact SHA-256, environment SHA-256, checker layers), minted when the checking gate ran; the store receipts attesting the bound build are listed with the layers each run credited. Nothing here asks the reader to trust the store: every layer is an independently runnable program over the quoted sources --- rebuild the module and re-run the checkers to reproduce the verdict. Receipts attesting superseded builds credit no layer and are omitted from this rendering (the count of omissions, if any, is noted below).}
\begin{receipttable}
\receiptrow{a9c0e34d3e3c9c8c}{2026-08-03T21:30:36Z}{build+decorrelated}{\badge{BadgeGreen}{current}}
\end{receipttable}
\sourcenote{21 further receipt(s) in the store attest superseded builds of this module and are omitted from this rendering.}
\sourcenote{module \texttt{BernsteinLogic.Core.Catalogue}; bound build \texttt{cbc92e976a1d\ldots} (bound from receipt 5abdb3589707b5df, 2026-08-03T21:30:33Z). A receipt records that named checking layers accepted a build: (module, compiled-artifact SHA-256, environment SHA-256, checker layers), minted when the checking gate ran; the store receipts attesting the bound build are listed with the layers each run credited. Nothing here asks the reader to trust the store: every layer is an independently runnable program over the quoted sources --- rebuild the module and re-run the checkers to reproduce the verdict. Receipts attesting superseded builds credit no layer and are omitted from this rendering (the count of omissions, if any, is noted below).}
\begin{receipttable}
\receiptrow{24804e66edba0bc9}{2026-08-03T20:25:09Z}{build+decorrelated}{\badge{BadgeGreen}{current}}
\receiptrow{5abdb3589707b5df}{2026-08-03T21:30:33Z}{build+decorrelated}{\badge{BadgeGreen}{current}}
\end{receipttable}
\sourcenote{18 further receipt(s) in the store attest superseded builds of this module and are omitted from this rendering.}
\layerline{\layeritem{build}{lake build (elaboration)}{\badge{BadgeGreen}{current}}\layeritem{replay}{leanchecker (official kernel replay)}{\badge{BadgeAmber}{stale only}}\layeritem{hardened}{safe\_verify (hardened screen)}{\badge{BadgeAmber}{stale only}}\layeritem{decorrelated}{lean4export + nanoda (decorrelated kernel)}{\badge{BadgeGreen}{current}}}
\sourcenote{\badge{BadgeGray}{comparator: none} no fixture in the three-implementation comparator harness (compiled / kernel-interpreted / reference-interpreted) exercises this rule.}
\ruleblock{HOL4 statement}
\sourcenote{theory \texttt{bern\_fam\_resid}, HOL4 development script \texttt{\texttt{bern\_fam\_residScript.sml}} \badge{BadgeGreen}{script bound}}
\holhead{Definition checkFam\_promise\_def}
\begin{holsrc}
Definition checkFam_promise_def:
  checkFam_promise (cols:columns) (ct:ctx) (i:num) ⇔
    case frSeq cols ct i of
    | NONE => F
    | SOME sq =>
      case frChildSeq cols ct i 0 of
      | NONE => F
      | SOME ss =>
        case frChildSeq cols ct i 1 of
        | NONE => F
        | SOME cs =>
          case frFormula cols ct i 0 of
          | NONE => F
          | SOME ctr =>
            case pathOf cols ct i of
            | NONE => F
            | SOME (stps, ifc) =>
              let pth = stepsToPath stps in
                ss.post = ctr ∧
                conjGet cs.pre pth = SOME ctr ∧
                conjReplace cs.pre pth ss.pre = SOME sq.pre ∧
                sq.post = cs.post ∧
                (case conjReplace cs.pre pth FTt of
                 | SOME cx0 =>
                     (EVERY (λw. MEM w ifc) (formulaFreeWires ctr) ∧
                      EVERY (λw. (MEM w (formulaFreeWires cx0) ∨
                                  MEM w (formulaFreeWires sq.post)) ⇒
                                 MEM w ifc)
                            (formulaFreeWires ss.pre))
                 | NONE => F)
End
\end{holsrc}
\holhead{Theorem famEq\_promise}
\begin{holsrc}
Theorem famEq_promise:
  wfColumns cols ct ∧ mkCtx cols = SOME ct ∧
  parseCert (cols.hdrImage ++ rest) = SOME (ws, pf) ∧
  nodeCorr cols ct i (CertPf sq nd) ∧
  FST (THE (tagOf cols i)) = 25 ⇒
  (checkFam_promise cols ct i ⇔ armPromise ws sq nd)
\end{holsrc}
\ruleblock{HOL verifier receipts}
\badge{BadgeBlue}{documentary} Documentary citations. Family validation: \texttt{bern\_fam\_resid} built GREEN with all ten \texttt{famEq} arms validated, \texttt{hol-strict} clean; spine census: \texttt{bern\_joins} GREEN under \texttt{hol-strict} (55 theorems, 0 offences, 208-theory ancestry including this theory), \texttt{checkJoins\_eq\_checkPf} verified verbatim in the built signature. The cited \texttt{famEq} lemma is the kernel-checked per-node equivalence: under well-formed columns (\texttt{wfColumns}/\texttt{mkCtx}/\texttt{nodeCorr}) at this rule\textquotesingle s wire tag, the columnar family check coincides with the flat reference walk\textquotesingle s arm for this rule. There is no signed HOL receipt store yet --- these citations are documentary.
Semantic stratum, also ported: the soundness twin
\texttt{sound\_promise} (theory \texttt{bern\_tr\_b14}: \texttt{hol-
strict} clean, 2 theorems, 0 offences, verification run of 2026-07-22 on
an isolated build host, content-addressed at the same sources), consumed
by the certified-fragment keystone (\texttt{checkPf\_sound\_frag},
theory \texttt{bern\_tr\_splice}, 20 theorems, 0 offences, same run): a
reference-checker accept inside the proven wire fragment yields sequent
validity.

\sourcenote{HOL receipts are documentary: quoted results of \texttt{hol-strict} runs
and census rows from the HOL4 development\textquotesingle s build
evidence, not minted receipt artefacts. A HOL receipt store (the Lean
store\textquotesingle s twin) is named future work.
}
\ruleblock{Proof}

\sourcenote{Genre: admissible rule (tier 2) --- a proof-tree-rewriting
recipe. Proof-theoretically \rulename{Promise} unpacks into core rules
(conjunct plumbing + \rulename{Cut}); in the realized calculus it is a
native constructor, kept for cost and for the locality-warrant payload.
The recipe shows the unpack; the semantic warrant is the kernel lemma
\texttt{sound\_promise}, which is the unpack's semantic shadow run in
one pass.}

\textbf{Setting.}\enspace The payload (the \texttt{PromiseStep}
structure quoted above) carries a supplier sequent $\Phi_1 \vdash Q$, a
consumer sequent whose premise contains the contract $Q$ at a declared
conjunct path $\pi$ (write it $\Phi_2 \wedge_{\pi} Q \vdash \Psi$ ---
the $\pi$-notation meaning $Q$ sits at path $\pi$ in the premise's
conjunction tree, $\Phi_2$ being the rest), a declared interface wire
list $Y$, and the composite sequent, whose premise is the consumer's
with $Q$ \emph{replaced} by $\Phi_1$ at $\pi$ and whose conclusion is
$\Psi$. Four checked hypotheses: \emph{(P1)} the path reads $Q$ in the
consumer's premise; \emph{(P2)} replacing $Q$ by $\Phi_1$ at $\pi$
yields exactly the composite premise; \emph{(P3)}
$\wires(Q) \subseteq Y$; \emph{(P4)} the supplier's wires meet the
consumer's and the goal's only inside $Y$. (P1)--(P2) are the splice
plumbing; (P3)--(P4) are the locality warrant.

\begin{brecipe}[the \rulename{Promise} instance --- cut through a
conjunct path]\label{rec:promise-instance}
The general admissible template (one rule node $\rightsquigarrow$ a
derivation fragment over core rules with the same conclusion, applied
leaf-first until no such node remains) instantiates as:
\[
\inferrule*[Right=\rulename{Promise}]{\;\mathcal{T}_s \rhd \Phi_1 \vdash Q \\ \mathcal{T}_c \rhd \Phi_2 \wedge_{\pi} Q \vdash \Psi\;}{\Phi_2 \wedge_{\pi} \Phi_1 \vdash \Psi}
\qquad\rightsquigarrow\qquad
\inferrule*[Right=\rulename{Cut}]{
  \inferrule*{\mathcal{T}_s}{\Phi_2 \wedge_{\pi} \Phi_1 \vdash Q}
  \\
  \inferrule*{\mathcal{T}_c}{\Phi_2 \wedge_{\pi} Q \vdash \Psi}
}{\Phi_2 \wedge_{\pi} \Phi_1 \vdash \Psi}
\]
where, writing $C$ for the composite premise
$\Phi_2 \wedge_{\pi} \Phi_1$: the left branch derives $C \vdash Q$ by
\rulename{Cut} from $C \vdash \Phi_1$ (an \rulename{Ax} leaf
$C \vdash C$, then a conclusion-side $\wedge$E chain along $\pi$)
against a \emph{premise-transported} copy of $\mathcal{T}_s$, and
\rulename{Cut}'s right premise proper is $C \wedge Q \vdash \Psi$, a
premise-transported copy of $\mathcal{T}_c$ --- the display elides both
transports. Premise transport (re-rooting a derivation $X \vdash \Psi$
at a premise $Y$ whose addressed conjuncts entail $X$) is \emph{not} a
local chain in this calculus: \rulename{Weaken} and \rulename{Cut} are
deliberately head-positional --- each grafts one conjunct at the
premise root's right --- so the premises reachable from an opaque
sub-proof $X \vdash \Psi$ by insertion alone are exactly $X$ with
right-appended conjuncts (\rulename{Cut} can also \emph{consume} an
already-present root-right conjunct when a proof of it is supplied ---
removal, never interior insertion or permutation), and even
$Y \wedge X \vdash \Psi$ is locally out of reach.
The transport is instead an induction over the sub-derivation itself;
its principal cases: rules that copy their premise tree verbatim from
their sub-proofs re-instantiate on the transported children;
premise-free leaves conclude at an arbitrary premise by the
\rulename{Norm} convention (an \rulename{Ax} leaf $\Phi \vdash \Phi$
re-derives at $Y$ as $Y \vdash Y$ plus a conclusion-side
$\wedge$E/$\wedge$I rebuild); and the two head-positional rules
are their own inductive cases, because re-rooting commutes with the
head graft: a transported \rulename{Weaken} grafts its conjunct onto
the new root, and a transported \rulename{Cut}'s right premise
re-forms as the new root with the cut formula appended. This list is
not exhaustive over the catalogue: the premise-\emph{forming} rules
($\vee$E's conclusion premise is the disjunction of its branch
premises; \rulename{Subdivide}'s branch premises conjoin the region
data) and the context-carrying rules (\rulename{Promise},
\rulename{Syzygy}, whose path plumbing must compose with the splice
--- the $\vee$E section's machinery) are neither verbatim copies nor
leaves; each needs its own commutation case, and those cases belong to
the same transport metatheorem flagged as unmechanized below.
Rewriting the premise field node by node under that case analysis
yields the transported copies of $\mathcal{T}_s$ and $\mathcal{T}_c$. Elimination costs
the sizes of $\mathcal{T}_s$ and $\mathcal{T}_c$ plus the
$\wedge$E chain along $\pi$ at the auxiliary derivation plus the
per-leaf \rulename{Ax} rebuilds --- and each transported leaf pays its
own conclusion-side rebuild chain (the proof-object grammar has no
sharing), so the total is bounded by the certificate sizes plus the
number of transported leaves times the addressing depth, linear in the
certificate data only when that depth is bounded ---
and consumes \rulename{Ax} leaves throughout; conservativity holds
with no new kernel semantic content, but as a global rewrite warranted
by the transport induction (stated here as prose --- the transport
metatheorem is not yet mechanized), not as insertion of a
constant-size fragment.
\end{brecipe}

\begin{bprop}[semantic warrant]\label{prop:promise-sound}
If the supplier and consumer sequents are semantically valid and
(P1)--(P2) hold, the composite sequent is semantically valid.
\end{bprop}
\begin{proof}
Fix an admissible valuation $\sigma$ satisfying the composite premise.
Reading the composite premise at $\pi$ returns $\Phi_1$: (P1) says
$\pi$ addresses $Q$ in the consumer's premise, so the replacement (P2)
performs succeeded at $\pi$, and the read of a successful replacement
returns the replacement (the round-trip lemma
\texttt{Formula.conjGet\_of\_conjReplace}); an addressed conjunct
of a holding conjunction holds (\texttt{Formula.holds\_of\_conjGet});
so $\sigma \models \Phi_1$, and supplier validity gives
$\sigma \models Q$. Replacing $\Phi_1$ back by $Q$ at $\pi$ recovers
the consumer's premise exactly
(\texttt{Formula.conjReplace\_back}, using (P1)--(P2)), and replacing
the addressed conjunct of a holding conjunction by a formula that
itself holds at $\sigma$ preserves holding
(\texttt{Formula.holds\_conjReplace}, applied at $\pi$ with $Q$ for
$\Phi_1$); so
$\sigma \models \Phi_2 \wedge_{\pi} Q$, and consumer validity gives
$\sigma \models \Psi$. This is \texttt{sound\_promise} verbatim: pure
predicate logic over the semantics --- admissibility of $\sigma$ is
passed through untouched, and neither confinement hypothesis is
consulted.
\end{proof}

\begin{bremark}[the confinement premises are warrant, not soundness]
\label{rem:promise-warrant}
(P3)--(P4) appear in no step of Proposition~\ref{prop:promise-sound},
and dropping them alone --- (P1)--(P2) retained --- admits no false
sequent: cut is sound regardless of where the contract's wires live. What their removal \emph{does} admit
is a degenerate reading of the rule --- a ``contract'' mentioning
supplier-private wires, i.e.\ a composition whose interface is as wide
as the components themselves. Everything the rule is \emph{for} then
fails: the contract stops being a datasheet (a public artefact whose
vocabulary is the declared separator), certificate size stops scaling
with the interface, and any bounded-disclosure reading of ``what the
consumer learns about the supplier'' becomes vacuous. Those
consequences are theorems proved elsewhere, \emph{about certificates
that carry the checks} --- they quantify over exactly the certificates
these hypotheses admit --- which is why the checks are checked --- a hypothesis whose
necessity is economic and epistemic rather than semantic is exactly
the kind a recipe must label, or the next reader deletes it as
redundant.
\end{bremark}

\begin{bremark}[relation to the literature]
Assume--guarantee composition descends from composing-specifications
calculi \cite{abadilamport1993}; the subtlety those calculi negotiate
is \emph{circularity} --- each component assuming what the other
guarantees --- which needs an induction over time or a fixpoint
argument \cite{mcmillan1999}. \rulename{Promise} deliberately takes
the acyclic case, where plain cut suffices and the whole rule is
bookkeeping around it; the circular case enters this calculus through
the feedback rules --- \rulename{SmallGain} and \rulename{SmallGainN}
at the sequent layer and the feedback cells of the satisfaction-layer
composition rule, where the guarded-fixpoint structure is explicit,
and \rulename{QWireComposeMarks} for liveness marks, where cyclic
wiring needs no well-foundedness device at all and the wiring
condition is a feasibility walk (schedules and pinned values must
close around cycles) rather than a causality walk. Separating the
cases --- cut where cut suffices, guarded fixpoints where feedback
threatens a safety face, a schedule walk where a cycle only has to
deliver a mark --- is the design's answer to the classical observation
that \emph{unguarded} circular assume--guarantee reasoning about
safety is unsound --- there, some well-foundedness device is required.
\end{bremark}

\begin{anchors}
Recipe \ref{rec:promise-instance} payload \& plumbing & \texttt{PromiseStep} (fields \texttt{path}, \texttt{pathReads}, \texttt{prePlumb}, \texttt{postEq})\\
path read/replace round trips & \texttt{Formula.conjGet\_of\_conjReplace}, \texttt{Formula.conjReplace\_back}\\
addressed conjunct entailed; replacement preserves holding & \texttt{Formula.holds\_of\_conjGet}, \texttt{Formula.holds\_conjReplace}\\
Proposition \ref{prop:promise-sound} & \texttt{sound\_promise}; wired as core \texttt{RuleData.promise} into \texttt{ProofObject.sound}\\
confinement checks (warrant only) & \texttt{formulaWiresWithin}, \texttt{wiresSeparated} (payload fields \texttt{contractConfined}, \texttt{separation})\\
\end{anchors}


\rulesection{33}{Syzygy}{core rule}
\label{blcatrule-Syzygy}
\ruleblock{The rule (Steele metanotation)}
\sourcenote{as realized in the mechanized development; source row: \texttt{\texttt{rules.toml}} --- \emph{Syzygy}. The metanotation statement is typeset by hand for notational quality; its mechanized ground truth is the Lean statement below, extracted from the sources at every build of this book and hash-bound. The rule \emph{roster} is machine-checked (every mechanized rule tag round-trips to exactly one catalogue row); the statement \emph{text} is editorial, checked by eye against the extracted Lean beside it.}
\begin{mathpar}
\inferrule*[Right=\rulename{Syzygy}]{
  \bigl\{\Phi_i \vdash e_i \leq 0\bigr\}_{i=1}^{n} \\
  \bigl\{g_k = 0\bigr\}_k \ \text{equality sub-conjuncts of } \Phi \text{ in wire variables} \\
  \lambda_i \in \Q_{\geq 0} \\
  \nu_k \in \Q \\
  c \equiv \textstyle\sum_i \lambda_i e_i + \sum_k \nu_k g_k \ \text{(exact canonical polynomial)} \\
  \Psi \in \bigl\{\, c \leq 0, \;\; e \geq 0 \text{ where } c = -e, \;\; l \leq r \text{ where } c = l - r \,\bigr\}
}{
  \Phi \vdash \Psi
}
\end{mathpar}
$\Phi \supseteq$ each $\Phi_i$ conjunct-wise (checked); the $\nu_k$ are free-sign exactly as \rulename{Farkas} frees $\lambda_j$ on $=$ atoms.

\ruleflavor{free-sign $\nu_k$ is Lemma A\textquotesingle s CommRing sign-blindness, one theorem for both dialects. Lifts \rulename{Farkas} in exactly two ways: atoms from sub-\emph{proofs} (not sub-conjuncts), atoms of arbitrary polynomial degree. The conclusion is exposed in any of the three \texttt{syzygyConclusionForm} readings --- one exact $\Q$ identity under a sign flip; the transport is the ge-bridge ring lemma, applied once inside \texttt{sound\_syzygy}. A Bernstein Logic rule outright: the baked base theory lists it core (its wire node decoded by the certified checker, its combination relation pinned); \textquotedblleft assume--guarantee\textquotedblright{} names its purpose, not a separate layer.}
\par\medskip
\rulename{Syzygy} is the calculus\textquotesingle s composition rule: it
combines already-proven component bounds into a bound on a weighted sum.
Each component sequent $\Phi_i \vdash e_i \leq 0$ is a sub-proof ---
typically ``this subsystem dissipates its storage function'' --- and the
composite premise $\Phi$ carries, among its conjuncts,
\emph{equality atoms} $g_k = 0$ describing how the subsystems are wired
together (power-preserving interconnection: what one port supplies,
another consumes). A producer reaches for the rule wherever a
system-level argument says ``add the component certificates with weights
$\lambda_i$ and let the internal supplies
cancel'': the weights are found by whatever
method the producer likes (an LP, typically), and only the resulting
exact identity ships in the certificate.

What the checker recomputes, given the payload (combination term,
component list, conclusion term): that every entry weight is
${}\geq 0$ --- premise \emph{equality} atoms expose both polarities to
the matcher, so a free-sign multiplier on an equality is realized as a
nonnegative weight on the appropriate polarity, and that freedom is what
makes them syzygies; that every entry
of the combination really is a component conclusion or a premise equality
sub-conjunct, matched by exact canonical polynomial; and that the
weighted sum of those polynomials is \emph{identically} the conclusion
polynomial as canonical forms over $\Q$ --- coefficient-for-coefficient
equality, not evaluation at test points. It further checks the
bookkeeping hypotheses: each component\textquotesingle s premise is
contained conjunct-wise in $\Phi$ (\texttt{preContains}), the conclusion
term has a determinate finite shape (\texttt{conclShape},
\texttt{conclFin}), and every wire in sight has positive size
(\texttt{wiresPos}, \texttt{compPostWiresPos}).

Each hypothesis is load-bearing. Drop the sign condition on the
inequality weights and a producer could multiply a bound by $-1$ and
conclude its negation. Drop the entry-sourcing check and the
``combination'' could smuggle in atoms
nobody proved. Drop the exact-identity check and the conclusion would be
unrelated to the premises. Drop premise containment and a component
proved over its own operating box could be applied outside it. Drop the
positive-size conditions and a component conclusion over a zero-size wire
--- vacuously true, since it quantifies over no tensor entries --- could
enter the combination with full algebraic weight and push the composite
to a falsehood. The proof section walks these in order.

\cataloguenote{Port-Hamiltonian reading: with $e_i = \mathcal{L}_i V_i -
s_i^{\mathrm{ext}} - s_i^{\mathrm{int}}$, the equality atoms $g_k = 0$
cancel the internal supplies and the conclusion is a dissipation
inequality for $\sum_i \lambda_i V_i$ --- Tellegen\textquotesingle s
theorem in certificate form. The free-body diagram is the junction tree;
internal forces at the cuts are the separator algebra.
}
\ruleblock{Lean statement}
\declhead{BernsteinLogic.Soundness}{SyzygyStep}{\badge{BadgeGreen}{statement bound}}
\begin{leansrc}
/-- Applicability data of the Syzygy rule over a hypernet — field-for-field
the reference checker's obligations, and field-for-field `PSyzygyData`
minus the component-proof list (validity hypotheses of `sound_syzygy`)
and `denBoundLog2` (the Q7 family; zero proof-theoretic weight in v1).
B13-shaped conclusion fields and B14 (`wiresPos`) are inherited verbatim
from `PFarkasData`; `compPostWiresPos` is the honest EXTENSION of B14 to
component-conclusion wires (Composition.lean header: without it a
vacuously-held component conclusion over a zero-size wire poisons the
combination — counterexample kernel-checkable). -/
structure SyzygyStep (H : Hypernet) where
  seq : Sequent
  components : List Sequent
  combination : Term
  conclusion : Term
  preContains : ∀ c ∈ components, premiseContains seq.pre c.pre
  conclShape : ∃ s, conclusion.bshape H = some s
  conclFin : conclusion.finWires H
  wiresPos : ∀ w ∈ seq.pre.freeWires ++ conclusion.freeWires,
      ∃ τ, H.wireType w = some τ ∧ 0 < Shape.totalSize τ.shape
  compPostWiresPos : ∀ c ∈ components, ∀ w ∈ c.post.freeWires,
      ∃ τ, H.wireType w = some τ ∧ 0 < Shape.totalSize τ.shape
  valid : syzygyCombinationValid combination components seq.pre conclusion
  postForm : syzygyConclusionForm conclusion seq.post
\end{leansrc}
\declhead{BernsteinLogic.Soundness}{holds\_of\_syzygyConclusionForm}{\badge{BadgeGreen}{statement bound}}
\begin{leansrc}
/-- The ge-bridge ring lemma (the both-polarity core change): every
reading admitted by
`syzygyConclusionForm conclusion` holds whenever the zero-normal atom
`conclusion ≤ 0` holds — the readings are the same exact ℚ identity
re-presented under a sign flip. The zero-form is the identity transport,
the ge-polarity flip is `-e ≤ 0 → e ≥ 0`, and the literal atom is
`l - r ≤ 0 → l ≤ r` (ArithAtom's tensor-general sub-zero bridge).
Applied once inside `sound_syzygy`, so PSyzygy (tier-1), PBound's
compact spelling (tier-3), and any future ≤-0-shaped fact rule inherit
all readings with no separate polarity justification. -/
theorem holds_of_syzygyConclusionForm {H : Hypernet} {σ : H.Valuation}
    {conclusion : Term} {F : Formula}
    (hle : Formula.holds σ (.atom .le conclusion (.const 0)))
    (hform : syzygyConclusionForm conclusion F) :
    Formula.holds σ F
\end{leansrc}
\declhead{BernsteinLogic.Soundness}{sound\_syzygy}{\badge{BadgeGreen}{statement bound}}
\begin{leansrc}
/-- Soundness of Syzygy (= Lemma B, dialect D): valid components give a
valid composite. The argument is `sound_farkas`'s, with the per-entry
source discharge generalized: a component-sourced entry's atom holds by
component validity (through `premiseContains`), a wiring-sourced entry's
atom holds as a premise sub-conjunct; either way
`atomLeZeroPolys_sound_entryC` delivers per-entry nonpositivity at every
clamped per-entry valuation, and the nonnegative fold closes. The
zero-normal bound then reaches `seq.post` through whichever
`syzygyConclusionForm` reading the payload declared (both-polarity core
change, 2026-07-18). -/
theorem sound_syzygy {H : Hypernet} (st : SyzygyStep H)
    (hcomps : ∀ c ∈ st.components, c.semanticallyValid H) :
    st.seq.semanticallyValid H
\end{leansrc}
\declhead{BernsteinLogic.Soundness}{atomLeZeroPolys\_sound\_entryC}{\badge{BadgeGreen}{statement bound}}
\begin{leansrc}
/-- Every `(e ≤ 0)`-normal-form polynomial a held atom entails evaluates
nonpositively at EVERY clamped per-entry valuation, given positive-size
wire types (general shapes — heterogeneous atoms read through the
clamp; the positivity is PROPOSED-B14's honest content: zero-size atoms
hold vacuously and entail nothing). -/
theorem atomLeZeroPolys_sound_entryC {H : Hypernet} (σ : H.Valuation)
    {F : Formula} (hF : Formula.holds σ F)
    (hw : ∀ w ∈ F.freeWires, ∃ τ, H.wireType w = some τ ∧
      0 < Shape.totalSize τ.shape)
    (hix : F.atomSidesIdxFree = true) (i : ℕ) :
    ∀ pa ∈ atomLeZeroPolys F, QArith.Poly.eval (σ.entryValC i) pa ≤ 0
\end{leansrc}
\declhead{BernsteinLogic.Soundness}{wire\_discipline\_of\_bshape\_fin}{\badge{BadgeGreen}{statement bound}}
\begin{leansrc}
/-- The PROPOSED-B13 field pair (`bshape` + `finWires`) yields the wire
discipline that per-entry evaluation consumes: every free wire reads a
finite tensor of shape `[]` or `s`. -/
theorem wire_discipline_of_bshape_fin {H : Hypernet} (σ : H.Valuation)
    {t : Term} {s : Shape} (hb : t.bshape H = some s)
    (hf : t.finWires H) (hix : t.idxFree = true) :
    ∀ w ∈ t.freeWires, ∃ qw, σ.wireQ w = some qw ∧
      (qw.shape = [] ∨ qw.shape = s)
\end{leansrc}
\declhead{BernsteinLogic.Soundness}{Term.evalQ\_total\_singleton}{\badge{BadgeGreen}{statement bound}}
\begin{leansrc}
/-- Total singleton evaluation: if every free wire reads as a finite
singleton tensor, the term evaluates to a singleton tensor whose scalar
is the polynomial evaluation at the valuation read-out. (Definedness half
of the PROPOSED-B13 Farkas fields, and the workhorse behind every scalar
extraction below.) -/
theorem Term.evalQ_total_singleton {H : Hypernet} (σ : H.Valuation) :
    ∀ (t : Term),
    (∀ w ∈ t.freeWires, ∃ qw, σ.wireQ w = some qw ∧ SingletonShape qw.shape) →
    t.idxFree = true →
    ∃ q, Term.evalQ σ t = some q ∧ SingletonShape q.shape ∧
      q.scalar = QArith.Poly.eval σ.rho t.toPoly
\end{leansrc}
\par\medskip
\textbf{What \texttt{syzygyCombinationValid} means --- and what it does
not.}\enspace The rule\textquotesingle s side condition appears in the
payload structure above as a single opaque name, so this book binds its
meaning by displaying the definition chain in full, from the checker
sources. Reading the chain: \texttt{syzygyCombinationValid} is precisely
the proposition ``\texttt{syzygyCheck} returns
\texttt{true}'', and \texttt{syzygyCheck}
is the disjunction of three branches of one algorithm (a branch for
index-free terms, a wire-coordinate branch for indexed terms, and a
mixed tagged-union branch for heterogeneous
combinations of both kinds). All three
branches do the same three things: decode the combination term into
weighted entries (\texttt{farkasEntries} --- a fixed canonical encoding,
anything else is rejected); require each weight nonnegative and each
entry\textquotesingle s polynomial to occur among the
$({\cdot} \leq 0)$-normal polynomials of a \emph{component conclusion} or
of an \emph{equality sub-conjunct of the premise}
(\texttt{atomLeZeroPolys} --- note that equality atoms contribute both
polarities, which is exactly the free-sign multiplier discipline); and
require the weighted fold of those polynomials to equal the conclusion
term\textquotesingle s canonical polynomial.

What a skeptical reader should confirm by eye in the definitions below:
the conclusion term enters the check \emph{only} as the target of the
exact-identity comparison (and through a syntactic index-freedom gate);
the step\textquotesingle s posted formula \texttt{seq.post} does not
appear in the check at all; and nothing in the chain evaluates, assumes,
or otherwise consults the \emph{truth} of the conclusion --- the check
manipulates polynomials, never valuations, so it could not covertly
``assert the conclusion is
true''. \texttt{seq.post} connects to the
conclusion term only through the separate, purely syntactic field
\texttt{postForm} --- the three readings of $c \leq 0$ enumerated by
\texttt{syzygyConclusionForm} below --- whose semantic transport is the
separately proven ring lemma \texttt{holds\_of\_syzygyConclusionForm}
quoted above. Semantic truth is established only in
\texttt{sound\_syzygy}, from the component-validity hypotheses.

\declhead{BernsteinLogic.Core.Catalogue}{syzygyConclusionForm}{\badge{BadgeGreen}{definition bound}}
\begin{leansrc}
def syzygyConclusionForm (conclusion : Term) (F : Formula) : Prop :=
  F = .atom .le conclusion (.const 0)
  ∨ (∃ e, conclusion = .neg e ∧ F = .atom .ge e (.const 0))
  ∨ (∃ l r, conclusion = .sub l r ∧ F = .atom .le l r)
\end{leansrc}
\declhead{BernsteinLogic.Core.Catalogue}{atomLeZeroPolys}{\badge{BadgeGreen}{definition bound}}
\begin{leansrc}
/-- The `(e ≤ 0)`-normal-form polynomials entailed by one premise atom
(exact polynomial forms — never rescaled; F7 discipline). -/
def atomLeZeroPolys : Formula → List QArith.Poly
  | .atom .le l r => [QArith.Poly.sub l.toPoly r.toPoly]
  | .atom .lt l r => [QArith.Poly.sub l.toPoly r.toPoly]
  | .atom .ge l r => [QArith.Poly.sub r.toPoly l.toPoly]
  | .atom .gt l r => [QArith.Poly.sub r.toPoly l.toPoly]
  | .atom .eq l r => [QArith.Poly.sub l.toPoly r.toPoly,
                      QArith.Poly.sub r.toPoly l.toPoly]
  | _ => []
\end{leansrc}
\declhead{BernsteinLogic.Core.Catalogue}{farkasEntries}{\badge{BadgeGreen}{definition bound}}
\begin{leansrc}
/-- Decode the combination-coefficient Term (pinned canonical encoding:
right-nested `.add` chain of `.mul (.const λⱼ) eⱼ`, `.const 0` end). -/
def farkasEntries : Term → Option (List (Rat × Term))
  | .const c => if c = 0 then some [] else none
  | .add (.mul (.const l) e) rest =>
    (farkasEntries rest).map ((l, e) :: ·)
  | _ => none
\end{leansrc}
\declhead{BernsteinLogic.Core.Catalogue}{syzygyCheckCore}{\badge{BadgeGreen}{definition bound}}
\begin{leansrc}
/-- The Syzygy combination check — `farkasCheck`'s shape, lifted the
AG draft's two ways: entry atoms may come from COMPONENT CONCLUSIONS
(the sub-proofs) as well as from equality sub-conjuncts of the composite
premise; atom polynomials are unrestricted in degree. Entry encoding,
fold, exactness discipline, and the G1 idx-freedom gates are inherited
verbatim from `farkasCheck`. -/
def syzygyCheckCore (cc : Term) (components : List Sequent) (premise : Formula)
    (conclusion : Term) : Bool :=
  match farkasEntries cc with
  | none => false
  | some entries =>
    conclusion.idxFree &&
    entries.all (fun (l, e) =>
      decide (0 ≤ l) &&
      ((components.any fun c =>
          c.post.atomSidesIdxFree &&
          ((atomLeZeroPolys c.post).any fun pa => decide (pa = e.toPoly))) ||
       (premise.conjunctList.any fun a =>
          a.isEqAtom && a.atomSidesIdxFree &&
          ((atomLeZeroPolys a).any fun pa => decide (pa = e.toPoly))))) &&
    decide ((entries.foldr
        (fun (l, e) acc => QArith.Poly.add (QArith.Poly.scale l e.toPoly) acc)
        []) = conclusion.toPoly)
\end{leansrc}
\declhead{BernsteinLogic.Core.Catalogue}{syzygyCheckW}{\badge{BadgeGreen}{definition bound}}
\begin{leansrc}
/-- The idx-form branch of the Syzygy check: entries
match component conclusions or premise equality sub-conjuncts by their
wire-aware coordinate polynomials — the multi-wire composites (DP
Syzygy/SmallGain, distinct region wires per component) are exactly why
`wireCoordKey` keys variables by (wire, coordinate). -/
def syzygyCheckW (cc : Term) (components : List Sequent) (premise : Formula)
    (conclusion : Term) : Bool :=
  match farkasEntries cc with
  | none => false
  | some entries =>
    match farkasEntriesW entries with
    | none => false
    | some pentries =>
      match conclusion.asCoordPolyIdxW? with
      | none => false
      | some pc =>
        pentries.all (fun le =>
          decide (0 ≤ le.1) &&
          ((components.any fun c =>
              (atomLeZeroPolysW c.post).any fun pa => decide (pa = le.2)) ||
           (premise.conjunctList.any fun a =>
              a.isEqAtom &&
              ((atomLeZeroPolysW a).any fun pa => decide (pa = le.2))))) &&
        decide ((pentries.foldr
            (fun le acc => QArith.Poly.add (QArith.Poly.scale le.1 le.2) acc)
            []) = pc)
\end{leansrc}
\declhead{BernsteinLogic.Core.Catalogue}{syzygyCheckM}{\badge{BadgeGreen}{definition bound}}
\begin{leansrc}
/-- The mixed-variable branch of the Syzygy check: entries match
component conclusions or premise equality sub-conjuncts by their MIXED
tagged polynomials — heterogeneous composites (a chart wire's TIdx
coordinates and bare scalar interface wires in one combination) are
exactly why the tagged union space exists. -/
def syzygyCheckM (cc : Term) (components : List Sequent) (premise : Formula)
    (conclusion : Term) : Bool :=
  match farkasEntries cc with
  | none => false
  | some entries =>
    match farkasEntriesM entries with
    | none => false
    | some pentries =>
      match conclusion.asCoordPolyMixed? with
      | none => false
      | some pc =>
        pentries.all (fun le =>
          decide (0 ≤ le.1) &&
          ((components.any fun c =>
              (atomLeZeroPolysM c.post).any fun pa => decide (pa = le.2)) ||
           (premise.conjunctList.any fun a =>
              a.isEqAtom &&
              ((atomLeZeroPolysM a).any fun pa => decide (pa = le.2))))) &&
        decide ((pentries.foldr
            (fun le acc => QArith.Poly.add (QArith.Poly.scale le.1 le.2) acc)
            []) = pc)
\end{leansrc}
\declhead{BernsteinLogic.Core.Catalogue}{syzygyCheck}{\badge{BadgeGreen}{definition bound}}
\begin{leansrc}
/-- The Syzygy combination check: legacy idx-free branch OR the
wire-aware idx coordinate branch OR the mixed tagged branch (mirrors
`farkasCheck`'s three-way split). -/
def syzygyCheck (cc : Term) (components : List Sequent) (premise : Formula)
    (conclusion : Term) : Bool :=
  syzygyCheckCore cc components premise conclusion ||
  syzygyCheckW cc components premise conclusion ||
  syzygyCheckM cc components premise conclusion
\end{leansrc}
\declhead{BernsteinLogic.Core.Catalogue}{syzygyCombinationValid}{\badge{BadgeGreen}{definition bound}}
\begin{leansrc}
def syzygyCombinationValid (cc : Term) (components : List Sequent)
    (premise : Formula) (conclusion : Term) : Prop :=
  syzygyCheck cc components premise conclusion = true
\end{leansrc}
\ruleblock{Correspondence (Steele display $\leftrightarrow$ Lean)}
\badge{BadgeGreen}{reviewer-certified against the displayed texts}
\sourcenote{The dictionary below is a reviewer-certified statement of how the mechanized statements above interpret the displayed rule: display metavariable $\mapsto$ Lean binder, each display hypothesis $\mapsto$ the premise or side condition carrying it, the display conclusion $\mapsto$ the goal --- with any deliberate difference in width between display and mechanization named outright. Certification is recorded as a hash triple over (the displayed cell, the extracted statement texts, this dictionary\textquotesingle s own text); if any of the three changes, this block turns amber until a reviewer re-certifies. The correspondence claim itself remains editorial --- reviewer-grade, not machine-checked; its machine-checked successor types each soundness statement against the rule as a structured object, and supersedes this block per rule as it lands.}
Judgment frame as in \rulename{Ax}. The statement bundles its data as
\texttt{SyzygyStep} --- field-for-field the reference
checker\textquotesingle s obligations, and field-for-field the payload
minus the component-proof list (which enters \texttt{sound\_syzygy} as
the validity hypothesis \texttt{hcomps}) and minus the
denominator-bound declaration (resource accounting; zero
proof-theoretic weight, so rightly absent from the soundness
statement).
\begin{itemize}
\item the component judgments $\{\Phi_i \vdash e_i \leq 0\}_{i=1}^n$
  $\mapsto$ \texttt{st.components : List Sequent} with
  \texttt{hcomps}; the atoms $e_i \leq 0$ have no binders of their own
  --- each combination entry tagged \textsc{from-component} is checked
  against a component\textquotesingle s conclusion inside the
  combination check, which is where the $e_i$ live.
\item $\Phi \supseteq \Phi_i$ conjunct-wise $\mapsto$
  \texttt{st.preContains} $:$ \texttt{premiseContains seq.pre c.pre},
  per component.
\item the equality sub-conjuncts $\{g_k = 0\}$ of $\Phi$ $\mapsto$ the
  \textsc{from-wiring} entry tag: such an entry\textquotesingle s atom
  is checked to be a sub-conjunct of \texttt{seq.pre} directly.
\item $\lambda_i \in \Q_{\geq 0}$, $\nu_k \in \Q$, and $c \equiv
  \sum_i \lambda_i e_i + \sum_k \nu_k g_k$ $\mapsto$
  \texttt{st.combination} and \texttt{st.conclusion} under
  \texttt{st.valid} $:$ \texttt{syzygyCombinationValid} --- the
  canonical-polynomial identity is recomputed by one of \emph{three}
  certified routes (the idx-free \texttt{syzygyCheckCore} replay, the
  wire-coordinate \texttt{syzygyCheckW}, and the mixed
  \texttt{syzygyCheckM}), each discharged separately in
  \texttt{sound\_syzygy}.
\item the evaluability and positive-size packs $\mapsto$
  \texttt{conclShape}, \texttt{conclFin}, \texttt{wiresPos},
  \texttt{compPostWiresPos} --- displayed inline at
  \rulename{PBound}\textquotesingle s cell (this page\textquotesingle
  s cell leaves them to the caption); \texttt{compPostWiresPos} is
  the honest extension of the positive-size discipline to component
  conclusions, whose absence has a kernel-checkable
  vacuous-conclusion counterexample (the source docstring).
\item $\Psi$\textquotesingle s three readings $\mapsto$
  \texttt{st.postForm} $:$ \texttt{syzygyConclusionForm}, literally
  the displayed three-membered disjunction; the transport of the
  zero-normal bound to every reading is the certified-alongside
  ge-bridge lemma \texttt{holds\_of\_syzygyConclusionForm}, applied
  once here, inherited by \rulename{PBound} and the
  \rulename{PsdCert} trio.
\end{itemize}
Width note: nothing in the statement consumes the port-Hamiltonian
reading --- supply rates, storage functions, and the wiring identity
are producer-side interpretations of which the certified content is
exactly the polynomial identity and its component sourcing; and the
statement is proven at every hypernet and component list, with every
hypothesis consumed.

\ruleblock{Lean verifier receipts}
\sourcenote{module \texttt{BernsteinLogic.Soundness}; bound build \texttt{a20d90e66c13\ldots} (bound from receipt a9c0e34d3e3c9c8c, 2026-08-03T21:30:36Z). A receipt records that named checking layers accepted a build: (module, compiled-artifact SHA-256, environment SHA-256, checker layers), minted when the checking gate ran; the store receipts attesting the bound build are listed with the layers each run credited. Nothing here asks the reader to trust the store: every layer is an independently runnable program over the quoted sources --- rebuild the module and re-run the checkers to reproduce the verdict. Receipts attesting superseded builds credit no layer and are omitted from this rendering (the count of omissions, if any, is noted below).}
\begin{receipttable}
\receiptrow{a9c0e34d3e3c9c8c}{2026-08-03T21:30:36Z}{build+decorrelated}{\badge{BadgeGreen}{current}}
\end{receipttable}
\sourcenote{21 further receipt(s) in the store attest superseded builds of this module and are omitted from this rendering.}
\sourcenote{module \texttt{BernsteinLogic.Core.Catalogue}; bound build \texttt{cbc92e976a1d\ldots} (bound from receipt 5abdb3589707b5df, 2026-08-03T21:30:33Z). A receipt records that named checking layers accepted a build: (module, compiled-artifact SHA-256, environment SHA-256, checker layers), minted when the checking gate ran; the store receipts attesting the bound build are listed with the layers each run credited. Nothing here asks the reader to trust the store: every layer is an independently runnable program over the quoted sources --- rebuild the module and re-run the checkers to reproduce the verdict. Receipts attesting superseded builds credit no layer and are omitted from this rendering (the count of omissions, if any, is noted below).}
\begin{receipttable}
\receiptrow{24804e66edba0bc9}{2026-08-03T20:25:09Z}{build+decorrelated}{\badge{BadgeGreen}{current}}
\receiptrow{5abdb3589707b5df}{2026-08-03T21:30:33Z}{build+decorrelated}{\badge{BadgeGreen}{current}}
\end{receipttable}
\sourcenote{18 further receipt(s) in the store attest superseded builds of this module and are omitted from this rendering.}
\layerline{\layeritem{build}{lake build (elaboration)}{\badge{BadgeGreen}{current}}\layeritem{replay}{leanchecker (official kernel replay)}{\badge{BadgeAmber}{stale only}}\layeritem{hardened}{safe\_verify (hardened screen)}{\badge{BadgeAmber}{stale only}}\layeritem{decorrelated}{lean4export + nanoda (decorrelated kernel)}{\badge{BadgeGreen}{current}}}
\sourcenote{\badge{BadgeGray}{comparator: none for this rule} no fixture in the
three-implementation comparator harness (compiled / kernel-interpreted /
reference-interpreted) exercises a \rulename{Syzygy} node yet; current
fixtures cover \rulename{BernCone} and \rulename{Norm} shapes. The
end-to-end \rulename{Syzygy} certificate
\texttt{DpCompositeSyzygy.lean} is a kernel-checked instance but rides
the handshake harness, not the comparator.
}
\ruleblock{HOL4 statement}
\sourcenote{theory \texttt{bern\_catalogue}, HOL4 development script \texttt{\texttt{bern\_catalogueScript.sml}} \badge{BadgeGreen}{script bound}}
\holhead{Definition syzygyCheckCore\_def}
\begin{holsrc}
Definition syzygyCheckCore_def:
  syzygyCheckCore cc components premise conclusion =
    case farkasEntries cc of
    | NONE => F
    | SOME entries =>
        termIdxFree conclusion ∧
        EVERY (λ(l,e).
          0 ≤ l ∧
          (EXISTS (λc. atomSidesIdxFree c.post ∧
                       EXISTS (λpa. pa = termToPoly e) (atomLeZeroPolys c.post))
                  components ∨
           EXISTS (λa. isEqAtom a ∧ atomSidesIdxFree a ∧
                       EXISTS (λpa. pa = termToPoly e) (atomLeZeroPolys a))
                  (conjunctList premise))) entries ∧
        (FOLDR (λ(l,e) acc. polyAdd (polyScale l (termToPoly e)) acc)
               [] entries =
         termToPoly conclusion)
End
\end{holsrc}
\holhead{Definition syzygyCheck\_def}
\begin{holsrc}
Definition syzygyCheck_def:
  syzygyCheck cc components premise conclusion ⇔
    syzygyCheckCore cc components premise conclusion ∨
    syzygyCheckW cc components premise conclusion
End
\end{holsrc}
\ruleblock{HOL verifier receipts}
\badge{BadgeBlue}{documentary} \texttt{hol-strict} CLEAN over the ported checker chain: theory
\texttt{bern\_catalogue}, 472 theorems, 0 cheat-tagged, 0 nonstandard
axioms (whitelist: \texttt{DISK\_THM}; the four \texttt{bool} axioms) ---
verification run of 2026-07-22 on an isolated build host, scripts
sha-matched verbatim to the sources cited here (git blob
\texttt{e41b9ed0\ldots}). Downstream, the same chain carries the checker
to a CakeML binary with the \texttt{machine\_code\_sound} tag
\texttt{(["DISK\_THM"], [])} (theory \texttt{bernCertProofARM8}).
Family-level bridge lemmas for the \rulename{Syzygy} node
(\texttt{famEq\_syz}, \texttt{bern\_fam\_syz}) ride an in-flight assembly
and are not yet in the census table. As on the \rulename{Farkas} and
\rulename{PBound} pages: the Lean check has since grown a third (mixed
tagged-union) branch, added 2026-07-19 after this port run --- the
divergence is refusal-direction only, and its port is named residue.
Semantic stratum, also ported: the soundness twins
\texttt{sound\_syzygy} (the proof-object form) and
\texttt{sound\_pBound} (theory \texttt{bern\_tr\_soundsyz}: \texttt{hol-
strict} clean, 7 theorems, 0 offences, verification run of 2026-07-22 on
an isolated build host, content-addressed at the same sources); the
certified-fragment keystone (\texttt{checkPf\_sound\_frag}, theory
\texttt{bern\_tr\_splice}, 20 theorems, 0 offences, same run) discharges
this rule\textquotesingle s wire arm through the latter.

\sourcenote{HOL receipts are documentary: quoted results of \texttt{hol-strict} runs
and census rows from the HOL4 development\textquotesingle s build
evidence, not minted receipt artefacts. A HOL receipt store (the Lean
store\textquotesingle s twin) is named future work.
}
\ruleblock{Proof}

\sourcenote{Genre: core rule (tier 1) --- an ordinary mathematical proof at
Putnam-exam-solution detail. Every hypothesis of the rule is used below and
remarked on; the formal-anchor table maps each step to the Lean declaration
that carries it.}

\textbf{Setting and claim.}\enspace Fix a hypernet $H$. A \emph{valuation}
$\sigma$ assigns each wire a rational tensor of its declared shape; a
sequent $\Phi \vdash \Psi$ is \emph{semantically valid} when every
admissible $\sigma$ satisfying $\Phi$ satisfies $\Psi$. Atoms
$e \leq 0$ over tensor-shaped terms hold per-entry: every entry of the
evaluated tensor is $\leq 0$. The rule's data (the \texttt{SyzygyStep}
payload quoted above) supplies component sequents
$\Phi_i \vdash e_i \leq 0$ for $i = 1,\dots,n$, a combination certificate
with weights $\lambda_i \in \Q_{\geq 0}$ on the components and free
weights $\nu_k \in \Q$ on equality sub-conjuncts $g_k = 0$ of $\Phi$, a
conclusion term $c$, and a checked \emph{syntactic} identity
\begin{equation}
c \;\equiv\; \sum_{i} \lambda_i\, e_i \;+\; \sum_{k} \nu_k\, g_k
\qquad\text{(equality of exact canonical polynomials over } \Q\text{).}
\tag{$\ast$}
\end{equation}
The claim (\texttt{sound\_syzygy}): if every component sequent is
semantically valid, the composite sequent $\Phi \vdash \Psi$ is
semantically valid, where $\Psi$ is any reading of $c \leq 0$ admitted by
\texttt{syzygyConclusionForm} (Lemma~\ref{lem:syz-transport}).

\begin{blemma}[nonnegative combination]\label{lem:syz-comb}
In an ordered field, if $a_i \leq 0$ and $\lambda_i \geq 0$ for
$i = 1,\dots,n$, and $b_k = 0$ for $k = 1,\dots,m$ with $\nu_k$ arbitrary,
then $\sum_i \lambda_i a_i + \sum_k \nu_k b_k \leq 0$.
\end{blemma}
\begin{proof}
Each $\lambda_i a_i \leq 0$: multiplication by a nonnegative element
preserves $\leq$ (applied to $a_i \leq 0$), so
$\lambda_i a_i \leq \lambda_i \cdot 0 = 0$. Each $\nu_k b_k = \nu_k \cdot 0
= 0$ \emph{regardless of the sign of $\nu_k$} --- this sign-blindness on
equality atoms is precisely the freedom classical Farkas variants grant
multipliers of equality constraints \cite[\S7.3]{schrijver1986}. A finite
sum of nonpositive elements is nonpositive, by induction on $n + m$.
\end{proof}

\begin{blemma}[the certified identity evaluates]\label{lem:syz-eval}
Let $(\ast)$ hold as an identity of canonical polynomials, and let
$\sigma$ be a valuation at which $c$ and every $e_i$, $g_k$ evaluate
(per-entry, through the clamped per-entry valuation for tensor shapes).
Then the rational identity
$c(\sigma) = \sum_i \lambda_i\, e_i(\sigma) + \sum_k \nu_k\, g_k(\sigma)$
holds at every entry.
\end{blemma}
\begin{proof}
First fix what ``per-entry'' means. For an entry index $j$ of the
conclusion's shape, the \emph{clamped per-entry valuation} reads each
wire at position $j$ when $j$ is in range for that wire's shape and at
its entry $0$ otherwise --- total on every wire in sight, because the
B14 positivity pack guarantees each carries at least one entry, and
agreeing with the plain positional read wherever $j$ is in range. No
shape-agreement hypothesis is needed beyond that: a sourced
polynomial's wires share its own atom's shape (rank-0 broadcast
aside), so they clamp in step, and its clamped value at $j$ is an
honest entry of its own evaluated tensor --- which is where Step 4's
per-entry case split picks it up. Evaluation at this scalar assignment
is a ring homomorphism $\Q[\text{wires}] \to \Q$; a syntactic
identity of canonical forms is an identity of polynomials, and ring
homomorphisms preserve polynomial identities. Two remarks make this step deserve its own lemma. First,
$(\ast)$ is checked by \emph{exact canonical-form equality} over $\Q$
(\texttt{syzygyCombinationValid}, replayed through \texttt{toPoly}), not by
evaluation at test points --- so no completeness, degree bound, or
interpolation argument is needed anywhere in this proof; soundness is pure
arithmetic. Second, the checker admits three certified routes to the same
fact (the idx-free branch, the wire-coordinate branch, and the mixed
tagged-union branch:
\texttt{syzygyCheck = Core $\vee$ W $\vee$ M}); whichever branch
fires, the certified fact consumed downstream is the same --- the
identity $(\ast)$, under the same per-entry reading. The branches
differ in how terms are decoded, never in what is certified.
\end{proof}

\begin{blemma}[reading transport]\label{lem:syz-transport}
Suppose $\sigma$ satisfies the zero-normal atom $c \leq 0$. Then $\sigma$
satisfies every formula $F$ admitted by
$\mathrm{syzygyConclusionForm}(c, F)$, namely:
\emph{(Z$\leq$)} $F = (c \leq 0)$ itself;
\emph{(G)} $F = (e \geq 0)$ where $c$ is \emph{literally} the
negation term $-e$;
\emph{(L$\leq$)} $F = (l \leq r)$ where $c$ is literally the
difference term $l - r$ --- structural identities of terms, not merely
canonical-polynomial ones.
\end{blemma}
\begin{proof}
(Z$\leq$) is the identity transport. (G): per entry,
$-e(\sigma) \leq 0 \iff e(\sigma) \geq 0$ (negate and flip, an ordered-field
axiom). (L$\leq$): per entry, $l(\sigma) - r(\sigma) \leq 0 \iff l(\sigma)
\leq r(\sigma)$ (add $r(\sigma)$ to both sides). Structurality is load-bearing:
because $e$, $l$, $r$ are subterms of $c$, the evaluability that the
zero-normal atom's holding certifies for $c$ passes to them --- a
merely canonical identity would not carry evaluation (a cancelled
ill-typed subterm vanishes from the polynomial yet still poisons the
term), which is why the admitted readings are stated on term
structure. This is the ge-bridge ring
lemma \texttt{holds\_of\_syzygyConclusionForm} (the both-polarity
conclusion-form design): the readings are one exact $\Q$ identity
re-presented under a sign flip, so the polarity bridge lives \emph{once},
here, and every $\leq 0$-shaped rule (\rulename{Syzygy} itself,
\rulename{PBound}'s compact spelling, the PsdCert family) inherits all
readings with no separate justification.
\end{proof}

\begin{bthm}[soundness of \rulename{Syzygy} = Lemma B, dialect D]
\label{thm:syz-sound}
Under the payload's checked hypotheses, if every component sequent
$\Phi_i \vdash e_i \leq 0$ is semantically valid, then so is
$\Phi \vdash \Psi$.
\end{bthm}
\begin{proof}
Fix an admissible valuation $\sigma$ with $\sigma \models \Phi$; we must
show $\sigma \models \Psi$.

\emph{Step 1 (components hold at $\sigma$).} For each $i$, the payload
checks $\Phi \supseteq \Phi_i$ conjunct-wise (\texttt{premiseContains}
--- an ordered-tree containment certificate, and satisfaction descends
along it by one $\wedge$-projection per step); hence
$\sigma \models \Phi_i$, and by the assumed semantic validity of the
$i$-th component, $\sigma \models e_i \leq 0$: every entry of
$e_i(\sigma)$ is $\leq 0$.

\emph{Step 2 (wiring identities hold at $\sigma$).} Each $g_k = 0$ is an
equality sub-conjunct of $\Phi$, so $\sigma \models g_k = 0$ directly:
$g_k(\sigma) = 0$ at every entry.

\emph{Step 3 (the conclusion evaluates).} The B13 shape fields
(\texttt{conclShape}, \texttt{conclFin}) guarantee $c$ has a determinate
finite Bernstein shape over $H$, and the B14 positivity fields
(\autoref{rem:syz-b14}) guarantee every wire in sight has positive total
size; together these give a total per-entry evaluation of $c$ at $\sigma$
--- the two-case shape analysis (\,singleton vs.\ indexed\,) in
\texttt{sound\_syzygy}'s body.

\emph{Step 4 (combine).} At each entry, Lemma~\ref{lem:syz-eval} turns the
certified identity $(\ast)$ into the rational identity
$c(\sigma) = \sum_i \lambda_i e_i(\sigma) + \sum_k \nu_k g_k(\sigma)$; the
per-entry source discharge is the case split of Steps 1--2 (a
component-sourced entry's atom holds by component validity, a
wiring-sourced entry's atom holds as a premise sub-conjunct ---
\texttt{atomLeZeroPolys\_sound\_entryC} delivers per-entry nonpositivity in
both cases), and Lemma~\ref{lem:syz-comb} closes the fold:
$c(\sigma) \leq 0$ at every entry.

\emph{Step 5 (read off).} Lemma~\ref{lem:syz-transport} transports the
zero-normal bound $c \leq 0$ to whichever reading
$\mathrm{syzygyConclusionForm}(c, \Psi)$ the payload declared
(\texttt{postForm}). Hence $\sigma \models \Psi$.
\end{proof}

\begin{bremark}[why the B14 hypotheses are not decoration]
\label{rem:syz-b14}
\texttt{wiresPos} and its honest extension \texttt{compPostWiresPos}
require every wire appearing in $\Phi$'s free wires, the conclusion, and
each component conclusion to have positive total size. Drop
\texttt{compPostWiresPos} and the rule is \emph{unsound}: a component
conclusion over a zero-size wire holds \emph{vacuously} (its evaluated
tensor has no entries, so ``every entry $\leq 0$'' is empty), yet its
polynomial enters the combination $(\ast)$ with full algebraic weight ---
a vacuously-held component can then push the composite conclusion to
something false. Concretely: over a hypernet with one zero-size wire $z$, take the
single component $\top \vdash z - z + 1 \leq 0$ --- vacuously valid,
its evaluated tensor having no entries --- with weight $1$ and empty
$\nu$-row; the combination's canonical polynomial is the constant
$1$, so the checked identity admits the conclusion $1 \leq 0$, a
rank-0 falsehood concluded from a vacuous truth. (The counterexample
is also kernel-checkable, recorded at the field's introduction.) The
B13 shape fields have the matching witness one step earlier: with $c$
the term $w - w$ for an untyped wire $w$, the canonical polynomial is
$0$ --- identical to the empty combination's --- yet the conclusion
atom never evaluates and so never holds; \texttt{conclShape} and
\texttt{conclFin} exclude exactly this. This is the classic mismatch between
empty-product semantics and nonempty algebra, caught by making the
degenerate case unrepresentable rather than discouraged.
\end{bremark}

\begin{bremark}[relation to the literature]
The rule is the \emph{sound} (easy) direction of an affine Farkas-type
combination, lifted twice: atoms come from sub-\emph{proofs} rather than
sub-conjuncts, and atoms may have arbitrary polynomial degree. The linear
ancestor is Farkas' lemma \cite{farkas1902}; the textbook treatment of
multiplier sign conventions ($\lambda \geq 0$ on inequalities, free on
equalities) is \cite[\S7]{schrijver1986}. With polynomial atoms the
combination sits in the Krivine--Handelman family of positivity
certificates \cite{krivine1964,handelman1988} --- of which, again, only
the trivial direction is claimed or needed: no Positivstellensatz
completeness enters the kernel. The wiring reading --- components as
open subsystems, $g_k = 0$ as power-preserving interconnection, the
composite as dissipation of $\sum_i \lambda_i V_i$ --- is Tellegen's
theorem \cite{tellegen1952} in certificate form, i.e.\ compositional
dissipativity in the sense of Willems \cite{willems1972} and
Moylan--Hill \cite{moylanhill1978}; the port-Hamiltonian formulation is
surveyed in \cite{vanderschaftjeltsema2014}. Within the accompanying
development the rule is known as Lemma~B (dialect~D); a quantitative
effectivity result for tree-structured interconnections --- sparse
Handelman certificates with running-intersection support --- is developed
in a companion note.
\end{bremark}

\begin{anchors}
Step 1: conjunct-wise premise containment & \texttt{premiseContains} (payload field \texttt{preContains})\\
Steps 1--2, per-entry atom discharge & \texttt{atomLeZeroPolys\_sound\_entryC}; source tags \textsc{from-component}/\textsc{from-wiring}\\
Step 3: conclusion evaluates & \texttt{conclShape}, \texttt{conclFin} (B13); \texttt{wiresPos}, \texttt{compPostWiresPos} (B14); \texttt{wire\_discipline\_of\_bshape\_fin}, \texttt{Term.evalQ\_total\_singleton}\\
Step 4: certified identity, all three checker branches & \texttt{syzygyCombinationValid}, \texttt{syzygyCheck} = \texttt{syzygyCheckCore} $\vee$ \texttt{syzygyCheckW} $\vee$ \texttt{syzygyCheckM}\\
Lemma \ref{lem:syz-comb} fold & the nonnegative fold in \texttt{sound\_syzygy}'s body\\
Lemma \ref{lem:syz-transport} & \texttt{holds\_of\_syzygyConclusionForm}\\
Theorem \ref{thm:syz-sound} & \texttt{sound\_syzygy}; wired as core \texttt{RuleData.syzygy} into \texttt{ProofObject.sound}\\
\end{anchors}


\rulesection{34}{SmallGain}{alias for \textcolor{DiffBlue}{\rulename{Syzygy}}}
\label{blcatrule-SmallGain}
\ruleblock{The rule (Steele metanotation)}
\sourcenote{as realized in the mechanized development; source row: \texttt{\texttt{rules.toml}} --- \emph{SmallGain}. The metanotation statement is typeset by hand for notational quality; its mechanized ground truth is the Lean statement below, extracted from the sources at every build of this book and hash-bound. The rule \emph{roster} is machine-checked (every mechanized rule tag round-trips to exactly one catalogue row); the statement \emph{text} is editorial, checked by eye against the extracted Lean beside it.}
\begin{mathpar}
\inferrule*[Right=\rulename{SmallGain}]{
  \Phi_1 \vdash e_1 \leq 0 \\
  \Phi_2 \vdash e_2 \leq 0 \\
  \Phi_w \vdash \lambda\sigma_2 - (1{-}\theta)\,\alpha_1 \leq 0 \\
  \Phi_w \vdash \sigma_1 - \lambda(1{-}\theta)\,\alpha_2 \leq 0 \\
  \Phi \supseteq \Phi_1,\ \Phi_2,\ \Phi_w \ \text{conjunct-wise (checked)} \\
  \lambda \in \Q_{>0} \\
  \theta \in \Q \cap (0,1)
}{
  \Phi \vdash \mathcal{L}(V_1 + \lambda V_2) + \theta\,(\alpha_1 + \lambda\alpha_2) \leq 0
}
\end{mathpar}
$e_i \equiv \mathcal{L}_i V_i + \alpha_i - \sigma_i$; \quad $m_1 \equiv \lambda\sigma_2 - (1-\theta)\,\alpha_1$, \quad $m_2 \equiv \sigma_1 - \lambda(1-\theta)\,\alpha_2$; \quad conclusion identity, checked as exact canonical polynomial: $e_1 + \lambda e_2 + m_1 + m_2 = \mathcal{L}(V_1 + \lambda V_2) + \theta(\alpha_1 + \lambda\alpha_2)$.

\ruleflavor{the two membership sequents are ordinary BL sequents over the declared operating box $\Phi_w$, their \rulename{BernCone} multipliers \emph{are} the LP witness; in the conclusion identity the cross-supplies cancel by construction --- the checker never sees a fixpoint. In the baked base theory this rule sits in the admissible roster --- no proof-term sort, no wire tag: it elaborates to one \rulename{Syzygy} node at weights $(1, \lambda, 1, 1)$, which this catalogue renders as the alias tier.}
\par\medskip
\rulename{SmallGain} is the calculus\textquotesingle s
feedback-composition rule in its classical costume: two subsystems, each
dissipating its own storage function except for a supply fed by the
other, compose into a single dissipation certificate for
$V_1 + \lambda V_2$ when the loop gain is strictly less than one
($\theta$ is the margin). It is an \emph{alias}: the wire format never
carries a \rulename{SmallGain} node. A producer that finds
$(\lambda, \theta)$ --- an LP feasibility problem over the membership
certificates --- elaborates the step to the single \rulename{Syzygy}
node displayed in the recipe below before emission, so what the checker
recomputes is exactly \rulename{Syzygy}\textquotesingle s obligation
set: the exact conclusion identity at weights $(1, \lambda, 1, 1)$, in
which the cross-supplies cancel by construction. Nothing fixpoint-shaped
survives to checking time; the classical contraction argument lives
entirely on the producer\textquotesingle s side, in the choice of
$(\lambda, \theta)$.

What would break without each hypothesis is inherited from
\rulename{Syzygy} through the unpack --- that
section\textquotesingle s commentary applies verbatim to the unpacked
node. The two conditions native to this rule: $\lambda > 0$, of which
soundness uses only $\lambda \geq 0$ (the strictness backs the
rule\textquotesingle s published meaning), and $\theta \in (0,1)$, which
is well-formedness for $\theta$\textquotesingle s published meaning as a
robustness margin rather than a soundness need --- the membership
sequents self-police, and an infeasible pair simply yields unprovable
memberships.

\cataloguenote{The four component sequents are ordinary sequents over the declared
operating box; their cone-membership multipliers \emph{are} the LP
witness. Signed polynomial contracts (rather than norm gains) are
constitutive, not stylistic: swing-up analyses in the
{\AA}str\"om--Furuta family exploit the \emph{sign} of
$\dot\vartheta\cos\vartheta$, which no norm-gain contract can express.
$(\lambda,\theta)$-existence stays on the producer side, deliberately
outside the checker; the $n$-component generalization is the
\rulename{SmallGainN} section --- the same alias discipline over an
arbitrary interconnection graph, elaborating to one \rulename{Syzygy}
node, of which this rule is the $n = 2$ instance.
}
\ruleblock{Lean statement}
\declhead{BernsteinLogic.Composition}{SmallGainStep}{\badge{BadgeGreen}{statement bound}}
\begin{leansrc}
/-- Applicability data of the SmallGain rule — field-for-field the
draft's `smallGainShapes` obligations (§3.3):

* (a) the dissipation conclusions are `(eᵢ ≤ 0)` atoms — the residual
  reading eᵢ = 𝓛ᵢVᵢ + αᵢ − σᵢ is definitional at the checker level
  (`𝓛ᵢVᵢ := eᵢ − αᵢ + σᵢ`), so no separate field pins it;
* (b) the membership conclusions are `(mᵢ ≤ 0)` atoms whose polynomials
  equal their schema instances at (λ, θ) — exact canonical polynomial;
* (c) λ > 0 and 0 < θ < 1 — three rational comparisons. The θ range is
  NOT consumed by soundness (the memberships self-police; the range backs
  the published guard-modulus semantics — draft §3.2's honesty note);
* (d) the conclusion identity: the (1, λ, 1, 1) fold of the four posted
  polynomials equals the conclusion schema 𝓛(V₁+λV₂) + θ(α₁+λα₂) — the
  checker replays the cancellation (`smallGain_is_syzygy` proves the
  schema coherent: given (b) and the residual reading, (d) is a theorem,
  not a search).

B13/B14 fields as in `SyzygyStep`, component-post coverage included. -/
structure SmallGainStep (H : Hypernet) where
  seq : Sequent
  diss1 : Sequent
  diss2 : Sequent
  mem1 : Sequent
  mem2 : Sequent
  lambda : ℚ
  theta : ℚ
  alpha1 : Term
  alpha2 : Term
  sigma1 : Term
  sigma2 : Term
  lyapTerm : Term
  e1 : Term
  e2 : Term
  m1 : Term
  m2 : Term
  concl : Term
  -- (a) dissipation shapes
  diss1Post : diss1.post = .atom .le e1 (.const 0)
  diss2Post : diss2.post = .atom .le e2 (.const 0)
  -- (b) membership shapes and schemas
  mem1Post : mem1.post = .atom .le m1 (.const 0)
  mem2Post : mem2.post = .atom .le m2 (.const 0)
  mem1Poly : m1.toPoly = Poly.sub (Poly.scale lambda sigma2.toPoly)
      (Poly.scale (1 - theta) alpha1.toPoly)
  mem2Poly : m2.toPoly = Poly.sub sigma1.toPoly
      (Poly.scale (lambda * (1 - theta)) alpha2.toPoly)
  -- G1 rider: the posted residuals and the conclusion are idx-free —
  -- their canonical polynomials carry the rule's meaning (`Term.toPoly`'s
  -- G1 invariant; mirrors `farkasCheck`'s gates)
  e1IdxFree : e1.idxFree = true
  e2IdxFree : e2.idxFree = true
  m1IdxFree : m1.idxFree = true
  m2IdxFree : m2.idxFree = true
  conclIdxFree : concl.idxFree = true
  -- (c) the three rational comparisons
  lambdaPos : 0 < lambda
  thetaLo : 0 < theta
  thetaHi : theta < 1
  -- (d) the conclusion identity (checker-replayed cancellation), and the
  -- conclusion term's tie to the schema
  foldEq : Poly.add e1.toPoly (Poly.add (Poly.scale lambda e2.toPoly)
        (Poly.add m1.toPoly m2.toPoly))
      = Poly.add lyapTerm.toPoly
          (Poly.scale theta
            (Poly.add alpha1.toPoly (Poly.scale lambda alpha2.toPoly)))
  conclPoly : concl.toPoly = Poly.add lyapTerm.toPoly
      (Poly.scale theta
        (Poly.add alpha1.toPoly (Poly.scale lambda alpha2.toPoly)))
  postEq : seq.post = .atom .le concl (.const 0)
  -- premise containment for all four components
  pre1 : premiseContains seq.pre diss1.pre
  pre2 : premiseContains seq.pre diss2.pre
  preM1 : premiseContains seq.pre mem1.pre
  preM2 : premiseContains seq.pre mem2.pre
  -- B13/B14 (component-post coverage per the §R2 finding)
  conclShape : ∃ s, concl.bshape H = some s
  conclFin : concl.finWires H
  wiresPos : ∀ w ∈ seq.pre.freeWires ++ concl.freeWires,
      ∃ τ, H.wireType w = some τ ∧ 0 < Shape.totalSize τ.shape
  compPostWiresPos : ∀ c ∈ [diss1, diss2, mem1, mem2],
      ∀ w ∈ c.post.freeWires,
      ∃ τ, H.wireType w = some τ ∧ 0 < Shape.totalSize τ.shape
\end{leansrc}
\declhead{BernsteinLogic.Composition}{toSyzygy}{\badge{BadgeGreen}{statement bound}}
\begin{leansrc}
/-- The SmallGain node, AS the Syzygy instance it is (weights
(1, λ, 1, 1), components in combination order). -/
def toSyzygy (st : SmallGainStep H) : SyzygyStep H
\end{leansrc}
\declhead{BernsteinLogic.Composition}{sound\_smallGain}{\badge{BadgeGreen}{statement bound}}
\begin{leansrc}
/-- Soundness of SmallGain (= Lemma B, dialect G), as the Syzygy
corollary: the four component validities feed `sound_syzygy` on the
collapsed instance. The three rational comparisons (λ > 0, 0 < θ < 1)
enter only through λ ≥ 0 inside the combination check; θ's range is
well-formedness, exactly as the draft's honesty note states. -/
theorem sound_smallGain {H : Hypernet} (st : SmallGainStep H)
    (h1 : st.diss1.semanticallyValid H)
    (h2 : st.diss2.semanticallyValid H)
    (hm1 : st.mem1.semanticallyValid H)
    (hm2 : st.mem2.semanticallyValid H) :
    st.seq.semanticallyValid H
\end{leansrc}
\declhead{BernsteinLogic.Composition}{toSyzygy\_components}{\badge{BadgeGreen}{statement bound}}
\begin{leansrc}
theorem toSyzygy_components (st : SmallGainStep H) :
    st.toSyzygy.components = st.components
\end{leansrc}
\ruleblock{Correspondence (Steele display $\leftrightarrow$ Lean)}
\badge{BadgeGreen}{reviewer-certified against the displayed texts}
\sourcenote{The dictionary below is a reviewer-certified statement of how the mechanized statements above interpret the displayed rule: display metavariable $\mapsto$ Lean binder, each display hypothesis $\mapsto$ the premise or side condition carrying it, the display conclusion $\mapsto$ the goal --- with any deliberate difference in width between display and mechanization named outright. Certification is recorded as a hash triple over (the displayed cell, the extracted statement texts, this dictionary\textquotesingle s own text); if any of the three changes, this block turns amber until a reviewer re-certifies. The correspondence claim itself remains editorial --- reviewer-grade, not machine-checked; its machine-checked successor types each soundness statement against the rule as a structured object, and supersedes this block per rule as it lands.}
The alias\textquotesingle s Lean vehicle is the payload structure
\texttt{SmallGainStep H} (quoted in full above --- its fields carry the
display\textquotesingle s metavariables), the verified unpack
\texttt{toSyzygy}, and the soundness corollary
\texttt{sound\_smallGain}.
\begin{itemize}
\item $\lambda \mapsto$ \texttt{lambda} $: \Q$ with \texttt{lambdaPos}
  $: 0 < \lambda$; $\theta \mapsto$ \texttt{theta} with
  \texttt{thetaLo}, \texttt{thetaHi} $: 0 < \theta < 1$; $\alpha_i,
  \sigma_i \mapsto$ \texttt{alpha1}/\texttt{alpha2},
  \texttt{sigma1}/\texttt{sigma2}; the composite storage $V_1 + \lambda
  V_2 \mapsto$ \texttt{lyapTerm}.
\item the dissipation premises $\Phi_i \vdash e_i \leq 0$ $\mapsto$
  \texttt{diss1}/\texttt{diss2} with the shape ties
  \texttt{diss1Post}/\texttt{diss2Post} (each $e_i$ is a field; its
  published spelling $\mathcal{L}_i V_i + \alpha_i - \sigma_i$ lives on
  the producer\textquotesingle s side). The operating-box membership
  premises $\mapsto$ \texttt{mem1}/\texttt{mem2} with
  \texttt{mem1Poly}/\texttt{mem2Poly} tying their residuals to the
  displayed schemas $\lambda\sigma_2 - (1{-}\theta)\alpha_1$ and
  $\sigma_1 - \lambda(1{-}\theta)\alpha_2$ as canonical polynomials.
\item the display\textquotesingle s side text --- the exact conclusion
  identity (the cross-supply-cancellation gloss is the flavor
  text\textquotesingle s) --- $\mapsto$ \texttt{foldEq}
  and \texttt{conclPoly} (the checker-replayed cancellation, exact over
  canonical polynomials) plus \texttt{postEq} (the conclusion tie);
  the displayed containment premise ($\Phi \supseteq \Phi_1, \Phi_2,
  \Phi_w$ conjunct-wise) $\mapsto$
  \texttt{pre1}/\texttt{pre2}/\texttt{preM1}/\texttt{preM2}.
\item conclusion $\Phi \vdash \mathcal{L}(V_1 + \lambda V_2) +
  \theta\,(\alpha_1 + \lambda\alpha_2) \leq 0$ $\mapsto$ the goal
  \texttt{st.seq.semanticallyValid H} of \texttt{sound\_smallGain},
  whose four hypotheses are exactly the component validities.
\item the alias artefacts: the unpack function is \texttt{toSyzygy}
  (weights $(1, \lambda, 1, 1)$, components in combination order); the
  soundness corollary routes through the target rule\textquotesingle s
  lemma on the collapsed instance --- the displayed rule never reaches
  the wire.
\end{itemize}
Width notes, both directions: soundness uses only $\lambda \geq 0$ of
\texttt{lambdaPos} (the strictness backs the rule\textquotesingle s
published meaning), and $\theta$\textquotesingle s range is
well-formedness rather than a soundness need --- both as the
commentary states. Conversely the structure carries idx-freeness and
wire-shape fields (\texttt{e1IdxFree} through
\texttt{compPostWiresPos}) that the display leaves to the ambient
calculus\textquotesingle s conventions: the mechanization is narrower
there, naming every side condition the checker replays.

\ruleblock{Lean verifier receipts}
\sourcenote{module \texttt{BernsteinLogic.Composition}; bound build \texttt{9116f454beb0\ldots} (bound from receipt 548ea3dda8ed73a2, 2026-08-03T21:30:29Z). A receipt records that named checking layers accepted a build: (module, compiled-artifact SHA-256, environment SHA-256, checker layers), minted when the checking gate ran; the store receipts attesting the bound build are listed with the layers each run credited. Nothing here asks the reader to trust the store: every layer is an independently runnable program over the quoted sources --- rebuild the module and re-run the checkers to reproduce the verdict. Receipts attesting superseded builds credit no layer and are omitted from this rendering (the count of omissions, if any, is noted below).}
\begin{receipttable}
\receiptrow{548ea3dda8ed73a2}{2026-08-03T21:30:29Z}{build+decorrelated}{\badge{BadgeGreen}{current}}
\end{receipttable}
\sourcenote{10 further receipt(s) in the store attest superseded builds of this module and are omitted from this rendering.}
\layerline{\layeritem{build}{lake build (elaboration)}{\badge{BadgeGreen}{current}}\layeritem{replay}{leanchecker (official kernel replay)}{\badge{BadgeAmber}{stale only}}\layeritem{hardened}{safe\_verify (hardened screen)}{\badge{BadgeAmber}{stale only}}\layeritem{decorrelated}{lean4export + nanoda (decorrelated kernel)}{\badge{BadgeGreen}{current}}}
\sourcenote{\badge{BadgeGray}{comparator: structurally n/a} the alias unpacks
\emph{before} emission (no rule tag; producers elaborate a
\rulename{SmallGain} step to its \rulename{Syzygy} node), so no wire
certificate --- and hence no comparator fixture --- can ever contain a
\rulename{SmallGain} node. Comparator coverage for this rule \emph{is}
\rulename{Syzygy}\textquotesingle s. The end-to-end instance
\texttt{DpCompositeSmallGain.lean} is kernel-checked upstream of the
unpack.
}
\ruleblock{HOL4 statement}
\absentbox{OUT OF SCOPE BY DESIGN}{Out of scope for the HOL side \emph{by design}, not by backlog: the wire
format never carries a \rulename{SmallGain} node (no tag --- producers
elaborate to the $(1,\lambda,1,1)$ \rulename{Syzygy} node before
emission, warrant \texttt{SmallGainStep.toSyzygy}), and the HOL4
development\textquotesingle s object is the wire checker. The HOL
statement that verifies a composed small-gain certificate is therefore
\rulename{Syzygy}\textquotesingle s (\texttt{syzygyCheck}, previous
section); an independent HOL port of the alias\textquotesingle s unpack
function would be new assurance surface only if HOL-side producers ever
emit \rulename{SmallGain} steps directly.
}
\ruleblock{HOL verifier receipts}
\absentbox{NONE}{no HOL verification receipt exists for this rule.}
\ruleblock{Proof}

\sourcenote{Genre: alias rule (tier 3) --- a respelling template. An alias
owes the book five artefacts: the verified 1$\to$1 unpack function, the
filled-in application of the target rule, the payload-transport table, the
exact algebraic identity that makes the respelling correct, and the
soundness corollary through the target rule's lemma. All five follow.
(The tier criterion counts derivation nodes, not instantiation work: the
unpack \emph{computes} the weighting, but its image is a single
\rulename{Syzygy} node, so the rule is an alias.)}

\begin{brecipe}[the alias template]\label{rec:sg-template}
An alias rule $A$ is warranted by a verified function
$\mathrm{unpack} : \mathrm{Step}_A \to \mathrm{Step}_T$ into the payload
of a \emph{single} target node $T$ (core or admissible), such that the
$A$-node and its image conclude the same sequent from the same premises:
one node $\to$ one node, payload transported, nothing derived. Soundness
is then a corollary, $\mathrm{sound}_A = \mathrm{sound}_T \circ
\mathrm{unpack}$, and the wire format need never carry $A$ at all ---
producers elaborate before emission, which is why \rulename{SmallGain} has
no \texttt{RuleTag}.
\end{brecipe}

\begin{brecipe}[the filled-in application --- \rulename{SmallGain}'s
premises rewritten \emph{into} the \rulename{Syzygy} instance]
\label{rec:sg-filled}
A \rulename{SmallGain} step carries four component proofs: the two
dissipation sequents (over their own boxes $\Phi_1, \Phi_2$) and the two
gain-membership sequents (over the declared operating box $\Phi_w$).
Define the four atom terms and the substitution into
\rulename{Syzygy}'s metavariables:
\[
e_1 \equiv \mathcal{L}_1 V_1 + \alpha_1 - \sigma_1, \qquad
e_2 \equiv \mathcal{L}_2 V_2 + \alpha_2 - \sigma_2, \qquad
m_1 \equiv \lambda\sigma_2 - (1{-}\theta)\,\alpha_1, \qquad
m_2 \equiv \sigma_1 - \lambda(1{-}\theta)\,\alpha_2,
\]
\[
n := 4, \qquad
(e_1, e_2, e_3, e_4) := (e_1, e_2, m_1, m_2), \qquad
(\Phi_1, \Phi_2, \Phi_3, \Phi_4) := (\Phi_1, \Phi_2, \Phi_w, \Phi_w),
\]
\[
(\lambda_1, \lambda_2, \lambda_3, \lambda_4) := (1, \lambda, 1, 1),
\qquad
\{g_k\}_k := \varnothing \ \text{(no equality entries)}, \qquad
c := \mathcal{L}(V_1 + \lambda V_2) + \theta\,(\alpha_1 + \lambda\alpha_2),
\]
with $\lambda \in \Q_{>0}$ and $\theta \in \Q \cap (0,1)$ from the
payload. Under this substitution the \rulename{SmallGain} node \emph{is}
the following \rulename{Syzygy} application (compare the rule display in
the \rulename{Syzygy} section, field for field):
\begin{mathpar}
\inferrule*[Right=\rulename{Syzygy}]{
  \Phi_1 \vdash e_1 \leq 0 \\
  \Phi_2 \vdash e_2 \leq 0 \\
  \Phi_w \vdash m_1 \leq 0 \\
  \Phi_w \vdash m_2 \leq 0 \\\\
  (1, \lambda, 1, 1) \in \Q_{\geq 0}^{\,4} \\
  c \;\equiv\; 1\cdot e_1 + \lambda\cdot e_2 + 1\cdot m_1 + 1\cdot m_2
  \ \text{(exact canonical polynomial)}
}{
  \Phi \vdash c \leq 0
}
\end{mathpar}
with $\Phi \supseteq \Phi_1, \Phi_2, \Phi_w$ conjunct-wise (the
four-way \texttt{preContains} split). Discharging \rulename{Syzygy}'s
obligations at this instance: the nonnegativity checks are $1 \geq 0$
(three times) and $\lambda \geq 0$ (from $\lambda > 0$ --- the
\emph{only} use soundness makes of $\lambda$'s sign); the sourcing check
is immediate (each of the four entries is the corresponding component's
posted atom, verbatim); and the exact-identity check is
Proposition~\ref{prop:sg-identity} below; the five index-freeness
gates (the four posted residuals and the composite conclusion carry no
free indices --- recomputed syntactic checks, the G1 rider that lets
canonical polynomials carry the rule's meaning) discharge by
recomputation on the five terms; the remaining obligations
--- the B13 shape pack, the B14 positivity pack, and the zero-normal
\texttt{postForm} declaration --- transport verbatim, per
Recipe~\ref{rec:sg-transport}. The dead premise slots of
\rulename{Syzygy} are honestly dead: no equality sub-conjuncts of $\Phi$
enter (the $\nu$-row is empty), because the cross-supply cancellation is
performed by the membership atoms themselves. The verified function
computing this rewriting is \texttt{SmallGainStep.toSyzygy}, quoted
above.
\end{brecipe}

\begin{bprop}[the conclusion identity]\label{prop:sg-identity}
As an identity of exact canonical polynomials,
\[
e_1 + \lambda e_2 + m_1 + m_2
\;=\;
\mathcal{L}(V_1 + \lambda V_2) \;+\; \theta\,(\alpha_1 + \lambda\alpha_2).
\]
\end{bprop}
\begin{proof}
Expand and collect --- two lines, which is the point:
\begin{align*}
e_1 + \lambda e_2 + m_1 + m_2
&= \bigl(\mathcal{L}_1 V_1 + \lambda\,\mathcal{L}_2 V_2\bigr)
 + \alpha_1 + \lambda\alpha_2
 \underbrace{-\,\sigma_1 - \lambda\sigma_2
 + \lambda\sigma_2 + \sigma_1}_{=\,0\ \text{(cross-supplies cancel)}}
 - (1{-}\theta)\alpha_1 - \lambda(1{-}\theta)\alpha_2\\
&= \mathcal{L}(V_1 + \lambda V_2)
 + \bigl(1 - (1{-}\theta)\bigr)\alpha_1
 + \lambda\bigl(1 - (1{-}\theta)\bigr)\alpha_2
\;=\; \mathcal{L}(V_1 + \lambda V_2) + \theta\,(\alpha_1 + \lambda\alpha_2).
\end{align*}
The final rewrite reads $\mathcal{L}(V_1 + \lambda V_2)$ for
$\mathcal{L}_1 V_1 + \lambda\,\mathcal{L}_2 V_2$, and that is a
definition, not a step: over disjoint variable sets the composite
generator acts coordinate-wise, so the composite storage-rate term
--- the payload's \texttt{lyapTerm}, pinned to exactly this sum by
the exact canonical ties (\texttt{foldEq} $\cdot$
\texttt{conclPoly}) --- has $\mathcal{L}(V_1 + \lambda V_2)$ as its
published spelling. The cross-supply cancellation is \emph{by
construction of the multipliers}: the membership sequents were chosen so that their
polynomials supply exactly $+\lambda\sigma_2$ and $+\sigma_1$. The checker
replays this as \rulename{Syzygy}'s exact canonical-polynomial identity
$(\ast)$ and \textbf{never sees a fixpoint} --- the classical small-gain
contraction argument is dissolved into two lines of ring arithmetic at
certificate-checking time.
\end{proof}

\begin{brecipe}[payload transport --- \texttt{toSyzygy}, field for field]
\label{rec:sg-transport}
\begin{center}\small
\begin{tabular}{@{}ll@{}}
\textbf{\texttt{SmallGainStep}} & \textbf{\texttt{SyzygyStep}}\\
\midrule
\texttt{seq} & \texttt{seq}\\
\texttt{diss1, diss2, mem1, mem2} & \texttt{components} (that order = combination order)\\
\texttt{combTerm} \textnormal{(weights $(1,\lambda,1,1)$)} & \texttt{combination}\\
\texttt{concl} & \texttt{conclusion}\\
\texttt{pre1, pre2, preM1, preM2} & \texttt{preContains} (four-way case split)\\
\texttt{conclShape, conclFin} & \texttt{conclShape, conclFin} (verbatim)\\
\texttt{wiresPos, compPostWiresPos} & \texttt{wiresPos, compPostWiresPos} (verbatim)\\
\texttt{e1/e2/m1/m2/conclIdxFree} & \textnormal{consumed by} \texttt{combValid} $\mapsto$ \texttt{valid} \textnormal{(G1 gates)}\\
\texttt{combValid} & \texttt{valid}\\
\texttt{postEq} & \texttt{postForm} $=$ \textnormal{(Z$\leq$), the zero-normal reading}\\
\end{tabular}
\end{center}
The alias declares the \emph{zero-normal} reading of the both-polarity
conclusion form: \rulename{SmallGain}'s own contract stays
$\leq$-only, and any polarity flexibility is inherited from
\rulename{Syzygy}'s side of the fence (the reading-transport lemma
\texttt{holds\_of\_syzygyConclusionForm}, proved in the \rulename{Syzygy}
section), not re-justified here.
\end{brecipe}

\begin{bprop}[soundness corollary]\label{prop:sg-sound}
If the four component sequents are semantically valid, the composite
sequent
$\Phi \vdash \mathcal{L}(V_1 + \lambda V_2) + \theta(\alpha_1 +
\lambda\alpha_2) \leq 0$ is semantically valid.
\end{bprop}
\begin{proof}
\texttt{sound\_smallGain} $=$ \texttt{sound\_syzygy} applied to
\texttt{st.toSyzygy}, feeding the four validities through the four-way
membership case split. Beyond that routing there is no further
argument --- the alias genre working as intended.
\end{proof}

\begin{bremark}[$\theta$'s range is well-formedness, not soundness]
The identity of Proposition~\ref{prop:sg-identity} holds at \emph{any}
$\theta$; the memberships self-police: an infeasible $(\lambda, \theta)$
makes the required membership inequalities false over the operating
box, so no sound certificate can discharge them. The $(0,1)$ range check
backs the \emph{published meaning} of $\theta$ as a guard modulus --- a
first-class robustness radius travelling in the proof term, per-node in
the payload, with the certificate header publishing the composite
modulus as the min over nodes --- and $\lambda > 0$ enters soundness
only as $\lambda \geq 0$.
\end{bremark}

\begin{bremark}[relation to the literature]
The rule certificate-izes the Lyapunov small-gain construction: two ISS-style
dissipation inequalities with gain interconnection conditions yield a
composite Lyapunov function $V_1 + \lambda V_2$ --- the nonlinear
small-gain theorem of Jiang--Teel--Praly \cite{jiangteelpraly1994} in its
Lyapunov formulation \cite{jiangmareelswang1996}, sitting inside the
dissipativity-composition tradition \cite{willems1972,moylanhill1978}.
Where the classical theorems close a fixpoint/contraction argument, the
certificate form closes a ring identity: the analytic content is moved
into the \emph{choice} of $(\lambda, \theta)$ and the membership
certificates, and the kernel checks only arithmetic. The contracts are
\emph{signed polynomial} inequalities rather than norm gains for
cause: {\AA}str\"om--Furuta's swing-up analysis \cite{astromfuruta2000}
exploits the sign of $\dot\vartheta\cos\vartheta$, information a
pure norm-gain contract does not record. $(\lambda,\theta)$-existence
results are deliberately outside the checker; the $n$-component
generalization over arbitrary interconnection graphs
\cite{dashkovskiyrufferwirth2007} is catalogued as
\rulename{SmallGainN} --- the same alias discipline, one
\rulename{Syzygy} node, with this rule as its $n = 2$ instance.
\end{bremark}

\begin{anchors}
Recipe \ref{rec:sg-transport} unpack & \texttt{SmallGainStep.toSyzygy} (Composition.lean); \texttt{toSyzygy\_components}\\
Proposition \ref{prop:sg-identity} & \texttt{combValid}'s fold identity (\texttt{foldEq} $\cdot$ \texttt{conclPoly}), replayed as \rulename{Syzygy}'s $(\ast)$\\
Proposition \ref{prop:sg-sound} & \texttt{sound\_smallGain} (a \texttt{sound\_syzygy} corollary)\\
no wire tag, by design & \texttt{RuleTag} comment block (Core/Catalogue.lean): producers elaborate before emission\\
end-to-end instance & \texttt{Handshake/DpCompositeSmallGain.lean} (kernel-checked upstream of the unpack)\\
\end{anchors}


\rulesection{35}{SmallGainN}{alias for \textcolor{DiffBlue}{\rulename{Syzygy}}}
\label{blcatrule-SmallGainN}
\ruleblock{The rule (Steele metanotation)}
\sourcenote{as realized in the mechanized development; source row: \texttt{\texttt{rules.toml}} --- \emph{SmallGainN}. The metanotation statement is typeset by hand for notational quality; its mechanized ground truth is the Lean statement below, extracted from the sources at every build of this book and hash-bound. The rule \emph{roster} is machine-checked (every mechanized rule tag round-trips to exactly one catalogue row); the statement \emph{text} is editorial, checked by eye against the extracted Lean beside it.}
\begin{mathpar}
\inferrule*[Right=\rulename{SmallGainN}]{
  \bigl\{\Phi_i \vdash e_i \leq 0\bigr\}_{i=1}^{n} \\
  \bigl\{\Phi_w \vdash m_i \leq 0\bigr\}_{i=1}^{n} \\
  \Phi \supseteq \Phi_1, \ldots, \Phi_n, \Phi_w \ \text{conjunct-wise (checked)} \\
  \lambda \in \Q_{>0}^{n} \\
  B \in \Q_{\geq 0}^{n \times n} \\
  \textstyle\sum_i B_{ij} = \lambda_j \ (\forall j) \\
  \theta \in \Q \cap (0,1)
}{
  \Phi \vdash \mathcal{L}\bigl(\textstyle\sum_i \lambda_i V_i\bigr) + \theta \textstyle\sum_j \lambda_j \alpha_j \leq 0
}
\end{mathpar}
$e_i \equiv \mathcal{L}_i V_i + \alpha_i - \sigma_i$; \quad $m_i \equiv \lambda_i \sigma_i - (1-\theta) \textstyle\sum_j B_{ij}\, \alpha_j$ (the repayment matrix\textquotesingle s membership rows); \quad conclusion identity, checked as exact canonical polynomial: $\textstyle\sum_i \lambda_i e_i + \sum_i m_i = \mathcal{L}(\sum_i \lambda_i V_i) + \theta \sum_j \lambda_j \alpha_j$ under the column condition.

\ruleflavor{the $n$ membership sequents are ordinary BL sequents over the declared operating box $\Phi_w$, their \rulename{BernCone} multipliers \emph{are} the LP witness; the repayment matrix $B$ routes each component\textquotesingle s supply through its neighbours, and the column condition $\sum_i B_{ij} = \lambda_j$ makes every cross-supply cancel in the conclusion identity --- the checker never sees a fixpoint, a spectral radius, or the interconnection graph. Like the binary rule it is an alias: no proof-term sort, no wire tag; it elaborates to one \rulename{Syzygy} node at weights $(\lambda_1, \ldots, \lambda_n, 1, \ldots, 1)$, and $B$ itself never reaches the wire (the membership rows carry it; the exact identity is the whole obligation). The binary rule is the $n = 2$ instance on the nose: $\lambda = (1, \lambda)$, $B_{21} = 1$, $B_{12} = \lambda$ reproduce \rulename{SmallGain}\textquotesingle s membership rows and its $(1, \lambda, 1, 1)$ weights.}
\par\medskip
\rulename{SmallGainN} is the $n$-component quantitative small-gain
rule: $n$ subsystems on an arbitrary directed interconnection graph,
each dissipating its own storage function except for supplies fed by
its neighbours, compose into a single dissipation certificate for the
weighted sum $\sum_i \lambda_i V_i$ when a weighted loop-gain condition
holds with margin $\theta$. The graph structure enters through the
\emph{repayment matrix} $B$: the membership row
$m_i \equiv \lambda_i \sigma_i - (1{-}\theta)\sum_j B_{ij}\alpha_j$
charges component $i$\textquotesingle s incoming supply against a
$(1{-}\theta)$-discounted draw on its neighbours\textquotesingle{}
decay budgets, and the column condition $\sum_i B_{ij} = \lambda_j$
says every component\textquotesingle s budget is drawn on at exactly
its weight in the composite --- which is what makes the cross-supplies
cancel identically in the conclusion.

Like the binary rule, it is an \emph{alias}: the wire format never
carries a \rulename{SmallGainN} node, and the checker gains no new
obligation kind. A producer that finds $(\lambda, B, \theta)$ --- an LP
feasibility problem over the membership certificates; when the gain
matrix has spectral radius below $1 - \theta$, a Neumann-series weight
vector witnesses feasibility, so the classical spectral small-gain
condition is exactly an existence theory for this rule\textquotesingle s
side data --- elaborates the step to a single \rulename{Syzygy} node at
weights $(\lambda_1, \ldots, \lambda_n, 1, \ldots, 1)$ before emission.
What the checker recomputes is exactly
\rulename{Syzygy}\textquotesingle s obligation set: the exact
conclusion identity over canonical polynomials. Neither the graph, nor
$B$, nor anything fixpoint- or spectrum-shaped survives to checking
time.

What would break without each hypothesis is inherited from
\rulename{Syzygy} through the unpack --- that
section\textquotesingle s commentary applies verbatim to the unpacked
node. The conditions native to this rule: $\lambda_i > 0$, of which
soundness uses only $\lambda_i \geq 0$ (strictness backs the published
meaning: every component genuinely participates); $\theta \in (0,1)$,
well-formedness for $\theta$\textquotesingle s published meaning as a
robustness margin (the identity holds at any $\theta$; the memberships
self-police); and the column condition, which is well-formedness at the
\emph{schema} level --- it is what makes the membership rows deliver
the stated conclusion for every decay budget $\alpha$, while at any
concrete instance the exact identity subsumes it: an emitted node whose
identity checks is sound regardless of how the producer bookkept $B$.

The binary rule is the $n = 2$ instance on the nose: $\lambda = (1,
\lambda)$, $B_{21} = 1$, $B_{12} = \lambda$ reproduce
\rulename{SmallGain}\textquotesingle s membership rows (as a set ---
the two rows arrive in the opposite order, a bookkeeping difference in
combination order) and its $(1, \lambda, 1, 1)$ weights.

\cataloguenote{The $n$-component small-gain theorem for general interconnection
topologies is Dashkovskiy--R\"uffer--Wirth
\cite{dashkovskiyrufferwirth2007}; this rule is its certificate form,
with the cycle/spectral conditions of that tradition dissolved into the
producer\textquotesingle s choice of $(\lambda, B, \theta)$ and the
checker replaying only ring arithmetic. Signed polynomial contracts
rather than norm gains, for the same cause as the binary section.
The dedicated $n$-ary payload vehicle is mechanized as the binary
rule\textquotesingle s pattern at arbitrary $n$: the exact identity
(\texttt{smallGainN\_is\_syzygy} --- the cross-supplies cancel under the
column condition, by canonical-polynomial extensionality), the payload
structure with its verified collapse
(\texttt{SmallGainNStep.toSyzygy}, weights
$(\lambda_1, \ldots, \lambda_n, 1, \ldots, 1)$), the soundness
corollary (\texttt{sound\_smallGainN}), and the column
condition\textquotesingle s necessity as a machine-checked mutant
(\texttt{smallGainN\_colSum\_mutant\_refuted}: delete it and a
one-component instance separates fold from conclusion, by
\texttt{decide}). The semantic face lives beside the binary
one\textquotesingle s in the discrete-kernel module:
\texttt{smallGainN\_drift} with its assembled avoid bound, and the
binary recovery on the nose (\texttt{smallGain\_drift\_of\_N}: $\lambda
= (1, \lambda)$, $B_{21} = 1$, $B_{12} = \lambda$ reproduce the binary
rows --- order-swapped, set-equal, as the commentary states). Every
\emph{emitted} instance remains checked and warranted as
\rulename{Syzygy}\textquotesingle s node: the alias unpacks before
emission.
}
\ruleblock{Lean statement}
\declhead{BernsteinLogic.Composition}{SmallGainNStep}{\badge{BadgeGreen}{statement bound}}
\begin{leansrc}
/-- The n-ary payload vehicle: `SmallGainStep`'s pattern at arbitrary `kk`.
Every field is the checker's obligation for the rule; the collapse target is
one `Syzygy` node (`toSyzygy`). `lambdaPos`, `thetaLo/Hi`, and `colSum` are
rule-level well-formedness exactly as in the binary structure (soundness
uses only `0 ≤ λᵢ` inside the combination check; an emitted node whose
`foldEq` checks is sound however the producer bookkept `B` — the identity
`smallGainN_is_syzygy` is what guarantees the row shapes always discharge
`foldEq` under the column condition). -/
structure SmallGainNStep (H : Hypernet) (kk : ℕ) where
  seq : Sequent
  diss : Fin kk → Sequent
  mem : Fin kk → Sequent
  lambda : Fin kk → ℚ
  B : Fin kk → Fin kk → ℚ
  theta : ℚ
  e : Fin kk → Term
  m : Fin kk → Term
  concl : Term

  dissPost : ∀ i, (diss i).post = .atom .le (e i) (.const 0)
  memPost : ∀ i, (mem i).post = .atom .le (m i) (.const 0)

  eIdxFree : ∀ i, (e i).idxFree = true
  mIdxFree : ∀ i, (m i).idxFree = true
  conclIdxFree : concl.idxFree = true

  lambdaPos : ∀ i, 0 < lambda i
  thetaLo : 0 < theta
  thetaHi : theta < 1
  colSum : ∀ j, ((List.finRange kk).map (fun i => B i j)).sum = lambda j

  foldEq : (((List.finRange kk).map (fun i => ((lambda i : ℚ), e i))
      ++ (List.finRange kk).map (fun i => ((1 : ℚ), m i))).foldr
      (fun pr acc => Poly.add (Poly.scale pr.1 pr.2.toPoly) acc) [])
    = concl.toPoly
  postEq : seq.post = .atom .le concl (.const 0)

  preDiss : ∀ i, premiseContains seq.pre (diss i).pre
  preMem : ∀ i, premiseContains seq.pre (mem i).pre

  conclShape : ∃ s, concl.bshape H = some s
  conclFin : concl.finWires H
  wiresPos : ∀ w ∈ seq.pre.freeWires ++ concl.freeWires,
      ∃ τ, H.wireType w = some τ ∧ 0 < Shape.totalSize τ.shape
  compPostWiresPos : ∀ c ∈ (List.finRange kk).map diss
        ++ (List.finRange kk).map mem,
      ∀ w ∈ c.post.freeWires,
      ∃ τ, H.wireType w = some τ ∧ 0 < Shape.totalSize τ.shape
\end{leansrc}
\declhead{BernsteinLogic.Composition}{smallGainN\_is\_syzygy}{\badge{BadgeGreen}{statement bound}}
\begin{leansrc}
/-- **The n-ary exact identity** (the producer warrant that `foldEq` is
always dischargeable from the row shapes): under the column condition, the
`(λ₁, …, λ_kk, 1, …, 1)`-weighted entry fold over the dissipation and
membership polynomials equals the conclusion polynomial
`∑ᵢ λᵢ·LVᵢ + θ·∑ⱼ λⱼ·Aⱼ` — the cross-supplies cancel identically. Reduces
by canonical-polynomial extensionality to per-key rational arithmetic:
distribute, exchange the double sum, apply the column condition. -/
theorem smallGainN_is_syzygy (kk : ℕ) (LV A Sg : Fin kk → Poly)
    (lam : Fin kk → ℚ) (B : Fin kk → Fin kk → ℚ) (th : ℚ)
    (hcol : ∀ j, ((List.finRange kk).map (fun i => B i j)).sum = lam j) :
    (((List.finRange kk).map
        (fun i => (lam i, Poly.sub (Poly.add (LV i) (A i)) (Sg i)))
      ++ (List.finRange kk).map
        (fun i => ((1 : ℚ), smallGainNMemPoly kk A Sg lam B th i))).foldr
      (fun pr acc => Poly.add (Poly.scale pr.1 pr.2) acc) [])
    = Poly.add
        ((List.finRange kk).foldr
          (fun i acc => Poly.add (Poly.scale (lam i) (LV i)) acc) [])
        (Poly.scale th
          ((List.finRange kk).foldr
            (fun i acc => Poly.add (Poly.scale (lam i) (A i)) acc) []))
\end{leansrc}
\declhead{BernsteinLogic.Composition}{sound\_smallGainN}{\badge{BadgeGreen}{statement bound}}
\begin{leansrc}
/-- **Soundness of SmallGainN**, as the Syzygy corollary: the `2·kk`
component validities feed `sound_syzygy` on the collapsed instance. The
rational comparisons (`λᵢ > 0`, `0 < θ < 1`) enter only through `λᵢ ≥ 0`
inside the combination check; `θ`'s range and the column condition are
well-formedness, exactly as the catalogue commentary states. -/
theorem sound_smallGainN {H : Hypernet} {kk : ℕ} (st : SmallGainNStep H kk)
    (hdiss : ∀ i, (st.diss i).semanticallyValid H)
    (hmem : ∀ i, (st.mem i).semanticallyValid H) :
    st.seq.semanticallyValid H
\end{leansrc}
\declhead{BernsteinLogic.Composition}{smallGainN\_colSum\_mutant\_refuted}{\badge{BadgeGreen}{statement bound}}
\begin{leansrc}
/-- **M-N2 (the column condition deleted): the identity is REFUTED** — by
the B-swap, computation-free: the mutant at `B = 0` and the PROVED identity
at `B = 1` (column sums check: `1 = λ`) share every other datum
(`kk = 1`, `LV = Sg = 0`, `A = const 1`, `λ = 1`, `θ = 0`), so their entry
folds would be equal; reading coefficients through the fold lemmas, the two
membership polynomials differ by exactly the decay budget's coefficient —
forcing `coeffSum (const 1) = 0` everywhere, against canonicity. The column
condition is what makes the membership rows deliver the stated conclusion
for every decay budget. -/
theorem smallGainN_colSum_mutant_refuted :
    ¬ (∀ (kk : ℕ) (LV A Sg : Fin kk → Poly) (lam : Fin kk → ℚ)
        (B : Fin kk → Fin kk → ℚ) (th : ℚ),
        (((List.finRange kk).map
            (fun i => (lam i, Poly.sub (Poly.add (LV i) (A i)) (Sg i)))
          ++ (List.finRange kk).map
            (fun i => ((1 : ℚ), smallGainNMemPoly kk A Sg lam B th i))).foldr
          (fun pr acc => Poly.add (Poly.scale pr.1 pr.2) acc) [])
        = Poly.add
            ((List.finRange kk).foldr
              (fun i acc => Poly.add (Poly.scale (lam i) (LV i)) acc) [])
            (Poly.scale th
              ((List.finRange kk).foldr
                (fun i acc => Poly.add (Poly.scale (lam i) (A i)) acc) [])))
\end{leansrc}
\declhead{BernsteinLogicAux.QtslCompose.AGComposeDiscrete}{smallGainN\_drift}{\badge{BadgeGreen}{statement bound}}
\begin{leansrc}
/-- **B3′, the cancellation row at any n** (the semantic face of the
checker-level `(λ₁, …, λₙ, 1, …, 1)` syzygy fold): per-component dissipation
rows with cross-supplies on the joint closed-loop kernel, under the n-ary
membership rows and the exact column condition, yield the `θ`-margin drift
row for `W = ∑ i, lam i * V i`. The cross-supplies cancel by construction:
`∑ i, lam i * sigma i ≤ (1-θ) ∑ i ∑ j, B i j * alpha j = (1-θ) ∑ j, lam j * alpha j`
(swap the double sum, then the column condition). -/
theorem smallGainN_drift (hlam : ∀ i, 0 ≤ lam i)
    (hcol : ∀ j, ∑ i, B i j = lam j)
    (hint : ∀ i, ∀ y, y ∉ A → Integrable (V i) (K y))
    (hd : ∀ i, ∀ y, y ∉ A →
        ∫ z, V i z ∂K y ≤ V i y - alpha i y + sigma i y)
    (hm : ∀ i, ∀ y, y ∉ A →
        lam i * sigma i y ≤ (1 - theta) * ∑ j, B i j * alpha j y) :
    ∀ y, y ∉ A → ∫ z, (∑ i, lam i * V i z) ∂K y
      ≤ (∑ i, lam i * V i y) - theta * ∑ j, lam j * alpha j y
\end{leansrc}
\declhead{BernsteinLogicAux.QtslCompose.AGComposeDiscrete}{smallGain\_drift\_of\_N}{\badge{BadgeGreen}{statement bound}}
\begin{leansrc}
/-- **Binary recovery, on the nose**: `smallGain_drift`'s exact statement
re-derived from `smallGainN_drift` at `m = 2`, `lam = ![1, λ]`,
`B = ![![0, λ], ![1, 0]]` (column sums `0 + 1 = 1` and `λ + 0 = λ`). One
bookkeeping disanalogy, inherited from the catalogue commentary: the n-ary
membership row at `i = 0` reads `1 * σ₁ ≤ (1-θ)·λ·α₂` — the binary `hm₂`
with its factors reassociated — and at `i = 1` reads `λ * σ₂ ≤ (1-θ)·α₁` =
the binary `hm₁`; the two rows arrive order-swapped vs the binary labels,
set-equal. -/
theorem smallGain_drift_of_N {V₁ V₂ alpha₁ alpha₂ sigma₁ sigma₂ : S → ℝ}
    {lam theta : ℝ} (hlam : 0 ≤ lam)
    (hint₁ : ∀ y, y ∉ A → Integrable V₁ (K y))
    (hint₂ : ∀ y, y ∉ A → Integrable V₂ (K y))
    (hd₁ : ∀ y, y ∉ A → ∫ z, V₁ z ∂K y ≤ V₁ y - alpha₁ y + sigma₁ y)
    (hd₂ : ∀ y, y ∉ A → ∫ z, V₂ z ∂K y ≤ V₂ y - alpha₂ y + sigma₂ y)
    (hm₁ : ∀ y, y ∉ A → lam * sigma₂ y ≤ (1 - theta) * alpha₁ y)
    (hm₂ : ∀ y, y ∉ A → sigma₁ y ≤ lam * ((1 - theta) * alpha₂ y)) :
    ∀ y, y ∉ A → ∫ z, (V₁ z + lam * V₂ z) ∂K y
      ≤ (V₁ y + lam * V₂ y) - theta * (alpha₁ y + lam * alpha₂ y)
\end{leansrc}
\declhead{BernsteinLogicAux.QtslCompose.AGComposeDiscrete}{avoid\_smallGainN}{\badge{BadgeGreen}{statement bound}}
\begin{leansrc}
/-- **B3′, assembled**: on the closed loop, the n-ary small-gain data yields
the composite avoid bound for the weighted-sum certificate (the discrete face
of the catalogue rule end-to-end; Foster–Lyapunov metatheory through
`Battery.measure_reach_le`, exactly as the binary `avoid_smallGain`). -/
theorem avoid_smallGainN (hlam : ∀ i, 0 ≤ lam i) (htheta : 0 ≤ theta)
    {C : Fin m → ℝ} (hA : MeasurableSet A)
    (hcol : ∀ j, ∑ i, B i j = lam j)
    (hV : ∀ i, Measurable (V i))
    (hV0 : ∀ i y, 0 ≤ V i y) (hVC : ∀ i y, V i y ≤ C i)
    (hbar : ∀ y ∈ A, 1 ≤ ∑ i, lam i * V i y)
    (halpha : ∀ i, ∀ y, y ∉ A → 0 ≤ alpha i y)
    (hd : ∀ i, ∀ y, y ∉ A →
        ∫ z, V i z ∂K y ≤ V i y - alpha i y + sigma i y)
    (hm : ∀ i, ∀ y, y ∉ A →
        lam i * sigma i y ≤ (1 - theta) * ∑ j, B i j * alpha j y) (x : S) :
    pathMeasure K x {ω | ∃ n, ω n ∈ A}
      ≤ ENNReal.ofReal (∑ i, lam i * V i x)
\end{leansrc}
\ruleblock{Correspondence (Steele display $\leftrightarrow$ Lean)}
\absentbox{NO CORRESPONDENCE DICTIONARY YET}{no reviewer-certified dictionary maps the displayed rule onto the mechanized statements above; until one is written and certified, that mapping rests on this section\textquotesingle s editorial care alone.}
\ruleblock{Lean verifier receipts}
\sourcenote{module \texttt{BernsteinLogic.Composition}; bound build \texttt{9116f454beb0\ldots} (bound from receipt 548ea3dda8ed73a2, 2026-08-03T21:30:29Z). A receipt records that named checking layers accepted a build: (module, compiled-artifact SHA-256, environment SHA-256, checker layers), minted when the checking gate ran; the store receipts attesting the bound build are listed with the layers each run credited. Nothing here asks the reader to trust the store: every layer is an independently runnable program over the quoted sources --- rebuild the module and re-run the checkers to reproduce the verdict. Receipts attesting superseded builds credit no layer and are omitted from this rendering (the count of omissions, if any, is noted below).}
\begin{receipttable}
\receiptrow{548ea3dda8ed73a2}{2026-08-03T21:30:29Z}{build+decorrelated}{\badge{BadgeGreen}{current}}
\end{receipttable}
\sourcenote{10 further receipt(s) in the store attest superseded builds of this module and are omitted from this rendering.}
\sourcenote{module \texttt{BernsteinLogicAux.QtslCompose.AGComposeDiscrete}; bound build \texttt{d5233632f95a\ldots} (bound from receipt 88ff74180adf5b95, 2026-08-03T20:13:53Z). A receipt records that named checking layers accepted a build: (module, compiled-artifact SHA-256, environment SHA-256, checker layers), minted when the checking gate ran; the store receipts attesting the bound build are listed with the layers each run credited. Nothing here asks the reader to trust the store: every layer is an independently runnable program over the quoted sources --- rebuild the module and re-run the checkers to reproduce the verdict. Receipts attesting superseded builds credit no layer and are omitted from this rendering (the count of omissions, if any, is noted below).}
\begin{receipttable}
\receiptrow{88ff74180adf5b95}{2026-08-03T20:13:53Z}{build+decorrelated}{\badge{BadgeGreen}{current}}
\end{receipttable}
\sourcenote{11 further receipt(s) in the store attest superseded builds of this module and are omitted from this rendering.}
\layerline{\layeritem{build}{lake build (elaboration)}{\badge{BadgeGreen}{current}}\layeritem{replay}{leanchecker (official kernel replay)}{\badge{BadgeAmber}{stale only}}\layeritem{hardened}{safe\_verify (hardened screen)}{\badge{BadgeAmber}{stale only}}\layeritem{decorrelated}{lean4export + nanoda (decorrelated kernel)}{\badge{BadgeGreen}{current}}}
\sourcenote{\badge{BadgeGray}{comparator: structurally n/a} the alias unpacks
\emph{before} emission (no rule tag; producers elaborate a
\rulename{SmallGainN} step to its \rulename{Syzygy} node), so no wire
certificate --- and hence no comparator fixture --- can ever contain a
\rulename{SmallGainN} node. Comparator coverage for this rule \emph{is}
\rulename{Syzygy}\textquotesingle s.
}
\ruleblock{HOL4 statement}
\absentbox{OUT OF SCOPE BY DESIGN}{Out of scope for the HOL side \emph{by design}, not by backlog, and for
exactly the binary rule\textquotesingle s reason: the wire format never
carries a \rulename{SmallGainN} node --- producers elaborate to the
$(\lambda_1, \ldots, \lambda_n, 1, \ldots, 1)$ \rulename{Syzygy} node
before emission --- and the HOL4 development\textquotesingle s object
is the wire checker. The HOL statement that verifies a composed
$n$-component small-gain certificate is therefore
\rulename{Syzygy}\textquotesingle s (\texttt{syzygyCheck}).
}
\ruleblock{HOL verifier receipts}
\absentbox{NONE}{no HOL verification receipt exists for this rule.}
\ruleblock{Proof}

\sourcenote{Genre: alias rule (tier 3) --- a respelling template. The
five artefacts an alias owes the book: the 1$\to$1 unpack, the
filled-in application of the target rule, the payload-transport table,
the exact algebraic identity that makes the respelling correct, and the
soundness corollary through the target rule's lemma. All five follow;
the unpack and transport table are stated here as the producer's
contract, their dedicated mechanized vehicle being the one artefact not
yet in the Lean tree (stated loudly below and in the section's Lean
cell) --- soundness of every emitted node meanwhile rests on the target
rule's already-mechanized lemma, because the alias unpacks before
emission.}

\begin{brecipe}[the unpack --- one node, any $n$]\label{rec:sgn-unpack}
A \rulename{SmallGainN} step carries $2n$ component proofs: the $n$
dissipation sequents (each over its own box $\Phi_i$) and the $n$
membership sequents (over the declared operating box $\Phi_w$), with
side data $\lambda \in \Q_{>0}^{n}$, $B \in \Q_{\geq 0}^{n \times n}$
with $\sum_i B_{ij} = \lambda_j$ for every $j$, and $\theta \in \Q \cap
(0,1)$. Define the atom terms
\[
e_i \equiv \mathcal{L}_i V_i + \alpha_i - \sigma_i,
\qquad
m_i \equiv \lambda_i \sigma_i - (1{-}\theta) \sum_j B_{ij}\, \alpha_j ,
\]
and the substitution into \rulename{Syzygy}'s metavariables:
\[
N := 2n, \qquad
(e_1, \ldots, e_N) := (e_1, \ldots, e_n, m_1, \ldots, m_n), \qquad
(\Phi_1, \ldots, \Phi_N) := (\Phi_1, \ldots, \Phi_n,
  \underbrace{\Phi_w, \ldots, \Phi_w}_{n}),
\]
\[
(\lambda_1, \ldots, \lambda_N) := (\lambda_1, \ldots, \lambda_n,
  1, \ldots, 1), \qquad
\{g_k\}_k := \varnothing, \qquad
c := \mathcal{L}\bigl(\textstyle\sum_i \lambda_i V_i\bigr)
  + \theta \textstyle\sum_j \lambda_j \alpha_j .
\]
The repayment matrix $B$ is consumed here and does not appear in the
image: it is the producer's bookkeeping for \emph{finding} membership
rows that make the identity of
Proposition~\ref{prop:sgn-identity} true, and the identity is the only
trace of it the checker ever replays.
\end{brecipe}

\begin{brecipe}[the filled-in application]\label{rec:sgn-filled}
Under the substitution of Recipe~\ref{rec:sgn-unpack} the
\rulename{SmallGainN} node \emph{is} the following \rulename{Syzygy}
application:
\begin{mathpar}
\inferrule*[Right=\rulename{Syzygy}]{
  \bigl\{\Phi_i \vdash e_i \leq 0\bigr\}_{i=1}^{n} \\
  \bigl\{\Phi_w \vdash m_i \leq 0\bigr\}_{i=1}^{n} \\\\
  (\lambda_1, \ldots, \lambda_n, 1, \ldots, 1) \in \Q_{\geq 0}^{\,2n} \\
  c \;\equiv\; \textstyle\sum_i \lambda_i\, e_i + \sum_i 1 \cdot m_i
  \ \text{(exact canonical polynomial)}
}{
  \Phi \vdash c \leq 0
}
\end{mathpar}
with $\Phi \supseteq \Phi_1, \ldots, \Phi_n, \Phi_w$ conjunct-wise (the
$2n$-way containment split). Discharging \rulename{Syzygy}'s
obligations at this instance: the nonnegativity checks are $\lambda_i
\geq 0$ (from $\lambda_i > 0$ --- the only use soundness makes of the
signs) and $1 \geq 0$, $n$ times each; the sourcing check is immediate
(each of the $2n$ entries is the corresponding component's posted atom,
verbatim); the exact-identity check is
Proposition~\ref{prop:sgn-identity}; the index-freeness gates discharge
by recomputation on the $2n + 1$ terms; and the shape, positivity, and
zero-normal \texttt{postForm} obligations transport exactly as in the
binary section. The $\nu$-row is empty for the same reason as there:
the cross-supply cancellation is performed by the membership atoms
themselves, so no equality sub-conjuncts of $\Phi$ enter.
\end{brecipe}

\begin{bprop}[the conclusion identity]\label{prop:sgn-identity}
Under the column condition $\sum_i B_{ij} = \lambda_j$ (all $j$), as an
identity of exact canonical polynomials,
\[
\sum_i \lambda_i e_i + \sum_i m_i
\;=\;
\mathcal{L}\bigl(\textstyle\sum_i \lambda_i V_i\bigr)
\;+\; \theta \sum_j \lambda_j \alpha_j .
\]
\end{bprop}
\begin{proof}
Expand and collect --- still two lines at any $n$, which is the point:
\begin{align*}
\sum_i \lambda_i e_i + \sum_i m_i
&= \sum_i \lambda_i \mathcal{L}_i V_i
 + \sum_i \lambda_i \alpha_i
 \underbrace{- \sum_i \lambda_i \sigma_i
 + \sum_i \lambda_i \sigma_i}_{=\,0\ \text{(cross-supplies cancel)}}
 - (1{-}\theta) \sum_j \Bigl(\sum_i B_{ij}\Bigr) \alpha_j \\
&= \mathcal{L}\bigl(\textstyle\sum_i \lambda_i V_i\bigr)
 + \sum_j \bigl(1 - (1{-}\theta)\bigr) \lambda_j \alpha_j
\;=\; \mathcal{L}\bigl(\textstyle\sum_i \lambda_i V_i\bigr)
 + \theta \sum_j \lambda_j \alpha_j ,
\end{align*}
where the second line swaps the summation order in the $B$-term and
applies the column condition $\sum_i B_{ij} = \lambda_j$ --- the one
place it is used, and exactly the discount that makes the neighbour
draws rejoin the decay budgets at weight $\lambda_j$. As in the binary
section, reading $\mathcal{L}(\sum_i \lambda_i V_i)$ for $\sum_i
\lambda_i \mathcal{L}_i V_i$ is a definition, not a step: over disjoint
variable sets the composite generator acts coordinate-wise. The checker
replays this as \rulename{Syzygy}'s exact canonical-polynomial identity
and never sees a fixpoint, a spectral radius, or the interconnection
graph --- the $n$-component small-gain argument is dissolved into ring
arithmetic at certificate-checking time.
\end{proof}

\begin{bprop}[binary recovery, on the nose]\label{prop:sgn-binary}
At $n = 2$ with $\lambda = (1, \lambda)$, $B_{21} = 1$, $B_{12} =
\lambda$ (and $B_{11} = B_{22} = 0$), the column condition holds, the
membership rows are exactly the binary section's
$\{\,\lambda\sigma_2 - (1{-}\theta)\alpha_1,\;
\sigma_1 - \lambda(1{-}\theta)\alpha_2\,\}$, the conclusion is exactly
$\mathcal{L}(V_1 + \lambda V_2) + \theta(\alpha_1 + \lambda\alpha_2)
\leq 0$, and the \rulename{Syzygy} weights are $(1, \lambda, 1, 1)$ ---
i.e.\ \rulename{SmallGain} is the $n = 2$ instance of this rule.
\end{bprop}
\begin{proof}
Column sums: $B_{11} + B_{21} = 1 = \lambda_1$ and $B_{12} + B_{22} =
\lambda = \lambda_2$. Membership rows: $m_1 = \sigma_1 -
\lambda(1{-}\theta)\alpha_2$ and $m_2 = \lambda\sigma_2 -
(1{-}\theta)\alpha_1$ --- the binary section's rows as a \emph{set},
arriving in the opposite order (its $m_1$ is this rule's $m_2$): a
bookkeeping difference in combination order, visible in the node bytes
and irrelevant to validity, since \rulename{Syzygy}'s identity is a
sum. Conclusion and weights read off directly.
\end{proof}

\begin{brecipe}[payload transport --- the producer's contract]
\label{rec:sgn-transport}
The $n$-ary payload transports into \texttt{SyzygyStep} exactly as the
binary payload does, with the four fixed slots generalized to $n$-vectors:
\begin{center}\small
\begin{tabular}{@{}ll@{}}
\textbf{\rulename{SmallGainN} step (contract)} & \textbf{\texttt{SyzygyStep}}\\
\midrule
\texttt{seq} & \texttt{seq}\\
$\texttt{diss}_1, \ldots, \texttt{diss}_n, \texttt{mem}_1, \ldots, \texttt{mem}_n$ & \texttt{components} (that order = combination order)\\
weights $(\lambda_1, \ldots, \lambda_n, 1, \ldots, 1)$ & \texttt{combination}\\
\texttt{concl} & \texttt{conclusion}\\
$2n$-way containment split & \texttt{preContains}\\
shape/positivity/idx-freeness packs & verbatim, as in the binary table\\
$(\lambda, B, \theta)$ & consumed by the unpack; $B$ leaves no image field\\
\texttt{postEq} & \texttt{postForm} $=$ the zero-normal reading\\
\end{tabular}
\end{center}
This table is the contract of the dedicated mechanized vehicle --- the
$n$-ary payload structure with its verified unpack and soundness
corollary, the binary section's \texttt{SmallGainStep} apparatus at
arbitrary $n$ --- which is not yet in the Lean tree (the section's Lean
cell states the absence). Nothing checked on the wire waits on it: the
unpack happens producer-side, and the emitted node is
\rulename{Syzygy}'s.
\end{brecipe}

\begin{bprop}[soundness corollary]\label{prop:sgn-sound}
If the $2n$ component sequents are semantically valid, the composite
sequent $\Phi \vdash \mathcal{L}(\sum_i \lambda_i V_i) + \theta \sum_j
\lambda_j \alpha_j \leq 0$ is semantically valid.
\end{bprop}
\begin{proof}
The target rule's kernel lemma applied to the unpacked instance of
Recipe~\ref{rec:sgn-filled}, feeding the $2n$ validities through the
containment split; the exact-identity obligation is
Proposition~\ref{prop:sgn-identity}. Beyond that routing there is no
further argument --- the alias genre working as intended, at every $n$
at once.
\end{proof}

\begin{bremark}[what stays outside the checker]
$(\lambda, B, \theta)$-\emph{existence} is the producer's problem, and
the classical spectral conditions are its solution theory: when the
gain matrix $\Gamma$ has spectral radius $\rho(\Gamma) < 1 - \theta$,
the Neumann-series vector $\lambda := \sum_k (1-\theta)^{-k}
(\Gamma^{\mathsf{T}})^k \mathbb{1}$ converges and witnesses
$\Gamma^{\mathsf{T}} \lambda \leq (1-\theta) \lambda$, from which a
feasible $B$ is assembled column by column (under-allocation is
repaired by diagonal padding, which only weakens the membership row it
feeds). The $n$-component ISS small-gain theorem for general networks
is Dashkovskiy--R\"uffer--Wirth \cite{dashkovskiyrufferwirth2007};
where that tradition closes a monotone-operator fixpoint argument over
the interconnection graph, the certificate form closes the ring
identity of Proposition~\ref{prop:sgn-identity}, and the graph
combinatorics survive only as the producer's search space. The signed
polynomial contract discipline is the binary section's, unchanged.
\end{bremark}

\begin{anchors}
Recipe \ref{rec:sgn-unpack} unpack & producer-side elaboration; dedicated mechanized vehicle = named successor (binary vehicle: \texttt{SmallGainStep.toSyzygy})\\
Proposition \ref{prop:sgn-identity} & \texttt{combValid}'s fold identity, replayed as \rulename{Syzygy}'s $(\ast)$ on the emitted node\\
Proposition \ref{prop:sgn-binary} & the binary section's \texttt{SmallGainStep} apparatus at $n = 2$ (weights $(1, \lambda, 1, 1)$)\\
Proposition \ref{prop:sgn-sound} & \texttt{sound\_syzygy} on the unpacked instance ($n$-ary corollary = named successor)\\
no wire tag, by design & \texttt{RuleTag} comment block (Core/Catalogue.lean): producers elaborate before emission\\
\end{anchors}


\rulesection{36}{ProductLattice}{admissible rule}
\label{blcatrule-ProductLattice}
\ruleblock{The rule (Steele metanotation)}
\sourcenote{as realized in the mechanized development; source row: \texttt{\texttt{rules.toml}} --- \emph{ProductLattice}. The metanotation statement is typeset by hand for notational quality; its mechanized ground truth is the Lean statement below, extracted from the sources at every build of this book and hash-bound. The rule \emph{roster} is machine-checked (every mechanized rule tag round-trips to exactly one catalogue row); the statement \emph{text} is editorial, checked by eye against the extracted Lean beside it.}
\begin{mathpar}
\inferrule*[Right=\rulename{ProductLattice}]{
  B_1 \ \text{a certified attracting-block lattice for } \varphi_1 \ \text{over } H_1
  \\
  B_2 \ \text{likewise for } \varphi_2 \ \text{over } H_2
  \\
  \mathrm{wires}(H_1) \cap \mathrm{wires}(H_2) = \emptyset
}{
  \{\, b_1 \times b_2 \mid b_1 \in B_1,\ b_2 \in B_2 \,\} \ \text{(the meet-generators) certified for } \varphi_1 \times \varphi_2 \ \text{over } H_1 \sqcup H_2
}
\end{mathpar}
the displayed family is the list of pairwise-product meet-generators.

\ruleflavor{the full lattice $\langle B_1 \times B_2 \rangle_{\vee}$ (join-closure of the pairwise products) is the paper-strength reading --- the bound theorem quantifies over the meet-generators, and the finite-union/boundary-decomposition step remains paper-strength, exactly as the correspondence concedes. Every composite facet sequent is discharged BY CITATION: a composite facet is normal to one factor\textquotesingle s axis, its inwardness polynomial IS that factor\textquotesingle s facet polynomial with the other factor\textquotesingle s coordinates riding along as free box variables (the rider, conjoined onto the premise). Bernstein-wise the composite facet grid is the factor grid tensored with all-ones --- a coefficient identity; nothing for a checker to recompute beyond the ordinary \rulename{BernCone} recompute pattern. Declared in the baked base theory itself (its first certificate-family rule), not in any separate composition layer.}
\par\medskip
\rulename{ProductLattice} composes at the certificate-\emph{family}
level, one judgment up from the sequent rules: given certified
attracting-block lattices for two systems over wire-disjoint hypernets,
it certifies the join-closure of pairwise products as an
attracting-block lattice for the product system over the disjoint
union. The geometric fact that powers it: every internal facet of an
axis-aligned product block is normal to one factor\textquotesingle s
axis, so its inwardness polynomial \emph{is} that
factor\textquotesingle s facet polynomial, with the other
factor\textquotesingle s coordinates riding along as free box
variables. Each composite facet sequent is therefore discharged
\emph{by citation}: conclusion kept verbatim, the other
factor\textquotesingle s box conjoined onto the premise as a rider.

What the checker recomputes: no new Bernstein arithmetic ---
Bernstein-wise the composite facet grid is the factor grid tensored
with all-ones, a coefficient identity, so the per-facet discharge is
syntactic identity of the cited conclusion. What remains checked is
the citation plumbing itself: wire-disjointness of the two hypernets,
factor scopedness, and the rider conjunction on each cited premise.

Wire-disjointness is load-bearing twice. Structurally, the disjoint
union of the two hypernets is only well-formed without shared wires
(linearity and acyclicity of the union are proven \emph{from}
disjointness). Semantically, no single transported facet sequent\textquotesingle s
validity fails under a shared wire --- the weaken lemma conjoins any
rider --- but the assembled object stops being a product of the
factors: the identification of every composite facet with a cited
factor facet, the other factor entering only as a rider over fresh
coordinates, is proven from disjointness, and with a shared wire the
recomputed readout misdescribes the composite. A single shared wire is
precisely an interface, and composition along an interface is
\rulename{Promise}/small-gain territory, not this
rule\textquotesingle s.

\cataloguenote{The soundness statement is dynamics-free --- sound for flows and maps
alike, because it certifies a \emph{lower bound} family (\textquotedblleft
these blocks are certified\textquotedblright), never \textquotedblleft
these are all the attractors\textquotedblright. The completeness
reading --- that the generated lattice is the whole attractor lattice
of the product --- is a flows-only theorem of the source literature
and lives paper-side; the proof section exhibits the map-dynamics
counterexample that forces the fence. No wire tag is allocated yet:
soundness lands first, executable checking follows the family-object
wire-format work, and the emit path refuses the rule until then.
}
\ruleblock{Lean statement}
\par\medskip
\textbf{What \textquotedblleft discharge by
citation\textquotedblright{} means, stated in the flavor text and
bound by the declarations below.}\enspace \texttt{citeFacet} keeps the cited
factor facet\textquotesingle s conclusion verbatim and conjoins the
rider onto the premise; \texttt{prodBlock} assembles a product
block\textquotesingle s facet list as exactly the two citation
families; \texttt{prodBlocks} enumerates only the meet-generators of
the join-closure --- the exponentially many join blocks are never
materialized; their sub-facets discharge through the paper-strength
boundary decomposition (which selects a factor facet and rider), whose
final transport-and-weaken step is the same citation lemmas with
smaller riders. A skeptical reader should confirm by eye
that no definition below evaluates dynamics anywhere: the rule\textquotesingle s
semantic content is entirely the transport-and-weaken pair quoted
above.

\declhead{BernsteinLogic.ProductLattice}{BlockCert}{\badge{BadgeGreen}{statement bound}}
\begin{leansrc}
/-- One certified attracting block, certificate view (the recon's Route-2
shape): its region — the compact axis-aligned box(-union), as the formula
the producer wrote it with — and its internal-facet inwardness sequents
(strict inwardness via non-strict predicates with explicit positive
rational margins; block-hood is elementary for strictly inward box unions,
no combinatorial index data claimed — note §1.1). -/
structure BlockCert where
  region : Formula
  facets : List Sequent
\end{leansrc}
\declhead{BernsteinLogic.ProductLattice}{CertifiedBlockLattice}{\badge{BadgeGreen}{statement bound}}
\begin{leansrc}
/-- A certified attracting-block lattice over a hypernet — the catalogue's
first certificate-FAMILY object (premise object of ProductLattice).
The family carries its lattice structure explicitly at the index level:
`join`/`meet` with the six standard lattice laws (idempotence is
derivable). These are the certificate-level shadow of the note §1.1
closure ("closed under ∪ and regularized ∩"); that a join index denotes
the set-union of the denoted blocks is producer-side geometry in v1 —
kernel teeth for the denotation arrive with the family-object wire format
(C-4/format cycle). Wire confinement makes the family transportable into
any extension (§3) — it is what `sound_productLattice` consumes. -/
structure CertifiedBlockLattice (H : Hypernet) where
  blocks : List BlockCert
  join : Fin blocks.length → Fin blocks.length → Fin blocks.length
  meet : Fin blocks.length → Fin blocks.length → Fin blocks.length
  join_comm : ∀ i j, join i j = join j i
  join_assoc : ∀ i j k, join (join i j) k = join i (join j k)
  meet_comm : ∀ i j, meet i j = meet j i
  meet_assoc : ∀ i j k, meet (meet i j) k = meet i (meet j k)
  absorb_join : ∀ i j, join i (meet i j) = i
  absorb_meet : ∀ i j, meet i (join i j) = i
  regionConfined : ∀ b ∈ blocks, ∀ w ∈ b.region.freeWires, (H.wireType w).isSome
  facetsScoped : ∀ b ∈ blocks, ∀ s ∈ b.facets, s.wellScoped H
\end{leansrc}
\declhead{BernsteinLogic.ProductLattice}{PProductLatticeData}{\badge{BadgeGreen}{statement bound}}
\begin{leansrc}
/-- Applicability data of the ProductLattice rule — the two factor
families and the side condition. Field-for-field what a checker would
verify: the factor certificates (whose facet-sequent validity enters
`sound_productLattice` as hypotheses, Step-style) and wire-disjointness.
No first-class `RuleTag`/wire constructor in v1 (note §1.4: emission-layer
posture; BINARY.md §6.2 reserves tag 27 for the family-object format
cycle); correspondingly NOT in `v0frag` — like the ReLU family
pre-enactment, soundness lands first and executable checking follows the
format work. -/
structure PProductLatticeData (H1 H2 : Hypernet) where
  disjoint : H1.WireDisjoint H2
  left : CertifiedBlockLattice H1
  right : CertifiedBlockLattice H2
\end{leansrc}
\declhead{BernsteinLogic.ProductLattice}{Sequent.semanticallyValid\_extend}{\badge{BadgeGreen}{statement bound}}
\begin{leansrc}
/-- TRANSPORT: semantic validity is monotone under hypernet extension —
a sequent valid over H, well-scoped in H, is valid over any H' that
extends H's wire typing and contains H's applications. (For ProductLattice
this is "the factor's certified sign is the composite's inwardness": the
facet polynomial reads only the factor's wires, so the other factor's
coordinates ride along invisibly.) -/
theorem Sequent.semanticallyValid_extend {H H' : Hypernet} {seq : Sequent}
    (hext : H.WireExtends H')
    (happs : ∀ app ∈ H.applications, app ∈ H'.applications)
    (hws : seq.wellScoped H) (hval : seq.semanticallyValid H) :
    seq.semanticallyValid H'
\end{leansrc}
\declhead{BernsteinLogic.ProductLattice}{Sequent.semanticallyValid\_weaken}{\badge{BadgeGreen}{statement bound}}
\begin{leansrc}
/-- WEAKEN, semantic and standalone: any conjunct may join a valid
sequent's premise (PWeaken's semantic content — the rider needs no
side condition). -/
theorem Sequent.semanticallyValid_weaken {H : Hypernet}
    {pre extra post : Formula}
    (h : Sequent.semanticallyValid H ⟨pre, post⟩) :
    Sequent.semanticallyValid H ⟨.and pre extra, post⟩
\end{leansrc}
\declhead{BernsteinLogic.ProductLattice}{sound\_productLattice\_citeLeft}{\badge{BadgeGreen}{statement bound}}
\begin{leansrc}
/-- The standing citation discharge, left factor: ANY sequent citing a
valid, well-scoped left-factor facet — conclusion verbatim, premise
= factor premise ∧ an arbitrary rider — is valid over the union.
Soundness is transport (§3) then weaken: the note §1.2's one-line
argument, in two lines. A join block's sub-facet, once the paper-strength
boundary decomposition has selected its factor facet and rider, is
discharged by this lemma (rider = the sub-facet's free-coordinate cell);
the decomposition itself is not mechanized here. -/
theorem sound_productLattice_citeLeft {H1 H2 : Hypernet}
    (hd : H1.WireDisjoint H2) {s0 : Sequent} (rider : Formula)
    (hws : s0.wellScoped H1) (hval : s0.semanticallyValid H1) :
    (citeFacet s0 rider).semanticallyValid (H1.disjUnion H2 hd)
\end{leansrc}
\declhead{BernsteinLogic.ProductLattice}{sound\_productLattice\_citeRight}{\badge{BadgeGreen}{statement bound}}
\begin{leansrc}
/-- The standing citation discharge, right factor. -/
theorem sound_productLattice_citeRight {H1 H2 : Hypernet}
    (hd : H1.WireDisjoint H2) {s0 : Sequent} (rider : Formula)
    (hws : s0.wellScoped H2) (hval : s0.semanticallyValid H2) :
    (citeFacet s0 rider).semanticallyValid (H1.disjUnion H2 hd)
\end{leansrc}
\declhead{BernsteinLogic.ProductLattice}{sound\_productLattice}{\badge{BadgeGreen}{statement bound}}
\begin{leansrc}
/-- SOUNDNESS of ProductLattice (note §1.2, written properly): given two
certified block lattices over wire-disjoint hypernets, every facet
sequent of every product block in the conclusion family is semantically
valid over the disjoint union. Each facet is a factor facet discharged by
citation; validity = transport of the factor's validity along the union
injection (§3), then weakening in the rider. Dynamics-free — sound for
flows AND maps; the flows-only fence (header) concerns completeness,
which this statement deliberately does not express: the conclusion is a
certified LOWER BOUND family, and "the whole attractor lattice" is a
paper-side, flows-only reading (KW `cor_flows_att_iso`). -/
theorem sound_productLattice {H1 H2 : Hypernet}
    (d : PProductLatticeData H1 H2)
    (hL : d.left.Valid) (hR : d.right.Valid) :
    ∀ b ∈ d.prodBlocks, ∀ s ∈ b.facets,
      s.semanticallyValid (H1.disjUnion H2 d.disjoint)
\end{leansrc}
\declhead{BernsteinLogic.ProductLattice}{Hypernet.disjUnion}{\badge{BadgeGreen}{statement bound}}
\begin{leansrc}
/-- The disjoint union of two wire-disjoint hypernets: union typing,
concatenated applications; the three well-formedness predicates are
proven from the factors' plus disjointness. This is the composite system
φ₁ × φ₂'s hypernet, H₁ ⊔ H₂ (note §1.1). -/
def Hypernet.disjUnion (H1 H2 : Hypernet) (hd : H1.WireDisjoint H2) : Hypernet
\end{leansrc}
\declhead{BernsteinLogic.ProductLattice}{Hypernet.Valuation.restrict}{\badge{BadgeGreen}{statement bound}}
\begin{leansrc}
/-- Restrict a valuation of an extending hypernet to the extended one. -/
def Hypernet.Valuation.restrict {H H' : Hypernet} (σ : H'.Valuation)
    (hext : H.WireExtends H') : H.Valuation
\end{leansrc}
\declhead{BernsteinLogic.ProductLattice}{Formula.holds\_restrict}{\badge{BadgeGreen}{statement bound}}
\begin{leansrc}
/-- Formulas scoped to the extended hypernet's wires hold identically
under restriction. -/
theorem Formula.holds_restrict {H H' : Hypernet} (σ : H'.Valuation)
    (hext : H.WireExtends H') :
    ∀ (F : Formula), (∀ w ∈ F.freeWires, (H.wireType w).isSome) →
      (Formula.holds (σ.restrict hext) F ↔ Formula.holds σ F)
\end{leansrc}
\declhead{BernsteinLogic.ProductLattice}{PrimApp.semOk\_restrict}{\badge{BadgeGreen}{statement bound}}
\begin{leansrc}
/-- Per-application semantics survive restriction: every entry SemOk reads
is at a port wire, typed in the extended hypernet by type consistency. -/
theorem PrimApp.semOk_restrict {H H' : Hypernet} (σ : H'.Valuation)
    (hext : H.WireExtends H') {app : PrimApp}
    (hTC : app.TypeConsistent H.wireType) (h : app.SemOk σ) :
    app.SemOk (σ.restrict hext)
\end{leansrc}
\declhead{BernsteinLogic.ProductLattice}{prodBlocks\_confined}{\badge{BadgeGreen}{statement bound}}
\begin{leansrc}
/-- The conclusion family is itself wire-confined over the union — the
confinement half of what a NESTED application (product of products,
AG-inside/product-outside, note §2) consumes from its premise object. -/
theorem prodBlocks_confined {H1 H2 : Hypernet}
    (d : PProductLatticeData H1 H2) :
    ∀ b ∈ d.prodBlocks,
      (∀ w ∈ b.region.freeWires,
        ((H1.disjUnion H2 d.disjoint).wireType w).isSome) ∧
      ∀ s ∈ b.facets, s.wellScoped (H1.disjUnion H2 d.disjoint)
\end{leansrc}
\declhead{BernsteinLogic.ProductLattice}{Hypernet.WireDisjoint}{\badge{BadgeGreen}{definition bound}}
\begin{leansrc}
/-- No wire is typed in both hypernets — ProductLattice's side condition
(note §1.1: wires(H₁) ∩ wires(H₂) = ∅). A single shared wire demotes the
composition to the interface axis (Promise/SmallGain — note §2). -/
def Hypernet.WireDisjoint (H1 H2 : Hypernet) : Prop :=
  ∀ w : Wire, (H1.wireType w).isSome → (H2.wireType w).isSome → False
\end{leansrc}
\declhead{BernsteinLogic.ProductLattice}{CertifiedBlockLattice.Valid}{\badge{BadgeGreen}{definition bound}}
\begin{leansrc}
/-- The family is certified when every block's every facet sequent is
semantically valid. Validity enters `sound_productLattice` as a
hypothesis — the Step pattern: sub-proof warrants are supplied at the
use site (by `sound_v0` on the factor certificates, or manually). -/
def CertifiedBlockLattice.Valid {H : Hypernet}
    (L : CertifiedBlockLattice H) : Prop :=
  ∀ b ∈ L.blocks, ∀ s ∈ b.facets, s.semanticallyValid H
\end{leansrc}
\declhead{BernsteinLogic.ProductLattice}{citeFacet}{\badge{BadgeGreen}{definition bound}}
\begin{leansrc}
/-- Discharge BY CITATION (note §1.1): the composite facet sequent keeps
the cited factor facet's conclusion VERBATIM — polynomial and predicate
unchanged; the Bernstein grid of the composite facet is the factor's grid
replicated along the free coordinates, a coefficient identity, so
syntactic identity of the conclusion is the whole discharge — and
conjoins the rider (the other factor's box, riding along as free box
variables) onto the premise. -/
def citeFacet (s : Sequent) (rider : Formula) : Sequent :=
  ⟨.and s.pre rider, s.post⟩
\end{leansrc}
\declhead{BernsteinLogic.ProductLattice}{prodBlock}{\badge{BadgeGreen}{definition bound}}
\begin{leansrc}
/-- The product of two blocks: region = the conjunction of boxes; facets =
each factor's facets, cited with the other factor's region as rider.
(Every internal facet of an axis-aligned product is normal to one
factor's axis — note §1.2.) -/
def prodBlock (b1 b2 : BlockCert) : BlockCert :=
  ⟨.and b1.region b2.region,
   b1.facets.map (fun s => citeFacet s b2.region)
     ++ b2.facets.map (fun s => citeFacet s b1.region)⟩
\end{leansrc}
\declhead{BernsteinLogic.ProductLattice}{PProductLatticeData.prodBlocks}{\badge{BadgeGreen}{definition bound}}
\begin{leansrc}
/-- The conclusion family's meet-generators: all pairwise products.
⟨B₁ × B₂⟩∨ = the join-closure of THIS list (KW Prop
`prop:Lwedgevee-generated`); the joins — exponentially many — are never
enumerated: their sub-facets discharge via the paper-strength boundary
decomposition (unmechanized — it selects the factor facet and rider),
whose final step is `sound_productLattice_citeLeft/Right` below with
smaller riders. -/
def PProductLatticeData.prodBlocks {H1 H2 : Hypernet}
    (d : PProductLatticeData H1 H2) : List BlockCert :=
  d.left.blocks.flatMap fun b1 => d.right.blocks.map fun b2 => prodBlock b1 b2
\end{leansrc}
\ruleblock{Correspondence (Steele display $\leftrightarrow$ Lean)}
\badge{BadgeGreen}{reviewer-certified against the displayed texts}
\sourcenote{The dictionary below is a reviewer-certified statement of how the mechanized statements above interpret the displayed rule: display metavariable $\mapsto$ Lean binder, each display hypothesis $\mapsto$ the premise or side condition carrying it, the display conclusion $\mapsto$ the goal --- with any deliberate difference in width between display and mechanization named outright. Certification is recorded as a hash triple over (the displayed cell, the extracted statement texts, this dictionary\textquotesingle s own text); if any of the three changes, this block turns amber until a reviewer re-certifies. The correspondence claim itself remains editorial --- reviewer-grade, not machine-checked; its machine-checked successor types each soundness statement against the rule as a structured object, and supersedes this block per rule as it lands.}
The one rule of this book whose judgment sits a layer up: the
conclusion is not a sequent but a certified \emph{family}, so the
\rulename{Ax} frame applies one level in --- per facet sequent.
\begin{itemize}
\item $B_1, B_2 \mapsto$ \texttt{d.left}, \texttt{d.right} $:$
  \texttt{CertifiedBlockLattice} over \texttt{H1}, \texttt{H2};
  ``certified for $\varphi_i$\textquotesingle\textquotesingle{}
  $\mapsto$ \texttt{hL}, \texttt{hR} $:$ \texttt{Valid} --- every
  block\textquotesingle s every facet sequent semantically valid.
  Note what does \emph{not} appear: any semiflow. The certified
  statement is dynamics-free (the note above) --- the
  display\textquotesingle s $\varphi_1 \times \varphi_2$ is the
  intended reading of the facet sequents, not a consumed hypothesis.
\item $\mathrm{wires}(H_1) \cap \mathrm{wires}(H_2) = \emptyset$
  $\mapsto$ \texttt{d.disjoint} $:$ \texttt{H1.WireDisjoint H2}, and
  the composite carrier $H_1 \sqcup H_2 \mapsto$
  \texttt{Hypernet.disjUnion}.
\item the displayed conclusion (the meet-generator family
  certified) $\mapsto$ \texttt{sound\_productLattice}, which
  discharges every facet
  of every \emph{pairwise product} block (\texttt{prodBlocks} --- the
  join-closure\textquotesingle s meet-generators) --- the displayed
  family and the bound quantification now coincide. The full
  join-closure $\langle B_1 \times B_2 \rangle_{\vee}$ is the flavor
  text\textquotesingle s paper-strength reading: the joins,
  never enumerated, discharge facet-by-facet through the standing
  citation lemmas \texttt{sound\_productLattice\_citeLeft/Right}
  (certified alongside, at \emph{arbitrary} rider), with the boundary
  decomposition itself paper-strength --- no separate join statement
  is bound.
\item with no displayed row: well-scopedness --- the factor
  lattices\textquotesingle{} facets and regions are confined to their
  own hypernets (\texttt{facetsScoped}/\texttt{regionConfined}
  fields of \texttt{CertifiedBlockLattice}), which is what makes
  transport along the union injection sound; the undisplayed-pack
  class of the \rulename{PsdCert} evaluability rows.
\end{itemize}
Width note: the citation lemmas accept \emph{any} rider formula where
the rule\textquotesingle s conclusion only ever conjoins the other
factor\textquotesingle s region --- wider machinery in the safe
direction, and exactly the generality the join-block sub-facets
consume.

\ruleblock{Lean verifier receipts}
\sourcenote{module \texttt{BernsteinLogic.ProductLattice}; bound build \texttt{84b7018f79da\ldots} (bound from receipt 9ebc3d9af324d1f7, 2026-07-21T04:48:38Z). A receipt records that named checking layers accepted a build: (module, compiled-artifact SHA-256, environment SHA-256, checker layers), minted when the checking gate ran; the store receipts attesting the bound build are listed with the layers each run credited. Nothing here asks the reader to trust the store: every layer is an independently runnable program over the quoted sources --- rebuild the module and re-run the checkers to reproduce the verdict. Receipts attesting superseded builds credit no layer and are omitted from this rendering (the count of omissions, if any, is noted below).}
\begin{receipttable}
\receiptrow{224af3a8595b550f}{2026-07-21T01:49:09Z}{build+decorrelated}{\badge{BadgeGreen}{current}}
\receiptrow{7187ff3aa085cd58}{2026-07-21T01:58:41Z}{build+decorrelated}{\badge{BadgeGreen}{current}}
\receiptrow{9ebc3d9af324d1f7}{2026-07-21T04:48:38Z}{build+decorrelated}{\badge{BadgeGreen}{current}}
\receiptrow{fd5ec354b89369c0}{2026-07-21T05:28:29Z}{build+replay+hardened+decorrelated}{\badge{BadgeGreen}{current}}
\end{receipttable}
\sourcenote{7 further receipt(s) in the store attest superseded builds of this module and are omitted from this rendering.}
\layerline{\layeritem{build}{lake build (elaboration)}{\badge{BadgeGreen}{current}}\layeritem{replay}{leanchecker (official kernel replay)}{\badge{BadgeGreen}{current}}\layeritem{hardened}{safe\_verify (hardened screen)}{\badge{BadgeGreen}{current}}\layeritem{decorrelated}{lean4export + nanoda (decorrelated kernel)}{\badge{BadgeGreen}{current}}}
\sourcenote{\badge{BadgeGray}{comparator: none} no fixture in the three-implementation comparator harness (compiled / kernel-interpreted / reference-interpreted) exercises this rule.}
\ruleblock{HOL4 statement}
\absentbox{NO HOL STATEMENT YET}{The certificate-object stratum of this rule now has HOL twins, with
kernel-evaluated transcription receipts, across five ported theories:
the certified-block record and the citation constructor, the
product-block assembly, the certified block lattice with its
wellformedness and validity predicates, the applicability payload
record with its wellformedness predicate, the meet-generator
enumeration, and the valuation-restriction lemmas the transport
argument reads (verified against the scripts this writing;
\texttt{hol-strict} clean at the same 2026-07-22 run as the ported
entries). What has no HOL twin yet is exactly the
rule\textquotesingle s semantic content --- the transport-and-weaken
citation lemmas and the soundness theorem they assemble into (the
commentary\textquotesingle s own point: nothing in the certificate
stratum evaluates dynamics anywhere, so the semantic layer IS the
rule). The rule also has no arm in the ported proof-object family and
no wire tag on either side. Absent until the citation/soundness
stratum ports; the certificate-object groundwork above is in place
for that seating.
}
\ruleblock{HOL verifier receipts}
\absentbox{NONE}{no HOL verification receipt exists for this rule.}
\ruleblock{Proof}

\sourcenote{Genre: admissible rule (tier 2) --- a proof-tree-rewriting
recipe, at the certificate-family layer: one lattice node unpacks into
a family of cited facet sequents. Like \rulename{PConstExact}, part of
the discharge is by reduction to verified machinery rather than to
catalogue nodes --- the transport step is a semantic monotonicity
lemma with no catalogue counterpart, because hypernet extension is
invisible to the sequent calculus --- and the recipe says so.}

\textbf{Setting.}\enspace A \emph{certified attracting-block lattice}
over a hypernet $H$ (the \texttt{CertifiedBlockLattice} structure
quoted above) is a finite family of blocks --- each a region box
formula plus its internal-facet inwardness sequents --- carrying
explicit index-level lattice structure and wire confinement; the family
is \emph{valid} when every facet sequent is semantically valid over
$H$. The rule's premises: valid families $B_1$ over $H_1$ and $B_2$
over $H_2$, and wire-disjointness
$\wires(H_1) \cap \wires(H_2) = \emptyset$. Conclusion: the
join-closure $\langle B_1 \times B_2 \rangle_{\vee}$ is certified over
the disjoint union $H_1 \sqcup H_2$.

\begin{brecipe}[the \rulename{ProductLattice} instance --- one node,
a citation family]\label{rec:pl-instance}
The rewrite replaces the lattice node by one derivation fragment
\emph{per facet of per product block}, each of the same two-move shape:
\[
\inferrule*[Right=cite]{\;s_0 \rhd \text{a factor facet sequent (cited sub-proof)}\;}{\bigl\langle \mathrm{pre}(s_0) \wedge \mathrm{rider},\ \mathrm{post}(s_0) \bigr\rangle \ \text{over } H_1 \sqcup H_2}
\]
where the rider is the \emph{other} factor's box. The conclusion is
kept verbatim --- the polynomial and predicate of the cited facet,
unchanged --- because a composite facet of an axis-aligned product is
normal to one factor's axis: its inwardness polynomial \emph{is} that
factor's, the other factor's coordinates riding along as free box
variables. Only the meet-generators $\{b_1 \times b_2\}$ are
enumerated ($O(|B_1| \cdot |B_2|)$ blocks, $O(|B_1| + |B_2|)$
distinct cited factor sub-proofs, each reused across its products
with varying riders); the join-closure's remaining elements ---
exponentially many --- are never materialized. They discharge by the
same mechanism: a boundary point of a finite union of axis-aligned
blocks lies on the boundary of some member containing it (member
interiors are interior to the union), so a union block's boundary
decomposes into sub-facets, each a subset of some generator's facet
and hence again normal to one factor's axis; inwardness for the member
implies inwardness for the union (monotonicity under $\subseteq$), so
the sub-facet's sequent is the same factor citation with its premise
further constrained to the exposed portion --- and the citation
lemmas accept these further-constrained premises because they accept
an \emph{arbitrary} rider.
\end{brecipe}

\begin{blemma}[transport: validity is monotone under hypernet
extension]\label{lem:pl-transport}
Let $H'$ extend $H$'s wire typing and contain $H$'s applications, and
let $s$ be semantically valid over $H$ and well-scoped in $H$. Then
$s$ is semantically valid over $H'$.
\end{blemma}
\begin{proof}
An admissible valuation $\sigma$ of $H'$ restricts to a valuation of
$H$ (read each $H$-typed wire through the extension), and the
restriction is admissible: every entry a per-application semantic
predicate reads is at a port wire, typed in $H$ by type consistency,
where the restriction agrees with $\sigma$ by construction. Formulas
whose free wires are $H$-typed therefore evaluate identically under
$\sigma$ and its restriction --- by induction on terms and formulas,
the only base case being a wire read, which is where well-scopedness
is consumed. So from $\sigma \models \mathrm{pre}(s)$: the restriction
satisfies $\mathrm{pre}(s)$, validity over $H$ gives the restriction
satisfies $\mathrm{post}(s)$, and evaluation-agreement carries it back
to $\sigma$. This is \texttt{Sequent.semanticallyValid\_extend};
the well-scopedness hypothesis is not decoration --- a sequent reading
a wire \emph{untyped} in $H$ evaluates that read differently in an
extension that types it, and the induction's base case fails exactly
there.
\end{proof}

\begin{blemma}[weaken, standalone]\label{lem:pl-weaken}
If $\langle \mathrm{pre}, \mathrm{post} \rangle$ is semantically valid
over $H$, so is
$\langle \mathrm{pre} \wedge \mathrm{extra}, \mathrm{post} \rangle$,
for any formula $\mathrm{extra}$.
\end{blemma}
\begin{proof}
Project the premise (\texttt{Sequent.semanticallyValid\_weaken}). The
rider needs no side condition --- in particular it may mention wires
the cited factor has never heard of, which is the whole point.
\end{proof}

\begin{bprop}[soundness of \rulename{ProductLattice}]
\label{prop:pl-sound}
Given valid families over wire-disjoint hypernets, every facet sequent
of every product block is semantically valid over
$H_1 \sqcup H_2$.
\end{bprop}
\begin{proof}
A product block's facet list is, by construction (\texttt{prodBlock}),
the left factor's facets cited with the right factor's region as
rider, concatenated with the symmetric family. For a left-cited facet:
the disjoint union extends $H_1$'s wire typing and contains its
applications (disjointness makes the union's linearity, acyclicity,
and type consistency inherit from the factors --- this is where the
side condition does structural work), so
Lemma~\ref{lem:pl-transport} transports the factor's validity to the
union, and Lemma~\ref{lem:pl-weaken} conjoins the rider. Symmetrically
on the right. That is the entire proof
(\texttt{sound\_productLattice}, through the two standing citation
lemmas \texttt{sound\_productLattice\_citeLeft}/\texttt{citeRight}).
For the un-enumerated join blocks' sub-facets the same two lemmas
supply the per-generator transport and arbitrary-rider steps, with the
exposure constraints as riders --- but the join-closure extension
itself (the boundary decomposition of Recipe~\ref{rec:pl-instance}'s
argument) is paper-strength: no bound declaration represents finite
unions, their boundary decomposition, or the union-inwardness step,
and the mechanized theorem quantifies over the generator blocks only.
\end{proof}

\begin{bremark}[why wire-disjointness is not decoration]
\label{rem:pl-disjoint}
With a shared wire the recipe fails twice. Structurally, the disjoint
union is no longer well-formed as built (linearity of the concatenated
application lists is proven \emph{from} disjointness). Semantically, no
single transported sequent's validity would fail ---
Lemma~\ref{lem:pl-weaken} conjoins \emph{any} rider, shared wires
included --- but the assembled \emph{readout} would misdescribe the
composite: the identification of every product-block facet with a
cited factor facet (the other factor entering only as a rider) is
proven from disjointness, and with a shared wire the built object is
simply not a product of the factors.
A single shared wire is precisely an \emph{interface}, and composition
along an interface is \rulename{Promise}/small-gain territory --- the
rule's applicability condition is the boundary between the two
composition regimes, not a technicality.
\end{bremark}

\begin{bremark}[the completeness fence: flows only, and the
map counterexample]\label{rem:pl-fence}
Proposition~\ref{prop:pl-sound} is dynamics-free --- no premise
mentions a flow or a map, so the rule is sound for both. What is
\emph{flows-only} is the completeness reading from the source
literature \cite{kalieswehbe2026}: for flows, the attractor lattice of
a product is the coproduct of the factors' lattices, so the generated
family is everything. For maps it is not: with
$f(x) = -\tanh(2x)$ on $[-p, p]$ (where $\tanh(2p) = p$), the factor
attractor lattice is $\{\emptyset, P, X\}$ with $P = \{p, -p\}$, yet
$f \times f$ has the phase-locked period-two attractors
$\{(p,p), (-p,-p)\}$ and $\{(p,-p), (-p,p)\}$ --- neither generated by
factor attractors; the coproduct is a proper sublattice. (The
arithmetic: $\tanh(2x)$ leaves the diagonal with slope $2$, so a
fixed point $p \approx 0.96$ of $\tanh(2p) = p$ exists; $0$ repels,
$|f'(0)| = 2$; each product two-cycle attracts,
$(f^{2})'(\pm p) = 4(1 - p^{2})^{2} < 1$; and that the factor lattice
has nothing beyond $\emptyset, P, X$ is the source example's own
analysis \cite{kalieswehbe2026}.) The Lean
statement is shaped so this cannot overreach: it certifies a lower
bound family, and any decomposition-level readout from a map instance
must be stated as ``at least this structure,'' never ``exactly this.''
\end{bremark}

\begin{anchors}
Recipe \ref{rec:pl-instance}: citation \& product blocks & \texttt{citeFacet}, \texttt{prodBlock}, \texttt{PProductLatticeData.prodBlocks}\\
Lemma \ref{lem:pl-transport} & \texttt{Sequent.semanticallyValid\_extend}; restriction machinery \texttt{Hypernet.Valuation.restrict}, \texttt{Formula.holds\_restrict}, \texttt{PrimApp.semOk\_restrict}\\
Lemma \ref{lem:pl-weaken} & \texttt{Sequent.semanticallyValid\_weaken}\\
union well-formedness from disjointness & \texttt{Hypernet.disjUnion} (fields \texttt{linearity}, \texttt{acyclic}, \texttt{typeConsistent}); \texttt{Hypernet.WireDisjoint}\\
Proposition \ref{prop:pl-sound} & \texttt{sound\_productLattice}; \texttt{sound\_productLattice\_citeLeft}/\texttt{citeRight}\\
nested-application confinement & \texttt{prodBlocks\_confined}\\
family validity (premise form) & \texttt{CertifiedBlockLattice.Valid}\\
\end{anchors}


\blcatsection{Judgment Structure and Batteries}
\rulesection{37}{QSatFromRegion}{alias for \textcolor{DiffBlue}{\rulename{PConstExact}}}
\label{blcatrule-QSatFromRegion}
\ruleblock{The rule (Steele metanotation)}
\sourcenote{as realized in the mechanized development; source row: \texttt{\texttt{QTSL-rules.toml}} --- \emph{QSatFromRegion}. The metanotation statement is typeset by hand for notational quality; its mechanized ground truth is the Lean statement below, extracted from the sources at every build of this book and hash-bound. The rule \emph{roster} is machine-checked (every mechanized rule tag round-trips to exactly one catalogue row); the statement \emph{text} is editorial, checked by eye against the extracted Lean beside it.}
\begin{mathpar}
\inferrule*{
  \mathrm{child} : \mathsf{QRegionSat}(\varphi, R) \\
  \mathrm{eval} : \vdash x_0 \in R \ \text{(BL leaf; sublevel} \to \rulename{PConstExact}\text{, box} \to \text{registry interval check)} \\
  \mathsf{SInit}(\varphi) \in s.\forall
}{
  \mathcal{M} \vDash s
}
\end{mathpar}

\par\medskip
\rulename{QSatFromRegion} is the extension theory\textquotesingle s
closing move: it turns a region-level certificate into the root
satisfaction verdict. The workhorse judgment of the extension page
concludes region-level facts --- every point of a region $R$ satisfies
the formula $\varphi$, with $R$ under-approximating the
formula\textquotesingle s denotation --- and this rule is the single
bridge from that shape to the top-level claim
$\mathcal{M} \vDash s$: cite the child, check that the
specification\textquotesingle s initial point $x_0$ lies in $R$ by
exact rational arithmetic, and require that the warranted safety
component of the specification is the initial-state injection
$\mathsf{SInit}(\varphi)$ of the child\textquotesingle s formula. A producer uses
it exactly once per initialized certificate, as the last row.

What the checker recomputes: the membership leaf --- for a sublevel
atom $\{x : V(x) \le y\}$, exact substitution of $x_0$ into
$V$ and one closed comparison, which is precisely the ground-fact rule
\rulename{PConstExact}\textquotesingle s own check; for a box atom,
per-coordinate containment recomputed by the checker\textquotesingle s
side-condition registry --- and the component tie: the warrant row names which safety
component it warrants, anything other than the $\mathsf{SInit}$ image
of the child\textquotesingle s formula refuses at load (a menu of
one), and that the named component is among the
specification\textquotesingle s consumed components is admission
conformance.

What would break without each hypothesis: drop the child and $R$ is an
arbitrary set nobody certified --- the conclusion is unmoored from any
proof; drop the eval leaf and $x_0$ may lie outside $R$, so the
child\textquotesingle s universal never instantiates (a certificate
could certify a region the system never starts in); drop the component tie
and a certificate about some subformula --- or an unrelated formula
--- could claim a component it never certified.

\cataloguenote{The soundness content is one line --- instantiation of the
child\textquotesingle s universal at $x_0$ --- which is why the rule
prices as an alias even though it has two premise fields: the child is
a judgment \emph{citation} (an edge under the tower discipline, the
\rulename{ProofWellFormedness} section), and the sole derivation node the
respelling produces is the \rulename{PConstExact} leaf.
Mechanization state, plainly: the bound statements below are that one
line and its assembly. \texttt{sound\_QSatFromRegion} (module
\texttt{BernsteinLogicAux.QtslAliases}) is the instantiation itself,
generic in the satisfaction predicate and profile-free --- axiom-free,
even: the kernel certifies it depends on no axioms at all ---
and \texttt{sound\_QSatFromRegion\_sublevel} composes it over the
exact-$\Q$ membership leaf (the \rulename{RootEval} identities, this
page), which is the sole derivation node named above. The
payload-level elaborator is bound beside them (module
\texttt{BernsteinLogic.QtslElaborators}, the batch
template\textquotesingle s unpack made concrete):
\texttt{qsatFromRegionElab} takes the rule payload to the one
\rulename{PConstExact} payload --- an inhabitant of the
target\textquotesingle s own payload type, comparator confinement
discharged at elaboration --- with
\texttt{qsatFromRegionElab\_sound} the factoring through
\texttt{sound\_pConstExact} and \texttt{qsatFromRegionElab\_sat} the
judgment-level composition over the elaborated leaf; the semantic
content carries no profile residue, so the entry carries no
conditional flag. Inherited flag, travelling with the alias: the target
excludes the $\neq$ comparator from its admissible set (a
second-soundness-invariant exclusion, unpacked in the
target\textquotesingle s section); the membership leaf uses $\le$, and
any port must preserve the exclusion.
}
\ruleblock{Lean statement}
\declhead{BernsteinLogicAux.QtslAliases}{sound\_QSatFromRegion}{\badge{BadgeGreen}{statement bound}}
\begin{leansrc}
/-- **`sound_QSatFromRegion`** (QTSL-rules.toml row `QSatFromRegion`; the
catalogue's "Soundness: one line — instantiation of the universal"): the
child certificate concludes region-level satisfaction (`R` under-approximates
the truth set: every point of `R` satisfies), the eval leaf places the
specification's initial point in `R`, and the root claim follows by
instantiation. Generic in the satisfaction predicate: for the probabilistic
judgments the extension page concludes, `sat x` reads as the concluded
bound at start `x` (e.g. `sound_QSafetyBattery`'s conclusion), and the rule
is profile-free. The root tie `s.root = φ` is a syntactic side condition
recomputed by the checker; it selects *which* predicate this lemma is
instantiated at, and carries no semantic content of its own. -/
theorem sound_QSatFromRegion {X : Type*} {R : Set X} {sat : X → Prop} {x₀ : X}
    (hchild : ∀ x ∈ R, sat x) (heval : x₀ ∈ R) : sat x₀
\end{leansrc}
\declhead{BernsteinLogicAux.QtslAliases}{sound\_QSatFromRegion\_sublevel}{\badge{BadgeGreen}{statement bound}}
\begin{leansrc}
/-- **`QSatFromRegion` assembled over its `RootEval` leaf, sublevel shape**:
the child's region-level universal on the cast sublevel atom, instantiated
at the cast initial point through the exact-ℚ evaluation — the sole
derivation node the respelling produces, composed. -/
theorem sound_QSatFromRegion_sublevel {σ : Type*} {sat : (σ → ℝ) → Prop}
    {P : MvPolynomial σ ℚ} {thr : ℚ} {xq : σ → ℚ}
    (hchild : ∀ x ∈ {x : σ → ℝ |
      MvPolynomial.eval x (MvPolynomial.map (Rat.castHom ℝ) P) ≤ (thr : ℝ)}, sat x)
    (heval : MvPolynomial.eval xq P ≤ thr) :
    sat (fun i => (xq i : ℝ))
\end{leansrc}
\par\medskip
\textbf{The payload-level elaborator, displayed whole.}\enspace The
unpack builds the target\textquotesingle s payload from the
rule\textquotesingle s wire data --- the eval field \emph{is} the
\rulename{RootEval} leaf (that section binds the closed-term builder
and the wire-level exactness identity it rides).

\declhead{BernsteinLogic.QtslElaborators}{qsatFromRegionElab\_sound}{\badge{BadgeGreen}{statement bound}}
\begin{leansrc}
/-- The factoring, BL face — through `sound_pConstExact`, as the amber
box demanded. -/
theorem qsatFromRegionElab_sound {H : Hypernet} {seq : Sequent}
    {x0 : Nat → Rat} {V : Poly} {thr : Rat}
    (hEval : (rootEvalFormula x0 V thr).evalClosed = some true)
    (hpost : seq.post = rootEvalFormula x0 V thr) :
    seq.semanticallyValid H
\end{leansrc}
\declhead{BernsteinLogic.QtslElaborators}{qsatFromRegionElab\_sat}{\badge{BadgeGreen}{statement bound}}
\begin{leansrc}
/-- The judgement-level composition over the elaborated leaf
(`prop:qsr-sound` at the wire): instantiation of the child's universal
at x₀, the membership hypothesis supplied by the checked leaf through
the ℚ-face exactness identity. `sound_QSatFromRegion` is generic in the
satisfaction predicate and the carrier, so the composition consumes the
bound lemma verbatim — sound_A = (instantiation) ∘ (elaborated leaf). -/
theorem qsatFromRegionElab_sat {x0 : Nat → Rat} {V : Poly} {thr : Rat}
    {sat : (Nat → Rat) → Prop}
    (hchild : ∀ x ∈ {x : Nat → Rat | Poly.eval x V ≤ thr}, sat x)
    (hEval : (rootEvalFormula x0 V thr).evalClosed = some true) :
    sat x0
\end{leansrc}
\declhead{BernsteinLogic.QtslElaborators}{qsatFromRegionElab}{\badge{BadgeGreen}{definition bound}}
\begin{leansrc}
/-- **The QSatFromRegion payload elaborator**, sublevel case: the
rule's wire payload elaborates to exactly the RootEval leaf's payload —
the leaf-level respelling has its own row (`RootEval`), and this rule's
eval field IS that leaf (the consumption-site anchor of the RootEval
section). -/
def qsatFromRegionElab (H : Hypernet) (seq : Sequent) (x0 : Nat → Rat)
    (V : Poly) (thr : Rat)
    (hEval : (rootEvalFormula x0 V thr).evalClosed = some true)
    (hpost : seq.post = rootEvalFormula x0 V thr) :
    PConstExactData H seq :=
  rootEvalElab H seq x0 V thr hEval hpost
\end{leansrc}
\ruleblock{Correspondence (Steele display $\leftrightarrow$ Lean)}
\badge{BadgeGreen}{reviewer-certified against the displayed texts}
\sourcenote{The dictionary below is a reviewer-certified statement of how the mechanized statements above interpret the displayed rule: display metavariable $\mapsto$ Lean binder, each display hypothesis $\mapsto$ the premise or side condition carrying it, the display conclusion $\mapsto$ the goal --- with any deliberate difference in width between display and mechanization named outright. Certification is recorded as a hash triple over (the displayed cell, the extracted statement texts, this dictionary\textquotesingle s own text); if any of the three changes, this block turns amber until a reviewer re-certifies. The correspondence claim itself remains editorial --- reviewer-grade, not machine-checked; its machine-checked successor types each soundness statement against the rule as a structured object, and supersedes this block per rule as it lands.}
The displayed premises map onto \texttt{sound\_QSatFromRegion}
verbatim: the child $\mathsf{QRegionSat}(\varphi, R)$ $\mapsto$
\texttt{hchild : ∀ x ∈ R, sat x} (region-level satisfaction,
compiled); the eval leaf $\vdash x_0 \in R$ $\mapsto$ \texttt{heval :
x₀ ∈ R}; the conclusion $\mathcal{M} \vDash s$ $\mapsto$ \texttt{sat
x₀} --- universal instantiation, one line, axiom-free. The component tie
($\mathsf{SInit}(\varphi) \in s.\forall$, the refuse-mode warrant
naming) has deliberately no Lean counterpart: it
is the syntactic side condition that selects \emph{which} predicate
\texttt{sat} the lemma is instantiated at, recomputed by the checker,
carrying no semantic content --- the statement\textquotesingle s
docstring says exactly this. \texttt{sound\_QSatFromRegion\_sublevel}
is the assembled form over the exact-$\Q$ \rulename{RootEval} leaf:
child universal on the cast sublevel atom + exact-$\Q$ evaluation
\texttt{heval : MvPolynomial.eval xq P ≤ thr} $\Rightarrow$
satisfaction at the cast initial point --- the sole derivation node
the respelling produces, composed. Width note: \texttt{sat : X →
Prop} is generic --- probabilistic conclusions read \texttt{sat x} as
the concluded bound at start \texttt{x}, so the alias is
profile-free by binder shape; wider than any single displayed
instantiation, safe direction.

\ruleblock{Lean verifier receipts}
\sourcenote{module \texttt{BernsteinLogicAux.QtslAliases}; bound build \texttt{570599779c29\ldots} (bound from receipt 5a60f4d01fd81ac5, 2026-07-21T05:39:02Z). A receipt records that named checking layers accepted a build: (module, compiled-artifact SHA-256, environment SHA-256, checker layers), minted when the checking gate ran; the store receipts attesting the bound build are listed with the layers each run credited. Nothing here asks the reader to trust the store: every layer is an independently runnable program over the quoted sources --- rebuild the module and re-run the checkers to reproduce the verdict. Receipts attesting superseded builds credit no layer and are omitted from this rendering (the count of omissions, if any, is noted below).}
\begin{receipttable}
\receiptrow{0a3819a978e31127}{2026-07-21T04:51:36Z}{build+decorrelated}{\badge{BadgeGreen}{current}}
\receiptrow{0bb1074df705d863}{2026-07-21T05:21:58Z}{build+decorrelated}{\badge{BadgeGreen}{current}}
\receiptrow{5a60f4d01fd81ac5}{2026-07-21T05:39:02Z}{build+replay+hardened+decorrelated}{\badge{BadgeGreen}{current}}
\end{receipttable}
\sourcenote{11 further receipt(s) in the store attest superseded builds of this module and are omitted from this rendering.}
\sourcenote{module \texttt{BernsteinLogic.QtslElaborators}; bound build \texttt{0205ec213b91\ldots} (hashed from the built olean at bind time). A receipt records that named checking layers accepted a build: (module, compiled-artifact SHA-256, environment SHA-256, checker layers), minted when the checking gate ran; the store receipts attesting the bound build are listed with the layers each run credited. Nothing here asks the reader to trust the store: every layer is an independently runnable program over the quoted sources --- rebuild the module and re-run the checkers to reproduce the verdict. Receipts attesting superseded builds credit no layer and are omitted from this rendering (the count of omissions, if any, is noted below).}
\begin{receipttable}
\receiptrow{a237f4cff27cb15a}{2026-07-22T10:06:47Z}{build+decorrelated}{\badge{BadgeGreen}{current}}
\end{receipttable}
\layerline{\layeritem{build}{lake build (elaboration)}{\badge{BadgeGreen}{current}}\layeritem{replay}{leanchecker (official kernel replay)}{\badge{BadgeGreen}{current}}\layeritem{hardened}{safe\_verify (hardened screen)}{\badge{BadgeGreen}{current}}\layeritem{decorrelated}{lean4export + nanoda (decorrelated kernel)}{\badge{BadgeGreen}{current}}}
\sourcenote{\badge{BadgeGray}{comparator: none} no fixture in the three-implementation comparator harness (compiled / kernel-interpreted / reference-interpreted) exercises this rule.}
\ruleblock{HOL4 statement}
\absentbox{NO HOL STATEMENT YET}{The HOL4 development\textquotesingle s object is the base
calculus\textquotesingle s wire-format checker (theory
\texttt{bern\_catalogue} ports the base catalogue\textquotesingle s
check-function cone); no HOL development of the extension theory
exists, so there is no HOL statement for any extension rule to bind.
Absent rather than out-of-scope: the bridge\textquotesingle s baked
realization is the extension theory\textquotesingle s warrant seam
(\texttt{qsat\_warrant}), so a future HOL port of an
extension-theory checker would state it. The one leaf the alias adds
checks through the base ground-fact rule \rulename{PConstExact}, whose
own HOL statement is not yet writable either (it postdates the HOL4
port and its wire tag is unallocated --- see that
rule\textquotesingle s section).
}
\ruleblock{HOL verifier receipts}
\absentbox{NONE}{no HOL verification receipt exists for this rule.}
\ruleblock{Proof}

\sourcenote{Genre: alias rule (tier 3) --- a respelling template: the
1$\to$1 unpack, the filled-in target application, the payload-transport
table, the exactness identity, and the soundness corollary through the
target\textquotesingle s lemma. Two riders. First, the tier criterion
counts derivation \emph{nodes}, not premise fields: this rule has two
premise fields, but the child is a judgment \emph{citation} --- an edge
under the tower discipline (the \rulename{ProofWellFormedness} section), not
a node any unpack must produce --- so the rule\textquotesingle s proof
content per instance is at most one node: in the sublevel branch, the
single \rulename{PConstExact} application discharging its membership
leaf; in the box branch, no node at all (the leaf is registry
recomputation). The child instantiation and the component tie are
definitional transports along edges, not nodes any unpack produces ---
that, and nothing stronger, is the sense in which the rule is an
alias.
Second, the semantic composition is mechanized and bound in the entry,
and so --- now that the extension certificate schema exists --- is the
payload-level elaborator (the batch template\textquotesingle s unpack,
at the wire level; the remark below names the bound functions).}

\begin{brecipe}[the alias template, one extends-edge up]\label{rec:qsr-template}
An alias at the satisfaction level is warranted by a 1$\to$1 respelling
into a single application of a target rule, with judgment citations
transported as \emph{edges} and the concluded judgment recovered
definitionally from the cited child. Here the target is the base
calculus\textquotesingle s ground-fact rule \rulename{PConstExact}; the
one new node is the membership leaf; and the $\vDash$-level conclusion
is an instantiation of the child\textquotesingle s universal
(Proposition~\ref{prop:qsr-inst}), not a derived fact needing catalogue
structure.
\end{brecipe}

\begin{brecipe}[the filled-in application]\label{rec:qsr-filled}
The rule\textquotesingle s fields: a child judgment $\mathrm{child} :
\mathsf{QRegionSat}(\varphi, R)$ --- every point of $R$ satisfies
$\varphi$, with $R$ under-approximating the formula\textquotesingle s
denotation --- a membership leaf $\mathrm{eval}$, and the component tie
$\mathsf{SInit}(\varphi) \in s.\forall$ (the warranted safety
component is the initial-state injection of the child\textquotesingle s
formula --- the record carries component lists, not a single root, and
the warrant row names its component). For a sublevel region atom
$R = \{x : V(x) \le y\}$, with $V$ a polynomial over $\Q$
and $x_0$ the specification\textquotesingle s rational initial point,
the eval leaf \emph{is} the \rulename{PConstExact} node
\begin{mathpar}
\inferrule*[Right=\rulename{PConstExact}]{
  V(x_0) \le y \ \text{holds under exact evaluation in } \Q
}{
  \Phi \vdash V(x_0) \le y
}
\end{mathpar}
with arbitrary premise $\Phi$, as the target licenses. $V(x_0)$ is the
closed term obtained by exact substitution of $x_0$\textquotesingle s
coordinates --- wire-free, so the target\textquotesingle s verified
evaluator decides it, with comparator $\le$ inside the
target\textquotesingle s admissible set. The leaf-level respelling has
its own row on this page (\rulename{RootEval}), whose section carries
the transport and exactness detail; this section carries the
judgment-level composition. For a \emph{box} atom the leaf is not a
\rulename{PConstExact} node: per-coordinate interval containment is
recomputed directly by the checker\textquotesingle s side-condition
registry --- the alias reading covers the sublevel branch, and the box
branch is stated here so the respelling never overclaims.
\end{brecipe}

\begin{brecipe}[payload transport --- field for field]\label{rec:qsr-transport}
\begin{center}\small
\begin{tabular}{@{}ll@{}}
\textbf{\rulename{QSatFromRegion} field} & \textbf{destination}\\
\midrule
$\mathrm{child}$ & not transported --- a judgment citation (edge; the tower discipline)\\
$\mathrm{eval}$ & the \rulename{PConstExact} payload: closed comparison $V(x_0) \le y$\\
component tie $\mathsf{SInit}(\varphi) \in s.\forall$ & the concluded judgment\textquotesingle s formula\\
concluded $\mathcal{M} \vDash s$ & instantiation of the child\textquotesingle s universal at $x_0$\\
\end{tabular}
\end{center}
\end{brecipe}

\begin{bprop}[the instantiation identity]\label{prop:qsr-inst}
If $R \subseteq \{x : \mathcal{M}, x \vDash \varphi\}$ and
$x_0 \in R$, then $\mathcal{M}, x_0 \vDash \varphi$; with
the component tie, this is $\mathcal{M} \vDash s$. Moreover,
for the sublevel atom the membership hypothesis is \emph{exactly} the
checked leaf, as an identity of exact rationals:
\[
x_0 \in R \iff V(x_0) \le y,
\]
because substituting a rational point into a rational polynomial
evaluates with no rounding --- the checker\textquotesingle s
$\mathsf{true}$ verdict \emph{is} membership, not an approximation to
it (the \rulename{RootEval} section proves this as its exactness
identity).
\end{bprop}
\begin{proof}
Instantiate the universal at $x_0$; with the component tie
the instance \emph{is} $\mathcal{M} \vDash s$. For the moreover clause,
two equivalences compose: $x_0 \in R \iff V(x_0) \le y$ is
the sublevel region\textquotesingle s membership predicate read at
$x_0$ --- a definition, not a computation --- and the identification of
the checker\textquotesingle s $\mathsf{true}$ verdict with that exact
inequality is the \rulename{RootEval} exactness identity, cited rather
than re-proved. The instantiation is the rule\textquotesingle s entire
own content --- one node\textquotesingle s worth, which is why it
prices as an alias.
\end{proof}

\begin{bprop}[soundness corollary]
\label{prop:qsr-sound}
In the sublevel branch,
$\mathrm{sound}_{\mathsf{QSatFromRegion}} = $ \texttt{sound\_pConstExact}
(applied to the eval leaf) $\fatsemi$ (instantiation at $x_0$); in the
box branch the recomputed containment supplies the membership fact and
the same instantiation finishes. In both branches the
child\textquotesingle s validity is supplied by the induction hypothesis
of the tower discipline\textquotesingle s induction
(\rulename{ProofWellFormedness} section).
\end{bprop}
\begin{proof}
Sublevel branch: \texttt{sound\_pConstExact} (the target\textquotesingle s
kernel-verified lemma, quoted in the \rulename{PConstExact} section)
turns the $\mathsf{true}$ evaluation into the validity of
$V(x_0) \le y$; Proposition~\ref{prop:qsr-inst} finishes.
Box branch: the recomputed containment plays the leaf\textquotesingle s
role --- per-coordinate comparison of exact rationals, the box face of
the same no-rounding exactness the \rulename{RootEval} section proves
for sublevel leaves, so the recompute\textquotesingle s verdict is
membership itself; Proposition~\ref{prop:qsr-inst}\textquotesingle s
instantiation, which never asked how membership was obtained, finishes
identically.
\end{proof}

\begin{bremark}[the payload-level elaborator, mechanized]
\label{rem:qsr-elab}
The \emph{semantic} composition above is bound in the entry:
\texttt{sound\_QSatFromRegion} is the instantiation (axiom-free), and
\texttt{sound\_QSatFromRegion\_sublevel} is
Proposition~\ref{prop:qsr-sound}\textquotesingle s statement over the
exact-$\Q$ leaf (the \rulename{RootEval} identities). The payload-level
elaborator --- previously this page\textquotesingle s one absent
artefact, held back until the extension certificate schema existed ---
is bound with it: \texttt{qsatFromRegionElab} (module
\texttt{BernsteinLogic.QtslElaborators}) takes the rule\textquotesingle
s wire payload to (in the sublevel case) exactly the one
\rulename{PConstExact} payload of Recipe~\ref{rec:qsr-transport}, as an
inhabitant of the target\textquotesingle s own payload type --- the
comparator confinement discharged at elaboration, once ---
with \texttt{qsatFromRegionElab\_sound} the factoring through
\texttt{sound\_pConstExact} and \texttt{qsatFromRegionElab\_sat} the
judgment-level composition over the elaborated leaf, the
checker\textquotesingle s verdict entering as membership through the
wire-level exactness identity \texttt{rootEvalFormula\_true\_iff} (the
$\Q$ face of the \rulename{RootEval} identities). The box branch
remains checker-side registry recompute, by design --- nothing to
elaborate.
\end{bremark}

\begin{bremark}[the inherited comparator confinement]
\rulename{PConstExact} excludes one comparator ($\neq$) from its
admissible set, for reasons of a second soundness invariant unpacked in
that rule\textquotesingle s section. The membership leaf uses $\le$
only --- comfortably inside --- but any port or mechanization of this
rule must \emph{preserve} the exclusion rather than re-derive the
comparator set from the semantic lemma alone; the flag travels with
the alias.
\end{bremark}

\begin{anchors}
target lemma (Prop.~\ref{prop:qsr-sound}) & \texttt{sound\_pConstExact}; \texttt{Formula.holds\_of\_evalClosed\_true} (quoted in the \rulename{PConstExact} section)\\
leaf-level transport + exactness & the \rulename{RootEval} section (this page)\\
child validity & the tower discipline\textquotesingle s induction (\rulename{ProofWellFormedness} section)\\
alias-level corollary & \texttt{sound\_QSatFromRegion} + \texttt{sound\_QSatFromRegion\_sublevel} (bound in the entry)\\
payload-level elaborator (\autoref{rem:qsr-elab}) & \texttt{qsatFromRegionElab} + \texttt{qsatFromRegionElab\_sound} + \texttt{qsatFromRegionElab\_sat} (module \texttt{BernsteinLogic.QtslElaborators}, bound in the entry)\\
box-atom branch & checker-side registry recompute, by design --- no rule anchor\\
\end{anchors}


\rulesection{38}{QGuardIntro}{core rule}
\label{blcatrule-QGuardIntro}
\ruleblock{The rule (Steele metanotation)}
\sourcenote{as realized in the mechanized development; source row: \texttt{\texttt{QTSL-rules.toml}} --- \emph{QGuardIntro}. The metanotation statement is typeset by hand for notational quality; its mechanized ground truth is the Lean statement below, extracted from the sources at every build of this book and hash-bound. The rule \emph{roster} is machine-checked (every mechanized rule tag round-trips to exactly one catalogue row); the statement \emph{text} is editorial, checked by eye against the extracted Lean beside it.}
\begin{mathpar}
\inferrule*{
  \mathrm{child} : \mathsf{QRegionSat}(\varphi, R) \\
  \mathsf{SInit}(\varphi) \in s.\forall
}{
  \mathcal{M} \vDash \mathsf{Guard}(R,\ s)
}
\end{mathpar}

\par\medskip
\rulename{QGuardIntro} is implication introduction at the judgment
stratum: from a region-level certificate alone --- every state of $R$
satisfies $\varphi$ --- it concludes the \emph{guarded} claim
$\mathcal{M} \vDash \mathsf{Guard}(R, s)$, with no initial-membership
premise anywhere. The hypothesis that the unguarded route discharges
at the ground seam ($X_0 \subseteq R$, checked by exact rational
arithmetic) is instead discharged \emph{into the specification}: a
guarded specification demands its body\textquotesingle s guarantee
only from initial states inside $R$ and promises nothing off-region.
The guard region imports the certificate-side region grammar
(sublevel, box, affine) into the specification record at exactly one
slot; the record still binds no model and fixes no initial condition
--- the guard is a predicate \emph{on} the subject\textquotesingle s
initial data, not initial data.

What the checker recomputes: nothing beyond the child\textquotesingle s
own obligations --- that is the point. The two ties are load-time
conformance, not computation: the guarded warrant row names its
safety component, and anything other than the initial-state injection
of the warrant\textquotesingle s formula refuses at load (a menu of
one); the guard region stored on the record is derived from the
warrant row itself (store-and-check: a disagreeing producer row is a
second region on a one-region slot, refused at conformance).

What would break without each hypothesis: drop the child and the
guarded claim is unmoored from any proof; drop the component tie and
a certificate about an unrelated formula could claim the guarded
specification; drop the disjointness from the unguarded warrant
relation and the initial-membership demand would fire on a guarded
introduction --- exactly the hypothesis the former exists to
internalize.

\cataloguenote{The semantics is \emph{pointwise} over the credal stack: a behavior
object satisfies $\mathsf{Guard}(R, s)$ iff every admitted law whose
initial marginal is supported in $\llbracket R \rrbracket$ meets
$s$\textquotesingle s demand --- at demonic-box initial data this is
exactly conditioning the belief evaluation on initial membership.
Pointwise is \emph{strictly stronger} than the blanket conditional
(``if all of $X_0$ lies in $R$ then\ldots\textquotedbl), and the
strictness is a theorem, not a taste: the blanket reading would
vacuously discharge any guard the initial set merely straddles. The
class algebra stays a complement-free distributive lattice --- the
guard\textquotesingle s antecedent is a ground condition on system
data, not a behavior class, which is why the conditional lives at the
record stratum and no formula-level implication former exists or is
introduced.
}
\ruleblock{Theory rows (geolog)}
\sourcenote{The rule\textquotesingle s certificate content as it is declared in the geolog theory the checker loads --- payload schema, side-condition demands, and axiom rows, quoted from the theory source at every build of this catalogue (a stale quotation cannot exist). Axiom modes: \texttt{chase} = the engine derives; \texttt{report} = a violation is a loud judgment; \texttt{refuse} = a violation gates admission. Side-condition relations are discharged by bounded evaluation at the engine\textquotesingle s ground seam, never by derivation.}
\par\medskip
The baked realization: the guard slot on the specification record,
the guarded-route warrant relation (disjoint-by-relation from the
unguarded one, so the initial-membership axiom structurally never
fires on a guarded introduction), and the ties.

\declhead{theory Qtsl}{spec\_guard}{\badge{BadgeGreen}{quoted from the theory source at render}}
\begin{leansrc}
spec_guard : {spec: Spec, region: QRegionAtom} -> Prop;
\end{leansrc}
\declhead{theory Qtsl}{qsat\_guard\_warrant}{\badge{BadgeGreen}{quoted from the theory source at render}}
\begin{leansrc}
qsat_guard_warrant : {sat: QSat, warrant: QForallSat} -> Prop;
\end{leansrc}
\declhead{theory Qtsl}{ax/qguard/region\_tie}{\badge{BadgeGreen}{quoted from the theory source at render}}
\begin{leansrc}
chase ax/qguard/region_tie : forall r : QSat, c : QForallSat,
    sp : Spec.
  {sat: r, warrant: c} qsat_guard_warrant,
  r QSat/spec = sp
  |- {spec: sp, region: c QForallSat/region} spec_guard;
\end{leansrc}
\declhead{theory Qtsl}{ax/wf/qguard\_warrant\_member}{\badge{BadgeGreen}{quoted from the theory source at render}}
\begin{leansrc}
report ax/wf/qguard_warrant_member : forall r : QSat, sp : Spec,
    c : QForallSat.
  {sat: r, warrant: c} qsat_guard_warrant,
  r QSat/spec = sp
  |- {spec: sp, f: c QForallSat/component} qtsl_warrant_member;
\end{leansrc}
\ruleblock{Lean statement}
\par\medskip
\textbf{The guarded denotation and the introduction\textquotesingle s
soundness, displayed whole.}\enspace \texttt{GuardSat} is the
pointwise reading; the strictness witness against the blanket
conditional and the wrong-region countermodel bind beside the
soundness statements.

\declhead{QtslAux.QGuard}{sound\_QGuardIntro}{\badge{BadgeGreen}{statement bound}}
\begin{leansrc}
/-- Soundness of the introduction rule: a region warrant — the body's
demand holds from every state of `R` — yields the guarded claim with no
initial-membership premise, for any initial set. The membership hypothesis
that region-warranted satisfaction discharges at the ground seam is here
discharged into the former (implication introduction). -/
theorem sound_QGuardIntro {init R : Set X} {sat : X → Prop}
    (hchild : ∀ x ∈ R, sat x) : GuardSat init R sat
\end{leansrc}
\declhead{QtslAux.QGuard}{sound\_QGuardIntro\_sublevel}{\badge{BadgeGreen}{statement bound}}
\begin{leansrc}
/-- The rational-sublevel guard (a certificate family's sublevel set at an
exact-rational threshold, single-piece form): a warrant on the real
sublevel set of the coefficient-cast polynomial yields the guarded claim
with that sublevel set in the guard slot, for any initial set. -/
theorem sound_QGuardIntro_sublevel {σ : Type*} {init : Set (σ → ℝ)}
    {sat : (σ → ℝ) → Prop} {P : MvPolynomial σ ℚ} {thr : ℚ}
    (hchild : ∀ x ∈ {x : σ → ℝ |
        MvPolynomial.eval x (MvPolynomial.map (Rat.castHom ℝ) P) ≤ (thr : ℝ)},
      sat x) :
    GuardSat init
      {x : σ → ℝ |
        MvPolynomial.eval x (MvPolynomial.map (Rat.castHom ℝ) P) ≤ (thr : ℝ)}
      sat
\end{leansrc}
\declhead{QtslAux.QGuard}{sound\_QGuardIntro\_prob}{\badge{BadgeGreen}{statement bound}}
\begin{leansrc}
/-- Kernel-form corollary in the battery idiom: a per-initial-state
probability floor warranted on the region yields the guarded floor. -/
theorem sound_QGuardIntro_prob [MeasurableSpace X] {init R : Set X}
    {μ : X → Measure (ℕ → X)} {p : ENNReal} {A : Set (ℕ → X)}
    (hchild : ∀ x ∈ R, p ≤ μ x A) :
    GuardSat init R fun x => p ≤ μ x A
\end{leansrc}
\declhead{QtslAux.QGuard}{exists\_blanket\_not\_guardSat}{\badge{BadgeGreen}{statement bound}}
\begin{leansrc}
/-- Strictness witness for the semantics adjudication, and the permanent
twin of the dropped-membership mutant: the blanket reading does not imply
the pointwise one. On `Bool` with the full initial set, guard region
`{true}`, and an everywhere-false demand, the blanket conditional holds
vacuously (the initial set is not contained in the region) while the
pointwise guard is refuted at `true`. A blanket semantics would license
claims the pointwise semantics rejects; the former adopts the pointwise
reading. -/
theorem exists_blanket_not_guardSat :
    ∃ (init R : Set Bool) (sat : Bool → Prop),
      GuardSatBlanket init R sat ∧ ¬ GuardSat init R sat
\end{leansrc}
\declhead{QtslAux.QGuard}{exists\_not\_guardSat\_wrong\_region}{\badge{BadgeGreen}{statement bound}}
\begin{leansrc}
/-- Wrong-region countermodel (the permanent twin of the region-substitution
mutant): a warrant on one region does not license a guard on an unrelated
one — an everywhere-false demand is vacuously warranted on the empty
region, and the guard on the full region is refuted. -/
theorem exists_not_guardSat_wrong_region :
    ∃ (init R R' : Set Bool) (sat : Bool → Prop),
      (∀ x ∈ R, sat x) ∧ ¬ GuardSat init R' sat
\end{leansrc}
\declhead{QtslAux.QGuard}{GuardSat}{\badge{BadgeGreen}{definition bound}}
\begin{leansrc}
/-- The guarded-satisfaction denotation, pointwise (demonic-initial-set)
form: every admitted initial point inside the guard region meets the body's
demand; initial points outside the region are unconstrained. -/
def GuardSat (init R : Set X) (sat : X → Prop) : Prop :=
  ∀ x₀ ∈ init, x₀ ∈ R → sat x₀
\end{leansrc}
\ruleblock{Correspondence (Steele display $\leftrightarrow$ Lean)}
\badge{BadgeGreen}{reviewer-certified against the displayed texts}
\sourcenote{The dictionary below is a reviewer-certified statement of how the mechanized statements above interpret the displayed rule: display metavariable $\mapsto$ Lean binder, each display hypothesis $\mapsto$ the premise or side condition carrying it, the display conclusion $\mapsto$ the goal --- with any deliberate difference in width between display and mechanization named outright. Certification is recorded as a hash triple over (the displayed cell, the extracted statement texts, this dictionary\textquotesingle s own text); if any of the three changes, this block turns amber until a reviewer re-certifies. The correspondence claim itself remains editorial --- reviewer-grade, not machine-checked; its machine-checked successor types each soundness statement against the rule as a structured object, and supersedes this block per rule as it lands.}
The display maps onto \texttt{sound\_QGuardIntro} verbatim: the child
$\mathsf{QRegionSat}(\varphi, R)$ $\mapsto$ \texttt{h : ∀ x ∈ R, sat
x}; the conclusion $\mathcal{M} \vDash \mathsf{Guard}(R, s)$ $\mapsto$
\texttt{GuardSat init R sat} (pointwise: every admitted initial point
inside $R$ meets the body\textquotesingle s demand), generic in the
initial set \texttt{init} --- the rule is sound for \emph{any}
initial data, which is what ``no membership premise\textquotedbl{}
means formally. The component tie ($\mathsf{SInit}(\varphi) \in
s.\forall$) has deliberately no Lean counterpart: it is the load-time
refusal selecting \emph{which} predicate \texttt{sat} the lemma is
instantiated at. Width note: \texttt{sat : X → Prop} is generic, so
probabilistic bodies read \texttt{sat x} as the concluded bound at
start \texttt{x} --- wider than any single displayed instantiation,
safe direction. The specializations bind beside it: sublevel-atom
guards over exact-$\Q$ polynomial evaluation
(\texttt{sound\_QGuardIntro\_sublevel}, \texttt{guard\_eval\_sublevel})
and the probability-threshold body form
(\texttt{sound\_QGuardIntro\_prob}).

\ruleblock{Lean verifier receipts}
\sourcenote{module \texttt{QtslAux.QGuard}; bound build \texttt{3eb84fb9b559\ldots} (bound from receipt bd99b3a9109c4daf, 2026-07-28T18:45:47Z). A receipt records that named checking layers accepted a build: (module, compiled-artifact SHA-256, environment SHA-256, checker layers), minted when the checking gate ran; the store receipts attesting the bound build are listed with the layers each run credited. Nothing here asks the reader to trust the store: every layer is an independently runnable program over the quoted sources --- rebuild the module and re-run the checkers to reproduce the verdict. Receipts attesting superseded builds credit no layer and are omitted from this rendering (the count of omissions, if any, is noted below).}
\begin{receipttable}
\receiptrow{bd99b3a9109c4daf}{2026-07-28T18:45:47Z}{build+decorrelated}{\badge{BadgeGreen}{current}}
\end{receipttable}
\layerline{\layeritem{build}{lake build (elaboration)}{\badge{BadgeGreen}{current}}\layeritem{replay}{leanchecker (official kernel replay)}{\badge{BadgeGray}{none}}\layeritem{hardened}{safe\_verify (hardened screen)}{\badge{BadgeGray}{none}}\layeritem{decorrelated}{lean4export + nanoda (decorrelated kernel)}{\badge{BadgeGreen}{current}}}
\sourcenote{\badge{BadgeGray}{comparator: none} no fixture in the three-implementation comparator harness (compiled / kernel-interpreted / reference-interpreted) exercises this rule.}
\ruleblock{HOL4 statement}
\absentbox{NO HOL STATEMENT YET}{No HOL development of the extension theory exists (see the
region-to-satisfaction bridge\textquotesingle s section); the guarded
former\textquotesingle s realization is extension-theory record
structure, so a future HOL port of an extension-theory checker would
state it.
}
\ruleblock{HOL verifier receipts}
\absentbox{NONE}{no HOL verification receipt exists for this rule.}
\ruleblock{Proof}

\sourcenote{Genre: core rule (tier 1), structural group. The proof is
deliberately one step --- the rule exists to \emph{relocate} a
hypothesis, not to add content --- and the substantive work is the
semantics adjudication in the second remark, certified by the
strictness witness and the wrong-region countermodel bound at this
section.}

\textbf{Claim (\texttt{sound\_QGuardIntro}).}\enspace Let $R$ be the
warrant's region, $s$ the specification whose warranted safety
component is $\mathsf{SInit}(\varphi)$, and suppose the region
warrant: every state of $R$ satisfies $\varphi$. Then for \emph{any}
initial set $X_0$, the guarded claim holds: every admitted initial
point inside $R$ meets the body's demand.

\begin{proof}
Let $x \in X_0 \cap R$. The warrant instantiates at $x$: $x \in R$,
so $x \vDash \varphi$, which is the body's demand at $x$. Since $x$
was an arbitrary admitted initial point inside the guard region, the
pointwise guarded denotation holds. No hypothesis on $X_0$ was used
--- the rule is sound for arbitrary initial data, which is exactly
what licenses dropping the membership premise.
\end{proof}

\begin{bremark}[why pointwise, not blanket]
The blanket conditional --- \emph{if all of $X_0$ lies in $R$, then
the body holds from all of $X_0$} --- is strictly weaker: on an
initial set that merely straddles $R$, the blanket antecedent is
false and the blanket conditional vacuously true, while the pointwise
reading still demands the body from the inside points. The
strictness is mechanized (\texttt{exists\_blanket\_not\_guardSat});
the pointwise reading is the adopted semantics, and it is the one the
introduction rule is sound for.
\end{bremark}

\begin{bremark}[wrong region refuted]
Substituting any region other than the warrant's own breaks
soundness, by countermodel
(\texttt{exists\_not\_guardSat\_wrong\_region}); the load-time tie
(the stored guard region is derived from the warrant row,
store-and-check) is what makes the substitution unrepresentable
rather than detected.
\end{bremark}


\rulesection{39}{QGuardElim}{core rule}
\label{blcatrule-QGuardElim}
\ruleblock{The rule (Steele metanotation)}
\sourcenote{as realized in the mechanized development; source row: \texttt{\texttt{QTSL-rules.toml}} --- \emph{QGuardElim}. The metanotation statement is typeset by hand for notational quality; its mechanized ground truth is the Lean statement below, extracted from the sources at every build of this book and hash-bound. The rule \emph{roster} is machine-checked (every mechanized rule tag round-trips to exactly one catalogue row); the statement \emph{text} is editorial, checked by eye against the extracted Lean beside it.}
\begin{mathpar}
\inferrule*{
  \mathcal{M} \vDash \mathsf{Guard}(R,\ s) \\
  \mathrm{eval} : \vdash X_0 \subseteq R
}{
  \mathcal{M} \vDash s
}
\end{mathpar}

\par\medskip
\rulename{QGuardElim} is the elimination: a guarded claim plus initial
containment yields the body\textquotesingle s claim. The membership
premise is the same ground computation the unguarded route uses ---
per region-constructor pair, a degenerate box against a sublevel atom
is the exact evaluation leaf, a fat box against a sublevel atom is an
enclosure family, and box against box is interval containment ---
re-entering at exactly one place. Introduction followed by
elimination recovers the dissolved region-warranted route whole, so
the guarded former adds expressive structure without adding any new
soundness surface: everything the pair can conclude, the old route
could conclude given the same membership evidence.

What would break without the containment premise is the consumption
fence\textquotesingle s whole reason: a guarded claim consumed as if
unguarded is an unsound \emph{strengthening} --- the guarded claim
promises nothing off-region, so treating it as the body\textquotesingle
s claim smuggles in the very hypothesis the former internalized.

\cataloguenote{Realized \emph{at consumption}, deliberately --- there is no
standalone elimination node. Guard-presence triggers the membership
demand on any consumer of the claim: the composition rule walks its
premise list and demands initial membership for every guarded premise
(the quoted consumer-fence axiom below), and the verdict-compile and
equivalence evaluators are guard-aware. A consumer that cannot
discharge the membership defers loudly rather than assuming it.
}
\ruleblock{Theory rows (geolog)}
\sourcenote{The rule\textquotesingle s certificate content as it is declared in the geolog theory the checker loads --- payload schema, side-condition demands, and axiom rows, quoted from the theory source at every build of this catalogue (a stale quotation cannot exist). Axiom modes: \texttt{chase} = the engine derives; \texttt{report} = a violation is a loud judgment; \texttt{refuse} = a violation gates admission. Side-condition relations are discharged by bounded evaluation at the engine\textquotesingle s ground seam, never by derivation.}
\par\medskip
The baked realization of elimination-at-consumption: list membership
over the composition rule\textquotesingle s premises, the
consumer-fence axiom demanding initial membership for guarded
premises, and the ground membership relation it demands.

\declhead{theory Qtsl}{qsat\_list\_member}{\badge{BadgeGreen}{quoted from the theory source at render}}
\begin{leansrc}
qsat_list_member : {list: QSatList, sat: QSat} -> Prop;
\end{leansrc}
\declhead{theory Qtsl}{ax/wf/qwirecompose\_guard\_elim}{\badge{BadgeGreen}{quoted from the theory source at render}}
\begin{leansrc}
report ax/wf/qwirecompose_guard_elim : forall n : QWireCompose,
    l : QSatList, q : QSat, y : System, sp : Spec,
    g : QRegionAtom.
  n QWireCompose/components = l,
  {list: l, sat: q} qsat_list_member,
  q QSat/sys = y,
  q QSat/spec = sp,
  {spec: sp, region: g} spec_guard
  |- {region: g, sys: y} qtsl_init_member;
\end{leansrc}
\declhead{theory Qtsl}{qtsl\_init\_member}{\badge{BadgeGreen}{quoted from the theory source at render}}
\begin{leansrc}
qtsl_init_member : {region: QRegionAtom, sys: System} -> Prop;
\end{leansrc}
\ruleblock{Lean statement}
\declhead{QtslAux.QGuard}{sound\_QGuardElim}{\badge{BadgeGreen}{statement bound}}
\begin{leansrc}
/-- Soundness of the elimination rule: the initial-membership evaluation
re-enters exactly here — a guarded claim together with containment of the
initial set in the guard region yields the unguarded claim. -/
theorem sound_QGuardElim {init R : Set X} {sat : X → Prop}
    (hg : GuardSat init R sat) (heval : init ⊆ R) : ∀ x₀ ∈ init, sat x₀
\end{leansrc}
\declhead{QtslAux.QGuard}{qsatFromRegion\_derived}{\badge{BadgeGreen}{statement bound}}
\begin{leansrc}
/-- Coherence: introduction followed by elimination recovers the
region-warranted record route (warrant plus initial containment yields
satisfaction from every initial point). -/
theorem qsatFromRegion_derived {init R : Set X} {sat : X → Prop}
    (hchild : ∀ x ∈ R, sat x) (heval : init ⊆ R) : ∀ x₀ ∈ init, sat x₀
\end{leansrc}
\declhead{QtslAux.QGuard}{exists\_guardSat\_not\_unguarded}{\badge{BadgeGreen}{statement bound}}
\begin{leansrc}
/-- The permanent twin of the dropped-membership mutant, stated directly:
elimination WITHOUT its containment premise is false — a guarded claim
does not yield the unguarded one. On `Bool` with the full initial set and
guard region `{true}`, the demand "the state is `true`" is guarded-satisfied
(vacuously off-region, directly on it) yet fails from the initial point
`false`. -/
theorem exists_guardSat_not_unguarded :
    ∃ (init R : Set Bool) (sat : Bool → Prop),
      GuardSat init R sat ∧ ¬ ∀ x₀ ∈ init, sat x₀
\end{leansrc}
\declhead{QtslAux.QGuard}{guardSat\_singleton\_iff}{\badge{BadgeGreen}{statement bound}}
\begin{leansrc}
/-- At a singleton initial set the guarded denotation is the pointwise
conditional — the shape `Qtsl.Aliases.sound_QSatFromRegion` consumes. -/
theorem guardSat_singleton_iff {R : Set X} {sat : X → Prop} {x₀ : X} :
    GuardSat {x₀} R sat ↔ (x₀ ∈ R → sat x₀)
\end{leansrc}
\ruleblock{Correspondence (Steele display $\leftrightarrow$ Lean)}
\badge{BadgeGreen}{reviewer-certified against the displayed texts}
\sourcenote{The dictionary below is a reviewer-certified statement of how the mechanized statements above interpret the displayed rule: display metavariable $\mapsto$ Lean binder, each display hypothesis $\mapsto$ the premise or side condition carrying it, the display conclusion $\mapsto$ the goal --- with any deliberate difference in width between display and mechanization named outright. Certification is recorded as a hash triple over (the displayed cell, the extracted statement texts, this dictionary\textquotesingle s own text); if any of the three changes, this block turns amber until a reviewer re-certifies. The correspondence claim itself remains editorial --- reviewer-grade, not machine-checked; its machine-checked successor types each soundness statement against the rule as a structured object, and supersedes this block per rule as it lands.}
The display maps onto \texttt{sound\_QGuardElim} verbatim: the guarded
premise $\mapsto$ \texttt{hg : GuardSat init R sat}; the containment
premise $\vdash X_0 \subseteq R$ $\mapsto$ \texttt{hsub : init ⊆ R};
the conclusion $\mathcal{M} \vDash s$ $\mapsto$ \texttt{∀ x ∈ init,
sat x} (the unguarded pointwise demand). The dropped-membership
variant is refuted by countermodel
(\texttt{exists\_guardSat\_not\_unguarded}); the coherence statement
\texttt{qsatFromRegion\_derived} is the factoring of the dissolved
region-warranted route through the pair.

\ruleblock{Lean verifier receipts}
\sourcenote{module \texttt{QtslAux.QGuard}; bound build \texttt{3eb84fb9b559\ldots} (bound from receipt bd99b3a9109c4daf, 2026-07-28T18:45:47Z). A receipt records that named checking layers accepted a build: (module, compiled-artifact SHA-256, environment SHA-256, checker layers), minted when the checking gate ran; the store receipts attesting the bound build are listed with the layers each run credited. Nothing here asks the reader to trust the store: every layer is an independently runnable program over the quoted sources --- rebuild the module and re-run the checkers to reproduce the verdict. Receipts attesting superseded builds credit no layer and are omitted from this rendering (the count of omissions, if any, is noted below).}
\begin{receipttable}
\receiptrow{bd99b3a9109c4daf}{2026-07-28T18:45:47Z}{build+decorrelated}{\badge{BadgeGreen}{current}}
\end{receipttable}
\layerline{\layeritem{build}{lake build (elaboration)}{\badge{BadgeGreen}{current}}\layeritem{replay}{leanchecker (official kernel replay)}{\badge{BadgeGray}{none}}\layeritem{hardened}{safe\_verify (hardened screen)}{\badge{BadgeGray}{none}}\layeritem{decorrelated}{lean4export + nanoda (decorrelated kernel)}{\badge{BadgeGreen}{current}}}
\sourcenote{\badge{BadgeGray}{comparator: none} no fixture in the three-implementation comparator harness (compiled / kernel-interpreted / reference-interpreted) exercises this rule.}
\ruleblock{HOL4 statement}
\absentbox{NO HOL STATEMENT YET}{As at the introduction rule: no HOL development of the extension
theory exists; the consumption fence is extension-theory walk
structure.
}
\ruleblock{HOL verifier receipts}
\absentbox{NONE}{no HOL verification receipt exists for this rule.}
\ruleblock{Proof}

\sourcenote{Genre: core rule (tier 1), structural group. One-step
proof; the load-bearing content is the consumption fence (a guarded
claim consumed as if unguarded is an unsound strengthening), certified
by the dropped-membership countermodel.}

\textbf{Claim (\texttt{sound\_QGuardElim}).}\enspace If the guarded
claim holds (every admitted initial point inside $R$ meets the body's
demand) and $X_0 \subseteq R$, then the body's claim holds from every
admitted initial point.

\begin{proof}
Let $x \in X_0$. By containment $x \in R$, so $x \in X_0 \cap R$, and
the guarded claim yields the body's demand at $x$.
\end{proof}

\begin{bremark}[the fence, and its countermodel]
Dropping the containment premise is refuted by countermodel
(\texttt{exists\_guardSat\_not\_unguarded}): a guarded claim promises
nothing off-region, so an initial point outside $R$ can violate the
body while the guarded claim holds. This is why the elimination is
realized \emph{at consumption} --- guard-presence triggers the
membership demand on any consumer --- rather than trusting producers
to remember it.
\end{bremark}

\begin{bremark}[coherence]
Introduction followed by elimination recovers the dissolved
region-warranted route whole
(\texttt{qsatFromRegion\_derived}): warrant in, guarded claim out,
membership in, body's claim out --- the same conclusion the old
single rule drew from the same two pieces of evidence. The former
adds a place to \emph{stand} between them, not a new way to conclude.
\end{bremark}


\rulesection{40}{QSafetyBattery}{core rule}
\label{blcatrule-QSafetyBattery}
\ruleblock{The rule (Steele metanotation)}
\sourcenote{as realized in the mechanized development; source row: \texttt{\texttt{QTSL-rules.toml}} --- \emph{QSafetyBattery}. The metanotation statement is typeset by hand for notational quality; its mechanized ground truth is the Lean statement below, extracted from the sources at every build of this book and hash-bound. The rule \emph{roster} is machine-checked (every mechanized rule tag round-trips to exactly one catalogue row); the statement \emph{text} is editorial, checked by eye against the extracted Lean beside it.}
\begin{mathpar}
\inferrule*{
  V : \mathsf{QCertFam} \\
  \bigl\{D_m \vdash \mathcal{A}_m V_m \le 0\bigr\}_{m} \ \text{(interiors)} \\
  \bigl\{\mathrm{slab}_j \vdash g_j \fatsemi V_{m'_j} \le V_{m_j}\bigr\}_{j} \ \text{(guards, per jump $j$: source mode $m_j$, target $m'_j$, reset $g_j$)} \\
  \bigl\{U \cap \bar{D}_m \vdash V_m \ge 1\bigr\}_{m} \ \text{(floors)} \\
  \bigl\{\bar{D}_m \vdash V_m \ge 0\bigr\}_{m} \ \text{(nonnegs)} \\
  \mathrm{scope\_warrant} : \mathsf{QInvScope} \ \text{(optional — the dissolved ScopeRule)} \\
  p + y \le 1 \ \text{(threshold tie)}
}{
  \mathcal{M},\ x_0 \in \{V \le y\} \vDash \mathsf{SProbGE}\bigl(p,\ \mathsf{PAlways}\ \varphi\bigr)
}
\end{mathpar}

\ruleflavor{a region written left of the turnstile scopes the row: $R \vdash P$ asserts $P$ at every point of $R$ (each such row is an inherited BL \rulename{BernCone}-family sequent). The certificate family $V$ enters as certificate data --- a shared reference with no claim content of its own; the rows impose every claim --- never as a derivation edge. The threshold tie $p + y \le 1$ is the soundness theorem\textquotesingle s own hypothesis, displayed as a side condition: unlike the deadline and until nodes (whose baked sorts carry a \texttt{tie} proof field), the baked safety node carries no tie field yet --- a named schema divergence, not a checked leaf. The conclusion is the theory row\textquotesingle s own judgment, in constructor vocabulary: the region-conditioned probability formula ($\mathsf{SProbGE}(p, \mathsf{PAlways}\ \varphi)$ on the region $\mathsf{qraSublevel}(V, y)$). The record\textquotesingle s licensed alias tie is $\mathsf{SProbGE}(p, \psi) = \mathsf{SReturnGE}(p, \mathsf{retPath}\ \psi)$ --- a return respelling, which does not touch the state-region antecedent; the initial-set hypothesis $x_0 \in \{V \le y\}$ left of the judgment is the satisfaction binding\textquotesingle s own guard (the theory demands the subject\textquotesingle s initial data lie in the concluded region), rendered in the same hypothesis convention as every \rulename{BernCone}-family sequent. The internalized variant LANDED at the grid re-sort revision: the guarded-spec former ($\mathsf{Guard}$, a certificate-stratum record slot carrying the initial-region antecedent) with its introduction rule \rulename{QGuardIntro} (this catalogue, the structural group). The hypothesis form remains this catalogue\textquotesingle s rendering convention (one convention catalogue-wide); the guarded former is its internalized spelling, one \rulename{QGuardIntro} application away, and no belief-gauge reading is displayed.}
\par\medskip
\rulename{QSafetyBattery} is the probabilistic safety battery: from a
family of purely polynomial premise rows about a certificate $V$, it
concludes that with probability at least $p$ the system satisfies
$\mathsf{PAlways}\ \varphi$ (surface syntax $\mathsf{G}\,\varphi$) ---
\emph{forever} --- from every start state in
the sublevel region $\{V \le y\}$. The rule concludes a
satisfaction judgment ($\vDash$), one theory-extension edge above the
base calculus, and every premise row is an inherited base-calculus
sequent: \emph{interior} rows say the certificate does not increase
along the in-mode dynamics, \emph{guard} rows say jumps do not increase
it, \emph{floor} rows say $V \ge 1$ on the avoid set (the compiled
over-approximation of where $\varphi$ fails), and \emph{nonneg} rows say
$V \ge 0$ on the domain hull. The battery adds zero kernel arithmetic;
its entire content is the supermartingale assembly one level up: $V$
read along the process and stopped at the first failure is a nonnegative
supermartingale, and Ville\textquotesingle s inequality converts its
initial value into a bound on the probability of ever failing.

A producer reaches for the rule wherever a quantitative safety guarantee
over a hybrid stochastic system is wanted. The certificate is found by
whatever synthesis method the producer likes (typically linear or
sum-of-squares programming over the same polynomial cones the premise
rows are checked in); only the rows ship in the certificate.

What the checker recomputes are the side conditions: the domain-hull
condition (one step from any in-scope state stays inside the hull
$\bar{D}$, which is exactly where the floor and nonneg rows were
checked), the fail-closed containment (one-step mass stays in
$D \cup U$: leaving the modelled domain is counted as \emph{failure},
never as silent success), and the floor-scope resolution (the floor
rows\textquotesingle{} scopes resolve to an exact cover of the avoid set
within the hull). The threshold tie $p + y \le 1$ is displayed as a
side condition, not a checked leaf: it is the soundness
theorem\textquotesingle s own hypothesis, and the baked safety node
carries no tie proof field yet (contrast the deadline and until nodes)
--- a named schema divergence, as the rule\textquotesingle s flavor
text states.

Each hypothesis is load-bearing. Drop the floors and failure is free of
charge: a certificate identically $0$ satisfies every other row and
would certify almost-sure safety of a chain that fails at its first
step. Drop the nonnegs and the maximal inequality itself fails --- a
signed process can pay for a barrier visit with mass borrowed from
below. Drop the fail-closed containment and probability mass could exit
the modelled domain unaccounted. Drop the threshold tie and the
conclusion would assert success probability $p$ alongside a failure
budget $y$ with $p + y > 1$. The proof section
walks each of these, with counterexamples.

The dashboard row\textquotesingle s mark states the status honestly:
the hash-bound theorems warrant the rule at the
\emph{discrete-kernel profile} --- system semantics a one-step Markov
kernel on the flattened hybrid state space, interior rows read as
one-step expected decrease (the Foster--Lyapunov reading) --- with the
PDMP sub-profile of the continuous-time semantics additionally
machine-checked end-to-end. The general continuous-time profile
(jump--diffusion) is \emph{proven}, at the martingale-problem
presentation: the system enters through its generator drift data,
$\mathcal{A}V$ as data with the martingale-problem hypothesis for the
pair, and the battery\textquotesingle s conclusion follows end-to-end
(\texttt{sound\_GSHS\_safety\_mp},
\texttt{sound\_GSHS\_safety\_battery\_mp}, module
\texttt{QtslAux.LayerD.MPDischarge} of the semantic-universe
strand\textquotesingle s auxiliary package) --- the same epistemic
shape as the staged-row theorem\textquotesingle s (\texttt{sound\_QSafetyBattery\_staged}) mixture bridge: an
explicit presentation hypothesis, quantifying over exactly the
systems whose presentation satisfies it. The martingale-problem
theorems are hash-bound below like every other block --- their
package, the semantic-universe strand\textquotesingle s auxiliary
package in the public tree, is registered with this
catalogue\textquotesingle s binder --- so the dashboard row carries no
conditional mark: each named profile is proven, and its theorems are
bound.

At the general stratum the battery\textquotesingle s shape is also
stated over the semantic universe itself, in the semantic-universe
functor package (hash-bound below): trajectories live as closed record
sets under the Fell topology, and the undisplayed regularity pack ---
the measurability the profile proofs \emph{assume} of their region
events --- restates there as a theorem of the site\textquotesingle s
topology: the confinement event of a closed safe region is Borel over
every trajectory window, with no compactness and no
$\sigma$-compactness anywhere
(\texttt{measurableSet\_confinedEvent\_cylinder}). On that carrier the
safety conclusion is a \emph{wall} --- a $\forall$-face spec:
per-window confinement classes are closed under record restriction,
their persistence follows from window-calculus transitivity alone, and
the judgment bridge is an equivalence --- satisfaction \emph{equals}
almost-sure confinement at the window
(\texttt{satAtM\_safetySpecF\_iff}), with the indicator calculus
admitted exactly by the measurability above (\texttt{pull\_indPenM}).
Two further site-level facts sharpen the story. Record confinement of
a trajectory\textquotesingle s closed-up graph to a closed wall is
\emph{pointwise} safety at every stage, unconditional in the
trajectory (\texttt{graphCl\_mem\_confinedEvent\_iff}) --- the
segment-constancy side condition a skeleton-based route would carry is
absent by construction. And composing two wall demands along a window
compositor yields exactly the wall demand of the glued record class
(\texttt{mboxplusWF\_indPenM}): wall specs are closed under
composition, and persistence survives it. These statements bind the
site-general stratum --- the acceptance events and spec classes that
every profile\textquotesingle s conclusion instantiates; the displayed
judgment\textquotesingle s warrant remains with the profile theorems
above.

\cataloguenote{Fail-closed as a design reading: the underlying theorem\textquotesingle s
concluded event is \emph{stronger} than the safety
formula\textquotesingle s truth set --- never fail \emph{and} never
leave the modelled domain --- and the formula tie then weakens it. The
node presents its premises per mode and per jump (mode-indexed
certificate families); the bound theorem for that shape assembles the
flattened generator row through a sub-convex mixture decomposition of the
one-step kernel --- the mode\textquotesingle s flow step plus the
enabled guard jumps --- which is the presentation-to-semantics bridge,
carried as an explicit hypothesis and quoted in the
proof\textquotesingle s final remark.
}
\ruleblock{Theory rows (geolog)}
\sourcenote{The rule\textquotesingle s certificate content as it is declared in the geolog theory the checker loads --- payload schema, side-condition demands, and axiom rows, quoted from the theory source at every build of this catalogue (a stale quotation cannot exist). Axiom modes: \texttt{chase} = the engine derives; \texttt{report} = a violation is a loud judgment; \texttt{refuse} = a violation gates admission. Side-condition relations are discharged by bounded evaluation at the engine\textquotesingle s ground seam, never by derivation.}
\par\medskip
In the geolog theory the checker loads, every object this rule touches
is a declared row, and the display\textquotesingle s remaining ambient
notation resolves as follows. The concluded formula
$\mathsf{SProbGE}(p, \mathsf{PAlways}\ \varphi)$ (surface syntax
$\mathsf{P}_{\ge p}[\mathsf{G}\,\varphi]$) \emph{is} the safety-face
cap former \texttt{SProbGE} (concluding into the record\textquotesingle s
$\forall$-component sort) applied at the path former \texttt{PAlways}
--- the probability-floor former whose threshold canonicity is the
range axiom quoted below. The displayed $\mathcal{M}$ is the paired
subject: a \texttt{System} element, bound to the claim by the
\texttt{QSat} pairing (the specification consumes this formula among
its safety components; the warrant row names it). The generator reads $\mathcal{A}_m V_m$ are not
free-floating operator applications: \texttt{qtsl\_gen\_form} is the
derived generator-form seam --- given the subject system and a
certificate component (a hypernet from the state space to the
nonnegative extended reals), the row names the hypernet saying how the
certificate can change with duration, recomputed from the
subject\textquotesingle s content rows and never trusted as stored
data. The battery\textquotesingle s E2 payload is the
\texttt{QSafetyBattery} entity below, field for field the
display\textquotesingle s premise kit; the conclusion region
$\mathsf{qraSublevel}(V, y)$ is the region-atom constructor
of the same name.

\declhead{theory Qtsl}{SProbGE}{\badge{BadgeGreen}{quoted from the theory source at render}}
\begin{leansrc}
SProbGE      : {p: Q, path: PathPred} -> ForallSpec;
\end{leansrc}
\declhead{theory Qtsl}{ax/wf/prob\_range}{\badge{BadgeGreen}{quoted from the theory source at render}}
\begin{leansrc}
report ax/wf/prob_range : forall p : Q, ps : PathPred, f : ForallSpec.
  f = {p: p, path: ps} SProbGE
  |- {l: 0, r: p} q_lt, {l: p, r: 1} q_le;
\end{leansrc}
\declhead{theory Qtsl}{PAlways}{\badge{BadgeGreen}{quoted from the theory source at render}}
\begin{leansrc}
PAlways : StatePred -> PathPred;                          // G phi
\end{leansrc}
\declhead{theory Qtsl}{QSafetyBattery}{\badge{BadgeGreen}{quoted from the theory source at render}}
\begin{leansrc}
QSafetyBattery : Set;
\end{leansrc}
\declhead{theory Qtsl}{QSafetyBattery/*}{\badge{BadgeGreen}{quoted from the theory source at render}}
\begin{leansrc}
QSafetyBattery/proof : QSafetyBattery -> QProof;
QSafetyBattery/p : QSafetyBattery -> Q;
QSafetyBattery/arg : QSafetyBattery -> StatePred;
QSafetyBattery/cert : QSafetyBattery -> QCertFam;
QSafetyBattery/interiors : QSafetyBattery -> QProofList;
QSafetyBattery/guards : QSafetyBattery -> QProofList;
QSafetyBattery/lowbounds : QSafetyBattery -> QProofList;
QSafetyBattery/nonnegs : QSafetyBattery -> QProofList;
QSafetyBattery/avoid : QSafetyBattery -> QRegionAtom;
QSafetyBattery/scope_warrant : QSafetyBattery -> QInvScope;
\end{leansrc}
\declhead{theory Qtsl}{qraSublevel}{\badge{BadgeGreen}{quoted from the theory source at render}}
\begin{leansrc}
qraSublevel : {cert: QCertFam, thr: Q} -> QRegionAtom;
\end{leansrc}
\declhead{theory Qtsl}{qtsl\_gen\_form}{\badge{BadgeGreen}{quoted from the theory source at render}}
\begin{leansrc}
qtsl_gen_form : {sys: System, cert: HypernetRef, gen: HypernetRef}
  -> Prop;   // side-condition (oracle discharge / checker
\end{leansrc}
\ruleblock{Lean statement}
\declhead{BernsteinLogicAux.Battery}{sound\_QSafetyBattery}{\badge{BadgeGreen}{statement bound}}
\begin{leansrc}
/-- **`sound_QSafetyBattery`, discrete profile** (QTSL-rules.toml row `QSafetyBattery`
at the discrete-kernel profile; the concluded judgment reads as region
under-approximation — every start in the region satisfies the formula). Over the
fail-closed form, the formula tie (F) — the avoid region over-approximates the falsity
set of the `𝖦`-argument `φ`, compiled as `U := cmpl(φ)` — weakens the concluded event to
the truth of `𝖦 φ` along the path:
`Pr_{x}[∀ t, X t ⊨ φ] ≥ p` for every start `x` in the concluded region
`D ∩ {V ≤ thr}`. No measurability is asked of the satisfaction predicate: the `𝖦 φ`
event is bounded below through `measure_mono` on the (measurable) fail-closed event. -/
theorem sound_QSafetyBattery {sat : S → Prop}
    (hD : MeasurableSet D) (hDbar : MeasurableSet Dbar) (hU : MeasurableSet U)
    (hV : Measurable V)
    (hF : ∀ y, ¬ sat y → y ∈ U)
    (hΘ : p + thr ≤ 1)
    (hDsub : D ⊆ Dbar)
    (hR3 : ∀ y ∈ D, y ∉ U → K y Dbar = 1)
    (hR4 : ∀ y ∈ D, y ∉ U → K y (D ∪ U) = 1)
    (hB12int : ∀ y ∈ D, y ∉ U → Integrable V (K y))
    (hB12 : ∀ y ∈ D, y ∉ U → ∫ z, V z ∂K y ≤ V y)
    (hB3 : ∀ y ∈ U, y ∈ Dbar → 1 ≤ V y)
    (hB4 : ∀ y ∈ Dbar, 0 ≤ V y)
    (hx : x ∈ D) (hthr : V x ≤ thr) :
    ENNReal.ofReal p ≤ pathMeasure K x {ω | ∀ n, sat (ω n)}
\end{leansrc}
\declhead{BernsteinLogicAux.Battery}{sound\_QSafetyBattery\_failClosed}{\badge{BadgeGreen}{statement bound}}
\begin{leansrc}
/-- **`sound_QSafetyBattery`, discrete profile, fail-closed form** (the discrete-kernel
profile: the hybrid model flattened to one measurable state space with a single one-step
Markov kernel `K`). Hypotheses
are the checked rows (B1)∧(B2) as the kernel drift row, (B3) floors, (B4) nonnegs, the
side conditions (R3) (hull: `D ⊆ D̄` and one-step mass in `D̄`) and (R4) (fail-closed:
one-step mass in `D ∪ U`), the threshold arithmetic (Θ) `p + thr ≤ 1` (the staged
`q_add`/`q_le` vocabulary, in exact ℚ), and a start state in the concluded region
`D ∩ {V ≤ thr}` (`qraSublevel`). Conclusion: with probability at least `p` the chain
remains in `D \ U` **forever** — never reaching the avoid region and never leaving the
modeled domain. This is the strong (fail-closed) reading of the concluded `QRegionSat`;
`sound_QSafetyBattery` weakens it through the formula tie (F).

The `Integrable V (K y)` hypothesis is the discrete reading of
"domain membership of `V` in the extended generator's domain"; for the intended
instances (`V` polynomial, one-step images in the compact
`D̄` by (R3)) it is automatic. -/
theorem sound_QSafetyBattery_failClosed
    (hD : MeasurableSet D) (hDbar : MeasurableSet Dbar) (hU : MeasurableSet U)
    (hV : Measurable V)
    (hΘ : p + thr ≤ 1)
    (hDsub : D ⊆ Dbar)
    (hR3 : ∀ y ∈ D, y ∉ U → K y Dbar = 1)
    (hR4 : ∀ y ∈ D, y ∉ U → K y (D ∪ U) = 1)
    (hB12int : ∀ y ∈ D, y ∉ U → Integrable V (K y))
    (hB12 : ∀ y ∈ D, y ∉ U → ∫ z, V z ∂K y ≤ V y)
    (hB3 : ∀ y ∈ U, y ∈ Dbar → 1 ≤ V y)
    (hB4 : ∀ y ∈ Dbar, 0 ≤ V y)
    (hx : x ∈ D) (hthr : V x ≤ thr) :
    ENNReal.ofReal p ≤ pathMeasure K x {ω | ∀ n, ω n ∈ D \ U}
\end{leansrc}
\declhead{BernsteinLogicAux.LayerD.Skeleton}{sound\_GSHS\_safety\_battery}{\badge{BadgeGreen}{statement bound}}
\begin{leansrc}
/-- **`sound_QSafetyBattery`, GSHS instance** (T5; lifetime-relative in the
value-presentation reading — and unconditional for W-Z-free models,
those a.s. free of the two stall phenomena, Zeno accumulation and
vanishing into a never-terminating segment): with the threshold arithmetic `p + thr ≤ 1` and a
start whose certificate is under the threshold, the chain remains outside
the avoid region — at every alive entry *and* every halt value — with
probability at least `p`. -/
theorem sound_GSHS_safety_battery [IsMarkovKernel M.seg] [IsMarkovKernel M.jump]
    {V : State → ℝ} {U : Set State} {p thr : ℝ}
    (hV : Measurable V) (hU : MeasurableSet U)
    (hSegD : SegmentDynkin M 0 V) (hForce : ForcedNonexpansion M V)
    (hNonneg : NonnegV V) {C : ℝ} (hC0 : 0 ≤ C) (hVC : ∀ y, V y ≤ C)
    (hFloor : ∀ y ∈ U, 1 ≤ V y)
    (hΘ : p + thr ≤ 1) (x : State) (hthr : V x ≤ thr) :
    ENNReal.ofReal p ≤ Battery.pathMeasure (skel M) (SkelState.alive x)
      {ω | ∀ n, ω n ∉ skelBad U}
\end{leansrc}
\par\medskip
\textbf{Discharging the segment field, corner-immune route.}\enspace
The per-model segment-Dynkin field the statement above consumes is
itself a theorem for piecewise-deterministic segments: the statement
below derives it from upper-Dini rate rows alone --- no
differentiability of the certificate along the flow is assumed
anywhere, so corner points, absolutely-continuous solutions, and
Filippov-style trajectories are in scope exactly as the
module\textquotesingle s header fences state.

\declhead{BernsteinLogicAux.LayerD.DiniDischarge}{segmentDynkin\_of\_pdmp\_dini}{\badge{BadgeGreen}{statement bound}}
\begin{leansrc}
/-- **[SegD] discharged at Dini strength** for models whose entry states
carry terminating Dini flow data. Mixing with the landed smooth data is
free — the field is a conjunction over states, so discharge is pointwise;
this theorem states the all-Dini case, and mixed cases compose by
`rcases` at the call site exactly as `segmentDynkin_of_pdmp` does. -/
theorem segmentDynkin_of_pdmp_dini (M : SegmentModel State) (θ : ℝ)
    (V : State → ℝ) [IsSFiniteKernel M.jump] (hV : Measurable V)
    (h : ∀ x : State, Nonempty (PDMPFlowFinDini M θ V x)) :
    SegmentDynkin M θ V
\end{leansrc}
\declhead{BernsteinLogicAux.LayerD.DiniDischarge}{pdmp\_dynkin\_core\_dini}{\badge{BadgeGreen}{statement bound}}
\begin{leansrc}
/-- **The corner-immune [SegD] core**: the B1 row at upper-right-Dini
strength — stated in the eventual/liminf-slope form the mean-value fence
consumes, so it survives switch corners, Zeno accumulation of corners,
and (through the AC/Filippov reading) sliding — still bounds the
θ-discounted, survival-weighted terminal read plus the accumulated jump
read by the entry value. No derivative of `v` is hypothesized anywhere. -/
theorem pdmp_dynkin_core_dini {v r j : ℝ → ℝ} {θ T : ℝ} (hT : 0 ≤ T)
    (hr : Continuous r) (hj : Continuous j)
    (hv_cont : ContinuousOn v (Icc 0 T))
    (hrow : ∀ t ∈ Ico 0 T, ∀ ρ : ℝ, (r t - θ) * v t - j t < ρ →
      ∃ᶠ z in 𝓝[>] t, slope v t z < ρ) :
    jint r j θ T + survWeight r θ T * v T ≤ v 0
\end{leansrc}
\par\medskip
\textbf{Discharging the segment field, diffusion route.}\enspace For
jump-diffusion segments the same field follows from a single quarantined
per-model input --- the sampled-supermartingale inequality --- stated
one stratum below as an interface structure; the theorem below converts
any inhabitant into the segment-Dynkin field, with a deterministic
positive control included in the module.

\declhead{BernsteinLogicAux.LayerD.ItoInterface}{segmentDynkin\_of\_ito}{\badge{BadgeGreen}{statement bound}}
\begin{leansrc}
/-- **[SegD] discharged for the diffusion profile against the stated
interface** (assault doc §3 row 1(a)): a model all of whose entry states
carry Itô-interface data satisfies the `SegmentDynkin` hypothesis field.
Mixing with PDMP/Dini data is free at call sites, exactly as for
`segmentDynkin_of_pdmp_dini` — the field is a conjunction over states,
discharged pointwise. -/
theorem segmentDynkin_of_ito (M : SegmentModel State) (θ : ℝ)
    (V : State → ℝ) [IsSFiniteKernel M.jump] (hV : Measurable V)
    (h : ∀ x : State, Nonempty (ItoInterface M θ V x)) :
    SegmentDynkin M θ V
\end{leansrc}
\par\medskip
\textbf{The interface\textquotesingle s own discharge.}\enspace The
quarantined inequality is in turn delivered by optional sampling: from a
supermartingale property of the certificate process (plus the
d\'ebut-measurability family), the assembly below inhabits the full
interface --- leaving, per diffusion entry state, exactly the
supermartingale property and the d\'ebut family. The blocks that follow
discharge each in turn.

\declhead{BernsteinLogicAux.LayerD.ItoSampling}{ItoInterface.ofSupermartingale}{\badge{BadgeGreen}{statement bound}}
\begin{leansrc}
/-- **The interface, discharged down to the supermartingale property**:
a full `ItoInterface` instance from the non-analytic pathwise data plus
`Supermartingale` + the début family. Composing with
`segmentDynkin_of_ito`, the chain
`Supermartingale + début family ⟹ [SegD]` is Lean end-to-end; what a
concrete diffusion model still owes is exactly the supermartingale
property (theorem 6's shape / upstream Itô). -/
noncomputable def ItoInterface.ofSupermartingale
    (M : SegmentModel State) [IsProbabilityMeasure P]
    (hT : ∀ ω, 0 ≤ T ω)
    (hk : ∀ ω, k ω = ExitKind.guard ∨ k ω = ExitKind.halt
      ∨ k ω = ExitKind.live)
    (hγ : ∀ ω, Measurable (γ ω))
    (hvc : ∀ ω, Continuous fun u => V (γ ω u))
    (hVγ : ∀ ω u, 0 ≤ V (γ ω u))
    (hr : ∀ ω, Continuous (r ω)) (hr0 : ∀ ω u, 0 ≤ r ω u)
    (hjr : ∀ ω, Continuous (jr ω)) (hjr0 : ∀ ω u, 0 ≤ jr ω u)
    (hjr_couple : ∀ ω, ∀ u ∈ segDom (T ω) (k ω),
      ENNReal.ofReal (jr ω u)
        = ENNReal.ofReal (r ω u)
          * ∫⁻ z, ENNReal.ofReal (V z) ∂(M.jump (γ ω u, ExitKind.alarm)))
    (hread : ∫⁻ o, outcomeRead M θ V o ∂(M.seg x)
      ≤ ∫⁻ ω, ∫⁻ o, outcomeRead M θ V o
          ∂(segMeasureIto (γ ω) (r ω) (T ω) (k ω)) ∂P)
    (hγ0 : ∀ ω, γ ω 0 = x)
    (hZ : Supermartingale (certProcess θ V γ r jr) ℱ P)
    (hτ : ∀ s : ℝ≥0, MeasurableSet[ℱ s]
      {ω | exitDebut T k ω ≤ (s : ℝ≥0∞)}) :
    ItoInterface M θ V x
\end{leansrc}
\declhead{BernsteinLogicAux.LayerD.ItoSampling}{sampled\_le\_of\_supermartingale}{\badge{BadgeGreen}{statement bound}}
\begin{leansrc}
/-- **Partial discharge of (A.1)** from the banked optional-sampling
suite: the supermartingale property of the certificate process (the
[WR-2]/[WR-3] residue) plus the exit début's measurability family imply
the `hsampled` field, with the start pinned at `x` and `P` a probability
measure supplying `Y₀ = V x`. -/
theorem sampled_le_of_supermartingale [IsProbabilityMeasure P]
    (hT : ∀ ω, 0 ≤ T ω)
    (hvc : ∀ ω, Continuous fun u => V (γ ω u))
    (hVγ : ∀ ω u, 0 ≤ V (γ ω u))
    (hr : ∀ ω, Continuous (r ω)) (hjr : ∀ ω, Continuous (jr ω))
    (hjr0 : ∀ ω u, 0 ≤ jr ω u)
    (hγ0 : ∀ ω, γ ω 0 = x)
    (hZ : Supermartingale (certProcess θ V γ r jr) ℱ P)
    (hτ : ∀ s : ℝ≥0, MeasurableSet[ℱ s]
      {ω | exitDebut T k ω ≤ (s : ℝ≥0∞)}) :
    ∀ t : ℝ, 0 ≤ t →
      ∫⁻ ω, sampledRead θ V (γ ω) (r ω) (jr ω) (T ω) (k ω) t ∂P
        ≤ ENNReal.ofReal (V x)
\end{leansrc}
\par\medskip
\textbf{The residue, bundled.}\enspace The per-model inputs the two
routes leave open are packaged as a structure whose only analytic
fields are the certificate-process supermartingale property and the
exit d\'ebut\textquotesingle s measurability family; the theorems below
convert per-entry-state families of such bundles (all-diffusion, or a
per-state diffusion/Dini mixture) into the segment field, and the class
constructor builds the battery\textquotesingle s Ville class with that
field proven rather than assumed.

\declhead{BernsteinLogicAux.GshsItoDischarge}{ItoSupermartingaleData}{\badge{BadgeGreen}{statement bound}}
\begin{leansrc}
/-- **The supermartingale-floor bundle**: exactly
`ItoInterface.ofSupermartingale`'s inputs as a structure, so that
per-entry-state families can be quantified. The two analytic fields are
`hZ` and `hτ`; everything else is the pathwise bookkeeping the interface
already carried. -/
structure ItoSupermartingaleData (M : SegmentModel State) (θ : ℝ)
    (V : State → ℝ) (x : State) where
  /-- the driving path-data space -/
  Ω : Type*
  /-- its σ-algebra -/
  [mΩ : MeasurableSpace Ω]
  /-- the law of the driving data issued from `x` -/
  P : Measure Ω
  /-- `P` is a probability measure (`Y₀ = V x` is the model's business —
  the optional-sampling route consumes this) -/
  hprob : IsProbabilityMeasure P
  /-- the filtration the certificate process rides -/
  ℱ : Filtration ℝ≥0 mΩ
  /-- exit time of each path (unread on live paths) -/
  T : Ω → ℝ
  hT : ∀ ω, 0 ≤ T ω
  /-- terminator species of each path -/
  k : Ω → ExitKind
  hk : ∀ ω, k ω = ExitKind.guard ∨ k ω = ExitKind.halt
    ∨ k ω = ExitKind.live
  /-- the trajectory of each path (clamped extension off the segment) -/
  γ : Ω → ℝ → State
  hγ : ∀ ω, Measurable (γ ω)
  hvc : ∀ ω, Continuous fun u => V (γ ω u)
  hVγ : ∀ ω u, 0 ≤ V (γ ω u)
  /-- paths start at the entry state -/
  hγ0 : ∀ ω, γ ω 0 = x
  /-- alarm intensity along each path -/
  r : Ω → ℝ → ℝ
  hr : ∀ ω, Continuous (r ω)
  hr0 : ∀ ω u, 0 ≤ r ω u
  /-- λ-weighted post-jump read along each path -/
  jr : Ω → ℝ → ℝ
  hjr : ∀ ω, Continuous (jr ω)
  hjr0 : ∀ ω u, 0 ≤ jr ω u
  hjr_couple : ∀ ω, ∀ u ∈ segDom (T ω) (k ω),
    ENNReal.ofReal (jr ω u)
      = ENNReal.ofReal (r ω u)
        * ∫⁻ z, ENNReal.ofReal (V z) ∂(M.jump (γ ω u, ExitKind.alarm))
  /-- the read-level disintegration of `M.seg x` along paths
  (definitional for `Rep`-constructed models) -/
  hread : ∫⁻ o, outcomeRead M θ V o ∂(M.seg x)
    ≤ ∫⁻ ω, ∫⁻ o, outcomeRead M θ V o
        ∂(segMeasureIto (γ ω) (r ω) (T ω) (k ω)) ∂P
  /-- **analytic field 1**: the certificate process is a supermartingale
  — the exact [WR-2]/[WR-3] residue -/
  hZ : Supermartingale (certProcess θ V γ r jr) ℱ P
  /-- **analytic field 2**: the exit début's measurability family -/
  hτ : ∀ s : ℝ≥0, MeasurableSet[ℱ s]
    {ω | exitDebut T k ω ≤ (s : ℝ≥0∞)}
\end{leansrc}
\declhead{BernsteinLogicAux.GshsItoDischarge}{segmentDynkin\_of\_supermartingale}{\badge{BadgeGreen}{statement bound}}
\begin{leansrc}
/-- **[SegD] from the supermartingale floor, all-diffusion form**: a
model all of whose entry states carry supermartingale-floor bundles
satisfies the `SegmentDynkin` hypothesis field. -/
theorem segmentDynkin_of_supermartingale (M : SegmentModel State) (θ : ℝ)
    (V : State → ℝ) [IsSFiniteKernel M.jump] (hV : Measurable V)
    (h : ∀ x : State, Nonempty (ItoSupermartingaleData M θ V x)) :
    SegmentDynkin M θ V
\end{leansrc}
\declhead{BernsteinLogicAux.GshsItoDischarge}{segmentDynkin\_of\_mixed}{\badge{BadgeGreen}{statement bound}}
\begin{leansrc}
/-- **The pointwise diffusion/Dini mixture, enacted** (both discharge
docstrings promised this "free at call sites"; here it is as one
theorem): each entry state carries either a supermartingale-floor bundle
(diffusion segments) or terminating Dini flow data (deterministic/PDMP
segments, switch corners included) — mixed freely, [SegD] follows. Two
disjuncts only: Dini is the deterministic floor (smooth data restates at
it, or lifts by `dini_row_of_deriv_row`). -/
theorem segmentDynkin_of_mixed (M : SegmentModel State) (θ : ℝ)
    (V : State → ℝ) [IsSFiniteKernel M.jump] (hV : Measurable V)
    (h : ∀ x : State, Nonempty (ItoSupermartingaleData M θ V x)
      ∨ Nonempty (PDMPFlowFinDini M θ V x)) :
    SegmentDynkin M θ V
\end{leansrc}
\declhead{BernsteinLogicAux.GshsItoDischarge}{GshsSegmentVille.ofDischarge}{\badge{BadgeGreen}{statement bound}}
\begin{leansrc}
/-- **The discharge constructor — `gshsGen:segment-dynkin`'s diffusion
sub-case retired to the floor**: the rung's class with its quarantined
field PROVEN from row-conditional per-state discharge data, never
assumed. `hdata`'s conditional shape mirrors the field's own: the
certified θ-discounted generator rows are where a concrete model's
analytic side takes its input (Itô's formula consumes `𝒜V ≤ 0` on the
cells to prove `hZ`; the Dini row data is the compiled rows themselves).
What a model owes per diffusion entry state is exactly `hZ` +`hτ` — the
`sdeGen:dynkin-supermartingale` route-(i) floor. -/
noncomputable def GshsSegmentVille.ofDischarge
    [MeasurableSpace Mode] [DecidableEq Mode]
    {P : GshsProfile n m T Mode} {V : Mode → State n → ℝ}
    {gradV : Mode → State n → State n}
    {hessV : Mode → State n → Fin n → Fin n → ℝ} {θ : ℝ}
    (M : LayerD.SegmentModel (Flat n Mode)) [IsSFiniteKernel M.jump]
    (rd : Flat n Mode → Flat n Mode)
    (jump_guard_det : ∀ p, M.jump (p, LayerD.ExitKind.guard)
      = Measure.dirac (rd p))
    (dispatch_cover : ∀ p : Flat n Mode,
      (∃ e : Fin T, (P.trans e).forced = true ∧ (P.trans e).src = p.1
        ∧ p.2 ∈ (P.trans e).guard
        ∧ rd p = ((P.trans e).dst, (P.trans e).reset p.2))
      ∨ rd p = p)
    (hV : Measurable (flatCert V))
    (hdata : (∀ mo : Mode, ∀ z ∈ P.cell mo,
        θ * V mo z + generator P V gradV hessV mo z ≤ 0) →
      ∀ p : Flat n Mode,
        Nonempty (LayerD.ItoSupermartingaleData M θ (flatCert V) p)
        ∨ Nonempty (LayerD.PDMPFlowFinDini M θ (flatCert V) p)) :
    GshsSegmentVille P V gradV hessV θ
\end{leansrc}
\par\medskip
\textbf{Discharging the d\'ebut family.}\enspace The second analytic
field is a library application: exit from a closed certified cell is
the d\'ebut of the open complement, and the d\'ebut of an open set
under right-continuous paths is a stopping time of the
right-continuous filtration --- no completion, no metric, no per-model
measure theory.

\declhead{BernsteinLogicAux.Debut}{isStoppingTime\_debut}{\badge{BadgeGreen}{statement bound}}
\begin{leansrc}
/-- The début of an open set by a right-continuous adapted process is a
bona-fide mathlib **stopping time** for any right-continuous filtration
(`ℝ≥0∞ := WithTop ℝ≥0` definitionally, so no translation layer): the pin's
`isStoppingTime_of_measurableSet_lt_of_isRightContinuous` upgrades the
`<`-family. This is the wide-sense-to-strict step the census's G3 entry
priced; it costs an instance, never a per-model proof — models on a raw
filtration `𝓕` use `𝓕.rightCont` (`𝓕₊`), whose `IsRightContinuous`
instance mathlib provides. -/
theorem isStoppingTime_debut [ℱ.IsRightContinuous]
    (hX : ∀ t : ℝ≥0, Measurable[ℱ t] (X t))
    (hrc : ∀ ω, ∀ s : ℝ≥0, ContinuousWithinAt (fun u => X u ω) (Set.Ici s) s)
    (hU : IsOpen U) :
    IsStoppingTime ℱ (debut X U)
\end{leansrc}
\declhead{BernsteinLogicAux.Debut}{exitDebut\_le\_family\_of\_isClosed\_cell}{\badge{BadgeGreen}{statement bound}}
\begin{leansrc}
/-- The certified-cell instantiation: exit from a CLOSED cell (the
box-union geometry — finite unions of closed boxes are closed by
`isClosed_Icc`/`isClosed_set_pi` + `Set.Finite.isClosed_biUnion`, one line
per model) is the début of the open complement, so the `hτ` field is a
library application. -/
theorem exitDebut_le_family_of_isClosed_cell [ℱ.IsRightContinuous]
    {T : Ω → ℝ} {k : Ω → ExitKind} {X : ℝ≥0 → Ω → E} {cell : Set E}
    (hcell : IsClosed cell)
    (hEq : ∀ ω, exitDebut T k ω = debut X cellᶜ ω)
    (hX : ∀ t : ℝ≥0, Measurable[ℱ t] (X t))
    (hrc : ∀ ω, ∀ s : ℝ≥0, ContinuousWithinAt (fun u => X u ω) (Set.Ici s) s) :
    ∀ s : ℝ≥0, MeasurableSet[ℱ s] {ω | exitDebut T k ω ≤ (s : ℝ≥0∞)}
\end{leansrc}
\par\medskip
\textbf{Discharging the supermartingale property, generic half.}\enspace
The first analytic field\textquotesingle s martingale theory is also a
library application: the exponential-discounting theorem produces the
required supermartingale from a martingale-problem hypothesis (the
compensated certificate process is a martingale --- what It\^o
\textquotesingle s formula and the interior rows establish per model,
with no stochastic integration in the certified statement), the
bounded-horizon and horizon-free optional-sampling inequalities are the
theorems the interface discharge above consumes, and the continuous-time
Doob/Ville maximal inequality converts supermartingale level crossings
into probability bounds. What remains per model is exactly the
martingale-problem hypothesis and the uniform bounds it consumes.

\declhead{BernsteinLogicAux.OptionalStoppingRC}{supermartingale\_exp\_discount}{\badge{BadgeGreen}{statement bound}}
\begin{leansrc}
/-- **Exponential discounting of a drift-dominated process.** Suppose
`W : ℝ≥0 → Ω → ℝ` is adapted with continuous paths, `g : ℝ≥0 → Ω → ℝ` is
jointly measurable, both are uniformly bounded, the compensated process
`fun t ω => W t ω - ∫ s in (0 : ℝ)..(t : ℝ), g s.toNNReal ω` is a
martingale (the martingale-problem formulation of `g` being the generator
drift of `W`), and `g t ω ≤ c * W t ω` pointwise. Then

`fun t ω => Real.exp (-c * t) * W t ω` is a supermartingale.

This is the discounting step of Dynkin-formula arguments, carried out with
pathwise Riemann–Stieltjes integration by parts against the `C¹` function
`s ↦ exp (-c * s)` and a conditional Fubini argument — no stochastic
integral is involved. -/
theorem supermartingale_exp_discount
    [IsProbabilityMeasure μ]
    (W : ℝ≥0 → Ω → ℝ) (g : ℝ≥0 → Ω → ℝ) (c : ℝ) (C : ℝ)
    (hWadp : Adapted ℱ W)
    (hWcont : ∀ ω, Continuous fun t => W t ω)
    (hWb : ∀ t ω, |W t ω| ≤ C)
    (hgmeas : Measurable (Function.uncurry g))
    (hgb : ∀ t ω, |g t ω| ≤ C)
    (hmart : Martingale
      (fun (t : ℝ≥0) (ω : Ω) =>
        W t ω - ∫ s in (0 : ℝ)..(t : ℝ), g s.toNNReal ω) ℱ μ)
    (hdom : ∀ t ω, g t ω ≤ c * W t ω) :
    Supermartingale
      (fun (t : ℝ≥0) (ω : Ω) => Real.exp (-c * (t : ℝ)) * W t ω) ℱ μ
\end{leansrc}
\declhead{BernsteinLogicAux.OptionalStoppingRC}{lintegral\_stoppedValue\_min\_le}{\badge{BadgeGreen}{statement bound}}
\begin{leansrc}
/-- The upper dyadic approximation to a nonnegative real. -/
noncomputable def dyadicCeil (n : ℕ) (x : ℝ≥0) : ℕ :=
  ⌈x * (2 ^ n : ℝ≥0)⌉₊

lemma dyadicCeil_time_tendsto (x : ℝ≥0) :
    Filter.Tendsto (fun n => dyadicTime n (dyadicCeil n x)) Filter.atTop (nhds x) := by
  -- Use the fact that $x \leq \lceil x \cdot 2^n \rceil / 2^n \leq x + 1 / 2^n$.
  have h_bounds : ∀ n : ℕ, x ≤ (dyadicCeil n x : ℝ≥0) / 2^n ∧ (dyadicCeil n x : ℝ≥0) / 2^n ≤ x + 1 / 2^n := by
    intro n
    simp [dyadicCeil, dyadicTime];
    rw [ le_div_iff₀, div_le_iff₀ ] <;> norm_num;
    exact ⟨ Nat.le_ceil _, by rw [ add_mul, inv_mul_cancel₀ ( by positivity ) ] ; exact_mod_cast Nat.ceil_lt_add_one ( by positivity ) |> le_of_lt ⟩;
  refine' tendsto_of_tendsto_of_tendsto_of_le_of_le' tendsto_const_nhds _ ( Filter.Eventually.of_forall fun n => h_bounds n |>.1 ) ( Filter.Eventually.of_forall fun n => h_bounds n |>.2 );
  convert tendsto_const_nhds.add ( tendsto_inv_atTop_zero.comp ( tendsto_pow_atTop_atTop_of_one_lt ( show ( 1 : ℝ≥0 ) < 2 by norm_num ) ) ) using 2 ; norm_num;
  exacts [ rfl, by rw [ add_zero ] ]

lemma dyadicCeil_time_mem_Ici (n : ℕ) (x : ℝ≥0) :
    dyadicTime n (dyadicCeil n x) ∈ Set.Ici x := by
  unfold dyadicTime dyadicCeil;
  exact Set.mem_Ici.mpr ( by rw [ le_div_iff₀ ( by positivity ) ] ; exact Nat.le_ceil _ )

lemma dyadicCeil_le_iff (n k : ℕ) (x : ℝ≥0) :
    dyadicCeil n x ≤ k ↔ x ≤ dyadicTime n k := by
  rw [dyadicCeil, Nat.ceil_le, dyadicTime, le_div_iff₀ (by positivity)]

lemma dyadicCeil_stoppingTime
    (τ : Ω → ℝ≥0) (hτ : ∀ t : ℝ≥0, MeasurableSet[ℱ t] {ω | τ ω ≤ t})
    (T : ℝ≥0) (n : ℕ) :
    IsStoppingTime (dyadicFiltration (ℱ := ℱ) n)
      (fun ω => (dyadicCeil n (min (τ ω) T) : ℕ∞)) := by
  intro k
  by_cases hT : T ≤ dyadicTime n k
  · refine MeasurableSet.univ.congr ?_
    refine (Set.eq_univ_of_forall fun ω => ?_).symm
    show (dyadicCeil n (min (τ ω) T) : ℕ∞) ≤ (k : ℕ∞)
    exact_mod_cast (dyadicCeil_le_iff n k (min (τ ω) T)).mpr ((min_le_right _ _).trans hT)
  · refine (hτ (dyadicTime n k)).congr ?_
    ext ω
    simp only [Set.mem_setOf_eq]
    constructor
    · intro h
      show (dyadicCeil n (min (τ ω) T) : ℕ∞) ≤ (k : ℕ∞)
      exact_mod_cast (dyadicCeil_le_iff n k (min (τ ω) T)).mpr ((min_le_left _ _).trans h)
    · intro h
      have h' : (dyadicCeil n (min (τ ω) T) : ℕ∞) ≤ (k : ℕ∞) := h
      have hnk : dyadicCeil n (min (τ ω) T) ≤ k := by exact_mod_cast h'
      have := (dyadicCeil_le_iff n k (min (τ ω) T)).mp hnk
      exact ((min_le_iff.mp this).resolve_right hT)

lemma dyadicCeil_bounded (τ : Ω → ℝ≥0) (T : ℝ≥0) (n : ℕ) (ω : Ω) :
    (dyadicCeil n (min (τ ω) T) : ℕ∞) ≤ dyadicCeil n T := by
  exact_mod_cast Nat.ceil_mono ( mul_le_mul_of_nonneg_right ( min_le_right _ _ ) ( by positivity ) )

lemma integral_dyadic_stopped_le
    [IsProbabilityMeasure μ]
    (Z : ℝ≥0 → Ω → ℝ) (hZ : Supermartingale Z ℱ μ)
    (τ : Ω → ℝ≥0)
    (hτ : ∀ t : ℝ≥0, MeasurableSet[ℱ t] {ω | τ ω ≤ t})
    (T : ℝ≥0) (n : ℕ) :
    ∫ ω, Z (dyadicTime n (dyadicCeil n (min (τ ω) T))) ω ∂μ
      ≤ ∫ ω, Z 0 ω ∂μ := by
  have h_submartingale :
      Submartingale (fun k => -Z (dyadicTime n k)) (dyadicFiltration (ℱ := ℱ) n) μ :=
    (supermartingale_dyadic Z hZ n).neg
  have h_stopping_time : IsStoppingTime (dyadicFiltration (ℱ := ℱ) n)
      (fun ω => (dyadicCeil n (min (τ ω) T) : ℕ∞)) := dyadicCeil_stoppingTime τ hτ T n
  have h0 : IsStoppingTime (dyadicFiltration (ℱ := ℱ) n) (fun _ : Ω => ((0 : ℕ) : ℕ∞)) := by
    intro k
    refine MeasurableSet.univ.congr ?_
    refine (Set.eq_univ_of_forall fun ω => ?_).symm
    show ((0 : ℕ) : ℕ∞) ≤ (k : ℕ∞)
    exact_mod_cast Nat.zero_le k
  have key := h_submartingale.expected_stoppedValue_mono h0 h_stopping_time
    (fun ω => by
      show ((0 : ℕ) : ℕ∞) ≤ (dyadicCeil n (min (τ ω) T) : ℕ∞)
      exact_mod_cast Nat.zero_le _)
    (N := dyadicCeil n T) (fun ω => by exact_mod_cast dyadicCeil_bounded τ T n ω)
  have hL : stoppedValue (fun k => -Z (dyadicTime n k)) (fun _ : Ω => ((0 : ℕ) : ℕ∞))
      = fun ω => -Z (dyadicTime n 0) ω := rfl
  have hd0 : dyadicTime n 0 = 0 := by simp [dyadicTime]
  have hR : stoppedValue (fun k => -Z (dyadicTime n k))
      (fun ω => (dyadicCeil n (min (τ ω) T) : ℕ∞))
      = fun ω => -Z (dyadicTime n (dyadicCeil n (min (τ ω) T))) ω := rfl
  rw [hL, hR, hd0] at key
  simp only [integral_neg] at key
  exact neg_le_neg_iff.mp key

lemma integrable_dyadic_stopped
    [IsProbabilityMeasure μ]
    (Z : ℝ≥0 → Ω → ℝ) (hZ : Supermartingale Z ℱ μ)
    (τ : Ω → ℝ≥0)
    (hτ : ∀ t : ℝ≥0, MeasurableSet[ℱ t] {ω | τ ω ≤ t})
    (T : ℝ≥0) (n : ℕ) :
    Integrable (fun ω => Z (dyadicTime n (dyadicCeil n (min (τ ω) T))) ω) μ := by
  have h_sub : Submartingale (fun k => -Z (dyadicTime n k)) (dyadicFiltration (ℱ := ℱ) n) μ :=
    (supermartingale_dyadic Z hZ n).neg
  have h_int := h_sub.integrable_stoppedValue (dyadicCeil_stoppingTime τ hτ T n)
    (N := dyadicCeil n T) (fun ω => by exact_mod_cast dyadicCeil_bounded τ T n ω)
  have hR : stoppedValue (fun k => -Z (dyadicTime n k))
      (fun ω => (dyadicCeil n (min (τ ω) T) : ℕ∞))
      = fun ω => -Z (dyadicTime n (dyadicCeil n (min (τ ω) T))) ω := rfl
  rw [hR] at h_int
  simpa using h_int.neg

lemma dyadic_stopped_tendsto
    (Z : ℝ≥0 → Ω → ℝ)
    (hZrc : ∀ ω, ∀ t : ℝ≥0,
      ContinuousWithinAt (fun s => Z s ω) (Set.Ici t) t)
    (τ : Ω → ℝ≥0) (T : ℝ≥0) (ω : Ω) :
    Filter.Tendsto
      (fun n => Z (dyadicTime n (dyadicCeil n (min (τ ω) T))) ω)
      Filter.atTop (nhds (Z (min (τ ω) T) ω)) := by
  refine' Filter.Tendsto.comp ( hZrc ω ( min ( τ ω ) T ) ) _;
  rw [ tendsto_nhdsWithin_iff ];
  exact ⟨ dyadicCeil_time_tendsto _, Filter.Eventually.of_forall fun n => dyadicCeil_time_mem_Ici _ _ ⟩

/-
**Optional sampling inequality at a bounded horizon.** For a nonnegative
supermartingale `Z` over `ℝ≥0` with right-continuous paths, a stopping time
`τ` (given by its measurability family) and a deterministic horizon `T`,

`∫⁻ ω, ENNReal.ofReal (Z (min (τ ω) T) ω) ∂μ ≤ ∫⁻ ω, ENNReal.ofReal (Z 0 ω) ∂μ`.

Stated with the lower Lebesgue integral so that no integrability of the
stopped value needs to be assumed.
-/
theorem lintegral_stoppedValue_min_le
    [IsProbabilityMeasure μ]
    (Z : ℝ≥0 → Ω → ℝ) (hZ : Supermartingale Z ℱ μ)
    (hZnn : ∀ t ω, 0 ≤ Z t ω)
    (hZrc : ∀ ω, ∀ t : ℝ≥0,
      ContinuousWithinAt (fun s => Z s ω) (Set.Ici t) t)
    (τ : Ω → ℝ≥0)
    (hτ : ∀ t : ℝ≥0, MeasurableSet[ℱ t] {ω | τ ω ≤ t})
    (T : ℝ≥0) :
    ∫⁻ ω, ENNReal.ofReal (Z (min (τ ω) T) ω) ∂μ
      ≤ ∫⁻ ω, ENNReal.ofReal (Z 0 ω) ∂μ
\end{leansrc}
\declhead{BernsteinLogicAux.OptionalStoppingRC}{lintegral\_stoppedValue\_le}{\badge{BadgeGreen}{statement bound}}
\begin{leansrc}
/-- Truncate a wide-sense nonnegative stopping time at a finite horizon. -/
noncomputable def ennrealTrunc (τ : Ω → ℝ≥0∞) (T : ℝ≥0) (ω : Ω) : ℝ≥0 :=
  if τ ω ≤ (T : ℝ≥0∞) then (τ ω).toNNReal else T

lemma ennrealTrunc_le_iff (τ : Ω → ℝ≥0∞) (T t : ℝ≥0) (ht : t < T) :
    {ω | ennrealTrunc τ T ω ≤ t} = {ω | τ ω ≤ (t : ℝ≥0∞)} := by
  ext ω;
  by_cases h : τ ω ≤ T <;> simp +decide [ h, ennrealTrunc ];
  · cases h' : τ ω <;> aesop;
  · exact iff_of_false ( not_le_of_gt ht ) ( fun h' => h <| le_trans h' <| mod_cast ht.le )

lemma ennrealTrunc_eq_of_le (τ : Ω → ℝ≥0∞) (T : ℝ≥0) (ω : Ω)
    (h : τ ω ≤ (T : ℝ≥0∞)) :
    ennrealTrunc τ T ω = (τ ω).toNNReal := by
  exact if_pos h

/-
**Optional sampling inequality, horizon-free form.** For a nonnegative
right-continuous supermartingale `Z` and a wide-sense stopping time `τ` with
values in `ℝ≥0∞` (given by its measurability family), the contribution of
the event `{τ < ∞}` satisfies

`∫⁻ ω, (if τ ω = ∞ then 0 else ENNReal.ofReal (Z (τ ω).toNNReal ω)) ∂μ
   ≤ ∫⁻ ω, ENNReal.ofReal (Z 0 ω) ∂μ`.

Follows from `lintegral_stoppedValue_min_le` by monotone convergence along
an increasing sequence of horizons, using nonnegativity to discard the
`{τ > T}` remainder.
-/
theorem lintegral_stoppedValue_le
    [IsProbabilityMeasure μ]
    (Z : ℝ≥0 → Ω → ℝ) (hZ : Supermartingale Z ℱ μ)
    (hZnn : ∀ t ω, 0 ≤ Z t ω)
    (hZrc : ∀ ω, ∀ t : ℝ≥0,
      ContinuousWithinAt (fun s => Z s ω) (Set.Ici t) t)
    (τ : Ω → ℝ≥0∞)
    (hτ : ∀ t : ℝ≥0, MeasurableSet[ℱ t] {ω | τ ω ≤ (t : ℝ≥0∞)}) :
    ∫⁻ ω, (if τ ω = ∞ then 0
           else ENNReal.ofReal (Z (τ ω).toNNReal ω)) ∂μ
      ≤ ∫⁻ ω, ENNReal.ofReal (Z 0 ω) ∂μ
\end{leansrc}
\declhead{BernsteinLogicAux.OptionalStoppingRC}{measure\_sup\_ge\_le\_of\_supermartingale}{\badge{BadgeGreen}{statement bound}}
\begin{leansrc}
/-- Restrict a supermartingale along a monotone map of time indices. -/
private lemma Supermartingale.comp_monotone
    {ι κ : Type*} [Preorder ι] [Preorder κ]
    {ℱι : Filtration ι m0} {Z : ι → Ω → ℝ}
    (hZ : Supermartingale Z ℱι μ) (f : κ → ι) (hf : Monotone f) :
    Supermartingale (fun k ↦ Z (f k))
      ({ seq := fun k ↦ ℱι (f k)
         mono' := fun _ _ h ↦ ℱι.mono (hf h)
         le' := fun k ↦ ℱι.le (f k) } : Filtration κ m0) μ := by
  exact ⟨fun k ↦ hZ.stronglyAdapted (f k),
    fun i j hij ↦ hZ.condExp_ae_le (hf hij), fun k ↦ hZ.integrable (f k)⟩

/-
A continuous path hitting `c` on `[0,T]` is approximated on uniform grids
at every strictly smaller level.
-/
private lemma exists_grid_hit_of_continuous
    (z : ℝ≥0 → ℝ) (hz : Continuous z) (T : ℝ≥0) (c d : ℝ) (hdc : d < c)
    (h : ∃ s : ℝ≥0, s ≤ T ∧ c ≤ z s) :
    ∃ m : ℕ, 0 < m ∧ ∃ k ≤ m,
      d ≤ z (⟨(k : ℝ) * T / m, by positivity⟩ : ℝ≥0) := by
  by_contra h_contra;
  obtain ⟨ s, hs₁, hs₂ ⟩ := h;
  -- By the continuity of $z$, for any $\epsilon > 0$, there exists a $\delta > 0$ such that $|z(t) - z(s)| < \epsilon$ whenever $|t - s| < \delta$.
  obtain ⟨δ, hδ_pos, hδ⟩ : ∃ δ > 0, ∀ t : ℝ≥0, |(t : ℝ) - s| < δ → |z t - z s| < (c - d) / 2 := by
    obtain ⟨δ, hδ, hδ'⟩ :=
      Metric.continuous_iff.mp hz s ((c - d) / 2) (half_pos (sub_pos.mpr hdc))
    refine ⟨δ, hδ, fun t ht ↦ ?_⟩
    simpa only [Real.dist_eq] using hδ' t (by simpa only [NNReal.dist_eq] using ht);
  -- Choose $m$ large enough such that $T/m < \delta$.
  obtain ⟨m, hm_pos, hm⟩ : ∃ m : ℕ, 0 < m ∧ (T : ℝ) / m < δ := by
    exact ⟨ ⌊T / δ⌋₊ + 1, Nat.succ_pos _, by rw [ div_lt_iff₀ ] <;> push_cast <;> nlinarith [ Nat.lt_floor_add_one ( T / δ ), mul_div_cancel₀ ( T : ℝ ) hδ_pos.ne' ] ⟩;
  -- Choose $k$ such that $k * T / m$ is within $\delta$ of $s$.
  obtain ⟨k, hk⟩ : ∃ k : ℕ, k ≤ m ∧ |(k * T / m : ℝ) - s| < δ := by
    refine' ⟨ ⌊ ( s : ℝ ) * m / T⌋₊, _, _ ⟩;
    · exact Nat.floor_le_of_le ( div_le_of_le_mul₀ ( by positivity ) ( by positivity ) ( by nlinarith [ show ( s : ℝ ) ≤ T from hs₁, show ( m : ℝ ) ≥ 1 by norm_cast ] ) );
    · by_cases hT : T = 0 <;> simp_all +decide [div_lt_iff₀];
      rw [ abs_lt ] ; constructor <;> nlinarith [ Nat.floor_le ( show 0 ≤ ( s : ℝ ) * m / T by positivity ), Nat.lt_floor_add_one ( ( s : ℝ ) * m / T ), show ( T : ℝ ) > 0 by positivity, mul_div_cancel₀ ( ( s : ℝ ) * m ) ( show ( T : ℝ ) ≠ 0 by positivity ), mul_div_cancel₀ ( ( ⌊ ( s : ℝ ) * m / T⌋₊ : ℝ ) * T ) ( show ( m : ℝ ) ≠ 0 by positivity ) ];
  exact h_contra ⟨ m, hm_pos, k, hk.1, by linarith [ abs_lt.mp ( hδ ⟨ k * T / m, by positivity ⟩ hk.2 ) ] ⟩

/-
**Doob-type maximal inequality, continuous time.** For a nonnegative
supermartingale `Z` over the index set `ℝ≥0` with continuous paths, the
probability that the running supremum over `[0, T]` reaches a level `c > 0`
is at most `(∫ ω, Z 0 ω ∂μ) / c`.
-/
theorem measure_sup_ge_le_of_supermartingale
    [IsProbabilityMeasure μ]
    (Z : ℝ≥0 → Ω → ℝ) (hZ : Supermartingale Z ℱ μ)
    (hZnn : ∀ t ω, 0 ≤ Z t ω)
    (hZcont : ∀ ω, Continuous fun t => Z t ω)
    (T : ℝ≥0) (c : ℝ) (hc : 0 < c) :
    μ {ω | ∃ s : ℝ≥0, s ≤ T ∧ c ≤ Z s ω}
      ≤ ENNReal.ofReal ((∫ ω, Z 0 ω ∂μ) / c)
\end{leansrc}
\par\medskip
\textbf{The martingale-problem profile, end to end.}\enspace The
general continuous-time theorems: the inter-jump motion enters through
its generator --- $\mathcal{A}V$ as data, tied to the process only by
the martingale hypothesis for the discounted compensated certificate
process --- the segment field is derived rather than assumed, and the
battery conclusion follows for every system whose presentation
satisfies the martingale-problem hypothesis.

\declhead{QtslAux.LayerD.MPDischarge}{sound\_GSHS\_safety\_mp}{\badge{BadgeGreen}{statement bound}}
\begin{leansrc}
/-- Safety (avoidance) at the martingale-problem profile: the probability
that the segment skeleton ever enters the bad set is at most `V x`. -/
theorem sound_GSHS_safety_mp [IsMarkovKernel M.seg] [IsMarkovKernel M.jump]
    {V : State → ℝ} {U : Set State} (hV : Measurable V) (hU : MeasurableSet U)
    (hmp : ∀ x : State, Nonempty (MPFlow M 0 V x))
    (hForce : ForcedNonexpansion M V)
    (hNonneg : NonnegV V) {C : ℝ} (hC0 : 0 ≤ C) (hVC : ∀ y, V y ≤ C)
    (hFloor : ∀ y ∈ U, 1 ≤ V y) (x : State) :
    Battery.pathMeasure (skel M) (SkelState.alive x)
        {ω | ∃ n, ω n ∈ skelBad U}
      ≤ ENNReal.ofReal (V x)
\end{leansrc}
\declhead{QtslAux.LayerD.MPDischarge}{sound\_GSHS\_safety\_battery\_mp}{\badge{BadgeGreen}{statement bound}}
\begin{leansrc}
/-- The safety battery at the martingale-problem profile. -/
theorem sound_GSHS_safety_battery_mp
    [IsMarkovKernel M.seg] [IsMarkovKernel M.jump]
    {V : State → ℝ} {U : Set State} {p thr : ℝ}
    (hV : Measurable V) (hU : MeasurableSet U)
    (hmp : ∀ x : State, Nonempty (MPFlow M 0 V x))
    (hForce : ForcedNonexpansion M V)
    (hNonneg : NonnegV V) {C : ℝ} (hC0 : 0 ≤ C) (hVC : ∀ y, V y ≤ C)
    (hFloor : ∀ y ∈ U, 1 ≤ V y)
    (hΘ : p + thr ≤ 1) (x : State) (hthr : V x ≤ thr) :
    ENNReal.ofReal p ≤ Battery.pathMeasure (skel M) (SkelState.alive x)
      {ω | ∀ n, ω n ∉ skelBad U}
\end{leansrc}
\par\medskip
The staged-row form binds the rule\textquotesingle s premise shapes as
the schema realizes them: per-mode interior rows against the
mode\textquotesingle s flow mechanism, per-jump guard rows, per-mode
floors and nonnegs, and the exact-$\Q$ threshold row cast once ---
assembled onto the flattened battery through the mixture bridge
discussed in the proof section\textquotesingle s final remark.

\declhead{BernsteinLogicAux.QtslBattery}{sound\_QSafetyBattery\_staged}{\badge{BadgeGreen}{statement bound}}
\begin{leansrc}
/-- **`sound_QSafetyBattery_staged`** — the `QSafetyBattery` node's staged
right column (QTSL-rules.toml row `QSafetyBattery`,
amended side) as a bound theorem, discrete-kernel profile. Hypothesis
families mirror the staged rosters:

* `hB1` **interiors** — per mode `m`, the flow chart's one-step drift
  `∫ V_m d(flow_m z) ≤ V_m z` on `D_m \ U_m` (the refined admissible scope;
  the B1 row), with `hint` its domain-membership integrability;
* `hB2` **guards** — per jump `j`, the boundary row
  `V_{dst j}(g_j z) ≤ V_{src j} z` on the guard slab (the B2 row);
* `hB3` **floors** and `hB4` **nonnegs** — per mode on `U ∩ D̄_m` / `D̄_m`
  (the B3/B4 rows);
* `hΘ` — the threshold row `p + thr ≤ 1` in **exact ℚ**, the staged
  `q_add`/`q_le` vocabulary, cast once here;
* `hR3`/`hR4` — the registry side conditions at the compiled kernel;
* `hmix` — **the System bridge**: the presentation's
  charts compile, at this profile, to the step kernel `K` — pointwise, the
  one-step measure from any in-scope state is a sub-convex mixture of the
  mode's flow step and the jumps enabled there (weights supported on
  `src j = m ∧ z ∈ G j`). A System whose compilation satisfies `hmix` is
  exactly what stands left of the concluded `⊨`.

Conclusion: region-level, as the concluded `QRegionSat` with the
`qraSublevel(V, thr)` atom reads — from every start in
`D_m ∩ {V_m ≤ thr}`, the chain satisfies `𝖦 φ` with probability at least
`p`. Assembled onto `sound_QSafetyBattery` through `mixture_drift` /
`mixture_integrable`; the GSHS-profile transport is the standing
residue. -/
theorem sound_QSafetyBattery_staged
    {K : Kernel (Md × X) (Md × X)} [IsMarkovKernel K]
    {Dm Dbarm Um : Md → Set X} {p thr : ℚ} {sat : Md × X → Prop}
    (hDm : ∀ m, MeasurableSet (Dm m)) (hDbarm : ∀ m, MeasurableSet (Dbarm m))
    (hUm : ∀ m, MeasurableSet (Um m)) (hVmeas : ∀ m, Measurable (V m))
    (hF : ∀ m z, ¬ sat (m, z) → z ∈ Um m)
    (hΘ : p + thr ≤ 1)
    (hDsub : ∀ m, Dm m ⊆ Dbarm m)
    (hR3 : ∀ m z, z ∈ Dm m → z ∉ Um m → K (m, z) {y | y.2 ∈ Dbarm y.1} = 1)
    (hR4 : ∀ m z, z ∈ Dm m → z ∉ Um m →
      K (m, z) ({y | y.2 ∈ Dm y.1} ∪ {y | y.2 ∈ Um y.1}) = 1)
    (hint : ∀ m z, z ∈ Dm m → z ∉ Um m → Integrable (V m) ((flow m) z))
    (hB1 : ∀ m z, z ∈ Dm m → z ∉ Um m → ∫ ζ, V m ζ ∂(flow m) z ≤ V m z)
    (hB2 : ∀ j, ∀ z ∈ G j, V (dst j) (g j z) ≤ V (src j) z)
    (hB3 : ∀ m z, z ∈ Um m → z ∈ Dbarm m → 1 ≤ V m z)
    (hB4 : ∀ m z, z ∈ Dbarm m → 0 ≤ V m z)
    (hmix : ∀ m z, z ∈ Dm m → z ∉ Um m →
      ∃ (w₀ : ℝ≥0) (w : J → ℝ≥0), w₀ + ∑ j, w j = 1 ∧
        (∀ j, w j ≠ 0 → src j = m ∧ z ∈ G j) ∧
        K (m, z) = (w₀ : ℝ≥0∞) • ((flow m) z).map (Prod.mk m)
          + ∑ j, (w j : ℝ≥0∞) • Measure.dirac (dst j, g j z)) :
    ∀ m z, z ∈ Dm m → V m z ≤ (thr : ℝ) →
      ENNReal.ofReal (p : ℝ)
        ≤ pathMeasure K (m, z) {ω | ∀ n, sat (ω n)}
\end{leansrc}
\par\medskip
\textbf{The general stratum: acceptance events on the record
carrier.}\enspace Over the semantic universe\textquotesingle s Fell
carrier the safety acceptance event is measurable by the topology of
the site rather than by assumption: confinement events of closed
walls are Fell-closed, hence Borel under the Effros structure ---
over every trajectory window, with no compactness hypothesis
anywhere.

\declhead{SuFunctor.Site.FellAcceptance}{measurableSet\_confinedEvent}{\badge{BadgeGreen}{statement bound}}
\begin{leansrc}
/-- The safety acceptance event of a closed wall is Borel — no
compactness, no σ-compactness, no Hausdorffness. -/
theorem measurableSet_confinedEvent {A : Set α} (hA : IsClosed A) :
    MeasurableSet (confinedEvent A)
\end{leansrc}
\declhead{SuFunctor.Site.FellAcceptance}{measurableSet\_confinedEvent\_cylinder}{\badge{BadgeGreen}{statement bound}}
\begin{leansrc}
/-- **The safety acceptance event is Borel over every window**: the
confinement event of a closed safe region's cylinder, with no
compactness and no σ-compactness — the A8 discharge of the safety
battery's measurability pack. -/
theorem measurableSet_confinedEvent_cylinder {D : Set X}
    (hD : IsClosed D) :
    MeasurableSet
      (confinedEvent ((Set.univ : Set ↥W) ×ˢ D) : Set (FTraj W X))
\end{leansrc}
\par\medskip
\textbf{Walls over the record carrier.}\enspace The
battery\textquotesingle s conclusion shape at the general stratum:
per-window confinement classes with the full judgment bridge
(satisfaction equals almost-sure confinement), the indicator
transport law that admits them, the closed-graph strengthening
(record confinement is pointwise safety, unconditional in the
trajectory), and the composition seam (the window compositor of two
wall demands is the wall demand of the glued class --- persistence
survives composition).

\declhead{SuFunctor.Site.FellBehaviorSpec}{indPenM}{\badge{BadgeGreen}{statement bound}}
\begin{leansrc}
/-- The indicator penalty of a class at the general stratum: permitted
outright iff the joint's mass avoids the class's complement, forbidden
outright otherwise.  Total in the class (mathlib's outer-measure
evaluation); its transport law below is where measurability binds. -/
def indPenM (S : Set K) : MeasPen K
\end{leansrc}
\declhead{SuFunctor.Site.FellBehaviorSpec}{pull\_indPenM}{\badge{BadgeGreen}{statement bound}}
\begin{leansrc}
/-- **Transport of the indicator along a measurable map is the
indicator of the preimage** — the general-stratum form of
`pull_indPen`, with the two measurability hypotheses that stratum
charges: the map's, and the class's.  The latter is the A8 admission
ticket: only Borel acceptance events enter the indicator calculus. -/
theorem pull_indPenM (h : K → L) (hm : Measurable h) {S : Set L}
    (hS : MeasurableSet S) :
    MeasPen.pull h (indPenM S) = indPenM (h ⁻¹' S)
\end{leansrc}
\declhead{SuFunctor.Site.FellBehaviorSpec}{isPersistentMF\_clipDemandsF}{\badge{BadgeGreen}{statement bound}}
\begin{leansrc}
/-- **Persistence of clip demands**: window-calculus transitivity
alone (`penExtMF_penExtMF`), with openness of the site's windows as
the only hypothesis.  No condition on the bases whatsoever. -/
theorem isPersistentMF_clipDemandsF (U : ι → Set T)
    (hU : ∀ i, IsOpen (U i))
    (base : ∀ i, MeasPen (FTraj (U i) X)) :
    IsPersistentMF U (clipDemandsF U base)
\end{leansrc}
\declhead{SuFunctor.Site.FellBehaviorSpec}{satAtM\_specOfBehaviorF\_iff}{\badge{BadgeGreen}{statement bound}}
\begin{leansrc}
/-- **The judgment bridge at the general stratum**: a joint satisfies
the spec of a Fell behavior subtype at a window iff its mass lies in
the class there.  Forward consumes the reflexive demand; backward
consumes the subpresheaf condition, clip measurability (openness of
the site's windows), and class measurability (the A8 ticket) — each
exactly once. -/
theorem satAtM_specOfBehaviorF_iff (U : ι → Set T)
    (hU : ∀ i, IsOpen (U i)) (B : WBehaviorF U X)
    (hB : ∀ i, MeasurableSet (B.cls i)) (j : ι)
    (μ : MeasJoint (FTraj (U j) X)) :
    SatAtM μ (specOfBehaviorF U B j) ↔
      (μ.1 : Measure (FTraj (U j) X)) (B.cls j)ᶜ = 0
\end{leansrc}
\declhead{SuFunctor.Site.FellBehaviorSpec}{satAtM\_safetySpecF\_iff}{\badge{BadgeGreen}{statement bound}}
\begin{leansrc}
/-- The safety wall's judgment bridge, with its measurability paid by
A8: for a closed safe region, satisfaction of the safety spec at a
window is almost-sure confinement there — the SU-general form of the
safety battery's acceptance reading, with the measurability pack
discharged by the topology of the site. -/
theorem satAtM_safetySpecF_iff (U : ι → Set T)
    (hU : ∀ i, IsOpen (U i)) {D : Set X} (hD : IsClosed D) (j : ι)
    (μ : MeasJoint (FTraj (U j) X)) :
    SatAtM μ (specOfBehaviorF U (safetyBehaviorF U D) j) ↔
      (μ.1 : Measure (FTraj (U j) X))
        (confinedEvent ((Set.univ : Set ↥(U j)) ×ˢ D))ᶜ = 0
\end{leansrc}
\declhead{SuFunctor.Site.FellBehaviorSpec}{graphCl\_mem\_confinedEvent\_iff}{\badge{BadgeGreen}{statement bound}}
\begin{leansrc}
/-- **The A5 strengthening**: record confinement of a trajectory's
closed-up graph to a closed wall is pointwise confinement at every
stage — unconditional in the trajectory.  The skeleton route's
segment-constancy side condition is absent by construction: the record
constrains every `(t, x)` of the closed graph, mid-segment stages
included, and the closure completion costs nothing against a closed
wall. -/
theorem graphCl_mem_confinedEvent_iff {D : Set X} (hD : IsClosed D)
    (f : Traj W X) :
    graphCl f ∈ confinedEvent ((Set.univ : Set ↥W) ×ˢ D) ↔
      ∀ w, f w ∈ D
\end{leansrc}
\declhead{SuFunctor.Site.FellBehaviorSpec}{mboxplusWF\_indPenM}{\badge{BadgeGreen}{statement bound}}
\begin{leansrc}
/-- **The record-carrier composition seam on indicators** (the
general-stratum sibling of `boxplusW_indPen`): the window compositor
of two wall demands is the wall demand of the glued record class —
the pullback intersection along the two state transports.  Properness
feeds the transports' measurability, and the classes' measurability
is the A8 ticket again.  This is the A4 transport core: composing
wall specs along the compositor never leaves the wall fragment, so
the composed conclusion is again a per-window class demand, and
persistence survives by `IsPersistentMF.image2_mboxplusWF`.  The
per-profile grid cells of the wire-composition rule (ODE, sampled
data, gauge-form feedback) are drift-assembly content at the
presentation profiles, deliberately not restated here. -/
theorem mboxplusWF_indPenM {W : Set T} {π₁ : Z → U'} {π₂ : Z → V'}
    (hp₁ : IsProperMap π₁) (hp₂ : IsProperMap π₂)
    {S₁ : Set (FTraj W U')} {S₂ : Set (FTraj W V')}
    (hS₁ : MeasurableSet S₁) (hS₂ : MeasurableSet S₂) :
    mboxplusWF π₁ π₂ (indPenM S₁) (indPenM S₂)
      = indPenM (fellMap (W
\end{leansrc}
\ruleblock{Correspondence (Steele display $\leftrightarrow$ Lean)}
\badge{BadgeAmber}{a bound Lean statement CHANGED since certification --- re-review and re-certify}
\sourcenote{The dictionary below is a reviewer-certified statement of how the mechanized statements above interpret the displayed rule: display metavariable $\mapsto$ Lean binder, each display hypothesis $\mapsto$ the premise or side condition carrying it, the display conclusion $\mapsto$ the goal --- with any deliberate difference in width between display and mechanization named outright. Certification is recorded as a hash triple over (the displayed cell, the extracted statement texts, this dictionary\textquotesingle s own text); if any of the three changes, this block turns amber until a reviewer re-certifies. The correspondence claim itself remains editorial --- reviewer-grade, not machine-checked; its machine-checked successor types each soundness statement against the rule as a structured object, and supersedes this block per rule as it lands.}
Model frame as at \rulename{QBasin}, probabilistically read: the
displayed $\mathcal{M},\ x_0 \in \{V \le y\} \vDash
\mathsf{SProbGE}\bigl(p,\ \mathsf{PAlways}\ \varphi\bigr)$ $\mapsto$
the goal $\mathrm{ofReal}\ p \le$
\texttt{pathMeasure K x \{ω | ∀ n, sat (ω n)\}} for every start in the
concluded region --- \texttt{sat} the $\varphi$-satisfaction
predicate, of which \emph{no measurability is asked} (the fail-closed
event is the measurable one; the tie weakens by
\texttt{measure\_mono}). The staged-row theorem
\texttt{sound\_QSafetyBattery\_staged} is the amended
display\textquotesingle s binding; the flattened
\texttt{sound\_QSafetyBattery} is its engine.
\begin{itemize}
\item profile, stated plainly: the map splits in two. The bullets
  below marked \emph{discrete map} bind the rule at the
  \emph{discrete-kernel profile} --- $\mathcal{M}$\textquotesingle s
  step semantics is a Markov kernel chain, its trajectory measure
  rendered as \texttt{pathMeasure} in the bound statements --- via
  \texttt{sound\_QSafetyBattery\_staged}, with the PDMP sub-profile
  carried end-to-end by \texttt{sound\_GSHS\_safety\_battery}. The
  martingale-problem theorems (\texttt{sound\_GSHS\_safety\_mp},
  \texttt{sound\_GSHS\_safety\_battery\_mp} --- hash-bound below) are
  the \emph{continuous-time avoidance engine}, mapped separately in
  the \emph{MP map} bullet; they do not themselves bind the full
  displayed judgment.
\item \emph{MP map}, its own binders and its own goal:
  the generator presentation and the interior rows $\mapsto$
  \texttt{hmp} $:$ $\forall x$, \texttt{Nonempty (MPFlow M 0 V x)}
  (the martingale-problem presentation hypothesis, carried as data);
  guards $\mapsto$ \texttt{hForce} $:$ \texttt{ForcedNonexpansion M
  V}; nonnegs $\mapsto$ \texttt{hNonneg} $:$ \texttt{NonnegV V};
  floors $\mapsto$ \texttt{hFloor} $:$ $\forall y \in U$, $1 \le V\,
  y$; the threshold row $\mapsto$ \texttt{hΘ}, and the start
  $\mapsto$ \texttt{x}/\texttt{hthr}. Undisplayed MP regularity, named
  outright: \texttt{hV}/\texttt{hU} (measurability),
  \texttt{hC0}/\texttt{hVC} (a global certificate bound $C$), and the
  \texttt{IsMarkovKernel} instances on \texttt{M.seg}/\texttt{M.jump}.
  Honesty about the goal: the MP battery\textquotesingle s literal
  conclusion is skeleton avoidance ---
  $\mathrm{ofReal}\ p \le$ \texttt{Battery.pathMeasure (skel M)
  (SkelState.alive x) \{ω | ∀ n, ω n ∉ skelBad U\}} --- avoidance of
  the $U$-bad set along the embedded skeleton, not the displayed
  $\mathsf{SProbGE}(p, \mathsf{PAlways}\ \varphi)$ judgment verbatim; the
  formula tie from $\varphi$ to $U$ is the discrete
  map\textquotesingle s \texttt{hF} business, and only the discrete
  staged theorem binds the displayed judgment itself.
\item \emph{discrete map} --- what stands left of the $\vDash$
  $\mapsto$ \texttt{hmix}, the
  System bridge: the presentation\textquotesingle s charts compile to
  the kernel as a sub-convex mixture of the mode\textquotesingle s
  flow step and the enabled jumps --- an explicit hypothesis, so the
  certified statement quantifies over exactly the Systems whose
  compilation satisfies it.
\item \emph{discrete map} --- the four premise families map
  one-to-one: interiors $\mapsto$
  \texttt{hB1} (per-mode one-step flow drift on $D_m \smallsetminus
  U_m$) with the undisplayed integrability rider \texttt{hint};
  guards $\mapsto$ \texttt{hB2} (per-jump boundary rows on the guard
  slab); floors $\mapsto$ \texttt{hB3} (on $U \cap \bar{D}_m$, the
  ruled fail-closed scope); nonnegs $\mapsto$ \texttt{hB4}.
\item \emph{discrete map} --- the threshold tie $p +
  y \le 1$ (a side condition in the display) $\mapsto$ \texttt{hΘ},
  in exact $\Q$, cast once inside the statement.
\item \emph{discrete map} --- side conditions R3/R4 $\mapsto$
  \texttt{hR3}/\texttt{hR4} at
  the compiled kernel (hull containment; fail-closed reset
  containment); the formula tie $U := \mathrm{cmpl}(\varphi)$
  $\mapsto$ \texttt{hF} $:$ $\neg\,\texttt{sat}\ y \to y \in U$ ---
  note the certified hypothesis asks only that $U$
  \emph{over-approximate} the falsity set: strictly wider than the
  displayed equation, in the safe direction.
\item \emph{discrete map} --- conclusion region
  $\mathsf{qraSublevel}(V, y)$
  $\mapsto$ the start conditions \texttt{z ∈ Dm m} and \texttt{V m z
  ≤ thr} (the Lean source\textquotesingle s own identifier for the
  display\textquotesingle s $y$), region-level as the judgment reads.
\end{itemize}
Width note: measurability of the region/certificate data
(\texttt{hDm}/\texttt{hDbarm}/\texttt{hUm}/\texttt{hVmeas}) and the
interior integrability rider are profile-level regularity with no
displayed rows --- the probabilistic sibling of the
\rulename{PsdCert} evaluability pack; every mandatory displayed row is
consumed --- the display\textquotesingle s \emph{optional}
$\mathrm{scope\_warrant}$ slot (the dissolved ScopeRule\textquotesingle
s interface) is taken by neither certified statement and has no entry
in this map --- and the two widenings (over-approximating $U$,
sub-convex \texttt{hmix}) both enlarge the certified reach beyond the
display\textquotesingle s letter.

\ruleblock{Lean verifier receipts}
\sourcenote{module \texttt{BernsteinLogicAux.Battery}; bound build \texttt{2fd1c031c80b\ldots} (bound from receipt 8d9ccb68993ea0ec, 2026-08-04T01:56:55Z). A receipt records that named checking layers accepted a build: (module, compiled-artifact SHA-256, environment SHA-256, checker layers), minted when the checking gate ran; the store receipts attesting the bound build are listed with the layers each run credited. Nothing here asks the reader to trust the store: every layer is an independently runnable program over the quoted sources --- rebuild the module and re-run the checkers to reproduce the verdict. Receipts attesting superseded builds credit no layer and are omitted from this rendering (the count of omissions, if any, is noted below).}
\begin{receipttable}
\receiptrow{8d9ccb68993ea0ec}{2026-08-04T01:56:55Z}{build+decorrelated}{\badge{BadgeGreen}{current}}
\receiptrow{5075586e4813a7c5}{2026-08-04T03:11:30Z}{build+decorrelated}{\badge{BadgeGreen}{current}}
\end{receipttable}
\sourcenote{13 further receipt(s) in the store attest superseded builds of this module and are omitted from this rendering.}
\sourcenote{module \texttt{BernsteinLogicAux.LayerD.Skeleton}; bound build \texttt{9a4f5ca8bc26\ldots} (bound from receipt f3d40eda96c4513b, 2026-08-04T01:57:01Z). A receipt records that named checking layers accepted a build: (module, compiled-artifact SHA-256, environment SHA-256, checker layers), minted when the checking gate ran; the store receipts attesting the bound build are listed with the layers each run credited. Nothing here asks the reader to trust the store: every layer is an independently runnable program over the quoted sources --- rebuild the module and re-run the checkers to reproduce the verdict. Receipts attesting superseded builds credit no layer and are omitted from this rendering (the count of omissions, if any, is noted below).}
\begin{receipttable}
\receiptrow{f3d40eda96c4513b}{2026-08-04T01:57:01Z}{build+decorrelated}{\badge{BadgeGreen}{current}}
\receiptrow{5966aa764026dde9}{2026-08-04T03:11:31Z}{build+decorrelated}{\badge{BadgeGreen}{current}}
\end{receipttable}
\sourcenote{13 further receipt(s) in the store attest superseded builds of this module and are omitted from this rendering.}
\sourcenote{module \texttt{BernsteinLogicAux.LayerD.DiniDischarge}; bound build \texttt{db41370b5e77\ldots} (bound from receipt 2a294d6f4ffc9a7b, 2026-07-21T04:51:51Z). A receipt records that named checking layers accepted a build: (module, compiled-artifact SHA-256, environment SHA-256, checker layers), minted when the checking gate ran; the store receipts attesting the bound build are listed with the layers each run credited. Nothing here asks the reader to trust the store: every layer is an independently runnable program over the quoted sources --- rebuild the module and re-run the checkers to reproduce the verdict. Receipts attesting superseded builds credit no layer and are omitted from this rendering (the count of omissions, if any, is noted below).}
\begin{receipttable}
\receiptrow{073080d556f0cf08}{2026-07-20T22:02:26Z}{build+decorrelated}{\badge{BadgeGreen}{current}}
\receiptrow{6f4929afcf219d67}{2026-07-20T22:52:48Z}{build+decorrelated}{\badge{BadgeGreen}{current}}
\receiptrow{f258f966d21ec502}{2026-07-21T00:03:10Z}{build+decorrelated}{\badge{BadgeGreen}{current}}
\receiptrow{2a294d6f4ffc9a7b}{2026-07-21T04:51:51Z}{build+decorrelated}{\badge{BadgeGreen}{current}}
\receiptrow{277ec2c0fa8227a9}{2026-07-21T05:38:41Z}{build+replay+hardened+decorrelated}{\badge{BadgeGreen}{current}}
\receiptrow{a142e9bdad29ed20}{2026-07-22T18:27:03Z}{build+decorrelated}{\badge{BadgeGreen}{current}}
\receiptrow{6aeb017a3939a647}{2026-07-31T21:59:58Z}{build+decorrelated}{\badge{BadgeGreen}{current}}
\receiptrow{85b98a8efabbfc19}{2026-07-31T22:28:52Z}{build+decorrelated}{\badge{BadgeGreen}{current}}
\receiptrow{fcd8973bd3d8aa0b}{2026-08-02T23:31:02Z}{build+decorrelated}{\badge{BadgeGreen}{current}}
\receiptrow{45fcba5d9746e88b}{2026-08-03T16:27:06Z}{build+decorrelated}{\badge{BadgeGreen}{current}}
\receiptrow{5a035508867866dd}{2026-08-03T22:45:49Z}{build+decorrelated}{\badge{BadgeGreen}{current}}
\end{receipttable}
\sourcenote{module \texttt{BernsteinLogicAux.LayerD.ItoInterface}; bound build \texttt{9b18e4e710f1\ldots} (bound from receipt 8b9184a6df93a83e, 2026-07-21T04:51:52Z). A receipt records that named checking layers accepted a build: (module, compiled-artifact SHA-256, environment SHA-256, checker layers), minted when the checking gate ran; the store receipts attesting the bound build are listed with the layers each run credited. Nothing here asks the reader to trust the store: every layer is an independently runnable program over the quoted sources --- rebuild the module and re-run the checkers to reproduce the verdict. Receipts attesting superseded builds credit no layer and are omitted from this rendering (the count of omissions, if any, is noted below).}
\begin{receipttable}
\receiptrow{4e3a32599ecb1605}{2026-07-21T00:03:12Z}{build+decorrelated}{\badge{BadgeGreen}{current}}
\receiptrow{8b9184a6df93a83e}{2026-07-21T04:51:52Z}{build+decorrelated}{\badge{BadgeGreen}{current}}
\receiptrow{04aa3a28f81ea9f9}{2026-07-21T05:38:43Z}{build+replay+hardened+decorrelated}{\badge{BadgeGreen}{current}}
\receiptrow{d95924f2dd7f318f}{2026-07-22T18:27:04Z}{build+decorrelated}{\badge{BadgeGreen}{current}}
\receiptrow{c78dba97e2eee2ba}{2026-07-31T21:59:59Z}{build+decorrelated}{\badge{BadgeGreen}{current}}
\receiptrow{ca54011b021b532e}{2026-07-31T22:28:52Z}{build+decorrelated}{\badge{BadgeGreen}{current}}
\receiptrow{bdf792a36abf7da8}{2026-08-02T23:31:02Z}{build+decorrelated}{\badge{BadgeGreen}{current}}
\receiptrow{e8478cde1b56df65}{2026-08-03T16:27:06Z}{build+decorrelated}{\badge{BadgeGreen}{current}}
\receiptrow{e987ba886a91fc22}{2026-08-03T22:45:50Z}{build+decorrelated}{\badge{BadgeGreen}{current}}
\end{receipttable}
\sourcenote{module \texttt{BernsteinLogicAux.LayerD.ItoSampling}; bound build \texttt{fb8de8b804d7\ldots} (bound from receipt 467820f314b0750c, 2026-07-21T08:16:24Z). A receipt records that named checking layers accepted a build: (module, compiled-artifact SHA-256, environment SHA-256, checker layers), minted when the checking gate ran; the store receipts attesting the bound build are listed with the layers each run credited. Nothing here asks the reader to trust the store: every layer is an independently runnable program over the quoted sources --- rebuild the module and re-run the checkers to reproduce the verdict. Receipts attesting superseded builds credit no layer and are omitted from this rendering (the count of omissions, if any, is noted below).}
\begin{receipttable}
\receiptrow{467820f314b0750c}{2026-07-21T08:16:24Z}{build+decorrelated}{\badge{BadgeGreen}{current}}
\end{receipttable}
\sourcenote{9 further receipt(s) in the store attest superseded builds of this module and are omitted from this rendering.}
\sourcenote{module \texttt{BernsteinLogicAux.GshsItoDischarge}; bound build \texttt{858d3234fc4f\ldots} (bound from receipt d6088699caa41da2, 2026-08-03T20:13:46Z). A receipt records that named checking layers accepted a build: (module, compiled-artifact SHA-256, environment SHA-256, checker layers), minted when the checking gate ran; the store receipts attesting the bound build are listed with the layers each run credited. Nothing here asks the reader to trust the store: every layer is an independently runnable program over the quoted sources --- rebuild the module and re-run the checkers to reproduce the verdict. Receipts attesting superseded builds credit no layer and are omitted from this rendering (the count of omissions, if any, is noted below).}
\begin{receipttable}
\receiptrow{1d9003db41305809}{2026-08-02T12:53:42Z}{build+decorrelated}{\badge{BadgeGreen}{current}}
\receiptrow{615e37d28395ee85}{2026-08-02T23:30:51Z}{build+decorrelated}{\badge{BadgeGreen}{current}}
\receiptrow{306da0b697e1ba52}{2026-08-03T16:26:54Z}{build+decorrelated}{\badge{BadgeGreen}{current}}
\receiptrow{d6088699caa41da2}{2026-08-03T20:13:46Z}{build+decorrelated}{\badge{BadgeGreen}{current}}
\receiptrow{b13efcc7c507c534}{2026-08-03T22:45:38Z}{build+decorrelated}{\badge{BadgeGreen}{current}}
\end{receipttable}
\sourcenote{7 further receipt(s) in the store attest superseded builds of this module and are omitted from this rendering.}
\sourcenote{module \texttt{BernsteinLogicAux.Debut}; bound build \texttt{cc670db616db\ldots} (bound from receipt 38571a839f9da48c, 2026-07-24T09:28:37Z). A receipt records that named checking layers accepted a build: (module, compiled-artifact SHA-256, environment SHA-256, checker layers), minted when the checking gate ran; the store receipts attesting the bound build are listed with the layers each run credited. Nothing here asks the reader to trust the store: every layer is an independently runnable program over the quoted sources --- rebuild the module and re-run the checkers to reproduce the verdict. Receipts attesting superseded builds credit no layer and are omitted from this rendering (the count of omissions, if any, is noted below).}
\begin{receipttable}
\receiptrow{a7c2ce15b3ea5ce7}{2026-07-22T16:23:09Z}{build+decorrelated}{\badge{BadgeGreen}{current}}
\receiptrow{e8af72cc3335497f}{2026-07-22T18:26:50Z}{build+decorrelated}{\badge{BadgeGreen}{current}}
\receiptrow{38571a839f9da48c}{2026-07-24T09:28:37Z}{build+decorrelated}{\badge{BadgeGreen}{current}}
\receiptrow{d0e0aa26ac8f8d34}{2026-07-31T21:59:41Z}{build+decorrelated}{\badge{BadgeGreen}{current}}
\receiptrow{b5cc2fcbabda9723}{2026-07-31T22:28:35Z}{build+decorrelated}{\badge{BadgeGreen}{current}}
\receiptrow{f494a0acd9a534b4}{2026-08-02T23:30:43Z}{build+decorrelated}{\badge{BadgeGreen}{current}}
\receiptrow{72ed3556543f0ada}{2026-08-03T16:26:47Z}{build+decorrelated}{\badge{BadgeGreen}{current}}
\receiptrow{aea75952b3723556}{2026-08-03T22:45:30Z}{build+decorrelated}{\badge{BadgeGreen}{current}}
\end{receipttable}
\sourcenote{module \texttt{BernsteinLogicAux.OptionalStoppingRC}; bound build \texttt{3f566c19afe8\ldots} (bound from receipt 704016ef1d4bd23d, 2026-07-24T09:28:43Z). A receipt records that named checking layers accepted a build: (module, compiled-artifact SHA-256, environment SHA-256, checker layers), minted when the checking gate ran; the store receipts attesting the bound build are listed with the layers each run credited. Nothing here asks the reader to trust the store: every layer is an independently runnable program over the quoted sources --- rebuild the module and re-run the checkers to reproduce the verdict. Receipts attesting superseded builds credit no layer and are omitted from this rendering (the count of omissions, if any, is noted below).}
\begin{receipttable}
\receiptrow{704016ef1d4bd23d}{2026-07-24T09:28:43Z}{build+decorrelated}{\badge{BadgeGreen}{current}}
\end{receipttable}
\sourcenote{6 further receipt(s) in the store attest superseded builds of this module and are omitted from this rendering.}
\sourcenote{module \texttt{QtslAux.LayerD.MPDischarge}; bound build \texttt{baded74dd943\ldots} (bound from receipt 7228834dea10e5f9, 2026-07-24T11:04:04Z). A receipt records that named checking layers accepted a build: (module, compiled-artifact SHA-256, environment SHA-256, checker layers), minted when the checking gate ran; the store receipts attesting the bound build are listed with the layers each run credited. Nothing here asks the reader to trust the store: every layer is an independently runnable program over the quoted sources --- rebuild the module and re-run the checkers to reproduce the verdict. Receipts attesting superseded builds credit no layer and are omitted from this rendering (the count of omissions, if any, is noted below).}
\begin{receipttable}
\receiptrow{7228834dea10e5f9}{2026-07-24T11:04:04Z}{build+decorrelated}{\badge{BadgeGreen}{current}}
\end{receipttable}
\sourcenote{module \texttt{BernsteinLogicAux.QtslBattery}; bound build \texttt{5c75f9923f27\ldots} (bound from receipt 3fd830b4914bdeda, 2026-08-04T01:56:58Z). A receipt records that named checking layers accepted a build: (module, compiled-artifact SHA-256, environment SHA-256, checker layers), minted when the checking gate ran; the store receipts attesting the bound build are listed with the layers each run credited. Nothing here asks the reader to trust the store: every layer is an independently runnable program over the quoted sources --- rebuild the module and re-run the checkers to reproduce the verdict. Receipts attesting superseded builds credit no layer and are omitted from this rendering (the count of omissions, if any, is noted below).}
\begin{receipttable}
\receiptrow{3fd830b4914bdeda}{2026-08-04T01:56:58Z}{build+decorrelated}{\badge{BadgeGreen}{current}}
\receiptrow{96b18d8572451f65}{2026-08-04T03:11:31Z}{build+decorrelated}{\badge{BadgeGreen}{current}}
\end{receipttable}
\sourcenote{13 further receipt(s) in the store attest superseded builds of this module and are omitted from this rendering.}
\sourcenote{module \texttt{SuFunctor.Site.FellAcceptance}; bound build \texttt{651fb557bbfe\ldots} (bound from receipt aa2169560f8f571a, 2026-08-03T20:19:04Z). A receipt records that named checking layers accepted a build: (module, compiled-artifact SHA-256, environment SHA-256, checker layers), minted when the checking gate ran; the store receipts attesting the bound build are listed with the layers each run credited. Nothing here asks the reader to trust the store: every layer is an independently runnable program over the quoted sources --- rebuild the module and re-run the checkers to reproduce the verdict. Receipts attesting superseded builds credit no layer and are omitted from this rendering (the count of omissions, if any, is noted below).}
\begin{receipttable}
\receiptrow{7ad85720de6e419e}{2026-08-03T10:13:50Z}{build+decorrelated}{\badge{BadgeGreen}{current}}
\receiptrow{94a408109996d1bf}{2026-08-03T13:41:21Z}{build+decorrelated}{\badge{BadgeGreen}{current}}
\receiptrow{580a7e551fab03b6}{2026-08-03T15:39:44Z}{build+decorrelated}{\badge{BadgeGreen}{current}}
\receiptrow{aa2169560f8f571a}{2026-08-03T20:19:04Z}{build+decorrelated}{\badge{BadgeGreen}{current}}
\end{receipttable}
\sourcenote{2 further receipt(s) in the store attest superseded builds of this module and are omitted from this rendering.}
\sourcenote{module \texttt{SuFunctor.Site.FellBehaviorSpec}; bound build \texttt{345f9caf9453\ldots} (bound from receipt 72f2c3b62b6771fc, 2026-08-03T20:19:08Z). A receipt records that named checking layers accepted a build: (module, compiled-artifact SHA-256, environment SHA-256, checker layers), minted when the checking gate ran; the store receipts attesting the bound build are listed with the layers each run credited. Nothing here asks the reader to trust the store: every layer is an independently runnable program over the quoted sources --- rebuild the module and re-run the checkers to reproduce the verdict. Receipts attesting superseded builds credit no layer and are omitted from this rendering (the count of omissions, if any, is noted below).}
\begin{receipttable}
\receiptrow{d7b846e785c1d7fb}{2026-08-03T10:13:53Z}{build+decorrelated}{\badge{BadgeGreen}{current}}
\receiptrow{c0e3dc15419765aa}{2026-08-03T13:41:24Z}{build+decorrelated}{\badge{BadgeGreen}{current}}
\receiptrow{9a62b29c08b63b4f}{2026-08-03T15:39:47Z}{build+decorrelated}{\badge{BadgeGreen}{current}}
\receiptrow{72f2c3b62b6771fc}{2026-08-03T20:19:08Z}{build+decorrelated}{\badge{BadgeGreen}{current}}
\end{receipttable}
\sourcenote{2 further receipt(s) in the store attest superseded builds of this module and are omitted from this rendering.}
\layerline{\layeritem{build}{lake build (elaboration)}{\badge{BadgeGreen}{current}}\layeritem{replay}{leanchecker (official kernel replay)}{\badge{BadgeGreen}{current}}\layeritem{hardened}{safe\_verify (hardened screen)}{\badge{BadgeGreen}{current}}\layeritem{decorrelated}{lean4export + nanoda (decorrelated kernel)}{\badge{BadgeGreen}{current}}}
\sourcenote{\badge{BadgeGray}{comparator: none} no fixture in the three-implementation comparator harness (compiled / kernel-interpreted / reference-interpreted) exercises this rule.}
\ruleblock{HOL4 statement}
\absentbox{NO HOL STATEMENT YET}{The HOL4 development ports the base calculus\textquotesingle s checker
chain (theory \texttt{bern\_catalogue}, developed side by side with the
base catalogue sources, stopping below the proof-object layer); it
contains no probability or measure theory.
\rulename{QSafetyBattery} lives one theory-extension edge above the base
calculus --- it concludes a satisfaction judgment, not a sequent --- and
its soundness content is path-space measure theory (trajectory kernels,
stopped supermartingales, Ville\textquotesingle s inequality), for which
no counterpart exists anywhere in the ported development. A HOL
statement of this rule would be a new measure-theoretic campaign, not an
extension of the ported checker chain; until one exists, this box
records the current state of the edifice.
}
\ruleblock{HOL verifier receipts}
\absentbox{NONE}{no HOL verification receipt exists for this rule.}
\ruleblock{Proof}

\sourcenote{Genre: core rule (tier 1) --- an ordinary mathematical proof
at Putnam-exam-solution detail. This rule concludes a \emph{satisfaction
judgment} ($\vDash$), one theory-extension edge above the base calculus:
every premise row is an inherited base-calculus sequent, and soundness
is proved \emph{given} base soundness (the extension-edge convention).
The theorems quoted above warrant the rule at the
\textbf{discrete-kernel profile}; the profile remark at the end of this
section states exactly what remains a hypothesis beyond it. Lemmas
\ref{lem:qsb-stopped}--\ref{lem:qsb-ville} and
Proposition~\ref{prop:qsb-pattern} are the battery family\textquotesingle s
shared probabilistic preamble.}

\textbf{Setting and claim.}\enspace Fix a measurable state space $S$ ---
the flattened hybrid state space, i.e.\ the disjoint sum of the per-mode
domains --- and a Markov one-step kernel $K$ on $S$, carrying both
discretized flow steps and guarded jumps. For $x \in S$ write $\Pr_x$
for the path measure on $S^{\N}$: the trajectory law of the
time-homogeneous chain driven by $K$, started deterministically at $x$
(the Ionescu--Tulcea construction, with the constant kernel family).
$(\mathcal{F}_n)_n$ is the canonical filtration, $\mathcal{F}_n$
generated by the coordinates $\le n$.

The rule\textquotesingle s data: a modelled domain $D$, a hull
$\bar{D} \supseteq D$ (the region on which the polynomial rows were
actually checked), an avoid set $U$ compiled to over-approximate the
falsity set of the $\mathsf{G}$-argument --- $D$, $\bar{D}$, $U$ all
measurable subsets of $S$, as is every region a row speaks on, here
and in the sibling rules riding this substrate --- a measurable
certificate $V : S \to \R$, and rationals $p, y$. The checked rows and
side conditions, read semantically at this profile:
\begin{itemize}
\item[(B1$\wedge$B2)] \emph{generator}: for every $z \in D \smallsetminus U$,
  $V$ is integrable against $K\,z$ and $\int V \,dK\,z \le V(z)$. (At
  this profile the interior and guard rows land \emph{together}: the
  one-step kernel carries flow and enabled jumps in a single step, so
  the two row families compile to one kernel-level generator row.)
\item[(B3)] \emph{floors}: $V \ge 1$ on $U \cap \bar{D}$.
\item[(B4)] \emph{nonnegs}: $V \ge 0$ on $\bar{D}$.
\item[(R3)] \emph{hull}: for every $z \in D \smallsetminus U$,
  $K\,z\,(\bar{D}) = 1$ --- one-step mass stays where the rows were
  checked.
\item[(R4)] \emph{fail-closed containment}: for every
  $z \in D \smallsetminus U$, $K\,z\,(D \cup U) = 1$.
\item[($\Theta$)] \emph{threshold}: $p + y \le 1$, checked in
  exact rational arithmetic.
\item[(F)] \emph{formula tie}: $\lnot \mathrm{sat}(z) \Rightarrow z \in U$,
  where $\mathrm{sat}$ is the state-satisfaction predicate of the
  $\mathsf{G}$-argument $\varphi$.
\end{itemize}
The claim (\texttt{sound\_QSafetyBattery}): for every $x \in D$ with
$V(x) \le y$,
\[
\Pr_x^{\,*}\bigl[\forall n,\ \mathrm{sat}(\omega_n)\bigr] \;\ge\; p
\]
--- outer measure, because no measurability is ever asked of the
satisfaction predicate (Lemma~\ref{lem:qsb-ville} states the
convention); it coincides with the ordinary probability whenever the
satisfaction event is measurable.

\begin{blemma}[the stopped certificate]\label{lem:qsb-stopped}
Let $f : S \to \R$ and let $A \subseteq S$ be a stop set. Define the
stopped certificate process $X_n(\omega) := f(\omega_{\,n \wedge \tau_A})$,
where $\tau_A$ is the first visit to $A$ --- equivalently, by the
primitive recursion that freezes the value once $A$ has been visited.
Then: (i) if $A$ is measurable and $f$ is measurable, $X$ is adapted to
$(\mathcal{F}_n)$; (ii) any pointwise property of $f$\textquotesingle s
values holds of $X$\textquotesingle s values; (iii) \emph{barrier
transport}: if $f \ge 1$ on $A$, then $X_n \ge 1$ on the event
$\{\tau_A \le n\}$; (iv) $X_n$ depends only on the coordinates $\le n$.
\end{blemma}
\begin{proof}
All four are elementary inductions on $n$ over the two branches of the
recursion (surviving: $X_{n+1} = f(\omega_{n+1})$; frozen:
$X_{n+1} = X_n$). No stopping-time API is needed anywhere --- the
recursion \emph{is} the stopped process, and (iii) is proved by chasing
the first index at which $A$ was visited.
\end{proof}

\begin{blemma}[supermartingale bridge]\label{lem:qsb-bridge}
Let $f$ be measurable with $0 \le f \le C$, let $A$ be measurable, and
suppose the one-step generator row holds \emph{off} $A$: for all
$z \notin A$, $\int f\,dK\,z \le f(z)$. Then for every start $x$, the
stopped certificate $X$ of Lemma~\ref{lem:qsb-stopped} is a
supermartingale with respect to $(\mathcal{F}_n, \Pr_x)$.
\end{blemma}
\begin{proof}
Integrability is the bound $0 \le X_n \le C$ (Lemma
\ref{lem:qsb-stopped}(ii)). For the supermartingale inequality,
condition on $\mathcal{F}_n$: by the Markov property of the trajectory
construction, the conditional expectation of a bounded
$\mathcal{F}_{n+1}$-measurable function given $\mathcal{F}_n$ is its
integral against the trajectory kernel continued from the $n$-prefix.
Split on the survival event
$S_n := \{\omega : \omega_k \notin A \ \text{for all } k \le n\}$ ---
$\mathcal{F}_n$-measurable, $A$ being measurable, and determined by the
coordinates $\le n$, so paths patched to agree with $\omega$ up to $n$
have the same survival status. On the surviving branch the integrand is
$f(\text{next coordinate})$, whose one-step marginal is $K(\omega_n)$;
since survival means $\omega_n \notin A$, the generator row applies and the
conditional expectation is $\le f(\omega_n) = X_n(\omega)$. On the
frozen branch the integrand is the constant $X_n(\omega)$. Note that
\emph{every} hypothesis here is pointwise --- there is no
almost-everywhere invariance argument anywhere in the bridge; this is
what the fail-closed design buys, and the reachability sibling
(\rulename{QDeadlineBattery}) will have to pay precisely this cost when
it can no longer enlarge its stop set.
\end{proof}

\begin{blemma}[Ville\textquotesingle s inequality, barrier form]
\label{lem:qsb-ville}
Let $X$ be a nonnegative supermartingale with $X_n \ge 1$ on the event
$\{$some coordinate $\le n$ lies in $A\}$. Then
$\Pr_x^{\,*}[\exists n,\ \omega_n \in A] \le \mathbb{E}\,X_0$, where
$\Pr_x^{\,*}$ is the outer measure --- equal to the ordinary
probability whenever the visit event is measurable, as it is at every
consumption in this family ($A$ is measurable there).
\end{blemma}
\begin{proof}
This is the maximal inequality for nonnegative supermartingales at
barrier $1$ \cite{ville1939} (textbook treatment: \cite{doob1953}): the
probability that a nonnegative supermartingale ever reaches level $a$ is
at most $\mathbb{E}\,X_0 / a$. The mechanized form deliberately asks no
measurability of the visit event itself --- the bad event is bounded by
inclusion in the \emph{measurable} event $\{\exists n,\ X_n \ge 1\}$,
whose probability the maximal inequality bounds, and only the process
needs to be adapted; the outer-measure phrasing of the conclusion is
exactly this inclusion, spelled.
\end{proof}

\begin{bprop}[\textbf{the battery pattern}]\label{prop:qsb-pattern}
Let $A$ be measurable and let $f$ be measurable with
$0 \le f \le C$ globally, $f \ge 1$ on $A$, and one-step drift off $A$.
Then for every $x$,
\[
\Pr_x\bigl[\exists n,\ \omega_n \in A\bigr] \;\le\; f(x).
\]
\end{bprop}
\begin{proof}
Lemmas \ref{lem:qsb-stopped}--\ref{lem:qsb-ville} assembled; the initial
value is $\mathbb{E}\,X_0 = f(x)$ because coordinate $0$ is
deterministically $x$. This proposition --- \emph{nonnegative
supermartingale $+$ barrier $+$ initial value $\Rightarrow$ probability
bound} --- is the shape every rule in the battery family consumes; the
siblings cite it as ``the battery pattern
(\rulename{QSafetyBattery}, Proposition~\ref{prop:qsb-pattern})''.
\end{proof}

\begin{blemma}[the clamp: from checked rows to global hypotheses]
\label{lem:qsb-clamp}
The rows (B3), (B4) speak only on $\bar{D}$, while
Proposition~\ref{prop:qsb-pattern} wants its bounds globally. Set
$A := U \cup D^{\mathrm{c}}$ (the fail-closed stop set),
$C := \max(y, 1)$, and let $g$ be the clamped certificate:
on $D$, clamp $V$ into $[0, C]$; off $D$, clamp $\max(V, 1)$ into
$[0,C]$ (so the floor is \emph{forced} off the modelled domain). Then:
(i) $0 \le g \le C$ globally; (ii) $g \ge 1$ on $A$; (iii) on
$\bar{D} \cap (D \cup U)$ the clamp is invisible:
$g = \min(V, C)$ there; (iv) under (R3) $+$ (R4), $g$ satisfies the
one-step generator row off $A$; (v) $g(x) = V(x)$ at any
$x \in D$ with $V(x) \le y$. Throughout, $g$ is measurable:
a case split between measurable pieces ($V$, $D$ measurable) composed
with clamps --- so $(g, A)$ is an admissible input to
Proposition~\ref{prop:qsb-pattern}.
\end{blemma}
\begin{proof}
(i), (ii), (v) are case checks from (B3), (B4) and the definition ---
(ii) uses (B3) on $U \cap \bar D$ and the forced floor off $D$. (iii):
on $D$ the lower clamp is inert by (B4); on
$(\bar{D} \cap U) \smallsetminus D$ the forced floor is inert by (B3).
For (iv), fix $z \in D \smallsetminus U$ (i.e.\ $z \notin A$). By (R3)
and (R4) the one-step measure $K\,z$ is carried by
$\bar{D} \cap (D \cup U)$ --- the complement has measure
$\le K z (\bar D^{\mathrm c}) + K z ((D \cup U)^{\mathrm c}) = 0$ ---
so by (iii) the integrand may be replaced by $\min(V, C)$. Then
$\int \min(V, C)\,dK\,z \le \min\bigl(\int V\,dK\,z,\; C\bigr)
\le \min(V(z), C) = g(z)$, using monotonicity of the integral, the
generator row (B1$\wedge$B2), and (iii) at $z$ itself.
\end{proof}

\begin{bthm}[fail-closed form]\label{thm:qsb-failclosed}
Under (B1)--(B4), (R3), (R4), ($\Theta$): for every $x \in D$ with
$V(x) \le y$,
\[
\Pr_x\bigl[\forall n,\ \omega_n \in D \smallsetminus U\bigr] \;\ge\; p .
\]
\end{bthm}
\begin{proof}
Apply Proposition~\ref{prop:qsb-pattern} to $(g, A)$ of
Lemma~\ref{lem:qsb-clamp}:
$\Pr_x[\text{reach } A] \le g(x) = V(x) \le y$. The
complement of ever visiting $A = U \cup D^{\mathrm{c}}$ is exactly
$\{\forall n,\ \omega_n \in D \smallsetminus U\}$, so its probability is
at least $1 - y \ge p$ by ($\Theta$). Note the concluded
event is \emph{stronger} than the safety formula\textquotesingle s truth
set: the chain never fails \emph{and never leaves the modelled domain}.
Leaving counts as failure --- that is the fail-closed reading, and it is
what (R4) is recomputed for.
\end{proof}

\begin{bprop}[soundness of \rulename{QSafetyBattery}]\label{prop:qsb-sound}
Under the hypotheses above including the formula tie (F): for every
$x \in D$ with $V(x) \le y$,
$\Pr_x^{\,*}[\forall n,\ \mathrm{sat}(\omega_n)] \ge p$ --- the concluded
satisfaction judgment
$\mathcal{M},\ x_0 \in \{V \le y\} \vDash
\mathsf{SProbGE}\bigl(p,\ \mathsf{PAlways}\ \varphi\bigr)$, read on the
sublevel region $\{V \le y\}$.
\end{bprop}
\begin{proof}
By the contrapositive of (F), $\omega_n \in D \smallsetminus U$ implies
$\mathrm{sat}(\omega_n)$; so the event of
Theorem~\ref{thm:qsb-failclosed} is contained in
$\{\forall n,\ \mathrm{sat}(\omega_n)\}$, and monotonicity of the
(outer) measure finishes. No measurability is ever asked of the
satisfaction predicate --- the smaller, fail-closed event is the
measurable one.
\end{proof}

\begin{bremark}[why the floors are not decoration]\label{rem:qsb-floors}
Drop (B3) and take $V \equiv 0$: drift, nonneg, and threshold rows all
hold, with $V(x) = 0 \le y$ for any $y \ge 0$, so
with $y = 0$, ($\Theta$) permits $p = 1$: the rule would
certify \emph{almost-sure} safety --- of any chain at all, including the
two-state chain $D = \{d\}$, $U = \{u\}$, $K\,d = \delta_u$, which fails
at its first step. With the floor present, this certificate dies
immediately: (B3) forces $V(u) \ge 1$, and then the generator row at $d$
demands $\int V\,dK\,d = V(u) \le V(d) = 0$ --- contradiction. Failure
must cost at least one unit of certificate value; the floor is the
price list.
\end{bremark}

\begin{bremark}[why the nonnegs are not decoration]\label{rem:qsb-nonnegs}
The maximal inequality of Lemma~\ref{lem:qsb-ville} is \emph{false} for
signed supermartingales, and a two-branch chain isolates the sign
hypothesis exactly --- every other row intact, boundedness included.
Take $D = \{s, b\}$, $U = \{a\}$, $\bar{D} = \{s, a, b\}$, with
$K\,s = \tfrac{1}{2}\delta_a + \tfrac{1}{2}\delta_b$,
$K\,b = \delta_b$, and the signed certificate $V(s) = 0$, $V(a) = 1$,
$V(b) = -1$. The generator row holds on $D \smallsetminus U$
($\int V\,dK\,s = 0 \le 0$; $\int V\,dK\,b = -1 \le -1$), the floor
holds on $U \cap \bar{D}$, the hull and fail-closed rows hold, and
$|V| \le 1$. Were Proposition~\ref{prop:qsb-pattern} available without
nonnegativity, the certificate would conclude
$\Pr_s[\text{reach } a] \le V(s) = 0$ --- but the chain reaches $a$
with probability $\tfrac{1}{2}$: the $a$-branch's barrier visit is
subsidized by the mass parked at $V = -1$ on the other branch, and
(B4) is what forbids the loan. The classical unbounded picture is the
simple symmetric random walk on $\mathbb{Z}$ with $V(z) = z$ and
$A = \{1\}$ --- same phenomenon, though as a drop-one witness it is
impure: it violates the boundedness hypothesis along with the sign,
which is why the two-branch chain is the recorded counterexample.
\end{bremark}

\begin{bremark}[the hull and fail-closed conditions (R3), (R4)]
\label{rem:qsb-r34}
Both are recomputed by the checker, never taken from the producer, and
each blocks a distinct hole. Without (R3), the one-step integral in the
generator row could consult states outside $\bar{D}$ --- states where the
floor and nonneg rows were \emph{never checked} --- and
Lemma~\ref{lem:qsb-clamp}(iv)\textquotesingle s clamp-invisibility
computation collapses. Without (R4), one-step mass could leave
$D \cup U$ entirely: the concluded event counts leaving the modelled
domain as failure, and (R4) is exactly what makes the clamped
certificate\textquotesingle s drift hold at the \emph{enlarged} stop set
$U \cup D^{\mathrm{c}}$ --- the design that let every hypothesis in
Lemma~\ref{lem:qsb-bridge} stay pointwise. A companion condition
(floor-scope resolution) is discharged at row-compilation time: it is
what makes the checked floor rows imply (B3) as stated here, and it has
no semantic residue at this layer.
\end{bremark}

\begin{bremark}[threshold and integrability]\label{rem:qsb-theta}
($\Theta$) is bookkeeping, and it is sharp: the certificate buys
$\Pr[\mathrm{fail}] \le y$, so asserting success
$\ge p$ requires $p \le 1 - y$, and
$y = 1 - p$ is the tightest admissible choice --- producer
weakening ($y < 1 - p$) is sound for free. The integrability
half of (B1$\wedge$B2) is the discrete-profile reading of
``the certificate lies in the generator\textquotesingle s
domain''; for the intended instances --- polynomial $V$, one-step
images inside the compact hull by (R3) --- it is automatic.
\end{bremark}

\begin{bremark}[\textbf{profile} --- what is proven, at which
presentations]
\label{rem:qsb-profile}
Proven, and quoted above this proof: the discrete-kernel profile ---
one-step Markov kernel semantics over the trajectory construction, the
theorems \texttt{sound\_QSafetyBattery} and
\texttt{sound\_QSafetyBattery\_failClosed}, sorry-free. Also proven,
end-to-end: the \emph{PDMP sub-profile} of the continuous-time semantics
--- piecewise-deterministic flow with guarded jumps, no diffusion
(\texttt{sound\_GSHS\_safety\_battery}). And proven, at the profile
that carries the general continuous-time semantics, jump--diffusions
included: the \emph{martingale-problem profile} --- the system
presented by its generator drift data, $\mathcal{A}V$ carried as data
beside the certificate with the martingale-problem hypothesis for the
pair, c\`adl\`ag paths per the presentation --- the theorems
\texttt{sound\_GSHS\_safety\_mp} and
\texttt{sound\_GSHS\_safety\_battery\_mp}, machine-checked sorry-free
in the semantic-universe strand\textquotesingle s auxiliary package
(module \texttt{QtslAux.LayerD.MPDischarge}). The three ingredients an
earlier edition of this remark filed as missing are each discharged
there: continuous-time generator domains ride the presentation ---
for the calculus\textquotesingle s own denoted class, compositional
symbolic piecewise-polynomial hypernets, $\mathcal{A}V$ is again
polynomial data and domain bookkeeping dissolves; fine path
properties ride the presentation\textquotesingle s c\`adl\`ag
clause; and optional stopping for right-continuous nonnegative
supermartingales is a library theorem, consumed at the exit
d\'ebut (\texttt{OptionalSamplingRC}, with
\texttt{supermartingale\_of\_compensator\_nonpos} and the
exponential-discount construction). The epistemic shape is the
discrete profile\textquotesingle s own: exactly as the staged theorem
quantifies over systems whose compilation satisfies the mixture
bridge, the martingale-problem theorems quantify over systems whose
presentation satisfies the martingale-problem hypothesis --- an
explicit bridge carried as data, never a silent assumption, so the
argument this section walks is now the theorem\textquotesingle s own
at every profile. The martingale-problem theorems live in the
semantic-universe strand\textquotesingle s auxiliary package, which is
registered with this catalogue\textquotesingle s binder: their
statements are extracted and hash-bound like every other block in
this section, so the dashboard row carries no conditional mark ---
each named profile is proven, and its theorems are bound.
\end{bremark}

\begin{bremark}[the staged premise shapes and the system bridge]
\label{rem:qsb-staged}
The rule\textquotesingle s node presents its premises per mode and per
jump: mode-indexed certificate family $(V_m)$, interior rows against
each mode\textquotesingle s flow mechanism, guard rows per jump, floors
and nonnegs per mode, and the threshold row in exact rational
arithmetic, cast once. The bound theorem for that shape
(\texttt{sound\_QSafetyBattery\_staged}) closes the gap to the flattened
hypotheses used above through an explicit bridge hypothesis: at every
in-scope state, the one-step kernel decomposes as a sub-convex mixture
of the mode\textquotesingle s flow step and the enabled guard-reset
jumps (weights supported on enabled jumps). Under that decomposition the
mode\textquotesingle s interior row and the enabled guard rows
\emph{are} the flattened generator row (\texttt{mixture\_drift}): the
one-step integral splits over the mixture, each jump term is a guard row
evaluation, and the weights sum to one. What stands left of the
concluded $\vDash$ is any system presentation whose compilation at this
profile satisfies the bridge.
\end{bremark}

\begin{bremark}[relation to the literature]\label{rem:qsb-lit}
Supermartingale safety certificates for stochastic systems go back to
Kushner \cite{kushner1967}; the modern certificate-form idiom ---
synthesize a barrier function, check polynomial rows, conclude a
probability bound --- is Prajna--Jadbabaie--Pappas
\cite{prajnajadbabaiepappas2007}. The probability engine is
Ville\textquotesingle s inequality \cite{ville1939}, with the systematic
supermartingale treatment in Doob \cite{doob1953}. The reading of the
interior rows as one-step expected decrease is the Foster--Lyapunov
drift idiom \cite{meyntweedie2009}, whose geometric form powers the
deadline sibling. What this rule adds to that literature is not
mathematics but \emph{trust structure}: every analytic ingredient lives
in a verified lemma, every polynomial ingredient is an inherited
base-calculus row replayed by the kernel, and the checker recomputes the
side conditions rather than trusting the producer.
\end{bremark}

\begin{anchors}
Lemma \ref{lem:qsb-stopped}: stopped certificate, adaptedness, barrier transport & \texttt{stopped}, \texttt{survive}, \texttt{stronglyAdapted\_stopped}, \texttt{one\_le\_stopped}, \texttt{stopped\_congr}\\
Lemma \ref{lem:qsb-bridge}: supermartingale bridge & \texttt{supermartingale\_stopped} (Markov property via \texttt{Kernel.condExp\_traj}; prefix patching via \texttt{Kernel.integral\_traj}; one-step marginal via \texttt{Kernel.map\_traj\_succ\_self})\\
Lemma \ref{lem:qsb-ville}: maximal inequality at barrier $1$ & \texttt{BernsteinLogic.Ville.reach\_avoid}, consumed at $a = 1$\\
Proposition \ref{prop:qsb-pattern}: the battery pattern & \texttt{measure\_reach\_le}\\
Lemma \ref{lem:qsb-clamp}: the clamp & the construction $g$, \texttt{hclamp}, \texttt{hgdrift} inside \texttt{sound\_QSafetyBattery\_failClosed}\\
Theorem \ref{thm:qsb-failclosed} & \texttt{sound\_QSafetyBattery\_failClosed}\\
Proposition \ref{prop:qsb-sound}: formula tie, monotonicity & \texttt{sound\_QSafetyBattery}\\
Remark \ref{rem:qsb-staged}: staged shapes, mixture bridge & \texttt{mixture\_integrable}, \texttt{mixture\_drift}, \texttt{sound\_QSafetyBattery\_staged}\\
Remark \ref{rem:qsb-profile}: PDMP sub-profile & \texttt{sound\_GSHS\_safety\_battery}\\
Remark \ref{rem:qsb-profile}: martingale-problem profile & \texttt{sound\_GSHS\_safety\_mp}, \texttt{sound\_GSHS\_safety\_battery\_mp} (\texttt{QtslAux.LayerD.MPDischarge}; hash-bound via the qtsl-aux package registration)\\
\end{anchors}


\rulesection{41}{QDeadlineBattery}{core rule}
\label{blcatrule-QDeadlineBattery}
\ruleblock{The rule (Steele metanotation)}
\sourcenote{as realized in the mechanized development; source row: \texttt{\texttt{QTSL-rules.toml}} --- \emph{QDeadlineBattery}. The metanotation statement is typeset by hand for notational quality; its mechanized ground truth is the Lean statement below, extracted from the sources at every build of this book and hash-bound. The rule \emph{roster} is machine-checked (every mechanized rule tag round-trips to exactly one catalogue row); the statement \emph{text} is editorial, checked by eye against the extracted Lean beside it.}
\begin{mathpar}
\inferrule*{
  V : \mathsf{QCertFam} \\
  R \subseteq [\![\varphi]\!] \ \text{(target formula tie)} \\
  \bigl\{D_m \smallsetminus R \vdash \mathcal{A}_m V_m \le -\theta V_m\bigr\}_{m} \ \text{(decay)} \\
  \bigl\{\mathrm{slab}_j \vdash g_j \fatsemi V_{m'_j} \le V_{m_j}\bigr\}_{j} \ \text{(guards, per jump $j$: source mode $m_j$, target $m'_j$, reset $g_j$)} \\
  \bigl\{\bar{D}_m \smallsetminus R \vdash V_m \ge 1\bigr\}_{m} \ \text{(floors, off-target only)} \\
  \bigl\{\bar{D}_m \vdash V_m \ge 0\bigr\}_{m},\ \bigl\{\bar{D}_m \vdash V_m \le M\bigr\}_{m} \ \text{(nonnegs; bound)} \\
  \vdash q \le \textstyle\sum_{k \le n} (\theta T)^k / k! \ \text{(one \rulename{TaylorBound} leaf)} \\
  0 \le \beta,\ \beta \le e^{-\theta},\ \theta \ge 0,\ 0 \le q,\ p \le 1 \ \text{(drift tie; sign gates)}
}{
  \mathcal{M},\ x_0 \in \{V \le (1-p)\,q\} \vDash \mathsf{SProbGE}\bigl(p,\ \mathsf{PDeadline}(T, \varphi, b)\bigr)
}
\end{mathpar}

\ruleflavor{rule sort baked in \texttt{theory Qtsl} (\texttt{QDeadlineBattery}; the tabular tier\textquotesingle s \texttt{QDeadlineTable} beside it) --- the display is the realized form, as the discrete-profile theorems bind it. A region left of the turnstile scopes the row ($R \vdash P$: $P$ at every point of $R$); the certificate family $V$ enters as certificate data --- a shared reference with no claim content of its own; the rows impose every claim --- never as a derivation edge. In the drift tie and sign gates, both transcendental sides ($\beta \le e^{-\theta}$ and the Taylor leaf\textquotesingle s $e^{\theta T}$ enclosure) live in the soundness lemma, never in the checker. The conclusion is the region-conditioned judgment in constructor vocabulary, as at \rulename{QSafetyBattery}; $b$ is \texttt{PDeadline}\textquotesingle s producer-chosen-horizon mark.}
\par\medskip
\rulename{QDeadlineBattery} is the bounded-reachability battery: it
concludes that with probability at least $p$ the system reaches the
target --- satisfies $\mathsf{PDeadline}(T, \varphi, b)$ (surface
syntax $\mathsf{F}_{\le T}\,\varphi$) --- from every start
in the sublevel region $\{V \le (1-p)\,q\}$. Its engine is exponential
decay: decay rows say the certificate shrinks by a factor $\beta$ per
step off the target, so the discounted process $\beta^{-t} V$ is a
supermartingale off the target, while the off-target floor $V \ge 1$
makes \emph{missing} the deadline cost at least one unit of discounted
certificate value; Markov\textquotesingle s inequality at the horizon
gives $\Pr[\mathrm{miss}] \le \beta^{T} V(x)$.

The rule\textquotesingle s division of labour between checker and lemma
is its most instructive feature. The checker replays only exact rational
comparisons: the one arithmetic leaf is the Taylor comparison
$q \le \sum_{k \le n} (\theta T)^k / k!$ --- a closed, wire-free
inequality discharged by the ground-fact rule --- while both
transcendental facts ($\sum_{k \le n} (\theta T)^k/k! \le e^{\theta T}$
and $e^{-\theta} e^{\theta} = 1$) live inside the soundness lemma. The
checker never needs to know $e$ exists.

What the checker recomputes beyond that leaf: the domain-hull and
fail-closed side conditions (as in the safety battery), and the
threshold arithmetic in the exact rational vocabulary
($\mathrm{thr}_{\mathsf F} = \mathrm{thr} \cdot q$ atop
$p + \mathrm{thr} \le 1$). Soundness is independent of the producer
knobs $\theta$ and $q$: any discrete rate $0 \le \beta \le e^{-\theta}$
and any $q$ below the Taylor sum are admissible.

Each hypothesis is load-bearing. Drop the off-target floor and a
certificate identically $0$ would certify certain reachability of an
unreachable target. Drop the boundedness row and the decay row no longer transfers to the
clamped certificate --- the discount multiplies the certificate by
$\beta^{-t}$, so the ceiling must be invisible wherever one-step mass
reaches, which is exactly what $V \le M$ buys (and, in the mechanized
engine, the integrability of the discounted certificate rides the same
bound). Drop the
fail-closed containment and escape from the modelled domain would count
as success: a reach guarantee must show the escape event \emph{null},
where the safety battery could simply enlarge its stop set (leaving
counts as failing is conservative for safety, but would be unsound
generosity for reachability). Drop the threshold arithmetic and the
concluded region over-budgets. The proof section walks each.

The dashboard mark states the status: the hash-bound theorems warrant
the rule at the discrete-kernel profile (geometric drift,
$\beta = e^{-\theta}$ discretized); the continuous-time profile is
proven at the martingale-problem presentation --- the discounted
supermartingale built from the undiscounted martingale-problem
hypothesis by the exponential-discount construction, stopped at the
d\'ebut by the right-continuous optional-sampling library theorem
(\texttt{sound\_GSHS\_deadline\_mp},
\texttt{sound\_GSHS\_deadline\_battery\_mp},
\texttt{QtslAux.LayerD.MPDischarge}) --- hash-bound below via the
qtsl-aux package registration, so the dashboard row carries no
conditional mark: each named profile is proven, and its theorems are
bound.

At the general stratum, over the semantic universe\textquotesingle s
record carrier (the semantic-universe functor package, hash-bound
below), the deadline conclusion is a \emph{mark} --- an
$\exists$-face spec: the per-window family of acceptance-event
demands is persistent, and satisfaction forces almost-sure acceptance
at the window (\texttt{satAtM\_acceptanceSpecF\_ae}). The backward
direction is deliberately absent, and the absence is the semantics: a
record that reaches its target only in a superwindow is invisible to
the subwindow, so a subwindow\textquotesingle s reach demand is a
genuinely stronger demand --- a deadline --- and extension is exactly
how a met deadline persists to every superwindow. The measurability
the mark face charges is issued by the site\textquotesingle s
topology rather than assumed: hit-events commute with countable
unions in the set argument (\texttt{hitEvent\_iUnion}), turning grid
decompositions into event decompositions, and a compact target
cell\textquotesingle s visit event over a $\sigma$-compact stage
window is Borel
(\texttt{measurableSet\_hitEvent\_univ\_prod\_of\_isCompact}). One
typing note rides the transposition: the site-general grids are
indexed by bare naturals; at clocked continuum strata the legs
re-type to rational-duration stage cuts --- the tick-indexed reading
of a windowed horizon is a forcing of the finitary record site, not
of the grammar.

\cataloguenote{The checked floor is \emph{off-target only}: on the target the
certificate may take any value, since success needs no budget. This is
deliberately weaker than a global floor, and the proven statement binds
the weaker form; a catalogued rendering that states a global floor is
sound but relaxable to match. The concluded threshold $(1-p)\,q$ is the
compile\textquotesingle s choice; the underlying theorem consumes only
the single inequality $\beta^{T} \cdot \mathrm{thr} \le 1 - p$. At this
profile, guard rows land inside the one-step decay row (the kernel
carries flow and jumps together); the instantaneous-jump reading, where
the discount factor does not tick during a jump, is carried at the
continuous-time profile by the martingale-problem
discharge\textquotesingle s no-alarm segment assembly.
}
\ruleblock{Lean statement}
\declhead{BernsteinLogicAux.Battery}{sound\_QDeadline}{\badge{BadgeGreen}{statement bound}}
\begin{leansrc}
/-- **`sound_QDeadline`, discrete profile** (T6): over the target-hit form, the deadline
formula tie — the target region *under*-approximates the truth set of the
`𝖥_{≤T}`-argument — weakens the concluded event to the truth of the reach formula:
`Pr_x[∃ t ≤ T, X t ⊨ φ] ≥ p` on the concluded region `D ∩ {V ≤ thr}`. No measurability
is asked of the satisfaction predicate (`measure_mono` on the measurable hit event). -/
theorem sound_QDeadline {T : ℕ} {M : ℝ} {sat : S → Prop}
    (hD : MeasurableSet D) (hDbar : MeasurableSet Dbar) (hB : MeasurableSet B)
    (hV : Measurable V)
    (hF : ∀ y ∈ B, sat y)
    (hβ : 0 ≤ β)
    (hΘF : β ^ T * thr ≤ 1 - p)
    (hDsub : D ⊆ Dbar)
    (hR3 : ∀ y ∈ D, y ∉ B → K y Dbar = 1)
    (hR4 : ∀ y ∈ D, y ∉ B → K y (D ∪ B) = 1)
    (hint : ∀ y ∈ D, y ∉ B → Integrable V (K y))
    (hdecay : ∀ y ∈ D, y ∉ B → ∫ z, V z ∂K y ≤ β * V y)
    (hfloor : ∀ y ∈ Dbar, y ∉ B → 1 ≤ V y)
    (hnonneg : ∀ y ∈ Dbar, 0 ≤ V y)
    (hM : ∀ y ∈ Dbar, V y ≤ M)
    (hx : x ∈ D) (hthr : V x ≤ thr) :
    ENNReal.ofReal p ≤ pathMeasure K x {ω | ∃ k ≤ T, sat (ω k)}
\end{leansrc}
\declhead{BernsteinLogicAux.Battery}{sound\_QDeadline\_hit}{\badge{BadgeGreen}{statement bound}}
\begin{leansrc}
/-- **`sound_QDeadline`, discrete profile, target-hit form** (T6 — the safety battery's
deadline successor — at the discrete-kernel profile). Hypotheses are the
decay rows (`∫ V ∂K y ≤ β·V y` off the target `B`, the discrete `𝓛V ≤ -θV` with
`β = e^{-θ}`), floors **off** the target, nonnegs, a certificate bound `M` on the cover,
side conditions (R3)/(R4), and the deadline threshold arithmetic
`(Θ_F): β ^ T * thr ≤ 1 - p` — the inequality form of the compile's `thr = (1-p)·q` with
`q ≤ e^{θT}`. Conclusion: from any start in `D ∩ {V ≤ thr}`, the chain visits `B` by time
`T` with probability at least `p`. Consumes the landed geometric-drift engine
`BernsteinLogic.Deadline.deadline_tail` at the enlarged target `B ∪ Dᶜ`,
then removes the escape branch with `ae_no_escape`. -/
theorem sound_QDeadline_hit {T : ℕ} {M : ℝ}
    (hD : MeasurableSet D) (hDbar : MeasurableSet Dbar) (hB : MeasurableSet B)
    (hV : Measurable V)
    (hβ : 0 ≤ β)
    (hΘF : β ^ T * thr ≤ 1 - p)
    (hDsub : D ⊆ Dbar)
    (hR3 : ∀ y ∈ D, y ∉ B → K y Dbar = 1)
    (hR4 : ∀ y ∈ D, y ∉ B → K y (D ∪ B) = 1)
    (hint : ∀ y ∈ D, y ∉ B → Integrable V (K y))
    (hdecay : ∀ y ∈ D, y ∉ B → ∫ z, V z ∂K y ≤ β * V y)
    (hfloor : ∀ y ∈ Dbar, y ∉ B → 1 ≤ V y)
    (hnonneg : ∀ y ∈ Dbar, 0 ≤ V y)
    (hM : ∀ y ∈ Dbar, V y ≤ M)
    (hx : x ∈ D) (hthr : V x ≤ thr) :
    ENNReal.ofReal p ≤ pathMeasure K x {ω | ∃ k ≤ T, ω k ∈ B}
\end{leansrc}
\par\medskip
The catalogue form: the $(\theta, q)$ parameterization, concluding on
$\{V \le (1-p)\,q\}$ from the Taylor leaf and the drift tie
$\beta \le e^{-\theta}$, with both transcendental facts inside the
lemma. It is a wrapper over \texttt{sound\_QDeadline}, which carries the
general inequality form $\beta^{T} \cdot \mathrm{thr} \le 1 - p$.

\declhead{BernsteinLogicAux.QtslBattery}{sound\_QDeadlineBattery}{\badge{BadgeGreen}{statement bound}}
\begin{leansrc}
/-- **`sound_QDeadlineBattery`** (QTSL-rules.toml row `QDeadlineBattery` —
the catalogue's `(θ, q)` parameterization of T6, discrete-kernel profile).
The concluded region is `{V ≤ (1 − p)·q}`: the deadline threshold
`thr_F = thr · q` in the staged `q_mul` vocabulary atop `thr = 1 − p`.
Its arithmetic premises are

* the **T7 Taylor leaf** `q ≤ ∑_{k ≤ n} (θT)^k / k!` — the one closed
  wire-free comparison the checker replays in exact ℚ (rule row
  `TaylorBound (T7)`, alias → `PConstExact`), and
* the **drift tie** `β ≤ e^{−θ}` — the producer's discrete geometric rate
  matching the catalogue's `β = e^{−θ}` (any smaller rate is sound).

The transcendental facts — `∑_{k ≤ n} (θT)^k/k! ≤ e^{θT}`
(`Real.sum_le_exp_of_nonneg`) and `e^{−θ}·e^{θ} = 1` — live in this lemma,
never in the checker (the catalogue's "the checker never needs to know `e`
exists"). Wrapper over `sound_QDeadline`, which carries the general
inequality form `β ^ T · thr ≤ 1 − p`.

System bridge: for any System whose chart compilation at
the discrete-kernel profile yields `K`, this is the concluded
`QRegionSat`'s reading on that System's records. -/
theorem sound_QDeadlineBattery {n : ℕ} {sat : S → Prop}
    (hD : MeasurableSet D) (hDbar : MeasurableSet Dbar) (hB : MeasurableSet B)
    (hV : Measurable V)
    (hF : ∀ y ∈ B, sat y)
    (hβ : 0 ≤ β) (hβθ : β ≤ Real.exp (-θ)) (hθ : 0 ≤ θ)
    (hq0 : 0 ≤ q)
    (hq : q ≤ ∑ k ∈ range (n + 1), (θ * T) ^ k / (Nat.factorial k : ℝ))
    (hp : p ≤ 1)
    (hDsub : D ⊆ Dbar)
    (hR3 : ∀ y ∈ D, y ∉ B → K y Dbar = 1)
    (hR4 : ∀ y ∈ D, y ∉ B → K y (D ∪ B) = 1)
    (hint : ∀ y ∈ D, y ∉ B → Integrable V (K y))
    (hdecay : ∀ y ∈ D, y ∉ B → ∫ z, V z ∂K y ≤ β * V y)
    (hfloor : ∀ y ∈ Dbar, y ∉ B → 1 ≤ V y)
    (hnonneg : ∀ y ∈ Dbar, 0 ≤ V y)
    (hM : ∀ y ∈ Dbar, V y ≤ M)
    (hx : x ∈ D) (hthr : V x ≤ (1 - p) * q) :
    ENNReal.ofReal p ≤ pathMeasure K x {ω | ∃ k ≤ T, sat (ω k)}
\end{leansrc}
\par\medskip
\textbf{The general-stochastic instance, assembled.}\enspace The
deadline bound on the timed skeleton of the segment model: the
$\theta$-discounted generator rows enter per segment, the clock is a
coordinate of the timed chain, and the conclusion is the displayed
deadline-overrun bound on that chain\textquotesingle s events ---
headline, off-target, and battery forms.

\declhead{BernsteinLogicAux.LayerD.Deadline}{sound\_GSHS\_deadline}{\badge{BadgeGreen}{statement bound}}
\begin{leansrc}
/-- **The GSHS deadline tail bound** (T6, headline form). Inherited
premises, stated (the leg's discipline — never silent):

* `hSegD : SegmentDynkin M θ V` at the discount `θ` — **the θ > 0 premise**:
  the B1 row compiled at positive discount; the
  statement is typed at `0 ≤ θ` and degenerates to the trivial floor bound
  at `θ = 0`.
* `hFloor : ∀ y, ε ≤ V y` — **the F7 floor premise**: the
  positive floor on the working region, conclusion scaled by `1/ε`.
  Without it the conclusion is unsound (the review's two-point witness: `V ≡ 0`
  off-target satisfies every remaining hypothesis while the deadline fails
  surely).

Conclusion: the probability that the timed skeleton chain ever presents a
value stamped past the deadline `T` is at most `e^{−θT} V(x)/ε` — the T6
display `Pr[τ_B > T] ≤ e^{−θT} V(x₀)/ε` in the value-presentation reading:
instantiated at the stopped presentation `Rep(P^{†B})` with halts into the
target ("the stop IS the target" — a design choice
made here and visible), a path outside this event either reached the
target by `T` or stalled (Zeno / vanished — the residues emptied a.s. for
in-scope models by the prose W-Z-freedom theorem).
Consumes `Battery.measure_reach_le` (= `Ville.reach_avoid`) over the
timed chain: the deadline battery rides literally the same engine as T5,
with the capped certificate `min(e^{θs}V, εe^{θT})/εe^{θT}` — and no
global bound `V ≤ C` is demanded (the cap manufactures its own). -/
theorem sound_GSHS_deadline [IsMarkovKernel M.seg] [IsMarkovKernel M.jump]
    {V : State → ℝ} {θ ε : ℝ} {T : ℝ≥0}
    (hV : Measurable V) (hθ : 0 ≤ θ)
    (hSegD : SegmentDynkin M θ V) (hForce : ForcedNonexpansion M V)
    (hNonneg : NonnegV V) (hε : 0 < ε) (hFloor : ∀ y, ε ≤ V y) (x : State) :
    Battery.pathMeasure (tskel M) (TSkelState.alive (0, x))
        {ω | ∃ n, T < tclock (ω n)}
      ≤ ENNReal.ofReal (Real.exp (-(θ * T)) * V x / ε)
\end{leansrc}
\declhead{BernsteinLogicAux.LayerD.Deadline}{sound\_GSHS\_deadline\_offTarget}{\badge{BadgeGreen}{statement bound}}
\begin{leansrc}
/-- **T6, off-target-floor variant** (the sharper premise, the narrower
event): the floor is demanded off the target only — the F7 floor premise
read narrowly — and the conclusion event asks for a past-deadline
presentation *outside* `B`. What this variant gives up is the crossing-jump
corner: a path whose first past-`T` presentation is the target hit itself
(the dispatch that crossed the deadline landed in `B`) misses the deadline
without entering this event — `exceed_cases` makes the dichotomy explicit.
The headline `sound_GSHS_deadline` closes the corner by flooring the whole
working region instead. -/
theorem sound_GSHS_deadline_offTarget [IsMarkovKernel M.seg] [IsMarkovKernel M.jump]
    {V : State → ℝ} {θ ε : ℝ} {T : ℝ≥0} {B : Set State}
    (hV : Measurable V) (hθ : 0 ≤ θ) (hB : MeasurableSet B)
    (hSegD : SegmentDynkin M θ V) (hForce : ForcedNonexpansion M V)
    (hNonneg : NonnegV V) (hε : 0 < ε) (hFloorB : ∀ y ∉ B, ε ≤ V y)
    (x : State) :
    Battery.pathMeasure (tskel M) (TSkelState.alive (0, x))
        {ω | ∃ n, T < tclock (ω n) ∧ ω n ∉ tbadIn B}
      ≤ ENNReal.ofReal (Real.exp (-(θ * T)) * V x / ε)
\end{leansrc}
\declhead{BernsteinLogicAux.LayerD.Deadline}{sound\_GSHS\_deadline\_battery}{\badge{BadgeGreen}{statement bound}}
\begin{leansrc}
/-- **The deadline battery form** (T6 as `sound_QDeadline`'s GSHS instance,
complement arithmetic): with the threshold accounting
`p + e^{−θT}·thr/ε ≤ 1` and a start under the threshold, the chain's
presented clock stays `≤ T` at every index with probability at least `p`.
On a stopped-at-target instantiation with halts into the target, that event
reads *reached the target by the deadline* — up to the two stall residues
(forced-jump Zeno accumulation; vanishing into a never-terminating
segment), which are inside the event, named, and emptied a.s. for in-scope
models by the prose W-Z-freedom theorem: the honest
one-sided residue, visible rather than erased. -/
theorem sound_GSHS_deadline_battery [IsMarkovKernel M.seg] [IsMarkovKernel M.jump]
    {V : State → ℝ} {θ ε p thr : ℝ} {T : ℝ≥0}
    (hV : Measurable V) (hθ : 0 ≤ θ)
    (hSegD : SegmentDynkin M θ V) (hForce : ForcedNonexpansion M V)
    (hNonneg : NonnegV V) (hε : 0 < ε) (hFloor : ∀ y, ε ≤ V y)
    (hΘ : p + Real.exp (-(θ * T)) * thr / ε ≤ 1) (x : State) (hthr : V x ≤ thr) :
    ENNReal.ofReal p ≤ Battery.pathMeasure (tskel M) (TSkelState.alive (0, x))
      {ω | ∀ n, tclock (ω n) ≤ T}
\end{leansrc}
\par\medskip
\textbf{The martingale-problem profile, end to end.}\enspace The
general continuous-time theorems: the discounted supermartingale is
built from the undiscounted martingale-problem hypothesis by the
exponential-discount construction, stopped at the d\'ebut by the
right-continuous optional-sampling library theorem, and the deadline
conclusion follows for every system whose presentation satisfies the
martingale-problem hypothesis.

\declhead{QtslAux.LayerD.MPDischarge}{sound\_GSHS\_deadline\_mp}{\badge{BadgeGreen}{statement bound}}
\begin{leansrc}
/-- The deadline reach bound at the martingale-problem profile. -/
theorem sound_GSHS_deadline_mp [IsMarkovKernel M.seg] [IsMarkovKernel M.jump]
    {V : State → ℝ} {θ ε : ℝ} {T : ℝ≥0}
    (hV : Measurable V) (hθ : 0 ≤ θ)
    (hmp : ∀ x : State, Nonempty (MPFlow M θ V x))
    (hForce : ForcedNonexpansion M V)
    (hNonneg : NonnegV V) (hε : 0 < ε) (hFloor : ∀ y, ε ≤ V y) (x : State) :
    Battery.pathMeasure (tskel M) (TSkelState.alive (0, x))
        {ω | ∃ n, T < tclock (ω n)}
      ≤ ENNReal.ofReal (Real.exp (-(θ * T)) * V x / ε)
\end{leansrc}
\declhead{QtslAux.LayerD.MPDischarge}{sound\_GSHS\_deadline\_battery\_mp}{\badge{BadgeGreen}{statement bound}}
\begin{leansrc}
/-- The deadline battery at the martingale-problem profile. -/
theorem sound_GSHS_deadline_battery_mp
    [IsMarkovKernel M.seg] [IsMarkovKernel M.jump]
    {V : State → ℝ} {θ ε p thr : ℝ} {T : ℝ≥0}
    (hV : Measurable V) (hθ : 0 ≤ θ)
    (hmp : ∀ x : State, Nonempty (MPFlow M θ V x))
    (hForce : ForcedNonexpansion M V)
    (hNonneg : NonnegV V) (hε : 0 < ε) (hFloor : ∀ y, ε ≤ V y)
    (hΘ : p + Real.exp (-(θ * T)) * thr / ε ≤ 1) (x : State)
    (hthr : V x ≤ thr) :
    ENNReal.ofReal p ≤ Battery.pathMeasure (tskel M) (TSkelState.alive (0, x))
      {ω | ∀ n, tclock (ω n) ≤ T}
\end{leansrc}
\par\medskip
\textbf{The general stratum: grid tickets by topology.}\enspace The
countable-combination engine and the reach leg\textquotesingle s
measurability ticket over the record carrier: hitting commutes with
unions in the set argument, and a compact target cell\textquotesingle s
visit event over a $\sigma$-compact stage window is Borel --- issued
by the site\textquotesingle s topology, not assumed.

\declhead{SuFunctor.Site.FellAcceptance}{hitEvent\_iUnion}{\badge{BadgeGreen}{statement bound}}
\begin{leansrc}
/-- **Hitting commutes with unions in the set argument** — the engine
that turns countable set decompositions into countable event
decompositions.  Arbitrary index, pure set algebra. -/
theorem hitEvent_iUnion {ι : Sort*} (A : ι → Set α) :
    hitEvent (⋃ i, A i) = ⋃ i, hitEvent (A i)
\end{leansrc}
\declhead{SuFunctor.Site.FellAcceptance}{measurableSet\_hitEvent\_iUnion}{\badge{BadgeGreen}{statement bound}}
\begin{leansrc}
/-- Hit events pass countable unions of measurable pieces. -/
theorem measurableSet_hitEvent_iUnion {ι : Sort*} [Countable ι]
    {A : ι → Set α} (h : ∀ i, MeasurableSet (hitEvent (A i))) :
    MeasurableSet (hitEvent (⋃ i, A i))
\end{leansrc}
\declhead{SuFunctor.Site.FellAcceptance}{measurableSet\_hitEvent\_univ\_prod\_of\_isCompact}{\badge{BadgeGreen}{statement bound}}
\begin{leansrc}
/-- **Hit-events of certified cells are Effros-measurable for free**
(the memo's headline, cashed): a compact cell's visit event over a
σ-compact stage window is a countable union of Fell-closed hit
events. -/
theorem measurableSet_hitEvent_univ_prod_of_isCompact
    [SigmaCompactSpace ↥W] {R : Set X} (hR : IsCompact R) :
    MeasurableSet
      (hitEvent ((Set.univ : Set ↥W) ×ˢ R) : Set (FTraj W X))
\end{leansrc}
\par\medskip
\textbf{Marks over the record carrier.}\enspace The deadline
conclusion\textquotesingle s general-stratum shape: acceptance-event
spec families, persistent by window-calculus transitivity, with the
forward judgment lemma; the backward direction is deliberately absent
--- deadline semantics, as the commentary states.

\declhead{SuFunctor.Site.FellBehaviorSpec}{acceptanceSpecF}{\badge{BadgeGreen}{statement bound}}
\begin{leansrc}
/-- The spec family of a per-window acceptance event (deadline, until,
Streett): clip demands of the events' indicators. -/
def acceptanceSpecF (U : ι → Set T)
    (E : ∀ i, Set (FTraj (U i) X)) (j : ι) :
    Set (MeasPen (FTraj (U j) X))
\end{leansrc}
\declhead{SuFunctor.Site.FellBehaviorSpec}{isPersistentMF\_acceptanceSpecF}{\badge{BadgeGreen}{statement bound}}
\begin{leansrc}
theorem isPersistentMF_acceptanceSpecF (U : ι → Set T)
    (hU : ∀ i, IsOpen (U i)) (E : ∀ i, Set (FTraj (U i) X)) :
    IsPersistentMF U (acceptanceSpecF U E)
\end{leansrc}
\declhead{SuFunctor.Site.FellBehaviorSpec}{satAtM\_acceptanceSpecF\_ae}{\badge{BadgeGreen}{statement bound}}
\begin{leansrc}
/-- The forward judgment: satisfying an acceptance spec at a window
forces almost-sure acceptance there.  The backward direction is
deliberately absent — liveness events do not clip-restrict, so
subwindow demands are genuinely stronger (a reach demand within a
subwindow is a deadline); the module docstring carries the reading. -/
theorem satAtM_acceptanceSpecF_ae (U : ι → Set T)
    (E : ∀ i, Set (FTraj (U i) X)) (j : ι)
    {μ : MeasJoint (FTraj (U j) X)}
    (h : SatAtM μ (acceptanceSpecF U E j)) :
    (μ.1 : Measure (FTraj (U j) X)) (E j)ᶜ = 0
\end{leansrc}
\ruleblock{Correspondence (Steele display $\leftrightarrow$ Lean)}
\badge{BadgeAmber}{a bound Lean statement CHANGED since certification --- re-review and re-certify}
\sourcenote{The dictionary below is a reviewer-certified statement of how the mechanized statements above interpret the displayed rule: display metavariable $\mapsto$ Lean binder, each display hypothesis $\mapsto$ the premise or side condition carrying it, the display conclusion $\mapsto$ the goal --- with any deliberate difference in width between display and mechanization named outright. Certification is recorded as a hash triple over (the displayed cell, the extracted statement texts, this dictionary\textquotesingle s own text); if any of the three changes, this block turns amber until a reviewer re-certifies. The correspondence claim itself remains editorial --- reviewer-grade, not machine-checked; its machine-checked successor types each soundness statement against the rule as a structured object, and supersedes this block per rule as it lands.}
Model frame and profile as at \rulename{QSafetyBattery}: hash-bound at
the discrete-kernel profile (the amended cell\textquotesingle s own
caveat), with the continuous-time profile proven at the
martingale-problem presentation (hash-bound below, per the profile
remark). The catalogue form binds
\texttt{sound\_QDeadlineBattery}; \texttt{sound\_QDeadline} carries
the general inequality it wraps.
\begin{itemize}
\item the displayed judgment $\mathcal{M}
  ,\ x_0 \in \{V \le (1-p)\,q\} \vDash
  \mathsf{SProbGE}\bigl(p,\ \mathsf{PDeadline}(T, \varphi, b)\bigr)$
  $\mapsto$ the goal
  $\mathrm{ofReal}\ p \le$ \texttt{pathMeasure K x \{ω | ∃ k ≤ T, sat
  (ω k)\}} from every start with \texttt{x ∈ D} and \texttt{V x ≤ (1
  − p)·q}; the target tie $\mapsto$ \texttt{hF} $:$ $\forall y \in B,
  \texttt{sat}\ y$ --- $B$ \emph{under}-approximates the target truth
  set (the polarity opposite to \rulename{QSafetyBattery}\textquotesingle
  s avoid tie, both safe directions), and again no measurability is
  asked of \texttt{sat}.
\item decay rows $\mapsto$ \texttt{hdecay} (one-step $\int V \le
  \beta V$ on $D \smallsetminus B$; at this profile guard rows land
  inside it, per the note) with integrability rider \texttt{hint};
  floors \emph{off-target only} $\mapsto$ \texttt{hfloor} (on
  $\bar{D} \smallsetminus B$ --- the display\textquotesingle s
  weaker, deliberately relaxable form is exactly what is bound);
  nonnegs $\mapsto$ \texttt{hnonneg}; the bound row $V \le M$
  $\mapsto$ \texttt{hM}.
\item the TaylorBound leaf $q \le \sum_{k \le n} (\theta T)^k / k!$
  $\mapsto$ \texttt{hq}, verbatim --- the one closed comparison a
  checker replays in exact $\Q$ --- and the drift tie $\beta \le
  e^{-\theta}$ $\mapsto$ \texttt{hβθ}, with both transcendental
  facts (\texttt{Real.sum\_le\_exp\_of\_nonneg}; $e^{-\theta}
  e^{\theta} = 1$) discharged \emph{inside} the lemma: the certified
  statement is the witness that the checker never needs to know $e$
  exists.
\item sign gates $0 \le \beta$, $\theta \ge 0$, $0 \le q$, $p \le 1$
  $\mapsto$ \texttt{hβ}/\texttt{hθ}/\texttt{hq0}/\texttt{hp}, verbatim
  --- the base sign gate \texttt{hβ} is consumed by the power-order
  step; side conditions and region
  measurability as at
  \rulename{QSafetyBattery} (\texttt{hR3}/\texttt{hR4};
  \texttt{hD}/\texttt{hDbar}/\texttt{hB}/\texttt{hV}, the
  undisplayed regularity pack).
\end{itemize}
Width note: the wrapped \texttt{sound\_QDeadline} consumes only
$\beta^T \cdot \mathrm{thr} \le 1 - p$ --- every $(\theta, q)$
satisfying the displayed leaves yields it, so the catalogue
parameterization is an instance of a strictly wider certified core;
producer knobs, exactly as the roster note says.

\ruleblock{Lean verifier receipts}
\sourcenote{module \texttt{BernsteinLogicAux.Battery}; bound build \texttt{2fd1c031c80b\ldots} (bound from receipt 8d9ccb68993ea0ec, 2026-08-04T01:56:55Z). A receipt records that named checking layers accepted a build: (module, compiled-artifact SHA-256, environment SHA-256, checker layers), minted when the checking gate ran; the store receipts attesting the bound build are listed with the layers each run credited. Nothing here asks the reader to trust the store: every layer is an independently runnable program over the quoted sources --- rebuild the module and re-run the checkers to reproduce the verdict. Receipts attesting superseded builds credit no layer and are omitted from this rendering (the count of omissions, if any, is noted below).}
\begin{receipttable}
\receiptrow{8d9ccb68993ea0ec}{2026-08-04T01:56:55Z}{build+decorrelated}{\badge{BadgeGreen}{current}}
\receiptrow{5075586e4813a7c5}{2026-08-04T03:11:30Z}{build+decorrelated}{\badge{BadgeGreen}{current}}
\end{receipttable}
\sourcenote{13 further receipt(s) in the store attest superseded builds of this module and are omitted from this rendering.}
\sourcenote{module \texttt{BernsteinLogicAux.QtslBattery}; bound build \texttt{5c75f9923f27\ldots} (bound from receipt 3fd830b4914bdeda, 2026-08-04T01:56:58Z). A receipt records that named checking layers accepted a build: (module, compiled-artifact SHA-256, environment SHA-256, checker layers), minted when the checking gate ran; the store receipts attesting the bound build are listed with the layers each run credited. Nothing here asks the reader to trust the store: every layer is an independently runnable program over the quoted sources --- rebuild the module and re-run the checkers to reproduce the verdict. Receipts attesting superseded builds credit no layer and are omitted from this rendering (the count of omissions, if any, is noted below).}
\begin{receipttable}
\receiptrow{3fd830b4914bdeda}{2026-08-04T01:56:58Z}{build+decorrelated}{\badge{BadgeGreen}{current}}
\receiptrow{96b18d8572451f65}{2026-08-04T03:11:31Z}{build+decorrelated}{\badge{BadgeGreen}{current}}
\end{receipttable}
\sourcenote{13 further receipt(s) in the store attest superseded builds of this module and are omitted from this rendering.}
\sourcenote{module \texttt{BernsteinLogicAux.LayerD.Deadline}; bound build \texttt{ae6767f0b59d\ldots} (bound from receipt 1c87bcb62534ed03, 2026-08-04T01:57:04Z). A receipt records that named checking layers accepted a build: (module, compiled-artifact SHA-256, environment SHA-256, checker layers), minted when the checking gate ran; the store receipts attesting the bound build are listed with the layers each run credited. Nothing here asks the reader to trust the store: every layer is an independently runnable program over the quoted sources --- rebuild the module and re-run the checkers to reproduce the verdict. Receipts attesting superseded builds credit no layer and are omitted from this rendering (the count of omissions, if any, is noted below).}
\begin{receipttable}
\receiptrow{1c87bcb62534ed03}{2026-08-04T01:57:04Z}{build+decorrelated}{\badge{BadgeGreen}{current}}
\receiptrow{ea6411e23e04ba3c}{2026-08-04T03:11:30Z}{build+decorrelated}{\badge{BadgeGreen}{current}}
\end{receipttable}
\sourcenote{14 further receipt(s) in the store attest superseded builds of this module and are omitted from this rendering.}
\sourcenote{module \texttt{QtslAux.LayerD.MPDischarge}; bound build \texttt{baded74dd943\ldots} (bound from receipt 7228834dea10e5f9, 2026-07-24T11:04:04Z). A receipt records that named checking layers accepted a build: (module, compiled-artifact SHA-256, environment SHA-256, checker layers), minted when the checking gate ran; the store receipts attesting the bound build are listed with the layers each run credited. Nothing here asks the reader to trust the store: every layer is an independently runnable program over the quoted sources --- rebuild the module and re-run the checkers to reproduce the verdict. Receipts attesting superseded builds credit no layer and are omitted from this rendering (the count of omissions, if any, is noted below).}
\begin{receipttable}
\receiptrow{7228834dea10e5f9}{2026-07-24T11:04:04Z}{build+decorrelated}{\badge{BadgeGreen}{current}}
\end{receipttable}
\sourcenote{module \texttt{SuFunctor.Site.FellAcceptance}; bound build \texttt{651fb557bbfe\ldots} (bound from receipt aa2169560f8f571a, 2026-08-03T20:19:04Z). A receipt records that named checking layers accepted a build: (module, compiled-artifact SHA-256, environment SHA-256, checker layers), minted when the checking gate ran; the store receipts attesting the bound build are listed with the layers each run credited. Nothing here asks the reader to trust the store: every layer is an independently runnable program over the quoted sources --- rebuild the module and re-run the checkers to reproduce the verdict. Receipts attesting superseded builds credit no layer and are omitted from this rendering (the count of omissions, if any, is noted below).}
\begin{receipttable}
\receiptrow{7ad85720de6e419e}{2026-08-03T10:13:50Z}{build+decorrelated}{\badge{BadgeGreen}{current}}
\receiptrow{94a408109996d1bf}{2026-08-03T13:41:21Z}{build+decorrelated}{\badge{BadgeGreen}{current}}
\receiptrow{580a7e551fab03b6}{2026-08-03T15:39:44Z}{build+decorrelated}{\badge{BadgeGreen}{current}}
\receiptrow{aa2169560f8f571a}{2026-08-03T20:19:04Z}{build+decorrelated}{\badge{BadgeGreen}{current}}
\end{receipttable}
\sourcenote{2 further receipt(s) in the store attest superseded builds of this module and are omitted from this rendering.}
\sourcenote{module \texttt{SuFunctor.Site.FellBehaviorSpec}; bound build \texttt{345f9caf9453\ldots} (bound from receipt 72f2c3b62b6771fc, 2026-08-03T20:19:08Z). A receipt records that named checking layers accepted a build: (module, compiled-artifact SHA-256, environment SHA-256, checker layers), minted when the checking gate ran; the store receipts attesting the bound build are listed with the layers each run credited. Nothing here asks the reader to trust the store: every layer is an independently runnable program over the quoted sources --- rebuild the module and re-run the checkers to reproduce the verdict. Receipts attesting superseded builds credit no layer and are omitted from this rendering (the count of omissions, if any, is noted below).}
\begin{receipttable}
\receiptrow{d7b846e785c1d7fb}{2026-08-03T10:13:53Z}{build+decorrelated}{\badge{BadgeGreen}{current}}
\receiptrow{c0e3dc15419765aa}{2026-08-03T13:41:24Z}{build+decorrelated}{\badge{BadgeGreen}{current}}
\receiptrow{9a62b29c08b63b4f}{2026-08-03T15:39:47Z}{build+decorrelated}{\badge{BadgeGreen}{current}}
\receiptrow{72f2c3b62b6771fc}{2026-08-03T20:19:08Z}{build+decorrelated}{\badge{BadgeGreen}{current}}
\end{receipttable}
\sourcenote{2 further receipt(s) in the store attest superseded builds of this module and are omitted from this rendering.}
\layerline{\layeritem{build}{lake build (elaboration)}{\badge{BadgeGreen}{current}}\layeritem{replay}{leanchecker (official kernel replay)}{\badge{BadgeAmber}{stale only}}\layeritem{hardened}{safe\_verify (hardened screen)}{\badge{BadgeAmber}{stale only}}\layeritem{decorrelated}{lean4export + nanoda (decorrelated kernel)}{\badge{BadgeGreen}{current}}}
\sourcenote{\badge{BadgeGray}{comparator: none} no fixture in the three-implementation comparator harness (compiled / kernel-interpreted / reference-interpreted) exercises this rule.}
\ruleblock{HOL4 statement}
\absentbox{NO HOL STATEMENT YET}{The HOL4 development ports the base calculus\textquotesingle s checker
chain (theory \texttt{bern\_catalogue}, developed side by side with the
base catalogue sources, stopping below the proof-object layer); it
contains no probability or measure theory.
\rulename{QDeadlineBattery} lives one theory-extension edge above the
base calculus and its soundness content is path-space measure theory
(discounted supermartingales, optional stopping, an escape-null
argument), for which no counterpart exists anywhere in the ported
development. A HOL statement of this rule would be a new
measure-theoretic campaign, not an extension of the ported checker
chain; until one exists, this box is the honest state of the edifice.
}
\ruleblock{HOL verifier receipts}
\absentbox{NONE}{no HOL verification receipt exists for this rule.}
\ruleblock{Proof}

\sourcenote{Genre: core rule (tier 1) --- an ordinary mathematical proof
at Putnam-exam-solution detail. Substrate, stopped-process machinery,
and the battery pattern are stated with \rulename{QSafetyBattery}
(Lemmas \ref{lem:qsb-stopped}--\ref{lem:qsb-ville},
Proposition~\ref{prop:qsb-pattern}) and cited here by reference. The
theorems quoted above warrant the rule at the \textbf{discrete-kernel
profile}; the profile remark at the end states what remains a
hypothesis beyond it.}

\textbf{Setting and claim.}\enspace The substrate of
\rulename{QSafetyBattery}: state space $S$, one-step Markov kernel $K$,
path measure $\Pr_x$, canonical filtration. The rule\textquotesingle s
data: modelled domain $D$, hull $\bar{D} \supseteq D$, a \emph{target}
$B$ compiled to \emph{under}-approximate the truth set of the
$\mathsf{F}_{\le T}$-argument --- note the orientation flip against the
safety battery\textquotesingle s avoid set: the formula tie here is
(F$'$)\ $\forall y \in B,\ \mathrm{sat}(y)$ --- a measurable certificate
$V$, a discrete decay factor $\beta \ge 0$, a horizon $T \in \N$, and
$p, \mathrm{thr}$. The checked rows, read semantically:
\begin{itemize}
\item[(D1)] \emph{decay}: for $y \in D \smallsetminus B$, $V$ integrable
  against $K\,y$ and $\int V \,dK\,y \le \beta\, V(y)$ --- the discrete
  reading of $\mathcal{A} V \le -\theta V$ with $\beta = e^{-\theta}$.
  (At this profile guard rows land inside this one row; see
  \autoref{rem:qdb-guards}.)
\item[(D2)] \emph{floor, off-target}: $V \ge 1$ on
  $\bar{D} \smallsetminus B$.
\item[(D3)] \emph{nonneg}: $V \ge 0$ on $\bar{D}$; \emph{bound}:
  $V \le M$ on $\bar{D}$.
\item[(R3)/(R4)] hull and fail-closed containment at target $B$: from
  $y \in D \smallsetminus B$, one-step mass stays in $\bar{D}$ and in
  $D \cup B$.
\item[($\Theta_{\mathsf F}$)] \emph{deadline threshold}:
  $\beta^{T} \cdot \mathrm{thr} \le 1 - p$.
\end{itemize}
The claim (\texttt{sound\_QDeadline}): for every $x \in D$ with
$V(x) \le \mathrm{thr}$,
\[
\Pr_x^{\,*}\bigl[\exists k \le T,\ \mathrm{sat}(\omega_k)\bigr] \;\ge\; p
\]
--- outer measure, as everywhere in the family: no measurability is
asked of the satisfaction predicate, and the mark is inert whenever its
truth set is measurable.

\begin{blemma}[discounted tail bound at the enlarged target]
\label{lem:qdb-tail}
Let $B' := B \cup D^{\mathrm{c}}$. Under (D1)--(D3), (R3), (R4): for
every $x \in D \smallsetminus B$,
\[
\Pr_x\bigl[\forall k \le T,\ \omega_k \notin B'\bigr]
\;\le\; \beta^{T}\, V(x).
\]
\end{blemma}
\begin{proof}
Clamp as in the safety battery, but with the floor forced on
$D \smallsetminus B$ rather than off $D$: let
$C := \max(M, 1)$ and let $g$ clamp $\max(V,1)$ into $[0,C]$ on
$D \smallsetminus B$ and $\max(V,0)$ into $[0,C]$ elsewhere. As in
Lemma~\ref{lem:qsb-clamp}: $0 \le g \le C$; $g \ge 1$ \emph{off} $B'$
(here the floor (D2) does the work on
$\bar{D} \smallsetminus B$, and the forced branch covers the rest); on
$\bar{D} \cap (D \cup B)$ --- which carries all one-step mass from
$D \smallsetminus B$, by (R3)$+$(R4) --- the clamp is invisible,
$g = \min(V, C)$. And by the bound row (D3), $V \le M \le C$ on the
hull, so the ceiling never truncates there: $g = V$ on the whole
mass-carrying set and at $y$ itself. The decay row therefore transfers,
$\int g \,dK\,y = \int V \,dK\,y \le \beta\, V(y) = \beta\, g(y)$ for
$y \in D \smallsetminus B$ --- and this is where the bound row is
load-bearing, not an integrability convenience: without it the transfer
would need $\min(\beta V(y), C) \le \beta \min(V(y), C)$, false for
$0 < \beta < 1$ wherever $V > C$ (at $\beta = 0$ both sides vanish;
that branch is handled separately below). Likewise $g(x) = V(x)$. Assume first
$\beta > 0$. Now the
discounted process
$M_n := \beta^{-(n \wedge \tau_{B'})}\, g(\omega_{n \wedge \tau_{B'}})$
is a nonnegative supermartingale (the decay row supplies exactly the
extra factor $\beta$ the discount consumes at each surviving step), and
on the miss event $\{\forall k \le T,\ \omega_k \notin B'\}$ we have
$\tau_{B'} > T$ and $g(\omega_T) \ge 1$, so $M_T \ge \beta^{-T}$.
Markov\textquotesingle s inequality at time $T$ gives
$\Pr[\text{miss}] \le \beta^{T}\, \mathbb{E} M_0
= \beta^{T} g(x) = \beta^{T} V(x)$. For $\beta = 0$ the inverse
discount is unavailable and the bound holds directly: at $T = 0$ the
right side is $V(x) = g(x) \ge 1 \ge \Pr[\text{miss}]$ (the floor off
$B'$), while for $T \ge 1$ the decay row at $\beta = 0$ forces
$\int g \,dK_y = 0$ for every $y$ off $B'$, so --- $g \ge 1$ off $B'$
--- the kernel from any off-$B'$ point charges $B'$ with full mass and
the miss event is null, by the one-step marginal of the trajectory
law --- the same prefix--step disintegration
Lemma~\ref{lem:qdb-noescape} rides. The mechanized route packages this
computation as the geometric-drift engine (the discounted
supermartingale and its hitting-time moment bound), consumed at the
enlarged target $B'$.
\end{proof}

\begin{blemma}[no escape]\label{lem:qdb-noescape}
Under (R4) at target $B$: for $x \in D$,
\[
\Pr_x\Bigl[\,\exists n,\ \bigl(\forall k \le n,\ \omega_k \notin B\bigr)
\wedge \omega_n \notin D \Bigr] \;=\; 0 .
\]
\end{blemma}
\begin{proof}
Write $E_n$ for the event that some escape as displayed occurs by time
$n$. Induction on $n$ gives $\Pr_x[E_n] = 0$ for every $n$. Base:
coordinate $0$ is deterministically $x \in D$. Step: an escape at time
$n+1$ that has not yet visited $B$ either already escaped by time $n$
(null by the inductive hypothesis) or steps from a state in
$D \smallsetminus B$ to a state outside $D \cup B$ --- and that step is
null, because the joint law of the $n$-prefix and the next step
disintegrates through the one-step kernel, at which (R4) says the
section $(D \cup B)^{\mathrm{c}}$ has measure zero from every surviving
prefix endpoint. The displayed event is $\bigcup_{n} E_n$, null as a
countable union of null events. This lemma is the almost-everywhere
invariance content of (R4) that a \emph{reach} conclusion cannot avoid:
the safety battery got fail-closure for free by enlarging its stop set
--- leaving counts as failing is conservative there --- but a reach
guarantee cannot count escape as success, so the escape event must be
shown null outright.
\end{proof}

\begin{bthm}[target-hit form]\label{thm:qdb-hit}
Under (D1)--(D3), (R3), (R4), ($\Theta_{\mathsf F}$): for every
$x \in D$ with $V(x) \le \mathrm{thr}$,
$\Pr_x[\exists k \le T,\ \omega_k \in B] \ge p$.
\end{bthm}
\begin{proof}
If $x \in B$, the deadline is met at time $0$ almost surely (coordinate
$0$ is deterministic) and $p \le 1$ suffices --- the degenerate corner,
handled first so the clamp may assume $x \notin B$. Otherwise decompose
the miss event: $\{$miss $B$ through $T\} \subseteq \{$miss $B'$ through
$T\} \cup \{$escape before visiting $B\}$. The second is null
(Lemma~\ref{lem:qdb-noescape}); the first has probability
$\le \beta^{T} V(x) \le \beta^{T}\,\mathrm{thr} \le 1 - p$
(Lemma~\ref{lem:qdb-tail} and ($\Theta_{\mathsf F}$)). Complement.
\end{proof}

\begin{bprop}[soundness of \rulename{QDeadlineBattery}]
\label{prop:qdb-sound}
With the formula tie (F$'$): $\Pr_x^{\,*}[\exists k \le T,\
\mathrm{sat}(\omega_k)] \ge p$ for every $x \in D$ with
$V(x) \le \mathrm{thr}$ --- the concluded judgment read on the sublevel
region.
\end{bprop}
\begin{proof}
$B$ under-approximates the truth set, so the hit event of
Theorem~\ref{thm:qdb-hit} is contained in the reach-formula event;
monotonicity of the (outer) measure. As always, no measurability is
asked of the satisfaction predicate.
\end{proof}

\begin{bprop}[catalogue form: the $(\theta, q)$ parameterization]
\label{prop:qdb-catalogue}
Suppose additionally $\beta \le e^{-\theta}$ with $\theta \ge 0$ (the
drift tie), $0 \le q \le \sum_{k \le n} (\theta T)^k / k!$ for some $n$
(the Taylor leaf --- the one closed, wire-free comparison the checker
replays in exact rational arithmetic, discharged as a ground-fact-rule
application), and $p \le 1$. Then the conclusion of
Proposition~\ref{prop:qdb-sound} holds on the region
$\{V \le (1-p)\,q\}$.
\end{bprop}
\begin{proof}
$q \le \sum_{k \le n} (\theta T)^k/k! \le e^{\theta T}$ --- the first
transcendental fact, living here and not in the checker, applied at the
nonnegative exponent $\theta T$ --- and $\beta^{T} \le e^{-\theta T}$,
powers preserving order between nonnegative bases ($0 \le \beta \le
e^{-\theta}$, $T \in \N$); so
$\beta^{T} q \le e^{-\theta T} e^{\theta T} = 1$ --- the second. Hence
$\beta^{T} \cdot (1-p)\,q \le 1 - p$: hypothesis
($\Theta_{\mathsf F}$) holds at $\mathrm{thr} = (1-p)\,q$, and
Proposition~\ref{prop:qdb-sound} applies. The checker never needs to
know $e$ exists; soundness is indifferent to the producer knobs
$\theta$ and $q$ --- any smaller rate or budget is admissible.
\end{proof}

\begin{bremark}[why the off-target floor is not decoration]
\label{rem:qdb-floor}
Drop (D2) and take $V \equiv 0$ with the identity kernel
$K\,y = \delta_y$ at some $x \in D \smallsetminus B$: decay holds
($0 \le \beta \cdot 0$), nonneg and bound hold, and
($\Theta_{\mathsf F}$) permits any $p \le 1$ since
$\beta^T \cdot 0 = 0$. The rule would then certify that a chain frozen
forever outside the target reaches it with probability one. With the
floor present the certificate dies: $V \ge 1$ on
$\bar{D} \smallsetminus B \ni x$ contradicts $V(x) = 0$. Missing the
deadline must cost at least one unit of discounted certificate value ---
that is precisely the step $M_T \ge \beta^{-T}$ in
Lemma~\ref{lem:qdb-tail}. And the floor is asked \emph{off the target
only}: on $B$ the certificate may take any value, because success needs
no budget. This checked form is deliberately weaker than a global
floor, and the proven statement binds the weaker form.
\end{bremark}

\begin{bremark}[the two intrinsic asymmetries against the safety battery]
\label{rem:qdb-asym}
Both are consequences of concluding \emph{reach} rather than
\emph{avoid}, and neither is an artefact of the mechanization. First,
the boundedness row $V \le M$: the safety battery's \emph{additive}
drift survives its ceiling for free ---
$\int \min(V, C)\,dK\,y \le \min(\int V\,dK\,y,\, C) \le
\min(V(y), C)$ --- so a ceiling at
$\max(\mathrm{thr}, 1)$ cost nothing; here the discount multiplies the
certificate by $\beta^{-n}$, and the transfer of the decay row to the
clamped certificate needs the ceiling to be invisible wherever one-step
mass reaches --- exactly (D3)'s bound, consumed inside
Lemma~\ref{lem:qdb-tail} (and in the mechanized geometric-drift engine
it additionally supplies integrability of the discounted certificate)
--- for polynomial certificates over a compact hull it is free. Second, escape: the safety battery counts
leaving the modelled domain as failure by enlarging its stop set, which
is conservative for an avoid conclusion but would be unsound generosity
for a reach conclusion --- hence Lemma~\ref{lem:qdb-noescape}, whose
almost-everywhere argument is the one place this
section\textquotesingle s hypotheses stop being pointwise.
\end{bremark}

\begin{bremark}[guards at this profile]\label{rem:qdb-guards}
The catalogued rule lists guard rows (jumps instantaneous). At the
discrete-kernel profile the one-step kernel carries flow and enabled
jumps together, so the guard rows land inside the single decay row (D1)
--- exactly as the safety battery\textquotesingle s interior and guard
rows land together. The instantaneous-jump reading --- the discount
factor does not tick during a jump, so the $e^{\theta t}$ factor cancels
across it --- belongs to the continuous-time profile, where it is
carried by the martingale-problem discharge of
\autoref{rem:qdb-profile}.
\end{bremark}

\begin{bremark}[\textbf{profile} --- what is proven, at which
presentations]
\label{rem:qdb-profile}
Proven: the discrete-kernel profile, sorry-free
(\texttt{sound\_QDeadline\_hit}, \texttt{sound\_QDeadline},
\texttt{sound\_QDeadlineBattery}). Also proven: the continuous-time
profile at the martingale-problem presentation
(\texttt{sound\_GSHS\_deadline\_mp},
\texttt{sound\_GSHS\_deadline\_battery\_mp}, module
\texttt{QtslAux.LayerD.MPDischarge} of the semantic-universe
strand\textquotesingle s auxiliary package) --- and the three
ingredients this remark once filed are exactly what that discharge
supplies: the $e^{\theta t} V$ supermartingale over continuous time
is the exponential-discount construction from the undiscounted
martingale-problem hypothesis; optional stopping for right-continuous
nonnegative supermartingales is the library theorem, consumed at the
d\'ebut; and the instantaneous-jump reading rides the
presentation\textquotesingle s no-alarm segment assembly. The
martingale-problem theorems are hash-bound by this
catalogue\textquotesingle s binder via the qtsl-aux package
registration, as for the safety sibling, so the dashboard row carries
no conditional mark --- each named profile is proven, and its
theorems are bound.
\end{bremark}

\begin{bremark}[relation to the literature]\label{rem:qdb-lit}
The engine is the geometric drift condition of the Foster--Lyapunov
family \cite{meyntweedie2009}: expected decrease by a factor per step
off the target, yielding exponential hitting-time tails. The discounted
supermartingale and its optional-stopping consumption are classical
\cite{doob1953}; the barrier mechanics are the same
Ville-type argument \cite{ville1939} the whole family rides. The
certificate-form packaging --- exact rational rows, transcendental facts
quarantined inside a verified lemma --- is this
calculus\textquotesingle s contribution rather than the
literature\textquotesingle s.
\end{bremark}

\begin{anchors}
Lemma \ref{lem:qdb-tail}: discounted engine, clamp for the deadline reading & \texttt{BernsteinLogic.Deadline.deadline\_tail} (via \texttt{supermartingale\_inv\_pow\_mul}, \texttt{expMoment\_hitting}); the clamp $g$ inside \texttt{sound\_QDeadline\_hit}\\
Lemma \ref{lem:qdb-noescape}: escape-null induction & \texttt{ae\_no\_escape} (file-local; joint prefix--step law via \texttt{Kernel.partialTraj\_compProd\_eq\_map\_traj}); base case \texttt{ae\_coord\_zero\_mem}\\
Theorem \ref{thm:qdb-hit}: assembly, degenerate corner & \texttt{sound\_QDeadline\_hit}\\
Proposition \ref{prop:qdb-sound}: formula tie & \texttt{sound\_QDeadline}\\
Proposition \ref{prop:qdb-catalogue}: Taylor leaf, drift tie, transcendental quarantine & \texttt{sound\_QDeadlineBattery} (via \texttt{Real.sum\_le\_exp\_of\_nonneg})\\
\end{anchors}


\rulesection{42}{QUnionBound}{core rule}
\label{blcatrule-QUnionBound}
\ruleblock{The rule (Steele metanotation)}
\sourcenote{as realized in the mechanized development; source row: \texttt{\texttt{QTSL-rules.toml}} --- \emph{QUnionBound}. The metanotation statement is typeset by hand for notational quality; its mechanized ground truth is the Lean statement below, extracted from the sources at every build of this book and hash-bound. The rule \emph{roster} is machine-checked (every mechanized rule tag round-trips to exactly one catalogue row); the statement \emph{text} is editorial, checked by eye against the extracted Lean beside it.}
\begin{mathpar}
\inferrule*{
  \mu(E_1) \le a \\ \mu(E_2) \le b \\ \mathrm{fail} \subseteq E_1 \cup E_2
}{
  \mu(\mathrm{fail}) \le a + b
}
\end{mathpar}

\ruleflavor{rule sort baked in \texttt{theory Qtsl} --- realized measure-generically (arbitrary measure; no profile residue). Packaged instance (the d\'ebut decomposition, pure path-set algebra over sequence space): failure of $(\varphi_s\,\mathsf{U}_{\le T}\,\varphi)$, at stop set $\mathit{Stop}$ and good-exit set $R$, is contained in $\{\text{some \emph{first} } \mathit{Stop}\text{-visit presents a bad state}\} \cup \{\text{no } \mathit{Stop}\text{-visit through } T\}$ --- the containment the until battery cites as its $\mathrm{fail} \subseteq E_1 \cup E_2$ premise.}
\par\medskip
\rulename{QUnionBound} is the battery family\textquotesingle s glue:
from two probability bounds and an event decomposition it concludes that
the sum of the budgets bounds the failure event ---
$\Pr[E_1] \le a$, $\Pr[E_2] \le b$, and
$\mathrm{fail} \subseteq E_1 \cup E_2$ give
$\Pr[\mathrm{fail}] \le a + b$. Its semantic content is monotonicity
plus subadditivity (Boole\textquotesingle s inequality), stated over an
\emph{arbitrary} measure: the rule is profile-generic, and its instance
at any system-semantics profile is the same lemma applied at that
profile\textquotesingle s path measure. Its dashboard row accordingly
carries an unconditional check with no profile discharge even needed
--- there was never a profile residue to discharge (the battery
siblings\textquotesingle{} rows earned theirs through the
martingale-problem discharge; this one is measure-generic from
birth).

The hypothesis $\mathrm{fail} \subseteq E_1 \cup E_2$ is where all the
content concentrates. The rule is sound for \emph{any} decomposition;
the citing rule must justify the containment. The packaged instance ---
the one the bounded-until battery cites --- is the d\'ebut
decomposition: failure of ``stay safe until the target, by the
deadline'' is contained in $\{$some \emph{first} stop-set visit presents
a bad state$\}$ $\cup$ $\{$no stop-set visit by the horizon$\}$. That
containment is pure path-set algebra over sequence space: no measure, no
measurability, only the well-ordering of the time axis.

A producer rarely writes this node by hand; it appears where a composite
temporal property splits its failure modes across two certificates, and
what the checker recomputes at such a node is precisely the containment
side of the decomposition (or it discharges it by citing the packaged
path-set lemma). The same additivity doubles as the slack-composition
device in assume--guarantee composition, where interface slack budgets
add across a wiring --- one lemma, two consumers.

Why core rather than admissible: there is nothing to unpack into.
Subadditivity of measure is not derivable from the other battery rules;
it is new --- deliberately small --- semantic content at the measure
level, and keeping it a separate node keeps every consumer\textquotesingle s
soundness proof honest about where its additivity comes from.

\cataloguenote{Stated in two forms: an extended-nonneg form (the primitive), and a
real-valued form in the shape battery bounds actually arrive
(certificate initial values), whose nonnegativity hypotheses are free at
every use site.
}
\ruleblock{Lean statement}
\declhead{BernsteinLogicAux.QtslBattery}{sound\_QUnionBound}{\badge{BadgeGreen}{statement bound}}
\begin{leansrc}
/-- **`sound_QUnionBound`** (ℝ≥0∞ form). Subadditivity through an event
decomposition — the entire semantic content of the `QUnionBound` rule
(QTSL-rules.toml row `QUnionBound`). -/
theorem sound_QUnionBound {Ω : Type*} [MeasurableSpace Ω] (μ : Measure Ω)
    {E₁ E₂ F : Set Ω} {a b : ℝ≥0∞}
    (h₁ : μ E₁ ≤ a) (h₂ : μ E₂ ≤ b) (hsub : F ⊆ E₁ ∪ E₂) :
    μ F ≤ a + b
\end{leansrc}
\declhead{BernsteinLogicAux.QtslBattery}{sound\_QUnionBound\_ofReal}{\badge{BadgeGreen}{statement bound}}
\begin{leansrc}
/-- **`sound_QUnionBound`** (`ENNReal.ofReal` form), in the shape the battery
bounds arrive: real-valued certificate bounds `a`, `b` combine to `a + b`.
The nonnegativity hypotheses are free at every use site (certificate initial
values, nonneg rows). -/
theorem sound_QUnionBound_ofReal {Ω : Type*} [MeasurableSpace Ω] (μ : Measure Ω)
    {E₁ E₂ F : Set Ω} {a b : ℝ} (ha : 0 ≤ a) (hb : 0 ≤ b)
    (h₁ : μ E₁ ≤ ENNReal.ofReal a) (h₂ : μ E₂ ≤ ENNReal.ofReal b)
    (hsub : F ⊆ E₁ ∪ E₂) :
    μ F ≤ ENNReal.ofReal (a + b)
\end{leansrc}
\declhead{BernsteinLogicAux.QtslBattery}{compl\_firstExit\_subset}{\badge{BadgeGreen}{statement bound}}
\begin{leansrc}
/-- **The first-exit event decomposition** (the `QUnionBound` premise's
début decomposition): over
sequence space, the complement of the first-exit-success event — some time
`n ≤ T` whose state is in the good exit set `Good` with no earlier visit to
the stop set `Stop` — is contained in the union of {some `Stop`-visit is a
*first* visit presenting a bad state (`Stop \ Good`)} and {no `Stop`-visit
through `T`}. Pure path-set algebra, no measure and no measurability; the
until battery instantiates `Stop := (A ∪ R) ∪ Dᶜ`, `Good := R`. This is the
lemma the `QUnionBound` premise `fail ⊆ E₁ ∪ E₂` cites for T8. -/
theorem compl_firstExit_subset {S : Type*} (Stop Good : Set S) (T : ℕ) :
    {ω : ℕ → S | ∃ n ≤ T, ω n ∈ Good ∧ ∀ j < n, ω j ∉ Stop}ᶜ ⊆
      {ω : ℕ → S | ∃ n, ω n ∈ Stop \ Good ∧ ∀ j < n, ω j ∉ Stop} ∪
        {ω : ℕ → S | ∀ k ≤ T, ω k ∉ Stop}
\end{leansrc}
\ruleblock{Correspondence (Steele display $\leftrightarrow$ Lean)}
\badge{BadgeAmber}{a bound Lean statement CHANGED since certification --- re-review and re-certify}
\sourcenote{The dictionary below is a reviewer-certified statement of how the mechanized statements above interpret the displayed rule: display metavariable $\mapsto$ Lean binder, each display hypothesis $\mapsto$ the premise or side condition carrying it, the display conclusion $\mapsto$ the goal --- with any deliberate difference in width between display and mechanization named outright. Certification is recorded as a hash triple over (the displayed cell, the extracted statement texts, this dictionary\textquotesingle s own text); if any of the three changes, this block turns amber until a reviewer re-certifies. The correspondence claim itself remains editorial --- reviewer-grade, not machine-checked; its machine-checked successor types each soundness statement against the rule as a structured object, and supersedes this block per rule as it lands.}
The cleanest dictionary in the Q-family --- the certified statement is
the displayed rule at full measure-generality, with no profile
residue.
\begin{itemize}
\item $\mu \mapsto$ an arbitrary \texttt{Measure μ} on an arbitrary
  measurable space --- the amended display\textquotesingle s own
  generality (\emph{wider} than the original display\textquotesingle s
  path-probability reading, in the safe direction): the rule\textquotesingle
  s GSHS instance and the AG slack-composition use are the same lemma
  at other measures (the roster\textquotesingle s ``one item, two
  consumers\textquotesingle\textquotesingle).
\item the three premises $\mu(E_1) \le a$, $\mu(E_2) \le b$,
  $\mathrm{fail} \subseteq E_1 \cup E_2$ $\mapsto$ \texttt{h₁},
  \texttt{h₂}, \texttt{hsub}, verbatim; conclusion $\mu(\mathrm{fail})
  \le a + b$ $\mapsto$ the goal, verbatim. No hypothesis is unused
  and nothing undisplayed is demanded --- not even measurability of
  the events (monotonicity and subadditivity hold for outer measure).
\item the \texttt{ofReal} companion carries the same rule in the
  shape battery bounds arrive (real-valued $a$, $b$), its two
  nonnegativity riders free at every use site, per the note.
\item the packaged instance --- the flavor text\textquotesingle s
  d\'ebut decomposition
  --- $\mapsto$ \texttt{compl\_firstExit\_subset}, certified
  alongside: pure path-set algebra over sequence space (no measure,
  no measurability), delivering exactly the $\mathrm{fail} \subseteq
  E_1 \cup E_2$ premise the until battery cites at $\mathit{Stop} :=
  (A \cup R) \cup D^{\mathrm{c}}$, $\mathit{Good} := R$.
\end{itemize}

\ruleblock{Lean verifier receipts}
\sourcenote{module \texttt{BernsteinLogicAux.QtslBattery}; bound build \texttt{5c75f9923f27\ldots} (bound from receipt 3fd830b4914bdeda, 2026-08-04T01:56:58Z). A receipt records that named checking layers accepted a build: (module, compiled-artifact SHA-256, environment SHA-256, checker layers), minted when the checking gate ran; the store receipts attesting the bound build are listed with the layers each run credited. Nothing here asks the reader to trust the store: every layer is an independently runnable program over the quoted sources --- rebuild the module and re-run the checkers to reproduce the verdict. Receipts attesting superseded builds credit no layer and are omitted from this rendering (the count of omissions, if any, is noted below).}
\begin{receipttable}
\receiptrow{3fd830b4914bdeda}{2026-08-04T01:56:58Z}{build+decorrelated}{\badge{BadgeGreen}{current}}
\receiptrow{96b18d8572451f65}{2026-08-04T03:11:31Z}{build+decorrelated}{\badge{BadgeGreen}{current}}
\end{receipttable}
\sourcenote{13 further receipt(s) in the store attest superseded builds of this module and are omitted from this rendering.}
\layerline{\layeritem{build}{lake build (elaboration)}{\badge{BadgeGreen}{current}}\layeritem{replay}{leanchecker (official kernel replay)}{\badge{BadgeAmber}{stale only}}\layeritem{hardened}{safe\_verify (hardened screen)}{\badge{BadgeAmber}{stale only}}\layeritem{decorrelated}{lean4export + nanoda (decorrelated kernel)}{\badge{BadgeGreen}{current}}}
\sourcenote{\badge{BadgeGray}{comparator: none} no fixture in the three-implementation comparator harness (compiled / kernel-interpreted / reference-interpreted) exercises this rule.}
\ruleblock{HOL4 statement}
\absentbox{NO HOL STATEMENT YET}{The HOL4 development ports the base calculus\textquotesingle s checker
chain (theory \texttt{bern\_catalogue}, developed side by side with the
base catalogue sources, stopping below the proof-object layer); it
contains no probability or measure theory. \rulename{QUnionBound} lives
one theory-extension edge above the base calculus, and even its small
soundness content (measure subadditivity, a path-set decomposition over
sequence space) has no counterpart anywhere in the ported development. A
HOL statement of this rule would open a measure-theory development, not
extend the ported checker chain; until one exists, this box is the
honest state of the edifice.
}
\ruleblock{HOL verifier receipts}
\absentbox{NONE}{no HOL verification receipt exists for this rule.}
\ruleblock{Proof}

\sourcenote{Genre: core rule (tier 1) --- the battery
family\textquotesingle s smallest member, and its only
\emph{measure-generic} one: the theorems below are stated over an
arbitrary measure, so the rule carries no profile conditionality at all
--- its instance at any system-semantics profile is the same lemma
applied at that profile\textquotesingle s path measure. A Putnam grader
would accept the one-line proof and then ask where the content went; the
honest answer --- into the decomposition hypothesis --- is
\autoref{rem:qub-content}.}

\textbf{Setting and claim.}\enspace Let $(\Omega, \mu)$ be any
measurable space with a measure, read throughout as the induced outer
measure on arbitrary subsets (coinciding with $\mu$ on measurable sets)
--- no measurability is asked of any event below, and both monotonicity
and finite subadditivity hold in that reading. The rule: from
$\mu(E_1) \le a$, $\mu(E_2) \le b$, and a decomposition
$\mathrm{fail} \subseteq E_1 \cup E_2$, conclude
$\mu(\mathrm{fail}) \le a + b$. In its battery consumption, $\mu$ is the
path measure $\Pr_x$ of the profile at hand, $\mathrm{fail}$ is the
failure event of a composite temporal property, and $a, b$ are the
initial values of two certificates.

\begin{bthm}[soundness of \rulename{QUnionBound}]\label{thm:qub-sound}
$\mu(E_1) \le a$, $\mu(E_2) \le b$, and
$\mathrm{fail} \subseteq E_1 \cup E_2$ imply
$\mu(\mathrm{fail}) \le a + b$.
\end{bthm}
\begin{proof}
Monotonicity, then subadditivity (Boole\textquotesingle s inequality):
$\mu(\mathrm{fail}) \le \mu(E_1 \cup E_2) \le \mu(E_1) + \mu(E_2)
\le a + b$. Note what is \emph{not} used: no measurability of
\emph{any} of the three events (the Setting\textquotesingle s
outer-measure reading supplies monotonicity and finite subadditivity on
arbitrary subsets), no
disjointness of $E_1, E_2$, no probability normalization --- the
statement holds for any measure, which is exactly what makes the rule
profile-generic.
\end{proof}

\begin{bprop}[real-valued form]\label{prop:qub-ofreal}
With $a, b \ge 0$ real and the bounds given as
$\mu(E_i) \le \mathrm{ofReal}$-embedded budgets, the conclusion is
$\mu(\mathrm{fail}) \le a + b$ under the same embedding.
\end{bprop}
\begin{proof}
The embedding of nonnegative reals into extended nonnegative reals is
additive on nonnegatives; apply Theorem~\ref{thm:qub-sound}. This is the
shape battery bounds actually arrive in --- certificate initial values
--- and its nonnegativity hypotheses are free at every use site
(certificate nonneg rows).
\end{proof}

\begin{blemma}[the d\'ebut decomposition --- the packaged instance]
\label{lem:qub-debut}
Over sequence space $S^{\N}$, fix a stop set
$\mathit{Stop} \subseteq S$, a good-exit set $\mathit{Good} \subseteq S$,
and a horizon $T$. Then the complement of the first-exit-success event
--- some $n \le T$ with $\omega_n \in \mathit{Good}$ and no earlier
visit to $\mathit{Stop}$ --- is contained in
\[
\bigl\{\exists n,\ \omega_n \in \mathit{Stop} \smallsetminus \mathit{Good}
\ \wedge\ \forall j < n,\ \omega_j \notin \mathit{Stop}\bigr\}
\;\cup\;
\bigl\{\forall k \le T,\ \omega_k \notin \mathit{Stop}\bigr\},
\]
i.e.\ in $\{$some \emph{first} $\mathit{Stop}$-visit presents a bad
state$\}$ $\cup$ $\{$no $\mathit{Stop}$-visit through $T\}$.
\end{blemma}
\begin{proof}
Suppose the path is not a first-exit success. If it never visits
$\mathit{Stop}$ through $T$, the second disjunct holds. Otherwise some
$k \le T$ has $\omega_k \in \mathit{Stop}$; let $n$ be the \emph{least}
index with $\omega_n \in \mathit{Stop}$ (the well-ordering of $\N$ ---
this is the entire ``d\'ebut'' content; at continuous time the
first-entry construction is a genuine theorem, while over $\N$ it is a
minimum). Then $n \le k \le T$ and no $j < n$ visits $\mathit{Stop}$;
were $\omega_n \in \mathit{Good}$, the path would be a first-exit
success at $n$ --- contradiction. So $\omega_n \in \mathit{Stop}
\smallsetminus \mathit{Good}$, and $n$ is a first visit: the first
disjunct holds. No measure, no measurability, no topology was consulted
--- the lemma is pure path-set algebra.
\end{proof}

\begin{bremark}[where the content lives]\label{rem:qub-content}
The rule is sound for \emph{any} events $E_1, E_2$; every gram of
content sits in the hypothesis
$\mathrm{fail} \subseteq E_1 \cup E_2$, which the citing rule must
justify and the checker must see justified. Drop it and the conclusion
is unrelated to the premises --- one could ``bound'' any event by any
two budgets. In the bounded-until battery\textquotesingle s use,
Lemma~\ref{lem:qub-debut} \emph{is} the justification, at
$\mathit{Stop} := (A \cup R) \cup D^{\mathrm{c}}$ and
$\mathit{Good} := R$: failure of ``stay safe until the target, by the
deadline'' splits into a bad first exit (bounded by the avoid
certificate) and no exit by the horizon (bounded by the deadline
certificate). The split covers every failure: a failing path either
never visits $\mathit{Stop}$ through the horizon (the second disjunct)
or has a first $\mathit{Stop}$-visit, which either is or is not a good
landing --- and a good landing, since $R \subseteq \mathit{Stop}$ in
this instance, is a success, so no third thing remains.
\end{bremark}

\begin{bremark}[measure-genericity, stated sibling-aware]
\label{rem:qub-generic}
Within this book\textquotesingle s battery family --- safety, deadline,
until, and this rule --- the other three carry conditional marks (now
recording, for safety, deadline, and until, that their
continuous-time martingale-problem theorems are cited rather than
hash-bound; see each profile remark); this rule alone is warranted
unconditionally, because its statement never mentions a kernel, a path
construction, or a filtration. The prediction this remark has carried
since its first edition held: the continuous-time discharges of the
siblings consume these same two theorems verbatim --- nothing here
needed to be touched. (The d\'ebut
lemma is the sequence-space instance of its decomposition; a
continuous-time consumer would supply its own first-entry construction
--- a genuine theorem there, as the lemma\textquotesingle s proof notes
--- through those same transports.)
\end{bremark}

\begin{bremark}[the second consumer]\label{rem:qub-second}
The same additivity is the slack-composition device in assume--guarantee
composition: when interface contracts compose across a wiring, their
slack budgets add, and the addition is this rule\textquotesingle s
theorem applied to the composed failure events. One lemma, two consumers
--- which is the reason the rule is kept as a separate node rather than
inlined into the until battery: every consumer\textquotesingle s
soundness proof stays honest about where its additivity comes from.
\end{bremark}

\begin{bremark}[why core, and why this small]\label{rem:qub-tier}
There is nothing to unpack into: subadditivity of measure is not
derivable from the other battery rules --- each of them concludes a
bound on a single certificate-shaped event from kernel rows, and none
mentions two arbitrary events, let alone their union --- so the rule
cannot be admissible: it is new semantic content, deliberately minimal. Its smallness is a
design decision inherited from the base
calculus\textquotesingle s taste: keep kernel-level content atomic, let
composite rules be admissible unpackings over the atoms.
\end{bremark}

\begin{anchors}
Theorem \ref{thm:qub-sound} & \texttt{sound\_QUnionBound}\\
Proposition \ref{prop:qub-ofreal} & \texttt{sound\_QUnionBound\_ofReal}\\
Lemma \ref{lem:qub-debut} & \texttt{compl\_firstExit\_subset}\\
the until-battery consumption & the \texttt{sound\_QUnionBound\_ofReal} application inside \texttt{sound\_QUntilBattery\_firstHit}\\
\end{anchors}


\rulesection{43}{QUntilBattery}{admissible rule}
\label{blcatrule-QUntilBattery}
\ruleblock{The rule (Steele metanotation)}
\sourcenote{as realized in the mechanized development; source row: \texttt{\texttt{QTSL-rules.toml}} --- \emph{QUntilBattery}. The metanotation statement is typeset by hand for notational quality; its mechanized ground truth is the Lean statement below, extracted from the sources at every build of this book and hash-bound. The rule \emph{roster} is machine-checked (every mechanized rule tag round-trips to exactly one catalogue row); the statement \emph{text} is editorial, checked by eye against the extracted Lean beside it.}
\begin{mathpar}
\inferrule*{
  E_{\mathrm{bad}} := \{\text{the first exit from } D \smallsetminus (A \cup R) \ \text{lands in } A \cup D^{\mathrm{c}}\}, \quad E_{\mathrm{miss}} := \{A \cup R \ \text{not reached by } T\} \\
  [\![\neg\varphi_s]\!] \subseteq A, \qquad R \subseteq [\![\varphi]\!] \ \text{(formula ties)} \\
  \mu_x(E_{\mathrm{bad}}) \le V_a(x) \ \text{(the \rulename{QSafetyBattery} slice for } V_a \ \text{at stop set } (A \cup R) \cup D^{\mathrm{c}}\text{: floor on } A \ \text{only)} \\
  \mu_x(E_{\mathrm{miss}}) \le \beta^{T}\, V_d(x) \ \text{(the \rulename{QDeadlineBattery} slice for } V_d \ \text{at target } A \cup R\text{)} \\
  [\![\neg(\varphi_s\,\mathsf{U}_{\le T}\,\varphi)]\!] \subseteq E_{\mathrm{bad}} \cup E_{\mathrm{miss}} \ \text{(the d\'ebut decomposition; consumed by \rulename{QUnionBound})} \\
  0 \le \beta, \qquad 0 \le V_d(x), \qquad \beta^{T} \le c \ \text{(sign gates; rate tie)}
}{
  \mathcal{M},\ x_0 \in \{V_a + c\,V_d \le 1 - p\} \vDash \mathsf{SProbGE}\bigl(p,\ \mathsf{PUntil}(T, \varphi_s, \varphi)\bigr)
}
\end{mathpar}

\ruleflavor{rule sort baked in \texttt{theory Qtsl} (\texttt{QUntilBattery}, avoidance and target certificates as separate family fields), the concluded affine region atom (\texttt{qraAffine}) baked beside it --- the display is the realized form, as the discrete-profile assembly binds it. The two slice premises are the sibling batteries\textquotesingle{} kits re-aimed --- the avoid slice is the safety kit at the composite stop set with its floor asked on the bad exit $A$ only, the deadline slice drives into $A \cup R$ by the horizon --- and each premise is the measure bound that slice\textquotesingle s rows conclude; the rows themselves are the sibling entries\textquotesingle{} premises. Catalogue instance of the rate tie: $c = q^{-1}$ with $q \le e^{\theta T}$, $\beta = e^{-\theta}$ --- fully rational at the discrete profile, no transcendental leaf. The conclusion is the region-conditioned judgment in constructor vocabulary, its composite region the baked \texttt{qraAffine} atom over $V_a + c\,V_d$ at threshold $1 - p$.}
\par\medskip
\rulename{QUntilBattery} concludes the bounded-until judgment
$\mathsf{SProbGE}\bigl(p,\ \mathsf{PUntil}(T, \varphi_s, \varphi)\bigr)$
(surface syntax $\mathsf{P}_{\ge p}[\varphi_s\, \mathsf{U}_{\le T}\,
\varphi]$) on a composite \emph{affine} region
$\mathsf{qraAffine}([(1,\, V_a),\, (c,\, V_d)],\ 1 - p)$ --- the
sublevel set $\{V_a + c\, V_d \le 1 - p\}$ --- from two certificates doing different
work: an avoid certificate $V_a$ guarding the safe-side formula
$\varphi_s$ until the target is reached, and a deadline certificate
$V_d$ driving the process into the target by the horizon. It is this
catalogue\textquotesingle s exemplar admissible rule at the satisfaction
stratum: one node unpacks into two battery slices and one
\rulename{QUnionBound} application, and the mechanized soundness proof
\emph{is} that unpacking --- it consumes the deadline
battery\textquotesingle s theorem as a black box, applies the union
bound once, and never re-proves either battery.

The two floors do different work, and the asymmetry is the
rule\textquotesingle s design. $V_a$\textquotesingle s floor is asked on
the bad-exit set $A$ \emph{only} --- the avoid certificate is silent on
$R \smallsetminus A$, so a landing there costs no certificate
budget: the rule does not charge for winning (a landing in
$A \cap R$, when the formula ties permit one, succeeds but is still
priced as avoid-territory). $V_d$\textquotesingle s
floor is on all of the hull off the target, exactly as its parent
deadline battery
demands. Fail-closed: leaving the modelled domain joins the stop set on
the \emph{bad} side, while an $R$-landing succeeds wherever it lands
(the target\textquotesingle s formula tie is global).

What the checker recomputes: the side conditions inherited from both
battery slices, plus the one arithmetic hypothesis binding the slices
together --- the rate tie $\beta^{T} \le c$ on the affine
atom\textquotesingle s second coefficient. The catalogued instance
$c = q^{-1}$ with $q \le e^{\theta T}$ and $\beta = e^{-\theta}$
satisfies the tie; at this profile the whole certificate is exact
rational arithmetic, with no transcendental leaf (one becomes available
by the deadline battery\textquotesingle s Taylor-leaf pattern when
$\theta$-rates are wanted). Drop the tie and the affine atom
under-budgets the miss mass; drop either slice\textquotesingle s rows
and that slice\textquotesingle s own counterexamples apply.

Tier, argued: admissible, not alias --- three premise sub-derivations,
not a respelling of one. And deliberately not monolithic core: a third
self-contained lemma would duplicate two batteries\textquotesingle{}
soundness; the mechanized assembly vindicates the choice by factoring
through the named machinery instead. The dashboard mark states the
status: hash-bound at the discrete-kernel profile; proven at the
continuous-time martingale-problem profile
(\texttt{sound\_GSHS\_until\_mp},
\texttt{sound\_GSHS\_until\_battery\_mp},
\texttt{QtslAux.LayerD.MPDischarge}) --- hash-bound below via the
qtsl-aux package registration, so the row carries no conditional
mark: each named profile is proven, and its theorems are bound.

At the general stratum the until acceptance atom is stated over the
semantic universe\textquotesingle s record carrier directly (the
semantic-universe functor package, hash-bound below): at some grid
index the reach leg is hit and the avoid leg is missed
(\texttt{untilEvent}) --- a countable Boolean combination of
hit-events, measurable from per-component tickets alone
(\texttt{measurableSet\_untilEvent}). The clocked instantiation shape
is recorded at the corollary: reach legs closed (a $\le$-cut of
stages times a closed target), avoid legs open (a $<$-cut times an
open unsafe region) --- and miss-events of open sets are measurable
with no compactness anywhere. As at the deadline battery, the until
conclusion is a mark ($\exists$-face): persistent by the same
window-calculus engine, with the forward judgment lemma and the same
deliberate absence of a backward direction --- the deadline semantics
the two batteries share.

\cataloguenote{The concluded region is the affine region atom with rows
$[(1, V_a), (c, V_d)]$ and threshold $1 - p$ --- an affine sublevel
region over certificate values, the region-vocabulary generalization
whose $k = 1$ case is the ordinary sublevel atom.
}
\ruleblock{Lean statement}
\declhead{BernsteinLogicAux.QtslBattery}{sound\_QUntilBattery}{\badge{BadgeGreen}{statement bound}}
\begin{leansrc}
/-- **`sound_QUntilBattery`, discrete profile** (T8): over the first-hit
form, the two formula ties — `A` over-approximates the falsity set of the
safe-side formula (`¬ sat_s y → y ∈ A`, the T13 OUTWARD direction), `R`
under-approximates the truth set of the target formula (`∀ y ∈ R, sat y`) —
weaken the concluded event to the truth of `φ_s U_{≤T} φ` along the path.
No measurability is asked of either satisfaction predicate (`measure_mono`
on the measurable first-hit event). -/
theorem sound_QUntilBattery {sat_s sat : S → Prop}
    (hD : MeasurableSet D) (hDbar : MeasurableSet Dbar)
    (hA : MeasurableSet A) (hR : MeasurableSet R)
    (hVa : Measurable Va) (hVd : Measurable Vd)
    (hFs : ∀ y, ¬ sat_s y → y ∈ A) (hFt : ∀ y ∈ R, sat y)
    (hβ : 0 ≤ β) (hc : β ^ T ≤ c)
    (hDsub : D ⊆ Dbar)
    (hR3 : ∀ y ∈ D, y ∉ A ∪ R → K y Dbar = 1)
    (hR4 : ∀ y ∈ D, y ∉ A ∪ R → K y (D ∪ (A ∪ R)) = 1)
    (hVaint : ∀ y ∈ D, y ∉ A ∪ R → Integrable Va (K y))
    (hVadrift : ∀ y ∈ D, y ∉ A ∪ R → ∫ z, Va z ∂K y ≤ Va y)
    (hVafloor : ∀ y ∈ A, y ∈ Dbar → 1 ≤ Va y)
    (hVanonneg : ∀ y ∈ Dbar, 0 ≤ Va y)
    (hVdint : ∀ y ∈ D, y ∉ A ∪ R → Integrable Vd (K y))
    (hVddecay : ∀ y ∈ D, y ∉ A ∪ R → ∫ z, Vd z ∂K y ≤ β * Vd y)
    (hVdfloor : ∀ y ∈ Dbar, y ∉ A ∪ R → 1 ≤ Vd y)
    (hVdnonneg : ∀ y ∈ Dbar, 0 ≤ Vd y)
    (hVdM : ∀ y ∈ Dbar, Vd y ≤ M)
    (hx : x ∈ D) (hthr : Va x + c * Vd x ≤ 1 - p) :
    ENNReal.ofReal p ≤ pathMeasure K x
      {ω | ∃ n ≤ T, sat (ω n) ∧ ∀ j < n, sat_s (ω j)}
\end{leansrc}
\declhead{BernsteinLogicAux.QtslBattery}{sound\_QUntilBattery\_firstHit}{\badge{BadgeGreen}{statement bound}}
\begin{leansrc}
/-- **`sound_QUntilBattery`, discrete profile, first-hit form** (the
catalogue's T8 row, fail-closed reading). The rule is
admissible over {T5, T6, `QUnionBound`}, and this proof *is* that unpacking:

* the T5 slice is `measure_firstBad_le` on the composite stop set
  `Stop := (A ∪ R) ∪ Dᶜ` with the clamped avoid certificate `V_a` — floor
  rows on `A` only, drift rows off `A ∪ R` (the process is stopped at the
  target, so `V_a` is silent on `R`);
* the T6 slice is `measure_missTarget_le` at target `A ∪ R` with the
  deadline certificate `V_d` — global floor off the target;
* the glue is one `sound_QUnionBound_ofReal` application whose event
  decomposition is `compl_firstExit_subset`.

The concluded region `{V_a + c·V_d ≤ 1 − p}` with `β ^ T ≤ c` is the
`qraAffine` atom (row `[(1, V_a), (c, V_d)]`, threshold `1 − p`);
the catalogue's `c = q⁻¹` with `q ≤ e^{θT}`, `β = e^{−θ}` satisfies
`β ^ T ≤ c` — the tie is dischargeable by the same T7 leaf pattern as
`sound_QDeadlineBattery`. Fail-closed: leaving the modeled domain `D`
before reaching `R` counts as failure (`Dᶜ` joins the stop set on the bad
side); reaching `R` counts as success wherever it lands, since the formula
tie for `R` is global.

System bridge: for any System whose chart compilation at
the discrete-kernel profile yields `K`, this is the concluded
`QRegionSat`'s reading on that System's records, fail-closed form. -/
theorem sound_QUntilBattery_firstHit
    (hD : MeasurableSet D) (hDbar : MeasurableSet Dbar)
    (hA : MeasurableSet A) (hR : MeasurableSet R)
    (hVa : Measurable Va) (hVd : Measurable Vd)
    (hβ : 0 ≤ β) (hc : β ^ T ≤ c)
    (hDsub : D ⊆ Dbar)
    (hR3 : ∀ y ∈ D, y ∉ A ∪ R → K y Dbar = 1)
    (hR4 : ∀ y ∈ D, y ∉ A ∪ R → K y (D ∪ (A ∪ R)) = 1)
    (hVaint : ∀ y ∈ D, y ∉ A ∪ R → Integrable Va (K y))
    (hVadrift : ∀ y ∈ D, y ∉ A ∪ R → ∫ z, Va z ∂K y ≤ Va y)
    (hVafloor : ∀ y ∈ A, y ∈ Dbar → 1 ≤ Va y)
    (hVanonneg : ∀ y ∈ Dbar, 0 ≤ Va y)
    (hVdint : ∀ y ∈ D, y ∉ A ∪ R → Integrable Vd (K y))
    (hVddecay : ∀ y ∈ D, y ∉ A ∪ R → ∫ z, Vd z ∂K y ≤ β * Vd y)
    (hVdfloor : ∀ y ∈ Dbar, y ∉ A ∪ R → 1 ≤ Vd y)
    (hVdnonneg : ∀ y ∈ Dbar, 0 ≤ Vd y)
    (hVdM : ∀ y ∈ Dbar, Vd y ≤ M)
    (hx : x ∈ D) (hthr : Va x + c * Vd x ≤ 1 - p) :
    ENNReal.ofReal p ≤ pathMeasure K x
      {ω | ∃ n ≤ T, ω n ∈ R ∧ ∀ j < n, ω j ∉ (A ∪ R) ∪ Dᶜ}
\end{leansrc}
\declhead{BernsteinLogicAux.QtslBattery}{measure\_firstBad\_le}{\badge{BadgeGreen}{statement bound}}
\begin{leansrc}
/-- **The stopped avoid bound** (the until battery's T5 slice, stopped at
`τ_{A∪R}`): a certificate `f`, bounded in `[0, C]`,
one-step non-increasing off the stop set `Stop`, and `≥ 1` on a bad-exit set
`Bad`, bounds the probability that some *first* `Stop`-visit presents a
`Bad` state: `μ {∃ n, ω n ∈ Bad ∧ ∀ j < n, ω j ∉ Stop} ≤ f x`.
`Supermartingale.ville_ineq` consumed on the `Battery.stopped` process at
barrier `1`. Unlike `measure_reach_le` (whose floor is on the stop set
itself), the floor here is asked only on `Bad` — which is what lets the
until battery's `V_a` be silent on the target region; `Bad ⊆ Stop` is not
required by the proof. -/
theorem measure_firstBad_le {K : Kernel S S} [IsMarkovKernel K]
    {f : S → ℝ} {Stop Bad : Set S} {C : ℝ}
    (hStop : MeasurableSet Stop) (hf : Measurable f)
    (hf0 : ∀ y, 0 ≤ f y) (hfC : ∀ y, f y ≤ C)
    (hf1 : ∀ y ∈ Bad, 1 ≤ f y)
    (hdrift : ∀ y, y ∉ Stop → ∫ z, f z ∂K y ≤ f y) (x : S) :
    pathMeasure K x {ω | ∃ n, ω n ∈ Bad ∧ ∀ j < n, ω j ∉ Stop}
      ≤ ENNReal.ofReal (f x)
\end{leansrc}
\declhead{BernsteinLogicAux.QtslBattery}{measure\_missTarget\_le}{\badge{BadgeGreen}{statement bound}}
\begin{leansrc}
/-- **The deadline miss bound** (the until battery's T6 slice, at target
`A ∪ R`): under the deadline battery's rows at a target
`B`, the probability of missing `B` through time `T` is at most
`β ^ T · V x`. This is `sound_QDeadline_hit` consumed as a black box and
complemented once — the Lean proof structure mirroring the until rule's
admissible-tier unpacking, which cites the T6 battery rather than re-proving
it. The deliberately weak conclusion shape (a bound on the miss event) is
what the `QUnionBound` application consumes. -/
theorem measure_missTarget_le {K : Kernel S S} [IsMarkovKernel K]
    {D Dbar B : Set S} {V : S → ℝ} {β : ℝ} {x : S} {T : ℕ} {M : ℝ}
    (hD : MeasurableSet D) (hDbar : MeasurableSet Dbar) (hB : MeasurableSet B)
    (hV : Measurable V) (hβ : 0 ≤ β)
    (hDsub : D ⊆ Dbar)
    (hR3 : ∀ y ∈ D, y ∉ B → K y Dbar = 1)
    (hR4 : ∀ y ∈ D, y ∉ B → K y (D ∪ B) = 1)
    (hint : ∀ y ∈ D, y ∉ B → Integrable V (K y))
    (hdecay : ∀ y ∈ D, y ∉ B → ∫ z, V z ∂K y ≤ β * V y)
    (hfloor : ∀ y ∈ Dbar, y ∉ B → 1 ≤ V y)
    (hnonneg : ∀ y ∈ Dbar, 0 ≤ V y)
    (hM : ∀ y ∈ Dbar, V y ≤ M)
    (hx : x ∈ D) :
    pathMeasure K x {ω | ∀ k ≤ T, ω k ∉ B} ≤ ENNReal.ofReal (β ^ T * V x)
\end{leansrc}
\par\medskip
\textbf{The general-stochastic instance, assembled.}\enspace The
until-failure bound on the timed skeleton of the stopped presentation:
failure decomposes into an unsafe presentation or a deadline overrun,
both events of the one timed chain --- the safety half runs the
avoidance engine at zero discount (the clock drops out pointwise), the
deadline half is the deadline instance below, and the glue is measure
algebra. Headline and battery forms.

\declhead{BernsteinLogicAux.LayerD.Until}{sound\_GSHS\_until}{\badge{BadgeGreen}{statement bound}}
\begin{leansrc}
/-- **The GSHS until failure bound** (T8, headline form — the K4 event
decomposition on one chain). Hypotheses are exactly the union of the two
halves' (each inherited premise stated, never silent — the safety half's
global bound `C`, the deadline half's `θ`-row and F7 floor `ε ≤ V_d`):
the probability that the timed skeleton of the stopped presentation ever
presents an unsafe value **or** overruns the deadline is at most
`V_a(x) + e^{−θT}·V_d(x)/ε`. -/
theorem sound_GSHS_until [IsMarkovKernel M.seg] [IsMarkovKernel M.jump]
    {Va Vd : State → ℝ} {Us : Set State} {θ ε : ℝ} {T : ℝ≥0}
    (hVa : Measurable Va) (hUs : MeasurableSet Us)
    (hSegDa : SegmentDynkin M 0 Va) (hForcea : ForcedNonexpansion M Va)
    (hNonnega : NonnegV Va) {C : ℝ} (hC0 : 0 ≤ C) (hVC : ∀ y, Va y ≤ C)
    (hFloora : ∀ y ∈ Us, 1 ≤ Va y)
    (hVd : Measurable Vd) (hθ : 0 ≤ θ)
    (hSegDd : SegmentDynkin M θ Vd) (hForced : ForcedNonexpansion M Vd)
    (hNonnegd : NonnegV Vd) (hε : 0 < ε) (hFloord : ∀ y, ε ≤ Vd y)
    (x : State) :
    Battery.pathMeasure (tskel M) (TSkelState.alive (0, x))
        {ω | ∃ n, ω n ∈ tbadIn Us ∨ T < tclock (ω n)}
      ≤ ENNReal.ofReal (Va x)
        + ENNReal.ofReal (Real.exp (-(θ * T)) * Vd x / ε)
\end{leansrc}
\declhead{BernsteinLogicAux.LayerD.Until}{sound\_GSHS\_until\_battery}{\badge{BadgeGreen}{statement bound}}
\begin{leansrc}
/-- **The until battery form** (T8 as `QUntilBattery`'s GSHS instance,
complement arithmetic): with the additive threshold accounting
`p + thr_a + e^{−θT}·thr_d/ε ≤ 1` and a start under both thresholds, the
chain presents only safe values with the clock never past `T`, with
probability at least `p`. On the stop-IS-target instantiation that event
reads *`φ_s` held until the target, reached by the deadline* — up to the
stall residues (inside the event, named, emptied a.s. in-scope by
W-Z-freedom). The composite-atom rendering `{V_a + (1/(qε))·V_d ≤ 1−p}`
of the catalogue's T8 entry maps onto (`thr_a`, `thr_d`) with
`q ≤ e^{θT}`; the `qraAffine` staging decision is
schema-side and does not gate this lemma. -/
theorem sound_GSHS_until_battery [IsMarkovKernel M.seg] [IsMarkovKernel M.jump]
    {Va Vd : State → ℝ} {Us : Set State} {θ ε p thra thrd : ℝ} {T : ℝ≥0}
    (hVa : Measurable Va) (hUs : MeasurableSet Us)
    (hSegDa : SegmentDynkin M 0 Va) (hForcea : ForcedNonexpansion M Va)
    (hNonnega : NonnegV Va) {C : ℝ} (hC0 : 0 ≤ C) (hVC : ∀ y, Va y ≤ C)
    (hFloora : ∀ y ∈ Us, 1 ≤ Va y)
    (hVd : Measurable Vd) (hθ : 0 ≤ θ)
    (hSegDd : SegmentDynkin M θ Vd) (hForced : ForcedNonexpansion M Vd)
    (hNonnegd : NonnegV Vd) (hε : 0 < ε) (hFloord : ∀ y, ε ≤ Vd y)
    (hΘ : p + (thra + Real.exp (-(θ * T)) * thrd / ε) ≤ 1)
    (x : State) (hthra : Va x ≤ thra) (hthrd : Vd x ≤ thrd) :
    ENNReal.ofReal p ≤ Battery.pathMeasure (tskel M) (TSkelState.alive (0, x))
      {ω | ∀ n, ω n ∉ tbadIn Us ∧ tclock (ω n) ≤ T}
\end{leansrc}
\declhead{BernsteinLogicAux.LayerD.Until}{sound\_GSHS\_avoid\_timed}{\badge{BadgeGreen}{statement bound}}
\begin{leansrc}
/-- **The safety half on the timed chain**: the probability that the timed
skeleton ever presents a value in `Us` (alive entries *and* value-carrying
halts — kills into `Us` included) is bounded by the initial safety
certificate. This is `sound_GSHS_safety` transported to the timed chain —
at `θ = 0` the clock drops out of the certificate pointwise, so the same
engine (`Battery.measure_reach_le` = K1) runs unchanged. -/
theorem sound_GSHS_avoid_timed [IsMarkovKernel M.seg] [IsMarkovKernel M.jump]
    {Va : State → ℝ} {Us : Set State}
    (hVa : Measurable Va) (hUs : MeasurableSet Us)
    (hSegD : SegmentDynkin M 0 Va) (hForce : ForcedNonexpansion M Va)
    (hNonneg : NonnegV Va) {C : ℝ} (hC0 : 0 ≤ C) (hVC : ∀ y, Va y ≤ C)
    (hFloor : ∀ y ∈ Us, 1 ≤ Va y) (x : State) :
    Battery.pathMeasure (tskel M) (TSkelState.alive (0, x))
        {ω | ∃ n, ω n ∈ tbadIn Us}
      ≤ ENNReal.ofReal (Va x)
\end{leansrc}
\par\medskip
\textbf{The martingale-problem profile, end to end.}\enspace The
general continuous-time theorems: both slices\textquotesingle{}
formerly per-model constituents --- the safety slice\textquotesingle s
stopping theory, the deadline slice\textquotesingle s discounted
supermartingale --- are supplied by the same martingale-problem
discharge that serves the parents, composite stop set included, and
the until conclusion follows for every system whose presentation
satisfies the martingale-problem hypothesis.

\declhead{QtslAux.LayerD.MPDischarge}{sound\_GSHS\_until\_mp}{\badge{BadgeGreen}{statement bound}}
\begin{leansrc}
/-- The timed until reach bound at the martingale-problem profile: both the
avoidance certificate (rate `0`) and the deadline certificate (rate `θ`)
are presented by martingale problems. -/
theorem sound_GSHS_until_mp [IsMarkovKernel M.seg] [IsMarkovKernel M.jump]
    {Va Vd : State → ℝ} {Us : Set State} {θ ε : ℝ} {T : ℝ≥0}
    (hVa : Measurable Va) (hUs : MeasurableSet Us)
    (hmpa : ∀ x : State, Nonempty (MPFlow M 0 Va x))
    (hForcea : ForcedNonexpansion M Va)
    (hNonnega : NonnegV Va) {C : ℝ} (hC0 : 0 ≤ C) (hVC : ∀ y, Va y ≤ C)
    (hFloora : ∀ y ∈ Us, 1 ≤ Va y)
    (hVd : Measurable Vd) (hθ : 0 ≤ θ)
    (hmpd : ∀ x : State, Nonempty (MPFlow M θ Vd x))
    (hForced : ForcedNonexpansion M Vd)
    (hNonnegd : NonnegV Vd) (hε : 0 < ε) (hFloord : ∀ y, ε ≤ Vd y)
    (x : State) :
    Battery.pathMeasure (tskel M) (TSkelState.alive (0, x))
        {ω | ∃ n, ω n ∈ tbadIn Us ∨ T < tclock (ω n)}
      ≤ ENNReal.ofReal (Va x)
        + ENNReal.ofReal (Real.exp (-(θ * T)) * Vd x / ε)
\end{leansrc}
\declhead{QtslAux.LayerD.MPDischarge}{sound\_GSHS\_until\_battery\_mp}{\badge{BadgeGreen}{statement bound}}
\begin{leansrc}
/-- The until battery at the martingale-problem profile. -/
theorem sound_GSHS_until_battery_mp
    [IsMarkovKernel M.seg] [IsMarkovKernel M.jump]
    {Va Vd : State → ℝ} {Us : Set State} {θ ε p thra thrd : ℝ} {T : ℝ≥0}
    (hVa : Measurable Va) (hUs : MeasurableSet Us)
    (hmpa : ∀ x : State, Nonempty (MPFlow M 0 Va x))
    (hForcea : ForcedNonexpansion M Va)
    (hNonnega : NonnegV Va) {C : ℝ} (hC0 : 0 ≤ C) (hVC : ∀ y, Va y ≤ C)
    (hFloora : ∀ y ∈ Us, 1 ≤ Va y)
    (hVd : Measurable Vd) (hθ : 0 ≤ θ)
    (hmpd : ∀ x : State, Nonempty (MPFlow M θ Vd x))
    (hForced : ForcedNonexpansion M Vd)
    (hNonnegd : NonnegV Vd) (hε : 0 < ε) (hFloord : ∀ y, ε ≤ Vd y)
    (hΘ : p + (thra + Real.exp (-(θ * T)) * thrd / ε) ≤ 1)
    (x : State) (hthra : Va x ≤ thra) (hthrd : Vd x ≤ thrd) :
    ENNReal.ofReal p ≤ Battery.pathMeasure (tskel M) (TSkelState.alive (0, x))
      {ω | ∀ n, ω n ∉ tbadIn Us ∧ tclock (ω n) ≤ T}
\end{leansrc}
\par\medskip
\textbf{The general stratum: the until atom on the record
carrier.}\enspace The deadline-grid until event --- some grid index
at which the reach leg is hit and the avoid leg missed --- as a
countable Boolean combination of hit-events, measurable from
per-component tickets; the clocked corollary takes closed reach legs
and open avoid legs, with no compactness on the avoid side.

\declhead{SuFunctor.Site.FellAcceptance}{untilEvent}{\badge{BadgeGreen}{statement bound}}
\begin{leansrc}
/-- **The deadline-grid until event**: at some grid index, the reach
leg is hit and the avoid leg is missed.  The countable-combination
skeleton of safety-until-a-deadline at the record tier: the clocked
instantiation takes `Reach n` = (stages ≤ nth cut) × target and
`Avoid n` = (stages < nth cut) × unsafe. -/
def untilEvent (Reach Avoid : ℕ → Set α) : Set (FellCloseds α)
\end{leansrc}
\declhead{SuFunctor.Site.FellAcceptance}{mem\_untilEvent\_iff}{\badge{BadgeGreen}{statement bound}}
\begin{leansrc}
theorem mem_untilEvent_iff {Reach Avoid : ℕ → Set α}
    {F : FellCloseds α} :
    F ∈ untilEvent Reach Avoid ↔
      ∃ n, ((F : Set α) ∩ Reach n).Nonempty ∧
        Disjoint (F : Set α) (Avoid n)
\end{leansrc}
\declhead{SuFunctor.Site.FellAcceptance}{measurableSet\_untilEvent}{\badge{BadgeGreen}{statement bound}}
\begin{leansrc}
theorem measurableSet_untilEvent {Reach Avoid : ℕ → Set α}
    (h₁ : ∀ n, MeasurableSet (hitEvent (Reach n)))
    (h₂ : ∀ n, MeasurableSet (missEvent (Avoid n))) :
    MeasurableSet (untilEvent Reach Avoid)
\end{leansrc}
\declhead{SuFunctor.Site.FellAcceptance}{measurableSet\_untilEvent\_of\_isClosed\_isOpen}{\badge{BadgeGreen}{statement bound}}
\begin{leansrc}
/-- The clocked until shape: closed reach legs (a `≤`-cut of stages
times a closed target), open avoid legs (a `<`-cut times an open
unsafe region).  The avoid legs' measurability costs nothing at all —
miss events of opens are Fell-closed at every stratum. -/
theorem measurableSet_untilEvent_of_isClosed_isOpen [SigmaCompactSpace α]
    {Reach Avoid : ℕ → Set α} (h₁ : ∀ n, IsClosed (Reach n))
    (h₂ : ∀ n, IsOpen (Avoid n)) :
    MeasurableSet (untilEvent Reach Avoid)
\end{leansrc}
\par\medskip
\textbf{Marks over the record carrier.}\enspace As at the deadline
battery: the until conclusion\textquotesingle s general-stratum spec
family is persistent with the forward judgment lemma, and the
deliberately absent backward direction is the shared deadline
semantics.

\declhead{SuFunctor.Site.FellBehaviorSpec}{acceptanceSpecF}{\badge{BadgeGreen}{statement bound}}
\begin{leansrc}
/-- The spec family of a per-window acceptance event (deadline, until,
Streett): clip demands of the events' indicators. -/
def acceptanceSpecF (U : ι → Set T)
    (E : ∀ i, Set (FTraj (U i) X)) (j : ι) :
    Set (MeasPen (FTraj (U j) X))
\end{leansrc}
\declhead{SuFunctor.Site.FellBehaviorSpec}{isPersistentMF\_acceptanceSpecF}{\badge{BadgeGreen}{statement bound}}
\begin{leansrc}
theorem isPersistentMF_acceptanceSpecF (U : ι → Set T)
    (hU : ∀ i, IsOpen (U i)) (E : ∀ i, Set (FTraj (U i) X)) :
    IsPersistentMF U (acceptanceSpecF U E)
\end{leansrc}
\declhead{SuFunctor.Site.FellBehaviorSpec}{satAtM\_acceptanceSpecF\_ae}{\badge{BadgeGreen}{statement bound}}
\begin{leansrc}
/-- The forward judgment: satisfying an acceptance spec at a window
forces almost-sure acceptance there.  The backward direction is
deliberately absent — liveness events do not clip-restrict, so
subwindow demands are genuinely stronger (a reach demand within a
subwindow is a deadline); the module docstring carries the reading. -/
theorem satAtM_acceptanceSpecF_ae (U : ι → Set T)
    (E : ∀ i, Set (FTraj (U i) X)) (j : ι)
    {μ : MeasJoint (FTraj (U j) X)}
    (h : SatAtM μ (acceptanceSpecF U E j)) :
    (μ.1 : Measure (FTraj (U j) X)) (E j)ᶜ = 0
\end{leansrc}
\ruleblock{Correspondence (Steele display $\leftrightarrow$ Lean)}
\badge{BadgeAmber}{a bound Lean statement CHANGED since certification --- re-review and re-certify}
\sourcenote{The dictionary below is a reviewer-certified statement of how the mechanized statements above interpret the displayed rule: display metavariable $\mapsto$ Lean binder, each display hypothesis $\mapsto$ the premise or side condition carrying it, the display conclusion $\mapsto$ the goal --- with any deliberate difference in width between display and mechanization named outright. Certification is recorded as a hash triple over (the displayed cell, the extracted statement texts, this dictionary\textquotesingle s own text); if any of the three changes, this block turns amber until a reviewer re-certifies. The correspondence claim itself remains editorial --- reviewer-grade, not machine-checked; its machine-checked successor types each soundness statement against the rule as a structured object, and supersedes this block per rule as it lands.}
Model frame and profile as at \rulename{QSafetyBattery}
(discrete-kernel; continuous time at the martingale-problem
presentation, hash-bound below per the profile remark). One
structural honesty first:
the display presents the two battery slices by their \emph{conclusions}
--- the measure bounds $\mu_x(E_{\mathrm{bad}}) \le V_a(x)$ and
$\mu_x(E_{\mathrm{miss}}) \le \beta^T V_d(x)$ --- plus the d\'ebut
containment that \rulename{QUnionBound} consumes, while the certified
statement \texttt{sound\_QUntilBattery} takes the two slices\textquotesingle{}
row families \emph{inline} (flattened hypothesis lists over one
kernel) and derives those bounds internally (the witness names the
waypoints: \texttt{measure\_firstBad\_le},
\texttt{measure\_missTarget\_le}, with the containment bound as
\texttt{compl\_firstExit\_subset}); the admissible-tier
unpacking is realized by the \emph{proof}, which consumes the QDeadlineBattery
theorem as a black box, not by the statement\textquotesingle s shape.
\begin{itemize}
\item the concluded judgment $\mathcal{M}
  ,\ x_0 \in \{V_a + c\,V_d \le 1 - p\} \vDash
  \mathsf{SProbGE}\bigl(p,\ \mathsf{PUntil}(T, \varphi_s,
  \varphi)\bigr)$ $\mapsto$ the goal
  $\mathrm{ofReal}\ p \le$ \texttt{pathMeasure K x \{ω | ∃ n ≤ T, sat
  (ω n) ∧ ∀ j < n, sat\_s (ω j)\}} --- the until event verbatim, both
  satisfaction predicates measurability-free.
\item the two formula ties, opposite polarities (the QClosure directions),
  displayed as the inclusions $[\![\neg\varphi_s]\!] \subseteq A$ and
  $R \subseteq [\![\varphi]\!]$: $A$ \emph{over}-approximates the
  falsity set of $\varphi_s$
  (\texttt{hFs}) and $R$ \emph{under}-approximates the truth set of
  $\varphi$ (\texttt{hFt}) --- the displayed inclusions verbatim, safe
  direction on both sides.
\item the \rulename{QSafetyBattery} slice for $V_a$ (rows behind the
  displayed bound $\mu_x(E_{\mathrm{bad}}) \le V_a(x)$) $\mapsto$
  \texttt{hVadrift}/\texttt{hVaint}
  (drift at the composite stop set $(A \cup R) \cup D^{\mathrm{c}}$,
  encoded as $y \in D,\ y \notin A \cup R$),
  \texttt{hVafloor} --- on $A$ \emph{only}, the display\textquotesingle
  s ``floor on $A$ only\textquotesingle\textquotesingle{}
  --- and \texttt{hVanonneg}.
\item the \rulename{QDeadlineBattery} slice for $V_d$ (rows behind the
  displayed bound $\mu_x(E_{\mathrm{miss}}) \le \beta^T V_d(x)$) $\mapsto$
  \texttt{hVddecay}/\texttt{hVdint} ($\beta$-decay off $A \cup R$),
  \texttt{hVdfloor} (global off-target), \texttt{hVdnonneg},
  \texttt{hVdM}.
\item the rate tie $\beta^T \le c$ and its displayed sign gates
  $0 \le \beta$, $0 \le V_d(x)$ $\mapsto$ \texttt{hc}, \texttt{hβ},
  and the nonneg row read at the start point
  (\texttt{hVdnonneg} with \texttt{hx}/\texttt{hDsub}) --- the
  $V_d(x)$ gate is what lets the outer assembly pass from the
  $\beta^T V_d(x)$ miss bound to the affine row\textquotesingle s
  $c\,V_d(x)$; fully
  rational at this profile; the catalogue\textquotesingle s
  $c = q^{-1}$ is one instance.
\item the concluded \texttt{qraAffine} atom $[(1, V_a), (c, V_d)]$ at
  threshold $1 - p$ $\mapsto$ the start condition \texttt{hthr} $:$
  \texttt{Va x + c·Vd x ≤ 1 − p}, verbatim affine row.
\item side conditions and regularity as at the siblings
  (\texttt{hR3}/\texttt{hR4} at the composite stop set; the
  measurability pack undisplayed).
\end{itemize}
Width note: beyond the two safe-direction ties, neither side is wider
--- every displayed row lands on a named hypothesis, and the
statement\textquotesingle s flattening loses nothing the display
promises (each inline family is exactly its slice\textquotesingle s
row family, stopped where the display stops it).

\ruleblock{Lean verifier receipts}
\sourcenote{module \texttt{BernsteinLogicAux.QtslBattery}; bound build \texttt{5c75f9923f27\ldots} (bound from receipt 3fd830b4914bdeda, 2026-08-04T01:56:58Z). A receipt records that named checking layers accepted a build: (module, compiled-artifact SHA-256, environment SHA-256, checker layers), minted when the checking gate ran; the store receipts attesting the bound build are listed with the layers each run credited. Nothing here asks the reader to trust the store: every layer is an independently runnable program over the quoted sources --- rebuild the module and re-run the checkers to reproduce the verdict. Receipts attesting superseded builds credit no layer and are omitted from this rendering (the count of omissions, if any, is noted below).}
\begin{receipttable}
\receiptrow{3fd830b4914bdeda}{2026-08-04T01:56:58Z}{build+decorrelated}{\badge{BadgeGreen}{current}}
\receiptrow{96b18d8572451f65}{2026-08-04T03:11:31Z}{build+decorrelated}{\badge{BadgeGreen}{current}}
\end{receipttable}
\sourcenote{13 further receipt(s) in the store attest superseded builds of this module and are omitted from this rendering.}
\sourcenote{module \texttt{BernsteinLogicAux.LayerD.Until}; bound build \texttt{5abcb1cb41f6\ldots} (bound from receipt dbbf5ae370de0748, 2026-08-04T01:57:07Z). A receipt records that named checking layers accepted a build: (module, compiled-artifact SHA-256, environment SHA-256, checker layers), minted when the checking gate ran; the store receipts attesting the bound build are listed with the layers each run credited. Nothing here asks the reader to trust the store: every layer is an independently runnable program over the quoted sources --- rebuild the module and re-run the checkers to reproduce the verdict. Receipts attesting superseded builds credit no layer and are omitted from this rendering (the count of omissions, if any, is noted below).}
\begin{receipttable}
\receiptrow{dbbf5ae370de0748}{2026-08-04T01:57:07Z}{build+decorrelated}{\badge{BadgeGreen}{current}}
\receiptrow{e2dab1e692fdab41}{2026-08-04T03:11:31Z}{build+decorrelated}{\badge{BadgeGreen}{current}}
\end{receipttable}
\sourcenote{11 further receipt(s) in the store attest superseded builds of this module and are omitted from this rendering.}
\sourcenote{module \texttt{QtslAux.LayerD.MPDischarge}; bound build \texttt{baded74dd943\ldots} (bound from receipt 7228834dea10e5f9, 2026-07-24T11:04:04Z). A receipt records that named checking layers accepted a build: (module, compiled-artifact SHA-256, environment SHA-256, checker layers), minted when the checking gate ran; the store receipts attesting the bound build are listed with the layers each run credited. Nothing here asks the reader to trust the store: every layer is an independently runnable program over the quoted sources --- rebuild the module and re-run the checkers to reproduce the verdict. Receipts attesting superseded builds credit no layer and are omitted from this rendering (the count of omissions, if any, is noted below).}
\begin{receipttable}
\receiptrow{7228834dea10e5f9}{2026-07-24T11:04:04Z}{build+decorrelated}{\badge{BadgeGreen}{current}}
\end{receipttable}
\sourcenote{module \texttt{SuFunctor.Site.FellAcceptance}; bound build \texttt{651fb557bbfe\ldots} (bound from receipt aa2169560f8f571a, 2026-08-03T20:19:04Z). A receipt records that named checking layers accepted a build: (module, compiled-artifact SHA-256, environment SHA-256, checker layers), minted when the checking gate ran; the store receipts attesting the bound build are listed with the layers each run credited. Nothing here asks the reader to trust the store: every layer is an independently runnable program over the quoted sources --- rebuild the module and re-run the checkers to reproduce the verdict. Receipts attesting superseded builds credit no layer and are omitted from this rendering (the count of omissions, if any, is noted below).}
\begin{receipttable}
\receiptrow{7ad85720de6e419e}{2026-08-03T10:13:50Z}{build+decorrelated}{\badge{BadgeGreen}{current}}
\receiptrow{94a408109996d1bf}{2026-08-03T13:41:21Z}{build+decorrelated}{\badge{BadgeGreen}{current}}
\receiptrow{580a7e551fab03b6}{2026-08-03T15:39:44Z}{build+decorrelated}{\badge{BadgeGreen}{current}}
\receiptrow{aa2169560f8f571a}{2026-08-03T20:19:04Z}{build+decorrelated}{\badge{BadgeGreen}{current}}
\end{receipttable}
\sourcenote{2 further receipt(s) in the store attest superseded builds of this module and are omitted from this rendering.}
\sourcenote{module \texttt{SuFunctor.Site.FellBehaviorSpec}; bound build \texttt{345f9caf9453\ldots} (bound from receipt 72f2c3b62b6771fc, 2026-08-03T20:19:08Z). A receipt records that named checking layers accepted a build: (module, compiled-artifact SHA-256, environment SHA-256, checker layers), minted when the checking gate ran; the store receipts attesting the bound build are listed with the layers each run credited. Nothing here asks the reader to trust the store: every layer is an independently runnable program over the quoted sources --- rebuild the module and re-run the checkers to reproduce the verdict. Receipts attesting superseded builds credit no layer and are omitted from this rendering (the count of omissions, if any, is noted below).}
\begin{receipttable}
\receiptrow{d7b846e785c1d7fb}{2026-08-03T10:13:53Z}{build+decorrelated}{\badge{BadgeGreen}{current}}
\receiptrow{c0e3dc15419765aa}{2026-08-03T13:41:24Z}{build+decorrelated}{\badge{BadgeGreen}{current}}
\receiptrow{9a62b29c08b63b4f}{2026-08-03T15:39:47Z}{build+decorrelated}{\badge{BadgeGreen}{current}}
\receiptrow{72f2c3b62b6771fc}{2026-08-03T20:19:08Z}{build+decorrelated}{\badge{BadgeGreen}{current}}
\end{receipttable}
\sourcenote{2 further receipt(s) in the store attest superseded builds of this module and are omitted from this rendering.}
\layerline{\layeritem{build}{lake build (elaboration)}{\badge{BadgeGreen}{current}}\layeritem{replay}{leanchecker (official kernel replay)}{\badge{BadgeAmber}{stale only}}\layeritem{hardened}{safe\_verify (hardened screen)}{\badge{BadgeAmber}{stale only}}\layeritem{decorrelated}{lean4export + nanoda (decorrelated kernel)}{\badge{BadgeGreen}{current}}}
\sourcenote{\badge{BadgeGray}{comparator: none} no fixture in the three-implementation comparator harness (compiled / kernel-interpreted / reference-interpreted) exercises this rule.}
\ruleblock{HOL4 statement}
\absentbox{NO HOL STATEMENT YET}{The HOL4 development ports the base calculus\textquotesingle s checker
chain (theory \texttt{bern\_catalogue}, developed side by side with the
base catalogue sources, stopping below the proof-object layer); it
contains no probability or measure theory. \rulename{QUntilBattery}
lives one theory-extension edge above the base calculus and its
soundness content is path-space measure theory (a stopped-supermartingale
slice, a discounted-supermartingale slice, and a union bound over a
d\'ebut decomposition), for which no counterpart exists anywhere in the
ported development. A HOL statement of this rule would be a new
measure-theoretic campaign, not an extension of the ported checker
chain; until one exists, this box is the honest state of the edifice.
}
\ruleblock{HOL verifier receipts}
\absentbox{NONE}{no HOL verification receipt exists for this rule.}
\ruleblock{Proof}

\sourcenote{Genre: admissible rule (tier 2) --- a proof-tree-rewriting
recipe. The general admissible template (one node $\to$ a derivation
fragment over core and previously-justified rules, applied leaf-first)
is stated in the \rulename{Norm} section, Recipe~\ref{rec:norm-template},
and not restated here; this section gives the rule\textquotesingle s
instance. The verified
theorems quoted above warrant the rule at the \textbf{discrete-kernel
profile}; the profile remark at the end states what remains a
hypothesis beyond it.}

\begin{brecipe}[the \rulename{QUntilBattery} rewrite]\label{rec:qutb}
A \rulename{QUntilBattery} node carries: an avoid certificate $V_a$ with
its rows (drift off $A \cup R$, floor \emph{on $A$ only}, nonneg), a
deadline certificate $V_d$ with its rows (decay factor $\beta$ off
$A \cup R$, floor on all of the hull off the target, nonneg, bound),
the shared side
conditions (hull; fail-closed at the composite exit set $A \cup R$), the
sign gate $\beta \ge 0$ with the rate tie $\beta^{T} \le c$, and the
threshold row
$V_a(x) + c\,V_d(x) \le 1 - p$ on the concluded affine region. Here $A$
over-approximates the falsity set of the safe-side formula $\varphi_s$
and $R$ under-approximates the truth set of the target $\varphi$. The
node rewrites to a derivation fragment over three constituents:
\[
\inferrule*[Right=\rulename{QUntilBattery}]{
  \text{$V_a$ rows} \\ \text{$V_d$ rows} \\ \beta^{T} \le c \\
  V_a + c\,V_d \le 1 - p \ \text{on the region}
}{\mathcal{M},\ x_0 \in \{V_a + c\,V_d \le 1 - p\} \vDash
  \mathsf{SProbGE}\bigl(p,\ \mathsf{PUntil}(T, \varphi_s, \varphi)\bigr)}
\;\rightsquigarrow\;
\begin{array}{l}
\text{(i) a \emph{QSafetyBattery slice}: the battery pattern
  (\rulename{QSafetyBattery}, Prop.~\ref{prop:qsb-pattern})}\\
\quad\text{at the composite stop set
  $\mathit{Stop} := (A \cup R) \cup D^{\mathrm{c}}$, certificate $V_a$;}\\
\text{(ii) a \emph{QDeadlineBattery slice}: the deadline miss bound
  (\rulename{QDeadlineBattery}, Lem.~\ref{lem:qdb-tail} via
  Thm.~\ref{thm:qdb-hit})}\\
\quad\text{at target $A \cup R$, certificate $V_d$;}\\
\text{(iii) one \rulename{QUnionBound} node on the d\'ebut
  decomposition (Lem.~\ref{lem:qub-debut}).}
\end{array}
\]
Honesty note, in the \rulename{PConstExact} tradition: the two slices
are consumed through their \emph{verified slice lemmas}, not by literal
re-application of full battery nodes --- the QSafetyBattery
slice\textquotesingle s floor family is deliberately weaker than a full
safety node\textquotesingle s (floor on the bad-exit set only; see
\autoref{rem:qutb-floors}), and the QDeadlineBattery slice is consumed in
complemented (miss-bound) form. The shipped verified object is therefore
the assembled soundness lemma, whose Lean structure \emph{is} this
rewrite: it consumes the deadline theorem as a black box, applies the
union bound once, and re-proves neither battery.
\end{brecipe}

\begin{bprop}[semantic warrant, first-hit form]\label{prop:qutb-firsthit}
On the substrate of the shared preamble, suppose: the $V_a$ rows ---
drift off $A \cup R$ within $D$, floor $V_a \ge 1$ on $A \cap \bar{D}$,
nonneg on $\bar{D}$; the $V_d$ rows --- decay
$\int V_d\,dK\,y \le \beta V_d(y)$ off $A \cup R$ within $D$, floor
$V_d \ge 1$ on $\bar{D} \smallsetminus (A \cup R)$, nonneg and bound $M$
on $\bar{D}$; hull and fail-closed side conditions at exit set
$A \cup R$; $\beta \ge 0$ and the rate tie $\beta^{T} \le c$. Then for
every $x \in D$ with $V_a(x) + c\, V_d(x) \le 1 - p$,
\[
\Pr_x\bigl[\exists n \le T,\ \omega_n \in R \ \wedge\
\forall j < n,\ \omega_j \notin \mathit{Stop}\bigr] \;\ge\; p,
\qquad \mathit{Stop} = (A \cup R) \cup D^{\mathrm{c}} .
\]
\end{bprop}
\begin{proof}
Three bounds, then arithmetic.
\emph{(i, the QSafetyBattery slice.)} Clamp $V_a$ as in the safety battery, with the
floor forced only off $D \cup R$ (the certificate is silent on the
target away from $A$). The stopped process at $\mathit{Stop}$ is a
nonnegative
supermartingale, $\ge 1$ whenever a \emph{first} $\mathit{Stop}$-visit
presents a bad state: a bad state lies in
$\mathit{Stop} \smallsetminus R = (A \smallsetminus R) \cup
(D^{\mathrm{c}} \smallsetminus R)$, and every reachable one is covered
--- in $A \smallsetminus R$ it lies in $\bar D$ (a first visit is either
the start, in $D$, or one kernel step from $D \smallsetminus
\mathit{Stop}$, confined to $\bar D$ by the hull row), where the floor
row on $A \cap \bar D$ applies; in $D^{\mathrm{c}} \smallsetminus R$,
which \emph{is} the region off $D \cup R$, the forced clamp applies.
Ville\textquotesingle s inequality gives
\[
E_1 := \bigl\{\text{some first } \mathit{Stop}\text{-visit is bad}\bigr\},
\qquad \Pr_x[E_1] \le V_a(x).
\]
The floor is asked on $A$ and off $D \cup R$ only --- in particular
never on $R \smallsetminus A$, which is what lets a landing there cost
no $V_a$-budget (a landing in $A \cap R$, if the ties permit one, still
succeeds, but the avoid floor prices it as avoid-territory; the
no-charge-for-winning story is the $R \smallsetminus A$ story ---
\autoref{rem:qutb-floors}).
\emph{(ii, the QDeadlineBattery slice.)} The deadline battery\textquotesingle s miss
bound at target $A \cup R$, consumed as a black box and weakened from
missing $A \cup R$ to missing the larger $\mathit{Stop}$ --- avoiding
$\mathit{Stop} \supseteq A \cup R$ through $T$ entails avoiding
$A \cup R$ through $T$, so $E_2$ below is contained in the
black box\textquotesingle s miss event:
\[
E_2 := \bigl\{\forall k \le T,\ \omega_k \notin \mathit{Stop}\bigr\},
\qquad \Pr_x[E_2] \le \beta^{T} V_d(x) \le c\, V_d(x),
\]
the last step by the rate tie and $V_d(x) \ge 0$.
\emph{(iii, the glue.)} The d\'ebut decomposition
(Lemma~\ref{lem:qub-debut}) at $(\mathit{Stop}, R, T)$ says the
complement of the concluded first-hit event is contained in
$E_1 \cup E_2$; one \rulename{QUnionBound} application in real-valued
form (Prop.~\ref{prop:qub-ofreal}, its nonnegativity riders supplied by
the nonneg rows and by $c \ge \beta^{T} \ge 0$) gives
$\Pr_x[\text{fail}] \le V_a(x) + c\,V_d(x) \le 1 - p$. Complement ---
licit because the first-hit event is a region event over measurable
regions (the profile\textquotesingle s regularity pack), so its measure
and its complement\textquotesingle s sum to one; the weakening to
possibly-nonmeasurable satisfaction events happens only in the next
proposition, in the monotone direction.
\end{proof}

\begin{bprop}[soundness of \rulename{QUntilBattery}]\label{prop:qutb-sound}
With the two formula ties --- $A$ over-approximating where $\varphi_s$
fails, $R$ under-approximating where $\varphi$ holds --- the first-hit
event weakens to the truth of the bounded-until formula along the path:
for $x$ in the affine region,
$\Pr_x[\exists n \le T,\ \mathrm{sat}(\omega_n) \wedge
\forall j < n,\ \mathrm{sat}_s(\omega_j)] \ge p$ --- the concluded
judgment
$\mathcal{M},\ x_0 \in \{V_a + c\,V_d \le 1 - p\} \vDash
\mathsf{SProbGE}\bigl(p,\ \mathsf{PUntil}(T, \varphi_s, \varphi)\bigr)$,
its region the sublevel set $\{V_a + c\,V_d \le 1-p\}$.
\end{bprop}
\begin{proof}
On the first-hit event: the landing state is in $R$, hence satisfies
$\varphi$; every earlier state avoided $\mathit{Stop} \supseteq A$,
hence satisfies $\varphi_s$ by the contrapositive of the safe-side tie.
Monotonicity of the (outer) measure; neither satisfaction predicate is
required measurable.
\end{proof}

\begin{bremark}[two floors doing different work]\label{rem:qutb-floors}
$V_a$\textquotesingle s floor is on $A$ only: the avoid certificate is
silent on $R \smallsetminus A$, so a landing there costs no
certificate budget --- the rule does not charge for winning (a state of
$A \cap R$, when the ties permit one, is a successful landing that the
avoid floor nonetheless prices as avoid-territory; the free wins are
exactly the $R \smallsetminus A$ ones). Demand the
floor on all of $R$ too and the concluded region shrinks pointlessly:
every
reachable success state would need $V_a \ge 1$, pushing
$V_a(x) + c V_d(x)$ past any useful threshold. $V_d$\textquotesingle s
floor is on all of $\bar D$ off the target, exactly as its parent
deadline battery
requires --- weaken \emph{it} and \autoref{rem:qdb-floor}\textquotesingle s
counterexample (the frozen chain certified to arrive) imports after one
padding: take $A = \varnothing$ and $V_a \equiv 0$ (all avoid rows
trivially satisfied), and the until node \emph{is} its deadline slice,
frozen chain and all.
The asymmetry is not an optimization; it is the two
slices\textquotesingle{} different jobs made visible in their price
lists.
\end{bremark}

\begin{bremark}[fail-closed, and the rate tie]\label{rem:qutb-failclosed}
$D^{\mathrm{c}}$ joins the stop set on the \emph{bad} side: leaving the
modelled domain before reaching the target counts as failure, while an
$R$-landing succeeds wherever it lands (the target tie is global). Drop
the rate tie $\beta^{T} \le c$ and step (ii)\textquotesingle s budget
transfer fails --- the affine atom would under-budget the
\emph{certified} miss bound by the gap $(\beta^{T} - c)\, V_d(x)$
(the actual miss mass sits somewhere below the bound; it is the
transfer, not the mass, that the tie protects). The catalogued instance
$c = q^{-1}$ with $0 < q \le e^{\theta T}$, $\beta = e^{-\theta}$
satisfies the tie, and at this profile the whole certificate is exact rational
arithmetic --- no transcendental leaf anywhere; one becomes available
through the deadline battery\textquotesingle s Taylor-leaf pattern
(Prop.~\ref{prop:qdb-catalogue}) when $\theta$-rates are wanted.
\end{bremark}

\begin{bremark}[tier, argued]\label{rem:qutb-tier}
Admissible, not alias: three premise sub-derivations (two battery
slices and a union bound), not a respelling of one. And deliberately
not monolithic core: a third self-contained soundness lemma would
duplicate two batteries\textquotesingle{} arguments line by line ---
the assembled proof instead consumes the deadline theorem as a black
box and the union bound as one application, which is this
catalogue\textquotesingle s admissibility discipline enforced at the
Lean level. The concluded region is the affine atom with rows
$[(1, V_a), (c, V_d)]$ and threshold $1 - p$ --- the
region-vocabulary generalization whose $k = 1$ case is the ordinary
sublevel atom.
\end{bremark}

\begin{bremark}[\textbf{profile} --- what is proven, at which
presentations]
\label{rem:qutb-profile}
Proven: the discrete-kernel profile, sorry-free
(\texttt{sound\_QUntilBattery\_firstHit},
\texttt{sound\_QUntilBattery}), with both slice lemmas and the union
bound as named constituents. Also proven: the continuous-time profile
at the martingale-problem presentation
(\texttt{sound\_GSHS\_until\_mp},
\texttt{sound\_GSHS\_until\_battery\_mp}, module
\texttt{QtslAux.LayerD.MPDischarge} of the semantic-universe
strand\textquotesingle s auxiliary package) --- the recipe above is
profile-independent in \emph{shape}, and both
parents\textquotesingle{} formerly filed constituents (the safety
slice\textquotesingle s stopping theory, the deadline
slice\textquotesingle s discounted supermartingale) are supplied by
the same martingale-problem discharge that serves the parents,
composite stop set included. The martingale-problem theorems are
hash-bound by this catalogue\textquotesingle s binder via the
qtsl-aux package registration, family-wide, so the dashboard row
carries no conditional mark --- each named profile is proven, and
its theorems are bound.
\end{bremark}

\begin{anchors}
Recipe \ref{rec:qutb}, slice (i) & \texttt{measure\_firstBad\_le} (Ville via \texttt{Supermartingale.ville\_ineq} on the stopped process; floor asked on the bad-exit set only)\\
Recipe \ref{rec:qutb}, slice (ii) & \texttt{measure\_missTarget\_le} (black-boxing \texttt{sound\_QDeadline\_hit}, complemented once)\\
Recipe \ref{rec:qutb}, glue (iii) & \texttt{sound\_QUnionBound\_ofReal} on \texttt{compl\_firstExit\_subset}\\
Proposition \ref{prop:qutb-firsthit}: the assembly & \texttt{sound\_QUntilBattery\_firstHit}\\
Proposition \ref{prop:qutb-sound}: formula ties & \texttt{sound\_QUntilBattery}\\
the concluded affine atom & region rows $[(1, V_a), (c, V_d)]$, threshold $1-p$; rate-tie hypothesis \texttt{hc} : $\beta^{T} \le c$\\
\end{anchors}


\rulesection{44}{QStreettBattery}{core rule}
\label{blcatrule-QStreettBattery}
\ruleblock{The rule (Steele metanotation)}
\sourcenote{as realized in the mechanized development; source row: \texttt{\texttt{QTSL-rules.toml}} --- \emph{QStreettBattery}. The metanotation statement is typeset by hand for notational quality; its mechanized ground truth is the Lean statement below, extracted from the sources at every build of this book and hash-bound. The rule \emph{roster} is machine-checked (every mechanized rule tag round-trips to exactly one catalogue row); the statement \emph{text} is editorial, checked by eye against the extracted Lean beside it.}
\begin{mathpar}
\inferrule*{
  V_0,\ (W_i)_{i \le k},\ (V_i)_{i \in A} : \mathsf{QCertFam}; \quad A \subseteq \{1, \ldots, k\} \ \text{(the pairs posting absorbing certificates)}; \quad \varepsilon_0 \in \Q_{>0}; \quad (\gamma_i)_{i \in A} \in \Q_{>0} \\
  \text{the \rulename{QSafetyBattery} rows for } V_0 \ \text{at } U := \mathrm{cmpl}(I) \ \text{(interior, guard, floor on } U \cap \bar{D}\text{, nonneg on } \bar{D}\text{)} \\
  \bigl\{S \smallsetminus [\![\rho_i]\!] \vdash \mathcal{A}_m V_{i,m} \le 0\bigr\}_{i \in A} \ \text{(absorbing drift: off-request, on all of } S\text{)} \\
  \bigl\{S \vdash 0 \le V_i \le C^{a}_i\bigr\}_{i \in A} \ \text{(absorbing bounds, on all of } S\text{)}, \quad \bigl\{[\![\rho_i]\!] \vdash V_i \ge 1\bigr\}_{i \in A} \ \text{(absorbing floors)} \\
  \bigl\{([\![\varphi_i]\!] \cup J_i)^{\complement} \vdash \mathcal{A}_m W_{i,m} \le -\varepsilon_0\bigr\}_{i \le k} \ \text{(ranking decay)}, \quad J_i := \{V_i \le 1 - \gamma_i\} \cap (I \smallsetminus [\![\rho_i]\!]) \ \text{for } i \in A; \quad J_i := \emptyset \ \text{for } i \notin A \\
  W_i \equiv 0 \ \text{on target cells (containment)}; \quad \bigl\{\bar{D} \vdash W_i \ge 0\bigr\}_{i \le k} \ \text{(ranking nonnegs)} \\
  \bigl\{\bar{D} \vdash W_i \le C_i\bigr\}_{i \le k}, \quad C_i \in \Q \\
  p + y \le 1 \ \text{(threshold tie)}
}{
  \mathcal{M},\ x_0 \in \{V_0 \le y\} \vDash \mathsf{SProbGE}\bigl(p,\ \mathsf{PDwellStreett}(\mathsf{PairBase}(\rho_1, \varphi_1), \ldots, \mathsf{PairBase}(\rho_k, \varphi_k))\bigr)
}
\end{mathpar}

\ruleflavor{a region written left of the turnstile scopes the row ($R \vdash P$: $P$ at every point of $R$); the certificate families enter as certificate data --- shared references with no claim content; the rows impose every claim. $J_i$ is \emph{defined, not supplied}: the mixed soundness theorem defines it from the posted absorbing data, with the $J_i := \emptyset$ case (absorbing family omitted, $i \notin A$) the conservative scope --- the baked node carries no absorbing-family fields yet, so this optional face is not yet wire-representable (a named schema divergence). $V_0$\textquotesingle s own nonnegativity rides the \rulename{QSafetyBattery} kit row, stated once there. The margin premise $\gamma_i > 0$ together with the absorbing floor forces $J_i \cap [\![\rho_i]\!] = \emptyset$; no disjointness row is asked. The target-cell containment $W_i \equiv 0$ is recomputed by the checker. The threshold tie $p + y \le 1$ is the soundness theorem\textquotesingle s own hypothesis, displayed as a side condition: the baked Streett node carries no tie proof field yet (contrast the deadline and until nodes) --- a named schema divergence, not a checked leaf. The conclusion is the region-conditioned judgment in constructor vocabulary, as at \rulename{QSafetyBattery}; the base-exact request consumption is the $\mathsf{PairBase}$ constructor.}
\par\medskip
\rulename{QStreettBattery} concludes a probabilistic Streett-acceptance
judgment: with probability at least $p$, a path of the modelled system
satisfies every pair of the Streett condition
$\mathsf{PDwellStreett}(\mathsf{PairBase}(\rho_1, \varphi_1), \ldots,
\mathsf{PairBase}(\rho_k, \varphi_k))$ (surface syntax
$\mathsf{Streett}((\rho_1, \varphi_1), \ldots, (\rho_k, \varphi_k))$) ---
for each pair $i$, if \emph{request} states (those satisfying $\rho_i$)
occur infinitely often along the path, then \emph{response} states
(those satisfying $\varphi_i$) also occur infinitely often. This is the
shape of liveness-under-fairness obligations, and it subsumes recurrence
and stability claims as degenerate pairs (the sub-rule sections that
follow this one). A producer discharges such a specification by
synthesizing the certificate family
$(V_0;\ (V_i, \gamma_i, W_i)_{i \le k};\ \varepsilon_0)$: one safety
certificate $V_0$ confining probability mass to the operating region,
and, per pair, a ranking certificate $W_i$ that pays out a fixed rational
$\varepsilon_0 > 0$ at every un-responded request --- plus, optionally,
absorbing certificates $(V_i, \gamma_i)$ licensing the sharper ranking
scope discussed in the proof. Every premise row is an inherited
base-calculus sequent (a Bernstein-cone-family polynomial inequality the
base kernel replays), so the bundle adds \emph{zero} kernel arithmetic:
its entire content is the supermartingale assembly one judgment level
up, which is where its soundness lemma lives.

What the checker recomputes rather than trusts: the target-cell
containment behind the $W_i \equiv 0$ convention (the collapse
side-condition of the \rulename{QStreettTargetCellCollapse} section), and
--- when the absorbing face becomes wire-representable --- the carved-out
set $J_i := \{V_i \le 1 - \gamma_i\} \cap
(I \smallsetminus [\![\rho_i]\!])$, $I$ the display's confinement
region $\{V_0 < 1\}$: $J_i$ is \emph{defined}, never supplied by the
producer, but today it is defined by the soundness theorem from posted
theorem data --- the baked node carries no absorbing-family fields, so no
enacted recomputation exists yet. Also recomputed: the
region-hull mass side-conditions that make the certificate scopes
fail-closed. The threshold tie $p + y \le 1$ is displayed as a side
condition, not a checked leaf: it is the soundness
theorem\textquotesingle s own hypothesis, and the baked Streett node
carries no tie proof field yet (contrast the deadline and until
nodes) --- a named schema divergence, as the rule\textquotesingle s
flavor text states. What would break without the hypotheses: the nonnegativity
rows for all three certificate families are load-bearing (a signed
certificate can subsidize probability-mass leakage --- the recorded
two-branch counterexample at the quantitative-safety rule); the
strictness $\varepsilon_0 > 0$ is load-bearing ($W \equiv 1$ satisfies
every row at $\varepsilon_0 = 0$ while no response is ever forced); the
floor and hull rows pair up fail-closed. Each is remarked, with its
counterexample, in the proof.

Scope, stated plainly: the statements bound and receipt-credited in
this edition are the five
excursion-cluster theorems quoted below --- the analytic engine of the
typeset proof --- together with the end-to-end
assembly \texttt{sound\_QStreettBattery} (the conservative-scope bundle
theorem over \texttt{Kernel.traj}) and its discharged carved-scope
corollary \texttt{sound\_QStreettBattery\_carved\_discharged} (absorbing
certificates supplied as kernel-level rows, composed through the
absorption discharge), both in the auxiliary module
\texttt{BernsteinLogicAux.QStreett}. These cover the rule at the
\emph{discrete-kernel profile} (a time-homogeneous Markov chain
presented by a single step kernel); the complement-request tie
\texttt{sound\_QStreettBattery\_tie} is likewise bound in this edition,
under the \rulename{QStreettCmplRequest} entry that carries its
event-level kernel; the parameterized
\texttt{sound\_QStreettBattery\_carved} and the absorption discharge
\texttt{absorptionBound\_of\_kernel} remain machine-checked and cited,
not bound. The continuous-time transport that used to sit beyond that
--- the residue this rule shared with its safety and deadline
siblings, and this entry\textquotesingle s former conditional mark ---
is discharged and bound: at the martingale-problem profile the whole
bundle is proven, confinement and zero-decrement ranking rows
(\texttt{sound\_GSHS\_streett\_bundle\_mp}) together with the
strict-decrement rows from strict presentations
(\texttt{sound\_GSHS\_streett\_bundle\_mp\_strict}, whose ranking-row
input is produced by \texttt{rankingSegmentDynkin\_of\_mp\_strict}
from the pathwise progress row of \texttt{MPFlowStrict}), leaving no
analytic per-model hypothesis in the bundle. All four are hash-bound
below via the registered qtsl-aux package, so the dashboard row
carries no conditional mark: each named profile is proven, and its
theorems are bound.

\cataloguenote{Two design notes travel with the rule. First, schema: the bundle
rule\textquotesingle s rows are baked in the extension theory
(per-pair rank certificates and nonneg proof lists); the source
catalogue\textquotesingle s old MISSING-ROW finding is discharged. Second, why Streett is the
\emph{primitive} acceptance condition here: a B\"uchi condition is
exactly a one-pair Streett condition with the trivially-true request,
and the degenerate-pair specialization erases the resulting tautological
rows --- so a Streett-primitive compile pays no premium on B\"uchi-shaped
workloads, while a B\"uchi-primitive compile would pay an
exponential-in-$k$ automaton re-encoding on Streett-shaped ones. The
sub-rules of the following sections (degenerate request, complement
request, target-cell collapse) are the compile-level faces of this
design.
}
\ruleblock{Lean statement}
\declhead{BernsteinLogicAux.Excursions}{Supermartingale.exists\_ae\_tendsto\_of\_nonneg}{\badge{BadgeGreen}{statement bound}}
\begin{leansrc}
/-- **Almost-everywhere convergence of nonnegative supermartingales** (K5a): a
nonnegative supermartingale is automatically bounded in `L¹` (its expectation is
nonincreasing), so Doob's convergence theorem applies with no further hypothesis. Dual
of `MeasureTheory.Submartingale.exists_ae_tendsto_of_bdd` via `Supermartingale.neg`. -/
theorem Supermartingale.exists_ae_tendsto_of_nonneg [IsFiniteMeasure μ] {f : ℕ → Ω → ℝ}
    (hf : Supermartingale f 𝒢 μ) (hnonneg : 0 ≤ f) :
    ∀ᵐ ω ∂μ, ∃ c, Tendsto (fun n => f n ω) atTop (𝓝 c)
\end{leansrc}
\declhead{BernsteinLogicAux.Excursions}{ranking\_hitting\_sum\_le}{\badge{BadgeGreen}{statement bound}}
\begin{leansrc}
/-- **Ranking hitting-time bound** (K5b, layer form): `ε₀ · 𝔼[τ_B] ≤ 𝔼[V 0]`, the
expected hitting time decomposed by the layer-cake formula as
`𝔼[τ_B] = ∑' k, μ {τ_B > k}` over the strict-survival events
`{τ_B > k} = {∀ j < k + 1, X j ∉ B}`. This is the certificate-facing form: the ranking
rows certify `hdrift`, the nonneg rows certify `hnonneg`, and the conclusion is a bound
on the expected number of steps before reaching `B`. -/
theorem ranking_hitting_sum_le (hε : 0 ≤ ε₀)
    (hint : ∀ n, Integrable (V n) μ) (hnonneg : 0 ≤ V)
    (hmB : ∀ n, MeasurableSet[𝒢 n] {ω | X n ω ∈ B})
    (hdrift : ∀ n, ∀ᵐ ω ∂μ, X n ω ∉ B → (μ[V (n + 1) | 𝒢 n]) ω ≤ V n ω - ε₀) :
    ENNReal.ofReal ε₀ * ∑' k, μ {ω | ∀ j < k + 1, X j ω ∉ B}
      ≤ ENNReal.ofReal (∫ ω, V 0 ω ∂μ)
\end{leansrc}
\declhead{BernsteinLogicAux.Excursions}{ae\_hits\_of\_ranking}{\badge{BadgeGreen}{statement bound}}
\begin{leansrc}
/-- **Almost-sure hitting from a ranking certificate** (K5b, qualitative form): a
strictly positive ranking drift off `B` forces the state process to reach `B` almost
surely — the summability of the survival probabilities from `ranking_hitting_sum_le`
sends them to zero. -/
theorem ae_hits_of_ranking (hε : 0 < ε₀)
    (hint : ∀ n, Integrable (V n) μ) (hnonneg : 0 ≤ V)
    (hmB : ∀ n, MeasurableSet[𝒢 n] {ω | X n ω ∈ B})
    (hdrift : ∀ n, ∀ᵐ ω ∂μ, X n ω ∉ B → (μ[V (n + 1) | 𝒢 n]) ω ≤ V n ω - ε₀) :
    ∀ᵐ ω ∂μ, ∃ n, X n ω ∈ B
\end{leansrc}
\declhead{BernsteinLogicAux.Excursions}{streett\_window\_sum\_le}{\badge{BadgeGreen}{statement bound}}
\begin{leansrc}
/-- **Summability of un-responded requests** (K5c, quantitative core): for each window
start `N`, `ε₀ · ∑_{j} μ(no response in [N, N+j] ∧ request at N+j) ≤ 𝔼[V N]` — the form
the Borel–Cantelli lemma consumes. -/
theorem streett_window_sum_le (hε : 0 ≤ ε₀)
    (hint : ∀ n, Integrable (V n) μ) (hnonneg : 0 ≤ V)
    (hmR : ∀ n, MeasurableSet[𝒢 n] {ω | X n ω ∈ R})
    (hmG : ∀ n, MeasurableSet[𝒢 n] {ω | X n ω ∈ G})
    (hdrift : ∀ n, ∀ᵐ ω ∂μ, X n ω ∉ G → (μ[V (n + 1) | 𝒢 n]) ω ≤ V n ω)
    (hdriftR : ∀ n, ∀ᵐ ω ∂μ,
      X n ω ∈ R → X n ω ∉ G → (μ[V (n + 1) | 𝒢 n]) ω ≤ V n ω - ε₀)
    (N : ℕ) :
    ENNReal.ofReal ε₀ * ∑' j : ℕ,
        μ ({ω | ∀ k, N ≤ k → k < N + j + 1 → X k ω ∉ G} ∩ {ω | X (N + j) ω ∈ R})
      ≤ ENNReal.ofReal (∫ ω, V N ω ∂μ)
\end{leansrc}
\declhead{BernsteinLogicAux.Excursions}{streett\_ae\_frequently\_response}{\badge{BadgeGreen}{statement bound}}
\begin{leansrc}
/-- **Streett excursion assembly** (K5c): if `V` is nonincreasing in conditional
expectation off the response set `G` and decreases by `ε₀ > 0` on requests (`R`, off
`G`), then almost surely infinitely many requests imply infinitely many responses —
the probabilistic content of a Streett pair `(R, G)`. Proof: by
`streett_window_sum_le` the un-responded-request events of each window are summable, so
their limsup is null (`MeasureTheory.measure_limsup_atTop_eq_zero`, Borel–Cantelli);
a path with infinitely many requests and no response after `N` lies in window `N`'s
limsup, and the union over `N` is null. -/
theorem streett_ae_frequently_response (hε : 0 < ε₀)
    (hint : ∀ n, Integrable (V n) μ) (hnonneg : 0 ≤ V)
    (hmR : ∀ n, MeasurableSet[𝒢 n] {ω | X n ω ∈ R})
    (hmG : ∀ n, MeasurableSet[𝒢 n] {ω | X n ω ∈ G})
    (hdrift : ∀ n, ∀ᵐ ω ∂μ, X n ω ∉ G → (μ[V (n + 1) | 𝒢 n]) ω ≤ V n ω)
    (hdriftR : ∀ n, ∀ᵐ ω ∂μ,
      X n ω ∈ R → X n ω ∉ G → (μ[V (n + 1) | 𝒢 n]) ω ≤ V n ω - ε₀) :
    ∀ᵐ ω ∂μ, (∃ᶠ n in atTop, X n ω ∈ R) → ∃ᶠ n in atTop, X n ω ∈ G
\end{leansrc}
\declhead{BernsteinLogicAux.QStreett}{sound\_QStreettBattery}{\badge{BadgeGreen}{statement bound}}
\begin{leansrc}
/-- **`sound_QStreettBattery`, discrete profile** (catalogue T9, rev 3
`thm:bundle-sound` transported; the assembly the receipt ledger carried as
FILED). Hypotheses: the `V₀` battery rows — quantitative safety of the
confinement region `D \ U` (the catalogue's `I := {V₀ < 1}` is the instance
`U := {y | 1 ≤ V₀ y}`, whose floor row is tautological) — and, per pair
`i : ι` (countable; a Streett bundle's finite `k` is an instance), the ranking
rows for `Wᵢ` at the **conservative scope** of the design note §4.1's ranking
row ("conservative on `I \ ⟦φᵢ⟧`"): nonincrease off the response region,
`εᵢ`-decrease at requests, on the confinement region. Conclusion:
`Pr[⋂ᵢ (Fin(Rqᵢ) ∪ Inf(Gᵢ))] ≥ p` from the start region `D ∩ {V₀ ≤ thr}` with
`p + thr ≤ 1` — the compile's atom `{V₀ ≤ 1 − p}` is `thr = 1 − p`.

The absorbing certificates `(Vᵢ, γᵢ)` of the full bundle datum do not appear:
at the conservative ranking scope they are not consumed (they license only the
`Jᵢ` carve-out — see `sound_QStreettBattery_carved`). Presentation→semantics
bridge and GSHS residue: module docstring. -/
theorem sound_QStreettBattery {Rq G : ι → Set S} {W : ι → S → ℝ} {ε C : ι → ℝ}
    (hD : MeasurableSet D) (hDbar : MeasurableSet Dbar) (hU : MeasurableSet U)
    (hV₀ : Measurable V₀)
    (hRq : ∀ i, MeasurableSet (Rq i)) (hG : ∀ i, MeasurableSet (G i))
    (hW : ∀ i, Measurable (W i))
    (hΘ : p + thr ≤ 1)
    (hDsub : D ⊆ Dbar)
    (hR3 : ∀ y ∈ D, y ∉ U → K y Dbar = 1)
    (hR4 : ∀ y ∈ D, y ∉ U → K y (D ∪ U) = 1)
    (hV₀int : ∀ y ∈ D, y ∉ U → Integrable V₀ (K y))
    (hV₀drift : ∀ y ∈ D, y ∉ U → ∫ z, V₀ z ∂K y ≤ V₀ y)
    (hV₀floor : ∀ y ∈ U, y ∈ Dbar → 1 ≤ V₀ y)
    (hV₀nonneg : ∀ y ∈ Dbar, 0 ≤ V₀ y)
    (hε : ∀ i, 0 < ε i)
    (hW0 : ∀ i, ∀ y ∈ Dbar, 0 ≤ W i y) (hWC : ∀ i, ∀ y ∈ Dbar, W i y ≤ C i)
    (hWint : ∀ i, ∀ y ∈ D, y ∉ U → y ∉ G i → Integrable (W i) (K y))
    (hWdrift : ∀ i, ∀ y ∈ D, y ∉ U → y ∉ G i → ∫ z, W i z ∂K y ≤ W i y)
    (hWdriftR : ∀ i, ∀ y ∈ D, y ∉ U → y ∈ Rq i → y ∉ G i →
      ∫ z, W i z ∂K y ≤ W i y - ε i)
    (hx : x ∈ D) (hthr : V₀ x ≤ thr) :
    ENNReal.ofReal p ≤ pathMeasure K x
      {ω | ∀ i, (∃ᶠ n in atTop, ω n ∈ Rq i) → ∃ᶠ n in atTop, ω n ∈ G i}
\end{leansrc}
\declhead{BernsteinLogicAux.QStreett}{sound\_QStreettBattery\_carved\_discharged}{\badge{BadgeGreen}{statement bound}}
\begin{leansrc}
/-- **`sound_QStreettBattery_carved`, discharged form** — the unconditional
corollary: the full bundle datum with the absorbing certificates `(Vᵢ, γᵢ)`
supplied as kernel-level rows (global scope, per the section docstring)
instead of the named hypothesis. The parameterized
`sound_QStreettBattery_carved` remains the citable assembly; this composes it
with `absorptionBound_of_kernel` per pair. -/
theorem sound_QStreettBattery_carved_discharged {D Dbar U : Set S} {V₀ : S → ℝ}
    {p thr : ℝ} {x : S} {ι : Type*} [Countable ι] {Rq G J : ι → Set S}
    {W Va : ι → S → ℝ} {ε C γ Ca : ι → ℝ}
    (hD : MeasurableSet D) (hDbar : MeasurableSet Dbar) (hU : MeasurableSet U)
    (hV₀ : Measurable V₀)
    (hRq : ∀ i, MeasurableSet (Rq i)) (hG : ∀ i, MeasurableSet (G i))
    (hJ : ∀ i, MeasurableSet (J i)) (hW : ∀ i, Measurable (W i))
    (hΘ : p + thr ≤ 1)
    (hDsub : D ⊆ Dbar)
    (hR3 : ∀ y ∈ D, y ∉ U → K y Dbar = 1)
    (hR4 : ∀ y ∈ D, y ∉ U → K y (D ∪ U) = 1)
    (hV₀int : ∀ y ∈ D, y ∉ U → Integrable V₀ (K y))
    (hV₀drift : ∀ y ∈ D, y ∉ U → ∫ z, V₀ z ∂K y ≤ V₀ y)
    (hV₀floor : ∀ y ∈ U, y ∈ Dbar → 1 ≤ V₀ y)
    (hV₀nonneg : ∀ y ∈ Dbar, 0 ≤ V₀ y)
    (hε : ∀ i, 0 < ε i) (hγ : ∀ i, 0 < γ i)
    (hW0 : ∀ i, ∀ y ∈ Dbar, 0 ≤ W i y) (hWC : ∀ i, ∀ y ∈ Dbar, W i y ≤ C i)
    (hWint : ∀ i, ∀ y ∈ D, y ∉ U → y ∉ G i → y ∉ J i → Integrable (W i) (K y))
    (hWdrift : ∀ i, ∀ y ∈ D, y ∉ U → y ∉ G i → y ∉ J i →
      ∫ z, W i z ∂K y ≤ W i y)
    (hWdriftR : ∀ i, ∀ y ∈ D, y ∉ U → y ∈ Rq i → y ∉ G i → y ∉ J i →
      ∫ z, W i z ∂K y ≤ W i y - ε i)
    (hVa : ∀ i, Measurable (Va i))
    (hVa0 : ∀ i y, 0 ≤ Va i y) (hVaC : ∀ i y, Va i y ≤ Ca i)
    (hVafloor : ∀ i, ∀ y ∈ Rq i, 1 ≤ Va i y)
    (hVadrift : ∀ i y, y ∉ Rq i → ∫ z, Va i z ∂K y ≤ Va i y)
    (hVaJ : ∀ i, ∀ y ∈ J i, Va i y ≤ 1 - γ i)
    (hx : x ∈ D) (hthr : V₀ x ≤ thr) :
    ENNReal.ofReal p ≤ pathMeasure K x
      {ω | ∀ i, (∃ᶠ n in atTop, ω n ∈ Rq i) → ∃ᶠ n in atTop, ω n ∈ G i}
\end{leansrc}
\par\medskip
\textbf{The general-stochastic instance, assembled.}\enspace The
bundle on the segment model\textquotesingle s plain skeleton chain,
obtained by instantiating the discrete bundle above at the skeleton
kernel: alive-only region transport, scoped per-pair ranking fields
with additive decrement (subtraction-free extended-nonnegative form),
and the three fail-closed containment rows enter as per-model fields;
the discrete engine\textquotesingle s localization and windowed
telescoping do the rest, restart-free. Plain and confined forms.

\declhead{BernsteinLogicAux.LayerD.Streett}{sound\_GSHS\_streett\_bundle}{\badge{BadgeGreen}{statement bound}}
\begin{leansrc}
/-- **`sound_QStreettBattery`, GSHS instance** (T9, conservative ranking
scope) — the discrete-statement-shaped corollary: the chain-level Streett
event alone. The `_confined` form is the semantically loaded one; this one
matches `QStreett.sound_QStreettBattery` line for line. -/
theorem sound_GSHS_streett_bundle
    [IsMarkovKernel M.seg] [IsMarkovKernel M.jump]
    (hD : MeasurableSet D) (hDbar : MeasurableSet Dbar) (hU : MeasurableSet U)
    (hV₀ : Measurable V₀)
    (hRq : ∀ i, MeasurableSet (Rq i)) (hG : ∀ i, MeasurableSet (G i))
    (hW : ∀ i, Measurable (W i))
    (hΘ : p + thr ≤ 1) (hDsub : D ⊆ Dbar)
    (hLive : NoVanish M (D \ U)) (hHalt : HaltsCaught M (D \ U) (U ∩ Dbar))
    (hReset : ResetsInto M (D \ U) ((D ∪ U) ∩ Dbar))
    (hSegD₀ : RankingSegmentDynkin M (D \ U) V₀ 0)
    (hForce₀ : ForcedNonexpansion M V₀) (hNonneg₀ : NonnegV V₀)
    (hC₀0 : 0 ≤ C₀) (hV₀C : ∀ y, V₀ y ≤ C₀)
    (hFloor₀ : ∀ y ∈ U, y ∈ Dbar → 1 ≤ V₀ y)
    (hε : ∀ i, 0 < ε i)
    (hSegDW : ∀ i, RankingSegmentDynkin M ((D \ U) \ G i) (W i) 0)
    (hSegDWR : ∀ i, RankingSegmentDynkin M (((D \ U) ∩ Rq i) \ G i) (W i) (ε i))
    (hForceW : ∀ i, ForcedNonexpansion M (W i)) (hNonnegW : ∀ i, NonnegV (W i))
    (hC0 : ∀ i, 0 ≤ C i) (hWC : ∀ i, ∀ y, W i y ≤ C i)
    (hx : x ∈ D) (hthr : V₀ x ≤ thr) :
    ENNReal.ofReal p ≤ Battery.pathMeasure (skel M) (SkelState.alive x)
      {ω | ∀ i, (∃ᶠ n in atTop, ω n ∈ skelIn (Rq i)) →
        ∃ᶠ n in atTop, ω n ∈ skelIn (G i)}
\end{leansrc}
\declhead{BernsteinLogicAux.LayerD.Streett}{sound\_GSHS\_streett\_bundle\_confined}{\badge{BadgeGreen}{statement bound}}
\begin{leansrc}
/-- **`sound_QStreettBattery`, GSHS instance, confined form** (T9; assault
doc §2.7 Theorem G at the conservative ranking scope, mechanized
restart-free — module docstring for why G1/G2's epoch machinery is not
consumed). Hypothesis fields are segment-level and per-model: the three
containment rows, the θ = 0 confinement field for `V₀`, and per pair the
scoped ranking fields for `W i` (nonincrease off the response region,
`ε i`-decrease at request entries — per-step shape, module docstring for
the rate-form conversion). Conclusion: with probability at least `p` the
skeleton chain of `Rep(P)` stays alive in `D \ U` forever **and** every
certified pair satisfies the chain-level Streett condition. Semantic
transport of the chain events: module docstring (T9b monotonicity for
requests, entry reads for responses, W-Z for time-cofinality). -/
theorem sound_GSHS_streett_bundle_confined
    [IsMarkovKernel M.seg] [IsMarkovKernel M.jump]
    (hD : MeasurableSet D) (hDbar : MeasurableSet Dbar) (hU : MeasurableSet U)
    (hV₀ : Measurable V₀)
    (hRq : ∀ i, MeasurableSet (Rq i)) (hG : ∀ i, MeasurableSet (G i))
    (hW : ∀ i, Measurable (W i))
    (hΘ : p + thr ≤ 1) (hDsub : D ⊆ Dbar)
    (hLive : NoVanish M (D \ U)) (hHalt : HaltsCaught M (D \ U) (U ∩ Dbar))
    (hReset : ResetsInto M (D \ U) ((D ∪ U) ∩ Dbar))
    (hSegD₀ : RankingSegmentDynkin M (D \ U) V₀ 0)
    (hForce₀ : ForcedNonexpansion M V₀) (hNonneg₀ : NonnegV V₀)
    (hC₀0 : 0 ≤ C₀) (hV₀C : ∀ y, V₀ y ≤ C₀)
    (hFloor₀ : ∀ y ∈ U, y ∈ Dbar → 1 ≤ V₀ y)
    (hε : ∀ i, 0 < ε i)
    (hSegDW : ∀ i, RankingSegmentDynkin M ((D \ U) \ G i) (W i) 0)
    (hSegDWR : ∀ i, RankingSegmentDynkin M (((D \ U) ∩ Rq i) \ G i) (W i) (ε i))
    (hForceW : ∀ i, ForcedNonexpansion M (W i)) (hNonnegW : ∀ i, NonnegV (W i))
    (hC0 : ∀ i, 0 ≤ C i) (hWC : ∀ i, ∀ y, W i y ≤ C i)
    (hx : x ∈ D) (hthr : V₀ x ≤ thr) :
    ENNReal.ofReal p ≤ Battery.pathMeasure (skel M) (SkelState.alive x)
      {ω | (∀ n, ω n ∈ skelIn D \ skelBad U) ∧
        ∀ i, (∃ᶠ n in atTop, ω n ∈ skelIn (Rq i)) →
          ∃ᶠ n in atTop, ω n ∈ skelIn (G i)}
\end{leansrc}
\par\medskip
\textbf{The martingale-problem profile, end to end.}\enspace The
general continuous-time theorems: the inter-jump motion enters through
its generator --- $\mathcal{A}V$ as data, tied to the process only by
the martingale hypothesis for the discounted compensated certificate
process. The mixed-profile bundle discharges the confinement half and
every zero-decrement ranking row; \texttt{MPFlowStrict} adds the one
pathwise progress row (accumulated generator drift reaching
$-\varepsilon$ by a finite sample time), from which the strict
ranking rows are derived at exactly the requested-but-ungranted
states --- and the strict bundle theorem then carries no analytic
per-model hypothesis at all.

\declhead{QtslAux.LayerD.MPDischarge}{sound\_GSHS\_streett\_bundle\_mp}{\badge{BadgeGreen}{statement bound}}
\begin{leansrc}
/-- The Streett/recurrence bundle at the martingale-problem profile,
mixed ranking scope: the confinement certificate `V₀` and every pair's
ranking certificate `W i` are presented by martingale problems, which
discharges all the `ε = 0` ranking rows. The strict-decrement rows
`hSegDWR` (`ε i > 0`) stay hypotheses *of this theorem*, for models that
discharge them by other means; `sound_GSHS_streett_bundle_mp_strict` below
discharges them too, from strict presentations (`MPFlowStrict`) on the
requested-but-ungranted region. -/
theorem sound_GSHS_streett_bundle_mp
    {ι : Type*} [Countable ι]
    {D Dbar U : Set State} {V₀ : State → ℝ} {Rq G : ι → Set State}
    {W : ι → State → ℝ} {p thr C₀ : ℝ} {ε C : ι → ℝ} {x : State}
    [IsMarkovKernel M.seg] [IsMarkovKernel M.jump]
    (hD : MeasurableSet D) (hDbar : MeasurableSet Dbar) (hU : MeasurableSet U)
    (hV₀ : Measurable V₀)
    (hRq : ∀ i, MeasurableSet (Rq i)) (hG : ∀ i, MeasurableSet (G i))
    (hW : ∀ i, Measurable (W i))
    (hΘ : p + thr ≤ 1) (hDsub : D ⊆ Dbar)
    (hLive : NoVanish M (D \ U)) (hHalt : HaltsCaught M (D \ U) (U ∩ Dbar))
    (hReset : ResetsInto M (D \ U) ((D ∪ U) ∩ Dbar))
    (hmp₀ : ∀ x : State, Nonempty (MPFlow M 0 V₀ x))
    (hForce₀ : ForcedNonexpansion M V₀) (hNonneg₀ : NonnegV V₀)
    (hC₀0 : 0 ≤ C₀) (hV₀C : ∀ y, V₀ y ≤ C₀)
    (hFloor₀ : ∀ y ∈ U, y ∈ Dbar → 1 ≤ V₀ y)
    (hε : ∀ i, 0 < ε i)
    (hmpW : ∀ i, ∀ x : State, Nonempty (MPFlow M 0 (W i) x))
    (hSegDWR : ∀ i, RankingSegmentDynkin M (((D \ U) ∩ Rq i) \ G i) (W i) (ε i))
    (hForceW : ∀ i, ForcedNonexpansion M (W i)) (hNonnegW : ∀ i, NonnegV (W i))
    (hC0 : ∀ i, 0 ≤ C i) (hWC : ∀ i, ∀ y, W i y ≤ C i)
    (hx : x ∈ D) (hthr : V₀ x ≤ thr) :
    ENNReal.ofReal p ≤ Battery.pathMeasure (skel M) (SkelState.alive x)
      {ω | ∀ i, (∃ᶠ n in atTop, ω n ∈ skelIn (Rq i)) →
        ∃ᶠ n in atTop, ω n ∈ skelIn (G i)}
\end{leansrc}
\declhead{QtslAux.LayerD.MPDischarge}{MPFlowStrict}{\badge{BadgeGreen}{statement bound}}
\begin{leansrc}
/-- **The strict martingale-problem presentation of a segment.** An
`MPFlow` together with the pathwise progress row: along every path, the
accumulated generator drift of the discounted certificate reaches `≤ -ε`
by some finite sample time. For exit paths the natural witness is the exit
time itself (`stopClip` caps the accumulation there), so the row reads
`∫₀^{T ω} mpDrift ≤ -ε`; the intended certificate-side derivation is a
strict drift bound `≤ -δ` on the segment together with a minimum dwell
time `t_min`, giving `ε = δ · t_min`. -/
structure MPFlowStrict (M : SegmentModel State) (θ : ℝ) (V : State → ℝ)
    (x : State) (ε : ℝ) extends MPFlow M θ V x where
  /-- the pathwise progress row -/
  hprog : ∀ ω, ∃ t : ℝ, 0 ≤ t ∧
    ∫ s in (0:ℝ)..(stopClip (T ω) (k ω) t),
      mpDrift θ V γ r jr AV s ω ≤ -ε
\end{leansrc}
\declhead{QtslAux.LayerD.MPDischarge}{rankingSegmentDynkin\_of\_mp\_strict}{\badge{BadgeGreen}{statement bound}}
\begin{leansrc}
/-- **The strict-decrement ranking rows at the martingale-problem
profile.** Entry states of the ranking region carrying strict
presentations at level `ε` satisfy `RankingSegmentDynkin M Sc V ε` — the
row family that `rankingSegmentDynkin_of_mp` could not produce. Note the
quantifier: only states in `Sc` need presentations (the `ε = 0` route
quantified over all states because it factored through `SegmentDynkin`). -/
theorem rankingSegmentDynkin_of_mp_strict (M : SegmentModel State)
    (V : State → ℝ) [IsSFiniteKernel M.jump] (hV : Measurable V)
    {Sc : Set State} {ε : ℝ} (hε : 0 ≤ ε)
    (h : ∀ x ∈ Sc, Nonempty (MPFlowStrict M 0 V x ε)) :
    RankingSegmentDynkin M Sc V ε
\end{leansrc}
\declhead{QtslAux.LayerD.MPDischarge}{sound\_GSHS\_streett\_bundle\_mp\_strict}{\badge{BadgeGreen}{statement bound}}
\begin{leansrc}
/-- The Streett/recurrence bundle at the **fully martingale-problem
profile**: the confinement certificate `V₀` and every pair's ranking
certificate `W i` are presented by martingale problems, with the
requested-but-ungranted region additionally carrying *strict*
presentations (`MPFlowStrict` at level `ε i`). This discharges the
`hSegDWR` hypothesis of `sound_GSHS_streett_bundle_mp` — at this profile
no analytic per-model hypothesis remains. -/
theorem sound_GSHS_streett_bundle_mp_strict
    {ι : Type*} [Countable ι]
    {D Dbar U : Set State} {V₀ : State → ℝ} {Rq G : ι → Set State}
    {W : ι → State → ℝ} {p thr C₀ : ℝ} {ε C : ι → ℝ} {x : State}
    [IsMarkovKernel M.seg] [IsMarkovKernel M.jump]
    (hD : MeasurableSet D) (hDbar : MeasurableSet Dbar) (hU : MeasurableSet U)
    (hV₀ : Measurable V₀)
    (hRq : ∀ i, MeasurableSet (Rq i)) (hG : ∀ i, MeasurableSet (G i))
    (hW : ∀ i, Measurable (W i))
    (hΘ : p + thr ≤ 1) (hDsub : D ⊆ Dbar)
    (hLive : NoVanish M (D \ U)) (hHalt : HaltsCaught M (D \ U) (U ∩ Dbar))
    (hReset : ResetsInto M (D \ U) ((D ∪ U) ∩ Dbar))
    (hmp₀ : ∀ x : State, Nonempty (MPFlow M 0 V₀ x))
    (hForce₀ : ForcedNonexpansion M V₀) (hNonneg₀ : NonnegV V₀)
    (hC₀0 : 0 ≤ C₀) (hV₀C : ∀ y, V₀ y ≤ C₀)
    (hFloor₀ : ∀ y ∈ U, y ∈ Dbar → 1 ≤ V₀ y)
    (hε : ∀ i, 0 < ε i)
    (hmpW : ∀ i, ∀ x : State, Nonempty (MPFlow M 0 (W i) x))
    (hmpWR : ∀ i, ∀ x ∈ ((D \ U) ∩ Rq i) \ G i,
      Nonempty (MPFlowStrict M 0 (W i) x (ε i)))
    (hForceW : ∀ i, ForcedNonexpansion M (W i)) (hNonnegW : ∀ i, NonnegV (W i))
    (hC0 : ∀ i, 0 ≤ C i) (hWC : ∀ i, ∀ y, W i y ≤ C i)
    (hx : x ∈ D) (hthr : V₀ x ≤ thr) :
    ENNReal.ofReal p ≤ Battery.pathMeasure (skel M) (SkelState.alive x)
      {ω | ∀ i, (∃ᶠ n in atTop, ω n ∈ skelIn (Rq i)) →
        ∃ᶠ n in atTop, ω n ∈ skelIn (G i)}
\end{leansrc}
\ruleblock{Correspondence (Steele display $\leftrightarrow$ Lean)}
\badge{BadgeGreen}{reviewer-certified against the displayed texts}
\sourcenote{The dictionary below is a reviewer-certified statement of how the mechanized statements above interpret the displayed rule: display metavariable $\mapsto$ Lean binder, each display hypothesis $\mapsto$ the premise or side condition carrying it, the display conclusion $\mapsto$ the goal --- with any deliberate difference in width between display and mechanization named outright. Certification is recorded as a hash triple over (the displayed cell, the extracted statement texts, this dictionary\textquotesingle s own text); if any of the three changes, this block turns amber until a reviewer re-certifies. The correspondence claim itself remains editorial --- reviewer-grade, not machine-checked; its machine-checked successor types each soundness statement against the rule as a structured object, and supersedes this block per rule as it lands.}
The displayed judgment is concluded by
\texttt{sound\_QStreettBattery} (second block): from the $V_0$ battery
rows on the confinement region and, per pair, the conservative-scope
ranking rows (nonincrease off the response region, $\varepsilon_i$-decay
at requests), it concludes
$\Pr[\bigcap_i (\mathrm{Fin}(Rq_i) \cup \mathrm{Inf}(G_i))] \ge p$
from the start region at \texttt{p + thr ≤ 1} --- the
display\textquotesingle s conclusion at the discrete-kernel profile
over \texttt{Kernel.traj}; \texttt{\_carved\_discharged} adds the
carved scope with the absorbing certificates consumed as kernel-level
rows. (Notation: $\mathrm{Fin}(A)$ / $\mathrm{Inf}(A)$ denote the path
events that a trajectory visits $A$ finitely / infinitely often ---
the acceptance atoms of the Hanoi Omega-Automata
format~\cite{babiak2015hoa} and of Emerson--Lei
conditions~\cite{emersonlei1987}; equivalently $\mathrm{Inf}(A) =
\{\rho \mid \mathrm{Inf}(\rho) \cap A \neq \emptyset\}$ in the
infinity-set convention, $\mathrm{Inf}(\rho)$ being the set of states
visited infinitely often. The upright $\mathrm{Fin}$ is distinct from
the sans-serif type former $\mathsf{Fin}$ of the mechanization.) The five bound excursion theorems are the K5 analytic engine the
assembly consumes, stated over an arbitrary filtered probability space
\texttt{(Ω, 𝒢, μ)} with an adapted state process \texttt{X} ---
\emph{wider} than the discrete-kernel chain. What the
display\textquotesingle s rows reach in the engine\textquotesingle s
bound set, per pair:
\begin{itemize}
\item ranking interior rows $\mathcal{A}_m W_{i,m} \le -\varepsilon_0$
  off $[\![\varphi_i]\!] \cup J_i$ $\mapsto$ the conditional-shadow
  pair in \texttt{streett\_ae\_frequently\_response}: \texttt{hdrift}
  (conditional non-increase off the response set \texttt{G}) and
  \texttt{hdriftR} (decay by \texttt{ε₀} at requests off \texttt{G})
  --- per-time a.e.\ conditional-expectation clauses, the
  probabilistic sibling of the integrated semantic-shadow forms on
  the deterministic pages; the BernCone rows reach them through the
  kernel bridge in the unbound assembly. Conclusion: \texttt{∀ᵐ ω,
  (∃ᶠ n, X n ω ∈ R) → ∃ᶠ n, X n ω ∈ G} --- the probabilistic content
  of one Streett pair, exactly.
\item nonneg rows $\mapsto$ \texttt{hnonneg : 0 ≤ V} --- and the
  roster\textquotesingle s two-branch signed-certificate
  counterexample is why the binder is there.
\item $\varepsilon_0 \in \Q_{>0}$ $\mapsto$ \texttt{hε : 0 < ε₀}
  (real, the $\Q$ cast upstream); the quantitative window/hitting
  forms (\texttt{streett\_window\_sum\_le},
  \texttt{ranking\_hitting\_sum\_le}) take the nonstrict \texttt{0 ≤
  ε₀} --- strictness is consumed exactly where Borel--Cantelli or
  summability-to-hitting needs it (\texttt{ae\_hits\_of\_ranking},
  the frequently-response theorem).
\item the measurability pack \texttt{hmR}/\texttt{hmG}
  (request/response events adapted at each time) is undisplayed
  regularity --- the probabilistic sibling of the PsdCert
  evaluability class, as on the battery pages.
\item \texttt{Supermartingale.exists\_ae\_tendsto\_of\_nonneg} (K5a)
  is engine stock: Doob convergence with the $L^1$ bound derived
  from nonnegativity, no display row corresponds.
\end{itemize}
Width note: the $V_0$ safety kit, the $W_i \le C_i$
upper-bound rows, the threshold tie $p + y \le 1$, and the
carved $J_i$ with the absorbing families $(V_i, \gamma_i)$ as
kernel-level rows are all consumed by the two bound assembly
statements; the G5 target-cell containment binds in the
\rulename{QStreettTargetCellCollapse} entry, and the complement-request
tie remains cited, not bound. The martingale-problem profile, stated
at its true extent: the confinement half and all zero-decrement
ranking rows are discharged by
\texttt{sound\_GSHS\_streett\_bundle\_mp}, and the strict-decrement
rows ($\varepsilon_i > 0$) by strict presentations ---
\texttt{MPFlowStrict} extends \texttt{MPFlow} by one pathwise progress
row (along every path the accumulated generator drift of the
discounted certificate reaches $-\varepsilon$ by a finite sample time;
the certificate-side derivation is $\varepsilon = \delta \cdot
t_{\min}$ from a strict segment drift bound and a minimum dwell time),
consumed by \texttt{rankingSegmentDynkin\_of\_mp\_strict}, which needs
presentations only at requested-but-ungranted states --- so
\texttt{sound\_GSHS\_streett\_bundle\_mp\_strict} leaves no analytic
per-model hypothesis in the bundle; all four hash-bound below. The
honest scope note that replaces the retired residue sentence: the
strict route\textquotesingle s progress row is \emph{pathwise}
(per-path accumulation, not in-expectation), the per-state bound
$\varepsilon \le V(x)$ is derived by the theorem rather than assumed
as a row, and the strict theorem is a companion to the mixed-profile
\texttt{\_mp} form, not an in-place replacement --- the landed
$\varepsilon = 0$ signature survives for citation stability. Width
note, deliberate strengthening
stated: the displayed ranking row asks strict $-\varepsilon_0$ decay
everywhere off $[\![\varphi_i]\!] \cup J_i$, where the bound
statements need nonincrease off the response and strict decrease only
at requests --- the displayed single-row form is stronger than
consumed, safe direction, and the producer emits it as one family.

\ruleblock{Lean verifier receipts}
\sourcenote{module \texttt{BernsteinLogicAux.Excursions}; bound build \texttt{ad62a846f92e\ldots} (hashed from the built olean at bind time). A receipt records that named checking layers accepted a build: (module, compiled-artifact SHA-256, environment SHA-256, checker layers), minted when the checking gate ran; the store receipts attesting the bound build are listed with the layers each run credited. Nothing here asks the reader to trust the store: every layer is an independently runnable program over the quoted sources --- rebuild the module and re-run the checkers to reproduce the verdict. Receipts attesting superseded builds credit no layer and are omitted from this rendering (the count of omissions, if any, is noted below).}
\begin{receipttable}
\receiptrow{495e3739c82a865d}{2026-07-18T13:23:35Z}{build+decorrelated}{\badge{BadgeGreen}{current}}
\receiptrow{fd91ed9c49139bb3}{2026-07-19T18:23:53Z}{build+decorrelated}{\badge{BadgeGreen}{current}}
\receiptrow{40087fade64fc483}{2026-07-20T01:21:05Z}{build+decorrelated}{\badge{BadgeGreen}{current}}
\receiptrow{483b08e5a60a003d}{2026-07-20T20:46:29Z}{build+decorrelated}{\badge{BadgeGreen}{current}}
\receiptrow{7142f9529d463930}{2026-07-20T20:58:02Z}{build+replay+hardened+decorrelated}{\badge{BadgeGreen}{current}}
\receiptrow{a477dcd09a788edc}{2026-07-21T00:02:42Z}{build+decorrelated}{\badge{BadgeGreen}{current}}
\receiptrow{30b972191d8805c2}{2026-07-21T05:38:35Z}{build+replay+hardened+decorrelated}{\badge{BadgeGreen}{current}}
\receiptrow{c6df511b240a81db}{2026-07-22T18:26:52Z}{build+decorrelated}{\badge{BadgeGreen}{current}}
\receiptrow{7d4ae5654cfa6cc8}{2026-07-31T21:59:44Z}{build+decorrelated}{\badge{BadgeGreen}{current}}
\receiptrow{50671021df35b6cc}{2026-07-31T22:28:38Z}{build+decorrelated}{\badge{BadgeGreen}{current}}
\receiptrow{153a2e7124bfd8d7}{2026-08-02T23:30:45Z}{build+decorrelated}{\badge{BadgeGreen}{current}}
\receiptrow{069303d7f7ac0945}{2026-08-03T16:26:49Z}{build+decorrelated}{\badge{BadgeGreen}{current}}
\receiptrow{b0dd60fee0daa2eb}{2026-08-03T22:45:33Z}{build+decorrelated}{\badge{BadgeGreen}{current}}
\end{receipttable}
\sourcenote{module \texttt{BernsteinLogicAux.QStreett}; bound build \texttt{469881beb770\ldots} (bound from receipt 9d1d848e0ba1718a, 2026-07-29T01:03:31Z). A receipt records that named checking layers accepted a build: (module, compiled-artifact SHA-256, environment SHA-256, checker layers), minted when the checking gate ran; the store receipts attesting the bound build are listed with the layers each run credited. Nothing here asks the reader to trust the store: every layer is an independently runnable program over the quoted sources --- rebuild the module and re-run the checkers to reproduce the verdict. Receipts attesting superseded builds credit no layer and are omitted from this rendering (the count of omissions, if any, is noted below).}
\begin{receipttable}
\receiptrow{e3990fe544d4708d}{2026-07-28T23:48:35Z}{build+decorrelated}{\badge{BadgeGreen}{current}}
\receiptrow{68e9374fa1b418fe}{2026-07-29T00:00:34Z}{build+decorrelated}{\badge{BadgeGreen}{current}}
\receiptrow{e5895c9706854cd3}{2026-07-29T00:49:04Z}{build+replay+hardened+decorrelated}{\badge{BadgeGreen}{current}}
\receiptrow{9d1d848e0ba1718a}{2026-07-29T01:03:31Z}{build+replay+hardened+decorrelated}{\badge{BadgeGreen}{current}}
\receiptrow{4a8d377b721b8523}{2026-07-31T22:00:09Z}{build+decorrelated}{\badge{BadgeGreen}{current}}
\receiptrow{d251a122540ab741}{2026-07-31T22:29:03Z}{build+decorrelated}{\badge{BadgeGreen}{current}}
\receiptrow{333635b8e047113a}{2026-08-02T23:31:14Z}{build+decorrelated}{\badge{BadgeGreen}{current}}
\receiptrow{9e96f80d63e4ad82}{2026-08-03T16:27:18Z}{build+decorrelated}{\badge{BadgeGreen}{current}}
\receiptrow{932b4b9e0b0fd323}{2026-08-03T22:46:02Z}{build+decorrelated}{\badge{BadgeGreen}{current}}
\end{receipttable}
\sourcenote{9 further receipt(s) in the store attest superseded builds of this module and are omitted from this rendering.}
\sourcenote{module \texttt{BernsteinLogicAux.LayerD.Streett}; bound build \texttt{7393065bbc7f\ldots} (bound from receipt abac471866cfcfd1, 2026-07-21T05:38:48Z). A receipt records that named checking layers accepted a build: (module, compiled-artifact SHA-256, environment SHA-256, checker layers), minted when the checking gate ran; the store receipts attesting the bound build are listed with the layers each run credited. Nothing here asks the reader to trust the store: every layer is an independently runnable program over the quoted sources --- rebuild the module and re-run the checkers to reproduce the verdict. Receipts attesting superseded builds credit no layer and are omitted from this rendering (the count of omissions, if any, is noted below).}
\begin{receipttable}
\receiptrow{0f3b50b4d3f0d33c}{2026-07-20T22:52:50Z}{build+decorrelated}{\badge{BadgeGreen}{current}}
\receiptrow{109669f848453ec2}{2026-07-21T00:03:19Z}{build+decorrelated}{\badge{BadgeGreen}{current}}
\receiptrow{7e17aa32b276623b}{2026-07-21T04:51:49Z}{build+decorrelated}{\badge{BadgeGreen}{current}}
\receiptrow{abac471866cfcfd1}{2026-07-21T05:38:48Z}{build+replay+hardened+decorrelated}{\badge{BadgeGreen}{current}}
\end{receipttable}
\sourcenote{6 further receipt(s) in the store attest superseded builds of this module and are omitted from this rendering.}
\sourcenote{module \texttt{QtslAux.LayerD.MPDischarge}; bound build \texttt{baded74dd943\ldots} (bound from receipt 7228834dea10e5f9, 2026-07-24T11:04:04Z). A receipt records that named checking layers accepted a build: (module, compiled-artifact SHA-256, environment SHA-256, checker layers), minted when the checking gate ran; the store receipts attesting the bound build are listed with the layers each run credited. Nothing here asks the reader to trust the store: every layer is an independently runnable program over the quoted sources --- rebuild the module and re-run the checkers to reproduce the verdict. Receipts attesting superseded builds credit no layer and are omitted from this rendering (the count of omissions, if any, is noted below).}
\begin{receipttable}
\receiptrow{7228834dea10e5f9}{2026-07-24T11:04:04Z}{build+decorrelated}{\badge{BadgeGreen}{current}}
\end{receipttable}
\layerline{\layeritem{build}{lake build (elaboration)}{\badge{BadgeGreen}{current}}\layeritem{replay}{leanchecker (official kernel replay)}{\badge{BadgeGreen}{current}}\layeritem{hardened}{safe\_verify (hardened screen)}{\badge{BadgeGreen}{current}}\layeritem{decorrelated}{lean4export + nanoda (decorrelated kernel)}{\badge{BadgeGreen}{current}}}
\sourcenote{\badge{BadgeGray}{comparator: none} no fixture in the three-implementation comparator harness (compiled / kernel-interpreted / reference-interpreted) exercises this rule.}
\ruleblock{HOL4 statement}
\absentbox{NO HOL STATEMENT YET}{Absent, with the gap located precisely: the HOL4 development ports the
\emph{base-calculus} checker cone (the \texttt{bern\_catalogue} theory,
built side-by-side with the Lean catalogue module) --- its object is the
wire-format certificate checker of the sequent stratum. The extension
theory\textquotesingle s satisfaction stratum has no HOL artefact at
all, so there is no HOL-side wire object a checker could
consume today. A HOL statement becomes meaningful
once the extension stratum is ported; it would then join the proof-object
layer the HOL port\textquotesingle s own coordination notes name as
successor territory. (Checked against the HOL sources: no Streett-,
satisfaction-, or NNF-shaped definitions exist there.)
}
\ruleblock{HOL verifier receipts}
\absentbox{NONE}{no HOL verification receipt exists for this rule.}
\ruleblock{Proof}

\sourcenote{Genre: core rule (tier 1) --- an ordinary mathematical proof at
Putnam-exam-solution detail, at the rule's \emph{proven scope}: discrete-kernel
profile, conservative ranking scope (\autoref{rem:qstb-scope} states the
scope loudly, including what is and is not hash-bound in this edition). Every
hypothesis is used and remarked on; the anchor table maps steps to the Lean
declarations that carry them.}

\textbf{Setting and claim.}\enspace The semantic object at the discrete-kernel
profile is a time-homogeneous Markov chain on a measurable state space $S$,
presented by a single step kernel $K$ (modes, guards, and resets of the
underlying hybrid presentation flattened into $K$), with path measure $\mu_x$
on $S^{\mathbb{N}}$ from a start state $x$ (the Ionescu--Tulcea trajectory
construction). Region data: the operating region $D$, its checker-recomputed
hull $\bar{D} \supseteq D$, and the exit set $U$ for the confinement
certificate; per Streett pair $i$ (finitely or countably many), a certified
request region $R_i$ and response region $G_i$. Certificate data: the
confinement certificate $V_0 : S \to \R$ with threshold
$y$, and per pair a ranking certificate $W_i : S \to \R$ with decay
rate $\varepsilon_i \in \Q_{>0}$ and bound $C_i$. The premise rows, each an
inherited base-calculus sequent:

\begin{itemize}
\item[(H1)] threshold arithmetic: $p + y \le 1$ (exact rational).
\item[(H2)] hull and mass rows, for every $z \in D \smallsetminus U$:
  $K(z)(\bar{D}) = 1$ and $K(z)(D \cup U) = 1$ --- one-step mass stays in the
  hull, and stays in operating-or-exit.
\item[(H3)] confinement rows: one-step generator row
  $\int V_0 \, dK(z) \le V_0(z)$ on $D \smallsetminus U$ (with one-step
  integrability); floor $V_0 \ge 1$ on $U \cap \bar{D}$; nonnegativity
  $V_0 \ge 0$ on $\bar{D}$.
\item[(H4)] ranking rows, per pair, on the confinement region
  $D \smallsetminus U$: bounds $0 \le W_i \le C_i$ on $\bar{D}$; nonincrease
  $\int W_i \, dK(z) \le W_i(z)$ off the response region $G_i$; request decay
  $\int W_i \, dK(z) \le W_i(z) - \varepsilon_i$ on $R_i \smallsetminus G_i$.
\item[(H5)] initialization: $x \in D$ and $V_0(x) \le y$ --- i.e.\
  $x$ lies in the concluded region atom
  $\mathsf{qraSublevel}(V_0, y)$.
\end{itemize}

The claim: under (H1)--(H5),
\[
\mu_x\Bigl\{\omega \;\Bigm|\; \forall i,\ \bigl(\exists^\infty n,\ \omega_n
\in R_i\bigr) \Rightarrow \bigl(\exists^\infty n,\ \omega_n \in G_i\bigr)
\Bigr\} \;\ge\; p ,
\]
which is the Streett-acceptance reading of the concluded judgment
$\mathcal{M},\ x_0 \in \{V_0 \le y\} \vDash
\mathsf{SProbGE}\bigl(p,\ \mathsf{PDwellStreett}(\mathsf{PairBase}(\rho_1,
\varphi_1), \ldots, \mathsf{PairBase}(\rho_k, \varphi_k))\bigr)$ once the
certified regions are tied to the formulas (\autoref{rem:qstb-tie}).

\begin{blemma}[confinement]\label{lem:qstb-confine}
$\mu_x\{\omega \mid \forall n,\ \omega_n \in D \smallsetminus U\} \ge
1 - V_0(x) \ge 1 - y \ge p$.
\end{blemma}
\begin{proof}
This is the quantitative-safety rule consumed whole: by (H2)--(H3), $V_0$
read along the path and stopped at the first $U$-visit is a nonnegative
supermartingale with barrier $\ge 1$ on the stopped set, so Ville's maximal
inequality bounds the exit probability by $V_0(x)$; the mass rows (H2) keep
every reachable state inside $\bar{D}$, where the nonnegativity and floor
rows actually hold --- the fail-closed pairing. The full argument, with the
two-branch counterexample showing the nonnegativity row is not decoration,
is the \rulename{QSafetyBattery} section's; the displayed intermediate
bound is the chain inside that argument --- Proposition~\ref{prop:qsb-pattern}
applied to the clamped certificate of Lemma~\ref{lem:qsb-clamp} bounds
the exit probability by $g(x) = V_0(x)$ --- consumed here through the
same end-to-end statement (anchor table), whose packaged form
Theorem~\ref{thm:qsb-failclosed} then states as the $\ge p$ complement.
The last two inequalities are (H5) and (H1).
\end{proof}

\begin{blemma}[windowed telescoping invariant]\label{lem:qstb-window}
Fix one pair $(R, G)$ with certificate $W$, rate $\varepsilon$, and suppose
the generator rows hold \emph{globally}: $\mathbb{E}[W_{n+1} \mid \mathcal{F}_n]
\le W_n$ a.e.\ off $G$, and $\le W_n - \varepsilon$ a.e.\ on requests off
$G$, for the process $W_n := W(X_n)$ adapted to a filtration
$(\mathcal{F}_n)$, with $W_n \ge 0$ and integrable. For a window start $N$
let $A_{N,n} := \{\text{no response in } [N, n)\}$. Then for all $n \ge N$,
\[
\int_{A_{N,n}} W_n \, d\mu \;+\; \varepsilon \sum_{k = N}^{n-1}
\mu\bigl(A_{N,k+1} \cap \{X_k \in R\}\bigr) \;\le\; \int W_N \, d\mu .
\]
\end{blemma}
\begin{proof}
Induction on $n$ from $n = N$, where $A_{N,N}$ is everything and the sum is
empty. For the step: $A_{N,n+1}$ is $\mathcal{F}_n$-measurable (it
constrains coordinates $< n + 1$), so the tower property converts
$\int_{A_{N,n+1}} W_{n+1}$ into $\int_{A_{N,n+1}}
\mathbb{E}[W_{n+1} \mid \mathcal{F}_n]$. Split $A_{N,n+1}$ at time $n$ into
its request and non-request parts: on the request part the decay row pays
$\varepsilon \cdot \mu(A_{N,n+1} \cap \{X_n \in R\})$ --- exactly the new
summand --- and on the non-request part the nonincrease row holds (both
parts lie off $G$ at time $n$, since $A_{N,n+1}$ forbids a response at
$n$). Finally $A_{N,n+1} \subseteq A_{N,n}$ and $W_n \ge 0$ shrink the
integral domain monotonically, closing the induction. Two features deserve
remark. First, \emph{no sign condition on $\varepsilon$ is used here}: the
decrement paid at step $k$ is exactly the summand recorded at index $k$,
an exact cancellation --- positivity of $\varepsilon$ enters only in
Lemma~\ref{lem:qstb-bc}. Second, the two generator rows do genuinely different
work and neither subsumes the other: nonincrease carries the (possibly
long) non-request stretches; decay is charged only at requests.
\end{proof}

\begin{blemma}[summability and Borel--Cantelli]\label{lem:qstb-bc}
Under the hypotheses of Lemma~\ref{lem:qstb-window} with $\varepsilon > 0$
and $W$ bounded: almost surely, infinitely many requests imply infinitely
many responses.
\end{blemma}
\begin{proof}
Fix a window start $N$. Lemma~\ref{lem:qstb-window} bounds every
partial sum of the series below by $\mathbb{E}[W_N]$ uniformly (the
leading integral it retains is nonnegative), and a nonnegative series
is the supremum of its partial sums, so
\[
\varepsilon \sum_{j \ge 0} \mu\bigl(\{\text{no response in } [N, N+j]\}
\cap \{X_{N+j} \in R\}\bigr) \;\le\; \mathbb{E}[W_N] \;<\; \infty ,
\]
so the un-responded-request events of window $N$ are summable, and by the
Borel--Cantelli lemma their limsup is null. A path with infinitely many
requests but only finitely many responses has a finite (possibly empty)
set of response times; for any window start $N$ beyond all of them
($N = 0$ when there are none), the path lies in infinitely many
of window $N$'s un-responded-request events --- i.e.\ in that null limsup.
The union over the countably many window starts is null, which is the
claim. (Here $\varepsilon > 0$ is finally consumed, dividing the bound;
boundedness of $W$ supplies the finite right-hand side and, upstream,
integrability of every $W_n$.)
\end{proof}

\begin{blemma}[localization to the confinement region]\label{lem:qstb-loc}
The bundle's ranking rows (H4) hold only on $D \smallsetminus U$, while
Lemma~\ref{lem:qstb-window} wants global rows. Define the \emph{localized}
pair process: flag the state with the survival indicator of the event
``no visit to $U \cup D^{\mathrm{c}}$ so far'', zero the certificate at the
first exit, and count exited states as \emph{responded}. Then the flagged
process satisfies the global hypotheses of Lemmas
\ref{lem:qstb-window}--\ref{lem:qstb-bc}, and the conclusion, restricted to
paths confined forever, descends to the raw process: almost surely on
$\{\forall n,\ \omega_n \in D \smallsetminus U\}$, infinitely many
$R_i$-visits imply infinitely many $G_i$-visits.
\end{blemma}
\begin{proof}
Write it out. Let $\sigma := \inf\{n : \omega_n \in U \cup
D^{\mathrm{c}}\}$ (a stopping time --- the regions are measurable), and
localize by
\[
W'_n := W(\omega_n)\,\mathbf{1}[\sigma > n],
\qquad
G' \ \text{holds at time } n \ :\iff\ \omega_n \in G \ \text{or}\
\sigma \le n ,
\]
the zeroed certificate and the enlarged (exited-counts-as-responded)
response. Check the global hypotheses of
Lemma~\ref{lem:qstb-window} for $(W', G')$ by cases on the state class
at time $n$:

\emph{(surviving, non-request)} $\sigma > n$ and $\omega_n \notin R$:
off $G'$ means also $\omega_n \notin G$, so $\omega_n \in
(D \smallsetminus U) \smallsetminus G$; the mass row (H2) carries the
next step into $\bar{D}$, where $0 \le W \le C$ (H4), so
$W'_{n+1} = W(\omega_{n+1})\mathbf{1}[\sigma > n+1] \le W(\omega_{n+1})$
almost surely on this branch --- the masked certificate is dominated by
the unmasked one --- and
$\mathbb{E}[W'_{n+1} \mid \mathcal{F}_n] \le \int W \, dK(\omega_n)
\le W(\omega_n) = W'_n$ by the nonincrease row.

\emph{(surviving, request)} the same domination, closed by the decay
row: $\mathbb{E}[W'_{n+1} \mid \mathcal{F}_n] \le W'_n - \varepsilon$.

\emph{(exited)} $\sigma \le n$: the state is in the enlarged response
$G'$ by construction, so no generator row is demanded of it --- the global
hypotheses hold vacuously, and indeed $W'_n = 0$ there.

Bounds and adaptedness: $0 \le W'_n \le C$ everywhere (H4 on $\bar{D}$
while surviving --- reachable surviving states lie in $\bar D$ by (H2)
and the deterministic start --- zero after), and $W'$ is adapted since
$\sigma$ is a stopping time. (To consume
Lemma~\ref{lem:qstb-window} literally in its certificate reading, run
the chain on the flagged state space $S \times \{0,1\}$: the flag is a
running conjunction updated by membership of each new coordinate in
$U \cup D^{\mathrm{c}}$, so the flagged chain is again time-homogeneous
Markov, $W'$ and the enlarged response are state functions there ---
$W'(z, f) = f \cdot W(z)$, response $\iff$ $f = 0$ or $z \in G$ --- and
the flagged trajectory law projects to $\mu_x$.) Lemmas
\ref{lem:qstb-window}--\ref{lem:qstb-bc} then give: almost surely,
infinitely many requests-while-surviving imply infinitely many
$G'$-responses. Restricted to confined paths the flag never trips
($\sigma = \infty$), $W' = W$, the enlarged response set coincides with
$G_i$, and the conclusion is the raw one. The
exited-counts-as-responded device is not a weakening of the claim: the
theorem below asserts the pair implications only on the confinement event
of Lemma~\ref{lem:qstb-confine}, exactly where the device is invisible.
\end{proof}

\begin{bthm}[soundness of \rulename{QStreettBattery}, discrete profile,
conservative ranking scope]\label{thm:qstb-sound}
Under (H1)--(H5),
$\mu_x\{\omega \mid \forall i,\ (\exists^\infty n,\ \omega_n \in R_i)
\Rightarrow (\exists^\infty n,\ \omega_n \in G_i)\} \ge p$.
\end{bthm}
\begin{proof}
Per pair $i$, Lemmas \ref{lem:qstb-window}--\ref{lem:qstb-loc} give the pair
implication almost surely on the confinement event. The pairs are
countably many, so the almost-sure statements intersect (countable
intersection of co-null sets); the Streett event therefore contains the
confinement event up to a null set, and Lemma~\ref{lem:qstb-confine} bounds
its measure below by $p$. $\qedhere$
\end{proof}

\begin{bremark}[every hypothesis earns its seat]\label{rem:qstb-hyp}
\emph{(H1)}: at $p + y > 1$ the concluded atom admits start
states with $1 - V_0(x) < p$ and Lemma~\ref{lem:qstb-confine} yields less
than $p$ --- the sublevel atom and the probability bound are coupled by
exactly this row. \emph{(H2)}: drop the hull row and probability mass can
reach $U \smallsetminus \bar{D}$, where no floor row holds, voiding
Ville's barrier in Lemma~\ref{lem:qstb-confine}; drop the
operating-or-exit row and the stopped-process reading of the confinement
certificate is unsound (mass escapes the region where a generator row was ever
asserted) --- the two recomputed side-conditions are the fail-closed
pairing. \emph{(H3) nonnegativity}: load-bearing, not hygiene --- a signed
confinement certificate can go negative on one branch and thereby
subsidize exit mass on another; the recorded two-branch counterexample
lives in the \rulename{QSafetyBattery} section. \emph{(H4) strictness}:
at $\varepsilon_i = 0$, take $W_i \equiv 1$ on a one-state chain with
$K(s) = \delta_s$, $R_i = \{s\}$, $G_i = \varnothing$, $V_0 \equiv 0$:
every row holds, the path requests forever, and no response ever occurs
--- the conclusion fails at any $p > 0$. \emph{(H4) bound rows}: the
constant $C_i$ gives integrability of the localized certificate
(Lemma~\ref{lem:qstb-loc}) and the finite right-hand side of
Lemma~\ref{lem:qstb-bc}; nonnegativity of $W_i$ additionally powers the
monotone domain-shrink step inside Lemma~\ref{lem:qstb-window}.
\emph{(H5)}: initialization is what makes the region atom an
\emph{under}-approximation of the satisfaction set --- the workhorse
orientation of the satisfaction judgment.
\end{bremark}

\begin{bremark}[scope, loudly]\label{rem:qstb-scope}
Three ledger entries: the discrete-profile warrant, the cited
testimony, and the continuous-time profile. \emph{(i) Bound in
this edition}: the five excursion-cluster declarations (anchor table) ---
the windowed telescoping engine (Lemmas
\ref{lem:qstb-window}--\ref{lem:qstb-bc}), its additive-drift sibling for
plain reachability (the ranking hitting-time bound and almost-sure
hitting), and the nonnegative-supermartingale convergence theorem (part
of the cluster, available for the general excursion picture; the
presented proof path does not consume it) --- together with the
end-to-end assembly at the conservative scope
(\texttt{sound\_QStreettBattery} --- Lemmas \ref{lem:qstb-confine},
\ref{lem:qstb-loc} and Theorem~\ref{thm:qstb-sound} as one statement)
and the discharged carved-scope corollary
(\texttt{sound\_QStreettBattery\_carved\_discharged}), both in the
auxiliary module \texttt{BernsteinLogicAux.QStreett}. These warrant the
rule at the discrete-kernel profile. Also bound in this edition, under
the \rulename{QStreettCmplRequest} entry whose section carries its
event-level kernel: the formula tie
(\texttt{sound\_QStreettBattery\_tie}, with
\texttt{streett\_event\_mono}). \emph{(ii) Machine-checked and
cited, not bound}: the localization lemma under its
own name (\texttt{streett\_pair\_localized}), and the parameterized
route to the sharper \emph{carved} ranking scope, where decay is
demanded only off $G_i \cup J_i$ with
$J_i := \{V_i \le 1 - \gamma_i\} \cap
(I \smallsetminus [\![\rho_i]\!])$ --- $I$ is the display's confinement
region $\{V_0 < 1\}$; this section's exit set instances it as
$U = \{1 \le V_0\}$ --- and the absorbing certificates
$(V_i, \gamma_i)$ repay the carve-out: a L\'evy upward-convergence
argument shows frequent $J_i$-visits preclude frequent requests, the
absorption bound being discharged from kernel-level rows by a
time-shifted supermartingale bridge plus a conditional Ville inequality
(\texttt{sound\_QStreettBattery\_carved},
\texttt{absorptionBound\_of\_kernel}). Binding those in a later edition
is the named follow-up; until then this page claims them as
testimony, not as hash-bound warrant. \emph{(iii) The
martingale-problem profile, discharged whole}: the confinement half
and every zero-decrement ranking row ($\varepsilon_i = 0$:
nonincrease off the response region) are discharged by
\texttt{sound\_GSHS\_streett\_bundle\_mp}, and the strict-decrement
rows ($\varepsilon_i > 0$) by strict presentations ---
\texttt{MPFlowStrict} extends the martingale-problem presentation by
one pathwise progress row (along every path the accumulated generator
drift of the discounted certificate reaches $-\varepsilon$ by a
finite sample time; the intended certificate-side derivation is
$\varepsilon = \delta \cdot t_{\min}$ from a strict segment drift
bound and a minimum dwell time), and
\texttt{rankingSegmentDynkin\_of\_mp\_strict} derives from it the
ranking row family the $\varepsilon = 0$ route could not produce,
needing presentations only at requested-but-ungranted states. The
proof mechanism, for the record, is \emph{not} the lower-bound
optional-sampling inequality an earlier edition of this remark
predicted: the stopped certificate splits as martingale plus drift
accumulation, the existing upper-bound optional sampling serves the
martingale part, and dominated convergence on
$\varepsilon$-truncated drift accumulations supplies the strict term
(the lower-bound lemma was built anyway, completing the
optional-sampling module). The assembled
\texttt{sound\_GSHS\_streett\_bundle\_mp\_strict} therefore carries
no analytic per-model hypothesis --- all four statements live in
module \texttt{QtslAux.LayerD.MPDischarge} and are hash-bound in
this edition via the registered qtsl-aux package, so, like its
safety and deadline siblings, this rule carries no conditional
mark: each named profile is proven, and its theorems are bound.
\end{bremark}

\begin{bremark}[the formula tie]\label{rem:qstb-tie}
The theorem is stated over certified regions $R_i, G_i$; the concluded
judgment speaks of formulas $\rho_i, \varphi_i$. The sound directions are
asymmetric: certified \emph{request} regions may \emph{over}-approximate
the request denotation, certified \emph{response} regions must
\emph{under}-approximate the response denotation --- weakening the
hypothesis and strengthening the conclusion of each pair implication.
This is the content of the complement-request sub-rule
(\rulename{QStreettCmplRequest}, two sections on), whose event-level kernel
is two monotonicity steps.
\end{bremark}

\begin{bremark}[degenerate pairs, and why Streett is primitive]
A B\"uchi acceptance condition is the one-pair Streett condition with
trivially-true request; on that pair the bundle degenerates to the
classical recurrence certificate (the \rulename{QStreettDegenerate}
section). The converse embedding does not exist at parity of cost:
deterministic B\"uchi automata recognize only limit languages
\cite{landweber1969}, and re-encoding a $k$-pair Streett condition
B\"uchi-primitively costs an exponential-in-$k$ automaton product, where
the bundle pays $2k + 1$ certificates. Streett conditions enter the
literature as the complemented-pairs acceptance of \cite{streett1982},
the linear-time face of Rabin's pairs \cite{rabin1969}; the standard
survey is \cite{thomas1990}.
\end{bremark}

\begin{bremark}[relation to the literature]
The ranking row is a Foster--Lyapunov drift condition
\cite{foster1953,meyntweedie1993} in certificate form; its use to force
almost-sure hitting is the ranking-supermartingale argument of the
probabilistic-programs tradition \cite{chakarov2013}. The confinement
certificate is Ville's inequality \cite{ville1939}; the convergence item
of the bound cluster is Doob's supermartingale convergence theorem
\cite{doob1953}; the summability-to-limsup step is Borel--Cantelli (see
e.g.\ \cite[Thm.~2.3.1]{durrett2019}). As with the polynomial rules of
the base calculus, the certificate form moves all analytic content into
the \emph{choice} of $(V_0, W_i, \varepsilon_0)$: what the checker
replays is row arithmetic, and what this proof adds is only the
assembly. No completeness is claimed or needed here; the source
catalogue records a completeness theorem for the bundle relative to its
termination seam, outside this book's kernel-facing scope.
\end{bremark}

\begin{anchors}
Lemma \ref{lem:qstb-confine} & \texttt{sound\_QSafetyBattery\_failClosed} (bound under the \rulename{QSafetyBattery} entry; consumed whole)\\
Lemma \ref{lem:qstb-window} + no-sign-condition remark & \texttt{streett\_window\_invariant} (the engine's inner induction; subsumed by the bound \texttt{streett\_window\_sum\_le})\\
Lemma \ref{lem:qstb-bc}, summability form & \texttt{streett\_window\_sum\_le} (bound)\\
Lemma \ref{lem:qstb-bc}, Borel--Cantelli assembly & \texttt{streett\_ae\_frequently\_response} (bound)\\
additive-drift sibling (reachability) & \texttt{ranking\_hitting\_sum\_le}, \texttt{ae\_hits\_of\_ranking} (bound)\\
supermartingale convergence (cluster item) & \texttt{Supermartingale.exists\_ae\_tendsto\_of\_nonneg} (bound; not consumed by the presented path)\\
Lemma \ref{lem:qstb-loc} & \texttt{streett\_pair\_localized} (\texttt{BernsteinLogicAux.QStreett}; machine-checked, binding pending)\\
Theorem \ref{thm:qstb-sound} & \texttt{sound\_QStreettBattery} (same module, \textbf{bound in this edition})\\
Remark \ref{rem:qstb-tie} & \texttt{streett\_event\_mono}, \texttt{sound\_QStreettBattery\_tie} (same module; \textbf{bound in this edition} under the \rulename{QStreettCmplRequest} entry)\\
carved scope, discharged corollary & \texttt{sound\_QStreettBattery\_carved\_discharged} (same module, \textbf{bound in this edition})\\
Remark \ref{rem:qstb-scope}(ii), carved scope route & \texttt{AbsorptionBound}, \texttt{ae\_not\_frequently\_of\_absorption}, \texttt{supermartingale\_stopped\_shift}, \texttt{measure\_reach\_after\_le}, \texttt{absorptionBound\_of\_kernel}, \texttt{sound\_QStreettBattery\_carved} (same module, binding pending)\\
Remark \ref{rem:qstb-scope}(iii), martingale-problem profile & \texttt{sound\_GSHS\_streett\_bundle\_mp}, \texttt{MPFlowStrict}, \texttt{rankingSegmentDynkin\_of\_mp\_strict}, \texttt{sound\_GSHS\_streett\_bundle\_mp\_strict} (\texttt{QtslAux.LayerD.MPDischarge}; \textbf{bound in this edition} via the qtsl-aux package registration)\\
\end{anchors}


\rulesection{45}{QAvgBattery}{core rule}
\label{blcatrule-QAvgBattery}
\ruleblock{The rule (Steele metanotation)}
\sourcenote{as realized in the mechanized development; source row: \texttt{\texttt{QTSL-rules.toml}} --- \emph{QAvgBattery}. The metanotation statement is typeset by hand for notational quality; this rule cites no Lean statements yet (the absence is stated loudly below). The rule \emph{roster} is machine-checked (every mechanized rule tag round-trips to exactly one catalogue row); the statement \emph{text} is editorial.}
\begin{mathpar}
\inferrule*{
  \overline{\mathrm{drift}_m : R_m \vdash g - r_m - \mathcal{A}_m B_m \ \ge\ 0} \\
  \overline{\mathrm{guard}_j} \quad \overline{\mathrm{floor}_m} \quad \overline{\mathrm{nonneg}_m} \\
  \mathrm{tie} : \vdash g \bowtie c
}{
  \mathcal{M},\ x_0 \in R \vDash \mathsf{SReturn}{\bowtie}\bigl(c,\ \mathsf{poAvg}(o)\bigr)
}
\end{mathpar}

\par\medskip
\rulename{QAvgBattery} concludes long-run-average claims from the
certificate pair of the Poisson equation: a rational \emph{gain} $g$
and a \emph{bias} certificate family $B$, with the drift rows
carrying the constant term inside each row\textquotesingle s sequent
polynomial --- so the discharge species is the same sign-sequent
genre as every other battery, with no new checker architecture. On
the upper face the drifts certify gain minus reward minus
generator-of-bias nonnegative on scope, and the tie bounds the gain
above the claimed constant; on the lower face the orientations flip,
and the floors bound the bias above so the correction term dies in
the limit. The walk recomputes the required orientation from the
concluded root --- a producer never declares a face.

\cataloguenote{Entered at the roster census: the node was baked --- including its
tie proof field, which the safety and acceptance batteries still lack
--- but had never been catalogued. No Lean statement binds it yet;
the display is the schema-faithful schematic form, and the soundness
obligation (the Poisson-pair argument: telescoping the bias along the
trajectory, dividing by the horizon, the floors killing the
correction term) is stated in the proof section as an obligation,
not discharged.
}
\ruleblock{Theory rows (geolog)}
\sourcenote{The rule\textquotesingle s certificate content as it is declared in the geolog theory the checker loads --- payload schema, side-condition demands, and axiom rows, quoted from the theory source at every build of this catalogue (a stale quotation cannot exist). Axiom modes: \texttt{chase} = the engine derives; \texttt{report} = a violation is a loud judgment; \texttt{refuse} = a violation gates admission. Side-condition relations are discharged by bounded evaluation at the engine\textquotesingle s ground seam, never by derivation.}
\declhead{theory Qtsl}{QAvgBattery}{\badge{BadgeGreen}{quoted from the theory source at render}}
\begin{leansrc}
QAvgBattery : Set;
\end{leansrc}
\declhead{theory Qtsl}{QAvgBattery/*}{\badge{BadgeGreen}{quoted from the theory source at render}}
\begin{leansrc}
QAvgBattery/proof : QAvgBattery -> QProof;
QAvgBattery/gain : QAvgBattery -> Q;
QAvgBattery/arg : QAvgBattery -> StatePred;
QAvgBattery/cert : QAvgBattery -> QCertFam;
QAvgBattery/drifts : QAvgBattery -> QProofList;
QAvgBattery/guards : QAvgBattery -> QProofList;
QAvgBattery/lowbounds : QAvgBattery -> QProofList;
QAvgBattery/nonnegs : QAvgBattery -> QProofList;
QAvgBattery/tie : QAvgBattery -> Proof;
QAvgBattery/scope_warrant : QAvgBattery -> QInvScope;
\end{leansrc}
\ruleblock{Lean statement}
\badge{BadgeAmber}{not yet mechanized} this rule cites no Lean statements: nothing on this page is hash-bound for it, and its dashboard row shows the same absence across the statement, receipt, and correspondence columns.
\ruleblock{HOL4 statement}
\absentbox{NO HOL STATEMENT YET}{No HOL development of the extension theory exists.
}
\ruleblock{HOL verifier receipts}
\absentbox{NONE}{no HOL verification receipt exists for this rule.}
\ruleblock{Proof}

\sourcenote{Genre: core rule (tier 1), battery group. This section states the rule's soundness
obligation exactly, per the catalogue's completeness-over-polish
charter (a rule that can only be named-and-gapped is still a row);
the discharge is open debt on the mechanization column, and the
schema quoted in the entry is the baked ground truth.}

\textbf{Soundness obligation.}\enspace The Poisson-equation pair argument: with drift rows certifying
$g - r_m - \mathcal{A}_m B_m \ge 0$ on scope per mode (upper face),
guards at the seams, floors and nonnegativity on the bias family,
telescoping the compensated bias along any admitted trajectory bounds
the accumulated reward's time average by the gain plus a correction
term of order bias-over-horizon; the floors kill the correction in
the limit, and the tie transfers the gain bound to the claimed
constant. The lower face is the same argument with orientations
flipped. The obligation is stated over the same behavioural models as
the safety battery's, with the walk recomputing the required
orientation from the concluded root.

\begin{bremark}[status]
No Lean statement binds this node. The discharge species is the
landed sign-sequent genre --- the expected mechanization reuses the
battery kit's supermartingale machinery with the compensated bias in
the certificate slot --- and the node's tie field is already baked,
so no schema motion is expected at mechanization time.
\end{bremark}


\rulesection{46}{QRateFromClock}{core rule}
\label{blcatrule-QRateFromClock}
\ruleblock{The rule (Steele metanotation)}
\sourcenote{as realized in the mechanized development; source row: \texttt{\texttt{QTSL-rules.toml}} --- \emph{QRateFromClock}. The metanotation statement is typeset by hand for notational quality; this rule cites no Lean statements yet (the absence is stated loudly below). The rule \emph{roster} is machine-checked (every mechanized rule tag round-trips to exactly one catalogue row); the statement \emph{text} is editorial.}
\begin{mathpar}
\inferrule*{
  \mathsf{TrMeter}(t, d) \\
  \rho \cdot d \ge 1 \\
  \sigma \ge 1 \\
  \mathsf{qtsl\_rate\_clock}(n)
}{
  \mathcal{M} \vDash \mathsf{SRateLE}(\{t\},\ \rho,\ \sigma)
}
\end{mathpar}

\par\medskip
\rulename{QRateFromClock} discharges the rate-envelope specification
former from clock data alone: a transition whose clock carries a dwell
floor $d$ (demonic, forced, or carved --- never the instantaneous
race clock, which has no floor) satisfies
$\mathsf{SRateLE}(\{t\}, \rho, \sigma)$ for any $\rho \cdot d \ge 1$
and $\sigma \ge 1$. The rate lane\textquotesingle s other feed is the
supermartingale kit --- a different evidence family, deliberately not
conflated; this rule is the clock seam.

The premise is a premise \emph{class}, not a node:
$\mathsf{TrMeter}(t, d)$ names three things together --- the
transition\textquotesingle s clock-content conformance (exactly-one
content presentation, wing positivity; already checker-discharged
rows), a POSITIVE lower witness $0 < d \le d'$ below the declared
floor $d'$ (positivity is the class\textquotesingle s own formation:
a nonpositive witness would inhabit the class while its division law
is vacuous), and the meter law as a \emph{named obligation} --- the
conformance rows alone do not prove elapsed-time separation. The law:
under the declared meter semantics (dwell accrues per transition;
guard exit and own firing reset it), successive firings of $t$ are
separated by at least $d$ of elapsed time, hence
$N_t(T) \le T/d + 1$ on any window under every resolution. The
checker\textquotesingle s trusted contribution is two exact rational
inequalities ($\rho \cdot d \ge 1$; $\sigma \ge 1$) over conformance
rows it already checks --- with $d$ resolved deterministically as the
transition\textquotesingle s declared finite floor (no producer
freedom; the class is monotone in $d$, so the declared floor is the
tightest lawful witness).

The concluded envelope is an upper-prevision claim; the obligation
actually proved is pathwise and surely --- strictly stronger --- and
the entry records that strength so a future pathwise rate former can
strengthen the conclusion without touching the node. Singleton
rosters compose by the rate lane\textquotesingle s documented
disjoint-roster sum law (envelopes add on disjoint rosters); this
rule is that law\textquotesingle s first in-catalogue consumer.

\cataloguenote{The schema is baked in the extension theory: the node, its walk, the
race-refusal, and the per-wing meter-demand rows are quoted below
from the theory source; the per-transition conformance rows the
premise class reads were checker-discharged before the node landed.
What remains owed is the meter law\textquotesingle s first
formalization (the shared keystone of this rule family\textquotesingle s
continuous-time obligations --- the discrete-profile induction is
elementary). No Lean statement binds the node yet; the mechanization
column is honest open debt.
}
\ruleblock{Theory rows (geolog)}
\sourcenote{The rule\textquotesingle s certificate content as it is declared in the geolog theory the checker loads --- payload schema, side-condition demands, and axiom rows, quoted from the theory source at every build of this catalogue (a stale quotation cannot exist). Axiom modes: \texttt{chase} = the engine derives; \texttt{report} = a violation is a loud judgment; \texttt{refuse} = a violation gates admission. Side-condition relations are discharged by bounded evaluation at the engine\textquotesingle s ground seam, never by derivation.}
\par\medskip
The baked rows: the node with its two exact-rational envelope
reports, the race-refusal (the instantaneous wing carries no floor
and is refused as subject), the walk, and the per-wing meter-demand
rows that mint the premise class\textquotesingle s
conformance-plus-floor reading at the declared dwell floor.

\declhead{theory Qtsl}{QRateFromClock}{\badge{BadgeGreen}{quoted from the theory source at render}}
\begin{leansrc}
QRateFromClock : Set;
\end{leansrc}
\declhead{theory Qtsl}{QRateFromClock/*}{\badge{BadgeGreen}{quoted from the theory source at render}}
\begin{leansrc}
QRateFromClock/proof : QRateFromClock -> QProof;
QRateFromClock/tr    : QRateFromClock -> Transition;
QRateFromClock/rho   : QRateFromClock -> Q;
QRateFromClock/sigma : QRateFromClock -> Q;
\end{leansrc}
\declhead{theory Qtsl}{ax/wf/qrfc1\_rho\_pos}{\badge{BadgeGreen}{quoted from the theory source at render}}
\begin{leansrc}
report ax/wf/qrfc1_rho_pos : forall n : QRateFromClock, x : Q.
  n QRateFromClock/rho = x
  |- {l: 0, r: x} q_lt;
\end{leansrc}
\declhead{theory Qtsl}{ax/wf/qrfc1\_sigma\_ge\_one}{\badge{BadgeGreen}{quoted from the theory source at render}}
\begin{leansrc}
report ax/wf/qrfc1_sigma_ge_one : forall n : QRateFromClock, x : Q.
  n QRateFromClock/sigma = x
  |- {l: 1, r: x} q_le;
\end{leansrc}
\declhead{theory Qtsl}{ax/qtsl/rate\_clock\_race}{\badge{BadgeGreen}{quoted from the theory source at render}}
\begin{leansrc}
refuse ax/qtsl/rate_clock_race : forall n : QRateFromClock,
    t : Transition, e : EnnRat, k : TrClock.
  n QRateFromClock/tr = t,
  t Transition/clock = k,
  k = {majorant: e} tkRace
  |- false;
\end{leansrc}
\declhead{theory Qtsl}{qtsl\_rate\_clock}{\badge{BadgeGreen}{quoted from the theory source at render}}
\begin{leansrc}
qtsl_rate_clock : {node: QRateFromClock} -> Prop;
\end{leansrc}
\declhead{theory Qtsl}{ax/tf/qrfc1\_walk}{\badge{BadgeGreen}{quoted from the theory source at render}}
\begin{leansrc}
report ax/tf/qrfc1_walk : forall n : QRateFromClock.
  |- {node: n} qtsl_rate_clock;
\end{leansrc}
\declhead{theory Qtsl}{ax/tf/qrfc1\_meter\_forced}{\badge{BadgeGreen}{quoted from the theory source at render}}
\begin{leansrc}
report ax/tf/qrfc1_meter_forced : forall n : QRateFromClock,
    t : Transition, d : Q, k : TrClock.
  n QRateFromClock/tr = t,
  t Transition/clock = k,
  k = {dwell: d} tkForced
  |- {tr: t, d: d} qtsl_tr_meter;
\end{leansrc}
\declhead{theory Qtsl}{ax/tf/qrfc1\_meter\_demonic}{\badge{BadgeGreen}{quoted from the theory source at render}}
\begin{leansrc}
report ax/tf/qrfc1_meter_demonic : forall n : QRateFromClock,
    t : Transition, d : Q, e : EnnRat, k : TrClock.
  n QRateFromClock/tr = t,
  t Transition/clock = k,
  k = {dwellMin: e} tkDemonic,
  e = {q: d} enFin
  |- {tr: t, d: d} qtsl_tr_meter;
\end{leansrc}
\declhead{theory Qtsl}{ax/tf/qrfc1\_meter\_carved}{\badge{BadgeGreen}{quoted from the theory source at render}}
\begin{leansrc}
report ax/tf/qrfc1_meter_carved : forall n : QRateFromClock,
    t : Transition, fl : Q, ce : EnnRat, k : TrClock.
  n QRateFromClock/tr = t,
  t Transition/clock = k,
  k = {floor: fl, ceil: ce} tkCarved
  |- {tr: t, d: fl} qtsl_tr_meter;
\end{leansrc}
\declhead{theory Qtsl}{qtsl\_tr\_meter}{\badge{BadgeGreen}{quoted from the theory source at render}}
\begin{leansrc}
qtsl_tr_meter : {tr: Transition, d: Q} -> Prop;
\end{leansrc}
\declhead{theory Qtsl}{qtsl\_tr\_content}{\badge{BadgeGreen}{quoted from the theory source at render}}
\begin{leansrc}
qtsl_tr_content : {tr: Transition} -> Prop;             // side-condition
\end{leansrc}
\ruleblock{Lean statement}
\badge{BadgeAmber}{not yet mechanized} this rule cites no Lean statements: nothing on this page is hash-bound for it, and its dashboard row shows the same absence across the statement, receipt, and correspondence columns.
\ruleblock{HOL4 statement}
\absentbox{NO HOL STATEMENT YET}{No HOL development of the extension theory exists; this
node\textquotesingle s content is walk structure over baked rows.
}
\ruleblock{HOL verifier receipts}
\absentbox{NONE}{no HOL verification receipt exists for this rule.}
\ruleblock{Proof}

\sourcenote{Genre: core rule (tier 1), the timers family (the
rate-envelope former's clock seam). This section states the rule's
soundness obligation exactly; the rule's schema and walk are baked in
the extension theory (rows quoted in this section), and the meter
law's formalization is the named obligation of the mechanization
column.}

\textbf{Soundness obligation (meter law + pathwise domination).}\enspace
Under the declared meter semantics (dwell accrues per transition;
guard exit and own firing reset it), a transition $t$ whose clock
carries floor $d > 0$ fires with successive firings separated by at
least $d$ of elapsed time, hence $N_t(T) \le T/d + 1$ on any window
of length $T$, under every resolution. The witness $d$ is the
transition's declared finite floor, resolved deterministically by
the walk (the premise class demands a POSITIVE witness at most the
declared floor, and is monotone, so the declared floor is the
tightest lawful instance; no producer freedom enters). With the
checker-verified exact inequalities $\rho \cdot d \ge 1$ and
$\sigma \ge 1$,
$N_t(T) \le T/d + 1 \le \rho T + \sigma$ pathwise and surely; the
upper prevision of the roster firing count therefore satisfies the
envelope --- pathwise domination, strictly stronger than the
concluded prevision-level claim. The instantaneous race clock is
refused as subject: it carries no floor, and its rate story is a
different evidence family (the thinning majorant), deliberately not
conflated with this seam.

\begin{bremark}[status]
The schema is baked: the node, its walk, the race-refusal, and the
per-wing meter-demand rows are landed theory rows (quoted in this
section), and the premise class's per-transition conformance rows
were checker-discharged before the node landed. No Lean statement
binds the node yet; the meter LAW is the named obligation ---
elementary at the discrete profile, and the shared keystone of this
rule family's continuous-time lifts.
\end{bremark}


\rulesection{47}{QInvRelativize}{core rule}
\label{blcatrule-QInvRelativize}
\ruleblock{The rule (Steele metanotation)}
\sourcenote{as realized in the mechanized development; source row: \texttt{\texttt{QTSL-rules.toml}} --- \emph{QInvRelativize}. The metanotation statement is typeset by hand for notational quality; its mechanized ground truth is the Lean statement below, extracted from the sources at every build of this book and hash-bound. The rule \emph{roster} is machine-checked (every mechanized rule tag round-trips to exactly one catalogue row); the statement \emph{text} is editorial, checked by eye against the extracted Lean beside it.}
\begin{mathpar}
\inferrule*{
  \mathsf{BoxInv}(B, \nu, x)
}{
  \mathsf{AService}(B, T, \chi \cap \nu, x) \;\leftrightarrow\; \mathsf{AService}(B, T, \chi, x)
}
\end{mathpar}
$B$ an \emph{arbitrary} behaviour --- no closure field consumed; $\mathsf{AService}$ is the d\'ebut-as-infimum reading unpacked without infima (a visit at some $t < T + \varepsilon$ for every $\varepsilon > 0$), $\mathsf{BoxInv}$ the $\Box$ reading (every point of every solution from $x$).

\ruleflavor{certificate content out of scope of the current baked theory --- realized in Lean over the behavioural floor.}
\par\medskip
One of the three timed-core basis schemata proved non-admissible in
the gnand completeness development (machine independence witnesses;
see the proof section's anchors), entering as tier-1 core: its
soundness obligation is semantic, discharged at Putnam-exam-solution
detail in the proof section and mechanized in
BernsteinLogic.QtslTimedCore. Sound over \emph{all} behavioural models
--- no side condition; the proof consumes none of the closure laws.

\ruleblock{Lean statement}
\declhead{BernsteinLogic.QtslTimedCore}{sound\_QInvRelativize}{\badge{BadgeGreen}{statement bound}}
\begin{leansrc}
/-- **QInvRelativize** (obligation relativization through invariants), the
biconditional core: under `□ν`, the visit sets of `χ ∩ ν` and `χ`
coincide along every solution from `x`, so the two service claims
agree. Sound for ARBITRARY behaviors — none of B1–B3 is consumed.
Catalogue: tier-1 core; independence machine-checked — a 3-state
deadline-blind structure validates the remaining basis and every rule
while refuting this schema, so it is underivable. -/
theorem sound_QInvRelativize (B : Behavior X) (T : ℝ≥0) (ν χ : Set X) (x : X)
    (hbox : B.BoxInv ν x) :
    B.AService T (χ ∩ ν) x ↔ B.AService T χ x
\end{leansrc}
\ruleblock{Correspondence (Steele display $\leftrightarrow$ Lean)}
\absentbox{NO CORRESPONDENCE DICTIONARY YET}{no reviewer-certified dictionary maps the displayed rule onto the mechanized statements above; until one is written and certified, that mapping rests on this section\textquotesingle s editorial care alone.}
\ruleblock{Lean verifier receipts}
\sourcenote{module \texttt{BernsteinLogic.QtslTimedCore}; bound build \texttt{e6370ca0a2ac\ldots} (bound from receipt 6460a6dd4d227c9a, 2026-08-04T01:16:03Z). A receipt records that named checking layers accepted a build: (module, compiled-artifact SHA-256, environment SHA-256, checker layers), minted when the checking gate ran; the store receipts attesting the bound build are listed with the layers each run credited. Nothing here asks the reader to trust the store: every layer is an independently runnable program over the quoted sources --- rebuild the module and re-run the checkers to reproduce the verdict. Receipts attesting superseded builds credit no layer and are omitted from this rendering (the count of omissions, if any, is noted below).}
\begin{receipttable}
\receiptrow{6460a6dd4d227c9a}{2026-08-04T01:16:03Z}{build+decorrelated}{\badge{BadgeGreen}{current}}
\receiptrow{659fe3396a48b674}{2026-08-04T02:23:12Z}{build+decorrelated}{\badge{BadgeGreen}{current}}
\end{receipttable}
\sourcenote{4 further receipt(s) in the store attest superseded builds of this module and are omitted from this rendering.}
\layerline{\layeritem{build}{lake build (elaboration)}{\badge{BadgeGreen}{current}}\layeritem{replay}{leanchecker (official kernel replay)}{\badge{BadgeAmber}{stale only}}\layeritem{hardened}{safe\_verify (hardened screen)}{\badge{BadgeAmber}{stale only}}\layeritem{decorrelated}{lean4export + nanoda (decorrelated kernel)}{\badge{BadgeGreen}{current}}}
\sourcenote{\badge{BadgeGray}{comparator: none} no fixture in the three-implementation comparator harness (compiled / kernel-interpreted / reference-interpreted) exercises this rule.}
\ruleblock{HOL4 statement}
\sourcenote{theory \texttt{bern\_tr\_t5qdyn}, HOL4 development script \texttt{\texttt{bern\_tr\_t5qdynScript.sml}} \badge{BadgeGreen}{script bound}}
\holhead{Theorem sound\_QInvRelativize}
\begin{holsrc}
Theorem sound_QInvRelativize:
  ∀fwd bud nu chi (x:'x).
    qdBoxInv fwd nu x ⇒
    (qdAService fwd bud (chi INTER nu) x ⇔ qdAService fwd bud chi x)
\end{holsrc}
\ruleblock{HOL verifier receipts}
\badge{BadgeBlue}{documentary} \texttt{hol-strict} CLEAN: theory \texttt{bern\_tr\_t5qdyn}, 22
theorems, 0 cheat-tagged, 0 nonstandard axioms (whitelist:
\texttt{DISK\_THM}; the four \texttt{bool} axioms) --- verification
run of 2026-07-22 on an isolated build host, against a
content-addressed checkout of the exact sources cited here (git blob
\texttt{d64e1954\ldots}); an independent acceptance host ran the same
build green at the porting seating (2026-07-21). The timed core this
statement quantifies over is ported with an inventoried seam:
trajectories ride the full real line with $0 \le t$ guards at exactly
the positions where the nonnegative-time binders sit, and the
behaviour structure\textquotesingle s closure proof fields flatten to
an explicit wellformedness conjunction. This rule\textquotesingle s
twin is closure-free exactly as its Lean statement is --- stated over
a bare forward-solution set, strictly stronger than a
hypothesis-for-hypothesis transliteration.

\sourcenote{HOL receipts are documentary: quoted results of \texttt{hol-strict} runs
and census rows from the HOL4 development\textquotesingle s build
evidence, not minted receipt artefacts. A HOL receipt store (the Lean
store\textquotesingle s twin) is named future work.
}
\ruleblock{Proof}

\sourcenote{Genre: core rule (tier 1) --- an ordinary mathematical
proof at Putnam-exam-solution detail. This section also fixes the
semantic setting shared by the three timed-core sections
(\rulename{QInvRelativize}, \rulename{QFirstService},
\rulename{QFirstServiceSplit}); the later two abbreviate it as ``the
timed-core setting''.}

\textbf{Setting (shared by the timed-core sections).}\enspace A
\emph{behavioural model} is a state set $X$ with atom denotations
$\llbracket\cdot\rrbracket \subseteq X$, extended over $\neg, \wedge,
\vee$ as set operations (so $\llbracket\chi\wedge\nu\rrbracket =
\llbracket\chi\rrbracket \cap \llbracket\nu\rrbracket$, etc.), and,
for each $x \in X$, a set
$B_x$ of trajectories $\tau : \mathbb{R}_{\ge 0} \to X$, subject to
three closure conditions and nothing else: \textbf{(B1)}
$\tau \in B_x \Rightarrow \tau(0) = x$; \textbf{(B2)}
$\tau \in B_x \Rightarrow (s \mapsto \tau(t+s)) \in B_{\tau(t)}$;
\textbf{(B3)} $\tau \in B_x$ and $\tau' \in B_{\tau(t)}$ imply
$\tau \cdot_t \tau' \in B_x$, where $\tau \cdot_t \tau'$ follows
$\tau$ on $[0,t)$ and $\tau'$ (re-anchored) from $t$ on. Empty $B_x$
is permitted. The \emph{début} of $\tau$ into $E \subseteq X$ is
$\tau_E := \inf\{t : \tau(t) \in E\}$ ($\inf\emptyset = \infty$);
worst-case service and invariance are
$$x \vDash \mathsf{A}_T\varphi \iff \forall \tau \in B_x :
\tau_{\llbracket\varphi\rrbracket} \le T, \qquad
x \vDash \Box\varphi \iff \forall \tau \in B_x\ \forall t :
\tau(t) \in \llbracket\varphi\rrbracket,$$
where deadlines $T$ range over \emph{finite} nonnegative reals
(the mechanization types them so); in particular $\tau_E \le T$
forces the visit set nonempty.
For $S \subseteq \mathbb{R}_{\ge 0}$ and finite $T \ge 0$,
$\inf S \le T$ holds iff for
every $\varepsilon > 0$ some $s \in S$ has $s < T + \varepsilon$ ---
the infimum-free reading used throughout, and the form the mechanized
statements take. A début need \emph{not} be attained: nothing in B1--B3
forces a visit set to contain its infimum. The three sections split
exactly along that line: \rulename{QInvRelativize} is sound over all behavioural
models; the two first-service laws demand \emph{attained débuts} for
their obligation formulas (their sections state the attainment lemma
and the countermodel making the demand honest).

\textbf{Claim (\texttt{sound\_QInvRelativize}).}\enspace
$\Box\nu \wedge \mathsf{A}_T(\chi \wedge \nu) \;\leftrightarrow\;
\Box\nu \wedge \mathsf{A}_T\chi$ is valid in every behavioural model.

\begin{bthm}[soundness of \rulename{QInvRelativize}]\label{thm:QInvRelativize-sound}
At every state $x$ of every behavioural model:
$x \vDash \Box\nu \wedge \mathsf{A}_T(\chi\wedge\nu)$ iff
$x \vDash \Box\nu \wedge \mathsf{A}_T\chi$.
\end{bthm}
\begin{proof}
It suffices to show, for $x \vDash \Box\nu$, that
$x \vDash \mathsf{A}_T(\chi\wedge\nu) \iff x \vDash \mathsf{A}_T\chi$;
conjoining $\Box\nu$ on both sides then gives the displayed
biconditional (both sides carry it as a conjunct).

Let $\tau \in B_x$. By $\Box\nu$, every point of $\tau$ lies in
$\llbracket\nu\rrbracket$, so
$$\{t : \tau(t) \in \llbracket\chi \wedge \nu\rrbracket\}
= \{t : \tau(t) \in \llbracket\chi\rrbracket\} \cap
  \{t : \tau(t) \in \llbracket\nu\rrbracket\}
= \{t : \tau(t) \in \llbracket\chi\rrbracket\},$$
the second equality because the second intersecting set is all of
$\mathbb{R}_{\ge 0}$. Equal visit sets have equal infima:
$\tau_{\llbracket\chi\wedge\nu\rrbracket} =
\tau_{\llbracket\chi\rrbracket}$. So the two service conditions agree
trajectory by trajectory, hence at $x$.
\end{proof}

\begin{bremark}[why the invariant conjunct is everything]
Drop $\Box\nu$ and the right-to-left direction fails in a two-location
clocked model whose $\chi$-location lies outside $\nu$: the service
fact about $\chi$ says nothing about $\chi \wedge \nu$. (Left-to-right
needs no invariant --- every $\chi\wedge\nu$-visit is a $\chi$-visit,
so that direction is valid outright.) The proof uses
$\Box\nu$ once, exactly to collapse the two visit sets.
\end{bremark}

\begin{bremark}[no closure is consumed]
The proof never builds a trajectory: B1--B3 are not used, so the
schema is sound over arbitrary trajectory assignments. Its
independence witness must therefore deviate in the \emph{operator}
rather than the trajectory space, and does: the independence witness
is a machine-checked, deadline-blind three-state structure whose
service-operator interpretation deviates at an unreachable state; it
validates the seed basis --- the axiom set prior to the corrections
these sections record --- while refuting \rulename{QInvRelativize} directly.
Nothing else in the basis ties finite worst-case débuts to
reachability, which is precisely what this axiom asserts.
\end{bremark}

\begin{anchors}
Theorem \ref{thm:QInvRelativize-sound} & \texttt{sound\_QInvRelativize} (\texttt{BernsteinLogic/QtslTimedCore.lean}); semantics \texttt{Behavior.AService}, \texttt{Behavior.BoxInv}\\
independence witness & machine-checked three-state deviant structure (see the remark above)\\
\end{anchors}


\rulesection{48}{QFirstService}{core rule}
\label{blcatrule-QFirstService}
\ruleblock{The rule (Steele metanotation)}
\sourcenote{as realized in the mechanized development; source row: \texttt{\texttt{QTSL-rules.toml}} --- \emph{QFirstService}. The metanotation statement is typeset by hand for notational quality; its mechanized ground truth is the Lean statement below, extracted from the sources at every build of this book and hash-bound. The rule \emph{roster} is machine-checked (every mechanized rule tag round-trips to exactly one catalogue row); the statement \emph{text} is editorial, checked by eye against the extracted Lean beside it.}
\begin{mathpar}
\inferrule*{
  B.\mathsf{FusionClosed} \\ 0 < n \\
  \forall i.\ \mathsf{FirstVisits}(B, \alpha_i) \\
  \forall i.\ \mathsf{AService}(B, T_i, \alpha_i, x)
}{
  \mathsf{AService}\big(B,\ {\textstyle\min_i T_i},\ \mathsf{fsTarget}(\alpha, T),\ x\big)
}
\end{mathpar}
$\mathsf{fsTarget}(\alpha, T) = \bigcup_i \big(\alpha_i \cap \{y \mid \forall j \ne i.\ \mathsf{AService}(B, T_j, \alpha_j, y)\}\big)$ --- some obligation served \emph{right here} while every other still binds at its ORIGINAL budget; $\mathsf{FirstVisits}$ = nonempty visit sets have least elements (attained d\'ebuts).

\ruleflavor{certificate content out of scope of the current baked theory --- realized in Lean over fusion-closed behaviours with attained d\'ebuts.}
\par\medskip
One of the three timed-core basis schemata proved non-admissible in
the gnand completeness development (machine independence witnesses;
see the proof section's anchors), entering as tier-1 core: its
soundness obligation is semantic, discharged at Putnam-exam-solution
detail in the proof section and mechanized in
BernsteinLogic.QtslTimedCore. Sound over behavioural models with
\emph{attained d\'ebuts} for the obligation formulas --- a real side
condition: the proof section's closing remark carries the
countermodel (mechanized) witnessing failure without it.

\ruleblock{Lean statement}
\declhead{BernsteinLogic.QtslTimedCore}{sound\_QFirstService}{\badge{BadgeGreen}{statement bound}}
\begin{leansrc}
/-- **FS_n (first-service, EQ-FS)** — mechanized soundness. Over a
fusion-closed behavior with attained débuts for the obligations: if
every obligation binds at `x`, then within the tightest budget every
solution reaches a state where some obligation is served and every
other still binds at its ORIGINAL budget. Catalogue: tier-1 core
(schema family, n ≥ 2; the n = 1 instance is degenerate-derivable);
independence machine-checked — a blind 3-state structure validates the
FS-less basis over all definable sets while refuting a 2-obligation
instance of FS. Class boundary:
`qFirstService_attainment_necessary` below. -/
theorem sound_QFirstService (hB3 : B.FusionClosed) (hn : 0 < n)
    (α : Fin n → Set X) (T : Fin n → ℝ≥0) (x : X)
    (hatt : ∀ i, B.FirstVisits (α i))
    (hx : ∀ i, B.AService (T i) (α i) x) :
    B.AService (Finset.univ.inf' (Finset.univ_nonempty_iff.mpr ⟨⟨0, hn⟩⟩) T)
      (B.fsTarget α T) x
\end{leansrc}
\declhead{BernsteinLogic.QtslTimedCore}{qFirstService\_attainment\_necessary}{\badge{BadgeGreen}{statement bound}}
\begin{leansrc}
/-- **The theorem the master doc's countermodel E certifies**: there is
a fusion-closed behavior and a state satisfying the FS₂ premise at
budgets (2, 2) whose FS-target conclusion FAILS — so the attainment
hypothesis of `sound_QFirstService` is not removable, and FS_n is
invalid over raw C₀. Witness trajectory for the failure: `stay 0`;
along it every point of the FS-target would be a spine point at some
parameter in (1, 2) at which the OTHER atom still binds — but from a
rational point the escape `esc t t` never meets `Eb`, and from an
irrational point it never meets `Ea`. -/
theorem qFirstService_attainment_necessary :
    ∃ (S : Type) (B : Behavior S) (α : Fin 2 → Set S) (x : S),
      B.FusionClosed ∧
      (∀ i, B.AService 2 (α i) x) ∧
      ¬ B.AService 2 (B.fsTarget α (fun _ => 2)) x
\end{leansrc}
\ruleblock{Correspondence (Steele display $\leftrightarrow$ Lean)}
\absentbox{NO CORRESPONDENCE DICTIONARY YET}{no reviewer-certified dictionary maps the displayed rule onto the mechanized statements above; until one is written and certified, that mapping rests on this section\textquotesingle s editorial care alone.}
\ruleblock{Lean verifier receipts}
\sourcenote{module \texttt{BernsteinLogic.QtslTimedCore}; bound build \texttt{e6370ca0a2ac\ldots} (bound from receipt 6460a6dd4d227c9a, 2026-08-04T01:16:03Z). A receipt records that named checking layers accepted a build: (module, compiled-artifact SHA-256, environment SHA-256, checker layers), minted when the checking gate ran; the store receipts attesting the bound build are listed with the layers each run credited. Nothing here asks the reader to trust the store: every layer is an independently runnable program over the quoted sources --- rebuild the module and re-run the checkers to reproduce the verdict. Receipts attesting superseded builds credit no layer and are omitted from this rendering (the count of omissions, if any, is noted below).}
\begin{receipttable}
\receiptrow{6460a6dd4d227c9a}{2026-08-04T01:16:03Z}{build+decorrelated}{\badge{BadgeGreen}{current}}
\receiptrow{659fe3396a48b674}{2026-08-04T02:23:12Z}{build+decorrelated}{\badge{BadgeGreen}{current}}
\end{receipttable}
\sourcenote{4 further receipt(s) in the store attest superseded builds of this module and are omitted from this rendering.}
\layerline{\layeritem{build}{lake build (elaboration)}{\badge{BadgeGreen}{current}}\layeritem{replay}{leanchecker (official kernel replay)}{\badge{BadgeAmber}{stale only}}\layeritem{hardened}{safe\_verify (hardened screen)}{\badge{BadgeAmber}{stale only}}\layeritem{decorrelated}{lean4export + nanoda (decorrelated kernel)}{\badge{BadgeGreen}{current}}}
\sourcenote{\badge{BadgeGray}{comparator: none} no fixture in the three-implementation comparator harness (compiled / kernel-interpreted / reference-interpreted) exercises this rule.}
\ruleblock{HOL4 statement}
\sourcenote{theory \texttt{bern\_tr\_t5qdyn}, HOL4 development script \texttt{\texttt{bern\_tr\_t5qdynScript.sml}} \badge{BadgeGreen}{script bound}}
\holhead{Theorem sound\_QFirstService}
\begin{holsrc}
Theorem sound_QFirstService:
  ∀fwd al (Tb:num->real) n (x:'x).
    qdFusionClosed fwd ∧ 0 < n ∧
    (∀i. i < n ⇒ qdFirstVisits fwd (al i)) ∧
    (∀i. i < n ⇒ qdAService fwd (Tb i) (al i) x) ⇒
    qdAService fwd (rminN Tb n) (qdFsTarget fwd al Tb n) x
\end{holsrc}
\ruleblock{HOL verifier receipts}
\badge{BadgeBlue}{documentary} \texttt{hol-strict} CLEAN: theory \texttt{bern\_tr\_t5qdyn}, 22
theorems, 0 cheat-tagged, 0 nonstandard axioms (whitelist:
\texttt{DISK\_THM}; the four \texttt{bool} axioms) --- verification
run of 2026-07-22 on an isolated build host, against a
content-addressed checkout of the exact sources cited here (git blob
\texttt{d64e1954\ldots}); an independent acceptance host ran the same
build green at the porting seating (2026-07-21). The timed core this
statement quantifies over is ported with an inventoried seam:
trajectories ride the full real line with $0 \le t$ guards at exactly
the positions where the nonnegative-time binders sit, and the
behaviour structure\textquotesingle s closure proof fields flatten to
an explicit wellformedness conjunction; the fusion-closure hypothesis
is consumed exactly as displayed. The d\'ebut/attainment engine
lemmas behind the statement (first-visit and served-point existence)
are ported alongside as local instruments, mirroring their private
status on the Lean side.

\sourcenote{HOL receipts are documentary: quoted results of \texttt{hol-strict} runs
and census rows from the HOL4 development\textquotesingle s build
evidence, not minted receipt artefacts. A HOL receipt store (the Lean
store\textquotesingle s twin) is named future work.
}
\ruleblock{Proof}

\sourcenote{Genre: core rule (tier 1), in the timed-core setting fixed
at the \rulename{QInvRelativize} section. This section proves the attainment
lemma shared by the timed-core sections; the class boundary (attained
débuts are a genuine side condition) is certified by an explicit
countermodel, sketched in the final remark and mechanized alongside
the soundness lemma.}

\textbf{Claim (\texttt{sound\_QFirstService}).}\enspace For $n \ge 2$
obligations,
$$\bigwedge_{i \le n} \mathsf{A}_{T_i}\alpha_i \;\sqsubseteq\;
\mathsf{A}_{\min_i T_i}\Big(\underbrace{\bigvee_{i}\big(\alpha_i \wedge
\bigwedge_{i' \ne i}\mathsf{A}_{T_{i'}}\alpha_{i'}\big)}_{=:~\Phi,
\text{ the FS-target}}\Big)$$
is valid in every behavioural model with \emph{attained débuts} for the
$\alpha_i$: along every trajectory, a nonempty
$\alpha_i$-visit set has a least element.

\begin{blemma}[attainment]\label{lem:fs-attain}
Attained débuts hold: (i) for closed
$\llbracket\alpha_i\rrbracket$ in any model whose trajectories are
right-continuous into a topological state space; (ii) in every clocked
model, for every formula; (iii) in discrete time, for every formula.
\end{blemma}
\begin{proof}
(i) Let $E$ be closed, $\tau$ right-continuous, $V := \{t : \tau(t)
\in E\}$ nonempty, $t_0 := \inf V$. If $t_0 \notin V$, pick visits
$v_n \downarrow t_0$ with $v_n > t_0$; right-continuity at $t_0$ gives
$\tau(v_n) \to \tau(t_0)$, and since each $\tau(v_n) \in E$, every
neighbourhood of $\tau(t_0)$ meets $E$, so $\tau(t_0) \in \overline{E}
= E$ --- contradiction. (No Hausdorff hypothesis: only ``closed'' and
``converges'' are used.) (ii) A clocked trajectory is a
right-continuous step function through finitely many locations on any
bounded window, with the post-jump location taken \emph{at} the jump time,
and \emph{every} formula denotation is a union of locations --- a
subset of the finite discrete location set, hence closed; so (i)
applies with the discrete topology. Concretely: given any visit at
time $R$, the visit set restricted to $[0,R]$ is a finite union of
left-closed intervals, whose least point is the global least point.
(iii) Nonempty subsets
of $\mathbb{N}$ have least elements.
\end{proof}

\begin{bthm}[soundness of \rulename{QFirstService}]\label{thm:fs-sound}
In every behavioural model with attained débuts for
$\alpha_1,\dots,\alpha_n$, at every state $x$: if
$x \vDash \mathsf{A}_{T_i}\alpha_i$ for all $i$, then
$x \vDash \mathsf{A}_{\min_i T_i}\Phi$.
\end{bthm}
\begin{proof}
Let $\tau \in B_x$.

\emph{First services exist.} For each $i$,
$\tau_{\llbracket\alpha_i\rrbracket} \le T_i < \infty$ forces the
visit set nonempty; attainment gives its least element
$t_i = \tau_{\llbracket\alpha_i\rrbracket} \le T_i$.

\emph{The first service overall.} Put $t^* := \min_i t_i$, attained
at some $i_0$, and $y := \tau(t^*) \in
\llbracket\alpha_{i_0}\rrbracket$. Choosing $j^* \in \arg\min_i T_i$:
$t^* \le t_{j^*} \le T_{j^*} = \min_i T_i$.

\emph{Co-obligations persist at $y$, at original budgets.} Fix
$i' \ne i_0$, $\sigma \in B_y$, and $\varepsilon > 0$. The fusion
$\rho := \tau \cdot_{t^*} \sigma$ lies in $B_x$ by B3 (typed by B1:
$\sigma(0) = y = \tau(t^*)$), so
$\rho_{\llbracket\alpha_{i'}\rrbracket} \le T_{i'}$, giving a
$\rho$-visit of $\alpha_{i'}$ at some $s < T_{i'} + \varepsilon$. If
$s < t^*$ then $\rho(s) = \tau(s)$, so $t_{i'} \le s < t^* \le
t_{i'}$ --- impossible. Hence $s \ge t^*$ and $\rho(s) =
\sigma(s - t^*)$: a $\sigma$-visit at time $s - t^* \le s <
T_{i'} + \varepsilon$. As $\varepsilon$ was arbitrary,
$\sigma_{\llbracket\alpha_{i'}\rrbracket} \le T_{i'}$. (Indeed
$\le T_{i'} - t^*$; the schema states the weaker original-budget
bound, which is what survives grid transport --- the residual-split
law \rulename{QFirstServiceSplit} is where the true residual is
recovered.)

\emph{Assembly.} $\sigma$ was arbitrary, so $y \vDash
\mathsf{A}_{T_{i'}}\alpha_{i'}$ for every $i' \ne i_0$; with
$y \in \llbracket\alpha_{i_0}\rrbracket$, $y \in
\llbracket\Phi\rrbracket$, and $\tau$ visits it at
$t^* \le \min_i T_i$.
\end{proof}

\begin{bremark}[each hypothesis, remarked]
B1 types the fusion; B3 is the engine (facts about all of $B_y$ reach
back to the premise at $x$ only through fusion); B2 is not consumed.
Attainment is consumed exactly once, and essentially --- to produce
the service point $y = \tau(t^*)$: an infimum is not a visit. (The
placement of $s$ then runs on the ordinary lower-bound property of
$t_{i'}$.) It is not removable: see the next remark. At
$n = 1$ the instance degenerates and is derivable; the schema family
starts at $n = 2$.
\end{bremark}

\begin{bremark}[the class boundary is real]\label{rem:fs-cm}
Over bare B1--B3 models the schema is \emph{false}. Countermodel (mechanized
as \texttt{qFirstService\_attainment\_necessary}): a spine
$\{y_p : p \ge 0\}$ with $\llbracket a \rrbracket$ the rational
$p \in (1,2)$ and $\llbracket b \rrbracket$ the irrational ones; from
every $y_p$, the trajectories are the spine-ride and, for each
$s \ge p$ with $s > 1$, a ride-to-$s$-then-exit to an atom-free sink.
The family is fusion-closed; every trajectory from $y_0$ meets both
atoms before time $2$ (density), so the premise
$\mathsf{A}_2 a \wedge \mathsf{A}_2 b$ holds; but at every
$y_q \in \llbracket a \rrbracket$ (so $q \in (1,2)$, making the
immediate exit legal) that exit kills $\mathsf{A}_2 b$, at every
$y_\iota \in \llbracket b \rrbracket$ it kills $\mathsf{A}_2 a$, and
states outside the two atoms fail both disjuncts outright, so the
FS-target denotes
$\emptyset$ and the conclusion fails. Dense interleaving of disjoint
atoms along one trajectory is exactly what right-continuity into
closed denotations forbids (Lemma \ref{lem:fs-attain}(i)) --- the
countermodel lives strictly outside the attained class, as it must.
\end{bremark}

\begin{anchors}
Theorem \ref{thm:fs-sound} & \texttt{sound\_QFirstService}; class mixin \texttt{Behavior.FusionClosed}, \texttt{Behavior.FirstVisits}; target \texttt{Behavior.fsTarget}\\
Lemma \ref{lem:fs-attain} & prose (this section); clocked/discrete instances discharge \texttt{FirstVisits} per model\\
Remark \ref{rem:fs-cm} & \texttt{qFirstService\_attainment\_necessary} (the countermodel, mechanized)\\
independence witness & machine-checked structure validating the corrected basis with this schema removed while refuting the schema itself\\
\end{anchors}


\rulesection{49}{QFirstServiceSplit}{core rule}
\label{blcatrule-QFirstServiceSplit}
\ruleblock{The rule (Steele metanotation)}
\sourcenote{as realized in the mechanized development; source row: \texttt{\texttt{QTSL-rules.toml}} --- \emph{QFirstServiceSplit}. The metanotation statement is typeset by hand for notational quality; its mechanized ground truth is the Lean statement below, extracted from the sources at every build of this book and hash-bound. The rule \emph{roster} is machine-checked (every mechanized rule tag round-trips to exactly one catalogue row); the statement \emph{text} is editorial, checked by eye against the extracted Lean beside it.}
\begin{mathpar}
\inferrule*{
  B.\mathsf{FusionClosed} \\ 0 < n \\ \forall i.\ c < T_i \\
  \forall i.\ \mathsf{FirstVisits}(B, \alpha_i) \\
  \forall i.\ \mathsf{AService}(B, T_i, \alpha_i, x)
}{
  \mathsf{AService}\big(B,\ c,\ \mathsf{fsTarget}(\alpha, T) \cup \mathsf{residualCell}(\alpha, T, c),\ x\big)
}
\end{mathpar}
$\mathsf{residualCell}(\alpha, T, c) = \{y \mid \forall i.\ \mathsf{AService}(B, T_i - c, \alpha_i, y)\}$ --- the no-service-by-$c$ branch pins elapsed time by the cutoff, recovering TRUE residuals $T_i - c$ where FS$_n$ can only re-issue original budgets.

\ruleflavor{certificate content out of scope of the current baked theory --- realized in Lean over fusion-closed behaviours with attained d\'ebuts.}
\par\medskip
One of the three timed-core basis schemata proved non-admissible in
the gnand completeness development --- for this law the machine
independence witnesses cover integer and half-integer deadline menus,
with finer denominators open (see the proof section's closing remark
and anchors) --- entering as tier-1 core: its soundness obligation is
semantic, discharged at Putnam-exam-solution detail in the proof
section and mechanized in BernsteinLogic.QtslTimedCore. Sound over
behavioural models with \emph{attained d\'ebuts} for the obligation
formulas --- a real side condition; its witnessing countermodel is
prose-only (not yet mechanized; the QFirstService countermodel is the
mechanized one).

\ruleblock{Lean statement}
\declhead{BernsteinLogic.QtslTimedCore}{sound\_QFirstServiceSplit}{\badge{BadgeGreen}{statement bound}}
\begin{leansrc}
/-- **FS-R (residual-split first-service)** — mechanized soundness,
same class as FS_n. For any cutoff `c` strictly below every budget:
within `c`, every solution reaches EITHER a first-service state
(original co-budgets — the FS-target) OR the residual cell, where
every obligation holds at its true residual `T i - c`. Genuinely new
on top of FS_n; independence machine-checked — a minimal 6-state
witness settles every all-integer deadline menu, and two further
certified 6-state witnesses extend the underivability to all
half-integer menus. -/
theorem sound_QFirstServiceSplit (hB3 : B.FusionClosed) (hn : 0 < n)
    (α : Fin n → Set X) (T : Fin n → ℝ≥0) (c : ℝ≥0) (x : X)
    (hc : ∀ i, c < T i)
    (hatt : ∀ i, B.FirstVisits (α i))
    (hx : ∀ i, B.AService (T i) (α i) x) :
    B.AService c (B.fsTarget α T ∪ B.residualCell α T c) x
\end{leansrc}
\ruleblock{Correspondence (Steele display $\leftrightarrow$ Lean)}
\absentbox{NO CORRESPONDENCE DICTIONARY YET}{no reviewer-certified dictionary maps the displayed rule onto the mechanized statements above; until one is written and certified, that mapping rests on this section\textquotesingle s editorial care alone.}
\ruleblock{Lean verifier receipts}
\sourcenote{module \texttt{BernsteinLogic.QtslTimedCore}; bound build \texttt{e6370ca0a2ac\ldots} (bound from receipt 6460a6dd4d227c9a, 2026-08-04T01:16:03Z). A receipt records that named checking layers accepted a build: (module, compiled-artifact SHA-256, environment SHA-256, checker layers), minted when the checking gate ran; the store receipts attesting the bound build are listed with the layers each run credited. Nothing here asks the reader to trust the store: every layer is an independently runnable program over the quoted sources --- rebuild the module and re-run the checkers to reproduce the verdict. Receipts attesting superseded builds credit no layer and are omitted from this rendering (the count of omissions, if any, is noted below).}
\begin{receipttable}
\receiptrow{6460a6dd4d227c9a}{2026-08-04T01:16:03Z}{build+decorrelated}{\badge{BadgeGreen}{current}}
\receiptrow{659fe3396a48b674}{2026-08-04T02:23:12Z}{build+decorrelated}{\badge{BadgeGreen}{current}}
\end{receipttable}
\sourcenote{4 further receipt(s) in the store attest superseded builds of this module and are omitted from this rendering.}
\layerline{\layeritem{build}{lake build (elaboration)}{\badge{BadgeGreen}{current}}\layeritem{replay}{leanchecker (official kernel replay)}{\badge{BadgeAmber}{stale only}}\layeritem{hardened}{safe\_verify (hardened screen)}{\badge{BadgeAmber}{stale only}}\layeritem{decorrelated}{lean4export + nanoda (decorrelated kernel)}{\badge{BadgeGreen}{current}}}
\sourcenote{\badge{BadgeGray}{comparator: none} no fixture in the three-implementation comparator harness (compiled / kernel-interpreted / reference-interpreted) exercises this rule.}
\ruleblock{HOL4 statement}
\sourcenote{theory \texttt{bern\_tr\_t5qdyn}, HOL4 development script \texttt{\texttt{bern\_tr\_t5qdynScript.sml}} \badge{BadgeGreen}{script bound}}
\holhead{Theorem sound\_QFirstServiceSplit}
\begin{holsrc}
Theorem sound_QFirstServiceSplit:
  ∀fwd al (Tb:num->real) n c (x:'x).
    qdFusionClosed fwd ∧ 0 < n ∧ 0 ≤ c ∧
    (∀i. i < n ⇒ c < Tb i) ∧
    (∀i. i < n ⇒ qdFirstVisits fwd (al i)) ∧
    (∀i. i < n ⇒ qdAService fwd (Tb i) (al i) x) ⇒
    qdAService fwd c
      (qdFsTarget fwd al Tb n UNION qdResidualCell fwd al Tb n c) x
\end{holsrc}
\ruleblock{HOL verifier receipts}
\badge{BadgeBlue}{documentary} \texttt{hol-strict} CLEAN: theory \texttt{bern\_tr\_t5qdyn}, 22
theorems, 0 cheat-tagged, 0 nonstandard axioms (whitelist:
\texttt{DISK\_THM}; the four \texttt{bool} axioms) --- verification
run of 2026-07-22 on an isolated build host, against a
content-addressed checkout of the exact sources cited here (git blob
\texttt{d64e1954\ldots}); an independent acceptance host ran the same
build green at the porting seating (2026-07-21). The timed core this
statement quantifies over is ported with an inventoried seam:
trajectories ride the full real line with $0 \le t$ guards at exactly
the positions where the nonnegative-time binders sit, and the
behaviour structure\textquotesingle s closure proof fields flatten to
an explicit wellformedness conjunction. The residual
budget\textquotesingle s truncated subtraction is rendered as the
max-$0$ coercion --- agreeing with the nonnegative-reals reading on
the whole consulted domain, not only under the split
theorem\textquotesingle s guard.

\sourcenote{HOL receipts are documentary: quoted results of \texttt{hol-strict} runs
and census rows from the HOL4 development\textquotesingle s build
evidence, not minted receipt artefacts. A HOL receipt store (the Lean
store\textquotesingle s twin) is named future work.
}
\ruleblock{Proof}

\sourcenote{Genre: core rule (tier 1), in the timed-core setting fixed
at the \rulename{QInvRelativize} section, with the attainment lemma of the
\rulename{QFirstService} section. This is the fourth and final basis
correction from the gnand completeness analysis: the law that makes
the decision
procedure's deadline-tightening step derivable.}

\textbf{Claim (\texttt{sound\_QFirstServiceSplit}).}\enspace For $n \ge 1$
obligations and any cutoff $0 < c < \min_i T_i$,
$$\bigwedge_{i \le n} \mathsf{A}_{T_i}\alpha_i \;\sqsubseteq\;
\mathsf{A}_c\Big(\Phi \;\vee\; \underbrace{\bigwedge_i
\mathsf{A}_{T_i - c}\,\alpha_i}_{=:~P,\text{ the residual cell}}\Big)$$
is valid in every behavioural model with attained débuts for the
$\alpha_i$, where $\Phi$ is the FS-target of \rulename{QFirstService}
(read at $n = 1$ with the empty co-obligation conjunction as $\top$,
so $\Phi = \alpha_1$).
Reading: within the cutoff, either some first service has happened
(and the co-obligations persist at original budgets), or nothing has
been served and every obligation holds at its \emph{true} residual.

\begin{bthm}[soundness of \rulename{QFirstServiceSplit}]\label{thm:fsr-sound}
In every behavioural model with attained débuts for
$\alpha_1,\dots,\alpha_n$, at every state $x$: if
$x \vDash \mathsf{A}_{T_i}\alpha_i$ for all $i$ and
$0 < c < \min_i T_i$, then $x \vDash \mathsf{A}_c(\Phi \vee P)$.
\end{bthm}
\begin{proof}
Let $\tau \in B_x$. As in the \rulename{QFirstService} proof,
attainment gives each obligation a least visit time
$t_i \le T_i$ along $\tau$; put $t^* := \min_i t_i$, attained at
$i_0$. Split on the cutoff.

\emph{Case $t^* \le c$.} The \rulename{QFirstService} argument,
verbatim, puts $y := \tau(t^*)$ in $\llbracket\Phi\rrbracket$: it
serves $\alpha_{i_0}$, and for $i' \ne i_0$ every $\sigma \in B_y$
inherits $\sigma_{\llbracket\alpha_{i'}\rrbracket} \le T_{i'}$ through
the fusion $\tau \cdot_{t^*} \sigma \in B_x$, the prefix being
$\alpha_{i'}$-free below $t^*$ by minimality. So $\tau$ visits
$\llbracket\Phi \vee P\rrbracket$ at $t^* \le c$.

\emph{Case $t^* > c$.} Let $y := \tau(c)$; we show
$y \in \llbracket P \rrbracket$, so that $\tau$ visits the disjunction
at time $c$. Fix $i$, $\sigma \in B_y$, $\varepsilon > 0$. Fuse:
$\rho := \tau \cdot_c \sigma \in B_x$ by B3, so there is a
$\rho$-visit of $\alpha_i$ at some $s < T_i + \varepsilon$. The prefix
is service-free through $c$: for $s' < c$, $\rho(s') = \tau(s')$, and
a visit there would give $t_i \le s' < c < t^* \le t_i$ ---
impossible; and the point at $c$ itself is $\rho(c) = \sigma(0) =
\tau(c)$, likewise unvisited (a visit at $c$ would give
$t_i \le c < t^*$). So $s \ge c$ and $\rho(s) = \sigma(s - c)$: a
$\sigma$-visit at time $s - c < (T_i - c) + \varepsilon$ (subtract
$c$ across $s < T_i + \varepsilon$, using $c < T_i$ to keep the
residual budget a genuine nonnegative quantity). As $\varepsilon$ was
arbitrary, $\sigma_{\llbracket\alpha_i\rrbracket} \le T_i - c$.
\end{proof}

\begin{bremark}[each hypothesis, remarked]
$c < \min_i T_i$ keeps every residual budget strictly positive ---
the regime the law is stated for; at $c = \min_i T_i$ the split
degenerates (under attained débuts the served case always fires,
since $t^* \le \min_i T_i$) and the residual row would carry
$\mathsf{A}_0$. Attainment
is consumed only through $t^*$ --- and only the \emph{served} case consumes
it: the no-service case uses the visit $s$ directly, never its
infimum. That asymmetry is honest and sharp: over bare B1--B3 models
the residual branch alone survives, and the schema as a whole fails
--- the countermodel (prose; not mechanized) retunes the
\rulename{QFirstService} countermodel --- the dense cluster moved, the
escape-legality threshold lowered to match --- and adds ``lazy
wings'': branches
postponing every service to just inside the original budgets, legal
precisely because the premise pins only original budgets before the
cluster, defeating
target and residual simultaneously along the spine. Attained débuts
restore the dichotomy the case split runs on.
\end{bremark}

\begin{bremark}[what the law buys]
\rulename{QFirstService} can only re-issue original budgets at the
service point; \rulename{QFirstServiceSplit} recovers the true residual
$T_i - c$ on the no-service branch, where the elapsed time is pinned
by the cutoff rather than by an unknowable service instant. Iterated
at successive cutoffs this derives the decision procedure's
deadline-tightening
step exactly (its descent lemma), which is what closed the
derivation side of the gnand completeness theorem. Independence: a
six-state witness proves the law underivable from the
FS-corrected basis for integer deadline menus, extended to
half-integer menus by two certified witnesses; independence at finer
denominators is open, and citation is bounded accordingly.
\end{bremark}

\begin{anchors}
Theorem \ref{thm:fsr-sound} & \texttt{sound\_QFirstServiceSplit}; residual \texttt{Behavior.residualCell}\\
countermodel (class boundary) & prose only (not yet mechanized); sketched in the first remark\\
independence witnesses & six-state minimal witness + two half-integer-menu witnesses, machine-checked (see the closing remark for the citation bound)\\
\end{anchors}


\blcatsection{Streett Rules and the Specification Surface}
\rulesection{50}{QStreettDegenerate}{alias for \textcolor{DiffBlue}{\rulename{QStreettBattery}}}
\label{blcatrule-QStreettDegenerate}
\ruleblock{The rule (Steele metanotation)}
\sourcenote{as realized in the mechanized development; source row: \texttt{\texttt{QTSL-rules.toml}} --- \emph{QStreettDegenerate}. The metanotation statement is typeset by hand for notational quality; its mechanized ground truth is the Lean statement below, extracted from the sources at every build of this book and hash-bound. The rule \emph{roster} is machine-checked (every mechanized rule tag round-trips to exactly one catalogue row); the statement \emph{text} is editorial, checked by eye against the extracted Lean beside it.}
\begin{mathpar}
\inferrule*{
  \text{one \rulename{QStreettBattery} node at } \rho_i := \mathsf{STrue} \ \text{for every pair} \\
  \text{(absorbing families and } J_i \ \text{syntactically empty; ranking decay off } [\![\varphi_i]\!] \ \text{alone)}
}{
  \mathcal{M},\ x_0 \in \{V_0 \le y\} \vDash \mathsf{SProbGE}\bigl(p,\ \mathsf{PDwellStreett}(\mathsf{PairBase}(\mathsf{STrue}, \varphi_1), \ldots, \mathsf{PairBase}(\mathsf{STrue}, \varphi_k))\bigr)
}
\end{mathpar}

\ruleflavor{the pair at request $\top$ means $\mathrm{Inf}([\![\varphi_i]\!])$ --- the classical recurrence certificate; the degenerate-pair specialization erases the tautological rows before emission, so the emitted kit is the recurrence kit row for row.}
\par\medskip
\rulename{QStreettDegenerate} is the recurrence rule in its compact
spelling: to conclude, with probability at least $p$, that every response
region is visited infinitely often ($\mathsf{GF}$-shaped claims, the
classical B\"uchi/recurrence certificate), instantiate the Streett bundle
at the trivially-true request. It is an \emph{alias} (n\'e StreettDegenerate --- the first sub-rule of the battery row): one
\rulename{QStreettBattery} application, at request regions equal to the
whole state space, is the entire unpacking --- the emptiness it exploits
is syntactic, not semantic. When every state requests, the
finitely-many-requests disjunct of each pair is empty along the tail
filter, the absorbing-certificate family and its carved-out sets
$J_i$ are vacuous, and the ranking decay row is demanded off the
response region alone --- so the degenerate-pair specialization erases
exactly the tautological rows, and what the checker replays is the
recurrence certificate, row for row. A producer uses it wherever the
specification is recurrence rather than
request-response fairness; nothing about the certificate synthesis
changes except that the request machinery never materializes.

Scope, stated plainly: the
bound statement below is the degenerate-instance witness
\texttt{sound\_QStreettDegenerate} (module
\texttt{BernsteinLogicAux.QStreett} --- one application of
\texttt{sound\_QStreettBattery} at request set \texttt{univ}, whose
finitely-many-requests disjunct is empty along the tail filter; its
conclusion is the recurrence event
$\Pr[\forall i,\ \mathrm{Inf}(G_i)] \ge p$, exactly the amended
display\textquotesingle s claim). The witness covers the rule at the
\emph{discrete-kernel profile} (a time-homogeneous Markov chain
presented by a single step kernel); this entry once carried a
conditional mark inherited from the bundle\textquotesingle s
continuous-time residue, and it goes the way the
bundle\textquotesingle s went: at the martingale-problem profile the
bundle is discharged whole --- strict-decrement rows included,
hash-bound under the \rulename{QStreettBattery} entry --- and the
alias inherits that discharge the same way it inherits everything
else, as one bundle application at request set \texttt{univ}, per the
bundle section\textquotesingle s profile remark. The payload-level
unpack function is not among this edition\textquotesingle s bound
artefacts: the alias deliberately has no wire sort of its own
(producers elaborate before emission), and the transport table in the
proof is stated at the catalogued-field level.

\cataloguenote{The degenerate pair is also the honest cost anchor of the
Streett-primitive design: on recurrence-shaped workloads the bundle
emits the same rows a dedicated recurrence rule would, so the general
rule costs nothing where the special case suffices --- while the reverse
specialization (recurrence-primitive, Streett by automaton product)
would pay exponentially in the number of pairs.
}
\ruleblock{Lean statement}
\declhead{BernsteinLogicAux.QStreett}{sound\_QStreettDegenerate}{\badge{BadgeGreen}{statement bound}}
\begin{leansrc}
/-- **`QStreettDegenerate` (T9a)**: request `= ⊤` degenerates the pair to
`Inf(Gᵢ)` — the classical recurrence certificate, one `sound_QStreettBattery`
application with the request set `univ` (whose `Fin` is empty along `atTop`).
Alias tier: the statement is the bundle's, specialized and simplified. -/
theorem sound_QStreettDegenerate {G : ι → Set S} {W : ι → S → ℝ} {ε C : ι → ℝ}
    (hD : MeasurableSet D) (hDbar : MeasurableSet Dbar) (hU : MeasurableSet U)
    (hV₀ : Measurable V₀)
    (hG : ∀ i, MeasurableSet (G i)) (hW : ∀ i, Measurable (W i))
    (hΘ : p + thr ≤ 1)
    (hDsub : D ⊆ Dbar)
    (hR3 : ∀ y ∈ D, y ∉ U → K y Dbar = 1)
    (hR4 : ∀ y ∈ D, y ∉ U → K y (D ∪ U) = 1)
    (hV₀int : ∀ y ∈ D, y ∉ U → Integrable V₀ (K y))
    (hV₀drift : ∀ y ∈ D, y ∉ U → ∫ z, V₀ z ∂K y ≤ V₀ y)
    (hV₀floor : ∀ y ∈ U, y ∈ Dbar → 1 ≤ V₀ y)
    (hV₀nonneg : ∀ y ∈ Dbar, 0 ≤ V₀ y)
    (hε : ∀ i, 0 < ε i)
    (hW0 : ∀ i, ∀ y ∈ Dbar, 0 ≤ W i y) (hWC : ∀ i, ∀ y ∈ Dbar, W i y ≤ C i)
    (hWint : ∀ i, ∀ y ∈ D, y ∉ U → y ∉ G i → Integrable (W i) (K y))
    (hWdriftR : ∀ i, ∀ y ∈ D, y ∉ U → y ∉ G i → ∫ z, W i z ∂K y ≤ W i y - ε i)
    (hx : x ∈ D) (hthr : V₀ x ≤ thr) :
    ENNReal.ofReal p ≤ pathMeasure K x
      {ω | ∀ i, ∃ᶠ n in atTop, ω n ∈ G i}
\end{leansrc}
\par\medskip
\textbf{The general-stochastic instance.}\enspace The degenerate
reading needs no separate statement at the segment-model tier: it rides
the confined form of the bundle instance below with the vacuous premise
sets --- the same instantiation-at-the-skeleton the bundle
section\textquotesingle s general-stochastic block describes.

\declhead{BernsteinLogicAux.LayerD.Streett}{sound\_GSHS\_streett\_bundle\_confined}{\badge{BadgeGreen}{statement bound}}
\begin{leansrc}
/-- **`sound_QStreettBattery`, GSHS instance, confined form** (T9; assault
doc §2.7 Theorem G at the conservative ranking scope, mechanized
restart-free — module docstring for why G1/G2's epoch machinery is not
consumed). Hypothesis fields are segment-level and per-model: the three
containment rows, the θ = 0 confinement field for `V₀`, and per pair the
scoped ranking fields for `W i` (nonincrease off the response region,
`ε i`-decrease at request entries — per-step shape, module docstring for
the rate-form conversion). Conclusion: with probability at least `p` the
skeleton chain of `Rep(P)` stays alive in `D \ U` forever **and** every
certified pair satisfies the chain-level Streett condition. Semantic
transport of the chain events: module docstring (T9b monotonicity for
requests, entry reads for responses, W-Z for time-cofinality). -/
theorem sound_GSHS_streett_bundle_confined
    [IsMarkovKernel M.seg] [IsMarkovKernel M.jump]
    (hD : MeasurableSet D) (hDbar : MeasurableSet Dbar) (hU : MeasurableSet U)
    (hV₀ : Measurable V₀)
    (hRq : ∀ i, MeasurableSet (Rq i)) (hG : ∀ i, MeasurableSet (G i))
    (hW : ∀ i, Measurable (W i))
    (hΘ : p + thr ≤ 1) (hDsub : D ⊆ Dbar)
    (hLive : NoVanish M (D \ U)) (hHalt : HaltsCaught M (D \ U) (U ∩ Dbar))
    (hReset : ResetsInto M (D \ U) ((D ∪ U) ∩ Dbar))
    (hSegD₀ : RankingSegmentDynkin M (D \ U) V₀ 0)
    (hForce₀ : ForcedNonexpansion M V₀) (hNonneg₀ : NonnegV V₀)
    (hC₀0 : 0 ≤ C₀) (hV₀C : ∀ y, V₀ y ≤ C₀)
    (hFloor₀ : ∀ y ∈ U, y ∈ Dbar → 1 ≤ V₀ y)
    (hε : ∀ i, 0 < ε i)
    (hSegDW : ∀ i, RankingSegmentDynkin M ((D \ U) \ G i) (W i) 0)
    (hSegDWR : ∀ i, RankingSegmentDynkin M (((D \ U) ∩ Rq i) \ G i) (W i) (ε i))
    (hForceW : ∀ i, ForcedNonexpansion M (W i)) (hNonnegW : ∀ i, NonnegV (W i))
    (hC0 : ∀ i, 0 ≤ C i) (hWC : ∀ i, ∀ y, W i y ≤ C i)
    (hx : x ∈ D) (hthr : V₀ x ≤ thr) :
    ENNReal.ofReal p ≤ Battery.pathMeasure (skel M) (SkelState.alive x)
      {ω | (∀ n, ω n ∈ skelIn D \ skelBad U) ∧
        ∀ i, (∃ᶠ n in atTop, ω n ∈ skelIn (Rq i)) →
          ∃ᶠ n in atTop, ω n ∈ skelIn (G i)}
\end{leansrc}
\ruleblock{Correspondence (Steele display $\leftrightarrow$ Lean)}
\badge{BadgeGreen}{reviewer-certified against the displayed texts}
\sourcenote{The dictionary below is a reviewer-certified statement of how the mechanized statements above interpret the displayed rule: display metavariable $\mapsto$ Lean binder, each display hypothesis $\mapsto$ the premise or side condition carrying it, the display conclusion $\mapsto$ the goal --- with any deliberate difference in width between display and mechanization named outright. Certification is recorded as a hash triple over (the displayed cell, the extracted statement texts, this dictionary\textquotesingle s own text); if any of the three changes, this block turns amber until a reviewer re-certifies. The correspondence claim itself remains editorial --- reviewer-grade, not machine-checked; its machine-checked successor types each soundness statement against the rule as a structured object, and supersedes this block per rule as it lands.}
The displayed conclusion $\mathcal{M}
,\ x_0 \in \{V_0 \le y\} \vDash
\mathsf{SProbGE}\bigl(p,\ \mathsf{PDwellStreett}(\mathsf{PairBase}(\mathsf{STrue},
\varphi_1), \ldots, \mathsf{PairBase}(\mathsf{STrue}, \varphi_k))\bigr)$
$\mapsto$ \texttt{ofReal p ≤ pathMeasure K x \{ω | ∀ i, ∃ᶠ n in atTop,
ω n ∈ G i\}} for every start with \texttt{x ∈ D} and \texttt{V₀ x ≤
thr} (the Lean source\textquotesingle s identifier for the
display\textquotesingle s $y$) --- the recurrence event, with the
$\mathsf{STrue}$-requests \emph{erased
at the statement level}, not merely in the proof: no request binder
exists in \texttt{sound\_QStreettDegenerate} at all, which is the
display\textquotesingle s ``absorbing families and $J_i$ syntactically
empty'' carried into the type. The alias tier is visible in the proof
term: one \texttt{sound\_QStreettBattery} application at \texttt{Rq :=
fun \_ => Set.univ}, whose frequently-clause is discharged by
\texttt{univ}-membership.
\begin{itemize}
\item the inherited $V_0$ QSafetyBattery kit $\mapsto$ the four
  kernel-integral rows verbatim: \texttt{hV₀drift} (interior),
  \texttt{hV₀floor} (floor on $U \cap \bar{D}$), \texttt{hV₀nonneg}
  (on $\bar{D}$), with \texttt{hV₀int} the undisplayed integrability
  rider, plus the hull/fail-closed side conditions
  \texttt{hR3}/\texttt{hR4} and the threshold \texttt{hΘ : p + thr ≤
  1} --- the battery pattern, unchanged.
\item ``ranking decay off $[\![\varphi_i]\!]$ alone'' $\mapsto$
  \texttt{hWdriftR : ∀ i, ∀ y ∈ D, y ∉ U → y ∉ G i → ∫ z, W i z ∂K y
  ≤ W i y - ε i} --- one generator row where the bundle carries two: the
  request-conditioned decay and the off-response non-increase
  coincide when every state is a request, and the statement asks
  only the merged row (the proof re-splits it for the bundle).
\item nonneg + upper-bound rows $\mapsto$ \texttt{hW0}/\texttt{hWC}
  on $\bar{D}$; strictness $\mapsto$ \texttt{hε : ∀ i, 0 < ε i} ---
  note the specialized statement takes per-pair \texttt{ε i}, mildly
  wider than the display\textquotesingle s single $\varepsilon_0$.
\item the measurability pack
  (\texttt{hD}/\texttt{hDbar}/\texttt{hU}/\texttt{hV₀}/\texttt{hG}/
  \texttt{hW}) is undisplayed profile regularity, per the battery
  precedent.
\end{itemize}
Width note: this alias\textquotesingle s warrant is inherited from
the bundle at both profiles (hash-bound at the discrete-kernel
profile; the martingale-problem discharge, strict rows included,
hash-bound under the bundle entry); the producer-side erasure of
tautological rows before emission is compile behaviour with no Lean
counterpart --- correctly, since the certified statement never
receives the tautological rows to begin with.

\ruleblock{Lean verifier receipts}
\sourcenote{module \texttt{BernsteinLogicAux.QStreett}; bound build \texttt{469881beb770\ldots} (bound from receipt 9d1d848e0ba1718a, 2026-07-29T01:03:31Z). A receipt records that named checking layers accepted a build: (module, compiled-artifact SHA-256, environment SHA-256, checker layers), minted when the checking gate ran; the store receipts attesting the bound build are listed with the layers each run credited. Nothing here asks the reader to trust the store: every layer is an independently runnable program over the quoted sources --- rebuild the module and re-run the checkers to reproduce the verdict. Receipts attesting superseded builds credit no layer and are omitted from this rendering (the count of omissions, if any, is noted below).}
\begin{receipttable}
\receiptrow{e3990fe544d4708d}{2026-07-28T23:48:35Z}{build+decorrelated}{\badge{BadgeGreen}{current}}
\receiptrow{68e9374fa1b418fe}{2026-07-29T00:00:34Z}{build+decorrelated}{\badge{BadgeGreen}{current}}
\receiptrow{e5895c9706854cd3}{2026-07-29T00:49:04Z}{build+replay+hardened+decorrelated}{\badge{BadgeGreen}{current}}
\receiptrow{9d1d848e0ba1718a}{2026-07-29T01:03:31Z}{build+replay+hardened+decorrelated}{\badge{BadgeGreen}{current}}
\receiptrow{4a8d377b721b8523}{2026-07-31T22:00:09Z}{build+decorrelated}{\badge{BadgeGreen}{current}}
\receiptrow{d251a122540ab741}{2026-07-31T22:29:03Z}{build+decorrelated}{\badge{BadgeGreen}{current}}
\receiptrow{333635b8e047113a}{2026-08-02T23:31:14Z}{build+decorrelated}{\badge{BadgeGreen}{current}}
\receiptrow{9e96f80d63e4ad82}{2026-08-03T16:27:18Z}{build+decorrelated}{\badge{BadgeGreen}{current}}
\receiptrow{932b4b9e0b0fd323}{2026-08-03T22:46:02Z}{build+decorrelated}{\badge{BadgeGreen}{current}}
\end{receipttable}
\sourcenote{9 further receipt(s) in the store attest superseded builds of this module and are omitted from this rendering.}
\sourcenote{module \texttt{BernsteinLogicAux.LayerD.Streett}; bound build \texttt{7393065bbc7f\ldots} (bound from receipt abac471866cfcfd1, 2026-07-21T05:38:48Z). A receipt records that named checking layers accepted a build: (module, compiled-artifact SHA-256, environment SHA-256, checker layers), minted when the checking gate ran; the store receipts attesting the bound build are listed with the layers each run credited. Nothing here asks the reader to trust the store: every layer is an independently runnable program over the quoted sources --- rebuild the module and re-run the checkers to reproduce the verdict. Receipts attesting superseded builds credit no layer and are omitted from this rendering (the count of omissions, if any, is noted below).}
\begin{receipttable}
\receiptrow{0f3b50b4d3f0d33c}{2026-07-20T22:52:50Z}{build+decorrelated}{\badge{BadgeGreen}{current}}
\receiptrow{109669f848453ec2}{2026-07-21T00:03:19Z}{build+decorrelated}{\badge{BadgeGreen}{current}}
\receiptrow{7e17aa32b276623b}{2026-07-21T04:51:49Z}{build+decorrelated}{\badge{BadgeGreen}{current}}
\receiptrow{abac471866cfcfd1}{2026-07-21T05:38:48Z}{build+replay+hardened+decorrelated}{\badge{BadgeGreen}{current}}
\end{receipttable}
\sourcenote{6 further receipt(s) in the store attest superseded builds of this module and are omitted from this rendering.}
\layerline{\layeritem{build}{lake build (elaboration)}{\badge{BadgeGreen}{current}}\layeritem{replay}{leanchecker (official kernel replay)}{\badge{BadgeGreen}{current}}\layeritem{hardened}{safe\_verify (hardened screen)}{\badge{BadgeGreen}{current}}\layeritem{decorrelated}{lean4export + nanoda (decorrelated kernel)}{\badge{BadgeGreen}{current}}}
\sourcenote{\badge{BadgeGray}{comparator: none} no fixture in the three-implementation comparator harness (compiled / kernel-interpreted / reference-interpreted) exercises this rule.}
\ruleblock{HOL4 statement}
\absentbox{OUT OF SCOPE BY DESIGN}{Out of scope for the HOL side \emph{by design}, not by backlog: the rule
is a producer-side respelling that unpacks to a single
\rulename{QStreettBattery} node before any certificate is emitted, so no
wire object can ever contain a \rulename{QStreettDegenerate} node --- the
HOL statement that would verify a degenerate-request bundle is the
bundle\textquotesingle s own (itself absent today, for the reason given
in that section: the satisfaction stratum has no HOL artefact).
Checked against the HOL sources: the
\texttt{bern\_catalogue} theory is the base-calculus checker cone and
carries nothing Streett-shaped.
}
\ruleblock{HOL verifier receipts}
\absentbox{NONE}{no HOL verification receipt exists for this rule.}
\ruleblock{Proof}

\sourcenote{Genre: alias rule (tier 3) --- a respelling template. This
section additionally states, once for the whole batch of
satisfaction-level aliases, the general template those sections share;
the \rulename{QStreettCmplRequest}, \rulename{QStreettTargetCellCollapse},
\rulename{QNnfRewrite}, \rulename{QStreettRequestNnf},
\rulename{QAvoid}, \rulename{QStreettGF}, \rulename{QStreettFG},
and \rulename{QNegFAlias} sections reference it by
this section's rule name instead of restating it.}

\begin{brecipe}[the satisfaction-level alias template]\label{rec:sd-template}
An alias rule $A$ at the satisfaction stratum is warranted by a $1 \to 1$
unpack into a single application of a target rule $T$ (or, for the
spec-surface rewrites, into a single definitional denotation equality):
the $A$-node and its image conclude the same judgment from the same
premise data, payload transported, nothing derived. Soundness is then a
corollary, $\mathrm{sound}_A = \mathrm{unpack} \fatsemi \mathrm{sound}_T$,
and the certificate format need never carry $A$: producers elaborate
before emission. The five artefacts an alias owes this book are the same
as at the base stratum (the \rulename{SmallGain} section is the model):
the unpack, the filled-in target application, the payload-transport
table, the exactness identity, and the soundness corollary. One honest
divergence from the base-stratum model, stated once here: no alias of
this batch has a \emph{verified payload function} among this
edition\textquotesingle s bound artefacts. The
transport tables of this batch are therefore stated at the
catalogued-field level, and the unpack-function slot of each alias
says so plainly; where a discrete-profile
\emph{semantic} witness exists in the auxiliary Lean development it is
named (the soundness content is machine-checked even though the
payload-level function is not), with its statement pin
riding the bindings record wherever the per-section text says
hash-bound, and where none exists the status is stated as filed. In
every case it is the payload-level slot, never the semantic witness,
that remains outside the bound warrant; the per-section text says which
case holds.
\end{brecipe}

\begin{brecipe}[the filled-in application --- \rulename{QStreettDegenerate}
rewritten into its \rulename{QStreettBattery} instance]\label{rec:sd-filled}
A degenerate node claims, from the confinement certificate $V_0$ and
per-region ranking certificates $(W_i, \varepsilon_i)$, that with
probability $\ge p$ every response region $G_i$ is visited infinitely
often. Its image is the single bundle node at request regions
$R_i := S$ (every state requests):
\begin{mathpar}
\inferrule*[Right=\rulename{QStreettBattery}]{
  \text{confinement rows for } V_0 \\
  \bigl\{\vdash \text{nonincrease of } W_i \ \text{off } G_i\bigr\} \\
  \bigl\{\vdash \varepsilon_i\text{-decay of } W_i \ \text{off } G_i\bigr\}
  \ \text{(the request row at } R_i = S\text{)} \\\\
  p + y \le 1 \ \text{(threshold tie)}
}{
  \mathcal{M},\ x_0 \in \{V_0 \le y\} \vDash
  \mathsf{SProbGE}\bigl(p,\ \mathsf{PDwellStreett}(\mathsf{PairBase}(\mathsf{STrue},
  \varphi_1), \ldots, \mathsf{PairBase}(\mathsf{STrue}, \varphi_k))\bigr)
}
\end{mathpar}
Discharging the bundle's obligations at this instance: the request
regions are measurable trivially; the decay row \emph{subsumes} the
nonincrease row (off $G_i$, every state is a request state, so
nonincrease follows from decay with the kit\textquotesingle s
$\varepsilon_i > 0$ --- nonnegativity is all the subsumption spends); the
absorbing families $(V_i, \gamma_i)$ and carved sets $J_i$ are vacuous
--- each $J_i$ lives inside $I \smallsetminus [\![\rho_i]\!]$, with $I$
the confinement region of the bundle\textquotesingle s carved-scope
formula, and at $\rho_i = \top$ that base is
$I \smallsetminus S = \varnothing$, so
$J_i = \varnothing$ syntactically, the rows carved to $J_i$ and to
off-request regions are rows over the empty region, and the
specialization targets the bundle's conservative (absorber-free)
scope --- the scope its ranking-route soundness theorem consumes ---
so no request-region absorber obligation survives. The dead premise slots are honestly dead, and the
degenerate-pair specialization pass erases them before emission.
\end{brecipe}

\begin{brecipe}[payload transport, catalogued-field level]\label{rec:sd-transport}
\begin{center}\small
\begin{tabular}{@{}ll@{}}
\textbf{degenerate node} & \textbf{\rulename{QStreettBattery} node}\\
\midrule
$V_0$ rows, threshold, init & verbatim\\
response regions $G_i$, certificates $(W_i, \varepsilon_i)$ & verbatim (bound rows $0 \le W_i \le C_i$ and strict $\varepsilon_i > 0$ included)\\
shared side conditions (hull, fail-closed) + measurability pack & verbatim (the battery pattern\textquotesingle s)\\
--- & request regions $:= S$ (each pair; measurable trivially)\\
--- & absorbing families: empty; $J_i = \varnothing$ (syntactic)\\
decay rows off $G_i$ & request-decay rows (at $R_i = S$); nonincrease subsumed\\
\end{tabular}
\end{center}
Unpack-function slot: no verified payload function in this edition
(Recipe~\ref{rec:sd-template}); the semantic witness exists ---
\texttt{sound\_QStreettDegenerate}, one \texttt{sound\_QStreettBattery}
application at request set \texttt{univ} (auxiliary module; hash-bound
in this edition).
\end{brecipe}

\begin{bprop}[the exactness identity]\label{prop:sd-exact}
For every path $\omega$ and pair $i$, at $R_i = S$:
\[
\bigl(\exists^\infty n,\ \omega_n \in S\bigr) \Rightarrow
\bigl(\exists^\infty n,\ \omega_n \in G_i\bigr)
\quad\Longleftrightarrow\quad
\exists^\infty n,\ \omega_n \in G_i .
\]
\end{bprop}
\begin{proof}
The hypothesis $\exists^\infty n,\ \omega_n \in S$ holds of every path
(every time is a request time), so the implication and its conclusion are
equivalent --- one line, which is the point. Equivalently, in acceptance
vocabulary: the finitely-many-requests disjunct
$\mathrm{Fin}(S)$ is empty along the tail filter, so the pair condition
$\mathrm{Fin}(R_i) \vee \mathrm{Inf}(G_i)$ collapses to
$\mathrm{Inf}(G_i)$.
\end{proof}

\begin{bprop}[soundness corollary]\label{prop:sd-sound}
Under the degenerate node's rows,
$\mu_x\{\omega \mid \forall i,\ \exists^\infty n,\ \omega_n \in G_i\}
\ge p$.
\end{bprop}
\begin{proof}
Apply the bundle's soundness theorem (the \rulename{QStreettBattery}
section's --- at its proven scope, the discrete-kernel profile and
conservative ranking scope; the bundle\textquotesingle s
martingale-problem discharge transfers to this alias the same way,
one application at the trivially-true request) at $R_i = S$ and
rewrite the event by
Proposition~\ref{prop:sd-exact}. There is no other content --- the alias
genre working as intended. (Machine-checked as
\texttt{sound\_QStreettDegenerate}; hash-bound in this edition.)
\end{proof}

\begin{bremark}[this is the B\"uchi case]
A B\"uchi condition \emph{is} the one-pair Streett condition with
trivially-true request, so this alias is where the classical
recurrence certificate lives inside the Streett-primitive design.
Deterministic B\"uchi automata recognize only limit languages
\cite{landweber1969}, so the reverse direction --- expressing a general
Streett condition B\"uchi-primitively --- is not a respelling at all but
an automaton re-encoding, exponential in the number of pairs for the
classical translations; the asymmetry is why the bundle,
not the recurrence rule, is the primitive (\cite{streett1982,thomas1990};
see the bundle section's remark).
\end{bremark}

\begin{anchors}
Recipe \ref{rec:sd-filled} target & the \rulename{QStreettBattery} section, Theorem \& hypotheses (H1)--(H5)\\
Proposition \ref{prop:sd-exact} & syntactic; no formal anchor needed (one line of filter algebra)\\
Proposition \ref{prop:sd-sound} & \texttt{sound\_QStreettDegenerate} (\texttt{BernsteinLogicAux.QStreett}; machine-checked, hash-bound in this edition)\\
no wire tag, by design & producers elaborate before emission (Recipe~\ref{rec:sd-template})\\
\end{anchors}


\rulesection{51}{QStreettCmplRequest}{admissible rule}
\label{blcatrule-QStreettCmplRequest}
\ruleblock{The rule (Steele metanotation)}
\sourcenote{as realized in the mechanized development; source row: \texttt{\texttt{QTSL-rules.toml}} --- \emph{QStreettCmplRequest}. The metanotation statement is typeset by hand for notational quality; its mechanized ground truth is the Lean statement below, extracted from the sources at every build of this book and hash-bound. The rule \emph{roster} is machine-checked (every mechanized rule tag round-trips to exactly one catalogue row); the statement \emph{text} is editorial, checked by eye against the extracted Lean beside it.}
\begin{mathpar}
\inferrule*{
  \text{a \rulename{QStreettBattery} kit at certified regions } (R_i, G_i) \\
  \bigl\{ [\![\mathrm{cmpl}\,\rho_i]\!] \subseteq R_i \bigr\} \ \text{(requests over-approximated)} \\
  \bigl\{ G_i \subseteq [\![\varphi_i]\!] \bigr\} \ \text{(responses under-approximated)}
}{
  \mathcal{M},\ x_0 \in \{V_0 \le y\} \vDash \mathsf{SProbGE}\bigl(p,\ \mathsf{PDwellStreett}(\mathsf{PairCmpl}(\rho_1, \varphi_1), \ldots)\bigr)
}
\end{mathpar}

\ruleflavor{requests sit on the hypothesis side of each pair implication, so the orientation swap of the complement (the workhorse judgment certifies under-approximations; complements arrive as over-approximations) is exactly the sound direction, for free.}
\par\medskip
\rulename{QStreettCmplRequest} lets a Streett pair carry its request as a
\emph{complement}: the pair $(\mathrm{cmpl}\,\rho_i, \varphi_i)$ asks
that if states \emph{outside} $[\![\rho_i]\!]$ recur forever, responses
recur forever. The satisfaction judgment\textquotesingle s workhorse
orientation certifies regions that \emph{under}-approximate a
formula\textquotesingle s denotation --- one orientation, deliberately:
the complement of an under-approximation is an
\emph{over}-approximation, so a complement-typed request arrives with
the opposite orientation for free. The rule (n\'e StreettCmplRequest) is \emph{admissible}, not alias: its unpacking is one
\rulename{QStreettBattery} node \emph{plus} region-weakening premises ---
the certified request regions over-approximate the complement
denotations, the certified response regions under-approximate the
response denotations --- and the direction check that makes this the
sound weakening is genuine content, not transport. Requests sit on the
hypothesis side of each pair implication and responses on the conclusion
side, so over-approximating requests and under-approximating responses
each \emph{weaken the claimed event}: any path satisfying the certified
pair satisfies the semantic pair. What the checker recomputes: nothing
new in kind --- the weakening premises are ordinary region-containment
rows it already knows how to verify, and
the complement itself is consumed at the region-atom level (the atom
dualities of the \rulename{QNnfRewrite} section), never as a formula
node.

Scope, stated plainly (this entry is marked \emph{conditional}): the
bound statements below are the discrete-profile witnesses --- the
event-level direction lemma \texttt{streett\_event\_mono} (two
monotonicity steps, exactly the two lines the recipe below typesets) and
the assembled bundle-through-formula-tie
\texttt{sound\_QStreettBattery\_tie} (requests over-approximated,
responses under-approximated, concluded over the semantic satisfaction
predicates), both in module \texttt{BernsteinLogicAux.QStreett}, plus
the atom-level complement step \texttt{denot\_elimNeg} in
\texttt{BernsteinLogicAux.NNFDuality}. One named residue keeps the
flag --- and it is this entry\textquotesingle s own, not the
bundle\textquotesingle s (whose continuous-time residue is discharged
whole at the martingale-problem profile, strict rows included,
hash-bound under the \rulename{QStreettBattery} entry): the
continuous-time face of the complement step remains prose. At the
general-stochastic-hybrid profile, consuming a complemented
request requires fine-topology regularity of the region (that the
region is nearly Borel and contains its regular points), which is
argued in the design record.

\cataloguenote{Why one orientation of the workhorse judgment suffices, in one line: an
under-approximating region calculus is closed under the connectives the
rules consume, and where a complement appears, the swap to
over-approximation is exactly what the hypothesis side of an implication
wants --- demanding a second, over-approximating judgment kind would
duplicate every rule for zero content.
}
\ruleblock{Lean statement}
\declhead{BernsteinLogicAux.QStreett}{streett\_event\_mono}{\badge{BadgeGreen}{statement bound}}
\begin{leansrc}
/-- **`QStreettCmplRequest` (T9b), the event-level kernel** — the discrete face
of request/response event monotonicity: requests **over**-approximated and
responses **under**-approximated weaken the per-pair Streett event in the
sound direction. Two `Frequently.mono`s — the "two lines" the catalogue prices. -/
theorem streett_event_mono {Rq Rqhat G Ghat : Set S} (hR : Rq ⊆ Rqhat)
    (hG : Ghat ⊆ G) (ω : ℕ → S)
    (h : (∃ᶠ n in atTop, ω n ∈ Rqhat) → ∃ᶠ n in atTop, ω n ∈ Ghat) :
    (∃ᶠ n in atTop, ω n ∈ Rq) → ∃ᶠ n in atTop, ω n ∈ G
\end{leansrc}
\declhead{BernsteinLogicAux.QStreett}{sound\_QStreettBattery\_tie}{\badge{BadgeGreen}{statement bound}}
\begin{leansrc}
/-- **`QStreettCmplRequest` (T9b), assembled** — the bundle consumed through the
formula tie, direction-checked by request/response event monotonicity and the
antitone region-complement kernel: the certified request region
`Rq i` over-approximates the truth set of the semantic request `reqSat i` (for
`cmpl`-typed requests this is atom-level complement elimination's
`⟦cmpl φ⟧ ⊇ ⟦φ⟧ᶜ = ⟦¬φ⟧`), and the certified response region
under-approximates the semantic response. Conclusion over the semantic
satisfaction predicates; no measurability is asked of them. -/
theorem sound_QStreettBattery_tie {Rq G : ι → Set S} {W : ι → S → ℝ} {ε C : ι → ℝ}
    {reqSat respSat : ι → S → Prop}
    (hD : MeasurableSet D) (hDbar : MeasurableSet Dbar) (hU : MeasurableSet U)
    (hV₀ : Measurable V₀)
    (hRq : ∀ i, MeasurableSet (Rq i)) (hG : ∀ i, MeasurableSet (G i))
    (hW : ∀ i, Measurable (W i))
    (hreq : ∀ i y, reqSat i y → y ∈ Rq i)
    (hresp : ∀ i, ∀ y ∈ G i, respSat i y)
    (hΘ : p + thr ≤ 1)
    (hDsub : D ⊆ Dbar)
    (hR3 : ∀ y ∈ D, y ∉ U → K y Dbar = 1)
    (hR4 : ∀ y ∈ D, y ∉ U → K y (D ∪ U) = 1)
    (hV₀int : ∀ y ∈ D, y ∉ U → Integrable V₀ (K y))
    (hV₀drift : ∀ y ∈ D, y ∉ U → ∫ z, V₀ z ∂K y ≤ V₀ y)
    (hV₀floor : ∀ y ∈ U, y ∈ Dbar → 1 ≤ V₀ y)
    (hV₀nonneg : ∀ y ∈ Dbar, 0 ≤ V₀ y)
    (hε : ∀ i, 0 < ε i)
    (hW0 : ∀ i, ∀ y ∈ Dbar, 0 ≤ W i y) (hWC : ∀ i, ∀ y ∈ Dbar, W i y ≤ C i)
    (hWint : ∀ i, ∀ y ∈ D, y ∉ U → y ∉ G i → Integrable (W i) (K y))
    (hWdrift : ∀ i, ∀ y ∈ D, y ∉ U → y ∉ G i → ∫ z, W i z ∂K y ≤ W i y)
    (hWdriftR : ∀ i, ∀ y ∈ D, y ∉ U → y ∈ Rq i → y ∉ G i →
      ∫ z, W i z ∂K y ≤ W i y - ε i)
    (hx : x ∈ D) (hthr : V₀ x ≤ thr) :
    ENNReal.ofReal p ≤ pathMeasure K x
      {ω | ∀ i, (∃ᶠ n in atTop, reqSat i (ω n)) → ∃ᶠ n in atTop, respSat i (ω n)}
\end{leansrc}
\declhead{BernsteinLogicAux.NNFDuality}{denot\_elimNeg}{\badge{BadgeGreen}{statement bound}}
\begin{leansrc}
/-- **Denotation preservation of dual-atom elimination**: where each dual atom
denotes the complement (`hd` — the per-duality definitional axioms, e.g. the
`compl_le_atom` family above), elimination preserves the region. This is the
lawfulness of consuming a `cmpl`-typed request through swapped atoms
(atom-level complement elimination's atom step, the T9b input shape). -/
theorem denot_elimNeg {σ : Type*} (den : α → Set σ) (d : α → α)
    (hd : ∀ a, den (d a) = (den a)ᶜ) :
    ∀ φ : NForm α, (φ.elimNeg d).denot den = φ.denot den
\end{leansrc}
\par\medskip
\textbf{The general-stochastic instance, carved requests
discharged.}\enspace The complemented-request form at the segment-model
tier: the carved variants below instantiate the bundle on the skeleton
chain with the complement-carved request regions supplied and their
side conditions discharged --- plain and confined forms.

\declhead{BernsteinLogicAux.LayerD.Streett}{sound\_GSHS\_streett\_bundle\_carved\_discharged}{\badge{BadgeGreen}{statement bound}}
\begin{leansrc}
/-- **The full bundle datum, GSHS instance** (T9, carved scope, absorbing
certificates as segment-level rows) — the discrete-statement-shaped
corollary of the `_confined` form; matches
`QStreett.sound_QStreettBattery_carved_discharged` line for line. -/
theorem sound_GSHS_streett_bundle_carved_discharged
    [IsMarkovKernel M.seg] [IsMarkovKernel M.jump]
    (hD : MeasurableSet D) (hDbar : MeasurableSet Dbar) (hU : MeasurableSet U)
    (hV₀ : Measurable V₀)
    (hRq : ∀ i, MeasurableSet (Rq i)) (hG : ∀ i, MeasurableSet (G i))
    (hJ : ∀ i, MeasurableSet (J i)) (hW : ∀ i, Measurable (W i))
    (hΘ : p + thr ≤ 1) (hDsub : D ⊆ Dbar)
    (hLive : NoVanish M (D \ U)) (hHalt : HaltsCaught M (D \ U) (U ∩ Dbar))
    (hReset : ResetsInto M (D \ U) ((D ∪ U) ∩ Dbar))
    (hSegD₀ : RankingSegmentDynkin M (D \ U) V₀ 0)
    (hForce₀ : ForcedNonexpansion M V₀) (hNonneg₀ : NonnegV V₀)
    (hC₀0 : 0 ≤ C₀) (hV₀C : ∀ y, V₀ y ≤ C₀)
    (hFloor₀ : ∀ y ∈ U, y ∈ Dbar → 1 ≤ V₀ y)
    (hε : ∀ i, 0 < ε i) (hγ : ∀ i, 0 < γ i)
    (hSegDW : ∀ i, RankingSegmentDynkin M ((D \ U) \ (G i ∪ J i)) (W i) 0)
    (hSegDWR : ∀ i,
      RankingSegmentDynkin M (((D \ U) ∩ Rq i) \ (G i ∪ J i)) (W i) (ε i))
    (hForceW : ∀ i, ForcedNonexpansion M (W i)) (hNonnegW : ∀ i, NonnegV (W i))
    (hC0 : ∀ i, 0 ≤ C i) (hWC : ∀ i, ∀ y, W i y ≤ C i)
    (hVa : ∀ i, Measurable (Va i)) (hNonnegVa : ∀ i, NonnegV (Va i))
    (hCa0 : ∀ i, 0 ≤ Ca i) (hVaC : ∀ i, ∀ y, Va i y ≤ Ca i)
    (hVaFloor : ∀ i, ∀ y ∈ Rq i, 1 ≤ Va i y)
    (hSegDVa : ∀ i, RankingSegmentDynkin M ((Rq i)ᶜ) (Va i) 0)
    (hForceVa : ∀ i, ForcedNonexpansion M (Va i))
    (hVaJ : ∀ i, ∀ y ∈ J i, Va i y ≤ 1 - γ i)
    (hx : x ∈ D) (hthr : V₀ x ≤ thr) :
    ENNReal.ofReal p ≤ Battery.pathMeasure (skel M) (SkelState.alive x)
      {ω | ∀ i, (∃ᶠ n in atTop, ω n ∈ skelIn (Rq i)) →
        ∃ᶠ n in atTop, ω n ∈ skelIn (G i)}
\end{leansrc}
\declhead{BernsteinLogicAux.LayerD.Streett}{sound\_GSHS\_streett\_bundle\_carved\_discharged\_confined}{\badge{BadgeGreen}{statement bound}}
\begin{leansrc}
/-- **The full bundle datum, GSHS instance, confined form** (T9 with the
absorbing certificates; assault doc §2.7 Theorem G proper, at the
hypothesis-carrying ranking scope `I \ (⟦φᵢ⟧ ∪ Jᵢ)`): the ranking rows are
demanded only off the absorbed region `J i`, and in exchange each pair
carries an absorbing certificate `Va i` — nonnegative, bounded, `≥ 1` on
the request region, segment-nonincreasing off it, `≤ 1 − γ i` on `J i`.
The absorption bound is discharged on the skeleton by
`QStreett.absorptionBound_of_kernel` (the shifted supermartingale bridge +
conditional Ville — the mechanized replacement for prose-G2's epoch
argument); its off-request drift row is demanded at **every alive state**
(scope `(Rq i)ᶜ`), the global-sharp reading of the discrete discharge —
frozen states satisfy it for free. -/
theorem sound_GSHS_streett_bundle_carved_discharged_confined
    [IsMarkovKernel M.seg] [IsMarkovKernel M.jump]
    (hD : MeasurableSet D) (hDbar : MeasurableSet Dbar) (hU : MeasurableSet U)
    (hV₀ : Measurable V₀)
    (hRq : ∀ i, MeasurableSet (Rq i)) (hG : ∀ i, MeasurableSet (G i))
    (hJ : ∀ i, MeasurableSet (J i)) (hW : ∀ i, Measurable (W i))
    (hΘ : p + thr ≤ 1) (hDsub : D ⊆ Dbar)
    (hLive : NoVanish M (D \ U)) (hHalt : HaltsCaught M (D \ U) (U ∩ Dbar))
    (hReset : ResetsInto M (D \ U) ((D ∪ U) ∩ Dbar))
    (hSegD₀ : RankingSegmentDynkin M (D \ U) V₀ 0)
    (hForce₀ : ForcedNonexpansion M V₀) (hNonneg₀ : NonnegV V₀)
    (hC₀0 : 0 ≤ C₀) (hV₀C : ∀ y, V₀ y ≤ C₀)
    (hFloor₀ : ∀ y ∈ U, y ∈ Dbar → 1 ≤ V₀ y)
    (hε : ∀ i, 0 < ε i) (hγ : ∀ i, 0 < γ i)
    (hSegDW : ∀ i, RankingSegmentDynkin M ((D \ U) \ (G i ∪ J i)) (W i) 0)
    (hSegDWR : ∀ i,
      RankingSegmentDynkin M (((D \ U) ∩ Rq i) \ (G i ∪ J i)) (W i) (ε i))
    (hForceW : ∀ i, ForcedNonexpansion M (W i)) (hNonnegW : ∀ i, NonnegV (W i))
    (hC0 : ∀ i, 0 ≤ C i) (hWC : ∀ i, ∀ y, W i y ≤ C i)
    (hVa : ∀ i, Measurable (Va i)) (hNonnegVa : ∀ i, NonnegV (Va i))
    (hCa0 : ∀ i, 0 ≤ Ca i) (hVaC : ∀ i, ∀ y, Va i y ≤ Ca i)
    (hVaFloor : ∀ i, ∀ y ∈ Rq i, 1 ≤ Va i y)
    (hSegDVa : ∀ i, RankingSegmentDynkin M ((Rq i)ᶜ) (Va i) 0)
    (hForceVa : ∀ i, ForcedNonexpansion M (Va i))
    (hVaJ : ∀ i, ∀ y ∈ J i, Va i y ≤ 1 - γ i)
    (hx : x ∈ D) (hthr : V₀ x ≤ thr) :
    ENNReal.ofReal p ≤ Battery.pathMeasure (skel M) (SkelState.alive x)
      {ω | (∀ n, ω n ∈ skelIn D \ skelBad U) ∧
        ∀ i, (∃ᶠ n in atTop, ω n ∈ skelIn (Rq i)) →
          ∃ᶠ n in atTop, ω n ∈ skelIn (G i)}
\end{leansrc}
\ruleblock{Correspondence (Steele display $\leftrightarrow$ Lean)}
\badge{BadgeGreen}{reviewer-certified against the displayed texts}
\sourcenote{The dictionary below is a reviewer-certified statement of how the mechanized statements above interpret the displayed rule: display metavariable $\mapsto$ Lean binder, each display hypothesis $\mapsto$ the premise or side condition carrying it, the display conclusion $\mapsto$ the goal --- with any deliberate difference in width between display and mechanization named outright. Certification is recorded as a hash triple over (the displayed cell, the extracted statement texts, this dictionary\textquotesingle s own text); if any of the three changes, this block turns amber until a reviewer re-certifies. The correspondence claim itself remains editorial --- reviewer-grade, not machine-checked; its machine-checked successor types each soundness statement against the rule as a structured object, and supersedes this block per rule as it lands.}
Three bound statements, three display constituents. The
orientation-swap kernel --- the flavor text\textquotesingle s
``exactly the sound direction, for free'' --- $\mapsto$
\texttt{streett\_event\_mono}: \texttt{Rq ⊆ Rqhat} (requests
over-approximated) and \texttt{Ghat ⊆ G} (responses
under-approximated) transport the per-pair Streett implication
pathwise; two \texttt{Frequently.mono}s, the priced two lines,
axiom-light and \emph{deterministic} --- no measure appears.
\begin{itemize}
\item the assembled rule $\mapsto$
  \texttt{sound\_QStreettBattery\_tie}: the full bundle kit (the
  battery-pattern binder block, as on the degenerate-request page)
  plus the two tie hypotheses, and the polarity pair is carried
  exactly --- \texttt{hreq : ∀ i y, reqSat i y → y ∈ Rq i} (the
  certified region \emph{over}-approximates the semantic request
  truth set: for $\mathrm{cmpl}$-typed requests this is the
  atom-level complement-elimination inclusion
  $[\![\mathrm{cmpl}\,\varphi]\!] \supseteq
  [\![\varphi]\!]^{\mathsf{c}}$) and \texttt{hresp : ∀ i, ∀ y ∈ G i,
  respSat i y} (responses \emph{under}-approximated). The conclusion
  event quantifies over the semantic predicates
  \texttt{reqSat}/\texttt{respSat}, of which \emph{no measurability
  is asked} (the measurable objects are the certified regions ---
  the QSafetyBattery \texttt{sat} pattern), and both tie hypotheses
  are safe-direction widenings in the QClosure polarity sense.
\item the atom-level complement step $\mapsto$ \texttt{denot\_elimNeg}
  (the NNF module): where each dual atom denotes the complement,
  consuming a $\mathrm{cmpl}$-typed request through swapped atoms
  preserves the region --- the lawfulness clause cited at
  QStreettRequestNnf\textquotesingle s consumption boundary (its
  flavor text).
\end{itemize}
Width note: the conditional mark\textquotesingle s one remaining
residue, exactly as the roster carries it --- the fine-topology
content of the complement step at the
general-stochastic-hybrid profile (nearly-Borel regions,
$E^r \subseteq E$), which remains design-note prose with no bound
counterpart; the bundle side\textquotesingle s former continuous-time
residue is discharged and hash-bound under the bundle entry. \texttt{streett\_event\_mono} is
itself profile-free (pure path logic); the profile enters only with
the assembled tie.

\ruleblock{Lean verifier receipts}
\sourcenote{module \texttt{BernsteinLogicAux.QStreett}; bound build \texttt{469881beb770\ldots} (bound from receipt 9d1d848e0ba1718a, 2026-07-29T01:03:31Z). A receipt records that named checking layers accepted a build: (module, compiled-artifact SHA-256, environment SHA-256, checker layers), minted when the checking gate ran; the store receipts attesting the bound build are listed with the layers each run credited. Nothing here asks the reader to trust the store: every layer is an independently runnable program over the quoted sources --- rebuild the module and re-run the checkers to reproduce the verdict. Receipts attesting superseded builds credit no layer and are omitted from this rendering (the count of omissions, if any, is noted below).}
\begin{receipttable}
\receiptrow{e3990fe544d4708d}{2026-07-28T23:48:35Z}{build+decorrelated}{\badge{BadgeGreen}{current}}
\receiptrow{68e9374fa1b418fe}{2026-07-29T00:00:34Z}{build+decorrelated}{\badge{BadgeGreen}{current}}
\receiptrow{e5895c9706854cd3}{2026-07-29T00:49:04Z}{build+replay+hardened+decorrelated}{\badge{BadgeGreen}{current}}
\receiptrow{9d1d848e0ba1718a}{2026-07-29T01:03:31Z}{build+replay+hardened+decorrelated}{\badge{BadgeGreen}{current}}
\receiptrow{4a8d377b721b8523}{2026-07-31T22:00:09Z}{build+decorrelated}{\badge{BadgeGreen}{current}}
\receiptrow{d251a122540ab741}{2026-07-31T22:29:03Z}{build+decorrelated}{\badge{BadgeGreen}{current}}
\receiptrow{333635b8e047113a}{2026-08-02T23:31:14Z}{build+decorrelated}{\badge{BadgeGreen}{current}}
\receiptrow{9e96f80d63e4ad82}{2026-08-03T16:27:18Z}{build+decorrelated}{\badge{BadgeGreen}{current}}
\receiptrow{932b4b9e0b0fd323}{2026-08-03T22:46:02Z}{build+decorrelated}{\badge{BadgeGreen}{current}}
\end{receipttable}
\sourcenote{9 further receipt(s) in the store attest superseded builds of this module and are omitted from this rendering.}
\sourcenote{module \texttt{BernsteinLogicAux.NNFDuality}; bound build \texttt{97564f8e3950\ldots} (bound from receipt 21f04961748c0302, 2026-07-29T01:03:33Z). A receipt records that named checking layers accepted a build: (module, compiled-artifact SHA-256, environment SHA-256, checker layers), minted when the checking gate ran; the store receipts attesting the bound build are listed with the layers each run credited. Nothing here asks the reader to trust the store: every layer is an independently runnable program over the quoted sources --- rebuild the module and re-run the checkers to reproduce the verdict. Receipts attesting superseded builds credit no layer and are omitted from this rendering (the count of omissions, if any, is noted below).}
\begin{receipttable}
\receiptrow{1417cdf27f69842a}{2026-07-28T23:48:35Z}{build+decorrelated}{\badge{BadgeGreen}{current}}
\receiptrow{65fc1b51055d38bd}{2026-07-29T00:00:33Z}{build+decorrelated}{\badge{BadgeGreen}{current}}
\receiptrow{2abbbc61b6065811}{2026-07-29T00:49:05Z}{build+replay+hardened+decorrelated}{\badge{BadgeGreen}{current}}
\receiptrow{21f04961748c0302}{2026-07-29T01:03:33Z}{build+replay+hardened+decorrelated}{\badge{BadgeGreen}{current}}
\receiptrow{8879c5b9721acf0c}{2026-07-31T22:00:03Z}{build+decorrelated}{\badge{BadgeGreen}{current}}
\receiptrow{e47a2c2d7f152711}{2026-07-31T22:28:57Z}{build+decorrelated}{\badge{BadgeGreen}{current}}
\receiptrow{ce1bb97f7cbcaf4f}{2026-08-02T23:31:07Z}{build+decorrelated}{\badge{BadgeGreen}{current}}
\receiptrow{8755b28f77990ba0}{2026-08-03T16:27:11Z}{build+decorrelated}{\badge{BadgeGreen}{current}}
\receiptrow{64948e60a56ae3bc}{2026-08-03T22:45:55Z}{build+decorrelated}{\badge{BadgeGreen}{current}}
\end{receipttable}
\sourcenote{8 further receipt(s) in the store attest superseded builds of this module and are omitted from this rendering.}
\sourcenote{module \texttt{BernsteinLogicAux.LayerD.Streett}; bound build \texttt{7393065bbc7f\ldots} (bound from receipt abac471866cfcfd1, 2026-07-21T05:38:48Z). A receipt records that named checking layers accepted a build: (module, compiled-artifact SHA-256, environment SHA-256, checker layers), minted when the checking gate ran; the store receipts attesting the bound build are listed with the layers each run credited. Nothing here asks the reader to trust the store: every layer is an independently runnable program over the quoted sources --- rebuild the module and re-run the checkers to reproduce the verdict. Receipts attesting superseded builds credit no layer and are omitted from this rendering (the count of omissions, if any, is noted below).}
\begin{receipttable}
\receiptrow{0f3b50b4d3f0d33c}{2026-07-20T22:52:50Z}{build+decorrelated}{\badge{BadgeGreen}{current}}
\receiptrow{109669f848453ec2}{2026-07-21T00:03:19Z}{build+decorrelated}{\badge{BadgeGreen}{current}}
\receiptrow{7e17aa32b276623b}{2026-07-21T04:51:49Z}{build+decorrelated}{\badge{BadgeGreen}{current}}
\receiptrow{abac471866cfcfd1}{2026-07-21T05:38:48Z}{build+replay+hardened+decorrelated}{\badge{BadgeGreen}{current}}
\end{receipttable}
\sourcenote{6 further receipt(s) in the store attest superseded builds of this module and are omitted from this rendering.}
\layerline{\layeritem{build}{lake build (elaboration)}{\badge{BadgeGreen}{current}}\layeritem{replay}{leanchecker (official kernel replay)}{\badge{BadgeGreen}{current}}\layeritem{hardened}{safe\_verify (hardened screen)}{\badge{BadgeGreen}{current}}\layeritem{decorrelated}{lean4export + nanoda (decorrelated kernel)}{\badge{BadgeGreen}{current}}}
\sourcenote{\badge{BadgeGray}{comparator: none} no fixture in the three-implementation comparator harness (compiled / kernel-interpreted / reference-interpreted) exercises this rule.}
\ruleblock{HOL4 statement}
\absentbox{NO HOL STATEMENT YET}{Absent, precisely located: the rule\textquotesingle s realization rides
the typed request field of the Streett bundle --- and the HOL4
development (\texttt{bern\_catalogue}, the base-calculus checker cone)
carries nothing at the satisfaction stratum, so there is no HOL-side
wire object for a checker to consume. A HOL statement becomes
meaningful only after the extension stratum is ported; the weakening rows
themselves are ordinary containment sequents the base checker already
covers. Checked against the HOL sources this edition.
}
\ruleblock{HOL verifier receipts}
\absentbox{NONE}{no HOL verification receipt exists for this rule.}
\ruleblock{Proof}

\sourcenote{Genre: admissible rule (tier 2) --- a proof-tree-rewriting
recipe (the general admissible template is stated in the
\rulename{PConstExact} section; the batch-level honesty clause about
absent payload functions is Recipe~\ref{rec:sd-template} of the
\rulename{QStreettDegenerate}
section and applies here verbatim).}

\begin{brecipe}[the \rulename{QStreettCmplRequest} instance --- one node,
one bundle, per-pair weakening premises]\label{rec:scr-instance}
A node concluding
$\mathcal{M},\ x_0 \in \{V_0 \le y\} \vDash
\mathsf{SProbGE}\bigl(p,\ \mathsf{PDwellStreett}(\mathsf{PairCmpl}(\rho_1,
\varphi_1), \ldots)\bigr)$ rewrites to the derivation fragment
\begin{mathpar}
\inferrule*[Right=\rulename{QStreettBattery}]{
  \text{the bundle kit at certified regions } (R_i, G_i) \\\\
  \bigl\{\, [\![\mathrm{cmpl}\,\rho_i]\!] \subseteq R_i \,\bigr\}_i
  \ \text{(request weakening)} \\
  \bigl\{\, G_i \subseteq [\![\varphi_i]\!] \,\bigr\}_i
  \ \text{(response weakening)}
}{
  \mathcal{M},\ x_0 \in \{V_0 \le y\} \vDash
  \mathsf{SProbGE}\bigl(p,\ \mathsf{PDwellStreett}(\mathsf{PairCmpl}(\rho_1,
  \varphi_1), \ldots)\bigr)
}
\end{mathpar}
where the request-weakening premises are supplied for free by the
orientation swap: a certified under-approximation $A_i \subseteq
[\![\rho_i]\!]$ yields the over-approximation $A_i^{\mathrm{c}} \supseteq
[\![\rho_i]\!]^{\mathrm{c}}$, and $R_i := A_i^{\mathrm{c}}$ is a region
atom by the atom dualities (\rulename{QNnfRewrite} section) wherever
$A_i$\textquotesingle s atom class carries a dual --- the
ordered-comparison classes, which are the shapes normalized request
bodies produce; an over-approximation of any other shape is simply
supplied as its own certified region --- and the
complement is consumed at the atom level, never as a formula node. The
weakening premises are ordinary containment rows the checker already
knows how to verify; nothing new \emph{in kind} is
recomputed.
\end{brecipe}

\begin{bprop}[the direction lemma --- two lines]\label{prop:scr-mono}
Fix a path $\omega$ and one pair. If the certified regions satisfy
$\mathrm{Rq} \subseteq \widehat{\mathrm{Rq}}$ (requests
over-approximated when read semantic-to-certified) and
$\widehat{G} \subseteq G$ (responses under-approximated), then
\[
\Bigl[\bigl(\exists^\infty n,\ \omega_n \in \widehat{\mathrm{Rq}}\bigr)
\Rightarrow \bigl(\exists^\infty n,\ \omega_n \in \widehat{G}\bigr)\Bigr]
\;\Longrightarrow\;
\Bigl[\bigl(\exists^\infty n,\ \omega_n \in \mathrm{Rq}\bigr)
\Rightarrow \bigl(\exists^\infty n,\ \omega_n \in G\bigr)\Bigr] .
\]
\end{bprop}
\begin{proof}
Frequent membership is monotone in the set: frequent
$\mathrm{Rq}$-visits are frequent $\widehat{\mathrm{Rq}}$-visits
(hypothesis side, first monotonicity), and frequent $\widehat{G}$-visits
are frequent $G$-visits (conclusion side, second monotonicity). Two
applications of monotonicity --- which is the entire lemma, and why the
weakening premises are the only new obligations of the rewrite.
(Machine-checked as \texttt{streett\_event\_mono}; hash-bound in this
edition.)
\end{proof}

\begin{bprop}[semantic warrant of the rewrite]\label{prop:scr-sound}
Under the bundle kit's rows and the weakening premises of
Recipe~\ref{rec:scr-instance}, the concluded judgment holds: with
probability $\ge p$, every semantic pair implication
$(\exists^\infty\ \text{states outside } [\![\rho_i]\!]) \Rightarrow
(\exists^\infty\ \text{states in } [\![\varphi_i]\!])$ is satisfied.
\end{bprop}
\begin{proof}
The bundle theorem gives the certified-region Streett event probability
$\ge p$; Proposition~\ref{prop:scr-mono}, applied per pair with
$\mathrm{Rq} := [\![\mathrm{cmpl}\,\rho_i]\!] \subseteq R_i$ and
$G_i \subseteq [\![\varphi_i]\!]$, includes that event in the semantic
one; the certified event is measurable (its regions are the
kit\textquotesingle s), and the path measure, read as outer measure on
the possibly-nonmeasurable semantic event, is monotone --- the lower
bound transfers. (Machine-checked, assembled, as
\texttt{sound\_QStreettBattery\_tie}; hash-bound in this edition.)
\end{proof}

\begin{bremark}[why the orientations are not symmetric, and not optional]
Swap either direction and the rule is unsound. If a certified request
region \emph{under}-approximated the complement, a path could request
infinitely often in the semantic sense while the certificate sees only
finitely many requests --- the certified pair implication holds
vacuously, the semantic one fails. Dually, a response region
\emph{over}-approximating $[\![\varphi_i]\!]$ lets certified
``responses'' land outside the actual response set. The hypothesis-side
/ conclusion-side asymmetry of implication is the entire content of the
direction check, and it is why the complement orientation swap --- which
produces exactly over-approximations --- is fitted to requests and to
nothing else.
\end{bremark}

\begin{bremark}[profile residue]
At the discrete-kernel profile the lemma chain above is complete,
machine-checked, and hash-bound in this edition. At the
continuous-time martingale-problem profile the bundle side is
discharged whole per the bundle section\textquotesingle s profile
remark (confinement, zero-decrement, and strict-decrement rows;
hash-bound there); what keeps this entry\textquotesingle s
conditional mark is the complement consumption\textquotesingle s own
fine-topology content (regular points of the
complement region), which remains prose in the design record
--- a residue of this entry\textquotesingle s own, not of the
bundle\textquotesingle s.
\end{bremark}

\begin{anchors}
Recipe \ref{rec:scr-instance} target & the \rulename{QStreettBattery} section\\
atom-level complement step & \texttt{denot\_elimNeg} (\texttt{BernsteinLogicAux.NNFDuality}; machine-checked, hash-bound)\\
Proposition \ref{prop:scr-mono} & \texttt{streett\_event\_mono} (\texttt{BernsteinLogicAux.QStreett}; machine-checked, hash-bound)\\
Proposition \ref{prop:scr-sound} & \texttt{sound\_QStreettBattery\_tie} (same module, hash-bound)\\
no wire object in this edition & producers rewrite before emission (Recipe~\ref{rec:scr-instance})\\
\end{anchors}


\rulesection{52}{QStreettTargetCellCollapse}{alias for \textcolor{DiffBlue}{\rulename{QStreettBattery}}}
\label{blcatrule-QStreettTargetCellCollapse}
\ruleblock{The rule (Steele metanotation)}
\sourcenote{as realized in the mechanized development; source row: \texttt{\texttt{QTSL-rules.toml}} --- \emph{QStreettTargetCellCollapse}. The metanotation statement is typeset by hand for notational quality; its mechanized ground truth is the Lean statement below, extracted from the sources at every build of this book and hash-bound. The rule \emph{roster} is machine-checked (every mechanized rule tag round-trips to exactly one catalogue row); the statement \emph{text} is editorial, checked by eye against the extracted Lean beside it.}
\begin{mathpar}
\inferrule*[right={$g(\mathrm{slab}) \subseteq T$,\ $W_i \equiv 0$ on $T$\ (G5, recomputed)}]{
  \mathcal{K},\ \mathrm{slab} \vdash W_i \ge 0
}{
  \mathcal{K},\ \mathrm{slab} \vdash g \fatsemi W_i \le W_i
}
\end{mathpar}
$\mathcal{K}$ ranges over \rulename{QStreettBattery} kits; the slab region left of the turnstile scopes both sequents.

\ruleflavor{a deduction rule between sequents: the kits demand the nonnegativity premise anyway, and under the side conditions (checker-recomputed, never a trusted annotation) the reset lands where the ranking certificate is identically zero, so the conclusion follows pointwise from the premise. On slabs whose resets leave the zero region the side condition fails and the full guard row is owed as usual.}
\par\medskip
\rulename{QStreettTargetCellCollapse} is a deduction rule between
sequents of one \rulename{QStreettBattery} kit context (n\'e StreettTargetCellCollapse): with $\mathcal{K}$ the premise metavariable
ranging over bundle kits, it derives the ranking-guard row
$g \fatsemi W_i \le W_i$ on a guard slab from the nonnegativity row
$W_i \ge 0$ there --- a row the kit demands anyway --- under two side
conditions: the containment of the reset image in the zero-convention
region ($g(\mathrm{slab}) \subseteq T$, with $W_i \equiv 0$ on $T$).
The rule is sound because under the side conditions the
conclusion\textquotesingle s left side is literally zero on the slab.
Both side conditions are recomputed geometric checks, never trusted
annotations --- which is also why the rule cannot be abused: on a slab
whose resets leave the zero region, the side condition fails and the
full guard row is owed as usual.

Mechanization state, plainly: the bound statement below is the one-line
witness \texttt{ranking\_guard\_collapse} (module
\texttt{BernsteinLogicAux.QStreett}) --- under the zero-on-target
convention and the recomputed containment, the guard row follows from
the nonneg row on the slab, which is the entirety of the
rule\textquotesingle s semantic claim (pure state-space set algebra: no
kernel, no profile, no probability enters, so no profile residue
attaches and the entry carries no conditional flag). The payload-level
unpack owes no sort of its own --- the extension certificate
schema\textquotesingle s adjudication of the batch family
(\rulename{QStreettDegenerate} section, Recipe~\ref{rec:sd-template}):
the collapse is a citation (the guard row cites the nonneg
row\textquotesingle s proof), and the containment is a recomputed side
condition of the battery walk, never a stored row --- so there is
nothing to elaborate, by design.

\cataloguenote{The checker\textquotesingle s implementation is citation-shaped,
deliberately decoupled from the deduction-rule presentation: no wire
object ever carries a collapse node --- a producer that has already
discharged the nonneg row cites it where the guard row was owed and
emits nothing new, one row to one row. Tier flag, carried honestly from
the source catalogue: that one-to-one citation shape is the alias
tiering, but the licensing side condition is a recomputed containment
--- if side-condition weight is ever ruled heavier than the alias
criterion intends, the rule re-tiers admissible with machinery
unchanged. Flagged, not forced.
}
\ruleblock{Lean statement}
\declhead{BernsteinLogicAux.QStreett}{ranking\_guard\_collapse}{\badge{BadgeGreen}{statement bound}}
\begin{leansrc}
/-- **`QStreettTargetCellCollapse` (T9c)**: when the reset of a guard transition
lands inside a region where the ranking certificate is identically zero (the
`Wᵢ ≡ 0`-on-target-cells convention, G5), the ranking-guard row
`W (g y) ≤ W y` collapses to the nonneg row `0 ≤ W y` — one row to one row,
under the recomputed containment side-condition `g '' slab ⊆ T`. -/
theorem ranking_guard_collapse {W : S → ℝ} {T slab : Set S} {g : S → S}
    (hzero : ∀ y ∈ T, W y = 0) (himg : ∀ y ∈ slab, g y ∈ T) :
    ∀ y ∈ slab, 0 ≤ W y → W (g y) ≤ W y
\end{leansrc}
\ruleblock{Correspondence (Steele display $\leftrightarrow$ Lean)}
\badge{BadgeGreen}{reviewer-certified against the displayed texts}
\sourcenote{The dictionary below is a reviewer-certified statement of how the mechanized statements above interpret the displayed rule: display metavariable $\mapsto$ Lean binder, each display hypothesis $\mapsto$ the premise or side condition carrying it, the display conclusion $\mapsto$ the goal --- with any deliberate difference in width between display and mechanization named outright. Certification is recorded as a hash triple over (the displayed cell, the extracted statement texts, this dictionary\textquotesingle s own text); if any of the three changes, this block turns amber until a reviewer re-certifies. The correspondence claim itself remains editorial --- reviewer-grade, not machine-checked; its machine-checked successor types each soundness statement against the rule as a structured object, and supersedes this block per rule as it lands.}
One bound statement, and the displayed deduction rule is its literal
shape: \texttt{ranking\_guard\_collapse} concludes \texttt{∀ y ∈ slab,
0 ≤ W y → W (g y) ≤ W y} --- the premise sequent as the antecedent, the
concluded ranking-guard row as the consequent, per point of the slab. The two
licensing conditions map as: $W_i \equiv 0$ on target cells
$\mapsto$ \texttt{hzero : ∀ y ∈ T, W y = 0} (the G5 convention as a
hypothesis), and the recomputed containment $g(\mathrm{slab})
\subseteq T$ $\mapsto$ \texttt{himg : ∀ y ∈ slab, g y ∈ T} ---
consumed as a hypothesis, recomputed checker-side, which is the
tier-ambiguity flag\textquotesingle s exact locus rendered as a
binder. Width note: the statement is one line over arbitrary sets
and an arbitrary reset map \texttt{g : S → S} --- no slab geometry,
no polynomiality of $W$, no guard structure --- strictly wider than
the displayed instance, safe direction; the failure branch the
flavor text names (resets leaving the zero region owe the full guard
row) is representable precisely as \texttt{himg} failing, whereupon
the theorem offers nothing --- refusal-shaped, as it should be.

\ruleblock{Lean verifier receipts}
\sourcenote{module \texttt{BernsteinLogicAux.QStreett}; bound build \texttt{469881beb770\ldots} (bound from receipt 9d1d848e0ba1718a, 2026-07-29T01:03:31Z). A receipt records that named checking layers accepted a build: (module, compiled-artifact SHA-256, environment SHA-256, checker layers), minted when the checking gate ran; the store receipts attesting the bound build are listed with the layers each run credited. Nothing here asks the reader to trust the store: every layer is an independently runnable program over the quoted sources --- rebuild the module and re-run the checkers to reproduce the verdict. Receipts attesting superseded builds credit no layer and are omitted from this rendering (the count of omissions, if any, is noted below).}
\begin{receipttable}
\receiptrow{e3990fe544d4708d}{2026-07-28T23:48:35Z}{build+decorrelated}{\badge{BadgeGreen}{current}}
\receiptrow{68e9374fa1b418fe}{2026-07-29T00:00:34Z}{build+decorrelated}{\badge{BadgeGreen}{current}}
\receiptrow{e5895c9706854cd3}{2026-07-29T00:49:04Z}{build+replay+hardened+decorrelated}{\badge{BadgeGreen}{current}}
\receiptrow{9d1d848e0ba1718a}{2026-07-29T01:03:31Z}{build+replay+hardened+decorrelated}{\badge{BadgeGreen}{current}}
\receiptrow{4a8d377b721b8523}{2026-07-31T22:00:09Z}{build+decorrelated}{\badge{BadgeGreen}{current}}
\receiptrow{d251a122540ab741}{2026-07-31T22:29:03Z}{build+decorrelated}{\badge{BadgeGreen}{current}}
\receiptrow{333635b8e047113a}{2026-08-02T23:31:14Z}{build+decorrelated}{\badge{BadgeGreen}{current}}
\receiptrow{9e96f80d63e4ad82}{2026-08-03T16:27:18Z}{build+decorrelated}{\badge{BadgeGreen}{current}}
\receiptrow{932b4b9e0b0fd323}{2026-08-03T22:46:02Z}{build+decorrelated}{\badge{BadgeGreen}{current}}
\end{receipttable}
\sourcenote{9 further receipt(s) in the store attest superseded builds of this module and are omitted from this rendering.}
\layerline{\layeritem{build}{lake build (elaboration)}{\badge{BadgeGreen}{current}}\layeritem{replay}{leanchecker (official kernel replay)}{\badge{BadgeGreen}{current}}\layeritem{hardened}{safe\_verify (hardened screen)}{\badge{BadgeGreen}{current}}\layeritem{decorrelated}{lean4export + nanoda (decorrelated kernel)}{\badge{BadgeGreen}{current}}}
\sourcenote{\badge{BadgeGray}{comparator: none} no fixture in the three-implementation comparator harness (compiled / kernel-interpreted / reference-interpreted) exercises this rule.}
\ruleblock{HOL4 statement}
\absentbox{OUT OF SCOPE BY DESIGN}{Out of scope for the HOL side \emph{by design}: the rule rewrites one
premise row of a bundle kit into another \emph{before}
emission --- no wire object ever contains a collapse node, and the
surviving row is an ordinary base-calculus sequent the base checker
already covers. The HOL statement that would matter is the
bundle\textquotesingle s own, absent for the reasons given in that
section. Checked against the HOL sources: the \texttt{bern\_catalogue}
theory carries nothing Streett-shaped.
}
\ruleblock{HOL verifier receipts}
\absentbox{NONE}{no HOL verification receipt exists for this rule.}
\ruleblock{Proof}

\sourcenote{Genre: alias rule (tier 3); the batch's shared template,
including the honesty clause about absent payload functions, is
Recipe~\ref{rec:sd-template} of the
\rulename{QStreettDegenerate} section. This section's rule is stated in
deduction-rule form --- the respelling content is carried by the
payload-transport recipe and the checker's citation-shaped
implementation, noted there.}

\begin{brecipe}[the rule --- a deduction between kit sequents]
\label{rec:tcc-filled}
Fix a \rulename{QStreettBattery} kit context $\mathcal{K}$, a pair $i$, a
guard transition with reset map $g$ and guard slab $\mathrm{slab}$, and
let $T$ be a region on which the kit's convention $W_i \equiv 0$ holds
(the target cells). The rule derives the ranking-guard row from the
nonnegativity row the kit demands anyway --- the kit's bound rows live
on the hull $\bar{D}$, so on guard slabs inside the hull the premise
below is a row already owed, while a slab outside them simply owes its
premise sequent separately; the deduction as displayed takes the
on-slab sequent itself as premise either way:
\begin{mathpar}
\inferrule*[right={$g(\mathrm{slab}) \subseteq T$,\ $W_i \equiv 0$ on $T$\ (G5, recomputed)}]{
  \mathcal{K},\ \mathrm{slab} \vdash W_i \ge 0
}{
  \mathcal{K},\ \mathrm{slab} \vdash g \fatsemi W_i \le W_i
}
\end{mathpar}
Both side conditions are checker-recomputed (the G5-convention check),
never producer-supplied. The premise metavariable does the honest work:
premise and conclusion are sequents of the \emph{same} kit context, so
admitting the rule claims exactly that any kit admitting the premise
admits the conclusion --- Proposition~\ref{prop:tcc-sound}, and nothing
else.
\end{brecipe}

\begin{brecipe}[payload transport, catalogued-field level]\label{rec:tcc-transport}
\begin{center}\small
\begin{tabular}{@{}ll@{}}
\textbf{collapse node} & \textbf{bundle kit row}\\
\midrule
guard row $(W_i, g, \mathrm{slab})$ & the kit's nonneg row for $W_i$ on $\mathrm{slab}$, cited (licit where the slab lies in the cited row's region)\\
containment $g(\mathrm{slab}) \subseteq T$ & recomputed side-condition (never producer-supplied)\\
zero convention $W_i \equiv 0$ on $T$ & the kit's target-cell convention, verbatim (that it holds on $T$ is the recomputed G5 check, not a trusted annotation)\\
\end{tabular}
\end{center}
The checker's implementation is citation-shaped: a producer that has
already discharged the nonnegativity row cites it where the guard row was
owed, and emits nothing new --- one row to one row, which is the alias
tiering. Presentation and implementation are deliberately decoupled: the
rule \emph{is} the deduction above; the citation is how a certificate
spells an application of it. Unpack-function slot: nothing to
elaborate, by design --- the collapse is a citation (the guard row
cites the nonneg row\textquotesingle s proof) and the containment is a
recomputed side condition, never a stored row; the semantic witness
\texttt{ranking\_guard\_collapse} (auxiliary module) is hash-bound in
this edition.
\end{brecipe}

\begin{bprop}[the exactness identity --- one line]\label{prop:tcc-exact}
For every $y \in \mathrm{slab}$: if $g(y) \in T$ and $W_i \equiv 0$ on
$T$ and $0 \le W_i(y)$, then $W_i(g(y)) = 0 \le W_i(y)$.
\end{bprop}
\begin{proof}
Substitution --- the displayed chain is the proof. The guard row's left
side is identically zero on the slab, so the guard row and the nonneg
row assert the same inequality there.
\end{proof}

\begin{bprop}[soundness corollary --- admissibility of the rule]
\label{prop:tcc-sound}
Any bundle kit admitting the nonneg row on the slab admits the
ranking-guard row on the slab --- the displayed rule, read as a claim;
hence every conclusion drawn through \rulename{QStreettBattery} from the
collapsed kit is drawn from the full kit, with the bundle's own
soundness theorem as warrant.
\end{bprop}
\begin{proof}
Proposition~\ref{prop:tcc-exact} produces the guard row pointwise from
the nonneg row under the side conditions --- that collapse is the
machine-checked content (\texttt{ranking\_guard\_collapse}; hash-bound
in this edition), and it is the whole of what this section proves. The
guard row is consumed \emph{upstream} of the flattened kernel: for a
hybrid presentation it is one of the base-calculus sequents from which
the kernel-integral ranking rows --- the bundle kit's H4, which the
bundle theorem is what actually consumes --- are assembled by the
presentation-to-kernel compile bridge (the safety
battery\textquotesingle s staged form carries that bridge as its
explicit \texttt{hmix} hypothesis; the bundle theorem takes the
compiled kernel and its H4 rows directly), never re-proven here. The corollary therefore lives at the kit stratum: the
collapsed kit discharges every base-calculus obligation the full
kit\textquotesingle s guard row was owed for, and every conclusion
drawn through \rulename{QStreettBattery} whose compile bridge owed that
guard row is drawn from the full kit, the bundle\textquotesingle s own
soundness theorem --- its assembly hypothesis unchanged --- supplying
the rest.
\end{proof}

\begin{bremark}[why the side-condition is not decoration]
Drop the containment and the collapse is unsound: a reset landing at a
state with $W_i(g(y)) > W_i(y)$ (outside the zero region) breaks the
upstream derivation of the ranking generator rows that the bundle's
telescoping lemma consumes --- the certificate could be pumped back up
at each guard
crossing, and the summability of un-responded requests
(\rulename{QStreettBattery} section) fails. The containment is
checker-recomputed for precisely this reason; its weight is the tier
flag recorded in the commentary.
\end{bremark}

\begin{anchors}
Recipe \ref{rec:tcc-filled} premise context $\mathcal{K}$ & the \rulename{QStreettBattery} section's kit (H4 rows)\\
Proposition \ref{prop:tcc-exact} & \texttt{ranking\_guard\_collapse} (\texttt{BernsteinLogicAux.QStreett}; machine-checked, hash-bound)\\
side conditions & the zero-on-target containment, recomputed (the bundle's G5-convention check)\\
\end{anchors}


\rulesection{53}{QNnfRewrite}{alias rule}
\label{blcatrule-QNnfRewrite}
\ruleblock{The rule (Steele metanotation)}
\sourcenote{as realized in the mechanized development; source row: \texttt{\texttt{QTSL-rules.toml}} --- \emph{QNnfRewrite}. The metanotation statement is typeset by hand for notational quality; its mechanized ground truth is the Lean statement below, extracted from the sources at every build of this book and hash-bound. The rule \emph{roster} is machine-checked (every mechanized rule tag round-trips to exactly one catalogue row); the statement \emph{text} is editorial, checked by eye against the extracted Lean beside it.}
\begin{mathpar}
\inferrule*{ }{ [\![\neg\neg\varphi]\!] = [\![\varphi]\!] }
\and
\inferrule*{ }{ [\![\neg(\varphi \wedge \psi)]\!] = [\![\neg\varphi \vee \neg\psi]\!] }
\and
\inferrule*{ }{ [\![\neg(\varphi \vee \psi)]\!] = [\![\neg\varphi \wedge \neg\psi]\!] }
\and
\inferrule*{ }{ [\![\neg\{f \le c\}]\!] = \{f > c\} }
\and
\inferrule*{ }{ [\![\neg\{f < c\}]\!] = \{f \ge c\} }
\and
\inferrule*{ }{ [\![\neg\{f \ge c\}]\!] = \{f < c\} }
\and
\inferrule*{ }{ [\![\neg\{f > c\}]\!] = \{f \le c\} }
\end{mathpar}

\ruleflavor{seven axioms --- deduction rules with nothing above the line --- at the spec surface, each concluding its region-denotation equality; a satisfaction judgment rewrites across any of them in either direction with certified region and rows transported verbatim, and no negation node survives to the baked representation (surviving negative literals are the typing rejection\textquotesingle s business, not a rule\textquotesingle s).}
\par\medskip
\rulename{QNnfRewrite} is the spec-surface negation calculus (n\'e NNFDuality): double negation, De Morgan through
$\wedge/\vee$, and the four atom/strict swaps
$\neg\{f \le c\} \leftrightarrow \{f > c\}$,
$\neg\{f < c\} \leftrightarrow \{f \ge c\}$ --- both orientations of
both comparison directions. The seven dualities are \emph{axioms}:
deduction rules with nothing above the line, each concluding its
region-denotation equality, stated once per duality. Each licenses a
judgment-level $1 \to 1$ rewrite of a formula inside a satisfaction
judgment: the rewritten formula denotes the same set of states, so any
region under-approximating one denotation under-approximates the other,
and the judgment transports unchanged. A
producer uses the dualities to push negations down to atoms (the
negation-normal-form pass of the \rulename{QStreettRequestNnf} section)
so that every negative position is consumed either by an atom swap or by
the typing pass\textquotesingle s rejection --- never by a negation
node in the certificate. The checker recomputes nothing: the dualities
are definitional set identities over the state space, and the
strict-inequality sides interact with the closure discipline (rounding
of strict literals at positive positions) only through the separate
closure rules, deliberately out of this rule\textquotesingle s scope.

There is no single target rule, and hence no \emph{alias-for} header
here: each axiom stands on its own definitional equality.
Mechanization state, plainly: the bound
statements below are the four atom swaps, machine-checked as set-level
denotation equalities in both orientations of both comparison
directions (module \texttt{BernsteinLogicAux.NNFDuality}:
\texttt{compl\_le\_atom}, \texttt{compl\_lt\_atom},
\texttt{compl\_ge\_atom}, \texttt{compl\_gt\_atom}), together with
the normalization pass\textquotesingle s denotation-preservation
theorem \texttt{denot\_nnfAux}, inside which double negation and De
Morgan are consumed from the ambient set library --- each applied once
per duality, so the displayed equalities are covered jointly by the
swaps (atoms) and the pass theorem (connectives). These are definitional
set identities over the state space, warranted in full; the entry
carries no conditional flag.

\cataloguenote{Both orientations of both directions are stated deliberately: an early
one-sided draft of the atom swaps was caught in review, and the repaired
statement is the four-lemma family. The strict-atom sides are stated
here as \emph{denotation} facts only; whether a strict literal may
appear in a concluded region is the closure discipline\textquotesingle s
question, answered by polarity-directed rounding in its own catalogue
rows.
}
\ruleblock{Lean statement}
\declhead{BernsteinLogicAux.NNFDuality}{compl\_le\_atom}{\badge{BadgeGreen}{statement bound}}
\begin{leansrc}
/-- T12 atom swap: `¬{f ≤ c} ↔ {f > c}`. -/
theorem compl_le_atom : {x | f x ≤ c}ᶜ = {x | c < f x}
\end{leansrc}
\declhead{BernsteinLogicAux.NNFDuality}{compl\_lt\_atom}{\badge{BadgeGreen}{statement bound}}
\begin{leansrc}
/-- T12 atom swap: `¬{f < c} ↔ {f ≥ c}`. -/
theorem compl_lt_atom : {x | f x < c}ᶜ = {x | c ≤ f x}
\end{leansrc}
\declhead{BernsteinLogicAux.NNFDuality}{compl\_ge\_atom}{\badge{BadgeGreen}{statement bound}}
\begin{leansrc}
/-- T12 atom swap, the other orientation: `¬{f ≥ c} ↔ {f < c}`. -/
theorem compl_ge_atom : {x | c ≤ f x}ᶜ = {x | f x < c}
\end{leansrc}
\declhead{BernsteinLogicAux.NNFDuality}{compl\_gt\_atom}{\badge{BadgeGreen}{statement bound}}
\begin{leansrc}
/-- T12 atom swap, the other orientation: `¬{f > c} ↔ {f ≤ c}`. -/
theorem compl_gt_atom : {x | c < f x}ᶜ = {x | f x ≤ c}
\end{leansrc}
\declhead{BernsteinLogicAux.NNFDuality}{denot\_nnfAux}{\badge{BadgeGreen}{statement bound}}
\begin{leansrc}
/-- **Denotation preservation of the NNF pass, both polarities** — the T12
dualities (double negation `compl_compl`, De Morgan `Set.compl_inter`/
`Set.compl_union`) are exactly the cases of this induction; each is applied
once, which is the catalogue's "stated once as a definitional axiom per
duality" realized as proof structure. -/
theorem denot_nnfAux : ∀ (φ : Form α) (b : Bool),
    (φ.nnfAux b).denot den = if b then φ.denot den else (φ.denot den)ᶜ
\end{leansrc}
\ruleblock{Correspondence (Steele display $\leftrightarrow$ Lean)}
\badge{BadgeGreen}{reviewer-certified against the displayed texts}
\sourcenote{The dictionary below is a reviewer-certified statement of how the mechanized statements above interpret the displayed rule: display metavariable $\mapsto$ Lean binder, each display hypothesis $\mapsto$ the premise or side condition carrying it, the display conclusion $\mapsto$ the goal --- with any deliberate difference in width between display and mechanization named outright. Certification is recorded as a hash triple over (the displayed cell, the extracted statement texts, this dictionary\textquotesingle s own text); if any of the three changes, this block turns amber until a reviewer re-certifies. The correspondence claim itself remains editorial --- reviewer-grade, not machine-checked; its machine-checked successor types each soundness statement against the rule as a structured object, and supersedes this block per rule as it lands.}
The displayed atom/strict swaps map one-to-one onto the four bound
equalities, each a region-denotation identity:
$[\![\neg\{f \le c\}]\!] = \{f > c\}$ $\mapsto$
\texttt{compl\_le\_atom : \{x | f x ≤ c\}ᶜ = \{x | c < f x\}}, and
likewise \texttt{compl\_lt\_atom} ($\neg\{f < c\} = \{f \ge c\}$),
\texttt{compl\_ge\_atom}, \texttt{compl\_gt\_atom} (the reversed
orientations --- both directions bound separately, honouring the
recorded one-sidedness catch). Binder note, wider-safe: \texttt{f : σ
→ ℝ} is an arbitrary function on an arbitrary carrier --- the
polynomial atom structure of the certificate format is erased, so
the certified equalities cover every instance the rule will ever be
applied to. The propositional dualities ($\neg\neg\varphi$, De
Morgan both ways) $\mapsto$ the \texttt{neg}/\texttt{and}/\texttt{or}
cases of \texttt{denot\_nnfAux}\textquotesingle s both-polarity
induction, where \texttt{compl\_compl} /
\texttt{Set.compl\_inter} / \texttt{Set.compl\_union} are each
consumed from mathlib once --- the entry\textquotesingle s
``stated once per duality'' realized as proof structure rather than
as four more standalone statements. Width note: the entry is
unconditional --- set-level witnesses cover the catalogued claim in
full (axioms: \texttt{propext} alone) --- and the
grammar-level guarantee the flavor text mentions (no negation node
survives to the baked representation) is certified not here but by
the QStreettRequestNnf entry\textquotesingle s normal-form \emph{type}; this
page\textquotesingle s statements are the denotation half only, as
its note already rules.

\ruleblock{Lean verifier receipts}
\sourcenote{module \texttt{BernsteinLogicAux.NNFDuality}; bound build \texttt{97564f8e3950\ldots} (bound from receipt 21f04961748c0302, 2026-07-29T01:03:33Z). A receipt records that named checking layers accepted a build: (module, compiled-artifact SHA-256, environment SHA-256, checker layers), minted when the checking gate ran; the store receipts attesting the bound build are listed with the layers each run credited. Nothing here asks the reader to trust the store: every layer is an independently runnable program over the quoted sources --- rebuild the module and re-run the checkers to reproduce the verdict. Receipts attesting superseded builds credit no layer and are omitted from this rendering (the count of omissions, if any, is noted below).}
\begin{receipttable}
\receiptrow{1417cdf27f69842a}{2026-07-28T23:48:35Z}{build+decorrelated}{\badge{BadgeGreen}{current}}
\receiptrow{65fc1b51055d38bd}{2026-07-29T00:00:33Z}{build+decorrelated}{\badge{BadgeGreen}{current}}
\receiptrow{2abbbc61b6065811}{2026-07-29T00:49:05Z}{build+replay+hardened+decorrelated}{\badge{BadgeGreen}{current}}
\receiptrow{21f04961748c0302}{2026-07-29T01:03:33Z}{build+replay+hardened+decorrelated}{\badge{BadgeGreen}{current}}
\receiptrow{8879c5b9721acf0c}{2026-07-31T22:00:03Z}{build+decorrelated}{\badge{BadgeGreen}{current}}
\receiptrow{e47a2c2d7f152711}{2026-07-31T22:28:57Z}{build+decorrelated}{\badge{BadgeGreen}{current}}
\receiptrow{ce1bb97f7cbcaf4f}{2026-08-02T23:31:07Z}{build+decorrelated}{\badge{BadgeGreen}{current}}
\receiptrow{8755b28f77990ba0}{2026-08-03T16:27:11Z}{build+decorrelated}{\badge{BadgeGreen}{current}}
\receiptrow{64948e60a56ae3bc}{2026-08-03T22:45:55Z}{build+decorrelated}{\badge{BadgeGreen}{current}}
\end{receipttable}
\sourcenote{8 further receipt(s) in the store attest superseded builds of this module and are omitted from this rendering.}
\layerline{\layeritem{build}{lake build (elaboration)}{\badge{BadgeGreen}{current}}\layeritem{replay}{leanchecker (official kernel replay)}{\badge{BadgeGreen}{current}}\layeritem{hardened}{safe\_verify (hardened screen)}{\badge{BadgeGreen}{current}}\layeritem{decorrelated}{lean4export + nanoda (decorrelated kernel)}{\badge{BadgeGreen}{current}}}
\sourcenote{\badge{BadgeGray}{comparator: none} no fixture in the three-implementation comparator harness (compiled / kernel-interpreted / reference-interpreted) exercises this rule.}
\ruleblock{HOL4 statement}
\absentbox{OUT OF SCOPE BY DESIGN}{Out of scope for the HOL side \emph{by design}: the dualities are
producer-side formula rewrites that unpack to definitional denotation
equalities before any certificate is emitted --- no wire object carries
a negation node (the typing discipline rejects surviving negative
literals), so there is nothing for a wire checker, HOL or otherwise, to
verify. Checked against the HOL sources: the \texttt{bern\_catalogue}
theory is the base-calculus checker cone; formula-level negation of the
extension theory\textquotesingle s spec surface does not appear in it.
}
\ruleblock{HOL verifier receipts}
\absentbox{NONE}{no HOL verification receipt exists for this rule.}
\ruleblock{Proof}

\sourcenote{Genre: axioms at the spec surface (roster tier 3). The seven
dualities are deduction rules with nothing above the line, each
concluding a definitional region-denotation equality --- there is no
single target rule being respelled, and no premise to carry: the
equality is the entire content. The payload-transport honesty clause of
the batch is Recipe~\ref{rec:sd-template} of the
\rulename{QStreettDegenerate} section, and applies here vacuously (no
payload ever exists to transport --- the dualities are applied before
emission).}

\begin{brecipe}[the seven axioms]
\label{rec:nnf-filled}
Over a state space $S$, with $[\![\cdot]\!]$ the region denotation of
spec-surface formulas and $f : S \to \R$ ranging over atom left-hand
sides:
\begin{mathpar}
\inferrule*{ }{ [\![\neg\neg\varphi]\!] = [\![\varphi]\!] }
\and
\inferrule*{ }{ [\![\neg(\varphi \wedge \psi)]\!] = [\![\neg\varphi \vee \neg\psi]\!] }
\and
\inferrule*{ }{ [\![\neg(\varphi \vee \psi)]\!] = [\![\neg\varphi \wedge \neg\psi]\!] }
\and
\inferrule*{ }{ [\![\neg\{f \le c\}]\!] = \{f > c\} }
\and
\inferrule*{ }{ [\![\neg\{f < c\}]\!] = \{f \ge c\} }
\and
\inferrule*{ }{ [\![\neg\{f \ge c\}]\!] = \{f < c\} }
\and
\inferrule*{ }{ [\![\neg\{f > c\}]\!] = \{f \le c\} }
\end{mathpar}
Each axiom licenses a judgment-level $1 \to 1$ rewrite in either
direction: a satisfaction judgment whose formula contains one side
concludes identically with the other side substituted, the certified
region unchanged.
\end{brecipe}

\begin{brecipe}[payload transport]\label{rec:nnf-transport}
\begin{center}\small
\begin{tabular}{@{}ll@{}}
\textbf{rewritten node} & \textbf{image}\\
\midrule
formula (one duality applied at one position) & formula (the dual form at that position)\\
certified region, judgment kind, all rows & verbatim\\
\end{tabular}
\end{center}
Unpack-function slot: no verified payload function in this edition
(batch template) ---
and here the slot is moot by design: the typing
discipline rejects surviving negative literals, so the dualities exist
to be applied \emph{before} any payload exists.
\end{brecipe}

\begin{bprop}[soundness of the axioms --- two lines each]
\label{prop:nnf-exact}
All seven concluded equalities hold.
\end{bprop}
\begin{proof}
The propositional three are complement involution and De Morgan for
sets. The atom four: pointwise over $S$, the negation of $f(y) \le c$ is
$c < f(y)$ and the negation of $f(y) < c$ is $c \le f(y)$ --- the total
order on $\R$ --- and the remaining two are the same facts read in the
other orientation. (Machine-checked: the four atom swaps as set-level
equalities, and the propositional three consumed from the ambient set
library inside the normalization pass's denotation theorem; hash-bound
in this edition.)
\end{proof}

\begin{bprop}[soundness corollary]\label{prop:nnf-sound}
If $R \subseteq [\![\varphi]\!]$ and $\varphi'$ arises from $\varphi$ by
one duality, then $R \subseteq [\![\varphi']\!]$; hence every
satisfaction judgment transports across a duality unchanged, in both
directions.
\end{bprop}
\begin{proof}
Denotations are equal (Proposition~\ref{prop:nnf-exact}), and
substitution preserves that: region denotation is compositional ---
each connective's denotation is a fixed set operation on its
children's --- so by induction on the surrounding formula, replacing a
subformula by one of equal denotation leaves every enclosing
denotation unchanged. The machine-checked cousin of this one-position
congruence is the normalization pass's denotation-preservation
theorem, which pushes negation through \emph{all} eligible positions
in one traversal. The pass consumes the propositional facts --- the
complement of a conjunction or disjunction, once per constructor
case --- and treats atoms opaquely: a negatively-occurring atom
survives as a negated literal whose denotation is the complement, and
the four atom dualities are separate bound facts that swap that
complement for its closed-form partner. That division realizes what
the display means by ``stated once per duality'': each equality is
one definitional fact, cited where it is used, never re-derived.
\end{proof}

\begin{bremark}[what this rule deliberately does not say]
The strict-atom swaps are \emph{denotation} facts; whether a strict
inequality may appear in a concluded region atom is a different
question, answered by the closure discipline (polarity-directed
rounding) in its own catalogue rows. And the probability operator is
absent from the display on purpose: a negated probability bound at even
polarity is not rewritten but \emph{rejected} by the typing pass --- the
polarity theorem of the \rulename{QStreettRequestNnf} section is what
makes that rejection a surface-level check.
\end{bremark}

\begin{anchors}
atom swaps, both orientations & \texttt{compl\_le\_atom}, \texttt{compl\_lt\_atom}, \texttt{compl\_ge\_atom}, \texttt{compl\_gt\_atom} (\texttt{BernsteinLogicAux.NNFDuality}; machine-checked, hash-bound)\\
Proposition \ref{prop:nnf-sound} induction & \texttt{denot\_nnfAux} (same module, hash-bound; the whole-pass form --- the one-position congruence is the compositionality induction in the proof)\\
no wire object, by design & negation nodes are rejected by typing; producers rewrite before emission\\
\end{anchors}


\rulesection{54}{QStreettRequestNnf}{alias for \textcolor{DiffBlue}{\rulename{QStreettBattery}}}
\label{blcatrule-QStreettRequestNnf}
\ruleblock{The rule (Steele metanotation)}
\sourcenote{as realized in the mechanized development; source row: \texttt{\texttt{QTSL-rules.toml}} --- \emph{QStreettRequestNnf}. The metanotation statement is typeset by hand for notational quality; its mechanized ground truth is the Lean statement below, extracted from the sources at every build of this book and hash-bound. The rule \emph{roster} is machine-checked (every mechanized rule tag round-trips to exactly one catalogue row); the statement \emph{text} is editorial, checked by eye against the extracted Lean beside it.}
\begin{mathpar}
\inferrule*{
  \mathcal{M},\ x_0 \in R \vDash \mathsf{SProbGE}\bigl(p,\ \mathsf{PDwellStreett}(\ldots,\ \mathsf{PairCmpl}(\mathrm{elimNeg}(\mathrm{nnf}(\rho_i)),\ \varphi_i),\ \ldots)\bigr)
}{
  \mathcal{M},\ x_0 \in R \vDash \mathsf{SProbGE}\bigl(p,\ \mathsf{PDwellStreett}(\ldots,\ \mathsf{PairCmpl}(\rho_i,\ \varphi_i),\ \ldots)\bigr)
}
\and
\inferrule*{ }{ [\![\mathrm{elimNeg}(\mathrm{nnf}(\rho))]\!] = [\![\rho]\!] }
\end{mathpar}

\ruleflavor{per pair $(\rho_i, \varphi_i)$, applied at each $\mathrm{cmpl}$-typed pair while base-typed requests pass unchanged: the request \emph{body} is rewritten to literal normal form by the QNnfRewrite dualities --- negation-freedom of the output holds \emph{by construction} of the normal-form grammar (no negation node exists to build) --- with the $\mathrm{cmpl}$ tag RETAINED: the complement applies at the consumption boundary (QStreettCmplRequest), never as a literal. The displayed axiom is the warrant: complements of equal denotations are equal, so the retained tag consumes identically. Polarity theorem: a surviving negative literal is exactly an odd-polarity atom occurrence of the input --- so the neg-P rejection is a surface polarity check, stable under normalization.}
\par\medskip
\rulename{QStreettRequestNnf} normalizes each Streett pair into the
bundle\textquotesingle s expected input shape (n\'e StreettRequestNNF): base-typed requests pass through unchanged, and a
complement-typed request has its formula rewritten to negation normal
form --- negations pushed to the atoms by the \rulename{QNnfRewrite}
rewrites --- so that the complement is applied at the \emph{consumption
boundary} (the atom swaps, or the complement-request rule), never
carried as a formula literal. It is an alias: the rewrite is a $1 \to 1$
respelling of one pair inside one \rulename{QStreettBattery} node, and
its warrant is denotation preservation of the normalization pass ---
the region the pair constrains is unchanged, so the judgment transports
verbatim.

The load-bearing extra is the \emph{polarity theorem}: the surviving
negative literals of the normal form are exactly the odd-polarity atom
occurrences of the input formula. Its consequence is that the typing
pass\textquotesingle s rejection of negated probability literals ---
which guards against smuggling a probability \emph{upper} bound through
an even number of negations --- is equivalent to a surface-level
polarity check on the un-normalized formula: normalization can never
manufacture a rejected literal that was not already there at odd
polarity, and never loses one. What would break without it: a
normalization pass without the polarity invariant could present a
well-typed surface whose normal form hides a rejected literal, and the
consumption boundary would silently consume a bound in the wrong
direction.

Mechanization state, plainly: the pass and both theorems are
machine-checked structurally, and the bound statements below are their
statements (module \texttt{BernsteinLogicAux.NNFDuality}): the
normalizer \texttt{Form.nnf} lands in a literal normal form whose
\emph{type} has no negation node --- negation-freedom by construction,
the strongest form of the well-formedness claim --- with denotation
preservation \texttt{denot\_nnf} (the displayed
$[\![\mathrm{nnf}(\rho_i)]\!] = [\![\rho_i]\!]$), the polarity
theorem \texttt{negLits\_nnf} (surviving negative literals $=$
odd-polarity atom occurrences) and its rejection face
\texttt{nnf\_no\_negLit\_iff}, and the consumption-boundary
elimination \texttt{denot\_elimNeg}. The rule\textquotesingle s scope
is the propositional skeleton, deliberately --- the path-operator nodes
of the spec surface are handled by their own ruled aliases and typing
rejections, never by this pass --- and the witnesses cover that scope in
full, so the entry carries no conditional flag.

\cataloguenote{The pass is a reference implementation on the producer side; nothing
checker-facing depends on running it. What the certificate format sees
is only the normalized result, and the typing pass independently
rejects any surviving negative probability literal --- so a defective
producer normalization can produce rejections, never unsoundness.
}
\ruleblock{Lean statement}
\declhead{BernsteinLogicAux.NNFDuality}{denot\_nnf}{\badge{BadgeGreen}{statement bound}}
\begin{leansrc}
/-- Denotation preservation of `Form.nnf` (T22's soundness face). -/
theorem denot_nnf (φ : Form α) : (φ.nnf).denot den = φ.denot den
\end{leansrc}
\declhead{BernsteinLogicAux.NNFDuality}{negLits\_nnf}{\badge{BadgeGreen}{statement bound}}
\begin{leansrc}
/-- **The T22 polarity theorem** ("a surviving neg-P literal is ALWAYS the
rejected even-polarity form", catalogue T22 — the wf theorem): the surviving
negative literals of the normal form are exactly the **odd-polarity** atom
occurrences of the input. Consequently the typing pass's neg-literal rejection
(`qtsl_typing/v1`, catalogue R6) is equivalent to a surface-level polarity
check on the un-normalized formula. -/
theorem negLits_nnf (φ : Form α) : (φ.nnf).negLits = φ.occs false
\end{leansrc}
\declhead{BernsteinLogicAux.NNFDuality}{nnf\_no\_negLit\_iff}{\badge{BadgeGreen}{statement bound}}
\begin{leansrc}
/-- The rejection face of the polarity theorem: the normal form is free of
`P`-flagged negative literals iff no `P`-atom occurs at odd polarity. -/
theorem nnf_no_negLit_iff (φ : Form α) (P : α → Prop) :
    (∀ a ∈ (φ.nnf).negLits, ¬ P a) ↔ ∀ a ∈ φ.occs false, ¬ P a
\end{leansrc}
\declhead{BernsteinLogicAux.NNFDuality}{denot\_elimNeg}{\badge{BadgeGreen}{statement bound}}
\begin{leansrc}
/-- **Denotation preservation of dual-atom elimination**: where each dual atom
denotes the complement (`hd` — the per-duality definitional axioms, e.g. the
`compl_le_atom` family above), elimination preserves the region. This is the
lawfulness of consuming a `cmpl`-typed request through swapped atoms
(atom-level complement elimination's atom step, the T9b input shape). -/
theorem denot_elimNeg {σ : Type*} (den : α → Set σ) (d : α → α)
    (hd : ∀ a, den (d a) = (den a)ᶜ) :
    ∀ φ : NForm α, (φ.elimNeg d).denot den = φ.denot den
\end{leansrc}
\ruleblock{Correspondence (Steele display $\leftrightarrow$ Lean)}
\badge{BadgeGreen}{reviewer-certified against the displayed texts}
\sourcenote{The dictionary below is a reviewer-certified statement of how the mechanized statements above interpret the displayed rule: display metavariable $\mapsto$ Lean binder, each display hypothesis $\mapsto$ the premise or side condition carrying it, the display conclusion $\mapsto$ the goal --- with any deliberate difference in width between display and mechanization named outright. Certification is recorded as a hash triple over (the displayed cell, the extracted statement texts, this dictionary\textquotesingle s own text); if any of the three changes, this block turns amber until a reviewer re-certifies. The correspondence claim itself remains editorial --- reviewer-grade, not machine-checked; its machine-checked successor types each soundness statement against the rule as a structured object, and supersedes this block per rule as it lands.}
The flavor text\textquotesingle s ``negation-freedom holds by
construction of the normal-form grammar'' is certified \emph{by
type}: the target grammar \texttt{NForm} (positive/negative literals
under $\wedge$/$\vee$) has no negation constructor, so no theorem
about negation-freedom exists to state --- the inductive definition
is the certificate. The displayed constituents --- the judgment-transport rule and its
warranting axiom --- map onto the
four bound statements:
\begin{itemize}
\item $[\![\mathrm{elimNeg}(\mathrm{nnf}(\rho_i))]\!] =
  [\![\rho_i]\!]$ $\mapsto$ the composition of \texttt{denot\_nnf}
  (\texttt{(φ.nnf).denot den = φ.denot den} --- the normalizer
  preserves the region, QNnfRewrite\textquotesingle s dualities consumed as
  the induction) with \texttt{denot\_elimNeg}, whose hypothesis
  \texttt{hd : ∀ a, den (d a) = (den a)ᶜ} is the per-duality
  definitional-axiom pack: elimination is lawful exactly where a
  dual atom exists, which is the flavor text\textquotesingle s
  ``the complement applies at the consumption boundary, never as a
  literal'' --- the retained $\mathrm{cmpl}$ tag consumes at QStreettCmplRequest
  through this same lemma, and dropping the tag would change the
  pair\textquotesingle s truth value (a one-state counterexample
  witnesses it), which is why the rewrite retains it.
\item the polarity theorem $\mapsto$ \texttt{negLits\_nnf :
  (φ.nnf).negLits = φ.occs false} --- surviving negative literals
  are \emph{exactly} the odd-polarity atom occurrences --- with the
  rejection face \texttt{nnf\_no\_negLit\_iff} certifying the
  flavor text\textquotesingle s ``the neg-P rejection is a surface
  polarity check, stable under normalization.''
\end{itemize}
Width note: atoms and denotations are fully abstract (\texttt{α},
\texttt{den : α → Set σ}) --- the per-pair Streett context, the
base-typed-request pass-through, and the probability flavour of the
neg-P check are all consumers of these statements, not content of
them; unconditional, \texttt{propext} alone, set-level witnesses
covering the catalogued claim in full.

\ruleblock{Lean verifier receipts}
\sourcenote{module \texttt{BernsteinLogicAux.NNFDuality}; bound build \texttt{97564f8e3950\ldots} (bound from receipt 21f04961748c0302, 2026-07-29T01:03:33Z). A receipt records that named checking layers accepted a build: (module, compiled-artifact SHA-256, environment SHA-256, checker layers), minted when the checking gate ran; the store receipts attesting the bound build are listed with the layers each run credited. Nothing here asks the reader to trust the store: every layer is an independently runnable program over the quoted sources --- rebuild the module and re-run the checkers to reproduce the verdict. Receipts attesting superseded builds credit no layer and are omitted from this rendering (the count of omissions, if any, is noted below).}
\begin{receipttable}
\receiptrow{1417cdf27f69842a}{2026-07-28T23:48:35Z}{build+decorrelated}{\badge{BadgeGreen}{current}}
\receiptrow{65fc1b51055d38bd}{2026-07-29T00:00:33Z}{build+decorrelated}{\badge{BadgeGreen}{current}}
\receiptrow{2abbbc61b6065811}{2026-07-29T00:49:05Z}{build+replay+hardened+decorrelated}{\badge{BadgeGreen}{current}}
\receiptrow{21f04961748c0302}{2026-07-29T01:03:33Z}{build+replay+hardened+decorrelated}{\badge{BadgeGreen}{current}}
\receiptrow{8879c5b9721acf0c}{2026-07-31T22:00:03Z}{build+decorrelated}{\badge{BadgeGreen}{current}}
\receiptrow{e47a2c2d7f152711}{2026-07-31T22:28:57Z}{build+decorrelated}{\badge{BadgeGreen}{current}}
\receiptrow{ce1bb97f7cbcaf4f}{2026-08-02T23:31:07Z}{build+decorrelated}{\badge{BadgeGreen}{current}}
\receiptrow{8755b28f77990ba0}{2026-08-03T16:27:11Z}{build+decorrelated}{\badge{BadgeGreen}{current}}
\receiptrow{64948e60a56ae3bc}{2026-08-03T22:45:55Z}{build+decorrelated}{\badge{BadgeGreen}{current}}
\end{receipttable}
\sourcenote{8 further receipt(s) in the store attest superseded builds of this module and are omitted from this rendering.}
\layerline{\layeritem{build}{lake build (elaboration)}{\badge{BadgeGreen}{current}}\layeritem{replay}{leanchecker (official kernel replay)}{\badge{BadgeGreen}{current}}\layeritem{hardened}{safe\_verify (hardened screen)}{\badge{BadgeGreen}{current}}\layeritem{decorrelated}{lean4export + nanoda (decorrelated kernel)}{\badge{BadgeGreen}{current}}}
\sourcenote{\badge{BadgeGray}{comparator: none} no fixture in the three-implementation comparator harness (compiled / kernel-interpreted / reference-interpreted) exercises this rule.}
\ruleblock{HOL4 statement}
\absentbox{OUT OF SCOPE BY DESIGN}{Out of scope for the HOL side \emph{by design}: the normalization is a
producer-side rewrite completed before emission --- certificates carry
only its output, and the output\textquotesingle s well-formedness is
enforced by the typing pass, not by a checker replay of the rewrite. No
wire object exists for a HOL statement to be about. Checked against the
HOL sources: the \texttt{bern\_catalogue} theory carries no spec-surface
formula grammar.
}
\ruleblock{HOL verifier receipts}
\absentbox{NONE}{no HOL verification receipt exists for this rule.}
\ruleblock{Proof}

\sourcenote{Genre: alias rule (tier 3) --- a respelling template; the
batch's shared template is Recipe~\ref{rec:sd-template} of the
\rulename{QStreettDegenerate}
section.}

\begin{brecipe}[the filled-in application --- one pair, normalized]
\label{rec:rnnf-filled}
Fix a Streett pair $(\rho, \varphi)$ inside a \rulename{QStreettBattery}
node, with $\rho$ complement-typed. Write $\mathrm{nnf}$ for the
negation-normal-form pass on the propositional skeleton (atoms,
$\neg$, $\wedge$, $\vee$): negations pushed inward by the
\rulename{QNnfRewrite} rewrites until only literals remain. The rewrite
is
\[
\bigl(\mathrm{cmpl}\,\rho,\ \varphi\bigr)
\quad\rightsquigarrow\quad
\bigl(\mathrm{cmpl}\,\bigl(\mathrm{elimNeg}(\mathrm{nnf}(\rho))\bigr),\ \varphi\bigr)
\]
--- the body normalized, the complement tag \emph{retained}: per the
source discipline the complement applies at the consumption boundary
(the complement-request rule), never as a literal. After the rewrite
every literal of the request body is positive (atom-swapped where it
was negative), and the pair is in exactly the input shape the
bundle and the complement-request rule expect. A base-typed pair is
already in that shape: on it the rewrite is the identity, the transport
trivial, and the exactness identity reflexivity --- stated once here so
the corollary\textquotesingle s case split below is complete. One pair
to one pair,
inside one node: the alias shape.
\end{brecipe}

\begin{brecipe}[payload transport]\label{rec:rnnf-transport}
\begin{center}\small
\begin{tabular}{@{}ll@{}}
\textbf{un-normalized pair} & \textbf{normalized pair}\\
\midrule
request body $\rho$ ($\mathrm{cmpl}$ tag retained) & $\mathrm{elimNeg}(\mathrm{nnf}(\rho))$, negative literals eliminated through dual atoms\\
response formula, certificates, rows & verbatim\\
polarity bookkeeping & the polarity theorem (Proposition~\ref{prop:rnnf-polarity}); nothing stored\\
\end{tabular}
\end{center}
Unpack-function slot: no verified payload function in this edition
(batch template); the
structural mechanization exists (the normalizer and its theorems, named
below; the four theorems hash-bound in this edition).
\end{brecipe}

\begin{bprop}[the exactness identity --- denotation preservation]
\label{prop:rnnf-exact}
$[\![\mathrm{nnf}(\rho)]\!] = [\![\rho]\!]$, and eliminating negative
literals through dual atoms preserves the denotation.
\end{bprop}
\begin{proof}
Induction on $\rho$, with the \rulename{QNnfRewrite} equalities as the
cases --- double negation and De Morgan drive the recursion, the atom
swaps close the leaves. The elimination step substitutes, for each
negative literal, the dual atom denoting the complement; equality is
preserved literal by literal. (Machine-checked: \texttt{denot\_nnf} and
\texttt{denot\_elimNeg}, hash-bound in this edition; the eliminated form
has no negative literals at
all, \texttt{negLits\_elimNeg} --- an unpinned auxiliary of the same
module.) Complements of equal
sets are equal, so the retained $\mathrm{cmpl}$ tag's
consumption-boundary denotation $S \smallsetminus [\![\rho]\!]$ is
preserved as well --- the rewrite itself never discharges the
complement.
\end{proof}

\begin{bprop}[the polarity theorem]\label{prop:rnnf-polarity}
The surviving negative literals of $\mathrm{nnf}(\rho)$ are exactly the
atom occurrences of $\rho$ at odd polarity (an odd number of enclosing
negations). Consequently the normal form is free of negated occurrences
of a given atom class if and only if the input has no odd-polarity
occurrence of that class.
\end{bprop}
\begin{proof}
Strengthen to the polarity-indexed pass (normalizing $\rho$ or
$\neg\rho$ according to a polarity bit) and induct: at an atom the bit
decides the literal's sign; at a negation the bit flips; at $\wedge$ and
$\vee$ the bit distributes (De Morgan swapping the connective when
negative). The negative-literal set on one side and the odd-polarity
occurrence set on the other satisfy the same recursion. The rejection
face is a rewrite of the same equality. (Machine-checked:
\texttt{negLits\_nnf}, via the unpinned auxiliary
\texttt{negLits\_nnfAux}, and
\texttt{nnf\_no\_negLit\_iff}; both theorems hash-bound in this edition.
Terminology, once: a \emph{negative literal} is a signed-literal
constructor of the normal-form grammar, not an occurrence of the
formula negation node --- the grammar carries literals of both signs
and no negation constructor, which is the sense in which
negation-freedom of the output holds by \emph{type}; the elimination
pass then removes the negative-signed literals through dual atoms.)
\end{proof}

\begin{bprop}[soundness corollary]\label{prop:rnnf-sound}
A bundle node whose pair is normalized concludes identically to the
un-normalized node: through \rulename{QStreettBattery} directly for
base-typed requests, and through the \rulename{QStreettCmplRequest}
weakening for the consumed complement --- with
Proposition~\ref{prop:rnnf-exact} guaranteeing both nodes constrain the
same regions.
\end{bprop}
\begin{proof}
Immediate from Proposition~\ref{prop:rnnf-exact} and the target rules'
theorems, each at its own proven scope --- the rewrite itself is
propositional and profile-free; whatever conditional mark the
concluding bundle or weakening node carries is inherited there, not
here; the alias adds no semantic content of its own.
\end{proof}

\begin{bremark}[what the polarity theorem is protecting]
The typing discipline rejects probability-atom occurrences at odd
polarity --- an odd number of enclosing negations, whether worn as a
bare negated literal in even context or as a positive literal under an
odd stack --- because a negated lower bound is an upper bound, and the
certificate calculus deliberately has no rule concluding probability
upper bounds through this pass. The theorem makes the rejection
\emph{stable under normalization}: no sequence of dualities can
manufacture or hide a rejected literal, so the check may run on the
surface formula once, before normalization, with the same verdict. A
guard that had to chase literals through rewrites would be a moving
target; a polarity count is not.
\end{bremark}

\begin{anchors}
normalizer, negation-freedom by type & \texttt{Form.nnf}, \texttt{NForm} (\texttt{BernsteinLogicAux.NNFDuality}; definitions --- no statement pins, their theorems below carry the warrant)\\
Proposition \ref{prop:rnnf-exact} & \texttt{denot\_nnf}, \texttt{denot\_elimNeg} (same module, hash-bound); \texttt{negLits\_elimNeg} (unpinned auxiliary)\\
Proposition \ref{prop:rnnf-polarity} & \texttt{negLits\_nnf}, \texttt{nnf\_no\_negLit\_iff} (same module, hash-bound)\\
Proposition \ref{prop:rnnf-sound} targets & the \rulename{QStreettBattery} and \rulename{QStreettCmplRequest} sections\\
\end{anchors}


\rulesection{55}{QAvoid}{alias rule}
\label{blcatrule-QAvoid}
\ruleblock{The rule (Steele metanotation)}
\sourcenote{as realized in the mechanized development; source row: \texttt{\texttt{QTSL-rules.toml}} --- \emph{QAvoid}. The metanotation statement is typeset by hand for notational quality; its mechanized ground truth is the Lean statement below, extracted from the sources at every build of this book and hash-bound. The rule \emph{roster} is machine-checked (every mechanized rule tag round-trips to exactly one catalogue row); the statement \emph{text} is editorial, checked by eye against the extracted Lean beside it.}
\begin{mathpar}
\inferrule*{
  \mathcal{M},\ x_0 \in R \vDash \mathsf{SProbGE}\bigl(p,\ \mathsf{PAlways}(\mathsf{SNot}\,\beta)\bigr)
}{
  \mathcal{M},\ x_0 \in R \vDash \mathsf{P}_{\ge p}[\mathsf{avoid}\,\beta]
}
\end{mathpar}

\ruleflavor{$\mathsf{avoid}\,\beta := \mathsf{G}\neg\beta$ --- the unfolding is the rule (the premise is the target judgment at the expansion), so the alias is sound by definition: the equation is definitional at the set level (\texttt{rfl}-grade), and the negation atom is ruled in by the \rulename{QNnfRewrite} swaps. The surface constructor builds its expansion rather than a new node kind — witnessed by the producer-side concrete-syntax canonicality battery; no rule sort exists and none is needed.}
\par\medskip
\rulename{QAvoid} is the avoidance sugar of the spec surface:
$\mathsf{avoid}\,\beta := \mathsf{G}\neg\beta$ --- avoidance is safety
of the negation, exact, the negation being ruled in by the
\rulename{QNnfRewrite} atom swaps. An alias unpacking to \emph{one}
application of an existing rule (the safety battery at the avoid set),
and the strongest-fidelity member of the sugar trio (its siblings:
\rulename{QStreettGF}, \rulename{QStreettFG}): the identity the
definition rides is machine-checked as a path-event identity in the
target rule\textquotesingle s concluded shape (module
\texttt{BernsteinLogicAux.QtslAliases}), \texttt{sugarAvoid\_denot},
proof \texttt{rfl}, axiom-free --- definitional exactness, certified as
such by the kernel. A producer, and a reader of specifications, may use
the sugar constructor freely: it builds its expansion rather than a new
node kind, so nothing downstream --- typing, compilation, checking ---
ever sees a sugar node, and no checker recomputation exists for this
rule at all. There is no soundness lemma here because the equation is a
definition; the witness that the surface really does build the
expansion remains the reference implementation\textquotesingle s
concrete-syntax canonicality test battery --- an implementation
property, not a semantic claim, so it stays a test rather than a
theorem. No profile enters; the entry carries no conditional flag.

\cataloguenote{The avoid sugar is also the compatibility layer to plain safety
specification: avoidance properties state directly against the safety
kernel with no automaton machinery, which is the degenerate (and most
common) case of the classical-surface translation the trio provides.
}
\ruleblock{Lean statement}
\declhead{BernsteinLogicAux.QtslAliases}{sugarAvoid\_denot}{\badge{BadgeGreen}{statement bound}}
\begin{leansrc}
/-- **`SugarBatch` (T10), first sugar**: `avoid β := G ¬β` — avoidance *is*
safety of the negation, definitionally (the proof is `rfl`; the negation is
ruled in by the `NNFDuality.lean` atom swaps (rule QNnfRewrite)). Exact: no direction is lost. -/
theorem sugarAvoid_denot {S : Type*} (B : Set S) :
    {ω : ℕ → S | ∀ n, ω n ∉ B} = {ω : ℕ → S | ∀ n, ω n ∈ Bᶜ}
\end{leansrc}
\ruleblock{Correspondence (Steele display $\leftrightarrow$ Lean)}
\badge{BadgeGreen}{reviewer-certified against the displayed texts}
\sourcenote{The dictionary below is a reviewer-certified statement of how the mechanized statements above interpret the displayed rule: display metavariable $\mapsto$ Lean binder, each display hypothesis $\mapsto$ the premise or side condition carrying it, the display conclusion $\mapsto$ the goal --- with any deliberate difference in width between display and mechanization named outright. Certification is recorded as a hash triple over (the displayed cell, the extracted statement texts, this dictionary\textquotesingle s own text); if any of the three changes, this block turns amber until a reviewer re-certifies. The correspondence claim itself remains editorial --- reviewer-grade, not machine-checked; its machine-checked successor types each soundness statement against the rule as a structured object, and supersedes this block per rule as it lands.}
The displayed rule transports the target judgment across the
defining equation $\mathrm{avoid}\,\beta := \mathsf{G}\neg\beta$
(premise: the expansion\textquotesingle s judgment; conclusion: the
sugared form). \texttt{sugarAvoid\_denot} certifies that equation and
is \texttt{rfl} --- definitional at the set level, the
strongest possible fidelity, with the negation-atom step ruled in by
the \rulename{QNnfRewrite} swaps as the docstring notes. Width note: a
per-path equivalence over arbitrary state types --- no measure enters,
so the identity transports under every profile\textquotesingle s path
measure by congruence; the remaining accompanying claim (the entry\textquotesingle s flavor
text: the certificate
surface \emph{builds} the expansion rather than minting a node kind) is
a property of the concrete-syntax canonicality battery --- a test, not
a theorem, deliberately outside this bound set, exactly as the roster
witness says.

\ruleblock{Lean verifier receipts}
\sourcenote{module \texttt{BernsteinLogicAux.QtslAliases}; bound build \texttt{570599779c29\ldots} (bound from receipt 5a60f4d01fd81ac5, 2026-07-21T05:39:02Z). A receipt records that named checking layers accepted a build: (module, compiled-artifact SHA-256, environment SHA-256, checker layers), minted when the checking gate ran; the store receipts attesting the bound build are listed with the layers each run credited. Nothing here asks the reader to trust the store: every layer is an independently runnable program over the quoted sources --- rebuild the module and re-run the checkers to reproduce the verdict. Receipts attesting superseded builds credit no layer and are omitted from this rendering (the count of omissions, if any, is noted below).}
\begin{receipttable}
\receiptrow{0a3819a978e31127}{2026-07-21T04:51:36Z}{build+decorrelated}{\badge{BadgeGreen}{current}}
\receiptrow{0bb1074df705d863}{2026-07-21T05:21:58Z}{build+decorrelated}{\badge{BadgeGreen}{current}}
\receiptrow{5a60f4d01fd81ac5}{2026-07-21T05:39:02Z}{build+replay+hardened+decorrelated}{\badge{BadgeGreen}{current}}
\end{receipttable}
\sourcenote{11 further receipt(s) in the store attest superseded builds of this module and are omitted from this rendering.}
\layerline{\layeritem{build}{lake build (elaboration)}{\badge{BadgeGreen}{current}}\layeritem{replay}{leanchecker (official kernel replay)}{\badge{BadgeGreen}{current}}\layeritem{hardened}{safe\_verify (hardened screen)}{\badge{BadgeGreen}{current}}\layeritem{decorrelated}{lean4export + nanoda (decorrelated kernel)}{\badge{BadgeGreen}{current}}}
\sourcenote{\badge{BadgeGray}{comparator: none} no fixture in the three-implementation comparator harness (compiled / kernel-interpreted / reference-interpreted) exercises this rule.}
\ruleblock{HOL4 statement}
\absentbox{OUT OF SCOPE BY DESIGN}{Out of scope for the HOL side \emph{by design}, in the strongest form
an alias can have: the sugar is a definitional constructor that builds
its expansion at the surface --- no sugar node exists in any
representation downstream of parsing, wire format included, so there is
no artefact for any checker port to be about. Checked against the HOL
sources: the \texttt{bern\_catalogue} theory carries no spec-surface
grammar.
}
\ruleblock{HOL verifier receipts}
\absentbox{NONE}{no HOL verification receipt exists for this rule.}
\ruleblock{Proof}

\sourcenote{Genre: alias rule (tier 3) --- a respelling template; the
shared template is Recipe~\ref{rec:sd-template} of the
\rulename{QStreettDegenerate} section. This rule is the degenerate case
of even that template: the unpack is a \emph{definition}, so the five
artefacts collapse to the definitional equation, its reading, and the
target section.}

\begin{brecipe}[the filled-in application]\label{rec:qaa-filled}
\[
\mathsf{avoid}\,\beta \;:=\; \mathsf{G}\neg\beta
\]
The concluding judgment discharges through one existing rule:
avoidance through the safety battery (\rulename{QSafetyBattery}) at the
avoid set $[\![\beta]\!]$, with $\neg\beta$ read by the semantic
negation clause $[\![\neg\beta]\!] = S \smallsetminus [\![\beta]\!]$
(the \rulename{QNnfRewrite} atom dualities supply the syntactic normal
form --- the atom swap alone when $\beta$ is an atom, De Morgan
propagation when it is compound).
\end{brecipe}

\begin{brecipe}[payload transport]\label{rec:qaa-transport}
The expansion $\mathsf{G}\neg\beta$ carries the safety kit at
$U := [\![\beta]\!]$, verbatim. No unpack function is owed here,
because none will ever exist: the constructor \emph{is} its expansion,
so there is no sugar payload to transport from --- the strongest
possible form of producers-elaborate-before-emission.
\end{brecipe}

\begin{bprop}[the exactness identity]\label{prop:qaa-exact}
A path always avoids $[\![\beta]\!]$ iff it always satisfies
$\neg\beta$ --- definitional, by the negation clause
$[\![\neg\beta]\!] = S \smallsetminus [\![\beta]\!]$, with
\rulename{QNnfRewrite} normalizing the syntax where needed.
\end{bprop}
\begin{proof}
Definitional (the bound Lean identity is \texttt{rfl}, axiom-free ---
the kernel certifies the two sides as one term).
\end{proof}

\begin{bprop}[soundness corollary]\label{prop:qaa-sound}
The sugared judgment is concluded by the target rule's theorem, with
no additional obligation: the sugar constructor builds the expansion,
so the sugared judgment \emph{is} the target judgment.
\end{bprop}
\begin{proof}
Definitional unfolding plus the target section's theorem. The witness
that the surface implementation honours the definition --- that the
constructor builds its expansion rather than introducing a new node
kind --- is the reference implementation's concrete-syntax canonicality
test battery; no checker-facing artefact depends on it, since no
downstream representation contains a sugar node.
\end{proof}

\begin{anchors}
target & the \rulename{QSafetyBattery} section; negation via the \rulename{QNnfRewrite} atom swaps\\
the identity, event-level & \texttt{sugarAvoid\_denot} (\texttt{rfl}, axiom-free; bound in the entry)\\
siblings (the sugar trio) & \rulename{QStreettGF}, \rulename{QStreettFG}\\
constructor builds expansion & the reference implementation\textquotesingle s concrete-syntax canonicality battery (documentary; no wire artefact exists to check)\\
\end{anchors}


\rulesection{56}{QStreettGF}{alias rule}
\label{blcatrule-QStreettGF}
\ruleblock{The rule (Steele metanotation)}
\sourcenote{as realized in the mechanized development; source row: \texttt{\texttt{QTSL-rules.toml}} --- \emph{QStreettGF}. The metanotation statement is typeset by hand for notational quality; its mechanized ground truth is the Lean statement below, extracted from the sources at every build of this book and hash-bound. The rule \emph{roster} is machine-checked (every mechanized rule tag round-trips to exactly one catalogue row); the statement \emph{text} is editorial, checked by eye against the extracted Lean beside it.}
\begin{mathpar}
\inferrule*{
  \mathcal{M},\ x_0 \in R \vDash \mathsf{SProbGE}\bigl(p,\ \mathsf{PDwellStreett}(\mathsf{PairBase}(\mathsf{STrue}, \varphi))\bigr) \\
  \text{paths total over the state space}
}{
  \mathcal{M},\ x_0 \in R \vDash \mathsf{SProbGE}(p,\ \mathsf{GF}\,\varphi)
}
\end{mathpar}

\ruleflavor{$\mathsf{GF}\varphi := \mathsf{Streett}((\top, \varphi))$ --- recurrence is the one-pair Streett condition at the trivially-true request, \rulename{QStreettDegenerate}\textquotesingle s concluded shape, so the sugar lands on that rule\textquotesingle s bound statement without adjustment. The totality premise is the bound chart, exactly the certified statement\textquotesingle s scope. Killed-process classes are the named unbound residue: the identity extends whenever the cemetery is adjoined inside the region vocabulary \emph{and} killed paths are continued as total paths at \(\dagger\) (\(\top = V \cup \{\dagger\}\)) --- that \(\dagger\)-extended restatement is not yet part of the certified semantics, so it stays here rather than in a premise; with the cemetery outside the region vocabulary the rewrite direction over-claims, and that class is kept closed by the batteries\textquotesingle{} full-mass gates. The surface constructor builds its expansion rather than a new node kind — witnessed by the producer-side concrete-syntax canonicality battery; no rule sort exists and none is needed.}
\par\medskip
\rulename{QStreettGF} is the recurrence sugar of the spec surface:
$\mathsf{GF}\varphi := \mathsf{Streett}((\top, \varphi))$ --- recurrence
is the degenerate Streett pair, the \rulename{QStreettDegenerate}
section\textquotesingle s concluded shape, so the sugar lands on that
rule\textquotesingle s bound statement without adjustment. An alias
unpacking to \emph{one} application of an existing rule; with its
persistence twin \rulename{QStreettFG} it is the compatibility layer to
the classical temporal-logic surface ---
$\mathsf{GF}$/$\mathsf{FG}$ specifications translate to the
Streett-primitive kernel without automaton products, the practical
payoff of B\"uchi acceptance being a one-pair Streett condition. The
identity the definition rides is machine-checked (module
\texttt{BernsteinLogicAux.QtslAliases}): \texttt{sugarGF\_denot} ---
the per-pair Streett event at the trivially-true request \emph{is} the
recurrence event.

Scope, stated in the open because this member alone has one: the
mechanized identity lives on the total path type, so it consumes the
no-cemetery chart --- exact on total-path state spaces and on
killed-process classes with the cemetery adjoined inside the region
vocabulary ($\top = V \cup \{\dagger\}$); with the cemetery outside the
region vocabulary the rewrite direction over-claims, and that class is
kept closed by the batteries\textquotesingle{} full-mass gates. The
named repair (the entry\textquotesingle s residue): restate the
identity on the $\dagger$-extended state type, matching the promotion
lemma. The surface constructor builds its expansion rather than a new
node kind; the witness for that remains the reference
implementation\textquotesingle s concrete-syntax canonicality test
battery --- an implementation property, a test rather than a theorem.

\cataloguenote{Why the factoring put the honesty here: as a batch row, one PROVEN
badge covered three warrants of different semantic depths ---
\texttt{rfl}-grade, chart-consuming, and pure-logic --- and the
chart-consuming member\textquotesingle s scope restriction was a
memo-level footnote. As its own row it is display-level text: the
exactness chart is part of what this rule \emph{says}.
}
\ruleblock{Lean statement}
\declhead{BernsteinLogicAux.QtslAliases}{sugarGF\_denot}{\badge{BadgeGreen}{statement bound}}
\begin{leansrc}
/-- **`SugarBatch` (T10), second sugar**: `GF φ := Streett((⊤, φ))` — the
per-pair Streett event at the trivially-true request (every state requests,
and along `atTop` over `ℕ` the request is frequent) *is* the recurrence
event. The expansion's target is `QStreettDegenerate` (T9a), whose bound
statement `sound_QStreettDegenerate` concludes exactly the right-hand
event. -/
theorem sugarGF_denot {S : Type*} (G : Set S) (ω : ℕ → S) :
    ((∃ᶠ n in atTop, ω n ∈ (Set.univ : Set S)) → ∃ᶠ n in atTop, ω n ∈ G)
      ↔ ∃ᶠ n in atTop, ω n ∈ G
\end{leansrc}
\ruleblock{Correspondence (Steele display $\leftrightarrow$ Lean)}
\badge{BadgeGreen}{reviewer-certified against the displayed texts}
\sourcenote{The dictionary below is a reviewer-certified statement of how the mechanized statements above interpret the displayed rule: display metavariable $\mapsto$ Lean binder, each display hypothesis $\mapsto$ the premise or side condition carrying it, the display conclusion $\mapsto$ the goal --- with any deliberate difference in width between display and mechanization named outright. Certification is recorded as a hash triple over (the displayed cell, the extracted statement texts, this dictionary\textquotesingle s own text); if any of the three changes, this block turns amber until a reviewer re-certifies. The correspondence claim itself remains editorial --- reviewer-grade, not machine-checked; its machine-checked successor types each soundness statement against the rule as a structured object, and supersedes this block per rule as it lands.}
\texttt{sugarGF\_denot} ($\mathsf{GF}\varphi :=
\mathsf{Streett}((\top, \varphi))$) certifies the pair event at
request \texttt{Set.univ} $\leftrightarrow$ the bare recurrence event
--- exactly \texttt{sound\_QStreettDegenerate}\textquotesingle s
concluded form, so the sugar lands on
\rulename{QStreettDegenerate}\textquotesingle s bound statement without
adjustment. Fidelity caveat, bound honestly: the Lean statement lives
on the total path type, which IS the displayed rule\textquotesingle s
totality premise (``paths total over the state space'') --- the premise
and the bound scope now coincide exactly. The flavor
text\textquotesingle s $\dagger$-extended case (cemetery adjoined
inside the vocabulary with killed paths continued as total paths at
$\dagger$: exact; cemetery outside: the rewrite over-claims) is the
named unbound residue the bound statement does not and
cannot witness; it is carried by the roster note\textquotesingle s
residue line, and the $\dagger$-extended restatement named there would
move it into this bound set. No measure enters; the identity transports
under every profile\textquotesingle s path measure by congruence. The
builds-expansions claim rides the canonicality battery, deliberately
outside this bound set, exactly as the roster witness says.

\ruleblock{Lean verifier receipts}
\sourcenote{module \texttt{BernsteinLogicAux.QtslAliases}; bound build \texttt{570599779c29\ldots} (bound from receipt 5a60f4d01fd81ac5, 2026-07-21T05:39:02Z). A receipt records that named checking layers accepted a build: (module, compiled-artifact SHA-256, environment SHA-256, checker layers), minted when the checking gate ran; the store receipts attesting the bound build are listed with the layers each run credited. Nothing here asks the reader to trust the store: every layer is an independently runnable program over the quoted sources --- rebuild the module and re-run the checkers to reproduce the verdict. Receipts attesting superseded builds credit no layer and are omitted from this rendering (the count of omissions, if any, is noted below).}
\begin{receipttable}
\receiptrow{0a3819a978e31127}{2026-07-21T04:51:36Z}{build+decorrelated}{\badge{BadgeGreen}{current}}
\receiptrow{0bb1074df705d863}{2026-07-21T05:21:58Z}{build+decorrelated}{\badge{BadgeGreen}{current}}
\receiptrow{5a60f4d01fd81ac5}{2026-07-21T05:39:02Z}{build+replay+hardened+decorrelated}{\badge{BadgeGreen}{current}}
\end{receipttable}
\sourcenote{11 further receipt(s) in the store attest superseded builds of this module and are omitted from this rendering.}
\layerline{\layeritem{build}{lake build (elaboration)}{\badge{BadgeGreen}{current}}\layeritem{replay}{leanchecker (official kernel replay)}{\badge{BadgeGreen}{current}}\layeritem{hardened}{safe\_verify (hardened screen)}{\badge{BadgeGreen}{current}}\layeritem{decorrelated}{lean4export + nanoda (decorrelated kernel)}{\badge{BadgeGreen}{current}}}
\sourcenote{\badge{BadgeGray}{comparator: none} no fixture in the three-implementation comparator harness (compiled / kernel-interpreted / reference-interpreted) exercises this rule.}
\ruleblock{HOL4 statement}
\absentbox{OUT OF SCOPE BY DESIGN}{Out of scope for the HOL side \emph{by design}: the sugar is a
definitional constructor that builds its expansion at the surface ---
no sugar node exists in any representation downstream of parsing, wire
format included, so there is no artefact for any checker port to be
about. Checked against the HOL sources: the \texttt{bern\_catalogue}
theory carries no spec-surface grammar. (The chart restriction is a
property of the Lean statement\textquotesingle s ambient path type, not
of any wire artefact --- there is nothing HOL-side for it to attach to.)
}
\ruleblock{HOL verifier receipts}
\absentbox{NONE}{no HOL verification receipt exists for this rule.}
\ruleblock{Proof}

\sourcenote{Genre: alias rule (tier 3) --- a respelling template; the
shared template is Recipe~\ref{rec:sd-template} of the
\rulename{QStreettDegenerate} section. The unpack is a \emph{definition},
so the five artefacts collapse to the definitional equation, its
reading, its exactness chart, and the target section.}

\begin{brecipe}[the filled-in application]\label{rec:qgf-filled}
\[
\mathsf{GF}\varphi \;:=\; \mathsf{Streett}((\top, \varphi))
\]
The concluding judgment discharges through one existing rule:
recurrence through the degenerate bundle
(\rulename{QStreettDegenerate}).
\end{brecipe}

\begin{brecipe}[payload transport]\label{rec:qgf-transport}
The expansion $\mathsf{Streett}((\top, \varphi))$ carries the
degenerate kit, verbatim. No unpack function is owed here, because none
will ever exist: the constructor \emph{is} its expansion, so there is
no sugar payload to transport from --- the strongest possible form of
producers-elaborate-before-emission.
\end{brecipe}

\begin{bprop}[the exactness identity, and its chart]\label{prop:qgf-exact}
The pair $(\top, \varphi)$ holds iff $\varphi$-states occur infinitely
often (the degeneration identity, proved as the exactness identity,
Proposition~\ref{prop:sd-exact}, of the \rulename{QStreettDegenerate}
section). As mechanized, the identity is stated on the total path type;
it therefore holds on total-path state spaces, and on killed-process
classes whenever the cemetery is adjoined inside the region vocabulary
($\top = V \cup \{\dagger\}$). With the cemetery outside the region
vocabulary the rewrite direction over-claims --- a killed path visits no
state cofinally, yet the vocabulary cannot say so --- and that class is
kept closed by the batteries' full-mass gates. The named repair is to
restate the identity on the $\dagger$-extended state type, matching the
promotion lemma; until then the chart restriction is carried in the
rule's own display, in the open.
\end{bprop}
\begin{proof}
The identity: by the degeneration argument of the target section,
mechanized as the bound Lean statement. The chart restriction: read off
the bound statement's ambient path type, as displayed.
\end{proof}

\begin{bprop}[soundness corollary]\label{prop:qgf-sound}
On the stated chart, the sugared judgment is concluded by the target
rule's theorem, with no additional obligation: the sugar constructor
builds the expansion, so the sugared judgment \emph{is} the target
judgment.
\end{bprop}
\begin{proof}
Definitional unfolding plus the target section's theorem, on the chart
of Proposition~\ref{prop:qgf-exact}. The witness that the surface
implementation honours the definition --- that the constructor builds
its expansion rather than introducing a new node kind --- is the
reference implementation's concrete-syntax canonicality test battery;
no checker-facing artefact depends on it, since no downstream
representation contains a sugar node.
\end{proof}

\begin{anchors}
target & the \rulename{QStreettDegenerate} section (degeneration identity proved there)\\
the identity, event-level & \texttt{sugarGF\_denot} (bound in the entry; total path type = the chart)\\
the chart's named repair & the $\dagger$-extended restatement, $\top = V \cup \{\dagger\}$ (the promotion lemma's form)\\
twin & \rulename{QStreettFG} (the safety-polarized half; consumes no chart)\\
constructor builds expansion & the reference implementation\textquotesingle s concrete-syntax canonicality battery (documentary; no wire artefact exists to check)\\
\end{anchors}


\rulesection{57}{QStreettFG}{alias rule}
\label{blcatrule-QStreettFG}
\ruleblock{The rule (Steele metanotation)}
\sourcenote{as realized in the mechanized development; source row: \texttt{\texttt{QTSL-rules.toml}} --- \emph{QStreettFG}. The metanotation statement is typeset by hand for notational quality; its mechanized ground truth is the Lean statement below, extracted from the sources at every build of this book and hash-bound. The rule \emph{roster} is machine-checked (every mechanized rule tag round-trips to exactly one catalogue row); the statement \emph{text} is editorial, checked by eye against the extracted Lean beside it.}
\begin{mathpar}
\inferrule*{
  \mathcal{M},\ x_0 \in R \vDash \mathsf{SProbGE}\bigl(p,\ \mathsf{PDwellStreett}(\mathsf{PairBase}(\beta, \mathsf{SFalse}))\bigr)
}{
  \mathcal{M},\ x_0 \in R \vDash \mathsf{SProbGE}(p,\ \mathsf{FG}\,\neg\beta)
}
\end{mathpar}

\ruleflavor{$\mathsf{FG}\neg\beta := \mathsf{Streett}((\beta, \bot))$ --- eventual permanent avoidance is the one-pair Streett condition that forbids infinitely many \(\beta\)-visits: with an unsatisfiable response, the pair condition says exactly that \(\beta\)-states occur only finitely often. Pure per-path logic — the safety-polarized half of the recurrence/persistence pair, consuming no conservativity. The rule\textquotesingle s honest scope is this polarized form (\(\mathsf{FG}\) applied to \(\neg\beta\)): general persistence \(\mathsf{FG}\varphi\) is not this rule. The surface constructor builds its expansion rather than a new node kind — witnessed by the producer-side concrete-syntax canonicality battery; no rule sort exists and none is needed.}
\par\medskip
\rulename{QStreettFG} is the persistence-of-avoidance sugar of the spec
surface: $\mathsf{FG}\neg\beta := \mathsf{Streett}((\beta, \bot))$ ---
eventual permanent avoidance is the one-pair Streett condition that
forbids infinitely many $\beta$-visits: with an unsatisfiable response,
the pair condition says exactly that $\beta$-states occur only finitely
often. An alias unpacking to \emph{one} application of an existing rule
(one \rulename{QStreettBattery} node at the pair $(\beta, \bot)$); the
safety-polarized half of the recurrence/persistence pair --- its twin
is \rulename{QStreettGF}, and the two land at different points of the
same Streett rule. The honest scope is the polarized form
($\mathsf{FG}$ applied to $\neg\beta$): general persistence
$\mathsf{FG}\varphi$ is not this rule. The identity the definition
rides is machine-checked (module
\texttt{BernsteinLogicAux.QtslAliases}): \texttt{sugarFG\_denot} ---
the per-pair Streett event at an unsatisfiable response \emph{is}
eventual permanent avoidance, the finitely-often reading; pure per-path
logic, no conservativity and no chart hypothesis consumed. The surface
constructor builds its expansion rather than a new node kind; the
witness for that remains the reference implementation\textquotesingle s
concrete-syntax canonicality test battery --- an implementation
property, a test rather than a theorem. No profile enters; the entry
carries no conditional flag.

\cataloguenote{The two dual sugars land at different points of the same rule:
$\mathsf{GF}\varphi$ requests always and so erases the request
machinery, while $\mathsf{FG}\neg\beta$ makes $\beta$ itself the
request and forbids response --- the ranking certificate must then pay
out at every $\beta$-visit with nothing ever refunding it, which is
exactly a proof that $\beta$-visits are finitely many. That both
directions are single nodes of one rule --- where automaton-primitive
designs pay a product construction for one of them --- is the practical
content of the Streett-primitive choice.
}
\ruleblock{Lean statement}
\declhead{BernsteinLogicAux.QtslAliases}{sugarFG\_denot}{\badge{BadgeGreen}{statement bound}}
\begin{leansrc}
/-- **`SugarBatch` (T10), third sugar**: `FG ¬β := Streett((β, ⊥))` — the
per-pair Streett event at an unsatisfiable response says the request region
is visited at most finitely often, which over `ℕ` *is* eventual permanent
avoidance. The pair condition forbids infinitely many `β`-visits: with
response `⊥`, "infinitely often `β` implies infinitely often `False`" is
exactly "eventually always `¬β`". -/
theorem sugarFG_denot {S : Type*} (B : Set S) (ω : ℕ → S) :
    ((∃ᶠ n in atTop, ω n ∈ B) → ∃ᶠ n in atTop, ω n ∈ (∅ : Set S))
      ↔ ∀ᶠ n in atTop, ω n ∉ B
\end{leansrc}
\ruleblock{Correspondence (Steele display $\leftrightarrow$ Lean)}
\badge{BadgeGreen}{reviewer-certified against the displayed texts}
\sourcenote{The dictionary below is a reviewer-certified statement of how the mechanized statements above interpret the displayed rule: display metavariable $\mapsto$ Lean binder, each display hypothesis $\mapsto$ the premise or side condition carrying it, the display conclusion $\mapsto$ the goal --- with any deliberate difference in width between display and mechanization named outright. Certification is recorded as a hash triple over (the displayed cell, the extracted statement texts, this dictionary\textquotesingle s own text); if any of the three changes, this block turns amber until a reviewer re-certifies. The correspondence claim itself remains editorial --- reviewer-grade, not machine-checked; its machine-checked successor types each soundness statement against the rule as a structured object, and supersedes this block per rule as it lands.}
\texttt{sugarFG\_denot} ($\mathsf{FG}\neg\beta :=
\mathsf{Streett}((\beta, \bot))$) certifies the pair event at response
$\emptyset$ $\leftrightarrow$ eventual permanent avoidance
(\texttt{∀ᶠ}), the finitely-often reading --- stated, like its
siblings, in the target rule\textquotesingle s concluded shape, which
is what makes it consumable as a rewrite rather than a mere fact. Width
note: a per-path equivalence over arbitrary state types --- no measure
enters, so the identity transports under every profile\textquotesingle s
path measure by congruence; the remaining accompanying claim (the
entry\textquotesingle s flavor text: the certificate
surface \emph{builds} the expansion rather than minting a node kind) is
a property of the concrete-syntax canonicality battery --- a test, not
a theorem, deliberately outside this bound set, exactly as the roster
witness says.

\ruleblock{Lean verifier receipts}
\sourcenote{module \texttt{BernsteinLogicAux.QtslAliases}; bound build \texttt{570599779c29\ldots} (bound from receipt 5a60f4d01fd81ac5, 2026-07-21T05:39:02Z). A receipt records that named checking layers accepted a build: (module, compiled-artifact SHA-256, environment SHA-256, checker layers), minted when the checking gate ran; the store receipts attesting the bound build are listed with the layers each run credited. Nothing here asks the reader to trust the store: every layer is an independently runnable program over the quoted sources --- rebuild the module and re-run the checkers to reproduce the verdict. Receipts attesting superseded builds credit no layer and are omitted from this rendering (the count of omissions, if any, is noted below).}
\begin{receipttable}
\receiptrow{0a3819a978e31127}{2026-07-21T04:51:36Z}{build+decorrelated}{\badge{BadgeGreen}{current}}
\receiptrow{0bb1074df705d863}{2026-07-21T05:21:58Z}{build+decorrelated}{\badge{BadgeGreen}{current}}
\receiptrow{5a60f4d01fd81ac5}{2026-07-21T05:39:02Z}{build+replay+hardened+decorrelated}{\badge{BadgeGreen}{current}}
\end{receipttable}
\sourcenote{11 further receipt(s) in the store attest superseded builds of this module and are omitted from this rendering.}
\layerline{\layeritem{build}{lake build (elaboration)}{\badge{BadgeGreen}{current}}\layeritem{replay}{leanchecker (official kernel replay)}{\badge{BadgeGreen}{current}}\layeritem{hardened}{safe\_verify (hardened screen)}{\badge{BadgeGreen}{current}}\layeritem{decorrelated}{lean4export + nanoda (decorrelated kernel)}{\badge{BadgeGreen}{current}}}
\sourcenote{\badge{BadgeGray}{comparator: none} no fixture in the three-implementation comparator harness (compiled / kernel-interpreted / reference-interpreted) exercises this rule.}
\ruleblock{HOL4 statement}
\absentbox{OUT OF SCOPE BY DESIGN}{Out of scope for the HOL side \emph{by design}: the sugar is a
definitional constructor that builds its expansion at the surface ---
no sugar node exists in any representation downstream of parsing, wire
format included, so there is no artefact for any checker port to be
about. Checked against the HOL sources: the \texttt{bern\_catalogue}
theory carries no spec-surface grammar.
}
\ruleblock{HOL verifier receipts}
\absentbox{NONE}{no HOL verification receipt exists for this rule.}
\ruleblock{Proof}

\sourcenote{Genre: alias rule (tier 3) --- a respelling template; the
shared template is Recipe~\ref{rec:sd-template} of the
\rulename{QStreettDegenerate} section. The unpack is a \emph{definition},
so the five artefacts collapse to the definitional equation, its
reading, and the target section.}

\begin{brecipe}[the filled-in application]\label{rec:qfg-filled}
\[
\mathsf{FG}\neg\beta \;:=\; \mathsf{Streett}((\beta, \bot))
\]
The concluding judgment discharges through one existing rule: eventual
permanent avoidance through one \rulename{QStreettBattery} node at the
pair $(\beta, \bot)$. The honest scope is this polarized form
($\mathsf{FG}$ applied to $\neg\beta$): general persistence
$\mathsf{FG}\varphi$ is not this rule.
\end{brecipe}

\begin{brecipe}[payload transport]\label{rec:qfg-transport}
The expansion $\mathsf{Streett}((\beta, \bot))$ carries one bundle
pair: requests $[\![\beta]\!]$, response empty. No unpack function is
owed here, because none will ever exist: the constructor \emph{is} its
expansion, so there is no sugar payload to transport from --- the
strongest possible form of producers-elaborate-before-emission.
\end{brecipe}

\begin{bprop}[the exactness identity]\label{prop:qfg-exact}
The pair $(\beta, \bot)$ holds iff $\beta$-states occur only finitely
often, i.e.\ iff eventually the path permanently satisfies $\neg\beta$:
\[
\Bigl(
\bigl(\exists^\infty n,\ \omega_n \in [\![\beta]\!]\bigr) \Rightarrow
\bigl(\exists^\infty n,\ \omega_n \in \varnothing\bigr)
\Bigr)
\;\Longleftrightarrow\;
\neg\bigl(\exists^\infty n,\ \omega_n \in [\![\beta]\!]\bigr) ,
\]
since the consequent is false outright --- two lines: an implication
into falsehood is the negation of its hypothesis, and finitely many
$\beta$-visits along a discrete path is precisely a tail of
$\neg\beta$-states. Pure per-path logic; no conservativity and no chart
hypothesis is consumed.
\end{bprop}
\begin{proof}
As displayed; the two lines are shown.
\end{proof}

\begin{bprop}[soundness corollary]\label{prop:qfg-sound}
The sugared judgment is concluded by the target rule's theorem, with
no additional obligation: the sugar constructor builds the expansion,
so the sugared judgment \emph{is} the target judgment.
\end{bprop}
\begin{proof}
Definitional unfolding plus the target section's theorem. The witness
that the surface implementation honours the definition --- that the
constructor builds its expansion rather than introducing a new node
kind --- is the reference implementation's concrete-syntax canonicality
test battery; no checker-facing artefact depends on it, since no
downstream representation contains a sugar node.
\end{proof}

\begin{bremark}[why $\mathsf{FG}$ costs a bundle and $\mathsf{GF}$ costs
only its degenerate case]
The two dual sugars land at different points of the same rule.
$\mathsf{GF}\varphi$ (the twin, \rulename{QStreettGF}) requests always
and so erases the request machinery; $\mathsf{FG}\neg\beta$ makes
$\beta$ itself the request and forbids response --- the ranking
certificate must then pay out at every $\beta$-visit with nothing ever
refunding it, which is exactly a proof that $\beta$-visits are finitely
many. That both directions are single nodes of one rule --- where
automaton-primitive designs pay a product construction for one of them
--- is the practical content of the Streett-primitive choice (see the
\rulename{QStreettBattery} section's remark).
\end{bremark}

\begin{anchors}
target & the \rulename{QStreettBattery} section, one pair $(\beta, \bot)$\\
the identity, event-level & \texttt{sugarFG\_denot} (bound in the entry)\\
twin & \rulename{QStreettGF} (recurrence; the member that consumes a chart)\\
constructor builds expansion & the reference implementation\textquotesingle s concrete-syntax canonicality battery (documentary; no wire artefact exists to check)\\
\end{anchors}


\rulesection{58}{QNegFAlias}{alias for \textcolor{DiffBlue}{\rulename{QSafetyBattery}}}
\label{blcatrule-QNegFAlias}
\ruleblock{The rule (Steele metanotation)}
\sourcenote{as realized in the mechanized development; source row: \texttt{\texttt{QTSL-rules.toml}} --- \emph{QNegFAlias}. The metanotation statement is typeset by hand for notational quality; its mechanized ground truth is the Lean statement below, extracted from the sources at every build of this book and hash-bound. The rule \emph{roster} is machine-checked (every mechanized rule tag round-trips to exactly one catalogue row); the statement \emph{text} is editorial, checked by eye against the extracted Lean beside it.}
\begin{mathpar}
\inferrule*{
  \mathrm{child} : \mathsf{QRegionSat}\bigl(\mathsf{SProbGE}(p, \mathsf{PAlways}(\mathsf{SNot}\,\varphi)),\ R\bigr)
}{
  \mathcal{M},\ x_0 \in R \vDash \mathsf{SProbGE}(p, \mathsf{PNotPath}(\mathsf{PDeadline}(T, \varphi, \mathrm{false})))
}
\end{mathpar}

\par\medskip
\rulename{QNegFAlias} discharges a negated deadline: to certify, with
probability at least $p$, that $\varphi$ does \emph{not} occur within
horizon $T$ ($\neg\mathsf{F}_{\le T}\,\varphi$), it suffices to certify
the \emph{unbounded} safety claim --- $\varphi$ never occurs at all ---
over the same region, and the alias rewrites the one judgment into the
other: the child is a satisfaction judgment for
always-not-$\varphi$, the conclusion the negated-deadline judgment,
region verbatim. It is an alias to the safety battery
(\rulename{QSafetyBattery}): one node to one node, and the entire
mathematical content is a two-line denotation inclusion --- a path on
which $\varphi$ never holds certainly does not reach $\varphi$ within
$T$, and not-reaching-within-$T$ \emph{is} the bounded-always of the
negation. The rewrite is choice-free (no producer parameter beyond the
horizon already present in the goal), and the checker recomputes
nothing new: the child is an ordinary safety-battery conclusion.

The rule is deliberately lossy in one direction --- the unbounded claim
is strictly stronger than the bounded one --- which is the honest price
of reusing the safety kernel with no horizon bookkeeping. A producer
who needs the sharp bounded claim (because $\varphi$ \emph{does} occur
eventually, only late) uses the deadline machinery instead; this alias
serves the common case where the avoidance is permanent.

Mechanization state, plainly: the bound statements below are the two
lines, as priced. \texttt{negF\_denot\_subset} (module
\texttt{BernsteinLogicAux.QtslAliases}) is the denotation inclusion
--- a path on which $\varphi$ never holds does not reach $\varphi$
within $T$ --- pure path-set algebra, axiom-free (the kernel certifies
it depends on no axioms at all); and \texttt{sound\_QNegFAlias} is the
alias\textquotesingle s transport: a lower bound on the unbounded
safety event (the target\textquotesingle s conclusion shape, with
$\mathrm{sat} := \neg\varphi$) is a lower bound on the
negated-deadline event, for \emph{any} measure on path space ---
profile-generic, so the discrete-kernel and GSHS records read it
alike, and the entry carries no conditional flag.

\cataloguenote{Two flags travel with the row, both from the source catalogue. Tier
(flagged, not forced): alias as sorted, but if the folded denotation
inclusion must itself be a rule application, the tier moves to
admissible --- machinery unchanged either way. And the mirror is a
deliberate hole: there is no dual rule discharging a negated
\emph{always} by a bounded operator, because the conclusion it would
need ($\mathsf{F}\varphi$ unbounded) is structurally unrepresentable in
the spec grammar --- the unbounded-eventually constructor was excluded
by testable-negative design, unbounded eventuality surviving only as a
producer-choice provenance flag on deadlines. The pair must never be
ported as one duality entry: one side is a statable choice-free alias,
the other hides a horizon choice and a metatheory line inside the
deadline battery.
}
\ruleblock{Lean statement}
\declhead{BernsteinLogicAux.QtslAliases}{negF\_denot\_subset}{\badge{BadgeGreen}{statement bound}}
\begin{leansrc}
/-- **The denotation inclusion** (QTSL-rules.toml row `QNegFAlias` —
"the denotation inclusion `⟦G¬φ⟧ ⊆ ⟦G_{≤T}¬φ⟧ = ⟦¬F_{≤T}φ⟧` folds into the
rewrite"; the catalogue's "two lines"): a path on which `φ` never holds
certainly does not reach `φ` within `T` — and not-reaching-within-`T` *is*
the negated deadline event. Pure path-set algebra; no measure, no
measurability. -/
theorem negF_denot_subset {S : Type*} (sat : S → Prop) (T : ℕ) :
    {ω : ℕ → S | ∀ n, ¬ sat (ω n)} ⊆ {ω : ℕ → S | ¬ ∃ k ≤ T, sat (ω k)}
\end{leansrc}
\declhead{BernsteinLogicAux.QtslAliases}{sound\_QNegFAlias}{\badge{BadgeGreen}{statement bound}}
\begin{leansrc}
/-- **`sound_QNegFAlias`** (rule row `QNegFAlias`): the alias's
transport — a lower bound on the probability of the unbounded safety event
(`QSafetyBattery`'s conclusion shape, e.g. `sound_QSafetyBattery` /
`sound_QSafetyBattery_staged` with `sat := ¬φ`) is a lower bound on the
probability of the negated-deadline event. Profile-generic (any measure on
path space — discrete-kernel and GSHS records alike), by monotonicity
through `negF_denot_subset`. The rewrite is choice-free: no producer
parameter beyond the horizon `T` already present in the goal. -/
theorem sound_QNegFAlias {S : Type*} [MeasurableSpace S]
    (μ : Measure (ℕ → S)) {sat : S → Prop} {b : ENNReal} {T : ℕ}
    (h : b ≤ μ {ω | ∀ n, ¬ sat (ω n)}) :
    b ≤ μ {ω | ¬ ∃ k ≤ T, sat (ω k)}
\end{leansrc}
\ruleblock{Correspondence (Steele display $\leftrightarrow$ Lean)}
\badge{BadgeGreen}{reviewer-certified against the displayed texts}
\sourcenote{The dictionary below is a reviewer-certified statement of how the mechanized statements above interpret the displayed rule: display metavariable $\mapsto$ Lean binder, each display hypothesis $\mapsto$ the premise or side condition carrying it, the display conclusion $\mapsto$ the goal --- with any deliberate difference in width between display and mechanization named outright. Certification is recorded as a hash triple over (the displayed cell, the extracted statement texts, this dictionary\textquotesingle s own text); if any of the three changes, this block turns amber until a reviewer re-certifies. The correspondence claim itself remains editorial --- reviewer-grade, not machine-checked; its machine-checked successor types each soundness statement against the rule as a structured object, and supersedes this block per rule as it lands.}
The displayed rewrite --- child
$\mathsf{QRegionSat}(\mathsf{SProbGE}(p, \mathsf{PAlways}(\mathsf{SNot}\,
\varphi)), R)$ concluding
$\mathsf{SProbGE}(p, \mathsf{PNotPath}(\mathsf{PDeadline}(T, \varphi,
\mathrm{false})))$ --- $\mapsto$ the two-statement chain:
\texttt{negF\_denot\_subset} (the denotation inclusion \texttt{\{ω |
∀ n, ¬ sat (ω n)\} ⊆ \{ω | ¬ ∃ k ≤ T, sat (ω k)\}} --- pure path-set
algebra, axiom-free, no measure, no measurability) and
\texttt{sound\_QNegFAlias}, its transport: a lower bound \texttt{b ≤
μ} on the unbounded-safety event is a lower bound on the
negated-deadline event, for \emph{any} measure on path space ---
profile-generic by binder shape (\texttt{μ : Measure (ℕ → S)}
arbitrary), the \rulename{QUnionBound} pattern; discrete-kernel and
GSHS records alike are instances. The
display\textquotesingle s choice-freeness is a binder fact: no
producer parameter appears beyond the horizon \texttt{T} already in
the goal. Width note: the child\textquotesingle s
$\mathsf{SProbGE}$ shape enters only as the abstract bound \texttt{b ≤
μ \{…\}} --- the certified transport is wider than the displayed
probabilistic judgment (any \texttt{ENNReal} bound transports), safe
direction; paths are step-indexed (\texttt{ℕ → S}), the deadline
read discretely, as at \rulename{TaylorBound}.

\ruleblock{Lean verifier receipts}
\sourcenote{module \texttt{BernsteinLogicAux.QtslAliases}; bound build \texttt{570599779c29\ldots} (bound from receipt 5a60f4d01fd81ac5, 2026-07-21T05:39:02Z). A receipt records that named checking layers accepted a build: (module, compiled-artifact SHA-256, environment SHA-256, checker layers), minted when the checking gate ran; the store receipts attesting the bound build are listed with the layers each run credited. Nothing here asks the reader to trust the store: every layer is an independently runnable program over the quoted sources --- rebuild the module and re-run the checkers to reproduce the verdict. Receipts attesting superseded builds credit no layer and are omitted from this rendering (the count of omissions, if any, is noted below).}
\begin{receipttable}
\receiptrow{0a3819a978e31127}{2026-07-21T04:51:36Z}{build+decorrelated}{\badge{BadgeGreen}{current}}
\receiptrow{0bb1074df705d863}{2026-07-21T05:21:58Z}{build+decorrelated}{\badge{BadgeGreen}{current}}
\receiptrow{5a60f4d01fd81ac5}{2026-07-21T05:39:02Z}{build+replay+hardened+decorrelated}{\badge{BadgeGreen}{current}}
\end{receipttable}
\sourcenote{11 further receipt(s) in the store attest superseded builds of this module and are omitted from this rendering.}
\layerline{\layeritem{build}{lake build (elaboration)}{\badge{BadgeGreen}{current}}\layeritem{replay}{leanchecker (official kernel replay)}{\badge{BadgeGreen}{current}}\layeritem{hardened}{safe\_verify (hardened screen)}{\badge{BadgeGreen}{current}}\layeritem{decorrelated}{lean4export + nanoda (decorrelated kernel)}{\badge{BadgeGreen}{current}}}
\sourcenote{\badge{BadgeGray}{comparator: none} no fixture in the three-implementation comparator harness (compiled / kernel-interpreted / reference-interpreted) exercises this rule.}
\ruleblock{HOL4 statement}
\absentbox{OUT OF SCOPE BY DESIGN}{Out of scope for the HOL side \emph{by design}: the alias unpacks to a
single safety-battery node before emission --- the wire never carries a
negated-deadline node, so the HOL statement that would verify the
discharged certificate is the safety battery\textquotesingle s own
(absent today: the satisfaction stratum has no HOL artefact, per that
section). Checked against the HOL sources: the \texttt{bern\_catalogue}
theory is the base-calculus checker cone and carries no spec-surface
path operators.
}
\ruleblock{HOL verifier receipts}
\absentbox{NONE}{no HOL verification receipt exists for this rule.}
\ruleblock{Proof}

\sourcenote{Genre: alias rule (tier 3) --- a respelling template; the
batch's shared template is Recipe~\ref{rec:sd-template} of the
\rulename{QStreettDegenerate} section. The soundness content is two lines, recorded as such in the
source catalogue, and both are now mechanized and bound in the entry:
the inclusion (\texttt{negF\_denot\_subset}, axiom-free) and its
measure transport (\texttt{sound\_QNegFAlias}, profile-generic).}

\begin{brecipe}[the filled-in application]\label{rec:nf-filled}
\begin{mathpar}
\inferrule*[Right=\rulename{QNegFAlias}]{
  \mathrm{child} : \mathsf{QRegionSat}\bigl(\mathsf{SProbGE}(p,
  \mathsf{PAlways}(\mathsf{SNot}\,\varphi)),\ R\bigr)
  \ \text{(one \rulename{QSafetyBattery} conclusion)}
}{
  \mathcal{M},\ x_0 \in R \vDash
  \mathsf{SProbGE}\bigl(p,\ \mathsf{PNotPath}(\mathsf{PDeadline}(T,
  \varphi, \mathrm{false}))\bigr)
}
\end{mathpar}
One node to one node: the child is the safety battery's conclusion at
avoid set $[\![\varphi]\!]$, the concluded judgment re-reads the same
certified region against the negated deadline, horizon $T$ taken from
the goal. In the staged spec vocabulary the conclusion's formula is the
negated-path reading of a deadline node whose met-flag is false; the
child's is the always of the negation --- the rewrite touches the
formula only.
\end{brecipe}

\begin{brecipe}[payload transport]\label{rec:nf-transport}
\begin{center}\small
\begin{tabular}{@{}ll@{}}
\textbf{\rulename{QNegFAlias} node} & \textbf{\rulename{QSafetyBattery} node}\\
\midrule
certified region $R$, probability $p$ & verbatim\\
goal formula $\neg\mathsf{F}_{\le T}\,\varphi$ & child formula $\mathsf{G}\,\neg\varphi$ (horizon dropped)\\
horizon $T$ & not consumed (\autoref{rem:nf-lossy})\\
\end{tabular}
\end{center}
Unpack-function slot: no verified payload function in this edition
(batch template); the
semantic witness is bound in the entry (\texttt{negF\_denot\_subset} +
\texttt{sound\_QNegFAlias}).
\end{brecipe}

\begin{bprop}[the exactness inclusion --- two lines]\label{prop:nf-exact}
$[\![\mathsf{G}\,\neg\varphi]\!] \subseteq
[\![\mathsf{G}_{\le T}\,\neg\varphi]\!] =
[\![\neg\mathsf{F}_{\le T}\,\varphi]\!]$.
\end{bprop}
\begin{proof}
A path that never satisfies $\varphi$ in particular does not satisfy
$\varphi$ at any time $\le T$ --- the inclusion, by restriction of a
universal quantifier. And not satisfying $\varphi$ at any time $\le T$
is the definition of not reaching $\varphi$ within $T$ --- the equality,
definitional. Hence every path in the child's path set lies in the
conclusion's, and the probability bound transports monotonically.
\end{proof}

\begin{bprop}[soundness corollary]\label{prop:nf-sound}
If the child judgment holds --- from every state of $R$, with
probability $\ge p$ the path always avoids $[\![\varphi]\!]$ --- then
from every state of $R$, with probability $\ge p$ the path does not
reach $[\![\varphi]\!]$ within $T$.
\end{bprop}
\begin{proof}
Proposition~\ref{prop:nf-exact} plus monotonicity of measures, applied
under the safety battery's theorem. There is no other content --- and
this two-step corollary is precisely what the bound statements are:
the inclusion and the transport, one decl each.
\end{proof}

\begin{bremark}[the deliberate lossiness]\label{rem:nf-lossy}
The alias proves the bounded claim from the unbounded one, discarding
$T$. When $\varphi$ genuinely occurs eventually --- only after the
horizon --- the child judgment is unprovable and this alias is the wrong
tool; the deadline machinery (the deadline battery, with its geometric
decay certificates) proves the sharp bounded claim instead. The alias
exists because permanent avoidance is the common case, and for it the
safety kernel is both cheaper and already ported.
\end{bremark}

\begin{bremark}[the mirror that is not a rule]\label{rem:nf-mirror}
There is deliberately no dual entry discharging a negated
\emph{always}: its conclusion would need unbounded eventuality
($\mathsf{F}\varphi$), and the spec grammar has no unbounded-eventually
constructor --- an exclusion by testable-negative design, unbounded
eventuality surviving only as a provenance flag recording that a
deadline's horizon was producer-chosen. What replaces the missing rule
is three separate pieces: a negation-normal-form duality
(\rulename{QNnfRewrite}), a horizon instantiation (a producer choice,
flagged), and a weakening line \emph{inside} the deadline battery's
lemma where the flag is discharged. The asymmetry is worth this remark
because a port that treated the pair as one duality entry would smuggle
a choice and a metatheory line into what looks like a definitional
rewrite.
\end{bremark}

\begin{bremark}[tier flag]
Alias as sorted; if the folded inclusion of
Proposition~\ref{prop:nf-exact} must itself be a rule application
rather than a lemma of the unpack, the tier moves to admissible ---
flagged in the source catalogue, machinery unchanged either way.
\end{bremark}

\begin{anchors}
Recipe \ref{rec:nf-filled} target & the \rulename{QSafetyBattery} section\\
Proposition \ref{prop:nf-exact} + corollary & \texttt{negF\_denot\_subset} + \texttt{sound\_QNegFAlias} (two lines, as priced; bound in the entry)\\
Remark \ref{rem:nf-mirror} & the spec grammar\textquotesingle s constructor census (no unbounded-eventually node)\\
\end{anchors}


\blcatsection{Discharge Kits, Closure, and Arithmetic Leaves}
\rulesection{59}{HypMultiplier}{admissible rule}
\label{blcatrule-HypMultiplier}
\ruleblock{The rule (Steele metanotation)}
\sourcenote{as realized in the mechanized development; source row: \texttt{\texttt{QTSL-rules.toml}} --- \emph{HypMultiplier}. The metanotation statement is typeset by hand for notational quality; its mechanized ground truth is the Lean statement below, extracted from the sources at every build of this book and hash-bound. The rule \emph{roster} is machine-checked (every mechanized rule tag round-trips to exactly one catalogue row); the statement \emph{text} is editorial, checked by eye against the extracted Lean beside it.}
\begin{mathpar}
\inferrule*{
  \vdash \lambda \ge 0 \ \text{on } \mathrm{cell} \\
  \vdash N + \lambda\, g_q \le 0 \ \text{on } \mathrm{cell}
}{
  \vdash N \le 0 \ \text{on } \mathrm{cell} \cap \{g_q \ge 0\}
}
\end{mathpar}
the child atom $\{g_q \ge 0\}$ is a concluded child certificate\textquotesingle s region.

\ruleflavor{the child atom enters through the presentation bridge; both premises are \rulename{BernCone}-family rows, replayed by the base kernel — realized as the verified pointwise warrant \texttt{hypMultiplier\_sound}.}
\par\medskip
\rulename{HypMultiplier} is the discharge kit\textquotesingle s
multiplier move, and the only member of the kit that manufactures a
new row rather than reshaping the scope of an existing one. The
situation it serves: a premise row of a satisfaction-judgment rule
must hold not on a whole cell but on the part of the cell where a
child certificate\textquotesingle s region atom holds --- $N \le 0$ on
$\mathrm{cell} \cap \{g_q \ge 0\}$, with $\{g_q \ge 0\}$ the region
concluded by a nested certificate. The base
calculus\textquotesingle s rows are scoped to cells, not to
cell-with-polynomial-atom intersections, so the obligation cannot be
posted as a base row directly. The rule buys the restriction with one
multiplier: exhibit $\lambda$ and post two ordinary cone-family rows
--- $\lambda \ge 0$ on the cell, and $N + \lambda\, g_q \le 0$ on the
cell.

What the checker recomputes: nothing beyond the two premise rows,
which are cone-membership sequents the base kernel already replays;
the combination step is the degree-one corner of the
positivstellensatz family, the same multiplier discipline as the
calculus\textquotesingle s Farkas machinery. The producer typically
finds $\lambda$ by linear programming; only the two rows ship.

Both hypotheses are load-bearing, in opposite ways. Drop the sign row
and a negative multiplier flips the direction of the trade: with
$N \equiv 1$, $\lambda \equiv -1$, $g_q \equiv 1$, the combined row
$N + \lambda g_q \equiv 0 \le 0$ holds while the conclusion fails
everywhere on the child atom. Widen the conclusion to the whole cell
and the rule is unsound wherever the child atom fails: where
$g_q(x) < 0$ the term $\lambda(x)\, g_q(x)$ is a subsidy rather than
a tax, and the combined row bounds nothing. The proof section walks
both, counterexamples attached.

\cataloguenote{Before this rule\textquotesingle s mechanization, the move lived in
the producer toolkit as unchecked metatheory --- the discharge step
performed the combination and asserted the scoped row. The verified
lemma closes that trust boundary. The certificate row that would let
the wire format carry the move natively is out of scope of the current
baked theory, so today the lemma warrants the toolkit\textquotesingle s
elaboration rather than a kernel-replayed node.
}
\ruleblock{Lean statement}
\par\medskip
The auxiliary library states the extension theory\textquotesingle s
structural lemmas over the compiled semantic class: a bare state
type, scopes and cells as sets, observables as real-valued functions.
The presentation-level child atom enters through the named bridge
hypothesis (\texttt{hbridge}), which the presentation layer
discharges definitionally at compile time; nothing else about the
system enters at all.

\declhead{BernsteinLogicAux.Structural}{hypMultiplier\_sound}{\badge{BadgeGreen}{statement bound}}
\begin{leansrc}
/-- **T2 HypMultiplier**, soundness. `childAtom` is the presentation-level
child atom; its bridge names the compiled superlevel reading `{x | 0 ≤ g x}`.
The premises are the two compiled BernCone-family rows. -/
theorem hypMultiplier_sound (cell childAtom : Set X) (N lam g : X → ℝ)
    (hbridge : childAtom = {x | 0 ≤ g x})
    (hlam : ∀ x ∈ cell, 0 ≤ lam x)
    (hcomb : ∀ x ∈ cell, N x + lam x * g x ≤ 0) :
    ∀ x ∈ cell ∩ childAtom, N x ≤ 0
\end{leansrc}
\ruleblock{Correspondence (Steele display $\leftrightarrow$ Lean)}
\badge{BadgeGreen}{reviewer-certified against the displayed texts}
\sourcenote{The dictionary below is a reviewer-certified statement of how the mechanized statements above interpret the displayed rule: display metavariable $\mapsto$ Lean binder, each display hypothesis $\mapsto$ the premise or side condition carrying it, the display conclusion $\mapsto$ the goal --- with any deliberate difference in width between display and mechanization named outright. Certification is recorded as a hash triple over (the displayed cell, the extracted statement texts, this dictionary\textquotesingle s own text); if any of the three changes, this block turns amber until a reviewer re-certifies. The correspondence claim itself remains editorial --- reviewer-grade, not machine-checked; its machine-checked successor types each soundness statement against the rule as a structured object, and supersedes this block per rule as it lands.}
One bound statement, transcription-close. The displayed premises
$\vdash \lambda \ge 0$ on cell and $\vdash N + \lambda g_q \le 0$ on
cell $\mapsto$ \texttt{hlam}/\texttt{hcomb}, pointwise on
\texttt{cell}; the conclusion $\vdash N \le 0$ on $\mathrm{cell} \cap
\{g_q \ge 0\}$ $\mapsto$ the goal \texttt{∀ x ∈ cell ∩ childAtom, N x
≤ 0}, with the child atom entering through the System-bridge
hypothesis \texttt{hbridge : childAtom = \{x | 0 ≤ g x\}} --- the
presentation-level region named by an equality the presentation layer
will discharge definitionally; until then it is an honest binder, not
an assertion. Width note: \texttt{N}, \texttt{lam}, \texttt{g} are
arbitrary real-valued functions and \texttt{cell} an arbitrary set ---
the BernCone row structure and polynomiality are consumed one layer
down where the rows are replayed; the certified warrant is the
pointwise Farkas step alone, degree-1 positivstellensatz as displayed,
strictly wider than the polynomial instance, safe direction.

\ruleblock{Lean verifier receipts}
\sourcenote{module \texttt{BernsteinLogicAux.Structural}; bound build \texttt{1222cdcc326f\ldots} (bound from receipt cb0c3820c6877f22, 2026-07-29T01:03:36Z). A receipt records that named checking layers accepted a build: (module, compiled-artifact SHA-256, environment SHA-256, checker layers), minted when the checking gate ran; the store receipts attesting the bound build are listed with the layers each run credited. Nothing here asks the reader to trust the store: every layer is an independently runnable program over the quoted sources --- rebuild the module and re-run the checkers to reproduce the verdict. Receipts attesting superseded builds credit no layer and are omitted from this rendering (the count of omissions, if any, is noted below).}
\begin{receipttable}
\receiptrow{0b7e68907f45b27f}{2026-07-28T23:48:36Z}{build+decorrelated}{\badge{BadgeGreen}{current}}
\receiptrow{e97d5614d4f08c7e}{2026-07-29T00:00:35Z}{build+decorrelated}{\badge{BadgeGreen}{current}}
\receiptrow{27de21f047e6b194}{2026-07-29T00:49:05Z}{build+replay+hardened+decorrelated}{\badge{BadgeGreen}{current}}
\receiptrow{cb0c3820c6877f22}{2026-07-29T01:03:36Z}{build+replay+hardened+decorrelated}{\badge{BadgeGreen}{current}}
\receiptrow{3530841119a9a607}{2026-07-31T22:00:22Z}{build+decorrelated}{\badge{BadgeGreen}{current}}
\receiptrow{b6cc23eb3e5b7f97}{2026-07-31T22:29:16Z}{build+decorrelated}{\badge{BadgeGreen}{current}}
\receiptrow{bf3dacc3ab7f20ea}{2026-08-02T23:31:31Z}{build+decorrelated}{\badge{BadgeGreen}{current}}
\receiptrow{384a8569d2005df7}{2026-08-03T16:27:35Z}{build+decorrelated}{\badge{BadgeGreen}{current}}
\receiptrow{717e07acf9acac8a}{2026-08-03T22:46:19Z}{build+decorrelated}{\badge{BadgeGreen}{current}}
\end{receipttable}
\sourcenote{11 further receipt(s) in the store attest superseded builds of this module and are omitted from this rendering.}
\layerline{\layeritem{build}{lake build (elaboration)}{\badge{BadgeGreen}{current}}\layeritem{replay}{leanchecker (official kernel replay)}{\badge{BadgeGreen}{current}}\layeritem{hardened}{safe\_verify (hardened screen)}{\badge{BadgeGreen}{current}}\layeritem{decorrelated}{lean4export + nanoda (decorrelated kernel)}{\badge{BadgeGreen}{current}}}
\sourcenote{\badge{BadgeGray}{comparator: none} no fixture in the three-implementation comparator harness (compiled / kernel-interpreted / reference-interpreted) exercises this rule.}
\ruleblock{HOL4 statement}
\absentbox{NO HOL STATEMENT YET}{The HOL4 development\textquotesingle s object is the base
calculus\textquotesingle s wire-format checker (theory
\texttt{bern\_catalogue} ports the catalogue check-function cone).
This rule belongs to the extension theory\textquotesingle s discharge
kit: it postdates that port, and no row for the move exists in the
baked theory --- there is nothing for a HOL statement to be
about. The box becomes fillable when a revision carries the
move\textquotesingle s certificate rows and the corresponding check is
ported.
}
\ruleblock{HOL verifier receipts}
\absentbox{NONE}{no HOL verification receipt exists for this rule.}
\ruleblock{Proof}

\sourcenote{Genre: admissible rule (tier 2) --- a proof-tree-rewriting
recipe; the general admissible template, both species included, is
stated in the \rulename{Norm} section and instantiated here in its
boundary species (discharge against a named verified theorem). The rule
is sequent-side machinery: premises and conclusion are
all $\vdash$ rows, and what the extension theory contributes is the
\emph{scope} vocabulary --- the child atom $\{g_q \ge 0\}$ names a
region concluded by a nested certificate, a shape no base cell can
express. The rule closes a named trust boundary: before its
mechanization, the move was performed producer-side and asserted.}

\begin{brecipe}[the \rulename{HypMultiplier} instance]\label{rec:hm-instance}
\[
\inferrule*[Right=\rulename{HypMultiplier}]{
  \mathcal{T}_\lambda \rhd\ \vdash \lambda \ge 0 \ \text{on } \mathrm{cell} \\
  \mathcal{T}_c \rhd\ \vdash N + \lambda\, g_q \le 0 \ \text{on } \mathrm{cell}
}{
  \vdash N \le 0 \ \text{on } \mathrm{cell} \cap \{g_q \ge 0\}
}
\qquad\rightsquigarrow\qquad
\begin{array}{l}
\text{\upshape the premise subtrees stand unchanged;}\\
\text{\texttt{hypMultiplier\_sound} discharges the scoped row.}
\end{array}
\]
The rewrite is degenerate in the same honest sense as the base
catalogue\textquotesingle s ground-fact rule: the conclusion is not
derivable from catalogue nodes, because the base grammar has no row
form scoped to a cell-with-polynomial-atom intersection --- that scope
exists one layer up, where the child certificate that concluded
$\{g_q \ge 0\}$ lives. The two premise subtrees are ordinary
cone-family derivations, checked by the base kernel unchanged; the
node itself discharges against the verified pointwise warrant below,
which consumes the two premises plus the one bridge hypothesis naming
the child atom's compiled reading (the closing remark) --- and nothing
else.
\end{brecipe}

\begin{bprop}[semantic warrant]\label{prop:hm-sound}
If $\lambda \ge 0$ everywhere on the cell and $N + \lambda g_q \le 0$
everywhere on the cell, then $N \le 0$ everywhere on
$\mathrm{cell} \cap \{g_q \ge 0\}$.
\end{bprop}
\begin{proof}
Fix $x$ in the intersection. Then $g_q(x) \ge 0$ (membership in the
child atom) and $\lambda(x) \ge 0$ (the sign row), so
$\lambda(x)\, g_q(x) \ge 0$, and
\[
N(x) \;\le\; N(x) + \lambda(x)\, g_q(x) \;\le\; 0,
\]
the first inequality by nonnegativity of the added term, the second by
the combined row at $x$. That is the whole proof ---
\texttt{hypMultiplier\_sound} is the two displayed inequalities.
\end{proof}

\begin{bremark}[the sign row is half the rule]
Drop $\vdash \lambda \ge 0$ and the rule is unsound. Counterexample:
$N \equiv 1$, $\lambda \equiv -1$, $g_q \equiv 1$ on any nonempty cell
(a singleton suffices). The
combined row holds identically ($1 - 1 = 0 \le 0$), yet $N \equiv 1$
fails everywhere on $\mathrm{cell} \cap \{g_q \ge 0\} = \mathrm{cell}$.
A free-signed multiplier converts the tax the child atom levies into a
subsidy: the combination can cancel exactly the positivity the
conclusion denies.
\end{bremark}

\begin{bremark}[the scope restriction is the other half]
The conclusion cannot claim the whole cell. With $N \equiv 1$,
$\lambda \equiv 1$, $g_q \equiv -1$ on the same nonempty cell, both
premises hold
($\lambda \ge 0$; $1 - 1 = 0 \le 0$) while $N > 0$ everywhere --- and
the rule is nonetheless sound, because
$\mathrm{cell} \cap \{g_q \ge 0\} = \varnothing$ and the scoped claim
is vacuous. Where $g_q(x) < 0$ the term $\lambda(x) g_q(x)$ subsidizes
$N$ instead of taxing it, and the combined row bounds nothing: the
scoped conclusion is exactly as strong as the premises pay for, no
stronger.
\end{bremark}

\begin{bremark}[degree-one positivstellensatz, and where the search lives]
The move is the degree-one corner of certifying positivity over a
constrained set by nonnegative combinations: Farkas\textquotesingle{}
multiplier discipline \cite{farkas1902} (see \cite{schrijver1986} for
the sign conventions) in the affine case, sitting at the bottom of the
polynomial cone hierarchy of Krivine \cite{krivine1964} and Handelman
\cite{handelman1988} that the calculus\textquotesingle s cone rule
draws on. Nothing beyond degree one is consumed here: one multiplier,
one product, no certificate machinery beyond the two posted rows. The
\emph{search} for $\lambda$ (typically a linear program over the
cone-membership witnesses) is producer-side; only its result ships,
and the checker replays the two rows without ever seeing the search.
\end{bremark}

\begin{bremark}[the bridge hypothesis]
In the verified statement the child atom is an abstract set with one
bridge hypothesis naming its compiled superlevel reading
($\{x \mid 0 \le g_q(x)\}$). When the presentation layer of the
extension theory lands, the hypothesis discharges definitionally from
the compile --- the scoped conclusion itself is unchanged; only the
bridge binder becomes definitional. Until then it is the honest
record of exactly what the warrant assumes about the child
certificate\textquotesingle s conclusion --- its region is the
superlevel set it claims to be, and nothing more.
\end{bremark}

\begin{anchors}
Recipe \ref{rec:hm-instance} warrant & \texttt{hypMultiplier\_sound} (BernsteinLogicAux/Structural.lean)\\
premise rows & cone-family sequent machinery of the base calculus (the \rulename{BernCone} and \rulename{Farkas} sections)\\
child-atom reading & the \texttt{hbridge} hypothesis of \texttt{hypMultiplier\_sound}\\
trust boundary closed & the producer-side discharge step this lemma now warrants\\
\end{anchors}


\rulesection{60}{QFullDomainWeaken}{alias rule}
\label{blcatrule-QFullDomainWeaken}
\ruleblock{The rule (Steele metanotation)}
\sourcenote{as realized in the mechanized development; source row: \texttt{\texttt{QTSL-rules.toml}} --- \emph{QFullDomainWeaken}. The metanotation statement is typeset by hand for notational quality; its mechanized ground truth is the Lean statement below, extracted from the sources at every build of this book and hash-bound. The rule \emph{roster} is machine-checked (every mechanized rule tag round-trips to exactly one catalogue row); the statement \emph{text} is editorial, checked by eye against the extracted Lean beside it.}
\begin{mathpar}
\inferrule*{
  \vdash P \ \text{on } D_m \\
  \mathrm{scope} \subseteq D_m \ \text{(compile fact, not a posted premise)}
}{
  \vdash P \ \text{on scope}
}
\end{mathpar}
emitted shape: $\mathrm{scope} = D_m \smallsetminus \mathrm{Off}$ (the off-scoped row).

\ruleflavor{payload verbatim, scope replaced — one application at the region layer, realized as \texttt{fullDomainWeaken\_sound} (axiom-free) $+$ \texttt{\_offScope}.}
\par\medskip
\rulename{QFullDomainWeaken} is the cheapest move in the discharge
kit: an off-scoped row --- an obligation $P$ on
$D_m \smallsetminus \mathrm{Off}$, the part of a
mode\textquotesingle s domain outside an off-region --- is discharged
by proving $P$ on all of $D_m$ and forgetting the carve-out. It is an
alias: one scope-weakening application at the region layer,
transporting the row payload verbatim and replacing the scope by the
full domain.

A producer reaches for it first. When the row happens to hold
globally, the off-region was doing no work for this row, and the
discharge costs nothing beyond the containment of the scope in the
domain --- which is a compile guarantee, not a posted premise. When
the global row fails (the off-region was earning its keep), the move
is simply unavailable, and the kit\textquotesingle s working move,
\rulename{QVerifiedCover}, localizes the effort instead.

What the checker recomputes: nothing. Soundness is the monotonicity
of universal quantification --- a row asserted on a superset
restricts to any subset --- and the verified lemma is a one-line
restriction along an inclusion, provable without any classical
reasoning at all.

\cataloguenote{Two deliberate non-identifications. First, the alias\textquotesingle s
target is not a catalogue node: no region-weakening rule exists to
apply, and the unpacking discharges directly against the verified
monotonicity lemma --- the alias-tier twin of the admissible
tier\textquotesingle s recorded reduction-to-verified-machinery
boundary case. Second, the rule is deliberately NOT the base
calculus\textquotesingle s sequent weakening: that rule manages a
sequent\textquotesingle s premise list; this one manages the region a
row is quantified over, one extension layer up. Unifying them would
conflate two different quantifiers that merely rhyme.
}
\ruleblock{Lean statement}
\par\medskip
Stated over the compiled semantic class (sets over a bare state
type). The off-scoped shape the discharge plan actually emits is the
second declaration, an instance of the first; the scope-in-domain
containment enters as the bridge hypothesis the compile guarantees.

\declhead{BernsteinLogicAux.Structural}{fullDomainWeaken\_sound}{\badge{BadgeGreen}{statement bound}}
\begin{leansrc}
/-- **T3a FullDomainWeaken**, soundness. `scope` is the presentation-level
off-scoped row's scope; the bridge names the one fact the compile
guarantees about it: it lies inside the full domain. -/
theorem fullDomainWeaken_sound (Dm scope : Set X) (P : X → Prop)
    (hbridge : scope ⊆ Dm)
    (hrow : ∀ x ∈ Dm, P x) :
    ∀ x ∈ scope, P x
\end{leansrc}
\declhead{BernsteinLogicAux.Structural}{fullDomainWeaken\_offScope}{\badge{BadgeGreen}{statement bound}}
\begin{leansrc}
/-- The off-scope instance: the scope compiled as `D_m` minus an off-region.
(The shape the discharge plan actually emits.) -/
theorem fullDomainWeaken_offScope (Dm Off : Set X) (P : X → Prop)
    (hrow : ∀ x ∈ Dm, P x) :
    ∀ x ∈ Dm \ Off, P x
\end{leansrc}
\ruleblock{Correspondence (Steele display $\leftrightarrow$ Lean)}
\badge{BadgeGreen}{reviewer-certified against the displayed texts}
\sourcenote{The dictionary below is a reviewer-certified statement of how the mechanized statements above interpret the displayed rule: display metavariable $\mapsto$ Lean binder, each display hypothesis $\mapsto$ the premise or side condition carrying it, the display conclusion $\mapsto$ the goal --- with any deliberate difference in width between display and mechanization named outright. Certification is recorded as a hash triple over (the displayed cell, the extracted statement texts, this dictionary\textquotesingle s own text); if any of the three changes, this block turns amber until a reviewer re-certifies. The correspondence claim itself remains editorial --- reviewer-grade, not machine-checked; its machine-checked successor types each soundness statement against the rule as a structured object, and supersedes this block per rule as it lands.}
The displayed premise $\vdash P$ on $D_m$ $\mapsto$ \texttt{hrow : ∀ x
∈ Dm, P x}; the compile fact $\mathrm{scope} \subseteq D_m$ ---
displayed as ``not a posted premise'' --- $\mapsto$ the System-bridge
hypothesis \texttt{hbridge : scope ⊆ Dm}, an explicit binder here
because lean-aux lacks the compile layer that guarantees it; the
conclusion $\mapsto$ \texttt{∀ x ∈ scope, P x}. The emitted shape
$\mathrm{scope} = D_m \smallsetminus \mathrm{Off}$ is bound separately
as \texttt{fullDomainWeaken\_offScope}, where the bridge discharges by
\texttt{Set.diff\_subset} --- zero hypotheses beyond the row.
Width note: \texttt{P : X → Prop} arbitrary (payload verbatim, as
displayed --- the certified statement cannot even see the payload,
which is the strongest possible form of ``payload verbatim, scope
replaced''); the root statement is axiom-free pure
$\forall$-monotonicity, exactly the alias tier\textquotesingle s
price.

\ruleblock{Lean verifier receipts}
\sourcenote{module \texttt{BernsteinLogicAux.Structural}; bound build \texttt{1222cdcc326f\ldots} (bound from receipt cb0c3820c6877f22, 2026-07-29T01:03:36Z). A receipt records that named checking layers accepted a build: (module, compiled-artifact SHA-256, environment SHA-256, checker layers), minted when the checking gate ran; the store receipts attesting the bound build are listed with the layers each run credited. Nothing here asks the reader to trust the store: every layer is an independently runnable program over the quoted sources --- rebuild the module and re-run the checkers to reproduce the verdict. Receipts attesting superseded builds credit no layer and are omitted from this rendering (the count of omissions, if any, is noted below).}
\begin{receipttable}
\receiptrow{0b7e68907f45b27f}{2026-07-28T23:48:36Z}{build+decorrelated}{\badge{BadgeGreen}{current}}
\receiptrow{e97d5614d4f08c7e}{2026-07-29T00:00:35Z}{build+decorrelated}{\badge{BadgeGreen}{current}}
\receiptrow{27de21f047e6b194}{2026-07-29T00:49:05Z}{build+replay+hardened+decorrelated}{\badge{BadgeGreen}{current}}
\receiptrow{cb0c3820c6877f22}{2026-07-29T01:03:36Z}{build+replay+hardened+decorrelated}{\badge{BadgeGreen}{current}}
\receiptrow{3530841119a9a607}{2026-07-31T22:00:22Z}{build+decorrelated}{\badge{BadgeGreen}{current}}
\receiptrow{b6cc23eb3e5b7f97}{2026-07-31T22:29:16Z}{build+decorrelated}{\badge{BadgeGreen}{current}}
\receiptrow{bf3dacc3ab7f20ea}{2026-08-02T23:31:31Z}{build+decorrelated}{\badge{BadgeGreen}{current}}
\receiptrow{384a8569d2005df7}{2026-08-03T16:27:35Z}{build+decorrelated}{\badge{BadgeGreen}{current}}
\receiptrow{717e07acf9acac8a}{2026-08-03T22:46:19Z}{build+decorrelated}{\badge{BadgeGreen}{current}}
\end{receipttable}
\sourcenote{11 further receipt(s) in the store attest superseded builds of this module and are omitted from this rendering.}
\layerline{\layeritem{build}{lake build (elaboration)}{\badge{BadgeGreen}{current}}\layeritem{replay}{leanchecker (official kernel replay)}{\badge{BadgeGreen}{current}}\layeritem{hardened}{safe\_verify (hardened screen)}{\badge{BadgeGreen}{current}}\layeritem{decorrelated}{lean4export + nanoda (decorrelated kernel)}{\badge{BadgeGreen}{current}}}
\sourcenote{\badge{BadgeGray}{comparator: none} no fixture in the three-implementation comparator harness (compiled / kernel-interpreted / reference-interpreted) exercises this rule.}
\ruleblock{HOL4 statement}
\absentbox{NO HOL STATEMENT YET}{The HOL4 development\textquotesingle s object is the base
calculus\textquotesingle s wire-format checker (theory
\texttt{bern\_catalogue} ports the catalogue check-function cone).
This rule is scope machinery of the extension
theory\textquotesingle s discharge kit: it postdates that port, and
no row for the move exists in the baked theory --- there is
nothing for a HOL statement to be about. The box becomes fillable
when a revision carries the move\textquotesingle s certificate rows
and the corresponding check is ported.
}
\ruleblock{HOL verifier receipts}
\absentbox{NONE}{no HOL verification receipt exists for this rule.}
\ruleblock{Proof}

\sourcenote{Genre: alias rule (tier 3) --- a respelling template. An
alias owes the book five artefacts: the verified 1$\to$1 unpack, the
filled-in application of the target, the payload-transport table, the
exactness identity, and the soundness corollary through the
target\textquotesingle s lemma. This alias is the
tier\textquotesingle s degenerate case --- its target is a verified
lemma rather than a catalogue node, because no region-weakening rule
exists to apply --- and the five artefacts below say so honestly
rather than inventing a node.}

\begin{brecipe}[the alias template, degenerate-target form]\label{rec:fdw-template}
An alias rule $A$ is warranted by a verified function carrying its
payload to the payload of a \emph{single} target application --- one
node $\to$ one node, nothing derived. Where the target is not a
catalogue node but a verified theorem, the transport carries the
payload into the theorem\textquotesingle s hypotheses instead, and the
soundness corollary is the theorem itself: the alias-tier twin of the
admissible tier\textquotesingle s recorded
reduction-to-verified-machinery boundary case.
\end{brecipe}

\begin{brecipe}[the filled-in application]\label{rec:fdw-filled}
\[
\inferrule*[Right=\rulename{QFullDomainWeaken}]{
  \mathcal{T} \rhd\ \vdash P \ \text{on } D_m
}{
  \vdash P \ \text{on } D_m \smallsetminus \mathrm{Off}
}
\qquad\rightsquigarrow\qquad
\texttt{fullDomainWeaken\_offScope}\bigl[\mathcal{T}\bigr]
\;\rhd\; \vdash P \ \text{on } D_m \smallsetminus \mathrm{Off}
\]
with the general form --- any scope the compile certifies inside the
domain --- carried by \texttt{fullDomainWeaken\_sound}: from $P$ on
$D_m$ and $\mathrm{scope} \subseteq D_m$, conclude $P$ on
$\mathrm{scope}$.
\end{brecipe}

\begin{brecipe}[payload transport]\label{rec:fdw-transport}
\begin{center}\small
\begin{tabular}{@{}ll@{}}
\textbf{off-scoped row} & \textbf{full-domain row}\\
\midrule
row payload $P$ & $P$ (verbatim)\\
scope $D_m \smallsetminus \mathrm{Off}$ & $D_m$ (the off-region dropped)\\
$\mathrm{scope} \subseteq D_m$ & compile guarantee (\texttt{hbridge}), not a posted premise\\
\end{tabular}
\end{center}
\end{brecipe}

\begin{bprop}[exactness identity]\label{prop:fdw-identity}
$D_m \smallsetminus \mathrm{Off} \subseteq D_m$; more generally, every
off-scoped row\textquotesingle s scope is contained in its
mode\textquotesingle s domain.
\end{bprop}
\begin{proof}
A set difference is contained in its left operand --- one line
(\texttt{Set.diff\_subset}). The general containment is the compile
fact the bridge hypothesis names, supplied by the presentation layer,
never proven per-instance.
\end{proof}

\begin{bprop}[soundness corollary]\label{prop:fdw-sound}
If $P$ holds everywhere on $D_m$, it holds everywhere on any
$\mathrm{scope} \subseteq D_m$ --- in particular on
$D_m \smallsetminus \mathrm{Off}$.
\end{bprop}
\begin{proof}
Restriction of a universally quantified statement along an inclusion:
for $x \in \mathrm{scope}$, the inclusion gives $x \in D_m$, and the
full-domain row gives $P(x)$. \texttt{fullDomainWeaken\_sound} is this
one line, and it is axiom-free --- no classical reasoning, no choice:
pure monotonicity of $\forall$.
\end{proof}

\begin{bremark}[what is spent]
Nothing, soundness-wise --- and possibly everything,
feasibility-wise: if $P$ fails anywhere in $D_m \cap \mathrm{Off}$
(points of the off-region outside $D_m$ never enter the row), the
full-domain row is simply unprovable and the move is unavailable. The
off-region existing at all signals that some row somewhere needed it;
this move is the discharge for the rows that did not, and its
availability is itself a diagnosis: the carve-out was free for this
row. When it fails, the kit\textquotesingle s working move
(\rulename{QVerifiedCover}) localizes the work instead of abandoning
the carve-out.
\end{bremark}

\begin{bremark}[deliberately not the sequent weakening]
The base calculus has a weakening rule; this is deliberately not it.
Sequent weakening grows a premise \emph{list} --- what a sequent
assumes; this move shrinks a quantification \emph{region} --- where a
row is asserted. The two live one extension layer apart and neither
subsumes the other: a base sequent has no scope field to weaken, and
this move never touches a premise list. The kinship is real, and the
identification would be a type error.
\end{bremark}

\begin{anchors}
Recipe \ref{rec:fdw-filled} unpack & \texttt{fullDomainWeaken\_offScope} (the emitted shape)\\
general form & \texttt{fullDomainWeaken\_sound} (axiom-free)\\
Proposition \ref{prop:fdw-identity} & \texttt{Set.diff\_subset}; the \texttt{hbridge} inclusion\\
region-layer monotonicity, generic & \texttt{regionSat\_mono} (same module)\\
\end{anchors}


\rulesection{61}{QVerifiedCover}{admissible rule}
\label{blcatrule-QVerifiedCover}
\ruleblock{The rule (Steele metanotation)}
\sourcenote{as realized in the mechanized development; source row: \texttt{\texttt{QTSL-rules.toml}} --- \emph{QVerifiedCover}. The metanotation statement is typeset by hand for notational quality; its mechanized ground truth is the Lean statement below, extracted from the sources at every build of this book and hash-bound. The rule \emph{roster} is machine-checked (every mechanized rule tag round-trips to exactly one catalogue row); the statement \emph{text} is editorial, checked by eye against the extracted Lean beside it.}
\begin{mathpar}
\inferrule*{
  \vdash P \ \text{on } C \ \text{(the cover row)} \\
  \bigl\{\vdash K_i \subseteq \mathrm{Off}\bigr\}_i \ \text{(per-component containments)} \\
  D_m \smallsetminus C \subseteq \textstyle\bigcup_i K_i \ \text{(residual decomposition)}
}{
  \vdash P \ \text{on } D_m \smallsetminus \mathrm{Off}
}
\end{mathpar}

\ruleflavor{the per-component containments discharge by exact Bernstein upper bounds, and the residual decomposition is recomputed by the checker. Dimension-free as realized (\texttt{verifiedCover\_components\_sound}, arbitrary component index); the cover bookkeeping $C \subseteq D_m$ is checker-side well-formedness, not consumed by soundness.}
\par\medskip
\rulename{QVerifiedCover} is the discharge kit\textquotesingle s
working move: an off-scoped row --- $P$ on
$D_m \smallsetminus \mathrm{Off}$ --- is discharged by proving the
row on a cover $C$ and certifying that everything of $D_m$ the cover
misses lies inside the off-region. The un-covered residual is
decomposed into components $K_i$, and each component ships one
containment premise $K_i \subseteq \mathrm{Off}$, discharged by exact
Bernstein upper bounds --- ordinary sequent rows the base kernel
replays.

A producer reaches for it when \rulename{QFullDomainWeaken} fails ---
the row genuinely needs its carve-out --- but the on-region can be
covered cleanly. Cover \emph{selection} is producer-side search (the
current toolkit selects one-dimensional covers); the verified
statement is dimension-free, so richer covers need no new warrant.

What the checker recomputes: the residual decomposition
$D_m \smallsetminus C \subseteq \bigcup_i K_i$ (set bookkeeping over
the posted components) and the per-component containment rows. The
case analysis is exhaustive by construction: a point of the scope is
either covered, and the row-on-$C$ premise applies, or un-covered, in
which case it lies in some $K_i \subseteq \mathrm{Off}$ ---
contradicting membership in the scope. Drop the per-component
containments and an un-covered point outside the off-region escapes
both cases; drop the decomposition and the component list could
silently miss part of the residual.

\cataloguenote{One premise of the catalogued statement is deliberately absent from
the verified one: the cover bookkeeping $C \subseteq D_m$. The proof
never consumes it --- a cover reaching outside the domain proves more
rows than needed, harmlessly --- so the mechanized statement omits
it, which is strictly stronger in the sound direction; the checker
keeps the containment as cover well-formedness. The neighbouring
scope-preparation steps (hull closure, cell covers, off-scope
splitting, reset containment, floor boxes) are deliberately not rules
at all: they are side conditions the checker recomputes, and rows for
them would misstate that.
}
\ruleblock{Lean statement}
\par\medskip
Stated over the compiled semantic class (sets over a bare state
type), dimension-free. The residual form is the kernel statement; the
per-component form is the shape the discharge plan emits, with an
arbitrary component index --- nothing one-dimensional survives into
the warrant.

\declhead{BernsteinLogicAux.Structural}{verifiedCover\_sound}{\badge{BadgeGreen}{statement bound}}
\begin{leansrc}
/-- **T3b VerifiedCover**, soundness (residual form). `scope` is the
off-scoped row's compiled scope `D_m \ Off`; the row is proven on the cover
`C`, and the un-covered residual `D_m \ C` is certified inside the
off-region. -/
theorem verifiedCover_sound (Dm C Off scope : Set X) (P : X → Prop)
    (hbridge : scope = Dm \ Off)
    (hrowC : ∀ x ∈ C, P x)
    (hresidual : Dm \ C ⊆ Off) :
    ∀ x ∈ scope, P x
\end{leansrc}
\declhead{BernsteinLogicAux.Structural}{verifiedCover\_components\_sound}{\badge{BadgeGreen}{statement bound}}
\begin{leansrc}
/-- **T3b VerifiedCover**, per-component form: the residual is decomposed
into components `K i`, each separately certified inside the off-region (the
"exact Bernstein upper bounds" premises, one per component). -/
theorem verifiedCover_components_sound {ι : Type*} (Dm C Off scope : Set X)
    (K : ι → Set X) (P : X → Prop)
    (hbridge : scope = Dm \ Off)
    (hrowC : ∀ x ∈ C, P x)
    (hdecomp : Dm \ C ⊆ ⋃ i, K i)
    (hcomp : ∀ i, K i ⊆ Off) :
    ∀ x ∈ scope, P x
\end{leansrc}
\ruleblock{Correspondence (Steele display $\leftrightarrow$ Lean)}
\badge{BadgeGreen}{reviewer-certified against the displayed texts}
\sourcenote{The dictionary below is a reviewer-certified statement of how the mechanized statements above interpret the displayed rule: display metavariable $\mapsto$ Lean binder, each display hypothesis $\mapsto$ the premise or side condition carrying it, the display conclusion $\mapsto$ the goal --- with any deliberate difference in width between display and mechanization named outright. Certification is recorded as a hash triple over (the displayed cell, the extracted statement texts, this dictionary\textquotesingle s own text); if any of the three changes, this block turns amber until a reviewer re-certifies. The correspondence claim itself remains editorial --- reviewer-grade, not machine-checked; its machine-checked successor types each soundness statement against the rule as a structured object, and supersedes this block per rule as it lands.}
The displayed premises map onto \texttt{verifiedCover\_components\_sound}
one for one: the cover row $\vdash P$ on $C$ $\mapsto$ \texttt{hrowC};
the per-component containments $K_i \subseteq \mathrm{Off}$ $\mapsto$
\texttt{hcomp : ∀ i, K i ⊆ Off} (the exact-Bernstein-upper-bound
premises, one per component, over an \emph{arbitrary} index type ---
the dimension-freedom the flavor text claims is
a binder fact); the recomputed residual decomposition $D_m
\smallsetminus C \subseteq \bigcup_i K_i$ $\mapsto$ \texttt{hdecomp},
an explicit hypothesis (recompute rendered as binder, the
\rulename{QStreettTargetCellCollapse} pattern); the conclusion scope
$D_m \smallsetminus \mathrm{Off}$ $\mapsto$ the System-bridge equality
\texttt{hbridge : scope = Dm \textbackslash{} Off}. The residual form
\texttt{verifiedCover\_sound} is the same statement with the
decomposition pre-collapsed. Width note, divergence D1 as certified:
the catalogued well-formedness $C \subseteq D_m$ appears in NO binder
--- soundness demonstrably never consumes it, which the module header
banks as the divergence and the flavor text
carries; payload \texttt{P} arbitrary, per the QFullDomainWeaken pattern.

\ruleblock{Lean verifier receipts}
\sourcenote{module \texttt{BernsteinLogicAux.Structural}; bound build \texttt{1222cdcc326f\ldots} (bound from receipt cb0c3820c6877f22, 2026-07-29T01:03:36Z). A receipt records that named checking layers accepted a build: (module, compiled-artifact SHA-256, environment SHA-256, checker layers), minted when the checking gate ran; the store receipts attesting the bound build are listed with the layers each run credited. Nothing here asks the reader to trust the store: every layer is an independently runnable program over the quoted sources --- rebuild the module and re-run the checkers to reproduce the verdict. Receipts attesting superseded builds credit no layer and are omitted from this rendering (the count of omissions, if any, is noted below).}
\begin{receipttable}
\receiptrow{0b7e68907f45b27f}{2026-07-28T23:48:36Z}{build+decorrelated}{\badge{BadgeGreen}{current}}
\receiptrow{e97d5614d4f08c7e}{2026-07-29T00:00:35Z}{build+decorrelated}{\badge{BadgeGreen}{current}}
\receiptrow{27de21f047e6b194}{2026-07-29T00:49:05Z}{build+replay+hardened+decorrelated}{\badge{BadgeGreen}{current}}
\receiptrow{cb0c3820c6877f22}{2026-07-29T01:03:36Z}{build+replay+hardened+decorrelated}{\badge{BadgeGreen}{current}}
\receiptrow{3530841119a9a607}{2026-07-31T22:00:22Z}{build+decorrelated}{\badge{BadgeGreen}{current}}
\receiptrow{b6cc23eb3e5b7f97}{2026-07-31T22:29:16Z}{build+decorrelated}{\badge{BadgeGreen}{current}}
\receiptrow{bf3dacc3ab7f20ea}{2026-08-02T23:31:31Z}{build+decorrelated}{\badge{BadgeGreen}{current}}
\receiptrow{384a8569d2005df7}{2026-08-03T16:27:35Z}{build+decorrelated}{\badge{BadgeGreen}{current}}
\receiptrow{717e07acf9acac8a}{2026-08-03T22:46:19Z}{build+decorrelated}{\badge{BadgeGreen}{current}}
\end{receipttable}
\sourcenote{11 further receipt(s) in the store attest superseded builds of this module and are omitted from this rendering.}
\layerline{\layeritem{build}{lake build (elaboration)}{\badge{BadgeGreen}{current}}\layeritem{replay}{leanchecker (official kernel replay)}{\badge{BadgeGreen}{current}}\layeritem{hardened}{safe\_verify (hardened screen)}{\badge{BadgeGreen}{current}}\layeritem{decorrelated}{lean4export + nanoda (decorrelated kernel)}{\badge{BadgeGreen}{current}}}
\sourcenote{\badge{BadgeGray}{comparator: none} no fixture in the three-implementation comparator harness (compiled / kernel-interpreted / reference-interpreted) exercises this rule.}
\ruleblock{HOL4 statement}
\absentbox{NO HOL STATEMENT YET}{The HOL4 development\textquotesingle s object is the base
calculus\textquotesingle s wire-format checker (theory
\texttt{bern\_catalogue} ports the catalogue check-function cone).
This rule is scope machinery of the extension
theory\textquotesingle s discharge kit: it postdates that port, and
no row for the move exists in the baked theory --- there is
nothing for a HOL statement to be about. The box becomes fillable
when the extension\textquotesingle s certificate rows land in a wire
format and the corresponding check is ported.
}
\ruleblock{HOL verifier receipts}
\absentbox{NONE}{no HOL verification receipt exists for this rule.}
\ruleblock{Proof}

\sourcenote{Genre: admissible rule (tier 2) --- a proof-tree-rewriting
recipe. The rule discharges an off-scoped row by an exhaustive case
analysis: prove the row on a cover, certify everything un-covered into
the off-region. Its premises are ordinary sequent rows; the case
analysis is the verified content.}

The admissible template, general form, is stated once for the whole
catalogue as Recipe~\ref{rec:norm-template} of the \rulename{Norm}
section; this section instantiates it.

\begin{brecipe}[the \rulename{QVerifiedCover} instance]\label{rec:vc-instance}
\[
\inferrule*[Right=\rulename{QVerifiedCover}]{
  \mathcal{T}_C \rhd\ \vdash P \ \text{on } C \\
  \bigl\{\, \mathcal{T}_i \rhd\ \vdash K_i \subseteq \mathrm{Off} \,\bigr\}_i \\
  D_m \smallsetminus C \subseteq \textstyle\bigcup_i K_i
}{
  \vdash P \ \text{on } D_m \smallsetminus \mathrm{Off}
}
\qquad\rightsquigarrow\qquad
\begin{array}{l}
\text{\upshape premise subtrees stand;}\\
\text{\texttt{verifiedCover\_components\_sound}}\\
\text{\upshape discharges the scoped row.}
\end{array}
\]
The per-component containments $K_i \subseteq \mathrm{Off}$ are exact
Bernstein upper-bound rows --- ordinary sequents the base kernel
replays; the residual decomposition is checker-recomputed set
bookkeeping over the posted components, not a proof obligation. The
node discharges against the verified case analysis below.
\end{brecipe}

\begin{bprop}[semantic warrant, residual form]\label{prop:vc-sound}
If $P$ holds everywhere on $C$ and
$D_m \smallsetminus C \subseteq \mathrm{Off}$, then $P$ holds
everywhere on $D_m \smallsetminus \mathrm{Off}$.
\end{bprop}
\begin{proof}
Fix $x \in D_m \smallsetminus \mathrm{Off}$ and split on membership in
the cover. If $x \in C$, the cover row gives $P(x)$. If $x \notin C$,
then $x \in D_m \smallsetminus C \subseteq \mathrm{Off}$ ---
contradicting $x \notin \mathrm{Off}$. The two cases are exhaustive,
so the second is vacuous and the first decides
(\texttt{verifiedCover\_sound}).
\end{proof}

\begin{bprop}[semantic warrant, per-component form]\label{prop:vc-comp}
If $P$ holds everywhere on $C$, the residual decomposes as
$D_m \smallsetminus C \subseteq \bigcup_i K_i$, and each component
satisfies $K_i \subseteq \mathrm{Off}$, then $P$ holds everywhere on
$D_m \smallsetminus \mathrm{Off}$.
\end{bprop}
\begin{proof}
The union of the components is contained in $\mathrm{Off}$ (each is),
so $D_m \smallsetminus C \subseteq \mathrm{Off}$ and
Proposition~\ref{prop:vc-sound} applies
(\texttt{verifiedCover\_components\_sound}). The component index is an
arbitrary type: no finiteness, no dimension, no ordering of components
is consumed.
\end{proof}

\begin{bremark}[each premise family is load-bearing]
Drop the cover row and any $x \in C \cap (D_m \smallsetminus
\mathrm{Off})$ with $P(x)$ false refutes the conclusion while both
set-bookkeeping premises hold. Drop the per-component containments
and an un-covered point outside the off-region escapes both cases:
$x \in D_m \smallsetminus C$ with
$x \notin \mathrm{Off}$ and $P(x)$ false refutes the conclusion while
the remaining premises hold. Drop the residual decomposition and the
component list can silently miss part of the residual --- the same
point, now lying in no $K_i$, escapes again. The exhaustiveness of the
case analysis is exactly the conjunction of the last two; the cover
row is what it decides to.
\end{bremark}

\begin{bremark}[the omitted premise --- deliberate, and in the sound direction]
The catalogued statement carries a cover-bookkeeping premise
$C \subseteq D_m$; the verified statement omits it, because the proof
never consumes it: a cover reaching outside the domain proves more
rows than needed, harmlessly. Omitting an un-consumed hypothesis makes
the mechanized statement strictly stronger, in the sound direction ---
the checker retains the containment as cover \emph{well-formedness},
which is where it belongs.
\end{bremark}

\begin{bremark}[the kit as a spectrum]
The corner $C = D_m$ (empty residual, no components) is
\rulename{QFullDomainWeaken}; the opposite corner $C = \varnothing$
demands the entire domain certified into the off-region, at which
point the scope is empty and the row vacuous. Between the corners, the
cover is the dial trading global proof effort against per-component
containment certificates --- and cover \emph{selection} is
producer-side search (one-dimensional in the current toolkit), spent
before anything ships. The verified statement is dimension-free:
richer covers change the search, never the warrant.
\end{bremark}

\begin{anchors}
Recipe \ref{rec:vc-instance} warrant & \texttt{verifiedCover\_components\_sound}\\
Proposition \ref{prop:vc-sound} & \texttt{verifiedCover\_sound} (the case split)\\
containment premises & exact Bernstein upper-bound rows (the \rulename{BernCone} section)\\
scope reading & the \texttt{hbridge} hypothesis ($\mathrm{scope} = D_m \smallsetminus \mathrm{Off}$)\\
degenerate corner & the \rulename{QFullDomainWeaken} section\\
\end{anchors}


\rulesection{62}{QPositiveCombinationRoot}{admissible rule}
\label{blcatrule-QPositiveCombinationRoot}
\ruleblock{The rule (Steele metanotation)}
\sourcenote{as realized in the mechanized development; source row: \texttt{\texttt{QTSL-rules.toml}} --- \emph{QPositiveCombinationRoot}. The metanotation statement is typeset by hand for notational quality; its mechanized ground truth is the Lean statement below, extracted from the sources at every build of this book and hash-bound. The rule \emph{roster} is machine-checked (every mechanized rule tag round-trips to exactly one catalogue row); the statement \emph{text} is editorial, checked by eye against the extracted Lean beside it.}
\begin{mathpar}
\inferrule*{
  \mathrm{child}_1 : \mathsf{QRegionSat}(\varphi_1, R_1) \\
  \mathrm{child}_2 : \mathsf{QRegionSat}(\varphi_2, R_2)
}{
  \mathcal{M},\ x_0 \in R_1 \cap R_2 \vDash \varphi_1 \wedge \varphi_2
}
\and
\inferrule*{
  \mathrm{child}_1 : \mathsf{QRegionSat}(\varphi_1, R_1) \\
  \mathrm{child}_2 : \mathsf{QRegionSat}(\varphi_2, R_2)
}{
  \mathcal{M},\ x_0 \in R_1 \cup R_2 \vDash \varphi_1 \vee \varphi_2
}
\end{mathpar}

\ruleflavor{$n$-ary forms $\bigcap_i R_i$ / $\bigcup_i R_i$ over a block family (\texttt{positiveCombination\{And,Or,Inter,Union\}\_sound}); positive connectives only — negation rides the normal-form and closure machinery, never this rule.}
\par\medskip
\rulename{QPositiveCombinationRoot} assembles a certificate for a root
boolean combination out of certificates for its blocks, at zero new
row cost: if the children conclude $\mathsf{QRegionSat}(\varphi_1,
R_1)$ and $\mathsf{QRegionSat}(\varphi_2, R_2)$, the conjunction is
certified on $R_1 \cap R_2$ and the disjunction on $R_1 \cup R_2$ ---
and the $n$-ary forms take $\bigcap_i R_i$ and $\bigcup_i R_i$ over a
whole block family. A producer uses it at the top of every
specification whose root is a positive boolean shape: each block is
certified by its own sub-derivation (a battery, an invariance
certificate, a nested combination), and the root costs two set
operations.

What the checker recomputes: the region intersection or union itself.
The semantic clauses --- that the conjunction denotes the
intersection of the denotations, the disjunction the union --- enter
as definitional bridges supplied by the specification layer, never as
posted premises.

The word \emph{positive} in the name is load-bearing.
Under-approximation is preserved by intersection and union: a region
inside each denotation is inside their intersection; a union of
regions inside denotations is inside the union of the denotations. It
is REVERSED by complement: the complement of an under-approximation
over-approximates --- the empty region under-approximates every
denotation, and its complement, the whole space, over-approximates
every complement. A negation passed through this rule would silently
flip the judgment kind\textquotesingle s one orientation guarantee.
Negation therefore never passes through it; it rides the normal-form
and closure machinery instead, where the direction swap is handled
deliberately.

\cataloguenote{Admissible, not alias: several children and two lemma directions ---
more than a single application. The node shape (a pair of binary
connective nodes versus one $n$-ary combination node) is an enactment
decision still open on the schema side; the verified lemmas cover
both forms, so neither choice waits on new mathematics.
}
\ruleblock{Lean statement}
\par\medskip
\texttt{RegionSat} is the compiled reading of the workhorse judgment
kind --- the region under-approximates the denotation --- quoted in
full because the four lemmas\textquotesingle{} meaning turns on it.
The connective denotations enter through bridge hypotheses
(\texttt{denAnd} $=$ intersection of the children\textquotesingle s
denotations, \texttt{denOr} $=$ union), supplied by the specification
layer at compile time.

\declhead{BernsteinLogicAux.Structural}{positiveCombinationAnd\_sound}{\badge{BadgeGreen}{statement bound}}
\begin{leansrc}
/-- **T11**, conjunction: intersection of under-approximations
under-approximates the conjunction. -/
theorem positiveCombinationAnd_sound (R₁ R₂ den₁ den₂ denAnd : Set X)
    (hbridge : denAnd = den₁ ∩ den₂)
    (h₁ : RegionSat R₁ den₁) (h₂ : RegionSat R₂ den₂) :
    RegionSat (R₁ ∩ R₂) denAnd
\end{leansrc}
\declhead{BernsteinLogicAux.Structural}{positiveCombinationOr\_sound}{\badge{BadgeGreen}{statement bound}}
\begin{leansrc}
/-- **T11**, disjunction: union of under-approximations under-approximates
the disjunction. -/
theorem positiveCombinationOr_sound (R₁ R₂ den₁ den₂ denOr : Set X)
    (hbridge : denOr = den₁ ∪ den₂)
    (h₁ : RegionSat R₁ den₁) (h₂ : RegionSat R₂ den₂) :
    RegionSat (R₁ ∪ R₂) denOr
\end{leansrc}
\declhead{BernsteinLogicAux.Structural}{positiveCombinationInter\_sound}{\badge{BadgeGreen}{statement bound}}
\begin{leansrc}
/-- **T11**, n-ary conjunction (the staged n-ary node shape). -/
theorem positiveCombinationInter_sound {ι : Type*} (R den : ι → Set X)
    (denAll : Set X)
    (hbridge : denAll = ⋂ i, den i)
    (h : ∀ i, RegionSat (R i) (den i)) :
    RegionSat (⋂ i, R i) denAll
\end{leansrc}
\declhead{BernsteinLogicAux.Structural}{positiveCombinationUnion\_sound}{\badge{BadgeGreen}{statement bound}}
\begin{leansrc}
/-- **T11**, n-ary disjunction. -/
theorem positiveCombinationUnion_sound {ι : Type*} (R den : ι → Set X)
    (denAll : Set X)
    (hbridge : denAll = ⋃ i, den i)
    (h : ∀ i, RegionSat (R i) (den i)) :
    RegionSat (⋃ i, R i) denAll
\end{leansrc}
\declhead{BernsteinLogicAux.Structural}{RegionSat}{\badge{BadgeGreen}{definition bound}}
\begin{leansrc}
/-- The satisfaction judgment's workhorse kind, compiled: region `R`
under-approximates denotation `den`. -/
def RegionSat (R den : Set X) : Prop := R ⊆ den
\end{leansrc}
\ruleblock{Correspondence (Steele display $\leftrightarrow$ Lean)}
\badge{BadgeGreen}{reviewer-certified against the displayed texts}
\sourcenote{The dictionary below is a reviewer-certified statement of how the mechanized statements above interpret the displayed rule: display metavariable $\mapsto$ Lean binder, each display hypothesis $\mapsto$ the premise or side condition carrying it, the display conclusion $\mapsto$ the goal --- with any deliberate difference in width between display and mechanization named outright. Certification is recorded as a hash triple over (the displayed cell, the extracted statement texts, this dictionary\textquotesingle s own text); if any of the three changes, this block turns amber until a reviewer re-certifies. The correspondence claim itself remains editorial --- reviewer-grade, not machine-checked; its machine-checked successor types each soundness statement against the rule as a structured object, and supersedes this block per rule as it lands.}
The judgment kind itself is bound by display: \texttt{RegionSat R den
= R ⊆ den} --- $\mathsf{QRegionSat}$ compiled, the formula appearing
only through its denotation (the System bridge, module \S0). The two
displayed binary rules $\mapsto$
\texttt{positiveCombinationAnd\_sound} /
\texttt{Or\_sound}: children $\mathsf{QRegionSat}(\varphi_i, R_i)$
$\mapsto$ \texttt{h₁}/\texttt{h₂ : RegionSat Rᵢ denᵢ}; conclusions at
$R_1 \cap R_2$ / $R_1 \cup R_2$ verbatim; and the semantic clauses
$[\![\varphi_1 \wedge \varphi_2]\!] = [\![\varphi_1]\!] \cap
[\![\varphi_2]\!]$ (resp.\ $\cup$) $\mapsto$ the bridge hypotheses
\texttt{hbridge : denAnd = den₁ ∩ den₂} --- supplied by the future
spec layer, never asserted in lean-aux, exactly as the roster witness
says. The $n$-ary forms (\texttt{Inter\_sound}/\texttt{Union\_sound},
arbitrary index) cover the baked node shapes (\texttt{QAndRegion}/
\texttt{QOrRegion}, each n-ary over a children list). Width note: pure
lattice monotonicity over arbitrary sets --- no formula grammar, no
positivity syntax; the ``positive connectives only'' restriction is
enforced by which bridges exist to supply (negation has no
intersection/union clause), not by any binder here --- refusal lives
in the spec layer, the certified statements are the two monotone
directions the display prices.

\ruleblock{Lean verifier receipts}
\sourcenote{module \texttt{BernsteinLogicAux.Structural}; bound build \texttt{1222cdcc326f\ldots} (bound from receipt cb0c3820c6877f22, 2026-07-29T01:03:36Z). A receipt records that named checking layers accepted a build: (module, compiled-artifact SHA-256, environment SHA-256, checker layers), minted when the checking gate ran; the store receipts attesting the bound build are listed with the layers each run credited. Nothing here asks the reader to trust the store: every layer is an independently runnable program over the quoted sources --- rebuild the module and re-run the checkers to reproduce the verdict. Receipts attesting superseded builds credit no layer and are omitted from this rendering (the count of omissions, if any, is noted below).}
\begin{receipttable}
\receiptrow{0b7e68907f45b27f}{2026-07-28T23:48:36Z}{build+decorrelated}{\badge{BadgeGreen}{current}}
\receiptrow{e97d5614d4f08c7e}{2026-07-29T00:00:35Z}{build+decorrelated}{\badge{BadgeGreen}{current}}
\receiptrow{27de21f047e6b194}{2026-07-29T00:49:05Z}{build+replay+hardened+decorrelated}{\badge{BadgeGreen}{current}}
\receiptrow{cb0c3820c6877f22}{2026-07-29T01:03:36Z}{build+replay+hardened+decorrelated}{\badge{BadgeGreen}{current}}
\receiptrow{3530841119a9a607}{2026-07-31T22:00:22Z}{build+decorrelated}{\badge{BadgeGreen}{current}}
\receiptrow{b6cc23eb3e5b7f97}{2026-07-31T22:29:16Z}{build+decorrelated}{\badge{BadgeGreen}{current}}
\receiptrow{bf3dacc3ab7f20ea}{2026-08-02T23:31:31Z}{build+decorrelated}{\badge{BadgeGreen}{current}}
\receiptrow{384a8569d2005df7}{2026-08-03T16:27:35Z}{build+decorrelated}{\badge{BadgeGreen}{current}}
\receiptrow{717e07acf9acac8a}{2026-08-03T22:46:19Z}{build+decorrelated}{\badge{BadgeGreen}{current}}
\end{receipttable}
\sourcenote{11 further receipt(s) in the store attest superseded builds of this module and are omitted from this rendering.}
\layerline{\layeritem{build}{lake build (elaboration)}{\badge{BadgeGreen}{current}}\layeritem{replay}{leanchecker (official kernel replay)}{\badge{BadgeGreen}{current}}\layeritem{hardened}{safe\_verify (hardened screen)}{\badge{BadgeGreen}{current}}\layeritem{decorrelated}{lean4export + nanoda (decorrelated kernel)}{\badge{BadgeGreen}{current}}}
\sourcenote{\badge{BadgeGray}{comparator: none} no fixture in the three-implementation comparator harness (compiled / kernel-interpreted / reference-interpreted) exercises this rule.}
\ruleblock{HOL4 statement}
\absentbox{NO HOL STATEMENT YET}{The HOL4 development\textquotesingle s object is the base
calculus\textquotesingle s wire-format checker (theory
\texttt{bern\_catalogue} ports the catalogue check-function cone).
This rule concludes a satisfaction judgment of the extension theory:
it postdates that port, and although the combination nodes are baked
in the extension theory (\texttt{QAndRegion}/\texttt{QOrRegion}), no
HOL development consumes the extension theory\textquotesingle s rows
--- there is no ported check for a HOL statement to be about. The box
becomes fillable when a checker consuming the extension\textquotesingle s
certificate rows is ported.
}
\ruleblock{HOL verifier receipts}
\absentbox{NONE}{no HOL verification receipt exists for this rule.}
\ruleblock{Proof}

\sourcenote{Genre: admissible rule (tier 2) --- a proof-tree-rewriting
recipe. Admissible and not alias by count: several children, and two
lemma directions (one per connective). The rule adds zero sequent
rows; its entire content is two one-line lattice-monotonicity facts
about under-approximation.}

The admissible template, general form, is stated once for the whole
catalogue as Recipe~\ref{rec:norm-template} of the \rulename{Norm}
section; this section instantiates it.

\begin{brecipe}[the \rulename{QPositiveCombinationRoot} instance]\label{rec:pcr-instance}
\begin{mathpar}
\inferrule*[Right=$\wedge$-form]{
  \mathcal{T}_1 \rhd\ \mathsf{QRegionSat}(\varphi_1, R_1) \\
  \mathcal{T}_2 \rhd\ \mathsf{QRegionSat}(\varphi_2, R_2)
}{
  \mathcal{M},\ x_0 \in R_1 \cap R_2 \vDash \varphi_1 \wedge \varphi_2
}
\and
\inferrule*[Right=$\vee$-form]{
  \mathcal{T}_1 \rhd\ \mathsf{QRegionSat}(\varphi_1, R_1) \\
  \mathcal{T}_2 \rhd\ \mathsf{QRegionSat}(\varphi_2, R_2)
}{
  \mathcal{M},\ x_0 \in R_1 \cup R_2 \vDash \varphi_1 \vee \varphi_2
}
\end{mathpar}
Each node rewrites to its children --- which stand unchanged --- plus
one application of the corresponding lattice-monotonicity warrant
(Propositions \ref{prop:pcr-and} and \ref{prop:pcr-or}); the $n$-ary
forms replace the pair by a family and the binary operation by
$\bigcap_i$ or $\bigcup_i$. No sequent row is added anywhere: the
whole rewrite happens at the judgment layer.
\end{brecipe}

\begin{bprop}[conjunction: intersection preserves under-approximation]\label{prop:pcr-and}
If $R_1$ under-approximates the denotation of $\varphi_1$ and $R_2$
that of $\varphi_2$, then $R_1 \cap R_2$ under-approximates the
denotation of $\varphi_1 \wedge \varphi_2$ (which is the intersection
of the two denotations, by the specification layer\textquotesingle s
definitional clause).
\end{bprop}
\begin{proof}
$x \in R_1 \cap R_2$ gives $x \in R_1 \subseteq
[\![\varphi_1]\!]$ and $x \in R_2 \subseteq [\![\varphi_2]\!]$, so
$x \in [\![\varphi_1]\!] \cap [\![\varphi_2]\!] =
[\![\varphi_1 \wedge \varphi_2]\!]$. One line
(\texttt{positiveCombinationAnd\_sound}); the final equality is the
\texttt{denAnd} bridge, supplied at compile time, never posted.
\end{proof}

\begin{bprop}[disjunction: union preserves under-approximation]\label{prop:pcr-or}
Under the same hypotheses, $R_1 \cup R_2$ under-approximates the
denotation of $\varphi_1 \vee \varphi_2$ (the union of the
denotations).
\end{bprop}
\begin{proof}
$x \in R_1 \cup R_2$ lies in some $R_j \subseteq [\![\varphi_j]\!]
\subseteq [\![\varphi_1]\!] \cup [\![\varphi_2]\!] =
[\![\varphi_1 \vee \varphi_2]\!]$. One line
(\texttt{positiveCombinationOr\_sound}), with the \texttt{denOr}
bridge closing it.
\end{proof}

\begin{bremark}[the $n$-ary forms, and the empty corners]
The $n$-ary lemmas (\texttt{positiveCombinationInter\_sound},
\texttt{positiveCombinationUnion\_sound}) quantify over an arbitrary
index type --- no finiteness is consumed --- so a single $n$-ary
combination node brings no new mathematical idea beyond the same
monotonicity principle; the $n$-ary lemmas themselves are its warrant
(an arbitrary-index family is not a fold of binary steps, which is why
they are separate statements). The
empty corners are sound instances rather than exclusions: an empty
conjunction denotes everything and $\bigcap_{i \in \varnothing} R_i$
is the whole space, which under-approximates it; an empty disjunction
denotes nothing and the empty union is empty, which
under-approximates that. Neither degenerate certificate claims
anything false, and each is worthless in its own honest way: the
disjunction corner certifies the empty region, the conjunction corner
certifies a tautology.
\end{bremark}

\begin{bremark}[the polarity fence, with the two-line counterexample]
Complement REVERSES the direction this rule preserves: from
$R \subseteq [\![\varphi]\!]$ nothing at all follows about
$\mathrm{cmpl}\,R$ under-approximating
$\mathrm{cmpl}\,[\![\varphi]\!]$ --- the containment flips. Sharpest
instance: $R = \varnothing$ under-approximates every denotation, and
its complement is the whole space, which under-approximates only the
trivial one. A negation routed through this rule would silently
convert the judgment kind\textquotesingle s under-approximation
guarantee into its opposite. The name \emph{positive} is that fence:
negation rides the normal-form rewrites and the closure discipline
(the \rulename{QClosureExact} section and its siblings),
where the direction swap is handled deliberately, one atom at a time.
\end{bremark}

\begin{bremark}[the bridge hypotheses]
The semantic clauses $[\![\varphi_1 \wedge \varphi_2]\!] =
[\![\varphi_1]\!] \cap [\![\varphi_2]\!]$ and its $\vee$/union twin
enter the verified statements as named bridge hypotheses
(\texttt{denAnd}, \texttt{denOr}), to be discharged definitionally by
the specification layer\textquotesingle s compile. The warrant
deliberately does not assert them: what the connectives denote is the
specification layer\textquotesingle s to define, and this
rule\textquotesingle s content is only that under-approximation
respects whatever lattice structure that definition induces.
\end{bremark}

\begin{anchors}
Recipe \ref{rec:pcr-instance}, $\wedge$-form & \texttt{positiveCombinationAnd\_sound}\\
Recipe \ref{rec:pcr-instance}, $\vee$-form & \texttt{positiveCombinationOr\_sound}\\
$n$-ary forms & \texttt{positiveCombinationInter\_sound}, \texttt{positiveCombinationUnion\_sound}\\
the judgment reading & \texttt{RegionSat} (quoted in full above)\\
the polarity fence\textquotesingle s other side & \texttt{compl\_over\_approx} (the \rulename{QClosureOutward} section)\\
\end{anchors}


\rulesection{63}{QClosureExact}{alias rule}
\label{blcatrule-QClosureExact}
\ruleblock{The rule (Steele metanotation)}
\sourcenote{as realized in the mechanized development; source row: \texttt{\texttt{QTSL-rules.toml}} --- \emph{QClosureExact}. The metanotation statement is typeset by hand for notational quality; its mechanized ground truth is the Lean statement below, extracted from the sources at every build of this book and hash-bound. The rule \emph{roster} is machine-checked (every mechanized rule tag round-trips to exactly one catalogue row); the statement \emph{text} is editorial, checked by eye against the extracted Lean beside it.}
\begin{mathpar}
\inferrule*{ }{ \mathrm{cmpl}\{p < c\} = \{c \le p\} }
\and
\inferrule*{ }{ \mathrm{cmpl}\{c < p\} = \{p \le c\} }
\end{mathpar}

\ruleflavor{two axioms --- deduction rules with nothing above the line, each an exact set identity, both orientations stated (\texttt{closureExactLt\_sound}/\texttt{Gt}) — no margin spent, no side condition; the closed form is a closed set for continuous observables (\texttt{closureExact\_isClosed\_ge}/\texttt{le}).}
\par\medskip
The closure discipline is the calculus\textquotesingle s answer to
strict inequalities: rows quantify over closed cells and certificates
consume closed atoms, so strict literals --- $\{p < c\}$, $\{c < p\}$
--- never reach a row. Three moves resolve them, keyed to where the
literal sits. \rulename{QClosureExact}, this rule, is the
free one: the \emph{complement} of a strict literal is already a
closed atom, exactly --- $\mathrm{cmpl}\{p < c\} = \{c \le p\}$ ---
so a complemented strict literal is replaced by its closed form with
no loss, no margin, and no side condition. The sibling moves pay:
\rulename{QClosureOutward} closes the complement of a
\emph{closed} atom by an orientation swap that over-approximates (the
sound direction at negative positions), and
\rulename{QClosureInwardEps} closes a strict literal at a
positive position by rounding inward a producer-chosen
$\varepsilon > 0$ (the sound direction there). This rule is the alias
of the trio; the other two are admissible.

What the checker recomputes: nothing --- the rewrite is a
definitional set identity, stated in both orientations because both
strict forms occur in the atom grammar. The closedness that names the
move is a fact about continuous observables: the closed form is a
preimage of a closed ray, and the calculus\textquotesingle s
observables are polynomials.

\cataloguenote{The identity spends no margin --- contrast the sibling moves, which
each pay in one currency (a boundary cut, or an $\varepsilon$-band).
Like \rulename{QFullDomainWeaken}, the alias\textquotesingle s target
is a verified identity rather than a catalogue node. The
$\varepsilon$-accounting for the inward move travels as producer
annotation data on the certificate, never as a rule.
}
\ruleblock{Lean statement}
\par\medskip
The strict atom enters through the named bridge hypothesis
(\texttt{hbridge}) fixing its compiled sublevel reading; the
closedness lemmas ask only continuity of the observable.

\declhead{BernsteinLogicAux.Structural}{closureExactLt\_sound}{\badge{BadgeGreen}{statement bound}}
\begin{leansrc}
/-- **T13 EXACT**, `<`-orientation: the complement of `{p < c}` IS the
closed atom `{c ≤ p}` — no rounding, no margin spent. -/
theorem closureExactLt_sound (p : X → ℝ) (c : ℝ) (strictAtom : Set X)
    (hbridge : strictAtom = {x | p x < c}) :
    strictAtomᶜ = {x | c ≤ p x}
\end{leansrc}
\declhead{BernsteinLogicAux.Structural}{closureExactGt\_sound}{\badge{BadgeGreen}{statement bound}}
\begin{leansrc}
/-- **T13 EXACT**, `>`-orientation. -/
theorem closureExactGt_sound (p : X → ℝ) (c : ℝ) (strictAtom : Set X)
    (hbridge : strictAtom = {x | c < p x}) :
    strictAtomᶜ = {x | p x ≤ c}
\end{leansrc}
\declhead{BernsteinLogicAux.Structural}{closureExact\_isClosed\_ge}{\badge{BadgeGreen}{statement bound}}
\begin{leansrc}
/-- The closedness that names the move: `{c ≤ p}` is closed for continuous
`p`. -/
theorem closureExact_isClosed_ge [TopologicalSpace X] {p : X → ℝ}
    (hp : Continuous p) (c : ℝ) : IsClosed {x | c ≤ p x}
\end{leansrc}
\declhead{BernsteinLogicAux.Structural}{closureExact\_isClosed\_le}{\badge{BadgeGreen}{statement bound}}
\begin{leansrc}
/-- Closedness, the other orientation: `{p ≤ c}` is closed for continuous
`p`. -/
theorem closureExact_isClosed_le [TopologicalSpace X] {p : X → ℝ}
    (hp : Continuous p) (c : ℝ) : IsClosed {x | p x ≤ c}
\end{leansrc}
\ruleblock{Correspondence (Steele display $\leftrightarrow$ Lean)}
\badge{BadgeGreen}{reviewer-certified against the displayed texts}
\sourcenote{The dictionary below is a reviewer-certified statement of how the mechanized statements above interpret the displayed rule: display metavariable $\mapsto$ Lean binder, each display hypothesis $\mapsto$ the premise or side condition carrying it, the display conclusion $\mapsto$ the goal --- with any deliberate difference in width between display and mechanization named outright. Certification is recorded as a hash triple over (the displayed cell, the extracted statement texts, this dictionary\textquotesingle s own text); if any of the three changes, this block turns amber until a reviewer re-certifies. The correspondence claim itself remains editorial --- reviewer-grade, not machine-checked; its machine-checked successor types each soundness statement against the rule as a structured object, and supersedes this block per rule as it lands.}
Transcription-close, both orientations bound separately (the recorded
one-sidedness lesson as bind structure): $\mathrm{cmpl}\{p < c\} =
\{c \le p\}$ $\mapsto$ \texttt{closureExactLt\_sound}, whose goal is
the set identity \texttt{strictAtomᶜ = \{x | c ≤ p x\}} under the
System-bridge equality \texttt{hbridge : strictAtom = \{x | p x <
c\}}; the $>$-orientation $\mapsto$ \texttt{closureExactGt\_sound}.
The closedness that names the move $\mapsto$
\texttt{closureExact\_isClosed\_ge}/\texttt{\_le}, and their binder
row carries the one consumed regularity fact: \texttt{Continuous p}
--- the observable\textquotesingle s continuity, which for polynomial
atoms is free downstream but is named here, plus the
\texttt{[TopologicalSpace X]} instance those two lemmas alone need.
Width note: \texttt{p : X → ℝ} arbitrary --- the polynomial atom
grammar is erased, safe direction; the identity pair is margin-free
and side-condition-free exactly as displayed (no hypothesis beyond
the bridge), and no completeness claim exists to state --- an
identity loses nothing by construction.

\ruleblock{Lean verifier receipts}
\sourcenote{module \texttt{BernsteinLogicAux.Structural}; bound build \texttt{1222cdcc326f\ldots} (bound from receipt cb0c3820c6877f22, 2026-07-29T01:03:36Z). A receipt records that named checking layers accepted a build: (module, compiled-artifact SHA-256, environment SHA-256, checker layers), minted when the checking gate ran; the store receipts attesting the bound build are listed with the layers each run credited. Nothing here asks the reader to trust the store: every layer is an independently runnable program over the quoted sources --- rebuild the module and re-run the checkers to reproduce the verdict. Receipts attesting superseded builds credit no layer and are omitted from this rendering (the count of omissions, if any, is noted below).}
\begin{receipttable}
\receiptrow{0b7e68907f45b27f}{2026-07-28T23:48:36Z}{build+decorrelated}{\badge{BadgeGreen}{current}}
\receiptrow{e97d5614d4f08c7e}{2026-07-29T00:00:35Z}{build+decorrelated}{\badge{BadgeGreen}{current}}
\receiptrow{27de21f047e6b194}{2026-07-29T00:49:05Z}{build+replay+hardened+decorrelated}{\badge{BadgeGreen}{current}}
\receiptrow{cb0c3820c6877f22}{2026-07-29T01:03:36Z}{build+replay+hardened+decorrelated}{\badge{BadgeGreen}{current}}
\receiptrow{3530841119a9a607}{2026-07-31T22:00:22Z}{build+decorrelated}{\badge{BadgeGreen}{current}}
\receiptrow{b6cc23eb3e5b7f97}{2026-07-31T22:29:16Z}{build+decorrelated}{\badge{BadgeGreen}{current}}
\receiptrow{bf3dacc3ab7f20ea}{2026-08-02T23:31:31Z}{build+decorrelated}{\badge{BadgeGreen}{current}}
\receiptrow{384a8569d2005df7}{2026-08-03T16:27:35Z}{build+decorrelated}{\badge{BadgeGreen}{current}}
\receiptrow{717e07acf9acac8a}{2026-08-03T22:46:19Z}{build+decorrelated}{\badge{BadgeGreen}{current}}
\end{receipttable}
\sourcenote{11 further receipt(s) in the store attest superseded builds of this module and are omitted from this rendering.}
\layerline{\layeritem{build}{lake build (elaboration)}{\badge{BadgeGreen}{current}}\layeritem{replay}{leanchecker (official kernel replay)}{\badge{BadgeGreen}{current}}\layeritem{hardened}{safe\_verify (hardened screen)}{\badge{BadgeGreen}{current}}\layeritem{decorrelated}{lean4export + nanoda (decorrelated kernel)}{\badge{BadgeGreen}{current}}}
\sourcenote{\badge{BadgeGray}{comparator: none} no fixture in the three-implementation comparator harness (compiled / kernel-interpreted / reference-interpreted) exercises this rule.}
\ruleblock{HOL4 statement}
\absentbox{NO HOL STATEMENT YET}{The HOL4 development\textquotesingle s object is the base
calculus\textquotesingle s wire-format checker (theory
\texttt{bern\_catalogue} ports the catalogue check-function cone).
The closure discipline is a specification-surface rewrite layer of
the extension theory: it postdates that port, and no row for any
of the three moves exists in the baked theory --- there is nothing
for a HOL statement to be about. The box becomes fillable when a
revision carries the closure moves\textquotesingle{} certificate rows
and the corresponding check is ported.
}
\ruleblock{HOL verifier receipts}
\absentbox{NONE}{no HOL verification receipt exists for this rule.}
\ruleblock{Proof}

\sourcenote{Genre: alias rule (tier 3) --- a respelling template, here
in the degenerate-target form: the target is a verified set identity,
not a catalogue node. This section also carries the closure
discipline\textquotesingle s shared preamble; the
\rulename{QClosureOutward} and
\rulename{QClosureInwardEps} sections build on it by
reference.}

\begin{brecipe}[the closure discipline --- shared preamble for the trio]
\label{rec:cd-preamble}
The kernel\textquotesingle s rows quantify over closed cells; its
certificates consume closed atoms; its exact arithmetic decides
non-strict comparisons. Strict literals therefore never reach a row:
every strict atom is resolved, before rows are emitted, by one of
three moves keyed to the literal\textquotesingle s position.
\begin{center}\small
\begin{tabular}{@{}lll@{}}
\textbf{literal} & \textbf{move} & \textbf{cost}\\
\midrule
complement of a strict atom & \rulename{QClosureExact}: replace by the closed form & none (set identity)\\
complement of a closed atom & \rulename{QClosureOutward}: swap orientation, over-approximate & the cut $\{p = c\}$\\
strict atom, positive position & \rulename{QClosureInwardEps}: round inward by $\varepsilon > 0$ & the band $(c - \varepsilon, c)$; mirror $(c, c + \varepsilon)$\\
\end{tabular}
\end{center}
The polarity logic: at a \emph{positive} position --- the atom
contributes to what the certified region must lie inside ---
under-approximation is sound, so shrinking the atom (the inward move)
is the safe direction. At a \emph{negative} position --- the atom
arrives complemented --- the consuming side wants the complement
over-approximated, so growing it by the boundary cut (the outward
move) is the safe direction. The exact move is direction-free: an
identity is sound everywhere. Soundness-only is the
discipline\textquotesingle s contract: no move claims to preserve
provability, and a specification tight at a critical threshold fails
honestly --- the certificate is refused --- rather than being
silently widened.
\end{brecipe}

\begin{brecipe}[the \rulename{QClosureExact} instance]
\label{rec:ce-instance}
\[
\mathrm{cmpl}\{p < c\} \;\rightsquigarrow\; \{c \le p\}
\qquad\qquad
\mathrm{cmpl}\{c < p\} \;\rightsquigarrow\; \{p \le c\}
\]
Payload transport: observable $p$ and threshold $c$ verbatim;
orientation swapped; strictness dropped; the complement consumed. One
literal $\to$ one closed atom, and the respelling is \emph{exact}:
Proposition~\ref{prop:ce-identity} is an equality, not a containment,
so the rewrite is sound at every position, positive or negative ---
the only member of the trio with that property.
\end{brecipe}

\begin{bprop}[the exactness identity]\label{prop:ce-identity}
$\mathrm{cmpl}\{x \mid p(x) < c\} = \{x \mid c \le p(x)\}$, and
symmetrically
$\mathrm{cmpl}\{x \mid c < p(x)\} = \{x \mid p(x) \le c\}$.
\end{bprop}
\begin{proof}
Pointwise, $\neg\bigl(p(x) < c\bigr) \iff c \le p(x)$ --- totality of
the order on $\R$, nothing more.
\texttt{closureExactLt\_sound} and \texttt{closureExactGt\_sound} are
this equivalence under one set-extensionality step; there is no
analytic content anywhere in the move.
\end{proof}

\begin{bprop}[the closedness that names the move]\label{prop:ce-closed}
For continuous $p$, the atoms $\{c \le p\}$ and $\{p \le c\}$ are
closed sets.
\end{bprop}
\begin{proof}
Each is the preimage under $p$ of a closed ray
(\texttt{closureExact\_isClosed\_ge}, \texttt{\_le}). The
calculus\textquotesingle s observables are polynomials, hence
continuous; the lemmas ask continuity only, so any continuous
extension of the observable grammar inherits the move unchanged.
\end{proof}

\begin{bremark}[both orientations, deliberately]
The identity is stated twice because both strict forms occur in the
atom grammar, and a one-sided statement of a complement lemma is a
standing trap for author and porter alike. Stating both
orientations costs two lines and removes the trap.
\end{bremark}

\begin{bremark}[soundness corollary, and what the siblings inherit]
Substituting equals for equals needs no direction analysis, so the
soundness corollary is immediate: any judgment consuming
$\mathrm{cmpl}\{p < c\}$ may consume $\{c \le p\}$ with identical
denotation, and no margin is spent. The sibling moves are
containments, not identities, and their soundness is
position-dependent --- which is exactly why the discipline routes
each literal to the move whose direction its position wants
(Recipe~\ref{rec:cd-preamble}).
\end{bremark}

\begin{anchors}
Proposition \ref{prop:ce-identity} & \texttt{closureExactLt\_sound}, \texttt{closureExactGt\_sound}\\
Proposition \ref{prop:ce-closed} & \texttt{closureExact\_isClosed\_ge}, \texttt{closureExact\_isClosed\_le}\\
strict-atom reading & the \texttt{hbridge} hypothesis of the identity lemmas\\
the sibling moves & the \rulename{QClosureOutward} and \rulename{QClosureInwardEps} sections\\
\end{anchors}


\rulesection{64}{QClosureOutward}{admissible rule}
\label{blcatrule-QClosureOutward}
\ruleblock{The rule (Steele metanotation)}
\sourcenote{as realized in the mechanized development; source row: \texttt{\texttt{QTSL-rules.toml}} --- \emph{QClosureOutward}. The metanotation statement is typeset by hand for notational quality; its mechanized ground truth is the Lean statement below, extracted from the sources at every build of this book and hash-bound. The rule \emph{roster} is machine-checked (every mechanized rule tag round-trips to exactly one catalogue row); the statement \emph{text} is editorial, checked by eye against the extracted Lean beside it.}
\begin{mathpar}
\inferrule*{ }{ \mathrm{cmpl}\{p \le c\} \subseteq \{c \le p\} }
\and
\inferrule*{ }{ \mathrm{cmpl}\{c \le p\} \subseteq \{p \le c\} }
\and
\inferrule*{
  R \subseteq [\![\varphi]\!]
}{
  \mathrm{cmpl}\,[\![\varphi]\!] \subseteq \mathrm{cmpl}\,R
}
\end{mathpar}

\ruleflavor{the closed swap over-approximates the true complement, losing only the cut $\{p = c\}$ — the direction negative positions consume; atom face both orientations (\texttt{closureOutward\_sound}/\texttt{Ge}) $+$ the antitone region face (\texttt{compl\_over\_approx}).}
\par\medskip
\rulename{QClosureOutward} is the closure
discipline\textquotesingle s move for complements of \emph{closed}
atoms (the discipline itself is set out in the
\rulename{QClosureExact} section). The true complement of
$\{p \le c\}$ is the strict set $\{c < p\}$ --- not consumable by a
row --- and the move consumes the orientation swap $\{c \le p\}$
instead, which contains it: an over-approximation of the true
complement, losing only the cut $\{p = c\}$. Over-approximation is
exactly what the consuming side wants at negative positions: there
the atom describes a set to be excluded or avoided, and a certificate
against a superset is a fortiori a certificate against the set.

The warrant has two faces, both verified. The \emph{atom} face is the
pointwise containment above, in both orientations. The \emph{region}
face is the antitone kernel: a certified region under-approximates
its denotation, so its complement over-approximates the
denotation\textquotesingle s complement --- the direction transport
that lets a negative position consume a child
certificate\textquotesingle s complement soundly. What the checker
recomputes: the orientation swap itself, which is syntactic; the
containments are the verified lemmas\textquotesingle{} content, spent
once, globally.

The move is one-directional by design: at a \emph{positive} position
the same swap would overclaim --- take $p$ constant at $c$: the swap
$\{c \le p\}$ is the whole space while the true complement of
$\{p \le c\}$ is empty, the maximal over-claim --- so the
discipline\textquotesingle s routing table, not a side condition,
keeps this move on the negative side.

\cataloguenote{The conditional flag this entry once carried recorded the composite
normal-form induction that pushes complements through connectives ---
deliberately outside this rule\textquotesingle s own lemmas, since
that induction is stated once, for all consumers, with the
normal-form rewriting rules. It has since been mechanized where it
lives, so the row carries no conditional flag: \texttt{denot\_nnf}
(\texttt{BernsteinLogicAux.NNFDuality}, bound at the
\rulename{QStreettRequestNnf} entry) is the denotation-preserving
normal-form pass whose induction cases are exactly the
\rulename{QNnfRewrite} axioms, with \texttt{NForm.elimNeg} /
\texttt{denot\_elimNeg} the consumption-boundary elimination. A reader
tracing a complement through a full specification follows this rule
for the atom step and the normal-form sections for the traversal; the
composition-as-applied stays producer-side, which is why the traversal
is cited here, not bound.
}
\ruleblock{Lean statement}
\par\medskip
\texttt{RegionSat} (quoted in full) is the compiled reading of the
workhorse judgment kind, consumed by the region face; the atom faces
fix the consumed closed swap through the bridge hypothesis
(\texttt{hbridge}).

\declhead{BernsteinLogicAux.Structural}{closureOutward\_sound}{\badge{BadgeGreen}{statement bound}}
\begin{leansrc}
/-- **T13 OUTWARD**, `≤`-atom: the consumed closed swap `{c ≤ p}` contains
the true complement of `{p ≤ c}`. No margin is spent (the containment loses
only the measure-zero cut `{p = c}`). -/
theorem closureOutward_sound (p : X → ℝ) (c : ℝ) (cmplAtom : Set X)
    (hbridge : cmplAtom = {x | c ≤ p x}) :
    {x | p x ≤ c}ᶜ ⊆ cmplAtom
\end{leansrc}
\declhead{BernsteinLogicAux.Structural}{closureOutwardGe\_sound}{\badge{BadgeGreen}{statement bound}}
\begin{leansrc}
/-- **T13 OUTWARD**, `≥`-atom orientation. -/
theorem closureOutwardGe_sound (p : X → ℝ) (c : ℝ) (cmplAtom : Set X)
    (hbridge : cmplAtom = {x | p x ≤ c}) :
    {x | c ≤ p x}ᶜ ⊆ cmplAtom
\end{leansrc}
\declhead{BernsteinLogicAux.Structural}{compl\_over\_approx}{\badge{BadgeGreen}{statement bound}}
\begin{leansrc}
/-- The antitone region-complement kernel (region face): an
under-approximating certified region's complement OVER-approximates the
semantic complement — the direction negative positions consume. -/
theorem compl_over_approx {R den : Set X} (h : RegionSat R den) :
    denᶜ ⊆ Rᶜ
\end{leansrc}
\declhead{BernsteinLogicAux.Structural}{RegionSat}{\badge{BadgeGreen}{definition bound}}
\begin{leansrc}
/-- The satisfaction judgment's workhorse kind, compiled: region `R`
under-approximates denotation `den`. -/
def RegionSat (R den : Set X) : Prop := R ⊆ den
\end{leansrc}
\ruleblock{Correspondence (Steele display $\leftrightarrow$ Lean)}
\badge{BadgeGreen}{reviewer-certified against the displayed texts}
\sourcenote{The dictionary below is a reviewer-certified statement of how the mechanized statements above interpret the displayed rule: display metavariable $\mapsto$ Lean binder, each display hypothesis $\mapsto$ the premise or side condition carrying it, the display conclusion $\mapsto$ the goal --- with any deliberate difference in width between display and mechanization named outright. Certification is recorded as a hash triple over (the displayed cell, the extracted statement texts, this dictionary\textquotesingle s own text); if any of the three changes, this block turns amber until a reviewer re-certifies. The correspondence claim itself remains editorial --- reviewer-grade, not machine-checked; its machine-checked successor types each soundness statement against the rule as a structured object, and supersedes this block per rule as it lands.}
Three displayed clauses, three bound statements. The atom face, both
orientations: $\mathrm{cmpl}\{p \le c\} \subseteq \{c \le p\}$
$\mapsto$ \texttt{closureOutward\_sound} (goal \texttt{\{x | p x ≤
c\}ᶜ ⊆ cmplAtom} under the bridge \texttt{cmplAtom = \{x | c ≤ p
x\}}), mirrored by \texttt{closureOutwardGe\_sound} --- containments,
not identities: the flavor text\textquotesingle s ``losing only the
cut $\{p = c\}$'' is
commentary on the gap, and the certified statement claims exactly the
over-approximation direction, nothing finer. The region face $R
\subseteq [\![\varphi]\!] \Rightarrow \mathrm{cmpl}[\![\varphi]\!]
\subseteq \mathrm{cmpl}\,R$ $\mapsto$ \texttt{compl\_over\_approx :
denᶜ ⊆ Rᶜ} from \texttt{RegionSat R den} --- the antitone kernel,
the region face of the closure discipline, two tokens of proof. Width
note: this row once carried a conditional mark for a composite
residue --- the De Morgan \emph{traversal} of a full specification ---
which is discharged: the traversal lives with the
QNnfRewrite/QStreettRequestNnf normal-form machinery, deliberately not
here, and is machine-checked there (\texttt{denot\_nnf} +
\texttt{denot\_elimNeg}, cross-cited in the note), so the row carries
no conditional flag and what remains is only that division of binding
responsibility. The certified atom + region faces are unconditional
set-level facts over arbitrary observables and regions; no topology,
no continuity, no polynomiality is consumed anywhere in this bound
set.

\ruleblock{Lean verifier receipts}
\sourcenote{module \texttt{BernsteinLogicAux.Structural}; bound build \texttt{1222cdcc326f\ldots} (bound from receipt cb0c3820c6877f22, 2026-07-29T01:03:36Z). A receipt records that named checking layers accepted a build: (module, compiled-artifact SHA-256, environment SHA-256, checker layers), minted when the checking gate ran; the store receipts attesting the bound build are listed with the layers each run credited. Nothing here asks the reader to trust the store: every layer is an independently runnable program over the quoted sources --- rebuild the module and re-run the checkers to reproduce the verdict. Receipts attesting superseded builds credit no layer and are omitted from this rendering (the count of omissions, if any, is noted below).}
\begin{receipttable}
\receiptrow{0b7e68907f45b27f}{2026-07-28T23:48:36Z}{build+decorrelated}{\badge{BadgeGreen}{current}}
\receiptrow{e97d5614d4f08c7e}{2026-07-29T00:00:35Z}{build+decorrelated}{\badge{BadgeGreen}{current}}
\receiptrow{27de21f047e6b194}{2026-07-29T00:49:05Z}{build+replay+hardened+decorrelated}{\badge{BadgeGreen}{current}}
\receiptrow{cb0c3820c6877f22}{2026-07-29T01:03:36Z}{build+replay+hardened+decorrelated}{\badge{BadgeGreen}{current}}
\receiptrow{3530841119a9a607}{2026-07-31T22:00:22Z}{build+decorrelated}{\badge{BadgeGreen}{current}}
\receiptrow{b6cc23eb3e5b7f97}{2026-07-31T22:29:16Z}{build+decorrelated}{\badge{BadgeGreen}{current}}
\receiptrow{bf3dacc3ab7f20ea}{2026-08-02T23:31:31Z}{build+decorrelated}{\badge{BadgeGreen}{current}}
\receiptrow{384a8569d2005df7}{2026-08-03T16:27:35Z}{build+decorrelated}{\badge{BadgeGreen}{current}}
\receiptrow{717e07acf9acac8a}{2026-08-03T22:46:19Z}{build+decorrelated}{\badge{BadgeGreen}{current}}
\end{receipttable}
\sourcenote{11 further receipt(s) in the store attest superseded builds of this module and are omitted from this rendering.}
\layerline{\layeritem{build}{lake build (elaboration)}{\badge{BadgeGreen}{current}}\layeritem{replay}{leanchecker (official kernel replay)}{\badge{BadgeGreen}{current}}\layeritem{hardened}{safe\_verify (hardened screen)}{\badge{BadgeGreen}{current}}\layeritem{decorrelated}{lean4export + nanoda (decorrelated kernel)}{\badge{BadgeGreen}{current}}}
\sourcenote{\badge{BadgeGray}{comparator: none} no fixture in the three-implementation comparator harness (compiled / kernel-interpreted / reference-interpreted) exercises this rule.}
\ruleblock{HOL4 statement}
\absentbox{NO HOL STATEMENT YET}{The HOL4 development\textquotesingle s object is the base
calculus\textquotesingle s wire-format checker (theory
\texttt{bern\_catalogue} ports the catalogue check-function cone).
The closure discipline is a specification-surface rewrite layer of
the extension theory: it postdates that port, and no row for any
of the three moves exists in the baked theory --- there is nothing
for a HOL statement to be about. The box becomes fillable when a
revision carries the closure moves\textquotesingle{} certificate rows
and the corresponding check is ported.
}
\ruleblock{HOL verifier receipts}
\absentbox{NONE}{no HOL verification receipt exists for this rule.}
\ruleblock{Proof}

\sourcenote{Genre: admissible rule (tier 2) --- a proof-tree-rewriting
recipe. The closure discipline this move belongs to --- why strict
literals never reach a row, and how the three moves divide the
positions --- is set out in full in the
\rulename{QClosureExact} section; this section states only
what is native to the outward move: the two verified faces of its
warrant, and the direction analysis that confines it to negative
positions.}

The admissible template, general form, is stated once for the whole
catalogue as Recipe~\ref{rec:norm-template} of the \rulename{Norm}
section; this section instantiates it.

\begin{brecipe}[the \rulename{QClosureOutward} instance --- two faces]
\label{rec:co-instance}
\emph{Atom face.} At a negative position, the complement of a closed
atom is consumed as its orientation swap:
\[
\mathrm{cmpl}\{p \le c\} \;\rightsquigarrow\; \{c \le p\}
\qquad\qquad
\mathrm{cmpl}\{c \le p\} \;\rightsquigarrow\; \{p \le c\}
\]
warranted by the containment
$\mathrm{cmpl}\{p \le c\} \subseteq \{c \le p\}$
(Proposition~\ref{prop:co-atom}): the consumed atom
\emph{over-approximates} the true complement, losing only the cut
$\{p = c\}$.

\emph{Region face.} A negative position consuming the complement of a
\emph{certified region} $R$ (a child certificate\textquotesingle s
conclusion, not a syntactic atom) is warranted by the antitone kernel
(Proposition~\ref{prop:co-region}): certified under-approximation of
the denotation becomes over-approximation of the
denotation\textquotesingle s complement. Both faces discharge against
their verified containment; the orientation swap itself is syntactic
and recomputed by the checker.
\end{brecipe}

\begin{bprop}[atom face]\label{prop:co-atom}
$\mathrm{cmpl}\{x \mid p(x) \le c\} \subseteq \{x \mid c \le p(x)\}$,
and symmetrically
$\mathrm{cmpl}\{x \mid c \le p(x)\} \subseteq \{x \mid p(x) \le c\}$.
\end{bprop}
\begin{proof}
If $\neg(p(x) \le c)$ then $c < p(x)$, hence $c \le p(x)$ --- the
strict fact weakened to its non-strict form
(\texttt{closureOutward\_sound}; the mirror is
\texttt{closureOutwardGe\_sound}). The two sides differ exactly by
the cut $\{p = c\}$, which the swap adjoins.
\end{proof}

\begin{bprop}[region face --- the antitone kernel]\label{prop:co-region}
If a certified region under-approximates a denotation,
$R \subseteq [\![\varphi]\!]$, then
$\mathrm{cmpl}\,[\![\varphi]\!] \subseteq \mathrm{cmpl}\,R$: the
complement of the certificate over-approximates the complement of the
denotation.
\end{bprop}
\begin{proof}
Contrapositive, one line: $x \notin \mathrm{cmpl}\,R$ means
$x \in R \subseteq [\![\varphi]\!]$, so
$x \notin \mathrm{cmpl}\,[\![\varphi]\!]$
(\texttt{compl\_over\_approx}). This is the direction transport that
lets negative positions consume child certificates at all: the
judgment kind fixes one orientation, and complementation flips it in
precisely the way the consuming side wants.
\end{proof}

\begin{bremark}[what is lost, and why the direction is right]
Only the cut $\{p = c\}$ --- and in the sound direction. At a
negative position the atom names a set to be excluded or avoided;
certifying against a superset of it is a fortiori certifying against
it. The move never claims the cut is negligible --- for a nondegenerate
\emph{polynomial} observable ($p - c$ not identically zero on a
full-dimensional cell) it is a
measure-zero level set, but soundness nowhere consumes that: the
containment alone carries the warrant, so even a degenerate observable
whose cut is fat (at the extreme $p \equiv c$, the whole cell) only
makes the certificate more conservative, never wrong.
\end{bremark}

\begin{bremark}[one-directional by design]
At a \emph{positive} position the same swap would overclaim. Take $p$
constant at $c$: the swap $\{c \le p\}$ is the whole space, while the
true complement of $\{p \le c\}$ is empty --- the maximal over-claim.
Nothing in the move\textquotesingle s own premises excludes this; the
closure discipline\textquotesingle s routing table
(the preamble, \rulename{QClosureExact} section) is the
guard --- the recipe\textquotesingle s ``at a negative position''
names that routing assignment, not a premise the move checks for
itself --- sending positive-position strict literals to the inward
move instead. A porter implementing the discipline ports the routing with
it, or re-learns this remark the hard way.
\end{bremark}

\begin{bremark}[what is deliberately not warranted here]
The verified faces are the atom step and the region step. The
composite induction --- pushing complements through conjunctions,
disjunctions, and nested operators until they reach atoms --- is
stated once, for all consumers, with the normal-form rewriting rules
--- this rule\textquotesingle s warrant is deliberately partial in
exactly that sense: a complement traversing a full specification
rides the normal-form sections for the traversal and this rule for
the atoms it arrives at.
\end{bremark}

\begin{anchors}
Proposition \ref{prop:co-atom} & \texttt{closureOutward\_sound}, \texttt{closureOutwardGe\_sound}\\
Proposition \ref{prop:co-region} & \texttt{compl\_over\_approx}\\
the judgment reading & \texttt{RegionSat} (quoted in full above)\\
the routing guard & the closure-discipline preamble (\rulename{QClosureExact} section)\\
the composite traversal & the normal-form rewriting sections (deliberately not this rule)\\
\end{anchors}


\rulesection{65}{QClosureInwardEps}{admissible rule}
\label{blcatrule-QClosureInwardEps}
\ruleblock{The rule (Steele metanotation)}
\sourcenote{as realized in the mechanized development; source row: \texttt{\texttt{QTSL-rules.toml}} --- \emph{QClosureInwardEps}. The metanotation statement is typeset by hand for notational quality; its mechanized ground truth is the Lean statement below, extracted from the sources at every build of this book and hash-bound. The rule \emph{roster} is machine-checked (every mechanized rule tag round-trips to exactly one catalogue row); the statement \emph{text} is editorial, checked by eye against the extracted Lean beside it.}
\begin{mathpar}
\inferrule*{
  \varepsilon \in \Q_{>0}
}{
  \{p \le c - \varepsilon\} \subseteq \{p < c\}
}
\and
\inferrule*{
  \varepsilon \in \Q_{>0}
}{
  \{c + \varepsilon \le p\} \subseteq \{c < p\}
}
\end{mathpar}

\ruleflavor{the producer-chosen $\varepsilon$ elaborates a strict atom at a positive position to its rounded closed form ($\{p < c\} \rightsquigarrow \{p \le c - \varepsilon\}$, $\{c < p\} \rightsquigarrow \{c + \varepsilon \le p\}$), and the displayed containments are why that is sound: the rounded closed atom strengthens the claim (\texttt{closureInwardEpsLt\_sound}/\texttt{Gt}, any real $\varepsilon > 0$ — rationality subsumed by cast); soundness-only, no completeness preservation claimed — the honest failure at a critical threshold is refusal. Closedness rider: \texttt{closureInwardEps\_isClosed}.}
\par\medskip
\rulename{QClosureInwardEps} is the closure
discipline\textquotesingle s move for strict literals at positive
positions (the discipline is set out in the
\rulename{QClosureExact} section): the literal $\{p < c\}$
is replaced by the closed atom $\{p \le c - \varepsilon\}$ for a
producer-chosen rational $\varepsilon > 0$. The rounded atom is
contained in the literal, so at a position where the certified region
must lie \emph{inside} the denotation, the substitution strengthens
the obligation --- sound by construction --- and hands the kernel a
closed atom it can consume.

The $\varepsilon$ is real spend: the band between $c - \varepsilon$
and $c$ is surrendered (between $c$ and $c + \varepsilon$ in the
mirror orientation). The choice of $\varepsilon$ lives in the
producer toolkit --- searchable, like a multiplier or a cover --- and
its accounting travels as annotation data on the certificate, not as
a rule. What the checker recomputes: nothing beyond the rewritten
atom; the containment is the verified lemma, spent once, globally.

The contract is soundness-only, and the honest failure mode is
refusal: a specification exactly tight at its threshold --- the
intended witness satisfies $p < c$ everywhere, but $p$ approaches $c$
on the region --- admits no $\varepsilon$, and the certificate is
refused rather than silently widened. No completeness preservation is
claimed, deliberately.

\cataloguenote{The verified statement takes any real $\varepsilon > 0$; the
producer\textquotesingle s $\varepsilon$ being rational is a
representation fact about certificates, subsumed by the cast ---
soundness for every positive real is a fortiori soundness for every
positive rational. The closedness of the rounded atom is an instance
of the exact move\textquotesingle s closedness lemma at the shifted
threshold.
}
\ruleblock{Lean statement}
\par\medskip
Both orientations of the rounding, plus the closedness rider; the
rounded atom enters through the bridge hypothesis (\texttt{hbridge})
fixing its compiled sublevel reading, and $\varepsilon$ enters as a
real with $0 < \varepsilon$ (rationality of the
producer\textquotesingle s choice is subsumed by the cast).

\declhead{BernsteinLogicAux.Structural}{closureInwardEpsLt\_sound}{\badge{BadgeGreen}{statement bound}}
\begin{leansrc}
/-- **T13 INWARD-EPS**, `<`-orientation: `{p ≤ c - ε} ⊆ {p < c}` for any
`ε > 0`. Substituting the rounded closed atom at a positive position
strengthens the certified claim. -/
theorem closureInwardEpsLt_sound (p : X → ℝ) (c ε : ℝ) (hε : 0 < ε)
    (roundedAtom : Set X)
    (hbridge : roundedAtom = {x | p x ≤ c - ε}) :
    roundedAtom ⊆ {x | p x < c}
\end{leansrc}
\declhead{BernsteinLogicAux.Structural}{closureInwardEpsGt\_sound}{\badge{BadgeGreen}{statement bound}}
\begin{leansrc}
/-- **T13 INWARD-EPS**, `>`-orientation: `{c + ε ≤ p} ⊆ {c < p}`. -/
theorem closureInwardEpsGt_sound (p : X → ℝ) (c ε : ℝ) (hε : 0 < ε)
    (roundedAtom : Set X)
    (hbridge : roundedAtom = {x | c + ε ≤ p x}) :
    roundedAtom ⊆ {x | c < p x}
\end{leansrc}
\declhead{BernsteinLogicAux.Structural}{closureInwardEps\_isClosed}{\badge{BadgeGreen}{statement bound}}
\begin{leansrc}
/-- The inward-rounding closedness rider: the rounded atom is closed
(instance of the EXACT move's closedness at the shifted threshold). -/
theorem closureInwardEps_isClosed [TopologicalSpace X] {p : X → ℝ}
    (hp : Continuous p) (c ε : ℝ) : IsClosed {x | p x ≤ c - ε}
\end{leansrc}
\ruleblock{Correspondence (Steele display $\leftrightarrow$ Lean)}
\badge{BadgeGreen}{reviewer-certified against the displayed texts}
\sourcenote{The dictionary below is a reviewer-certified statement of how the mechanized statements above interpret the displayed rule: display metavariable $\mapsto$ Lean binder, each display hypothesis $\mapsto$ the premise or side condition carrying it, the display conclusion $\mapsto$ the goal --- with any deliberate difference in width between display and mechanization named outright. Certification is recorded as a hash triple over (the displayed cell, the extracted statement texts, this dictionary\textquotesingle s own text); if any of the three changes, this block turns amber until a reviewer re-certifies. The correspondence claim itself remains editorial --- reviewer-grade, not machine-checked; its machine-checked successor types each soundness statement against the rule as a structured object, and supersedes this block per rule as it lands.}
Both displayed containments bound, orientation for orientation:
$\{p \le c - \varepsilon\}
\subseteq \{p < c\}$ $\mapsto$ \texttt{closureInwardEpsLt\_sound}
(the rounded atom entering through the bridge \texttt{roundedAtom =
\{x | p x ≤ c − ε\}}, the strict margin as \texttt{hε : 0 < ε}), and
the $>$-orientation $\mapsto$ \texttt{closureInwardEpsGt\_sound};
the inward-rounding closedness rider (the rounded atom is closed)
$\mapsto$
\texttt{closureInwardEps\_isClosed} (continuity consumed there
alone, the QClosureExact pattern). Two width facts, one per direction of
the display: the certified statements take \emph{any real}
$\varepsilon > 0$ where the display asks $\varepsilon \in \Q_{>0}$
--- producer rationality subsumed by cast, divergence D3, statement
wider in the safe direction; and no completeness-preservation
statement exists in the module --- the flavor text\textquotesingle s
``soundness-only'' is certified by absence, with total loss at a
critical cut remaining the honest refusal mode, exactly as the
ledger inherits it.

\ruleblock{Lean verifier receipts}
\sourcenote{module \texttt{BernsteinLogicAux.Structural}; bound build \texttt{1222cdcc326f\ldots} (bound from receipt cb0c3820c6877f22, 2026-07-29T01:03:36Z). A receipt records that named checking layers accepted a build: (module, compiled-artifact SHA-256, environment SHA-256, checker layers), minted when the checking gate ran; the store receipts attesting the bound build are listed with the layers each run credited. Nothing here asks the reader to trust the store: every layer is an independently runnable program over the quoted sources --- rebuild the module and re-run the checkers to reproduce the verdict. Receipts attesting superseded builds credit no layer and are omitted from this rendering (the count of omissions, if any, is noted below).}
\begin{receipttable}
\receiptrow{0b7e68907f45b27f}{2026-07-28T23:48:36Z}{build+decorrelated}{\badge{BadgeGreen}{current}}
\receiptrow{e97d5614d4f08c7e}{2026-07-29T00:00:35Z}{build+decorrelated}{\badge{BadgeGreen}{current}}
\receiptrow{27de21f047e6b194}{2026-07-29T00:49:05Z}{build+replay+hardened+decorrelated}{\badge{BadgeGreen}{current}}
\receiptrow{cb0c3820c6877f22}{2026-07-29T01:03:36Z}{build+replay+hardened+decorrelated}{\badge{BadgeGreen}{current}}
\receiptrow{3530841119a9a607}{2026-07-31T22:00:22Z}{build+decorrelated}{\badge{BadgeGreen}{current}}
\receiptrow{b6cc23eb3e5b7f97}{2026-07-31T22:29:16Z}{build+decorrelated}{\badge{BadgeGreen}{current}}
\receiptrow{bf3dacc3ab7f20ea}{2026-08-02T23:31:31Z}{build+decorrelated}{\badge{BadgeGreen}{current}}
\receiptrow{384a8569d2005df7}{2026-08-03T16:27:35Z}{build+decorrelated}{\badge{BadgeGreen}{current}}
\receiptrow{717e07acf9acac8a}{2026-08-03T22:46:19Z}{build+decorrelated}{\badge{BadgeGreen}{current}}
\end{receipttable}
\sourcenote{11 further receipt(s) in the store attest superseded builds of this module and are omitted from this rendering.}
\layerline{\layeritem{build}{lake build (elaboration)}{\badge{BadgeGreen}{current}}\layeritem{replay}{leanchecker (official kernel replay)}{\badge{BadgeGreen}{current}}\layeritem{hardened}{safe\_verify (hardened screen)}{\badge{BadgeGreen}{current}}\layeritem{decorrelated}{lean4export + nanoda (decorrelated kernel)}{\badge{BadgeGreen}{current}}}
\sourcenote{\badge{BadgeGray}{comparator: none} no fixture in the three-implementation comparator harness (compiled / kernel-interpreted / reference-interpreted) exercises this rule.}
\ruleblock{HOL4 statement}
\absentbox{NO HOL STATEMENT YET}{The HOL4 development\textquotesingle s object is the base
calculus\textquotesingle s wire-format checker (theory
\texttt{bern\_catalogue} ports the catalogue check-function cone).
The closure discipline is a specification-surface rewrite layer of
the extension theory: it postdates that port, and no row for any
of the three moves exists in the baked theory --- there is nothing
for a HOL statement to be about. The box becomes fillable when a
revision carries the closure moves\textquotesingle{} certificate rows
and the corresponding check is ported.
}
\ruleblock{HOL verifier receipts}
\absentbox{NONE}{no HOL verification receipt exists for this rule.}
\ruleblock{Proof}

\sourcenote{Genre: admissible rule (tier 2) --- a proof-tree-rewriting
recipe. The closure discipline this move belongs to is set out in full
in the \rulename{QClosureExact} section; this section states
what is native to the inward move: the rounding containment in both
orientations, the closedness rider, and the two honesty clauses ---
strictness of $\varepsilon > 0$, and the deliberate absence of any
completeness claim.}

The admissible template, general form, is stated once for the whole
catalogue as Recipe~\ref{rec:norm-template} of the \rulename{Norm}
section; this section instantiates it.

\begin{brecipe}[the \rulename{QClosureInwardEps} instance]
\label{rec:cie-instance}
At a positive position, a strict literal rounds inward by a
producer-chosen $\varepsilon > 0$:
\[
\{p < c\} \;\rightsquigarrow\; \{p \le c - \varepsilon\}
\qquad\qquad
\{c < p\} \;\rightsquigarrow\; \{c + \varepsilon \le p\}
\]
warranted by the containment
$\{p \le c - \varepsilon\} \subseteq \{p < c\}$
(Proposition~\ref{prop:cie-sound}): the rounded closed atom
\emph{strengthens} the certified claim, which is the sound direction
where the certified region must lie inside the denotation. Payload
transport: observable and threshold verbatim, threshold shifted by
$\varepsilon$, strictness dropped; the $\varepsilon$ itself travels
as certificate annotation data, never as a rule premise.
\end{brecipe}

\begin{bprop}[the rounding containment]\label{prop:cie-sound}
For every $\varepsilon > 0$:
$\{x \mid p(x) \le c - \varepsilon\} \subseteq \{x \mid p(x) < c\}$,
and symmetrically
$\{x \mid c + \varepsilon \le p(x)\} \subseteq \{x \mid c < p(x)\}$.
\end{bprop}
\begin{proof}
$p(x) \le c - \varepsilon < c$, the second inequality because
$\varepsilon > 0$ --- one line of linear arithmetic
(\texttt{closureInwardEpsLt\_sound}; mirror
\texttt{closureInwardEpsGt\_sound}). The verified statements take any
\emph{real} $\varepsilon > 0$; the producer\textquotesingle s
$\varepsilon$ being rational is a representation fact about
certificates, subsumed by the cast.
\end{proof}

\begin{bprop}[closedness rider]\label{prop:cie-closed}
For continuous $p$, both rounded atoms are closed sets:
$\{p \le c - \varepsilon\}$ and $\{c + \varepsilon \le p\}$.
\end{bprop}
\begin{proof}
An instance of the exact move\textquotesingle s closedness at the
shifted threshold $c - \varepsilon$
(\texttt{closureInwardEps\_isClosed}, one application of
\texttt{closureExact\_isClosed\_le}); the mirrored orientation is
$\{c + \varepsilon \le p\} = p^{-1}\bigl([\,c + \varepsilon,
\infty)\bigr)$, closed by the same continuity --- the exact
move\textquotesingle s $\ge$-orientation closedness at threshold
$c + \varepsilon$ (the cited binding covers the $\le$ orientation;
the mirror discharges against the $\ge$ twin). The move thus
delivers exactly what the discipline asked of it: a closed atom the
kernel can consume.
\end{proof}

\begin{bremark}[$\varepsilon > 0$ strictly --- the boundary counterexample]
At $\varepsilon = 0$ the containment fails: $\{p \le c\}$ is not
contained in $\{p < c\}$, and the failure can be total --- take $p$
constant at $c$: the un-rounded atom is the whole space while the
literal is empty. The strict positivity of $\varepsilon$ is therefore
the whole rule, in the same sense that the sign row is half of
\rulename{HypMultiplier}: remove it and the move manufactures the
maximal over-claim from nothing.
\end{bremark}

\begin{bremark}[soundness-only, and the refusal posture]
No completeness preservation is claimed, deliberately. The
$\varepsilon$-band is surrendered --- $(c - \varepsilon, c)$ in the
first orientation, $(c, c + \varepsilon)$ in the mirror --- and a
specification exactly tight at its threshold --- the intended witness
satisfies $p < c$ everywhere while $p$ approaches $c$ on the region
--- admits no $\varepsilon$ at all: every rounding excludes points the
witness needs. The honest outcome is refusal (no certificate exists
through this move), never silent widening; a producer meeting refusal
here either finds slack elsewhere in the specification or reports the
tightness, and the discipline\textquotesingle s contract is that the
checker is never the party doing the forgiving.
\end{bremark}

\begin{bremark}[why not under a complement]
Rounding an atom that will be complemented is doubly wrong as a
discharge of \emph{this} obligation: the complement of the rounded
atom, $\{c - \varepsilon < p\}$, is strict again --- the closedness
goal is defeated, and the $\varepsilon$ was spent for nothing.
Complemented literals belong to the exact move (free, when the inner
literal is strict) or the outward move (when the inner atom is
closed); the routing table in the closure-discipline preamble
(\rulename{QClosureExact} section) assigns every case
exactly one move, and the assignment is what makes the trio a
discipline rather than three tricks.
\end{bremark}

\begin{anchors}
Proposition \ref{prop:cie-sound} & \texttt{closureInwardEpsLt\_sound}, \texttt{closureInwardEpsGt\_sound}\\
Proposition \ref{prop:cie-closed} & \texttt{closureInwardEps\_isClosed} (via \texttt{closureExact\_isClosed\_le})\\
rounded-atom reading & the \texttt{hbridge} hypothesis of the rounding lemmas\\
$\varepsilon$ accounting & certificate annotation data, producer-side (not a rule)\\
the routing guard & the closure-discipline preamble (\rulename{QClosureExact} section)\\
\end{anchors}


\rulesection{66}{TaylorBound}{alias for \textcolor{DiffBlue}{\rulename{PConstExact}}}
\label{blcatrule-TaylorBound}
\ruleblock{The rule (Steele metanotation)}
\sourcenote{as realized in the mechanized development; source row: \texttt{\texttt{QTSL-rules.toml}} --- \emph{TaylorBound}. The metanotation statement is typeset by hand for notational quality; its mechanized ground truth is the Lean statement below, extracted from the sources at every build of this book and hash-bound. The rule \emph{roster} is machine-checked (every mechanized rule tag round-trips to exactly one catalogue row); the statement \emph{text} is editorial, checked by eye against the extracted Lean beside it.}
\begin{mathpar}
\inferrule*[Right=\rulename{TaylorBound}]{
  q \in \Q; \quad \theta \in \Q_{>0}; \quad T \in \N; \quad n \in \N \\
  q \le \textstyle\sum_{k \le n} (\theta T)^k / k! \ \text{holds under exact evaluation in } \Q
}{
  \Phi \vdash q \le \textstyle\sum_{k \le n} (\theta T)^k / k!
}
\end{mathpar}

\ruleflavor{one \rulename{PConstExact} application, displayed at the deadline battery\textquotesingle s Taylor comparison — every term is wire-free (closed), so the comparison is decided by the verified closed evaluator; no QTSL sort exists and none is needed. The transcendental fact $\sum_{k \le n} (\theta T)^k / k! \le e^{\theta T}$ belongs to the consuming lemma, never the checker: the checker never needs to know $e$ exists.}
\par\medskip
\rulename{TaylorBound} supplies the deadline battery\textquotesingle s
rational surrogate for a transcendental bound. The deadline
rule\textquotesingle s soundness argument wants $q \le e^{\theta T}$
for a producer-chosen rational $q$, with decay rate
$\theta \in \Q_{>0}$ and horizon $T \in \Q_{\ge 0}$ from its payload:
its concluded sublevel region $\{V \le (1-p)\,q\}$ grows with $q$, so
the producer wants the sharpest $q$ the arithmetic affords. No checker
doing exact rational arithmetic can evaluate $e^{\theta T}$, and none
needs to: the certificate posts the single wire-free comparison
$q \le \sum_{k \le n} (\theta T)^k / k!$ --- a finite sum of
rationals, checked in exact $\Q$ --- and the consuming lemma carries
the enclosure $\sum_{k \le n} (\theta T)^k / k! \le e^{\theta T}$
(nonnegative Taylor remainder at nonnegative argument), so
$q \le e^{\theta T}$ follows without the checker ever knowing $e$
exists. A producer picks the truncation order $n$ and any convenient
$q$ at or below the partial sum.

What the checker recomputes is the target rule\textquotesingle s
entire check: exact closed evaluation of one comparison. No
certificate ever carries a \rulename{TaylorBound} node --- no sort
exists and none is needed; the producer emits the ground-fact row
directly, and this row is the catalogue-level name for that
elaboration.

What would break without each hypothesis: with rounded rather than
exact evaluation, $q$ could creep above the true partial sum and the
surrogate above $e^{\theta T}$; without wire-freedom the
``constant'' could depend on state and the enclosure is meaningless;
and the sign condition $\theta T \ge 0$ is load-bearing in the
consuming lemma --- at negative arguments the Taylor remainder
alternates in sign and the enclosure direction fails at alternate
truncation orders (the proof section carries the counterexample).

\cataloguenote{Deliberately \emph{not} aliased onto the interval-enclosure rule:
that would import witness-bracket machinery for a metatheory
one-liner, and the enclosure rule today lives only in a separate
reference-checker implementation besides. The seam this rule
instantiates is a design invariant of the whole extension page:
transcendental facts live in soundness lemmas (the Lean-side deadline
lemma consumes the standard library\textquotesingle s
partial-sum-below-exponential fact), arithmetic lives in the checker.
Mechanization state, plainly: the bound statements below are the
enclosure and its exact-$\Q$ face. \texttt{taylorLeaf\_le\_exp}
(module \texttt{BernsteinLogicAux.QtslAliases}) is the enclosure step
--- $q$ at or below the truncated sum is below $e^{x}$ at nonnegative
argument, the nonnegative Taylor remainder --- and
\texttt{taylorLeaf\_rat\_le\_exp} is the row end to end: the wire-free
comparison exactly as the checker replays it, in $\Q$, with the
$(\theta, T)$ payload shape, concluding $q \le e^{\theta T}$. The
consuming site is \texttt{sound\_QDeadlineBattery}\textquotesingle s
$hq$ premise (the target\textquotesingle s entry), whose internal
enclosure step these statements name as a rule-level fact. The sign
condition $0 \le \theta T$ is a stated hypothesis, exactly as
load-bearing as the commentary\textquotesingle s counterexample says
it is. No profile enters; the entry carries no conditional flag. The
target\textquotesingle s comparator exclusion ($\neq$ inadmissible)
travels with the alias to any port.
}
\ruleblock{Lean statement}
\declhead{BernsteinLogicAux.QtslAliases}{taylorLeaf\_le\_exp}{\badge{BadgeGreen}{statement bound}}
\begin{leansrc}
/-- **The enclosure step** (QTSL-rules.toml row `TaylorBound (T7)`): a
producer-chosen `q` at or below the order-`n` Taylor partial sum of
`exp` at a nonnegative argument is below `exp` itself — the nonnegative
Taylor remainder (`Real.sum_le_exp_of_nonneg`). The sign condition
`0 ≤ x` is load-bearing: at negative arguments the remainder alternates
in sign, so the enclosure fails at alternate truncation orders (at
`x = -1` it fails at `n = 0`, since `1 > e⁻¹`, and holds at `n = 1`) —
there is no all-`n` enclosure to be had (counter-witnesses in the
book's proof section). -/
theorem taylorLeaf_le_exp {q x : ℝ} {n : ℕ} (hx : 0 ≤ x)
    (hq : q ≤ ∑ k ∈ range (n + 1), x ^ k / (Nat.factorial k : ℝ)) :
    q ≤ Real.exp x
\end{leansrc}
\declhead{BernsteinLogicAux.QtslAliases}{taylorLeaf\_rat\_le\_exp}{\badge{BadgeGreen}{statement bound}}
\begin{leansrc}
/-- **`TaylorBound`, the exact-ℚ face, end to end** (rule row
`TaylorBound (T7)` — the single wire-free closed comparison, checked in
exact ℚ, that stands in for `q ≤ e^{θT}`): decay rate `θ ∈ ℚ≥0` and
horizon `T ∈ ℕ` from the deadline payload, truncation order `n`
producer-chosen. The checker replays the hypothesis; the conclusion —
which mentions `e` — lives only here (the catalogue's "the checker never
needs to know `e` exists"). Consuming site: `sound_QDeadlineBattery`'s
`hq` premise. -/
theorem taylorLeaf_rat_le_exp {qq θ : ℚ} {T n : ℕ} (hθ : 0 ≤ θ)
    (hleaf : qq ≤ ∑ k ∈ range (n + 1), (θ * T) ^ k / (Nat.factorial k : ℚ)) :
    (qq : ℝ) ≤ Real.exp (θ * T)
\end{leansrc}
\par\medskip
\textbf{The payload-level elaborator, displayed whole.}\enspace
$(q, \theta, T, n)$ to the one \rulename{PConstExact} payload, the
right side expanded producer-side in exact $\Q$; the checker replays
one closed comparison, and the transcendental conclusion appears only
in \texttt{taylorElab\_le\_exp}\textquotesingle s goal.

\declhead{BernsteinLogic.QtslElaborators}{taylorFormula\_true\_iff}{\badge{BadgeGreen}{statement bound}}
\begin{leansrc}
/-- The wire-level replay identity: the checker's `true` verdict is
exactly the ℚ comparison against the partial sum. -/
theorem taylorFormula_true_iff (q θ : Rat) (T n : Nat) :
    (taylorFormula q θ T n).evalClosed = some true ↔ q ≤ taylorPartialSum θ T n
\end{leansrc}
\declhead{BernsteinLogic.QtslElaborators}{taylorElab\_sound}{\badge{BadgeGreen}{statement bound}}
\begin{leansrc}
/-- The factoring, BL face: `sound_pConstExact` at the elaborated
payload (`prop:tb-sound`'s first half — validity of the posted
comparison at every valuation). -/
theorem taylorElab_sound {H : Hypernet} {seq : Sequent} {q θ : Rat} {T n : Nat}
    (hEval : (taylorFormula q θ T n).evalClosed = some true)
    (hpost : seq.post = taylorFormula q θ T n) :
    seq.semanticallyValid H
\end{leansrc}
\declhead{BernsteinLogic.QtslElaborators}{taylorElab\_le\_exp}{\badge{BadgeGreen}{statement bound}}
\begin{leansrc}
/-- The semantic face, end to end (`prop:tb-sound` composed with
`prop:tb-enclosure`): the checked wire comparison bounds `q` by
`e^{θT}` through the bound enclosure — the elaborator factoring through
the kernel replay into `taylorLeaf_rat_le_exp`. The sign condition
`0 ≤ θ` is load-bearing (the proof section's negative-argument
counterexample); the deadline rule's typing supplies it. -/
theorem taylorElab_le_exp {q θ : Rat} {T n : Nat} (hθ : 0 ≤ θ)
    (hEval : (taylorFormula q θ T n).evalClosed = some true) :
    (q : ℝ) ≤ Real.exp (θ * T)
\end{leansrc}
\declhead{BernsteinLogic.QtslElaborators}{taylorPartialSum}{\badge{BadgeGreen}{definition bound}}
\begin{leansrc}
/-- The order-`n` Taylor partial sum of `exp (θ·T)`, computed in exact
ℚ — the producer-side expansion of `rec:tb-transport`'s right side. -/
def taylorPartialSum (θ : Rat) (T n : Nat) : Rat :=
  ∑ k ∈ Finset.range (n + 1), (θ * (T : Rat)) ^ k / (Nat.factorial k : Rat)
\end{leansrc}
\declhead{BernsteinLogic.QtslElaborators}{taylorElab}{\badge{BadgeGreen}{definition bound}}
\begin{leansrc}
/-- **The TaylorBound payload elaborator** (`rec:tb-transport`): (q, θ,
T, n) plus the checked comparison, to the one `PConstExactData`
inhabitant. No `TaylorBound` node sort exists and none is needed
(`rec:tb-template`) — this function is the catalogue-level name for the
producer's elaboration, mechanized. -/
def taylorElab (H : Hypernet) (seq : Sequent) (q θ : Rat) (T n : Nat)
    (hEval : (taylorFormula q θ T n).evalClosed = some true)
    (hpost : seq.post = taylorFormula q θ T n) :
    PConstExactData H seq :=
  .mk H seq (.const q) (.const (taylorPartialSum θ T n)) .le (by decide)
    hEval hpost
\end{leansrc}
\ruleblock{Correspondence (Steele display $\leftrightarrow$ Lean)}
\badge{BadgeGreen}{reviewer-certified against the displayed texts}
\sourcenote{The dictionary below is a reviewer-certified statement of how the mechanized statements above interpret the displayed rule: display metavariable $\mapsto$ Lean binder, each display hypothesis $\mapsto$ the premise or side condition carrying it, the display conclusion $\mapsto$ the goal --- with any deliberate difference in width between display and mechanization named outright. Certification is recorded as a hash triple over (the displayed cell, the extracted statement texts, this dictionary\textquotesingle s own text); if any of the three changes, this block turns amber until a reviewer re-certifies. The correspondence claim itself remains editorial --- reviewer-grade, not machine-checked; its machine-checked successor types each soundness statement against the rule as a structured object, and supersedes this block per rule as it lands.}
The displayed wire-free comparison $q \le \sum_{k \le n}
(\theta T)^k / k!$ standing in for $q \le e^{\theta T}$ $\mapsto$ the
bound pair: \texttt{taylorLeaf\_le\_exp} (the enclosure step ---
partial sum below \texttt{Real.exp} at nonnegative argument, the
nonneg-remainder citation) and \texttt{taylorLeaf\_rat\_le\_exp}, the
end-to-end exact-$\Q$ face: \texttt{qq θ : ℚ}, \texttt{T n : ℕ},
hypothesis the exact-$\Q$ partial-sum comparison the checker replays,
conclusion \texttt{(qq : ℝ) ≤ Real.exp (θ * T)} --- the
flavor text\textquotesingle s ``the checker never needs to know $e$
exists'' as a binder fact: $e$ appears in the goal alone, never in a
hypothesis. The sign condition $\mapsto$ \texttt{hθ : 0 ≤ θ}
(load-bearing --- the alternating-remainder counterexample in the
proof section is why the binder exists); the truncation order
\texttt{n} is producer-chosen, unconstrained. Width note: the
consuming site (\texttt{sound\_QDeadlineBattery}\textquotesingle s
\texttt{hq} premise) is cited, not bound here --- correct
package-boundary behaviour, the deadline battery\textquotesingle s own
entry carries it; the displayed horizon $T \in \N$ matches the bound end-to-end
lemma\textquotesingle s \texttt{T : ℕ} exactly (the rational-horizon
generalization is mathematically available from
\texttt{taylorLeaf\_le\_exp} but is neither bound nor displayed); and
the displayed $\theta \in \Q_{>0}$ is a deliberately stronger surface
restriction than the bound \texttt{hθ : 0 ≤ θ} --- safe direction, the
deadline battery\textquotesingle s own decay-rate positivity.

\ruleblock{Lean verifier receipts}
\sourcenote{module \texttt{BernsteinLogicAux.QtslAliases}; bound build \texttt{570599779c29\ldots} (bound from receipt 5a60f4d01fd81ac5, 2026-07-21T05:39:02Z). A receipt records that named checking layers accepted a build: (module, compiled-artifact SHA-256, environment SHA-256, checker layers), minted when the checking gate ran; the store receipts attesting the bound build are listed with the layers each run credited. Nothing here asks the reader to trust the store: every layer is an independently runnable program over the quoted sources --- rebuild the module and re-run the checkers to reproduce the verdict. Receipts attesting superseded builds credit no layer and are omitted from this rendering (the count of omissions, if any, is noted below).}
\begin{receipttable}
\receiptrow{0a3819a978e31127}{2026-07-21T04:51:36Z}{build+decorrelated}{\badge{BadgeGreen}{current}}
\receiptrow{0bb1074df705d863}{2026-07-21T05:21:58Z}{build+decorrelated}{\badge{BadgeGreen}{current}}
\receiptrow{5a60f4d01fd81ac5}{2026-07-21T05:39:02Z}{build+replay+hardened+decorrelated}{\badge{BadgeGreen}{current}}
\end{receipttable}
\sourcenote{11 further receipt(s) in the store attest superseded builds of this module and are omitted from this rendering.}
\sourcenote{module \texttt{BernsteinLogic.QtslElaborators}; bound build \texttt{0205ec213b91\ldots} (hashed from the built olean at bind time). A receipt records that named checking layers accepted a build: (module, compiled-artifact SHA-256, environment SHA-256, checker layers), minted when the checking gate ran; the store receipts attesting the bound build are listed with the layers each run credited. Nothing here asks the reader to trust the store: every layer is an independently runnable program over the quoted sources --- rebuild the module and re-run the checkers to reproduce the verdict. Receipts attesting superseded builds credit no layer and are omitted from this rendering (the count of omissions, if any, is noted below).}
\begin{receipttable}
\receiptrow{a237f4cff27cb15a}{2026-07-22T10:06:47Z}{build+decorrelated}{\badge{BadgeGreen}{current}}
\end{receipttable}
\layerline{\layeritem{build}{lake build (elaboration)}{\badge{BadgeGreen}{current}}\layeritem{replay}{leanchecker (official kernel replay)}{\badge{BadgeGreen}{current}}\layeritem{hardened}{safe\_verify (hardened screen)}{\badge{BadgeGreen}{current}}\layeritem{decorrelated}{lean4export + nanoda (decorrelated kernel)}{\badge{BadgeGreen}{current}}}
\sourcenote{\badge{BadgeGray}{comparator: none} no fixture in the three-implementation comparator harness (compiled / kernel-interpreted / reference-interpreted) exercises this rule.}
\ruleblock{HOL4 statement}
\absentbox{OUT OF SCOPE BY DESIGN}{Out of scope for the HOL side \emph{by design}, not by backlog: the
wire format never carries a \rulename{TaylorBound} node (no sort
exists and none is needed --- the producer emits the ground-fact
comparison row directly), so no wire certificate a HOL checker reads
could ever contain one; HOL coverage for this rule \emph{is} the
target rule\textquotesingle s. Today even that statement is not
writable: the HOL4 development (theory \texttt{bern\_catalogue}, the
base catalogue\textquotesingle s check-function cone) predates
\rulename{PConstExact}, whose wire tag is unallocated --- see that
rule\textquotesingle s section for the honest state of the
target\textquotesingle s own box.
}
\ruleblock{HOL verifier receipts}
\absentbox{NONE}{no HOL verification receipt exists for this rule.}
\ruleblock{Proof}

\sourcenote{Genre: alias rule (tier 3) --- a respelling template. The
five artefacts an alias owes the book: the 1$\to$1 unpack, the filled-in
target application, the payload-transport table, the exactness identity,
and the soundness corollary through the target\textquotesingle s lemma.
All five follow, with one honest asymmetry marked where it occurs: this
rule\textquotesingle s exactness identity is not a ring identity but an
\emph{enclosure} --- the Taylor-remainder inequality --- and it lives in
the consuming lemma, never in the checker. Mechanization state: the
target\textquotesingle s lemma is kernel-verified; the alias-level
semantic content --- the enclosure and its exact-$\Q$ face --- is
mechanized and bound in the entry, and so --- now that the extension
certificate schema exists --- is the payload-level elaborator (the
remark below names the bound functions).}

\begin{brecipe}[the template instance --- no node sort, by design]
\label{rec:tb-template}
The deadline rule (\rulename{QDeadlineBattery}) wants the fact
$q \le e^{\theta T}$ for a producer-chosen rational surrogate $q$, with
decay rate $\theta \in \Q_{> 0}$ and horizon $T \in \Q_{\ge 0}$ from its
payload: its concluded sublevel region $\{V \le (1-p)\,q\}$ grows with
$q$, so the producer wants the sharpest $q$ the arithmetic affords. No
checker doing exact rational arithmetic can evaluate $e^{\theta T}$, and
none needs to: the certificate posts the single wire-free comparison
\[
q \;\le\; \textstyle\sum_{k \le n} (\theta T)^k / k!
\]
for a producer-chosen truncation order $n$ --- a finite sum of
rationals, one closed comparison --- and the consuming lemma carries the
enclosure to $e^{\theta T}$ (Proposition~\ref{prop:tb-enclosure}). The
unpack image is a single \rulename{PConstExact} node; indeed no
\rulename{TaylorBound} node sort exists and none is needed --- the
producer emits the ground-fact row directly, and this rule is the
catalogue-level name for that elaboration. One node $\to$ one node,
nothing derived.
\end{brecipe}

\begin{brecipe}[the filled-in application]\label{rec:tb-filled}
\begin{mathpar}
\inferrule*[Right=\rulename{PConstExact}]{
  q \le \textstyle\sum_{k \le n} (\theta T)^k / k! \ \text{holds under exact evaluation in } \Q
}{
  \Phi \vdash q \le \textstyle\sum_{k \le n} (\theta T)^k / k!
}
\end{mathpar}
with arbitrary premise $\Phi$, as the target licenses (the rule is
premise-free). Both sides are closed rational terms --- the sum is
finite, every summand exact --- so the target\textquotesingle s verified
evaluator decides the comparison in exact $\Q$, and the comparator
($\le$) lies inside the target\textquotesingle s admissible comparator
set.
\end{brecipe}

\begin{brecipe}[payload transport --- field for field]\label{rec:tb-transport}
\begin{center}\small
\begin{tabular}{@{}ll@{}}
\textbf{\rulename{TaylorBound} data} & \textbf{\rulename{PConstExact} payload}\\
\midrule
$q$ \textnormal{(the rational surrogate)} & left side of the comparison\\
$\theta, T, n$ \textnormal{(rate, horizon, order)} & the right side $\textstyle\sum_{k \le n} (\theta T)^k/k!$, expanded producer-side\\
comparator & $\le$ (inside the admissible set)\\
\textnormal{wire-freedom} & implied, not declared: the evaluator returns $\mathsf{none}$ on any wire\\
\end{tabular}
\end{center}
The checker sees one closed comparison and nothing else; $\theta$, $T$,
and $n$ survive only through the value of the right-hand term. The
transcendental target $e^{\theta T}$ appears in \emph{no} field --- that
is the point.
\end{brecipe}

\begin{bprop}[the enclosure --- the exactness identity, honest form]
\label{prop:tb-enclosure}
For $\theta T \ge 0$ and any $n$,
\[
\textstyle\sum_{k \le n} (\theta T)^k / k! \;\le\; e^{\theta T},
\qquad\text{hence}\qquad
q \le \textstyle\sum_{k \le n} (\theta T)^k / k! \;\Longrightarrow\; q \le e^{\theta T}.
\]
\end{bprop}
\begin{proof}
Two lines, though of a different species than a ring identity: at
$x := \theta T \ge 0$ every term of the exponential series is
nonnegative, so each partial sum is dominated by the limit --- in the
classical Taylor form, $e^{x} = \sum_{k \le n} x^k/k! + R_n$ with the
Lagrange remainder $R_n = e^{\xi} x^{n+1}/(n+1)! \ge 0$ for some
$\xi \in [0, x]$ \cite{rudin1976}. The sign condition is load-bearing:
at negative arguments the remainder alternates in sign and the
enclosure direction fails at alternate truncation orders ($n = 0$,
$x = -1$: the partial sum is $1 > e^{-1}$; at $n = 1$ it happens to
hold again --- there is no all-$n$ enclosure to be had) --- the
deadline rule\textquotesingle s typing ($\theta > 0$, $T \ge 0$) is
what entitles the lemma to the inequality, via $\theta T \ge 0$. Unlike the
base-calculus exemplar alias, whose two-line identity is checked by the
kernel as exact polynomial arithmetic, this identity is consumed
\emph{inside the deadline rule\textquotesingle s soundness lemma} (the
Lean side consumes the standard library\textquotesingle s
partial-sum-below-exponential fact,
\texttt{Real.sum\_le\_exp\_of\_nonneg}, there) and never reaches the
checker: the checker\textquotesingle s whole obligation is the exact
rational comparison of Recipe~\ref{rec:tb-filled}. The checker never
needs to know $e$ exists.
\end{proof}

\begin{bprop}[soundness corollary]\label{prop:tb-sound}
If the posted comparison checks and $\theta T \ge 0$ (the sign
hypothesis of Proposition~\ref{prop:tb-enclosure}), then
$q \le e^{\theta T}$, which is the fact the deadline
rule\textquotesingle s lemma consumes.
\end{bprop}
\begin{proof}
\texttt{sound\_pConstExact} (the target\textquotesingle s
kernel-verified lemma, quoted in the \rulename{PConstExact} section)
turns the $\mathsf{true}$ evaluation into the validity of
$q \le \sum_{k \le n} (\theta T)^k/k!$ at every valuation;
Proposition~\ref{prop:tb-enclosure} extends to $e^{\theta T}$ inside
the consuming lemma. There is no other content --- the alias genre
working as intended.
\end{proof}

\begin{bremark}[the payload-level elaborator, mechanized]
\label{rem:tb-elab}
The \emph{semantic} content of the respelling is bound in the entry:
\texttt{taylorLeaf\_le\_exp} is
Proposition~\ref{prop:tb-enclosure}\textquotesingle s enclosure step,
and \texttt{taylorLeaf\_rat\_le\_exp} is
Proposition~\ref{prop:tb-sound}\textquotesingle s exact-$\Q$ face end
to end. The payload-level elaborator --- previously this
page\textquotesingle s one absent artefact, held back until the
extension certificate schema existed --- is bound with it:
\texttt{taylorElab} (module \texttt{BernsteinLogic.QtslElaborators})
takes $(q, \theta, T, n)$ to exactly the one \rulename{PConstExact}
payload of Recipe~\ref{rec:tb-transport}, the right side expanded
producer-side in exact $\Q$ (\texttt{taylorPartialSum}), as an
inhabitant of the target\textquotesingle s own payload type --- the
comparator confinement discharged at elaboration, once --- with
\texttt{taylorElab\_sound} the factoring through
\texttt{sound\_pConstExact}, \texttt{taylorFormula\_true\_iff} the
wire-level replay identity, and \texttt{taylorElab\_le\_exp} the
semantic face end to end (the checked comparison bounding $q$ by
$e^{\theta T}$ through the bound enclosure --- the transcendental
conclusion still appears in no field, only in the goal). Stated at the
discrete horizon, $T : \mathbb{N}$, matching the bound exact-$\Q$ face
--- the entry\textquotesingle s width note carries the divergence from
the recipe\textquotesingle s rational-time payload; the enclosure of
Proposition~\ref{prop:tb-enclosure} itself needs only
$\theta T \ge 0$.
\end{bremark}

\begin{bremark}[deliberately not an enclosure-rule alias]
The calculus also has interval-enclosure machinery with witness
brackets. Respelling this leaf onto it was considered and declined:
that would import witness-bracket payload for a metatheory one-liner,
and the enclosure rule today lives only in a separate
reference-checker implementation besides. The seam this rule
instantiates is a design invariant of the whole extension page:
transcendental facts live in soundness lemmas, arithmetic lives in the
checker.
\end{bremark}

\begin{bremark}[the inherited comparator confinement]
\rulename{PConstExact} excludes one comparator ($\neq$) from its
admissible set, for reasons of a second soundness invariant unpacked in
that rule\textquotesingle s section. This leaf uses $\le$ only ---
comfortably inside --- but any port or mechanization of this rule must
\emph{preserve} the exclusion rather than re-derive the comparator set
from the semantic lemma alone; the flag travels with the alias.
\end{bremark}

\begin{anchors}
target lemma (Prop.~\ref{prop:tb-sound}) & \texttt{sound\_pConstExact}; \texttt{Formula.holds\_of\_evalClosed\_true} (quoted in the \rulename{PConstExact} section)\\
wire-freedom implied at the leaf & \texttt{Formula.evalClosed} ($\mathsf{none}$ on any wire; same section)\\
the enclosure, mechanized where consumed & \texttt{Real.sum\_le\_exp\_of\_nonneg} inside the deadline lemma \texttt{sound\_QDeadlineBattery} (see that rule\textquotesingle s entry)\\
the enclosure + exact-$\Q$ face, rule-level & \texttt{taylorLeaf\_le\_exp}, \texttt{taylorLeaf\_rat\_le\_exp} (bound in the entry)\\
payload-level elaborator (\autoref{rem:tb-elab}) & \texttt{taylorElab} + \texttt{taylorElab\_sound} + \texttt{taylorElab\_le\_exp}, over \texttt{taylorPartialSum} + \texttt{taylorFormula\_true\_iff} (module \texttt{BernsteinLogic.QtslElaborators}, bound in the entry)\\
\end{anchors}


\rulesection{67}{RootEval}{alias for \textcolor{DiffBlue}{\rulename{PConstExact}}}
\label{blcatrule-RootEval}
\ruleblock{The rule (Steele metanotation)}
\sourcenote{as realized in the mechanized development; source row: \texttt{\texttt{QTSL-rules.toml}} --- \emph{RootEval}. The metanotation statement is typeset by hand for notational quality; its mechanized ground truth is the Lean statement below, extracted from the sources at every build of this book and hash-bound. The rule \emph{roster} is machine-checked (every mechanized rule tag round-trips to exactly one catalogue row); the statement \emph{text} is editorial, checked by eye against the extracted Lean beside it.}
\begin{mathpar}
\inferrule*{
  R = \{x : V(x) \le y\} \ \text{(sublevel root atom)} \\
  \vdash V(x_0) \le y \ \text{(one \rulename{PConstExact} leaf: the substituted comparison is closed)}
}{
  \vdash x_0 \in R
}
\and
\inferrule*{
  R = \textstyle\prod_i [a_i, b_i] \ \text{(box root atom)} \\
  \forall i.\ a_i \le (x_0)_i \le b_i \ \text{(recomputed in exact } \Q\text{)}
}{
  \vdash x_0 \in R
}
\end{mathpar}

\ruleflavor{catalogued as the eval leaf of the region-to-satisfaction bridge (a BL \texttt{Proof} leaf; the bridge\textquotesingle s node sort is out of scope of the current baked theory — its baked warrant seam is \texttt{qsat\_warrant}); sublevel atoms discharge through \rulename{PConstExact} (the first form is that rule\textquotesingle s application at the substituted point), box atoms through the registry interval check (the second form: conformance recomputation, not a rule application).}
\par\medskip
\rulename{RootEval} is the membership leaf itself, catalogued as a
rule of its own because every initialized certificate ends with one:
membership of the specification\textquotesingle s initial point in the
\emph{root atom} --- the region atom concluded at the
certificate\textquotesingle s root (``root'' names the
certificate\textquotesingle s root claim, not roots of polynomials)
--- by exact rational evaluation. In the certificate schema it
is not a node sort at all: it is the eval leaf of the region-to-
satisfaction bridge (\rulename{QSatFromRegion} --- whose own node sort
is out of scope of the current baked theory; the baked warrant seam is
\texttt{qsat\_warrant}), a base-proof leaf. For a sublevel atom
$\{x : V(x) \le y\}$ the leaf is one \rulename{PConstExact}
application: substitute the rational point into the polynomial ---
exact, no rounding anywhere --- and decide one closed comparison. For
a box atom, the checker recomputes per-coordinate containment directly
(the side-condition registry: conformance recomputation, not a rule
application).

What the checker recomputes: everything --- the leaf \emph{is} a
recompute; there is no payload beyond the point and the atom. What
would break without exactness: under naive (non-directed) rounding a
point lying $\varepsilon$ outside the region could certify
membership, and the bridge rule would conclude satisfaction from a
false premise --- directed outward rounding keeps that direction
sound but surrenders the equivalence at the boundary --- the exactness identity in the
proof section is a literal identity of values, and it is the whole
rule. The comparator exclusion inherited from the target ($\neq$
inadmissible, a second-soundness-invariant exclusion unpacked in the
target\textquotesingle s section) must survive any port.

\cataloguenote{The judgment-level composition --- how the checked leaf feeds the
instantiation that concludes root satisfaction --- lives in the
\rulename{QSatFromRegion} section; this row owns the leaf.
Mechanization state, plainly: the bound statements below are the
exactness identity itself. \texttt{rootEval\_eval\_cast} (module
\texttt{BernsteinLogicAux.QtslAliases}) is the literal identity of
values --- evaluating the $\Q$-coefficient polynomial at the rational
point and casting equals evaluating the cast polynomial at the cast
point --- and \texttt{rootEval\_sublevel\_iff} /
\texttt{rootEval\_box\_iff} are its two atom shapes, each an
\emph{equivalence}: the exact-$\Q$ comparison decides the
real-semantics membership in both directions, which is precisely the
content a rounded evaluator would lose. The statements are over the
generic polynomial-evaluation vocabulary (\texttt{MvPolynomial} over
$\Q$, cast to $\R$); the wire-term evaluator that produces the
$\Q$-comparison is the target\textquotesingle s own kernel-verified
theorem, quoted in the \rulename{PConstExact} section, and the
box-atom branch\textquotesingle s registry recompute is checker-side
by design. No profile enters; the entry carries no conditional
flag.
}
\ruleblock{Lean statement}
\declhead{BernsteinLogicAux.QtslAliases}{rootEval\_eval\_cast}{\badge{BadgeGreen}{statement bound}}
\begin{leansrc}
/-- **The exactness identity** (QTSL-rules.toml row `RootEval (T19)` — "a
literal identity of values, and it is the whole rule"): evaluating a
ℚ-coefficient polynomial at a rational point and casting the result equals
evaluating the cast polynomial at the cast point. This is why the checker's
exact rational arithmetic (`PConstExact`'s replay) *is* the real-semantics
evaluation, with no rounding gap anywhere. -/
theorem rootEval_eval_cast {σ : Type*} (P : MvPolynomial σ ℚ) (xq : σ → ℚ) :
    ((MvPolynomial.eval xq P : ℚ) : ℝ)
      = MvPolynomial.eval (fun i => (xq i : ℝ)) (MvPolynomial.map (Rat.castHom ℝ) P)
\end{leansrc}
\declhead{BernsteinLogicAux.QtslAliases}{rootEval\_sublevel\_iff}{\badge{BadgeGreen}{statement bound}}
\begin{leansrc}
/-- **`RootEval`, sublevel atom, as an equivalence**: the exact ℚ comparison
at the initial point *decides* membership of the cast point in the cast
sublevel atom `{x : V(x) ≤ thr}` — both directions, which is the content a
rounded evaluator would lose (one direction would survive rounding; the
identity does not). The forward direction is the leaf as the checker
discharges it (rule row `RootEval (T19)`, sublevel branch →
`PConstExact`). -/
theorem rootEval_sublevel_iff {σ : Type*} (P : MvPolynomial σ ℚ) (xq : σ → ℚ)
    (thr : ℚ) :
    MvPolynomial.eval xq P ≤ thr
      ↔ MvPolynomial.eval (fun i => (xq i : ℝ))
          (MvPolynomial.map (Rat.castHom ℝ) P) ≤ (thr : ℝ)
\end{leansrc}
\declhead{BernsteinLogicAux.QtslAliases}{rootEval\_box\_iff}{\badge{BadgeGreen}{statement bound}}
\begin{leansrc}
/-- **`RootEval`, box atom, as an equivalence**: per-coordinate exact ℚ
containment decides containment of the cast point in the cast box (rule row
`RootEval (T19)`, box branch — the side-condition registry's per-coordinate
recompute). -/
theorem rootEval_box_iff {ι : Type*} (aq bq xq : ι → ℚ) :
    (∀ i, aq i ≤ xq i ∧ xq i ≤ bq i)
      ↔ ∀ i, ((aq i : ℝ) ≤ (xq i : ℝ) ∧ (xq i : ℝ) ≤ (bq i : ℝ))
\end{leansrc}
\par\medskip
\textbf{The payload-level elaborator, displayed whole.}\enspace The
closed left-hand term is built by verified substitution
(\texttt{polyClosedTerm}; its evaluation lemma says the
checker\textquotesingle s exact replay computes the value $V(x_0)$
itself), and the elaborated payload inhabits the
target\textquotesingle s own payload type --- comparator confinement
discharged at elaboration.

\declhead{BernsteinLogic.QtslElaborators}{polyClosedTerm\_evalQ}{\badge{BadgeGreen}{statement bound}}
\begin{leansrc}
@[simp] theorem polyClosedTerm_evalQ (x0 : Nat → Rat) :
    ∀ p : Poly, (polyClosedTerm x0 p).evalQ? = some (Poly.eval x0 p)
  | [] => by simp [polyClosedTerm, Poly.eval, Term.evalQ?]
  | mc :: rest => by
    have ih
\end{leansrc}
\declhead{BernsteinLogic.QtslElaborators}{rootEvalFormula\_true\_iff}{\badge{BadgeGreen}{statement bound}}
\begin{leansrc}
/-- The wire-level exactness identity, ℚ face (`prop:re-exact` at the
payload: the checker's `true` verdict IS the exact membership
inequality — both directions; the ℝ face stays
`Aliases.rootEval_sublevel_iff`). -/
theorem rootEvalFormula_true_iff (x0 : Nat → Rat) (V : Poly) (thr : Rat) :
    (rootEvalFormula x0 V thr).evalClosed = some true ↔ Poly.eval x0 V ≤ thr
\end{leansrc}
\declhead{BernsteinLogic.QtslElaborators}{rootEvalElab\_sound}{\badge{BadgeGreen}{statement bound}}
\begin{leansrc}
/-- The factoring, BL face: the elaborated payload is sound through the
target's kernel-verified lemma — one application of
`sound_pConstExact`, which is itself one step to
`Formula.holds_of_evalClosed_true` (`prop:re-sound`). -/
theorem rootEvalElab_sound {H : Hypernet} {seq : Sequent} {x0 : Nat → Rat}
    {V : Poly} {thr : Rat}
    (hEval : (rootEvalFormula x0 V thr).evalClosed = some true)
    (hpost : seq.post = rootEvalFormula x0 V thr) :
    seq.semanticallyValid H
\end{leansrc}
\declhead{BernsteinLogic.QtslElaborators}{polyClosedTerm}{\badge{BadgeGreen}{definition bound}}
\begin{leansrc}
/-- The closed polynomial term at a rational point. Wire-free by
construction, and its exact evaluation is `Poly.eval` at that point
(`polyClosedTerm_evalQ`) — the checker computes the value `V(x₀)`
itself, not an approximation to it. -/
def polyClosedTerm (x0 : Nat → Rat) (p : Poly) : Term :=
  p.foldr (fun mc acc => .add (.mul (.const mc.2) (monoClosedTerm x0 mc.1)) acc)
    (.const 0)
\end{leansrc}
\declhead{BernsteinLogic.QtslElaborators}{rootEvalFormula}{\badge{BadgeGreen}{definition bound}}
\begin{leansrc}
/-- The elaborated leaf formula: `V(x₀) ≤ thr` as one closed exact
comparison. -/
def rootEvalFormula (x0 : Nat → Rat) (V : Poly) (thr : Rat) : Formula :=
  .atom .le (polyClosedTerm x0 V) (.const thr)
\end{leansrc}
\declhead{BernsteinLogic.QtslElaborators}{rootEvalElab}{\badge{BadgeGreen}{definition bound}}
\begin{leansrc}
/-- **The RootEval payload elaborator** (`rec:re-transport`, sublevel
branch): the alias's data — point, atom — plus its own checked leaf, to
the one `PConstExactData` inhabitant. The comparator confinement
travels by construction (`.le ≠ .ne` discharged here, once). -/
def rootEvalElab (H : Hypernet) (seq : Sequent) (x0 : Nat → Rat) (V : Poly)
    (thr : Rat)
    (hEval : (rootEvalFormula x0 V thr).evalClosed = some true)
    (hpost : seq.post = rootEvalFormula x0 V thr) :
    PConstExactData H seq :=
  .mk H seq (polyClosedTerm x0 V) (.const thr) .le (by decide) hEval hpost
\end{leansrc}
\ruleblock{Correspondence (Steele display $\leftrightarrow$ Lean)}
\badge{BadgeGreen}{reviewer-certified against the displayed texts}
\sourcenote{The dictionary below is a reviewer-certified statement of how the mechanized statements above interpret the displayed rule: display metavariable $\mapsto$ Lean binder, each display hypothesis $\mapsto$ the premise or side condition carrying it, the display conclusion $\mapsto$ the goal --- with any deliberate difference in width between display and mechanization named outright. Certification is recorded as a hash triple over (the displayed cell, the extracted statement texts, this dictionary\textquotesingle s own text); if any of the three changes, this block turns amber until a reviewer re-certifies. The correspondence claim itself remains editorial --- reviewer-grade, not machine-checked; its machine-checked successor types each soundness statement against the rule as a structured object, and supersedes this block per rule as it lands.}
The displayed leaf rules --- exact rational evaluation concluding
root-atom membership, one form per atom shape --- $\mapsto$ three statements, and the
strength gradation is the dictionary\textquotesingle s content.
\texttt{rootEval\_eval\_cast} is the literal identity of values:
$\Q$-evaluation, cast, equals $\R$-evaluation of the cast polynomial
at the cast point (\texttt{MvPolynomial} over any variable type ---
dimension-free). The two atom shapes are then \emph{equivalences},
not implications: \texttt{rootEval\_sublevel\_iff} (exact-$\Q$
comparison $\leftrightarrow$ membership in the cast sublevel atom)
and \texttt{rootEval\_box\_iff} (per-coordinate $\Q$ containment
$\leftrightarrow$ cast-box membership) --- two-directional exactness,
which is precisely what a rounded evaluator would lose (under
directed outward rounding the sound forward direction survives and
the iff does not; undirected rounding loses even soundness): the entry\textquotesingle s
``no rounding gap anywhere'' as bind structure. Width note: the
checker consumes only the forward directions (leaf discharge via
\rulename{PConstExact} / the registry interval recompute); the
certified equivalences are strictly stronger, safe direction, and no
profile, measure, or topology enters any binder.

\ruleblock{Lean verifier receipts}
\sourcenote{module \texttt{BernsteinLogicAux.QtslAliases}; bound build \texttt{570599779c29\ldots} (bound from receipt 5a60f4d01fd81ac5, 2026-07-21T05:39:02Z). A receipt records that named checking layers accepted a build: (module, compiled-artifact SHA-256, environment SHA-256, checker layers), minted when the checking gate ran; the store receipts attesting the bound build are listed with the layers each run credited. Nothing here asks the reader to trust the store: every layer is an independently runnable program over the quoted sources --- rebuild the module and re-run the checkers to reproduce the verdict. Receipts attesting superseded builds credit no layer and are omitted from this rendering (the count of omissions, if any, is noted below).}
\begin{receipttable}
\receiptrow{0a3819a978e31127}{2026-07-21T04:51:36Z}{build+decorrelated}{\badge{BadgeGreen}{current}}
\receiptrow{0bb1074df705d863}{2026-07-21T05:21:58Z}{build+decorrelated}{\badge{BadgeGreen}{current}}
\receiptrow{5a60f4d01fd81ac5}{2026-07-21T05:39:02Z}{build+replay+hardened+decorrelated}{\badge{BadgeGreen}{current}}
\end{receipttable}
\sourcenote{11 further receipt(s) in the store attest superseded builds of this module and are omitted from this rendering.}
\sourcenote{module \texttt{BernsteinLogic.QtslElaborators}; bound build \texttt{0205ec213b91\ldots} (hashed from the built olean at bind time). A receipt records that named checking layers accepted a build: (module, compiled-artifact SHA-256, environment SHA-256, checker layers), minted when the checking gate ran; the store receipts attesting the bound build are listed with the layers each run credited. Nothing here asks the reader to trust the store: every layer is an independently runnable program over the quoted sources --- rebuild the module and re-run the checkers to reproduce the verdict. Receipts attesting superseded builds credit no layer and are omitted from this rendering (the count of omissions, if any, is noted below).}
\begin{receipttable}
\receiptrow{a237f4cff27cb15a}{2026-07-22T10:06:47Z}{build+decorrelated}{\badge{BadgeGreen}{current}}
\end{receipttable}
\layerline{\layeritem{build}{lake build (elaboration)}{\badge{BadgeGreen}{current}}\layeritem{replay}{leanchecker (official kernel replay)}{\badge{BadgeGreen}{current}}\layeritem{hardened}{safe\_verify (hardened screen)}{\badge{BadgeGreen}{current}}\layeritem{decorrelated}{lean4export + nanoda (decorrelated kernel)}{\badge{BadgeGreen}{current}}}
\sourcenote{\badge{BadgeGray}{comparator: none} no fixture in the three-implementation comparator harness (compiled / kernel-interpreted / reference-interpreted) exercises this rule.}
\ruleblock{HOL4 statement}
\absentbox{OUT OF SCOPE BY DESIGN}{Out of scope for the HOL side \emph{by design}, not by backlog: this
rule is a \emph{field} of the bridge rule
(\rulename{QSatFromRegion}\textquotesingle s eval slot), never a
node, so no wire certificate a HOL checker reads could ever contain a
\rulename{RootEval} node; HOL coverage for this rule \emph{is} the
target rule\textquotesingle s, and the box-atom branch is loader-side
registry recomputation, likewise not a wire-format rule. Today even
the target\textquotesingle s statement is not writable: the HOL4
development (theory \texttt{bern\_catalogue}, the base
catalogue\textquotesingle s check-function cone) predates
\rulename{PConstExact}, whose wire tag is unallocated --- see that
rule\textquotesingle s section for the honest state of the
target\textquotesingle s own box.
}
\ruleblock{HOL verifier receipts}
\absentbox{NONE}{no HOL verification receipt exists for this rule.}
\ruleblock{Proof}

\sourcenote{Genre: alias rule (tier 3) --- a respelling template: the
1$\to$1 unpack, the filled-in target application, the payload-transport
table, the exactness identity, and the soundness corollary through the
target\textquotesingle s lemma. This row is the membership leaf itself,
catalogued as a rule of its own because every initialized certificate
ends with one; in the extension certificate schema it is a \emph{field} of
the bridge rule (\rulename{QSatFromRegion}\textquotesingle s eval slot),
not a node sort. The judgment-level story lives in the bridge
rule\textquotesingle s section --- this section owns the leaf.
Mechanization state: the target\textquotesingle s lemma is
kernel-verified; the exactness identities are mechanized and bound in
the entry, and so --- now that the extension certificate schema exists
--- is the payload-level elaborator (the remark below names the bound
functions).}

\begin{brecipe}[the template instance]\label{rec:re-template}
One leaf $\to$ one node: membership of the initial point $x_0$ in the
certificate\textquotesingle s \emph{root atom} --- the region atom
concluded at the root of the certificate (naming note: ``root'' refers
to the certificate\textquotesingle s root claim, not to roots of
polynomials) --- by exact rational evaluation. For a \emph{sublevel}
atom $\{x : V(x) \le y\}$ the unpack image is a single
\rulename{PConstExact} application
(Recipe~\ref{rec:re-filled}). For a \emph{box} atom the leaf is not a
rule application at all: per-coordinate interval containment is
recomputed directly by the checker\textquotesingle s side-condition
registry (conformance recomputation) --- the alias reading covers the
sublevel branch, and the box branch is stated so the respelling never
overclaims. Nothing else is derived: the rule is a compact spelling of
one ground-fact row.
\end{brecipe}

\begin{brecipe}[the filled-in application]\label{rec:re-filled}
Sublevel atom $R = \{x : V(x) \le y\}$ with $V$ a polynomial
over $\Q$, threshold $y \in \Q$, and initial point
$x_0$ with rational coordinates:
\begin{mathpar}
\inferrule*[Right=\rulename{PConstExact}]{
  V(x_0) \le y \ \text{holds under exact evaluation in } \Q
}{
  \Phi \vdash V(x_0) \le y
}
\end{mathpar}
with arbitrary premise $\Phi$, as the target licenses (the rule is
premise-free). $V(x_0)$ denotes the closed term obtained by
substituting every coordinate of $x_0$ into $V$ --- closed by
construction, so the target\textquotesingle s verified evaluator
decides the comparison in exact $\Q$, and a $\mathsf{true}$ verdict
certifies closedness as a byproduct (the target section unpacks why).
\end{brecipe}

\begin{brecipe}[payload transport --- field for field]\label{rec:re-transport}
\begin{center}\small
\begin{tabular}{@{}ll@{}}
\textbf{\rulename{RootEval} data} & \textbf{\rulename{PConstExact} payload}\\
\midrule
$x_0$ \textnormal{(rational coordinates)} & substituted into the closed left-hand term\\
sublevel atom $(V, y)$ & the comparison\textquotesingle s two sides\\
comparator & $\le$ (inside the admissible set)\\
box atom \textnormal{(per-axis bounds)} & \emph{not} transported --- registry recompute (Recipe~\ref{rec:re-template})\\
\end{tabular}
\end{center}
\end{brecipe}

\begin{bprop}[the exactness identity]\label{prop:re-exact}
\[
x_0 \in \{x : V(x) \le y\}
\quad\Longleftrightarrow\quad
\mathrm{evalClosed}\bigl(V(x_0) \le y\bigr) = \mathsf{true}.
\]
\end{bprop}
\begin{proof}
Two lines, which is the point. Substituting a rational point into a
rational polynomial is exact arithmetic --- no rounding enters at any
step --- so the evaluator computes the value $V(x_0)$ itself, not an
approximation to it; and the decidable comparison of two exact
rationals is the membership predicate verbatim. This is the degenerate
--- and best --- case of the exactness identity an alias owes: not an
algebraic rearrangement but a literal identity of values. It is also
the whole rule: under naively rounded (non-directed) evaluation, a
point lying $\varepsilon$ outside the region could certify
membership, and the bridge rule\textquotesingle s instantiation (its
section, the extension page) would conclude satisfaction from a false
premise; directed outward rounding would keep that direction sound
but surrender the equivalence --- boundary points, where
$V(x_0) = y$ exactly, typically become undecidable (an
outward enclosure that straddles the threshold cannot decide; an
evaluation whose enclosure does not cross it --- an exact zero-width
value, or a one-sided enclosure touching the threshold --- still
decides).
\end{proof}

\begin{bprop}[soundness corollary]\label{prop:re-sound}
If the leaf checks, $x_0$ belongs to the root atom.
\end{bprop}
\begin{proof}
\texttt{Formula.holds\_of\_evalClosed\_true} (the target\textquotesingle s
kernel-verified evaluator theorem, quoted in the \rulename{PConstExact}
section): a closed comparison whose exact evaluation returns
$\mathsf{true}$ holds at every valuation. By
Proposition~\ref{prop:re-exact} that is membership: substitution
closed the formula, so $V(x_0) \le y$ mentions no wires,
its truth is valuation-independent, and that truth is literally the
sublevel atom\textquotesingle s membership predicate at $x_0$. There is no other
content --- the alias genre working as intended. (For the box branch
the recomputed per-coordinate containment plays the same role, warranted
by the recompute\textquotesingle s definition rather than a rule lemma.)
\end{proof}

\begin{bremark}[the payload-level elaborator, mechanized]
\label{rem:re-elab}
The \emph{semantic} content of the respelling is bound in the entry:
\texttt{rootEval\_eval\_cast} is
Proposition~\ref{prop:re-exact}\textquotesingle s literal identity of
values, and \texttt{rootEval\_sublevel\_iff} /
\texttt{rootEval\_box\_iff} are the two atom shapes as equivalences.
The payload-level elaborator --- previously this page\textquotesingle s
one absent artefact, held back until the extension certificate schema
existed --- is bound with it: \texttt{rootEvalElab} (module
\texttt{BernsteinLogic.QtslElaborators}) takes $(x_0, \text{atom})$ to
exactly the one \rulename{PConstExact} payload of
Recipe~\ref{rec:re-transport} (sublevel branch), as an inhabitant of
the target\textquotesingle s own payload type --- the comparator
confinement discharged at elaboration, once. The closed left-hand term
is built by verified substitution (\texttt{polyClosedTerm}, whose
evaluation lemma \texttt{polyClosedTerm\_evalQ} says the
checker\textquotesingle s exact replay computes
\emph{the value $V(x_0)$ itself}), \texttt{rootEvalFormula\_true\_iff}
is Proposition~\ref{prop:re-exact}\textquotesingle s identity at the
wire (the $\Q$ face --- both directions, as the proposition demands),
and \texttt{rootEvalElab\_sound} is the factoring through
\texttt{sound\_pConstExact}, itself one step to
\texttt{Formula.holds\_of\_evalClosed\_true}. The box branch is not
elaborated, by design (Recipe~\ref{rec:re-template}: registry
recompute, not a rule application).
\end{bremark}

\begin{bremark}[the inherited comparator confinement]
\rulename{PConstExact} excludes one comparator ($\neq$) from its
admissible set, for reasons of a second soundness invariant unpacked in
that rule\textquotesingle s section. The membership leaf uses $\le$
only --- comfortably inside --- but any port or mechanization of this
rule must \emph{preserve} the exclusion rather than re-derive the
comparator set from the semantic lemma alone; the flag travels with
the alias.
\end{bremark}

\begin{anchors}
target lemma (Prop.~\ref{prop:re-sound}) & \texttt{Formula.holds\_of\_evalClosed\_true}; \texttt{sound\_pConstExact} (quoted in the \rulename{PConstExact} section)\\
closedness implied at the leaf & \texttt{Formula.evalClosed} ($\mathsf{none}$ on any wire; same section)\\
exactness identities & \texttt{rootEval\_eval\_cast}, \texttt{rootEval\_sublevel\_iff}, \texttt{rootEval\_box\_iff} (bound in the entry)\\
payload-level elaborator (\autoref{rem:re-elab}) & \texttt{rootEvalElab} + \texttt{rootEvalElab\_sound}, over \texttt{polyClosedTerm} + \texttt{polyClosedTerm\_evalQ} + \texttt{rootEvalFormula\_true\_iff} (module \texttt{BernsteinLogic.QtslElaborators}, bound in the entry)\\
box-atom branch & checker-side registry recompute, by design --- no rule anchor\\
consumption site & the bridge rule\textquotesingle s eval field (\rulename{QSatFromRegion} section)\\
\end{anchors}


\rulesection{68}{BoxBern}{admissible rule}
\label{blcatrule-BoxBern}
\ruleblock{The rule (Steele metanotation)}
\sourcenote{as realized in the mechanized development; source row: \texttt{\texttt{QTSL-rules.toml}} --- \emph{BoxBern}. The metanotation statement is typeset by hand for notational quality; its mechanized ground truth is the Lean statement below, extracted from the sources at every build of this book and hash-bound. The rule \emph{roster} is machine-checked (every mechanized rule tag round-trips to exactly one catalogue row); the statement \emph{text} is editorial, checked by eye against the extracted Lean beside it.}
\begin{mathpar}
\inferrule*[Right=\rulename{BoxBern}]{
  N \in \Q[z_1, \ldots, z_n]; \quad D = \textstyle\prod_{v} [a_v, b_v] \subseteq \Q^n
  \\
  \varphi = (A,\ A^{-1},\ b): \quad A,\ A^{-1} : \Q^{[n] \times [n]},\ \ b : \Q^{[n]}; \qquad \varphi(w) = A\,w + b \ \text{(affine-rung witness; } \varphi \ \text{omitted reads as } \mathrm{id}\text{)}
  \\
  \vdash A^{-1} A = 1_{[n]} \ \text{(the one certified product, exact)}
  \\
  \hat N = \varphi \fatsemi N \ \text{(exact substitution)}
  \\
  \bigl\{b_\alpha \le 0\bigr\}_{\alpha} \ \text{(every tensor-product Bernstein coefficient of } \hat N \ \text{over } D\text{)}
}{
  D \vdash \hat N \le 0
}
\end{mathpar}

\ruleflavor{the witness is data with index structure --- two square rational tensors and an offset vector, the declared inverse carried rather than computed, its sole obligation the exact product (over a field, one side suffices) --- and the ladder above the affine rung is graded by each rung\textquotesingle s own grader: an unimplemented rung is a named refusing variant, never a partial check. The pullback $\hat N = \varphi \fatsemi N$ is recomputed by the checker and never rides the certificate, so a laundered pullback is unrepresentable rather than detected. The enacted rule is stated from its landed realization (\texttt{check\_box\_row}/\texttt{check\_wire\_family} in the workshop checker, over the wire-family grammar; refusal witness $(\alpha, b_\alpha)$) --- its certificate content is out of scope of the current baked theory: a box-cell row family is the named successor, and this cell moves onto the theory row when a revision carries it.}
\par\medskip
\rulename{BoxBern} is the tensor-product (box) sibling of
\rulename{BernCone}: where \rulename{BernCone} certifies a polynomial
predicate over a simplex-shaped region by exact Bernstein enclosure,
\rulename{BoxBern} certifies $\hat N \le 0$ over an axis-aligned
rational box, with the polynomial optionally pulled back through a
certified reparameterization witness first. A producer reaches for it
wherever certificate cells are boxes rather than simplices --- grid
covers, dyadic bisection cells, per-cell contraction rows for discrete
maps --- because the obligation count is the product of per-axis
degrees rather than the simplex count times a multinomial: at the
worked fixture\textquotesingle s shape (four variables, per-axis
degrees $(2,4,2,4)$, total degree $8$), one box costs $225$ exact
coefficients where the Kuhn-simplex decomposition of the same box
costs $24 \cdot 495 = 11{,}880$ --- a structural gap that compounds
with per-coefficient cost, measured at three orders of magnitude
wall-clock on the acceptance fixture.

What the checker recomputes, given the row $(N, \varphi?, D)$: that
the reparameterization witness certifies (for an affine witness,
shapes and the exact left-inverse identity $A^{-1}A = 1$; later ladder
rungs by their own graders; an unimplemented rung is a \emph{named
refusing variant}, never a partial check); the pullback
$\hat N = \varphi \fatsemi N$ by exact substitution --- no pullback ever
rides the certificate, so a laundered $\hat N$ is unrepresentable
rather than detected; and then the enclosure itself: substitute
$x_v = a_v + (b_v - a_v)\,t_v$ per axis, change basis per axis at the
polynomial\textquotesingle s actual degrees, and test every
tensor-product Bernstein coefficient against $\le 0$, refusing with
the first offending multi-index and its exact value. Everything is
exact rational arithmetic; nothing consults a float.

Each hypothesis is load-bearing. Drop the slot certification and an
uninvertible ``witness\textquotesingle\textquotesingle{} could smuggle
a coordinate change under which the claim means something else
entirely. Drop the recompute-the-pullback discipline and the checked
polynomial decouples from the claimed one. Drop the per-coefficient
sweep and the conclusion is unrelated to the box. The conclusion is
deliberately narrow --- one sign fact on one box --- because the row
\emph{species} that consumes it (coverage of the operative region,
S-procedure multipliers $\lambda \ge 0$, init/containment side rows)
owns every regional and compositional obligation; this rule is the
arithmetic leaf they all stand on.

\cataloguenote{The enclosure bound is conservative --- coefficients overshoot the
true range except at corners, where they interpolate --- and
completeness is recovered by subdivision exactly as for the simplex
rule: \rulename{Subdivide} remains the completeness partner. Degrees
enter the certified statement as \emph{data}: any degree family
admitting a valid tensor-Bernstein expansion encloses, so ``use the
actual degrees\textquotesingle\textquotesingle{} is a checker cost
policy, not a soundness premise. Degenerate axes ($a_v = b_v$) are
inside the statement --- the substitution sends the axis to a
constant and the enclosure is unaffected.
}
\ruleblock{Lean statement}
\declhead{BernsteinLogicAux.BoxBern}{boxBern\_enclosure}{\badge{BadgeGreen}{statement bound}}
\begin{leansrc}
/-- **The tensor-basis enclosure lemma** (soundness part (ii) of the
box-Bernstein rule): if the box-shifted claim polynomial
`N ∘ (lo + (hi − lo)·u)` has tensor-Bernstein coefficients `b`, then any
bounds on `b` bound `N` on the box. `hexp` is the checker's exact-`ℚ`
replay obligation (the shift and basis-conversion passes); the conclusion
quantifies over the landed arena's box set. -/
theorem boxBern_enclosure (lo hi : Fin n → ℚ) (N : MvPolynomial (Fin n) ℚ)
    (d : Fin n → ℕ) (b : (∀ v, Fin (d v + 1)) → ℚ) (m M : ℚ)
    (hexp : boxShift lo hi N = ∑ α, C (b α) * tensorBern d α)
    (hm : ∀ α, m ≤ b α) (hM : ∀ α, b α ≤ M) :
    ∀ x ∈ (ratBox lo hi).set, (m : ℝ) ≤ aeval x N ∧ aeval x N ≤ (M : ℝ)
\end{leansrc}
\declhead{BernsteinLogicAux.BoxBern}{boxBern\_enclosure\_minmax}{\badge{BadgeGreen}{statement bound}}
\begin{leansrc}
/-- The doc-verbatim min/max form (part (ii) as displayed):
`min_α b_α ≤ N(x) ≤ max_α b_α` on the box. -/
theorem boxBern_enclosure_minmax (lo hi : Fin n → ℚ)
    (N : MvPolynomial (Fin n) ℚ)
    (d : Fin n → ℕ) (b : (∀ v, Fin (d v + 1)) → ℚ)
    (hexp : boxShift lo hi N = ∑ α, C (b α) * tensorBern d α) :
    ∀ x ∈ (ratBox lo hi).set,
      ((Finset.univ.inf' Finset.univ_nonempty b : ℚ) : ℝ) ≤ aeval x N ∧
        aeval x N ≤ ((Finset.univ.sup' Finset.univ_nonempty b : ℚ) : ℝ)
\end{leansrc}
\declhead{BernsteinLogicAux.BoxBern}{boxBern\_nonpos}{\badge{BadgeGreen}{statement bound}}
\begin{leansrc}
/-- **The checker-facing sign corollary** (the rule's accept condition:
"accept iff every coefficient `b_α ≤ 0`"): nonpositive coefficients certify
`N ≤ 0` on the box — exactly one `hcells` premise of `Repar.t0_soundness` /
`t05_soundness_gridcx`, over the same `GridCx.Box` object. -/
theorem boxBern_nonpos (lo hi : Fin n → ℚ) (N : MvPolynomial (Fin n) ℚ)
    (d : Fin n → ℕ) (b : (∀ v, Fin (d v + 1)) → ℚ)
    (hexp : boxShift lo hi N = ∑ α, C (b α) * tensorBern d α)
    (hcoeff : ∀ α, b α ≤ 0) :
    ∀ x ∈ (ratBox lo hi).set, aeval x N ≤ 0
\end{leansrc}
\declhead{BernsteinLogicAux.BoxBern}{aeval\_boxShift}{\badge{BadgeGreen}{statement bound}}
\begin{leansrc}
/-- **Per-axis affine substitution exactness** (doc (ii) proof-shape item
"per-axis affine substitution exactness — a diagonal T0 map"): evaluating
the shifted polynomial at `u` is evaluating the original at the shifted
point. -/
theorem aeval_boxShift (lo hi : Fin n → ℚ) (u : Fin n → ℝ)
    (N : MvPolynomial (Fin n) ℚ) :
    aeval u (boxShift lo hi N)
      = aeval (fun v => (lo v : ℝ) + ((hi v : ℝ) - (lo v : ℝ)) * u v) N
\end{leansrc}
\declhead{BernsteinLogicAux.BoxBern}{sum\_aeval\_tensorBern}{\badge{BadgeGreen}{statement bound}}
\begin{leansrc}
/-- Sub-lemma 2c (partition of unity): the tensor basis sums to `1` on the
unit cube — "products of partitions of unity are partitions of unity"
(doc (ii)); univariate `bernsteinPolynomial.sum` distributed over the
product by `Fintype.prod_sum`. -/
theorem sum_aeval_tensorBern (d : Fin n → ℕ) (u : Fin n → ℝ) :
    ∑ α : ∀ v, Fin (d v + 1), aeval u (tensorBern d α) = 1
\end{leansrc}
\ruleblock{Correspondence (Steele display $\leftrightarrow$ Lean)}
\badge{BadgeGreen}{reviewer-certified against the displayed texts}
\sourcenote{The dictionary below is a reviewer-certified statement of how the mechanized statements above interpret the displayed rule: display metavariable $\mapsto$ Lean binder, each display hypothesis $\mapsto$ the premise or side condition carrying it, the display conclusion $\mapsto$ the goal --- with any deliberate difference in width between display and mechanization named outright. Certification is recorded as a hash triple over (the displayed cell, the extracted statement texts, this dictionary\textquotesingle s own text); if any of the three changes, this block turns amber until a reviewer re-certifies. The correspondence claim itself remains editorial --- reviewer-grade, not machine-checked; its machine-checked successor types each soundness statement against the rule as a structured object, and supersedes this block per rule as it lands.}
The displayed rule --- certified slot, exact pullback, coefficient
sweep, conclusion $\hat N \le 0$ on $D$ --- maps onto the bound
statements as follows. $D \mapsto$ \texttt{lo hi : Fin n → ℚ} through
\texttt{ratBox} (the same \texttt{GridCx.Box} object the transport
theorems quantify over); $\hat N \mapsto$ \texttt{N : MvPolynomial
(Fin n) ℚ} --- the \emph{pulled-back} polynomial, because the checker
recomputes the pullback and hands the enclosure lemma its result: the
slot legs (certification, replay) are consumed upstream by the
transport theorems (\texttt{t0\_soundness},
\texttt{t05\_soundness\_gridcx}, and the T1/TC rung lemmas), quoted
and bound in their own sections, never re-stated here. One
orientation divergence in the slot leg is named rather than left to
the reader: the enacted checker verifies the \emph{displayed} product
--- its single exact multiplication computes $A^{-1} A$ and compares
it with $1_{[n]}$ (\texttt{affine\_left\_certified}) --- whereas the
Lean witness structure stores the \emph{opposite} orientation
(\texttt{AffineWitness.certified} $: A \cdot A^{-1} = 1$) and derives
the displayed one as the theorem \texttt{left\_certified}
$: A^{-1} \cdot A = 1$ (via \texttt{mul\_eq\_one\_comm}; the matrices
are square over the commutative ring $\Q$, so either product
determines the other and the two orientations are equivalent). The
displayed coefficient table $\mapsto$ the expansion hypothesis of
\texttt{boxBern\_enclosure}: an explicit representation of the
\texttt{boxShift}-substituted polynomial as
$\sum_\alpha b_\alpha \prod_v B^{d_v}_{\alpha_v}$ --- the
u-space/x-space transport is a theorem (\texttt{aeval\_boxShift}),
not notation, and the degree family $d$ is a binder, not a fixed
choice. The displayed accept condition (every $b_\alpha \le 0$)
$\mapsto$ the hypothesis of \texttt{boxBern\_nonpos}, whose conclusion
is verbatim one \texttt{hcells} premise of the transport theorems ---
the two compose without glue, which is the certified content of
``exactly one premise of the landed transport
theorems.\textquotesingle\textquotesingle{} Width note, two-sided:
\texttt{boxBern\_enclosure\_minmax} is strictly stronger than the
accept condition needs (both range ends, where the rule consumes only
the max side), and nothing certified is wider than the display ---
the slot legs live upstream by construction, and no coverage or
multiplier obligation enters any binder here (the row species owns
them, as the display says).

\ruleblock{Lean verifier receipts}
\sourcenote{module \texttt{BernsteinLogicAux.BoxBern}; bound build \texttt{110ba4bda664\ldots} (bound from receipt 6935ac926dca961f, 2026-08-03T20:13:44Z). A receipt records that named checking layers accepted a build: (module, compiled-artifact SHA-256, environment SHA-256, checker layers), minted when the checking gate ran; the store receipts attesting the bound build are listed with the layers each run credited. Nothing here asks the reader to trust the store: every layer is an independently runnable program over the quoted sources --- rebuild the module and re-run the checkers to reproduce the verdict. Receipts attesting superseded builds credit no layer and are omitted from this rendering (the count of omissions, if any, is noted below).}
\begin{receipttable}
\receiptrow{63a3e12578528f22}{2026-08-02T02:06:52Z}{build+decorrelated}{\badge{BadgeGreen}{current}}
\receiptrow{c2fe5c957feb2507}{2026-08-02T23:30:40Z}{build+decorrelated}{\badge{BadgeGreen}{current}}
\receiptrow{665f63f0c0026da8}{2026-08-03T16:26:44Z}{build+decorrelated}{\badge{BadgeGreen}{current}}
\receiptrow{6935ac926dca961f}{2026-08-03T20:13:44Z}{build+decorrelated}{\badge{BadgeGreen}{current}}
\receiptrow{a3e7c09b848c5808}{2026-08-03T22:45:28Z}{build+decorrelated}{\badge{BadgeGreen}{current}}
\end{receipttable}
\sourcenote{8 further receipt(s) in the store attest superseded builds of this module and are omitted from this rendering.}
\layerline{\layeritem{build}{lake build (elaboration)}{\badge{BadgeGreen}{current}}\layeritem{replay}{leanchecker (official kernel replay)}{\badge{BadgeAmber}{stale only}}\layeritem{hardened}{safe\_verify (hardened screen)}{\badge{BadgeAmber}{stale only}}\layeritem{decorrelated}{lean4export + nanoda (decorrelated kernel)}{\badge{BadgeGreen}{current}}}
\sourcenote{\badge{BadgeGray}{comparator: none} no fixture in the three-implementation comparator harness (compiled / kernel-interpreted / reference-interpreted) exercises this rule.}
\ruleblock{HOL4 statement}
\absentbox{NO HOL STATEMENT YET}{The HOL family port (\texttt{checkFam\_boxbern} + a \texttt{famEq}
arm, the A1/A3 mechanism) is a named successor of the kernel leg that
landed this rule --- absence recorded at landing, with the mechanism
that will discharge it, so this reason cannot silently go stale the
way pre-wave absence reasons did.
}
\ruleblock{HOL verifier receipts}
\absentbox{NONE}{no HOL verification receipt exists for this rule.}
\ruleblock{Proof}

\sourcenote{Genre: admissible rule (tier 2) --- admissible by
\emph{reduction to verified checking machinery}, the taxonomy\textquotesingle s
recorded boundary case. The reduction shape is stated first; the
enclosure proposition it stands on is then proved in full, because it
is the section\textquotesingle s one piece of genuinely new
mathematics --- the tensor twin of the simplex grounding that
\rulename{BernCone} already carries.}

\begin{brecipe}[the reduction --- what the checker replays]
\label{rec:bb-reduction}
A \rulename{BoxBern} row is $(N,\ \varphi?,\ D)$: a claim polynomial
$N \in \Q[z_1, \dots, z_n]$, an optional reparameterization witness
$\varphi$ from the ladder ($\varnothing$ = identity), and an
axis-aligned box $D = \prod_v [a_v, b_v] \subseteq \Q^n$ in the
witness\textquotesingle s working coordinates. The discharge is three
replays, all exact $\Q$, each standing on a named verified theorem
rather than on catalogue structure:
\begin{enumerate}
\item \emph{slot certification} --- for an affine witness, shapes and
  the exact left-inverse identity $A^{-1}A = 1$ (a left inverse
  suffices over a field); later ladder rungs certify by their own
  graders, and a rung without a grader is a \emph{named refusing
  variant} --- refused loudly, never half-checked.
\item \emph{substitution replay} --- $\hat N = \varphi \fatsemi N$ by
  exact substitution. No pullback field exists in the row, so a forged
  $\hat N$ is unrepresentable rather than detected.
\item \emph{enclosure} --- the tensor-product Bernstein coefficients
  of $\hat N$ over $D$ (Proposition~\ref{prop:bb-enclosure}); accept
  iff every coefficient is $\le 0$, refusing with the witness
  $(\alpha, b_\alpha)$ otherwise.
\end{enumerate}
The conclusion, $\hat N \le 0$ on $D$, is deliberately exactly one
\texttt{hcells} premise of the landed transport theorems: coverage of
the operative region, S-procedure multipliers, and init/containment
side rows are owned by the consuming row species, never by this leaf.
\end{brecipe}

\begin{bprop}[tensor-product Bernstein enclosure]\label{prop:bb-enclosure}
Let $p \in \Q[x_1, \dots, x_n]$, let $D = \prod_v [a_v, b_v]$ with
rational endpoints ($a_v \le b_v$, degenerate axes allowed), and
suppose that after the per-axis affine substitution
$x_v = a_v + (b_v - a_v)\,u_v$ the resulting polynomial admits the
tensor-product Bernstein representation
\[
p\bigl(a + (b - a)\odot u\bigr) \;=\; \sum_{\alpha}\, b_\alpha
\prod_{v=1}^{n} B^{d_v}_{\alpha_v}(u_v),
\qquad
B^{d}_{i}(t) = \binom{d}{i}\, t^{i} (1-t)^{d-i},
\]
for some degree family $(d_v)$ and coefficients
$b_\alpha \in \Q$, $\alpha \in \prod_v \{0, \dots, d_v\}$. Then for
every $x \in D$,
\[
\min_\alpha b_\alpha \;\le\; p(x) \;\le\; \max_\alpha b_\alpha .
\]
In particular, if every $b_\alpha \le 0$ then $p \le 0$ on $D$.
\end{bprop}
\begin{proof}
Fix $x \in D$. Each coordinate lies in its interval, so there is a
$u \in [0,1]^n$ with $x_v = a_v + (b_v - a_v)u_v$ for every axis: on a
non-degenerate axis take $u_v = (x_v - a_v)/(b_v - a_v)$; on a
degenerate axis ($a_v = b_v$) the substitution is the constant $a_v$
regardless of $u_v$, and $u_v = 0$ serves. Substitution exactness ---
that evaluating $p$ at this $x$ equals evaluating the substituted
polynomial at $u$ --- is a per-axis affine replay, not a notational
identification (the bound statement keeps the two polynomials
distinct and transports along the substitution lemma).

On $[0,1]$, each Bernstein basis polynomial is nonnegative
($t^i \ge 0$, $(1-t)^{d-i} \ge 0$, binomial coefficients positive) and
each degree-$d$ family sums to $1$: $\sum_i B^d_i(t) =
\bigl(t + (1-t)\bigr)^d = 1$ by the binomial theorem. Products of
partitions of unity are partitions of unity, so the tensor family
$\bigl\{\prod_v B^{d_v}_{\alpha_v}(u_v)\bigr\}_\alpha$ is a family of
nonnegative weights summing to $1$ at every $u \in [0,1]^n$.
Evaluating the representation at $u$ therefore exhibits $p(x)$ as a
convex combination of the $b_\alpha$, and a convex combination lies
between the minimum and the maximum of its generators. The accept
condition is the corollary: if every generator is $\le 0$, so is
every convex combination.
\end{proof}

\begin{bremark}[degrees are data; degenerate axes are inside the statement]
Nothing in Proposition~\ref{prop:bb-enclosure} privileges the
polynomial\textquotesingle s actual per-axis degrees: \emph{any}
degree family admitting a valid representation encloses, so the
checker\textquotesingle s ``use the actual degrees after the
shift\textquotesingle\textquotesingle{} is a cost policy --- smallest
coefficient table --- and never a soundness premise. Likewise a
degenerate axis needs no exclusion: the substitution collapses it to
a constant, degrees drop, and the convex-combination argument is
untouched. Both points are formalization deltas made \emph{explicit}
in the bound statements, in the conservative direction --- the
expansion hypothesis speaks about the substituted polynomial (the
u-space/x-space transport is a theorem), and the degree family is a
binder.
\end{bremark}

\begin{bremark}[conservatism, and where completeness comes from]
The bound is not tight: coefficients overshoot the true range except
at the corners of $D$, where the representation interpolates
($B^d_i(0) = [i = 0]$, $B^d_i(1) = [i = d]$, so corner coefficients
\emph{equal} corner values). A true-but-refused row is repaired by
subdivision --- Bernstein enclosures converge under subdivision and
degree elevation --- so \rulename{Subdivide} is the completeness
partner here exactly as for the simplex rule, and refusal is always
the failure direction.
\end{bremark}

\begin{bremark}[why a box rule at all --- the obligation-count structure]
At per-axis degrees $(d_v)$, one box costs $\prod_v (d_v + 1)$
coefficients; the kernel-checkable simplex shape for the same box
costs $n!$ Kuhn simplices at $\binom{d + n}{n}$ coefficients each
(total degree $d$). At the acceptance fixture\textquotesingle s shape
--- $n = 4$, degrees $(2,4,2,4)$, $d = 8$ --- that is $3 \cdot 5
\cdot 3 \cdot 5 = 225$ against $24 \cdot \binom{12}{4} = 11{,}880$
per cell, and the per-coefficient work also drops (one-dimensional
passes against multinomial walks). The measured same-harness gap on
the fixture\textquotesingle s 26 cells is three orders of magnitude
wall-clock. None of this changes what is certified --- the rule is an
optimization of checker cost, not a widening of any claim.
\end{bremark}

\begin{anchors}
Recipe \ref{rec:bb-reduction} slot leg & \texttt{t0\_soundness} / \texttt{t05\_soundness\_gridcx} (their sections); later rungs per their graders, refusing variants named\\
Recipe \ref{rec:bb-reduction} replay leg & pullback recomputed, never carried --- laundering unrepresentable by row shape\\
Proposition \ref{prop:bb-enclosure} & \texttt{boxBern\_enclosure} (bound form), \texttt{boxBern\_enclosure\_minmax} (verbatim min/max form)\\
substitution exactness & \texttt{aeval\_boxShift} (per-axis affine replay; degenerate axes via the unit-scalar witness)\\
partition of unity & \texttt{sum\_aeval\_tensorBern} (tensor family sums to 1)\\
the accept corollary & \texttt{boxBern\_nonpos} --- conclusion verbatim one \texttt{hcells} premise of the transport theorems\\
\end{anchors}


\rulesection{69}{PExpTypeEnclosure}{admissible rule}
\label{blcatrule-PExpTypeEnclosure}
\ruleblock{The rule (Steele metanotation)}
\sourcenote{as realized in the mechanized development; source row: \texttt{\texttt{QTSL-rules.toml}} --- \emph{PExpTypeEnclosure}. The metanotation statement is typeset by hand for notational quality; its mechanized ground truth is the Lean statement below, extracted from the sources at every build of this book and hash-bound. The rule \emph{roster} is machine-checked (every mechanized rule tag round-trips to exactly one catalogue row); the statement \emph{text} is editorial, checked by eye against the extracted Lean beside it.}
\begin{mathpar}
\inferrule*[Right=\rulename{PExpTypeEnclosure}]{
  f \in F \quad n \in \mathbb{N} \quad r, R \in \Q
  \\
  a_k = a^f_k \ \text{for } k = 0..n \ \text{(exact coefficient recurrence)}
  \\
  0 \le r < n + 2
  \qquad
  R \;\ge\; \frac{r^{n+1}}{(n+1)!}\cdot\frac{n+2}{\,n+2-r\,}
}{
  \forall x,\ |x| \le r \;\Rightarrow\; \Bigl|f(x) - \textstyle\sum_{k \le n} a_k x^k\Bigr| \le R
}
\end{mathpar}

\ruleflavor{the coefficients are recomputed by the checker (never carried); the check is pure $\Q$-arithmetic and never evaluates $f$; the real functions enter only in the kernel-checked soundness theorem. Scope of $F$: entire functions with exact-$\Q$ Maclaurin coefficients dominated by $|a_k| \le 1/k!$ (entire of exponential type, normalized) --- $\exp$, $\sin$, $\cos$, and (qtsl tower) the reduced error-function kernel; a new in-class function costs one coefficient recurrence + one domination case, the tail argument untouched. Finite-radius functions ($\log$, rational powers) are certified by sibling routes, not by this rule; every downstream consumption decomposes into existing polynomial sequents (recentering, Taylor-model combination, endpoint substitution).}
\par\medskip
\rulename{PExpTypeEnclosure} is the calculus\textquotesingle s one
transcendental leaf: a remainder-bound rule that brackets a
non-polynomial function between two polynomials with all data rational,
so that every downstream obligation --- recentering, Taylor-model
combination, endpoint substitution into a generator inequality ---
decomposes into polynomial sequents the existing rules already check.
A producer reaches for it wherever certified dynamics carry
transcendental drift terms (pendulum-family $\sin/\cos$, rate and decay
$\exp$, Gaussian-kernel reductions): the rule certifies a Maclaurin
truncation once, at a radius, with an explicit remainder, and from
there the trusted arithmetic is exact rational computation exactly as
for every other leaf.

The witness is $(f, n, r, R)$: a function tag $f$ from the registered
library, a truncation degree $n$, a radius $r \in \Q_{\ge 0}$, and a
remainder budget $R \in \Q$. The check is pure rational arithmetic ---
the coefficients $a^f_k$ are recomputed from $(f, n)$ by exact per-tag
recurrence rather than carried (a laundered coefficient list is
unrepresentable on the wire), the radius must satisfy $r < n+2$, and
$R$ must dominate the closed-form tail bound
$\smash{\frac{r^{n+1}}{(n+1)!}\cdot\frac{n+2}{n+2-r}}$. The checker
never evaluates $f$; the real functions enter only in the
kernel-checked soundness theorem, which concludes the two-sided
bracket $|f(x) - \sum_{k \le n} a_k x^k| \le R$ on $[-r, r]$.

The scope is deliberately exact: \emph{entire functions with
exact-rational Maclaurin coefficients dominated by
$|a_k| \le 1/k!$} --- entire of exponential type, normalized so
$e^{|x|}$ dominates coefficientwise. All three conditions are
load-bearing: exact-$\Q$ recurrences make clause (i) checkable
term-by-term; the $1/k!$ domination is what lets ONE tail formula
serve every tag; entirety means any rational radius works, subject
only to the formula\textquotesingle s own guard $r < n+2$. This is
strictly narrower than ``anything with a closed-form Taylor
series\textquotesingle\textquotesingle: $\log$ and rational powers are
analytic but not entire (a finite-radius variant would need a
radius-ratio geometric tail and a domain side condition; they are
certified today by sibling routes built on their own bounds), and
$\tan$ is meromorphic (handled as $\sin/\cos$ plus division at the
consumer). Extending the library to a new in-class function costs one
coefficient recurrence and one domination case --- the tail argument
is untouched. The error-function kernel exercised exactly this clause:
under the $y = z^2$ reduction its coefficients are
$(-1)^k / (k!\,(2k+1))$, dominated by $1/k!$ verbatim, and the
soundness statement needed no edits beyond the per-tag cases.

\cataloguenote{Admissible by reduction to verified checking machinery (the same
boundary clause as \rulename{BoxBern}): the checker replays
\texttt{encCheck}, and the trust lands on the kernel-checked meaning
theorems, never on catalogue structure. Conservativity: no new sequent
theorems --- the rule concludes into a sibling judgment kind, and its
consumers re-enter the calculus as ordinary polynomial rows. Stated
outright rather than smoothed: this is the one rule where
transcendental meaning enters the trust story --- a deliberate,
priced, single-rule growth of the trusted base, carried by the
soundness theorems cited below. Slack is the producer\textquotesingle s
business (pick $n$, recenter, subdivide): subdivision shrinks bracket
width first-order while degree shrinks it factorially, and the
measured budgets on the acceptance systems ran at degree $9$ with
remainders near $10^{-7}$ --- orders below the certificate margins
they feed.
}
\ruleblock{Lean statement}
\declhead{BernsteinLogicAux.Enclosure}{encCheck\_sound}{\badge{BadgeGreen}{statement bound}}
\begin{leansrc}
/-- **The remainder-bound rule is sound** (note §1's boxed rule): a witness passing
`encCheck` certifies `|f(x) − P(x)| ≤ R` for every `|x| ≤ r`, where `f` is the tag's
real function and `P` the witness polynomial. -/
theorem encCheck_sound (w : EncWitness) (h : encCheck w = true) :
    ∀ x : ℝ, |x| ≤ (w.r : ℝ) →
      |w.tag.eval x - evalPoly w.coeffs x| ≤ (w.R : ℝ)
\end{leansrc}
\declhead{BernsteinLogicAux.Enclosure}{encCheck\_sound\_bracket}{\badge{BadgeGreen}{statement bound}}
\begin{leansrc}
/-- Two-sided bracket form: `P(x) − R ≤ f(x) ≤ P(x) + R` on `[−r, r]` — the pair of
polynomial brackets the certificate's sequents consume (note §3). -/
theorem encCheck_sound_bracket (w : EncWitness) (h : encCheck w = true) :
    ∀ x : ℝ, |x| ≤ (w.r : ℝ) →
      evalPoly w.coeffs x - (w.R : ℝ) ≤ w.tag.eval x ∧
      w.tag.eval x ≤ evalPoly w.coeffs x + (w.R : ℝ)
\end{leansrc}
\declhead{BernsteinLogicAux.Enclosure}{encCheck\_point}{\badge{BadgeGreen}{statement bound}}
\begin{leansrc}
/-- Point-bracket corollary (producer recentering, note §3(1)): evaluating the witness at
a rational point `c` with `|c| ≤ r` brackets the transcendental constant `f(c)`. -/
theorem encCheck_point (w : EncWitness) (h : encCheck w = true) (c : ℚ)
    (hc : |c| ≤ w.r) :
    evalPoly w.coeffs (c : ℝ) - (w.R : ℝ) ≤ w.tag.eval (c : ℝ) ∧
    w.tag.eval (c : ℝ) ≤ evalPoly w.coeffs (c : ℝ) + (w.R : ℝ)
\end{leansrc}
\declhead{BernsteinLogicAux.Enclosure}{linear\_interp\_le}{\badge{BadgeGreen}{statement bound}}
\begin{leansrc}
/-- Generator VCs are linear in each transcendental leaf occurrence: if `A + C·w ≤ 0`
holds at both endpoints of `[l, h]`, it holds for every `w ∈ [l, h]`. With this, per-cell
interval hypotheses for the enclosed leaf discharge as TWO ordinary polynomial sequents —
no new wire dimension. -/
theorem linear_interp_le {A C l h w : ℝ} (hlw : l ≤ w) (hwh : w ≤ h)
    (hl : A + C * l ≤ 0) (hh : A + C * h ≤ 0) : A + C * w ≤ 0
\end{leansrc}
\declhead{QtslAux.Enclosure}{encCheck\_sound}{\badge{BadgeGreen}{statement bound}}
\begin{leansrc}
theorem encCheck_sound (w : EncWitness) (h : encCheck w = true) :
    ∀ x : ℝ, |x| ≤ (w.r : ℝ) →
      |w.tag.eval x - evalPoly w.coeffs x| ≤ (w.R : ℝ)
\end{leansrc}
\declhead{QtslAux.Enclosure}{encCheck\_sound\_bracket}{\badge{BadgeGreen}{statement bound}}
\begin{leansrc}
theorem encCheck_sound_bracket (w : EncWitness) (h : encCheck w = true) :
    ∀ x : ℝ, |x| ≤ (w.r : ℝ) →
      evalPoly w.coeffs x - (w.R : ℝ) ≤ w.tag.eval x ∧
      w.tag.eval x ≤ evalPoly w.coeffs x + (w.R : ℝ)
\end{leansrc}
\declhead{QtslAux.Enclosure}{integral\_exp\_neg\_sq}{\badge{BadgeGreen}{statement bound}}
\begin{leansrc}
/-- The Gaussian integral reduces to the error-function kernel by the
substitution `t = z * s`: `∫ t in (0:ℝ)..z, exp (-(t ^ 2)) = z * T (z ^ 2)`.
Both sides are odd in `z`, and the oriented interval integral makes the
identity hold for every real `z` without a sign case split. Together with a
bracket for `2 / √π` (see `two_div_sqrt_pi_mem`) this carries polynomial
enclosures of `T` to enclosures of the error function. -/
theorem integral_exp_neg_sq (z : ℝ) :
    ∫ t in (0:ℝ)..z, Real.exp (-(t ^ 2)) = z * FnTag.erfT.eval (z ^ 2)
\end{leansrc}
\ruleblock{Correspondence (Steele display $\leftrightarrow$ Lean)}
\badge{BadgeGreen}{reviewer-certified against the displayed texts}
\sourcenote{The dictionary below is a reviewer-certified statement of how the mechanized statements above interpret the displayed rule: display metavariable $\mapsto$ Lean binder, each display hypothesis $\mapsto$ the premise or side condition carrying it, the display conclusion $\mapsto$ the goal --- with any deliberate difference in width between display and mechanization named outright. Certification is recorded as a hash triple over (the displayed cell, the extracted statement texts, this dictionary\textquotesingle s own text); if any of the three changes, this block turns amber until a reviewer re-certifies. The correspondence claim itself remains editorial --- reviewer-grade, not machine-checked; its machine-checked successor types each soundness statement against the rule as a structured object, and supersedes this block per rule as it lands.}
The displayed rule maps onto the bound statements as follows. The
witness $(f, n, r, R) \mapsto$ \texttt{EncWitness} (fields
\texttt{func}, \texttt{deg}, \texttt{radius}, \texttt{rem}, plus the
coefficient list governed by the displayed recurrence clause,
recompute semantics per the flavor text); the three displayed check clauses $\mapsto$ the definition
of \texttt{encCheck}, whose acceptance is the single hypothesis of
every soundness statement. The displayed conclusion $\mapsto$
\texttt{encCheck\_sound} (the absolute-value form) and
\texttt{encCheck\_sound\_bracket} (the same content as a two-sided
polynomial bracket --- equal width, different packaging);
\texttt{encCheck\_point} instantiates the bracket at a rational point
(the recentering step\textquotesingle s constant legs), and
\texttt{linear\_interp\_le} carries the endpoint-substitution step the
consumers use. One deliberate width difference, named outright: on the
wire the coefficient vector is never serialized --- the checker
recomputes it --- while the Lean witness carries the list explicitly
and clause (i) checks it against the recurrence; the wire grain is
therefore strictly narrower than the mechanized statement, and the
coefficient-mutation receipts live kernel-side precisely because those
mutations are unrepresentable on the wire. The two Lean packages prove
the same statement shape over different tag libraries (the qtsl
package adds the error-function kernel, with its meaning stated
against the Gaussian integral and bridged by
\texttt{integral\_exp\_neg\_sq}); the HOL statements are the
side-by-side port against \texttt{transcTheory}\textquotesingle s
natively per-index series --- reindexing-free exactly where the Lean
proof needed injection reindexing.

\ruleblock{Lean verifier receipts}
\sourcenote{module \texttt{BernsteinLogicAux.Enclosure}; bound build \texttt{aeef11c4ffdf\ldots} (bound from receipt c61f02086268c593, 2026-07-21T00:02:40Z). A receipt records that named checking layers accepted a build: (module, compiled-artifact SHA-256, environment SHA-256, checker layers), minted when the checking gate ran; the store receipts attesting the bound build are listed with the layers each run credited. Nothing here asks the reader to trust the store: every layer is an independently runnable program over the quoted sources --- rebuild the module and re-run the checkers to reproduce the verdict. Receipts attesting superseded builds credit no layer and are omitted from this rendering (the count of omissions, if any, is noted below).}
\begin{receipttable}
\receiptrow{ec9147f09c619209}{2026-07-17T21:43:40Z}{build+replay+hardened+decorrelated}{\badge{BadgeGreen}{current}}
\receiptrow{ef188baecdf2d5cc}{2026-07-19T18:23:51Z}{build+decorrelated}{\badge{BadgeGreen}{current}}
\receiptrow{74824f999e9c6c1a}{2026-07-20T01:21:03Z}{build+decorrelated}{\badge{BadgeGreen}{current}}
\receiptrow{c61f02086268c593}{2026-07-21T00:02:40Z}{build+decorrelated}{\badge{BadgeGreen}{current}}
\receiptrow{5f18c0fcd7885f0e}{2026-07-22T18:26:52Z}{build+decorrelated}{\badge{BadgeGreen}{current}}
\receiptrow{0b776138b6a8c8d7}{2026-07-31T21:59:44Z}{build+decorrelated}{\badge{BadgeGreen}{current}}
\receiptrow{2ee240defc1ac528}{2026-07-31T22:28:37Z}{build+decorrelated}{\badge{BadgeGreen}{current}}
\receiptrow{1fb846e38f09ba02}{2026-08-02T23:30:45Z}{build+decorrelated}{\badge{BadgeGreen}{current}}
\receiptrow{2f17504ab4610d57}{2026-08-03T16:26:49Z}{build+decorrelated}{\badge{BadgeGreen}{current}}
\receiptrow{c53cac58e2ed77d9}{2026-08-03T22:45:33Z}{build+decorrelated}{\badge{BadgeGreen}{current}}
\end{receipttable}
\sourcenote{module \texttt{QtslAux.Enclosure}; bound build \texttt{7de2b931c19f\ldots} (bound from receipt 840791fb258b3984, 2026-07-26T15:38:06Z). A receipt records that named checking layers accepted a build: (module, compiled-artifact SHA-256, environment SHA-256, checker layers), minted when the checking gate ran; the store receipts attesting the bound build are listed with the layers each run credited. Nothing here asks the reader to trust the store: every layer is an independently runnable program over the quoted sources --- rebuild the module and re-run the checkers to reproduce the verdict. Receipts attesting superseded builds credit no layer and are omitted from this rendering (the count of omissions, if any, is noted below).}
\begin{receipttable}
\receiptrow{969915618955be8f}{2026-07-25T11:27:08Z}{build+decorrelated}{\badge{BadgeGreen}{current}}
\receiptrow{840791fb258b3984}{2026-07-26T15:38:06Z}{build+replay+hardened+decorrelated}{\badge{BadgeGreen}{current}}
\end{receipttable}
\layerline{\layeritem{build}{lake build (elaboration)}{\badge{BadgeGreen}{current}}\layeritem{replay}{leanchecker (official kernel replay)}{\badge{BadgeGreen}{current}}\layeritem{hardened}{safe\_verify (hardened screen)}{\badge{BadgeGreen}{current}}\layeritem{decorrelated}{lean4export + nanoda (decorrelated kernel)}{\badge{BadgeGreen}{current}}}
\sourcenote{\badge{BadgeGray}{comparator: none} no fixture in the three-implementation comparator harness (compiled / kernel-interpreted / reference-interpreted) exercises this rule.}
\ruleblock{HOL4 statement}
\sourcenote{theory \texttt{bern\_enclosure}, HOL4 development script \texttt{\texttt{bern\_enclosureScript.sml}} \badge{BadgeGreen}{script bound}}
\holhead{Definition encCheck\_def}
\begin{holsrc}
Definition encCheck_def:
  encCheck w ⇔
    LENGTH w.coeffs = w.n + 1 ∧
    w.coeffs = GENLIST (maclaurin w.tag) (w.n + 1) ∧
    0 ≤ w.r ∧ w.r < &w.n + 2 ∧
    tailBound w.n w.r ≤ w.R
End
\end{holsrc}
\holhead{Theorem encCheck\_sound}
\begin{holsrc}
Theorem encCheck_sound:
  ∀w x.
    encCheck w ∧ abs x ≤ real_of_rat w.r ⇒
    abs (fnEval w.tag x - evalPolyR w.coeffs x) ≤ real_of_rat w.R
\end{holsrc}
\holhead{Theorem encCheck\_sound\_bracket}
\begin{holsrc}
Theorem encCheck_sound_bracket:
  ∀w x.
    encCheck w ∧ abs x ≤ real_of_rat w.r ⇒
    evalPolyR w.coeffs x - real_of_rat w.R ≤ fnEval w.tag x ∧
    fnEval w.tag x ≤ evalPolyR w.coeffs x + real_of_rat w.R
\end{holsrc}
\ruleblock{HOL verifier receipts}
\badge{BadgeBlue}{documentary} Honest partial, stated as such: the port deliberately carries the
MEANING theorem (\texttt{encCheck\_sound} against
\texttt{transcTheory}\textquotesingle s $\exp/\sin/\cos$), unlike the
neighbouring slices whose meaning stays Lean-side. Of the
theory\textquotesingle s fifteen theorems, the eleven foundation
theorems (datatypes through the factorial bound, including the full
rational-to-real cast bridge) were machine-verified in a batch build;
the four analysis theorems (\texttt{TERM\_LE},
\texttt{MAJORANT\_SUMS}, \texttt{TAIL\_MAIN},
\texttt{encCheck\_sound} with corollaries) are assembled from
interactively-verified tactic pieces, and a full in-file batch
re-verification currently exceeds the memory of every build host used
to date ($\sim$117\,GB resident observed at failure; a
$\ge$256\,GB host is the named path). Until that build lands, the HOL
half is a stated port with interactive verification, not a
batch-verified theory --- the badge here says exactly that. The
\texttt{encReceipts} table (seven rows, expected \texttt{TTTFFFF})
replays the same seven witnesses that appear byte-identically in the
Lean kernel receipts and the reference checker\textquotesingle s
battery: three accepts (including one with $R$ exactly at the tail
bound) and four single-mutation rejects.

\sourcenote{HOL receipts are documentary: quoted results of \texttt{hol-strict} runs
and census rows from the HOL4 development\textquotesingle s build
evidence, not minted receipt artefacts. A HOL receipt store (the Lean
store\textquotesingle s twin) is named future work.
}
\ruleblock{Proof}

\sourcenote{Genre: admissible rule (tier 2) --- admissible by
\emph{reduction to verified checking machinery}, the
taxonomy\textquotesingle s recorded boundary case. The reduction is
stated first; the tail bound it stands on is then proved in full,
because it is the section\textquotesingle s one piece of genuinely new
mathematics --- and the one place in the whole calculus where the
meaning of a transcendental function enters the trusted base.}

\begin{brecipe}[the reduction --- what the checker replays]
\label{rec:enc-reduction}
A \rulename{PExpTypeEnclosure} row is $(f,\ n,\ r,\ R)$: a function
tag $f$ from the registered library $F$, a truncation degree
$n \in \N$, a radius $r \in \Q_{\ge 0}$, and a remainder budget
$R \in \Q$. The discharge is three replays, all exact $\Q$:
\begin{enumerate}
\item \emph{coefficient recompute} --- the Maclaurin prefix
  $a_0, \dots, a_n$ is regenerated from $(f, n)$ by the
  tag\textquotesingle s exact rational recurrence; no coefficient ever
  rides the certificate, so a laundered coefficient list is
  unrepresentable rather than detected.
\item \emph{radius guard} --- $0 \le r < n + 2$, the side condition
  under which the closed-form tail bound below is finite and valid.
\item \emph{remainder comparison} --- $R \ge \mathrm{tail}(n, r)$
  where
  $\mathrm{tail}(n, r) = \dfrac{r^{n+1}}{(n+1)!}\cdot
   \dfrac{n+2}{\,n+2-r\,}$,
  a comparison of two explicit rationals.
\end{enumerate}
Nothing here evaluates $f$. The conclusion --- for all real $x$ with
$|x| \le r$, $\bigl|f(x) - \sum_{k \le n} a_k x^k\bigr| \le R$, i.e.\
the two-sided polynomial bracket $P - R \le f \le P + R$ on
$[-r, r]$ --- is the statement of the kernel-checked soundness
theorem, and that theorem is where the real $\exp$, $\sin$, $\cos$
(as the proof assistant\textquotesingle s own library defines them)
enter, once, for the whole calculus.
\end{brecipe}

\begin{bprop}[uniform Maclaurin tail bound]\label{prop:enc-tail}
Let $f$ be entire with Maclaurin coefficients satisfying
$|a_k| \le 1/k!$ for all $k$ (rationality of the $a_k$ is what makes
the \emph{check} exact; the bound itself needs only the domination).
For $n \in \N$ and $0 \le r < n + 2$,
\[
\sup_{|x| \le r}\ \Bigl| f(x) - \sum_{k \le n} a_k x^k \Bigr|
\;\le\;
\frac{r^{n+1}}{(n+1)!}\cdot\frac{n+2}{\,n+2-r\,}
\;=\; \mathrm{tail}(n, r).
\]
\end{bprop}

\begin{proof}
For $|x| \le r$ the omitted tail is dominated termwise:
\[
\Bigl|\sum_{k > n} a_k x^k\Bigr|
\;\le\; \sum_{k > n} \frac{r^k}{k!}.
\]
Write $N = n + 1$ and index the tail by $j \ge 0$ as $k = N + j$. The
factorial lower bound $(N + j)! \ge N!\,(N+1)^j$ --- induction on $j$:
the base is trivial, and the step multiplies by $N + j + 1 \ge N + 1$
--- gives
\[
\frac{r^{N+j}}{(N+j)!}
\;\le\; \frac{r^N}{N!}\Bigl(\frac{r}{N+1}\Bigr)^{j}.
\]
The guard $r < n + 2 = N + 1$ makes the ratio $q = r/(N+1)$ satisfy
$0 \le q < 1$, so the geometric series sums to
$1/(1 - q) = (N+1)/(N+1-r)$, and
\[
\sum_{k > n} \frac{r^k}{k!}
\;\le\; \frac{r^{N}}{N!}\cdot\frac{N+1}{\,N+1-r\,}
\;=\; \frac{r^{n+1}}{(n+1)!}\cdot\frac{n+2}{\,n+2-r\,}. \qedhere
\]
\end{proof}

Soundness of the rule is now the two-line composition the mechanized
statements carry: clause (1) certifies that the recomputed prefix
agrees with $f$\textquotesingle s true coefficients (a per-tag
recurrence identity, proven once per library member against the
library definition of $f$); every registered tag satisfies the
domination $|a_k| \le 1/k!$ (a per-tag lemma --- for $\exp$ an
identity, for $\sin/\cos$ the odd/even zero pattern, for the reduced
error-function kernel $|a_k| = 1/(k!\,(2k+1))$); and
Proposition~\ref{prop:enc-tail} turns clauses (2)--(3) into the
bracket. Accepting a witness with $R$ \emph{exactly} at
$\mathrm{tail}(n, r)$ is sound --- the receipt battery pins one such
row on purpose --- and every strengthening of the check (a smaller
$R$, a larger $r$, a shortened prefix, a flipped coefficient sign) is
refused by construction, which is what the four mutation receipts
witness.

\begin{bremark}[scope --- what the class is, and is not]
The hypotheses of Proposition~\ref{prop:enc-tail} are the exact scope
of the rule: entire, exact-$\Q$ coefficients, dominated by $1/k!$
(entire of exponential type, normalized so $e^{|x|}$ dominates
coefficientwise). Each condition excludes something familiar.
Finite-radius functions --- $\log$, rational powers --- fail
entirety: a variant for them would carry a radius-ratio geometric
tail \emph{and} a domain side condition, and they are certified today
by sibling routes built on their own bounds rather than by stretching
this rule. $\tan$ is meromorphic and out entirely; consumers build it
from $\sin/\cos$ and division. And $e^{x^2}$ is entire but not of
exponential type --- its coefficients escape the domination, so no
choice of $n$ rescues the one uniform tail formula. Conversely,
admitting a new in-class function costs exactly one coefficient
recurrence and one domination case; the tail argument, the check, and
the soundness statement shape are untouched --- the error-function
kernel entered this way, under a $y = z^2$ reduction chosen precisely
to land its coefficients inside the domination.
\end{bremark}

\begin{bremark}[slack is the producer\textquotesingle s business]
The rule certifies one truncation at one radius; tightness is policy,
not soundness. Recentering (the point brackets of the mechanized
\texttt{encCheck\_point}, combined by Taylor-model arithmetic over
$\Q$) and subdivision shrink bracket width first-order, while raising
the degree shrinks it factorially --- so producers buy accuracy with
$n$ and locality with cells, and every combination step re-enters the
calculus as ordinary polynomial rows. On the measured acceptance
systems, degree $9$ at unit radius gave remainders near $10^{-7}$,
orders below the certificate margins consuming them.
\end{bremark}


\blcatsection{Morse Topology}
\rulesection{70}{QMorseDischarge}{core rule}
\label{blcatrule-QMorseDischarge}
\ruleblock{The rule (Steele metanotation)}
\sourcenote{as realized in the mechanized development; source row: \texttt{\texttt{QTSL-rules.toml}} --- \emph{QMorseDischarge}. The metanotation statement is typeset by hand for notational quality; its mechanized ground truth is the Lean statement below, extracted from the sources at every build of this book and hash-bound. The rule \emph{roster} is machine-checked (every mechanized rule tag round-trips to exactly one catalogue row); the statement \emph{text} is editorial, checked by eye against the extracted Lean beside it.}
\begin{mathpar}
\inferrule*{
  L : \mathsf{QBlockLattice} \ \text{(blocks over one rational grid)} \\
  \bigl\{F_f \vdash \sigma_f\, \mathcal{A}\ell_f > 0\bigr\}_{f \in \mathrm{facets}(L)} \ \text{(facet-sign rows: } \ell_f \ \text{the facet\textquotesingle s normal coordinate, } \sigma_f \in \{\pm 1\} \ \text{the declared crossing sign, } F_f \ \text{the facet)} \\
  h : \mathsf{QCellMap} \ (\text{cell} \to \text{vertex} \sqcup \{\text{transit}\}) \\
  \bigl\{\vdash T(b) \subseteq \rho(h(b))\bigr\}_{h(b) \ne \text{transit}} \ \text{(containments, per vertex cell)} \\
  \bigl\{T(b) \vdash \mathcal{A} V_b + \varepsilon_b \le 0\bigr\}_{h(b) = \text{transit}}, \quad \varepsilon_b \in \Q_{>0} \ \text{(transits, per transit cell)} \\
  \bigl\{\vdash \rho(u) \cap \rho(v) = \emptyset\bigr\}_{u \ne v \in G}, \quad \bigl\{\vdash \rho(v) \subseteq D\bigr\}_{v \in G} \ \text{(} \rho\text{-wf: rows where not box-aligned, recomputed where box-aligned)} \\
  \mathrm{CR}(\mathcal{M}|_D) \subseteq \textstyle\bigcup_{h(b) \ne \text{transit}} T(b) \ \text{(cell-level chain-recurrence cover)} \\
  \mathsf{MorseAssembly}(L, h, G, \rho) \ \text{(the assembly ledger, defined below)}
}{
  \mathcal{M} \vDash \mathsf{Morse}(D, G, \rho)
}
\end{mathpar}

\ruleflavor{the lattice payload carries name, grid, arena, and blocks, and enters as certificate data (a shared reference; the rows impose the claims); its per-facet sign rows are displayed inline as the second premise --- each certifies a uniform strict sign of the flow across one grid facet ($\mathcal{A}\ell_f$ is the vector field\textquotesingle s facet-normal component on the facet, frozen at the gridline), and the transit certificates $V_b$ ride the same payload. The $\rho$-wf obligations split by alignment: box-aligned disjointness/containment is bounded interval arithmetic, recomputed; non-box-aligned pairs are per-pair BL rows. The displayed $\mathsf{MorseAssembly}$ package comprises the assembly ledger --- arena compactness, per-cell-pair crossing certificates on shared facets, orbit-interval order rows --- plus profile identification and the ambient metric structure chain recurrence needs (model-class); the chain-recurrence cover is displayed as its own premise. Checker-recomputed: lattice/cell-map presentation conformance (closure recompute, join-irreducible readout, the acyclicity check on the spec graph, the arena inside the top member\textquotesingle s realization --- the nonemptiness row).}
\par\medskip
\rulename{QMorseDischarge} concludes the global Morse judgment: all
recurrence of the model inside the arena lives in the declared regions
$\rho(v)$, and transitions between them run only down the declared
order. It is reduction-to-base in its cleanest form: the certificate is
a lattice of blocks over one rational grid whose every internal facet
carries a strictly-inward sign row of the base calculus, a map $h$
sending the lattice\textquotesingle s cells to spec vertices or to
``transit'', containment rows placing each vertex cell inside its
region, and per-transit-cell decrease rows $\mathcal{A}V + \varepsilon
\le 0$ --- every leaf a sequent the base kernel already checks, and
zero new kernel rules.

A producer reaches for it from the refinement loop: synthesize the
block family (splitting cells where a facet sign fails to certify,
merging where refinement is not needed), pick transit observables from
low-degree templates, and match cells to the spec graph. Failure mode
is refusal with a pointer: a spec whose resolution exceeds what the
certificate can contain reports exactly which cells obstruct, and that
report is itself the useful artefact.

What the checker recomputes rather than trusts: the lattice closure
under union and grid-level intersection; the readout --- the
join-irreducible members, their unique lower covers, the cells
$T(b) = \mathrm{cl}(b \smallsetminus b^{-})$, and the induced cell
order; convexity of $h$\textquotesingle s fibres and monotonicity of
$h$ on vertex cells; acyclicity of the spec graph; and the box-aligned
well-formedness checks on $\rho$. The witness the rule assembles is
$M(v) := \mathrm{Inv}\bigl(\bigcup_{h(b)=v} T(b)\bigr)$ --- the
invariant part of the fibre\textquotesingle s cell-union, not the
union of per-cell invariant parts; the proof section shows the
distinction is load-bearing.

Each hypothesis earns its place. Drop the containment rows and the
location clause fails. Drop the transit rows and a transit cell could
harbour recurrence the $\rho$-cover misses. Drop fibre convexity and
the judgment could assert a strict order the certificate does not
enforce. Drop facet strictness and three arguments die at once:
no-facet-riding, fibre disjointness, and the chain-margin bound
(``strictness pays a second time'').

The mechanization is conditional, honestly: the bound theorems prove
the witness assembly with a ledger of named hypothesis families ---
G3 (compactness of the invariant parts), discharged outright from
model-class facts; G4 (disjointness of distinct fibres), reduced to
transit certificates on fibre overlaps; G5 (the chain-recurrence
cover), reduced to the cell-level cover; and G6 (the order clause),
which remains a hypothesis of the statements bound here, its content
being the block lattice\textquotesingle s orbit combinatorics. The
proof section renders each as mathematics, together with G2, the
pointwise-to-integrated bridge that every certificate-shaped
hypothesis consumes.

\cataloguenote{The judgment is deliberately an ``at most'' claim: empty Morse sets
are permitted, so the rule certifies confinement of recurrence and of
transition order, never existence of structure --- the existence half
(Conley indices, connection witnesses) is a separately priced
annotation tier, absent by design. End-to-end certificate receipts
exist in the development\textquotesingle s exact-arithmetic
certificate-checker lane, which accepts both direct and
product-assembled Morse-family certificates against a mutant battery;
those receipts ride the checker\textquotesingle s own test suite, not
this book\textquotesingle s receipt store.
}
\ruleblock{Theory rows (geolog)}
\sourcenote{The rule\textquotesingle s certificate content as it is declared in the geolog theory the checker loads --- payload schema, side-condition demands, and axiom rows, quoted from the theory source at every build of this catalogue (a stale quotation cannot exist). Axiom modes: \texttt{chase} = the engine derives; \texttt{report} = a violation is a loud judgment; \texttt{refuse} = a violation gates admission. Side-condition relations are discharged by bounded evaluation at the engine\textquotesingle s ground seam, never by derivation.}
\par\medskip
The rule\textquotesingle s official certificate content in the geolog
theory the checker loads. \texttt{QMorse} is the concluded judgment
kind --- the graph-over-regions claim, with its arena --- and
\texttt{QMorseDischarge} is this rule\textquotesingle s payload
entity, field for field the display\textquotesingle s premise kit
(lattice, cell map, containment/transit/well-formedness proof-leaf
lists, and the per-transit-cell certificate polynomials in
\texttt{transit\_certs}). The arena tie is the quoted axiom: a
discharged claim\textquotesingle s arena \emph{is} the
lattice\textquotesingle s, as a stated equation between hash-consed
values. The R1 side condition is the demand relation
\texttt{qtsl\_morse\_readout} --- lattice closure and readout
recompute, discharged by bounded evaluation at the
engine\textquotesingle s ground seam. And the transit rows\textquotesingle
$\mathcal{A}V$ is not an ambient operator: \texttt{qtsl\_gen\_form} is
the derived generator-form seam tying the subject system and the
certificate component $V$ to the hypernet that says how $V$ can change
with duration --- recomputed, never trusted as stored data.

\declhead{theory Qtsl}{QMorse}{\badge{BadgeGreen}{quoted from the theory source at render}}
\begin{leansrc}
QMorse          : Set;   // entity: graph-over-regions claim
\end{leansrc}
\declhead{theory Qtsl}{QMorse/*}{\badge{BadgeGreen}{quoted from the theory source at render}}
\begin{leansrc}
QMorse/judgment : QMorse -> QJudgment;
QMorse/graph : QMorse -> QDag;
QMorse/rho   : QMorse -> QVertexRegionList;
QMorse/arena : QMorse -> QArena;
\end{leansrc}
\declhead{theory Qtsl}{QMorseDischarge}{\badge{BadgeGreen}{quoted from the theory source at render}}
\begin{leansrc}
QMorseDischarge : Set;   // entity: the lattice-reading discharge
\end{leansrc}
\declhead{theory Qtsl}{QMorseDischarge/*}{\badge{BadgeGreen}{quoted from the theory source at render}}
\begin{leansrc}
QMorseDischarge/proof : QMorseDischarge -> QProof;
QMorseDischarge/lattice : QMorseDischarge -> QBlockLattice;
QMorseDischarge/h : QMorseDischarge -> QCellMap;
QMorseDischarge/containments : QMorseDischarge -> QProofList;
QMorseDischarge/transits : QMorseDischarge -> QProofList;
QMorseDischarge/transit_certs : QMorseDischarge -> QCellPolyList;
QMorseDischarge/rho_wf : QMorseDischarge -> QProofList;
\end{leansrc}
\declhead{theory Qtsl}{ax/qmorsedischarge/arena}{\badge{BadgeGreen}{quoted from the theory source at render}}
\begin{leansrc}
chase ax/qmorsedischarge/arena : forall r : QMorseDischarge,
    p : QProof, j : QMorse, lat : QBlockLattice, k : QJudgment,
    ac : Shape, gc : QGridCx.
  r QMorseDischarge/proof = p,
  p QProof/concludes = k,
  j QMorse/judgment = k,
  r QMorseDischarge/lattice = lat,
  lat QBlockLattice/arena = ac,
  lat QBlockLattice/cx = gc
  |- j QMorse/arena = {cells: ac, cx: gc} qarMk;
\end{leansrc}
\declhead{theory Qtsl}{qtsl\_morse\_readout}{\badge{BadgeGreen}{quoted from the theory source at render}}
\begin{leansrc}
qtsl_morse_readout : {dag: QDag, h: QCellMap, lat: QBlockLattice,
                      rho: QVertexRegionList} -> Prop;
\end{leansrc}
\declhead{theory Qtsl}{qtsl\_gen\_form}{\badge{BadgeGreen}{quoted from the theory source at render}}
\begin{leansrc}
qtsl_gen_form : {sys: System, cert: HypernetRef, gen: HypernetRef}
  -> Prop;   // side-condition (oracle discharge / checker
\end{leansrc}
\ruleblock{Lean statement}
\declhead{BernsteinLogic.QtslSoundness}{qMorseDischarge\_sound}{\badge{BadgeGreen}{statement bound}}
\begin{leansrc}
/-- QMorseDischarge soundness — Lemma M(v), transcribed as the witness
assembly `M(v) := MaxInv(⋃_{h(b)=v} T(b))` (the note's own definition,
"the invariant part of the fiber's cell-union, NOT the union of per-cell
invariant parts" — its self-caught first-draft defect). PROVED outright:
the location clause (from the containment rows) and the invariance
clause (`maxInv_isInvariant`). INTERFACE (gaps G3–G6, each the named
prose paragraph of the note's M(v)/M(ii)/M(iii) proofs): compactness,
fiber disjointness, the chain-recurrence cover — into which the transit
rows enter via `transit_maxInv_empty` (M(iv), mechanized above) — and
the order clause. A successor leg mechanizing a gap discharges the
corresponding hypothesis; the statement is already the roster's. -/
theorem qMorseDischarge_sound (φ : Semiflow X) (D : Set X)
    (ρ : Vtx → Set X) (cells : ι → Set X) (hmap : ι → Option Vtx)
    (hContain : ∀ i v, hmap i = some v → cells i ⊆ ρ v)
    (hCompact : ∀ v, IsCompact (φ.MaxInv (fiberUnion cells hmap v)))
    (hDisj : ∀ v w, v ≠ w →
      Disjoint (φ.MaxInv (fiberUnion cells hmap v))
        (φ.MaxInv (fiberUnion cells hmap w)))
    (hCR : ∀ x ∈ D, ChainRecurrentIn φ D x →
      ∃ v, x ∈ φ.MaxInv (fiberUnion cells hmap v))
    (hOrder : ∀ γ : ℝ → X, φ.CompleteOrbit γ → Set.range γ ⊆ D →
      (∀ v, ¬ Set.range γ ⊆ φ.MaxInv (fiberUnion cells hmap v)) →
      ∃ v w, w < v ∧ alphaSet γ ⊆ φ.MaxInv (fiberUnion cells hmap v) ∧
        omegaSet γ ⊆ φ.MaxInv (fiberUnion cells hmap w)) :
    MorseSat φ D ρ
\end{leansrc}
\declhead{BernsteinLogic.QtslSoundness}{MorseWitness}{\badge{BadgeGreen}{statement bound}}
\begin{leansrc}
/-- A witness for `M ⊨_D Morse(G, ρ)` — the graded Morse-decomposition
semantics, clause for clause. `Vtx` carries the QDag's reachability
order (the edge set is presentation; v1 semantics consumes only the
reachability order); empty `M v` are permitted (the deliberate v1
"at most this structure" polarity). Order convention: source-to-sink,
`w < v` = connecting orbits may run from `M v` down to `M w`. -/
structure MorseWitness (φ : Semiflow X) (D : Set X) [PartialOrder Vtx]
    (ρ : Vtx → Set X) where
  /-- the family `{M(v)}` of Morse sets. -/
  M : Vtx → Set X
  /-- each compact… -/
  isCompact : ∀ v, IsCompact (M v)
  /-- …invariant… -/
  invariant : ∀ v, φ.IsInvariantIn (M v)
  /-- clause 1 (location): `M(v) ⊆ ρ(v)`. -/
  located : ∀ v, M v ⊆ ρ v
  /-- …pairwise disjoint. -/
  disjoint : ∀ v w, v ≠ w → Disjoint (M v) (M w)
  /-- clause 2 (completeness of the cover): every chain-recurrent point
  of `φ|_D` lies in some `M(v)`. -/
  crCover : ∀ x ∈ D, ChainRecurrentIn φ D x → ∃ v, x ∈ M v
  /-- clause 3 (order): a complete orbit in `D` contained in no `M(v)`
  connects a strictly higher Morse set to a strictly lower one. -/
  order : ∀ γ : ℝ → X, φ.CompleteOrbit γ → Set.range γ ⊆ D →
    (∀ v, ¬ Set.range γ ⊆ M v) →
    ∃ v w, w < v ∧ alphaSet γ ⊆ M v ∧ omegaSet γ ⊆ M w
\end{leansrc}
\declhead{BernsteinLogic.QtslSoundness}{MorseSat}{\badge{BadgeGreen}{statement bound}}
\begin{leansrc}
/-- The `QMorse` judgment's semantics: some family witnesses it. -/
def MorseSat (φ : Semiflow X) (D : Set X) [PartialOrder Vtx]
    (ρ : Vtx → Set X) : Prop
\end{leansrc}
\declhead{BernsteinLogic.QtslDischarge}{qMorseDischarge\_sound\_modelClass}{\badge{BadgeGreen}{statement bound}}
\begin{leansrc}
/-- `qMorseDischarge_sound` with the gap ledger at its post-discharge
state. Same conclusion, same rows; the interface hypotheses are now:

* model-class facts (the ODE-doctrine profile's warrants, per the
  System notion): spacewise-continuous flow, compact cells —
  G3 is GONE, these replace it;
* `hOverlap` — G4's residue: shared-facet overlaps carry transit
  certificates ("strictness pays a second time");
* `hCell` — G5's residue: the CELL-level chain-recurrence cover
  (M(iii), KMV-III, literature strength);
* `hTransitCert` — M(iv)'s input, now in certificate form;
* `hOrder` — G6, unchanged (wall named in the module header).

The witness assembly and the proved clauses (location, invariance) are
inherited from `qMorseDischarge_sound` untouched. -/
theorem qMorseDischarge_sound_modelClass (φ : Semiflow X) (D : Set X)
    (ρ : Vtx → Set X) (cells : ι → Set X) (hmap : ι → Option Vtx)
    (hflow : ∀ t, Continuous (φ.flow t))
    (hcells : ∀ i, IsCompact (cells i))
    (hContain : ∀ i v, hmap i = some v → cells i ⊆ ρ v)
    (hOverlap : ∀ v w, v ≠ w →
      φ.HasTransitCert
        (fiberUnion cells hmap v ∩ fiberUnion cells hmap w))
    (hCell : ∀ x ∈ D, ChainRecurrentIn φ D x → ∃ i, x ∈ φ.MaxInv (cells i))
    (hTransitCert : ∀ i, hmap i = none → φ.HasTransitCert (cells i))
    (hOrder : ∀ γ : ℝ → X, φ.CompleteOrbit γ → Set.range γ ⊆ D →
      (∀ v, ¬ Set.range γ ⊆ φ.MaxInv (fiberUnion cells hmap v)) →
      ∃ v w, w < v ∧ alphaSet γ ⊆ φ.MaxInv (fiberUnion cells hmap v) ∧
        omegaSet γ ⊆ φ.MaxInv (fiberUnion cells hmap w)) :
    MorseSat φ D ρ
\end{leansrc}
\declhead{BernsteinLogic.QtslDischarge}{maxInv\_isCompact}{\badge{BadgeGreen}{statement bound}}
\begin{leansrc}
/-- G3 discharged: `MaxInv` of a compact region is compact, given a
Hausdorff state space and spacewise-continuous flow maps — the "flow
continuity" the gap ledger names, and nothing else. Route: the set of
complete orbits inside `C` is a closed subset of the compact product
`∀ _ : ℝ, C` (Tychonoff; orbit constraints are equalizers of
continuous maps, closed in T2), and `MaxInv C` is its continuous image
under evaluation at 0. -/
theorem maxInv_isCompact [T2Space X] {φ : Semiflow X} {C : Set X}
    (hflow : ∀ t, Continuous (φ.flow t)) (hC : IsCompact C) :
    IsCompact (φ.MaxInv C)
\end{leansrc}
\declhead{BernsteinLogic.QtslDischarge}{fiberUnion\_isCompact}{\badge{BadgeGreen}{statement bound}}
\begin{leansrc}
/-- Finitely many compact cells have compact fiber unions. -/
theorem fiberUnion_isCompact {Vtx ι : Type*} [Finite ι]
    {cells : ι → Set X} {hmap : ι → Option Vtx}
    (hcells : ∀ i, IsCompact (cells i)) (v : Vtx) :
    IsCompact (fiberUnion cells hmap v)
\end{leansrc}
\declhead{BernsteinLogic.QtslDischarge}{HasTransitCert.disjoint\_maxInv}{\badge{BadgeGreen}{statement bound}}
\begin{leansrc}
/-- The G4 reduction ("strictness pays a second time"): if the overlap of two regions carries a transit
certificate, their maximal invariant sets are disjoint — a common
point would be forward-confined to the overlap by
`maxInv_forward_confined` applied on both sides. No backward
uniqueness, no injectivity of the flow maps. -/
theorem HasTransitCert.disjoint_maxInv {φ : Semiflow X} {C₁ C₂ : Set X}
    (h : φ.HasTransitCert (C₁ ∩ C₂)) :
    Disjoint (φ.MaxInv C₁) (φ.MaxInv C₂)
\end{leansrc}
\declhead{BernsteinLogic.QtslDischarge}{HasTransitCert.maxInv\_empty}{\badge{BadgeGreen}{statement bound}}
\begin{leansrc}
/-- Lemma M(iv) through the certificate interface: a transit-certified
region has empty maximal invariant set. (Shorter than the original
route: forward confinement suffices.) -/
theorem HasTransitCert.maxInv_empty {φ : Semiflow X} {C : Set X}
    (h : φ.HasTransitCert C) : φ.MaxInv C = ∅
\end{leansrc}
\declhead{BernsteinLogic.QtslDischarge}{crCover\_assemble}{\badge{BadgeGreen}{statement bound}}
\begin{leansrc}
/-- The cover-completeness paragraph of Lemma M(v), mechanized: the
fiber-level chain-recurrence cover follows from the CELL-level cover
(M(iii), the KMV-III margin argument — the residue, deliberately at
literature strength) plus transit certificates on transit cells
(M(iv), through `HasTransitCert.maxInv_empty`). -/
theorem crCover_assemble {φ : Semiflow X} {D : Set X}
    {cells : ι → Set X} {hmap : ι → Option Vtx}
    (hCell : ∀ x ∈ D, ChainRecurrentIn φ D x → ∃ i, x ∈ φ.MaxInv (cells i))
    (hTransit : ∀ i, hmap i = none → φ.HasTransitCert (cells i)) :
    ∀ x ∈ D, ChainRecurrentIn φ D x →
      ∃ v, x ∈ φ.MaxInv (fiberUnion cells hmap v)
\end{leansrc}
\declhead{BernsteinLogic.QtslOrderAssembly}{qMorseDischarge\_sound\_orderAssembled}{\badge{BadgeGreen}{statement bound}}
\begin{leansrc}
/-- `qMorseDischarge_sound_modelClass` with G6 at its residue: `hOrder`
is DERIVED (`orderClause_assemble`), and the interface carries instead
arena compactness (`hD`, a model-class fact: the arena is a compact
box-union), the two R1 shadows (`hConv`, `hMono`), and **G6a**
(`hInterval` — M(ii)'s orbit-interval conclusion, prose tier A; see
`orderClause_assemble`). No statement in `QtslSoundness.lean` or the
leg-5 module changes shape. -/
theorem qMorseDischarge_sound_orderAssembled (φ : Semiflow X) (D : Set X)
    (ρ : Vtx → Set X) (cells : ι → Set X) (hmap : ι → Option Vtx)
    (hD : IsCompact D)
    (hflow : ∀ t, Continuous (φ.flow t))
    (hcells : ∀ i, IsCompact (cells i))
    (hContain : ∀ i v, hmap i = some v → cells i ⊆ ρ v)
    (hOverlap : ∀ v w, v ≠ w →
      φ.HasTransitCert
        (fiberUnion cells hmap v ∩ fiberUnion cells hmap w))
    (hCell : ∀ x ∈ D, ChainRecurrentIn φ D x → ∃ i, x ∈ φ.MaxInv (cells i))
    (hTransitCert : ∀ i, hmap i = none → φ.HasTransitCert (cells i))
    (hConv : ∀ i j k v, i ≤ j → j ≤ k →
      hmap i = some v → hmap k = some v → hmap j = some v)
    (hMono : ∀ i j v w, i ≤ j →
      hmap i = some v → hmap j = some w → v ≤ w)
    (hInterval : ∀ γ : ℝ → X, φ.CompleteOrbit γ → Set.range γ ⊆ D →
      ∃ bα bω : ι, bω ≤ bα ∧
        alphaSet γ ⊆ cells bα ∧ omegaSet γ ⊆ cells bω ∧
        ∀ t : ℝ, ∃ i, bω ≤ i ∧ i ≤ bα ∧ γ t ∈ cells i) :
    MorseSat φ D ρ
\end{leansrc}
\declhead{BernsteinLogic.QtslEventualMember}{qMorseDischarge\_sound\_eventualMember}{\badge{BadgeGreen}{statement bound}}
\begin{leansrc}
/-- `qMorseDischarge_sound_orderAssembled` with G6a DERIVED: the
`hInterval` hypothesis is replaced by the block-family interface — the
five R1-recompute combinatorial rows, per-member forward invariance
(G1's class), closedness of join-irreducible realizations, and the
strict-inwardness shadows (F1's meet-confinement clause and the two
F2 absorption facts). No statement in `QtslSoundness.lean` or the
leg-5/leg-6 modules changes shape. What remains of G6 after this leg:
nothing under the name G6a — the residue classes are those already on
the ledger (G1b facet-row bridges per member, G2-class
pointwise-to-integrated for the shadows, R1 recompute conformance). -/
theorem qMorseDischarge_sound_eventualMember (φ : Semiflow X) (D : Set X)
    (ρ : Vtx → Set X) (cells : ι → Set X) (hmap : ι → Option Vtx)
    (blocks : β → Set X) (emb : ι → β) (cover : ι → β) (bD : β)
    (hD : IsCompact D)
    (hflow : ∀ t, Continuous (φ.flow t))
    (hcells : ∀ i, IsCompact (cells i))
    (hContain : ∀ i v, hmap i = some v → cells i ⊆ ρ v)
    (hOverlap : ∀ v w, v ≠ w →
      φ.HasTransitCert
        (fiberUnion cells hmap v ∩ fiberUnion cells hmap w))
    (hCell : ∀ x ∈ D, ChainRecurrentIn φ D x → ∃ i, x ∈ φ.MaxInv (cells i))
    (hTransitCert : ∀ i, hmap i = none → φ.HasTransitCert (cells i))
    (hConv : ∀ i j k v, i ≤ j → j ≤ k →
      hmap i = some v → hmap k = some v → hmap j = some v)
    (hMono : ∀ i j v w, i ≤ j →
      hmap i = some v → hmap j = some w → v ≤ w)
    (hbD : D ⊆ blocks bD)
    (hJoin : ∀ b : β, ∀ x ∈ blocks b,
      ∃ i : ι, x ∈ blocks (emb i) ∧ blocks (emb i) ⊆ blocks b)
    (hcover : ∀ i : ι, blocks (cover i) ⊂ blocks (emb i))
    (hcellsSup : ∀ i : ι, blocks (emb i) \ blocks (cover i) ⊆ cells i)
    (horder : ∀ i j : ι, blocks (emb i) ⊆ blocks (emb j) → i ≤ j)
    (hclosed : ∀ i : ι, IsClosed (blocks (emb i)))
    (hInv : ∀ b : β, φ.Invariant (blocks b))
    (hMeet : ∀ a b : β, ∃ c : β, blocks c ⊆ blocks a ∩ blocks b ∧
      ∀ γ : ℝ → X, φ.CompleteOrbit γ → Set.range γ ⊆ D →
        ∀ s t : ℝ, (∀ u, s ≤ u → u ≤ t → γ u ∈ blocks a ∩ blocks b) →
          ∀ u, s ≤ u → u ≤ t → γ u ∈ blocks c)
    (hAbsorbOmega : ∀ γ : ℝ → X, φ.CompleteOrbit γ → Set.range γ ⊆ D →
      ∀ b : β, (omegaSet γ ∩ blocks b).Nonempty →
        ∃ T : ℝ, ∀ t, T ≤ t → γ t ∈ blocks b)
    (hAbsorbAlpha : ∀ γ : ℝ → X, φ.CompleteOrbit γ → Set.range γ ⊆ D →
      ∀ b : β, (alphaSet γ ∩ blocks b).Nonempty →
        ∃ T : ℝ, ∀ t, t ≤ T → γ t ∈ blocks b) :
    MorseSat φ D ρ
\end{leansrc}
\declhead{BernsteinLogic.QtslCrossing}{qMorseDischarge\_sound\_cellCover}{\badge{BadgeGreen}{statement bound}}
\begin{leansrc}
/-- `qMorseDischarge_sound_modelClass` with G4 at its sharpened residue:
the fiber-overlap hypothesis `hOverlap` (one transit certificate on each
whole fiber-intersection — not the row shape) is replaced by per-cell-
pair crossing certificates `hCross` (the row shape: one certified facet
functional per shared facet). New model-class fact consumed: trajectory
continuity `htraj` — the same hypothesis G1's
`invariant_of_boundaryDwell` already takes. Everything else is
inherited unchanged; G5's and G6's residues are exactly as in
`qMorseDischarge_sound_modelClass`. -/
theorem qMorseDischarge_sound_cellCover (φ : Semiflow X) (D : Set X)
    (ρ : Vtx → Set X) (cells : ι → Set X) (hmap : ι → Option Vtx)
    (hflow : ∀ t, Continuous (φ.flow t))
    (htraj : ∀ x, Continuous fun t : ℝ≥0 => φ.flow t x)
    (hcells : ∀ i, IsCompact (cells i))
    (hContain : ∀ i v, hmap i = some v → cells i ⊆ ρ v)
    (hCross : ∀ v w, v ≠ w → ∀ i j, hmap i = some v → hmap j = some w →
      φ.HasCrossingCert (cells i ∩ cells j))
    (hCell : ∀ x ∈ D, ChainRecurrentIn φ D x → ∃ i, x ∈ φ.MaxInv (cells i))
    (hTransitCert : ∀ i, hmap i = none → φ.HasTransitCert (cells i))
    (hOrder : ∀ γ : ℝ → X, φ.CompleteOrbit γ → Set.range γ ⊆ D →
      (∀ v, ¬ Set.range γ ⊆ φ.MaxInv (fiberUnion cells hmap v)) →
      ∃ v w, w < v ∧ alphaSet γ ⊆ φ.MaxInv (fiberUnion cells hmap v) ∧
        omegaSet γ ⊆ φ.MaxInv (fiberUnion cells hmap w)) :
    MorseSat φ D ρ
\end{leansrc}
\ruleblock{Correspondence (Steele display $\leftrightarrow$ Lean)}
\badge{BadgeAmber}{a bound Lean statement CHANGED since certification --- re-review and re-certify}
\sourcenote{The dictionary below is a reviewer-certified statement of how the mechanized statements above interpret the displayed rule: display metavariable $\mapsto$ Lean binder, each display hypothesis $\mapsto$ the premise or side condition carrying it, the display conclusion $\mapsto$ the goal --- with any deliberate difference in width between display and mechanization named outright. Certification is recorded as a hash triple over (the displayed cell, the extracted statement texts, this dictionary\textquotesingle s own text); if any of the three changes, this block turns amber until a reviewer re-certifies. The correspondence claim itself remains editorial --- reviewer-grade, not machine-checked; its machine-checked successor types each soundness statement against the rule as a structured object, and supersedes this block per rule as it lands.}
The displayed judgment $\mathcal{M} \vDash \mathsf{Morse}(D, G, \rho)$
$\mapsto$ the goal \texttt{MorseSat φ D ρ} = \texttt{Nonempty
(MorseWitness φ D ρ)} --- the Morse note \S2.2 clauses (compactness,
invariance, location, disjointness, chain-recurrence cover, order)
packaged as a structure, so the certified conclusion is existential
exactly as the judgment reads; the witness assembly $M(v) :=
\mathrm{MaxInv}\bigl(\bigcup_{h(b)=v} T(b)\bigr)$ lives inside the
proof term of \texttt{qMorseDischarge\_sound}, never in its statement.
The display\textquotesingle s spec graph $G$ enters as the binder pair
\texttt{Vtx} + \texttt{[PartialOrder Vtx]}: the semantics consumes
only the reachability order (the \texttt{MorseWitness} docstring says
so in words), and the dag check on the presentation\textquotesingle s
edge set is R1 recompute with no Lean counterpart --- correctly.
Order polarity carried exactly: \texttt{hOrder}\textquotesingle s
conclusion places $\alpha$-limits in the strictly \emph{higher} set
and $\omega$-limits in the lower (\texttt{w < v ∧ alphaSet γ ⊆ \ldots
v ∧ omegaSet γ ⊆ \ldots w}), the display\textquotesingle s
``transitions run only down the declared order.''
\begin{itemize}
\item the E2 lattice $L$ and cell map $h$ $\mapsto$ the abstract
  binders \texttt{cells : ι → Set X} and \texttt{hmap : ι → Option
  Vtx} (\texttt{none} = transit, the display\textquotesingle s
  $\sqcup\,\{\text{transit}\}$): the whole lattice stratum ---
  closure recompute, join-irreducible readout, $T(b)$, fibre
  convexity and monotonicity of $h$, box-aligned $\rho$-wf --- is
  presentation the checker recomputes (side condition R1), and the
  bound statements never see it; they quantify over arbitrary cell
  families, strictly wider than the display, in the safe direction.
\item containment rows $\mapsto$ \texttt{hContain} (\texttt{hmap i =
  some v → cells i ⊆ ρ v}), one-to-one at the semantic stratum.
\item transit rows $\mathcal{A}V + \varepsilon \le 0$ $\mapsto$
  \texttt{hTransitCert i} (per \texttt{hmap i = none}) in the
  semantic-shadow form \texttt{HasTransitCert}: some observable $V$,
  bounded below on the cell, with explicit margin $\varepsilon > 0$
  and the \emph{integrated} decrease clause $V(\Phi_t x) \le V(x) -
  \varepsilon t$ along confined prefixes. The pointwise-to-integrated
  bridge (gap G2; the \texttt{FlowDerivOn} profile identification) is
  outside this bound set --- the roster witness cites its
  mechanization (\texttt{QtslMeanValue.lean}), not bound here.
\item $\rho$-wf rows $\mapsto$ nothing: no bound hypothesis consumes
  them (fibre disjointness arrives through \texttt{hOverlap}
  instead) --- R1 territory, like the lattice.
\item the conditional ledger, exactly as bound:
  \texttt{qMorseDischarge\_sound} carries G3--G6 as hypotheses
  (\texttt{hCompact}/\texttt{hDisj}/\texttt{hCR}/\texttt{hOrder});
  \texttt{\_modelClass} discharges G3 outright
  (\texttt{maxInv\_isCompact}: Tychonoff, consuming \texttt{[T2Space
  X]}, spacewise flow continuity \texttt{hflow}, \texttt{[Finite ι]}
  + compact cells via \texttt{fiberUnion\_isCompact} --- model-class
  facts all), reduces G4 to overlap transit certificates
  (\texttt{HasTransitCert.disjoint\_maxInv}: forward confinement on
  both sides, no backward uniqueness), reduces G5 to the cell-level
  cover \texttt{hCell} (M(iii), KMV-III, literature strength) through
  \texttt{crCover\_assemble} + \texttt{HasTransitCert.maxInv\_empty}
  (M(iv)), and carries G6 (\texttt{hOrder}) untouched.
\end{itemize}
Width note, both directions stated. Three assembly statements extend
the bound set (the three assembly blocks below): the
order assembly (\texttt{qMorseDischarge\_sound\_orderAssembled},
module \texttt{QtslOrderAssembly}) derives \texttt{hOrder} --- G6
proper --- carrying instead arena compactness, the two R1 shadows, and
G6a, the orbit-interval row; the eventual-member assembly
(\texttt{qMorseDischarge\_sound\_eventualMember}, module
\texttt{QtslEventualMember}) derives G6a in turn from the block-family
interface (the five R1-recompute combinatorial rows, per-member
forward invariance, closedness, strict-inwardness shadows), leaving
nothing of G6 under its own name; and the crossing assembly
(\texttt{qMorseDischarge\_sound\_cellCover}, module
\texttt{QtslCrossing}) sharpens G4 to the row shape --- per-cell-pair
crossing certificates on shared facets in place of whole-fibre-overlap
transit certificates. All three are hash-bound, so the certified
reach extends past \texttt{\_modelClass}\textquotesingle s ledger and
the conditional residue shrinks to the ledger\textquotesingle s
G1b/G2-class profile identifications (the
$\mathrm{FlowDerivOn}\,V\,(\mathcal{A}V)$ chain-rule bridges --- the
mean-value and Dini profile strata remain cited theorems, not bound),
the cell-level chain-recurrence cover --- now a displayed premise
($\mathrm{CR}(\mathcal{M}|_D) \subseteq \bigcup T(b)$ over vertex
cells) --- $\mapsto$ \texttt{hCell} (the fibre-to-cell margin argument
is the proof\textquotesingle s stated residue), and the displayed
$\mathsf{MorseAssembly}$ package $\mapsto$ the remaining
crossing/order/interface hypotheses plus recompute
conformance. Binder cargo: \texttt{[PseudoMetricSpace
X]} is consumed by chain recurrence itself --- $(\varepsilon,T)$-chains
need distances --- ambient structure the assembly package carries.

\ruleblock{Lean verifier receipts}
\sourcenote{module \texttt{BernsteinLogic.QtslSoundness}; bound build \texttt{77005991d61d\ldots} (bound from receipt 8130a7a7a679cef4, 2026-08-04T01:16:06Z). A receipt records that named checking layers accepted a build: (module, compiled-artifact SHA-256, environment SHA-256, checker layers), minted when the checking gate ran; the store receipts attesting the bound build are listed with the layers each run credited. Nothing here asks the reader to trust the store: every layer is an independently runnable program over the quoted sources --- rebuild the module and re-run the checkers to reproduce the verdict. Receipts attesting superseded builds credit no layer and are omitted from this rendering (the count of omissions, if any, is noted below).}
\begin{receipttable}
\receiptrow{8130a7a7a679cef4}{2026-08-04T01:16:06Z}{build+decorrelated}{\badge{BadgeGreen}{current}}
\receiptrow{9eb3706354082144}{2026-08-04T02:23:12Z}{build+decorrelated}{\badge{BadgeGreen}{current}}
\end{receipttable}
\sourcenote{5 further receipt(s) in the store attest superseded builds of this module and are omitted from this rendering.}
\sourcenote{module \texttt{BernsteinLogic.QtslDischarge}; bound build \texttt{4aa46910a600\ldots} (bound from receipt dd2946805f790f99, 2026-08-04T01:16:09Z). A receipt records that named checking layers accepted a build: (module, compiled-artifact SHA-256, environment SHA-256, checker layers), minted when the checking gate ran; the store receipts attesting the bound build are listed with the layers each run credited. Nothing here asks the reader to trust the store: every layer is an independently runnable program over the quoted sources --- rebuild the module and re-run the checkers to reproduce the verdict. Receipts attesting superseded builds credit no layer and are omitted from this rendering (the count of omissions, if any, is noted below).}
\begin{receipttable}
\receiptrow{dd2946805f790f99}{2026-08-04T01:16:09Z}{build+decorrelated}{\badge{BadgeGreen}{current}}
\receiptrow{e1bc68890e2aebf1}{2026-08-04T02:23:11Z}{build+decorrelated}{\badge{BadgeGreen}{current}}
\end{receipttable}
\sourcenote{14 further receipt(s) in the store attest superseded builds of this module and are omitted from this rendering.}
\sourcenote{module \texttt{BernsteinLogic.QtslOrderAssembly}; bound build \texttt{1fef28f78168\ldots} (bound from receipt f6e3628b18a7ed89, 2026-07-20T08:27:03Z). A receipt records that named checking layers accepted a build: (module, compiled-artifact SHA-256, environment SHA-256, checker layers), minted when the checking gate ran; the store receipts attesting the bound build are listed with the layers each run credited. Nothing here asks the reader to trust the store: every layer is an independently runnable program over the quoted sources --- rebuild the module and re-run the checkers to reproduce the verdict. Receipts attesting superseded builds credit no layer and are omitted from this rendering (the count of omissions, if any, is noted below).}
\begin{receipttable}
\receiptrow{b066c7d1be7303f1}{2026-07-19T04:29:30Z}{build+decorrelated}{\badge{BadgeGreen}{current}}
\receiptrow{ef9ccb70d0449d3b}{2026-07-19T09:07:23Z}{build+decorrelated}{\badge{BadgeGreen}{current}}
\receiptrow{f6e3628b18a7ed89}{2026-07-20T08:27:03Z}{build+decorrelated}{\badge{BadgeGreen}{current}}
\receiptrow{08fb682d6ccdcbc7}{2026-07-20T20:44:50Z}{build+decorrelated}{\badge{BadgeGreen}{current}}
\receiptrow{62d115665ebe5578}{2026-07-20T21:01:19Z}{build+replay+hardened+decorrelated}{\badge{BadgeGreen}{current}}
\receiptrow{b2fd112afe983043}{2026-07-21T05:28:46Z}{build+replay+hardened+decorrelated}{\badge{BadgeGreen}{current}}
\end{receipttable}
\sourcenote{module \texttt{BernsteinLogic.QtslEventualMember}; bound build \texttt{9282e7682307\ldots} (bound from receipt 839bcd47d53458bf, 2026-07-20T08:27:05Z). A receipt records that named checking layers accepted a build: (module, compiled-artifact SHA-256, environment SHA-256, checker layers), minted when the checking gate ran; the store receipts attesting the bound build are listed with the layers each run credited. Nothing here asks the reader to trust the store: every layer is an independently runnable program over the quoted sources --- rebuild the module and re-run the checkers to reproduce the verdict. Receipts attesting superseded builds credit no layer and are omitted from this rendering (the count of omissions, if any, is noted below).}
\begin{receipttable}
\receiptrow{d9e8d15537b8c931}{2026-07-19T11:09:44Z}{build+decorrelated}{\badge{BadgeGreen}{current}}
\receiptrow{dcd494c897af133b}{2026-07-19T11:29:34Z}{build+decorrelated}{\badge{BadgeGreen}{current}}
\receiptrow{839bcd47d53458bf}{2026-07-20T08:27:05Z}{build+decorrelated}{\badge{BadgeGreen}{current}}
\receiptrow{0b348bff26bc835c}{2026-07-20T20:44:43Z}{build+decorrelated}{\badge{BadgeGreen}{current}}
\receiptrow{fbf234c5b555fbd6}{2026-07-20T21:01:12Z}{build+replay+hardened+decorrelated}{\badge{BadgeGreen}{current}}
\receiptrow{a67d63ba90d91fb5}{2026-07-21T05:28:37Z}{build+replay+hardened+decorrelated}{\badge{BadgeGreen}{current}}
\end{receipttable}
\sourcenote{2 further receipt(s) in the store attest superseded builds of this module and are omitted from this rendering.}
\sourcenote{module \texttt{BernsteinLogic.QtslCrossing}; bound build \texttt{29a1be816112\ldots} (bound from receipt 3dec8934e00e8465, 2026-07-20T08:27:06Z). A receipt records that named checking layers accepted a build: (module, compiled-artifact SHA-256, environment SHA-256, checker layers), minted when the checking gate ran; the store receipts attesting the bound build are listed with the layers each run credited. Nothing here asks the reader to trust the store: every layer is an independently runnable program over the quoted sources --- rebuild the module and re-run the checkers to reproduce the verdict. Receipts attesting superseded builds credit no layer and are omitted from this rendering (the count of omissions, if any, is noted below).}
\begin{receipttable}
\receiptrow{58226ade60f5d15c}{2026-07-19T04:33:17Z}{build+decorrelated}{\badge{BadgeGreen}{current}}
\receiptrow{8fa273e016e7de3b}{2026-07-19T09:07:20Z}{build+decorrelated}{\badge{BadgeGreen}{current}}
\receiptrow{3dec8934e00e8465}{2026-07-20T08:27:06Z}{build+decorrelated}{\badge{BadgeGreen}{current}}
\receiptrow{204c60d7cf1cf318}{2026-07-20T20:44:39Z}{build+decorrelated}{\badge{BadgeGreen}{current}}
\receiptrow{574f67e1c43258ae}{2026-07-20T21:01:08Z}{build+replay+hardened+decorrelated}{\badge{BadgeGreen}{current}}
\receiptrow{5e4d769e0abef0a1}{2026-07-21T05:28:33Z}{build+replay+hardened+decorrelated}{\badge{BadgeGreen}{current}}
\end{receipttable}
\layerline{\layeritem{build}{lake build (elaboration)}{\badge{BadgeGreen}{current}}\layeritem{replay}{leanchecker (official kernel replay)}{\badge{BadgeGreen}{current}}\layeritem{hardened}{safe\_verify (hardened screen)}{\badge{BadgeGreen}{current}}\layeritem{decorrelated}{lean4export + nanoda (decorrelated kernel)}{\badge{BadgeGreen}{current}}}
\sourcenote{\badge{BadgeGray}{comparator: none} no fixture in the three-implementation comparator harness (compiled / kernel-interpreted / reference-interpreted) exercises this rule.}
\ruleblock{HOL4 statement}
\absentbox{NO HOL STATEMENT YET}{Absent, and honestly so at two levels. The HOL4
development\textquotesingle s object is the base-kernel checker cone
--- the first-order computational surface of the sequent-stratum
certificate contract (dtype factorization, exact rational arithmetic,
the catalogue check functions) --- and the satisfaction-judgment
stratum this rule concludes in has no HOL twin: no Morse-family or
judgment-stratum definition exists in the ported scripts (verified
against the port\textquotesingle s definition inventory). The base
leaves this rule consumes --- facet-sign, containment, and decrease
rows --- ARE covered by the ported checker; what has no HOL
counterpart is the judgment-stratum soundness development itself,
which today is Lean-only. A port would be new assurance surface over
the semiflow stratum, not a translation of existing checker code.
}
\ruleblock{HOL verifier receipts}
\absentbox{NONE}{no HOL verification receipt exists for this rule.}
\ruleblock{Proof}

\sourcenote{Genre: core rule (tier 1) --- an ordinary mathematical proof
at Putnam-exam-solution detail; the largest mechanized lemma set in this
book. The rule's mechanization is \emph{conditional}: the soundness
theorem is proved with a ledger of four named hypothesis families, of
which one is discharged outright below, two are reduced to
transit-certificate form, and the fourth --- the order clause --- is
carried as an explicit hypothesis by the discharged-floor theorem and
discharged by the separately bound assembly legs (ledger remark). The
ledger remark near the end states each as mathematics. Every hypothesis is used and remarked on; the formal-anchor
table maps each step to the Lean declaration that carries it.}

\textbf{Setting.}\enspace Fix a forward semiflow $\varphi$ on a metric
state space $X$ ($\Phi_0 = \mathrm{id}$,
$\Phi_{s+t} = \Phi_t \circ \Phi_s$), and a compact forward-invariant
arena $D \subseteq X$. Four standard notions, exactly as the mechanized
statements carry them:

\emph{(complete orbits)} $\gamma : \R \to X$ is a complete orbit when
$\gamma(s + t) = \Phi_t(\gamma(s))$ for all $s \in \R$, $t \geq 0$ ---
a two-sided trajectory compatible with the forward flow.

\emph{(maximal invariant part)} for $C \subseteq X$,
\[
\mathrm{Inv}(C) \;:=\; \{\, x \ :\ \exists\, \gamma \ \text{complete
orbit},\ \gamma(0) = x,\ \mathrm{range}\,\gamma \subseteq C \,\}.
\]
An \emph{invariant set} $S$ (the sense fit for a semiflow) is one that
is forward-invariant and satisfies $S \subseteq \mathrm{Inv}(S)$ ---
every point rides a complete orbit inside.

\emph{(limit sets)} $\omega(\gamma) := \bigcap_{T}
\mathrm{cl}\,\gamma([T, \infty))$ and $\alpha(\gamma) := \bigcap_{T}
\mathrm{cl}\,\gamma((-\infty, T])$.

\emph{(chains)} an $(\varepsilon, T)$-chain from $x$ to $y$ inside $D$
is a finite concatenation of orbit segments of duration $\geq T$ with
jumps of size $< \varepsilon$ at points of $D$ (generated by: a single
step $x, y \in D$, $t \geq T$,
$\mathrm{dist}(\Phi_t x, y) < \varepsilon$; and transitivity). A point
is \emph{chain-recurrent in $D$} when for every $\varepsilon > 0$ and
$T$ there is an $(\varepsilon, T)$-chain from $x$ to itself.

\textbf{The judgment.}\enspace
$\mathcal{M} \vDash \mathsf{Morse}(D, G, \rho)$ holds when some family
$\{M(v)\}_{v \in V(G)}$ --- indexed by the poset the graph $G$
presents by reachability --- satisfies:
each $M(v)$ compact and invariant; the family pairwise disjoint;
\emph{(location)} $M(v) \subseteq \rho(v)$;
\emph{(completeness of the cover)} every chain-recurrent point of
$\varphi|_D$ lies in some $M(v)$; and
\emph{(order)} every complete orbit in $D$ contained in no single
$M(v)$ has $\alpha(\gamma) \subseteq M(v)$ and
$\omega(\gamma) \subseteq M(w)$ for some $w < v$. The order runs
source-to-sink: $w < v$ means connecting orbits may run from $M(v)$
down to $M(w)$, so the $\alpha$-side vertex is the higher one. Empty
$M(v)$ are permitted --- the judgment is an ``at most this structure''
claim (recurrence confined to the $\rho(v)$, transitions confined to
the order), never an existence claim; the existence half is a
separately priced annotation tier (Conley indices, connection
matrices), deliberately absent here.

\textbf{The rule's data and the checker's recompute.}\enspace The
certificate supplies: the family object --- a finite lattice of
box-union blocks over one declared rational grid, each internal facet
carrying a strictly-inward sign row (the \rulename{QBlockInv} section's
row species; the same rows serve both rules); a cell map $h$ from the
lattice's cells to $V(G) \sqcup \{\mathsf{transit}\}$; per vertex cell
a containment row $T(b) \subseteq \rho(h(b))$; per transit cell a
polynomial $V$ with decrease rows $\mathcal{A}V + \varepsilon \leq 0$;
and well-formedness rows for $\rho$. The checker \emph{recomputes},
rather than trusts, the lattice readout: closure under union and
grid-level intersection, the join-irreducible members $b$ with unique
lower covers $b^-$, the cells $T(b) := \mathrm{cl}(b \smallsetminus
b^-)$, the induced cell order, convexity of $h$'s fibres and
monotonicity of $h$ on vertex cells, and acyclicity of $G$. Writing
$\mathrm{fib}(v) := \bigcup_{h(b) = v} T(b)$ for the fibre's
cell-union, the witness the rule assembles is
\[
M(v) \;:=\; \mathrm{Inv}\big(\mathrm{fib}(v)\big)
\]
--- the invariant part \emph{of the fibre's cell-union}, not the union
of per-cell invariant parts. The distinction is load-bearing
(\autoref{rem:qmd-mv}).

\begin{blemma}[$\mathrm{Inv}$ is located and invariant]
\label{lem:qmd-inv}
\emph{(i)} $\mathrm{Inv}(C) \subseteq C$.
\emph{(ii)} $\mathrm{Inv}(C)$ is an invariant set.
\end{blemma}
\begin{proof}
(i) The witnessing orbit passes through $x$ at time $0$ with range in
$C$. (ii) Two facts, both by time-shifting: shifting a complete orbit
by any $s$ ($u \mapsto \gamma(u + s)$) is again a complete orbit, and
every point $\gamma(s)$ of an orbit with range in $C$ therefore lies in
$\mathrm{Inv}(C)$ with the shifted orbit as witness. Forward
invariance: for $x \in \mathrm{Inv}(C)$ with orbit $\gamma$,
$\Phi_t x = \gamma(0 + t)$ is such a point. The complete-orbit-through
clause: the same statement at every shift shows
$\mathrm{range}\,\gamma \subseteq \mathrm{Inv}(C)$, so the witnessing
orbit itself runs inside $\mathrm{Inv}(C)$.
\end{proof}

\begin{blemma}[location]\label{lem:qmd-loc}
Under the containment rows, $\mathrm{fib}(v) \subseteq \rho(v)$, hence
$M(v) \subseteq \rho(v)$ by Lemma~\ref{lem:qmd-inv}(i).
\end{blemma}
\begin{proof}
A union of cells each inside $\rho(v)$ is inside $\rho(v)$.
\end{proof}

\begin{blemma}[forward confinement]\label{lem:qmd-conf}
If $x \in \mathrm{Inv}(C)$ then $\Phi_t x \in C$ for every
$t \geq 0$.
\end{blemma}
\begin{proof}
The witnessing orbit's forward half \emph{is} the forward orbit:
$\Phi_t x = \Phi_t(\gamma(0)) = \gamma(t) \in C$. This one-line lemma
is the section's quiet workhorse: a semiflow is deterministic
\emph{forward}, so membership in an $\mathrm{Inv}$ set constrains the
actual forward trajectory --- and nothing here assumes backward
uniqueness (a semiflow may branch backward; complete orbits through a
point need not be unique). The disjointness reduction below survives
precisely because it only ever uses forward halves.
\end{proof}

\begin{blemma}[transit certificates kill invariance]\label{lem:qmd-transit}
Call $C$ \emph{transit-certified} when there are $V : X \to \R$ and
$\varepsilon > 0$ with $V$ bounded below on $C$ and, along every
forward trajectory confined to $C$,
$V(\Phi_t x) \leq V(x) - \varepsilon\, t$. Then no point is
forward-confined to $C$; in particular
$\mathrm{Inv}(C) = \emptyset$.
\end{blemma}
\begin{proof}
Suppose $\Phi_t x \in C$ for all $t$, and let $m$ bound $V$ below on
$C$. By the Archimedean property choose $n$ with
$n\,\varepsilon > V(x) - m$; the decrease clause at horizon $n$ gives
$V(\Phi_n x) \leq V(x) - \varepsilon n < m$, contradicting
$\Phi_n x \in C$. The $\mathrm{Inv}$ statement follows by
Lemma~\ref{lem:qmd-conf}. (The certified rows are pointwise ---
$\mathcal{A}V + \varepsilon \leq 0$ on the cell, for the deterministic
generator $\mathcal{A} = f \cdot \nabla$ --- and reach the integrated
clause used here by integrating the derivative along the trajectory;
that bridge is ledger item G2 in \autoref{rem:qmd-ledger}, stated
once and consumed by every certificate-shaped hypothesis in this
section.)
\end{proof}

\begin{blemma}[disjointness from overlap certificates]
\label{lem:qmd-disj}
If $C_1 \cap C_2$ is transit-certified, then
$\mathrm{Inv}(C_1)$ and $\mathrm{Inv}(C_2)$ are disjoint.
\end{blemma}
\begin{proof}
A common point is forward-confined to $C_1$ and to $C_2$
(Lemma~\ref{lem:qmd-conf} twice), hence to $C_1 \cap C_2$ --- which
Lemma~\ref{lem:qmd-transit} forbids. No backward uniqueness, no
injectivity of the flow maps.
\end{proof}

\begin{bremark}[why the certificate sits on the overlap ---
``strictness pays a second time'']\label{rem:qmd-overlap}
Distinct fibres' cell-unions intersect only in shared boundary faces
(cell interiors are disjoint by the readout's construction ---
generically facets, but possibly lower-dimensional: two cells may meet
in an edge or a corner alone). On a certified facet incident to a
shared face, the facet polynomial supplies the observable ---
\emph{oriented to decrease}: the rows certify inward growth of $\ell$,
so the transit observable Lemma~\ref{lem:qmd-transit} consumes is
$V = -\ell$, bounded below on the compact overlap, with the certified
strict rate as its $\varepsilon$; the shared face lies in $\ell$'s
zero set, and confinement to the face would force $\ell \equiv 0$.
Where the overlap is a union of several faces, one signed observable
per face does not by itself assemble into a single global certificate
--- the row-shaped assembly is per-cell-pair crossing certificates on
shared facets, the development's G4 sharpening, bound in this edition
(anchor table) --- and the hypothesis
Theorem~\ref{thm:qmd-sound} takes is the assembled overlap
certificate. The strict margin of the \rulename{QBlockInv} rows is
thus consumed a second time, as the transit raw material for facet
regions. One might
hope to weaken the hypothesis to per-piece certificates --- each
$C_i$ separately transit-certified --- and conclude something about
the union; that is false. On the circle rotation, two overlapping
closed arcs each carry a monotone angular observable (a transit
certificate for the arc), yet the union is the whole circle, every
point recurrent: certificates on pieces do not compose across an
overlapping union, because no global observable exists. The overlap
form is the correct, and certifiable, hypothesis.
\end{bremark}

\begin{blemma}[compactness of the invariant parts]\label{lem:qmd-cpt}
Suppose $X$ is Hausdorff, each flow map $\Phi_t : X \to X$ is
continuous, and $C$ is compact. Then $\mathrm{Inv}(C)$ is compact.
Consequently each $\mathrm{fib}(v)$, a finite union of compact cells,
has compact invariant part.
\end{blemma}
\begin{proof}
Let $K \subseteq X^{\R}$ (product topology) be the set of complete
orbits with range in $C$:
\[
K \;=\; \Big(\bigcap_{s \in \R,\ t \geq 0}
\{\gamma : \gamma(s+t) = \Phi_t(\gamma(s))\}\Big)
\;\cap\; \prod_{r \in \R} C .
\]
Each orbit constraint is the equalizer of two continuous maps ---
evaluation at $s + t$, and $\Phi_t \circ (\text{evaluation at } s)$
--- hence closed, $X$ being Hausdorff; the product
$\prod_{\R} C$ is compact by Tychonoff (the one genuinely infinitary
step of this section); a closed subset of a compact set is compact. And
$\mathrm{Inv}(C)$ is exactly the image of $K$ under evaluation at $0$,
a continuous map, hence compact. The fibre statement: a finite union
of compact cells is compact, and the above applies to it.
\end{proof}

\begin{blemma}[the cover assembles from the cell level]
\label{lem:qmd-cover}
Suppose \emph{(cell-level cover)} every chain-recurrent point of
$\varphi|_D$ lies in $\mathrm{Inv}(T(b))$ for some cell $b$, and
\emph{(transit rows)} every transit cell is transit-certified. Then
every chain-recurrent point of $\varphi|_D$ lies in some
$M(v) = \mathrm{Inv}(\mathrm{fib}(v))$.
\end{blemma}
\begin{proof}
Let $x$ be chain-recurrent, $x \in \mathrm{Inv}(T(b))$. If $h(b) =
\mathsf{transit}$, Lemma~\ref{lem:qmd-transit} makes
$\mathrm{Inv}(T(b))$ empty --- contradiction; so $h(b)$ is a vertex
$v$. Then $T(b) \subseteq \mathrm{fib}(v)$, and $\mathrm{Inv}$ is
monotone in its region (a complete orbit inside the cell is inside the
fibre's union), so $x \in \mathrm{Inv}(\mathrm{fib}(v))$.
\end{proof}

\begin{bthm}[soundness of \rulename{QMorseDischarge}, at the
discharged floor]\label{thm:qmd-sound}
Suppose:
\emph{(model class)} $X$ is a Hausdorff metric space, each flow map
$\Phi_t$ is continuous, and every cell is compact;
\emph{(containment rows)} $T(b) \subseteq \rho(h(b))$ for every vertex
cell;
\emph{(overlap certificates)} for $v \neq w$, the overlap
$\mathrm{fib}(v) \cap \mathrm{fib}(w)$ is transit-certified;
\emph{(cell-level cover)} every chain-recurrent point of $\varphi|_D$
lies in some cell's invariant part;
\emph{(transit rows)} every transit cell is transit-certified; and
\emph{(order clause)} every complete orbit in $D$ contained in no
single $M(v)$ connects a strictly higher fibre's invariant part to a
strictly lower one, in the sense of the judgment's order clause with
$M(v) = \mathrm{Inv}(\mathrm{fib}(v))$.
Then $\mathcal{M} \vDash \mathsf{Morse}(D, G, \rho)$.
\end{bthm}
\begin{proof}
Take the family $M(v) := \mathrm{Inv}(\mathrm{fib}(v))$ and check the
witness clauses. \emph{Location}: Lemma~\ref{lem:qmd-loc}.
\emph{Invariance}: Lemma~\ref{lem:qmd-inv}(ii). \emph{Compactness}:
Lemma~\ref{lem:qmd-cpt} under the model-class facts.
\emph{Disjointness}: Lemma~\ref{lem:qmd-disj} applied to the overlap
certificates. \emph{Cover}: Lemma~\ref{lem:qmd-cover} from the
cell-level cover and the transit rows. \emph{Order}: the hypothesis,
verbatim. All six clauses hold, so the family witnesses the judgment.
\end{proof}

\begin{bremark}[the gap ledger, as mathematics]\label{rem:qmd-ledger}
The rule was first mechanized as the bare witness assembly: location
and invariance proved outright, and four clauses --- labelled G3--G6 in
the development, the labels this book adopts as the rule's hypothesis
names --- taken as opaque interface hypotheses. The discharge stratum
bound by this section moves the ledger to the state
Theorem~\ref{thm:qmd-sound} exhibits:

\emph{G3 (compactness of the invariant parts): proven.}
Lemma~\ref{lem:qmd-cpt} derives it from model-class facts alone ---
Hausdorff state space, spacewise-continuous flow, compact cells. No
compactness hypothesis survives at the fibre level.

\emph{G4 (disjointness of distinct fibres' invariant parts): reduced,
and sharpened to rows.} Lemma~\ref{lem:qmd-disj} trades it for transit
certificates on pairwise overlaps; the residue is the overlap
observable of \autoref{rem:qmd-overlap} --- shared-facet
functionals whose certified strict signs are the raw material ---
together with the G2 bridge below; and the row-shaped assembly
(per-cell-pair crossing certificates on shared facets in place of
whole-overlap certificates) is mechanized and bound
(\texttt{qMorseDischarge\_sound\_cellCover}, anchor table).

\emph{G5 (completeness of the chain-recurrence cover): reduced.}
Lemma~\ref{lem:qmd-cover} trades the fibre-level cover for the
cell-level one. The residue is the margin argument: a chain-recurrent
point's chains stay within its own cell, because the strict facet
margin is a uniform bound on admissible chain jumps near facets --- an
$\varepsilon$-chain with $\varepsilon$ below the margin can no more
exit a block than an orbit can. The classical home of this
containment is Conley's theory \cite{conley1978}; its algorithmic
form is the chain-recurrence algorithm of
Kalies--Mischaikow--van der Vorst \cite{kmv2005}.

\emph{G6 (the order clause): a hypothesis of
Theorem~\ref{thm:qmd-sound}, discharged by the bound assembly legs.}
Its mathematical content is the orbit
combinatorics of the block lattice: the \emph{eventual member} of a
forward orbit --- the minimal lattice member containing a forward
tail, well-defined because the tail-containing family is finite,
closed under the lattice meet, and nonempty: the arena sits inside the
top member's realization (one of the combinatorial interface rows the
assembly legs consume), so the top member contains every forward
tail --- descends the cell
order and stabilizes, giving every complete orbit in $D$ extreme
cells $b_\omega \leq b_\alpha$ that bracket every visited cell; the
fibre-convexity case analysis then converts bracketing into the
witness's order clause (the extremes are not transit, by
Lemma~\ref{lem:qmd-transit} against their nonempty invariant limit
sets; a same-fibre pair traps the whole orbit in its fibre by
convexity, contradicting the ``no single $M(v)$'' hypothesis; distinct
fibres yield the strict order by monotonicity of $h$). The
development's assembly legs mechanize exactly this --- the cell order
entering as theorem-interface data, the eventual-member construction
proved against combinatorial row hypotheses --- and are bound in this
edition (\texttt{qMorseDischarge\_sound\_orderAssembled},
\texttt{\_eventualMember}; anchor table), so nothing of G6 survives
under its own name: the residue is the R1-recompute combinatorial rows
and strict-inwardness shadows those legs consume, plus the G2-class
profile identifications below. Theorem~\ref{thm:qmd-sound} is stated
at the discharged floor and takes the clause as an explicit hypothesis
there; the rule's conditional badge is judged against the full bound
set.

\emph{G2 (the pointwise-to-integrated bridge), named for
completeness:} every ``transit-certified'' hypothesis above is the
integrated semantic shadow of pointwise rows
$\mathcal{A}V + \varepsilon \leq 0$; the bridge is one-variable
calculus along trajectories, and its surviving analytic content is
the identification of the certified polynomial rate with the
trajectory-wise derivative on the region --- the same model-class
profile warrant as ledger item (a) of the \rulename{QBlockInv}
section, whose G1 (block invariance) completes the shared ledger.
\end{bremark}

\begin{bremark}[the witness is $\mathrm{Inv}$ of the union, not the
union of $\mathrm{Inv}$s]\label{rem:qmd-mv}
Define instead $M'(v) := \bigcup_{h(b) = v} \mathrm{Inv}(T(b))$ and
the order clause fails: a complete orbit crossing between two cells of
the \emph{same} fibre lies in no per-cell invariant part (its range is
in neither cell alone), so it belongs to no $M'(v)$ while also
connecting nothing strictly ordered --- the ``orbit contained in no
Morse set'' case triggers with no higher/lower pair to name. The
correct definition absorbs such orbits into their fibre's $M(v)$,
which is exactly the role of fibre convexity: every cell between the
orbit's extremes maps to the same vertex, so the whole orbit lies in
the fibre's union and hence in its invariant part. Without convexity
the judgment could assert a strict order the certificate does not
enforce --- convexity of $h$'s fibres is well-formedness, not a
nicety, and the checker recomputes it.
\end{bremark}

\begin{bremark}[relation to the literature]
The judgment is Conley's Morse decomposition \cite{conley1978} in
certified form: a decomposition indexed by a two-element order is an
attractor--repeller pair, and the general form is the lattice-refined
iteration of that idea; the completeness clause is the containment of
chain recurrence in the Morse sets, and the certified staircase of
lattice members is a finite shard of a complete Lyapunov function
(Conley's fundamental theorem). The block-lattice route --- attracting
blocks closed under union and intersection, join-irreducibles, the
tessellation into cells $T(b)$, and Birkhoff duality
\cite{birkhoff1937} between the attractor lattice and the Morse-graph
poset --- is the program of Kalies--Mischaikow--van der Vorst's lattice
series \cite{kmv2016,kmv2022}; the chain-recurrence cover's
algorithmic form is their earlier chain-recurrence algorithm
\cite{kmv2005}; and the computational tradition of certifying global
dynamics by index machinery is surveyed in Mischaikow--Mrozek
\cite{mischaikowmrozek2002}. What the rule adds to that tradition is
the artefact discipline: those pipelines end in trusted computation,
while here every analytic claim rides certified sign rows --- base
kernel sequents in exact $\Q$ --- plus combinatorics the checker
recomputes, with the semantic warrant being the lemma chain above.
\end{bremark}

\begin{anchors}
witness assembly, location + invariance proved & \texttt{qMorseDischarge\_sound} (QtslSoundness.lean); \texttt{maxInv\_subset}, \texttt{maxInv\_isInvariant}\\
Lemma \ref{lem:qmd-inv} & \texttt{maxInv\_subset}, \texttt{maxInv\_invariant}, \texttt{maxInv\_isInvariant}\\
Lemma \ref{lem:qmd-loc} & \texttt{fiberUnion\_subset}; rows \texttt{hContain}\\
Lemma \ref{lem:qmd-conf} & \texttt{maxInv\_forward\_confined} (QtslDischarge.lean)\\
Lemma \ref{lem:qmd-transit} & \texttt{HasTransitCert}, \texttt{HasTransitCert.not\_forward\_confined}, \texttt{HasTransitCert.maxInv\_empty}; the bare form \texttt{transit\_maxInv\_empty} (QtslSoundness.lean)\\
Lemma \ref{lem:qmd-disj} & \texttt{HasTransitCert.disjoint\_maxInv}; hypothesis \texttt{hOverlap}\\
Lemma \ref{lem:qmd-cpt} & \texttt{maxInv\_isCompact} (Tychonoff route), \texttt{fiberUnion\_isCompact}\\
Lemma \ref{lem:qmd-cover} & \texttt{crCover\_assemble}; \texttt{maxInv\_mono}, \texttt{subset\_fiberUnion}; hypotheses \texttt{hCell}, \texttt{hTransitCert}\\
Theorem \ref{thm:qmd-sound} & \texttt{qMorseDischarge\_sound\_modelClass}; order hypothesis \texttt{hOrder} (ledger G6)\\
judgment semantics & \texttt{MorseWitness}, \texttt{MorseSat}, \texttt{EpsChainIn}, \texttt{ChainRecurrentIn}, \texttt{omegaSet}, \texttt{alphaSet}\\
G6's assembly legs (\textbf{bound in this edition}) & \texttt{orderClause\_assemble}, \texttt{eventualMember\_interval}, \texttt{qMorseDischarge\_sound\_orderAssembled}, \texttt{\_eventualMember} (QtslOrderAssembly, QtslEventualMember)\\
G4's row sharpening (\textbf{bound in this edition}) & \texttt{qMorseDischarge\_sound\_cellCover} (QtslCrossing)\\
\end{anchors}


\rulesection{71}{QMorseProduct}{core rule}
\label{blcatrule-QMorseProduct}
\ruleblock{The rule (Steele metanotation)}
\sourcenote{as realized in the mechanized development; source row: \texttt{\texttt{QTSL-rules.toml}} --- \emph{QMorseProduct}. The metanotation statement is typeset by hand for notational quality; its mechanized ground truth is the Lean statement below, extracted from the sources at every build of this book and hash-bound. The rule \emph{roster} is machine-checked (every mechanized rule tag round-trips to exactly one catalogue row); the statement \emph{text} is editorial, checked by eye against the extracted Lean beside it.}
\begin{mathpar}
\inferrule*{
  \mathcal{M}_1 \vDash \mathsf{Morse}(D_1, G_1, \rho_1) \\
  \mathcal{M}_2 \vDash \mathsf{Morse}(D_2, G_2, \rho_2) \\
  L_1, L_2 : \mathsf{QBlockLattice} \ \text{(the factor lattices)}
}{
  \mathcal{M}_1 \times \mathcal{M}_2 \vDash \mathsf{Morse}(D_1 \times D_2,\ G_1 \times G_2,\ \rho_1 \times \rho_2)
}
\end{mathpar}

\ruleflavor{the two factor judgments enter as derivation premises (judgment-kind edges); the factor lattices enter as certificate data (shared references --- the premise judgments\textquotesingle{} own certificates impose their claims), consulted by the checker\textquotesingle s product tie.}
\par\medskip
\rulename{QMorseProduct} composes two global judgments along the
state-space axis: disjoint variable sets, product arena, product spec
graph, product regions. A producer reaches for it exactly where global
structure is engineered compositionally --- certify each factor
subsystem\textquotesingle s Morse graph, then take the product ---
because the alternative (re-discharging the composite through
\rulename{QMorseDischarge} on the composite lattice) recomputes an
exponentially larger readout and re-derives the cell map.

The rule is fenced to the deterministic-flow class, and the fence is
the rule\textquotesingle s own mathematical content: the composite
certificate\textquotesingle s readout --- join-irreducibles, cells,
cell order --- must be the product of the factor readouts, and for
flows that is a theorem (the attractor lattice of a product flow is
the coproduct of the factor attractor lattices), while for maps it is
false: products of maps exhibit emergent, phase-locked attractors
outside the coproduct, so a composite readout computed from factor
certificates would misrepresent the product\textquotesingle s global
structure. For maps the rule must not be derived at all.

What the checker recomputes: wire-disjointness of the factors, the
product surgery on graphs and regions, and --- on the certificate
route --- the composite lattice\textquotesingle s readout. Facet and
transit rows ride along by citation: every composite facet is normal
to one factor\textquotesingle s axis and its inwardness polynomial is
that factor\textquotesingle s, the other factor\textquotesingle s
coordinates entering as free box variables; a factor transit
observable still decreases on cell $\times$ anything, the other
factor\textquotesingle s drift differentiating it to zero.

The witness-level lemma chain is mechanized in full: the product
construction on the semiflow stratum (\texttt{Semiflow.prod}), the
factoring lemmas for orbits, limit sets, and chains, the six-clause
witness transport (\texttt{MorseWitness.prod}), and the soundness
theorem (\texttt{qMorseProduct\_sound}) --- unconditional, with no
interface hypotheses: from the two factor judgments, the product
judgment, on the product arena, product order, and product regions.
The bound statements consume neither compactness nor invariance of
the arenas (the proof never needed them), and Hausdorffness exactly
where the transport says ``compact, hence closed.''

\cataloguenote{The proof section separates two strata deliberately: the
witness-transport argument (compactness, invariance, location,
disjointness, cover, order --- all coordinate-wise) is elementary and
does not consume continuity of time, while the readout-transport
theorem that warrants the certificate route is flow-specific. The
fence binds the rule as catalogued; the decomposition is recorded so
the mechanization and any map-side successor judgment know where the
flow hypothesis actually bites. The mechanization respects the split
exactly: the bound statements carry the witness-level entailment in
full and import nothing of Kalies--Wehbe --- the fence's warrant
remains at literature strength, where it belongs, since it certifies
the canonicity of the certificate-route composite readout rather than
the satisfaction-level soundness bound here.
}
\ruleblock{Lean statement}
\declhead{BernsteinLogic.QtslProduct}{qMorseProduct\_sound}{\badge{BadgeGreen}{statement bound}}
\begin{leansrc}
/-- QMorseProduct soundness, witness level (thm:qmp-sound; the bind
target). From the two factor judgments — the rule's E1 premises `left,
right : QMorse` — conclude the product judgment on the product arena,
product reachability order, and product regions. UNCONDITIONAL: no
interface hypotheses (contrast the `\dcond` family); the analytic
content is the witness transport above. The E2 lattice premises and the
`flows_only` fence's readout-transport warrant (Kalies–Wehbe) live at
the certificate stratum and are deliberately not consumed — see the
module header and the book page's fence remark. -/
theorem qMorseProduct_sound
    (h₁ : MorseSat φ₁ D₁ ρ₁) (h₂ : MorseSat φ₂ D₂ ρ₂) :
    MorseSat (φ₁.prod φ₂) (D₁ ×ˢ D₂)
      (fun p : Vtx₁ × Vtx₂ => ρ₁ p.1 ×ˢ ρ₂ p.2)
\end{leansrc}
\declhead{BernsteinLogic.QtslProduct}{MorseWitness.prod}{\badge{BadgeGreen}{statement bound}}
\begin{leansrc}
/-- Witness transport (lem:qmp-transport): the product family
`M(v, w) := M₁(v) × M₂(w)`, indexed by the product order, witnesses the
product judgment. Six clauses, each the book paragraph:

* compactness — a product of compacts is compact;
* invariance — coordinate-wise forward invariance, and a product point
  rides the pairing of the factor witnesses' complete orbits;
* location — the product of the factor locations;
* disjointness — distinct product vertices differ in some coordinate,
  and that factor's disjointness kills the intersection;
* cover — a chain-recurrent product point has chain-recurrent
  coordinates (§4), each covered by its factor family;
* order — a product orbit contained in no `M(v, w)` has an uncontained
  coordinate; three cases (both uncontained / first only / second
  only), the contained coordinate handled by compact-hence-closed and
  the closure bound (§3), the α-side vertex strictly higher in the
  product order exactly as the source-to-sink convention demands. -/
def MorseWitness.prod (W₁ : MorseWitness φ₁ D₁ ρ₁)
    (W₂ : MorseWitness φ₂ D₂ ρ₂) :
    MorseWitness (φ₁.prod φ₂) (D₁ ×ˢ D₂)
      (fun p : Vtx₁ × Vtx₂ => ρ₁ p.1 ×ˢ ρ₂ p.2)
\end{leansrc}
\declhead{BernsteinLogic.QtslProduct}{completeOrbit\_prod\_iff}{\badge{BadgeGreen}{statement bound}}
\begin{leansrc}
/-- Orbits factor (lem:qmp-factor(i)): `γ` is a complete orbit of the
product iff each coordinate is a complete orbit of its factor. -/
theorem completeOrbit_prod_iff {γ : ℝ → X₁ × X₂} :
    (φ₁.prod φ₂).CompleteOrbit γ ↔
      φ₁.CompleteOrbit (fun s => (γ s).1) ∧
      φ₂.CompleteOrbit (fun s => (γ s).2)
\end{leansrc}
\declhead{BernsteinLogic.QtslProduct}{omegaSet\_prod\_subset}{\badge{BadgeGreen}{statement bound}}
\begin{leansrc}
/-- ω-limit sets factor (lem:qmp-factor(ii)):
`ω(γ) ⊆ ω(γ₁) × ω(γ₂)`. Book proof verbatim: the image of `[T, ∞)`
under `γ` sits inside the product of the coordinate images; closure of
a product is the product of closures. -/
theorem omegaSet_prod_subset (γ : ℝ → X₁ × X₂) :
    omegaSet γ ⊆
      (omegaSet fun s => (γ s).1) ×ˢ (omegaSet fun s => (γ s).2)
\end{leansrc}
\declhead{BernsteinLogic.QtslProduct}{alphaSet\_prod\_subset}{\badge{BadgeGreen}{statement bound}}
\begin{leansrc}
/-- α-limit sets factor (lem:qmp-factor(ii), the `likewise for α`). -/
theorem alphaSet_prod_subset (γ : ℝ → X₁ × X₂) :
    alphaSet γ ⊆
      (alphaSet fun s => (γ s).1) ×ˢ (alphaSet fun s => (γ s).2)
\end{leansrc}
\declhead{BernsteinLogic.QtslProduct}{EpsChainIn.fst}{\badge{BadgeGreen}{statement bound}}
\begin{leansrc}
/-- Chains project, first coordinate (lem:qmp-factor(iii)): under the
max metric (mathlib's product metric, `Prod.dist_eq`) factor distance
is ≤ product distance, so a chain step projects to a chain step — jump
points project into `D₁`, durations unchanged — and concatenations
project to concatenations. -/
theorem EpsChainIn.fst {D₁ : Set X₁} {D₂ : Set X₂} {ε : ℝ} {T : ℝ≥0}
    {x y : X₁ × X₂}
    (h : EpsChainIn (φ₁.prod φ₂) (D₁ ×ˢ D₂) ε T x y) :
    EpsChainIn φ₁ D₁ ε T x.1 y.1
\end{leansrc}
\declhead{BernsteinLogic.QtslProduct}{EpsChainIn.snd}{\badge{BadgeGreen}{statement bound}}
\begin{leansrc}
/-- Chains project, second coordinate. -/
theorem EpsChainIn.snd {D₁ : Set X₁} {D₂ : Set X₂} {ε : ℝ} {T : ℝ≥0}
    {x y : X₁ × X₂}
    (h : EpsChainIn (φ₁.prod φ₂) (D₁ ×ˢ D₂) ε T x y) :
    EpsChainIn φ₂ D₂ ε T x.2 y.2
\end{leansrc}
\declhead{BernsteinLogic.QtslProduct}{ChainRecurrentIn.fst}{\badge{BadgeGreen}{statement bound}}
\begin{leansrc}
/-- Chain recurrence projects, first coordinate (lem:qmp-factor(iii),
"hence a point chain-recurrent in `D₁ × D₂` has both coordinates
chain-recurrent in their arenas"). -/
theorem ChainRecurrentIn.fst {D₁ : Set X₁} {D₂ : Set X₂} {x : X₁ × X₂}
    (h : ChainRecurrentIn (φ₁.prod φ₂) (D₁ ×ˢ D₂) x) :
    ChainRecurrentIn φ₁ D₁ x.1
\end{leansrc}
\declhead{BernsteinLogic.QtslProduct}{ChainRecurrentIn.snd}{\badge{BadgeGreen}{statement bound}}
\begin{leansrc}
/-- Chain recurrence projects, second coordinate. -/
theorem ChainRecurrentIn.snd {D₁ : Set X₁} {D₂ : Set X₂} {x : X₁ × X₂}
    (h : ChainRecurrentIn (φ₁.prod φ₂) (D₁ ×ˢ D₂) x) :
    ChainRecurrentIn φ₂ D₂ x.2
\end{leansrc}
\declhead{BernsteinLogic.QtslProduct}{CompleteOrbit.prodMk}{\badge{BadgeGreen}{statement bound}}
\begin{leansrc}
/-- Pairing complete orbits (lem:qmp-factor(i), ⇐): factor orbits ride
together as a product orbit. -/
theorem CompleteOrbit.prodMk {γ₁ : ℝ → X₁} {γ₂ : ℝ → X₂}
    (h₁ : φ₁.CompleteOrbit γ₁) (h₂ : φ₂.CompleteOrbit γ₂) :
    (φ₁.prod φ₂).CompleteOrbit fun s => (γ₁ s, γ₂ s)
\end{leansrc}
\declhead{BernsteinLogic.QtslProduct}{omegaSet\_subset\_closure\_range}{\badge{BadgeGreen}{statement bound}}
\begin{leansrc}
/-- `ω(γ) ⊆ cl(range γ)` — one intersect and one image bound; feeds the
order clause's contained cases (book:
`α(γ₂), ω(γ₂) ⊆ cl(range γ₂) ⊆ M₂(w₀)`). -/
theorem omegaSet_subset_closure_range (γ : ℝ → X) :
    omegaSet γ ⊆ closure (Set.range γ)
\end{leansrc}
\declhead{BernsteinLogic.QtslProduct}{alphaSet\_subset\_closure\_range}{\badge{BadgeGreen}{statement bound}}
\begin{leansrc}
/-- `α(γ) ⊆ cl(range γ)` — the α twin. -/
theorem alphaSet_subset_closure_range (γ : ℝ → X) :
    alphaSet γ ⊆ closure (Set.range γ)
\end{leansrc}
\ruleblock{Correspondence (Steele display $\leftrightarrow$ Lean)}
\absentbox{NO CORRESPONDENCE DICTIONARY YET}{no reviewer-certified dictionary maps the displayed rule onto the mechanized statements above; until one is written and certified, that mapping rests on this section\textquotesingle s editorial care alone.}
\ruleblock{Lean verifier receipts}
\sourcenote{module \texttt{BernsteinLogic.QtslProduct}; bound build \texttt{1b85f883dfdc\ldots} (bound from receipt d5b482c9d6ed9b6c, 2026-07-20T03:30:32Z). A receipt records that named checking layers accepted a build: (module, compiled-artifact SHA-256, environment SHA-256, checker layers), minted when the checking gate ran; the store receipts attesting the bound build are listed with the layers each run credited. Nothing here asks the reader to trust the store: every layer is an independently runnable program over the quoted sources --- rebuild the module and re-run the checkers to reproduce the verdict. Receipts attesting superseded builds credit no layer and are omitted from this rendering (the count of omissions, if any, is noted below).}
\begin{receipttable}
\receiptrow{d5b482c9d6ed9b6c}{2026-07-20T03:30:32Z}{build+decorrelated}{\badge{BadgeGreen}{current}}
\receiptrow{15fa0566d44e84fd}{2026-07-20T20:44:52Z}{build+decorrelated}{\badge{BadgeGreen}{current}}
\receiptrow{e8214fce26d91980}{2026-07-20T21:01:21Z}{build+replay+hardened+decorrelated}{\badge{BadgeGreen}{current}}
\receiptrow{b87a989a1475042c}{2026-07-21T05:28:48Z}{build+replay+hardened+decorrelated}{\badge{BadgeGreen}{current}}
\end{receipttable}
\layerline{\layeritem{build}{lake build (elaboration)}{\badge{BadgeGreen}{current}}\layeritem{replay}{leanchecker (official kernel replay)}{\badge{BadgeGreen}{current}}\layeritem{hardened}{safe\_verify (hardened screen)}{\badge{BadgeGreen}{current}}\layeritem{decorrelated}{lean4export + nanoda (decorrelated kernel)}{\badge{BadgeGreen}{current}}}
\sourcenote{\badge{BadgeGray}{comparator: none} no fixture in the three-implementation comparator harness (compiled / kernel-interpreted / reference-interpreted) exercises this rule.}
\ruleblock{HOL4 statement}
\absentbox{NO HOL STATEMENT YET}{Absent at the judgment stratum: the HOL4 development\textquotesingle s
object is the base-kernel checker cone, and no Morse-family definition
exists in the ported scripts (verified against the
port\textquotesingle s definition inventory) --- so the
satisfaction-judgment semantics this rule\textquotesingle s soundness
is stated in has no HOL twin to mirror the now-complete Lean
mechanization into. A future port has its interface fixed twice over:
the witness clauses quoted in the \rulename{QMorseDischarge} section
and the product lemma chain bound here.
}
\ruleblock{HOL verifier receipts}
\absentbox{NONE}{no HOL verification receipt exists for this rule.}
\ruleblock{Proof}

\sourcenote{Genre: core rule (tier 1) --- an ordinary mathematical proof
at Putnam-exam-solution detail, mechanized \emph{in full}: the
accompanying development proves the product construction, the factoring
lemmas, the witness transport, and the soundness theorem with no
interface hypotheses (contrast the \rulename{QMorseDischarge} ledger;
within the Morse family only \rulename{QBasin} shares this fully-proved
status). The lemma chain below fixed the mechanization's interface in
advance, and the formal-anchor table maps each step to the declaration
that carries it. The readout-transport theorem of the fence remark is
deliberately \emph{not} mechanized: it warrants the canonicity of the
certificate-route composite readout, not the witness-level entailment
proved here, and remains at literature strength.}

\textbf{Setting and claim.}\enspace Let $\varphi_1, \varphi_2$ be
forward semiflows on Hausdorff metric state spaces $X_1, X_2$, with
compact forward-invariant arenas $D_1, D_2$. The \emph{product
semiflow} $\varphi_1 \times \varphi_2$ on $X_1 \times X_2$ (with the
max metric) flows coordinate-wise:
$\Phi_t(x_1, x_2) = (\Phi^1_t x_1, \Phi^2_t x_2)$; the semiflow laws
hold coordinate-wise. The rule: from
$\mathcal{M}_1 \vDash \mathsf{Morse}(D_1, G_1, \rho_1)$ and
$\mathcal{M}_2 \vDash \mathsf{Morse}(D_2, G_2, \rho_2)$, with
disjoint variable sets, conclude
\[
\mathcal{M}_1 \times \mathcal{M}_2 \ \vDash\
\mathsf{Morse}\big(D_1 \times D_2,\ G_1 \times G_2,\ \rho_1 \times \rho_2\big),
\]
where $G_1 \times G_2$ presents the product order
($(v_1, v_2) \geq (w_1, w_2)$ iff $v_i \geq w_i$ in both coordinates)
and $(\rho_1 \times \rho_2)(v_1, v_2) := \rho_1(v_1) \times
\rho_2(v_2)$.

\begin{blemma}[orbits, limit sets, and chains factor]
\label{lem:qmp-factor}
\emph{(i)} $\gamma = (\gamma_1, \gamma_2)$ is a complete orbit of
$\varphi_1 \times \varphi_2$ iff each $\gamma_i$ is a complete orbit
of $\varphi_i$.
\emph{(ii)} $\omega(\gamma) \subseteq \omega(\gamma_1) \times
\omega(\gamma_2)$, and likewise for $\alpha$.
\emph{(iii)} the projection of an $(\varepsilon, T)$-chain inside
$D_1 \times D_2$ is an $(\varepsilon, T)$-chain inside $D_i$; hence a
point chain-recurrent in $D_1 \times D_2$ has both coordinates
chain-recurrent in their arenas.
\end{blemma}
\begin{proof}
(i) is the coordinate-wise flow law read in both directions. (ii):
$\gamma([T, \infty)) \subseteq \gamma_1([T,\infty)) \times
\gamma_2([T,\infty))$; closure of a product is the product of
closures; and an intersection of product sets is the product of the
intersections --- pointwise, $(x_1, x_2)$ lies in every
$A_T \times B_T$ iff each $x_i$ lies in every factor set. (iii): under
the max metric, factor distance is $\leq$ product distance, so a
single chain step projects to a chain step (jump points project into
$D_i$, durations are unchanged), and concatenations project to
concatenations.
\end{proof}

\begin{blemma}[witness transport]\label{lem:qmp-transport}
Let $\{M_1(v)\}_{v \in V(G_1)}$ and $\{M_2(w)\}_{w \in V(G_2)}$
witness the factor judgments. Then the product family
\[
M(v, w) \;:=\; M_1(v) \times M_2(w),
\]
indexed by the product order, witnesses the product judgment.
\end{blemma}
\begin{proof}
Six clauses.

\emph{Compactness}: a product of two compact sets is compact.

\emph{Invariance}: forward invariance holds coordinate-wise; and a
point $(x_1, x_2)$ of the product rides the pairing of the factor
witnesses' complete orbits through $x_1$ and $x_2$
(Lemma~\ref{lem:qmp-factor}(i)), which stays inside the product set.

\emph{Location}: $M_1(v) \times M_2(w) \subseteq \rho_1(v) \times
\rho_2(w)$, the product of the factor locations.

\emph{Disjointness}: if $(v, w) \neq (v', w')$ they differ in some
coordinate, say $v \neq v'$; then $M_1(v) \cap M_1(v') = \emptyset$
makes the products disjoint regardless of the other coordinate.

\emph{Cover}: let $(x_1, x_2)$ be chain-recurrent in $D_1 \times D_2$.
By Lemma~\ref{lem:qmp-factor}(iii) each $x_i$ is chain-recurrent in
$D_i$, hence covered by its factor family: $x_1 \in M_1(v)$,
$x_2 \in M_2(w)$ for some $v, w$; so
$(x_1, x_2) \in M(v, w)$.

\emph{Order}: let $\gamma = (\gamma_1, \gamma_2)$ be a complete orbit
in $D_1 \times D_2$ contained in no single $M(v, w)$. Containment in a
product set is coordinate-wise, so
--- negating ``some $(v,w)$ contains $\gamma$'' ---
either $\gamma_1$ is contained in no $M_1(v)$, or $\gamma_2$ in no
$M_2(w)$. Three cases.
\emph{(both uncontained)}: the factor order clauses give
$v > v'$ with $\alpha(\gamma_1) \subseteq M_1(v)$,
$\omega(\gamma_1) \subseteq M_1(v')$, and $w > w'$ likewise for
$\gamma_2$; by Lemma~\ref{lem:qmp-factor}(ii),
$\alpha(\gamma) \subseteq M(v, w)$ and
$\omega(\gamma) \subseteq M(v', w')$, with
$(v, w) > (v', w')$ strictly in the product order.
\emph{($\gamma_1$ uncontained, $\gamma_2 \subseteq M_2(w_0)$)}: the
first factor gives $v > v'$ as above; for the second,
$\mathrm{range}\,\gamma_2 \subseteq M_2(w_0)$ with $M_2(w_0)$ compact
in a Hausdorff space, hence closed, gives
$\alpha(\gamma_2), \omega(\gamma_2) \subseteq
\mathrm{cl}(\mathrm{range}\,\gamma_2) \subseteq M_2(w_0)$; so
$\alpha(\gamma) \subseteq M(v, w_0)$,
$\omega(\gamma) \subseteq M(v', w_0)$, and
$(v, w_0) > (v', w_0)$ strictly.
\emph{(the symmetric case)}: exchange the roles of the factors.
In each case the $\alpha$-side vertex is the strictly higher one, as
the source-to-sink convention demands.
\end{proof}

\begin{bthm}[soundness of \rulename{QMorseProduct}, witness level]
\label{thm:qmp-sound}
If both factor judgments hold, the product judgment holds.
\end{bthm}
\begin{proof}
Choose factor witnesses and apply Lemma~\ref{lem:qmp-transport}.
\end{proof}

\begin{bremark}[where the flow hypothesis actually binds --- the rule's
fence]\label{rem:qmp-fence}
Nothing in Lemmas~\ref{lem:qmp-factor}--\ref{lem:qmp-transport}
consumes continuity of time: read over $\mathbb{N}$-time systems, the
same six clauses transport verbatim, at the same ``at most'' polarity.
The rule is nonetheless fenced to the deterministic-\emph{flow} class,
for reasons that live above the witness stratum. The certificate-level
discharge route does not consume semantic witnesses: it composes the
factor \emph{lattices} (the join-closure of the products of factor
blocks, every composite facet's inwardness polynomial being a factor's
with the other factor's coordinates riding along as free box
variables) and must know that the composite lattice's recomputed
readout --- its join-irreducibles, cells, and cell order --- is the
product of the factor readouts, so that the composite tessellation
represents the product dynamics at the resolution the factors
certified. For flows this is a theorem: the attractor lattice of a
product flow is the coproduct of the factor attractor lattices
(Kalies--Wehbe \cite{kalieswehbe2026}), whose Birkhoff dual
\cite{birkhoff1937} is exactly ``composite Morse graph $=$ product of
factor graphs.'' For maps it is false: products of maps exhibit
emergent, phase-locked attractors outside the coproduct
(Kalies--Wehbe's counterexample --- already visible in the product of
two period-two orbits, whose diagonal and anti-diagonal are attractors
of the product generated by no product of factor attractors), so a
composite readout computed from factor certificates would silently
misrepresent the product's global structure, and the judgment's
ecosystem --- which reads a certified Morse graph as the canonical
decomposition at the declared resolution --- would be poisoned at
exactly the canonicity it relies on. Hence the fence: for maps the
rule must not be derived at all; for flows the readout-transport
theorem is the rule's semantic import, and it is the reason this core
rule carries its own lemma rather than reducing to lattice
combinatorics.
\end{bremark}

\begin{bremark}[what rides along for free]\label{rem:qmp-rides}
Two certificate-level transports cost nothing beyond citation, and are
worth recording with the rule even though the witness-level proof
above does not touch them. \emph{Facet rows}: every internal facet of
a composite block is normal to one factor's axis, and its inwardness
polynomial \emph{is} that factor's facet polynomial --- the other
factor's coordinates enter as free box variables, so the composite
row's enclosure data is the factor's, replicated along the free
coordinates. \emph{Transit rows}: if $\mathcal{A}_1 V \leq
-\varepsilon$ on a factor cell, the same $V$ read on the product
satisfies $\mathcal{A}_{1 \times 2} V = \mathcal{A}_1 V \leq
-\varepsilon$ on cell $\times$ anything --- the other factor's drift
differentiates $V$ to zero. Both ride the same disjointness of the variable sets --- the facet
rows by weakening alone, the transit rows by weakening plus the
product-generator identity
$\mathcal{A}_{1\times 2}(V \circ \pi_i) = (\mathcal{A}_i V) \circ \pi_i$,
which is the precise content of ``differentiates $V$ to zero.''
\end{bremark}

\begin{bremark}[the mechanization consumes no arena hypotheses]
\label{rem:qmp-arena}
The setting above posits compact forward-invariant arenas because that
is what arenas are in this catalogue's ecosystem; the transport itself
never uses either property, and the mechanized statements quantify over
bare sets $D_1, D_2$ --- strictly more general, and faithful to the
witness clauses, which constrain the arena only through the judgments.
Hausdorffness is consumed exactly once, where the proof of
Lemma~\ref{lem:qmp-transport} says ``compact in a Hausdorff space,
hence closed''; the mechanization takes it as a topological hypothesis
over a pseudometric base, which every metric state space satisfies.
\end{bremark}

\begin{bremark}[relation to the literature]
The rule is the certified face of compositionality results for global
dynamics: Kalies--Wehbe \cite{kalieswehbe2026} prove the
attractor-lattice coproduct isomorphism for product flows and delimit
it (maps fail; skew products degrade to one-way results), inside the
lattice-theoretic program of Kalies--Mischaikow--van der Vorst
\cite{kmv2016,kmv2022}. The witness-level transport
(Lemma~\ref{lem:qmp-transport}) is elementary Conley theory
\cite{conley1978} over products; the observation that its ``at most''
clauses survive even where the canonical readout does not
(\autoref{rem:qmp-fence}) is exactly why the judgment's fence is
stated at the rule, not at the witness.
\end{bremark}

\begin{anchors}
product semiflow (Setting) & \texttt{Semiflow.prod}, \texttt{prod\_flow} (QtslProduct.lean); the max metric is the product metric, \texttt{Prod.dist\_eq}\\
Lemma \ref{lem:qmp-factor}(i) & \texttt{completeOrbit\_prod\_iff}, \texttt{CompleteOrbit.prodMk}\\
Lemma \ref{lem:qmp-factor}(ii) & \texttt{omegaSet\_prod\_subset}, \texttt{alphaSet\_prod\_subset}\\
Lemma \ref{lem:qmp-factor}(iii) & \texttt{EpsChainIn.fst}/\texttt{.snd}, \texttt{ChainRecurrentIn.fst}/\texttt{.snd}\\
Lemma \ref{lem:qmp-transport} & \texttt{MorseWitness.prod} (all six clauses; the contained cases of the order clause via \texttt{alphaSet\_subset\_closure\_range}, \texttt{omegaSet\_subset\_closure\_range})\\
Theorem \ref{thm:qmp-sound} & \texttt{qMorseProduct\_sound} --- unconditional, no interface hypotheses\\
witness clauses (the interface targeted) & \texttt{MorseWitness}, \texttt{MorseSat} (QtslSoundness.lean, bound in the \rulename{QMorseDischarge} section)\\
Remark \ref{rem:qmp-fence} (readout transport, flows) & \emph{deliberately not mechanized}: Kalies--Wehbe \cite{kalieswehbe2026}, attractor-lattice coproduct isomorphism; Birkhoff duality \cite{birkhoff1937}\\
interim discharge path & \rulename{QMorseDischarge} on the composite lattice (that section's theorem chain) --- no longer needed as warrant; retained as the certificate-route alternative\\
\end{anchors}


\rulesection{72}{QBlockInv}{core rule}
\label{blcatrule-QBlockInv}
\ruleblock{The rule (Steele metanotation)}
\sourcenote{as realized in the mechanized development; source row: \texttt{\texttt{QTSL-rules.toml}} --- \emph{QBlockInv}. The metanotation statement is typeset by hand for notational quality; its mechanized ground truth is the Lean statement below, extracted from the sources at every build of this book and hash-bound. The rule \emph{roster} is machine-checked (every mechanized rule tag round-trips to exactly one catalogue row); the statement \emph{text} is editorial, checked by eye against the extracted Lean beside it.}
\begin{mathpar}
\inferrule*{
  N : \mathsf{QRegionAtom} \ \text{(the block)} \\
  \bigl\{F_f \vdash \mathcal{A}\ell_f \ge \varepsilon_f\bigr\}_{f \in \mathrm{facets}(N)}, \quad \varepsilon_f \in \Q_{>0} \ \text{(} \ell_f \ \text{the facet\textquotesingle s inward affine coordinate; } F_f \ \text{its facet region)} \\
  \mathsf{FacetBridge}(\mathcal{M}, D, N, (F_f, \varepsilon_f)_f) \ \text{(the bridge package, defined below)} \\
  \vdash X_0 \subseteq N \ \text{(\texttt{init\_containment}, a single BL leaf)}
}{
  \mathcal{M} \vDash \mathsf{QInvScope}(N, X_0)
}
\end{mathpar}

\ruleflavor{each facet-inward row bounds the flow derivative of the facet\textquotesingle s inward coordinate ($\ell_f$: affine, vanishing on the facet hyperplane, nonnegative on the block\textquotesingle s side) below by a strictly positive rational margin, on the facet region $F_f$. The displayed $\mathsf{FacetBridge}$ package comprises: the block realized closed in its arena; each facet region a neighbourhood of its entire facet box, corners included (tangential margin, not merely normal thickening); the per-facet wedge test (one-directional by design); arena invariance ($N \subseteq D$, $D$ invariant); trajectory continuity and the pointwise-to-integrated row bridge --- checker-recomputed or interface-carried per profile (model-class).}
\par\medskip
\rulename{QBlockInv} concludes the invariance-scope judgment: the block
$N$ is forward-invariant and encloses the declared initial region. It
is the Morse family\textquotesingle s workhorse producer --- the
judgment it mints is consumed by \rulename{QBasin}\textquotesingle s
safety factory, by scope-restriction of reach obligations, and as the
arena warrant of nested certificates --- and its premise kit is
minimal: per internal facet of the box-union block, one strictly-inward
sign row with explicit positive rational margin (facets on the arena
boundary owe nothing), plus one containment row $X_0 \subseteq N$.

What the checker recomputes: the block\textquotesingle s grid geometry
--- that the bitmap realizes a closed box-union inside its arena, and
the facet-list well-formedness --- plus the replay of every row through
the base kernel. What would break without each hypothesis: without the
sign rows, exit through a facet; without strictness-by-margin, a strict
field sign could not be certified through the non-strict enclosure
predicates of exact-$\Q$ arithmetic (and the Morse rules that cite the
same rows would lose the arguments where strictness is genuinely
load-bearing); without the containment row, a start outside the block;
without closedness of the realization, the first-exit
argument\textquotesingle s sup-time step fails.

The mechanization is conditional at a named interface: the bound
theorems descend the invariance obligation through a first-exit
skeleton (closed block, continuous trajectories, boundary dwell) and a
facet-row bridge (integrated growth of the facet observables through a
one-directional wedge test), leaving exactly three hypothesis families
--- the integrated row clause (reachable from the pointwise rows
through a mean-value bridge whose surviving content is the
profile-level identification of the certified rate with the
trajectory-wise derivative), the grid facts of the box-union, and
arena invariance (nested citation or model-class warrant, terminating
at the whole space). The proof section states the ledger, and why the
wedge test must be one-directional (reentrant corners of box-unions).

\cataloguenote{The same facet rows serve three consumers --- this rule, the lattice
membership of the family object, and the Morse
discharge\textquotesingle s facet arguments --- which is the design
reason strict margins are mandatory here even though this
rule\textquotesingle s own conclusion would survive non-strict signs
(the proof section\textquotesingle s strictness remark).
}
\ruleblock{Lean statement}
\declhead{BernsteinLogic.QtslSoundness}{qBlockInv\_sound}{\badge{BadgeGreen}{statement bound}}
\begin{leansrc}
/-- QBlockInv soundness, the M(i) corollary (roster: "2 lines"). The
facet rows deliver `hInv` through Lemma M(i) — strict facet-inwardness
of the compact box-union block implies forward invariance (gap G1, tier
A prose); `init_containment` delivers `hInit`. The rule-level content is
the assembly of the `QInvScope` claim. -/
theorem qBlockInv_sound {φ : Semiflow X} {N X0 : Set X}
    (hInv : φ.Invariant N) (hInit : X0 ⊆ N) :
    φ.InvScopeSat N X0
\end{leansrc}
\declhead{BernsteinLogic.QtslDischarge}{qBlockInv\_sound\_of\_dwell}{\badge{BadgeGreen}{statement bound}}
\begin{leansrc}
/-- QBlockInv soundness with G1 at its residue: the interface fact
drops from "the block is invariant" to "boundary points dwell"
(G1b — the facet-row bridge). -/
theorem qBlockInv_sound_of_dwell {φ : Semiflow X} {N X0 : Set X}
    (hN : IsClosed N)
    (htraj : ∀ x, Continuous fun t : ℝ≥0 => φ.flow t x)
    (hdwell : ∀ x ∈ N, x ∈ frontier N →
      ∃ δ : ℝ≥0, 0 < δ ∧ ∀ s ≤ δ, φ.flow s x ∈ N)
    (hInit : X0 ⊆ N) : φ.InvScopeSat N X0
\end{leansrc}
\declhead{BernsteinLogic.QtslDischarge}{invariant\_of\_boundaryDwell}{\badge{BadgeGreen}{statement bound}}
\begin{leansrc}
/-- G1's tier-A "first-exit argument", mechanized. A closed region with
continuous trajectories is forward-invariant as soon as every BOUNDARY
point dwells (stays in the region for some positive time) — interior
points dwell by continuity alone. The residue G1b (certified strict
facet-inward sign rows ⇒ the boundary dwell) is the polynomial-field
bridge, prose tier A, same class as G2.

Proof shape = the note's first-exit: the set of times whose whole
prefix stays in `N` is nonempty, closed under the sup (closedness +
continuity), and extends past its sup by the dwell at the sup point —
so it is unbounded, i.e. all of `ℝ≥0`. -/
theorem invariant_of_boundaryDwell {φ : Semiflow X} {N : Set X}
    (hN : IsClosed N)
    (htraj : ∀ x, Continuous fun t : ℝ≥0 => φ.flow t x)
    (hdwell : ∀ x ∈ N, x ∈ frontier N →
      ∃ δ : ℝ≥0, 0 < δ ∧ ∀ s ≤ δ, φ.flow s x ∈ N) :
    φ.Invariant N
\end{leansrc}
\declhead{BernsteinLogic.QtslDischarge}{HasFacetRows.boundaryDwell}{\badge{BadgeGreen}{statement bound}}
\begin{leansrc}
/-- G1b mechanized at its honest stratum: facet rows + block-in-arena +
a forward-invariant arena ⇒ the boundary dwell that
`invariant_of_boundaryDwell` consumes. Proof shape = the entry half of
the interior-entry argument ("interior entry is the same sign fact
read forward"): at a boundary point, trajectory continuity confines a short
prefix to the wedge neighborhood and every active certified region at
once (finitely many); on that prefix each active observable grows from
`0` at rate `ε i` — nonnegative — arena invariance supplies the `D`
conjunct, and the wedge test re-admits every prefix point. Inactive
facets never appear: the wedge test quotes only active ones. -/
theorem HasFacetRows.boundaryDwell {φ : Semiflow X} {D N : Set X}
    (h : φ.HasFacetRows D N) (hND : N ⊆ D) (hD : φ.Invariant D)
    (htraj : ∀ x, Continuous fun t : ℝ≥0 => φ.flow t x) :
    ∀ x ∈ N, x ∈ frontier N →
      ∃ δ : ℝ≥0, 0 < δ ∧ ∀ s ≤ δ, φ.flow s x ∈ N
\end{leansrc}
\declhead{BernsteinLogic.QtslDischarge}{qBlockInv\_sound\_of\_facetRows}{\badge{BadgeGreen}{statement bound}}
\begin{leansrc}
/-- QBlockInv soundness from the certificate's facet rows — the
composed G1 corollary at the current floor. Standing interface facts,
each at its named honest stratum (module header, G1 entry): `hN` and
`hND` are grid facts of the block (compact box-union in its arena),
`hD` is the model-class arena fact (`invariant_univ` when the block
has no ∂D-faces), `hrows` carries the G2-class integrated row residue
and the R1-class wedge geometry, `htraj` is flow continuity (G3's
model-class clause). -/
theorem qBlockInv_sound_of_facetRows {φ : Semiflow X} {D N X0 : Set X}
    (hN : IsClosed N)
    (htraj : ∀ x, Continuous fun t : ℝ≥0 => φ.flow t x)
    (hND : N ⊆ D) (hD : φ.Invariant D) (hrows : φ.HasFacetRows D N)
    (hInit : X0 ⊆ N) : φ.InvScopeSat N X0
\end{leansrc}
\declhead{BernsteinLogic.QtslMeanValue}{qBlockInv\_sound\_of\_rows}{\badge{BadgeGreen}{statement bound}}
\begin{leansrc}
/-- QBlockInv soundness from pointwise rows — G1's chain at its new
floor: per-facet row bounds + profile + wedge geometry + arena
invariance + model-class regularity ⇒ `InvScopeSat`. Composes
`hasFacetRows_of_rows` into `qBlockInv_sound_of_facetRows`; every
hypothesis is now either R1-recomputable grid geometry, a model-class
fact, exact row arithmetic, or the profile identification. -/
theorem qBlockInv_sound_of_rows {φ : Semiflow X} {D N X0 : Set X}
    {n : ℕ} {ℓ : Fin n → X → ℝ} {ε : Fin n → ℝ} {F : Fin n → Set X}
    (W : Fin n → X → ℝ)
    (hN : IsClosed N)
    (htraj : ∀ x, Continuous fun t : ℝ≥0 => φ.flow t x)
    (hND : N ⊆ D) (hD : φ.Invariant D)
    (hε : ∀ i, 0 < ε i)
    (hℓ : ∀ i, ContinuousOn (ℓ i) (F i))
    (hW : ∀ i, φ.FlowDerivOn (ℓ i) (W i) (F i))
    (hrow : ∀ i, ∀ y ∈ F i, ε i ≤ W i y)
    (hlocal : ∀ x ∈ frontier N, ∃ (A : Set (Fin n)) (U : Set X), U ∈ 𝓝 x ∧
      (∀ i ∈ A, ℓ i x = 0 ∧ F i ∈ 𝓝 x) ∧
      ∀ y ∈ U, y ∈ D → (∀ i ∈ A, 0 ≤ ℓ i y) → y ∈ N)
    (hInit : X0 ⊆ N) : φ.InvScopeSat N X0
\end{leansrc}
\declhead{BernsteinLogic.QtslGridGeometry}{qBlockInv\_sound\_of\_grid}{\badge{BadgeGreen}{statement bound}}
\begin{leansrc}
/-- QBlockInv soundness on the grid model, end to end: grid geometry
discharged outright; the remaining hypotheses are exactly the named
residues — the arena invariance interface `hArena` (residue (c), see
the two reduction theorems below) and the integrated row clause
(residue (a), G2's class) — plus flow continuity (G3's model-class
clause). -/
theorem qBlockInv_sound_of_grid (φ : Semiflow (Fin d → ℝ)) (B : Finset G.CellIdx)
    {X0 : Set (Fin d → ℝ)}
    (htraj : ∀ x, Continuous fun t : ℝ≥0 => φ.flow t x)
    (hArena : φ.Invariant G.arena)
    (ε : G.FacetData → ℝ) (F : G.FacetData → Set (Fin d → ℝ))
    (hε : ∀ p, G.IsFacet B p → 0 < ε p)
    (hF : ∀ p, G.IsFacet B p → ∀ x ∈ G.facetBox p, F p ∈ 𝓝 x)
    (hrow : ∀ p, G.IsFacet B p → ∀ x ∈ F p, ∀ t : ℝ≥0,
      (∀ s ≤ t, φ.flow s x ∈ F p) →
      G.obs p x + ε p * (t : ℝ) ≤ G.obs p (φ.flow t x))
    (hInit : X0 ⊆ G.blockRegion B) :
    φ.InvScopeSat (G.blockRegion B) X0
\end{leansrc}
\declhead{BernsteinLogic.QtslGridGeometry}{qBlockInv\_sound\_of\_pointwiseRows}{\badge{BadgeGreen}{statement bound}}
\begin{leansrc}
/-- CAPSTONE of the same-day gap-fill pair (this file × the interval
mean-value bridge, `QtslMeanValue.lean`): grid block + POINTWISE row
bounds, consumed through the profile predicate `FlowDerivOn`, ⇒
`InvScopeSat`. Every remaining hypothesis is a model-class regularity
fact (`htraj`, `hV`), exact row arithmetic (`hε`, `hrow`), the arena
interface (`hArena`, residue (c)'s reductions), the certified-region
neighborhood fact (`hF`), or the profile identification (`hW` —
`W p = 𝒜(obs p)` for the ODE doctrine, the sole analytic remainder of
the whole G1 chain). -/
theorem qBlockInv_sound_of_pointwiseRows (φ : Semiflow (Fin d → ℝ))
    (B : Finset G.CellIdx) {X0 : Set (Fin d → ℝ)}
    (htraj : ∀ x, Continuous fun t : ℝ≥0 => φ.flow t x)
    (hArena : φ.Invariant G.arena)
    (ε : G.FacetData → ℝ) (F : G.FacetData → Set (Fin d → ℝ))
    (W : G.FacetData → (Fin d → ℝ) → ℝ)
    (hε : ∀ p, G.IsFacet B p → 0 < ε p)
    (hF : ∀ p, G.IsFacet B p → ∀ x ∈ G.facetBox p, F p ∈ 𝓝 x)
    (hV : ∀ p, G.IsFacet B p → ContinuousOn (G.obs p) (F p))
    (hW : ∀ p, G.IsFacet B p → φ.FlowDerivOn (G.obs p) (W p) (F p))
    (hrow : ∀ p, G.IsFacet B p → ∀ y ∈ F p, ε p ≤ W p y)
    (hInit : X0 ⊆ G.blockRegion B) :
    φ.InvScopeSat (G.blockRegion B) X0
\end{leansrc}
\ruleblock{Correspondence (Steele display $\leftrightarrow$ Lean)}
\badge{BadgeAmber}{a bound Lean statement CHANGED since certification --- re-review and re-certify}
\sourcenote{The dictionary below is a reviewer-certified statement of how the mechanized statements above interpret the displayed rule: display metavariable $\mapsto$ Lean binder, each display hypothesis $\mapsto$ the premise or side condition carrying it, the display conclusion $\mapsto$ the goal --- with any deliberate difference in width between display and mechanization named outright. Certification is recorded as a hash triple over (the displayed cell, the extracted statement texts, this dictionary\textquotesingle s own text); if any of the three changes, this block turns amber until a reviewer re-certifies. The correspondence claim itself remains editorial --- reviewer-grade, not machine-checked; its machine-checked successor types each soundness statement against the rule as a structured object, and supersedes this block per rule as it lands.}
The displayed judgment $\mathcal{M} \vDash \mathsf{QInvScope}(N, X_0)$
$\mapsto$ the goal \texttt{InvScopeSat N X0} = \texttt{φ.Invariant N ∧
X0 ⊆ N} --- the claim record\textquotesingle s denotation verbatim.
The five bound statements form one discharge chain:
\texttt{qBlockInv\_sound} (the roster\textquotesingle s ``two lines'':
invariance + containment $\Rightarrow$ the claim, M(i) as hypothesis),
\texttt{\_of\_dwell} riding \texttt{invariant\_of\_boundaryDwell} (G1
mechanized: closed block + trajectory continuity + boundary dwell
$\Rightarrow$ invariance, the first-exit sup-time argument), and
\texttt{\_of\_facetRows} riding \texttt{HasFacetRows.boundaryDwell},
the deepest bound form, whose kit is the honest ledger:
\begin{itemize}
\item the block $N : \mathsf{QRegionAtom}$ $\mapsto$ an abstract
  \texttt{N : Set X} with \texttt{hN : IsClosed N}: the grid stratum
  (bitmap realizes a closed box-union in its arena; facet-list
  well-formedness) is checker-side recompute, and closedness is the
  one grid fact the bound statements consume --- as an explicit
  hypothesis, not a derived fact.
\item facet-inward sign rows $\mapsto$ \texttt{hrows :
  φ.HasFacetRows D N}: finitely many facet observables $\ell_i$ with
  explicit positive margins (\texttt{∀ i, 0 < ε i} --- the
  display\textquotesingle s strictness-by-margin, carried exactly),
  each growing at \emph{integrated} rate $\ge \varepsilon_i$ along
  confined prefixes --- the pointwise row $\mathcal{A}\ell \ge
  \varepsilon$ reaches that form through the G2-class mean-value
  bridge, outside this bound set --- plus the local wedge test,
  one-directional by design (reentrant corners of box-unions defeat
  the $\leftrightarrow$ form), with ``$\partial D$-faces owe
  nothing'' landing as the wedge test\textquotesingle s \texttt{y ∈
  D} conjunct.
\item \texttt{init\_containment} $\mapsto$ \texttt{hInit : X0 ⊆ N},
  one-to-one, unchanged at every link of the chain.
\item the displayed \textsf{FacetBridge} package (its contents
  defined in the flavor text) $\mapsto$
  \texttt{htraj} (trajectory continuity --- model-class), and the
  arena pair \texttt{hND : N ⊆ D} + \texttt{hD : φ.Invariant D} ---
  the certified statement makes the invariance regress explicit (the
  roster witness names its termination at \texttt{invariant\_univ},
  which is in the module but not this bound set).
\end{itemize}
Width note: the conditional flag\textquotesingle s three named
families are visible verbatim in the facet-row
statement\textquotesingle s binders --- the integrated row clause
(inside \texttt{HasFacetRows}), the grid facts (\texttt{hN}, plus the
wedge-test geometry carried \emph{inside} \texttt{hrows} rather than
discharged), and arena invariance (\texttt{hD}). The bound set
extends past the two-line core (the two blocks below): the
row-to-shadow discharge \texttt{qBlockInv\_sound\_of\_rows} (module
\texttt{QtslMeanValue} --- pointwise rows reach the integrated shadow
form through the interval mean-value bridge, leaving the profile
identification \texttt{hW} as the sole analytic remainder of the G1
chain), the grid-model discharge \texttt{qBlockInv\_sound\_of\_grid}
(module \texttt{QtslGridGeometry} --- grid geometry discharged
outright, residues exactly the arena interface and the integrated row
clause), and the capstone \texttt{qBlockInv\_sound\_of\_pointwiseRows}
(grid block + pointwise rows through \texttt{FlowDerivOn} to the
judgment) are hash-bound, so the conditional residue shrinks to
the profile identification plus the arena-invariance interface
(its reduction theorems terminate at \texttt{invariant\_univ}), the
facet-neighbourhood interface clause (each certified region $F_i$ a
neighbourhood of every point of its facet box, corners included ---
the capstone consumes it as an explicit hypothesis) and
model-class regularity. One deliberate width inversion: the
display\textquotesingle s conclusion tie is an unconditional judgment
about the model, while \texttt{qBlockInv\_sound} alone is wider than
the rule (any invariant set qualifies); the facet-row form narrows
back to certificate-shaped input. No hypothesis in the bound set
demands the block be a box-union --- the geometry enters only through
what the wedge test can actually be discharged for.

\ruleblock{Lean verifier receipts}
\sourcenote{module \texttt{BernsteinLogic.QtslSoundness}; bound build \texttt{77005991d61d\ldots} (bound from receipt 8130a7a7a679cef4, 2026-08-04T01:16:06Z). A receipt records that named checking layers accepted a build: (module, compiled-artifact SHA-256, environment SHA-256, checker layers), minted when the checking gate ran; the store receipts attesting the bound build are listed with the layers each run credited. Nothing here asks the reader to trust the store: every layer is an independently runnable program over the quoted sources --- rebuild the module and re-run the checkers to reproduce the verdict. Receipts attesting superseded builds credit no layer and are omitted from this rendering (the count of omissions, if any, is noted below).}
\begin{receipttable}
\receiptrow{8130a7a7a679cef4}{2026-08-04T01:16:06Z}{build+decorrelated}{\badge{BadgeGreen}{current}}
\receiptrow{9eb3706354082144}{2026-08-04T02:23:12Z}{build+decorrelated}{\badge{BadgeGreen}{current}}
\end{receipttable}
\sourcenote{5 further receipt(s) in the store attest superseded builds of this module and are omitted from this rendering.}
\sourcenote{module \texttt{BernsteinLogic.QtslDischarge}; bound build \texttt{4aa46910a600\ldots} (bound from receipt dd2946805f790f99, 2026-08-04T01:16:09Z). A receipt records that named checking layers accepted a build: (module, compiled-artifact SHA-256, environment SHA-256, checker layers), minted when the checking gate ran; the store receipts attesting the bound build are listed with the layers each run credited. Nothing here asks the reader to trust the store: every layer is an independently runnable program over the quoted sources --- rebuild the module and re-run the checkers to reproduce the verdict. Receipts attesting superseded builds credit no layer and are omitted from this rendering (the count of omissions, if any, is noted below).}
\begin{receipttable}
\receiptrow{dd2946805f790f99}{2026-08-04T01:16:09Z}{build+decorrelated}{\badge{BadgeGreen}{current}}
\receiptrow{e1bc68890e2aebf1}{2026-08-04T02:23:11Z}{build+decorrelated}{\badge{BadgeGreen}{current}}
\end{receipttable}
\sourcenote{14 further receipt(s) in the store attest superseded builds of this module and are omitted from this rendering.}
\sourcenote{module \texttt{BernsteinLogic.QtslMeanValue}; bound build \texttt{ed04b606934b\ldots} (bound from receipt c2f57d21d0b1c270, 2026-07-20T08:27:08Z). A receipt records that named checking layers accepted a build: (module, compiled-artifact SHA-256, environment SHA-256, checker layers), minted when the checking gate ran; the store receipts attesting the bound build are listed with the layers each run credited. Nothing here asks the reader to trust the store: every layer is an independently runnable program over the quoted sources --- rebuild the module and re-run the checkers to reproduce the verdict. Receipts attesting superseded builds credit no layer and are omitted from this rendering (the count of omissions, if any, is noted below).}
\begin{receipttable}
\receiptrow{160082fd7894faf1}{2026-07-19T07:39:26Z}{build+decorrelated}{\badge{BadgeGreen}{current}}
\receiptrow{260eab5039951515}{2026-07-19T09:07:22Z}{build+decorrelated}{\badge{BadgeGreen}{current}}
\receiptrow{c2f57d21d0b1c270}{2026-07-20T08:27:08Z}{build+decorrelated}{\badge{BadgeGreen}{current}}
\receiptrow{f75d120f9128f80c}{2026-07-20T20:44:49Z}{build+decorrelated}{\badge{BadgeGreen}{current}}
\receiptrow{891ea7ab9c4f79ff}{2026-07-20T21:01:17Z}{build+replay+hardened+decorrelated}{\badge{BadgeGreen}{current}}
\receiptrow{af025f8810d8a6cc}{2026-07-21T05:28:44Z}{build+replay+hardened+decorrelated}{\badge{BadgeGreen}{current}}
\end{receipttable}
\sourcenote{module \texttt{BernsteinLogic.QtslGridGeometry}; bound build \texttt{9060a661920e\ldots} (bound from receipt 64d6a21041b4042f, 2026-07-20T08:27:10Z). A receipt records that named checking layers accepted a build: (module, compiled-artifact SHA-256, environment SHA-256, checker layers), minted when the checking gate ran; the store receipts attesting the bound build are listed with the layers each run credited. Nothing here asks the reader to trust the store: every layer is an independently runnable program over the quoted sources --- rebuild the module and re-run the checkers to reproduce the verdict. Receipts attesting superseded builds credit no layer and are omitted from this rendering (the count of omissions, if any, is noted below).}
\begin{receipttable}
\receiptrow{bf51898b80b37884}{2026-07-19T12:35:27Z}{build+decorrelated}{\badge{BadgeGreen}{current}}
\receiptrow{64d6a21041b4042f}{2026-07-20T08:27:10Z}{build+decorrelated}{\badge{BadgeGreen}{current}}
\receiptrow{c7deda94ed05f8e5}{2026-07-20T20:44:45Z}{build+decorrelated}{\badge{BadgeGreen}{current}}
\receiptrow{6bab7a61aa41db61}{2026-07-20T21:01:14Z}{build+replay+hardened+decorrelated}{\badge{BadgeGreen}{current}}
\receiptrow{5e72e577b5f436e5}{2026-07-21T05:28:39Z}{build+replay+hardened+decorrelated}{\badge{BadgeGreen}{current}}
\end{receipttable}
\sourcenote{2 further receipt(s) in the store attest superseded builds of this module and are omitted from this rendering.}
\layerline{\layeritem{build}{lake build (elaboration)}{\badge{BadgeGreen}{current}}\layeritem{replay}{leanchecker (official kernel replay)}{\badge{BadgeGreen}{current}}\layeritem{hardened}{safe\_verify (hardened screen)}{\badge{BadgeGreen}{current}}\layeritem{decorrelated}{lean4export + nanoda (decorrelated kernel)}{\badge{BadgeGreen}{current}}}
\sourcenote{\badge{BadgeGray}{comparator: none} no fixture in the three-implementation comparator harness (compiled / kernel-interpreted / reference-interpreted) exercises this rule.}
\ruleblock{HOL4 statement}
\sourcenote{theory \texttt{bern\_tr\_t5qdyn}, HOL4 development script \texttt{\texttt{bern\_tr\_t5qdynScript.sml}} \badge{BadgeGreen}{script bound}}
\holhead{Definition sfInvariant\_def}
\begin{holsrc}
Definition sfInvariant_def:
  sfInvariant (fl:real->'x->'x) (sv:'x->bool) ⇔
    ∀x. x IN sv ⇒ ∀t. 0 ≤ t ⇒ fl t x IN sv
End
\end{holsrc}
\holhead{Definition sfInvScopeSat\_def}
\begin{holsrc}
Definition sfInvScopeSat_def:
  sfInvScopeSat (fl:real->'x->'x) (sv:'x->bool) (X0:'x->bool) ⇔
    sfInvariant fl sv ∧ X0 SUBSET sv
End
\end{holsrc}
\holhead{Theorem sfQBlockInv\_sound}
\begin{holsrc}
Theorem sfQBlockInv_sound:
  ∀fl (N:'x->bool) X0.
    sfInvariant fl N ∧ X0 SUBSET N ⇒ sfInvScopeSat fl N X0
\end{holsrc}
\ruleblock{HOL verifier receipts}
\badge{BadgeBlue}{documentary} \texttt{hol-strict} CLEAN: theory \texttt{bern\_tr\_t5qdyn}, 22
theorems, 0 cheat-tagged, 0 nonstandard axioms (whitelist:
\texttt{DISK\_THM}; the four \texttt{bool} axioms) --- verification
run of 2026-07-22 on an isolated build host, against a
content-addressed checkout of the exact sources cited here (git blob
\texttt{d64e1954\ldots}); an independent acceptance host ran the same
build green at the porting seating (2026-07-21). What is ported is
the semiflow-stratum root: the flow law, the invariance and
invariance-scope-satisfaction predicates, and the root soundness
theorem concluding the scope judgment from invariance plus
containment --- the exact root the Lean development\textquotesingle s
satisfaction stratum supplies for this rule. The facet-row bridge
(the first-exit skeleton and integrated-growth descent that discharge
invariance from the certificate\textquotesingle s sign rows) and the
grid-geometry recomputation remain Lean-side; the conditional reading
of this section\textquotesingle s mechanization carries over to the
HOL twin unchanged. The facet-sign and containment rows the rule
consumes are covered by the ported base checker as before.

\sourcenote{HOL receipts are documentary: quoted results of \texttt{hol-strict} runs
and census rows from the HOL4 development\textquotesingle s build
evidence, not minted receipt artefacts. A HOL receipt store (the Lean
store\textquotesingle s twin) is named future work.
}
\ruleblock{Proof}

\sourcenote{Genre: core rule (tier 1) --- an ordinary mathematical proof
at Putnam-exam-solution detail. The rule's mechanization is
\emph{conditional}: the section descends the invariance obligation from
an opaque interface fact to three named hypothesis families ---
integrated facet-row clauses, grid geometry of the block, arena
invariance --- and the ledger remark at the end states exactly what
remains and why. Every hypothesis is used and remarked on; the
formal-anchor table maps each step to the Lean declaration that carries
it.}

\textbf{Setting and claim.}\enspace Fix a forward semiflow $\varphi$ on
a state space $X$ ($\Phi_0 = \mathrm{id}$,
$\Phi_{s+t} = \Phi_t \circ \Phi_s$; the \rulename{QBasin} section's
setting, which also fixes the $\mathsf{QInvScope}$ denotation: $S$
forward-invariant and $X_0 \subseteq S$). The rule's data: a
\emph{block} $N$ --- the realization of a box-union over the
certificate's declared rational grid, a compact closed set --- sitting
in an \emph{arena} $D \supseteq N$; per internal facet of $N$, one
strictly-inward sign row of the base calculus (a
\rulename{BernCone}-family sequent asserting the field's normal
component points into the block, with an explicit positive rational
margin; facets on $\partial D$ owe nothing); and one containment row
$X_0 \subseteq N$. The claim: these data warrant
$\mathsf{QInvScope}(N, X_0)$ --- established below at the discharged
floor, conditional on the named ledger of
\autoref{rem:qbi-ledger}.

\begin{blemma}[rule-level assembly]\label{lem:qbi-assembly}
If $N$ is forward-invariant and $X_0 \subseteq N$, then
$\mathsf{QInvScope}(N, X_0)$ holds.
\end{blemma}
\begin{proof}
The judgment \emph{is} that conjunction. The lemma is stated to fix
where the rule-level content ends: everything below descends the
invariance conjunct toward the certificate's rows.
\end{proof}

\begin{blemma}[the first-exit skeleton: closed $+$ continuous $+$
boundary dwell $\Rightarrow$ invariant]\label{lem:qbi-dwell}
Suppose $N$ is closed, every trajectory $t \mapsto \Phi_t x$ is
continuous, and every boundary point dwells: for each
$x \in N \cap \partial N$ there is $\delta > 0$ with
$\Phi_s x \in N$ for all $0 \leq s \leq \delta$. Then $N$ is
forward-invariant.
\end{blemma}
\begin{proof}
First, \emph{every} point of $N$ dwells: an interior point $x$ has
$\Phi_0 x = x \in \mathrm{int}\,N$, so continuity of the trajectory at
time $0$ confines a short prefix to $\mathrm{int}\,N$ --- the dwell
hypothesis is needed only at the frontier, which is where the
certificate's rows live.

Now fix $x \in N$ and $t \geq 0$, and suppose toward a contradiction
that $\Phi_t x \notin N$. Consider the set of \emph{prefix-confined
times}
\[
S \;:=\; \{\, u \geq 0 \ :\ \Phi_s x \in N \ \text{for all } s \leq u
\,\}.
\]
$S$ contains $0$, and every element of $S$ is $\leq t$ (an element
$> t$ would confine time $t$ itself); throughout, times $s, u$ range
over $[0, \infty)$. Let $T := \sup S \leq t$. Every
time $s < T$ is confined (it is below some element of $S$). The sup
itself is confined: for $T > 0$ the preimage of $N$ under the
continuous trajectory is closed and contains $[0, T)$, hence contains
$T \in \mathrm{cl}\,[0,T)$; for $T = 0$ it is the hypothesis
$x \in N$. So $T \in S$. But now the dwell at the point $\Phi_T x \in N$
supplies $\delta > 0$ with $\Phi_s(\Phi_T x) \in N$ for
$0 \leq s \leq \delta$,
i.e. (semiflow law) $\Phi_u x \in N$ for $T \leq u \leq T + \delta$;
combined with $T \in S$ this puts $T + \delta \in S$, contradicting
$T = \sup S$ with $\delta > 0$.
\end{proof}

\begin{bremark}[each skeleton hypothesis is load-bearing]
\emph{Closedness}: on $X = \R$ with $\Phi_t x = x + t$ and the
half-open $N = [0, 1)$, every point of $N$ dwells (interior points by
continuity; the one frontier point of $N$, namely $0$, dwells since
$\Phi_s 0 = s \in [0,1)$ for small $s$ --- the other frontier point
$1$ is not in $N$), yet the trajectory from $0$ exits: the proof fails
exactly at
``the sup itself is confined'', which is where closedness is consumed.
\emph{Trajectory continuity}: with a discontinuous time-dependence the
preimage argument dies and a trajectory may jump across the frontier
without ever presenting a boundary point to the dwell hypothesis; in
this stratum continuity is a model-class fact (polynomial-drift flows),
not a theorem. \emph{Dwell at the frontier}: dropping it re-opens exit
through the boundary --- it is precisely the obligation the facet rows
discharge (next lemma).
\end{bremark}

\begin{blemma}[the facet-row bridge: certified rows $\Rightarrow$
boundary dwell]\label{lem:qbi-bridge}
Suppose the certificate supplies, for block $N$ in arena $D$: finitely
many \emph{facet observables} $\ell_1, \dots, \ell_n : X \to \R$ (the
affine facet polynomials), margins $\varepsilon_i > 0$, and certified
regions $F_i$ (boxes around the facets, thickened to full dimension
\emph{with tangential margin}: each $F_i$ must be a neighbourhood of
every point of its facet, corners included --- a box thickened only in
the normal direction is a neighbourhood of no corner point, and the
wedge test below demands one), such that

\emph{(row, integrated form)} along any trajectory segment confined to
$F_i$, $\ell_i$ grows at rate at least $\varepsilon_i$:
if $\Phi_s x \in F_i$ for all $s \leq t$, then
$\ell_i(\Phi_t x) \geq \ell_i(x) + \varepsilon_i\, t$; and

\emph{(local wedge test)} every frontier point $x$ of $N$ has an
\emph{active set} $A \subseteq \{1, \dots, n\}$ --- each active
observable vanishing at $x$, each active $F_i$ a neighbourhood of $x$
--- and a neighbourhood $U$ of $x$ such that every $y \in U$ with
$y \in D$ and $\ell_i(y) \geq 0$ for all $i \in A$ lies in $N$.

If moreover $N \subseteq D$, $D$ is forward-invariant, and trajectories
are continuous, then every boundary point of $N$ dwells.
\end{blemma}
\begin{proof}
Fix $x \in N \cap \partial N$ with active set $A$ and wedge
neighbourhood $U$. The set $U \cap \bigcap_{i \in A} F_i$ is a finite
intersection of neighbourhoods of $x$, hence a neighbourhood; continuity
of the trajectory at time $0$ (where it sits at $x$) confines a prefix
$[0, \delta]$ to it, some $\delta > 0$. Now take any
$0 \leq s \leq \delta$
and check the wedge test's three conditions at $y = \Phi_s x$:
membership in $U$ is the confinement; membership in $D$ is arena
invariance applied to $x \in N \subseteq D$; and for each active $i$,
the whole prefix up to $s$ is confined to $F_i$, so the integrated row
gives $\ell_i(\Phi_s x) \geq \ell_i(x) + \varepsilon_i s
= 0 + \varepsilon_i s \geq 0$, using exactly the activation fact
$\ell_i(x) = 0$. The wedge test re-admits $\Phi_s x \in N$. Inactive
facets never appear: the test quotes only active ones.
\end{proof}

\begin{bthm}[soundness of \rulename{QBlockInv}, at the discharged
floor]\label{thm:qbi-sound}
Suppose: $N$ is closed with $N \subseteq D$; $D$ is forward-invariant;
trajectories are continuous; the certificate's facet rows satisfy
Lemma~\ref{lem:qbi-bridge}'s two clauses; and $X_0 \subseteq N$. Then
$\mathsf{QInvScope}(N, X_0)$ holds.
\end{bthm}
\begin{proof}
Lemma~\ref{lem:qbi-bridge} delivers the boundary dwell;
Lemma~\ref{lem:qbi-dwell} converts it to forward invariance of $N$;
Lemma~\ref{lem:qbi-assembly} assembles the judgment with the
containment row.
\end{proof}

\begin{bremark}[the wedge test is one-directional by design]
\label{rem:qbi-wedge}
At a \emph{reentrant} corner of a box-union block --- an inner corner
of an L-shape --- the block is locally the \emph{union} of two
half-spaces $\{\ell_1 \geq 0\} \cup \{\ell_2 \geq 0\}$, while the
wedge test's hypothesis set is their \emph{intersection}: the
implication ``both $\ell_i \geq 0$ $\Rightarrow$ $y \in N$'' holds
(the intersection sits inside the union), but its converse is false
there --- points of $N$ with one active observable negative exist on
the other branch. The biconditional form of the test would therefore
be unsatisfiable at reentrant corners, while the one-directional form
is exactly what the first-exit argument consumes
(Lemma~\ref{lem:qbi-bridge} never needs to \emph{characterize} the
frontier, only to re-admit prefix points). A certificate whose rows
are jointly unsatisfiable at a boundary point --- checkerboard blocks
touching at a corner and demanding opposite strict signs there ---
simply makes the bridge's hypotheses unfulfillable: soundness only
consumes what is given.
\end{bremark}

\begin{bremark}[$\partial D$-faces owe nothing; the arena regress
terminates]
Facets of $N$ on the arena boundary carry no rows: in
Lemma~\ref{lem:qbi-bridge} they enter only through the conjunct
$y \in D$, which arena invariance supplies. The arena's own invariance
is either a nested $\mathsf{QInvScope}$ citation (a chart sub-block
certified by its own facet rows --- this same rule, one level up) or
the model-class fact that a chart's denotation confines trajectories to
its domain; the regress terminates at the whole space, which is
invariant vacuously.
\end{bremark}

\begin{bremark}[strictness of the margins: where it is and is not
consumed]\label{rem:qbi-strict}
The invariance conclusion of this section survives
$\varepsilon_i = 0$: the integrated clause then says $\ell_i$ is
non-decreasing along confined segments, and growth-from-zero still
yields $\ell_i \geq 0$ --- consistent with the classical invariance
theorem (Nagumo \cite{nagumo1942}), which needs only subtangentiality.
The explicit positive rational margin is nonetheless mandatory in the
certificate convention, for two reasons that live outside this
section. First, exactness: a \emph{strict} field sign cannot be
certified directly by the non-strict enclosure predicates of the base
calculus; ``non-strict row on $\mathcal{A}\ell - \varepsilon$ with
explicit $\varepsilon > 0$'' is how strictness becomes checkable in
exact $\Q$. Second, the same rows are consumed by the Morse-family
rules, where strictness is genuinely load-bearing --- no orbit rides a
facet, distinct fibres' invariant parts are disjoint, and the facet
margin bounds admissible chain jumps (``strictness pays a second
time''; see the \rulename{QMorseDischarge} section's ledger). A margin
quietly dropped to zero would leave this rule's conclusion intact and
break those.
\end{bremark}

\begin{bremark}[the hypothesis ledger --- what remains, named]
\label{rem:qbi-ledger}
The chain above discharges the rule's invariance obligation down to
three hypothesis families, and the rule's Lean column renders with the
hypothesis-flavoured check because of exactly these:

\emph{(a) the integrated row clause: discharged, and bound.} The
certificate's rows are pointwise: per facet box, a base-calculus
sequent bounding the field's certified rate,
$\mathcal{A}\ell_i \geq \varepsilon_i$ (the generator
$\mathcal{A} = f \cdot \nabla$ of the deterministic model class).
Lemma~\ref{lem:qbi-bridge} consumes the \emph{integrated} clause along
trajectory segments. The bridge between them is one-variable calculus
--- a mean-value argument along the trajectory, mechanized and bound
in this edition (anchor table) --- and its surviving analytic content
is the identification of the certified polynomial rate with the
trajectory-wise derivative on the region: a model-class profile
warrant (for smooth flows the identification is literal; switched and
hybrid profiles replace it by one-sided derivative bounds, which is
their own seam, deliberately not this section's).

\emph{(b) grid geometry of the block: closed for grid blocks, and
bound --- except one interface clause.} Closedness of $N$,
$N \subseteq D$, the vanishing of each facet observable on its facet
box, and the wedge re-admission implication are
finite geometry of the box-union over its declared rational grid ---
recomputable combinatorics, proved outright for grid blocks in the
grid-geometry module and bound in this edition (anchor table),
end-to-end through the capstone that takes grid block plus pointwise
row bounds to the concluded judgment. The one geometric clause that
stays interface-level even there: each certified region $F_i$ must be
a neighbourhood of every point of its facet box, corners included ---
the tangential-margin requirement of Lemma~\ref{lem:qbi-bridge} ---
which the capstone consumes as a hypothesis on the certificate's
regions rather than proving from the boxes. The interface-level
theorems above remain stated with (b) as hypotheses, which is what the
general (non-grid) reading consumes.

\emph{(c) arena invariance.} As in the regress remark: nested
citation, or model-class warrant, terminating at the whole space.

None of the three is an opaque dynamical assumption; each names a
specific fact with a specific owner --- (a) the profile stratum, (b)
the checker's recompute (and the grid module), (c) the arena's own
certificate. That is the sense in which this core rule is
\emph{conditionally} mechanized: the cited theorems are sorry-free,
and their hypotheses are the ledger above.
\end{bremark}

\begin{bremark}[relation to the literature]
The block-with-inward-facets certificate is the exact-arithmetic
descendant of Conley theory's isolating blocks
\cite{conleyeaston1971}: compact sets whose boundary behaviour is
transverse, so that dynamical claims reduce to boundary sign checks.
The invariance argument itself is Nagumo's theorem
\cite{nagumo1942} in first-exit form, specialized to polyhedral sets,
where the subtangentiality condition becomes finitely many polynomial
sign rows --- the form set-invariance control has long exploited
\cite{blanchini1999}. What the certificate stratum adds is the
discipline: strict signs via non-strict predicates with exact rational
margins (\autoref{rem:qbi-strict}), and a one-directional wedge
test chosen so that box-union (nonconvex) blocks are admissible
without any frontier characterization
(\autoref{rem:qbi-wedge}).
\end{bremark}

\begin{anchors}
Lemma \ref{lem:qbi-assembly} & \texttt{qBlockInv\_sound} (QtslSoundness.lean)\\
Lemma \ref{lem:qbi-dwell}, first-exit skeleton & \texttt{invariant\_of\_boundaryDwell} (QtslDischarge.lean); composed as \texttt{qBlockInv\_sound\_of\_dwell}\\
Lemma \ref{lem:qbi-bridge}, row shadow + wedge test & \texttt{HasFacetRows} (the definition); \texttt{HasFacetRows.boundaryDwell}\\
Theorem \ref{thm:qbi-sound}, end-to-end & \texttt{qBlockInv\_sound\_of\_facetRows}\\
arena regress terminus & \texttt{invariant\_univ}; nested-citation path \texttt{invariant\_of\_invScopeSat} (development)\\
ledger item (a), the mean-value bridge & \texttt{integrated\_ge\_of\_flowDeriv\_ge}, composed as \texttt{qBlockInv\_sound\_of\_rows} (QtslMeanValue; \textbf{bound in this edition})\\
ledger item (b), grid facts & \texttt{Grid.wedge\_containment}, \texttt{Grid.qBlockInv\_sound\_of\_grid}; capstone \texttt{Grid.qBlockInv\_sound\_of\_pointwiseRows} (QtslGridGeometry; \textbf{bound in this edition})\\
facet rows & strictly-inward \rulename{BernCone}-family sequents with explicit rational margin (base kernel)\\
\end{anchors}


\rulesection{73}{QLatticeInv}{alias for \textcolor{DiffBlue}{\rulename{QBlockInv}}}
\label{blcatrule-QLatticeInv}
\ruleblock{The rule (Steele metanotation)}
\sourcenote{as realized in the mechanized development; source row: \texttt{\texttt{QTSL-rules.toml}} --- \emph{QLatticeInv}. The metanotation statement is typeset by hand for notational quality; its mechanized ground truth is the Lean statement below, extracted from the sources at every build of this book and hash-bound. The rule \emph{roster} is machine-checked (every mechanized rule tag round-trips to exactly one catalogue row); the statement \emph{text} is editorial, checked by eye against the extracted Lean beside it.}
\begin{mathpar}
\inferrule*{
  L : \mathsf{QBlockLattice} \\
  C \ \text{(\texttt{downset}, a cell-set; membership and down-closedness recomputed)} \\
  \mathcal{M} \vDash \mathsf{QInvScope}\bigl(\textstyle\bigcup_{i \in C} \mathrm{blocks}(\mathrm{emb}\ i),\ X_0\bigr) \ \text{(the cited \rulename{QBlockInv} conclusion)} \\
  \vdash X_0 \subseteq \textstyle\bigcup\downarrow C \ \text{(\texttt{init\_containment})}
}{
  \mathcal{M} \vDash \mathsf{QInvScope}\bigl(\textstyle\bigcup\downarrow C,\ X_0\bigr)
}
\end{mathpar}

\par\medskip
\rulename{QLatticeInv} is the family object\textquotesingle s
invariance alias: any down-closed cell-set $C$ of the recomputed
readout names an invariant region $\bigcup\downarrow C$, and the rule
concludes its invariance-scope judgment by citation --- the unpacking
is a single \rulename{QBlockInv} application whose block is the
lattice member realizing the down-set (existence by join-closure) and
whose facet rows are that member\textquotesingle s own, already
present in the family object. No row is re-proven; the producer merely
selects which down-set encloses the initial region away from the avoid
sets.

What the checker recomputes: membership of the realization in the
lattice, and down-closedness of $C$ in the recomputed cell order ---
readout facts, not dynamics. What would break without them: a cell-set
that is not down-closed names a region orbits can leave (an orbit
descends the cell order out of it), and the exactness identity between
the down-set\textquotesingle s cell union and the
member\textquotesingle s realization fails --- already on a two-member
chain: the top member alone is not down-closed, and its cell omits
everything the bottom member contributes.

Mechanization state, plainly: the bound statements below are the
exactness identity and the soundness corollaries, at both semantic
strata (module \texttt{BernsteinLogic.QtslLatticeInv}).
\texttt{qLatticeInv\_exact} proves the identity that makes the
respelling exact --- for down-closed $C$, the union of $C$\textquotesingle
s cells equals the union of its join-irreducible realizations --- by
the discharge chain\textquotesingle s own well-founded descent through
lower covers (\texttt{exists\_cellMem\_subset},
over the same descent rows the eventual-member interface consumes ---
rows that assume no distributivity,
plus one row this alias adds: cell inside member, the readout closure
half\textquotesingle s set-level shadow). One sharp point, because the
page\textquotesingle s classical route consumes distributivity: the
mechanized route consumes none, and the row pack does not encode
join-primality --- absorption-positive families satisfy every row, and
the bound (union-form) identity holds in them regardless --- because
the descent runs inside a single witnessed member of $C$ and never
forms the join. The join-form reading (the union as the member
$N_C$\textquotesingle s realization) costs exactly
join-realizes-as-union, the checker\textquotesingle s recomputed
membership of the down-set\textquotesingle s realization, strictly
weaker than distributivity; that recompute is also what refuses the
diamond-lattice witness when its top is realized past the
atoms\textquotesingle{} union. The corollaries
\texttt{Semiflow.qLatticeInv\_sound} /
\texttt{Behavior.qLatticeInv\_sound} (and their
\texttt{\_of\_blockInv} end-to-end forms) are
$\mathrm{sound}_A = \mathrm{sound}_T \circ \mathrm{unpack}$ through
the cited \rulename{QBlockInv} conclusion, delivered as judgment-level
cargo. Unconditional: the lattice rows are the checker\textquotesingle
s recomputed readout facts entering as theorem parameters (R1 shadows),
and join-closure membership of the down-set\textquotesingle s
realization is deliberately \emph{not} consumed --- with
down-closedness consumed and the cited member judgment supplied, the
entailment holds for the bare cell union, so the rule is sound a
fortiori with the checker\textquotesingle s extra side condition. The
payload-level unpack elaborator is bound beside the corollaries (as
with the sibling extension rules; module
\texttt{BernsteinLogic.QtslElaborators}):
\texttt{LatticeInvPayload} reifies the wire payload --- the leg-7 row
pack, the down-set, the containment row ---
\texttt{LatticeInvPayload.unpack} is the verified transport onto the
\rulename{QBlockInv} payload (block $:=$ the down-set\textquotesingle
s realization; the containment row rewritten along the exactness
identity; facets a citation, never payload data), and the
\texttt{qLatticeInv\_sound\_elab} pair is $\mathrm{sound}_A =
\mathrm{sound}_T \circ \mathrm{unpack}$ with
\texttt{qBlockInv\_sound} as the literal $\mathrm{sound}_T$, at both
semantic strata. The tier is flagged, not forced:
if registry-backed citation of the member\textquotesingle s rows is
ruled heavier than the alias criterion allows, the rule reclassifies
as admissible with all machinery unchanged.

\cataloguenote{The warrant for the reading --- once an orbit enters a lattice member
it never leaves, so visited cells descend the cell order and down-sets
are forward-invariant --- is the same orbit combinatorics the Morse
discharge\textquotesingle s order clause rests on, consumed here at
its easy end.
}
\ruleblock{Lean statement}
\declhead{BernsteinLogic.QtslLatticeInv}{qLatticeInv\_exact}{\badge{BadgeGreen}{statement bound}}
\begin{leansrc}
/-- **prop:qli-exact, the exactness identity.** For a down-closed cell-set
`C`, the union of the cells of `C` equals the union of the realizations of
`C`'s join-irreducible members — the realization of the display's member
`N_C = ⋁ C` (join = union in the union-closed family object). The ⊇ half
is leg 7's well-founded descent through lower covers
(`Semiflow.exists_cellMem_subset`): a member point lands in the cell of a
join-irreducible below it, and down-closedness (through `horder`) pins
that cell inside `C`. The ⊆ half is `hcellSub`, cell inside member. This
is what makes the alias's respelling exact rather than merely `⊇`. -/
theorem qLatticeInv_exact
    (hdc : IsLowerSet C)
    (hJoin : ∀ b : β, ∀ x ∈ blocks b,
      ∃ i : ι, x ∈ blocks (emb i) ∧ blocks (emb i) ⊆ blocks b)
    (hcover : ∀ i : ι, blocks (cover i) ⊂ blocks (emb i))
    (hcellsSup : ∀ i : ι, blocks (emb i) \ blocks (cover i) ⊆ cells i)
    (horder : ∀ i j : ι, blocks (emb i) ⊆ blocks (emb j) → i ≤ j)
    (hcellSub : ∀ i : ι, cells i ⊆ blocks (emb i)) :
    (⋃ i ∈ C, cells i) = ⋃ i ∈ C, blocks (emb i)
\end{leansrc}
\declhead{BernsteinLogic.QtslLatticeInv}{member\_subset\_cells\_le}{\badge{BadgeGreen}{statement bound}}
\begin{leansrc}
/-- **"cells tile members"** (the page proof's first readout fact), at the
principal down-set: every point of a join-irreducible's realization lies
in the cell of some join-irreducible at-or-below it in the cell order.
The identity's ⊇ half, exposed as its own statement because the book page
names it. -/
theorem member_subset_cells_le
    (hJoin : ∀ b : β, ∀ x ∈ blocks b,
      ∃ i : ι, x ∈ blocks (emb i) ∧ blocks (emb i) ⊆ blocks b)
    (hcover : ∀ i : ι, blocks (cover i) ⊂ blocks (emb i))
    (hcellsSup : ∀ i : ι, blocks (emb i) \ blocks (cover i) ⊆ cells i)
    (horder : ∀ i j : ι, blocks (emb i) ⊆ blocks (emb j) → i ≤ j)
    (i : ι) :
    blocks (emb i) ⊆ ⋃ j ∈ Set.Iic i, cells j
\end{leansrc}
\declhead{BernsteinLogic.QtslLatticeInv}{Semiflow.qLatticeInv\_sound}{\badge{BadgeGreen}{statement bound}}
\begin{leansrc}
/-- **prop:qli-sound, through the target** — the alias soundness corollary
in its genre-native shape `sound_A = sound_T ∘ unpack`: the premise is the
cited `QBlockInv` node's conclusion on the member `N_C` (its realization:
the union of the down-set's join-irreducible realizations), delivered as
judgment-level cargo (the `qBasin_sound` precedent — the facet rows live
in the cited node, not here); the conclusion is the alias's displayed
judgment on `⋃↓C`. The bridge is `qLatticeInv_exact`, and there is no
other content — the alias genre working as intended. -/
theorem Semiflow.qLatticeInv_sound {φ : Semiflow X} {X0 : Set X}
    (hdc : IsLowerSet C)
    (hJoin : ∀ b : β, ∀ x ∈ blocks b,
      ∃ i : ι, x ∈ blocks (emb i) ∧ blocks (emb i) ⊆ blocks b)
    (hcover : ∀ i : ι, blocks (cover i) ⊂ blocks (emb i))
    (hcellsSup : ∀ i : ι, blocks (emb i) \ blocks (cover i) ⊆ cells i)
    (horder : ∀ i j : ι, blocks (emb i) ⊆ blocks (emb j) → i ≤ j)
    (hcellSub : ∀ i : ι, cells i ⊆ blocks (emb i))
    (hMember : φ.InvScopeSat (⋃ i ∈ C, blocks (emb i)) X0) :
    φ.InvScopeSat (⋃ i ∈ C, cells i) X0
\end{leansrc}
\declhead{BernsteinLogic.QtslLatticeInv}{Semiflow.qLatticeInv\_sound\_of\_blockInv}{\badge{BadgeGreen}{statement bound}}
\begin{leansrc}
/-- **rec:qli-filled, composed end-to-end**: the filled-in `QBlockInv`
application. Premises: forward invariance of the member's realization
(the cited facet-row chain's output — dischargeable per member via
`qBlockInv_sound_of_facetRows` and kin) and the rule's own containment
row, stated on the alias's displayed region and transported verbatim
across the identity. -/
theorem Semiflow.qLatticeInv_sound_of_blockInv {φ : Semiflow X} {X0 : Set X}
    (hdc : IsLowerSet C)
    (hJoin : ∀ b : β, ∀ x ∈ blocks b,
      ∃ i : ι, x ∈ blocks (emb i) ∧ blocks (emb i) ⊆ blocks b)
    (hcover : ∀ i : ι, blocks (cover i) ⊂ blocks (emb i))
    (hcellsSup : ∀ i : ι, blocks (emb i) \ blocks (cover i) ⊆ cells i)
    (horder : ∀ i j : ι, blocks (emb i) ⊆ blocks (emb j) → i ≤ j)
    (hcellSub : ∀ i : ι, cells i ⊆ blocks (emb i))
    (hInv : φ.Invariant (⋃ i ∈ C, blocks (emb i)))
    (hInit : X0 ⊆ ⋃ i ∈ C, cells i) :
    φ.InvScopeSat (⋃ i ∈ C, cells i) X0
\end{leansrc}
\declhead{BernsteinLogic.QtslLatticeInv}{Behavior.qLatticeInv\_sound}{\badge{BadgeGreen}{statement bound}}
\begin{leansrc}
/-- `Semiflow.qLatticeInv_sound` at the behavioral floor — the identity is
dynamics-free, so the alias transports over solution-trajectory-set
semantics verbatim (the two-directions carve: `QInvScope` is
safety-directional, sound here without existence hypotheses). -/
theorem Behavior.qLatticeInv_sound {B : Behavior X} {X0 : Set X}
    (hdc : IsLowerSet C)
    (hJoin : ∀ b : β, ∀ x ∈ blocks b,
      ∃ i : ι, x ∈ blocks (emb i) ∧ blocks (emb i) ⊆ blocks b)
    (hcover : ∀ i : ι, blocks (cover i) ⊂ blocks (emb i))
    (hcellsSup : ∀ i : ι, blocks (emb i) \ blocks (cover i) ⊆ cells i)
    (horder : ∀ i j : ι, blocks (emb i) ⊆ blocks (emb j) → i ≤ j)
    (hcellSub : ∀ i : ι, cells i ⊆ blocks (emb i))
    (hMember : B.InvScopeSat (⋃ i ∈ C, blocks (emb i)) X0) :
    B.InvScopeSat (⋃ i ∈ C, cells i) X0
\end{leansrc}
\declhead{BernsteinLogic.QtslLatticeInv}{Behavior.qLatticeInv\_sound\_of\_blockInv}{\badge{BadgeGreen}{statement bound}}
\begin{leansrc}
/-- `Semiflow.qLatticeInv_sound_of_blockInv` at the behavioral floor. -/
theorem Behavior.qLatticeInv_sound_of_blockInv {B : Behavior X} {X0 : Set X}
    (hdc : IsLowerSet C)
    (hJoin : ∀ b : β, ∀ x ∈ blocks b,
      ∃ i : ι, x ∈ blocks (emb i) ∧ blocks (emb i) ⊆ blocks b)
    (hcover : ∀ i : ι, blocks (cover i) ⊂ blocks (emb i))
    (hcellsSup : ∀ i : ι, blocks (emb i) \ blocks (cover i) ⊆ cells i)
    (horder : ∀ i j : ι, blocks (emb i) ⊆ blocks (emb j) → i ≤ j)
    (hcellSub : ∀ i : ι, cells i ⊆ blocks (emb i))
    (hInv : B.Invariant (⋃ i ∈ C, blocks (emb i)))
    (hInit : X0 ⊆ ⋃ i ∈ C, cells i) :
    B.InvScopeSat (⋃ i ∈ C, cells i) X0
\end{leansrc}
\par\medskip
\textbf{The payload-level unpack, displayed whole.}\enspace
Recipe~\ref{rec:qli-filled} as a function: the wire payload reified
field-for-field against the register, the containment row transported
along the exactness identity, the target\textquotesingle s soundness
consumed literally in the factoring pair.

\declhead{BernsteinLogic.QtslElaborators}{LatticeInvPayload}{\badge{BadgeGreen}{statement bound}}
\begin{leansrc}
/-- **The QLatticeInv wire payload, Lean face** — fields mirror the
geolog register (section J's adjudication: "registers already staged;
the elaborator's domains are complete"): the family object with its
recomputed readout (`QLatticeInv/lattice`, as the leg-7 R1-shadow row
pack — no lattice structure minted, the leg-7 discipline), the
down-set (`QLatticeInv/downset`, with down-closedness the checker's
recomputed side condition), and the rule's own containment row
(`QLatticeInv/init_containment`). Membership of the down-set's
realization in the lattice is deliberately NOT a field — checker-side,
not consumed (soundness holds for the bare union, a fortiori with the
extra condition). -/
structure LatticeInvPayload (X ι β : Type*) [Finite ι] [PartialOrder ι] where
  /-- Member realizations (the family object's blocks). -/
  blocks : β → Set X
  /-- Join-irreducible embedding over the cell index. -/
  emb : ι → β
  /-- Lower covers. -/
  cover : ι → β
  /-- Cells `T(b) = cl(b ∖ b⁻)`, as realized sets. -/
  cells : ι → Set X
  /-- The down-set `C` (wire field `downset`). -/
  downset : Set ι
  /-- The initial region `X₀`. -/
  initRegion : Set X
  /-- Down-closedness in the recomputed cell order (R1 recompute). -/
  hdc : IsLowerSet downset
  /-- Birkhoff join-decomposition, realization form (R1 shadow). -/
  hJoin : ∀ b : β, ∀ x ∈ blocks b,
    ∃ i : ι, x ∈ blocks (emb i) ∧ blocks (emb i) ⊆ blocks b
  /-- Strict lower covers (R1 shadow). -/
  hcover : ∀ i : ι, blocks (cover i) ⊂ blocks (emb i)
  /-- Cell ⊇ member-minus-cover (R1 shadow). -/
  hcellsSup : ∀ i : ι, blocks (emb i) \ blocks (cover i) ⊆ cells i
  /-- Order reflects realization inclusion (R1 shadow). -/
  horder : ∀ i j : ι, blocks (emb i) ⊆ blocks (emb j) → i ≤ j
  /-- Cell ⊆ its member's realization (the row this alias adds). -/
  hcellSub : ∀ i : ι, cells i ⊆ blocks (emb i)
  /-- The rule's own containment row (wire field `init_containment`),
  stated on the alias's displayed region. -/
  hInit : initRegion ⊆ ⋃ i ∈ downset, cells i
\end{leansrc}
\declhead{BernsteinLogic.QtslElaborators}{BlockInvPayload}{\badge{BadgeGreen}{statement bound}}
\begin{leansrc}
/-- **The QBlockInv payload, Lean face** (`rec:qli-transport`'s
codomain): the block region and the containment row. The facet rows are
a CITATION (judgement-level cargo, the `qBasin_sound` precedent) — they
enter the factoring theorems as the cited node's invariance, never as
payload data. -/
structure BlockInvPayload (X : Type*) where
  /-- The block `N_C` — here its realization, the bare union (the
  join realizes as the union in the union-closed family object; that
  the union is a lattice MEMBER is the checker's recomputed side
  condition, deliberately not consumed). -/
  block : Set X
  /-- The initial region `X₀`, transported verbatim. -/
  initRegion : Set X
  /-- `init_containment`, rewritten along `qLatticeInv_exact`. -/
  hInit : initRegion ⊆ block
\end{leansrc}
\declhead{BernsteinLogic.QtslElaborators}{Semiflow.qLatticeInv\_sound\_elab}{\badge{BadgeGreen}{statement bound}}
\begin{leansrc}
/-- **sound_A = sound_T ∘ unpack, semiflow stratum** — the target's
soundness lemma (`qBlockInv_sound`, literally sound_T) applied to the
UNPACKED payload, concluding the alias's displayed judgement. The cited
member's facet-row chain enters as `hInv` (judgement-level cargo). -/
theorem Semiflow.qLatticeInv_sound_elab {φ : Semiflow X}
    (p : LatticeInvPayload X ι β) (hInv : φ.Invariant p.unpack.block) :
    φ.InvScopeSat (⋃ i ∈ p.downset, p.cells i) p.initRegion
\end{leansrc}
\declhead{BernsteinLogic.QtslElaborators}{Behavior.qLatticeInv\_sound\_elab}{\badge{BadgeGreen}{statement bound}}
\begin{leansrc}
/-- **sound_A = sound_T ∘ unpack, behavioural floor** — the identity is
dynamics-free, so the unpack transports at every semantics for free. -/
theorem Behavior.qLatticeInv_sound_elab {B : Behavior X}
    (p : LatticeInvPayload X ι β) (hInv : B.Invariant p.unpack.block) :
    B.InvScopeSat (⋃ i ∈ p.downset, p.cells i) p.initRegion
\end{leansrc}
\declhead{BernsteinLogic.QtslElaborators}{LatticeInvPayload.unpack}{\badge{BadgeGreen}{definition bound}}
\begin{leansrc}
/-- **The verified unpack** — `rec:qli-filled` as a function, the first
of the five alias artefacts (`unpack : Step_A → Step_T`): block :=
the down-set's realization `N_C`, facets cited (not payload),
`init_containment` transported along the exactness identity
`qLatticeInv_exact` (`⋃↓C = N_C` as regions — the alias's conclusion is
the target's on the nose). -/
def LatticeInvPayload.unpack (p : LatticeInvPayload X ι β) : BlockInvPayload X where
  block := ⋃ i ∈ p.downset, p.blocks (p.emb i)
  initRegion := p.initRegion
  hInit := by
    rw [← qLatticeInv_exact p.hdc p.hJoin p.hcover p.hcellsSup p.horder p.hcellSub]
    exact p.hInit
\end{leansrc}
\ruleblock{Correspondence (Steele display $\leftrightarrow$ Lean)}
\badge{BadgeGreen}{reviewer-certified against the displayed texts}
\sourcenote{The dictionary below is a reviewer-certified statement of how the mechanized statements above interpret the displayed rule: display metavariable $\mapsto$ Lean binder, each display hypothesis $\mapsto$ the premise or side condition carrying it, the display conclusion $\mapsto$ the goal --- with any deliberate difference in width between display and mechanization named outright. Certification is recorded as a hash triple over (the displayed cell, the extracted statement texts, this dictionary\textquotesingle s own text); if any of the three changes, this block turns amber until a reviewer re-certifies. The correspondence claim itself remains editorial --- reviewer-grade, not machine-checked; its machine-checked successor types each soundness statement against the rule as a structured object, and supersedes this block per rule as it lands.}
The displayed judgment $\mathcal{M} \vDash
\mathsf{QInvScope}(\bigcup\downarrow C, X_0)$ $\mapsto$
\texttt{InvScopeSat (⋃ i ∈ C, cells i) X0}, at BOTH semantic strata
(\texttt{Semiflow.}/\texttt{Behavior.qLatticeInv\_sound} --- the
exactness identity is dynamics-free, so the behavioural transport is
verbatim). The alias\textquotesingle s entire content is
\texttt{qLatticeInv\_exact}: for down-closed \texttt{C} (\texttt{hdc
: IsLowerSet C} --- the R1 down-closedness recompute as a binder),
the union of $C$\textquotesingle s cells EQUALS the union of its
join-irreducible realizations --- the member $N_C$\textquotesingle s
realization exactly where the checker\textquotesingle s membership
recompute holds (join realizes as union; the bound statement is on the
bare union and consumes no distributivity) --- so the respelling is
exact, not merely one-sided.
\begin{itemize}
\item the E2 lattice $L$ $\mapsto$ the Birkhoff row pack, abstract:
  \texttt{hJoin} (every block point lands in a join-irreducible below
  it), \texttt{hcover} (strict lower covers), \texttt{hcellsSup}
  (cell $\supseteq$ member minus cover), \texttt{horder} (order
  reflects realization inclusion) --- byte-identical to the Morse
  discharge\textquotesingle s eventual-member interface: no lattice
  structure is minted, the rows ARE the interface --- plus the one row
  new to this alias,
  \texttt{hcellSub} (cell $\subseteq$ member: the readout closure
  half\textquotesingle s set-level shadow).
\item the cited \rulename{QBlockInv} premise $\mapsto$
  \texttt{hMember : InvScopeSat (⋃ i ∈ C, blocks (emb i)) X0} ---
  judgment-level cargo on the member\textquotesingle s realization
  (the \texttt{qBasin\_sound} precedent: facet rows live in the cited
  node, never here); \texttt{\_of\_blockInv} composes the target in
  end-to-end (its \texttt{hInv}/\texttt{hInit} pair is
  \texttt{qBlockInv\_sound}\textquotesingle s kit, the containment
  row transported across the identity).
\item \texttt{init\_containment} $\mapsto$ \texttt{hInit : X0 ⊆ ⋃ i ∈
  C, cells i}, stated on the alias\textquotesingle s displayed region.
\item \texttt{member\_subset\_cells\_le} exposes the
  page\textquotesingle s ``cells tile members'' fact at the principal
  down-set --- bound because the proof section names it.
\end{itemize}
Width note: the module is topology-free and dynamics-free except
through the cited judgment --- \texttt{cells}/\texttt{blocks}
arbitrary set families over an abstract ordered index; the grid
realization of the lattice is R1 territory, strictly wider in the
safe direction. Module row bound to an attested build (the receipt
provenance is recorded in the bindings).

\ruleblock{Lean verifier receipts}
\sourcenote{module \texttt{BernsteinLogic.QtslLatticeInv}; bound build \texttt{467dc57e1c30\ldots} (bound from receipt 98755f4d8dbefb10, 2026-07-21T04:48:40Z). A receipt records that named checking layers accepted a build: (module, compiled-artifact SHA-256, environment SHA-256, checker layers), minted when the checking gate ran; the store receipts attesting the bound build are listed with the layers each run credited. Nothing here asks the reader to trust the store: every layer is an independently runnable program over the quoted sources --- rebuild the module and re-run the checkers to reproduce the verdict. Receipts attesting superseded builds credit no layer and are omitted from this rendering (the count of omissions, if any, is noted below).}
\begin{receipttable}
\receiptrow{98755f4d8dbefb10}{2026-07-21T04:48:40Z}{build+decorrelated}{\badge{BadgeGreen}{current}}
\receiptrow{210205ac3ea031f4}{2026-07-21T05:17:21Z}{build+decorrelated}{\badge{BadgeGreen}{current}}
\receiptrow{9a870584fbb58e05}{2026-07-21T05:28:41Z}{build+replay+hardened+decorrelated}{\badge{BadgeGreen}{current}}
\end{receipttable}
\sourcenote{3 further receipt(s) in the store attest superseded builds of this module and are omitted from this rendering.}
\sourcenote{module \texttt{BernsteinLogic.QtslElaborators}; bound build \texttt{0205ec213b91\ldots} (hashed from the built olean at bind time). A receipt records that named checking layers accepted a build: (module, compiled-artifact SHA-256, environment SHA-256, checker layers), minted when the checking gate ran; the store receipts attesting the bound build are listed with the layers each run credited. Nothing here asks the reader to trust the store: every layer is an independently runnable program over the quoted sources --- rebuild the module and re-run the checkers to reproduce the verdict. Receipts attesting superseded builds credit no layer and are omitted from this rendering (the count of omissions, if any, is noted below).}
\begin{receipttable}
\receiptrow{a237f4cff27cb15a}{2026-07-22T10:06:47Z}{build+decorrelated}{\badge{BadgeGreen}{current}}
\end{receipttable}
\layerline{\layeritem{build}{lake build (elaboration)}{\badge{BadgeGreen}{current}}\layeritem{replay}{leanchecker (official kernel replay)}{\badge{BadgeGreen}{current}}\layeritem{hardened}{safe\_verify (hardened screen)}{\badge{BadgeGreen}{current}}\layeritem{decorrelated}{lean4export + nanoda (decorrelated kernel)}{\badge{BadgeGreen}{current}}}
\sourcenote{\badge{BadgeGray}{comparator: none} no fixture in the three-implementation comparator harness (compiled / kernel-interpreted / reference-interpreted) exercises this rule.}
\ruleblock{HOL4 statement}
\absentbox{NO HOL STATEMENT YET}{Absent at the judgment stratum: no judgment-stratum definition exists
in the ported base-kernel checker scripts (verified against the
port\textquotesingle s definition inventory), so the satisfaction-level
witnesses --- now mechanized on the Lean side (module
\texttt{QtslLatticeInv}, the entry\textquotesingle s bound statements)
--- have no HOL twin to mirror them. The target rule\textquotesingle s
rows (facet signs, containment) ARE covered by the ported checker.
}
\ruleblock{HOL verifier receipts}
\absentbox{NONE}{no HOL verification receipt exists for this rule.}
\ruleblock{Proof}

\sourcenote{Genre: alias rule (tier 3) --- a respelling template. An
alias owes the book five artefacts: the verified 1$\to$1 unpack
function, the filled-in application of the target rule, the
payload-transport table, the exact algebraic identity that makes the
respelling correct, and the soundness corollary through the target
rule's lemma. All five are present --- the filled-in application, the
transport table, the exactness identity and the soundness corollary
mechanized and bound in the entry
(\texttt{qLatticeInv\_exact}, Proposition~\ref{prop:qli-exact}; the
\texttt{qLatticeInv\_sound} family, Proposition~\ref{prop:qli-sound},
at the semiflow \emph{and} behavioural strata), and --- now that the
extension certificate schema exists --- the verified payload-level
unpack function itself, \texttt{LatticeInvPayload.unpack} with its
factoring pair (the remark below names the bound functions). (Tier
flag, carried honestly: the
unpack cites facet rows through the checker's recomputed lattice
readout --- if registry-backed citation is ruled heavier than the alias
criterion allows, the rule reclassifies as admissible with all
machinery unchanged.)}

\begin{bremark}[the payload-level elaborator, mechanized]
\label{rem:qli-elab}
The \emph{semantic} content of this page is bound in the entry:
\texttt{qLatticeInv\_exact} is Proposition~\ref{prop:qli-exact}
(over the discharge chain's own Birkhoff rows --- no lattice structure
is minted), and the \texttt{qLatticeInv\_sound}
family is Proposition~\ref{prop:qli-sound} in both its
through-the-target and end-to-end forms, at both semantic strata. The
payload-level elaborator --- previously this page's one absent
artefact, held back until the extension certificate schema existed ---
is bound with it, as with the sibling extension rules:
\texttt{LatticeInvPayload} (module
\texttt{BernsteinLogic.QtslElaborators}) reifies the rule's wire
payload --- family object as the leg-7 row pack, down-set, containment
row --- and \texttt{LatticeInvPayload.unpack} takes it to the single
\rulename{QBlockInv} payload of Recipe~\ref{rec:qli-filled}
(\texttt{BlockInvPayload}: block $:=$ the down-set's realization, the
containment row transported along Proposition~\ref{prop:qli-exact};
facets are a citation, entering the factoring theorems as the cited
node's invariance, never as payload data). The factoring pair
\texttt{Semiflow.qLatticeInv\_sound\_elab} /
\texttt{Behavior.qLatticeInv\_sound\_elab} is
$\mathrm{sound}_A = \mathrm{sound}_T \circ \mathrm{unpack}$ with
\texttt{qBlockInv\_sound} as the literal $\mathrm{sound}_T$, applied to
the unpacked payload. The checker-side ingredients (membership of the
down-set's realization in the lattice, down-closedness of the cell-set)
are recomputed side conditions, already specified with the rule; the
mechanization deliberately consumes the down-closedness and leaves the
membership fact checker-side (with down-closedness consumed and the
cited member judgment supplied, soundness holds for the bare cell
union, a fortiori with the extra condition).
\end{bremark}

\begin{brecipe}[the alias template]\label{rec:qli-template}
An alias rule $A$ is warranted by a verified function
$\mathrm{unpack} : \mathrm{Step}_A \to \mathrm{Step}_T$ into the
payload of a \emph{single} target node $T$, such that the $A$-node and
its image conclude the same judgment from the same premises: one node
$\to$ one node, payload transported, nothing derived. Soundness is then
a corollary, $\mathrm{sound}_A = \mathrm{sound}_T \circ
\mathrm{unpack}$.
\end{brecipe}

\textbf{The rule.}\enspace The certificate supplies the family object
(the block lattice of the \rulename{QMorseDischarge} section, with its
recomputed readout: join-irreducible members $b$, lower covers $b^-$,
cells $T(b) := \mathrm{cl}(b \smallsetminus b^-)$, the induced cell
order), a cell-set $C$ of join-irreducibles that is \emph{down-closed}
in the cell order, and one containment row
$X_0 \subseteq \bigcup\downarrow C$. It concludes
$\mathsf{QInvScope}\big(\bigcup\downarrow C,\ X_0\big)$, writing
$\bigcup\downarrow C$ for the union of the cells of $C$.

\begin{brecipe}[the filled-in application --- \rulename{QLatticeInv}
rewritten \emph{into} the \rulename{QBlockInv} instance]
\label{rec:qli-filled}
Let $N_C := \bigvee_{b \in C} b$ --- the lattice join of the down-set's
members, itself a member of the family object by join-closure. The
\rulename{QLatticeInv} node rewrites to the single \rulename{QBlockInv}
application
\begin{mathpar}
\inferrule*[Right=\rulename{QBlockInv}]{
  N_C \ \text{(the member: block data from the family object)} \\
  \bigl\{\vdash \text{facet-inward sign rows of } N_C\bigr\}
  \ \text{(cited from the family object, byte-identical)} \\
  \vdash X_0 \subseteq N_C \ \text{(the rule's own containment row,
  rewritten along Proposition~\ref{prop:qli-exact})}
}{
  \mathcal{M} \vDash \mathsf{QInvScope}(N_C,\ X_0)
}
\end{mathpar}
and the conclusion is the alias's conclusion \emph{on the nose} by
Proposition~\ref{prop:qli-exact}: $N_C = \bigcup\downarrow C$ as
regions. Every internal facet of $N_C$ already carries its sign row in
the family object --- membership of $N_C$ in the lattice is exactly
the checker's recomputed side condition, and no row is re-proven.
\end{brecipe}

\begin{bprop}[the exactness identity]\label{prop:qli-exact}
For a down-closed set $C$ of join-irreducibles,
\[
\bigcup_{b \in C} T(b) \;=\; \bigvee_{b \in C} b
\qquad \text{(equality of realizations).}
\]
\end{bprop}
\begin{proof}
Two readout facts, one lattice fact. \emph{(cells tile members)} Every
member $N$ of the lattice is the union of the cells of the
join-irreducibles below it, $N = \bigcup \{T(b) : b \in \mathsf{J},\
b \subseteq N\}$: downward induction along lower covers --- $T(b)$ is
by definition what $b$ adds over $b^-$, so peeling cells descends
through the lattice and exhausts $N$. \emph{(Birkhoff exactness)} In a
finite \emph{distributive} lattice (a lattice of regions closed under
union and intersection --- a ring of sets --- qualifies) the map
$N \mapsto \{b \in \mathsf{J} : b
\subseteq N\}$ is an isomorphism onto down-closed sets of
join-irreducibles \cite{birkhoff1937}; for down-closed
$C$ the join $N_C$ then satisfies $\{b \in \mathsf{J} : b \subseteq N_C\} =
C$. Distributivity --- join-primality of the irreducibles --- is what
forbids the join from absorbing new join-irreducibles, and it is not
decoration: in the diamond lattice $M_3$, two of the three atoms form
a down-closed set of join-irreducibles whose join absorbs the third. Substituting the second fact into the
first: $N_C = \bigcup_{b \in C} T(b)$. (Where each half lives: the
tiling is the readout's own construction, recomputed by the checker
per certificate; the isomorphism is the one classical fact, and
down-closedness of $C$ is the recomputed side condition that makes the
identity exact rather than merely $\supseteq$.) \emph{Mechanization
note, the honest disanalogy, adjudicated against the descent:}
\texttt{qLatticeInv\_exact} binds the identity in \emph{union form} ---
$\bigcup_{b \in C} T(b) = \bigcup_{b \in C} b$ as realizations, the
right side the bare union over the down-set --- and proves it from the
row pack alone: \emph{no distributivity hypothesis is consumed}, and
none is encoded, since the rows are satisfiable in absorption-positive
families. The route never forms the join: each point of the union
arrives with a witnessing member $b \in C$; the well-founded descent
through lower covers (the discharge chain's
\texttt{exists\_cellMem\_subset}) runs \emph{inside that one member},
landing the point in the cell of a join-irreducible $b' \subseteq b$;
and the pin-step is two rows --- the induced-order binding
\texttt{horder} turns $b' \subseteq b$ into $b' \leq b$ in the cell
order, and down-closedness closes $b' \in C$. Join-primality forbids
absorption below the \emph{join}; the descent never visits the join,
so the prohibition is sidestepped, not re-proven. What distributivity
carries is therefore only the display seam --- reading the bare union
as the realization of the member $N_C = \bigvee C$ --- and the seam
consumes less than the full hypothesis: it needs exactly
join-realizes-as-union, which is the checker's recomputed membership
of the down-set's realization (a member that is a union of members is
their least upper bound), a consequence of the ring-of-sets structure
strictly weaker than distributivity (union-closed absorption-positive
families exist: there the identity holds in \emph{both} forms while
the classifying map still fails). The $M_3$ witness above, realized as
sets --- atom realizations $\{1\}, \{2\}, \{3\}$, top $\{1,2,3\}$,
$C$ the first two atoms --- separates the strata cleanly: every row
holds, the union form holds ($\{1,2\} = \{1,2\}$), the join form
fails ($\{1,2\} \neq \{1,2,3\}$), and it is the membership recompute
that refuses. Nothing mechanized rests on the classical citation; it
remains as warrant genealogy for this proof's route.
\end{proof}

\begin{brecipe}[payload transport --- field for field]
\label{rec:qli-transport}
\begin{center}\small
\begin{tabular}{@{}ll@{}}
\textbf{\rulename{QLatticeInv} payload} & \textbf{\rulename{QBlockInv} payload}\\
\midrule
family object $+$ down-set $C$ & block $N_C = \bigvee C$ (member, by join-closure)\\
\emph{(no own rows)} & \texttt{facets} $:=$ $N_C$'s rows, cited from the family object\\
\texttt{init\_containment} & \texttt{init\_containment} (rewritten along Proposition~\ref{prop:qli-exact}: $\bigcup\downarrow C = N_C$)\\
conclusion region $\bigcup\downarrow C$ & conclusion region $N_C$ (Proposition~\ref{prop:qli-exact})\\
\end{tabular}
\end{center}
The checker's recomputed side conditions at the alias: $C$ down-closed
in the recomputed cell order, and $N_C$ a member whose facet rows are
all present --- both facts about the readout, not about the dynamics.
\end{brecipe}

\begin{bprop}[soundness corollary, through the target]
\label{prop:qli-sound}
If the family object's rows are valid and
$X_0 \subseteq \bigcup\downarrow C$, then --- under the
\rulename{QBlockInv} section's standing hypotheses (its conditional
ledger, applied to the member $N_C$) ---
$\mathsf{QInvScope}\big(\bigcup\downarrow C,\ X_0\big)$ holds.
\end{bprop}
\begin{proof}
The member $N_C$ is a strictly facet-inward compact block of the
family object, so the \rulename{QBlockInv} section's theorem chain
applies to it verbatim (its rows \emph{are} rows of the family
object); the conclusion transports along
Proposition~\ref{prop:qli-exact}. Mechanized, this proposition is one
composition, $\mathrm{sound}_{\textsc{QLatticeInv}} =
\mathrm{sound}_{\textsc{QBlockInv}} \circ \mathrm{unpack}$ --- there
is no other content, which is the alias genre working as intended: the
through-the-target form (\texttt{Semiflow.qLatticeInv\_sound} /
\texttt{Behavior.qLatticeInv\_sound}) takes the cited node's
$\mathsf{QInvScope}$ conclusion as judgment-level cargo and rewrites
along the identity; the end-to-end form (\texttt{\_of\_blockInv})
composes through \texttt{qBlockInv\_sound} verbatim. Both strata are
covered because the identity is dynamics-free --- the alias transports
at every semantics for free.
\end{proof}

\begin{bremark}[why the rule earns its keep]
Down-closed cell-sets are the invariant enclosures the Morse readout
offers for free: once an orbit enters a member it never leaves, so the
sequence of cells it visits is non-increasing in the cell order, and
the down-set of any initial cell-set is forward-invariant and encloses
the reachable set. The alias turns each such enclosure into a one-node
citation --- the member's rows already exist in the family object ---
where a direct proof would re-derive facet rows the certificate
already carries. Downstream, the minted $\mathsf{QInvScope}$ judgment
feeds \rulename{QBasin}'s safety factory and scope-restriction of
reach obligations, which is where the row-count payoff compounds.
\end{bremark}

\begin{anchors}
target soundness & \texttt{qBlockInv\_sound} (QtslSoundness.lean) and the \rulename{QBlockInv} section's discharge chain\\
Proposition \ref{prop:qli-exact} & \texttt{qLatticeInv\_exact} (QtslLatticeInv.lean, bound in the entry; union form, no distributivity consumed) via the \texttt{exists\_cellMem\_subset} descent; Birkhoff representation \cite{birkhoff1937} as warrant genealogy (nothing mechanized rests on it)\\
``cells tile members'' & \texttt{member\_subset\_cells\_le} (principal down-set instance)\\
alias-level corollary & \texttt{Semiflow.qLatticeInv\_sound} $+$ \texttt{\_of\_blockInv} and the behavioural-floor pair (bound in the entry)\\
payload-level elaborator (\autoref{rem:qli-elab}) & \texttt{LatticeInvPayload.unpack} + \texttt{Semiflow.qLatticeInv\_sound\_elab} + \texttt{Behavior.qLatticeInv\_sound\_elab} (module \texttt{BernsteinLogic.QtslElaborators}, bound in the entry)\\
side conditions & down-closedness: consumed (\texttt{IsLowerSet} $+$ \texttt{horder}); membership of $N_C$ in the lattice: checker-side, deliberately not consumed\\
\end{anchors}


\rulesection{74}{QBasin}{core rule}
\label{blcatrule-QBasin}
\ruleblock{The rule (Steele metanotation)}
\sourcenote{as realized in the mechanized development; source row: \texttt{\texttt{QTSL-rules.toml}} --- \emph{QBasin}. The metanotation statement is typeset by hand for notational quality; its mechanized ground truth is the Lean statement below, extracted from the sources at every build of this book and hash-bound. The rule \emph{roster} is machine-checked (every mechanized rule tag round-trips to exactly one catalogue row); the statement \emph{text} is editorial, checked by eye against the extracted Lean beside it.}
\begin{mathpar}
\inferrule*{
  \mathcal{M} \vDash \mathsf{QInvScope}(S, X_0) \\
  \bigl\{\vdash S \cap [\![\beta_j]\!] = \emptyset\bigr\}_{j} \ \text{(\texttt{avoid\_disjoint} rows)}
}{
  \mathcal{M},\ x_0 \in X_0 \vDash \mathsf{SProbGE}\bigl(1,\ \mathsf{PAlways}(\mathsf{SNot}\,\beta_j)\bigr)
}
\end{mathpar}

\ruleflavor{the invariance premise is deliberately the \\textsf{QInvScope} judgment and not the Morse-graph judgment: basin avoidance consumes invariance of $S$ from $X_0$ alone, and asking the stronger judgment would couple the rule to machinery its soundness never uses.}
\par\medskip
\rulename{QBasin} is the campaign rule: safety atoms by citation. From
one invariance-scope judgment and one disjointness fact per avoid
region, it concludes the safety-shaped region judgment
$\mathsf{SProbGE}\bigl(1,\ \mathsf{PAlways}(\mathsf{SNot}\,\beta)\bigr)$
(surface syntax $\mathsf{P}_{\ge 1}[\mathsf{G}\,\neg\beta]$) on the
scope\textquotesingle s init region --- the deterministic reading, with
no probabilistic machinery: for the deterministic model class the
certified claim is universal avoidance, and no supermartingale
certificate is synthesized or needed. The economics are the point: the
invariance judgment is proved once (typically by \rulename{QBlockInv}
or \rulename{QLatticeInv}, whose facet rows already live in the family
object), and each additional avoid region costs exactly its own
disjointness rows and nothing else --- the invariance side is cited,
never re-certified, even when the regions are taken jointly (the batch
theorem).

What the checker recomputes: the conclusion tie (the cited
judgment\textquotesingle s region and init fields equal the
rule\textquotesingle s own), plus the per-box disjointness sign rows,
or the grid bookkeeping when the avoid region is cell-aligned. What
would break without each hypothesis: without invariance, a trajectory
exits the scope and reaches the avoid set (the proof
section\textquotesingle s drift counterexample); without the
containment, failure at time zero; without disjointness, failure at
time zero inside the scope.

The premise is deliberately the invariance judgment and NOT the global
Morse judgment: recurrence structure alone implies no enclosure of the
init region --- an orbit may transit the avoid set en route to its
Morse set. Basin-minted and battery-minted safety atoms are
type-identical conclusions, so downstream citations consume either
producer indifferently; keying on provenance is forbidden by design.

This is the one Morse-family rule whose mechanization is
unconditional: the cited theorems are proved outright, with no
interface hypotheses.

\cataloguenote{The probability-one costume ($\mathsf{SProbGE}(1, \cdot)$) matters for
composition: the
conclusion lands in the same judgment form the probabilistic batteries
mint, so tower citations, boolean combination at the root, and the
assume--guarantee spine consume deterministic and stochastic safety
atoms through one interface.
}
\ruleblock{Lean statement}
\declhead{BernsteinLogic.QtslSoundness}{qBasin\_sound}{\badge{BadgeGreen}{statement bound}}
\begin{leansrc}
/-- QBasin soundness (roster: "3 lines against the Morse note's lemma
corpus"; axioms note §3.5) — FULLY PROVED, no interface facts. Premises:
the cited `QInvScope` judgment (the E1 edge `QBasin/inv`) and the
avoid-disjointness fact (`QBasin/avoid_disjoint`'s rows). Conclusion:
the safety-shaped `QRegionSat` on the scope's init region. Note the
invariance premise arrives as judgment-level cargo — the certificate-
level fact `QInvScope` carries — which is why no dynamics beyond the
semiflow algebra is consumed here (the roster's §4.6 correction: the
premise is NOT the `QMorse` judgment). -/
theorem qBasin_sound {φ : Semiflow X} {S X0 B : Set X}
    (hscope : φ.InvScopeSat S X0) (hAvoid : S ∩ B = ∅) :
    φ.SafetySat X0 B
\end{leansrc}
\declhead{BernsteinLogic.QtslSoundness}{qBasin\_batch}{\badge{BadgeGreen}{statement bound}}
\begin{leansrc}
/-- The batch-factory payoff, mechanized ("one family object, many
G-judgments, no new rows" — zero new rows per additional `β` after the
first containment): one invariant scope
discharges every avoid-region at the price of its own disjointness fact
alone — even jointly. -/
theorem qBasin_batch {φ : Semiflow X} {S X0 : Set X} {J : Type*}
    {B : J → Set X} (hscope : φ.InvScopeSat S X0)
    (hAvoid : ∀ j, S ∩ B j = ∅) :
    φ.SafetySat X0 (⋃ j, B j)
\end{leansrc}
\ruleblock{Correspondence (Steele display $\leftrightarrow$ Lean)}
\badge{BadgeAmber}{a bound Lean statement CHANGED since certification --- re-review and re-certify}
\sourcenote{The dictionary below is a reviewer-certified statement of how the mechanized statements above interpret the displayed rule: display metavariable $\mapsto$ Lean binder, each display hypothesis $\mapsto$ the premise or side condition carrying it, the display conclusion $\mapsto$ the goal --- with any deliberate difference in width between display and mechanization named outright. Certification is recorded as a hash triple over (the displayed cell, the extracted statement texts, this dictionary\textquotesingle s own text); if any of the three changes, this block turns amber until a reviewer re-certifies. The correspondence claim itself remains editorial --- reviewer-grade, not machine-checked; its machine-checked successor types each soundness statement against the rule as a structured object, and supersedes this block per rule as it lands.}
Both statements are read against the abstract model class: the
display\textquotesingle s model $\mathcal{M}$ is any forward semiflow
$\varphi$ (\texttt{Semiflow X}) on a bare state type; the v1 model
class realizes it, and nothing in either statement consumes the
realization.
\begin{itemize}
\item $S$, $X_0$ (the scope\textquotesingle s region and init region)
  $\mapsto$ the binders \texttt{S X0 : Set X}; the avoid
  predicate\textquotesingle s extent $[\![\beta_j]\!]$ $\mapsto$
  \texttt{B : Set X} in \texttt{qBasin\_sound}, and the indexed family
  \texttt{B : J $\to$ Set X} in \texttt{qBasin\_batch} (the
  display\textquotesingle s $j$-indexed premise row, taken jointly).
\item premise $\mathrm{inv} : \mathsf{QInvScope}(S, X_0)$ $\mapsto$
  \texttt{hscope} $: \varphi.\texttt{InvScopeSat}\ S\ X_0$ ---
  definitionally the conjunction of forward invariance of $S$ and
  $X_0 \subseteq S$, the judgment\textquotesingle s two clauses.
\item each \texttt{avoid\_disjoint} row $\vdash S \cap [\![\beta_j]\!]
  = \emptyset$ $\mapsto$ \texttt{hAvoid} $: S \cap B = \emptyset$
  (batch: $\forall j,\ S \cap B\,j = \emptyset$).
\item conclusion $\mathcal{M},\ x_0 \in X_0 \vDash
  \mathsf{SProbGE}\bigl(1,\ \mathsf{PAlways}(\mathsf{SNot}\,\beta_j)\bigr)$
  $\mapsto$ the goal
  $\varphi.\texttt{SafetySat}\ X_0\ B$: every trajectory from $X_0$
  avoids $B$ at every forward time --- the deterministic
  probability-one reading of the displayed judgment.
  \texttt{qBasin\_batch} concludes
  $\varphi.\texttt{SafetySat}\ X_0\ (\bigcup_j B\,j)$, which yields the
  displayed per-$j$ family of conclusions by monotonicity.
\end{itemize}
Width note: the display\textquotesingle s E1 annotation, glossed by
the flavor text (the premise is
the invariance judgment, \emph{not} the Morse judgment), is enforced by
type rather than by hypothesis --- \texttt{qBasin\_sound} consumes an
\texttt{InvScopeSat} fact and nothing of Morse shape fits that premise
slot. No hypothesis of either statement is unused, and neither
statement is wider than the displayed rule.

\ruleblock{Lean verifier receipts}
\sourcenote{module \texttt{BernsteinLogic.QtslSoundness}; bound build \texttt{77005991d61d\ldots} (bound from receipt 8130a7a7a679cef4, 2026-08-04T01:16:06Z). A receipt records that named checking layers accepted a build: (module, compiled-artifact SHA-256, environment SHA-256, checker layers), minted when the checking gate ran; the store receipts attesting the bound build are listed with the layers each run credited. Nothing here asks the reader to trust the store: every layer is an independently runnable program over the quoted sources --- rebuild the module and re-run the checkers to reproduce the verdict. Receipts attesting superseded builds credit no layer and are omitted from this rendering (the count of omissions, if any, is noted below).}
\begin{receipttable}
\receiptrow{8130a7a7a679cef4}{2026-08-04T01:16:06Z}{build+decorrelated}{\badge{BadgeGreen}{current}}
\receiptrow{9eb3706354082144}{2026-08-04T02:23:12Z}{build+decorrelated}{\badge{BadgeGreen}{current}}
\end{receipttable}
\sourcenote{5 further receipt(s) in the store attest superseded builds of this module and are omitted from this rendering.}
\layerline{\layeritem{build}{lake build (elaboration)}{\badge{BadgeGreen}{current}}\layeritem{replay}{leanchecker (official kernel replay)}{\badge{BadgeAmber}{stale only}}\layeritem{hardened}{safe\_verify (hardened screen)}{\badge{BadgeAmber}{stale only}}\layeritem{decorrelated}{lean4export + nanoda (decorrelated kernel)}{\badge{BadgeGreen}{current}}}
\sourcenote{\badge{BadgeGray}{comparator: none} no fixture in the three-implementation comparator harness (compiled / kernel-interpreted / reference-interpreted) exercises this rule.}
\ruleblock{HOL4 statement}
\absentbox{NO HOL STATEMENT YET}{Absent: the HOL4 development\textquotesingle s object is the
base-kernel checker cone, and the satisfaction-judgment stratum this
rule concludes in has no HOL twin --- no judgment-stratum definition
exists in the ported scripts (verified against the
port\textquotesingle s definition inventory). The disjointness rows
this rule consumes ARE covered by the ported checker; the three-line
judgment-stratum soundness lemma is Lean-only today.
}
\ruleblock{HOL verifier receipts}
\absentbox{NONE}{no HOL verification receipt exists for this rule.}
\ruleblock{Proof}

\sourcenote{Genre: core rule (tier 1) --- an ordinary mathematical proof at
Putnam-exam-solution detail. Every hypothesis of the rule is used below
and remarked on; the formal-anchor table maps each step to the Lean
declaration that carries it. This rule concludes a judgment one stratum
above the sequent calculus: premises typed $\vdash$ are ordinary sequents
the base kernel replays, and the rule's own soundness lives at the
judgment stratum, over a dynamical model. Uniquely among the Morse-family
rules in this book, its mechanization is \emph{unconditional}: the cited
theorems are proved outright, with no interface hypotheses.}

\textbf{Setting and claim.}\enspace Fix a \emph{forward semiflow} $\varphi$
on a state space $X$: a map $\Phi : \R_{\geq 0} \times X \to X$, written
$\Phi_t$, with $\Phi_0 = \mathrm{id}$ and $\Phi_{s+t} = \Phi_s \fatsemi
\Phi_t$ (the abstract shadow of the model class: a deterministic
polynomial-drift flow on its chart domains; the two semiflow laws are
model-frame data --- the proofs below consume only the rule premises
(a)--(c)). Regions denote subsets of
$X$. Two judgment forms participate:

\emph{(invariance scope)}\enspace $\mathsf{QInvScope}(S, X_0)$ asserts the
conjunction: $S$ is forward-invariant ($x \in S \Rightarrow \Phi_t x \in S$
for all $t \geq 0$) and $X_0 \subseteq S$.

\emph{(safety)}\enspace the region judgment
$\mathcal{M},\ x_0 \in X_0 \vDash
\mathsf{SProbGE}\bigl(1,\ \mathsf{PAlways}(\mathsf{SNot}\,\beta)\bigr)$
asserts, in the deterministic reading, that every point of $X_0$ avoids
$B := \llbracket\beta\rrbracket$ at every forward time:
$\forall x \in X_0,\ \forall t \geq 0,\ \Phi_t x \notin B$. The
probability-one former $\mathsf{SProbGE}(1, \cdot)$
is the degenerate probabilistic costume of this
universal statement --- no measure, no stochastic machinery, and no
supermartingale certificate participates anywhere in this section (the
model-class fence: deterministic flows only).

The rule's premises: the cited invariance judgment (an internal edge ---
the claim is consumed, not its derivation), and one disjointness fact
$S \cap \llbracket\beta_j\rrbracket = \emptyset$ per avoid region,
discharged by per-box sign rows of the base calculus (or by cell
bookkeeping the checker recomputes, when $\beta_j$ is aligned to the
certificate's grid).

\begin{bthm}[soundness of \rulename{QBasin}]\label{thm:qbasin-sound}
Let $S, X_0, B \subseteq X$ with
\emph{(a)} $S$ forward-invariant,
\emph{(b)} $X_0 \subseteq S$
(jointly: the cited $\mathsf{QInvScope}(S, X_0)$ judgment), and
\emph{(c)} $S \cap B = \emptyset$.
Then every $x \in X_0$ satisfies $\Phi_t x \notin B$ for all
$t \geq 0$.
\end{bthm}
\begin{proof}
Three sentences. By (b), $x \in S$; by (a), $\Phi_t x \in S$ for every
$t$. Were $\Phi_t x \in B$, it would lie in $S \cap B$, which (c)
declares empty. Hence $\Phi_t x \notin B$.
\end{proof}

\begin{bprop}[the batch form]\label{prop:qbasin-batch}
Let $\{B_j\}_{j \in J}$ be any family of avoid regions with, for each
$j$, the single fact $S \cap B_j = \emptyset$. Under (a) and (b) of
Theorem~\ref{thm:qbasin-sound}, every $x \in X_0$ avoids
$\bigcup_j B_j$ at every forward time --- equivalently, each $B_j$
separately, and all of them jointly.
\end{bprop}
\begin{proof}
A point of the union lies in some $B_j$; apply
Theorem~\ref{thm:qbasin-sound} to that $j$. Nothing about $J$ is used ---
the index family may be infinite, though a certificate of course carries
finitely many rows.
\end{proof}

\begin{bremark}[each hypothesis is load-bearing]\label{rem:qbasin-hyp}
Drop (a), keeping (b) and (c): on $X = \R$ with the unit-speed drift
$\Phi_t x = x + t$, take $S = X_0 = [0,1]$ and $B = [2,3]$. Then
$X_0 \subseteq S$ and $S \cap B = \emptyset$, yet the trajectory from
$x = 0$ enters $B$ at $t = 2$ --- invariance is the entire dynamical
content of the rule, and it arrives wholesale through the cited
judgment. Drop (b): a point $x \in X_0 \smallsetminus S$ may sit in $B$
already, violating the conclusion at $t = 0$ --- concretely, the
identity flow with $S = \{0\}$ and $X_0 = B = \{1\}$: (a) and (c) hold,
and the start state is the violation. Drop (c): a point of
$S \cap B$ reachable from $X_0$ (at worst a point of
$X_0 \cap B$ itself, again at $t = 0$) violates it --- concretely, the
identity flow with $X_0 = S = B = \{0\}$: (a) and (b) hold, and the
start state sits in the avoid set. The two set-theoretic
hypotheses fail at time zero when dropped; the dynamical one fails at
positive time --- which is why (a) is a judgment with its own producing
rules while (b) and (c) are single containment/disjointness rows.
\end{bremark}

\begin{bremark}[the premise is the invariance judgment, not the Morse
judgment]\label{rem:qbasin-premise}
A natural mis-design would take the global Morse judgment as the
premise: it is the stronger-looking claim. But
$\mathsf{Morse}(G, \rho)$ alone implies no particular invariant
enclosure of $X_0$ --- an orbit from $X_0$ may transit
$\llbracket\beta\rrbracket$ \emph{en route} to its Morse set
(concretely: $\dot{x} = 3 - x$ on $\R$, Morse structure the attractor
$\{3\}$, $X_0 = \{0\}$, $B = (1, 2)$ --- the orbit crosses $B$ on its
way to the attractor), so
recurrence structure without an enclosure proves no safety at all. The
sound premise is the certificate-level invariance fact, and the Morse
machinery contributes through the \emph{factory} instead: every lattice
member of the family object (and every down-closed cell-set, via
\rulename{QLatticeInv}) yields a $\mathsf{QInvScope}$ judgment whose
facet rows already exist in the family object. This rule then mints one
safety atom per avoid region at the price of that region\textquotesingle s
own disjointness rows and \emph{nothing else} --- the invariance side is
cited, never re-certified, so no invariance or facet row is ever added
however many avoid regions arrive --- which is the row-count payoff the
batch form
(Proposition~\ref{prop:qbasin-batch}) states as mathematics.
\end{bremark}

\begin{bremark}[what the checker recomputes]
At this rule the checker verifies: that the cited judgment's region and
init fields equal the rule's own scope data (the conclusion tie), and
the disjointness rows --- per-box, per-atom sign sequents of the base
calculus, or grid bookkeeping when $\beta_j$ is cell-aligned. Nothing
else: soundness is pure logic over the judgment's semantic content,
which is why the mechanized statement is three lines and
hypothesis-free. Basin-minted and battery-minted safety atoms are
\emph{type-identical} conclusions, deliberately: downstream citations
consume the judgment and may not key on its producer.
\end{bremark}

\begin{bremark}[relation to the literature]
Safety-by-invariance is the classical frame of set-invariance control
--- Nagumo's subtangentiality theorem \cite{nagumo1942} is the
ancestral invariance certificate, and Blanchini's survey
\cite{blanchini1999} the modern map of the field --- and of barrier
certificates \cite{prajnajadbabaie2004}, which in their usual packaging
synthesize one function
per specification, coupling the invariance argument to the safety
argument. The rule's point is the \emph{decoupling}: invariance is
certified once, as a judgment (by strictly inward facet rows, at
\rulename{QBlockInv}), and each safety claim is thereafter a citation
plus disjointness rows. The batch corollary is the certificate-side
face of the elementary observation that one invariant set separates its
contents from every disjoint avoid set at once.
\end{bremark}

\begin{anchors}
Theorem \ref{thm:qbasin-sound} & \texttt{qBasin\_sound} (QtslSoundness.lean); premise cargo \texttt{hscope}, \texttt{hAvoid}\\
Proposition \ref{prop:qbasin-batch} & \texttt{qBasin\_batch}\\
invariance-scope denotation & \texttt{Semiflow.InvScopeSat}\\
safety denotation ($\mathsf{SProbGE}(1, \cdot)$, deterministic reading) & \texttt{Semiflow.SafetySat}\\
disjointness rows & per-box \rulename{BernCone}-family sequents (base kernel); cell-aligned case = recomputed bookkeeping\\
\end{anchors}


\blcatsection{Existence Certificates}
\rulesection{75}{QZeroByMiranda}{core rule}
\label{blcatrule-QZeroByMiranda}
\ruleblock{The rule (Steele metanotation)}
\sourcenote{as realized in the mechanized development; source row: \texttt{\texttt{QTSL-rules.toml}} --- \emph{QZeroByMiranda}. The metanotation statement is typeset by hand for notational quality; this rule cites no Lean statements yet (the absence is stated loudly below). The rule \emph{roster} is machine-checked (every mechanized rule tag round-trips to exactly one catalogue row); the statement \emph{text} is editorial.}
\begin{mathpar}
\inferrule*{
  b : \mathsf{QMirandaBattery} \\
  \mathsf{qtsl\_miranda\_battery}(b,\ C,\ g,\ \Pi)
}{
  \vDash \mathsf{QZeroExists}(g,\ C)
}
\end{mathpar}

\par\medskip
\rulename{QZeroByMiranda} is Miranda\textquotesingle s theorem as a
certificate: $2m$ face rows over an equidimensional polynomial view,
each row a sign sequent claiming the signed component restricted to
one face of the cell nonpositive on the lower face and nonnegative on
the upper --- and the theorem concludes a zero of the vector field
inside the cell. Existence only: nothing in this family can conclude
uniqueness, count, or non-existence. Face geometry is derived from
each row\textquotesingle s indices, never stored --- a
producer-supplied face would be a trusted row --- and the
face-restriction recompute (exact-rational substitution) is the
walk\textquotesingle s derived-sequent authority.

\cataloguenote{The battery walk checks completeness (each (axis, face) exactly
once), per-axis component agreement, component bijectivity, and index
ranges as decided-false violations; non-box atoms, non-canonical or
over-wide polynomials, list ambiguity, and budget overruns refuse
with no verdict. Entered at the roster census (baked, never
catalogued); no Lean statement binds it yet.
}
\ruleblock{Theory rows (geolog)}
\sourcenote{The rule\textquotesingle s certificate content as it is declared in the geolog theory the checker loads --- payload schema, side-condition demands, and axiom rows, quoted from the theory source at every build of this catalogue (a stale quotation cannot exist). Axiom modes: \texttt{chase} = the engine derives; \texttt{report} = a violation is a loud judgment; \texttt{refuse} = a violation gates admission. Side-condition relations are discharged by bounded evaluation at the engine\textquotesingle s ground seam, never by derivation.}
\declhead{theory Qtsl}{QZeroByMiranda}{\badge{BadgeGreen}{quoted from the theory source at render}}
\begin{leansrc}
QZeroByMiranda : Set;
\end{leansrc}
\declhead{theory Qtsl}{QZeroByMiranda/*}{\badge{BadgeGreen}{quoted from the theory source at render}}
\begin{leansrc}
QZeroByMiranda/proof : QZeroByMiranda -> QProof;
QZeroByMiranda/g : QZeroByMiranda -> PolyVec;
QZeroByMiranda/cell : QZeroByMiranda -> QRegionAtom;
QZeroByMiranda/battery : QZeroByMiranda -> QMirandaBattery;
\end{leansrc}
\declhead{theory Qtsl}{qtsl\_miranda\_battery}{\badge{BadgeGreen}{quoted from the theory source at render}}
\begin{leansrc}
qtsl_miranda_battery : {battery: QMirandaBattery,
                        cell: QRegionAtom, g: PolyVec,
                        params: QRegionAtom} -> Prop;
\end{leansrc}
\ruleblock{Lean statement}
\badge{BadgeAmber}{not yet mechanized} this rule cites no Lean statements: nothing on this page is hash-bound for it, and its dashboard row shows the same absence across the statement, receipt, and correspondence columns.
\ruleblock{HOL4 statement}
\absentbox{NO HOL STATEMENT YET}{No HOL development of the extension theory exists.
}
\ruleblock{HOL verifier receipts}
\absentbox{NONE}{no HOL verification receipt exists for this rule.}
\ruleblock{Proof}

\sourcenote{Genre: core rule (tier 1), existence group. This section states the rule's soundness
obligation exactly, per the catalogue's completeness-over-polish
charter (a rule that can only be named-and-gapped is still a row);
the discharge is open debt on the mechanization column, and the
schema quoted in the entry is the baked ground truth.}

\textbf{Soundness obligation.}\enspace Miranda's theorem, certified: if each signed component of the
equidimensional polynomial view is nonpositive on its lower face and
nonnegative on its upper face of the cell (the $2m$ face rows, with
face geometry derived from row indices and each restriction
recomputed by exact-rational substitution), then the view has a zero
in the cell --- for every parameter point of the parameter cell.
Existence only; uniqueness, count, and non-existence are outside the
family by design.

\begin{bremark}[status]
No Lean statement binds this node. Miranda's theorem
(Poincar\'e--Miranda) is classical; the mechanization route would
state it over the boxed cell with the walk's derived face
restrictions as hypotheses.
\end{bremark}


\rulesection{76}{QEquilibriumFromInward}{admissible rule}
\label{blcatrule-QEquilibriumFromInward}
\ruleblock{The rule (Steele metanotation)}
\sourcenote{as realized in the mechanized development; source row: \texttt{\texttt{QTSL-rules.toml}} --- \emph{QEquilibriumFromInward}. The metanotation statement is typeset by hand for notational quality; this rule cites no Lean statements yet (the absence is stated loudly below). The rule \emph{roster} is machine-checked (every mechanized rule tag round-trips to exactly one catalogue row); the statement \emph{text} is editorial.}
\begin{mathpar}
\inferrule*{
  \mathrm{inv} : \mathsf{QBlockInv}\ \text{(strictly-inward facets, box block)}
}{
  \vDash \mathsf{QEquilibriumExists}(f,\ C)
}
\end{mathpar}

\par\medskip
\rulename{QEquilibriumFromInward} is a derived rule: a strictly-inward
facet battery on a box block --- every facet\textquotesingle s sign
sequent points the field inward --- is already the existence battery
for the negated field with identity pairing, so an equilibrium
inside the block follows by citation, with zero new oracle demand.
The premises are the block-invariant evidence rows themselves; the
unpacking instantiates the zero-existence obligation.

\cataloguenote{For fixed points the battery\textquotesingle s subject is the map
minus the identity; the elaboration recompute tying the two is a
named deferral at the ground seam. Entered at the roster census
(baked, never catalogued); no Lean statement binds it yet.
}
\ruleblock{Theory rows (geolog)}
\sourcenote{The rule\textquotesingle s certificate content as it is declared in the geolog theory the checker loads --- payload schema, side-condition demands, and axiom rows, quoted from the theory source at every build of this catalogue (a stale quotation cannot exist). Axiom modes: \texttt{chase} = the engine derives; \texttt{report} = a violation is a loud judgment; \texttt{refuse} = a violation gates admission. Side-condition relations are discharged by bounded evaluation at the engine\textquotesingle s ground seam, never by derivation.}
\declhead{theory Qtsl}{QEquilibriumFromInward}{\badge{BadgeGreen}{quoted from the theory source at render}}
\begin{leansrc}
QEquilibriumFromInward : Set;
\end{leansrc}
\declhead{theory Qtsl}{QEquilibriumFromInward/*}{\badge{BadgeGreen}{quoted from the theory source at render}}
\begin{leansrc}
QEquilibriumFromInward/proof : QEquilibriumFromInward -> QProof;
QEquilibriumFromInward/inv : QEquilibriumFromInward -> QBlockInv;
\end{leansrc}
\ruleblock{Lean statement}
\badge{BadgeAmber}{not yet mechanized} this rule cites no Lean statements: nothing on this page is hash-bound for it, and its dashboard row shows the same absence across the statement, receipt, and correspondence columns.
\ruleblock{HOL4 statement}
\absentbox{NO HOL STATEMENT YET}{No HOL development of the extension theory exists.
}
\ruleblock{HOL verifier receipts}
\absentbox{NONE}{no HOL verification receipt exists for this rule.}
\ruleblock{Proof}

\sourcenote{Genre: admissible rule (tier 2), existence group. This section states the rule's soundness
obligation exactly, per the catalogue's completeness-over-polish
charter (a rule that can only be named-and-gapped is still a row);
the discharge is open debt on the mechanization column, and the
schema quoted in the entry is the baked ground truth.}

\textbf{Soundness obligation.}\enspace A strictly-inward facet battery on a box block is the existence
battery for the negated field with identity pairing: inwardness on
every facet gives exactly the sign pattern Miranda's theorem demands
for $-f$, so the unpacking cites the block-invariant evidence rows
and instantiates the zero-existence obligation --- premises by
citation, zero new oracle demand. Soundness is therefore inherited
from the zero-existence rule's obligation.

\begin{bremark}[status]
No Lean statement binds this node; its obligation factors
entirely through the zero-existence rule's. For fixed points the
battery's subject is the map minus the identity, and the elaboration
recompute tying the two is a named deferral at the ground seam.
\end{bremark}


\rulesection{77}{QCoverByMiranda}{core rule}
\label{blcatrule-QCoverByMiranda}
\ruleblock{The rule (Steele metanotation)}
\sourcenote{as realized in the mechanized development; source row: \texttt{\texttt{QTSL-rules.toml}} --- \emph{QCoverByMiranda}. The metanotation statement is typeset by hand for notational quality; this rule cites no Lean statements yet (the absence is stated loudly below). The rule \emph{roster} is machine-checked (every mechanized rule tag round-trips to exactly one catalogue row); the statement \emph{text} is editorial.}
\begin{mathpar}
\inferrule*{
  b : \mathsf{QMirandaBattery} \\
  \mathsf{qtsl\_miranda\_battery}\bigl(b,\ W,\ \sigma - \mathrm{id}_{Z},\ Z\bigr)
}{
  \vDash Z \subseteq \sigma(W)
}
\end{mathpar}

\par\medskip
\rulename{QCoverByMiranda} is the chart-image discharge at
per-(patch cell, target cell) granularity: the battery\textquotesingle
s subject is chart minus target in joint variables, with the target
cell as parameters --- a certified zero of $\sigma(w) - z$ in the
patch cell for every $z$ in the target cell is exactly the
containment of the target cell in the chart image. Union to a region
cover rides the existing cover machinery; a cover is never inferred
from an image enclosure --- enclosures bound images from outside,
and only the zero-existence argument bounds them from inside.

\cataloguenote{Entered at the roster census (baked, never catalogued); no Lean
statement binds it yet. The battery walk is shared with the
zero-existence rule.
}
\ruleblock{Theory rows (geolog)}
\sourcenote{The rule\textquotesingle s certificate content as it is declared in the geolog theory the checker loads --- payload schema, side-condition demands, and axiom rows, quoted from the theory source at every build of this catalogue (a stale quotation cannot exist). Axiom modes: \texttt{chase} = the engine derives; \texttt{report} = a violation is a loud judgment; \texttt{refuse} = a violation gates admission. Side-condition relations are discharged by bounded evaluation at the engine\textquotesingle s ground seam, never by derivation.}
\declhead{theory Qtsl}{QCoverByMiranda}{\badge{BadgeGreen}{quoted from the theory source at render}}
\begin{leansrc}
QCoverByMiranda : Set;
\end{leansrc}
\declhead{theory Qtsl}{QCoverByMiranda/*}{\badge{BadgeGreen}{quoted from the theory source at render}}
\begin{leansrc}
QCoverByMiranda/proof : QCoverByMiranda -> QProof;
QCoverByMiranda/chart : QCoverByMiranda -> PolyVec;
QCoverByMiranda/wcell : QCoverByMiranda -> QRegionAtom;
QCoverByMiranda/zcell : QCoverByMiranda -> QRegionAtom;
QCoverByMiranda/battery : QCoverByMiranda -> QMirandaBattery;
\end{leansrc}
\ruleblock{Lean statement}
\badge{BadgeAmber}{not yet mechanized} this rule cites no Lean statements: nothing on this page is hash-bound for it, and its dashboard row shows the same absence across the statement, receipt, and correspondence columns.
\ruleblock{HOL4 statement}
\absentbox{NO HOL STATEMENT YET}{No HOL development of the extension theory exists.
}
\ruleblock{HOL verifier receipts}
\absentbox{NONE}{no HOL verification receipt exists for this rule.}
\ruleblock{Proof}

\sourcenote{Genre: core rule (tier 1), existence group. This section states the rule's soundness
obligation exactly, per the catalogue's completeness-over-polish
charter (a rule that can only be named-and-gapped is still a row);
the discharge is open debt on the mechanization column, and the
schema quoted in the entry is the baked ground truth.}

\textbf{Soundness obligation.}\enspace For each target point $z$ of the target cell, the battery (subject:
chart minus target in joint variables, parameters the target cell)
certifies a zero of $\sigma(w) - z$ in the patch cell --- i.e.\ a
preimage of $z$ under the chart --- so the target cell is contained
in the chart image. Union to a region cover rides the existing cover
machinery; a cover is never inferred from an image enclosure.

\begin{bremark}[status]
No Lean statement binds this node; the obligation factors through
the zero-existence rule's, applied parametrically in the target
point.
\end{bremark}


\blcatsection{Composition and Aggregation}
\rulesection{78}{QWireCompose}{core rule}
\label{blcatrule-QWireCompose}
\ruleblock{The rule (Steele metanotation)}
\sourcenote{as realized in the mechanized development; source row: \texttt{\texttt{QTSL-rules.toml}} --- \emph{QWireCompose}. The metanotation statement is typeset by hand for notational quality; its mechanized ground truth is the Lean statement below, extracted from the sources at every build of this book and hash-bound. The rule \emph{roster} is machine-checked (every mechanized rule tag round-trips to exactly one catalogue row); the statement \emph{text} is editorial, checked by eye against the extracted Lean beside it.}
\begin{mathpar}
\inferrule*[Right=\rulename{QWireCompose}]{
  \bigl\{\Sigma_i \vDash (I_i, \alpha_i, \gamma_i[, \kappa_i])\bigr\}_{i=1}^{n} \\
  \mathsf{QWireCert}\bigl(w,\ (\mathit{spec}_i)_{i} \Rightarrow \mathit{spec}_{\mathrm{out}},\ \kappa\bigr) \\
  \mathrm{rank} : \{1, \ldots, n\} \to \N; \quad \bigl\{\mathrm{rank}_i < \mathrm{rank}_j\bigr\}_{(i \to j) \in w} \ \text{(the causality certificate: one rank per component, strictly increasing along every wire)} \\
  \mathrm{W0}\!: \ \text{the composite martingale problem admits a c\`adl\`ag strong-Markov solution, unique in law, from every initial state (stochastic mechanisms; the martingale-problem sense --- distinct from the aggregation language\textquotesingle s own W0)} \\
  \mathrm{W1}\!: \ \text{transition data factor through per-place colour multisets (SDCPN mechanisms; typing-discharged in the aggregation language)} \\
  \mathrm{W2}\!: \ \text{the within-instant immediate-jump relation is well-founded, of uniformly bounded depth} \\
  \mathrm{W3}\!: \ \text{simultaneous enablings carry a declared serialization or a certified reset commutation (a.s.-null collision row for stochastic timers)} \\
  \kappa_{\mathrm{out}} := \textstyle\bigoplus_{\mathrm{legs}} \kappa_{\mathrm{leg}}, \quad \kappa_a \oplus \kappa_b := \kappa_a + \kappa_b \ \text{(slack addition along lens composition; checker-recomputed, never stored)}
}{
  \mathsf{SysWire}(w, [\Sigma_1, \ldots, \Sigma_n]) \vDash \mathit{spec}_{\mathrm{out}}
}
\end{mathpar}

\ruleflavor{satisfaction of an interface-bearing spec is plain \texttt{QSat} --- behaviours whose inputs conform to the assumption deliver the guarantee; closed $\mathcal{M} \vDash s$ is the empty-interface instance (defined once at the judgment-kinds home; the former subscripted display $\vDash_{\mathrm{I}}$ and the \texttt{QSatI} sort are both dissolved). The first two premises are derivation edges (judgment-kind premises); the well-posedness families enter as proof leaves on the node\textquotesingle s own well-posedness list, each owed only by the mechanism profiles that name it --- W0 by stochastic profiles, W1 by SDCPN profiles, W2/W3 the syntactic class (their four-way recompute relation is not yet enacted; the leaves enter as posted proof rows) --- and the causality ranks ride the node as one natural per component. The causality certificate and the slack composite are recomputed by the checker. Tier split derived from mechanism field-family population, never a stored tag: \textbf{admissible} where every mechanism derives ODE/discrete --- the composite certificate $\sum_i \lambda_i V_i$ is syntactically constructible and the node unpacks through the sequent-layer transport (\rulename{Promise}/\rulename{Syzygy}) --- \textbf{core} where any mechanism is stochastic (SDE/GSHS): the warrant is the per-profile semantic theorems. This display is the acyclic $n$-ary form (the causality certificate makes a directed cycle unrepresentable); feedback rides the binary small-gain cells, sum-form and gauge-form.}
\par\medskip
\rulename{QWireCompose} is the satisfaction-judgment
stratum\textquotesingle s composition rule --- assume--guarantee
transport over a structured system presentation. From
interface-indexed satisfaction of each component
($\mathsf{QSatI}(\Sigma_i, \mathit{spec}_i)$: every behaviour of
$\Sigma_i$ whose inputs satisfy the assumption $\alpha_i$ delivers the
guarantee $\gamma_i$, with quantitative slack where declared) and a
wiring certificate (the lens square transporting delivered guarantees
into downstream assumptions, with slack
$\kappa \in \mathcal{K}^0_\infty$), it concludes satisfaction of the
wired composite at the output spec. A producer uses it to certify a
system of interacting components --- each certified once, against its
own datasheet --- without ever building the monolithic certificate.

The soundness warrant factors, and the factorization organizes
everything on this page: $\mathrm{sound} = $ (transport arithmetic)
$\circ$ (semantic reading). The transport half is profile-generic
exact polynomial arithmetic --- the $(1, \lambda, 1, 1)$ cancellation
of cross-supplies, the sequent-layer machinery of \rulename{Promise}
and \rulename{Syzygy} --- and consumes only linearity of the
obligation operator over the wiring. The semantic half is
per-profile, and it is where this entry\textquotesingle s conditional
check lives. What the cited theorems prove, per profile and shape:
for general stochastic hybrid systems (GSHS), \emph{series}
composition (the conditioning tower, qualitative and product-bound
forms, plus the certified pipeline through the maximal-inequality
kernel), \emph{parallel} composition (the weighted-sum
supermartingale --- the semantic face of the syzygy weights), and
\emph{feedback in sum-form} (the cancellation run under conditional
expectation along the joint jump skeleton, with non-expansive
boundary rows) are proven; the $n$-ary \emph{acyclic} assembly is
proven at any $n$ by strong induction on the causality certificate,
with union-bound slack composition; and the per-profile tier split is
itself a proven, derived function of the presentation. \emph{Feedback
in gauge-form} (jumps that expand the certificate, repaid by flow
decay) is partial: its reduction theorem and the max-form
piecewise-deterministic instance are machine-checked end-to-end, while
the finite-index price of the diffusion term and the admissible-window
pinch remain open --- with the diffusion cell\textquotesingle s
parking reason retired: within-segment decay for gauged certificates
is now available at the martingale-problem profile, so the open part
is the accounting, not the stochastic calculus. The
deterministic-ODE row is at statement level;
the discrete-kernel row\textquotesingle s assembly is proven --- B1
$n$-ary acyclic, B2 parallel, and B3 sum-form small-gain, end-to-end
over the trajectory kernel --- while the sampled-data row remains
stated, not yet assembled. The discrete row carries a substrate
asymmetry worth stating with the grid: it consumes \emph{no}
interface hypothesis fields (contrast the GSHS cells\textquotesingle{}
D0--D4 data --- Dynkin, d\'ebut sampling, skeleton adaptedness),
because its execution semantics is mathlib\textquotesingle s own
\texttt{Kernel.traj}. The commentary states
this grid plainly because the dashboard\textquotesingle s conditional
mark warrants exactly the proven cells, never the whole catalogued
rule.

Two fences are part of the statement, not omissions. Undirected
stochastic wire sharing (common noise) is out of scope: stochastic
components sharing a wire need exactly one determining side, and the
proven parallel cell covers directed coupling and disjoint tensor
only. And the $n$-ary theorems are acyclic by construction --- the
presentation\textquotesingle s causality certificate makes a directed
cycle unrepresentable --- so feedback rides the binary small-gain
cells, where the coupling of jump times between the two components is
handled by the joint-skeleton argument.

What a checker recomputes, the rule\textquotesingle s row sorts now
baked: the causality certificate (a rank function strictly
increasing along wires), the slack composite (lens composition plus
slack addition, with the class membership rows discharged by
\rulename{QWireCertDischarge}), and the well-posedness row families
for the stochastic profiles. The component satisfactions and the
wiring squares are cited judgments and inherited sequent leaves ---
the rule adds no kernel arithmetic of its own.

\cataloguenote{The tier split is derived from mechanism field-family population,
never stored as a tag: admissible where every mechanism in the
presentation derives the ODE or discrete profile (the composite
certificate $\sum_i \lambda_i V_i$ is syntactically constructible and
the node unpacks through the sequent-layer transport), core where any
mechanism is stochastic (no unpacking exists; the semantic theorems
are the warrant). One rule node either way. The judgment sorts the
rule concludes into are baked in the extension theory; the checker
chain that would consume them is not yet ported --- the rule is
proven ahead of being wireable, and the dashboard row
shows exactly that state.
}
\ruleblock{Lean statement}
\declhead{BernsteinLogicAux.AGCompose.GshsSeries}{certified}{\badge{BadgeGreen}{statement bound}}
\begin{leansrc}
/-- **C1, certified form**: the full conditional-on-interface pipeline. Component
certificates enter as I-bridge outputs on the composite law — nonnegative
supermartingales along the stopped joint skeleton [D1][D2], with barrier values on
their bad sets (floor-row transport [D2]) and violation-witness inclusions [D3]
(the stopped skeleton samples the début of a violation; for component 2 the
certificate process is understood α₂-conformance-stopped, which is why its witness
hypothesis is relative to `conform₂`). The landed [K1] (`reach_avoid_div`,
`Ville.lean`) converts each certificate to a failure bound; `fail_bound` composes
them. Conclusion: the composite QSatI failure bound with certificate-valued ε's,

`μ (conform₁ ∩ satOutᶜ) ≤ ofReal (∫ Y₁ 0) / ofReal a₁ + ofReal (∫ Y₂ 0) / ofReal a₂`. -/
theorem certified [IsFiniteMeasure μ]
    {Y₁ Y₂ : ℕ → Ω → ℝ} {X₁ : ℕ → Ω → S₁} {X₂ : ℕ → Ω → S₂}
    {B₁ : Set S₁} {B₂ : Set S₂} {a₁ a₂ : ℝ}
    -- [D1]+[D2]: I-bridge outputs on the composite law
    (hY₁ : Supermartingale Y₁ 𝒢 μ) (hY₁0 : 0 ≤ Y₁)
    (hY₂ : Supermartingale Y₂ 𝒢 μ) (hY₂0 : 0 ≤ Y₂)
    -- floor-row barriers at skeleton samples [D2]
    (ha₁ : 0 < a₁) (ha₂ : 0 < a₂)
    (hbar₁ : ∀ n ω, X₁ n ω ∈ B₁ → a₁ ≤ Y₁ n ω)
    (hbar₂ : ∀ n ω, X₂ n ω ∈ B₂ → a₂ ≤ Y₂ n ω)
    -- violation witnesses at skeleton samples [D3]
    (hwit₁ : conform₁ ∩ sat₁ᶜ ⊆ {ω | ∃ n, X₁ n ω ∈ B₁})
    (hwit₂ : conform₂ ∩ sat₂ᶜ ⊆ {ω | ∃ n, X₂ n ω ∈ B₂})
    -- the wiring square (κ-slack folded) and the output reading
    (hw : conform₁ ∩ sat₁ ⊆ conform₂)
    (hout : conform₁ ∩ sat₁ ∩ sat₂ ⊆ satOut) :
    μ (conform₁ ∩ satOutᶜ) ≤
      ENNReal.ofReal (∫ ω, Y₁ 0 ω ∂μ) / ENNReal.ofReal a₁ +
      ENNReal.ofReal (∫ ω, Y₂ 0 ω ∂μ) / ENNReal.ofReal a₂
\end{leansrc}
\declhead{BernsteinLogicAux.AGCompose.GshsParallel}{certified}{\badge{BadgeGreen}{statement bound}}
\begin{leansrc}
/-- **C2, certified form** (n-ary, sum-form): component certificates — nonnegative
supermartingales `Y i` w.r.t. the joint skeleton filtration [D1][D2], syzygy
weights `λᵢ ≥ 0`, and per-component floor rows giving a common barrier level `a` on
each bad set (boundary rows transport pointwise through sums, inventory §3.C
claim 1: on `Bad i` the term `λᵢ·Yᵢ` alone clears `a`, the other terms being
nonnegative) — yield the composite avoid bound via the landed [K1]:

`ofReal a * μ {∃ n, X n ∈ ⋃ i ∈ s, Bad i} ≤ ofReal (∑ i ∈ s, λᵢ · ∫ Yᵢ 0)`.

The right side is the syzygy-weighted initial certificate mass — the semantic
reading of the composite certificate `∑ λᵢVᵢ` at the initial skeleton time. To
recover per-component barrier levels `aᵢ`, instantiate `a := min over s of λᵢ·aᵢ`. -/
theorem certified (s : Finset ι) [IsFiniteMeasure μ]
    {Y : ι → ℕ → Ω → ℝ} {lam : ι → ℝ} {X : ℕ → Ω → St} {Bad : ι → Set St} {a : ℝ}
    (hY : ∀ i ∈ s, Supermartingale (Y i) 𝒢 μ)
    (hY0 : ∀ i ∈ s, 0 ≤ Y i)
    (hlam : ∀ i ∈ s, 0 ≤ lam i)
    (hbar : ∀ i ∈ s, ∀ n ω, X n ω ∈ Bad i → a ≤ lam i * Y i n ω) :
    ENNReal.ofReal a * μ {ω | ∃ n, X n ω ∈ ⋃ i ∈ s, Bad i} ≤
      ENNReal.ofReal (∑ i ∈ s, lam i * ∫ ω, Y i 0 ω ∂μ)
\end{leansrc}
\declhead{BernsteinLogicAux.AGCompose.GshsParallel}{certified\_pair}{\badge{BadgeGreen}{statement bound}}
\begin{leansrc}
/-- **C2, certified form, two components** — the inventory's `V₁ + λV₂` shape,
directly. Barrier reading: component 1's floor clears `a` on `B₁`, component 2's
λ-weighted floor clears `a` on `B₂`; the composite avoid bound is

`ofReal a * μ {∃ n, X n ∈ B₁ ∪ B₂} ≤ ofReal (∫ Y₁ 0 + λ · ∫ Y₂ 0)`. -/
theorem certified_pair [IsFiniteMeasure μ]
    {Y₁ Y₂ : ℕ → Ω → ℝ} {X : ℕ → Ω → St} {B₁ B₂ : Set St} {lam a : ℝ}
    (hY₁ : Supermartingale Y₁ 𝒢 μ) (hY₁0 : 0 ≤ Y₁)
    (hY₂ : Supermartingale Y₂ 𝒢 μ) (hY₂0 : 0 ≤ Y₂) (hlam : 0 ≤ lam)
    (hbar₁ : ∀ n ω, X n ω ∈ B₁ → a ≤ Y₁ n ω)
    (hbar₂ : ∀ n ω, X n ω ∈ B₂ → a ≤ lam * Y₂ n ω) :
    ENNReal.ofReal a * μ {ω | ∃ n, X n ω ∈ B₁ ∪ B₂} ≤
      ENNReal.ofReal (∫ ω, Y₁ 0 ω ∂μ + lam * ∫ ω, Y₂ 0 ω ∂μ)
\end{leansrc}
\declhead{BernsteinLogicAux.GshsFeedbackSum}{supermartingale}{\badge{BadgeGreen}{statement bound}}
\begin{leansrc}
/-- **The composite certificate is a supermartingale** for the joint-skeleton
filtration — the C3a conclusion's engine. Follows from `drift` by discarding the
nonnegative θ-margin. -/
theorem supermartingale : Supermartingale F.W 𝒢 μ
\end{leansrc}
\declhead{BernsteinLogicAux.GshsFeedbackSum}{avoid\_bound}{\badge{BadgeGreen}{statement bound}}
\begin{leansrc}
/-- **The composite avoid bound** (K1, consumed as the landed
`BernsteinLogic.Ville.reach_avoid`): a barrier `a` for `W` on a sampled bad set `B`
bounds the probability that any joint-skeleton sample reaches `B`. -/
theorem avoid_bound {S : Type*} {X : ℕ → Ω → S} {B : Set S} {a : ℝ}
    (hbarrier : ∀ k ω, X k ω ∈ B → a ≤ F.W k ω) :
    ENNReal.ofReal a * μ {ω | ∃ k, X k ω ∈ B} ≤ ENNReal.ofReal (∫ ω, F.W 0 ω ∂μ)
\end{leansrc}
\declhead{BernsteinLogicAux.MaxFormPDMP}{sound\_GSHS\_safety\_barrier}{\badge{BadgeGreen}{statement bound}}
\begin{leansrc}
/-- **The T5 consumer at a general positive barrier** (Statement II's
Step 7 barrier `min(1, λ/λ')` needs this): with floor `m > 0` on the avoid
region, the reach probability is bounded by `V x / m`. Proved by scaling
the certificate to barrier `1` and consuming the landed
`sound_GSHS_safety`; no new engine plumbing. -/
theorem sound_GSHS_safety_barrier [IsMarkovKernel M.seg] [IsMarkovKernel M.jump]
    {W : State → ℝ} {U : Set State} {m : ℝ} (hm : 0 < m)
    (hW : Measurable W) (hU : MeasurableSet U)
    (hSegD : SegmentDynkin M 0 W) (hForce : ForcedNonexpansion M W)
    (hNonneg : NonnegV W) {C : ℝ} (hC0 : 0 ≤ C) (hWC : ∀ y, W y ≤ C)
    (hFloor : ∀ y ∈ U, m ≤ W y) (x : State) :
    Battery.pathMeasure (skel M) (SkelState.alive x) {ω | ∃ n, ω n ∈ skelBad U}
      ≤ ENNReal.ofReal (W x / m)
\end{leansrc}
\par\medskip
\textbf{The presentation layer, bound by display.}\enspace The
\texttt{System} structure is the abstraction the composition theorems
consume: the field-family population pattern of each mechanism (from
which the doctrine profile is \emph{derived}, never stored) and the
wiring order with its causality certificate carried as data --- a rank
function strictly increasing along wires, which is precisely what the
checker-side causality condition recomputes. \texttt{System.qagTier}
derives the rule\textquotesingle s tier from the same data, with both
unfoldings proven; \texttt{admissible\_constructs} makes the
admissible face\textquotesingle s \textquotedblleft syntactically
constructible\textquotedblright{} clause a checkable statement rather
than prose.

\declhead{BernsteinLogicAux.QtslCompose.Presentation}{System}{\badge{BadgeGreen}{statement bound}}
\begin{leansrc}
/-- A System presentation, at the composition-relevant abstraction: `n` charts
with their field-family population patterns, the wiring DAG (`upstream j` =
the charts whose outputs wire into `j`'s inputs), and the causality
certificate carried as **data** — a rank function strictly increasing along
wires, exactly what the checker's causality-DAG side condition recomputes.
Feedback (a directed cycle) is *unrepresentable* in this structure; the
feedback transport shape rides the landed binary cells
(`AGCompose/GshsFeedbackSum.lean`, `GaugeSmallGain.lean`, `MaxFormPDMP.lean`)
— see the AGCompose module docstring. -/
structure System (n : ℕ) : Type where
  charts : Fin n → ChartShape
  upstream : Fin n → Finset (Fin n)
  rank : Fin n → ℕ
  causality : ∀ j, ∀ i ∈ upstream j, rank i < rank j
\end{leansrc}
\declhead{BernsteinLogicAux.QtslCompose.Presentation}{System.qagTier}{\badge{BadgeGreen}{statement bound}}
\begin{leansrc}
/-- The `QWireCompose` tier at this System, **derived from the derived
profiles** (never a stored tag): admissible iff every chart's profile admits a
syntactically constructible composite certificate. -/
def System.qagTier {n : ℕ} (S : System n) : Tier
\end{leansrc}
\declhead{BernsteinLogicAux.QtslCompose.Presentation}{System.qagTier\_admissible\_iff}{\badge{BadgeGreen}{statement bound}}
\begin{leansrc}
/-- The split, unfolded to the profile characterization: admissible iff every
chart derives ODE or discrete. -/
theorem System.qagTier_admissible_iff {n : ℕ} (S : System n) :
    S.qagTier = .admissible ↔
      ∀ i, (S.charts i).profile = .ode ∨ (S.charts i).profile = .discrete
\end{leansrc}
\declhead{BernsteinLogicAux.QtslCompose.Presentation}{System.qagTier\_core\_iff}{\badge{BadgeGreen}{statement bound}}
\begin{leansrc}
/-- The core face of the split: core iff some chart is stochastic (SDE or
GSHS). -/
theorem System.qagTier_core_iff {n : ℕ} (S : System n) :
    S.qagTier = .core ↔
      ∃ i, (S.charts i).profile = .sde ∨ (S.charts i).profile = .gshs
\end{leansrc}
\declhead{BernsteinLogicAux.QtslCompose.Presentation}{admissible\_constructs}{\badge{BadgeGreen}{statement bound}}
\begin{leansrc}
/-- The core face of the split: core iff some chart is stochastic (SDE or
GSHS). -/
theorem System.qagTier_core_iff {n : ℕ} (S : System n) :
    S.qagTier = .core ↔
      ∃ i, (S.charts i).profile = .sde ∨ (S.charts i).profile = .gshs := by
  have hsplit : S.qagTier = .admissible ∨ S.qagTier = .core := by
    unfold System.qagTier
    split <;> simp
  constructor
  · intro h
    by_contra hno
    push_neg at hno
    have : S.qagTier = .admissible := by
      rw [S.qagTier_admissible_iff]
      intro i
      have := hno i
      cases hp : (S.charts i).profile <;> simp_all
    rw [this] at h
    exact absurd h (by simp)
  · intro ⟨i, hi⟩
    rcases hsplit with h | h
    · rw [S.qagTier_admissible_iff] at h
      rcases h i with ho | hd <;> rcases hi with hs | hs <;> simp_all
    · exact h

/-! ## The admissible face's constructibility, pointwise

The composite certificate for the admissible tier is the weighted sum
`∑ i, λ i • V i`; its composite obligation row follows *pointwise* from the
component dissipation rows (`e i ≤ 0` on the operating region) and the wiring
identification rows (`g k = 0` there), at nonnegative certificate weights and
free-sign wiring multipliers. This is the pointwise shadow of the landed
sequent-layer theorem (`sound_syzygy`, main package — the (T) layer of the
inventory §0 factorization); it lives here so that the tier definition's
"syntactically constructible" clause is a statement, not prose. -/
theorem admissible_constructs {ι κIdx X : Type*} [Fintype ι] [Fintype κIdx]
    {R : Set X} (lam : ι → ℝ) (e : ι → X → ℝ) (nu : κIdx → ℝ) (g : κIdx → X → ℝ)
    (hlam : ∀ i, 0 ≤ lam i)
    (he : ∀ i, ∀ x ∈ R, e i x ≤ 0)
    (hg : ∀ k, ∀ x ∈ R, g k x = 0) :
    ∀ x ∈ R, (∑ i, lam i * e i x) + (∑ k, nu k * g k x) ≤ 0
\end{leansrc}
\declhead{BernsteinLogicAux.QtslCompose.AGCompose}{sat\_of\_conform}{\badge{BadgeGreen}{statement bound}}
\begin{leansrc}
/-- **Conformance propagation** (pointwise core). If the wiring squares
transport external conformance + upstream deliveries into each component's
conformance, and each component delivers on conforming inputs (at this `ω`),
then every component delivers. Strong induction on the causality rank — the
presentation's causality certificate is exactly what makes this well-founded. -/
theorem sat_of_conform
    (hw : ∀ j, conformExt ∩ (⋂ i ∈ S.upstream j, sat i) ⊆ conform j)
    {ω : Ω} (h : ∀ j, ω ∈ conform j → ω ∈ sat j)
    (hext : ω ∈ conformExt) : ∀ j, ω ∈ sat j
\end{leansrc}
\declhead{BernsteinLogicAux.QtslCompose.AGCompose}{qualitative}{\badge{BadgeGreen}{statement bound}}
\begin{leansrc}
/-- **QWireCompose, qualitative form** (a.e.; the n-ary generalization of
`AGCompose.Series.qualitative`): component QSatI's as a.e. implications +
wiring squares yield the composite a.e. delivery of every guarantee, hence of
any output spec read from them (`hout`). -/
theorem qualitative [MeasurableSpace Ω] {μ : Measure Ω}
    (hw : ∀ j, conformExt ∩ (⋂ i ∈ S.upstream j, sat i) ⊆ conform j)
    (h : ∀ j, ∀ᵐ ω ∂μ, ω ∈ conform j → ω ∈ sat j)
    {satOut : Set Ω} (hout : conformExt ∩ (⋂ j, sat j) ⊆ satOut) :
    ∀ᵐ ω ∂μ, ω ∈ conformExt → ω ∈ satOut
\end{leansrc}
\declhead{BernsteinLogicAux.QtslCompose.AGCompose}{fail\_subset}{\badge{BadgeGreen}{statement bound}}
\begin{leansrc}
/-- **The n-ary failure decomposition** (pure set algebra; the rule's
union-bound side of slack composition): a composite failure under conforming
external inputs localizes at some component that failed *while conforming* —
namely a rank-minimal failing one, whose upstream all delivered. -/
theorem fail_subset
    (hw : ∀ j, conformExt ∩ (⋂ i ∈ S.upstream j, sat i) ⊆ conform j) :
    conformExt ∩ (⋂ j, sat j)ᶜ ⊆ ⋃ j, conform j ∩ (sat j)ᶜ
\end{leansrc}
\declhead{BernsteinLogicAux.QtslCompose.AGCompose}{fail\_bound}{\badge{BadgeGreen}{statement bound}}
\begin{leansrc}
/-- **QWireCompose, probabilistic form: n-ary union-bound slack composition** —
the inventory §3.B1 display `ℙ(composite fails) ≤ ∑ i (1 − p_i)` at any `n`,
in failure-measure form (`ε j` plays `1 − p_j`). -/
theorem fail_bound [MeasurableSpace Ω] {μ : Measure Ω} {ε : Fin n → ℝ≥0∞}
    (hw : ∀ j, conformExt ∩ (⋂ i ∈ S.upstream j, sat i) ⊆ conform j)
    (hε : ∀ j, μ (conform j ∩ (sat j)ᶜ) ≤ ε j) :
    μ (conformExt ∩ (⋂ j, sat j)ᶜ) ≤ ∑ j, ε j
\end{leansrc}
\declhead{BernsteinLogicAux.QtslCompose.AGCompose}{fail\_bound\_out}{\badge{BadgeGreen}{statement bound}}
\begin{leansrc}
/-- `fail_bound` read at an output spec (`hout` = the output reading of the
delivered guarantees, as in the binary cells). -/
theorem fail_bound_out [MeasurableSpace Ω] {μ : Measure Ω} {ε : Fin n → ℝ≥0∞}
    {satOut : Set Ω}
    (hw : ∀ j, conformExt ∩ (⋂ i ∈ S.upstream j, sat i) ⊆ conform j)
    (hout : conformExt ∩ (⋂ j, sat j) ⊆ satOut)
    (hε : ∀ j, μ (conform j ∩ (sat j)ᶜ) ≤ ε j) :
    μ (conformExt ∩ satOutᶜ) ≤ ∑ j, ε j
\end{leansrc}
\par\medskip
\textbf{The B-row (discrete-kernel) assembly.}\enspace The row the
commentary\textquotesingle s grid records as proven: B1
$n$-ary acyclic union-bound composition through the first-failure
device, B2 parallel weighted-sum on the product kernel, B3 sum-form
small-gain cancellation --- machine-checked end-to-end over
\texttt{Kernel.traj} through the battery substrate, with no semantic
hypothesis fields (the profile identification enters as the choice of
substrate, which is the conditional mark\textquotesingle s honest
residue for this row).

\declhead{BernsteinLogicAux.QtslCompose.AGComposeDiscrete}{firstFail}{\badge{BadgeGreen}{statement bound}}
\begin{leansrc}
/-- The event that component `j` is a **first** violator: at some time the path
is in `j`'s violation set, no component's violation set having been visited
before. Distinct components' first-failure events may overlap (simultaneous
first violations); the union bound is indifferent. -/
def firstFail (bad : Fin m → Set S) (j : Fin m) : Set (ℕ → S)
\end{leansrc}
\declhead{BernsteinLogicAux.QtslCompose.AGComposeDiscrete}{reach\_subset\_iUnion\_firstFail}{\badge{BadgeGreen}{statement bound}}
\begin{leansrc}
/-- **The first-failure decomposition** (the discrete-profile face of K4's event
decomposition; kin of the event-stratum `AGCompose.fail_subset`, with
time-minimality replacing rank-minimality): a path that ever visits any
violation set has a first violation, and that violation belongs to some
component. Pure path-set algebra. -/
theorem reach_subset_iUnion_firstFail (bad : Fin m → Set S) :
    {ω : ℕ → S | ∃ n, ω n ∈ ⋃ i, bad i} ⊆ ⋃ j, firstFail bad j
\end{leansrc}
\declhead{BernsteinLogicAux.QtslCompose.AGComposeDiscrete}{measure\_firstFail\_le}{\badge{BadgeGreen}{statement bound}}
\begin{leansrc}
/-- **The stopped-process device** (inventory §3.B1), isolated: stopping the
composite chain at the first violation by *any* component, component `j`'s
certificate — with drift live off the **joint** violation set and barrier only
on `j`'s **own** violation set — bounds the probability that `j` fails first
by its initial value. K1 (`Supermartingale.ville_ineq`) at barrier `1`;
contrast `Battery.measure_reach_le`, whose barrier must cover the whole stop
set and which therefore cannot see a multi-component decomposition. -/
theorem measure_firstFail_le {bad : Fin m → Set S} {V : S → ℝ} {C : ℝ} {j : Fin m}
    (hA : MeasurableSet (⋃ i, bad i)) (hV : Measurable V)
    (hV0 : ∀ y, 0 ≤ V y) (hVC : ∀ y, V y ≤ C)
    (hbar : ∀ y ∈ bad j, 1 ≤ V y)
    (hdrift : ∀ y, y ∉ (⋃ i, bad i) → ∫ z, V z ∂K y ≤ V y) (x : S) :
    pathMeasure K x (firstFail bad j) ≤ ENNReal.ofReal (V x)
\end{leansrc}
\declhead{BernsteinLogicAux.QtslCompose.AGComposeDiscrete}{fail\_bound}{\badge{BadgeGreen}{statement bound}}
\begin{leansrc}
/-- **QWireCompose B-row, probabilistic form** (engine form; inventory §3.B1's
display `ℙ(composite fails) ≤ ∑ᵢ (1−pᵢ)` at any `n`, semantically discharged
over `Kernel.traj`): per-component bounded certificates with drift off the
joint violation set and barriers on their own violation sets bound the
composite failure probability by the sum of initial certificate values. The
ε-shaped quantities the event-stratum `AGCompose.fail_bound` consumes as
hypotheses are here established from kernel-level data. -/
theorem fail_bound {bad : Fin m → Set S} {V : Fin m → S → ℝ} {C : Fin m → ℝ}
    (hA : MeasurableSet (⋃ i, bad i))
    (hV : ∀ j, Measurable (V j)) (hV0 : ∀ j y, 0 ≤ V j y)
    (hVC : ∀ j y, V j y ≤ C j)
    (hbar : ∀ j, ∀ y ∈ bad j, 1 ≤ V j y)
    (hdrift : ∀ j, ∀ y, y ∉ (⋃ i, bad i) → ∫ z, V j z ∂K y ≤ V j y) (x : S) :
    pathMeasure K x {ω | ∃ n, ω n ∈ ⋃ i, bad i}
      ≤ ∑ j, ENNReal.ofReal (V j x)
\end{leansrc}
\declhead{BernsteinLogicAux.QtslCompose.AGComposeDiscrete}{fail\_bound\_wired}{\badge{BadgeGreen}{statement bound}}
\begin{leansrc}
/-- **QWireCompose B-row, wired form** (the catalogued rule's hypothesis shape):
per-component **datasheet** rows — drift valid only on the component's own
conformance region — plus the state-level wiring square (external conformance
folded; joint upstream delivery implies conformance, the I-bridge reading of
`QWireCert` at this stratum) yield the same union-bound conclusion. The
`System`'s wiring order supplies the square's index set; its rank/causality
data is consumed by the event-stratum assembly, not here (time-minimality
replaces the rank induction — see the module header, and `fail_bound_upstream`
for the honest-width form this now specializes). -/
theorem fail_bound_wired (Sys : System m) {conf bad : Fin m → Set S}
    {V : Fin m → S → ℝ} {C : Fin m → ℝ}
    (hA : MeasurableSet (⋃ i, bad i))
    (hw : ∀ j, (⋂ i ∈ Sys.upstream j, (bad i)ᶜ) ⊆ conf j)
    (hV : ∀ j, Measurable (V j)) (hV0 : ∀ j y, 0 ≤ V j y)
    (hVC : ∀ j y, V j y ≤ C j)
    (hbar : ∀ j, ∀ y ∈ bad j, 1 ≤ V j y)
    (hdrift : ∀ j, ∀ y ∈ conf j, y ∉ bad j → ∫ z, V j z ∂K y ≤ V j y) (x : S) :
    pathMeasure K x {ω | ∃ n, ω n ∈ ⋃ i, bad i}
      ≤ ∑ j, ENNReal.ofReal (V j x)
\end{leansrc}
\declhead{BernsteinLogicAux.QtslCompose.AGComposeDiscrete}{ae\_avoid\_of\_zero}{\badge{BadgeGreen}{statement bound}}
\begin{leansrc}
/-- **The qualitative reading** (inventory §3.B1's "a.s.-invariance-shaped"
half): started strictly inside every certificate's zero set, the composite
chain a.s. never violates any component's guarantee — the conditioning-tower
display degenerates, at the discrete profile, to the union bound at total
mass zero. -/
theorem ae_avoid_of_zero {bad : Fin m → Set S} {V : Fin m → S → ℝ}
    {C : Fin m → ℝ}
    (hA : MeasurableSet (⋃ i, bad i))
    (hV : ∀ j, Measurable (V j)) (hV0 : ∀ j y, 0 ≤ V j y)
    (hVC : ∀ j y, V j y ≤ C j)
    (hbar : ∀ j, ∀ y ∈ bad j, 1 ≤ V j y)
    (hdrift : ∀ j, ∀ y, y ∉ (⋃ i, bad i) → ∫ z, V j z ∂K y ≤ V j y) {x : S}
    (hzero : ∀ j, V j x = 0) :
    ∀ᵐ ω ∂pathMeasure K x, ∀ n, ω n ∉ ⋃ i, bad i
\end{leansrc}
\declhead{BernsteinLogicAux.QtslCompose.AGComposeDiscrete}{drift\_parallelComp}{\badge{BadgeGreen}{statement bound}}
\begin{leansrc}
/-- **B2, the linearity row** (inventory §3.B2's one-paragraph argument,
machine-checked): component drift rows transport to the weighted-sum
certificate on the synchronous product — the syzygy combination at weights
`(1, λ)`, verbatim. Off the lifted joint violation set both component rows
are live; linearity of the integral over `K₁ ∥ₖ K₂` does the rest. -/
theorem drift_parallelComp {bad₁ : Set S₁} {bad₂ : Set S₂}
    {V₁ : S₁ → ℝ} {V₂ : S₂ → ℝ} {lam : ℝ} (hlam : 0 ≤ lam)
    (hint₁ : ∀ y, y ∉ bad₁ → Integrable V₁ (K₁ y))
    (hint₂ : ∀ y, y ∉ bad₂ → Integrable V₂ (K₂ y))
    (hd₁ : ∀ y, y ∉ bad₁ → ∫ z, V₁ z ∂K₁ y ≤ V₁ y)
    (hd₂ : ∀ y, y ∉ bad₂ → ∫ z, V₂ z ∂K₂ y ≤ V₂ y) :
    ∀ p : S₁ × S₂, p ∉ (Prod.fst ⁻¹' bad₁ ∪ Prod.snd ⁻¹' bad₂) →
      ∫ q, (V₁ q.1 + lam * V₂ q.2) ∂(K₁ ∥ₖ K₂) p
        ≤ V₁ p.1 + lam * V₂ p.2
\end{leansrc}
\declhead{BernsteinLogicAux.QtslCompose.AGComposeDiscrete}{avoid\_parallelComp}{\badge{BadgeGreen}{statement bound}}
\begin{leansrc}
/-- **B2, assembled** (the discrete face of `GshsParallel.certified`'s
weighted-sum supermartingale claim, end-to-end): on the synchronous product
chain, the weighted-sum certificate bounds the probability that *either*
component ever violates by its initial value. The weight's role is visible in
the barrier hypotheses: component 2's barrier is carried at level `1/λ`
(stated as `1 ≤ λ·V₂` — how the syzygy weights meet the common barrier). -/
theorem avoid_parallelComp {bad₁ : Set S₁} {bad₂ : Set S₂}
    {V₁ : S₁ → ℝ} {V₂ : S₂ → ℝ} {C₁ C₂ lam : ℝ} (hlam : 0 ≤ lam)
    (hbad₁ : MeasurableSet bad₁) (hbad₂ : MeasurableSet bad₂)
    (hV₁ : Measurable V₁) (hV₂ : Measurable V₂)
    (hV₁0 : ∀ y, 0 ≤ V₁ y) (hV₂0 : ∀ y, 0 ≤ V₂ y)
    (hV₁C : ∀ y, V₁ y ≤ C₁) (hV₂C : ∀ y, V₂ y ≤ C₂)
    (hbar₁ : ∀ y ∈ bad₁, 1 ≤ V₁ y) (hbar₂ : ∀ y ∈ bad₂, 1 ≤ lam * V₂ y)
    (hd₁ : ∀ y, y ∉ bad₁ → ∫ z, V₁ z ∂K₁ y ≤ V₁ y)
    (hd₂ : ∀ y, y ∉ bad₂ → ∫ z, V₂ z ∂K₂ y ≤ V₂ y) (x : S₁ × S₂) :
    pathMeasure (K₁ ∥ₖ K₂) x
        {ω | ∃ n, ω n ∈ (Prod.fst ⁻¹' bad₁ ∪ Prod.snd ⁻¹' bad₂ : Set (S₁ × S₂))}
      ≤ ENNReal.ofReal (V₁ x.1 + lam * V₂ x.2)
\end{leansrc}
\declhead{BernsteinLogicAux.QtslCompose.AGComposeDiscrete}{smallGain\_drift}{\badge{BadgeGreen}{statement bound}}
\begin{leansrc}
/-- **B3, the cancellation row** (inventory §3.B3; the semantic face of the
checker-level `(1, λ, 1, 1)` syzygy fold `sound_smallGain`, at the kernel
level): per-component dissipation rows with cross-supplies on the joint
closed-loop kernel, under the small-gain membership rows, yield the
`θ`-margin drift row for `W = V₁ + λV₂`. Only `0 ≤ λ`, `0 ≤ θ` are consumed;
the rule-level ranges `λ > 0`, `θ ∈ (0,1)` are well-formedness, never
hypotheses (the width find of the C3a cell, reproduced at this profile). -/
theorem smallGain_drift (hlam : 0 ≤ lam)
    (hint₁ : ∀ y, y ∉ A → Integrable V₁ (K y))
    (hint₂ : ∀ y, y ∉ A → Integrable V₂ (K y))
    (hd₁ : ∀ y, y ∉ A → ∫ z, V₁ z ∂K y ≤ V₁ y - alpha₁ y + sigma₁ y)
    (hd₂ : ∀ y, y ∉ A → ∫ z, V₂ z ∂K y ≤ V₂ y - alpha₂ y + sigma₂ y)
    (hm₁ : ∀ y, y ∉ A → lam * sigma₂ y ≤ (1 - theta) * alpha₁ y)
    (hm₂ : ∀ y, y ∉ A → sigma₁ y ≤ lam * ((1 - theta) * alpha₂ y)) :
    ∀ y, y ∉ A → ∫ z, (V₁ z + lam * V₂ z) ∂K y
      ≤ (V₁ y + lam * V₂ y) - theta * (alpha₁ y + lam * alpha₂ y)
\end{leansrc}
\declhead{BernsteinLogicAux.QtslCompose.AGComposeDiscrete}{avoid\_smallGain}{\badge{BadgeGreen}{statement bound}}
\begin{leansrc}
/-- **B3, assembled** (the discrete face of the C3a conclusion, end-to-end):
on the closed loop, the small-gain data yields the composite avoid bound for
the weighted-sum certificate — Foster–Lyapunov metatheory
(`ville_ineq`/optional stopping, through `Battery.measure_reach_le`) closing
the cell exactly as the inventory's §3.B3 statement names it. The `α`-sign
hypotheses convert the `θ`-margin into a plain drift row; the barrier is the
caller's (how `1 ≤ W` distributes between components is model data). -/
theorem avoid_smallGain (hlam : 0 ≤ lam) (htheta : 0 ≤ theta)
    {C₁ C₂ : ℝ} (hA : MeasurableSet A)
    (hV₁ : Measurable V₁) (hV₂ : Measurable V₂)
    (hV₁0 : ∀ y, 0 ≤ V₁ y) (hV₂0 : ∀ y, 0 ≤ V₂ y)
    (hV₁C : ∀ y, V₁ y ≤ C₁) (hV₂C : ∀ y, V₂ y ≤ C₂)
    (hbar : ∀ y ∈ A, 1 ≤ V₁ y + lam * V₂ y)
    (halpha₁ : ∀ y, y ∉ A → 0 ≤ alpha₁ y)
    (halpha₂ : ∀ y, y ∉ A → 0 ≤ alpha₂ y)
    (hd₁ : ∀ y, y ∉ A → ∫ z, V₁ z ∂K y ≤ V₁ y - alpha₁ y + sigma₁ y)
    (hd₂ : ∀ y, y ∉ A → ∫ z, V₂ z ∂K y ≤ V₂ y - alpha₂ y + sigma₂ y)
    (hm₁ : ∀ y, y ∉ A → lam * sigma₂ y ≤ (1 - theta) * alpha₁ y)
    (hm₂ : ∀ y, y ∉ A → sigma₁ y ≤ lam * ((1 - theta) * alpha₂ y)) (x : S) :
    pathMeasure K x {ω | ∃ n, ω n ∈ A}
      ≤ ENNReal.ofReal (V₁ x + lam * V₂ x)
\end{leansrc}
\ruleblock{Correspondence (Steele display $\leftrightarrow$ Lean)}
\badge{BadgeAmber}{a bound Lean statement CHANGED since certification --- re-review and re-certify}
\sourcenote{The dictionary below is a reviewer-certified statement of how the mechanized statements above interpret the displayed rule: display metavariable $\mapsto$ Lean binder, each display hypothesis $\mapsto$ the premise or side condition carrying it, the display conclusion $\mapsto$ the goal --- with any deliberate difference in width between display and mechanization named outright. Certification is recorded as a hash triple over (the displayed cell, the extracted statement texts, this dictionary\textquotesingle s own text); if any of the three changes, this block turns amber until a reviewer re-certifies. The correspondence claim itself remains editorial --- reviewer-grade, not machine-checked; its machine-checked successor types each soundness statement against the rule as a structured object, and supersedes this block per rule as it lands.}
The display is one $n$-ary rule; the bound set is a \emph{grid}, and
the dictionary\textquotesingle s first duty is to say which display
constituent lands in which theorem. All statements live over the
compiled behaviour space (\texttt{Ω}, the composite execution
law\textquotesingle s carrier), reached from the presentation through
bridge hypotheses --- the events \texttt{conform j} / \texttt{sat j}
(component $j$\textquotesingle s input-conformance and
guarantee-delivery, interpreted over composite behaviours: the
I-bridge of the kernel-semantics interface) --- so the certified
statements quantify over exactly the Systems whose compilation
discharges those bridges.
\begin{itemize}
\item the E1 premises $\Sigma_i \vDash (\ldots)$ (the
  \texttt{QSatI} judgments) $\mapsto$,
  $n$-ary face: \texttt{h : ∀ j, ∀ᵐ ω ∂μ, ω ∈ conform j → ω ∈ sat j}
  (\texttt{qualitative}) and quantitatively \texttt{hε : μ (conform j
  ∩ (sat j)ᶜ) ≤ ε j} (\texttt{fail\_bound}) --- per-component
  satisfaction as a.e.\ implication / failure-measure bound.
\item the $\mathsf{QWireCert}$ premise $\mapsto$ \texttt{hw : ∀ j,
  conformExt ∩ (⋂ i ∈ S.upstream j, sat i) ⊆ conform j}: the wiring
  squares with the $\kappa$-slack \emph{folded into the predicates}.
  Consequence, stated plainly: the displayed $\kappa_{\mathrm{out}}$
  recompute row has no counterpart in the $n$-ary bound statements
  --- quantitative slack composition appears instead as the
  union-bound $\sum_j \varepsilon_j$ (\texttt{fail\_bound\_out}, the
  inventory\textquotesingle s $\sum_i (1 - p_i)$ display at any $n$).
\item the causality certificate $\mapsto$ \texttt{System.rank} +
  \texttt{System.causality} as \emph{data} (the structure is bound by
  display): rank strictly increasing along wires is a field, the
  induction of \texttt{sat\_of\_conform} runs on it, and a directed
  cycle is unrepresentable --- the Lean face of ``recomputed.''
\item W0--W3 well-posedness rows $\mapsto$ nothing bound ---
  schema-stratum families, honestly absent.
\item the conclusion $\mathsf{SysWire}(w, \vec{\Sigma}) \vDash
  \mathit{spec}_{\mathrm{out}}$ $\mapsto$ \texttt{qualitative}\textquotesingle s
  goal (a.e.\ delivery of \texttt{satOut} on conforming inputs) and
  \texttt{fail\_bound\_out}\textquotesingle s (failure measure at the
  output spec), with \texttt{hout} the output reading of the
  delivered guarantees; \texttt{fail\_subset} is the pure
  set-algebra localization engine (a composite failure localizes at a
  rank-minimal failing component, conforming there).
\item the tier split (the display\textquotesingle s
  admissible/core tier split, stated in the flavor text) $\mapsto$ \texttt{System.qagTier}, a
  \emph{derived function of the presentation} with both unfoldings
  proven (\texttt{qagTier\_admissible\_iff}: every chart ODE or
  discrete; \texttt{qagTier\_core\_iff}: some chart stochastic), and
  \texttt{admissible\_constructs} making ``syntactically
  constructible'' a statement: composite row of $\sum_i \lambda_i
  V_i$ from component dissipation rows + wiring identification rows,
  pointwise, at nonnegative certificate weights and free-sign wiring
  multipliers --- the pointwise shadow of the main
  package\textquotesingle s \texttt{sound\_syzygy} (witness prose,
  package boundary).
\item the binary GSHS cells, premise shapes as bound:
  \texttt{Series.certified} consumes two supermartingale certificates
  (\texttt{Supermartingale Yᵢ 𝒢 μ}, nonneg, positive barriers,
  floor-row transport \texttt{hbarᵢ}, violation-witness inclusions
  \texttt{hwitᵢ}) + one wiring square \texttt{conform₁ ∩ sat₁ ⊆
  conform₂}, concluding the composite failure bound with
  certificate-valued $\varepsilon$\textquotesingle s (initial mass
  over barrier, per component); \texttt{Parallel.certified} /
  \texttt{certified\_pair} conclude the avoid bound for the
  weighted-sum certificate (the semantic face of the syzygy
  weights); \texttt{FeedbackSum.supermartingale} /
  \texttt{avoid\_bound} operate on the bundled
  \texttt{SumFormFeedback} structure --- and its binder rows carry a
  width find: soundness consumes only \texttt{0 ≤ lambda}, \texttt{0
  ≤ theta}; the rule-level ranges $\lambda > 0$, $\theta \in (0,1)$
  are well-formedness, never hypotheses --- with the $(1, \lambda, 1,
  1)$ cancellation entering as the two integrated a.e.\ membership
  rows \texttt{mem₁}/\texttt{mem₂};
  \texttt{sound\_GSHS\_safety\_barrier} is the $\kappa = 0$ max-form
  PDMP instance at a general positive floor \texttt{m}, consuming
  measurability, kernel instances, Dynkin/non-expansion data, and a
  global bound \texttt{C} on the certificate.
\item the B-row (discrete-kernel) assembly
  (\texttt{AGComposeDiscrete}; end-to-end over \texttt{Kernel.traj}
  through the battery substrate, with \emph{no} D-layer hypothesis
  fields --- the semantic hypotheses of the GSHS cells above have no
  counterpart here because the discrete substrate is in Lean): the E1
  premises $\mapsto$ per-component bounded certificates \texttt{V j}
  with datasheet drift on \texttt{conf j} and barrier \texttt{1 ≤ V j}
  on \texttt{bad j} (state-level); the $\mathsf{QWireCert}$ premise
  $\mapsto$ \texttt{hw : ⋂ i ∈ upstream j, (bad i)ᶜ ⊆ conf j} --- the
  square pre-folded at the state stratum, while rank and causality
  stay event-stratum: time-minimality (the first-failure device,
  \texttt{firstFail} / \texttt{measure\_firstFail\_le}, the
  K1-consumption primitive the multi-component setting needs) replaces
  the rank induction, so \texttt{System.rank}/\texttt{causality} are
  honestly unused in \texttt{fail\_bound\_wired}; the conclusion
  $\mapsto$ the failure-measure bound $\sum_j
  \mathrm{ofReal}(V_j\,x)$, the
  $\varepsilon$\textquotesingle s of the event-stratum
  \texttt{fail\_bound} \emph{established} rather than assumed. K2
  (optional stopping) is consumed transitively --- it is bridge step 1
  inside \texttt{ville\_ineq} --- and no direct optional-stopping
  consumption is needed for avoid-shaped conclusions. B2/B3 transpose
  the binary GSHS cells to the kernel level:
  \texttt{drift\_parallelComp} / \texttt{avoid\_parallelComp} (the
  weighted-sum row on \texttt{K₁ ∥ₖ K₂}, the discrete face of the
  parallel cell) and \texttt{smallGain\_drift} /
  \texttt{avoid\_smallGain} (the $(1, \lambda, 1, 1)$ cancellation,
  consuming only \texttt{0 ≤ lam} and \texttt{0 ≤ theta} --- the
  sum-form cell\textquotesingle s width find reproduced at this
  stratum).
\end{itemize}
Width note: the conditional mark\textquotesingle s content is the
grid-vs-rule gap the commentary already states (the deterministic-ODE row at statement level, the sampled-data row stated but not yet assembled, feedback in gauge-form partial; the discrete-kernel row's assembly is discharged, so the gap is one row lighter) --- the dictionary adds
that within the \emph{proven} cells the mappings above are exact, that
the $n$-ary statements are acyclic by type (not by side condition),
and that no bound statement consumes the $\mathcal{K}^0_\infty$
structure of $\kappa$ itself (that lives in
\rulename{QWireCertDischarge}\textquotesingle s bound set; here slack
is pre-folded or union-bounded).

\ruleblock{Lean verifier receipts}
\sourcenote{module \texttt{BernsteinLogicAux.AGCompose.GshsSeries}; bound build \texttt{ff5cb858ef68\ldots} (bound from receipt 911818331ed966fc, 2026-07-20T01:42:13Z). A receipt records that named checking layers accepted a build: (module, compiled-artifact SHA-256, environment SHA-256, checker layers), minted when the checking gate ran; the store receipts attesting the bound build are listed with the layers each run credited. Nothing here asks the reader to trust the store: every layer is an independently runnable program over the quoted sources --- rebuild the module and re-run the checkers to reproduce the verdict. Receipts attesting superseded builds credit no layer and are omitted from this rendering (the count of omissions, if any, is noted below).}
\begin{receipttable}
\receiptrow{cb0f28f9ca6b1c44}{2026-07-19T18:23:37Z}{build+decorrelated}{\badge{BadgeGreen}{current}}
\receiptrow{c973aa90108c94d9}{2026-07-20T01:20:47Z}{build+decorrelated}{\badge{BadgeGreen}{current}}
\receiptrow{911818331ed966fc}{2026-07-20T01:42:13Z}{build+decorrelated}{\badge{BadgeGreen}{current}}
\receiptrow{eada292d784c9848}{2026-07-20T20:46:24Z}{build+decorrelated}{\badge{BadgeGreen}{current}}
\receiptrow{1a58b3dfc83a9cb5}{2026-07-20T20:57:57Z}{build+replay+hardened+decorrelated}{\badge{BadgeGreen}{current}}
\receiptrow{3358b4d03fa77d7c}{2026-07-21T00:02:20Z}{build+decorrelated}{\badge{BadgeGreen}{current}}
\receiptrow{d4f9f9c67009d52f}{2026-07-21T05:38:28Z}{build+replay+hardened+decorrelated}{\badge{BadgeGreen}{current}}
\receiptrow{f86a2ccbe493f911}{2026-07-22T18:26:47Z}{build+decorrelated}{\badge{BadgeGreen}{current}}
\receiptrow{6d17867e1f4f7cdc}{2026-07-31T21:59:37Z}{build+decorrelated}{\badge{BadgeGreen}{current}}
\receiptrow{f4d14ab85f9d9554}{2026-07-31T22:28:31Z}{build+decorrelated}{\badge{BadgeGreen}{current}}
\receiptrow{90881889f02ba9fe}{2026-08-02T23:30:38Z}{build+decorrelated}{\badge{BadgeGreen}{current}}
\receiptrow{684b3dcbc3c2e336}{2026-08-03T16:26:42Z}{build+decorrelated}{\badge{BadgeGreen}{current}}
\receiptrow{782c96e76de07ffa}{2026-08-03T22:45:26Z}{build+decorrelated}{\badge{BadgeGreen}{current}}
\end{receipttable}
\sourcenote{2 further receipt(s) in the store attest superseded builds of this module and are omitted from this rendering.}
\sourcenote{module \texttt{BernsteinLogicAux.AGCompose.GshsParallel}; bound build \texttt{8b8d0ee45b1a\ldots} (hashed from the built olean at bind time). A receipt records that named checking layers accepted a build: (module, compiled-artifact SHA-256, environment SHA-256, checker layers), minted when the checking gate ran; the store receipts attesting the bound build are listed with the layers each run credited. Nothing here asks the reader to trust the store: every layer is an independently runnable program over the quoted sources --- rebuild the module and re-run the checkers to reproduce the verdict. Receipts attesting superseded builds credit no layer and are omitted from this rendering (the count of omissions, if any, is noted below).}
\begin{receipttable}
\receiptrow{f455dcef7549cd27}{2026-07-18T13:22:39Z}{build+decorrelated}{\badge{BadgeGreen}{current}}
\end{receipttable}
\sourcenote{12 further receipt(s) in the store attest superseded builds of this module and are omitted from this rendering.}
\sourcenote{module \texttt{BernsteinLogicAux.GshsFeedbackSum}; bound build \texttt{a0ce53c03e8d\ldots} (bound from receipt 45ff739d5d3e0306, 2026-07-20T07:45:49Z). A receipt records that named checking layers accepted a build: (module, compiled-artifact SHA-256, environment SHA-256, checker layers), minted when the checking gate ran; the store receipts attesting the bound build are listed with the layers each run credited. Nothing here asks the reader to trust the store: every layer is an independently runnable program over the quoted sources --- rebuild the module and re-run the checkers to reproduce the verdict. Receipts attesting superseded builds credit no layer and are omitted from this rendering (the count of omissions, if any, is noted below).}
\begin{receipttable}
\receiptrow{45ff739d5d3e0306}{2026-07-20T07:45:49Z}{build+decorrelated}{\badge{BadgeGreen}{current}}
\receiptrow{eb766d96123685ba}{2026-07-20T20:46:31Z}{build+decorrelated}{\badge{BadgeGreen}{current}}
\receiptrow{f0ef6207375847ca}{2026-07-20T20:58:04Z}{build+replay+hardened+decorrelated}{\badge{BadgeGreen}{current}}
\receiptrow{cf3cd505d8c77450}{2026-07-21T00:02:49Z}{build+decorrelated}{\badge{BadgeGreen}{current}}
\receiptrow{51f1db6c22c527e0}{2026-07-21T05:38:37Z}{build+replay+hardened+decorrelated}{\badge{BadgeGreen}{current}}
\receiptrow{474b9edb3b9dbacd}{2026-07-22T18:26:54Z}{build+decorrelated}{\badge{BadgeGreen}{current}}
\receiptrow{b0859917a7e80e56}{2026-07-31T21:59:48Z}{build+decorrelated}{\badge{BadgeGreen}{current}}
\receiptrow{e0b2e5fdbdb04064}{2026-07-31T22:28:42Z}{build+decorrelated}{\badge{BadgeGreen}{current}}
\receiptrow{57b040e29c83ab7a}{2026-08-02T23:30:50Z}{build+decorrelated}{\badge{BadgeGreen}{current}}
\receiptrow{d425bb06e18d9d85}{2026-08-03T16:26:54Z}{build+decorrelated}{\badge{BadgeGreen}{current}}
\receiptrow{30583333e955b957}{2026-08-03T22:45:37Z}{build+decorrelated}{\badge{BadgeGreen}{current}}
\end{receipttable}
\sourcenote{3 further receipt(s) in the store attest superseded builds of this module and are omitted from this rendering.}
\sourcenote{module \texttt{BernsteinLogicAux.MaxFormPDMP}; bound build \texttt{426c38790122\ldots} (bound from receipt 0a8c057e8cbf9aab, 2026-07-20T01:42:14Z). A receipt records that named checking layers accepted a build: (module, compiled-artifact SHA-256, environment SHA-256, checker layers), minted when the checking gate ran; the store receipts attesting the bound build are listed with the layers each run credited. Nothing here asks the reader to trust the store: every layer is an independently runnable program over the quoted sources --- rebuild the module and re-run the checkers to reproduce the verdict. Receipts attesting superseded builds credit no layer and are omitted from this rendering (the count of omissions, if any, is noted below).}
\begin{receipttable}
\receiptrow{42fadeaf4c4d0d5b}{2026-07-19T18:24:26Z}{build+decorrelated}{\badge{BadgeGreen}{current}}
\receiptrow{81de1969c895bf9f}{2026-07-20T01:21:39Z}{build+decorrelated}{\badge{BadgeGreen}{current}}
\receiptrow{0a8c057e8cbf9aab}{2026-07-20T01:42:14Z}{build+decorrelated}{\badge{BadgeGreen}{current}}
\receiptrow{12151b1f2bcd3a91}{2026-07-20T20:46:35Z}{build+decorrelated}{\badge{BadgeGreen}{current}}
\receiptrow{5821a7ee1f92823e}{2026-07-20T20:58:08Z}{build+replay+hardened+decorrelated}{\badge{BadgeGreen}{current}}
\receiptrow{c1facca742c8a57f}{2026-07-21T00:03:24Z}{build+decorrelated}{\badge{BadgeGreen}{current}}
\receiptrow{b37c4c71b35d3ed7}{2026-07-21T05:38:51Z}{build+replay+hardened+decorrelated}{\badge{BadgeGreen}{current}}
\receiptrow{e32a92fac2e848c6}{2026-07-22T18:27:07Z}{build+decorrelated}{\badge{BadgeGreen}{current}}
\receiptrow{2ed08aae31ac4f48}{2026-07-31T22:00:02Z}{build+decorrelated}{\badge{BadgeGreen}{current}}
\receiptrow{1c0dd63c4cc8d828}{2026-07-31T22:28:56Z}{build+decorrelated}{\badge{BadgeGreen}{current}}
\receiptrow{47a3e55f8aa07fc7}{2026-08-02T23:31:07Z}{build+decorrelated}{\badge{BadgeGreen}{current}}
\receiptrow{5ada5a799af57e12}{2026-08-03T16:27:10Z}{build+decorrelated}{\badge{BadgeGreen}{current}}
\receiptrow{c2de6b6e1b9e676c}{2026-08-03T22:45:54Z}{build+decorrelated}{\badge{BadgeGreen}{current}}
\end{receipttable}
\sourcenote{2 further receipt(s) in the store attest superseded builds of this module and are omitted from this rendering.}
\sourcenote{module \texttt{BernsteinLogicAux.QtslCompose.Presentation}; bound build \texttt{efd8b570583d\ldots} (bound from receipt c2304cb9372c9e10, 2026-07-21T05:39:15Z). A receipt records that named checking layers accepted a build: (module, compiled-artifact SHA-256, environment SHA-256, checker layers), minted when the checking gate ran; the store receipts attesting the bound build are listed with the layers each run credited. Nothing here asks the reader to trust the store: every layer is an independently runnable program over the quoted sources --- rebuild the module and re-run the checkers to reproduce the verdict. Receipts attesting superseded builds credit no layer and are omitted from this rendering (the count of omissions, if any, is noted below).}
\begin{receipttable}
\receiptrow{0d24201fcc020c14}{2026-07-20T07:45:50Z}{build+decorrelated}{\badge{BadgeGreen}{current}}
\receiptrow{c6d195b51ddfe44f}{2026-07-21T00:04:00Z}{build+decorrelated}{\badge{BadgeGreen}{current}}
\receiptrow{c2304cb9372c9e10}{2026-07-21T05:39:15Z}{build+replay+hardened+decorrelated}{\badge{BadgeGreen}{current}}
\end{receipttable}
\sourcenote{11 further receipt(s) in the store attest superseded builds of this module and are omitted from this rendering.}
\sourcenote{module \texttt{BernsteinLogicAux.QtslCompose.AGCompose}; bound build \texttt{88a891d2d053\ldots} (bound from receipt b967c74819cd1516, 2026-07-21T05:39:06Z). A receipt records that named checking layers accepted a build: (module, compiled-artifact SHA-256, environment SHA-256, checker layers), minted when the checking gate ran; the store receipts attesting the bound build are listed with the layers each run credited. Nothing here asks the reader to trust the store: every layer is an independently runnable program over the quoted sources --- rebuild the module and re-run the checkers to reproduce the verdict. Receipts attesting superseded builds credit no layer and are omitted from this rendering (the count of omissions, if any, is noted below).}
\begin{receipttable}
\receiptrow{91f9e1cce3d69343}{2026-07-19T03:38:55Z}{build+decorrelated}{\badge{BadgeGreen}{current}}
\receiptrow{ccc58392dbd207a9}{2026-07-19T18:24:44Z}{build+decorrelated}{\badge{BadgeGreen}{current}}
\receiptrow{875fc85d5bcb76ca}{2026-07-20T01:21:56Z}{build+decorrelated}{\badge{BadgeGreen}{current}}
\receiptrow{23e1df54aafc0165}{2026-07-20T20:46:49Z}{build+decorrelated}{\badge{BadgeGreen}{current}}
\receiptrow{02e33e58a762ea11}{2026-07-20T20:58:22Z}{build+replay+hardened+decorrelated}{\badge{BadgeGreen}{current}}
\receiptrow{9690cee413e67035}{2026-07-21T00:03:48Z}{build+decorrelated}{\badge{BadgeGreen}{current}}
\receiptrow{b967c74819cd1516}{2026-07-21T05:39:06Z}{build+replay+hardened+decorrelated}{\badge{BadgeGreen}{current}}
\end{receipttable}
\sourcenote{6 further receipt(s) in the store attest superseded builds of this module and are omitted from this rendering.}
\sourcenote{module \texttt{BernsteinLogicAux.QtslCompose.AGComposeDiscrete}; bound build \texttt{d5233632f95a\ldots} (bound from receipt 88ff74180adf5b95, 2026-08-03T20:13:53Z). A receipt records that named checking layers accepted a build: (module, compiled-artifact SHA-256, environment SHA-256, checker layers), minted when the checking gate ran; the store receipts attesting the bound build are listed with the layers each run credited. Nothing here asks the reader to trust the store: every layer is an independently runnable program over the quoted sources --- rebuild the module and re-run the checkers to reproduce the verdict. Receipts attesting superseded builds credit no layer and are omitted from this rendering (the count of omissions, if any, is noted below).}
\begin{receipttable}
\receiptrow{88ff74180adf5b95}{2026-08-03T20:13:53Z}{build+decorrelated}{\badge{BadgeGreen}{current}}
\end{receipttable}
\sourcenote{11 further receipt(s) in the store attest superseded builds of this module and are omitted from this rendering.}
\layerline{\layeritem{build}{lake build (elaboration)}{\badge{BadgeGreen}{current}}\layeritem{replay}{leanchecker (official kernel replay)}{\badge{BadgeGreen}{current}}\layeritem{hardened}{safe\_verify (hardened screen)}{\badge{BadgeGreen}{current}}\layeritem{decorrelated}{lean4export + nanoda (decorrelated kernel)}{\badge{BadgeGreen}{current}}}
\sourcenote{\badge{BadgeGray}{comparator: none} no fixture in the three-implementation comparator harness (compiled / kernel-interpreted / reference-interpreted) exercises this rule.}
\ruleblock{HOL4 statement}
\absentbox{NO HOL STATEMENT YET}{The satisfaction-judgment extension postdates the HOL4 port entirely:
the ported cone is the base calculus\textquotesingle s first-order
check-function surface (theory \texttt{bern\_catalogue}, stopping
before even the base proof-object layer), and no extension-theory
definition exists anywhere in the HOL development. A HOL statement of
this rule presupposes a HOL port of the satisfaction stratum, which is
its own campaign; until then this box is the honest state of the
edifice.
}
\ruleblock{HOL verifier receipts}
\absentbox{NONE}{no HOL verification receipt exists for this rule.}
\ruleblock{Proof}

\sourcenote{Genre: core rule (tier 1), with a per-profile tier split
derived from the presentation. The coverage is deliberately
conditional: this section proves the cells the mechanization proves,
names the bridge hypotheses under which they apply, and states what
remains open --- never claiming past the mechanized grid.}

\textbf{Setting and claim.}\enspace A \emph{system} is a structured
presentation: finitely many \emph{mechanisms} (primitive dynamics
declarations --- drift, diffusion, jump, and guard--reset field
families over typed wires), plus \emph{wirings} (variable
identifications) carrying a \emph{causality certificate} --- a rank
function strictly increasing along wires, so a directed cycle is
unrepresentable in the assembled structure. Each mechanism's
\emph{doctrine profile} is derived from which field families are
populated: drift only $=$ ODE; $+$diffusion $=$ SDE; $+$real-time
jumps $=$ general stochastic hybrid (GSHS); counter-time jumps only
$=$ discrete kernel. An interface-indexed satisfaction judgment
$\mathsf{QSatI}(\Sigma, (I, \alpha, \gamma[, \kappa]))$ asserts: every
behaviour of $\Sigma$ whose input trajectory satisfies the assumption
$\alpha$ pointwise delivers the guarantee $\gamma$ (with quantitative
margin per the slack $\kappa$ where declared). A wiring certificate
$\mathsf{QWireCert}(w, (\mathit{spec}_i) \Rightarrow
\mathit{spec}_{\mathrm{out}}, \kappa)$ asserts the lens square
\cite{capuccimyers2026}: delivered upstream guarantees, degraded by at
most $\kappa$, imply downstream assumptions through $w$, and the
output spec reads off the delivered guarantees (the square's discharge
is \rulename{QWireCertDischarge}'s business). The rule: from
$\Sigma_i \vDash \mathit{spec}_i$ for $i = 1, \dots, n$ and
the wiring certificate, conclude
$\mathsf{SysWire}(w, [\Sigma_1, \ldots, \Sigma_n]) \vDash
\mathit{spec}_{\mathrm{out}}$.

The warrant factors as
$\mathrm{sound} = \text{(T)} \circ \text{(S)}$: \emph{transport
arithmetic} --- that composite obligation \emph{rows} follow from
component rows plus wiring rows --- which is profile-generic exact
polynomial arithmetic; and the \emph{semantic reading} --- that the
checked row family implies the behavioural claim --- which is
per-profile, and is where the tier split and every open cell live.

\begin{blemma}[transport is profile-silent]\label{lem:qag-transport}
With residuals $e_i := \mathcal{L}_i V_i + \alpha_i - \sigma_i$ and
membership rows $m_1 := \lambda\sigma_2 - (1{-}\theta)\alpha_1$,
$m_2 := \sigma_1 - \lambda(1{-}\theta)\alpha_2$,
\[
e_1 + \lambda e_2 + m_1 + m_2
\;=\;
\mathcal{L}(V_1 + \lambda V_2) + \theta\,(\alpha_1 + \lambda\alpha_2),
\]
an identity of polynomials --- the cross-supplies cancel by
construction of the multipliers. More generally, for nonnegative
certificate weights $\lambda_i$ and free-sign wiring multipliers
$\nu_k$: if each $e_i \leq 0$ and each wiring identity $g_k = 0$ holds
on the operating region, then
$\sum_i \lambda_i e_i + \sum_k \nu_k g_k \leq 0$ there.
\end{blemma}
\begin{proof}
The displayed identity is a two-line expansion; the only property of
the obligation operator $\mathcal{L}$ it consumes is \emph{linearity
over the wiring} --- which every profile's operator has:
$\nabla(\cdot) \cdot f$ for ODE, one-step conditional expectation
minus identity for kernels, and the extended generator's interior
terms for GSHS. The pointwise fold is
\texttt{admissible\_constructs}: each $\lambda_i e_i \leq 0$
(nonnegative times nonpositive), each $\nu_k g_k = 0$ regardless of
sign, finite sums of nonpositives are nonpositive. At the sequent
layer this is the \rulename{Promise}/\rulename{Syzygy} machinery, and
the small-gain instance is exactly the $(1, \lambda, 1, 1)$-weighted
\rulename{Syzygy} node --- which is why the split of
\autoref{rem:qag-tier} can define ``admissible'' as
\emph{syntactically constructible}: the composite certificate for the
non-stochastic profiles \emph{is} $\sum_i \lambda_i V_i$, and its row
is this lemma. The break in uniformity, wherever it appears below, is
never here: GSHS forced jumps contribute boundary rows, not
$\mathcal{L}$-terms, so part of the composite obligation family
bypasses linearity entirely and must be handled semantically.
\end{proof}

\begin{blemma}[conformance propagation --- the acyclic engine]
\label{lem:qag-conform}
Let a system presentation $S$ have $n$ components, and work in the
composite behaviour space. Write $\mathrm{conform}_j$ for the event
that component $j$'s inputs conform to $\alpha_j$, $\mathrm{sat}_j$
for delivery of $\gamma_j$, and $\mathrm{ext}$ for external-input
conformance. Suppose the wiring squares transport, pointwise:
$\mathrm{ext} \cap \bigcap_{i \in \mathrm{upstream}(j)} \mathrm{sat}_i
\subseteq \mathrm{conform}_j$ for every $j$. If at a behaviour $\omega
\in \mathrm{ext}$ every component delivers on conformance
($\omega \in \mathrm{conform}_j \Rightarrow \omega \in
\mathrm{sat}_j$), then every component delivers at $\omega$.
\end{blemma}
\begin{proof}
Strong induction on the causality rank (\texttt{sat\_of\_conform}):
if every component of rank $< N$ delivers, then for $j$ of rank $< N{+}1$
all of $j$'s upstream components have strictly smaller rank (the
causality certificate), hence deliver; the wiring square then puts
$\omega \in \mathrm{conform}_j$, and the per-component implication
closes. The causality certificate is precisely what makes the
induction well-founded, and its necessity is a two-component
counterexample: components in a cycle, each delivering \emph{only if}
the other does, satisfy every wiring square while both fail --- the
implication chain has no base case. That failure is the classical
unsoundness of naive circular assume--guarantee
\cite{abadilamport1993, mcmillan1999}; this calculus fences it
structurally (cycles unrepresentable here) and readmits feedback only
through the guarded small-gain cells of
Lemma~\ref{lem:qag-feedback}, where the well-foundedness device is the
$\theta$-margin.
\end{proof}

\begin{blemma}[failure decomposition and slack composition]
\label{lem:qag-union}
Under Lemma~\ref{lem:qag-conform}'s wiring hypothesis, a composite
failure under conforming external inputs localizes at some component
that failed \emph{while conforming}:
\[
\mathrm{ext} \cap \Bigl(\bigcap_j \mathrm{sat}_j\Bigr)^{\mathsf{c}}
\;\subseteq\;
\bigcup_j \bigl(\mathrm{conform}_j \cap \mathrm{sat}_j^{\mathsf{c}}\bigr),
\]
hence for any measure with per-component conforming-failure bounds
$\mu(\mathrm{conform}_j \cap \mathrm{sat}_j^{\mathsf{c}}) \leq
\varepsilon_j$, the composite fails with mass at most
$\sum_j \varepsilon_j$ --- the $\sum_j (1 - p_j)$ display, at any $n$.
\end{blemma}
\begin{proof}
Pick a \emph{rank-minimal} failing component $j_m$ (finitely many
components; the failing set is nonempty). Every upstream component of
$j_m$ has strictly smaller rank, hence delivers (minimality), so the
wiring square puts $\omega \in \mathrm{conform}_{j_m}$: a conforming
failure (\texttt{fail\_subset}). Monotonicity and finite subadditivity
of the measure close (\texttt{fail\_bound}; read at an output spec,
\texttt{fail\_bound\_out}). The per-component event
$\mathrm{conform}_j \cap \mathrm{sat}_j^{\mathsf{c}}$ --- not bare
failure --- is forced: component $j$'s certificate priced failure
\emph{only under conforming inputs}; charging it for
non-conforming failures would demand a guarantee nobody certified,
and the rank-minimality argument is exactly what shows the composite
never needs one.
\end{proof}

\begin{blemma}[the stochastic-hybrid cells: series, parallel,
feedback sum-form]\label{lem:qag-feedback}
Over the composite execution law of a GSHS composite, with component
certificates entering as nonnegative supermartingales along the
(stopped) joint jump skeleton, with barrier values on their bad sets
and violation-witness inclusions:
\emph{(i, series)} the conditioning tower composes --- qualitatively,
delivery a.e.\ on conforming inputs chains through the wiring square;
quantitatively, an upstream success mass $\geq p_1$ and a downstream
conditional success probability $\geq p_2$ --- the conditional bound
holding almost everywhere \emph{on the upstream success event}, itself
measurable in the conditioning $\sigma$-algebra --- compose to
composite success mass $\geq p_2 \cdot p_1$ (stronger than the
union-bound shape by exactly $(1-p_1)(1-p_2) \ge 0$);
and the certified pipeline bounds the composite conforming-failure
mass by $\int Y_1(0) / a_1 + \int Y_2(0) / a_2$ through the maximal
inequality.
\emph{(ii, parallel)} for nonnegative weights, the weighted sum
$\sum_i \lambda_i V_i$ of supermartingales w.r.t.\ one (joint)
filtration is a supermartingale --- the semantic face of the syzygy
weights --- and the composite avoid bound follows through the same
maximal inequality.
\emph{(iii, feedback, sum-form)} for two components in feedback whose
certificates satisfy, along the \emph{joint} skeleton, the
between-jumps conditional decrease (pre-jump value plus integrated
assumption-minus-supply increment conditionally below the current
value) and \emph{non-expansive} jump accounting (post-jump value a.e.\
at most pre-jump value), with the integrated membership rows of the
small-gain pair at $\lambda, \theta \geq 0$: the composite
$W = V_1 + \lambda V_2$ is a nonnegative supermartingale for the
joint-skeleton filtration, with the $\theta$-margin banked by the
compensated process, and the composite avoid bound follows.
\end{blemma}
\begin{proof}
(i) is \texttt{AGCompose.Series}: the qualitative form is an a.e.\
implication chain through the pointwise wiring square
(\texttt{qualitative}); the product form integrates the downstream
conditional bound over the upstream event --- the set-integral of the
conditional expectation of the indicator over an upstream-measurable
set is the joint mass (\texttt{tower\_step}, \texttt{tower}); the
certified form (\texttt{certified}) converts each certificate to a
failure bound via the maximal inequality for nonnegative
supermartingales (Ville's inequality, mechanized in the accompanying
development) and composes by the failure decomposition. Exogeneity is
where acyclicity earns its keep semantically: in a cascade each
component's input process is measurable w.r.t.\ upstream trajectories,
so the component's certificate argument runs pathwise in its input.

(ii) is \texttt{AGCompose.Parallel}: induction on the finite index
set, scalar-multiple and sum closure of supermartingales
(\texttt{supermartingale\_weighted\_sum}), then the barrier argument
as in (i) (\texttt{certified}, \texttt{certified\_pair}).

(iii) is the file quoted above (\texttt{SumFormFeedback}): folding
jump non-expansion into the between-jumps decrease gives the
per-component skeleton step; the weighted combination of the two steps
cancels the cross-supplies \emph{under conditional expectation} ---
Lemma~\ref{lem:qag-transport} run pointwise inside the conditional ---
leaving the $\theta$-weighted assumption margin (\texttt{drift});
nonnegativity and the supermartingale property assemble
(\texttt{supermartingale}), and the avoid bound is the maximal
inequality again (\texttt{avoid\_bound}, with the continuous-time
reading conditional on the d\'ebut-sampling inclusion,
\texttt{debut\_avoid\_bound}). Two hypotheses deserve their remark.
\emph{The joint skeleton is not a convenience}: in feedback no
component's input is exogenous --- the jump chain interleaves both
components' guard hits, each component's between-jumps argument is
valid only between \emph{joint} jumps, and forced jumps, having no
intensity, enter only through the jump-accounting field at those
times. \emph{Non-expansion is the sum-form's load-bearing wall}: only
$\lambda, \theta \geq 0$ and the non-expansive accounting are consumed
--- no dwell-time or schedule hypothesis appears anywhere in the
descent, which is a checkable structural fact about the statements ---
and the rule-level ranges $\lambda > 0$, $\theta \in (0,1)$ are
well-formedness backing the published guard-modulus reading, not
soundness needs.
\end{proof}

\begin{bthm}[the $n$-ary acyclic assembly]\label{thm:qag-nary}
Let $S$ be an acyclic system presentation with wiring squares
transporting pointwise as in Lemma~\ref{lem:qag-conform}. If each
component delivers a.e.\ on conformance, the composite delivers its
output spec a.e.\ on external conformance
(\texttt{QtslCompose.qualitative}); and with per-component
conforming-failure bounds $\varepsilon_j$, the composite output-spec
failure mass is at most $\sum_j \varepsilon_j$
(\texttt{fail\_bound\_out}) --- at any $n$, by
Lemmas~\ref{lem:qag-conform} and~\ref{lem:qag-union}.
\end{bthm}
\begin{proof}
Immediate from the two lemmas: the a.e.\ form collects the countable
family of per-component a.e.\ implications, applies the pointwise
propagation on the common full-measure set, and reads the output spec
off the delivered guarantees; the probabilistic form is the failure
decomposition plus subadditivity. The per-profile content enters
entirely through the \emph{bridge hypotheses} --- that the
component-certificate events and their bounds exist over the composite
law is the per-profile discharge of Lemma~\ref{lem:qag-feedback} (GSHS
acyclic cells) or its lighter analogues; this theorem is the
profile-generic event algebra on top.
\end{proof}

\begin{bremark}[the tier split, derived --- never a stored tag]
\label{rem:qag-tier}
The rule node's tier is a function of the presentation
(\texttt{System.qagTier}): \emph{admissible} iff every mechanism's
derived profile admits a syntactically constructible composite
certificate --- proven equivalent to ``every mechanism derives ODE or
discrete'' (\texttt{qagTier\_admissible\_iff}) --- and \emph{core} iff
some mechanism is stochastic (\texttt{qagTier\_core\_iff}). On the
admissible face the node unpacks: the composite certificate is
$\sum_i \lambda_i V_i$ (Lemma~\ref{lem:qag-transport}), transported at
the sequent layer by the \rulename{Promise}/\rulename{Syzygy}
machinery, and the rule adds no semantic content. On the core face no
unpacking exists --- the warrant \emph{is}
Lemmas~\ref{lem:qag-conform}--\ref{lem:qag-feedback} --- which is the
taxonomy's definition of core applied honestly, one rule node either
way. Deriving the tier from field-family population (rather than
storing a system-level tag) is forced by mixed composites: a
sampled-data plant-plus-controller would be ill-typed across its own
components under any stored per-system tag.
\end{bremark}

\begin{bremark}[what is \emph{not} proven --- the open cells, loud]
\label{rem:qag-gaps}
The conditional mark on this rule's dashboard row warrants exactly the
cells above, and the honest remainder is:
\emph{(a)} the deterministic-ODE row's semantic reading is at
statement level (its literature anchor is the certified-lens
composition theorem \cite{capuccimyers2026}; the classical comparison
argument is standard but not yet written in-house);
\emph{(b)} the discrete-kernel row's assembly is discharged ---
$n$-ary acyclic union-bound, parallel weighted-sum, and sum-form
small-gain, end-to-end over the trajectory kernel (the width note's
dictionary names the statements) --- with the \emph{sampled-data} row
the one still stated, not yet assembled (item (d));
\emph{(c)} feedback in \emph{gauge-form} --- jumps that expand the
certificate ($V(g(x)) \leq c\,V(x)$, $c > 1$), repaid by flow decay
under schedule hypotheses --- is partial: the gauge-then-sum reduction
and the machine-checked max-form piecewise-deterministic instance are
landed sorry-free, while the finite-index price of the diffusion term
and the admissible-window pinch remain open research --- though the
reason the diffusion cell was \emph{parked} is retired: within-segment
supermartingale decay for gauged certificates is now available at the
continuous-time martingale-problem profile (the battery
family\textquotesingle s discharge), so what remains is the finite-index
accounting and the window pinch themselves, not a missing stochastic
calculus; sum-form
composites provably do not cover this regime, which is exactly why the
core tier exists. Non-expansive-reset models live entirely in
Lemma~\ref{lem:qag-feedback}(iii), so nothing below gates on (c);
\emph{(d)} the mixed sampled-data row is stated, not proven ---
deterministic throughout, the natural first boundary-row exercise;
\emph{(e)} undirected stochastic wire sharing (common noise) is
\emph{fenced}, not missing: stochastic components sharing a wire need
one determining side, and for fenced models no discharge of the
joint-filtration hypothesis exists --- the lemmas are correctly
unusable there;
\emph{(f)} the $n$-ary \emph{product-form} (tower) composition is
binary only --- its $n$-ary form needs a filtration chain along the
wiring order and is real further work;
\emph{(g)} the judgment sorts the rule concludes into are not yet
staged in the certificate schema, so no end-to-end checkable
certificate exercises the rule yet. Each bridge hypothesis
(the interface layers carrying a component certificate onto the
composite law) is a named hypothesis discharged per model class ---
parameters, never silent assumptions.
\end{bremark}

\begin{bremark}[relation to the literature]
The acyclic discipline is the composing-specifications tradition
\cite{abadilamport1993}, with circularity readmitted only under a
well-foundedness device \cite{mcmillan1999} --- here the
$\theta$-margin of the small-gain pair, whose Lyapunov formulation
descends from the nonlinear small-gain theorems
\cite{jiangteelpraly1994, jiangmareelswang1996}. The parallel cell is
compositional dissipativity \cite{willems1972, moylanhill1978} run at
the satisfaction layer --- the weighted-sum supermartingale is the
stochastic face of the storage-function sum. The wiring square, its
slack algebra, and the certified-lens composition law are
\cite{capuccimyers2026}. What the mechanized grid adds over all of
these is the certificate form: exact-arithmetic transport rows plus
per-profile semantic lemmas with their hypotheses reified as checkable
data, rather than a metatheorem quantified over models nobody
replays.
\end{bremark}

\begin{anchors}
Lemma \ref{lem:qag-transport}: cancellation, pointwise fold & \texttt{admissible\_constructs}; sequent layer \texttt{sound\_promise}, \texttt{sound\_syzygy}, \texttt{sound\_smallGain} (main package)\\
Lemma \ref{lem:qag-conform}: rank induction & \texttt{sat\_of\_conform}; \texttt{System} (fields \texttt{rank}, \texttt{causality})\\
Lemma \ref{lem:qag-union}: rank-minimal localization, union bound & \texttt{fail\_subset}, \texttt{fail\_bound}, \texttt{fail\_bound\_out}\\
Lemma \ref{lem:qag-feedback}(i): series & \texttt{AGCompose.Series.qualitative}, \texttt{tower\_step}, \texttt{tower}, \texttt{fail\_bound}, \texttt{certified}\\
Lemma \ref{lem:qag-feedback}(ii): parallel & \texttt{AGCompose.Parallel.supermartingale\_weighted\_sum}, \texttt{certified}, \texttt{certified\_pair}\\
Lemma \ref{lem:qag-feedback}(iii): feedback sum-form & \texttt{SkeletonCertificate} (fields \texttt{dynkin}, \texttt{jump\_nonexpansive}), \texttt{SkeletonCertificate.step}, \texttt{SumFormFeedback.drift}, \texttt{supermartingale}, \texttt{avoid\_bound}, \texttt{debut\_avoid\_bound}, \texttt{compensated\_supermartingale}, \texttt{margin\_budget}\\
gauge-form partial (max-form instance) & \texttt{sound\_GSHS\_safety\_barrier} (MaxFormPDMP)\\
Theorem \ref{thm:qag-nary} & \texttt{QtslCompose.qualitative}, \texttt{fail\_bound\_out}\\
Remark \ref{rem:qag-tier}: derived split & \texttt{System.qagTier}, \texttt{qagTier\_admissible\_iff}, \texttt{qagTier\_core\_iff}; \texttt{ChartShape.profile}\\
\end{anchors}


\rulesection{79}{QWireCertDischarge}{core rule}
\label{blcatrule-QWireCertDischarge}
\ruleblock{The rule (Steele metanotation)}
\sourcenote{as realized in the mechanized development; source row: \texttt{\texttt{QTSL-rules.toml}} --- \emph{QWireCertDischarge}. The metanotation statement is typeset by hand for notational quality; its mechanized ground truth is the Lean statement below, extracted from the sources at every build of this book and hash-bound. The rule \emph{roster} is machine-checked (every mechanized rule tag round-trips to exactly one catalogue row); the statement \emph{text} is editorial, checked by eye against the extracted Lean beside it.}
\begin{mathpar}
\inferrule*[Right=\rulename{QWireCertDischarge}]{
  \bigl\{R_\ell,\ \varphi(s) \wedge \alpha(v(s), a) \vdash \mathrm{Lift}_s\, \varphi(u(s, a))\bigr\}_{\ell},\ \ \bigl\{R_\ell,\ \varphi(s) \vdash \gamma(v(s))\bigr\}_{\ell} \ \text{(mechanism-certification legs; } s \ \text{a mechanism state, } a \ \text{an admissible input)} \\
  \bigl\{R_\ell \vdash \alpha_2(w(o_1), a_2) + \kappa(\gamma_1(o_1)) \ge \alpha_1(w^{\sharp}(o_1, a_2))\bigr\}_{\ell} \ \text{(assumption legs)} \\
  \bigl\{R_\ell \vdash \gamma_1(o_1) \ge \gamma_2(w(o_1))\bigr\}_{\ell} \ \text{(guarantee legs)} \\
  \vdash \kappa(0) = 0 \ \text{(closed evaluation)} \\
  [0, \infty) \vdash \kappa' \ge 0 \ \text{(derivative-nonnegativity row)} \\
  \vdash \deg \kappa > 0,\ \mathrm{lead}(\kappa) \ge 0 \ \text{(leading-coefficient row)}
}{
  \mathsf{QWireCert}\bigl(w,\ (\mathit{spec}_i) \Rightarrow \mathit{spec}_{\mathrm{out}},\ \kappa\bigr)
}
\end{mathpar}
every leg family indexes over its posted proof leaves $\ell$, and each leaf binds its own region $R_\ell$ as the antecedent of its sequent; the three arithmetic rows certify $\kappa \in \mathcal{K}^0_\infty$ (zero at zero, monotone on the ray, unbounded).

\ruleflavor{$\mathcal{K}^0_\infty$ is the class under which slacks add and lenses compose; the premises are pointwise implications over polynomial data, one inherited BL sequent per posted leaf, the leaf\textquotesingle s antecedent carrying its region --- no region cover is stored on the node, and the obligation that the posted regions jointly cover the consuming composition\textquotesingle s interface domain lives one row up, at the consumption site. The two Capucci--Myers condition families are the mechanism-certification legs (n\'e chart certification) and the wiring-lens legs. No new checker machinery: the rule is leaf assembly, and its semantic content lives one row up in the composition rule.}
\par\medskip
\rulename{QWireCertDischarge} is the wiring-certificate
judgment\textquotesingle s own discharge: it assembles a
$\mathsf{QWireCert}$ --- the license \rulename{QWireCompose} consumes
--- entirely from inherited BL sequent leaves. Four condition families
are checked, each indexing over its posted proof leaves $\ell$, and
each leaf binds its own region $R_\ell$ as the antecedent of its own
sequent. The two \emph{mechanism-certification legs} certify the
component\textquotesingle s contract at its mechanism:
$R_\ell,\ \varphi(s) \wedge \alpha(v(s), a) \vdash
\mathrm{Lift}_s\, \varphi(u(s, a))$ --- the certified invariant,
together with the assumption evaluated at the mechanism\textquotesingle
s readout and an admissible input, is preserved by the update --- and
$R_\ell,\ \varphi(s) \vdash \gamma(v(s))$ --- the invariant delivers
the guarantee at the readout --- where $s$ ranges over the
mechanism\textquotesingle s state carrier and $a$ over its admissible
input carrier. The two \emph{lens legs} certify the wiring square: the
\emph{assumption leg} (delivered upstream guarantees, degraded by at
most the slack $\kappa$, imply the downstream assumption through the
wiring map) and the \emph{guarantee leg} (the upstream guarantee
dominates the composite\textquotesingle s output reading). For
polynomial data all four are pointwise implications --- ordinary
sequents the base calculus
already checks. Three further arithmetic rows certify the
slack\textquotesingle s class membership,
$\kappa \in \mathcal{K}^0_\infty$: the constant coefficient is zero
(closed evaluation), the derivative is nonnegative on the ray (a
derivative-nonnegativity row), and the leading coefficient is positive
at positive degree.

What the cited theorems warrant --- and the scope of this
entry\textquotesingle s conditional check: that the three row families
genuinely certify class membership (a polynomial passing them is zero
at zero, monotone on $[0, \infty)$, and unbounded), and that the class
supports the composition rule\textquotesingle s slack algebra ---
slacks are closed under pointwise sum, closed under composition, the
two-slack contract chain composes to the summed slack, and the
guarantee legs chain along a cascade. That algebra is the
quantitative content of \textquotedblleft compose by lens composition
and by addition of the slacks\textquotedblright.

Each membership row is load-bearing for the algebra. Drop zero-at-zero
and a constant slack certifies a square that degrades even exact
guarantee delivery --- and the nonnegativity that anchors the additive
algebra (a class slack is nonnegative on the ray \emph{because} it is
zero at zero and monotone) fails with it. Drop monotonicity and
composition of slacks fails outright: the outer slack of a two-lens
cascade need no longer bound the chained degradation, since a
non-monotone $\kappa_1$ can reward a worse input. Drop unboundedness
and bounded candidates enter ($\kappa \equiv 0$ passes the other two
rows): the discharge then certifies membership in a class the slack
does not belong to, and the published reading of $\kappa$ as a
reusable comparison-function-style gain is lost. The proof section
walks each closure with the field it consumes.

\cataloguenote{Core-as-glue: the rule is leaf assembly, and the semantic content of
what a wiring certificate \emph{buys} lives one row up, in
\rulename{QWireCompose}\textquotesingle s semantic reading. The
non-strict monotone reading of the class follows the
catalogue\textquotesingle s own row roster (a derivative-nonnegativity
row certifies monotone, not strictly increasing); the zero slack is
thereby not a member --- a slack-free square is a square with the
$\kappa$-term absent, not with $\kappa = 0$. The judgment sort itself
is baked in the extension theory; the mechanized warrant was
proven ahead of the bake, and the square leaves are landed
main-package sequent machinery, cited rather than re-bound here.
}
\ruleblock{Lean statement}
\par\medskip
\textbf{The slack class, bound by display.}\enspace
\texttt{IsClassK0Inf} is the catalogue\textquotesingle s
$\mathcal{K}^0_\infty$: zero at zero, monotone on the nonnegative ray,
unbounded at infinity. Its closure theorems (\texttt{add} --- slacks
add; \texttt{comp} --- lenses compose) are proven in the same module
inside the class namespace; \texttt{polynomial\_discharge} maps the
three checked row families onto the three fields, and
\texttt{slack\_triangle}/\texttt{guarantee\_leg\_comp} are the
quantitative and guarantee-side composition faces.

\declhead{BernsteinLogicAux.QtslCompose.WireCertDischarge}{polynomial\_discharge}{\badge{BadgeGreen}{statement bound}}
\begin{leansrc}
/-- **The polynomial discharge** — the rule note's three membership rows, each
mapped to the class field it certifies:
1. zero-at-zero = closed evaluation (`p.coeff 0 = 0`, checked by evaluating);
2. monotone = the derivative-nonnegativity row (`p' ≥ 0` on the ray);
3. unbounded = positive leading coefficient (with positive degree).
A polynomial passing the three BL row families is genuinely a 𝒦⁰∞ slack. -/
theorem polynomial_discharge (p : Polynomial ℝ)
    (h0 : p.coeff 0 = 0)
    (hderiv : ∀ x : ℝ, 0 ≤ x → 0 ≤ p.derivative.eval x)
    (hdeg : 0 < p.degree) (hlead : 0 ≤ p.leadingCoeff) :
    IsClassK0Inf (fun x => p.eval x)
\end{leansrc}
\declhead{BernsteinLogicAux.QtslCompose.WireCertDischarge}{slack\_triangle}{\badge{BadgeGreen}{statement bound}}
\begin{leansrc}
/-- **The additive slack triangle** (the quantitative face of "compose by
lens composition and by addition of the slacks"): two contract inequalities
with slack compose to one with the summed slack. Stated pointwise over an
arbitrary interface carrier — the per-region reading is `X` = the region. -/
theorem slack_triangle {X : Type*} {f g h : X → ℝ} {c₁ c₂ : ℝ}
    (h₁ : ∀ x, f x ≤ g x + c₁) (h₂ : ∀ x, g x ≤ h x + c₂) :
    ∀ x, f x ≤ h x + (c₁ + c₂)
\end{leansrc}
\declhead{BernsteinLogicAux.QtslCompose.WireCertDischarge}{guarantee\_leg\_comp}{\badge{BadgeGreen}{statement bound}}
\begin{leansrc}
/-- The guarantee-side lens leg composes (the κ-free leg of the CM square:
`γ₁ ≥ γ₂ ∘ w` chains along a cascade). -/
theorem guarantee_leg_comp {A B C : Type*} {w₁ : A → B} {w₂ : B → C}
    {γ₁ : A → ℝ} {γ₂ : B → ℝ} {γ₃ : C → ℝ}
    (h₁ : ∀ a, γ₂ (w₁ a) ≤ γ₁ a) (h₂ : ∀ b, γ₃ (w₂ b) ≤ γ₂ b) :
    ∀ a, γ₃ (w₂ (w₁ a)) ≤ γ₁ a
\end{leansrc}
\declhead{BernsteinLogicAux.QtslCompose.WireCertDischarge}{nonneg}{\badge{BadgeGreen}{statement bound}}
\begin{leansrc}
/-- A 𝒦⁰∞ slack is nonnegative on the nonnegative ray. -/
theorem nonneg (h : IsClassK0Inf κ) {x : ℝ} (hx : 0 ≤ x) : 0 ≤ κ x
\end{leansrc}
\declhead{BernsteinLogicAux.QtslCompose.WireCertDischarge}{add}{\badge{BadgeGreen}{statement bound}}
\begin{leansrc}
/-- **Slacks add** ("by addition of the slacks"): the class is closed under
pointwise sum. -/
theorem add (h₁ : IsClassK0Inf κ₁) (h₂ : IsClassK0Inf κ₂) :
    IsClassK0Inf (fun x => κ₁ x + κ₂ x)
\end{leansrc}
\declhead{BernsteinLogicAux.QtslCompose.WireCertDischarge}{comp}{\badge{BadgeGreen}{statement bound}}
\begin{leansrc}
/-- **Lenses compose**: the class is closed under composition (the slack of a
two-lens cascade). -/
theorem comp (h₁ : IsClassK0Inf κ₁) (h₂ : IsClassK0Inf κ₂) :
    IsClassK0Inf (κ₁ ∘ κ₂)
\end{leansrc}
\declhead{BernsteinLogicAux.QtslCompose.WireCertDischarge}{IsClassK0Inf}{\badge{BadgeGreen}{definition bound}}
\begin{leansrc}
/-- Class 𝒦⁰∞ membership (the `QWireCert` slack class): zero at zero,
monotone on the nonnegative ray, unbounded at infinity. -/
structure IsClassK0Inf (κ : ℝ → ℝ) : Prop where
  zero_at_zero : κ 0 = 0
  mono : MonotoneOn κ (Ici 0)
  unbounded : Tendsto κ atTop atTop
\end{leansrc}
\ruleblock{Correspondence (Steele display $\leftrightarrow$ Lean)}
\badge{BadgeGreen}{reviewer-certified against the displayed texts}
\sourcenote{The dictionary below is a reviewer-certified statement of how the mechanized statements above interpret the displayed rule: display metavariable $\mapsto$ Lean binder, each display hypothesis $\mapsto$ the premise or side condition carrying it, the display conclusion $\mapsto$ the goal --- with any deliberate difference in width between display and mechanization named outright. Certification is recorded as a hash triple over (the displayed cell, the extracted statement texts, this dictionary\textquotesingle s own text); if any of the three changes, this block turns amber until a reviewer re-certifies. The correspondence claim itself remains editorial --- reviewer-grade, not machine-checked; its machine-checked successor types each soundness statement against the rule as a structured object, and supersedes this block per rule as it lands.}
The displayed conclusion $\mathsf{QWireCert}(w, (\mathit{spec}_i)
\Rightarrow \mathit{spec}_{\mathrm{out}}, \kappa)$ is concluded by
\emph{no} bound statement --- deliberately: the judgment sort lives
theory-side (the note), the square leaves are landed sequent machinery
cited as witness prose, and what this bound set certifies is the
\emph{semantic warrant behind the arithmetic rows} plus the class
algebra the rule\textquotesingle s composition claims lean on. The
class itself is bound by display: \texttt{IsClassK0Inf} =
zero-at-zero + \texttt{MonotoneOn κ (Ici 0)} + \texttt{Tendsto κ
atTop atTop} --- $\mathcal{K}^0_\infty$ as the catalogue reads it,
non-strict monotone (the note\textquotesingle s ruling, visible in
the field type).
\begin{itemize}
\item the three membership rows $\mapsto$
  \texttt{polynomial\_discharge}, one hypothesis per row per field:
  closed evaluation $\mapsto$ \texttt{p.coeff 0 = 0} $\Rightarrow$
  \texttt{zero\_at\_zero}; derivative-nonnegativity on the ray
  $\mapsto$ \texttt{hderiv} $\Rightarrow$ \texttt{mono} (through the
  mean-value machinery \texttt{monotoneOn\_of\_deriv\_nonneg});
  leading-coefficient row $\mapsto$ \texttt{0 < p.degree} +
  \texttt{0 ≤ p.leadingCoeff} $\Rightarrow$ \texttt{unbounded}. The
  displayed row asks only $\mathrm{lead}(\kappa) \ge 0$ and so does
  the statement --- nonstrict suffices because a degree-witnessing
  leading coefficient is nonzero by definition; the
  commentary\textquotesingle s ``positive'' is the derived fact, not
  an extra demand.
\item ``slacks add, lenses compose'' $\mapsto$
  \texttt{slack\_triangle} (two contract inequalities chain to the
  summed slack, stated pointwise over an arbitrary carrier --- the
  per-region reading is the carrier choice) and
  \texttt{guarantee\_leg\_comp} (the $\kappa$-free guarantee leg
  chains along a cascade). The class \emph{closure} lemmas
  \texttt{IsClassK0Inf.add} and \texttt{IsClassK0Inf.comp} (with the
  nonnegativity anchor \texttt{IsClassK0Inf.nonneg}) are bound with
  the entry and carry the commentary\textquotesingle s closure claims
  (sum, composition), alongside the triangle and the guarantee leg.
\item polynomiality of $\kappa$ $\mapsto$ \texttt{p : Polynomial ℝ}
  in the discharge theorem only; the class and its algebra are stated
  for arbitrary \texttt{ℝ → ℝ} functions --- wider than the
  display\textquotesingle s polynomial slacks, safe direction.
\end{itemize}
Width note: the four leaf families of the display map onto the baked
node\textquotesingle s two proof lists --- the mechanism-certification
legs (both shapes) ride \texttt{QWireCertDischarge/mech\_rows}, and
the assumption and guarantee lens legs ride
\texttt{QWireCertDischarge/lens\_rows} --- with every leaf\textquotesingle
s own $R_\ell$ as its sequent\textquotesingle s antecedent. In the
bound set the lens legs appear
only through their compositional shadows (\texttt{slack\_triangle},
\texttt{guarantee\_leg\_comp}); the per-leaf discharge of all four
families is
inherited BL sequents, main package, outside this page\textquotesingle s
bindings --- the package-boundary rule, stated in the module header.

\ruleblock{Lean verifier receipts}
\sourcenote{module \texttt{BernsteinLogicAux.QtslCompose.WireCertDischarge}; bound build \texttt{cec5cc1749e1\ldots} (bound from receipt 90dcbae3ace744d7, 2026-07-19T18:24:58Z). A receipt records that named checking layers accepted a build: (module, compiled-artifact SHA-256, environment SHA-256, checker layers), minted when the checking gate ran; the store receipts attesting the bound build are listed with the layers each run credited. Nothing here asks the reader to trust the store: every layer is an independently runnable program over the quoted sources --- rebuild the module and re-run the checkers to reproduce the verdict. Receipts attesting superseded builds credit no layer and are omitted from this rendering (the count of omissions, if any, is noted below).}
\begin{receipttable}
\receiptrow{7aeb2249883cf96c}{2026-07-19T03:38:57Z}{build+decorrelated}{\badge{BadgeGreen}{current}}
\receiptrow{90dcbae3ace744d7}{2026-07-19T18:24:58Z}{build+decorrelated}{\badge{BadgeGreen}{current}}
\receiptrow{164a0a543f6637cb}{2026-07-20T01:22:11Z}{build+decorrelated}{\badge{BadgeGreen}{current}}
\receiptrow{100eb93aa5f473af}{2026-07-20T20:47:01Z}{build+decorrelated}{\badge{BadgeGreen}{current}}
\receiptrow{1f74787842459411}{2026-07-20T20:58:34Z}{build+replay+hardened+decorrelated}{\badge{BadgeGreen}{current}}
\receiptrow{807e980b9db27d7a}{2026-07-21T00:04:07Z}{build+decorrelated}{\badge{BadgeGreen}{current}}
\receiptrow{efa933bf3665661d}{2026-07-21T05:39:20Z}{build+replay+hardened+decorrelated}{\badge{BadgeGreen}{current}}
\receiptrow{64dad6f82e052f30}{2026-07-22T18:27:19Z}{build+decorrelated}{\badge{BadgeGreen}{current}}
\receiptrow{2a1434c39627845c}{2026-07-31T22:00:17Z}{build+decorrelated}{\badge{BadgeGreen}{current}}
\receiptrow{d12854bca5ce90c3}{2026-07-31T22:29:11Z}{build+decorrelated}{\badge{BadgeGreen}{current}}
\receiptrow{b0d9534dd5d7039b}{2026-08-02T23:31:22Z}{build+decorrelated}{\badge{BadgeGreen}{current}}
\receiptrow{bf1f56225eddc2cf}{2026-08-03T16:27:26Z}{build+decorrelated}{\badge{BadgeGreen}{current}}
\receiptrow{e839dbc9244ca3d2}{2026-08-03T22:46:10Z}{build+decorrelated}{\badge{BadgeGreen}{current}}
\end{receipttable}
\layerline{\layeritem{build}{lake build (elaboration)}{\badge{BadgeGreen}{current}}\layeritem{replay}{leanchecker (official kernel replay)}{\badge{BadgeGreen}{current}}\layeritem{hardened}{safe\_verify (hardened screen)}{\badge{BadgeGreen}{current}}\layeritem{decorrelated}{lean4export + nanoda (decorrelated kernel)}{\badge{BadgeGreen}{current}}}
\sourcenote{\badge{BadgeGray}{comparator: none} no fixture in the three-implementation comparator harness (compiled / kernel-interpreted / reference-interpreted) exercises this rule.}
\ruleblock{HOL4 statement}
\absentbox{NO HOL STATEMENT YET}{The satisfaction-judgment extension postdates the HOL4 port entirely:
the ported cone is the base calculus\textquotesingle s check-function
surface (theory \texttt{bern\_catalogue}), and no extension-theory
definition exists in the HOL development --- in particular nothing
shaped like a wiring-certificate discharge. The rule\textquotesingle s
sequent \emph{leaves} would already be expressible in the ported
vocabulary, but the judgment they assemble would not; the box waits on
a HOL port of the satisfaction stratum.
}
\ruleblock{HOL verifier receipts}
\absentbox{NONE}{no HOL verification receipt exists for this rule.}
\ruleblock{Proof}

\sourcenote{Genre: core rule (tier 1), core-as-glue --- the rule
assembles inherited sequent leaves into the wiring-certificate
judgment, and its own mathematical content is deliberately thin: that
the checked arithmetic rows genuinely certify slack-class membership,
and that the class supports the composition algebra the composition
rule recomputes. The section proves exactly that and marks the
assembly gap loudly.}

\textbf{Setting and claim.}\enspace A wiring certificate
$\mathsf{QWireCert}(w, (\mathit{spec}_i) \Rightarrow
\mathit{spec}_{\mathrm{out}}, \kappa)$ carries two per-region families
of pointwise implications --- the lens square's assumption leg
$\alpha_2(w(o_1), a_2) + \kappa(\gamma_1(o_1)) \geq
\alpha_1(w^{\sharp}(o_1, a_2))$ and guarantee leg
$\gamma_1(o_1) \geq \gamma_2(w(o_1))$
\cite{capuccimyers2026} --- which for polynomial data are ordinary
sequents the base calculus checks, plus a slack
$\kappa \in \mathcal{K}^0_\infty$: zero at zero, monotone on
$[0, \infty)$, unbounded at infinity (the \texttt{IsClassK0Inf}
structure quoted above; the non-strict monotone reading is the
catalogued one). The rule's own obligations, proven here: the three
checked rows certify membership for polynomial slacks
(Lemma~\ref{lem:qwc-poly}), and the class supports ``compose by lens
composition and by addition of the slacks''
(Lemma~\ref{lem:qwc-algebra}).

\begin{blemma}[the polynomial discharge]\label{lem:qwc-poly}
Let $p$ be a real polynomial with \emph{(R1)} constant coefficient
zero, \emph{(R2)} derivative nonnegative on $[0, \infty)$, and
\emph{(R3)} positive degree with nonnegative leading coefficient. Then
$x \mapsto p(x)$ is in $\mathcal{K}^0_\infty$.
\end{blemma}
\begin{proof}
(\texttt{polynomial\_discharge}.) Zero at zero: $p(0)$ \emph{is} the
constant coefficient, so (R1) is the field it certifies --- and note
the checker's form of (R1) is one closed exact evaluation, the
ground-fact pattern. Monotone: $[0, \infty)$ is convex and $p$ is
everywhere continuous and differentiable, so nonnegativity of $p'$ on
the interior gives monotonicity on the ray (the mean-value
inequality); (R2) is a single-variable nonnegativity row of exactly
the calculus's cone-checkable shape. Unbounded: a polynomial of
positive degree with nonnegative leading coefficient tends to
$+\infty$ --- (R3) is one syntactic degree check plus one sign row.
Each row maps onto exactly one class field; none is spare.
\end{proof}

\begin{blemma}[the slack algebra]\label{lem:qwc-algebra}
$\mathcal{K}^0_\infty$ is closed under pointwise sum and under
composition; two contract inequalities with constant slacks chain to
one with the summed slack
($f \leq g + c_1$ and $g \leq h + c_2$ pointwise give
$f \leq h + (c_1 + c_2)$); and guarantee legs chain along a cascade
($\gamma_2 \circ w_1 \leq \gamma_1$ and $\gamma_3 \circ w_2 \leq
\gamma_2$ give $\gamma_3 \circ w_2 \circ w_1 \leq \gamma_1$).
\end{blemma}
\begin{proof}
Sum (\texttt{IsClassK0Inf.add}): zeros add; monotonicity is termwise;
for unboundedness, eventually the second summand is nonnegative (it
tends to infinity), so the sum eventually dominates the first, which
tends to infinity. Composition (\texttt{IsClassK0Inf.comp}):
$\kappa_1(\kappa_2(0)) = \kappa_1(0) = 0$; monotone-composed-monotone
is monotone \emph{provided the inner map lands in the outer's domain
of monotonicity} --- which is where nonnegativity of a class slack on
the ray (zero at zero plus monotone) is consumed, not decoration;
unbounded-composed-unbounded is unbounded. The triangle
(\texttt{slack\_triangle}) and the guarantee chain
(\texttt{guarantee\_leg\_comp}) are one-line pointwise transitivity,
stated over an arbitrary interface carrier so the per-region reading
is instantiation, not generalization. In application the
``constants'' are the slack values at the delivered guarantees --- the
triangle is applied pointwise per region --- while composing the slack
\emph{functions} of a cascade is the closure pair's job.
\end{proof}

\begin{bremark}[why each membership row is load-bearing]
\label{rem:qwc-rows}
Drop (R1) and $\kappa(x) = x + 1$ enters the class candidates: a
wiring square then certifies with degradation even on \emph{exact}
upstream delivery ($\kappa(0) = 1$); and the anchor of the algebra ---
nonnegativity of a class slack on the ray, proved from zero-at-zero
plus monotonicity --- fails too, witness $\kappa(x) = x - 1$. Drop (R2) and
$\kappa(x) = x^2 - x$ enters (zero at zero, positive degree, positive
lead): negative on $(0, 1)$, so the nonnegativity anchor fails, and
non-monotone, so composition closure fails at the monotonicity step
(Lemma~\ref{lem:qwc-algebra}'s proof consumes it twice: once as the
outer map's field, once through inner-map nonnegativity), and the
quantitative reading --- a weaker delivered guarantee never
\emph{shrinks} the licensed degradation --- inverts. Drop (R3) and
bounded slacks enter ($\kappa \equiv 0$ passes (R1)--(R2) alone); the
discharge then no longer certifies unboundedness, and $\kappa$ loses
its comparison-function-style gain reading. The zero polynomial is excluded by (R3) deliberately: a
slack-free square is a square with the $\kappa$-term \emph{absent},
never one with $\kappa = 0$, so the additive algebra needs no zero
element.
\end{bremark}

\begin{bremark}[the assembly gap, loud]\label{rem:qwc-gap}
What this section does \emph{not} contain: the square leaves
themselves (landed sequent machinery of the base calculus, cited by
the certificate --- pointwise polynomial implications per region,
checked by the cone machinery), and the end-to-end assembly of the
judgment object, which awaits the staging of the judgment sort in the
certificate schema. The dashboard's conditional mark scopes the
warrant to exactly Lemmas~\ref{lem:qwc-poly}
and~\ref{lem:qwc-algebra}: the arithmetic rows mean what they claim,
and the class composes. The semantic content of what a wiring
certificate \emph{buys} a composite is deliberately not here --- it is
the composition rule's semantic reading, one section over.
\end{bremark}

\begin{anchors}
class definition (bound by display) & \texttt{IsClassK0Inf} (fields \texttt{zero\_at\_zero}, \texttt{mono}, \texttt{unbounded})\\
Lemma \ref{lem:qwc-poly} & \texttt{polynomial\_discharge}\\
Lemma \ref{lem:qwc-algebra}: sum, composition & \texttt{IsClassK0Inf.add}, \texttt{IsClassK0Inf.comp}; nonnegativity anchor \texttt{IsClassK0Inf.nonneg}\\
Lemma \ref{lem:qwc-algebra}: triangle, guarantee chain & \texttt{slack\_triangle}, \texttt{guarantee\_leg\_comp}\\
square legs (cited, not re-proved) & inherited \rulename{BernCone}-family sequents (main package)\\
\end{anchors}


\rulesection{80}{QWireComposeMarks}{core rule}
\label{blcatrule-QWireComposeMarks}
\ruleblock{The rule (Steele metanotation)}
\sourcenote{as realized in the mechanized development; source row: \texttt{\texttt{QTSL-rules.toml}} --- \emph{QWireComposeMarks}. The metanotation statement is typeset by hand for notational quality; its mechanized ground truth is the Lean statement below, extracted from the sources at every build of this book and hash-bound. The rule \emph{roster} is machine-checked (every mechanized rule tag round-trips to exactly one catalogue row); the statement \emph{text} is editorial, checked by eye against the extracted Lean beside it.}
\begin{mathpar}
\inferrule*[Right=\rulename{QWireComposeMarks}]{
  \bigl\{\Sigma_i \vDash (I_i,\ A_i,\ E_i)\bigr\}_{i=1}^{n}
  \quad \text{where } A_i \equiv \textstyle\bigwedge_{j \in \mathrm{in}(i)} \mathrm{ev}(g_j)\ \mathsf{PathPred},
  \ \ E_i \equiv g_i\ \mathsf{ExistsSpec}
  \ \ \text{(conditional delivery: } \forall\text{-conforming-}\exists\text{-delivering, properness presupposed)} \\
  \bigl\{\ell_i\ \text{latency rows}\bigr\}_{i=1}^{n}
  \ \text{(each premise ships its delivery-stage arithmetic + seriality, at its own profile tier)} \\
  \mathsf{QWireCert}\bigl(w,\ (\mathit{spec}_i)_i \Rightarrow \mathit{spec}_{\mathrm{out}},\ -\bigr)
  \qquad
  \bigwedge_i g_i \vdash \mathrm{marks}(\mathit{spec}_{\mathrm{out}})
  \ \text{(mark-entailment cut)} \\
  \mathsf{MarkFlow}\bigl(w,\ (t_k)_k,\ (P_k)_{k \in \mathrm{cyc}_0(w)}\bigr)
  \ \text{(the schedule walk, two columns: stage rows --- difference-constraint feasibility at the labels; carve rows --- containment $P_k \subseteq \mathrm{mark}_k$ + carved-delivery citations + one species-tagged diagonal certificate per zero-lag cycle, exact values as the degenerate species)} \\
  \mathrm{W0}\ \text{(composite properness --- the existing family; the liveness construction\textquotesingle s tier discharge is stepwise, schedule-relative)}
}{
  \mathsf{SysWire}(w, [\Sigma_1, \ldots, \Sigma_n]) \vDash \mathit{spec}_{\mathrm{out}}
}
\end{mathpar}

\ruleflavor{the $\exists$-face sibling of \rulename{QWireCompose}: same wiring premise vocabulary, opposite polarity economics. \textbf{No causality certificate, no ranks, no guard column} --- a dip under a mark survives every extension, so a cyclic wiring conspires \emph{for} the conclusion, and the well-foundedness kit appears precisely in the rule whose conclusion face is $\mathsf{ForallSpec}$, never here. In its place, one composite-level ingredient in two columns: the $\mathsf{MarkFlow}$ walk --- the $\exists$-dual of the causality walk, same arithmetic species (difference constraints over the declaration closure, bounded recompute), opposite verdict role: the $\forall$-walk demands strict increase along every wire, refusing all cycles; this walk demands \emph{consistency}, admitting every cycle whose schedule closes and refusing time-travel. Its stage arithmetic forces every dependency cycle to total lag zero with one shared label --- so the cycles the rule admits are exactly the instantaneous self-supporting loops, and there the walk\textquotesingle s carve column takes over: producer-supplied \emph{carves} (one predicate per cycle mark, presented as a scalar-codomain gate hypernet cited by digest --- its sign-fiber the carve; a plated aggregate gate for token-population members; never a trace, and never a bare point of a continuous state space) checked by containment arithmetic, carved-delivery citations, and one diagonal certificate per cycle asserting the instant loop closes jointly \emph{inside} the carve product --- the discriminant carve-invariance alone provably misses. Exact value pins survive as the degenerate diagonal species, honest exactly where states are finitely presentable. Off the cycles no carve is needed: the $\forall\exists$ premises chain deliveries in topological order, producer-side --- and the premises themselves are the maximal-carve rows, so the whole column is the rule\textquotesingle s own premise form run at tighter predicates. The conclusion is root satisfaction with the demanded marks in the record\textquotesingle s $\mathsf{ExistsSpec}$ list --- one turnstile, no subscript; the face lives in the argument. The joint witness (record, joint law, PDE solution) is constructed once, inside the soundness metatheorem, and never appears on the wire. Guarded premises are not exempt: the composition consumer eliminates each component guard at consumption --- initial membership discharged before the premise is read as a guarantee --- exactly as at \rulename{QWireCompose} (the mirrored elimination row).}
\par\medskip
\rulename{QWireComposeMarks} is the satisfaction stratum\textquotesingle s
composition rule for cyclic wiring at the liveness face --- the polarity
complement of \rulename{QWireCompose}, whose causality certificate makes
a directed cycle unrepresentable. The polarity fact underneath: a safety
violation and a mark realization are both finite-stage events with
opposite persistence. A violation poisons every extension, so circular
reasoning about guarantees-as-floors needs a well-foundedness device; a
dip under a mark survives every extension, so a cycle\textquotesingle s
self-support is productive --- the two-relay loop that refutes the
unguarded circular safety rule is, read at the mark face, exactly the
witness that realizes ``both outputs hit $1$.'' Accordingly this rule
carries \emph{no causality certificate, no ranks, no guards}: from
per-component conditional-delivery judgments (every admitted resolution
whose interface conforms to the assumed in-neighbour marks admits some
resolution through it realizing the own mark --- with a properness
presupposition on the conforming sub-body), their latency rows, the
wiring plumbing, one mark-entailment sequent, the $\mathsf{MarkFlow}$
walk, and the properness family, it concludes satisfaction of the wired
composite at an output spec whose demanded content is its
record\textquotesingle s $\mathsf{ExistsSpec}$ list.

The one composite-level ingredient is the walk, and it runs in two
columns of the same arithmetic species --- bounded recompute over the
declaration closure, never chase derivation, no well-foundedness
anywhere. The \emph{stage column} is the causality walk\textquotesingle s
$\exists$-dual: producer-supplied stage labels per mark, checked against
the difference-constraint system the latency rows induce --- where the
$\forall$-walk demands strict increase along every wire (refusing all
cycles), this walk demands consistency, admitting every cycle whose
schedule closes and refusing time-travel (a relay against an initialized
unit delay demands $t_x = t_y$ and $t_y = t_x + 1$ at once). Its own
arithmetic forces every admitted cycle to total lag zero and one shared
label --- so the admitted cycles are precisely the instantaneous
self-supporting loops. There the \emph{carve column} takes over:
producer-supplied carves --- one \emph{predicate} per cycle mark,
presented as a scalar-codomain gate hypernet cited by digest, its
sign-fiber the carve; a plated aggregate gate where the members are
token-population systems; never a trace, and never a bare point of a
continuous state space --- checked by three row families: containment
(the carve refines its mark), carved-delivery citations of the very
premises the rule already holds (instantiated at the carve events),
and one \emph{diagonal} certificate per cycle, asserting the instant
loop closes jointly inside the carve product. The diagonal is the
load-bearing novelty: carve-invariance alone provably misses it, and
exact value pins survive as its degenerate species --- honest exactly
where states are finitely presentable, which is the reason the column
generalized. The premises themselves are the maximal-carve rows
(machine-checked as an equivalence), so the whole column is the
rule\textquotesingle s own premise form run at tighter predicates ---
a dial from the sketch\textquotesingle s judgments down to points. The
proof section derives both columns from countermodels: delete the
stage column and the relay--delay loop concludes falsely; delete the
diagonal and the swap loop with a swap-closed mark --- the relay
loop\textquotesingle s evil twin, distinguishable by delivered content
alone --- concludes falsely while every premise holds, \emph{including
carve-invariance at every carve}. Both refutations, the refuse/admit
pairs, the discriminant pair, and the subsumption of pins by carves
are machine-checked.

The soundness warrant is a metatheorem whose joint object --- record,
joint law, or solution, per profile --- is constructed once in the
proof and never appears on the wire. The mechanized instance is the
discrete-relational tier: the assembly engine in carved form (carve
rows + schedule + seriality + diagonal closure produce a joint run
realizing every mark at its label, by stage recursion --- the
walk\textquotesingle s rows align every in-neighbour of a marked
component at its own label, where the carve-at-label invariant holds,
so the carved-delivery rows discharge the marked constraints; the
pinned engine is re-derived through it at singleton carves), and the
chaining bridge at both ends of the dial (on acyclic-instant networks
the conditional-delivery premises alone conclude --- via derived
points, and equally via the maximal carve with no diagonal presented
--- so the carve column enters the wire only where a cycle genuinely
closed). The diagonal\textquotesingle s per-species discharge splits
the labor the way the stochastic tier already does: the checker
verifies finite species rows (a witness point, a bounded enumeration,
Lipschitz/gain rows with cycle product below one, convexity or
order-interval flags), and the fixpoint theorem that converts them
into joint solvability --- exhaustion, Banach, Kakutani, Tarski ---
lives in the metatheorem, per profile, proved once. Properness earns a tier
note: the liveness construction consumes \emph{stepwise} properness
(the composite stays live at every stage, given the schedule\textquotesingle s
pins), because bare nonemptiness of the behaviour set is witnessed by
runs the construction cannot use --- blocking components and
pin-sensitive instant subsystems are the countermodels, named in the
proof section.

Two fences are part of the statement. Marks are single-dip
graded-reach demands --- finitely realizable, extension-persistent;
recurrence-shaped liveness (Streett-class acceptance around cycles)
stays outside both cyclic variants. And the rule deliberately trades
completeness for finite presentation: a component delivering only on
its favourite conforming input can witness a joint behaviour that no
conditional-delivery premise certifies --- certifying such composites
without presenting behaviours is the strategy form, future work.

\cataloguenote{The classical anchors sit on the safety side: circular
assume--guarantee is unsound for safety without a well-foundedness
device \cite{abadilamport1993}, and the standard repairs induct over
time \cite{mcmillan1999}. This rule is the observation that the
liveness-mark fragment of the specification language needs no such
device at all --- greatest-fixpoint territory, where the walk replaces
the ranking: schedules must close around cycles (and deliveries must
agree around the instantaneous ones --- agreement certified over
carved predicates, with exact values as the degenerate species) rather
than cycles being forbidden to close. The certificate carries
judgments, stage labels, gate-hypernet digests with their containment,
delivery, and diagonal rows, and one sequent; nothing behaviour-sized.
Mechanization posture, stated loudly: the discrete-relational tier is
proven end-to-end in carved form with the pinned form subsumed (the
Lean cells below, sorry-free, standard axiom closures); stochastic and
PDE instances of the metatheorem, the per-species diagonal devices at
continuous tiers (Banach for the contraction species, Kakutani for the
convex, Tarski for the order-theoretic), the history-dependent
mechanism tier, and the general difference-constraint walk are open,
and the schema realization (an additive sibling node kind beside the
acyclic rule\textquotesingle s, carrying premise citations, latency
rows, the two-column walk as a side condition, and the reused
properness family) is staged for the extension theory\textquotesingle s
next revision, which is why this entry\textquotesingle s dashboard row
shows a catalogued rule ahead of its wire format.
}
\ruleblock{Lean statement}
\par\medskip
\textbf{The discrete-relational tier, bound whole.}\enspace The engine,
the chaining bridge, and the battery --- the hypotheses-necessity
refutations are first-class bound statements, not test scaffolding: the
rule\textquotesingle s two walk columns each carry their forcing
countermodel in the same module that proves the repaired kit sound.

\declhead{BernsteinLogicAux.QtslCompose.AGComposeMarks}{marks\_of\_pinned\_delivery}{\badge{BadgeGreen}{statement bound}}
\begin{leansrc}
/-- **The assembly engine** (the soundness metatheorem's discrete-relational
instance, pinned form): pin rows + schedule walk + seriality + schedule-relative
instant closure produce a joint run realizing every mark at its label. The
joint behaviour exists only under this conclusion's binder. -/
theorem marks_of_pinned_delivery (hser : Serial N) (hfeas : MarkFlowFeasible N t)
    (hpin : PinnedDelivery N mark a) (hclose : InstantClosure N a t) :
    ∃ ω : ℕ → Fin n → V, JointRun N ω ∧ ∀ i, ω (t i) i ∈ mark i
\end{leansrc}
\declhead{BernsteinLogicAux.QtslCompose.AGComposeMarks}{pins\_of\_conditional\_delivery}{\badge{BadgeGreen}{statement bound}}
\begin{leansrc}
/-- **Chaining**: on an acyclic network (any topological order of the
dependency graph), the ∀∃ conditional-delivery premises alone produce a pin —
each component's delivered value is chosen against its in-neighbours' already-
chosen deliveries, in topological order. This is why the premise-shaped rule
needs no value column on the wire in the acyclic case: the pin is derivable,
producer-side, from the premises it already cites. -/
theorem pins_of_conditional_delivery (hloc : DepLocal N) (ord : Fin n → ℕ)
    (hord : ∀ i, ∀ j ∈ N.dep i, ord j < ord i)
    (hdel : ∀ i, ConditionalDelivery N mark i) :
    ∃ a : Fin n → V, PinnedDelivery N mark a
\end{leansrc}
\declhead{BernsteinLogicAux.QtslCompose.AGComposeMarks}{instantClosure\_of\_acyclic}{\badge{BadgeGreen}{statement bound}}
\begin{leansrc}
/-- Acyclicity dissolves the joint instant solve: with a topological order,
seriality alone closes every instant, pins included — coordinate by
coordinate, in order. (Holds for ANY pin vector `a`; the pinned coordinates'
own admission is the pointed rows' business, not this closure's.) -/
theorem instantClosure_of_acyclic (hser : Serial N) (hloc : DepLocal N)
    (ord : Fin n → ℕ) (hord : ∀ i, ∀ j ∈ N.dep i, ord j < ord i)
    (a : Fin n → V) (t : Fin n → ℕ) :
    InstantClosure N a t
\end{leansrc}
\declhead{BernsteinLogicAux.QtslCompose.AGComposeMarks}{marks\_of\_conditional\_delivery\_acyclic}{\badge{BadgeGreen}{statement bound}}
\begin{leansrc}
/-- **The premise-shaped rule, acyclic instance**: on a network with a
topological order, the ∀∃ conditional-delivery premises + the schedule walk +
seriality conclude — no pin appears among the hypotheses (it is derived).
This is the discrete face of the design's claim that the value column is
forced only by instant cycles. -/
theorem marks_of_conditional_delivery_acyclic (hser : Serial N) (hloc : DepLocal N)
    (ord : Fin n → ℕ) (hord : ∀ i, ∀ j ∈ N.dep i, ord j < ord i)
    (hfeas : MarkFlowFeasible N t) (hdel : ∀ i, ConditionalDelivery N mark i) :
    ∃ ω : ℕ → Fin n → V, JointRun N ω ∧ ∀ i, ω (t i) i ∈ mark i
\end{leansrc}
\declhead{BernsteinLogicAux.QtslCompose.AGComposeMarks}{z1\_walk\_infeasible}{\badge{BadgeGreen}{statement bound}}
\begin{leansrc}
/-- The schedule walk REFUSES the relay–delay loop: its rows demand
`t 0 = t 1 ∧ t 1 = t 0 + 1`. -/
theorem z1_walk_infeasible : ¬ ∃ t : Fin 2 → ℕ, MarkFlowFeasible z1Net t
\end{leansrc}
\declhead{BernsteinLogicAux.QtslCompose.AGComposeMarks}{allOnes\_walk\_feasible}{\badge{BadgeGreen}{statement bound}}
\begin{leansrc}
/-- The schedule walk ADMITS the all-ones loop at the shared label `0`. -/
theorem allOnes_walk_feasible : MarkFlowFeasible allOnesNet ![0, 0]
\end{leansrc}
\declhead{BernsteinLogicAux.QtslCompose.AGComposeMarks}{z1\_refutes\_walkless}{\badge{BadgeGreen}{statement bound}}
\begin{leansrc}
/-- **M1 (the walk deleted): the pin-rowed rule WITHOUT the schedule walk is
unsound** — the relay–delay loop satisfies seriality, the pin rows, and the
instant closure, yet realizes no mark on any run at any stage. (The design
memo's forcing argument for MarkFlow, machine-checked; `z1_walk_infeasible`
shows the walk indeed refuses this instance.) -/
theorem z1_refutes_walkless :
    ¬ (∀ (m : ℕ) (W : Type) (_ : Nonempty W) (N : MarkNet m W)
        (mark : Fin m → Set W) (a : Fin m → W) (t : Fin m → ℕ),
        Serial N → PinnedDelivery N mark a → InstantClosure N a t →
        ∃ ω : ℕ → Fin m → W, JointRun N ω ∧ ∀ i, ∃ k, ω k i ∈ mark i)
\end{leansrc}
\declhead{BernsteinLogicAux.QtslCompose.AGComposeMarks}{z5\_refutes\_pinless}{\badge{BadgeGreen}{statement bound}}
\begin{leansrc}
/-- **M2 (the pins deleted): the premise-shaped rule without a value column is
unsound** — conditional delivery (∀∃) at every component, dep-locality,
a feasible schedule, unpinned joint instant closure, AND a nonempty composite
all hold for the swap loop with swap-closed marks, yet no joint run realizes
any mark: each component maps the mark set into the mark set (the premises
are true), but the loop's instant fixpoint lands outside it. The one
composite-level ingredient the stage arithmetic cannot see is the VALUE
consistency — the pin rows. -/
theorem z5_refutes_pinless :
    ¬ (∀ (m : ℕ) (W : Type) (_ : Nonempty W) (N : MarkNet m W)
        (mark : Fin m → Set W) (t : Fin m → ℕ),
        Serial N → DepLocal N → (∀ i, ConditionalDelivery N mark i) →
        MarkFlowFeasible N t → W0Serial N → (∃ ω, JointRun N ω) →
        ∃ ω : ℕ → Fin m → W, JointRun N ω ∧ ∀ i, ∃ k, ω k i ∈ mark i)
\end{leansrc}
\declhead{BernsteinLogicAux.QtslCompose.AGComposeMarks}{z5\_no\_pin}{\badge{BadgeGreen}{statement bound}}
\begin{leansrc}
/-- **The repair refuses Z5**: no pin exists — the pin rows demand a value
assignment with `a 0 = swap3 (a 1)` and `a 1 = a 0` inside the marks, and the
swap has no fixpoint there. (With `z5_refutes_pinless`: the pin rows are
exactly the ingredient that separates the swap loop from the all-ones loop,
which `allOnes_marks_realized` below certifies through the same engine.) -/
theorem z5_no_pin : ¬ ∃ a : Fin 2 → Fin 3, PinnedDelivery z5Net z5Mark a
\end{leansrc}
\declhead{BernsteinLogicAux.QtslCompose.AGComposeMarks}{allOnes\_marks\_realized}{\badge{BadgeGreen}{statement bound}}
\begin{leansrc}
/-- **Positive control**: the assembly engine's full hypothesis kit is
satisfiable on the all-ones relay loop and concludes — both marks realized at
the shared label `0` (the design's Role-2 example, end-to-end: self-support
productive under the ∃-face). -/
theorem allOnes_marks_realized :
    ∃ ω : ℕ → Fin 2 → Bool, JointRun allOnesNet ω ∧
      ∀ i, ω ((![0, 0] : Fin 2 → ℕ) i) i ∈ boolMark i
\end{leansrc}
\par\medskip
\textbf{The carve generalization (r5), bound whole.}\enspace Pins as
predicates: the carved engine at general lags, the subsumption of value
pins as singleton carves (degeneracy as a theorem), the dial identity
(the rule\textquotesingle s $\forall\exists$ premise IS the
maximal-carve row), the acyclic bridge at the maximal carve, the
discriminant pair (carve-invariance holds on the swap loop; no carve
inside the marks closes its diagonal), the carve-level mutants, and the
fat-carve positive control.

\declhead{BernsteinLogicAux.QtslCompose.AGComposeCarves}{marks\_of\_carved\_delivery}{\badge{BadgeGreen}{statement bound}}
\begin{leansrc}
/-- **The generalized assembly engine** (C1): carve rows + schedule walk +
seriality + diagonal closure produce a joint run realizing every mark at its
label. No hypothesis names a value; the joint behaviour — and every delivered
point — exists only under this conclusion's binder. -/
theorem marks_of_carved_delivery (hser : Serial N) (hfeas : MarkFlowFeasible N t)
    (hpin : CarvedDelivery N mark P) (hclose : CarveClosure N P t) :
    ∃ ω : ℕ → Fin n → V, JointRun N ω ∧ ∀ i, ω (t i) i ∈ mark i
\end{leansrc}
\declhead{BernsteinLogicAux.QtslCompose.AGComposeCarves}{carvedDelivery\_of\_pinned}{\badge{BadgeGreen}{statement bound}}
\begin{leansrc}
/-- The r4 pin rows are the singleton instance of the carve rows. -/
theorem carvedDelivery_of_pinned (hpin : PinnedDelivery N mark a) :
    CarvedDelivery N mark (fun i => {a i})
\end{leansrc}
\declhead{BernsteinLogicAux.QtslCompose.AGComposeCarves}{carveClosure\_of\_pinned}{\badge{BadgeGreen}{statement bound}}
\begin{leansrc}
/-- The r4 closure discharges the singleton diagonal. The walk is
load-bearing: every in-neighbour of a pinned instant coordinate is itself
labeled `k` (zero-lag row), so it is either an instant coordinate held at its
pin by the closure or a pipelined coordinate held at its pin by conformance —
and then the pointed row fires. -/
theorem carveClosure_of_pinned (hfeas : MarkFlowFeasible N t)
    (hpin : PinnedDelivery N mark a) (hclose : InstantClosure N a t) :
    CarveClosure N (fun i => {a i}) t
\end{leansrc}
\declhead{BernsteinLogicAux.QtslCompose.AGComposeCarves}{marks\_of\_pinned\_delivery\_via\_carve}{\badge{BadgeGreen}{statement bound}}
\begin{leansrc}
/-- **Degeneracy is a theorem, not a coincidence**: the r4 engine's exact
statement, re-derived through the generalized engine at singleton carves. -/
theorem marks_of_pinned_delivery_via_carve (hser : Serial N)
    (hfeas : MarkFlowFeasible N t) (hpin : PinnedDelivery N mark a)
    (hclose : InstantClosure N a t) :
    ∃ ω : ℕ → Fin n → V, JointRun N ω ∧ ∀ i, ω (t i) i ∈ mark i
\end{leansrc}
\declhead{BernsteinLogicAux.QtslCompose.AGComposeCarves}{carvedDelivery\_mark\_iff\_conditional}{\badge{BadgeGreen}{statement bound}}
\begin{leansrc}
/-- At `P = mark`, the carve rows ARE the conditional-delivery premises: his
sketch's judgments and the r4 pin rows are the two ends of one dial. -/
theorem carvedDelivery_mark_iff_conditional :
    CarvedDelivery N mark mark ↔ ∀ i, ConditionalDelivery N mark i
\end{leansrc}
\declhead{BernsteinLogicAux.QtslCompose.AGComposeCarves}{carveClosure\_of\_acyclic}{\badge{BadgeGreen}{statement bound}}
\begin{leansrc}
/-- Acyclicity dissolves the diagonal at ANY delivery-closed carve: with a
topological order, pinned instant coordinates choose from `F ∩ P` one
`ord`-level at a time (their in-neighbours are all labeled `k` by the walk:
pipelined ones held in-carve by conformance, instant ones settled at earlier
levels), and unpinned coordinates close by seriality as in r4. -/
theorem carveClosure_of_acyclic (hser : Serial N) (hloc : DepLocal N)
    (ord : Fin n → ℕ) (hord : ∀ i, ∀ j ∈ N.dep i, ord j < ord i)
    (hfeas : MarkFlowFeasible N t) (hpin : CarvedDelivery N mark P) :
    CarveClosure N P t
\end{leansrc}
\declhead{BernsteinLogicAux.QtslCompose.AGComposeCarves}{marks\_acyclic\_via\_maximal\_carve}{\badge{BadgeGreen}{statement bound}}
\begin{leansrc}
/-- **The dial's other end, machine-checked** (the second derivation of r4's
acyclic corollary): on acyclic networks, his sketch's ∀∃ premises conclude
THROUGH the carve engine at the MAXIMAL carve `P = mark` — no point derived,
no diagonal presented. With `marks_of_conditional_delivery_acyclic` (r4,
via derived points) this brackets the dial from both ends. -/
theorem marks_acyclic_via_maximal_carve (hser : Serial N) (hloc : DepLocal N)
    (ord : Fin n → ℕ) (hord : ∀ i, ∀ j ∈ N.dep i, ord j < ord i)
    (hfeas : MarkFlowFeasible N t)
    (hdel : ∀ i, ConditionalDelivery N mark i) :
    ∃ ω : ℕ → Fin n → V, JointRun N ω ∧ ∀ i, ω (t i) i ∈ mark i
\end{leansrc}
\declhead{BernsteinLogicAux.QtslCompose.AGComposeCarves}{z5\_carvedDelivery\_mark}{\badge{BadgeGreen}{statement bound}}
\begin{leansrc}
/-- **C5 — invariance is NOT the discriminant**: the swap loop's carve rows
HOLD at the maximal carve (swap maps the mark set into the mark set; the
relay copies it). What Z5 loses is the diagonal, not the delivery. -/
theorem z5_carvedDelivery_mark : CarvedDelivery z5Net z5Mark z5Mark
\end{leansrc}
\declhead{BernsteinLogicAux.QtslCompose.AGComposeCarves}{z5\_no\_carve\_closure}{\badge{BadgeGreen}{statement bound}}
\begin{leansrc}
/-- **C6 — the diagonal IS the discriminant**: NO carve system inside the
marks closes the swap loop's instant diagonal — containment and closure alone
contradict (the delivery rows are not even needed). The r4 refusal
`z5_no_pin` is this fact's singleton shadow. -/
theorem z5_no_carve_closure :
    ¬ ∃ P : Fin 2 → Set (Fin 3), (∀ i, P i ⊆ z5Mark i) ∧
      CarveClosure z5Net P ![0, 0]
\end{leansrc}
\declhead{BernsteinLogicAux.QtslCompose.AGComposeCarves}{z5\_refutes\_closureless}{\badge{BadgeGreen}{statement bound}}
\begin{leansrc}
/-- **M2′ (the diagonal deleted): the carve rule without `CarveClosure` is
unsound** — carve rows at the maximal carve, dep-locality, a feasible
schedule, unpinned joint instant closure, and a nonempty composite all hold
on the swap loop, yet no joint run realizes any mark. The r4 mutant
`z5_refutes_pinless` restated in r5 vocabulary: what the stage arithmetic
cannot see, carve-invariance cannot see either. -/
theorem z5_refutes_closureless :
    ¬ (∀ (m : ℕ) (W : Type) (_ : Nonempty W) (N : MarkNet m W)
        (mark P : Fin m → Set W) (t : Fin m → ℕ),
        Serial N → DepLocal N → CarvedDelivery N mark P →
        MarkFlowFeasible N t → W0Serial N → (∃ ω, JointRun N ω) →
        ∃ ω : ℕ → Fin m → W, JointRun N ω ∧ ∀ i, ∃ k, ω k i ∈ mark i)
\end{leansrc}
\declhead{BernsteinLogicAux.QtslCompose.AGComposeCarves}{z1\_refutes\_walkless\_carve}{\badge{BadgeGreen}{statement bound}}
\begin{leansrc}
/-- **M1′ (the walk deleted): the carve rule without `MarkFlowFeasible` is
unsound** — the relay–delay loop carries carve rows AND a discharging
diagonal at the singleton carve `{true}` (the closure's conformance
hypothesis hands the delay's base value in-carve, and the relay copies it),
yet realizes no mark on any run. The walk is forced at carve level exactly
as at value level. -/
theorem z1_refutes_walkless_carve :
    ¬ (∀ (m : ℕ) (W : Type) (_ : Nonempty W) (N : MarkNet m W)
        (mark P : Fin m → Set W) (t : Fin m → ℕ),
        Serial N → CarvedDelivery N mark P → CarveClosure N P t →
        ∃ ω : ℕ → Fin m → W, JointRun N ω ∧ ∀ i, ∃ k, ω k i ∈ mark i)
\end{leansrc}
\declhead{BernsteinLogicAux.QtslCompose.AGComposeCarves}{fatCarve\_marks\_realized}{\badge{BadgeGreen}{statement bound}}
\begin{leansrc}
/-- **The measure-zero objection, answered at theorem level**: the relay loop
realizes its marks through the engine with carve = mark = `{v ≠ 0}` — a
NON-singleton carve; no hypothesis names a value. (The diagonal closes at a
point the PROOF chooses; the kit never presents one. Z5's swap loop, with the
same carve and the same walk, has no such point — `z5_no_carve_closure`.) -/
theorem fatCarve_marks_realized :
    ∃ ω : ℕ → Fin 2 → Fin 3, JointRun relay3Net ω ∧
      ∀ i, ω ((![0, 0] : Fin 2 → ℕ) i) i ∈ relay3Mark i
\end{leansrc}
\ruleblock{Correspondence (Steele display $\leftrightarrow$ Lean)}
\badge{BadgeAmber}{the Steele display and a bound Lean statement and the dictionary text CHANGED since certification --- re-review and re-certify}
\sourcenote{The dictionary below is a reviewer-certified statement of how the mechanized statements above interpret the displayed rule: display metavariable $\mapsto$ Lean binder, each display hypothesis $\mapsto$ the premise or side condition carrying it, the display conclusion $\mapsto$ the goal --- with any deliberate difference in width between display and mechanization named outright. Certification is recorded as a hash triple over (the displayed cell, the extracted statement texts, this dictionary\textquotesingle s own text); if any of the three changes, this block turns amber until a reviewer re-certifies. The correspondence claim itself remains editorial --- reviewer-grade, not machine-checked; its machine-checked successor types each soundness statement against the rule as a structured object, and supersedes this block per rule as it lands.}
The display is the rule at the satisfaction stratum; the bound set is
its soundness metatheorem\textquotesingle s discrete-relational
instance, over an explicit network model (`MarkNet`: in-neighbour sets,
per-component delivery lags, memoryless admitted-output relations over
wire snapshots, initial segments; `JointRun` the per-stage constraint
satisfaction). The dictionary, display constituent by bound statement:
\begin{itemize}
\item the conditional-delivery premises $\Sigma_i \vDash (I_i, A_i,
  E_i)$ $\mapsto$ \texttt{ConditionalDelivery} (every mark-conforming
  snapshot admits an own-mark output through it) --- the memoryless
  instantiation; the latency rows\textquotesingle{} seriality
  obligation $\mapsto$ \texttt{Serial}; dep-locality (a component
  reads only its declared in-neighbours) $\mapsto$ \texttt{DepLocal}.
\item the $\mathsf{MarkFlow}$ stage column $\mapsto$
  \texttt{MarkFlowFeasible} --- the walk\textquotesingle s Lean
  \emph{specification} at the exact-lag row shape ($t_i = t_j +
  \mathrm{lag}_i$ per in-edge, $\mathrm{lag}_i \leq t_i$); the general
  difference-constraint species (windows, rational stages) is the wire
  format\textquotesingle s and is honestly absent here.
\item the carve column\textquotesingle s row families $\mapsto$
  \texttt{CarvedDelivery} (containment + carved delivery: every
  snapshot with in-neighbour reads inside the carves admits an output
  in the own carve) and \texttt{CarveClosure} (the diagonal: at every
  stage, for every carve-conforming base, the instant subsystem closes
  jointly with labeled coordinates in-carve and admitted) --- the
  latter consumed \emph{abstractly}: its per-species certificate
  dischargers are the wire format\textquotesingle s, its fixpoint
  devices the metatheorem\textquotesingle s. The r4 value column
  $\mapsto$ \texttt{PinnedDelivery} + \texttt{InstantClosure},
  subsumed as the singleton case by
  \texttt{carvedDelivery\_of\_pinned} / \texttt{carveClosure\_of\_pinned} /
  \texttt{marks\_of\_pinned\_delivery\_via\_carve}; the premises
  themselves are the maximal-carve rows,
  \texttt{carvedDelivery\_mark\_iff\_conditional}.
\item the W0-species rows $\mapsto$ \texttt{InstantClosure} /
  \texttt{CarveClosure} (schedule-relative stepwise properness) ---
  deliberately STRONGER than bare nonemptiness, the tier discharge the
  construction consumes; the unpinned form \texttt{W0Serial} appears
  only in the refuted premise-only statements, to make those
  refutations as strong as possible.
\item the conclusion $\mapsto$ the existential in
  \texttt{marks\_of\_carved\_delivery} (engine) /
  \texttt{marks\_of\_conditional\_delivery\_acyclic} +
  \texttt{marks\_acyclic\_via\_maximal\_carve} (both dial ends of the
  acyclic case): some \texttt{JointRun} realizes every mark at its
  label (the joint witness under the binder, never a hypothesis).
\item the countermodel battery, as prereg\textquotesingle d mutants:
  \texttt{z1\_refutes\_walkless} / \texttt{z1\_refutes\_walkless\_carve}
  (stage column deleted, at value and carve level),
  \texttt{z5\_refutes\_pinless} / \texttt{z5\_refutes\_closureless}
  (pin/diagonal column deleted --- against kits with \texttt{W0Serial}
  and composite nonemptiness granted), the discriminant pair
  \texttt{z5\_carvedDelivery\_mark} (invariance holds on the swap
  loop) / \texttt{z5\_no\_carve\_closure} (no carve inside the marks
  closes its diagonal), \texttt{z5\_no\_pin} (its singleton shadow),
  \texttt{z1\_walk\_infeasible} / \texttt{allOnes\_walk\_feasible}
  (the walk\textquotesingle s refuse/admit pair),
  \texttt{allOnes\_marks\_realized} (positive control through the
  pinned engine) + \texttt{allOnes\_direct} (independent witness) +
  \texttt{fatCarve\_marks\_realized} (positive control through the
  carved engine at a non-singleton carve --- no hypothesis names a
  value).
\end{itemize}
Width note: the conditional mark\textquotesingle s content is the
tier gap stated in the commentary --- within the mechanized tier the
mappings above are exact, the engine is proven at general lags, and
the mark-entailment sequent is consumed at its identity instance (the
conclusion realizes the premises\textquotesingle{} own marks; the cut
against a distinct output-spec mark list is sequent-layer machinery,
not this module\textquotesingle s).

\ruleblock{Lean verifier receipts}
\sourcenote{module \texttt{BernsteinLogicAux.QtslCompose.AGComposeMarks}; no attested build yet; source-bound only. The statements above are hash-bound to the source file, but no build attestation exists for this module yet, so no receipt can be judged to attest the quoted build.}
\badge{BadgeAmber}{NO ATTESTED BUILD BOUND for this module}
\sourcenote{7 receipt(s) in the store attest superseded builds of this module and are omitted from this rendering.}
\sourcenote{module \texttt{BernsteinLogicAux.QtslCompose.AGComposeCarves}; no attested build yet; source-bound only. The statements above are hash-bound to the source file, but no build attestation exists for this module yet, so no receipt can be judged to attest the quoted build.}
\badge{BadgeAmber}{NO ATTESTED BUILD BOUND for this module}
\sourcenote{6 receipt(s) in the store attest superseded builds of this module and are omitted from this rendering.}
\layerline{\layeritem{build}{lake build (elaboration)}{\badge{BadgeAmber}{stale only}}\layeritem{replay}{leanchecker (official kernel replay)}{\badge{BadgeGray}{none}}\layeritem{hardened}{safe\_verify (hardened screen)}{\badge{BadgeGray}{none}}\layeritem{decorrelated}{lean4export + nanoda (decorrelated kernel)}{\badge{BadgeAmber}{stale only}}}
\sourcenote{\badge{BadgeGray}{comparator: none} no fixture in the three-implementation comparator harness (compiled / kernel-interpreted / reference-interpreted) exercises this rule.}
\ruleblock{HOL4 statement}
\absentbox{NO HOL STATEMENT YET}{The satisfaction-judgment extension postdates the HOL4 port entirely
(the ported cone is the base calculus\textquotesingle s first-order
check-function surface), and this rule\textquotesingle s wire format is
itself staged ahead of the schema — there is no wire checker for a HOL
statement to mirror yet. The box inherits \rulename{QWireCompose}\textquotesingle s
honest state one rule further out.
}
\ruleblock{HOL verifier receipts}
\absentbox{NONE}{no HOL verification receipt exists for this rule.}
\ruleblock{Proof}

\sourcenote{Genre: core rule (tier 1). The soundness warrant is a
metatheorem whose joint witness never appears on the wire; this section
proves it at the discrete-relational tier --- the tier the mechanization
proves --- and derives every certificate row from a countermodel that
breaks the rule without it. Per-profile instances beyond this tier are
stated as open, never claimed.}

\textbf{Setting and claim.}\enspace The acyclic composition rule's
causality certificate makes a directed cycle unrepresentable; this rule
is its polarity complement, existing precisely to admit cyclic wiring
--- for \emph{liveness marks only}. The polarity fact that makes this
sound where the safety analogue fails: a safety violation and a mark
realization are both finite-stage events with opposite persistence. A
violation of a prefix-closed guarantee poisons every extension, so under
a $\forall$-reading a cyclic wiring can conspire to admit a bad
behaviour and a well-foundedness device is mandatory. A dip under a
mark survives every extension, so under an $\exists$-reading the same
conspiracy works \emph{for} the conclusion: circular self-support is a
post-fixpoint, productive rather than vicious. The rule concludes
$\mathsf{SysWire}(w, [\Sigma_1, \ldots, \Sigma_n]) \vDash
\mathit{spec}_{\mathrm{out}}$ where the output spec's demanded content
is its record's $\mathsf{ExistsSpec}$ list --- one turnstile, no
subscript: the face lives in the argument. Premises: per-component
\emph{conditional delivery} (for every admitted resolution whose
interface conforms to the assumed in-neighbour marks, some admitted
resolution through it realizes the own mark --- the
$\forall$-conforming-$\exists$-delivering reading, with a properness
presupposition on the conforming sub-body), each shipping its
\emph{latency row} (delivery-stage arithmetic at its own profile tier);
the wiring certificate's plumbing; one mark-entailment sequent
$\bigwedge_i g_i \vdash \mathrm{marks}(\mathit{spec}_{\mathrm{out}})$;
the $\mathsf{MarkFlow}$ walk in two columns (below); and the
W0-species properness rows.

\begin{bremark}[the walk, stage column: what breaks without it]
\label{rem:qwcm-z1}
Take a relay ($\mathrm{out} \equiv \mathrm{in}$, instantaneous) closed
against a unit delay initialized to $0$, marks ``hits $1$'' on both
wires. Both conditional-delivery premises are true --- a relay fed a
trace hitting $1$ emits one; a delay shifts a hit by one stage --- and
the composite is proper. But its unique behaviour is all-zeros: the
initialization propagates around the cycle forever, and the conclusion
is false. Premise shape alone is therefore unsound; something
composite-level is forced. The forced ingredient is \emph{schedule
arithmetic}: the producer supplies one stage label $t_k$ per mark, and
the checker verifies the finite difference-constraint system assembled
from the latency rows --- label equality across instantaneous wires,
shifts through delay media. The relay--delay loop is refused exactly
here: its rows demand $t_x = t_y$ and $t_y = t_x + 1$, infeasible ---
while the loop of two relays passes at the shared label $0$. This is
the $\exists$-dual of the causality walk: the same arithmetic species
(difference constraints over the declaration closure, bounded
recompute, never chase derivation), the opposite verdict role --- the
$\forall$-walk demands strict increase along every wire, refusing all
cycles; this walk demands consistency, admitting every cycle whose
schedule closes. Feasibility of difference constraints is a
Bellman--Ford-class check, not an induction device: no
well-foundedness enters. Mechanized: \texttt{z1\_refutes\_walkless}
(every other hypothesis of the pinned engine true, conclusion false),
with the refuse/admit pair \texttt{z1\_walk\_infeasible} /
\texttt{allOnes\_walk\_feasible}.
\end{bremark}

\begin{bremark}[the walk, value column: what breaks without it]
\label{rem:qwcm-z5}
The stage column has a blind spot, and it is exactly where this rule
lives. Around any dependency cycle the latency rows sum, so a feasible
schedule forces every cycle to total lag zero and one shared label:
the admitted cycles are the instantaneous self-supporting loops, whose
marks all land at one instant. There the $\exists$-chosen delivery
\emph{values} must agree around the loop --- and no stage arithmetic
can see values. Witness: a pointwise swap gate over $\{0, 1, 1'\}$
($1 \mapsto 1'$, $1' \mapsto 1$, $0 \mapsto 0$) closed against a relay,
with mark ``hits $\{1, 1'\}$'' on both wires --- a mark set closed
under the swap. Every remaining hypothesis holds: both conditional
deliveries (the swap maps the mark set into itself; the relay copies
it), the schedule ($t_x = t_y = 0$), properness (all-zeros --- $0$ is
the swap's fixpoint), even joint instant seriality. But every joint
run satisfies $x = \mathrm{swap}(x)$ around the loop, hence is
all-zeros, and the conclusion is false: the swap loop is the relay
loop's evil twin, distinguished by values alone. The forced repair is
the walk's \emph{value column}: on each zero-lag cycle the producer
supplies a \emph{pin} $\alpha_k$ per mark --- the delivered value,
spec-representable and label-sized, the value analogue of $t_k$, never
a trace --- and the checker verifies mark membership per pinned wire
plus \emph{pointed delivery}: each pinned component's
conditional-delivery judgment cited at the pins (through every
conforming input agreeing with the neighbouring pins, the own pin is
admitted). The relay loop passes ($\alpha_x = \alpha_y = 1$); the swap
loop is refused ($\alpha_x = \mathrm{swap}(\alpha_y)$,
$\alpha_y = \alpha_x$ has no solution inside the marks). Mechanized:
\texttt{z5\_refutes\_pinless} (the premise-only kit accepts a false
conclusion) and \texttt{z5\_no\_pin} (the pin rows refuse the
witness). Off the cycles no pin is needed --- see the chaining lemma.
An exact value is, however, the wrong general datum --- the next
remark generalizes the column and keeps this one as its degenerate
species.
\end{bremark}

\begin{bremark}[the carve column: pins generalized to predicates]
\label{rem:qwcm-carve}
An exact pin is measure-zero data: honest for a finite alphabet,
unpresentable for members whose wire states are continuous or
infinite-dimensional. The column therefore generalizes: per zero-lag-
cycle mark the producer supplies a \emph{carve} --- a predicate on the
wire\textquotesingle s state carrier, presented as a scalar-codomain
gate hypernet cited by digest (its nonnegativity locus the carve), or
a plated aggregate gate where the member is a token-population system.
The soundness question this raises is the design\textquotesingle s
crux: a point composes trivially around an instant cycle (the joint
assignment \emph{is} the pin vector), while a carve leaves each
delivery free within a set, so joint existence becomes a nonemptiness
obligation over the carve product --- the \emph{diagonal}. Carve-
invariance does not settle it: on the swap loop the carve rows hold at
the maximal carve --- the swap maps the mark set into itself
(\texttt{z5\_carvedDelivery\_mark}) --- yet \emph{no} carve inside the
marks closes the diagonal (\texttt{z5\_no\_carve\_closure}): around
the loop any jointly-admitted value is a swap fixpoint, and the only
one lies outside the marks. So the carve rows come in three families:
\emph{containment} (carve refines mark), \emph{carved delivery} (the
rule\textquotesingle s own conditional-delivery judgment, cited at the
carve events --- the premises are literally the maximal-carve rows,
\texttt{carvedDelivery\_mark\_iff\_conditional}, so the column is one
dial from the sketch\textquotesingle s judgments down to r4\textquotesingle s
points), and one \emph{diagonal certificate} per cycle, species-tagged:
a witness point or bounded enumeration where states are finite (the
degenerate species --- and the swap loop is refused mechanically
there); contraction rows with exact-rational cycle gain below one
(Banach closes the diagonal --- small-gain arithmetic reused at the
$\exists$-face); convexity or order-interval flags (Kakutani, Tarski).
The checker verifies only the species\textquotesingle{} finite rows;
each fixpoint theorem lives in the soundness metatheorem, per profile
--- exactly where Ionescu--Tulcea already lives for the stochastic
assembly. Mechanized at this tier: the carved engine
(\texttt{marks\_of\_carved\_delivery}, the diagonal consumed
abstractly as \texttt{CarveClosure}), the subsumption of pins as
singleton carves (\texttt{marks\_of\_pinned\_delivery\_via\_carve} ---
degeneracy as a theorem), the carve-level mutants
(\texttt{z5\_refutes\_closureless},
\texttt{z1\_refutes\_walkless\_carve}), and a fat-carve positive
control in which no hypothesis names a value
(\texttt{fatCarve\_marks\_realized}).
\end{bremark}

\begin{bremark}[the reading of the premise, the properness
presupposition, and W0: what breaks without each]
\label{rem:qwcm-rest}
Three further hypotheses earn their places by countermodels at the
design tier. (i) Reading the component premise as co-realization
(``some behaviour realizes assumptions and guarantee together'') is
unsound: a component can co-realize only in a mode the closed loop
never enters (a gate that cooperates only on inputs the wiring cannot
produce), so the conditioning must do real work ---
$\forall$-conforming-$\exists$-delivering is the reading under which
the metatheorem is true. (ii) The properness presupposition on each
premise's conforming sub-body blocks vacuous conditioning: a component
admitting no conforming input satisfies the conditional vacuously
while delivering nothing. (iii) The W0-species rows block the empty
composite (an inconsistent loop --- a negation gate against a relay
--- satisfies every safety spec and no liveness spec), and for the
liveness \emph{construction} their tier discharge is stepwise: the
composite must stay live at every stage given the schedule's pins,
because bare nonemptiness of the behaviour set is witnessed by runs
the construction cannot use (a component may block after delivering;
an instant subsystem's solvability may flip on a pinned coordinate's
value). Stepwise properness is what the mechanized engine consumes.
\end{bremark}

\begin{blemma}[the assembly engine --- soundness at the
discrete-relational tier]\label{lem:qwcm-engine}
Fix $n$ components over a wire alphabet, each with in-neighbours, a
delivery lag ($0$ = instantaneous), a memoryless admitted-output
relation reading its lagged snapshot, and an initial segment. Suppose:
(a) seriality (initial sets and output relations nonempty --- the
mechanism-presentation obligation of the latency rows); (b) a feasible
schedule; (c) pins with their membership and pointed-delivery rows;
(d) schedule-relative instant closure (at every stage, with pinned
instantaneous coordinates held at their pins and pipelined coordinates
at any given values, the remaining instantaneous coordinates close
jointly). Then some joint run realizes every mark at its label.

\emph{Proof (as mechanized, \texttt{marks\_of\_pinned\_delivery}).}
Build the run by stage recursion. At stage $k$, pipelined coordinates
choose: their initial value before the pipe fills; \emph{the pin} at
their own label (legal because $t_i \geq \mathrm{lag}_i$ is a walk
row); any admitted output otherwise (seriality). Instantaneous
coordinates close jointly by (d), with the schedule's pinned ones held
at their pins. Two facts discharge the pinned constraints. First, the
\emph{pin invariant}: at its label every coordinate carries its pin ---
by construction for pipelined coordinates, by (d)'s base clause for
instantaneous ones. Second, the walk's alignment: for a marked
component $i$, every in-neighbour $j$ satisfies $t_i = t_j +
\mathrm{lag}_i$, so the snapshot $i$ reads at delivery time is taken at
stage $t_j$ --- where, by the invariant, $j$ carries its pin. The
pointed-delivery row therefore applies verbatim and admits the pin.
Extension past the last label is seriality plus (d) at unpinned
stages. The realized-at-label conclusion is the invariant plus the
membership rows. \emph{Carved form (as mechanized,
\texttt{marks\_of\_carved\_delivery}):} the same recursion with the
invariant weakened from ``carries its pin'' to ``lands inside its
carve'' --- pipelined labeled coordinates now \emph{choose} from the
intersection of their admitted set with their carve (nonempty by the
carved-delivery row, since the walk aligns their in-neighbours at
labels where the invariant holds), and instantaneous labeled
coordinates close through the diagonal certificate, in-carve and
admitted at once. The pinned statement above is re-derived from the
carved engine at singleton carves
(\texttt{marks\_of\_pinned\_delivery\_via\_carve}). \qedhere
\end{blemma}

\begin{blemma}[the chaining bridge --- no pins where no cycle bites]
\label{lem:qwcm-chain}
On a network whose dependency graph admits a topological order, the
conditional-delivery premises alone produce pins satisfying the rows
of Lemma~\ref{lem:qwcm-engine}: choose each component's delivered
value by its premise, applied to the already-chosen deliveries of its
in-neighbours, in topological order; dep-locality (a component reads
only its declared in-neighbours) turns the chosen witness into the
pointed row. Moreover acyclicity dissolves the joint instant solve:
seriality alone furnishes the instant closure, coordinate by
coordinate in order. Hence the \emph{premise-shaped rule with no value
column is sound on acyclic-instant networks}
(\texttt{marks\_of\_conditional\_delivery\_acyclic}) --- the pins are
producer-derivable search data there, entering the wire only where the
walk's own arithmetic says a cycle closed. \emph{Proof (as
mechanized, \texttt{pins\_of\_conditional\_delivery} +
\texttt{instantClosure\_of\_acyclic}).} Induction on topological
levels: at each level, apply the premise at the snapshot assembling
the lower levels' pins (conforming by their membership conjuncts);
transport the chosen admission along dep-locality to every
dep-agreeing snapshot. The instant closure by the same induction, with
seriality supplying unpinned choices. The bridge holds at the dial's
other end too: the same premises \emph{are} the maximal-carve rows,
and acyclicity dissolves the diagonal at any delivery-closed carve
(\texttt{carveClosure\_of\_acyclic}), so the carved engine also
concludes with nothing pinned and nothing carved on the wire
(\texttt{marks\_acyclic\_via\_maximal\_carve}). \qedhere
\end{blemma}

\begin{bremark}[the metatheorem beyond this tier]\label{rem:qwcm-tiers}
The joint object lives in the soundness proof and only there: at the
discrete-relational tier it is the recursion above; at stochastic
tiers the joint law is assembled by Ionescu--Tulcea/Kolmogorov
extension along the wiring's stagewise dependency structure, with the
mark-escaping resolutions selected per the premises; for PDE profiles,
whatever solution object the profile's semantic theorem provides. The
rule's premise kit does not change per profile --- only the bridge
lemmas do, exactly as in the acyclic rule's grid. Those instances, the
per-species diagonal devices at continuous tiers (Banach for the
contraction species, Kakutani/Brouwer for the convex, Knaster--Tarski
for the order-theoretic --- each converting its species' finite rows
into the joint instant solve the carved engine consumes abstractly),
history-dependent mechanisms, and the general difference-constraint
walk (windows, rational stages) are open; the mechanized tier is the
one this section proves, and the conditional mark warrants exactly it.
Two scope fences are part of the statement: marks are single-dip
graded-reach demands (finitely realizable, extension-persistent) ---
recurrence-shaped liveness stays outside; and the rule trades
completeness for finite presentation --- a component delivering only
on its favourite conforming input can witness a joint behaviour no
conditional-delivery premise certifies.
\end{bremark}


\rulesection{81}{QSimTransport}{core rule}
\label{blcatrule-QSimTransport}
\ruleblock{The rule (Steele metanotation)}
\sourcenote{as realized in the mechanized development; source row: \texttt{\texttt{QTSL-rules.toml}} --- \emph{QSimTransport}. The metanotation statement is typeset by hand for notational quality; its mechanized ground truth is the Lean statement below, extracted from the sources at every build of this book and hash-bound. The rule \emph{roster} is machine-checked (every mechanized rule tag round-trips to exactly one catalogue row); the statement \emph{text} is editorial, checked by eye against the extracted Lean beside it.}
\begin{mathpar}
\inferrule*[Right=\rulename{QSimTransport}]{
  S' \vDash \mathit{spec}' \\
  \mathsf{QSim}(\sigma, S, S')
}{
  S \vDash \sigma^{*}\mathit{spec}'
}
\end{mathpar}

\ruleflavor{certification transport along a coarse-graining: $\sigma$ is an exact simulation witness, $S'$ coarse-grains $S$ --- if $S'$ is certified and $\mathsf{QSim}(\sigma, S, S')$ holds ($\sigma$ maps every behaviour of $S$ to a behaviour of $S'$), then $S$ is certified for the pulled-back spec ($\sigma^{*}$ pulls back the whole interface spec, assumptions and guarantees). Expectation specs pull back with the bound intact; region/invariance specs pull back to the preimage region. The former side condition (per-mode polynomial map + compatibility squares as BL leaves) is the kind\textquotesingle s discharge, one node down: \rulename{QSimDischarge} consumes exactly those squares ($\sigma$-strict tier; the kernel tier joins with its own bounded coverage walk), \rulename{QSimCompose} composes witnesses. Exactness is what $\mathsf{QSim}$ asserts --- the approximate $\sigma$-$\varepsilon$ class concludes $\mathsf{QDist}$, never this premise. One-sided functional refinement --- a simulation witness is never a two-sided coupling. satisfaction of an interface-bearing spec is plain \texttt{QSat}, as at \rulename{QWireCompose}.}
\par\medskip
\rulename{QSimTransport} transports certification along a
coarse-graining: if the abstract system $S'$ is certified, and the
simulation judgment $\mathsf{QSim}(\sigma, S, S')$ holds --- its
discharge, one node down, checks the per-mode polynomial
map\textquotesingle s compatibility squares as inherited BL leaves ---
then the concrete system $S$ is certified for the pulled-back spec
$\sigma^{*}\mathit{spec}'$. A producer uses it in the direction the
rule states: certify the abstraction \emph{once} (small state space,
cheap rows), then reuse that certificate under every refinement that
simulates it, paying only the simulation judgment\textquotesingle s
discharge per refinement.

What the cited theorems warrant is three strata. The
\emph{behavioural} form: over compiled behaviour classes, if $\sigma$
maps every behaviour of $S$ to a behaviour of $S'$ (the semantic content
of the squares) and $S'$\textquotesingle s behaviours all satisfy the
spec, then $S$\textquotesingle s behaviours satisfy the pullback --- two
lines, which is the honest size of the rule\textquotesingle s semantic
content; everything substantive in an application is the sequent-level
discharge of the squares. The \emph{finite-kernel} instance, where the
development has machine-checked simulation machinery: a functional simulation
witness at the finite Markov tier is exactly a pushforward-kernel
identity (per state, the pushed concrete kernel row equals the abstract
kernel row at the image state), the marginal-process lemma then
transports laws at every time, and both spec-transport faces follow ---
expectation specs pull back along $\sigma$ with the bound intact, and
region/invariance specs pull back to the preimage region. And the
\emph{general measurable-kernel} tier, mirroring the finite instance
name-for-name over Markov kernels and the Giry monad: the same
pushforward identity (now a \texttt{Measure.map} equation per state),
the same law transport at every horizon, both spec-transport faces ---
with the finite tier\textquotesingle s nonnegativity hypotheses gone,
true by type.

The witness is one-sided by design: a simulation map is functional
refinement, not a two-sided relational coupling. A functional
simulation induces a degenerate coupling, but the transport theorems
never route through it --- they are the pushforward identity plus
summation identities, and the relational tier\textquotesingle s
coupling machinery (joint kernels, absorption hypotheses) is
deliberately not shared. What breaks without the hypotheses: without
the compatibility squares the pushed law and the abstract law decouple
at the first step, and the pullback claims anything; without the
nonnegativity hypotheses of the region form, signed
\textquotedblleft mass\textquotedblright{} could cancel and hide
probability outside the preimage region.

\cataloguenote{The general measurable-kernel tier is discharged (third stratum,
bound below). The simulation premise is the judgment kind
$\mathsf{QSim}$, baked in the extension theory together with its
discharge and composition nodes: the former side condition (map data +
compatibility squares) is that kind\textquotesingle s discharge, so
nothing here renders as a hypothesis mark. Exactness is what the kind
asserts --- the approximate class concludes distance judgments
instead. What remains beyond the three proven strata: the
categorical strengthening under which the composition rule could cite
the lens-composition theorem as a lemma \emph{template} rather than
proving its cells directly --- the direct proofs never pull back along
strict weakenings, which is why the strengthening is not a dependency
of anything on this page. No end-to-end certificate exercises the
rule yet. Related reading, not
dependency: an exchangeability quotient is a simulation map, so
aggregate-space reduction can be \emph{read} as an instance of this
rule --- it stands on its own proof.
}
\ruleblock{Lean statement}
\declhead{BernsteinLogicAux.QtslCompose.SimTransport}{pullback}{\badge{BadgeGreen}{statement bound}}
\begin{leansrc}
/-- **QSimTransport, behavioral form.** `Beh`/`Beh'` are the compiled
behavior classes of the concrete and abstract Systems; `hsim` is the
semantic content of the simulation witness (σ maps behaviors to behaviors);
`hsat` is `QSatI(S', spec')`. Conclusion: `QSatI(S, σ* spec')`. -/
theorem pullback {Ω Ω' : Type*} {σ : Ω → Ω'} {Beh : Set Ω} {Beh' : Set Ω'}
    {spec' : Set Ω'} (hsim : σ '' Beh ⊆ Beh') (hsat : Beh' ⊆ spec') :
    Beh ⊆ σ ⁻¹' spec'
\end{leansrc}
\declhead{BernsteinLogicAux.QtslCompose.SimTransport}{expect\_push}{\badge{BadgeGreen}{statement bound}}
\begin{leansrc}
/-- Observable pullback: the expectation of `g` under a pushed law is the
expectation of the pulled-back observable `g ∘ σ` under the original law.
(The `σ* spec` of the rule display, at expectation-spec granularity.) -/
theorem expect_push (σ : α → β) (μ : α → ℝ) (g : β → ℝ) :
    ∑ b, push σ μ b * g b = ∑ s, μ s * g (σ s)
\end{leansrc}
\declhead{BernsteinLogicAux.QtslCompose.SimTransport}{sim\_expect}{\badge{BadgeGreen}{statement bound}}
\begin{leansrc}
/-- **QSimTransport, expectation form** (finite tier). Along a simulation
witness `PushKernel σ K Q`, the time-`n` expectation of a pulled-back
observable on the concrete chain *equals* the abstract chain's expectation —
`push_law` composed with `expect_push`. -/
theorem sim_expect {σ : α → β} {K : α → α → ℝ} {Q : β → β → ℝ}
    (hσ : PushKernel σ K Q) (μ₀ : α → ℝ) (g : β → ℝ) (n : ℕ) :
    ∑ s, law K μ₀ n s * g (σ s) = ∑ b, law Q (push σ μ₀) n b * g b
\end{leansrc}
\declhead{BernsteinLogicAux.QtslCompose.SimTransport}{sim\_expect\_bound}{\badge{BadgeGreen}{statement bound}}
\begin{leansrc}
/-- Certified expectation bounds transport: a bound proven against the
abstract (coarse) chain holds for the concrete chain at the pulled-back
observable. This is CM's rule 2 at the finite tier, in the direction the
rule states it (certify the abstraction once, reuse under refinement). -/
theorem sim_expect_bound {σ : α → β} {K : α → α → ℝ} {Q : β → β → ℝ}
    (hσ : PushKernel σ K Q) (μ₀ : α → ℝ) (g : β → ℝ) {n : ℕ} {b₀ : ℝ}
    (hb : ∑ b, law Q (push σ μ₀) n b * g b ≤ b₀) :
    ∑ s, law K μ₀ n s * g (σ s) ≤ b₀
\end{leansrc}
\declhead{BernsteinLogicAux.QtslCompose.SimTransport}{sim\_region}{\badge{BadgeGreen}{statement bound}}
\begin{leansrc}
/-- **QSimTransport, region form** (finite tier): if the abstract law at time
`n` puts no mass outside `R'`, the concrete law puts no mass outside
`σ ⁻¹' R'` — invariance specs pull back to the preimage region. -/
theorem sim_region {σ : α → β} {K : α → α → ℝ} {Q : β → β → ℝ}
    (hσ : PushKernel σ K Q) {μ₀ : α → ℝ}
    (hK : ∀ s t, 0 ≤ K s t) (hμ₀ : ∀ s, 0 ≤ μ₀ s) {n : ℕ} {R' : Set β}
    (hR' : ∀ b, b ∉ R' → law Q (push σ μ₀) n b = 0) :
    ∀ s, σ s ∉ R' → law K μ₀ n s = 0
\end{leansrc}
\declhead{BernsteinLogicAux.QtslCompose.SimTransportMeas}{push\_law}{\badge{BadgeGreen}{statement bound}}
\begin{leansrc}
/-- **QSimTransport, measurable-kernel tier — law transport.** Along a
simulation witness, the pushforward of the time-`n` concrete law is the
time-`n` abstract law of the pushed start. (The measurable-tier face of
`QRelChain.push_law`; both spec-transport faces below are corollaries.) -/
theorem push_law {σ : α → β} {K : Kernel α α} {Q : Kernel β β}
    (hσ : Measurable σ) (h : PushKernel σ K Q) (μ : Measure α) :
    ∀ n, (law K μ n).map σ = law Q (μ.map σ) n
  | 0 => rfl
  | n + 1 => by
    rw [law_succ, map_bind _ _ hσ, funext h, ← bind_map _ _ hσ,
      push_law hσ h μ n, law_succ]

/-! ## The two spec-transport faces -/

/-- Expectation specs pull back along `σ` with the bound intact —
`ℝ≥0∞`-observable face, an *equality* as at the finite tier. -/
theorem sim_expect_lintegral {σ : α → β} {K : Kernel α α} {Q : Kernel β β}
    (hσ : Measurable σ) (h : PushKernel σ K Q) (μ : Measure α)
    {g : β → ℝ≥0∞} (hg : Measurable g) (n : ℕ) :
    ∫⁻ s, g (σ s) ∂law K μ n = ∫⁻ b, g b ∂law Q (μ.map σ) n
\end{leansrc}
\declhead{BernsteinLogicAux.QtslCompose.SimTransportMeas}{sim\_expect\_bound}{\badge{BadgeGreen}{statement bound}}
\begin{leansrc}
/-- Certified expectation bounds transport in the direction the rule states:
certify the abstraction once, reuse under refinement. -/
theorem sim_expect_bound {σ : α → β} {K : Kernel α α} {Q : Kernel β β}
    (hσ : Measurable σ) (h : PushKernel σ K Q) (μ : Measure α)
    {g : β → ℝ≥0∞} (hg : Measurable g) {n : ℕ} {b₀ : ℝ≥0∞}
    (hb : ∫⁻ b, g b ∂law Q (μ.map σ) n ≤ b₀) :
    ∫⁻ s, g (σ s) ∂law K μ n ≤ b₀
\end{leansrc}
\declhead{BernsteinLogicAux.QtslCompose.SimTransportMeas}{sim\_region}{\badge{BadgeGreen}{statement bound}}
\begin{leansrc}
/-- Region/invariance specs pull back to the preimage region. The finite
tier's nonnegativity hypotheses (`hK`, `hμ₀` — "signed mass could cancel")
are absent: at this tier they hold by type. -/
theorem sim_region {σ : α → β} {K : Kernel α α} {Q : Kernel β β}
    (hσ : Measurable σ) (h : PushKernel σ K Q) (μ : Measure α)
    {R' : Set β} (hR' : MeasurableSet R') {n : ℕ}
    (hz : law Q (μ.map σ) n R'ᶜ = 0) :
    law K μ n (σ ⁻¹' R')ᶜ = 0
\end{leansrc}
\declhead{BernsteinLogicAux.QtslCompose.SimTransportMeas}{pushKernel\_iff\_comp}{\badge{BadgeGreen}{statement bound}}
\begin{leansrc}
/-- A functional simulation witness is exactly commutation of the step
kernels with the deterministic kernel of `σ`: the square the lens/fibred
strengthening would functor over. -/
theorem pushKernel_iff_comp {σ : α → β} (hσ : Measurable σ)
    {K : Kernel α α} {Q : Kernel β β} :
    PushKernel σ K Q ↔
      Kernel.deterministic σ hσ ∘ₖ K = Q ∘ₖ Kernel.deterministic σ hσ
\end{leansrc}
\ruleblock{Correspondence (Steele display $\leftrightarrow$ Lean)}
\badge{BadgeGreen}{reviewer-certified against the displayed texts}
\sourcenote{The dictionary below is a reviewer-certified statement of how the mechanized statements above interpret the displayed rule: display metavariable $\mapsto$ Lean binder, each display hypothesis $\mapsto$ the premise or side condition carrying it, the display conclusion $\mapsto$ the goal --- with any deliberate difference in width between display and mechanization named outright. Certification is recorded as a hash triple over (the displayed cell, the extracted statement texts, this dictionary\textquotesingle s own text); if any of the three changes, this block turns amber until a reviewer re-certifies. The correspondence claim itself remains editorial --- reviewer-grade, not machine-checked; its machine-checked successor types each soundness statement against the rule as a structured object, and supersedes this block per rule as it lands.}
Three strata, as the commentary\textquotesingle s conditional scope
states; the display maps onto each at its own granularity.
\emph{Behavioural}: the displayed conclusion $S \vDash
\sigma^{*}\mathit{spec}'$ $\mapsto$ \texttt{pullback}\textquotesingle s
goal \texttt{Beh ⊆ σ ⁻¹\textquotesingle{} spec\textquotesingle} over
compiled behaviour classes; the E1 premise $S' \vDash
\mathit{spec}'$ $\mapsto$ \texttt{hsat : Beh\textquotesingle{} ⊆
spec\textquotesingle}; the judgment premise
$\mathsf{QSim}(\sigma, S, S')$ $\mapsto$
\texttt{hsim : σ \textquotesingle\textquotesingle{} Beh ⊆
Beh\textquotesingle} --- the kind\textquotesingle s \emph{semantic
content} exactly: the per-mode polynomial structure of $\sigma$ is
erased (\texttt{σ} is any function), and the BL discharge of the
compatibility squares lives in the kind\textquotesingle s own
discharge node, outside this bound set. The
statement is axiom-free set algebra --- two lines, the honest size of
the rule\textquotesingle s semantic content.
\begin{itemize}
\item \emph{finite-kernel instance}: the simulation witness $\mapsto$
  \texttt{PushKernel σ K Q} --- the per-state marginal identity
  \texttt{push σ (K s) = Q (σ s)}, which is what a functional
  simulation IS at this tier.
\item expectation specs pull back ``with the bound intact'' $\mapsto$
  \texttt{sim\_expect}, and the certified statement is
  \emph{stronger} than the display\textquotesingle s reading: an
  \emph{equality} of expectations at every horizon (\texttt{∑ s, law
  K μ₀ n s * g (σ s) = ∑ b, law Q (push σ μ₀) n b * g b}), with
  \texttt{sim\_expect\_bound} the displayed inequality direction as
  its corollary (certify the abstraction once, reuse under
  refinement).
\item region/invariance specs pull back to the preimage $\mapsto$
  \texttt{sim\_region}: abstract law puts no mass outside
  $R'$ $\Rightarrow$ concrete law puts none outside
  $\sigma^{-1} R'$ --- and this face alone consumes
  nonnegativity (\texttt{hK}, \texttt{hμ₀}): the commentary\textquotesingle s
  ``signed mass could cancel'' hypothesis, visible as exactly two
  binder rows.
\item \texttt{expect\_push} is the observable-pullback engine
  ($\sigma^{*}$ at expectation granularity), fibre-summation algebra.
\end{itemize}
\emph{Measurable-kernel tier} (second bound block, module
\texttt{SimTransportMeas}, namespace \texttt{…QtslCompose.Meas}): the
simulation witness $\mapsto$
\texttt{Meas.PushKernel σ K Q} --- per state, \texttt{(K s).map σ = Q
(σ s)} over \texttt{ProbabilityTheory.Kernel}s; \texttt{push\_law} is
the law transport at every horizon over the \texttt{n}-fold Giry bind;
\texttt{sim\_expect\_bound} and \texttt{sim\_region} are the two
spec-transport faces in measure vocabulary (the addressability
equalities \texttt{sim\_expect\_lintegral}/\texttt{\_integral} live in
the module, cited not bound); and \texttt{pushKernel\_iff\_comp} is
the intertwining square --- \texttt{Kernel.deterministic σ hσ ∘ₖ K = Q
∘ₖ Kernel.deterministic σ hσ}, push-after-step equals step-after-push
--- the shape the lens/fibred strengthening would functor over, so
that residue names a bound \emph{object}, not a blank. The finite
tier\textquotesingle s nonnegativity binders (\texttt{hK},
\texttt{hμ₀}) vanish at this tier by type, which is the ``signed mass
could cancel'' clause discharged by construction.
Width note: one-sidedness is structural in the bound statements ---
nothing consumes or constructs a joint kernel; \texttt{push\_law} (the
shared marginal-process lemma) is genuinely the same lemma the
relational tier rides, at both tiers, while the coupling machinery is
deliberately absent from this bound set. The conditional mark is
dropped: the former side condition is the baked $\mathsf{QSim}$
premise; what remains named --- the lex/fibred (categorical)
strengthening --- is a dependency of nothing on this page.

\ruleblock{Lean verifier receipts}
\sourcenote{module \texttt{BernsteinLogicAux.QtslCompose.SimTransport}; bound build \texttt{b2a09651f9a0\ldots} (bound from receipt dc0250279a3a2d0d, 2026-07-19T18:24:55Z). A receipt records that named checking layers accepted a build: (module, compiled-artifact SHA-256, environment SHA-256, checker layers), minted when the checking gate ran; the store receipts attesting the bound build are listed with the layers each run credited. Nothing here asks the reader to trust the store: every layer is an independently runnable program over the quoted sources --- rebuild the module and re-run the checkers to reproduce the verdict. Receipts attesting superseded builds credit no layer and are omitted from this rendering (the count of omissions, if any, is noted below).}
\begin{receipttable}
\receiptrow{edfcab6a34d8774b}{2026-07-19T03:38:57Z}{build+decorrelated}{\badge{BadgeGreen}{current}}
\receiptrow{dc0250279a3a2d0d}{2026-07-19T18:24:55Z}{build+decorrelated}{\badge{BadgeGreen}{current}}
\receiptrow{714abbd46c61ec9c}{2026-07-20T01:22:07Z}{build+decorrelated}{\badge{BadgeGreen}{current}}
\receiptrow{28aeb5d1779c6cdc}{2026-07-20T20:46:58Z}{build+decorrelated}{\badge{BadgeGreen}{current}}
\receiptrow{1bb39a35738bbfe3}{2026-07-20T20:58:31Z}{build+replay+hardened+decorrelated}{\badge{BadgeGreen}{current}}
\receiptrow{1ae45ad709a68e43}{2026-07-21T00:04:02Z}{build+decorrelated}{\badge{BadgeGreen}{current}}
\receiptrow{5b96bfe1e468ce59}{2026-07-21T05:39:17Z}{build+replay+hardened+decorrelated}{\badge{BadgeGreen}{current}}
\receiptrow{016e87c2921cd31d}{2026-07-22T18:27:17Z}{build+decorrelated}{\badge{BadgeGreen}{current}}
\receiptrow{e776a27adcccb860}{2026-07-31T22:00:16Z}{build+decorrelated}{\badge{BadgeGreen}{current}}
\receiptrow{cf62752ca2ea53ab}{2026-07-31T22:29:10Z}{build+decorrelated}{\badge{BadgeGreen}{current}}
\receiptrow{3c2ed3fc0dce1ac5}{2026-08-02T23:31:21Z}{build+decorrelated}{\badge{BadgeGreen}{current}}
\receiptrow{6121f42e978829a4}{2026-08-03T16:27:25Z}{build+decorrelated}{\badge{BadgeGreen}{current}}
\receiptrow{3a4a61220d562456}{2026-08-03T22:46:09Z}{build+decorrelated}{\badge{BadgeGreen}{current}}
\end{receipttable}
\sourcenote{module \texttt{BernsteinLogicAux.QtslCompose.SimTransportMeas}; bound build \texttt{9530b0b3c465\ldots} (bound from receipt 3c92f1e3a2986972, 2026-07-20T08:24:45Z). A receipt records that named checking layers accepted a build: (module, compiled-artifact SHA-256, environment SHA-256, checker layers), minted when the checking gate ran; the store receipts attesting the bound build are listed with the layers each run credited. Nothing here asks the reader to trust the store: every layer is an independently runnable program over the quoted sources --- rebuild the module and re-run the checkers to reproduce the verdict. Receipts attesting superseded builds credit no layer and are omitted from this rendering (the count of omissions, if any, is noted below).}
\begin{receipttable}
\receiptrow{3c92f1e3a2986972}{2026-07-20T08:24:45Z}{build+decorrelated}{\badge{BadgeGreen}{current}}
\receiptrow{06e7f179212d370d}{2026-07-20T20:47:00Z}{build+decorrelated}{\badge{BadgeGreen}{current}}
\receiptrow{c314290e315215f2}{2026-07-20T20:58:33Z}{build+replay+hardened+decorrelated}{\badge{BadgeGreen}{current}}
\receiptrow{c471feb792abf799}{2026-07-21T00:04:04Z}{build+decorrelated}{\badge{BadgeGreen}{current}}
\receiptrow{b18bbdfab8a60463}{2026-07-21T05:39:18Z}{build+replay+hardened+decorrelated}{\badge{BadgeGreen}{current}}
\receiptrow{16821f18670ae2eb}{2026-07-22T18:27:18Z}{build+decorrelated}{\badge{BadgeGreen}{current}}
\receiptrow{fb024d6489c5da75}{2026-07-31T22:00:16Z}{build+decorrelated}{\badge{BadgeGreen}{current}}
\receiptrow{ed5047f4fc559336}{2026-07-31T22:29:10Z}{build+decorrelated}{\badge{BadgeGreen}{current}}
\receiptrow{e618e69852b924d8}{2026-08-02T23:31:21Z}{build+decorrelated}{\badge{BadgeGreen}{current}}
\receiptrow{19820df5b08bd97a}{2026-08-03T16:27:25Z}{build+decorrelated}{\badge{BadgeGreen}{current}}
\receiptrow{682c99d152518997}{2026-08-03T22:46:09Z}{build+decorrelated}{\badge{BadgeGreen}{current}}
\end{receipttable}
\layerline{\layeritem{build}{lake build (elaboration)}{\badge{BadgeGreen}{current}}\layeritem{replay}{leanchecker (official kernel replay)}{\badge{BadgeGreen}{current}}\layeritem{hardened}{safe\_verify (hardened screen)}{\badge{BadgeGreen}{current}}\layeritem{decorrelated}{lean4export + nanoda (decorrelated kernel)}{\badge{BadgeGreen}{current}}}
\sourcenote{\badge{BadgeGray}{comparator: none} no fixture in the three-implementation comparator harness (compiled / kernel-interpreted / reference-interpreted) exercises this rule.}
\ruleblock{HOL4 statement}
\absentbox{NO HOL STATEMENT YET}{The satisfaction-judgment extension postdates the HOL4 port entirely:
the ported cone is the base calculus\textquotesingle s check-function
surface (theory \texttt{bern\_catalogue}), and no extension-theory
definition --- nothing shaped like simulation transport --- exists in
the HOL development. The box waits on a HOL port of the satisfaction
stratum, which is its own campaign.
}
\ruleblock{HOL verifier receipts}
\absentbox{NONE}{no HOL verification receipt exists for this rule.}
\ruleblock{Proof}

\sourcenote{Genre: core rule (tier 1). The rule's semantic content is
honestly small --- the behavioural pullback is two lines --- and the
section says so rather than inflating it; the substance of any
application is the sequent-level discharge of the compatibility
squares, plus the finite-kernel instance where the transport is proven
against actual chain machinery. The strata beyond the two walked here,
and the genuine gaps, are marked loudly.}

\textbf{Setting and claim.}\enspace A simulation witness
$\sigma : S \to S'$ is a per-mode polynomial map whose
\emph{compatibility squares} --- commutation of $\sigma$ with the two
systems' dynamics declarations, mode by mode --- are checked as
inherited sequent leaves; the rule consumes the witness as the
judgment $\mathsf{QSim}(\sigma, S, S')$, whose own discharge node
carries those squares. The rule: from
$S' \vDash \mathit{spec}'$ and the judgment, conclude
$S \vDash \sigma^{*}\mathit{spec}'$ --- certify the
abstraction once, reuse under refinement \cite{lynchvaandrager1995,
capuccimyers2026}.

\begin{bthm}[behavioral pullback]\label{thm:qst-pullback}
Let $\mathrm{Beh}$, $\mathrm{Beh}'$ be the compiled behaviour classes
of $S$, $S'$, and suppose $\sigma$ maps behaviours to behaviours
($\sigma[\mathrm{Beh}] \subseteq \mathrm{Beh}'$ --- the semantic
content of the compatibility squares). If every behaviour of $S'$
satisfies $\mathit{spec}'$, then every behaviour of $S$ satisfies
$\sigma^{-1}[\mathit{spec}']$.
\end{bthm}
\begin{proof}
$\omega \in \mathrm{Beh}$ gives
$\sigma(\omega) \in \mathrm{Beh}' \subseteq \mathit{spec}'$, i.e.\
$\omega \in \sigma^{-1}[\mathit{spec}']$
(\texttt{pullback}). Two lines is the honest size: the rule's
\emph{semantic} content is inclusion-chasing, and everything
substantive in an application is establishing
$\sigma[\mathrm{Beh}] \subseteq \mathrm{Beh}'$ from the checked
squares --- per-mode polynomial identities, discharged by the base
calculus. The hypothesis is not droppable to
``$\sigma$ is any map'': without the squares the image of a concrete
behaviour need not be an abstract behaviour, and the pullback claims
anything at all.
\end{proof}

\textbf{The finite-kernel instance.}\enspace At the finite Markov
tier the witness has teeth the behavioural form leaves abstract: a
\emph{functional simulation} is per-state marginal identity ---
pushing the concrete kernel row at $s$ along $\sigma$ yields the
abstract kernel row at $\sigma(s)$ (a pushforward-kernel witness) ---
and the transport theorems are proven against the chain machinery.

\begin{blemma}[observable pullback]\label{lem:qst-expect}
For any law $\mu$ on the concrete states and observable $g$ on the
abstract states,
$\sum_b (\sigma_{*}\mu)(b) \, g(b) = \sum_s \mu(s) \, g(\sigma(s))$;
and along a pushforward-kernel witness, the time-$n$ expectation of a
pulled-back observable on the concrete chain \emph{equals} the
abstract chain's expectation of $g$ at time $n$ (started from the
pushed initial law). Hence any certified expectation bound proven
against the abstract chain holds for the concrete chain at the
pulled-back observable.
\end{blemma}
\begin{proof}
The first identity is summation over fibres of $\sigma$
(\texttt{expect\_push}). The second composes it with the
marginal-process lemma --- the per-state marginal identities
propagate, by induction on time, to equality of pushed laws at every
time --- giving \texttt{sim\_expect}; the bound transports by
rewriting (\texttt{sim\_expect\_bound}). Note the direction: the
producer proves one bound on the small chain and every refinement
inherits it; nothing is re-proven per refinement beyond the witness
identities themselves. The interface reading rides the same
equalities: an assumption event $A'$ pulls back with its mass intact
(the pushed law at time $n$ \emph{equals} the abstract law, so
$\nu_n(\sigma^{-1}A') = \lambda_n(A')$), hence assumption-weighted
bounds transport, and conditional expectations transport whenever the
conditioning mass is positive.
\end{proof}

\begin{blemma}[region pullback]\label{lem:qst-region}
Along a pushforward-kernel witness with entrywise-nonnegative kernel
and initial law: if the abstract law at time $n$ puts no mass outside
$R'$, the concrete law at time $n$ puts no mass outside
$\sigma^{-1}[R']$.
\end{blemma}
\begin{proof}
For $s$ with $\sigma(s) \notin R'$: the concrete mass at $s$ is at
most the pushed mass at $\sigma(s)$ (one nonnegative term of a
fibre sum), which equals the abstract mass at $\sigma(s)$ (the
marginal-process lemma), which is zero; nonnegativity closes the
squeeze (\texttt{sim\_region}). The nonnegativity hypotheses are
load-bearing, not hygiene: with signed ``mass'' the fibre sum could
cancel --- a negative sibling hiding positive mass at $s$ outside the
preimage region --- and the single-term bound fails. Nonnegativity is
what makes ``no mass outside $R'$'' transfer pointwise through a
many-to-one map.
\end{proof}

\begin{bremark}[refinement is one-sided; a simulation is not a
coupling]\label{rem:qst-onesided}
The witness machinery genuinely shared with the relational
(two-chain) tier is exactly the pushforward-kernel identity and the
marginal-process lemma --- a functional simulation \emph{is} such a
witness. Nothing else transfers: the relational tier's certificates
are joint-kernel data over a product space with absorption and drift
hypotheses, and none of that is simulation transport. A functional
simulation induces a degenerate coupling, but the theorems above
never route through it --- they are pushforward plus summation
identities, full stop. Conflating the two would import hypotheses
(coupling emptiness, absorption) that this rule neither needs nor
checks.
\end{bremark}

\begin{bremark}[the gaps, loud]\label{rem:qst-gaps}
The section walks two strata --- behavioural, and finite-kernel. The
general measurable-kernel tier is mechanized and bound in this
rule's entry, mirroring the finite instance name-for-name over
Markov kernels (the same pushforward identity per state, now a
\texttt{Measure.map} equation; law transport at every horizon; both
spec-transport faces --- with the finite tier's nonnegativity
hypotheses gone, true by type); it is cited there rather than
re-walked here. Not yet mechanized: the categorical strengthening
under which the composition rule could cite the lens-composition
theorem of
\cite{capuccimyers2026} as a lemma \emph{template} --- the direct
composition proofs never pull back along strict weakenings, so
nothing in this book depends on that strengthening; it is named so
that its absence is a recorded choice, not an oversight. The
judgment kinds the rule consumes and concludes into are staged in
the certificate schema (the simulation judgment together with its
discharge and composition nodes), but no end-to-end certificate yet
exercises the rule --- loader and end-to-end work, not mathematics. Related
reading, not dependency: an exchangeability quotient is a simulation
map, so aggregate-space reduction can be \emph{read} as an instance
of this rule; it stands on its own proof and cites nothing here.
\end{bremark}

\begin{anchors}
Theorem \ref{thm:qst-pullback} & \texttt{pullback}\\
Lemma \ref{lem:qst-expect} & \texttt{expect\_push}, \texttt{sim\_expect}, \texttt{sim\_expect\_bound}; marginal-process lemma \texttt{push\_law} (chain substrate module)\\
Lemma \ref{lem:qst-region} & \texttt{sim\_region}; law nonnegativity \texttt{law\_nonneg} (chain substrate module)\\
witness form & \texttt{PushKernel} (chain substrate module) --- the per-state marginal identity\\
compatibility squares (cited, not re-proved) & inherited sequent leaves, per mode (main package)\\
\end{anchors}


\rulesection{82}{QFlatten}{admissible rule}
\label{blcatrule-QFlatten}
\ruleblock{The rule (Steele metanotation)}
\sourcenote{as realized in the mechanized development; source row: \texttt{\texttt{QTSL-rules.toml}} --- \emph{QFlatten}. The metanotation statement is typeset by hand for notational quality; its mechanized ground truth is the Lean statement below, extracted from the sources at every build of this book and hash-bound. The rule \emph{roster} is machine-checked (every mechanized rule tag round-trips to exactly one catalogue row); the statement \emph{text} is editorial, checked by eye against the extracted Lean beside it.}
\begin{mathpar}
\inferrule*[Right=\rulename{QFlatten}]{
  \mathsf{QFlatEq}(\mathrm{sys}, \mathit{flat}) \\
  \mathcal{M}_{\mathit{flat}},\ x_0 \in R \vDash \varphi \ \text{(the closed-system pipeline\textquotesingle s judgment)}
}{
  \mathrm{sys},\ x_0 \in R \vDash \varphi
}
\end{mathpar}

\ruleflavor{$\mathit{flat}$ is the flattening of the composite $\mathrm{sys}$; the substitution rule it always was: satisfaction transports across the flattening equation. The former side condition (polynomial substitution + guard-product bookkeeping) is the kind\textquotesingle s discharge --- oracle-shaped by construction (\texttt{qtsl\_flatten}, a bounded deterministic recompute over declaration structure), with the refusal boundary stated positively as the judgment\textquotesingle s scope: shared-carrier confinement, and binary wiring with $n$-ary/directed cases reducing by currying = the checker\textquotesingle s substitution recompute (refusal-shaped completeness cost, never a soundness hypothesis). The guarantee the rule was minted to state --- $\mathsf{SysWire}$ adds structure, never expressiveness --- is carried by the kind\textquotesingle s existence: the flattened manifest\textquotesingle s compiled semantics is the composite\textquotesingle s (theorem at both kernel tiers on the synchronous-product face), and the projections remain free \rulename{QSimTransport} witnesses under sibling well-formedness (unit row sums). The clock rows (schedule hypotheses, simultaneity, non-Zeno) are premise classes cited by every composition rule, deliberately not rows here.}
\par\medskip
\rulename{QFlatten} transports satisfaction across the flattening
equation: given the judgment $\mathsf{QFlatEq}(\mathrm{sys},
\mathit{flat})$ --- $\mathit{flat}$ is the flattening of the composite
presentation, discharged by the checker\textquotesingle s substitution
recompute --- a verdict on the monolithic manifest is a verdict on the
composite. A
producer reaches for it as the bridge to every closed-system pipeline:
compile nodes, certificate synthesis, and the whole battery machinery
operate on one flat manifest, while modelling happens compositionally.
The rule\textquotesingle s two guarantees pull in opposite directions
and are both load-bearing: wiring adds \emph{structure}, never
\emph{expressiveness}. Never expressiveness: the flattened
manifest\textquotesingle s compiled semantics is exactly the
composite\textquotesingle s --- at the compiled class this is
definitional (a presentation\textquotesingle s semantics IS its
flattened object\textquotesingle s), and at the finite-kernel tier it
is a theorem: the monolithic chain\textquotesingle s law is, at every
horizon, the product of the component laws. Added structure: the
flattened manifest remembers its components --- the projections are
exact simulation witnesses for free, so every component-level
specification remains addressable on the flattened object through
\rulename{QSimTransport}.

What the judgment\textquotesingle s discharge recomputes: the
substitution itself --- polynomial
substitution plus guard-product bookkeeping --- which is why the rule is
admissible rather than core: it discharges against verified
recomputation and the semantic witness pair, adding no new kernel
content. Directed substitution wirings do \emph{not} reduce to the
independent synchronous product of the witness pair --- wiring an
upstream output into a downstream input creates conditional dependence
the product law does not carry; they rest on the checker\textquotesingle
s substitution recompute itself (currying the upstream trajectory into
a dependent downstream kernel), definitional in the compiled class.

What would break: component well-formedness (unit row sums) is what
makes the projections exact witnesses --- against a defective second
factor, the first projection\textquotesingle s pushed row picks up the
leftover mass as a spurious scale factor, and the component-spec
addressability claim fails. The scheduling hypotheses (clock rows,
simultaneity, non-Zeno families) are deliberately premise classes cited
by every composition rule rather than rows of this one.

\cataloguenote{The flattening premise is the judgment kind $\mathsf{QFlatEq}$,
baked in the extension theory; its discharge is the flatten
recompute behind a registry seam, with the shared-carrier and
binary-wiring confinements as the judgment\textquotesingle s stated
scope --- refusal-shaped, never a hypothesis of this rule.
The proven fragment spans two tiers: the finite-kernel tier, on the
same substrate as the relational chain machinery --- the tier where
\textquotedblleft law\textquotedblright, \textquotedblleft
product\textquotedblright, and \textquotedblleft
pushforward\textquotedblright{} are finite exact-arithmetic objects ---
and the general measurable-kernel tier, where they
are the Giry-monad objects (second bound block). The remaining
confinement is the \emph{independent binary product} witness pair --- a
design fence, not a gap: $n$-ary independent wirings iterate it, while
directed wirings rest on the checker\textquotesingle s substitution
recompute (a dependent downstream kernel, not the independent product).
The general compiled-class statement is substitution correctness,
checker recomputation by design.
}
\ruleblock{Lean statement}
\declhead{BernsteinLogicAux.QtslCompose.Flatten}{law\_prodKernel}{\badge{BadgeGreen}{statement bound}}
\begin{leansrc}
/-- **The flatten witness** (`QFlatten`'s "never expressiveness" face): the
flattened monolithic chain's law is, at every time, the product of the
component chains' laws — the composite presentation and its monolithic
manifest compile to the same semantics. -/
theorem law_prodKernel (K₁ : α → α → ℝ) (K₂ : β → β → ℝ)
    (μ₁ : α → ℝ) (μ₂ : β → ℝ) :
    ∀ n, law (prodKernel K₁ K₂) (prodDist μ₁ μ₂) n
      = prodDist (law K₁ μ₁ n) (law K₂ μ₂ n)
  | 0 => rfl
  | n + 1 => by
    rw [law_succ, law_prodKernel K₁ K₂ μ₁ μ₂ n, step_prodKernel, law_succ,
      law_succ]

section Projections

variable [DecidableEq α] [DecidableEq β]

/-- The added *structure*: the first projection is a simulation witness
(`QSimTransport.PushKernel`) from the flattened manifest onto component 1 —
for free, whenever component 2 is well-formed. -/
theorem pushKernel_fst {K₁ : α → α → ℝ} {K₂ : β → β → ℝ}
    (h₂ : ∀ b, ∑ b', K₂ b b' = 1) :
    PushKernel Prod.fst (prodKernel K₁ K₂) K₁
\end{leansrc}
\declhead{BernsteinLogicAux.QtslCompose.Flatten}{pushKernel\_fst}{\badge{BadgeGreen}{statement bound}}
\begin{leansrc}
/-- The added *structure*: the first projection is a simulation witness
(`QSimTransport.PushKernel`) from the flattened manifest onto component 1 —
for free, whenever component 2 is well-formed. -/
theorem pushKernel_fst {K₁ : α → α → ℝ} {K₂ : β → β → ℝ}
    (h₂ : ∀ b, ∑ b', K₂ b b' = 1) :
    PushKernel Prod.fst (prodKernel K₁ K₂) K₁
\end{leansrc}
\declhead{BernsteinLogicAux.QtslCompose.Flatten}{pushKernel\_snd}{\badge{BadgeGreen}{statement bound}}
\begin{leansrc}
/-- Mirror: the second projection simulates onto component 2. -/
theorem pushKernel_snd {K₁ : α → α → ℝ} {K₂ : β → β → ℝ}
    (h₁ : ∀ a, ∑ a', K₁ a a' = 1) :
    PushKernel Prod.snd (prodKernel K₁ K₂) K₂
\end{leansrc}
\declhead{BernsteinLogicAux.QtslCompose.FlattenMeas}{law\_parallelComp}{\badge{BadgeGreen}{statement bound}}
\begin{leansrc}
/-- **QFlatten, measurable-kernel tier — the flatten witness** ("never
expressiveness" face): the flattened monolithic chain's law is, at every
horizon, the product of the component chains' laws — the composite
presentation and its monolithic manifest compile to the same semantics.
(The measurable-tier face of `law_prodKernel`.) -/
theorem law_parallelComp (μ : Measure α) (ν : Measure β)
    (K : Kernel α α) (Q : Kernel β β)
    [SFinite μ] [SFinite ν] [IsSFiniteKernel K] [IsSFiniteKernel Q] :
    ∀ n, law (K ∥ₖ Q) (μ.prod ν) n = (law K μ n).prod (law Q ν n)
  | 0 => rfl
  | n + 1 => by
    rw [law_succ, law_parallelComp μ ν K Q n, bind_parallelComp,
      ← law_succ, ← law_succ]

/-! ## The added structure: projections are simulation witnesses -/

/-- The added *structure* ("adds structure" face): the first projection is a
simulation witness (`SimTransportMeas.PushKernel`) from the flattened
manifest onto component 1 — for free, consuming exactly the *sibling's*
Markov property (`[IsMarkovKernel Q]`: unit total mass per row) and nothing
of `K`'s own well-formedness. Margin-zero boundary, displayed by
`Measure.map_fst_prod`: against a defective sibling the pushed row is
`(Q p.2 univ) • K p.1` — the leak coefficient is the sibling's total mass.
(The measurable-tier face of the finite tier's `pushKernel_fst`.) -/
theorem pushKernel_fst {K : Kernel α α} {Q : Kernel β β}
    [IsSFiniteKernel K] [IsMarkovKernel Q] :
    PushKernel Prod.fst (K ∥ₖ Q) K
\end{leansrc}
\declhead{BernsteinLogicAux.QtslCompose.FlattenMeas}{pushKernel\_fst}{\badge{BadgeGreen}{statement bound}}
\begin{leansrc}
/-- The added *structure* ("adds structure" face): the first projection is a
simulation witness (`SimTransportMeas.PushKernel`) from the flattened
manifest onto component 1 — for free, consuming exactly the *sibling's*
Markov property (`[IsMarkovKernel Q]`: unit total mass per row) and nothing
of `K`'s own well-formedness. Margin-zero boundary, displayed by
`Measure.map_fst_prod`: against a defective sibling the pushed row is
`(Q p.2 univ) • K p.1` — the leak coefficient is the sibling's total mass.
(The measurable-tier face of the finite tier's `pushKernel_fst`.) -/
theorem pushKernel_fst {K : Kernel α α} {Q : Kernel β β}
    [IsSFiniteKernel K] [IsMarkovKernel Q] :
    PushKernel Prod.fst (K ∥ₖ Q) K
\end{leansrc}
\declhead{BernsteinLogicAux.QtslCompose.FlattenMeas}{pushKernel\_snd}{\badge{BadgeGreen}{statement bound}}
\begin{leansrc}
/-- Mirror: the second projection simulates onto component 2, consuming
exactly the first component's Markov property. -/
theorem pushKernel_snd {K : Kernel α α} {Q : Kernel β β}
    [IsMarkovKernel K] [IsSFiniteKernel Q] :
    PushKernel Prod.snd (K ∥ₖ Q) Q
\end{leansrc}
\declhead{BernsteinLogicAux.QtslCompose.FlattenMeas}{bind\_parallelComp}{\badge{BadgeGreen}{statement bound}}
\begin{leansrc}
/-- One step of the flattened (synchronous-product) chain factors as the
product of the component steps. (The measurable-tier face of
`step_prodKernel`.) Proof is pure `∘ₘ`-algebra:
`Measure.prod_comp_right`/`prod_comp_left` are the two one-sided
measure-level product-bind identities, fused by
`Kernel.parallelComp_comp_parallelComp`. -/
theorem bind_parallelComp (μ : Measure α) (ν : Measure β)
    (K : Kernel α α) (Q : Kernel β β)
    [SFinite μ] [SFinite ν] [IsSFiniteKernel K] [IsSFiniteKernel Q] :
    (μ.prod ν).bind (K ∥ₖ Q) = (μ.bind K).prod (ν.bind Q)
\end{leansrc}
\ruleblock{Correspondence (Steele display $\leftrightarrow$ Lean)}
\badge{BadgeGreen}{reviewer-certified against the displayed texts}
\sourcenote{The dictionary below is a reviewer-certified statement of how the mechanized statements above interpret the displayed rule: display metavariable $\mapsto$ Lean binder, each display hypothesis $\mapsto$ the premise or side condition carrying it, the display conclusion $\mapsto$ the goal --- with any deliberate difference in width between display and mechanization named outright. Certification is recorded as a hash triple over (the displayed cell, the extracted statement texts, this dictionary\textquotesingle s own text); if any of the three changes, this block turns amber until a reviewer re-certifies. The correspondence claim itself remains editorial --- reviewer-grade, not machine-checked; its machine-checked successor types each soundness statement against the rule as a structured object, and supersedes this block per rule as it lands.}
The displayed judgment premise $\mathsf{QFlatEq}(\mathrm{sys},
\mathit{flat})$ --- the flattened manifest has the
\emph{same} compiled semantics, each projection $\pi_i$ a free
\rulename{QSimTransport} witness --- $\mapsto$ the theorem pair at the
finite-kernel tier: \texttt{law\_prodKernel} (never expressiveness)
and \texttt{pushKernel\_fst}/\texttt{\_snd} (added structure). The
kind\textquotesingle s discharge (the substitution recompute) has
\emph{no} Lean
counterpart by design: at the compiled class the
presentation\textquotesingle s semantics IS its flattened
object\textquotesingle s, so that discharge is checker recomputation,
and what is theorem-shaped is exactly what is bound. The
rule\textquotesingle s own transport step (a verdict on
$\mathcal{M}_{\mathit{flat}}$ is a verdict on $\mathrm{sys}$) is
substitution of equals along the judgment --- no separate statement
to bind.
\begin{itemize}
\item $\mathsf{SysWire}(w, \mathit{Sys}_1, \ldots, \mathit{Sys}_n)$
  $\mapsto$ the \emph{binary synchronous product}: kernels
  \texttt{K₁ : α → α → ℝ}, \texttt{K₂ : β → β → ℝ} over finite types,
  flattened as \texttt{prodKernel} on \texttt{α × β}. Two confinements
  at this block, both design fences: binary-independent (directed and
  $n$-ary substitution wirings rest on the checker\textquotesingle s
  substitution recompute --- currying the upstream trajectory into a
  \emph{dependent} downstream kernel, which is not the independent
  product bound here; per the module header), and
  finite-kernel tier (\texttt{[Fintype α] [Fintype β]}) --- the latter
  discharged by the measurable-tier block below, leaving the
  independent binary product as this bound set\textquotesingle s
  remaining confinement.
\item same compiled semantics $\mapsto$ \texttt{law\_prodKernel}:
  \texttt{law (prodKernel K₁ K₂) (prodDist μ₁ μ₂) n = prodDist (law K₁
  μ₁ n) (law K₂ μ₂ n)} at every \texttt{n} --- and its binder row is
  the width fact: the identity is proven for \emph{arbitrary} kernel
  and law data, no Markov/distribution well-formedness consumed ---
  pure bilinear algebra, strictly wider than the display\textquotesingle s
  well-formed-system reading, safe direction.
\item projections as free \rulename{QSimTransport} witnesses $\mapsto$
  \texttt{pushKernel\_fst : PushKernel Prod.fst (prodKernel K₁ K₂) K₁}
  --- consuming only \texttt{∀ b, ∑ b', K₂ b b' = 1}: the
  \emph{sibling} component\textquotesingle s unit row sums, nothing of
  its own, and not even nonnegativity --- the certified statement
  isolates exactly which well-formedness fragment the
  commentary\textquotesingle s ``what would break'' paragraph names
  (a defective second factor leaks scale into the pushed row).
  \texttt{\_snd} mirrors. \texttt{[DecidableEq]} binders enter only
  here (the pushforward\textquotesingle s fibre sums).
\end{itemize}
\emph{Measurable-kernel tier} (second bound block, module
\texttt{FlattenMeas}, namespace \texttt{…QtslCompose.Meas}): the
product is mathlib\textquotesingle s own pair
--- \texttt{K ∥ₖ Q} (\texttt{Kernel.parallelComp}) over
\texttt{Measure.prod} --- and the theorem pair transposes
name-for-name:
\texttt{law\_parallelComp} (never expressiveness: \texttt{law (K ∥ₖ Q)
(μ.prod ν) n = (law K μ n).prod (law Q ν n)} at every horizon, over
\texttt{SimTransportMeas}\textquotesingle s Giry-bind \texttt{law}) and
\texttt{pushKernel\_fst}/\texttt{\_snd} (added structure: projections
as \texttt{Meas.PushKernel} witnesses), with \texttt{bind\_parallelComp}
the one-step glue. Binder floor, the honest analogue of the finite
tier\textquotesingle s ``arbitrary data'' width fact:
\texttt{SFinite}/\texttt{IsSFiniteKernel} --- product-measure theory is
s-finite theory --- still strictly wider than the
display\textquotesingle s Markov/probability reading, same safe
direction. And the ``what would break'' margin-zero witness is
displayed by mathlib itself: \texttt{Measure.map\_fst\_prod} gives the
pushed row as \texttt{(ν univ) • μ} --- leak coefficient the
sibling\textquotesingle s total mass --- with \texttt{pushKernel\_fst}
consuming exactly \texttt{[IsMarkovKernel Q]} and nothing of
\texttt{K}\textquotesingle s own well-formedness, not even finiteness.
Width note: the clock rows (schedule/simultaneity/non-Zeno) are
premise classes deliberately absent from both display and bound
statements; the closure facts at both tiers live in their modules but
outside this bound set --- cited, not certified here: finite
\texttt{isMarkov\_prodKernel}, \texttt{isDist\_prodDist},
\texttt{push\_fst\_prodDist}; measurable \texttt{sfinite\_law},
\texttt{isProbabilityMeasure\_law}, the addressability compositions
\texttt{sim\_expect\_fst}/\texttt{\_snd} (component specs consumed
end-to-end through \rulename{QSimTransport}\textquotesingle s
measurable engine), the bridge \texttt{pushKernel\_iff\_stochSim}, and
the mathlib instances \texttt{IsMarkovKernel (K ∥ₖ Q)} /
\texttt{prod.instIsProbabilityMeasure}.

\ruleblock{Lean verifier receipts}
\sourcenote{module \texttt{BernsteinLogicAux.QtslCompose.Flatten}; bound build \texttt{30ed43793772\ldots} (bound from receipt 8cd9a31cf76c1595, 2026-07-19T18:24:49Z). A receipt records that named checking layers accepted a build: (module, compiled-artifact SHA-256, environment SHA-256, checker layers), minted when the checking gate ran; the store receipts attesting the bound build are listed with the layers each run credited. Nothing here asks the reader to trust the store: every layer is an independently runnable program over the quoted sources --- rebuild the module and re-run the checkers to reproduce the verdict. Receipts attesting superseded builds credit no layer and are omitted from this rendering (the count of omissions, if any, is noted below).}
\begin{receipttable}
\receiptrow{e4496ce5107a3441}{2026-07-19T03:38:56Z}{build+decorrelated}{\badge{BadgeGreen}{current}}
\receiptrow{8cd9a31cf76c1595}{2026-07-19T18:24:49Z}{build+decorrelated}{\badge{BadgeGreen}{current}}
\receiptrow{62e3b3d615f78f37}{2026-07-20T01:22:02Z}{build+decorrelated}{\badge{BadgeGreen}{current}}
\receiptrow{b45a913927ac5488}{2026-07-20T20:46:52Z}{build+decorrelated}{\badge{BadgeGreen}{current}}
\receiptrow{d52960f0b18dc76a}{2026-07-20T20:58:26Z}{build+replay+hardened+decorrelated}{\badge{BadgeGreen}{current}}
\receiptrow{0eabb0d2fdfbc00c}{2026-07-21T00:03:55Z}{build+decorrelated}{\badge{BadgeGreen}{current}}
\receiptrow{1f0f373238ef2c2b}{2026-07-21T05:39:11Z}{build+replay+hardened+decorrelated}{\badge{BadgeGreen}{current}}
\receiptrow{35d0f3c7c71856fc}{2026-07-22T18:27:16Z}{build+decorrelated}{\badge{BadgeGreen}{current}}
\receiptrow{53dde5ffb738ab3a}{2026-07-31T22:00:14Z}{build+decorrelated}{\badge{BadgeGreen}{current}}
\receiptrow{50f8cfda2d25e116}{2026-07-31T22:29:08Z}{build+decorrelated}{\badge{BadgeGreen}{current}}
\receiptrow{edc0e93a8721d195}{2026-08-02T23:31:19Z}{build+decorrelated}{\badge{BadgeGreen}{current}}
\receiptrow{a6fae671c06c5ebb}{2026-08-03T16:27:23Z}{build+decorrelated}{\badge{BadgeGreen}{current}}
\receiptrow{0427fd20db259dcd}{2026-08-03T22:46:07Z}{build+decorrelated}{\badge{BadgeGreen}{current}}
\end{receipttable}
\sourcenote{module \texttt{BernsteinLogicAux.QtslCompose.FlattenMeas}; bound build \texttt{40d14da6aacd\ldots} (bound from receipt 4747724a289473cd, 2026-07-20T08:24:46Z). A receipt records that named checking layers accepted a build: (module, compiled-artifact SHA-256, environment SHA-256, checker layers), minted when the checking gate ran; the store receipts attesting the bound build are listed with the layers each run credited. Nothing here asks the reader to trust the store: every layer is an independently runnable program over the quoted sources --- rebuild the module and re-run the checkers to reproduce the verdict. Receipts attesting superseded builds credit no layer and are omitted from this rendering (the count of omissions, if any, is noted below).}
\begin{receipttable}
\receiptrow{4747724a289473cd}{2026-07-20T08:24:46Z}{build+decorrelated}{\badge{BadgeGreen}{current}}
\receiptrow{2a11a368eea6adc5}{2026-07-20T20:46:54Z}{build+decorrelated}{\badge{BadgeGreen}{current}}
\receiptrow{d11ebe966df6a113}{2026-07-20T20:58:27Z}{build+replay+hardened+decorrelated}{\badge{BadgeGreen}{current}}
\receiptrow{ba9ffebde141f18b}{2026-07-21T00:03:56Z}{build+decorrelated}{\badge{BadgeGreen}{current}}
\receiptrow{13a1681b1015f8c2}{2026-07-21T05:39:13Z}{build+replay+hardened+decorrelated}{\badge{BadgeGreen}{current}}
\receiptrow{89eb1271b74a3f48}{2026-07-22T18:27:16Z}{build+decorrelated}{\badge{BadgeGreen}{current}}
\receiptrow{d86334f5a6f3694d}{2026-07-31T22:00:15Z}{build+decorrelated}{\badge{BadgeGreen}{current}}
\receiptrow{4aceb786dcde76f2}{2026-07-31T22:29:09Z}{build+decorrelated}{\badge{BadgeGreen}{current}}
\receiptrow{e040a1347e8cd2a0}{2026-08-02T23:31:20Z}{build+decorrelated}{\badge{BadgeGreen}{current}}
\receiptrow{52c94a6a75909deb}{2026-08-03T16:27:23Z}{build+decorrelated}{\badge{BadgeGreen}{current}}
\receiptrow{0a269083e637c33e}{2026-08-03T22:46:07Z}{build+decorrelated}{\badge{BadgeGreen}{current}}
\end{receipttable}
\layerline{\layeritem{build}{lake build (elaboration)}{\badge{BadgeGreen}{current}}\layeritem{replay}{leanchecker (official kernel replay)}{\badge{BadgeGreen}{current}}\layeritem{hardened}{safe\_verify (hardened screen)}{\badge{BadgeGreen}{current}}\layeritem{decorrelated}{lean4export + nanoda (decorrelated kernel)}{\badge{BadgeGreen}{current}}}
\sourcenote{\badge{BadgeGray}{comparator: none} no fixture in the three-implementation comparator harness (compiled / kernel-interpreted / reference-interpreted) exercises this rule.}
\ruleblock{HOL4 statement}
\absentbox{NO HOL STATEMENT YET}{The HOL4 development checks the base calculus\textquotesingle s wire
format (theory \texttt{bern\_catalogue}). \rulename{QFlatten} operates
on system presentations of the extension theory --- baked, System
stratum and all; its semantic witnesses live in the
auxiliary Lean package. No HOL development consumes the extension
theory\textquotesingle s rows; the box becomes writable when a
flatten recompute joins a ported checker chain.
}
\ruleblock{HOL verifier receipts}
\absentbox{NONE}{no HOL verification receipt exists for this rule.}
\ruleblock{Proof}

\sourcenote{Genre: admissible rule (tier 2) --- a proof-tree-rewriting
recipe of the reduction-to-verified-machinery kind: the node discharges
against the checker's substitution recompute plus a semantic witness
pair, not against catalogue nodes. Scope, stated loudly: the witness
pair is proven at the finite-kernel tier, on the synchronous-product
wiring; the general compiled-class statement is substitution
correctness, definitional by the presentation semantics.}

\begin{brecipe}[the \rulename{QFlatten} rewrite]\label{rec:qfl-rewrite}
A \rulename{QFlatten} node transports satisfaction across the
flattening judgment --- $\mathit{flat}$ the monolithic manifest of a
composite presentation:
\[
\inferrule*[Right=\rulename{QFlatten}]{\;\mathsf{QFlatEq}(\mathrm{sys}, \mathit{flat}) \\ \mathcal{M}_{\mathit{flat}},\ x_0 \in R \vDash \varphi\;}{\mathrm{sys},\ x_0 \in R \vDash \varphi}
\qquad\rightsquigarrow\qquad
\text{\upshape the checker replays the substitution.}
\]
The rewrite target contains zero catalogue nodes: a presentation's
compiled semantics \emph{is} its flattened object's, so at the compiled
class the judgment's discharge is a recomputation --- polynomial
substitution
plus guard-product bookkeeping --- and the rule adds no kernel semantic
content (the admissibility criterion). What deserves proof, and gets it
below, is the slogan the rule is trusted for: \textbf{wiring adds
structure, never expressiveness}. Directed substitution wirings do
\emph{not} reduce to the independent synchronous product treated below
--- wiring an upstream output into a downstream input creates
conditional dependence the product law does not carry --- so the
flatten witness below attests the independent-product case, and the
directed case rests on the substitution recompute (the judgment's
discharge above) alone.
\end{brecipe}

\begin{bprop}[never expressiveness: the flatten witness]\label{prop:qfl-law}
Let $K_1, K_2$ be finite kernels with initial laws $\mu_1, \mu_2$, and
let the flattened manifest be the product kernel
$K_{\otimes}\bigl((a,b), (a',b')\bigr) = K_1(a, a')\, K_2(b, b')$ with
initial law $(\mu_1 \otimes \mu_2)(a, b) = \mu_1(a)\, \mu_2(b)$. Then
at every horizon $n$,
\[
\mathrm{law}_{K_\otimes}\bigl(\mu_1 \otimes \mu_2,\ n\bigr)
\;=\; \mathrm{law}_{K_1}(\mu_1, n) \,\otimes\, \mathrm{law}_{K_2}(\mu_2, n):
\]
the monolithic chain's law is, at every time, exactly the product of
the component laws --- the composite presentation and its flattened
manifest compile to the same semantics.
\end{bprop}
\begin{proof}
One step factors, by Fubini for finite sums and bilinearity:
\[
\mathrm{step}_{K_\otimes}(\mu_1 \otimes \mu_2)(a', b')
= \sum_{a, b} \mu_1(a) \mu_2(b)\, K_1(a, a') K_2(b, b')
= \Bigl(\sum_a \mu_1(a) K_1(a, a')\Bigr)\Bigl(\sum_b \mu_2(b) K_2(b, b')\Bigr),
\]
i.e.\ $\mathrm{step}_{K_\otimes}(\mu_1 \otimes \mu_2) =
\mathrm{step}_{K_1}\mu_1 \otimes \mathrm{step}_{K_2}\mu_2$. Induction
on $n$ closes. Closure rides along: the product of kernels is a kernel
and the product of laws is a law (nonnegativity termwise; unit masses
by the same bilinear factorization ---
$\sum_{a,b} \mu_1(a)\,\mu_2(b) = 1 \cdot 1$ for the product law,
$\sum_{a',b'} K_1(a, a')\,K_2(b, b') = 1 \cdot 1$ per product row).
\end{proof}

\begin{bprop}[added structure: the projections are free simulation
witnesses]\label{prop:qfl-proj}
If $K_2$ is row-stochastic, the first projection is an \emph{exact}
intertwining witness from $K_\otimes$ onto $K_1$:
$\pi_{1\#} K_\otimes\bigl((a,b), \cdot\bigr) = K_1(a, \cdot)$ for every
product state $(a, b)$; symmetrically for $\pi_2$ under
row-stochastic $K_1$; and $\pi_{1\#}(\mu_1 \otimes \mu_2) = \mu_1$
whenever $\mu_2$ has unit mass.
\end{bprop}
\begin{proof}
$\pi_{1\#} K_\otimes\bigl((a,b), \cdot\bigr)(a') =
\sum_{b'} K_1(a, a') K_2(b, b') = K_1(a, a') \sum_{b'} K_2(b, b') =
K_1(a, a')$, the last step by the unit row sum. The initial-law
statement is the same computation at $\mu$. Consequence: the flattened
manifest \emph{remembers} its components --- each projection is a free
exact-intertwining (\texttt{PushKernel}) witness, the
\rulename{QSimTransport} witness form, instantiating the strict
(zero-degradation) column of the \rulename{QDistTransport} matrix ---
so every component-level specification remains addressable on the
flattened object with zero degradation. $\mathsf{SysWire}$ adds these
handles and nothing else: Proposition~\ref{prop:qfl-law} caps
expressiveness, this proposition delivers the structure.
\end{proof}

\begin{bremark}[why the well-formedness hypotheses are not decoration]
Both directions consume them. In Proposition~\ref{prop:qfl-proj}, drop
the unit row sum of the \emph{other} factor and the projection witness
fails: the pushed row picks up $\sum_{b'} K_2(b, b')$ as a spurious
scale factor, so a sub-stochastic (leaky) second component would make
every component-1 judgment transport against a deflated law --- the
component-spec addressability claim silently weakens. In
Proposition~\ref{prop:qfl-law}, unit masses are what make the product
law a probability law at every time. The scheduling hypotheses of
composite dynamics --- clock rows, simultaneity, non-Zeno families ---
are deliberately \emph{premise classes} cited by every composition rule
rather than rows of this one: flattening is about state and kernel
structure, and importing the clock discipline here would duplicate its
single authoritative statement.
\end{bremark}

\begin{bremark}[the tier boundary, honestly]
The propositions above are the synchronous-product instance at the
finite-kernel tier --- the tier where law, product, and pushforward are
finite exact-arithmetic objects, shared with the relational-chain
machinery. General wirings (directed substitution, guard products)
reduce by the checker's own recompute; the reduction is definitional in
the compiled class and is not re-proven as a theorem here. Nothing
above claims the measure/kernel tier; its flatten-commutes statements
--- the same theorem pair over Markov kernels and product measures ---
are mechanized and bound in this rule's entry, cited there rather than
re-walked here.
\end{bremark}

\begin{anchors}
Recipe \ref{rec:qfl-rewrite} judgment discharge & the substitution recompute (polynomial substitution + guard-product bookkeeping; checker-side, behind the flattening judgment's registry seam)\\
Proposition \ref{prop:qfl-law} & \texttt{prodKernel}; \texttt{prodDist}; \texttt{step\_prodKernel}; \texttt{law\_prodKernel}; closure \texttt{isMarkov\_prodKernel}, \texttt{isDist\_prodDist}\\
Proposition \ref{prop:qfl-proj} & \texttt{pushKernel\_fst}; \texttt{pushKernel\_snd}; \texttt{push\_fst\_prodDist}\\
consumed by & \rulename{QSimTransport} (the \texttt{PushKernel} witness form); the strict column of \rulename{QDistTransport}'s matrix (\texttt{strict\_cell\_exact})\\
\end{anchors}


\rulesection{83}{QLinGrowthDischarge}{core rule}
\label{blcatrule-QLinGrowthDischarge}
\ruleblock{The rule (Steele metanotation)}
\sourcenote{as realized in the mechanized development; source row: \texttt{\texttt{QTSL-rules.toml}} --- \emph{QLinGrowthDischarge}. The metanotation statement is typeset by hand for notational quality; this rule cites no Lean statements yet (the absence is stated loudly below). The rule \emph{roster} is machine-checked (every mechanized rule tag round-trips to exactly one catalogue row); the statement \emph{text} is editorial.}
\begin{mathpar}
\inferrule*{
  \overline{\mathrm{env}_i : B \vdash\ \pm f_i(x) \ \le\ a + b\,\lVert x \rVert}
}{
  \vDash \mathsf{QLinGrowth}(f,\ a,\ b,\ B)
}
\end{mathpar}

\par\medskip
\rulename{QLinGrowthDischarge} concludes the linear-growth-envelope
kind with the constants carried: membership of a polynomial vector
field in the envelope $(a, b)$ over a scope box, from per-coordinate
affine-envelope sign sequents --- exact-rational, checker-replayed.
Downstream cells consume the pair $(a, b)$ arithmetically, never a
membership bit: the kind exists precisely so that growth constants
compose.

\cataloguenote{Entered at the roster census (baked, never catalogued); no Lean
statement binds it yet. The displayed norm form is schematic --- the
envelope rows\textquotesingle{} own sequent shape is the authority.
}
\ruleblock{Theory rows (geolog)}
\sourcenote{The rule\textquotesingle s certificate content as it is declared in the geolog theory the checker loads --- payload schema, side-condition demands, and axiom rows, quoted from the theory source at every build of this catalogue (a stale quotation cannot exist). Axiom modes: \texttt{chase} = the engine derives; \texttt{report} = a violation is a loud judgment; \texttt{refuse} = a violation gates admission. Side-condition relations are discharged by bounded evaluation at the engine\textquotesingle s ground seam, never by derivation.}
\declhead{theory Qtsl}{QLinGrowth}{\badge{BadgeGreen}{quoted from the theory source at render}}
\begin{leansrc}
QLinGrowth    : Set;
\end{leansrc}
\declhead{theory Qtsl}{QLinGrowth/*}{\badge{BadgeGreen}{quoted from the theory source at render}}
\begin{leansrc}
QLinGrowth/judgment : QLinGrowth -> QJudgment;
QLinGrowth/f : QLinGrowth -> PolyVec;
QLinGrowth/a : QLinGrowth -> Q;
QLinGrowth/b : QLinGrowth -> Q;
QLinGrowth/scope : QLinGrowth -> Box;
\end{leansrc}
\declhead{theory Qtsl}{QLinGrowthDischarge}{\badge{BadgeGreen}{quoted from the theory source at render}}
\begin{leansrc}
QLinGrowthDischarge : Set;
\end{leansrc}
\declhead{theory Qtsl}{QLinGrowthDischarge/*}{\badge{BadgeGreen}{quoted from the theory source at render}}
\begin{leansrc}
QLinGrowthDischarge/proof : QLinGrowthDischarge -> QProof;
QLinGrowthDischarge/envelopes : QLinGrowthDischarge -> QProofList;
\end{leansrc}
\ruleblock{Lean statement}
\badge{BadgeAmber}{not yet mechanized} this rule cites no Lean statements: nothing on this page is hash-bound for it, and its dashboard row shows the same absence across the statement, receipt, and correspondence columns.
\ruleblock{HOL4 statement}
\absentbox{NO HOL STATEMENT YET}{No HOL development of the extension theory exists.
}
\ruleblock{HOL verifier receipts}
\absentbox{NONE}{no HOL verification receipt exists for this rule.}
\ruleblock{Proof}

\sourcenote{Genre: core rule (tier 1), the discharge species. This section states the rule's soundness
obligation exactly, per the catalogue's completeness-over-polish
charter (a rule that can only be named-and-gapped is still a row);
the discharge is open debt on the mechanization column, and the
schema quoted in the entry is the baked ground truth.}

\textbf{Soundness obligation.}\enspace The per-coordinate affine-envelope sign sequents on the scope box
jointly certify membership of the field in the linear-growth envelope
with constants $(a, b)$; the kind carries the constants so that
downstream consumers compose them arithmetically. The obligation is
that the envelope rows' sequent semantics entail the envelope
inequality coordinatewise on the scope box --- a per-row sign-sequent
instantiation.

\begin{bremark}[status]
No Lean statement binds this node. The displayed norm form is
schematic; the envelope rows' own sequent shape (exact-rational,
checker-replayed) is the authority.
\end{bremark}


\rulesection{84}{QMarginalCompat}{core rule}
\label{blcatrule-QMarginalCompat}
\ruleblock{The rule (Steele metanotation)}
\sourcenote{as realized in the mechanized development; source row: \texttt{\texttt{QTSL-rules.toml}} --- \emph{QMarginalCompat}. The metanotation statement is typeset by hand for notational quality; this rule cites no Lean statements yet (the absence is stated loudly below). The rule \emph{roster} is machine-checked (every mechanized rule tag round-trips to exactly one catalogue row); the statement \emph{text} is editorial.}
\begin{mathpar}
\inferrule*{
  \overline{\mathrm{dom}_k}\ \ \text{(per-contributor dominance rows)} \\
  \mathsf{qtsl\_stochwire\_shape}(\mathit{cl},\ \mathit{sp})
}{
  \vDash \mathsf{QMarginalCompat}(\mathit{sp},\ \mathit{cl})
}
\end{mathpar}

\par\medskip
\rulename{QMarginalCompat} packages marginal compatibility for
stochastic wires: a belief names convex-analytic content
(the indicator of a named cell set, or a registered function class),
a stochastic wire carries one with its convexity,
lower-semicontinuity, and properness row kits, and this node collects
the per-contributor dominance rows over a shared sample space. The
contributor list\textquotesingle s shape agreement against the space
is an admission demand, discharged by the shape walk quoted below.

\cataloguenote{Entered at the roster census (baked, never catalogued); no Lean
statement binds it yet; display is schema-faithful schematic.
}
\ruleblock{Theory rows (geolog)}
\sourcenote{The rule\textquotesingle s certificate content as it is declared in the geolog theory the checker loads --- payload schema, side-condition demands, and axiom rows, quoted from the theory source at every build of this catalogue (a stale quotation cannot exist). Axiom modes: \texttt{chase} = the engine derives; \texttt{report} = a violation is a loud judgment; \texttt{refuse} = a violation gates admission. Side-condition relations are discharged by bounded evaluation at the engine\textquotesingle s ground seam, never by derivation.}
\declhead{theory Qtsl}{QMarginalCompat}{\badge{BadgeGreen}{quoted from the theory source at render}}
\begin{leansrc}
QMarginalCompat : Set;
\end{leansrc}
\declhead{theory Qtsl}{QMarginalCompat/*}{\badge{BadgeGreen}{quoted from the theory source at render}}
\begin{leansrc}
QMarginalCompat/proof : QMarginalCompat -> QProof;
QMarginalCompat/space : QMarginalCompat -> Shape;
QMarginalCompat/contributors : QMarginalCompat -> QStochContribList;
QMarginalCompat/dominance_rows : QMarginalCompat -> QProofList;
\end{leansrc}
\declhead{theory Qtsl}{qtsl\_stochwire\_shape}{\badge{BadgeGreen}{quoted from the theory source at render}}
\begin{leansrc}
qtsl_stochwire_shape : {contribs: QStochContribList, space: Shape}
  -> Prop;   // side-condition
\end{leansrc}
\declhead{theory Qtsl}{ax/wf/stochwire\_shape}{\badge{BadgeGreen}{quoted from the theory source at render}}
\begin{leansrc}
report ax/wf/stochwire_shape : forall m : QMarginalCompat,
    sp : Shape, cl : QStochContribList.
  m QMarginalCompat/space = sp,
  m QMarginalCompat/contributors = cl
  |- {contribs: cl, space: sp} qtsl_stochwire_shape;
\end{leansrc}
\ruleblock{Lean statement}
\badge{BadgeAmber}{not yet mechanized} this rule cites no Lean statements: nothing on this page is hash-bound for it, and its dashboard row shows the same absence across the statement, receipt, and correspondence columns.
\ruleblock{HOL4 statement}
\absentbox{NO HOL STATEMENT YET}{No HOL development of the extension theory exists.
}
\ruleblock{HOL verifier receipts}
\absentbox{NONE}{no HOL verification receipt exists for this rule.}
\ruleblock{Proof}

\sourcenote{Genre: core rule (tier 1), stochastic wires. This section states the rule's soundness
obligation exactly, per the catalogue's completeness-over-polish
charter (a rule that can only be named-and-gapped is still a row);
the discharge is open debt on the mechanization column, and the
schema quoted in the entry is the baked ground truth.}

\textbf{Soundness obligation.}\enspace The per-contributor dominance rows, over a shared sample space whose
shape agreement with the contributor list is an admission demand,
jointly certify marginal compatibility of the contributors'
beliefs --- the packaging judgment consumed by stochastic wire
composition. The obligation is that dominance per contributor entails
compatibility of the marginals on the shared space.

\begin{bremark}[status]
No Lean statement binds this node; the convex-analytic kits
(convexity, lower semicontinuity, properness) ride the stochastic
wire's own rows and are consumed here by citation.
\end{bremark}


\rulesection{85}{QAggReduce}{core rule}
\label{blcatrule-QAggReduce}
\ruleblock{The rule (Steele metanotation)}
\sourcenote{as realized in the mechanized development; source row: \texttt{\texttt{QTSL-rules.toml}} --- \emph{QAggReduce}. The metanotation statement is typeset by hand for notational quality; its mechanized ground truth is the Lean statement below, extracted from the sources at every build of this book and hash-bound. The rule \emph{roster} is machine-checked (every mechanized rule tag round-trips to exactly one catalogue row); the statement \emph{text} is editorial, checked by eye against the extracted Lean beside it.}
\begin{mathpar}
\inferrule*{
  (\vec{s}, \vec{k}) \in K \vdash G(\vec{s}, \vec{k}) \le 0 \ \text{(one \rulename{BernCone}-family sequent over the aggregate coordinates)} \\
  s^P_j := \textstyle\sum_{t:P} \phi_j(c_t), \qquad k_P := \#P \ \text{(the aggregate coordinates, defined)} \\
  \bigl\{\vdash l_j\, k_P \le s^P_j \le u_j\, k_P\bigr\}_{j, P} \ \text{(coupling rows; } [l_j, u_j] \ \text{a Bernstein enclosure of } \phi_j \ \text{on } B_P\text{)} \\
  \bigl\{\vdash 0 \le k_P \le M_P\bigr\}_{P} \ \text{(count rows)} \\
  \bigl\{\vdash y^{\top} \vec{k} = y^{\top} \vec{k}(x_0)\bigr\}_{y \in Y} \ \text{(P-invariant conservation rows: each cited witness } y \ge 0,\ y^{\top} C = 0 \ \text{verified against the uncoloured skeleton\textquotesingle s incidence matrix } C\text{; } x_0 \ \text{the initial marking)} \\
  \bigl\{\vdash (s^P_\alpha)^2 \le k_P\, s^P_{2\alpha}\bigr\}_{\alpha, P} \ \text{(moment rows, optional tightening)} \\
  \forall x \in \mathrm{scope},\ \forall t{:}P \ \text{in } x:\ c_t \in B_P \ \text{(colour-box scope)}
}{
  x \in \mathrm{scope} \vdash G\bigl(\vec{s}(x), \vec{k}(x)\bigr) \le 0
}
\end{mathpar}

\ruleflavor{the defined aggregate coordinates replace the certificate\textquotesingle s feature dictionary in the display: the dictionary is certificate data (a shared reference) whose entire displayed content is the two defining equations, with $\phi_j$ cited by the coupling rows. The enclosure box discharge is a side condition. The checker never sees a marking: the multiset quotient is compiled away producer-side, aggregate dimension independent of the marking bound $M$; every row family is valid at every reachable in-scope marking under its stated premises (the moment and count-lower rows unconditionally; the count-upper row under the marking bound $M$; conservation rows by their verified witness vectors, the conserved value pinned at the initial marking (the rows ride the count-row proof list --- no dedicated conservation field exists on the baked node, a named schema divergence); coupling rows under their per-token bounds), so adding in-scope rows only shrinks $K$ and never cuts a reachable in-scope aggregate point.}
\par\medskip
\rulename{QAggReduce} is the aggregation rule for nets whose state is a
\emph{marking} --- a multiset of coloured tokens distributed over places,
with variable dimension and a mandatory exchangeability quotient
(tokens within a place are interchangeable, so any legitimate
observable is a symmetric function of each place\textquotesingle s
colour multiset). It replaces an obligation quantified over all
in-scope markings by a single polynomial obligation over finitely many
\emph{aggregate} coordinates: token counts $k_P$ and pooled features
$s^P_j = \sum_{t:P} \phi_j(c_t)$ for a finite feature dictionary. A
producer reaches for it because the alternative --- flattening marking
shapes into modes --- explodes combinatorially in the token bound and
breaks the symmetry; here obligation count and dimension are
independent of the marking bound, which enters only as a constant in
count rows.

What the checker recomputes: one hypothesis-carrying polynomial sequent
$G \le 0$ over the aggregate box, whose hypothesis rows assemble the
containment certificate for a polytope $K$ covering every in-scope
marking\textquotesingle s aggregate image --- coupling rows
$l_j k_P \le s^P_j \le u_j k_P$ from enclosures of each feature on its
place\textquotesingle s colour box, count rows $0 \le k_P \le M_P$ plus
linear conservation laws $y^{\top} \vec{k} = y^{\top} \vec{k}(x_0)$
(place invariants \cite{murata1989}, the conserved value pinned at the
initial marking; each cited witness $y \ge 0$, $y^{\top} C = 0$
checker-verified against the uncoloured skeleton\textquotesingle s
incidence matrix), and
optional moment rows $(s^P_\alpha)^2 \le k_P\, s^P_{2\alpha}$
(Cauchy--Schwarz over the bag). Soundness of the reduction is
evaluation: if the aggregate image lands in $K$ and $G \le 0$ on $K$,
the marking-quantified obligation holds pointwise. The whole content is
the containment certificate, with each row family proven valid at its
stated scope: nonnegativity and moment rows for \emph{every} bag
outright; coupling rows for every bag under their colour-scope premise
(the box discharge); count bounds and conservation rows as certified
reachability facts. Under those side conditions, adding a row shrinks
$K$ without ever cutting a reachable aggregate point.

What would break: a row NOT valid for every reachable marking makes $K$
fail containment and the conclusion simply false --- the central
instance being the coupling rows\textquotesingle{} scope premise, since
an enclosure of $\phi_j$ on the colour box says nothing about tokens
outside it; the box must be \emph{discharged} (exit-guard transitions,
an invariance proof, or stopped-process semantics), never assumed. The
converse failure is honest: not every polytope point is a realizable
marking, so $K$ over-approximates --- a completeness cost, refusal-shaped
and never soundness-shaped, with the row menu as its dial.

\cataloguenote{The operative aggregate is the measurable-query observable
$\langle \mathrm{marking}, f \rangle$ over bags of colours; counts and
pools are the $f \equiv 1$ and $f = \mathrm{id}$ degenerate queries,
and a query is a semiring contraction of the count tensor against $f$
along the colour axis. Semantically the pools are commutative-monoid
folds along the tokens-to-places classifying map, with the multiset as
the universal aggregate \cite{spivak2021agg}; the dictionary\textquotesingle s
expressive ceiling per shape is the classical generation of
multisymmetric invariants by polarized power sums \cite{weyl1939}. The
coupling rows here are deterministic pool--count arithmetic within one
marking --- \emph{not} probabilistic couplings; the identically-named
rows of \rulename{QCoupleDrift} are joint-kernel marginal data over a
product of chains. Same word, disjoint mathematics, neither cites the
other.
}
\ruleblock{Lean statement}
\declhead{BernsteinLogicAux.QtslCompose.AggReduce}{aggReduce}{\badge{BadgeGreen}{statement bound}}
\begin{leansrc}
/-- **QAggReduce, semantic form** ("proof: evaluation"): if `G ≤ 0` on `K`
and every in-scope point's aggregate vector lands in `K`, then the
marking-quantified obligation holds. `agg` is the compiled aggregate map
`x ↦ (s⃗(x), k⃗(x))` — its pointwise containment (`hin`) is what the
coupling/count/moment rows + the P-invariant conservation laws certify. -/
theorem aggReduce {X A : Type*} (agg : X → A) (scope : Set X) (K : Set A)
    (G : A → ℝ) (hG : ∀ a ∈ K, G a ≤ 0) (hin : ∀ x ∈ scope, agg x ∈ K) :
    ∀ x ∈ scope, G (agg x) ≤ 0
\end{leansrc}
\declhead{BernsteinLogicAux.QtslCompose.AggReduce}{coupling\_row}{\badge{BadgeGreen}{statement bound}}
\begin{leansrc}
/-- **The coupling row is valid for every marking** (NOT a probabilistic
coupling — see the module docstring's fence): if every token's `f`-value
lies in `[l, u]`, the pool–count coupling `l · k ≤ s ≤ u · k` holds. Adding
these rows to the polytope `K` never cuts a reachable aggregate point. -/
theorem coupling_row {m : Multiset C} {f : C → ℝ} {l u : ℝ}
    (hl : ∀ c ∈ m, l ≤ f c) (hu : ∀ c ∈ m, f c ≤ u) :
    l * m.card ≤ query m f ∧ query m f ≤ u * m.card
\end{leansrc}
\declhead{BernsteinLogicAux.QtslCompose.AggReduce}{count\_row\_le}{\badge{BadgeGreen}{statement bound}}
\begin{leansrc}
/-- The count row's upper half, from any marking bound (the `M_P` bound is a
reachability/P-invariant fact, consumed as a hypothesis — the box-discharge
side condition [D10], not proven here). -/
theorem count_row_le {m : Multiset C} {M : ℕ} (h : m.card ≤ M) :
    (m.card : ℝ) ≤ M
\end{leansrc}
\declhead{BernsteinLogicAux.QtslCompose.AggReduce}{count\_row\_nonneg}{\badge{BadgeGreen}{statement bound}}
\begin{leansrc}
/-- The count row's lower half: counts are nonnegative. -/
theorem count_row_nonneg (m : Multiset C) : (0 : ℝ) ≤ m.card
\end{leansrc}
\declhead{BernsteinLogicAux.QtslCompose.AggReduce}{moment\_row}{\badge{BadgeGreen}{statement bound}}
\begin{leansrc}
/-- **The moment row is valid for every marking**: `⟨m, f⟩² ≤ k · ⟨m, f²⟩` —
Cauchy–Schwarz over the bag (the rule display's `(s_α)² ≤ k · s_{2α}`,
with `s_α = ⟨m, x^α⟩` instantiated at `f = x^α`). -/
theorem moment_row (m : Multiset C) (f : C → ℝ) :
    (query m f) ^ 2 ≤ m.card * query m (fun c => f c ^ 2)
\end{leansrc}
\declhead{BernsteinLogicAux.QtslCompose.AggReduce}{query\_const\_one}{\badge{BadgeGreen}{statement bound}}
\begin{leansrc}
/-- Counts are the `f ≡ 1` degenerate query (the 14:20 correction's
normal form: `k_P = ⟨m, 1⟩`). -/
theorem query_const_one (m : Multiset C) :
    query m (fun _ => (1 : ℝ)) = m.card
\end{leansrc}
\declhead{BernsteinLogicAux.QtslCompose.AggReduce}{query\_count}{\badge{BadgeGreen}{statement bound}}
\begin{leansrc}
/-- **The contraction-row form** (the γ-trace): the query observable is the
count tensor contracted against `f` along the color axis. -/
theorem query_count [DecidableEq C] (m : Multiset C) (f : C → ℝ) :
    query m f = ∑ c ∈ m.toFinset, (m.count c : ℝ) * f c
\end{leansrc}
\ruleblock{Correspondence (Steele display $\leftrightarrow$ Lean)}
\badge{BadgeGreen}{reviewer-certified against the displayed texts}
\sourcenote{The dictionary below is a reviewer-certified statement of how the mechanized statements above interpret the displayed rule: display metavariable $\mapsto$ Lean binder, each display hypothesis $\mapsto$ the premise or side condition carrying it, the display conclusion $\mapsto$ the goal --- with any deliberate difference in width between display and mechanization named outright. Certification is recorded as a hash triple over (the displayed cell, the extracted statement texts, this dictionary\textquotesingle s own text); if any of the three changes, this block turns amber until a reviewer re-certifies. The correspondence claim itself remains editorial --- reviewer-grade, not machine-checked; its machine-checked successor types each soundness statement against the rule as a structured object, and supersedes this block per rule as it lands.}
The displayed conclusion $\forall x$ in scope: $G(\vec{s}(x),
\vec{k}(x)) \le 0$ $\mapsto$ \texttt{aggReduce}\textquotesingle s goal
\texttt{∀ x ∈ scope, G (agg x) ≤ 0}, with the displayed premise
$\vdash G \le 0$ on $K$ $\mapsto$ \texttt{hG} and the entire row
apparatus $\mapsto$ \texttt{hin : ∀ x ∈ scope, agg x ∈ K} --- the
containment certificate, which the coupling/count/moment/conservation
rows exist to assemble. The binder row is the biggest width fact on
this page: \texttt{aggReduce} is stated for \emph{any} types \texttt{X
A}, any \texttt{agg : X → A}, any \texttt{K : Set A}, any \texttt{G :
A → ℝ} --- no polynomiality of $G$, no polytopality of $K$, no
aggregate box --- so the certified semantic core is strictly wider
than the displayed polynomial instance, safe direction (proof by
evaluation needs none of the structure the checker needs).
\begin{itemize}
\item coupling rows $l_j k_P \le s^P_j \le u_j k_P$ $\mapsto$
  \texttt{coupling\_row}: hypotheses \texttt{∀ c ∈ m, l ≤ f c} and
  \texttt{f c ≤ u} --- the enclosure quantified over tokens
  \emph{in the bag}, which is precisely where the box-discharge side
  condition lands in the certified statement (an enclosure on
  the colour box reaches the hypothesis only for markings confined to
  the box; the flavor text carries this as the box-discharge side
  condition, the statement as the membership guard).
\item count rows $0 \le k_P \le M_P$ $\mapsto$
  \texttt{count\_row\_nonneg} (unconditional) +
  \texttt{count\_row\_le}, whose bound \texttt{m.card ≤ M} is a
  \emph{hypothesis} --- the $M_P$ bound is a reachability/P-invariant
  fact, consumed, never proven. The displayed conservation rows
  $y^{\top} \vec{k} = y^{\top} \vec{k}(x_0)$ (P-invariants, the
  conserved value pinned at the initial marking; witnesses $y \ge 0$,
  $y^{\top} C = 0$ checker-verified) have no bound Lean counterpart at
  all: they are
  reachability facts of the net, outside this module\textquotesingle s
  every-bag warrants. On the baked node they ride the count-row proof
  list --- no dedicated conservation field exists, a named schema
  divergence.
\item moment rows $(s^P_\alpha)^2 \le k_P\, s^P_{2\alpha}$ $\mapsto$
  \texttt{moment\_row}, unconditional for every bag (Cauchy--Schwarz
  by induction), instantiated at \texttt{f = x\^{}α}.
\item the defined aggregate coordinates $s^P_j := \sum_{t:P}
  \phi_j(c_t)$ and $k_P := \#P$ $\mapsto$ per-query functions
  \texttt{f : C → ℝ}, arbitrary --- polynomiality of the $\phi_j$ is
  producer-side; \texttt{query\_count} certifies the
  contraction-row form $\langle m, f\rangle = \sum_c
  \mathrm{count}(c) \cdot f(c)$ (the $\gamma$-trace) under
  \texttt{[DecidableEq C]}, and \texttt{query\_const\_one} pins counts
  as the $f \equiv 1$ degenerate query.
\end{itemize}
Width note: every bound row warrant is per-bag arithmetic --- the
multiset quotient is already taken (\texttt{Multiset C} IS the
exchangeability-quotiented object), so the display\textquotesingle s
``the checker never sees a marking'' has as its Lean face that no
statement here mentions places, nets, or reachability; scope
containment \texttt{hin} is where all of that enters, as one
hypothesis.
The conditional mark\textquotesingle s content, at its exact extent:
the count-bound and conservation rows are reachability facts of the
net --- consumed as hypotheses, never proven here --- and a coupling
row\textquotesingle s colour-box confinement must be discharged by
the producer (exit-guard transitions, an invariance certificate, or
stopped-process semantics), never assumed. Every other row family is
warranted for every bag outright.

\ruleblock{Lean verifier receipts}
\sourcenote{module \texttt{BernsteinLogicAux.QtslCompose.AggReduce}; bound build \texttt{be9dbd259553\ldots} (bound from receipt bd4d8aca0400b822, 2026-07-21T04:51:40Z). A receipt records that named checking layers accepted a build: (module, compiled-artifact SHA-256, environment SHA-256, checker layers), minted when the checking gate ran; the store receipts attesting the bound build are listed with the layers each run credited. Nothing here asks the reader to trust the store: every layer is an independently runnable program over the quoted sources --- rebuild the module and re-run the checkers to reproduce the verdict. Receipts attesting superseded builds credit no layer and are omitted from this rendering (the count of omissions, if any, is noted below).}
\begin{receipttable}
\receiptrow{5f6de8c6a1ff5236}{2026-07-21T01:51:28Z}{build+decorrelated}{\badge{BadgeGreen}{current}}
\receiptrow{bd4d8aca0400b822}{2026-07-21T04:51:40Z}{build+decorrelated}{\badge{BadgeGreen}{current}}
\receiptrow{c37acfaf586210ee}{2026-07-21T05:39:09Z}{build+replay+hardened+decorrelated}{\badge{BadgeGreen}{current}}
\receiptrow{24c5f02176c0ee6a}{2026-07-22T18:27:15Z}{build+decorrelated}{\badge{BadgeGreen}{current}}
\receiptrow{1eaf40d323232bf8}{2026-07-31T22:00:14Z}{build+decorrelated}{\badge{BadgeGreen}{current}}
\receiptrow{47734200303b43f0}{2026-07-31T22:29:07Z}{build+decorrelated}{\badge{BadgeGreen}{current}}
\receiptrow{712878ff235d6c76}{2026-08-02T23:31:19Z}{build+decorrelated}{\badge{BadgeGreen}{current}}
\receiptrow{36b853e33e482b54}{2026-08-03T16:27:22Z}{build+decorrelated}{\badge{BadgeGreen}{current}}
\receiptrow{cb8f537f8267cf3e}{2026-08-03T22:46:06Z}{build+decorrelated}{\badge{BadgeGreen}{current}}
\end{receipttable}
\sourcenote{6 further receipt(s) in the store attest superseded builds of this module and are omitted from this rendering.}
\layerline{\layeritem{build}{lake build (elaboration)}{\badge{BadgeGreen}{current}}\layeritem{replay}{leanchecker (official kernel replay)}{\badge{BadgeGreen}{current}}\layeritem{hardened}{safe\_verify (hardened screen)}{\badge{BadgeGreen}{current}}\layeritem{decorrelated}{lean4export + nanoda (decorrelated kernel)}{\badge{BadgeGreen}{current}}}
\sourcenote{\badge{BadgeGray}{comparator: none} no fixture in the three-implementation comparator harness (compiled / kernel-interpreted / reference-interpreted) exercises this rule.}
\ruleblock{HOL4 statement}
\absentbox{NO HOL STATEMENT YET}{The HOL4 development checks the base calculus\textquotesingle s wire
format (theory \texttt{bern\_catalogue}). \rulename{QAggReduce} is an
extension-edge semantic rule whose node is baked in the extension
theory; what reaches the base checker after the
compile is ordinary hypothesis-carrying \rulename{BernCone} sequents,
which the HOL side already covers as sequents without knowing their
aggregation provenance. A HOL statement \emph{about the rule} becomes
writable when a checker consuming the extension theory\textquotesingle s
rows is ported.
}
\ruleblock{HOL verifier receipts}
\absentbox{NONE}{no HOL verification receipt exists for this rule.}
\ruleblock{Proof}

\sourcenote{Genre: core rule (tier 1). Scope, stated loudly: the
semantic core (the aggregate-space reduction) is proven, and the
row-warrant families are proven at their stated scopes ---
nonnegativity and moment rows for every bag, coupling rows under their
colour-scope premise, while count upper bounds and place-invariant
rows enter as certified in-scope reachability facts, not bagwise
identities; the marking-space compile itself
--- producing the aggregate map and its rows from a net --- is checker
and compiler recomputation, outside these statements. Every hypothesis
is used and remarked on.}

\textbf{Setting and claim.}\enspace The state of a coloured net is a
\emph{marking}: for each place $P$, a finite multiset of colours drawn
from a compact rational box $B_P \subset \R^{n_P}$ --- so the state
space is a disjoint sum, over count vectors, of symmetric powers, with
variable dimension and a mandatory exchangeability quotient (a
``function of the state'' that reads token identities is not a function
on the quotient at all). The certificate language works in
\emph{aggregate coordinates}: counts $k_P = \#P$ and pooled features
$s^P_j = \sum_{t : P} \phi_j(c_t)$ for a finite dictionary
$\Phi = \{\phi_j\}$ of polynomials. An obligation arrives as
``for every in-scope marking $x$:
$G\bigl(\vec{s}(x), \vec{k}(x)\bigr) \leq 0$'' with $G$ polynomial. The
rule replaces marking quantification by one polynomial sequent over a
polytope $K$ in aggregate space, whose hypothesis rows are the
containment certificate. The claim (\texttt{aggReduce} + the row
warrants): the reduction is sound, and every row family below --- each
under its stated scope premise --- may be added to $K$'s description
without ever cutting a reachable aggregate point.

\begin{blemma}[queries, counts, and the contraction form]
\label{lem:qar-query}
For a bag $m$ of colours and $f : \text{colour} \to \R$, write
$\langle m, f \rangle = \sum_{c \in m} f(c)$ (with multiplicity). Then
$\langle m, \mathbf{1} \rangle = \#m$, and
$\langle m, f \rangle = \sum_{c \in \mathrm{supp}(m)}
\mathrm{count}_m(c) \cdot f(c)$ --- the count tensor contracted against
$f$ along the colour axis.
\end{blemma}
\begin{proof}
Both are fibre decompositions of a finite sum: group the terms of
$\sum_{c \in m} f(c)$ by value of $c$. Counts and pools are thus the
$f \equiv 1$ and coordinate-projection instances of one observable
class, and every axis of the contraction is either summed out (colour)
or an output (the query index) --- there is no un-read axis, which is
why, for the chosen finite dictionary, aggregate dimension is the
dictionary's index set: set by the dictionary, never by the marking.
\end{proof}

\begin{blemma}[coupling rows are valid for every bag]\label{lem:qar-couple}
If $l \leq f(c) \leq u$ for every colour $c$ occurring in $m$, then
\[
l \cdot \#m \;\leq\; \langle m, f \rangle \;\leq\; u \cdot \#m .
\]
\end{blemma}
\begin{proof}
A sum of $\#m$ terms, each in $[l, u]$ (induction on the bag). In use,
$[l, u]$ is a certified enclosure of $\phi_j$ on the colour box $B_P$
--- polynomial-enclosure machinery the calculus already owns --- so the
row couples the pool wire to the count wire. Two remarks. The bounds
hypothesis quantifies over colours \emph{in the bag}: the row is
therefore only as good as the guarantee that tokens stay in $B_P$ ---
the scope premise of \autoref{rem:qar-box}. And the degenerate case
is load-bearing: at $k_P = 0$ the two inequalities pin $s^P_j = 0$ ---
emptiness is polytope geometry, not a case split.
\end{proof}

\begin{blemma}[count rows]\label{lem:qar-count}
$0 \leq \#m$ always; and $\#m \leq M$ for any certified marking bound
$M$. Linear conservation laws of the uncoloured skeleton (place
invariants \cite{murata1989}) are further valid count rows.
\end{blemma}
\begin{proof}
Cardinalities are nonnegative integers; the bound and the invariants
are reachability facts of the net, consumed as hypotheses (they are
certified upstream as reachability/place-invariant facts, never
re-proven per obligation). Integer-valuedness of $k_P$ is what
the counting closure discipline (\rulename{QCountAtomRounding})
exploits at complemented count atoms.
\end{proof}

\begin{blemma}[moment rows are valid for every bag]\label{lem:qar-moment}
For every bag $m$ and every $f$:
$\langle m, f \rangle^2 \leq \#m \cdot \langle m, f^2 \rangle$.
\end{blemma}
\begin{proof}
Cauchy--Schwarz over the counting measure on the bag:
$\bigl(\sum_{c \in m} 1 \cdot f(c)\bigr)^2 \leq
\bigl(\sum_{c \in m} 1^2\bigr)\bigl(\sum_{c \in m} f(c)^2\bigr)$.
Instantiated at monomial features this is the family
$(s^P_\alpha)^2 \leq k_P\, s^P_{2\alpha}$ --- optional tightenings:
each one added shrinks the relaxation $K$, and by this lemma never
below the reachable image.
\end{proof}

\begin{bthm}[soundness of \rulename{QAggReduce}]\label{thm:qar-sound}
Let $\mathrm{agg}$ be the aggregate map
$x \mapsto (\vec{s}(x), \vec{k}(x))$, let
$K$ satisfy $\mathrm{agg}(x) \in K$ for every in-scope $x$, and suppose
the sequent $G \leq 0$ holds on $K$. Then
$G(\mathrm{agg}(x)) \leq 0$ for every in-scope $x$.
\end{bthm}
\begin{proof}
Evaluation: $\mathrm{agg}(x) \in K$ and $G \leq 0$ on $K$. That the
proof is one line \emph{is} the honest size of the semantic core ---
the entire content of a \rulename{QAggReduce} node is the containment
certificate for $K$, which Lemmas
\ref{lem:qar-couple}--\ref{lem:qar-moment} assemble from rows each
valid at every reachable marking: a conjunction of per-marking-valid
constraints contains the whole reachable aggregate image, so the single
polynomial sequent over $K$ (discharged by the hypothesis-carrying
cone machinery of the base calculus) covers every in-scope marking at
once. Obligation count and dimension are independent of the marking
bound, which enters only as a constant in Lemma~\ref{lem:qar-count}'s
rows --- the reason this rule exists instead of a mode explosion over
marking shapes.
\end{proof}

\begin{bremark}[the box-discharge scope premise]\label{rem:qar-box}
Lemma~\ref{lem:qar-couple}'s hypothesis --- every occurring colour in
$B_P$ --- is a claim about the \emph{process}, and with nondegenerate
diffusion a process exits every compact box with positive probability.
An enclosure row quantified over states the process leaves is
unsoundness, not incompleteness; so the box must be \emph{discharged},
by one of three named routes: model box exit as a forced transition
(the boundary becomes part of the model, so no in-scope state ever
lies outside the box and the exit's successor region supplies its own
rows), prove invariance of the box for every reachable state
(vanishing diffusion at the boundary plus inward drift --- polynomial,
box-supported conditions landing in the existing row machinery), or
certify the process stopped at exit --- the stopped state kept inside
the box, with the exit event carried in the specification's
probability budget. Which route is a per-model
declaration; that \emph{some} route is declared is a side condition of
the rule. Dropping it admits the counterexample: a generator row verified
on $B_P$ says nothing at the first token that leaves, and the
concluded obligation can be false on reachable states.
\end{bremark}

\begin{bremark}[what the relaxation costs, and what it never costs]
\label{rem:qar-gap}
$K$ over-approximates the reachable aggregate image --- not every
polytope point is a realizable marking. The cost is
\emph{completeness}: a true obligation may fail on a spurious point of
$K$ and the sequent then refuses; it is never soundness, by
Theorem~\ref{thm:qar-sound}. The row menu is the dial: integer
rounding at count atoms, conservation laws, and moment rows each cut
spurious regions. The drain example, worked to its numbers in
\autoref{rem:qar-drain}, is the sharpest small witness: its
interior row is \emph{infeasible} over the real-relaxed
off-target scope $k > 0$ (it evaluates to $+\tfrac14$ at
$k = \tfrac12$, $s_2 = 0$, a point no reachable marking occupies) and
discharges, with an exact linear-multiplier witness, once
\rulename{QCountAtomRounding} tightens the scope to $k \geq 1$ ---
the tightening menu earning its keep on the first toy. Conversely a
row that fails at some reachable in-scope marking (a naively-cleared
mean at $k = 0$, say) would cut reachable points and break
Theorem~\ref{thm:qar-sound}'s containment hypothesis --- validity at
every reachable in-scope marking, under the row family's stated
premises, is the admission criterion for the menu: the nonnegativity
and moment rows meet it for every bag outright, the coupling rows
meet it once their colour-scope premise is discharged, and the count
bounds and invariants meet it as certified reachability facts
consumed as hypotheses.
\end{bremark}

\begin{bremark}[the drain example, worked to its numbers]
\label{rem:qar-drain}
One place, colour $x \in [-2, 2]$, one completion transition consuming
a single token at rate $\lambda(x) = 1 + x^2$, no source; the deadline
certificate is $V = 1 + k$ with $\theta = \tfrac12$ (counts alone
suffice for $V$; the pooled feature $s_2 = \sum_t x_t^2$ enters the
interior row through the rate, so its coupling row joins the scope).
The generator: the continuous part vanishes ($V$ reads no colours),
and each completion removes one token, so
$\mathcal{L}V = \sum_t \lambda(x_t)\bigl(V(k{-}1) - V(k)\bigr)
= -\sum_t (1 + x_t^2) = -(k + s_2)$. The interior row is
\[
\theta V + \mathcal{L}V
\;=\; \tfrac12(1 + k) - (k + s_2)
\;=\; -\tfrac12(k - 1) - s_2 .
\]
Over the real-relaxed off-target scope
$\{\, k > 0,\ 0 \le s_2 \le 4k \,\}$ --- the coupling row
(Lemma~\ref{lem:qar-couple}) at the enclosure $x^2 \in [0, 4]$ on the
colour box --- the row obligation fails: $(k, s_2) = (\tfrac12, 0)$ is
a scope point but no marking (counts are integers), and there the row
evaluates to $\tfrac12 \cdot \tfrac32 - \tfrac12 = +\tfrac14$. Once
\rulename{QCountAtomRounding} rounds the off-target scope
$\mathrm{cmpl}\{k \le 0\}$ to $k \ge 1$, the rewriting displayed above
\emph{is} the exact linear-multiplier witness: multiplier $\tfrac12$
on the scope row $k - 1 \ge 0$ and $1$ on the coupling floor
$s_2 \ge 0$ reproduce the row identically --- a polynomial identity,
no residue, no numerics --- and the certificate is tight at
$(k, s_2) = (1, 0)$: the single-cold-token state, the true bottleneck,
since the slowest completion rate is $\lambda(0) = 1$. Every number
cited in \autoref{rem:qar-gap} recomputes from this display; the
example's other rows (floor, atom, root) are not the story here and
are not staged.
\end{bremark}

\begin{bremark}[relation to the literature]
The aggregate coordinates are the classical multisymmetric power sums:
per fixed shape, polarized power sums generate the full invariant
algebra in characteristic $0$ \cite{weyl1939}, which locates the
dictionary's expressiveness ceiling (generation degree grows with the
shape --- a resolution bound, priced as completeness, never claimed
away). Degree-one certificates in counts alone recover classical
place-invariant analysis \cite{murata1989}. Semantically, pools are
commutative-monoid folds along the tokens-to-places classifying map,
with the multiset as the universal aggregate every other one factors
through \cite{spivak2021agg} --- exchangeability and
factoring-through-the-free-commutative-monoid are the same statement
from two sides. The identically-named ``coupling rows'' of the
relational tier (\rulename{QCoupleDrift}) are probability couplings
over a product of chains; the rows here are deterministic pool--count
arithmetic within one marking. Same word, disjoint mathematics,
deliberately uncited in both directions.
\end{bremark}

\begin{anchors}
Lemma \ref{lem:qar-query} & \texttt{query}; \texttt{query\_const\_one}; \texttt{query\_count}\\
Lemma \ref{lem:qar-couple} & \texttt{coupling\_row} (degenerate pin at $k = 0$: its two halves)\\
Lemma \ref{lem:qar-count} & \texttt{count\_row\_nonneg}; \texttt{count\_row\_le}\\
Lemma \ref{lem:qar-moment} & \texttt{moment\_row}\\
Theorem \ref{thm:qar-sound} & \texttt{aggReduce}\\
integrality at count atoms & \rulename{QCountAtomRounding} section (\texttt{countAtomRounding\_sound}/\texttt{\_eq})\\
\end{anchors}


\rulesection{86}{QCountAtomRounding}{alias rule}
\label{blcatrule-QCountAtomRounding}
\ruleblock{The rule (Steele metanotation)}
\sourcenote{as realized in the mechanized development; source row: \texttt{\texttt{QTSL-rules.toml}} --- \emph{QCountAtomRounding}. The metanotation statement is typeset by hand for notational quality; its mechanized ground truth is the Lean statement below, extracted from the sources at every build of this book and hash-bound. The rule \emph{roster} is machine-checked (every mechanized rule tag round-trips to exactly one catalogue row); the statement \emph{text} is editorial, checked by eye against the extracted Lean beside it.}
\begin{mathpar}
\inferrule*{
  k_P \ \text{a count observable, integer-valued at every reachable marking} \\
  j \in \mathbb{Z}
}{
  \mathrm{cmpl}\,\{k_P \le j\} \ \rightsquigarrow \ \{k_P \ge j + 1\} \quad \text{(hypothesis positions, reachable scope)}
}
\end{mathpar}

\ruleflavor{integrality is the compile\textquotesingle s guarantee. Exact over the reachable locus (the complement over reachable markings \emph{is} $\{k_P \ge j + 1\}$, not $\{k_P > j\}$), and unconditionally never wider ($\{k_P \ge j+1\} \subseteq \{k_P > j\}$) --- rounding to the next integer is free and strictly tighter than the real complement, the counting analogue of the strict-atom closure discipline.}
\par\medskip
\rulename{QCountAtomRounding} is the counting closure discipline: when a
count observable $k_P$ is complemented in a hypothesis position, the
real-arithmetic complement of $\{k_P \le j\}$ --- the open condition
$\{k_P > j\}$ --- is replaced by the closed integer-rounded atom
$\{k_P \ge j + 1\}$. Over reachable markings the two are \emph{equal}
(counts are integers there), and everywhere the rounded atom is
contained in the open one, so the rewrite is free and never wider
than the real complement --- strictly tighter wherever the slab is
realized: substituting it in hypothesis position never widens the
constraint set, and the gained slab $j < k_P < j + 1$
(unreachable, but visible to real-arithmetic relaxations) is often
exactly what makes a downstream combination certificate feasible. It is
the counting analogue of the strict-atom closure discipline for
continuous observables, with integrality replacing topology.

There is nothing for the checker to recompute beyond the syntactic
rewrite itself: producers apply it before emission, and what the wire
carries is the already-tightened atom inside an aggregation
rule\textquotesingle s hypothesis rows. The one genuine hypothesis is
integrality --- that $k_P$ takes integer values at every reachable
marking. That is a guarantee of the marking-space compile (counts are
cardinalities), not of the rewrite; in the mechanization it enters as a
named hypothesis, discharged definitionally when the net layer lands.
Without it the exactness identity fails (a fractional
\textquotedblleft count\textquotedblright{} at $k_P = j + \tfrac12$
satisfies the true complement but not the rounded atom), and only the
containment direction survives --- which is not enough in hypothesis
position: the rounded scope is \emph{smaller}, so a row verified only
on it silently skips exactly such fractional points, and a misapplied
rounding can silently accept rather than refuse (the proof
section\textquotesingle s integrality remark carries the witness). The
integrality warrant is load-bearing for soundness, not a convenience.

\cataloguenote{Load-bearing on the first worked aggregation example: a drain
net\textquotesingle s interior row is infeasible over the real-relaxed
off-target scope $k > 0$ (it fails at $k = \tfrac12$) and discharges
with the rounded scope $k \ge 1$, via an exact linear-multiplier
witness. The rewrite rides with \rulename{QAggReduce}\textquotesingle s
family: its target is an atom-denotation identity consumed by count
rows, not a separate rule node --- one atom in, one atom out.
}
\ruleblock{Lean statement}
\declhead{BernsteinLogicAux.Structural}{countAtomRounding\_sound}{\badge{BadgeGreen}{statement bound}}
\begin{leansrc}
/-- **P3.3 CountAtomRounding**, the soundness side (hypothesis position):
every reachable point of the true complement of `{k ≤ j}` satisfies the
tightened atom `{j + 1 ≤ k}`. This is the side the drain example's interior
row consumes — the aggregate polytope `K` must contain every in-scope
reachable point, and it does. -/
theorem countAtomRounding_sound (k : X → ℝ) (j : ℤ) (Reach : Set X)
    (hbridge : ∀ x ∈ Reach, ∃ n : ℤ, k x = (n : ℝ)) :
    ∀ x ∈ Reach, x ∈ {x | k x ≤ (j : ℝ)}ᶜ → (j : ℝ) + 1 ≤ k x
\end{leansrc}
\declhead{BernsteinLogicAux.Structural}{countAtomRounding\_tighter}{\badge{BadgeGreen}{statement bound}}
\begin{leansrc}
/-- **P3.3**, the free-tightening direction: the rounded atom is contained
in the real complement unconditionally — substituting it never widens the
constraint set, so the rewrite is strictly tighter (by the cut
`(j, j+1)`), which is what buys the drain example its Farkas
feasibility. Stated at any real threshold `j`: the containment needs no
integrality — only exactness (`countAtomRounding_eq`) and soundness
(`countAtomRounding_sound`) ride the integer-threshold bridge. -/
theorem countAtomRounding_tighter (k : X → ℝ) (j : ℝ) :
    {x | j + 1 ≤ k x} ⊆ {x | j < k x}
\end{leansrc}
\declhead{BernsteinLogicAux.Structural}{countAtomRounding\_eq}{\badge{BadgeGreen}{statement bound}}
\begin{leansrc}
/-- **P3.3**, exactness on the reachable locus: over integer-valued points
the rewrite loses nothing — the catalogue's "the complement ... over
reachable markings IS `{k_P ≥ j + 1}`", verbatim as a set identity. -/
theorem countAtomRounding_eq (k : X → ℝ) (j : ℤ) (Reach : Set X)
    (hbridge : ∀ x ∈ Reach, ∃ n : ℤ, k x = (n : ℝ)) :
    Reach ∩ {x | k x ≤ (j : ℝ)}ᶜ = Reach ∩ {x | (j : ℝ) + 1 ≤ k x}
\end{leansrc}
\ruleblock{Correspondence (Steele display $\leftrightarrow$ Lean)}
\badge{BadgeGreen}{reviewer-certified against the displayed texts}
\sourcenote{The dictionary below is a reviewer-certified statement of how the mechanized statements above interpret the displayed rule: display metavariable $\mapsto$ Lean binder, each display hypothesis $\mapsto$ the premise or side condition carrying it, the display conclusion $\mapsto$ the goal --- with any deliberate difference in width between display and mechanization named outright. Certification is recorded as a hash triple over (the displayed cell, the extracted statement texts, this dictionary\textquotesingle s own text); if any of the three changes, this block turns amber until a reviewer re-certifies. The correspondence claim itself remains editorial --- reviewer-grade, not machine-checked; its machine-checked successor types each soundness statement against the rule as a structured object, and supersedes this block per rule as it lands.}
The entry\textquotesingle s two claims (exactness and never-wider,
now the flavor text) and the displayed integrality premise
map onto the three bound statements. The compile\textquotesingle s
integrality guarantee --- the rule\textquotesingle s System bridge,
as the roster says in words --- $\mapsto$ \texttt{hbridge : ∀ x ∈
Reach, ∃ n : ℤ, k x = (n : ℝ)}: counts integer-valued at every
reachable point, a named hypothesis because lean-aux lacks the net
layer that guarantees it. Then: the hypothesis-position rewrite
$\mapsto$ \texttt{countAtomRounding\_sound} (every reachable point of
the true complement of $\{k \le j\}$ satisfies $\{j + 1 \le k\}$ ---
the side the drain example\textquotesingle s interior row consumes);
``unconditionally never wider'' $\mapsto$
\texttt{countAtomRounding\_tighter}, whose binder row is the width
find: NO integrality, NO reachable scope --- the containment $\{j+1
\le k\} \subseteq \{j < k\}$ holds for arbitrary real observables,
so the substitution never \emph{widens} anywhere; soundness of the
hypothesis-position substitution itself rides the exactness side ---
the integrality hypothesis --- certified separately as
\texttt{countAtomRounding\_eq} --- the
catalogue\textquotesingle s ``the complement over reachable markings
IS $\{k_P \ge j+1\}$'' verbatim as a set identity on the reachable
locus. Width note: \texttt{k : X → ℝ} throughout;
\texttt{\_tighter} now takes $j : \mathbb{R}$ outright (the
containment needs no integrality, so the anchor is stated at the
proposition\textquotesingle s full width), while \texttt{\_sound} and
\texttt{\_eq} keep the $j : \mathbb{Z}$ cast their integer-threshold
bridge demands ---
the count-observable grammar (the $f \equiv 1$ query) is upstream;
divergence D2\textquotesingle s question (which side is load-bearing)
is answered by the bind structure itself: soundness rides
\texttt{\_sound}, the free tightening rides \texttt{\_tighter}, and
each is consumable independently.

\ruleblock{Lean verifier receipts}
\sourcenote{module \texttt{BernsteinLogicAux.Structural}; bound build \texttt{1222cdcc326f\ldots} (bound from receipt cb0c3820c6877f22, 2026-07-29T01:03:36Z). A receipt records that named checking layers accepted a build: (module, compiled-artifact SHA-256, environment SHA-256, checker layers), minted when the checking gate ran; the store receipts attesting the bound build are listed with the layers each run credited. Nothing here asks the reader to trust the store: every layer is an independently runnable program over the quoted sources --- rebuild the module and re-run the checkers to reproduce the verdict. Receipts attesting superseded builds credit no layer and are omitted from this rendering (the count of omissions, if any, is noted below).}
\begin{receipttable}
\receiptrow{0b7e68907f45b27f}{2026-07-28T23:48:36Z}{build+decorrelated}{\badge{BadgeGreen}{current}}
\receiptrow{e97d5614d4f08c7e}{2026-07-29T00:00:35Z}{build+decorrelated}{\badge{BadgeGreen}{current}}
\receiptrow{27de21f047e6b194}{2026-07-29T00:49:05Z}{build+replay+hardened+decorrelated}{\badge{BadgeGreen}{current}}
\receiptrow{cb0c3820c6877f22}{2026-07-29T01:03:36Z}{build+replay+hardened+decorrelated}{\badge{BadgeGreen}{current}}
\receiptrow{3530841119a9a607}{2026-07-31T22:00:22Z}{build+decorrelated}{\badge{BadgeGreen}{current}}
\receiptrow{b6cc23eb3e5b7f97}{2026-07-31T22:29:16Z}{build+decorrelated}{\badge{BadgeGreen}{current}}
\receiptrow{bf3dacc3ab7f20ea}{2026-08-02T23:31:31Z}{build+decorrelated}{\badge{BadgeGreen}{current}}
\receiptrow{384a8569d2005df7}{2026-08-03T16:27:35Z}{build+decorrelated}{\badge{BadgeGreen}{current}}
\receiptrow{717e07acf9acac8a}{2026-08-03T22:46:19Z}{build+decorrelated}{\badge{BadgeGreen}{current}}
\end{receipttable}
\sourcenote{11 further receipt(s) in the store attest superseded builds of this module and are omitted from this rendering.}
\layerline{\layeritem{build}{lake build (elaboration)}{\badge{BadgeGreen}{current}}\layeritem{replay}{leanchecker (official kernel replay)}{\badge{BadgeGreen}{current}}\layeritem{hardened}{safe\_verify (hardened screen)}{\badge{BadgeGreen}{current}}\layeritem{decorrelated}{lean4export + nanoda (decorrelated kernel)}{\badge{BadgeGreen}{current}}}
\sourcenote{\badge{BadgeGray}{comparator: none} no fixture in the three-implementation comparator harness (compiled / kernel-interpreted / reference-interpreted) exercises this rule.}
\ruleblock{HOL4 statement}
\absentbox{NO HOL STATEMENT YET}{The HOL4 development checks the base calculus\textquotesingle s wire
format (theory \texttt{bern\_catalogue}). This rewrite is applied
producer-side, before emission: even with the aggregation node baked,
the wire is expected to carry the already-tightened atom, so the HOL
statement that would matter is the aggregation
sequent\textquotesingle s --- and no checker consuming the extension
theory\textquotesingle s rows is ported yet.
Nothing exists for a HOL statement to be about; if a future format ever
carries the un-rounded atom, the rewrite\textquotesingle s exactness
identity becomes the portable statement.
}
\ruleblock{HOL verifier receipts}
\absentbox{NONE}{no HOL verification receipt exists for this rule.}
\ruleblock{Proof}

\sourcenote{Genre: alias rule (tier 3) --- a respelling template. One
variation from the alias norm, stated up front: the unpack's target is
an \emph{atom-denotation identity} consumed by
\rulename{QAggReduce}'s count rows (the pattern of the strict-atom
closure moves), not a separate rule node --- one atom in, one atom out,
producers rewrite before emission. The five alias artefacts follow in
that key: the unpack, the filled-in application, the payload transport,
the exactness identity, and the soundness corollary.}

\begin{brecipe}[the respelling]\label{rec:car-unpack}
The unpack acts on a complemented count atom in hypothesis position:
\[
\mathrm{cmpl}\,\{\, k_P \leq j \,\}
\qquad\rightsquigarrow\qquad
\{\, k_P \geq j + 1 \,\},
\qquad j \in \mathbb{Z},
\]
where $k_P$ is a count observable --- integer-valued at every reachable
marking, because counts are cardinalities (the compile's integrality
guarantee, entering the mechanization as a named hypothesis until the
net layer discharges it definitionally). Payload transport, field for
field:
\begin{center}\small
\begin{tabular}{@{}lll@{}}
\textbf{field} & \textbf{before} & \textbf{after}\\
\midrule
observable & $k_P$ & $k_P$ (verbatim)\\
threshold & $j$ & $j + 1$ (integer bump)\\
orientation & complemented $\leq$ (open $>$) & closed $\geq$\\
scope & reachable markings & reachable markings (verbatim)\\
\end{tabular}
\end{center}
The filled-in application is a scope row of an aggregation obligation:
where a compile node scopes a row ``off the target $\{k_P \leq j\}$'',
the row's hypothesis list carries $k_P \geq j + 1$ in place of the
real-arithmetic complement.
\end{brecipe}

\begin{bprop}[exactness over the reachable locus]\label{prop:car-eq}
On integer-valued points, for $j \in \mathbb{Z}$,
$\mathrm{Reach} \cap \mathrm{cmpl}\{k_P \leq j\}
= \mathrm{Reach} \cap \{k_P \geq j + 1\}$. (Both restrictions are
load-bearing: at a non-integer threshold the integer $n = \lceil j
\rceil$ would satisfy $n > j$ but not $n \geq j + 1$.)
\end{bprop}
\begin{proof}
At a reachable point, $k_P = n$ for some integer $n$; and for integers,
$n > j \iff n \geq j + 1$. Two lines, which is the point: the
complement of $\{k_P \leq j\}$ \emph{over reachable markings} is
$\{k_P \geq j+1\}$, not $\{k_P > j\}$ --- nothing is lost.
\end{proof}

\begin{bprop}[free tightening, unconditionally]\label{prop:car-tight}
$\{k_P \geq j + 1\} \subseteq \{k_P > j\}$ on \emph{all} points,
integer-valued or not.
\end{bprop}
\begin{proof}
$j + 1 \leq k_P$ implies $j < k_P$. So substituting the rounded atom in
hypothesis position never widens the constraint set; what it removes
is exactly the slab $j < k_P < j + 1$ --- strictly, whenever that
slab's preimage is nonempty: unreachable by any marking, but fully
visible to a real-arithmetic relaxation, which is where the gain
lands.
\end{proof}

\begin{bprop}[soundness corollary]\label{prop:car-sound}
Every reachable point of the true complement of $\{k_P \leq j\}$
satisfies the tightened atom; hence any row valid on
$\mathrm{Reach} \cap \{k_P \geq j + 1\}$ is valid on the reachable part
of the true complement --- the direction hypothesis positions consume,
and by Proposition~\ref{prop:car-eq} nothing more is being assumed than
the complement itself, reachably.
\end{bprop}
\begin{proof}
Immediate from Proposition~\ref{prop:car-eq}, left-to-right. The
asymmetry of positions is worth one sentence: in \emph{hypothesis}
position a tighter scope is free (fewer points to verify, no reachable
point dropped, by exactness); in \emph{conclusion} position a tightened
atom would be a stronger claim and is not licensed by this rule ---
the alias is a hypothesis-position discipline, exactly like its
continuous sibling.
\end{proof}

\begin{bremark}[why integrality is not decoration]\label{rem:car-int}
Drop it and exactness fails: a fractional ``count'' at
$k_P = j + \tfrac12$ satisfies the true complement but not the rounded
atom, so a row verified only on the rounded scope would silently skip
it --- if such points were reachable, that is unsoundness in hypothesis
position. The integrality hypothesis is precisely the statement that
they are not: counts are cardinalities of token multisets. Note there is \emph{no} graceful failure mode when the rewrite is
misapplied to a non-count observable: in hypothesis (scope) position
the rounded atom \emph{shrinks} the region a row must be verified on,
so misapplication can silently accept --- the half-integer point above
is exactly the witness. The integrality warrant is load-bearing, not a
convenience.
\end{bremark}

\begin{bremark}[load-bearing on the first worked example]
The drain net's deadline certificate needs its interior row off the
target $\{k \leq 0\}$. Over the real-relaxed scope $k > 0$ the row is
infeasible --- it evaluates to $+\tfrac14$ at $(k, s_2) =
(\tfrac12, 0)$, a point no marking occupies --- while over the rounded
scope $k \geq 1$ it discharges with an exact linear-multiplier witness
(both numbers recompute in the \rulename{QAggReduce} section's
drain-example remark, where the row and its multipliers are displayed).
The rounding is the counting analogue of the strict-atom closure
discipline (polarity-directed rounding of strict literals), with
integrality playing the role topology plays there: both convert an
open condition into the closed atom the row machinery wants, at zero
semantic cost on the states that matter.
\end{bremark}

\begin{anchors}
Proposition \ref{prop:car-eq} & \texttt{countAtomRounding\_eq}\\
Proposition \ref{prop:car-tight} & \texttt{countAtomRounding\_tighter}\\
Proposition \ref{prop:car-sound} & \texttt{countAtomRounding\_sound}\\
integrality bridge & the named hypothesis in \texttt{countAtomRounding\_sound} and \texttt{\_eq} (compile-side guarantee); deliberately absent from \texttt{\_tighter}, whose containment is unconditional\\
consumed by & \rulename{QAggReduce} scope rows (count atoms); the drain example's interior row\\
\end{anchors}


\rulesection{87}{QPAlwaysFromFold}{core rule}
\label{blcatrule-QPAlwaysFromFold}
\ruleblock{The rule (Steele metanotation)}
\sourcenote{as realized in the mechanized development; source row: \texttt{\texttt{QTSL-rules.toml}} --- \emph{QPAlwaysFromFold}. The metanotation statement is typeset by hand for notational quality; this rule cites no Lean statements yet (the absence is stated loudly below). The rule \emph{roster} is machine-checked (every mechanized rule tag round-trips to exactly one catalogue row); the statement \emph{text} is editorial.}
\begin{mathpar}
\inferrule*{
  \mathrm{upper} : \mathsf{qpaf\_upper}(n,\ c_{\mathrm{u}}) \\
  \mathrm{lower} : \mathsf{qpaf\_lower}(n,\ c_{\mathrm{l}}) \\
  \mathsf{qtsl\_palways\_fold}(n)
}{
  \mathcal{M} \vDash \mathsf{SProbGE}\bigl(p,\ \mathsf{PAlways}(\mathsf{SAnd}[\mathsf{SAtomLE}(c, V),\ \mathsf{SAtomGE}(c, V)])\bigr)
}
\end{mathpar}

\par\medskip
\rulename{QPAlwaysFromFold} ties plated-fold safety evidence to the
judgment stratum: it concludes a surely-invariant claim --- a
conserved-charge invariant, held with probability one --- over the
demonic composite of a token-population subject, from certificates
about a plated fold of the population. Coverage telescopes footprint
rows \emph{pathwise} across every resolution of every clock wing: the
telescoping argument reads firings, not clocks, so race, demonic, and
forced transitions are covered uniformly and no scheduler resolution
escapes. Initial membership is checked inside the evaluator\textquotesingle
s walk (the deadline-table precedent): there is no warrant row, no
region-level child, and no boxed region --- the subject\textquotesingle
s marking is unbounded, and a boxed region for it would be a type
error, not a convenience.

A producer cites one certificate per direction: an upper-bound
citation for the $\le$ face, a lower-bound citation for the $\ge$
face --- at least one, at most one row each; the lower\textquotesingle
s features are recomputed by the walk, never trusted. The accepted
conclusion shapes are one atom or the displayed two-sided conjunction
with equal polynomials; anything wider refuses loudly (a narrow walk
mis-verifies nothing).

\cataloguenote{The verification content lives in the evaluator\textquotesingle s
walk, whose semantics are normative in the plated-transfer design:
conclusion resolution, spec shape, atom decomposition, cert
citations, subject agreement including the plate-to-population token
bindings, fold scale, polynomial correspondence against the
count-wire enumeration, init membership from the bag marking,
sibling-claim presence, and the polarity roster --- with the
pathwise-species demand on covered transitions: an expectation-species
footprint row never warrants a surely-invariant. No Lean statement
binds the node yet; the schema and walk are baked, the walk mutant is
red in the theory battery, and the mechanization column is honest
open debt.
}
\ruleblock{Theory rows (geolog)}
\sourcenote{The rule\textquotesingle s certificate content as it is declared in the geolog theory the checker loads --- payload schema, side-condition demands, and axiom rows, quoted from the theory source at every build of this catalogue (a stale quotation cannot exist). Axiom modes: \texttt{chase} = the engine derives; \texttt{report} = a violation is a loud judgment; \texttt{refuse} = a violation gates admission. Side-condition relations are discharged by bounded evaluation at the engine\textquotesingle s ground seam, never by derivation.}
\par\medskip
The baked rows: the node, its direction citations, the walk, and the
plate-population binding the walk\textquotesingle s subject-agreement
step reads.

\declhead{theory Qtsl}{QPAlwaysFromFold}{\badge{BadgeGreen}{quoted from the theory source at render}}
\begin{leansrc}
QPAlwaysFromFold : Set;
\end{leansrc}
\declhead{theory Qtsl}{QPAlwaysFromFold/*}{\badge{BadgeGreen}{quoted from the theory source at render}}
\begin{leansrc}
QPAlwaysFromFold/proof : QPAlwaysFromFold -> QProof;
\end{leansrc}
\declhead{theory Qtsl}{qpaf\_upper}{\badge{BadgeGreen}{quoted from the theory source at render}}
\begin{leansrc}
qpaf_upper : {node: QPAlwaysFromFold, cert: HCert} -> Prop;
\end{leansrc}
\declhead{theory Qtsl}{qpaf\_lower}{\badge{BadgeGreen}{quoted from the theory source at render}}
\begin{leansrc}
qpaf_lower : {node: QPAlwaysFromFold, cert: HCert} -> Prop;
\end{leansrc}
\declhead{theory Qtsl}{ax/wf/qpalwaysfold\_walk}{\badge{BadgeGreen}{quoted from the theory source at render}}
\begin{leansrc}
report ax/wf/qpalwaysfold_walk : forall n : QPAlwaysFromFold.
  |- {node: n} qtsl_palways_fold;
\end{leansrc}
\declhead{theory Qtsl}{hplate\_token}{\badge{BadgeGreen}{quoted from the theory source at render}}
\begin{leansrc}
hplate_token : {plate: HPlate, token: Token} -> Prop;
\end{leansrc}
\ruleblock{Lean statement}
\badge{BadgeAmber}{not yet mechanized} this rule cites no Lean statements: nothing on this page is hash-bound for it, and its dashboard row shows the same absence across the statement, receipt, and correspondence columns.
\ruleblock{HOL4 statement}
\absentbox{NO HOL STATEMENT YET}{No HOL development of the extension theory exists; this
node\textquotesingle s content is walk structure over baked rows.
}
\ruleblock{HOL verifier receipts}
\absentbox{NONE}{no HOL verification receipt exists for this rule.}
\ruleblock{Proof}

\sourcenote{Genre: core rule (tier 1), the plated-transfer genus (safety face). This section states the rule's soundness
obligation exactly, per the catalogue's completeness-over-polish
charter (a rule that can only be named-and-gapped is still a row);
the discharge is open debt on the mechanization column, and the
schema quoted in the entry is the baked ground truth.}

\textbf{Soundness obligation.}\enspace With the walk's ten steps discharged --- conclusion resolution, spec
shape, atom decomposition, direction citations, subject agreement
including plate-to-population token bindings, fold scale, polynomial
correspondence against the count-wire enumeration, initial membership
from the bag marking, sibling-claim presence, and the polarity roster
with its pathwise-species demand --- the cited fold certificates'
footprint rows telescope pathwise across every resolution of every
clock wing, so the conserved-charge invariant holds along every run:
the concluded claim holds surely, and any probability threshold in
$(0,1]$ is sound.

\begin{bremark}[status]
No Lean statement binds this node. The walk semantics are
normative in the plated-transfer design; the walk mutant is red in
the theory battery, which certifies the walk's presence, not the
transfer theorem --- the transfer theorem's mechanization is the open
debt.
\end{bremark}


\rulesection{88}{QFireBoundFromFold}{core rule}
\label{blcatrule-QFireBoundFromFold}
\ruleblock{The rule (Steele metanotation)}
\sourcenote{as realized in the mechanized development; source row: \texttt{\texttt{QTSL-rules.toml}} --- \emph{QFireBoundFromFold}. The metanotation statement is typeset by hand for notational quality; this rule cites no Lean statements yet (the absence is stated loudly below). The rule \emph{roster} is machine-checked (every mechanized rule tag round-trips to exactly one catalogue row); the statement \emph{text} is editorial.}
\begin{mathpar}
\inferrule*{
  \mathrm{cert} : \mathsf{HCert} \\
  d_{\min} > 0 \\
  \mathsf{qtsl\_firebound\_fold}(n)
}{
  \mathcal{M} \vDash \mathsf{SReturnLE}\bigl(B,\ \mathsf{retFireCount}(\mathit{trs})\bigr)
}
\end{mathpar}

\par\medskip
\rulename{QFireBoundFromFold} is the ranking face of the same genus:
from \emph{one} gap-rowed plated certificate it concludes a roster
firing-count ceiling --- the specification\textquotesingle s warranted
component is the count bound over the roster firing-count run
functional. The transfer theorem\textquotesingle s pathwise face,
applied: with footprint, coverage, floor, and flow discharged, every
roster firing drops the fold by at least $d_{\min}$ in the reading
scale, the fold never rises anywhere else, and the fold is floored at
zero --- so roster firings number at most
$\mathrm{fold}_\varphi(m_0)/d_{\min}$ along \emph{every} run,
uniformly in token count and interleaving.

What the checker recomputes: two exact inequalities --- each roster
row\textquotesingle s gap at least $d_{\min}$, and $B \cdot d_{\min}$
at least the fold\textquotesingle s initial value --- and never an
argmin: $d_{\min}$ is pinned by the producer, the same witness
species as a pinned tolerance. Counting is firing-metered per roster
member --- deliberately not region-entry counting (the crossing-count
acceptance device) and not the instant-crossing count (the
crossing-budget family\textquotesingle s compensator-mass species):
three counting species, three devices, none interchangeable.

\cataloguenote{Closed-judgment route, deliberately: no warrant row, no region-level
child, no boxed certificate-family table for an unbounded marking;
init arithmetic lives inside the walk. The cert citation is a
function, not a relation --- a ranking consequence cites exactly one
certificate, and exactly-one is semantic here (totality is the
sort\textquotesingle s own obligation); contrast the safety
sibling\textquotesingle s two-direction relation shape. No Lean
statement binds the node yet; the schema, the positivity report row,
and the walk are baked with red mutants in the theory battery, and
the mechanization column is honest open debt.
}
\ruleblock{Theory rows (geolog)}
\sourcenote{The rule\textquotesingle s certificate content as it is declared in the geolog theory the checker loads --- payload schema, side-condition demands, and axiom rows, quoted from the theory source at every build of this catalogue (a stale quotation cannot exist). Axiom modes: \texttt{chase} = the engine derives; \texttt{report} = a violation is a loud judgment; \texttt{refuse} = a violation gates admission. Side-condition relations are discharged by bounded evaluation at the engine\textquotesingle s ground seam, never by derivation.}
\par\medskip
The baked rows: the node with its pinned decrement, the positivity
report, the walk, and the roster relation the concluded run
functional reads.

\declhead{theory Qtsl}{QFireBoundFromFold}{\badge{BadgeGreen}{quoted from the theory source at render}}
\begin{leansrc}
QFireBoundFromFold : Set;
\end{leansrc}
\declhead{theory Qtsl}{QFireBoundFromFold/*}{\badge{BadgeGreen}{quoted from the theory source at render}}
\begin{leansrc}
QFireBoundFromFold/proof : QFireBoundFromFold -> QProof;
QFireBoundFromFold/cert  : QFireBoundFromFold -> HCert;
QFireBoundFromFold/dmin  : QFireBoundFromFold -> Q;
\end{leansrc}
\declhead{theory Qtsl}{ax/wf/qfirebound\_dmin\_pos}{\badge{BadgeGreen}{quoted from the theory source at render}}
\begin{leansrc}
report ax/wf/qfirebound_dmin_pos : forall n : QFireBoundFromFold,
    d : Q.
  n QFireBoundFromFold/dmin = d
  |- {l: 0, r: d} q_lt;
\end{leansrc}
\declhead{theory Qtsl}{ax/wf/qfirebound\_walk}{\badge{BadgeGreen}{quoted from the theory source at render}}
\begin{leansrc}
report ax/wf/qfirebound_walk : forall n : QFireBoundFromFold.
  |- {node: n} qtsl_firebound_fold;
\end{leansrc}
\declhead{theory Qtsl}{retFireCount}{\badge{BadgeGreen}{quoted from the theory source at render}}
\begin{leansrc}
retFireCount : {trs: TrList} -> PathObservable;
\end{leansrc}
\declhead{theory Qtsl}{qtsl\_firecount\_roster}{\badge{BadgeGreen}{quoted from the theory source at render}}
\begin{leansrc}
qtsl_firecount_roster : {sys: System, trs: TrList} -> Prop;
\end{leansrc}
\ruleblock{Lean statement}
\badge{BadgeAmber}{not yet mechanized} this rule cites no Lean statements: nothing on this page is hash-bound for it, and its dashboard row shows the same absence across the statement, receipt, and correspondence columns.
\ruleblock{HOL4 statement}
\absentbox{NO HOL STATEMENT YET}{No HOL development of the extension theory exists; this
node\textquotesingle s content is walk structure over baked rows.
}
\ruleblock{HOL verifier receipts}
\absentbox{NONE}{no HOL verification receipt exists for this rule.}
\ruleblock{Proof}

\sourcenote{Genre: core rule (tier 1), the plated-transfer genus (ranking face). This section states the rule's soundness
obligation exactly, per the catalogue's completeness-over-polish
charter (a rule that can only be named-and-gapped is still a row);
the discharge is open debt on the mechanization column, and the
schema quoted in the entry is the baked ground truth.}

\textbf{Soundness obligation.}\enspace With footprint, coverage, floor, and flow discharged for the one
gap-rowed certificate, every roster firing drops the fold by at least
$d_{\min}$ in the reading scale, the fold never rises anywhere else,
and the fold is floored at zero; hence roster firings number at most
$\mathrm{fold}_\varphi(m_0)/d_{\min}$ along every run, uniformly in
token count and interleaving --- the concluded firing-count ceiling.
The checker's contribution is two exact inequalities (each roster
row's gap at least $d_{\min}$; the ceiling times $d_{\min}$ at least
the fold's initial value); $d_{\min}$ is pinned by the producer and
never derived.

\begin{bremark}[status]
No Lean statement binds this node. The walk semantics are
normative in the transfer-former design (twelve steps); the theory
battery's mutants are red at the walk and positivity loci.
\end{bremark}


\rulesection{89}{QPAlwaysTFromFold}{core rule}
\label{blcatrule-QPAlwaysTFromFold}
\ruleblock{The rule (Steele metanotation)}
\sourcenote{as realized in the mechanized development; source row: \texttt{\texttt{QTSL-rules.toml}} --- \emph{QPAlwaysTFromFold}. The metanotation statement is typeset by hand for notational quality; this rule cites no Lean statements yet (the absence is stated loudly below). The rule \emph{roster} is machine-checked (every mechanized rule tag round-trips to exactly one catalogue row); the statement \emph{text} is editorial.}
\begin{mathpar}
\inferrule*{
  n\,\mathsf{cert} = c \\ n\,\mathsf{goalfold} = f_g \\ n\,\mathsf{horizon} = T \\
  \mathsf{coverage\_lb}(c) \\ \mathsf{flow\_lb}(c) \\
  \forall r.\ u_r > 0 \Rightarrow \mathsf{TrMeter}(\mathit{tr}(r), d_r) \\
  \mathsf{qtsl\_palwayst\_fold}(n)
}{
  \mathcal{M} \vDash \mathsf{SProbGE}\bigl(p,\ \mathsf{PAlwaysT}(T,\ \mathrm{cmpl}(\mathit{goal}))\bigr)
}
\end{mathpar}

\par\medskip
\rulename{QPAlwaysTFromFold} is the timed-floor face of the
plated-transfer genus --- the third sibling beside
\rulename{QPAlwaysFromFold} (safety) and \rulename{QFireBoundFromFold}
(ranking): from one certificate whose footprint rows carry the
\emph{drop-at-most} polarity (per firing, the fold decreases by at
most $u_r$ --- including $u_r = 0$, never decreases), together with
per-row clock floors, it concludes that no run of the system reaches
a designated goal region within a horizon $T$ --- a surely-satisfied
time-bounded invariance claim over the goal\textquotesingle s
complement.

The soundness assembly is the transfer telescope\textquotesingle s
mirror face plus a meter law: with drop-at-most coverage of every
may-touch transition, mirror flow, per-row footprint sequents, and
--- for conjugated reading kinds --- the conjugation-parameter
conformance discharged (integrality of conjugated powers and the
order-preservation domain; the transformed-space telescope is
unjustified without it), every run satisfies
$\mathrm{fold}_\varphi(m_t) \ge \mathrm{fold}_\varphi(m_0) - \sum_r
u_r N_r(t)$; with clock floors $d_r > 0$ on every counted row,
$N_r(t) \le \max(0, (t - o_r)/d_r) + 1$ under every resolution. A run
reaching the goal region (where the fold is at most $f_g$) by time
$T$ would force $\mathrm{fold}_\varphi(m_0) - f_g \le \sum_{r : u_r >
0} u_r \cdot (\max(0, (T - o_r)/d_r) + 1)$ --- and the
checker\textquotesingle s whole trusted contribution is ONE exact
rational verification of the strict reverse inequality. $T$, $f_g$,
and every $u_r$ are pinned by the producer; the checker never
computes an argmin or an optimum.

Counting from below stays deliberately \emph{inside} the walk: the
specification language refuses a standing ``at least $n$ events
occur\textquotesingle\textquotesingle{} former (a liveness claim
consuming survival mass), and this rule honors that refusal --- the
conclusion is safety-polarity, violated by finite prefixes. The
conclusion is Zeno-robust without any crossing budget: rows with
$u_r = 0$ (instantaneous rows included) may fire unboundedly, since
they cannot drop the fold; an instantaneous row with $u_r > 0$
refuses outright (it has no time floor to meter).

\cataloguenote{The schema is baked in the extension theory: the node, the
fourteen-step walk (normative in the judgment-tie design), the two
mirror side-conditions, the per-wing meter demands, and the
drop-at-most footprint polarity are all landed rows, quoted below
from the theory source. The polarity\textquotesingle s formation
carried its guard obligations into the bake --- most critically that
the \emph{existing} upper-coverage side condition counts exactly
rows of the two elder polarities (safety, gap), stated as a closed
positive membership rather than a prohibition: a drop-at-most row
bounds nothing above, and an evaluator that counted one would unsound
the two elder siblings.
The mechanization ladder names the discrete-unfolding profile first
(round-indexed runs); the continuous lift waits on a first
formalization of the clock meter semantics, shared with
\rulename{QRateFromClock} and \rulename{QPAlwaysTFromLineage}. No
Lean statement binds the node yet; the mechanization column is
honest open debt.
}
\ruleblock{Theory rows (geolog)}
\sourcenote{The rule\textquotesingle s certificate content as it is declared in the geolog theory the checker loads --- payload schema, side-condition demands, and axiom rows, quoted from the theory source at every build of this catalogue (a stale quotation cannot exist). Axiom modes: \texttt{chase} = the engine derives; \texttt{report} = a violation is a loud judgment; \texttt{refuse} = a violation gates admission. Side-condition relations are discharged by bounded evaluation at the engine\textquotesingle s ground seam, never by derivation.}
\par\medskip
The baked rows: the node with its pinned horizon and goal-fold
fields, the optional lineage-offsets citation, the walk, the two
mirror side-conditions on the cited certificate, the per-wing meter
demands, and the drop-at-most polarity with its displacement bound.

\declhead{theory Qtsl}{QPAlwaysTFromFold}{\badge{BadgeGreen}{quoted from the theory source at render}}
\begin{leansrc}
QPAlwaysTFromFold : Set;
\end{leansrc}
\declhead{theory Qtsl}{QPAlwaysTFromFold/*}{\badge{BadgeGreen}{quoted from the theory source at render}}
\begin{leansrc}
QPAlwaysTFromFold/proof : QPAlwaysTFromFold -> QProof;
QPAlwaysTFromFold/cert  : QPAlwaysTFromFold -> HCert;
QPAlwaysTFromFold/goalfold : QPAlwaysTFromFold -> Q;   // f_g, pinned
QPAlwaysTFromFold/horizon  : QPAlwaysTFromFold -> Q;   // T, pinned
\end{leansrc}
\declhead{theory Qtsl}{ax/wf/qpatf\_horizon\_nonneg}{\badge{BadgeGreen}{quoted from the theory source at render}}
\begin{leansrc}
report ax/wf/qpatf_horizon_nonneg : forall n : QPAlwaysTFromFold,
    x : Q.
  n QPAlwaysTFromFold/horizon = x
  |- {l: 0, r: x} q_le;
\end{leansrc}
\declhead{theory Qtsl}{qpatf\_offsets}{\badge{BadgeGreen}{quoted from the theory source at render}}
\begin{leansrc}
qpatf_offsets : {node: QPAlwaysTFromFold,
                 lineage: QPAlwaysTFromLineage} -> Prop;
\end{leansrc}
\declhead{theory Qtsl}{qtsl\_palwayst\_fold}{\badge{BadgeGreen}{quoted from the theory source at render}}
\begin{leansrc}
qtsl_palwayst_fold : {node: QPAlwaysTFromFold} -> Prop;
\end{leansrc}
\declhead{theory Qtsl}{ax/tf/qpatf\_walk}{\badge{BadgeGreen}{quoted from the theory source at render}}
\begin{leansrc}
report ax/tf/qpatf_walk : forall n : QPAlwaysTFromFold.
  |- {node: n} qtsl_palwayst_fold;
\end{leansrc}
\declhead{theory Qtsl}{ax/tf/qpatf\_coverage\_lb}{\badge{BadgeGreen}{quoted from the theory source at render}}
\begin{leansrc}
report ax/tf/qpatf_coverage_lb : forall n : QPAlwaysTFromFold,
    c : HCert.
  n QPAlwaysTFromFold/cert = c
  |- {cert: c} qtsl_hplate_coverage_lb;
\end{leansrc}
\declhead{theory Qtsl}{ax/tf/qpatf\_flow\_lb}{\badge{BadgeGreen}{quoted from the theory source at render}}
\begin{leansrc}
report ax/tf/qpatf_flow_lb : forall n : QPAlwaysTFromFold,
    c : HCert.
  n QPAlwaysTFromFold/cert = c
  |- {cert: c} qtsl_hplate_flow_lb;
\end{leansrc}
\declhead{theory Qtsl}{ax/tf/qpatf\_meter\_forced}{\badge{BadgeGreen}{quoted from the theory source at render}}
\begin{leansrc}
report ax/tf/qpatf_meter_forced : forall n : QPAlwaysTFromFold,
    c : HCert, r : HFootprintRow, t : Transition, d : Q,
    k : TrClock.
  n QPAlwaysTFromFold/cert = c,
  r HFootprintRow/cert = c,
  r HFootprintRow/polarity = hrpDropAtMost,
  r HFootprintRow/tr = t,
  t Transition/clock = k,
  k = {dwell: d} tkForced
  |- {tr: t, d: d} qtsl_tr_meter;
\end{leansrc}
\declhead{theory Qtsl}{ax/tf/qpatf\_meter\_demonic}{\badge{BadgeGreen}{quoted from the theory source at render}}
\begin{leansrc}
report ax/tf/qpatf_meter_demonic : forall n : QPAlwaysTFromFold,
    c : HCert, r : HFootprintRow, t : Transition, d : Q,
    e : EnnRat, k : TrClock.
  n QPAlwaysTFromFold/cert = c,
  r HFootprintRow/cert = c,
  r HFootprintRow/polarity = hrpDropAtMost,
  r HFootprintRow/tr = t,
  t Transition/clock = k,
  k = {dwellMin: e} tkDemonic,
  e = {q: d} enFin
  |- {tr: t, d: d} qtsl_tr_meter;
\end{leansrc}
\declhead{theory Qtsl}{ax/tf/qpatf\_meter\_carved}{\badge{BadgeGreen}{quoted from the theory source at render}}
\begin{leansrc}
report ax/tf/qpatf_meter_carved : forall n : QPAlwaysTFromFold,
    c : HCert, r : HFootprintRow, t : Transition, fl : Q,
    ce : EnnRat, k : TrClock.
  n QPAlwaysTFromFold/cert = c,
  r HFootprintRow/cert = c,
  r HFootprintRow/polarity = hrpDropAtMost,
  r HFootprintRow/tr = t,
  t Transition/clock = k,
  k = {floor: fl, ceil: ce} tkCarved
  |- {tr: t, d: fl} qtsl_tr_meter;
\end{leansrc}
\declhead{theory Qtsl}{hrpDropAtMost}{\badge{BadgeGreen}{quoted from the theory source at render}}
\begin{leansrc}
hrpDropAtMost : HRowPolarity;
\end{leansrc}
\declhead{theory Qtsl}{qtsl\_hplate\_coverage\_lb}{\badge{BadgeGreen}{quoted from the theory source at render}}
\begin{leansrc}
qtsl_hplate_coverage_lb : {cert: HCert} -> Prop;
\end{leansrc}
\declhead{theory Qtsl}{qtsl\_hplate\_flow\_lb}{\badge{BadgeGreen}{quoted from the theory source at render}}
\begin{leansrc}
qtsl_hplate_flow_lb : {cert: HCert} -> Prop;
\end{leansrc}
\ruleblock{Lean statement}
\badge{BadgeAmber}{not yet mechanized} this rule cites no Lean statements: nothing on this page is hash-bound for it, and its dashboard row shows the same absence across the statement, receipt, and correspondence columns.
\ruleblock{HOL4 statement}
\absentbox{NO HOL STATEMENT YET}{No HOL development of the extension theory exists; this
node\textquotesingle s content is walk structure over baked rows.
}
\ruleblock{HOL verifier receipts}
\absentbox{NONE}{no HOL verification receipt exists for this rule.}
\ruleblock{Proof}

\sourcenote{Genre: core rule (tier 1), the plated-transfer genus
(timed-floor face). This section states the rule's soundness
obligation exactly, per the catalogue's completeness-over-polish
charter; the rule's schema and walk are baked in the extension
theory (rows quoted in this section), and the discharge below is the
named obligation of the mechanization column.}

\textbf{Soundness obligation (mirror telescope + meter law).}\enspace
Let $F := \mathrm{fold}_\varphi(m_0)$, recomputed by the walk from the
declared initial marking. (Telescope.) With drop-at-most coverage of
every may-touch transition, per-row footprint sequents (produced-side
fold contribution at least consumed side minus $u_r$, in the reading
scale, under the guard hypotheses, for all colors in the boxes),
mirror flow (the fold does not decrease between firings), and --- for
conjugated reading kinds --- the conjugation-parameter conformance
(integrality of conjugated powers; the order-preservation domain,
even powers requiring certified nonnegative feature ranges: without
it the transformed-space reading of the telescope is unjustified)
discharged,
every run satisfies
$\mathrm{fold}_\varphi(m_t) \ge F - \sum_r u_r N_r(t)$, where
$N_r(t)$ counts firings of row $r$ by time $t$. (Meter.) With every
counted row ($u_r > 0$) carrying a clock floor $d_r > 0$ and
first-firing offset $o_r$ (from a cited lineage node, else the
one-dwell default $o_r := d_r$),
$N_r(t) \le \max(0, (t - o_r)/d_r) + 1$ under every resolution.
(Assembly.) The goal region carries a certified fold ceiling $f_g$;
a run reaching it at $t \le T$ forces
$F - f_g \;\le\; \sum_r u_r N_r(t) \;\le\; \sum_{r : u_r > 0} u_r
\cdot \bigl(\max(0, \tfrac{T - o_r}{d_r}) + 1\bigr)$,
using $u_r \ge 0$ throughout and monotonicity in $t$; the
checker-verified strict inequality
$\sum_{r : u_r > 0} u_r \cdot (\max(0, \tfrac{T - o_r}{d_r}) + 1) <
F - f_g$ refutes it, so no admitted run reaches the goal region
within $T$ --- the concluded time-bounded invariance, surely.

The conclusion is safety-polarity and consumes no crossing budget:
rows with $u_r = 0$ (instantaneous rows included) may fire
unboundedly without touching the telescope's right side; an
instantaneous row with $u_r > 0$ refuses at the walk's meter step.

\begin{bremark}[status]
The schema is baked: the node, its fourteen-step walk, the mirror
side-conditions, the per-wing meter demands, and the drop-at-most
polarity are landed theory rows (quoted in this section). No Lean
statement binds the node yet: the named obligation is the
discrete-unfolding profile first (round-indexed runs), with the
continuous lift behind the clock meter semantics' first
formalization --- the keystone shared across this rule family.
\end{bremark}


\rulesection{90}{QPAlwaysTFromLineage}{core rule}
\label{blcatrule-QPAlwaysTFromLineage}
\ruleblock{The rule (Steele metanotation)}
\sourcenote{as realized in the mechanized development; source row: \texttt{\texttt{QTSL-rules.toml}} --- \emph{QPAlwaysTFromLineage}. The metanotation statement is typeset by hand for notational quality; this rule cites no Lean statements yet (the absence is stated loudly below). The rule \emph{roster} is machine-checked (every mechanized rule tag round-trips to exactly one catalogue row); the statement \emph{text} is editorial.}
\begin{mathpar}
\inferrule*{
  \mathsf{unit\_lineage}(c) \\
  \{\mathsf{hrpGap}(1), \mathsf{hrpDropAtMost}(1)\}\ \text{per advancing row on } \chi = -\psi\ \text{(pathwise species)} \\
  \mathsf{TrMeter}(\mathit{tr}(r), d)\ \text{per advancing row} \\
  \forall \tau \in m_0.\ \psi(\tau) \le 0 \\
  T < k \cdot d \\
  \mathsf{qtsl\_palwayst\_lineage}(n)
}{
  \mathcal{M} \vDash \mathsf{SProbGE}\bigl(p,\ \mathsf{PAlwaysT}(T,\ \textstyle\bigwedge_{q\,:\,\mathrm{depth}(q) \ge k} \mathrm{count}_q = 0)\bigr)
}
\end{mathpar}

\par\medskip
\rulename{QPAlwaysTFromLineage} is a causal-depth meter: tokens whose
production chain has depth $k$ cannot exist strictly before time
$k \cdot d$. It concludes plate-emptiness --- no token of declared
causal depth at least $k$ exists at any time $t \le T$, for a pinned
horizon $T$ checked STRICTLY below $k \cdot d$ (the induction proves
emptiness through round $k - 1$ and permits a depth-$k$ token in the
round-$k$ marking, so the inclusive endpoint is deliberately not
concluded; on the round grid $T = (k-1) \cdot d$ is the tightest
warranted choice, and $k$ is a positive integer) --- from four
ingredients: \emph{unit lineage} (every covered transition consumes
at most one and produces at most one token across the cited plates,
so per-token lineage is a path: for a row consuming a cited token and
producing a cited successor, the fold-level displacement of the depth
dictionary is the successor\textquotesingle s depth minus the
consumed token\textquotesingle s; birth rows root at depth at most
one, and a death row on a positive-depth token satisfies neither
bracket and is refused --- a completeness cost, never an
unsoundness); a
\emph{bracket} of footprint rows on the negated depth dictionary
pinning each chain firing\textquotesingle s advance to exactly one
(the drop-at-most row and the gap row at equal witness --- both sides
load-bearing, since a decrease bound alone would admit one firing
advancing depth by ten --- and every bracket row of the pathwise
deterministic-reset species: a stochastic moment row could satisfy
both sequents at expected advance one while individual outcomes jump
farther, defeating the surely conclusion); clock floors $d$ on every
advancing row; and
initial depth at most zero, recomputed by the walk from the declared
initial marking (universally over its tokens --- vacuous when empty).

The rule has a second face: a \rulename{QPAlwaysTFromFold} node may
cite it as the source of its per-row first-firing offsets ($o_r$ :=
the row\textquotesingle s consumed-tuple minimum declared depth over
the guard box, times $d$ --- recomputed by the CONSUMING
walk, never read from stored producer data, because an inflated
offset errs on the unsound side; a row whose guard box admits no
consumed cited token takes the consumer\textquotesingle s one-dwell
default instead). This face is unaffected by the strict endpoint: a
firing that consumes a depth-$k$ token happens at round $k$ or
later, time at least $k \cdot d$, which is exactly the offset claim.

The rule mints \emph{profile-gated}, and the gate is stated honestly
rather than smoothed over: per-token chain spacing is NOT derivable
from the per-transition meter law when several tokens share a guard
--- a transition\textquotesingle s accrued dwell can be spent on a
freshly produced token. This is the classical timed-Petri-net
memory-policy fork (single-server versus per-token timers). The
profile, exactly: every covered row carries the one period $d$ as its
dwell parameter (demonic and forced wings admitted; a carved wing at
floor equal to ceiling $d$, so firings cannot fall off the round
grid); the subject is declared at the round-indexed step reading ---
firings at round $r$ read the round-$r$ marking and produce into
round $r + 1$, rounds at times $r \cdot d$, the theory\textquotesingle s
own discrete-time-unfolding reading, whose canonical case is the
all-demonic uniform-dwell net; and the cited plates\textquotesingle{}
token colors do not evolve between firings (the brackets constrain
firing displacement only, so between-round color drift would evade
them). Under that profile the per-token spacing IS the step
law\textquotesingle s own content and the soundness induction is
clean; the continuous-time lift is a genuine semantics decision the
proposal prices without baking.

\cataloguenote{The schema is baked in the extension theory: the node, its walk, and
its structural side-conditions are quoted below from the theory
source. Profile-gated from birth, per the catalogue\textquotesingle s
precedent of rules proven at a named profile and filed beyond it; the
concluded specification former itself is profile-independent --- the
gate lives in the walk\textquotesingle s side conditions, never in
the conclusion\textquotesingle s meaning. A worked instance of the
discrete induction (a fabrication-line model with a 370-deep
production chain) is already kernel-checked; the oracle obligation
that generalizes it to the profile rule remains owed, and no Lean
statement binds the node yet --- the mechanization column is honest
open debt.
}
\ruleblock{Theory rows (geolog)}
\sourcenote{The rule\textquotesingle s certificate content as it is declared in the geolog theory the checker loads --- payload schema, side-condition demands, and axiom rows, quoted from the theory source at every build of this catalogue (a stale quotation cannot exist). Axiom modes: \texttt{chase} = the engine derives; \texttt{report} = a violation is a loud judgment; \texttt{refuse} = a violation gates admission. Side-condition relations are discharged by bounded evaluation at the engine\textquotesingle s ground seam, never by derivation.}
\par\medskip
The baked rows: the node with its positive depth, positive period,
and nonnegative horizon reports, the unit-lineage structural
side-condition on the cited certificate, and the causal-depth walk.

\declhead{theory Qtsl}{QPAlwaysTFromLineage}{\badge{BadgeGreen}{quoted from the theory source at render}}
\begin{leansrc}
QPAlwaysTFromLineage : Set;
\end{leansrc}
\declhead{theory Qtsl}{QPAlwaysTFromLineage/*}{\badge{BadgeGreen}{quoted from the theory source at render}}
\begin{leansrc}
QPAlwaysTFromLineage/proof : QPAlwaysTFromLineage -> QProof;
QPAlwaysTFromLineage/cert  : QPAlwaysTFromLineage -> HCert;
QPAlwaysTFromLineage/depth   : QPAlwaysTFromLineage -> Nat;
QPAlwaysTFromLineage/period  : QPAlwaysTFromLineage -> Q;
QPAlwaysTFromLineage/horizon : QPAlwaysTFromLineage -> Q;
\end{leansrc}
\declhead{theory Qtsl}{ax/wf/qpatl\_depth\_pos}{\badge{BadgeGreen}{quoted from the theory source at render}}
\begin{leansrc}
report ax/wf/qpatl_depth_pos : forall n : QPAlwaysTFromLineage,
    k : Nat.
  n QPAlwaysTFromLineage/depth = k
  |- {l: 0, r: k} nat_lt;
\end{leansrc}
\declhead{theory Qtsl}{ax/wf/qpatl\_period\_pos}{\badge{BadgeGreen}{quoted from the theory source at render}}
\begin{leansrc}
report ax/wf/qpatl_period_pos : forall n : QPAlwaysTFromLineage,
    d : Q.
  n QPAlwaysTFromLineage/period = d
  |- {l: 0, r: d} q_lt;
\end{leansrc}
\declhead{theory Qtsl}{ax/wf/qpatl\_horizon\_nonneg}{\badge{BadgeGreen}{quoted from the theory source at render}}
\begin{leansrc}
report ax/wf/qpatl_horizon_nonneg : forall n : QPAlwaysTFromLineage,
    x : Q.
  n QPAlwaysTFromLineage/horizon = x
  |- {l: 0, r: x} q_le;
\end{leansrc}
\declhead{theory Qtsl}{qtsl\_unit\_lineage}{\badge{BadgeGreen}{quoted from the theory source at render}}
\begin{leansrc}
qtsl_unit_lineage : {cert: HCert} -> Prop;
\end{leansrc}
\declhead{theory Qtsl}{qtsl\_palwayst\_lineage}{\badge{BadgeGreen}{quoted from the theory source at render}}
\begin{leansrc}
qtsl_palwayst_lineage : {node: QPAlwaysTFromLineage} -> Prop;
\end{leansrc}
\declhead{theory Qtsl}{ax/tf/qpatl\_walk}{\badge{BadgeGreen}{quoted from the theory source at render}}
\begin{leansrc}
report ax/tf/qpatl_walk : forall n : QPAlwaysTFromLineage.
  |- {node: n} qtsl_palwayst_lineage;
\end{leansrc}
\declhead{theory Qtsl}{ax/tf/qpatl\_unit\_lineage}{\badge{BadgeGreen}{quoted from the theory source at render}}
\begin{leansrc}
report ax/tf/qpatl_unit_lineage : forall n : QPAlwaysTFromLineage,
    c : HCert.
  n QPAlwaysTFromLineage/cert = c
  |- {cert: c} qtsl_unit_lineage;
\end{leansrc}
\ruleblock{Lean statement}
\badge{BadgeAmber}{not yet mechanized} this rule cites no Lean statements: nothing on this page is hash-bound for it, and its dashboard row shows the same absence across the statement, receipt, and correspondence columns.
\ruleblock{HOL4 statement}
\absentbox{NO HOL STATEMENT YET}{No HOL development of the extension theory exists; this
node\textquotesingle s content is walk structure over baked rows.
}
\ruleblock{HOL verifier receipts}
\absentbox{NONE}{no HOL verification receipt exists for this rule.}
\ruleblock{Proof}

\sourcenote{Genre: core rule (tier 1), the plated-transfer genus
(causal-depth face), profile-gated from birth. This section states
the rule's soundness obligation exactly, gate included; the rule's
schema and walk are baked in the extension theory (rows quoted in
this section), and the discharge below is the named obligation of
the mechanization column.}

\textbf{Soundness obligation (depth induction at the
discrete-unfolding profile).}\enspace At the gated profile (every
covered row carrying the one period $d$ as its dwell parameter ---
demonic and forced wings admitted, carved wings at floor equal to
ceiling $d$; the subject declared at the round-indexed step reading,
firings at round $r$ reading the round-$r$ marking and producing
into round $r + 1$, rounds at times $r \cdot d$; cited
plates' token colors static between firings), by
induction over rounds: unit lineage makes each token's production
history a path through covered firings; the bracket rows on the
negated depth dictionary --- every one of the pathwise
deterministic-reset species --- pin each chain firing's depth advance
to exactly one (the gap row bounds the advance from below, the
drop-at-most row from above --- both sides load-bearing); initial
depth is at most zero, universally over the declared initial
marking's tokens (vacuous when empty), by the walk's recomputation.
Hence a token present at round $m$ arose through at
most $m$ chain firings in strictly increasing rounds and has declared
depth at most $m$ --- so plates of declared depth at least $k$
($k \ge 1$) are empty at rounds $0, \ldots, k-1$: emptiness at every
time $t \le T$ for the pinned horizon $T < k \cdot d$, surely. The
endpoint is strict by construction --- the induction permits a
depth-$k$ token in the round-$k$ marking itself, so the inclusive
horizon $k \cdot d$ is deliberately not concluded.

\textbf{The gate, exactly.}\enspace Beyond the profile, per-token
chain spacing is NOT a consequence of the per-transition meter law
when tokens share a guard: a transition's accrued dwell can be spent
on a freshly produced token, collapsing the chain's temporal spacing
to one dwell for arbitrarily many chain steps. This is the classical
timed-Petri-net memory-policy fork (single-server versus per-token
timers); the continuous-time lift is a semantics decision the
proposal prices without baking, and the rule refuses outside the
profile rather than concluding from an unproven spacing premise.

\begin{bremark}[status]
The schema is baked: the node, its walk, and its structural
side-conditions are landed theory rows (quoted in this section). No
Lean statement binds the node yet. A worked instance of the discrete
induction (a fabrication-line model with a 370-deep production
chain) is already kernel-checked; the named obligation generalizes
it to the profile rule.
\end{bremark}


\blcatsection{Distance Transport}
\rulesection{91}{QCoupleDrift}{core rule}
\label{blcatrule-QCoupleDrift}
\ruleblock{The rule (Steele metanotation)}
\sourcenote{as realized in the mechanized development; source row: \texttt{\texttt{QTSL-rules.toml}} --- \emph{QCoupleDrift}. The metanotation statement is typeset by hand for notational quality; its mechanized ground truth is the Lean statement below, extracted from the sources at every build of this book and hash-bound. The rule \emph{roster} is machine-checked (every mechanized rule tag round-trips to exactly one catalogue row); the statement \emph{text} is editorial, checked by eye against the extracted Lean beside it.}
\begin{mathpar}
\inferrule*{
  \bigl\{\vdash K(p, \cdot) \ge 0,\ \textstyle\sum_q K(p, q) = 1,\ \pi_{i\#} K(p, \cdot) = \widehat{K}_i(p_i, \cdot)\bigr\} \ \text{(coupling rows, exact } \Q\text{; } \widehat{K}_i \ \text{the completed component kernels)} \\
  \bigl\{\vdash K(p, p) = 1 \ \text{for both-absorbed } p\bigr\} \ \text{(freeze rows)} \\
  \bigl\{\vdash \textstyle\sum_q K(p, q)\, V(q) \le V(p)\bigr\} \ \text{(relational supermartingale generator row, every product state)} \\
  \vdash V \ge 0 \\
  \bigl\{\vdash V \ge 1 \ \text{(disagreeing), resp.\ } V \ge |v_1 - v_2|, \ \text{on both-absorbed states}\bigr\} \ \text{(terminal domination)} \\
  \text{absorption: per-component reach criterion}; \quad \mu_0 \ \text{a probability law with } \pi_{i\#}\mu_0 = \text{the component initial laws (initial marginal rows)}; \quad \vdash \textstyle\sum_p \mu_0(p)\, V(p) \le b
}{
  \mathsf{QDist}[d](\mathcal{M}_1, \mathcal{M}_2) \le b \qquad d \in \{\mathrm{TV},\ \text{gap-pseudometrics}\}
}
\end{mathpar}

\ruleflavor{the coupling marginals are tied to the \emph{completed} component kernels $\widehat{K}_i$ (absorbing states frozen --- the completion carries the component freezes), and the absorption criterion is recomputed by the checker (a per-component graph reach). Coupling emptiness refuses, never lies; the $\mathbb{W}_1$ registry entry and the mixture-polynomial tier are designed, not yet realized --- the finite-chain tier above is the proven fragment.}
\par\medskip
\rulename{QCoupleDrift} is the coupling-witnessed discharge of the
distance judgment: from a certified coupling of two absorbing Markov
chains and a relational supermartingale on the product, it concludes a
bound on how far the two systems\textquotesingle{} final readouts can
differ --- total-variation distance of the final readout laws, or the
gap of final expected values. A producer reaches for it to compare two
models of the same system (a nominal chain against a faulted one, a
fine model against a simplification): the producer chooses a coupling
--- a product-space kernel whose marginals are the two chains --- that
keeps the pair together where the chains agree, and a nonnegative
function $V$ that decreases along the coupled dynamics and dominates
the disagreement indicator at absorption. The certified bound is
$V$\textquotesingle s initial expectation.

What the checker recomputes, all in exact rational arithmetic: that the
coupling rows form a genuine kernel (entrywise nonnegative, unit row
sums) whose row marginals are exactly the two completed component
kernels; that both-absorbed states are frozen; the drift inequality
$\sum_q K(p,q)\, V(q) \le V(p)$ at every product state; nonnegativity
of $V$; the terminal-domination rows ($V \ge 1$ on disagreeing
both-absorbed states for the total-variation form, $V \ge |v_1 - v_2|$
for the expectation-gap form); the per-component absorption criterion
as a finite-graph reachability check; and the initial bound
$\sum_p \mu_0(p)\, V(p) \le b$.

Each hypothesis is load-bearing, and the proof section walks the
omission counterexamples: broken marginals sever the product chain from
the systems being compared; without nonnegativity the drift induction
cannot be localized to the absorbing set; without terminal domination
the coupling inequality has nothing to consume; without absorption the
final pair law is not a coupling of the final component laws (mass that
never both-absorbs escapes the accounting). A coupling that cannot be
completed refuses --- the coupling is an input premise, so when none
exists no certificate can be formed: the rule never lies by emptiness.

\cataloguenote{The judgment-primary form ($\mathsf{QDist}[d] \le \varepsilon$, one
conclusion per certificate) deliberately differs from pre/post-style
relational program logics \cite{hsu2017}: it composes at judgment level
(triangle, weaken, symmetry where flagged), and expectation-gap
readouts enter as gap-pseudometric registry entries rather than as a
second judgment form. The proven fragment is the finite-chain
certificate tier, and its conclusion now reads in general vocabulary:
the certified bound dominates the Jordan-decomposition total-variation
distance of the final readout laws as measures (see the width note);
the mixture-polynomial tier and
the $\mathbb{W}_1$ registry entry are designed, not yet realized. A
pilot study ran four exact-rational certificates against an independent
checker with a mutation battery; on correlated faulting the relational
route certified a total-variation bound roughly three times tighter
than the best single-chain route --- the reason the tier exists.
}
\ruleblock{Lean statement}
\declhead{BernsteinLogicAux.TvDist}{exists\_hahn}{\badge{BadgeGreen}{statement bound}}
\begin{leansrc}
/-- **The Hahn-set bridge**: for finite measures of equal total mass there
is a measurable set `i` on which the difference attains `tvDist`, and which
dominates the difference on every measurable set. This is the single lemma
through which all Jordan-decomposition bookkeeping flows; everything
downstream (metric axioms, coupling inequality, finite-tier seam) consumes
only its statement. -/
theorem exists_hahn (μ ν : Measure X) [IsFiniteMeasure μ] [IsFiniteMeasure ν]
    (hmass : μ.real univ = ν.real univ) :
    ∃ i : Set X, MeasurableSet i ∧ μ.real i - ν.real i = tvDist μ ν ∧
      ∀ A : Set X, MeasurableSet A → μ.real A - ν.real A ≤ μ.real i - ν.real i
\end{leansrc}
\declhead{BernsteinLogicAux.TvDist}{tvDist\_le\_coupling}{\badge{BadgeGreen}{statement bound}}
\begin{leansrc}
/-- **The coupling inequality**: if `π` on `X × X` has marginals `μ` and
`ν`, then `d_TV(μ,ν) ≤ π {r | r.1 ≠ r.2}`. The off-diagonal is evaluated by
the *outer* measure, so no measurability of the diagonal (and no separation
hypothesis on `X`) is needed; the marginal conditions themselves force the
total masses to agree. -/
theorem tvDist_le_coupling {μ ν : Measure X} [IsFiniteMeasure μ] [IsFiniteMeasure ν]
    (π : Measure (X × X)) (hπ₁ : π.map Prod.fst = μ) (hπ₂ : π.map Prod.snd = ν) :
    tvDist μ ν ≤ (π {r : X × X | r.1 ≠ r.2}).toReal
\end{leansrc}
\declhead{BernsteinLogicAux.TvDist}{tvDist\_ofFn}{\badge{BadgeGreen}{statement bound}}
\begin{leansrc}
/-- **The finite-tier compatibility lemma**: on a finite discrete space the
general TV distance of the induced measures is the checker's own
`(1/2)·L¹` total variation. Requires only nonnegativity and equal totals
(sub-probability pairs of equal mass included). -/
theorem tvDist_ofFn {p q : L → ℝ} (hp0 : ∀ l, 0 ≤ p l) (hq0 : ∀ l, 0 ≤ q l)
    (htot : ∑ l, p l = ∑ l, q l) :
    tvDist (ofFn p) (ofFn q) = QRelChain.tvDist p q
\end{leansrc}
\declhead{BernsteinLogicAux.TvDist}{sound\_QRelTV\_measure}{\badge{BadgeGreen}{statement bound}}
\begin{leansrc}
/-- **`sound_QRelTV`, conclusion re-scoped to the general object**: under
exactly the hypotheses of `QRelChain.sound_QRelTV` — the checker's
obligations C1–C9, nothing added — the certified bound dominates the
*general* total-variation distance (`TV.tvDist`, the Jordan/coupling
object of this module) between the final readout laws, viewed as measures
on the readout space. `Fintype` survives only where the certificate
arithmetic lives (the chains and the readout space); the conclusion is the
unrestricted object. -/
theorem sound_QRelTV_measure {L : Type*} [Fintype L] [DecidableEq L]
    [MeasurableSpace L] [MeasurableSingletonClass L] {f : σ₁ → L} {g : σ₂ → L}
    {V : σ₁ × σ₂ → ℝ} {b : ℝ}
    (hK : QRelChain.IsMarkov K) (hμ₀ : QRelChain.IsDist μ₀)
    (hK₁ : QRelChain.IsMarkov K₁) (hK₂ : QRelChain.IsMarkov K₂)
    (hm₁ : QRelChain.PushKernel Prod.fst K K₁)
    (hm₂ : QRelChain.PushKernel Prod.snd K K₂)
    (hfreeze : ∀ p ∈ A₁ ×ˢ A₂, K p p = 1)
    (hfr₁ : ∀ a ∈ A₁, K₁ a a = 1) (hfr₂ : ∀ a ∈ A₂, K₂ a a = 1)
    (habs₁ : Filter.Tendsto
      (fun n => ∑ a ∈ A₁, QRelChain.law K₁ (QRelChain.push Prod.fst μ₀) n a)
      Filter.atTop (nhds 1))
    (habs₂ : Filter.Tendsto
      (fun n => ∑ a ∈ A₂, QRelChain.law K₂ (QRelChain.push Prod.snd μ₀) n a)
      Filter.atTop (nhds 1))
    (hV0 : ∀ p, 0 ≤ V p) (hdrift : ∀ p, ∑ r, K p r * V r ≤ V p)
    (hterm : ∀ p ∈ A₁ ×ˢ A₂, f p.1 ≠ g p.2 → 1 ≤ V p)
    (hb : ∑ p, μ₀ p * V p ≤ b) :
    tvDist
      (ofFn (QRelChain.push f (QRelChain.finalLaw K₁ (QRelChain.push Prod.fst μ₀) A₁)))
      (ofFn (QRelChain.push g (QRelChain.finalLaw K₂ (QRelChain.push Prod.snd μ₀) A₂)))
      ≤ b
\end{leansrc}
\declhead{BernsteinLogicAux.QRelChain}{sound\_QRelTV}{\badge{BadgeGreen}{statement bound}}
\begin{leansrc}
/-- **Soundness of `qrel-explicit-v0` `tv_le` certificates** (README
§"Soundness chain", assembled; the `sound_QCoupleDrift` obligation of the
doctrine sketch at the explicit finite tier, TV readout): under the checker's
obligations — C1/C3 (`IsMarkov`), C4 (`PushKernel`, both faces), the
completed-kernel freeze, C5 (`0 ≤ V`), C6 (drift), C7 (`V ≥ 1` on
disagreeing both-absorbed states), C8 (absorption, semantic form, per
component), C9 (`E_{μ₀}[V] ≤ b`) — the certified bound dominates the true
total-variation distance of the two final readout laws. -/
theorem sound_QRelTV [Fintype L] [DecidableEq L] {f : σ₁ → L} {g : σ₂ → L}
    {V : σ₁ × σ₂ → ℝ} {b : ℝ}
    (hK : IsMarkov K) (hμ₀ : IsDist μ₀) (hK₁ : IsMarkov K₁) (hK₂ : IsMarkov K₂)
    (hm₁ : PushKernel Prod.fst K K₁) (hm₂ : PushKernel Prod.snd K K₂)
    (hfreeze : ∀ p ∈ A₁ ×ˢ A₂, K p p = 1)
    (hfr₁ : ∀ a ∈ A₁, K₁ a a = 1) (hfr₂ : ∀ a ∈ A₂, K₂ a a = 1)
    (habs₁ : Tendsto (fun n => ∑ a ∈ A₁, law K₁ (push Prod.fst μ₀) n a)
      atTop (𝓝 1))
    (habs₂ : Tendsto (fun n => ∑ a ∈ A₂, law K₂ (push Prod.snd μ₀) n a)
      atTop (𝓝 1))
    (hV0 : ∀ p, 0 ≤ V p) (hdrift : ∀ p, ∑ q, K p q * V q ≤ V p)
    (hterm : ∀ p ∈ A₁ ×ˢ A₂, f p.1 ≠ g p.2 → 1 ≤ V p)
    (hb : ∑ p, μ₀ p * V p ≤ b) :
    tvDist (push f (finalLaw K₁ (push Prod.fst μ₀) A₁))
      (push g (finalLaw K₂ (push Prod.snd μ₀) A₂)) ≤ b
\end{leansrc}
\declhead{BernsteinLogicAux.QRelChain}{sound\_QRelExpGap}{\badge{BadgeGreen}{statement bound}}
\begin{leansrc}
/-- **Soundness of `qrel-explicit-v0` `expgap_le` certificates** (README
§"Soundness chain", assembled; Jensen replaces the coupling lemma): under the
same obligations with C7 in its expectation-gap form
(`|v₁ − v₂| ≤ V` at both-absorbed states), the certified bound dominates the
true gap of final expected readout values. -/
theorem sound_QRelExpGap {v₁ : σ₁ → ℝ} {v₂ : σ₂ → ℝ} {V : σ₁ × σ₂ → ℝ} {b : ℝ}
    (hK : IsMarkov K) (hμ₀ : IsDist μ₀) (hK₁ : IsMarkov K₁) (hK₂ : IsMarkov K₂)
    (hm₁ : PushKernel Prod.fst K K₁) (hm₂ : PushKernel Prod.snd K K₂)
    (hfreeze : ∀ p ∈ A₁ ×ˢ A₂, K p p = 1)
    (hfr₁ : ∀ a ∈ A₁, K₁ a a = 1) (hfr₂ : ∀ a ∈ A₂, K₂ a a = 1)
    (habs₁ : Tendsto (fun n => ∑ a ∈ A₁, law K₁ (push Prod.fst μ₀) n a)
      atTop (𝓝 1))
    (habs₂ : Tendsto (fun n => ∑ a ∈ A₂, law K₂ (push Prod.snd μ₀) n a)
      atTop (𝓝 1))
    (hV0 : ∀ p, 0 ≤ V p) (hdrift : ∀ p, ∑ q, K p q * V q ≤ V p)
    (hterm : ∀ p ∈ A₁ ×ˢ A₂, |v₁ p.1 - v₂ p.2| ≤ V p)
    (hb : ∑ p, μ₀ p * V p ≤ b) :
    |(∑ a₁ ∈ A₁, finalMass K₁ (push Prod.fst μ₀) a₁ * v₁ a₁) -
      ∑ a₂ ∈ A₂, finalMass K₂ (push Prod.snd μ₀) a₂ * v₂ a₂| ≤ b
\end{leansrc}
\ruleblock{Correspondence (Steele display $\leftrightarrow$ Lean)}
\badge{BadgeGreen}{reviewer-certified against the displayed texts}
\sourcenote{The dictionary below is a reviewer-certified statement of how the mechanized statements above interpret the displayed rule: display metavariable $\mapsto$ Lean binder, each display hypothesis $\mapsto$ the premise or side condition carrying it, the display conclusion $\mapsto$ the goal --- with any deliberate difference in width between display and mechanization named outright. Certification is recorded as a hash triple over (the displayed cell, the extracted statement texts, this dictionary\textquotesingle s own text); if any of the three changes, this block turns amber until a reviewer re-certifies. The correspondence claim itself remains editorial --- reviewer-grade, not machine-checked; its machine-checked successor types each soundness statement against the rule as a structured object, and supersedes this block per rule as it lands.}
The displayed $d \in \{\mathrm{TV}, \text{gap-pseudometrics}\}$ is two
theorems: $\mathsf{QDist}[\mathrm{TV}](\mathcal{M}_1, \mathcal{M}_2)
\le b$ $\mapsto$ \texttt{sound\_QRelTV}\textquotesingle s goal
\texttt{tvDist (push f (finalLaw K₁ …)) (push g (finalLaw K₂ …)) ≤ b}
--- the true total-variation distance of the two final \emph{readout}
laws (the readout maps \texttt{f}/\texttt{g} are the displayed
pseudometric\textquotesingle s observation structure, explicit in the
binder) --- and the gap form $\mapsto$
\texttt{sound\_QRelExpGap}\textquotesingle s absolute-difference goal.
The obligation roster C1--C9 maps row family by row family:
\begin{itemize}
\item coupling rows $\mapsto$ \texttt{hK : IsMarkov K} (entrywise
  nonneg, unit sums) + \texttt{hm₁}/\texttt{hm₂} (\texttt{PushKernel
  Prod.fst/snd K Kᵢ} --- marginals equal the completed component
  kernels, both faces).
\item freeze rows $\mapsto$ \texttt{hfreeze} (both-absorbed diagonal
  $= 1$) + \texttt{hfr₁}/\texttt{hfr₂} (component freezes).
\item generator row $\mapsto$ \texttt{hdrift : ∀ p, ∑ q, K p q * V q ≤ V
  p}, every product state, displayed vocabulary carried exactly;
  $V \ge 0$ $\mapsto$ \texttt{hV0}.
\item terminal domination $\mapsto$ \texttt{hterm}, per distance flavour, exactly as displayed: the
  indicator form (\texttt{f p.1 ≠ g p.2 → 1 ≤ V p}) and the gap form
  (\texttt{|v₁ p.1 − v₂ p.2| ≤ V p}), both on both-absorbed states
  only.
\item the absorption criterion carries this
  entry\textquotesingle s one representation delta: the display asks the
  per-component reach criterion (recomputed, per the flavor text); the BOUND
  statements consume absorption in \emph{semantic} form ---
  \texttt{habs₁}/\texttt{habs₂}, absorbed mass tending to one. The
  \texttt{\_of\_reach} variants consuming the checker\textquotesingle s
  own graph criterion (via \texttt{tendsto\_absorbedMass\_one})
  exist in the module and are cited by the roster witness, but are
  not in this bound set --- machine-checked and cited, not bound, in
  this edition.
\item initial marginal rows $\mapsto$ carried \emph{by construction}:
  the component initial laws in the goal ARE \texttt{push Prod.fst/snd
  μ₀} --- no separate hypothesis exists because the statement never
  admits mismatched initials; $\vdash \sum_p \mu_0(p) V(p) \le b$
  $\mapsto$ \texttt{hb}.
\end{itemize}
Width note: the \emph{conclusion} is the general object ---
\texttt{TV.sound\_QRelTV\_measure} (module
\texttt{BernsteinLogicAux.TvDist}) restates
\texttt{sound\_QRelTV}\textquotesingle s bound against the
Jordan-decomposition total-variation distance of the final readout
laws viewed as measures, with zero new hypotheses, through the general
coupling inequality (\texttt{tvDist\_le\_coupling}: any coupling,
off-diagonal weighed by outer measure, no diagonal-measurability
assumption) and the finite-tier seam \texttt{tvDist\_ofFn}
(\texttt{QRelChain.tvDist} \emph{is} the general distance restricted);
TV is packaged as a metric on probability measures (\texttt{TVProb}).
The conditional mark therefore re-scopes to the certificate-arithmetic
tier: \texttt{Fintype} chains with exact real arithmetic standing in
for exact $\Q$ --- a fragment fence, by design. The countable tier now carries kernel-level twins
(\texttt{BernsteinLogicAux.CountableTier}: the drift-induction reach
bound, the optional-stopping reach lower bound, and geometric
absorption, over \texttt{Countable} state spaces with row hypotheses as
\texttt{tsum} series, summability supplied by the unit-row-sum and
boundedness rows), machine-checked and cited, not bound, like the
\texttt{\_of\_reach} variants --- the $\mathbb{W}_1$
registry entry and the mixture-polynomial tier remain designed, not
realized, exactly as the flavor text says;
coupling emptiness has no Lean face because the theorems consume a
supplied coupling --- refusal-shaped completeness cost lives
producer-side.

\ruleblock{Lean verifier receipts}
\sourcenote{module \texttt{BernsteinLogicAux.TvDist}; bound build \texttt{54f7e48568aa\ldots} (bound from receipt 1c592f5cc199fad5, 2026-07-24T09:28:46Z). A receipt records that named checking layers accepted a build: (module, compiled-artifact SHA-256, environment SHA-256, checker layers), minted when the checking gate ran; the store receipts attesting the bound build are listed with the layers each run credited. Nothing here asks the reader to trust the store: every layer is an independently runnable program over the quoted sources --- rebuild the module and re-run the checkers to reproduce the verdict. Receipts attesting superseded builds credit no layer and are omitted from this rendering (the count of omissions, if any, is noted below).}
\begin{receipttable}
\receiptrow{a908c952d3907ab4}{2026-07-22T16:31:55Z}{build+decorrelated}{\badge{BadgeGreen}{current}}
\receiptrow{4f8b7d7adcc8f5d2}{2026-07-22T18:27:35Z}{build+decorrelated}{\badge{BadgeGreen}{current}}
\receiptrow{1c592f5cc199fad5}{2026-07-24T09:28:46Z}{build+decorrelated}{\badge{BadgeGreen}{current}}
\receiptrow{2cc8a0435e486a23}{2026-07-31T22:00:34Z}{build+decorrelated}{\badge{BadgeGreen}{current}}
\receiptrow{b78489856a2622df}{2026-07-31T22:29:28Z}{build+decorrelated}{\badge{BadgeGreen}{current}}
\receiptrow{10be3d04a0098b77}{2026-08-02T23:31:44Z}{build+decorrelated}{\badge{BadgeGreen}{current}}
\receiptrow{9bcf8f97503e6168}{2026-08-03T16:27:47Z}{build+decorrelated}{\badge{BadgeGreen}{current}}
\receiptrow{c41aeb2d063d9b00}{2026-08-03T22:46:32Z}{build+decorrelated}{\badge{BadgeGreen}{current}}
\end{receipttable}
\sourcenote{module \texttt{BernsteinLogicAux.QRelChain}; bound build \texttt{64ae4b299db4\ldots} (bound from receipt efcd1abd35a1b7f0, 2026-07-19T18:24:37Z). A receipt records that named checking layers accepted a build: (module, compiled-artifact SHA-256, environment SHA-256, checker layers), minted when the checking gate ran; the store receipts attesting the bound build are listed with the layers each run credited. Nothing here asks the reader to trust the store: every layer is an independently runnable program over the quoted sources --- rebuild the module and re-run the checkers to reproduce the verdict. Receipts attesting superseded builds credit no layer and are omitted from this rendering (the count of omissions, if any, is noted below).}
\begin{receipttable}
\receiptrow{54e0e5134c7a203f}{2026-07-19T01:11:59Z}{build+decorrelated}{\badge{BadgeGreen}{current}}
\receiptrow{1ac5f362df2230fb}{2026-07-19T01:58:14Z}{build+decorrelated}{\badge{BadgeGreen}{current}}
\receiptrow{efcd1abd35a1b7f0}{2026-07-19T18:24:37Z}{build+decorrelated}{\badge{BadgeGreen}{current}}
\receiptrow{5a18ef8a1d5d481d}{2026-07-20T01:21:49Z}{build+decorrelated}{\badge{BadgeGreen}{current}}
\receiptrow{1119d63191cbca38}{2026-07-20T20:46:42Z}{build+decorrelated}{\badge{BadgeGreen}{current}}
\receiptrow{53ac120ea74296cc}{2026-07-20T20:58:15Z}{build+replay+hardened+decorrelated}{\badge{BadgeGreen}{current}}
\receiptrow{1189ecab6f8e092f}{2026-07-21T00:03:39Z}{build+decorrelated}{\badge{BadgeGreen}{current}}
\receiptrow{4a715ffa478d6acd}{2026-07-21T05:38:59Z}{build+replay+hardened+decorrelated}{\badge{BadgeGreen}{current}}
\receiptrow{ab2a676a4de29f86}{2026-07-22T18:27:12Z}{build+decorrelated}{\badge{BadgeGreen}{current}}
\receiptrow{a1b101422ff8770b}{2026-07-31T22:00:09Z}{build+decorrelated}{\badge{BadgeGreen}{current}}
\receiptrow{6aed590f9ba0b241}{2026-07-31T22:29:03Z}{build+decorrelated}{\badge{BadgeGreen}{current}}
\receiptrow{b83a01eca4bc0784}{2026-08-02T23:31:14Z}{build+decorrelated}{\badge{BadgeGreen}{current}}
\receiptrow{2bbbc458342cfb9e}{2026-08-03T16:27:17Z}{build+decorrelated}{\badge{BadgeGreen}{current}}
\receiptrow{04d2eae82194b82e}{2026-08-03T22:46:01Z}{build+decorrelated}{\badge{BadgeGreen}{current}}
\end{receipttable}
\layerline{\layeritem{build}{lake build (elaboration)}{\badge{BadgeGreen}{current}}\layeritem{replay}{leanchecker (official kernel replay)}{\badge{BadgeGreen}{current}}\layeritem{hardened}{safe\_verify (hardened screen)}{\badge{BadgeGreen}{current}}\layeritem{decorrelated}{lean4export + nanoda (decorrelated kernel)}{\badge{BadgeGreen}{current}}}
\sourcenote{\badge{BadgeGray}{comparator: none} no fixture in the three-implementation comparator harness (compiled / kernel-interpreted / reference-interpreted) exercises this rule.}
\ruleblock{HOL4 statement}
\absentbox{NO HOL STATEMENT YET}{The HOL4 development checks the base calculus\textquotesingle s wire
format (theory \texttt{bern\_catalogue}, the first-order check-function
cone of the core catalogue). \rulename{QCoupleDrift} concludes a
distance judgment of the extension theory: its rule sort is out of
scope of the current baked theory, and the soundness spine lives in the
auxiliary Lean package. Nothing exists yet for a HOL statement to be
about; the box becomes writable when a revision carries the relational
certificate rows and its checker chain is ported.
}
\ruleblock{HOL verifier receipts}
\absentbox{NONE}{no HOL verification receipt exists for this rule.}
\ruleblock{Proof}

\sourcenote{Genre: core rule (tier 1). Scope, stated loudly: the
finite-chain tier --- finite state spaces, distribution-evolution
semantics, exact-rational certificate data instantiated over $\R$ by
cast. The mixture-polynomial tier and the $\mathbb{W}_1$ conclusion
family are designed, not proven, and nothing below claims them. Every
hypothesis of the rule is used and remarked on.}

\textbf{Setting and claim.}\enspace Two finite chains are compared:
state spaces $\sigma_1, \sigma_2$, row-stochastic kernels $K_1, K_2$,
absorbing target sets $A_i$ (each $a \in A_i$ frozen: $K_i(a, a) = 1$),
and readouts $f : \sigma_1 \to L$, $g : \sigma_2 \to L$ (resp.\ values
$v_1, v_2 : \sigma_i \to \R$). The certificate supplies a
\emph{coupling} kernel $K$ on $\sigma_1 \times \sigma_2$ with initial
law $\mu_0$, and a function $V \geq 0$ with bound $b$. The chain's law
evolves by $\mu_{n+1}(t) = \sum_s \mu_n(s) K(s, t)$; the \emph{final
law} of an absorbing state is its monotone limit mass
$\mathrm{fin}(a) = \sup_n \mu_n(a)$. Hypotheses (the checker's
obligation list, in the order the proof consumes them): the coupling
rows make $K$ a kernel whose row marginals are exactly $K_1, K_2$
(\emph{marginal identities}); both-absorbed states
$A_1 \times A_2$ are frozen; $V$ is a supermartingale for $K$
($\sum_q K(p, q) V(q) \leq V(p)$, every $p$); $V \geq 0$; terminal
domination ($V(p) \geq 1$ on both-absorbed $p$ with
$f(p_1) \neq g(p_2)$, resp.\ $V(p) \geq |v_1(p_1) - v_2(p_2)|$); each
component chain absorbs almost surely (certified by a finite-graph
reach criterion); and $\sum_p \mu_0(p) V(p) \leq b$. The claims
(\texttt{sound\_QRelTV}, \texttt{sound\_QRelExpGap}): the
total-variation distance of the two final readout laws, resp.\ the gap
of final expected values, is at most $b$.

\begin{blemma}[marginal-process lemma]\label{lem:qcd-marginal}
If every row of $K$ pushes along $\pi_i$ to the corresponding row of
$K_i$, then at every time $n$ the pushforward of the coupled law is the
component law: $\pi_{i\#}\mu_n = \mathrm{law}_i(n)$, where
$\mathrm{law}_i$ evolves $\pi_{i\#}\mu_0$ under $K_i$.
\end{blemma}
\begin{proof}
Induction on $n$; the base is definitional. For the step it suffices
that the row identities intertwine the two Markov operators:
\[
(\pi_{\#}\,\mathrm{step}_K\,\mu)(c)
= \sum_{\pi(t) = c} \sum_{s} \mu(s)\, K(s, t)
= \sum_{s} \mu(s) \sum_{\pi(t) = c} K(s, t)
= \sum_{s} \mu(s)\, K_i(\pi(s), c)
= (\mathrm{step}_{K_i}\, \pi_{\#}\mu)(c),
\]
where the third equality is the pushed-row identity at $s$ and the last
regroups the outer sum over the fibres of $\pi$. This is the coupling's
first duty: the product chain really \emph{is} a joint presentation of
the two systems being compared.
\end{proof}

\begin{blemma}[supermartingale induction]\label{lem:qcd-drift}
$\sum_p \mu_n(p)\, V(p) \leq \sum_p \mu_0(p)\, V(p)$ for all $n$.
\end{blemma}
\begin{proof}
One step: $\sum_t (\mathrm{step}\,\mu)(t) V(t) = \sum_s \mu(s)
\sum_t K(s, t) V(t) \leq \sum_s \mu(s) V(s)$, using the generator row at
each $s$ and $\mu(s) \geq 0$. Iterate.
\end{proof}

\begin{blemma}[the finite optional-stopping bound]\label{lem:qcd-stop}
Let $A \subseteq$ states be frozen, and let $0 \leq \gamma \leq V$ on
$A$. Then $\sum_{a \in A} \mathrm{fin}(a)\, \gamma(a) \leq b$.
\end{blemma}
\begin{proof}
At every finite time,
\[
\sum_{a \in A} \mu_n(a)\, \gamma(a)
\;\leq\; \sum_{a \in A} \mu_n(a)\, V(a)
\;\leq\; \sum_{\text{all } s} \mu_n(s)\, V(s)
\;\leq\; b,
\]
the middle step because $V \geq 0$ lets the sum extend off $A$, the
last by Lemma~\ref{lem:qcd-drift} and the initial bound. Frozen states
have monotone mass ($\mu_{n+1}(a) \geq \mu_n(a) K(a,a) = \mu_n(a)$), so
each term converges upward to $\mathrm{fin}(a)\gamma(a)$; the limit of
the finite sum obeys the uniform bound. This is optional stopping with
monotone convergence of finite sums doing Fatou's work --- no measure
theory enters at this tier.
\end{proof}

\begin{blemma}[coupling inequality]\label{lem:qcd-couple}
For any nonnegative weight $\gamma$ on $\sigma_1 \times \sigma_2$,
\[
\mathrm{TV}\bigl((f \circ \pi_1)_{\#}\gamma,\ (g \circ \pi_2)_{\#}\gamma\bigr)
\;\leq\; \gamma\{\, p : f(p_1) \neq g(p_2) \,\}.
\]
\end{blemma}
\begin{proof}
Fix a label $l$. Splitting each pushed mass over the agreement region
$\{f = g\}$ and its complement, the agreement contributions to the two
sides are \emph{equal} (there, $f(p_1) = l \iff g(p_2) = l$), so
\[
\bigl| (f\pi_1)_{\#}\gamma(l) - (g\pi_2)_{\#}\gamma(l) \bigr|
\;\leq\; \gamma\{f \neq g,\ f(p_1) = l\} + \gamma\{f \neq g,\ g(p_2) = l\}.
\]
Summing over $l$ counts the disagreement mass twice, and
$\mathrm{TV} = \tfrac12 \sum_l |\cdot|$ halves it. This is the
classical coupling lemma \cite[\S I.2]{lindvall1992}, in Doeblin's
tradition \cite{doeblin1938}; stated for arbitrary nonnegative
$\gamma$ because normalization is never consumed --- Step 2 of the
theorem applies it to the final pair law, whose unit mass
Lemma~\ref{lem:qcd-final} supplies separately.
\end{proof}

\begin{blemma}[final marginal identities]\label{lem:qcd-final}
Suppose additionally each component chain absorbs almost surely
(its absorbed mass tends to $1$). Then for $a_1 \in A_1$:
$\sum_{a_2 \in A_2} \mathrm{fin}_K(a_1, a_2) =
\mathrm{fin}_{K_1}(a_1)$, and symmetrically. Consequently the final
pair law on $A_1 \times A_2$ is a genuine \emph{coupling of the two
final component laws}, of total mass $1$.
\end{blemma}
\begin{proof}
Let $S_n = \sum_{a_2 \in A_2} \mu_n(a_1, a_2)$ and
$T_n = \sum_{s_2} \mu_n(a_1, s_2)$; by Lemma~\ref{lem:qcd-marginal},
$T_n$ is the component-1 law at $a_1$. Then $S_n \leq T_n$, and
\[
T_n - S_n \;=\; \textstyle\sum_{s_2 \notin A_2} \mu_n(a_1, s_2)
\;\leq\; \textstyle\sum_{s_2 \notin A_2} (\pi_{2\#}\mu_n)(s_2)
\;=\; 1 - (\text{component-2 absorbed mass at } n) \;\longrightarrow\; 0 :
\]
the \emph{one-absorbed-loner} mass --- pairs whose first coordinate has
frozen while the second still wanders --- is dominated by the second
chain's unabsorbed mass, which its absorption hypothesis kills.
Squeeze; both-absorbed pairs are frozen, so $S_n$ also converges to the
final pair mass. Total mass $1$ follows by summing over $A_1$ and
applying component-1 absorption.
\end{proof}

\begin{bthm}[soundness, total-variation form]\label{thm:qcd-tv}
Under all hypotheses,
$\mathrm{TV}\bigl(f_{\#}\mathrm{fin}_{K_1},\ g_{\#}\mathrm{fin}_{K_2}\bigr)
\leq b$.
\end{bthm}
\begin{proof}
\emph{Step 1.} By Lemma~\ref{lem:qcd-final} the final pair law $\gamma$
on $A_1 \times A_2$ has marginals $\mathrm{fin}_{K_1}$,
$\mathrm{fin}_{K_2}$, so the two final readout laws are the two
pushforwards of the single weight $\gamma$.
\emph{Step 2.} Lemma~\ref{lem:qcd-couple} bounds their total-variation
distance by $\gamma$'s disagreement mass.
\emph{Step 3.} Terminal domination says the disagreement indicator is
$\leq V$ on the frozen set $A_1 \times A_2$ (and $\geq 0$ everywhere),
so Lemma~\ref{lem:qcd-stop} with $\gamma = \mathbf{1}[f \neq g]$ bounds
the disagreement mass by $b$.
\end{proof}

\begin{bthm}[soundness, expectation-gap form]\label{thm:qcd-gap}
With terminal domination in the form
$V(p) \geq |v_1(p_1) - v_2(p_2)|$ on $A_1 \times A_2$,
\[
\Bigl| \textstyle\sum_{a_1} \mathrm{fin}_{K_1}(a_1)\, v_1(a_1)
- \sum_{a_2} \mathrm{fin}_{K_2}(a_2)\, v_2(a_2) \Bigr| \;\leq\; b .
\]
\end{bthm}
\begin{proof}
Rewrite both final expectations over the final pair law via
Lemma~\ref{lem:qcd-final}; the difference becomes
$\sum_p \mathrm{fin}_K(p)\,(v_1(p_1) - v_2(p_2))$ over
$A_1 \times A_2$. The triangle inequality moves the absolute value
inside, and Lemma~\ref{lem:qcd-stop} with
$\gamma = |v_1 - v_2|$ closes. The coupling lemma is not consumed here
--- domination of the value gap replaces domination of the disagreement
indicator; two conclusion families, one engine.
\end{proof}

\begin{bremark}[the absorption criterion is discharged, not assumed]
\label{rem:qcd-reach}
The checker certifies absorption per component as a finite-graph fact:
from every state, the frozen set is reachable through
positive-probability edges. Finiteness uniformizes this into one
horizon $N$ and one margin $\varepsilon > 0$ valid for all states; each
$N$-step block then absorbs at least an $\varepsilon$-fraction of the
still-loose mass (frozen mass never returns), so the unabsorbed mass
decays geometrically and the absorbed mass tends to $1$
(\texttt{tendsto\_absorbedMass\_one}). The \texttt{\_of\_reach}
theorem forms consume the graph criterion directly.
\end{bremark}

\begin{bremark}[why each hypothesis is needed]\label{rem:qcd-omit}
\emph{Marginal identities}: drop them and $\gamma$ in
Theorem~\ref{thm:qcd-tv} no longer has the two systems' final laws as
marginals --- the bound then constrains an unrelated pair (couple
$\mathcal{M}_1$ with a copy of \emph{itself} instead of
$\mathcal{M}_2$: every absorbed pair agrees, so $V = 0$ on the
reachable diagonal --- with $V = 1$ on the unreachable disagreeing
absorbed pairs --- passes every other row at $b = 0$, whatever
$\mathcal{M}_2$ is). \emph{$V \geq 0$}: the extension step in
Lemma~\ref{lem:qcd-stop} ($\sum_A \leq \sum_{\text{all}}$) fails, and
negative values can subsidize the generator row: from an initial state of
value $0$, branch half-half into a frozen \emph{disagreeing} state of
value $1$ and a frozen \emph{agreeing} state of value $-1$ (domination
constrains only disagreeing states). The generator row holds at every state
($\tfrac12 \cdot 1 + \tfrac12 \cdot (-1) = 0 \leq 0$), the initial
bound passes at $b = 0$ --- and the disagreement mass is $\tfrac12$. \emph{Terminal domination}: with
$V \equiv 0$ and disagreeing finals, generator and initial rows pass at
$b = 0$ while the true distance is $1$. \emph{The freeze}: monotone
convergence of frozen mass (Lemmas \ref{lem:qcd-stop},
\ref{lem:qcd-final}) is what makes ``final law'' well-defined;
unfrozen, the two-state swap chain puts mass $1$ on the target at even
times and $0$ at odd --- $\sup_n \mu_n(a)$ is no longer a limit, and
certified mass drains back out.
\emph{Absorption}: without it the one-absorbed-loner mass survives ---
take $K_2$ the identity off $A_2$: the pair never both-absorbs, the
final pair law is $0$, and Lemma~\ref{lem:qcd-couple} holds vacuously
about sub-laws that no longer represent the systems; the marginal
identity of Lemma~\ref{lem:qcd-final} is exactly what fails --- and
the bound itself fails with it: component 1 absorbing at once,
component 2 frozen nowhere it reaches, $V = 0$ on the reachable ray
and $1$ on the (unreachable) disagreeing both-absorbed pairs passes
every other row at $b = 0$, while the final readout laws sit at
total-variation distance $\tfrac12$.
\emph{Initial bound}: the conclusion's $b$ is
$V$'s initial expectation; without it nothing quantitative survives.
A coupling that cannot satisfy the marginal rows simply does not
exist as a certificate --- the coupling is an input premise, so its
emptiness refuses (a completeness cost), never lies.
\end{bremark}

\begin{bremark}[relation to the literature]
The engine is the coupling method for Markov chains
\cite{doeblin1938,lindvall1992}: total variation bounded by
disagreement probability of a coupling, here in distribution-evolution
form over finite spaces \cite[Ch.~4--5]{levinperes2017}, with a
relational supermartingale supplying the quantitative disagreement
bound --- the certificate-side sibling of coupling-based relational
program logics \cite{hsu2017}, transposed from pre/post judgments to a
single distance judgment that composes by metric axioms (triangle,
weakening) at judgment level. The expectation-gap form replaces the
coupling lemma by direct domination --- the gap-pseudometric readout.
On correlated perturbations the relational route is strictly sharper
than composing two one-sided certificates, which is the rule's reason
to exist: the coupling sees the correlation, unary certificates
cannot.
\end{bremark}

\begin{anchors}
Lemma \ref{lem:qcd-marginal} & \texttt{PushKernel}; \texttt{push\_step}; \texttt{push\_law}\\
Lemma \ref{lem:qcd-drift} & \texttt{expect\_step\_le}; \texttt{expect\_law\_le}\\
Lemma \ref{lem:qcd-stop} & \texttt{law\_mono\_absorbing}; \texttt{finalMass}; \texttt{finalExpect\_le}\\
Lemma \ref{lem:qcd-couple} & \texttt{tvDist}; \texttt{tvDist\_push\_le\_disagree}\\
Lemma \ref{lem:qcd-final} & \texttt{sum\_finalMass\_snd\_eq}; \texttt{sum\_finalMass\_fst\_eq}; \texttt{sum\_finalMass\_prod\_eq\_one}; \texttt{push\_finalLaw\_fst}/\texttt{\_snd}\\
Theorem \ref{thm:qcd-tv} & \texttt{sound\_QRelTV} (graph-criterion form \texttt{sound\_QRelTV\_of\_reach})\\
Theorem \ref{thm:qcd-gap} & \texttt{sound\_QRelExpGap} (graph-criterion form \texttt{sound\_QRelExpGap\_of\_reach})\\
Remark \ref{rem:qcd-reach} & \texttt{exists\_uniform\_reach}; \texttt{absorbed\_decay}; \texttt{tendsto\_absorbedMass\_one}\\
\end{anchors}


\rulesection{92}{QFixBall}{core rule}
\label{blcatrule-QFixBall}
\ruleblock{The rule (Steele metanotation)}
\sourcenote{as realized in the mechanized development; source row: \texttt{\texttt{QTSL-rules.toml}} --- \emph{QFixBall}. The metanotation statement is typeset by hand for notational quality; its mechanized ground truth is the Lean statement below, extracted from the sources at every build of this book and hash-bound. The rule \emph{roster} is machine-checked (every mechanized rule tag round-trips to exactly one catalogue row); the statement \emph{text} is editorial, checked by eye against the extracted Lean beside it.}
\begin{mathpar}
\inferrule*{
  d : \mathsf{QDistRegistry}(X_d); \qquad (X_d, d) \ \text{complete} \\
  \vdash \rho < 1 \ \text{two-point on } \overline{B}(\hat\mu, r) \ \text{(contraction row)} \\
  \vdash d(T\hat\mu, \hat\mu) \le \delta \ \text{(residual row)} \\
  \vdash \delta + \rho\, r \le r \ \text{(radius row, division-free)}
}{
  \vdash \mathsf{QDist}[d](\hat\mu,\ \mathrm{fix}\,T) \le r
}
\end{mathpar}

\ruleflavor{every hypothesis anchored at the known snapshot $\hat\mu$, never at the unknown fixed point: the contraction row discharges by operator-norm/psdCert rows at the affine tier, the residual row is one concrete sweep at the known point, and completeness is mechanized at the affine/Euclidean instantiation ($\mathbb{W}_2$ completeness by instance on $(\mathcal{P}_2(X), W_2)$). The conclusion\textquotesingle s definite description needs no separate license premise: this rule is the only derivation that makes $\mathrm{fix}\,T$ speakable, and it discharges the description\textquotesingle s well-formedness itself --- exactly one fixed point exists in $\overline{B}(\hat\mu, r)$, manufactured by the contraction (the $\exists!$ of the soundness statement; no existence stratum consumed). The registry typing has no baked sort yet (the promoted $\mathsf{QDist}$ kind awaits its row species); completeness is the instance hypothesis of the Lean statement. Refusals: $\rho \ge 1$, and registry entries the subject sort does not carry.}
\par\medskip
\rulename{QFixBall} is the contraction-witnessed discharge of the
distance judgment: it certifies that a computed snapshot $\hat\mu$ lies
within distance $r$ of the fixed point of an update map $T$ --- and, in
the same breath, that \textquotedblleft the fixed point of
$T$\textquotedblright{} is a well-defined description at all (exactly
one fixed point exists in the certified ball). A producer reaches for
it wherever an iterative solver has been run to approximate a
fixed-point object --- a stationary law, a variational posterior, a
self-consistent field --- and the deliverable is not the iterate but a
certified \emph{ball} around it. Its three rows realize the anchoring
inversion: every hypothesis is stated at the \emph{known} snapshot,
never at the unknown fixed point --- the contraction row bounds
$T$\textquotesingle s two-point ratio on the closed ball around
$\hat\mu$, the residual row measures one concrete application
$d(T\hat\mu, \hat\mu) \le \delta$, and the radius row
$\delta + \rho\, r \le r$ is division-free arithmetic.

What the checker recomputes: the contraction bound $\rho < 1$ (at the
affine tier, an operator-norm row discharged by positive-semidefinite
certificate machinery on $\rho^2 I - A^{\top}\!A$-shaped data), the residual
evaluation (one sweep of the map at the snapshot, exact rational
arithmetic), and the radius inequality. Everything else --- ball
self-mapping, existence, uniqueness, the anchored bound
$d(\hat\mu, \mathrm{fix}\,T) \le \delta/(1-\rho) \le r$ --- is
manufactured by the mathematics from those three rows; no existence
stratum is consumed, so the judgment is existential without vacuity.

What would break: with $\rho \ge 1$ the fixed point need not be unique
or exist and the description license fails --- the rule refuses (the
identity map has residual $0$ at every point and fixes all of them).
Without the radius row the ball need not be self-mapped and can miss
the fixed point entirely, even with a genuine contraction. Without the
residual row nothing ties the ball to the dynamics. The completeness of
the subject space is a genuine hypothesis, carried honestly --- and both
named instantiations now discharge it by instance: the affine/Euclidean
case directly, and the Wasserstein case on $(\mathcal{P}_2(X), W_2)$
over a complete separable base (the instantiation block below) --- never
an axiom smuggled in.

\cataloguenote{The rule concludes onto a definite-description subject
($\mathrm{fix}\,T$) and therefore carries the description
well-formedness discharge inside the rule --- the one new typing
obligation of the distance-judgment grammar. An independent
exact-arithmetic checker replays a pilot certificate\textquotesingle s
eleven obligations; its mutation battery refuses each mutilated
certificate at a named obligation. Producer strategy (proposing
$\rho$ from iterate data, pinning the affine reduction) is deliberately
outside the kernel.
}
\ruleblock{Lean statement}
\declhead{BernsteinLogicAux.QFixBall}{sound\_QFixBall}{\badge{BadgeGreen}{statement bound}}
\begin{leansrc}
/-- **`sound_QFixBall`** — the rule's full conclusion: the certified
rows produce EXACTLY ONE fixed point of `f` in `closedBall x₀ r`.
Existence is manufactured by the contraction itself (the third case of
the two-direction carve: no existence stratum is consumed, and the
judgment is existential without vacuity). `[CompleteSpace α]` is the
honest boundary — see the header. -/
theorem sound_QFixBall [CompleteSpace α] (hK : K < 1)
    (hlip : LipschitzOnWith K f (closedBall x₀ r))
    (hres : dist (f x₀) x₀ ≤ δ) (hrad : δ + K * r ≤ r) :
    ∃! y, y ∈ closedBall x₀ r ∧ IsFixedPt f y
\end{leansrc}
\declhead{BernsteinLogicAux.QFixBall}{sound\_QFixBall\_affine}{\badge{BadgeGreen}{statement bound}}
\begin{leansrc}
/-- **`sound_QFixBall_affine`** — the caviwball-v0 instantiation shape:
snapshot mean vector in `EuclideanSpace ℝ (Fin M)` (completeness is an
instance — the `[CompleteSpace]` boundary discharges for free here),
affine sweep map, operator-norm + residual + radius rows ⇒ unique
fixed point in the certified ball. -/
theorem sound_QFixBall_affine {M : ℕ}
    (A : EuclideanSpace ℝ (Fin M) →L[ℝ] EuclideanSpace ℝ (Fin M))
    (c x₀ : EuclideanSpace ℝ (Fin M)) {K : ℝ≥0} {δ r : ℝ}
    (hK : K < 1) (hA : ‖A‖₊ ≤ K)
    (hres : dist (A x₀ + c) x₀ ≤ δ) (hrad : δ + K * r ≤ r) :
    ∃! y, y ∈ closedBall x₀ r ∧ IsFixedPt (fun x => A x + c) y
\end{leansrc}
\declhead{BernsteinLogicAux.QFixBall}{dist\_fixedPoint\_le}{\badge{BadgeGreen}{statement bound}}
\begin{leansrc}
/-- **The anchored a-priori bound** — for ANY fixed point in the ball
(uniqueness not consumed): `dist x₀ μ⋆ ≤ δ / (1 - K)`. Proved by one
triangle inequality through the snapshot's image; the literature's
one-point-at-the-limit anchoring never appears. -/
theorem dist_fixedPoint_le (hK : K < 1)
    (hlip : LipschitzOnWith K f (closedBall x₀ r))
    (hres : dist (f x₀) x₀ ≤ δ) (hrad : δ + K * r ≤ r)
    {y : α} (hy : y ∈ closedBall x₀ r) (hfy : IsFixedPt f y) :
    dist x₀ y ≤ δ / (1 - K)
\end{leansrc}
\declhead{BernsteinLogicAux.QFixBall}{anchored\_le\_radius}{\badge{BadgeGreen}{statement bound}}
\begin{leansrc}
/-- The radius row also caps the anchored bound by `r` itself
(`δ/(1-K) ≤ r`) — the form the checker's ball-radius obligation uses. -/
theorem anchored_le_radius (hK : K < 1)
    (hrad : δ + K * r ≤ r) : δ / (1 - K) ≤ r
\end{leansrc}
\par\medskip
\textbf{The Wasserstein instantiation, discharged by instance.}\enspace
The completeness hypothesis named above is closed on the space where
this rule\textquotesingle s stochastic clients live: the statement
below runs \texttt{sound\_QFixBall} on $(\mathcal{P}_2(X), W_2)$ ---
second-moment laws under quadratic Wasserstein distance over a
complete separable metric space --- with the \texttt{CompleteSpace}
side condition supplied by instance resolution rather than hypothesis.
No completeness assumption survives in the statement; the affine and
Wasserstein instantiations now stand on the same footing.

\declhead{BernsteinLogicAux.Wasserstein.Complete}{sound\_QFixBall\_wasserstein}{\badge{BadgeGreen}{statement bound}}
\begin{leansrc}
/-- **`sound_QFixBall`, instantiated at the Wasserstein-2 space.** The
`[CompleteSpace]` hypothesis — the honest boundary of
`BernsteinLogicAux/QFixBall.lean` (header item 1), and the side
condition that kept the rule's coverage marked conditional — is
discharged here **by instance**: for any complete separable metric `X`,
certified contraction/residual/radius rows on `(𝒫₂(X), W₂)` produce
the unique fixed point in the certified ball. -/
theorem sound_QFixBall_wasserstein {T : P2 X → P2 X} {μ₀ : P2 X} {K : ℝ≥0}
    {δ r : ℝ} (hK : K < 1)
    (hlip : LipschitzOnWith K T (closedBall μ₀ r))
    (hres : dist (T μ₀) μ₀ ≤ δ) (hrad : δ + K * r ≤ r) :
    ∃! ν, ν ∈ closedBall μ₀ r ∧ IsFixedPt T ν
\end{leansrc}
\ruleblock{Correspondence (Steele display $\leftrightarrow$ Lean)}
\badge{BadgeGreen}{reviewer-certified against the displayed texts}
\sourcenote{The dictionary below is a reviewer-certified statement of how the mechanized statements above interpret the displayed rule: display metavariable $\mapsto$ Lean binder, each display hypothesis $\mapsto$ the premise or side condition carrying it, the display conclusion $\mapsto$ the goal --- with any deliberate difference in width between display and mechanization named outright. Certification is recorded as a hash triple over (the displayed cell, the extracted statement texts, this dictionary\textquotesingle s own text); if any of the three changes, this block turns amber until a reviewer re-certifies. The correspondence claim itself remains editorial --- reviewer-grade, not machine-checked; its machine-checked successor types each soundness statement against the rule as a structured object, and supersedes this block per rule as it lands.}
The displayed conclusion splits, and so do the bound statements: the
description license --- $\mathrm{fix}\,T$ well-defined (unique) in
$\overline{B}(\hat\mu, r)$ --- $\mapsto$ \texttt{sound\_QFixBall :
∃! y, y ∈ closedBall x₀ r ∧ IsFixedPt f y} (existence manufactured by
the contraction itself; no existence stratum consumed, exactly as the
roster says); the distance clause $\mathsf{QDist}[d](\hat\mu,
\mathrm{fix}\,T) \le r$ is then ball membership, with the sharper
anchored bound $\delta/(1-\rho)$ carried by
\texttt{dist\_fixedPoint\_le} --- whose binder row shows uniqueness is
NOT consumed (any fixed point in the ball obeys it) --- and capped by
the radius row via \texttt{anchored\_le\_radius}.
\begin{itemize}
\item contraction row, two-point on the ball $\mapsto$ \texttt{hlip :
  LipschitzOnWith K f (closedBall x₀ r)} with \texttt{K : ℝ≥0},
  \texttt{hK : K < 1} --- the $\rho \ge 1$ refusal is the strict
  binder; at the affine tier the operator-norm row $\mapsto$
  \texttt{hA : ‖A‖₊ ≤ K} in \texttt{sound\_QFixBall\_affine}
  (EuclideanSpace, the caviwball-v0 pinned mean map).
\item residual row $\mapsto$ \texttt{hres : dist (f x₀) x₀ ≤ δ} ---
  anchored at the snapshot, and no hypothesis anywhere mentions the
  fixed point: the flavor text\textquotesingle s ``never at the unknown
  fixed point'' is a checkable property of the binder list.
\item radius row $\mapsto$ \texttt{hrad : δ + K * r ≤ r},
  division-free verbatim.
\item registry completeness $\mapsto$ the instance binder
  \texttt{[CompleteSpace α]} --- a named hypothesis of the general
  statement, never an axiom, and discharged by instance resolution in
  both named instantiations: directly for the affine/Euclidean case,
  and for $\mathbb{W}_2$ by the Wasserstein block
  (\texttt{sound\_QFixBall\_wasserstein}) below.
\end{itemize}
Width note: the subject space is an arbitrary metric space --- the
registry-entry structure of $d$ is erased, so the certified reach
covers every registry instantiation whose completeness discharge
arrives (both named instantiations\textquotesingle{} discharges are now
in hand), strictly wider than the displayed registry framing, safe
direction; the means-only (variance-conclusion) refusal and the
Gaussian-pinning producer strategy are kernel-external, no Lean
counterpart, correctly.

\ruleblock{Lean verifier receipts}
\sourcenote{module \texttt{BernsteinLogicAux.QFixBall}; bound build \texttt{68796bc2f9db\ldots} (bound from receipt 060c7093ff4899bf, 2026-07-21T04:51:42Z). A receipt records that named checking layers accepted a build: (module, compiled-artifact SHA-256, environment SHA-256, checker layers), minted when the checking gate ran; the store receipts attesting the bound build are listed with the layers each run credited. Nothing here asks the reader to trust the store: every layer is an independently runnable program over the quoted sources --- rebuild the module and re-run the checkers to reproduce the verdict. Receipts attesting superseded builds credit no layer and are omitted from this rendering (the count of omissions, if any, is noted below).}
\begin{receipttable}
\receiptrow{a88bac06d271526e}{2026-07-21T00:27:07Z}{build+decorrelated}{\badge{BadgeGreen}{current}}
\receiptrow{060c7093ff4899bf}{2026-07-21T04:51:42Z}{build+decorrelated}{\badge{BadgeGreen}{current}}
\receiptrow{ae5851604e101130}{2026-07-21T05:38:57Z}{build+replay+hardened+decorrelated}{\badge{BadgeGreen}{current}}
\receiptrow{85c6bac865d222ba}{2026-07-22T18:27:12Z}{build+decorrelated}{\badge{BadgeGreen}{current}}
\receiptrow{bf722ef76b576c08}{2026-07-31T22:00:08Z}{build+decorrelated}{\badge{BadgeGreen}{current}}
\receiptrow{0caae5e60f19815a}{2026-07-31T22:29:02Z}{build+decorrelated}{\badge{BadgeGreen}{current}}
\receiptrow{9780a725a44c924b}{2026-08-02T23:31:13Z}{build+decorrelated}{\badge{BadgeGreen}{current}}
\receiptrow{490335613ee08847}{2026-08-03T16:27:17Z}{build+decorrelated}{\badge{BadgeGreen}{current}}
\receiptrow{e68b74ba7bf3e8de}{2026-08-03T22:46:01Z}{build+decorrelated}{\badge{BadgeGreen}{current}}
\end{receipttable}
\sourcenote{5 further receipt(s) in the store attest superseded builds of this module and are omitted from this rendering.}
\sourcenote{module \texttt{BernsteinLogicAux.Wasserstein.Complete}; bound build \texttt{06292c010c42\ldots} (bound from receipt fc2b784b4f0c347e, 2026-07-29T08:09:21Z). A receipt records that named checking layers accepted a build: (module, compiled-artifact SHA-256, environment SHA-256, checker layers), minted when the checking gate ran; the store receipts attesting the bound build are listed with the layers each run credited. Nothing here asks the reader to trust the store: every layer is an independently runnable program over the quoted sources --- rebuild the module and re-run the checkers to reproduce the verdict. Receipts attesting superseded builds credit no layer and are omitted from this rendering (the count of omissions, if any, is noted below).}
\begin{receipttable}
\receiptrow{f660300e93359aad}{2026-07-28T23:54:07Z}{build+decorrelated}{\badge{BadgeGreen}{current}}
\receiptrow{ddcacf582a03fe61}{2026-07-29T05:43:06Z}{build+decorrelated}{\badge{BadgeGreen}{current}}
\receiptrow{fc2b784b4f0c347e}{2026-07-29T08:09:21Z}{build+replay+hardened+decorrelated}{\badge{BadgeGreen}{current}}
\receiptrow{439956f3f0160d8a}{2026-07-31T22:00:35Z}{build+decorrelated}{\badge{BadgeGreen}{current}}
\receiptrow{d52523d643b3462b}{2026-07-31T22:29:29Z}{build+decorrelated}{\badge{BadgeGreen}{current}}
\receiptrow{a5c8597e9187735d}{2026-08-02T23:31:45Z}{build+decorrelated}{\badge{BadgeGreen}{current}}
\receiptrow{9f4e8bbd7042154f}{2026-08-03T16:27:49Z}{build+decorrelated}{\badge{BadgeGreen}{current}}
\receiptrow{c0e6026aa1939c74}{2026-08-03T22:46:33Z}{build+decorrelated}{\badge{BadgeGreen}{current}}
\end{receipttable}
\sourcenote{4 further receipt(s) in the store attest superseded builds of this module and are omitted from this rendering.}
\layerline{\layeritem{build}{lake build (elaboration)}{\badge{BadgeGreen}{current}}\layeritem{replay}{leanchecker (official kernel replay)}{\badge{BadgeGreen}{current}}\layeritem{hardened}{safe\_verify (hardened screen)}{\badge{BadgeGreen}{current}}\layeritem{decorrelated}{lean4export + nanoda (decorrelated kernel)}{\badge{BadgeGreen}{current}}}
\sourcenote{\badge{BadgeGray}{comparator: none} no fixture in the three-implementation comparator harness (compiled / kernel-interpreted / reference-interpreted) exercises this rule.}
\ruleblock{HOL4 statement}
\absentbox{NO HOL STATEMENT YET}{The HOL4 development\textquotesingle s object is the base
calculus\textquotesingle s wire-format checker (theory
\texttt{bern\_catalogue} ports the first-order check-function cone of
the core catalogue). \rulename{QFixBall} lives one extension edge above
that stratum --- it concludes a distance judgment of the extension
theory, its rule sort is out of scope of the current baked theory, and
its mechanization lives in the auxiliary Lean package. There is nothing
checker-shaped for a HOL statement to be about yet; the box becomes
writable when a revision of the theory carries the certificate
objects and the corresponding checker
chain is ported.
}
\ruleblock{HOL verifier receipts}
\absentbox{NONE}{no HOL verification receipt exists for this rule.}
\ruleblock{Proof}

\sourcenote{Genre: core rule (tier 1) --- Banach's fixed-point theorem
\cite{banach1922} in certificate form. Scope, stated loudly: the
theorems are parametric in a complete metric space; the affine/Euclidean
instantiation is mechanized, and completeness of the Wasserstein space
is a \emph{named hypothesis}, not a discharged one. Every hypothesis of
the rule is used below and remarked on.}

\textbf{Setting and claim.}\enspace Fix a metric space $(X, d)$, a map
$T : X \to X$, a snapshot $\hat\mu \in X$ (the object the producer
actually computed), and rational certificate data
$\rho \in \Q \cap [0, 1)$, $\delta, r \in \Q$. The three certified rows:
\begin{align*}
&\text{(H-\textsc{anchor})}\quad d(Tx, Ty) \leq \rho\, d(x, y)
  \ \text{ for all } x, y \in \overline{B}(\hat\mu, r)
  && \text{(two-point contraction \emph{on the ball})},\\
&\text{(H-\textsc{resid})}\quad d(T\hat\mu, \hat\mu) \leq \delta
  && \text{(one concrete sweep at the snapshot)},\\
&\text{(H-\textsc{radius})}\quad \delta + \rho\, r \leq r
  && \text{(division-free; equivalent to } r \geq \delta/(1-\rho)\text{)}.
\end{align*}
The claim (\texttt{sound\_QFixBall}): if $X$ is complete, there is
\emph{exactly one} fixed point of $T$ in $\overline{B}(\hat\mu, r)$ ---
so the description $\mathrm{fix}\,T$ is licensed on the ball, and the
distance judgment $d(\hat\mu, \mathrm{fix}\,T) \leq r$ holds; moreover
the sharper anchored bound $d(\hat\mu, \mathrm{fix}\,T) \leq
\delta/(1-\rho)$ holds (\texttt{dist\_fixedPoint\_le},
\texttt{anchored\_le\_radius}). Every hypothesis is anchored at the
\emph{known} point $\hat\mu$; the unknown fixed point appears only in
conclusions --- the anchoring inversion that makes the statement
checkable at all.

\begin{blemma}[the rows are sign-coherent]\label{lem:qfb-signs}
H-\textsc{resid} forces $\delta \geq 0$; together with $\rho < 1$,
H-\textsc{radius} forces $r \geq 0$, so the ball is honest.
\end{blemma}
\begin{proof}
$0 \leq d(T\hat\mu, \hat\mu) \leq \delta$. From H-\textsc{radius},
$(1 - \rho)\, r \geq \delta \geq 0$, and $1 - \rho > 0$.
\end{proof}

\begin{blemma}[self-mapping is derived, never assumed]\label{lem:qfb-maps}
$T$ maps $\overline{B}(\hat\mu, r)$ into itself.
\end{blemma}
\begin{proof}
Let $d(x, \hat\mu) \leq r$. Both $x$ and $\hat\mu$ lie in the ball
(Lemma~\ref{lem:qfb-signs} for $\hat\mu$), so H-\textsc{anchor} applies
to the pair:
\[
d(Tx, \hat\mu) \;\leq\; d(Tx, T\hat\mu) + d(T\hat\mu, \hat\mu)
\;\leq\; \rho\, d(x, \hat\mu) + \delta
\;\leq\; \rho\, r + \delta \;\leq\; r . \qedhere
\]
\end{proof}

\begin{blemma}[in-ball uniqueness]\label{lem:qfb-unique}
If $y, z \in \overline{B}(\hat\mu, r)$ are both fixed by $T$, then
$y = z$.
\end{blemma}
\begin{proof}
$d(y, z) = d(Ty, Tz) \leq \rho\, d(y, z)$ with $\rho < 1$ forces
$d(y, z) \leq 0$, hence $= 0$. Note the proof consumes only the
\emph{two-point} row, on the ball --- no iteration, no completeness.
\end{proof}

\begin{blemma}[the anchored a-priori bound]\label{lem:qfb-anchor}
Any fixed point $y \in \overline{B}(\hat\mu, r)$ satisfies
$d(\hat\mu, y) \leq \delta/(1 - \rho) \leq r$.
\end{blemma}
\begin{proof}
One triangle inequality through the snapshot's image:
\[
d(\hat\mu, y) \;\leq\; d(\hat\mu, T\hat\mu) + d(T\hat\mu, Ty)
\;\leq\; \delta + \rho\, d(\hat\mu, y),
\]
so $(1 - \rho)\, d(\hat\mu, y) \leq \delta$. The cap by $r$ is
H-\textsc{radius} rearranged: $\delta \leq (1 - \rho)\, r$
(\texttt{anchored\_le\_radius}). Uniqueness is not consumed --- the
bound holds for \emph{any} fixed point in the ball, which is what lets
the rule state it before existence is settled.
\end{proof}

\begin{bthm}[soundness of \rulename{QFixBall}]\label{thm:qfb-sound}
Let $X$ be complete and let the three rows hold. Then $T$ has exactly
one fixed point $\mu^{*}$ in $\overline{B}(\hat\mu, r)$, and
$d(\hat\mu, \mu^{*}) \leq \delta/(1-\rho) \leq r$.
\end{bthm}
\begin{proof}
\emph{Existence.} By Lemma~\ref{lem:qfb-maps} the closed ball is
$T$-invariant; it is a closed subset of a complete space, hence
complete, and $T$ restricted to it is a $\rho$-contraction
(H-\textsc{anchor}). Banach's theorem \cite{banach1922} applies to the
restriction: the Picard iterates $T^n \hat\mu$ form a Cauchy sequence
($d(T^{n+1}\hat\mu, T^n\hat\mu) \leq \rho^n \delta$, summable) whose
limit $\mu^{*}$ lies in the ball and is fixed. Existence is thus
\emph{manufactured} by the contraction --- the rule consumes no
separate existence certificate, so the concluded description is
existential without vacuity.

\emph{Uniqueness} is Lemma~\ref{lem:qfb-unique}; the \emph{bound} is
Lemma~\ref{lem:qfb-anchor}. Together: $\mathrm{fix}\,T$ is well-defined
on the ball (the description license the judgment carries) and the
distance judgment holds at radius $r$, with $\delta/(1-\rho)$ as the
sharper certified value.
\end{proof}

\begin{bprop}[the affine instantiation]\label{prop:qfb-affine}
On a finite-dimensional Euclidean space, let $T x = A x + c$ with $A$
linear and an operator-norm row $\lVert A \rVert \leq \rho < 1$. Then
H-\textsc{anchor} holds globally (a fortiori on the ball), completeness
holds by instance, and Theorem~\ref{thm:qfb-sound} applies verbatim
(\texttt{sound\_QFixBall\_affine}).
\end{bprop}
\begin{proof}
$d(Tx, Ty) = \lVert A(x - y) \rVert \leq \lVert A \rVert\,
\lVert x - y \rVert \leq \rho\, d(x, y)$. The operator-norm row is the
affine tier's discharge of H-\textsc{anchor}: it is certified by
positive-semidefiniteness of $\rho^2 I - A^{\top} A$, which is
\rulename{PSDCert}-family machinery on rational data.
\end{proof}

\begin{bremark}[why each hypothesis is needed]\label{rem:qfb-omit}
\emph{Drop $\rho < 1$}: the identity map has residual $0$ at every
snapshot and every radius row holds with $\rho = 1$, yet every point is
fixed --- the description license (uniqueness) fails, so the rule
\emph{refuses} $\rho \geq 1$ rather than concluding.
\emph{Drop H-\textsc{radius}}: on $\R$, take $Tx = \tfrac{1}{2}x + 2$,
$\hat\mu = 0$, so $\rho = \tfrac12$, $\delta = 2$; the unique fixed
point is $4$. With $r = 1$ the ball $[-1, 1]$ contains no fixed point
--- and indeed $\delta + \rho r = \tfrac52 > 1$: the radius row fails,
correctly. \emph{Drop H-\textsc{resid}}: nothing ties the ball to the
dynamics --- any $r$ would pass on any map with small $\rho$, concluding
a distance bound to a fixed point that can sit arbitrarily far away.
\emph{Weaken completeness}: on $X = \Q \cap [1, 2]$ (incomplete),
$Tx = \tfrac{x}{2} + \tfrac{1}{x}$ maps $X$ into itself with
$|T'| \le \tfrac12$ and contracts toward the irrational $\sqrt{2}$ ---
iterates Cauchy, no fixed point in $X$, and the certificate rows even
pass ($\hat\mu = \tfrac32$, $r = \tfrac12$, $\rho = \tfrac12$,
$\delta = \tfrac{1}{12}$, with $\delta + \rho r = \tfrac13 \le
\tfrac12$); the conclusion's existential half is exactly what
completeness buys. The two-point row being
required only \emph{on the ball} is deliberate: locality is what a
finite certificate can attest (globally, $\lVert A \rVert < 1$ may even
be false while the ball-restricted bound holds).
\end{bremark}

\begin{bremark}[relation to the literature]
The theorem is Banach's \cite{banach1922}; the closed-ball variant with
the condition $d(T\hat\mu, \hat\mu) \leq (1 - \rho)\, r$ and the
a-priori bound $\delta/(1-\rho)$ is the textbook local form
\cite[\S I.1]{granasdugundji2003}. The certificate-form contribution is
the \emph{anchoring inversion}: classical statements estimate the
distance from an iterate to the limit; the rule re-expresses every
hypothesis as a row checkable at the known snapshot in exact rational
arithmetic --- contraction as an operator-norm (or two-point) row,
residual as one evaluated sweep, radius as division-free arithmetic ---
so the analytic content survives while everything the checker touches
is decidable. The completeness gap for measure-space subjects
(Wasserstein distance over a Euclidean base) is carried as a named
hypothesis; the pilot certificate discharges it by pinning the subject
to the affine/Euclidean reduction of
Proposition~\ref{prop:qfb-affine}.
\end{bremark}

\begin{anchors}
Lemma \ref{lem:qfb-signs} & \texttt{delta\_nonneg}, \texttt{radius\_nonneg}\\
Lemma \ref{lem:qfb-maps} & \texttt{mapsTo\_closedBall}\\
Lemma \ref{lem:qfb-unique} & \texttt{fixedPoint\_unique}\\
Lemma \ref{lem:qfb-anchor} & \texttt{dist\_fixedPoint\_le}; \texttt{anchored\_le\_radius}\\
Theorem \ref{thm:qfb-sound} & \texttt{sound\_QFixBall} (existence via the contraction on the complete invariant ball)\\
Proposition \ref{prop:qfb-affine} & \texttt{lipschitzWith\_affine}; \texttt{sound\_QFixBall\_affine}\\
\end{anchors}


\rulesection{93}{QDistTransport}{core rule}
\label{blcatrule-QDistTransport}
\ruleblock{The rule (Steele metanotation)}
\sourcenote{as realized in the mechanized development; source row: \texttt{\texttt{QTSL-rules.toml}} --- \emph{QDistTransport}. The metanotation statement is typeset by hand for notational quality; its mechanized ground truth is the Lean statement below, extracted from the sources at every build of this book and hash-bound. The rule \emph{roster} is machine-checked (every mechanized rule tag round-trips to exactly one catalogue row); the statement \emph{text} is editorial, checked by eye against the extracted Lean beside it.}
\begin{mathpar}
\inferrule*{
  J = \bigl(\mathcal{M}_c \vDash s\bigr) \ \text{(strict, law-level)} \quad \text{or} \quad J = \bigl(\mathsf{QDist}[\mathrm{TV}](\mu_c, \nu) \le \delta\bigr), \ \text{on the coarse side} \\
  \bigl\{\vdash \mathrm{tv}\bigl(\pi_{\#}\, K(s, \cdot),\ Q(\pi(s), \cdot)\bigr) \le \varepsilon\bigr\}_{s \in S_{\mathrm{f}}} \ \text{(} \sigma\text{-witness rows, one per fine state; } \varepsilon = 0 \ \text{is exact intertwining)} \\
  \vdash \mu^{0}_{c} = \pi_{\#}\, \mu^{0}_{\mathrm{f}} \ \text{(initial row)} \\
  n \in \N \ \text{(declared horizon; } \varepsilon\text{-cell only)}
}{
  \text{exact cell: } J \ \text{transports unchanged} \qquad \varepsilon\text{-cell: } \mathsf{QDist}[\mathrm{TV}]\bigl(\pi_{\#}(\text{fine law at } n),\ \text{coarse law at } n\bigr) \le n\,\varepsilon
}
\end{mathpar}

\ruleflavor{$\pi : S_{\mathrm{f}} \to S_c$ is the state projection, $K$ and $Q$ the fine and coarse kernels, $\mu^0$ the initial laws; each $\sigma$-witness row asserts that one row of the fine kernel pushes along $\pi$ to within total-variation $\varepsilon$ of the coarse kernel\textquotesingle s row at the projected state. Unbounded horizons refuse without a contraction row. The displayed $J$ is restricted to the built judgment classes (strict law-level specs and $\mathsf{QDist}[\mathrm{TV}]$); the successor cells widen it as they land. Transport is a matrix (judgment class $\times$ witness class), each built cell carrying its own lemma and refusal condition; sequential witnesses compose under the certified budget $\varepsilon_1 + \varepsilon_2$ (an upper bound --- the actual composite defect may be smaller), and the strict column is the $\varepsilon = 0$ instantiation of the family. Built at the finite-kernel tier; the Lipschitz-pushforward ($\mathbb{W}_p$) and description-subject cells are named successors.}
\par\medskip
\rulename{QDistTransport} moves judgments across a simulation map. It
is a \emph{matrix} of transport cells --- judgment class (strict spec
vs.\ distance judgment) $\times$ witness class (exact intertwining
vs.\ $\varepsilon$-approximate intertwining) --- each cell carrying its
own conclusion form, its own lemma, and its own refusal condition. A
producer reaches for it to certify a fine model through a coarse one:
prove the spec on the abstraction, exhibit the simulation witness, and
transport. The two built kernel-tier cells: under an \emph{exact}
witness (every row of the fine kernel pushes along the readout map
exactly onto the coarse kernel\textquotesingle s row) laws transport
unchanged at every horizon, so any judgment about the coarse law pulls
back with zero degradation; under an $\varepsilon$-witness (every row
pushes to within total-variation distance $\varepsilon$) a strict
judgment does NOT transport --- what transports is a \emph{distance}
judgment, with $\varepsilon$ inherited and horizon-compounded: the
pushed fine law tracks the coarse law to $n\,\varepsilon$ at declared
horizon $n$. The degrading cell is what makes approximate simulation a
distance-judgment \emph{producer}.

What the checker recomputes: the per-row intertwining identities (exact
case) or per-row total-variation defects against the declared
$\varepsilon$ (approximate case), well-formedness of both kernels, and
the horizon arithmetic. The compounding law $n\,\varepsilon$ is proved,
never guessed: one-step defect by mixture convexity, then the coarse
Markov operator\textquotesingle s total-variation nonexpansiveness
(data processing --- the coarse chain cannot amplify tracking error),
then telescoping. Sequential witnesses compose with certified budget
$\varepsilon_1 + \varepsilon_2$ --- an upper bound with no additional
bookkeeping slack (plain addition is the honest lattice at this tier),
though the composite\textquotesingle s actual defect may be smaller.

What would break: without the coarse kernel\textquotesingle s row
stochasticity the operator could amplify differences and the telescoping
collapses; without a genuine probability law at the start, mixture
convexity fails; and without a contraction row the linear growth in $n$
is sharp, which is why the $\varepsilon$-cell refuses unbounded
horizons rather than extrapolating.

\cataloguenote{The strict column is the $\varepsilon = 0$ instantiation of the family
--- one biconditional, so the exact cell re-exports the marginal-process
lemma rather than restating it. Analytic lineage: approximate
(bi)simulation with quantitative degradation \cite{girardpappas2007}
and behavioural pseudometrics for probabilistic systems
\cite{desharnais2004}. Built at the finite-kernel tier (the scope the
Lean statements carry); the Lipschitz-pushforward cells for
$\mathbb{W}_p$ registry entries, the description-subject cells, and the
per-cell composition-coherence squares are named successors.
}
\ruleblock{Lean statement}
\declhead{BernsteinLogicAux.QDistTransport}{sound\_QDistTransport\_eps}{\badge{BadgeGreen}{statement bound}}
\begin{leansrc}
/-- **`sound_QDistTransport_eps` — the degrading cell's transport
theorem** (matrix cell: strict judgment × ε-witness ⇒ distance
judgment): at horizon `n`, the pushed fine law and the coarse law are
within TV-distance `n·ε`. Telescoping: one-step defect (`ε`, above) +
TV-nonexpansiveness of the coarse operator (§3) + triangle. This is
the theorem that makes "ε inherited, horizon-compounded" a proved
compounding law rather than a guess; with no contraction row the
linear growth is the honest rate (the refusal condition for unbounded
horizons lives spec-side). -/
theorem sound_QDistTransport_eps [DecidableEq β] {π : α → β}
    {K : α → α → ℝ} {Q : β → β → ℝ} {ε : ℝ} {μ₀ : α → ℝ}
    (hK : IsMarkov K) (hQ : IsMarkov Q) (hμ₀ : IsDist μ₀)
    (h : EpsPushKernel π K Q ε) (n : ℕ) :
    tvDist (push π (law K μ₀ n)) (law Q (push π μ₀) n) ≤ n * ε
\end{leansrc}
\declhead{BernsteinLogicAux.QDistTransport}{epsPushKernel\_comp}{\badge{BadgeGreen}{statement bound}}
\begin{leansrc}
/-- **ε-arithmetic composes** (derived, not guessed): transport along
`σ₁` then `σ₂` carries the certified budget `ε₁ + ε₂` — an upper bound
(triangle through the mid-tier pushforward + data processing; the actual
composite defect may be smaller). Plain addition is the right
composition law at the TV tier. -/
theorem epsPushKernel_comp [DecidableEq β] [DecidableEq γ]
    {π₁ : α → β} {π₂ : β → γ} {K₁ : α → α → ℝ} {K₂ : β → β → ℝ}
    {K₃ : γ → γ → ℝ} {ε₁ ε₂ : ℝ}
    (h₁ : EpsPushKernel π₁ K₁ K₂ ε₁) (h₂ : EpsPushKernel π₂ K₂ K₃ ε₂) :
    EpsPushKernel (π₂ ∘ π₁) K₁ K₃ (ε₁ + ε₂)
\end{leansrc}
\declhead{BernsteinLogicAux.QDistTransport}{sound\_QDistSelfTransport}{\badge{BadgeGreen}{statement bound}}
\begin{leansrc}
/-- **Matrix cell (distance judgment × strict witness), TV entry**: a
TV distance judgment transports along any simulation/readout map
unchanged (registry reading: TV → TV at Lipschitz constant 1). -/
theorem sound_QDistSelfTransport [DecidableEq β] (π : α → β)
    {μ ν : α → ℝ} {ε : ℝ} (h : tvDist μ ν ≤ ε) :
    tvDist (push π μ) (push π ν) ≤ ε
\end{leansrc}
\declhead{BernsteinLogicAux.QDistTransport}{strict\_cell\_exact}{\badge{BadgeGreen}{statement bound}}
\begin{leansrc}
/-- **Matrix cell (strict judgment × strict/kernel witness)**: under
the exact witness, laws transport EXACTLY at every horizon
(`QRelChain.push_law` re-exported as the matrix's strict kernel cell)
— so any judgment about the coarse law pulls back with zero
degradation, and the degrading cell's `n·ε` bound specializes to `0`. -/
theorem strict_cell_exact [DecidableEq β] {π : α → β} {K : α → α → ℝ}
    {Q : β → β → ℝ} (h : PushKernel π K Q) (μ₀ : α → ℝ) (n : ℕ) :
    push π (law K μ₀ n) = law Q (push π μ₀) n
\end{leansrc}
\declhead{BernsteinLogicAux.QDistTransport}{pushKernel\_iff\_eps\_zero}{\badge{BadgeGreen}{statement bound}}
\begin{leansrc}
/-- The exact intertwining is precisely the `ε = 0` witness. -/
theorem pushKernel_iff_eps_zero [DecidableEq β] {π : α → β}
    {K : α → α → ℝ} {Q : β → β → ℝ} :
    PushKernel π K Q ↔ EpsPushKernel π K Q 0
\end{leansrc}
\ruleblock{Correspondence (Steele display $\leftrightarrow$ Lean)}
\badge{BadgeGreen}{reviewer-certified against the displayed texts}
\sourcenote{The dictionary below is a reviewer-certified statement of how the mechanized statements above interpret the displayed rule: display metavariable $\mapsto$ Lean binder, each display hypothesis $\mapsto$ the premise or side condition carrying it, the display conclusion $\mapsto$ the goal --- with any deliberate difference in width between display and mechanization named outright. Certification is recorded as a hash triple over (the displayed cell, the extracted statement texts, this dictionary\textquotesingle s own text); if any of the three changes, this block turns amber until a reviewer re-certifies. The correspondence claim itself remains editorial --- reviewer-grade, not machine-checked; its machine-checked successor types each soundness statement against the rule as a structured object, and supersedes this block per rule as it lands.}
The display is a matrix and the bound set is its built column, cell
by cell. The $\varepsilon$-cell\textquotesingle s conclusion
$\mathsf{QDist}[\mathrm{TV}](\text{fine pushed law}, \text{coarse
law}) \le n\varepsilon$ at horizon $n$ $\mapsto$
\texttt{sound\_QDistTransport\_eps : tvDist (push π (law K μ₀ n))
(law Q (push π μ₀) n) ≤ n * ε} --- the compounding law as a proved
induction (one-step defect + TV-nonexpansiveness of the coarse
operator + triangle), never a guess.
\begin{itemize}
\item the $\sigma$-witness rows $\mapsto$ \texttt{EpsPushKernel π K Q
  ε} (every fine row pushes to within TV $\varepsilon$ of the coarse
  row at the projected state), with the exact-intertwining column
  certified as its $\varepsilon = 0$ face:
  \texttt{pushKernel\_iff\_eps\_zero}, and the strict cell
  \texttt{strict\_cell\_exact} re-exporting \texttt{push\_law} ---
  laws transport \emph{exactly} at every horizon, zero degradation.
\item the initial row (coarse initial law $= \pi_{\#}$ fine initial)
  $\mapsto$ carried by construction: the coarse arguments in every
  goal ARE \texttt{push π μ₀} --- mismatched initials are
  unrepresentable in the statement, so no hypothesis row exists.
\item the declared horizon $\mapsto$ the explicit \texttt{n : ℕ}
  binder; the refusal condition for unbounded horizons without a
  contraction row is spec-side with no Lean counterpart --- the
  honest linear rate IS the certified statement.
\item sequential accumulation $\mapsto$ \texttt{epsPushKernel\_comp}:
  certified composite budget $\varepsilon_1 + \varepsilon_2$, plain
  addition as the displayed lattice --- an upper bound (the actual
  composite defect may be smaller); the distance-judgment $\times$
  strict-witness cell $\mapsto$ \texttt{sound\_QDistSelfTransport}
  (TV data processing at $L = 1$).
\item consumed regularity: the $\varepsilon$-cell alone demands
  well-formedness (\texttt{IsMarkov} both kernels, \texttt{IsDist
  μ₀}) --- the mixture-convexity and nonexpansiveness steps use it
  --- while the strict cell and the composition law are
  hypothesis-free beyond their witnesses: visible in the binder
  rows.
\end{itemize}
Width note: finite-kernel tier (\texttt{Fintype},
\texttt{DecidableEq} for fibre sums), as the flavor text
confines it; the unbuilt cells ($\mathbb{W}_p$ Lipschitz
pushforward, description subjects, per-cell coherence squares) are
the conditional residue, named in the note, bound by nothing --- and
the judgment-class abstraction of the display\textquotesingle s $J$
(``or a strict spec\textquotesingle\textquotesingle) is realized
concretely: strict specs ride \texttt{strict\_cell\_exact}, distance
judgments ride the TV theorems; no single parameterized $J$ exists in
Lean, the matrix\textquotesingle s rows are separate theorems.

\ruleblock{Lean verifier receipts}
\sourcenote{module \texttt{BernsteinLogicAux.QDistTransport}; bound build \texttt{4bd1c5103f38\ldots} (bound from receipt 6ff1345414960d8f, 2026-07-21T05:38:55Z). A receipt records that named checking layers accepted a build: (module, compiled-artifact SHA-256, environment SHA-256, checker layers), minted when the checking gate ran; the store receipts attesting the bound build are listed with the layers each run credited. Nothing here asks the reader to trust the store: every layer is an independently runnable program over the quoted sources --- rebuild the module and re-run the checkers to reproduce the verdict. Receipts attesting superseded builds credit no layer and are omitted from this rendering (the count of omissions, if any, is noted below).}
\begin{receipttable}
\receiptrow{6ff1345414960d8f}{2026-07-21T05:38:55Z}{build+replay+hardened+decorrelated}{\badge{BadgeGreen}{current}}
\end{receipttable}
\sourcenote{13 further receipt(s) in the store attest superseded builds of this module and are omitted from this rendering.}
\layerline{\layeritem{build}{lake build (elaboration)}{\badge{BadgeGreen}{current}}\layeritem{replay}{leanchecker (official kernel replay)}{\badge{BadgeGreen}{current}}\layeritem{hardened}{safe\_verify (hardened screen)}{\badge{BadgeGreen}{current}}\layeritem{decorrelated}{lean4export + nanoda (decorrelated kernel)}{\badge{BadgeGreen}{current}}}
\sourcenote{\badge{BadgeGray}{comparator: none} no fixture in the three-implementation comparator harness (compiled / kernel-interpreted / reference-interpreted) exercises this rule.}
\ruleblock{HOL4 statement}
\absentbox{NO HOL STATEMENT YET}{The HOL4 development checks the base calculus\textquotesingle s wire
format (theory \texttt{bern\_catalogue}). The transport
matrix\textquotesingle s cells conclude judgments of the extension
theory, whose rule sorts are out of scope of the current baked theory;
the mechanized cells live in the auxiliary Lean package. Nothing exists
yet for a HOL statement to be about; the box becomes writable when a
revision carries the distance-judgment rows and its checker chain is
ported.
}
\ruleblock{HOL verifier receipts}
\absentbox{NONE}{no HOL verification receipt exists for this rule.}
\ruleblock{Proof}

\sourcenote{Genre: core rule (tier 1). Scope, stated loudly: the
finite-kernel tier --- finite state spaces, distribution evolution,
total-variation entries of the registry. The Lipschitz-pushforward
cells for Wasserstein entries, the description-subject cells, and the
per-cell composition-coherence squares are named successors, not
claimed. Every hypothesis is used and remarked on.}

\textbf{Setting and claim.}\enspace A \emph{fine} chain (states
$\alpha$, kernel $K$, initial law $\mu_0$), a \emph{coarse} chain
(states $\beta$, kernel $Q$), and a readout map
$\pi : \alpha \to \beta$. Write $\pi_{\#}$ for pushforward,
$\mathrm{law}_K(n)$ for the time-$n$ law, and
$\mathrm{TV}(p, q) = \tfrac12 \sum_l |p(l) - q(l)|$. The witness comes
in two grades: \emph{exact intertwining} --- every kernel row pushes
exactly, $\pi_{\#} K(s, \cdot) = Q(\pi(s), \cdot)$ --- and
\emph{$\varepsilon$-intertwining} ---
$\mathrm{TV}\bigl(\pi_{\#} K(s, \cdot),\ Q(\pi(s), \cdot)\bigr) \leq
\varepsilon$ for every $s$. The transport matrix's built cells, which
are the claims proved below: \textbf{(strict $\times$ exact)} laws
transport unchanged at every horizon, so any judgment about the coarse
law pulls back with zero degradation
(\texttt{strict\_cell\_exact}); \textbf{(distance $\times$ strict
witness)} a total-variation distance judgment moves along any readout
map unchanged (\texttt{sound\_QDistSelfTransport});
\textbf{(strict $\times$ $\varepsilon$-witness --- the degrading
cell)} what transports is a \emph{distance} judgment:
$\mathrm{TV}\bigl(\pi_{\#}\mathrm{law}_K(n),\
\mathrm{law}_Q(n)\bigr) \leq n\,\varepsilon$ at declared horizon $n$
(\texttt{sound\_QDistTransport\_eps}); and witnesses compose with
$\varepsilon_1 + \varepsilon_2$ (\texttt{epsPushKernel\_comp}).

\begin{blemma}[data processing: pushforward is TV-nonexpansive]
\label{lem:qdt-dpi}
For any map $\pi$ and any $p, q$:
$\mathrm{TV}(\pi_{\#} p, \pi_{\#} q) \leq \mathrm{TV}(p, q)$.
\end{blemma}
\begin{proof}
Per label $b$, $|\sum_{\pi(s) = b} (p - q)(s)| \leq
\sum_{\pi(s) = b} |p - q|(s)$; summing over $b$ regroups the right side
into $\sum_s |p - q|(s)$. Merging states can only cancel differences,
never create them \cite[Ch.~4]{levinperes2017}. This lemma \emph{is}
the (distance $\times$ strict-witness) cell: a certified
$\mathrm{TV} \leq \varepsilon$ on the fine side transports to the
coarse side unchanged --- the registry reading of total variation as a
Lipschitz-$1$ entry under pushforward.
\end{proof}

\begin{blemma}[the coarse operator cannot amplify]\label{lem:qdt-markov}
If $Q$ is row-stochastic with nonnegative entries, then
$\mathrm{TV}(\mathrm{step}_Q\, p,\ \mathrm{step}_Q\, q) \leq
\mathrm{TV}(p, q)$.
\end{blemma}
\begin{proof}
$|\mathrm{step}_Q(p - q)(t)| \leq \sum_s |p - q|(s)\, Q(s, t)$ by the
triangle inequality and $Q \geq 0$; summing over $t$ and swapping sums,
the inner $\sum_t Q(s, t) = 1$ leaves $\sum_s |p - q|(s)$. Both
hypotheses on $Q$ are consumed: nonnegativity to pull the absolute
value through, unit row sums to collapse the mass.
\end{proof}

\begin{blemma}[one-step defect]\label{lem:qdt-defect}
Let $\mu$ be a probability law and let the $\varepsilon$-intertwining
rows hold. Then
$\mathrm{TV}\bigl(\pi_{\#}(\mathrm{step}_K \mu),\
\mathrm{step}_Q(\pi_{\#}\mu)\bigr) \leq \varepsilon$.
\end{blemma}
\begin{proof}
Both sides are $\mu$-mixtures of per-state rows:
$\pi_{\#}(\mathrm{step}_K \mu) = \sum_s \mu(s)\,
\pi_{\#}K(s, \cdot)$ and $\mathrm{step}_Q(\pi_{\#}\mu) =
\sum_s \mu(s)\, Q(\pi(s), \cdot)$ (the second by regrouping the coarse
step over the fibres of $\pi$). Total variation is jointly convex ---
a mixture of pairs is bounded by the mixture of pairwise distances ---
so the defect is at most $\sum_s \mu(s) \cdot \varepsilon =
\varepsilon$, using that $\mu$ is nonnegative with total mass $1$.
\end{proof}

\begin{bthm}[the degrading cell: horizon-compounded transport]
\label{thm:qdt-eps}
With $K, Q$ row-stochastic, $\mu_0$ a probability law, and the
$\varepsilon$-intertwining rows: for every $n$,
\[
\mathrm{TV}\bigl(\pi_{\#}\,\mathrm{law}_K(n),\ \mathrm{law}_Q(n)\bigr)
\;\leq\; n\, \varepsilon,
\]
where the coarse chain starts from $\pi_{\#}\mu_0$.
\end{bthm}
\begin{proof}
Induction on $n$; the base is $\mathrm{TV}(x, x) = 0$. Step, by the
triangle inequality through the hybrid term
$\mathrm{step}_Q\bigl(\pi_{\#}\,\mathrm{law}_K(n)\bigr)$:
\[
\mathrm{TV}\bigl(\pi_{\#}\,\mathrm{law}_K(n{+}1),\ \mathrm{law}_Q(n{+}1)\bigr)
\;\leq\;
\underbrace{\mathrm{TV}\bigl(\pi_{\#}\mathrm{step}_K\,\mathrm{law}_K(n),\
\mathrm{step}_Q\,\pi_{\#}\mathrm{law}_K(n)\bigr)}_{\leq\ \varepsilon
\ \text{(Lemma \ref{lem:qdt-defect}, at the time-}n\text{ law)}}
\;+\;
\underbrace{\mathrm{TV}\bigl(\mathrm{step}_Q\,\pi_{\#}\mathrm{law}_K(n),\
\mathrm{step}_Q\,\mathrm{law}_Q(n)\bigr)}_{\leq\ n \varepsilon
\ \text{(Lemma \ref{lem:qdt-markov} + induction)}} .
\]
The compounding law is thus \emph{proved}, not guessed: one fresh
defect per step, prior error never amplified. A strict judgment does
NOT transport along an $\varepsilon$-witness; what transports is a
distance judgment with $\varepsilon$ inherited and horizon-compounded
--- precisely this cell.
\end{proof}

\begin{bprop}[witness composition]\label{prop:qdt-comp}
If $\sigma_1 = (\pi_1, \varepsilon_1)$ intertwines $K_1$ with $K_2$ and
$\sigma_2 = (\pi_2, \varepsilon_2)$ intertwines $K_2$ with $K_3$, then
$\pi_2 \circ \pi_1$ is an $(\varepsilon_1 + \varepsilon_2)$-witness
from $K_1$ to $K_3$.
\end{bprop}
\begin{proof}
At any state $s$, triangle through the mid-tier row
$\pi_{2\#} K_2(\pi_1(s), \cdot)$: the first leg is
$\mathrm{TV}\bigl(\pi_{2\#}\pi_{1\#}K_1(s,\cdot),\
\pi_{2\#}K_2(\pi_1 s, \cdot)\bigr) \leq \varepsilon_1$ by
Lemma~\ref{lem:qdt-dpi} applied to $\sigma_1$'s row defect; the second
leg is $\sigma_2$'s row at $\pi_1(s)$. Plain addition records the honest
\emph{upper bound} at this tier --- $\varepsilon_1 + \varepsilon_2$ is
what the triangle gives, with no minimality claim: data processing can
gain slack (a constant second projection has composite defect $0$
whatever $\varepsilon_1$ was).
\end{proof}

\begin{bprop}[the strict column is the $\varepsilon = 0$ instantiation]
\label{prop:qdt-strict}
Exact intertwining holds iff the rows are $0$-intertwining; under it,
$\pi_{\#}\,\mathrm{law}_K(n) = \mathrm{law}_Q(n)$ for every $n$, and
every judgment about the coarse law pulls back with zero degradation.
\end{bprop}
\begin{proof}
($\Rightarrow$) equal rows have distance $0$. ($\Leftarrow$)
$\mathrm{TV}(p, q) \leq 0$ forces $p = q$ entrywise (a sum of absolute
values is zero only if each vanishes). The exact law transport is then
the marginal-process induction (the same pushed-row-identity induction
as the \rulename{QCoupleDrift} section's marginal-process lemma, there
instantiated at product projections, here at an arbitrary readout
map), or equivalently Theorem~\ref{thm:qdt-eps} at $\varepsilon = 0$.
\end{proof}

\begin{bremark}[why each hypothesis is needed]\label{rem:qdt-omit}
\emph{Row-stochasticity of $Q$} (Lemma~\ref{lem:qdt-markov}): a
non-stochastic coarse operator amplifies --- with $Q = 2\,\mathrm{Id}$
each step doubles any existing difference and the telescoping
collapses; with signed entries the absolute value cannot be pulled
through. \emph{$\mu$ a probability law}
(Lemma~\ref{lem:qdt-defect}): mixture convexity needs nonnegative
weights of total mass one --- weights summing to $2$ double the
defect. \emph{The declared horizon}: the linear rate is sharp to first order
without further rows --- on the two-state fine chain that leaks mass
$\varepsilon$ toward an absorbing label each step, against a static
coarse chain, the distance at horizon $n$ is
$1 - (1 - \varepsilon)^n = n\varepsilon - O(n^2\varepsilon^2)$, so no
$o(n\varepsilon)$ bound holds, and over unbounded horizons the error
climbs to order one --- the $\varepsilon$-cell therefore
\emph{refuses} unbounded horizons rather than extrapolating; a
contraction row is the named mechanism that would buy a
uniform-in-$n$ bound, and it is a different certificate. \emph{Per-row quantification} of the intertwining: a
defect bound only at states charged by $\mu_0$ fails at later times,
since the law spreads --- the rows are quantified over all states
precisely so Lemma~\ref{lem:qdt-defect} applies at \emph{every}
time-$n$ law.
\end{bremark}

\begin{bremark}[relation to the literature]
The degrading cell is the kernel-tier, exact-arithmetic face of
approximate simulation with quantitative degradation
\cite{girardpappas2007}: there, an $\varepsilon$-approximate
(bi)simulation transports properties up to a metric slack; here the
slack is total variation, the systems are finite kernels, and the
horizon-compounding law is proved from mixture convexity plus data
processing. Behavioural pseudometrics for probabilistic systems
\cite{desharnais2004} are the semantic tradition the distance-judgment
registry discretizes; the coupling toolkit is classical
\cite{lindvall1992,levinperes2017}. The matrix framing --- judgment
class $\times$ witness class, one lemma and one refusal condition per
cell --- is the rule's own contribution: it makes ``what survives
transport along a defective witness'' a typed question with a proved
answer per cell, instead of a convention.
\end{bremark}

\begin{anchors}
Lemma \ref{lem:qdt-dpi} & \texttt{tvDist\_push\_le}; cell form \texttt{sound\_QDistSelfTransport}\\
Lemma \ref{lem:qdt-markov} & \texttt{tvDist\_step\_le}\\
Lemma \ref{lem:qdt-defect} & \texttt{EpsPushKernel}; \texttt{push\_step\_defect}\\
Theorem \ref{thm:qdt-eps} & \texttt{sound\_QDistTransport\_eps} (triangle: \texttt{tvDist\_triangle})\\
Proposition \ref{prop:qdt-comp} & \texttt{epsPushKernel\_comp}; \texttt{push\_comp}\\
Proposition \ref{prop:qdt-strict} & \texttt{pushKernel\_iff\_eps\_zero}; \texttt{eq\_of\_tvDist\_le\_zero}; \texttt{strict\_cell\_exact}\\
\end{anchors}


\blcatsection{Stored-Chain Certificate Species}
\rulesection{94}{DtmcAbsorbBracket}{core rule}
\label{blcatrule-DtmcAbsorbBracket}
\ruleblock{The rule (Steele metanotation)}
\sourcenote{as realized in the mechanized development; source row: \texttt{\texttt{QTSL-rules.toml}} --- \emph{DtmcAbsorbBracket}. The metanotation statement is typeset by hand for notational quality; this rule cites no Lean statements yet (the absence is stated loudly below). The rule \emph{roster} is machine-checked (every mechanized rule tag round-trips to exactly one catalogue row); the statement \emph{text} is editorial.}
\begin{mathpar}
\inferrule*{
  q : \mathsf{DAbsQuery} \\
  \mathsf{qtsl\_dabs\_wf}(q) \\
  \text{stored rows of } q.\mathsf{chain}\ (\mathsf{dprob}, \mathsf{dinit}, \mathsf{dclass}, \mathsf{dreject})
}{
  \mathsf{qtsl\_dabs\_absorb\_bracket}(q,\ c,\ lo,\ hi)
}
\end{mathpar}

\par\medskip
\rulename{DtmcAbsorbBracket} is a certificate species rather than a
proof-node rule: a checker-recomputed claim over the stored rows of a
finite chain, introducing zero new judgment sorts and zero new
specification formers. It certifies guaranteed two-sided brackets on
the eventual-absorption measure of one class of a finite stored-row
chain --- and, through its normalized sibling
\rulename{DtmcPosteriorBracket}, on normalized conditional measures
(posteriors of discrete probabilistic programs, where the compilation
of a program to a chain is an untrusted translation outside the trust
path). The species is stated at full generality: any finite chain, any
class labeling, any horizon.

The evaluation semantics, once: $m_0$ is the chain\textquotesingle s
stored initial measure; $m_{t+1} = m_t K$, where $K$ is the stored
transition-row matrix COMPLETED by a self-loop at every state with no
stored outgoing row (\emph{rowless-is-absorbing}). For a query with
horizon $T$, class labeling, and rejection labeling: $N_c$ is the
class-$c$ mass of $m_T$; $S = \sum_c N_c$; $R = 1 - S -
(\text{rejected mass of } m_T)$; and $R_c$ is the
\emph{reachability-sharpened} per-class residual --- the part of $R$
sitting at states from which some class-$c$ state is reachable
through strictly positive stored rows (a backward walk seeded at the
class, run under the same work budget; budget exhaustion during the
walk refuses the whole claim, never verdicts). The claim holds iff
$lo \le N_c$ and $N_c + R_c \le hi$. Mass that can never arrive in
the class is never charged to it: a class unreachable from every
alive state has $R_c = 0$ and a width-zero bracket at any $T \ge 1$,
with no termination assumption; and $R_c \le R$ always, so any claim
admissible under the uniform-residual reading stays admissible.

The completion convention is the species\textquotesingle{}
load-bearing design move: a truncated row set can only WIDEN brackets
(unclassified mass accrues to the residual, never to any class), so
row-set completeness is a tightness concern rather than a soundness
obligation --- there is no closure side condition, deliberately, and
under-expansion is unrepresentable as an attack. Derived quantities
($N_c$, $S$, $R$, $R_c$, ratio denominators) are never stored: the
evaluator recomputes them from the chain rows, so forging an
intermediate is unrepresentable; claims may be looser than canonical,
never tighter.

Soundness has two sides. Below, absorption is monotone: classed states
are demanded stored-absorbing (the well-formedness side condition),
so class mass never decreases past the horizon. Above, the
class\textquotesingle s co-reachable set $B_C$ (states with a
strictly-positive-support path into the class; $C \subseteq B_C$) has
a forward-closed complement --- no positive-support row leads from
outside $B_C$ into it --- so the mass of $B_C$ is nonincreasing and
later class gain is bounded by $\mathrm{mass}_T(B_C \setminus C) =
R_c$. The uniform-residual bound ($R_c$ relaxed to $1 -
\mathrm{mass}_T(\mathrm{Abs})$) is the pre-sharpening corollary. The
argument is mechanized in the accompanying Lean development
(\texttt{SuFunctor/Adequacy/DtmcAbsorb.lean}; sorry-free, standard
axiom closure, axiom-audited on every build): the sharpened upper
side \texttt{dtmcAbsorb\_mass\_le\_reach} is notably HYPOTHESIS-FREE
on the class --- it holds for any class of any finite stochastic
chain, the well-formedness demands protecting only the monotone lower
side; beside it sit \texttt{dtmcAbsorbBracket\_reach\_sound} (the
bracket, both sides), \texttt{dtmcAbsorb\_reach\_residual\_le} (the
never-loosens warrant $R_c \le 1 - \mathrm{mass}_T(\mathrm{Abs})$,
itself needing only that absorbing states absorb),
\texttt{dtmcAbsorb\_eventual\_bracket\_reach} (the claim shape), the
forward-closed generalization (\texttt{FinDTMC.ForwardClosed},
\texttt{dtmc\_mass\_mono\_of\_forwardClosed}) from which the
incumbent monotonicity lemmas re-derive in one line, and a
kernel-checked strict-improvement witness
(\texttt{idBool\_strict\_improvement}: an unreachable class where the
uniform-residual bracket has width one and the sharpened bracket
width zero --- a standing regression tooth, re-checked every build).
The mechanized statements are deliberately universally quantified
($\forall\, t' \ge T$) rather than supremum-shaped: over $\Q$ the
supremum of a monotone bounded mass sequence need not exist as a
rational, and the universal form is the citable content --- any
eventual-absorption reading (sup or limit, in any completion of
$\Q$) inherits the claimed bracket.

One clause of the species is graded prose by design: the
\emph{truncation-transfer} reading --- stored rows agreeing with an
external program\textquotesingle s chain on rowed states implies the
bracket transfers to the program --- is a coupling argument about a
translation that lives OUTSIDE the trust path, and its assurance
instrument is a battery (vendored source, Monte-Carlo differential,
closed-form containment), not a kernel theorem. The certificate
itself never claims more than the stored rows warrant.

\cataloguenote{The schema is baked in the certificate theory as a species sibling
extending the chain grain, purely additive --- the query sort with
anchor-mediated chain binding (the query\textquotesingle s chain IS
its anchor\textquotesingle s chain, labels bound by chase rows), the
horizon field, the class and rejection labelings, and the three
side-condition Props with their report, refusal, and chase rows, all
quoted below from the theory source. The doubly-classed state is
unrepresentable at the attribute grain (the class labeling is a
keyed function row). The mechanized soundness lemmas are named in
the commentary; no Lean statement is bound into this book yet ---
the binding column is honest open debt, and joins when the
statement-binding machinery registers the development\textquotesingle s
package.
}
\ruleblock{Theory rows (geolog)}
\sourcenote{The rule\textquotesingle s certificate content as it is declared in the geolog theory the checker loads --- payload schema, side-condition demands, and axiom rows, quoted from the theory source at every build of this catalogue (a stale quotation cannot exist). Axiom modes: \texttt{chase} = the engine derives; \texttt{report} = a violation is a loud judgment; \texttt{refuse} = a violation gates admission. Side-condition relations are discharged by bounded evaluation at the engine\textquotesingle s ground seam, never by derivation.}
\par\medskip
The chain grain this species reads: states, stored transition rows,
stored initial rows, and the row-mass admission demand. The species
adds no rows here; it consumes them.

\declhead{theory Dtmc}{DState}{\badge{BadgeGreen}{quoted from the theory source at render}}
\begin{leansrc}
DState : Set;   // one state of one chain
\end{leansrc}
\declhead{theory Dtmc}{dprob}{\badge{BadgeGreen}{quoted from the theory source at render}}
\begin{leansrc}
dprob : {src: DState, dst: DState} -> Q;
\end{leansrc}
\declhead{theory Dtmc}{dinit}{\badge{BadgeGreen}{quoted from the theory source at render}}
\begin{leansrc}
dinit : {chain: Dtmc, state: DState} -> Q;
\end{leansrc}
\declhead{theory Dtmc}{qtsl\_dtmc\_rowmass}{\badge{BadgeGreen}{quoted from the theory source at render}}
\begin{leansrc}
qtsl_dtmc_rowmass  : {chain: Dtmc} -> Prop;   // side-condition
\end{leansrc}
\par\medskip
The baked rows of the species: the query sort with its anchor and
horizon, the class and rejection labelings, the disjointness
refusal, the labels-bind-own-chain chase rows, the well-formedness
demand, and the two bracket claims with their interval-order
reports.

\declhead{theory DtmcAbsorb}{DAbsQuery}{\badge{BadgeGreen}{quoted from the theory source at render}}
\begin{leansrc}
DAbsQuery : Set;
\end{leansrc}
\declhead{theory DtmcAbsorb}{DAbsQuery/*}{\badge{BadgeGreen}{quoted from the theory source at render}}
\begin{leansrc}
DAbsQuery/anchor  : DAbsQuery -> DState;
DAbsQuery/horizon : DAbsQuery -> Nat;
\end{leansrc}
\declhead{theory DtmcAbsorb}{dclass}{\badge{BadgeGreen}{quoted from the theory source at render}}
\begin{leansrc}
dclass : {q: DAbsQuery, state: DState} -> Nat;
\end{leansrc}
\declhead{theory DtmcAbsorb}{dreject}{\badge{BadgeGreen}{quoted from the theory source at render}}
\begin{leansrc}
dreject : {q: DAbsQuery, state: DState} -> Prop;
\end{leansrc}
\declhead{theory DtmcAbsorb}{ax/wf/dabs\_class\_reject\_disjoint}{\badge{BadgeGreen}{quoted from the theory source at render}}
\begin{leansrc}
refuse ax/wf/dabs_class_reject_disjoint : forall q : DAbsQuery,
    s : DState, c : Nat.
  c = {q: q, state: s} dclass,
  {q: q, state: s} dreject
  |- false;
\end{leansrc}
\declhead{theory DtmcAbsorb}{ax/wf/dabs\_class\_same\_dtmc}{\badge{BadgeGreen}{quoted from the theory source at render}}
\begin{leansrc}
chase ax/wf/dabs_class_same_dtmc : forall q : DAbsQuery, s : DState,
    c : Nat.
  c = {q: q, state: s} dclass
  |- s DState/dtmc = q DAbsQuery/anchor DState/dtmc;
\end{leansrc}
\declhead{theory DtmcAbsorb}{ax/wf/dabs\_reject\_same\_dtmc}{\badge{BadgeGreen}{quoted from the theory source at render}}
\begin{leansrc}
chase ax/wf/dabs_reject_same_dtmc : forall q : DAbsQuery, s : DState.
  {q: q, state: s} dreject
  |- s DState/dtmc = q DAbsQuery/anchor DState/dtmc;
\end{leansrc}
\declhead{theory DtmcAbsorb}{qtsl\_dabs\_wf}{\badge{BadgeGreen}{quoted from the theory source at render}}
\begin{leansrc}
qtsl_dabs_wf : {q: DAbsQuery} -> Prop;   // side-condition
\end{leansrc}
\declhead{theory DtmcAbsorb}{ax/wf/dabs\_wf}{\badge{BadgeGreen}{quoted from the theory source at render}}
\begin{leansrc}
report ax/wf/dabs_wf : forall q : DAbsQuery.
  |- {q: q} qtsl_dabs_wf;
\end{leansrc}
\declhead{theory DtmcAbsorb}{qtsl\_dabs\_absorb\_bracket}{\badge{BadgeGreen}{quoted from the theory source at render}}
\begin{leansrc}
qtsl_dabs_absorb_bracket : {q: DAbsQuery, class: Nat, lo: Q, hi: Q}
  -> Prop;   // side-condition
\end{leansrc}
\declhead{theory DtmcAbsorb}{ax/wf/dabs\_absorb\_order}{\badge{BadgeGreen}{quoted from the theory source at render}}
\begin{leansrc}
report ax/wf/dabs_absorb_order : forall q : DAbsQuery, c : Nat,
    l : Q, h : Q.
  {q: q, class: c, lo: l, hi: h} qtsl_dabs_absorb_bracket
  |- {l: l, r: h} q_le;
\end{leansrc}
\ruleblock{Lean statement}
\badge{BadgeAmber}{not yet mechanized} this rule cites no Lean statements: nothing on this page is hash-bound for it, and its dashboard row shows the same absence across the statement, receipt, and correspondence columns.
\ruleblock{HOL4 statement}
\absentbox{NO HOL STATEMENT YET}{No HOL development of the extension theory exists; this
species\textquotesingle{} content is an exact-rational evaluator
recompute over baked rows, and its soundness theorems live in the
Lean development named in the commentary.
}
\ruleblock{HOL verifier receipts}
\absentbox{NONE}{no HOL verification receipt exists for this rule.}
\ruleblock{Proof}

\sourcenote{Genre: certificate species (tier 1 semantic soundness;
checker recomputation, no proof node). The proof below is the
mechanized argument restated at textbook detail; the named Lean
declarations carry it, sorry-free at standard axiom closure.}

\textbf{Setting.}\enspace A finite stored-row chain: state set $X$,
stored transition rows $K_0 \subseteq X \times X \to \Q_{\ge 0}$
satisfying the row-mass admission demand (each rowed source's stored
masses total exactly one), stored initial measure $m_0$, horizon
$T \in \N$, class labeling and rejection labeling over $X$. Let $K$
be $K_0$ completed by a unit self-loop at every state with no stored
outgoing row (\emph{rowless-is-absorbing}); $K$ is then a stochastic
matrix on $X$, and $m_{t+1} := m_t K$ defines the exact-rational
evolution. Write $C$ for the set of class-$c$ states,
$\mathrm{Abs}$ for the labeled (classed or rejected) states --- each
demanded stored-absorbing by the well-formedness side condition:
rowless, or exactly the unit self-loop --- and
$N_c := \mathrm{mass}_T(C)$, $S := \sum_{c'} N_{c'}$,
$R := 1 - S - (\text{rejected mass of } m_T)$.

\textbf{Lemma 1 (monotone absorption --- the lower side).}\enspace
For $t' \ge T$: $\mathrm{mass}_{t'}(C) \ge \mathrm{mass}_T(C)$.
\emph{Proof.} Each state of $C$ is stored-absorbing, so under $K$ its
mass never leaves; class mass is nondecreasing step by step.
(Removing the well-formedness hypothesis admits the counterexample
the mutant battery pins: a classed state with a stored mass-leaking
outgoing row loses class mass at the very next step.)\hfill$\square$

\textbf{Lemma 2 (forward-closed complement --- the upper side).}\enspace
Let $B_C$ be the co-reachable set of $C$: states admitting a path
into $C$ through strictly positive stored rows ($C \subseteq B_C$).
Then $X \setminus B_C$ is forward-closed --- no positive-support row
of $K$ leads from outside $B_C$ into $B_C$ (such a row would itself
witness co-reachability of its source) --- hence
$\mathrm{mass}_{t'}(B_C)$ is nonincreasing in $t'$.
\emph{Proof.} For a forward-closed set $F$, $\mathbf{1}_F K^\top
\ge \mathbf{1}_F$ pointwise on $F$'s complement contributes nothing,
and mass on $F$ can only stay or grow; apply to
$F = X \setminus B_C$ and complement. This is the mechanized
generalization (\texttt{FinDTMC.ForwardClosed},
\texttt{dtmc\_mass\_mono\_of\_forwardClosed}), from which Lemma~1
re-derives in one line: a stored-absorbing set is forward-closed
around each of its points.\hfill$\square$

\textbf{Theorem (the bracket).}\enspace Define
$R_c := \mathrm{mass}_T(B_C \setminus C)$ (the
reachability-sharpened per-class residual, recomputed by a backward
walk over strictly positive stored rows, seeded at $C$, under the
same work budget --- exhaustion refuses the whole claim). For every
$t' \ge T$:
$$\mathrm{mass}_{t'}(C) \in [\,\mathrm{mass}_T(C),\
\mathrm{mass}_T(C) + R_c\,].$$
\emph{Proof.} Lower: Lemma 1. Upper: class mass at $t'$ sits inside
$B_C$ (mass outside $B_C$ can never arrive in $C$, by Lemma~2
applied along the trajectory), and $C$'s gain since $T$ is at most
$B_C$'s non-$C$ mass at $T$, which is $R_c$, since $B_C$'s total is
nonincreasing while its $C$-part is nondecreasing.\hfill$\square$

The mechanized upper side (\texttt{dtmcAbsorb\_mass\_le\_reach}) is
HYPOTHESIS-FREE on the class: it holds for any class of any finite
stochastic chain --- the well-formedness demands protect only
Lemma~1. The never-loosens warrant
(\texttt{dtmcAbsorb\_reach\_residual\_le}) gives $R_c \le
1 - \mathrm{mass}_T(\mathrm{Abs})$, needing only that labeled states
absorb, so every claim admissible under the uniform residual remains
admissible under the sharpened one; and the strict-improvement tooth
(\texttt{idBool\_strict\_improvement}) exhibits an unreachable class
where the uniform bracket has width one and the sharpened bracket
width zero --- a standing kernel-checked regression witness. A class
unreachable from every alive state has $R_c = 0$: a width-zero
bracket at any $T \ge 1$, with no termination assumption anywhere.

\textbf{Corollary (eventual readings).}\enspace The statements are
universally quantified over $t' \ge T$ deliberately: over $\Q$ the
supremum of the monotone bounded sequence $\mathrm{mass}_{t'}(C)$
need not exist as a rational. Any eventual-absorption reading ---
supremum or limit, in any completion of $\Q$ --- inherits the
claimed bracket from the universal form
(\texttt{dtmcAbsorb\_eventual\_bracket\_reach}).

\textbf{Why each hypothesis.}\enspace Row-mass admission: without
it, doubled stored masses inflate $N_c$ (a battery locus). The
well-formedness demand: Lemma 1's counterexample above. The budget
refusal on the backward walk: a partial walk could under-build
$B_C$, dropping states whose mass belongs in $R_c$ --- so exhaustion
refuses rather than verdicts. Interval order $lo \le hi$ is a report
row; inversion refuses. And the completion convention
(rowless-is-absorbing) makes truncation \emph{widen} brackets:
omitted rows leave their states absorbing and their mass in the
residual, never in a class --- there is no closure side condition to
attack, and claims may be looser than canonical, never tighter.

\begin{bremark}[status]
The schema is baked (species rows quoted in this section); the
argument above is mechanized sorry-free at standard axiom closure,
axiom-audited on every build, in the module named in the commentary.
No Lean statement is bound into this book yet: the binding column is
honest open debt until the statement-binding machinery registers the
development's package.
\end{bremark}


\rulesection{95}{DtmcPosteriorBracket}{core rule}
\label{blcatrule-DtmcPosteriorBracket}
\ruleblock{The rule (Steele metanotation)}
\sourcenote{as realized in the mechanized development; source row: \texttt{\texttt{QTSL-rules.toml}} --- \emph{DtmcPosteriorBracket}. The metanotation statement is typeset by hand for notational quality; this rule cites no Lean statements yet (the absence is stated loudly below). The rule \emph{roster} is machine-checked (every mechanized rule tag round-trips to exactly one catalogue row); the statement \emph{text} is editorial.}
\begin{mathpar}
\inferrule*{
  q : \mathsf{DAbsQuery} \\
  \mathsf{qtsl\_dabs\_wf}(q) \\
  \text{stored rows of } q.\mathsf{chain}\ (\mathsf{dprob}, \mathsf{dinit}, \mathsf{dclass}, \mathsf{dreject}) \\
  S > 0\ \text{(recomputed; refusal otherwise)}
}{
  \mathsf{qtsl\_dabs\_posterior\_bracket}(q,\ c,\ lo,\ hi)
}
\end{mathpar}

\par\medskip
\rulename{DtmcPosteriorBracket} is the normalized face of
\rulename{DtmcAbsorbBracket}: a bracket on the \emph{posterior}
eventual-absorption measure of one class, with normalizer the total
classed mass. Hard conditioning is a CLASS, not a score: rejected
mass (failed conditioning) and unlabeled or nonterminating mass are
excluded from the normalizer by the same labeling mechanism, so every
stored row remains a probability row and the chain grain\textquotesingle s
row-mass admission demand is untouched. This is what lets posteriors
of discrete probabilistic programs ride the species without any
weighted-row machinery.

The claim holds iff $lo \le N_c/(S+R)$ and $(N_c+R)/(S+R) \le hi$,
with $N_c$, $S$, $R$ recomputed by the same $T$-step exact-rational
evaluation as the sibling; it REFUSES at $S = 0$ (no
positive-evidence witness). The mathematical content beyond the
sibling\textquotesingle s conservation bracket is the correlated-ratio
lemma --- pure algebra, no chain in sight: from $N_c \le M_c$
classwise, $\sum_c (M_c - N_c) \le R$, and $S > 0$, the true
posterior lies in $[\,N_c/(S+R),\ (N_c+R)/(S+R)\,]$. The correlation
is the point: numerator and denominator draw their errors from ONE
budget, so both endpoints use the single normalizer $S + R$ rather
than the looser interval quotient --- never worse, often sharper
(measured $3.0\times$ tighter than the interval-quotient baseline at
matched truncation on the species\textquotesingle{} first worked
example). Mechanized as \texttt{dtmcPosteriorBracket\_sound}, over
any finite class alphabet, in the same module and with the same
receipts as the sibling.

An honesty note carried from the design: the bracket width is exactly
$R/(S+R)$, class-uniform --- the tightness of the whole judgment is
one number, which is what keeps comparison-table claim shapes honest.
The posterior face deliberately keeps the UNIFORM residual $R$ even
though the unnormalized sibling mints with the reachability-sharpened
$R_c$: the single-normalizer argument requires one shared budget
across all classes, and the class-uniform width claim stays honest
only under it. An asymmetric refinement (uniform lower endpoint,
class-specific upper endpoint $(N_c+R_c)/(S+R_c)$) is sound and is
deliberately NOT minted --- it boards a future revision of the
theory if and when a use case demands it.

\cataloguenote{The schema is baked in the certificate theory alongside the sibling:
same query sort, labelings, and well-formedness demand; its own claim
relation with its own interval-order report; the $S > 0$ refusal is
evaluator-side (a recompute refusal, not a stored row). Quoted below
from the theory source. No Lean statement is bound into this book
yet --- the binding column is honest open debt, as the sibling.
}
\ruleblock{Theory rows (geolog)}
\sourcenote{The rule\textquotesingle s certificate content as it is declared in the geolog theory the checker loads --- payload schema, side-condition demands, and axiom rows, quoted from the theory source at every build of this catalogue (a stale quotation cannot exist). Axiom modes: \texttt{chase} = the engine derives; \texttt{report} = a violation is a loud judgment; \texttt{refuse} = a violation gates admission. Side-condition relations are discharged by bounded evaluation at the engine\textquotesingle s ground seam, never by derivation.}
\par\medskip
The baked rows beyond the sibling\textquotesingle s: the posterior
claim relation and its interval-order report. The query sort,
labelings, and well-formedness demand are the sibling\textquotesingle s
own rows, quoted in its section.

\declhead{theory DtmcAbsorb}{qtsl\_dabs\_posterior\_bracket}{\badge{BadgeGreen}{quoted from the theory source at render}}
\begin{leansrc}
qtsl_dabs_posterior_bracket : {q: DAbsQuery, class: Nat, lo: Q,
                               hi: Q} -> Prop;   // side-condition
\end{leansrc}
\declhead{theory DtmcAbsorb}{ax/wf/dabs\_posterior\_order}{\badge{BadgeGreen}{quoted from the theory source at render}}
\begin{leansrc}
report ax/wf/dabs_posterior_order : forall q : DAbsQuery, c : Nat,
    l : Q, h : Q.
  {q: q, class: c, lo: l, hi: h} qtsl_dabs_posterior_bracket
  |- {l: l, r: h} q_le;
\end{leansrc}
\ruleblock{Lean statement}
\badge{BadgeAmber}{not yet mechanized} this rule cites no Lean statements: nothing on this page is hash-bound for it, and its dashboard row shows the same absence across the statement, receipt, and correspondence columns.
\ruleblock{HOL4 statement}
\absentbox{NO HOL STATEMENT YET}{No HOL development of the extension theory exists; this
species\textquotesingle{} content is an exact-rational evaluator
recompute over baked rows plus one algebraic lemma, mechanized in
the Lean development named in the commentary.
}
\ruleblock{HOL verifier receipts}
\absentbox{NONE}{no HOL verification receipt exists for this rule.}
\ruleblock{Proof}

\sourcenote{Genre: certificate species (tier 1 semantic soundness;
checker recomputation, no proof node). The content beyond the
unnormalized sibling is one algebraic lemma, proved here in full;
the mechanized form is \texttt{dtmcPosteriorBracket\_sound}, over
any finite class alphabet.}

\textbf{Setting.}\enspace As in \rulename{DtmcAbsorbBracket}: from
the $T$-step evaluation, per-class truncation masses $N_c \ge 0$,
their total $S = \sum_c N_c$, and the residual $R \ge 0$. Write
$M_c$ for the true eventual class-$c$ mass (any completion-valued
reading). The sibling's conservation bracket gives, classwise,
$N_c \le M_c$ and $\sum_c (M_c - N_c) \le R$.

\textbf{Lemma (correlated ratio).}\enspace If additionally $S > 0$,
then for every class $c$:
$$\frac{M_c}{\sum_{c'} M_{c'}} \;\in\;
\Bigl[\,\frac{N_c}{S+R},\ \frac{N_c+R}{S+R}\,\Bigr].$$
\emph{Proof.} Let $\delta_{c'} := M_{c'} - N_{c'} \ge 0$ and
$\Delta := \sum_{c'} \delta_{c'} \le R$. The denominator is
$\sum_{c'} M_{c'} = S + \Delta \in [S,\ S+R]$, and $S > 0$ makes
every quotient well-defined. Upper endpoint:
$$\frac{M_c}{S+\Delta} = \frac{N_c + \delta_c}{S + \Delta}
\le \frac{N_c + \Delta}{S + \Delta} \le \frac{N_c + R}{S + R},$$
the first step since $\delta_c \le \Delta$, the second since
$x \mapsto (N_c + x)/(S + x)$ is nondecreasing for $x \ge 0$ (as
$N_c \le S$) and $\Delta \le R$. Lower endpoint:
$$\frac{M_c}{S+\Delta} \ge \frac{N_c}{S+\Delta} \ge
\frac{N_c}{S+R}.$$
\hfill$\square$

The correlation is the content: numerator and denominator draw their
errors from the ONE budget $\Delta \le R$, so both endpoints use the
single normalizer $S + R$. The uncorrelated alternative --- bracket
numerator and denominator separately and take the interval quotient
$[\,N_c/(S+R),\ (N_c+R)/S\,]$ --- is never better and usually worse;
at matched truncation on the species' first worked example the
correlated form measured $3.0\times$ tighter.

\textbf{Why each hypothesis.}\enspace $S > 0$: at $S = 0$ the
posterior is $0/0$-shaped and the species REFUSES (no
positive-evidence witness) --- a refusal, never a verdict.
Nonnegativity of the $\delta_{c'}$ is exactly the sibling's monotone
lower side; the one-budget bound $\Delta \le R$ is its conservation
upper side, summed. Hard conditioning enters as a CLASS (rejected
mass leaves the normalizer by labeling), so every stored row remains
a probability row and no weighted-row machinery is consumed.

\textbf{Width honesty.}\enspace Both endpoints differ by exactly
$R/(S+R)$, class-uniformly: the tightness of the whole judgment is
one number. The posterior face deliberately keeps the uniform $R$
(the one-budget argument above requires a single normalizer shared
by all classes); the asymmetric per-class refinement with upper
endpoint $(N_c+R_c)/(S+R_c)$ is sound and deliberately not part of
this revision --- it boards a future one if a use case demands it.

\begin{bremark}[status]
The schema is baked (claim relation and interval-order report quoted
in this section; query sort and labelings in the sibling's). The
lemma above is mechanized sorry-free at standard axiom closure over
any finite class alphabet. No Lean statement is bound into this book
yet: the binding column is honest open debt, as the sibling's.
\end{bremark}


\rulesection{96}{QGenFormContraction}{core rule}
\label{blcatrule-QGenFormContraction}
\ruleblock{The rule (Steele metanotation)}
\sourcenote{as realized in the mechanized development; source row: \texttt{\texttt{QTSL-rules.toml}} --- \emph{QGenFormContraction}. The metanotation statement is typeset by hand for notational quality; this rule cites no Lean statements yet (the absence is stated loudly below). The rule \emph{roster} is machine-checked (every mechanized rule tag round-trips to exactly one catalogue row); the statement \emph{text} is editorial.}
\begin{mathpar}
\inferrule*{
  w : \mathsf{DGenContraction} \\
  \text{stored rows: } \mathsf{dweight}\ (A = \mathrm{supp}\,V,\ \text{nonempty}),\ \mathsf{dprob},\ \mathsf{dinit} \\
  \textstyle\sum_{y \in A} K(x,y)\,V(y) \le c\,V(x)\ \ \forall x \in A \\
  K(x,y) = 0\ \text{for stored rows } x \notin A,\ y \in A \\
  \min_A V > 0 \\
  \mathsf{qtsl\_dtmc\_genform\_contraction}(w)
}{
  \text{each claimed row of } w:\ \mathsf{qtsl\_dgc\_partial\_lo},\ \mathsf{qtsl\_dgc\_ev\_bracket},\ \mathsf{qtsl\_dgc\_survival\_tail},\ \mathsf{qtsl\_dgc\_z\_one}
}
\end{mathpar}

\par\medskip
\rulename{QGenFormContraction} strengthens the chain-content
generator-form seam from the incumbent\textquotesingle s boolean
claim --- transfer of the certificate is nonpositive-drift --- to a
score-carrying, regional, \emph{quantitative} contraction whose
conclusions include the incumbent\textquotesingle s and add geometric
tails, first-moment tails, and a termination-mass-one claim. It is a
certificate species, not a proof-node rule: checker recomputation
over stored rows, zero new judgment sorts, zero new formers. The
mechanized carrier is an arbitrary nonnegative weight matrix (the
score factorization $s(x)K(x,y)$ is presentation, not mathematics),
so stochastic, sub-probability, and soft-score row sets ride ONE
induction.

This is the fold-don\textquotesingle t-sibling reading of the change
discipline enacted: the would-be separate weighted-tail species folds
into the incumbent seam, and the incumbent survives as the
degenerate face (score $\equiv 1$, $c = 1$, region = all states) ---
its rows, wire shape, and consuming lemmas untouched. Subsumption is
mechanized twice over, as the strengthening discipline demands: the
instance chain (the weighted core at $s \equiv 1$ on a stochastic
kernel re-derives the previously proven stochastic core\textquotesingle s
exact statement:
\texttt{massAt\_expect\_le\_of\_contraction\_of\_weighted}) and the
composed theorem
(\texttt{qGenFormContraction\_subsumes\_genform}: from the
incumbent\textquotesingle s recomputed nonpositive generator form,
the strengthened engine derives the incumbent
consumption\textquotesingle s exact path-law conclusion, never by
citing the incumbent\textquotesingle s own proof) --- so every
certificate admitted under the incumbent discharges under the
strengthening with an unchanged conclusion.

The evaluator recomputes, per stored row: the region-restricted
weighted contraction $\sum_{y \in A} K(x,y)V(y) \le c\,V(x)$ on the
region $A$; no re-entry ($K(x,y) = 0$ for stored rows from outside
$A$ into $A$ --- equivalently, the complement is closed); and the
strictly positive region minimum $\mu = \min_A V > 0$. Conclusions,
each bound exactly where its claim row is spent: weighted alive-mass
decay $\langle m_t|_A, V\rangle \le c^t \langle m_0|_A, V\rangle$;
for $c < 1$, window exit-weight truncation and the time-weighted
first-moment tail (the expected-value rows\textquotesingle{}
engine); the exact partial survival sum (the lower claim); and ---
through the DERIVED conversion $1/\mu$ rescaling the mechanized
unit-floor lemma by scale-invariance --- the bare alive-mass tail
$\mathrm{mass}_t(A) \le c^t \langle m_0|_A, V\rangle/\mu$, the
stored-row form of a hitting-time tail
$\Pr[\text{hitting time} > t] \le M c^t$ (a self-stabilization
time\textquotesingle s region is the not-yet-stabilized
configurations, whose complement\textquotesingle s closedness is
exactly the no-re-entry recompute).

The enacted schema\textquotesingle s deviations from the drafted
proposal are each load-bearing. There is NO stored exit or
conversion constant: the constant is derived ($1/\min_A V$, with the
minimum demanded strictly positive) because a stored constant would
be a forgeable derived quantity --- the species\textquotesingle{}
own never-store-what-you-can-recompute discipline. The region IS the
weight relation\textquotesingle s support (a stored weight row is
the stored membership; no separate region relation to fall out of
sync). Four NAMED claim relations replace a flavor datum --- partial
lower bound, expected-exit-time bracket, survival tail, and
termination-mass-one --- making the wire self-documenting and
letting each side condition bind exactly where its conclusion is
spent (the $c < 1$ demand binds exactly at tail-claiming rows; at
$c = 1$ the incumbent decay-only conclusion remains). A PRESENT
score field REFUSES at this revision (deferred-safe-fail): the
mechanized core is already score-carrying, and the evaluator face
for present scores boards its own future revision rather than
shipping unexercised. Higher moments ($k \ge 2$) are deliberately
absent --- the same per-term reduction with $\sum t^k c^t$ closed
forms, filed to the first leg that needs them.

The soundness engine is the weighted-transfer induction (one-step
transfer, induction on $t$, geometric series --- no stochastic
analysis anywhere), mechanized sorry-free at standard axiom closure
and axiom-audited on every build
(\texttt{SuFunctor/Adequacy/DtmcContraction.lean}): \texttt{wmassAt}
with the instance tie \texttt{massAt\_eq\_wmassAt}; the core
\texttt{wmassAt\_expect\_le\_of\_contraction} (certificate
nonnegativity DROPPED --- strictly stronger than the drafted
specification); the score face; the regional packaging via the
region-supported observable; exit domination; the exact
arithmetico-geometric identity and series tails; and the two tail
assemblies, in universally-quantified partial-sum form (over $\Q$
the limit need not exist; any completion-valued reading inherits the
bound). The alive-mass face lives on its own module
(\texttt{SuFunctor/Adequacy/DtmcAliveTail.lean}):
\texttt{massAt\_expect\_on\_le\_of\_regionContraction} and
\texttt{dtmcAliveTail\_of\_regionContraction}, whose two added
hypotheses --- the floor and the closed complement --- are each
witnessed load-bearing by kernel-checked countermodels COMMITTED
beside the lemmas (\texttt{aliveTail\_needs\_floor},
\texttt{aliveTail\_needs\_no\_reentry}), re-checked on every build.
The species\textquotesingle{} mutant battery was run as committed
throwaway edits, each red at its predicted load-bearing proof step.

Two further faces are deliberately NOT part of this revision. The
TEMPLATE face (geometric certificate families $V(n) = M\gamma^n$
over offset families, collapsing region instances to single rational
inequalities) and the PATTERN face (finitely-presented infinite
chains and product-profile templates) are deferred to their own
design work: the pattern face is underdetermined at this grain (no
canonical state-coordinate binding; a pattern\textquotesingle s
generated kernel and the stored completion\textquotesingle s
rowless-is-absorbing reading are two different kernels and must not
share one chain reading; and the finite-state induction does not by
itself warrant infinite-sum exchange --- a carrier theorem is owed),
and no admitted certificate demands the one-dimensional template at
the stored-weights grain, where the finite face stores every weight.
Their collapse warrants are already mechanized and banked for that
future work (\texttt{genFormTemplate\_collapse} for one-dimensional
offset families; \texttt{aliveStep\_factor},
\texttt{transferRatio\_of\_ge}, \texttt{productProfile\_collapse}
and its window form,
\texttt{SuFunctor/Adequacy/DtmcProductProfile.lean}, proven at full
profile generality), together with a kernel-checked necessity
witness (\texttt{rampProfile\_needs\_window\_reach}) showing the
class-check family alone cannot cover ramp-resident exterior states
--- so the face\textquotesingle s eventual schema must carry a
profile-length demand beside window coverage. What remains for that
face is the infinite-carrier bridge, the kernel-object tie, and the
stored/template composition glue; until then, template- and
pattern-shaped claims live outside the certified lane, labeled as
such.

\cataloguenote{The FINITE FACE is baked in the certificate theory as a species
sibling extending the chain grain, purely additive: the witness
entity (anchor-mediated chain binding, absolute horizon, the
deterministic-prologue offset datum, the contraction constant, the
optional score reference whose presence refuses at this revision),
the keyed weight relation whose support is the region (nonempty
demanded), the engine-decided demand minted per witness by a report
row, and the four named claim relations with their range and
interval-order reports --- all quoted below from the theory source.
The incumbent generator-form rows are untouched (additive delta
only), and the sibling species\textquotesingle{} deliberately-absent
notes carry their pointer flips from the same revision. The
mechanized soundness lemmas are named in the commentary; no Lean
statement is bound into this book yet --- the binding column is
honest open debt.
}
\ruleblock{Theory rows (geolog)}
\sourcenote{The rule\textquotesingle s certificate content as it is declared in the geolog theory the checker loads --- payload schema, side-condition demands, and axiom rows, quoted from the theory source at every build of this catalogue (a stale quotation cannot exist). Axiom modes: \texttt{chase} = the engine derives; \texttt{report} = a violation is a loud judgment; \texttt{refuse} = a violation gates admission. Side-condition relations are discharged by bounded evaluation at the engine\textquotesingle s ground seam, never by derivation.}
\par\medskip
The incumbent seam this species strengthens --- and survives as the
degenerate face of: the chain-content generator-form demand.

\declhead{theory Dtmc}{qtsl\_dtmc\_genform}{\badge{BadgeGreen}{quoted from the theory source at render}}
\begin{leansrc}
qtsl_dtmc_genform : {sys: System, cert: HypernetRef,
                     gen: HypernetRef} -> Prop;
\end{leansrc}
\par\medskip
The baked rows of the finite face: the witness entity, the weight
relation with its same-chain chase row, the engine-decided demand,
and the four named claim relations with their wf reports (constant
range, interval order on expected-value brackets, tail-coefficient
ranges, and the strict-contraction demand exactly where an
expected-value bracket is claimed).

\declhead{theory DtmcContraction}{DGenContraction}{\badge{BadgeGreen}{quoted from the theory source at render}}
\begin{leansrc}
DGenContraction : Set;
\end{leansrc}
\declhead{theory DtmcContraction}{DGenContraction/*}{\badge{BadgeGreen}{quoted from the theory source at render}}
\begin{leansrc}
DGenContraction/anchor  : DGenContraction -> DState;
DGenContraction/horizon : DGenContraction -> Nat;
DGenContraction/offset  : DGenContraction -> Nat;
DGenContraction/c       : DGenContraction -> Q;
DGenContraction/score   : DGenContraction -> HypernetRef;
\end{leansrc}
\declhead{theory DtmcContraction}{dweight}{\badge{BadgeGreen}{quoted from the theory source at render}}
\begin{leansrc}
dweight : {cw: DGenContraction, state: DState} -> Q;
\end{leansrc}
\declhead{theory DtmcContraction}{ax/wf/dgc\_weight\_same\_dtmc}{\badge{BadgeGreen}{quoted from the theory source at render}}
\begin{leansrc}
chase ax/wf/dgc_weight_same_dtmc : forall w : DGenContraction,
    s : DState, v : Q.
  v = {cw: w, state: s} dweight
  |- s DState/dtmc = w DGenContraction/anchor DState/dtmc;
\end{leansrc}
\declhead{theory DtmcContraction}{qtsl\_dtmc\_genform\_contraction}{\badge{BadgeGreen}{quoted from the theory source at render}}
\begin{leansrc}
qtsl_dtmc_genform_contraction : {cw: DGenContraction} -> Prop;
\end{leansrc}
\declhead{theory DtmcContraction}{ax/wf/dgc\_demand}{\badge{BadgeGreen}{quoted from the theory source at render}}
\begin{leansrc}
report ax/wf/dgc_demand : forall w : DGenContraction.
  |- {cw: w} qtsl_dtmc_genform_contraction;
\end{leansrc}
\declhead{theory DtmcContraction}{qtsl\_dgc\_partial\_lo}{\badge{BadgeGreen}{quoted from the theory source at render}}
\begin{leansrc}
qtsl_dgc_partial_lo : {cw: DGenContraction, lo: Q} -> Prop;
\end{leansrc}
\declhead{theory DtmcContraction}{qtsl\_dgc\_ev\_bracket}{\badge{BadgeGreen}{quoted from the theory source at render}}
\begin{leansrc}
qtsl_dgc_ev_bracket : {cw: DGenContraction, lo: Q, hi: Q} -> Prop;
\end{leansrc}
\declhead{theory DtmcContraction}{qtsl\_dgc\_survival\_tail}{\badge{BadgeGreen}{quoted from the theory source at render}}
\begin{leansrc}
qtsl_dgc_survival_tail : {cw: DGenContraction, coeff: Q, rate: Q}
  -> Prop;   // side-condition (R21)
\end{leansrc}
\declhead{theory DtmcContraction}{qtsl\_dgc\_z\_one}{\badge{BadgeGreen}{quoted from the theory source at render}}
\begin{leansrc}
qtsl_dgc_z_one : {cw: DGenContraction} -> Prop;
\end{leansrc}
\declhead{theory DtmcContraction}{ax/wf/dgc\_c\_range}{\badge{BadgeGreen}{quoted from the theory source at render}}
\begin{leansrc}
report ax/wf/dgc_c_range : forall w : DGenContraction, c : Q.
  c = w DGenContraction/c
  |- {l: 0, r: c} q_le, {l: c, r: 1} q_le;
\end{leansrc}
\declhead{theory DtmcContraction}{ax/wf/dgc\_ev\_order}{\badge{BadgeGreen}{quoted from the theory source at render}}
\begin{leansrc}
report ax/wf/dgc_ev_order : forall w : DGenContraction, l : Q, h : Q.
  {cw: w, lo: l, hi: h} qtsl_dgc_ev_bracket
  |- {l: l, r: h} q_le;
\end{leansrc}
\declhead{theory DtmcContraction}{ax/wf/dgc\_tail\_range}{\badge{BadgeGreen}{quoted from the theory source at render}}
\begin{leansrc}
report ax/wf/dgc_tail_range : forall w : DGenContraction, b : Q,
    r : Q.
  {cw: w, coeff: b, rate: r} qtsl_dgc_survival_tail
  |- {l: 0, r: b} q_le, {l: 0, r: r} q_le, {l: r, r: 1} q_lt;
\end{leansrc}
\declhead{theory DtmcContraction}{ax/wf/dgc\_ev\_strict}{\badge{BadgeGreen}{quoted from the theory source at render}}
\begin{leansrc}
report ax/wf/dgc_ev_strict : forall w : DGenContraction, l : Q,
    h : Q, c : Q.
  {cw: w, lo: l, hi: h} qtsl_dgc_ev_bracket,
  c = w DGenContraction/c
  |- {l: c, r: 1} q_lt;
\end{leansrc}
\ruleblock{Lean statement}
\badge{BadgeAmber}{not yet mechanized} this rule cites no Lean statements: nothing on this page is hash-bound for it, and its dashboard row shows the same absence across the statement, receipt, and correspondence columns.
\ruleblock{HOL4 statement}
\absentbox{NO HOL STATEMENT YET}{No HOL development of the extension theory exists; this
species\textquotesingle{} content is an exact-rational evaluator
recompute over baked rows, and its soundness theorems live in the
Lean development named in the commentary.
}
\ruleblock{HOL verifier receipts}
\absentbox{NONE}{no HOL verification receipt exists for this rule.}
\ruleblock{Proof}

\sourcenote{Genre: certificate species (tier 1 semantic soundness;
checker recomputation strengthening an existing seam). The proof is
the weighted-transfer induction plus exact series arithmetic --- no
stochastic analysis anywhere; the mechanized lemma roster is named
inline.}

\textbf{Setting.}\enspace A finite stored-row chain with completed
kernel $K$ (rowless-is-absorbing, as in the sibling species); a
stored weight function $V$ on the region $A := \mathrm{supp}\,V$
(the region IS the support of the stored weight rows --- nonempty
demanded); a contraction constant $c \in [0,1]$; and the recomputed
side conditions: regional contraction
$\sum_{y \in A} K(x,y)\,V(y) \le c\,V(x)$ for every $x \in A$;
no re-entry, $K(x,y) = 0$ for stored rows with $x \notin A$,
$y \in A$ (equivalently, the complement of $A$ is closed); and
strictly positive region minimum $\mu := \min_A V > 0$. The carrier
of the mechanized core is an arbitrary nonnegative weight matrix ---
the score factorization $s(x)K(x,y)$ is presentation, not
mathematics --- so stochastic, sub-probability, and soft-score row
sets ride one induction; certificate nonnegativity is NOT assumed by
the core (\texttt{wmassAt\_expect\_le\_of\_contraction}), which is
strictly stronger than the drafted specification.

\textbf{Lemma 1 (one-step transfer).}\enspace With $w_t :=
\langle m_t|_A,\ V\rangle$ the region-supported weighted mass:
$w_{t+1} \le c\, w_t$. \emph{Proof.} Expand $m_{t+1} = m_t K$;
mass entering $A$ from outside contributes nothing (no re-entry),
and each $x \in A$ contributes $m_t(x) \sum_{y \in A} K(x,y)V(y)
\le c\, m_t(x) V(x)$ by the recomputed contraction, using
nonnegativity of kernel entries and of the evolving mass (the two
mutant loci at this step). Sum over $x \in A$.\hfill$\square$

\textbf{Theorem (decay, tails, termination).}\enspace By induction,
$w_t \le c^t w_0$ (weighted alive-mass decay --- at $c = 1$ this is
exactly the incumbent seam's nonpositive-drift conclusion). Per-step
exit weight is dominated through the region values (exit
domination), so for $c < 1$ and any window $[T, T')$, the geometric
series gives the truncation bound
$\kappa\, c^T \langle m_0, V\rangle/(1-c)$ and --- via the exact
arithmetico-geometric identity
(\texttt{sum\_range\_mul\_pow\_mul\_eq}) and its tail estimates
(\texttt{sum\_Ico\_pow\_le}, \texttt{sum\_Ico\_mul\_pow\_le}) --- the
time-weighted first-moment tail
$\kappa\, \langle m_0, V\rangle\, c^T (T(1-c)+c)/(1-c)^2$, both in
universally-quantified partial-sum form: over $\Q$ the limit need
not exist, and any completion-valued normalizer, expected-value
numerator, or moment row inherits the bound
(\texttt{wmassAt\_exit\_tail\_le},
\texttt{wmassAt\_exit\_moment\_tail\_le}). Dropping $c < 1$ from a
tail claim is unsound, not vacuous --- at $c = 1$ the rational
division convention would make the assembled bound false (the
battery's witness) --- which is why the strict-contraction demand
binds exactly where a tail-shaped claim row is spent.

\textbf{Corollary (bare alive-mass tail --- the derived
conversion).}\enspace Where $\mu = \min_A V > 0$, apply the
unit-floor lemma to the rescaled certificate $V/\mu$ (contraction
and no-re-entry are scale-invariant; the floor $V/\mu \ge 1$ on $A$
holds by definition of $\mu$):
$$\mathrm{mass}_t(A) \;\le\; c^t\, \langle m_0|_A,\ V\rangle/\mu$$
(\texttt{massAt\_expect\_on\_le\_of\_regionContraction},
\texttt{dtmcAliveTail\_of\_regionContraction}) --- the stored-row
form of the hitting-time tail $\Pr[\text{hitting time} > t] \le
M c^t$. The conversion constant $1/\mu$ is DERIVED, never stored: a
stored constant would be a forgeable derived quantity. Both added
hypotheses are witnessed load-bearing by kernel-checked
countermodels committed beside the lemmas and re-checked every
build: without the floor, one state at $V = 1/2$, $c = 1$ makes the
bound read $1 \le 1/2$ (\texttt{aliveTail\_needs\_floor}); without
the closed complement, two states with the exterior feeding the
region at probability one make it read $1 \le 0$ at $t = 1$
(\texttt{aliveTail\_needs\_no\_reentry}). With stochastic rows at
every region state, an initial probability measure, $0 \le c < 1$,
no re-entry, and the positive minimum, weighted decay converts to
bare survival decay and forces termination mass exactly one --- the
fourth claim relation; a sign-free weighted contraction alone would
NOT suffice (the negative-weight self-loop countermodel is a battery
tooth).

\textbf{Subsumption (the strengthening's extra obligation).}\enspace
Mechanized twice over, as a strengthening of a landed seam must be:
the instance chain --- the weighted core at $s \equiv 1$ on a
stochastic kernel re-derives the previously proven stochastic core's
exact statement
(\texttt{massAt\_expect\_le\_of\_contraction\_of\_weighted}) --- and
the composed theorem
(\texttt{qGenFormContraction\_subsumes\_genform}): from the
incumbent's recomputed nonpositive generator form, the strengthened
engine at the degenerate face (score $\equiv 1$, $c = 1$, region =
all states) derives the incumbent consumption's exact path-law
conclusion through the duality bridge, never by citing the
incumbent's own proof. Every certificate admitted under the
incumbent therefore discharges under the strengthening with an
unchanged conclusion; the incumbent's rows, wire shape, and
consuming lemmas are untouched.

\textbf{Prologue offsets.}\enspace The offset datum $k$ (absent = 0)
excludes a deterministic prologue from the claim window: partial
survival sums index from $k$, and where a partial or expected-value
claim spends it, the well-formedness demand recomputes prologue
determinism --- killed-on-exit evolution from the full stored
initial measure with total surviving mass exactly one at step
$k - 1$ ($k = 0$ vacuous). An unearned offset would shift the tail
window onto mass the induction never controlled; the recompute
refuses it.

\begin{bremark}[status]
The finite face's schema is baked (entity, weight relation, demand,
and the four named claim relations quoted in this section); the
whole species' soundness, the subsumption pair, and the alive-mass
face are mechanized sorry-free at standard axiom closure,
axiom-audited on every build, with the species' mutant battery run
as committed throwaway edits, each red at its predicted load-bearing
step. No Lean statement is bound into this book yet: the binding
column is honest open debt. The template and pattern faces are NOT
part of this revision --- their collapse warrants and the necessity
witness for the eventual profile-length demand are mechanized and
banked, and the faces board the revision that follows their own
design work; until then template- and pattern-shaped claims live
outside the certified lane, labeled as such.
\end{bremark}


\blcatsection{The Soundness Theorem}
{\itshape Its content is three conformance obligations on derivation edges, and together they are what turns the page\textquotesingle s per-rule \emph{relative} soundness lemmas --- each proved on the hypothesis that the base calculus is sound --- into a single verdict about a whole certificate, unconditional once those per-rule lemmas, base soundness, and the loader\textquotesingle s conformance passes are supplied.}

\rulesection{97}{ProofWellFormedness}{edge discipline (glue)}
\label{blcatrule-ProofWellFormedness}
\ruleflavor{baked as \emph{typing}, not as a node sort: E1 — derivation premises are judgment kinds XOR inherited $\mathsf{Proof}$, never \texttt{QProof}, never a rule sort (the baked proof-node sorts\textquotesingle{} premise fields carry exactly these types); E2 — certificate-data references are not derivation edges; E3 — the admittedness + acyclicity pass over the certificate\textquotesingle s own row graph (loader conformance, deliberately a check rather than a theory axiom: a cyclic certificate would chase-satisfy every axiom while proving nothing).}
\ruleblock{Lean statement}
\par\medskip
The discipline, mechanized at exactly the certificate-object level
the recipes state it. The certificate shape binds by display --- the four definitions
ARE the E1/E2 typing claims --- and the theorems bind as statements:
the E3 bridge, the induction template in well-founded and
finite-acyclic forms with root corollary, and the consumption-site
transport pair. The induction theorem proper
(\texttt{towerCitation\_sound\_of\_wf}) closes under no axioms at all.

\declhead{BernsteinLogicAux.Acceptance}{acyclic\_wellFounded}{\badge{BadgeGreen}{statement bound}}
\begin{leansrc}
/-- **The E3 bridge**: on a FINITE certificate, the loader's checkable
acyclicity pass yields exactly the well-foundedness the induction needs.
This is why E3 can be loader conformance rather than an axiom of the
theory: well-foundedness is an induction principle the logical framework
deliberately lacks, but on finite row sets it is equivalent to the
decidable absence of citation cycles. -/
theorem acyclic_wellFounded [Finite ι] (C : ι → Row σ J ι)
    (hacyc : Acyclic C) : WellFounded (cites C)
\end{leansrc}
\declhead{BernsteinLogicAux.Acceptance}{towerCitation\_sound\_of\_wf}{\badge{BadgeGreen}{statement bound}}
\begin{leansrc}
/-- **`rec:tc-template`, well-founded form** — the extension-edge
induction, in the shape the reflection theorem's induction consumes.
Hypotheses are the template's clauses: clause 1 = `hbase` (base
soundness; in deployment, the BL checker-soundness frontier — residue
R4); clause 2 = `hside` + `hstep` (the acceptance relation's local-check
clause as its own parameter, then the per-rule conditional-soundness
step in the cookie-cutter shape — the per-rule lemmas are CONSUMED here,
never proven: residue R1); clause 3 = `hwf` (the well-foundedness the E3
pass buys via `acyclic_wellFounded`). Conclusion: every admitted row's
conclusion-denotation holds — at each row the induction hypothesis
supplies exactly what the per-rule lemma consumes, and clause 1 grounds
the leaves. -/
theorem towerCitation_sound_of_wf
    (C : ι → Row σ J ι) (baseTrue : σ → Prop) (sem : J → Prop)
    (SideOK : ι → Prop)
    (hbase : ∀ i s, C i = Row.base s → baseTrue s)
    (hside : ∀ i, SideOK i)
    (hstep : ∀ i prems concl, C i = Row.ext prems concl → SideOK i →
      (∀ j ∈ prems, Row.holds baseTrue sem (C j)) → sem concl)
    (hwf : WellFounded (cites C)) :
    ∀ i, Row.holds baseTrue sem (C i)
\end{leansrc}
\declhead{BernsteinLogicAux.Acceptance}{towerCitation\_sound}{\badge{BadgeGreen}{statement bound}}
\begin{leansrc}
/-- **`rec:tc-template`, headline form** — the discipline as the checker
consumes it: a FINITE certificate whose citation graph passes the E3
acyclicity check, base soundness, local checks, and the per-rule
conditional lemmas jointly warrant every admitted row's
conclusion-denotation. This is the theorem that turns the book's
per-rule RELATIVE soundness lemmas into an unconditional verdict about
a whole certificate. -/
theorem towerCitation_sound [Finite ι]
    (C : ι → Row σ J ι) (baseTrue : σ → Prop) (sem : J → Prop)
    (SideOK : ι → Prop)
    (hbase : ∀ i s, C i = Row.base s → baseTrue s)
    (hside : ∀ i, SideOK i)
    (hstep : ∀ i prems concl, C i = Row.ext prems concl → SideOK i →
      (∀ j ∈ prems, Row.holds baseTrue sem (C j)) → sem concl)
    (hacyc : Acyclic C) :
    ∀ i, Row.holds baseTrue sem (C i)
\end{leansrc}
\declhead{BernsteinLogicAux.Acceptance}{towerCitation\_root}{\badge{BadgeGreen}{statement bound}}
\begin{leansrc}
/-- **Root soundness** (the acceptance relation's clause-4 shape): if
some admitted row of an accepted certificate concludes the root
judgment, that judgment's denotation holds. The abstract skeleton of
the reflection theorem's statement. -/
theorem towerCitation_root [Finite ι]
    (C : ι → Row σ J ι) (baseTrue : σ → Prop) (sem : J → Prop)
    (SideOK : ι → Prop)
    (hbase : ∀ i s, C i = Row.base s → baseTrue s)
    (hside : ∀ i, SideOK i)
    (hstep : ∀ i prems concl, C i = Row.ext prems concl → SideOK i →
      (∀ j ∈ prems, Row.holds baseTrue sem (C j)) → sem concl)
    (hacyc : Acyclic C)
    {Jroot : J} (hroot : ∃ i prems, C i = Row.ext prems Jroot) :
    sem Jroot
\end{leansrc}
\declhead{BernsteinLogicAux.Acceptance}{citePos\_sound}{\badge{BadgeGreen}{statement bound}}
\begin{leansrc}
/-- **`prop:tc-consume` (a), positive citations**: the child's concluded
region `R` under-approximates its formula's denotation `D`, and any
MONOTONE consuming claim proved over `R` transfers to `D`. Which
operators are monotone in which positions is each consuming rule's own
lemma's business; this lemma fixes only the direction. A certificate can
lose completeness at a citation, never soundness. -/
theorem citePos_sound {X : Type*} {R D : Set X} {P : Set X → Prop}
    (hmono : Monotone P) (hsub : R ⊆ D) (hR : P R) : P D
\end{leansrc}
\declhead{BernsteinLogicAux.Acceptance}{citeCmpl\_subset}{\badge{BadgeGreen}{statement bound}}
\begin{leansrc}
/-- **`prop:tc-consume` (b), the orientation swap**: an
under-approximating certified region's complement OVER-approximates the
denotation's complement — the conservative direction where regions
appear negatively. (Set-inclusion twin of
`BernsteinLogic.Structural.compl_over_approx`, which states the same
direction through the compiled region-judgment reading.) -/
theorem citeCmpl_subset {X : Type*} {R D : Set X} (hsub : R ⊆ D) :
    Dᶜ ⊆ Rᶜ
\end{leansrc}
\declhead{BernsteinLogicAux.Acceptance}{citeCmpl\_sound}{\badge{BadgeGreen}{statement bound}}
\begin{leansrc}
/-- **`prop:tc-consume` (b), consumer form**: an ANTITONE claim proved on
the citable complement `Rᶜ` covers the true violation set `Dᶜ` — a
premise row proved on, or a probability bound proved for reaching, the
over-approximation covers the truth. -/
theorem citeCmpl_sound {X : Type*} {R D : Set X} {P : Set X → Prop}
    (hanti : Antitone P) (hsub : R ⊆ D) (hR : P Rᶜ) : P Dᶜ
\end{leansrc}
\declhead{BernsteinLogicAux.Acceptance}{Row}{\badge{BadgeGreen}{definition bound}}
\begin{leansrc}
/-- **E1, as a type**: one row of a decoded certificate. A `base` row is a
base-calculus proof leaf concluding the sequent `s` (the `⊢` side); an
`ext` row is an extension row carrying its premise slots — references to
other rows, the derivation edges, and nothing else — and its concluded
satisfaction judgment (the `⊨` side). A premise slot can land on nothing
but another admitted row, consumed as that row's conclusion. -/
inductive Row (σ J ι : Type*) where
  | base (s : σ)
  | ext (prems : List ι) (concl : J)
\end{leansrc}
\declhead{BernsteinLogicAux.Acceptance}{Row.cited}{\badge{BadgeGreen}{definition bound}}
\begin{leansrc}
/-- **E2, in one definition**: the derivation edges out of a row are its
premise slots and nothing else — the induction metric counts these only.
Certificate-data references (shared certificate families, lattice
objects) are unrepresented here by design; their erasure is the loader's
conformance obligation, not this module's. -/
def Row.cited : Row σ J ι → List ι
  | .base _ => []
  | .ext prems _ => prems
\end{leansrc}
\declhead{BernsteinLogicAux.Acceptance}{Row.holds}{\badge{BadgeGreen}{definition bound}}
\begin{leansrc}
/-- The denotation of a row's CONCLUSION, under a kernel semantics
`baseTrue` for base sequents and a judgment semantics `sem` (the
judgment-denotation interface, kept abstract — zero measure theory in
this module by design). The two-case shape is E1's XOR: an edge lands on
kernel truth or on a judgment denotation; there is no third thing to
cite. -/
def Row.holds (baseTrue : σ → Prop) (sem : J → Prop) : Row σ J ι → Prop
  | .base s => baseTrue s
  | .ext _ concl => sem concl
\end{leansrc}
\declhead{BernsteinLogicAux.Acceptance}{cites}{\badge{BadgeGreen}{definition bound}}
\begin{leansrc}
/-- Row `j` is cited by row `i`: a derivation edge from `i` down to `j`.
The well-founded induction descends along this relation. -/
def cites (C : ι → Row σ J ι) (j i : ι) : Prop :=
  j ∈ (C i).cited
\end{leansrc}
\declhead{BernsteinLogicAux.Acceptance}{Acyclic}{\badge{BadgeGreen}{definition bound}}
\begin{leansrc}
/-- **E3's checkable half**, as the induction consumes it: no row
transitively cites itself. On a finite certificate this is precisely
acyclicity of the citation graph — the decidable pass the loader runs. -/
def Acyclic (C : ι → Row σ J ι) : Prop :=
  ∀ i, ¬ Relation.TransGen (cites C) i i
\end{leansrc}
\ruleblock{Correspondence (Steele display $\leftrightarrow$ Lean)}
\absentbox{NO CORRESPONDENCE DICTIONARY YET}{no reviewer-certified dictionary maps the displayed rule onto the mechanized statements above; until one is written and certified, that mapping rests on this section\textquotesingle s editorial care alone.}
\ruleblock{Lean verifier receipts}
\sourcenote{module \texttt{BernsteinLogicAux.Acceptance}; bound build \texttt{e78746b9f98a\ldots} (bound from receipt 3d1fca564d6e11dc, 2026-07-20T03:41:07Z). A receipt records that named checking layers accepted a build: (module, compiled-artifact SHA-256, environment SHA-256, checker layers), minted when the checking gate ran; the store receipts attesting the bound build are listed with the layers each run credited. Nothing here asks the reader to trust the store: every layer is an independently runnable program over the quoted sources --- rebuild the module and re-run the checkers to reproduce the verdict. Receipts attesting superseded builds credit no layer and are omitted from this rendering (the count of omissions, if any, is noted below).}
\begin{receipttable}
\receiptrow{3d1fca564d6e11dc}{2026-07-20T03:41:07Z}{build+decorrelated}{\badge{BadgeGreen}{current}}
\receiptrow{3546fbecf9c242a2}{2026-07-20T20:46:26Z}{build+decorrelated}{\badge{BadgeGreen}{current}}
\receiptrow{e684b09bdf01a640}{2026-07-20T20:57:59Z}{build+replay+hardened+decorrelated}{\badge{BadgeGreen}{current}}
\receiptrow{fe82a04ec14e7cee}{2026-07-21T00:02:22Z}{build+decorrelated}{\badge{BadgeGreen}{current}}
\receiptrow{f0f540e41c59285c}{2026-07-21T05:38:30Z}{build+replay+hardened+decorrelated}{\badge{BadgeGreen}{current}}
\receiptrow{69a2515828782208}{2026-07-22T18:26:47Z}{build+decorrelated}{\badge{BadgeGreen}{current}}
\receiptrow{bd7c9443b7f8649a}{2026-07-31T21:59:38Z}{build+decorrelated}{\badge{BadgeGreen}{current}}
\receiptrow{652943c163b17f25}{2026-07-31T22:28:31Z}{build+decorrelated}{\badge{BadgeGreen}{current}}
\receiptrow{dcd3934b54c0db35}{2026-08-02T23:30:39Z}{build+decorrelated}{\badge{BadgeGreen}{current}}
\receiptrow{5d9475b9d8bbe3c2}{2026-08-03T16:26:43Z}{build+decorrelated}{\badge{BadgeGreen}{current}}
\receiptrow{cbbad707861e8740}{2026-08-03T22:45:26Z}{build+decorrelated}{\badge{BadgeGreen}{current}}
\end{receipttable}
\layerline{\layeritem{build}{lake build (elaboration)}{\badge{BadgeGreen}{current}}\layeritem{replay}{leanchecker (official kernel replay)}{\badge{BadgeGreen}{current}}\layeritem{hardened}{safe\_verify (hardened screen)}{\badge{BadgeGreen}{current}}\layeritem{decorrelated}{lean4export + nanoda (decorrelated kernel)}{\badge{BadgeGreen}{current}}}
\sourcenote{\badge{BadgeGray}{comparator: none} no fixture in the three-implementation comparator harness (compiled / kernel-interpreted / reference-interpreted) exercises this rule.}
\ruleblock{HOL4 statement}
\absentbox{OUT OF SCOPE BY DESIGN}{Out of scope by structure: there is no rule statement to port --- the
discipline concludes nothing and has no check function of its own. Its
one checkable obligation, the admittedness + acyclicity pass, is
certificate-loader conformance, which would arrive with a HOL-side
extension-theory loader as conformance code, not as a rule check.
Independently, the HOL4 development\textquotesingle s object today is
the base calculus\textquotesingle s wire-format checker (theory
\texttt{bern\_catalogue}, the base catalogue\textquotesingle s
check-function cone); no extension-theory HOL development exists for
such a pass to live in.
}
\ruleblock{HOL verifier receipts}
\absentbox{NONE}{no HOL verification receipt exists for this rule.}
\ruleblock{Proof}

\sourcenote{Genre: edge discipline (glue) --- the extension
page\textquotesingle s fourth category, deliberately outside the three
tiers. This entry is not a rule: it has no conclusion judgment, so
``semantic soundness of the conclusion'' does not type-check against
it, and it has nothing to unpack, so neither rewriting genre applies
--- it \emph{is} the unpacking mechanism. What it owes the book
instead is the warrant every other section of the extension page
silently leans on: each rule\textquotesingle s soundness lemma is
proved \emph{relative} to base soundness, and this discipline is how
those conditional lemmas compose into a single verdict about a whole
certificate --- unconditional once the per-rule lemmas, base
soundness, and the loader\textquotesingle s conformance passes are in
hand; the closing remark names what remains conditional. The presentation apes the recipe form: the general
induction template, then the discipline\textquotesingle s obligations
as the instance that makes the template apply.}

\begin{brecipe}[what the discipline licenses]\label{rec:tc-license}
A concluded child judgment
$\mathcal{M},\ x_0 \in R \vDash \varphi_{\mathrm{child}}$ (the region-level
satisfaction kind: every point of $R$ satisfies
$\varphi_{\mathrm{child}}$, with $R$ under-approximating the
formula\textquotesingle s denotation) makes $R$ \emph{citable} as a
region atom in any parent row\textquotesingle s scope or target. The
working example: a nested probability former
$\mathsf{SProbGE}(p, \psi)$ (surface syntax $\mathsf{P}_{\ge p}[\psi]$)
inside a state formula compiles its child
first and cites the child\textquotesingle s concluded region as the
sublevel atom $\{V_{\mathrm{child}} \le 1 - p\}$. The license is an
\emph{edge} between rows of one certificate, never a node: the
citation itself concludes nothing.
\end{brecipe}

\begin{brecipe}[the induction template]\label{rec:tc-template}
Let a certificate be a finite set of admitted rows: base-calculus
proof rows (leaves, concluding sequents on the $\vdash$ side) and
extension rows (concluding satisfaction judgments on the $\vDash$
side), each extension row\textquotesingle s premise fields referencing
other rows; call a row \emph{admitted} when its local checks pass ---
its premise fields are well typed and its side conditions recompute
true. Suppose:
\begin{enumerate}
\item \emph{base soundness} --- every admitted base row\textquotesingle s
  sequent is true in kernel semantics;
\item \emph{per-rule conditional soundness} --- for every extension
  rule sort $R$ a lemma $\mathrm{sound}_R$: if the referenced child
  judgments\textquotesingle{} denotations hold, the base leaves are
  true, and the side conditions recompute true, then the concluded
  judgment\textquotesingle s denotation holds;
\item \emph{acyclicity} --- the certificate\textquotesingle s citation
  graph is acyclic.
\end{enumerate}
Then every admitted row\textquotesingle s judgment denotation holds:
induct over the citation order, which clause 3 (with finiteness)
makes well-founded; at each row the induction hypothesis supplies the
child denotations, admittedness supplies the recomputed side
conditions, and clause 1 grounds the leaves --- together exactly what
clause 2\textquotesingle s lemma consumes.
The conditional lemmas are the per-rule sections of this book; the
induction is this section.
\end{brecipe}

\begin{brecipe}[the instance --- the discipline\textquotesingle s three
obligations]\label{rec:tc-instance}
The discipline is staged as \emph{typing and conformance}, not as a
node sort, in three obligations, each of which is one premise of
Recipe~\ref{rec:tc-template} made checkable:
\begin{itemize}
\item \textbf{E1 (edge typing).} Every premise field of an extension
  row is typed as a satisfaction-judgment kind or as an inherited
  base-proof leaf, exclusively --- never as a raw rule row of either
  theory. This makes the induction\textquotesingle s case analysis
  exhaustive: every edge lands on something clause 1 or the induction
  hypothesis covers; no edge can cite an object that has no soundness
  account.
\item \textbf{E2 (data/edge separation).} Certificate-\emph{data}
  references --- a certificate family reused across rows, a lattice
  object --- are not derivation edges. The induction metric counts
  derivation edges only, so shared data can neither smuggle a cycle
  past the acyclicity pass nor impose ordering the derivation
  structure does not have.
\item \textbf{E3 (admittedness + acyclicity).} A loader pass over the
  certificate\textquotesingle s own row graph checks that every
  referenced row is admitted and that the citation graph is acyclic
  --- clause 3 verbatim. This is \emph{conformance}, deliberately not
  an axiom of the theory: an existence-of-concluding-proof axiom would
  let the theory\textquotesingle s own saturation fabricate proof
  rows, and well-foundedness is an induction principle, which the
  logical framework deliberately lacks. A cyclic certificate would
  satisfy every local check while proving nothing --- or worse: take
  two rows, each citing the other\textquotesingle s concluded
  judgment as its sole premise, with per-rule lemmas of the shape
  $J_B \Rightarrow J_A$ and $J_A \Rightarrow J_B$. In an instance
  where both judgments are false, both implications hold, every
  per-row obligation and the admittedness pass discharge, and only
  acyclicity objects; accepting the pair would not merely warrant
  nothing --- it would assert falsehoods.
\end{itemize}
\end{brecipe}

\begin{bprop}[consumption-site soundness]\label{prop:tc-consume}
Let the child conclude $\mathcal{M},\ x_0 \in R \vDash \varphi$ with
$R \subseteq [\![\varphi]\!]$. Then at every consumption site the
citation is used in the sound direction:
\begin{enumerate}
\item[(a)] a \emph{positive} citation passes $R$ exactly, and any
  monotone claim over $R$ (``the state reaches / remains in $R$'')
  entails the same claim for $[\![\varphi]\!]$;
\item[(b)] a \emph{complement} citation swaps orientation:
  $\mathrm{cmpl}(R) \supseteq \mathrm{cmpl}([\![\varphi]\!])$, so a
  premise row proved on $\mathrm{cmpl}(R)$ --- or an \emph{upper}
  probability bound proved for reaching it --- covers the true
  violation set.
\end{enumerate}
\end{bprop}
\begin{proof}
(a) is monotonicity of the consuming operator in
$R \subseteq [\![\varphi]\!]$; (b) is the same inclusion complemented.
The discipline contributes exactly this inclusion and its swap;
\emph{which} operators are monotone in which positions is each
consuming rule\textquotesingle s own lemma\textquotesingle s business,
stated and proved in that rule\textquotesingle s section.
Under-approximation on the positive side and over-approximation on the
complement side are both conservative: a certificate can lose
completeness at a citation, never soundness.
\end{proof}

\begin{bremark}[why this is not a rule, and what the mechanization is]
The three tiers each predicate on a conclusion: core is the concluded
judgment\textquotesingle s semantic soundness; admissible and alias
are verified unpackings of a concluded node. A discipline with no
conclusion satisfies none of the three definitions --- it is the
induction step the tiers\textquotesingle{} lemmas plug into, which is
why it renders under its own category head rather than wearing a tier
that does not type-check. Its one \emph{checkable} obligation (E3)
lands as certificate-loader conformance code. The warrant itself is
now mechanized at exactly the certificate-object level this
section\textquotesingle s recipes state it
(\texttt{BernsteinLogicAux/Acceptance.lean}), anchor by anchor:
Recipe~\ref{rec:tc-template} is \texttt{towerCitation\_sound} ---
well-founded induction over admitted rows, consuming the per-rule
lemmas as its step hypothesis --- in finite-acyclic (headline) and
well-founded forms, with root corollary
\texttt{towerCitation\_root}; the E1 typing of
Recipe~\ref{rec:tc-instance} is the row type itself, so the
exhaustiveness of the induction\textquotesingle s case analysis is
type-checked (two consumptions exist: kernel truth of a base leaf,
denotation of a concluded judgment --- no third case to cite); E2 is
the edge metric \texttt{Row.cited} (premise slots only); E3\textquotesingle s
sufficiency is the bridge \texttt{acyclic\_wellFounded} --- on a
finite certificate the checkable acyclicity pass yields exactly the
well-foundedness the induction needs, which is why conformance
suffices where the logical framework deliberately lacks an induction
axiom; and Proposition~\ref{prop:tc-consume} is the transport pair
\texttt{citePos\_sound} / \texttt{citeCmpl\_subset} /
\texttt{citeCmpl\_sound}. The induction theorem proper
(\texttt{towerCitation\_sound\_of\_wf}) depends on no axioms at all;
the bridge and transport lemmas close under the standard three.
\end{bremark}

\begin{bremark}[what remains conditional, named]
This discipline\textquotesingle s mechanization is conditional on
named residues rather than a vague gap.
\textbf{R1} --- the per-rule lemmas enter the induction as a
hypothesis: that is the discipline\textquotesingle s literal content
(it composes conditional lemmas; it does not prove them), but the
type-level identification of that hypothesis\textquotesingle s
instances with this book\textquotesingle s bound per-rule statement
family awaits the deep embedding and the per-rule schema wrappers;
until then the identification is editorial, rule by rule.
\textbf{R2} --- E2 is discharged by \emph{representation}: the
abstract row type carries derivation edges only, and the faithfulness
of the concrete certificate format\textquotesingle s erasure onto it
(no derivation edge lost, no data reference kept) is loader/decode
conformance, outside the module. \textbf{R3} --- admittedness
(E3\textquotesingle s first half) is absorbed by the totality of the
abstract certificate\textquotesingle s indexing: every reference IS an
admitted row, by type; the loader\textquotesingle s admittedness check
is what establishes that totality. (Clause~1, base soundness, is the
base calculus\textquotesingle s own soundness frontier and enters as
the named hypothesis it always was --- a premise of the template, not
a residue of the mechanization.)
\end{bremark}



\begingroup
\renewcommand{\bibname}{References}
\renewcommand{\chapter}{\section}
\pdfbookmark[1]{References}{blcat-refs}\clearpage
\begin{thebibliography}{99}
\footnotesize

\bibitem{farkas1902}
J.~Farkas.
\emph{Theorie der einfachen Ungleichungen}.
Journal f\"ur die reine und angewandte Mathematik 124 (1902), 1--27.

\bibitem{schrijver1986}
A.~Schrijver.
\emph{Theory of Linear and Integer Programming}.
Wiley, 1986. (Farkas variants and multiplier sign conventions: \S7.)

\bibitem{krivine1964}
J.-L.~Krivine.
\emph{Anneaux pr\'eordonn\'es}.
Journal d'Analyse Math\'ematique 12 (1964), 307--326.

\bibitem{handelman1988}
D.~Handelman.
\emph{Representing polynomials by positive linear functions on compact
convex polyhedra}.
Pacific Journal of Mathematics 132 (1988), no.~1, 35--62.

\bibitem{tellegen1952}
B.~D.~H.~Tellegen.
\emph{A general network theorem, with applications}.
Philips Research Reports 7 (1952), 259--269.

\bibitem{willems1972}
J.~C.~Willems.
\emph{Dissipative dynamical systems, Part I: General theory}.
Archive for Rational Mechanics and Analysis 45 (1972), 321--351.

\bibitem{moylanhill1978}
P.~J.~Moylan and D.~J.~Hill.
\emph{Stability criteria for large-scale systems}.
IEEE Transactions on Automatic Control 23 (1978), no.~2, 143--149.

\bibitem{vanderschaftjeltsema2014}
A.~van~der~Schaft and D.~Jeltsema.
\emph{Port-Hamiltonian Systems Theory: An Introductory Overview}.
Foundations and Trends in Systems and Control 1 (2014), no.~2--3, 173--378.

\bibitem{jiangteelpraly1994}
Z.-P.~Jiang, A.~R.~Teel, and L.~Praly.
\emph{Small-gain theorem for ISS systems and applications}.
Mathematics of Control, Signals and Systems 7 (1994), 95--120.

\bibitem{jiangmareelswang1996}
Z.-P.~Jiang, I.~M.~Y.~Mareels, and Y.~Wang.
\emph{A Lyapunov formulation of the nonlinear small-gain theorem for
interconnected ISS systems}.
Automatica 32 (1996), no.~8, 1211--1215.

\bibitem{astromfuruta2000}
K.~J.~{\AA}str\"om and K.~Furuta.
\emph{Swinging up a pendulum by energy control}.
Automatica 36 (2000), no.~2, 287--295.

\bibitem{steele2017}
G.~L.~Steele~Jr.
\emph{It's Time for a New Old Language}.
Keynote, ACM SIGPLAN Symposium on Principles and Practice of Parallel
Programming (PPoPP), 2017. (The ``Computer Science Metanotation'' this
book's rule cells are set in.)


\bibitem{gentzen1935}
G.~Gentzen.
\emph{Untersuchungen \"uber das logische Schlie{\ss}en I, II}.
Mathematische Zeitschrift 39 (1935), 176--210 and 405--431. English
translation in M.~E.~Szabo (ed.), \emph{The Collected Papers of Gerhard
Gentzen}, North-Holland, 1969.

\bibitem{girard1987}
J.-Y.~Girard.
\emph{Linear logic}.
Theoretical Computer Science 50 (1987), no.~1, 1--101. (Source of the
additive/multiplicative vocabulary for context-sharing vs.\
context-splitting rule formats.)

\bibitem{birkhoff1967}
G.~Birkhoff.
\emph{Lattice Theory}.
3rd ed., American Mathematical Society Colloquium Publications 25, 1967.
(M$_3$, the diamond, as the canonical modular non-distributive lattice.)

\bibitem{rudin1976}
W.~Rudin.
\emph{Principles of Mathematical Analysis}.
3rd ed., McGraw--Hill, 1976. (Taylor's theorem with the Lagrange form
of the remainder: Theorem 5.15; the nonnegative-remainder enclosure at
nonnegative argument used by the \rulename{TaylorBound} section.)

\bibitem{banach1922}
S.~Banach.
\emph{Sur les op\'erations dans les ensembles abstraits et leur
application aux \'equations int\'egrales}.
Fundamenta Mathematicae 3 (1922), 133--181.

\bibitem{granasdugundji2003}
A.~Granas and J.~Dugundji.
\emph{Fixed Point Theory}.
Springer Monographs in Mathematics, Springer, 2003. (The closed-ball
condition $d(Tx_0, x_0) \leq (1 - \rho)\,r$ and the a-priori bound
$d(x_0, x^*) \leq \delta/(1-\rho)$: \S I.1.)

\bibitem{doeblin1938}
W.~Doeblin.
\emph{Expos\'e de la th\'eorie des cha\^ines simples constantes de
Markov \`a un nombre fini d'\'etats}.
Revue Math\'ematique de l'Union Interbalkanique 2 (1938), 77--105.
(The origin of the coupling method for Markov chains.)

\bibitem{lindvall1992}
T.~Lindvall.
\emph{Lectures on the Coupling Method}.
Wiley, 1992; corrected reprint, Dover, 2002. (The coupling inequality
$\mathrm{TV} \leq \Pr[X \neq Y]$: \S I.2.)

\bibitem{levinperes2017}
D.~A.~Levin and Y.~Peres, with E.~L.~Wilmer.
\emph{Markov Chains and Mixing Times}.
2nd edition, American Mathematical Society, 2017. (Total-variation
distance, couplings, and the data-processing/projection inequality:
Ch.~4--5.)

\bibitem{hsu2017}
J.~Hsu.
\emph{Probabilistic Couplings for Probabilistic Reasoning}.
PhD thesis, University of Pennsylvania, 2017. arXiv:1710.09951.
(Couplings as certificates in relational program logics --- the
pre/post-style sibling of the judgment form used here.)

\bibitem{girardpappas2007}
A.~Girard and G.~J.~Pappas.
\emph{Approximation metrics for discrete and continuous systems}.
IEEE Transactions on Automatic Control 52 (2007), no.~5, 782--798.
(Approximate simulation with quantitative degradation --- the analytic
ancestor of the $\varepsilon$-witness transport cell.)

\bibitem{desharnais2004}
J.~Desharnais, V.~Gupta, R.~Jagadeesan, and P.~Panangaden.
\emph{Metrics for labelled Markov processes}.
Theoretical Computer Science 318 (2004), no.~3, 323--354.
(Behavioral pseudometrics for probabilistic systems.)

\bibitem{weyl1939}
H.~Weyl.
\emph{The Classical Groups: Their Invariants and Representations}.
Princeton University Press, 1939. (Polarized power sums generate the
multisymmetric invariants in characteristic 0.)

\bibitem{murata1989}
T.~Murata.
\emph{Petri nets: Properties, analysis and applications}.
Proceedings of the IEEE 77 (1989), no.~4, 541--580. (Place invariants
--- the linear conservation laws entering as free count rows.)

\bibitem{spivak2021agg}
D.~I.~Spivak, R.~Garner, and A.~D.~Fairbanks.
\emph{Functorial aggregation}.
arXiv:2111.10968. (Aggregation as commutative-monoid folds along a
classifying map; the multiset as the universal aggregate.)

\bibitem{dashkovskiyrufferwirth2007}
S.~Dashkovskiy, B.~S.~R\"uffer, F.~R.~Wirth.
An ISS small gain theorem for general networks.
\emph{Mathematics of Control, Signals, and Systems} 19(2):93--122,
2007. DOI 10.1007/s00498-007-0014-8. (The $n$-component ISS small-gain
theorem over arbitrary interconnection graphs whose certificate form
the \rulename{SmallGainN} section catalogues; its monotone-operator
fixpoint machinery lives producer-side as existence theory for the
$(\lambda, B, \theta)$ side data.)

\bibitem{nelsonoppen1980}
G.~Nelson and D.~C.~Oppen.
\emph{Fast decision procedures based on congruence closure}.
Journal of the ACM 27 (1980), no.~2, 356--364.

\bibitem{nieuwenhuisoliveras2007}
R.~Nieuwenhuis and A.~Oliveras.
\emph{Fast congruence closure and extensions}.
Information and Computation 205 (2007), no.~4, 557--574.

\bibitem{bernstein1912}
S.~N.~Bernstein.
\emph{D\'emonstration du th\'eor\`eme de Weierstrass fond\'ee sur le
calcul des probabilit\'es}.
Communications de la Soci\'et\'e Math\'ematique de Kharkov (2) 13
(1912), 1--2.

\bibitem{cargoshisha1966}
G.~T.~Cargo and O.~Shisha.
\emph{The Bernstein form of a polynomial}.
Journal of Research of the National Bureau of Standards 70B (1966),
79--81. (Bernstein coefficients enclose the range on the simplex ---
the sound direction this calculus's kernel uses.)

\bibitem{farouki2012}
R.~T.~Farouki.
\emph{The Bernstein polynomial basis: a centennial retrospective}.
Computer Aided Geometric Design 29 (2012), no.~6, 379--419.

\bibitem{decasteljau1959}
P.~de~Casteljau.
\emph{Outillages m\'ethodes calcul}.
Technical report, A.~Citro\"en, Paris, 1959. (The founding
subdivision/evaluation algorithm for polynomials in Bernstein form.)

\bibitem{laneriesenfeld1980}
J.~M.~Lane and R.~F.~Riesenfeld.
\emph{A theoretical development for the computer generation and display
of piecewise polynomial surfaces}.
IEEE Transactions on Pattern Analysis and Machine Intelligence 2
(1980), no.~1, 35--46.

\bibitem{garloff1986}
J.~Garloff.
\emph{Convergent bounds for the range of multivariate polynomials}.
In: Interval Mathematics 1985 (K.~Nickel, ed.), Lecture Notes in
Computer Science 212, Springer, 1986, 37--56. (Subdivision sharpens
Bernstein enclosures --- the completeness partner of the enclosure
soundness used here.)

\bibitem{ehlers2017}
R.~Ehlers.
\emph{Formal verification of piece-wise linear feed-forward neural
networks}.
Proceedings of ATVA 2017 (Automated Technology for Verification and
Analysis), 2017. (The ``triangle'' linear relaxation of ReLU.)

\bibitem{wongkolter2018}
E.~Wong and J.~Z.~Kolter.
\emph{Provable defenses against adversarial examples via the convex
outer adversarial polytope}.
Proceedings of ICML 2018, PMLR 80, 2018.

\bibitem{zhang2018crown}
H.~Zhang, T.-W.~Weng, P.-Y.~Chen, C.-J.~Hsieh, and L.~Daniel.
\emph{Efficient neural network robustness certification with general
activation functions}.
Advances in Neural Information Processing Systems 31 (NeurIPS 2018).
(CROWN: adaptive two-sided linear bounds.)

\bibitem{singh2019deeppoly}
G.~Singh, T.~Gehr, M.~P\"uschel, and M.~Vechev.
\emph{An abstract domain for certifying neural networks}.
Proceedings of the ACM on Programming Languages 3 (POPL 2019),
Article 41. (DeepPoly.)

\bibitem{xu2020lirpa}
K.~Xu, Z.~Shi, H.~Zhang, Y.~Wang, K.-W.~Chang, M.~Huang,
B.~Kailkhura, X.~Lin, and C.-J.~Hsieh.
\emph{Automatic perturbation analysis for scalable certified
robustness and beyond}.
Advances in Neural Information Processing Systems 33 (NeurIPS 2020).
(The LiRPA framework and name.)

\bibitem{katz2017reluplex}
G.~Katz, C.~Barrett, D.~L.~Dill, K.~Julian, and M.~J.~Kochenderfer.
\emph{Reluplex: an efficient SMT solver for verifying deep neural
networks}.
Proceedings of CAV 2017 (Computer Aided Verification), 2017.

\bibitem{bunel2020branch}
R.~Bunel, J.~Lu, I.~Turkaslan, P.~H.~S.~Torr, P.~Kohli, and
M.~P.~Kumar.
\emph{Branch and bound for piecewise linear neural network
verification}.
Journal of Machine Learning Research 21 (2020).

\bibitem{dalrymple2026psd}
David ``davidad'' Dalrymple.
\emph{Certifying strict positive definiteness via an approximate
triangular inverse} (post). X (formerly Twitter), 2026, status
2022316806094913669. (The triangular-inverse method:
sign-and-pattern checks on $A \cdot Z$ for an untrusted triangular
witness $Z$; no divisions, witness accuracy never load-bearing.)

\bibitem{irving2026certificate}
G.~Irving.
\emph{Answer to: Certificate for strict positive definiteness of a
sparse matrix}.
Computational Science Stack Exchange, 2026 (question 45365, answer
45400). (The rounding-tolerant integer variant and its certificate-size
analysis.)

\bibitem{irvingaks}
G.~Irving.
\emph{aks} (Lean formalization repository): the \texttt{Random}
library's diagonal-dominance PSD bridge and a certified
$65536$-node spectral expander.
\texttt{github.com/girving/aks}, 2026.

\bibitem{gershgorin1931}
S.~Gerschgorin.
\emph{\"Uber die Abgrenzung der Eigenwerte einer Matrix}.
Bulletin de l'Acad\'emie des Sciences de l'URSS, Classe des sciences
math\'ematiques et naturelles 6 (1931), 749--754.

\bibitem{taussky1949}
O.~Taussky.
\emph{A recurring theorem on determinants}.
American Mathematical Monthly 56 (1949), no.~10, 672--676. (The
Levy--Desplanques strict-dominance nonsingularity theorem and its
history.)

\bibitem{boman2005}
E.~G.~Boman, D.~Chen, O.~Parekh, and S.~Toledo.
\emph{On factor width and symmetric H-matrices}.
Linear Algebra and its Applications 405 (2005), 239--248.
(Factor-width-2 = scaled diagonal dominance --- the class the /DD
tier's congruence certificates inhabit.)

\bibitem{hornjohnson2013}
R.~A.~Horn and C.~R.~Johnson.
\emph{Matrix Analysis}, 2nd edition.
Cambridge University Press, 2013. (Sylvester's criterion, congruence,
and Gershgorin discs, as cited in the certificate proofs.)

\bibitem{abadilamport1993}
M.~Abadi and L.~Lamport.
\emph{Composing Specifications}.
ACM Transactions on Programming Languages and Systems 15 (1993),
no.~1, 73--132.

\bibitem{mcmillan1999}
K.~L.~McMillan.
\emph{Circular Compositional Reasoning about Liveness}.
In: Correct Hardware Design and Verification Methods (CHARME 1999),
Lecture Notes in Computer Science 1703, Springer, 1999, 342--346.

\bibitem{lynchvaandrager1995}
N.~Lynch and F.~Vaandrager.
\emph{Forward and Backward Simulations, Part I: Untimed Systems}.
Information and Computation 121 (1995), no.~2, 214--233.

\bibitem{capuccimyers2026}
M.~Capucci and D.~J.~Myers.
Certified lenses and compositional system certification.
arXiv:2604.03017 (2026). (The wiring-square conditions, the
$\mathcal{K}^0_\infty$ slack discipline, and the
compose-by-lens-composition-and-slack-addition law cited in the
composition sections.)


\bibitem{ville1939}
J.~Ville.
\emph{\'Etude critique de la notion de collectif}.
Monographies des Probabilit\'es, Gauthier-Villars, Paris, 1939.
(The maximal inequality for nonnegative supermartingales --- the
probability engine of every battery rule.)

\bibitem{doob1953}
J.~L.~Doob.
\emph{Stochastic Processes}.
Wiley, New York, 1953. (Systematic supermartingale theory; maximal
inequalities and optional sampling.)

\bibitem{kushner1967}
H.~J.~Kushner.
\emph{Stochastic Stability and Control}.
Mathematics in Science and Engineering 33, Academic Press, New York,
1967. (Supermartingale safety certificates for stochastic systems.)

\bibitem{meyntweedie2009}
S.~P.~Meyn and R.~L.~Tweedie.
\emph{Markov Chains and Stochastic Stability}.
2nd edition, Cambridge University Press, 2009. (Foster--Lyapunov drift
conditions; the geometric drift form behind the deadline battery.)

\bibitem{prajnajadbabaiepappas2007}
S.~Prajna, A.~Jadbabaie, and G.~J.~Pappas.
\emph{A framework for worst-case and stochastic safety verification
using barrier certificates}.
IEEE Transactions on Automatic Control 52 (2007), no.~8, 1415--1428.

\bibitem{streett1982}
R.~S.~Streett.
\emph{Propositional dynamic logic of looping and converse is elementarily
decidable}.
Information and Control 54 (1982), no.~1--2, 121--141. (The
complemented-pairs acceptance condition.)

\bibitem{rabin1969}
M.~O.~Rabin.
\emph{Decidability of second-order theories and automata on infinite
trees}.
Transactions of the American Mathematical Society 141 (1969), 1--35.

\bibitem{landweber1969}
L.~H.~Landweber.
\emph{Decision problems for $\omega$-automata}.
Mathematical Systems Theory 3 (1969), no.~4, 376--384. (Deterministic
B\"uchi automata recognize exactly the limit languages.)

\bibitem{thomas1990}
W.~Thomas.
\emph{Automata on infinite objects}.
In: Handbook of Theoretical Computer Science, vol.~B, 133--191.
Elsevier, 1990.

\bibitem{foster1953}
F.~G.~Foster.
\emph{On the stochastic matrices associated with certain queuing
processes}.
Annals of Mathematical Statistics 24 (1953), no.~3, 355--360. (The
additive-drift criterion.)

\bibitem{meyntweedie1993}
S.~P.~Meyn and R.~L.~Tweedie.
\emph{Markov Chains and Stochastic Stability}.
Springer, 1993. (Foster--Lyapunov drift conditions, systematized.)

\bibitem{chakarov2013}
A.~Chakarov and S.~Sankaranarayanan.
\emph{Probabilistic program analysis with martingales}.
In: Computer Aided Verification (CAV 2013), Lecture Notes in Computer
Science 8044, 511--526. Springer, 2013. (Ranking supermartingales.)

\bibitem{durrett2019}
R.~Durrett.
\emph{Probability: Theory and Examples}.
5th edition, Cambridge University Press, 2019. (Borel--Cantelli.)

\bibitem{conley1978}
C.~Conley.
\emph{Isolated Invariant Sets and the Morse Index}.
CBMS Regional Conference Series in Mathematics 38, American Mathematical
Society, 1978.

\bibitem{conleyeaston1971}
C.~Conley and R.~Easton.
\emph{Isolated invariant sets and isolating blocks}.
Transactions of the American Mathematical Society 158 (1971), 35--61.

\bibitem{birkhoff1937}
G.~Birkhoff.
\emph{Rings of sets}.
Duke Mathematical Journal 3 (1937), no.~3, 443--454. (The representation
theorem for finite distributive lattices via join-irreducibles.)

\bibitem{nagumo1942}
M.~Nagumo.
\emph{\"Uber die Lage der Integralkurven gew\"ohnlicher
Differentialgleichungen}.
Proceedings of the Physico-Mathematical Society of Japan (3) 24 (1942),
551--559.

\bibitem{blanchini1999}
F.~Blanchini.
\emph{Set invariance in control}.
Automatica 35 (1999), no.~11, 1747--1767.

\bibitem{prajnajadbabaie2004}
S.~Prajna and A.~Jadbabaie.
\emph{Safety verification of hybrid systems using barrier certificates}.
Hybrid Systems: Computation and Control (HSCC 2004), LNCS 2993,
Springer, 2004, 477--492.

\bibitem{kmv2005}
W.~D.~Kalies, K.~Mischaikow, and R.~C.~A.~M.~van~der~Vorst.
\emph{An algorithmic approach to chain recurrence}.
Foundations of Computational Mathematics 5 (2005), no.~4, 409--449.

\bibitem{kmv2016}
W.~D.~Kalies, K.~Mischaikow, and R.~C.~A.~M.~van~der~Vorst.
\emph{Lattice structures for attractors II}.
Foundations of Computational Mathematics 16 (2016), 1151--1191.

\bibitem{kmv2022}
W.~D.~Kalies, K.~Mischaikow, and R.~C.~A.~M.~van~der~Vorst.
\emph{Lattice structures for attractors III}.
Journal of Dynamics and Differential Equations 34 (2022), 1729--1768.
arXiv:1911.09382.

\bibitem{mischaikowmrozek2002}
K.~Mischaikow and M.~Mrozek.
\emph{Conley index}.
Handbook of Dynamical Systems, Vol.~2, North-Holland, 2002, 393--460.

\bibitem{kalieswehbe2026}
W.~D.~Kalies and T.~Wehbe.
\emph{Compositionality of Global Dynamics in Product and Skew-Product
Systems}.
arXiv:2606.30925, 2026.

\bibitem{statman1979}
R.~Statman.
\emph{Lower bounds on Herbrand's theorem}.
Proceedings of the American Mathematical Society 75 (1979), no.~1,
104--107. (Non-elementary lower bounds on quantifier-inference
elimination; with \cite{orevkov1979}, the classical witness for the
cut-elimination blow-up.)

\bibitem{orevkov1979}
V.~P.~Orevkov.
\emph{Lower bounds for lengthening of proofs after cut-elimination}.
Zapiski Nauchnykh Seminarov LOMI 88 (1979), 137--162; English
translation in Journal of Soviet Mathematics (1982).
\end{thebibliography}

\endgroup
\endgroup
 \restoregeometry
 \chapter{Partial Completeness}
\label{ch:qtsl-partial-completeness}

Every rule in the catalogue of \bookref{ch:qtsl-manual}{the QTSL
manual} carries its own soundness story: a statement, a proof, and,
where mechanization has caught up, a machine-checked realization.
Soundness can therefore be audited locally. Each rule earns its keep
on its own, and a user who trusts the kernel need trust nothing else.

Completeness makes a different kind of claim, and no amount of local
auditing can establish it. A completeness theorem says that the rule
set as a whole is not missing anything, so that every true entailment
expressible in some class of specifications is actually derivable.
Claims of that shape must be won fragment by fragment, by exhibiting
either a derivation or a countermodel for every specification the
fragment can express. The prize is a change in the meaning of
refusal. Without completeness, a tool that fails to derive a
specification reports only that a rule might be missing; within a
complete fragment, the same failure delivers a definite verdict about
the specification itself.

This chapter records what has been proved. A reader who wants only
the shape of the result can stop after the completeness theorem
(Theorem~\ref{thm:gnand-completeness}); the axiom basis in
section~\ref{sec:gnand-proof-theory} shows what completeness costs,
and the closing sections account for the evidence and draw the
boundary of the claim. The most surprising idea in the chapter, for
a reader browsing, lives in Remark~\ref{rem:delayer-wings} and in
Figure~\ref{fig:delayer-wings}: the models that witness completeness
must branch, even though nothing in the specifications mentions
branching.

\section{A complete fragment: hard deadlines under staggering}
\label{sec:qtsl-gnand}

The fragment lives in the timed core of the calculus and uses two
modalities over propositional logic.

\begin{definition}[clocked model]
\label{def:clocked-model}
A \emph{clocked model} has finitely many locations, each carrying a
positive maximal dwell time and a fixed valuation of the atoms. A
\emph{trajectory} is a timed path through locations that respects the
dwell bounds.
\end{definition}

The two modalities read as follows.
\begin{itemize}
\item $\Box\varphi$, \emph{invariance}: $\varphi$ holds at every
  point reachable from here. The formula $D := \Box\bot$
  (``deadness'') holds exactly at the states with no future, so
  $\neg D$ asserts liveness.
\item $\mathsf{A}_T\varphi$, \emph{worst-case service}: \emph{every}
  trajectory from here reaches a $\varphi$-point within time at most
  $T$. The quantifier ranges over all branches, so the operator
  reports a worst case rather than a per-path promise. That one word,
  ``every'', is the source of everything difficult below.
\end{itemize}

The characteristic constraint of the fragment, and the origin of its
working name, is the \emph{gnand}\footnote{``Graded NAND'': a
prohibited conjunction, as in NAND, whose conjuncts are graded by
time bounds. The name arose as a working label during the
development and proved too convenient to retire.}, a globally held
prohibition of joint imminence:
\[
\mathrm{gnand}(a, d_a, b, d_b) \;:=\;
\Box\neg\big(\mathsf{A}_{d_a}a \wedge \mathsf{A}_{d_b}b\big),
\]
read: at no reachable moment is $a$ due within $d_a$ while $b$ is due
within $d_b$. It serves as a natural mutual-exclusion primitive for
services that must stay temporally staggered, two operations that may
never both be imminent, in either order.

\begin{example}[a single-boiler espresso machine]
\label{ex:espresso}
A single-boiler espresso machine cannot brew and steam at the same
time. Suppose an espresso order $e$ must be served within $90$
seconds on every branch of the schedule, a milk order $m$ within
$120$, and the boiler forbids any reachable moment at which both are
due within $30$ seconds. The specification comprises the pending
obligations $\neg e \wedge \mathsf{A}_{90}\,e$ and $\neg m \wedge
\mathsf{A}_{120}\,m$, the staggering constraint
$\mathrm{gnand}(e, 30, m, 30)$, and liveness $\neg D$. Suppose the
caf\'e also advertises that the $90$-second figure is tight, so that
some branch really does take the full $90$. That assertion,
$\neg\mathsf{A}_{89}\,e$, already forces the schedule to
\emph{branch}: a barista who always serves promptly would make the
worst case better than declared, so some branch of the schedule must
dawdle while another serves. This innocent observation grows into
the central construction of the completeness proof
(Remark~\ref{rem:delayer-wings}).
\end{example}

The example's shape, in general, defines the fragment.

\begin{definition}[gnand core]
\label{def:gnand-core}
A \emph{gnand core} is a specification asserting, over $n$ atoms:
\begin{itemize}
\item \textbf{pending obligations} $\neg a \wedge \mathsf{A}_{s_a}a$,
  meaning $a$ has not yet been served and must be served within $s_a$
  on every branch;
\item optionally, \textbf{exactness} $\neg\mathsf{A}_{s_a-1}a$ (the
  deadline is tight: some branch delays past $s_a - 1$) or an
  \textbf{avoidance fact} $\neg\mathsf{A}_h a$ (some branch withholds
  $a$ through time $h$);
\item any finite set of gnand constraints over pairs of atoms; and
  liveness $\neg D$.
\end{itemize}
Deadlines may be taken to be integers with no loss of
generality.\footnote{Rational rescaling and the decomposition of a
deadline into iterated unit deadlines are derivable rewrites in the
proof theory, so a fractional menu of deadlines reduces to an
integer one.}
\end{definition}

The fragment deliberately excludes triggered obligations
$\Box(x \to \mathsf{A}_d y)$, the atom-level variant of the
staggering constraint, and possibility assertions $\neg\Box\nu$ with
their associated time-locks. The boundary is principled rather than
incidental; it follows a \emph{sign}. Every constraint in the
fragment either forbids joint imminence or states a fact at the
initial state, and the model construction below is protected by a
monotonicity lemma that holds for such negative-side checks.
The excluded formers \emph{demand} positive worst-case facts at
interior states of a model. Constructions that meet such demands
exist in the development record, but their completeness theory is
open, and this section claims nothing about them.

\subsection{The proof theory}
\label{sec:gnand-proof-theory}

The axiom basis was not designed in advance; it was
\emph{discovered}, by iterating between attempted completeness proofs
and machine searches for countermodels. Three axiom schemata were
added to the seed basis this way, each certified genuinely new by an
explicit independence witness before being adopted, and several
others were pruned as derivable. The final basis comprises:
\begin{itemize}
\item four axioms governing the invariance modality;
\item reflexive service, $\varphi \to \mathsf{A}_T\varphi$;
\item budget composition,
  $\mathsf{A}_{T_1}\mathsf{A}_{T_2}\varphi \leftrightarrow
  \mathsf{A}_{T_1+T_2}\varphi$;
\item the graded absurdity law, $\mathsf{A}_T\bot \to D$;
\item relativization of obligations through invariants,
  $\Box\nu \wedge \mathsf{A}_T(\chi \wedge \nu) \leftrightarrow
  \Box\nu \wedge \mathsf{A}_T\chi$;
\item a first-service law: at the first service among $n$ co-pending
  obligations, the remaining obligations still bind; and
\item the last and decisive addition, a \emph{residual-split}
  first-service law. For any cutoff $c$ strictly below all deadlines,
  \[
  \bigwedge_{i \le n} \mathsf{A}_{T_i}\alpha_i \;\to\;
  \mathsf{A}_c\Big(\bigvee_{i}\big(\alpha_i \wedge
  \bigwedge_{i' \ne i} \mathsf{A}_{T_{i'}}\alpha_{i'}\big) \;\vee\;
  \bigwedge_i \mathsf{A}_{T_i - c}\alpha_i\Big),
  \]
  meaning that within $c$, either some first service has occurred
  (and the co-obligations persist), or nothing has been served and
  every obligation holds at its true residual $T_i - c$.
\end{itemize}
Derivation rules are propositional calculus, modus ponens, import of
a background theory, monotonicity for both modalities, and rational
rescaling.

\subsection{The completeness theorem}
\label{sec:gnand-theorem}

\begin{theorem}[completeness for the gnand fragment]
\label{thm:gnand-completeness}
For gnand cores, the following are equivalent:
\begin{enumerate}
\item the core is satisfiable in a clocked model;
\item an effective arithmetic on the core (a fixpoint computation on
  per-atom effective deadlines and service-precedence edges, with
  one additional branch-conditional check) fires none of its
  refutation conditions;
\item an explicit finite model of the core exists on the uniform
  grid with dwell $\tfrac{1}{n+1}$, where $n$ is the number of
  pending atoms, and the theorem's proof constructs it;
\item the core is consistent: $\not\vdash D$ from the core in the
  basis above.
\end{enumerate}
\end{theorem}

In particular, derivability and clocked validity coincide on the
fragment. Since soundness is proved over a much broader class of
trajectory spaces, of which clocked models are a special case,
validity over every intermediate semantic class coincides with both.
Satisfiability and validity of cores are therefore
\emph{decidable}, not merely semi-decidable; the fixpoint arithmetic
of item~(2) decides both.

Within its fragment the theorem is unconditional; its statement
carries no unproven hypotheses. Its proof, however, forms a chain
with several links, and the accounting must proceed link by link.
That every arithmetic refutation in item~(2) is matched by a formal
derivation of $D$ is a theorem, with the residual-split law doing
the work its adoption promised: the arithmetic's
deadline-tightening step is derivable by iterating the law at
successive cutoffs, descending step by step to the arithmetic's own
value. That every arithmetic-feasible core yields a model at grid
$\tfrac{1}{n+1}$ combines a proved monotone-protection lemma with one
combinatorial schema whose verification is per-instance and
population-scale (section~\ref{sec:gnand-evidence}) rather than
symbolic.

\begin{remark}[delayer wings]
\label{rem:delayer-wings}
The model construction deserves a name of its own. The emitted
models are \emph{schedule trees}: an earliest-deadline-first spine
together with \emph{delayer wings}, auxiliary branches whose only
role is to keep each declared worst case at its declared value, so
that the negative graded facts remain true even though every
individual path serves early (Figure~\ref{fig:delayer-wings}). That
branching structure can satisfy staggering constraints which every
branching-free model provably violates. This was the central
semantic discovery of the development, and it is why the fragment's
models are trees rather than runs.
\end{remark}

\begin{figure}[b]
\centering
\begin{tikzpicture}[
    >=Stealth,
    state/.style={draw, rounded corners=2pt, inner sep=4pt,
                  minimum height=2.6em,
                  font=\footnotesize, align=center},
    spine/.style={->, thick},
    wing/.style={->, thick, dashed, gray!55!black},
    lbl/.style={font=\footnotesize\itshape, midway, above=1pt},
    winglbl/.style={font=\footnotesize\itshape, midway, below=2pt,
                    sloped, gray!55!black}]
  \node[state] (root) {$\neg e,\ \neg m$\\ $e$ due $90$, $m$ due $120$};
  \node[state, right=30mm of root] (se) {$e$ served\\ $m$ due $75$};
  \node[state, right=16mm of se] (sm) {$m$ served};
  \node[state, below right=16mm and 6mm of root] (wait)
    {still $\neg e,\ \neg m$\\ $e$ due $30$, $m$ due $60$};
  \node[state, right=22mm of wait] (we) {$e$ served late\\ ($t = 90$)};
  \draw[spine] (root) -- node[lbl] {serve $e$ at $t=45$} (se);
  \draw[spine] (se) -- node[lbl] {serve $m$} (sm);
  \draw[wing] (root.south) -- node[winglbl] {dawdle to $t=60$} (wait.west);
  \draw[wing] (wait) -- node[winglbl] {serve $e$} (we);
\end{tikzpicture}
\caption{A schedule tree for the espresso core of
Example~\ref{ex:espresso}. The solid spine serves early, in
earliest-deadline-first order. The dashed delayer wing dawdles, so
that the worst-case service time of $e$ really attains the declared
$90$ seconds and the exactness assertion $\neg\mathsf{A}_{89}\,e$
comes out true. No single run can do both jobs; the tree does.}
\label{fig:delayer-wings}
\end{figure}

\subsection{Evidence and assurance status}
\label{sec:gnand-evidence}

The proofs are prose mathematics, developed against an unusually
dense instrument battery. The completeness chain itself is not yet
mechanized in a proof assistant. The three axiom schemata the search
added, however, now carry sorry-free machine-checked soundness
lemmas (\texttt{sound\_u10f}, \texttt{sound\_firstService}, and
\texttt{sound\_fsResidual}, in \texttt{Qtsl/TimedCore.lean} under
\texttt{lean/qtsl} in \path{maitria-coalition/qtsl}; their
statements are reproduced in Appendix~\ref{app:lean-cone}). A
machine-checked countermodel,
\texttt{firstService\_attainment\_necessary}, accompanies them. The
two first-service lemmas assume two properties of the trajectory space,
fusion closure and attainment of first service; clocked models
satisfy both. The countermodel shows that the attainment hypothesis
cannot be dropped, by exhibiting a fusion-closed space in which the
first-service schema's premises hold and its conclusion fails.
Rule-level mechanization otherwise tracks the
catalogue's own realization discipline (\bookref{ch:qtsl-manual}{the
QTSL manual}). The machine evidence behind the chain, all of it
replayable from the development record as it migrates to
\path{maitria-coalition/qtsl}:

\begin{center}
\begin{tabular}{@{}lll@{}}
\toprule
instrument & what it checks & result \\
\midrule
derivation sweep & arithmetic-refuted cores derive $D$ &
  $64{,}786$ cells, $0$ underivable \\
cross-pair sweep & the descent lemma's coverage &
  $4{,}116/4{,}116$ derived \\
branch engine & branch-conditional refutations derive $D$ &
  $1{,}000/1{,}000$; cycle mode $200/200$ \\
model emitter & feasible cores yield grid models &
  $122/122$ named; $541/541$ random \\
soundness controls & no false derivations of $D$ &
  silent on every satisfiable core \\
\bottomrule
\end{tabular}
\end{center}

The controls in the last row did real work during the development.
They caught two draft derivation engines asserting more than their
machinery guaranteed, both retired before any proof consumed them,
and the same battery falsified three earlier conjectured axiom bases
outright. A rule set that has survived that regime differs, as an
epistemic object, from one that has merely not yet failed.

\subsection{What is not claimed}
\label{sec:gnand-not-claimed}

Completeness is claimed for gnand cores only. The composition of
this fragment with triggered obligations and possibility constraints
is open theory; so is the calculus's next expressive tier, in which
worst-case service generalizes to a time-bounded \emph{until} that
carries safety clauses along the way. Minimality of the basis, a
separate question from sufficiency, is settled for integer deadline
menus (the residual-split law is genuinely underivable from the
rest: there is a six-state witness, and six states is minimal) and
open for arbitrarily fine fractional menus. No answer to the
minimality question can weaken the theorem, since a redundant axiom
would only mean that the same completeness holds from fewer
assumptions.

\part{MTK}
\chapter{Tutorial: A Re-entrant Semiconductor Fab}
\label{ch:qtsl-minifab}

\definecolor{wkA}{HTML}{332288} \definecolor{wkB}{HTML}{CC6677} \definecolor{wkC}{HTML}{117733} \definecolor{wkD}{HTML}{882255} \definecolor{wkE}{HTML}{44AA99} \definecolor{wkF}{HTML}{AA4499} \definecolor{czero}{RGB}{108,108,108}
\colorlet{hnaccent}{accent}
\colorlet{hnmid}{accentmid}
\colorlet{hndeep}{accentdeep}
\definecolor{hngray}{RGB}{92,92,92}
\pgfdeclarelayer{holebg}
\pgfdeclarelayer{boxbg}
\pgfdeclarelayer{plates}
\pgfsetlayers{holebg,boxbg,plates,main}

\noindent\draftmark\medskip

The preceding chapters introduced the stack one layer at a time:
mechanisms and their composition (chapter~\ref{ch:qtsl-mathematical-models}),
the pattern inventory that composition ships as
(chapter~\ref{ch:qtsl-composition-patterns}), the semantics that
specifications are stated in
(chapter~\ref{ch:qtsl-semantic-universe}), the plated-certificate
discipline for token-population systems (section~\ref{sec:hyp-plates}
and section~\ref{sec:mt-sdcpn}), and the producer stack that finds
certificates for a checker to replay (chapter~\ref{ch:mtk-cookbook}).
This chapter runs the whole stack over a single model, end to end:
describe a factory line, declare it as database rows --- an SDCPN, a
stochastically and dynamically coloured Petri net in the modelling
catalogue's frame (section~\ref{sec:mt-sdcpn}), which is the native
habitat of everything that follows --- recompose it along its natural
seams, state what should be certified about it, and read the
certificates. The certificates come in two families whose
division of labour is stated rather than smoothed over. The primary family
(section~\ref{sec:minifab-exact}) certifies the discrete line itself,
at its full combinatorial state space: exact statements holding along
every run, at every lot count, under every resolution of the
scheduling freedom --- the variadic-certificate rows of
section~\ref{sec:mt-sdcpn} at work on their first industrial
subject, and then re-read on the wall clock, where the dispatch
seat is priced end to end: a certified floor over every policy, the
exact optimum and the supervisor family attaining it, a widened
full-dispatch seat, an untrusted filler inside a certified
envelope, and one certificate instantiated across six orders of
magnitude of lot count.
The second family (section~\ref{sec:minifab-certs}) certifies the
\emph{fluid relaxation} of the same line: policy-indexed stability
under sustained arrivals, at loads within a tenth of a lot per week
of capacity --- claims the discrete family does not yet price, which
is why the fluid lane is retained in full as the comparison lane,
with the contrast stated where the two families meet. A closing
section then builds the fluid lane's model on the toolkit's
TypeScript surface and drives its release grid back to checked
certificates, in a listing that runs.

The model is the \emph{Mini-Fab}: five machines in three stations,
six route steps, every station visited twice per lot. It was posed by
Kempf at Intel in 1994 as a deliberately minimal semiconductor-fab
exercise and has circulated as a community benchmark since; the
complete quantitative restatement that this chapter's constants are
witnessed against is van den Berk's 2004 Eindhoven master's thesis \cite{vandenberk04}
(the systems engineering group of Rooda and Lefeber), whose second
chapter restates the entire case. The spec digest recording that
witnessing, decision point by decision point, is committed beside the
certificates in the experiments repository
(section~\ref{sec:exp-repro}).

Small as it is, the line exhibits the domain's whole hard structure:
re-entrant flow, a single-machine bottleneck, batching with mix
constraints, sequence-dependent setups, personnel shared across
stations, maintenance calendars, and machine failures that destroy
work in progress. And its headline claim family has teeth, because it
is famously falsifiable: re-entrant lines can destabilize strictly
below nominal capacity under plausible dispatch policies (the
Kumar--Seidman phenomenon \cite{kumar1990dynamic}), so ``this policy
keeps the line stable at this load'' is a claim the world is entirely
capable of refuting.

\section{The line}
\label{sec:minifab-line}

Three workstations process lots of wafers. \emph{Diffusion} has two
furnaces, a buffer for 18 lots, and consumes its input in batches of
exactly three lots. \emph{Implantation} has two machines and a buffer
for 12. \emph{Lithography} has one machine and a buffer for 12. Every
lot follows the same six-step route --- diffusion, implantation,
lithography, implantation, diffusion, lithography --- so each station
is visited twice, running a ``low'' recipe on a lot's first visit and
a ``high'' recipe on its second:

\begin{center}\small
\begin{tabular}{@{}clcrrr@{}}
\toprule
step & station & recipe & process (min) & load (min) & unload (min)\\
\midrule
1 & diffusion     & low  & 225 & 20 & 40\\
2 & implantation  & low  &  30 & 15 & 15\\
3 & lithography   & low  &  55 & 10 & 10\\
4 & implantation  & high &  50 & 15 & 15\\
5 & diffusion     & high & 255 & 20 & 40\\
6 & lithography   & high &  10 & 10 & 10\\
\bottomrule
\end{tabular}
\end{center}

\noindent Loading and unloading hold the machine; at diffusion they
are paid per \emph{batch} of three, not per lot. There is no
preemption and no rework.

\begin{figure}[t]
\centering
\begin{tikzpicture}[x=1mm, y=1mm,
  lot/.style={-{Stealth[length=2.4mm]}, line width=0.6pt, draw=black!70},
  buf/.style={draw=black!60, fill=black!4, rounded corners=1.2pt,
              minimum width=16mm, minimum height=10.5mm,
              align=center, font=\scriptsize, inner sep=1.2pt},
  bufdiff/.style={buf, fill=accent!55},
  buflith/.style={buf, fill=accent!22, draw=accentdeep, line width=0.9pt},
  stn/.style={draw=black!60, fill=white, rounded corners=1.2pt,
              minimum width=34mm, minimum height=12.5mm,
              align=center, font=\scriptsize, inner sep=1.5pt},
  serve/.style={black!45, line width=0.5pt},
  servelith/.style={accentdeep!75, line width=0.7pt}]
  \node[bufdiff] (s1) at (8,30)   {step 1\\low\\225\,min};
  \node[buf]     (s2) at (28,30)  {step 2\\low\\30\,min};
  \node[buflith] (s3) at (48,30)  {step 3\\low\\55\,min};
  \node[buf]     (s4) at (68,30)  {step 4\\high\\50\,min};
  \node[bufdiff] (s5) at (88,30)  {step 5\\high\\255\,min};
  \node[buflith] (s6) at (108,30) {step 6\\high\\10\,min};
  \draw[lot] (-7,30) -- (s1.west);
  \node[align=center, font=\scriptsize, anchor=south] at (-7,33)
    {release\\84 lots/wk};
  \draw[lot] (s1.east) -- (s2.west);
  \draw[lot] (s2.east) -- (s3.west);
  \draw[lot] (s3.east) -- (s4.west);
  \draw[lot] (s4.east) -- (s5.west);
  \draw[lot] (s5.east) -- (s6.west);
  \draw[lot] (s6.east) -- (123,30);
  \node[align=center, font=\scriptsize, anchor=south] at (123,33)
    {finished\\lots};
  \node[stn, fill=accent!55] (diff) at (20,2)
    {\textbf{diffusion}\\2 machines $\cdot$ buffer 18 $\cdot$ batches of 3\\
     83.3\% busy at 84 lots/wk};
  \node[stn] (impl) at (64,2)
    {\textbf{implantation}\\2 machines $\cdot$ buffer 12\\
     58.3\% busy at 84 lots/wk};
  \node[stn, fill=accent!22, draw=accentdeep, line width=1pt] (lith) at (108,2)
    {\textbf{lithography}\\\textcolor{accentdeep}{\textbf{1 machine}}
     $\cdot$ buffer 12\\
     \textcolor{accentdeep}{\textbf{87.5\% busy --- the bottleneck}}};
  \draw[serve] (diff.north) -- (s1.south);
  \draw[serve] (diff.north) -- (s5.south);
  \draw[serve] (impl.north) -- (s2.south);
  \draw[serve] (impl.north) -- (s4.south);
  \draw[servelith] (lith.north) -- (s3.south);
  \draw[servelith] (lith.north) -- (s6.south);
\end{tikzpicture}
\caption{The Mini-Fab's re-entrant route. Buffers are drawn in route
order (top); the three stations (bottom) each serve two
non-consecutive steps, which is what \emph{re-entrant} means --- the
crossing connectors are the route folding back on itself. Lithography
runs both of its steps on a single machine and is the line's
bottleneck. Utilizations are the deterministic machine-occupancy
figures at the nominal release of 84 lots/week.}
\label{fig:minifab-route}
\end{figure}

Three product types are released on constant per-type intervals:
3 test lots, 51 of type A, and 30 of type B per week of $7 \times
1440 = 10080$ minutes --- 84 lots/week in all. One release interval
is worth pausing on: type A arrives every $10080/51 = 3360/17$
minutes. The parameters are all integers, and still the event times
of the system are not, so exact rational time arithmetic is a
necessity rather than a nicety, and the declared model keeps
every duration in $\Q$. Arrivals are deterministic; the only
stochastic elements anywhere in the specification are implantation
failures and repairs (below).

Around the route sit the resources that make the model interesting.
One \emph{transporter} (capacity one lot, 4 minutes per move between
adjacent areas, 1 minute to load or unload) carries every lot between
stations. Two \emph{operators} perform all machine loading and
unloading, non-preemptively: operator~0 serves diffusion and
implantation, operator~1 serves lithography and implantation, so
implantation is contested and every station's handling competes with
another station's. Lithography additionally pays a sequence-dependent
\emph{setup}, performed by an operator with the machine empty,
whenever its input changes against the previous job: 5 minutes for a
product-type change, 10 for a step change, 12 for both. At the
two-machine stations, a \emph{test} lot's second visit must use the
machine it did not use on its first (so both machines stay
monitored); a diffusion batch contains at most one test lot, and a
high-step batch must not mix the two commercial types. Each operator
takes three 60-minute off-blocks per 720-minute shift, placed freely.
One \emph{technician} performs scheduled maintenance (75 minutes per
diffusion furnace per day; 30 and 120 minutes per shift-window at
lithography and implantation respectively), and implantation machines
\emph{fail} during processing --- time between failures uniform on
$[1440, 4560]$ minutes of processing time, repair uniform on $[360,
480]$ --- destroying the lot in process.

The capacity arithmetic that organizes everything downstream: per
lot, lithography is occupied $(10 + 55 + 10) + (10 + 10 + 10) = 105$
minutes on its single machine; diffusion $95 + 105 = 200$ minutes
spread over two furnaces (batch handling divided by three); and
implantation $60 + 80 = 140$ minutes over two machines. At 84
lots/week the three stations run at $87.5\%$, $83.3\%$, and $58.3\%$
respectively --- before setups and maintenance, and the margin
shrinks from there. The single-machine lithography station is the
bottleneck, and the line's deterministic fluid capacity is
$10080/105 = 96$ lots/week, lithography-pinned
(Figure~\ref{fig:minifab-route}).

\section{The model as rows}
\label{sec:minifab-rows}

The declared model is a set of database rows --- the Geolog substrate
of chapters~\ref{ch:geolog-quickstart}--\ref{ch:geolog-reference} ---
over the System vocabulary of chapter~\ref{ch:qtsl-mathematical-models}. The
ground layer is nothing but the tables of the previous section,
recorded as entities with attributes, in the entity-per-row idiom:

\begin{Verbatim}[fontsize=\small,frame=leftline,rulecolor=\color{accentmid},framesep=3mm,xleftmargin=3mm]
Station : Set;
Station/name     : Station -> Str;
Station/machines : Station -> Nat;
Station/bufcap   : Station -> Nat;
Station/batch    : Station -> Nat;   // machine-input batch size (1 = none)

// st_diff: "diffusion" 2 18 3;   st_impl: "implantation" 2 12 1;
// st_lith: "lithography" 1 12 1

RouteStep : Set;
RouteStep/idx     : RouteStep -> Nat;
RouteStep/station : RouteStep -> Station;
RouteStep/recipe  : RouteStep -> Str;   // "low" | "high"
RouteStep/proc    : RouteStep -> Q;     // minutes
RouteStep/load    : RouteStep -> Q;     // handling; the machine is held
RouteStep/unload  : RouteStep -> Q;

// (1, diff, low, 225, 20, 40)   (2, impl, low,   30, 15, 15)
// (3, lith, low,  55, 10, 10)   (4, impl, high,  50, 15, 15)
// (5, diff, high, 255, 20, 40)  (6, lith, high,  10, 10, 10)

SetupRule : Set;                     // lithography, sequence-dependent
SetupRule/typeChange : SetupRule -> B;
SetupRule/stepChange : SetupRule -> B;
SetupRule/minutes    : SetupRule -> Q;
// (t,f,5)  (f,t,10)  (t,t,12)  (f,f,0)

FailRule : Set;                      // implantation only
FailRule/tbfLo : FailRule -> Q;   FailRule/tbfHi : FailRule -> Q;
FailRule/repLo : FailRule -> Q;   FailRule/repHi : FailRule -> Q;
// time between failures 1440..4560 (processing time); repair 360..480
\end{Verbatim}

\noindent Rows are data rather than code: nothing about these
declarations commits to a decomposition, and both of the readings
below compile their dynamics from the same ground tables.

The \emph{monolithic reading} renders the whole fab as one mechanism
--- an SDCPN, and the class name is a parts list for exactly this
specification, worth unpacking on first contact. \emph{Places} hold
coloured tokens: waiting lots (coloured by product type, release
time, and first-visit machine), machine phases (setup, loading, busy,
awaiting unload, unloading, maintenance, down), personnel and
transporter phases, and a sink of finished lots. \emph{Dynamically
coloured}: every duration in the tables is a token colour in motion
--- a timer drifting at $\dot\tau = 1$ inside its phase and forcing a
jump at its threshold --- so time lives in the colours, never in a
side-calendar. \emph{Stochastically}: implantation's failure and
repair durations are random firing measures on the reset, drawn from
the declared uniform family --- and nothing else in this
specification is random. The transitions carry the guards and firing
laws, and the thesis's own release idiom --- per-type signal
generators firing at constant intervals --- \emph{is} the
timer-mechanism idiom, so the model's front door needs no vocabulary
at all beyond what chapter~\ref{ch:qtsl-mathematical-models} already
provides. The net's interface exposes \emph{aggregate} coordinates
only: counts and pools --- work in progress, per-step queue lengths,
busy-machine counts per station, a completed-lot count, and a pool
summing realized flow times.

Three features of the specification look, at first sight, like they
should resist an aggregate-count rendering. None of them does, and
each dissolution is informative.

\paragraph{Batching survives count coordinates.} A diffusion load
consumes three same-step lots subject to joint colour constraints (at
most one test lot; high-step batches commercially pure). At the
aggregate stratum, ``a feasible batch exists'' is a finite
disjunction of count inequalities --- for the high step, with
$k_A, k_B, k_{\mathrm{test}}$ the per-type queue counts,
\[
  k_A \ge 3 \;\vee\; k_B \ge 3 \;\vee\;
  (k_{\mathrm{test}} \ge 1 \wedge k_A \ge 2) \;\vee\;
  (k_{\mathrm{test}} \ge 1 \wedge k_B \ge 2),
\]
and the reset decrements the counts by the chosen composition.
\emph{Which} three lots form the batch is invisible to every
aggregate observable, by exchangeability. What the guard does add is
a \emph{choice} among feasible compositions --- which belongs to
dispatch, below, not to a new row species. Machine dedication for
test lots is likewise just a colour coordinate (the first-visit
machine), with a mismatch guard on the second visit.

\paragraph{FIFO dissolves into class-level dispatch.} The count
grammar deliberately has no expression for ``the $i$-th token,'' so
first-in-first-out order within a colour class is structurally
invisible --- and for this model the invisibility is semantically
exact, not a compromise: service law depends only on a lot's type and
step, so permuting same-colour lots moves no count trajectory, no
utilization or throughput observable, and (by Little's law, which
relates only means) not even mean flow time. What FIFO \emph{does}
determine that matters is which colour class a freed machine serves
next --- and that residue is a dispatch policy: the index of every
certificate in the fluid lane (section~\ref{sec:minifab-certs}), and
a freedom the combinatorial certificates of
section~\ref{sec:minifab-exact} instead quantify over ---
adversarially, every policy at once. The
dissolution does have a boundary: per-lot flow-time
\emph{distribution tails} do depend on within-class order, so a
specification demanding per-lot quantile bounds would reopen the
question (an age coordinate --- still a colour --- is the natural
extension).

\paragraph{Simultaneity is the norm, not a corner case.} Every
duration in the tables is an integer number of minutes measured from
a common origin, so distinct timers reach their thresholds at the
same instant constantly --- every shift boundary fires off-block
guards on three resources at once. The declared model therefore
carries explicit serialization data for simultaneous forced jumps
(a priority order on transition families, with commutation as the
alternative discharge), and the reference simulator's deterministic
tie-break --- time, then declared family order, then entity identity
--- is that serialization made executable. Two correct simulators
that disagree under tie-reordering are answering the commutation
question the serialization row exists to settle.

\section{The compositional reading}
\label{sec:minifab-composite}

The monolithic reading is complete and faithful, and it deliberately
gives up the one thing the composition axis exists to buy:
decomposition. A per-station argument, a per-resource invariant, a
certificate that transports across a boundary --- all of these want
the fab as a \emph{composite}: release generators, three station
mechanisms, two operators, a technician, a transporter, and a sink,
wired together. The seams come in four kinds --- directed lot
handoffs, resource-pool memberships, fused event ports, and the
dispatch freedom the fusions create --- and the pattern inventory
they draw on is chapter~\ref{ch:qtsl-composition-patterns}'s.

What the seams are, mathematically, pattern by pattern:

\paragraph{Directed lot flow.} A station's completions, as a counted
output, wired into the next station's arrival input --- token flow
carried as signal flow, with the lot's colours riding as a tuple. This
is ordinary directed wiring. It has one limitation worth naming:
conservation (a lot leaves diffusion if and only if it arrives at
implantation) is a semantic property the wire itself does not state;
in the compositional reading it is recovered where it matters, as the
pool pattern's invariant.

\paragraph{Resource pools.} Operator~1 must be exclusively held by
either a lithography task or an implantation task, never both; so
must the technician, across three stations, and the transporter,
across every handoff. Continuous signal-sharing carries none of this
--- there is no signal, only a conserved unit. The resource-pool
pattern of chapter~\ref{ch:qtsl-composition-patterns} renders each
resource as a pool with \emph{take}/\emph{give} ports and a crossing
discipline (each user declares which of its modes count as holding
the unit), and its payoff is a theorem rather than a convention: for
a well-formed pool share, every jump of the composite preserves
\[
  k_{\mathrm{pool}} \;+\; \sum_{u \,\in\, \mathrm{users}} \#H_u \;=\; C,
\]
the count-row conservation invariant --- free to the checker, since
it is a linear invariant of the count coordinates. For this model,
$C = 1$ four times over: each of the four pools is a single human or
vehicle.

\paragraph{Event fusion.} ``Load starts'' is one event touching three
parties: the machine (empty, set up), the buffer (a feasible batch
present), and the operator (idle, at the station). The fusion pattern
makes it one \emph{joint jump}: enabling is the conjunction of the
parties' local guards evaluated on the pre-jump snapshot, each party
writes only its own state (so reset commutation within the fusion is
a theorem), and one party --- here the station's timer --- owns the
timing. Without fusion, every synchronized cluster collapses into a
single mechanism, which is the monolithic reading by forced march;
fusion is what lets the composite reading exist at all.

\paragraph{Dispatch.} Fusion creates contention: a freed lithography
machine with both step-3 and step-6 lots waiting; an idle operator
with two stations bidding; a feasible-batch guard satisfiable by
several compositions. Undispatched, the composite is
nondeterministic, and the semantics resolves the freedom
\emph{demonically} --- which is the correct default precisely because
of Kumar--Seidman: ``stable at this load, whatever resolves the
contention'' is typically false below capacity, and should be
unprovable when it is false. A dispatch row installs a named policy
--- a chooser over the enabled-conflict set, reading aggregate
observables. The two certificate families divide over exactly this
freedom: each fluid certificate of section~\ref{sec:minifab-certs}
is indexed by the $(\text{system}, \text{policy})$ pair, while the
combinatorial certificates of section~\ref{sec:minifab-exact} leave
the freedom demonic and hold under every resolution of it. The off-block
placement freedom (operators schedule their own breaks) sits in the
same slot: it is control data, not plant data.

\paragraph{Distribution-valued resets.} Implantation's failure and
repair durations are uniform on rational intervals --- noisy timer
thresholds, neither deterministic resets nor exponential clocks. The
declared form is a reset kernel from a named family
($\mathrm{uniform}[a,b]$ here), consumed by the certification
interface only through its exact rational moments. Fixed and
uniformly-distributed durations make the count process semi-Markov;
the timer-as-state discipline (the supplementary-variable
construction) is what keeps the presentation Markovian in the
extended state, so this is vocabulary, not new metatheory.

Here is where the two readings stand. The monolithic reading uses
mechanism vocabulary alone --- timers, coloured counts, guarded
jumps; the combinatorial certificates of
section~\ref{sec:minifab-exact} are stated against its
token-population refinement (places as plates, lots as tokens), and
the fluid certificates of section~\ref{sec:minifab-certs} against
its aggregate projection. The compositional
reading exercises the fusion, pool, dispatch, and reset-kernel
patterns exactly as chapter~\ref{ch:qtsl-composition-patterns}
catalogues them; the fusion and pool patterns had been reserved as a
future axis until this model arrived as their first customer and
ended the deferral, and the composite row set is mid-flip from
design to artefact --- the fusion stratum's producer surface has
landed, the pool packaging has not, and
section~\ref{sec:minifab-ts} draws that boundary where the code
makes it concrete. Equally informative is what the model does
\emph{not} demand: no nesting, no mode-switched macro-locations
(shifts are timer guards inside mechanisms), no credal wire content,
and no delayed trace --- until some future dispatch chooser reads
\emph{lagged} observables (yesterday's throughput, smoothed work in
progress), at which point the loop becomes observation-delayed
feedback and lands in the delayed-trace pattern. A discrete-event
benchmark exercises a largely disjoint slice of the vocabulary from
the robotics and hybrid-dynamics exemplars the patterns were designed
against, and it still lands on the catalogued inventory; that is
evidence the inventory's shape is holding.

One number quantifies what the composite seams own. The reference
discrete-event simulation committed beside the model tables
(\path{semicon/minifab/sim} in the experiments repository) implements
the deterministic layers only --- machines, batching, setups,
dedication, first-in-first-out within colour --- on an exact
$\tfrac{1}{17}$-minute integer tick, and measures a
deterministic-core throughput ceiling of 87 lots/week. The thesis's
full model --- adding the personnel, transport, maintenance, and
failure layers --- measures a maximal FIFO throughput of $57.5$
lots/week. The roughly thirty lots per week between those two numbers
is carried by \emph{exactly} the layers that live on the pool, fusion,
and reset-kernel seams: a quantitative reading of how much of this
line's behaviour the composition patterns are responsible for.

\section{The specification}
\label{sec:minifab-spec}

What should be certified about a fab divides into two claim
families, and the division is structural rather than stylistic: the
two families quantify differently, are checked by different
machinery, and fail differently.

The first family is \emph{exact and pathwise}: claims about the
discrete line itself, holding along every run, at every lot count,
under every resolution of the scheduling freedom. Conservation ---
machines neither minted nor lost; a lot census that nothing but a
completed lot can decrease; a congestion measure that no firing
increases; and hard progress --- from any work-in-progress profile,
a stated budget of transition firings that no scheduler can exceed.
These are the apparent infinity ``one claim per marking, one row per
transition instance'' collapsed to finitely many symbolic rows by
the plate discipline (section~\ref{sec:hyp-plates}), and
section~\ref{sec:minifab-exact} certifies them. In the released
specification vocabulary
(chapter~\ref{ch:qtsl-specification-language}) every member of this
family now renders. The safety members render as $\mathsf{PAlways}$
atoms over the model's aggregate wires, and a worked render appears
beside its certificate below. The hard firing-count member renders
through the count-ceiling spelling,
$\mathsf{SReturnLE}(B, \mathsf{retFireCount}(R))$: a sure bound on
an event count was a catalogued deliberate absence of the language
until the transition-count functional landed
(section~\ref{sec:spec-lang-devcount}), and the witness class that
entry names for the demand --- structurally bounded
resource-sharing net families --- is instantiated by this model.
The judgment stratum arrived in the same motion: the Rule Catalogue
(chapter~\ref{ch:qtsl-manual}) now carries one plated-transfer rule
per face, each concluding a formed $S \models \varphi$ row --- a
surely-held invariant from one certificate citation per direction,
and a roster firing-count ceiling from one gap-rowed certificate.
What still separates this chapter's rows from formed judgments is
narrower than a missing former, and it is stated with the
certificates (section~\ref{sec:minifab-exact-status}).

The second family is \emph{long-run}, and the claim class matters
(section~\ref{sec:exp-claims}: these are the infinite-horizon-average
rows of the taxonomy, and the class label travels with every result). Over
the monolithic reading's aggregate interface, with $\mathrm{wip}(t)$
the total work in progress and $N(T)$ the completed-lot count, the
specification atoms are:
\begin{itemize}
\item \textbf{stability}: there is a finite $W$ with
  $\limsup_{T \to \infty} \frac{1}{T}\int_0^T \mathrm{wip}(t)\,dt
  \le W$ --- the line does not drown;
\item \textbf{throughput}: $\liminf_{T \to \infty} N(T)/T \ge
  \delta$ --- the release rate is actually sustained;
\item \textbf{cycle time}: mean flow time at most $\varphi$ ---
  which Little's law, $\bar\varphi = \bar{W}/\lambda$, derives from
  the other two rather than asking anyone to certify it separately.
  Certify two, derive the third.
\end{itemize}
In the released vocabulary this family now splits down its middle,
and the split is a fact about the language worth stating precisely.
The stability atom \emph{renders}: the long-run average is an
accumulator of the specification vocabulary --- $\mathsf{poAvg}$,
the Ces\`aro average of a state observable
(section~\ref{sec:spec-lang-returnstratum}) --- and the atom is its
expectation cap,
\[
  \mathsf{SReturnLE}\bigl(W,\ \mathsf{poAvg}(\mathrm{wip})\bigr),
\]
the upper prevision of the averaged work in progress over every law
the dispatched composite admits, with the residual freedoms (setup
sequencing, allocation slack, tie-breaks) resolved adversarially
--- exactly the quantifier this family wants; the level faces
(section~\ref{sec:spec-lang-level}) put threshold probes on the
same functional, and the occupancy siblings render the same way:
the long-run occupied fraction of a station is the average of an
indicator observable, so ``lithography is busy at least
seven-eighths of the time, in the long run'' is
$\mathsf{SReturnGE}(7/8,\ \mathsf{poAvg}(\cdot))$ over the
busy-count indicator --- the vocabulary's own compositional
spelling of occupation fractions. These renders form and build
today (section~\ref{sec:minifab-ts} carries the live witness), and
what stands behind a formed long-run claim is stated with the same
care as everything else in this chapter: the average former's
certificate species --- the gain/bias pair, a Poisson-equation
drift row --- is named by the rule catalogue and not yet claimed,
so these are formed claims awaiting their battery, where the
safety renders of the first family already sit on discharged
certificates.

The throughput atom is the other half, and it now renders --- as an
envelope, which is the only honest shape for it.
$\liminf_{T \to \infty} N(T)/T$ as a bare long-run limit is
deliberately not the vocabulary's object --- a limit claim carries
zero finite-window content, so a consumer sized against it can be
overwhelmed on any finite horizon --- and the event-rate family's
primitive is instead the affine envelope from start
(section~\ref{sec:spec-lang-rate}). The throughput claim the list
above records is the envelope's lower face over the
completion-transition roster, at the release rate $\delta$ and a
certified burst budget $\sigma$,
\[
  \mathsf{SRateGE}\bigl(\delta,\ \sigma,\ R_{\mathrm{sink}}\bigr):
\]
along every admitted law, the expected completed-lot count
dominates $\delta t - \sigma$ at \emph{every} horizon --- which
yields the displayed limit reading as a corollary
($\underline{E}[N(T)]/T \ge \delta - \sigma/T$), never the other
way around. What stands behind the formed claim is stated with the
same care as the stability half's: its certificate species --- the
lower-envelope drift rows, the Foster--Lyapunov shape whose
potential's forced descent pays out completions --- is named by
the rule catalogue and not yet claimed for this exhibit, so the
throughput render joins the long-run family's formed claims
awaiting their battery. The argument the earlier edition recorded
beside the vocabulary --- census certified
(section~\ref{sec:minifab-conservation}), drain rates bounded by
the machine roster, a finite average-work-in-progress bound
forcing $\mathrm{wip}(T)/T \to 0$, the census pinning completion
rate to release rate --- is now the recipe for that certificate's
potential rather than a measurement of an absence. And the
cycle-time atom never wanted a former at all: Little's law derives
it, as the list above already said.

Every claim of this second family is about a $(\text{system},
\text{policy})$ pair, never about the system alone. The undispatched
claim --- stable under \emph{every} resolution of contention --- is
demonic, and for re-entrant lines it is typically \emph{false} below
capacity; the specification's refusal to let per-component stability
compose across an undispatched boundary is correctness rather than
caution. The first family, by contrast, needs no policy index at all
--- its claims are demonic by construction --- and that is precisely
why it can say nothing yet about sustained arrivals: a released lot
\emph{raises} every potential of
section~\ref{sec:minifab-exact}, and pricing arrivals is a different
certificate species (the crossing budgets of the hybrid-process
correspondence, section~\ref{sec:mt-gshs}), deliberately not
interchangeable with ranking rows. Today the arrivals story is the
fluid lane's, the exactness story is the combinatorial lane's, and
the frontier itself is the seam between them.

\section{The combinatorial certificates}
\label{sec:minifab-exact}

What has been certified about the discrete line, with no fluid
scaling and no aggregate projection anywhere in the statement: the
lot-flow core of the Mini-Fab --- the exact Kempf process times, the
station machine counts $(2, 2, 1)$, the full re-entrant route,
simplified so that at most one lot per process step may start,
receive a unit of service, or finish per tick (the model-boundary
paragraph, page~\pageref{par:minifab-model-boundary}) --- as an SDCPN of the modelling catalogue's frame,
a token-population System at its full
combinatorial state space, carrying twelve plated certificates whose
$150$ footprint sequents are discharged by the producer in exact
rational arithmetic (section~\ref{sec:minifab-exact-status} says
what checks them), one row set covering \emph{every} lot count,
\emph{every} scheduler, \emph{every} run. The model is small enough
to print whole, and this section does --- the net on its own page in
Petri notation (Figure~\ref{fig:minifab-sdcpn-petri}), the
classification tables beside it, and the certificate hypernets in
plate notation (Figure~\ref{fig:minifab-cert-hypernets}); every
constant below
re-derives from the campaign committed beside the toolkit's
semiconductor strategy, whose test battery replays the entire
verification (chapter~\ref{ch:mtk-cookbook} for the producer stack;
the strategy source and its battery ship in the toolkit repository).

\subsection{The Mini-Fab as an SDCPN}
\label{sec:minifab-sdcpn}
\label{sec:minifab-plated}

The exhibit is declared in the SDCPN frame's own compositional
reading --- \emph{places are Bags, arcs are patterns}
(section~\ref{sec:mt-sdcpn}) --- and this subsection presents it the
way the container carries it. Sixteen places, each a Bag of coloured
tokens; as certificate subjects they are sixteen \emph{plates} ---
symbolic token populations in the sense of
section~\ref{sec:hyp-plates}:

\begin{center}\small
\begin{tabular}{@{}llll@{}}
\toprule
places (as plates) & count & tokens are & colour\\
\midrule
$q_0$ & 1 & lots not yet released & none\\
$q_1, \dots, q_6$ & 6 & lots waiting at step $k$ & none\\
$p_1, \dots, p_6$ & 6 & lots in process at step $k$ &
  remaining work $w \in [0, T_k]$\\
$m_{\mathrm{diff}}, m_{\mathrm{impl}}, m_{\mathrm{lith}}$ & 3 &
  idle machines & none\\
\bottomrule
\end{tabular}
\end{center}

\noindent with $T = (225, 30, 55, 50, 255, 10)$ the process times of
section~\ref{sec:minifab-line}, read as unit-work quanta: a lot in
process at step $k$ carries the integer minutes of work it has left,
in the declared colour box $[0, T_k]$. Nineteen Petri transitions,
in five families, each declaring its clock species from the frame's
own transition vocabulary:

\begin{center}\small
\begin{tabular}{@{}llll@{}}
\toprule
transition & count & firing & clock\\
\midrule
$\mathrm{release}$ & 1 & $q_0 \to q_1$ & carved $[1,1]$, demonic seat\\
$\mathrm{start}_k$ & 6 & $q_k + m_{\mathrm{st}(k)} \to p_k$ at
  $w = T_k$ & carved $[1,1]$, demonic seat\\
$\mathrm{tick}_k$ & 6 & $p_k$ at $w \ge 1 \to p_k$ at $w - 1$ &
  forced, dwell $1$\\
$\mathrm{finish}_k$ ($k < 6$) & 5 & $p_k$ at $w = 0 \to q_{k+1} +
  m_{\mathrm{st}(k)}$ & carved $[1,1]$, demonic seat\\
$\mathrm{ship}$ & 1 & $p_6$ at $w = 0 \to m_{\mathrm{lith}}$ &
  carved $[1,1]$, demonic seat\\
\bottomrule
\end{tabular}
\end{center}

\noindent The clock declarations are the SDCPN's own. Processing
quanta ride the \emph{forced} species at unit dwell (deterministic
unit work); every scheduling choice --- when to release, which
waiting lot a free machine starts, when a finished lot moves on ---
rides a clock carved to the must-fire window $[1, 1]$, with the
choice itself left to the demonic scheduler seat: timing owned by
the clocks, choice by the seat. Nothing else is pinned --- no
dispatch policy, no tie-break, no priority; both stochastic
transition species stand available and unexercised, every firing
law is Dirac, and no arc reads without consuming: the all-demonic,
unit-dwell, Dirac corner of the frame, whose unfolding is the
faithful ordering-and-counting embedding of an untimed net. Two
derived time coordinates run through everything that follows, both
read off these same clock rows: \emph{counting time} --- the
unfolding's step index, makespans in firings --- carries every
certificate statement from here to the carved clock, and
\emph{metric time} --- the tick meter the carved windows define ---
re-prices the same rows on the wall clock
(section~\ref{sec:minifab-carved}). Whatever a certificate asserts
below, it asserts against every scheduler at once.

Figure~\ref{fig:minifab-sdcpn-petri} draws the net whole, in Petri
notation, initial marking and supervisor seat included; the tables
that follow it print the classification: sixteen places as Bag
populations, nineteen Petri transitions with every arc typed against
the founding taxonomy of Everdij and Blom \cite{everdij-blom-2010},
the arc-type census with the pattern each element flows through, and
the twelve certificates with their folds and transfer polarities.
The rows of these tables are the rows of the emitted container ---
one table line per declared row, in the container's own species
vocabulary --- so the isomorphism is checkable by diffing the page
against the bytes, and the places where a declared row home is
still a named successor are marked on the page rather than smoothed
over. The pattern inventory the classification lands on is
chapter~\ref{ch:qtsl-composition-patterns}'s: the destroy and
create edges are the mode-switching machinery of
section~\ref{sec:cp-switch} in rendezvous with the transition's
timing seat (section~\ref{sec:cp-eventport}), and the enabling-arc
row --- unexercised by the plant, exercised by the supervisor
seat's read-and-return wires
(section~\ref{sec:minifab-synthesis}) --- is the per-tuple
parameter-setting wire of section~\ref{sec:cp-set}. The statics
square behind the ``mechanized'' status entries is
Problem~\ref{prob:mt-sdcpn}, closed in the semantic-universe
development.


\newgeometry{hmargin=1cm,vmargin=1.2cm}
\begin{landscape}
\setlength{\parskip}{0pt}\vspace*{\fill}
\centering
\begin{tikzpicture}[x=1cm, y=0.97cm,
  place/.style={circle, draw=black!80, line width=0.7pt,
    minimum size=8.5mm, inner sep=0pt},
  trans/.style={rectangle, draw=black!80, line width=0.7pt, fill=white,
    minimum width=9.5mm, minimum height=4.5mm, inner sep=1.5pt,
    font=\footnotesize},
  every label/.style={font=\footnotesize, label distance=1.5pt},
  arc/.style={line width=0.55pt, draw=black!80,
    arrows={Bar[width=2.1mm]}-{Stealth[length=2.4mm]}},
  carc/.style={arc, draw=czero},
  rarc/.style={line width=0.55pt, draw=black!80,
    arrows={Stealth[length=2.4mm]}-{Stealth[length=2.4mm]}}]

\node[place, label=above:{$q_0$}] (q0) at (0,0) {};
\node[trans] (rel) at (1.7,0) {$\mathrm{release}$};
\foreach \i in {1,...,6}
  \node[place, label=above:{$q_\i$}] (q\i) at ({3.4*\i},0) {};

\foreach \i/\T/\c/\a in {1/225/wkA/0, 2/30/wkB/0, 3/55/wkC/0,
                      4/50/wkD/0, 5/255/wkE/0, 6/10/wkF/45}{
  \node[place, draw=\c, line width=0.9pt,
        label={[font=\footnotesize]\a:{$p_\i$}}]
        (p\i) at ({3.4*\i+1.0},-4.2) {};
  \node[font=\scriptsize, align=left, anchor=north west, xshift=-1ex]
        at ({3.4*\i+1.5},-4.32)
        {colour:\\[-1.5pt] $\textcolor{\c}{\T} \to \textcolor{czero}{0}$};
}

\foreach \i in {1,...,6}
  \node[trans] (s\i) at ({3.4*\i+0.5},-2.1) {$\mathrm{start}_\i$};
\foreach \i in {1,...,5}
  \node[trans] (f\i) at ({3.4*\i+2.3},-2.1) {$\mathrm{finish}_\i$};
\foreach \i in {1,...,6}
  \node[trans] (t\i) at ({3.4*\i+1.0},-6.3) {$\mathrm{tick}_\i$};
\node[trans] (ship) at (23.3,-4.2) {$\mathrm{ship}$};

\node[place, label=below:{$m_{\mathrm{diff}}$}] (mdiff) at (9.5,-9.3) {};
\node[place, label=below:{$m_{\mathrm{impl}}$}] (mimpl) at (12.5,-9.3) {};
\node[place, label=below:{$m_{\mathrm{lith}}$}] (mlith) at (15.5,-9.3) {};

\fill[black] (q0.center) circle (0.065);
\foreach \a in {0,60,120,180,240,300}
  \fill[black] ([shift={(\a:0.235)}]q0.center) circle (0.065);
\fill[black] ([xshift=-1.2mm]mdiff.center) circle (0.065);
\fill[black] ([xshift= 1.2mm]mdiff.center) circle (0.065);
\fill[black] ([xshift=-1.2mm]mimpl.center) circle (0.065);
\fill[black] ([xshift= 1.2mm]mimpl.center) circle (0.065);
\fill[black] (mlith.center) circle (0.065);

\draw[arc] (q0) -- (rel);
\draw[arc] (rel) -- (q1);
\foreach \i in {1,...,6} \draw[arc] (q\i) -- (s\i);
\foreach \i/\j in {1/2, 2/3, 3/4, 4/5, 5/6} \draw[arc] (f\i) -- (q\j);

\foreach \i/\c in {1/wkA, 2/wkB, 3/wkC, 4/wkD, 5/wkE, 6/wkF}{
  \draw[arc, draw=\c] (s\i) -- (p\i);
  \draw[arc, draw=\c] (t\i) to[bend right=30] (p\i);
}
\foreach \i in {1,...,6} \draw[carc] (p\i) to[bend right=30] (t\i);
\foreach \i in {1,...,5} \draw[carc] (p\i) -- (f\i);
\draw[carc] (p6) -- (ship);

\draw[arc, rounded corners=6pt] (mdiff.115) |- (3.65,-7.35)
      -- ([xshift=-2.5mm]s1.south);
\draw[arc, rounded corners=6pt] (mdiff.65)  |- (17.25,-7.35)
      -- ([xshift=-2.5mm]s5.south);
\draw[arc, rounded corners=6pt] (mimpl.115) |- (7.05,-7.55)
      -- ([xshift=-2.5mm]s2.south);
\draw[arc, rounded corners=6pt] (mimpl.65)  |- (13.85,-7.55)
      -- ([xshift=-2.5mm]s4.south);
\draw[arc, rounded corners=6pt] (mlith.115) |- (10.45,-7.75)
      -- ([xshift=-2.5mm]s3.south);
\draw[arc, rounded corners=6pt] (mlith.65)  |- (20.65,-7.75)
      -- ([xshift=-2.5mm]s6.south);
\draw[arc, rounded corners=6pt] ([xshift=2.5mm]f1.south)
      -- (5.95,-7.95)  -- (9.43,-7.95)  -- (mdiff.100);
\draw[arc, rounded corners=6pt] ([xshift=2.5mm]f5.south)
      -- (19.55,-7.95) -- (9.57,-7.95)  -- (mdiff.80);
\draw[arc, rounded corners=6pt] ([xshift=2.5mm]f2.south)
      -- (9.35,-8.15)  -- (12.43,-8.15) -- (mimpl.100);
\draw[arc, rounded corners=6pt] ([xshift=2.5mm]f4.south)
      -- (16.15,-8.15) -- (12.57,-8.15) -- (mimpl.80);
\draw[arc, rounded corners=6pt] ([xshift=2.5mm]f3.south)
      -- (12.75,-8.35) -- (15.43,-8.35) -- (mlith.100);
\draw[arc, rounded corners=6pt] (ship.south) |- (15.57,-8.35)
      -- (mlith.80);

\begin{pgfonlayer}{holebg}
  \draw[dashed, line width=0.7pt, draw=black!75]
        (-0.4,-11.6) rectangle (24.4,-10.35);
  \node[font=\footnotesize\itshape] at (12.0,-11.15)
        {Dispatcher policy};
\end{pgfonlayer}
\node[place, fill=white,
      label={[font=\footnotesize]160:{$g_1$}}] (g1) at (3.0,-10.35) {};
\node[place, fill=white,
      label={[font=\footnotesize]200:{$c_1$}}] (c1) at (5.45,-10.35) {};
\node[place, fill=white,
      label={[font=\footnotesize]200:{$g_5$}}] (g5) at (16.7,-10.35) {};
\node[rectangle, draw=black!80, line width=0.7pt, fill=white,
      minimum size=3.5mm, inner sep=0pt,
      label={[font=\footnotesize]200:{$p_5$ pool}}] (tp5) at (19.3,-10.35) {};
\node[place, fill=white,
      label={[font=\footnotesize]340:{$c_5$}}] (c5) at (20.25,-10.35) {};
\node[rectangle, draw=black!80, line width=0.7pt, fill=white,
      minimum size=3.5mm, inner sep=0pt,
      label={[font=\footnotesize]200:{$q_5$ pool}}] (tq5) at (24.05,-10.35) {};
\draw[rarc, draw=wkA, rounded corners=6pt] (g1) |- (s1.west);
\draw[arc,  draw=wkA] ([xshift=-2.5mm]f1.south) -- (c1.north);
\draw[rarc, draw=wkE, rounded corners=6pt] (g5) |- (s5.west);
\draw[rarc, rounded corners=6pt] (p5.80) |- (19.9,-3.45)
      -- (19.9,-9.7) -- (19.3,-9.7) -- (tp5.north);
\draw[arc,  draw=wkE, rounded corners=6pt] (f5.east) -| (c5.north);
\draw[rarc, rounded corners=6pt] (q5.east) -- (17.7,0) -- (17.7,-1.2)
      -- (24.05,-1.2) -- (tq5.north);
\end{tikzpicture}\par\medskip
\begin{minipage}{0.94\linewidth}
\platecaption{The Mini-Fab exhibit of this section, in Petri
notation: circles are places, rectangles Petri transitions; every arc
runs place-to-transition or transition-to-place, bar at its source,
head at its target --- $16$ places, $19$ transitions, $49$ arcs, the
census of the classification tables of
section~\ref{sec:minifab-sdcpn}. The initial marking is drawn: $N = 7$
lots pending in $q_0$ (the hexagonal cluster; $N$ appears nowhere else
in the model), and $(2, 2, 1)$ idle machines. Colour coding: each
in-process place $p_k$ declares its colour type, tagged
\emph{colour:} $T_k \to 0$ with $T_k$ the step's work quantum; arcs
\emph{producing} into $p_k$ (its start; its tick's return, paying one
unit of $w$ per firing) wear that type's hue, and every arc
\emph{consuming} from an in-process place wears the shared grey of
the terminal colour $0$ --- the value each countdown drains to, and
the guard at which finishes and shipment fire; the tag numerals wear
the same coding. Machine take and give-back ride
the rectilinear arcs to the pool places at the bottom --- the
two-steps-per-machine-group wiring ($1{,}5$ diffusion; $2{,}4$
implantation; $3{,}6$ lithography) is the re-entrancy. The dashed box
at the foot is the open system's supervisor hole --- the one seat the
plant leaves open, labelled by the policy that fills it --- drawn at
the boundary the synthesized dispatch policy of
section~\ref{sec:minifab-optimum} exercises, one grammar with
Figure~\ref{fig:minifab-pistar}, which draws that filler seated: two
gate places $g_1$, $g_5$, whose read-and-return enabling arcs (a head
at both ends) gate the two diffusion starts --- the seat's act wires;
two firing counters $c_1$, $c_5$, one token per
$\mathrm{finish}_1$/$\mathrm{finish}_5$ firing; and two pooled
read-and-return taps on $q_5$ and $p_5$ --- four observables; with the
gates, the strict subset of the full-dispatch interface
(section~\ref{sec:minifab-fulldispatch}) that $\pi^{\ast}$ needs.
Boundary wires wear the gated step's hue; the boundary is drawn
unfilled --- gate tokens are exactly what a filler mints. The seat's \emph{first}
filler (section~\ref{sec:minifab-synthesis}) sealed this same seat at
a narrower interface --- one weight-$N$ cumulative-arrival tap on
$q_6$, a gate on $\mathrm{start}_6$, and the armed plate's one-shot
key, its conservation pair $k + g = 1$ the certified Boolean gate
discipline. The hole apparatus is the comb's, beyond the plant census
above.}{fig:minifab-sdcpn-petri}
\end{minipage}
\par
\vspace*{\fill}
\thispagestyle{empty}
\end{landscape}
\restoregeometry


{\footnotesize
\setlength{\tabcolsep}{3.5pt}\renewcommand{\arraystretch}{0.92}

\noindent
\begin{minipage}[t]{0.40\textwidth}
\begin{center}
\begin{tabular}{@{}rlllr@{}}
\multicolumn{5}{@{}l@{}}{\textbf{Places} (species \texttt{HPlate})}\\
\toprule
\# & plate & colour & Bag carrier & init\\
\midrule
0 & $q_0$ & --- & $\mathrm{Bag}(\mathbf{1}) \cong \mathbb{N}$ & $N$\\
1 & $q_1$ & --- & $\mathrm{Bag}(\mathbf{1})$ & 0\\
2 & $q_2$ & --- & $\mathrm{Bag}(\mathbf{1})$ & 0\\
3 & $q_3$ & --- & $\mathrm{Bag}(\mathbf{1})$ & 0\\
4 & $q_4$ & --- & $\mathrm{Bag}(\mathbf{1})$ & 0\\
5 & $q_5$ & --- & $\mathrm{Bag}(\mathbf{1})$ & 0\\
6 & $q_6$ & --- & $\mathrm{Bag}(\mathbf{1})$ & 0\\
7 & $p_1$ & $[0, 225]$ & $\mathrm{Bag}([0,225])$ & 0\\
8 & $p_2$ & $[0, 30]$ & $\mathrm{Bag}([0,30])$ & 0\\
9 & $p_3$ & $[0, 55]$ & $\mathrm{Bag}([0,55])$ & 0\\
10 & $p_4$ & $[0, 50]$ & $\mathrm{Bag}([0,50])$ & 0\\
11 & $p_5$ & $[0, 255]$ & $\mathrm{Bag}([0,255])$ & 0\\
12 & $p_6$ & $[0, 10]$ & $\mathrm{Bag}([0,10])$ & 0\\
13 & $m_{\mathrm{diff}}$ & --- & $\mathrm{Bag}(\mathbf{1})$ & 2\\
14 & $m_{\mathrm{impl}}$ & --- & $\mathrm{Bag}(\mathbf{1})$ & 2\\
15 & $m_{\mathrm{lith}}$ & --- & $\mathrm{Bag}(\mathbf{1})$ & 1\\
\bottomrule
\end{tabular}
\end{center}
\end{minipage}\hfill
\begin{minipage}[t]{0.57\textwidth}
\begin{center}
\begin{tabular}{@{}lrl@{}}
\multicolumn{3}{@{}l@{}}{\textbf{The arc-type census, classified}}\\
\toprule
SDCPN element & \# & pattern instance\\
\midrule
ordinary arc, input face & 25 & Bag destroy edge (rendezvous)\\
ordinary arc, output face & 24 & Bag create edge$^{\dagger}$ (spawn)\\
enabledness count & 19 & monotone scope rows\\
colour guard slab & 12 & guard content on the edge\\
enabling arc & 0 & per-tuple parameter-setting\\
inhibitor arc & 0 & excluded by formation\\
delay (rate-raced) transition & 0 & race clock, finite majorant\\
forced (deterministic delay) & 6 & forced clock, dwell $1$\\
immediate transition & 0 & race clock, infinite majorant\\
demonic scheduling & 13 & carved clock $[1,1]$; choice stays demonic\\
firing measure & 19 & Dirac reset legs\\
initial marking & 1 & root init rows\\
\bottomrule
\end{tabular}
\end{center}
\end{minipage}

\smallskip

\begin{center}
\begin{tabular}{@{}rlllll@{}}
\multicolumn{6}{@{}l@{}}{\textbf{Petri transitions} (species
\texttt{Transition}: grain, clock, attachment)}\\
\toprule
\# & transition & consumes [guard slab] & produces [colour law] & clock &
reset rows\\
\midrule
0 & $\mathrm{release}$ & $q_0$ & $q_1$ &
  carved $[1,1]$ & exact tier$^{\dagger}$\\
1 & $\mathrm{start}_1$ & $q_1,\; m_{\mathrm{diff}}$ & $p_1\,[w := 225]$ &
  carved $[1,1]$ & exact tier$^{\dagger}$\\
2 & $\mathrm{start}_2$ & $q_2,\; m_{\mathrm{impl}}$ & $p_2\,[w := 30]$ &
  carved $[1,1]$ & exact tier$^{\dagger}$\\
3 & $\mathrm{start}_3$ & $q_3,\; m_{\mathrm{lith}}$ & $p_3\,[w := 55]$ &
  carved $[1,1]$ & exact tier$^{\dagger}$\\
4 & $\mathrm{start}_4$ & $q_4,\; m_{\mathrm{impl}}$ & $p_4\,[w := 50]$ &
  carved $[1,1]$ & exact tier$^{\dagger}$\\
5 & $\mathrm{start}_5$ & $q_5,\; m_{\mathrm{diff}}$ & $p_5\,[w := 255]$ &
  carved $[1,1]$ & exact tier$^{\dagger}$\\
6 & $\mathrm{start}_6$ & $q_6,\; m_{\mathrm{lith}}$ & $p_6\,[w := 10]$ &
  carved $[1,1]$ & exact tier$^{\dagger}$\\
7 & $\mathrm{tick}_1$ & $p_1\,[1 \le w \le 225]$ & $p_1\,[w := w-1]$ &
  forced & \texttt{tr.scale}\\
8 & $\mathrm{tick}_2$ & $p_2\,[1 \le w \le 30]$ & $p_2\,[w := w-1]$ &
  forced & \texttt{tr.scale}\\
9 & $\mathrm{tick}_3$ & $p_3\,[1 \le w \le 55]$ & $p_3\,[w := w-1]$ &
  forced & \texttt{tr.scale}\\
10 & $\mathrm{tick}_4$ & $p_4\,[1 \le w \le 50]$ & $p_4\,[w := w-1]$ &
  forced & \texttt{tr.scale}\\
11 & $\mathrm{tick}_5$ & $p_5\,[1 \le w \le 255]$ & $p_5\,[w := w-1]$ &
  forced & \texttt{tr.scale}\\
12 & $\mathrm{tick}_6$ & $p_6\,[1 \le w \le 10]$ & $p_6\,[w := w-1]$ &
  forced & \texttt{tr.scale}\\
13 & $\mathrm{finish}_1$ & $p_1\,[w = 0]$ & $q_2,\; m_{\mathrm{diff}}$ &
  carved $[1,1]$ & exact tier$^{\dagger}$\\
14 & $\mathrm{finish}_2$ & $p_2\,[w = 0]$ & $q_3,\; m_{\mathrm{impl}}$ &
  carved $[1,1]$ & exact tier$^{\dagger}$\\
15 & $\mathrm{finish}_3$ & $p_3\,[w = 0]$ & $q_4,\; m_{\mathrm{lith}}$ &
  carved $[1,1]$ & exact tier$^{\dagger}$\\
16 & $\mathrm{finish}_4$ & $p_4\,[w = 0]$ & $q_5,\; m_{\mathrm{impl}}$ &
  carved $[1,1]$ & exact tier$^{\dagger}$\\
17 & $\mathrm{finish}_5$ & $p_5\,[w = 0]$ & $q_6,\; m_{\mathrm{diff}}$ &
  carved $[1,1]$ & exact tier$^{\dagger}$\\
18 & $\mathrm{ship}$ & $p_6\,[w = 0]$ & $m_{\mathrm{lith}}$ &
  carved $[1,1]$ & exact tier$^{\dagger}$\\
\bottomrule
\end{tabular}
\end{center}

\smallskip

\begin{center}
\begin{tabular}{@{}lllll@{}}
\multicolumn{5}{@{}l@{}}{\textbf{Certificates} (species
\texttt{HCert}; $150$ footprint rows, species
\texttt{HFootprintRow})}\\
\toprule
certificate & claim family & fold & features & footprint rows\\
\midrule
work-potential & hard progress (ranking) & sum & $v = F$ &
  $19 \times \text{gap } 1$\\
lot-count & census upper bound & sum & $1$ on lot plates &
  $18$ safety $+$ gap at ship\\
machines-diff-ub & place invariant, $\le$ & sum &
  $+1$ on $m_{\mathrm{diff}}, p_1, p_5$ & $6$ safety\\
machines-diff-lb & place invariant, $\ge$ & sum &
  $-1$ on $m_{\mathrm{diff}}, p_1, p_5$ & $6$ safety\\
machines-impl-ub & place invariant, $\le$ & sum &
  $+1$ on $m_{\mathrm{impl}}, p_2, p_4$ & $6$ safety\\
machines-impl-lb & place invariant, $\ge$ & sum &
  $-1$ on $m_{\mathrm{impl}}, p_2, p_4$ & $6$ safety\\
machines-lith-ub & place invariant, $\le$ & sum &
  $+1$ on $m_{\mathrm{lith}}, p_3, p_6$ & $6$ safety\\
machines-lith-lb & place invariant, $\ge$ & sum &
  $-1$ on $m_{\mathrm{lith}}, p_3, p_6$ & $6$ safety\\
congestion-$\ell_2$ & congestion non-increase & conj.\ pow.\ $p{=}2$ &
  $v = R$ & $6 \times \varphi$-gap $1$ $+$ $13$ safety\\
work-ceiling-max & ceiling safety & max & $v = R$ &
  $19$ safety (gap refused)\\
arrival-potential & deadline ranking (reach $q_6$) & sum & $v = D$ &
  $16 \times \text{gap } 1$ $+$ $3$ safety\\
step6-budget & post-arrival budget & sum & $v = F - D$ &
  $3 \times \text{gap } 1$ $+$ $16$ safety\\
\bottomrule
\end{tabular}
\end{center}

\medskip
\noindent\textbf{Places.} A place is a Bag of coloured tokens ---
remaining work $w$ in its declared box for the in-process places,
unit tokens (so: a counter) everywhere else. The
machine plates are the resource-pool pattern in SDCPN dress: take
and give are ordinary arcs, and the pool's conservation law is
certified (the machine pairs above), not assumed.

\smallskip
\noindent\textbf{Transitions.} Every arc is an \emph{ordinary} arc:
each consumed entry is one input-face arc, each produced entry one
output-face arc, every firing law Dirac --- a deterministic colour
polynomial in the consumed colours. There are no enabling arcs
(every emitted selection row carries $\text{read} = 0$, a fact a
reader can grep the container for) and no inhibitor arcs:
enabledness everywhere is a monotone count-and-slab condition, so
adding tokens never disables a step. Guard slabs are the footprint
sequents' hypothesis slabs; every dwell is $1$.
$^{\dagger}$The six same-place ticks carry their reset as a declared
scale row (\texttt{tr.scale}: image and reset polynomial in the
container). The thirteen place-crossing transitions discharge their
produced side directly in the exact-rational tier; the declared
reset-row home arrives with the Bag reset-row presentation, staged
for the next theory revision.

\smallskip
\noindent\textbf{The census.} Counts are arcs, transitions, or slabs.
The zero rows read two ways: enabling arcs and both stochastic
transition species are \emph{available and unexercised} --- this
exhibit is the demonic-choice, carved-unit-dwell, Dirac corner of
the frame (every clock carved to the window $[1,1]$: timing owned
by the clocks, choice by the seat), the discrete-time embedding of
an untimed net, faithful for ordering and counting --- while the
inhibitor row is
\emph{unavailable by design}, warranted by the upward-closure
theorem. Since every certificate quantifies over the demonic
scheduler, its claims dominate every rate-raced refinement of the
same scheduling freedom: the stochastic species, when they arrive,
only resolve freedom the certificates have already priced. The
founding taxonomy's conflict-resolution element is likewise priced
rather than modelled: no priority or randomization rows exist, and
every consequence holds by dominance over all of them at once.

\smallskip
\noindent\textbf{Certificates.} $F$ is the remaining-firings
dictionary ($F_k = \sum_{j \ge k} (T_j + 2)$;
$v(q_0) = F_1 + 1 = 638$, $v(q_k) = F_k$,
$v(p_k, w) = w + 1 + F_{k+1}$, machines uncited) and $R$ the
remaining-work dictionary ($R_k = \sum_{j \ge k} T_j$;
$v(q_0) = R_1 = 625$, $v(q_k) = R_k$,
$v(p_k, w) = w + R_{k+1}$). Transfer polarity reads off the row
profile: an all-gap sum certificate is a ranking (every firing pays,
so runs are budgeted --- $638$ firings per lot, tight); a
safety-rowed pair in both signs is a conserved place invariant
($\#m_s + \#p_i + \#p_j = C_s$); the conjugated-power gap rows
budget the ticks in $\varphi$-space; and the max fold carries safety
only --- a gap row on an idempotent fold is a formation error the
vocabulary load-rejects, exercised by the mutant battery.

\smallskip
\noindent\textbf{Deadline.} $D$ is the arrival dictionary: $F$ with
the step-6 remainder $F_6 = T_6 + 2 = 12$ subtracted on every
position before $q_6$, and $0$ from $q_6$ on --- $v_D$ vanishes
exactly where a lot has entered $q_6$, so the fold's zero-set
\emph{is} the reach goal. The pointwise split $D + (F - D) = F$
retells the ranking as a deadline argument: pre-arrival firings pay
$D$ (at most $626N$ from $N$ pending lots), post-arrival firings pay
$F - D$ (at most $12$ per arrived lot), so along every maximal run
every lot has entered $q_6$ within $T(N) = 638N - 12$ steps ---
counting time, the discrete-time unfolding's step index (the
metric-time note below prices the same bound on the tick meter).
Read as a Spec this is
the deadline lower bound
$\mathrm{P}_{\ge 1}\!\left[\mathrm{F}_{\le T(N)}\,
(\text{all $N$ lots have entered } q_6)\right]$ with demonic timing:
in the Dirac corner every scheduler resolution is a point mass, the
dictionaries quantify over the whole demonic seat, and the bound is
tight --- the schedule that defers the last $\mathrm{finish}_5$
behind all other work attains it exactly (and typically-scheduled
runs sit at it too: the step-6 tail is $12$ firings against step
5's $255$-tick pole). Three honesty lines the arithmetic enforces:
the goal is \emph{cumulative} arrival, never simultaneous occupancy
($\#q_6 = N$ at once is demonically avoidable --- start step 6 as
lots arrive and $q_6$ stays drained); the claim reads at the
deadline wall (a record that stops existing is not late; the
completion conjunct that upgrades the wall to full worst-case reach
is the liveness face, priced separately); and the counting-time
reading is primary but no longer alone --- every clock in this
encoding is carved to a finite must-fire window, and the
metric-time note below prices the same bound with steps replaced by
window widths. The judgment row binding this Spec to the container
is the completion face's declared row home --- not smoothed over:
the deadline content shipped here is the two certificates' $38$
footprint rows.

\smallskip
\noindent\textbf{Deadline, metric time.} With every clock carved to
a positive finite must-fire window (dwell exactly one at these
declarations), the counting deadline becomes a wall-clock statement:
all $N$ lots enter the finishing queue within $638N - 12$ time units
along every maximal schedule ($626N$ under the latest-safe unlock)
--- probability $1$, uniformly over the scheduler's remaining
freedom. The certificate content is exactly the dictionary pair
already presented; the transfer multiplies by the largest window
ceiling (here $1$), read off the model bytes by a must-fire census
that refuses any idling member. Three clauses carry the reading,
each at its stated tier: the dictionary rows are geolog-checked
in-loop; the window-top law (a finite ceiling is a deadline) is the
specification's declared clock semantics, its meter arithmetic owed
at the proof tier; and maximality with deadlock-freedom-below-goal
remain the same two argued clauses the counting statement already
carries --- the checkable species for the latter exists (the
angelic-selection witness family) and awaits its active-region
staging. The bound is firing-tight but not time-tight:
simultaneously due windows fire in batches, and the round replay
brackets the attained wall well below the certified one ---
structure, not slack in the rows.

\smallskip
\noindent\textbf{Open system and synthesis.} The plant opens into a
comb: one supervisor seat whose declared interface reads the
cumulative-arrival observable and drives the $\mathrm{start}_6$
gate wire. A supervision-pattern filler closes it inside the model
class --- per gate, a gate place $g$, a one-shot key $k$ (initial
$1$), and an unlock jump consuming $k$, reading $q_6$ at weight $N$
(returning what it reads) and minting $g$; the gated Petri
transition gains a read-and-return enabling arc, the census row the
plant alone leaves at zero. Closed roster: $18$ places, $20$ Petri
transitions, $14$ certificates, $160$ footprint rows --- the
plant's twelve certificates extend (the ranking pays the unlock's
single firing through the key: total completing-run firings
$= 638N + 1$), and the Boolean gate discipline is \emph{certified},
not declared, by the conservation pair $k + g = 1$. Joint
policy+certificate synthesis runs lexicographically: certified
completion first (the gate that never opens ties the arrival
optimum, then deadlocks with every lot in $q_6$ shipping nothing
--- refused, starved plate named; so is the over-gated policy whose
claimed horizon $370N$ beats the optimum by starving arrival work
out of its own constraint set), then minimize the certified
horizon, where the least dictionary-certifiable $T$ is an
exact-rational linear program over dictionary space (gap rows on
the admitted pre-goal roster, boxwise nonnegativity, goal pins,
goal dominance) whose deterministic optimal vertex \emph{is} the
arrival dictionary $D$ restricted to the closed roster. The winner
is the latest-safe unlock (reopen exactly at the $N$th arrival):
under it no pre-goal firing exists but the $626N$ arrival-paying
ones, so every schedule of the closed composite reaches
all-arrivals at exactly step $626N$ --- against the open plant's
$638N - 12$, a certified throughput-floor gain of $12(N-1)$ steps
($+1.6\%$ at $N = 7$; $+1.9\%$ at the weekly $N = 84$: $52584$ vs
$53580$). Honesty lines: the geolog-checked rows carry the
\emph{within}-$626N$ direction (gap rows in the loop); the
\emph{not-before} direction is the producer tier's exact-drop
identities (each gap row's displacement equals its gap identically,
both directions by Bernstein) plus the replay bracket, its in-loop
row species a named absence; the $626N$ floor over \emph{all}
fillers is plant arithmetic ($626$ lot-specific firings per
arrival, one firing per composite step, policy firings only add);
completion in every encoding is two-tier (ranking geolog-checked;
enabledness below the ship goal argued); the seat's weight-$N$
observation read arc is the one arc-separation clause the flattened
encoding cannot pass --- carried as the staged observation
stratum's named debt, never smoothed; and this armed-plate encoding
has a native-read sibling ($17$ places, $20$ Petri transitions, no
key) whose completion carries one clause more --- its unlock
consumes nothing, so no ranking row can gap it, and
unlock-fires-finitely-often rides as a stated convention --- each
variant citing its own roster and tier, neither borrowing the
other's numbers.

}
Two structural findings the classification surfaces, both
informative rather than embarrassing. First, the deliberate corner
the clock table pins --- all freedom demonic, all processing
deterministic unit work, every firing law Dirac --- is load-bearing
for transport: the certificates are stated against that corner's
adversary, which is the \emph{largest} one, so every rate-raced or
policy-refined elaboration of the same line --- the stochastic
species the frame holds ready, the failure and repair kernels of
the fluid lane's fuller model --- only resolves freedom these rows
have already priced, and the classification doubles as a transport
argument. Second, the machine plates exhibit the resource-pool
pattern and the inhibitor exclusion in one stroke: mutual exclusion
of a conserved unit needs only ordinary arcs plus the certified
invariant (the machine-conservation pairs of
section~\ref{sec:minifab-conservation}), and nowhere in a
re-entrant line with true capacities does an absence test arise ---
the model never misses the excluded primitive, which is the
exclusion argument of section~\ref{sec:mt-sdcpn} read off a working
industrial subject.

\subsection{One row per Petri transition}
\label{sec:minifab-footprint}

A certificate here is a plated certificate exactly as
section~\ref{sec:mt-sdcpn} catalogues the species: a per-token
\emph{feature} --- a polynomial in the representative token's colour
block --- on each cited plate; a commutative-monoid \emph{fold}
(sum, max, or the conjugated sums between them); and one
\emph{footprint sequent} per (certificate, Petri transition) pair,
asserting that the firing's produced contribution is bounded by its
consumed contribution --- minus a pinned rational gap, on the rows
that claim progress. Grammar equivariance is what makes the row
finite: no expression names ``the $i$-th token,'' so the
transition-level row \emph{is} the family, at every marking, at
every lot count. The
twelve certificates below carry $150$ footprint sequents in all ---
$19$ each for the six whole-net certificates, $6$ for each of the
six machine-conservation certificates --- and every sequent is a
polynomial inequality of degree at most two over the declared colour
boxes, discharged in the producer's exact tier by
Bernstein-coefficient nonnegativity, in exact $\Q$, with no
floating-point number representable in the pipeline. Coverage (every
transition that may touch a cited plate carries a row), floor rows
(gap-carrying folds bounded below), gap positivity, and
scale-parameter conformance are checked in the same pass. The
certificates themselves are finite drawn objects --- plated feature
cells, folds, four shapes carrying the whole roster --- and
Figure~\ref{fig:minifab-cert-hypernets} draws them in plate
notation.

\subsection{The work potential, and the run-length budget}
\label{sec:minifab-ranking}

The flagship certificate is a shaped potential over the sum fold.
Its feature dictionary is ``remaining firings'': with
$F_k = \sum_{j \ge k} (T_j + 2)$ the firings a lot still owes from
queue $k$ onward (one start, $T_j$ ticks, one finish per remaining
step), assign
\[
  v(q_0) = F_1 + 1, \qquad
  v(q_k) = F_k, \qquad
  v(p_k, w) = w + 1 + F_{k+1}, \qquad
  v(m) = 0 .
\]
Every one of the nineteen transitions carries a gap row with gap $1$,
and each is one line of arithmetic: a release moves a token from
$F_1 + 1$ to $F_1$; a start consumes $F_k = T_k + 2 + F_{k+1}$ and
produces $T_k + 1 + F_{k+1}$; a tick pays one unit of $w$; a finish
consumes $1 + F_{k+1}$ and produces $F_{k+1}$; the final shipment
consumes $1$ and produces nothing the dictionary counts. Coverage is
total, the floor is nonnegativity of the features, and the transfer
consequence (section~\ref{sec:mt-sdcpn}) is the specification's hard
progress claim: \textbf{along every run from any marking $m_0$, the
total number of transition firings is at most the fold
$V(m_0) = \sum v$} --- uniformly in the lot count and in the
scheduler. At the Kempf constants the per-lot budget is
\[
  1 + \sum_{k=1}^{6} (T_k + 2) \;=\; 1 + 625 + 12 \;=\; 638
\]
firings, so $N$ pending lots drain in at most $638N$ firings ---
$53{,}592$ at the nominal weekly release of $84$ lots. And because
this particular dictionary telescopes \emph{exactly} --- every
transition drops the fold by exactly one --- the bound is tight: every
complete drain attains it, which the independent replay below
measures at every tested lot count. This is the claim the fluid
relaxation structurally cannot state (a continuous level supports no
firing count), just as the arrivals-and-load claims of
section~\ref{sec:minifab-certs} are the ones this certificate
deliberately does not make.

The budget now has its render in the released vocabulary: for a run
from marking $m_0$, the claim is the count ceiling
$\mathsf{SReturnLE}\bigl(V(m_0),\, \mathsf{retFireCount}(R)\bigr)$
with $R$ the full nineteen-transition roster
(section~\ref{sec:spec-lang-devcount}), and the catalogue's
ranking-face plated-transfer rule concludes exactly that judgment
from one gap-rowed certificate --- this one. What stands between
this render and a formed row is recorded in
section~\ref{sec:minifab-exact-status}.

\subsection{Conservation}
\label{sec:minifab-conservation}

Three certificate groups carry the bookkeeping claims, each over the
sum fold with constant features.

\emph{Lot census.} Feature $1$ on every lot population (pending,
queues, in-process; machines uncited): eighteen safety rows assert
that no transition mints a lot, and the shipment transition carries
a gap row --- the census strictly drops exactly there. The transfer
consequence: shipments along any run number at most the initial lot
count. The dual lower bound is deliberately \emph{not} certified,
because it is false --- lots genuinely leave at shipment --- and the
refusal is the honest content: this vocabulary cannot be talked into
asserting conservation of a quantity the model does not conserve.

\emph{Machine populations.} For each station, a pair of sum
certificates --- features $+1$ and $-1$ on the idle plate and the
station's two in-process plates, safety rows on the six incident
transitions --- pins both directions of
\[
  \#m_s \;+\; \#p_i \;+\; \#p_j \;=\; C_s ,
  \qquad (C_{\mathrm{diff}}, C_{\mathrm{impl}}, C_{\mathrm{lith}})
  = (2, 2, 1) :
\]
the classical place invariant, as certificate rows: a start takes an
idle machine and opens an in-process slot, a finish does the
reverse, a tick touches neither, and no other transition can touch the
three plates at all --- machines are neither minted nor lost, under
any scheduler, at any lot count. In the specification vocabulary
this pair renders as
$\mathsf{PAlways}(\mathsf{SAnd}(\mathsf{SAtomLE}(C_s, \kappa_s),
\mathsf{SAtomGE}(C_s, \kappa_s)))$ for the station's population
observable $\kappa_s$ --- a const-$1$ aggregate wire over the three
cited plates, the same row species as the queue-length wires below;
the judgment row that consumes the render is the catalogue's
safety-face plated-transfer rule, whose conclusion is this
conjunction under the sure cap $\mathsf{SProbGE}(1, \cdot)$ and
whose premise shape --- one certificate citation per direction ---
is this pair's own (section~\ref{sec:minifab-spec}).

\subsection{The bottleneck pair: two folds, one dictionary}
\label{sec:minifab-bottleneck}

The next two certificates share one feature dictionary ---
\emph{remaining total work}: with $R_k = \sum_{j \ge k} T_j$ (so
$R_1 = 625$), assign $v(q_0) = R_1$, $v(q_k) = R_k$,
$v(p_k, w) = w + R_{k+1}$ --- and differ only in the fold, which is
exactly the point.

\emph{Congestion, $\ell_2$.} The conjugated-sum fold at $p = 2$
(the scale-parametric family of section~\ref{sec:mt-sdcpn},
strictly monotone at every finite $p$, interpolating from sum
toward max). Read in $\varphi$-space, per that section's normative
discipline --- the rows constrain $\sum_t v_t^2$, and the checker
never computes a root: every tick transition carries a gap row with
$\varphi$-space gap $1$ (from $w \ge 1$, the drop
$(w + R)^2 - (w - 1 + R)^2 = 2(w + R) - 1 \ge 1$), and every other
transition carries a safety row, because the dictionary is arranged so
that moves preserve remaining work exactly: a start converts
$R_k = T_k + R_{k+1}$ into $w + R_{k+1}$ at $w = T_k$, a finish
converts $0 + R_{k+1}$ into $R_{k+1}$, and the shipment retires
$R_7 = 0$. The even-power order-preservation side condition is the
nonnegativity of the feature range on the declared boxes ---
discharged, not assumed. Consequences: the $\ell_2$ work-congestion
measure never increases along any run, and processing quanta number
at most the initial $\varphi$-mass $\sum v(m_0)^2$. Stated with the
modulus honesty the modelling chapter demands: the $\varphi$-space
budget is loose in aggregate units (at the weekly release the
$52{,}500$ ticks of the drain sit far below the initial
$\varphi$-mass of $84 \times 625^2 = 32{,}812{,}500$), which is
exactly why the $\varphi$-space statement --- not a root of it ---
is the certified one; the sharp tick count is the work potential's
job, one certificate over.

In the released specification vocabulary the non-increase claim has
a worked render today. The emitted aggregate wires include the
$\varphi$-space work masses of the six in-process populations and
the six queue lengths, so the congestion observable is the linear
read-out
\[
  s_2 \;=\; \sum_{k} \mathrm{l2work}_{p_k}
       \;+\; \sum_{k} R_k^2 \cdot \mathrm{len}_{q_k}
       \;+\; R_1^2 \cdot \mathrm{len}_{q_0} ,
\]
and for a run from a marking with congestion value $S_0$ the render
is $\mathsf{PAlways}(\mathsf{SAtomLE}(S_0, s_2))$ --- formation per
chapter~\ref{ch:qtsl-specification-language}, the atom's polynomial
canonical over the wire coordinates.

\emph{Ceiling, max.} The same dictionary under the max fold: every
transition a safety row, no gap row anywhere --- the maximum remaining
total work over lots never increases. The asymmetry with the
congestion certificate is enforced, not stylistic:
section~\ref{sec:mt-sdcpn} establishes that idempotent folds are
progress-incapable in every scale (an untouched dominating token
absorbs any local drop), so a gap row on a max certificate is a
\emph{formation error}, and the vocabulary load-rejects it --- the
constructor refuses the row, and the mutant battery below exercises
the refusal. Safety is what a max fold can carry, and safety is what
this certificate claims.


\newgeometry{hmargin=1cm,vmargin=1.2cm}
\begin{landscape}
\setlength{\parskip}{0pt}\vspace*{\fill}
\centering
\begin{tikzpicture}[x=1cm, y=1cm,
  cell/.style={draw=hndeep, fill=hnaccent!50, rounded corners=1pt,
    font=\footnotesize, inner sep=3pt, minimum height=4.2mm},
  fold/.style={draw=hndeep, fill=hnaccent!50, rounded corners=1pt,
    font=\footnotesize, inner sep=3pt, minimum height=4.2mm,
    align=center},
  pbox/.style={draw=hngray, fill=hngray!12, rounded corners=1pt,
    font=\footnotesize\ttfamily, inner sep=2.5pt, minimum height=4.2mm},
  plate/.style={draw=hnmid, line width=0.7pt, rounded corners=2pt,
    fill=hnaccent!18, inner sep=5pt},
  platelab/.style={font=\tiny\ttfamily, text=hnmid, inner sep=1pt},
  wire/.style={-, line width=0.5pt, draw=black!75},
  wlab/.style={font=\scriptsize, inner sep=1pt},
  port/.style={font=\small\itshape, inner sep=1.5pt},
  ptitle/.style={font=\footnotesize\bfseries, anchor=west, inner sep=0pt},
  ann/.style={font=\scriptsize, text=black!70, align=left,
    anchor=north west, inner sep=0pt}]

\begin{scope}[shift={(0,0)}]
\node[ptitle] at (-0.4,1.0)
  {(a) work-potential --- hard progress (ranking); fold $\Sigma$;
   dictionary $F$};
\node[port] (a-in0) at (0,0)    {$q_0$};
\node[port] (a-in1) at (0,-1.35) {$q_k$};
\node[port] (a-in2) at (0,-2.7)  {$p_k$};
\node[cell] (a-c0) at (2.7,0)     {$F_1 + 1$};
\node[cell] (a-c1) at (2.7,-1.35) {$F_k$};
\node[cell] (a-c2) at (2.7,-2.7)  {$w + 1 + F_{k+1}$};
\begin{pgfonlayer}{plates}
  \node[plate, fit=(a-c0)] (a-pl0) {};
  \node[plate, fit=(a-c1)] (a-pl1) {};
  \node[plate, fit=(a-c2)] (a-pl2) {};
\end{pgfonlayer}
\node[platelab, anchor=north east] at (a-pl0.south east) {plate $q_0$};
\node[platelab, anchor=north east] at (a-pl1.south east) {plate $q_k$};
\node[platelab, anchor=north east] at (a-pl2.south east) {plate $p_k$};
\node[fold] (a-f) at (6.4,-1.35) {fold $\Sigma$};
\node[port] (a-V) at (8.1,-1.35) {$V$};
\draw[wire] (a-in0) -- node[wlab, above, pos=0.4] {$[q_0]$} (a-c0);
\draw[wire] (a-in1) -- node[wlab, above, pos=0.4] {$[q_k]$} (a-c1);
\draw[wire] (a-in2) -- node[wlab, above, pos=0.4] {$[p_k]$} (a-c2);
\draw[wire] (a-c0.east) -- (a-f.west);
\draw[wire] (a-c1.east) -- (a-f.west);
\draw[wire] (a-c2.east) -- (a-f.west);
\draw[wire] (a-f) -- (a-V);
\node[ann, text width=9.4cm] at (-0.4,-3.7)
  {(one $q_k$ and one $p_k$ row per step, $k = 1, \dots, 6$; all
   plate indices bound by the fold) \\[2pt]
   $F_k = \sum_{j \ge k} (T_j + 2)$, the firings a lot still owes;
   $v(m) = 0$: machines uncited. One gap row per transition, gap
   $1$ --- $19$ rows; along every run, from any marking $m_0$, total
   firings $\le V(m_0)$, tight at $638$ per lot.};
\end{scope}

\begin{scope}[shift={(0,-7.3)}]
\node[ptitle] at (-0.4,1.0)
  {(b) machine conservation --- place invariants; fold $\Sigma$;
   $\pm$ twins};
\node[port] (b-in0) at (0,0)     {$m_s$};
\node[port] (b-in1) at (0,-1.35) {$p_i$};
\node[port] (b-in2) at (0,-2.7)  {$p_j$};
\node[cell] (b-c0) at (2.7,0)     {$+1$};
\node[cell] (b-c1) at (2.7,-1.35) {$+1$};
\node[cell] (b-c2) at (2.7,-2.7)  {$+1$};
\begin{pgfonlayer}{plates}
  \node[plate, fit=(b-c0)] (b-pl0) {};
  \node[plate, fit=(b-c1)] (b-pl1) {};
  \node[plate, fit=(b-c2)] (b-pl2) {};
\end{pgfonlayer}
\node[platelab, anchor=north east] at (b-pl0.south east) {plate $m_s$};
\node[platelab, anchor=north east] at (b-pl1.south east) {plate $p_i$};
\node[platelab, anchor=north east] at (b-pl2.south east) {plate $p_j$};
\node[fold] (b-f) at (6.4,-1.35) {fold $\Sigma$};
\node[port] (b-V) at (8.1,-1.35) {$\kappa_s$};
\draw[wire] (b-in0) -- node[wlab, above, pos=0.4] {$[m_s]$} (b-c0);
\draw[wire] (b-in1) -- node[wlab, above, pos=0.4] {$[p_i]$} (b-c1);
\draw[wire] (b-in2) -- node[wlab, above, pos=0.4] {$[p_j]$} (b-c2);
\draw[wire] (b-c0.east) -- (b-f.west);
\draw[wire] (b-c1.east) -- (b-f.west);
\draw[wire] (b-c2.east) -- (b-f.west);
\draw[wire] (b-f) -- (b-V);
\node[ann, text width=9.4cm] at (-0.4,-3.7)
  {the $-1$ twin certifies the reverse inequality, pinning
   $\kappa_s = \#m_s + \#p_i + \#p_j = C_s$ under every scheduler, at
   every lot count; $(s;\, i, j) \in \{(\mathrm{diff};\, 1, 5),\,
   (\mathrm{impl};\, 2, 4),\, (\mathrm{lith};\, 3, 6)\}$ with
   $(C_{\mathrm{diff}}, C_{\mathrm{impl}}, C_{\mathrm{lith}}) =
   (2, 2, 1)$ --- $6$ safety rows per certificate, $36$ in all.};
\end{scope}

\begin{scope}[shift={(12.3,0)}]
\node[ptitle] at (-0.4,1.0)
  {(c) the bottleneck pair --- two folds, one dictionary $R$};
\node[port] (c-in0) at (0,0)     {$q_0$};
\node[port] (c-in1) at (0,-1.35) {$q_k$};
\node[port] (c-in2) at (0,-2.7)  {$p_k$};
\node[cell] (c-c0) at (2.6,0)     {$R_1$};
\node[cell] (c-c1) at (2.6,-1.35) {$R_k$};
\node[cell] (c-c2) at (2.6,-2.7)  {$w + R_{k+1}$};
\begin{pgfonlayer}{plates}
  \node[plate, fit=(c-c0)] (c-pl0) {};
  \node[plate, fit=(c-c1)] (c-pl1) {};
  \node[plate, fit=(c-c2)] (c-pl2) {};
\end{pgfonlayer}
\node[platelab, anchor=north east] at (c-pl0.south east) {plate $q_0$};
\node[platelab, anchor=north east] at (c-pl1.south east) {plate $q_k$};
\node[platelab, anchor=north east] at (c-pl2.south east) {plate $p_k$};
\node[pbox] (c-cp0) at (5.1,0)     {copy$_2$};
\node[pbox] (c-cp1) at (5.1,-1.35) {copy$_2$};
\node[pbox] (c-cp2) at (5.1,-2.7)  {copy$_2$};
\node[fold] (c-f1) at (7.6,-0.55) {fold $\Sigma_\varphi$,\;
  $\varphi = ({\cdot})^2$};
\node[fold] (c-f2) at (7.6,-3.05) {fold $\max$};
\node[port] (c-V1) at (10.1,-0.55) {$V_2$};
\node[port] (c-V2) at (10.1,-3.05) {$V_{\max}$};
\draw[wire] (c-in0) -- node[wlab, above, pos=0.4] {$[q_0]$} (c-c0);
\draw[wire] (c-in1) -- node[wlab, above, pos=0.4] {$[q_k]$} (c-c1);
\draw[wire] (c-in2) -- node[wlab, above, pos=0.4] {$[p_k]$} (c-c2);
\draw[wire] (c-c0.east) -- (c-cp0.west);
\draw[wire] (c-c1.east) -- (c-cp1.west);
\draw[wire] (c-c2.east) -- (c-cp2.west);
\draw[wire] ([yshift= 1.1mm]c-cp0.east) -- (c-f1.west);
\draw[wire] ([yshift= 1.1mm]c-cp1.east) -- (c-f1.west);
\draw[wire] ([yshift= 1.1mm]c-cp2.east) -- (c-f1.west);
\draw[wire] ([yshift=-1.1mm]c-cp0.east) -- (c-f2.west);
\draw[wire] ([yshift=-1.1mm]c-cp1.east) -- (c-f2.west);
\draw[wire] ([yshift=-1.1mm]c-cp2.east) -- (c-f2.west);
\draw[wire] (c-f1) -- (c-V1);
\draw[wire] (c-f2) -- (c-V2);
\node[ann, text width=10.6cm] at (-0.4,-4.2)
  {$R_k = \sum_{j \ge k} T_j$, the work a lot still owes ($R_1 =
   625$); rows repeat per step as in (a). Congestion-$\ell_2$
   (upper fold): obligations stated in $\varphi$-scale --- $\sum
   v^2$, no root ever computed, verdicts transporting monotonely ---
   $6$ $\varphi$-gap rows (the ticks) $+$ $13$ safety rows.
   Work-ceiling (lower fold): safety rows only ($19$); its gap row
   is \emph{refused at load} --- a gap row on an idempotent fold is a
   formation error, not a lint (section~\ref{sec:mt-sdcpn}).};
\end{scope}
\end{tikzpicture}\par\medskip
\begin{minipage}{0.94\linewidth}
\platecaption{The exhibit's certificate hypernets, in plate
notation. Each tinted enclosing box is a functorial box declared per
node: the enclosed feature cell is applied once per token of the plate
named at the box corner (its whiskered lift,
section~\ref{sec:hyp-plates}), and its wires carry the plate's
symbolic extent, marked $[q_k]$, $[p_k]$; a fold node consumes the
plated wires into one scalar. The drawn artefact is therefore a single
finite object denoting a function of the whole variable-dimensional
marking --- the grammar has no expression for ``the $i$-th token,''
which is why one footprint row per (certificate, transition) pair
covers every instance at every marking
(section~\ref{sec:minifab-footprint}). The sum fold is contraction
with the plate indices bound; constant features compile to the
standard count wire, so (b) reads
$\kappa_s = \#m_s + \#p_i + \#p_j$. Four shapes carry all twelve
certificates of the classification tables
(section~\ref{sec:minifab-sdcpn}): (a) the ranking; (b) the six
machine invariants; (c) the two folds sharing the remaining-work
dictionary $R$; the lot census is (b)'s wiring with feature $1$ on the
nine lot populations ($q_0$, $q_k$, $p_k$; machines uncited), $18$
safety rows and a gap row exactly at shipment; and the deadline pair
(section~\ref{sec:minifab-deadline}) is (a)'s wiring twice more, its
dictionary split pointwise into $D$ and $F - D$ ($16$ gap $+$ $3$
safety, and $3$ gap $+$ $16$ safety). Row total:
$19 + 19 + 36 + 19 + 19 + 38 = 150$ footprint rows, the tables'
census.}{fig:minifab-cert-hypernets}
\end{minipage}
\par
\vspace*{\fill}
\thispagestyle{empty}
\end{landscape}
\restoregeometry

\subsection{The deadline pair: the flagship dictionary, split
pointwise}
\label{sec:minifab-deadline}

The last two certificates price a deadline, and they cost no new
dictionary: they split the work potential's $F$ pointwise. Let $D$
be the \emph{arrival dictionary} --- $F$ with the step-6 remainder
$F_6 = T_6 + 2 = 12$ subtracted on every position before $q_6$, and
$0$ from $q_6$ on --- so $v_D$ vanishes exactly where a lot has
entered $q_6$: the fold's zero-set \emph{is} the reach goal.
\emph{Arrival-potential} carries $D$ over the sum fold: sixteen gap
rows with gap $1$, and safety rows on the three transitions at or
past arrival, which pay nothing. \emph{Step6-budget} carries the
complement $F - D$: three gap rows on the post-arrival transitions,
sixteen safety. The pointwise identity $D + (F - D) = F$ retells the
run-length budget as a deadline argument --- pre-arrival firings pay
$D$ (at most $626N$ from $N$ pending lots), post-arrival firings pay
$F - D$ (at most $12$ per arrived lot) --- so along every maximal
run, every lot has entered $q_6$ within
\[
  T(N) \;=\; 638N - 12
\]
steps: counting time, the discrete-time unfolding's step index. As a
specification this is the deadline lower bound
$\mathrm{P}_{\ge 1}\!\left[\mathrm{F}_{\le T(N)}\,
(\text{all $N$ lots have entered } q_6)\right]$ with demonic timing
--- in the Dirac corner every scheduler resolution is a point mass,
and the dictionaries quantify over the whole demonic seat. The
bound is tight: the schedule that defers the last
$\mathrm{finish}_5$ behind all other work attains it exactly, and
typically-scheduled runs sit at it too, the step-6 tail being $12$
firings against step 5's $255$-tick pole --- the independent replay
below witnesses both dynamically.

Three honesty lines the arithmetic enforces, stated rather than
smoothed over. The goal is \emph{cumulative} arrival, never
simultaneous occupancy: $\#q_6 = N$ at once is demonically avoidable
(start step 6 as lots arrive and $q_6$ stays drained), so no
certificate of this pair claims it. The claim reads at the deadline
wall --- a record that stops existing is not late; the completion
conjunct that upgrades the wall to full worst-case reach is the
liveness face, priced separately. And the deadline above is
stated in counting time; its metric-time twin rides the carved
windows the clock rows already declare --- a dwell floor with a
finite must-fire ceiling --- under which these same two
dictionaries price the same bound with steps replaced by ticks
(section~\ref{sec:minifab-carved}). The
judgment row binding the deadline specification to the container is
the completion face's declared row home --- a named successor; the
deadline content shipped today is the pair's $38$ footprint
sequents.

\subsection{The open system: the policy synthesized jointly with
its certificate}
\label{sec:minifab-synthesis}

The deadline pair prices the line against every dispatch at once,
and that universality is also the bound's ceiling: $638N - 12$ is
what no scheduler can exceed, not what a good one achieves. The
specification section already named the shape of the next claim
family --- everything long-run is about a $(\text{system},
\text{policy})$ pair, never the system alone
(section~\ref{sec:minifab-spec}) --- and here the exact pathwise
family acquires the same index, as \emph{syntax}: the plant opens
into a comb. Same sixteen places, same nineteen Petri transitions,
plus one \emph{hole} --- the dashed box at the foot of
Figure~\ref{fig:minifab-sdcpn-petri}, a first-class term of the
System grammar carrying nothing but a declared interface: read the
cumulative-arrival observable, drive the $\mathrm{start}_6$ gate
wire. The hole is not a figure of speech; the emitted container
carries the comb's own constructor rows, and the claim that exactly
one hole occurs is decided in the production loop by an occurrence
walk over those rows. What fills it is the supervision pattern of
section~\ref{sec:cp-supervise}, arriving exactly as that entry
promised: the filler is an ordinary jump-only System owning a gate
place, the gated transition carries a weight-1 read-no-consume
enabling arc --- the census row the plant alone leaves at zero ---
and the closed composite is again a coloured net of the same class.
Supervision is composition inside the model class, so filling the
hole is no semantic ceremony: it is one relation between three
System terms, and the closed composite re-enters the same
certificate machinery as the plant itself.

The synthesis is joint: policy and certificate are searched for
together, the certificate machinery serving as the search
instrument rather than a post-hoc audit. The objective is
lexicographic, completion first. A candidate policy must first
carry a certified completion claim --- a ranking over the closed
composite whose rows the checker re-derives --- and only then
competes on the certified arrival horizon, where the least
dictionary-certifiable $T$ is an exact-rational linear program over
dictionary space: gap rows on the transitions the policy admits
pre-goal, boxwise nonnegativity, zero pinned on the goal region,
dominance everywhere a pending lot sits. Refutation is the
improvement signal: exact replay of the closed composite under
adversarial schedules either attains the certified horizon ---
closing the bracket pointwise --- or returns a trajectory whose
step ledger shows where the adversary spends its firings, which is
precisely the gate the next candidate should own. On this model the
search terminates at a theorem rather than a timeout. The winner is
the \emph{latest-safe unlock} --- hold $\mathrm{start}_6$ closed
until the $N$th lot enters $q_6$, then open --- and under it no
pre-goal firing exists except the $626N$ that pay the arrival
dictionary, so every schedule of the closed composite reaches
all-arrivals at exactly step $626N$. The linear program agrees from
the other side: its deterministic optimal vertex \emph{is} the
arrival dictionary $D$ of the deadline pair, restricted to the
closed roster. Against the open plant's $638N - 12$ the synthesized
policy saves $12(N-1)$ steps --- a certified throughput-floor gain
of $+1.6\%$ at $N = 7$ and $+1.9\%$ at the weekly $N = 84$ ($52584$
against $53580$) --- and $626N$ is the floor over \emph{all}
fillers: each arrival needs $626$ lot-specific plant firings,
firings interleave one per composite step, and policy firings only
add, so no filler of this comb does better in counting time. The
synthesized policy is not merely better; it is optimal, with both
endpoints of the bracket certified objects.

The lexicographic order is not pedantry, and the mutant battery
proves it the vivid way. Drop the completion clause and the
``best'' policy is the gate that never opens: it ties the arrival
optimum --- with the post-arrival work locked out, every firing
pays the arrival dictionary, and all $N$ lots reach $q_6$ at
exactly step $626N$ --- and then deadlocks with every lot in
$q_6$, shipping nothing, a policy the deadline alone would score
perfect. Its sharper sibling gates $\mathrm{start}_5$ instead and
claims a horizon of $370N$, ``beating'' the true optimum by
starving the arrival work out of its own constraint set: gap rows
are demanded only on transitions the policy admits, so a policy
that forbids the work escapes the rows that would price it. Both
ride the battery with expectations committed in advance, and both
are refused at the completion certificate, the starved queue named
in the refusal. That is the completion-first discipline earning its
place in the objective, and it is the synthesis story in miniature:
the certificate engine does not merely bless the winner --- it is
what catches the degenerate candidates a deadline-only score would
crown.

The campaign ships the closed composite in two encodings, and the
difference between them is a certificate-tier fact, stated rather
than smoothed. The \emph{native-read} encoding adds one gate place
and one unlock jump that reads $q_6$ at weight $N$ and returns what
it reads --- seventeen places, twenty Petri transitions, the roster
the comb container carries --- but its unlock consumes nothing, so
no ranking row can gap it: a maximal run may spin the unlock
forever instead of shipping, and ``the unlock fires finitely
often'' rides as a stated convention, argued rather than checked,
one clause more than the completion claim otherwise needs. The \emph{armed-plate} encoding
spends one more place to delete the clause: a one-shot key $k$ with
initial marking $1$, consumed by the unlock --- eighteen places,
twenty Petri transitions --- makes the extended ranking all-gap
including the unlock (the completing run pays $638N + 1$ firings,
the key buying the single unlock) and turns the Boolean gate
discipline from a declaration into a certified conservation pair
$k + g = 1$, the machine-pair idiom aimed at the policy's own
state. Each variant cites its own roster and its own tier; neither
borrows the other's numbers. In both, the completion claim itself
is two-tier --- the ranking is geolog-checked; enabledness below the
ship goal (some transition stays live while lots remain: in-process
lots always tick or finish, machine conservation supplies an
in-process lot whenever a station's idle plate is empty, queues are
unbounded) is argued, the liveness face the deadline subsection
already priced separately. The remaining honesty lines are the
deadline pair's, inherited and extended. The geolog-checked rows
carry the
within-$626N$ direction; the not-before direction is the producer
tier's exact-drop identities --- each gap row's displacement equals
its gap identically, both directions by Bernstein --- together with
the replay bracket, pointwise $626N$ at every schedule tested, its
in-loop row species a named absence (the lower-drift polarity the
row vocabulary does not yet carry). The supervisor's weight-$N$
observation read arc fails exactly one clause of the supervision
census --- a policy arc reading a plant place; the packaged
observation legs are the staged stratum that owns it --- carried as
a named debt in the container's own notes. And the
filler-interface tie on a wired comb is refused producer-side, the
substitution binding behind it argued until the checker-side
evaluator lands. The two successors earlier editions of this
chapter named here have both since landed against this comb, and
the closing subsections of this section walk them: the metric rung ---
the carved clock, under which these same dictionaries price
wall-clock deadlines and the dispatch seat is priced end to end
(sections~\ref{sec:minifab-carved}--\ref{sec:minifab-fulldispatch})
--- and the envelope arm, an uncertified dispatcher run inside a
certified switching envelope whose downside the same family-level
safety rows bound (section~\ref{sec:minifab-envelope}). That both
were designed against this comb is the architectural point: the
open system is one comb, and everything after it is fillers.

\subsection{What checks these rows --- today, and next}
\label{sec:minifab-exact-status}

The trust story is staged, and each stage is named rather than
blurred.

\begin{itemize}
\item \textbf{Checked today by the producer (trusted Rust, not
  formally verified):} the toolkit's producer-side verifier,
  \texttt{verify\_campaign} in mtk, checks all $150$ footprint
  sequents, coverage, the certificates' floor conditions (every cited
  feature nonnegative over its plate's declared colour box --- a
  side condition of each certificate, not a lower bound on the
  optimum), gap positivity, and scale-parameter
  conformance, in exact rational arithmetic (Bernstein-coefficient
  nonnegativity, exact $\Q$). The campaign --- model, dictionaries,
  rows --- and the tests that run this verification are public, beside
  the toolkit's semiconductor strategy, at the mtk revision the
  experiments repository pins (\texttt{8904467}). Two
  separate facts limit what that buys a reader. First, the verifier
  is ordinary Rust that we trust, not formally verified. Second, it cannot
  be rebuilt from public sources today: the snapshot depends on
  development revisions of sibling crates that are not published, so
  a reader cannot reproduce the producer check from public sources
  alone. The checker of the next item, which
  reads the sealed container, independently re-evaluates only $36$ of the $150$ footprint rows,
  those of the six tick transitions; the other $114$, including the
  rows of the release, start, finish and ship transitions on which
  the deadline depends, are checked by this producer alone. The
  sealed bytes do not identify the net: a variant in which lots skip
  step~2 seals to the same bytes. So the sealed checker's $36$-row
  evaluation does not by itself authenticate the net printed in this
  chapter, nor the $114$ producer-only rows; ``the net'' in every
  claim here is the one defined by the pinned mtk source.
\item \textbf{Discharged today, in the production loop:} emitted
  certificate containers are sealed deterministically (double
  emission is byte-identical), and the engine-side theory provider
  discharges the canonicity claims --- one per distinct feature
  polynomial, forty at this model --- while refusing
  containers that claim nothing (a vacuous accept is structurally
  unavailable, and the battery below includes the probe). Since the
  evaluator family reached the coloured tier, the same in-loop walk
  also discharges the plated side-condition claims the container
  now carries: footprint on the thirty-six tick rows --- the
  flagship's six, whose colour-typed feature dictionaries are read
  and decided from the container's own bytes, the lot count's six,
  the machine conservations' twelve, and the deadline pair's twelve
  --- flow on all twelve
  certificates, and the five floor claims the sign analysis reaches,
  with zero refusals and zero failures on the real emission; a
  standalone checker binary consuming the same sealed bytes reports
  identical tallies.
\item \textbf{Landed at a named fragment, refused loudly beyond it:}
  the theory demands its own side-condition families for every
  plated certificate --- footprint, coverage, floor, flow, and
  parameter conformance --- and the engine-side evaluator family
  that discharges them in-loop is now decided at the coloured
  place/transition fragment: same-place colour dictionaries, sum
  folds, displacement settled where it is constant or of
  single-signed even shape over the colour space, exact $\Q$. What
  still lies beyond is named per row rather than blurred: coverage
  and the thirteen place-crossing transitions' footprint rows await the
  produced-side reset-row presentation; the congestion ticks, the
  flagship's floor, and the conjugated power's parameter conformance
  await the declared colour-box rows (the Bernstein-over-box core
  that will decide them is written and exercised against exactly
  these shapes); the machine lower bounds' floors await a coupling
  route for negative features; and the max fold is an idempotent
  kind, outside the polynomial walk by design. A claim placed on
  any of these is refused with the unrealized case named --- the
  battery below exercises two such refusals --- so their discharge
  remains the exact tier's, replayed independently below, and the
  demand rows travel with the container either way: nothing is
  silently waived.
\item \textbf{The warrant, and the judgment:} the run-length
  consequence rests on the transfer argument of
  section~\ref{sec:mt-sdcpn} --- elementary telescoping, each
  dropped hypothesis refuted there by a finite countermodel. The
  judgment stratum has landed since this campaign first ran: the
  catalogue's plated-transfer pair concludes formed
  $S \models \varphi$ rows for the safety and ranking faces, with
  the verification content normative in the evaluators' walks
  (section~\ref{sec:minifab-spec}). The catalogue states what
  remains open, loudly: no Lean statement binds either rule node
  yet, so the formal lemma pack is still the named successor. The
  gap is not hidden in the artefact: the certificate carries rows
  and demands, and states its consequence as a consequence.
\item \textbf{Jurisdiction, measured:} extraction toward the
  certified-binary checker reports this family's containers as
  carrying no content in that checker's fragment --- correctly, and
  by design: plated rows are certificate-stratum vocabulary outside
  the fragment the machine-code soundness theorem binds, and their
  discharge route is the evaluator family and the lemma pack above,
  not the certified binary. The measured refusal is the correct
  answer, and it is the division of labour the stack's consumers
  were built to make explicit: one set of bytes, checkers that fail
  differently, none asked to vouch beyond its fragment.
\item \textbf{The reproduction spine:} this campaign's committed
  home in the experiments repository, with a make target per number,
  arrives with the spine's uniform regeneration --- the motion that
  has already reached the fluid lane next door
  (section~\ref{sec:minifab-certs}) --- and the chapter's numbers
  move in lock-step with it.
\end{itemize}

\subsection{The mutant battery and the independent replay}
\label{sec:minifab-exact-mutants}

Twelve deliberately corrupted campaigns, expectations committed
before any run, each aimed at a specific gate and each in fact
refused there: a broken gap (a tick that returns its token
unchanged --- the exact tier's footprint discharge fails); a gap row
on the max fold (refused at construction --- the formation error
above, exercised); a dropped footprint row (the coverage gate); a
nonpositive gap (the gap-positivity gate); a conjugated power at
$p = 0$ (the parameter floor); a non-canonical feature polynomial
(killed in the production loop by the engine provider's canonicity
check --- a real in-loop kill, not a producer-side courtesy); a
container claiming nothing (refused in-loop as a vacuous accept);
now that the claims themselves ride in-loop, three aimed at
the coloured walk: a broken gap in the emitted bytes (its footprint
claim \emph{fails} in the loop --- lying bytes under a live claim);
a footprint claim placed on a row the walk cannot yet decide
(refused, the missing colour-domain stratum named); and a floor
claim on a negative-feature certificate (refused, the unbuilt
coupling route named); and, with the deadline pair landed, two
aimed at its rows: a corrupted arrival dictionary (the exact tier's
footprint discharge fails at the mutated transition, the locus
named), and a deadline gap inflated in the sealed bytes under a
live claim (the in-loop footprint claim fails on the corrupted row
--- the checker refuses a lying deadline).

Beside the battery, an independent replay closes the artefact pair:
an exact-rational drain simulation --- written against the model
tables, not the certificate code, with its expectations committed
before the run --- executes seeded demonic schedules at lot counts
$1$, $2$, $5$, and $84$, and confirms: every drain terminates
cleanly with all lots shipped and machine populations restored to
$(2,2,1)$; total firings equal $638N$ \emph{exactly} at every lot
count and every schedule tested (the budget is tight, as the
telescoping predicts); every firing satisfies every touched
certificate row in its reading scale; the place invariants pin;
the congestion fold is monotone with ticks inside the
$\varphi$-space budget; every seeded schedule meets the deadline,
with the arrival fold reaching zero exactly at the $N$th arrival;
the adversarial schedule that defers the last $\mathrm{finish}_5$
attains $638N - 12$ \emph{exactly} at every tested lot count (the
deadline bound is tight dynamically, as the split predicts); and
the broken-gap mutant violates its row
at the \emph{first} mutated firing --- the dynamic kill agreeing
with the static refusal. The replay rides the toolkit's workshop
today and lands in the reproduction spine with the vendoring motion
above.

\subsection{The carved clock: the same rows in real time}
\label{sec:minifab-carved}

The clock rows of section~\ref{sec:minifab-sdcpn} support two
readings, and everything so far in this section has used the first:
\emph{counting time} --- makespans measured in firings, the
discrete-time unfolding's step index --- where total work is
scheduler-invariant ($638N$ firings, always: the work potential's
telescoping) and \emph{when} the work happens is the only
nontrivial question. This subsection and the four that follow read
the same declared clocks the other way --- \emph{metric time}, the
tick meter --- and push the open-seat story to its natural end: what the
wall clock rewards, exactly how much, a policy that attains the
optimum for $N \le 7$ and the conjectured value at the replayed
larger counts $N = 8, \dots, 12$ and $N = 84$, a machine-checked sense in which ``optimal'' is the
decision-theoretically right word, a deliberately untrusted policy
run inside a certified envelope, and one certificate that spans lot
counts from $7$ to $10^{7}$. The subject is unchanged --- the
lot-flow core, with the model boundary stated at this section's end
inherited whole.

The metric-time semantics is already in the declarations: every
transition's clock is carved to the window $[1,1]$, so each row
dwells exactly one tick and fires at most once per tick --- the
per-row meter law; simultaneous instances resolve through the demonic scheduler
seat, so the only freedom left is genuinely dispatch freedom. Both
closed-composite encodings of section~\ref{sec:minifab-synthesis}
reseal under this reading with row tallies identical to the
counting-time containers --- the certificates themselves never
change.

Two transfers move the counting-time results onto the wall clock.
First, upper bounds transfer by a \emph{no-idle} argument: before
the goal, a tick with an empty due set --- no enabled transition
instance whose one-tick dwell has elapsed --- would freeze the
state forever, so along every goal-reaching schedule each tick
fires at least one row, and ticks-to-all-arrivals is bounded by
pre-goal firings: at most $626N$ ticks under the synthesized gate
and $638N - 12$ open, carried by the \emph{same} geolog-checked
rows plus two named glue clauses, argued --- the meter/no-idle
lemma and
enabledness below the goal. Second, there is a wall-clock
\emph{floor} over all dispatch policies of any kind, supervisory or
never-idle:
\[
\text{ticks to all arrivals} \;\ge\; 255N + 371
\qquad (21{,}791 \text{ at the weekly } N = 84),
\]
because step~5's service tick is one row --- at most one firing per
tick, $255$ firings per lot --- behind an uncompressible
$370$-event first-lot pipeline fill, plus the final transfer into
the arrival queue. The bound is Z3-verified below $N = 8$, and
its machines-deleted variant is refused too (checked at $N \le 2$):
machine capacity is not the binding constraint; the clock law and
per-lot precedence are. (Tight at $N = 1$: the singleton trajectory
takes exactly $626$ ticks.)

Measured against that floor at $N = 84$, the dispatch levers
separate cleanly. The synthesized arrival gate, whose counting-time
value is certified (section~\ref{sec:minifab-synthesis}), moves
wall-clock arrivals $0.0\%$: greedy dispatch reaches all-arrivals
in $22{,}019$ ticks with the gate on or off --- the gate's value
was always the certified worst-case floor, a guarantee over every
schedule, not a typical-case speedup. Pull order is a $1.83\times$
lever (an adversarial pull priority needs $40{,}406$ ticks,
$+83.5\%$); and greedy itself sits $228$ ticks above the floor.
Real-time value lives almost entirely in the dispatch seat ---
which is why the optimum is a policy story, next, and why the seat
then widens (section~\ref{sec:minifab-fulldispatch}).

\subsection{Optimal throughput: exact for up to seven lots}
\label{sec:minifab-optimum}

For $N = 1, \dots, 7$, the optimum, measured as the tick of the
$N$-th arrival into the finishing queue $q_6$ (not the tick at which
all lots have shipped), is
\[
T_{\mathrm{opt}}(N) \;=\; 255N + 371 + d(N),
\qquad d = 0\ (N \le 2),\quad 2\ (N = 3),\quad 4\ (N \ge 4),
\]
and is established from both sides. For $N \ge 8$ the same expression
is only the conjectured optimum: $\pi^{\ast}$ attains it at the
replayed counts below, but optimality on the net is open.
\emph{Erratum (2026-10-19): optimality is established for $N \le 7$
only.} For $N = 1, \dots, 7$ it is closed from both sides: no schedule
meets $T_{\mathrm{opt}} - 1$ on a stated relaxation of the net that
also imposes a queue-order reduction (unsatisfiable: LRAT proofs are
checked by the CakeML-verified checker \texttt{cake\_lpr}, and an
independent integer SMT encoding provides a separate check), and a schedule at
$T_{\mathrm{opt}}$ exists and is re-validated. Two pen-and-paper
arguments carry the lower side to every policy: that the relaxation
contains every run of the net, and an exchange lemma justifying the
queue-order reduction; for $N=1,\dots,5$ there are also UNSAT
certificates for the encoding without the queue-order reduction. All
of this is about the simplified net of this section, in which at
most one lot per process step may start, receive a unit of service,
or finish per tick, even at two-machine stations (model note at this
section's end).

On the upper side an explicit supervisor family $\pi^{\ast}$,
drawn seated in its interface as Figure~\ref{fig:minifab-pistar}
(hold the first few diffusion starts on one machine with two-tick
handovers, then alternate pre-positioned; run the step-1 services
back to back), attains the displayed expression in replay at
$N = 1, \dots, 12$ and at $N = 84$; at each replayed $N \ge 4$ that
value is $255N + 375$: $21{,}795$ ticks at the weekly count, within
$4$ ticks ($0.018\%$) of the machine-free floor above; the displayed
expression divided by $N$ tends to $255$ ticks per lot, which is an
algebraic limit, not an all-$N$ replay result. The attainment is a \emph{replay}:
$\pi^{\ast}$'s explicit firing table replays to arrival at exactly
$21{,}795$ against the decoded container plus a named transcription
of the net rows the container does not seal (the transcription is
\texttt{semicon/minifab-floor/glc-replay/src/replay.rs} in the
experiments repository; section~\ref{sec:minifab-exact-status}),
under the meter law, every
hold accounted. At $N = 1, \dots, 12$ the replay runs on the same
decoded net with exactly two sealed facts changed, the pending-lot
count and the unlock threshold, both sealed as $84$; it is not a
replay of a sealed $N$-lot container. (The original archive also
replayed $N = 30$, at $8{,}025$ ticks, against the producer's
in-memory model rather than the decoded container; that receipt is
kept but not re-run.) For $3 \le N \le 7$ the two-sided optimum also
shows that the machine-free floor $255N + 371$ is not attained. A
machine-contiguity floor of $255N + 375$ for $N \ge 4$ is proved in
Lean~4 over mathlib --- a five-case analysis of the first four
diffusion entrants, showing any schedule in that abstract model
faster than $255N + 375$ leaves the third or fourth lot's service without a contiguous
machine window, three of the five cases dying inside the first
pigeonhole handover the tick-meter forces, with adversarial mutants
refused by the kernel --- but it is proved for an \emph{abstract
schedule model} onto which the net does not project exactly: the net
also allows a tick in a lot's own start round and a same-round
handover of a diffusion machine. So for $N \ge 8$ the floor is not
yet proved for the net, and optimality there is open, although
the replays at $N = 8, \dots, 12$ and $N = 84$ attain $255N + 375$.
Those replays establish the upper side at those counts, not an
all-policy lower bound. The Lean analysis accounts for the $+4$ in
the abstract lower bound by diffusion-machine handovers; it does not
transfer that bound to the net. An earlier edition of this
subsection claimed a Lean-checked floor for the net and a two-sided
optimum at every $N$; both overclaimed.
Separately, the counter/meter bookkeeping behind
the machine-free floor landed as geolog-checked rows --- this
container family's first checker-side clock-law rows.

Two consequences deserve their own sentences. \emph{Deliberate
idling pays, among the policies tested}: at the weekly count the
best tested never-idle (work-conserving) policy needs $255N + 455$
ticks, $80$ behind $\pi^{\ast}$, with its never-idle law verified at
every pre-goal tick of the replay, and a measured greedy baseline
runs $224$ ticks behind; whether some untested work-conserving policy
closes that gap is open. Gating the \emph{last} step moves arrivals
$0.0\%$, while holding the \emph{diffusion starts} recovers the full
gap, and replaying $\pi^{\ast}$ with holds disallowed is
\emph{refused} by the engine at its first delayed diffusion start:
where the supervisor holds is what its last $80$ ticks are made of.
And \emph{the decision-theoretic reading is discharged, not
gestured at}: with deterministic closed dynamics the belief state
is a point mass, and imprecise-decision-theoretic selection at the
widest-uncertainty locus provably collapses to makespan
minimization --- the reduction Lean-checked against the
mechanized imprecise-decision-theory development, with a concrete
refutation cell showing the collapse fails without determinism.
``The synthesized policy is decision-theoretically correct''
therefore means optimal throughput --- which the receipts above
establish for $N \le 7$ and leave open beyond, and the seat-regret
face of the specification vocabulary
(section~\ref{sec:spec-lang-regret}) is where the language states
exactly this shape: the regret of a filler against its own seat,
here $0$ for $N \le 7$ (established two-sided as above) and
unproved beyond, where the Lean floor covers only the abstract
schedule model.

Figure~\ref{fig:minifab-pistar} closes the subsection by drawing
$\pi^{\ast}$ as what it is: a filler for the supervisor seat, the
Petri page's dashed box grown to nearly the whole page, one
enormous Petri transition inside --- the policy's firing kernel,
its interior a plate-notation hypernet over the seat's four
observables. The page and its caption carry the policy's entire
content: no code, no lookup table, no appeal to the replay --- two
threshold cells and two piecewise-constant schedule decoders,
readable in one look.


\newgeometry{hmargin=1cm,vmargin=1.2cm}
\begin{landscape}
\setlength{\parskip}{0pt}\vspace*{\fill}
\centering
\begin{tikzpicture}[x=1cm, y=0.86cm,
  place/.style={circle, draw=black!80, line width=0.7pt,
    minimum size=8.5mm, inner sep=0pt},
  trans/.style={rectangle, draw=black!80, line width=0.7pt, fill=white,
    minimum width=9.5mm, minimum height=4.5mm, inner sep=1.5pt,
    font=\footnotesize},
  every label/.style={font=\footnotesize, label distance=1.5pt},
  arc/.style={line width=0.55pt, draw=black!80,
    arrows={Bar[width=2.1mm]}-{Stealth[length=2.4mm]}},
  carc/.style={arc, draw=czero},
  rarc/.style={line width=0.55pt, draw=black!80,
    arrows={Stealth[length=2.4mm]}-{Stealth[length=2.4mm]}},
  elide/.style={line width=0.55pt, draw=czero, dotted,
    -{Stealth[length=2.4mm]}},
  cell/.style={draw=hndeep, fill=hnaccent!50, rounded corners=1pt,
    font=\footnotesize, inner sep=3pt, minimum height=4.2mm},
  fold/.style={draw=hndeep, fill=hnaccent!50, rounded corners=1pt,
    font=\footnotesize, inner sep=3pt, minimum height=4.2mm,
    align=center},
  defcell/.style={cell, dashed},
  plate/.style={draw=hnmid, line width=0.7pt, rounded corners=2pt,
    fill=hnaccent!18, inner sep=5pt},
  platelab/.style={font=\tiny\ttfamily, text=hnmid, inner sep=1pt},
  wire/.style={-, line width=0.5pt, draw=black!75},
  wlab/.style={font=\scriptsize, inner sep=1pt},
  port/.style={font=\small\itshape, inner sep=1.5pt},
  ann/.style={font=\scriptsize, text=black!70, align=left,
    anchor=north west, inner sep=0pt}]

\node[ann] at (-0.2,1.55) {\emph{the plant of the Petri page, elided
  to the diffusion re-entrancy (steps 1 and 5 share the machine
  pool); everything ungated is omitted}};

\node[place, label=above:{$q_1$}] (q1) at (1.0,0) {};
\node[trans] (s1) at (3.0,0) {$\mathrm{start}_1$};
\node[place, draw=wkA, line width=0.9pt, label=above:{$p_1$}] (p1) at (4.8,0) {};
\node[trans] (f1) at (6.6,0) {$\mathrm{finish}_1$};
\node[place, label=above:{$q_5$}] (q5) at (12.6,0) {};
\node[trans] (s5) at (14.4,0) {$\mathrm{start}_5$};
\node[place, draw=wkE, line width=0.9pt, label=above:{$p_5$}] (p5) at (16.2,0) {};
\node[trans] (f5) at (18.0,0) {$\mathrm{finish}_5$};
\node[place, label=above:{$q_6$}] (q6) at (19.8,0) {};

\draw[arc] (q1) -- (s1);
\draw[arc, draw=wkA] (s1) -- (p1);
\draw[carc] (p1) -- (f1);
\draw[elide] (f1) to[bend left=14]
  node[wlab, above, text=czero] {steps 2--4 (implantation, lithography) --- ungated} (q5);
\draw[arc] (q5) -- (s5);
\draw[arc, draw=wkE] (s5) -- (p5);
\draw[carc] (p5) -- (f5);
\draw[arc] (f5) -- (q6);
\draw[elide] (q6.east) -- ++(1.45,0)
  node[wlab, right, text=czero, align=left]
  {step 6, ship\\(the arrival gate of the Petri page)};

\node[place, label={[font=\footnotesize]180:{$m_{\mathrm{diff}}$}}]
  (mdiff) at (8.7,-1.05) {};
\fill[black] ([xshift=-1.2mm]mdiff.center) circle (0.065);
\fill[black] ([xshift= 1.2mm]mdiff.center) circle (0.065);
\draw[arc, rounded corners=5pt] (mdiff.170) -| ([xshift=-2mm]s1.south);
\draw[arc, rounded corners=5pt] (mdiff.10)  -| ([xshift=-2mm]s5.south);
\draw[arc, rounded corners=5pt] ([xshift=2mm]f1.south) |- (mdiff.140);
\draw[arc, rounded corners=5pt] ([xshift=2mm]f5.south) -- ++(0,-0.62)
  -| (mdiff.40);

\begin{pgfonlayer}{holebg}
  \draw[dashed, line width=0.7pt, draw=black!75]
        (-0.2,-12.1) rectangle (24.6,-2.05);
\end{pgfonlayer}
\node[font=\footnotesize\itshape] at (12.2,-11.8)
  {Supervisor seat --- filled by the synthesized policy $\pi^{\ast}$};

\node[place, fill=white,
  label={[font=\footnotesize]160:{$g_1$}}] (g1) at (3.0,-2.05) {};
\fill[black] (g1.center) circle (0.065);
\node[place, fill=white,
  label={[font=\footnotesize]160:{$g_5$}}] (g5) at (14.4,-2.05) {};
\draw[rarc, draw=wkA] (g1) -- ([xshift=2mm]s1.south);
\draw[rarc, draw=wkE] (g5) -- ([xshift=2mm]s5.south);
\node[place, fill=white,
  label={[font=\footnotesize]200:{$c_1$}}] (c1) at (6.6,-2.05) {};
\node[place, fill=white,
  label={[font=\footnotesize]340:{$c_5$}}] (c5) at (18.0,-2.05) {};
\draw[arc, draw=wkA] (f1.south) -- (c1.north);
\draw[arc, draw=wkE] ([xshift=-2mm]f5.south) -- ++(0,-0.35) -| (c5.north);
\node[rectangle, draw=black!80, line width=0.7pt, fill=white,
  minimum size=3.5mm, inner sep=0pt,
  label={[font=\footnotesize]200:{$q_5$ pool}}] (tq5) at (12.0,-2.05) {};
\node[rectangle, draw=black!80, line width=0.7pt, fill=white,
  minimum size=3.5mm, inner sep=0pt,
  label={[font=\footnotesize]200:{$p_5$ pool}}] (tp5) at (16.8,-2.05) {};
\draw[rarc] (q5.south west) -- (tq5.north);
\draw[rarc] (p5.south east) -- (tp5.north);

\begin{pgfonlayer}{boxbg}
  \draw[line width=0.9pt, draw=black!80, fill=white]
        (1.2,-11.35) rectangle (24.2,-3.35);
\end{pgfonlayer}
\node[font=\footnotesize\bfseries, anchor=north east, align=left,
  text width=5.35cm] at (24.05,-3.55)
  {$\pi^{\ast}$ --- one transition, fired every tick: the kernel
   reads the four boundary observables and rewrites the two gates};

\draw[rarc] (c1.south) -- (c1.south |- 0,-3.35);
\draw[rarc] (c5.south) -- (c5.south |- 0,-3.35);
\draw[rarc] (tq5.south) -- (tq5.south |- 0,-3.35);
\draw[rarc] (tp5.south) -- (tp5.south |- 0,-3.35);
\draw[arc, draw=wkA] ([xshift=-1.2mm]g1.south) -- ++(0,-1.0);
\draw[arc, draw=wkA] ([xshift=1.2mm]g1.south) ++(0,-1.0) -- ++(0,1.0);
\draw[arc, draw=wkE] ([xshift=-1.2mm]g5.south) -- ++(0,-1.0);
\draw[arc, draw=wkE] ([xshift=1.2mm]g5.south) ++(0,-1.0) -- ++(0,1.0);

\node[port] (Pg1) at (3.0,-3.75) {$g_1$};
\node[port] (Pc1) at (6.6,-3.75) {$c_1$};
\node[port] (Pq5) at (12.0,-3.75) {$q_5$};
\node[port] (Pg5) at (14.4,-3.75) {$g_5$};
\node[port] (Pp5) at (16.8,-3.75) {$p_5$};
\node[port] (Pc5) at (18.0,-3.75) {$c_5$};

\node[cell] (sc1) at (11.2,-5.7) {$1$};
\node[fold] (sf1) at (13.0,-5.7) {fold $\Sigma$};
\node[cell] (scw) at (16.9,-5.7) {$w$};
\node[fold] (sfw) at (16.9,-7.3) {fold $\Sigma$};
\begin{pgfonlayer}{plates}
  \node[plate, fit=(sc1)] (pl1) {};
  \node[plate, fit=(scw)] (plw) {};
\end{pgfonlayer}
\node[platelab, anchor=north east] at (pl1.south east) {plate $q_5$};
\node[platelab, anchor=north east] at (plw.south east) {plate $p_5$};
\draw[wire, rounded corners=4pt] (Pq5.south) -- ++(0,-0.55)
  -- node[wlab, above left, pos=0.6] {$[q_5]$} (sc1.north);
\draw[wire, rounded corners=4pt] (Pp5.south)
  -- node[wlab, right, pos=0.45] {$[p_5]$} (scw.north);
\draw[wire] (sc1.east) -- (sf1.west);
\draw[wire] (scw.south) -- (sfw.north);

\node[defcell] (thc) at (21.6,-7.3) {margin $\theta(c_5)$};
\node[defcell] (phc) at (7.0,-8.6) {quota $\varphi(c_1, c_5)$};
\draw[wire, rounded corners=4pt] (Pc5.south) -- (18.0,-7.3)
  -- (thc.west);
\draw[wire, rounded corners=4pt] (Pc1.south) -- (6.6,-7.0)
  -- node[wlab, left, pos=0.85] {$c_1$} (6.6,-8.35);
\draw[wire, rounded corners=4pt] (18.0,-7.3) -- (18.0,-9.6)
  -- node[wlab, above, pos=0.65] {$c_5$} (7.0,-9.6) -- (phc.south);

\node[cell] (G5) at (15.4,-8.6)
  {$\ell \ge 1 \;\wedge\; u \le \theta$};
\node[cell] (G1) at (3.9,-6.6) {$c_1 < \varphi$};
\draw[wire, rounded corners=4pt] (sf1.east) -- ++(0.35,0)
  -- node[wlab, left, pos=0.55] {$\ell$} (14.05,-8.35)
  -- (G5.west |- 0,-8.5) -- (G5.west);
\draw[wire, rounded corners=4pt] (sfw.south)
  -- node[wlab, right, pos=0.35] {$u$} (16.9,-8.0)
  -- (16.1,-8.0) -- (16.1,-8.25);
\draw[wire, rounded corners=4pt] (thc.south)
  |- node[wlab, above, pos=0.78] {$\theta$} (G5.east);
\draw[wire, rounded corners=4pt] (phc.west)
  -- node[wlab, above, pos=0.6] {$\varphi$} (3.9,-8.6) -- (G1.south);
\draw[wire, rounded corners=4pt] (6.6,-5.2) -- (3.9,-5.2) -- (G1.north);
\draw[wire, rounded corners=4pt] (G5.north) -- (15.4,-4.6)
  -- node[wlab, above, pos=0.35] {open/hold} (14.4,-4.6) -- (Pg5.south);
\draw[wire, rounded corners=4pt] (G1.west) -- (3.0,-6.6)
  -- node[wlab, left, pos=0.6] {open/hold} (Pg1.south);

\node[ann, text width=21.6cm] at (1.7,-10.15)
  {closed forms of the two dashed defined-symbol cells ---
   $\theta$ (diffusion-slot margin): the first $1 + \min(N-2,\,2)$
   step-5 slots run consecutively on one machine --- $\theta = 0$, the
   two-tick handover; every later slot alternates machines,
   pre-positioned one tick early --- $\theta = 1$.\quad
   $\varphi$ (feed quota): unbounded while $c_1 < \min(N,\,4)$
   (step-1 services back-to-back); afterwards one admission per slot
   window (services 5, 6 into the first freed windows, then
   alternating).};
\end{tikzpicture}\par\medskip
\begin{minipage}{0.94\linewidth}
\platecaption{The synthesized policy $\pi^{\ast}$ (optimal for $N \le 7$), drawn
as what it is: a filler for the supervisor seat. The dashed seat of
Figure~\ref{fig:minifab-sdcpn-petri} is here nearly the whole page,
and the plant survives only as the gated neighbourhood along the top
--- the diffusion re-entrancy, where $\mathrm{start}_1$ and
$\mathrm{start}_5$ compete for the same two machines. Straddling the
boundary, the seat's declared interface: the two pooled reads, the two
firing counters, and the two gate places --- a strict subset of the
fifteen-observable, four-act full-dispatch interface of
section~\ref{sec:minifab-fulldispatch}, and the only wires
$\pi^{\ast}$ needs. Inside sits one enormous Petri transition: the
policy's firing kernel, applied once per tick, reading the four
observables and rewriting the two gates. Its interior is a hypernet in the plate notation of
Figure~\ref{fig:minifab-cert-hypernets}: the pooled observables are
two plated folds ($\ell = \Sigma\,1$ over $[q_5]$, the queue length;
$u = \Sigma\,w$ over $[p_5]$, the in-service residual ---
pre-positioning keeps at most the incumbent in $p_5$, so the sum reads
its residual), the counters decode the schedule phase through the two
\emph{dashed} cells ($\theta$, $\varphi$ --- defined symbols, not
primitive generators: piecewise-constant functions of the counters,
closed forms in the box), and two threshold cells (region-membership)
emit the gate bits. Holding $\mathrm{start}_5$ until $u \le \theta$
pre-positions each diffusion slot; holding $\mathrm{start}_1$ at the
quota keeps a diffusion machine free for step~5 --- deliberate idling,
the one capability the never-idle class lacks. What this filler
attains, each claim at its tier (section~\ref{sec:minifab-optimum}):
the $N$-th arrival into the last queue at $255N + 371 + d(N)$
ticks ($d=0,0,2$, then $4$) for $N=1,\dots,12$ and $N=84$.
The weekly replay reaches $21{,}795$ against the decoded
$N=84$ container plus a transcription of the rows it does not
seal; the $N=1,\dots,12$ replays use that model with the
pending-lot count and unlock threshold substituted. Optimality
on the simplified net is established for $N=1,\dots,7$
and is open beyond; the
Lean machine-contiguity floor is for an abstract schedule model. The best tested never-idle policy needs
$255N + 455$ at $N = 84$ --- these two gates are worth $80$ ticks over it
(whether an untested work-conserving policy closes that gap is open), and replaying $\pi^{\ast}$ with
holds disallowed is refused at its first delayed diffusion start.
Under deterministic closed dynamics, imprecise-decision-theoretic
selection provably collapses to makespan minimization (Lean-checked)
--- so where $\pi^{\ast}$ is optimal ($N \le 7$), this page is also
the decision-theoretic optimum. Initial
marking drawn: $g_1$ open ($c_1 = 0 < \varphi$), $g_5$ shut
($\ell = 0$).}{fig:minifab-pistar}
\end{minipage}
\par
\vspace*{\fill}
\thispagestyle{empty}
\end{landscape}
\restoregeometry

\subsection{Widening the seat: full dispatch as a filler}
\label{sec:minifab-fulldispatch}

The supervisor seat of section~\ref{sec:minifab-synthesis} reads
one observable and drives one gate wire. That narrowness was an
encoding shim, not a formalism limit: an interface always was a
pair of port lists, and the same comb construction re-seals the
same single seat with a \emph{full-dispatch} interface --- fifteen
observables (six queue lengths, three machine-idle flags, six
in-process work counts) and four act ports (three per-station pull
decisions plus token selection) --- around the unchanged plant. The
closed net is the plant net, so the exhibit's twelve-certificate
family binds every filler \emph{a fortiori}, geolog-checked: the
sealed full-dispatch container re-verifies with the full row
roster, and the certificate never mentions the policy.

Two findings from exercising it. Token selection is a \emph{second}
dispatch freedom that determinism requires: a pull priority alone
does not close the system; fixing both yields one trajectory per
policy, verified by replaying residual-order permutations to
identical traces. And across the $16$-policy deterministic grid
($8$ pull priorities $\times$ $2$ selection rules, $N \in
\{1, 7, 84\}$; Table~\ref{tab:minifab-dispatchgrid}), only the
diffusion-station pull preference binds the outcome --- the
implantation and lithography columns of the table never move a
number --- while token selection swings roughly $7{,}000$ ticks and
\emph{crosses} with pull preference; the best cell lands on
$255N + 455$ exactly at $N \in \{7, 84\}$ --- the best tested
work-conserving value of the previous subsection, consistent with it,
since grid policies never deliberately idle. Threshold-style gate supervisors provably
cannot express per-station pull preference --- argued: a gate place
maintained by fixed-weight fused arcs is an affine functional of
the marking, read through upward-closed thresholds, and a pull
preference's acceptance region is a union no affine threshold
carves. Dispatch is a genuine seat species, not sugar over gates
--- enablement carving and seat selection are different composition
operations.

\begin{table}[t]
\centering
\footnotesize
\setlength{\tabcolsep}{2.2pt}\newcommand{\bestcell}[1]{\cellcolor{accent!30}\textbf{#1}}\begin{tabular}{@{}cccl rrrr rr@{}}
\toprule
\multicolumn{4}{@{}c}{dispatch policy} &
\multicolumn{4}{c}{ticks per arrival} &
\multicolumn{2}{c@{}}{reference / listed value (\%)} \\
\cmidrule(r){1-4}\cmidrule(lr){5-8}\cmidrule(l){9-10}
$m_{\mathrm{diff}}$ pulls & $m_{\mathrm{impl}}$ pulls &
$m_{\mathrm{lith}}$ pulls & selection &
$N{=}1$ & $N{=}7$ & $N{=}84$ & $N{=}10^{6}$ & $N{=}84$ & $N{=}10^{6}$ \\
\midrule
\multicolumn{4}{@{}l}{\textbf{ours: $\pi^{\ast}$} (synthesized holds --- not never-idle)} &
626.0 & \bestcell{308.6} & \bestcell{259.5} & \bestcell{255.000375} &
\bestcell{100.00$^{\dagger}$} & \bestcell{100.00$^{\dagger}$} \\
\midrule
$q_5$ & $q_2$ & $q_3$ & most-advanced  & 626.0 & 320.0 & 260.4 & --- & 99.63 & --- \\
$q_5$ & $q_2$ & $q_6$ & most-advanced  & 626.0 & 320.0 & 260.4 & --- & 99.63 & --- \\
$q_5$ & $q_4$ & $q_3$ & most-advanced  & 626.0 & 320.0 & 260.4 & --- & 99.63 & --- \\
$q_5$ & $q_4$ & $q_6$ & most-advanced  & 626.0 & 320.0 & 260.4 & --- & 99.63 & --- \\
$q_5$ & $q_2$ & $q_3$ & least-advanced & 626.0 & 384.7 & 343.2 & --- & 75.61 & --- \\
$q_5$ & $q_2$ & $q_6$ & least-advanced & 626.0 & 384.7 & 343.2 & --- & 75.61 & --- \\
$q_5$ & $q_4$ & $q_3$ & least-advanced & 626.0 & 384.7 & 343.2 & --- & 75.61 & --- \\
$q_5$ & $q_4$ & $q_6$ & least-advanced & 626.0 & 384.7 & 343.2 & --- & 75.61 & --- \\
$q_1$ & $q_2$ & $q_3$ & least-advanced & 626.0 & 445.0 & 397.1 & --- & 65.34 & --- \\
$q_1$ & $q_2$ & $q_6$ & least-advanced & 626.0 & 445.0 & 397.1 & --- & 65.34 & --- \\
$q_1$ & $q_4$ & $q_3$ & least-advanced & 626.0 & 445.0 & 397.1 & --- & 65.34 & --- \\
$q_1$ & $q_4$ & $q_6$ & least-advanced & 626.0 & 445.0 & 397.1 & --- & 65.34 & --- \\
$q_1$ & $q_2$ & $q_3$ & most-advanced  & 626.0 & 448.6 & 477.4 & --- & 54.35 & --- \\
$q_1$ & $q_2$ & $q_6$ & most-advanced  & 626.0 & 448.6 & 477.4 & --- & 54.35 & --- \\
$q_1$ & $q_4$ & $q_3$ & most-advanced  & 626.0 & 448.6 & 477.4 & --- & 54.35 & --- \\
$q_1$ & $q_4$ & $q_6$ & most-advanced  & 626.0 & 448.6 & 477.4 & --- & 54.35 & --- \\
\midrule
\multicolumn{4}{@{}l}{\emph{reference ($255N{+}375$ for $N \ge 4$; proved optimal for $N \le 7$, not proved a floor for $N \ge 8$)}} &
\emph{626.0} & \emph{308.6} & \emph{259.5} & \emph{255.000375} &
\emph{100} & \emph{100} \\
\bottomrule
\end{tabular}
\caption{All seventeen policies of these subsections, benchmark-style
(the synthesized policy itself is drawn as the seated filler in
Figure~\ref{fig:minifab-pistar}): the sixteen never-idle full-dispatch
cells (three binary pull preferences $\times$ token selection;
deterministic replay, verified by residual-order permutations) and the
synthesized $\pi^{\ast}$, which is not a row of the grid because it
deliberately idles --- the one capability the never-idle class lacks,
worth $80$ ticks over the best tested never-idle cell. Read the parameter columns: only
the diffusion column ever moves a number; the implantation and
lithography preferences are provably irrelevant at this roster. Token
selection crosses with pull preference (most-advanced wins under
$q_5$, loses under $q_1$). Units: ticks per arrival ($T/N$, the
throughput cost of one lot; the asymptotic ideal is $255$); at $N = 1$ every tabulated policy attains the optimum of $626$
ticks. At $N = 7$, $\pi^{\ast}$ attains the proved optimum of
$2{,}160$ ticks ($308.6$ ticks per arrival, rounded). The
percentage columns divide the reference time by the listed policy
time at the same $N$, before rounding. $^{\dagger}$At $N = 84$
the denominator for $\pi^{\ast}$ is its replayed time of
$21{,}795$ ticks. At $N = 10^{6}$ the only policy entry is a
formula evaluation: $(255N+375)/N = 255.000375$ ticks per
arrival, not a replay. Its $100.00\%$ compares the formula with
itself. Neither percentage column certifies optimality on the net;
an earlier edition called these columns certified at every $N$.
The sixteen never-idle cells are deterministic replays and have
no million-lot entries (Table~\ref{tab:minifab-nscaling}). The single-supervisor lanes of section~\ref{sec:minifab-carved}
show the same pull-order lever at the same magnitude ($40{,}406$ ticks
under an adversarial greedy pull there).}
\label{tab:minifab-dispatchgrid}
\end{table}

\subsection{An untrusted policy in a certified envelope}
\label{sec:minifab-envelope}

The envelope certificates --- safety, floors, the deadline ---
never read the policy: they quantify over every filler of the seat,
which is why the policy may be \emph{anything}, including a black
box nobody should trust. (For replay-based performance claims the dispatcher need not be
trusted: what is checked is its trace. For $\pi^{\ast}$ the replay
model and its relation to the sealed instance are specified in
section~\ref{sec:minifab-optimum}. The million-lot entry is a
formula evaluation, not a checked trace. Of the never-idle rows of
Table~\ref{tab:minifab-dispatchgrid}, only the best value at
$N = 84$, $21{,}875$ ticks, is also reached by a work-conserving
schedule replayed against the decoded container; the other cells
are replays by the producer's strategy tier, which is trusted Rust
and cannot be reproduced from public sources alone.) A third closure of the same seat
demonstrates exactly that:
the open comb filled by an uncertified lookahead heuristic (a
stand-in for a learned or LLM dispatcher) composed with a certified
switching envelope --- a decision-budget bench with a certified
recovery region. In replay: the heuristic ran free, was never
benched, completed, and attained the envelope's arrival bound
$638N - 11$ \emph{exactly} at the tested $N = 84$ --- one tick
above the open plant's $638N - 12$, the closed composite's own
unlock firing --- so the bound is witnessed tight, not merely
sound; a deliberate idler hit the budget boundary and was benched,
with recovery at exactly $1{,}009$ firings; a rogue policy was
benched after eight named in-loop refusals with the plant
untouched; and disabling the bench stalls the line --- the switch
is load-bearing, not decorative. After all runs the sealed
container re-verified \emph{byte-identical}: the certificate is
carried entirely by the envelope, geolog-checked, and no behaviour
of
the filler --- competent, idle, or hostile --- can touch it.

\subsection{One certificate, every lot count}
\label{sec:minifab-nscaling}

The deadline certificates of this section are stated once and
instantiated at any lot count~$N$: the sealed container's row
obligations (footprint, flow, floor, canonicity) never mention $N$
--- they quantify over \emph{all} markings --- and $N$ enters only
as three integers in the container's own data (the initial lot
count, the unlock guard's threshold multiplicity, and hence the
instantiated horizon, which the reader recovers from the marking
row by one dot product with the arrival dictionary).
Table~\ref{tab:minifab-nscaling} instantiates, reseals, and
re-verifies the pair of exhibits at lot counts spanning six orders
of magnitude. The certificate's byte size moves only with the
decimal width of those integers (six bytes across the whole sweep);
the checker's wall time is flat at $\sim$2\,ms; and the objects the
claims quantify over grow combinatorially --- at $N = 10^{6}$ the
model has at least $2\times10^{38}$ reachable markings and at least
$10^{2\times10^{7}}$ distinct maximal schedules, and the checker
never visits any of them.

\begin{table}[t]
\centering
\small
\begin{tabular}{@{}rrrrrll@{}}
\toprule
$N$ & $T^{*}=626N$ & $638N-12$ & cert.\ (B) & check (ms) &
markings $\geq$ & schedules $\geq$ \\
\midrule
$7$      & $4{,}382$         & $4{,}454$         & $3{,}105$ & $1.9$ & $3.4\times10^{3}$  & $10^{123}$ \\
$84$     & $52{,}584$        & $53{,}580$        & $3{,}105$ & $1.9$ & $8.1\times10^{9}$  & $10^{1{,}719}$ \\
$10^{3}$ & $626{,}000$       & $637{,}988$       & $3{,}107$ & $1.9$ & $2.0\times10^{17}$ & $10^{20{,}469}$ \\
$10^{6}$ & $626{,}000{,}000$ & $637{,}999{,}988$ & $3{,}109$ & $1.8$ & $2.0\times10^{38}$ & $10^{2.0\times10^{7}}$ \\
$10^{7}$ & $6{,}260{,}000{,}000$ & $6{,}379{,}999{,}988$ & $3{,}111$ & $1.9$ & $2.0\times10^{45}$ & $10^{2.0\times10^{8}}$ \\
\bottomrule
\end{tabular}
\caption{The synthesized-policy deadline certificate instantiated
at lot count $N$. Columns: the certified all-arrivals horizon under
the synthesized gate ($T^{*} = 626N$ steps, exact on every
completing run); the no-policy worst case ($638N-12$); the sealed
container's size in bytes (supervised exhibit; the no-policy
container runs $2{,}963$--$2{,}966$\,B over the same sweep); the
checker's wall time (median of 7 re-discharges of the sealed bytes
by the production engine's in-loop evaluator walk, single core of
an NVIDIA GB10 (Grace, aarch64) under concurrent load; all samples
across all rows fell in 1.7--2.0\,ms); a reachable-marking lower
bound; and a lower bound on distinct maximal schedules the claims
quantify over. Intermediate lot counts $10^{4}$ and $10^{5}$ were
measured identically (containers $3{,}107$ and $3{,}109$\,B;
medians $1.9$\,ms).}
\label{tab:minifab-nscaling}
\end{table}

\paragraph{What is $N$-parametric and what is not --- the honest
split.} The \emph{certificate} and the \emph{checker} are
$N$-parametric in the strong sense: the same $36$ footprint rows
(of the $150$ footprint sequents the producer checks, the sealed
checker evaluates only these $36$, the six tick transitions' rows),
the same dictionaries, the same claim roster discharge at every $N$
(identical tally vectors across the sweep), and at the landed weekly
count the resealed container reproduces the landed exhibit exactly.
Re-instantiation at a new $N$ does re-run the \emph{producer}: it
reseals the container ($0.002$\,s at $N=84$, $106$\,s at $N=10^{7}$
--- the producer's in-memory form expands the weight-$N$ guard
unarily, a linear-cost implementation artifact of the untrusted
side that never touches the sealed bytes or the checker), and
re-derives the exact-rational producer receipts (the
$=$-direction's drop identities; $0.002$\,s at $N=84$, $9.8$\,s at
$N=10^{6}$). Dynamic attainment witnesses (replays arriving at
exactly $626N$) were re-run through $N=10^{4}$, i.e.\
$6{,}260{,}000$-step runs; beyond that the replay grinder's
quadratic cost makes simulation the bottleneck --- which is the
point of certifying.

\paragraph{Claim grades at every $N$.} The within-$T^{*}$ direction
rides geolog-checked rows at every $N$ (the row obligations are
$N$-free; the horizon instantiates from the container's marking
row). The not-before direction and exact attainment remain
producer-tier receipts, and completion keeps its graded split
(ranking geolog-checked, enabledness glue argued), exactly as stated with the
certificates above --- instantiation at large $N$ upgrades nothing
and degrades nothing.

\paragraph{The contrast columns, derived not simulated.} Both
bounds are conservative combinatorics on the net. \emph{Markings:}
lots distribute over pending, the six queues, and shipped; with
machines idle every such distribution is reachable (advance one lot
at a time) and distinct, giving $\binom{N+7}{7}$ reachable markings
before counting any in-process or machine state. \emph{Schedules:}
two lots processing concurrently on distinct stations (one in step
1, one in step 2) admit $\binom{259}{32} \approx 10^{40.9}$
interleavings of their firing chains, all distinct as transition
words; scheduling $\lfloor N/2\rfloor$ such pairs in blocks gives
$\binom{259}{32}^{\lfloor N/2\rfloor}$ distinct maximal runs. Every
one of them is covered by the same $\sim$3\,kB of certificate.

\paragraph{Reproduction.} The reproduction kit rebuilds the independent checker at a pinned
revision and re-checks the sealed containers, with per-relation
claim counts and integrity mutants as gates. Of the Mini-Fab
campaign's $150$ footprint rows, this check evaluates the $36$
tick-transition rows; it neither checks the other $114$ nor
identifies the net from the sealed bytes. The kit also includes
the producer-tier mutant tests, but that leg needs unpublished
dependency revisions and cannot be reproduced from public sources
alone. For the dispatch grids, the small-$N$ optimality receipts
(including the LRAT proofs and the independent SMT check), and
the lot-count sweep, see section~\ref{sec:exp-repro}.

\paragraph{The model boundary, honestly.}\label{par:minifab-model-boundary} Everything in this
section --- counting time and carved clock alike --- certifies
the lot-flow core, as a simplified discrete net, and each exclusion
is a named refinement rather than fine print. The first is
parallelism: under the wall-clock semantics each transition fires at
most once per tick, so at a given process step at most one lot can
start, receive a unit of service, or finish in any one tick; a
station's two machines can work in the same tick on lots at
different steps, never on two lots at the same step. Kempf's
specification allows more than one lot at a process step to start,
receive service, or finish in a tick; that parallelism \emph{adds}
runs, so the a-fortiori argument below does not cover it: whether
the optimality results of this section carry over to a net with
Kempf's parallelism, or to his full specification, is open. The
others: diffusion's batching-of-three and its mix rules,
product types and test-lot dedication, the declared buffer
capacities, the transporter, the operators and technician, the
maintenance calendars and failures, and the timed weekly release
schedule. The certificates are arranged to survive the refinements
that only restrict runs --- a conservation or safety row over all
schedulers holds \emph{a fortiori} over fewer --- while refinements
that add Petri transitions re-run the same discharge over the
extended roster; and arrivals-as-births are, as stated in
section~\ref{sec:minifab-spec}, a different certificate species
with its own budget vocabulary, not a missing row here. Nothing in
this section claims throughput under sustained release; that claim
lives, for now, one section over --- about the relaxation.

\section{The fluid certificates: the comparison lane}
\label{sec:minifab-certs}

The second certificate family predates the first and certifies a
different object: the \emph{fluid model} of the Mini-Fab's
deterministic core skeleton, under two named dispatch policies, with
every residual freedom resolved adversarially. It is retained in
full --- as the comparison lane --- because it answers exactly the
question the combinatorial family does not yet price: sustained
arrivals, at stated loads, policy by policy. The contrast, stated
once: a fluid certificate asserts stability of the
\emph{relaxation} --- continuous levels, aggregate drain rates,
indexed by a dispatch policy --- at releases within a tenth of a lot
per week of capacity; a combinatorial certificate
(section~\ref{sec:minifab-exact}) asserts exact statements about the
\emph{discrete system} --- every run, every lot count, scheduler-uniform --- with no continuum limit anywhere. Neither family
subsumes the other today, and the seam between them is measured
rather than smoothed over: the granularity conversion still open on
the fluid side is named at this section's end, and the arrivals
species still open on the combinatorial side is named at
section~\ref{sec:minifab-exact}'s. All numbers in this section
re-derive from the committed artefacts in the experiments repository
--- the certificates and their mutant battery under
\path{semicon/out/minifab}, replayed by the independent checker
through the repository's semiconductor make target --- and are
displayed under the conventions of
section~\ref{sec:exp-conventions}. The campaign summary lives in
section~\ref{sec:exp-semicon}; this section walks the same artefact
at reading depth.

\subsection{The fluid model, with its decisions priced}

Fluid scaling replaces lot counts by continuous levels $q \in
\R_{\ge 0}^6$, one per route step, drained at the stations'
machine-occupancy rates. From the tables of
section~\ref{sec:minifab-line}, with handling folded into occupancy
(the machine is held during load and unload, per-batch at
diffusion):

\begin{center}\small
\begin{tabular}{@{}clrr@{}}
\toprule
step & station (machines) & occupied min/lot & per-machine rate (lots/min)\\
\midrule
1 & diffusion (2)    & $285/3 = 95$  & $1/95$\\
2 & implantation (2) & $60$          & $1/60$\\
3 & lithography (1)  & $75$          & $1/75$\\
4 & implantation (2) & $80$          & $1/80$\\
5 & diffusion (2)    & $315/3 = 105$ & $1/105$\\
6 & lithography (1)  & $30$          & $1/30$\\
\bottomrule
\end{tabular}
\end{center}

\noindent The declared-model ledger is explicit, in the spirit of the
cookbook's declared-prior pattern
(chapter~\ref{ch:mtk-cookbook}): operators are folded into machine
occupancy rather than modelled as a contended resource (that layer is
the pool seam of section~\ref{sec:minifab-composite}); products share
a single class per step (the mix enters only through total release;
dedication and batch-mix rules are granularity effects that vanish in
fluid scaling); batching becomes a rate ($3$ lots per $285$ or $315$
machine-minutes); arrivals are deterministic at exact rate $\lambda$;
and failures and maintenance are excluded from the core skeleton ---
their sound entry point is capacity intervals from availability
moments, named successor work, not fine print smuggled into these
certificates. Setups get two sound variants: a \texttt{core} variant
with no setups (matching the capacity arithmetic), and a
\texttt{setup-adversarial} variant in which every lithography job may
pay up to the full worst-case 12-minute setup, encoded as service-rate
\emph{intervals} $\mu_3 \in [1/87, 1/75]$ and $\mu_6 \in [1/42,
1/30]$ resolved adversarially at every instant --- so the certificate
quantifies over every setup sequencing with no monotonicity
assumption anywhere (queueing networks are not monotone in service
rates; intervals sidestep the issue rather than arguing it).
Worst-case-inline setups put lithography's per-lot occupancy at $129$
minutes, an adversarial ceiling of $10080/129 \approx 78.1$
lots/week.

The two named policies are the classical static priorities for
re-entrant lines: \emph{last-buffer-first-served} (LBFS --- later
steps preempt earlier ones for their station's capacity: step 5 over
step 1, 4 over 2, 6 over 3) and \emph{first-buffer-first-served}
(FBFS, the reverse). These are the disciplines the Kumar--Seidman
destabilization literature is about, which is what gives ``stable at
$\rho$'' its falsifiable content. FIFO is not directly certified --- its
fluid content is class-level dispatch, per
section~\ref{sec:minifab-rows} --- but the setup-adversarial rows
soundly cover FIFO's setup burden for \emph{any} within-class order
at loads up to 78 lots/week. Everything the policy does not pin ---
allocation slack on empty classes, idling below the top positive
priority, rate realization inside the declared intervals --- is
inside the adversary.

\subsection{The certified statement and its checker}

Each certificate exhibits a linear Lyapunov function $V(q) = \langle
w, q \rangle$ with $w \in \Q_{>0}^6$ and an exact rational
$\varepsilon > 0$ such that along \emph{every} admissible fluid
trajectory under the named policy, $\frac{d}{dt} V(q(t)) \le
-\varepsilon$ for almost every $t$ with $q(t) \ne 0$. Consequently
any work-in-progress profile $q_0$ drains to the empty line within
$V(q_0)/\varepsilon$ minutes, stays drained, and the line sustains
its release --- the stability and throughput atoms of
section~\ref{sec:minifab-spec} at once, with a quantitative drain
rate as the bonus.

The checker (\texttt{fluidcheck}, a small Rust program) trusts nothing about how the certificate
was found. It rebuilds every polytope itself from the declared
network data --- routes, rates, capacities, policy --- and replays
one exact-rational Farkas dual per \emph{sign mode}: for each of the
$2^6 - 1 = 63$ nonempty patterns of which queues are positive, the
admissible drift vectors form a polytope, and the certificate's dual
for that mode is a nonnegative combination of the polytope's
constraints witnessing $\langle w, \dot q \rangle \le -\varepsilon$
over all of it. Soundness rests on weak linear-programming duality
and one measure-theory lemma --- the next subsection --- and on
nothing else; in particular, no solver, no floating-point arithmetic,
and no completeness claim about the producer participates in the
trust story. Each certificate carries $819$ obligations and
re-verifies in $1.5$--$3.0$ milliseconds on one CPU core. The
producer that found the certificates (floating-point proposal,
exact-rational certification, self-replay before emission) took
about $0.2$ seconds per certificate --- and is, per the conventions,
freely sloppy: being wrong costs retries, never soundness.

The bespoke container is, since this campaign first shipped, no
longer the only committed form. Ten of the twelve certificates below
are re-emitted as sealed geolog-commit containers on the unified
certificate schema, in two committed lanes --- one against the
reproduction vocabulary, one pinned to the baked theory pair ---
each lane carrying its own mutant battery and re-verified, sign mode
by sign mode, through the engine's exact-rational ground seam;
section~\ref{sec:exp-repro} records the transition's terms, and the
two tree-cover certificates await their species' re-emission. The
original containers remain the record of the campaign that produced
this section's numbers, and the checker above still replays them ---
a committed lane never dies.

\subsection{The boundary-equality lemma, the hinge}

One semantic decision carries the entire result. What does the fluid
model say about an \emph{empty} buffer? The lazy encoding says only
$\dot q_k \ge 0$ when $q_k = 0$ (levels stay nonnegative). The right
encoding says more:

\begin{lemma}[boundary equality]
\label{lem:minifab-boundary}
Let $q_k : [0,\infty) \to \R_{\ge 0}$ be absolutely continuous. Then
at almost every $t$ with $q_k(t) = 0$, the derivative exists and
$\dot q_k(t) = 0$.
\end{lemma}

\noindent (At almost every point of the level set $\{q_k = 0\}$ ---
discard the null set where $q_k$ is not differentiable --- the point
is a Lebesgue density point of the set, hence an accumulation point
of other zeros, and a derivative along a sequence of zeros is zero.)

The measured difference between the two encodings is not a
refinement at the margin. Under the inequality-only semantics, the
adversary may pour phantom flow into empty buffers whose station
idles, and \emph{no linear Lyapunov function exists at any tested
load} --- the certificate family is empty. Under the equality, linear
certificates exist to within a tenth of a lot per week of fluid
capacity. The
equality also makes many sign modes vacuous: at the flagship
configuration, $41$ of the $63$ mode polytopes are empty (a pattern
like ``only $q_5$ positive'' forces full diffusion output through an
empty downstream line that lithography cannot pass). Vacuity needs no
case split and no emptiness branch in the checker: an infeasible
polytope has an unbounded dual, so the same Farkas row shape that
certifies a drift bound on a nonempty mode certifies \emph{any} bound
on an empty one. The checker has one code path, on purpose.

\subsection{The ladder}

Twelve certificates and two deliberate failures, all replayed by the
checker from the committed artefacts:

\begin{center}\small
\begin{tabular}{@{}llccrl@{}}
\toprule
certificate & policy & variant & lots/week & $\rho_{\mathrm{lith}}$ &
$\varepsilon$ ($V$-units/min)\\
\midrule
\texttt{lbfs\_core\_84wk} & LBFS & core & 84 & $0.875$ & $1/3600$\\
\texttt{lbfs\_core\_90wk} & LBFS & core & 90 & $0.938$ & $1/7200$\\
\texttt{lbfs\_core\_94wk} & LBFS & core & 94 & $0.979$ & $1/21600$\\
\texttt{lbfs\_core\_95wk} & LBFS & core & 95 & $0.990$ & $1/43200$\\
\texttt{lbfs\_core\_191\_2wk}\,\textsuperscript{t} & LBFS & core &
  95.5 & $0.995$ & $1157/10^{8}$\\
\texttt{lbfs\_core\_959\_10wk}\,\textsuperscript{t} & LBFS & core &
  95.9 & $0.999$ & $231/10^{8}$\\
--- probe & LBFS & core & $96$ = capacity & $1.0$ & none (correctly)\\
\texttt{fbfs\_core\_84wk} & FBFS & core & 84 & $0.875$ &
  $\approx 2.8 \times 10^{-4}$\,\textsuperscript{a}\\
\texttt{fbfs\_core\_95wk} & FBFS & core & 95 & $0.990$ & $1/38880$\\
\texttt{lbfs\_setupadv\_72wk} & LBFS & setup-adv & 72 & $0.921^{*}$ &
  $11/74760$\\
\texttt{lbfs\_setupadv\_76wk} & LBFS & setup-adv & 76 & $0.973^{*}$ &
  $23/448560$\\
\texttt{lbfs\_setupadv\_78wk} & LBFS & setup-adv & 78 & $0.998^{*}$ &
  $1/299040$\\
--- probe & LBFS & setup-adv & $79 >$ ceiling & $>1^{*}$ & none (correctly)\\
\texttt{fbfs\_setupadv\_78wk} & FBFS & setup-adv & 78 & $0.998^{*}$ &
  $1/287280$\\
\bottomrule
\end{tabular}
\end{center}

\noindent {\small ($*$) setup-adversarial utilization is measured
against worst-case-inline lithography occupancy, $129$ min/lot.
(a)~exact value $6923209477/24747380706600$; the display is rounded
in the claim-weakening direction. (t)~emitted as fully-refined
sign-split \emph{tree covers} --- the certificate format's refinement
species, adjudicated by the checker's tree walk rather than the mode-mask
walk --- with the weight search and all cell duals solved as one
sparse linear program.} \medskip

$V$-units are normalized so $\sum_k w_k = 1$: $V$ is a weighted lot
count. At the flagship configuration the optimizer's own weights came
out as $w = (7, 7, 7, 3.5, 3.5, 2)/30$ --- remaining-lithography-work
flavoured, found by the linear program, not designed. A worked
reading: a ten-lot backlog at step 1, at the nominal 84 lots/week
release, has $V = 7/3$ and is certified to clear within $8400$
minutes --- under six days --- against every adversarial resolution
of the residual freedoms.

The table carries three findings. First, the certified load meets
capacity: certificates exist at $95.9$ of $96$ lots/week ($99.90\%$
of fluid capacity) and at $78$ of the $78.1$-lot adversarial ceiling
($99.8\%$); producer-side bisection brackets the last certifiable
release within $6 \times 10^{-5}$ lots/week of the fluid capacity
(the one-lot gap in the original release grid was grid granularity,
not lost capacity); and the probes at $96$ and $79$ fail
\emph{because the load is genuinely at or over capacity} --- so the
demonic generosity (everything unpinned handed to the adversary) is
measurably free here, costing less than a tenth of a lot per week of
certified capacity. Second, the exact drain rates
separate the policies: FBFS edges LBFS near capacity ($1/38880$
versus $1/43200$ at 95 lots/week; $1/287280$ versus $1/299040$ at
the adversarial 78) --- both stable everywhere tested, consistent
with the classical fluid-limit picture for re-entrant lines, and the
kind of exact policy comparison the claim class was chosen to
produce. Third, the certificates and the reference simulation
triangulate: certified setup-adversarial stability up to $78$
lots/week $<$ simulated FIFO stability at $84$ $<$ the simulated
deterministic-core ceiling at $87$ $<$ certified zero-setup capacity
at $96$. The gap between $78$ and $87$ is the certificate's
conservatism against worst-case setups ($24$ minutes of setup per
lot charged by the adversary, against $9.4$ measured for FIFO in
simulation) --- and closing most of it is a policy question
(run-length policies amortize setups), not a checker limitation.

The mutant battery, per the standing convention
(section~\ref{sec:exp-conventions}): seven deliberately corrupted
certificates --- inflated $\varepsilon$, zeroed weight vector,
negated dual, broken policy identity, dropped mode, swapped modes,
over-capacity release --- each designed to be refused at a specific
check, and each in fact refused.

\paragraph{What this section does not claim.} These are claims about
the declared fluid model of the deterministic core, not about the
discrete line --- whose own certificates are
section~\ref{sec:minifab-exact}'s, and are of a deliberately
different species: exact run-length budgets and invariants, not rate
stability under arrivals. Of the granularity work that converts a
fluid drain rate into discrete-line guarantees for the count chain,
the Little's-law (exact-moments) half has since been discharged: a
discrete-lot certificate species replays lot-level cycles
deterministically, certifying rate stability at the fluid capacity
itself --- a statement no fluid certificate can make --- together
with exact steady work-in-progress and cycle-time moments from
declared initial families, with its own designed-locus mutant
battery. The Foster--Lyapunov half --- universal coverage over all
mixed work-in-progress states --- remains open, as does the sound
entry of failures and maintenance via availability intervals, and
the personnel, transport, and batching-granularity layers live on
the composite seams of section~\ref{sec:minifab-composite} until
those patterns carry certificates of their own. Nothing here claims either policy is
\emph{optimal} --- only that its stability region is certified, to
within a tenth of a lot per week of the model's capacity, by an
artefact a reader can re-refuse and re-accept on their own machine in
milliseconds.

\section{The construction in TypeScript}
\label{sec:minifab-ts}

One orientation before the listing: the construction below is the
fluid lane's (section~\ref{sec:minifab-certs}) --- the model it
assembles is the aggregate reading, and the certificates it drives
are the release-grid drain certificates of that section's ladder.
The combinatorial campaign of section~\ref{sec:minifab-exact} is
produced today on the toolkit's Rust strategy surface; its
TypeScript twin awaits the production-strategy surface on the
toolkit's TypeScript side (the specification surface's revision has
since landed, and this section's listings ride it), and this
section will grow that second listing when the surface lands,
rather than simulate one now.

This section was staged until the toolkit's TypeScript construction
surface shipped, with its acceptance test fixed in advance: a single
executable listing entering the ground tables of
section~\ref{sec:minifab-rows} as rows, declaring the mechanisms as
manifests, assembling the composite, forming specification content
against the aggregate interface, handing the obligation set to the
session workspace of chapter~\ref{ch:qtsl-tutorial}, and --- driven
over the same release grid --- producing certificates the
independent checker accepts at every load certified in
section~\ref{sec:minifab-certs} and none at the two probe loads that
correctly failed there. The surface shipped; the listing runs; the
test passed, on the terms recorded at the end of this section. The
complete file is committed as \path{qtsl/book/listings/minifab.ts}
(with its emission seam \path{fluidseam.ts} beside it), against the
\texttt{@maitria/mtk} and \texttt{@maitria/mtk-session} packages of
section~\ref{sec:qs-typescript} and
section~\ref{sec:tut-session}; the excerpts below are verbatim
slices of it, and the transcripts are the run's own output.

The ground tables enter exactly as
section~\ref{sec:minifab-rows} declared them --- stations, route
steps, release generators, the lithography setup rules, the
implantation failure brackets --- as data the rest of the file
compiles from, so they are not reprinted here. The mechanisms are:
each station at the aggregate stratum (per-visit queue levels, a
busy-machine count, and the one device the platform renders every
duration with --- a timer drifting at rate one), with the declared
buffer and machine capacities carried as domain constraints:

\begin{Verbatim}[fontsize=\small,frame=leftline,rulecolor=\color{accentmid},framesep=3mm,xleftmargin=3mm]
function station(st: number): system.System {
  const visits = ROUTE_STEP.filter((r) => r.st === st); // two per station
  const n = visits.length;
  const width = n + 2; // queues, busy, timer
  return system.buildSystem({
    name: STATION[st].name,
    state: [
      ...visits.map((r, v) => ({ id: v, name: `q${r.idx}`, sort: "Nat" })),
      { id: n, name: "busy", sort: "Nat" },
      { id: n + 1, name: "tau" },
    ],
    interface: {
      inputs: visits.map((r, v) => ({ id: v, name: `arrive${r.idx}` })),
      outputs: visits.map((r, v) => ({ id: v, name: `complete${r.idx}` })),
    },
    modes: [
      {
        id: 0,
        name: "serving",
        domain: [
          ...visits.map((_, v) => ({ poly: mono(width, v), op: "ge", rhs: 0 })),
          ...visits.map((_, v) => ({
            poly: mono(width, v),
            op: "le",
            rhs: STATION[st].bufcap,
          })),
          { poly: mono(width, n), op: "ge", rhs: 0 },
          { poly: mono(width, n), op: "le", rhs: STATION[st].machines },
        ],
      },
    ],
    drift: [{ mode: 0, var: n + 1, poly: constant(width, "1") }], // tau' = 1
    init: { kind: "point", coords: Array.from({ length: width }, () => 0) },
  });
}
\end{Verbatim}

\noindent The release mechanism is three timers (the thesis's
signal-generator idiom is the timer-mechanism idiom, as
section~\ref{sec:minifab-rows} observed), the sink counts finished
lots and pools their flow times, and each conserved unit ---
operators, technician, transporter --- is a unit-capacity pool
mechanism whose conservation range is likewise a domain constraint.
None of these manifests matches a catalogued profile, and none needs
to: the free constructor never refuses, and the preflight's
\texttt{noKnownProfile} notes in the transcript below are
classification, not complaint --- a mechanism whose populated family
set matches no preset is a first-class citizen of
chapter~\ref{ch:qtsl-mathematical-models}.

The composite wires the re-entrant lot flow, completion wire into
arrival wire along the route:

\begin{Verbatim}[fontsize=\small,frame=leftline,rulecolor=\color{accentmid},framesep=3mm,xleftmargin=3mm]
const fab = system.composite({
  name: "minifab",
  components: [release, diff, impl, lith, sink, ...pools],
  wiring: [
    flow([0, 0], [1, 0]), // release      -> diffusion    (step 1)
    flow([1, 0], [2, 0]), // diffusion    -> implantation (step 2)
    flow([2, 0], [3, 0]), // implantation -> lithography  (step 3)
    flow([3, 0], [2, 1]), // lithography  -> implantation (step 4)
    flow([2, 1], [1, 1]), // implantation -> diffusion    (step 5)
    flow([1, 1], [3, 1]), // diffusion    -> lithography  (step 6)
    flow([3, 1], [4, 0]), // lithography  -> finished lots
  ],
});
\end{Verbatim}

\noindent One boundary deserves stating here, where the code makes it
concrete rather than abstract: the directed lot-flow seams of
section~\ref{sec:minifab-composite} are constructible today, as the
wiring rows above; the \emph{pool-membership} and
\emph{event-fusion} seams are
the event-port stratum, whose fusion primitive now carries a
producer surface (typed fusion rows over the stations' population
edges; section~\ref{sec:cp-eventport}) while the pool-membership
packaging is still design-tense, which is why the pools stand here
as constructed mechanisms whose seize and release bindings await
the packaged rows. And the dispatch freedom enters precisely where
section~\ref{sec:minifab-certs} put it: as the policy index of each
certificate, carried below by the release-grid driver, not as a row
on this composite.

The specification content formable in the released vocabulary is the
certificates' own quantitative consequence --- the worked reading of
section~\ref{sec:minifab-certs}, as a typed record over the
six-level aggregate interface:

\begin{Verbatim}[fontsize=\small,frame=leftline,rulecolor=\color{accentmid},framesep=3mm,xleftmargin=3mm]
const drainSpec = spec.build({
  name: "minifab drain: LBFS core 84/wk, ten lots at step 1",
  forall: [spec.probGe(1, spec.deadline(8400, spec.polyLe(0, wip)))],
  iface: {
    invars: [],
    outvars: ROUTE_STEP.map((r) => ({ name: `q${r.idx}`, shape: [] })),
  },
});
\end{Verbatim}

\noindent with \texttt{wip} the sum of the six coordinate
polynomials. The long-run family of
section~\ref{sec:minifab-spec} enters the same file at its formable
half: the stability atom, through the average former. For the fluid
skeleton under this policy at this release the honest cap is zero
--- the line drains and stays drained, so the transient washes out
of the time average:

\begin{Verbatim}[fontsize=\small,frame=leftline,rulecolor=\color{accentmid},framesep=3mm,xleftmargin=3mm]
const longRunStable = spec.build({
  name: "minifab long-run: LBFS core 84/wk, average wip",
  forall: [spec.returnLe(0, spec.poAvg(spec.oPoly(wip)))],
  iface: {
    invars: [],
    outvars: ROUTE_STEP.map((r) => ({ name: `q${r.idx}`, shape: [] })),
  },
});
\end{Verbatim}

\noindent The record preflights, builds, encodes, and round-trips
on the released surface (the receipt below carries the run); its
resting state is the one section~\ref{sec:minifab-spec} assigns it:
a formed claim whose certificate species is named and not yet
claimed, so in the workspace's books the obligation would stand
open --- which is where a claim without its battery belongs. The
throughput atom stays out of this listing for a narrower reason
than it once had: its former is staged in the theory
(section~\ref{sec:spec-lang-rate}), but the producer surface
quoted here does not construct it yet, and this listing does not
simulate constructions the toolkit does not have.

The producer core is the Cookbook's propose-then-dispose discipline
(chapter~\ref{ch:mtk-cookbook}) on the checker's own vocabulary: the
listing rebuilds the same constraint rows the checker rebuilds, asks
one floating-point robust linear program for candidate weights
(purely a proposal), then certifies one exact dual per sign mode and
replays every obligation in exact arithmetic before emitting a byte.
The equality rows are where two of this chapter's decisions become
five lines of algebra --- the boundary-equality lemma
(Lemma~\ref{lem:minifab-boundary}) as flow conservation on each
empty buffer, and the policy's entire fluid content as saturation
above each nonempty priority level:

\begin{Verbatim}[fontsize=\small,frame=leftline,rulecolor=\color{accentmid},framesep=3mm,xleftmargin=3mm]
function eqRows(net: Net, mask: number): Row[] {
  const k = net.k;
  const rows: Row[] = [];
  for (let i = 0; i < k; i++) {
    if (!((mask >> i) & 1)) {
      // buffer i empty: inflow = outflow (a.e. zero derivative)
      const r = zeros(2 * k);
      r[i] = Q.from(-1);
      net.route.forEach((succ, j) => {
        if (succ === i + 1) r[j] = r[j].add(Q.from(1));
      });
      rows.push({ r, rhs: net.arrivals[i].neg() });
    }
  }
  net.prio.forEach((order, s) => {
    for (let pos = 0; pos < order.length; pos++) {
      if ((mask >> (order[pos] - 1)) & 1) {
        // a nonempty buffer at priority: everything at or above it
        // saturates the station
        const r = zeros(2 * k);
        for (const j of order.slice(0, pos + 1)) r[k + j - 1] = Q.from(1);
        rows.push({ r, rhs: net.caps[s] });
      }
    }
  });
  return rows;
}
\end{Verbatim}

Finally the obligation set, handed to the session workspace: one
obligation per grid point, probes included, under a statement kind
whose single move runs the producer, emits through the canonical
encoder, and lets the checker adjudicate. The kind's failure mode is
part of its design: at an infeasible load it returns
\emph{inconclusive} --- ``the load is at or over the variant's
capacity; absence of a certificate is not a refutation'' --- because
no certificate is not a disproof, and the workspace's books are only
as good as the outcomes recorded in them.

\begin{Verbatim}[fontsize=\small,frame=leftline,rulecolor=\color{accentmid},framesep=3mm,xleftmargin=3mm]
const GRID: [string, string, number][] = [
  ["LBFS", "core", 84],
  ["LBFS", "core", 90],
  ["LBFS", "core", 94],
  ["LBFS", "core", 95],
  ["LBFS", "core", 96], // probe: fluid capacity
  ["FBFS", "core", 84],
  ["FBFS", "core", 95],
  ["LBFS", "setup-adversarial", 72],
  ["LBFS", "setup-adversarial", 76],
  ["LBFS", "setup-adversarial", 78],
  ["LBFS", "setup-adversarial", 79], // probe: over the adversarial ceiling
  ["FBFS", "setup-adversarial", 78],
];

const ws = new Workspace([fluidStable]);
for (const [policy, variant, lotsPerWeek] of GRID) {
  const ob = ws.pose({ kind: "minifab.fluidStable", policy, variant, lotsPerWeek });
  const r = ws.apply(ob.id, "farkas");
  const w = r.attempt.witness as { eps?: string; checker?: string } | undefined;
  console.log(
    `${ob.id} ${ws.describe(ob.id)}\n  -> ${r.attempt.outcome}` +
      (w ? ` eps=${w.eps} (${w.checker})` : ` (${r.attempt.note})`),
  );
}
\end{Verbatim}

The run, whole --- construction, specification, and the twelve
obligations --- takes under four seconds:

\begin{Verbatim}[fontsize=\small,frame=leftline,rulecolor=\color{accentmid},framesep=3mm,xleftmargin=3mm]
$ npx tsx minifab.ts
composite: minifab (composite), schema maitria.system.manifest.v0;
    9 preflight notes, kinds: noKnownProfile
spec: built ("minifab drain: LBFS core 84/wk, ten lots at step 1")
ob1 fluid skeleton drains under LBFS (core) at 84 lots/week
  -> discharged eps=1/3600 (ACCEPT modes=63 obligations=819 eps=1/3600)
ob2 fluid skeleton drains under LBFS (core) at 90 lots/week
  -> discharged eps=1/7200 (ACCEPT modes=63 obligations=819 eps=1/7200)
ob3 fluid skeleton drains under LBFS (core) at 94 lots/week
  -> discharged eps=1/21600 (ACCEPT modes=63 obligations=819 eps=1/21600)
ob4 fluid skeleton drains under LBFS (core) at 95 lots/week
  -> discharged eps=1/43200 (ACCEPT modes=63 obligations=819 eps=1/43200)
ob5 fluid skeleton drains under LBFS (core) at 96 lots/week
  -> inconclusive (the load is at or over the variant's capacity;
     absence of a certificate is not a refutation)
ob6 fluid skeleton drains under FBFS (core) at 84 lots/week
  -> discharged eps=6923209477/24747380706600 (ACCEPT modes=63
     obligations=819 eps=6923209477/24747380706600)
ob7 fluid skeleton drains under FBFS (core) at 95 lots/week
  -> discharged eps=1/38880 (ACCEPT modes=63 obligations=819 eps=1/38880)
ob8 fluid skeleton drains under LBFS (setup-adversarial) at 72 lots/week
  -> discharged eps=11/74760 (ACCEPT modes=63 obligations=819 eps=11/74760)
ob9 fluid skeleton drains under LBFS (setup-adversarial) at 76 lots/week
  -> discharged eps=23/448560 (ACCEPT modes=63 obligations=819 eps=23/448560)
ob10 fluid skeleton drains under LBFS (setup-adversarial) at 78 lots/week
  -> discharged eps=1/299040 (ACCEPT modes=63 obligations=819 eps=1/299040)
ob11 fluid skeleton drains under LBFS (setup-adversarial) at 79 lots/week
  -> inconclusive (the load is at or over the variant's capacity;
     absence of a certificate is not a refutation)
ob12 fluid skeleton drains under FBFS (setup-adversarial) at 78 lots/week
  -> discharged eps=1/287280 (ACCEPT modes=63 obligations=819 eps=1/287280)

12 obligations: 10 discharged, 2 open (96wk, 79wk)
flagship V weights: (7/30, 7/30, 7/30, 7/60, 7/60, 1/15)
\end{Verbatim}

\noindent And the acceptance test's other half, against the
standalone checker binary over the emitted files:

\begin{Verbatim}[fontsize=\small,frame=leftline,rulecolor=\color{accentmid},framesep=3mm,xleftmargin=3mm]
$ for f in out/*.cbor; do fluidcheck $f > /dev/null && pass=$((pass+1)); done
$ echo "standalone fluidcheck: $pass/10 ACCEPT"
standalone fluidcheck: 10/10 ACCEPT
\end{Verbatim}

Read the transcript against the ladder of
section~\ref{sec:minifab-certs}. The acceptance criterion was
checker-acceptance, deliberately --- a reimplemented producer may
find different valid certificates --- and what happened is a report,
not a promise: this producer landed the \emph{same} certificates.
Every drain rate matches the table's, down to the FBFS
exact value whose display there is rounded; both probes fail for the
right reason; and the weight vector $w = (7, 7, 7, 3.5, 3.5, 2)/30$
--- which that section reported as found by the linear program, not
designed --- has now been found twice, by two implementations
proposing over the same polytopes. The independence boundary is
drawn where it always is in this book: the \emph{producer} above is
a fresh implementation on the TypeScript surface, while the emission
seam links the checker crate itself, so that the bytes written and
the bytes judged pass through one canonical encoder (at this preview
the seam is a thin binding around the checker crate the reproduction
spine imports; its migration is tracked in section~\ref{sec:exp-repro}) --- and the standalone
binary replays every certificate from the bytes alone, trusting
nothing about how they were made.

 \chapter{Implementation Pearls}
\label{ch:mtk-pearls}

\noindent\draftmark\medskip

This chapter is written in the tradition of Bentley's \emph{Programming
Pearls} \cite{bentley2000pearls}: short, self-contained engineering essays, one per design
decision that earned its place in the producer stack. The genre has a
discipline of its own --- a pearl is not advice; it is a technique with
the measurement that decided it --- and this chapter keeps to it. Each
essay states the problem, the technique, and the receipt: numbers from
committed artefacts, reproducible from the public repositories where
the artefact has migrated, and identified as workshop-record
measurements where it has not yet (per the preview note at the front of
this volume). The common thread is that \emph{certified-producer
engineering} --- code whose every output faces an unforgiving checker
--- quietly inverts several classic engineering trade-offs, and each
pearl records one such inversion. Where the agents' cookbook
(\bookref{ch:mtk-cookbook}{the agents' cookbook}) states the moves of a
good untrusted producer, this chapter states the decisions that hold up
the machinery those moves run on.

\section{The Certifying Tier Is a GPU Workload Too}
\label{sec:pearl-exact-gpu}

\paragraph{The problem.} The producer stack's central split ---
a fast float tier proposes, an exact tier certifies
(\bookref{ch:mtk-cookbook}{the cookbook}, \emph{Sloppy Proposer, Exact
Disposer}) --- is usually read as a hardware split: floats on the
device, exact rationals on the host. The economics that force floats
out of the inner search loop seem, by the same logic, to force
exactness off the device. In the workshop prototype that seeded the
toolkit's kernels, one per-leaf basis-change GEMV at a realistic shape
cost $715\,\mu\mathrm{s}$ in exact-$\Q$ arithmetic on the CPU, against
$4.65\,\mu\mathrm{s}$ in directed \texttt{f64} on the same CPU and
$191$\,ns on the device --- three and a half orders of magnitude
between the exact path and the device path. Leave exactness host-side
and serial, and the moment search becomes device-fast, certificate
assembly becomes the pipeline's ceiling.

\paragraph{The technique.} Put the exact tier on the device as an
\emph{integer} workload. The certifying kernels in \path{crates/mtk}
represent a batch of rational tables as sign-magnitude integer
numerators in base-$2^{64}$ limb planes over \emph{one} positive shared
denominator per cell (the schoolbook common-denominator move,
industrialized), so that every transform is an integer recurrence
with an explicit denominator-bookkeeping rule: no division, no gcd, no
floating point anywhere in the certifying path. Reduction to canonical
form happens host-side at the boundary. Fixed-capacity limb arithmetic
is sound only if capacity exhaustion is \emph{detected}, never wrapped,
so every product writes into scratch of compile-time capacity with
stated headroom lemmas, every store back is fit-checked, and the
carry-out flag is captured and treated as overflow anyway --- the lemma
and the flag would both have to fail. On overflow the kernel poisons
the cell: outputs zeroed, routing flag set, the cell's work abandoned.
A poisoned cell's outputs are never read; the promotion ladder
re-derives it from its \emph{inputs} at the next tier --- a wider
device width, up to sixteen limbs, then unbounded host integers. Being
wrong about a width costs a retry one rung up, never a wrong table.
The kernels themselves ship as embedded CUDA source, compiled at
runtime for whatever device is present, which makes the development
loop one \texttt{cargo run} on any machine; an ahead-of-time path can join
later behind the same module seam (and when it does, the
pipeline-of-record law of section~\ref{sec:pearl-lane-law} applies to
it).

\paragraph{The receipt.} The exact-tier selftest runs eleven cases
spanning the census shapes, fused transform chains, a
$\sim$130-digit-coefficient regime, and engineered overflow-boundary
fixtures, demanding bit-exact equality against the host big-integer
reference for every reduced table entry, every verdict, and every
first-violation address --- with the promotion ladder's rungs each
demonstrably exercised. Replaying a committed certificate from its
bytes reproduces its certified threshold --- a 55-digit reduced
rational --- bit-for-bit across three substrates: Python
\texttt{fractions}, Rust \texttt{num-bigint}, and CUDA \texttt{u64}
limb planes. The full soundness burden, including what ``bit-exact''
means here and the residual risks, is stated in
\path{crates/mtk/BERNX-SOUNDNESS.md} in the
\path{maitria-coalition/mtk} repository.

\section{Round Outward --- to the Verifier's Resolution}
\label{sec:pearl-outward}

\paragraph{The problem.} The base law of float participation in a
certified system is directed rounding: on any verdict-relevant path,
every operation rounds toward the safe side --- toward $-\infty$ for
lower bounds, toward $+\infty$ for upper bounds --- so enclosures stay
enclosures. The law is necessary, and it is not sufficient. A
certifier campaign in the workshop lineage banked the counterexample:
its checker decides transcendental sign questions by bisection on a
dyadic lattice, deepening to a resolution of $2^{-256}$, and a
producer emission relying on a $10^{-70}$ \emph{relative} guard alone
can land a correctly-rounded margin \emph{below one lattice unit}.
Every rounding outward, every claim true in $\mathbb{R}$ --- and
undecidable at any depth the checker will ever spend: a correct
refusal of a true claim, forever.

\paragraph{The technique.} Treat the emission grid as part of the
interface. The producer quantizes every emitted bound outward onto a
dyadic grid coarse enough that each claimed margin clears the
verifier's decision lattice by whole grid steps --- in the campaign
above, one full $2^{-96}$ step outward keeps every margin decidable at
bisection depth 16--64. The doctrine gains a clause: \emph{exactness
lives in the direction of rounding, and the magnitude of rounding must
respect the verifier's resolution.} The cost is $2^{-96}$ of
tightness, unmeasurable at certificate scale; the benefit is that a
refusal now carries information --- it indicts the claim, not the
resolution. Two smaller dividends follow. Dyadic endpoints are
exact rationals --- an \texttt{f64} \emph{is} a dyadic rational --- so
the checking side compares them in exact arithmetic, and a certificate
that confines itself to the dyadic grid never smuggles a rounding into
its own statement. And grid-aligned emissions keep denominators
uniform across a certificate's rows, which the exact tier's
shared-denominator representation
(section~\ref{sec:pearl-exact-gpu}) rewards directly.

\paragraph{The receipt.} Workshop record: the continuous-latents
certifier behind the counterexample above landed with thirteen
artefacts accepted under truth containment, thirteen mutants refused
at their designed loci, and a twenty-six-case property layer --- on
the grid-clearing emission discipline. The lattice constants above
($2^{-256}$ checker resolution, $2^{-96}$ emission step, decision
depth 16--64) are that campaign's numbers; its artefacts migrate to
the public repositories with their campaign, per the preview note.

\section{A Lane May Change Cost, Never Verdicts}
\label{sec:pearl-lane-law}

\paragraph{The problem.} Multi-architecture fast paths rot in a
characteristic way: not silent breakage but \emph{silent divergence},
an optimized lane that is almost always equal to the reference.
Hand-vectorized code that must be faster by assumption --- and equal
by assumption --- is how fast-path crates decay. In a system whose
outputs face a checker, a divergent producer-side lane wastes
campaigns; a divergent screening lane on the checking side is a
soundness incident in waiting.

\paragraph{The technique.} The shared-kernels repository
(\path{maitria-coalition/kernels}) binds every kernel to one law:
\emph{an acceleration lane may change cost, never verdicts.} Three
structural consequences. \emph{Reference first}: every kernel ships a
portable scalar reference, written for obviousness, always compiled
--- the reference is the semantics; lanes are the economics.
\emph{Receipts before dispatch}: a property-based differential battery
gates every lane in CI, but passing the battery earns a lane only
existence; a lane becomes its architecture's default only with a
measured receipt on real hardware, committed to the repository --- and
a lane that loses to the reference on its own target stays available,
un-dispatched, and marked. \emph{Fit-detection, never approximation}:
fixed-width paths detect inputs they cannot represent exactly and
refuse them to the caller's promotion ladder; no rounding exists on
any verdict-relevant path. For JIT-compiled device lanes there is a
fourth: the same kernel source can compile branch-free under an
offline pipeline and branchy --- an order of magnitude slower ---
under a runtime pipeline, so every device lane names its \emph{pipeline
of record}, acceptance benches run that pipeline, and the disassembly
branch count is printed at build time and recorded in the receipt as a
watched number.

\paragraph{The receipt.} The founding receipts of the
coefficient-sweep family, committed under \path{receipts/} in the
kernels repository: the NEON lane at $1.47$--$3.00\times$ the
reference across shapes on a Grace-class aarch64 workstation, the AVX2
lane at $1.25$--$3.96\times$ on a Skylake-W Xeon --- with the
receipt's own candid reading attached, that the wins are modest at the
largest sizes because a streaming \texttt{i64} scan leaves cache and
becomes memory-bandwidth-bound. The dispatch code cites those numbers,
not the intuition they corrected. And the batteries earn their keep in
both directions: the GPU interval family's \emph{first} battery run
falsified the oracle, not the kernels --- a naive exact-rational
reference had missed IEEE~754's signed-zero rules \cite{ieee754} (an
exactly-cancelling round-toward-$-\infty$ addition yields $-0.0$; an
exact-zero product carries the XOR of the operand signs). The device
was right; the oracle now spells the exact-zero cases out, and the
repaired family ran sixteen-for-sixteen green at a hundred property
examples each --- a battery its committed receipts have since grown
to thirty-three properties, still all green at the same hundred. Exact rationals carry no signed zero, so an exactness oracle for
float behaviour must state IEEE's sign semantics explicitly --- a
differential battery is symmetric, and the symmetry is a feature.

\section{The Babysitter Thread}
\label{sec:pearl-babysitter}

\paragraph{The problem.} An exact producer grinding a large obligation
frontier can be legitimately silent for minutes: no intermediate
result is trustworthy enough to print, and the hot loop has better
things to do than I/O. But at an operator's terminal ---
increasingly, an operating \emph{agent's} terminal --- ``no output for
minutes'' is indistinguishable from ``hung,'' and the rational
response to a hung-looking grinder is to kill it, wasting exactly the
work the silence was hiding. The doubt is the cost: it can be resolved
only by waiting longer or killing sooner, and both resolutions are
expensive. Classic optimization-solver CLIs solved this decades ago:
an iteration line, about once a second.

\paragraph{The technique.} The toolkit's convention: any hot loop
expected to run beyond about a second reports progress to standard
output at least once per second --- iteration count plus whatever else
is cheaply at hand (objective, residual, cells retired). The
implementation keeps the reporting out of the loop's way: the hot loop
maintains \emph{only counters}, written with relaxed atomic stores ---
no I/O, no clock reads, no timing branches on the hot path --- while a
\emph{babysitter thread} wakes on a system timer about once a second,
reads the counters, and prints. Where the loop body is already chunky,
the degenerate inline form suffices: consult the clock only every
$N$th iteration and print when a second has elapsed --- the event-loop
engine in \path{crates/mtk-eventloop} does exactly this in its replay
path (every 64th application, one elapsed check) and rate-limits its
per-pass progress lines to the same once-per-second cadence, with the
final line always printed. Where a verb's
standard output is a machine-consumed data channel, the collision is
flagged and the progress line moves --- never silently dropped.

\paragraph{The receipt.} The convention entered the toolkit the usual
way: a long exact-tier run drew a ``no sign of progress'' report, and
settling that doubt cost more than the instrumentation ever would
have. The fix is as close to free as instrumentation gets --- a
relaxed store per iteration and one thread sleeping on a timer --- and
both forms are landed machinery in the \path{mtk} repository. The
exact solver loops carry the full form
(\path{crates/mtk/src/progress.rs}: two relaxed counters, threaded as
an optional argument so that opting out costs nothing, with the
\texttt{lp-solve} verb's exact path spawning the once-per-second
monitor); the
event-loop engine's replay path carries the inline form
(\path{crates/mtk-eventloop/src/engine.rs}). The surface also crosses
the language boundary. The toolkit's TypeScript bindings expose the
same counters as a polling handle over asynchronous solve variants;
JavaScript gives the babysitter and the baby one thread to share, so
the solve must leave it before any timer can read a counter mid-grind.
The printing monitor stays on the consumer's side of the boundary,
where a collision with a machine-consumed output channel is fixable at
the call site rather than at a process boundary. One detail is worth copying: a test in
each ticking module, on both sides of the language boundary, asserts
both that the counters move and that their presence changes no outcome
--- the same run, bit-identical, with and without the observer. The pearl is the
asymmetry: the counters cost nanoseconds per iteration; their absence
is priced in killed runs.

\section{A Certificate Is Database Content, Not a File Format}
\label{sec:pearl-cert-rows}

\paragraph{The problem.} Every bespoke certificate file format breeds
an ecosystem tax: one parser per consumer, version skew between them,
vendored readers drifting from the spec, and an interchange layer that
slowly comes to carry semantics nobody proved anything about.
The toolkit's development lineage paid this tax once --- an early
binary certificate format grew four independent readers before it was
deliberately retired --- and the replacement is a structural decision,
not a better format.

\paragraph{The technique.} Certificates are ordinary rows in the
logical database: columns of exact scalars, tensor-valued content
where the payload is genuinely tensor-shaped, stored under the
certificate theory that the \path{qtsl} repository defines as data
(\path{theories/Qtsl.geolog}, extending the Bernstein Logic proof
tower of \path{theories/BernsteinLogic.geolog}), in the
same content-addressed commits as everything else. There is no
intermediate certificate file format to version, so there is none to
drift; a checker's input is a sealed commit body, and the same bytes
serve every consumer --- the production engine's incremental
monitoring, the certified-checker route, the reflection route into a
proof assistant's kernel --- with no translation step left to audit.
Merge is set union plus re-derivation, and provenance questions
(\emph{who asserted this row, and did anyone assert its
contradiction?}) are structural properties of signed content-addressed
commits rather than per-format afterthoughts. What remains outside the
database is deliberately thin: the untrusted tiers exchange proposals
and refusals as CBOR shapes in a dependency-free types crate
(\path{CONTRACTS.md} \S1.2 in the \path{mtk} repository draws the
boundary: witness interchange is \emph{not a proof format}), so the
producer's conveniences can evolve freely without a byte of them
leaking into the checked surface's semantics.

\paragraph{The receipt.} The interchange layer's golden vectors decode
a committed certificate whose payload carries 55-digit big-integer
fractions and round-trip it losslessly through the v0 shapes
(\path{crates/witness-types/tests/golden.rs} in the \path{mtk} repository)
--- exactness surviving serialization is a test, not a hope. And the
next pearl depends on this one: canonical content is what makes
byte-identity a meaningful instrument in the first place.

\section{Byte-Identity Is a Conformance Instrument}
\label{sec:pearl-byte-identity}

\paragraph{The problem.} Cross-platform reproducibility for numerical
systems is conventionally stated in tolerances --- ``agrees to
$10^{-9}$ across machines'' --- which reintroduces judgment exactly
where an audit wants none: who chose the tolerance, and what drifted
inside it? A certificate-producing system can demand more. If the
certifying path is exact and the emission is canonical, the correct
cross-platform claim is not agreement within tolerance; it is
\emph{the same bytes}.

\paragraph{The technique.} Make the producer deterministic by
construction --- exact arithmetic on the certifying path, so no
platform's float behaviour can enter
(section~\ref{sec:pearl-exact-gpu}); fixed iteration and emission
orders, so no hash-map enumeration nondeterminism survives; canonical
serialization (section~\ref{sec:pearl-cert-rows}) --- and
cross-silicon conformance collapses to \texttt{cmp}. The split this
buys is exactly the split the system's trust story wants: the
\emph{verdict} column of any re-run table must be identical across
machines, while the \emph{time} column is allowed to vary and thereby
becomes a genuine hardware measurement instead of a confound.

\paragraph{The receipt.} The reproduction spine's semiconductor-fab
case study (the worked example of \bookref{ch:qtsl-minifab}{the
semiconductor chapter}; battery and receipts in the \path{semicon/}
tree of \path{maitria-coalition/qtsl-experiments}) was re-run
end-to-end on two boxes of different instruction sets --- a
Grace-class aarch64 workstation and an x86-64 EPYC cloud instance.
Every emitted certificate came back byte-identical to the committed
fixtures on both: ten of ten fluid-scheduling certificates, with the
certified slack values matching as exact rationals; six of six
station-cover certificates by SHA-256; the policy-separation numbers
reproduced as identical exact rationals ($\varepsilon = 1/43200$
against $1/38880$ at the near-capacity operating point); and the
deliberate failure probes --- the honest-stop cases --- still failing
correctly on both machines, producer face and checker-mutant face
alike. The checker battery ran all sixteen positive fixtures and all
thirteen mutants to their expected verdicts on both boxes. The timing
columns, meanwhile, differed by a uniform $1.7$--$1.9\times$
single-core factor across producer walls and check kernels --- a
bridge measurement that classifies every cross-box timing delta as
hardware economics, with evidence, rather than leaving it to argument.
Determinism turned an entire re-validation campaign into a column of
\emph{identical} plus one real hardware number.

\bigskip

\noindent Six decisions, one shape: in a system where every claim must
survive an independent checker, the engineering wins go to whoever
moves judgment out of the pipeline and into artefacts --- exactness
where floats were assumed, grids where guards were assumed, receipts
where intuition was assumed, bytes where tolerances were assumed. The
measurements above are what those trades cost and what they paid; the
repositories carry the artefacts to re-run them.
 \chapter{Cookbook for Agents}
\label{ch:mtk-cookbook}

\noindent\draftmark\medskip

This chapter is the toolkit's proof-production cookbook: moves in the
P\'olya shape --- the move, its recognition conditions, a worked example
with receipts, the failure mode it prevents, exercises --- re-derived
from \emph{How to Solve It} \cite{polya1945solve} for minds that
produce machine-checkable certificates. The intended reader is a large language model (or any
capable agent) doing formal verification and heavy engineering; a human
reads it profitably as the toolkit's operational philosophy stated
concretely. The register follows from the founding intent: a
meta-insight that could once only be pattern-matched through years of
accumulated practice can now simply be \emph{stated} to an agent, in
the agent's native language, and applied faster for being explained.

Every pattern carries the same five rubrics. \emph{The move}, stated
once at full generality. \emph{When it applies}: the recognition
conditions --- how you notice, mid-task, that this is the situation you
are in. \emph{Worked example from the record}: a real incident from the
workshop record that produced this toolkit, receipts included, so the
reader can go verify rather than take the chapter's word --- itself one
of the cookbook's moves: a claim worth acting on ships with an artefact
positioned to contradict it. \emph{The failure mode it prevents}: what
happens, concretely, on the path where the move is not made. And
\emph{exercises}: open legs from the producing campaigns' own backlogs
at this preview, so that working an exercise produces value, not
homework. Artefacts named in the examples migrate to the public
repositories as they graduate; per the preview note at the front of
this volume, treat a mention as a statement of what exists, not yet a
link. Status words (``shipped,'' ``proposed,'' ``in design'') are as of
this preview.

The workshop's cookbook has already grown past this chapter: two
further moves graduated after this preview's cut (one on charters as
falsifiable predictions, one a proof-assistant transparency lesson),
and a strong candidate --- adversarial review with planted doubts ---
stays deliberately unwritten here, because its authoritative statement
is a live operating document in the workshop, with measured hit rates
across the review campaigns that shaped it, and a book copy of a live
norm would drift from the norm. Each will be written from its receipts
when they migrate. Behind them, a harvest queue of roughly ninety
further candidate moves waits its turn.

\section{Sloppy Proposer, Exact Disposer}
\label{sec:ckbk-proposer-disposer}

\paragraph{The move.} Split every computation whose result must be
trusted into two roles with an explicit boundary:

\begin{itemize}
\item a \textbf{proposer} --- fast, approximate, unverified, and
  structurally \emph{untrustable}: a float library, a
  validated-numerics package, an ML search, an SMT heuristic, your own
  quick script. It produces candidates: a region, a bound, a witness, a
  decomposition, a step size.
\item a \textbf{disposer} --- exact, minimal, independently
  implemented. It re-derives, from the candidate data alone, everything
  the certificate asserts (exact-$\Q$ arithmetic, from-scratch identity
  checks) and accepts or rejects. Nothing the proposer computed appears
  in the trusted conclusion; only what the disposer re-derived.
\end{itemize}

Then police the boundary \textbf{in both directions}:

\begin{enumerate}
\item Proposer arithmetic never enters the certificate. The trusted
  computing base contains the disposer and nothing else; the proposer
  is trusted only \emph{as a search heuristic}, where being wrong costs
  retries, never soundness.
\item Disposer arithmetic never enters the proposer's inner loop. Exact
  computation inside a search loop is pure waste: the loop's job is to
  find a candidate worth checking, and it should be as cheap as
  correctness-irrelevance allows. The division of labour is the point
  --- a proposer made rigorous is as much a design error as a disposer
  made sloppy.
\end{enumerate}

This is the book's assurance stance (chapter~\ref{ch:exec-summary})
applied at the scale of a single computation, and the producer stack is
shaped by it end to end: every GPU kernel row returns witness-typed
output built for the exact tier's disposal
(chapter~\ref{ch:mtk-pearls}), and the checker at the bottom
(chapter~\ref{ch:qtsl-manual}) is the disposer of last resort.

\paragraph{When it applies.}
\begin{itemize}
\item Any time a fast-but-unverified computation must feed a
  certificate whose soundness has to be airtight. The recognition
  question: ``does the assurance story currently require this
  library/model/script to be bug-free?'' If yes, restructure so it only
  has to be \emph{useful}.
\item Any time you notice exact or expensive computation inside a loop
  that exists to search. The recognition sign is a profiler line: the
  search dominating wall-clock with arithmetic whose precision the
  final verdict never uses.
\item When choosing whether to ``just trust'' an established
  rigorous-numerics or solver library. The relevant question is not
  whether the library is good (it usually is) but whether its
  \emph{output representation} already matches something your checker
  can re-verify --- in which case de-trusting it is nearly free.
\end{itemize}

\paragraph{Worked example from the record.} The validated-ODE
de-trusting ladder: take enclosure claims from CAPD --- a well-regarded
C++ validated-numerics library for rigorous ODE integration --- and
re-found them on an exact checker, so that the assurance story never
rests on trusting C++ interval arithmetic. The architecture is the move
verbatim: three certificate species in which the proposer of every
candidate --- CAPD, an in-house float Picard iterator, anything --- is
untrusted by construction, and the checker re-derives all sequent
polynomials from the candidate data by exact $\Q$ arithmetic. Validity
of a proposed flow tube is a checkable fixed-point inclusion (the
Picard operator maps the candidate box into itself --- the
a-posteriori validation architecture formalized, kernel and all, by
Br\'ehard, Mahboubi, and Pous \cite{lipics-itp-2019-8}), so the
checker never integrates anything --- it verifies inclusion-invariance of what
the untrusted integrator proposed, with zero per-instance quantitative
side conditions: no Lipschitz constant, no contraction factor, no
step-ratio condition in any certificate.

Both directions of boundary-policing have receipts:

\begin{itemize}
\item \textbf{Direction 2, learned the expensive way.} The first
  proposer implementation ran exact Bernstein computations inside its
  candidate-search loop over radius and error-matrix parameters ---
  measured at roughly forty seconds per pass of pure waste, since the
  disposer re-derives everything anyway. The campaign's own
  retrospective names the division of labour as one of its surprises:
  the lesson was not ``be rigorous'' (the pipeline already was, at the
  disposer) but ``stop being rigorous in the wrong place.''
\item \textbf{Direction 1, cashed out when the real library arrived.}
  CAPD, vendored and driven through the frozen proposer interface,
  checker untouched: the exact-$\Q$ end box agreed with CAPD's
  uncertified claim to three decimals with a tube radius about
  $50\times$ tighter --- the cost of de-trusting was approximately zero
  enclosure quality (wall-clock about $5\times$, spent in the checker,
  which is the one place spending is correct). And the certified
  envelope \emph{extended}: a Van der Pol step size the in-house float
  proposer could not certify became certifiable with CAPD proposing. An
  isolation probe attributed the remaining failure wall to the in-house
  proposer's own search loop, not to CAPD --- a diagnosis only
  available because the interface made proposers interchangeable.
\end{itemize}

Note what interchangeability bought beyond soundness: upgrading the
proposer from quick in-house script to best-in-class library was a pure
plug-in, no re-verification of anything, because the trust boundary
never moved.

\paragraph{The failure mode it prevents.}
\begin{itemize}
\item \textbf{Trusted-base bloat}: the assurance story ``the theorem is
  only as trustworthy as the library plus the compiler run'' --- which
  is CAPD's own openly stated architecture, and precisely the thing
  being escaped. Every trusted component is a standing liability that
  must be re-audited whenever it updates; a proposer never needs
  auditing at all.
\item \textbf{The dual waste}: search loops that run at certificate
  precision, making exploration the bottleneck for zero soundness
  benefit. This one is insidious because it \emph{feels} virtuous ---
  rigour everywhere --- while actually signalling that the boundary was
  never drawn.
\item \textbf{Un-upgradeable pipelines}: when proposer and verifier
  share arithmetic (or code), swapping in a better search engine means
  re-validating the world. With the boundary drawn, proposers compete
  and the checker stands still.
\end{itemize}

\paragraph{Exercises.}
\begin{enumerate}
\item \textbf{The jet-transport adapter.} CAPD's \texttt{CnOdeSolver}
  computes derivatives of the flow map (jets). The adapter interface is
  frozen and the ladder's first-rung sequent layer is jet-ready; wire
  \texttt{CnOdeSolver} through the same seam, checker untouched.
  Success condition: goldens accept with CAPD-proposed jet enclosures;
  the checker diff is empty.
\item \textbf{Settle the flagged hypothesis.} The adapter's arrival
  recorded that one benchmark's certifiable-step wall moved
  \emph{inward} under CAPD, with a stated-but-unverified mechanism
  (centred forms losing a cancellation under the checker's crude
  absolute-value bound). Design the two-cell experiment that confirms
  or refutes the mechanism; if confirmed, the fix belongs in the
  disposer's bound, not the proposer.
\item \textbf{Generalize the isolation probe.} The probe that
  attributed the residual wall to the in-house search loop was built
  bespoke. Extract the pattern --- swap proposers across the frozen
  interface, hold the disposer fixed, bisect which side owns a failure
  --- into a reusable harness verb for any proposer/disposer pipeline.
\end{enumerate}

\section{Certify the Honest Stop}
\label{sec:ckbk-honest-stop}

\paragraph{The move.} When a line of attack ends --- the rate cannot be
improved, the constant cannot be sharpened, the method genuinely runs
out --- do not shrug and move on, and do not quietly ship the weaker
thing as if it had been the goal. \textbf{Prove the stop.} Convert ``we
couldn't push further'' into one of three certified endings, and record
which one you are in:

\begin{enumerate}
\item \textbf{The claim past the stop is false.} Produce the witness (a
  concrete counterexample instance) \emph{and}, where the machinery
  permits, a certificate-side proof that no valid certificate of the
  stronger claim exists (e.g.\ any candidate decomposition is forced
  into an inadmissible component). Then the stop is a theorem, not a
  fatigue report.
\item \textbf{The claim is true but this method cannot reach it.} Name
  the missing ingredient as a first-class object --- a new prior
  family, a stronger bridge lemma, a representation the format lacks
  --- so the stop becomes a scoped successor leg rather than a vague
  wall.
\item \textbf{The claim is true at a boundary value.} Certify the
  boundary from both sides (the bound holds at $c$; a witness attains
  $c$), so the constant ships as sharp rather than as ``the best we
  found.''
\end{enumerate}

In all three cases the artefact records the stop with the same rigour as
the positive results --- ideally \emph{inside} the same certificate and
test battery, so a future reader (or a future you, or a cheaper model
retrying the problem) cannot re-open the question without meeting the
proof.

\paragraph{When it applies.}
\begin{itemize}
\item Whenever a sweep, ladder, or optimization plateaus and you are
  about to write ``we achieved $X$'' for the $X$ you reached. The
  recognition question: \emph{is the plateau a fact about the problem
  or about your effort?} If you don't know, that is ending 1 to chase
  --- the answer is usually cheaper than the next ladder rung would
  have been.
\item Whenever you are tempted to soften a claim's statement so that
  what you proved matches it. Softening is sometimes right --- but only
  with the un-softened claim's status recorded (false? open?
  method-blocked?).
\item When banking machinery or priors for future use: an input that
  turned out to be \emph{unconsumable} (its only consumer is a false
  claim) is a result about the problem's structure, worth the two
  sentences it costs to record.
\end{itemize}

\paragraph{Worked example from the record.} The Kuramoto--Sivashinsky
energy certificate, fourth target of a campaign producing
exact-rational, independently replayed dissipation certificates for
PDEs. The campaign's earlier targets had certified decay \emph{rates}
--- viscous Burgers at the sharp rate
$\tfrac{d}{dt}E \le -2\nu E$. For Kuramoto--Sivashinsky on the marginal
domain $\mathbb{T}_{2\pi}$ the analogous hope --- any positive rate
$\alpha > 0$ --- is \textbf{false}, and the artefact treats that
falsity as a deliverable (ending 1, with ending 3 nested inside):

\begin{itemize}
\item \textbf{The witness}: at mode $k = 1$ ($u = \varepsilon\sin x$)
  the dissipation form $k^4 - k^2$ vanishes and the nonlinearity is
  energy-neutral, so $\tfrac{d}{dt}E \le -\alpha E$ fails for every
  $\alpha > 0$. One line, checkable by hand.
\item \textbf{The certificate-side impossibility}: the producer carries
  an assert that any nonnegative-weight decomposition of the
  $\alpha$-remainder forces weight $-\alpha$ on the $u_2^2$ square ---
  unconditional in the free weight. The stronger claim is not merely
  un-found; it is un-certifiable in the calculus, and the artefact
  proves that.
\item \textbf{The battery enforcement}: a mutant recomputes
  \emph{everything} consistently at a positive rate --- every jet
  identity, every replay green --- and is rejected \emph{only} by the
  negative weight. The stop is enforced in the executable spec
  (section~\ref{sec:ckbk-mutant-battery}): a future certificate
  claiming a positive rate on this domain will be refused by machine,
  not by memory.
\item \textbf{The named residue} (ending 2): the artefact's own
  stopping-point section states what the composed result is ---
  energy monotonicity at the marginal domain, sharp, the strongest true
  statement for this functional on this domain, with the stop
  \emph{provable} --- and then names what the real literature benchmark
  ($L > 2\pi$ mean-energy bounds, where energy genuinely grows through
  the unstable band) additionally requires: a background functional and
  a genuinely new prior family (Gagliardo--Nirenberg/interpolation
  rows). Its closing words are this section's thesis in miniature:
  ``named next leg, not tonight's overclaim.''
\item \textbf{The unconsumed input, recorded}: a banked spectral prior
  ($\int u_2^2 \ge \int u_0^2$) is left unconsumed, on the record ---
  the campaign notes its only would-be consumer is the false rate
  claim. Two sentences, and no successor will wonder why it sits
  unused.
\end{itemize}

A sibling specimen, same campaign, showing ending 3 in both directions:
the Fisher--KPP certificate proves the invariant region $[0,1]$ is
maximal by mutants --- enlarged above to $[0, \tfrac98]$, all
bookkeeping recomputed consistently, rejection comes only from a
genuinely negative Bernstein coefficient that no subdivision can repair;
extended below, the endpoint row itself fires. ``The certificate cannot
be strengthened past the true boundary in either direction'' is a
tested property of the checker, not a sentence in prose.

\paragraph{The failure mode it prevents.}
\begin{itemize}
\item \textbf{The shrugged dead-end}: an unrecorded stop is a question
  that reopens. Someone (often a cheaper agent, often you-next-week)
  retries the false claim, burns the same compute, and --- worse ---
  may ``succeed'' via a subtly unsound route precisely because nothing
  in the machine refuses it.
\item \textbf{The laundered stop}: shipping the weaker claim in the
  stronger claim's rhetorical clothing (``we certify decay'' when the
  hoped-for rate silently became $0$). The reader's trust is spent on
  the discrepancy when they find it --- and in adversarial review they
  will.
\item \textbf{Effort-plateau mistaken for truth-plateau}: without the
  falsity/impossibility half, ``sharp'' means ``we got tired.'' With
  it, sharpness is a two-sided certificate and the constant is
  settled.
\end{itemize}

\paragraph{Exercises.}
\begin{enumerate}
\item \textbf{The symbolic-domain carriage.} Carry the domain length
  $L$ symbolically through the Kuramoto--Sivashinsky certificate: for
  $L < 2\pi$ the same form yields a \emph{positive}-rate nonlinear
  decay; for $L > 2\pi$ it yields nothing --- the energy method
  genuinely ends there. Certify both sides of that boundary, making the
  method's own limit an ending-3 artefact.
\item \textbf{The background-functional leg.} The named ending-2
  residue: build the interpolation prior family (the first genuinely
  new prior species in the campaign) and attack the $L > 2\pi$
  mean-energy benchmark. The honest-stop record is the specification of
  exactly what is missing.
\item \textbf{Settle a flagged wall.} The CAPD adapter's arrival
  (section~\ref{sec:ckbk-proposer-disposer}) left one
  certifiable-step wall's inward move as a stated-but-unverified
  hypothesis. Determine which ending it is: a true wall of the bound
  (certify it), or an artefact of a crude inequality (fix the disposer
  and move it back). Either way the stop stops being a guess.
\end{enumerate}

\section{The Mutant Battery Is the Spec}
\label{sec:ckbk-mutant-battery}

\paragraph{The move.} A checker's meaning is not what it accepts ---
accepting the golden certificate is table stakes --- but what it
\textbf{rejects}. So specify the checker by building its rejections: a
battery of \emph{mutants}, each a certificate wrong in one deliberate,
named way, each asserted to be rejected \textbf{at its intended check
with its intended message}. The battery is then the executable
specification of the trust boundary: every mutant is a theorem of the
form ``a certificate lying in \emph{this} particular way cannot pass,''
and the battery as a whole delimits what an ACCEPT verdict actually
asserts.

Two disciplines make the battery a spec rather than a pile of negative
tests:

\begin{enumerate}
\item \textbf{Consistency mutants are the genre that carries the
  spec.} The strongest mutants recompute \emph{everything downstream
  of the lie consistently}, so that every identity, every replay, every side
  computation is green, and rejection can only come from the single row
  that carries the claim's truth. A mutant rejected by accidental
  earlier breakage (a malformed field, an arithmetic slip you didn't
  intend) tests nothing about the check you meant to test --- which is
  why each rejection's \emph{location and message} are asserted, not
  just the rejection.
\item \textbf{Cover the lie taxonomy, not just arithmetic.} Wrong
  numbers are the easy genre. The battery must also include: sharpness
  violations in \textbf{both} directions (claim strengthened past the
  true boundary / input weakened below it), laundered caveats (an
  obligation's status upgraded to one no receipt backs), provenance
  breaks (data whose justifying source is removed, dangling, or
  unlisted), structural smuggles (declaring only a slice of the term
  the side condition should see), and --- where claims are mechanism-
  or context-scoped --- transport violations (a true-in-context-A claim
  rewired into context B).
\end{enumerate}

The same discipline runs at every tier of this toolkit. The rule
catalogue (chapter~\ref{ch:qtsl-manual}) is its proof-theory face: the
roster is machine-checked, every mechanized rule carries receipts, and
the coverage dashboard up front is the candid map of what the assurance
edifice does and does not yet refuse. The composition-pattern inventory
(chapter~\ref{ch:qtsl-composition-patterns}) is its modelling face: the
refusals do real work there too, encoded structurally rather than
left to the modeller's vigilance.

\paragraph{When it applies.}
\begin{itemize}
\item Whenever you write a checker, validator, decoder, or replay tool
  whose ACCEPT someone will trust. If the mutant battery does not
  exist, the trust boundary exists only in your intentions.
\item Whenever a checker gains a new check: the check's first artefact
  is the mutant it refuses. A check without a firing mutant is
  unfalsifiable --- you cannot distinguish ``this check works'' from
  ``this check is dead code.''
\item When reviewing someone else's verification claim. ``10/10 checks
  passed'' is uninformative until you know what the 11th --- the one
  built to fail --- would have done. Ask for the battery before the
  benchmark.
\end{itemize}

\paragraph{Worked example from the record.} The PDE certificate
campaign of the previous section: exact-rational dissipation
certificates for Burgers, Fisher--KPP, and Kuramoto--Sivashinsky, each
with a Python producer and an independent exact-fraction checker
sharing zero polynomial code, each shipping a battery with per-mutant
rejection messages asserted. The taxonomy, one specimen each, all on
disk:

\begin{itemize}
\item \textbf{Consistency $+$ sharpness (above)}: a Burgers mutant
  claims decay rate $\alpha = 3\nu$ (true rate: $2\nu$) and recomputes
  all bookkeeping consistently --- all arithmetic passes and rejection
  comes only from the negative combination weight: the checker enforces
  the sharpness boundary, not mere self-consistency. The
  Kuramoto--Sivashinsky twin pushes a linear rate of $24$ to $25$ and
  is refused only by the forced weight $-\tfrac1{16}$.
\item \textbf{Sharpness (below) $+$ unrepairability}: the Fisher--KPP
  mutant widens the certified invariant region $[0,1]$ to
  $[0,\tfrac98]$; the endpoint rows still pass (any $[0,b]$ with
  $b \ge 1$ is inward at the ends) and the refusal comes from one
  genuinely negative Bernstein coefficient equal to $f(\tfrac98) < 0$
  --- no subdivision can repair it, so the battery proves the boundary
  is the true one rather than a resolution artefact.
\item \textbf{Structural smuggle}: advection-smuggle mutants add a term
  the semilinear/jet-linearity side condition must see; every
  arithmetic identity passes and only the structure check fires ---
  establishing that the side condition, not the arithmetic, is what
  excludes that class of false instance.
\item \textbf{Provenance edges}: the Fisher--KPP deficit-decay
  remainder is the campaign's first that genuinely needs a bounding
  box, and the box is supplied \emph{by another claim in the same
  certificate} with an explicit provenance edge. Three mutants break
  the edge three different ways --- source removed, dangling, supplier
  unlisted --- and each is refused by the provenance check
  specifically.
\item \textbf{Laundered caveat}: both laundering flavours --- a
  hypothesis quietly dropped from the declared list, and an
  obligation's status promoted to ``proven'' with no discharge receipt
  --- are refused by the ledger discipline
  (section~\ref{sec:ckbk-prior-ledgers}), tested both at the ledger and
  inside the discharge object.
\item \textbf{Transport (the campaign's centrepiece)}: a mutant takes
  the linear-mechanism spectral-gap rate claim --- true, certified,
  sharp \emph{for the linear mechanism} --- and rewires it onto the
  nonlinear flow, where mode-1 pumping makes it false. Every jet
  identity, weight, and spectral replay is green; rejection comes
  \textbf{only} from the mechanism-scoped premise-transport table. In
  the artefact's own words: ``the checker refusing a false theorem
  whose every checkable identity is true.'' That single mutant is the
  spec of what the transport table means.
\end{itemize}

Scale receipts, for calibration: the three batteries run 31, 18, and 29
tests (29, 17, and 28 mutants respectively), messages asserted, with a
combined regression replaying all three checkers green --- and each
battery was built within its target's single working session: hours,
not weeks.

\paragraph{The failure mode it prevents.}
\begin{itemize}
\item \textbf{The golden-only checker}: passes its examples, accepts a
  subtly-false certificate the day an adversary (or an optimizer, or a
  bug) produces one. Without mutants you have tested the happy path of
  a tool whose entire purpose is the unhappy path.
\item \textbf{Dead checks}: a check that nothing exercises can be
  broken by refactor and no test notices. The firing mutant is the
  check's liveness monitor.
\item \textbf{Accidental-rejection illusions}: a battery whose mutants
  die of malformedness rather than of their intended lie certifies
  nothing --- the dangerous mutant is the \emph{well-formed} one, and
  only location-asserted rejections distinguish the two.
\item \textbf{Unfalsifiable sharpness and scope claims}: ``sharp,''
  ``maximal,'' ``linear mechanism only'' as prose adjectives, versus as
  mutants the machine refuses. The former decay into folklore; the
  latter are enforced.
\end{itemize}

\paragraph{Exercises.}
\begin{enumerate}
\item \textbf{The spectral wire rows.} When the jet certificate
  schema's spectral wire family ships, write its battery \emph{before}
  its first consumer: wrong-constant discharge, missing witness mode,
  transport onto an unlicensed mechanism --- the taxonomy above,
  instantiated at the wire tier.
\item \textbf{Retrofit audit.} The neural-supermartingale certificate
  suite grew its mutant receipts incrementally across its campaign.
  Audit its batteries against the taxonomy: which genres are missing
  (transport? provenance? both-direction sharpness?), and add the
  cheapest missing genre with rejection locations asserted.
\item \textbf{The battery-of-the-battery.} The location-assertion
  discipline currently lives in test code by convention. Build the
  meta-check: a harness pass that runs each mutant, captures
  \emph{which} check fired, and fails the battery if any mutant is
  refused earlier than its declared target --- turning discipline 1 of
  this section from prose into a tool.
\end{enumerate}

\section{Declared-Prior Ledgers}
\label{sec:ckbk-prior-ledgers}

\paragraph{The move.} Every analytic fact your certificate
\textbf{consumes but does not prove} --- a functional inequality, a
comparison principle, a semantic bridge between the certificate's
algebra and the theorem's meaning --- is carried by an explicit
\textbf{ledger} inside the certificate, under a \textbf{closed status
vocabulary}.
Concretely:

\begin{enumerate}
\item \textbf{Declare, never absorb.} An unproven ingredient appears as
  a named entry (\texttt{assumed}, \texttt{named-unproven}, \ldots),
  not as an invisible premise folded into the prose or --- worse ---
  into the checker's silence. The certificate's ACCEPT then
  means: \emph{these claims hold, conditional on this list.}
\item \textbf{Close the vocabulary and back every status.} The checker
  enforces that statuses come from a fixed set and that any status
  implying a proof has its receipt present and replaying green. An
  obligation promoted to ``proven'' with no receipt is a \emph{rejected
  certificate} --- laundering is a machine-refused move, not a style
  violation.
\item \textbf{Discharge as an arc, with a tie.} Because the ledger is
  explicit, an assumption can be discharged \emph{later,
  incrementally}: a new certified object replaces \texttt{assumed} with
  a proof-backed status, and a \textbf{tie check} verifies the
  discharged object is the same fact the consuming row uses (same
  integrand, same constant) --- so an internally valid proof of a
  \emph{neighbouring} fact cannot silently license the claim.
\end{enumerate}

The ledger converts ``is this pipeline rigorous?'' from a judgement call into a query: read the list. And it converts rigour-improvement from a
rewrite into a diff: shrink the list, one discharge at a time. The rule
catalogue's coverage dashboard (chapter~\ref{ch:qtsl-manual}) is the
same move applied to the assurance edifice itself: a declared,
per-rule account of which correspondences are certified and which
remain pending, in data rather than in reputation.

\paragraph{When it applies.}
\begin{itemize}
\item Whenever a formal artefact meets analysis (or any outside
  theory): the bridge facts --- Parseval, Gr\"onwall, a comparison
  principle, a solution concept's energy inequality --- are exactly the
  ingredients a polynomial/arithmetic checker cannot re-derive. They
  must go \emph{somewhere}; the only choice is ledger or shadow.
\item Whenever you catch yourself writing ``standard'' or
  ``well-known'' next to a fact your checker never sees. Standard facts
  have names; name them in data, with a status.
\item When planning a campaign: rungs you cannot yet prove are not
  blockers if they can be \emph{declared} --- the pipeline ships
  sound-modulo-ledger on day one, and the ledger is the campaign's
  remaining-work list, machine-readable.
\end{itemize}

\paragraph{Worked example from the record.} The Wirtinger discharge arc
in the Burgers energy certificate. The sharp-rate claim
$\tfrac{d}{dt}E \le -2\nu E$ needs the Wirtinger/Poincar\'e inequality
$\int u_1^2 \ge \int u_0^2$ for mean-zero data on $\mathbb{T}_{2\pi}$
--- an analytic fact the exact-$\Q$ jet checker cannot see. It shipped
\textbf{declared-assumed}: a named prior in the certificate, with the
mean-zero hypothesis's \emph{invariance} itself certified (a pure
divergence row), and the declaration stated out loud in the campaign
record. The pipeline was already refusing laundering at this stage: the
Fisher--KPP sibling's battery included both laundering mutants
(hypothesis quietly dropped; status upgraded receipt-free),
machine-refused.

A later session, the same day, discharged the prior
\emph{in-certificate}: on the torus it diagonalizes exactly, so the
checkable kernel is the mode polynomial $q(k) = k^2 - 1 \ge 0$ on
$|k| \ge 1$ --- certified by an exact-$\Q$ shift row
($q(1+t) = t^2 + 2t$, nonnegative coefficients, covering every mode
uniformly), with sharpness as a width-zero eigenvalue enclosure
$\lambda_1 \in [1,1]$ (lower side the shift row, upper side the witness
mode $k = 1$). The ledger flips: \emph{no assumed priors remain}. Three
details carry the section:

\begin{itemize}
\item \textbf{The tie.} The checker verifies the consuming row's
  integrand \emph{is} $u_1^2 - c\,u_0^2$ at the discharged constant $c$
  --- an internally valid discharge of a different constant cannot
  silently license the claim, and a consistency mutant enforces
  exactly that. Discharge without a tie is a new laundering surface;
  the arc closed it the same session it opened.
\item \textbf{The residue is itself ledgered.} The discharge consumes
  four analytic bridge facts of its own (Parseval,
  derivative-coefficient, mean-is-coefficient-zero, termwise
  summation). They did not vanish into prose: each is a named
  \texttt{bridge:} entry citing its mathlib lemma at a pinned commit,
  citations checked live at writing time, with one remaining gap named
  (the $H^1$ extension of the derivative form, which holds as cited at
  $C^1$). The ledger got \emph{shorter and sharper}, not empty by
  fiat.
\item \textbf{Laundering mutants re-targeted post-discharge.} With a
  discharge object present, the new lie is tampering \emph{inside} it,
  or claiming spectral-certified status without it; the battery grew
  nine mutants aimed at the ledger and the discharge both. The status
  vocabulary stayed closed the whole way.
\end{itemize}

Cost calibration, because trained priors on verification cost run
years stale: the whole arc --- declare, ship sound-modulo-ledger,
discharge with sharpness, tie, re-target the battery --- completed
within one working day across three sessions. Declaring first did not
slow the campaign; it is what \emph{let} the campaign ship while the
discharge waited its turn. And the banked bonus paid forward the same
day: the discharge artefact pre-certified the higher-order siblings
($q = k^4 - k^2$, $k^4 - 1$), which the Kuramoto--Sivashinsky target
then consumed as its spectral priors.

\paragraph{The failure mode it prevents.}
\begin{itemize}
\item \textbf{The silent assumption import}: a ``verified'' pipeline
  resting on an unnamed analytic fact --- the reader believes the
  checker saw everything; the gap surfaces in adversarial review (or in
  production) as ``what exactly does ACCEPT mean here?'', and the
  answer costs the pipeline its credibility even where it was sound.
\item \textbf{Laundering under deadline}: statuses drifting upward
  (``assumed'' $\to$ ``known'' $\to$ ``proven'') across retellings,
  with no machine in the path to refuse the drift. This is the
  certificate-level twin of the documentation failure in which a
  distilled guess hardens into quoted law: the repair in both cases is
  provenance enforced at the artefact, not vigilance in the prose.
\item \textbf{The wrong-theorem discharge}: proving a true neighbouring
  fact (different constant, different function space) and letting it
  license the claim --- refutable only if the tie between discharge and
  consumer is checked, which is why the tie is part of the move and not
  a refinement of it.
\item \textbf{All-or-nothing rigour}: without a ledger, ``not fully
  proven'' and ``not shippable'' collapse into one, and campaigns stall
  at their hardest lemma. With one, soundness-modulo-a-shrinking-list
  is a publishable, checkable state.
\end{itemize}

\paragraph{Exercises.}
\begin{enumerate}
\item \textbf{The Lean bridge lemmas.} The four \texttt{bridge:}
  entries in the discharge (plus the linear-mechanism spectral-support
  obligation in the Kuramoto--Sivashinsky artefact) are named with
  mathlib citations at a pin. Formalize the statements as Lean lemmas
  against that pin, flipping each ledger entry from cited to
  proof-backed --- and add the battery mutant that tampers with the new
  receipts.
\item \textbf{The named gap.} The derivative-coefficient bridge holds
  at $C^1$ data; the $H^1$ extension is an explicitly named gap in the
  ledger. Close it --- part of the exercise is determining whether the
  gap closes at the pin or needs an upstream mathlib contribution.
\item \textbf{A new prior species, ledger-first.} The $L > 2\pi$
  Kuramoto--Sivashinsky benchmark needs interpolation
  (Gagliardo--Nirenberg-type) priors --- the first genuinely new prior
  family since Wirtinger. Introduce them the ledger-first way: declared
  species with closed statuses and laundering mutants \emph{from
  birth}, so the discharge arc has its rails before anyone attempts the
  discharge.
\end{enumerate}

\section{Cheap Before Expensive}
\label{sec:ckbk-cheap-first}

\paragraph{The move.} Before committing the expensive resource --- a
build, a rented-GPU run, a finer subdivision, a proof campaign, a
frontier-model token stream --- spend the cheap instrument that could
make the expensive spend unnecessary, radically smaller, or
better-aimed. The move has three forms, in ascending machinery:

\begin{enumerate}
\item \textbf{The cheap question.} Does the record, or a free
  qualitative argument, already settle it? Read the project log before
  the build; check for the monotonicity observation before assembling
  the sharp ladder. The cheapest ending to an expensive plan is
  discovering it already happened, or was never needed.
\item \textbf{The cheap pass.} When cost scales with a count (cells,
  claims, pieces, obligations), run the inexpensive consolidation
  \emph{before} paying the per-unit price --- and demand the
  consolidation be loss-free where the theory permits it (same bound,
  fewer pieces), so the saving is pure.
\item \textbf{The cheap measurement.} When two routes to the same
  certified quantity compete, race them at \emph{equal budget} on a
  real instance before scaling either. ``Which route wins'' is an
  empirical question with a cheap experiment, not a doctrine --- and
  the fair comparison is at matched spend, since anything beats
  anything given more budget.
\end{enumerate}

Two boundary clauses guard the move. First, it orders
instruments; it never licenses deferral --- a plan that stops at the
cheap instrument needs a stated mechanism for when the expensive one
fires, not a vibe; and when the expensive machinery gets built anyway,
it is \emph{banked, not wasted}: retained explicitly for the sibling
case where the cheap argument does not apply. Second, the currency is
latency and aim, not primarily money: in a regime of abundant
implementation capacity, what the move conserves is the time the
expensive spend takes, the attention it occupies, and the risk of
aiming it at a target the cheap instrument would have moved.

\paragraph{When it applies.}
\begin{itemize}
\item About to dispatch or build something whose success condition
  might already be satisfied. The project record is the cheapest
  instrument available, and the check extends to every dispatch, not
  only the ones that feel uncertain.
\item About to turn a budget dial --- subdivision depth, piece count,
  degree, machine-hours, review rounds --- because the current answer is
  not tight enough. First ask: is there a cheap restructuring (a
  coarsen pass, a representation change) that dominates dial-turning?
  The recognition sign is a cost that scales like
  $\mathcal{O}(h^{-k})$ against an error shrinking like
  $\mathcal{O}(h)$: the dial is the losing side of an exponent gap.
\item Holding a certified artefact whose slack is visible (margins,
  lift values, mutant-refusal exponents): slack is prepaid budget, and
  a cheap pass can often spend it on consolidation with zero loss.
\item Routing work to minds: at any date there is a quality/cost/speed
  frontier of available models, and selection is by context, not by
  rank --- the cheap capable point first, the frontier model where
  context-understanding or taste binds. The same logic governs how
  much deliberation to request from a given model: effort is a dial,
  and the cheap setting is the right default for work that does not
  need the expensive one.
\end{itemize}

\paragraph{Worked examples from the record.}

\textbf{The consumed dispatch.} An overnight agent was dispatched
toward what its brief framed as a live build: an exact certifying-tier
GPU kernel family, on a rented machine. The agent's first instrument was
one read of the project log --- which showed the tier had shipped
\emph{that morning} (fused exact Bernstein slot kernels, 11/11
device-versus-host-reference bit-exact, with two successor tasks
already consuming them). Verdict: dispatch adjudicated consumed; a
candidate follow-on was named without being fired; total agent runtime
132 seconds. What a rented-machine campaign would have re-built, a
two-minute read retired.

\textbf{The slack-funded coarsen pass.} A family of fixed-shift
projection-kernel certificates was born-certifiable with lift $0$ ---
the theorem's manufactured slack absorbed all float noise. That slack
is prepaid budget, and a cheap coarsen pass spent it: one
representative solution went from 640 uniform pieces to 26 nonuniform
dyadic pieces (89 claims), \emph{bound bit-identical}, lift still $0$
--- the receipts' own headline was ``coarsening pays $14\times$.''
Every downstream cost of that row (certify, reference-check, emit,
fast-check, per claim) shrank by the same factor before any expensive
tier ran. One detail worth copying: the coarsen pass's excess
ledger initially under-billed (an interpolant's convexity overshoot
went unbilled), and it was the pinned test suite that caught it --- the
cheap pass runs under the same verification discipline as everything else.

\textbf{Tube versus hull at equal claim budget.} A design note posed
its own pricing question --- a curved-boundary ``tube'' upgrade ``pays
exactly when hull conservatism is the binding loss --- measure first''
--- and the measurement ran both routes to a certified enclosure of the
same switching integral, same exact tier, budget counted in claims. At
about 10 claims the tube already matched the hull's 112-claim width; at
59 claims the tube was ${\sim}1300\times$ tighter than the hull at 237;
matching the tube's width by dial-turning the hull would extrapolate to
subdivision depth ${\sim}17$, $\mathcal{O}(10^5)$ claims. The
receipts' verdict --- on this benchmark the tube wins by two to four
orders of magnitude at every matched budget --- was scoped to the
benchmark, with both scaling laws measured. One cheap head-to-head, and
the budget dial was retired as the losing instrument for this loss
class \emph{before}, not after, a $10^5$-claim campaign.

A sibling specimen from the proof side: a closed-loop encounter
certificate needed ``stays separated forever after exit,'' and the
author built the full ray-clearance LP ladder --- then found the
certified exit condition makes the exit face \emph{recede}, so
post-exit eternity is free model algebra on already-certified facts.
The LP machinery was demoted as unnecessary for that class and
explicitly retained for the non-receding sibling classes: the cheap
question asked late still pays, and the expensive machinery, banked, is
not waste.

\paragraph{The failure mode it prevents.}
\begin{itemize}
\item \textbf{The duplicate build}: re-doing shipped work because the
  record went unread. The cost is not only the spend --- a re-built
  sibling diverges from the shipped one, and reconciling them later
  costs more than either.
\item \textbf{Dial-turning into an exponent gap}: paying
  $\mathcal{O}(h^{-k})$ for $\mathcal{O}(h)$ relief because scaling
  the incumbent felt like progress and restructuring felt like a
  detour. The receipts above put real numbers on how bad this gets:
  five orders of magnitude of claims to match a route change.
\item \textbf{Doctrine substituting for measurement}: ``the
  sophisticated route wins'' (or ``the simple route suffices'')
  adopted from taste, when the equal-budget race costs minutes. The
  tube-versus-hull brief explicitly allowed the inverse outcome
  --- the measurement declined it, and only the measurement could.
\item \textbf{Thrift theatre}: the inverted failure --- reading this
  section as ``spend less.'' The expensive spend is often correct; if
  you are an agent whose cost intuitions were trained on older
  tooling, your priors on verification cost in particular run stale
  and high (the receipts throughout this chapter say hours where the
  prior says months). The move is about what you learn \emph{before}
  aiming the big spend, never about flinching from it.
\end{itemize}

\paragraph{Exercises.}
\begin{enumerate}
\item \textbf{The agreement-matrix verb.} The consumed dispatch above
  named a follow-on without firing it: a consolidated
  float-versus-exact agreement-matrix verb over existing machinery,
  its content currently distributed across two campaigns' receipts.
  Build it as a standing cheap instrument (one CLI verb) so the cheap
  question ``do the tiers agree, where, at what exponent?'' costs one
  invocation before any future expensive exact-tier run. Success
  condition: the verb runs against the committed receipts and emits the
  matrix; its output for one known row matches the published receipt.
\item \textbf{Coarsening at degree 2.} The projection-kernel receipts
  state the boundary plainly: the degree-2 fixed-shift row is a-priori
  sizing only --- no monotone theorem at degree 2 (the Korovkin
  boundary, measured in the same receipts as a step change in
  monotonicity violation). The coarsen pass is theory-backed only at
  degree 1. Determine whether a-posteriori coarsening pays at degree
  2: coarsen, re-certify, compare bound and claim count against the
  uncoarsened row. Success condition: a new receipts row with
  pieces/claims/bound/lift columns, whichever way the verdict falls ---
  an inverse outcome would be a finding rather than a failure.
\item \textbf{The next binding-loss measurement.} The tube route won
  where hull conservatism was the binding loss on a two-dimensional
  benchmark with one curved boundary. Before extending tubes (higher
  dimension, multiple switching surfaces), run the pricing question
  forward: construct the next benchmark, measure the hull baseline's
  scaling first, and only build the tube generalization if the
  exponent gap is present there too. Success condition: a
  receipts-shaped table for the new benchmark with both routes at
  matched budgets --- or the measured absence of the gap, banked as
  the reason the generalization waits.
\end{enumerate}

\section{Scope and Frame}
\label{sec:ckbk-scope-frame}

\paragraph{The move.} Every claim lives in a frame --- a mechanism, an
interface, a semantics under which its sentences are read --- and most
catastrophic errors in certificate work are not wrong identities but
\emph{right identities read in the wrong frame}. The move is to pin the
frame as a first-class, machine-checked object before building, and
then enforce scope in three places:

\begin{enumerate}
\item \textbf{Premises transport only where licensed.} A fact proven in
  one frame (one mechanism, one solution regime, one data class) does
  not follow the notation to a neighbouring frame. Record, per premise,
  the frames where it holds and the mechanism that carries it there ---
  an invariance row, a bridge obligation --- as \emph{checker data}, so
  that using a premise outside its license is a refusal, not an
  oversight.
\item \textbf{Changes enter only through the seam.} When an untrusted
  component is upgraded or swapped, the modification's entire footprint
  is the declared interface; everything on the trusted side runs
  byte-for-byte unchanged, and that claim is grep-verifiable, not
  asserted.
\item \textbf{Out-of-scope reads go vacuous, never false.} When the
  frame itself is discovered ambiguous --- the same certificate
  admitting two readings, one of them unsound --- the repair is
  semantic before it is mechanical: declare the conditional reading
  normative (the certificate asserts its conclusion \emph{given} its
  declared environment), so a certificate whose hypotheses cannot be
  jointly met claims nothing, rather than something false. Then add the
  mechanical check that makes the declaration well-formed.
\end{enumerate}

The common discipline: the frame is stated where the machine can refuse
against it. A frame that lives in prose (``this only applies to the
linear case,'' ``don't touch the checker'') decays into folklore; a
frame that lives in a transport table, a frozen function signature, or
a normative semantics section with its checker pass is enforced by
machine.
The composition-pattern inventory
(chapter~\ref{ch:qtsl-composition-patterns}) runs the same discipline
at the model tier: which certificates transport across a composite
boundary, and which are refused, is part of each pattern's definition.

\paragraph{When it applies.}
\begin{itemize}
\item Any artefact with more than one mechanism, regime, or model of
  the same system --- linearizations beside nonlinear flows, discrete
  beside continuum, per-piece beside global. The recognition question:
  ``could every \emph{identity} in this proof be true while the
  \emph{theorem} is false, because a premise crossed a mechanism
  boundary?'' If the answer is not obviously no, the transport
  structure is missing.
\item Any planned upgrade of a component inside a verified pipeline.
  The recognition sign is the phrase ``and we'll just also need to
  tweak the checker'' --- the seam is about to leak; freeze the
  interface first and make the swap's footprint checkable.
\item Any review finding of the shape ``under reading A this is sound,
  under reading B it accepts a lie.'' That is not a bug in a check; it
  is an unpinned semantics, and patching the check without declaring
  the reading leaves the ambiguity for the next consumer.
\item At design time, before requirements: the frame question (``what
  kind of thing is this, really?'') outranks every detailed
  requirement, because a wrong frame silently poisons all of them. The
  workshop's own worst-case specimen was a requirements document
  written against an implicit wrong frame --- the modelling substrate
  read as an expressive logic rather than as a logical framework ---
  whose review verdict was that most of the work needed re-doing. The
  frame question costs a paragraph; its absence costs the document.
\end{itemize}

\paragraph{Worked examples from the record.}

\textbf{Mechanism-scoped premise transport.} The
Kuramoto--Sivashinsky certificate of
sections~\ref{sec:ckbk-honest-stop} and~\ref{sec:ckbk-mutant-battery}
certifies two decay claims on two mechanisms: the nonlinear flow
(energy monotone, sharp at $\alpha = 0$) and the linear mechanism
restricted to modes $|k| \ge 2$ (rate $-24$, sharp at mode 2). The
spectral-gap premise is licensed for the linear mechanism \emph{only}
--- under the nonlinear flow, mode 1 is pumped from pairs like
$(3, -2)$ and the premise is not invariant --- and that license is not
prose: every mode-exclusion premise carries a transport table keyed by
mechanism (invariance-row kind for mean-zero, bridge-obligation kind
for spectral support). The battery's centrepiece rewires the gap-sector
rate claim onto the nonlinear mechanism: every jet identity, weight,
and spectral replay is green, and rejection comes only from the
transport table. Without the transport structure, nothing else in the
certificate would have objected.

\textbf{The frozen seam under a proposer swap.} When the real
validated-numerics library replaced the in-house float proposer
(section~\ref{sec:ckbk-proposer-disposer}, seen here from the scope
side), the entire change entered through one frozen interface: the
adapter builds its Taylor-model polynomial and feeds it through the
frozen proposal entry point, and the exact fixed point, split-tree
certification, and checker run byte-for-byte unchanged. The claim is
grep-verifiable and was re-verified at this chapter's writing: the
override parameter appears in the proposer module (the seam's
definition) and the adapter's own files, and nowhere in the checker ---
the checker's text is untouched by construction, not by promise. The
payoff compounds: because the swap's footprint was the seam, an
isolation probe could later attribute a residual certification wall to
the in-house proposer's search loop rather than to the library --- a
diagnosis that only type-checks when proposers are interchangeable
behind a frame that never moved.

\textbf{Vacuous, not false.} An adversarial review round found that a
new integral-monotonicity rule in the jet certificate schema, read
under plain all-pre-valuations sequent semantics, admitted a false
ACCEPT: a precondition constraining a bound integral wire, plus
free-sign equality cancellation, licensed a conclusion no trajectory
satisfies. The repair came in the order this move prescribes ---
semantic half first: the contract declares the binding-environment
reading normative (a certificate asserts its sequents over environments
satisfying its preconditions \emph{and} every binding it declares ---
the declared-hypothesis pattern), under which the attacking certificate
makes the hypothesis set unsatisfiable and the certificate
\emph{vacuous} --- never a false assertion, because under this
semantics no environment is claimed. Then the mechanical half: a new
certificate-global binding-functionality pass (duplicate bindings of
one wire reject; the negative control fires in isolation) makes the
declared environment a well-formed function of wires. The residue is
scoped rather than hidden: until the registered-environment lane ships,
a consumer composing these certificates owns the hypothesis-set
satisfiability check the same way it owns assumed rows --- stated in
the normative section, where a composer will meet it.

\paragraph{The failure mode it prevents.}
\begin{itemize}
\item \textbf{Premise smuggling}: a true lemma licensing a false
  theorem across a mechanism boundary. This is the failure that green
  batteries cannot catch from below --- every constituent identity
  replays --- which is why the transport table has to exist as its own
  refusal channel.
\item \textbf{Seam leak}: an ``upgrade'' that touches the trusted side
  ``just a little,'' after which nothing about the pipeline is a
  plug-in and every future swap re-opens verification of the world.
  The byte-for-byte-unchanged claim, grep-backed, is what keeps
  upgrade cost at zero on the trusted side.
\item \textbf{The silently-false certificate}: an artefact sound under
  its author's intended reading and unsound under the reading a
  consumer will actually use. Declaring the semantics normative in the
  contract --- and making the unsound reading \emph{forbidden} rather
  than merely undiscussed --- converts a latent soundness hole into a
  documented conditional.
\item \textbf{Frame-blind requirements}: detailed specification
  against an implicit wrong frame, discovered at review, fixed by
  ``most of this needs re-doing.''
\end{itemize}

\paragraph{Exercises.}
\begin{enumerate}
\item \textbf{The registered environment.} The jet schema names the
  successor leg for its own residue: carry the declared binding set as
  first-class certificate data (a registry object) with a sound
  existential-elimination/composition rule, so consumers compose
  certificates without carrying the conditional semantics by
  convention --- the frame promoted from normative prose to checked
  data, this section's move applied to its own example. Success
  condition: two certificates composed through the rule, with a mutant
  pair (a composition whose environments conflict) refused at a named
  check.
\item \textbf{Transport tables as shared machinery.} Mechanism-scoped
  premise transport currently lives in the Kuramoto--Sivashinsky
  artefact's checker. Extract it into the PDE campaign's shared
  checker machinery so the Burgers and Fisher--KPP artefacts consume
  the same transport module --- the fan-out built it once where it was
  needed; the unification makes it the family's frame discipline.
  Success condition: all three batteries green through the shared
  module, plus one new transport mutant per retrofitted artefact.
\item \textbf{Unrepresentable out-of-frame states.} The modelling
  substrate's tensor-contraction definition
  (chapter~\ref{ch:geolog-reference}) was deliberately ratified in the
  form whose side condition --- every un-hit index is an output ---
  makes the value-relevant free choice (a dropped axis silently
  scaling results) unrepresentable rather than discouraged; the side
  condition is exactly the proof-irrelevance theorem. Realize
  contraction rows with that condition enforced by construction, and
  write the property test that witnesses proof-irrelevance (no
  auxiliary choice affects any value). Success condition: the lowering
  rejects ill-framed index maps at the type/row level, and the
  property test runs green against the reference semantics.
\end{enumerate}

\section{Separate the Timescales}
\label{sec:ckbk-timescales}

\paragraph{The move.} When a system's rates split into well-separated
bands, do not pay the stiffness tax in the checker --- restructure the
\emph{proposal}. Timescale separation is a \textbf{producer strategy}:
everything below lives in the untrusted lane and in certificate
\emph{structure}; the checker never learns what ``stiff'' means, and
the separation parameter $\varepsilon$ enters certificates only ever as
an exact rational coefficient in rows, never as a concept --- no new
rule family, no new trust surface. Four sub-moves, run in order:

\begin{enumerate}
\item \textbf{Diagnose} (cheap, float): decide whether you are in the
  situation at all (recognition signs below). The diagnosis routes; it
  is never a claim.
\item \textbf{Factor} --- reparameterize into fast/slow coordinates,
  as rungs of the toolkit's reparameterization ladder. The first rung
  (shipped) is eigenmode coordinates: the modal matrix of the
  linearization is float-proposed and dyadic-snapped --- the
  proposer/disposer pattern again --- with the declared inverse
  verified by exact multiplication and the pullback identities
  replayed, so the frame travels with the certificate as an interface
  hypothesis whose identities the checker replays. The eigenvector
  matrix is the standard fast/slow splitter, and it costs nothing to
  de-trust. Higher rungs cover patch-local slow sets (per-patch
  affine; chart maps with certified-positive denominators) up to the
  ladder's ceiling --- hypernet maps for genuinely nonlinear slow
  manifolds, where the reparameterization is itself certificate
  content cited by digest. The ceiling is frontier at this preview:
  named, wire-format slot reserved, semantics awaiting their first
  consumer.
\item \textbf{Witness} --- make ``separated'' a \emph{checkable
  claim}, not an assumption. Two readings, used per what your matrix
  admits. The quadratic-form/Lyapunov reading (shipped): a fast-block
  rate floor $\alpha_f$ via positive-semidefiniteness rows
  ($P_f \succ 0$, $Q_f - \alpha_f P_f \succ 0$), a slow-block rate
  ceiling, and the \textbf{gap row} $\alpha_f \ge 2 r \alpha_s$ as one
  exact comparison --- and this reading works for nonsymmetric
  linearizations, which is the generic case. The eigenvalue reading
  (proposed, not shipped): certified spectral windows and an inertia
  rule --- eigenvalue \emph{counting} across shifts, the one
  conclusion shape the shipped roster does not span. When those ship,
  ``the spectrum splits into
  $(-\infty, -1/\varepsilon] \cup [-c, 0)$'' becomes a specific
  certified sentence.
\item \textbf{Compose} --- certify the blocks separately, then buy the
  coupled conclusion with the classical two-timescale composite
  $W = w_s^{\mathsf T} P_s w_s + \lambda\, w_f^{\mathsf T} P_f w_f$
  (Khalil, \emph{Nonlinear Systems} \cite{khalil2002nonlinear},
  ch.~11, in linear-quadratic form). For polynomial and linear blocks the composite's rows are
  ordinary conclusions on the full coupled system --- the weighted sum
  \emph{is} the certificate, and the weight $\lambda$ is the
  producer's small-gain balancing act, float-searched and exactly
  disposed. The general rule family is in the catalogue:
  \rulename{SmallGain} (chapter~\ref{ch:qtsl-manual}), the calculus's
  feedback-composition rule, which elaborates to a single
  \rulename{Syzygy} node whose exact conclusion identity at weights
  $(1, \lambda, 1, 1)$ makes the cross-supplies cancel by construction
  --- nothing fixpoint-shaped survives to checking time. Timescale
  separation is the special case in which the gain asymmetry makes the
  small-gain product easy: the fast block's gain \emph{to} the slow
  one scales like $1/\alpha_f$, so a real gap buys the composition its
  slack. And carry \textbf{the full error term}: whatever the
  slow-manifold or modal-frame approximation leaves over is a
  \emph{certified remainder} inside the composite rows, never a
  dropped term. In the linear case this remainder is literally visible
  --- the exact off-diagonal entries the dyadic snap leaves in the
  modal frame. In the nonlinear case it is a certified enclosure of
  the approximation error: Taylor models today, with a
  fancier-approximants ladder (Chebyshev models
  \cite{ensl-00472509v2}, Remez with a-posteriori error
  certification, Pad\'e with certified-positive denominators) in
  design at this preview.
\end{enumerate}

\paragraph{When it applies.} Recognition signs, each a cheap
float-lane probe:
\begin{itemize}
\item \textbf{Eigenvalue-ratio probe}: eigenvalues of the
  linearization at a few representative points; a ratio
  ${\gtrsim}10\times$ between rate bands is the signature (one
  \texttt{eigvals} call).
\item \textbf{Enclosure anisotropy}: cover refinement concentrating
  splits along particular directions everywhere --- cells forced small
  in two axes at once by a skewed feature is exactly what a
  reparameterization fixes.
\item \textbf{Step-size collapse}: an explicit integrator's stable
  step is orders of magnitude below the timescale of the observables
  you care about --- classical stiffness, and the reason the
  certificate's cover would otherwise inherit the fast timescale
  globally.
\item \textbf{Search-space conditioning}: certificate searches
  starving because coordinates live at wildly different scales. A
  measured instance from the workshop record: an Adam-driven
  certificate search starving on value ranges spanning $1$ to $10^6$,
  fixed by a coordinate-scaling knob --- the same pathology one level
  up, in the proposal lane's own dynamics.
\item \textbf{The domain says so}: filters, controllers, chemical
  networks, neural dynamics --- fast transient plus slow drift is the
  modal architecture of engineered and biological systems both. The
  workshop's worn witness is a deployed interactive-tracking system
  whose fast loop runs at 72--120\,Hz against a 0.03--1\,Hz slow
  loop, three orders of magnitude apart, named in its own
  documentation as the decision everything else rests on.
\end{itemize}

\paragraph{Worked example from the record.} A deliberately linear toy,
taken end to end through all four moves on shipped machinery: a
two-fast/one-slow coupled linear system --- slow rate ${\sim}1$, fast
rates ${\sim}64, 80$, all-dyadic entries --- exact throughout:

\begin{itemize}
\item float \texttt{eig} proposes the modal matrix; dyadic snap at
  $2^{-12}$; exact inverse declared; the checker replays
  $T \cdot T^{-1} = I$ and $\tilde A = T^{-1} A T$ in $\Q$;
\item the snap's leftover off-diagonal coupling
  (${\sim}7 \times 10^{-3}$, exact) is the certified remainder,
  carried inside the composite;
\item fast rate floor $\alpha_f = 117097/1024$ through the
  \emph{unmodified} checker, slow rate $\alpha_s$ exact (a
  $1 \times 1$ block's rate is one number), gap row
  $\alpha_f \ge 2 \cdot 32 \cdot \alpha_s$ --- certified per-state
  ratio $57.9\times$;
\item the two-timescale composite with $\lambda = 8$ from a 12-point
  float search: composite rate $\alpha_c = 455/256 \approx 90\%$ of
  the slow quadratic bound --- the slow mode governs, which is the
  separation story's signature;
\item 0.81\,s end to end, a 554-byte certificate, 3/3 producer refusal
  probes and 8/8 checker mutants refused. The probe worth naming:
  cranking the back-coupling flips the slow mode \emph{unstable} (the
  Schur-complement correction crosses zero) and the pipeline refuses
  at the exact modal row --- the classical quasi-steady-state trap,
  caught by arithmetic rather than by vigilance.
\end{itemize}

The toy is deliberately linear --- every claim rests on rows that
exist at this preview. What the linear case cannot exhibit, and the
frontier behind it: hypernet slow-manifold reparameterization
(frontier), the eigenvalue-reading witness (proposed), and nonlinear
certified remainders (in design). Each item names work in flight.

\paragraph{The failure mode it prevents.}
\begin{itemize}
\item \textbf{Paying the stiffness tax in the cover}: without the
  factoring, a certificate over the coupled system inherits the fast
  timescale everywhere --- cell budgets explode along directions the
  slow claim never needed. The reparameterization moves the tax into
  one exact similarity replay.
\item \textbf{The quasi-steady-state trap}: classical
  singular-perturbation practice \emph{assumes} the fast subsystem
  settles and substitutes its equilibrium --- sound only under
  conditions (attractivity, weak back-coupling) that the classical
  workflow rarely certifies. The witness sub-move makes exactly those
  conditions the checkable content; the toy's refusal probe is the
  trap sprung on purpose.
\item \textbf{Monolithic non-reuse}: one big Lyapunov certificate for
  the coupled system is re-derived from scratch whenever either block
  changes. The composite factors the work: the fast block's
  certificate is a reusable component, and only the small-gain
  arithmetic re-runs when the slow side varies --- the same economy
  the assume--guarantee layer exists to buy
  (chapter~\ref{ch:qtsl-composition-patterns}), arriving here through
  the gap.
\end{itemize}

\paragraph{Exercises.}
\begin{enumerate}
\item \textbf{Land the eigenvalue reading}: the inertia and
  spectral-window rules are scoped with priced proof obligations in
  the workshop record. When they ship, upgrade the toy's gap row to a
  certified spectral-window sentence and re-run the mutant battery.
\item \textbf{Certify a worn witness}: take a production system you
  operate that already runs a fast/slow split as engineering practice
  (the tracking-filter shape above). State the two loops as a
  two-mechanism system, witness the rate gap, and compose --- the
  first deployed-system instance of this section.
\item \textbf{Van der Pol at large $\mu$}: the classic nonlinear
  two-timescale system. The slow manifold is a cubic --- mid-ladder
  territory, with the approximants ladder's certified remainder as the
  error term once it ships. Success condition: a certified
  relaxation-oscillation enclosure whose cover is coarse along the
  slow manifold.
\item \textbf{Timescale $\times$ parameter box}: compose this section
  with the campaign's parameter-boxed oscillator artefact (uniform
  decrease over a damping box) --- a fast/slow system whose
  \emph{rates} live in boxes, certifying separation uniformly over the
  box. One new fixture, no new rules.
\end{enumerate}

\section{Anchor at the Known Point}
\label{sec:ckbk-anchor}

\paragraph{The move.} When you must certify something about an object
you cannot compute --- a fixed point, a limit, an optimum --- do not
analyse the iteration that approaches it. \textbf{Invert the
anchoring: state every hypothesis at the concrete artefact you hold}
(the snapshot, the candidate, the last iterate), and let one
contraction-style row transport the conclusion to the unknown object.
Concretely, for a fixed point $\mu_\star$ of an operator $T$ and a
snapshot $\hat\mu$:

\begin{enumerate}
\item certify a \textbf{two-point contraction} on a ball around
  $\hat\mu$: $d(Tx, Ty) \le \rho\, d(x, y)$ with $\rho < 1$, as
  checkable rows;
\item certify the \textbf{one-step residual} at the snapshot:
  $d(T\hat\mu, \hat\mu) \le \delta$ --- one application of $T$ to a
  concrete object, always computable;
\item conclude $d(\hat\mu, \mu_\star) \le \delta/(1-\rho)$, with
  existence \emph{and} uniqueness of $\mu_\star$ in the ball
  manufactured by the contraction itself (Banach) --- no convergence
  analysis, no iteration count, no history of how $\hat\mu$ was
  produced.
\end{enumerate}

The inversion matters because the literature's theorems are usually
anchored at the unknown point (``the iteration contracts \emph{toward}
$\mu_\star$'' --- a one-point statement at a place you cannot evaluate
anything), which is useless as a certificate hypothesis. A
certificate's hypotheses must be checkable, and only the known point
supports checking. The two-point form costs more analysis once, in the
untrusted lane; every certificate after that is arithmetic at
$\hat\mu$. This section is the producer half of a seam: the checker
half is the catalogue's \rulename{QFixBall} rule
(chapter~\ref{ch:qtsl-manual}), whose hypotheses are anchored at the
known snapshot by construction and whose refusals include
$\rho \ge 1$.

The producer-side companions, from the same worked example:

\begin{itemize}
\item \textbf{Reduce before you certify.} If the operator provably
  maps everything into a small family (a pinning lemma: any output of
  $T$ is Gaussian with known precision, whatever the input), the
  measure-space problem collapses to a finite-dimensional one --- and
  the reduction is \emph{forced}, not assumed, so the certificate
  still speaks about the full space.
\item \textbf{Guess-and-check the spectral row.} The contraction
  factor $\rho$ needs a norm bound on a nonnegative matrix; propose a
  Perron row vector with floats, snap to rationals, and let the
  checker verify the row inequality. When bare power iteration
  oscillates (period-2 on a two-block system), iterate $(R+I)/2$
  instead --- same fixed vector, damped rotation. Both proposals are
  sound to check; the fix buys \emph{sharpness}, not soundness --- the
  checker is indifferent, and that indifference is the seam working.
\end{itemize}

\paragraph{When it applies.}
\begin{itemize}
\item the claim is ``the true object is within $\varepsilon$ of
  \emph{this}'' --- distance-to-fixpoint, distance-to-optimum, a
  posteriori error bounds --- rather than ``the algorithm
  converges'';
\item every hypothesis you have tried to state mentions the unknown
  object's location (a Lipschitz constant ``near $\mu_\star$,'' a
  curvature ``at the optimum'') --- the tell that the literature's
  anchoring is backwards for your purpose;
\item one application of the operator to a concrete object is cheap,
  even though finding its fixed point is not (one coordinate-ascent
  sweep versus the converged posterior; one Newton step versus the
  root; one Bellman backup versus the value function);
\item a self-mapped ball is available: contraction plus residual imply
  $T(\overline{B}(\hat\mu, r)) \subseteq \overline{B}(\hat\mu, r)$
  for $r \ge \delta/(1-\rho)$, so the hypotheses only need to hold on
  a set you can name in advance.
\end{itemize}

Interval-Newton and Krawczyk verification are this move's classical
ancestors in $\R^n$; the move generalizes them to any complete metric
space a reduction can reach.

\paragraph{Worked example from the record.} Certified Wasserstein
balls for variational inference. The task: certify that a CAVI
(coordinate-ascent variational inference) snapshot is within
$\mathbb{W}_2$-distance $r$ of the true self-consistent fixed point.

\begin{itemize}
\item The literature's theorems (contraction of the CAVI iteration
  \cite{caprio26}) are one-point statements at the unknown
  $\mu_\star$. The certificate
  inverts: all hypotheses at the snapshot $\hat\mu$.
\item The pinning reduction: every CAVI output is exactly Gaussian
  with precision $Q_{ii}$ regardless of input, so the fixed-point
  problem in measure space collapses to an affine map on means --- and
  the certificate quantifies over \emph{all} product measures in
  $\mathcal{P}_2$, not just Gaussian candidates.
\item Receipts: ACCEPT with certified $\rho = 96/181$ against true
  contraction factor $3/(4\sqrt{2}) \approx 0.53033$; ball radius
  $3.6 \times 10^{-7}$; ground truth (computed independently, never
  consumed by the certifier) inside every ball with slack
  ${\times}1.9$; 11/11 mutants refused at named obligations.
\item The period-2 Perron oscillation appeared live: the first run's
  bare power iteration proposed a loose $\rho = 3/4$; the $(R+I)/2$
  fix recovered sharpness. Both fixture generations were accepted by
  the checker --- a live specimen of the proposer being fallible and
  the trust boundary not caring
  (section~\ref{sec:ckbk-proposer-disposer}, observed end to end).
\end{itemize}

\paragraph{The failure mode it prevents.} Without the inversion, the
certificate route dead-ends in one of two ways. Either the producer
imports the literature theorem as-is and the certificate's hypotheses
quantify over a neighbourhood of a point nobody can locate ---
unfalsifiable rows, which a checker must either trust (trust-surface
creep) or refuse (no certificate at all). Or the producer certifies
the \emph{iteration history} (``$n$ steps, each contracting'') ---
which couples the certificate to the trajectory that produced the
snapshot, breaks on restarts and warm starts, and still needs the
contraction constant somewhere. The anchored form is history-free: any
snapshot from any source, certified on its own merits --- a corrupted
snapshot simply fails the residual row (refusal, not a lie).

\paragraph{Exercises.}
\begin{enumerate}
\item \textbf{Block-Gaussian instance} --- the worked example's blocks
  were scalars; run the same pattern with matrix blocks, exercising
  operator-norm rows proper. Everything is staged; only the fixture
  generator changes.
\item \textbf{Discrete instance (2-spin Ising)} --- the
  total-variation contraction is free at bounded couplings; this
  produces the first \emph{discrete} Wasserstein-ball fixture and
  tests whether the move survives leaving the Gaussian reduction.
\item \textbf{The covariance fold-in} --- fold covariance error into
  the ball (dropping the variance-exactness refusal); the disanalogy
  to price is that the mean-map affinity no longer collapses the whole
  problem.
\item \textbf{The transfer question}, worth a page before any build:
  which of your own stack's implicit objects (stationary laws of
  samplers, Bellman fixed points in value-iteration producers) already
  have a cheap one-step residual and are one two-point contraction row
  away from an anchored ball?
\end{enumerate}

\section{Safety Walls and Liveness Tilts}
\label{sec:ckbk-polarity}

\paragraph{The move.} A specification that asks a system to be
\emph{robustly safe while optimistically achieving a goal} has two
halves of opposite character, and most threshold-reading errors in
probabilistic specification come from applying one half's natural
reading to the other. The move is to give the two halves opposite
\textbf{polarities}, explicitly:

\begin{enumerate}
\item \textbf{Safety is a probability claim, read violation-form: a
  wall.} Budget the mass on the bad event ---
  $P(\mathrm{bad}) \le \beta$ --- and hold violation infinitely
  inconsistent: the belief (the level-of-inconsistency functional of
  chapter~\ref{ch:qtsl-semantic-universe}) is $0$ on the half-space
  where the budget holds and $\infty$ where mass on the bad event
  exceeds it, with nothing in between. In particular, the run that
  halts, idles, or leaks mass incurs no inconsistency --- \emph{a
  system that stops has not hit the wall}.
  This is not leniency but correctness: safety is
  closure-invariant --- adding stopping behaviours to a safe system
  cannot make it unsafe, and leaked mass can only lower the
  probability of reaching the bad event
  (chapter~\ref{ch:qtsl-semantic-universe} carries this asymmetry as
  a theorem of the partial-model semantics).
\item \textbf{Liveness is an expectation claim, read as a tilt: a
  hinge.} Price the shortfall of the good event's mass below its
  target $\tau$ at a finite rate ---
  $\lambda \cdot \bigl(\tau - \mu(\mathrm{goal})\bigr)_+$ --- and
  charge nothing at or above the target. Beliefs take values in
  $[0, \infty]$ (the Lawvere quantale: there are no negative levels
  of inconsistency), so the shortfall is a saturated difference, written
  $(\cdot)_+$, never a signed one. The hinge's two sides carry the
  reading: the flat side means overachievement is free and buys
  nothing --- in particular it cannot be traded against the wall ---
  and the sloped side means a near-miss is \emph{priced}, not merely
  refused, which is what ``optimistically'' means as mathematics.
\end{enumerate}

The two readings are one family, not two species: as
$\lambda \to \infty$ the hinge steepens around its kink at
$\mu(\mathrm{goal}) = \tau$ and converges --- in the epigraph sense,
the right sense for beliefs --- to the hard threshold: $0$ at or
above the target, $\infty$ below. A \emph{hard} liveness threshold is
therefore the $\lambda = \infty$ endpoint of the tilt,
satisfaction-form as it must be: the run that does nothing fails it,
and should --- that is what a hard liveness threshold \emph{means}. Choosing
$\lambda$ is choosing where on the family the claim lives; the trap
(below) is only ever reading a threshold at the wrong polarity, never
choosing a finite $\lambda$ on purpose.

Compose by maximum: $\max(\mathrm{wall}, \mathrm{hinge})$ is the
conjunction of beliefs, and it is \textbf{convex} --- the wall is
the $\{0, \infty\}$-indicator of a half-space, the hinge is a maximum
of an affine function with zero, and a pointwise maximum of convex
beliefs is convex. Certification inherits the cheapness: half-space
walls, \rulename{Farkas} combination, and the ReLU row family
(\rulename{ReLU-Pos}, \rulename{ReLU-Above} and kin,
chapter~\ref{ch:qtsl-manual}) --- whose native object is exactly
$\max(\mathrm{affine}, 0)$ --- already span the composite. The
payoff sentence: the standard reach--avoid specification, in its
robustly-safe-while-optimistically-live form, costs nothing the
toolkit does not already ship.

\paragraph{When it applies.}
\begin{itemize}
\item Any specification with a ``never / at most'' half and an
  ``eventually / at least'' half: reach--avoid, stabilize within
  bounds, achieve a return under a resource budget. The two halves
  are both thresholds, and they need opposite readings.
\item Whenever the modelled system can halt, diverge, exit scope, or
  otherwise leak mass --- partiality --- so that the two readings of
  a threshold genuinely differ. The recognition question is: ``what
  does this specification say about the run that does
  \emph{nothing}?'' If safety and liveness give the same answer,
  a polarity has been misassigned.
\item Whenever you meet a bare threshold --- yours or someone else's
  --- that does not carry its reading. The natural default flips sign
  between the two halves, so ``which polarity did the author mean?''
  is a question to settle before consuming the row, not after.
\item When composition or certification cost binds: the wall-plus-tilt
  composite stays convex, so the polarity discipline is also what
  keeps the composed specification in the cheap, Farkas-certifiable
  class.
\end{itemize}

\paragraph{Worked example.} Three outcomes suffice: \textsc{Reach}
(the goal), \textsc{Avoid} (the bad set), \textsc{Inertia} (the run
that never commits --- halts, idles, diverges). A distribution
$\mu = (\mu_R, \mu_A, \mu_I)$ over the three is a point of the
probability $2$-simplex, and the specification reads: \emph{reach with
probability at least $\tfrac34$, never more than $\tfrac14$ on the
bad set.} The safety wall is the violation half-space
$\mu_A > \tfrac14$ penalized at $\infty$; the liveness tilt is the
hinge $\lambda \cdot (\tfrac34 - \mu_R)_+$. Four distributions ---
the three vertices and one genuine mixture --- against the wall, the
hinge at two rates, and the hinge's hard limit:

\begin{center}
\begin{tabular}{@{}lcrrrr@{}}
\toprule
& $(\mu_R, \mu_A, \mu_I)$
& \multicolumn{1}{c}{wall}
& \multicolumn{1}{c}{$\lambda = 4$}
& \multicolumn{1}{c}{$\lambda = 40$}
& \multicolumn{1}{c}{$\lambda = \infty$} \\
\midrule
commits to the goal    & $(1, 0, 0)$ & $0$ & $0$ & $0$ & $0$ \\
commits to the bad set & $(0, 1, 0)$ & $\infty$ & $3$ & $30$ & $\infty$ \\
never commits          & $(0, 0, 1)$ & $0$ & $3$ & $30$ & $\infty$ \\
the mixed run          & $\bigl(\tfrac{7}{10}, \tfrac1{20}, \tfrac14\bigr)$
                       & $0$ & $\tfrac15$ & $2$ & $\infty$ \\
\bottomrule
\end{tabular}
\end{center}

\noindent Read the columns against the doctrine. The \textbf{wall}
distinguishes the two failures the hinge cannot: committing to the
bad set hits it, never committing does not --- failing dangerously
and failing idly are different verdicts, and only the
violation-form reading can tell them apart. The \textbf{hinge}
carries the orthogonal information: the mixed run, with
$\mu_R = \tfrac{7}{10}$ against the target $\tfrac34$, certifies at
\emph{every} finite $\lambda$, paying the priced shortfall
$\lambda \cdot \tfrac1{20}$ --- the optimistic reading, a near-miss
ranked rather than refused --- while any distribution at or above
the target pays exactly $0$ on the flat side. And the
\textbf{$\lambda = \infty$ column} is the hinge steepened into the
step: the mixed run and \textsc{Inertia} now both fail, which is
correct precisely when the claim intended is the hard one. The
composite $\max(\mathrm{wall}, \mathrm{hinge})$ needs both columns'
information: drop the wall and dangerous failure is merely priced
(a large enough achievement elsewhere could dominate it); drop the
hinge and every near-miss collapses into undifferentiated refusal.

Where this stands in the toolkit at this preview. The
probability-threshold former \textsc{SProbGE}
(chapter~\ref{ch:qtsl-specification-language}) is a mass floor --- lower
probability at least $p$ --- i.e.\ the $\lambda = \infty$ liveness
endpoint, and the credal semantics makes it feel every lost unit of
mass, exactly as a hard liveness claim should
(chapter~\ref{ch:qtsl-semantic-universe}: a liveness claim cannot be
laundered through a model that is allowed to quietly stop). Safety
already runs at the other polarity internally: the safety batteries' certified
object is an upper bound on the violation mass
($\overline{P}(\mathrm{bad}) \le V(x_0)$, the threshold row of
\rulename{QSafetyBattery} and kin), which is why those certificates
remain sound under leak --- the wall reading is already how the
machinery works. The expectation formers
(\textsc{SReturnGE}/\textsc{SReturnLE},
chapter~\ref{ch:qtsl-specification-language}) are the shipped
expectation-claim species that carry the tilt: a hinge belief over an
indicator channel is an expected-shortfall bound. And the polarity
itself is declared in the grammar --- as a sort, which is stronger
than a field: the specification cap
(chapter~\ref{ch:qtsl-specification-language}) types every threshold
component by its face, so a mis-faced premise is a formation error,
refused before any semantics runs, and each former's reading is
definitional in the certificate theory rather than a habit of its
author. One disambiguation earns its keep here, because that chapter
also says ``safety face'' and means a different axis by it. The cap's
polarity answers \emph{who quantifies over the admitted laws}: a
$\mathsf{ForallSpec}$ component binds every law, an
$\mathsf{ExistsSpec}$ component asserts some law attains. This
section's polarity answers \emph{which way a single threshold faces}:
wall on the violation mass, floor on the goal mass. The two axes are
independent, and one specimen shows how far apart they sit:
\textsc{SProbGE} --- this section's hard \emph{liveness} endpoint ---
concludes $\mathsf{ForallSpec}$, the credal \emph{safety} face,
because demanding the goal's mass from every admitted law is a
guarantee about a floor. What remains the author's is this section's
own discipline: the grammar refuses a mis-faced component, and no
grammar can refuse a wall aimed at the wrong event.

\paragraph{The failure mode it prevents.} One trap, two signs, plus
the composite's own failure:
\begin{itemize}
\item \textbf{Liveness read violation-form forgives doing nothing.}
  Two outcomes; a system that succeeds with probability $\tfrac25$
  and otherwise halts; the claim ``reach success with probability at
  least $\tfrac12$.'' The violation-form misreading --- refuse only
  if mass on the \emph{failure} event exceeds $\tfrac12$ --- sees
  zero mass on failure and passes, but $\mu(\mathrm{success}) =
  \tfrac25 < \tfrac12$: the mass floor fails, and the correct
  verdict is refusal. Liveness laundered through leak --- the
  quietly-stopping system certified for a goal it does not deliver.
\item \textbf{Safety read satisfaction-form punishes halting as if it
  were violation.} Spell ``avoid the bad set'' as a floor on the
  safe-completion event and the safely-idling system fails, though it
  never approaches the bad set. The downstream rot is worse than the
  false refusal: an author who knows the system may stop weakens the
  floor to tolerate it, and the weakened floor now also tolerates
  genuine violation mass. The repair was never a looser constant; it
  was the other polarity.
\item \textbf{The negotiable wall.} Price safety at a finite rate ---
  treat it as one more tilt --- and conjunction-by-maximum degrades
  into trade: enough achievement elsewhere dominates a bounded
  safety belief, and the optimizer will find that trade.
  ``Robustly'' means the wall is not for sale at any exchange rate,
  and only the violation form at infinite inconsistency says so.
\end{itemize}
Two polarities, no trap: liveness formers are satisfaction-form by
definition, hard at $\lambda = \infty$ and tilted at finite $\lambda$;
the safety machinery is violation-form by construction; and the
reading travels with the claim as data instead of living in the
reader's habits.

\paragraph{Exercises.}
\begin{enumerate}
\item \textbf{The face mutants, now that the sort has shipped.} The
  cap's polarity discipline --- a mis-faced premise is a formation
  error --- is a refusal, and this chapter's own doctrine applies to
  it (section~\ref{sec:ckbk-mutant-battery}): a refusal without a
  firing witness is indistinguishable from dead code. Build the
  formation battery, each mis-faced construction fed to the grammar
  with its refusal location asserted --- and include the one cell no
  formation check can own: the laundering-through-leak specimen
  above, a floor swapped for a wall on the complementary event, which
  is well-formed at both faces and wrong only about what the author
  meant. That last mutant's home is the semantic tier (the flipped
  reading asserted against the partial-model semantics and refused
  there), and placing it correctly is part of the exercise.
\item \textbf{The epi-limit, certified.} The one-family claim ---
  the hinge at rate $\lambda$ converges in the epigraph sense to the
  hard threshold as $\lambda \to \infty$ --- is stated here as
  belief mathematics. Formalize it against the belief semantics,
  so that a hard liveness former and the $\lambda = \infty$ tilt are
  interchangeable obligations in one worked certificate rather than
  two spellings the reader must trust to agree.
\item \textbf{The composite on a shipped witness.} The spec
  catalogue's hitting-shape witness (goal reached by deadline with
  probability at least $\tfrac34$) is a hard liveness claim beside a
  safety wall. Restate it composite-form --- wall on the bad set,
  hinge on the goal mass at two finite rates --- and produce the
  certificate with half-space walls, \rulename{Farkas}, and ReLU
  rows only; then check that the $\lambda = \infty$ run of the same
  pipeline recovers the original hard claim. Success condition: the
  composite certifies with no row species beyond the shipped roster,
  and the two finite-$\lambda$ certificates price the same shortfall
  the table above prices.
\end{enumerate}

\appendix
\part{Appendices}
\chapter{The Lean Cone}
\label{app:lean-cone}


Throughout this book, claims that lean on a mechanized result cite
the result by its Lean name, set in typewriter face. This appendix
collects those citations into one place: for every cited declaration,
it reproduces the declaration's statement from the \path{lean/}
trees of \path{maitria-coalition/qtsl}, with its file and line. We call this set the \emph{Lean cone} of the book: the set
of mechanized results the book's prose binds to, closed under
nothing --- what you see is what the text actually cites, re-derived
from the sources on every build, so the listing cannot drift from
the tree it describes.

A word on how to read the entries. Statements are reproduced at
signature level: a theorem up to its proof, a definition up to its
body, a structure up to its fields (an ellipsis marks the elision),
and long source lines re-wrapped to the measure of the page.
Each entry is shown in the namespace where it lives, and surrounding
context (\texttt{open} and \texttt{variable} declarations,
notation) lives in the cited file, so an entry here reads as a
faithful excerpt rather than a self-contained module. To verify an
entry, clone the repository and run \texttt{lake build} inside the
relevant Lake package under \path{lean/} (for example
\path{lean/qtsl} or \path{lean/su-functor}); the toolchain is pinned
by each package's \path{lean-toolchain} file, and a build that
succeeds has kernel-checked every statement listed here, since none
of the trees admits \texttt{sorry}.

Two boundaries of the cone, so that its edges are as legible as its
contents. First, the Rule Catalogue displays a Lean or HOL4
statement beside each rule already, extracted from the mechanized
rule sources and hash-bound at emission; this appendix lists the
catalogue's citations that resolve into the \path{lean/} trees of
this repository and defers to the catalogue's own apparatus for its
rule-source and HOL4 layers. Second, module citations (a dotted
name denoting a whole file or directory of results) are listed per
chapter without statements; the module's file is the reference.

The listing follows the book's chapter order, and each declaration
appears once, under the first chapter that cites it; later chapters
that cite the same result say so in a closing line rather than
repeating the statement.



\section{The Semantic Universe}
Citations from Chapter~\ref{ch:qtsl-semantic-universe}.

\par\medskip\noindent{\small \texttt{Qtsl.DiamondSite.}\texttt{\bfseries ChronologyKit}}\hfill{\footnotesize \path{lean/qtsl-aux/QtslAux/DiamondSite.lean}, line 79}\par\nopagebreak
\begin{leansrc}
structure ChronologyKit (X : Type*) [PartialOrder X] where …
\end{leansrc}
\par\medskip\noindent{\small \texttt{Qtsl.DiamondSite.}\texttt{\bfseries minkowskiKit}}\hfill{\footnotesize \path{lean/qtsl-aux/QtslAux/DiamondTopology.lean}, line 283}\par\nopagebreak
\begin{leansrc}
def minkowskiKit : ChronologyKit (Mink E) where …
\end{leansrc}
\par\medskip\noindent{\small \texttt{Qtsl.DiamondSite.}\texttt{\bfseries diamondTopology}}\hfill{\footnotesize \path{lean/qtsl-aux/QtslAux/DiamondTopology.lean}, line 114}\par\nopagebreak
\begin{leansrc}
def diamondTopology (K : ChronologyKit X) :
    GrothendieckTopology (Diamond X) where …
\end{leansrc}
\par\medskip\noindent{\small \texttt{Qtsl.DiamondSite.}\texttt{\bfseries mem\_of\_not\_thick}}\hfill{\footnotesize \path{lean/qtsl-aux/QtslAux/DiamondTopology.lean}, line 142}\par\nopagebreak
\begin{leansrc}
theorem mem_of_not_thick {D : Diamond X} (h : ¬ Thick K D)
    (S : Sieve D) : S ∈ diamondTopology K D
\end{leansrc}
\par\medskip\noindent{\small \texttt{Qtsl.DiamondSite.}\texttt{\bfseries mem\_diamondTopology\_degenerate\_iff}}\hfill{\footnotesize \path{lean/qtsl-aux/QtslAux/DiamondTopology.lean}, line 158}\par\nopagebreak
\begin{leansrc}
theorem mem_diamondTopology_degenerate_iff
    (hK : ∀ a b : X, K.chron a b ↔ a ≤ b) {D : Diamond X} {S : Sieve D} :
    S ∈ diamondTopology K D ↔ S = ⊤
\end{leansrc}
\par\medskip\noindent{\small \texttt{Qtsl.DiamondSite.}\texttt{\bfseries int\_lt\_no\_chronologyKit}}\hfill{\footnotesize \path{lean/qtsl-aux/QtslAux/DiamondTopology.lean}, line 383}\par\nopagebreak
\begin{leansrc}
theorem int_lt_no_chronologyKit :
    ¬ ∃ K : ChronologyKit ℤ, ∀ a b : ℤ, K.chron a b ↔ a < b
\end{leansrc}
\par\medskip\noindent{\small \texttt{SuFunctor.Site.TTT.}\texttt{\bfseries anchorQuot}}\hfill{\footnotesize \path{lean/su-functor/SuFunctor/Site/TemporalTypeTheory.lean}, line 156}\par\nopagebreak
\begin{leansrc}
def anchorQuot : Window ⥤ IntervalBase where …
\end{leansrc}
\par\medskip\noindent{\small \texttt{SuFunctor.Site.TTT.}\texttt{\bfseries anchorTrajIso}}\hfill{\footnotesize \path{lean/su-functor/SuFunctor/Site/TemporalTypeTheory.lean}, line 323}\par\nopagebreak
\begin{leansrc}
def anchorTrajIso : anchorQuot.op ⋙ ssTraj X ≅ windowTraj X
\end{leansrc}
\par\medskip\noindent{\small \texttt{SuFunctor.Site.}\texttt{\bfseries ttObject}}\hfill{\footnotesize \path{lean/su-functor/SuFunctor/Site/BehaviorSpec.lean}, line 246}\par\nopagebreak
\begin{leansrc}
def ttObject (U : ι → Set T) {A : WireSpan El} (B : WBehavior U (states A)) :
    SiteSpecTotal U El
\end{leansrc}
\par\medskip\noindent{\small \texttt{SuFunctor.Site.}\texttt{\bfseries satAt\_specOfBehavior\_iff}}\hfill{\footnotesize \path{lean/su-functor/SuFunctor/Site/BehaviorSpec.lean}, line 208}\par\nopagebreak
\begin{leansrc}
theorem satAt_specOfBehavior_iff (U : ι → Set T) (B : WBehavior U X) (j : ι)
    (ν : Joint (Traj (U j) X)) :
    SatAt ν (specOfBehavior U B j) ↔ (↑ν.wt.support : Set (Traj (U j) X)) ⊆ B.cls j
\end{leansrc}
\par\medskip\noindent{\small \texttt{SuFunctor.Site.}\texttt{\bfseries wSpecMap\_specOfBehavior}}\hfill{\footnotesize \path{lean/su-functor/SuFunctor/Site/BehaviorSpec.lean}, line 256}\par\nopagebreak
\begin{leansrc}
theorem wSpecMap_specOfBehavior (U : ι → Set T) {A B : WireSpan El}
    (φ : A ⟶ B) (Bh : WBehavior U (states A)) :
    wSpecMap U φ (specOfBehaviorW U Bh) =
      specOfBehaviorW U (Bh.pull (res φ))
\end{leansrc}
\par\medskip\noindent{\small \texttt{SuFunctor.Site.TTT.}\texttt{\bfseries mem\_interiorTopology\_iff\_derived}}\hfill{\footnotesize \path{lean/su-functor/SuFunctor/Site/CausalConcordance.lean}, line 110}\par\nopagebreak
\begin{leansrc}
theorem mem_interiorTopology_iff_derived {W : Window} (S : Sieve W) :
    S ∈ interiorTopology W ↔
      S.functorPushforward windowToDiamond ∈
        diamondTopology (linearOrderKit ℝ) W.toDiamond
\end{leansrc}
\par\medskip\noindent{\small \texttt{SuFunctor.Site.TTT.}\texttt{\bfseries wavy\_mem\_iff\_derived}}\hfill{\footnotesize \path{lean/su-functor/SuFunctor/Site/CausalConcordance.lean}, line 133}\par\nopagebreak
\begin{leansrc}
theorem wavy_mem_iff_derived {W : Window} (S : Sieve W) :
    S.functorPushforward anchorQuot ∈ wavyTopology ⟨W.len⟩ ↔
      S.functorPushforward windowToDiamond ∈
        diamondTopology (linearOrderKit ℝ) W.toDiamond
\end{leansrc}
\par\medskip\noindent{\small \texttt{Qtsl.QuantumStates.}\texttt{\bfseries recordShadowQ\_separates\_bell\_mixed}}\hfill{\footnotesize \path{lean/qtsl-aux/QtslAux/QuantumStates.lean}, line 939}\par\nopagebreak
\begin{leansrc}
theorem recordShadowQ_separates_bell_mixed :
    recordShadowQ zzRecord zzRecord_mem_top
        (stateClassMk qubitNet Dstar ⟨bell, bell_isState⟩) ≠
      recordShadowQ zzRecord zzRecord_mem_top
        (stateClassMk qubitNet Dstar ⟨mixed, mixed_isState⟩)
\end{leansrc}
\par\medskip\noindent{\small \texttt{Qtsl.DiamondSite.}\texttt{\bfseries minkowski\_diamond\_spacelike\_pair}}\hfill{\footnotesize \path{lean/qtsl-aux/QtslAux/DiamondTopology.lean}, line 245}\par\nopagebreak
\begin{leansrc}
theorem minkowski_diamond_spacelike_pair {x₀ : E} (hx : x₀ ≠ 0)
    (hx1 : ‖x₀‖ ≤ 1) :
    ∃ p q : Mink E,
      ((⟨0, 0⟩ : Mink E) ≤ p ∧ p ≤ ⟨2, 0⟩) ∧
      ((⟨0, 0⟩ : Mink E) ≤ q ∧ q ≤ ⟨2, 0⟩) ∧
      ¬ p ≤ q ∧ ¬ q ≤ p
\end{leansrc}
\par\medskip\noindent{\small \texttt{SuFunctor.Universe.FellCloseds.}\texttt{\bfseries compactSpace}}\hfill{\footnotesize \path{lean/su-functor/SuFunctor/Universe/FellTopology.lean}, line 268}\par\nopagebreak
\begin{leansrc}
instance compactSpace : CompactSpace (FellCloseds α)
\end{leansrc}
\par\medskip\noindent{\small \texttt{SuFunctor.Universe.}\texttt{\bfseries fell\_eq\_vietoris}}\hfill{\footnotesize \path{lean/su-functor/SuFunctor/Universe/FellTopology.lean}, line 181}\par\nopagebreak
\begin{leansrc}
theorem fell_eq_vietoris [CompactSpace α] [T2Space α] :
    SuFunctor.Universe.fell α = TopologicalSpace.vietoris α
\end{leansrc}
\par\medskip\noindent{\small \texttt{SuFunctor.Universe.fell.}\texttt{\bfseries continuous\_restrictOpen}}\hfill{\footnotesize \path{lean/su-functor/SuFunctor/Universe/FellRestrict.lean}, line 81}\par\nopagebreak
\begin{leansrc}
theorem continuous_restrictOpen (hU : IsOpen U) :
    Continuous (restrictOpen U)
\end{leansrc}
\par\medskip\noindent{\small \texttt{SuFunctor.Universe.fell.}\texttt{\bfseries continuous\_image\_of\_isProperMap}}\hfill{\footnotesize \path{lean/su-functor/SuFunctor/Universe/FellRestrict.lean}, line 130}\par\nopagebreak
\begin{leansrc}
theorem continuous_image_of_isProperMap (hf : IsProperMap f) :
    Continuous (Set.image f : Set α → Set β)
\end{leansrc}
\par\medskip\noindent{\small \texttt{SuFunctor.Universe.FellCloseds.}\texttt{\bfseries projLimHomeomorph}}\hfill{\footnotesize \path{lean/su-functor/SuFunctor/Universe/FellDescent.lean}, line 315}\par\nopagebreak
\begin{leansrc}
def projLimHomeomorph (U : ι → Set α) (hUo : ∀ i, IsOpen (U i))
    (hdir : Directed (· ⊆ ·) U) (hcov : ⋃ i, U i = univ) :
    FellCloseds α ≃ₜ ProjLim U where …
\end{leansrc}
\par\medskip\noindent{\small \texttt{SuFunctor.Universe.FellCloseds.IsLocalLocus.}\texttt{\bfseries isClosed}}\hfill{\footnotesize \path{lean/su-functor/SuFunctor/Universe/FellRecordLocus.lean}, line 99}\par\nopagebreak
\begin{leansrc}
theorem IsLocalLocus.isClosed (hloc : IsLocalLocus U R₀ R)
    (hUo : ∀ i, IsOpen (U i)) (hRcl : ∀ i, IsClosed (R i)) :
    IsClosed R₀
\end{leansrc}
\par\medskip\noindent{\small \texttt{SuFunctor.Universe.FellCloseds.IsLocalLocus.}\texttt{\bfseries isCompact}}\hfill{\footnotesize \path{lean/su-functor/SuFunctor/Universe/FellRecordLocus.lean}, line 113}\par\nopagebreak
\begin{leansrc}
theorem IsLocalLocus.isCompact (hloc : IsLocalLocus U R₀ R)
    (hUo : ∀ i, IsOpen (U i)) (hRcl : ∀ i, IsClosed (R i)) :
    IsCompact R₀
\end{leansrc}
\par\medskip\noindent{\small \texttt{SuFunctor.Universe.FellCloseds.IsLocalLocusEq.}\texttt{\bfseries glue\_mem}}\hfill{\footnotesize \path{lean/su-functor/SuFunctor/Adequacy/SolutionSetMesh.lean}, line 125}\par\nopagebreak
\begin{leansrc}
theorem IsLocalLocusEq.glue_mem {U : ι → Set α} {R₀ : Set (FellCloseds α)}
    {R : ∀ i, Set (FellCloseds ↥(U i))} (hloc : IsLocalLocusEq U R₀ R)
    (hUo : ∀ i, IsOpen (U i)) (hcov : ⋃ i, U i = univ)
    {p : Equalizer U} (hp : ∀ i, p.1 i ∈ R i) :
    glueEqualizer U hUo hcov p ∈ R₀
\end{leansrc}
\par\medskip\noindent{\small \texttt{SuFunctor.Universe.FellCloseds.}\texttt{\bfseries speedLocus}}\hfill{\footnotesize \path{lean/su-functor/SuFunctor/Universe/FellRecordLocus.lean}, line 234}\par\nopagebreak
\begin{leansrc}
def speedLocus : Set (FellCloseds β)
\end{leansrc}
\par\medskip\noindent{\small \texttt{SuFunctor.Universe.FellCloseds.}\texttt{\bfseries isClosed\_speedLocus}}\hfill{\footnotesize \path{lean/su-functor/SuFunctor/Universe/FellRecordLocus.lean}, line 259}\par\nopagebreak
\begin{leansrc}
theorem isClosed_speedLocus (hτ : Continuous τ) (hσ : Continuous σ)
    (c : ℝ≥0) : IsClosed (speedLocus τ σ c)
\end{leansrc}
\par\medskip\noindent{\small \texttt{SuFunctor.Universe.FellCloseds.}\texttt{\bfseries isSpeedGraph\_iff\_forall\_clip}}\hfill{\footnotesize \path{lean/su-functor/SuFunctor/Universe/FellRecordLocus.lean}, line 308}\par\nopagebreak
\begin{leansrc}
theorem isSpeedGraph_iff_forall_clip (hdir : Directed (· ⊆ ·) U)
    (hcov : ⋃ i, U i = univ) (F : FellCloseds β) :
    IsSpeedGraph τ σ c (F : Set β)
      ↔ ∀ i, FellCloseds.clip (U i) F
          ∈ speedLocus (τ ∘ (↑)) (σ ∘ (↑)) c
\end{leansrc}
\par\medskip\noindent{\small \texttt{SuFunctor.Universe.FellCloseds.}\texttt{\bfseries speedLocusHomeomorph}}\hfill{\footnotesize \path{lean/su-functor/SuFunctor/Universe/FellRecordLocus.lean}, line 346}\par\nopagebreak
\begin{leansrc}
noncomputable def speedLocusHomeomorph (hUo : ∀ i, IsOpen (U i))
    (hdir : Directed (· ⊆ ·) U) (hcov : ⋃ i, U i = univ) :
    ↥(speedLocus τ σ c)
      ≃ₜ ↥(FellCloseds.levelLocus U
          fun _ => speedLocus (τ ∘ (↑)) (σ ∘ (↑)) c)
\end{leansrc}
\par\medskip\noindent{\small \texttt{SuFunctor.Universe.FellCloseds.}\texttt{\bfseries bot\_mem\_speedLocus}}\hfill{\footnotesize \path{lean/su-functor/SuFunctor/Universe/FellRecordLocus.lean}, line 242}\par\nopagebreak
\begin{leansrc}
theorem bot_mem_speedLocus : (⊥ : FellCloseds β) ∈ speedLocus τ σ c
\end{leansrc}
\par\medskip\noindent{\small \texttt{SuFunctor.Universe.FellCloseds.}\texttt{\bfseries isLocallyClosed\_speedLocus}}\hfill{\footnotesize \path{lean/su-functor/SuFunctor/Universe/FellRecordLocus.lean}, line 294}\par\nopagebreak
\begin{leansrc}
theorem isLocallyClosed_speedLocus (hτ : Continuous τ)
    (hσ : Continuous σ) (c : ℝ≥0) :
    IsLocallyClosed (speedLocus τ σ c)
\end{leansrc}
\par\medskip\noindent{\small \texttt{SuFunctor.Universe.FellCloseds.}\texttt{\bfseries isSpeedGraph\_iff\_exists\_lipschitzOnWith}}\hfill{\footnotesize \path{lean/su-functor/SuFunctor/Universe/FellRecordLocus.lean}, line 440}\par\nopagebreak
\begin{leansrc}
theorem isSpeedGraph_iff_exists_lipschitzOnWith [CompleteSpace E]
    [Nonempty E] {F : FellCloseds (T × E)} :
    IsSpeedGraph Prod.fst Prod.snd c (F : Set (T × E))
      ↔ ∃ D : Set T, IsClosed D ∧ ∃ f : T → E, LipschitzOnWith c f D ∧
          (F : Set (T × E)) = (fun t => (t, f t)) '' D
\end{leansrc}
\par\medskip\noindent{\small \texttt{SuFunctor.Universe.FellCloseds.}\texttt{\bfseries topology\_onePoint\_eq\_vietoris}}\hfill{\footnotesize \path{lean/su-functor/SuFunctor/Universe/FellOnePoint.lean}, line 319}\par\nopagebreak
\begin{leansrc}
theorem topology_onePoint_eq_vietoris [WeaklyLocallyCompactSpace α] [T2Space α] :
    (FellCloseds.topology : TopologicalSpace (FellCloseds (OnePoint α)))
      = .induced (↑) (TopologicalSpace.vietoris (OnePoint α))
\end{leansrc}
\par\medskip\noindent{\small \texttt{SuFunctor.Universe.FellCloseds.}\texttt{\bfseries nablaGen}}\hfill{\footnotesize \path{lean/su-functor/SuFunctor/Universe/FellCapacity.lean}, line 75}\par\nopagebreak
\begin{leansrc}
def nablaGen (B : ι → Set Y) (hits misses : Finset ι) : Set (FellCloseds Y)
\end{leansrc}
\par\medskip\noindent{\small \texttt{SuFunctor.Universe.FellCloseds.}\texttt{\bfseries isTopologicalBasis\_nablaGen}}\hfill{\footnotesize \path{lean/su-functor/SuFunctor/Universe/FellCapacity.lean}, line 136}\par\nopagebreak
\begin{leansrc}
theorem isTopologicalBasis_nablaGen [LocallyCompactSpace Y] [T2Space Y]
    {B : ι → Set Y} (hB : IsTopologicalBasis (Set.range B))
    (hBcl : ∀ i, IsCompact (closure (B i))) :
    IsTopologicalBasis
      (Set.range fun p : Finset ι × Finset ι => nablaGen B p.1 p.2)
\end{leansrc}
\par\medskip\noindent{\small \texttt{SuFunctor.Universe.FellCloseds.}\texttt{\bfseries measure\_ext\_of\_nablaGen}}\hfill{\footnotesize \path{lean/su-functor/SuFunctor/Universe/FellCapacity.lean}, line 278}\par\nopagebreak
\begin{leansrc}
theorem measure_ext_of_nablaGen [LocallyCompactSpace Y] [T2Space Y] [Countable ι]
    {B : ι → Set Y} (hB : IsTopologicalBasis (Set.range B))
    (hBcl : ∀ i, IsCompact (closure (B i)))
    [MeasurableSpace (FellCloseds Y)] [BorelSpace (FellCloseds Y)]
    {μ ν : Measure (FellCloseds Y)} [IsFiniteMeasure μ]
    (h : ∀ hits misses : Finset ι,
      μ (nablaGen B hits misses) = ν (nablaGen B hits misses)) :
    μ = ν
\end{leansrc}
\par\medskip\noindent{\small \texttt{SuFunctor.Universe.FellCloseds.}\texttt{\bfseries exists\_nablaGen\_measure\_ne}}\hfill{\footnotesize \path{lean/su-functor/SuFunctor/Universe/FellCapacity.lean}, line 302}\par\nopagebreak
\begin{leansrc}
theorem exists_nablaGen_measure_ne [LocallyCompactSpace Y] [T2Space Y] [Countable ι]
    {B : ι → Set Y} (hB : IsTopologicalBasis (Set.range B))
    (hBcl : ∀ i, IsCompact (closure (B i)))
    [MeasurableSpace (FellCloseds Y)] [BorelSpace (FellCloseds Y)]
    {μ ν : Measure (FellCloseds Y)} [IsFiniteMeasure μ]
    (hμν : μ ≠ ν) :
    ∃ hits misses : Finset ι,
      μ (nablaGen B hits misses) ≠ ν (nablaGen B hits misses)
\end{leansrc}
\par\medskip\noindent{\small \texttt{SuFunctor.Universe.FellCloseds.}\texttt{\bfseries scottContinuous\_ext\_of\_nablaGen}}\hfill{\footnotesize \path{lean/su-functor/SuFunctor/Universe/FellCapacity.lean}, line 326}\par\nopagebreak
\begin{leansrc}
theorem scottContinuous_ext_of_nablaGen [LocallyCompactSpace Y] [T2Space Y]
    {B : ι → Set Y} (hB : IsTopologicalBasis (Set.range B))
    (hBcl : ∀ i, IsCompact (closure (B i)))
    {β : Type*} [PartialOrder β] {f g : Opens (FellCloseds Y) → β}
    (hf : ScottContinuous f) (hg : ScottContinuous g)
    (h : ∀ U : Opens (FellCloseds Y),
      (∃ s : Finset (Finset ι × Finset ι),
        (U : Set (FellCloseds Y)) = ⋃ p ∈ s, nablaGen B p.1 p.2) →
      f U = g U) :
    f = g
\end{leansrc}
\par\medskip\noindent{\small \texttt{SuFunctor.Universe.FellCloseds.}\texttt{\bfseries not\_isOpen\_disjoint\_univ}}\hfill{\footnotesize \path{lean/su-functor/SuFunctor/Universe/FellCapacity.lean}, line 375}\par\nopagebreak
\begin{leansrc}
theorem not_isOpen_disjoint_univ [T1Space Y] (hY : ¬CompactSpace Y) :
    ¬IsOpen {F : FellCloseds Y | Disjoint (F : Set Y) univ}
\end{leansrc}
\par\medskip\noindent{\small \texttt{SuFunctor.Universe.FellCloseds.}\texttt{\bfseries exists\_scottContinuous\_eq\_on\_lower\_ne}}\hfill{\footnotesize \path{lean/su-functor/SuFunctor/Universe/FellCapacity.lean}, line 401}\par\nopagebreak
\begin{leansrc}
theorem exists_scottContinuous_eq_on_lower_ne [T1Space Y] [Nonempty Y] :
    ∃ f g : Opens (FellCloseds Y) → ℝ≥0∞,
      ScottContinuous f ∧ ScottContinuous g ∧
      (∀ U : Opens (FellCloseds Y),
        IsOpen[TopologicalSpace.generateFrom
          ((fun V => {F : FellCloseds Y | ((F : Set Y) ∩ V).Nonempty})
            '' {V : Set Y | IsOpen V})]
          (U : Set (FellCloseds Y)) → f U = g U) ∧
      f ≠ g
\end{leansrc}
\par\medskip\noindent{\small \texttt{Pointfree.NablaData.Valuation.}\texttt{\bfseries ext\_of\_nabla}}\hfill{\footnotesize \path{lean/pointfree/Pointfree/Frame/HyperspacePresentation.lean}, line 259}\par\nopagebreak
\begin{leansrc}
theorem ext_of_nabla
    (hμS : μ.IsScottContinuous) (hνS : ν.IsScottContinuous)
    (hfin : ∀ p, μ (nabla geom p) ≠ ∞)
    (h : ∀ p, μ (nabla geom p) = ν (nabla geom p)) : μ = ν
\end{leansrc}
\par\medskip\noindent{\small \texttt{Pointfree.NablaData.Valuation.}\texttt{\bfseries exists\_nabla\_ne}}\hfill{\footnotesize \path{lean/pointfree/Pointfree/Frame/HyperspacePresentation.lean}, line 271}\par\nopagebreak
\begin{leansrc}
theorem exists_nabla_ne
    (hμS : μ.IsScottContinuous) (hνS : ν.IsScottContinuous)
    (hfin : ∀ p, μ (nabla geom p) ≠ ∞)
    (hne : μ ≠ ν) : ∃ p, μ (nabla geom p) ≠ ν (nabla geom p)
\end{leansrc}
\par\medskip\noindent{\small \texttt{Pointfree.Valuation.}\texttt{\bfseries ext\_of\_generators}}\hfill{\footnotesize \path{lean/pointfree/Pointfree/Valuation/Determination.lean}, line 131}\par\nopagebreak
\begin{leansrc}
theorem ext_of_generators
    (hμS : μ.IsScottContinuous) (hνS : ν.IsScottContinuous)
    (hGinf : ∀ i j, ∃ k, G i ⊓ G j = G k)
    (hGgen : ∀ x : L, ∃ S : Set κ, x = ⨆ i ∈ S, G i)
    (hfin : ∀ i, μ (G i) ≠ ∞)
    (h : ∀ i, μ (G i) = ν (G i)) : μ = ν
\end{leansrc}
\par\medskip\noindent{\small \texttt{Pointfree.Valuation.}\texttt{\bfseries ext\_of\_isTopologicalBasis}}\hfill{\footnotesize \path{lean/pointfree/Pointfree/Valuation/SpatialDetermination.lean}, line 60}\par\nopagebreak
\begin{leansrc}
theorem ext_of_isTopologicalBasis
    (hB : IsTopologicalBasis (Set.range fun i => (N i : Set X)))
    (hNinf : ∀ i j, ∃ k, N i ⊓ N j = N k)
    {μ ν : Valuation (Opens X)}
    (hμS : μ.IsScottContinuous) (hνS : ν.IsScottContinuous)
    (hfin : ∀ i, μ (N i) ≠ ∞)
    (h : ∀ i, μ (N i) = ν (N i)) : μ = ν
\end{leansrc}
\par\medskip\noindent{\small \texttt{Pointfree.SubProb.}\texttt{\bfseries promoteHomeomorph}}\hfill{\footnotesize \path{lean/pointfree/Pointfree/Valuation/Topology.lean}, line 230}\par\nopagebreak
\begin{leansrc}
noncomputable def promoteHomeomorph : SubProb L ≃ₜ Prob (L × Prop)
\end{leansrc}
\par\medskip\noindent{\small \texttt{Pointfree.SubProb.}\texttt{\bfseries induced\_promote\_eq}}\hfill{\footnotesize \path{lean/pointfree/Pointfree/Valuation/Topology.lean}, line 171}\par\nopagebreak
\begin{leansrc}
theorem induced_promote_eq :
    TopologicalSpace.induced (promote : SubProb L → Prob (L × Prop))
        (Prob.instTopologicalSpace (L × Prop))
      = SubProb.instTopologicalSpace L
\end{leansrc}
\par\medskip\noindent{\small \texttt{SuFunctor.Universe.General.}\texttt{\bfseries SubProbabilityMeasure}}\hfill{\footnotesize \path{lean/su-functor/SuFunctor/Universe/General/Carrier.lean}, line 51}\par\nopagebreak
\begin{leansrc}
def SubProbabilityMeasure : Type _
\end{leansrc}
\par\medskip\noindent{\small \texttt{Pointfree.SubProbabilityMeasure.}\texttt{\bfseries continuous\_toSubProb}}\hfill{\footnotesize \path{lean/pointfree/Pointfree/Valuation/BridgeContinuity.lean}, line 180}\par\nopagebreak
\begin{leansrc}
theorem SubProbabilityMeasure.continuous_toSubProb :
    Continuous (SubProbabilityMeasure.toSubProb (X := X))
\end{leansrc}
\par\medskip\noindent{\small \texttt{Pointfree.}\texttt{\bfseries daggerChartMeasure}}\hfill{\footnotesize \path{lean/pointfree/Pointfree/Measure/DaggerChart.lean}, line 390}\par\nopagebreak
\begin{leansrc}
noncomputable def daggerChartMeasure :
    SubProbabilityMeasure X ≃ₜ ProbabilityMeasure (X ⊕ Unit) where …
\end{leansrc}
\par\medskip\noindent{\small \texttt{Pointfree.}\texttt{\bfseries compactSpace\_of\_daggerChartMeasure}}\hfill{\footnotesize \path{lean/pointfree/Pointfree/Measure/DaggerChart.lean}, line 415}\par\nopagebreak
\begin{leansrc}
theorem compactSpace_of_daggerChartMeasure [T2Space X] [BorelSpace X]
    [CompactSpace X] : CompactSpace (SubProbabilityMeasure X)
\end{leansrc}
\par\medskip\noindent{\small \texttt{Pointfree.}\texttt{\bfseries metrizableSpace\_subProbabilityMeasure}}\hfill{\footnotesize \path{lean/pointfree/Pointfree/Measure/DaggerChart.lean}, line 425}\par\nopagebreak
\begin{leansrc}
theorem metrizableSpace_subProbabilityMeasure [MetrizableSpace X]
    [SeparableSpace X] [BorelSpace X] :
    MetrizableSpace (SubProbabilityMeasure X)
\end{leansrc}
\par\medskip\noindent{\small \texttt{SuFunctor.Universe.}\texttt{\bfseries belCarrier}}\hfill{\footnotesize \path{lean/su-functor/SuFunctor/Universe/Carrier.lean}, line 63}\par\nopagebreak
\begin{leansrc}
def belCarrier : Set (V → ℝ)
\end{leansrc}
\par\medskip\noindent{\small \texttt{SuFunctor.Universe.}\texttt{\bfseries BelState}}\hfill{\footnotesize \path{lean/su-functor/SuFunctor/Universe/Carrier.lean}, line 119}\par\nopagebreak
\begin{leansrc}
abbrev BelState : Type _
\end{leansrc}
\par\medskip\noindent{\small \texttt{SuFunctor.Universe.}\texttt{\bfseries convex\_belCarrier}}\hfill{\footnotesize \path{lean/su-functor/SuFunctor/Universe/Carrier.lean}, line 72}\par\nopagebreak
\begin{leansrc}
theorem convex_belCarrier : Convex ℝ (belCarrier V)
\end{leansrc}
\par\medskip\noindent{\small \texttt{SuFunctor.Universe.}\texttt{\bfseries isCompact\_belCarrier}}\hfill{\footnotesize \path{lean/su-functor/SuFunctor/Universe/Carrier.lean}, line 107}\par\nopagebreak
\begin{leansrc}
theorem isCompact_belCarrier : IsCompact (belCarrier V)
\end{leansrc}
\par\medskip\noindent{\small \texttt{SuFunctor.Universe.}\texttt{\bfseries probCarrier}}\hfill{\footnotesize \path{lean/su-functor/SuFunctor/Universe/ProbCarrier.lean}, line 45}\par\nopagebreak
\begin{leansrc}
def probCarrier : Set (V → ℝ)
\end{leansrc}
\par\medskip\noindent{\small \texttt{SuFunctor.Universe.}\texttt{\bfseries ProbState}}\hfill{\footnotesize \path{lean/su-functor/SuFunctor/Universe/ProbCarrier.lean}, line 85}\par\nopagebreak
\begin{leansrc}
abbrev ProbState : Type _
\end{leansrc}
\par\medskip\noindent{\small \texttt{SuFunctor.Universe.General.}\texttt{\bfseries sat\_sup\_iff}}\hfill{\footnotesize \path{lean/su-functor/SuFunctor/Universe/General/Faces.lean}, line 72}\par\nopagebreak
\begin{leansrc}
theorem sat_sup_iff {σ φ ψ : StatePenalty X} :
    Sat σ (φ ⊔ ψ) ↔ Sat σ φ ∧ Sat σ ψ
\end{leansrc}
\par\medskip\noindent{\small \texttt{SuFunctor.Universe.General.}\texttt{\bfseries entails\_inf\_iff}}\hfill{\footnotesize \path{lean/su-functor/SuFunctor/Universe/General/Faces.lean}, line 101}\par\nopagebreak
\begin{leansrc}
theorem entails_inf_iff {φ ψ ρ : StatePenalty X} :
    ρ ≤ φ ⊓ ψ ↔ ρ ≤ φ ∧ ρ ≤ ψ
\end{leansrc}
\par\medskip\noindent{\small \texttt{SuFunctor.Universe.General.}\texttt{\bfseries exists\_sat\_inf\_not\_sat}}\hfill{\footnotesize \path{lean/su-functor/SuFunctor/Universe/General/Faces.lean}, line 112}\par\nopagebreak
\begin{leansrc}
theorem exists_sat_inf_not_sat [Nonempty X] :
    ∃ σ φ ψ : StatePenalty X, Sat σ (φ ⊓ ψ) ∧ ¬ Sat σ φ ∧ ¬ Sat σ ψ
\end{leansrc}
\par\medskip\noindent{\small \texttt{SuFunctor.Universe.}\texttt{\bfseries iSup\_boxRows\_eq\_self}}\hfill{\footnotesize \path{lean/su-functor/SuFunctor/Universe/Generation.lean}, line 146}\par\nopagebreak
\begin{leansrc}
theorem iSup_boxRows_eq_self {f : ProbPenalty V}
    (hf : LowerSemicontinuous f) :
    ⨆ (q : ℚ) (l : V → ℚ) (u : V → ℚ) (_ : 0 < q)
        (_ : ∀ ν, boxRow q l u ν ≤ f ν), boxRow q l u = f
\end{leansrc}
\par\medskip\noindent{\small \texttt{SuFunctor.Universe.}\texttt{\bfseries iSup\_tents\_eq\_self}}\hfill{\footnotesize \path{lean/su-functor/SuFunctor/Universe/Tents.lean}, line 76}\par\nopagebreak
\begin{leansrc}
theorem iSup_tents_eq_self {σ : Penalty V} (hσ : LowerSemicontinuous σ) :
    ⨆ (μ₀ : BelState V) (L : ℝ) (M : ℝ)
        (_ : ∀ ν, tent μ₀ L M ν ≤ σ ν), tent μ₀ L M = σ
\end{leansrc}
\par\medskip\noindent{\small \texttt{SuFunctor.Universe.}\texttt{\bfseries ProbPenalty}}\hfill{\footnotesize \path{lean/su-functor/SuFunctor/Universe/ProbCarrier.lean}, line 157}\par\nopagebreak
\begin{leansrc}
abbrev ProbPenalty : Type _
\end{leansrc}
\par\medskip\noindent{\small \texttt{SuFunctor.Universe.General.}\texttt{\bfseries sat\_sup\_iff}}\hfill{\footnotesize \path{lean/su-functor/SuFunctor/Universe/General/Faces.lean}, line 72}\par\nopagebreak
\begin{leansrc}
theorem sat_sup_iff {σ φ ψ : StatePenalty X} :
    Sat σ (φ ⊔ ψ) ↔ Sat σ φ ∧ Sat σ ψ
\end{leansrc}
\par\medskip\noindent{\small \texttt{SuFunctor.Universe.General.}\texttt{\bfseries entails\_inf\_iff}}\hfill{\footnotesize \path{lean/su-functor/SuFunctor/Universe/General/Faces.lean}, line 101}\par\nopagebreak
\begin{leansrc}
theorem entails_inf_iff {φ ψ ρ : StatePenalty X} :
    ρ ≤ φ ⊓ ψ ↔ ρ ≤ φ ∧ ρ ≤ ψ
\end{leansrc}
\par\medskip\noindent{\small \texttt{SuFunctor.Universe.General.}\texttt{\bfseries exists\_sat\_inf\_not\_sat}}\hfill{\footnotesize \path{lean/su-functor/SuFunctor/Universe/General/Faces.lean}, line 112}\par\nopagebreak
\begin{leansrc}
theorem exists_sat_inf_not_sat [Nonempty X] :
    ∃ σ φ ψ : StatePenalty X, Sat σ (φ ⊓ ψ) ∧ ¬ Sat σ φ ∧ ¬ Sat σ ψ
\end{leansrc}
\par\medskip\noindent{\small \texttt{SuFunctor.Universe.General.}\texttt{\bfseries lowerSemicontinuous\_wall\_iff}}\hfill{\footnotesize \path{lean/su-functor/SuFunctor/Universe/General/Faces.lean}, line 222}\par\nopagebreak
\begin{leansrc}
theorem lowerSemicontinuous_wall_iff {C : Set (SubProbabilityMeasure X)} :
    LowerSemicontinuous (wall C) ↔ IsClosed C
\end{leansrc}
\par\medskip\noindent{\small \texttt{SuFunctor.Universe.General.}\texttt{\bfseries iSup\_gambleRows\_eq\_self}}\hfill{\footnotesize \path{lean/su-functor/SuFunctor/Universe/General/GenerationGeneral.lean}, line 248}\par\nopagebreak
\begin{leansrc}
theorem iSup_gambleRows_eq_self {ι : Y → FiniteMeasure X}
    (hι : Topology.IsInducing ι) {f : Y → ℝ≥0∞} (hf : LowerSemicontinuous f) :
    ⨆ (q : ℚ) (F : Finset ((X →ᵇ ℝ≥0) × ℚ × ℚ)) (_ : 0 < q)
        (_ : ∀ ν, gambleRow ι (ENNReal.ofReal q) F ν ≤ f ν),
      gambleRow ι (ENNReal.ofReal q) F = f
\end{leansrc}
\par\medskip\noindent{\small \texttt{SuFunctor.Universe.General.}\texttt{\bfseries isInducing\_testProfile}}\hfill{\footnotesize \path{lean/su-functor/SuFunctor/Universe/General/GenerationGeneral.lean}, line 176}\par\nopagebreak
\begin{leansrc}
theorem isInducing_testProfile {ι : Y → FiniteMeasure X}
    (hι : Topology.IsInducing ι) : Topology.IsInducing (testProfile ι)
\end{leansrc}
\par\medskip\noindent{\small \texttt{SuFunctor.Universe.General.}\texttt{\bfseries exists\_countable\_gambleRows\_eq\_self}}\hfill{\footnotesize \path{lean/su-functor/SuFunctor/Universe/General/GenerationGeneral.lean}, line 305}\par\nopagebreak
\begin{leansrc}
theorem exists_countable_gambleRows_eq_self [SecondCountableTopology Y]
    {ι : Y → FiniteMeasure X} (hι : Topology.IsInducing ι) {f : Y → ℝ≥0∞}
    (hf : LowerSemicontinuous f) :
    ∃ R : Set (ℚ × Finset ((X →ᵇ ℝ≥0) × ℚ × ℚ)), R.Countable
      ∧ (∀ p ∈ R, ∀ ν, gambleRow ι (ENNReal.ofReal p.1) p.2 ν ≤ f ν)
      ∧ ⨆ p ∈ R, gambleRow ι (ENNReal.ofReal p.1) p.2 = f
\end{leansrc}
\par\medskip\noindent{\small \texttt{SuFunctor.Universe.General.}\texttt{\bfseries secondCountableTopology\_probabilityMeasure}}\hfill{\footnotesize \path{lean/su-functor/SuFunctor/Universe/General/GenerationGeneral.lean}, line 433}\par\nopagebreak
\begin{leansrc}
theorem secondCountableTopology_probabilityMeasure
    [CompactSpace X] [TopologicalSpace.MetrizableSpace X] [BorelSpace X] :
    SecondCountableTopology (ProbabilityMeasure X)
\end{leansrc}
\par\medskip\noindent{\small \texttt{SuFunctor.Plain.Joint.}\texttt{\bfseries toMeas\_surjective}}\hfill{\footnotesize \path{lean/su-functor/SuFunctor/Plain/SeamDensity.lean}, line 333}\par\nopagebreak
\begin{leansrc}
theorem Joint.toMeas_surjective :
    Function.Surjective (Joint.toMeas : Joint X → MeasJoint X)
\end{leansrc}
\par\medskip\noindent{\small \texttt{Qtsl.Imprecise.}\texttt{\bfseries sum\_dominates\_sup}}\hfill{\footnotesize \path{lean/qtsl-aux/QtslAux/Imprecise/MarginalCompat.lean}, line 103}\par\nopagebreak
\begin{leansrc}
theorem sum_dominates_sup (m : ∀ i, FinPMF X → FinPMF (Xi i))
    (f : ∀ i, FinPMF (Xi i) → ENNRat) (ν : FinPMF X) :
    supPenalty m f ν ≤ compositePenalty m f ν
\end{leansrc}
\par\medskip\noindent{\small \texttt{Qtsl.Imprecise.}\texttt{\bfseries sublevel\_marginal\_of\_dominates}}\hfill{\footnotesize \path{lean/qtsl-aux/QtslAux/Imprecise/MarginalCompat.lean}, line 136}\par\nopagebreak
\begin{leansrc}
theorem sublevel_marginal_of_dominates (m : ∀ i, FinPMF X → FinPMF (Xi i))
    (f : ∀ i, FinPMF (Xi i) → ENNRat) {F : FinPMF X → ENNRat} {ν : FinPMF X}
    (hdom : supPenalty m f ν ≤ F ν) {t : ENNRat}
    (h : F ν ≤ t) (i : ι) : f i (m i ν) ≤ t
\end{leansrc}
\par\medskip\noindent{\small \texttt{SuFunctor.Universe.General.}\texttt{\bfseries compileEquiv}}\hfill{\footnotesize \path{lean/su-functor/SuFunctor/Universe/RadialUniverse.lean}, line 99}\par\nopagebreak
\begin{leansrc}
noncomputable def compileEquiv [Nonempty X] :
    {f : Penalty X // LowerSemicontinuous f} ≃ RadialUniverse X where …
\end{leansrc}
\par\medskip\noindent{\small \texttt{SuFunctor.Adequacy.}\texttt{\bfseries rwall\_ne\_top}}\hfill{\footnotesize \path{lean/su-functor/SuFunctor/Universe/RadialUniverse.lean}, line 349}\par\nopagebreak
\begin{leansrc}
theorem rwall_ne_top (c : SubPMF V → Prop) :
    rwall c ≠ (⊤ : Plain.Pen V)
\end{leansrc}
\par\medskip\noindent{\small \texttt{SuFunctor.Adequacy.}\texttt{\bfseries contradictory\_composite\_denotes\_bot}}\hfill{\footnotesize \path{lean/su-functor/SuFunctor/Universe/RadialUniverse.lean}, line 589}\par\nopagebreak
\begin{leansrc}
theorem contradictory_composite_denotes_bot :
    rwall xyEqC ⊔ rwall xFalseC ⊔ rwall yTrueC
      = rayPoint (SubPMF.zero (Bool × Bool))
\end{leansrc}
\par\medskip\noindent{\small \texttt{SuFunctor.Universe.General.}\texttt{\bfseries compile\_ne\_singleton}}\hfill{\footnotesize \path{lean/su-functor/SuFunctor/Universe/RadialUniverse.lean}, line 172}\par\nopagebreak
\begin{leansrc}
theorem compile_ne_singleton (f : Penalty X) {ν : SubProbabilityMeasure X}
    (hν : ν ≠ 0) : compile f ≠ ({ν} : Set (SubProbabilityMeasure X))
\end{leansrc}
\par\medskip\noindent{\small \texttt{SuFunctor.StateSpace.}\texttt{\bfseries sortedCone}}\hfill{\footnotesize \path{lean/su-functor/SuFunctor/StateSpace/FamCompletion.lean}, line 165}\par\nopagebreak
\begin{leansrc}
def sortedCone (n : ℕ) : OmegaSA
\end{leansrc}
\par\medskip\noindent{\small \texttt{SuFunctor.StateSpace.}\texttt{\bfseries mem\_sortedCone\_iff\_monotone}}\hfill{\footnotesize \path{lean/su-functor/SuFunctor/StateSpace/FamCompletion.lean}, line 132}\par\nopagebreak
\begin{leansrc}
theorem mem_sortedCone_iff_monotone {n : ℕ} (x : Fin n → ℝ) :
    x ∈ BasicSet (sortedConds n) ↔ Monotone x
\end{leansrc}
\par\medskip\noindent{\small \texttt{SuFunctor.BagFormer.}\texttt{\bfseries bagObj}}\hfill{\footnotesize \path{lean/su-functor/SuFunctor/BagFormer.lean}, line 99}\par\nopagebreak
\begin{leansrc}
def bagObj (A : WireSpan El) : WireSpan El
\end{leansrc}
\par\medskip\noindent{\small \texttt{SuFunctor.StateSpace.}\texttt{\bfseries bagIsColimit}}\hfill{\footnotesize \path{lean/su-functor/SuFunctor/StateSpace/FamCompletion.lean}, line 188}\par\nopagebreak
\begin{leansrc}
def bagIsColimit :
    IsColimit (famComponentsCofan bagObj)
\end{leansrc}
\par\medskip\noindent{\small \texttt{Pointfree.Prob.}\texttt{\bfseries lowerSemicontinuous\_eval}}\hfill{\footnotesize \path{lean/pointfree/Pointfree/Valuation/Topology.lean}, line 79}\par\nopagebreak
\begin{leansrc}
theorem lowerSemicontinuous_eval (a : M) :
    LowerSemicontinuous fun μ : Prob M => μ a
\end{leansrc}
\par\medskip\noindent{\small \texttt{Pointfree.SubProb.}\texttt{\bfseries upperSemicontinuous\_mass}}\hfill{\footnotesize \path{lean/pointfree/Pointfree/Valuation/Topology.lean}, line 122}\par\nopagebreak
\begin{leansrc}
theorem upperSemicontinuous_mass :
    UpperSemicontinuous fun μ : SubProb L => μ ⊤
\end{leansrc}
\par\medskip\noindent{\small \texttt{Pointfree.SubProb.}\texttt{\bfseries isClosed\_liveness\_wall}}\hfill{\footnotesize \path{lean/pointfree/Pointfree/Valuation/Topology.lean}, line 132}\par\nopagebreak
\begin{leansrc}
theorem isClosed_liveness_wall (p : ℝ≥0∞) :
    IsClosed {μ : SubProb L | p ≤ μ ⊤}
\end{leansrc}
\par\medskip\noindent{\small \texttt{SuFunctor.}\texttt{\bfseries satCredal\_lower\_rayC\_geC}}\hfill{\footnotesize \path{lean/su-functor/SuFunctor/Adequacy/Envelope.lean}, line 390}\par\nopagebreak
\begin{leansrc}
theorem satCredal_lower_rayC_geC {μ : SubPMF V} (hμ : μ ≠ SubPMF.zero V)
    (A : Set V) (q : ℚ) :
    SatCredal .lower (rayC μ) (geC A q) ↔ q ≤ μ.prob A
\end{leansrc}
\par\medskip\noindent{\small \texttt{SuFunctor.}\texttt{\bfseries maxSlice\_rayC\_empty\_iff}}\hfill{\footnotesize \path{lean/su-functor/SuFunctor/Adequacy/Envelope.lean}, line 216}\par\nopagebreak
\begin{leansrc}
theorem maxSlice_rayC_empty_iff {μ : SubPMF V} :
    (∀ ν, ¬ maxSlice (rayC μ) ν) ↔ μ = SubPMF.zero V
\end{leansrc}
\par\medskip\noindent{\small \texttt{SuFunctor.}\texttt{\bfseries not\_satCredal\_upper\_rayC\_geC}}\hfill{\footnotesize \path{lean/su-functor/SuFunctor/Adequacy/Envelope.lean}, line 381}\par\nopagebreak
\begin{leansrc}
theorem not_satCredal_upper_rayC_geC (μ : SubPMF V) {A : Set V} {q : ℚ}
    (hq : 0 < q) : ¬ SatCredal .upper (rayC μ) (geC A q)
\end{leansrc}
\par\medskip\noindent{\small \texttt{SuFunctor.Universe.SpecDnf.}\texttt{\bfseries top\_notMem\_modelClass}}\hfill{\footnotesize \path{lean/su-functor/SuFunctor/Universe/SpecDnf.lean}, line 216}\par\nopagebreak
\begin{leansrc}
theorem top_notMem_modelClass (D : SpecDnf K) :
    (⊤ : K → ℝ≥0∞) ∉ D.modelClass
\end{leansrc}
\par\medskip\noindent{\small \texttt{SuFunctor.Universe.General.}\texttt{\bfseries isClosed\_compile\_iff}}\hfill{\footnotesize \path{lean/su-functor/SuFunctor/Universe/General/Closedness.lean}, line 308}\par\nopagebreak
\begin{leansrc}
theorem isClosed_compile_iff [Nonempty X] {f : Penalty X} :
    IsClosed (compile f) ↔ LowerSemicontinuous f
\end{leansrc}
\par\medskip\noindent{\small \texttt{SuFunctor.Universe.General.}\texttt{\bfseries exists\_lsc\_compile\_eq\_iff}}\hfill{\footnotesize \path{lean/su-functor/SuFunctor/Universe/General/Closedness.lean}, line 294}\par\nopagebreak
\begin{leansrc}
theorem exists_lsc_compile_eq_iff [Nonempty X]
    {R : Set (SubProbabilityMeasure X)} :
    (∃ f : Penalty X, LowerSemicontinuous f ∧ compile f = R)
      ↔ R.Nonempty ∧ IsClosed R ∧ IsRadial R
\end{leansrc}
\par\medskip\noindent{\small \texttt{SuFunctor.Universe.General.}\texttt{\bfseries compile\_rayPenalty}}\hfill{\footnotesize \path{lean/su-functor/SuFunctor/Universe/General/Closedness.lean}, line 260}\par\nopagebreak
\begin{leansrc}
theorem compile_rayPenalty [Nonempty X] {R : Set (SubProbabilityMeasure X)}
    (hcl : IsClosed R) (hrad : IsRadial R) (hne : R.Nonempty) :
    compile (rayPenalty R) = R
\end{leansrc}
\par\medskip\noindent{\small \texttt{SuFunctor.Universe.General.}\texttt{\bfseries compile\_injective}}\hfill{\footnotesize \path{lean/su-functor/SuFunctor/Universe/General/Dictionary.lean}, line 238}\par\nopagebreak
\begin{leansrc}
theorem compile_injective [Nonempty X] :
    Function.Injective (compile (X := X))
\end{leansrc}
\par\medskip\noindent{\small \texttt{SuFunctor.Universe.General.}\texttt{\bfseries rayPenalty\_compile}}\hfill{\footnotesize \path{lean/su-functor/SuFunctor/Universe/General/Dictionary.lean}, line 231}\par\nopagebreak
\begin{leansrc}
theorem rayPenalty_compile [Nonempty X] (f : Penalty X) :
    rayPenalty (compile f) = f
\end{leansrc}
\par\medskip\noindent{\small \texttt{SuFunctor.Universe.General.}\texttt{\bfseries upperSemicontinuous\_rayGauge}}\hfill{\footnotesize \path{lean/su-functor/SuFunctor/Universe/General/Closedness.lean}, line 233}\par\nopagebreak
\begin{leansrc}
theorem upperSemicontinuous_rayGauge {R : Set (SubProbabilityMeasure X)}
    (hcl : IsClosed R) (hrad : IsRadial R) (hne : R.Nonempty) :
    UpperSemicontinuous (rayGauge R)
\end{leansrc}
\par\medskip\noindent{\small \texttt{SuFunctor.Universe.General.}\texttt{\bfseries le\_gradeN\_of\_forall\_lt}}\hfill{\footnotesize \path{lean/su-functor/SuFunctor/Universe/General/Carrier.lean}, line 326}\par\nopagebreak
\begin{leansrc}
theorem le_gradeN_of_forall_lt {a : ℝ≥0∞} {c : ℝ≥0} (hc : c ≤ 1)
    (h : ∀ y, y < a → c ≤ gradeN y) : c ≤ gradeN a
\end{leansrc}
\par\medskip\noindent{\small \texttt{SuFunctor.Universe.}\texttt{\bfseries compile}}\hfill{\footnotesize \path{lean/su-functor/SuFunctor/Universe/VietorisDictionary.lean}, line 168}\par\nopagebreak
\begin{leansrc}
def compile (f : ProbPenalty V) : Set (BelState V)
\end{leansrc}
\par\medskip\noindent{\small \texttt{SuFunctor.Universe.EpistemicStates.}\texttt{\bfseries condTarget\_mix}}\hfill{\footnotesize \path{lean/su-functor/SuFunctor/Universe/EpistemicStates.lean}, line 935}\par\nopagebreak
\begin{leansrc}
theorem condTarget_mix (k : B → ProbState A) {μ ν : ProbState (A × B)}
    {s t : ℝ} (hs : 0 ≤ s) (ht : 0 ≤ t) (hst : s + t = 1) :
    condTarget k (ProbState.mix s t hs ht hst μ ν)
      = ProbState.mix s t hs ht hst (condTarget k μ)
          (condTarget k ν)
\end{leansrc}
\par\medskip\noindent{\small \texttt{SuFunctor.Universe.EpistemicStates.}\texttt{\bfseries klq\_mix\_mix\_le}}\hfill{\footnotesize \path{lean/su-functor/SuFunctor/Universe/EpistemicStates.lean}, line 429}\par\nopagebreak
\begin{leansrc}
theorem klq_mix_mix_le {μ ν ρ σ : ProbState V} {s t : ℝ}
    (hs : 0 < s) (ht : 0 < t) (hst : s + t = 1) :
    klq (ProbState.mix s t hs.le ht.le hst μ ν)
        (ProbState.mix s t hs.le ht.le hst ρ σ)
      ≤ ENNReal.ofReal s * klq μ ρ + ENNReal.ofReal t * klq ν σ
\end{leansrc}
\par\medskip\noindent{\small \texttt{SuFunctor.Universe.EpistemicStates.}\texttt{\bfseries lowerSemicontinuous\_condPenalty}}\hfill{\footnotesize \path{lean/su-functor/SuFunctor/Universe/EpistemicStates.lean}, line 1056}\par\nopagebreak
\begin{leansrc}
theorem lowerSemicontinuous_condPenalty {β : ℝ≥0∞} (hβ : β ≠ ⊤)
    (k : B → ProbState A) :
    LowerSemicontinuous (condPenalty β k)
\end{leansrc}
\par\medskip\noindent{\small \texttt{SuFunctor.Universe.EpistemicStates.}\texttt{\bfseries klq\_condTarget\_eq\_zero\_iff}}\hfill{\footnotesize \path{lean/su-functor/SuFunctor/Universe/EpistemicStates.lean}, line 974}\par\nopagebreak
\begin{leansrc}
theorem klq_condTarget_eq_zero_iff (k : B → ProbState A)
    (μ : ProbState (A × B)) :
    klq μ (condTarget k μ) = 0
      ↔ ∀ b, 0 < (sndMarg μ : B → ℝ) b
          → ∀ a, (μ : A × B → ℝ) (a, b) / (sndMarg μ : B → ℝ) b
              = (k b : A → ℝ) a
\end{leansrc}
\par\medskip\noindent{\small \texttt{SuFunctor.Universe.EpistemicStates.}\texttt{\bfseries mutualInfoQ\_eq\_zero\_iff}}\hfill{\footnotesize \path{lean/su-functor/SuFunctor/Universe/EpistemicStates.lean}, line 844}\par\nopagebreak
\begin{leansrc}
theorem mutualInfoQ_eq_zero_iff (μ : ProbState (A × B)) :
    mutualInfoQ μ = 0 ↔ IsProductLaw μ
\end{leansrc}
\par\medskip\noindent{\small \texttt{SuFunctor.Universe.EpistemicStates.}\texttt{\bfseries continuous\_mutualInfoQ}}\hfill{\footnotesize \path{lean/su-functor/SuFunctor/Universe/EpistemicStates.lean}, line 882}\par\nopagebreak
\begin{leansrc}
theorem continuous_mutualInfoQ :
    Continuous (mutualInfoQ (A := A) (B := B))
\end{leansrc}
\par\medskip\noindent{\small \texttt{SuFunctor.Universe.EpistemicStates.Witness.}\texttt{\bfseries not\_penaltyConvex\_indepPenalty}}\hfill{\footnotesize \path{lean/su-functor/SuFunctor/Universe/EpistemicStates.lean}, line 1262}\par\nopagebreak
\begin{leansrc}
theorem not_penaltyConvex_indepPenalty {β : ℝ≥0∞} (hβ : β ≠ 0) :
    ¬ PenaltyConvex (indepPenalty (A := Fin 2) (B := Fin 2) β)
\end{leansrc}
\par\medskip\noindent{\small \texttt{SuFunctor.Universe.EpistemicStates.}\texttt{\bfseries indepBody\_credalBody\_not\_convex}}\hfill{\footnotesize \path{lean/su-functor/SuFunctor/Universe/EpistemicStates.lean}, line 1435}\par\nopagebreak
\begin{leansrc}
theorem indepBody_credalBody_not_convex :
    IsCredalBody (indepBody (Fin 2) (Fin 2))
      ∧ ¬ IsConvexCredalBody (indepBody (Fin 2) (Fin 2))
\end{leansrc}
\par\medskip\noindent{\small \texttt{SuFunctor.Universe.EpistemicStates.}\texttt{\bfseries compile\_marginalWallSum}}\hfill{\footnotesize \path{lean/su-functor/SuFunctor/Universe/EpistemicStates.lean}, line 1145}\par\nopagebreak
\begin{leansrc}
theorem compile_marginalWallSum {MX : Set (ProbState A)}
    {MY : Set (ProbState B)} (hX : MX.Nonempty) (hY : MY.Nonempty) :
    compile (fun μ : ProbState (A × B) =>
        probWall MX (fstMarg μ) + probWall MY (sndMarg μ))
      = {x : BelState (A × B) | ∃ μ ∈ frechetClass MX MY, ∃ γ : ℝ,
          0 ≤ γ ∧ γ ≤ 1 ∧ (x : A × B → ℝ) = γ • (μ : A × B → ℝ)}
\end{leansrc}
\par\medskip\noindent{\small \texttt{SuFunctor.Universe.EpistemicStates.}\texttt{\bfseries marginalSum\_eq\_zero\_iff}}\hfill{\footnotesize \path{lean/su-functor/SuFunctor/Universe/EpistemicStates.lean}, line 1078}\par\nopagebreak
\begin{leansrc}
theorem marginalSum_eq_zero_iff (g : ProbState A → ℝ≥0∞)
    (h : ProbState B → ℝ≥0∞) (μ : ProbState (A × B)) :
    g (fstMarg μ) + h (sndMarg μ) = 0
      ↔ g (fstMarg μ) = 0 ∧ h (sndMarg μ) = 0
\end{leansrc}
\par\medskip\noindent{\small \texttt{SuFunctor.Universe.EpistemicStates.Witness.}\texttt{\bfseries frechet\_does\_not\_imply\_independence}}\hfill{\footnotesize \path{lean/su-functor/SuFunctor/Universe/EpistemicStates.lean}, line 1321}\par\nopagebreak
\begin{leansrc}
theorem frechet_does_not_imply_independence :
    (probWall {uniform2} (fstMarg chi)
        + probWall {uniform2} (sndMarg chi) = 0)
      ∧ ¬ IsProductLaw chi ∧ mutualInfoQ chi ≠ 0
\end{leansrc}
\par\medskip\noindent{\small \texttt{SuFunctor.Universe.CredalPoset.}\texttt{\bfseries house\_includes\_all}}\hfill{\footnotesize \path{lean/su-functor/SuFunctor/Universe/CredalPoset.lean}, line 47}\par\nopagebreak
\begin{leansrc}
theorem house_includes_all {V : Type*} [Fintype V] [Nonempty V]
    (S : Set (BelState V)) :
    (IsConvexCredalBody S → IsCredalBody S) ∧
    (IsRichardsonBody S → IsCredalBody S) ∧
    (IsUltra S → IsCredalBody S) ∧
    (IsMVS S → IsCredalBody S) ∧
    (IsHL S → IsCredalBody S) ∧
    (IsCozman S → IsCredalBody S)
\end{leansrc}
\par\medskip\noindent{\small \texttt{SuFunctor.Universe.General.CredalPoset.}\texttt{\bfseries scaleSegment\_eq\_compile\_dirWall}}\hfill{\footnotesize \path{lean/su-functor/SuFunctor/Universe/General/CredalPoset.lean}, line 446}\par\nopagebreak
\begin{leansrc}
theorem scaleSegment_eq_compile_dirWall {K : Set (ProbabilityMeasure X)}
    (hne : K.Nonempty) : scaleSegment K = compile (dirWall K)
\end{leansrc}
\par\medskip\noindent{\small \texttt{SuFunctor.Universe.CredalPoset.}\texttt{\bfseries exists\_penaltyConvex\_not\_closed}}\hfill{\footnotesize \path{lean/su-functor/SuFunctor/Universe/CredalPoset/Adjudication.lean}, line 172}\par\nopagebreak
\begin{leansrc}
theorem exists_penaltyConvex_not_closed :
    ∃ f : ProbPenalty (Fin 2),
      PenaltyConvex f ∧ ¬ IsClosed (compile f)
\end{leansrc}
\par\medskip\noindent{\small \texttt{SuFunctor.Universe.CredalPoset.}\texttt{\bfseries exists\_penaltyConvex\_compile\_not\_richardson}}\hfill{\footnotesize \path{lean/su-functor/SuFunctor/Universe/CredalPoset/Adjudication.lean}, line 201}\par\nopagebreak
\begin{leansrc}
theorem exists_penaltyConvex_compile_not_richardson :
    ∃ f : ProbPenalty (Fin 2),
      PenaltyConvex f ∧ ¬ IsRichardsonBody (compile f)
\end{leansrc}
\par\medskip\noindent{\small \texttt{SuFunctor.Universe.CredalPoset.IsRichardsonBody.}\texttt{\bfseries isConvexCredalBody}}\hfill{\footnotesize \path{lean/su-functor/SuFunctor/Universe/CredalPoset/Inclusions.lean}, line 176}\par\nopagebreak
\begin{leansrc}
theorem IsRichardsonBody.isConvexCredalBody [Nonempty V] {S : Set (BelState V)}
    (h : IsRichardsonBody S) : IsConvexCredalBody S
\end{leansrc}
\par\medskip\noindent{\small \texttt{SuFunctor.Universe.General.CredalPoset.}\texttt{\bfseries GradeConvex}}\hfill{\footnotesize \path{lean/su-functor/SuFunctor/Universe/General/CredalPoset.lean}, line 165}\par\nopagebreak
\begin{leansrc}
def GradeConvex (f : Penalty X) : Prop
\end{leansrc}
\par\medskip\noindent{\small \texttt{SuFunctor.Universe.General.CredalPoset.}\texttt{\bfseries convexBody\_compile\_of\_gradeConvex}}\hfill{\footnotesize \path{lean/su-functor/SuFunctor/Universe/General/CredalPoset.lean}, line 252}\par\nopagebreak
\begin{leansrc}
theorem convexBody_compile_of_gradeConvex [Nonempty X] {f : Penalty X}
    (hgc : GradeConvex f) : ConvexBody (compile f)
\end{leansrc}
\par\medskip\noindent{\small \texttt{SuFunctor.Universe.CredalPoset.IsUltra.}\texttt{\bfseries isConvexCredalBody}}\hfill{\footnotesize \path{lean/su-functor/SuFunctor/Universe/CredalPoset/Defs.lean}, line 196}\par\nopagebreak
\begin{leansrc}
theorem IsUltra.isConvexCredalBody {S : Set (BelState V)}
    (h : IsUltra S) : IsConvexCredalBody S
\end{leansrc}
\par\medskip\noindent{\small \texttt{SuFunctor.Universe.CredalPoset.IsPrevisionBody.}\texttt{\bfseries isUltra}}\hfill{\footnotesize \path{lean/su-functor/SuFunctor/Universe/CredalPoset/Previsions.lean}, line 50}\par\nopagebreak
\begin{leansrc}
theorem IsPrevisionBody.isUltra {S : Set (BelState V)}
    (hP : IsPrevisionBody S) : IsUltra S
\end{leansrc}
\par\medskip\noindent{\small \texttt{SuFunctor.Universe.CredalPoset.}\texttt{\bfseries isPrevisionBody\_iff\_isUltra}}\hfill{\footnotesize \path{lean/su-functor/SuFunctor/Universe/CredalPoset/Previsions.lean}, line 202}\par\nopagebreak
\begin{leansrc}
theorem isPrevisionBody_iff_isUltra {S : Set (BelState V)} :
    IsPrevisionBody S ↔ IsUltra S
\end{leansrc}
\par\medskip\noindent{\small \texttt{SuFunctor.Universe.General.CredalPoset.}\texttt{\bfseries isClosed\_ptwiseDownHull}}\hfill{\footnotesize \path{lean/su-functor/SuFunctor/Universe/General/CredalPoset.lean}, line 588}\par\nopagebreak
\begin{leansrc}
theorem isClosed_ptwiseDownHull
    (horder : IsClosed {p : FiniteMeasure X × FiniteMeasure X |
      (p.1 : Measure X) ≤ (p.2 : Measure X)})
    (F : Finset (SubProbabilityMeasure X)) :
    IsClosed (ptwiseDownHull F)
\end{leansrc}
\par\medskip\noindent{\small \texttt{SuFunctor.Universe.CredalPoset.IsMVS.}\texttt{\bfseries isUltra}}\hfill{\footnotesize \path{lean/su-functor/SuFunctor/Universe/CredalPoset/Inclusions.lean}, line 270}\par\nopagebreak
\begin{leansrc}
theorem IsMVS.isUltra {S : Set (BelState V)} (h : IsMVS S) : IsUltra S
\end{leansrc}
\par\medskip\noindent{\small \texttt{SuFunctor.}\texttt{\bfseries maxSlice\_rayC}}\hfill{\footnotesize \path{lean/su-functor/SuFunctor/Adequacy/Envelope.lean}, line 176}\par\nopagebreak
\begin{leansrc}
theorem maxSlice_rayC {μ : SubPMF V} (hμ : μ ≠ SubPMF.zero V)
    (ν : SubPMF V) : maxSlice (rayC μ) ν ↔ singletonC μ ν
\end{leansrc}
\par\medskip\noindent{\small \texttt{SuFunctor.}\texttt{\bfseries satCredal\_lower\_rayC}}\hfill{\footnotesize \path{lean/su-functor/SuFunctor/Adequacy/Envelope.lean}, line 327}\par\nopagebreak
\begin{leansrc}
theorem satCredal_lower_rayC {μ : SubPMF V} (hμ : μ ≠ SubPMF.zero V)
    (c : SubPMF V → Prop) : SatCredal .lower (rayC μ) c ↔ c μ
\end{leansrc}
\par\medskip\noindent{\small \texttt{SuFunctor.}\texttt{\bfseries DownClosed}}\hfill{\footnotesize \path{lean/su-functor/SuFunctor/Adequacy/Envelope.lean}, line 305}\par\nopagebreak
\begin{leansrc}
def DownClosed (c : SubPMF V → Prop) : Prop
\end{leansrc}
\par\medskip\noindent{\small \texttt{SuFunctor.}\texttt{\bfseries rayC\_eq\_singletonC\_iff}}\hfill{\footnotesize \path{lean/su-functor/SuFunctor/Adequacy/Envelope.lean}, line 251}\par\nopagebreak
\begin{leansrc}
theorem rayC_eq_singletonC_iff {μ : SubPMF V} :
    (∀ ν, rayC μ ν ↔ singletonC μ ν) ↔ μ = SubPMF.zero V
\end{leansrc}
\par\medskip\noindent{\small \texttt{SuFunctor.}\texttt{\bfseries maxSlice\_unionRays\_subset}}\hfill{\footnotesize \path{lean/su-functor/SuFunctor/Adequacy/EnvelopeUnion.lean}, line 138}\par\nopagebreak
\begin{leansrc}
theorem maxSlice_unionRays_subset {μs : ι → SubPMF V} {ν : SubPMF V}
    (h : maxSlice (unionRays μs) ν) : unionSingletons μs ν
\end{leansrc}
\par\medskip\noindent{\small \texttt{SuFunctor.}\texttt{\bfseries maxSlice\_unionRays}}\hfill{\footnotesize \path{lean/su-functor/SuFunctor/Adequacy/EnvelopeUnion.lean}, line 183}\par\nopagebreak
\begin{leansrc}
theorem maxSlice_unionRays {μs : ι → SubPMF V} {m : ℚ}
    (hm : ∀ i, (μs i).mass = m) (hm0 : m ≠ 0) (ν : SubPMF V) :
    maxSlice (unionRays μs) ν ↔ unionSingletons μs ν
\end{leansrc}
\par\medskip\noindent{\small \texttt{Qtsl.Imprecise.}\texttt{\bfseries natExtG\_le\_lowerEnvG}}\hfill{\footnotesize \path{lean/qtsl-aux/QtslAux/Imprecise/NaturalExtensionGeneral.lean}, line 372}\par\nopagebreak
\begin{leansrc}
theorem natExtG_le_lowerEnvG {M' : Set (ProbabilityMeasure X)}
    (hM' : M'.Nonempty) (hdom : ∀ i, c i ≤ lowerEnvG M' (h i))
    (g : BddGamble X) : natExtG h c g ≤ lowerEnvG M' g
\end{leansrc}
\par\medskip\noindent{\small \texttt{Qtsl.Imprecise.}\texttt{\bfseries avoidsSureLossG\_of\_dominating}}\hfill{\footnotesize \path{lean/qtsl-aux/QtslAux/Imprecise/NaturalExtensionGeneral.lean}, line 381}\par\nopagebreak
\begin{leansrc}
theorem avoidsSureLossG_of_dominating {M' : Set (ProbabilityMeasure X)}
    (hM' : M'.Nonempty) (hdom : ∀ i, c i ≤ lowerEnvG M' (h i)) :
    AvoidsSureLossG h c
\end{leansrc}
\par\medskip\noindent{\small \texttt{Pointfree.}\texttt{\bfseries Prevision}}\hfill{\footnotesize \path{lean/pointfree/Pointfree/Prevision/Basic.lean}, line 286}\par\nopagebreak
\begin{leansrc}
structure Prevision (L : Type*) [Order.Frame L] where …
\end{leansrc}
\par\medskip\noindent{\small \texttt{SuFunctor.Universe.}\texttt{\bfseries singletons\_not\_definable}}\hfill{\footnotesize \path{lean/su-functor/SuFunctor/Universe/Points.lean}, line 165}\par\nopagebreak
\begin{leansrc}
theorem singletons_not_definable {σ : ProbPenalty V}
    (hσ : LowerSemicontinuous σ) (hproper : ∃ μ, σ μ ≠ ⊤)
    {n : ℕ} (u : Fin n → ProbPenalty V) (𝒱 : Fin n → Set (ProbPenalty V))
    (hfin : ∀ i, (𝒱 i).Finite) :
    ¬ ∀ τ : ProbPenalty V, LowerSemicontinuous τ →
        ((∃ i, CoupledSat τ (u i) (𝒱 i)) ↔ τ = σ)
\end{leansrc}
\par\medskip\noindent{\small \texttt{SuFunctor.Plain.}\texttt{\bfseries idealPoint\_le\_iff}}\hfill{\footnotesize \path{lean/su-functor/SuFunctor/Plain/Battery.lean}, line 117}\par\nopagebreak
\begin{leansrc}
theorem idealPoint_le_iff (σ : Pen (states A)) (𝒮 : SpecClass A) :
    ({σ} : SpecClass A) ⊆ 𝒮 ↔ σ ∈ 𝒮
\end{leansrc}
\par\medskip\noindent{\small \texttt{SuFunctor.Universe.}\texttt{\bfseries crescentClass}}\hfill{\footnotesize \path{lean/su-functor/SuFunctor/Universe/VietorisStratum.lean}, line 315}\par\nopagebreak
\begin{leansrc}
def crescentClass (u : ProbPenalty V) (𝒱 : Set (ProbPenalty V)) :
    Set (NonemptyCompacts (BelState V))
\end{leansrc}
\par\medskip\noindent{\small \texttt{SuFunctor.Universe.General.}\texttt{\bfseries isLocallyClosed\_crescentClass}}\hfill{\footnotesize \path{lean/su-functor/SuFunctor/Universe/General/CrescentStratum.lean}, line 246}\par\nopagebreak
\begin{leansrc}
theorem isLocallyClosed_crescentClass [Nonempty X] {u : Penalty X}
    {𝒱 : Set (Penalty X)} (hu : LowerSemicontinuous u)
    (h𝒱 : 𝒱.Finite) (hlsc : ∀ v ∈ 𝒱, LowerSemicontinuous v) :
    IsLocallyClosed (crescentClass u 𝒱)
\end{leansrc}
\par\medskip\noindent{\small \texttt{SuFunctor.Universe.SpecDnf.}\texttt{\bfseries compileNC\_mem\_crescentClass\_iff}}\hfill{\footnotesize \path{lean/su-functor/SuFunctor/Universe/SpecDnf.lean}, line 512}\par\nopagebreak
\begin{leansrc}
theorem SpecDnf.compileNC_mem_crescentClass_iff [Nonempty V]
    {σ : ProbPenalty V} (hσ : LowerSemicontinuous σ)
    (D : SpecDnf (ProbState V)) :
    compileNC σ hσ ∈ D.crescentClass ↔ σ ∈ D.modelClass
\end{leansrc}
\par\medskip\noindent{\small \texttt{SuFunctor.Universe.General.}\texttt{\bfseries compile\_subset\_compile\_iff}}\hfill{\footnotesize \path{lean/su-functor/SuFunctor/Universe/General/Dictionary.lean}, line 122}\par\nopagebreak
\begin{leansrc}
theorem compile_subset_compile_iff {f h : Penalty X} :
    compile f ⊆ compile h ↔ h ≤ f
\end{leansrc}
\par\medskip\noindent{\small \texttt{SuFunctor.Universe.}\texttt{\bfseries mem\_modelClass\_reachAvoid\_iff}}\hfill{\footnotesize \path{lean/su-functor/SuFunctor/Universe/ReachAvoid.lean}, line 94}\par\nopagebreak
\begin{leansrc}
theorem mem_modelClass_reachAvoid_iff {u ℓ σ : K → ℝ≥0∞} :
    σ ∈ (SpecAtom.reachAvoid u ℓ).modelClass ↔ u ≤ σ ∧ Dip σ ℓ
\end{leansrc}
\par\medskip\noindent{\small \texttt{SuFunctor.Universe.}\texttt{\bfseries coupledSat\_crispRA\_iff}}\hfill{\footnotesize \path{lean/su-functor/SuFunctor/Universe/ReachAvoid.lean}, line 494}\par\nopagebreak
\begin{leansrc}
theorem coupledSat_crispRA_iff {σ : K → ℝ≥0∞} {S T : Set K} :
    σ ∈ (SpecAtom.reachAvoid (wallOn S) (wallOn Tᶜ)).modelClass
      ↔ (∀ x, x ∉ S → σ x = ⊤) ∧ ∃ x ∈ T, σ x < ⊤
\end{leansrc}
\par\medskip\noindent{\small \texttt{SuFunctor.Universe.}\texttt{\bfseries exists\_genClosure\_reachAvoid\_not\_crispGenerated}}\hfill{\footnotesize \path{lean/su-functor/SuFunctor/Universe/ReachAvoid.lean}, line 563}\par\nopagebreak
\begin{leansrc}
theorem exists_genClosure_reachAvoid_not_crispGenerated :
    ∃ C : Set (Unit → ℝ≥0∞),
      GenClosure (reachAvoidClasses Unit) C
        ∧ ¬ GenClosure (crispReachAvoidClasses Unit) C
\end{leansrc}
\par\medskip\noindent{\small \texttt{SuFunctor.Universe.}\texttt{\bfseries genClosure\_reachAvoid\_iff\_exists\_specDnf}}\hfill{\footnotesize \path{lean/su-functor/SuFunctor/Universe/ReachAvoid.lean}, line 259}\par\nopagebreak
\begin{leansrc}
theorem genClosure_reachAvoid_iff_exists_specDnf {C : Set (K → ℝ≥0∞)} :
    GenClosure (reachAvoidClasses K) C
      ↔ ∃ D : SpecDnf K, C = D.modelClass
\end{leansrc}
\par\medskip\noindent{\small \texttt{SuFunctor.Universe.}\texttt{\bfseries genClosure\_reachAvoidCres\_iff\_exists\_specDnf}}\hfill{\footnotesize \path{lean/su-functor/SuFunctor/Universe/ReachAvoid.lean}, line 370}\par\nopagebreak
\begin{leansrc}
theorem genClosure_reachAvoidCres_iff_exists_specDnf
    {C : Set (NonemptyCompacts (BelState V))} :
    GenClosure (reachAvoidCrescents V) C
      ↔ ∃ D : SpecDnf (ProbState V), C = D.crescentClass
\end{leansrc}
\par\medskip\noindent{\small \texttt{SuFunctor.Universe.}\texttt{\bfseries isConstructibleUnion\_of\_genClosure\_reachAvoidCrescentsLsc}}\hfill{\footnotesize \path{lean/su-functor/SuFunctor/Universe/ReachAvoid.lean}, line 456}\par\nopagebreak
\begin{leansrc}
theorem isConstructibleUnion_of_genClosure_reachAvoidCrescentsLsc
    {C : Set (NonemptyCompacts (BelState V))}
    (h : GenClosure (reachAvoidCrescentsLsc V) C) :
    IsConstructibleUnion C
\end{leansrc}
\par\medskip\noindent{\small \texttt{SuFunctor.Universe.}\texttt{\bfseries modelClass\_eq\_iInter\_reachAvoid}}\hfill{\footnotesize \path{lean/su-functor/SuFunctor/Universe/ReachAvoid.lean}, line 122}\par\nopagebreak
\begin{leansrc}
theorem modelClass_eq_iInter_reachAvoid (a : SpecAtom K) :
    a.modelClass
      = ⋂ ℓ ∈ a.marks, (SpecAtom.reachAvoid a.floor ℓ).modelClass
\end{leansrc}
\par\medskip\noindent{\small \texttt{SuFunctor.Universe.}\texttt{\bfseries crescentClass\_eq\_iInter\_reachAvoid}}\hfill{\footnotesize \path{lean/su-functor/SuFunctor/Universe/ReachAvoid.lean}, line 296}\par\nopagebreak
\begin{leansrc}
theorem crescentClass_eq_iInter_reachAvoid (a : SpecAtom (ProbState V)) :
    crescentClass a.floor a.marks
      = ⋂ ℓ ∈ a.marks, crescentClass a.floor {ℓ, ⊤}
\end{leansrc}
\par\medskip\noindent{\small \texttt{SuFunctor.Universe.}\texttt{\bfseries exists\_modelClass\_ne\_iUnion\_reachAvoid}}\hfill{\footnotesize \path{lean/su-functor/SuFunctor/Universe/ReachAvoid.lean}, line 141}\par\nopagebreak
\begin{leansrc}
theorem exists_modelClass_ne_iUnion_reachAvoid :
    ∃ a : SpecAtom Unit,
      a.modelClass
        ≠ ⋃ ℓ ∈ a.marks, (SpecAtom.reachAvoid a.floor ℓ).modelClass
\end{leansrc}
\par\medskip\noindent{\small \texttt{SuFunctor.Universe.}\texttt{\bfseries dip\_sup\_iff}}\hfill{\footnotesize \path{lean/su-functor/SuFunctor/Universe/CoupledJudgment.lean}, line 226}\par\nopagebreak
\begin{leansrc}
theorem dip_sup_iff {σ v₁ v₂ : K → ℝ≥0∞} :
    Dip σ (v₁ ⊔ v₂) ↔ Dip σ v₁ ∨ Dip σ v₂
\end{leansrc}
\par\medskip\noindent{\small \texttt{SuFunctor.Universe.}\texttt{\bfseries IsConstructibleUnion}}\hfill{\footnotesize \path{lean/su-functor/SuFunctor/Universe/LocallyClosedLayer.lean}, line 406}\par\nopagebreak
\begin{leansrc}
def IsConstructibleUnion (s : Set X) : Prop
\end{leansrc}
\par\medskip\noindent{\small \texttt{SuFunctor.Universe.General.}\texttt{\bfseries crescentClass\_inter\_crescentClass}}\hfill{\footnotesize \path{lean/su-functor/SuFunctor/Universe/General/CrescentStratum.lean}, line 286}\par\nopagebreak
\begin{leansrc}
theorem crescentClass_inter_crescentClass {u₁ u₂ : Penalty X}
    {𝒱₁ 𝒱₂ : Set (Penalty X)} :
    crescentClass u₁ 𝒱₁ ∩ crescentClass u₂ 𝒱₂
      = crescentClass (u₁ ⊔ u₂) (𝒱₁ ∪ 𝒱₂)
\end{leansrc}
\par\medskip\noindent{\small \texttt{SuFunctor.Universe.}\texttt{\bfseries inter\_dnfClass\_append}}\hfill{\footnotesize \path{lean/su-functor/SuFunctor/Universe/DisjunctionElim.lean}, line 204}\par\nopagebreak
\begin{leansrc}
theorem inter_dnfClass_append {n m : ℕ}
    (γ : Set (NonemptyCompacts (BelState V)))
    (u : Fin n → ProbPenalty V) (𝒱 : Fin n → Set (ProbPenalty V))
    (u' : Fin m → ProbPenalty V) (𝒱' : Fin m → Set (ProbPenalty V)) :
    γ ∩ dnfClass (Fin.append u u') (Fin.append 𝒱 𝒱')
      = (γ ∩ dnfClass u 𝒱) ∪ (γ ∩ dnfClass u' 𝒱')
\end{leansrc}
\par\medskip\noindent{\small \texttt{SuFunctor.Universe.General.}\texttt{\bfseries compile\_pullPen}}\hfill{\footnotesize \path{lean/su-functor/SuFunctor/Universe/General/TransportGeneral.lean}, line 176}\par\nopagebreak
\begin{leansrc}
theorem compile_pullPen [Nonempty X] (hθ : Continuous θ) (f : Penalty Y) :
    compile (pullPen hθ f) = push θ ⁻¹' compile f
\end{leansrc}
\par\medskip\noindent{\small \texttt{SuFunctor.Universe.General.}\texttt{\bfseries transportClass\_union}}\hfill{\footnotesize \path{lean/su-functor/SuFunctor/Universe/General/TransportGeneral.lean}, line 278}\par\nopagebreak
\begin{leansrc}
theorem transportClass_union (θ : X → Y)
    (C₁ C₂ : Set (FellCloseds (SubProbabilityMeasure Y))) :
    transportClass θ (C₁ ∪ C₂)
      = transportClass θ C₁ ∪ transportClass θ C₂
\end{leansrc}
\par\medskip\noindent{\small \texttt{SuFunctor.Universe.}\texttt{\bfseries ctxOrE\_valid\_iff}}\hfill{\footnotesize \path{lean/su-functor/SuFunctor/Universe/DisjunctionElim.lean}, line 99}\par\nopagebreak
\begin{leansrc}
theorem ctxOrE_valid_iff :
    (∀ γ ψ : Set α, γ ∩ A ⊆ ψ → γ ∩ B ⊆ ψ → γ ∩ J ⊆ ψ) ↔ J ⊆ A ∪ B
\end{leansrc}
\par\medskip\noindent{\small \texttt{SuFunctor.Universe.}\texttt{\bfseries orI\_ctxOrE\_valid\_iff}}\hfill{\footnotesize \path{lean/su-functor/SuFunctor/Universe/DisjunctionElim.lean}, line 107}\par\nopagebreak
\begin{leansrc}
theorem orI_ctxOrE_valid_iff :
    (A ⊆ J ∧ B ⊆ J ∧
        ∀ γ ψ : Set α, γ ∩ A ⊆ ψ → γ ∩ B ⊆ ψ → γ ∩ J ⊆ ψ)
      ↔ J = A ∪ B
\end{leansrc}
\par\medskip\noindent{\small \texttt{SuFunctor.Universe.}\texttt{\bfseries exists\_ctxOrE\_countermodel\_of\_ssubset}}\hfill{\footnotesize \path{lean/su-functor/SuFunctor/Universe/DisjunctionElim.lean}, line 123}\par\nopagebreak
\begin{leansrc}
theorem exists_ctxOrE_countermodel_of_ssubset (h : A ∪ B ⊂ J) :
    ∃ γ ψ : Set α, γ ∩ A ⊆ ψ ∧ γ ∩ B ⊆ ψ ∧ ¬ γ ∩ J ⊆ ψ
\end{leansrc}
\par\medskip\noindent{\small \texttt{SuFunctor.Universe.}\texttt{\bfseries exists\_union\_crescentClass\_not\_isLocallyClosed}}\hfill{\footnotesize \path{lean/su-functor/SuFunctor/Universe/DisjunctionElim.lean}, line 403}\par\nopagebreak
\begin{leansrc}
theorem exists_union_crescentClass_not_isLocallyClosed :
    ∃ u v : ProbPenalty Bool,
      LowerSemicontinuous u ∧ LowerSemicontinuous v ∧
        ¬ IsLocallyClosed
          (crescentClass u ∅
            ∪ crescentClass (0 : ProbPenalty Bool) {v})
\end{leansrc}
\par\medskip\noindent{\small \texttt{SuFunctor.Universe.}\texttt{\bfseries dnfClass\_append}}\hfill{\footnotesize \path{lean/su-functor/SuFunctor/Universe/DisjunctionElim.lean}, line 146}\par\nopagebreak
\begin{leansrc}
theorem dnfClass_append {n m : ℕ} (u : Fin n → ProbPenalty V)
    (𝒱 : Fin n → Set (ProbPenalty V)) (u' : Fin m → ProbPenalty V)
    (𝒱' : Fin m → Set (ProbPenalty V)) :
    dnfClass (Fin.append u u') (Fin.append 𝒱 𝒱')
      = dnfClass u 𝒱 ∪ dnfClass u' 𝒱'
\end{leansrc}
\par\medskip\noindent{\small \texttt{SuFunctor.Universe.}\texttt{\bfseries isConstructibleUnion\_dnfClass\_append}}\hfill{\footnotesize \path{lean/su-functor/SuFunctor/Universe/DisjunctionElim.lean}, line 226}\par\nopagebreak
\begin{leansrc}
theorem isConstructibleUnion_dnfClass_append {n m : ℕ}
    {u : Fin n → ProbPenalty V} {𝒱 : Fin n → Set (ProbPenalty V)}
    {u' : Fin m → ProbPenalty V} {𝒱' : Fin m → Set (ProbPenalty V)}
    (hu : ∀ i, LowerSemicontinuous (u i)) (h𝒱 : ∀ i, (𝒱 i).Finite)
    (hlsc : ∀ i, ∀ v ∈ 𝒱 i, LowerSemicontinuous v)
    (hu' : ∀ i, LowerSemicontinuous (u' i)) (h𝒱' : ∀ i, (𝒱' i).Finite)
    (hlsc' : ∀ i, ∀ v ∈ 𝒱' i, LowerSemicontinuous v) :
    IsConstructibleUnion
      (dnfClass (Fin.append u u') (Fin.append 𝒱 𝒱'))
\end{leansrc}
\par\medskip\noindent{\small \texttt{SuFunctor.Universe.}\texttt{\bfseries ctxOrE\_dnfClass}}\hfill{\footnotesize \path{lean/su-functor/SuFunctor/Universe/DisjunctionElim.lean}, line 216}\par\nopagebreak
\begin{leansrc}
theorem ctxOrE_dnfClass {n m : ℕ} {u : Fin n → ProbPenalty V}
    {𝒱 : Fin n → Set (ProbPenalty V)} {u' : Fin m → ProbPenalty V}
    {𝒱' : Fin m → Set (ProbPenalty V)}
    {γ ψ : Set (NonemptyCompacts (BelState V))}
    (h₁ : γ ∩ dnfClass u 𝒱 ⊆ ψ) (h₂ : γ ∩ dnfClass u' 𝒱' ⊆ ψ) :
    γ ∩ dnfClass (Fin.append u u') (Fin.append 𝒱 𝒱') ⊆ ψ
\end{leansrc}
\par\medskip\noindent{\small \texttt{SuFunctor.Universe.}\texttt{\bfseries isClosed\_pairHull}}\hfill{\footnotesize \path{lean/su-functor/SuFunctor/Universe/HullClosure.lean}, line 270}\par\nopagebreak
\begin{leansrc}
theorem isClosed_pairHull {C D : Set (ProbabilityMeasure X)}
    (hC : IsClosed C) (hD : IsClosed D) : IsClosed (pairHull C D)
\end{leansrc}
\par\medskip\noindent{\small \texttt{SuFunctor.Universe.}\texttt{\bfseries isClosed\_binaryHull}}\hfill{\footnotesize \path{lean/su-functor/SuFunctor/Universe/HullClosure.lean}, line 395}\par\nopagebreak
\begin{leansrc}
theorem isClosed_binaryHull {C D : Set (ProbabilityMeasure X)}
    (hC : IsClosed C) (hD : IsClosed D) : IsClosed (binaryHull C D)
\end{leansrc}
\par\medskip\noindent{\small \texttt{SuFunctor.Universe.}\texttt{\bfseries sublevel\_qCoproduct\_eq}}\hfill{\footnotesize \path{lean/su-functor/SuFunctor/Universe/HullClosure.lean}, line 686}\par\nopagebreak
\begin{leansrc}
theorem sublevel_qCoproduct_eq {σ τ' : ProbabilityMeasure X → ℝ≥0∞}
    (hσ : LowerSemicontinuous σ) (hτ' : LowerSemicontinuous τ') (c : ℝ≥0∞) :
    {μ | qCoproduct σ τ' μ ≤ c}
      = binaryHull {ν | σ ν ≤ c} {ν | τ' ν ≤ c}
\end{leansrc}
\par\medskip\noindent{\small \texttt{SuFunctor.Universe.}\texttt{\bfseries qCoproduct\_le\_min}}\hfill{\footnotesize \path{lean/su-functor/SuFunctor/Universe/HullClosure.lean}, line 663}\par\nopagebreak
\begin{leansrc}
theorem qCoproduct_le_min (σ τ' : ProbabilityMeasure X → ℝ≥0∞)
    (μ : ProbabilityMeasure X) : qCoproduct σ τ' μ ≤ min (σ μ) (τ' μ)
\end{leansrc}
\par\medskip\noindent{\small \texttt{SuFunctor.Universe.}\texttt{\bfseries le\_qCoproduct\_of\_forall\_mixConvex}}\hfill{\footnotesize \path{lean/su-functor/SuFunctor/Universe/HullClosure.lean}, line 670}\par\nopagebreak
\begin{leansrc}
theorem le_qCoproduct_of_forall_mixConvex {σ τ' h : ProbabilityMeasure X → ℝ≥0∞}
    (hqc : ∀ c, MixConvex {ν | h ν ≤ c}) (h1 : ∀ ν, h ν ≤ σ ν)
    (h2 : ∀ ν, h ν ≤ τ' ν) (μ : ProbabilityMeasure X) :
    h μ ≤ qCoproduct σ τ' μ
\end{leansrc}
\par\medskip\noindent{\small \texttt{SuFunctor.Universe.}\texttt{\bfseries boxClass\_union\_ssubset}}\hfill{\footnotesize \path{lean/su-functor/SuFunctor/Universe/DisjunctionElim.lean}, line 275}\par\nopagebreak
\begin{leansrc}
theorem boxClass_union_ssubset {x y : α} (hxy : x ≠ y) :
    boxClass ({x} : Set α) ∪ boxClass {y} ⊂ boxClass ({x, y} : Set α)
\end{leansrc}
\par\medskip\noindent{\small \texttt{SuFunctor.Universe.}\texttt{\bfseries exists\_ctxOrE\_countermodel\_body\_reading\_carrier}}\hfill{\footnotesize \path{lean/su-functor/SuFunctor/Universe/DisjunctionElim.lean}, line 311}\par\nopagebreak
\begin{leansrc}
theorem exists_ctxOrE_countermodel_body_reading_carrier :
    ∃ γ ψ : Set (NonemptyCompacts (BelState Bool)),
      γ ∩ boxClass {BelState.dirac true} ⊆ ψ
        ∧ γ ∩ boxClass {BelState.dirac false} ⊆ ψ
        ∧ ¬ γ ∩ boxClass
            ({BelState.dirac true, BelState.dirac false}
              : Set (BelState Bool)) ⊆ ψ
\end{leansrc}
\par\medskip\noindent{\small \texttt{SuFunctor.Universe.}\texttt{\bfseries mem\_dnfClass\_append\_iff}}\hfill{\footnotesize \path{lean/su-functor/SuFunctor/Universe/DisjunctionElim.lean}, line 175}\par\nopagebreak
\begin{leansrc}
theorem mem_dnfClass_append_iff {n m : ℕ} {u : Fin n → ProbPenalty V}
    {𝒱 : Fin n → Set (ProbPenalty V)} {u' : Fin m → ProbPenalty V}
    {𝒱' : Fin m → Set (ProbPenalty V)} {K : NonemptyCompacts (BelState V)} :
    K ∈ dnfClass (Fin.append u u') (Fin.append 𝒱 𝒱')
      ↔ K ∈ dnfClass u 𝒱 ∨ K ∈ dnfClass u' 𝒱'
\end{leansrc}
\par\medskip\noindent{\small \texttt{SuFunctor.Universe.}\texttt{\bfseries coupled\_union\_refinement\_iff}}\hfill{\footnotesize \path{lean/su-functor/SuFunctor/Universe/LocallyClosedLayer.lean}, line 138}\par\nopagebreak
\begin{leansrc}
theorem coupled_union_refinement_iff {P : (K → ℝ≥0∞) → Prop} {u : K → ℝ≥0∞}
    {𝒱 : Set (K → ℝ≥0∞)} {u' : Fin m → K → ℝ≥0∞}
    {𝒱' : Fin m → Set (K → ℝ≥0∞)}
    (hPc : ∀ c, SelValid 𝒱' c → P (selSup u c)) :
    (∀ σ, P σ → CoupledSat σ u 𝒱 → ∃ j, CoupledSat σ (u' j) (𝒱' j))
      ↔ ∀ c, SelValid 𝒱' c → Blocked u 𝒱 u' c
\end{leansrc}
\par\medskip\noindent{\small \texttt{SuFunctor.Universe.}\texttt{\bfseries exists\_naive\_cover\_incomplete}}\hfill{\footnotesize \path{lean/su-functor/SuFunctor/Universe/LocallyClosedLayer.lean}, line 235}\par\nopagebreak
\begin{leansrc}
theorem exists_naive_cover_incomplete :
    ∃ u v : ProbPenalty (Fin 2),
      LowerSemicontinuous u ∧ LowerSemicontinuous v
        ∧ (∀ σ, CoupledSat σ u {u ⊔ v} ↔ CoupledSat σ u {v})
        ∧ ¬ u ⊔ v ≤ v
\end{leansrc}
\par\medskip\noindent{\small \texttt{SuFunctor.Universe.}\texttt{\bfseries exists\_essential\_floor\_raise}}\hfill{\footnotesize \path{lean/su-functor/SuFunctor/Universe/LocallyClosedLayer.lean}, line 294}\par\nopagebreak
\begin{leansrc}
theorem exists_essential_floor_raise :
    ∃ u₁ w : ProbPenalty (Fin 2),
      LowerSemicontinuous u₁ ∧ LowerSemicontinuous w
        ∧ (∀ σ : ProbPenalty (Fin 2), CoupledSat σ 0 {(⊤ : ProbPenalty (Fin 2))} →
            CoupledSat σ u₁ {(⊤ : ProbPenalty (Fin 2))} ∨ CoupledSat σ 0 {w})
        ∧ ¬ (∀ σ : ProbPenalty (Fin 2), LowerSemicontinuous σ →
              CoupledSat σ 0 {(⊤ : ProbPenalty (Fin 2))} →
              CoupledSat σ u₁ {(⊤ : ProbPenalty (Fin 2))})
        ∧ ¬ (∀ σ : ProbPenalty (Fin 2), LowerSemicontinuous σ →
              CoupledSat σ 0 {(⊤ : ProbPenalty (Fin 2))} →
              CoupledSat σ 0 {w})
        ∧ ¬ u₁ ≤ (0 : ProbPenalty (Fin 2))
\end{leansrc}
\par\medskip\noindent{\small \texttt{SuFunctor.Universe.}\texttt{\bfseries le\_infProj\_iff}}\hfill{\footnotesize \path{lean/su-functor/SuFunctor/Universe/Marginalization.lean}, line 79}\par\nopagebreak
\begin{leansrc}
theorem le_infProj_iff {p : Z → W} {σ : Z → ℝ≥0∞} {φ : W → ℝ≥0∞} :
    φ ≤ infProj p σ ↔ φ ∘ p ≤ σ
\end{leansrc}
\par\medskip\noindent{\small \texttt{SuFunctor.Universe.}\texttt{\bfseries le\_hide\_iff}}\hfill{\footnotesize \path{lean/su-functor/SuFunctor/Universe/Marginalization.lean}, line 195}\par\nopagebreak
\begin{leansrc}
theorem le_hide_iff {p : Z → W} {σ : Z → ℝ≥0∞} {φ : W → ℝ≥0∞}
    (hφ : LowerSemicontinuous φ) :
    φ ≤ hide p σ ↔ φ ∘ p ≤ σ
\end{leansrc}
\par\medskip\noindent{\small \texttt{SuFunctor.Universe.}\texttt{\bfseries hide\_sup\_le}}\hfill{\footnotesize \path{lean/su-functor/SuFunctor/Universe/HideFrobenius.lean}, line 156}\par\nopagebreak
\begin{leansrc}
theorem hide_sup_le (p : Z → W) (σ : Z → ℝ≥0∞) {g : W → ℝ≥0∞}
    (hg : LowerSemicontinuous g) :
    hide p σ ⊔ g ≤ hide p (σ ⊔ g ∘ p)
\end{leansrc}
\par\medskip\noindent{\small \texttt{SuFunctor.Universe.}\texttt{\bfseries hide\_add\_le}}\hfill{\footnotesize \path{lean/su-functor/SuFunctor/Universe/HideFrobenius.lean}, line 168}\par\nopagebreak
\begin{leansrc}
theorem hide_add_le (p : Z → W) (σ : Z → ℝ≥0∞) {g : W → ℝ≥0∞}
    (hg : LowerSemicontinuous g) :
    hide p σ + g ≤ hide p (σ + g ∘ p)
\end{leansrc}
\par\medskip\noindent{\small \texttt{SuFunctor.Universe.}\texttt{\bfseries exists\_lowerSemicontinuous\_hide\_add\_ne}}\hfill{\footnotesize \path{lean/su-functor/SuFunctor/Universe/HideFrobenius.lean}, line 309}\par\nopagebreak
\begin{leansrc}
theorem exists_lowerSemicontinuous_hide_add_ne :
    ∃ (σ : ProbabilityMeasure ℕ → ℝ≥0∞) (g : ℝ≥0∞ → ℝ≥0∞),
      LowerSemicontinuous σ ∧ LowerSemicontinuous g ∧
      hide liveMass (σ + g ∘ liveMass) ≠ hide liveMass σ + g
\end{leansrc}
\par\medskip\noindent{\small \texttt{SuFunctor.Universe.}\texttt{\bfseries lscEnv\_charFn}}\hfill{\footnotesize \path{lean/su-functor/SuFunctor/Universe/HideFrobenius.lean}, line 111}\par\nopagebreak
\begin{leansrc}
theorem lscEnv_charFn (C : Set Z) : lscEnv (charFn C) = charFn (closure C)
\end{leansrc}
\par\medskip\noindent{\small \texttt{SuFunctor.Universe.}\texttt{\bfseries exists\_judgment\_separating\_hide\_add}}\hfill{\footnotesize \path{lean/su-functor/SuFunctor/Universe/HideFrobenius.lean}, line 347}\par\nopagebreak
\begin{leansrc}
theorem exists_judgment_separating_hide_add :
    ∃ (σ : ProbabilityMeasure ℕ → ℝ≥0∞) (g φ : ℝ≥0∞ → ℝ≥0∞),
      LowerSemicontinuous σ ∧ LowerSemicontinuous g ∧
      LowerSemicontinuous φ ∧
      φ ≤ hide liveMass (σ + g ∘ liveMass) ∧
      ¬ φ ≤ hide liveMass σ + g
\end{leansrc}
\par\medskip\noindent{\small \texttt{SuFunctor.Universe.}\texttt{\bfseries hide\_sup\_comp}}\hfill{\footnotesize \path{lean/su-functor/SuFunctor/Universe/HideFrobenius.lean}, line 189}\par\nopagebreak
\begin{leansrc}
theorem hide_sup_comp [T2Space W] {p : Z → W} (hp : Continuous p)
    {σ : Z → ℝ≥0∞} (hσ : LowerSemicontinuous σ)
    (hcpt : ∀ c : ℝ≥0∞, c ≠ ⊤ → IsCompact {z | σ z ≤ c})
    {g : W → ℝ≥0∞} (hg : LowerSemicontinuous g) :
    hide p (σ ⊔ g ∘ p) = hide p σ ⊔ g
\end{leansrc}
\par\medskip\noindent{\small \texttt{SuFunctor.Universe.}\texttt{\bfseries hide\_add\_comp}}\hfill{\footnotesize \path{lean/su-functor/SuFunctor/Universe/HideFrobenius.lean}, line 209}\par\nopagebreak
\begin{leansrc}
theorem hide_add_comp [T2Space W] {p : Z → W} (hp : Continuous p)
    {σ : Z → ℝ≥0∞} (hσ : LowerSemicontinuous σ)
    (hcpt : ∀ c : ℝ≥0∞, c ≠ ⊤ → IsCompact {z | σ z ≤ c})
    {g : W → ℝ≥0∞} (hg : LowerSemicontinuous g) :
    hide p (σ + g ∘ p) = hide p σ + g
\end{leansrc}
\par\medskip\noindent{\small \texttt{Qtsl.Imprecise.}\texttt{\bfseries hurwicz\_ge\_iff\_exists\_chord}}\hfill{\footnotesize \path{lean/qtsl-aux/QtslAux/Imprecise/HurwiczCrescents.lean}, line 67}\par\nopagebreak
\begin{leansrc}
theorem hurwicz_ge_iff_exists_chord (α : ℚ) (hα0 : 0 ≤ α) (hα1 : α ≤ 1)
    (M : ι → FinPMF X) (f : X → ℚ) (lo : ℚ) :
    lo ≤ hurwicz α M f
      ↔ ∃ a b : ℚ, lo ≤ α * a + (1 - α) * b
          ∧ a ≤ upperE M f ∧ b ≤ lowerE M f
\end{leansrc}
\par\medskip\noindent{\small \texttt{Qtsl.Imprecise.}\texttt{\bfseries hurwiczSuperlevel\_eq\_iUnion\_crescents}}\hfill{\footnotesize \path{lean/qtsl-aux/QtslAux/Imprecise/HurwiczCrescents.lean}, line 81}\par\nopagebreak
\begin{leansrc}
theorem hurwiczSuperlevel_eq_iUnion_crescents (α : ℚ) (hα0 : 0 ≤ α)
    (hα1 : α ≤ 1) (M : ι → FinPMF X) (lo : ℚ) :
    {f : X → ℚ | lo ≤ hurwicz α M f}
      = ⋃ p : {q : ℚ × ℚ // lo ≤ α * q.1 + (1 - α) * q.2},
          ({f : X → ℚ | p.val.1 ≤ upperE M f}
            ∩ {f : X → ℚ | p.val.2 ≤ lowerE M f})
\end{leansrc}
\par\medskip\noindent{\small \texttt{Qtsl.Imprecise.}\texttt{\bfseries finite\_chords\_strict}}\hfill{\footnotesize \path{lean/qtsl-aux/QtslAux/Imprecise/HurwiczCrescents.lean}, line 112}\par\nopagebreak
\begin{leansrc}
theorem finite_chords_strict {α lo : ℚ} (hα0 : 0 < α) (hα1 : α < 1)
    (S : Finset (ℚ × ℚ)) (hS : ∀ p ∈ S, lo ≤ α * p.1 + (1 - α) * p.2) :
    ∃ (M : Bool → FinPMF Bool) (f : Bool → ℚ),
      lo ≤ hurwicz α M f ∧
        ∀ p ∈ S, ¬(p.2 ≤ lowerE M f ∧ p.1 ≤ upperE M f)
\end{leansrc}
\par\medskip\noindent{\small \texttt{Qtsl.Imprecise.}\texttt{\bfseries le\_hurwicz\_of\_chord}}\hfill{\footnotesize \path{lean/qtsl-aux/QtslAux/Imprecise/HurwiczCrescents.lean}, line 54}\par\nopagebreak
\begin{leansrc}
theorem le_hurwicz_of_chord {α lo a b : ℚ} (hα0 : 0 ≤ α) (hα1 : α ≤ 1)
    {M : ι → FinPMF X} {f : X → ℚ}
    (hchord : lo ≤ α * a + (1 - α) * b)
    (ha : a ≤ upperE M f) (hb : b ≤ lowerE M f) :
    lo ≤ hurwicz α M f
\end{leansrc}
\par\medskip\noindent{\small \texttt{Qtsl.Imprecise.}\texttt{\bfseries penUpperE}}\hfill{\footnotesize \path{lean/qtsl-aux/QtslAux/Imprecise/PenaltyEnvelope.lean}, line 100}\par\nopagebreak
\begin{leansrc}
def penUpperE (M : ι → FinPMF X) (c : ι → ℚ) (f : X → ℚ) : ℚ
\end{leansrc}
\par\medskip\noindent{\small \texttt{Qtsl.Imprecise.}\texttt{\bfseries penLowerE}}\hfill{\footnotesize \path{lean/qtsl-aux/QtslAux/Imprecise/PenaltyEnvelope.lean}, line 105}\par\nopagebreak
\begin{leansrc}
def penLowerE (M : ι → FinPMF X) (c : ι → ℚ) (f : X → ℚ) : ℚ
\end{leansrc}
\par\medskip\noindent{\small \texttt{Qtsl.Imprecise.}\texttt{\bfseries penUpperE\_zero\_cost}}\hfill{\footnotesize \path{lean/qtsl-aux/QtslAux/Imprecise/PenaltyEnvelope.lean}, line 138}\par\nopagebreak
\begin{leansrc}
theorem penUpperE_zero_cost (M : ι → FinPMF X) (f : X → ℚ) :
    penUpperE M 0 f = upperE M f
\end{leansrc}
\par\medskip\noindent{\small \texttt{Qtsl.Imprecise.}\texttt{\bfseries penUpperE\_eq\_sup\_levelProfile}}\hfill{\footnotesize \path{lean/qtsl-aux/QtslAux/Imprecise/PenaltyEnvelope.lean}, line 332}\par\nopagebreak
\begin{leansrc}
theorem penUpperE_eq_sup_levelProfile (M : ι → FinPMF X) (c : ι → ℚ)
    (f : X → ℚ) :
    penUpperE M c f
      = univ.sup' univ_nonempty (fun i₀ => levelProfile M c f i₀ - c i₀)
\end{leansrc}
\par\medskip\noindent{\small \texttt{Qtsl.Imprecise.}\texttt{\bfseries penUpperE\_compFam}}\hfill{\footnotesize \path{lean/qtsl-aux/QtslAux/Imprecise/PenaltyKernelCompose.lean}, line 71}\par\nopagebreak
\begin{leansrc}
theorem penUpperE_compFam (M : ι → FinPMF X) (c : ι → ℚ)
    (K : (x : X) → κ x → FinPMF Y) (d : (x : X) → κ x → ℚ)
    (f : X × Y → ℚ) :
    penUpperE (compFam M K) (penCompCost M c d) f
      = penUpperE M c fun x => penUpperE (K x) (d x) fun y => f (x, y)
\end{leansrc}
\par\medskip\noindent{\small \texttt{Qtsl.Imprecise.}\texttt{\bfseries penUpperE\_mix\_le}}\hfill{\footnotesize \path{lean/qtsl-aux/QtslAux/Imprecise/PenaltyEnvelope.lean}, line 219}\par\nopagebreak
\begin{leansrc}
theorem penUpperE_mix_le (M : ι → FinPMF X) (c : ι → ℚ) (f g : X → ℚ)
    {t : ℚ} (h₀ : 0 ≤ t) (h₁ : t ≤ 1) :
    penUpperE M c (fun x => t * f x + (1 - t) * g x)
      ≤ t * penUpperE M c f + (1 - t) * penUpperE M c g
\end{leansrc}
\par\medskip\noindent{\small \texttt{Qtsl.Imprecise.Mutants.}\texttt{\bfseries penUpperE\_not\_posHomogeneous}}\hfill{\footnotesize \path{lean/qtsl-aux/QtslAux/Imprecise/PenaltyEnvelope.lean}, line 362}\par\nopagebreak
\begin{leansrc}
theorem penUpperE_not_posHomogeneous :
    penUpperE (fun _ : Unit => FinPMF.dirac (default : Unit)) (fun _ => 1)
        (fun x => (2 : ℚ) * (fun _ : Unit => (0 : ℚ)) x)
      ≠ 2 * penUpperE (fun _ : Unit => FinPMF.dirac (default : Unit))
          (fun _ => 1) (fun _ : Unit => (0 : ℚ))
\end{leansrc}
\par\medskip\noindent{\small \texttt{Qtsl.Imprecise.}\texttt{\bfseries penUpperE\_dirac\_family}}\hfill{\footnotesize \path{lean/qtsl-aux/QtslAux/Imprecise/PenaltyEnvelope.lean}, line 292}\par\nopagebreak
\begin{leansrc}
theorem penUpperE_dirac_family (w : M → ℚ) (g : M → ℚ) :
    penUpperE (fun m => FinPMF.dirac m) w g
      = univ.sup' univ_nonempty (fun m => g m - w m)
\end{leansrc}
\par\medskip\noindent{\small \texttt{SuFunctor.Plain.}\texttt{\bfseries MeasJoint}}\hfill{\footnotesize \path{lean/su-functor/SuFunctor/Plain/MeasJoint.lean}, line 110}\par\nopagebreak
\begin{leansrc}
def MeasJoint (X : Type u) [MeasurableSpace X] : Type u
\end{leansrc}
\par\medskip\noindent{\small \texttt{SuFunctor.Plain.Joint.}\texttt{\bfseries toMeas\_push}}\hfill{\footnotesize \path{lean/su-functor/SuFunctor/Plain/GeneralSeam.lean}, line 131}\par\nopagebreak
\begin{leansrc}
theorem toMeas_push (h : X → Y) (hm : Measurable h) (μ : Joint X) :
    toMeas (μ.push h) = (toMeas μ).push h
\end{leansrc}
\par\medskip\noindent{\small \texttt{SuFunctor.Site.}\texttt{\bfseries penExtM}}\hfill{\footnotesize \path{lean/su-functor/SuFunctor/Site/GeneralWindow.lean}, line 86}\par\nopagebreak
\begin{leansrc}
abbrev penExtM {V W : Set T} (hVW : V ⊆ W) {X : Type u} [MeasurableSpace X] :
    MeasPen (Traj V X) → MeasPen (Traj W X)
\end{leansrc}
\par\medskip\noindent{\small \texttt{SuFunctor.Site.}\texttt{\bfseries penExtM\_comp\_toMeas}}\hfill{\footnotesize \path{lean/su-functor/SuFunctor/Site/GeneralWindow.lean}, line 265}\par\nopagebreak
\begin{leansrc}
theorem penExtM_comp_toMeas {V W : Set T} (hVW : V ⊆ W) {X : Type u}
    [MeasurableSpace X] (F : MeasPen (Traj V X)) :
    (penExtM hVW F) ∘ Joint.toMeas = penExt hVW (F ∘ Joint.toMeas)
\end{leansrc}
\par\medskip\noindent{\small \texttt{SuFunctor.Site.IsPersistentM.}\texttt{\bfseries seamRestrict}}\hfill{\footnotesize \path{lean/su-functor/SuFunctor/Site/GeneralWindow.lean}, line 273}\par\nopagebreak
\begin{leansrc}
theorem IsPersistentM.seamRestrict {U : ι → Set T} {X : Type u}
    [MeasurableSpace X] {Φ : ∀ i, Set (MeasPen (Traj (U i) X))}
    (hΦ : IsPersistentM U Φ) :
    IsPersistent U (fun i => (fun F => F ∘ Joint.toMeas) '' Φ i)
\end{leansrc}
\par\medskip\noindent{\small \texttt{SuFunctor.Plain.Joint.}\texttt{\bfseries toMeas\_injective}}\hfill{\footnotesize \path{lean/su-functor/SuFunctor/Plain/GeneralSeam.lean}, line 119}\par\nopagebreak
\begin{leansrc}
theorem toMeas_injective [MeasurableSingletonClass X] :
    Function.Injective (toMeas : Joint X → MeasJoint X)
\end{leansrc}
\par\medskip\noindent{\small \texttt{SuFunctor.Site.}\texttt{\bfseries penExt}}\hfill{\footnotesize \path{lean/su-functor/SuFunctor/Site/WindowPenalty.lean}, line 76}\par\nopagebreak
\begin{leansrc}
abbrev penExt {V W : Set T} (hVW : V ⊆ W) {X : Type u} :
    Pen (Traj V X) → Pen (Traj W X)
\end{leansrc}
\par\medskip\noindent{\small \texttt{SuFunctor.Plain.Joint.}\texttt{\bfseries dense\_ratSeam}}\hfill{\footnotesize \path{lean/su-functor/SuFunctor/Plain/SeamDensity.lean}, line 299}\par\nopagebreak
\begin{leansrc}
theorem Joint.dense_ratSeam
    (hd : DenseRange (Joint.toMeas : Joint X → MeasJoint X)) :
    Dense (Joint.toMeas '' {ν : Joint X | ν.IsRatWeighted})
\end{leansrc}
\par\medskip\noindent{\small \texttt{SuFunctor.Plain.Joint.}\texttt{\bfseries denseRange\_toMeas\_nat}}\hfill{\footnotesize \path{lean/su-functor/SuFunctor/Plain/SeamDensity.lean}, line 443}\par\nopagebreak
\begin{leansrc}
theorem Joint.denseRange_toMeas_nat :
    DenseRange (Joint.toMeas : Joint ℕ → MeasJoint ℕ)
\end{leansrc}
\par\medskip\noindent{\small \texttt{SuFunctor.Plain.MeasPen.}\texttt{\bfseries le\_of\_seam\_le}}\hfill{\footnotesize \path{lean/su-functor/SuFunctor/Plain/SeamDensity.lean}, line 149}\par\nopagebreak
\begin{leansrc}
theorem MeasPen.le_of_seam_le
    (hd : DenseRange (Joint.toMeas : Joint X → MeasJoint X))
    {F G : MeasPen X} (hF : LowerSemicontinuous F)
    (hG : UpperSemicontinuous G)
    (h : ∀ ν : Joint X, F ν.toMeas ≤ G ν.toMeas) : F ≤ G
\end{leansrc}
\par\medskip\noindent{\small \texttt{SuFunctor.Plain.MeasPen.}\texttt{\bfseries iSup\_seam\_eq}}\hfill{\footnotesize \path{lean/su-functor/SuFunctor/Plain/SeamDensity.lean}, line 176}\par\nopagebreak
\begin{leansrc}
theorem MeasPen.iSup_seam_eq
    (hd : DenseRange (Joint.toMeas : Joint X → MeasJoint X))
    {F : MeasPen X} (hF : LowerSemicontinuous F) :
    ⨆ ν : Joint X, F ν.toMeas = ⨆ μ : MeasJoint X, F μ
\end{leansrc}
\par\medskip\noindent{\small \texttt{SuFunctor.Plain.MeasPen.}\texttt{\bfseries eq\_of\_seam\_eq}}\hfill{\footnotesize \path{lean/su-functor/SuFunctor/Plain/SeamDensity.lean}, line 166}\par\nopagebreak
\begin{leansrc}
theorem MeasPen.eq_of_seam_eq
    (hd : DenseRange (Joint.toMeas : Joint X → MeasJoint X))
    {F G : MeasPen X} (hF : Continuous F) (hG : Continuous G)
    (h : ∀ ν : Joint X, F ν.toMeas = G ν.toMeas) : F = G
\end{leansrc}
\par\medskip\noindent{\small \texttt{SuFunctor.Plain.}\texttt{\bfseries dense\_seam\_inf\_gap}}\hfill{\footnotesize \path{lean/su-functor/SuFunctor/Plain/SeamDensity.lean}, line 635}\par\nopagebreak
\begin{leansrc}
theorem dense_seam_inf_gap :
    DenseRange (Joint.toMeas : Joint ℕ → MeasJoint ℕ) ∧
      ∃ F : MeasPen ℕ, LowerSemicontinuous F ∧
        (∀ ν : Joint ℕ, 1 ≤ F ν.toMeas) ∧
        (⨅ ν : Joint ℕ, F ν.toMeas) ≠ ⨅ μ : MeasJoint ℕ, F μ
\end{leansrc}
\par\medskip\noindent{\small \texttt{SuFunctor.Site.}\texttt{\bfseries IsMatching}}\hfill{\footnotesize \path{lean/su-functor/SuFunctor/Site/WindowPenalty.lean}, line 241}\par\nopagebreak
\begin{leansrc}
def IsMatching (U : ι → Set T) {X : Type u}
    (Φ : ∀ i, Set (Pen (Traj (U i) X))) : Prop
\end{leansrc}
\par\medskip\noindent{\small \texttt{SuFunctor.Site.}\texttt{\bfseries IsPersistent}}\hfill{\footnotesize \path{lean/su-functor/SuFunctor/Site/WindowPenalty.lean}, line 234}\par\nopagebreak
\begin{leansrc}
def IsPersistent (U : ι → Set T) {X : Type u}
    (Φ : ∀ i, Set (Pen (Traj (U i) X))) : Prop
\end{leansrc}
\par\medskip\noindent{\small \texttt{SuFunctor.Site.IsPersistent.}\texttt{\bfseries image2\_boxplusW}}\hfill{\footnotesize \path{lean/su-functor/SuFunctor/Site/WindowPenalty.lean}, line 269}\par\nopagebreak
\begin{leansrc}
theorem IsPersistent.image2_boxplusW {U : ι → Set T}
    {A B C : WireSpan El} {M : Pat A B} {N : Pat B C}
    {Φ : ∀ i, Set (Pen (Traj (U i) (Pat.Tot M)))}
    {Ψ : ∀ i, Set (Pen (Traj (U i) (Pat.Tot N)))}
    (hΦ : IsPersistent U Φ) (hΨ : IsPersistent U Ψ) :
    IsPersistent U (fun i => Set.image2 boxplusW (Φ i) (Ψ i))
\end{leansrc}
\par\medskip\noindent{\small \texttt{SuFunctor.Site.}\texttt{\bfseries classDescent}}\hfill{\footnotesize \path{lean/su-functor/SuFunctor/Site/ClassDescent.lean}, line 243}\par\nopagebreak
\begin{leansrc}
def classDescent (hUo : ∀ i, IsOpen (U i)) (hdir : Directed (· ⊆ ·) U)
    (hcov : ⋃ i, U i = univ) :
    Set (FellCloseds α) ≃o Set (ProjLim U) where …
\end{leansrc}
\par\medskip\noindent{\small \texttt{SuFunctor.Site.}\texttt{\bfseries isConstructibleUnion\_classDescent\_iff}}\hfill{\footnotesize \path{lean/su-functor/SuFunctor/Site/ClassDescent.lean}, line 268}\par\nopagebreak
\begin{leansrc}
theorem isConstructibleUnion_classDescent_iff (hUo : ∀ i, IsOpen (U i))
    (hdir : Directed (· ⊆ ·) U) (hcov : ⋃ i, U i = univ)
    {Φ : Set (FellCloseds α)} :
    IsConstructibleUnion (classDescent hUo hdir hcov Φ)
      ↔ IsConstructibleUnion Φ
\end{leansrc}
\par\medskip\noindent{\small \texttt{SuFunctor.Site.IsLocallyClosed.}\texttt{\bfseries classExt}}\hfill{\footnotesize \path{lean/su-functor/SuFunctor/Site/ClassDescent.lean}, line 167}\par\nopagebreak
\begin{leansrc}
theorem IsLocallyClosed.classExt (hV : IsOpen V) (hVW : V ⊆ W)
    {Φ : Set (FellCloseds ↥V)} (hΦ : IsLocallyClosed Φ) :
    IsLocallyClosed (classExt hVW Φ)
\end{leansrc}
\par\medskip\noindent{\small \texttt{SuFunctor.Site.IsConstructibleUnion.}\texttt{\bfseries classExt}}\hfill{\footnotesize \path{lean/su-functor/SuFunctor/Site/ClassDescent.lean}, line 174}\par\nopagebreak
\begin{leansrc}
theorem IsConstructibleUnion.classExt (hV : IsOpen V) (hVW : V ⊆ W)
    {Φ : Set (FellCloseds ↥V)} (hΦ : IsConstructibleUnion Φ) :
    IsConstructibleUnion (Site.classExt hVW Φ)
\end{leansrc}
\par\medskip\noindent{\small \texttt{SuFunctor.Site.}\texttt{\bfseries bijective\_toProjLim}}\hfill{\footnotesize \path{lean/su-functor/SuFunctor/Site/ClassDescent.lean}, line 281}\par\nopagebreak
\begin{leansrc}
theorem bijective_toProjLim (hUo : ∀ i, IsOpen (U i))
    (hdir : Directed (· ⊆ ·) U) (hcov : ⋃ i, U i = univ) :
    Function.Bijective (toProjLim U)
\end{leansrc}
\par\medskip\noindent{\small \texttt{SuFunctor.Site.}\texttt{\bfseries clip\_mem\_classExt\_iff}}\hfill{\footnotesize \path{lean/su-functor/SuFunctor/Site/ClassDescent.lean}, line 159}\par\nopagebreak
\begin{leansrc}
theorem clip_mem_classExt_iff (hVW : V ⊆ W) {Φ : Set (FellCloseds ↥V)}
    (F : FellCloseds α) :
    clip W F ∈ classExt hVW Φ ↔ clip V F ∈ Φ
\end{leansrc}
\par\medskip\noindent{\small \texttt{SuFunctor.}\texttt{\bfseries SingleSiteSemanticUniverse}}\hfill{\footnotesize \path{lean/su-functor/SuFunctor/SingleSiteSemanticUniverse.lean}, line 489}\par\nopagebreak
\begin{leansrc}
abbrev SingleSiteSemanticUniverse (El : Type u)
\end{leansrc}
\par\medskip\noindent{\small \texttt{SuFunctor.QtslSem.}\texttt{\bfseries SpecJudgLayer}}\hfill{\footnotesize \path{lean/su-functor/SuFunctor/QtslSem/SpecJudg.lean}, line 108}\par\nopagebreak
\begin{leansrc}
abbrev SpecJudgLayer (El : Type u)
\end{leansrc}
\par\medskip\noindent{\small \texttt{SuFunctor.QtslSem.}\texttt{\bfseries cellOfSubset}}\hfill{\footnotesize \path{lean/su-functor/SuFunctor/QtslSem/SpecJudg.lean}, line 128}\par\nopagebreak
\begin{leansrc}
def cellOfSubset {El : Type u} {a b : SpecJudgLayer El}
    {M : Pat a.un b.un} {𝒰 𝒱 : Cpl M} (h : 𝒰 ⊆ 𝒱) :
    Cell (C := SpecJudgLayer El)
      (⟨M, 𝒰⟩ : JudgLayer.JLoose (specJudgOrder El) a b) ⟨M, 𝒱⟩
      (𝟙 a) (𝟙 b)
\end{leansrc}
\par\medskip\noindent{\small \texttt{SuFunctor.Site.}\texttt{\bfseries wComp\_mono}}\hfill{\footnotesize \path{lean/su-functor/SuFunctor/Site/Battery.lean}, line 127}\par\nopagebreak
\begin{leansrc}
theorem wComp_mono {M : Pat A B} {N : Pat B C} {𝒰 𝒰' : WCpl U M}
    {𝒱 𝒱' : WCpl U N} (h₁ : ∀ i, 𝒰.1 i ⊆ 𝒰'.1 i)
    (h₂ : ∀ i, 𝒱.1 i ⊆ 𝒱'.1 i) :
    ∀ i, (wComp U 𝒰 𝒱).1 i ⊆ (wComp U 𝒰' 𝒱').1 i
\end{leansrc}
\par\medskip\noindent{\small \texttt{SuFunctor.Site.}\texttt{\bfseries wUnit\_transparent}}\hfill{\footnotesize \path{lean/su-functor/SuFunctor/Site/Battery.lean}, line 87}\par\nopagebreak
\begin{leansrc}
theorem wUnit_transparent {M : Pat A B} (𝒰 : WCpl U M) :
    wCellMap U (PseudoDoubleCategoryStruct.lLeftUnitor (C := WireSpan El) M)
      (wComp U (wUnit U A) 𝒰) = 𝒰
\end{leansrc}
\par\medskip\noindent{\small \texttt{CategoryTheory.}\texttt{\bfseries PseudoDoubleCategoryStruct}}\hfill{\footnotesize \path{lean/double-cat/DoubleCat/Basic.lean}, line 101}\par\nopagebreak
\begin{leansrc}
class PseudoDoubleCategoryStruct (C : Type u) [Category.{v} C] where …
\end{leansrc}
\par\medskip\noindent{\small \texttt{CategoryTheory.}\texttt{\bfseries PseudoDoubleCategory}}\hfill{\footnotesize \path{lean/double-cat/DoubleCat/Basic.lean}, line 322}\par\nopagebreak
\begin{leansrc}
class PseudoDoubleCategory (C : Type u) [Category.{v} C]
    [PseudoDoubleCategoryStruct C] : Prop where …
\end{leansrc}
\par\medskip\noindent{\small \texttt{CategoryTheory.PseudoDoubleCategoryStruct.Cell.}\texttt{\bfseries cast}}\hfill{\footnotesize \path{lean/double-cat/DoubleCat/Basic.lean}, line 178}\par\nopagebreak
\begin{leansrc}
def Cell.cast {a b c d : C} {m : Loose a b} {n : Loose c d} {f f' : a ⟶ c}
    {g g' : b ⟶ d} (hf : f = f') (hg : g = g') (α : Cell m n f g) :
    Cell m n f' g'
\end{leansrc}
\par\medskip\noindent{\small \texttt{CategoryTheory.}\texttt{\bfseries LaxDoubleFunctor}}\hfill{\footnotesize \path{lean/double-cat/DoubleCat/Functor/Lax.lean}, line 48}\par\nopagebreak
\begin{leansrc}
structure LaxDoubleFunctor extends C ⥤ D where …
\end{leansrc}
\par\medskip\noindent{\small \texttt{SuFunctor.Plain.MeasJoint.}\texttt{\bfseries push\_id}}\hfill{\footnotesize \path{lean/su-functor/SuFunctor/Plain/MeasJoint.lean}, line 156}\par\nopagebreak
\begin{leansrc}
@[simp] theorem push_id (μ : MeasJoint X) : μ.push id = μ
\end{leansrc}
\par\medskip\noindent{\small \texttt{SuFunctor.Plain.MeasJoint.}\texttt{\bfseries push\_push}}\hfill{\footnotesize \path{lean/su-functor/SuFunctor/Plain/MeasJoint.lean}, line 163}\par\nopagebreak
\begin{leansrc}
theorem push_push (h₁ : X → Y) (h₂ : Y → Z) (hm₁ : Measurable h₁)
    (hm₂ : Measurable h₂) (μ : MeasJoint X) :
    (μ.push h₁).push h₂ = μ.push (h₂ ∘ h₁)
\end{leansrc}
\par\medskip\noindent{\small \texttt{SuFunctor.Plain.MeasJoint.}\texttt{\bfseries mass\_push}}\hfill{\footnotesize \path{lean/su-functor/SuFunctor/Plain/MeasJoint.lean}, line 170}\par\nopagebreak
\begin{leansrc}
theorem mass_push (h : X → Y) (hm : Measurable h) (μ : MeasJoint X) :
    (μ.push h).1.mass = μ.1.mass
\end{leansrc}
\par\medskip\noindent{\small \texttt{SuFunctor.Plain.}\texttt{\bfseries continuous\_push}}\hfill{\footnotesize \path{lean/su-functor/SuFunctor/Plain/MeasJoint.lean}, line 336}\par\nopagebreak
\begin{leansrc}
theorem continuous_push {h : X → Y} (h_cont : Continuous h) :
    Continuous (fun μ : MeasJoint X => μ.push h)
\end{leansrc}
\par\medskip\noindent{\small \texttt{SuFunctor.Plain.MeasJoint.}\texttt{\bfseries push\_toMeasure}}\hfill{\footnotesize \path{lean/su-functor/SuFunctor/Plain/MeasJoint.lean}, line 153}\par\nopagebreak
\begin{leansrc}
@[simp] theorem push_toMeasure (h : X → Y) (μ : MeasJoint X) :
    ((μ.push h).1 : Measure Y) = (μ.1 : Measure X).map h
\end{leansrc}
\par\medskip\noindent{\small \texttt{SuFunctor.Universe.}\texttt{\bfseries exists\_lowerSemicontinuous\_not\_lowerSemicontinuous\_infProj}}\hfill{\footnotesize \path{lean/su-functor/SuFunctor/Universe/Marginalization.lean}, line 558}\par\nopagebreak
\begin{leansrc}
theorem exists_lowerSemicontinuous_not_lowerSemicontinuous_infProj :
    ∃ σ : ProbabilityMeasure ℕ → ℝ≥0∞, LowerSemicontinuous σ ∧
      ¬ LowerSemicontinuous (infProj liveMass σ)
\end{leansrc}
\par\medskip\noindent{\small \texttt{SuFunctor.Universe.}\texttt{\bfseries le\_lscEnv\_iff}}\hfill{\footnotesize \path{lean/su-functor/SuFunctor/Universe/Marginalization.lean}, line 173}\par\nopagebreak
\begin{leansrc}
theorem le_lscEnv_iff {g φ : Z → ℝ≥0∞} (hφ : LowerSemicontinuous φ) :
    φ ≤ lscEnv g ↔ φ ≤ g
\end{leansrc}
\par\medskip\noindent{\small \texttt{SuFunctor.Universe.}\texttt{\bfseries isCompact\_sublevel\_of\_isTightMeasureSet}}\hfill{\footnotesize \path{lean/su-functor/SuFunctor/Universe/Marginalization.lean}, line 317}\par\nopagebreak
\begin{leansrc}
theorem isCompact_sublevel_of_isTightMeasureSet
    {σ : ProbabilityMeasure 𝓧 → ℝ≥0∞} (hσ : LowerSemicontinuous σ)
    {c : ℝ≥0∞}
    (ht : IsTightMeasureSet
      {((μ : ProbabilityMeasure 𝓧) : Measure 𝓧) | μ ∈ {ν | σ ν ≤ c}}) :
    IsCompact {ν | σ ν ≤ c}
\end{leansrc}
\par\medskip\noindent{\small \texttt{SuFunctor.Universe.}\texttt{\bfseries sublevel\_infProj\_eq}}\hfill{\footnotesize \path{lean/su-functor/SuFunctor/Universe/Marginalization.lean}, line 213}\par\nopagebreak
\begin{leansrc}
theorem sublevel_infProj_eq [T1Space W] {p : Z → W} (hp : Continuous p)
    {σ : Z → ℝ≥0∞} (hσ : LowerSemicontinuous σ)
    (hcpt : ∀ c : ℝ≥0∞, c ≠ ⊤ → IsCompact {z | σ z ≤ c})
    {c : ℝ≥0∞} (hc : c ≠ ⊤) :
    {w | infProj p σ w ≤ c} = p '' {z | σ z ≤ c}
\end{leansrc}
\par\medskip\noindent{\small \texttt{SuFunctor.Universe.}\texttt{\bfseries exists\_infProj\_attained}}\hfill{\footnotesize \path{lean/su-functor/SuFunctor/Universe/Marginalization.lean}, line 270}\par\nopagebreak
\begin{leansrc}
theorem exists_infProj_attained [T1Space W] {p : Z → W} (hp : Continuous p)
    {σ : Z → ℝ≥0∞} (hσ : LowerSemicontinuous σ)
    (hcpt : ∀ c : ℝ≥0∞, c ≠ ⊤ → IsCompact {z | σ z ≤ c})
    {ν : W} (hν : infProj p σ ν ≠ ⊤) :
    ∃ z, p z = ν ∧ σ z = infProj p σ ν
\end{leansrc}
\par\medskip\noindent{\small \texttt{SuFunctor.Universe.}\texttt{\bfseries isCompact\_sublevel\_infProj}}\hfill{\footnotesize \path{lean/su-functor/SuFunctor/Universe/Marginalization.lean}, line 283}\par\nopagebreak
\begin{leansrc}
theorem isCompact_sublevel_infProj [T1Space W] {p : Z → W} (hp : Continuous p)
    {σ : Z → ℝ≥0∞} (hσ : LowerSemicontinuous σ)
    (hcpt : ∀ c : ℝ≥0∞, c ≠ ⊤ → IsCompact {z | σ z ≤ c})
    {c : ℝ≥0∞} (hc : c ≠ ⊤) :
    IsCompact {w | infProj p σ w ≤ c}
\end{leansrc}
\par\medskip\noindent{\small \texttt{SuFunctor.Universe.}\texttt{\bfseries lowerSemicontinuous\_infProj}}\hfill{\footnotesize \path{lean/su-functor/SuFunctor/Universe/Marginalization.lean}, line 292}\par\nopagebreak
\begin{leansrc}
theorem lowerSemicontinuous_infProj [T2Space W] {p : Z → W} (hp : Continuous p)
    {σ : Z → ℝ≥0∞} (hσ : LowerSemicontinuous σ)
    (hcpt : ∀ c : ℝ≥0∞, c ≠ ⊤ → IsCompact {z | σ z ≤ c}) :
    LowerSemicontinuous (infProj p σ)
\end{leansrc}
\par\medskip\noindent{\small \texttt{SuFunctor.Universe.}\texttt{\bfseries sublevel\_infProj\_margFst\_eq}}\hfill{\footnotesize \path{lean/su-functor/SuFunctor/Universe/Marginalization.lean}, line 347}\par\nopagebreak
\begin{leansrc}
theorem sublevel_infProj_margFst_eq
    {σ : ProbabilityMeasure (𝓧 × 𝓨) → ℝ≥0∞} (hσ : LowerSemicontinuous σ)
    (ht : ∀ c : ℝ≥0∞, c ≠ ⊤ → IsTightMeasureSet
      {((ρ : ProbabilityMeasure (𝓧 × 𝓨)) : Measure (𝓧 × 𝓨)) | ρ ∈ {ν | σ ν ≤ c}})
    {c : ℝ≥0∞} (hc : c ≠ ⊤) :
    {w | infProj margFst σ w ≤ c} = margFst '' {ρ | σ ρ ≤ c}
\end{leansrc}
\par\medskip\noindent{\small \texttt{SuFunctor.Plain.Pen.}\texttt{\bfseries pull}}\hfill{\footnotesize \path{lean/su-functor/SuFunctor/Plain/Penalty.lean}, line 100}\par\nopagebreak
\begin{leansrc}
def pull (h : X → Y) (f : Pen Y) : Pen X
\end{leansrc}
\par\medskip\noindent{\small \texttt{SuFunctor.Site.}\texttt{\bfseries siteSpecFunctor}}\hfill{\footnotesize \path{lean/su-functor/SuFunctor/Site/SiteFunctor.lean}, line 128}\par\nopagebreak
\begin{leansrc}
def siteSpecFunctor (U : ι → Set T) (El : Type u) :
    SpanValued.{u + 1, u, u, u + 1, u} (WireSpan El) where …
\end{leansrc}
\par\medskip\noindent{\small \texttt{SuFunctor.Plain.}\texttt{\bfseries measSpecObj}}\hfill{\footnotesize \path{lean/su-functor/SuFunctor/Plain/GeneralFunctor.lean}, line 56}\par\nopagebreak
\begin{leansrc}
def measSpecObj : MeasCat.{u}ᵒᵖ ⥤ Type u where …
\end{leansrc}
\par\medskip\noindent{\small \texttt{SuFunctor.Site.}\texttt{\bfseries SiteSpecTotal}}\hfill{\footnotesize \path{lean/su-functor/SuFunctor/Site/SiteFunctor.lean}, line 279}\par\nopagebreak
\begin{leansrc}
abbrev SiteSpecTotal (U : ι → Set T) (El : Type u)
\end{leansrc}
\par\medskip\noindent{\small \texttt{SuFunctor.Site.}\texttt{\bfseries siteSpecTotal\_pseudoDoubleCategory}}\hfill{\footnotesize \path{lean/su-functor/SuFunctor/Site/SiteFunctor.lean}, line 285}\par\nopagebreak
\begin{leansrc}
theorem siteSpecTotal_pseudoDoubleCategory (U : ι → Set T) (El : Type u) :
    ∃ _ : PseudoDoubleCategory (SiteSpecTotal U El), True
\end{leansrc}
\par\medskip\noindent{\small \texttt{SuFunctor.Site.}\texttt{\bfseries siteSpecProj\_isPseudo}}\hfill{\footnotesize \path{lean/su-functor/SuFunctor/Site/SiteFunctor.lean}, line 295}\par\nopagebreak
\begin{leansrc}
theorem siteSpecProj_isPseudo (U : ι → Set T) (El : Type u) :
    (siteSpecProj U El).IsPseudo
\end{leansrc}
\par\medskip\noindent{\small \texttt{SuFunctor.Site.}\texttt{\bfseries siteSpecProj\_discreteOpfibration}}\hfill{\footnotesize \path{lean/su-functor/SuFunctor/Site/SiteFunctor.lean}, line 304}\par\nopagebreak
\begin{leansrc}
theorem siteSpecProj_discreteOpfibration (U : ι → Set T) (El : Type u) :
    IsDiscreteDoubleOpfibration (SpanValued.proj (siteSpecFunctor U El))
\end{leansrc}
\par\medskip\noindent{\small \texttt{SuFunctor.}\texttt{\bfseries SingleSiteSemanticUniverse}}\hfill{\footnotesize \path{lean/su-functor/SuFunctor/SingleSiteSemanticUniverse.lean}, line 489}\par\nopagebreak
\begin{leansrc}
abbrev SingleSiteSemanticUniverse (El : Type u)
\end{leansrc}
\par\medskip\noindent{\small \texttt{SuFunctor.}\texttt{\bfseries dnfSpecFunctor}}\hfill{\footnotesize \path{lean/su-functor/SuFunctor/SingleSiteSemanticUniverse.lean}, line 365}\par\nopagebreak
\begin{leansrc}
def dnfSpecFunctor (El : Type u) :
    SpanValued.{u + 1, u, u, u + 1, u} (WireSpan El) where …
\end{leansrc}
\par\medskip\noindent{\small \texttt{SuFunctor.}\texttt{\bfseries singleSiteUniverseProj}}\hfill{\footnotesize \path{lean/su-functor/SuFunctor/SingleSiteSemanticUniverse.lean}, line 500}\par\nopagebreak
\begin{leansrc}
def singleSiteUniverseProj (El : Type u) :
    LaxDoubleFunctor (SingleSiteSemanticUniverse El) (WireSpan El)
\end{leansrc}
\par\medskip\noindent{\small \texttt{SuFunctor.}\texttt{\bfseries singleSiteUniverseProj\_isPseudo}}\hfill{\footnotesize \path{lean/su-functor/SuFunctor/SingleSiteSemanticUniverse.lean}, line 505}\par\nopagebreak
\begin{leansrc}
theorem singleSiteUniverseProj_isPseudo (El : Type u) :
    (singleSiteUniverseProj El).IsPseudo
\end{leansrc}
\par\medskip\noindent{\small \texttt{SuFunctor.}\texttt{\bfseries singleSiteUniverseProj\_discreteOpfibration}}\hfill{\footnotesize \path{lean/su-functor/SuFunctor/SingleSiteSemanticUniverse.lean}, line 514}\par\nopagebreak
\begin{leansrc}
theorem singleSiteUniverseProj_discreteOpfibration (El : Type u) :
    IsDiscreteDoubleOpfibration (singleSiteUniverseProj El)
\end{leansrc}
\par\medskip\noindent{\small \texttt{SuFunctor.}\texttt{\bfseries specMap\_modelClass\_eq}}\hfill{\footnotesize \path{lean/su-functor/SuFunctor/SingleSiteSemanticUniverse.lean}, line 533}\par\nopagebreak
\begin{leansrc}
theorem specMap_modelClass_eq (φ : A ⟶ B)
    (hφ : Function.Surjective (res φ)) (D : DnfSpec A) :
    specMap φ D.modelClass
      = (dnfPull (res φ) D).modelClass ∩ Set.range (Pen.pull (res φ))
\end{leansrc}
\par\medskip\noindent{\small \texttt{SuFunctor.}\texttt{\bfseries specMap\_modelClass\_subset}}\hfill{\footnotesize \path{lean/su-functor/SuFunctor/SingleSiteSemanticUniverse.lean}, line 540}\par\nopagebreak
\begin{leansrc}
theorem specMap_modelClass_subset (φ : A ⟶ B)
    (hφ : Function.Surjective (res φ)) (D : DnfSpec A) :
    specMap φ D.modelClass ⊆ (dnfPull (res φ) D).modelClass
\end{leansrc}
\par\medskip\noindent{\small \texttt{SuFunctor.Universe.SpecDnf.}\texttt{\bfseries isConstructibleUnion\_probRead}}\hfill{\footnotesize \path{lean/su-functor/SuFunctor/SingleSiteSemanticUniverse.lean}, line 594}\par\nopagebreak
\begin{leansrc}
theorem Universe.SpecDnf.isConstructibleUnion_probRead
    (D : SpecDnf (Joint V))
    (hfloor : ∀ a ∈ D.probRead.atoms, LowerSemicontinuous a.floor)
    (hmarks : ∀ a ∈ D.probRead.atoms, ∀ v ∈ a.marks, LowerSemicontinuous v) :
    IsConstructibleUnion D.probRead.crescentClass
\end{leansrc}
\par\medskip\noindent{\small \texttt{SuFunctor.Site.}\texttt{\bfseries singleton\_boxplusW}}\hfill{\footnotesize \path{lean/su-functor/SuFunctor/Site/Battery.lean}, line 103}\par\nopagebreak
\begin{leansrc}
theorem singleton_boxplusW {W : Set T} {M : Pat A B} {N : Pat B C}
    (u₁ : Pen (Traj W (Pat.Tot M))) (u₂ : Pen (Traj W (Pat.Tot N))) :
    Set.image2 boxplusW ({u₁} : Set (Pen (Traj W (Pat.Tot M))))
        ({u₂} : Set (Pen (Traj W (Pat.Tot N)))) =
      {boxplusW u₁ u₂}
\end{leansrc}
\par\medskip\noindent{\small \texttt{SuFunctor.Site.}\texttt{\bfseries transport\_singletonW}}\hfill{\footnotesize \path{lean/su-functor/SuFunctor/Site/Battery.lean}, line 113}\par\nopagebreak
\begin{leansrc}
theorem transport_singletonW {W : Set T} (φ : A ⟶ B)
    (σ : Pen (Traj W (states A))) :
    Pen.pull (trajMap (res φ)) '' ({σ} : Set (Pen (Traj W (states A)))) =
      {Pen.pull (trajMap (res φ)) σ}
\end{leansrc}
\par\medskip\noindent{\small \texttt{SuFunctor.Site.}\texttt{\bfseries wIdealPoint\_le\_iff}}\hfill{\footnotesize \path{lean/su-functor/SuFunctor/Site/Battery.lean}, line 121}\par\nopagebreak
\begin{leansrc}
theorem wIdealPoint_le_iff {W : Set T} (σ : Pen (Traj W (states A)))
    (S : Set (Pen (Traj W (states A)))) :
    ({σ} : Set (Pen (Traj W (states A)))) ⊆ S ↔ σ ∈ S
\end{leansrc}
\par\medskip\noindent{\small \texttt{SuFunctor.Universe.}\texttt{\bfseries exists\_coupledSat\_add\_fails}}\hfill{\footnotesize \path{lean/su-functor/SuFunctor/Universe/SpecDnf.lean}, line 420}\par\nopagebreak
\begin{leansrc}
theorem exists_coupledSat_add_fails :
    ∃ (σ₁ σ₂ u₁ u₂ : Unit → ℝ≥0∞) (𝒱₁ 𝒱₂ : Set (Unit → ℝ≥0∞)),
      ⊤ ∈ 𝒱₁ ∧ ⊤ ∈ 𝒱₂ ∧
      CoupledSat σ₁ u₁ 𝒱₁ ∧ CoupledSat σ₂ u₂ 𝒱₂ ∧
      ¬ CoupledSat (σ₁ + σ₂) (u₁ + u₂) (𝒱₁ ∪ 𝒱₂)
\end{leansrc}
\par\medskip\noindent{\small \texttt{SuFunctor.Universe.}\texttt{\bfseries add\_comp\_safety\_le}}\hfill{\footnotesize \path{lean/su-functor/SuFunctor/Universe/SpecDnf.lean}, line 405}\par\nopagebreak
\begin{leansrc}
theorem add_comp_safety_le (g₁ : K → K₁) (g₂ : K → K₂)
    {u₁ σ₁ : K₁ → ℝ≥0∞} {u₂ σ₂ : K₂ → ℝ≥0∞}
    (h₁ : u₁ ≤ σ₁) (h₂ : u₂ ≤ σ₂) :
    (u₁ ∘ g₁) + (u₂ ∘ g₂) ≤ (σ₁ ∘ g₁) + (σ₂ ∘ g₂)
\end{leansrc}
\par\medskip\noindent{\small \texttt{CategoryTheory.ImplDouble.}\texttt{\bfseries cell\_subsingleton}}\hfill{\footnotesize \path{lean/double-cat/DoubleCat/ImplDouble.lean}, line 68}\par\nopagebreak
\begin{leansrc}
theorem cell_subsingleton {a b c d : ImplDouble C} (m : Loose a b)
    (n : Loose c d) (f : a ⟶ c) (g : b ⟶ d) :
    Subsingleton (Cell m n f g)
\end{leansrc}
\par\medskip\noindent{\small \texttt{SuFunctor.Site.}\texttt{\bfseries mem\_wComp\_iff}}\hfill{\footnotesize \path{lean/su-functor/SuFunctor/Site/Battery.lean}, line 69}\par\nopagebreak
\begin{leansrc}
theorem mem_wComp_iff {M : Pat A B} {N : Pat B C} {𝒰 : WCpl U M}
    {𝒱 : WCpl U N} {i : ι}
    {w : Pen (Traj (U i) (Pat.Tot (Pat.comp M N)))} :
    w ∈ (wComp U 𝒰 𝒱).1 i ↔
      ∃ u₁ ∈ 𝒰.1 i, ∃ u₂ ∈ 𝒱.1 i, boxplusW u₁ u₂ = w
\end{leansrc}
\par\medskip\noindent{\small \texttt{SuFunctor.Site.}\texttt{\bfseries wComp\_nonempty}}\hfill{\footnotesize \path{lean/su-functor/SuFunctor/Site/Battery.lean}, line 78}\par\nopagebreak
\begin{leansrc}
theorem wComp_nonempty {M : Pat A B} {N : Pat B C} {𝒰 : WCpl U M}
    {𝒱 : WCpl U N} (h₁ : ∀ i, (𝒰.1 i).Nonempty)
    (h₂ : ∀ i, (𝒱.1 i).Nonempty) :
    ∀ i, ((wComp U 𝒰 𝒱).1 i).Nonempty
\end{leansrc}
\par\medskip\noindent{\small \texttt{SuFunctor.Site.}\texttt{\bfseries taxedBoxplusW\_breaks\_unit\_transparency}}\hfill{\footnotesize \path{lean/su-functor/SuFunctor/Site/Battery.lean}, line 152}\par\nopagebreak
\begin{leansrc}
theorem taxedBoxplusW_breaks_unit_transparency {W : Set T} (M : Pat A B)
    {c : ℝ≥0∞} (hc : c ≠ 0) :
    ¬ ∀ u : Pen (Traj W (Pat.Tot M)),
      Pen.pull (trajMap (PatCell.tot
          (PseudoDoubleCategoryStruct.lLeftUnitor (C := WireSpan El) M)))
          (taxedBoxplusW c (unitPenW A W) u) = u
\end{leansrc}
\par\medskip\noindent{\small \texttt{SuFunctor.Site.}\texttt{\bfseries funnelCompW\_breaks\_unit\_coherence}}\hfill{\footnotesize \path{lean/su-functor/SuFunctor/Site/Battery.lean}, line 176}\par\nopagebreak
\begin{leansrc}
theorem funnelCompW_breaks_unit_coherence {W : Set T} (M : Pat A B) :
    ¬ ∀ 𝒰 : Set (Pen (Traj W (Pat.Tot M))),
      Pen.pull (trajMap (PatCell.tot
          (PseudoDoubleCategoryStruct.lLeftUnitor (C := WireSpan El) M))) ''
          funnelCompW ({unitPenW A W} :
            Set (Pen (Traj W (Pat.Tot (Pat.id A))))) 𝒰 = 𝒰
\end{leansrc}
\par\medskip\noindent{\small \texttt{Qtsl.Imprecise.}\texttt{\bfseries sublevel\_compositeG\_subset\_glueSetG}}\hfill{\footnotesize \path{lean/qtsl-aux/QtslAux/Imprecise/ContainmentGeneral.lean}, line 111}\par\nopagebreak
\begin{leansrc}
theorem sublevel_compositeG_subset_glueSetG (m : ∀ i, J → Mi i)
    (f : ∀ i, Mi i → ENNRat) (F : ENNRat) :
    {ν | compositePenaltyG m f ν ≤ F} ⊆ glueSetG m f F
\end{leansrc}
\par\medskip\noindent{\small \texttt{Qtsl.Imprecise.}\texttt{\bfseries certificate\_transfer}}\hfill{\footnotesize \path{lean/qtsl-aux/QtslAux/Imprecise/ContainmentGeneral.lean}, line 121}\par\nopagebreak
\begin{leansrc}
theorem certificate_transfer (m : ∀ i, J → Mi i) (f : ∀ i, Mi i → ENNRat)
    (F : ENNRat) {P : J → Prop} (h : ∀ ν ∈ glueSetG m f F, P ν) :
    ∀ ν, compositePenaltyG m f ν ≤ F → P ν
\end{leansrc}
\par\medskip\noindent{\small \texttt{Qtsl.Imprecise.}\texttt{\bfseries sublevel\_pmComposite\_subset\_glue}}\hfill{\footnotesize \path{lean/qtsl-aux/QtslAux/Imprecise/ContainmentGeneral.lean}, line 307}\par\nopagebreak
\begin{leansrc}
theorem sublevel_pmComposite_subset_glue {proj : ∀ i, X → Xi i}
    (hproj : ∀ i, Continuous (proj i))
    (f : ∀ i, ProbabilityMeasure (Xi i) → ENNRat) (F : ENNRat) :
    {ν | pmCompositePenalty proj (fun i => (hproj i).measurable) f ν ≤ F}
      ⊆ glueSetG (pmMarginal proj fun i => (hproj i).measurable) f F
\end{leansrc}
\par\medskip\noindent{\small \texttt{Qtsl.Imprecise.}\texttt{\bfseries isClosed\_pmGlueSet}}\hfill{\footnotesize \path{lean/qtsl-aux/QtslAux/Imprecise/ContainmentGeneral.lean}, line 300}\par\nopagebreak
\begin{leansrc}
theorem isClosed_pmGlueSet {proj : ∀ i, X → Xi i}
    (hproj : ∀ i, Continuous (proj i))
    {f : ∀ i, ProbabilityMeasure (Xi i) → ENNRat}
    (hf : ∀ i, LowerSemicontinuous (f i)) (F : ENNRat) :
    IsClosed (glueSetG (pmMarginal proj fun i => (hproj i).measurable) f F)
\end{leansrc}
\par\medskip\noindent{\small \texttt{Qtsl.Imprecise.}\texttt{\bfseries forall\_sublevel\_eq\_iff}}\hfill{\footnotesize \path{lean/qtsl-aux/QtslAux/Imprecise/ContainmentGeneral.lean}, line 162}\par\nopagebreak
\begin{leansrc}
theorem forall_sublevel_eq_iff {g h : J → ENNRat} :
    (∀ F, {ν | g ν ≤ F} = {ν | h ν ≤ F}) ↔ g = h
\end{leansrc}
\par\medskip\noindent{\small \texttt{Qtsl.Imprecise.ENNRat.}\texttt{\bfseries add\_eq\_max\_iff}}\hfill{\footnotesize \path{lean/qtsl-aux/QtslAux/Imprecise/ContainmentGeneral.lean}, line 81}\par\nopagebreak
\begin{leansrc}
theorem ENNRat.add_eq_max_iff (a b : ENNRat) :
    a + b = max a b ↔ min a b = 0 ∨ max a b = ⊤
\end{leansrc}
\par\medskip\noindent{\small \texttt{Qtsl.Imprecise.}\texttt{\bfseries binaryComposite\_ne\_glue}}\hfill{\footnotesize \path{lean/qtsl-aux/QtslAux/Imprecise/ContainmentGeneral.lean}, line 191}\par\nopagebreak
\begin{leansrc}
theorem binaryComposite_ne_glue :
    (fun ν : Unit => (fun _ : Unit => (1 : ENNRat)) ν
        + (fun _ : Unit => (1 : ENNRat)) ν)
      ≠ fun ν : Unit => max ((fun _ : Unit => (1 : ENNRat)) ν)
          ((fun _ : Unit => (1 : ENNRat)) ν)
\end{leansrc}
\par\medskip\noindent{\small \texttt{Qtsl.Imprecise.}\texttt{\bfseries sublevel\_ssubset\_glue\_witness}}\hfill{\footnotesize \path{lean/qtsl-aux/QtslAux/Imprecise/ContainmentGeneral.lean}, line 206}\par\nopagebreak
\begin{leansrc}
theorem sublevel_ssubset_glue_witness :
    {ν : Unit | compositePenaltyG (fun _ : Bool => id)
        (fun _ : Bool => fun _ : Unit => (1 : ENNRat)) ν ≤ 1}
      ⊂ glueSetG (fun _ : Bool => id)
        (fun _ : Bool => fun _ : Unit => (1 : ENNRat)) 1
\end{leansrc}
\par\medskip\noindent{\small \texttt{Qtsl.Imprecise.}\texttt{\bfseries signed\_relaxation\_fails}}\hfill{\footnotesize \path{lean/qtsl-aux/QtslAux/Imprecise/MarginalCompat.lean}, line 179}\par\nopagebreak
\begin{leansrc}
theorem signed_relaxation_fails :
    ∃ f : Fin 2 → FinPMF (Fin 2) → ℚ,
      ¬ ∀ (i : Fin 2) (ν : FinPMF (Fin 2)),
          f i ν ≤ compositePenaltyQ (Xi := fun _ => Fin 2)
            (fun _ ν => ν) f ν
\end{leansrc}
\par\medskip\noindent{\small \texttt{Qtsl.Imprecise.}\texttt{\bfseries subunit\_weight\_fails}}\hfill{\footnotesize \path{lean/qtsl-aux/QtslAux/Imprecise/MarginalCompat.lean}, line 189}\par\nopagebreak
\begin{leansrc}
theorem subunit_weight_fails :
    ∃ (c : Fin 1 → ℚ) (f : Fin 1 → FinPMF (Fin 2) → ℚ),
      (∀ i ν, 0 ≤ f i ν) ∧ (∀ i, 0 ≤ c i) ∧
      ¬ ∀ (i : Fin 1) (ν : FinPMF (Fin 2)),
          f i ν ≤ compositePenaltyWQ (Xi := fun _ => Fin 2) c
            (fun _ ν => ν) f ν
\end{leansrc}
\par\medskip\noindent{\small \texttt{SuFunctor.Plain.}\texttt{\bfseries jointEquivBel}}\hfill{\footnotesize \path{lean/su-functor/SuFunctor/Plain/Seam.lean}, line 100}\par\nopagebreak
\begin{leansrc}
def jointEquivBel : Joint V ≃ Universe.BelState V where …
\end{leansrc}
\par\medskip\noindent{\small \texttt{SuFunctor.Plain.}\texttt{\bfseries jointEquivProb}}\hfill{\footnotesize \path{lean/su-functor/SuFunctor/Plain/Seam.lean}, line 147}\par\nopagebreak
\begin{leansrc}
def jointEquivProb :
    {μ : Joint V // μ.wt.sum (fun _ v => v) = 1} ≃ Universe.ProbState V where …
\end{leansrc}
\par\medskip\noindent{\small \texttt{SuFunctor.Plain.}\texttt{\bfseries penEquiv}}\hfill{\footnotesize \path{lean/su-functor/SuFunctor/Plain/Seam.lean}, line 200}\par\nopagebreak
\begin{leansrc}
def penEquiv : Pen V ≃o (Universe.BelState V → ℝ≥0∞) where …
\end{leansrc}
\par\medskip\noindent{\small \texttt{SuFunctor.Plain.Joint.}\texttt{\bfseries toMeas}}\hfill{\footnotesize \path{lean/su-functor/SuFunctor/Plain/GeneralSeam.lean}, line 76}\par\nopagebreak
\begin{leansrc}
def toMeas (μ : Joint X) : MeasJoint X
\end{leansrc}
\par\medskip\noindent{\small \texttt{SuFunctor.Plain.Joint.}\texttt{\bfseries pull\_toMeas}}\hfill{\footnotesize \path{lean/su-functor/SuFunctor/Plain/GeneralSeam.lean}, line 150}\par\nopagebreak
\begin{leansrc}
theorem pull_toMeas (h : X → Y) (hm : Measurable h) (F : MeasPen Y) :
    (MeasPen.pull h F) ∘ toMeas = Pen.pull h (F ∘ toMeas)
\end{leansrc}
\par\medskip\noindent{\small \texttt{SuFunctor.Plain.Joint.}\texttt{\bfseries image\_pull\_toMeas}}\hfill{\footnotesize \path{lean/su-functor/SuFunctor/Plain/GeneralSeam.lean}, line 159}\par\nopagebreak
\begin{leansrc}
theorem image_pull_toMeas (h : X → Y) (hm : Measurable h)
    (𝒮 : Set (MeasPen Y)) :
    (fun F => F ∘ toMeas) '' (MeasPen.pull h '' 𝒮) =
      Pen.pull h '' ((fun F => F ∘ toMeas) '' 𝒮)
\end{leansrc}
\par\medskip\noindent{\small \texttt{SuFunctor.Universe.General.}\texttt{\bfseries coupledSat\_pullPen\_iff}}\hfill{\footnotesize \path{lean/su-functor/SuFunctor/Universe/General/TransportGeneral.lean}, line 239}\par\nopagebreak
\begin{leansrc}
theorem coupledSat_pullPen_iff (hθ : Continuous θ)
    (hsurj : Function.Surjective (probPush hθ))
    {σ u : Penalty Y} {𝒱 : Set (Penalty Y)} :
    CoupledSat (pullPen hθ σ) (pullPen hθ u) (pullPen hθ '' 𝒱)
      ↔ CoupledSat σ u 𝒱
\end{leansrc}
\par\medskip\noindent{\small \texttt{SuFunctor.Universe.General.}\texttt{\bfseries transportClass\_crescentClass}}\hfill{\footnotesize \path{lean/su-functor/SuFunctor/Universe/General/TransportGeneral.lean}, line 418}\par\nopagebreak
\begin{leansrc}
theorem transportClass_crescentClass [Nonempty X] (hθ : Continuous θ)
    (hsurj : Function.Surjective
      (push θ : SubProbabilityMeasure X → SubProbabilityMeasure Y))
    (hmap : IsClosedMap
      (push θ : SubProbabilityMeasure X → SubProbabilityMeasure Y))
    (u : Penalty Y) (𝒱 : Set (Penalty Y)) :
    transportClass θ (crescentClass u 𝒱)
      = crescentClass (pullPen hθ u) (pullPen hθ '' 𝒱)
        ∩ {L | Saturated θ (L : Set (SubProbabilityMeasure X))}
\end{leansrc}
\par\medskip\noindent{\small \texttt{SuFunctor.Universe.General.}\texttt{\bfseries exists\_transport\_not\_subset\_pulled\_crescent}}\hfill{\footnotesize \path{lean/su-functor/SuFunctor/Universe/General/TransportGeneral.lean}, line 665}\par\nopagebreak
\begin{leansrc}
theorem exists_transport_not_subset_pulled_crescent :
    ∃ (θ : Unit → Bool) (hθ : Continuous θ) (u : Penalty Bool)
      (𝒱 : Set (Penalty Bool)),
      LowerSemicontinuous u ∧ (∀ v ∈ 𝒱, LowerSemicontinuous v) ∧ 𝒱.Finite ∧
      ¬ (transportClass θ (crescentClass u 𝒱)
          ⊆ crescentClass (pullPen hθ u) (pullPen hθ '' 𝒱))
\end{leansrc}
\par\medskip\noindent{\small \texttt{SuFunctor.Universe.General.}\texttt{\bfseries compile\_add\_radProd}}\hfill{\footnotesize \path{lean/su-functor/SuFunctor/Universe/General/TransportGeneral.lean}, line 478}\par\nopagebreak
\begin{leansrc}
theorem compile_add_radProd [Nonempty X] (f g : Penalty X) :
    compile (f + g) = radProd (compile f) (compile g)
\end{leansrc}
\par\medskip\noindent{\small \texttt{SuFunctor.Universe.General.}\texttt{\bfseries probPush\_surjective\_of\_section}}\hfill{\footnotesize \path{lean/su-functor/SuFunctor/Universe/General/TransportGeneral.lean}, line 208}\par\nopagebreak
\begin{leansrc}
theorem probPush_surjective_of_section (hθ : Continuous θ) {s : Y → X}
    (hs_mble : Measurable s) (hs : Function.RightInverse s θ) :
    Function.Surjective (probPush hθ)
\end{leansrc}
\par\medskip\noindent{\small \texttt{SuFunctor.Universe.}\texttt{\bfseries cross\_site\_union\_refinement\_iff}}\hfill{\footnotesize \path{lean/su-functor/SuFunctor/Universe/DemonicClock.lean}, line 262}\par\nopagebreak
\begin{leansrc}
theorem cross_site_union_refinement_iff {P : (K → ℝ≥0∞) → Prop}
    {u w : K → ℝ≥0∞} {𝒱 : Set (K → ℝ≥0∞)} {u' : Fin m → K → ℝ≥0∞}
    {𝒱' : Fin m → Set (K → ℝ≥0∞)}
    (hPc : ∀ c, SelValid 𝒱' c → P (selSup u c ⊔ w)) :
    (∀ σ, P σ → CoupledSat σ (u ⊔ w) 𝒱 → ∃ j, CoupledSat σ (u' j) (𝒱' j))
      ↔ ∀ c, SelValid 𝒱' c → Blocked (u ⊔ w) 𝒱 u' c
\end{leansrc}
\par\medskip\noindent{\small \texttt{SuFunctor.Universe.}\texttt{\bfseries lowerSemicontinuous\_selSup}}\hfill{\footnotesize \path{lean/su-functor/SuFunctor/Universe/DemonicClock.lean}, line 279}\par\nopagebreak
\begin{leansrc}
theorem lowerSemicontinuous_selSup {u : K → ℝ≥0∞}
    {c : Fin m → Option (K → ℝ≥0∞)} (hu : LowerSemicontinuous u)
    (hc : ∀ j w, c j = some w → LowerSemicontinuous w) :
    LowerSemicontinuous (selSup u c)
\end{leansrc}
\par\medskip\noindent{\small \texttt{SuFunctor.Universe.General.}\texttt{\bfseries probWall}}\hfill{\footnotesize \path{lean/su-functor/SuFunctor/Universe/General/DemonicGeneral.lean}, line 49}\par\nopagebreak
\begin{leansrc}
noncomputable def probWall (B : Set (ProbabilityMeasure X)) : Penalty X
\end{leansrc}
\par\medskip\noindent{\small \texttt{SuFunctor.Universe.}\texttt{\bfseries crescentClass\_inter\_wall}}\hfill{\footnotesize \path{lean/su-functor/SuFunctor/Universe/DemonicClock.lean}, line 323}\par\nopagebreak
\begin{leansrc}
theorem crescentClass_inter_wall (u w : ProbPenalty V)
    (𝒱 : Set (ProbPenalty V)) :
    crescentClass u 𝒱 ∩ crescentClass w ∅ = crescentClass (u ⊔ w) 𝒱
\end{leansrc}
\par\medskip\noindent{\small \texttt{SuFunctor.Universe.}\texttt{\bfseries cross\_site\_union\_refinement\_iff\_probWall}}\hfill{\footnotesize \path{lean/su-functor/SuFunctor/Universe/DemonicClock.lean}, line 332}\par\nopagebreak
\begin{leansrc}
theorem cross_site_union_refinement_iff_probWall {B : Set (ProbState V)}
    (hB : IsClosed B) {u : ProbPenalty V} {𝒱 : Set (ProbPenalty V)}
    {u' : Fin m → ProbPenalty V} {𝒱' : Fin m → Set (ProbPenalty V)}
    (hu : LowerSemicontinuous u)
    (h𝒱' : ∀ j, ∀ v ∈ 𝒱' j, LowerSemicontinuous v) :
    (∀ σ, LowerSemicontinuous σ → CoupledSat σ (u ⊔ probWall B) 𝒱 →
        ∃ j, CoupledSat σ (u' j) (𝒱' j))
      ↔ ∀ c, SelValid 𝒱' c → Blocked (u ⊔ probWall B) 𝒱 u' c
\end{leansrc}
\par\medskip\noindent{\small \texttt{SuFunctor.Universe.}\texttt{\bfseries exists\_no\_functional\_transport}}\hfill{\footnotesize \path{lean/su-functor/SuFunctor/Universe/DemonicClock.lean}, line 186}\par\nopagebreak
\begin{leansrc}
theorem exists_no_functional_transport :
    ∃ r : Bool → Bool → Prop, ¬ ∃ f : Bool → Bool, ∀ x y, r x y ↔ f x = y
\end{leansrc}
\par\medskip\noindent{\small \texttt{SuFunctor.Universe.}\texttt{\bfseries emLift\_comp}}\hfill{\footnotesize \path{lean/su-functor/SuFunctor/Universe/DemonicClock.lean}, line 143}\par\nopagebreak
\begin{leansrc}
theorem emLift_comp (r : X → Y → Prop) (s : Y → Z → Prop)
    (A : Set X) (C : Set Z) :
    emLift (relComp r s) A C ↔ ∃ B, emLift r A B ∧ emLift s B C
\end{leansrc}
\par\medskip\noindent{\small \texttt{SuFunctor.Spacetime.}\texttt{\bfseries spacetimeSpecFunctor\_fibre\_eq}}\hfill{\footnotesize \path{lean/su-functor/SuFunctor/Spacetime/Universe.lean}, line 141}\par\nopagebreak
\begin{leansrc}
theorem spacetimeSpecFunctor_fibre_eq (A : SpacetimeBase Clk Rec) :
    (spacetimeSpecFunctor Clk Rec).toFunctor.obj A
      = SpecDnf (Joint (A.wires → ClockedEl Clk Rec))
\end{leansrc}
\par\medskip\noindent{\small \texttt{SuFunctor.Spacetime.}\texttt{\bfseries profileLiftHom}}\hfill{\footnotesize \path{lean/su-functor/SuFunctor/Spacetime/Base.lean}, line 364}\par\nopagebreak
\begin{leansrc}
def profileLiftHom (S : Type u) (φ : S → B.wires) :
    profileLift (Rec := Rec) S φ ⟶ B
\end{leansrc}
\par\medskip\noindent{\small \texttt{SuFunctor.Spacetime.}\texttt{\bfseries profile\_unique\_of\_hom}}\hfill{\footnotesize \path{lean/su-functor/SuFunctor/Spacetime/Base.lean}, line 373}\par\nopagebreak
\begin{leansrc}
theorem profile_unique_of_hom {S : Type u} {κ : S → Clk}
    {φ : S → B.wires} (h : ∀ w, B.profile (φ w) = κ w) :
    κ = fun w => B.profile (φ w)
\end{leansrc}
\par\medskip\noindent{\small \texttt{SuFunctor.Spacetime.}\texttt{\bfseries mixedInterface}}\hfill{\footnotesize \path{lean/su-functor/SuFunctor/Spacetime/Base.lean}, line 382}\par\nopagebreak
\begin{leansrc}
def mixedInterface (c c' : Clk) : SpacetimeBase Clk Rec
\end{leansrc}
\par\medskip\noindent{\small \texttt{SuFunctor.Spacetime.}\texttt{\bfseries profiledEmbed}}\hfill{\footnotesize \path{lean/su-functor/SuFunctor/Spacetime/Universe.lean}, line 94}\par\nopagebreak
\begin{leansrc}
def profiledEmbed (Clk : Type u) (Rec : Clk → Type u) :
    LaxDoubleFunctor (SpacetimeBase Clk Rec)
      (WireSpan (ClockedEl Clk Rec)) where …
\end{leansrc}
\par\medskip\noindent{\small \texttt{SuFunctor.Spacetime.}\texttt{\bfseries SemanticUniverse}}\hfill{\footnotesize \path{lean/su-functor/SuFunctor/Spacetime/Universe.lean}, line 159}\par\nopagebreak
\begin{leansrc}
abbrev SemanticUniverse (Clk : Type u) (Rec : Clk → Type u)
\end{leansrc}
\par\medskip\noindent{\small \texttt{SuFunctor.Spacetime.}\texttt{\bfseries universeProj\_discreteOpfibration}}\hfill{\footnotesize \path{lean/su-functor/SuFunctor/Spacetime/Universe.lean}, line 185}\par\nopagebreak
\begin{leansrc}
theorem universeProj_discreteOpfibration
    (Clk : Type u) (Rec : Clk → Type u) :
    IsDiscreteDoubleOpfibration (universeProj Clk Rec)
\end{leansrc}
\par\medskip\noindent{\small \texttt{SuFunctor.Spacetime.}\texttt{\bfseries constHomEquiv}}\hfill{\footnotesize \path{lean/su-functor/SuFunctor/Spacetime/Universe.lean}, line 216}\par\nopagebreak
\begin{leansrc}
def constHomEquiv (S T : WireSpan (Rec c)) :
    (S ⟶ T) ≃ (constObj (Rec := Rec) c S.wires ⟶ constObj c T.wires) where …
\end{leansrc}
\par\medskip\noindent{\small \texttt{SuFunctor.Universe.}\texttt{\bfseries fibInf}}\hfill{\footnotesize \path{lean/su-functor/SuFunctor/Universe/CertificateLayer.lean}, line 79}\par\nopagebreak
\begin{leansrc}
def fibInf (g : K₁ → K₂) (τ : K₁ → ℝ≥0∞) : K₂ → ℝ≥0∞
\end{leansrc}
\par\medskip\noindent{\small \texttt{SuFunctor.Universe.}\texttt{\bfseries coupledSat\_comap\_iff\_fibInf}}\hfill{\footnotesize \path{lean/su-functor/SuFunctor/Universe/CertificateLayer.lean}, line 128}\par\nopagebreak
\begin{leansrc}
theorem coupledSat_comap_iff_fibInf {g : K₁ → K₂} {τ : K₁ → ℝ≥0∞}
    {u : K₂ → ℝ≥0∞} {𝒱 : Set (K₂ → ℝ≥0∞)} :
    CoupledSat τ (u ∘ g) ((· ∘ g) '' 𝒱) ↔ CoupledSat (fibInf g τ) u 𝒱
\end{leansrc}
\par\medskip\noindent{\small \texttt{SuFunctor.Universe.SpecDnf.}\texttt{\bfseries comap\_modelClass\_mono}}\hfill{\footnotesize \path{lean/su-functor/SuFunctor/Universe/CertificateLayer.lean}, line 181}\par\nopagebreak
\begin{leansrc}
theorem SpecDnf.comap_modelClass_mono {g : K₁ → K₂} {D E : SpecDnf K₂}
    (h : D.modelClass ⊆ E.modelClass) :
    (D.comap g).modelClass ⊆ (E.comap g).modelClass
\end{leansrc}
\par\medskip\noindent{\small \texttt{SuFunctor.Universe.}\texttt{\bfseries exists\_comp\_notMem\_modelClass\_comap}}\hfill{\footnotesize \path{lean/su-functor/SuFunctor/Universe/CertificateLayer.lean}, line 213}\par\nopagebreak
\begin{leansrc}
theorem exists_comp_notMem_modelClass_comap :
    ∃ (g : Unit → Bool) (a : SpecAtom Bool) (σ : Bool → ℝ≥0∞),
      σ ∈ a.modelClass ∧ σ ∘ g ∉ (a.comap g).modelClass
\end{leansrc}
\par\medskip\noindent{\small \texttt{SuFunctor.Universe.}\texttt{\bfseries coupledSat\_comp\_add\_of\_dips}}\hfill{\footnotesize \path{lean/su-functor/SuFunctor/Universe/CertificateLayer.lean}, line 258}\par\nopagebreak
\begin{leansrc}
theorem coupledSat_comp_add_of_dips (g₁ : K → K₁) (g₂ : K → K₂)
    {σ₁ u₁ : K₁ → ℝ≥0∞} {σ₂ u₂ : K₂ → ℝ≥0∞}
    {𝒱₁ : Set (K₁ → ℝ≥0∞)} {𝒱₂ : Set (K₂ → ℝ≥0∞)}
    (h₁ : CoupledSat σ₁ u₁ 𝒱₁) (h₂ : CoupledSat σ₂ u₂ 𝒱₂)
    (hcon : ∀ w ∈ (· ∘ g₁) '' 𝒱₁ ∪ (· ∘ g₂) '' 𝒱₂,
      Dip ((σ₁ ∘ g₁) + (σ₂ ∘ g₂)) w) :
    CoupledSat ((σ₁ ∘ g₁) + (σ₂ ∘ g₂)) ((u₁ ∘ g₁) + (u₂ ∘ g₂))
      ((· ∘ g₁) '' 𝒱₁ ∪ (· ∘ g₂) '' 𝒱₂)
\end{leansrc}
\par\medskip\noindent{\small \texttt{SuFunctor.Universe.}\texttt{\bfseries exists\_boxplus\_not\_semantically\_monotone}}\hfill{\footnotesize \path{lean/su-functor/SuFunctor/Universe/CertificateLayer.lean}, line 285}\par\nopagebreak
\begin{leansrc}
theorem exists_boxplus_not_semantically_monotone :
    ∃ (a e e₁ e₂ : SpecAtom Unit) (τ : Unit → ℝ≥0∞),
      e.modelClass = e₁.modelClass ∪ e₂.modelClass ∧
      CoupledSat τ (a.floor + e.floor) (a.marks ∪ e.marks) ∧
      ¬ CoupledSat τ (a.floor + e₁.floor) (a.marks ∪ e₁.marks) ∧
      ¬ CoupledSat τ (a.floor + e₂.floor) (a.marks ∪ e₂.marks)
\end{leansrc}
\par\medskip\noindent{\small \texttt{SuFunctor.Universe.SuFunctor.}\texttt{\bfseries dnfBoxplus\_mono}}\hfill{\footnotesize \path{lean/su-functor/SuFunctor/Universe/CertificateLayer.lean}, line 405}\par\nopagebreak
\begin{leansrc}
theorem dnfBoxplus_mono {A B C : WireSpan El} {M : Pat A B} {N : Pat B C}
    {D₁ E₁ : DnfCpl M} {D₂ E₂ : DnfCpl N}
    (h₁ : D₁.atoms ⊆ E₁.atoms) (h₂ : D₂.atoms ⊆ E₂.atoms) :
    (dnfBoxplus D₁ D₂).atoms ⊆ (dnfBoxplus E₁ E₂).atoms
\end{leansrc}
\par\medskip\noindent{\small \texttt{SuFunctor.Spacetime.}\texttt{\bfseries spacetimeSpecFunctor}}\hfill{\footnotesize \path{lean/su-functor/SuFunctor/Spacetime/Universe.lean}, line 134}\par\nopagebreak
\begin{leansrc}
noncomputable def spacetimeSpecFunctor (Clk : Type u) (Rec : Clk → Type u) :
    SpanValued.{u + 1, u, u, u + 1, u} (SpacetimeBase Clk Rec)
\end{leansrc}
\par\medskip\noindent{\small \texttt{SuFunctor.Spacetime.}\texttt{\bfseries semanticUniverse\_pseudoDoubleCategory}}\hfill{\footnotesize \path{lean/su-functor/SuFunctor/Spacetime/Universe.lean}, line 164}\par\nopagebreak
\begin{leansrc}
theorem semanticUniverse_pseudoDoubleCategory
    (Clk : Type u) (Rec : Clk → Type u) :
    ∃ _ : PseudoDoubleCategory (SemanticUniverse Clk Rec), True
\end{leansrc}
\par\medskip\noindent{\small \texttt{SuFunctor.Spacetime.}\texttt{\bfseries universeProj\_isPseudo}}\hfill{\footnotesize \path{lean/su-functor/SuFunctor/Spacetime/Universe.lean}, line 177}\par\nopagebreak
\begin{leansrc}
theorem universeProj_isPseudo (Clk : Type u) (Rec : Clk → Type u) :
    (universeProj Clk Rec).IsPseudo
\end{leansrc}
\par\medskip\noindent{\small \texttt{SuFunctor.Spacetime.}\texttt{\bfseries profiledEmbed\_faithful\_tight}}\hfill{\footnotesize \path{lean/su-functor/SuFunctor/Spacetime/Universe.lean}, line 122}\par\nopagebreak
\begin{leansrc}
theorem profiledEmbed_faithful_tight {A B : SpacetimeBase Clk Rec}
    {φ ψ : A ⟶ B}
    (h : (profiledEmbed Clk Rec).map φ = (profiledEmbed Clk Rec).map ψ) :
    φ = ψ
\end{leansrc}
\par\medskip\noindent{\small \texttt{SuFunctor.Spacetime.}\texttt{\bfseries constProfile}}\hfill{\footnotesize \path{lean/su-functor/SuFunctor/Spacetime/Universe.lean}, line 298}\par\nopagebreak
\begin{leansrc}
def constProfile :
    LaxDoubleFunctor (WireSpan (Rec c)) (SpacetimeBase Clk Rec) where …
\end{leansrc}
\par\medskip\noindent{\small \texttt{SuFunctor.Spacetime.}\texttt{\bfseries constProfile\_isPseudo}}\hfill{\footnotesize \path{lean/su-functor/SuFunctor/Spacetime/Universe.lean}, line 340}\par\nopagebreak
\begin{leansrc}
theorem constProfile_isPseudo :
    (constProfile (Rec := Rec) c).IsPseudo where
  laxComp_isInvertible {_ T _} M N
\end{leansrc}
\par\medskip\noindent{\small \texttt{SuFunctor.Spacetime.}\texttt{\bfseries siteRead}}\hfill{\footnotesize \path{lean/su-functor/SuFunctor/Spacetime/Universe.lean}, line 273}\par\nopagebreak
\begin{leansrc}
noncomputable def siteRead (c : Clk) {S : Type u}
    (D : SpecDnf (Joint (S → ClockedEl Clk Rec))) :
    SpecDnf (Joint (S → Rec c))
\end{leansrc}
\par\medskip\noindent{\small \texttt{SuFunctor.Spacetime.SuFunctor.Universe.SpecDnf.}\texttt{\bfseries mem\_modelClass\_withWall}}\hfill{\footnotesize \path{lean/su-functor/SuFunctor/Spacetime/Universe.lean}, line 408}\par\nopagebreak
\begin{leansrc}
theorem SpecDnf.mem_modelClass_withWall {D : SpecDnf K}
    {w σ : K → ℝ≥0∞} :
    σ ∈ (D.withWall w).modelClass ↔ w ≤ σ ∧ σ ∈ D.modelClass
\end{leansrc}
\par\medskip\noindent{\small \texttt{SuFunctor.Adequacy.}\texttt{\bfseries rwall\_singletonC}}\hfill{\footnotesize \path{lean/su-functor/SuFunctor/Universe/RadialUniverse.lean}, line 303}\par\nopagebreak
\begin{leansrc}
theorem rwall_singletonC (μ : SubPMF V) :
    rwall (singletonC μ) = rayPoint μ
\end{leansrc}
\par\medskip\noindent{\small \texttt{SuFunctor.Adequacy.}\texttt{\bfseries idealPoint}}\hfill{\footnotesize \path{lean/su-functor/SuFunctor/Adequacy/SemanticUniverseWiring.lean}, line 213}\par\nopagebreak
\begin{leansrc}
noncomputable def idealPoint (μ : SubPMF V) : Plain.Pen V
\end{leansrc}
\par\medskip\noindent{\small \texttt{SuFunctor.SubPMF.}\texttt{\bfseries toJoint}}\hfill{\footnotesize \path{lean/su-functor/SuFunctor/Adequacy/SemanticUniverseWiring.lean}, line 99}\par\nopagebreak
\begin{leansrc}
noncomputable def toJoint (μ : SubPMF V) : Plain.Joint V where …
\end{leansrc}
\par\medskip\noindent{\small \texttt{SuFunctor.Adequacy.}\texttt{\bfseries diagramDnf}}\hfill{\footnotesize \path{lean/su-functor/SuFunctor/Adequacy/SemanticUniverseWiring.lean}, line 385}\par\nopagebreak
\begin{leansrc}
noncomputable def diagramDnf (L : List (SubPMF V → Prop)) :
    SpecDnf (Plain.Joint V) where …
\end{leansrc}
\par\medskip\noindent{\small \texttt{SuFunctor.Adequacy.}\texttt{\bfseries coupledSat\_idealPoint\_cwall\_iff}}\hfill{\footnotesize \path{lean/su-functor/SuFunctor/Adequacy/SemanticUniverseWiring.lean}, line 295}\par\nopagebreak
\begin{leansrc}
theorem coupledSat_idealPoint_cwall_iff {c : SubPMF V → Prop}
    {μ : SubPMF V} :
    CoupledSat (idealPoint μ) (cwall c) {(⊤ : Plain.Pen V)} ↔ c μ
\end{leansrc}
\par\medskip\noindent{\small \texttt{SuFunctor.Adequacy.}\texttt{\bfseries idealPoint\_mem\_diagramDnf\_iff}}\hfill{\footnotesize \path{lean/su-functor/SuFunctor/Adequacy/SemanticUniverseWiring.lean}, line 394}\par\nopagebreak
\begin{leansrc}
theorem idealPoint_mem_diagramDnf_iff (L : List (SubPMF V → Prop))
    (μ : SubPMF V) :
    idealPoint μ ∈ (diagramDnf L).modelClass ↔ looseConj L μ
\end{leansrc}
\par\medskip\noindent{\small \texttt{SuFunctor.Plain.}\texttt{\bfseries semanticallyValid\_iff\_subset}}\hfill{\footnotesize \path{lean/su-functor/SuFunctor/Plain/CheckerFace.lean}, line 215}\par\nopagebreak
\begin{leansrc}
theorem semanticallyValid_iff_subset (H : Hypernet) (seq : Sequent) :
    seq.semanticallyValid H ↔ admCyl H seq ⊆ postEvt H seq
\end{leansrc}
\par\medskip\noindent{\small \texttt{SuFunctor.Plain.}\texttt{\bfseries upperProb\_eq\_zero\_iff}}\hfill{\footnotesize \path{lean/su-functor/SuFunctor/Plain/CheckerFace.lean}, line 177}\par\nopagebreak
\begin{leansrc}
theorem upperProb_eq_zero_iff (B E : Set X) :
    upperProb B E = 0 ↔ E ∩ B = ∅
\end{leansrc}
\par\medskip\noindent{\small \texttt{SuFunctor.Plain.}\texttt{\bfseries semanticallyValid\_iff\_face}}\hfill{\footnotesize \path{lean/su-functor/SuFunctor/Plain/CheckerFace.lean}, line 224}\par\nopagebreak
\begin{leansrc}
theorem semanticallyValid_iff_face (H : Hypernet) (seq : Sequent) :
    seq.semanticallyValid H ↔
      upperProb (admCyl H seq) (postEvt H seq)ᶜ = 0
\end{leansrc}
\par\medskip\noindent{\small \texttt{SuFunctor.Plain.}\texttt{\bfseries suppIn\_empty}}\hfill{\footnotesize \path{lean/su-functor/SuFunctor/Plain/CheckerFace.lean}, line 78}\par\nopagebreak
\begin{leansrc}
theorem suppIn_empty : suppIn (∅ : Set X) = {Joint.zero X}
\end{leansrc}
\par\medskip\noindent{\small \texttt{SuFunctor.Plain.}\texttt{\bfseries face\_ex\_falso}}\hfill{\footnotesize \path{lean/su-functor/SuFunctor/Plain/CheckerFace.lean}, line 243}\par\nopagebreak
\begin{leansrc}
theorem face_ex_falso (H : Hypernet) (seq : Sequent)
    (h : admCyl H seq = ∅) :
    upperProb (admCyl H seq) (postEvt H seq)ᶜ = 0
\end{leansrc}
\par\medskip\noindent{\small \texttt{SuFunctor.}\texttt{\bfseries hnetInclusion}}\hfill{\footnotesize \path{lean/su-functor/SuFunctor/HypernetDouble.lean}, line 559}\par\nopagebreak
\begin{leansrc}
def hnetInclusion : LaxDoubleFunctor (HNet ty) (SingleSiteSemanticUniverse ElVal)
\end{leansrc}
\par\medskip\noindent{\small \texttt{SuFunctor.Universe.General.}\texttt{\bfseries sat\_wall\_iff}}\hfill{\footnotesize \path{lean/su-functor/SuFunctor/Universe/General/CheckerAnchor.lean}, line 82}\par\nopagebreak
\begin{leansrc}
theorem sat_wall_iff {σ : StatePenalty X}
    {C : Set (SubProbabilityMeasure X)} :
    Sat σ (wall C) ↔ ∀ ν, σ ν ≠ ⊤ → ν ∈ C
\end{leansrc}
\par\medskip\noindent{\small \texttt{SuFunctor.Universe.General.}\texttt{\bfseries sat\_wall\_wall\_iff}}\hfill{\footnotesize \path{lean/su-functor/SuFunctor/Universe/General/CheckerAnchor.lean}, line 92}\par\nopagebreak
\begin{leansrc}
theorem sat_wall_wall_iff {S C : Set (SubProbabilityMeasure X)} :
    Sat (wall S) (wall C) ↔ S ⊆ C
\end{leansrc}
\par\medskip\noindent{\small \texttt{SuFunctor.Universe.General.}\texttt{\bfseries sat\_iff\_sublevel\_subset}}\hfill{\footnotesize \path{lean/su-functor/SuFunctor/Universe/General/CheckerAnchor.lean}, line 111}\par\nopagebreak
\begin{leansrc}
theorem sat_iff_sublevel_subset {σ φ : StatePenalty X} :
    Sat σ φ ↔ ∀ s : ℝ≥0∞,
      {ν : SubProbabilityMeasure X | σ ν ≤ s}
        ⊆ {ν : SubProbabilityMeasure X | φ ν ≤ s}
\end{leansrc}
\par\medskip\noindent{\small \texttt{PPLC.}\texttt{\bfseries OpenHypernet}}\hfill{\footnotesize \path{lean/qtsl/Qtsl/HypernetCospan.lean}, line 222}\par\nopagebreak
\begin{leansrc}
structure OpenHypernet where …
\end{leansrc}
\par\medskip\noindent{\small \texttt{PPLC.OpenHypernet.}\texttt{\bfseries glue\_admissible\_iff}}\hfill{\footnotesize \path{lean/qtsl/Qtsl/HypernetCospan.lean}, line 444}\par\nopagebreak
\begin{leansrc}
theorem glue_admissible_iff (σ : (glueNet hA).Valuation) :
    σ.Admissible ↔
      (∀ a ∈ H.net.applications, a.SemOk σ) ∧
      (∀ a ∈ G.net.applications, a.SemOk σ)
\end{leansrc}
\par\medskip\noindent{\small \texttt{PPLC.OpenHypernet.}\texttt{\bfseries unitOpen}}\hfill{\footnotesize \path{lean/qtsl/Qtsl/HypernetCospan.lean}, line 465}\par\nopagebreak
\begin{leansrc}
def unitOpen (X : List Wire) (wt₀ : Wire → Option TensorType)
    (hX : X.Nodup) : OpenHypernet where …
\end{leansrc}
\par\medskip\noindent{\small \texttt{PPLC.OpenHypernet.}\texttt{\bfseries glue\_assoc}}\hfill{\footnotesize \path{lean/qtsl/Qtsl/HypernetCospan.lean}, line 558}\par\nopagebreak
\begin{leansrc}
theorem glue_assoc {H G K : OpenHypernet}
    (hHG : Aligned H G) (hGK : Aligned G K)
    (hL : Aligned (glue hHG) K) (hR : Aligned H (glue hGK)) :
    glue hL = glue hR
\end{leansrc}
\par\medskip\noindent{\small \texttt{SuFunctor.}\texttt{\bfseries theta\_rename}}\hfill{\footnotesize \path{lean/su-functor/SuFunctor/HypernetInclusion.lean}, line 912}\par\nopagebreak
\begin{leansrc}
theorem theta_rename {m : AlignedLoose X Y} {m' : AlignedLoose X' Y'}
    (c : AlignedCell ty m m') (T : List PrimApp)
    (hTC : ∀ a ∈ T, a.TypeConsistent (tableOn ty m.apex))
    (hw : ∀ a ∈ T, ∀ w ∈ a.inputWires ++ a.outputWires, w ∈ m.apex) :
    dnfPull (renameTot ty c) (austereDnf ty m T)
      = austereDnf ty m' (T.map (PrimApp.rename c.map))
\end{leansrc}
\par\medskip\noindent{\small \texttt{PPLC.ProofObject.}\texttt{\bfseries sound}}\hfill{\footnotesize \path{lean/qtsl/Qtsl/Soundness.lean}, line 8330}\par\nopagebreak
\begin{leansrc}
theorem ProofObject.sound {H : Hypernet} {seq : Sequent}
    (pf : ProofObject H seq) : seq.semanticallyValid H
\end{leansrc}
\par\medskip\noindent{\small \texttt{PPLC.}\texttt{\bfseries check\_sound}}\hfill{\footnotesize \path{lean/qtsl/Qtsl/Soundness.lean}, line 8439}\par\nopagebreak
\begin{leansrc}
theorem check_sound {H : Hypernet} {seq : Sequent} (pf : ProofObject H seq)
    (_hcheck : check pf = true) : seq.semanticallyValid H
\end{leansrc}
\par\medskip\noindent Modules cited: \texttt{QtslAux.DiamondSite}, \texttt{QtslAux.DiamondTopology}, \texttt{SuFunctor.Site.TemporalTypeTheory}, \texttt{SuFunctor.Site.BehaviorSpec}, \texttt{SuFunctor.Site.CausalConcordance}, \texttt{SuFunctor.StateSpace}, \texttt{SuFunctor.StateSpace.FamCat}, \texttt{QtslAux.QuantumCovariance}, \texttt{QtslAux.OpenQuantum}, \texttt{QtslAux.QuantumRecords}, \texttt{QtslAux.QuantumDescent}, \texttt{QtslAux.QuantumStates}, \texttt{QtslAux.QuantumAssembly}, \texttt{QtslAux.QuantumDeterminacy}, \texttt{QtslAux.QuantumOuterStorey}, \texttt{SuFunctor.Universe.FellTopology}, \texttt{SuFunctor.Universe.FellRestrict}, \texttt{SuFunctor.Universe.FellSeamWitness}, \texttt{SuFunctor.Universe.FellDescent}, \texttt{SuFunctor.Universe.FellRecordLocus}, \texttt{SuFunctor.Universe.FellRecordWitness}, \texttt{SuFunctor.Universe.FellOnePoint}, \texttt{SuFunctor.Universe.FellCapacity}, \texttt{Pointfree.Frame.Presentation}, \texttt{Pointfree.Frame.HyperspacePresentation}, \texttt{Pointfree.Valuation.Basic}, \texttt{Pointfree.Valuation.Topology}, \texttt{Pointfree.Valuation.EscapeWitness}, \texttt{Pointfree.Valuation.Bridge}, \texttt{Pointfree.Measure.DaggerChart}, \texttt{SuFunctor.Universe.General}, \texttt{SuFunctor.Universe.Judgment}, \texttt{SuFunctor.Universe.Closure}, \texttt{SuFunctor.Universe.Generation}, \texttt{SuFunctor.Universe.Tents}, \texttt{QtslAux.Imprecise.MarginalCompat}, \texttt{QtslAux.Imprecise.MarginalCompatGeneral}, \texttt{SuFunctor.Universe.RadialUniverse}, \texttt{SuFunctor.StateSpace.FamCompletion}, \texttt{SuFunctor.Universe.VietorisStratum}, \texttt{SuFunctor.Universe.ImageChar}, \texttt{QtslAux.Imprecise.PenaltyCredalIso}, \texttt{SuFunctor.Universe.EpistemicStates}, \texttt{SuFunctor.Universe.General.CredalPoset}, \texttt{SuFunctor.Universe.General.Richardson}, \texttt{SuFunctor.Universe.CredalPoset}, \texttt{SuFunctor.Universe.HausdorffKL}, \texttt{SuFunctor.Adequacy.Envelope}, \texttt{SuFunctor.Adequacy.EnvelopeUnion}, \texttt{QtslAux.Imprecise.Credal}, \texttt{QtslAux.Imprecise.NaturalExtensionGeneral}, \texttt{QtslAux.Imprecise.NaturalExtension}, \texttt{QtslAux.LeakAware}, \texttt{SuFunctor.Universe.General.CrescentStratum}, \texttt{SuFunctor.Universe.HullClosure}, \texttt{SuFunctor.Site.GeneralWindow}, \texttt{SuFunctor.Plain.SeamDensity}, \texttt{Pointfree.MeasureAlgebra.Basic}, \texttt{Pointfree.Prevision.Basic}, \texttt{SuFunctor.QtslSem.SpecJudg}, \texttt{DoubleCat.Basic}, \texttt{DoubleCat.Functor.Lax}, \texttt{DoubleCat.DepSpan}, \texttt{SuFunctor.Universe.Marginalization}, \texttt{DoubleCat.WireSpan}, \texttt{SuFunctor.Plain.Penalty}, \texttt{SuFunctor.Plain.SpecFunctor}, \texttt{SuFunctor.Plain.Battery}, \texttt{SuFunctor.Site.WindowPenalty}, \texttt{SuFunctor.Site.SiteFunctor}, \texttt{SuFunctor.Site.Battery}, \texttt{SuFunctor.Site.MatchingFork}, \texttt{SuFunctor.Plain.GeneralFunctor}, \texttt{DoubleCat.Grothendieck}, \texttt{DoubleCat.Opfibration}, \texttt{DoubleCat.JudgLayer}, \texttt{QtslAux.Imprecise.Containment}, \texttt{QtslAux.Imprecise.ContainmentGeneral}, \texttt{QtslAux.Imprecise.CouplingGlue}, \texttt{QtslAux.Imprecise.FeasibilityGlue}, \texttt{SuFunctor.Plain.Seam}, \texttt{SuFunctor.Plain.GeneralSeam}, \texttt{SuFunctor.Universe.TransportCitizenship}, \texttt{SuFunctor.Universe.General.TransportGeneral}, \texttt{SuFunctor.Universe.DemonicClock}, \texttt{SuFunctor.Universe.General.DemonicGeneral}, \texttt{DoubleCat.CrossSite}, \texttt{SuFunctor.Spacetime.Base}, \texttt{SuFunctor.Universe.CertificateLayer}, \texttt{SuFunctor.Spacetime.Universe}, \texttt{SuFunctor.Target}, \texttt{SuFunctor.Adequacy.Inclusions}, \texttt{SuFunctor.Adequacy.Dtmc}, \texttt{SuFunctor.Adequacy.Pgm}, \texttt{SuFunctor.Adequacy.Mdp}, \texttt{SuFunctor.Adequacy.MdpDemonic}, \texttt{SuFunctor.Adequacy.DelayLine}, \texttt{SuFunctor.Adequacy.HeatTorus}, \texttt{SuFunctor.Adequacy.MarginalGlue}, \texttt{SuFunctor.Adequacy.HeatTorusInit}, \texttt{SuFunctor.Adequacy.SolutionSetMesh}, \texttt{SuFunctor.Plain.CheckerFace}, \texttt{SuFunctor.HypernetDouble}, \texttt{SuFunctor.Universe.General.CheckerAnchor}, \texttt{Qtsl.HypernetCospan}, \texttt{SuFunctor.HypernetInclusion}, \texttt{Qtsl.Soundness}.

\section{Hypernets: Primitive Functions}
Citations from Chapter~\ref{ch:qtsl-hypernets}.

\par\medskip\noindent{\small \texttt{PPLC.OpenHypernet.}\texttt{\bfseries glue\_unit\_left}}\hfill{\footnotesize \path{lean/qtsl/Qtsl/HypernetCospan.lean}, line 517}\par\nopagebreak
\begin{leansrc}
theorem glue_unit_left (H : OpenHypernet)
    (wt₀ : Wire → Option TensorType)
    (hsub : ∀ w τ, wt₀ w = some τ → H.net.wireType w = some τ)
    (hA : Aligned (unitOpen H.inputs wt₀ H.inputs_nodup) H) :
    glue hA = H
\end{leansrc}
\par\medskip\noindent{\small \texttt{PPLC.OpenHypernet.}\texttt{\bfseries glue\_unit\_right}}\hfill{\footnotesize \path{lean/qtsl/Qtsl/HypernetCospan.lean}, line 532}\par\nopagebreak
\begin{leansrc}
theorem glue_unit_right (H : OpenHypernet)
    (wt₀ : Wire → Option TensorType)
    (hsub : ∀ w τ, wt₀ w = some τ → H.net.wireType w = some τ)
    (hA : Aligned H (unitOpen H.outputs wt₀ H.outputs_nodup)) :
    glue hA = H
\end{leansrc}
\par\medskip\noindent{\small \texttt{SuFunctor.}\texttt{\bfseries theta\_unit}}\hfill{\footnotesize \path{lean/su-functor/SuFunctor/HypernetInclusion.lean}, line 842}\par\nopagebreak
\begin{leansrc}
theorem theta_unit (X : Finset Wire) :
    dnfPull (idTot ty X) (dnfUnit (ctxIface X))
      = austereDnf ty (AlignedLoose.unit X) []
\end{leansrc}
\par\medskip\noindent{\small \texttt{SuFunctor.}\texttt{\bfseries theta\_comp}}\hfill{\footnotesize \path{lean/su-functor/SuFunctor/HypernetInclusion.lean}, line 864}\par\nopagebreak
\begin{leansrc}
theorem theta_comp (m : AlignedLoose X Y) (n : AlignedLoose Y Z)
    (T₁ T₂ : List PrimApp)
    (hTC₁ : ∀ a ∈ T₁, a.TypeConsistent (tableOn ty m.apex))
    (hw₁ : ∀ a ∈ T₁, ∀ w ∈ a.inputWires ++ a.outputWires, w ∈ m.apex)
    (hTC₂ : ∀ a ∈ T₂, a.TypeConsistent (tableOn ty n.apex))
    (hw₂ : ∀ a ∈ T₂, ∀ w ∈ a.inputWires ++ a.outputWires, w ∈ n.apex) :
    dnfPull (splitTot ty m n)
        (dnfBoxplus (austereDnf ty m T₁) (austereDnf ty n T₂))
      = austereDnf ty (m.seq n) (T₁ ++ T₂)
\end{leansrc}
\par\medskip\noindent{\small \texttt{SuFunctor.}\texttt{\bfseries vacDnf\_dnfPull}}\hfill{\footnotesize \path{lean/su-functor/SuFunctor/HypernetInclusion.lean}, line 811}\par\nopagebreak
\begin{leansrc}
theorem vacDnf_dnfPull {A B : Type} (h : A → B) :
    dnfPull h (vacDnf (Plain.Joint B)) = vacDnf (Plain.Joint A)
\end{leansrc}
\par\medskip\noindent{\small \texttt{SuFunctor.}\texttt{\bfseries hnetInclusion\_pushGLoose}}\hfill{\footnotesize \path{lean/su-functor/SuFunctor/HypernetDouble.lean}, line 574}\par\nopagebreak
\begin{leansrc}
theorem hnetInclusion_pushGLoose {p q : HNet ty}
    (M : SpanValued.GLoose p q) {c d : AlignedCtx} {n : Loose c d}
    {f' : p.1 ⟶ c} {g' : q.1 ⟶ d} (γ : Cell M.base n f' g') :
    SpanValued.bmapLoose (stateBase ty) (hnetTheta ty)
        (SpanValued.pushGLoose M γ) ≍
      SpanValued.pushGLoose
        (SpanValued.bmapLoose (stateBase ty) (hnetTheta ty) M)
        ((stateBase ty).mapCell γ)
\end{leansrc}
\par\medskip\noindent{\small \texttt{SuFunctor.}\texttt{\bfseries austereDnf\_eq\_iff}}\hfill{\footnotesize \path{lean/su-functor/SuFunctor/HypernetInclusion.lean}, line 934}\par\nopagebreak
\begin{leansrc}
theorem austereDnf_eq_iff (m : AlignedLoose X Y) (T T' : List PrimApp) :
    austereDnf ty m T = austereDnf ty m T'
      ↔ AdmRegion ty m T = AdmRegion ty m T'
\end{leansrc}
\par\medskip\noindent{\small \texttt{PPLC.}\texttt{\bfseries substituteRaw\_eq\_glue}}\hfill{\footnotesize \path{lean/qtsl/Qtsl/SubstitutionBridge.lean}, line 1960}\par\nopagebreak
\begin{leansrc}
theorem substituteRaw_eq_glue :
    substituteRaw outer φ E inner =
      { wireTypes := bridgeTable outer φ E inner
        apps := (substituteOpen outer φ E inner
            hwfO hwfI hpo hti hto).net.applications
        ins := (substituteOpen outer φ E inner
            hwfO hwfI hpo hti hto).inputs
        outs := (substituteOpen outer φ E inner
            hwfO hwfI hpo hti hto).outputs }
\end{leansrc}
\par\medskip\noindent{\small \texttt{PPLC.}\texttt{\bfseries substitute\_eq\_glue\_renumber}}\hfill{\footnotesize \path{lean/qtsl/Qtsl/SubstitutionBridge.lean}, line 1985}\par\nopagebreak
\begin{leansrc}
theorem substitute_eq_glue_renumber :
    substitute outer φ E inner =
      (HypernetP.mk (bridgeTable outer φ E inner)
        ((substituteOpen outer φ E inner
            hwfO hwfI hpo hti hto).net.applications)
        ((substituteOpen outer φ E inner hwfO hwfI hpo hti hto).inputs)
        ((substituteOpen outer φ E inner
            hwfO hwfI hpo hti hto).outputs)).renumber
\end{leansrc}
\par\medskip\noindent{\small \texttt{PPLC.}\texttt{\bfseries substituteOpen\_admissible\_iff}}\hfill{\footnotesize \path{lean/qtsl/Qtsl/SubstitutionBridge.lean}, line 2001}\par\nopagebreak
\begin{leansrc}
theorem substituteOpen_admissible_iff
    (σ : (substituteOpen outer φ E inner
        hwfO hwfI hpo hti hto).net.Valuation) :
    σ.Admissible ↔
      ((∀ a ∈ inner.apps, a.SemOk σ) ∧
        (∀ a ∈ (substSeam outer φ E inner).repairs, a.SemOk σ)) ∧
      (∀ a ∈ outer.apps.map
          (PrimApp.remapWires (substOuterWire outer φ E inner)),
        a.SemOk σ)
\end{leansrc}
\par\medskip\noindent Modules cited: \texttt{Qtsl.SubstitutionBridge}, \texttt{Qtsl.PackingGap}.
\par\medskip\noindent Also cited in this chapter, listed under an earlier one: \texttt{Qtsl.HypernetCospan}, \texttt{SuFunctor.HypernetDouble}, \texttt{SuFunctor.HypernetInclusion}, \texttt{glue\_admissible\_iff}, \texttt{glue\_assoc}, \texttt{hnetInclusion}, \texttt{theta\_rename}.

\section{Hypernets: Markov Kernels}
Citations from Chapter~\ref{ch:qtsl-markov-kernels}.

\par\medskip\noindent{\small \texttt{SuFunctor.Stochastic.}\texttt{\bfseries sum\_noiseDenotFun}}\hfill{\footnotesize \path{lean/su-functor/SuFunctor/Stochastic/NoiseSemantics.lean}, line 117}\par\nopagebreak
\begin{leansrc}
theorem sum_noiseDenotFun (g : (Fin N → ℝ) → V) (hg : Measurable g) :
    ∑ v, noiseDenotFun g v = 1
\end{leansrc}
\par\medskip\noindent{\small \texttt{SuFunctor.Stochastic.}\texttt{\bfseries noiseDenot\_comp}}\hfill{\footnotesize \path{lean/su-functor/SuFunctor/Stochastic/NoiseSemantics.lean}, line 162}\par\nopagebreak
\begin{leansrc}
theorem noiseDenot_comp (g : (Fin N → ℝ) → V) (π : V → W)
    (hg : Measurable g) (hπg : Measurable (π ∘ g)) :
    noiseDenot (π ∘ g) hπg = ProbState.push π (noiseDenot g hg)
\end{leansrc}
\par\medskip\noindent{\small \texttt{SuFunctor.Stochastic.}\texttt{\bfseries noiseDenot\_seqComp}}\hfill{\footnotesize \path{lean/su-functor/SuFunctor/Stochastic/NoiseKleisli.lean}, line 232}\par\nopagebreak
\begin{leansrc}
theorem noiseDenot_seqComp (g : (Fin N₁ → ℝ) → V) (k : V → (Fin N₂ → ℝ) → W)
    (hg : Measurable g) (hk : ∀ v, Measurable (k v)) :
    noiseDenot (seqComp g k) (measurable_seqComp hg hk)
      = (noiseDenot g hg).bind fun v => noiseDenot (k v) (hk v)
\end{leansrc}
\par\medskip\noindent{\small \texttt{SuFunctor.Stochastic.}\texttt{\bfseries copy\_ne\_kleisli}}\hfill{\footnotesize \path{lean/su-functor/SuFunctor/Stochastic/NoiseKleisli.lean}, line 332}\par\nopagebreak
\begin{leansrc}
theorem copy_ne_kleisli :
    noiseDenot copyPair measurable_copyPair
      ≠ (noiseDenot bitStage measurable_bitStage).bind
          (fun v => noiseDenot (pairStage v) (measurable_pairStage v))
\end{leansrc}
\par\medskip\noindent{\small \texttt{SuFunctor.Stochastic.}\texttt{\bfseries ThinningForm}}\hfill{\footnotesize \path{lean/su-functor/SuFunctor/Stochastic/JumpMeasure.lean}, line 297}\par\nopagebreak
\begin{leansrc}
structure ThinningForm (V : Type u) [Fintype V] : Type u where …
\end{leansrc}
\par\medskip\noindent{\small \texttt{SuFunctor.Stochastic.JumpTerm.}\texttt{\bfseries elaborate\_mass\_partition}}\hfill{\footnotesize \path{lean/su-functor/SuFunctor/Stochastic/JumpMeasure.lean}, line 541}\par\nopagebreak
\begin{leansrc}
theorem elaborate_mass_partition (J : JumpTerm V) (x : V) :
    ((elaborate J).rate : ℝ) *
        ((∑ y, ((elaborate J).kernel x : Step V → ℝ) (Step.next y))
          + ((elaborate J).kernel x : Step V → ℝ) Step.stay
          + ((elaborate J).kernel x : Step V → ℝ) Step.cemetery)
      = (J.intensity : ℝ)
\end{leansrc}
\par\medskip\noindent{\small \texttt{SuFunctor.Stochastic.JumpTerm.}\texttt{\bfseries outflow\_le\_intensity}}\hfill{\footnotesize \path{lean/su-functor/SuFunctor/Stochastic/JumpMeasure.lean}, line 557}\par\nopagebreak
\begin{leansrc}
theorem outflow_le_intensity (J : JumpTerm V) (x : V) :
    ((elaborate J).rate : ℝ) *
        ((∑ y, ((elaborate J).kernel x : Step V → ℝ) (Step.next y))
          + ((elaborate J).kernel x : Step V → ℝ) Step.cemetery)
      ≤ (J.intensity : ℝ)
\end{leansrc}
\par\medskip\noindent{\small \texttt{SuFunctor.Stochastic.JumpTerm.}\texttt{\bfseries generator\_elaborate}}\hfill{\footnotesize \path{lean/su-functor/SuFunctor/Stochastic/JumpMeasure.lean}, line 381}\par\nopagebreak
\begin{leansrc}
theorem generator_elaborate (J : JumpTerm V) :
    (elaborate J).generator = J.generator
\end{leansrc}
\par\medskip\noindent{\small \texttt{SuFunctor.Stochastic.ThinningForm.}\texttt{\bfseries generator\_dilate}}\hfill{\footnotesize \path{lean/su-functor/SuFunctor/Stochastic/JumpMeasure.lean}, line 592}\par\nopagebreak
\begin{leansrc}
theorem generator_dilate (c : ℚ≥0) (hc : 1 ≤ c) (T : ThinningForm V) :
    (T.dilate c hc).generator = T.generator
\end{leansrc}
\par\medskip\noindent{\small \texttt{SuFunctor.Stochastic.ThinningForm.}\texttt{\bfseries generator\_absorbDiagonal}}\hfill{\footnotesize \path{lean/su-functor/SuFunctor/Stochastic/JumpMeasure.lean}, line 681}\par\nopagebreak
\begin{leansrc}
theorem generator_absorbDiagonal (T : ThinningForm V) :
    T.absorbDiagonal.generator = T.generator
\end{leansrc}
\par\medskip\noindent{\small \texttt{SuFunctor.Stochastic.ThinningForm.JumpTerm.}\texttt{\bfseries elaborate\_superpose}}\hfill{\footnotesize \path{lean/su-functor/SuFunctor/Stochastic/JumpAlgebra.lean}, line 534}\par\nopagebreak
\begin{leansrc}
theorem elaborate_superpose (J₁ J₂ : JumpTerm V) :
    (superpose J₁ J₂).elaborate
      = (J₁.elaborate).superpose (J₂.elaborate)
\end{leansrc}
\par\medskip\noindent{\small \texttt{SuFunctor.Stochastic.ThinningForm.JumpTerm.}\texttt{\bfseries generator\_elaborate\_interleave}}\hfill{\footnotesize \path{lean/su-functor/SuFunctor/Stochastic/JumpAlgebra.lean}, line 561}\par\nopagebreak
\begin{leansrc}
theorem generator_elaborate_interleave (J₁ : JumpTerm V₁)
    (J₂ : JumpTerm V₂) :
    (interleave J₁ J₂).elaborate.generator
      = ((J₁.elaborate).interleave (J₂.elaborate)).generator
\end{leansrc}
\par\medskip\noindent{\small \texttt{SuFunctor.Stochastic.}\texttt{\bfseries generator\_eq\_stepMix}}\hfill{\footnotesize \path{lean/su-functor/SuFunctor/Stochastic/JumpDiffusionThinning.lean}, line 187}\par\nopagebreak
\begin{leansrc}
theorem generator_eq_stepMix (P : JdProfile n m) (V : State n → ℝ)
    (gradV : State n → State n) (hessV : State n → Fin n → Fin n → ℝ)
    (z : State n) :
    Qtsl.JumpDiffusion.generator P V gradV hessV z
      = (∑ i, gradV z i * P.f z i)
        + (1 / 2) * (∑ i, ∑ j, P.gram z i j * hessV z i j)
        + P.kill z * stepMixGenerator (killMixture n) V z
        + (P.legs.map fun lm =>
            lm.1 z * stepMixGenerator (deficitToCemetery lm.2) V z).sum
\end{leansrc}
\par\medskip\noindent Modules cited: \texttt{SuFunctor.Stochastic}, \texttt{SuFunctor.Stochastic.JumpAlgebra}.

\section{Systems: Mathematical Models}
Citations from Chapter~\ref{ch:qtsl-mathematical-models}.

\par\medskip\noindent{\small \texttt{SuFunctor.Adequacy.}\texttt{\bfseries isOdeSol\_iff\_integralCurve}}\hfill{\footnotesize \path{lean/su-functor/SuFunctor/Adequacy/Ode.lean}, line 209}\par\nopagebreak
\begin{leansrc}
theorem isOdeSol_iff_integralCurve {f : E → E} {x₀ : E} {y : ℝ → E}
    (hf : Continuous f) :
    IsOdeSol f x₀ y ↔
      y 0 = x₀ ∧ IsIntegralCurveOn y (fun _ => f) (Ici 0)
\end{leansrc}
\par\medskip\noindent{\small \texttt{SuFunctor.Adequacy.}\texttt{\bfseries isOdeSol\_iff\_integral}}\hfill{\footnotesize \path{lean/su-functor/SuFunctor/Adequacy/Ode.lean}, line 180}\par\nopagebreak
\begin{leansrc}
theorem isOdeSol_iff_integral {f : E → E} {x₀ : E} {y : ℝ → E}
    (hf : Continuous f) :
    IsOdeSol f x₀ y ↔ ContinuousOn y (Ici 0) ∧
      ∀ t ∈ Ici (0 : ℝ), y t = x₀ + ∫ s in (0 : ℝ)..t, f (y s)
\end{leansrc}
\par\medskip\noindent{\small \texttt{SuFunctor.Adequacy.}\texttt{\bfseries ode\_adequacy}}\hfill{\footnotesize \path{lean/su-functor/SuFunctor/Adequacy/Ode.lean}, line 319}\par\nopagebreak
\begin{leansrc}
theorem ode_adequacy [CompleteSpace E] {f : E → E} (hf : Continuous f)
    (x₀ : E) {a : Crd (ℝ → E)} :
    ∃ α : Cell (looseConj (odeAtoms f x₀) : Loose a a)
        (suppC (integralCurves f x₀)) (𝟙 a) (𝟙 a),
      Cell.IsInvertible α
\end{leansrc}
\par\medskip\noindent{\small \texttt{SuFunctor.Adequacy.}\texttt{\bfseries ode\_box\_adequacy}}\hfill{\footnotesize \path{lean/su-functor/SuFunctor/Adequacy/Ode.lean}, line 414}\par\nopagebreak
\begin{leansrc}
theorem ode_box_adequacy [CompleteSpace E] {f : E → E} (hf : Continuous f)
    (B : Set E) {a : Crd (ℝ → E)} :
    ∃ α : Cell (looseConj (odeBoxAtoms f B) : Loose a a)
        (suppC (⋃ x₀ ∈ B, integralCurves f x₀)) (𝟙 a) (𝟙 a),
      Cell.IsInvertible α
\end{leansrc}
\par\medskip\noindent{\small \texttt{SuFunctor.Adequacy.}\texttt{\bfseries odeSols\_eqOn\_of\_lipschitzOnWith}}\hfill{\footnotesize \path{lean/su-functor/SuFunctor/Adequacy/Ode.lean}, line 234}\par\nopagebreak
\begin{leansrc}
theorem odeSols_eqOn_of_lipschitzOnWith {f : E → E} {K : ℝ≥0} {s : Set E}
    (hL : LipschitzOnWith K f s) {x₀ : E} {y₁ y₂ : ℝ → E}
    (h₁ : IsOdeSol f x₀ y₁) (h₂ : IsOdeSol f x₀ y₂)
    (hs₁ : ∀ t ∈ Ici (0 : ℝ), y₁ t ∈ s)
    (hs₂ : ∀ t ∈ Ici (0 : ℝ), y₂ t ∈ s) :
    EqOn y₁ y₂ (Ici 0)
\end{leansrc}
\par\medskip\noindent{\small \texttt{SuFunctor.Adequacy.}\texttt{\bfseries su\_ode\_adequacy}}\hfill{\footnotesize \path{lean/su-functor/SuFunctor/Adequacy/Ode.lean}, line 453}\par\nopagebreak
\begin{leansrc}
theorem su_ode_adequacy [CompleteSpace E] {f : E → E} (hf : Continuous f)
    (x₀ : E) (μ : SubPMF (ℝ → E)) :
    idealPoint μ ∈ (odeDnf f x₀).modelClass ↔
      suppC (integralCurves f x₀) μ
\end{leansrc}
\par\medskip\noindent{\small \texttt{SuFunctor.Adequacy.}\texttt{\bfseries heat\_mesh\_glues}}\hfill{\footnotesize \path{lean/su-functor/SuFunctor/Adequacy/HeatTorus.lean}, line 695}\par\nopagebreak
\begin{leansrc}
theorem heat_mesh_glues (hT : IsOpen T)
    (h1 : IsHeatSol D M.I1 T c1) (h2 : IsHeatSol D M.I2 T c2)
    (h3 : IsHeatSol D M.I3 T c3)
    (hs12 : ∀ x ∈ M.I1 ∩ M.I2, ∀ t ∈ T, c1 x t = c2 x t)
    (hs23 : ∀ x ∈ M.I2 ∩ M.I3, ∀ t ∈ T, c2 x t = c3 x t)
    (hw : WrapWall M T c1 c3) :
    ∃ u : ℝ → ℝ → E, IsHeatTorusSol D M.L T u ∧
      (∀ x ∈ M.I1, ∀ t ∈ T, u x t = c1 x t) ∧
      (∀ x ∈ M.I2, ∀ t ∈ T, u x t = c2 x t) ∧
      (∀ x ∈ M.I3, ∀ t ∈ T, u x t = c3 x t)
\end{leansrc}
\par\medskip\noindent{\small \texttt{SuFunctor.Adequacy.}\texttt{\bfseries heat\_adequacy}}\hfill{\footnotesize \path{lean/su-functor/SuFunctor/Adequacy/HeatTorus.lean}, line 1122}\par\nopagebreak
\begin{leansrc}
theorem heat_adequacy (M : RingMesh) (D : ℝ) {T : Set ℝ}
    (hT : IsOpen T) {a : Crd (HeatRec E)} :
    ∃ α : Cell (looseConj (heatAtoms M D T) : Loose a a)
        (suppC (heatClassicalRecs M D T)) (𝟙 a) (𝟙 a),
      Cell.IsInvertible α
\end{leansrc}
\par\medskip\noindent{\small \texttt{SuFunctor.Adequacy.}\texttt{\bfseries chainAtoms\_dirac\_lin}}\hfill{\footnotesize \path{lean/su-functor/SuFunctor/Adequacy/HeatTorus.lean}, line 919}\par\nopagebreak
\begin{leansrc}
theorem chainAtoms_dirac_lin (e : E) :
    looseConj (heatChainAtoms M D T) (SubPMF.dirac (linRec e))
\end{leansrc}
\par\medskip\noindent{\small \texttt{SuFunctor.Adequacy.}\texttt{\bfseries not\_wrap\_lin}}\hfill{\footnotesize \path{lean/su-functor/SuFunctor/Adequacy/HeatTorus.lean}, line 932}\par\nopagebreak
\begin{leansrc}
theorem not_wrap_lin {e : E} (he : e ≠ 0) {t₀ : ℝ} (ht₀ : t₀ ∈ T) :
    ¬ suppC (wrapC M T) (SubPMF.dirac (linRec e))
\end{leansrc}
\par\medskip\noindent{\small \texttt{SuFunctor.Adequacy.}\texttt{\bfseries heat\_chain\_glues}}\hfill{\footnotesize \path{lean/su-functor/SuFunctor/Adequacy/HeatTorus.lean}, line 484}\par\nopagebreak
\begin{leansrc}
theorem heat_chain_glues {M : RingMesh} {D : ℝ} {T : Set ℝ}
    (hT : IsOpen T) {c1 c2 c3 : ℝ → ℝ → E}
    (h1 : IsHeatSol D M.I1 T c1) (h2 : IsHeatSol D M.I2 T c2)
    (h3 : IsHeatSol D M.I3 T c3)
    (hs12 : ∀ x ∈ M.I1 ∩ M.I2, ∀ t ∈ T, c1 x t = c2 x t)
    (hs23 : ∀ x ∈ M.I2 ∩ M.I3, ∀ t ∈ T, c2 x t = c3 x t) :
    ∃ w : ℝ → ℝ → E, IsHeatSol D M.chart T w ∧
      (∀ x ∈ M.I1, ∀ t ∈ T, w x t = c1 x t) ∧
      (∀ x ∈ M.I2, ∀ t ∈ T, w x t = c2 x t) ∧
      (∀ x ∈ M.I3, ∀ t ∈ T, w x t = c3 x t)
\end{leansrc}
\par\medskip\noindent{\small \texttt{SuFunctor.}\texttt{\bfseries vorobev\_no\_joint}}\hfill{\footnotesize \path{lean/su-functor/SuFunctor/Adequacy/MarginalGlue.lean}, line 246}\par\nopagebreak
\begin{leansrc}
theorem vorobev_no_joint (μ : SubPMF (Bool × Bool × Bool)) :
    ¬ (μ.push proj12 = pairCorr ∧ μ.push proj23 = pairCorr ∧
        μ.push proj13 = pairIndep)
\end{leansrc}
\par\medskip\noindent{\small \texttt{SuFunctor.}\texttt{\bfseries credalCompile\_vorobevGlue}}\hfill{\footnotesize \path{lean/su-functor/SuFunctor/Adequacy/MarginalGlue.lean}, line 415}\par\nopagebreak
\begin{leansrc}
theorem credalCompile_vorobevGlue :
    credalCompile vorobevGlue = rayC (SubPMF.zero (Bool × Bool × Bool))
\end{leansrc}
\par\medskip\noindent{\small \texttt{SuFunctor.Adequacy.}\texttt{\bfseries heat\_ring\_glues\_vorobev\_ring\_not}}\hfill{\footnotesize \path{lean/su-functor/SuFunctor/Adequacy/HeatTorus.lean}, line 1038}\par\nopagebreak
\begin{leansrc}
theorem heat_ring_glues_vorobev_ring_not :
    (∃ μ : SubPMF (HeatRec E), looseConj (heatAtoms M D T) μ ∧
        μ ≠ SubPMF.zero (HeatRec E)) ∧
      ¬ ∃ μ : SubPMF (Bool × Bool × Bool), vorobevGlue μ
\end{leansrc}
\par\medskip\noindent{\small \texttt{SuFunctor.Adequacy.}\texttt{\bfseries heat\_init\_adequacy}}\hfill{\footnotesize \path{lean/su-functor/SuFunctor/Adequacy/HeatTorusInit.lean}, line 291}\par\nopagebreak
\begin{leansrc}
theorem heat_init_adequacy (M : RingMesh) (D : ℝ) {T : Set ℝ}
    (hT : IsOpen T) (g : ℝ → E) {a : Crd (HeatRec E)} :
    ∃ α : Cell (looseConj (heatInitAtoms M D T g) : Loose a a)
        (suppC (heatInitClassicalRecs M D T g)) (𝟙 a) (𝟙 a),
      Cell.IsInvertible α
\end{leansrc}
\par\medskip\noindent{\small \texttt{SuFunctor.Adequacy.}\texttt{\bfseries heatInitAtoms\_consistent\_of\_feed}}\hfill{\footnotesize \path{lean/su-functor/SuFunctor/Adequacy/HeatTorusInit.lean}, line 323}\par\nopagebreak
\begin{leansrc}
theorem heatInitAtoms_consistent_of_feed {M : RingMesh} {D : ℝ}
    {T : Set ℝ} (hT : IsOpen T) {g : ℝ → E}
    (hfeed : PortFeed M D T g) :
    ∃ μ : SubPMF (HeatRec E), looseConj (heatInitAtoms M D T g) μ ∧
      μ ≠ SubPMF.zero (HeatRec E)
\end{leansrc}
\par\medskip\noindent{\small \texttt{SuFunctor.Adequacy.}\texttt{\bfseries initFace\_wrap\_compat}}\hfill{\footnotesize \path{lean/su-functor/SuFunctor/Adequacy/HeatTorusInit.lean}, line 178}\par\nopagebreak
\begin{leansrc}
theorem initFace_wrap_compat {M : RingMesh} {D : ℝ} {T : Set ℝ}
    {g : ℝ → E} {p : HeatRec E} (h : IsInitMeshSol M D T g p)
    (hne : (𝓝[T] (0 : ℝ)).NeBot) {x : ℝ} (hx3 : x ∈ M.I3)
    (hx1 : x - M.L ∈ M.I1) : g x = g (x - M.L)
\end{leansrc}
\par\medskip\noindent{\small \texttt{SuFunctor.Adequacy.IsHeatTorusSolInit.}\texttt{\bfseries trace\_wrap\_compat}}\hfill{\footnotesize \path{lean/su-functor/SuFunctor/Adequacy/HeatTorusInit.lean}, line 208}\par\nopagebreak
\begin{leansrc}
theorem IsHeatTorusSolInit.trace_wrap_compat {M : RingMesh} {D : ℝ}
    {T : Set ℝ} {u : ℝ → ℝ → E} {g : ℝ → E}
    (h : IsHeatTorusSolInit M D T u g) (hne : (𝓝[T] (0 : ℝ)).NeBot)
    {x : ℝ} (hx : x ∈ M.chart) (hxL : x + M.L ∈ M.chart) :
    g (x + M.L) = g x
\end{leansrc}
\par\medskip\noindent{\small \texttt{SuFunctor.Adequacy.}\texttt{\bfseries portFeed\_const}}\hfill{\footnotesize \path{lean/su-functor/SuFunctor/Adequacy/HeatTorusInit.lean}, line 335}\par\nopagebreak
\begin{leansrc}
theorem portFeed_const (M : RingMesh) (D : ℝ) (T : Set ℝ) (k : E) :
    PortFeed M D T (fun _ => k)
\end{leansrc}
\par\medskip\noindent{\small \texttt{SuFunctor.Adequacy.}\texttt{\bfseries portFeed\_mode}}\hfill{\footnotesize \path{lean/su-functor/SuFunctor/Adequacy/HeatTorusInit.lean}, line 436}\par\nopagebreak
\begin{leansrc}
theorem portFeed_mode (M : RingMesh) (D : ℝ) (T : Set ℝ) (e : E)
    {ω : ℝ} (k : ℤ) (hω : ω * M.L = k * (2 * Real.pi)) :
    PortFeed M D T (modeDatum ω e)
\end{leansrc}
\par\medskip\noindent{\small \texttt{SuFunctor.Adequacy.}\texttt{\bfseries modeDatum\_ne\_const}}\hfill{\footnotesize \path{lean/su-functor/SuFunctor/Adequacy/HeatTorusInit.lean}, line 447}\par\nopagebreak
\begin{leansrc}
theorem modeDatum_ne_const {e : E} (he : e ≠ 0) {ω : ℝ} (hω : ω ≠ 0) :
    modeDatum ω e (ω⁻¹ * (Real.pi / 2)) ≠ modeDatum ω e 0
\end{leansrc}
\par\medskip\noindent{\small \texttt{SuFunctor.Adequacy.}\texttt{\bfseries initFace\_neBot}}\hfill{\footnotesize \path{lean/su-functor/SuFunctor/Adequacy/HeatTorusInit.lean}, line 99}\par\nopagebreak
\begin{leansrc}
theorem initFace_neBot {τ : ℝ} (hτ : 0 < τ) :
    (𝓝[Ioo (0 : ℝ) τ] (0 : ℝ)).NeBot
\end{leansrc}
\par\medskip\noindent{\small \texttt{SuFunctor.Universe.FellCloseds.}\texttt{\bfseries equalizerHomeomorph}}\hfill{\footnotesize \path{lean/su-functor/SuFunctor/Universe/FellDescentEqualizer.lean}, line 287}\par\nopagebreak
\begin{leansrc}
def equalizerHomeomorph [R1Space α] (U : ι → Set α) (hUo : ∀ i, IsOpen (U i))
    (hcov : ⋃ i, U i = univ) :
    FellCloseds α ≃ₜ Equalizer U where …
\end{leansrc}
\par\medskip\noindent{\small \texttt{SuFunctor.Adequacy.}\texttt{\bfseries meshDescent}}\hfill{\footnotesize \path{lean/su-functor/SuFunctor/Adequacy/SolutionSetMesh.lean}, line 235}\par\nopagebreak
\begin{leansrc}
noncomputable def meshDescent (M : RingMesh) (T : Set ℝ) :
    FellCloseds (MeshPt M T E) ≃ₜ
      FellCloseds.Equalizer (arcRegion M T E)
\end{leansrc}
\par\medskip\noindent{\small \texttt{SuFunctor.Adequacy.}\texttt{\bfseries setStratum\_holonomy}}\hfill{\footnotesize \path{lean/su-functor/SuFunctor/Adequacy/SolutionSetMesh.lean}, line 688}\par\nopagebreak
\begin{leansrc}
theorem setStratum_holonomy (D : ℝ) {t₀ : ℝ} (ht₀ : t₀ ∈ T) :
    ∃ p : FellCloseds.Equalizer (arcRegion M T ℝ),
      (∀ i, IsCellTube M D T i (p.1 i : Set ↥(arcRegion M T ℝ i))) ∧
      TubeWrapWall M T (p.1 0 : Set ↥(arcRegion M T ℝ 0))
        (p.1 2 : Set ↥(arcRegion M T ℝ 2)) ∧
      IsChartTube M D T (meshGlue M T p : Set (MeshPt M T ℝ)) ∧
      ¬ IsTorusTube M D T (meshGlue M T p : Set (MeshPt M T ℝ))
\end{leansrc}
\par\medskip\noindent{\small \texttt{SuFunctor.Adequacy.}\texttt{\bfseries orbitTube\_not\_torusTube}}\hfill{\footnotesize \path{lean/su-functor/SuFunctor/Adequacy/SolutionSetMesh.lean}, line 640}\par\nopagebreak
\begin{leansrc}
theorem orbitTube_not_torusTube (D : ℝ) {t₀ : ℝ} (ht₀ : t₀ ∈ T) :
    ¬ IsTorusTube M D T (orbitTube M T)
\end{leansrc}
\par\medskip\noindent{\small \texttt{SuFunctor.Adequacy.}\texttt{\bfseries not\_isLocalLocusEq\_torusTube}}\hfill{\footnotesize \path{lean/su-functor/SuFunctor/Adequacy/SolutionSetMesh.lean}, line 723}\par\nopagebreak
\begin{leansrc}
theorem not_isLocalLocusEq_torusTube (D : ℝ) {t₀ : ℝ} (ht₀ : t₀ ∈ T) :
    ¬ FellCloseds.IsLocalLocusEq (arcRegion M T ℝ)
        (torusTubeLocus M D T) (cellTubeLocus M D T)
\end{leansrc}
\par\medskip\noindent{\small \texttt{SuFunctor.Adequacy.}\texttt{\bfseries tubeAtoms\_not\_classical}}\hfill{\footnotesize \path{lean/su-functor/SuFunctor/Adequacy/SolutionSetMesh.lean}, line 863}\par\nopagebreak
\begin{leansrc}
theorem tubeAtoms_not_classical {t₀ : ℝ} (ht₀ : t₀ ∈ T) (D : ℝ) :
    ¬ ∀ μ : SubPMF (FellCloseds.Equalizer (arcRegion M T ℝ)),
        looseConj (tubeAtoms M D T) μ ↔ suppC (torusTubeEqs M D T) μ
\end{leansrc}
\par\medskip\noindent{\small \texttt{SuFunctor.Adequacy.}\texttt{\bfseries solset\_forward\_cell}}\hfill{\footnotesize \path{lean/su-functor/SuFunctor/Adequacy/SolutionSetMesh.lean}, line 818}\par\nopagebreak
\begin{leansrc}
theorem solset_forward_cell
    {a : Crd (FellCloseds.Equalizer (arcRegion M T E))} :
    ∃ _ : Cell (suppC (torusTubeEqs M D T) : Loose a a)
        (looseConj (tubeAtoms M D T)) (𝟙 a) (𝟙 a), True
\end{leansrc}
\par\medskip\noindent{\small \texttt{SuFunctor.Adequacy.}\texttt{\bfseries isTorusTube\_univ}}\hfill{\footnotesize \path{lean/su-functor/SuFunctor/Adequacy/SolutionSetMesh.lean}, line 372}\par\nopagebreak
\begin{leansrc}
theorem isTorusTube_univ :
    IsTorusTube M D T (univ : Set (MeshPt M T E))
\end{leansrc}
\par\medskip\noindent{\small \texttt{SuFunctor.Adequacy.}\texttt{\bfseries labelled\_ensembles\_glue}}\hfill{\footnotesize \path{lean/su-functor/SuFunctor/Adequacy/SolutionSetMesh.lean}, line 385}\par\nopagebreak
\begin{leansrc}
theorem labelled_ensembles_glue {ι' : Type*} (hT : IsOpen T)
    {c1 c2 c3 : ι' → ℝ → ℝ → E}
    (h1 : ∀ j, IsHeatSol D M.I1 T (c1 j))
    (h2 : ∀ j, IsHeatSol D M.I2 T (c2 j))
    (h3 : ∀ j, IsHeatSol D M.I3 T (c3 j))
    (hs12 : ∀ j, ∀ x ∈ M.I1 ∩ M.I2, ∀ t ∈ T, c1 j x t = c2 j x t)
    (hs23 : ∀ j, ∀ x ∈ M.I2 ∩ M.I3, ∀ t ∈ T, c2 j x t = c3 j x t)
    (hw : ∀ j, WrapWall M T (c1 j) (c3 j)) :
    ∃ u : ι' → ℝ → ℝ → E, ∀ j,
      IsHeatTorusSol D M.L T (u j) ∧
      (∀ x ∈ M.I1, ∀ t ∈ T, u j x t = c1 j x t) ∧
      (∀ x ∈ M.I2, ∀ t ∈ T, u j x t = c2 j x t) ∧
      (∀ x ∈ M.I3, ∀ t ∈ T, u j x t = c3 j x t)
\end{leansrc}
\par\medskip\noindent{\small \texttt{SuFunctor.Adequacy.}\texttt{\bfseries orbit\_enclosure\_forces\_zero}}\hfill{\footnotesize \path{lean/su-functor/SuFunctor/Adequacy/SolutionSetMesh.lean}, line 991}\par\nopagebreak
\begin{leansrc}
theorem orbit_enclosure_forces_zero (D : ℝ) {t₀ : ℝ} (ht₀ : t₀ ∈ T) :
    ∀ μ : SubPMF (HeatRec ℝ),
      looseConj (heatAtoms M D T (E := ℝ)
          ++ [suppC (enclose1C M T (orbitSet M)),
              suppC (enclose2C M T (orbitSet M)),
              suppC (enclose3C M T (orbitSet M))]) μ →
      μ = SubPMF.zero (HeatRec ℝ)
\end{leansrc}
\par\medskip\noindent{\small \texttt{SuFunctor.Adequacy.LocalRows.}\texttt{\bfseries mesh\_adequacy}}\hfill{\footnotesize \path{lean/su-functor/SuFunctor/Adequacy/Pde.lean}, line 400}\par\nopagebreak
\begin{leansrc}
theorem mesh_adequacy (R : LocalRows E) (M : RingMesh) {T : Set ℝ}
    (hT : IsOpen T) {a : Crd (HeatRec E)} :
    ∃ α : Cell (looseConj (R.meshAtoms M T) : Loose a a)
        (suppC (R.classicalRecs M T)) (𝟙 a) (𝟙 a),
      Cell.IsInvertible α
\end{leansrc}
\par\medskip\noindent{\small \texttt{SuFunctor.Adequacy.LocalRows.PdeAt.}\texttt{\bfseries congr}}\hfill{\footnotesize \path{lean/su-functor/SuFunctor/Adequacy/Pde.lean}, line 644}\par\nopagebreak
\begin{leansrc}
theorem PdeAt.congr {u v : ℝ → ℝ → ι → ℝ} {x t : ℝ}
    {Sx St : Set ℝ} (hSx : IsOpen Sx) (hSt : IsOpen St)
    (hx : x ∈ Sx) (ht : t ∈ St)
    (hagree : ∀ y ∈ Sx, ∀ s ∈ St, u y s = v y s)
    (h : PdeAt P v x t) : PdeAt P u x t
\end{leansrc}
\par\medskip\noindent{\small \texttt{SuFunctor.Adequacy.LocalRows.PdeAt.}\texttt{\bfseries comp\_sub\_const}}\hfill{\footnotesize \path{lean/su-functor/SuFunctor/Adequacy/Pde.lean}, line 683}\par\nopagebreak
\begin{leansrc}
theorem PdeAt.comp_sub_const {w : ℝ → ℝ → ι → ℝ} {x t c : ℝ}
    (h : PdeAt P w (x - c) t) :
    PdeAt P (fun y s => w (y - c) s) x t
\end{leansrc}
\par\medskip\noindent{\small \texttt{SuFunctor.Adequacy.LocalRows.}\texttt{\bfseries pde\_adequacy}}\hfill{\footnotesize \path{lean/su-functor/SuFunctor/Adequacy/Pde.lean}, line 768}\par\nopagebreak
\begin{leansrc}
theorem pde_adequacy (P : PdeMechanism ι k) (M : RingMesh) {T : Set ℝ}
    (hT : IsOpen T) {a : Crd (HeatRec (ι → ℝ))} :
    ∃ α : Cell (looseConj (P.rows.meshAtoms M T) : Loose a a)
        (suppC (P.rows.classicalRecs M T)) (𝟙 a) (𝟙 a),
      Cell.IsInvertible α
\end{leansrc}
\par\medskip\noindent{\small \texttt{SuFunctor.Adequacy.LocalRows.}\texttt{\bfseries heat\_adequacy\_of\_mesh}}\hfill{\footnotesize \path{lean/su-functor/SuFunctor/Adequacy/Pde.lean}, line 510}\par\nopagebreak
\begin{leansrc}
theorem heat_adequacy_of_mesh (M : RingMesh) (D : ℝ) {T : Set ℝ}
    (hT : IsOpen T) {a : Crd (HeatRec E)} :
    ∃ α : Cell (looseConj (heatAtoms M D T) : Loose a a)
        (suppC (heatClassicalRecs M D T)) (𝟙 a) (𝟙 a),
      Cell.IsInvertible α
\end{leansrc}
\par\medskip\noindent{\small \texttt{SuFunctor.Adequacy.LocalRows.}\texttt{\bfseries heatAt\_iff\_pdeAt}}\hfill{\footnotesize \path{lean/su-functor/SuFunctor/Adequacy/Pde.lean}, line 835}\par\nopagebreak
\begin{leansrc}
theorem heatAt_iff_pdeAt {D : ℝ} {u : ℝ → ℝ → ℝ} {x t : ℝ} :
    HeatAt D u x t ↔ PdeAt (heatMechanism D) (fun y s _ => u y s) x t
\end{leansrc}
\par\medskip\noindent{\small \texttt{SuFunctor.Adequacy.LocalRows.}\texttt{\bfseries burgers\_adequacy}}\hfill{\footnotesize \path{lean/su-functor/SuFunctor/Adequacy/Pde.lean}, line 918}\par\nopagebreak
\begin{leansrc}
theorem burgers_adequacy (ν : ℝ) (M : RingMesh) {T : Set ℝ}
    (hT : IsOpen T) {a : Crd (HeatRec (PUnit → ℝ))} :
    ∃ α : Cell
        (looseConj ((burgersMechanism ν).rows.meshAtoms M T) : Loose a a)
        (suppC ((burgersMechanism ν).rows.classicalRecs M T))
        (𝟙 a) (𝟙 a),
      Cell.IsInvertible α
\end{leansrc}
\par\medskip\noindent{\small \texttt{SuFunctor.Adequacy.LocalRows.}\texttt{\bfseries kpp\_adequacy}}\hfill{\footnotesize \path{lean/su-functor/SuFunctor/Adequacy/Pde.lean}, line 928}\par\nopagebreak
\begin{leansrc}
theorem kpp_adequacy (M : RingMesh) {T : Set ℝ}
    (hT : IsOpen T) {a : Crd (HeatRec (PUnit → ℝ))} :
    ∃ α : Cell
        (looseConj (kppMechanism.rows.meshAtoms M T) : Loose a a)
        (suppC (kppMechanism.rows.classicalRecs M T)) (𝟙 a) (𝟙 a),
      Cell.IsInvertible α
\end{leansrc}
\par\medskip\noindent{\small \texttt{SuFunctor.Adequacy.LocalRows.}\texttt{\bfseries coupled\_adequacy}}\hfill{\footnotesize \path{lean/su-functor/SuFunctor/Adequacy/Pde.lean}, line 937}\par\nopagebreak
\begin{leansrc}
theorem coupled_adequacy (M : RingMesh) {T : Set ℝ}
    (hT : IsOpen T) {a : Crd (HeatRec (Fin 2 → ℝ))} :
    ∃ α : Cell
        (looseConj (coupledMechanism.rows.meshAtoms M T) : Loose a a)
        (suppC (coupledMechanism.rows.classicalRecs M T)) (𝟙 a) (𝟙 a),
      Cell.IsInvertible α
\end{leansrc}
\par\medskip\noindent{\small \texttt{SuFunctor.Adequacy.LocalRows.}\texttt{\bfseries PdeRung}}\hfill{\footnotesize \path{lean/su-functor/SuFunctor/Adequacy/Pde.lean}, line 1010}\par\nopagebreak
\begin{leansrc}
structure PdeRung {k : ℕ} (P : PdeMechanism ι k) where …
\end{leansrc}
\par\medskip\noindent{\small \texttt{SuFunctor.Adequacy.LocalRows.}\texttt{\bfseries rung\_adequacy}}\hfill{\footnotesize \path{lean/su-functor/SuFunctor/Adequacy/Pde.lean}, line 1045}\par\nopagebreak
\begin{leansrc}
theorem rung_adequacy {k : ℕ} {P : PdeMechanism ι k} (B : PdeRung P)
    (M : RingMesh) {T : Set ℝ} (hT : IsOpen T)
    {a : Crd (HeatRec (ι → ℝ))} :
    ∃ α : Cell (looseConj (B.rows.meshAtoms M T) : Loose a a)
        (suppC (B.rows.classicalRecs M T)) (𝟙 a) (𝟙 a),
      Cell.IsInvertible α
\end{leansrc}
\par\medskip\noindent{\small \texttt{SuFunctor.Adequacy.LocalRows.}\texttt{\bfseries pde\_rung\_transfer}}\hfill{\footnotesize \path{lean/su-functor/SuFunctor/Adequacy/Pde.lean}, line 1070}\par\nopagebreak
\begin{leansrc}
theorem pde_rung_transfer {k : ℕ} {P : PdeMechanism ι k} (B : PdeRung P)
    (M : RingMesh) (T : Set ℝ) {a : Crd (HeatRec (ι → ℝ))} :
    ∃ _ : Cell (looseConj (P.rows.meshAtoms M T) : Loose a a)
        (looseConj (B.rows.meshAtoms M T)) (𝟙 a) (𝟙 a), True
\end{leansrc}
\par\medskip\noindent{\small \texttt{SuFunctor.Adequacy.LocalRows.}\texttt{\bfseries pdeAtoms\_consistent\_of\_equilibrium}}\hfill{\footnotesize \path{lean/su-functor/SuFunctor/Adequacy/Pde.lean}, line 795}\par\nopagebreak
\begin{leansrc}
theorem pdeAtoms_consistent_of_equilibrium {M : RingMesh} {T : Set ℝ}
    {e : ι → ℝ} (he : P.IsEquilibriumValue e) :
    ∃ μ : SubPMF (HeatRec (ι → ℝ)),
      looseConj (P.rows.meshAtoms M T) μ ∧
        μ ≠ SubPMF.zero (HeatRec (ι → ℝ))
\end{leansrc}
\par\medskip\noindent{\small \texttt{SuFunctor.Adequacy.LocalRows.}\texttt{\bfseries kpp\_not\_equilibrium\_half}}\hfill{\footnotesize \path{lean/su-functor/SuFunctor/Adequacy/Pde.lean}, line 975}\par\nopagebreak
\begin{leansrc}
theorem kpp_not_equilibrium_half :
    ¬ kppMechanism.IsEquilibriumValue (fun _ => (1 / 2 : ℝ))
\end{leansrc}
\par\medskip\noindent{\small \texttt{SuFunctor.Adequacy.LocalRows.}\texttt{\bfseries su\_pde\_adequacy}}\hfill{\footnotesize \path{lean/su-functor/SuFunctor/Adequacy/Pde.lean}, line 1132}\par\nopagebreak
\begin{leansrc}
theorem su_pde_adequacy {k : ℕ} (P : PdeMechanism ι k) (M : RingMesh)
    {T : Set ℝ} (hT : IsOpen T) (μ : SubPMF (HeatRec (ι → ℝ))) :
    idealPoint μ ∈ (P.rows.meshDnf M T).modelClass ↔
      suppC (P.rows.classicalRecs M T) μ
\end{leansrc}
\par\medskip\noindent{\small \texttt{SuFunctor.Adequacy.PdeRowSpec.}\texttt{\bfseries toMechanism}}\hfill{\footnotesize \path{lean/su-functor/SuFunctor/Adequacy/PdeRows.lean}, line 173}\par\nopagebreak
\begin{leansrc}
noncomputable def toMechanism [Inhabited ι] (S : PdeRowSpec ι k) :
    PdeMechanism ι k where …
\end{leansrc}
\par\medskip\noindent{\small \texttt{SuFunctor.Adequacy.}\texttt{\bfseries toMv\_congr}}\hfill{\footnotesize \path{lean/su-functor/SuFunctor/Adequacy/PdeRows.lean}, line 154}\par\nopagebreak
\begin{leansrc}
theorem toMv_congr {σ : Type} {b₁ b₂ : ℕ → σ} :
    ∀ rows : List RowTerm,
      (∀ t ∈ rows, ∀ j, j < t.exps.length → b₁ j = b₂ j) →
      toMv b₁ rows = toMv b₂ rows
  | [], _ => rfl
  | t :: ts, h => by
    have hterm : termProd b₁ 0 t.exps = termProd b₂ 0 t.exps
\end{leansrc}
\par\medskip\noindent{\small \texttt{SuFunctor.Adequacy.PdeRowSpec.}\texttt{\bfseries toMechanism\_congr\_inhabited}}\hfill{\footnotesize \path{lean/su-functor/SuFunctor/Adequacy/PdeRows.lean}, line 181}\par\nopagebreak
\begin{leansrc}
theorem toMechanism_congr_inhabited (S : PdeRowSpec ι k) (hWF : S.WF)
    (d₁ d₂ : Inhabited ι) :
    @toMechanism ι k d₁ S = @toMechanism ι k d₂ S
\end{leansrc}
\par\medskip\noindent{\small \texttt{SuFunctor.Adequacy.}\texttt{\bfseries heatRowSpec\_toMechanism}}\hfill{\footnotesize \path{lean/su-functor/SuFunctor/Adequacy/PdeRows.lean}, line 220}\par\nopagebreak
\begin{leansrc}
theorem heatRowSpec_toMechanism (c : ℚ) :
    (heatRowSpec c).toMechanism = heatMechanism (c : ℝ)
\end{leansrc}
\par\medskip\noindent{\small \texttt{SuFunctor.Adequacy.}\texttt{\bfseries heat\_rows\_adequacy}}\hfill{\footnotesize \path{lean/su-functor/SuFunctor/Adequacy/PdeRows.lean}, line 239}\par\nopagebreak
\begin{leansrc}
theorem heat_rows_adequacy (c : ℚ) (M : RingMesh) {T : Set ℝ}
    (hT : IsOpen T) {a : Crd (HeatRec (PUnit → ℝ))} :
    ∃ α : Cell
        (looseConj (((heatRowSpec c).toMechanism).rows.meshAtoms M T) : Loose a a)
        (suppC (((heatRowSpec c).toMechanism).rows.classicalRecs M T))
        (𝟙 a) (𝟙 a),
      Cell.IsInvertible α
\end{leansrc}
\par\medskip\noindent{\small \texttt{SuFunctor.Adequacy.}\texttt{\bfseries su\_pde\_rows\_adequacy}}\hfill{\footnotesize \path{lean/su-functor/SuFunctor/Adequacy/PdeRows.lean}, line 310}\par\nopagebreak
\begin{leansrc}
theorem su_pde_rows_adequacy {ι : Type} [Fintype ι] [Inhabited ι] {k : ℕ}
    (S : PdeRowSpec ι k) (M : RingMesh) {T : Set ℝ} (hT : IsOpen T)
    (μ : SubPMF (HeatRec (ι → ℝ))) :
    idealPoint μ ∈ ((S.toMechanism).rows.meshDnf M T).modelClass ↔
      suppC ((S.toMechanism).rows.classicalRecs M T) μ
\end{leansrc}
\par\medskip\noindent{\small \texttt{SuFunctor.Adequacy.}\texttt{\bfseries composite\_iff\_dde}}\hfill{\footnotesize \path{lean/su-functor/SuFunctor/Adequacy/DelayLine.lean}, line 442}\par\nopagebreak
\begin{leansrc}
theorem composite_iff_dde {f : E → E → E} {c τ : ℝ} {φ : ℝ → E}
    (hc : 0 < c) (hτ : 0 < τ) (hφ : ContinuousOn φ (Icc (-τ) 0))
    {y z : ℝ → E} {u : ℝ → ℝ → E} :
    IsCompositeSol f c τ φ y z u ↔
      ∃ Y, IsDdeSol f τ φ Y ∧ RecMatches c τ Y y z u
\end{leansrc}
\par\medskip\noindent{\small \texttt{SuFunctor.Adequacy.}\texttt{\bfseries delayline\_adequacy}}\hfill{\footnotesize \path{lean/su-functor/SuFunctor/Adequacy/DelayLine.lean}, line 592}\par\nopagebreak
\begin{leansrc}
theorem delayline_adequacy (f : E → E → E) {c τ : ℝ} (hc : 0 < c)
    (hτ : 0 < τ) {φ : ℝ → E} (hφ : ContinuousOn φ (Icc (-τ) 0))
    {a : Crd (DelayRec E)} :
    ∃ α : Cell (looseConj (delayAtoms f c τ φ) : Loose a a)
        (suppC (classicalRecs f c τ φ)) (𝟙 a) (𝟙 a),
      Cell.IsInvertible α
\end{leansrc}
\par\medskip\noindent{\small \texttt{SuFunctor.Adequacy.IsCompositeSol.}\texttt{\bfseries outflow\_delay}}\hfill{\footnotesize \path{lean/su-functor/SuFunctor/Adequacy/DelayLine.lean}, line 295}\par\nopagebreak
\begin{leansrc}
theorem outflow_delay (hc : 0 < c) (hτ : 0 < τ)
    (h : IsCompositeSol f c τ φ y z u) :
    ∀ t ∈ Ici (0 : ℝ), z t = concat φ y (t - τ)
\end{leansrc}
\par\medskip\noindent{\small \texttt{SuFunctor.Adequacy.IsCompositeSol.}\texttt{\bfseries field\_slice\_is\_history}}\hfill{\footnotesize \path{lean/su-functor/SuFunctor/Adequacy/DelayLine.lean}, line 310}\par\nopagebreak
\begin{leansrc}
theorem field_slice_is_history (hc : 0 < c)
    (h : IsCompositeSol f c τ φ y z u) {t : ℝ} (ht : t ∈ Ici (0 : ℝ)) :
    ∀ s ∈ Icc (t - τ) t, u (c * (t - s)) t = concat φ y s
\end{leansrc}
\par\medskip\noindent{\small \texttt{Qtsl.JumpDiffusion.}\texttt{\bfseries MpSolution}}\hfill{\footnotesize \path{lean/qtsl-aux/QtslAux/JumpDiffusionMartingaleProblem.lean}, line 191}\par\nopagebreak
\begin{leansrc}
structure MpSolution (P : JdProfile n m) (D : Set (State n))
    (Ω : Type*) [MeasureSpace Ω] where …
\end{leansrc}
\par\medskip\noindent{\small \texttt{Qtsl.JumpDiffusion.MpSolution.}\texttt{\bfseries supermartingale\_of\_gen\_nonpos}}\hfill{\footnotesize \path{lean/qtsl-aux/QtslAux/JumpDiffusionMartingaleProblem.lean}, line 234}\par\nopagebreak
\begin{leansrc}
theorem supermartingale_of_gen_nonpos
    [IsProbabilityMeasure (volume : Measure Ω)]
    (mp : MpSolution P D Ω) (hC2 : IsC2Fn V gradV hessV)
    (hgint : ∀ (ω : Ω) (t : ℝ), IntervalIntegrable
      (fun s => generator P V gradV hessV (mp.X (Real.toNNReal s) ω))
      volume 0 t)
    (hVint : ∀ t : ℝ≥0, Integrable (fun ω => V (mp.X t ω))
      (volume : Measure Ω))
    (hrow : ∀ z ∈ D, generator P V gradV hessV z ≤ 0) :
    Supermartingale (fun (t : ℝ≥0) (ω : Ω) => V (mp.X t ω)) mp.ℱ
      (volume : Measure Ω)
\end{leansrc}
\par\medskip\noindent{\small \texttt{Qtsl.SdeMartingaleProblem.MpSolution.}\texttt{\bfseries toSdeDynkinVille}}\hfill{\footnotesize \path{lean/qtsl-aux/QtslAux/SdeMartingaleProblem.lean}, line 409}\par\nopagebreak
\begin{leansrc}
noncomputable def toSdeDynkinVille
    [IsProbabilityMeasure (volume : Measure Ω)]
    (mp : MpSolution P D Ω) (τ : ℕ → ℝ≥0) (hτ : Monotone τ)
    (hC2 : IsC2Fn V gradV hessV)
    (hgint : ∀ (ω : Ω) (t : ℝ), IntervalIntegrable
      (fun s => generator P V gradV hessV (mp.X (Real.toNNReal s) ω))
      volume 0 t)
    (hVint : ∀ t : ℝ≥0, Integrable (fun ω => V (mp.X t ω))
      (volume : Measure Ω)) :
    SdeDynkinVille P D V gradV hessV Ω where …
\end{leansrc}
\par\medskip\noindent{\small \texttt{Qtsl.JumpDiffusion.}\texttt{\bfseries genPoly\_eval\_eq\_generator}}\hfill{\footnotesize \path{lean/qtsl-aux/QtslAux/JumpDiffusion.lean}, line 269}\par\nopagebreak
\begin{leansrc}
theorem genPoly_eval_eq_generator (P : PolyProfile n m)
    (V : MvPolynomial (Fin n) ℚ) (z : State n) :
    (aeval z (genPoly P V) : ℝ)
      = generator P.toJd
          (fun w => (aeval w V : ℝ))
          (fun w i => (aeval w (pderiv i V) : ℝ))
          (fun w i j => (aeval w (pderiv i (pderiv j V)) : ℝ)) z
\end{leansrc}
\par\medskip\noindent{\small \texttt{Qtsl.SdeMartingaleProblem.}\texttt{\bfseries isC2Fn\_poly}}\hfill{\footnotesize \path{lean/qtsl-aux/QtslAux/SdeMartingaleProblem.lean}, line 285}\par\nopagebreak
\begin{leansrc}
theorem isC2Fn_poly (V : MvPolynomial (Fin n) ℚ) :
    IsC2Fn (fun z : State n => (aeval z V : ℝ))
      (fun z i => (aeval z (pderiv i V) : ℝ))
      (fun z i j => (aeval z (pderiv i (pderiv j V)) : ℝ)) where
  hasFDerivAt_V z
\end{leansrc}
\par\medskip\noindent{\small \texttt{Qtsl.JumpDiffusion.}\texttt{\bfseries mp\_avoid\_row\_of\_genPoly}}\hfill{\footnotesize \path{lean/qtsl-aux/QtslAux/JumpDiffusionMartingaleProblem.lean}, line 326}\par\nopagebreak
\begin{leansrc}
theorem mp_avoid_row_of_genPoly
    [IsProbabilityMeasure (volume : Measure Ω)]
    {P₀ : PolyProfile n m} {D : Set (State n)}
    (V₀ : MvPolynomial (Fin n) ℚ)
    (mp : MpSolution P₀.toJd D Ω)
    (τ : ℕ → ℝ≥0) (hτ : Monotone τ)
    (hgint : ∀ (ω : Ω) (t : ℝ), IntervalIntegrable
      (fun s => (aeval (mp.X (Real.toNNReal s) ω) (genPoly P₀ V₀) : ℝ))
      volume 0 t)
    (hVint : ∀ t : ℝ≥0, Integrable
      (fun ω => (aeval (mp.X t ω) V₀ : ℝ)) (volume : Measure Ω))
    (hrow : ∀ z ∈ D, (aeval z (genPoly P₀ V₀) : ℝ) ≤ 0)
    (hnonneg : ∀ z ∈ D, 0 ≤ (aeval z V₀ : ℝ))
    {Bad : Set (State n)} {a : ℝ}
    (hfloor : ∀ z ∈ Bad, a ≤ (aeval z V₀ : ℝ)) :
    ENNReal.ofReal a * volume {ω | ∃ k, mp.X (τ k) ω ∈ Bad}
      ≤ ENNReal.ofReal (∫ ω, (aeval (mp.X (τ 0) ω) V₀ : ℝ))
\end{leansrc}
\par\medskip\noindent{\small \texttt{Qtsl.JumpDiffusion.}\texttt{\bfseries mp\_avoid\_init\_of\_genPoly}}\hfill{\footnotesize \path{lean/qtsl-aux/QtslAux/JumpDiffusionMartingaleProblem.lean}, line 370}\par\nopagebreak
\begin{leansrc}
theorem mp_avoid_init_of_genPoly
    [IsProbabilityMeasure (volume : Measure Ω)]
    {P₀ : PolyProfile n m} {D : Set (State n)}
    (V₀ : MvPolynomial (Fin n) ℚ)
    (mp : MpSolution P₀.toJd D Ω)
    (τ : ℕ → ℝ≥0) (hτ : Monotone τ)
    (hgint : ∀ (ω : Ω) (t : ℝ), IntervalIntegrable
      (fun s => (aeval (mp.X (Real.toNNReal s) ω) (genPoly P₀ V₀) : ℝ))
      volume 0 t)
    (hVint : ∀ t : ℝ≥0, Integrable
      (fun ω => (aeval (mp.X t ω) V₀ : ℝ)) (volume : Measure Ω))
    (hrow : ∀ z ∈ D, (aeval z (genPoly P₀ V₀) : ℝ) ≤ 0)
    (hnonneg : ∀ z ∈ D, 0 ≤ (aeval z V₀ : ℝ))
    {Bad : Set (State n)} {z0 : State n}
    (hstart : ∀ ω : Ω, mp.X (τ 0) ω = z0)
    (hfloor : ∀ z ∈ Bad, 1 ≤ (aeval z V₀ : ℝ)) :
    volume {ω | ∃ k, mp.X (τ k) ω ∈ Bad}
      ≤ ENNReal.ofReal (aeval z0 V₀ : ℝ)
\end{leansrc}
\par\medskip\noindent{\small \texttt{Qtsl.JumpDiffusion.}\texttt{\bfseries mp\_control}}\hfill{\footnotesize \path{lean/qtsl-aux/QtslAux/JumpDiffusionMartingaleProblem.lean}, line 543}\par\nopagebreak
\begin{leansrc}
theorem mp_control [IsProbabilityMeasure (volume : Measure Ω)]
    (x₀ : State 1) {Bad : Set (State 1)}
    (hfloor : ∀ z ∈ Bad,
      1 ≤ (aeval z (X 0 * X 0 : MvPolynomial (Fin 1) ℚ) : ℝ)) :
    volume {_ω : Ω | ∃ _k : ℕ, x₀ ∈ Bad}
      ≤ ENNReal.ofReal (aeval x₀ (X 0 * X 0 : MvPolynomial (Fin 1) ℚ) : ℝ)
\end{leansrc}
\par\medskip\noindent{\small \texttt{SuFunctor.Adequacy.SpectrumDecl.}\texttt{\bfseries plainTrace\_divergent\_none}}\hfill{\footnotesize \path{lean/su-functor/SuFunctor/Adequacy/PdeRows.lean}, line 549}\par\nopagebreak
\begin{leansrc}
theorem plainTrace_divergent_none (ds : List DriverDecl) :
    (SpectrumDecl.mk ds .divergent).plainTrace = none
\end{leansrc}
\par\medskip\noindent{\small \texttt{SuFunctor.Adequacy.}\texttt{\bfseries isRodeSol\_iff\_integralCurve}}\hfill{\footnotesize \path{lean/su-functor/SuFunctor/Adequacy/Rode.lean}, line 63}\par\nopagebreak
\begin{leansrc}
theorem isRodeSol_iff_integralCurve [CompleteSpace E] {f : Θ → E → E}
    {x₀ : E} {θ : Θ} {y : ℝ → E} (hf : Continuous (f θ)) :
    IsRodeSol f x₀ θ y ↔
      y 0 = x₀ ∧ IsIntegralCurveOn y (fun _ => f θ) (Ici 0)
\end{leansrc}
\par\medskip\noindent{\small \texttt{SuFunctor.Adequacy.}\texttt{\bfseries rode\_draw\_adequacy}}\hfill{\footnotesize \path{lean/su-functor/SuFunctor/Adequacy/Rode.lean}, line 71}\par\nopagebreak
\begin{leansrc}
theorem rode_draw_adequacy [CompleteSpace E] {f : Θ → E → E} {θ : Θ}
    (hf : Continuous (f θ)) (x₀ : E) {a : Crd (ℝ → E)} :
    ∃ α : Cell (looseConj (odeAtoms (f θ) x₀) : Loose a a)
        (suppC (integralCurves (f θ) x₀)) (𝟙 a) (𝟙 a),
      Cell.IsInvertible α
\end{leansrc}
\par\medskip\noindent{\small \texttt{SuFunctor.Adequacy.}\texttt{\bfseries rode\_adequacy}}\hfill{\footnotesize \path{lean/su-functor/SuFunctor/Adequacy/Rode.lean}, line 160}\par\nopagebreak
\begin{leansrc}
theorem rode_adequacy [CompleteSpace E] {f : Θ → E → E} {P : Set Θ}
    (hf : ∀ θ ∈ P, Continuous (f θ)) (x₀ : E) {a : Crd (ℝ → E)} :
    ∃ α : Cell (looseConj (rodeAtoms f P x₀) : Loose a a)
        (suppC (⋃ θ ∈ P, integralCurves (f θ) x₀)) (𝟙 a) (𝟙 a),
      Cell.IsInvertible α
\end{leansrc}
\par\medskip\noindent{\small \texttt{SuFunctor.Adequacy.}\texttt{\bfseries rode\_box\_adequacy}}\hfill{\footnotesize \path{lean/su-functor/SuFunctor/Adequacy/Rode.lean}, line 241}\par\nopagebreak
\begin{leansrc}
theorem rode_box_adequacy [CompleteSpace E] {f : Θ → E → E} {P : Set Θ}
    (hf : ∀ θ ∈ P, Continuous (f θ)) (B : Set E) {a : Crd (ℝ → E)} :
    ∃ α : Cell (looseConj (rodeBoxAtoms f P B) : Loose a a)
        (suppC (⋃ θ ∈ P, ⋃ x₀ ∈ B, integralCurves (f θ) x₀))
        (𝟙 a) (𝟙 a),
      Cell.IsInvertible α
\end{leansrc}
\par\medskip\noindent{\small \texttt{SuFunctor.Adequacy.}\texttt{\bfseries rodeSols\_eqOn\_of\_lipschitzOnWith}}\hfill{\footnotesize \path{lean/su-functor/SuFunctor/Adequacy/Rode.lean}, line 81}\par\nopagebreak
\begin{leansrc}
theorem rodeSols_eqOn_of_lipschitzOnWith {f : Θ → E → E} {K : ℝ≥0}
    {s : Set E} {P : Set Θ} (hL : ∀ θ ∈ P, LipschitzOnWith K (f θ) s)
    {x₀ : E} {θ : Θ} (hθ : θ ∈ P) {y₁ y₂ : ℝ → E}
    (h₁ : IsRodeSol f x₀ θ y₁) (h₂ : IsRodeSol f x₀ θ y₂)
    (hs₁ : ∀ t ∈ Ici (0 : ℝ), y₁ t ∈ s)
    (hs₂ : ∀ t ∈ Ici (0 : ℝ), y₂ t ∈ s) :
    EqOn y₁ y₂ (Ici 0)
\end{leansrc}
\par\medskip\noindent{\small \texttt{SuFunctor.Adequacy.}\texttt{\bfseries su\_rode\_adequacy}}\hfill{\footnotesize \path{lean/su-functor/SuFunctor/Adequacy/Rode.lean}, line 281}\par\nopagebreak
\begin{leansrc}
theorem su_rode_adequacy [CompleteSpace E] {f : Θ → E → E} {P : Set Θ}
    (hf : ∀ θ ∈ P, Continuous (f θ)) (x₀ : E) (μ : SubPMF (ℝ → E)) :
    idealPoint μ ∈ (rodeDnf f P x₀).modelClass ↔
      suppC (⋃ θ ∈ P, integralCurves (f θ) x₀) μ
\end{leansrc}
\par\medskip\noindent{\small \texttt{Qtsl.JumpDiffusion.}\texttt{\bfseries JdProfile}}\hfill{\footnotesize \path{lean/qtsl-aux/QtslAux/JumpDiffusion.lean}, line 121}\par\nopagebreak
\begin{leansrc}
structure JdProfile (n m : ℕ) where …
\end{leansrc}
\par\medskip\noindent{\small \texttt{Qtsl.JumpDiffusion.}\texttt{\bfseries generator\_one\_dirac}}\hfill{\footnotesize \path{lean/qtsl-aux/QtslAux/JumpDiffusion.lean}, line 178}\par\nopagebreak
\begin{leansrc}
theorem generator_one_dirac (P : JdProfile n m) (lam : State n → ℝ)
    (g : State n → State n) (hkill : ∀ z, P.kill z = 0)
    (hlegs : P.legs = [(lam, [((1 : ℝ), g)])])
    (V : State n → ℝ) (gradV : State n → State n)
    (hessV : State n → Fin n → Fin n → ℝ) (z : State n) :
    generator P V gradV hessV z
      = SdeGenerator.generator ⟨P.f, P.σ, lam, g⟩ V gradV hessV z
\end{leansrc}
\par\medskip\noindent{\small \texttt{Qtsl.JumpDiffusion.}\texttt{\bfseries toJd\_generator\_one\_dirac}}\hfill{\footnotesize \path{lean/qtsl-aux/QtslAux/JumpDiffusion.lean}, line 293}\par\nopagebreak
\begin{leansrc}
theorem toJd_generator_one_dirac (P : PolyProfile n m)
    (lam : MvPolynomial (Fin n) ℚ) (g : Fin n → MvPolynomial (Fin n) ℚ)
    (hkill : P.kill = 0) (hlegs : P.legs = [(lam, [((1 : ℚ), g)])])
    (V : State n → ℝ) (gradV : State n → State n)
    (hessV : State n → Fin n → Fin n → ℝ) (z : State n) :
    generator P.toJd V gradV hessV z
      = SdeGenerator.generator (P.toSde lam g) V gradV hessV z
\end{leansrc}
\par\medskip\noindent{\small \texttt{Qtsl.JumpDiffusion.}\texttt{\bfseries transferSum\_sub\_self}}\hfill{\footnotesize \path{lean/qtsl-aux/QtslAux/JumpDiffusion.lean}, line 101}\par\nopagebreak
\begin{leansrc}
theorem transferSum_sub_self (μ : MixtureR n) (V : State n → ℝ)
    (z : State n) :
    transferSum μ V z - V z
      = (μ.map fun wg => wg.1 * (V (wg.2 z) - V z)).sum
        - (1 - (μ.map fun wg => wg.1).sum) * V z
\end{leansrc}
\par\medskip\noindent{\small \texttt{Qtsl.JumpDiffusion.}\texttt{\bfseries two\_leg\_control}}\hfill{\footnotesize \path{lean/qtsl-aux/QtslAux/JumpDiffusion.lean}, line 329}\par\nopagebreak
\begin{leansrc}
theorem two_leg_control (z : State 1) :
    4 * generator twoLeg.toJd (fun w => (w 0) ^ 2)
        (fun w _ => 2 * w 0) (fun _ _ _ => 2) z
      = 4 + (z 0) ^ 2 - 3 * (z 0) ^ 4
\end{leansrc}
\par\medskip\noindent{\small \texttt{Qtsl.JumpDiffusion.}\texttt{\bfseries JdDynkinVille}}\hfill{\footnotesize \path{lean/qtsl-aux/QtslAux/JumpDiffusionMartingaleProblem.lean}, line 82}\par\nopagebreak
\begin{leansrc}
structure JdDynkinVille (P : JdProfile n m) (D : Set (State n))
    (V : State n → ℝ) (gradV : State n → State n)
    (hessV : State n → Fin n → Fin n → ℝ)
    (Ω : Type*) [MeasureSpace Ω] where …
\end{leansrc}
\par\medskip\noindent{\small \texttt{Qtsl.JumpDiffusion.}\texttt{\bfseries generator\_toJd}}\hfill{\footnotesize \path{lean/qtsl-aux/QtslAux/JumpDiffusionMartingaleProblem.lean}, line 431}\par\nopagebreak
\begin{leansrc}
theorem generator_toJd (P : SdeProfile n m) (V : State n → ℝ)
    (gradV : State n → State n) (hessV : State n → Fin n → Fin n → ℝ)
    (z : State n) :
    generator P.toJd V gradV hessV z
      = SdeGenerator.generator P V gradV hessV z
\end{leansrc}
\par\medskip\noindent{\small \texttt{Qtsl.JumpDiffusion.\_root\_.Qtsl.SdeMartingaleProblem.MpSolution.}\texttt{\bfseries toJd}}\hfill{\footnotesize \path{lean/qtsl-aux/QtslAux/JumpDiffusionMartingaleProblem.lean}, line 445}\par\nopagebreak
\begin{leansrc}
noncomputable def _root_.Qtsl.SdeMartingaleProblem.MpSolution.toJd
    {P : SdeProfile n m} {D : Set (State n)} {Ω : Type*} [MeasureSpace Ω]
    (mp : SdeMartingaleProblem.MpSolution P D Ω) :
    MpSolution P.toJd D Ω where …
\end{leansrc}
\par\medskip\noindent{\small \texttt{SuFunctor.Adequacy.}\texttt{\bfseries satAtoms\_iff}}\hfill{\footnotesize \path{lean/su-functor/SuFunctor/Adequacy/Dtmc.lean}, line 598}\par\nopagebreak
\begin{leansrc}
theorem satAtoms_iff (hR : R.WF) (μ : SubPMF (Traj S k)) :
    looseConj (atoms R k) μ ↔ μ = rowsLaw R hR k
\end{leansrc}
\par\medskip\noindent{\small \texttt{SuFunctor.Adequacy.}\texttt{\bfseries dtmc\_adequacy}}\hfill{\footnotesize \path{lean/su-functor/SuFunctor/Adequacy/Dtmc.lean}, line 622}\par\nopagebreak
\begin{leansrc}
theorem dtmc_adequacy (M : FinDTMC S) (k : ℕ) {a : Crd (Traj S k)} :
    ∃ α : Cell (looseConj (atoms (compile M) k) : Loose a a)
        (singletonC (M.pathLaw k)) (𝟙 a) (𝟙 a),
      Cell.IsInvertible α
\end{leansrc}
\par\medskip\noindent{\small \texttt{SuFunctor.Adequacy.FinDTMC.}\texttt{\bfseries pathLaw\_mass}}\hfill{\footnotesize \path{lean/su-functor/SuFunctor/Adequacy/Dtmc.lean}, line 259}\par\nopagebreak
\begin{leansrc}
theorem FinDTMC.pathLaw_mass (M : FinDTMC S) (k : ℕ) :
    (M.pathLaw k).mass = 1
\end{leansrc}
\par\medskip\noindent{\small \texttt{SuFunctor.Adequacy.FinDTMC.}\texttt{\bfseries trajPMF\_eq}}\hfill{\footnotesize \path{lean/su-functor/SuFunctor/Adequacy/Dtmc.lean}, line 743}\par\nopagebreak
\begin{leansrc}
theorem FinDTMC.trajPMF_eq :
    ∀ (k : ℕ) (ω : Traj S k),
      (M.trajPMF k) ω = q2e ((M.pathLaw k).wt ω)
\end{leansrc}
\par\medskip\noindent{\small \texttt{SuFunctor.Adequacy.}\texttt{\bfseries ctmc\_jump\_adequacy}}\hfill{\footnotesize \path{lean/su-functor/SuFunctor/Adequacy/CtmcJumpChain.lean}, line 371}\par\nopagebreak
\begin{leansrc}
theorem ctmc_jump_adequacy (M : FinCTMC S) (k : ℕ) {a : Crd (Traj S k)} :
    ∃ α : Cell
        (looseConj (atoms (jumpRows (compileCtmc M)) k) : Loose a a)
        (singletonC (M.jumpLaw k)) (𝟙 a) (𝟙 a),
      Cell.IsInvertible α
\end{leansrc}
\par\medskip\noindent{\small \texttt{SuFunctor.Adequacy.}\texttt{\bfseries compile\_jump\_comm}}\hfill{\footnotesize \path{lean/su-functor/SuFunctor/Adequacy/CtmcJumpChain.lean}, line 352}\par\nopagebreak
\begin{leansrc}
theorem compile_jump_comm (M : FinCTMC S) :
    jumpRows (compileCtmc M) = compile M.jumpChain
\end{leansrc}
\par\medskip\noindent{\small \texttt{SuFunctor.Adequacy.}\texttt{\bfseries jumpRows\_wf}}\hfill{\footnotesize \path{lean/su-functor/SuFunctor/Adequacy/CtmcJumpChain.lean}, line 338}\par\nopagebreak
\begin{leansrc}
theorem jumpRows_wf {R : CtmcRows S} (hR : R.WF) : (jumpRows R).WF where
  outcome_nonneg p
\end{leansrc}
\par\medskip\noindent{\small \texttt{SuFunctor.Adequacy.FinCTMC.}\texttt{\bfseries jumpChain\_step\_mul\_exit}}\hfill{\footnotesize \path{lean/su-functor/SuFunctor/Adequacy/CtmcJumpChain.lean}, line 210}\par\nopagebreak
\begin{leansrc}
theorem FinCTMC.jumpChain_step_mul_exit (M : FinCTMC S) (s s' : S) :
    M.jumpChain.step s s' * M.exitRate s = M.offRate s s'
\end{leansrc}
\par\medskip\noindent{\small \texttt{SuFunctor.Adequacy.}\texttt{\bfseries compileCtmc\_diag\_blind}}\hfill{\footnotesize \path{lean/su-functor/SuFunctor/Adequacy/CtmcJumpChain.lean}, line 385}\par\nopagebreak
\begin{leansrc}
theorem compileCtmc_diag_blind (M M' : FinCTMC S)
    (hoff : ∀ s s', s ≠ s' → M.rate s s' = M'.rate s s')
    (hinit : M.init = M'.init) : compileCtmc M = compileCtmc M'
\end{leansrc}
\par\medskip\noindent{\small \texttt{SuFunctor.Adequacy.FinCTMC.}\texttt{\bfseries jumpLaw\_diag\_blind}}\hfill{\footnotesize \path{lean/su-functor/SuFunctor/Adequacy/CtmcJumpChain.lean}, line 407}\par\nopagebreak
\begin{leansrc}
theorem FinCTMC.jumpLaw_diag_blind (M M' : FinCTMC S)
    (hoff : ∀ s s', s ≠ s' → M.rate s s' = M'.rate s s')
    (hinit : M.init = M'.init) (k : ℕ) :
    M.jumpLaw k = M'.jumpLaw k
\end{leansrc}
\par\medskip\noindent{\small \texttt{SuFunctor.Adequacy.FinCTMC.}\texttt{\bfseries jumpTrajPMF\_eq}}\hfill{\footnotesize \path{lean/su-functor/SuFunctor/Adequacy/CtmcJumpChain.lean}, line 440}\par\nopagebreak
\begin{leansrc}
theorem FinCTMC.jumpTrajPMF_eq (M : FinCTMC S) (k : ℕ) (ω : Traj S k) :
    M.jumpTrajPMF k ω = q2e ((M.jumpLaw k).wt ω)
\end{leansrc}
\par\medskip\noindent{\small \texttt{SuFunctor.Adequacy.}\texttt{\bfseries mdp\_adequacy}}\hfill{\footnotesize \path{lean/su-functor/SuFunctor/Adequacy/Mdp.lean}, line 284}\par\nopagebreak
\begin{leansrc}
theorem mdp_adequacy (M : FinMDP S A) (π : Policy S A) (k : ℕ)
    {a : Crd (Traj S k)} :
    ∃ α : Cell
        (looseConj (atoms (closeRows (compileMdp M) (compilePolicy π)) k) :
          Loose a a)
        (singletonC (M.pathLaw π k)) (𝟙 a) (𝟙 a),
      Cell.IsInvertible α
\end{leansrc}
\par\medskip\noindent{\small \texttt{SuFunctor.Adequacy.}\texttt{\bfseries compile\_close\_comm}}\hfill{\footnotesize \path{lean/su-functor/SuFunctor/Adequacy/Mdp.lean}, line 266}\par\nopagebreak
\begin{leansrc}
theorem compile_close_comm (M : FinMDP S A) (π : Policy S A) :
    closeRows (compileMdp M) (compilePolicy π) =
      compile (M.closedLoop π)
\end{leansrc}
\par\medskip\noindent{\small \texttt{SuFunctor.Adequacy.FinMDP.}\texttt{\bfseries closedLoop\_stepPMF}}\hfill{\footnotesize \path{lean/su-functor/SuFunctor/Adequacy/Mdp.lean}, line 334}\par\nopagebreak
\begin{leansrc}
theorem FinMDP.closedLoop_stepPMF (M : FinMDP S A) (π : Policy S A)
    (s : S) :
    (M.closedLoop π).stepPMF s =
      (π.actPMF s).bind (fun a => M.stepPMF s a)
\end{leansrc}
\par\medskip\noindent{\small \texttt{SuFunctor.}\texttt{\bfseries unionSingletons}}\hfill{\footnotesize \path{lean/su-functor/SuFunctor/Adequacy/EnvelopeUnion.lean}, line 104}\par\nopagebreak
\begin{leansrc}
def unionSingletons (μs : ι → SubPMF V) : SubPMF V → Prop
\end{leansrc}
\par\medskip\noindent{\small \texttt{SuFunctor.Adequacy.FinMDP.}\texttt{\bfseries demonicCredal}}\hfill{\footnotesize \path{lean/su-functor/SuFunctor/Adequacy/MdpDemonic.lean}, line 81}\par\nopagebreak
\begin{leansrc}
def FinMDP.demonicCredal (M : FinMDP S A) (k : ℕ) :
    SubPMF (Traj S k) → Prop
\end{leansrc}
\par\medskip\noindent{\small \texttt{SuFunctor.Adequacy.}\texttt{\bfseries mdp\_demonic\_adequacy}}\hfill{\footnotesize \path{lean/su-functor/SuFunctor/Adequacy/MdpDemonic.lean}, line 157}\par\nopagebreak
\begin{leansrc}
theorem mdp_demonic_adequacy (M : FinMDP S A) (k : ℕ)
    {a : Crd (Traj S k)} :
    ∃ α : Cell (demonicC (compileMdp M) k : Loose a a)
        (M.demonicLaws k) (𝟙 a) (𝟙 a),
      Cell.IsInvertible α
\end{leansrc}
\par\medskip\noindent{\small \texttt{SuFunctor.Adequacy.}\texttt{\bfseries demonicC}}\hfill{\footnotesize \path{lean/su-functor/SuFunctor/Adequacy/MdpDemonic.lean}, line 92}\par\nopagebreak
\begin{leansrc}
def demonicC (R : MdpRows S A) (k : ℕ) : SubPMF (Traj S k) → Prop
\end{leansrc}
\par\medskip\noindent{\small \texttt{SuFunctor.Adequacy.}\texttt{\bfseries demonicC\_iff}}\hfill{\footnotesize \path{lean/su-functor/SuFunctor/Adequacy/MdpDemonic.lean}, line 128}\par\nopagebreak
\begin{leansrc}
theorem demonicC_iff (M : FinMDP S A) (k : ℕ) (μ : SubPMF (Traj S k)) :
    demonicC (compileMdp M) k μ ↔ M.demonicLaws k μ
\end{leansrc}
\par\medskip\noindent{\small \texttt{SuFunctor.Adequacy.}\texttt{\bfseries demonic\_maxSlice}}\hfill{\footnotesize \path{lean/su-functor/SuFunctor/Adequacy/MdpDemonic.lean}, line 184}\par\nopagebreak
\begin{leansrc}
theorem demonic_maxSlice (M : FinMDP S A) (k : ℕ)
    (ν : SubPMF (Traj S k)) :
    maxSlice (M.demonicCredal k) ν ↔ M.demonicLaws k ν
\end{leansrc}
\par\medskip\noindent{\small \texttt{SuFunctor.Adequacy.}\texttt{\bfseries satCredal\_lower\_demonicCredal}}\hfill{\footnotesize \path{lean/su-functor/SuFunctor/Adequacy/MdpDemonic.lean}, line 205}\par\nopagebreak
\begin{leansrc}
theorem satCredal_lower_demonicCredal (M : FinMDP S A) (k : ℕ)
    (c : SubPMF (Traj S k) → Prop) :
    SatCredal .lower (M.demonicCredal k) c ↔
      ∀ π : Policy S A, c (M.pathLaw π k)
\end{leansrc}
\par\medskip\noindent{\small \texttt{SuFunctor.Adequacy.}\texttt{\bfseries satCredal\_upper\_demonicCredal}}\hfill{\footnotesize \path{lean/su-functor/SuFunctor/Adequacy/MdpDemonic.lean}, line 214}\par\nopagebreak
\begin{leansrc}
theorem satCredal_upper_demonicCredal (M : FinMDP S A) (k : ℕ)
    {c : SubPMF (Traj S k) → Prop} (hc : DownClosed c) :
    SatCredal .upper (M.demonicCredal k) c ↔
      ∀ π : Policy S A, c (M.pathLaw π k)
\end{leansrc}
\par\medskip\noindent{\small \texttt{SuFunctor.Adequacy.}\texttt{\bfseries not\_satCredal\_upper\_demonicCredal\_geC}}\hfill{\footnotesize \path{lean/su-functor/SuFunctor/Adequacy/MdpDemonic.lean}, line 265}\par\nopagebreak
\begin{leansrc}
theorem not_satCredal_upper_demonicCredal_geC (M : FinMDP S A) (k : ℕ)
    [Nonempty A] {E : Set (Traj S k)} {q : ℚ} (hq : 0 < q) :
    ¬ SatCredal .upper (M.demonicCredal k) (geC E q)
\end{leansrc}
\par\medskip\noindent{\small \texttt{SuFunctor.Adequacy.}\texttt{\bfseries isEmpty\_policy\_iff}}\hfill{\footnotesize \path{lean/su-functor/SuFunctor/Adequacy/MdpDemonic.lean}, line 290}\par\nopagebreak
\begin{leansrc}
theorem isEmpty_policy_iff (M : FinMDP S A) :
    IsEmpty (Policy S A) ↔ IsEmpty A
\end{leansrc}
\par\medskip\noindent{\small \texttt{SuFunctor.}\texttt{\bfseries not\_maxSlice\_unionRays\_halfPair}}\hfill{\footnotesize \path{lean/su-functor/SuFunctor/Adequacy/EnvelopeUnion.lean}, line 316}\par\nopagebreak
\begin{leansrc}
theorem not_maxSlice_unionRays_halfPair (μ : SubPMF V) :
    ¬ maxSlice (unionRays (halfPair μ)) (halfPair μ false)
\end{leansrc}
\par\medskip\noindent{\small \texttt{SuFunctor.Adequacy.}\texttt{\bfseries ma\_adequacy}}\hfill{\footnotesize \path{lean/su-functor/SuFunctor/Adequacy/MarkovAutomaton.lean}, line 839}\par\nopagebreak
\begin{leansrc}
theorem ma_adequacy (M : FinMA S I A) (σ : Policy I A) (r : I → ℕ)
    (hr : ZenoFree (M.imStep σ) r) (k : ℕ) {a : Crd (Traj S k)} :
    ∃ α : Cell
        (looseConj (atoms (jumpRows
          (closeMaRows (compileMa M) (compilePolicy σ) (absFuel r))) k) :
          Loose a a)
        (singletonC (M.pathLaw σ r hr k)) (𝟙 a) (𝟙 a),
      Cell.IsInvertible α
\end{leansrc}
\par\medskip\noindent{\small \texttt{SuFunctor.Adequacy.}\texttt{\bfseries compile\_closeMa\_comm}}\hfill{\footnotesize \path{lean/su-functor/SuFunctor/Adequacy/MarkovAutomaton.lean}, line 819}\par\nopagebreak
\begin{leansrc}
theorem compile_closeMa_comm (M : FinMA S I A) (σ : Policy I A)
    (r : I → ℕ) (hr : ZenoFree (M.imStep σ) r) :
    closeMaRows (compileMa M) (compilePolicy σ) (absFuel r)
      = compileCtmc (M.close σ r hr)
\end{leansrc}
\par\medskip\noindent{\small \texttt{SuFunctor.Adequacy.}\texttt{\bfseries absorbIter\_mass}}\hfill{\footnotesize \path{lean/su-functor/SuFunctor/Adequacy/MarkovAutomaton.lean}, line 246}\par\nopagebreak
\begin{leansrc}
theorem absorbIter_mass (f : I → S ⊕ I → ℚ) (r : I → ℕ)
    (hr : ZenoFree f r) (hrow : ∀ i, ∑ y, f i y = 1) :
    ∀ (n : ℕ) (i : I), r i < n →
      ∑ m, absorbIter f n (Sum.inr i) m = 1
  | 0, _, hn => absurd hn (Nat.not_lt_zero _)
  | n + 1, i, hn => by
    simp_rw [absorbIter_succ_inr]
    rw [Finset.sum_comm]
    simp_rw [← Finset.mul_sum]
    have hterm : ∀ y, f i y * (∑ m, absorbIter f n y m) = f i y * 1
\end{leansrc}
\par\medskip\noindent{\small \texttt{SuFunctor.Adequacy.}\texttt{\bfseries absorbIter\_unique}}\hfill{\footnotesize \path{lean/su-functor/SuFunctor/Adequacy/MarkovAutomaton.lean}, line 276}\par\nopagebreak
\begin{leansrc}
theorem absorbIter_unique (f : I → S ⊕ I → ℚ) (r : I → ℕ)
    (hr : ZenoFree f r) (g : S ⊕ I → S → ℚ)
    (hg_inl : ∀ m' m, g (Sum.inl m') m = hitMass (Sum.inl m' : S ⊕ I) m)
    (hg_step : ∀ i m, g (Sum.inr i) m = ∑ y, f i y * g y m) :
    ∀ (n : ℕ) (i : I) (m : S), r i < n →
      g (Sum.inr i) m = absorbIter f n (Sum.inr i) m
  | 0, _, _, hn => absurd hn (Nat.not_lt_zero _)
  | n + 1, i, m, hn => by
    rw [hg_step, absorbIter_succ_inr]
    refine Finset.sum_congr rfl fun y _ => ?_
    by_cases hz : f i y = 0
    · rw [hz, zero_mul, zero_mul]
    · congr 1
      cases y with
      | inl m' => rw [hg_inl, absorbIter_inl]
      | inr j =>
        exact absorbIter_unique f r hr g hg_inl hg_step n j m
          (lt_of_lt_of_le (hr i j hz) (Nat.lt_succ_iff.mp hn))

def absFuel (r : I → ℕ) : ℕ
\end{leansrc}
\par\medskip\noindent{\small \texttt{SuFunctor.Adequacy.}\texttt{\bfseries absorbIter\_stab}}\hfill{\footnotesize \path{lean/su-functor/SuFunctor/Adequacy/MarkovAutomaton.lean}, line 223}\par\nopagebreak
\begin{leansrc}
theorem absorbIter_stab (f : I → S ⊕ I → ℚ) (r : I → ℕ)
    (hr : ZenoFree f r) (i : I) (m : S) {n n' : ℕ}
    (h : r i < n) (h' : r i < n') :
    absorbIter f n (Sum.inr i) m = absorbIter f n' (Sum.inr i) m
\end{leansrc}
\par\medskip\noindent{\small \texttt{SuFunctor.Adequacy.FinMA.}\texttt{\bfseries close\_offRate\_factor}}\hfill{\footnotesize \path{lean/su-functor/SuFunctor/Adequacy/MarkovAutomaton.lean}, line 540}\par\nopagebreak
\begin{leansrc}
theorem FinMA.close_offRate_factor (M : FinMA S I A) (σ : Policy I A)
    (r : I → ℕ) (hr : ZenoFree (M.imStep σ) r) {s s' : S}
    (hne : s ≠ s') :
    (M.close σ r hr).offRate s s'
      = M.exitRate s * ∑ x, M.preChain.jumpChain.step (Sum.inl s) x
          * M.absorb σ r x s'
\end{leansrc}
\par\medskip\noindent{\small \texttt{SuFunctor.Adequacy.FinMA.}\texttt{\bfseries close\_exitRate\_thinning}}\hfill{\footnotesize \path{lean/su-functor/SuFunctor/Adequacy/MarkovAutomaton.lean}, line 562}\par\nopagebreak
\begin{leansrc}
theorem FinMA.close_exitRate_thinning (M : FinMA S I A)
    (σ : Policy I A) (r : I → ℕ) (hr : ZenoFree (M.imStep σ) r)
    (s : S) :
    (M.close σ r hr).exitRate s
      = M.exitRate s - ∑ x, M.offRate s x * M.absorb σ r x s
\end{leansrc}
\par\medskip\noindent{\small \texttt{SuFunctor.Adequacy.FinMA.}\texttt{\bfseries close\_exitRate\_le}}\hfill{\footnotesize \path{lean/su-functor/SuFunctor/Adequacy/MarkovAutomaton.lean}, line 594}\par\nopagebreak
\begin{leansrc}
theorem FinMA.close_exitRate_le (M : FinMA S I A) (σ : Policy I A)
    (r : I → ℕ) (hr : ZenoFree (M.imStep σ) r) (s : S) :
    (M.close σ r hr).exitRate s ≤ M.exitRate s
\end{leansrc}
\par\medskip\noindent{\small \texttt{SuFunctor.Adequacy.}\texttt{\bfseries compileMa\_diag\_blind}}\hfill{\footnotesize \path{lean/su-functor/SuFunctor/Adequacy/MarkovAutomaton.lean}, line 713}\par\nopagebreak
\begin{leansrc}
theorem compileMa_diag_blind (M M' : FinMA S I A)
    (hoff : ∀ s x, Sum.inl s ≠ x → M.rate s x = M'.rate s x)
    (hbr : M.branch = M'.branch) (hinit : M.init = M'.init) :
    compileMa M = compileMa M'
\end{leansrc}
\par\medskip\noindent{\small \texttt{SuFunctor.Adequacy.FinMA.}\texttt{\bfseries trajPMF\_eq}}\hfill{\footnotesize \path{lean/su-functor/SuFunctor/Adequacy/MarkovAutomaton.lean}, line 869}\par\nopagebreak
\begin{leansrc}
theorem FinMA.trajPMF_eq (M : FinMA S I A) (σ : Policy I A)
    (r : I → ℕ) (hr : ZenoFree (M.imStep σ) r) (k : ℕ)
    (ω : Traj S k) :
    M.trajPMF σ r hr k ω = q2e ((M.pathLaw σ r hr k).wt ω)
\end{leansrc}
\par\medskip\noindent{\small \texttt{SuFunctor.Adequacy.}\texttt{\bfseries ma\_demonic\_adequacy}}\hfill{\footnotesize \path{lean/su-functor/SuFunctor/Adequacy/MaDemonic.lean}, line 299}\par\nopagebreak
\begin{leansrc}
theorem ma_demonic_adequacy (M : FinMA S I A) {r : I → ℕ}
    (hb : M.BranchZenoFree r) (k : ℕ) {a : Crd (Traj S k)} :
    ∃ α : Cell (maDemonicC (compileMa M) (absFuel r) k : Loose a a)
        (M.demonicLaws r hb k) (𝟙 a) (𝟙 a),
      Cell.IsInvertible α
\end{leansrc}
\par\medskip\noindent{\small \texttt{SuFunctor.Adequacy.}\texttt{\bfseries maDemonicC}}\hfill{\footnotesize \path{lean/su-functor/SuFunctor/Adequacy/MaDemonic.lean}, line 238}\par\nopagebreak
\begin{leansrc}
def maDemonicC (R : MaRows S I A) (n k : ℕ) :
    SubPMF (Traj S k) → Prop
\end{leansrc}
\par\medskip\noindent{\small \texttt{SuFunctor.Adequacy.}\texttt{\bfseries maDemonicC\_iff}}\hfill{\footnotesize \path{lean/su-functor/SuFunctor/Adequacy/MaDemonic.lean}, line 262}\par\nopagebreak
\begin{leansrc}
theorem maDemonicC_iff (M : FinMA S I A) {r : I → ℕ}
    (hb : M.BranchZenoFree r) (k : ℕ) (μ : SubPMF (Traj S k)) :
    maDemonicC (compileMa M) (absFuel r) k μ ↔
      M.demonicLaws r hb k μ
\end{leansrc}
\par\medskip\noindent{\small \texttt{SuFunctor.Adequacy.FinMA.}\texttt{\bfseries BranchZenoFree}}\hfill{\footnotesize \path{lean/su-functor/SuFunctor/Adequacy/MaDemonic.lean}, line 122}\par\nopagebreak
\begin{leansrc}
def FinMA.BranchZenoFree (M : FinMA S I A) (r : I → ℕ) : Prop
\end{leansrc}
\par\medskip\noindent{\small \texttt{SuFunctor.Adequacy.}\texttt{\bfseries zenoFree\_of\_branch}}\hfill{\footnotesize \path{lean/su-functor/SuFunctor/Adequacy/MarkovAutomaton.lean}, line 411}\par\nopagebreak
\begin{leansrc}
theorem zenoFree_of_branch (M : FinMA S I A) (r : I → ℕ)
    (hb : ∀ i a j, M.branch i a (Sum.inr j) ≠ 0 → r j < r i)
    (σ : Policy I A) : ZenoFree (M.imStep σ) r
\end{leansrc}
\par\medskip\noindent{\small \texttt{SuFunctor.Adequacy.FinMA.}\texttt{\bfseries pathLaw\_rank\_irrel}}\hfill{\footnotesize \path{lean/su-functor/SuFunctor/Adequacy/MaDemonic.lean}, line 169}\par\nopagebreak
\begin{leansrc}
theorem FinMA.pathLaw_rank_irrel (M : FinMA S I A) (σ : Policy I A)
    {r r' : I → ℕ} (hr : ZenoFree (M.imStep σ) r)
    (hr' : ZenoFree (M.imStep σ) r') (k : ℕ) :
    M.pathLaw σ r hr k = M.pathLaw σ r' hr' k
\end{leansrc}
\par\medskip\noindent{\small \texttt{SuFunctor.Adequacy.FinMA.}\texttt{\bfseries demonicLaws\_rank\_irrel}}\hfill{\footnotesize \path{lean/su-functor/SuFunctor/Adequacy/MaDemonic.lean}, line 207}\par\nopagebreak
\begin{leansrc}
theorem FinMA.demonicLaws_rank_irrel (M : FinMA S I A) {r r' : I → ℕ}
    (hb : M.BranchZenoFree r) (hb' : M.BranchZenoFree r') (k : ℕ) :
    M.demonicLaws r hb k = M.demonicLaws r' hb' k
\end{leansrc}
\par\medskip\noindent{\small \texttt{SuFunctor.Adequacy.}\texttt{\bfseries ma\_demonic\_maxSlice}}\hfill{\footnotesize \path{lean/su-functor/SuFunctor/Adequacy/MaDemonic.lean}, line 326}\par\nopagebreak
\begin{leansrc}
theorem ma_demonic_maxSlice (M : FinMA S I A) {r : I → ℕ}
    (hb : M.BranchZenoFree r) (k : ℕ) (ν : SubPMF (Traj S k)) :
    maxSlice (M.demonicCredal r hb k) ν ↔ M.demonicLaws r hb k ν
\end{leansrc}
\par\medskip\noindent{\small \texttt{SuFunctor.Adequacy.}\texttt{\bfseries satCredal\_lower\_maDemonicCredal}}\hfill{\footnotesize \path{lean/su-functor/SuFunctor/Adequacy/MaDemonic.lean}, line 349}\par\nopagebreak
\begin{leansrc}
theorem satCredal_lower_maDemonicCredal (M : FinMA S I A) {r : I → ℕ}
    (hb : M.BranchZenoFree r) (k : ℕ) (c : SubPMF (Traj S k) → Prop) :
    SatCredal .lower (M.demonicCredal r hb k) c ↔
      ∀ σ : Policy I A, c (M.pathLaw σ r (hb.zenoFree σ) k)
\end{leansrc}
\par\medskip\noindent{\small \texttt{SuFunctor.Adequacy.}\texttt{\bfseries satCredal\_upper\_maDemonicCredal}}\hfill{\footnotesize \path{lean/su-functor/SuFunctor/Adequacy/MaDemonic.lean}, line 359}\par\nopagebreak
\begin{leansrc}
theorem satCredal_upper_maDemonicCredal (M : FinMA S I A) {r : I → ℕ}
    (hb : M.BranchZenoFree r) (k : ℕ) {c : SubPMF (Traj S k) → Prop}
    (hc : DownClosed c) :
    SatCredal .upper (M.demonicCredal r hb k) c ↔
      ∀ σ : Policy I A, c (M.pathLaw σ r (hb.zenoFree σ) k)
\end{leansrc}
\par\medskip\noindent{\small \texttt{SuFunctor.Adequacy.Policy.}\texttt{\bfseries nonempty\_of\_isEmpty}}\hfill{\footnotesize \path{lean/su-functor/SuFunctor/Adequacy/MaDemonic.lean}, line 406}\par\nopagebreak
\begin{leansrc}
theorem Policy.nonempty_of_isEmpty {X : Type} [IsEmpty X] :
    Nonempty (Policy X A)
\end{leansrc}
\par\medskip\noindent{\small \texttt{SuFunctor.Adequacy.Policy.}\texttt{\bfseries isEmpty\_iff}}\hfill{\footnotesize \path{lean/su-functor/SuFunctor/Adequacy/MaDemonic.lean}, line 418}\par\nopagebreak
\begin{leansrc}
theorem Policy.isEmpty_iff {X : Type} [Nonempty X] [DecidableEq A] :
    IsEmpty (Policy X A) ↔ IsEmpty A
\end{leansrc}
\par\medskip\noindent{\small \texttt{SuFunctor.Adequacy.}\texttt{\bfseries not\_satCredal\_upper\_maDemonicCredal\_geC}}\hfill{\footnotesize \path{lean/su-functor/SuFunctor/Adequacy/MaDemonic.lean}, line 438}\par\nopagebreak
\begin{leansrc}
theorem not_satCredal_upper_maDemonicCredal_geC (M : FinMA S I A)
    {r : I → ℕ} (hb : M.BranchZenoFree r) (k : ℕ)
    [hp : Nonempty (Policy I A)] {E : Set (Traj S k)} {q : ℚ}
    (hq : 0 < q) :
    ¬ SatCredal .upper (M.demonicCredal r hb k) (geC E q)
\end{leansrc}
\par\medskip\noindent{\small \texttt{SuFunctor.Adequacy.FinMA.}\texttt{\bfseries nonempty\_markovian}}\hfill{\footnotesize \path{lean/su-functor/SuFunctor/Adequacy/MaDemonic.lean}, line 474}\par\nopagebreak
\begin{leansrc}
theorem FinMA.nonempty_markovian (M : FinMA S I A) [Nonempty I]
    [Nonempty A] {r : I → ℕ} (hb : M.BranchZenoFree r) : Nonempty S
\end{leansrc}
\par\medskip\noindent{\small \texttt{SuFunctor.Adequacy.FinMA.}\texttt{\bfseries CertScheduler}}\hfill{\footnotesize \path{lean/su-functor/SuFunctor/Adequacy/MaDemonic.lean}, line 504}\par\nopagebreak
\begin{leansrc}
structure FinMA.CertScheduler (M : FinMA S I A) where …
\end{leansrc}
\par\medskip\noindent{\small \texttt{SuFunctor.Adequacy.FinMA.}\texttt{\bfseries demonicLawsCert}}\hfill{\footnotesize \path{lean/su-functor/SuFunctor/Adequacy/MaDemonic.lean}, line 515}\par\nopagebreak
\begin{leansrc}
def FinMA.demonicLawsCert (M : FinMA S I A) (k : ℕ) :
    SubPMF (Traj S k) → Prop
\end{leansrc}
\par\medskip\noindent{\small \texttt{SuFunctor.Adequacy.FinMA.}\texttt{\bfseries demonicLawsCert\_iff}}\hfill{\footnotesize \path{lean/su-functor/SuFunctor/Adequacy/MaDemonic.lean}, line 531}\par\nopagebreak
\begin{leansrc}
theorem FinMA.demonicLawsCert_iff (M : FinMA S I A) {r : I → ℕ}
    (hb : M.BranchZenoFree r) (k : ℕ) (ν : SubPMF (Traj S k)) :
    M.demonicLawsCert k ν ↔ M.demonicLaws r hb k ν
\end{leansrc}
\par\medskip\noindent{\small \texttt{SuFunctor.Adequacy.Bag.}\texttt{\bfseries sortTuple\_apply\_eq\_inf\_sup}}\hfill{\footnotesize \path{lean/su-functor/SuFunctor/Adequacy/Bag.lean}, line 128}\par\nopagebreak
\begin{leansrc}
theorem sortTuple_apply_eq_inf_sup (x : Fin n → B) (k : Fin n) :
    sortTuple x k =
      ((univ : Finset (Fin n)).powersetCard (k.val + 1)).attach.inf'
        ((powersetCard_succ_nonempty k).attach)
        (fun S => S.1.sup' (powersetCard_mem_nonempty S.2) x)
\end{leansrc}
\par\medskip\noindent{\small \texttt{SuFunctor.Adequacy.Bag.}\texttt{\bfseries measurable\_symSec}}\hfill{\footnotesize \path{lean/su-functor/SuFunctor/Adequacy/Bag.lean}, line 377}\par\nopagebreak
\begin{leansrc}
theorem measurable_symSec
    (hle : MeasurableSet {p : B × B | p.1 ≤ p.2}) :
    Measurable (symSec : Sym B n → (Fin n → B))
\end{leansrc}
\par\medskip\noindent{\small \texttt{SuFunctor.Adequacy.Bag.}\texttt{\bfseries monotone\_eq\_of\_toSym\_eq}}\hfill{\footnotesize \path{lean/su-functor/SuFunctor/Adequacy/Bag.lean}, line 335}\par\nopagebreak
\begin{leansrc}
theorem monotone_eq_of_toSym_eq {x y : Fin n → B}
    (hx : Monotone x) (hy : Monotone y) (h : toSym x = toSym y) :
    x = y
\end{leansrc}
\par\medskip\noindent{\small \texttt{SuFunctor.Adequacy.Bag.}\texttt{\bfseries quotKernel\_comp\_eq}}\hfill{\footnotesize \path{lean/su-functor/SuFunctor/Adequacy/Bag.lean}, line 411}\par\nopagebreak
\begin{leansrc}
theorem quotKernel_comp_eq (L : Kernel α T) (hs : Measurable s)
    (hL : ∀ x, L (s (q x)) = L x) (x : α) :
    quotKernel L hs (q x) = L x
\end{leansrc}
\par\medskip\noindent{\small \texttt{SuFunctor.Adequacy.Bag.}\texttt{\bfseries iter\_intertwine}}\hfill{\footnotesize \path{lean/su-functor/SuFunctor/Adequacy/Bag.lean}, line 451}\par\nopagebreak
\begin{leansrc}
theorem iter_intertwine (hq : Measurable q) (hs : Measurable s)
    (K : Kernel α α)
    (hK : ∀ x, (K.map q) (s (q x)) = (K.map q) x) (k : ℕ) :
    Kernel.deterministic q hq ∘ₖ iterK K k =
      iterK (quotKernel (K.map q) hs) k ∘ₖ Kernel.deterministic q hq
\end{leansrc}
\par\medskip\noindent{\small \texttt{SuFunctor.Adequacy.Bag.}\texttt{\bfseries comap\_section\_congr}}\hfill{\footnotesize \path{lean/su-functor/SuFunctor/Adequacy/Bag.lean}, line 419}\par\nopagebreak
\begin{leansrc}
theorem comap_section_congr (L : Kernel α T)
    {s' : ᾱ → α} (hs : Measurable s) (hs' : Measurable s')
    (hsec' : ∀ m, q (s' m) = m)
    (hL : ∀ x, L (s (q x)) = L x) :
    L.comap s' hs' = L.comap s hs
\end{leansrc}
\par\medskip\noindent{\small \texttt{SuFunctor.Adequacy.Bag.}\texttt{\bfseries bag\_transport}}\hfill{\footnotesize \path{lean/su-functor/SuFunctor/Adequacy/Bag.lean}, line 598}\par\nopagebreak
\begin{leansrc}
theorem bag_transport [LinearOrder B] [MeasurableSpace B]
    (hle : MeasurableSet {p : B × B | p.1 ≤ p.2})
    (K : Kernel (BagPoints B) (BagPoints B))
    (hK : ∀ (n : ℕ) (x : Fin n → B) (σ : Equiv.Perm (Fin n)),
      (K.map toSymSigma) ⟨n, x ∘ σ⟩ = (K.map toSymSigma) ⟨n, x⟩)
    (k : ℕ) (p : BagPoints B) :
    (iterK K k p).map toSymSigma =
      iterK (quotKernel (K.map toSymSigma) (measurable_symSecSigma hle)) k
        (toSymSigma p)
\end{leansrc}
\par\medskip\noindent{\small \texttt{SuFunctor.Adequacy.Bag.}\texttt{\bfseries invariance\_iff\_perm}}\hfill{\footnotesize \path{lean/su-functor/SuFunctor/Adequacy/Bag.lean}, line 572}\par\nopagebreak
\begin{leansrc}
theorem invariance_iff_perm [LinearOrder B] [MeasurableSpace B]
    {T : Type*} [MeasurableSpace T]
    (L : Kernel (BagPoints B) T) :
    (∀ p : BagPoints B, L (symSecSigma (toSymSigma p)) = L p) ↔
      (∀ (n : ℕ) (x : Fin n → B) (σ : Equiv.Perm (Fin n)),
        L ⟨n, x ∘ σ⟩ = L ⟨n, x⟩)
\end{leansrc}
\par\medskip\noindent{\small \texttt{SuFunctor.Adequacy.Bag.}\texttt{\bfseries measurableSet\_le\_lex}}\hfill{\footnotesize \path{lean/su-functor/SuFunctor/Adequacy/Bag.lean}, line 627}\par\nopagebreak
\begin{leansrc}
theorem measurableSet_le_lex (d : ℕ) :
    MeasurableSet {q : (Fin d → ℝ) × (Fin d → ℝ) |
      toLex q.1 ≤ toLex q.2}
\end{leansrc}
\par\medskip\noindent{\small \texttt{SuFunctor.Adequacy.FinCISPN.}\texttt{\bfseries not\_branchZenoFree\_of\_disabled}}\hfill{\footnotesize \path{lean/su-functor/SuFunctor/Adequacy/StochasticPetriColouredDemonic.lean}, line 121}\par\nopagebreak
\begin{leansrc}
theorem not_branchZenoFree_of_disabled {i : N.Vanishing} {a : N.IAct}
    (h : ¬ N.IEnabled a i.1) (r : N.Vanishing → ℕ) :
    ¬ N.toFinMA.BranchZenoFree r
\end{leansrc}
\par\medskip\noindent{\small \texttt{SuFunctor.Adequacy.FinCISPN.}\texttt{\bfseries RespScheduler}}\hfill{\footnotesize \path{lean/su-functor/SuFunctor/Adequacy/StochasticPetriColouredDemonic.lean}, line 139}\par\nopagebreak
\begin{leansrc}
def RespScheduler : Type
\end{leansrc}
\par\medskip\noindent{\small \texttt{SuFunctor.Adequacy.SpnDemonicControls.}\texttt{\bfseries arena\_plural}}\hfill{\footnotesize \path{lean/su-functor/SuFunctor/Adequacy/StochasticPetriDemonic.lean}, line 532}\par\nopagebreak
\begin{leansrc}
theorem arena_plural :
    (net₁.uniformResp : net₁.RespScheduler).1.choice i10 1
      ≠ discardResp.1.choice i10 1
\end{leansrc}
\par\medskip\noindent{\small \texttt{SuFunctor.Adequacy.FinCISPN.}\texttt{\bfseries RespectingRows}}\hfill{\footnotesize \path{lean/su-functor/SuFunctor/Adequacy/StochasticPetriColouredDemonic.lean}, line 213}\par\nopagebreak
\begin{leansrc}
structure RespectingRows where …
\end{leansrc}
\par\medskip\noindent{\small \texttt{SuFunctor.Adequacy.FinISPN.}\texttt{\bfseries ispn\_demonic\_adequacy}}\hfill{\footnotesize \path{lean/su-functor/SuFunctor/Adequacy/StochasticPetriDemonic.lean}, line 295}\par\nopagebreak
\begin{leansrc}
theorem ispn_demonic_adequacy {r : N.Vanishing → ℕ}
    (hr : N.CascadeRank r) (k : ℕ) {a : Crd (Traj N.Tangible k)} :
    ∃ α : Cell (N.ispnDemonicC (absFuel r) k : Loose a a)
        (N.demonicTangibleLaws hr k) (𝟙 a) (𝟙 a),
      Cell.IsInvertible α
\end{leansrc}
\par\medskip\noindent{\small \texttt{SuFunctor.Adequacy.FinISPN.}\texttt{\bfseries ispnDemonicC\_iff}}\hfill{\footnotesize \path{lean/su-functor/SuFunctor/Adequacy/StochasticPetriDemonic.lean}, line 250}\par\nopagebreak
\begin{leansrc}
theorem ispnDemonicC_iff {r : N.Vanishing → ℕ} (hr : N.CascadeRank r)
    (k : ℕ) (μ : SubPMF (Traj N.Tangible k)) :
    N.ispnDemonicC (absFuel r) k μ ↔ N.demonicTangibleLaws hr k μ
\end{leansrc}
\par\medskip\noindent{\small \texttt{SuFunctor.Adequacy.FinCISPN.}\texttt{\bfseries demonicTangibleLaws\_rank\_irrel}}\hfill{\footnotesize \path{lean/su-functor/SuFunctor/Adequacy/StochasticPetriColouredDemonic.lean}, line 177}\par\nopagebreak
\begin{leansrc}
theorem demonicTangibleLaws_rank_irrel {r r' : N.Vanishing → ℕ}
    (hr : N.CascadeRank r) (hr' : N.CascadeRank r') (k : ℕ) :
    N.demonicTangibleLaws hr k = N.demonicTangibleLaws hr' k
\end{leansrc}
\par\medskip\noindent{\small \texttt{SuFunctor.Adequacy.FinCISPN.}\texttt{\bfseries demonicTangibleCredal\_rank\_irrel}}\hfill{\footnotesize \path{lean/su-functor/SuFunctor/Adequacy/StochasticPetriColouredDemonic.lean}, line 191}\par\nopagebreak
\begin{leansrc}
theorem demonicTangibleCredal_rank_irrel {r r' : N.Vanishing → ℕ}
    (hr : N.CascadeRank r) (hr' : N.CascadeRank r') (k : ℕ) :
    N.demonicTangibleCredal hr k = N.demonicTangibleCredal hr' k
\end{leansrc}
\par\medskip\noindent{\small \texttt{SuFunctor.Adequacy.FinISPN.}\texttt{\bfseries ispn\_demonic\_maxSlice}}\hfill{\footnotesize \path{lean/su-functor/SuFunctor/Adequacy/StochasticPetriDemonic.lean}, line 321}\par\nopagebreak
\begin{leansrc}
theorem ispn_demonic_maxSlice {r : N.Vanishing → ℕ}
    (hr : N.CascadeRank r) (k : ℕ) (ν : SubPMF (Traj N.Tangible k)) :
    maxSlice (N.demonicTangibleCredal hr k) ν ↔
      N.demonicTangibleLaws hr k ν
\end{leansrc}
\par\medskip\noindent{\small \texttt{SuFunctor.Adequacy.FinISPN.}\texttt{\bfseries satCredal\_lower\_ispnDemonicCredal}}\hfill{\footnotesize \path{lean/su-functor/SuFunctor/Adequacy/StochasticPetriDemonic.lean}, line 346}\par\nopagebreak
\begin{leansrc}
theorem satCredal_lower_ispnDemonicCredal {r : N.Vanishing → ℕ}
    (hr : N.CascadeRank r) (k : ℕ)
    (c : SubPMF (Traj N.Tangible k) → Prop) :
    SatCredal .lower (N.demonicTangibleCredal hr k) c ↔
      ∀ σ : N.RespScheduler,
        c (N.tangiblePathLaw σ.1 r (N.respZenoFree hr σ) k)
\end{leansrc}
\par\medskip\noindent{\small \texttt{SuFunctor.Adequacy.FinISPN.}\texttt{\bfseries su\_ispn\_demonic\_adequacy}}\hfill{\footnotesize \path{lean/su-functor/SuFunctor/Adequacy/StochasticPetriDemonic.lean}, line 412}\par\nopagebreak
\begin{leansrc}
theorem su_ispn_demonic_adequacy {r : N.Vanishing → ℕ}
    (hr : N.CascadeRank r) (k : ℕ) (μ : SubPMF (Traj N.Tangible k)) :
    idealPoint μ ∈ (N.ispnDemonicDnf (absFuel r) k).modelClass
      ↔ N.demonicTangibleLaws hr k μ
\end{leansrc}
\par\medskip\noindent{\small \texttt{SuFunctor.Adequacy.FinISPN.}\texttt{\bfseries chainFire\_val}}\hfill{\footnotesize \path{lean/su-functor/SuFunctor/Adequacy/StochasticPetriSetFiring.lean}, line 196}\par\nopagebreak
\begin{leansrc}
theorem chainFire_val : ∀ (l : List U) (m : BMarking P B),
    N.ChainEnabled l m → ∀ p : P,
    (N.chainFire l m p : ℕ) + N.preSum ↑l p
      = (m p : ℕ) + N.postSum ↑l p
  | [], m, _, p => by simp
  | u :: l, m, h, p => by
    obtain ⟨hu, hl⟩
\end{leansrc}
\par\medskip\noindent{\small \texttt{SuFunctor.Adequacy.FinISPN.}\texttt{\bfseries chainFire\_eq\_of\_enabled}}\hfill{\footnotesize \path{lean/su-functor/SuFunctor/Adequacy/StochasticPetriSetFiring.lean}, line 213}\par\nopagebreak
\begin{leansrc}
theorem chainFire_eq_of_enabled {l₁ l₂ : List U} {m : BMarking P B}
    (h₁ : N.ChainEnabled l₁ m) (h₂ : N.ChainEnabled l₂ m)
    (hml : (↑l₁ : Multiset U) = ↑l₂) :
    N.chainFire l₁ m = N.chainFire l₂ m
\end{leansrc}
\par\medskip\noindent{\small \texttt{SuFunctor.Adequacy.FinISPN.}\texttt{\bfseries chainFire\_eq\_compositeFire}}\hfill{\footnotesize \path{lean/su-functor/SuFunctor/Adequacy/StochasticPetriSetFiring.lean}, line 262}\par\nopagebreak
\begin{leansrc}
theorem chainFire_eq_compositeFire {l : List U} {m : BMarking P B}
    (h : N.ChainEnabled l m) (hsup : N.Supplied ↑l m) :
    N.chainFire l m = N.compositeFire ↑l m
\end{leansrc}
\par\medskip\noindent{\small \texttt{SuFunctor.Adequacy.FinISPN.}\texttt{\bfseries JointReservation}}\hfill{\footnotesize \path{lean/su-functor/SuFunctor/Adequacy/StochasticPetriSetFiring.lean}, line 296}\par\nopagebreak
\begin{leansrc}
def JointReservation (S : Multiset U) (m : BMarking P B) : Prop
\end{leansrc}
\par\medskip\noindent{\small \texttt{SuFunctor.Adequacy.FinISPN.}\texttt{\bfseries setNonConflicting\_of\_headroom}}\hfill{\footnotesize \path{lean/su-functor/SuFunctor/Adequacy/StochasticPetriSetFiring.lean}, line 377}\par\nopagebreak
\begin{leansrc}
theorem setNonConflicting_of_headroom {S : Multiset U} {m : BMarking P B}
    (hJ : N.JointReservation S m)
    (hcap : ∀ p, (m p : ℕ) + N.postSum S p ≤ B) :
    N.SetNonConflicting S m
\end{leansrc}
\par\medskip\noindent{\small \texttt{SuFunctor.Adequacy.FinISPN.}\texttt{\bfseries r1r3\_composite\_adequacy}}\hfill{\footnotesize \path{lean/su-functor/SuFunctor/Adequacy/StochasticPetriSetFiring.lean}, line 391}\par\nopagebreak
\begin{leansrc}
theorem r1r3_composite_adequacy {S : Multiset U} {m : BMarking P B}
    (hJ : N.JointReservation S m)
    (hcap : ∀ p, (m p : ℕ) + N.postSum S p ≤ B) :
    ∀ l : List U, (↑l : Multiset U) = S →
      N.ChainEnabled l m ∧ N.chainFire l m = N.compositeFire S m
\end{leansrc}
\par\medskip\noindent{\small \texttt{SuFunctor.Adequacy.FinCISPN.}\texttt{\bfseries flatten\_triangle\_of\_agree}}\hfill{\footnotesize \path{lean/su-functor/SuFunctor/Adequacy/StochasticPetriColouredComposed.lean}, line 343}\par\nopagebreak
\begin{leansrc}
theorem flatten_triangle_of_agree (σ : Policy N.Vanishing N.IAct)
    (r : N.Vanishing → ℕ) (hr : ZenoFree (N.toFinMA.imStep σ) r)
    {Mode : Type} [Fintype Mode] (e : Mode ≃ N.Tangible)
    (flat : FinCTMC Mode)
    (hoff : ∀ s s' : N.Tangible, s ≠ s' →
      (flat.relabel e).rate s s'
        = (N.toFinMA.close σ r hr).rate s s')
    (hinit : (flat.relabel e).init = (N.toFinMA.close σ r hr).init)
    (k : ℕ) {a : Crd (Traj N.Tangible k)} :
    ∃ α : Cell
        (looseConj (atoms (jumpRows
          (closeMaRows N.compileCispn (compilePolicy σ) (absFuel r))) k) :
          Loose a a)
        (looseConj (atoms (jumpRows (compileCtmc (flat.relabel e))) k))
        (𝟙 a) (𝟙 a),
      Cell.IsInvertible α
\end{leansrc}
\par\medskip\noindent{\small \texttt{SuFunctor.}\texttt{\bfseries SatAngelic}}\hfill{\footnotesize \path{lean/su-functor/SuFunctor/Adequacy/StochasticPetriAngelic.lean}, line 121}\par\nopagebreak
\begin{leansrc}
def SatAngelic (C c : SubPMF V → Prop) : Prop
\end{leansrc}
\par\medskip\noindent{\small \texttt{SuFunctor.Adequacy.FinISPN.}\texttt{\bfseries satAngelic\_ispnDemonicCredal}}\hfill{\footnotesize \path{lean/su-functor/SuFunctor/Adequacy/StochasticPetriAngelic.lean}, line 196}\par\nopagebreak
\begin{leansrc}
theorem satAngelic_ispnDemonicCredal {r : N.Vanishing → ℕ}
    (hr : N.CascadeRank r) (k : ℕ)
    (c : SubPMF (Traj N.Tangible k) → Prop) :
    SatAngelic (N.demonicTangibleCredal hr k) c ↔
      ∃ σ : N.RespScheduler,
        c (N.tangiblePathLaw σ.1 r (N.respZenoFree hr σ) k)
\end{leansrc}
\par\medskip\noindent{\small \texttt{SuFunctor.Adequacy.FinISPN.}\texttt{\bfseries ispn\_demonic\_verdict\_to\_member}}\hfill{\footnotesize \path{lean/su-functor/SuFunctor/Adequacy/StochasticPetriAngelic.lean}, line 220}\par\nopagebreak
\begin{leansrc}
theorem ispn_demonic_verdict_to_member {r : N.Vanishing → ℕ}
    (hr : N.CascadeRank r) (σ : N.RespScheduler) (k : ℕ)
    {c : SubPMF (Traj N.Tangible k) → Prop}
    (h : SatCredal .lower (N.demonicTangibleCredal hr k) c) :
    c (N.tangiblePathLaw σ.1 r (N.respZenoFree hr σ) k)
\end{leansrc}
\par\medskip\noindent{\small \texttt{SuFunctor.Adequacy.FinISPN.}\texttt{\bfseries flatten\_angelic\_of\_agree}}\hfill{\footnotesize \path{lean/su-functor/SuFunctor/Adequacy/StochasticPetriAngelic.lean}, line 264}\par\nopagebreak
\begin{leansrc}
theorem flatten_angelic_of_agree {r : N.Vanishing → ℕ}
    (hrk : N.CascadeRank r) (σ : N.RespScheduler)
    {Mode : Type} [Fintype Mode] (e : Mode ≃ N.Tangible)
    (flat : FinCTMC Mode)
    (hoff : ∀ s s' : N.Tangible, s ≠ s' →
      (flat.relabel e).rate s s'
        = (N.toFinMA.close σ.1 r (N.respZenoFree hrk σ)).rate s s')
    (hinit : (flat.relabel e).init
        = (N.toFinMA.close σ.1 r (N.respZenoFree hrk σ)).init)
    (k : ℕ) {c : SubPMF (Traj N.Tangible k) → Prop}
    (hc : c ((flat.relabel e).jumpLaw k)) :
    SatAngelic (N.demonicTangibleCredal hrk k) c
\end{leansrc}
\par\medskip\noindent{\small \texttt{SuFunctor.Adequacy.PetriVanilla.}\texttt{\bfseries face\_eq\_reach}}\hfill{\footnotesize \path{lean/su-functor/SuFunctor/Adequacy/PetriVanilla.lean}, line 409}\par\nopagebreak
\begin{leansrc}
theorem face_eq_reach (hd : ∀ t, 0 < d t) :
    face N d m0 = {m | Reach N m0 m}
\end{leansrc}
\par\medskip\noindent{\small \texttt{SuFunctor.Adequacy.PetriVanilla.}\texttt{\bfseries face\_dwell\_uniform}}\hfill{\footnotesize \path{lean/su-functor/SuFunctor/Adequacy/PetriVanilla.lean}, line 418}\par\nopagebreak
\begin{leansrc}
theorem face_dwell_uniform {d' : T → ℚ} (hd : ∀ t, 0 < d t)
    (hd' : ∀ t, 0 < d' t) : face N d m0 = face N d' m0
\end{leansrc}
\par\medskip\noindent{\small \texttt{SuFunctor.Adequacy.PetriVanilla.}\texttt{\bfseries su\_vanilla\_demonic\_adequacy}}\hfill{\footnotesize \path{lean/su-functor/SuFunctor/Adequacy/PetriVanilla.lean}, line 570}\par\nopagebreak
\begin{leansrc}
theorem su_vanilla_demonic_adequacy (k : ℕ)
    (μ : SubPMF (Traj (Marking P) k)) :
    idealPoint μ ∈ (vanillaDemonicDnf N d m0 k).modelClass
      ↔ vanillaDemonicLaws N d m0 k μ
\end{leansrc}
\par\medskip\noindent{\small \texttt{SuFunctor.Adequacy.PetriVanilla.}\texttt{\bfseries vanillaDemonicC\_iff}}\hfill{\footnotesize \path{lean/su-functor/SuFunctor/Adequacy/PetriVanilla.lean}, line 515}\par\nopagebreak
\begin{leansrc}
theorem vanillaDemonicC_iff (k : ℕ) (μ : SubPMF (Traj (Marking P) k)) :
    vanillaDemonicC N d m0 k μ ↔ vanillaDemonicLaws N d m0 k μ
\end{leansrc}
\par\medskip\noindent{\small \texttt{SuFunctor.Adequacy.PetriVanilla.}\texttt{\bfseries recordVisited\_eq\_face}}\hfill{\footnotesize \path{lean/su-functor/SuFunctor/Adequacy/PetriVanilla.lean}, line 622}\par\nopagebreak
\begin{leansrc}
theorem recordVisited_eq_face :
    recordVisited N d m0 = face N d m0
\end{leansrc}
\par\medskip\noindent{\small \texttt{SuFunctor.Adequacy.PetriVanilla.}\texttt{\bfseries su\_vanilla\_reach\_iff}}\hfill{\footnotesize \path{lean/su-functor/SuFunctor/Adequacy/PetriVanilla.lean}, line 650}\par\nopagebreak
\begin{leansrc}
theorem su_vanilla_reach_iff (hd : ∀ t, 0 < d t) (m : Marking P) :
    (∃ (k : ℕ) (μ : SubPMF (Traj (Marking P) k)),
        idealPoint μ ∈ (vanillaDemonicDnf N d m0 k).modelClass
          ∧ ∃ ω, μ.wt ω ≠ 0 ∧ ∃ i, ω i = m)
      ↔ Reach N m0 m
\end{leansrc}
\par\medskip\noindent{\small \texttt{SuFunctor.Adequacy.PetriVanilla.}\texttt{\bfseries su\_vanilla\_dwell\_uniform}}\hfill{\footnotesize \path{lean/su-functor/SuFunctor/Adequacy/PetriVanilla.lean}, line 671}\par\nopagebreak
\begin{leansrc}
theorem su_vanilla_dwell_uniform {d' : T → ℚ} (hd : ∀ t, 0 < d t)
    (hd' : ∀ t, 0 < d' t) :
    recordVisited N d m0 = recordVisited N d' m0
\end{leansrc}
\par\medskip\noindent{\small \texttt{SuFunctor.Adequacy.PetriVanilla.}\texttt{\bfseries su\_vanilla\_safety}}\hfill{\footnotesize \path{lean/su-functor/SuFunctor/Adequacy/PetriVanilla.lean}, line 681}\par\nopagebreak
\begin{leansrc}
theorem su_vanilla_safety (hd : ∀ t, 0 < d t) (Φ : Marking P → Prop) :
    (∀ m ∈ recordVisited N d m0, Φ m) ↔
      ∀ m, Reach N m0 m → Φ m
\end{leansrc}
\par\medskip\noindent{\small \texttt{SuFunctor.Adequacy.Pgm.}\texttt{\bfseries pgm\_adequacy}}\hfill{\footnotesize \path{lean/su-functor/SuFunctor/Adequacy/Pgm.lean}, line 727}\par\nopagebreak
\begin{leansrc}
theorem pgm_adequacy (B : FinBN n Val) {a : Crd (Asn Val)} :
    ∃ α : Cell (looseConj (atoms (compile B)) : Loose a a)
        (singletonC B.jointLaw) (𝟙 a) (𝟙 a),
      Cell.IsInvertible α
\end{leansrc}
\par\medskip\noindent{\small \texttt{SuFunctor.Adequacy.Pgm.FinBN.}\texttt{\bfseries jointLaw\_mass}}\hfill{\footnotesize \path{lean/su-functor/SuFunctor/Adequacy/Pgm.lean}, line 374}\par\nopagebreak
\begin{leansrc}
theorem FinBN.jointLaw_mass (B : FinBN n Val) : (B.jointLaw).mass = 1
\end{leansrc}
\par\medskip\noindent{\small \texttt{SuFunctor.Adequacy.Pgm.FinBN.}\texttt{\bfseries ancestralPMF\_eq}}\hfill{\footnotesize \path{lean/su-functor/SuFunctor/Adequacy/Pgm.lean}, line 921}\par\nopagebreak
\begin{leansrc}
theorem FinBN.ancestralPMF_eq (σ : Asn Val) :
    (B.ancestralPMF) σ = q2e ((B.jointLaw).wt σ)
\end{leansrc}
\par\medskip\noindent{\small \texttt{SuFunctor.Adequacy.}\texttt{\bfseries noiseDenot\_progTraj}}\hfill{\footnotesize \path{lean/su-functor/SuFunctor/Adequacy/ProbProg.lean}, line 200}\par\nopagebreak
\begin{leansrc}
theorem noiseDenot_progTraj (P : ProbProgram S N₀ N) :
    ∀ k, noiseDenot (progTraj P k) (measurable_progTraj P k) = pathDist P k
  | 0 => by
    calc noiseDenot (progTraj P 0) (measurable_progTraj P 0)
        = noiseDenot (trajZero ∘ P.start)
            (measurable_trajZero.comp P.measurable_start)
\end{leansrc}
\par\medskip\noindent{\small \texttt{SuFunctor.Adequacy.}\texttt{\bfseries su\_pprog\_adequacy}}\hfill{\footnotesize \path{lean/su-functor/SuFunctor/Adequacy/ProbProg.lean}, line 398}\par\nopagebreak
\begin{leansrc}
theorem su_pprog_adequacy {P : ProbProgram S N₀ N} {M : FinDTMC S}
    (hPM : Presents P M) (k : ℕ) (μ : SubPMF (Traj S k)) :
    idealPoint μ ∈ (pprogDnf (compile M) k).modelClass
      ↔ μ.probFun = ((noiseDenot (progTraj P k) (measurable_progTraj P k) :
          ProbState (Traj S k)) : Traj S k → ℝ)
\end{leansrc}
\par\medskip\noindent{\small \texttt{SuFunctor.Adequacy.QNet.}\texttt{\bfseries quantum\_net\_adequacy}}\hfill{\footnotesize \path{lean/su-functor/SuFunctor/Adequacy/QuantumNet.lean}, line 426}\par\nopagebreak
\begin{leansrc}
theorem quantum_net_adequacy (P : QNetSystem X n ι T κ) {a : Crd κ} :
    ∃ α : Cell (looseConj (QNet.atoms (QNet.compile P)) : Loose a a)
        (singletonC P.shadowLaw) (𝟙 a) (𝟙 a),
      Cell.IsInvertible α
\end{leansrc}
\par\medskip\noindent{\small \texttt{SuFunctor.Adequacy.}\texttt{\bfseries QNetSystem}}\hfill{\footnotesize \path{lean/su-functor/SuFunctor/Adequacy/QuantumNet.lean}, line 156}\par\nopagebreak
\begin{leansrc}
structure QNetSystem (X : Type*) [PartialOrder X]
    (n : Type*) [Fintype n] [DecidableEq n]
    (ι T κ : Type*) [Fintype κ] where …
\end{leansrc}
\par\medskip\noindent{\small \texttt{SuFunctor.Adequacy.QNetSystem.}\texttt{\bfseries shadowPMF\_prob\_eq\_born}}\hfill{\footnotesize \path{lean/su-functor/SuFunctor/Adequacy/QuantumNet.lean}, line 204}\par\nopagebreak
\begin{leansrc}
theorem shadowPMF_prob_eq_born (α : κ) :
    ((P.shadowPMF.prob α : ℚ) : ℝ) =
      ((P.state * P.frag.records α).trace).re
\end{leansrc}
\par\medskip\noindent{\small \texttt{SuFunctor.Adequacy.QNetSystem.}\texttt{\bfseries shadowPMF\_eq\_fragment\_shadow}}\hfill{\footnotesize \path{lean/su-functor/SuFunctor/Adequacy/QuantumNet.lean}, line 213}\par\nopagebreak
\begin{leansrc}
theorem shadowPMF_eq_fragment_shadow (α : κ) :
    ((P.shadowPMF.prob α : ℚ) : ℝ) =
      P.frag.shadow
        (stateClassMk P.frag.net P.frag.ambient
          ⟨P.state, P.state_isState⟩) α
\end{leansrc}
\par\medskip\noindent{\small \texttt{SuFunctor.Adequacy.}\texttt{\bfseries shadowPMF\_class\_invariant}}\hfill{\footnotesize \path{lean/su-functor/SuFunctor/Adequacy/QuantumNet.lean}, line 227}\par\nopagebreak
\begin{leansrc}
theorem shadowPMF_class_invariant (F : LCQFTFragment X n ι T κ)
    {ρ σ : Matrix n n ℂ} (hρ : IsState ρ) (hσ : IsState σ)
    {q₁ q₂ : κ → ℚ}
    (h₁ : ∀ α, ((q₁ α : ℚ) : ℝ) = recordShadow F.records ρ α)
    (h₂ : ∀ α, ((q₂ α : ℚ) : ℝ) = recordShadow F.records σ α)
    (hcls : stateClassMk F.net F.ambient ⟨ρ, hρ⟩ =
      stateClassMk F.net F.ambient ⟨σ, hσ⟩) :
    (QNetSystem.mk F ρ hρ q₁ h₁).shadowPMF =
      (QNetSystem.mk F σ hσ q₂ h₂).shadowPMF
\end{leansrc}
\par\medskip\noindent{\small \texttt{SuFunctor.Adequacy.QNetSystem.}\texttt{\bfseries shadowLaw\_mass}}\hfill{\footnotesize \path{lean/su-functor/SuFunctor/Adequacy/QuantumNet.lean}, line 198}\par\nopagebreak
\begin{leansrc}
theorem shadowLaw_mass : P.shadowLaw.mass = 1
\end{leansrc}
\par\medskip\noindent{\small \texttt{SuFunctor.Adequacy.QNet.}\texttt{\bfseries rowsLaw\_mass\_of\_complete}}\hfill{\footnotesize \path{lean/su-functor/SuFunctor/Adequacy/QuantumNet.lean}, line 313}\par\nopagebreak
\begin{leansrc}
theorem rowsLaw_mass_of_complete (R : QNetRows κ) (hR : R.WF)
    (hC : R.Complete) : (rowsLaw R hR).mass = 1
\end{leansrc}
\par\medskip\noindent{\small \texttt{SuFunctor.Adequacy.QNet.}\texttt{\bfseries bellQNet\_shadowPMF}}\hfill{\footnotesize \path{lean/su-functor/SuFunctor/Adequacy/QuantumNet.lean}, line 714}\par\nopagebreak
\begin{leansrc}
theorem bellQNet_shadowPMF : bellQNet.shadowPMF = bellShadowPMF
\end{leansrc}
\par\medskip\noindent{\small \texttt{SuFunctor.Adequacy.QNet.}\texttt{\bfseries mixedQNet\_shadowPMF}}\hfill{\footnotesize \path{lean/su-functor/SuFunctor/Adequacy/QuantumNet.lean}, line 718}\par\nopagebreak
\begin{leansrc}
theorem mixedQNet_shadowPMF : mixedQNet.shadowPMF = mixedShadowPMF
\end{leansrc}
\par\medskip\noindent{\small \texttt{SuFunctor.Adequacy.QNet.}\texttt{\bfseries exists\_incomplete\_mass\_lt\_one}}\hfill{\footnotesize \path{lean/su-functor/SuFunctor/Adequacy/QuantumNet.lean}, line 726}\par\nopagebreak
\begin{leansrc}
theorem exists_incomplete_mass_lt_one :
    ∃ (P : QNetSystem (Mink ℝ) (Fin 2 × Fin 2) (Fin 2) ℝ
        (Fin 2 × Fin 2)) (J : Finset (Fin 2 × Fin 2)),
      (QNet.rowsLaw (QNet.compileSub P J)
        (QNet.compileSub_wf P J)).mass < 1
\end{leansrc}
\par\medskip\noindent{\small \texttt{SuFunctor.Adequacy.QNet.}\texttt{\bfseries exists\_nonPSD\_family\_negative\_shadow}}\hfill{\footnotesize \path{lean/su-functor/SuFunctor/Adequacy/QuantumNet.lean}, line 822}\par\nopagebreak
\begin{leansrc}
theorem exists_nonPSD_family_negative_shadow :
    ∃ (R : Fin 2 → Matrix (Fin 1) (Fin 1) ℂ)
      (ρ : Matrix (Fin 1) (Fin 1) ℂ),
      IsState ρ ∧ (∑ α, R α = 1) ∧ ∃ α, recordShadow R ρ α < 0
\end{leansrc}
\par\medskip\noindent{\small \texttt{SuFunctor.Adequacy.QNet.}\texttt{\bfseries su\_qnet\_adequacy}}\hfill{\footnotesize \path{lean/su-functor/SuFunctor/Adequacy/QuantumNet.lean}, line 895}\par\nopagebreak
\begin{leansrc}
theorem su_qnet_adequacy (P : QNetSystem X n ι T κ) (μ : SubPMF κ) :
    idealPoint (μ.map recState)
        ∈ (qnetSpecW (QNet.compile P)).modelClass
      ↔ μ = P.shadowLaw
\end{leansrc}
\par\medskip\noindent{\small \texttt{SuFunctor.Adequacy.QNet.}\texttt{\bfseries quantum\_credal\_adequacy}}\hfill{\footnotesize \path{lean/su-functor/SuFunctor/Adequacy/QuantumNet.lean}, line 546}\par\nopagebreak
\begin{leansrc}
theorem quantum_credal_adequacy (F : LCQFTFragment X n ι T κ)
    (S : Set {ρ : Matrix n n ℂ // IsState ρ}) {a : Crd κ} :
    ∃ α : Cell (QNet.qnetCredalC F S : Loose a a)
        (QNet.quantumLaws F S) (𝟙 a) (𝟙 a),
      Cell.IsInvertible α
\end{leansrc}
\par\medskip\noindent{\small \texttt{SuFunctor.Adequacy.QNet.QNet.}\texttt{\bfseries qnetCredalC}}\hfill{\footnotesize \path{lean/su-functor/SuFunctor/Adequacy/QuantumNet.lean}, line 514}\par\nopagebreak
\begin{leansrc}
def qnetCredalC (F : LCQFTFragment X n ι T κ)
    (S : Set {ρ : Matrix n n ℂ // IsState ρ}) : SubPMF κ → Prop
\end{leansrc}
\par\medskip\noindent{\small \texttt{SuFunctor.Adequacy.QNet.QNet.}\texttt{\bfseries qnetCredalC\_iff}}\hfill{\footnotesize \path{lean/su-functor/SuFunctor/Adequacy/QuantumNet.lean}, line 522}\par\nopagebreak
\begin{leansrc}
theorem qnetCredalC_iff (F : LCQFTFragment X n ι T κ)
    (S : Set {ρ : Matrix n n ℂ // IsState ρ}) (μ : SubPMF κ) :
    qnetCredalC F S μ ↔ quantumLaws F S μ
\end{leansrc}
\par\medskip\noindent{\small \texttt{SuFunctor.Adequacy.QNet.QNet.}\texttt{\bfseries RatShadowMember}}\hfill{\footnotesize \path{lean/su-functor/SuFunctor/Adequacy/QuantumNet.lean}, line 442}\par\nopagebreak
\begin{leansrc}
structure RatShadowMember (F : LCQFTFragment X n ι T κ)
    (S : Set {ρ : Matrix n n ℂ // IsState ρ}) where …
\end{leansrc}
\par\medskip\noindent{\small \texttt{SuFunctor.Adequacy.QNet.}\texttt{\bfseries quantum\_maxSlice}}\hfill{\footnotesize \path{lean/su-functor/SuFunctor/Adequacy/QuantumNet.lean}, line 556}\par\nopagebreak
\begin{leansrc}
theorem quantum_maxSlice (F : LCQFTFragment X n ι T κ)
    (S : Set {ρ : Matrix n n ℂ // IsState ρ}) (ν : SubPMF κ) :
    maxSlice (QNet.quantumCredal F S) ν ↔ QNet.quantumLaws F S ν
\end{leansrc}
\par\medskip\noindent{\small \texttt{SuFunctor.Adequacy.QNet.}\texttt{\bfseries quantum\_ray\_singleton\_reconciliation}}\hfill{\footnotesize \path{lean/su-functor/SuFunctor/Adequacy/QuantumNet.lean}, line 562}\par\nopagebreak
\begin{leansrc}
theorem quantum_ray_singleton_reconciliation
    (F : LCQFTFragment X n ι T κ)
    (S : Set {ρ : Matrix n n ℂ // IsState ρ}) {a : Crd κ} :
    ∃ α : Cell (maxSlice (QNet.quantumCredal F S) : Loose a a)
        (QNet.quantumLaws F S) (𝟙 a) (𝟙 a),
      Cell.IsInvertible α
\end{leansrc}
\par\medskip\noindent{\small \texttt{SuFunctor.Adequacy.QNet.}\texttt{\bfseries satCredal\_lower\_quantumCredal}}\hfill{\footnotesize \path{lean/su-functor/SuFunctor/Adequacy/QuantumNet.lean}, line 573}\par\nopagebreak
\begin{leansrc}
theorem satCredal_lower_quantumCredal (F : LCQFTFragment X n ι T κ)
    (S : Set {ρ : Matrix n n ℂ // IsState ρ})
    (c : SubPMF κ → Prop) :
    SatCredal .lower (QNet.quantumCredal F S) c ↔
      ∀ w : QNet.RatShadowMember F S, c w.law
\end{leansrc}
\par\medskip\noindent{\small \texttt{SuFunctor.Adequacy.QNet.}\texttt{\bfseries satCredal\_upper\_quantumCredal}}\hfill{\footnotesize \path{lean/su-functor/SuFunctor/Adequacy/QuantumNet.lean}, line 582}\par\nopagebreak
\begin{leansrc}
theorem satCredal_upper_quantumCredal (F : LCQFTFragment X n ι T κ)
    (S : Set {ρ : Matrix n n ℂ // IsState ρ})
    {c : SubPMF κ → Prop} (hc : DownClosed c) :
    SatCredal .upper (QNet.quantumCredal F S) c ↔
      ∀ w : QNet.RatShadowMember F S, c w.law
\end{leansrc}
\par\medskip\noindent{\small \texttt{SuFunctor.Adequacy.QNet.QNet.}\texttt{\bfseries quantumLaws\_mix}}\hfill{\footnotesize \path{lean/su-functor/SuFunctor/Adequacy/QuantumNet.lean}, line 660}\par\nopagebreak
\begin{leansrc}
theorem quantumLaws_mix (hS : MixClosed S)
    (w₁ w₂ : RatShadowMember F S) {t : ℚ} (h₀ : 0 ≤ t) (h₁ : t ≤ 1) :
    quantumLaws F S (FinPMF.mix w₁.pmf w₂.pmf t h₀ h₁).toSubPMF
\end{leansrc}
\par\medskip\noindent{\small \texttt{Qtsl.QuantumStates.}\texttt{\bfseries recordShadowPMF\_mix}}\hfill{\footnotesize \path{lean/qtsl-aux/QtslAux/QuantumStates.lean}, line 678}\par\nopagebreak
\begin{leansrc}
theorem recordShadowPMF_mix [Fintype ι] (R : ι → Matrix n n ℂ)
    (ρ σ : Matrix n n ℂ) (q₁ q₂ : ι → ℚ)
    (hq₁ : ∀ α, (q₁ α : ℝ) = recordShadow R ρ α)
    (hq₂ : ∀ α, (q₂ α : ℝ) = recordShadow R σ α)
    (hR : ∀ α, (R α).PosSemidef) (hexh : ∑ α, R α = 1)
    (hρ : IsState ρ) (hσ : IsState σ) {t : ℚ} (h₀ : 0 ≤ t) (h₁ : t ≤ 1) :
    recordShadowPMF R (((t : ℚ) : ℂ) • ρ + (((1 - t : ℚ)) : ℂ) • σ)
        (fun α => t * q₁ α + (1 - t) * q₂ α)
        (fun α => recordShadow_mix_rat_witness hq₁ hq₂ α)
        hR hexh (isState_mix hρ hσ h₀ h₁) =
      FinPMF.mix (recordShadowPMF R ρ q₁ hq₁ hR hexh hρ)
        (recordShadowPMF R σ q₂ hq₂ hR hexh hσ) t h₀ h₁
\end{leansrc}
\par\medskip\noindent{\small \texttt{SuFunctor.Adequacy.QNet.}\texttt{\bfseries mixLaw\_not\_mem\_pair}}\hfill{\footnotesize \path{lean/su-functor/SuFunctor/Adequacy/QuantumNet.lean}, line 775}\par\nopagebreak
\begin{leansrc}
theorem mixLaw_not_mem_pair :
    ¬ QNet.quantumLaws isingQubitFragment bellMixedSet
        (FinPMF.mix bellMember.pmf mixedMember.pmf (1 / 2)
          (by norm_num) (by norm_num)).toSubPMF
\end{leansrc}
\par\medskip\noindent{\small \texttt{SuFunctor.Adequacy.QNet.}\texttt{\bfseries not\_satCredal\_upper\_quantumCredal\_geC}}\hfill{\footnotesize \path{lean/su-functor/SuFunctor/Adequacy/QuantumNet.lean}, line 609}\par\nopagebreak
\begin{leansrc}
theorem not_satCredal_upper_quantumCredal_geC
    (F : LCQFTFragment X n ι T κ)
    (S : Set {ρ : Matrix n n ℂ // IsState ρ})
    [hne : Nonempty (QNet.RatShadowMember F S)]
    {E : Set κ} {q : ℚ} (hq : 0 < q) :
    ¬ SatCredal .upper (QNet.quantumCredal F S) (geC E q)
\end{leansrc}
\par\medskip\noindent{\small \texttt{SuFunctor.Adequacy.QNet.}\texttt{\bfseries satCredal\_quantum\_of\_isEmpty}}\hfill{\footnotesize \path{lean/su-functor/SuFunctor/Adequacy/QuantumNet.lean}, line 625}\par\nopagebreak
\begin{leansrc}
theorem satCredal_quantum_of_isEmpty (F : LCQFTFragment X n ι T κ)
    (S : Set {ρ : Matrix n n ℂ // IsState ρ})
    [h : IsEmpty (QNet.RatShadowMember F S)] (kind : EnvelopeKind)
    (c : SubPMF κ → Prop) :
    SatCredal kind (QNet.quantumCredal F S) c
\end{leansrc}
\par\medskip\noindent{\small \texttt{Qtsl.QuantumStates.}\texttt{\bfseries existsUnique\_isState\_traceRepr}}\hfill{\footnotesize \path{lean/qtsl-aux/QtslAux/QuantumStates.lean}, line 518}\par\nopagebreak
\begin{leansrc}
theorem existsUnique_isState_traceRepr
    (hpos : ∀ v : n → ℂ, 0 ≤ φ (Matrix.vecMulVec v (star v)))
    (hone : φ 1 = 1) :
    ∃! ρ : Matrix n n ℂ, IsState ρ ∧ ∀ a, φ a = (ρ * a).trace
\end{leansrc}
\par\medskip\noindent{\small \texttt{Qtsl.QuantumStates.}\texttt{\bfseries exists\_shadow\_lowerE\_lt\_upperE}}\hfill{\footnotesize \path{lean/qtsl-aux/QtslAux/QuantumStates.lean}, line 953}\par\nopagebreak
\begin{leansrc}
theorem exists_shadow_lowerE_lt_upperE :
    ∃ f : Fin 2 × Fin 2 → ℚ,
      lowerE ![bellShadowPMF, mixedShadowPMF] f <
        upperE ![bellShadowPMF, mixedShadowPMF] f
\end{leansrc}
\par\medskip\noindent{\small \texttt{Qtsl.QuantumCoarseGraining.}\texttt{\bfseries pushforward\_ofPresentation}}\hfill{\footnotesize \path{lean/qtsl-aux/QtslAux/QuantumCoarseGraining.lean}, line 474}\par\nopagebreak
\begin{leansrc}
theorem pushforward_ofPresentation (f : ι → κ) (C : ι → Matrix d d ℂ)
    (ψ : d → ℂ)
    (hnorm : ∑ α, ∑ β, decoherenceFunctional C ψ α β = 1) :
    (QuantumMeasure.ofPresentation C ψ hnorm).pushforward f
      = QuantumMeasure.ofPresentation (coarsenFamily f C) ψ
          (by
            simp only [decoherenceFunctional_coarsenFamily]
            have h := sum_sum_fiberEvent f Finset.univ Finset.univ
              fun a b => decoherenceFunctional C ψ a b
            rw [h, fiberEvent_univ]
            exact hnorm)
\end{leansrc}
\par\medskip\noindent{\small \texttt{SuFunctor.Adequacy.}\texttt{\bfseries tvDistQ}}\hfill{\footnotesize \path{lean/su-functor/SuFunctor/Adequacy/QuantumContainment.lean}, line 135}\par\nopagebreak
\begin{leansrc}
def tvDistQ {κ : Type u} [Fintype κ] (μ ν : SubPMF κ) : ℚ
\end{leansrc}
\par\medskip\noindent{\small \texttt{SuFunctor.Adequacy.}\texttt{\bfseries tvDistQ\_eq\_zero\_iff}}\hfill{\footnotesize \path{lean/su-functor/SuFunctor/Adequacy/QuantumContainment.lean}, line 166}\par\nopagebreak
\begin{leansrc}
theorem tvDistQ_eq_zero_iff {κ : Type u} [Fintype κ] {μ ν : SubPMF κ} :
    tvDistQ μ ν = 0 ↔ μ = ν
\end{leansrc}
\par\medskip\noindent{\small \texttt{SuFunctor.Adequacy.}\texttt{\bfseries credalInflate}}\hfill{\footnotesize \path{lean/su-functor/SuFunctor/Adequacy/QuantumContainment.lean}, line 214}\par\nopagebreak
\begin{leansrc}
def credalInflate {κ : Type u} [Fintype κ] (ε : ℚ)
    (L : SubPMF κ → Prop) : SubPMF κ → Prop
\end{leansrc}
\par\medskip\noindent{\small \texttt{SuFunctor.Adequacy.Containment.}\texttt{\bfseries mono}}\hfill{\footnotesize \path{lean/su-functor/SuFunctor/Adequacy/QuantumContainment.lean}, line 282}\par\nopagebreak
\begin{leansrc}
theorem Containment.mono {ι κ : Type u} [Fintype ι] [Fintype κ]
    {LUV : SubPMF ι → Prop} {f : ι → κ} {LIR : SubPMF κ → Prop}
    {ε ε' : ℚ} (h : ε ≤ ε')
    (hc : Containment LUV f LIR ε) : Containment LUV f LIR ε'
\end{leansrc}
\par\medskip\noindent{\small \texttt{SuFunctor.Adequacy.Containment.}\texttt{\bfseries comp}}\hfill{\footnotesize \path{lean/su-functor/SuFunctor/Adequacy/QuantumContainment.lean}, line 292}\par\nopagebreak
\begin{leansrc}
theorem Containment.comp {ι κ σ : Type u}
    [Fintype ι] [Fintype κ] [Fintype σ]
    {L₁ : SubPMF ι → Prop} {f : ι → κ} {L₂ : SubPMF κ → Prop} {ε₁ : ℚ}
    {g : κ → σ} {L₃ : SubPMF σ → Prop} {ε₂ : ℚ}
    (h₁ : Containment L₁ f L₂ ε₁) (h₂ : Containment L₂ g L₃ ε₂) :
    Containment L₁ (g ∘ f) L₃ (ε₁ + ε₂)
\end{leansrc}
\par\medskip\noindent{\small \texttt{SuFunctor.SubPMF.}\texttt{\bfseries map\_wt\_apply}}\hfill{\footnotesize \path{lean/su-functor/SuFunctor/Adequacy/QuantumContainment.lean}, line 102}\par\nopagebreak
\begin{leansrc}
theorem map_wt_apply {X Y : Type u} [Fintype X] [DecidableEq Y]
    (f : X → Y) (μ : SubPMF X) (c : Y) :
    (μ.map f).wt c
      = ∑ a ∈ Finset.univ.filter (fun a => f a = c), μ.wt a
\end{leansrc}
\par\medskip\noindent{\small \texttt{SuFunctor.Adequacy.}\texttt{\bfseries tvDistQ\_map\_le}}\hfill{\footnotesize \path{lean/su-functor/SuFunctor/Adequacy/QuantumContainment.lean}, line 193}\par\nopagebreak
\begin{leansrc}
theorem tvDistQ_map_le {X Y : Type u} [Fintype X] [Fintype Y]
    (f : X → Y) (μ ν : SubPMF X) :
    tvDistQ (μ.map f) (ν.map f) ≤ tvDistQ μ ν
\end{leansrc}
\par\medskip\noindent{\small \texttt{SuFunctor.Adequacy.}\texttt{\bfseries recordShadow\_coarsenFamily}}\hfill{\footnotesize \path{lean/su-functor/SuFunctor/Adequacy/QuantumContainment.lean}, line 327}\par\nopagebreak
\begin{leansrc}
theorem recordShadow_coarsenFamily {ι' κ' : Type*} [Fintype ι']
    [Fintype κ'] [DecidableEq ι'] [DecidableEq κ']
    {d : Type*} [Fintype d]
    (f : ι' → κ') (R : ι' → Matrix d d ℂ) (ρ : Matrix d d ℂ) (c : κ') :
    recordShadow (coarsenFamily f R) ρ c
      = ∑ a ∈ fiber f c, recordShadow R ρ a
\end{leansrc}
\par\medskip\noindent{\small \texttt{SuFunctor.Adequacy.}\texttt{\bfseries uv\_ir\_containment}}\hfill{\footnotesize \path{lean/su-functor/SuFunctor/Adequacy/QuantumContainment.lean}, line 763}\par\nopagebreak
\begin{leansrc}
theorem uv_ir_containment :
    Containment (QNet.quantumLaws isingQubitFragment uvSeg) Prod.fst
      (QNet.quantumLaws irFragment irSet) (49 / 625)
\end{leansrc}
\par\medskip\noindent{\small \texttt{SuFunctor.Adequacy.}\texttt{\bfseries uv\_ir\_containment\_lower}}\hfill{\footnotesize \path{lean/su-functor/SuFunctor/Adequacy/QuantumContainment.lean}, line 788}\par\nopagebreak
\begin{leansrc}
theorem uv_ir_containment_lower (ε : ℚ)
    (h : Containment (QNet.quantumLaws isingQubitFragment uvSeg)
      Prod.fst (QNet.quantumLaws irFragment irSet) ε) :
    49 / 625 ≤ ε
\end{leansrc}
\par\medskip\noindent{\small \texttt{SuFunctor.Adequacy.}\texttt{\bfseries not\_uv\_ir\_containment\_zero}}\hfill{\footnotesize \path{lean/su-functor/SuFunctor/Adequacy/QuantumContainment.lean}, line 817}\par\nopagebreak
\begin{leansrc}
theorem not_uv_ir_containment_zero :
    ¬ Containment (QNet.quantumLaws isingQubitFragment uvSeg) Prod.fst
      (QNet.quantumLaws irFragment irSet) 0
\end{leansrc}
\par\medskip\noindent{\small \texttt{SuFunctor.Adequacy.}\texttt{\bfseries onestep\_defect}}\hfill{\footnotesize \path{lean/su-functor/SuFunctor/Adequacy/QuantumContainment.lean}, line 1056}\par\nopagebreak
\begin{leansrc}
theorem onestep_defect (x : Fin 2) :
    tvDistQ (coarseLaw (x + 1)) (irStep (coarseLaw x)) ≤ 49 / 625
\end{leansrc}
\par\medskip\noindent{\small \texttt{SuFunctor.Adequacy.}\texttt{\bfseries horizon\_bound}}\hfill{\footnotesize \path{lean/su-functor/SuFunctor/Adequacy/QuantumContainment.lean}, line 1073}\par\nopagebreak
\begin{leansrc}
theorem horizon_bound (x : Fin 2) (k : ℕ) :
    tvDistQ (coarseLaw ((· + 1)^[k] x)) (irStep^[k] (coarseLaw x))
      ≤ (k : ℚ) * (49 / 625)
\end{leansrc}
\par\medskip\noindent{\small \texttt{SuFunctor.Adequacy.}\texttt{\bfseries lax\_horizon}}\hfill{\footnotesize \path{lean/su-functor/SuFunctor/Adequacy/QuantumContainment.lean}, line 395}\par\nopagebreak
\begin{leansrc}
theorem lax_horizon {A : Type*} {B : Type*}
    (push : A → B) (stepA : A → A) (stepB : B → B)
    (d : B → B → ℚ) (ε : ℚ) (Reach : A → Prop)
    (hself : ∀ x, d x x = 0)
    (htri : ∀ x y z, d x z ≤ d x y + d y z)
    (hnonexp : ∀ x y, d (stepB x) (stepB y) ≤ d x y)
    (hinv : ∀ a, Reach a → Reach (stepA a))
    (hdef : ∀ a, Reach a → d (push (stepA a)) (stepB (push a)) ≤ ε) :
    ∀ a, Reach a → ∀ k : ℕ,
      d (push (stepA^[k] a)) (stepB^[k] (push a)) ≤ (k : ℚ) * ε
\end{leansrc}
\par\medskip\noindent{\small \texttt{SuFunctor.Adequacy.}\texttt{\bfseries tvDistQ\_stochAct\_le}}\hfill{\footnotesize \path{lean/su-functor/SuFunctor/Adequacy/QuantumContainment.lean}, line 368}\par\nopagebreak
\begin{leansrc}
theorem tvDistQ_stochAct_le {κ : Type u} [Fintype κ] (Q : κ → κ → ℚ)
    (hQ0 : ∀ c c', 0 ≤ Q c c') (hQ1 : ∀ c, ∑ c', Q c c' = 1)
    (μ ν : SubPMF κ) :
    tvDistQ (stochAct Q hQ0 hQ1 μ) (stochAct Q hQ0 hQ1 ν)
      ≤ tvDistQ μ ν
\end{leansrc}
\par\medskip\noindent{\small \texttt{SuFunctor.Adequacy.}\texttt{\bfseries ir\_tot\_containment}}\hfill{\footnotesize \path{lean/su-functor/SuFunctor/Adequacy/QuantumContainment.lean}, line 1135}\par\nopagebreak
\begin{leansrc}
theorem ir_tot_containment :
    Containment (QNet.quantumLaws irFragment irSet) totalMerge
      (QNet.quantumLaws totFragment totSet) 0
\end{leansrc}
\par\medskip\noindent{\small \texttt{SuFunctor.Adequacy.}\texttt{\bfseries uv\_tot\_containment}}\hfill{\footnotesize \path{lean/su-functor/SuFunctor/Adequacy/QuantumContainment.lean}, line 1159}\par\nopagebreak
\begin{leansrc}
theorem uv_tot_containment :
    Containment (QNet.quantumLaws isingQubitFragment uvSeg)
      (totalMerge ∘ Prod.fst)
      (QNet.quantumLaws totFragment totSet) (49 / 625 + 0)
\end{leansrc}
\par\medskip\noindent{\small \texttt{SuFunctor.Adequacy.}\texttt{\bfseries uvSeg\_mixClosed}}\hfill{\footnotesize \path{lean/su-functor/SuFunctor/Adequacy/QuantumContainment.lean}, line 570}\par\nopagebreak
\begin{leansrc}
theorem uvSeg_mixClosed : QNet.MixClosed uvSeg
\end{leansrc}
\par\medskip\noindent{\small \texttt{SuFunctor.Adequacy.}\texttt{\bfseries uvLaws\_mix\_exercised}}\hfill{\footnotesize \path{lean/su-functor/SuFunctor/Adequacy/QuantumContainment.lean}, line 620}\par\nopagebreak
\begin{leansrc}
theorem uvLaws_mix_exercised :
    QNet.quantumLaws isingQubitFragment uvSeg
      ((FinPMF.mix upMember.pmf entMember.pmf (1 / 2)
        (by norm_num) (by norm_num)).toSubPMF)
\end{leansrc}
\par\medskip\noindent{\small \texttt{SuFunctor.Adequacy.}\texttt{\bfseries pushforward\_probRow\_bridge}}\hfill{\footnotesize \path{lean/su-functor/SuFunctor/Adequacy/QuantumContainment.lean}, line 1302}\par\nopagebreak
\begin{leansrc}
theorem pushforward_probRow_bridge (c : Fin 2) :
    (uvMeasure.pushforward Prod.fst).probRow c
      = (((entMember.law.map Prod.fst).wt c : ℚ) : ℝ)
\end{leansrc}
\par\medskip\noindent{\small \texttt{SuFunctor.Adequacy.}\texttt{\bfseries eftTower\_offDiagMass\_antitone}}\hfill{\footnotesize \path{lean/su-functor/SuFunctor/Adequacy/QuantumContainment.lean}, line 1276}\par\nopagebreak
\begin{leansrc}
theorem eftTower_offDiagMass_antitone :
    Antitone fun n => (eftStage n).offDiagMass
\end{leansrc}
\par\medskip\noindent{\small \texttt{SuFunctor.Adequacy.}\texttt{\bfseries eftStage\_classical}}\hfill{\footnotesize \path{lean/su-functor/SuFunctor/Adequacy/QuantumContainment.lean}, line 1268}\par\nopagebreak
\begin{leansrc}
theorem eftStage_classical (n : ℕ) :
    ∃ (p : eftTowerX n → ℝ) (hp : ∀ c, 0 ≤ p c) (hs : ∑ c, p c = 1),
      eftStage n = QuantumMeasure.classical p hp hs
\end{leansrc}
\par\medskip\noindent{\small \texttt{SuFunctor.Adequacy.}\texttt{\bfseries recordShadow\_kchanApply}}\hfill{\footnotesize \path{lean/su-functor/SuFunctor/Adequacy/QuantumCompression.lean}, line 317}\par\nopagebreak
\begin{leansrc}
theorem recordShadow_kchanApply (A : J → Matrix e d ℂ)
    (S : κ' → Matrix e e ℂ) (ρ : Matrix d d ℂ) (c : κ') :
    recordShadow S (kchanApply A ρ) c
      = recordShadow (fun c => kchanDual A (S c)) ρ c
\end{leansrc}
\par\medskip\noindent{\small \texttt{SuFunctor.Adequacy.}\texttt{\bfseries kchanApply\_ptrKraus\_eq\_ptR}}\hfill{\footnotesize \path{lean/su-functor/SuFunctor/Adequacy/QuantumCompression.lean}, line 450}\par\nopagebreak
\begin{leansrc}
theorem kchanApply_ptrKraus_eq_ptR (ρ : TwoQubits) :
    kchanApply ptrKraus ρ = ptR ρ
\end{leansrc}
\par\medskip\noindent{\small \texttt{SuFunctor.Adequacy.}\texttt{\bfseries uv\_smallIR\_containment}}\hfill{\footnotesize \path{lean/su-functor/SuFunctor/Adequacy/QuantumCompression.lean}, line 718}\par\nopagebreak
\begin{leansrc}
theorem uv_smallIR_containment :
    Containment (QNet.quantumLaws isingQubitFragment uvSeg) Prod.fst
      (QNet.quantumLaws singleQubitFragment smallSet) (49 / 625)
\end{leansrc}
\par\medskip\noindent{\small \texttt{SuFunctor.Adequacy.QNet.}\texttt{\bfseries QNetRows}}\hfill{\footnotesize \path{lean/su-functor/SuFunctor/Adequacy/QuantumNet.lean}, line 251}\par\nopagebreak
\begin{leansrc}
structure QNetRows (κ : Type*) where …
\end{leansrc}
\par\medskip\noindent Modules cited: \texttt{SuFunctor.Adequacy.QuantumNet}, \texttt{SuFunctor.Adequacy.Ode}, \texttt{SuFunctor.Adequacy.Pde}, \texttt{SuFunctor.Adequacy.PdeRows}, \texttt{QtslAux.SdeGenerator}, \texttt{QtslAux.SdeGeneratorPoly}, \texttt{QtslAux.SdeMartingaleProblem}, \texttt{SuFunctor.Adequacy.Rode}, \texttt{QtslAux.JumpDiffusion}, \texttt{QtslAux.JumpDiffusionMartingaleProblem}, \texttt{SuFunctor.Adequacy.CtmcJumpChain}, \texttt{SuFunctor.Adequacy.MarkovAutomaton}, \texttt{SuFunctor.Adequacy.MaDemonic}, \texttt{SuFunctor.Adequacy.Bag}, \texttt{SuFunctor.TokenGame}, \texttt{SuFunctor.Adequacy.StochasticPetri}, \texttt{SuFunctor.Adequacy.StochasticPetriCascade}, \texttt{SuFunctor.Adequacy.StochasticPetriDemonic}, \texttt{SuFunctor.Adequacy.StochasticPetriSetFiring}, \texttt{SuFunctor.Adequacy.StochasticPetriFlatten}, \texttt{SuFunctor.Adequacy.StochasticPetriFlattenAgreement}, \texttt{SuFunctor.Adequacy.StochasticPetriComposed}, \texttt{QtslAux.KilledTimeIndexed}, \texttt{SuFunctor.Adequacy.StochasticPetriColouredComposed}, \texttt{SuFunctor.Adequacy.StochasticPetriAngelic}, \texttt{SuFunctor.Adequacy.StochasticPetriColoured}, \texttt{SuFunctor.Adequacy.StochasticPetriColouredCascade}, \texttt{SuFunctor.Adequacy.StochasticPetriColouredDemonic}, \texttt{SuFunctor.Adequacy.PetriVanilla}, \texttt{Qtsl.ScmTwin}, \texttt{SuFunctor.Stochastic.NoiseKleisli}, \texttt{SuFunctor.Adequacy.ProbProg}, \texttt{QtslAux.QuantumCoarseGraining}, \texttt{SuFunctor.Adequacy.QuantumContainment}.
\par\medskip\noindent Also cited in this chapter, listed under an earlier one: \texttt{QtslAux.QuantumAssembly}, \texttt{QtslAux.QuantumRecords}, \texttt{SuFunctor.Adequacy.DelayLine}, \texttt{SuFunctor.Adequacy.Dtmc}, \texttt{SuFunctor.Adequacy.EnvelopeUnion}, \texttt{SuFunctor.Adequacy.HeatTorus}, \texttt{SuFunctor.Adequacy.HeatTorusInit}, \texttt{SuFunctor.Adequacy.Inclusions}, \texttt{SuFunctor.Adequacy.MarginalGlue}, \texttt{SuFunctor.Adequacy.Mdp}, \texttt{SuFunctor.Adequacy.MdpDemonic}, \texttt{SuFunctor.Adequacy.Pgm}, \texttt{SuFunctor.Adequacy.SolutionSetMesh}, \texttt{copy\_ne\_kleisli}, \texttt{maxSlice\_unionRays}, \texttt{maxSlice\_unionRays\_subset}, \texttt{noiseDenot\_seqComp}.

\section{Systems: Composition Patterns}
Citations from Chapter~\ref{ch:qtsl-composition-patterns}.

\par\medskip\noindent{\small \texttt{SuFunctor.AlignedLoose.}\texttt{\bfseries seq\_assoc}}\hfill{\footnotesize \path{lean/su-functor/SuFunctor/HypernetComposition.lean}, line 219}\par\nopagebreak
\begin{leansrc}
theorem AlignedLoose.seq_assoc (m : AlignedLoose X Y) (n : AlignedLoose Y Z)
    (p : AlignedLoose Z V) : (m.seq n).seq p = m.seq (n.seq p)
\end{leansrc}
\par\medskip\noindent{\small \texttt{SuFunctor.}\texttt{\bfseries austereDnf\_subst}}\hfill{\footnotesize \path{lean/su-functor/SuFunctor/HypernetComposition.lean}, line 350}\par\nopagebreak
\begin{leansrc}
theorem austereDnf_subst (m : AlignedLoose X Y) (n : AlignedLoose Y Z)
    {T₁ T₁' T₂ T₂' : List PrimApp}
    (h₁ : austereDnf ty m T₁ = austereDnf ty m T₁')
    (h₂ : austereDnf ty n T₂ = austereDnf ty n T₂')
    (hTC₁ : ∀ a ∈ T₁, a.TypeConsistent (tableOn ty m.apex))
    (hw₁ : ∀ a ∈ T₁, ∀ w ∈ a.inputWires ++ a.outputWires, w ∈ m.apex)
    (hTC₂ : ∀ a ∈ T₂, a.TypeConsistent (tableOn ty n.apex))
    (hw₂ : ∀ a ∈ T₂, ∀ w ∈ a.inputWires ++ a.outputWires, w ∈ n.apex)
    (hTC₁' : ∀ a ∈ T₁', a.TypeConsistent (tableOn ty m.apex))
    (hw₁' : ∀ a ∈ T₁', ∀ w ∈ a.inputWires ++ a.outputWires, w ∈ m.apex)
    (hTC₂' : ∀ a ∈ T₂', a.TypeConsistent (tableOn ty n.apex))
    (hw₂' : ∀ a ∈ T₂', ∀ w ∈ a.inputWires ++ a.outputWires, w ∈ n.apex) :
    austereDnf ty (m.seq n) (T₁ ++ T₂)
      = austereDnf ty (m.seq n) (T₁' ++ T₂')
\end{leansrc}
\par\medskip\noindent{\small \texttt{SuFunctor.}\texttt{\bfseries admRegion\_append\_empty\_left}}\hfill{\footnotesize \path{lean/su-functor/SuFunctor/HypernetComposition.lean}, line 373}\par\nopagebreak
\begin{leansrc}
theorem admRegion_append_empty_left (m : AlignedLoose X Y)
    (n : AlignedLoose Y Z) {T₁ T₂ : List PrimApp}
    (h : AdmRegion ty m T₁ = ∅)
    (hTC₁ : ∀ a ∈ T₁, a.TypeConsistent (tableOn ty m.apex))
    (hw₁ : ∀ a ∈ T₁, ∀ w ∈ a.inputWires ++ a.outputWires, w ∈ m.apex)
    (hTC₂ : ∀ a ∈ T₂, a.TypeConsistent (tableOn ty n.apex))
    (hw₂ : ∀ a ∈ T₂, ∀ w ∈ a.inputWires ++ a.outputWires, w ∈ n.apex) :
    AdmRegion ty (m.seq n) (T₁ ++ T₂) = ∅
\end{leansrc}
\par\medskip\noindent{\small \texttt{SuFunctor.}\texttt{\bfseries share\_contradiction}}\hfill{\footnotesize \path{lean/su-functor/SuFunctor/HypernetComposition.lean}, line 465}\par\nopagebreak
\begin{leansrc}
theorem share_contradiction {w : Wire} {τ : TensorType}
    (c₁ c₂ : TensorValue τ) {i₀ : ℕ}
    (hne : c₁.entryE i₀ ≠ c₂.entryE i₀) :
    AdmRegion ty
        ((AlignedLoose.unit {w}).seq (AlignedLoose.unit {w}))
        ([PrimApp.const τ c₁ w] ++ [PrimApp.const τ c₂ w]) = ∅
\end{leansrc}
\par\medskip\noindent{\small \texttt{SuFunctor.}\texttt{\bfseries theta\_rename\_seq}}\hfill{\footnotesize \path{lean/su-functor/SuFunctor/HypernetComposition.lean}, line 564}\par\nopagebreak
\begin{leansrc}
theorem theta_rename_seq {m : AlignedLoose X Y} {n : AlignedLoose Y Z}
    {m' : AlignedLoose X' Y'} {n' : AlignedLoose Y' Z'}
    (c₁ : AlignedCell ty m m') (c₂ : AlignedCell ty n n')
    (hmap : c₁.map = c₂.map) (T₁ T₂ : List PrimApp)
    (hTC₁ : ∀ a ∈ T₁, a.TypeConsistent (tableOn ty m.apex))
    (hw₁ : ∀ a ∈ T₁, ∀ w ∈ a.inputWires ++ a.outputWires, w ∈ m.apex)
    (hTC₂ : ∀ a ∈ T₂, a.TypeConsistent (tableOn ty n.apex))
    (hw₂ : ∀ a ∈ T₂, ∀ w ∈ a.inputWires ++ a.outputWires, w ∈ n.apex)
    (hTC₁' : ∀ a ∈ T₁.map (PrimApp.rename c₁.map),
      a.TypeConsistent (tableOn ty m'.apex))
    (hw₁' : ∀ a ∈ T₁.map (PrimApp.rename c₁.map),
      ∀ w ∈ a.inputWires ++ a.outputWires, w ∈ m'.apex)
    (hTC₂' : ∀ a ∈ T₂.map (PrimApp.rename c₂.map),
      a.TypeConsistent (tableOn ty n'.apex))
    (hw₂' : ∀ a ∈ T₂.map (PrimApp.rename c₂.map),
      ∀ w ∈ a.inputWires ++ a.outputWires, w ∈ n'.apex) :
    dnfPull (renameTot ty (AlignedCell.seqCell ty c₁ c₂ hmap))
        (austereDnf ty (m.seq n) (T₁ ++ T₂))
      = dnfPull (splitTot ty m' n')
          (dnfBoxplus
            (dnfPull (renameTot ty c₁) (austereDnf ty m T₁))
            (dnfPull (renameTot ty c₂) (austereDnf ty n T₂)))
\end{leansrc}
\par\medskip\noindent{\small \texttt{SuFunctor.}\texttt{\bfseries admRegion\_append}}\hfill{\footnotesize \path{lean/su-functor/SuFunctor/HypernetInclusion.lean}, line 728}\par\nopagebreak
\begin{leansrc}
theorem admRegion_append (m : AlignedLoose X Y) (n : AlignedLoose Y Z)
    (T₁ T₂ : List PrimApp)
    (hTC₁ : ∀ a ∈ T₁, a.TypeConsistent (tableOn ty m.apex))
    (hw₁ : ∀ a ∈ T₁, ∀ w ∈ a.inputWires ++ a.outputWires, w ∈ m.apex)
    (hTC₂ : ∀ a ∈ T₂, a.TypeConsistent (tableOn ty n.apex))
    (hw₂ : ∀ a ∈ T₂, ∀ w ∈ a.inputWires ++ a.outputWires, w ∈ n.apex) :
    AdmRegion ty (m.seq n) (T₁ ++ T₂)
      = splitTot ty m n ⁻¹'
          (Pat.compFst (alignedPat ty m) (alignedPat ty n) ⁻¹'
              AdmRegion ty m T₁
            ∩ Pat.compSnd (alignedPat ty m) (alignedPat ty n) ⁻¹'
              AdmRegion ty n T₂)
\end{leansrc}
\par\medskip\noindent{\small \texttt{SuFunctor.}\texttt{\bfseries share\_feasible}}\hfill{\footnotesize \path{lean/su-functor/SuFunctor/HypernetComposition.lean}, line 442}\par\nopagebreak
\begin{leansrc}
theorem share_feasible {w : Wire} {τ : TensorType} (hty : ty w = some τ)
    (c : TensorValue τ) :
    (AdmRegion ty (AlignedLoose.unit {w})
      [PrimApp.const τ c w]).Nonempty
\end{leansrc}
\par\medskip\noindent{\small \texttt{SuFunctor.Spacetime.}\texttt{\bfseries hybridTime\_square\_subset}}\hfill{\footnotesize \path{lean/su-functor/SuFunctor/Adequacy/HybridTime.lean}, line 489}\par\nopagebreak
\begin{leansrc}
theorem hybridTime_square_subset
    (D : SpecDnf (Joint (W → ClockedEl Clk (SiteRec κ X)))) :
    ((D.withWall (sampleWireWall κ X ports i i' φ)).withWall
        (possWall (perTickRowsSet κ X ports i i' φ Pc))).modelClass ⊆
      ((D.withWall (sampleWireWall κ X ports i i' φ)).withWall
        (possWall (coarseRowsSet κ X ports i Pc))).modelClass
\end{leansrc}
\par\medskip\noindent{\small \texttt{SuFunctor.Spacetime.}\texttt{\bfseries tickSeam\_hybridTime\_square\_eq}}\hfill{\footnotesize \path{lean/su-functor/SuFunctor/Adequacy/HybridTime.lean}, line 552}\par\nopagebreak
\begin{leansrc}
theorem tickSeam_hybridTime_square_eq {X : Type u}
    (Pc : SpacetimeRec (tickClkOf.{u} ⟨false⟩) X → Prop)
    (D : SpecDnf (Joint (ULift.{u} Bool →
      ClockedEl (ULift.{u} Bool) (SiteRec tickClkOf X)))) :
    ((D.withWall (sampleWireWall tickClkOf X tickPorts ⟨false⟩ ⟨true⟩
          tickSeamHom)).withWall
        (possWall (perTickRowsSet tickClkOf X tickPorts ⟨false⟩ ⟨true⟩
          tickSeamHom Pc))).modelClass =
      ((D.withWall (sampleWireWall tickClkOf X tickPorts ⟨false⟩ ⟨true⟩
          tickSeamHom)).withWall
        (possWall (coarseRowsSet tickClkOf X tickPorts ⟨false⟩
          Pc))).modelClass
\end{leansrc}
\par\medskip\noindent{\small \texttt{SuFunctor.Plain.}\texttt{\bfseries funnelComp\_breaks\_unit\_coherence}}\hfill{\footnotesize \path{lean/su-functor/SuFunctor/Plain/Battery.lean}, line 177}\par\nopagebreak
\begin{leansrc}
theorem funnelComp_breaks_unit_coherence (M : Pat A B) :
    ¬ ∀ 𝒰 : Cpl M,
      Pen.pull (PatCell.tot
          (PseudoDoubleCategoryStruct.lLeftUnitor (C := WireSpan El) M)) ''
          funnelComp ({unitPen A} : Cpl (Pat.id A)) 𝒰 = 𝒰
\end{leansrc}
\par\medskip\noindent{\small \texttt{SuFunctor.Plain.}\texttt{\bfseries unitClass\_transparent}}\hfill{\footnotesize \path{lean/su-functor/SuFunctor/Plain/Battery.lean}, line 86}\par\nopagebreak
\begin{leansrc}
theorem unitClass_transparent (M : Pat A B) (𝒰 : Cpl M) :
    Pen.pull (PatCell.tot
        (PseudoDoubleCategoryStruct.lLeftUnitor (C := WireSpan El) M)) ''
        Set.image2 boxplus ({unitPen A} : Cpl (Pat.id A)) 𝒰 = 𝒰
\end{leansrc}
\par\medskip\noindent{\small \texttt{SuFunctor.Adequacy.}\texttt{\bfseries funnel\_coupledSat\_of\_parts}}\hfill{\footnotesize \path{lean/su-functor/SuFunctor/Adequacy/Funnel.lean}, line 195}\par\nopagebreak
\begin{leansrc}
theorem funnel_coupledSat_of_parts (F : FunnelData K) {σ₁ σ₂ : K → ℝ≥0∞}
    (h₁ : σ₁ ∈ F.transit.modelClass) (h₂ : σ₂ ∈ F.hold.modelClass)
    (hcon : ∀ w ∈ F.composite.marks, Dip (σ₁ + σ₂) w) :
    σ₁ + σ₂ ∈ F.composite.modelClass
\end{leansrc}
\par\medskip\noindent{\small \texttt{SuFunctor.Adequacy.FunnelParts.}\texttt{\bfseries assemble\_isSome\_iff}}\hfill{\footnotesize \path{lean/su-functor/SuFunctor/Adequacy/Funnel.lean}, line 138}\par\nopagebreak
\begin{leansrc}
theorem assemble_isSome_iff {p : FunnelParts K} :
    p.assemble.isSome ↔ p.transitTarget = p.holdInvariant ∧ 0 < p.budget
\end{leansrc}
\par\medskip\noindent{\small \texttt{SuFunctor.}\texttt{\bfseries vorobev\_record\_glue}}\hfill{\footnotesize \path{lean/su-functor/SuFunctor/Adequacy/MarginalGlue.lean}, line 274}\par\nopagebreak
\begin{leansrc}
theorem vorobev_record_glue (v : Bool × Bool × Bool) :
    (0 < pairCorr.wt (proj12 v) ∧ 0 < pairCorr.wt (proj23 v) ∧
        0 < pairIndep.wt (proj13 v)) ↔
      v = (false, false, false) ∨ v = (true, true, true)
\end{leansrc}
\par\medskip\noindent Modules cited: \texttt{SuFunctor.HypernetComposition}, \texttt{SuFunctor.Adequacy.HybridTime}, \texttt{SuFunctor.Adequacy.Funnel}, \texttt{SuFunctor.Adequacy.StochasticPetriSupervise}.
\par\medskip\noindent Also cited in this chapter, listed under an earlier one: \texttt{Qtsl.ScmTwin}, \texttt{SuFunctor.Adequacy.MarginalGlue}, \texttt{SuFunctor.HypernetInclusion}, \texttt{SuFunctor.Plain.Battery}, \texttt{SuFunctor.SingleSiteSemanticUniverse}, \texttt{coupledSat\_comp\_add\_of\_dips}, \texttt{credalCompile\_vorobevGlue}, \texttt{glue\_admissible\_iff}, \texttt{theta\_comp}, \texttt{vorobev\_no\_joint}.

\section{Specification Language}
Citations from Chapter~\ref{ch:qtsl-specification-language}.

\par\medskip\noindent{\small \texttt{SuFunctor.Universe.}\texttt{\bfseries Sat}}\hfill{\footnotesize \path{lean/su-functor/SuFunctor/Universe/Judgment.lean}, line 54}\par\nopagebreak
\begin{leansrc}
def Sat (σ φ : Penalty V) : Prop
\end{leansrc}
\par\medskip\noindent{\small \texttt{SuFunctor.Universe.}\texttt{\bfseries wallOn\_le\_iff}}\hfill{\footnotesize \path{lean/su-functor/SuFunctor/Universe/SpecElab.lean}, line 120}\par\nopagebreak
\begin{leansrc}
theorem wallOn_le_iff {C : Set K} {σ : K → ℝ≥0∞} :
    wallOn C ≤ σ ↔ permitted σ ⊆ C
\end{leansrc}
\par\medskip\noindent{\small \texttt{SuFunctor.Universe.}\texttt{\bfseries dip\_top\_iff}}\hfill{\footnotesize \path{lean/su-functor/SuFunctor/Universe/SpecElab.lean}, line 137}\par\nopagebreak
\begin{leansrc}
theorem dip_top_iff {σ : K → ℝ≥0∞} :
    Dip σ (⊤ : K → ℝ≥0∞) ↔ (permitted σ).Nonempty
\end{leansrc}
\par\medskip\noindent{\small \texttt{SuFunctor.}\texttt{\bfseries SatCredal}}\hfill{\footnotesize \path{lean/su-functor/SuFunctor/Adequacy/Envelope.lean}, line 289}\par\nopagebreak
\begin{leansrc}
def SatCredal (k : EnvelopeKind) (C c : SubPMF V → Prop) : Prop
\end{leansrc}
\par\medskip\noindent{\small \texttt{SuFunctor.Universe.SpecRecord.}\texttt{\bfseries elabSpec\_top\_mem}}\hfill{\footnotesize \path{lean/su-functor/SuFunctor/Universe/SpecElab.lean}, line 935}\par\nopagebreak
\begin{leansrc}
theorem SpecRecord.elabSpec_top_mem {width : α → ℕ} {A : α → Set K}
    {r : SpecRecord α} {D : SpecDnf K}
    (_h : r.elabSpec width A = .ok D) :
    ∀ a ∈ D.atoms, (⊤ : K → ℝ≥0∞) ∈ a.marks
\end{leansrc}
\par\medskip\noindent{\small \texttt{SuFunctor.Universe.SpecRecord.}\texttt{\bfseries elabSpec}}\hfill{\footnotesize \path{lean/su-functor/SuFunctor/Universe/SpecElab.lean}, line 902}\par\nopagebreak
\begin{leansrc}
noncomputable def SpecRecord.elabSpec (width : α → ℕ) (A : α → Set K)
    (r : SpecRecord α) : Except (ConformanceError α) (SpecDnf K)
\end{leansrc}
\par\medskip\noindent{\small \texttt{SuFunctor.Universe.SpecRecord.}\texttt{\bfseries resolveArity\_isOk\_iff}}\hfill{\footnotesize \path{lean/su-functor/SuFunctor/Universe/SpecElab.lean}, line 742}\par\nopagebreak
\begin{leansrc}
theorem SpecRecord.resolveArity_isOk_iff {width : α → ℕ}
    {r : SpecRecord α} :
    IsOk (r.resolveArity width) ↔ r.Conformant width
\end{leansrc}
\par\medskip\noindent{\small \texttt{SuFunctor.Universe.SpecRecord.}\texttt{\bfseries resolveArity\_error\_names\_offender}}\hfill{\footnotesize \path{lean/su-functor/SuFunctor/Universe/SpecElab.lean}, line 835}\par\nopagebreak
\begin{leansrc}
theorem SpecRecord.resolveArity_error_names_offender {width : α → ℕ}
    {r : SpecRecord α} {a : α} {f e : ℕ} {b : Bool}
    (h : r.resolveArity width = .error (.widthMismatch a f e b)) :
    width a = f ∧ f ≠ e ∧
      ((b = false ∧ a ∈ r.root.atoms) ∨
        (b = true ∧ ∃ I, r.iface? = some I ∧
          ∃ φ, I.assumption? = some φ ∧ a ∈ φ.atoms))
\end{leansrc}
\par\medskip\noindent{\small \texttt{SuFunctor.Universe.SpecRecord.}\texttt{\bfseries elabSpec\_isOk\_iff}}\hfill{\footnotesize \path{lean/su-functor/SuFunctor/Universe/SpecElab.lean}, line 909}\par\nopagebreak
\begin{leansrc}
theorem SpecRecord.elabSpec_isOk_iff {width : α → ℕ} {A : α → Set K}
    {r : SpecRecord α} :
    IsOk (r.elabSpec width A) ↔ r.Conformant width
\end{leansrc}
\par\medskip\noindent{\small \texttt{SuFunctor.Universe.SpecRecord.}\texttt{\bfseries elabSpec\_ok\_modelClass}}\hfill{\footnotesize \path{lean/su-functor/SuFunctor/Universe/SpecElab.lean}, line 918}\par\nopagebreak
\begin{leansrc}
theorem SpecRecord.elabSpec_ok_modelClass {width : α → ℕ}
    {A : α → Set K} {r : SpecRecord α} {D : SpecDnf K}
    (h : r.elabSpec width A = .ok D) :
    D.modelClass = NnfForm.classSem A r.root.nnf
\end{leansrc}
\par\medskip\noindent{\small \texttt{SuFunctor.Universe.SpecDnf.}\texttt{\bfseries modelClass\_top}}\hfill{\footnotesize \path{lean/su-functor/SuFunctor/Universe/SpecElab.lean}, line 323}\par\nopagebreak
\begin{leansrc}
theorem SpecDnf.modelClass_top :
    (SpecDnf.top : SpecDnf K).modelClass = safeProper Set.univ
\end{leansrc}
\par\medskip\noindent{\small \texttt{SuFunctor.Universe.SpecDnf.}\texttt{\bfseries modelClass\_bot}}\hfill{\footnotesize \path{lean/su-functor/SuFunctor/Universe/SpecElab.lean}, line 328}\par\nopagebreak
\begin{leansrc}
theorem SpecDnf.modelClass_bot :
    (SpecDnf.bot : SpecDnf K).modelClass = ∅
\end{leansrc}
\par\medskip\noindent{\small \texttt{SuFunctor.Universe.}\texttt{\bfseries den\_compl\_le}}\hfill{\footnotesize \path{lean/su-functor/SuFunctor/Universe/SpecElab.lean}, line 601}\par\nopagebreak
\begin{leansrc}
theorem den_compl_le (f : K → ℝ) (c : ℚ) :
    {μ | f μ ≤ (c : ℝ)}ᶜ = {μ | (c : ℝ) < f μ}
\end{leansrc}
\par\medskip\noindent{\small \texttt{SuFunctor.Universe.}\texttt{\bfseries safetyAtom}}\hfill{\footnotesize \path{lean/su-functor/SuFunctor/Universe/SpecElab.lean}, line 176}\par\nopagebreak
\begin{leansrc}
noncomputable def safetyAtom (C : Set K) : SpecAtom K where …
\end{leansrc}
\par\medskip\noindent{\small \texttt{SuFunctor.Universe.}\texttt{\bfseries coupledSat\_safetyAtom\_iff}}\hfill{\footnotesize \path{lean/su-functor/SuFunctor/Universe/SpecElab.lean}, line 191}\par\nopagebreak
\begin{leansrc}
theorem coupledSat_safetyAtom_iff {C : Set K} {σ : K → ℝ≥0∞} :
    CoupledSat σ (wallOn C) {⊤} ↔
      permitted σ ⊆ C ∧ (permitted σ).Nonempty
\end{leansrc}
\par\medskip\noindent{\small \texttt{SuFunctor.Universe.}\texttt{\bfseries safetyAtom\_modelClass}}\hfill{\footnotesize \path{lean/su-functor/SuFunctor/Universe/SpecElab.lean}, line 205}\par\nopagebreak
\begin{leansrc}
theorem safetyAtom_modelClass (C : Set K) :
    (safetyAtom C).modelClass = safeProper C
\end{leansrc}
\par\medskip\noindent{\small \texttt{SuFunctor.Universe.SpecAtom.}\texttt{\bfseries modelClass\_merge}}\hfill{\footnotesize \path{lean/su-functor/SuFunctor/Universe/SpecElab.lean}, line 243}\par\nopagebreak
\begin{leansrc}
theorem SpecAtom.modelClass_merge (a b : SpecAtom K) :
    (a.merge b).modelClass = a.modelClass ∩ b.modelClass
\end{leansrc}
\par\medskip\noindent{\small \texttt{SuFunctor.Universe.SpecDnf.}\texttt{\bfseries modelClass\_conj}}\hfill{\footnotesize \path{lean/su-functor/SuFunctor/Universe/SpecElab.lean}, line 292}\par\nopagebreak
\begin{leansrc}
theorem SpecDnf.modelClass_conj (D E : SpecDnf K) :
    (D.conj E).modelClass = D.modelClass ∩ E.modelClass
\end{leansrc}
\par\medskip\noindent{\small \texttt{SuFunctor.Universe.SpecDnf.}\texttt{\bfseries modelClass\_disj}}\hfill{\footnotesize \path{lean/su-functor/SuFunctor/Universe/SpecElab.lean}, line 308}\par\nopagebreak
\begin{leansrc}
theorem SpecDnf.modelClass_disj (D E : SpecDnf K) :
    (D.disj E).modelClass = D.modelClass ∪ E.modelClass
\end{leansrc}
\par\medskip\noindent{\small \texttt{SuFunctor.Universe.NnfForm.}\texttt{\bfseries exists\_hedging\_gap}}\hfill{\footnotesize \path{lean/su-functor/SuFunctor/Universe/SpecElab.lean}, line 557}\par\nopagebreak
\begin{leansrc}
theorem exists_hedging_gap :
    ∃ (A : Bool → Set Bool) (ψ χ : NnfForm Bool) (σ : Bool → ℝ≥0∞),
      σ ∈ safeProper ((NnfForm.disj ψ χ).den A) ∧
        σ ∉ classSem A (NnfForm.disj ψ χ)
\end{leansrc}
\par\medskip\noindent{\small \texttt{SuFunctor.Universe.NnfForm.}\texttt{\bfseries elabNnf\_modelClass\_disjFree}}\hfill{\footnotesize \path{lean/su-functor/SuFunctor/Universe/SpecElab.lean}, line 546}\par\nopagebreak
\begin{leansrc}
theorem elabNnf_modelClass_disjFree (A : α → Set K) (ψ : NnfForm α)
    (h : ψ.DisjFree) :
    (elabNnf A ψ).modelClass = safeProper (ψ.den A)
\end{leansrc}
\par\medskip\noindent{\small \texttt{SuFunctor.Universe.StateFormula.}\texttt{\bfseries den\_nnf}}\hfill{\footnotesize \path{lean/su-functor/SuFunctor/Universe/SpecElab.lean}, line 462}\par\nopagebreak
\begin{leansrc}
theorem den_nnf (A : α → Set K) (φ : StateFormula α) :
    (φ.nnf).den A = φ.den A
\end{leansrc}
\par\medskip\noindent{\small \texttt{SuFunctor.Universe.}\texttt{\bfseries recordCon\_iff\_stages}}\hfill{\footnotesize \path{lean/su-functor/SuFunctor/Adequacy/SpecPathLens.lean}, line 720}\par\nopagebreak
\begin{leansrc}
theorem recordCon_iff_stages {p : NonemptyCompacts X × ProbState X} :
    p ∈ recordCon ↔ stages .dwell p ⊆ stages .touch p
\end{leansrc}
\par\medskip\noindent{\small \texttt{SuFunctor.Universe.}\texttt{\bfseries exists\_untilOn\_species\_incomparable}}\hfill{\footnotesize \path{lean/su-functor/SuFunctor/Adequacy/SpecPathLens.lean}, line 840}\par\nopagebreak
\begin{leansrc}
theorem exists_untilOn_species_incomparable :
    ∃ p ∈ recordCon (X := Bool),
      (untilOn (fun b => cond b 1 0) 1 ∅ {false} (stages .touch p) ∧
        ¬ untilOn (fun b => cond b 1 0) 1 ∅ {false} (stages .dwell p)) ∧
      (untilOn (fun b => cond b 1 0) 1 {true} {true} (stages .dwell p) ∧
        ¬ untilOn (fun b => cond b 1 0) 1 {true} {true}
          (stages .touch p))
\end{leansrc}
\par\medskip\noindent{\small \texttt{SuFunctor.Universe.}\texttt{\bfseries exists\_streettOn\_species\_incomparable}}\hfill{\footnotesize \path{lean/su-functor/SuFunctor/Adequacy/SpecPathLens.lean}, line 888}\par\nopagebreak
\begin{leansrc}
theorem exists_streettOn_species_incomparable :
    ∃ p ∈ recordCon (X := Bool),
      (streettOn [(({true} : Set Bool), ({false} : Set Bool))]
          (stages .touch p) ∧
        ¬ streettOn [(({true} : Set Bool), ({false} : Set Bool))]
          (stages .dwell p)) ∧
      (streettOn [(({false} : Set Bool), (∅ : Set Bool))]
          (stages .dwell p) ∧
        ¬ streettOn [(({false} : Set Bool), (∅ : Set Bool))]
          (stages .touch p))
\end{leansrc}
\par\medskip\noindent{\small \texttt{SuFunctor.Universe.}\texttt{\bfseries touchWall\_eq\_wallOn}}\hfill{\footnotesize \path{lean/su-functor/SuFunctor/Adequacy/SpecPathLens.lean}, line 744}\par\nopagebreak
\begin{leansrc}
theorem touchWall_eq_wallOn (Bad : Set X) :
    touchWall Bad
      = wallOn {p : NonemptyCompacts X × ProbState X |
          alwaysOn Badᶜ (stages .touch p)}
\end{leansrc}
\par\medskip\noindent{\small \texttt{SuFunctor.Universe.}\texttt{\bfseries dwellWall\_eq\_wallOn}}\hfill{\footnotesize \path{lean/su-functor/SuFunctor/Adequacy/SpecPathLens.lean}, line 761}\par\nopagebreak
\begin{leansrc}
theorem dwellWall_eq_wallOn (Bad : Set X) :
    dwellWall Bad
      = wallOn {p : NonemptyCompacts X × ProbState X |
          alwaysOn Badᶜ (stages .dwell p)}
\end{leansrc}
\par\medskip\noindent{\small \texttt{SuFunctor.Universe.}\texttt{\bfseries markOn}}\hfill{\footnotesize \path{lean/su-functor/SuFunctor/Adequacy/SpecPathLens.lean}, line 100}\par\nopagebreak
\begin{leansrc}
noncomputable def markOn (C : Set K) : K → ℝ≥0∞
\end{leansrc}
\par\medskip\noindent{\small \texttt{SuFunctor.Universe.}\texttt{\bfseries markOn\_univ}}\hfill{\footnotesize \path{lean/su-functor/SuFunctor/Adequacy/SpecPathLens.lean}, line 112}\par\nopagebreak
\begin{leansrc}
@[simp] theorem markOn_univ : markOn (Set.univ : Set K) = ⊤
\end{leansrc}
\par\medskip\noindent{\small \texttt{SuFunctor.Universe.}\texttt{\bfseries coupledSat\_pathAtom\_iff}}\hfill{\footnotesize \path{lean/su-functor/SuFunctor/Adequacy/SpecPathLens.lean}, line 167}\par\nopagebreak
\begin{leansrc}
theorem coupledSat_pathAtom_iff {C G : Set K} {σ : K → ℝ≥0∞} :
    CoupledSat σ (pathAtom C G).floor (pathAtom C G).marks ↔
      permitted σ ⊆ C ∧ (permitted σ ∩ G).Nonempty
\end{leansrc}
\par\medskip\noindent{\small \texttt{SuFunctor.Universe.}\texttt{\bfseries stages\_uniformLift}}\hfill{\footnotesize \path{lean/su-functor/SuFunctor/Adequacy/SpecPathLens.lean}, line 732}\par\nopagebreak
\begin{leansrc}
theorem stages_uniformLift (sp : RecordSpecies) (G : NonemptyCompacts X) :
    stages sp (G, uniformLift G) = (G : Set X)
\end{leansrc}
\par\medskip\noindent{\small \texttt{SuFunctor.Universe.}\texttt{\bfseries alwaysOn}}\hfill{\footnotesize \path{lean/su-functor/SuFunctor/Adequacy/SpecPathLens.lean}, line 266}\par\nopagebreak
\begin{leansrc}
def alwaysOn (Good : Set X) (S : Set X) : Prop
\end{leansrc}
\par\medskip\noindent{\small \texttt{SuFunctor.Universe.}\texttt{\bfseries isClosed\_recordClass\_always}}\hfill{\footnotesize \path{lean/su-functor/SuFunctor/Adequacy/SpecPathLens.lean}, line 1010}\par\nopagebreak
\begin{leansrc}
theorem isClosed_recordClass_always (sp : RecordSpecies) (τ : X → ℕ)
    (A : α → Set X) (φ : StateFormula α) :
    IsClosed (recordClass sp τ A (.always φ))
\end{leansrc}
\par\medskip\noindent{\small \texttt{SuFunctor.Universe.}\texttt{\bfseries windowAlwaysOn}}\hfill{\footnotesize \path{lean/su-functor/SuFunctor/Adequacy/SpecPathLens.lean}, line 270}\par\nopagebreak
\begin{leansrc}
def windowAlwaysOn (τ : X → ℕ) (n : ℕ) (Good : Set X) (S : Set X) :
    Prop
\end{leansrc}
\par\medskip\noindent{\small \texttt{SuFunctor.Universe.}\texttt{\bfseries isClosed\_recordClass\_alwaysT}}\hfill{\footnotesize \path{lean/su-functor/SuFunctor/Adequacy/SpecPathLens.lean}, line 1017}\par\nopagebreak
\begin{leansrc}
theorem isClosed_recordClass_alwaysT (sp : RecordSpecies) (τ : X → ℕ)
    (A : α → Set X) (n : ℕ) (φ : StateFormula α) :
    IsClosed (recordClass sp τ A (.alwaysT n φ))
\end{leansrc}
\par\medskip\noindent{\small \texttt{SuFunctor.Universe.}\texttt{\bfseries windowReachOn}}\hfill{\footnotesize \path{lean/su-functor/SuFunctor/Adequacy/SpecPathLens.lean}, line 276}\par\nopagebreak
\begin{leansrc}
def windowReachOn (τ : X → ℕ) (n : ℕ) (Goal : Set X) (S : Set X) :
    Prop
\end{leansrc}
\par\medskip\noindent{\small \texttt{SuFunctor.Universe.}\texttt{\bfseries isOpen\_recordClass\_deadline}}\hfill{\footnotesize \path{lean/su-functor/SuFunctor/Adequacy/SpecPathLens.lean}, line 1030}\par\nopagebreak
\begin{leansrc}
theorem isOpen_recordClass_deadline (sp : RecordSpecies) (τ : X → ℕ)
    (A : α → Set X) (n : ℕ) (φ : StateFormula α) (b : Bool) :
    IsOpen (recordClass sp τ A (.deadline n φ b))
\end{leansrc}
\par\medskip\noindent{\small \texttt{SuFunctor.Universe.}\texttt{\bfseries untilOn}}\hfill{\footnotesize \path{lean/su-functor/SuFunctor/Adequacy/SpecPathLens.lean}, line 283}\par\nopagebreak
\begin{leansrc}
def untilOn (τ : X → ℕ) (n : ℕ) (Safe Target : Set X) (S : Set X) :
    Prop
\end{leansrc}
\par\medskip\noindent{\small \texttt{SuFunctor.Universe.}\texttt{\bfseries recordClass\_untilF\_eq}}\hfill{\footnotesize \path{lean/su-functor/SuFunctor/Adequacy/SpecPathLens.lean}, line 1045}\par\nopagebreak
\begin{leansrc}
theorem recordClass_untilF_eq (sp : RecordSpecies) (τ : X → ℕ)
    (A : α → Set X) (n : ℕ) (sf tg : StateFormula α) :
    recordClass sp τ A (.untilF n sf tg)
      = ⋃ t : Fin (n + 1),
          ({p : NonemptyCompacts X × ProbState X |
              (stages sp p ∩ (tg.den A ∩ {x | τ x = (t : ℕ)})).Nonempty}
            ∩ {p : NonemptyCompacts X × ProbState X |
                stages sp p ⊆ sf.den A ∪ {x | ¬ τ x < (t : ℕ)}})
\end{leansrc}
\par\medskip\noindent{\small \texttt{SuFunctor.Universe.}\texttt{\bfseries isConstructibleUnion\_recordClass\_untilF}}\hfill{\footnotesize \path{lean/su-functor/SuFunctor/Adequacy/SpecPathLens.lean}, line 1062}\par\nopagebreak
\begin{leansrc}
theorem isConstructibleUnion_recordClass_untilF (sp : RecordSpecies)
    (τ : X → ℕ) (A : α → Set X) (n : ℕ) (sf tg : StateFormula α) :
    IsConstructibleUnion (recordClass sp τ A (.untilF n sf tg))
\end{leansrc}
\par\medskip\noindent{\small \texttt{Qtsl.QStreett.}\texttt{\bfseries sound\_QStreettBundle}}\hfill{\footnotesize \path{lean/qtsl-aux/QtslAux/QStreett.lean}, line 347}\par\nopagebreak
\begin{leansrc}
theorem sound_QStreettBundle {Rq G : ι → Set S} {W : ι → S → ℝ} {ε C : ι → ℝ}
    (hD : MeasurableSet D) (hDbar : MeasurableSet Dbar) (hU : MeasurableSet U)
    (hV₀ : Measurable V₀)
    (hRq : ∀ i, MeasurableSet (Rq i)) (hG : ∀ i, MeasurableSet (G i))
    (hW : ∀ i, Measurable (W i))
    (hΘ : p + thr ≤ 1)
    (hDsub : D ⊆ Dbar)
    (hR3 : ∀ y ∈ D, y ∉ U → K y Dbar = 1)
    (hR4 : ∀ y ∈ D, y ∉ U → K y (D ∪ U) = 1)
    (hV₀int : ∀ y ∈ D, y ∉ U → Integrable V₀ (K y))
    (hV₀drift : ∀ y ∈ D, y ∉ U → ∫ z, V₀ z ∂K y ≤ V₀ y)
    (hV₀floor : ∀ y ∈ U, y ∈ Dbar → 1 ≤ V₀ y)
    (hV₀nonneg : ∀ y ∈ Dbar, 0 ≤ V₀ y)
    (hε : ∀ i, 0 < ε i)
    (hW0 : ∀ i, ∀ y ∈ Dbar, 0 ≤ W i y) (hWC : ∀ i, ∀ y ∈ Dbar, W i y ≤ C i)
    (hWint : ∀ i, ∀ y ∈ D, y ∉ U → y ∉ G i → Integrable (W i) (K y))
    (hWdrift : ∀ i, ∀ y ∈ D, y ∉ U → y ∉ G i → ∫ z, W i z ∂K y ≤ W i y)
    (hWdriftR : ∀ i, ∀ y ∈ D, y ∉ U → y ∈ Rq i → y ∉ G i →
      ∫ z, W i z ∂K y ≤ W i y - ε i)
    (hx : x ∈ D) (hthr : V₀ x ≤ thr) :
    ENNReal.ofReal p ≤ pathMeasure K x
      {ω | ∀ i, (∃ᶠ n in atTop, ω n ∈ Rq i) → ∃ᶠ n in atTop, ω n ∈ G i}
\end{leansrc}
\par\medskip\noindent{\small \texttt{SuFunctor.Universe.}\texttt{\bfseries streettOn}}\hfill{\footnotesize \path{lean/su-functor/SuFunctor/Adequacy/SpecPathLens.lean}, line 292}\par\nopagebreak
\begin{leansrc}
def streettOn (prs : List (Set X × Set X)) (S : Set X) : Prop
\end{leansrc}
\par\medskip\noindent{\small \texttt{SuFunctor.Universe.}\texttt{\bfseries pathNnf\_notPath\_alwaysT}}\hfill{\footnotesize \path{lean/su-functor/SuFunctor/Adequacy/SpecPathLens.lean}, line 677}\par\nopagebreak
\begin{leansrc}
@[simp] theorem pathNnf_notPath_alwaysT (n : ℕ) (φ : StateFormula α) :
    pathNnf (PathFormula.notPath (.alwaysT n φ))
      = .ok (.deadline n φ.neg false)
\end{leansrc}
\par\medskip\noindent{\small \texttt{SuFunctor.Universe.}\texttt{\bfseries pathNnf\_notPath\_notPath\_deadline}}\hfill{\footnotesize \path{lean/su-functor/SuFunctor/Adequacy/SpecPathLens.lean}, line 683}\par\nopagebreak
\begin{leansrc}
@[simp] theorem pathNnf_notPath_notPath_deadline (n : ℕ)
    (φ : StateFormula α) (b : Bool) :
    pathNnf (PathFormula.notPath (.notPath (.deadline n φ b)))
      = .ok (.deadline n φ b)
\end{leansrc}
\par\medskip\noindent{\small \texttt{SuFunctor.Universe.}\texttt{\bfseries pathNnf\_notPath\_always}}\hfill{\footnotesize \path{lean/su-functor/SuFunctor/Adequacy/SpecPathLens.lean}, line 659}\par\nopagebreak
\begin{leansrc}
@[simp] theorem pathNnf_notPath_always (φ : StateFormula α) :
    pathNnf (PathFormula.notPath (.always φ))
      = .error .negatedAlways
\end{leansrc}
\par\medskip\noindent{\small \texttt{SuFunctor.Universe.}\texttt{\bfseries pathNnf\_notPath\_untilF}}\hfill{\footnotesize \path{lean/su-functor/SuFunctor/Adequacy/SpecPathLens.lean}, line 664}\par\nopagebreak
\begin{leansrc}
@[simp] theorem pathNnf_notPath_untilF (n : ℕ)
    (sf tg : StateFormula α) :
    pathNnf (PathFormula.notPath (.untilF n sf tg))
      = .error .negatedUntil
\end{leansrc}
\par\medskip\noindent{\small \texttt{SuFunctor.Universe.}\texttt{\bfseries pathNnf\_notPath\_streett}}\hfill{\footnotesize \path{lean/su-functor/SuFunctor/Adequacy/SpecPathLens.lean}, line 671}\par\nopagebreak
\begin{leansrc}
@[simp] theorem pathNnf_notPath_streett (prs : List (StreettPair α)) :
    pathNnf (PathFormula.notPath (.streett prs))
      = .error .negatedStreett
\end{leansrc}
\par\medskip\noindent{\small \texttt{SuFunctor.Universe.}\texttt{\bfseries not\_alwaysOn\_iff}}\hfill{\footnotesize \path{lean/su-functor/SuFunctor/Adequacy/SpecPathLens.lean}, line 422}\par\nopagebreak
\begin{leansrc}
theorem not_alwaysOn_iff {Good S : Set X} :
    ¬ alwaysOn Good S ↔ ∃ x ∈ S, x ∈ Goodᶜ
\end{leansrc}
\par\medskip\noindent{\small \texttt{SuFunctor.Universe.}\texttt{\bfseries satCredal\_lower\_singletonC\_zero}}\hfill{\footnotesize \path{lean/su-functor/SuFunctor/Adequacy/SpecBehavioral.lean}, line 143}\par\nopagebreak
\begin{leansrc}
theorem satCredal_lower_singletonC_zero (c : SubPMF V → Prop) :
    SatCredal .lower (singletonC (SubPMF.zero V)) c
\end{leansrc}
\par\medskip\noindent{\small \texttt{SuFunctor.Universe.}\texttt{\bfseries probThresholdSet}}\hfill{\footnotesize \path{lean/su-functor/SuFunctor/Adequacy/SpecBehavioral.lean}, line 168}\par\nopagebreak
\begin{leansrc}
def probThresholdSet (K : Kernel S S) [IsMarkovKernel K] (p : ℝ)
    (E : Set (ℕ → S)) : Set S
\end{leansrc}
\par\medskip\noindent{\small \texttt{SuFunctor.Universe.}\texttt{\bfseries safetyAtom\_modelClass\_mono}}\hfill{\footnotesize \path{lean/su-functor/SuFunctor/Adequacy/SpecBehavioral.lean}, line 92}\par\nopagebreak
\begin{leansrc}
theorem safetyAtom_modelClass_mono {C D : Set X} (h : C ⊆ D) :
    (safetyAtom C).modelClass ⊆ (safetyAtom D).modelClass
\end{leansrc}
\par\medskip\noindent{\small \texttt{SuFunctor.Universe.}\texttt{\bfseries sprobAtom\_modelClass\_mono}}\hfill{\footnotesize \path{lean/su-functor/SuFunctor/Adequacy/SpecBehavioral.lean}, line 210}\par\nopagebreak
\begin{leansrc}
theorem sprobAtom_modelClass_mono {p : ℝ} {E E' : Set (ℕ → S)}
    (hE : E ⊆ E') :
    (safetyAtom (probThresholdSet K p E)).modelClass ⊆
      (safetyAtom (probThresholdSet K p E')).modelClass
\end{leansrc}
\par\medskip\noindent{\small \texttt{SuFunctor.Universe.}\texttt{\bfseries sprob\_sound\_of\_streettBundle}}\hfill{\footnotesize \path{lean/su-functor/SuFunctor/Adequacy/SpecBehavioral.lean}, line 224}\par\nopagebreak
\begin{leansrc}
theorem sprob_sound_of_streettBundle
    {K : Kernel S S} [IsMarkovKernel K]
    {D Dbar U : Set S} {V₀ : S → ℝ} {p thr : ℝ} {x : S}
    {ι : Type*} [Countable ι]
    {Rq G : ι → Set S} {W : ι → S → ℝ} {ε C : ι → ℝ}
    (hD : MeasurableSet D) (hDbar : MeasurableSet Dbar)
    (hU : MeasurableSet U) (hV₀ : Measurable V₀)
    (hRq : ∀ i, MeasurableSet (Rq i)) (hG : ∀ i, MeasurableSet (G i))
    (hW : ∀ i, Measurable (W i))
    (hΘ : p + thr ≤ 1)
    (hDsub : D ⊆ Dbar)
    (hR3 : ∀ y ∈ D, y ∉ U → K y Dbar = 1)
    (hR4 : ∀ y ∈ D, y ∉ U → K y (D ∪ U) = 1)
    (hV₀int : ∀ y ∈ D, y ∉ U → Integrable V₀ (K y))
    (hV₀drift : ∀ y ∈ D, y ∉ U → ∫ z, V₀ z ∂K y ≤ V₀ y)
    (hV₀floor : ∀ y ∈ U, y ∈ Dbar → 1 ≤ V₀ y)
    (hV₀nonneg : ∀ y ∈ Dbar, 0 ≤ V₀ y)
    (hε : ∀ i, 0 < ε i)
    (hW0 : ∀ i, ∀ y ∈ Dbar, 0 ≤ W i y)
    (hWC : ∀ i, ∀ y ∈ Dbar, W i y ≤ C i)
    (hWint : ∀ i, ∀ y ∈ D, y ∉ U → y ∉ G i → Integrable (W i) (K y))
    (hWdrift : ∀ i, ∀ y ∈ D, y ∉ U → y ∉ G i →
      ∫ z, W i z ∂K y ≤ W i y)
    (hWdriftR : ∀ i, ∀ y ∈ D, y ∉ U → y ∈ Rq i → y ∉ G i →
      ∫ z, W i z ∂K y ≤ W i y - ε i)
    (hx : x ∈ D) (hthr : V₀ x ≤ thr)
    {E : Set (ℕ → S)}
    (hE : {ω | ∀ i, (∃ᶠ n in atTop, ω n ∈ Rq i) →
      ∃ᶠ n in atTop, ω n ∈ G i} ⊆ E) :
    x ∈ probThresholdSet K p E
\end{leansrc}
\par\medskip\noindent{\small \texttt{SuFunctor.Universe.}\texttt{\bfseries satCredal\_lower\_singletonC\_geC}}\hfill{\footnotesize \path{lean/su-functor/SuFunctor/Adequacy/SpecBehavioral.lean}, line 115}\par\nopagebreak
\begin{leansrc}
theorem satCredal_lower_singletonC_geC {μ : SubPMF V}
    (hμ : μ ≠ SubPMF.zero V) (A : Set V) (q : ℚ) :
    SatCredal .lower (singletonC μ) (geC A q) ↔ q ≤ μ.prob A
\end{leansrc}
\par\medskip\noindent{\small \texttt{SuFunctor.Universe.}\texttt{\bfseries returnGESet}}\hfill{\footnotesize \path{lean/su-functor/SuFunctor/Adequacy/SpecBehavioral.lean}, line 284}\par\nopagebreak
\begin{leansrc}
def returnGESet (K : Kernel S S) [IsMarkovKernel K] (r : S → ℝ)
    (b : ℝ) (n : ℕ) : Set S
\end{leansrc}
\par\medskip\noindent{\small \texttt{SuFunctor.Universe.}\texttt{\bfseries returnGESet\_mono\_reward}}\hfill{\footnotesize \path{lean/su-functor/SuFunctor/Adequacy/SpecBehavioral.lean}, line 299}\par\nopagebreak
\begin{leansrc}
theorem returnGESet_mono_reward {r r' : S → ℝ} {b : ℝ} {n : ℕ}
    (hr : ∀ s, r s ≤ r' s)
    (hint : ∀ x, Integrable (finReturn r n) (pathMeasure K x))
    (hint' : ∀ x, Integrable (finReturn r' n) (pathMeasure K x)) :
    returnGESet K r b n ⊆ returnGESet K r' b n
\end{leansrc}
\par\medskip\noindent{\small \texttt{SuFunctor.Universe.}\texttt{\bfseries returnGEAtom\_modelClass\_mono\_reward}}\hfill{\footnotesize \path{lean/su-functor/SuFunctor/Adequacy/SpecBehavioral.lean}, line 334}\par\nopagebreak
\begin{leansrc}
theorem returnGEAtom_modelClass_mono_reward {r r' : S → ℝ} {b : ℝ}
    {n : ℕ} (hr : ∀ s, r s ≤ r' s)
    (hint : ∀ x, Integrable (finReturn r n) (pathMeasure K x))
    (hint' : ∀ x, Integrable (finReturn r' n) (pathMeasure K x)) :
    (safetyAtom (returnGESet K r b n)).modelClass ⊆
      (safetyAtom (returnGESet K r' b n)).modelClass
\end{leansrc}
\par\medskip\noindent{\small \texttt{SuFunctor.Universe.}\texttt{\bfseries returnGE\_sound\_of\_deadline}}\hfill{\footnotesize \path{lean/su-functor/SuFunctor/Adequacy/SpecBehavioral.lean}, line 441}\par\nopagebreak
\begin{leansrc}
theorem returnGE_sound_of_deadline [IsProbabilityMeasure μ]
    (hβ : 0 ≤ β) (ha : 0 < a)
    (hint : ∀ n, Integrable (L n) μ) (hnonneg : 0 ≤ L)
    (hmB : ∀ n, MeasurableSet[𝒢 n] {ω | X n ω ∈ B})
    (hdrift : ∀ n, ∀ᵐ ω ∂μ, X n ω ∉ B →
      (μ[L (n + 1) | 𝒢 n]) ω ≤ β * L n ω)
    (hbarrier : ∀ n ω, X n ω ∉ B → a ≤ L n ω) (n : ℕ) :
    1 - β ^ n * (∫ ω, L 0 ω ∂μ) / a ≤
      ∫ ω, finReturn (B.indicator fun _ => (1 : ℝ)) (n + 1)
        (fun t => X t ω) ∂μ
\end{leansrc}
\par\medskip\noindent{\small \texttt{SuFunctor.Universe.}\texttt{\bfseries discountedReturn\_le\_of\_drift}}\hfill{\footnotesize \path{lean/su-functor/SuFunctor/Adequacy/SpecBehavioral.lean}, line 543}\par\nopagebreak
\begin{leansrc}
theorem discountedReturn_le_of_drift [IsFiniteMeasure μ]
    {γ : ℝ} {r : ℕ → Ω → ℝ}
    (hβ : 0 ≤ β) (hγ : 0 ≤ γ)
    (hint : ∀ n, Integrable (L n) μ) (hnonneg : 0 ≤ L)
    (hdrift : ∀ n, ∀ᵐ ω ∂μ, (μ[L (n + 1) | 𝒢 n]) ω ≤ β * L n ω)
    (hrint : ∀ n, Integrable (r n) μ)
    (hr : ∀ n ω, r n ω ≤ L n ω) (n : ℕ) :
    ∫ ω, discountedPartialReturn γ r n ω ∂μ ≤
      (∑ t ∈ Finset.range n, (γ * β) ^ t) * ∫ ω, L 0 ω ∂μ
\end{leansrc}
\par\medskip\noindent{\small \texttt{SuFunctor.Universe.}\texttt{\bfseries integral\_le\_pow\_mul}}\hfill{\footnotesize \path{lean/su-functor/SuFunctor/Adequacy/SpecBehavioral.lean}, line 511}\par\nopagebreak
\begin{leansrc}
theorem integral_le_pow_mul [IsFiniteMeasure μ] (hβ : 0 ≤ β)
    (hint : ∀ n, Integrable (L n) μ) (hnonneg : 0 ≤ L)
    (hdrift : ∀ n, ∀ᵐ ω ∂μ, (μ[L (n + 1) | 𝒢 n]) ω ≤ β * L n ω)
    (n : ℕ) :
    ∫ ω, L n ω ∂μ ≤ β ^ n * ∫ ω, L 0 ω ∂μ
\end{leansrc}
\par\medskip\noindent{\small \texttt{SuFunctor.}\texttt{\bfseries livenessTilt\_top}}\hfill{\footnotesize \path{lean/su-functor-legacy/SuFunctorLegacy/Threshold.lean}, line 275}\par\nopagebreak
\begin{leansrc}
theorem livenessTilt_top (goal : Set V) (p : ℚ) :
    livenessTilt goal p ⊤ = livenessWall goal p
\end{leansrc}
\par\medskip\noindent{\small \texttt{SuFunctor.Universe.}\texttt{\bfseries untilOn\_univ\_safe}}\hfill{\footnotesize \path{lean/su-functor/SuFunctor/Adequacy/SpecPathLens.lean}, line 392}\par\nopagebreak
\begin{leansrc}
theorem untilOn_univ_safe {τ : X → ℕ} {n : ℕ} {Target S : Set X} :
    untilOn τ n Set.univ Target S ↔ windowReachOn τ n Target S
\end{leansrc}
\par\medskip\noindent{\small \texttt{SuFunctor.Universe.}\texttt{\bfseries streettOn\_singleton\_univ\_iff}}\hfill{\footnotesize \path{lean/su-functor/SuFunctor/Universe/SpecRewrite.lean}, line 599}\par\nopagebreak
\begin{leansrc}
theorem streettOn_singleton_univ_iff {Y : Type*} {G S : Set Y} :
    streettOn [(Set.univ, G)] S ↔ (S.Nonempty → (S ∩ G).Nonempty)
\end{leansrc}
\par\medskip\noindent{\small \texttt{SuFunctor.Universe.}\texttt{\bfseries streettOn\_singleton\_univ\_empty}}\hfill{\footnotesize \path{lean/su-functor/SuFunctor/Universe/SpecRewrite.lean}, line 625}\par\nopagebreak
\begin{leansrc}
theorem streettOn_singleton_univ_empty {Y : Type*} {G : Set Y} :
    streettOn [(Set.univ, G)] (∅ : Set Y) ∧ ¬ ((∅ : Set Y) ∩ G).Nonempty
\end{leansrc}
\par\medskip\noindent{\small \texttt{Qtsl.TemporalRewrite.}\texttt{\bfseries no\_f\_and\_fusion}}\hfill{\footnotesize \path{lean/qtsl-aux/QtslAux/TemporalRewrite.lean}, line 274}\par\nopagebreak
\begin{leansrc}
theorem no_f_and_fusion :
    ∃ V : Nat → Nat → Prop,
      sat V 0 (.and (.F (.atom 0)) (.F (.atom 1))) ∧
      ¬ sat V 0 (.F (.and (.atom 0) (.atom 1)))
\end{leansrc}
\par\medskip\noindent{\small \texttt{Qtsl.TemporalRewrite.}\texttt{\bfseries no\_g\_or\_fusion}}\hfill{\footnotesize \path{lean/qtsl-aux/QtslAux/TemporalRewrite.lean}, line 285}\par\nopagebreak
\begin{leansrc}
theorem no_g_or_fusion :
    ∃ V : Nat → Nat → Prop,
      sat V 0 (.G (.or (.atom 0) (.atom 1))) ∧
      ¬ sat V 0 (.or (.G (.atom 0)) (.G (.atom 1)))
\end{leansrc}
\par\medskip\noindent{\small \texttt{Qtsl.TemporalRewrite.}\texttt{\bfseries no\_gf\_and\_fusion}}\hfill{\footnotesize \path{lean/qtsl-aux/QtslAux/TemporalRewrite.lean}, line 298}\par\nopagebreak
\begin{leansrc}
theorem no_gf_and_fusion :
    ∃ V : Nat → Nat → Prop,
      sat V 0 (.and (.G (.F (.atom 0))) (.G (.F (.atom 1)))) ∧
      ¬ sat V 0 (.G (.F (.and (.atom 0) (.atom 1))))
\end{leansrc}
\par\medskip\noindent{\small \texttt{Qtsl.TemporalRewrite.}\texttt{\bfseries no\_fg\_or\_fusion}}\hfill{\footnotesize \path{lean/qtsl-aux/QtslAux/TemporalRewrite.lean}, line 311}\par\nopagebreak
\begin{leansrc}
theorem no_fg_or_fusion :
    ∃ V : Nat → Nat → Prop,
      sat V 0 (.F (.G (.or (.atom 0) (.atom 1)))) ∧
      ¬ sat V 0 (.or (.F (.G (.atom 0))) (.F (.G (.atom 1))))
\end{leansrc}
\par\medskip\noindent{\small \texttt{SuFunctor.Universe.}\texttt{\bfseries exists\_denEq\_not\_classEq}}\hfill{\footnotesize \path{lean/su-functor/SuFunctor/Universe/SpecRewrite.lean}, line 206}\par\nopagebreak
\begin{leansrc}
theorem exists_denEq_not_classEq :
    ∃ (A : Fin 3 → Set Bool) (φ ψ : StateFormula (Fin 3)),
      DenEq A φ ψ ∧ ¬ ClassEq A φ ψ
\end{leansrc}
\par\medskip\noindent{\small \texttt{SuFunctor.Universe.}\texttt{\bfseries classEquiv\_of\_nnf\_eq}}\hfill{\footnotesize \path{lean/su-functor/SuFunctor/Universe/SpecRewrite.lean}, line 139}\par\nopagebreak
\begin{leansrc}
theorem classEquiv_of_nnf_eq {φ ψ : StateFormula α} (h : φ.nnf = ψ.nnf) :
    ClassEquiv.{u} φ ψ
\end{leansrc}
\par\medskip\noindent Modules cited: \texttt{SuFunctor.Universe.SpecElab}, \texttt{SuFunctor.Universe.SpecRewrite}, \texttt{SuFunctor.Universe.PairSequent}.
\par\medskip\noindent Also cited in this chapter, listed under an earlier one: \texttt{General.iSup\_gambleRows\_eq\_self}, \texttt{iSup\_boxRows\_eq\_self}, \texttt{sat\_sup\_iff}, \texttt{top\_notMem\_modelClass}.

\section{Rule Catalogue}
Citations from Chapter~\ref{ch:qtsl-manual}.

\par\medskip\noindent{\small \texttt{PPLC.}\texttt{\bfseries PAxData}}\hfill{\footnotesize \path{lean/qtsl/Qtsl/Core/Catalogue.lean}, line 1353}\par\nopagebreak
\begin{leansrc}
inductive PAxData : Hypernet → Sequent → Type where …
\end{leansrc}
\par\medskip\noindent{\small \texttt{SuFunctor.StateSpace.FamObj.}\texttt{\bfseries famEq}}\hfill{\footnotesize \path{lean/su-functor/SuFunctor/StateSpace/FamLimits.lean}, line 169}\par\nopagebreak
\begin{leansrc}
def famEq : FamObj C
\end{leansrc}
\par\medskip\noindent{\small \texttt{PPLC.}\texttt{\bfseries sound\_ax}}\hfill{\footnotesize \path{lean/qtsl/Qtsl/Soundness.lean}, line 4165}\par\nopagebreak
\begin{leansrc}
theorem sound_ax {H : Hypernet} {seq : Sequent}
    (h : seq.pre = seq.post) : seq.semanticallyValid H
\end{leansrc}
\par\medskip\noindent{\small \texttt{PPLC.Formula.}\texttt{\bfseries holds}}\hfill{\footnotesize \path{lean/qtsl/Qtsl/Catalogue.lean}, line 402}\par\nopagebreak
\begin{leansrc}
def Formula.holds {H : Hypernet} (σ : H.Valuation) : Formula → Prop
  | .atom b l r => ∃ a a' : QTensor,
      Term.evalQ σ l = some a ∧ Term.evalQ σ r = some a' ∧
      QTensor.rel2 (BowTie.holdsQ b) a a'
  | .inRegion t simplex S =>
      simplex = S.simplex ∧
      ∃ q : QTensor, Term.evalQ σ t = some q ∧
        q.shape = [S.simplex.complex.ambDim] ∧
        MaskRegion S.simplex.complex S.simplex.contains S.contains
          (fun j => q.entries.getD j.val 0)
  | .and l r => Formula.holds σ l ∧ Formula.holds σ r
  | .or l r => Formula.holds σ l ∨ Formula.holds σ r
  | .tt => True
  | .ff => False

namespace ContractPattern

variable {inputRanks : List Nat} {outputRank : Nat}

def idxShape (pat : ContractPattern inputRanks outputRank) : Shape
\end{leansrc}
\par\medskip\noindent{\small \texttt{PPLC.Sequent.}\texttt{\bfseries semanticallyValid}}\hfill{\footnotesize \path{lean/qtsl/Qtsl/Catalogue.lean}, line 527}\par\nopagebreak
\begin{leansrc}
def Sequent.semanticallyValid (H : Hypernet) (seq : Sequent) : Prop
\end{leansrc}
\par\medskip\noindent{\small \texttt{PPLC.M3.}\texttt{\bfseries le\_refl}}\hfill{\footnotesize \path{lean/qtsl/Qtsl/Admissible.lean}, line 112}\par\nopagebreak
\begin{leansrc}
theorem le_refl : ∀ x : M3, le x x = true
\end{leansrc}
\par\medskip\noindent{\small \texttt{PPLC.}\texttt{\bfseries PWeakenData}}\hfill{\footnotesize \path{lean/qtsl/Qtsl/Core/Catalogue.lean}, line 1357}\par\nopagebreak
\begin{leansrc}
inductive PWeakenData : Hypernet → Sequent → Type where …
\end{leansrc}
\par\medskip\noindent{\small \texttt{PPLC.}\texttt{\bfseries ProofObject}}\hfill{\footnotesize \path{lean/qtsl/Qtsl/Core/Catalogue.lean}, line 1316}\par\nopagebreak
\begin{leansrc}
inductive ProofObject : Hypernet → Sequent → Type where …
\end{leansrc}
\par\medskip\noindent{\small \texttt{PPLC.}\texttt{\bfseries sound\_weaken}}\hfill{\footnotesize \path{lean/qtsl/Qtsl/Soundness.lean}, line 4171}\par\nopagebreak
\begin{leansrc}
theorem sound_weaken {H : Hypernet} {seq premise : Sequent} {extra : Formula}
    (ih : premise.semanticallyValid H)
    (hpre : seq.pre = .and premise.pre extra)
    (hpost : seq.post = premise.post) : seq.semanticallyValid H
\end{leansrc}
\par\medskip\noindent{\small \texttt{PPLC.M3.}\texttt{\bfseries weaken\_kit}}\hfill{\footnotesize \path{lean/qtsl/Qtsl/Admissible.lean}, line 121}\par\nopagebreak
\begin{leansrc}
theorem weaken_kit : ∀ x y z : M3,
    le x z = true → le (inf x y) z = true
\end{leansrc}
\par\medskip\noindent{\small \texttt{PPLC.}\texttt{\bfseries sound\_cut}}\hfill{\footnotesize \path{lean/qtsl/Qtsl/Soundness.lean}, line 4180}\par\nopagebreak
\begin{leansrc}
theorem sound_cut {H : Hypernet} {seq leftPremise rightPremise : Sequent}
    {cutFormula : Formula}
    (ihl : leftPremise.semanticallyValid H)
    (ihr : rightPremise.semanticallyValid H)
    (hpre : seq.pre = leftPremise.pre)
    (hpost : seq.post = rightPremise.post)
    (hcutL : leftPremise.post = cutFormula)
    (hcutR : rightPremise.pre = .and leftPremise.pre cutFormula) :
    seq.semanticallyValid H
\end{leansrc}
\par\medskip\noindent{\small \texttt{PPLC.}\texttt{\bfseries PCutData}}\hfill{\footnotesize \path{lean/qtsl/Qtsl/Core/Catalogue.lean}, line 1365}\par\nopagebreak
\begin{leansrc}
inductive PCutData : Hypernet → Sequent → Type where …
\end{leansrc}
\par\medskip\noindent{\small \texttt{PPLC.M3.}\texttt{\bfseries cut\_kit}}\hfill{\footnotesize \path{lean/qtsl/Qtsl/Admissible.lean}, line 125}\par\nopagebreak
\begin{leansrc}
theorem cut_kit : ∀ x m y : M3,
    le x m = true → le (inf x m) y = true → le x y = true
\end{leansrc}
\par\medskip\noindent{\small \texttt{PPLC.}\texttt{\bfseries sound\_andI}}\hfill{\footnotesize \path{lean/qtsl/Qtsl/Soundness.lean}, line 4198}\par\nopagebreak
\begin{leansrc}
theorem sound_andI {H : Hypernet} {seq leftPremise rightPremise : Sequent}
    (ihl : leftPremise.semanticallyValid H)
    (ihr : rightPremise.semanticallyValid H)
    (hpreMatch : leftPremise.pre = rightPremise.pre)
    (hpre : seq.pre = leftPremise.pre)
    (hpost : seq.post = .and leftPremise.post rightPremise.post) :
    seq.semanticallyValid H
\end{leansrc}
\par\medskip\noindent{\small \texttt{PPLC.}\texttt{\bfseries PAndIData}}\hfill{\footnotesize \path{lean/qtsl/Qtsl/Core/Catalogue.lean}, line 1376}\par\nopagebreak
\begin{leansrc}
inductive PAndIData : Hypernet → Sequent → Type where …
\end{leansrc}
\par\medskip\noindent{\small \texttt{PPLC.M3.}\texttt{\bfseries andI\_kit}}\hfill{\footnotesize \path{lean/qtsl/Qtsl/Admissible.lean}, line 129}\par\nopagebreak
\begin{leansrc}
theorem andI_kit : ∀ x y z : M3,
    le x y = true → le x z = true → le x (inf y z) = true
\end{leansrc}
\par\medskip\noindent{\small \texttt{PPLC.}\texttt{\bfseries sound\_andE\_left}}\hfill{\footnotesize \path{lean/qtsl/Qtsl/Soundness.lean}, line 4212}\par\nopagebreak
\begin{leansrc}
theorem sound_andE_left {H : Hypernet} {seq premise : Sequent} {other : Formula}
    (ih : premise.semanticallyValid H)
    (hpre : seq.pre = premise.pre)
    (hpost : premise.post = .and seq.post other) :
    seq.semanticallyValid H
\end{leansrc}
\par\medskip\noindent{\small \texttt{PPLC.}\texttt{\bfseries sound\_andE\_right}}\hfill{\footnotesize \path{lean/qtsl/Qtsl/Soundness.lean}, line 4223}\par\nopagebreak
\begin{leansrc}
theorem sound_andE_right {H : Hypernet} {seq premise : Sequent} {other : Formula}
    (ih : premise.semanticallyValid H)
    (hpre : seq.pre = premise.pre)
    (hpost : premise.post = .and other seq.post) :
    seq.semanticallyValid H
\end{leansrc}
\par\medskip\noindent{\small \texttt{PPLC.M3.}\texttt{\bfseries andE\_l\_kit}}\hfill{\footnotesize \path{lean/qtsl/Qtsl/Admissible.lean}, line 133}\par\nopagebreak
\begin{leansrc}
theorem andE_l_kit : ∀ x y z : M3,
    le x (inf y z) = true → le x y = true
\end{leansrc}
\par\medskip\noindent{\small \texttt{PPLC.M3.}\texttt{\bfseries andE\_r\_kit}}\hfill{\footnotesize \path{lean/qtsl/Qtsl/Admissible.lean}, line 137}\par\nopagebreak
\begin{leansrc}
theorem andE_r_kit : ∀ x y z : M3,
    le x (inf y z) = true → le x z = true
\end{leansrc}
\par\medskip\noindent{\small \texttt{PPLC.}\texttt{\bfseries sound\_orI\_left}}\hfill{\footnotesize \path{lean/qtsl/Qtsl/Soundness.lean}, line 4234}\par\nopagebreak
\begin{leansrc}
theorem sound_orI_left {H : Hypernet} {seq premise : Sequent} {other : Formula}
    (ih : premise.semanticallyValid H)
    (hpre : seq.pre = premise.pre)
    (hpost : seq.post = .or premise.post other) :
    seq.semanticallyValid H
\end{leansrc}
\par\medskip\noindent{\small \texttt{PPLC.}\texttt{\bfseries sound\_orI\_right}}\hfill{\footnotesize \path{lean/qtsl/Qtsl/Soundness.lean}, line 4244}\par\nopagebreak
\begin{leansrc}
theorem sound_orI_right {H : Hypernet} {seq premise : Sequent} {other : Formula}
    (ih : premise.semanticallyValid H)
    (hpre : seq.pre = premise.pre)
    (hpost : seq.post = .or other premise.post) :
    seq.semanticallyValid H
\end{leansrc}
\par\medskip\noindent{\small \texttt{PPLC.M3.}\texttt{\bfseries orI\_l\_kit}}\hfill{\footnotesize \path{lean/qtsl/Qtsl/Admissible.lean}, line 141}\par\nopagebreak
\begin{leansrc}
theorem orI_l_kit : ∀ x y z : M3,
    le x y = true → le x (sup y z) = true
\end{leansrc}
\par\medskip\noindent{\small \texttt{PPLC.M3.}\texttt{\bfseries orI\_r\_kit}}\hfill{\footnotesize \path{lean/qtsl/Qtsl/Admissible.lean}, line 145}\par\nopagebreak
\begin{leansrc}
theorem orI_r_kit : ∀ x y z : M3,
    le x z = true → le x (sup y z) = true
\end{leansrc}
\par\medskip\noindent{\small \texttt{PPLC.}\texttt{\bfseries sound\_orE}}\hfill{\footnotesize \path{lean/qtsl/Qtsl/Soundness.lean}, line 4254}\par\nopagebreak
\begin{leansrc}
theorem sound_orE {H : Hypernet} {seq leftPremise rightPremise : Sequent}
    (ihl : leftPremise.semanticallyValid H)
    (ihr : rightPremise.semanticallyValid H)
    (hpostMatch : leftPremise.post = rightPremise.post)
    (hpre : seq.pre = .or leftPremise.pre rightPremise.pre)
    (hpost : seq.post = leftPremise.post) :
    seq.semanticallyValid H
\end{leansrc}
\par\medskip\noindent{\small \texttt{PPLC.}\texttt{\bfseries OrECtxStep}}\hfill{\footnotesize \path{lean/qtsl/Qtsl/Admissible.lean}, line 20}\par\nopagebreak
\begin{leansrc}
structure OrECtxStep where …
\end{leansrc}
\par\medskip\noindent{\small \texttt{PPLC.}\texttt{\bfseries sound\_orECtx}}\hfill{\footnotesize \path{lean/qtsl/Qtsl/Admissible.lean}, line 33}\par\nopagebreak
\begin{leansrc}
theorem sound_orECtx {H : Hypernet} (st : OrECtxStep)
    (ihl : st.leftPremise.semanticallyValid H)
    (ihr : st.rightPremise.semanticallyValid H) :
    st.seq.semanticallyValid H
\end{leansrc}
\par\medskip\noindent{\small \texttt{PPLC.}\texttt{\bfseries cx\_not\_derivable}}\hfill{\footnotesize \path{lean/qtsl/Qtsl/Admissible.lean}, line 657}\par\nopagebreak
\begin{leansrc}
theorem cx_not_derivable : IsEmpty (ProofObject cxH cxSeq)
\end{leansrc}
\par\medskip\noindent{\small \texttt{PPLC.}\texttt{\bfseries orECtx\_not\_elaborable}}\hfill{\footnotesize \path{lean/qtsl/Qtsl/Admissible.lean}, line 666}\par\nopagebreak
\begin{leansrc}
theorem orECtx_not_elaborable :
    ¬ (∀ (st : OrECtxStep) (H : Hypernet),
        Nonempty (ProofObject H st.leftPremise) →
        Nonempty (ProofObject H st.rightPremise) →
        Nonempty (ProofObject H st.seq))
\end{leansrc}
\par\medskip\noindent{\small \texttt{PPLC.ProofObject.}\texttt{\bfseries m3sound}}\hfill{\footnotesize \path{lean/qtsl/Qtsl/Admissible.lean}, line 406}\par\nopagebreak
\begin{leansrc}
theorem ProofObject.m3sound : ∀ {H : Hypernet} {seq : Sequent},
    (∀ w, H.wireType w = none) → ProofObject H seq →
    M3.le (mval seq.pre) (mval seq.post) = true
  | _, _, _, .mk _ _ _ (.ax _ _ (.mk _ _ hpp)) => by
      rw [hpp]; exact M3.le_refl _
  | _, _, hE, .mk _ _ _ (.weaken _ _ (.mk _ _ _ pproof _ hpre hpost)) => by
      rw [hpre, hpost]
      exact M3.weaken_kit _ _ _ (pproof.m3sound hE)
  | _, _, hE, .mk _ _ _ (.cut _ _
      (.mk _ _ _ _ lproof _ rproof hpre hpost hcl hcr)) => by
      have hl
\end{leansrc}
\par\medskip\noindent{\small \texttt{PPLC.}\texttt{\bfseries cxStep}}\hfill{\footnotesize \path{lean/qtsl/Qtsl/Admissible.lean}, line 632}\par\nopagebreak
\begin{leansrc}
def cxStep : OrECtxStep where …
\end{leansrc}
\par\medskip\noindent{\small \texttt{PPLC.OrECtxStep.}\texttt{\bfseries elabHere}}\hfill{\footnotesize \path{lean/qtsl/Qtsl/Admissible.lean}, line 58}\par\nopagebreak
\begin{leansrc}
def OrECtxStep.elabHere {H : Hypernet} (st : OrECtxStep)
    (hpath : st.path = .here)
    (lpf : ProofObject H st.leftPremise) (rpf : ProofObject H st.rightPremise) :
    ProofObject H st.seq
\end{leansrc}
\par\medskip\noindent{\small \texttt{PPLC.}\texttt{\bfseries POrEData}}\hfill{\footnotesize \path{lean/qtsl/Qtsl/Core/Catalogue.lean}, line 1413}\par\nopagebreak
\begin{leansrc}
inductive POrEData : Hypernet → Sequent → Type where …
\end{leansrc}
\par\medskip\noindent{\small \texttt{PPLC.M3.}\texttt{\bfseries orE\_kit}}\hfill{\footnotesize \path{lean/qtsl/Qtsl/Admissible.lean}, line 149}\par\nopagebreak
\begin{leansrc}
theorem orE_kit : ∀ x y z : M3,
    le x z = true → le y z = true → le (sup x y) z = true
\end{leansrc}
\par\medskip\noindent{\small \texttt{PPLC.}\texttt{\bfseries substituteInFormula}}\hfill{\footnotesize \path{lean/qtsl/Qtsl/Core/Catalogue.lean}, line 581}\par\nopagebreak
\begin{leansrc}
def substituteInFormula (F : Formula) (w : Wire) (t : Term) (F' : Formula) : Prop
\end{leansrc}
\par\medskip\noindent{\small \texttt{PPLC.}\texttt{\bfseries sound\_eqElim}}\hfill{\footnotesize \path{lean/qtsl/Qtsl/Soundness.lean}, line 4269}\par\nopagebreak
\begin{leansrc}
theorem sound_eqElim {H : Hypernet} {seq eqPremise bodyPremise : Sequent}
    {wire : Wire} {termT termU : Term} {bodyForm : Formula}
    (ihEq : eqPremise.semanticallyValid H)
    (ihBody : bodyPremise.semanticallyValid H)
    (hpre : seq.pre = eqPremise.pre)
    (hpreMatch : eqPremise.pre = bodyPremise.pre)
    (heqPost : eqPremise.post = .atom .eq termT termU)
    (htyEq : ∃ s, termT.bshape H = some s ∧ termU.bshape H = some s)
    (hbodySubT : substituteInFormula bodyForm wire termT bodyPremise.post)
    (hconclSubU : substituteInFormula bodyForm wire termU seq.post) :
    seq.semanticallyValid H
\end{leansrc}
\par\medskip\noindent{\small \texttt{PPLC.Term.}\texttt{\bfseries evalQ\_subst\_congr}}\hfill{\footnotesize \path{lean/qtsl/Qtsl/Soundness.lean}, line 270}\par\nopagebreak
\begin{leansrc}
theorem Term.evalQ_subst_congr {H : Hypernet} (σ : H.Valuation)
    {w : Wire} {u₁ u₂ : Term} (hu : Term.evalQ σ u₁ = Term.evalQ σ u₂) :
    ∀ (t : Term),
      Term.evalQ σ (Term.subst w u₁ t) = Term.evalQ σ (Term.subst w u₂ t)
\end{leansrc}
\par\medskip\noindent{\small \texttt{PPLC.Formula.}\texttt{\bfseries holds\_subst\_congr}}\hfill{\footnotesize \path{lean/qtsl/Qtsl/Soundness.lean}, line 288}\par\nopagebreak
\begin{leansrc}
theorem Formula.holds_subst_congr {H : Hypernet} (σ : H.Valuation)
    {w : Wire} {u₁ u₂ : Term} (hu : Term.evalQ σ u₁ = Term.evalQ σ u₂) :
    ∀ (F : Formula), (Formula.holds σ (F.subst w u₁) ↔ Formula.holds σ (F.subst w
      u₂))
\end{leansrc}
\par\medskip\noindent{\small \texttt{PPLC.Term.}\texttt{\bfseries evalQ\_bshape}}\hfill{\footnotesize \path{lean/qtsl/Qtsl/Soundness.lean}, line 479}\par\nopagebreak
\begin{leansrc}
theorem Term.evalQ_bshape {H : Hypernet} (σ : H.Valuation) :
    ∀ (t : Term) (s : Shape) (q : QTensor),
    t.bshape H = some s → Term.evalQ σ t = some q → q.shape = s
\end{leansrc}
\par\medskip\noindent{\small \texttt{PPLC.Hypernet.}\texttt{\bfseries Valuation}}\hfill{\footnotesize \path{lean/qtsl/Qtsl/Catalogue.lean}, line 268}\par\nopagebreak
\begin{leansrc}
def Hypernet.Valuation (H : Hypernet) : Type
\end{leansrc}
\par\medskip\noindent{\small \texttt{PPLC.Term.}\texttt{\bfseries evalQ}}\hfill{\footnotesize \path{lean/qtsl/Qtsl/Catalogue.lean}, line 352}\par\nopagebreak
\begin{leansrc}
def Term.evalQ {H : Hypernet} (σ : H.Valuation) : Term → Option QTensor
  | .var w => σ.wireQ w
  | .const c => some ⟨[], [c], rfl⟩
  | .add l r => do QTensor.map2 (· + ·) (← Term.evalQ σ l) (← Term.evalQ σ r)
  | .mul l r => do QTensor.map2 (· * ·) (← Term.evalQ σ l) (← Term.evalQ σ r)
  | .sub l r => do QTensor.map2 (· - ·) (← Term.evalQ σ l) (← Term.evalQ σ r)
  | .neg t => (Term.evalQ σ t).map (QTensor.map1 Neg.neg)
  | .idx t i => do

    let q ← Term.evalQ σ t
    match q.shape with
    | [e] => if i < e then some ⟨[], [q.entries.getD i 0], rfl⟩ else none
    | _ => none

def BowTie.holdsQ : BowTie → ℚ → ℚ → Prop
  | .lt => (· < ·)
  | .le => (· ≤ ·)
  | .eq => (· = ·)
  | .ge => fun a b => b ≤ a
  | .gt => fun a b => b < a
  | .ne => (· ≠ ·)

def finSum (n : ℕ) (f : Fin n → ℝ) : ℝ
\end{leansrc}
\par\medskip\noindent{\small \texttt{PPLC.QTensor.}\texttt{\bfseries rel2}}\hfill{\footnotesize \path{lean/qtsl/Qtsl/Catalogue.lean}, line 332}\par\nopagebreak
\begin{leansrc}
def rel2 (r : ℚ → ℚ → Prop) (a b : QTensor) : Prop
\end{leansrc}
\par\medskip\noindent{\small \texttt{PPLC.QTensor.}\texttt{\bfseries eq\_of\_rel2}}\hfill{\footnotesize \path{lean/qtsl/Qtsl/Soundness.lean}, line 2431}\par\nopagebreak
\begin{leansrc}
theorem QTensor.eq_of_rel2 {r : ℚ → ℚ → Prop} {a b : QTensor}
    (hr : ∀ x y, r x y → x = y) (hs : a.shape = b.shape)
    (h : QTensor.rel2 r a b) : a = b
\end{leansrc}
\par\medskip\noindent{\small \texttt{PPLC.Formula.}\texttt{\bfseries subst}}\hfill{\footnotesize \path{lean/qtsl/Qtsl/Core/Catalogue.lean}, line 573}\par\nopagebreak
\begin{leansrc}
def Formula.subst (w : Wire) (t : Term) : Formula → Formula
  | .atom b l r => .atom b (Term.subst w t l) (Term.subst w t r)
  | .inRegion u simplex S => .inRegion (Term.subst w t u) simplex S
  | .and l r => .and (Formula.subst w t l) (Formula.subst w t r)
  | .or l r => .or (Formula.subst w t l) (Formula.subst w t r)
  | .tt => .tt
  | .ff => .ff

def substituteInFormula (F : Formula) (w : Wire) (t : Term) (F' : Formula) : Prop
\end{leansrc}
\par\medskip\noindent{\small \texttt{PPLC.Term.}\texttt{\bfseries bshape}}\hfill{\footnotesize \path{lean/qtsl/Qtsl/Core/Catalogue.lean}, line 517}\par\nopagebreak
\begin{leansrc}
def Term.bshape (H : Hypernet) : Term → Option Shape
  | .var w => (H.wireType w).map (·.shape)
  | .const _ => some []
  | .add l r | .mul l r | .sub l r =>
    match l.bshape H, r.bshape H with
    | some s₁, some s₂ =>
      if s₁ = s₂ then some s₁
      else if s₁ = [] then some s₂
      else if s₂ = [] then some s₁
      else none
    | _, _ => none
  | .neg t => t.bshape H
  | .idx t i =>

    match t.bshape H with
    | some [e] => if i < e then some [] else none
    | _ => none

def Term.finWires (H : Hypernet) (t : Term) : Prop
\end{leansrc}
\par\medskip\noindent{\small \texttt{PPLC.Term.}\texttt{\bfseries finWires}}\hfill{\footnotesize \path{lean/qtsl/Qtsl/Core/Catalogue.lean}, line 535}\par\nopagebreak
\begin{leansrc}
def Term.finWires (H : Hypernet) (t : Term) : Prop
\end{leansrc}
\par\medskip\noindent{\small \texttt{PPLC.Term.}\texttt{\bfseries evalQ\_defined\_of\_bshape\_fin}}\hfill{\footnotesize \path{lean/qtsl/Qtsl/Soundness.lean}, line 3488}\par\nopagebreak
\begin{leansrc}
theorem Term.evalQ_defined_of_bshape_fin {H : Hypernet} (σ : H.Valuation) :
    ∀ (t : Term) (s : Shape), t.bshape H = some s → t.finWires H →
    ∃ q, Term.evalQ σ t = some q ∧ q.shape = s
\end{leansrc}
\par\medskip\noindent{\small \texttt{PPLC.QTensor.}\texttt{\bfseries rel2\_refl}}\hfill{\footnotesize \path{lean/qtsl/Qtsl/Soundness.lean}, line 6773}\par\nopagebreak
\begin{leansrc}
theorem QTensor.rel2_refl {r : ℚ → ℚ → Prop} (hr : ∀ x, r x x)
    (a : QTensor) : QTensor.rel2 r a a
\end{leansrc}
\par\medskip\noindent{\small \texttt{PPLC.}\texttt{\bfseries sound\_eqRefl}}\hfill{\footnotesize \path{lean/qtsl/Qtsl/Soundness.lean}, line 6853}\par\nopagebreak
\begin{leansrc}
theorem sound_eqRefl {H : Hypernet} {seq : Sequent} {t : Term}
    (htyDef : ∃ s, t.bshape H = some s) (hfin : t.finWires H)
    (hpost : seq.post = .atom .eq t t) :
    seq.semanticallyValid H
\end{leansrc}
\par\medskip\noindent{\small \texttt{PPLC.Hypernet.Valuation.}\texttt{\bfseries wireQ\_isSome}}\hfill{\footnotesize \path{lean/qtsl/Qtsl/Soundness.lean}, line 70}\par\nopagebreak
\begin{leansrc}
theorem Hypernet.Valuation.wireQ_isSome {H : Hypernet} (σ : H.Valuation)
    {w : Wire} {τ : TensorType} (hτ : H.wireType w = some τ)
    (hfin : τ.dtype.ext = Ext.fin) : ∃ q, σ.wireQ w = some q
\end{leansrc}
\par\medskip\noindent{\small \texttt{PPLC.QTensor.}\texttt{\bfseries rel2\_symm}}\hfill{\footnotesize \path{lean/qtsl/Qtsl/Soundness.lean}, line 6789}\par\nopagebreak
\begin{leansrc}
theorem QTensor.rel2_symm {r : ℚ → ℚ → Prop}
    (hr : ∀ x y, r x y → r y x) {a b : QTensor}
    (h : QTensor.rel2 r a b) : QTensor.rel2 r b a
\end{leansrc}
\par\medskip\noindent{\small \texttt{PPLC.}\texttt{\bfseries sound\_eqSym}}\hfill{\footnotesize \path{lean/qtsl/Qtsl/Soundness.lean}, line 6863}\par\nopagebreak
\begin{leansrc}
theorem sound_eqSym {H : Hypernet} {seq premise : Sequent} {tT tU : Term}
    (ih : premise.semanticallyValid H)
    (hpre : seq.pre = premise.pre)
    (hpremPost : premise.post = .atom .eq tT tU)
    (hpost : seq.post = .atom .eq tU tT) :
    seq.semanticallyValid H
\end{leansrc}
\par\medskip\noindent{\small \texttt{PPLC.}\texttt{\bfseries PEqSymData}}\hfill{\footnotesize \path{lean/qtsl/Qtsl/Core/Catalogue.lean}, line 1443}\par\nopagebreak
\begin{leansrc}
inductive PEqSymData : Hypernet → Sequent → Type where …
\end{leansrc}
\par\medskip\noindent{\small \texttt{PPLC.QTensor.}\texttt{\bfseries rel2\_conj}}\hfill{\footnotesize \path{lean/qtsl/Qtsl/Soundness.lean}, line 6834}\par\nopagebreak
\begin{leansrc}
theorem QTensor.rel2_conj {r s t : ℚ → ℚ → Prop}
    (hrst : ∀ x y, r x y → s x y → t x y) {a b : QTensor}
    (hr : QTensor.rel2 r a b) (hs : QTensor.rel2 s a b) :
    QTensor.rel2 t a b
\end{leansrc}
\par\medskip\noindent{\small \texttt{PPLC.}\texttt{\bfseries sound\_antisym}}\hfill{\footnotesize \path{lean/qtsl/Qtsl/Soundness.lean}, line 6888}\par\nopagebreak
\begin{leansrc}
theorem sound_antisym {H : Hypernet} {seq premiseLe premiseGe : Sequent}
    {tT tU : Term}
    (ihLe : premiseLe.semanticallyValid H)
    (ihGe : premiseGe.semanticallyValid H)
    (hpreLe : seq.pre = premiseLe.pre)
    (hpreGe : seq.pre = premiseGe.pre)
    (hpremLePost : premiseLe.post = .atom .le tT tU)
    (hpremGePost : premiseGe.post = .atom .ge tT tU)
    (hpost : seq.post = .atom .eq tT tU) :
    seq.semanticallyValid H
\end{leansrc}
\par\medskip\noindent{\small \texttt{PPLC.}\texttt{\bfseries baryCoeffs1D}}\hfill{\footnotesize \path{lean/qtsl/Qtsl/Core/Catalogue.lean}, line 584}\par\nopagebreak
\begin{leansrc}
def baryCoeffs1D (cx : SimpC) (dContains : Fin cx.nVertices → Bool) :
    Option (List (Fin cx.nVertices × Rat × Rat))
\end{leansrc}
\par\medskip\noindent{\small \texttt{PPLC.}\texttt{\bfseries primAffine}}\hfill{\footnotesize \path{lean/qtsl/Qtsl/Core/Catalogue.lean}, line 599}\par\nopagebreak
\begin{leansrc}
def primAffine (s c : Rat) : Rat × Rat
\end{leansrc}
\par\medskip\noindent{\small \texttt{PPLC.}\texttt{\bfseries f4VertexAtom}}\hfill{\footnotesize \path{lean/qtsl/Qtsl/Core/Catalogue.lean}, line 612}\par\nopagebreak
\begin{leansrc}
def f4VertexAtom (inS : Bool) (off slp : Rat) (t : Term) : Formula
\end{leansrc}
\par\medskip\noindent{\small \texttt{PPLC.}\texttt{\bfseries barycentricHalfspaceBuild}}\hfill{\footnotesize \path{lean/qtsl/Qtsl/Core/Catalogue.lean}, line 616}\par\nopagebreak
\begin{leansrc}
def barycentricHalfspaceBuild (t : Term) (simplex : Simplex)
    (S : ExtensionSet) : Option Formula
\end{leansrc}
\par\medskip\noindent{\small \texttt{PPLC.}\texttt{\bfseries barycentricHalfspaceForm}}\hfill{\footnotesize \path{lean/qtsl/Qtsl/Core/Catalogue.lean}, line 625}\par\nopagebreak
\begin{leansrc}
def barycentricHalfspaceForm (t : Term) (simplex : Simplex)
    (S : ExtensionSet) (F : Formula) : Prop
\end{leansrc}
\par\medskip\noindent{\small \texttt{PPLC.}\texttt{\bfseries barycentricHalfspaceBuild\_spec}}\hfill{\footnotesize \path{lean/qtsl/Qtsl/Soundness.lean}, line 634}\par\nopagebreak
\begin{leansrc}
theorem barycentricHalfspaceBuild_spec {t : Term} {simplex : Simplex}
    {S : ExtensionSet} {F : Formula}
    (h : barycentricHalfspaceBuild t simplex S = some F) :
    S.simplex = simplex ∧ simplex.complex.ambDim = 1 ∧
    ∃ (a b : Fin simplex.complex.nVertices) (V₀ V₁ : ℚ),
      (List.finRange simplex.complex.nVertices).filter
        (fun v => simplex.contains v) = [a, b] ∧
      (∀ j, simplex.complex.coord a j = V₀) ∧
      (∀ j, simplex.complex.coord b j = V₁) ∧
      V₀ ≠ V₁ ∧
      F = .and
        (f4VertexAtom (S.containsNat a.val)
          (V₁ / (V₁ - V₀)) (-1 / (V₁ - V₀)) t)
        (f4VertexAtom (S.containsNat b.val)
          (-V₀ / (V₁ - V₀)) (1 / (V₁ - V₀)) t)
\end{leansrc}
\par\medskip\noindent{\small \texttt{PPLC.}\texttt{\bfseries sound\_convElim}}\hfill{\footnotesize \path{lean/qtsl/Qtsl/Soundness.lean}, line 4672}\par\nopagebreak
\begin{leansrc}
theorem sound_convElim {H : Hypernet} {seq premise : Sequent}
    {term : Term} {simplex : Simplex} {S : ExtensionSet}
    (ih : premise.semanticallyValid H)
    (hpre : seq.pre = premise.pre)
    (hpremiseForm : premise.post = .inRegion term simplex S)
    (hhalf : barycentricHalfspaceForm term simplex S seq.post) :
    seq.semanticallyValid H
\end{leansrc}
\par\medskip\noindent{\small \texttt{PPLC.}\texttt{\bfseries f4VertexAtom\_holds\_iff}}\hfill{\footnotesize \path{lean/qtsl/Qtsl/Soundness.lean}, line 1393}\par\nopagebreak
\begin{leansrc}
theorem f4VertexAtom_holds_iff {H : Hypernet} {σ : H.Valuation} {t : Term}
    {q : QTensor} (hq : Term.evalQ σ t = some q) (hsq : SingletonShape q.shape)
    (inS : Bool) (off slp : ℚ) :
    Formula.holds σ (f4VertexAtom inS off slp t) ↔
      (if inS then off + slp * q.scalar ≤ 0
       else 0 ≤ off + slp * q.scalar)
\end{leansrc}
\par\medskip\noindent{\small \texttt{PPLC.}\texttt{\bfseries maskRegion\_pair\_elim}}\hfill{\footnotesize \path{lean/qtsl/Qtsl/Soundness.lean}, line 2526}\par\nopagebreak
\begin{leansrc}
theorem maskRegion_pair_elim {C : SimpC} {mask ext : Fin C.nVertices → Bool}
    {x : Fin C.ambDim → ℚ} {a b : Fin C.nVertices} {V₀ V₁ : ℚ}
    (hfil : (List.finRange C.nVertices).filter (fun v => mask v) = [a, b])
    (hVa : ∀ j, C.coord a j = V₀) (hVb : ∀ j, C.coord b j = V₁)
    (hne : V₀ ≠ V₁) (j0 : Fin C.ambDim)
    (h : MaskRegion C mask ext x) :
    (ext a = true → (V₁ - x j0) / (V₁ - V₀) ≤ 0) ∧
    (ext a = false → 0 ≤ (V₁ - x j0) / (V₁ - V₀)) ∧
    (ext b = true → (x j0 - V₀) / (V₁ - V₀) ≤ 0) ∧
    (ext b = false → 0 ≤ (x j0 - V₀) / (V₁ - V₀))
\end{leansrc}
\par\medskip\noindent{\small \texttt{PPLC.}\texttt{\bfseries sound\_convIntro}}\hfill{\footnotesize \path{lean/qtsl/Qtsl/Soundness.lean}, line 4593}\par\nopagebreak
\begin{leansrc}
theorem sound_convIntro {H : Hypernet} {seq premise : Sequent}
    {term : Term} {simplex : Simplex} {S : ExtensionSet}
    (ih : premise.semanticallyValid H)
    (hpre : seq.pre = premise.pre)
    (htermShape : term.bshape H = some [simplex.complex.ambDim])
    (hhalf : barycentricHalfspaceForm term simplex S premise.post)
    (hpost : seq.post = .inRegion term simplex S) :
    seq.semanticallyValid H
\end{leansrc}
\par\medskip\noindent{\small \texttt{PPLC.}\texttt{\bfseries f4VertexAtom\_holds\_inv}}\hfill{\footnotesize \path{lean/qtsl/Qtsl/Soundness.lean}, line 1442}\par\nopagebreak
\begin{leansrc}
theorem f4VertexAtom_holds_inv {H : Hypernet} {σ : H.Valuation} {t : Term}
    {inS : Bool} {off slp : ℚ}
    (h : Formula.holds σ (f4VertexAtom inS off slp t)) :
    ∃ q, Term.evalQ σ t = some q
\end{leansrc}
\par\medskip\noindent{\small \texttt{PPLC.}\texttt{\bfseries maskRegion\_pair\_intro}}\hfill{\footnotesize \path{lean/qtsl/Qtsl/Soundness.lean}, line 2575}\par\nopagebreak
\begin{leansrc}
theorem maskRegion_pair_intro {C : SimpC} {mask ext : Fin C.nVertices → Bool}
    {x : Fin C.ambDim → ℚ} {a b : Fin C.nVertices} {V₀ V₁ : ℚ}
    (hfil : (List.finRange C.nVertices).filter (fun v => mask v) = [a, b])
    (hVa : ∀ j, C.coord a j = V₀) (hVb : ∀ j, C.coord b j = V₁)
    (hne : V₀ ≠ V₁) (hdim : C.ambDim = 1) (j0 : Fin C.ambDim)
    (hw0t : ext a = true → (V₁ - x j0) / (V₁ - V₀) ≤ 0)
    (hw0f : ext a = false → 0 ≤ (V₁ - x j0) / (V₁ - V₀))
    (hw1t : ext b = true → (x j0 - V₀) / (V₁ - V₀) ≤ 0)
    (hw1f : ext b = false → 0 ≤ (x j0 - V₀) / (V₁ - V₀)) :
    MaskRegion C mask ext x
\end{leansrc}
\par\medskip\noindent{\small \texttt{PPLC.}\texttt{\bfseries sound\_arithAtom\_fold}}\hfill{\footnotesize \path{lean/qtsl/Qtsl/Soundness.lean}, line 4297}\par\nopagebreak
\begin{leansrc}
theorem sound_arithAtom_fold {H : Hypernet} {seq premise : Sequent}
    {termL termR : Term} {b : BowTie}
    (ih : premise.semanticallyValid H)
    (hpre : seq.pre = premise.pre)
    (hprem : atomicCompareZero b termL termR premise.post)
    (hconcl : atomicCompareForm b termL termR seq.post) :
    seq.semanticallyValid H
\end{leansrc}
\par\medskip\noindent{\small \texttt{PPLC.}\texttt{\bfseries sound\_arithAtom\_unfold}}\hfill{\footnotesize \path{lean/qtsl/Qtsl/Soundness.lean}, line 4318}\par\nopagebreak
\begin{leansrc}
theorem sound_arithAtom_unfold {H : Hypernet} {seq premise : Sequent}
    {termL termR : Term} {b : BowTie}
    (ih : premise.semanticallyValid H)
    (hpre : seq.pre = premise.pre)
    (hprem : atomicCompareForm b termL termR premise.post)
    (hconcl : atomicCompareZero b termL termR seq.post) :
    seq.semanticallyValid H
\end{leansrc}
\par\medskip\noindent{\small \texttt{PPLC.}\texttt{\bfseries atomicCompareForm}}\hfill{\footnotesize \path{lean/qtsl/Qtsl/Core/Catalogue.lean}, line 553}\par\nopagebreak
\begin{leansrc}
def atomicCompareForm (b : BowTie) (l r : Term) (F : Formula) : Prop
\end{leansrc}
\par\medskip\noindent{\small \texttt{PPLC.}\texttt{\bfseries atomicCompareZero}}\hfill{\footnotesize \path{lean/qtsl/Qtsl/Core/Catalogue.lean}, line 556}\par\nopagebreak
\begin{leansrc}
def atomicCompareZero (b : BowTie) (l r : Term) (F : Formula) : Prop
\end{leansrc}
\par\medskip\noindent{\small \texttt{PPLC.QTensor.}\texttt{\bfseries map2\_defined\_of\_rel2}}\hfill{\footnotesize \path{lean/qtsl/Qtsl/Soundness.lean}, line 3462}\par\nopagebreak
\begin{leansrc}
theorem QTensor.map2_defined_of_rel2 {r : ℚ → ℚ → Prop} {a b : QTensor}
    (h : QTensor.rel2 r a b) (f : ℚ → ℚ → ℚ) :
    ∃ c, QTensor.map2 f a b = some c
\end{leansrc}
\par\medskip\noindent{\small \texttt{PPLC.}\texttt{\bfseries BowTie}}\hfill{\footnotesize \path{lean/qtsl/Qtsl/Core/Catalogue.lean}, line 346}\par\nopagebreak
\begin{leansrc}
inductive BowTie where …
\end{leansrc}
\par\medskip\noindent{\small \texttt{PPLC.QTensor.}\texttt{\bfseries rel2\_sub\_zero\_iff}}\hfill{\footnotesize \path{lean/qtsl/Qtsl/Soundness.lean}, line 2466}\par\nopagebreak
\begin{leansrc}
theorem QTensor.rel2_sub_zero_iff {b : BowTie} {al ar d : QTensor}
    (hd : QTensor.map2 (· - ·) al ar = some d) :
    QTensor.rel2 (BowTie.holdsQ b) d ⟨[], [0], rfl⟩
      ↔ QTensor.rel2 (BowTie.holdsQ b) al ar
\end{leansrc}
\par\medskip\noindent{\small \texttt{PPLC.Term.}\texttt{\bfseries evalQ\_sub\_inv}}\hfill{\footnotesize \path{lean/qtsl/Qtsl/Soundness.lean}, line 1306}\par\nopagebreak
\begin{leansrc}
theorem Term.evalQ_sub_inv {H : Hypernet} (σ : H.Valuation) {l r : Term}
    {q : QTensor} (h : Term.evalQ σ (.sub l r) = some q) :
    ∃ a b, Term.evalQ σ l = some a ∧ Term.evalQ σ r = some b ∧
      QTensor.map2 (· - ·) a b = some q
\end{leansrc}
\par\medskip\noindent{\small \texttt{PPLC.}\texttt{\bfseries PArithAtomData}}\hfill{\footnotesize \path{lean/qtsl/Qtsl/Core/Catalogue.lean}, line 1471}\par\nopagebreak
\begin{leansrc}
inductive PArithAtomData : Hypernet → Sequent → Type where …
\end{leansrc}
\par\medskip\noindent{\small \texttt{PPLC.}\texttt{\bfseries allForwardDiffBernCoeffsSatisfyP}}\hfill{\footnotesize \path{lean/qtsl/Qtsl/Core/Catalogue.lean}, line 760}\par\nopagebreak
\begin{leansrc}
def allForwardDiffBernCoeffsSatisfyP
    (simplex : Simplex) (S : ExtensionSet) (degree : Nat)
    (polynomial : Term) (predicate : Formula) : Prop
\end{leansrc}
\par\medskip\noindent{\small \texttt{PPLC.}\texttt{\bfseries satisfiesPredQ}}\hfill{\footnotesize \path{lean/qtsl/Qtsl/Core/Catalogue.lean}, line 701}\par\nopagebreak
\begin{leansrc}
def satisfiesPredQ (P : Formula) (c : Rat) : Bool
\end{leansrc}
\par\medskip\noindent{\small \texttt{PPLC.}\texttt{\bfseries bernConeCoeffCheck}}\hfill{\footnotesize \path{lean/qtsl/Qtsl/Core/Catalogue.lean}, line 727}\par\nopagebreak
\begin{leansrc}
def bernConeCoeffCheck (simplex : Simplex) (S : ExtensionSet) (degree : Nat)
    (polynomial : Term) (predicate : Formula) : Bool
\end{leansrc}
\par\medskip\noindent{\small \texttt{PPLC.}\texttt{\bfseries sound\_bernCone}}\hfill{\footnotesize \path{lean/qtsl/Qtsl/Soundness.lean}, line 5498}\par\nopagebreak
\begin{leansrc}
theorem sound_bernCone {H : Hypernet} {seq regionPremise : Sequent}
    {termX : Term} {simplex : Simplex} {S : ExtensionSet} {degree : Nat}
    {polynomial : Term} {predicate : Formula}
    (ih : regionPremise.semanticallyValid H)
    (hpre : seq.pre = regionPremise.pre)
    (hregionPost : regionPremise.post = .inRegion termX simplex S)
    (hpolyTie : polynomial.freeWires = [] ∨
      ∃ w, termX = Term.var w ∧ polynomial.freeWires ⊆ [w])
    (hpredShape : (¬ isExtensionSetEmpty S → isConicPredicate predicate)
      ∧ (isExtensionSetEmpty S → isConvexPredicate predicate))
    (hcoeff : allForwardDiffBernCoeffsSatisfyP simplex S degree polynomial
      predicate)
    (hpostApply : applyPredicateToTerm predicate polynomial seq.post) :
    seq.semanticallyValid H
\end{leansrc}
\par\medskip\noindent{\small \texttt{PPLC.}\texttt{\bfseries sound\_bernCone\_empty}}\hfill{\footnotesize \path{lean/qtsl/Qtsl/Soundness.lean}, line 5058}\par\nopagebreak
\begin{leansrc}
theorem sound_bernCone_empty {H : Hypernet} {seq regionPremise : Sequent}
    {termX : Term} {simplex : Simplex} {S : ExtensionSet} {degree : Nat}
    {polynomial : Term} {predicate : Formula}
    (ih : regionPremise.semanticallyValid H)
    (hpre : seq.pre = regionPremise.pre)
    (hregionPost : regionPremise.post = .inRegion termX simplex S)
    (hpolyTie : polynomial.freeWires = [] ∨
      ∃ w, termX = Term.var w ∧ polynomial.freeWires ⊆ [w])
    (hconvex : isExtensionSetEmpty S → isConvexPredicate predicate)
    (hSempty : isExtensionSetEmpty S)
    (hcoeff : allForwardDiffBernCoeffsSatisfyP simplex S degree polynomial
      predicate)
    (hpostApply : applyPredicateToTerm predicate polynomial seq.post) :
    seq.semanticallyValid H
\end{leansrc}
\par\medskip\noindent{\small \texttt{PPLC.}\texttt{\bfseries MaskRegion}}\hfill{\footnotesize \path{lean/qtsl/Qtsl/Catalogue.lean}, line 376}\par\nopagebreak
\begin{leansrc}
def MaskRegion (C : SimpC) (mask ext : Fin C.nVertices → Bool)
    (x : Fin C.ambDim → ℚ) : Prop
\end{leansrc}
\par\medskip\noindent{\small \texttt{PPLC.Term.}\texttt{\bfseries evalQ\_brRead}}\hfill{\footnotesize \path{lean/qtsl/Qtsl/Soundness.lean}, line 1578}\par\nopagebreak
\begin{leansrc}
theorem Term.evalQ_brRead {H : Hypernet} (σ : H.Valuation) (s : Shape) :
    ∀ (t : Term),
    (∀ w ∈ t.freeWires, ∃ qw, σ.wireQ w = some qw ∧
      (qw.shape = [] ∨ qw.shape = s)) →
    t.idxFree = true →
    ∃ q, Term.evalQ σ t = some q ∧ (q.shape = [] ∨ q.shape = s) ∧
      (q.shape = [] → ∀ w ∈ t.freeWires, ∀ qw, σ.wireQ w = some qw →
        qw.shape = []) ∧
      ∀ i, i < Shape.totalSize s →
        q.brRead i = QArith.Poly.eval (σ.entryVal i) t.toPoly
\end{leansrc}
\par\medskip\noindent{\small \texttt{PPLC.}\texttt{\bfseries maskRegion\_rat\_weights}}\hfill{\footnotesize \path{lean/qtsl/Qtsl/Soundness.lean}, line 4798}\par\nopagebreak
\begin{leansrc}
theorem maskRegion_rat_weights {C : SimpC}
    {mask ext : Fin C.nVertices → Bool} {x : Fin C.ambDim → ℚ}
    (hmask : MaskRegion C mask ext x)
    (hext0 : ∀ v, ext v = false) :
    ∃ qw : Fin ((List.finRange C.nVertices).filter
        (fun v => mask v)).length → ℚ,
      (∀ i, 0 ≤ qw i) ∧ (∑ i, qw i) = 1 ∧
      ∀ j : Fin C.ambDim,
        (∑ i, qw i * C.coord (((List.finRange C.nVertices).filter
          (fun v => mask v)).get i) j) = x j
\end{leansrc}
\par\medskip\noindent{\small \texttt{PPLC.ConvexRat.}\texttt{\bfseries exists\_nonneg\_rat\_solution}}\hfill{\footnotesize \path{lean/qtsl/Qtsl/ConvexRat.lean}, line 269}\par\nopagebreak
\begin{leansrc}
theorem exists_nonneg_rat_solution (A : Fin k → Fin m → ℚ)
    (b : Fin m → ℚ) (c : Fin k → ℝ) (hc0 : ∀ i, 0 ≤ c i)
    (hcsol : ∀ j, (∑ i, c i * (A i j : ℝ)) = (b j : ℝ)) :
    ∃ q : Fin k → ℚ, (∀ i, 0 ≤ q i) ∧ ∀ j, (∑ i, q i * A i j) = b j
\end{leansrc}
\par\medskip\noindent{\small \texttt{PPLC.}\texttt{\bfseries classifyPred}}\hfill{\footnotesize \path{lean/qtsl/Qtsl/Core/Catalogue.lean}, line 667}\par\nopagebreak
\begin{leansrc}
def classifyPred : Formula → Option PredClass
  | .and l r => do
    let a ← classifyPred l
    let b ← classifyPred r
    pure ⟨a.conic && b.conic, a.upClosed && b.upClosed⟩
  | f => classifyPredAtom f

def isConicPredicate (P : Formula) : Prop
\end{leansrc}
\par\medskip\noindent{\small \texttt{PPLC.}\texttt{\bfseries satisfiesPredQ\_convex}}\hfill{\footnotesize \path{lean/qtsl/Qtsl/Soundness.lean}, line 2871}\par\nopagebreak
\begin{leansrc}
theorem satisfiesPredQ_convex {P : Formula} {cls : PredClass}
    (hP : classifyPred P = some cls) :
    ∀ a₁ b₁ l : ℚ, 0 ≤ l → l ≤ 1 →
      satisfiesPredQ P a₁ = true → satisfiesPredQ P b₁ = true →
      satisfiesPredQ P (l * a₁ + (1 - l) * b₁) = true
\end{leansrc}
\par\medskip\noindent{\small \texttt{PPLC.QArith.}\texttt{\bfseries fdBernstein\_sound\_dim1}}\hfill{\footnotesize \path{lean/qtsl/Qtsl/QArithKernel.lean}, line 2219}\par\nopagebreak
\begin{leansrc}
theorem fdBernstein_sound_dim1
    (verts : List (List ℚ)) (smask : List Bool) (d : Nat) (p : Poly)
    (t : FDTable) (pred : ℚ → Bool) (u₀ : ℚ)
    (hmask : smask = [true, false] ∨ smask = [false, true])
    (hfd : fdBernstein verts smask d p = some t)
    (hgrid : gridAll pred t = true)
    (hu : 0 ≤ u₀)
    (hzero : pred 0 = true)
    (hscale : ∀ a c : ℚ, 0 ≤ c → pred a = true → pred (c * a) = true)
    (hconv : ∀ a b l : ℚ, 0 ≤ l → l ≤ 1 → pred a = true → pred b = true →
      pred (l * a + (1 - l) * b) = true) :
    pred (Poly.eval (fun j =>
      ((List.range verts.length).map
        fun i => (verts.getD i []).getD j 0 *
          (if smask.getD i false then -u₀ else 1 + u₀)).sum) p) = true
\end{leansrc}
\par\medskip\noindent{\small \texttt{PPLC.QArith.}\texttt{\bfseries fdBernstein\_sound\_const}}\hfill{\footnotesize \path{lean/qtsl/Qtsl/QArithKernel.lean}, line 1733}\par\nopagebreak
\begin{leansrc}
theorem fdBernstein_sound_const
    (verts : List (List ℚ)) (smask : List Bool) (d : Nat) (c₀ : ℚ)
    (t : FDTable) (pred : ℚ → Bool)
    (hfd : fdBernstein verts smask d (Poly.const c₀) = some t)
    (hgrid : gridAll pred t = true) :
    pred c₀ = true
\end{leansrc}
\par\medskip\noindent{\small \texttt{PPLC.}\texttt{\bfseries satisfiesPredQ\_conic\_zero}}\hfill{\footnotesize \path{lean/qtsl/Qtsl/Soundness.lean}, line 3066}\par\nopagebreak
\begin{leansrc}
theorem satisfiesPredQ_conic_zero {P : Formula} {cls : PredClass}
    (hP : classifyPred P = some cls) (hc : cls.conic = true) :
    satisfiesPredQ P 0 = true
\end{leansrc}
\par\medskip\noindent{\small \texttt{PPLC.}\texttt{\bfseries satisfiesPredQ\_conic\_scale}}\hfill{\footnotesize \path{lean/qtsl/Qtsl/Soundness.lean}, line 3102}\par\nopagebreak
\begin{leansrc}
theorem satisfiesPredQ_conic_scale {P : Formula} {cls : PredClass}
    (hP : classifyPred P = some cls) (hc : cls.conic = true) :
    ∀ a c : ℚ, 0 ≤ c → satisfiesPredQ P a = true →
      satisfiesPredQ P (c * a) = true
\end{leansrc}
\par\medskip\noindent{\small \texttt{PPLC.QArith.}\texttt{\bfseries fdBernstein}}\hfill{\footnotesize \path{lean/qtsl/Qtsl/Core/QArith.lean}, line 158}\par\nopagebreak
\begin{leansrc}
def fdBernstein (verts : List (List Rat)) (smask : List Bool) (d : Nat)
    (p : Poly) : Option FDTable
\end{leansrc}
\par\medskip\noindent{\small \texttt{PPLC.QArith.}\texttt{\bfseries gridAll}}\hfill{\footnotesize \path{lean/qtsl/Qtsl/Core/QArith.lean}, line 190}\par\nopagebreak
\begin{leansrc}
def gridAll (pred : Rat → Bool) (t : FDTable) : Bool
\end{leansrc}
\par\medskip\noindent{\small \texttt{PPLC.Term.}\texttt{\bfseries idxBounded}}\hfill{\footnotesize \path{lean/qtsl/Qtsl/Core/Catalogue.lean}, line 478}\par\nopagebreak
\begin{leansrc}
def Term.idxBounded (n : Nat) : Term → Bool
  | .var _ | .const _ => true
  | .add l r | .mul l r | .sub l r => l.idxBounded n && r.idxBounded n
  | .neg t => t.idxBounded n
  | .idx t i => decide (i < n) && t.idxBounded n

def Formula.freeWires : Formula → List Wire
  | .atom _ l r => l.freeWires ++ r.freeWires
  | .inRegion t _ _ => t.freeWires
  | .and l r | .or l r => l.freeWires ++ r.freeWires
  | .tt | .ff => []

def Term.type (H : Hypernet) : Term → Option TensorType
  | .var w => H.wireType w
  | .const _ => some ⟨[], Dtype.Q⟩
  | .idx t i =>

    match Term.type H t with
    | some τ =>
      match τ.shape with
      | [e] => if i < e then some ⟨[], τ.dtype⟩ else none
      | _ => none
    | none => none
  | .add l r | .mul l r | .sub l r =>
    match Term.type H l, Term.type H r with
    | some τ₁, some τ₂ => if τ₁ = τ₂ then some τ₁ else none
    | _, _ => none
  | .neg t => Term.type H t

def Formula.WellTyped (H : Hypernet) : Formula → Prop
  | .atom _ l r => ∃ τ, l.type H = some τ ∧ r.type H = some τ
  | .inRegion t simplex S =>
      (∃ τ, t.type H = some τ ∧ τ.shape = [simplex.complex.ambDim]) ∧
      S.simplex = simplex
  | .and l r => l.WellTyped H ∧ r.WellTyped H
  | .or l r  => l.WellTyped H ∧ r.WellTyped H
  | .tt      => True
  | .ff      => True

def Term.bshape (H : Hypernet) : Term → Option Shape
  …
\end{leansrc}
\par\medskip\noindent{\small \texttt{PPLC.}\texttt{\bfseries applyPredicateToTerm}}\hfill{\footnotesize \path{lean/qtsl/Qtsl/Core/Catalogue.lean}, line 631}\par\nopagebreak
\begin{leansrc}
def applyPredicateToTerm (P : Formula) (arg : Term) (F : Formula) : Prop
\end{leansrc}
\par\medskip\noindent{\small \texttt{PPLC.}\texttt{\bfseries holds\_substHole\_of\_satisfies}}\hfill{\footnotesize \path{lean/qtsl/Qtsl/Soundness.lean}, line 3144}\par\nopagebreak
\begin{leansrc}
theorem holds_substHole_of_satisfies {H : Hypernet} (σ : H.Valuation)
    {P : Formula} {cls : PredClass} (hP : classifyPred P = some cls)
    {arg : Term} {qa : QTensor}
    (ha : Term.evalQ σ arg = some qa) (hsa : SingletonShape qa.shape)
    (hsat : satisfiesPredQ P qa.scalar = true) :
    Formula.holds σ (Formula.subst predicateHole arg P)
\end{leansrc}
\par\medskip\noindent{\small \texttt{PPLC.ConeLemma.}\texttt{\bfseries StarIdentity}}\hfill{\footnotesize \path{lean/qtsl/Qtsl/ConeLemma.lean}, line 18}\par\nopagebreak
\begin{leansrc}
def StarIdentity : Prop
\end{leansrc}
\par\medskip\noindent{\small \texttt{PPLC.}\texttt{\bfseries PNormData}}\hfill{\footnotesize \path{lean/qtsl/Qtsl/Core/Catalogue.lean}, line 1505}\par\nopagebreak
\begin{leansrc}
inductive PNormData : Hypernet → Sequent → Type where …
\end{leansrc}
\par\medskip\noindent{\small \texttt{PPLC.}\texttt{\bfseries sound\_norm}}\hfill{\footnotesize \path{lean/qtsl/Qtsl/Soundness.lean}, line 4494}\par\nopagebreak
\begin{leansrc}
theorem sound_norm {H : Hypernet} {seq : Sequent} {termL termR : Term}
    (hcanon : polynomialCanonicalEq termL termR)
    (htyEq : termL.bshape H = termR.bshape H)
    (htyDef : ∃ s, termL.bshape H = some s)
    (hfinL : termL.finWires H) (hfinR : termR.finWires H)
    (hpost : seq.post = .atom .eq termL termR) :
    seq.semanticallyValid H
\end{leansrc}
\par\medskip\noindent{\small \texttt{PPLC.}\texttt{\bfseries polynomialCanonicalEq}}\hfill{\footnotesize \path{lean/qtsl/Qtsl/Core/Catalogue.lean}, line 698}\par\nopagebreak
\begin{leansrc}
def polynomialCanonicalEq (l r : Term) : Prop
\end{leansrc}
\par\medskip\noindent{\small \texttt{PPLC.Term.}\texttt{\bfseries toPoly}}\hfill{\footnotesize \path{lean/qtsl/Qtsl/Core/Catalogue.lean}, line 414}\par\nopagebreak
\begin{leansrc}
def Term.toPoly : Term → QArith.Poly
  | .var w => QArith.Poly.var w
  | .const c => QArith.Poly.const c
  | .add l r => QArith.Poly.add l.toPoly r.toPoly
  | .mul l r => QArith.Poly.mul l.toPoly r.toPoly
  | .sub l r => QArith.Poly.sub l.toPoly r.toPoly
  | .neg t => QArith.Poly.neg t.toPoly

  | .idx _ _ => QArith.Poly.const 0

def Term.evalQ? : Term → Option Rat
  | .var _ => none
  | .const c => some c
  | .add l r => do pure ((← l.evalQ?) + (← r.evalQ?))
  | .mul l r => do pure ((← l.evalQ?) * (← r.evalQ?))
  | .sub l r => do pure ((← l.evalQ?) - (← r.evalQ?))
  | .neg t => do pure (-(← t.evalQ?))
  | .idx _ _ => none

def BowTie.evalQ : BowTie → Rat → Rat → Bool
  | .lt, x, y => decide (x < y)
  | .le, x, y => decide (x ≤ y)
  | .eq, x, y => decide (x = y)
  | .ge, x, y => decide (y ≤ x)
  | .gt, x, y => decide (y < x)
  | .ne, x, y => decide (¬ x = y)

def Formula.evalClosed : Formula → Option Bool
  | .atom b l r => do pure (b.evalQ (← l.evalQ?) (← r.evalQ?))
  | .inRegion _ _ _ => none
  | .and l r => do pure ((← l.evalClosed) && (← r.evalClosed))
  | .or l r => do pure ((← l.evalClosed) || (← r.evalClosed))
  | .tt => some true
  | .ff => some false

def baryLambdas1D (cx : SimpC) (dContains : Fin cx.nVertices → Bool)
    (t : Term) : Option (List (Fin cx.nVertices × Term))
\end{leansrc}
\par\medskip\noindent{\small \texttt{PPLC.Term.}\texttt{\bfseries idxFree}}\hfill{\footnotesize \path{lean/qtsl/Qtsl/Core/Catalogue.lean}, line 472}\par\nopagebreak
\begin{leansrc}
def Term.idxFree : Term → Bool
  | .var _ | .const _ => true
  | .add l r | .mul l r | .sub l r => l.idxFree && r.idxFree
  | .neg t => t.idxFree
  | .idx _ _ => false

def Term.idxBounded (n : Nat) : Term → Bool
  | .var _ | .const _ => true
  | .add l r | .mul l r | .sub l r => l.idxBounded n && r.idxBounded n
  | .neg t => t.idxBounded n
  | .idx t i => decide (i < n) && t.idxBounded n

def Formula.freeWires : Formula → List Wire
  | .atom _ l r => l.freeWires ++ r.freeWires
  | .inRegion t _ _ => t.freeWires
  | .and l r | .or l r => l.freeWires ++ r.freeWires
  | .tt | .ff => []

def Term.type (H : Hypernet) : Term → Option TensorType
  | .var w => H.wireType w
  | .const _ => some ⟨[], Dtype.Q⟩
  | .idx t i =>

    match Term.type H t with
    | some τ =>
      match τ.shape with
      | [e] => if i < e then some ⟨[], τ.dtype⟩ else none
      | _ => none
    | none => none
  | .add l r | .mul l r | .sub l r =>
    match Term.type H l, Term.type H r with
    | some τ₁, some τ₂ => if τ₁ = τ₂ then some τ₁ else none
    | _, _ => none
  | .neg t => Term.type H t

def Formula.WellTyped (H : Hypernet) : Formula → Prop
  | .atom _ l r => ∃ τ, l.type H = some τ ∧ r.type H = some τ
  | .inRegion t simplex S =>
      (∃ τ, t.type H = some τ ∧ τ.shape = [simplex.complex.ambDim]) ∧
      S.simplex = simplex
  …
\end{leansrc}
\par\medskip\noindent{\small \texttt{PPLC.}\texttt{\bfseries farkasCombinationValid}}\hfill{\footnotesize \path{lean/qtsl/Qtsl/Core/Catalogue.lean}, line 918}\par\nopagebreak
\begin{leansrc}
def farkasCombinationValid (cc : Term) (premise : Formula)
    (conclusion : Term) : Prop
\end{leansrc}
\par\medskip\noindent{\small \texttt{PPLC.}\texttt{\bfseries farkasEntries}}\hfill{\footnotesize \path{lean/qtsl/Qtsl/Core/Catalogue.lean}, line 778}\par\nopagebreak
\begin{leansrc}
def farkasEntries : Term → Option (List (Rat × Term))
  | .const c => if c = 0 then some [] else none
  | .add (.mul (.const l) e) rest =>
    (farkasEntries rest).map ((l, e) :: ·)
  | _ => none

def wireCoordKey (w i : Nat) : Nat
\end{leansrc}
\par\medskip\noindent{\small \texttt{PPLC.}\texttt{\bfseries atomLeZeroPolys}}\hfill{\footnotesize \path{lean/qtsl/Qtsl/Core/Catalogue.lean}, line 769}\par\nopagebreak
\begin{leansrc}
def atomLeZeroPolys : Formula → List QArith.Poly
  | .atom .le l r => [QArith.Poly.sub l.toPoly r.toPoly]
  | .atom .lt l r => [QArith.Poly.sub l.toPoly r.toPoly]
  | .atom .ge l r => [QArith.Poly.sub r.toPoly l.toPoly]
  | .atom .gt l r => [QArith.Poly.sub r.toPoly l.toPoly]
  | .atom .eq l r => [QArith.Poly.sub l.toPoly r.toPoly,
                      QArith.Poly.sub r.toPoly l.toPoly]
  | _ => []

def farkasEntries : Term → Option (List (Rat × Term))
  | .const c => if c = 0 then some [] else none
  | .add (.mul (.const l) e) rest =>
    (farkasEntries rest).map ((l, e) :: ·)
  | _ => none

def wireCoordKey (w i : Nat) : Nat
\end{leansrc}
\par\medskip\noindent{\small \texttt{PPLC.}\texttt{\bfseries farkasCheckCore}}\hfill{\footnotesize \path{lean/qtsl/Qtsl/Core/Catalogue.lean}, line 899}\par\nopagebreak
\begin{leansrc}
def farkasCheckCore (cc : Term) (premise : Formula) (conclusion : Term) : Bool
\end{leansrc}
\par\medskip\noindent{\small \texttt{PPLC.}\texttt{\bfseries farkasCheck}}\hfill{\footnotesize \path{lean/qtsl/Qtsl/Core/Catalogue.lean}, line 914}\par\nopagebreak
\begin{leansrc}
def farkasCheck (cc : Term) (premise : Formula) (conclusion : Term) : Bool
\end{leansrc}
\par\medskip\noindent{\small \texttt{PPLC.}\texttt{\bfseries sound\_farkas}}\hfill{\footnotesize \path{lean/qtsl/Qtsl/Soundness.lean}, line 4337}\par\nopagebreak
\begin{leansrc}
theorem sound_farkas {H : Hypernet} {seq : Sequent}
    {combinationCoeffs conclusion : Term}
    (hshape : ∃ s, conclusion.bshape H = some s)
    (hfin : conclusion.finWires H)
    (hwpos : ∀ w ∈ seq.pre.freeWires ++ conclusion.freeWires,
      ∃ τ, H.wireType w = some τ ∧ 0 < Shape.totalSize τ.shape)
    (hvalid : farkasCombinationValid combinationCoeffs seq.pre conclusion)
    (hpost : seq.post = .atom .le conclusion (.const 0)) :
    seq.semanticallyValid H
\end{leansrc}
\par\medskip\noindent{\small \texttt{PPLC.}\texttt{\bfseries atomLeZeroPolysW}}\hfill{\footnotesize \path{lean/qtsl/Qtsl/Core/Catalogue.lean}, line 800}\par\nopagebreak
\begin{leansrc}
def atomLeZeroPolysW : Formula → List QArith.Poly
  | .atom op l r =>
    match l.asCoordPolyIdxW? with
    | none => []
    | some pl =>
      match r.asCoordPolyIdxW? with
      | none => []
      | some pr =>
        match op with
        | .le | .lt => [QArith.Poly.sub pl pr]
        | .ge | .gt => [QArith.Poly.sub pr pl]
        | .eq => [QArith.Poly.sub pl pr, QArith.Poly.sub pr pl]
        | .ne => []
  | _ => []

def farkasEntriesW : List (Rat × Term) → Option (List (Rat × QArith.Poly))
  | [] => some []
  | (l, e) :: rest =>
    match e.asCoordPolyIdxW?, farkasEntriesW rest with
    | some pe, some prest => some ((l, pe) :: prest)
    | _, _ => none

def farkasCheckW (cc : Term) (premise : Formula) (conclusion : Term) : Bool
\end{leansrc}
\par\medskip\noindent{\small \texttt{PPLC.}\texttt{\bfseries atomLeZeroPolysM}}\hfill{\footnotesize \path{lean/qtsl/Qtsl/Core/Catalogue.lean}, line 854}\par\nopagebreak
\begin{leansrc}
def atomLeZeroPolysM : Formula → List QArith.Poly
  | .atom op l r =>
    match l.asCoordPolyMixed? with
    | none => []
    | some pl =>
      match r.asCoordPolyMixed? with
      | none => []
      | some pr =>
        match op with
        | .le | .lt => [QArith.Poly.sub pl pr]
        | .ge | .gt => [QArith.Poly.sub pr pl]
        | .eq => [QArith.Poly.sub pl pr, QArith.Poly.sub pr pl]
        | .ne => []
  | _ => []

def farkasEntriesM : List (Rat × Term) → Option (List (Rat × QArith.Poly))
  | [] => some []
  | (l, e) :: rest =>
    match e.asCoordPolyMixed?, farkasEntriesM rest with
    | some pe, some prest => some ((l, pe) :: prest)
    | _, _ => none

def farkasCheckM (cc : Term) (premise : Formula) (conclusion : Term) : Bool
\end{leansrc}
\par\medskip\noindent{\small \texttt{PPLC.}\texttt{\bfseries farkasCheckW}}\hfill{\footnotesize \path{lean/qtsl/Qtsl/Core/Catalogue.lean}, line 822}\par\nopagebreak
\begin{leansrc}
def farkasCheckW (cc : Term) (premise : Formula) (conclusion : Term) : Bool
\end{leansrc}
\par\medskip\noindent{\small \texttt{PPLC.}\texttt{\bfseries farkasCheckM}}\hfill{\footnotesize \path{lean/qtsl/Qtsl/Core/Catalogue.lean}, line 876}\par\nopagebreak
\begin{leansrc}
def farkasCheckM (cc : Term) (premise : Formula) (conclusion : Term) : Bool
\end{leansrc}
\par\medskip\noindent{\small \texttt{PPLC.}\texttt{\bfseries wireCoordKey}}\hfill{\footnotesize \path{lean/qtsl/Qtsl/Core/Catalogue.lean}, line 784}\par\nopagebreak
\begin{leansrc}
def wireCoordKey (w i : Nat) : Nat
\end{leansrc}
\par\medskip\noindent{\small \texttt{PPLC.}\texttt{\bfseries farkas\_fold\_eval\_nonpos}}\hfill{\footnotesize \path{lean/qtsl/Qtsl/Soundness.lean}, line 3445}\par\nopagebreak
\begin{leansrc}
theorem farkas_fold_eval_nonpos (ρ : Nat → ℚ) :
    ∀ (entries : List (ℚ × Term)),
    (∀ e ∈ entries, 0 ≤ e.1 ∧ QArith.Poly.eval ρ e.2.toPoly ≤ 0) →
    QArith.Poly.eval ρ
      (entries.foldr
        (fun (le : ℚ × Term) acc =>
          QArith.Poly.add (QArith.Poly.scale le.1 le.2.toPoly) acc) []) ≤ 0
\end{leansrc}
\par\medskip\noindent{\small \texttt{PPLC.}\texttt{\bfseries farkas\_fold\_eval\_nonpos\_polys}}\hfill{\footnotesize \path{lean/qtsl/Qtsl/Soundness.lean}, line 3830}\par\nopagebreak
\begin{leansrc}
theorem farkas_fold_eval_nonpos_polys (ρ : Nat → ℚ) :
    ∀ (pentries : List (ℚ × QArith.Poly)),
    (∀ e ∈ pentries, 0 ≤ e.1 ∧ QArith.Poly.eval ρ e.2 ≤ 0) →
    QArith.Poly.eval ρ
      (pentries.foldr
        (fun (le : ℚ × QArith.Poly) acc =>
          QArith.Poly.add (QArith.Poly.scale le.1 le.2) acc) []) ≤ 0
\end{leansrc}
\par\medskip\noindent{\small \texttt{PPLC.}\texttt{\bfseries PFarkasData}}\hfill{\footnotesize \path{lean/qtsl/Qtsl/Core/Catalogue.lean}, line 1517}\par\nopagebreak
\begin{leansrc}
inductive PFarkasData : Hypernet → Sequent → Type where …
\end{leansrc}
\par\medskip\noindent{\small \texttt{PPLC.}\texttt{\bfseries coverInfeasibilityValid}}\hfill{\footnotesize \path{lean/qtsl/Qtsl/Core/Catalogue.lean}, line 1151}\par\nopagebreak
\begin{leansrc}
def coverInfeasibilityValid
    (_termX : Term) (simplex : Simplex) (S : ExtensionSet)
    (nBranches : Nat)
    (branchSimplex : Fin nBranches → Simplex)
    (branchS : Fin nBranches → ExtensionSet)
    (_coverCert : Term) : Prop
\end{leansrc}
\par\medskip\noindent{\small \texttt{PPLC.}\texttt{\bfseries coverExtendPass}}\hfill{\footnotesize \path{lean/qtsl/Qtsl/Core/Catalogue.lean}, line 1127}\par\nopagebreak
\begin{leansrc}
def coverExtendPass (ivs : List (Rat × Rat)) (r : Rat) : Rat
\end{leansrc}
\par\medskip\noindent{\small \texttt{PPLC.}\texttt{\bfseries coverReach}}\hfill{\footnotesize \path{lean/qtsl/Qtsl/Core/Catalogue.lean}, line 1130}\par\nopagebreak
\begin{leansrc}
def coverReach : Nat → List (Rat × Rat) → Rat → Rat
  | 0, _, r => r
  | n + 1, ivs, r => coverReach n ivs (coverExtendPass ivs r)

def coverChainCheck (ivs : List (Rat × Rat)) (lo hi : Rat) : Bool
\end{leansrc}
\par\medskip\noindent{\small \texttt{PPLC.}\texttt{\bfseries coverChainCheck}}\hfill{\footnotesize \path{lean/qtsl/Qtsl/Core/Catalogue.lean}, line 1134}\par\nopagebreak
\begin{leansrc}
def coverChainCheck (ivs : List (Rat × Rat)) (lo hi : Rat) : Bool
\end{leansrc}
\par\medskip\noindent{\small \texttt{PPLC.}\texttt{\bfseries coverCheck}}\hfill{\footnotesize \path{lean/qtsl/Qtsl/Core/Catalogue.lean}, line 1137}\par\nopagebreak
\begin{leansrc}
def coverCheck (simplex : Simplex) (S : ExtensionSet) (nBranches : Nat)
    (branchSimplex : Fin nBranches → Simplex)
    (branchS : Fin nBranches → ExtensionSet) : Bool
\end{leansrc}
\par\medskip\noindent{\small \texttt{PPLC.}\texttt{\bfseries sound\_subdivide}}\hfill{\footnotesize \path{lean/qtsl/Qtsl/Soundness.lean}, line 7758}\par\nopagebreak
\begin{leansrc}
theorem sound_subdivide {H : Hypernet} {seq regionPremise : Sequent}
    {nBranches : Nat} {branchPremise : Fin nBranches → Sequent}
    {termX : Term} {simplex : Simplex} {S : ExtensionSet}
    {branchSimplex : Fin nBranches → Simplex}
    {branchS : Fin nBranches → ExtensionSet}
    {coverCertificate : Term}
    (ihRegion : regionPremise.semanticallyValid H)
    (ihBranch : ∀ k, (branchPremise k).semanticallyValid H)
    (hpre : seq.pre = regionPremise.pre)
    (hregionPost : regionPremise.post = .inRegion termX simplex S)
    (hbranchPre : ∀ k, (branchPremise k).pre
                   = .and regionPremise.pre
                          (.inRegion termX (branchSimplex k) (branchS k)))
    (hbranchPost : ∀ k, (branchPremise k).post = seq.post)
    (hcover : coverInfeasibilityValid termX simplex S nBranches
                branchSimplex branchS coverCertificate) :
    seq.semanticallyValid H
\end{leansrc}
\par\medskip\noindent{\small \texttt{PPLC.}\texttt{\bfseries interval1D\_elim}}\hfill{\footnotesize \path{lean/qtsl/Qtsl/Soundness.lean}, line 7642}\par\nopagebreak
\begin{leansrc}
theorem interval1D_elim {S : ExtensionSet} {lo hi : ℚ}
    (hint : S.interval1D? = some (lo, hi))
    {x : Fin S.simplex.complex.ambDim → ℚ}
    (hmask : MaskRegion S.simplex.complex S.simplex.contains S.contains x)
    (j0 : Fin S.simplex.complex.ambDim) :
    lo ≤ x j0 ∧ x j0 ≤ hi ∧ lo < hi
\end{leansrc}
\par\medskip\noindent{\small \texttt{PPLC.}\texttt{\bfseries interval1D\_intro}}\hfill{\footnotesize \path{lean/qtsl/Qtsl/Soundness.lean}, line 7676}\par\nopagebreak
\begin{leansrc}
theorem interval1D_intro {S : ExtensionSet} {lo hi : ℚ}
    (hint : S.interval1D? = some (lo, hi))
    {x : Fin S.simplex.complex.ambDim → ℚ}
    (j0 : Fin S.simplex.complex.ambDim)
    (hxlo : lo ≤ x j0) (hxhi : x j0 ≤ hi) :
    MaskRegion S.simplex.complex S.simplex.contains S.contains x
\end{leansrc}
\par\medskip\noindent{\small \texttt{PPLC.}\texttt{\bfseries ChainInv}}\hfill{\footnotesize \path{lean/qtsl/Qtsl/Soundness.lean}, line 7700}\par\nopagebreak
\begin{leansrc}
def ChainInv (ivs : List (ℚ × ℚ)) (lo r : ℚ) : Prop
\end{leansrc}
\par\medskip\noindent{\small \texttt{PPLC.}\texttt{\bfseries chainInv\_extendPass}}\hfill{\footnotesize \path{lean/qtsl/Qtsl/Soundness.lean}, line 7704}\par\nopagebreak
\begin{leansrc}
theorem chainInv_extendPass (ivs : List (ℚ × ℚ)) (lo : ℚ) :
    ∀ (sub : List (ℚ × ℚ)), (∀ iv ∈ sub, iv ∈ ivs) →
    ∀ (r : ℚ), ChainInv ivs lo r →
    ChainInv ivs lo (coverExtendPass sub r)
\end{leansrc}
\par\medskip\noindent{\small \texttt{PPLC.}\texttt{\bfseries chainInv\_reach}}\hfill{\footnotesize \path{lean/qtsl/Qtsl/Soundness.lean}, line 7735}\par\nopagebreak
\begin{leansrc}
theorem chainInv_reach (ivs : List (ℚ × ℚ)) (lo : ℚ) :
    ∀ (n : Nat) (r : ℚ), ChainInv ivs lo r →
    ChainInv ivs lo (coverReach n ivs r)
\end{leansrc}
\par\medskip\noindent{\small \texttt{PPLC.}\texttt{\bfseries coverChain\_covers}}\hfill{\footnotesize \path{lean/qtsl/Qtsl/Soundness.lean}, line 7745}\par\nopagebreak
\begin{leansrc}
theorem coverChain_covers {ivs : List (ℚ × ℚ)} {lo hi : ℚ} (hlt : lo < hi)
    (hcheck : coverChainCheck ivs lo hi = true) :
    ∀ x : ℚ, lo ≤ x → x ≤ hi → ∃ iv ∈ ivs, iv.1 ≤ x ∧ x ≤ iv.2
\end{leansrc}
\par\medskip\noindent{\small \texttt{PPLC.}\texttt{\bfseries PSubdivideData}}\hfill{\footnotesize \path{lean/qtsl/Qtsl/Core/Catalogue.lean}, line 1533}\par\nopagebreak
\begin{leansrc}
inductive PSubdivideData : Hypernet → Sequent → Type where …
\end{leansrc}
\par\medskip\noindent{\small \texttt{PPLC.}\texttt{\bfseries reluLinked}}\hfill{\footnotesize \path{lean/qtsl/Qtsl/Core/Catalogue.lean}, line 1223}\par\nopagebreak
\begin{leansrc}
def reluLinked (H : Hypernet) (preact post : Wire) : Prop
\end{leansrc}
\par\medskip\noindent{\small \texttt{PPLC.PrimApp.}\texttt{\bfseries isReluAt}}\hfill{\footnotesize \path{lean/qtsl/Qtsl/Core/Catalogue.lean}, line 1217}\par\nopagebreak
\begin{leansrc}
def PrimApp.isReluAt (preact post : Wire) : PrimApp → Bool
  | .relu s i o =>
      decide (i = preact) && decide (o = post) &&
      decide (0 < Shape.totalSize s)
  | _ => false

def reluLinked (H : Hypernet) (preact post : Wire) : Prop
\end{leansrc}
\par\medskip\noindent{\small \texttt{PPLC.}\texttt{\bfseries affineWireTerm}}\hfill{\footnotesize \path{lean/qtsl/Qtsl/Core/Catalogue.lean}, line 1159}\par\nopagebreak
\begin{leansrc}
def affineWireTerm (offset slope : Rat) (w : Wire) : Term
\end{leansrc}
\par\medskip\noindent{\small \texttt{PPLC.}\texttt{\bfseries lirpaConclusionForm}}\hfill{\footnotesize \path{lean/qtsl/Qtsl/Core/Catalogue.lean}, line 1162}\par\nopagebreak
\begin{leansrc}
def lirpaConclusionForm
    (preact post : Wire) (α β γ δ : Rat) (F : Formula) : Prop
\end{leansrc}
\par\medskip\noindent{\small \texttt{PPLC.}\texttt{\bfseries reluLinked\_facts}}\hfill{\footnotesize \path{lean/qtsl/Qtsl/Soundness.lean}, line 8057}\par\nopagebreak
\begin{leansrc}
theorem reluLinked_facts {H : Hypernet} {preact post : Wire}
    (hlink : reluLinked H preact post) (σ : H.Valuation)
    (hadm : σ.Admissible) :
    ∃ qp qy : QTensor,
      σ.wireQ preact = some qp ∧ σ.wireQ post = some qy ∧
      qp.shape = qy.shape ∧ 0 < qp.entries.length ∧
      qp.entries.length = qy.entries.length ∧
      ∀ i, i < qp.entries.length →
        (qp.entries.getD i 0 ≤ 0 ∧ qy.entries.getD i 0 = 0) ∨
        (0 ≤ qp.entries.getD i 0 ∧
          qy.entries.getD i 0 = qp.entries.getD i 0)
\end{leansrc}
\par\medskip\noindent{\small \texttt{PPLC.}\texttt{\bfseries sound\_lirpa}}\hfill{\footnotesize \path{lean/qtsl/Qtsl/Soundness.lean}, line 8235}\par\nopagebreak
\begin{leansrc}
theorem sound_lirpa {H : Hypernet}
    {seq lowerPremise upperPremise : Sequent}
    {preact post : Wire} {a b α β γ δ : ℚ}
    (ihl : lowerPremise.semanticallyValid H)
    (ihu : upperPremise.semanticallyValid H)
    (hlink : reluLinked H preact post)
    (hAneg : a < 0) (hBpos : 0 < b)
    (hLower1 : α * a + β ≤ 0) (hLower2 : β ≤ 0) (hLower3 : α * b + β ≤ b)
    (hUpper1 : 0 ≤ γ * a + δ) (hUpper2 : 0 ≤ δ) (hUpper3 : b ≤ γ * b + δ)
    (hpre : seq.pre = lowerPremise.pre)
    (hpreMatch : lowerPremise.pre = upperPremise.pre)
    (hlowerPost : lowerPremise.post = .atom .le (.const a) (.var preact))
    (hupperPost : upperPremise.post = .atom .le (.var preact) (.const b))
    (hpostForm : lirpaConclusionForm preact post α β γ δ seq.post) :
    seq.semanticallyValid H
\end{leansrc}
\par\medskip\noindent{\small \texttt{PPLC.ReluExpansion.}\texttt{\bfseries reluExpH}}\hfill{\footnotesize \path{lean/qtsl/Qtsl/ReluExpansion.lean}, line 156}\par\nopagebreak
\begin{leansrc}
def reluExpH : Hypernet where …
\end{leansrc}
\par\medskip\noindent{\small \texttt{PPLC.ReluExpansion.}\texttt{\bfseries relu\_projection}}\hfill{\footnotesize \path{lean/qtsl/Qtsl/ReluExpansion.lean}, line 462}\par\nopagebreak
\begin{leansrc}
theorem relu_projection (p y : ℚ) :
    (∃ σ : reluExpH.Valuation, σ.Admissible ∧
        σ.entryFlat 0 0 = some (.fin p) ∧
        σ.entryFlat 8 0 = some (.fin y)) ↔
    ((p ≤ 0 ∧ y = 0) ∨ (0 ≤ p ∧ y = p))
\end{leansrc}
\par\medskip\noindent{\small \texttt{PPLC.}\texttt{\bfseries lirpa\_lower\_neg}}\hfill{\footnotesize \path{lean/qtsl/Qtsl/Soundness.lean}, line 8014}\par\nopagebreak
\begin{leansrc}
theorem lirpa_lower_neg {a α β p : ℚ} (ha : a < 0)
    (h1 : α * a + β ≤ 0) (h2 : β ≤ 0) (hap : a ≤ p) (hp0 : p ≤ 0) :
    α * p + β ≤ 0
\end{leansrc}
\par\medskip\noindent{\small \texttt{PPLC.}\texttt{\bfseries lirpa\_lower\_pos}}\hfill{\footnotesize \path{lean/qtsl/Qtsl/Soundness.lean}, line 8025}\par\nopagebreak
\begin{leansrc}
theorem lirpa_lower_pos {b α β p : ℚ} (hb : 0 < b)
    (h2 : β ≤ 0) (h3 : α * b + β ≤ b) (hp0 : 0 ≤ p) (hpb : p ≤ b) :
    α * p + β ≤ p
\end{leansrc}
\par\medskip\noindent{\small \texttt{PPLC.}\texttt{\bfseries lirpa\_upper\_neg}}\hfill{\footnotesize \path{lean/qtsl/Qtsl/Soundness.lean}, line 8036}\par\nopagebreak
\begin{leansrc}
theorem lirpa_upper_neg {a γ δ p : ℚ} (ha : a < 0)
    (h1 : 0 ≤ γ * a + δ) (h2 : 0 ≤ δ) (hap : a ≤ p) (hp0 : p ≤ 0) :
    0 ≤ γ * p + δ
\end{leansrc}
\par\medskip\noindent{\small \texttt{PPLC.}\texttt{\bfseries lirpa\_upper\_pos}}\hfill{\footnotesize \path{lean/qtsl/Qtsl/Soundness.lean}, line 8047}\par\nopagebreak
\begin{leansrc}
theorem lirpa_upper_pos {b γ δ p : ℚ} (hb : 0 < b)
    (h2 : 0 ≤ δ) (h3 : b ≤ γ * b + δ) (hp0 : 0 ≤ p) (hpb : p ≤ b) :
    p ≤ γ * p + δ
\end{leansrc}
\par\medskip\noindent{\small \texttt{PPLC.}\texttt{\bfseries affineWireTerm\_evalQ}}\hfill{\footnotesize \path{lean/qtsl/Qtsl/Soundness.lean}, line 7994}\par\nopagebreak
\begin{leansrc}
theorem affineWireTerm_evalQ {H : Hypernet} (σ : H.Valuation) {w : Wire}
    {qw : QTensor} (h : σ.wireQ w = some qw) (off slope : ℚ) :
    ∃ qa, Term.evalQ σ (affineWireTerm off slope w) = some qa ∧
      qa.shape = qw.shape ∧
      qa.entries.length = qw.entries.length ∧
      ∀ i, i < qw.entries.length →
        qa.entries.getD i 0 = off + slope * qw.entries.getD i 0
\end{leansrc}
\par\medskip\noindent{\small \texttt{PPLC.}\texttt{\bfseries sound\_reluPos}}\hfill{\footnotesize \path{lean/qtsl/Qtsl/Soundness.lean}, line 8156}\par\nopagebreak
\begin{leansrc}
theorem sound_reluPos {H : Hypernet} {seq premise : Sequent}
    {preact post : Wire}
    (ih : premise.semanticallyValid H)
    (hlink : reluLinked H preact post)
    (hpre : seq.pre = premise.pre)
    (hpremPost : premise.post = .atom .ge (.var preact) (.const 0))
    (hpost : seq.post = .atom .eq (.var post) (.var preact)) :
    seq.semanticallyValid H
\end{leansrc}
\par\medskip\noindent{\small \texttt{PPLC.}\texttt{\bfseries PReLUPosData}}\hfill{\footnotesize \path{lean/qtsl/Qtsl/Core/Catalogue.lean}, line 1578}\par\nopagebreak
\begin{leansrc}
inductive PReLUPosData : Hypernet → Sequent → Type where …
\end{leansrc}
\par\medskip\noindent{\small \texttt{PPLC.QTensor.}\texttt{\bfseries rel2\_const\_right\_elim}}\hfill{\footnotesize \path{lean/qtsl/Qtsl/Soundness.lean}, line 7914}\par\nopagebreak
\begin{leansrc}
theorem QTensor.rel2_const_right_elim {r : ℚ → ℚ → Prop} {q : QTensor}
    {c : ℚ} (h : QTensor.rel2 r q ⟨[], [c], rfl⟩) :
    ∀ x ∈ q.entries, r x c
\end{leansrc}
\par\medskip\noindent{\small \texttt{PPLC.}\texttt{\bfseries sound\_reluNeg}}\hfill{\footnotesize \path{lean/qtsl/Qtsl/Soundness.lean}, line 8193}\par\nopagebreak
\begin{leansrc}
theorem sound_reluNeg {H : Hypernet} {seq premise : Sequent}
    {preact post : Wire}
    (ih : premise.semanticallyValid H)
    (hlink : reluLinked H preact post)
    (hpre : seq.pre = premise.pre)
    (hpremPost : premise.post = .atom .le (.var preact) (.const 0))
    (hpost : seq.post = .atom .eq (.var post) (.const 0)) :
    seq.semanticallyValid H
\end{leansrc}
\par\medskip\noindent{\small \texttt{PPLC.}\texttt{\bfseries PReLUNegData}}\hfill{\footnotesize \path{lean/qtsl/Qtsl/Core/Catalogue.lean}, line 1588}\par\nopagebreak
\begin{leansrc}
inductive PReLUNegData : Hypernet → Sequent → Type where …
\end{leansrc}
\par\medskip\noindent{\small \texttt{PPLC.}\texttt{\bfseries PReLUNNData}}\hfill{\footnotesize \path{lean/qtsl/Qtsl/Core/Catalogue.lean}, line 1598}\par\nopagebreak
\begin{leansrc}
inductive PReLUNNData : Hypernet → Sequent → Type where …
\end{leansrc}
\par\medskip\noindent{\small \texttt{PPLC.}\texttt{\bfseries sound\_reluNN}}\hfill{\footnotesize \path{lean/qtsl/Qtsl/Soundness.lean}, line 8116}\par\nopagebreak
\begin{leansrc}
theorem sound_reluNN {H : Hypernet} {seq : Sequent} {preact post : Wire}
    (hlink : reluLinked H preact post)
    (hpost : seq.post = .atom .ge (.var post) (.const 0)) :
    seq.semanticallyValid H
\end{leansrc}
\par\medskip\noindent{\small \texttt{PPLC.}\texttt{\bfseries sound\_pConstExact}}\hfill{\footnotesize \path{lean/qtsl/Qtsl/Soundness.lean}, line 7403}\par\nopagebreak
\begin{leansrc}
theorem sound_pConstExact {H : Hypernet} {seq : Sequent} {e₁ e₂ : Term}
    {b : BowTie}
    (hEval : (Formula.atom b e₁ e₂).evalClosed = some true)
    (hpost : seq.post = .atom b e₁ e₂) :
    seq.semanticallyValid H
\end{leansrc}
\par\medskip\noindent{\small \texttt{PPLC.}\texttt{\bfseries PReLUAboveData}}\hfill{\footnotesize \path{lean/qtsl/Qtsl/Core/Catalogue.lean}, line 1605}\par\nopagebreak
\begin{leansrc}
inductive PReLUAboveData : Hypernet → Sequent → Type where …
\end{leansrc}
\par\medskip\noindent{\small \texttt{PPLC.}\texttt{\bfseries sound\_reluAbove}}\hfill{\footnotesize \path{lean/qtsl/Qtsl/Soundness.lean}, line 8138}\par\nopagebreak
\begin{leansrc}
theorem sound_reluAbove {H : Hypernet} {seq : Sequent} {preact post : Wire}
    (hlink : reluLinked H preact post)
    (hpost : seq.post = .atom .ge (.var post) (.var preact)) :
    seq.semanticallyValid H
\end{leansrc}
\par\medskip\noindent{\small \texttt{PPLC.}\texttt{\bfseries sound\_pBound}}\hfill{\footnotesize \path{lean/qtsl/Qtsl/Soundness.lean}, line 7412}\par\nopagebreak
\begin{leansrc}
theorem sound_pBound {H : Hypernet} {seq : Sequent} {components : List Sequent}
    {combination conclusion : Term}
    (hcomps : ∀ c ∈ components, c.semanticallyValid H)
    (hpreC : ∀ c ∈ components, premiseContains seq.pre c.pre)
    (hshape : ∃ s, conclusion.bshape H = some s)
    (hfin : conclusion.finWires H)
    (hwpos : ∀ w ∈ seq.pre.freeWires ++ conclusion.freeWires,
      ∃ τ, H.wireType w = some τ ∧ 0 < Shape.totalSize τ.shape)
    (hcompPos : ∀ c ∈ components, ∀ w ∈ c.post.freeWires,
      ∃ τ, H.wireType w = some τ ∧ 0 < Shape.totalSize τ.shape)
    (hvalid : syzygyCombinationValid combination components seq.pre conclusion)
    (hform : syzygyConclusionForm conclusion seq.post) :
    seq.semanticallyValid H
\end{leansrc}
\par\medskip\noindent{\small \texttt{PPLC.}\texttt{\bfseries baryLambdas1D}}\hfill{\footnotesize \path{lean/qtsl/Qtsl/Core/Catalogue.lean}, line 449}\par\nopagebreak
\begin{leansrc}
def baryLambdas1D (cx : SimpC) (dContains : Fin cx.nVertices → Bool)
    (t : Term) : Option (List (Fin cx.nVertices × Term))
\end{leansrc}
\par\medskip\noindent{\small \texttt{PPLC.}\texttt{\bfseries MaskRegionStrict}}\hfill{\footnotesize \path{lean/qtsl/Qtsl/Catalogue.lean}, line 386}\par\nopagebreak
\begin{leansrc}
def MaskRegionStrict (C : SimpC) (mask ext : Fin C.nVertices → Bool)
    (x : Fin C.ambDim → ℚ) : Prop
\end{leansrc}
\par\medskip\noindent{\small \texttt{PPLC.SimpC.}\texttt{\bfseries CellIntDisjoint}}\hfill{\footnotesize \path{lean/qtsl/Qtsl/Catalogue.lean}, line 396}\par\nopagebreak
\begin{leansrc}
def SimpC.CellIntDisjoint (C : SimpC) : Prop
\end{leansrc}
\par\medskip\noindent{\small \texttt{PPLC.}\texttt{\bfseries PCellInIntData}}\hfill{\footnotesize \path{lean/qtsl/Qtsl/Core/Catalogue.lean}, line 1612}\par\nopagebreak
\begin{leansrc}
inductive PCellInIntData : Hypernet → Sequent → Type where …
\end{leansrc}
\par\medskip\noindent{\small \texttt{PPLC.}\texttt{\bfseries sound\_cellInInterior}}\hfill{\footnotesize \path{lean/qtsl/Qtsl/Soundness.lean}, line 6531}\par\nopagebreak
\begin{leansrc}
theorem sound_cellInInterior {H : Hypernet} {seq premise : Sequent}
    {schema : SimpC} {input : Wire} {outputs : List Wire}
    {whichI : Fin schema.cells.length}
    (ih : premise.semanticallyValid H)
    (happ : PrimApp.cell schema input outputs ∈ H.applications)
    (hinshape : ∃ τ, H.wireType input = some τ ∧
      τ.shape = [schema.ambDim])
    (houts : ∀ o ∈ outputs, ∃ τ, H.wireType o = some τ ∧
      Shape.totalSize τ.shape = 1)
    (hpre : seq.pre = premise.pre)
    (hint : cellStrictInteriorForm schema (.var input) whichI premise.post)
    (hpost : cellOneHotForm schema outputs whichI seq.post) :
    seq.semanticallyValid H
\end{leansrc}
\par\medskip\noindent{\small \texttt{PPLC.}\texttt{\bfseries cellStrictInteriorBuild}}\hfill{\footnotesize \path{lean/qtsl/Qtsl/Core/Catalogue.lean}, line 1167}\par\nopagebreak
\begin{leansrc}
def cellStrictInteriorBuild (schema : SimpC) (t : Term)
    (k : Fin schema.cells.length) : Option Formula
\end{leansrc}
\par\medskip\noindent{\small \texttt{PPLC.}\texttt{\bfseries cellOneHotBuild}}\hfill{\footnotesize \path{lean/qtsl/Qtsl/Core/Catalogue.lean}, line 1180}\par\nopagebreak
\begin{leansrc}
def cellOneHotBuild (schema : SimpC) (outputs : List Wire)
    (k : Fin schema.cells.length) : Option Formula
\end{leansrc}
\par\medskip\noindent{\small \texttt{PPLC.}\texttt{\bfseries maskRegionStrict\_pair\_intro}}\hfill{\footnotesize \path{lean/qtsl/Qtsl/Soundness.lean}, line 6218}\par\nopagebreak
\begin{leansrc}
theorem maskRegionStrict_pair_intro {C : SimpC}
    {mask ext : Fin C.nVertices → Bool}
    {x : Fin C.ambDim → ℚ} {a b : Fin C.nVertices} {V₀ V₁ : ℚ}
    (hfil : (List.finRange C.nVertices).filter (fun v => mask v) = [a, b])
    (hscope : ∀ v, ext v = true → mask v = true)
    (hVa : ∀ j, C.coord a j = V₀) (hVb : ∀ j, C.coord b j = V₁)
    (hne : V₀ ≠ V₁) (hdim : C.ambDim = 1) (j0 : Fin C.ambDim)
    (hw0t : ext a = true → (V₁ - x j0) / (V₁ - V₀) < 0)
    (hw0f : ext a = false → 0 < (V₁ - x j0) / (V₁ - V₀))
    (hw1t : ext b = true → (x j0 - V₀) / (V₁ - V₀) < 0)
    (hw1f : ext b = false → 0 < (x j0 - V₀) / (V₁ - V₀)) :
    MaskRegionStrict C mask ext x
\end{leansrc}
\par\medskip\noindent{\small \texttt{PPLC.}\texttt{\bfseries holds\_cellOneHot}}\hfill{\footnotesize \path{lean/qtsl/Qtsl/Soundness.lean}, line 6480}\par\nopagebreak
\begin{leansrc}
theorem holds_cellOneHot {H : Hypernet} (σ : H.Valuation) {schema : SimpC}
    {outputs : List Wire} {k : Fin schema.cells.length} {F : Formula}
    (hbuild : cellOneHotForm schema outputs k F)
    (houts : ∀ o ∈ outputs, ∃ τ, H.wireType o = some τ ∧
      Shape.totalSize τ.shape = 1)
    (hsem : ∀ j : Fin outputs.length,
      σ.entryFlat (outputs.get j) 0
        = some (.fin (if j.val = k.val then 1 else 0))) :
    Formula.holds σ F
\end{leansrc}
\par\medskip\noindent{\small \texttt{PPLC.}\texttt{\bfseries cellStrictInteriorBuild\_spec}}\hfill{\footnotesize \path{lean/qtsl/Qtsl/Soundness.lean}, line 6091}\par\nopagebreak
\begin{leansrc}
theorem cellStrictInteriorBuild_spec {schema : SimpC} {t : Term}
    {k : Fin schema.cells.length} {F : Formula}
    (h : cellStrictInteriorForm schema t k F) :
    schema.ambDim = 1 ∧
    ∃ (a b : Fin schema.nVertices) (V₀ V₁ : ℚ),
      (List.finRange schema.nVertices).filter
        (fun v => (schema.cells.get k).1 v) = [a, b] ∧
      (∀ j, schema.coord a j = V₀) ∧
      (∀ j, schema.coord b j = V₁) ∧
      V₀ ≠ V₁ ∧
      F = .and
        (.atom (if (schema.cells.get k).2 a then .lt else .gt)
          (.add (.const (V₁ / (V₁ - V₀))) (.mul (.const (-1 / (V₁ - V₀))) t))
          (.const 0))
        (.atom (if (schema.cells.get k).2 b then .lt else .gt)
          (.add (.const (-V₀ / (V₁ - V₀))) (.mul (.const (1 / (V₁ - V₀))) t))
          (.const 0))
\end{leansrc}
\par\medskip\noindent{\small \texttt{PPLC.PrimApp.}\texttt{\bfseries SemOk}}\hfill{\footnotesize \path{lean/qtsl/Qtsl/Catalogue.lean}, line 445}\par\nopagebreak
\begin{leansrc}
def PrimApp.SemOk {H : Hypernet} (σ : H.Valuation) : PrimApp → Prop
  | .const _ty c out =>
      ∀ i : ℕ, σ.entryFlat out i = c.entryE i
  | .contract δ r inputs output _inputRanks _outputRank pat =>
      ∀ ν₀ ∈ Shape.positions pat.idxShape,
        ∃ eOut eExact,
          σ.entryAt output (pat.outCoords ν₀) = some eOut ∧
          contractExact σ inputs pat (pat.outCoords ν₀) = some eExact ∧
          δ.RoundOkE r eExact eOut
  | .cast _δfrom δto r shape input output =>
      ∀ i : Fin (Shape.totalSize shape),
        ∃ eIn eOut, σ.entryFlat input i.val = some eIn ∧
          σ.entryFlat output i.val = some eOut ∧ δto.RoundOkE r eIn eOut
  | .cell schema input outputs =>

      schema.CellIntDisjoint ∧
      ∃ x : Fin schema.ambDim → ℚ,
        (∀ j : Fin schema.ambDim, σ.entryAt input [j.val] = some (.fin (x j))) ∧
        ∃ k : Fin schema.cells.length,
          MaskRegion schema (schema.cells.get k).1 (schema.cells.get k).2 x ∧
          ∀ j : Fin outputs.length,
            σ.entryFlat (outputs.get j) 0 =
              some (.fin (if j.val = k.val then 1 else 0))
  | .reshape _newShape input output =>
      ∀ i : ℕ, σ.entryFlat output i = σ.entryFlat input i
  | .broadcast toShape input output =>
      ∃ τ, H.wireType input = some τ ∧
        ∀ c ∈ Shape.positions toShape,
          σ.entryAt output c = σ.entryAt input (Shape.bcastBack τ.shape toShape c)
  | .copy _k input outputs =>
      ∀ o ∈ outputs, ∀ i : ℕ, σ.entryFlat o i = σ.entryFlat input i
  | .delete _ => True
  | .swap in1 in2 out1 out2 =>
      (∀ i : ℕ, σ.entryFlat out1 i = σ.entryFlat in2 i) ∧
      (∀ i : ℕ, σ.entryFlat out2 i = σ.entryFlat in1 i)
  | .stack _k shape inputs output =>

      ∀ i : Fin inputs.length, ∀ j : Fin (Shape.totalSize shape),
        σ.entryFlat output (i.val * Shape.totalSize shape + j.val)
          = σ.entryFlat inputs[i] j.val
  …
\end{leansrc}
\par\medskip\noindent{\small \texttt{PPLC.}\texttt{\bfseries PCellOutData}}\hfill{\footnotesize \path{lean/qtsl/Qtsl/Core/Catalogue.lean}, line 1628}\par\nopagebreak
\begin{leansrc}
inductive PCellOutData : Hypernet → Sequent → Type where …
\end{leansrc}
\par\medskip\noindent{\small \texttt{PPLC.}\texttt{\bfseries cellOutZeroBuild}}\hfill{\footnotesize \path{lean/qtsl/Qtsl/Core/Catalogue.lean}, line 1207}\par\nopagebreak
\begin{leansrc}
def cellOutZeroBuild (schema : SimpC) (outputs : List Wire)
    (k : Fin schema.cells.length) : Option Formula
\end{leansrc}
\par\medskip\noindent{\small \texttt{PPLC.}\texttt{\bfseries sound\_cellOut}}\hfill{\footnotesize \path{lean/qtsl/Qtsl/Soundness.lean}, line 6632}\par\nopagebreak
\begin{leansrc}
theorem sound_cellOut {H : Hypernet} {seq premise : Sequent}
    {schema : SimpC} {input : Wire} {outputs : List Wire}
    {whichJ : Fin schema.cells.length} {wv : Fin schema.nVertices}
    (ih : premise.semanticallyValid H)
    (happ : PrimApp.cell schema input outputs ∈ H.applications)
    (hinshape : ∃ τ, H.wireType input = some τ ∧
      τ.shape = [schema.ambDim])
    (houts : ∀ o ∈ outputs, ∃ τ, H.wireType o = some τ ∧
      Shape.totalSize τ.shape = 1)
    (hpre : seq.pre = premise.pre)
    (hwit : cellOutWitnessForm schema (.var input) whichJ wv premise.post)
    (hpost : cellOutZeroForm schema outputs whichJ seq.post) :
    seq.semanticallyValid H
\end{leansrc}
\par\medskip\noindent{\small \texttt{PPLC.}\texttt{\bfseries cellOutWitnessBuild}}\hfill{\footnotesize \path{lean/qtsl/Qtsl/Core/Catalogue.lean}, line 1193}\par\nopagebreak
\begin{leansrc}
def cellOutWitnessBuild (schema : SimpC) (t : Term)
    (k : Fin schema.cells.length) (w : Fin schema.nVertices) :
    Option Formula
\end{leansrc}
\par\medskip\noindent{\small \texttt{PPLC.}\texttt{\bfseries cellOutWitnessBuild\_spec}}\hfill{\footnotesize \path{lean/qtsl/Qtsl/Soundness.lean}, line 6145}\par\nopagebreak
\begin{leansrc}
theorem cellOutWitnessBuild_spec {schema : SimpC} {t : Term}
    {k : Fin schema.cells.length} {w : Fin schema.nVertices} {F : Formula}
    (h : cellOutWitnessForm schema t k w F) :
    schema.ambDim = 1 ∧
    (schema.cells.get k).1 w = true ∧
    ∃ (a b : Fin schema.nVertices) (V₀ V₁ : ℚ),
      (List.finRange schema.nVertices).filter
        (fun v => (schema.cells.get k).1 v) = [a, b] ∧
      (∀ j, schema.coord a j = V₀) ∧
      (∀ j, schema.coord b j = V₁) ∧
      V₀ ≠ V₁ ∧
      ((w = a ∧ F = .atom (if (schema.cells.get k).2 a then .gt else .lt)
          (.add (.const (V₁ / (V₁ - V₀))) (.mul (.const (-1 / (V₁ - V₀))) t))
          (.const 0)) ∨
       (w = b ∧ F = .atom (if (schema.cells.get k).2 b then .gt else .lt)
          (.add (.const (-V₀ / (V₁ - V₀))) (.mul (.const (1 / (V₁ - V₀))) t))
          (.const 0)))
\end{leansrc}
\par\medskip\noindent{\small \texttt{PPLC.Formula.}\texttt{\bfseries evalClosed}}\hfill{\footnotesize \path{lean/qtsl/Qtsl/Core/Catalogue.lean}, line 441}\par\nopagebreak
\begin{leansrc}
def Formula.evalClosed : Formula → Option Bool
  | .atom b l r => do pure (b.evalQ (← l.evalQ?) (← r.evalQ?))
  | .inRegion _ _ _ => none
  | .and l r => do pure ((← l.evalClosed) && (← r.evalClosed))
  | .or l r => do pure ((← l.evalClosed) || (← r.evalClosed))
  | .tt => some true
  | .ff => some false

def baryLambdas1D (cx : SimpC) (dContains : Fin cx.nVertices → Bool)
    (t : Term) : Option (List (Fin cx.nVertices × Term))
\end{leansrc}
\par\medskip\noindent{\small \texttt{PPLC.Formula.}\texttt{\bfseries holds\_of\_evalClosed\_true}}\hfill{\footnotesize \path{lean/qtsl/Qtsl/Soundness.lean}, line 2728}\par\nopagebreak
\begin{leansrc}
theorem Formula.holds_of_evalClosed_true {H : Hypernet} (σ : H.Valuation) :
    ∀ (F : Formula), F.evalClosed = some true → F.holds σ
\end{leansrc}
\par\medskip\noindent{\small \texttt{PPLC.Formula.}\texttt{\bfseries holds\_of\_evalClosed\_true}}\hfill{\footnotesize \path{lean/qtsl/Qtsl/Soundness.lean}, line 2728}\par\nopagebreak
\begin{leansrc}
theorem Formula.holds_of_evalClosed_true {H : Hypernet} (σ : H.Valuation) :
    ∀ (F : Formula), F.evalClosed = some true → F.holds σ
\end{leansrc}
\par\medskip\noindent{\small \texttt{PPLC.Formula.}\texttt{\bfseries evalClosed}}\hfill{\footnotesize \path{lean/qtsl/Qtsl/Core/Catalogue.lean}, line 441}\par\nopagebreak
\begin{leansrc}
def Formula.evalClosed : Formula → Option Bool
  | .atom b l r => do pure (b.evalQ (← l.evalQ?) (← r.evalQ?))
  | .inRegion _ _ _ => none
  | .and l r => do pure ((← l.evalClosed) && (← r.evalClosed))
  | .or l r => do pure ((← l.evalClosed) || (← r.evalClosed))
  | .tt => some true
  | .ff => some false

def baryLambdas1D (cx : SimpC) (dContains : Fin cx.nVertices → Bool)
    (t : Term) : Option (List (Fin cx.nVertices × Term))
\end{leansrc}
\par\medskip\noindent{\small \texttt{PPLC.}\texttt{\bfseries SyzygyStep}}\hfill{\footnotesize \path{lean/qtsl/Qtsl/Soundness.lean}, line 7150}\par\nopagebreak
\begin{leansrc}
structure SyzygyStep (H : Hypernet) where …
\end{leansrc}
\par\medskip\noindent{\small \texttt{PPLC.}\texttt{\bfseries sound\_syzygy}}\hfill{\footnotesize \path{lean/qtsl/Qtsl/Soundness.lean}, line 7205}\par\nopagebreak
\begin{leansrc}
theorem sound_syzygy {H : Hypernet} (st : SyzygyStep H)
    (hcomps : ∀ c ∈ st.components, c.semanticallyValid H) :
    st.seq.semanticallyValid H
\end{leansrc}
\par\medskip\noindent{\small \texttt{PPLC.}\texttt{\bfseries syzygyConclusionForm}}\hfill{\footnotesize \path{lean/qtsl/Qtsl/Core/Catalogue.lean}, line 559}\par\nopagebreak
\begin{leansrc}
def syzygyConclusionForm (conclusion : Term) (F : Formula) : Prop
\end{leansrc}
\par\medskip\noindent{\small \texttt{PPLC.}\texttt{\bfseries syzygyCheck}}\hfill{\footnotesize \path{lean/qtsl/Qtsl/Core/Catalogue.lean}, line 1026}\par\nopagebreak
\begin{leansrc}
def syzygyCheck (cc : Term) (components : List Sequent) (premise : Formula)
    (conclusion : Term) : Bool
\end{leansrc}
\par\medskip\noindent{\small \texttt{PPLC.}\texttt{\bfseries PBoundData}}\hfill{\footnotesize \path{lean/qtsl/Qtsl/Core/Catalogue.lean}, line 1690}\par\nopagebreak
\begin{leansrc}
inductive PBoundData : Hypernet → Sequent → Type where …
\end{leansrc}
\par\medskip\noindent{\small \texttt{PPLC.}\texttt{\bfseries syzygyCombinationValid}}\hfill{\footnotesize \path{lean/qtsl/Qtsl/Core/Catalogue.lean}, line 1032}\par\nopagebreak
\begin{leansrc}
def syzygyCombinationValid (cc : Term) (components : List Sequent)
    (premise : Formula) (conclusion : Term) : Prop
\end{leansrc}
\par\medskip\noindent{\small \texttt{PPLC.}\texttt{\bfseries holds\_of\_syzygyConclusionForm}}\hfill{\footnotesize \path{lean/qtsl/Qtsl/Soundness.lean}, line 7165}\par\nopagebreak
\begin{leansrc}
theorem holds_of_syzygyConclusionForm {H : Hypernet} {σ : H.Valuation}
    {conclusion : Term} {F : Formula}
    (hle : Formula.holds σ (.atom .le conclusion (.const 0)))
    (hform : syzygyConclusionForm conclusion F) :
    Formula.holds σ F
\end{leansrc}
\par\medskip\noindent{\small \texttt{PPLC.PsdCert.}\texttt{\bfseries conclCheck}}\hfill{\footnotesize \path{lean/qtsl/Qtsl/Core/Catalogue.lean}, line 1105}\par\nopagebreak
\begin{leansrc}
def conclCheck (n : Nat) (M : Fin n → Fin n → Rat) (ts : List Term)
    (conclusion : Term) : Bool
\end{leansrc}
\par\medskip\noindent{\small \texttt{PPLC.Term.}\texttt{\bfseries toPoly}}\hfill{\footnotesize \path{lean/qtsl/Qtsl/Core/Catalogue.lean}, line 414}\par\nopagebreak
\begin{leansrc}
def Term.toPoly : Term → QArith.Poly
  | .var w => QArith.Poly.var w
  | .const c => QArith.Poly.const c
  | .add l r => QArith.Poly.add l.toPoly r.toPoly
  | .mul l r => QArith.Poly.mul l.toPoly r.toPoly
  | .sub l r => QArith.Poly.sub l.toPoly r.toPoly
  | .neg t => QArith.Poly.neg t.toPoly

  | .idx _ _ => QArith.Poly.const 0

def Term.evalQ? : Term → Option Rat
  | .var _ => none
  | .const c => some c
  | .add l r => do pure ((← l.evalQ?) + (← r.evalQ?))
  | .mul l r => do pure ((← l.evalQ?) * (← r.evalQ?))
  | .sub l r => do pure ((← l.evalQ?) - (← r.evalQ?))
  | .neg t => do pure (-(← t.evalQ?))
  | .idx _ _ => none

def BowTie.evalQ : BowTie → Rat → Rat → Bool
  | .lt, x, y => decide (x < y)
  | .le, x, y => decide (x ≤ y)
  | .eq, x, y => decide (x = y)
  | .ge, x, y => decide (y ≤ x)
  | .gt, x, y => decide (y < x)
  | .ne, x, y => decide (¬ x = y)

def Formula.evalClosed : Formula → Option Bool
  | .atom b l r => do pure (b.evalQ (← l.evalQ?) (← r.evalQ?))
  | .inRegion _ _ _ => none
  | .and l r => do pure ((← l.evalClosed) && (← r.evalClosed))
  | .or l r => do pure ((← l.evalClosed) || (← r.evalClosed))
  | .tt => some true
  | .ff => some false

def baryLambdas1D (cx : SimpC) (dContains : Fin cx.nVertices → Bool)
    (t : Term) : Option (List (Fin cx.nVertices × Term))
\end{leansrc}
\par\medskip\noindent{\small \texttt{PPLC.PsdCert.}\texttt{\bfseries conclCheckCore}}\hfill{\footnotesize \path{lean/qtsl/Qtsl/Core/Catalogue.lean}, line 1092}\par\nopagebreak
\begin{leansrc}
def conclCheckCore (n : Nat) (M : Fin n → Fin n → Rat) (ts : List Term)
    (conclusion : Term) : Bool
\end{leansrc}
\par\medskip\noindent{\small \texttt{PPLC.PsdCert.}\texttt{\bfseries conclCheckW}}\hfill{\footnotesize \path{lean/qtsl/Qtsl/Core/Catalogue.lean}, line 1098}\par\nopagebreak
\begin{leansrc}
def conclCheckW (n : Nat) (M : Fin n → Fin n → Rat) (ts : List Term)
    (conclusion : Term) : Bool
\end{leansrc}
\par\medskip\noindent{\small \texttt{PPLC.}\texttt{\bfseries sound\_psdCertOf}}\hfill{\footnotesize \path{lean/qtsl/Qtsl/Soundness.lean}, line 7466}\par\nopagebreak
\begin{leansrc}
theorem sound_psdCertOf {H : Hypernet} {seq : Sequent} {a : Nat}
    {M : Fin a → Fin a → ℚ} {ts : List Term} {conclusion : Term}
    (hpsd : ∀ x : Fin a → ℚ, 0 ≤ PsdCert.finQSum a fun i =>
      PsdCert.finQSum a fun j => M i j * x i * x j)
    (hconcl : PsdCert.conclCheck a M ts conclusion = true)
    (hshape : ∃ s, conclusion.bshape H = some s)
    (hfin : conclusion.finWires H)
    (hwpos : ∀ w ∈ conclusion.freeWires,
      ∃ τ, H.wireType w = some τ ∧ 0 < Shape.totalSize τ.shape)
    (hform : syzygyConclusionForm conclusion seq.post) :
    seq.semanticallyValid H
\end{leansrc}
\par\medskip\noindent{\small \texttt{PPLC.PsdCert.}\texttt{\bfseries congruence\_eq\_diagonal}}\hfill{\footnotesize \path{lean/qtsl/Qtsl/PsdCert.lean}, line 41}\par\nopagebreak
\begin{leansrc}
theorem congruence_eq_diagonal {A Z : Matrix n n ℝ}
    (hA : ∀ i j, A i j = A j i)
    (hZu : ∀ ⦃i j⦄, j < i → Z i j = 0)
    (hMl : ∀ ⦃i j⦄, i < j → (A * Z) i j = 0) :
    Zᴴ * (A * Z) = diagonal fun i => Z i i * (A * Z) i i
\end{leansrc}
\par\medskip\noindent{\small \texttt{PPLC.PsdCert.}\texttt{\bfseries posDef\_of\_congruence\_posDef}}\hfill{\footnotesize \path{lean/qtsl/Qtsl/PsdCert.lean}, line 77}\par\nopagebreak
\begin{leansrc}
theorem posDef_of_congruence_posDef {A Z : Matrix n n ℝ}
    (hZu : ∀ ⦃i j⦄, j < i → Z i j = 0)
    (hZd : ∀ i, 0 < Z i i)
    (hBd : (Zᴴ * (A * Z)).PosDef) :
    A.PosDef
\end{leansrc}
\par\medskip\noindent{\small \texttt{PPLC.PsdCert.}\texttt{\bfseries posDef\_of\_triangular\_witness}}\hfill{\footnotesize \path{lean/qtsl/Qtsl/PsdCert.lean}, line 98}\par\nopagebreak
\begin{leansrc}
theorem posDef_of_triangular_witness {A Z : Matrix n n ℝ}
    (hA : ∀ i j, A i j = A j i)
    (hZu : ∀ ⦃i j⦄, j < i → Z i j = 0)
    (hZd : ∀ i, 0 < Z i i)
    (hMl : ∀ ⦃i j⦄, i < j → (A * Z) i j = 0)
    (hMd : ∀ i, 0 < (A * Z) i i) :
    A.PosDef
\end{leansrc}
\par\medskip\noindent{\small \texttt{PPLC.PsdCert.}\texttt{\bfseries quadForm\_nonneg\_of\_posSemidef}}\hfill{\footnotesize \path{lean/qtsl/Qtsl/PsdCert.lean}, line 431}\par\nopagebreak
\begin{leansrc}
theorem quadForm_nonneg_of_posSemidef {a : Nat} {M : Fin a → Fin a → ℚ}
    (h : (matR M).PosSemidef) (x : Fin a → ℚ) :
    0 ≤ finQSum a (fun i => finQSum a (fun j => M i j * x i * x j))
\end{leansrc}
\par\medskip\noindent{\small \texttt{PPLC.PsdCert.}\texttt{\bfseries symB}}\hfill{\footnotesize \path{lean/qtsl/Qtsl/Core/Catalogue.lean}, line 1047}\par\nopagebreak
\begin{leansrc}
def symB (n : Nat) (A : Fin n → Fin n → Rat) : Bool
\end{leansrc}
\par\medskip\noindent{\small \texttt{PPLC.PsdCert.}\texttt{\bfseries upperTriPosDiagB}}\hfill{\footnotesize \path{lean/qtsl/Qtsl/Core/Catalogue.lean}, line 1051}\par\nopagebreak
\begin{leansrc}
def upperTriPosDiagB (n : Nat) (Z : Fin n → Fin n → Rat) : Bool
\end{leansrc}
\par\medskip\noindent{\small \texttt{PPLC.PsdCert.}\texttt{\bfseries mulEntry}}\hfill{\footnotesize \path{lean/qtsl/Qtsl/Core/Catalogue.lean}, line 1044}\par\nopagebreak
\begin{leansrc}
def mulEntry (n : Nat) (A B : Fin n → Fin n → Rat) (i j : Fin n) : Rat
\end{leansrc}
\par\medskip\noindent{\small \texttt{PPLC.PsdCert.}\texttt{\bfseries conclCheck}}\hfill{\footnotesize \path{lean/qtsl/Qtsl/Core/Catalogue.lean}, line 1105}\par\nopagebreak
\begin{leansrc}
def conclCheck (n : Nat) (M : Fin n → Fin n → Rat) (ts : List Term)
    (conclusion : Term) : Bool
\end{leansrc}
\par\medskip\noindent{\small \texttt{PPLC.PsdCert.}\texttt{\bfseries negQuadFormPoly}}\hfill{\footnotesize \path{lean/qtsl/Qtsl/Core/Catalogue.lean}, line 1085}\par\nopagebreak
\begin{leansrc}
def negQuadFormPoly (n : Nat) (M : Fin n → Fin n → Rat)
    (ps : List QArith.Poly) : QArith.Poly
\end{leansrc}
\par\medskip\noindent{\small \texttt{PPLC.PsdCert.}\texttt{\bfseries posDef\_of\_dominant}}\hfill{\footnotesize \path{lean/qtsl/Qtsl/PsdCert.lean}, line 126}\par\nopagebreak
\begin{leansrc}
theorem posDef_of_dominant [DecidableEq n] {B : Matrix n n ℝ}
    (hsym : ∀ i j, B i j = B j i)
    (hdom : ∀ i, (∑ j ∈ Finset.univ.erase i, |B i j|) < B i i) :
    B.PosDef
\end{leansrc}
\par\medskip\noindent{\small \texttt{PPLC.PsdCert.}\texttt{\bfseries posDef\_of\_dd\_witness}}\hfill{\footnotesize \path{lean/qtsl/Qtsl/PsdCert.lean}, line 213}\par\nopagebreak
\begin{leansrc}
theorem posDef_of_dd_witness [DecidableEq n] {A Z : Matrix n n ℝ}
    (hA : ∀ i j, A i j = A j i)
    (hZu : ∀ ⦃i j⦄, j < i → Z i j = 0)
    (hZd : ∀ i, 0 < Z i i)
    (hdom : ∀ i, (∑ j ∈ Finset.univ.erase i, |(Zᴴ * (A * Z)) i j|)
              < (Zᴴ * (A * Z)) i i) :
    A.PosDef
\end{leansrc}
\par\medskip\noindent{\small \texttt{PPLC.PsdCert.}\texttt{\bfseries finQSum}}\hfill{\footnotesize \path{lean/qtsl/Qtsl/Core/Catalogue.lean}, line 1042}\par\nopagebreak
\begin{leansrc}
def finQSum (n : Nat) (f : Fin n → Rat) : Rat
\end{leansrc}
\par\medskip\noindent{\small \texttt{PPLC.PsdCert.}\texttt{\bfseries posDef\_of\_dominant}}\hfill{\footnotesize \path{lean/qtsl/Qtsl/PsdCert.lean}, line 126}\par\nopagebreak
\begin{leansrc}
theorem posDef_of_dominant [DecidableEq n] {B : Matrix n n ℝ}
    (hsym : ∀ i j, B i j = B j i)
    (hdom : ∀ i, (∑ j ∈ Finset.univ.erase i, |B i j|) < B i i) :
    B.PosDef
\end{leansrc}
\par\medskip\noindent{\small \texttt{PPLC.PsdCert.}\texttt{\bfseries congEntry}}\hfill{\footnotesize \path{lean/qtsl/Qtsl/Core/Catalogue.lean}, line 1062}\par\nopagebreak
\begin{leansrc}
def congEntry (n : Nat) (A Z : Fin n → Fin n → Rat) (i j : Fin n) : Rat
\end{leansrc}
\par\medskip\noindent{\small \texttt{PPLC.PsdCert.}\texttt{\bfseries posSemidef\_of\_factor}}\hfill{\footnotesize \path{lean/qtsl/Qtsl/PsdCert.lean}, line 233}\par\nopagebreak
\begin{leansrc}
theorem posSemidef_of_factor {m' n' : Type*} [Fintype m'] [Fintype n']
    {A : Matrix n' n' ℝ} {B : Matrix m' m' ℝ} {L : Matrix m' n' ℝ}
    (hB : B.PosSemidef) (hA : A = Lᴴ * B * L) : A.PosSemidef
\end{leansrc}
\par\medskip\noindent{\small \texttt{PPLC.PsdCert.}\texttt{\bfseries kernelExt}}\hfill{\footnotesize \path{lean/qtsl/Qtsl/Core/Catalogue.lean}, line 1071}\par\nopagebreak
\begin{leansrc}
def kernelExt (m k : Nat) (N : Fin m → Fin k → Rat) (a : Fin m)
    (j : Fin (m + k)) : Rat
\end{leansrc}
\par\medskip\noindent{\small \texttt{PPLC.PsdCert.}\texttt{\bfseries posSemidef\_of\_factor}}\hfill{\footnotesize \path{lean/qtsl/Qtsl/PsdCert.lean}, line 233}\par\nopagebreak
\begin{leansrc}
theorem posSemidef_of_factor {m' n' : Type*} [Fintype m'] [Fintype n']
    {A : Matrix n' n' ℝ} {B : Matrix m' m' ℝ} {L : Matrix m' n' ℝ}
    (hB : B.PosSemidef) (hA : A = Lᴴ * B * L) : A.PosSemidef
\end{leansrc}
\par\medskip\noindent{\small \texttt{PPLC.}\texttt{\bfseries sound\_reluSplit}}\hfill{\footnotesize \path{lean/qtsl/Qtsl/ReluSplit.lean}, line 15}\par\nopagebreak
\begin{leansrc}
theorem sound_reluSplit {H : Hypernet}
    {seq activePremise inactivePremise : Sequent}
    {preact post marker : Wire} {bodyForm : Formula}
    (ihActive : activePremise.semanticallyValid H)
    (ihInactive : inactivePremise.semanticallyValid H)
    (hlink : reluLinked H preact post)
    (hrank0 : ∃ τ, H.wireType post = some τ ∧ τ.shape = [])
    (hActPre : activePremise.pre
      = .and seq.pre (.atom .ge (.var preact) (.const 0)))
    (hInactPre : inactivePremise.pre
      = .and seq.pre (.atom .le (.var preact) (.const 0)))
    (hActPost : substituteInFormula bodyForm marker (.var preact)
      activePremise.post)
    (hInactPost : substituteInFormula bodyForm marker (.const 0)
      inactivePremise.post)
    (hconcl : substituteInFormula bodyForm marker (.var post) seq.post) :
    seq.semanticallyValid H
\end{leansrc}
\par\medskip\noindent{\small \texttt{PPLC.}\texttt{\bfseries reluSplit\_rank0\_necessary}}\hfill{\footnotesize \path{lean/qtsl/Qtsl/ReluSplit.lean}, line 153}\par\nopagebreak
\begin{leansrc}
theorem reluSplit_rank0_necessary :
    ¬ (∀ (H : Hypernet) (seq activePremise inactivePremise : Sequent)
        (preact post marker : Wire) (bodyForm : Formula),
        activePremise.semanticallyValid H →
        inactivePremise.semanticallyValid H →
        reluLinked H preact post →
        activePremise.pre
          = .and seq.pre (.atom .ge (.var preact) (.const 0)) →
        inactivePremise.pre
          = .and seq.pre (.atom .le (.var preact) (.const 0)) →
        substituteInFormula bodyForm marker (.var preact)
          activePremise.post →
        substituteInFormula bodyForm marker (.const 0)
          inactivePremise.post →
        substituteInFormula bodyForm marker (.var post) seq.post →
        seq.semanticallyValid H)
\end{leansrc}
\par\medskip\noindent{\small \texttt{PPLC.ReluSplitScope.}\texttt{\bfseries HSplit}}\hfill{\footnotesize \path{lean/qtsl/Qtsl/ReluSplit.lean}, line 95}\par\nopagebreak
\begin{leansrc}
def HSplit : Hypernet where …
\end{leansrc}
\par\medskip\noindent{\small \texttt{PPLC.}\texttt{\bfseries PromiseStep}}\hfill{\footnotesize \path{lean/qtsl/Qtsl/Soundness.lean}, line 7116}\par\nopagebreak
\begin{leansrc}
structure PromiseStep where …
\end{leansrc}
\par\medskip\noindent{\small \texttt{PPLC.}\texttt{\bfseries sound\_promise}}\hfill{\footnotesize \path{lean/qtsl/Qtsl/Soundness.lean}, line 7132}\par\nopagebreak
\begin{leansrc}
theorem sound_promise {H : Hypernet} (st : PromiseStep)
    (ihs : st.supplier.semanticallyValid H)
    (ihc : st.consumer.semanticallyValid H) :
    st.seq.semanticallyValid H
\end{leansrc}
\par\medskip\noindent{\small \texttt{PPLC.Formula.}\texttt{\bfseries conjGet\_of\_conjReplace}}\hfill{\footnotesize \path{lean/qtsl/Qtsl/Soundness.lean}, line 6994}\par\nopagebreak
\begin{leansrc}
theorem Formula.conjGet_of_conjReplace :
    ∀ {π : ConjPath} {Γ Γ' g : Formula},
    Γ.conjReplace? π g = some Γ' → Γ'.conjGet? π = some g
\end{leansrc}
\par\medskip\noindent{\small \texttt{PPLC.Formula.}\texttt{\bfseries holds\_of\_conjGet}}\hfill{\footnotesize \path{lean/qtsl/Qtsl/Soundness.lean}, line 6916}\par\nopagebreak
\begin{leansrc}
theorem Formula.holds_of_conjGet {H : Hypernet} {σ : H.Valuation} :
    ∀ {π : ConjPath} {Γ d : Formula},
    Γ.conjGet? π = some d → Γ.holds σ → d.holds σ
\end{leansrc}
\par\medskip\noindent{\small \texttt{PPLC.Formula.}\texttt{\bfseries conjReplace\_back}}\hfill{\footnotesize \path{lean/qtsl/Qtsl/Soundness.lean}, line 7042}\par\nopagebreak
\begin{leansrc}
theorem Formula.conjReplace_back :
    ∀ {π : ConjPath} {Γ Γ' q g : Formula},
    Γ.conjGet? π = some q → Γ.conjReplace? π g = some Γ' →
    Γ'.conjReplace? π q = some Γ
\end{leansrc}
\par\medskip\noindent{\small \texttt{PPLC.Formula.}\texttt{\bfseries holds\_conjReplace}}\hfill{\footnotesize \path{lean/qtsl/Qtsl/Soundness.lean}, line 6949}\par\nopagebreak
\begin{leansrc}
theorem Formula.holds_conjReplace {H : Hypernet} {σ : H.Valuation} :
    ∀ {π : ConjPath} {Γ Γ' g : Formula},
    Γ.conjReplace? π g = some Γ' → Γ.holds σ → g.holds σ → Γ'.holds σ
\end{leansrc}
\par\medskip\noindent{\small \texttt{PPLC.}\texttt{\bfseries formulaWiresWithin}}\hfill{\footnotesize \path{lean/qtsl/Qtsl/Core/Catalogue.lean}, line 953}\par\nopagebreak
\begin{leansrc}
def formulaWiresWithin (f : Formula) (iface : List Wire) : Prop
\end{leansrc}
\par\medskip\noindent{\small \texttt{PPLC.}\texttt{\bfseries wiresSeparated}}\hfill{\footnotesize \path{lean/qtsl/Qtsl/Core/Catalogue.lean}, line 956}\par\nopagebreak
\begin{leansrc}
def wiresSeparated (left ctx post : Formula) (iface : List Wire) : Prop
\end{leansrc}
\par\medskip\noindent{\small \texttt{PPLC.}\texttt{\bfseries syzygyCheckCore}}\hfill{\footnotesize \path{lean/qtsl/Qtsl/Core/Catalogue.lean}, line 964}\par\nopagebreak
\begin{leansrc}
def syzygyCheckCore (cc : Term) (components : List Sequent) (premise : Formula)
    (conclusion : Term) : Bool
\end{leansrc}
\par\medskip\noindent{\small \texttt{PPLC.}\texttt{\bfseries syzygyCheckW}}\hfill{\footnotesize \path{lean/qtsl/Qtsl/Core/Catalogue.lean}, line 982}\par\nopagebreak
\begin{leansrc}
def syzygyCheckW (cc : Term) (components : List Sequent) (premise : Formula)
    (conclusion : Term) : Bool
\end{leansrc}
\par\medskip\noindent{\small \texttt{PPLC.}\texttt{\bfseries syzygyCheckM}}\hfill{\footnotesize \path{lean/qtsl/Qtsl/Core/Catalogue.lean}, line 1004}\par\nopagebreak
\begin{leansrc}
def syzygyCheckM (cc : Term) (components : List Sequent) (premise : Formula)
    (conclusion : Term) : Bool
\end{leansrc}
\par\medskip\noindent{\small \texttt{PPLC.}\texttt{\bfseries premiseContains}}\hfill{\footnotesize \path{lean/qtsl/Qtsl/Core/Catalogue.lean}, line 950}\par\nopagebreak
\begin{leansrc}
def premiseContains (pre sub : Formula) : Prop
\end{leansrc}
\par\medskip\noindent{\small \texttt{PPLC.}\texttt{\bfseries atomLeZeroPolys\_sound\_entryC}}\hfill{\footnotesize \path{lean/qtsl/Qtsl/Soundness.lean}, line 3380}\par\nopagebreak
\begin{leansrc}
theorem atomLeZeroPolys_sound_entryC {H : Hypernet} (σ : H.Valuation)
    {F : Formula} (hF : Formula.holds σ F)
    (hw : ∀ w ∈ F.freeWires, ∃ τ, H.wireType w = some τ ∧
      0 < Shape.totalSize τ.shape)
    (hix : F.atomSidesIdxFree = true) (i : ℕ) :
    ∀ pa ∈ atomLeZeroPolys F, QArith.Poly.eval (σ.entryValC i) pa ≤ 0
\end{leansrc}
\par\medskip\noindent{\small \texttt{PPLC.}\texttt{\bfseries wire\_discipline\_of\_bshape\_fin}}\hfill{\footnotesize \path{lean/qtsl/Qtsl/Soundness.lean}, line 2265}\par\nopagebreak
\begin{leansrc}
theorem wire_discipline_of_bshape_fin {H : Hypernet} (σ : H.Valuation)
    {t : Term} {s : Shape} (hb : t.bshape H = some s)
    (hf : t.finWires H) (hix : t.idxFree = true) :
    ∀ w ∈ t.freeWires, ∃ qw, σ.wireQ w = some qw ∧
      (qw.shape = [] ∨ qw.shape = s)
\end{leansrc}
\par\medskip\noindent{\small \texttt{PPLC.Term.}\texttt{\bfseries evalQ\_total\_singleton}}\hfill{\footnotesize \path{lean/qtsl/Qtsl/Soundness.lean}, line 1009}\par\nopagebreak
\begin{leansrc}
theorem Term.evalQ_total_singleton {H : Hypernet} (σ : H.Valuation) :
    ∀ (t : Term),
    (∀ w ∈ t.freeWires, ∃ qw, σ.wireQ w = some qw ∧ SingletonShape qw.shape) →
    t.idxFree = true →
    ∃ q, Term.evalQ σ t = some q ∧ SingletonShape q.shape ∧
      q.scalar = QArith.Poly.eval σ.rho t.toPoly
\end{leansrc}
\par\medskip\noindent{\small \texttt{PPLC.SmallGainStep.}\texttt{\bfseries toSyzygy}}\hfill{\footnotesize \path{lean/qtsl/Qtsl/Composition.lean}, line 409}\par\nopagebreak
\begin{leansrc}
def toSyzygy (st : SmallGainStep H) : SyzygyStep H where …
\end{leansrc}
\par\medskip\noindent{\small \texttt{PPLC.}\texttt{\bfseries sound\_smallGain}}\hfill{\footnotesize \path{lean/qtsl/Qtsl/Composition.lean}, line 435}\par\nopagebreak
\begin{leansrc}
theorem sound_smallGain {H : Hypernet} (st : SmallGainStep H)
    (h1 : st.diss1.semanticallyValid H)
    (h2 : st.diss2.semanticallyValid H)
    (hm1 : st.mem1.semanticallyValid H)
    (hm2 : st.mem2.semanticallyValid H) :
    st.seq.semanticallyValid H
\end{leansrc}
\par\medskip\noindent{\small \texttt{PPLC.SmallGainStep.}\texttt{\bfseries toSyzygy}}\hfill{\footnotesize \path{lean/qtsl/Qtsl/Composition.lean}, line 409}\par\nopagebreak
\begin{leansrc}
def toSyzygy (st : SmallGainStep H) : SyzygyStep H where …
\end{leansrc}
\par\medskip\noindent{\small \texttt{PPLC.}\texttt{\bfseries RuleTag}}\hfill{\footnotesize \path{lean/qtsl/Qtsl/Core/Catalogue.lean}, line 1291}\par\nopagebreak
\begin{leansrc}
inductive RuleTag where …
\end{leansrc}
\par\medskip\noindent{\small \texttt{PPLC.}\texttt{\bfseries SmallGainStep}}\hfill{\footnotesize \path{lean/qtsl/Qtsl/Composition.lean}, line 253}\par\nopagebreak
\begin{leansrc}
structure SmallGainStep (H : Hypernet) where …
\end{leansrc}
\par\medskip\noindent{\small \texttt{PPLC.SmallGainStep.}\texttt{\bfseries combTerm}}\hfill{\footnotesize \path{lean/qtsl/Qtsl/Composition.lean}, line 322}\par\nopagebreak
\begin{leansrc}
def combTerm (st : SmallGainStep H) : Term
\end{leansrc}
\par\medskip\noindent{\small \texttt{PPLC.SmallGainStep.}\texttt{\bfseries combValid}}\hfill{\footnotesize \path{lean/qtsl/Qtsl/Composition.lean}, line 338}\par\nopagebreak
\begin{leansrc}
theorem combValid (st : SmallGainStep H) :
    syzygyCombinationValid st.combTerm st.components st.seq.pre st.concl
\end{leansrc}
\par\medskip\noindent{\small \texttt{PPLC.SmallGainStep.}\texttt{\bfseries toSyzygy\_components}}\hfill{\footnotesize \path{lean/qtsl/Qtsl/Composition.lean}, line 430}\par\nopagebreak
\begin{leansrc}
theorem toSyzygy_components (st : SmallGainStep H) :
    st.toSyzygy.components = st.components
\end{leansrc}
\par\medskip\noindent{\small \texttt{PPLC.}\texttt{\bfseries smallGainN\_is\_syzygy}}\hfill{\footnotesize \path{lean/qtsl/Qtsl/Composition.lean}, line 582}\par\nopagebreak
\begin{leansrc}
theorem smallGainN_is_syzygy (kk : ℕ) (LV A Sg : Fin kk → Poly)
    (lam : Fin kk → ℚ) (B : Fin kk → Fin kk → ℚ) (th : ℚ)
    (hcol : ∀ j, ((List.finRange kk).map (fun i => B i j)).sum = lam j) :
    (((List.finRange kk).map
        (fun i => (lam i, Poly.sub (Poly.add (LV i) (A i)) (Sg i)))
      ++ (List.finRange kk).map
        (fun i => ((1 : ℚ), smallGainNMemPoly kk A Sg lam B th i))).foldr
      (fun pr acc => Poly.add (Poly.scale pr.1 pr.2) acc) [])
    = Poly.add
        ((List.finRange kk).foldr
          (fun i acc => Poly.add (Poly.scale (lam i) (LV i)) acc) [])
        (Poly.scale th
          ((List.finRange kk).foldr
            (fun i acc => Poly.add (Poly.scale (lam i) (A i)) acc) []))
\end{leansrc}
\par\medskip\noindent{\small \texttt{PPLC.SmallGainNStep.}\texttt{\bfseries toSyzygy}}\hfill{\footnotesize \path{lean/qtsl/Qtsl/Composition.lean}, line 796}\par\nopagebreak
\begin{leansrc}
def toSyzygy (st : SmallGainNStep H kk) : SyzygyStep H where …
\end{leansrc}
\par\medskip\noindent{\small \texttt{PPLC.}\texttt{\bfseries sound\_smallGainN}}\hfill{\footnotesize \path{lean/qtsl/Qtsl/Composition.lean}, line 825}\par\nopagebreak
\begin{leansrc}
theorem sound_smallGainN {H : Hypernet} {kk : ℕ} (st : SmallGainNStep H kk)
    (hdiss : ∀ i, (st.diss i).semanticallyValid H)
    (hmem : ∀ i, (st.mem i).semanticallyValid H) :
    st.seq.semanticallyValid H
\end{leansrc}
\par\medskip\noindent{\small \texttt{PPLC.}\texttt{\bfseries smallGainN\_colSum\_mutant\_refuted}}\hfill{\footnotesize \path{lean/qtsl/Qtsl/Composition.lean}, line 857}\par\nopagebreak
\begin{leansrc}
theorem smallGainN_colSum_mutant_refuted :
    ¬ (∀ (kk : ℕ) (LV A Sg : Fin kk → Poly) (lam : Fin kk → ℚ)
        (B : Fin kk → Fin kk → ℚ) (th : ℚ),
        (((List.finRange kk).map
            (fun i => (lam i, Poly.sub (Poly.add (LV i) (A i)) (Sg i)))
          ++ (List.finRange kk).map
            (fun i => ((1 : ℚ), smallGainNMemPoly kk A Sg lam B th i))).foldr
          (fun pr acc => Poly.add (Poly.scale pr.1 pr.2) acc) [])
        = Poly.add
            ((List.finRange kk).foldr
              (fun i acc => Poly.add (Poly.scale (lam i) (LV i)) acc) [])
            (Poly.scale th
              ((List.finRange kk).foldr
                (fun i acc => Poly.add (Poly.scale (lam i) (A i)) acc) [])))
\end{leansrc}
\par\medskip\noindent{\small \texttt{Qtsl.QtslCompose.Discrete.}\texttt{\bfseries smallGainN\_drift}}\hfill{\footnotesize \path{lean/qtsl-aux/QtslAux/QtslCompose/AGComposeDiscrete.lean}, line 372}\par\nopagebreak
\begin{leansrc}
theorem smallGainN_drift (hlam : ∀ i, 0 ≤ lam i)
    (hcol : ∀ j, ∑ i, B i j = lam j)
    (hint : ∀ i, ∀ y, y ∉ A → Integrable (V i) (K y))
    (hd : ∀ i, ∀ y, y ∉ A →
        ∫ z, V i z ∂K y ≤ V i y - alpha i y + sigma i y)
    (hm : ∀ i, ∀ y, y ∉ A →
        lam i * sigma i y ≤ (1 - theta) * ∑ j, B i j * alpha j y) :
    ∀ y, y ∉ A → ∫ z, (∑ i, lam i * V i z) ∂K y
      ≤ (∑ i, lam i * V i y) - theta * ∑ j, lam j * alpha j y
\end{leansrc}
\par\medskip\noindent{\small \texttt{Qtsl.QtslCompose.Discrete.}\texttt{\bfseries smallGain\_drift\_of\_N}}\hfill{\footnotesize \path{lean/qtsl-aux/QtslAux/QtslCompose/AGComposeDiscrete.lean}, line 416}\par\nopagebreak
\begin{leansrc}
theorem smallGain_drift_of_N {V₁ V₂ alpha₁ alpha₂ sigma₁ sigma₂ : S → ℝ}
    {lam theta : ℝ} (hlam : 0 ≤ lam)
    (hint₁ : ∀ y, y ∉ A → Integrable V₁ (K y))
    (hint₂ : ∀ y, y ∉ A → Integrable V₂ (K y))
    (hd₁ : ∀ y, y ∉ A → ∫ z, V₁ z ∂K y ≤ V₁ y - alpha₁ y + sigma₁ y)
    (hd₂ : ∀ y, y ∉ A → ∫ z, V₂ z ∂K y ≤ V₂ y - alpha₂ y + sigma₂ y)
    (hm₁ : ∀ y, y ∉ A → lam * sigma₂ y ≤ (1 - theta) * alpha₁ y)
    (hm₂ : ∀ y, y ∉ A → sigma₁ y ≤ lam * ((1 - theta) * alpha₂ y)) :
    ∀ y, y ∉ A → ∫ z, (V₁ z + lam * V₂ z) ∂K y
      ≤ (V₁ y + lam * V₂ y) - theta * (alpha₁ y + lam * alpha₂ y)
\end{leansrc}
\par\medskip\noindent{\small \texttt{PPLC.}\texttt{\bfseries citeFacet}}\hfill{\footnotesize \path{lean/qtsl/Qtsl/ProductLattice.lean}, line 675}\par\nopagebreak
\begin{leansrc}
def citeFacet (s : Sequent) (rider : Formula) : Sequent
\end{leansrc}
\par\medskip\noindent{\small \texttt{PPLC.}\texttt{\bfseries prodBlock}}\hfill{\footnotesize \path{lean/qtsl/Qtsl/ProductLattice.lean}, line 678}\par\nopagebreak
\begin{leansrc}
def prodBlock (b1 b2 : BlockCert) : BlockCert
\end{leansrc}
\par\medskip\noindent{\small \texttt{PPLC.PProductLatticeData.}\texttt{\bfseries prodBlocks}}\hfill{\footnotesize \path{lean/qtsl/Qtsl/ProductLattice.lean}, line 688}\par\nopagebreak
\begin{leansrc}
def PProductLatticeData.prodBlocks {H1 H2 : Hypernet}
    (d : PProductLatticeData H1 H2) : List BlockCert
\end{leansrc}
\par\medskip\noindent{\small \texttt{PPLC.}\texttt{\bfseries CertifiedBlockLattice}}\hfill{\footnotesize \path{lean/qtsl/Qtsl/ProductLattice.lean}, line 658}\par\nopagebreak
\begin{leansrc}
structure CertifiedBlockLattice (H : Hypernet) where …
\end{leansrc}
\par\medskip\noindent{\small \texttt{PPLC.Hypernet.}\texttt{\bfseries disjUnion}}\hfill{\footnotesize \path{lean/qtsl/Qtsl/ProductLattice.lean}, line 252}\par\nopagebreak
\begin{leansrc}
def Hypernet.disjUnion (H1 H2 : Hypernet) (hd : H1.WireDisjoint H2) : Hypernet
      where …
\end{leansrc}
\par\medskip\noindent{\small \texttt{PPLC.}\texttt{\bfseries sound\_productLattice}}\hfill{\footnotesize \path{lean/qtsl/Qtsl/ProductLattice.lean}, line 712}\par\nopagebreak
\begin{leansrc}
theorem sound_productLattice {H1 H2 : Hypernet}
    (d : PProductLatticeData H1 H2)
    (hL : d.left.Valid) (hR : d.right.Valid) :
    ∀ b ∈ d.prodBlocks, ∀ s ∈ b.facets,
      s.semanticallyValid (H1.disjUnion H2 d.disjoint)
\end{leansrc}
\par\medskip\noindent{\small \texttt{PPLC.Sequent.}\texttt{\bfseries semanticallyValid\_extend}}\hfill{\footnotesize \path{lean/qtsl/Qtsl/ProductLattice.lean}, line 636}\par\nopagebreak
\begin{leansrc}
theorem Sequent.semanticallyValid_extend {H H' : Hypernet} {seq : Sequent}
    (hext : H.WireExtends H')
    (happs : ∀ app ∈ H.applications, app ∈ H'.applications)
    (hws : seq.wellScoped H) (hval : seq.semanticallyValid H) :
    seq.semanticallyValid H'
\end{leansrc}
\par\medskip\noindent{\small \texttt{PPLC.Sequent.}\texttt{\bfseries semanticallyValid\_weaken}}\hfill{\footnotesize \path{lean/qtsl/Qtsl/ProductLattice.lean}, line 648}\par\nopagebreak
\begin{leansrc}
theorem Sequent.semanticallyValid_weaken {H : Hypernet}
    {pre extra post : Formula}
    (h : Sequent.semanticallyValid H ⟨pre, post⟩) :
    Sequent.semanticallyValid H ⟨.and pre extra, post⟩
\end{leansrc}
\par\medskip\noindent{\small \texttt{PPLC.}\texttt{\bfseries sound\_productLattice\_citeLeft}}\hfill{\footnotesize \path{lean/qtsl/Qtsl/ProductLattice.lean}, line 692}\par\nopagebreak
\begin{leansrc}
theorem sound_productLattice_citeLeft {H1 H2 : Hypernet}
    (hd : H1.WireDisjoint H2) {s0 : Sequent} (rider : Formula)
    (hws : s0.wellScoped H1) (hval : s0.semanticallyValid H1) :
    (citeFacet s0 rider).semanticallyValid (H1.disjUnion H2 hd)
\end{leansrc}
\par\medskip\noindent{\small \texttt{PPLC.PProductLatticeData.}\texttt{\bfseries prodBlocks}}\hfill{\footnotesize \path{lean/qtsl/Qtsl/ProductLattice.lean}, line 688}\par\nopagebreak
\begin{leansrc}
def PProductLatticeData.prodBlocks {H1 H2 : Hypernet}
    (d : PProductLatticeData H1 H2) : List BlockCert
\end{leansrc}
\par\medskip\noindent{\small \texttt{PPLC.Hypernet.Valuation.}\texttt{\bfseries restrict}}\hfill{\footnotesize \path{lean/qtsl/Qtsl/ProductLattice.lean}, line 320}\par\nopagebreak
\begin{leansrc}
def Hypernet.Valuation.restrict {H H' : Hypernet} (σ : H'.Valuation)
    (hext : H.WireExtends H') : H.Valuation
\end{leansrc}
\par\medskip\noindent{\small \texttt{PPLC.Formula.}\texttt{\bfseries holds\_restrict}}\hfill{\footnotesize \path{lean/qtsl/Qtsl/ProductLattice.lean}, line 429}\par\nopagebreak
\begin{leansrc}
theorem Formula.holds_restrict {H H' : Hypernet} (σ : H'.Valuation)
    (hext : H.WireExtends H') :
    ∀ (F : Formula), (∀ w ∈ F.freeWires, (H.wireType w).isSome) →
      (Formula.holds (σ.restrict hext) F ↔ Formula.holds σ F)
\end{leansrc}
\par\medskip\noindent{\small \texttt{PPLC.PrimApp.}\texttt{\bfseries semOk\_restrict}}\hfill{\footnotesize \path{lean/qtsl/Qtsl/ProductLattice.lean}, line 485}\par\nopagebreak
\begin{leansrc}
theorem PrimApp.semOk_restrict {H H' : Hypernet} (σ : H'.Valuation)
    (hext : H.WireExtends H') {app : PrimApp}
    (hTC : app.TypeConsistent H.wireType) (h : app.SemOk σ) :
    app.SemOk (σ.restrict hext)
\end{leansrc}
\par\medskip\noindent{\small \texttt{PPLC.Hypernet.}\texttt{\bfseries WireDisjoint}}\hfill{\footnotesize \path{lean/qtsl/Qtsl/ProductLattice.lean}, line 17}\par\nopagebreak
\begin{leansrc}
def Hypernet.WireDisjoint (H1 H2 : Hypernet) : Prop
\end{leansrc}
\par\medskip\noindent{\small \texttt{PPLC.}\texttt{\bfseries prodBlocks\_confined}}\hfill{\footnotesize \path{lean/qtsl/Qtsl/ProductLattice.lean}, line 729}\par\nopagebreak
\begin{leansrc}
theorem prodBlocks_confined {H1 H2 : Hypernet}
    (d : PProductLatticeData H1 H2) :
    ∀ b ∈ d.prodBlocks,
      (∀ w ∈ b.region.freeWires,
        ((H1.disjUnion H2 d.disjoint).wireType w).isSome) ∧
      ∀ s ∈ b.facets, s.wellScoped (H1.disjUnion H2 d.disjoint)
\end{leansrc}
\par\medskip\noindent{\small \texttt{PPLC.CertifiedBlockLattice.}\texttt{\bfseries Valid}}\hfill{\footnotesize \path{lean/qtsl/Qtsl/ProductLattice.lean}, line 671}\par\nopagebreak
\begin{leansrc}
def CertifiedBlockLattice.Valid {H : Hypernet}
    (L : CertifiedBlockLattice H) : Prop
\end{leansrc}
\par\medskip\noindent{\small \texttt{Qtsl.Aliases.}\texttt{\bfseries sound\_QSatFromRegion}}\hfill{\footnotesize \path{lean/qtsl-aux/QtslAux/QtslAliases.lean}, line 35}\par\nopagebreak
\begin{leansrc}
theorem sound_QSatFromRegion {X : Type*} {R : Set X} {sat : X → Prop} {x₀ : X}
    (hchild : ∀ x ∈ R, sat x) (heval : x₀ ∈ R) : sat x₀
\end{leansrc}
\par\medskip\noindent{\small \texttt{Qtsl.Aliases.}\texttt{\bfseries sound\_QSatFromRegion\_sublevel}}\hfill{\footnotesize \path{lean/qtsl-aux/QtslAux/QtslAliases.lean}, line 58}\par\nopagebreak
\begin{leansrc}
theorem sound_QSatFromRegion_sublevel {σ : Type*} {sat : (σ → ℝ) → Prop}
    {P : MvPolynomial σ ℚ} {thr : ℚ} {xq : σ → ℚ}
    (hchild : ∀ x ∈ {x : σ → ℝ |
      MvPolynomial.eval x (MvPolynomial.map (Rat.castHom ℝ) P) ≤ (thr : ℝ)}, sat x)
    (heval : MvPolynomial.eval xq P ≤ thr) :
    sat (fun i => (xq i : ℝ))
\end{leansrc}
\par\medskip\noindent{\small \texttt{PPLC.}\texttt{\bfseries rootEvalFormula\_true\_iff}}\hfill{\footnotesize \path{lean/qtsl/Qtsl/Elaborators.lean}, line 137}\par\nopagebreak
\begin{leansrc}
theorem rootEvalFormula_true_iff (x0 : Nat → Rat) (V : Poly) (thr : Rat) :
    (rootEvalFormula x0 V thr).evalClosed = some true
      ↔ Poly.eval x0 V ≤ thr
\end{leansrc}
\par\medskip\noindent{\small \texttt{Qtsl.Guard.}\texttt{\bfseries GuardSat}}\hfill{\footnotesize \path{lean/qtsl-aux/QtslAux/QGuard.lean}, line 57}\par\nopagebreak
\begin{leansrc}
def GuardSat (init R : Set X) (sat : X → Prop) : Prop
\end{leansrc}
\par\medskip\noindent{\small \texttt{Qtsl.Guard.}\texttt{\bfseries sound\_QGuardIntro}}\hfill{\footnotesize \path{lean/qtsl-aux/QtslAux/QGuard.lean}, line 74}\par\nopagebreak
\begin{leansrc}
theorem sound_QGuardIntro {init R : Set X} {sat : X → Prop}
    (hchild : ∀ x ∈ R, sat x) : GuardSat init R sat
\end{leansrc}
\par\medskip\noindent{\small \texttt{Qtsl.Guard.}\texttt{\bfseries sound\_QGuardIntro\_sublevel}}\hfill{\footnotesize \path{lean/qtsl-aux/QtslAux/QGuard.lean}, line 171}\par\nopagebreak
\begin{leansrc}
theorem sound_QGuardIntro_sublevel {σ : Type*} {init : Set (σ → ℝ)}
    {sat : (σ → ℝ) → Prop} {P : MvPolynomial σ ℚ} {thr : ℚ}
    (hchild : ∀ x ∈ {x : σ → ℝ |
        MvPolynomial.eval x (MvPolynomial.map (Rat.castHom ℝ) P) ≤ (thr : ℝ)},
      sat x) :
    GuardSat init
      {x : σ → ℝ |
        MvPolynomial.eval x (MvPolynomial.map (Rat.castHom ℝ) P) ≤ (thr : ℝ)}
      sat
\end{leansrc}
\par\medskip\noindent{\small \texttt{Qtsl.Guard.}\texttt{\bfseries guard\_eval\_sublevel}}\hfill{\footnotesize \path{lean/qtsl-aux/QtslAux/QGuard.lean}, line 185}\par\nopagebreak
\begin{leansrc}
theorem guard_eval_sublevel {σ : Type*} {P : MvPolynomial σ ℚ} {thr : ℚ}
    {xq : σ → ℚ} (heval : MvPolynomial.eval xq P ≤ thr) :
    (fun i => (xq i : ℝ)) ∈ {x : σ → ℝ |
      MvPolynomial.eval x (MvPolynomial.map (Rat.castHom ℝ) P) ≤ (thr : ℝ)}
\end{leansrc}
\par\medskip\noindent{\small \texttt{Qtsl.Guard.}\texttt{\bfseries sound\_QGuardIntro\_prob}}\hfill{\footnotesize \path{lean/qtsl-aux/QtslAux/QGuard.lean}, line 193}\par\nopagebreak
\begin{leansrc}
theorem sound_QGuardIntro_prob [MeasurableSpace X] {init R : Set X}
    {μ : X → Measure (ℕ → X)} {p : ENNReal} {A : Set (ℕ → X)}
    (hchild : ∀ x ∈ R, p ≤ μ x A) :
    GuardSat init R fun x => p ≤ μ x A
\end{leansrc}
\par\medskip\noindent{\small \texttt{Qtsl.Guard.}\texttt{\bfseries exists\_blanket\_not\_guardSat}}\hfill{\footnotesize \path{lean/qtsl-aux/QtslAux/QGuard.lean}, line 137}\par\nopagebreak
\begin{leansrc}
theorem exists_blanket_not_guardSat :
    ∃ (init R : Set Bool) (sat : Bool → Prop),
      GuardSatBlanket init R sat ∧ ¬ GuardSat init R sat
\end{leansrc}
\par\medskip\noindent{\small \texttt{Qtsl.Guard.}\texttt{\bfseries exists\_not\_guardSat\_wrong\_region}}\hfill{\footnotesize \path{lean/qtsl-aux/QtslAux/QGuard.lean}, line 160}\par\nopagebreak
\begin{leansrc}
theorem exists_not_guardSat_wrong_region :
    ∃ (init R R' : Set Bool) (sat : Bool → Prop),
      (∀ x ∈ R, sat x) ∧ ¬ GuardSat init R' sat
\end{leansrc}
\par\medskip\noindent{\small \texttt{Qtsl.Guard.}\texttt{\bfseries sound\_QGuardElim}}\hfill{\footnotesize \path{lean/qtsl-aux/QtslAux/QGuard.lean}, line 81}\par\nopagebreak
\begin{leansrc}
theorem sound_QGuardElim {init R : Set X} {sat : X → Prop}
    (hg : GuardSat init R sat) (heval : init ⊆ R) : ∀ x₀ ∈ init, sat x₀
\end{leansrc}
\par\medskip\noindent{\small \texttt{Qtsl.Guard.}\texttt{\bfseries exists\_guardSat\_not\_unguarded}}\hfill{\footnotesize \path{lean/qtsl-aux/QtslAux/QGuard.lean}, line 150}\par\nopagebreak
\begin{leansrc}
theorem exists_guardSat_not_unguarded :
    ∃ (init R : Set Bool) (sat : Bool → Prop),
      GuardSat init R sat ∧ ¬ ∀ x₀ ∈ init, sat x₀
\end{leansrc}
\par\medskip\noindent{\small \texttt{Qtsl.Guard.}\texttt{\bfseries qsatFromRegion\_derived}}\hfill{\footnotesize \path{lean/qtsl-aux/QtslAux/QGuard.lean}, line 88}\par\nopagebreak
\begin{leansrc}
theorem qsatFromRegion_derived {init R : Set X} {sat : X → Prop}
    (hchild : ∀ x ∈ R, sat x) (heval : init ⊆ R) : ∀ x₀ ∈ init, sat x₀
\end{leansrc}
\par\medskip\noindent{\small \texttt{Qtsl.LayerD.}\texttt{\bfseries sound\_GSHS\_safety\_mp}}\hfill{\footnotesize \path{lean/qtsl-aux/QtslAux/LayerD/MPDischarge.lean}, line 511}\par\nopagebreak
\begin{leansrc}
theorem sound_GSHS_safety_mp [IsMarkovKernel M.seg] [IsMarkovKernel M.jump]
    {V : State → ℝ} {U : Set State} (hV : Measurable V) (hU : MeasurableSet U)
    (hmp : ∀ x : State, Nonempty (MPFlow M 0 V x))
    (hForce : ForcedNonexpansion M V)
    (hNonneg : NonnegV V) {C : ℝ} (hC0 : 0 ≤ C) (hVC : ∀ y, V y ≤ C)
    (hFloor : ∀ y ∈ U, 1 ≤ V y) (x : State) :
    Battery.pathMeasure (skel M) (SkelState.alive x)
        {ω | ∃ n, ω n ∈ skelBad U}
      ≤ ENNReal.ofReal (V x)
\end{leansrc}
\par\medskip\noindent{\small \texttt{Qtsl.LayerD.}\texttt{\bfseries sound\_GSHS\_safety\_battery\_mp}}\hfill{\footnotesize \path{lean/qtsl-aux/QtslAux/LayerD/MPDischarge.lean}, line 524}\par\nopagebreak
\begin{leansrc}
theorem sound_GSHS_safety_battery_mp
    [IsMarkovKernel M.seg] [IsMarkovKernel M.jump]
    {V : State → ℝ} {U : Set State} {p thr : ℝ}
    (hV : Measurable V) (hU : MeasurableSet U)
    (hmp : ∀ x : State, Nonempty (MPFlow M 0 V x))
    (hForce : ForcedNonexpansion M V)
    (hNonneg : NonnegV V) {C : ℝ} (hC0 : 0 ≤ C) (hVC : ∀ y, V y ≤ C)
    (hFloor : ∀ y ∈ U, 1 ≤ V y)
    (hΘ : p + thr ≤ 1) (x : State) (hthr : V x ≤ thr) :
    ENNReal.ofReal p ≤ Battery.pathMeasure (skel M) (SkelState.alive x)
      {ω | ∀ n, ω n ∉ skelBad U}
\end{leansrc}
\par\medskip\noindent{\small \texttt{Qtsl.Battery.}\texttt{\bfseries sound\_QSafetyBattery\_staged}}\hfill{\footnotesize \path{lean/qtsl-aux/QtslAux/QtslBattery.lean}, line 556}\par\nopagebreak
\begin{leansrc}
theorem sound_QSafetyBattery_staged
    {K : Kernel (Md × X) (Md × X)} [IsMarkovKernel K]
    {Dm Dbarm Um : Md → Set X} {p thr : ℚ} {sat : Md × X → Prop}
    (hDm : ∀ m, MeasurableSet (Dm m)) (hDbarm : ∀ m, MeasurableSet (Dbarm m))
    (hUm : ∀ m, MeasurableSet (Um m)) (hVmeas : ∀ m, Measurable (V m))
    (hF : ∀ m z, ¬ sat (m, z) → z ∈ Um m)
    (hΘ : p + thr ≤ 1)
    (hDsub : ∀ m, Dm m ⊆ Dbarm m)
    (hR3 : ∀ m z, z ∈ Dm m → z ∉ Um m → K (m, z) {y | y.2 ∈ Dbarm y.1} = 1)
    (hR4 : ∀ m z, z ∈ Dm m → z ∉ Um m →
      K (m, z) ({y | y.2 ∈ Dm y.1} ∪ {y | y.2 ∈ Um y.1}) = 1)
    (hint : ∀ m z, z ∈ Dm m → z ∉ Um m → Integrable (V m) ((flow m) z))
    (hB1 : ∀ m z, z ∈ Dm m → z ∉ Um m → ∫ ζ, V m ζ ∂(flow m) z ≤ V m z)
    (hB2 : ∀ j, ∀ z ∈ G j, V (dst j) (g j z) ≤ V (src j) z)
    (hB3 : ∀ m z, z ∈ Um m → z ∈ Dbarm m → 1 ≤ V m z)
    (hB4 : ∀ m z, z ∈ Dbarm m → 0 ≤ V m z)
    (hmix : ∀ m z, z ∈ Dm m → z ∉ Um m →
      ∃ (w₀ : ℝ≥0) (w : J → ℝ≥0), w₀ + ∑ j, w j = 1 ∧
        (∀ j, w j ≠ 0 → src j = m ∧ z ∈ G j) ∧
        K (m, z) = (w₀ : ℝ≥0∞) • ((flow m) z).map (Prod.mk m)
          + ∑ j, (w j : ℝ≥0∞) • Measure.dirac (dst j, g j z)) :
    ∀ m z, z ∈ Dm m → V m z ≤ (thr : ℝ) →
      ENNReal.ofReal (p : ℝ)
        ≤ pathMeasure K (m, z) {ω | ∀ n, sat (ω n)}
\end{leansrc}
\par\medskip\noindent{\small \texttt{SuFunctor.Site.}\texttt{\bfseries measurableSet\_confinedEvent\_cylinder}}\hfill{\footnotesize \path{lean/su-functor/SuFunctor/Site/FellAcceptance.lean}, line 413}\par\nopagebreak
\begin{leansrc}
theorem measurableSet_confinedEvent_cylinder {D : Set X}
    (hD : IsClosed D) :
    MeasurableSet
      (confinedEvent ((Set.univ : Set ↥W) ×ˢ D) : Set (FTraj W X))
\end{leansrc}
\par\medskip\noindent{\small \texttt{SuFunctor.Site.}\texttt{\bfseries satAtM\_safetySpecF\_iff}}\hfill{\footnotesize \path{lean/su-functor/SuFunctor/Site/FellBehaviorSpec.lean}, line 254}\par\nopagebreak
\begin{leansrc}
theorem satAtM_safetySpecF_iff (U : ι → Set T)
    (hU : ∀ i, IsOpen (U i)) {D : Set X} (hD : IsClosed D) (j : ι)
    (μ : MeasJoint (FTraj (U j) X)) :
    SatAtM μ (specOfBehaviorF U (safetyBehaviorF U D) j) ↔
      (μ.1 : Measure (FTraj (U j) X))
        (confinedEvent ((Set.univ : Set ↥(U j)) ×ˢ D))ᶜ = 0
\end{leansrc}
\par\medskip\noindent{\small \texttt{SuFunctor.Site.}\texttt{\bfseries pull\_indPenM}}\hfill{\footnotesize \path{lean/su-functor/SuFunctor/Site/FellBehaviorSpec.lean}, line 110}\par\nopagebreak
\begin{leansrc}
theorem pull_indPenM (h : K → L) (hm : Measurable h) {S : Set L}
    (hS : MeasurableSet S) :
    MeasPen.pull h (indPenM S) = indPenM (h ⁻¹' S)
\end{leansrc}
\par\medskip\noindent{\small \texttt{SuFunctor.Site.}\texttt{\bfseries graphCl\_mem\_confinedEvent\_iff}}\hfill{\footnotesize \path{lean/su-functor/SuFunctor/Site/FellBehaviorSpec.lean}, line 340}\par\nopagebreak
\begin{leansrc}
theorem graphCl_mem_confinedEvent_iff {D : Set X} (hD : IsClosed D)
    (f : Traj W X) :
    graphCl f ∈ confinedEvent ((Set.univ : Set ↥W) ×ˢ D) ↔
      ∀ w, f w ∈ D
\end{leansrc}
\par\medskip\noindent{\small \texttt{SuFunctor.Site.}\texttt{\bfseries mboxplusWF\_indPenM}}\hfill{\footnotesize \path{lean/su-functor/SuFunctor/Site/FellBehaviorSpec.lean}, line 280}\par\nopagebreak
\begin{leansrc}
theorem mboxplusWF_indPenM {W : Set T} {π₁ : Z → U'} {π₂ : Z → V'}
    (hp₁ : IsProperMap π₁) (hp₂ : IsProperMap π₂)
    {S₁ : Set (FTraj W U')} {S₂ : Set (FTraj W V')}
    (hS₁ : MeasurableSet S₁) (hS₂ : MeasurableSet S₂) :
    mboxplusWF π₁ π₂ (indPenM S₁) (indPenM S₂)
      = indPenM (fellMap (W := W) π₁ ⁻¹' S₁ ∩
          fellMap (W := W) π₂ ⁻¹' S₂)
\end{leansrc}
\par\medskip\noindent{\small \texttt{Qtsl.Battery.}\texttt{\bfseries sound\_QSafetyBattery}}\hfill{\footnotesize \path{lean/qtsl-aux/QtslAux/Battery.lean}, line 433}\par\nopagebreak
\begin{leansrc}
theorem sound_QSafetyBattery {sat : S → Prop}
    (hD : MeasurableSet D) (hDbar : MeasurableSet Dbar) (hU : MeasurableSet U)
    (hV : Measurable V)
    (hF : ∀ y, ¬ sat y → y ∈ U)
    (hΘ : p + thr ≤ 1)
    (hDsub : D ⊆ Dbar)
    (hR3 : ∀ y ∈ D, y ∉ U → K y Dbar = 1)
    (hR4 : ∀ y ∈ D, y ∉ U → K y (D ∪ U) = 1)
    (hB12int : ∀ y ∈ D, y ∉ U → Integrable V (K y))
    (hB12 : ∀ y ∈ D, y ∉ U → ∫ z, V z ∂K y ≤ V y)
    (hB3 : ∀ y ∈ U, y ∈ Dbar → 1 ≤ V y)
    (hB4 : ∀ y ∈ Dbar, 0 ≤ V y)
    (hx : x ∈ D) (hthr : V x ≤ thr) :
    ENNReal.ofReal p ≤ pathMeasure K x {ω | ∀ n, sat (ω n)}
\end{leansrc}
\par\medskip\noindent{\small \texttt{Qtsl.Battery.}\texttt{\bfseries pathMeasure}}\hfill{\footnotesize \path{lean/qtsl-aux/QtslAux/Battery.lean}, line 69}\par\nopagebreak
\begin{leansrc}
noncomputable def pathMeasure (K : Kernel S S) [IsMarkovKernel K] (x : S) :
    Measure (ℕ → S)
\end{leansrc}
\par\medskip\noindent{\small \texttt{Qtsl.LayerD.}\texttt{\bfseries sound\_GSHS\_safety\_battery}}\hfill{\footnotesize \path{lean/qtsl-aux/QtslAux/LayerD/Skeleton.lean}, line 253}\par\nopagebreak
\begin{leansrc}
theorem sound_GSHS_safety_battery [IsMarkovKernel M.seg] [IsMarkovKernel M.jump]
    {V : State → ℝ} {U : Set State} {p thr : ℝ}
    (hV : Measurable V) (hU : MeasurableSet U)
    (hSegD : SegmentDynkin M 0 V) (hForce : ForcedNonexpansion M V)
    (hNonneg : NonnegV V) {C : ℝ} (hC0 : 0 ≤ C) (hVC : ∀ y, V y ≤ C)
    (hFloor : ∀ y ∈ U, 1 ≤ V y)
    (hΘ : p + thr ≤ 1) (x : State) (hthr : V x ≤ thr) :
    ENNReal.ofReal p ≤ Battery.pathMeasure (skel M) (SkelState.alive x)
      {ω | ∀ n, ω n ∉ skelBad U}
\end{leansrc}
\par\medskip\noindent{\small \texttt{Qtsl.Battery.}\texttt{\bfseries sound\_QSafetyBattery\_failClosed}}\hfill{\footnotesize \path{lean/qtsl-aux/QtslAux/Battery.lean}, line 311}\par\nopagebreak
\begin{leansrc}
theorem sound_QSafetyBattery_failClosed
    (hD : MeasurableSet D) (hDbar : MeasurableSet Dbar) (hU : MeasurableSet U)
    (hV : Measurable V)
    (hΘ : p + thr ≤ 1)
    (hDsub : D ⊆ Dbar)
    (hR3 : ∀ y ∈ D, y ∉ U → K y Dbar = 1)
    (hR4 : ∀ y ∈ D, y ∉ U → K y (D ∪ U) = 1)
    (hB12int : ∀ y ∈ D, y ∉ U → Integrable V (K y))
    (hB12 : ∀ y ∈ D, y ∉ U → ∫ z, V z ∂K y ≤ V y)
    (hB3 : ∀ y ∈ U, y ∈ Dbar → 1 ≤ V y)
    (hB4 : ∀ y ∈ Dbar, 0 ≤ V y)
    (hx : x ∈ D) (hthr : V x ≤ thr) :
    ENNReal.ofReal p ≤ pathMeasure K x {ω | ∀ n, ω n ∈ D \ U}
\end{leansrc}
\par\medskip\noindent{\small \texttt{OptionalSamplingRC.}\texttt{\bfseries supermartingale\_of\_compensator\_nonpos}}\hfill{\footnotesize \path{lean/optional-sampling/OptionalSamplingRC/Basic.lean}, line 1150}\par\nopagebreak
\begin{leansrc}
theorem supermartingale_of_compensator_nonpos
    [IsProbabilityMeasure μ]
    (W : ℝ≥0 → Ω → ℝ) (g : ℝ → Ω → ℝ)
    (hmart : Martingale
      (fun (t : ℝ≥0) (ω : Ω) => W t ω - ∫ s in (0:ℝ)..(t:ℝ), g s ω) ℱ μ)
    (hGadp : Adapted ℱ
      (fun (t : ℝ≥0) (ω : Ω) => ∫ s in (0:ℝ)..(t:ℝ), g s ω))
    (hWint : ∀ t : ℝ≥0, Integrable (W t) μ)
    (hgint : ∀ ω, ∀ t : ℝ, IntervalIntegrable (fun s => g s ω) volume 0 t)
    (hg0 : ∀ ω, ∀ s : ℝ, 0 ≤ s → g s ω ≤ 0) :
    Supermartingale W ℱ μ
\end{leansrc}
\par\medskip\noindent{\small \texttt{Qtsl.Battery.}\texttt{\bfseries mixture\_drift}}\hfill{\footnotesize \path{lean/qtsl-aux/QtslAux/QtslBattery.lean}, line 488}\par\nopagebreak
\begin{leansrc}
lemma mixture_drift {m : Md} {z : X} {κ : Measure (Md × X)}
    {w₀ : ℝ≥0} {w : J → ℝ≥0}
    (hVmeas : ∀ m', Measurable (V m'))
    (hκ : κ = (w₀ : ℝ≥0∞) • ((flow m) z).map (Prod.mk m)
      + ∑ j, (w j : ℝ≥0∞) • Measure.dirac (dst j, g j z))
    (hsum : w₀ + ∑ j, w j = 1)
    (hguard : ∀ j, w j ≠ 0 → V (dst j) (g j z) ≤ V m z)
    (hflowint : Integrable (fun ζ => V m ζ) ((flow m) z))
    (hflowdrift : ∫ ζ, V m ζ ∂(flow m) z ≤ V m z) :
    ∫ y, V y.1 y.2 ∂κ ≤ V m z
\end{leansrc}
\par\medskip\noindent{\small \texttt{Qtsl.Battery.}\texttt{\bfseries stronglyAdapted\_stopped}}\hfill{\footnotesize \path{lean/qtsl-aux/QtslAux/Battery.lean}, line 187}\par\nopagebreak
\begin{leansrc}
lemma stronglyAdapted_stopped {f : S → ℝ} {A : Set S} (hA : MeasurableSet A)
    (hf : Measurable f) :
    StronglyAdapted (Filtration.piLE (X := fun _ : ℕ => S)) (stopped f A)
\end{leansrc}
\par\medskip\noindent{\small \texttt{Qtsl.Battery.}\texttt{\bfseries one\_le\_stopped}}\hfill{\footnotesize \path{lean/qtsl-aux/QtslAux/Battery.lean}, line 133}\par\nopagebreak
\begin{leansrc}
lemma one_le_stopped {f : S → ℝ} {A : Set S} (hf1 : ∀ y ∈ A, 1 ≤ f y) :
    ∀ (n : ℕ) (ω : ℕ → S), (∃ k ≤ n, ω k ∈ A) → 1 ≤ stopped f A n ω
\end{leansrc}
\par\medskip\noindent{\small \texttt{Qtsl.Battery.}\texttt{\bfseries stopped\_congr}}\hfill{\footnotesize \path{lean/qtsl-aux/QtslAux/Battery.lean}, line 154}\par\nopagebreak
\begin{leansrc}
lemma stopped_congr {f : S → ℝ} {A : Set S} : ∀ {n : ℕ} {ω ω' : ℕ → S},
    (∀ k ≤ n, ω k = ω' k) → stopped f A n ω = stopped f A n ω'
\end{leansrc}
\par\medskip\noindent{\small \texttt{Qtsl.Battery.}\texttt{\bfseries supermartingale\_stopped}}\hfill{\footnotesize \path{lean/qtsl-aux/QtslAux/Battery.lean}, line 263}\par\nopagebreak
\begin{leansrc}
theorem supermartingale_stopped {f : S → ℝ} {A : Set S} {C : ℝ}
    (hA : MeasurableSet A) (hf : Measurable f) (hf0 : ∀ y, 0 ≤ f y) (hfC : ∀ y, f y
      ≤ C)
    (hdrift : ∀ y, y ∉ A → ∫ z, f z ∂K y ≤ f y) (x : S) :
    Supermartingale (stopped f A) (Filtration.piLE (X := fun _ : ℕ => S))
      (pathMeasure K x)
\end{leansrc}
\par\medskip\noindent{\small \texttt{Qtsl.Battery.}\texttt{\bfseries measure\_reach\_le}}\hfill{\footnotesize \path{lean/qtsl-aux/QtslAux/Battery.lean}, line 279}\par\nopagebreak
\begin{leansrc}
theorem measure_reach_le {f : S → ℝ} {A : Set S} {C : ℝ}
    (hA : MeasurableSet A) (hf : Measurable f) (hf0 : ∀ y, 0 ≤ f y) (hfC : ∀ y, f y
      ≤ C)
    (hf1 : ∀ y ∈ A, 1 ≤ f y) (hdrift : ∀ y, y ∉ A → ∫ z, f z ∂K y ≤ f y) (x : S) :
    pathMeasure K x {ω | ∃ n, ω n ∈ A} ≤ ENNReal.ofReal (f x)
\end{leansrc}
\par\medskip\noindent{\small \texttt{Qtsl.Battery.}\texttt{\bfseries mixture\_integrable}}\hfill{\footnotesize \path{lean/qtsl-aux/QtslAux/QtslBattery.lean}, line 468}\par\nopagebreak
\begin{leansrc}
lemma mixture_integrable {m : Md} {z : X} {κ : Measure (Md × X)}
    {w₀ : ℝ≥0} {w : J → ℝ≥0}
    (hVmeas : ∀ m', Measurable (V m'))
    (hκ : κ = (w₀ : ℝ≥0∞) • ((flow m) z).map (Prod.mk m)
      + ∑ j, (w j : ℝ≥0∞) • Measure.dirac (dst j, g j z))
    (hflowint : Integrable (fun ζ => V m ζ) ((flow m) z)) :
    Integrable (fun y : Md × X => V y.1 y.2) κ
\end{leansrc}
\par\medskip\noindent{\small \texttt{Qtsl.LayerD.}\texttt{\bfseries sound\_GSHS\_deadline\_mp}}\hfill{\footnotesize \path{lean/qtsl-aux/QtslAux/LayerD/MPDischarge.lean}, line 539}\par\nopagebreak
\begin{leansrc}
theorem sound_GSHS_deadline_mp [IsMarkovKernel M.seg] [IsMarkovKernel M.jump]
    {V : State → ℝ} {θ ε : ℝ} {T : ℝ≥0}
    (hV : Measurable V) (hθ : 0 ≤ θ)
    (hmp : ∀ x : State, Nonempty (MPFlow M θ V x))
    (hForce : ForcedNonexpansion M V)
    (hNonneg : NonnegV V) (hε : 0 < ε) (hFloor : ∀ y, ε ≤ V y) (x : State) :
    Battery.pathMeasure (tskel M) (TSkelState.alive (0, x))
        {ω | ∃ n, T < tclock (ω n)}
      ≤ ENNReal.ofReal (Real.exp (-(θ * T)) * V x / ε)
\end{leansrc}
\par\medskip\noindent{\small \texttt{Qtsl.LayerD.}\texttt{\bfseries sound\_GSHS\_deadline\_battery\_mp}}\hfill{\footnotesize \path{lean/qtsl-aux/QtslAux/LayerD/MPDischarge.lean}, line 552}\par\nopagebreak
\begin{leansrc}
theorem sound_GSHS_deadline_battery_mp
    [IsMarkovKernel M.seg] [IsMarkovKernel M.jump]
    {V : State → ℝ} {θ ε p thr : ℝ} {T : ℝ≥0}
    (hV : Measurable V) (hθ : 0 ≤ θ)
    (hmp : ∀ x : State, Nonempty (MPFlow M θ V x))
    (hForce : ForcedNonexpansion M V)
    (hNonneg : NonnegV V) (hε : 0 < ε) (hFloor : ∀ y, ε ≤ V y)
    (hΘ : p + Real.exp (-(θ * T)) * thr / ε ≤ 1) (x : State)
    (hthr : V x ≤ thr) :
    ENNReal.ofReal p ≤ Battery.pathMeasure (tskel M) (TSkelState.alive (0, x))
      {ω | ∀ n, tclock (ω n) ≤ T}
\end{leansrc}
\par\medskip\noindent{\small \texttt{SuFunctor.Site.}\texttt{\bfseries satAtM\_acceptanceSpecF\_ae}}\hfill{\footnotesize \path{lean/su-functor/SuFunctor/Site/FellBehaviorSpec.lean}, line 180}\par\nopagebreak
\begin{leansrc}
theorem satAtM_acceptanceSpecF_ae (U : ι → Set T)
    (E : ∀ i, Set (FTraj (U i) X)) (j : ι)
    {μ : MeasJoint (FTraj (U j) X)}
    (h : SatAtM μ (acceptanceSpecF U E j)) :
    (μ.1 : Measure (FTraj (U j) X)) (E j)ᶜ = 0
\end{leansrc}
\par\medskip\noindent{\small \texttt{SuFunctor.Site.}\texttt{\bfseries hitEvent\_iUnion}}\hfill{\footnotesize \path{lean/su-functor/SuFunctor/Site/FellAcceptance.lean}, line 123}\par\nopagebreak
\begin{leansrc}
theorem hitEvent_iUnion {ι : Sort*} (A : ι → Set α) :
    hitEvent (⋃ i, A i) = ⋃ i, hitEvent (A i)
\end{leansrc}
\par\medskip\noindent{\small \texttt{SuFunctor.Site.}\texttt{\bfseries measurableSet\_hitEvent\_univ\_prod\_of\_isCompact}}\hfill{\footnotesize \path{lean/su-functor/SuFunctor/Site/FellAcceptance.lean}, line 397}\par\nopagebreak
\begin{leansrc}
theorem measurableSet_hitEvent_univ_prod_of_isCompact
    [SigmaCompactSpace ↥W] {R : Set X} (hR : IsCompact R) :
    MeasurableSet
      (hitEvent ((Set.univ : Set ↥W) ×ˢ R) : Set (FTraj W X))
\end{leansrc}
\par\medskip\noindent{\small \texttt{Qtsl.Battery.}\texttt{\bfseries sound\_QDeadline}}\hfill{\footnotesize \path{lean/qtsl-aux/QtslAux/Battery.lean}, line 749}\par\nopagebreak
\begin{leansrc}
theorem sound_QDeadline {T : ℕ} {M : ℝ} {sat : S → Prop}
    (hD : MeasurableSet D) (hDbar : MeasurableSet Dbar) (hB : MeasurableSet B)
    (hV : Measurable V)
    (hF : ∀ y ∈ B, sat y)
    (hβ : 0 ≤ β)
    (hΘF : β ^ T * thr ≤ 1 - p)
    (hDsub : D ⊆ Dbar)
    (hR3 : ∀ y ∈ D, y ∉ B → K y Dbar = 1)
    (hR4 : ∀ y ∈ D, y ∉ B → K y (D ∪ B) = 1)
    (hint : ∀ y ∈ D, y ∉ B → Integrable V (K y))
    (hdecay : ∀ y ∈ D, y ∉ B → ∫ z, V z ∂K y ≤ β * V y)
    (hfloor : ∀ y ∈ Dbar, y ∉ B → 1 ≤ V y)
    (hnonneg : ∀ y ∈ Dbar, 0 ≤ V y)
    (hM : ∀ y ∈ Dbar, V y ≤ M)
    (hx : x ∈ D) (hthr : V x ≤ thr) :
    ENNReal.ofReal p ≤ pathMeasure K x {ω | ∃ k ≤ T, sat (ω k)}
\end{leansrc}
\par\medskip\noindent{\small \texttt{Qtsl.Battery.}\texttt{\bfseries sound\_QDeadlineBattery}}\hfill{\footnotesize \path{lean/qtsl-aux/QtslAux/QtslBattery.lean}, line 413}\par\nopagebreak
\begin{leansrc}
theorem sound_QDeadlineBattery {n : ℕ} {sat : S → Prop}
    (hD : MeasurableSet D) (hDbar : MeasurableSet Dbar) (hB : MeasurableSet B)
    (hV : Measurable V)
    (hF : ∀ y ∈ B, sat y)
    (hβ : 0 ≤ β) (hβθ : β ≤ Real.exp (-θ)) (hθ : 0 ≤ θ)
    (hq0 : 0 ≤ q)
    (hq : q ≤ ∑ k ∈ range (n + 1), (θ * T) ^ k / (Nat.factorial k : ℝ))
    (hp : p ≤ 1)
    (hDsub : D ⊆ Dbar)
    (hR3 : ∀ y ∈ D, y ∉ B → K y Dbar = 1)
    (hR4 : ∀ y ∈ D, y ∉ B → K y (D ∪ B) = 1)
    (hint : ∀ y ∈ D, y ∉ B → Integrable V (K y))
    (hdecay : ∀ y ∈ D, y ∉ B → ∫ z, V z ∂K y ≤ β * V y)
    (hfloor : ∀ y ∈ Dbar, y ∉ B → 1 ≤ V y)
    (hnonneg : ∀ y ∈ Dbar, 0 ≤ V y)
    (hM : ∀ y ∈ Dbar, V y ≤ M)
    (hx : x ∈ D) (hthr : V x ≤ (1 - p) * q) :
    ENNReal.ofReal p ≤ pathMeasure K x {ω | ∃ k ≤ T, sat (ω k)}
\end{leansrc}
\par\medskip\noindent{\small \texttt{Qtsl.Battery.}\texttt{\bfseries sound\_QDeadline\_hit}}\hfill{\footnotesize \path{lean/qtsl-aux/QtslAux/Battery.lean}, line 568}\par\nopagebreak
\begin{leansrc}
theorem sound_QDeadline_hit {T : ℕ} {M : ℝ}
    (hD : MeasurableSet D) (hDbar : MeasurableSet Dbar) (hB : MeasurableSet B)
    (hV : Measurable V)
    (hβ : 0 ≤ β)
    (hΘF : β ^ T * thr ≤ 1 - p)
    (hDsub : D ⊆ Dbar)
    (hR3 : ∀ y ∈ D, y ∉ B → K y Dbar = 1)
    (hR4 : ∀ y ∈ D, y ∉ B → K y (D ∪ B) = 1)
    (hint : ∀ y ∈ D, y ∉ B → Integrable V (K y))
    (hdecay : ∀ y ∈ D, y ∉ B → ∫ z, V z ∂K y ≤ β * V y)
    (hfloor : ∀ y ∈ Dbar, y ∉ B → 1 ≤ V y)
    (hnonneg : ∀ y ∈ Dbar, 0 ≤ V y)
    (hM : ∀ y ∈ Dbar, V y ≤ M)
    (hx : x ∈ D) (hthr : V x ≤ thr) :
    ENNReal.ofReal p ≤ pathMeasure K x {ω | ∃ k ≤ T, ω k ∈ B}
\end{leansrc}
\par\medskip\noindent{\small \texttt{MeasureTheory.}\texttt{\bfseries supermartingale\_inv\_pow\_mul}}\hfill{\footnotesize \path{lean/qtsl-aux/QtslAux/Deadline.lean}, line 34}\par\nopagebreak
\begin{leansrc}
theorem supermartingale_inv_pow_mul [IsFiniteMeasure μ] {f : ℕ → Ω → ℝ} {β : ℝ}
    (hadp : StronglyAdapted 𝒢 f) (hint : ∀ n, Integrable (f n) μ) (hβ : 0 < β)
    (hdrift : ∀ n, μ[f (n + 1) | 𝒢 n] ≤ᵐ[μ] fun ω => β * f n ω) :
    Supermartingale (fun n ω => β⁻¹ ^ n * f n ω) 𝒢 μ
\end{leansrc}
\par\medskip\noindent{\small \texttt{Qtsl.Deadline.}\texttt{\bfseries expMoment\_hitting}}\hfill{\footnotesize \path{lean/qtsl-aux/QtslAux/Deadline.lean}, line 128}\par\nopagebreak
\begin{leansrc}
theorem expMoment_hitting (hβ : 0 < β)
    (hint : ∀ n, Integrable (V n) μ) (hnonneg : 0 ≤ V)
    (hmB : ∀ n, MeasurableSet[𝒢 n] {ω | X n ω ∈ B})
    (hdrift : ∀ n, ∀ᵐ ω ∂μ, X n ω ∉ B → (μ[V (n + 1) | 𝒢 n]) ω ≤ β * V n ω)
    (hbarrier : ∀ n ω, X n ω ∈ B → a ≤ V n ω) :
    ENNReal.ofReal a *
        ∑' k, ENNReal.ofReal (β⁻¹ ^ k) * μ {ω | X k ω ∈ B ∧ ∀ j < k, X j ω ∉ B}
      ≤ ENNReal.ofReal (∫ ω, V 0 ω ∂μ)
\end{leansrc}
\par\medskip\noindent{\small \texttt{Qtsl.Battery.}\texttt{\bfseries ae\_no\_escape}}\hfill{\footnotesize \path{lean/qtsl-aux/QtslAux/Battery.lean}, line 502}\par\nopagebreak
\begin{leansrc}
private lemma ae_no_escape (hD : MeasurableSet D) (hB : MeasurableSet B)
    (hR4 : ∀ y ∈ D, y ∉ B → K y (D ∪ B) = 1) (hx : x ∈ D) :
    pathMeasure K x (⋃ n, {ω : ℕ → S | (∀ k ≤ n, ω k ∉ B) ∧ ω n ∉ D}) = 0
\end{leansrc}
\par\medskip\noindent{\small \texttt{Qtsl.Battery.}\texttt{\bfseries ae\_coord\_zero\_mem}}\hfill{\footnotesize \path{lean/qtsl-aux/QtslAux/Battery.lean}, line 484}\par\nopagebreak
\begin{leansrc}
private lemma ae_coord_zero_mem {E : Set S} (hE : MeasurableSet E) (hxE : x ∈ E) :
    ∀ᵐ ω ∂pathMeasure K x, ω 0 ∈ E
\end{leansrc}
\par\medskip\noindent{\small \texttt{Qtsl.Battery.}\texttt{\bfseries compl\_firstExit\_subset}}\hfill{\footnotesize \path{lean/qtsl-aux/QtslAux/QtslBattery.lean}, line 70}\par\nopagebreak
\begin{leansrc}
theorem compl_firstExit_subset {S : Type*} (Stop Good : Set S) (T : ℕ) :
    {ω : ℕ → S | ∃ n ≤ T, ω n ∈ Good ∧ ∀ j < n, ω j ∉ Stop}ᶜ ⊆
      {ω : ℕ → S | ∃ n, ω n ∈ Stop \ Good ∧ ∀ j < n, ω j ∉ Stop} ∪
        {ω : ℕ → S | ∀ k ≤ T, ω k ∉ Stop}
\end{leansrc}
\par\medskip\noindent{\small \texttt{Qtsl.Battery.}\texttt{\bfseries sound\_QUnionBound}}\hfill{\footnotesize \path{lean/qtsl-aux/QtslAux/QtslBattery.lean}, line 50}\par\nopagebreak
\begin{leansrc}
theorem sound_QUnionBound {Ω : Type*} [MeasurableSpace Ω] (μ : Measure Ω)
    {E₁ E₂ F : Set Ω} {a b : ℝ≥0∞}
    (h₁ : μ E₁ ≤ a) (h₂ : μ E₂ ≤ b) (hsub : F ⊆ E₁ ∪ E₂) :
    μ F ≤ a + b
\end{leansrc}
\par\medskip\noindent{\small \texttt{Qtsl.Battery.}\texttt{\bfseries sound\_QUnionBound\_ofReal}}\hfill{\footnotesize \path{lean/qtsl-aux/QtslAux/QtslBattery.lean}, line 58}\par\nopagebreak
\begin{leansrc}
theorem sound_QUnionBound_ofReal {Ω : Type*} [MeasurableSpace Ω] (μ : Measure Ω)
    {E₁ E₂ F : Set Ω} {a b : ℝ} (ha : 0 ≤ a) (hb : 0 ≤ b)
    (h₁ : μ E₁ ≤ ENNReal.ofReal a) (h₂ : μ E₂ ≤ ENNReal.ofReal b)
    (hsub : F ⊆ E₁ ∪ E₂) :
    μ F ≤ ENNReal.ofReal (a + b)
\end{leansrc}
\par\medskip\noindent{\small \texttt{Qtsl.Battery.}\texttt{\bfseries sound\_QUntilBattery\_firstHit}}\hfill{\footnotesize \path{lean/qtsl-aux/QtslAux/QtslBattery.lean}, line 206}\par\nopagebreak
\begin{leansrc}
theorem sound_QUntilBattery_firstHit
    (hD : MeasurableSet D) (hDbar : MeasurableSet Dbar)
    (hA : MeasurableSet A) (hR : MeasurableSet R)
    (hVa : Measurable Va) (hVd : Measurable Vd)
    (hβ : 0 ≤ β) (hc : β ^ T ≤ c)
    (hDsub : D ⊆ Dbar)
    (hR3 : ∀ y ∈ D, y ∉ A ∪ R → K y Dbar = 1)
    (hR4 : ∀ y ∈ D, y ∉ A ∪ R → K y (D ∪ (A ∪ R)) = 1)
    (hVaint : ∀ y ∈ D, y ∉ A ∪ R → Integrable Va (K y))
    (hVadrift : ∀ y ∈ D, y ∉ A ∪ R → ∫ z, Va z ∂K y ≤ Va y)
    (hVafloor : ∀ y ∈ A, y ∈ Dbar → 1 ≤ Va y)
    (hVanonneg : ∀ y ∈ Dbar, 0 ≤ Va y)
    (hVdint : ∀ y ∈ D, y ∉ A ∪ R → Integrable Vd (K y))
    (hVddecay : ∀ y ∈ D, y ∉ A ∪ R → ∫ z, Vd z ∂K y ≤ β * Vd y)
    (hVdfloor : ∀ y ∈ Dbar, y ∉ A ∪ R → 1 ≤ Vd y)
    (hVdnonneg : ∀ y ∈ Dbar, 0 ≤ Vd y)
    (hVdM : ∀ y ∈ Dbar, Vd y ≤ M)
    (hx : x ∈ D) (hthr : Va x + c * Vd x ≤ 1 - p) :
    ENNReal.ofReal p ≤ pathMeasure K x
      {ω | ∃ n ≤ T, ω n ∈ R ∧ ∀ j < n, ω j ∉ (A ∪ R) ∪ Dᶜ}
\end{leansrc}
\par\medskip\noindent{\small \texttt{Qtsl.LayerD.}\texttt{\bfseries sound\_GSHS\_until\_mp}}\hfill{\footnotesize \path{lean/qtsl-aux/QtslAux/LayerD/MPDischarge.lean}, line 584}\par\nopagebreak
\begin{leansrc}
theorem sound_GSHS_until_mp [IsMarkovKernel M.seg] [IsMarkovKernel M.jump]
    {Va Vd : State → ℝ} {Us : Set State} {θ ε : ℝ} {T : ℝ≥0}
    (hVa : Measurable Va) (hUs : MeasurableSet Us)
    (hmpa : ∀ x : State, Nonempty (MPFlow M 0 Va x))
    (hForcea : ForcedNonexpansion M Va)
    (hNonnega : NonnegV Va) {C : ℝ} (hC0 : 0 ≤ C) (hVC : ∀ y, Va y ≤ C)
    (hFloora : ∀ y ∈ Us, 1 ≤ Va y)
    (hVd : Measurable Vd) (hθ : 0 ≤ θ)
    (hmpd : ∀ x : State, Nonempty (MPFlow M θ Vd x))
    (hForced : ForcedNonexpansion M Vd)
    (hNonnegd : NonnegV Vd) (hε : 0 < ε) (hFloord : ∀ y, ε ≤ Vd y)
    (x : State) :
    Battery.pathMeasure (tskel M) (TSkelState.alive (0, x))
        {ω | ∃ n, ω n ∈ tbadIn Us ∨ T < tclock (ω n)}
      ≤ ENNReal.ofReal (Va x)
        + ENNReal.ofReal (Real.exp (-(θ * T)) * Vd x / ε)
\end{leansrc}
\par\medskip\noindent{\small \texttt{Qtsl.LayerD.}\texttt{\bfseries sound\_GSHS\_until\_battery\_mp}}\hfill{\footnotesize \path{lean/qtsl-aux/QtslAux/LayerD/MPDischarge.lean}, line 644}\par\nopagebreak
\begin{leansrc}
theorem sound_GSHS_until_battery_mp
    [IsMarkovKernel M.seg] [IsMarkovKernel M.jump]
    {Va Vd : State → ℝ} {Us : Set State} {θ ε p thra thrd : ℝ} {T : ℝ≥0}
    (hVa : Measurable Va) (hUs : MeasurableSet Us)
    (hmpa : ∀ x : State, Nonempty (MPFlow M 0 Va x))
    (hForcea : ForcedNonexpansion M Va)
    (hNonnega : NonnegV Va) {C : ℝ} (hC0 : 0 ≤ C) (hVC : ∀ y, Va y ≤ C)
    (hFloora : ∀ y ∈ Us, 1 ≤ Va y)
    (hVd : Measurable Vd) (hθ : 0 ≤ θ)
    (hmpd : ∀ x : State, Nonempty (MPFlow M θ Vd x))
    (hForced : ForcedNonexpansion M Vd)
    (hNonnegd : NonnegV Vd) (hε : 0 < ε) (hFloord : ∀ y, ε ≤ Vd y)
    (hΘ : p + (thra + Real.exp (-(θ * T)) * thrd / ε) ≤ 1)
    (x : State) (hthra : Va x ≤ thra) (hthrd : Vd x ≤ thrd) :
    ENNReal.ofReal p ≤ Battery.pathMeasure (tskel M) (TSkelState.alive (0, x))
      {ω | ∀ n, ω n ∉ tbadIn Us ∧ tclock (ω n) ≤ T}
\end{leansrc}
\par\medskip\noindent{\small \texttt{SuFunctor.Site.}\texttt{\bfseries untilEvent}}\hfill{\footnotesize \path{lean/su-functor/SuFunctor/Site/FellAcceptance.lean}, line 343}\par\nopagebreak
\begin{leansrc}
def untilEvent (Reach Avoid : ℕ → Set α) : Set (FellCloseds α)
\end{leansrc}
\par\medskip\noindent{\small \texttt{SuFunctor.Site.}\texttt{\bfseries measurableSet\_untilEvent}}\hfill{\footnotesize \path{lean/su-functor/SuFunctor/Site/FellAcceptance.lean}, line 353}\par\nopagebreak
\begin{leansrc}
theorem measurableSet_untilEvent {Reach Avoid : ℕ → Set α}
    (h₁ : ∀ n, MeasurableSet (hitEvent (Reach n)))
    (h₂ : ∀ n, MeasurableSet (missEvent (Avoid n))) :
    MeasurableSet (untilEvent Reach Avoid)
\end{leansrc}
\par\medskip\noindent{\small \texttt{Qtsl.Battery.}\texttt{\bfseries sound\_QUntilBattery}}\hfill{\footnotesize \path{lean/qtsl-aux/QtslAux/QtslBattery.lean}, line 371}\par\nopagebreak
\begin{leansrc}
theorem sound_QUntilBattery {sat_s sat : S → Prop}
    (hD : MeasurableSet D) (hDbar : MeasurableSet Dbar)
    (hA : MeasurableSet A) (hR : MeasurableSet R)
    (hVa : Measurable Va) (hVd : Measurable Vd)
    (hFs : ∀ y, ¬ sat_s y → y ∈ A) (hFt : ∀ y ∈ R, sat y)
    (hβ : 0 ≤ β) (hc : β ^ T ≤ c)
    (hDsub : D ⊆ Dbar)
    (hR3 : ∀ y ∈ D, y ∉ A ∪ R → K y Dbar = 1)
    (hR4 : ∀ y ∈ D, y ∉ A ∪ R → K y (D ∪ (A ∪ R)) = 1)
    (hVaint : ∀ y ∈ D, y ∉ A ∪ R → Integrable Va (K y))
    (hVadrift : ∀ y ∈ D, y ∉ A ∪ R → ∫ z, Va z ∂K y ≤ Va y)
    (hVafloor : ∀ y ∈ A, y ∈ Dbar → 1 ≤ Va y)
    (hVanonneg : ∀ y ∈ Dbar, 0 ≤ Va y)
    (hVdint : ∀ y ∈ D, y ∉ A ∪ R → Integrable Vd (K y))
    (hVddecay : ∀ y ∈ D, y ∉ A ∪ R → ∫ z, Vd z ∂K y ≤ β * Vd y)
    (hVdfloor : ∀ y ∈ Dbar, y ∉ A ∪ R → 1 ≤ Vd y)
    (hVdnonneg : ∀ y ∈ Dbar, 0 ≤ Vd y)
    (hVdM : ∀ y ∈ Dbar, Vd y ≤ M)
    (hx : x ∈ D) (hthr : Va x + c * Vd x ≤ 1 - p) :
    ENNReal.ofReal p ≤ pathMeasure K x
      {ω | ∃ n ≤ T, sat (ω n) ∧ ∀ j < n, sat_s (ω j)}
\end{leansrc}
\par\medskip\noindent{\small \texttt{Qtsl.Battery.}\texttt{\bfseries measure\_firstBad\_le}}\hfill{\footnotesize \path{lean/qtsl-aux/QtslAux/QtslBattery.lean}, line 110}\par\nopagebreak
\begin{leansrc}
theorem measure_firstBad_le {K : Kernel S S} [IsMarkovKernel K]
    {f : S → ℝ} {Stop Bad : Set S} {C : ℝ}
    (hStop : MeasurableSet Stop) (hf : Measurable f)
    (hf0 : ∀ y, 0 ≤ f y) (hfC : ∀ y, f y ≤ C)
    (hf1 : ∀ y ∈ Bad, 1 ≤ f y)
    (hdrift : ∀ y, y ∉ Stop → ∫ z, f z ∂K y ≤ f y) (x : S) :
    pathMeasure K x {ω | ∃ n, ω n ∈ Bad ∧ ∀ j < n, ω j ∉ Stop}
      ≤ ENNReal.ofReal (f x)
\end{leansrc}
\par\medskip\noindent{\small \texttt{Qtsl.Battery.}\texttt{\bfseries measure\_missTarget\_le}}\hfill{\footnotesize \path{lean/qtsl-aux/QtslAux/QtslBattery.lean}, line 143}\par\nopagebreak
\begin{leansrc}
theorem measure_missTarget_le {K : Kernel S S} [IsMarkovKernel K]
    {D Dbar B : Set S} {V : S → ℝ} {β : ℝ} {x : S} {T : ℕ} {M : ℝ}
    (hD : MeasurableSet D) (hDbar : MeasurableSet Dbar) (hB : MeasurableSet B)
    (hV : Measurable V) (hβ : 0 ≤ β)
    (hDsub : D ⊆ Dbar)
    (hR3 : ∀ y ∈ D, y ∉ B → K y Dbar = 1)
    (hR4 : ∀ y ∈ D, y ∉ B → K y (D ∪ B) = 1)
    (hint : ∀ y ∈ D, y ∉ B → Integrable V (K y))
    (hdecay : ∀ y ∈ D, y ∉ B → ∫ z, V z ∂K y ≤ β * V y)
    (hfloor : ∀ y ∈ Dbar, y ∉ B → 1 ≤ V y)
    (hnonneg : ∀ y ∈ Dbar, 0 ≤ V y)
    (hM : ∀ y ∈ Dbar, V y ≤ M)
    (hx : x ∈ D) :
    pathMeasure K x {ω | ∀ k ≤ T, ω k ∉ B} ≤ ENNReal.ofReal (β ^ T * V x)
\end{leansrc}
\par\medskip\noindent{\small \texttt{MeasureTheory.Supermartingale.}\texttt{\bfseries ville\_ineq}}\hfill{\footnotesize \path{lean/qtsl-aux/QtslAux/Ville.lean}, line 133}\par\nopagebreak
\begin{leansrc}
theorem Supermartingale.ville_ineq [IsFiniteMeasure μ] (hsup : Supermartingale f 𝒢
      μ)
    (hnonneg : 0 ≤ f) (ε : ℝ≥0) :
    ε * μ {ω | ∃ k, (ε : ℝ) ≤ f k ω} ≤ ENNReal.ofReal (∫ ω, f 0 ω ∂μ)
\end{leansrc}
\par\medskip\noindent{\small \texttt{Qtsl.QStreett.}\texttt{\bfseries absorptionBound\_of\_kernel}}\hfill{\footnotesize \path{lean/qtsl-aux/QtslAux/QStreett.lean}, line 804}\par\nopagebreak
\begin{leansrc}
theorem absorptionBound_of_kernel {A J Rq : Set S} {V : S → ℝ} {C γ : ℝ} {x : S}
    (hRq : MeasurableSet Rq) (hV : Measurable V)
    (hV0 : ∀ y, 0 ≤ V y) (hVC : ∀ y, V y ≤ C)
    (hVfloor : ∀ y ∈ Rq, 1 ≤ V y)
    (hVdrift : ∀ y, y ∉ Rq → ∫ z, V z ∂K y ≤ V y)
    (hVJ : ∀ y ∈ J, V y ≤ 1 - γ) :
    AbsorptionBound K x A J Rq γ
\end{leansrc}
\par\medskip\noindent{\small \texttt{Qtsl.LayerD.}\texttt{\bfseries sound\_GSHS\_streett\_bundle\_mp}}\hfill{\footnotesize \path{lean/qtsl-aux/QtslAux/LayerD/MPDischarge.lean}, line 612}\par\nopagebreak
\begin{leansrc}
theorem sound_GSHS_streett_bundle_mp
    {ι : Type*} [Countable ι]
    {D Dbar U : Set State} {V₀ : State → ℝ} {Rq G : ι → Set State}
    {W : ι → State → ℝ} {p thr C₀ : ℝ} {ε C : ι → ℝ} {x : State}
    [IsMarkovKernel M.seg] [IsMarkovKernel M.jump]
    (hD : MeasurableSet D) (hDbar : MeasurableSet Dbar) (hU : MeasurableSet U)
    (hV₀ : Measurable V₀)
    (hRq : ∀ i, MeasurableSet (Rq i)) (hG : ∀ i, MeasurableSet (G i))
    (hW : ∀ i, Measurable (W i))
    (hΘ : p + thr ≤ 1) (hDsub : D ⊆ Dbar)
    (hLive : NoVanish M (D \ U)) (hHalt : HaltsCaught M (D \ U) (U ∩ Dbar))
    (hReset : ResetsInto M (D \ U) ((D ∪ U) ∩ Dbar))
    (hmp₀ : ∀ x : State, Nonempty (MPFlow M 0 V₀ x))
    (hForce₀ : ForcedNonexpansion M V₀) (hNonneg₀ : NonnegV V₀)
    (hC₀0 : 0 ≤ C₀) (hV₀C : ∀ y, V₀ y ≤ C₀)
    (hFloor₀ : ∀ y ∈ U, y ∈ Dbar → 1 ≤ V₀ y)
    (hε : ∀ i, 0 < ε i)
    (hmpW : ∀ i, ∀ x : State, Nonempty (MPFlow M 0 (W i) x))
    (hSegDWR : ∀ i, RankingSegmentDynkin M (((D \ U) ∩ Rq i) \ G i) (W i) (ε i))
    (hForceW : ∀ i, ForcedNonexpansion M (W i)) (hNonnegW : ∀ i, NonnegV (W i))
    (hC0 : ∀ i, 0 ≤ C i) (hWC : ∀ i, ∀ y, W i y ≤ C i)
    (hx : x ∈ D) (hthr : V₀ x ≤ thr) :
    ENNReal.ofReal p ≤ Battery.pathMeasure (skel M) (SkelState.alive x)
      {ω | ∀ i, (∃ᶠ n in atTop, ω n ∈ skelIn (Rq i)) →
        ∃ᶠ n in atTop, ω n ∈ skelIn (G i)}
\end{leansrc}
\par\medskip\noindent{\small \texttt{Qtsl.LayerD.}\texttt{\bfseries sound\_GSHS\_streett\_bundle\_mp\_strict}}\hfill{\footnotesize \path{lean/qtsl-aux/QtslAux/LayerD/MPDischarge.lean}, line 1074}\par\nopagebreak
\begin{leansrc}
theorem sound_GSHS_streett_bundle_mp_strict
    {ι : Type*} [Countable ι]
    {D Dbar U : Set State} {V₀ : State → ℝ} {Rq G : ι → Set State}
    {W : ι → State → ℝ} {p thr C₀ : ℝ} {ε C : ι → ℝ} {x : State}
    [IsMarkovKernel M.seg] [IsMarkovKernel M.jump]
    (hD : MeasurableSet D) (hDbar : MeasurableSet Dbar) (hU : MeasurableSet U)
    (hV₀ : Measurable V₀)
    (hRq : ∀ i, MeasurableSet (Rq i)) (hG : ∀ i, MeasurableSet (G i))
    (hW : ∀ i, Measurable (W i))
    (hΘ : p + thr ≤ 1) (hDsub : D ⊆ Dbar)
    (hLive : NoVanish M (D \ U)) (hHalt : HaltsCaught M (D \ U) (U ∩ Dbar))
    (hReset : ResetsInto M (D \ U) ((D ∪ U) ∩ Dbar))
    (hmp₀ : ∀ x : State, Nonempty (MPFlow M 0 V₀ x))
    (hForce₀ : ForcedNonexpansion M V₀) (hNonneg₀ : NonnegV V₀)
    (hC₀0 : 0 ≤ C₀) (hV₀C : ∀ y, V₀ y ≤ C₀)
    (hFloor₀ : ∀ y ∈ U, y ∈ Dbar → 1 ≤ V₀ y)
    (hε : ∀ i, 0 < ε i)
    (hmpW : ∀ i, ∀ x : State, Nonempty (MPFlow M 0 (W i) x))
    (hmpWR : ∀ i, ∀ x ∈ ((D \ U) ∩ Rq i) \ G i,
      Nonempty (MPFlowStrict M 0 (W i) x (ε i)))
    (hForceW : ∀ i, ForcedNonexpansion M (W i)) (hNonnegW : ∀ i, NonnegV (W i))
    (hC0 : ∀ i, 0 ≤ C i) (hWC : ∀ i, ∀ y, W i y ≤ C i)
    (hx : x ∈ D) (hthr : V₀ x ≤ thr) :
    ENNReal.ofReal p ≤ Battery.pathMeasure (skel M) (SkelState.alive x)
      {ω | ∀ i, (∃ᶠ n in atTop, ω n ∈ skelIn (Rq i)) →
        ∃ᶠ n in atTop, ω n ∈ skelIn (G i)}
\end{leansrc}
\par\medskip\noindent{\small \texttt{Qtsl.LayerD.}\texttt{\bfseries rankingSegmentDynkin\_of\_mp\_strict}}\hfill{\footnotesize \path{lean/qtsl-aux/QtslAux/LayerD/MPDischarge.lean}, line 1060}\par\nopagebreak
\begin{leansrc}
theorem rankingSegmentDynkin_of_mp_strict (M : SegmentModel State)
    (V : State → ℝ) [IsSFiniteKernel M.jump] (hV : Measurable V)
    {Sc : Set State} {ε : ℝ} (hε : 0 ≤ ε)
    (h : ∀ x ∈ Sc, Nonempty (MPFlowStrict M 0 V x ε)) :
    RankingSegmentDynkin M Sc V ε
\end{leansrc}
\par\medskip\noindent{\small \texttt{Qtsl.LayerD.}\texttt{\bfseries MPFlowStrict}}\hfill{\footnotesize \path{lean/qtsl-aux/QtslAux/LayerD/MPDischarge.lean}, line 1047}\par\nopagebreak
\begin{leansrc}
structure MPFlowStrict (M : SegmentModel State) (θ : ℝ) (V : State → ℝ)
    (x : State) (ε : ℝ) extends MPFlow M θ V x where …
\end{leansrc}
\par\medskip\noindent{\small \texttt{Qtsl.Excursions.}\texttt{\bfseries streett\_ae\_frequently\_response}}\hfill{\footnotesize \path{lean/qtsl-aux/QtslAux/Excursions.lean}, line 365}\par\nopagebreak
\begin{leansrc}
theorem streett_ae_frequently_response (hε : 0 < ε₀)
    (hint : ∀ n, Integrable (V n) μ) (hnonneg : 0 ≤ V)
    (hmR : ∀ n, MeasurableSet[𝒢 n] {ω | X n ω ∈ R})
    (hmG : ∀ n, MeasurableSet[𝒢 n] {ω | X n ω ∈ G})
    (hdrift : ∀ n, ∀ᵐ ω ∂μ, X n ω ∉ G → (μ[V (n + 1) | 𝒢 n]) ω ≤ V n ω)
    (hdriftR : ∀ n, ∀ᵐ ω ∂μ,
      X n ω ∈ R → X n ω ∉ G → (μ[V (n + 1) | 𝒢 n]) ω ≤ V n ω - ε₀) :
    ∀ᵐ ω ∂μ, (∃ᶠ n in atTop, X n ω ∈ R) → ∃ᶠ n in atTop, X n ω ∈ G
\end{leansrc}
\par\medskip\noindent{\small \texttt{Qtsl.Excursions.}\texttt{\bfseries streett\_window\_sum\_le}}\hfill{\footnotesize \path{lean/qtsl-aux/QtslAux/Excursions.lean}, line 319}\par\nopagebreak
\begin{leansrc}
theorem streett_window_sum_le (hε : 0 ≤ ε₀)
    (hint : ∀ n, Integrable (V n) μ) (hnonneg : 0 ≤ V)
    (hmR : ∀ n, MeasurableSet[𝒢 n] {ω | X n ω ∈ R})
    (hmG : ∀ n, MeasurableSet[𝒢 n] {ω | X n ω ∈ G})
    (hdrift : ∀ n, ∀ᵐ ω ∂μ, X n ω ∉ G → (μ[V (n + 1) | 𝒢 n]) ω ≤ V n ω)
    (hdriftR : ∀ n, ∀ᵐ ω ∂μ,
      X n ω ∈ R → X n ω ∉ G → (μ[V (n + 1) | 𝒢 n]) ω ≤ V n ω - ε₀)
    (N : ℕ) :
    ENNReal.ofReal ε₀ * ∑' j : ℕ,
        μ ({ω | ∀ k, N ≤ k → k < N + j + 1 → X k ω ∉ G} ∩ {ω | X (N + j) ω ∈ R})
      ≤ ENNReal.ofReal (∫ ω, V N ω ∂μ)
\end{leansrc}
\par\medskip\noindent{\small \texttt{Qtsl.Excursions.}\texttt{\bfseries ranking\_hitting\_sum\_le}}\hfill{\footnotesize \path{lean/qtsl-aux/QtslAux/Excursions.lean}, line 147}\par\nopagebreak
\begin{leansrc}
theorem ranking_hitting_sum_le (hε : 0 ≤ ε₀)
    (hint : ∀ n, Integrable (V n) μ) (hnonneg : 0 ≤ V)
    (hmB : ∀ n, MeasurableSet[𝒢 n] {ω | X n ω ∈ B})
    (hdrift : ∀ n, ∀ᵐ ω ∂μ, X n ω ∉ B → (μ[V (n + 1) | 𝒢 n]) ω ≤ V n ω - ε₀) :
    ENNReal.ofReal ε₀ * ∑' k, μ {ω | ∀ j < k + 1, X j ω ∉ B}
      ≤ ENNReal.ofReal (∫ ω, V 0 ω ∂μ)
\end{leansrc}
\par\medskip\noindent{\small \texttt{Qtsl.Excursions.}\texttt{\bfseries ae\_hits\_of\_ranking}}\hfill{\footnotesize \path{lean/qtsl-aux/QtslAux/Excursions.lean}, line 174}\par\nopagebreak
\begin{leansrc}
theorem ae_hits_of_ranking (hε : 0 < ε₀)
    (hint : ∀ n, Integrable (V n) μ) (hnonneg : 0 ≤ V)
    (hmB : ∀ n, MeasurableSet[𝒢 n] {ω | X n ω ∈ B})
    (hdrift : ∀ n, ∀ᵐ ω ∂μ, X n ω ∉ B → (μ[V (n + 1) | 𝒢 n]) ω ≤ V n ω - ε₀) :
    ∀ᵐ ω ∂μ, ∃ n, X n ω ∈ B
\end{leansrc}
\par\medskip\noindent{\small \texttt{MeasureTheory.Supermartingale.}\texttt{\bfseries exists\_ae\_tendsto\_of\_nonneg}}\hfill{\footnotesize \path{lean/qtsl-aux/QtslAux/Excursions.lean}, line 42}\par\nopagebreak
\begin{leansrc}
theorem Supermartingale.exists_ae_tendsto_of_nonneg [IsFiniteMeasure μ] {f : ℕ → Ω
      → ℝ}
    (hf : Supermartingale f 𝒢 μ) (hnonneg : 0 ≤ f) :
    ∀ᵐ ω ∂μ, ∃ c, Tendsto (fun n => f n ω) atTop (𝓝 c)
\end{leansrc}
\par\medskip\noindent{\small \texttt{Qtsl.QStreett.}\texttt{\bfseries streett\_event\_mono}}\hfill{\footnotesize \path{lean/qtsl-aux/QtslAux/QStreett.lean}, line 418}\par\nopagebreak
\begin{leansrc}
theorem streett_event_mono {Rq Rqhat G Ghat : Set S} (hR : Rq ⊆ Rqhat)
    (hG : Ghat ⊆ G) (ω : ℕ → S)
    (h : (∃ᶠ n in atTop, ω n ∈ Rqhat) → ∃ᶠ n in atTop, ω n ∈ Ghat) :
    (∃ᶠ n in atTop, ω n ∈ Rq) → ∃ᶠ n in atTop, ω n ∈ G
\end{leansrc}
\par\medskip\noindent{\small \texttt{Qtsl.QStreett.}\texttt{\bfseries streett\_pair\_localized}}\hfill{\footnotesize \path{lean/qtsl-aux/QtslAux/QStreett.lean}, line 244}\par\nopagebreak
\begin{leansrc}
theorem streett_pair_localized {D Dbar U Rq G : Set S} {W : S → ℝ} {ε C : ℝ} {x : S}
    (hD : MeasurableSet D) (hDbar : MeasurableSet Dbar) (hU : MeasurableSet U)
    (hRq : MeasurableSet Rq) (hG : MeasurableSet G) (hW : Measurable W)
    (hε : 0 < ε) (hDsub : D ⊆ Dbar)
    (hR3 : ∀ y ∈ D, y ∉ U → K y Dbar = 1)
    (hW0 : ∀ y ∈ Dbar, 0 ≤ W y) (hWC : ∀ y ∈ Dbar, W y ≤ C)
    (hWint : ∀ y ∈ D, y ∉ U → y ∉ G → Integrable W (K y))
    (hWdrift : ∀ y ∈ D, y ∉ U → y ∉ G → ∫ z, W z ∂K y ≤ W y)
    (hWdriftR : ∀ y ∈ D, y ∉ U → y ∈ Rq → y ∉ G → ∫ z, W z ∂K y ≤ W y - ε) :
    ∀ᵐ ω ∂pathMeasure K x, (∀ n, ω n ∈ D \ U) →
      (∃ᶠ n in atTop, ω n ∈ Rq) → ∃ᶠ n in atTop, ω n ∈ G
\end{leansrc}
\par\medskip\noindent{\small \texttt{Qtsl.Excursions.}\texttt{\bfseries streett\_window\_invariant}}\hfill{\footnotesize \path{lean/qtsl-aux/QtslAux/Excursions.lean}, line 201}\par\nopagebreak
\begin{leansrc}
theorem streett_window_invariant
    (hint : ∀ n, Integrable (V n) μ) (hnonneg : 0 ≤ V)
    (hmR : ∀ n, MeasurableSet[𝒢 n] {ω | X n ω ∈ R})
    (hmG : ∀ n, MeasurableSet[𝒢 n] {ω | X n ω ∈ G})
    (hdrift : ∀ n, ∀ᵐ ω ∂μ, X n ω ∉ G → (μ[V (n + 1) | 𝒢 n]) ω ≤ V n ω)
    (hdriftR : ∀ n, ∀ᵐ ω ∂μ,
      X n ω ∈ R → X n ω ∉ G → (μ[V (n + 1) | 𝒢 n]) ω ≤ V n ω - ε₀)
    (N : ℕ) : ∀ n, N ≤ n →
    ∫ ω in {ω | ∀ k, N ≤ k → k < n → X k ω ∉ G}, V n ω ∂μ
        + ε₀ * ∑ k ∈ Ico N n,
            μ.real ({ω | ∀ j, N ≤ j → j < k + 1 → X j ω ∉ G} ∩ {ω | X k ω ∈ R})
      ≤ ∫ ω, V N ω ∂μ
\end{leansrc}
\par\medskip\noindent{\small \texttt{Qtsl.QStreett.}\texttt{\bfseries AbsorptionBound}}\hfill{\footnotesize \path{lean/qtsl-aux/QtslAux/QStreett.lean}, line 479}\par\nopagebreak
\begin{leansrc}
def AbsorptionBound (K : Kernel S S) [IsMarkovKernel K] (x : S)
    (A J Rq : Set S) (γ : ℝ) : Prop
\end{leansrc}
\par\medskip\noindent{\small \texttt{Qtsl.QStreett.}\texttt{\bfseries ae\_not\_frequently\_of\_absorption}}\hfill{\footnotesize \path{lean/qtsl-aux/QtslAux/QStreett.lean}, line 490}\par\nopagebreak
\begin{leansrc}
theorem ae_not_frequently_of_absorption {A J Rq : Set S} {γ : ℝ} {x : S}
    (hRq : MeasurableSet Rq) (hγ : 0 < γ)
    (habs : AbsorptionBound K x A J Rq γ) :
    ∀ᵐ ω ∂pathMeasure K x,
      (∃ᶠ k in atTop, ω ∈ survive A k ∧ ω k ∈ J) →
        ¬ ∃ᶠ n in atTop, ω n ∈ Rq
\end{leansrc}
\par\medskip\noindent{\small \texttt{Qtsl.QStreett.}\texttt{\bfseries supermartingale\_stopped\_shift}}\hfill{\footnotesize \path{lean/qtsl-aux/QtslAux/QStreett.lean}, line 740}\par\nopagebreak
\begin{leansrc}
theorem supermartingale_stopped_shift {f : S → ℝ} {R : Set S} {C : ℝ}
    (hR : MeasurableSet R) (hf : Measurable f) (hf0 : ∀ y, 0 ≤ f y) (hfC : ∀ y, f y
      ≤ C)
    (hdrift : ∀ y, y ∉ R → ∫ z, f z ∂K y ≤ f y) (k : ℕ) (u : Π _i : Iic k, S) :
    Supermartingale (fun n ω => stopped f R n (shift k ω)) (piLEShift k)
      (Kernel.traj (X := fun _ : ℕ => S) (κ := mchain K) k u)
\end{leansrc}
\par\medskip\noindent{\small \texttt{Qtsl.QStreett.}\texttt{\bfseries measure\_reach\_after\_le}}\hfill{\footnotesize \path{lean/qtsl-aux/QtslAux/QStreett.lean}, line 760}\par\nopagebreak
\begin{leansrc}
theorem measure_reach_after_le {f : S → ℝ} {R : Set S} {C : ℝ}
    (hR : MeasurableSet R) (hf : Measurable f) (hf0 : ∀ y, 0 ≤ f y) (hfC : ∀ y, f y
      ≤ C)
    (hf1 : ∀ y ∈ R, 1 ≤ f y) (hdrift : ∀ y, y ∉ R → ∫ z, f z ∂K y ≤ f y)
    (k : ℕ) (u : Π _i : Iic k, S) :
    Kernel.traj (X := fun _ : ℕ => S) (κ := mchain K) k u {ω | ∃ n, k ≤ n ∧ ω n ∈ R}
      ≤ ENNReal.ofReal (f (u ⟨k, mem_Iic.2 le_rfl⟩))
\end{leansrc}
\par\medskip\noindent{\small \texttt{PPLC.Qtsl.Behavior.}\texttt{\bfseries AService}}\hfill{\footnotesize \path{lean/qtsl/Qtsl/TimedCore.lean}, line 32}\par\nopagebreak
\begin{leansrc}
def Behavior.AService (B : Behavior X) (T : ℝ≥0) (φ : Set X) (x : X) : Prop
\end{leansrc}
\par\medskip\noindent{\small \texttt{PPLC.Qtsl.Behavior.}\texttt{\bfseries BoxInv}}\hfill{\footnotesize \path{lean/qtsl/Qtsl/TimedCore.lean}, line 35}\par\nopagebreak
\begin{leansrc}
def Behavior.BoxInv (B : Behavior X) (φ : Set X) (x : X) : Prop
\end{leansrc}
\par\medskip\noindent{\small \texttt{PPLC.Qtsl.Behavior.}\texttt{\bfseries FusionClosed}}\hfill{\footnotesize \path{lean/qtsl/Qtsl/TimedCore.lean}, line 28}\par\nopagebreak
\begin{leansrc}
def Behavior.FusionClosed (B : Behavior X) : Prop
\end{leansrc}
\par\medskip\noindent{\small \texttt{PPLC.Qtsl.Behavior.}\texttt{\bfseries FirstVisits}}\hfill{\footnotesize \path{lean/qtsl/Qtsl/TimedCore.lean}, line 38}\par\nopagebreak
\begin{leansrc}
def Behavior.FirstVisits (B : Behavior X) (φ : Set X) : Prop
\end{leansrc}
\par\medskip\noindent{\small \texttt{PPLC.Qtsl.Behavior.}\texttt{\bfseries fsTarget}}\hfill{\footnotesize \path{lean/qtsl/Qtsl/TimedCore.lean}, line 41}\par\nopagebreak
\begin{leansrc}
def Behavior.fsTarget (B : Behavior X) {n : ℕ}
    (α : Fin n → Set X) (T : Fin n → ℝ≥0) : Set X
\end{leansrc}
\par\medskip\noindent{\small \texttt{PPLC.Qtsl.Behavior.}\texttt{\bfseries FirstVisits}}\hfill{\footnotesize \path{lean/qtsl/Qtsl/TimedCore.lean}, line 38}\par\nopagebreak
\begin{leansrc}
def Behavior.FirstVisits (B : Behavior X) (φ : Set X) : Prop
\end{leansrc}
\par\medskip\noindent{\small \texttt{PPLC.Qtsl.Behavior.}\texttt{\bfseries residualCell}}\hfill{\footnotesize \path{lean/qtsl/Qtsl/TimedCore.lean}, line 45}\par\nopagebreak
\begin{leansrc}
def Behavior.residualCell (B : Behavior X) {n : ℕ}
    (α : Fin n → Set X) (T : Fin n → ℝ≥0) (c : ℝ≥0) : Set X
\end{leansrc}
\par\medskip\noindent{\small \texttt{Qtsl.NNF.}\texttt{\bfseries denot\_elimNeg}}\hfill{\footnotesize \path{lean/qtsl-aux/QtslAux/NNFDuality.lean}, line 198}\par\nopagebreak
\begin{leansrc}
theorem denot_elimNeg {σ : Type*} (den : α → Set σ) (d : α → α)
    (hd : ∀ a, den (d a) = (den a)ᶜ) :
    ∀ φ : NForm α, (φ.elimNeg d).denot den = φ.denot den
\end{leansrc}
\par\medskip\noindent{\small \texttt{Qtsl.QStreett.}\texttt{\bfseries ranking\_guard\_collapse}}\hfill{\footnotesize \path{lean/qtsl-aux/QtslAux/QStreett.lean}, line 464}\par\nopagebreak
\begin{leansrc}
theorem ranking_guard_collapse {W : S → ℝ} {T slab : Set S} {g : S → S}
    (hzero : ∀ y ∈ T, W y = 0) (himg : ∀ y ∈ slab, g y ∈ T) :
    ∀ y ∈ slab, 0 ≤ W y → W (g y) ≤ W y
\end{leansrc}
\par\medskip\noindent{\small \texttt{Qtsl.NNF.}\texttt{\bfseries compl\_le\_atom}}\hfill{\footnotesize \path{lean/qtsl-aux/QtslAux/NNFDuality.lean}, line 36}\par\nopagebreak
\begin{leansrc}
theorem compl_le_atom : {x | f x ≤ c}ᶜ = {x | c < f x}
\end{leansrc}
\par\medskip\noindent{\small \texttt{Qtsl.NNF.}\texttt{\bfseries compl\_lt\_atom}}\hfill{\footnotesize \path{lean/qtsl-aux/QtslAux/NNFDuality.lean}, line 41}\par\nopagebreak
\begin{leansrc}
theorem compl_lt_atom : {x | f x < c}ᶜ = {x | c ≤ f x}
\end{leansrc}
\par\medskip\noindent{\small \texttt{Qtsl.NNF.}\texttt{\bfseries compl\_ge\_atom}}\hfill{\footnotesize \path{lean/qtsl-aux/QtslAux/NNFDuality.lean}, line 46}\par\nopagebreak
\begin{leansrc}
theorem compl_ge_atom : {x | c ≤ f x}ᶜ = {x | f x < c}
\end{leansrc}
\par\medskip\noindent{\small \texttt{Qtsl.NNF.}\texttt{\bfseries compl\_gt\_atom}}\hfill{\footnotesize \path{lean/qtsl-aux/QtslAux/NNFDuality.lean}, line 51}\par\nopagebreak
\begin{leansrc}
theorem compl_gt_atom : {x | c < f x}ᶜ = {x | f x ≤ c}
\end{leansrc}
\par\medskip\noindent{\small \texttt{Qtsl.NNF.}\texttt{\bfseries denot\_nnfAux}}\hfill{\footnotesize \path{lean/qtsl-aux/QtslAux/NNFDuality.lean}, line 111}\par\nopagebreak
\begin{leansrc}
theorem denot_nnfAux : ∀ (φ : Form α) (b : Bool),
    (φ.nnfAux b).denot den = if b then φ.denot den else (φ.denot den)ᶜ
\end{leansrc}
\par\medskip\noindent{\small \texttt{Qtsl.NNF.Form.}\texttt{\bfseries nnf}}\hfill{\footnotesize \path{lean/qtsl-aux/QtslAux/NNFDuality.lean}, line 87}\par\nopagebreak
\begin{leansrc}
def Form.nnf (φ : Form α) : NForm α
\end{leansrc}
\par\medskip\noindent{\small \texttt{Qtsl.NNF.}\texttt{\bfseries denot\_nnf}}\hfill{\footnotesize \path{lean/qtsl-aux/QtslAux/NNFDuality.lean}, line 128}\par\nopagebreak
\begin{leansrc}
theorem denot_nnf (φ : Form α) : (φ.nnf).denot den = φ.denot den
\end{leansrc}
\par\medskip\noindent{\small \texttt{Qtsl.NNF.}\texttt{\bfseries negLits\_nnf}}\hfill{\footnotesize \path{lean/qtsl-aux/QtslAux/NNFDuality.lean}, line 170}\par\nopagebreak
\begin{leansrc}
theorem negLits_nnf (φ : Form α) : (φ.nnf).negLits = φ.occs false
\end{leansrc}
\par\medskip\noindent{\small \texttt{Qtsl.NNF.}\texttt{\bfseries nnf\_no\_negLit\_iff}}\hfill{\footnotesize \path{lean/qtsl-aux/QtslAux/NNFDuality.lean}, line 175}\par\nopagebreak
\begin{leansrc}
theorem nnf_no_negLit_iff (φ : Form α) (P : α → Prop) :
    (∀ a ∈ (φ.nnf).negLits, ¬ P a) ↔ ∀ a ∈ φ.occs false, ¬ P a
\end{leansrc}
\par\medskip\noindent{\small \texttt{Qtsl.NNF.}\texttt{\bfseries negLits\_elimNeg}}\hfill{\footnotesize \path{lean/qtsl-aux/QtslAux/NNFDuality.lean}, line 188}\par\nopagebreak
\begin{leansrc}
theorem negLits_elimNeg (d : α → α) : ∀ φ : NForm α, (φ.elimNeg d).negLits = ∅
\end{leansrc}
\par\medskip\noindent{\small \texttt{Qtsl.NNF.}\texttt{\bfseries negLits\_nnfAux}}\hfill{\footnotesize \path{lean/qtsl-aux/QtslAux/NNFDuality.lean}, line 153}\par\nopagebreak
\begin{leansrc}
theorem negLits_nnfAux : ∀ (φ : Form α) (b : Bool),
    (φ.nnfAux b).negLits = φ.occs (!b)
\end{leansrc}
\par\medskip\noindent{\small \texttt{Qtsl.Aliases.}\texttt{\bfseries sugarAvoid\_denot}}\hfill{\footnotesize \path{lean/qtsl-aux/QtslAux/QtslAliases.lean}, line 95}\par\nopagebreak
\begin{leansrc}
theorem sugarAvoid_denot {S : Type*} (B : Set S) :
    {ω : ℕ → S | ∀ n, ω n ∉ B} = {ω : ℕ → S | ∀ n, ω n ∈ Bᶜ}
\end{leansrc}
\par\medskip\noindent{\small \texttt{Qtsl.Aliases.}\texttt{\bfseries sugarGF\_denot}}\hfill{\footnotesize \path{lean/qtsl-aux/QtslAux/QtslAliases.lean}, line 98}\par\nopagebreak
\begin{leansrc}
theorem sugarGF_denot {S : Type*} (G : Set S) (ω : ℕ → S) :
    ((∃ᶠ n in atTop, ω n ∈ (Set.univ : Set S)) → ∃ᶠ n in atTop, ω n ∈ G)
      ↔ ∃ᶠ n in atTop, ω n ∈ G
\end{leansrc}
\par\medskip\noindent{\small \texttt{Qtsl.Aliases.}\texttt{\bfseries sugarFG\_denot}}\hfill{\footnotesize \path{lean/qtsl-aux/QtslAux/QtslAliases.lean}, line 105}\par\nopagebreak
\begin{leansrc}
theorem sugarFG_denot {S : Type*} (B : Set S) (ω : ℕ → S) :
    ((∃ᶠ n in atTop, ω n ∈ B) → ∃ᶠ n in atTop, ω n ∈ (∅ : Set S))
      ↔ ∀ᶠ n in atTop, ω n ∉ B
\end{leansrc}
\par\medskip\noindent{\small \texttt{Qtsl.Aliases.}\texttt{\bfseries negF\_denot\_subset}}\hfill{\footnotesize \path{lean/qtsl-aux/QtslAux/QtslAliases.lean}, line 84}\par\nopagebreak
\begin{leansrc}
theorem negF_denot_subset {S : Type*} (sat : S → Prop) (T : ℕ) :
    {ω : ℕ → S | ∀ n, ¬ sat (ω n)} ⊆ {ω : ℕ → S | ¬ ∃ k ≤ T, sat (ω k)}
\end{leansrc}
\par\medskip\noindent{\small \texttt{Qtsl.Structural.}\texttt{\bfseries hypMultiplier\_sound}}\hfill{\footnotesize \path{lean/qtsl-aux/QtslAux/Structural.lean}, line 75}\par\nopagebreak
\begin{leansrc}
theorem hypMultiplier_sound (cell childAtom : Set X) (N lam g : X → ℝ)
    (hbridge : childAtom = {x | 0 ≤ g x})
    (hlam : ∀ x ∈ cell, 0 ≤ lam x)
    (hcomb : ∀ x ∈ cell, N x + lam x * g x ≤ 0) :
    ∀ x ∈ cell ∩ childAtom, N x ≤ 0
\end{leansrc}
\par\medskip\noindent{\small \texttt{Qtsl.Structural.}\texttt{\bfseries fullDomainWeaken\_sound}}\hfill{\footnotesize \path{lean/qtsl-aux/QtslAux/Structural.lean}, line 88}\par\nopagebreak
\begin{leansrc}
theorem fullDomainWeaken_sound (Dm scope : Set X) (P : X → Prop)
    (hbridge : scope ⊆ Dm)
    (hrow : ∀ x ∈ Dm, P x) :
    ∀ x ∈ scope, P x
\end{leansrc}
\par\medskip\noindent{\small \texttt{Qtsl.Structural.}\texttt{\bfseries fullDomainWeaken\_offScope}}\hfill{\footnotesize \path{lean/qtsl-aux/QtslAux/Structural.lean}, line 94}\par\nopagebreak
\begin{leansrc}
theorem fullDomainWeaken_offScope (Dm Off : Set X) (P : X → Prop)
    (hrow : ∀ x ∈ Dm, P x) :
    ∀ x ∈ Dm \ Off, P x
\end{leansrc}
\par\medskip\noindent{\small \texttt{Qtsl.Structural.}\texttt{\bfseries regionSat\_mono}}\hfill{\footnotesize \path{lean/qtsl-aux/QtslAux/Structural.lean}, line 62}\par\nopagebreak
\begin{leansrc}
theorem regionSat_mono {R S den : Set X} (h : RegionSat R den) (hsub : S ⊆ R) :
    RegionSat S den
\end{leansrc}
\par\medskip\noindent{\small \texttt{Qtsl.Structural.}\texttt{\bfseries verifiedCover\_components\_sound}}\hfill{\footnotesize \path{lean/qtsl-aux/QtslAux/Structural.lean}, line 116}\par\nopagebreak
\begin{leansrc}
theorem verifiedCover_components_sound {ι : Type*} (Dm C Off scope : Set X)
    (K : ι → Set X) (P : X → Prop)
    (hbridge : scope = Dm \ Off)
    (hrowC : ∀ x ∈ C, P x)
    (hdecomp : Dm \ C ⊆ ⋃ i, K i)
    (hcomp : ∀ i, K i ⊆ Off) :
    ∀ x ∈ scope, P x
\end{leansrc}
\par\medskip\noindent{\small \texttt{Qtsl.Structural.}\texttt{\bfseries verifiedCover\_sound}}\hfill{\footnotesize \path{lean/qtsl-aux/QtslAux/Structural.lean}, line 102}\par\nopagebreak
\begin{leansrc}
theorem verifiedCover_sound (Dm C Off scope : Set X) (P : X → Prop)
    (hbridge : scope = Dm \ Off)
    (hrowC : ∀ x ∈ C, P x)
    (hresidual : Dm \ C ⊆ Off) :
    ∀ x ∈ scope, P x
\end{leansrc}
\par\medskip\noindent{\small \texttt{Qtsl.Structural.}\texttt{\bfseries RegionSat}}\hfill{\footnotesize \path{lean/qtsl-aux/QtslAux/Structural.lean}, line 57}\par\nopagebreak
\begin{leansrc}
def RegionSat (R den : Set X) : Prop
\end{leansrc}
\par\medskip\noindent{\small \texttt{Qtsl.Structural.}\texttt{\bfseries positiveCombinationAnd\_sound}}\hfill{\footnotesize \path{lean/qtsl-aux/QtslAux/Structural.lean}, line 129}\par\nopagebreak
\begin{leansrc}
theorem positiveCombinationAnd_sound (R₁ R₂ den₁ den₂ denAnd : Set X)
    (hbridge : denAnd = den₁ ∩ den₂)
    (h₁ : RegionSat R₁ den₁) (h₂ : RegionSat R₂ den₂) :
    RegionSat (R₁ ∩ R₂) denAnd
\end{leansrc}
\par\medskip\noindent{\small \texttt{Qtsl.Structural.}\texttt{\bfseries positiveCombinationOr\_sound}}\hfill{\footnotesize \path{lean/qtsl-aux/QtslAux/Structural.lean}, line 135}\par\nopagebreak
\begin{leansrc}
theorem positiveCombinationOr_sound (R₁ R₂ den₁ den₂ denOr : Set X)
    (hbridge : denOr = den₁ ∪ den₂)
    (h₁ : RegionSat R₁ den₁) (h₂ : RegionSat R₂ den₂) :
    RegionSat (R₁ ∪ R₂) denOr
\end{leansrc}
\par\medskip\noindent{\small \texttt{Qtsl.Structural.}\texttt{\bfseries positiveCombinationInter\_sound}}\hfill{\footnotesize \path{lean/qtsl-aux/QtslAux/Structural.lean}, line 141}\par\nopagebreak
\begin{leansrc}
theorem positiveCombinationInter_sound {ι : Type*} (R den : ι → Set X)
    (denAll : Set X)
    (hbridge : denAll = ⋂ i, den i)
    (h : ∀ i, RegionSat (R i) (den i)) :
    RegionSat (⋂ i, R i) denAll
\end{leansrc}
\par\medskip\noindent{\small \texttt{Qtsl.Structural.}\texttt{\bfseries positiveCombinationUnion\_sound}}\hfill{\footnotesize \path{lean/qtsl-aux/QtslAux/Structural.lean}, line 148}\par\nopagebreak
\begin{leansrc}
theorem positiveCombinationUnion_sound {ι : Type*} (R den : ι → Set X)
    (denAll : Set X)
    (hbridge : denAll = ⋃ i, den i)
    (h : ∀ i, RegionSat (R i) (den i)) :
    RegionSat (⋃ i, R i) denAll
\end{leansrc}
\par\medskip\noindent{\small \texttt{Qtsl.Structural.}\texttt{\bfseries compl\_over\_approx}}\hfill{\footnotesize \path{lean/qtsl-aux/QtslAux/Structural.lean}, line 201}\par\nopagebreak
\begin{leansrc}
theorem compl_over_approx {R den : Set X} (h : RegionSat R den) :
    denᶜ ⊆ Rᶜ
\end{leansrc}
\par\medskip\noindent{\small \texttt{Qtsl.Structural.}\texttt{\bfseries closureExactLt\_sound}}\hfill{\footnotesize \path{lean/qtsl-aux/QtslAux/Structural.lean}, line 157}\par\nopagebreak
\begin{leansrc}
theorem closureExactLt_sound (p : X → ℝ) (c : ℝ) (strictAtom : Set X)
    (hbridge : strictAtom = {x | p x < c}) :
    strictAtomᶜ = {x | c ≤ p x}
\end{leansrc}
\par\medskip\noindent{\small \texttt{Qtsl.Structural.}\texttt{\bfseries closureExact\_isClosed\_ge}}\hfill{\footnotesize \path{lean/qtsl-aux/QtslAux/Structural.lean}, line 171}\par\nopagebreak
\begin{leansrc}
theorem closureExact_isClosed_ge [TopologicalSpace X] {p : X → ℝ}
    (hp : Continuous p) (c : ℝ) : IsClosed {x | c ≤ p x}
\end{leansrc}
\par\medskip\noindent{\small \texttt{Qtsl.Structural.}\texttt{\bfseries closureExactGt\_sound}}\hfill{\footnotesize \path{lean/qtsl-aux/QtslAux/Structural.lean}, line 164}\par\nopagebreak
\begin{leansrc}
theorem closureExactGt_sound (p : X → ℝ) (c : ℝ) (strictAtom : Set X)
    (hbridge : strictAtom = {x | c < p x}) :
    strictAtomᶜ = {x | p x ≤ c}
\end{leansrc}
\par\medskip\noindent{\small \texttt{Qtsl.Structural.}\texttt{\bfseries closureExact\_isClosed\_le}}\hfill{\footnotesize \path{lean/qtsl-aux/QtslAux/Structural.lean}, line 175}\par\nopagebreak
\begin{leansrc}
theorem closureExact_isClosed_le [TopologicalSpace X] {p : X → ℝ}
    (hp : Continuous p) (c : ℝ) : IsClosed {x | p x ≤ c}
\end{leansrc}
\par\medskip\noindent{\small \texttt{Qtsl.Structural.}\texttt{\bfseries closureOutward\_sound}}\hfill{\footnotesize \path{lean/qtsl-aux/QtslAux/Structural.lean}, line 182}\par\nopagebreak
\begin{leansrc}
theorem closureOutward_sound (p : X → ℝ) (c : ℝ) (cmplAtom : Set X)
    (hbridge : cmplAtom = {x | c ≤ p x}) :
    {x | p x ≤ c}ᶜ ⊆ cmplAtom
\end{leansrc}
\par\medskip\noindent{\small \texttt{Qtsl.NNF.NForm.}\texttt{\bfseries elimNeg}}\hfill{\footnotesize \path{lean/qtsl-aux/QtslAux/NNFDuality.lean}, line 181}\par\nopagebreak
\begin{leansrc}
def NForm.elimNeg (d : α → α) : NForm α → NForm α
  | .pos a => .pos a
  | .negLit a => .pos (d a)
  | .and φ ψ => .and (φ.elimNeg d) (ψ.elimNeg d)
  | .or φ ψ => .or (φ.elimNeg d) (ψ.elimNeg d)

theorem negLits_elimNeg (d : α → α) : ∀ φ : NForm α, (φ.elimNeg d).negLits = ∅
\end{leansrc}
\par\medskip\noindent{\small \texttt{Qtsl.Structural.}\texttt{\bfseries closureOutwardGe\_sound}}\hfill{\footnotesize \path{lean/qtsl-aux/QtslAux/Structural.lean}, line 190}\par\nopagebreak
\begin{leansrc}
theorem closureOutwardGe_sound (p : X → ℝ) (c : ℝ) (cmplAtom : Set X)
    (hbridge : cmplAtom = {x | p x ≤ c}) :
    {x | c ≤ p x}ᶜ ⊆ cmplAtom
\end{leansrc}
\par\medskip\noindent{\small \texttt{Qtsl.Structural.}\texttt{\bfseries closureInwardEpsLt\_sound}}\hfill{\footnotesize \path{lean/qtsl-aux/QtslAux/Structural.lean}, line 209}\par\nopagebreak
\begin{leansrc}
theorem closureInwardEpsLt_sound (p : X → ℝ) (c ε : ℝ) (hε : 0 < ε)
    (roundedAtom : Set X)
    (hbridge : roundedAtom = {x | p x ≤ c - ε}) :
    roundedAtom ⊆ {x | p x < c}
\end{leansrc}
\par\medskip\noindent{\small \texttt{Qtsl.Structural.}\texttt{\bfseries closureInwardEps\_isClosed}}\hfill{\footnotesize \path{lean/qtsl-aux/QtslAux/Structural.lean}, line 229}\par\nopagebreak
\begin{leansrc}
theorem closureInwardEps_isClosed [TopologicalSpace X] {p : X → ℝ}
    (hp : Continuous p) (c ε : ℝ) : IsClosed {x | p x ≤ c - ε}
\end{leansrc}
\par\medskip\noindent{\small \texttt{Qtsl.Structural.}\texttt{\bfseries closureInwardEpsGt\_sound}}\hfill{\footnotesize \path{lean/qtsl-aux/QtslAux/Structural.lean}, line 219}\par\nopagebreak
\begin{leansrc}
theorem closureInwardEpsGt_sound (p : X → ℝ) (c ε : ℝ) (hε : 0 < ε)
    (roundedAtom : Set X)
    (hbridge : roundedAtom = {x | c + ε ≤ p x}) :
    roundedAtom ⊆ {x | c < p x}
\end{leansrc}
\par\medskip\noindent{\small \texttt{Qtsl.Aliases.}\texttt{\bfseries taylorLeaf\_le\_exp}}\hfill{\footnotesize \path{lean/qtsl-aux/QtslAux/QtslAliases.lean}, line 66}\par\nopagebreak
\begin{leansrc}
theorem taylorLeaf_le_exp {q x : ℝ} {n : ℕ} (hx : 0 ≤ x)
    (hq : q ≤ ∑ k ∈ range (n + 1), x ^ k / (Nat.factorial k : ℝ)) :
    q ≤ Real.exp x
\end{leansrc}
\par\medskip\noindent{\small \texttt{Qtsl.Aliases.}\texttt{\bfseries taylorLeaf\_rat\_le\_exp}}\hfill{\footnotesize \path{lean/qtsl-aux/QtslAux/QtslAliases.lean}, line 71}\par\nopagebreak
\begin{leansrc}
theorem taylorLeaf_rat_le_exp {qq θ : ℚ} {T n : ℕ} (hθ : 0 ≤ θ)
    (hleaf : qq ≤ ∑ k ∈ range (n + 1), (θ * T) ^ k / (Nat.factorial k : ℚ)) :
    (qq : ℝ) ≤ Real.exp (θ * T)
\end{leansrc}
\par\medskip\noindent{\small \texttt{PPLC.}\texttt{\bfseries taylorElab\_le\_exp}}\hfill{\footnotesize \path{lean/qtsl/Qtsl/Elaborators.lean}, line 224}\par\nopagebreak
\begin{leansrc}
theorem taylorElab_le_exp {q θ : Rat} {T n : Nat} (hθ : 0 ≤ θ)
    (hEval : (taylorFormula q θ T n).evalClosed = some true) :
    (q : ℝ) ≤ Real.exp (θ * T)
\end{leansrc}
\par\medskip\noindent{\small \texttt{PPLC.}\texttt{\bfseries taylorElab}}\hfill{\footnotesize \path{lean/qtsl/Qtsl/Elaborators.lean}, line 199}\par\nopagebreak
\begin{leansrc}
def taylorElab (H : Hypernet) (seq : Sequent) (q θ : Rat) (T n : Nat)
    (hEval : (taylorFormula q θ T n).evalClosed = some true)
    (hpost : seq.post = taylorFormula q θ T n) :
    PConstExactData H seq
\end{leansrc}
\par\medskip\noindent{\small \texttt{PPLC.}\texttt{\bfseries taylorPartialSum}}\hfill{\footnotesize \path{lean/qtsl/Qtsl/Elaborators.lean}, line 180}\par\nopagebreak
\begin{leansrc}
def taylorPartialSum (θ : Rat) (T n : Nat) : Rat
\end{leansrc}
\par\medskip\noindent{\small \texttt{PPLC.}\texttt{\bfseries taylorElab\_sound}}\hfill{\footnotesize \path{lean/qtsl/Qtsl/Elaborators.lean}, line 208}\par\nopagebreak
\begin{leansrc}
theorem taylorElab_sound {H : Hypernet} {seq : Sequent} {q θ : Rat} {T n : Nat}
    (hEval : (taylorFormula q θ T n).evalClosed = some true)
    (hpost : seq.post = taylorFormula q θ T n) :
    seq.semanticallyValid H
\end{leansrc}
\par\medskip\noindent{\small \texttt{PPLC.}\texttt{\bfseries taylorFormula\_true\_iff}}\hfill{\footnotesize \path{lean/qtsl/Qtsl/Elaborators.lean}, line 190}\par\nopagebreak
\begin{leansrc}
theorem taylorFormula_true_iff (q θ : Rat) (T n : Nat) :
    (taylorFormula q θ T n).evalClosed = some true
      ↔ q ≤ taylorPartialSum θ T n
\end{leansrc}
\par\medskip\noindent{\small \texttt{Qtsl.Aliases.}\texttt{\bfseries rootEval\_eval\_cast}}\hfill{\footnotesize \path{lean/qtsl-aux/QtslAux/QtslAliases.lean}, line 39}\par\nopagebreak
\begin{leansrc}
theorem rootEval_eval_cast {σ : Type*} (P : MvPolynomial σ ℚ) (xq : σ → ℚ) :
    ((MvPolynomial.eval xq P : ℚ) : ℝ)
      = MvPolynomial.eval (fun i => (xq i : ℝ)) (MvPolynomial.map (Rat.castHom ℝ) P)
\end{leansrc}
\par\medskip\noindent{\small \texttt{Qtsl.Aliases.}\texttt{\bfseries rootEval\_sublevel\_iff}}\hfill{\footnotesize \path{lean/qtsl-aux/QtslAux/QtslAliases.lean}, line 45}\par\nopagebreak
\begin{leansrc}
theorem rootEval_sublevel_iff {σ : Type*} (P : MvPolynomial σ ℚ) (xq : σ → ℚ)
    (thr : ℚ) :
    MvPolynomial.eval xq P ≤ thr
      ↔ MvPolynomial.eval (fun i => (xq i : ℝ))
          (MvPolynomial.map (Rat.castHom ℝ) P) ≤ (thr : ℝ)
\end{leansrc}
\par\medskip\noindent{\small \texttt{Qtsl.Aliases.}\texttt{\bfseries rootEval\_box\_iff}}\hfill{\footnotesize \path{lean/qtsl-aux/QtslAux/QtslAliases.lean}, line 53}\par\nopagebreak
\begin{leansrc}
theorem rootEval_box_iff {ι : Type*} (aq bq xq : ι → ℚ) :
    (∀ i, aq i ≤ xq i ∧ xq i ≤ bq i)
      ↔ ∀ i, ((aq i : ℝ) ≤ (xq i : ℝ) ∧ (xq i : ℝ) ≤ (bq i : ℝ))
\end{leansrc}
\par\medskip\noindent{\small \texttt{PPLC.}\texttt{\bfseries polyClosedTerm}}\hfill{\footnotesize \path{lean/qtsl/Qtsl/Elaborators.lean}, line 100}\par\nopagebreak
\begin{leansrc}
def polyClosedTerm (x0 : Nat → Rat) (p : Poly) : Term
\end{leansrc}
\par\medskip\noindent{\small \texttt{PPLC.}\texttt{\bfseries rootEvalElab}}\hfill{\footnotesize \path{lean/qtsl/Qtsl/Elaborators.lean}, line 147}\par\nopagebreak
\begin{leansrc}
def rootEvalElab (H : Hypernet) (seq : Sequent) (x0 : Nat → Rat) (V : Poly)
    (thr : Rat)
    (hEval : (rootEvalFormula x0 V thr).evalClosed = some true)
    (hpost : seq.post = rootEvalFormula x0 V thr) :
    PConstExactData H seq
\end{leansrc}
\par\medskip\noindent{\small \texttt{PPLC.}\texttt{\bfseries polyClosedTerm\_evalQ}}\hfill{\footnotesize \path{lean/qtsl/Qtsl/Elaborators.lean}, line 104}\par\nopagebreak
\begin{leansrc}
@[simp] theorem polyClosedTerm_evalQ (x0 : Nat → Rat) :
    ∀ p : Poly, (polyClosedTerm x0 p).evalQ? = some (Poly.eval x0 p)
  | [] => by simp [polyClosedTerm, Poly.eval, Term.evalQ?]
  | mc :: rest => by
    have ih
\end{leansrc}
\par\medskip\noindent{\small \texttt{PPLC.}\texttt{\bfseries rootEvalElab\_sound}}\hfill{\footnotesize \path{lean/qtsl/Qtsl/Elaborators.lean}, line 158}\par\nopagebreak
\begin{leansrc}
theorem rootEvalElab_sound {H : Hypernet} {seq : Sequent} {x0 : Nat → Rat}
    {V : Poly} {thr : Rat}
    (hEval : (rootEvalFormula x0 V thr).evalClosed = some true)
    (hpost : seq.post = rootEvalFormula x0 V thr) :
    seq.semanticallyValid H
\end{leansrc}
\par\medskip\noindent{\small \texttt{SuFunctor.Universe.}\texttt{\bfseries ratBox}}\hfill{\footnotesize \path{lean/su-functor/SuFunctor/Universe/Generation.lean}, line 67}\par\nopagebreak
\begin{leansrc}
def ratBox (l u : V → ℚ) : Set (ProbState V)
\end{leansrc}
\par\medskip\noindent{\small \texttt{Qtsl.GridCx.}\texttt{\bfseries Box}}\hfill{\footnotesize \path{lean/qtsl-aux/QtslAux/GridCx.lean}, line 49}\par\nopagebreak
\begin{leansrc}
structure Box (d : ℕ) where …
\end{leansrc}
\par\medskip\noindent{\small \texttt{Qtsl.Repar.AffineWitness.}\texttt{\bfseries t0\_soundness}}\hfill{\footnotesize \path{lean/qtsl-aux/QtslAux/Repar.lean}, line 196}\par\nopagebreak
\begin{leansrc}
theorem t0_soundness (N Nhat : MvPolynomial (Fin n) ℚ)
    (X D : Set (Fin n → ℝ))
    (hreplay : Nhat = φ.subst N)
    (hcorr : φ.invMapA '' X ⊆ D)
    (hcells : ∀ w ∈ D, aeval w Nhat ≤ 0) :
    ∀ z ∈ X, aeval z N ≤ 0
\end{leansrc}
\par\medskip\noindent{\small \texttt{Qtsl.Repar.AffineWitness.}\texttt{\bfseries t05\_soundness\_gridcx}}\hfill{\footnotesize \path{lean/qtsl-aux/QtslAux/Repar.lean}, line 247}\par\nopagebreak
\begin{leansrc}
theorem t05_soundness_gridcx {n : ℕ} {ι : Type*} (φs : ι → AffineWitness n)
    (N : MvPolynomial (Fin n) ℚ) (Nhat : ι → MvPolynomial (Fin n) ℚ)
    (B : ι → GridCx.Box n) (T : ∀ i, GridCx.PatchTree n (B i))
    (X : Set (Fin n → ℝ)) (Xp : ι → Set (Fin n → ℝ))
    (hX : X ⊆ ⋃ i, Xp i)
    (hreplay : ∀ i, Nhat i = (φs i).subst N)
    (hcorr : ∀ i, (φs i).invMapA '' Xp i ⊆ (B i).set)
    (hcells : ∀ i (c : (T i).CellIdx), ∀ w ∈ (T i).cell c,
      aeval w (Nhat i) ≤ 0) :
    ∀ z ∈ X, aeval z N ≤ 0
\end{leansrc}
\par\medskip\noindent{\small \texttt{Qtsl.Repar.AffineWitness.}\texttt{\bfseries left\_certified}}\hfill{\footnotesize \path{lean/qtsl-aux/QtslAux/Repar.lean}, line 77}\par\nopagebreak
\begin{leansrc}
theorem left_certified : φ.Ainv * φ.A = 1
\end{leansrc}
\par\medskip\noindent{\small \texttt{Qtsl.BoxBern.}\texttt{\bfseries boxBern\_enclosure}}\hfill{\footnotesize \path{lean/qtsl-aux/QtslAux/BoxBern.lean}, line 170}\par\nopagebreak
\begin{leansrc}
theorem boxBern_enclosure (lo hi : Fin n → ℚ) (N : MvPolynomial (Fin n) ℚ)
    (d : Fin n → ℕ) (b : (∀ v, Fin (d v + 1)) → ℚ) (m M : ℚ)
    (hexp : boxShift lo hi N = ∑ α, C (b α) * tensorBern d α)
    (hm : ∀ α, m ≤ b α) (hM : ∀ α, b α ≤ M) :
    ∀ x ∈ (ratBox lo hi).set, (m : ℝ) ≤ aeval x N ∧ aeval x N ≤ (M : ℝ)
\end{leansrc}
\par\medskip\noindent{\small \texttt{Qtsl.BoxBern.}\texttt{\bfseries boxShift}}\hfill{\footnotesize \path{lean/qtsl-aux/QtslAux/BoxBern.lean}, line 122}\par\nopagebreak
\begin{leansrc}
noncomputable def boxShift (lo hi : Fin n → ℚ) :
    MvPolynomial (Fin n) ℚ →ₐ[ℚ] MvPolynomial (Fin n) ℚ
\end{leansrc}
\par\medskip\noindent{\small \texttt{Qtsl.BoxBern.}\texttt{\bfseries aeval\_boxShift}}\hfill{\footnotesize \path{lean/qtsl-aux/QtslAux/BoxBern.lean}, line 126}\par\nopagebreak
\begin{leansrc}
theorem aeval_boxShift (lo hi : Fin n → ℚ) (u : Fin n → ℝ)
    (N : MvPolynomial (Fin n) ℚ) :
    aeval u (boxShift lo hi N)
      = aeval (fun v => (lo v : ℝ) + ((hi v : ℝ) - (lo v : ℝ)) * u v) N
\end{leansrc}
\par\medskip\noindent{\small \texttt{Qtsl.BoxBern.}\texttt{\bfseries boxBern\_nonpos}}\hfill{\footnotesize \path{lean/qtsl-aux/QtslAux/BoxBern.lean}, line 195}\par\nopagebreak
\begin{leansrc}
theorem boxBern_nonpos (lo hi : Fin n → ℚ) (N : MvPolynomial (Fin n) ℚ)
    (d : Fin n → ℕ) (b : (∀ v, Fin (d v + 1)) → ℚ)
    (hexp : boxShift lo hi N = ∑ α, C (b α) * tensorBern d α)
    (hcoeff : ∀ α, b α ≤ 0) :
    ∀ x ∈ (ratBox lo hi).set, aeval x N ≤ 0
\end{leansrc}
\par\medskip\noindent{\small \texttt{Qtsl.BoxBern.}\texttt{\bfseries boxBern\_enclosure\_minmax}}\hfill{\footnotesize \path{lean/qtsl-aux/QtslAux/BoxBern.lean}, line 184}\par\nopagebreak
\begin{leansrc}
theorem boxBern_enclosure_minmax (lo hi : Fin n → ℚ)
    (N : MvPolynomial (Fin n) ℚ)
    (d : Fin n → ℕ) (b : (∀ v, Fin (d v + 1)) → ℚ)
    (hexp : boxShift lo hi N = ∑ α, C (b α) * tensorBern d α) :
    ∀ x ∈ (ratBox lo hi).set,
      ((Finset.univ.inf' Finset.univ_nonempty b : ℚ) : ℝ) ≤ aeval x N ∧
        aeval x N ≤ ((Finset.univ.sup' Finset.univ_nonempty b : ℚ) : ℝ)
\end{leansrc}
\par\medskip\noindent{\small \texttt{Qtsl.BoxBern.}\texttt{\bfseries sum\_aeval\_tensorBern}}\hfill{\footnotesize \path{lean/qtsl-aux/QtslAux/BoxBern.lean}, line 79}\par\nopagebreak
\begin{leansrc}
theorem sum_aeval_tensorBern (d : Fin n → ℕ) (u : Fin n → ℝ) :
    ∑ α : ∀ v, Fin (d v + 1), aeval u (tensorBern d α) = 1
\end{leansrc}
\par\medskip\noindent{\small \texttt{Qtsl.Enclosure.}\texttt{\bfseries encCheck}}\hfill{\footnotesize \path{lean/qtsl-aux/QtslAux/Enclosure.lean}, line 93}\par\nopagebreak
\begin{leansrc}
def encCheck (w : EncWitness) : Bool
\end{leansrc}
\par\medskip\noindent{\small \texttt{Qtsl.Enclosure.}\texttt{\bfseries EncWitness}}\hfill{\footnotesize \path{lean/qtsl-aux/QtslAux/Enclosure.lean}, line 82}\par\nopagebreak
\begin{leansrc}
structure EncWitness where …
\end{leansrc}
\par\medskip\noindent{\small \texttt{Qtsl.Enclosure.}\texttt{\bfseries encCheck\_sound}}\hfill{\footnotesize \path{lean/qtsl-aux/QtslAux/Enclosure.lean}, line 380}\par\nopagebreak
\begin{leansrc}
theorem encCheck_sound (w : EncWitness) (h : encCheck w = true) :
    ∀ x : ℝ, |x| ≤ (w.r : ℝ) →
      |w.tag.eval x - evalPoly w.coeffs x| ≤ (w.R : ℝ)
\end{leansrc}
\par\medskip\noindent{\small \texttt{Qtsl.Enclosure.}\texttt{\bfseries encCheck\_sound\_bracket}}\hfill{\footnotesize \path{lean/qtsl-aux/QtslAux/Enclosure.lean}, line 426}\par\nopagebreak
\begin{leansrc}
theorem encCheck_sound_bracket (w : EncWitness) (h : encCheck w = true) :
    ∀ x : ℝ, |x| ≤ (w.r : ℝ) →
      evalPoly w.coeffs x - (w.R : ℝ) ≤ w.tag.eval x ∧
      w.tag.eval x ≤ evalPoly w.coeffs x + (w.R : ℝ)
\end{leansrc}
\par\medskip\noindent{\small \texttt{Qtsl.Enclosure.}\texttt{\bfseries encCheck\_point}}\hfill{\footnotesize \path{lean/qtsl-aux/QtslAux/Enclosure.lean}, line 437}\par\nopagebreak
\begin{leansrc}
theorem encCheck_point (w : EncWitness) (h : encCheck w = true) (c : ℚ)
    (hc : |c| ≤ w.r) :
    evalPoly w.coeffs (c : ℝ) - (w.R : ℝ) ≤ w.tag.eval (c : ℝ) ∧
    w.tag.eval (c : ℝ) ≤ evalPoly w.coeffs (c : ℝ) + (w.R : ℝ)
\end{leansrc}
\par\medskip\noindent{\small \texttt{Qtsl.Enclosure.}\texttt{\bfseries linear\_interp\_le}}\hfill{\footnotesize \path{lean/qtsl-aux/QtslAux/Enclosure.lean}, line 445}\par\nopagebreak
\begin{leansrc}
theorem linear_interp_le {A C l h w : ℝ} (hlw : l ≤ w) (hwh : w ≤ h)
    (hl : A + C * l ≤ 0) (hh : A + C * h ≤ 0) : A + C * w ≤ 0
\end{leansrc}
\par\medskip\noindent{\small \texttt{Qtsl.Enclosure.}\texttt{\bfseries integral\_exp\_neg\_sq}}\hfill{\footnotesize \path{lean/qtsl-aux/QtslAux/Enclosure.lean}, line 459}\par\nopagebreak
\begin{leansrc}
theorem integral_exp_neg_sq (z : ℝ) :
    ∫ t in (0:ℝ)..z, Real.exp (-(t ^ 2)) = z * FnTag.erfT.eval (z ^ 2)
\end{leansrc}
\par\medskip\noindent{\small \texttt{PPLC.Qtsl.Behavior.}\texttt{\bfseries qMorseDischarge\_sound}}\hfill{\footnotesize \path{lean/qtsl/Qtsl/Behavioral.lean}, line 369}\par\nopagebreak
\begin{leansrc}
theorem qMorseDischarge_sound (B : Behavior X) (D : Set X)
    (ρ : Vtx → Set X) (cells : ι → Set X) (hmap : ι → Option Vtx)
    (hContain : ∀ i v, hmap i = some v → cells i ⊆ ρ v)
    (hCompact : ∀ v, IsCompact (B.MaxInv (fiberUnion cells hmap v)))
    (hDisj : ∀ v w, v ≠ w →
      Disjoint (B.MaxInv (fiberUnion cells hmap v))
        (B.MaxInv (fiberUnion cells hmap w)))
    (hCell : ∀ x ∈ D, B.ChainRecurrentIn D x → ∃ i, x ∈ B.MaxInv (cells i))
    (hTransit : ∀ i, hmap i = none → B.HasTransitCert (cells i))
    (hOrder : ∀ γ ∈ B.cpl, Set.range γ ⊆ D →
      (∀ v, ¬ Set.range γ ⊆ B.MaxInv (fiberUnion cells hmap v)) →
      ∃ v w, w < v ∧ alphaSet γ ⊆ B.MaxInv (fiberUnion cells hmap v) ∧
        omegaSet γ ⊆ B.MaxInv (fiberUnion cells hmap w)) :
    B.MorseSat D ρ
\end{leansrc}
\par\medskip\noindent{\small \texttt{PPLC.Qtsl.Behavior.}\texttt{\bfseries MorseWitness}}\hfill{\footnotesize \path{lean/qtsl/Qtsl/Behavioral.lean}, line 99}\par\nopagebreak
\begin{leansrc}
structure MorseWitness (B : Behavior X) (D : Set X) [PartialOrder Vtx]
    (ρ : Vtx → Set X) where …
\end{leansrc}
\par\medskip\noindent{\small \texttt{PPLC.Qtsl.Behavior.}\texttt{\bfseries HasTransitCert}}\hfill{\footnotesize \path{lean/qtsl/Qtsl/Behavioral.lean}, line 124}\par\nopagebreak
\begin{leansrc}
def HasTransitCert (B : Behavior X) (C : Set X) : Prop
\end{leansrc}
\par\medskip\noindent{\small \texttt{PPLC.Qtsl.Semiflow.}\texttt{\bfseries FlowDerivOn}}\hfill{\footnotesize \path{lean/qtsl/Qtsl/MeanValue.lean}, line 24}\par\nopagebreak
\begin{leansrc}
def FlowDerivOn (φ : Semiflow X) (V W : X → ℝ) (C : Set X) : Prop
\end{leansrc}
\par\medskip\noindent{\small \texttt{PPLC.Qtsl.Semiflow.}\texttt{\bfseries maxInv\_isCompact}}\hfill{\footnotesize \path{lean/qtsl/Qtsl/Discharge.lean}, line 231}\par\nopagebreak
\begin{leansrc}
theorem maxInv_isCompact [T2Space X] {φ : Semiflow X} {C : Set X}
    (hflow : ∀ t, Continuous (φ.flow t)) (hC : IsCompact C) :
    IsCompact (φ.MaxInv C)
\end{leansrc}
\par\medskip\noindent{\small \texttt{PPLC.Qtsl.Semiflow.}\texttt{\bfseries fiberUnion\_isCompact}}\hfill{\footnotesize \path{lean/qtsl/Qtsl/Discharge.lean}, line 261}\par\nopagebreak
\begin{leansrc}
theorem fiberUnion_isCompact {Vtx ι : Type*} [Finite ι]
    {cells : ι → Set X} {hmap : ι → Option Vtx}
    (hcells : ∀ i, IsCompact (cells i)) (v : Vtx) :
    IsCompact (fiberUnion cells hmap v)
\end{leansrc}
\par\medskip\noindent{\small \texttt{PPLC.Qtsl.Semiflow.HasTransitCert.}\texttt{\bfseries disjoint\_maxInv}}\hfill{\footnotesize \path{lean/qtsl/Qtsl/Discharge.lean}, line 81}\par\nopagebreak
\begin{leansrc}
theorem HasTransitCert.disjoint_maxInv {φ : Semiflow X} {C₁ C₂ : Set X}
    (h : φ.HasTransitCert (C₁ ∩ C₂)) :
    Disjoint (φ.MaxInv C₁) (φ.MaxInv C₂)
\end{leansrc}
\par\medskip\noindent{\small \texttt{PPLC.Qtsl.Behavior.}\texttt{\bfseries crCover\_assemble}}\hfill{\footnotesize \path{lean/qtsl/Qtsl/Behavioral.lean}, line 352}\par\nopagebreak
\begin{leansrc}
theorem crCover_assemble {B : Behavior X} {D : Set X}
    {cells : ι → Set X} {hmap : ι → Option Vtx}
    (hCell : ∀ x ∈ D, B.ChainRecurrentIn D x → ∃ i, x ∈ B.MaxInv (cells i))
    (hTransit : ∀ i, hmap i = none → B.HasTransitCert (cells i)) :
    ∀ x ∈ D, B.ChainRecurrentIn D x →
      ∃ v, x ∈ B.MaxInv (fiberUnion cells hmap v)
\end{leansrc}
\par\medskip\noindent{\small \texttt{PPLC.Qtsl.Behavior.HasTransitCert.}\texttt{\bfseries maxInv\_empty}}\hfill{\footnotesize \path{lean/qtsl/Qtsl/Behavioral.lean}, line 144}\par\nopagebreak
\begin{leansrc}
theorem HasTransitCert.maxInv_empty {B : Behavior X} {C : Set X}
    (h : B.HasTransitCert C) : B.MaxInv C = ∅
\end{leansrc}
\par\medskip\noindent{\small \texttt{PPLC.Qtsl.Semiflow.}\texttt{\bfseries qMorseDischarge\_sound\_orderAssembled}}\hfill{\footnotesize \path{lean/qtsl/Qtsl/OrderAssembly.lean}, line 250}\par\nopagebreak
\begin{leansrc}
theorem qMorseDischarge_sound_orderAssembled (φ : Semiflow X) (D : Set X)
    (ρ : Vtx → Set X) (cells : ι → Set X) (hmap : ι → Option Vtx)
    (hD : IsCompact D)
    (hflow : ∀ t, Continuous (φ.flow t))
    (hcells : ∀ i, IsCompact (cells i))
    (hContain : ∀ i v, hmap i = some v → cells i ⊆ ρ v)
    (hOverlap : ∀ v w, v ≠ w →
      φ.HasTransitCert
        (fiberUnion cells hmap v ∩ fiberUnion cells hmap w))
    (hCell : ∀ x ∈ D, ChainRecurrentIn φ D x → ∃ i, x ∈ φ.MaxInv (cells i))
    (hTransitCert : ∀ i, hmap i = none → φ.HasTransitCert (cells i))
    (hConv : ∀ i j k v, i ≤ j → j ≤ k →
      hmap i = some v → hmap k = some v → hmap j = some v)
    (hMono : ∀ i j v w, i ≤ j →
      hmap i = some v → hmap j = some w → v ≤ w)
    (hInterval : ∀ γ : ℝ → X, φ.CompleteOrbit γ → Set.range γ ⊆ D →
      ∃ bα bω : ι, bω ≤ bα ∧
        alphaSet γ ⊆ cells bα ∧ omegaSet γ ⊆ cells bω ∧
        ∀ t : ℝ, ∃ i, bω ≤ i ∧ i ≤ bα ∧ γ t ∈ cells i) :
    MorseSat φ D ρ
\end{leansrc}
\par\medskip\noindent{\small \texttt{PPLC.Qtsl.Semiflow.}\texttt{\bfseries qMorseDischarge\_sound\_eventualMember}}\hfill{\footnotesize \path{lean/qtsl/Qtsl/EventualMember.lean}, line 274}\par\nopagebreak
\begin{leansrc}
theorem qMorseDischarge_sound_eventualMember (φ : Semiflow X) (D : Set X)
    (ρ : Vtx → Set X) (cells : ι → Set X) (hmap : ι → Option Vtx)
    (blocks : β → Set X) (emb : ι → β) (cover : ι → β) (bD : β)
    (hD : IsCompact D)
    (hflow : ∀ t, Continuous (φ.flow t))
    (hcells : ∀ i, IsCompact (cells i))
    (hContain : ∀ i v, hmap i = some v → cells i ⊆ ρ v)
    (hOverlap : ∀ v w, v ≠ w →
      φ.HasTransitCert
        (fiberUnion cells hmap v ∩ fiberUnion cells hmap w))
    (hCell : ∀ x ∈ D, ChainRecurrentIn φ D x → ∃ i, x ∈ φ.MaxInv (cells i))
    (hTransitCert : ∀ i, hmap i = none → φ.HasTransitCert (cells i))
    (hConv : ∀ i j k v, i ≤ j → j ≤ k →
      hmap i = some v → hmap k = some v → hmap j = some v)
    (hMono : ∀ i j v w, i ≤ j →
      hmap i = some v → hmap j = some w → v ≤ w)
    (hbD : D ⊆ blocks bD)
    (hJoin : ∀ b : β, ∀ x ∈ blocks b,
      ∃ i : ι, x ∈ blocks (emb i) ∧ blocks (emb i) ⊆ blocks b)
    (hcover : ∀ i : ι, blocks (cover i) ⊂ blocks (emb i))
    (hcellsSup : ∀ i : ι, blocks (emb i) \ blocks (cover i) ⊆ cells i)
    (horder : ∀ i j : ι, blocks (emb i) ⊆ blocks (emb j) → i ≤ j)
    (hclosed : ∀ i : ι, IsClosed (blocks (emb i)))
    (hInv : ∀ b : β, φ.Invariant (blocks b))
    (hMeet : ∀ a b : β, ∃ c : β, blocks c ⊆ blocks a ∩ blocks b ∧
      ∀ γ : ℝ → X, φ.CompleteOrbit γ → Set.range γ ⊆ D →
        ∀ s t : ℝ, (∀ u, s ≤ u → u ≤ t → γ u ∈ blocks a ∩ blocks b) →
          ∀ u, s ≤ u → u ≤ t → γ u ∈ blocks c)
    (hAbsorbOmega : ∀ γ : ℝ → X, φ.CompleteOrbit γ → Set.range γ ⊆ D →
      ∀ b : β, (omegaSet γ ∩ blocks b).Nonempty →
        ∃ T : ℝ, ∀ t, T ≤ t → γ t ∈ blocks b)
    (hAbsorbAlpha : ∀ γ : ℝ → X, φ.CompleteOrbit γ → Set.range γ ⊆ D →
      ∀ b : β, (alphaSet γ ∩ blocks b).Nonempty →
        ∃ T : ℝ, ∀ t, t ≤ T → γ t ∈ blocks b) :
    MorseSat φ D ρ
\end{leansrc}
\par\medskip\noindent{\small \texttt{PPLC.Qtsl.Semiflow.}\texttt{\bfseries qMorseDischarge\_sound\_cellCover}}\hfill{\footnotesize \path{lean/qtsl/Qtsl/Crossing.lean}, line 120}\par\nopagebreak
\begin{leansrc}
theorem qMorseDischarge_sound_cellCover (φ : Semiflow X) (D : Set X)
    (ρ : Vtx → Set X) (cells : ι → Set X) (hmap : ι → Option Vtx)
    (hflow : ∀ t, Continuous (φ.flow t))
    (htraj : ∀ x, Continuous fun t : ℝ≥0 => φ.flow t x)
    (hcells : ∀ i, IsCompact (cells i))
    (hContain : ∀ i v, hmap i = some v → cells i ⊆ ρ v)
    (hCross : ∀ v w, v ≠ w → ∀ i j, hmap i = some v → hmap j = some w →
      φ.HasCrossingCert (cells i ∩ cells j))
    (hCell : ∀ x ∈ D, ChainRecurrentIn φ D x → ∃ i, x ∈ φ.MaxInv (cells i))
    (hTransitCert : ∀ i, hmap i = none → φ.HasTransitCert (cells i))
    (hOrder : ∀ γ : ℝ → X, φ.CompleteOrbit γ → Set.range γ ⊆ D →
      (∀ v, ¬ Set.range γ ⊆ φ.MaxInv (fiberUnion cells hmap v)) →
      ∃ v w, w < v ∧ alphaSet γ ⊆ φ.MaxInv (fiberUnion cells hmap v) ∧
        omegaSet γ ⊆ φ.MaxInv (fiberUnion cells hmap w)) :
    MorseSat φ D ρ
\end{leansrc}
\par\medskip\noindent{\small \texttt{PPLC.Qtsl.Behavior.}\texttt{\bfseries maxInv\_subset}}\hfill{\footnotesize \path{lean/qtsl/Qtsl/Behavioral.lean}, line 51}\par\nopagebreak
\begin{leansrc}
theorem maxInv_subset (B : Behavior X) (C : Set X) : B.MaxInv C ⊆ C
\end{leansrc}
\par\medskip\noindent{\small \texttt{PPLC.Qtsl.Behavior.}\texttt{\bfseries maxInv\_isInvariant}}\hfill{\footnotesize \path{lean/qtsl/Qtsl/Behavioral.lean}, line 69}\par\nopagebreak
\begin{leansrc}
theorem maxInv_isInvariant (B : Behavior X) (C : Set X) :
    B.IsInvariantIn (B.MaxInv C)
\end{leansrc}
\par\medskip\noindent{\small \texttt{PPLC.Qtsl.Semiflow.}\texttt{\bfseries maxInv\_invariant}}\hfill{\footnotesize \path{lean/qtsl/Qtsl/QtslSoundness.lean}, line 63}\par\nopagebreak
\begin{leansrc}
theorem maxInv_invariant (φ : Semiflow X) (C : Set X) :
    φ.Invariant (φ.MaxInv C)
\end{leansrc}
\par\medskip\noindent{\small \texttt{PPLC.Qtsl.Semiflow.}\texttt{\bfseries fiberUnion\_subset}}\hfill{\footnotesize \path{lean/qtsl/Qtsl/QtslSoundness.lean}, line 195}\par\nopagebreak
\begin{leansrc}
theorem fiberUnion_subset {cells : ι → Set X} {hmap : ι → Option Vtx}
    {ρ : Vtx → Set X} (hContain : ∀ i v, hmap i = some v → cells i ⊆ ρ v)
    (v : Vtx) : fiberUnion cells hmap v ⊆ ρ v
\end{leansrc}
\par\medskip\noindent{\small \texttt{PPLC.Qtsl.Semiflow.}\texttt{\bfseries maxInv\_forward\_confined}}\hfill{\footnotesize \path{lean/qtsl/Qtsl/Discharge.lean}, line 38}\par\nopagebreak
\begin{leansrc}
theorem maxInv_forward_confined {φ : Semiflow X} {C : Set X} {x : X}
    (hx : x ∈ φ.MaxInv C) : ∀ t : ℝ≥0, φ.flow t x ∈ C
\end{leansrc}
\par\medskip\noindent{\small \texttt{PPLC.Qtsl.Behavior.HasTransitCert.}\texttt{\bfseries not\_forward\_confined}}\hfill{\footnotesize \path{lean/qtsl/Qtsl/Behavioral.lean}, line 129}\par\nopagebreak
\begin{leansrc}
theorem HasTransitCert.not_forward_confined {B : Behavior X} {C : Set X}
    (h : B.HasTransitCert C) {γ : ℝ≥0 → X} (hγ : γ ∈ B.fwd) :
    ¬ ∀ t : ℝ≥0, γ t ∈ C
\end{leansrc}
\par\medskip\noindent{\small \texttt{PPLC.Qtsl.Semiflow.}\texttt{\bfseries transit\_maxInv\_empty}}\hfill{\footnotesize \path{lean/qtsl/Qtsl/QtslSoundness.lean}, line 114}\par\nopagebreak
\begin{leansrc}
theorem transit_maxInv_empty {φ : Semiflow X} {C : Set X} {V : X → ℝ}
    {ε : ℝ} (hε : 0 < ε) (hBdd : BddBelow (V '' C))
    (hDec : ∀ x ∈ C, ∀ t : ℝ≥0,
      (∀ s : ℝ≥0, s ≤ t → φ.flow s x ∈ C) → V (φ.flow t x) ≤ V x - ε * t) :
    φ.MaxInv C = ∅
\end{leansrc}
\par\medskip\noindent{\small \texttt{PPLC.Qtsl.Behavior.}\texttt{\bfseries maxInv\_mono}}\hfill{\footnotesize \path{lean/qtsl/Qtsl/Behavioral.lean}, line 55}\par\nopagebreak
\begin{leansrc}
theorem maxInv_mono {B : Behavior X} {C D : Set X} (h : C ⊆ D) :
    B.MaxInv C ⊆ B.MaxInv D
\end{leansrc}
\par\medskip\noindent{\small \texttt{PPLC.Qtsl.Semiflow.}\texttt{\bfseries subset\_fiberUnion}}\hfill{\footnotesize \path{lean/qtsl/Qtsl/Discharge.lean}, line 45}\par\nopagebreak
\begin{leansrc}
theorem subset_fiberUnion {Vtx ι : Type*} {cells : ι → Set X}
    {hmap : ι → Option Vtx} {i : ι} {v : Vtx} (hi : hmap i = some v) :
    cells i ⊆ fiberUnion cells hmap v
\end{leansrc}
\par\medskip\noindent{\small \texttt{PPLC.Qtsl.Semiflow.}\texttt{\bfseries qMorseDischarge\_sound\_modelClass}}\hfill{\footnotesize \path{lean/qtsl/Qtsl/Discharge.lean}, line 297}\par\nopagebreak
\begin{leansrc}
theorem qMorseDischarge_sound_modelClass (φ : Semiflow X) (D : Set X)
    (ρ : Vtx → Set X) (cells : ι → Set X) (hmap : ι → Option Vtx)
    (hflow : ∀ t, Continuous (φ.flow t))
    (hcells : ∀ i, IsCompact (cells i))
    (hContain : ∀ i v, hmap i = some v → cells i ⊆ ρ v)
    (hOverlap : ∀ v w, v ≠ w →
      φ.HasTransitCert
        (fiberUnion cells hmap v ∩ fiberUnion cells hmap w))
    (hCell : ∀ x ∈ D, ChainRecurrentIn φ D x → ∃ i, x ∈ φ.MaxInv (cells i))
    (hTransitCert : ∀ i, hmap i = none → φ.HasTransitCert (cells i))
    (hOrder : ∀ γ : ℝ → X, φ.CompleteOrbit γ → Set.range γ ⊆ D →
      (∀ v, ¬ Set.range γ ⊆ φ.MaxInv (fiberUnion cells hmap v)) →
      ∃ v w, w < v ∧ alphaSet γ ⊆ φ.MaxInv (fiberUnion cells hmap v) ∧
        omegaSet γ ⊆ φ.MaxInv (fiberUnion cells hmap w)) :
    MorseSat φ D ρ
\end{leansrc}
\par\medskip\noindent{\small \texttt{PPLC.Qtsl.Behavior.}\texttt{\bfseries MorseSat}}\hfill{\footnotesize \path{lean/qtsl/Qtsl/Behavioral.lean}, line 118}\par\nopagebreak
\begin{leansrc}
def MorseSat (B : Behavior X) (D : Set X) [PartialOrder Vtx]
    (ρ : Vtx → Set X) : Prop
\end{leansrc}
\par\medskip\noindent{\small \texttt{PPLC.Qtsl.Behavior.}\texttt{\bfseries EpsChainIn}}\hfill{\footnotesize \path{lean/qtsl/Qtsl/Behavioral.lean}, line 80}\par\nopagebreak
\begin{leansrc}
inductive EpsChainIn (B : Behavior X) (D : Set X) (ε : ℝ) (T : ℝ≥0) :
    X → X → Prop
  | step {x y : X} {γ : ℝ≥0 → X} {t : ℝ≥0} (hγ : γ ∈ B.fwd) (hγ0 : γ 0 = x)
      (hx : x ∈ D) (hy : y ∈ D) (ht : T ≤ t) (hjump : dist (γ t) y < ε) :
      EpsChainIn B D ε T x y
  | trans {x y z : X} (h1 : EpsChainIn B D ε T x y)
      (h2 : EpsChainIn B D ε T y z) : EpsChainIn B D ε T x z

def ChainRecurrentIn (B : Behavior X) (D : Set X) (x : X) : Prop
\end{leansrc}
\par\medskip\noindent{\small \texttt{PPLC.Qtsl.Behavior.}\texttt{\bfseries ChainRecurrentIn}}\hfill{\footnotesize \path{lean/qtsl/Qtsl/Behavioral.lean}, line 88}\par\nopagebreak
\begin{leansrc}
def ChainRecurrentIn (B : Behavior X) (D : Set X) (x : X) : Prop
\end{leansrc}
\par\medskip\noindent{\small \texttt{PPLC.Qtsl.Semiflow.}\texttt{\bfseries omegaSet}}\hfill{\footnotesize \path{lean/qtsl/Qtsl/QtslSoundness.lean}, line 142}\par\nopagebreak
\begin{leansrc}
def omegaSet [TopologicalSpace X] (γ : ℝ → X) : Set X
\end{leansrc}
\par\medskip\noindent{\small \texttt{PPLC.Qtsl.Semiflow.}\texttt{\bfseries alphaSet}}\hfill{\footnotesize \path{lean/qtsl/Qtsl/QtslSoundness.lean}, line 145}\par\nopagebreak
\begin{leansrc}
def alphaSet [TopologicalSpace X] (γ : ℝ → X) : Set X
\end{leansrc}
\par\medskip\noindent{\small \texttt{PPLC.Qtsl.Semiflow.}\texttt{\bfseries orderClause\_assemble}}\hfill{\footnotesize \path{lean/qtsl/Qtsl/OrderAssembly.lean}, line 183}\par\nopagebreak
\begin{leansrc}
theorem orderClause_assemble (φ : Semiflow X) (D : Set X)
    (hD : IsCompact D) (hflow : ∀ t, Continuous (φ.flow t))
    (cells : ι → Set X) (hmap : ι → Option Vtx)
    (hConv : ∀ i j k v, i ≤ j → j ≤ k →
      hmap i = some v → hmap k = some v → hmap j = some v)
    (hMono : ∀ i j v w, i ≤ j →
      hmap i = some v → hmap j = some w → v ≤ w)
    (hTransit : ∀ i, hmap i = none → φ.HasTransitCert (cells i))
    (hInterval : ∀ γ : ℝ → X, φ.CompleteOrbit γ → Set.range γ ⊆ D →
      ∃ bα bω : ι, bω ≤ bα ∧
        alphaSet γ ⊆ cells bα ∧ omegaSet γ ⊆ cells bω ∧
        ∀ t : ℝ, ∃ i, bω ≤ i ∧ i ≤ bα ∧ γ t ∈ cells i) :
    ∀ γ : ℝ → X, φ.CompleteOrbit γ → Set.range γ ⊆ D →
      (∀ v, ¬ Set.range γ ⊆ φ.MaxInv (fiberUnion cells hmap v)) →
      ∃ v w, w < v ∧
        alphaSet γ ⊆ φ.MaxInv (fiberUnion cells hmap v) ∧
        omegaSet γ ⊆ φ.MaxInv (fiberUnion cells hmap w)
\end{leansrc}
\par\medskip\noindent{\small \texttt{PPLC.Qtsl.Semiflow.}\texttt{\bfseries eventualMember\_interval}}\hfill{\footnotesize \path{lean/qtsl/Qtsl/EventualMember.lean}, line 145}\par\nopagebreak
\begin{leansrc}
theorem eventualMember_interval (φ : Semiflow X) (D : Set X)
    (blocks : β → Set X) (emb : ι → β) (cover : ι → β) (bD : β)
    (cells : ι → Set X)
    (hbD : D ⊆ blocks bD)
    (hJoin : ∀ b : β, ∀ x ∈ blocks b,
      ∃ i : ι, x ∈ blocks (emb i) ∧ blocks (emb i) ⊆ blocks b)
    (hcover : ∀ i : ι, blocks (cover i) ⊂ blocks (emb i))
    (hcellsSup : ∀ i : ι, blocks (emb i) \ blocks (cover i) ⊆ cells i)
    (horder : ∀ i j : ι, blocks (emb i) ⊆ blocks (emb j) → i ≤ j)
    (hclosed : ∀ i : ι, IsClosed (blocks (emb i)))
    (hInv : ∀ b : β, φ.Invariant (blocks b))
    (hMeet : ∀ a b : β, ∃ c : β, blocks c ⊆ blocks a ∩ blocks b ∧
      ∀ γ : ℝ → X, φ.CompleteOrbit γ → Set.range γ ⊆ D →
        ∀ s t : ℝ, (∀ u, s ≤ u → u ≤ t → γ u ∈ blocks a ∩ blocks b) →
          ∀ u, s ≤ u → u ≤ t → γ u ∈ blocks c)
    (hAbsorbOmega : ∀ γ : ℝ → X, φ.CompleteOrbit γ → Set.range γ ⊆ D →
      ∀ b : β, (omegaSet γ ∩ blocks b).Nonempty →
        ∃ T : ℝ, ∀ t, T ≤ t → γ t ∈ blocks b)
    (hAbsorbAlpha : ∀ γ : ℝ → X, φ.CompleteOrbit γ → Set.range γ ⊆ D →
      ∀ b : β, (alphaSet γ ∩ blocks b).Nonempty →
        ∃ T : ℝ, ∀ t, t ≤ T → γ t ∈ blocks b) :
    ∀ γ : ℝ → X, φ.CompleteOrbit γ → Set.range γ ⊆ D →
      ∃ bα bω : ι, bω ≤ bα ∧
        alphaSet γ ⊆ cells bα ∧ omegaSet γ ⊆ cells bω ∧
        ∀ t : ℝ, ∃ i, bω ≤ i ∧ i ≤ bα ∧ γ t ∈ cells i
\end{leansrc}
\par\medskip\noindent{\small \texttt{PPLC.Qtsl.Semiflow.}\texttt{\bfseries prod}}\hfill{\footnotesize \path{lean/qtsl/Qtsl/Product.lean}, line 21}\par\nopagebreak
\begin{leansrc}
def prod (φ₁ : Semiflow X₁) (φ₂ : Semiflow X₂) : Semiflow (X₁ × X₂) where …
\end{leansrc}
\par\medskip\noindent{\small \texttt{PPLC.Qtsl.Semiflow.MorseWitness.}\texttt{\bfseries prod}}\hfill{\footnotesize \path{lean/qtsl/Qtsl/Product.lean}, line 154}\par\nopagebreak
\begin{leansrc}
def MorseWitness.prod (W₁ : MorseWitness φ₁ D₁ ρ₁)
    (W₂ : MorseWitness φ₂ D₂ ρ₂) :
    MorseWitness (φ₁.prod φ₂) (D₁ ×ˢ D₂)
      (fun p : Vtx₁ × Vtx₂ => ρ₁ p.1 ×ˢ ρ₂ p.2) where …
\end{leansrc}
\par\medskip\noindent{\small \texttt{PPLC.Qtsl.Semiflow.}\texttt{\bfseries qMorseProduct\_sound}}\hfill{\footnotesize \path{lean/qtsl/Qtsl/Product.lean}, line 235}\par\nopagebreak
\begin{leansrc}
theorem qMorseProduct_sound
    (h₁ : MorseSat φ₁ D₁ ρ₁) (h₂ : MorseSat φ₂ D₂ ρ₂) :
    MorseSat (φ₁.prod φ₂) (D₁ ×ˢ D₂)
      (fun p : Vtx₁ × Vtx₂ => ρ₁ p.1 ×ˢ ρ₂ p.2)
\end{leansrc}
\par\medskip\noindent{\small \texttt{PPLC.Qtsl.Semiflow.}\texttt{\bfseries prod\_flow}}\hfill{\footnotesize \path{lean/qtsl/Qtsl/Product.lean}, line 26}\par\nopagebreak
\begin{leansrc}
@[simp] theorem prod_flow (φ₁ : Semiflow X₁) (φ₂ : Semiflow X₂)
    (t : ℝ≥0) (p : X₁ × X₂) :
    (φ₁.prod φ₂).flow t p = (φ₁.flow t p.1, φ₂.flow t p.2)
\end{leansrc}
\par\medskip\noindent{\small \texttt{PPLC.Qtsl.Semiflow.}\texttt{\bfseries completeOrbit\_prod\_iff}}\hfill{\footnotesize \path{lean/qtsl/Qtsl/Product.lean}, line 39}\par\nopagebreak
\begin{leansrc}
theorem completeOrbit_prod_iff {γ : ℝ → X₁ × X₂} :
    (φ₁.prod φ₂).CompleteOrbit γ ↔
      φ₁.CompleteOrbit (fun s => (γ s).1) ∧
      φ₂.CompleteOrbit (fun s => (γ s).2)
\end{leansrc}
\par\medskip\noindent{\small \texttt{PPLC.Qtsl.Semiflow.CompleteOrbit.}\texttt{\bfseries prodMk}}\hfill{\footnotesize \path{lean/qtsl/Qtsl/Product.lean}, line 32}\par\nopagebreak
\begin{leansrc}
theorem CompleteOrbit.prodMk {γ₁ : ℝ → X₁} {γ₂ : ℝ → X₂}
    (h₁ : φ₁.CompleteOrbit γ₁) (h₂ : φ₂.CompleteOrbit γ₂) :
    (φ₁.prod φ₂).CompleteOrbit fun s => (γ₁ s, γ₂ s)
\end{leansrc}
\par\medskip\noindent{\small \texttt{PPLC.Qtsl.Semiflow.}\texttt{\bfseries omegaSet\_prod\_subset}}\hfill{\footnotesize \path{lean/qtsl/Qtsl/Product.lean}, line 55}\par\nopagebreak
\begin{leansrc}
theorem omegaSet_prod_subset (γ : ℝ → X₁ × X₂) :
    omegaSet γ ⊆
      (omegaSet fun s => (γ s).1) ×ˢ (omegaSet fun s => (γ s).2)
\end{leansrc}
\par\medskip\noindent{\small \texttt{PPLC.Qtsl.Semiflow.}\texttt{\bfseries alphaSet\_prod\_subset}}\hfill{\footnotesize \path{lean/qtsl/Qtsl/Product.lean}, line 72}\par\nopagebreak
\begin{leansrc}
theorem alphaSet_prod_subset (γ : ℝ → X₁ × X₂) :
    alphaSet γ ⊆
      (alphaSet fun s => (γ s).1) ×ˢ (alphaSet fun s => (γ s).2)
\end{leansrc}
\par\medskip\noindent{\small \texttt{PPLC.Qtsl.Semiflow.EpsChainIn.}\texttt{\bfseries fst}}\hfill{\footnotesize \path{lean/qtsl/Qtsl/Product.lean}, line 111}\par\nopagebreak
\begin{leansrc}
theorem EpsChainIn.fst {D₁ : Set X₁} {D₂ : Set X₂} {ε : ℝ} {T : ℝ≥0}
    {x y : X₁ × X₂}
    (h : EpsChainIn (φ₁.prod φ₂) (D₁ ×ˢ D₂) ε T x y) :
    EpsChainIn φ₁ D₁ ε T x.1 y.1
\end{leansrc}
\par\medskip\noindent{\small \texttt{PPLC.Qtsl.Semiflow.ChainRecurrentIn.}\texttt{\bfseries fst}}\hfill{\footnotesize \path{lean/qtsl/Qtsl/Product.lean}, line 133}\par\nopagebreak
\begin{leansrc}
theorem ChainRecurrentIn.fst {D₁ : Set X₁} {D₂ : Set X₂} {x : X₁ × X₂}
    (h : ChainRecurrentIn (φ₁.prod φ₂) (D₁ ×ˢ D₂) x) :
    ChainRecurrentIn φ₁ D₁ x.1
\end{leansrc}
\par\medskip\noindent{\small \texttt{PPLC.Qtsl.Semiflow.}\texttt{\bfseries alphaSet\_subset\_closure\_range}}\hfill{\footnotesize \path{lean/qtsl/Qtsl/Product.lean}, line 100}\par\nopagebreak
\begin{leansrc}
theorem alphaSet_subset_closure_range (γ : ℝ → X) :
    alphaSet γ ⊆ closure (Set.range γ)
\end{leansrc}
\par\medskip\noindent{\small \texttt{PPLC.Qtsl.Semiflow.}\texttt{\bfseries omegaSet\_subset\_closure\_range}}\hfill{\footnotesize \path{lean/qtsl/Qtsl/Product.lean}, line 95}\par\nopagebreak
\begin{leansrc}
theorem omegaSet_subset_closure_range (γ : ℝ → X) :
    omegaSet γ ⊆ closure (Set.range γ)
\end{leansrc}
\par\medskip\noindent{\small \texttt{PPLC.Qtsl.Behavior.}\texttt{\bfseries qBlockInv\_sound}}\hfill{\footnotesize \path{lean/qtsl/Qtsl/Behavioral.lean}, line 323}\par\nopagebreak
\begin{leansrc}
theorem qBlockInv_sound {B : Behavior X} {N X0 : Set X}
    (hInv : B.Invariant N) (hInit : X0 ⊆ N) : B.InvScopeSat N X0
\end{leansrc}
\par\medskip\noindent{\small \texttt{PPLC.Qtsl.Behavior.}\texttt{\bfseries invariant\_of\_boundaryDwell}}\hfill{\footnotesize \path{lean/qtsl/Qtsl/Behavioral.lean}, line 213}\par\nopagebreak
\begin{leansrc}
theorem invariant_of_boundaryDwell {B : Behavior X} {N : Set X}
    (hN : IsClosed N)
    (hcont : ∀ γ ∈ B.fwd, Continuous γ)
    (hdwell : ∀ γ ∈ B.fwd, γ 0 ∈ N → γ 0 ∈ frontier N →
      ∃ δ : ℝ≥0, 0 < δ ∧ ∀ s ≤ δ, γ s ∈ N) :
    B.Invariant N
\end{leansrc}
\par\medskip\noindent{\small \texttt{PPLC.Qtsl.Behavior.HasFacetRows.}\texttt{\bfseries boundaryDwell}}\hfill{\footnotesize \path{lean/qtsl/Qtsl/Behavioral.lean}, line 283}\par\nopagebreak
\begin{leansrc}
theorem HasFacetRows.boundaryDwell {B : Behavior X} {D N : Set X}
    (h : B.HasFacetRows D N) (hND : N ⊆ D) (hD : B.Invariant D)
    (hcont : ∀ γ ∈ B.fwd, Continuous γ) :
    ∀ γ ∈ B.fwd, γ 0 ∈ N → γ 0 ∈ frontier N →
      ∃ δ : ℝ≥0, 0 < δ ∧ ∀ s ≤ δ, γ s ∈ N
\end{leansrc}
\par\medskip\noindent{\small \texttt{PPLC.Qtsl.Behavior.}\texttt{\bfseries invariant\_univ}}\hfill{\footnotesize \path{lean/qtsl/Qtsl/Behavioral.lean}, line 320}\par\nopagebreak
\begin{leansrc}
theorem invariant_univ (B : Behavior X) : B.Invariant Set.univ
\end{leansrc}
\par\medskip\noindent{\small \texttt{PPLC.Qtsl.Behavior.}\texttt{\bfseries HasFacetRows}}\hfill{\footnotesize \path{lean/qtsl/Qtsl/Behavioral.lean}, line 200}\par\nopagebreak
\begin{leansrc}
def HasFacetRows (B : Behavior X) [TopologicalSpace X] (D N : Set X) : Prop
\end{leansrc}
\par\medskip\noindent{\small \texttt{PPLC.Qtsl.Semiflow.}\texttt{\bfseries qBlockInv\_sound\_of\_rows}}\hfill{\footnotesize \path{lean/qtsl/Qtsl/MeanValue.lean}, line 166}\par\nopagebreak
\begin{leansrc}
theorem qBlockInv_sound_of_rows {φ : Semiflow X} {D N X0 : Set X}
    {n : ℕ} {ℓ : Fin n → X → ℝ} {ε : Fin n → ℝ} {F : Fin n → Set X}
    (W : Fin n → X → ℝ)
    (hN : IsClosed N)
    (htraj : ∀ x, Continuous fun t : ℝ≥0 => φ.flow t x)
    (hND : N ⊆ D) (hD : φ.Invariant D)
    (hε : ∀ i, 0 < ε i)
    (hℓ : ∀ i, ContinuousOn (ℓ i) (F i))
    (hW : ∀ i, φ.FlowDerivOn (ℓ i) (W i) (F i))
    (hrow : ∀ i, ∀ y ∈ F i, ε i ≤ W i y)
    (hlocal : ∀ x ∈ frontier N, ∃ (A : Set (Fin n)) (U : Set X), U ∈ 𝓝 x ∧
      (∀ i ∈ A, ℓ i x = 0 ∧ F i ∈ 𝓝 x) ∧
      ∀ y ∈ U, y ∈ D → (∀ i ∈ A, 0 ≤ ℓ i y) → y ∈ N)
    (hInit : X0 ⊆ N) : φ.InvScopeSat N X0
\end{leansrc}
\par\medskip\noindent{\small \texttt{PPLC.Qtsl.Grid.}\texttt{\bfseries qBlockInv\_sound\_of\_grid}}\hfill{\footnotesize \path{lean/qtsl/Qtsl/GridGeometry.lean}, line 381}\par\nopagebreak
\begin{leansrc}
theorem qBlockInv_sound_of_grid (φ : Semiflow (Fin d → ℝ)) (B : Finset G.CellIdx)
    {X0 : Set (Fin d → ℝ)}
    (htraj : ∀ x, Continuous fun t : ℝ≥0 => φ.flow t x)
    (hArena : φ.Invariant G.arena)
    (ε : G.FacetData → ℝ) (F : G.FacetData → Set (Fin d → ℝ))
    (hε : ∀ p, G.IsFacet B p → 0 < ε p)
    (hF : ∀ p, G.IsFacet B p → ∀ x ∈ G.facetBox p, F p ∈ 𝓝 x)
    (hrow : ∀ p, G.IsFacet B p → ∀ x ∈ F p, ∀ t : ℝ≥0,
      (∀ s ≤ t, φ.flow s x ∈ F p) →
      G.obs p x + ε p * (t : ℝ) ≤ G.obs p (φ.flow t x))
    (hInit : X0 ⊆ G.blockRegion B) :
    φ.InvScopeSat (G.blockRegion B) X0
\end{leansrc}
\par\medskip\noindent{\small \texttt{PPLC.Qtsl.Grid.}\texttt{\bfseries qBlockInv\_sound\_of\_pointwiseRows}}\hfill{\footnotesize \path{lean/qtsl/Qtsl/GridGeometry.lean}, line 397}\par\nopagebreak
\begin{leansrc}
theorem qBlockInv_sound_of_pointwiseRows (φ : Semiflow (Fin d → ℝ))
    (B : Finset G.CellIdx) {X0 : Set (Fin d → ℝ)}
    (htraj : ∀ x, Continuous fun t : ℝ≥0 => φ.flow t x)
    (hArena : φ.Invariant G.arena)
    (ε : G.FacetData → ℝ) (F : G.FacetData → Set (Fin d → ℝ))
    (W : G.FacetData → (Fin d → ℝ) → ℝ)
    (hε : ∀ p, G.IsFacet B p → 0 < ε p)
    (hF : ∀ p, G.IsFacet B p → ∀ x ∈ G.facetBox p, F p ∈ 𝓝 x)
    (hV : ∀ p, G.IsFacet B p → ContinuousOn (G.obs p) (F p))
    (hW : ∀ p, G.IsFacet B p → φ.FlowDerivOn (G.obs p) (W p) (F p))
    (hrow : ∀ p, G.IsFacet B p → ∀ y ∈ F p, ε p ≤ W p y)
    (hInit : X0 ⊆ G.blockRegion B) :
    φ.InvScopeSat (G.blockRegion B) X0
\end{leansrc}
\par\medskip\noindent{\small \texttt{PPLC.Qtsl.Semiflow.}\texttt{\bfseries qBlockInv\_sound\_of\_dwell}}\hfill{\footnotesize \path{lean/qtsl/Qtsl/Discharge.lean}, line 168}\par\nopagebreak
\begin{leansrc}
theorem qBlockInv_sound_of_dwell {φ : Semiflow X} {N X0 : Set X}
    (hN : IsClosed N)
    (htraj : ∀ x, Continuous fun t : ℝ≥0 => φ.flow t x)
    (hdwell : ∀ x ∈ N, x ∈ frontier N →
      ∃ δ : ℝ≥0, 0 < δ ∧ ∀ s ≤ δ, φ.flow s x ∈ N)
    (hInit : X0 ⊆ N) : φ.InvScopeSat N X0
\end{leansrc}
\par\medskip\noindent{\small \texttt{PPLC.Qtsl.Behavior.}\texttt{\bfseries qBlockInv\_sound\_of\_facetRows}}\hfill{\footnotesize \path{lean/qtsl/Qtsl/Behavioral.lean}, line 310}\par\nopagebreak
\begin{leansrc}
theorem qBlockInv_sound_of_facetRows {B : Behavior X} {D N X0 : Set X}
    (hN : IsClosed N)
    (hcont : ∀ γ ∈ B.fwd, Continuous γ)
    (hND : N ⊆ D) (hD : B.Invariant D) (hrows : B.HasFacetRows D N)
    (hInit : X0 ⊆ N) : B.InvScopeSat N X0
\end{leansrc}
\par\medskip\noindent{\small \texttt{PPLC.Qtsl.Semiflow.}\texttt{\bfseries invariant\_of\_invScopeSat}}\hfill{\footnotesize \path{lean/qtsl/Qtsl/GridGeometry.lean}, line 422}\par\nopagebreak
\begin{leansrc}
theorem invariant_of_invScopeSat {φ : Semiflow X} {D X0 : Set X}
    (h : φ.InvScopeSat D X0) : φ.Invariant D
\end{leansrc}
\par\medskip\noindent{\small \texttt{PPLC.Qtsl.Semiflow.}\texttt{\bfseries integrated\_ge\_of\_flowDeriv\_ge}}\hfill{\footnotesize \path{lean/qtsl/Qtsl/MeanValue.lean}, line 105}\par\nopagebreak
\begin{leansrc}
theorem integrated_ge_of_flowDeriv_ge {φ : Semiflow X} {C : Set X}
    {V W : X → ℝ} {c : ℝ}
    (htraj : ∀ x, Continuous fun t : ℝ≥0 => φ.flow t x)
    (hV : ContinuousOn V C)
    (hW : φ.FlowDerivOn V W C)
    (hrow : ∀ y ∈ C, c ≤ W y)
    {x : X} {t : ℝ≥0} (hconf : ∀ s ≤ t, φ.flow s x ∈ C) :
    V x + c * t ≤ V (φ.flow t x)
\end{leansrc}
\par\medskip\noindent{\small \texttt{PPLC.Qtsl.Grid.}\texttt{\bfseries wedge\_containment}}\hfill{\footnotesize \path{lean/qtsl/Qtsl/GridGeometry.lean}, line 106}\par\nopagebreak
\begin{leansrc}
theorem wedge_containment (B : Finset G.CellIdx) (x : Fin d → ℝ)
    (hx : x ∈ G.blockRegion B) :
    ∃ U ∈ 𝓝 x, ∀ y ∈ U, y ∈ G.arena →
      (∀ p : G.FacetData, G.IsFacet B p → x ∈ G.facetBox p → 0 ≤ G.obs p y) →
      y ∈ G.blockRegion B
\end{leansrc}
\par\medskip\noindent{\small \texttt{PPLC.Qtsl.Grid.}\texttt{\bfseries qBlockInv\_sound\_of\_grid}}\hfill{\footnotesize \path{lean/qtsl/Qtsl/GridGeometry.lean}, line 381}\par\nopagebreak
\begin{leansrc}
theorem qBlockInv_sound_of_grid (φ : Semiflow (Fin d → ℝ)) (B : Finset G.CellIdx)
    {X0 : Set (Fin d → ℝ)}
    (htraj : ∀ x, Continuous fun t : ℝ≥0 => φ.flow t x)
    (hArena : φ.Invariant G.arena)
    (ε : G.FacetData → ℝ) (F : G.FacetData → Set (Fin d → ℝ))
    (hε : ∀ p, G.IsFacet B p → 0 < ε p)
    (hF : ∀ p, G.IsFacet B p → ∀ x ∈ G.facetBox p, F p ∈ 𝓝 x)
    (hrow : ∀ p, G.IsFacet B p → ∀ x ∈ F p, ∀ t : ℝ≥0,
      (∀ s ≤ t, φ.flow s x ∈ F p) →
      G.obs p x + ε p * (t : ℝ) ≤ G.obs p (φ.flow t x))
    (hInit : X0 ⊆ G.blockRegion B) :
    φ.InvScopeSat (G.blockRegion B) X0
\end{leansrc}
\par\medskip\noindent{\small \texttt{PPLC.Qtsl.Grid.}\texttt{\bfseries qBlockInv\_sound\_of\_pointwiseRows}}\hfill{\footnotesize \path{lean/qtsl/Qtsl/GridGeometry.lean}, line 397}\par\nopagebreak
\begin{leansrc}
theorem qBlockInv_sound_of_pointwiseRows (φ : Semiflow (Fin d → ℝ))
    (B : Finset G.CellIdx) {X0 : Set (Fin d → ℝ)}
    (htraj : ∀ x, Continuous fun t : ℝ≥0 => φ.flow t x)
    (hArena : φ.Invariant G.arena)
    (ε : G.FacetData → ℝ) (F : G.FacetData → Set (Fin d → ℝ))
    (W : G.FacetData → (Fin d → ℝ) → ℝ)
    (hε : ∀ p, G.IsFacet B p → 0 < ε p)
    (hF : ∀ p, G.IsFacet B p → ∀ x ∈ G.facetBox p, F p ∈ 𝓝 x)
    (hV : ∀ p, G.IsFacet B p → ContinuousOn (G.obs p) (F p))
    (hW : ∀ p, G.IsFacet B p → φ.FlowDerivOn (G.obs p) (W p) (F p))
    (hrow : ∀ p, G.IsFacet B p → ∀ y ∈ F p, ε p ≤ W p y)
    (hInit : X0 ⊆ G.blockRegion B) :
    φ.InvScopeSat (G.blockRegion B) X0
\end{leansrc}
\par\medskip\noindent{\small \texttt{PPLC.Qtsl.}\texttt{\bfseries qLatticeInv\_exact}}\hfill{\footnotesize \path{lean/qtsl/Qtsl/LatticeInv.lean}, line 28}\par\nopagebreak
\begin{leansrc}
theorem qLatticeInv_exact
    (hdc : IsLowerSet C)
    (hJoin : ∀ b : β, ∀ x ∈ blocks b,
      ∃ i : ι, x ∈ blocks (emb i) ∧ blocks (emb i) ⊆ blocks b)
    (hcover : ∀ i : ι, blocks (cover i) ⊂ blocks (emb i))
    (hcellsSup : ∀ i : ι, blocks (emb i) \ blocks (cover i) ⊆ cells i)
    (horder : ∀ i j : ι, blocks (emb i) ⊆ blocks (emb j) → i ≤ j)
    (hcellSub : ∀ i : ι, cells i ⊆ blocks (emb i)) :
    (⋃ i ∈ C, cells i) = ⋃ i ∈ C, blocks (emb i)
\end{leansrc}
\par\medskip\noindent{\small \texttt{PPLC.Qtsl.Semiflow.}\texttt{\bfseries exists\_cellMem\_subset}}\hfill{\footnotesize \path{lean/qtsl/Qtsl/EventualMember.lean}, line 112}\par\nopagebreak
\begin{leansrc}
theorem exists_cellMem_subset
    {blocks : β → Set X} {emb : ι → β} {cover : ι → β} {cells : ι → Set X}
    (hJoin : ∀ b : β, ∀ x ∈ blocks b,
      ∃ i : ι, x ∈ blocks (emb i) ∧ blocks (emb i) ⊆ blocks b)
    (hcover : ∀ i : ι, blocks (cover i) ⊂ blocks (emb i))
    (hcellsSup : ∀ i : ι, blocks (emb i) \ blocks (cover i) ⊆ cells i)
    {x : X} (i : ι) (hmem : x ∈ blocks (emb i)) :
    ∃ j : ι, blocks (emb j) ⊆ blocks (emb i) ∧
      x ∈ blocks (emb j) ∧ x ∈ cells j
\end{leansrc}
\par\medskip\noindent{\small \texttt{PPLC.Qtsl.Semiflow.}\texttt{\bfseries qLatticeInv\_sound}}\hfill{\footnotesize \path{lean/qtsl/Qtsl/LatticeInv.lean}, line 65}\par\nopagebreak
\begin{leansrc}
theorem Semiflow.qLatticeInv_sound {φ : Semiflow X} {X0 : Set X}
    (hdc : IsLowerSet C)
    (hJoin : ∀ b : β, ∀ x ∈ blocks b,
      ∃ i : ι, x ∈ blocks (emb i) ∧ blocks (emb i) ⊆ blocks b)
    (hcover : ∀ i : ι, blocks (cover i) ⊂ blocks (emb i))
    (hcellsSup : ∀ i : ι, blocks (emb i) \ blocks (cover i) ⊆ cells i)
    (horder : ∀ i j : ι, blocks (emb i) ⊆ blocks (emb j) → i ≤ j)
    (hcellSub : ∀ i : ι, cells i ⊆ blocks (emb i))
    (hMember : φ.InvScopeSat (⋃ i ∈ C, blocks (emb i)) X0) :
    φ.InvScopeSat (⋃ i ∈ C, cells i) X0
\end{leansrc}
\par\medskip\noindent{\small \texttt{PPLC.Qtsl.Behavior.}\texttt{\bfseries qLatticeInv\_sound}}\hfill{\footnotesize \path{lean/qtsl/Qtsl/LatticeInv.lean}, line 93}\par\nopagebreak
\begin{leansrc}
theorem Behavior.qLatticeInv_sound {B : Behavior X} {X0 : Set X}
    (hdc : IsLowerSet C)
    (hJoin : ∀ b : β, ∀ x ∈ blocks b,
      ∃ i : ι, x ∈ blocks (emb i) ∧ blocks (emb i) ⊆ blocks b)
    (hcover : ∀ i : ι, blocks (cover i) ⊂ blocks (emb i))
    (hcellsSup : ∀ i : ι, blocks (emb i) \ blocks (cover i) ⊆ cells i)
    (horder : ∀ i j : ι, blocks (emb i) ⊆ blocks (emb j) → i ≤ j)
    (hcellSub : ∀ i : ι, cells i ⊆ blocks (emb i))
    (hMember : B.InvScopeSat (⋃ i ∈ C, blocks (emb i)) X0) :
    B.InvScopeSat (⋃ i ∈ C, cells i) X0
\end{leansrc}
\par\medskip\noindent{\small \texttt{PPLC.Qtsl.Behavior.}\texttt{\bfseries qBasin\_sound}}\hfill{\footnotesize \path{lean/qtsl/Qtsl/Behavioral.lean}, line 327}\par\nopagebreak
\begin{leansrc}
theorem qBasin_sound {B : Behavior X} {S X0 Bad : Set X}
    (hscope : B.InvScopeSat S X0) (hAvoid : S ∩ Bad = ∅) :
    B.SafetySat X0 Bad
\end{leansrc}
\par\medskip\noindent{\small \texttt{PPLC.Qtsl.}\texttt{\bfseries member\_subset\_cells\_le}}\hfill{\footnotesize \path{lean/qtsl/Qtsl/LatticeInv.lean}, line 44}\par\nopagebreak
\begin{leansrc}
theorem member_subset_cells_le
    (hJoin : ∀ b : β, ∀ x ∈ blocks b,
      ∃ i : ι, x ∈ blocks (emb i) ∧ blocks (emb i) ⊆ blocks b)
    (hcover : ∀ i : ι, blocks (cover i) ⊂ blocks (emb i))
    (hcellsSup : ∀ i : ι, blocks (emb i) \ blocks (cover i) ⊆ cells i)
    (horder : ∀ i j : ι, blocks (emb i) ⊆ blocks (emb j) → i ≤ j)
    (i : ι) :
    blocks (emb i) ⊆ ⋃ j ∈ Set.Iic i, cells j
\end{leansrc}
\par\medskip\noindent{\small \texttt{PPLC.Qtsl.Semiflow.}\texttt{\bfseries qLatticeInv\_sound}}\hfill{\footnotesize \path{lean/qtsl/Qtsl/LatticeInv.lean}, line 65}\par\nopagebreak
\begin{leansrc}
theorem Semiflow.qLatticeInv_sound {φ : Semiflow X} {X0 : Set X}
    (hdc : IsLowerSet C)
    (hJoin : ∀ b : β, ∀ x ∈ blocks b,
      ∃ i : ι, x ∈ blocks (emb i) ∧ blocks (emb i) ⊆ blocks b)
    (hcover : ∀ i : ι, blocks (cover i) ⊂ blocks (emb i))
    (hcellsSup : ∀ i : ι, blocks (emb i) \ blocks (cover i) ⊆ cells i)
    (horder : ∀ i j : ι, blocks (emb i) ⊆ blocks (emb j) → i ≤ j)
    (hcellSub : ∀ i : ι, cells i ⊆ blocks (emb i))
    (hMember : φ.InvScopeSat (⋃ i ∈ C, blocks (emb i)) X0) :
    φ.InvScopeSat (⋃ i ∈ C, cells i) X0
\end{leansrc}
\par\medskip\noindent{\small \texttt{PPLC.Qtsl.Behavior.}\texttt{\bfseries qBasin\_batch}}\hfill{\footnotesize \path{lean/qtsl/Qtsl/Behavioral.lean}, line 336}\par\nopagebreak
\begin{leansrc}
theorem qBasin_batch {B : Behavior X} {S X0 : Set X} {J : Type*}
    {Bad : J → Set X} (hscope : B.InvScopeSat S X0)
    (hAvoid : ∀ j, S ∩ Bad j = ∅) :
    B.SafetySat X0 (⋃ j, Bad j)
\end{leansrc}
\par\medskip\noindent{\small \texttt{PPLC.Qtsl.Behavior.}\texttt{\bfseries InvScopeSat}}\hfill{\footnotesize \path{lean/qtsl/Qtsl/Behavioral.lean}, line 45}\par\nopagebreak
\begin{leansrc}
def InvScopeSat (B : Behavior X) (S X0 : Set X) : Prop
\end{leansrc}
\par\medskip\noindent{\small \texttt{PPLC.Qtsl.Behavior.}\texttt{\bfseries SafetySat}}\hfill{\footnotesize \path{lean/qtsl/Qtsl/Behavioral.lean}, line 42}\par\nopagebreak
\begin{leansrc}
def SafetySat (B : Behavior X) (X0 Bad : Set X) : Prop
\end{leansrc}
\par\medskip\noindent{\small \texttt{PPLC.Qtsl.Semiflow.}\texttt{\bfseries InvScopeSat}}\hfill{\footnotesize \path{lean/qtsl/Qtsl/QtslSoundness.lean}, line 84}\par\nopagebreak
\begin{leansrc}
def InvScopeSat (φ : Semiflow X) (S X0 : Set X) : Prop
\end{leansrc}
\par\medskip\noindent{\small \texttt{PPLC.Qtsl.Semiflow.}\texttt{\bfseries SafetySat}}\hfill{\footnotesize \path{lean/qtsl/Qtsl/QtslSoundness.lean}, line 87}\par\nopagebreak
\begin{leansrc}
def SafetySat (φ : Semiflow X) (X0 B : Set X) : Prop
\end{leansrc}
\par\medskip\noindent{\small \texttt{Qtsl.QtslCompose.System.}\texttt{\bfseries qagTier}}\hfill{\footnotesize \path{lean/qtsl-aux/QtslAux/QtslCompose/Presentation.lean}, line 111}\par\nopagebreak
\begin{leansrc}
def System.qagTier {n : ℕ} (S : System n) : Tier
\end{leansrc}
\par\medskip\noindent{\small \texttt{Qtsl.QtslCompose.}\texttt{\bfseries admissible\_constructs}}\hfill{\footnotesize \path{lean/qtsl-aux/QtslAux/QtslCompose/Presentation.lean}, line 163}\par\nopagebreak
\begin{leansrc}
theorem admissible_constructs {ι κIdx X : Type*} [Fintype ι] [Fintype κIdx]
    {R : Set X} (lam : ι → ℝ) (e : ι → X → ℝ) (nu : κIdx → ℝ) (g : κIdx → X → ℝ)
    (hlam : ∀ i, 0 ≤ lam i)
    (he : ∀ i, ∀ x ∈ R, e i x ≤ 0)
    (hg : ∀ k, ∀ x ∈ R, g k x = 0) :
    ∀ x ∈ R, (∑ i, lam i * e i x) + (∑ k, nu k * g k x) ≤ 0
\end{leansrc}
\par\medskip\noindent{\small \texttt{Qtsl.AGCompose.Parallel.}\texttt{\bfseries fail\_bound}}\hfill{\footnotesize \path{lean/qtsl-aux/QtslAux/AGCompose/GshsParallel.lean}, line 128}\par\nopagebreak
\begin{leansrc}
theorem fail_bound {ι : Type*} (s : Finset ι) {sat : ι → Set Ω} {satOut : Set Ω}
    {ε : ι → ℝ≥0∞}
    (h : ∀ i ∈ s, μ (sat i)ᶜ ≤ ε i)
    (hout : (⋂ i ∈ s, sat i) ⊆ satOut) :
    μ satOutᶜ ≤ ∑ i ∈ s, ε i
\end{leansrc}
\par\medskip\noindent{\small \texttt{Qtsl.QtslCompose.}\texttt{\bfseries fail\_bound\_out}}\hfill{\footnotesize \path{lean/qtsl-aux/QtslAux/QtslCompose/AGCompose.lean}, line 131}\par\nopagebreak
\begin{leansrc}
theorem fail_bound_out [MeasurableSpace Ω] {μ : Measure Ω} {ε : Fin n → ℝ≥0∞}
    {satOut : Set Ω}
    (hw : ∀ j, conformExt ∩ (⋂ i ∈ S.upstream j, sat i) ⊆ conform j)
    (hout : conformExt ∩ (⋂ j, sat j) ⊆ satOut)
    (hε : ∀ j, μ (conform j ∩ (sat j)ᶜ) ≤ ε j) :
    μ (conformExt ∩ satOutᶜ) ≤ ∑ j, ε j
\end{leansrc}
\par\medskip\noindent{\small \texttt{Qtsl.QtslCompose.}\texttt{\bfseries sat\_of\_conform}}\hfill{\footnotesize \path{lean/qtsl-aux/QtslAux/QtslCompose/AGCompose.lean}, line 48}\par\nopagebreak
\begin{leansrc}
theorem sat_of_conform
    (hw : ∀ j, conformExt ∩ (⋂ i ∈ S.upstream j, sat i) ⊆ conform j)
    {ω : Ω} (h : ∀ j, ω ∈ conform j → ω ∈ sat j)
    (hext : ω ∈ conformExt) : ∀ j, ω ∈ sat j
\end{leansrc}
\par\medskip\noindent{\small \texttt{Qtsl.AGCompose.Series.}\texttt{\bfseries fail\_subset}}\hfill{\footnotesize \path{lean/qtsl-aux/QtslAux/AGCompose/GshsSeries.lean}, line 45}\par\nopagebreak
\begin{leansrc}
theorem fail_subset
    (hw : conform₁ ∩ sat₁ ⊆ conform₂)
    (hout : conform₁ ∩ sat₁ ∩ sat₂ ⊆ satOut) :
    conform₁ ∩ satOutᶜ ⊆ (conform₁ ∩ sat₁ᶜ) ∪ (conform₂ ∩ sat₂ᶜ)
\end{leansrc}
\par\medskip\noindent{\small \texttt{Qtsl.QtslCompose.System.}\texttt{\bfseries qagTier\_admissible\_iff}}\hfill{\footnotesize \path{lean/qtsl-aux/QtslAux/QtslCompose/Presentation.lean}, line 116}\par\nopagebreak
\begin{leansrc}
theorem System.qagTier_admissible_iff {n : ℕ} (S : System n) :
    S.qagTier = .admissible ↔
      ∀ i, (S.charts i).profile = .ode ∨ (S.charts i).profile = .discrete
\end{leansrc}
\par\medskip\noindent{\small \texttt{Qtsl.QtslCompose.System.}\texttt{\bfseries qagTier\_core\_iff}}\hfill{\footnotesize \path{lean/qtsl-aux/QtslAux/QtslCompose/Presentation.lean}, line 135}\par\nopagebreak
\begin{leansrc}
theorem System.qagTier_core_iff {n : ℕ} (S : System n) :
    S.qagTier = .core ↔
      ∃ i, (S.charts i).profile = .sde ∨ (S.charts i).profile = .gshs
\end{leansrc}
\par\medskip\noindent{\small \texttt{Qtsl.AGCompose.Series.}\texttt{\bfseries certified}}\hfill{\footnotesize \path{lean/qtsl-aux/QtslAux/AGCompose/GshsSeries.lean}, line 132}\par\nopagebreak
\begin{leansrc}
theorem certified [IsFiniteMeasure μ]
    {Y₁ Y₂ : ℕ → Ω → ℝ} {X₁ : ℕ → Ω → S₁} {X₂ : ℕ → Ω → S₂}
    {B₁ : Set S₁} {B₂ : Set S₂} {a₁ a₂ : ℝ}

    (hY₁ : Supermartingale Y₁ 𝒢 μ) (hY₁0 : 0 ≤ Y₁)
    (hY₂ : Supermartingale Y₂ 𝒢 μ) (hY₂0 : 0 ≤ Y₂)

    (ha₁ : 0 < a₁) (ha₂ : 0 < a₂)
    (hbar₁ : ∀ n ω, X₁ n ω ∈ B₁ → a₁ ≤ Y₁ n ω)
    (hbar₂ : ∀ n ω, X₂ n ω ∈ B₂ → a₂ ≤ Y₂ n ω)

    (hwit₁ : conform₁ ∩ sat₁ᶜ ⊆ {ω | ∃ n, X₁ n ω ∈ B₁})
    (hwit₂ : conform₂ ∩ sat₂ᶜ ⊆ {ω | ∃ n, X₂ n ω ∈ B₂})

    (hw : conform₁ ∩ sat₁ ⊆ conform₂)
    (hout : conform₁ ∩ sat₁ ∩ sat₂ ⊆ satOut) :
    μ (conform₁ ∩ satOutᶜ) ≤
      ENNReal.ofReal (∫ ω, Y₁ 0 ω ∂μ) / ENNReal.ofReal a₁ +
      ENNReal.ofReal (∫ ω, Y₂ 0 ω ∂μ) / ENNReal.ofReal a₂
\end{leansrc}
\par\medskip\noindent{\small \texttt{Qtsl.AGCompose.Parallel.}\texttt{\bfseries certified}}\hfill{\footnotesize \path{lean/qtsl-aux/QtslAux/AGCompose/GshsParallel.lean}, line 58}\par\nopagebreak
\begin{leansrc}
theorem certified (s : Finset ι) [IsFiniteMeasure μ]
    {Y : ι → ℕ → Ω → ℝ} {lam : ι → ℝ} {X : ℕ → Ω → St} {Bad : ι → Set St} {a : ℝ}
    (hY : ∀ i ∈ s, Supermartingale (Y i) 𝒢 μ)
    (hY0 : ∀ i ∈ s, 0 ≤ Y i)
    (hlam : ∀ i ∈ s, 0 ≤ lam i)
    (hbar : ∀ i ∈ s, ∀ n ω, X n ω ∈ Bad i → a ≤ lam i * Y i n ω) :
    ENNReal.ofReal a * μ {ω | ∃ n, X n ω ∈ ⋃ i ∈ s, Bad i} ≤
      ENNReal.ofReal (∑ i ∈ s, lam i * ∫ ω, Y i 0 ω ∂μ)
\end{leansrc}
\par\medskip\noindent{\small \texttt{Qtsl.AGCompose.Parallel.}\texttt{\bfseries certified\_pair}}\hfill{\footnotesize \path{lean/qtsl-aux/QtslAux/AGCompose/GshsParallel.lean}, line 87}\par\nopagebreak
\begin{leansrc}
theorem certified_pair [IsFiniteMeasure μ]
    {Y₁ Y₂ : ℕ → Ω → ℝ} {X : ℕ → Ω → St} {B₁ B₂ : Set St} {lam a : ℝ}
    (hY₁ : Supermartingale Y₁ 𝒢 μ) (hY₁0 : 0 ≤ Y₁)
    (hY₂ : Supermartingale Y₂ 𝒢 μ) (hY₂0 : 0 ≤ Y₂) (hlam : 0 ≤ lam)
    (hbar₁ : ∀ n ω, X n ω ∈ B₁ → a ≤ Y₁ n ω)
    (hbar₂ : ∀ n ω, X n ω ∈ B₂ → a ≤ lam * Y₂ n ω) :
    ENNReal.ofReal a * μ {ω | ∃ n, X n ω ∈ B₁ ∪ B₂} ≤
      ENNReal.ofReal (∫ ω, Y₁ 0 ω ∂μ + lam * ∫ ω, Y₂ 0 ω ∂μ)
\end{leansrc}
\par\medskip\noindent{\small \texttt{Qtsl.AGCompose.SumFormFeedback.}\texttt{\bfseries avoid\_bound}}\hfill{\footnotesize \path{lean/qtsl-aux/QtslAux/GshsFeedbackSum.lean}, line 166}\par\nopagebreak
\begin{leansrc}
theorem avoid_bound {S : Type*} {X : ℕ → Ω → S} {B : Set S} {a : ℝ}
    (hbarrier : ∀ k ω, X k ω ∈ B → a ≤ F.W k ω) :
    ENNReal.ofReal a * μ {ω | ∃ k, X k ω ∈ B} ≤ ENNReal.ofReal (∫ ω, F.W 0 ω ∂μ)
\end{leansrc}
\par\medskip\noindent{\small \texttt{Qtsl.AGCompose.}\texttt{\bfseries SumFormFeedback}}\hfill{\footnotesize \path{lean/qtsl-aux/QtslAux/GshsFeedbackSum.lean}, line 79}\par\nopagebreak
\begin{leansrc}
structure SumFormFeedback (μ : Measure Ω) (𝒢 : Filtration ℕ m0) where …
\end{leansrc}
\par\medskip\noindent{\small \texttt{Qtsl.LayerD.}\texttt{\bfseries sound\_GSHS\_safety\_barrier}}\hfill{\footnotesize \path{lean/qtsl-aux/QtslAux/MaxFormPDMP.lean}, line 608}\par\nopagebreak
\begin{leansrc}
theorem sound_GSHS_safety_barrier [IsMarkovKernel M.seg] [IsMarkovKernel M.jump]
    {W : State → ℝ} {U : Set State} {m : ℝ} (hm : 0 < m)
    (hW : Measurable W) (hU : MeasurableSet U)
    (hSegD : SegmentDynkin M 0 W) (hForce : ForcedNonexpansion M W)
    (hNonneg : NonnegV W) {C : ℝ} (hC0 : 0 ≤ C) (hWC : ∀ y, W y ≤ C)
    (hFloor : ∀ y ∈ U, m ≤ W y) (x : State) :
    Battery.pathMeasure (skel M) (SkelState.alive x) {ω | ∃ n, ω n ∈ skelBad U}
      ≤ ENNReal.ofReal (W x / m)
\end{leansrc}
\par\medskip\noindent{\small \texttt{Qtsl.QtslCompose.Discrete.}\texttt{\bfseries firstFail}}\hfill{\footnotesize \path{lean/qtsl-aux/QtslAux/QtslCompose/AGComposeDiscrete.lean}, line 52}\par\nopagebreak
\begin{leansrc}
def firstFail (bad : Fin m → Set S) (j : Fin m) : Set (ℕ → S)
\end{leansrc}
\par\medskip\noindent{\small \texttt{Qtsl.QtslCompose.Discrete.}\texttt{\bfseries measure\_firstFail\_le}}\hfill{\footnotesize \path{lean/qtsl-aux/QtslAux/QtslCompose/AGComposeDiscrete.lean}, line 86}\par\nopagebreak
\begin{leansrc}
theorem measure_firstFail_le {bad : Fin m → Set S} {V : S → ℝ} {C : ℝ} {j : Fin m}
    (hA : MeasurableSet (⋃ i, bad i)) (hV : Measurable V)
    (hV0 : ∀ y, 0 ≤ V y) (hVC : ∀ y, V y ≤ C)
    (hbar : ∀ y ∈ bad j, 1 ≤ V y)
    (hdrift : ∀ y, y ∉ (⋃ i, bad i) → ∫ z, V z ∂K y ≤ V y) (x : S) :
    pathMeasure K x (firstFail bad j) ≤ ENNReal.ofReal (V x)
\end{leansrc}
\par\medskip\noindent{\small \texttt{Qtsl.QtslCompose.Discrete.}\texttt{\bfseries fail\_bound\_wired}}\hfill{\footnotesize \path{lean/qtsl-aux/QtslAux/QtslCompose/AGComposeDiscrete.lean}, line 156}\par\nopagebreak
\begin{leansrc}
theorem fail_bound_wired (Sys : System m) {conf bad : Fin m → Set S}
    {V : Fin m → S → ℝ} {C : Fin m → ℝ}
    (hA : MeasurableSet (⋃ i, bad i))
    (hw : ∀ j, (⋂ i ∈ Sys.upstream j, (bad i)ᶜ) ⊆ conf j)
    (hV : ∀ j, Measurable (V j)) (hV0 : ∀ j y, 0 ≤ V j y)
    (hVC : ∀ j y, V j y ≤ C j)
    (hbar : ∀ j, ∀ y ∈ bad j, 1 ≤ V j y)
    (hdrift : ∀ j, ∀ y ∈ conf j, y ∉ bad j → ∫ z, V j z ∂K y ≤ V j y) (x : S) :
    pathMeasure K x {ω | ∃ n, ω n ∈ ⋃ i, bad i}
      ≤ ∑ j, ENNReal.ofReal (V j x)
\end{leansrc}
\par\medskip\noindent{\small \texttt{MeasureTheory.Supermartingale.}\texttt{\bfseries ville\_ineq}}\hfill{\footnotesize \path{lean/qtsl-aux/QtslAux/Ville.lean}, line 133}\par\nopagebreak
\begin{leansrc}
theorem Supermartingale.ville_ineq [IsFiniteMeasure μ] (hsup : Supermartingale f 𝒢
      μ)
    (hnonneg : 0 ≤ f) (ε : ℝ≥0) :
    ε * μ {ω | ∃ k, (ε : ℝ) ≤ f k ω} ≤ ENNReal.ofReal (∫ ω, f 0 ω ∂μ)
\end{leansrc}
\par\medskip\noindent{\small \texttt{Qtsl.QtslCompose.Discrete.}\texttt{\bfseries drift\_parallelComp}}\hfill{\footnotesize \path{lean/qtsl-aux/QtslAux/QtslCompose/AGComposeDiscrete.lean}, line 207}\par\nopagebreak
\begin{leansrc}
theorem drift_parallelComp {bad₁ : Set S₁} {bad₂ : Set S₂}
    {V₁ : S₁ → ℝ} {V₂ : S₂ → ℝ} {lam : ℝ} (hlam : 0 ≤ lam)
    (hint₁ : ∀ y, y ∉ bad₁ → Integrable V₁ (K₁ y))
    (hint₂ : ∀ y, y ∉ bad₂ → Integrable V₂ (K₂ y))
    (hd₁ : ∀ y, y ∉ bad₁ → ∫ z, V₁ z ∂K₁ y ≤ V₁ y)
    (hd₂ : ∀ y, y ∉ bad₂ → ∫ z, V₂ z ∂K₂ y ≤ V₂ y) :
    ∀ p : S₁ × S₂, p ∉ (Prod.fst ⁻¹' bad₁ ∪ Prod.snd ⁻¹' bad₂) →
      ∫ q, (V₁ q.1 + lam * V₂ q.2) ∂(K₁ ∥ₖ K₂) p
        ≤ V₁ p.1 + lam * V₂ p.2
\end{leansrc}
\par\medskip\noindent{\small \texttt{Qtsl.QtslCompose.Discrete.}\texttt{\bfseries avoid\_parallelComp}}\hfill{\footnotesize \path{lean/qtsl-aux/QtslAux/QtslCompose/AGComposeDiscrete.lean}, line 233}\par\nopagebreak
\begin{leansrc}
theorem avoid_parallelComp {bad₁ : Set S₁} {bad₂ : Set S₂}
    {V₁ : S₁ → ℝ} {V₂ : S₂ → ℝ} {C₁ C₂ lam : ℝ} (hlam : 0 ≤ lam)
    (hbad₁ : MeasurableSet bad₁) (hbad₂ : MeasurableSet bad₂)
    (hV₁ : Measurable V₁) (hV₂ : Measurable V₂)
    (hV₁0 : ∀ y, 0 ≤ V₁ y) (hV₂0 : ∀ y, 0 ≤ V₂ y)
    (hV₁C : ∀ y, V₁ y ≤ C₁) (hV₂C : ∀ y, V₂ y ≤ C₂)
    (hbar₁ : ∀ y ∈ bad₁, 1 ≤ V₁ y) (hbar₂ : ∀ y ∈ bad₂, 1 ≤ lam * V₂ y)
    (hd₁ : ∀ y, y ∉ bad₁ → ∫ z, V₁ z ∂K₁ y ≤ V₁ y)
    (hd₂ : ∀ y, y ∉ bad₂ → ∫ z, V₂ z ∂K₂ y ≤ V₂ y) (x : S₁ × S₂) :
    pathMeasure (K₁ ∥ₖ K₂) x
        {ω | ∃ n, ω n ∈ (Prod.fst ⁻¹' bad₁ ∪ Prod.snd ⁻¹' bad₂ : Set (S₁ × S₂))}
      ≤ ENNReal.ofReal (V₁ x.1 + lam * V₂ x.2)
\end{leansrc}
\par\medskip\noindent{\small \texttt{Qtsl.QtslCompose.Discrete.}\texttt{\bfseries smallGain\_drift}}\hfill{\footnotesize \path{lean/qtsl-aux/QtslAux/QtslCompose/AGComposeDiscrete.lean}, line 278}\par\nopagebreak
\begin{leansrc}
theorem smallGain_drift (hlam : 0 ≤ lam)
    (hint₁ : ∀ y, y ∉ A → Integrable V₁ (K y))
    (hint₂ : ∀ y, y ∉ A → Integrable V₂ (K y))
    (hd₁ : ∀ y, y ∉ A → ∫ z, V₁ z ∂K y ≤ V₁ y - alpha₁ y + sigma₁ y)
    (hd₂ : ∀ y, y ∉ A → ∫ z, V₂ z ∂K y ≤ V₂ y - alpha₂ y + sigma₂ y)
    (hm₁ : ∀ y, y ∉ A → lam * sigma₂ y ≤ (1 - theta) * alpha₁ y)
    (hm₂ : ∀ y, y ∉ A → sigma₁ y ≤ lam * ((1 - theta) * alpha₂ y)) :
    ∀ y, y ∉ A → ∫ z, (V₁ z + lam * V₂ z) ∂K y
      ≤ (V₁ y + lam * V₂ y) - theta * (alpha₁ y + lam * alpha₂ y)
\end{leansrc}
\par\medskip\noindent{\small \texttt{Qtsl.QtslCompose.Discrete.}\texttt{\bfseries avoid\_smallGain}}\hfill{\footnotesize \path{lean/qtsl-aux/QtslAux/QtslCompose/AGComposeDiscrete.lean}, line 310}\par\nopagebreak
\begin{leansrc}
theorem avoid_smallGain (hlam : 0 ≤ lam) (htheta : 0 ≤ theta)
    {C₁ C₂ : ℝ} (hA : MeasurableSet A)
    (hV₁ : Measurable V₁) (hV₂ : Measurable V₂)
    (hV₁0 : ∀ y, 0 ≤ V₁ y) (hV₂0 : ∀ y, 0 ≤ V₂ y)
    (hV₁C : ∀ y, V₁ y ≤ C₁) (hV₂C : ∀ y, V₂ y ≤ C₂)
    (hbar : ∀ y ∈ A, 1 ≤ V₁ y + lam * V₂ y)
    (halpha₁ : ∀ y, y ∉ A → 0 ≤ alpha₁ y)
    (halpha₂ : ∀ y, y ∉ A → 0 ≤ alpha₂ y)
    (hd₁ : ∀ y, y ∉ A → ∫ z, V₁ z ∂K y ≤ V₁ y - alpha₁ y + sigma₁ y)
    (hd₂ : ∀ y, y ∉ A → ∫ z, V₂ z ∂K y ≤ V₂ y - alpha₂ y + sigma₂ y)
    (hm₁ : ∀ y, y ∉ A → lam * sigma₂ y ≤ (1 - theta) * alpha₁ y)
    (hm₂ : ∀ y, y ∉ A → sigma₁ y ≤ lam * ((1 - theta) * alpha₂ y)) (x : S) :
    pathMeasure K x {ω | ∃ n, ω n ∈ A}
      ≤ ENNReal.ofReal (V₁ x + lam * V₂ x)
\end{leansrc}
\par\medskip\noindent{\small \texttt{Qtsl.AGCompose.Series.}\texttt{\bfseries tower\_step}}\hfill{\footnotesize \path{lean/qtsl-aux/QtslAux/AGCompose/GshsSeries.lean}, line 81}\par\nopagebreak
\begin{leansrc}
theorem tower_step [IsFiniteMeasure μ] {mG : MeasurableSpace Ω} (hG : mG ≤ m0)
    {S : Set Ω} (hS : MeasurableSet[mG] S) {E : Set Ω} (hE : MeasurableSet[m0] E)
    {p : ℝ}
    (hcond : ∀ᵐ ω ∂μ, ω ∈ S → p ≤ (μ[E.indicator (fun _ => (1 : ℝ)) | mG]) ω) :
    ENNReal.ofReal p * μ S ≤ μ (S ∩ E)
\end{leansrc}
\par\medskip\noindent{\small \texttt{Qtsl.AGCompose.Parallel.}\texttt{\bfseries supermartingale\_weighted\_sum}}\hfill{\footnotesize \path{lean/qtsl-aux/QtslAux/AGCompose/GshsParallel.lean}, line 31}\par\nopagebreak
\begin{leansrc}
theorem supermartingale_weighted_sum {ι : Type*} (s : Finset ι) [IsFiniteMeasure μ]
    {Y : ι → ℕ → Ω → ℝ} {lam : ι → ℝ}
    (hY : ∀ i ∈ s, Supermartingale (Y i) 𝒢 μ) (hlam : ∀ i ∈ s, 0 ≤ lam i) :
    Supermartingale (fun n ω => ∑ i ∈ s, lam i * Y i n ω) 𝒢 μ
\end{leansrc}
\par\medskip\noindent{\small \texttt{Qtsl.AGCompose.SumFormFeedback.}\texttt{\bfseries debut\_avoid\_bound}}\hfill{\footnotesize \path{lean/qtsl-aux/QtslAux/GshsFeedbackSum.lean}, line 171}\par\nopagebreak
\begin{leansrc}
theorem debut_avoid_bound {S : Type*} {X : ℕ → Ω → S} {B : Set S} {a : ℝ}
    (hbarrier : ∀ k ω, X k ω ∈ B → a ≤ F.W k ω) {E : Set Ω}
    (hdebut : E ⊆ {ω | ∃ k, X k ω ∈ B}) :
    ENNReal.ofReal a * μ E ≤ ENNReal.ofReal (∫ ω, F.W 0 ω ∂μ)
\end{leansrc}
\par\medskip\noindent{\small \texttt{Qtsl.QtslCompose.}\texttt{\bfseries qualitative}}\hfill{\footnotesize \path{lean/qtsl-aux/QtslAux/QtslCompose/AGCompose.lean}, line 69}\par\nopagebreak
\begin{leansrc}
theorem qualitative [MeasurableSpace Ω] {μ : Measure Ω}
    (hw : ∀ j, conformExt ∩ (⋂ i ∈ S.upstream j, sat i) ⊆ conform j)
    (h : ∀ j, ∀ᵐ ω ∂μ, ω ∈ conform j → ω ∈ sat j)
    {satOut : Set Ω} (hout : conformExt ∩ (⋂ j, sat j) ⊆ satOut) :
    ∀ᵐ ω ∂μ, ω ∈ conformExt → ω ∈ satOut
\end{leansrc}
\par\medskip\noindent{\small \texttt{Qtsl.AGCompose.Series.}\texttt{\bfseries qualitative}}\hfill{\footnotesize \path{lean/qtsl-aux/QtslAux/AGCompose/GshsSeries.lean}, line 35}\par\nopagebreak
\begin{leansrc}
theorem qualitative
    (h₁ : ∀ᵐ ω ∂μ, ω ∈ conform₁ → ω ∈ sat₁)
    (hw : conform₁ ∩ sat₁ ⊆ conform₂)
    (h₂ : ∀ᵐ ω ∂μ, ω ∈ conform₂ → ω ∈ sat₂)
    (hout : conform₁ ∩ sat₁ ∩ sat₂ ⊆ satOut) :
    ∀ᵐ ω ∂μ, ω ∈ conform₁ → ω ∈ satOut
\end{leansrc}
\par\medskip\noindent{\small \texttt{Qtsl.AGCompose.Parallel.}\texttt{\bfseries supermartingale\_weighted\_sum}}\hfill{\footnotesize \path{lean/qtsl-aux/QtslAux/AGCompose/GshsParallel.lean}, line 31}\par\nopagebreak
\begin{leansrc}
theorem supermartingale_weighted_sum {ι : Type*} (s : Finset ι) [IsFiniteMeasure μ]
    {Y : ι → ℕ → Ω → ℝ} {lam : ι → ℝ}
    (hY : ∀ i ∈ s, Supermartingale (Y i) 𝒢 μ) (hlam : ∀ i ∈ s, 0 ≤ lam i) :
    Supermartingale (fun n ω => ∑ i ∈ s, lam i * Y i n ω) 𝒢 μ
\end{leansrc}
\par\medskip\noindent{\small \texttt{Qtsl.AGCompose.}\texttt{\bfseries SkeletonCertificate}}\hfill{\footnotesize \path{lean/qtsl-aux/QtslAux/GshsFeedbackSum.lean}, line 39}\par\nopagebreak
\begin{leansrc}
structure SkeletonCertificate (μ : Measure Ω) (𝒢 : Filtration ℕ m0) where …
\end{leansrc}
\par\medskip\noindent{\small \texttt{Qtsl.AGCompose.SkeletonCertificate.}\texttt{\bfseries step}}\hfill{\footnotesize \path{lean/qtsl-aux/QtslAux/GshsFeedbackSum.lean}, line 65}\par\nopagebreak
\begin{leansrc}
theorem step (C : SkeletonCertificate μ 𝒢) (k : ℕ) :
    μ[C.val (k + 1) + C.alphaInt k - C.sigmaInt k | 𝒢 k] ≤ᵐ[μ] C.val k
\end{leansrc}
\par\medskip\noindent{\small \texttt{Qtsl.AGCompose.SumFormFeedback.}\texttt{\bfseries drift}}\hfill{\footnotesize \path{lean/qtsl-aux/QtslAux/GshsFeedbackSum.lean}, line 124}\par\nopagebreak
\begin{leansrc}
theorem drift (k : ℕ) : μ[F.W (k + 1) + F.margin k | 𝒢 k] ≤ᵐ[μ] F.W k
\end{leansrc}
\par\medskip\noindent{\small \texttt{Qtsl.AGCompose.SumFormFeedback.}\texttt{\bfseries compensated\_supermartingale}}\hfill{\footnotesize \path{lean/qtsl-aux/QtslAux/GshsFeedbackSum.lean}, line 196}\par\nopagebreak
\begin{leansrc}
theorem compensated_supermartingale
    (hα₁ : ∀ k, StronglyMeasurable[𝒢 (k + 1)] (F.cert₁.alphaInt k))
    (hα₂ : ∀ k, StronglyMeasurable[𝒢 (k + 1)] (F.cert₂.alphaInt k)) :
    Supermartingale F.compensated 𝒢 μ
\end{leansrc}
\par\medskip\noindent{\small \texttt{Qtsl.AGCompose.SumFormFeedback.}\texttt{\bfseries margin\_budget}}\hfill{\footnotesize \path{lean/qtsl-aux/QtslAux/GshsFeedbackSum.lean}, line 226}\par\nopagebreak
\begin{leansrc}
theorem margin_budget
    (hα₁ : ∀ k, StronglyMeasurable[𝒢 (k + 1)] (F.cert₁.alphaInt k))
    (hα₂ : ∀ k, StronglyMeasurable[𝒢 (k + 1)] (F.cert₂.alphaInt k)) (n : ℕ) :
    ∫ ω, (∑ j ∈ Finset.range n, F.margin j) ω ∂μ ≤ ∫ ω, F.W 0 ω ∂μ
\end{leansrc}
\par\medskip\noindent{\small \texttt{Qtsl.QtslCompose.ChartShape.}\texttt{\bfseries profile}}\hfill{\footnotesize \path{lean/qtsl-aux/QtslAux/QtslCompose/Presentation.lean}, line 57}\par\nopagebreak
\begin{leansrc}
def ChartShape.profile (c : ChartShape) : Profile
\end{leansrc}
\par\medskip\noindent{\small \texttt{Qtsl.QtslCompose.}\texttt{\bfseries IsClassK0Inf}}\hfill{\footnotesize \path{lean/qtsl-aux/QtslAux/QtslCompose/WireCertDischarge.lean}, line 38}\par\nopagebreak
\begin{leansrc}
structure IsClassK0Inf (κ : ℝ → ℝ) : Prop where …
\end{leansrc}
\par\medskip\noindent{\small \texttt{Qtsl.QtslCompose.}\texttt{\bfseries polynomial\_discharge}}\hfill{\footnotesize \path{lean/qtsl-aux/QtslAux/QtslCompose/WireCertDischarge.lean}, line 75}\par\nopagebreak
\begin{leansrc}
theorem polynomial_discharge (p : Polynomial ℝ)
    (h0 : p.coeff 0 = 0)
    (hderiv : ∀ x : ℝ, 0 ≤ x → 0 ≤ p.derivative.eval x)
    (hdeg : 0 < p.degree) (hlead : 0 ≤ p.leadingCoeff) :
    IsClassK0Inf (fun x => p.eval x) where
  zero_at_zero
\end{leansrc}
\par\medskip\noindent{\small \texttt{Qtsl.QtslCompose.}\texttt{\bfseries slack\_triangle}}\hfill{\footnotesize \path{lean/qtsl-aux/QtslAux/QtslCompose/WireCertDischarge.lean}, line 94}\par\nopagebreak
\begin{leansrc}
theorem slack_triangle {X : Type*} {f g h : X → ℝ} {c₁ c₂ : ℝ}
    (h₁ : ∀ x, f x ≤ g x + c₁) (h₂ : ∀ x, g x ≤ h x + c₂) :
    ∀ x, f x ≤ h x + (c₁ + c₂)
\end{leansrc}
\par\medskip\noindent{\small \texttt{Qtsl.QtslCompose.}\texttt{\bfseries guarantee\_leg\_comp}}\hfill{\footnotesize \path{lean/qtsl-aux/QtslAux/QtslCompose/WireCertDischarge.lean}, line 100}\par\nopagebreak
\begin{leansrc}
theorem guarantee_leg_comp {A B C : Type*} {w₁ : A → B} {w₂ : B → C}
    {γ₁ : A → ℝ} {γ₂ : B → ℝ} {γ₃ : C → ℝ}
    (h₁ : ∀ a, γ₂ (w₁ a) ≤ γ₁ a) (h₂ : ∀ b, γ₃ (w₂ b) ≤ γ₂ b) :
    ∀ a, γ₃ (w₂ (w₁ a)) ≤ γ₁ a
\end{leansrc}
\par\medskip\noindent{\small \texttt{Qtsl.QtslCompose.IsClassK0Inf.}\texttt{\bfseries add}}\hfill{\footnotesize \path{lean/qtsl-aux/QtslAux/QtslCompose/WireCertDischarge.lean}, line 52}\par\nopagebreak
\begin{leansrc}
theorem add (h₁ : IsClassK0Inf κ₁) (h₂ : IsClassK0Inf κ₂) :
    IsClassK0Inf (fun x => κ₁ x + κ₂ x) where
  zero_at_zero
\end{leansrc}
\par\medskip\noindent{\small \texttt{Qtsl.QtslCompose.IsClassK0Inf.}\texttt{\bfseries comp}}\hfill{\footnotesize \path{lean/qtsl-aux/QtslAux/QtslCompose/WireCertDischarge.lean}, line 63}\par\nopagebreak
\begin{leansrc}
theorem comp (h₁ : IsClassK0Inf κ₁) (h₂ : IsClassK0Inf κ₂) :
    IsClassK0Inf (κ₁ ∘ κ₂) where
  zero_at_zero
\end{leansrc}
\par\medskip\noindent{\small \texttt{Qtsl.QtslCompose.IsClassK0Inf.}\texttt{\bfseries nonneg}}\hfill{\footnotesize \path{lean/qtsl-aux/QtslAux/QtslCompose/WireCertDischarge.lean}, line 47}\par\nopagebreak
\begin{leansrc}
theorem nonneg (h : IsClassK0Inf κ) {x : ℝ} (hx : 0 ≤ x) : 0 ≤ κ x
\end{leansrc}
\par\medskip\noindent{\small \texttt{BernsteinLogic.QtslCompose.Marks.}\texttt{\bfseries ConditionalDelivery}}\hfill{\footnotesize \path{lean/qtsl-aux/QtslAux/QtslCompose/AGComposeMarks.lean}, line 157}\par\nopagebreak
\begin{leansrc}
def ConditionalDelivery (N : MarkNet n V) (mark : Fin n → Set V) (i : Fin n) : Prop
\end{leansrc}
\par\medskip\noindent{\small \texttt{BernsteinLogic.QtslCompose.Marks.}\texttt{\bfseries DepLocal}}\hfill{\footnotesize \path{lean/qtsl-aux/QtslAux/QtslCompose/AGComposeMarks.lean}, line 109}\par\nopagebreak
\begin{leansrc}
def DepLocal (N : MarkNet n V) : Prop
\end{leansrc}
\par\medskip\noindent{\small \texttt{BernsteinLogic.QtslCompose.Marks.}\texttt{\bfseries MarkFlowFeasible}}\hfill{\footnotesize \path{lean/qtsl-aux/QtslAux/QtslCompose/AGComposeMarks.lean}, line 123}\par\nopagebreak
\begin{leansrc}
def MarkFlowFeasible (N : MarkNet n V) (t : Fin n → ℕ) : Prop
\end{leansrc}
\par\medskip\noindent{\small \texttt{BernsteinLogic.QtslCompose.Marks.}\texttt{\bfseries CarvedDelivery}}\hfill{\footnotesize \path{lean/qtsl-aux/QtslAux/QtslCompose/AGComposeCarves.lean}, line 69}\par\nopagebreak
\begin{leansrc}
def CarvedDelivery (N : MarkNet n V) (mark P : Fin n → Set V) : Prop
\end{leansrc}
\par\medskip\noindent{\small \texttt{BernsteinLogic.QtslCompose.Marks.}\texttt{\bfseries CarveClosure}}\hfill{\footnotesize \path{lean/qtsl-aux/QtslAux/QtslCompose/AGComposeCarves.lean}, line 76}\par\nopagebreak
\begin{leansrc}
def CarveClosure (N : MarkNet n V) (P : Fin n → Set V) (t : Fin n → ℕ) : Prop
\end{leansrc}
\par\medskip\noindent{\small \texttt{BernsteinLogic.QtslCompose.Marks.}\texttt{\bfseries PinnedDelivery}}\hfill{\footnotesize \path{lean/qtsl-aux/QtslAux/QtslCompose/AGComposeMarks.lean}, line 130}\par\nopagebreak
\begin{leansrc}
def PinnedDelivery (N : MarkNet n V) (mark : Fin n → Set V) (a : Fin n → V) : Prop
\end{leansrc}
\par\medskip\noindent{\small \texttt{BernsteinLogic.QtslCompose.Marks.}\texttt{\bfseries InstantClosure}}\hfill{\footnotesize \path{lean/qtsl-aux/QtslAux/QtslCompose/AGComposeMarks.lean}, line 148}\par\nopagebreak
\begin{leansrc}
def InstantClosure (N : MarkNet n V) (a : Fin n → V) (t : Fin n → ℕ) : Prop
\end{leansrc}
\par\medskip\noindent{\small \texttt{BernsteinLogic.QtslCompose.Marks.}\texttt{\bfseries carvedDelivery\_of\_pinned}}\hfill{\footnotesize \path{lean/qtsl-aux/QtslAux/QtslCompose/AGComposeCarves.lean}, line 276}\par\nopagebreak
\begin{leansrc}
theorem carvedDelivery_of_pinned (hpin : PinnedDelivery N mark a) :
    CarvedDelivery N mark (fun i => {a i})
\end{leansrc}
\par\medskip\noindent{\small \texttt{BernsteinLogic.QtslCompose.Marks.}\texttt{\bfseries carveClosure\_of\_pinned}}\hfill{\footnotesize \path{lean/qtsl-aux/QtslAux/QtslCompose/AGComposeCarves.lean}, line 290}\par\nopagebreak
\begin{leansrc}
theorem carveClosure_of_pinned (hfeas : MarkFlowFeasible N t)
    (hpin : PinnedDelivery N mark a) (hclose : InstantClosure N a t) :
    CarveClosure N (fun i => {a i}) t
\end{leansrc}
\par\medskip\noindent{\small \texttt{BernsteinLogic.QtslCompose.Marks.}\texttt{\bfseries marks\_of\_pinned\_delivery\_via\_carve}}\hfill{\footnotesize \path{lean/qtsl-aux/QtslAux/QtslCompose/AGComposeCarves.lean}, line 315}\par\nopagebreak
\begin{leansrc}
theorem marks_of_pinned_delivery_via_carve (hser : Serial N)
    (hfeas : MarkFlowFeasible N t) (hpin : PinnedDelivery N mark a)
    (hclose : InstantClosure N a t) :
    ∃ ω : ℕ → Fin n → V, JointRun N ω ∧ ∀ i, ω (t i) i ∈ mark i
\end{leansrc}
\par\medskip\noindent{\small \texttt{BernsteinLogic.QtslCompose.Marks.}\texttt{\bfseries carvedDelivery\_mark\_iff\_conditional}}\hfill{\footnotesize \path{lean/qtsl-aux/QtslAux/QtslCompose/AGComposeCarves.lean}, line 332}\par\nopagebreak
\begin{leansrc}
theorem carvedDelivery_mark_iff_conditional :
    CarvedDelivery N mark mark ↔ ∀ i, ConditionalDelivery N mark i
\end{leansrc}
\par\medskip\noindent{\small \texttt{BernsteinLogic.QtslCompose.Marks.}\texttt{\bfseries marks\_of\_carved\_delivery}}\hfill{\footnotesize \path{lean/qtsl-aux/QtslAux/QtslCompose/AGComposeCarves.lean}, line 260}\par\nopagebreak
\begin{leansrc}
theorem marks_of_carved_delivery (hser : Serial N) (hfeas : MarkFlowFeasible N t)
    (hpin : CarvedDelivery N mark P) (hclose : CarveClosure N P t) :
    ∃ ω : ℕ → Fin n → V, JointRun N ω ∧ ∀ i, ω (t i) i ∈ mark i
\end{leansrc}
\par\medskip\noindent{\small \texttt{BernsteinLogic.QtslCompose.Marks.}\texttt{\bfseries marks\_of\_conditional\_delivery\_acyclic}}\hfill{\footnotesize \path{lean/qtsl-aux/QtslAux/QtslCompose/AGComposeMarks.lean}, line 421}\par\nopagebreak
\begin{leansrc}
theorem marks_of_conditional_delivery_acyclic (hser : Serial N) (hloc : DepLocal N)
    (ord : Fin n → ℕ) (hord : ∀ i, ∀ j ∈ N.dep i, ord j < ord i)
    (hfeas : MarkFlowFeasible N t) (hdel : ∀ i, ConditionalDelivery N mark i) :
    ∃ ω : ℕ → Fin n → V, JointRun N ω ∧ ∀ i, ω (t i) i ∈ mark i
\end{leansrc}
\par\medskip\noindent{\small \texttt{BernsteinLogic.QtslCompose.Marks.}\texttt{\bfseries marks\_acyclic\_via\_maximal\_carve}}\hfill{\footnotesize \path{lean/qtsl-aux/QtslAux/QtslCompose/AGComposeCarves.lean}, line 475}\par\nopagebreak
\begin{leansrc}
theorem marks_acyclic_via_maximal_carve (hser : Serial N) (hloc : DepLocal N)
    (ord : Fin n → ℕ) (hord : ∀ i, ∀ j ∈ N.dep i, ord j < ord i)
    (hfeas : MarkFlowFeasible N t)
    (hdel : ∀ i, ConditionalDelivery N mark i) :
    ∃ ω : ℕ → Fin n → V, JointRun N ω ∧ ∀ i, ω (t i) i ∈ mark i
\end{leansrc}
\par\medskip\noindent{\small \texttt{BernsteinLogic.QtslCompose.Marks.}\texttt{\bfseries JointRun}}\hfill{\footnotesize \path{lean/qtsl-aux/QtslAux/QtslCompose/AGComposeMarks.lean}, line 104}\par\nopagebreak
\begin{leansrc}
def JointRun (N : MarkNet n V) (ω : ℕ → Fin n → V) : Prop
\end{leansrc}
\par\medskip\noindent{\small \texttt{BernsteinLogic.QtslCompose.Marks.}\texttt{\bfseries z1\_refutes\_walkless}}\hfill{\footnotesize \path{lean/qtsl-aux/QtslAux/QtslCompose/AGComposeMarks.lean}, line 505}\par\nopagebreak
\begin{leansrc}
theorem z1_refutes_walkless :
    ¬ (∀ (m : ℕ) (W : Type) (_ : Nonempty W) (N : MarkNet m W)
        (mark : Fin m → Set W) (a : Fin m → W) (t : Fin m → ℕ),
        Serial N → PinnedDelivery N mark a → InstantClosure N a t →
        ∃ ω : ℕ → Fin m → W, JointRun N ω ∧ ∀ i, ∃ k, ω k i ∈ mark i)
\end{leansrc}
\par\medskip\noindent{\small \texttt{BernsteinLogic.QtslCompose.Marks.}\texttt{\bfseries z1\_refutes\_walkless\_carve}}\hfill{\footnotesize \path{lean/qtsl-aux/QtslAux/QtslCompose/AGComposeCarves.lean}, line 590}\par\nopagebreak
\begin{leansrc}
theorem z1_refutes_walkless_carve :
    ¬ (∀ (m : ℕ) (W : Type) (_ : Nonempty W) (N : MarkNet m W)
        (mark P : Fin m → Set W) (t : Fin m → ℕ),
        Serial N → CarvedDelivery N mark P → CarveClosure N P t →
        ∃ ω : ℕ → Fin m → W, JointRun N ω ∧ ∀ i, ∃ k, ω k i ∈ mark i)
\end{leansrc}
\par\medskip\noindent{\small \texttt{BernsteinLogic.QtslCompose.Marks.}\texttt{\bfseries z5\_refutes\_pinless}}\hfill{\footnotesize \path{lean/qtsl-aux/QtslAux/QtslCompose/AGComposeMarks.lean}, line 597}\par\nopagebreak
\begin{leansrc}
theorem z5_refutes_pinless :
    ¬ (∀ (m : ℕ) (W : Type) (_ : Nonempty W) (N : MarkNet m W)
        (mark : Fin m → Set W) (t : Fin m → ℕ),
        Serial N → DepLocal N → (∀ i, ConditionalDelivery N mark i) →
        MarkFlowFeasible N t → W0Serial N → (∃ ω, JointRun N ω) →
        ∃ ω : ℕ → Fin m → W, JointRun N ω ∧ ∀ i, ∃ k, ω k i ∈ mark i)
\end{leansrc}
\par\medskip\noindent{\small \texttt{BernsteinLogic.QtslCompose.Marks.}\texttt{\bfseries z5\_refutes\_closureless}}\hfill{\footnotesize \path{lean/qtsl-aux/QtslAux/QtslCompose/AGComposeCarves.lean}, line 530}\par\nopagebreak
\begin{leansrc}
theorem z5_refutes_closureless :
    ¬ (∀ (m : ℕ) (W : Type) (_ : Nonempty W) (N : MarkNet m W)
        (mark P : Fin m → Set W) (t : Fin m → ℕ),
        Serial N → DepLocal N → CarvedDelivery N mark P →
        MarkFlowFeasible N t → W0Serial N → (∃ ω, JointRun N ω) →
        ∃ ω : ℕ → Fin m → W, JointRun N ω ∧ ∀ i, ∃ k, ω k i ∈ mark i)
\end{leansrc}
\par\medskip\noindent{\small \texttt{BernsteinLogic.QtslCompose.Marks.}\texttt{\bfseries z5\_carvedDelivery\_mark}}\hfill{\footnotesize \path{lean/qtsl-aux/QtslAux/QtslCompose/AGComposeCarves.lean}, line 495}\par\nopagebreak
\begin{leansrc}
theorem z5_carvedDelivery_mark : CarvedDelivery z5Net z5Mark z5Mark
\end{leansrc}
\par\medskip\noindent{\small \texttt{BernsteinLogic.QtslCompose.Marks.}\texttt{\bfseries z5\_no\_carve\_closure}}\hfill{\footnotesize \path{lean/qtsl-aux/QtslAux/QtslCompose/AGComposeCarves.lean}, line 509}\par\nopagebreak
\begin{leansrc}
theorem z5_no_carve_closure :
    ¬ ∃ P : Fin 2 → Set (Fin 3), (∀ i, P i ⊆ z5Mark i) ∧
      CarveClosure z5Net P ![0, 0]
\end{leansrc}
\par\medskip\noindent{\small \texttt{BernsteinLogic.QtslCompose.Marks.}\texttt{\bfseries z5\_no\_pin}}\hfill{\footnotesize \path{lean/qtsl-aux/QtslAux/QtslCompose/AGComposeMarks.lean}, line 663}\par\nopagebreak
\begin{leansrc}
theorem z5_no_pin : ¬ ∃ a : Fin 2 → Fin 3, PinnedDelivery z5Net z5Mark a
\end{leansrc}
\par\medskip\noindent{\small \texttt{BernsteinLogic.QtslCompose.Marks.}\texttt{\bfseries z1\_walk\_infeasible}}\hfill{\footnotesize \path{lean/qtsl-aux/QtslAux/QtslCompose/AGComposeMarks.lean}, line 460}\par\nopagebreak
\begin{leansrc}
theorem z1_walk_infeasible : ¬ ∃ t : Fin 2 → ℕ, MarkFlowFeasible z1Net t
\end{leansrc}
\par\medskip\noindent{\small \texttt{BernsteinLogic.QtslCompose.Marks.}\texttt{\bfseries allOnes\_walk\_feasible}}\hfill{\footnotesize \path{lean/qtsl-aux/QtslAux/QtslCompose/AGComposeMarks.lean}, line 469}\par\nopagebreak
\begin{leansrc}
theorem allOnes_walk_feasible : MarkFlowFeasible allOnesNet ![0, 0]
\end{leansrc}
\par\medskip\noindent{\small \texttt{BernsteinLogic.QtslCompose.Marks.}\texttt{\bfseries allOnes\_marks\_realized}}\hfill{\footnotesize \path{lean/qtsl-aux/QtslAux/QtslCompose/AGComposeMarks.lean}, line 677}\par\nopagebreak
\begin{leansrc}
theorem allOnes_marks_realized :
    ∃ ω : ℕ → Fin 2 → Bool, JointRun allOnesNet ω ∧
      ∀ i, ω ((![0, 0] : Fin 2 → ℕ) i) i ∈ boolMark i
\end{leansrc}
\par\medskip\noindent{\small \texttt{BernsteinLogic.QtslCompose.Marks.}\texttt{\bfseries allOnes\_direct}}\hfill{\footnotesize \path{lean/qtsl-aux/QtslAux/QtslCompose/AGComposeMarks.lean}, line 711}\par\nopagebreak
\begin{leansrc}
theorem allOnes_direct : JointRun allOnesNet fun _ => ![true, true]
\end{leansrc}
\par\medskip\noindent{\small \texttt{BernsteinLogic.QtslCompose.Marks.}\texttt{\bfseries fatCarve\_marks\_realized}}\hfill{\footnotesize \path{lean/qtsl-aux/QtslAux/QtslCompose/AGComposeCarves.lean}, line 666}\par\nopagebreak
\begin{leansrc}
theorem fatCarve_marks_realized :
    ∃ ω : ℕ → Fin 2 → Fin 3, JointRun relay3Net ω ∧
      ∀ i, ω ((![0, 0] : Fin 2 → ℕ) i) i ∈ relay3Mark i
\end{leansrc}
\par\medskip\noindent{\small \texttt{BernsteinLogic.QtslCompose.Marks.}\texttt{\bfseries marks\_of\_pinned\_delivery}}\hfill{\footnotesize \path{lean/qtsl-aux/QtslAux/QtslCompose/AGComposeMarks.lean}, line 287}\par\nopagebreak
\begin{leansrc}
theorem marks_of_pinned_delivery (hser : Serial N) (hfeas : MarkFlowFeasible N t)
    (hpin : PinnedDelivery N mark a) (hclose : InstantClosure N a t) :
    ∃ ω : ℕ → Fin n → V, JointRun N ω ∧ ∀ i, ω (t i) i ∈ mark i
\end{leansrc}
\par\medskip\noindent{\small \texttt{BernsteinLogic.QtslCompose.Marks.}\texttt{\bfseries pins\_of\_conditional\_delivery}}\hfill{\footnotesize \path{lean/qtsl-aux/QtslAux/QtslCompose/AGComposeMarks.lean}, line 308}\par\nopagebreak
\begin{leansrc}
theorem pins_of_conditional_delivery (hloc : DepLocal N) (ord : Fin n → ℕ)
    (hord : ∀ i, ∀ j ∈ N.dep i, ord j < ord i)
    (hdel : ∀ i, ConditionalDelivery N mark i) :
    ∃ a : Fin n → V, PinnedDelivery N mark a
\end{leansrc}
\par\medskip\noindent{\small \texttt{BernsteinLogic.QtslCompose.Marks.}\texttt{\bfseries instantClosure\_of\_acyclic}}\hfill{\footnotesize \path{lean/qtsl-aux/QtslAux/QtslCompose/AGComposeMarks.lean}, line 357}\par\nopagebreak
\begin{leansrc}
theorem instantClosure_of_acyclic (hser : Serial N) (hloc : DepLocal N)
    (ord : Fin n → ℕ) (hord : ∀ i, ∀ j ∈ N.dep i, ord j < ord i)
    (a : Fin n → V) (t : Fin n → ℕ) :
    InstantClosure N a t
\end{leansrc}
\par\medskip\noindent{\small \texttt{BernsteinLogic.QtslCompose.Marks.}\texttt{\bfseries carveClosure\_of\_acyclic}}\hfill{\footnotesize \path{lean/qtsl-aux/QtslAux/QtslCompose/AGComposeCarves.lean}, line 357}\par\nopagebreak
\begin{leansrc}
theorem carveClosure_of_acyclic (hser : Serial N) (hloc : DepLocal N)
    (ord : Fin n → ℕ) (hord : ∀ i, ∀ j ∈ N.dep i, ord j < ord i)
    (hfeas : MarkFlowFeasible N t) (hpin : CarvedDelivery N mark P) :
    CarveClosure N P t
\end{leansrc}
\par\medskip\noindent{\small \texttt{Qtsl.QtslCompose.}\texttt{\bfseries sim\_expect}}\hfill{\footnotesize \path{lean/qtsl-aux/QtslAux/QtslCompose/SimTransport.lean}, line 61}\par\nopagebreak
\begin{leansrc}
theorem sim_expect {σ : α → β} {K : α → α → ℝ} {Q : β → β → ℝ}
    (hσ : PushKernel σ K Q) (μ₀ : α → ℝ) (g : β → ℝ) (n : ℕ) :
    ∑ s, law K μ₀ n s * g (σ s) = ∑ b, law Q (push σ μ₀) n b * g b
\end{leansrc}
\par\medskip\noindent{\small \texttt{Qtsl.QtslCompose.}\texttt{\bfseries sim\_expect\_bound}}\hfill{\footnotesize \path{lean/qtsl-aux/QtslAux/QtslCompose/SimTransport.lean}, line 68}\par\nopagebreak
\begin{leansrc}
theorem sim_expect_bound {σ : α → β} {K : α → α → ℝ} {Q : β → β → ℝ}
    (hσ : PushKernel σ K Q) (μ₀ : α → ℝ) (g : β → ℝ) {n : ℕ} {b₀ : ℝ}
    (hb : ∑ b, law Q (push σ μ₀) n b * g b ≤ b₀) :
    ∑ s, law K μ₀ n s * g (σ s) ≤ b₀
\end{leansrc}
\par\medskip\noindent{\small \texttt{Qtsl.QtslCompose.}\texttt{\bfseries sim\_region}}\hfill{\footnotesize \path{lean/qtsl-aux/QtslAux/QtslCompose/SimTransport.lean}, line 77}\par\nopagebreak
\begin{leansrc}
theorem sim_region {σ : α → β} {K : α → α → ℝ} {Q : β → β → ℝ}
    (hσ : PushKernel σ K Q) {μ₀ : α → ℝ}
    (hK : ∀ s t, 0 ≤ K s t) (hμ₀ : ∀ s, 0 ≤ μ₀ s) {n : ℕ} {R' : Set β}
    (hR' : ∀ b, b ∉ R' → law Q (push σ μ₀) n b = 0) :
    ∀ s, σ s ∉ R' → law K μ₀ n s = 0
\end{leansrc}
\par\medskip\noindent{\small \texttt{Qtsl.QtslCompose.}\texttt{\bfseries expect\_push}}\hfill{\footnotesize \path{lean/qtsl-aux/QtslAux/QtslCompose/SimTransport.lean}, line 46}\par\nopagebreak
\begin{leansrc}
theorem expect_push (σ : α → β) (μ : α → ℝ) (g : β → ℝ) :
    ∑ b, push σ μ b * g b = ∑ s, μ s * g (σ s)
\end{leansrc}
\par\medskip\noindent{\small \texttt{Qtsl.QRelChain.}\texttt{\bfseries push\_law}}\hfill{\footnotesize \path{lean/qtsl-aux/QtslAux/QRelChain.lean}, line 152}\par\nopagebreak
\begin{leansrc}
theorem push_law [Fintype β] {π : α → β} {K : α → α → ℝ} {Q : β → β → ℝ}
    (h : PushKernel π K Q) (μ₀ : α → ℝ) :
    ∀ n, push π (law K μ₀ n) = law Q (push π μ₀) n
  | 0 => rfl
  | n + 1 => by rw [law_succ, push_step h, push_law h μ₀ n, law_succ]

end Push

section Drift

variable [Fintype α] {K : α → α → ℝ} {V μ₀ μ : α → ℝ}

theorem expect_step_le (hμ : ∀ s, 0 ≤ μ s)
    (hdrift : ∀ s, ∑ t, K s t * V t ≤ V s) :
    ∑ t, step K μ t * V t ≤ ∑ s, μ s * V s
\end{leansrc}
\par\medskip\noindent{\small \texttt{Qtsl.QtslCompose.Meas.}\texttt{\bfseries sim\_expect\_lintegral}}\hfill{\footnotesize \path{lean/qtsl-aux/QtslAux/QtslCompose/SimTransportMeas.lean}, line 91}\par\nopagebreak
\begin{leansrc}
theorem sim_expect_lintegral {σ : α → β} {K : Kernel α α} {Q : Kernel β β}
    (hσ : Measurable σ) (h : PushKernel σ K Q) (μ : Measure α)
    {g : β → ℝ≥0∞} (hg : Measurable g) (n : ℕ) :
    ∫⁻ s, g (σ s) ∂law K μ n = ∫⁻ b, g b ∂law Q (μ.map σ) n
\end{leansrc}
\par\medskip\noindent{\small \texttt{Qtsl.QtslCompose.Meas.}\texttt{\bfseries pushKernel\_iff\_comp}}\hfill{\footnotesize \path{lean/qtsl-aux/QtslAux/QtslCompose/SimTransportMeas.lean}, line 132}\par\nopagebreak
\begin{leansrc}
theorem pushKernel_iff_comp {σ : α → β} (hσ : Measurable σ)
    {K : Kernel α α} {Q : Kernel β β} :
    PushKernel σ K Q ↔
      Kernel.deterministic σ hσ ∘ₖ K = Q ∘ₖ Kernel.deterministic σ hσ
\end{leansrc}
\par\medskip\noindent{\small \texttt{Qtsl.QRelChain.}\texttt{\bfseries law\_nonneg}}\hfill{\footnotesize \path{lean/qtsl-aux/QtslAux/QRelChain.lean}, line 85}\par\nopagebreak
\begin{leansrc}
theorem law_nonneg (hK : ∀ s t, 0 ≤ K s t) (hμ₀ : ∀ s, 0 ≤ μ₀ s) :
    ∀ n s, 0 ≤ law K μ₀ n s
  | 0, s => hμ₀ s
  | n + 1, t => step_nonneg hK (law_nonneg hK hμ₀ n) t

theorem sum_step (hrow : ∀ s, ∑ t, K s t = 1) : ∑ t, step K μ t = ∑ s, μ s
\end{leansrc}
\par\medskip\noindent{\small \texttt{Qtsl.QRelChain.}\texttt{\bfseries PushKernel}}\hfill{\footnotesize \path{lean/qtsl-aux/QtslAux/QRelChain.lean}, line 130}\par\nopagebreak
\begin{leansrc}
def PushKernel (π : α → β) (K : α → α → ℝ) (Q : β → β → ℝ) : Prop
\end{leansrc}
\par\medskip\noindent{\small \texttt{Qtsl.QtslCompose.}\texttt{\bfseries law\_prodKernel}}\hfill{\footnotesize \path{lean/qtsl-aux/QtslAux/QtslCompose/Flatten.lean}, line 94}\par\nopagebreak
\begin{leansrc}
theorem law_prodKernel (K₁ : α → α → ℝ) (K₂ : β → β → ℝ)
    (μ₁ : α → ℝ) (μ₂ : β → ℝ) :
    ∀ n, law (prodKernel K₁ K₂) (prodDist μ₁ μ₂) n
      = prodDist (law K₁ μ₁ n) (law K₂ μ₂ n)
  | 0 => rfl
  | n + 1 => by
    rw [law_succ, law_prodKernel K₁ K₂ μ₁ μ₂ n, step_prodKernel, law_succ,
      law_succ]

section Projections

variable [DecidableEq α] [DecidableEq β]

theorem pushKernel_fst {K₁ : α → α → ℝ} {K₂ : β → β → ℝ}
    (h₂ : ∀ b, ∑ b', K₂ b b' = 1) :
    PushKernel Prod.fst (prodKernel K₁ K₂) K₁
\end{leansrc}
\par\medskip\noindent{\small \texttt{Qtsl.QtslCompose.}\texttt{\bfseries pushKernel\_fst}}\hfill{\footnotesize \path{lean/qtsl-aux/QtslAux/QtslCompose/Flatten.lean}, line 109}\par\nopagebreak
\begin{leansrc}
theorem pushKernel_fst {K₁ : α → α → ℝ} {K₂ : β → β → ℝ}
    (h₂ : ∀ b, ∑ b', K₂ b b' = 1) :
    PushKernel Prod.fst (prodKernel K₁ K₂) K₁
\end{leansrc}
\par\medskip\noindent{\small \texttt{Qtsl.QtslCompose.}\texttt{\bfseries prodKernel}}\hfill{\footnotesize \path{lean/qtsl-aux/QtslAux/QtslCompose/Flatten.lean}, line 34}\par\nopagebreak
\begin{leansrc}
def prodKernel (K₁ : α → α → ℝ) (K₂ : β → β → ℝ) :
    α × β → α × β → ℝ
\end{leansrc}
\par\medskip\noindent{\small \texttt{Qtsl.QtslCompose.Meas.}\texttt{\bfseries law\_parallelComp}}\hfill{\footnotesize \path{lean/qtsl-aux/QtslAux/QtslCompose/FlattenMeas.lean}, line 75}\par\nopagebreak
\begin{leansrc}
theorem law_parallelComp (μ : Measure α) (ν : Measure β)
    (K : Kernel α α) (Q : Kernel β β)
    [SFinite μ] [SFinite ν] [IsSFiniteKernel K] [IsSFiniteKernel Q] :
    ∀ n, law (K ∥ₖ Q) (μ.prod ν) n = (law K μ n).prod (law Q ν n)
  | 0 => rfl
  | n + 1 => by
    rw [law_succ, law_parallelComp μ ν K Q n, bind_parallelComp,
      ← law_succ, ← law_succ]

theorem pushKernel_fst {K : Kernel α α} {Q : Kernel β β}
    [IsSFiniteKernel K] [IsMarkovKernel Q] :
    PushKernel Prod.fst (K ∥ₖ Q) K
\end{leansrc}
\par\medskip\noindent{\small \texttt{Qtsl.QtslCompose.Meas.}\texttt{\bfseries PushKernel}}\hfill{\footnotesize \path{lean/qtsl-aux/QtslAux/QtslCompose/SimTransportMeas.lean}, line 43}\par\nopagebreak
\begin{leansrc}
def PushKernel (σ : α → β) (K : Kernel α α) (Q : Kernel β β) : Prop
\end{leansrc}
\par\medskip\noindent{\small \texttt{Qtsl.QtslCompose.Meas.}\texttt{\bfseries bind\_parallelComp}}\hfill{\footnotesize \path{lean/qtsl-aux/QtslAux/QtslCompose/FlattenMeas.lean}, line 58}\par\nopagebreak
\begin{leansrc}
theorem bind_parallelComp (μ : Measure α) (ν : Measure β)
    (K : Kernel α α) (Q : Kernel β β)
    [SFinite μ] [SFinite ν] [IsSFiniteKernel K] [IsSFiniteKernel Q] :
    (μ.prod ν).bind (K ∥ₖ Q) = (μ.bind K).prod (ν.bind Q)
\end{leansrc}
\par\medskip\noindent{\small \texttt{Qtsl.QtslCompose.}\texttt{\bfseries isMarkov\_prodKernel}}\hfill{\footnotesize \path{lean/qtsl-aux/QtslAux/QtslCompose/Flatten.lean}, line 45}\par\nopagebreak
\begin{leansrc}
theorem isMarkov_prodKernel {K₁ : α → α → ℝ} {K₂ : β → β → ℝ}
    (h₁ : IsMarkov K₁) (h₂ : IsMarkov K₂) :
    IsMarkov (prodKernel K₁ K₂) where
  nonneg
\end{leansrc}
\par\medskip\noindent{\small \texttt{Qtsl.QtslCompose.}\texttt{\bfseries isDist\_prodDist}}\hfill{\footnotesize \path{lean/qtsl-aux/QtslAux/QtslCompose/Flatten.lean}, line 61}\par\nopagebreak
\begin{leansrc}
theorem isDist_prodDist {μ₁ : α → ℝ} {μ₂ : β → ℝ}
    (h₁ : IsDist μ₁) (h₂ : IsDist μ₂) : IsDist (prodDist μ₁ μ₂) where
  nonneg
\end{leansrc}
\par\medskip\noindent{\small \texttt{Qtsl.QtslCompose.}\texttt{\bfseries push\_fst\_prodDist}}\hfill{\footnotesize \path{lean/qtsl-aux/QtslAux/QtslCompose/Flatten.lean}, line 135}\par\nopagebreak
\begin{leansrc}
theorem push_fst_prodDist {μ₁ : α → ℝ} {μ₂ : β → ℝ} (h₂ : IsDist μ₂) :
    push Prod.fst (prodDist μ₁ μ₂) = μ₁
\end{leansrc}
\par\medskip\noindent{\small \texttt{Qtsl.QtslCompose.Meas.}\texttt{\bfseries sfinite\_law}}\hfill{\footnotesize \path{lean/qtsl-aux/QtslAux/QtslCompose/FlattenMeas.lean}, line 38}\par\nopagebreak
\begin{leansrc}
instance sfinite_law (K : Kernel α α) [IsSFiniteKernel K]
    (μ : Measure α) [SFinite μ] : ∀ n, SFinite (law K μ n)
  | 0 => inferInstanceAs (SFinite μ)
  | n + 1 => by
    have
\end{leansrc}
\par\medskip\noindent{\small \texttt{Qtsl.QtslCompose.Meas.}\texttt{\bfseries isProbabilityMeasure\_law}}\hfill{\footnotesize \path{lean/qtsl-aux/QtslAux/QtslCompose/FlattenMeas.lean}, line 47}\par\nopagebreak
\begin{leansrc}
instance isProbabilityMeasure_law (K : Kernel α α) [IsMarkovKernel K]
    (μ : Measure α) [IsProbabilityMeasure μ] :
    ∀ n, IsProbabilityMeasure (law K μ n)
  | 0 => inferInstanceAs (IsProbabilityMeasure μ)
  | n + 1 => by
    have
\end{leansrc}
\par\medskip\noindent{\small \texttt{Qtsl.QtslCompose.Meas.}\texttt{\bfseries sim\_expect\_fst}}\hfill{\footnotesize \path{lean/qtsl-aux/QtslAux/QtslCompose/FlattenMeas.lean}, line 106}\par\nopagebreak
\begin{leansrc}
theorem sim_expect_fst {K : Kernel α α} {Q : Kernel β β}
    [IsSFiniteKernel K] [IsMarkovKernel Q]
    (μ : Measure α) (ν : Measure β) [IsProbabilityMeasure ν]
    {g : α → ℝ≥0∞} (hg : Measurable g) (n : ℕ) :
    ∫⁻ p, g p.1 ∂law (K ∥ₖ Q) (μ.prod ν) n = ∫⁻ a, g a ∂law K μ n
\end{leansrc}
\par\medskip\noindent{\small \texttt{Qtsl.QtslCompose.Meas.}\texttt{\bfseries pushKernel\_iff\_stochSim}}\hfill{\footnotesize \path{lean/qtsl-aux/QtslAux/QtslCompose/FlattenMeas.lean}, line 128}\par\nopagebreak
\begin{leansrc}
theorem pushKernel_iff_stochSim {σ : α → β} {K : Kernel α α}
    {Q : Kernel β β} : PushKernel σ K Q ↔ StochSim σ K Q
\end{leansrc}
\par\medskip\noindent{\small \texttt{Qtsl.QtslCompose.}\texttt{\bfseries prodDist}}\hfill{\footnotesize \path{lean/qtsl-aux/QtslAux/QtslCompose/Flatten.lean}, line 40}\par\nopagebreak
\begin{leansrc}
def prodDist (μ₁ : α → ℝ) (μ₂ : β → ℝ) : α × β → ℝ
\end{leansrc}
\par\medskip\noindent{\small \texttt{Qtsl.QtslCompose.}\texttt{\bfseries step\_prodKernel}}\hfill{\footnotesize \path{lean/qtsl-aux/QtslAux/QtslCompose/Flatten.lean}, line 75}\par\nopagebreak
\begin{leansrc}
theorem step_prodKernel (K₁ : α → α → ℝ) (K₂ : β → β → ℝ)
    (μ₁ : α → ℝ) (μ₂ : β → ℝ) :
    step (prodKernel K₁ K₂) (prodDist μ₁ μ₂)
      = prodDist (step K₁ μ₁) (step K₂ μ₂)
\end{leansrc}
\par\medskip\noindent{\small \texttt{Qtsl.QtslCompose.}\texttt{\bfseries pushKernel\_snd}}\hfill{\footnotesize \path{lean/qtsl-aux/QtslAux/QtslCompose/Flatten.lean}, line 122}\par\nopagebreak
\begin{leansrc}
theorem pushKernel_snd {K₁ : α → α → ℝ} {K₂ : β → β → ℝ}
    (h₁ : ∀ a, ∑ a', K₁ a a' = 1) :
    PushKernel Prod.snd (prodKernel K₁ K₂) K₂
\end{leansrc}
\par\medskip\noindent{\small \texttt{Qtsl.QDistTransport.}\texttt{\bfseries strict\_cell\_exact}}\hfill{\footnotesize \path{lean/qtsl-aux/QtslAux/QDistTransport.lean}, line 193}\par\nopagebreak
\begin{leansrc}
theorem strict_cell_exact [DecidableEq β] {π : α → β} {K : α → α → ℝ}
    {Q : β → β → ℝ} (h : PushKernel π K Q) (μ₀ : α → ℝ) (n : ℕ) :
    push π (law K μ₀ n) = law Q (push π μ₀) n
\end{leansrc}
\par\medskip\noindent{\small \texttt{Qtsl.QtslCompose.}\texttt{\bfseries aggReduce}}\hfill{\footnotesize \path{lean/qtsl-aux/QtslAux/QtslCompose/AggReduce.lean}, line 37}\par\nopagebreak
\begin{leansrc}
theorem aggReduce {X A : Type*} (agg : X → A) (scope : Set X) (K : Set A)
    (G : A → ℝ) (hG : ∀ a ∈ K, G a ≤ 0) (hin : ∀ x ∈ scope, agg x ∈ K) :
    ∀ x ∈ scope, G (agg x) ≤ 0
\end{leansrc}
\par\medskip\noindent{\small \texttt{Qtsl.QtslCompose.}\texttt{\bfseries coupling\_row}}\hfill{\footnotesize \path{lean/qtsl-aux/QtslAux/QtslCompose/AggReduce.lean}, line 80}\par\nopagebreak
\begin{leansrc}
theorem coupling_row {m : Multiset C} {f : C → ℝ} {l u : ℝ}
    (hl : ∀ c ∈ m, l ≤ f c) (hu : ∀ c ∈ m, f c ≤ u) :
    l * m.card ≤ query m f ∧ query m f ≤ u * m.card
\end{leansrc}
\par\medskip\noindent{\small \texttt{Qtsl.QtslCompose.}\texttt{\bfseries count\_row\_nonneg}}\hfill{\footnotesize \path{lean/qtsl-aux/QtslAux/QtslCompose/AggReduce.lean}, line 96}\par\nopagebreak
\begin{leansrc}
theorem count_row_nonneg (m : Multiset C) : (0 : ℝ) ≤ m.card
\end{leansrc}
\par\medskip\noindent{\small \texttt{Qtsl.QtslCompose.}\texttt{\bfseries count\_row\_le}}\hfill{\footnotesize \path{lean/qtsl-aux/QtslAux/QtslCompose/AggReduce.lean}, line 99}\par\nopagebreak
\begin{leansrc}
theorem count_row_le {m : Multiset C} {M : ℕ} (h : m.card ≤ M) :
    (m.card : ℝ) ≤ M
\end{leansrc}
\par\medskip\noindent{\small \texttt{Qtsl.QtslCompose.}\texttt{\bfseries moment\_row}}\hfill{\footnotesize \path{lean/qtsl-aux/QtslAux/QtslCompose/AggReduce.lean}, line 105}\par\nopagebreak
\begin{leansrc}
theorem moment_row (m : Multiset C) (f : C → ℝ) :
    (query m f) ^ 2 ≤ m.card * query m (fun c => f c ^ 2)
\end{leansrc}
\par\medskip\noindent{\small \texttt{Qtsl.QtslCompose.}\texttt{\bfseries query\_count}}\hfill{\footnotesize \path{lean/qtsl-aux/QtslAux/QtslCompose/AggReduce.lean}, line 66}\par\nopagebreak
\begin{leansrc}
theorem query_count [DecidableEq C] (m : Multiset C) (f : C → ℝ) :
    query m f = ∑ c ∈ m.toFinset, (m.count c : ℝ) * f c
\end{leansrc}
\par\medskip\noindent{\small \texttt{Qtsl.QtslCompose.}\texttt{\bfseries query\_const\_one}}\hfill{\footnotesize \path{lean/qtsl-aux/QtslAux/QtslCompose/AggReduce.lean}, line 55}\par\nopagebreak
\begin{leansrc}
theorem query_const_one (m : Multiset C) :
    query m (fun _ => (1 : ℝ)) = m.card
\end{leansrc}
\par\medskip\noindent{\small \texttt{Qtsl.Structural.}\texttt{\bfseries countAtomRounding\_sound}}\hfill{\footnotesize \path{lean/qtsl-aux/QtslAux/Structural.lean}, line 236}\par\nopagebreak
\begin{leansrc}
theorem countAtomRounding_sound (k : X → ℝ) (j : ℤ) (Reach : Set X)
    (hbridge : ∀ x ∈ Reach, ∃ n : ℤ, k x = (n : ℝ)) :
    ∀ x ∈ Reach, x ∈ {x | k x ≤ (j : ℝ)}ᶜ → (j : ℝ) + 1 ≤ k x
\end{leansrc}
\par\medskip\noindent{\small \texttt{Qtsl.Structural.}\texttt{\bfseries countAtomRounding\_tighter}}\hfill{\footnotesize \path{lean/qtsl-aux/QtslAux/Structural.lean}, line 250}\par\nopagebreak
\begin{leansrc}
theorem countAtomRounding_tighter (k : X → ℝ) (j : ℤ) :
    {x | (j : ℝ) + 1 ≤ k x} ⊆ {x | (j : ℝ) < k x}
\end{leansrc}
\par\medskip\noindent{\small \texttt{Qtsl.Structural.}\texttt{\bfseries countAtomRounding\_eq}}\hfill{\footnotesize \path{lean/qtsl-aux/QtslAux/Structural.lean}, line 259}\par\nopagebreak
\begin{leansrc}
theorem countAtomRounding_eq (k : X → ℝ) (j : ℤ) (Reach : Set X)
    (hbridge : ∀ x ∈ Reach, ∃ n : ℤ, k x = (n : ℝ)) :
    Reach ∩ {x | k x ≤ (j : ℝ)}ᶜ = Reach ∩ {x | (j : ℝ) + 1 ≤ k x}
\end{leansrc}
\par\medskip\noindent{\small \texttt{Qtsl.QRelChain.}\texttt{\bfseries sound\_QRelTV}}\hfill{\footnotesize \path{lean/qtsl-aux/QtslAux/QRelChain.lean}, line 736}\par\nopagebreak
\begin{leansrc}
theorem sound_QRelTV [Fintype L] [DecidableEq L] {f : σ₁ → L} {g : σ₂ → L}
    {V : σ₁ × σ₂ → ℝ} {b : ℝ}
    (hK : IsMarkov K) (hμ₀ : IsDist μ₀) (hK₁ : IsMarkov K₁) (hK₂ : IsMarkov K₂)
    (hm₁ : PushKernel Prod.fst K K₁) (hm₂ : PushKernel Prod.snd K K₂)
    (hfreeze : ∀ p ∈ A₁ ×ˢ A₂, K p p = 1)
    (hfr₁ : ∀ a ∈ A₁, K₁ a a = 1) (hfr₂ : ∀ a ∈ A₂, K₂ a a = 1)
    (habs₁ : Tendsto (fun n => ∑ a ∈ A₁, law K₁ (push Prod.fst μ₀) n a)
      atTop (𝓝 1))
    (habs₂ : Tendsto (fun n => ∑ a ∈ A₂, law K₂ (push Prod.snd μ₀) n a)
      atTop (𝓝 1))
    (hV0 : ∀ p, 0 ≤ V p) (hdrift : ∀ p, ∑ q, K p q * V q ≤ V p)
    (hterm : ∀ p ∈ A₁ ×ˢ A₂, f p.1 ≠ g p.2 → 1 ≤ V p)
    (hb : ∑ p, μ₀ p * V p ≤ b) :
    tvDist (push f (finalLaw K₁ (push Prod.fst μ₀) A₁))
      (push g (finalLaw K₂ (push Prod.snd μ₀) A₂)) ≤ b
\end{leansrc}
\par\medskip\noindent{\small \texttt{Qtsl.QRelChain.}\texttt{\bfseries sound\_QRelExpGap}}\hfill{\footnotesize \path{lean/qtsl-aux/QtslAux/QRelChain.lean}, line 777}\par\nopagebreak
\begin{leansrc}
theorem sound_QRelExpGap {v₁ : σ₁ → ℝ} {v₂ : σ₂ → ℝ} {V : σ₁ × σ₂ → ℝ} {b : ℝ}
    (hK : IsMarkov K) (hμ₀ : IsDist μ₀) (hK₁ : IsMarkov K₁) (hK₂ : IsMarkov K₂)
    (hm₁ : PushKernel Prod.fst K K₁) (hm₂ : PushKernel Prod.snd K K₂)
    (hfreeze : ∀ p ∈ A₁ ×ˢ A₂, K p p = 1)
    (hfr₁ : ∀ a ∈ A₁, K₁ a a = 1) (hfr₂ : ∀ a ∈ A₂, K₂ a a = 1)
    (habs₁ : Tendsto (fun n => ∑ a ∈ A₁, law K₁ (push Prod.fst μ₀) n a)
      atTop (𝓝 1))
    (habs₂ : Tendsto (fun n => ∑ a ∈ A₂, law K₂ (push Prod.snd μ₀) n a)
      atTop (𝓝 1))
    (hV0 : ∀ p, 0 ≤ V p) (hdrift : ∀ p, ∑ q, K p q * V q ≤ V p)
    (hterm : ∀ p ∈ A₁ ×ˢ A₂, |v₁ p.1 - v₂ p.2| ≤ V p)
    (hb : ∑ p, μ₀ p * V p ≤ b) :
    |(∑ a₁ ∈ A₁, finalMass K₁ (push Prod.fst μ₀) a₁ * v₁ a₁) -
      ∑ a₂ ∈ A₂, finalMass K₂ (push Prod.snd μ₀) a₂ * v₂ a₂| ≤ b
\end{leansrc}
\par\medskip\noindent{\small \texttt{Qtsl.QRelChain.}\texttt{\bfseries tendsto\_absorbedMass\_one}}\hfill{\footnotesize \path{lean/qtsl-aux/QtslAux/QRelChain.lean}, line 392}\par\nopagebreak
\begin{leansrc}
theorem tendsto_absorbedMass_one (hK : IsMarkov K) (hμ₀ : IsDist μ₀)
    (hfr : ∀ a ∈ A, K a a = 1)
    (hreach : ∀ s, ∃ k, 0 < ∑ a ∈ A, law K (pointMass s) k a) :
    Tendsto (fun n => ∑ a ∈ A, law K μ₀ n a) atTop (𝓝 1)
\end{leansrc}
\par\medskip\noindent{\small \texttt{Qtsl.QRelChain.}\texttt{\bfseries tvDist}}\hfill{\footnotesize \path{lean/qtsl-aux/QtslAux/QRelChain.lean}, line 453}\par\nopagebreak
\begin{leansrc}
def tvDist [Fintype L] (p q : L → ℝ) : ℝ
\end{leansrc}
\par\medskip\noindent{\small \texttt{Qtsl.QRelChain.}\texttt{\bfseries push\_step}}\hfill{\footnotesize \path{lean/qtsl-aux/QtslAux/QRelChain.lean}, line 133}\par\nopagebreak
\begin{leansrc}
theorem push_step [Fintype β] {π : α → β} {K : α → α → ℝ} {Q : β → β → ℝ}
    (h : PushKernel π K Q) (μ : α → ℝ) :
    push π (step K μ) = step Q (push π μ)
\end{leansrc}
\par\medskip\noindent{\small \texttt{Qtsl.QRelChain.}\texttt{\bfseries expect\_step\_le}}\hfill{\footnotesize \path{lean/qtsl-aux/QtslAux/QRelChain.lean}, line 164}\par\nopagebreak
\begin{leansrc}
theorem expect_step_le (hμ : ∀ s, 0 ≤ μ s)
    (hdrift : ∀ s, ∑ t, K s t * V t ≤ V s) :
    ∑ t, step K μ t * V t ≤ ∑ s, μ s * V s
\end{leansrc}
\par\medskip\noindent{\small \texttt{Qtsl.QRelChain.}\texttt{\bfseries expect\_law\_le}}\hfill{\footnotesize \path{lean/qtsl-aux/QtslAux/QRelChain.lean}, line 175}\par\nopagebreak
\begin{leansrc}
theorem expect_law_le (hK : ∀ s t, 0 ≤ K s t) (hμ₀ : ∀ s, 0 ≤ μ₀ s)
    (hdrift : ∀ s, ∑ t, K s t * V t ≤ V s) :
    ∀ n, ∑ s, law K μ₀ n s * V s ≤ ∑ s, μ₀ s * V s
  | 0 => le_rfl
  | n + 1 =>
    (expect_step_le (law_nonneg hK hμ₀ n) hdrift).trans
      (expect_law_le hK hμ₀ hdrift n)

end Drift

section Absorb

variable [Fintype α] {K : α → α → ℝ} {μ₀ : α → ℝ}

theorem law_mono_absorbing (hK : ∀ s t, 0 ≤ K s t) (hμ₀ : ∀ s, 0 ≤ μ₀ s)
    {a : α} (haa : K a a = 1) :
    Monotone fun n => law K μ₀ n a
\end{leansrc}
\par\medskip\noindent{\small \texttt{Qtsl.QRelChain.}\texttt{\bfseries law\_mono\_absorbing}}\hfill{\footnotesize \path{lean/qtsl-aux/QtslAux/QRelChain.lean}, line 189}\par\nopagebreak
\begin{leansrc}
theorem law_mono_absorbing (hK : ∀ s t, 0 ≤ K s t) (hμ₀ : ∀ s, 0 ≤ μ₀ s)
    {a : α} (haa : K a a = 1) :
    Monotone fun n => law K μ₀ n a
\end{leansrc}
\par\medskip\noindent{\small \texttt{Qtsl.QRelChain.}\texttt{\bfseries finalMass}}\hfill{\footnotesize \path{lean/qtsl-aux/QtslAux/QRelChain.lean}, line 198}\par\nopagebreak
\begin{leansrc}
def finalMass (K : α → α → ℝ) (μ₀ : α → ℝ) (a : α) : ℝ
\end{leansrc}
\par\medskip\noindent{\small \texttt{Qtsl.QRelChain.}\texttt{\bfseries finalExpect\_le}}\hfill{\footnotesize \path{lean/qtsl-aux/QtslAux/QRelChain.lean}, line 226}\par\nopagebreak
\begin{leansrc}
theorem finalExpect_le {V gv : α → ℝ} {A : Finset α} {b : ℝ}
    (hK : IsMarkov K) (hμ₀ : IsDist μ₀) (hfr : ∀ a ∈ A, K a a = 1)
    (hgV : ∀ a ∈ A, gv a ≤ V a) (hV0 : ∀ s, 0 ≤ V s)
    (hdrift : ∀ s, ∑ t, K s t * V t ≤ V s) (hb : ∑ s, μ₀ s * V s ≤ b) :
    ∑ a ∈ A, finalMass K μ₀ a * gv a ≤ b
\end{leansrc}
\par\medskip\noindent{\small \texttt{Qtsl.QRelChain.}\texttt{\bfseries tvDist}}\hfill{\footnotesize \path{lean/qtsl-aux/QtslAux/QRelChain.lean}, line 453}\par\nopagebreak
\begin{leansrc}
def tvDist [Fintype L] (p q : L → ℝ) : ℝ
\end{leansrc}
\par\medskip\noindent{\small \texttt{Qtsl.QRelChain.}\texttt{\bfseries tvDist\_push\_le\_disagree}}\hfill{\footnotesize \path{lean/qtsl-aux/QtslAux/QRelChain.lean}, line 455}\par\nopagebreak
\begin{leansrc}
theorem tvDist_push_le_disagree {σ₁ σ₂ : Type*} [Fintype σ₁] [Fintype σ₂]
    [Fintype L] [DecidableEq L] (γ : σ₁ × σ₂ → ℝ) (hγ : ∀ p, 0 ≤ γ p)
    (f : σ₁ → L) (g : σ₂ → L) :
    tvDist (push (fun p : σ₁ × σ₂ => f p.1) γ) (push (fun p : σ₁ × σ₂ => g p.2) γ) ≤
      ∑ p ∈ univ.filter fun p : σ₁ × σ₂ => f p.1 ≠ g p.2, γ p
\end{leansrc}
\par\medskip\noindent{\small \texttt{Qtsl.QRelChain.}\texttt{\bfseries sum\_finalMass\_snd\_eq}}\hfill{\footnotesize \path{lean/qtsl-aux/QtslAux/QRelChain.lean}, line 532}\par\nopagebreak
\begin{leansrc}
theorem sum_finalMass_snd_eq (hK : IsMarkov K) (hμ₀ : IsDist μ₀)
    (hK₁ : IsMarkov K₁) (hK₂ : IsMarkov K₂)
    (hm₁ : PushKernel Prod.fst K K₁) (hm₂ : PushKernel Prod.snd K K₂)
    (hfreeze : ∀ p ∈ A₁ ×ˢ A₂, K p p = 1) (hfr₁ : ∀ a ∈ A₁, K₁ a a = 1)
    (habs₂ : Tendsto (fun n => ∑ a ∈ A₂, law K₂ (push Prod.snd μ₀) n a)
      atTop (𝓝 1))
    {a₁ : σ₁} (ha₁ : a₁ ∈ A₁) :
    ∑ a₂ ∈ A₂, finalMass K μ₀ (a₁, a₂) = finalMass K₁ (push Prod.fst μ₀) a₁
\end{leansrc}
\par\medskip\noindent{\small \texttt{Qtsl.QRelChain.}\texttt{\bfseries sum\_finalMass\_fst\_eq}}\hfill{\footnotesize \path{lean/qtsl-aux/QtslAux/QRelChain.lean}, line 585}\par\nopagebreak
\begin{leansrc}
theorem sum_finalMass_fst_eq (hK : IsMarkov K) (hμ₀ : IsDist μ₀)
    (hK₁ : IsMarkov K₁) (hK₂ : IsMarkov K₂)
    (hm₁ : PushKernel Prod.fst K K₁) (hm₂ : PushKernel Prod.snd K K₂)
    (hfreeze : ∀ p ∈ A₁ ×ˢ A₂, K p p = 1) (hfr₂ : ∀ a ∈ A₂, K₂ a a = 1)
    (habs₁ : Tendsto (fun n => ∑ a ∈ A₁, law K₁ (push Prod.fst μ₀) n a)
      atTop (𝓝 1))
    {a₂ : σ₂} (ha₂ : a₂ ∈ A₂) :
    ∑ a₁ ∈ A₁, finalMass K μ₀ (a₁, a₂) = finalMass K₂ (push Prod.snd μ₀) a₂
\end{leansrc}
\par\medskip\noindent{\small \texttt{Qtsl.QRelChain.}\texttt{\bfseries sum\_finalMass\_prod\_eq\_one}}\hfill{\footnotesize \path{lean/qtsl-aux/QtslAux/QRelChain.lean}, line 638}\par\nopagebreak
\begin{leansrc}
theorem sum_finalMass_prod_eq_one (hK : IsMarkov K) (hμ₀ : IsDist μ₀)
    (hK₁ : IsMarkov K₁) (hK₂ : IsMarkov K₂)
    (hm₁ : PushKernel Prod.fst K K₁) (hm₂ : PushKernel Prod.snd K K₂)
    (hfreeze : ∀ p ∈ A₁ ×ˢ A₂, K p p = 1) (hfr₁ : ∀ a ∈ A₁, K₁ a a = 1)
    (habs₁ : Tendsto (fun n => ∑ a ∈ A₁, law K₁ (push Prod.fst μ₀) n a)
      atTop (𝓝 1))
    (habs₂ : Tendsto (fun n => ∑ a ∈ A₂, law K₂ (push Prod.snd μ₀) n a)
      atTop (𝓝 1)) :
    ∑ p ∈ A₁ ×ˢ A₂, finalMass K μ₀ p = 1
\end{leansrc}
\par\medskip\noindent{\small \texttt{Qtsl.QRelChain.}\texttt{\bfseries push\_finalLaw\_fst}}\hfill{\footnotesize \path{lean/qtsl-aux/QtslAux/QRelChain.lean}, line 667}\par\nopagebreak
\begin{leansrc}
theorem push_finalLaw_fst [Fintype L] [DecidableEq L]
    (hK : IsMarkov K) (hμ₀ : IsDist μ₀) (hK₁ : IsMarkov K₁) (hK₂ : IsMarkov K₂)
    (hm₁ : PushKernel Prod.fst K K₁) (hm₂ : PushKernel Prod.snd K K₂)
    (hfreeze : ∀ p ∈ A₁ ×ˢ A₂, K p p = 1) (hfr₁ : ∀ a ∈ A₁, K₁ a a = 1)
    (habs₂ : Tendsto (fun n => ∑ a ∈ A₂, law K₂ (push Prod.snd μ₀) n a)
      atTop (𝓝 1))
    (f : σ₁ → L) :
    push (fun p : σ₁ × σ₂ => f p.1) (finalLaw K μ₀ (A₁ ×ˢ A₂)) =
      push f (finalLaw K₁ (push Prod.fst μ₀) A₁)
\end{leansrc}
\par\medskip\noindent{\small \texttt{Qtsl.QRelChain.}\texttt{\bfseries sound\_QRelTV\_of\_reach}}\hfill{\footnotesize \path{lean/qtsl-aux/QtslAux/QRelChain.lean}, line 815}\par\nopagebreak
\begin{leansrc}
theorem sound_QRelTV_of_reach [Fintype L] [DecidableEq L] {f : σ₁ → L}
    {g : σ₂ → L} {V : σ₁ × σ₂ → ℝ} {b : ℝ}
    (hK : IsMarkov K) (hμ₀ : IsDist μ₀) (hK₁ : IsMarkov K₁) (hK₂ : IsMarkov K₂)
    (hm₁ : PushKernel Prod.fst K K₁) (hm₂ : PushKernel Prod.snd K K₂)
    (hfreeze : ∀ p ∈ A₁ ×ˢ A₂, K p p = 1)
    (hfr₁ : ∀ a ∈ A₁, K₁ a a = 1) (hfr₂ : ∀ a ∈ A₂, K₂ a a = 1)
    (hreach₁ : ∀ s, ∃ k, 0 < ∑ a ∈ A₁, law K₁ (pointMass s) k a)
    (hreach₂ : ∀ s, ∃ k, 0 < ∑ a ∈ A₂, law K₂ (pointMass s) k a)
    (hV0 : ∀ p, 0 ≤ V p) (hdrift : ∀ p, ∑ q, K p q * V q ≤ V p)
    (hterm : ∀ p ∈ A₁ ×ˢ A₂, f p.1 ≠ g p.2 → 1 ≤ V p)
    (hb : ∑ p, μ₀ p * V p ≤ b) :
    tvDist (push f (finalLaw K₁ (push Prod.fst μ₀) A₁))
      (push g (finalLaw K₂ (push Prod.snd μ₀) A₂)) ≤ b
\end{leansrc}
\par\medskip\noindent{\small \texttt{Qtsl.QRelChain.}\texttt{\bfseries sound\_QRelExpGap\_of\_reach}}\hfill{\footnotesize \path{lean/qtsl-aux/QtslAux/QRelChain.lean}, line 833}\par\nopagebreak
\begin{leansrc}
theorem sound_QRelExpGap_of_reach {v₁ : σ₁ → ℝ} {v₂ : σ₂ → ℝ}
    {V : σ₁ × σ₂ → ℝ} {b : ℝ}
    (hK : IsMarkov K) (hμ₀ : IsDist μ₀) (hK₁ : IsMarkov K₁) (hK₂ : IsMarkov K₂)
    (hm₁ : PushKernel Prod.fst K K₁) (hm₂ : PushKernel Prod.snd K K₂)
    (hfreeze : ∀ p ∈ A₁ ×ˢ A₂, K p p = 1)
    (hfr₁ : ∀ a ∈ A₁, K₁ a a = 1) (hfr₂ : ∀ a ∈ A₂, K₂ a a = 1)
    (hreach₁ : ∀ s, ∃ k, 0 < ∑ a ∈ A₁, law K₁ (pointMass s) k a)
    (hreach₂ : ∀ s, ∃ k, 0 < ∑ a ∈ A₂, law K₂ (pointMass s) k a)
    (hV0 : ∀ p, 0 ≤ V p) (hdrift : ∀ p, ∑ q, K p q * V q ≤ V p)
    (hterm : ∀ p ∈ A₁ ×ˢ A₂, |v₁ p.1 - v₂ p.2| ≤ V p)
    (hb : ∑ p, μ₀ p * V p ≤ b) :
    |(∑ a₁ ∈ A₁, finalMass K₁ (push Prod.fst μ₀) a₁ * v₁ a₁) -
      ∑ a₂ ∈ A₂, finalMass K₂ (push Prod.snd μ₀) a₂ * v₂ a₂| ≤ b
\end{leansrc}
\par\medskip\noindent{\small \texttt{Qtsl.QRelChain.}\texttt{\bfseries exists\_uniform\_reach}}\hfill{\footnotesize \path{lean/qtsl-aux/QtslAux/QRelChain.lean}, line 334}\par\nopagebreak
\begin{leansrc}
theorem exists_uniform_reach (hK : IsMarkov K) (hfr : ∀ a ∈ A, K a a = 1)
    (hreach : ∀ s, ∃ k, 0 < ∑ a ∈ A, law K (pointMass s) k a) :
    ∃ N ε, 0 < ε ∧ 1 ≤ N ∧ ∀ s, ε ≤ ∑ a ∈ A, law K (pointMass s) N a
\end{leansrc}
\par\medskip\noindent{\small \texttt{Qtsl.QRelChain.}\texttt{\bfseries absorbed\_decay}}\hfill{\footnotesize \path{lean/qtsl-aux/QtslAux/QRelChain.lean}, line 354}\par\nopagebreak
\begin{leansrc}
theorem absorbed_decay {μ : α → ℝ} (hK : IsMarkov K) (hμ : IsDist μ)
    (hfr : ∀ a ∈ A, K a a = 1) {N : ℕ} {ε : ℝ}
    (hε : ∀ s, ε ≤ ∑ a ∈ A, law K (pointMass s) N a) :
    1 - ∑ a ∈ A, law K μ N a ≤
      (1 - ε) * (1 - ∑ s ∈ A, μ s)
\end{leansrc}
\par\medskip\noindent{\small \texttt{Qtsl.QFixBall.}\texttt{\bfseries sound\_QFixBall}}\hfill{\footnotesize \path{lean/qtsl-aux/QtslAux/QFixBall.lean}, line 88}\par\nopagebreak
\begin{leansrc}
theorem sound_QFixBall [CompleteSpace α] (hK : K < 1)
    (hlip : LipschitzOnWith K f (closedBall x₀ r))
    (hres : dist (f x₀) x₀ ≤ δ) (hrad : δ + K * r ≤ r) :
    ∃! y, y ∈ closedBall x₀ r ∧ IsFixedPt f y
\end{leansrc}
\par\medskip\noindent{\small \texttt{Qtsl.QFixBall.}\texttt{\bfseries dist\_fixedPoint\_le}}\hfill{\footnotesize \path{lean/qtsl-aux/QtslAux/QFixBall.lean}, line 59}\par\nopagebreak
\begin{leansrc}
theorem dist_fixedPoint_le (hK : K < 1)
    (hlip : LipschitzOnWith K f (closedBall x₀ r))
    (hres : dist (f x₀) x₀ ≤ δ) (hrad : δ + K * r ≤ r)
    {y : α} (hy : y ∈ closedBall x₀ r) (hfy : IsFixedPt f y) :
    dist x₀ y ≤ δ / (1 - K)
\end{leansrc}
\par\medskip\noindent{\small \texttt{Qtsl.QFixBall.}\texttt{\bfseries anchored\_le\_radius}}\hfill{\footnotesize \path{lean/qtsl-aux/QtslAux/QFixBall.lean}, line 103}\par\nopagebreak
\begin{leansrc}
theorem anchored_le_radius (hK : K < 1)
    (hrad : δ + K * r ≤ r) : δ / (1 - K) ≤ r
\end{leansrc}
\par\medskip\noindent{\small \texttt{Qtsl.QFixBall.}\texttt{\bfseries sound\_QFixBall\_affine}}\hfill{\footnotesize \path{lean/qtsl-aux/QtslAux/QFixBall.lean}, line 120}\par\nopagebreak
\begin{leansrc}
theorem sound_QFixBall_affine {M : ℕ}
    (A : EuclideanSpace ℝ (Fin M) →L[ℝ] EuclideanSpace ℝ (Fin M))
    (c x₀ : EuclideanSpace ℝ (Fin M)) {K : ℝ≥0} {δ r : ℝ}
    (hK : K < 1) (hA : ‖A‖₊ ≤ K)
    (hres : dist (A x₀ + c) x₀ ≤ δ) (hrad : δ + K * r ≤ r) :
    ∃! y, y ∈ closedBall x₀ r ∧ IsFixedPt (fun x => A x + c) y
\end{leansrc}
\par\medskip\noindent{\small \texttt{Qtsl.Wasserstein.}\texttt{\bfseries sound\_QFixBall\_wasserstein}}\hfill{\footnotesize \path{lean/qtsl-aux/QtslAux/Wasserstein/Complete.lean}, line 449}\par\nopagebreak
\begin{leansrc}
theorem sound_QFixBall_wasserstein {T : P2 X → P2 X} {μ₀ : P2 X} {K : ℝ≥0}
    {δ r : ℝ} (hK : K < 1)
    (hlip : LipschitzOnWith K T (closedBall μ₀ r))
    (hres : dist (T μ₀) μ₀ ≤ δ) (hrad : δ + K * r ≤ r) :
    ∃! ν, ν ∈ closedBall μ₀ r ∧ IsFixedPt T ν
\end{leansrc}
\par\medskip\noindent{\small \texttt{Qtsl.QFixBall.}\texttt{\bfseries delta\_nonneg}}\hfill{\footnotesize \path{lean/qtsl-aux/QtslAux/QFixBall.lean}, line 34}\par\nopagebreak
\begin{leansrc}
theorem delta_nonneg (hres : dist (f x₀) x₀ ≤ δ) : 0 ≤ δ
\end{leansrc}
\par\medskip\noindent{\small \texttt{Qtsl.QFixBall.}\texttt{\bfseries radius\_nonneg}}\hfill{\footnotesize \path{lean/qtsl-aux/QtslAux/QFixBall.lean}, line 37}\par\nopagebreak
\begin{leansrc}
theorem radius_nonneg (hK : K < 1) (hres : dist (f x₀) x₀ ≤ δ)
    (hrad : δ + K * r ≤ r) : 0 ≤ r
\end{leansrc}
\par\medskip\noindent{\small \texttt{Qtsl.QFixBall.}\texttt{\bfseries mapsTo\_closedBall}}\hfill{\footnotesize \path{lean/qtsl-aux/QtslAux/QFixBall.lean}, line 45}\par\nopagebreak
\begin{leansrc}
theorem mapsTo_closedBall (hK : K < 1)
    (hlip : LipschitzOnWith K f (closedBall x₀ r))
    (hres : dist (f x₀) x₀ ≤ δ) (hrad : δ + K * r ≤ r) :
    Set.MapsTo f (closedBall x₀ r) (closedBall x₀ r)
\end{leansrc}
\par\medskip\noindent{\small \texttt{Qtsl.QFixBall.}\texttt{\bfseries fixedPoint\_unique}}\hfill{\footnotesize \path{lean/qtsl-aux/QtslAux/QFixBall.lean}, line 76}\par\nopagebreak
\begin{leansrc}
theorem fixedPoint_unique (hK : K < 1)
    (hlip : LipschitzOnWith K f (closedBall x₀ r))
    {y z : α} (hy : y ∈ closedBall x₀ r) (hz : z ∈ closedBall x₀ r)
    (hfy : IsFixedPt f y) (hfz : IsFixedPt f z) : y = z
\end{leansrc}
\par\medskip\noindent{\small \texttt{Qtsl.QFixBall.}\texttt{\bfseries lipschitzWith\_affine}}\hfill{\footnotesize \path{lean/qtsl-aux/QtslAux/QFixBall.lean}, line 113}\par\nopagebreak
\begin{leansrc}
theorem lipschitzWith_affine {E : Type*} [NormedAddCommGroup E]
    [NormedSpace ℝ E] (A : E →L[ℝ] E) (c : E) (hA : ‖A‖₊ ≤ K) :
    LipschitzWith K fun x => A x + c
\end{leansrc}
\par\medskip\noindent{\small \texttt{Qtsl.QDistTransport.}\texttt{\bfseries pushKernel\_iff\_eps\_zero}}\hfill{\footnotesize \path{lean/qtsl-aux/QtslAux/QDistTransport.lean}, line 183}\par\nopagebreak
\begin{leansrc}
theorem pushKernel_iff_eps_zero [DecidableEq β] {π : α → β}
    {K : α → α → ℝ} {Q : β → β → ℝ} :
    PushKernel π K Q ↔ EpsPushKernel π K Q 0
\end{leansrc}
\par\medskip\noindent{\small \texttt{Qtsl.QDistTransport.}\texttt{\bfseries epsPushKernel\_comp}}\hfill{\footnotesize \path{lean/qtsl-aux/QtslAux/QDistTransport.lean}, line 168}\par\nopagebreak
\begin{leansrc}
theorem epsPushKernel_comp [DecidableEq β] [DecidableEq γ]
    {π₁ : α → β} {π₂ : β → γ} {K₁ : α → α → ℝ} {K₂ : β → β → ℝ}
    {K₃ : γ → γ → ℝ} {ε₁ ε₂ : ℝ}
    (h₁ : EpsPushKernel π₁ K₁ K₂ ε₁) (h₂ : EpsPushKernel π₂ K₂ K₃ ε₂) :
    EpsPushKernel (π₂ ∘ π₁) K₁ K₃ (ε₁ + ε₂)
\end{leansrc}
\par\medskip\noindent{\small \texttt{Qtsl.QDistTransport.}\texttt{\bfseries sound\_QDistSelfTransport}}\hfill{\footnotesize \path{lean/qtsl-aux/QtslAux/QDistTransport.lean}, line 74}\par\nopagebreak
\begin{leansrc}
theorem sound_QDistSelfTransport [DecidableEq β] (π : α → β)
    {μ ν : α → ℝ} {ε : ℝ} (h : tvDist μ ν ≤ ε) :
    tvDist (push π μ) (push π ν) ≤ ε
\end{leansrc}
\par\medskip\noindent{\small \texttt{Qtsl.QRelChain.}\texttt{\bfseries IsMarkov}}\hfill{\footnotesize \path{lean/qtsl-aux/QtslAux/QRelChain.lean}, line 55}\par\nopagebreak
\begin{leansrc}
structure IsMarkov [Fintype α] (K : α → α → ℝ) : Prop where …
\end{leansrc}
\par\medskip\noindent{\small \texttt{Qtsl.QDistTransport.}\texttt{\bfseries sound\_QDistTransport\_eps}}\hfill{\footnotesize \path{lean/qtsl-aux/QtslAux/QDistTransport.lean}, line 152}\par\nopagebreak
\begin{leansrc}
theorem sound_QDistTransport_eps [DecidableEq β] {π : α → β}
    {K : α → α → ℝ} {Q : β → β → ℝ} {ε : ℝ} {μ₀ : α → ℝ}
    (hK : IsMarkov K) (hQ : IsMarkov Q) (hμ₀ : IsDist μ₀)
    (h : EpsPushKernel π K Q ε) (n : ℕ) :
    tvDist (push π (law K μ₀ n)) (law Q (push π μ₀) n) ≤ n * ε
\end{leansrc}
\par\medskip\noindent{\small \texttt{Qtsl.QDistTransport.}\texttt{\bfseries tvDist\_push\_le}}\hfill{\footnotesize \path{lean/qtsl-aux/QtslAux/QDistTransport.lean}, line 60}\par\nopagebreak
\begin{leansrc}
theorem tvDist_push_le [DecidableEq β] (π : α → β) (p q : α → ℝ) :
    tvDist (push π p) (push π q) ≤ tvDist p q
\end{leansrc}
\par\medskip\noindent{\small \texttt{Qtsl.QDistTransport.}\texttt{\bfseries tvDist\_step\_le}}\hfill{\footnotesize \path{lean/qtsl-aux/QtslAux/QDistTransport.lean}, line 86}\par\nopagebreak
\begin{leansrc}
theorem tvDist_step_le {K : α → α → ℝ} (hK : IsMarkov K) (p q : α → ℝ) :
    tvDist (step K p) (step K q) ≤ tvDist p q
\end{leansrc}
\par\medskip\noindent{\small \texttt{Qtsl.QDistTransport.}\texttt{\bfseries EpsPushKernel}}\hfill{\footnotesize \path{lean/qtsl-aux/QtslAux/QDistTransport.lean}, line 103}\par\nopagebreak
\begin{leansrc}
def EpsPushKernel [DecidableEq β] (π : α → β) (K : α → α → ℝ)
    (Q : β → β → ℝ) (ε : ℝ) : Prop
\end{leansrc}
\par\medskip\noindent{\small \texttt{Qtsl.QDistTransport.}\texttt{\bfseries push\_step\_defect}}\hfill{\footnotesize \path{lean/qtsl-aux/QtslAux/QDistTransport.lean}, line 107}\par\nopagebreak
\begin{leansrc}
theorem push_step_defect [DecidableEq β] {π : α → β} {K : α → α → ℝ}
    {Q : β → β → ℝ} {ε : ℝ} {μ : α → ℝ} (hμ : IsDist μ)
    (h : EpsPushKernel π K Q ε) :
    tvDist (push π (step K μ)) (step Q (push π μ)) ≤ ε
\end{leansrc}
\par\medskip\noindent{\small \texttt{Qtsl.QDistTransport.}\texttt{\bfseries tvDist\_triangle}}\hfill{\footnotesize \path{lean/qtsl-aux/QtslAux/QDistTransport.lean}, line 40}\par\nopagebreak
\begin{leansrc}
theorem tvDist_triangle {L : Type*} [Fintype L] (p q r : L → ℝ) :
    tvDist p r ≤ tvDist p q + tvDist q r
\end{leansrc}
\par\medskip\noindent{\small \texttt{SuFunctor.Universe.BelState.}\texttt{\bfseries push\_comp}}\hfill{\footnotesize \path{lean/su-functor/SuFunctor/Universe/TotalCitizenship.lean}, line 89}\par\nopagebreak
\begin{leansrc}
theorem BelState.push_comp (f : V → U) (g : U → W) (x : BelState V) :
    BelState.push (g ∘ f) x = BelState.push g (BelState.push f x)
\end{leansrc}
\par\medskip\noindent{\small \texttt{Qtsl.QDistTransport.}\texttt{\bfseries eq\_of\_tvDist\_le\_zero}}\hfill{\footnotesize \path{lean/qtsl-aux/QtslAux/QDistTransport.lean}, line 48}\par\nopagebreak
\begin{leansrc}
theorem eq_of_tvDist_le_zero {L : Type*} [Fintype L] {p q : L → ℝ}
    (h : tvDist p q ≤ 0) : p = q
\end{leansrc}
\par\medskip\noindent{\small \texttt{SuFunctor.Adequacy.}\texttt{\bfseries dtmcAbsorb\_mass\_le\_reach}}\hfill{\footnotesize \path{lean/su-functor/SuFunctor/Adequacy/DtmcAbsorb.lean}, line 484}\par\nopagebreak
\begin{leansrc}
theorem dtmcAbsorb_mass_le_reach (M : FinDTMC S) {C : Finset S}
    {T t' : ℕ} (h : T ≤ t') :
    massOn M t' C ≤ massOn M T C + massOn M T (M.reachSet C \ C)
\end{leansrc}
\par\medskip\noindent{\small \texttt{SuFunctor.Adequacy.}\texttt{\bfseries dtmcAbsorbBracket\_reach\_sound}}\hfill{\footnotesize \path{lean/su-functor/SuFunctor/Adequacy/DtmcAbsorb.lean}, line 512}\par\nopagebreak
\begin{leansrc}
theorem dtmcAbsorbBracket_reach_sound (M : FinDTMC S)
    {C Abs : Finset S} (hAbs : ∀ s ∈ Abs, M.Absorbing s)
    (hCA : C ⊆ Abs) {T t' : ℕ} (h : T ≤ t') :
    massOn M T C ≤ massOn M t' C ∧
      massOn M t' C ≤ massOn M T C + massOn M T (M.reachSet C \ C)
\end{leansrc}
\par\medskip\noindent{\small \texttt{SuFunctor.Adequacy.}\texttt{\bfseries dtmcAbsorb\_reach\_residual\_le}}\hfill{\footnotesize \path{lean/su-functor/SuFunctor/Adequacy/DtmcAbsorb.lean}, line 528}\par\nopagebreak
\begin{leansrc}
theorem dtmcAbsorb_reach_residual_le (M : FinDTMC S)
    {C Abs : Finset S} (hAbs : ∀ s ∈ Abs, M.Absorbing s) (T : ℕ) :
    massOn M T (M.reachSet C \ C) ≤ 1 - massOn M T Abs
\end{leansrc}
\par\medskip\noindent{\small \texttt{SuFunctor.Adequacy.}\texttt{\bfseries dtmcAbsorb\_eventual\_bracket\_reach}}\hfill{\footnotesize \path{lean/su-functor/SuFunctor/Adequacy/DtmcAbsorb.lean}, line 550}\par\nopagebreak
\begin{leansrc}
theorem dtmcAbsorb_eventual_bracket_reach (M : FinDTMC S)
    {C Abs : Finset S} (hAbs : ∀ s ∈ Abs, M.Absorbing s)
    (hCA : C ⊆ Abs) {T : ℕ} {lo hi : ℚ}
    (hlo : lo ≤ massOn M T C)
    (hhi : massOn M T C + massOn M T (M.reachSet C \ C) ≤ hi) :
    ∀ t', T ≤ t' → lo ≤ massOn M t' C ∧ massOn M t' C ≤ hi
\end{leansrc}
\par\medskip\noindent{\small \texttt{SuFunctor.Adequacy.FinDTMC.}\texttt{\bfseries ForwardClosed}}\hfill{\footnotesize \path{lean/su-functor/SuFunctor/Adequacy/DtmcAbsorb.lean}, line 229}\par\nopagebreak
\begin{leansrc}
def FinDTMC.ForwardClosed (M : FinDTMC S) (D : Finset S) : Prop
\end{leansrc}
\par\medskip\noindent{\small \texttt{SuFunctor.Adequacy.}\texttt{\bfseries dtmc\_mass\_mono\_of\_forwardClosed}}\hfill{\footnotesize \path{lean/su-functor/SuFunctor/Adequacy/DtmcAbsorb.lean}, line 276}\par\nopagebreak
\begin{leansrc}
theorem dtmc_mass_mono_of_forwardClosed (M : FinDTMC S)
    {D : Finset S} (hD : M.ForwardClosed D) {T t' : ℕ} (h : T ≤ t') :
    massOn M T D ≤ massOn M t' D
\end{leansrc}
\par\medskip\noindent{\small \texttt{SuFunctor.Adequacy.}\texttt{\bfseries idBool\_strict\_improvement}}\hfill{\footnotesize \path{lean/su-functor/SuFunctor/Adequacy/DtmcAbsorb.lean}, line 769}\par\nopagebreak
\begin{leansrc}
theorem idBool_strict_improvement {T t' : ℕ} (h : T ≤ t') :
    massOn idBool t' {true} = 0
\end{leansrc}
\par\medskip\noindent{\small \texttt{SuFunctor.Adequacy.}\texttt{\bfseries dtmcPosteriorBracket\_sound}}\hfill{\footnotesize \path{lean/su-functor/SuFunctor/Adequacy/DtmcAbsorb.lean}, line 574}\par\nopagebreak
\begin{leansrc}
theorem dtmcPosteriorBracket_sound {ι : Type} [Fintype ι]
    (N Mv : ι → ℚ) (R : ℚ) (hR : 0 ≤ R)
    (hN : ∀ c, 0 ≤ N c)
    (hS : 0 < ∑ c', N c')
    (hle : ∀ c, N c ≤ Mv c)
    (hgap : (∑ c', (Mv c' - N c')) ≤ R) (c : ι) :
    N c / ((∑ c', N c') + R) ≤ Mv c / (∑ c', Mv c') ∧
      Mv c / (∑ c', Mv c') ≤ (N c + R) / ((∑ c', N c') + R)
\end{leansrc}
\par\medskip\noindent{\small \texttt{SuFunctor.Adequacy.}\texttt{\bfseries massAt\_expect\_le\_of\_contraction\_of\_weighted}}\hfill{\footnotesize \path{lean/su-functor/SuFunctor/Adequacy/DtmcContraction.lean}, line 164}\par\nopagebreak
\begin{leansrc}
theorem massAt_expect_le_of_contraction_of_weighted (M : FinDTMC S)
    (V : S → ℚ) {c : ℚ} (hc : 0 ≤ c)
    (hcontr : ∀ s, (∑ s', M.step s s' * V s') ≤ c * V s) (t : ℕ) :
    (∑ s, massAt M t s * V s) ≤ c ^ t * ∑ s, massAt M 0 s * V s
\end{leansrc}
\par\medskip\noindent{\small \texttt{SuFunctor.Adequacy.}\texttt{\bfseries qGenFormContraction\_subsumes\_genform}}\hfill{\footnotesize \path{lean/su-functor/SuFunctor/Adequacy/DtmcContraction.lean}, line 181}\par\nopagebreak
\begin{leansrc}
theorem qGenFormContraction_subsumes_genform [DecidableEq S]
    (M : FinDTMC S) (V : S → ℚ)
    (hg : ∀ s, genRow (compile M) V s ≤ 0) (t : ℕ) :
    expectQ (M.pathLaw t) (atLast V) ≤
      expectQ (M.pathLaw 0) (atLast V)
\end{leansrc}
\par\medskip\noindent{\small \texttt{SuFunctor.Adequacy.}\texttt{\bfseries wmassAt}}\hfill{\footnotesize \path{lean/su-functor/SuFunctor/Adequacy/DtmcContraction.lean}, line 63}\par\nopagebreak
\begin{leansrc}
def wmassAt (init : S → ℚ) (W : S → S → ℚ) : ℕ → S → ℚ
  | 0 => init
  | t + 1 => fun s' => ∑ s, wmassAt init W t s * W s s'

@[simp] theorem wmassAt_zero (init : S → ℚ) (W : S → S → ℚ) :
    wmassAt init W 0 = init
\end{leansrc}
\par\medskip\noindent{\small \texttt{SuFunctor.Adequacy.}\texttt{\bfseries massAt\_eq\_wmassAt}}\hfill{\footnotesize \path{lean/su-functor/SuFunctor/Adequacy/DtmcContraction.lean}, line 85}\par\nopagebreak
\begin{leansrc}
theorem massAt_eq_wmassAt (M : FinDTMC S) (t : ℕ) :
    massAt M t = wmassAt M.init M.step t
\end{leansrc}
\par\medskip\noindent{\small \texttt{SuFunctor.Adequacy.}\texttt{\bfseries wmassAt\_expect\_le\_of\_contraction}}\hfill{\footnotesize \path{lean/su-functor/SuFunctor/Adequacy/DtmcContraction.lean}, line 116}\par\nopagebreak
\begin{leansrc}
theorem wmassAt_expect_le_of_contraction {init : S → ℚ} {W : S → S → ℚ}
    (V : S → ℚ) (hinit : ∀ s, 0 ≤ init s) (hW : ∀ s s', 0 ≤ W s s')
    {c : ℚ} (hc : 0 ≤ c)
    (hcontr : ∀ s, (∑ s', W s s' * V s') ≤ c * V s) (t : ℕ) :
    (∑ s, wmassAt init W t s * V s) ≤ c ^ t * ∑ s, init s * V s
\end{leansrc}
\par\medskip\noindent{\small \texttt{SuFunctor.Adequacy.}\texttt{\bfseries massAt\_expect\_on\_le\_of\_regionContraction}}\hfill{\footnotesize \path{lean/su-functor/SuFunctor/Adequacy/DtmcAliveTail.lean}, line 60}\par\nopagebreak
\begin{leansrc}
theorem massAt_expect_on_le_of_regionContraction (M : FinDTMC S)
    (A : Finset S) (V : S → ℚ) {c : ℚ} (hc : 0 ≤ c)
    (hclosed : ∀ x ∉ A, ∀ y ∈ A, M.step x y = 0)
    (hcontr : ∀ x ∈ A, (∑ y ∈ A, M.step x y * V y) ≤ c * V x) (t : ℕ) :
    (∑ x ∈ A, massAt M t x * V x) ≤ c ^ t * ∑ x ∈ A, M.init x * V x
\end{leansrc}
\par\medskip\noindent{\small \texttt{SuFunctor.Adequacy.}\texttt{\bfseries dtmcAliveTail\_of\_regionContraction}}\hfill{\footnotesize \path{lean/su-functor/SuFunctor/Adequacy/DtmcAliveTail.lean}, line 93}\par\nopagebreak
\begin{leansrc}
theorem dtmcAliveTail_of_regionContraction (M : FinDTMC S)
    (A : Finset S) (V : S → ℚ) {c : ℚ} (hc : 0 ≤ c)
    (hfloor : ∀ x ∈ A, 1 ≤ V x)
    (hclosed : ∀ x ∉ A, ∀ y ∈ A, M.step x y = 0)
    (hcontr : ∀ x ∈ A, (∑ y ∈ A, M.step x y * V y) ≤ c * V x) (t : ℕ) :
    massOn M t A ≤ c ^ t * ∑ x ∈ A, M.init x * V x
\end{leansrc}
\par\medskip\noindent{\small \texttt{SuFunctor.Adequacy.}\texttt{\bfseries aliveTail\_needs\_floor}}\hfill{\footnotesize \path{lean/su-functor/SuFunctor/Adequacy/DtmcAliveTail.lean}, line 128}\par\nopagebreak
\begin{leansrc}
theorem aliveTail_needs_floor :
    ¬ (massOn unitChain 1 Finset.univ ≤
        (1 : ℚ) ^ 1 * ∑ x ∈ Finset.univ, unitChain.init x * ((1 : ℚ) / 2))
\end{leansrc}
\par\medskip\noindent{\small \texttt{SuFunctor.Adequacy.}\texttt{\bfseries aliveTail\_needs\_no\_reentry}}\hfill{\footnotesize \path{lean/su-functor/SuFunctor/Adequacy/DtmcAliveTail.lean}, line 159}\par\nopagebreak
\begin{leansrc}
theorem aliveTail_needs_no_reentry :
    ¬ (massOn reentryChain 1 {true} ≤
        (1 : ℚ) ^ 1 * ∑ x ∈ ({true} : Finset Bool),
          reentryChain.init x * (1 : ℚ))
\end{leansrc}
\par\medskip\noindent{\small \texttt{SuFunctor.Adequacy.}\texttt{\bfseries genFormTemplate\_collapse}}\hfill{\footnotesize \path{lean/su-functor/SuFunctor/Adequacy/DtmcTemplate.lean}, line 35}\par\nopagebreak
\begin{leansrc}
theorem genFormTemplate_collapse {ι : Type} [Fintype ι]
    (w : ι → ℚ) (d : ι → ℤ) {M γ c : ℚ} (hM : 0 ≤ M) (hγ : 0 < γ)
    (h : (∑ i, w i * γ ^ d i) ≤ c) (n : ℤ) :
    (∑ i, w i * (M * γ ^ (n + d i))) ≤ c * (M * γ ^ n)
\end{leansrc}
\par\medskip\noindent{\small \texttt{SuFunctor.Adequacy.}\texttt{\bfseries aliveStep\_factor}}\hfill{\footnotesize \path{lean/su-functor/SuFunctor/Adequacy/DtmcProductProfile.lean}, line 264}\par\nopagebreak
\begin{leansrc}
theorem aliveStep_factor {g : List ℚ} {ρ : ℚ} (hg : ∀ q ∈ g, 0 < q)
    (hρ : 0 < ρ) (p : ℚ) {d : ℕ} (x : Fin d → ℕ)
    (hx : ∀ i, downMove x i ≠ 0) :
    aliveStep g ρ p x
      = (∑ i, transferRatio g ρ p (x i)) / (d : ℚ)
          * prodProfileV g ρ x
\end{leansrc}
\par\medskip\noindent{\small \texttt{SuFunctor.Adequacy.}\texttt{\bfseries transferRatio\_of\_ge}}\hfill{\footnotesize \path{lean/su-functor/SuFunctor/Adequacy/DtmcProductProfile.lean}, line 180}\par\nopagebreak
\begin{leansrc}
theorem transferRatio_of_ge {g : List ℚ} {ρ : ℚ} (hg : ∀ q ∈ g, 0 < q)
    (hρ : 0 < ρ) (p : ℚ) {v : ℕ} (hv : g.length + 1 ≤ v) :
    transferRatio g ρ p v = rInf ρ p
\end{leansrc}
\par\medskip\noindent{\small \texttt{SuFunctor.Adequacy.}\texttt{\bfseries productProfile\_collapse}}\hfill{\footnotesize \path{lean/su-functor/SuFunctor/Adequacy/DtmcProductProfile.lean}, line 297}\par\nopagebreak
\begin{leansrc}
theorem productProfile_collapse {g : List ℚ} {ρ : ℚ}
    (hg : ∀ q ∈ g, 0 < q) (hρ : 0 < ρ) (p : ℚ) {d : ℕ} {c : ℚ}
    (hclass : ∀ k ∈ Finset.Icc 1 d,
      ((d : ℚ) - (k : ℚ)) * rMax g ρ p + (k : ℚ) * rInf ρ p
        ≤ (d : ℚ) * c)
    (x : Fin d → ℕ) (j : Fin d) (hgeo : g.length + 1 ≤ x j)
    (h2 : 2 ≤ x j) :
    aliveStep g ρ p x ≤ c * prodProfileV g ρ x
\end{leansrc}
\par\medskip\noindent{\small \texttt{SuFunctor.Adequacy.}\texttt{\bfseries rampProfile\_needs\_window\_reach}}\hfill{\footnotesize \path{lean/su-functor/SuFunctor/Adequacy/DtmcProductProfile.lean}, line 370}\par\nopagebreak
\begin{leansrc}
theorem rampProfile_needs_window_reach :
    (∀ k ∈ Finset.Icc (1 : ℕ) 3,
      ((3 : ℚ) - (k : ℚ)) * rMax rampG (7/4) (1/4)
        + (k : ℚ) * rInf (7/4) (1/4) ≤ (3 : ℚ) * (193/200)) ∧
    (3 : ℚ) * (193/200)
      < 2 * transferRatio rampG (7/4) (1/4) 1
          + transferRatio rampG (7/4) (1/4) 5
\end{leansrc}
\par\medskip\noindent{\small \texttt{SuFunctor.Adequacy.}\texttt{\bfseries sum\_range\_mul\_pow\_mul\_eq}}\hfill{\footnotesize \path{lean/su-functor/SuFunctor/Adequacy/DtmcContraction.lean}, line 279}\par\nopagebreak
\begin{leansrc}
theorem sum_range_mul_pow_mul_eq (c : ℚ) (n : ℕ) :
    (∑ i ∈ Finset.range n, (i : ℚ) * c ^ i) * (1 - c) ^ 2 =
      c - (n : ℚ) * c ^ n + ((n : ℚ) - 1) * c ^ (n + 1)
\end{leansrc}
\par\medskip\noindent{\small \texttt{SuFunctor.Adequacy.}\texttt{\bfseries sum\_Ico\_pow\_le}}\hfill{\footnotesize \path{lean/su-functor/SuFunctor/Adequacy/DtmcContraction.lean}, line 318}\par\nopagebreak
\begin{leansrc}
theorem sum_Ico_pow_le {c : ℚ} (hc : 0 ≤ c) (hc1 : c < 1) (T T' : ℕ) :
    (∑ t ∈ Finset.Ico T T', c ^ t) ≤ c ^ T / (1 - c)
\end{leansrc}
\par\medskip\noindent{\small \texttt{SuFunctor.Adequacy.}\texttt{\bfseries sum\_Ico\_mul\_pow\_le}}\hfill{\footnotesize \path{lean/su-functor/SuFunctor/Adequacy/DtmcContraction.lean}, line 335}\par\nopagebreak
\begin{leansrc}
theorem sum_Ico_mul_pow_le {c : ℚ} (hc : 0 ≤ c) (hc1 : c < 1)
    (T T' : ℕ) :
    (∑ t ∈ Finset.Ico T T', (t : ℚ) * c ^ t) ≤
      c ^ T * ((T : ℚ) * (1 - c) + c) / (1 - c) ^ 2
\end{leansrc}
\par\medskip\noindent{\small \texttt{SuFunctor.Adequacy.}\texttt{\bfseries wmassAt\_exit\_tail\_le}}\hfill{\footnotesize \path{lean/su-functor/SuFunctor/Adequacy/DtmcContraction.lean}, line 374}\par\nopagebreak
\begin{leansrc}
theorem wmassAt_exit_tail_le {init : S → ℚ} {W : S → S → ℚ}
    (V e : S → ℚ) (hinit : ∀ s, 0 ≤ init s) (hW : ∀ s s', 0 ≤ W s s')
    {c : ℚ} (hc : 0 ≤ c) (hc1 : c < 1)
    (hcontr : ∀ s, (∑ s', W s s' * V s') ≤ c * V s)
    {κ : ℚ} (hκ : 0 ≤ κ) (hexit : ∀ s, e s ≤ κ * V s)
    (hIV : 0 ≤ ∑ s, init s * V s) (T T' : ℕ) :
    (∑ t ∈ Finset.Ico T T', ∑ s, wmassAt init W t s * e s) ≤
      κ * c ^ T * (∑ s, init s * V s) / (1 - c)
\end{leansrc}
\par\medskip\noindent{\small \texttt{SuFunctor.Adequacy.}\texttt{\bfseries wmassAt\_exit\_moment\_tail\_le}}\hfill{\footnotesize \path{lean/su-functor/SuFunctor/Adequacy/DtmcContraction.lean}, line 403}\par\nopagebreak
\begin{leansrc}
theorem wmassAt_exit_moment_tail_le {init : S → ℚ} {W : S → S → ℚ}
    (V e : S → ℚ) (hinit : ∀ s, 0 ≤ init s) (hW : ∀ s s', 0 ≤ W s s')
    {c : ℚ} (hc : 0 ≤ c) (hc1 : c < 1)
    (hcontr : ∀ s, (∑ s', W s s' * V s') ≤ c * V s)
    {κ : ℚ} (hκ : 0 ≤ κ) (hexit : ∀ s, e s ≤ κ * V s)
    (hIV : 0 ≤ ∑ s, init s * V s) (T T' : ℕ) :
    (∑ t ∈ Finset.Ico T T', (t : ℚ) * ∑ s, wmassAt init W t s * e s) ≤
      κ * (∑ s, init s * V s) * c ^ T * ((T : ℚ) * (1 - c) + c) /
        (1 - c) ^ 2
\end{leansrc}
\par\medskip\noindent{\small \texttt{Qtsl.Acceptance.}\texttt{\bfseries towerCitation\_sound\_of\_wf}}\hfill{\footnotesize \path{lean/qtsl-aux/QtslAux/Acceptance.lean}, line 56}\par\nopagebreak
\begin{leansrc}
theorem towerCitation_sound_of_wf
    (C : ι → Row σ J ι) (baseTrue : σ → Prop) (sem : J → Prop)
    (SideOK : ι → Prop)
    (hbase : ∀ i s, C i = Row.base s → baseTrue s)
    (hside : ∀ i, SideOK i)
    (hstep : ∀ i prems concl, C i = Row.ext prems concl → SideOK i →
      (∀ j ∈ prems, Row.holds baseTrue sem (C j)) → sem concl)
    (hwf : WellFounded (cites C)) :
    ∀ i, Row.holds baseTrue sem (C i)
\end{leansrc}
\par\medskip\noindent{\small \texttt{Qtsl.Acceptance.}\texttt{\bfseries towerCitation\_sound}}\hfill{\footnotesize \path{lean/qtsl-aux/QtslAux/Acceptance.lean}, line 74}\par\nopagebreak
\begin{leansrc}
theorem towerCitation_sound [Finite ι]
    (C : ι → Row σ J ι) (baseTrue : σ → Prop) (sem : J → Prop)
    (SideOK : ι → Prop)
    (hbase : ∀ i s, C i = Row.base s → baseTrue s)
    (hside : ∀ i, SideOK i)
    (hstep : ∀ i prems concl, C i = Row.ext prems concl → SideOK i →
      (∀ j ∈ prems, Row.holds baseTrue sem (C j)) → sem concl)
    (hacyc : Acyclic C) :
    ∀ i, Row.holds baseTrue sem (C i)
\end{leansrc}
\par\medskip\noindent{\small \texttt{Qtsl.Acceptance.}\texttt{\bfseries towerCitation\_root}}\hfill{\footnotesize \path{lean/qtsl-aux/QtslAux/Acceptance.lean}, line 86}\par\nopagebreak
\begin{leansrc}
theorem towerCitation_root [Finite ι]
    (C : ι → Row σ J ι) (baseTrue : σ → Prop) (sem : J → Prop)
    (SideOK : ι → Prop)
    (hbase : ∀ i s, C i = Row.base s → baseTrue s)
    (hside : ∀ i, SideOK i)
    (hstep : ∀ i prems concl, C i = Row.ext prems concl → SideOK i →
      (∀ j ∈ prems, Row.holds baseTrue sem (C j)) → sem concl)
    (hacyc : Acyclic C)
    {Jroot : J} (hroot : ∃ i prems, C i = Row.ext prems Jroot) :
    sem Jroot
\end{leansrc}
\par\medskip\noindent{\small \texttt{Qtsl.Acceptance.Row.}\texttt{\bfseries cited}}\hfill{\footnotesize \path{lean/qtsl-aux/QtslAux/Acceptance.lean}, line 36}\par\nopagebreak
\begin{leansrc}
def Row.cited : Row σ J ι → List ι
  | .base _ => []
  | .ext prems _ => prems

def Row.holds (baseTrue : σ → Prop) (sem : J → Prop) : Row σ J ι → Prop
  | .base s => baseTrue s
  | .ext _ concl => sem concl

def cites (C : ι → Row σ J ι) (j i : ι) : Prop
\end{leansrc}
\par\medskip\noindent{\small \texttt{Qtsl.Acceptance.}\texttt{\bfseries acyclic\_wellFounded}}\hfill{\footnotesize \path{lean/qtsl-aux/QtslAux/Acceptance.lean}, line 50}\par\nopagebreak
\begin{leansrc}
theorem acyclic_wellFounded [Finite ι] (C : ι → Row σ J ι)
    (hacyc : Acyclic C) : WellFounded (cites C)
\end{leansrc}
\par\medskip\noindent{\small \texttt{Qtsl.Acceptance.}\texttt{\bfseries citePos\_sound}}\hfill{\footnotesize \path{lean/qtsl-aux/QtslAux/Acceptance.lean}, line 101}\par\nopagebreak
\begin{leansrc}
theorem citePos_sound {X : Type*} {R D : Set X} {P : Set X → Prop}
    (hmono : Monotone P) (hsub : R ⊆ D) (hR : P R) : P D
\end{leansrc}
\par\medskip\noindent{\small \texttt{Qtsl.Acceptance.}\texttt{\bfseries citeCmpl\_subset}}\hfill{\footnotesize \path{lean/qtsl-aux/QtslAux/Acceptance.lean}, line 105}\par\nopagebreak
\begin{leansrc}
theorem citeCmpl_subset {X : Type*} {R D : Set X} (hsub : R ⊆ D) :
    Dᶜ ⊆ Rᶜ
\end{leansrc}
\par\medskip\noindent{\small \texttt{Qtsl.Acceptance.}\texttt{\bfseries citeCmpl\_sound}}\hfill{\footnotesize \path{lean/qtsl-aux/QtslAux/Acceptance.lean}, line 109}\par\nopagebreak
\begin{leansrc}
theorem citeCmpl_sound {X : Type*} {R D : Set X} {P : Set X → Prop}
    (hanti : Antitone P) (hsub : R ⊆ D) (hR : P Rᶜ) : P Dᶜ
\end{leansrc}
\par\medskip\noindent Modules cited: \texttt{QtslAux.QGuard}, \texttt{QtslAux.LayerD.MPDischarge}, \texttt{SuFunctor.Site.FellAcceptance}, \texttt{SuFunctor.Site.FellBehaviorSpec}, \texttt{QtslAux.Enclosure}.
\par\medskip\noindent Also cited in this chapter, listed under an earlier one: \texttt{ProofObject.sound}, \texttt{SuFunctor.Plain.CheckerFace}, \texttt{semanticallyValid\_iff\_face}.

\section{Partial Completeness}
Citations from Chapter~\ref{ch:qtsl-partial-completeness}.

\par\medskip\noindent{\small \texttt{PPLC.Qtsl.}\texttt{\bfseries sound\_u10f}}\hfill{\footnotesize \path{lean/qtsl/Qtsl/TimedCore.lean}, line 49}\par\nopagebreak
\begin{leansrc}
theorem sound_u10f (B : Behavior X) (T : ℝ≥0) (ν χ : Set X) (x : X)
    (hbox : B.BoxInv ν x) :
    B.AService T (χ ∩ ν) x ↔ B.AService T χ x
\end{leansrc}
\par\medskip\noindent{\small \texttt{PPLC.Qtsl.}\texttt{\bfseries sound\_firstService}}\hfill{\footnotesize \path{lean/qtsl/Qtsl/TimedCore.lean}, line 138}\par\nopagebreak
\begin{leansrc}
theorem sound_firstService (hB3 : B.FusionClosed) (hn : 0 < n)
    (α : Fin n → Set X) (T : Fin n → ℝ≥0) (x : X)
    (hatt : ∀ i, B.FirstVisits (α i))
    (hx : ∀ i, B.AService (T i) (α i) x) :
    B.AService (Finset.univ.inf' (Finset.univ_nonempty_iff.mpr ⟨⟨0, hn⟩⟩) T)
      (B.fsTarget α T) x
\end{leansrc}
\par\medskip\noindent{\small \texttt{PPLC.Qtsl.}\texttt{\bfseries sound\_fsResidual}}\hfill{\footnotesize \path{lean/qtsl/Qtsl/TimedCore.lean}, line 149}\par\nopagebreak
\begin{leansrc}
theorem sound_fsResidual (hB3 : B.FusionClosed) (hn : 0 < n)
    (α : Fin n → Set X) (T : Fin n → ℝ≥0) (c : ℝ≥0) (x : X)
    (hc : ∀ i, c < T i)
    (hatt : ∀ i, B.FirstVisits (α i))
    (hx : ∀ i, B.AService (T i) (α i) x) :
    B.AService c (B.fsTarget α T ∪ B.residualCell α T c) x
\end{leansrc}
\par\medskip\noindent{\small \texttt{PPLC.Qtsl.TimedCoreCountermodel.}\texttt{\bfseries firstService\_attainment\_necessary}}\hfill{\footnotesize \path{lean/qtsl/Qtsl/TimedCore.lean}, line 496}\par\nopagebreak
\begin{leansrc}
theorem firstService_attainment_necessary :
    ∃ (S : Type) (B : Behavior S) (α : Fin 2 → Set S) (x : S),
      B.FusionClosed ∧
      (∀ i, B.AService 2 (α i) x) ∧
      ¬ B.AService 2 (B.fsTarget α (fun _ => 2)) x
\end{leansrc}

\printbibliography[heading=bibintoc,title={References}]

\end{document}
